\documentclass[11pt,a4paper,leqno,oneside,openany]{book}

\usepackage[T1]{fontenc}
\usepackage[utf8]{inputenc}
\usepackage[provide=*,french,english]{babel}
\usepackage{lmodern}
\usepackage{amsmath,amssymb,amsfonts,mathrsfs}
\usepackage{mathtools}
\usepackage{tikz}
\usetikzlibrary{cd,arrows.meta}
\usepackage{geometry}
\geometry{left=21mm,right=21mm,top=22mm,bottom=23mm,marginparwidth=15mm,marginparsep=3mm}
\usepackage{enumitem}
\usepackage[normalem]{ulem}
\let\mathunderline\underline
\renewcommand{\underline}[1]{\ifmmode\mathunderline{#1}\else\uline{#1}\fi}
\emergencystretch=3em

\DeclareUnicodeCharacter{2192}{\ensuremath{\rightarrow}}
\DeclareUnicodeCharacter{2297}{\ensuremath{\otimes}}
\DeclareUnicodeCharacter{2295}{\ensuremath{\oplus}}
\DeclareUnicodeCharacter{2228}{\ensuremath{\vee}}
\DeclareUnicodeCharacter{223C}{\ensuremath{\sim}}
\DeclareUnicodeCharacter{2248}{\ensuremath{\approx}}
\DeclareUnicodeCharacter{00D7}{\ensuremath{\times}}
\DeclareUnicodeCharacter{03B1}{\ensuremath{\alpha}}
\DeclareUnicodeCharacter{03B2}{\ensuremath{\beta}}
\DeclareUnicodeCharacter{00A7}{\S}
\DeclareUnicodeCharacter{010C}{\v{C}}
\DeclareUnicodeCharacter{00C3}{\~{A}}
\DeclareUnicodeCharacter{02D9}{\.{}}

\providecommand{\Zl}{\mathbb{Z}_{\ell}}
\providecommand{\Ql}{\mathbb{Q}_{\ell}}
\providecommand{\Gal}{\operatorname{Gal}}
\providecommand{\Spec}{\operatorname{Spec}}
\providecommand{\Isom}{\operatorname{Isom}}
\providecommand{\Aut}{\operatorname{Aut}}
\providecommand{\Hom}{\operatorname{Hom}}
\providecommand{\Ext}{\operatorname{Ext}}
\providecommand{\Biext}{\operatorname{Biext}}
\providecommand{\Coker}{\operatorname{Coker}}
\providecommand{\Ker}{\operatorname{Ker}}
\providecommand{\Im}{\operatorname{Im}}
\providecommand{\Tor}{\operatorname{Tor}}
\providecommand{\prof}{\operatorname{prof}}
\providecommand{\codim}{\operatorname{codim}}
\providecommand{\pr}{\operatorname{pr}}
\providecommand{\id}{\mathrm{id}}
\providecommand{\sym}{\mathrm{sym}}
\providecommand{\Corr}{\operatorname{Corr}}
\providecommand{\Proj}{\operatorname{Proj}}
\providecommand{\Comp}{\operatorname{Comp}}
\providecommand{\cC}{\mathcal{C}}
\providecommand{\cY}{\mathcal{Y}}
\providecommand{\cT}{\mathcal{T}}
\providecommand{\cPsi}{\Psi}
\providecommand{\Dcat}{D}
\providecommand{\La}{\Lambda}
\providecommand{\Ghat}{\widehat{G}}
\providecommand{\uHom}{\underline{\operatorname{Hom}}}
\providecommand{\cl}{\operatorname{c\ell}}
\providecommand{\e}{\mathrm{e}}




\usepackage{etoolbox,adjustbox,marginnote,xurl,hyperref}
\tikzcdset{ampersand replacement=\&}
\tikzset{ampersand replacement=\&}
\hypersetup{hidelinks,hypertexnames=false,pdfauthor={GPT-6 Astra; source authors credited in the edition},pdftitle={SGA 7 I and II - English reader}}
\BeforeBeginEnvironment{tikzcd}{\begin{adjustbox}{max width=\linewidth}}
\AfterEndEnvironment{tikzcd}{\end{adjustbox}}
\setcounter{tocdepth}{1}
\setlength{\emergencystretch}{4em}
\newcount\SGAform
\SGAform=0
\ifnum\SGAform=0
  \long\def\SGAvar#1#2#3{#1}
\else\ifnum\SGAform=1
  \long\def\SGAvar#1#2#3{\SGAvmark{#3}#2}
\else
  \long\def\SGAvar#1#2#3{#2}
\fi\fi
\newcommand{\SGAvmark}[1]{\ifstrempty{#1}{}{\leavevmode\textsuperscript{\tiny\hyperlink{app-#1.1}{[#1]}}\ignorespaces}}
\newcommand{\SGApage}[2]{\marginnote{\normalfont\scriptsize #1\,\textbf{#2}}}
\newcommand{\SGAtoc}[1]{\phantomsection\addcontentsline{toc}{chapter}{#1}}
\newcommand{\SGAtocunderline}[1]{\ifmmode\mathunderline{#1}\else #1\fi}
\newcommand{\SGAsectiontoc}[2]{\phantomsection\addcontentsline{toc}{section}{\texorpdfstring{#1}{Section #2}}}

\providecommand{\Pic}{\operatorname{Pic}}
\providecommand{\uExt}{\underline{\operatorname{Ext}}}
\providecommand{\uTor}{\underline{\operatorname{Tor}}}
\makeatletter
\DeclareRobustCommand{\SGARefTarget}[1]{\mbox{\phantomsection\label{#1}}}
\DeclareRobustCommand{\SGARefAliasTarget}[1]{\mbox{\phantomsection\label{#1}}}
\DeclareRobustCommand{\SGARefLink}[2]{\@ifundefined{r@#1}{#2}{\hyperref[#1]{#2}}}
\let\SGAoriginaltextcircled\textcircled
\DeclareRobustCommand{\textcircled}[1]{\ifmmode\text{\SGAoriginaltextcircled{#1}}\else\SGAoriginaltextcircled{#1}\fi}
\renewcommand{\SGAtoc}[1]{\phantomsection\addcontentsline{toc}{chapter}{#1}\markboth{#1}{#1}}
\makeatother
\tikzcdset{ampersand replacement=\&}
\tikzset{ampersand replacement=\&}

\begin{document}\selectlanguage{english}\frontmatter
\begin{titlepage}\centering
{\Large S\'eminaire de g\'eom\'etrie alg\'ebrique du Bois Marie\par}\vspace{20mm}
{\Huge\bfseries SGA 7\par}\vspace{8mm}
{\LARGE Monodromy Groups in Algebraic Geometry\par}\vspace{8mm}
{\Large Volumes I and II\par}\vspace{16mm}
{\Large\ifcase\SGAform Diplomatic text\or Corrected text with apparatus\or Plain reading text\fi\par}\vfill
{\large Edition of 12 November 2026\par}
{\large GPT-6 Astra (OpenAI, ChatGPT Pro mode)\par}
\end{titlepage}
\section*{About this form}
\ifcase\SGAform
The diplomatic form follows the printed readings established by the paired source witnesses. Transcription repairs common to the two witnesses are retained. Expos\'e VI reproduces the original English directly.
\or
The corrected form follows the released English, the shared original English of Expos\'e VI, and the explicit editorial decisions. Every passage different from the diplomatic form is marked. The apparatus records each difference and its source. Inherited editorial notes and questions retained for review are presented separately.
\or
This is the corrected text without variant marks, critical apparatus, inherited editorial explanations, or unresolved-reading notes. The authorial text, authorial notes, bibliographies, and supplied publisher catalogue are retained. Printed folios appear in the margins only.
\fi
The margin labels refer to the original print, independently of this book's page numbers. An approximation sign marks an estimated English alignment between mathematical anchors; exact numbered-item matches are unqualified. The original English of Expos\'e VI retains every source folio marker. Expos\'es III--V do not exist as separate written texts.
\section*{Bibliographic identity}

Alexander Grothendieck, with M. Raynaud and D. S. Rim, Groupes de monodromie en géométrie algébrique I, Séminaire de géométrie algébrique du Bois-Marie 1967–1969 (SGA 7 I), Lecture Notes in Mathematics 288, Springer, 1972: Exposés I, II, VI–IX. Pierre Deligne and Nicholas Katz, Groupes de monodromie en géométrie algébrique II (SGA 7 II), Lecture Notes in Mathematics 340, Springer, 1973: Exposés X–XXII. The detailed authorship credits in the exposés are retained. The original-language Exposé VI by D. S. Rim is already English and remains English in both language readers.

\section*{Credits and provenance}

This build is credited to GPT-6 Astra (OpenAI, ChatGPT Pro mode). The French source is AI-transcribed; the English is AI-translated, except for Exposé VI, whose original English is retained. The edition contains AI-produced corrections. These descriptions do not transfer any human contribution to AI. The English authority is the supplied sga-en release v2026-11-12-english-edition, commit 4a4cbac; the French diplomatic and current layers, correction logs and decisions are the supplied November 2026 records. The August French layer is used only as a change-log baseline. The following source excerpts preserve the recorded credits, including contributions to other volumes of the wider SGA project; those wider credits do not assert a contribution to SGA 7.

\subsection*{Recorded human contributions}

Sources and authorship

The original SGA volumes retain their authors, seminar contributors and editors, named in the corresponding texts. The collection includes work of Alexander Grothendieck, Michèle Raynaud, Michel Raynaud, Pierre Deligne, Nicholas Katz, Michel Demazure, Michael Artin, Jean-Louis Verdier, Luc Illusie and their collaborators; this list does not replace the detailed credits in each volume.

SGA 1 retains the French retypesetting associated with Bas Edixhoven’s volunteer project, the Société mathématique de France and Michel Raynaud’s notes, from arXiv math/0206203v2 (\url{https://arxiv.org/abs/math/0206203v2}).

SGA 2 is AI-translated from the corrected French TeX of the re-edition edited by Yves Laszlo, Cohomologie locale des faisceaux cohérents et théorèmes de Lefschetz locaux et globaux (SGA 2), Documents Mathématiques 4, Société mathématique de France, 2005 (arXiv math/0511279v1 (\url{https://arxiv.org/abs/math/0511279v1})), and follows its numbering and labels. Most of the initial French typesetting was done by volunteers of Bas Edixhoven’s project (L. Bayle, N. Borne, O. Brinon, J. Buresi, M. Chardin, F. Ducrot, P. Graftiaux, F. Han, P. Karwasz, L. Koelblen, D. Madore, S. Morel, D. Naie, B. Osserman, J. Riou and V. Sécherre); C. Sabbah adapted it to the SMF format.

The SGA 2 and SGA 3 English sources acknowledge Jacob C. Reinhold’s English translation (\url{https://github.com/jcreinhold/sga/tree/e7a259f3f8608ad3edf9bf6eead3fd504dd2d23e}), snapshot \url{e7a259f3f8608ad3edf9bf6eead3fd504dd2d23e}, used for comparison and drafting. His stated CC BY 4.0 notice applies to that contribution and is preserved; it does not license the underlying French works or every contribution here. Inherited English wording remains attributed to its human source. In SGA 3, lemmas, proofs and corrections added in the Philippe Gille–Patrick Polo re-edition are given in Translator's notes and Additions that cite it.

SGA 4 is AI-translated from the French LaTeX edition of SGA 4, Tomes 1–3, distributed by Fabrice Orgogozo (normalesup.org/\textasciitilde{}forgogozo/SGA4 (\url{https://www.normalesup.org/~forgogozo/SGA4/}); source snapshot at git commit 71766d9), and uses its TeX configuration files. Its footnotes headed “N.D.E.” are that edition's editorial notes, translated.

SGA 4½ preserves Daniel Miller’s LaTeX typesetting (\url{https://github.com/dkmiller/sga4.5}) and its original typesetter’s note. It is not attributed to AI.

Additional English translation is AI-translated; subsequent corrections, editing, reconstruction and typesetting are AI-produced. The source-specific human credits and genuine third-party notices remain applicable; attribution does not establish permission or endorsement.


\subsection*{Recorded production history}

Jacob C. Reinhold's English translation was used for comparison and drafting in SGA 2 and SGA 3. Recorded AI work includes OpenAI Codex (GPT-5.6 Sol: translation; GPT-6 Astra: corrections and assembly, September–October 2026) and Claude Opus 5.5 (Anthropic: the SGA 2 and SGA 3 sources of 26 October 2026, and the translation of SGA 3 revised against the 1970 printing in November 2026). GPT-6.1 Sol (OpenAI API) proposed the edits that carry the corrections of the scan check of the French text into SGA 4 (November 2026); each edit was checked by hand. The release notes record the correction rounds: spectrum notation in SGA 1 (25 September 2026), mathematical and reference corrections in SGA 1 and SGA 7 II (3 October 2026), references in SGA 3 (24 October 2026), the scan-check corrections in SGA 4 (12 November 2026).


\subsection*{SGA 7 I source authority}

SGA 7 I is complete as the published volume stands. The Springer volume is explicitly Exposés I, II, VI, VII, VIII, IX. Its introductory account explains that Grothendieck's oral Exposés I–V were not separately written: their substance was summarized in the written Exposé I, and there are no separately written Exposés III–V. The jump from II to VI is therefore historical, not a translation or packaging gap. The current English component includes every exposé in the published volume. Its existing release filename is retained to preserve stable asset links.


\subsection*{Inherited English witness}

The inherited witness documented in \url{en/evidence/supplemental_working_decisions/sga7i/INHERITED_ENGLISH_PROVENANCE_AND_ADVERSE_USE.md} has 278,864 bytes and SHA-256 \url{CFCB4C17081E16F963DCBC3FA04A43C4D4065477CB37BD931233840D287AA2E9}. Its exact blob first appears in the public-mirror history at commit \url{cb420190bbc6311cff96a08a5a11a4caf483fdff}, dated 2026-05-29 12:24:48+02:00, under the generic identity “Manuscript Typesetting Project”; the Git blob is \url{2ac760f78b6107d2d71bd2e513e5baa4e3be418d}. The local working paths and address are omitted for privacy.

That service identity and the surviving package metadata do not identify the translating model or human. No defensible evidence presently attributes the text specifically to Claude, Codex, ChatGPT, or a human translator. Its author/model lineage is therefore recorded as unresolved.

This witness is retained unchanged for provenance, terminology comparison, and adverse-history purposes. It must not be merged continuously, treated as cumulative coverage, or cited as source-aligned. Individual phrases may be reused only after direct matching to the fresh French authority. The active English reader is source-first; Expos\textbackslash{}'e VI will be taken from its already-English frozen source after a structural check.


\subsection*{French corpus provenance}

SGA7 I French editions

This is an exact no-overwrite binding of two separate source-language editions through Tome-I EOF: o is the printed-diplomatic layer and c is the separately corrected scholarly layer. They are not a bilingual or interleaved edition. The English edition remains separate.

The upstream layer proves complete 528/528 source-folio coverage and stable 267-page readers. Its own audit record is intentionally narrower than a critical edition: representative pages and consequential ambiguities were checked directly, but no exhaustive high-detail page-by-page source audit is claimed. The scan remains controlling for future adjudication.


SGA7 II French editions

This is an exact no-overwrite binding of two separate source-language editions through Tome-II EOF: o is the printed-diplomatic layer and c is the separately corrected scholarly layer. They are not a bilingual or interleaved edition. The English edition remains separate.

The bound layer covers scan indices 8--445, with 438 markers / 438 unique body pages and no missing or duplicate index, across Exposes X--XXII. Gate-00 on 2026-08-20 repaired the omitted scan index 407 / printed folio 400 (the terminal Expose-XXI bibliography) in both French layers and rebuilt both readers. The immutable closure receipt is gate00\_idx407\_sourcekeeper\_closure\_20260820.json. Changed loci and terminal surfaces have direct review evidence; this binding does not silently promote that bounded review to an exhaustive high-detail audit of every unchanged printed page.


\section*{Printed front matter and contents}

The tables of contents in these readers are newly generated from the transcribed text. They do not purport to transcribe the printed tables. The SGA 7 II supplied layers begin with Exposé X; its preceding printed title, publication and contents pages are unavailable in these layers. SGA 7 I preserves its transcribed general introduction but has no separately transcribed printed contents table; no absent printed title or publication pages have been reconstructed. Exposés III–V are not missing transcriptions: the source-authority record explains their historical status.

\section*{Rights}

The rights in the original work belong to its author or the author's representatives or heirs. Where we hold rights, in the AI-produced transcription, translation, corrections and typesetting, they are dedicated to the public domain under CC0 1.0. This edition is published for educational, non-profit purposes, and we claim no rights we do not hold. If you hold rights in this work and object to its publication here, please open a GitHub issue and we will take it down.

\par\noindent\url{https://github.com/KokunoYumeto/sga-en/issues}

\par\noindent\url{https://creativecommons.org/publicdomain/zero/1.0/}

\clearpage

\tableofcontents
\mainmatter
\part{SGA 7 I}
\clearpage\SGAtoc{General introduction}
\selectlanguage{english}
\SGApage{I}{V}

\thispagestyle{empty}
\begin{center}
\underline{INTRODUCTION}
\end{center}

Let $S$ be a scheme and $f:X\longrightarrow S$ a morphism of schemes. If $f$ is
proper and smooth, and if the prime number $\ell$ is invertible on $S$, the
$\ell$-adic cohomology groups $H^i(X_{\bar s},\mathbb Z_\ell)$ of the geometric
fibres of $f$ form an $\ell$-adic local system on $S$. When $f$ is only assumed
proper, these groups are the fibres of an $\ell$-adic sheaf
$R^i f_*\mathbb Z_\ell$ on $S$. The \underline{theory of vanishing cycles}
relates the ramification of this sheaf on $S$ to the singularities of $f$.

We shall consider only the case where $S$ is a henselian trait (= the spectrum
of a henselian discrete valuation ring). In practice, this is the essential
case. I refer to (I 2.1) for a heuristic description of the theory. Let
$S=\Spec(V)$, let $k(\bar\eta)$ be an algebraic closure of the fraction field
$k(\eta)$ of $V$, and let $k(\bar s)$ be the corresponding algebraic closure of
the residue field $k(s)$ ($k(\bar s)$ is the residue field of the normalization
$\bar V$ of $V$ in $k(\bar\eta)$). In the special case where $X$ is proper and
flat over $S$, of relative dimension $n$, and where $f$ has only one point of
non-smoothness $x\in X_s$, one defines $\Gal(\bar\eta/\eta)$-modules $\phi^i$,
of a purely local nature in a neighbourhood of $x$, vanishing for
$i\notin[0,n]$ (I 4.2), and one constructs a long exact sequence of
$\Gal(\bar\eta/\eta)$-modules

\[
\cdots \longrightarrow H^i(X_{\bar s},\mathbb Z_\ell)
\xrightarrow{\mathrm{sp}} H^i(X_{\bar\eta},\mathbb Z_\ell)
\longrightarrow \phi^i
\longrightarrow H^{i+1}(X_{\bar s},\mathbb Z_\ell)
\longrightarrow \cdots
\]

(the $\Gal(\bar s/s)$-module $H^i(X_{\bar s},\mathbb Z_\ell)$ is regarded as a
$\Gal(\bar\eta/\eta)$-module on which inertia acts trivially; $\mathrm{sp}$ is
the specialization map).

\SGAvar{Criteria are also given, valid for the moment only in characteristic $0$, for
the $\phi^i$ with small $i$ to vanish (for example, if $X$ is moreover locally
a complete intersection, $\phi^i=0$ for $i\ne n$ (I. 4.5)).

}{Criteria are also given, valid for the moment only in characteristic $0$, for
the $\phi^i$ with small $i$ to vanish (for example, if $X$ is moreover locally
a complete intersection, $\phi^i=0$ for $i\ne n$ (I. 4.6)).

}{E7I-Introduction-001}\SGApage{I}{VI}

The \underline{monodromy theorem} asserts that the action of the inertia
subgroup $I$ of $\Gal(\bar\eta/\eta)$ is quasi-unipotent: for $T\in I$, there
exist integers $N>0$, $M>0$ such that the endomorphism $(T^M-I)^N$ of
$H^i(X_{\bar\eta},\mathbb Z_\ell)$ is zero. Two proofs of this theorem are
given. The first (I.1.2), which is arithmetic in nature, applies as soon as the
cohomology groups under consideration have reasonable finiteness properties.
The key point is that, when $k(s)$ is of finite type over its prime subfield,
the action of $I$ is quasi-unipotent for every $\ell$-adic representation of
$\Gal(\bar\eta/\eta)$. The second proof, which is more geometric, requires
purity and resolution of singularities: for the moment it is valid only in
characteristic $0$ or for an $H^1$. It provides valuable information about the
nilpotence exponent $N$ ($N\leq i+1$ for an $H^i$). As will be seen later, over
$\mathbb C$ it is well suited to comparison with transcendental theory.

Applied to the $H^1$ of abelian varieties, i.e. to their Tate module, these
results make it possible to study their reduction modulo $p$. Thus two proofs
are given of the stable reduction theorem for abelian varieties, according to
which, after ramification, the connected special fibre of the Néron model is
an extension of an abelian variety by a torus. The first (I. 6) follows the
method of the arithmetic proof of the monodromy theorem. The second (IX 3.6)
rests on a much finer analysis of the Néron model and of its polarization
properties.

Grothendieck's Exposés I through V were not written up. They were summarized in
one Exposé I. The results stated there are proved succinctly, but essentially
completely.

Exposé II applies the method of Lefschetz pencils and (I 5.3) to the study of
the fundamental group. If $S$ is a surface over an algebraically closed field
$k$ of characteristic exponent $p$, it is shown that the profinite group
$\pi_1^{(p)}(S,s)$, the completion away from $p$ of the fundamental group,
admits a finite pro-$(p)$ presentation.

As explained above, Exposés III through V do not exist.

\SGApage{I}{VII}

Exposé VI contains, with a few additions, Schlessinger's theory of
deformations. Care is taken there to account for the infinitesimal
automorphisms of the objects being classified, and not to pass too abruptly to
the functor of isomorphism classes of objects. It also gives a new construction
of $\underline{\operatorname{Ext}}^{i}(L_{X/S},-)$ ($L_{X/S}$ being the
relative cotangent complex of $X/S$). In the rest of the seminar this Exposé is
used chiefly through its study of deformations of ordinary quadratic
singularities.

Exposés VII and VIII are devoted to the theory of biextensions. This theory
plays an essential role in Exposé IX, in expressing what becomes of a
polarization when one passes from an abelian variety to its Néron model. Exposé
IX contains the proof of the stable reduction theorem for abelian varieties,
as well as various applications.

The continuation of this seminar, SGA 7 II, by P. Deligne and N. Katz, will
appear later.

\begin{flushright}
Bures-sur-Yvette, May 1972,\\
P. DELIGNE
\end{flushright}


\markboth{General introduction}{General introduction}

\ifnum\SGAform=1
\clearpage\section*{Critical apparatus}\begingroup\small\raggedright\sloppy
Readings are given in the exact TeX notation of the two English witnesses. Decimal suffixes identify the individual differences in each marked passage.

\par\medskip\noindent\hypertarget{app-E7I-Introduction-001.1}{}\textbf{E7I-Introduction-001.1}\quad printed folio V.\par
\noindent Printed reading: {\ttfamily\footnotesize 5\allowbreak{}}\par
\noindent Corrected reading: {\ttfamily\footnotesize 6\allowbreak{}}\par
\noindent Source: {\ttfamily\footnotesize ED-EN-INTRO-I4.6;\allowbreak{} F7I-INTRO-003.2;\allowbreak{} fr/sga7i/\{diplomatic\allowbreak{},current\}/general\_in\allowbreak{}troduction\_body.tex\allowbreak{}}\par

\par\endgroup\normalsize\fi
\clearpage\SGAtoc{Exposé I}
\selectlanguage{english}
\SGApage{I}{1}


\SGARefAliasTarget{sga7.sga7i.ei.nroot}SGA 7

\underline{Expos\'e I}

\underline{\underline{Summary of the First Lectures by A. Grothendieck}}

\underline{written up by P. Deligne}

\underline{Contents}

0. Preliminaries

1. Arithmetic proof of the monodromy theorem

2. Vanishing cycles

3. Geometric proof of the monodromy theorem

4. Vanishing criteria for the sheaves of vanishing cycles

5. Action of monodromy on the $\pi_1$

6. Appendix by P. Deligne: arithmetic proof of the stable reduction theorem

0. \underline{Preliminaries}

\SGARefTarget{sga7.sga7i.ei.n0-0}{0.0}. \underline{Terminology}. Recall the following definitions:

(\SGARefTarget{sga7.sga7i.ei.n0-0-1}{0.0.1}) \underline{trait}: the spectrum of a discrete valuation ring. The generic and closed points of a trait are often denoted by $\eta$ and $s$.

(\SGARefTarget{sga7.sga7i.ei.n0-0-2}{0.0.2}) \underline{strictly local}: for a ring, local henselian with separably closed residue field; for a scheme, the spectrum of such a ring.

(\SGARefTarget{sga7.sga7i.ei.n0-0-3}{0.0.3}) \underline{geometric point of $S$} (in this expos\'e): a morphism $\bar{x}:\operatorname{Spec}(\bar{k})\longrightarrow S$ with image $x\in S$, where $k(\bar{x})\underset{\mathrm{dfn}}{=}\bar{k}$ is a separable closure of $k(x)$. For a henselian trait $S=\operatorname{Spec}(V)$, a generic geometric point of $S$ (that is, one with image $\eta$) is often denoted by $\bar{\eta}$. The residue field of the valuation ring $\bar{V}$, the integral closure of $V$ in $k(\bar{\eta})$, is a separably closed algebraic extension of $k(s)$ and defines a geometric point $\bar{s}$ of $S$ centered at $s$.


\SGApage{I}{\(\approx\)2}
(\SGARefTarget{sga7.sga7i.ei.n0-0-4}{0.0.4}) \underline{strict localization of $X$ at the geometric point $\bar{x}$}: SGA 4 VIII 4.4.

(\SGARefTarget{sga7.sga7i.ei.n0-0-5}{0.0.5}) \underline{essentially of finite type over $S$}: for a local (respectively, henselian local; respectively, strictly local) $S$-scheme, the spectrum of a local ring (respectively, the henselization of a local ring; respectively, a strict localization) of a scheme of finite type over $S$.

(\SGARefTarget{sga7.sga7i.ei.n0-0-6}{0.0.6}) $F\longmapsto F(n)$: Tate twist.

(\SGARefTarget{sga7.sga7i.ei.n0-0-7}{0.0.7}) An endomorphism $T$ of a finite-dimensional vector space is \underline{quasi-unipotent} if its eigenvalues are roots of unity, that is, if there exist $N\geqslant 1$ and $M\geqslant 1$ such that $(T^N-1)^M=0$.



\SGARefTarget{sga7.sga7i.ei.n0-1}{0.1}. \underline{Review of $\ell$-adic cohomology}

Let $X$ be a scheme of finite type over an algebraically closed field $k$ of characteristic exponent $p$, and let $\ell$ be a prime different from $p$.

The $\ell$-adic cohomology groups of $X$ were defined in SGA 4 and SGA 5.

a) $H^i(X,\mathbb{Z}/\ell^n)$ is the $i$-th cohomology group of the \`etale site $X_{\mathrm{et}}$ of $X$ with coefficients in $\mathbb{Z}/\ell^n$;

b) by definition,
\[
H^i(X,\mathbb{Z}_\ell)=\varprojlim_n H^i(X,\mathbb{Z}/\ell^n);
\]

c) by definition,
\[
H^i(X,\mathbb{Q}_\ell)=H^i(X,\mathbb{Z}_\ell)\otimes_{\mathbb{Z}_\ell}\mathbb{Q}_\ell.
\]

The $\ell$-adic cohomology groups of $X$ with proper supports are defined in the same way from the groups $H^i_c(X,\mathbb{Z}/\ell^n)$ of SGA 4 XVII.

In each of the following cases, one knows how to prove that these groups $H^i(X)$ have reasonable finiteness properties:

(\SGARefTarget{sga7.sga7i.ei.n0-1-1}{0.1.1}) for cohomology with proper supports;

(\SGARefTarget{sga7.sga7i.ei.n0-1-2}{0.1.2}) when $X$ is proper over $k$;

(\SGARefTarget{sga7.sga7i.ei.n0-1-3}{0.1.3}) when $X$ is smooth over $k$ (by (\SGARefLink{sga7.sga7i.ei.n0-1-1}{0.1.1}) and Poincar\'e duality, SGA 4 XVIII);

\SGAvar{(\SGARefTarget{sga7.sga7i.ei.n0-1-4}{0.1.4}) when Hironoka's resolution of singularities is available for schemes of finite type over $k$ of dimension at most $\dim(X)$;

}{(\SGARefTarget{sga7.sga7i.ei.n0-1-4}{0.1.4}) when Hironaka's resolution of singularities is available for schemes of finite type over $k$ of dimension at most $\dim(X)$;

}{E7I-I-001}(\SGARefTarget{sga7.sga7i.ei.n0-1-5}{0.1.5}) when $i=0,1$.


\SGApage{I}{\(\approx\)3}
In all these cases, the groups $H^i(X,\mathbb{Z}/\ell^n)$ are finite, the groups $H^i(X,\mathbb{Z}_\ell)$ are finitely generated, the groups $H^i(X,\mathbb{Q}_\ell)$ are finite-dimensional, and (except perhaps in case (\SGARefLink{sga7.sga7i.ei.n0-1-5}{0.1.5})) there are split short exact sequences

\begin{equation*}
0 \longrightarrow H^i(X,\mathbb{Z}_\ell)\otimes \mathbb{Z}/\ell^n
\longrightarrow H^i(X,\mathbb{Z}/\ell^n)
\longrightarrow \Tor_1(H^{i+1}(X,\mathbb{Z}_\ell),\mathbb{Z}/\ell^n)
\longrightarrow 0.
\end{equation*}
Moreover, in each case the proof of the finiteness theorem also proves a generic constructibility theorem: if $k$ is the field of fractions of an integral noetherian scheme $S$, and $X$ is the generic fibre of a morphism $f_1:X_1\longrightarrow S$ of finite type, then there exists a nonempty open subset $U\subset S$ such that over $U$ the sheaves $R^i f_*\mathbb{Z}/\ell^n$ are locally constant and their formation commutes with every base change.

\SGARefTarget{sga7.sga7i.ei.n0-2}{0.2}. We no longer suppose that $k$ is algebraically closed. Let $\bar{k}$ be an algebraic closure of $k$ (or a separable closure, which amounts to the same thing here). We consider the ``geometric'' cohomology groups $H^i(X_{\bar{k}})$, where $X_{\bar{k}}$ is obtained from $X$ by extension of scalars from $k$ to $\bar{k}$. By transport of structure, the Galois group $\Gal(\bar{k}/k)$ acts on these groups. In $\mathbb{Q}_\ell$-cohomology, in any of the cases (\SGARefLink{sga7.sga7i.ei.n0-1-1}{0.1.1})--(\SGARefLink{sga7.sga7i.ei.n0-1-5}{0.1.5}), this gives continuous representations of $\Gal(\bar{k}/k)$ in a linear group $\mathrm{GL}(n,\mathbb{Q}_\ell)$.

\SGARefTarget{sga7.sga7i.ei.n0-3}{0.3}. \underline{Review of traits} (\SGARefLink{sga7.sga7i.ei.n0-0-1}{0.0.1})

For the proofs, we refer to [4]. Let $S$ be a henselian trait, and let $\eta$, $s$, $V$, $\bar V$, $\bar\eta$, and $\bar s$ be as in (\SGARefLink{sga7.sga7i.ei.n0-0-1}{0.0.1}) and (\SGARefLink{sga7.sga7i.ei.n0-0-3}{0.0.3}). Let $p$ denote the characteristic exponent of $k(s)$, and let $\Gal(\bar\eta/\eta)$ (respectively, $\Gal(\bar s/s)$) denote the Galois group $\Gal(k(\bar\eta)/k(\eta))$ (respectively, the corresponding group for $\bar s/s$). By transport of structure, $\Gal(\bar\eta/\eta)$ acts on $k(\bar s)$; hence there is a homomorphism from $\Gal(\bar\eta/\eta)$ onto $\Gal(\bar s/s)$ whose kernel is the \underline{inertia group} $I$. The subfield of $k(\bar\eta)$ fixed by $I$ is a maximal unramified extension of $k(\eta)$.

The profinite group $I$ has a largest pro-$p$ subgroup $P$. The quotient $I/P$ is the \underline{tame inertia group}. Canonically,

\begin{equation*}
I/P\simeq \varprojlim_n\mu_n(k(\bar s))
\underset{\mathrm{dfn}}{=}\widehat{\mathbb{Z}}(1)(k(\bar s)).
\end{equation*}


\SGApage{I}{\(\approx\)4}
If $\phi_n$ is the canonical map from $I/P$ to $\mu_n(\bar k)$ and $u$ is an element of $\bar V$ of valuation $1/n$, then, for $\sigma\in I$,

\SGAvar{\begin{equation*}
\sigma(u)\equiv\phi_n(u)\,u
\qquad (\text{modulo elements of valuation greater than }1/n).
\end{equation*}
In summary,

}{\begin{equation*}
\sigma(u)\equiv\phi_n(\sigma)\,u
\qquad (\text{modulo elements of valuation greater than }1/n).
\end{equation*}
In summary,

}{E7I-I-002}\begin{equation}\tag{\SGARefTarget{sga7.sga7i.ei.n0-3-1}{0.3.1}}
\begin{aligned}
\Gal(\bar\eta/\eta)&\supset I\supset P,\\
\Gal(\bar\eta/\eta)/I&\simeq\Gal(\bar s/s),\\
I/P&\simeq\widehat{\mathbb{Z}}(1)(k(\bar s)),\\
P&:\ \text{a pro-$p$ group}\quad(\{e\}\text{ if }p=1),
\end{aligned}
\end{equation}
and the action (by inner automorphisms in $\Gal(\bar\eta/\eta)$) of $\Gal(\bar\eta/\eta)/I$ on $I/P$ is identified with the natural action of $\Gal(\bar s/s)$ on $\widehat{\mathbb{Z}}(1)(k(\bar s))$.

\SGAvar{\SGARefTarget{sga7.sga7i.ei.n0-4}{0.4}. More generally, let $S$ be strictly local and regular, let $D$ be a regular divisor in $S$, and let $p'$ be the characteristic exponent of the generic point of $D$. Let $\eta$ be a geometric generic point of $S$, and let $\bar s$ be the resulting geometric point centered at the closed point $s$. By Abhyankar's lemma, the fundamental group $\pi_1(S-D,\bar\eta)$ has a filtration analogous to (\SGARefLink{sga7.sga7i.ei.n0-3-1}{0.3.1}):

}{\SGARefTarget{sga7.sga7i.ei.n0-4}{0.4}. More generally, let $S$ be strictly local and regular, let $D$ be a regular divisor in $S$, and let $p'$ be the characteristic exponent of the generic point of $D$. Let $\bar\eta$ be a geometric generic point of $S$, and let $\bar s$ be the resulting geometric point centered at the closed point $s$. By Abhyankar's lemma, the fundamental group $\pi_1(S-D,\bar\eta)$ has a filtration analogous to (\SGARefLink{sga7.sga7i.ei.n0-3-1}{0.3.1}):

}{E7I-I-003}\begin{equation}\tag{\SGARefTarget{sga7.sga7i.ei.n0-4-1}{0.4.1}}
\begin{aligned}
\pi_1(S-D,\bar\eta)&\supset I\supset P,\\
\pi_1(S-D,\bar\eta)/I&\simeq\Gal(\bar s/s),\\
I/P&\simeq\widehat{\mathbb{Z}}(1)(k(\bar s)),\\
P&:\ \text{with no quotient of order prime to }p'.
\end{aligned}
\end{equation}
The action by inner automorphisms of $\pi_1(S-D,\bar\eta)/I$ on $I/P$ is identified with the natural action of $\Gal(\bar s/s)$ on $\widehat{\mathbb{Z}}(1)(k(\bar s))$.

\SGARefTarget{sga7.sga7i.ei.n0-5}{0.5}. \underline{N\'eron desingularization}

The method of N\'eron desingularization gives the following result.

\underline{Lemma} \SGARefTarget{sga7.sga7i.ei.n0-5-1}{0.5.1}. \underline{Let $S\longrightarrow S_o$ be a morphism of henselian traits, with $S=\Spec(V)$ and $S_o=\Spec(V_o)$, and let $\pi$ be a uniformizer of $V_o$. Suppose that $\pi$ is a uniformizer of $V$, and that the extensions $k(s)/k(s_o)$ and $k(\eta)/k(\eta_o)$ are separable. Then $V$ is a filtered direct limit of smooth henselian $V_o$-algebras $B_i$ essentially of finite type over $V_o$.}


\SGApage{I}{\(\approx\)5}
To construct the desired $V_o$-algebras, one takes $V_o\subset B\subset V$, with $B$ of finite type, and desingularizes $\Spec(B)$ by N\'eron's method ([1], \S~4).

Let $S=\Spec(V)$ be a henselian trait.

(\SGARefTarget{sga7.sga7i.ei.n0-5-2}{0.5.2}) If $S$ is of equal characteristic, with uniformizer $\pi$, one may apply (\SGARefLink{sga7.sga7i.ei.n0-5-1}{0.5.1}) with $S_o=\mathbf{F}_p[\pi]_{(\pi)}$. The extension $s/s_o$ is separable because $\mathbf{F}_p$ is perfect; the extension $\eta/\eta_o$ is separable because $\pi$, being a uniformizer, is not a $p$-th power (if $p\neq1$), and hence belongs to a $p$-basis. It follows that $S$ is a limit of spectra of smooth henselian $\mathbf{F}_p[\pi]_{(\pi)}$-algebras essentially of finite type.

(\SGARefTarget{sga7.sga7i.ei.n0-5-3}{0.5.3}) If $S$ is complete and of mixed characteristic, with perfect residue field $k$, then $V$ is an Eisenstein extension of the ring of Witt vectors (that is, the Cohen ring) $W(k)$:

\begin{equation*}
V\simeq W(k)[\pi]/\left(\pi^n+\sum_{i=0}^{n-1}a_i\pi^i\right).
\end{equation*}
If the $a_i$ are perturbed slightly, one obtains an isomorphic extension. We may therefore suppose that the $a_i$ belong to $W(k')$, where $k'$ is of finite type over $\mathbf{F}_p$. Let $k_o$ be the perfect closure of $k'$ in $k$, and set

\SGApage{I}{4}
\begin{equation*}
V_o=W(k_o)[\pi]/\left(\pi^n+\sum_{i=0}^{n-1}a_i\pi^i\right).
\end{equation*}
Lemma \SGARefLink{sga7.sga7i.ei.n0-5-1}{0.5.1} applies to $V/V_o$, and the residue field of $V_o$ is a radicial extension of a field of finite type over $\mathbf{F}_p$.

\SGApage{I}{\(\approx\)5}
1. \underline{Arithmetic proof of the monodromy theorem}

The reader will find in [5], Appendix, a proof of the following result of A. Grothendieck.

\underline{Proposition} \SGARefTarget{sga7.sga7i.ei.n1-1}{1.1}. \underline{Let $S$ be a henselian trait, let $\ell$ be a prime different from the characteristic exponent $p$ of $k(s)$, and let $\rho$ be a continuous representation of $\Gal(\bar\eta/\eta)$ in $\mathrm{GL}(n,\mathbb{Q}_\ell)$}. \underline{Suppose that the following condition holds:}

$(*_\ell)$ \quad \underline{No finite extension of $k(s)$ contains all roots of unity of $\ell$-power order.}


\SGApage{I}{\(\approx\)6}
\underline{Then there exists an open subgroup $I_1$ of the inertia group $I$ such that $\rho(\sigma)$ is unipotent for every $\sigma\in I_1$.}

Condition $(*_\ell)$ is automatically satisfied when $k(s)$ is of finite type over the prime field (or is radicial over such a field).

The following theorem follows immediately from \SGARefLink{sga7.sga7i.ei.n1-1}{1.1}.

\SGAvar{\underline{Monodromy theorem} \SGARefTarget{sga7.sga7i.ei.n1-2}{1.2}. \underline{With the preceding notation, let $X$ be a scheme of finite type over $\eta$, and let $H$ be a cohomology space of $X_{\bar\eta}$ with coefficients in $\mathbb{Q}_\ell$, of one of the types (\SGARefLink{sga7.sga7i.ei.n0-1-1}{1.1.1})--(\SGARefLink{sga7.sga7i.ei.n0-1-5}{1.1.5}). If condition $(*_\ell)$ is satisfied, the action on $H$ of every element of $I$ is quasi-unipotent.}

}{\underline{Monodromy theorem} \SGARefTarget{sga7.sga7i.ei.n1-2}{1.2}. \underline{With the preceding notation, let $X$ be a scheme of finite type over $\eta$, and let $H$ be a cohomology space of $X_{\bar\eta}$ with coefficients in $\mathbb{Q}_\ell$, of one of the types (\SGARefLink{sga7.sga7i.ei.n0-1-1}{0.1.1})--(\SGARefLink{sga7.sga7i.ei.n0-1-5}{0.1.5}). If condition $(*_\ell)$ is satisfied, the action on $H$ of every element of $I$ is quasi-unipotent.}

}{E7I-I-004}The limit arguments of \SGARefLink{sga7.sga7i.ei.n0-5}{0.5} allow us to dispense with hypothesis $(*_\ell)$.

\underline{Variant} \SGARefTarget{sga7.sga7i.ei.n1-3}{1.3}. \underline{Condition $(*_\ell)$ in \SGARefLink{sga7.sga7i.ei.n1-2}{1.2} is superfluous.}

Let $H'_o$ be a $\mathbb{Z}_\ell$-cohomology space of $X_{\bar\eta}$ giving rise to $H$, and put $H_o=H'_o/\text{torsion}$. By hypothesis,
$H=H_o\otimes_{\mathbb{Z}_\ell}\mathbb{Q}_\ell$, and the action of $\Gal(\bar\eta/\eta)$ is induced by

\begin{equation*}
\rho:\Gal(\bar\eta/\eta)\longrightarrow\mathrm{GL}(H_o).
\end{equation*}
Let $I_{\ell'}$ be the largest prime-to-$\ell$ subgroup of $I$. Its image in $\mathrm{GL}(H_o)$ is an $\ell$-adic Lie group of order prime to $\ell$, and is therefore finite.

Make an extension of traits $S'\longrightarrow S$. This does not change $H_o$ (SGA 4 XVI 6.6), and $\rho(I')$ has finite index in $\rho(I)$, because $I'/P'$ has finite index in $I/P$, by what precedes. It therefore suffices to prove the theorem over $S'$; we may suppose that

a) $\Gal(\bar\eta/\eta)$ acts trivially on $H_o/\ell H_o$;

b) if $S$ is of mixed characteristic, then $S$ is complete with perfect residue field (EGA $\mathrm{O}_{III}$ 10.3.1).

Put $S=\Spec(V)$, define $V_o$ as in (\SGARefLink{sga7.sga7i.ei.n0-5-2}{0.5.2}) and (\SGARefLink{sga7.sga7i.ei.n0-5-3}{0.5.3}), and apply (\SGARefLink{sga7.sga7i.ei.n0-5-1}{0.5.1}). We have

\SGApage{I}{\(\approx\)7}
\begin{equation*}
V=\varinjlim_i B_i.
\end{equation*}


Here $S_i=\Spec(B_i)$ is essentially of finite type and smooth over $V_o$. Let $D_i\subset S_i$ be defined by the equation $\pi=0$. For sufficiently large $i$, $H_o$ comes from a twisted constant $\mathbb{Z}_\ell$-sheaf $\mathcal H_i$ on $S_i-D_i$ (by ``generic constructibility,'' (\SGARefLink{sga7.sga7i.ei.n0-1}{0.1})), and $\mathcal H_i/\ell\mathcal H_i$ is constant.

The filtrations (\SGARefLink{sga7.sga7i.ei.n0-3-1}{0.3.1}) and (\SGARefLink{sga7.sga7i.ei.n0-4-1}{0.4.1}) give rise to commutative diagrams:

\begin{equation}\tag{\SGARefTarget{sga7.sga7i.ei.n1-3-1}{1.3.1}}
\begin{tikzcd}[column sep=small,row sep=large]
\Gal(\bar\eta/\eta) \arrow[d] \& \supset \& I \arrow[d] \& \supset \& P \arrow[d] \\
\pi_1(S_i-D_i,\bar\eta) \& \supset \& I_i \& \supset \& P_i
\end{tikzcd}
\qquad\text{and}\qquad
\begin{tikzcd}[row sep=large]
I/P \arrow[r,dash,"\sim"] \arrow[d] \& \widehat{\mathbb Z}(1) \arrow[d,equal] \\
I_i/P_i \arrow[r,dash,"\sim"] \& \widehat{\mathbb Z}(1)
\end{tikzcd}
\end{equation}

\SGAvar{The representation on $H_o$ of $\Gal(\bar\eta/\eta)$ is induced by an action of $\pi_1(S_i-D_i,\bar\eta)$ on $H_o$ that is trivial on $H_o/\ell H_o$. The group
\[
\Ker\bigl(\mathrm{GL}(H)\longrightarrow\mathrm{GL}(H/\ell H)\bigr)
\]
is a pro-$\ell$ group. Thus $P_i$ acts trivially, and we apply the following variant of \SGARefLink{sga7.sga7i.ei.n1-1}{1.1}.

\underline{Variant} \SGARefTarget{sga7.sga7i.ei.n1-4}{1.4}. \underline{Retain the notation of \SGARefLink{sga7.sga7i.ei.n0-4}{0.4}, and let $\rho$ be a continuous representation of $\Gal(\bar\eta/\eta)$ in $\mathrm{Gal}(n,\mathbb{Q}_\ell)$. Suppose that $k(s)$ satisfies $(*_\ell)$ and that $\rho(P)$ is finite. Then the action of $I$ is quasi-unipotent.}



}{The representation on $H_o$ of $\Gal(\bar\eta/\eta)$ is induced by an action of $\pi_1(S_i-D_i,\bar\eta)$ on $H_o$ that is trivial on $H_o/\ell H_o$. The group
\[
\Ker\bigl(\mathrm{GL}(H_o)\longrightarrow\mathrm{GL}(H_o/\ell H_o)\bigr)
\]
is a pro-$\ell$ group. Thus $P_i$ acts trivially, and we apply the following variant of \SGARefLink{sga7.sga7i.ei.n1-1}{1.1}.

\underline{Variant} \SGARefTarget{sga7.sga7i.ei.n1-4}{1.4}. \underline{Retain the notation of \SGARefLink{sga7.sga7i.ei.n0-4}{0.4}, and let $\rho$ be a continuous representation of $\Gal(\bar\eta/\eta)$ in $\mathrm{GL}(n,\mathbb{Q}_\ell)$. Suppose that $k(s)$ satisfies $(*_\ell)$ and that $\rho(P)$ is finite. Then the action of $I$ is quasi-unipotent.}



}{E7I-I-005}2. \underline{Vanishing cycles}

\SGARefTarget{sga7.sga7i.ei.n2-1}{2.1}. Let $S$ be a strictly local trait, and retain the notation of \SGARefLink{sga7.sga7i.ei.n0-3}{0.3}. Thus, by hypothesis, $s=\bar s$. Let $f:X\longrightarrow S$ be a scheme of finite type over $S$. The formalism of vanishing cycles is a tool for studying the relation between the cohomology of $X_s$ and that of $X_{\bar\eta}$, as well as the action of the inertia group on the latter, using information about $f$ that is \underline{local} on $X$.

A key result of this kind is the specialization theorem (SGA 4 XVI 2.2), according to which $X_s$ and $X_{\bar\eta}$ have the same cohomology when $f$ is proper and smooth.

In the transcendental theory, if $f:X\longrightarrow D$ is a proper morphism from an analytic space $X$ to the disk $D$, vanishing cycles arise as follows:

a) After shrinking $D$ if necessary, one may suppose that $X$ is a topological fibre bundle over the punctured disk $D^*$ (Thom's isotopy theorem). In particular, the fibres


\SGApage{I}{\(\approx\)8}
$X_t$ $(t\in D^*)$ are all isomorphic.

b) It is not unreasonable to picture the special fibre $X_s$ as obtained from \lq\lq the\rq\rq\ general fibre $X_t$ by contracting polyhedra in $X_t$; one would then have a map $\pi:X_t\longrightarrow X_s$. The method of vanishing cycles consists in studying the difference between the cohomology of $X_s$ and that of $X_t$ by means of the Leray spectral sequence for $\pi$.

The geometric theory explained below is very close to this transcendental theory: the general point $t\in D$ must be replaced by $\bar\eta$, the geometric generic point of $S$. Conversely, the transcendental theory can be modeled on the geometric theory by replacing $t$ or $\bar\eta$ with the universal cover $\widetilde D^*$ of $D^*$, and $X_t$ or $X_{\bar\eta}$ with $X\times_D\widetilde D^*$. This eliminates many $\varepsilon$'s and $\eta$'s, but sometimes loses information.

\SGARefTarget{sga7.sga7i.ei.n2-2}{2.2}. Fix a torsion ring $\Lambda$ in which $p$ is invertible (for example, $\Lambda=\mathbb Z/\ell^n\mathbb Z$, with $\ell$ prime to $p$). Consider the morphisms

\SGAvar{\begin{center}
\begin{tikzcd}[column sep=large,row sep=huge]
X_s \arrow[r,"i"] \arrow[d,"f_s"] \& X \arrow[d,"f"] \& \& X_{\bar\eta} \arrow[ll,"\bar j"] \arrow[d,"f_{\bar\eta}"] \\
s \arrow[r] \& S \& \eta \arrow[l] \& \bar\eta \arrow[l]
\end{tikzcd}
\end{center}
The sheaves of vanishing cycles $\psi_i$ on $X_s$ are

}{\begin{center}
\begin{tikzcd}[column sep=large,row sep=huge]
X_s \arrow[r,"i"] \arrow[d,"f_s"] \& X \arrow[d,"f"] \& \& X_{\bar\eta} \arrow[ll,"\bar j"] \arrow[d,"f_{\bar\eta}"] \\
s \arrow[r] \& S \& \eta \arrow[l] \& \bar\eta \arrow[l]
\end{tikzcd}
\end{center}
The sheaves of vanishing cycles $\psi^i$ on $X_s$ are

}{E7I-I-006}\begin{equation*}\tag{\SGARefTarget{sga7.sga7i.ei.n2-2-1}{2.2.1}}
\psi^i=i^*R^i\bar j_*\Lambda.
\end{equation*}
If $f$ is \underline{proper}, then (SGA 4 XII 5.5)

\SGAvar{\begin{equation*}\tag{\SGARefTarget{sga7.sga7i.ei.n2-2-2}{2.2.2}}
i^*:H^p(X,R^q\bar j_*\Lambda)\xrightarrow{\ \sim\ }
H^p(X_s,i^*R^q j_*\Lambda),
\end{equation*}
and the Leray spectral sequence for $\bar j$ becomes

}{\begin{equation*}\tag{\SGARefTarget{sga7.sga7i.ei.n2-2-2}{2.2.2}}
i^*:H^p(X,R^q\bar j_*\Lambda)\xrightarrow{\ \sim\ }
H^p(X_s,i^*R^q\bar j_*\Lambda),
\end{equation*}
and the Leray spectral sequence for $\bar j$ becomes

}{E7I-I-007}\begin{equation*}\tag{\SGARefTarget{sga7.sga7i.ei.n2-2-3}{2.2.3}}
H^p(X_s,\psi^q)\Longrightarrow H^{p+q}(X_{\bar\eta},\Lambda).
\end{equation*}
When nothing happens at infinity, this spectral sequence can sometimes be obtained without a properness hypothesis. That some hypothesis at infinity is necessary is already clear from the trivial case $X_s=\varnothing$.


\SGApage{I}{\(\approx\)9}
The definitions immediately give:

\underline{Proposition} \SGARefTarget{sga7.sga7i.ei.n2-3}{2.3}. \underline{Let $\bar x$ be a geometric point of $X_s$ (for example, a closed point), let $X_{(\bar x)}$ be the strict henselization of $X$ at $\bar x$, and let $X_{(\bar x)\bar\eta}$ be the geometric generic fibre of the projection $X_{(\bar x)}\longrightarrow S$. Then}

\begin{equation*}
\psi^i_{\bar x}=H^i(X_{(\bar x)\bar\eta},\Lambda).
\end{equation*}

\underline{Corollary} \SGARefTarget{sga7.sga7i.ei.n2-4}{2.4}. \underline{Wherever $f$ is smooth, $\psi^i=0$ for $i>0$ and $\psi^0=\Lambda$.}

This is SGA 4 XV 2.1.

The group $\Gal(\bar\eta/\eta)$ (which, by hypothesis, equals the inertia group) acts on the $\psi^i$ by transport of structure, and

\underline{Scholium} \SGARefTarget{sga7.sga7i.ei.n2-5}{2.5}. \underline{The spectral sequence (\SGARefLink{sga7.sga7i.ei.n2-2-3}{2.2.3}) is $\Gal(\bar\eta/\eta)$-equivariant.}

\underline{Variant} \SGARefTarget{sga7.sga7i.ei.n2-6}{2.6}. Put $\bar S=\Spec(\bar V)$ and $\bar X=X\times_S\bar S$. The scheme $X_{\bar\eta}$ is the complement of $X_s=X_{\bar s}$ in $\bar X$. Put

\begin{equation*}
\phi^i=\underline{H}^{\,i+1}_{X_s}(\bar X,\Lambda)
\end{equation*}
(the local-cohomology sheaves). There is an exact sequence of sheaves

\begin{equation*}
0\longrightarrow\phi^{-1}\longrightarrow\Lambda\longrightarrow\psi^0
\longrightarrow\phi^0\longrightarrow0,
\end{equation*}
and, for $i>0$, isomorphisms $\psi^i\xrightarrow{\ \sim\ }\phi^i$. If $f$ is smooth, the $\phi^i$ vanish.

\SGAvar{Put $\phi^i_{\mathrm{gl}}=H^{i}_{X_s}(\bar X,\Lambda)$, the global analogue of $\phi^i$. There is a spectral sequence

}{Put $\phi^i_{\mathrm{gl}}=H^{i+1}_{X_s}(\bar X,\Lambda)$, the global analogue of $\phi^i$. There is a spectral sequence

}{E7I-I-008}\begin{equation*}\tag{\SGARefTarget{sga7.sga7i.ei.n2-6-1}{2.6.1}}
H^p(X_s,\phi^q)\Longrightarrow\phi^{p+q}_{\mathrm{gl}},
\end{equation*}
and a long exact cohomology sequence which, when $f$ is proper, becomes

\begin{equation*}\tag{\SGARefTarget{sga7.sga7i.ei.n2-6-2}{2.6.2}}
\cdots\longrightarrow H^i(X_s,\Lambda)\longrightarrow H^i(X_{\bar\eta},\Lambda)
\longrightarrow\phi^i_{\mathrm{gl}}\longrightarrow\cdots.
\end{equation*}

\underline{Variant} \SGARefTarget{sga7.sga7i.ei.n2-7}{2.7}. Let $K_t$ be the maximal tamely ramified extension of $K$ in $\bar K$. Put $X_t=X\times_S\Spec(K_t)$, and let $\bar j_t$ be the projection from $X_t$ to $X$. Define tame variants of the $\psi^i$ by


\SGApage{I}{\(\approx\)10}
\begin{equation*}
\psi^i_t=R^i\bar j_{t*}\Lambda.
\end{equation*}
Since $P$ is a pro-$p$ group and $p$ is invertible in $\Lambda$, one obtains

\begin{equation}\tag{\SGARefTarget{sga7.sga7i.ei.n2-7-1}{2.7.1}}
H^*(X_t,\Lambda)=H^*(X_{\bar\eta},\Lambda)^P,
\end{equation}
\begin{equation}\tag{\SGARefTarget{sga7.sga7i.ei.n2-7-2}{2.7.2}}
\psi^i_t=(\psi^i)^P.
\end{equation}
When $f$ is proper, the Leray spectral sequence for $\bar j_t$ is obtained from (\SGARefLink{sga7.sga7i.ei.n2-2-3}{2.2.3}) by passing to $P$-invariants:

\begin{equation}\tag{\SGARefTarget{sga7.sga7i.ei.n2-7-3}{2.7.3}}
H^p(X_s,\psi^q_t)\Longrightarrow H^{p+q}(X_t,\Lambda)
=H^{p+q}(X_{\bar\eta},\Lambda)^P.
\end{equation}



3. \underline{Geometric proof of the monodromy theorem}

Let $\Lambda$ be as in \SGARefLink{sga7.sga7i.ei.n2-2}{2.2}.

\SGAvar{\underline{Conjecture} \SGARefTarget{sga7.sga7i.ei.n3-1}{3.1} (\underline{purity}). \underline{Let $A$ be an
excellent regular strictly local ring} (\SGARefLink{sga7.sga7i.ei.n0-0-2}{0.2}) \underline{and let $D$ be a
regular divisor in $\Spec(A)$ defined by a parameter $x$. Put
$U=\Spec(A)-D=\Spec(A[1/X])$. Then}

}{\underline{Conjecture} \SGARefTarget{sga7.sga7i.ei.n3-1}{3.1} (\underline{purity}). \underline{Let $A$ be an
excellent regular strictly local ring} (\SGARefLink{sga7.sga7i.ei.n0-0-2}{0.0.2}) \underline{and let $D$ be a
regular divisor in $\Spec(A)$ defined by a parameter $x$. Put
$U=\Spec(A)-D=\Spec(A[1/x])$. Then}

}{E7I-I-009}\begin{equation}\tag{\SGARefTarget{sga7.sga7i.ei.n3-1-1}{3.1.1}}
H^{q}(U,\Lambda)=
\begin{cases}
\Lambda & \text{if }q=0,\\[4pt]
\Lambda(-1) & \text{if }q=1\quad\text{(notation (\SGARefLink{sga7.sga7i.ei.n0-0-6}{0.0.6})),}\\[4pt]
0 & \text{if }q>1.
\end{cases}
\end{equation}

Here are some cases in which this conjecture is known (SGA 4 XIX): for
$q=0$ or $1$; when $A$ has characteristic $0$; when $A$ has equal
characteristic $p$ and resolution of singularities is available in
characteristic $p$ in dimensions $\leqslant\dim(A)$; or when
$(\Spec(A),D)$ is a smooth pair (the relative purity theorem, SGA 4 XVI
3.7).

When purity is available and
$D=\sum_{i\in B}D_i$ is a normal-crossings divisor in $\Spec(A)$ defined
by part of a regular system of parameters, the cohomology of
$U=\Spec(A)-D$ is given by


\begin{equation}\tag{\SGARefTarget{sga7.sga7i.ei.n3-1-2}{3.1.2}}
\begin{aligned}
H^0(U,\Lambda)&=\Lambda,\\[4pt]
H^1(U,\Lambda)&=\Lambda(-1)^B,\\[4pt]
H^q(U,\Lambda)&=\bigwedge\nolimits^q H^1(U,\Lambda)
  \qquad\text{(via the cup product).}
\end{aligned}
\end{equation}
These formulas are identical to those which, over $\mathbb C$, give the
cohomology of $D^{*B}\times D^k$.

\SGApage{I}{11}
\SGARefTarget{sga7.sga7i.ei.n3-2}{3.2}. Let $S$ be a strictly local trait, and retain the notation of
(\SGARefLink{sga7.sga7i.ei.n0-0-1}{0.0.1}) and (\SGARefLink{sga7.sga7i.ei.n0-0-3}{0.0.3}). Let $f:X\longrightarrow S$ be a morphism of finite
type. Suppose that $X$ is regular, that $X_\eta$ is smooth, and that
$X_s$ is a normal-crossings divisor in $X$. We shall suppose that this
divisor is, globally and not merely locally for the \'{e}tale topology, a
sum of regular divisors:
\[
(X_s)_{\mathrm{red}}=\sum_{i\in B}D_i
\]
(the general case can be recovered by localization). Denote their
multiplicities by $m_i$, so that
$X_s=\sum_{i\in B}m_iD_i$.

Let $x$ be a point of $X_s$, let $\bar x$ be a geometric point of $X$
lying over $x$ (for example, $x$ may be closed and equal to $\bar x$),
put $C=\{i\mid x\in D_i\}$, and let $K$ be the complex concentrated in
degrees $0$ and $-1$
\begin{equation*}
K:\mathbf Z^C\xrightarrow{\ d\ }\mathbf Z,
\qquad d((n_j))=\sum_jm_jn_j.
\end{equation*}

\underline{Theorem} \SGARefTarget{sga7.sga7i.ei.n3-3}{3.3} (\underline{uses purity}). \underline{The groups
of tame vanishing cycles $\psi^i_{t\bar x}$ are as follows:}

(i) \underline{There is a canonical isomorphism of graded
$\Lambda$-algebras}
\begin{equation*}
\psi^*_{t\bar x}\simeq
\bigwedge\nolimits^*\bigl(H_1(K)\otimes\Lambda(-1)\bigr)
\otimes_\Lambda\psi^0_{t\bar x}.
\end{equation*}

(ii) $\psi^0_{t\bar x}$ \underline{is the} $\Lambda$-\underline{algebra}
$\Lambda^E$, \underline{where $E$ is a set on which the tame inertia
group} $I_t=I/P=\widehat{\mathbf Z}(1)(k(s))$
\underline{acts transitively, and whose cardinality is the largest
prime-to-$p$ divisor of $\#H_0(K)$.}

Here is the method of proof.

\SGAvar{a) Let $X_{(\bar x)}$ be the strict localization of $X$ at $\bar x$,
let $U=X_{(\bar x)}-X_{(\bar x)s}$ (the complement of a
normal-crossings divisor), let $t_j$ be a local equation for $D_j$, put
\[
X_n=X_{( x)}\bigl[(t_j^{1/n})\bigr]
\qquad ((n,p)=1),
\]
which is a regular scheme, and let $U_n$ be the inverse image of $U$ in
$X_n$. By (\SGARefLink{sga7.sga7i.ei.n3-1-2}{3.1.2}),
\[
H^q(U_n,\Lambda)=\bigwedge\nolimits^q\bigl(\Lambda(-1)^C\bigr).
\]
For $m=xd$, the pullback map from $H^q(U_n,\Lambda)$ to
$H^q(U_m,\Lambda)$ is multiplication


}{a) Let $X_{(\bar x)}$ be the strict localization of $X$ at $\bar x$,
let $U=X_{(\bar x)}-X_{(\bar x)s}$ (the complement of a
normal-crossings divisor), let $t_j$ be a local equation for $D_j$, put
\[
X_n=X_{(\bar x)}\bigl[(t_j^{1/n})\bigr]
\qquad ((n,p)=1),
\]
which is a regular scheme, and let $U_n$ be the inverse image of $U$ in
$X_n$. By (\SGARefLink{sga7.sga7i.ei.n3-1-2}{3.1.2}),
\[
H^q(U_n,\Lambda)=\bigwedge\nolimits^q\bigl(\Lambda(-1)^C\bigr).
\]
For $m=nd$, the pullback map from $H^q(U_n,\Lambda)$ to
$H^q(U_m,\Lambda)$ is multiplication


}{E7I-I-010}\SGApage{I}{\(\approx\)12}
by $d^q$. Putting
$\widetilde U=\varprojlim_{(n,p)=1}U_n$, one therefore has
\begin{equation}\tag{\SGARefTarget{sga7.sga7i.ei.n3-3-1}{3.3.1}}
\begin{cases}
H^0(\widetilde U,\Lambda)=\Lambda,\\[6pt]
H^q(\widetilde U,\Lambda)=\displaystyle\varinjlim_n
  H^q(U_n,\Lambda)=0 & \text{if }q\ne0.
\end{cases}
\end{equation}

b) $\widetilde U$ is a pro-covering of $U$ with group
$\widehat{\mathbf Z}(1)(k(s))^C$. By definition,
\[
\psi^*_{t\bar x}=H^*(U_{\eta_t},\Lambda).
\]
$\widetilde U$ is a pro-covering, with group
$H_1(K)\otimes\widehat{\mathbf Z}(1)(k(s))$, of each connected component
of $U_{\eta_t}$, and these components are permuted transitively by the
tame inertia group. Formula (\SGARefLink{sga7.sga7i.ei.n3-3-1}{3.3.1}) and the Hochschild--Serre spectral
sequence now give the formulas in \SGARefLink{sga7.sga7i.ei.n3-3}{3.3}.

Pictorially, one may say that $U_{\bar\eta}$ is the fibre of a morphism
$\delta$ from the cohomological
$K(\widehat{\mathbf Z}(1)(k(s))^C,1)$ represented by $U$ to the
cohomological $K(\widehat{\mathbf Z}(1)(k(s)),1)$ represented by $\eta$;
the morphism $\delta$ induces
\[
d:(n_j)\longmapsto\sum_jm_jn_j:
\widehat{\mathbf Z}(1)(k(s))^C
\longrightarrow\widehat{\mathbf Z}(1)(k(s)).
\]

\SGAvar{\underline{Corollary} \SGARefTarget{sga7.sga7i.ei.n3-4}{3.4} (\underline{uses purity}). \underline{Suppose
that $f$ is proper. Let $N$ be the least common multiple of the
multiplicities $m_j$ $(j\in B)$. Let $q\geqslant0$, and let $q'$ be the
minimum of $q$ and the greatest number of divisors $D_j$ passing through
one point. For $T$ in the tame inertia group $I_t=I/P$, one has}
\begin{equation*}
(T^N-1)^{q'}=0
\quad\underline{\text{on}}\quad
H^q(X_{\bar\eta},\Lambda)^P.
\end{equation*}
This follows at once from \SGARefLink{sga7.sga7i.ei.n3-3}{3.3}, Scholium \SGARefLink{sga7.sga7i.ei.n2-5}{2.5}, and Variant \SGARefLink{sga7.sga7i.ei.n2-7}{2.7}.

}{\underline{Corollary} \SGARefTarget{sga7.sga7i.ei.n3-4}{3.4} (\underline{uses purity}). \underline{Suppose
that $f$ is proper. Let $N$ be the least common multiple of the
multiplicities $m_j$ $(j\in B)$. Let $q\geqslant0$, and let $q'$ be the
minimum of $q$ and the greatest number of divisors $D_j$ passing through
one point. For $T$ in the tame inertia group $I_t=I/P$, one has}
\begin{equation*}
(T^N-1)^{q'+1}=0
\quad\underline{\text{on}}\quad
H^q(X_{\bar\eta},\Lambda)^P.
\end{equation*}
This follows at once from \SGARefLink{sga7.sga7i.ei.n3-3}{3.3}, Scholium \SGARefLink{sga7.sga7i.ei.n2-5}{2.5}, and Variant \SGARefLink{sga7.sga7i.ei.n2-7}{2.7}.

}{E7I-I-011}The required purity results are available for $H^1$, in characteristic
$0$, or when $S$ is essentially of finite type over a field.

\underline{Theorem} \SGARefTarget{sga7.sga7i.ei.n3-5}{3.5}. \underline{Let $C$ be a smooth projective curve
over the fraction field $k(\eta)$ of a discrete valuation ring $V$. The
inertia group $I$ acts on $H^1(C_{\bar\eta},\mathbf Q_\ell)$ by
quasi-unipotent automorphisms of level $2$: there exists $N$ such that,
for $T\in I$,}
\begin{equation*}
(T^N-1)^2=0
\quad\text{on}\quad H^1(C_{\bar\eta},\mathbf Q_\ell).
\end{equation*}
We reduce to the case in which $V$ is strictly local. The assertion is
that there exists an open subgroup $I_1$ of $I$ such that
$(T-1)^2=0$ on $H^1$ for $T\in I_1$.


\SGApage{I}{\(\approx\)13}
We may extend the scalars of $k(\eta)$ to a finite extension, replacing
$V$ by its normalization in that extension. Since the image of the
pro-$p$ group $P$ in an $\ell$-adic Lie group
$\mathrm{GL}(n,\mathbf Q_\ell)$ is finite, we may in particular suppose
that $P$ acts trivially.

Let $C_o$ be a minimal model of $C$ over $V$ (for the existence of
$C_o$, which is based on Abhyankar's theorem, see the discussion in [3],
\S~2, p.~87; one could instead reduce to the case where $V$ is excellent
by passing to its completion and invoking SGA 4 XVI 1.6).

Let $C_1$ be a regular model of $C$, obtained from $C_o$ by blowups, whose
special fibre is a normal-crossings divisor. The arguments of \SGARefLink{sga7.sga7i.ei.n3-3}{3.3} and
\SGARefLink{sga7.sga7i.ei.n3-4}{3.4} apply to $C_1$, $\psi_t^0$, $\psi_t^1$, and
\[
H^1(C_{\bar\eta},\mathbb Z/\ell^n)^P
=H^1(C_{\bar\eta},\mathbb Z/\ell^n),
\]
with $N$ independent of $n$; Theorem \SGARefLink{sga7.sga7i.ei.n3-5}{3.5} follows.

\SGAvar{\underline{Remark} \SGARefTarget{sga7.sga7i.ei.n3-6}{3.6}. Let $A$ be an abelian variety over $k(\eta)$.
Since $A$ is a quotient of a Jacobian, it follows again that the action
of $I$ on $H^1(A_{\eta},\mathbb Q_\ell)$, or on $T_\ell(A)$, is
quasi-unipotent of level $2$.

\underline{Remark} \SGARefTarget{sga7.sga7i.ei.n3-7}{3.7}. The action of $I$ on
$H^1(X_{\bar\eta},\mathbb Q)$ is in fact quasi-unipotent of level
$2$ for every scheme $X$ of finite type over $\eta$.

}{\underline{Remark} \SGARefTarget{sga7.sga7i.ei.n3-6}{3.6}. Let $A$ be an abelian variety over $k(\eta)$.
Since $A$ is a quotient of a Jacobian, it follows again that the action
of $I$ on $H^1(A_{\bar\eta},\mathbb Q_\ell)$, or on $T_\ell(A)$, is
quasi-unipotent of level $2$.

\underline{Remark} \SGARefTarget{sga7.sga7i.ei.n3-7}{3.7}. The action of $I$ on
$H^1(X_{\bar\eta},\mathbb Q_\ell)$ is in fact quasi-unipotent of level
$2$ for every scheme $X$ of finite type over $\eta$.

}{E7I-I-012}\SGApage{I}{12}
Theorem \SGARefLink{sga7.sga7i.ei.n3-5}{3.5} (or Remark \SGARefLink{sga7.sga7i.ei.n3-6}{3.6}) is the basis for the proof of the stable
reduction theorem for abelian varieties that will be given in Expos\'{e}
IX, \S~3.

\SGApage{I}{13}
\SGARefTarget{sga7.sga7i.ei.n3-8}{3.8}. It is clear that whenever purity and resolution of singularities
are available (for example, in characteristic $0$), the preceding method
proves the monodromy theorem and gives bounds for the nilpotence
exponent. Thus, in equal characteristic $0$, if $X$ is proper and smooth
over $\eta$ and $T$ belongs to an open subgroup $I_1$ of $I$, one has
\begin{equation*}
(T-1)^{q+1}=0
\quad\text{on}\quad H^q(X_{\bar\eta},\mathbb Z_\ell)
\qquad(T\in I_1\subset I).
\end{equation*}



4. \underline{Vanishing criteria for the sheaves of vanishing cycles}

\SGARefTarget{sga7.sga7i.ei.n4-1}{4.1}. Let $S$ be a strictly local trait and let
$f:X\longrightarrow S$ be a scheme of finite type over $S$. Put
$n=\dim(X_\eta)$.

If $X'$ is the scheme-theoretic closure of $X_\eta$ in $X$, then:

a) On $X_s-X'_s$, one has $\phi^{-1}=\Lambda$, and the other
$\phi^i$ vanish;


\SGApage{I}{\(\approx\)14}
b) on $X'_s$, the $\phi^i$ and $\psi^i$ relative to $X/S$ or to
$X'/S$ coincide, and $\phi^{-1}=0$.

Since $X'$ is flat over $S$, this often permits a reduction to the case
in which $X$ is flat over $S$.

\underline{Theorem} \SGARefTarget{sga7.sga7i.ei.n4-2}{4.2}.

(i) \underline{One has $\psi^i=0$ for $i>n$.}

(ii) \underline{More precisely}, \underline{let $x$ be a point of $X_s$ whose
closure has dimension $k$}
($k=\operatorname{trdeg}_{k(s)}k(x)$),
\underline{and let $\bar x$ be a geometric point lying over $x$. Then}
\begin{equation*}
\psi^i_{\bar x}=0
\quad\underline{\text{for}}\quad i>n-k,
\qquad\text{i.e.}\qquad
d(\psi^i)\leqslant n-i
\quad(\underline{\text{notation of SGA 4 XIV 2.1}}).
\end{equation*}
Indeed,
\[
\psi^i_{\bar x}=H^i(X_{(\bar x)\bar\eta},\Lambda),
\]
and $X_{(\bar x)\bar\eta}$ is an inverse limit of affine schemes of
dimension at most $n$ over $k(\bar\eta)$. Assertion (i) therefore follows
from the affine Lefschetz theorem (SGA 4 XIV 3.2) by passage to the
limit.

To prove (ii), we may suppose that $X$ is flat over $S$ (\SGARefLink{sga7.sga7i.ei.n4-1}{4.1}) and that
there is a quasi-finite flat $S$-morphism
\[
g:X\longrightarrow\mathbb A^k_S
\]
such that $g(x)$ is the generic point of $\mathbb A^k_s$. Let $S'$ be
the strict localization of $\mathbb A^k_S$ at $\bar x$. Then:

a) $S'_{\bar\eta}$ is the spectrum of a field, and
$\Gal(\bar\eta'/s'_{\bar\eta})$ is a $p$-group $Q$;

b) by (i), applied to $g$,
\[
H^i(X_{(\bar x)\bar\eta'},\Lambda)=0
\qquad\text{for }i>n-k;
\]

c) one has
\[
H^i(X_{(\bar x)\bar\eta},\Lambda)
=H^i(X_{(\bar x)\bar\eta'},\Lambda)^Q,
\]
which proves the theorem.

\underline{Corollary} \SGARefTarget{sga7.sga7i.ei.n4-3}{4.3}. \underline{If $f$ is proper and flat and the
non-smooth locus of $f$ in $X_s$ has dimension at most $d$, the
specialization morphism}
\[
H^i(X_s)\longrightarrow H^i(X_{\bar\eta})
\]
\underline{is an isomorphism for $i>n+d+1$ and an epimorphism for
$i=n+d+1$.}

\SGARefTarget{sga7.sga7i.ei.n4-4}{4.4}. To prove, in certain cases, that the $\phi^i$ vanish for small
$i$, we shall use local duality, either directly or through the local
Lefschetz theorems of SGA 2 XIV. This tool is available only in
characteristic $0$, or in equal characteristic $p$ when resolution of
singularities \`a la Hironaka is available in the dimensions under
consideration.


\SGApage{I}{15}
\underline{Theorem} \SGARefTarget{sga7.sga7i.ei.n4-5}{4.5}. \underline{Suppose that $S$ has characteristic
$0$. Let $f:X\longrightarrow S$ be a morphism of finite type and let
$x$ be a closed point of $X_s$. Suppose that:}

(a) \underline{$x$ is an isolated non-smooth point in its fibre;}

(b) \underline{$X$ has relative geometric depth at least $n$ at $x$:
in a neighbourhood of $x$, the scheme $X$ is identified with a
subscheme defined by $k$ equations in a smooth scheme of relative
dimension $N$ at $x$, with $n\leqslant N-k$.}

\underline{Then $\phi^i_x=0$ for $i<n$.}

\SGAvar{Let $X_x$ be the strict localization of $X$ at $x$, and put
$V=X_x-\{x\}$. For every trait $S'$ finite over $S$, put similarly
\[
X'_x=X_x\times_S S',\qquad V'=V\times_S S'.
\]
Let $ S$ be the normalization of $S$ in $\bar\eta$, and put
\[
\bar X_x=X_x\times_S\bar S,
\qquad \bar V=V\times_S\bar S.
\]
We have:

(1) If $\phi^i=0$ on $(X_x)_s$ (i.e. at every generization of $x$ in
$X_s$) for $i\leqslant m$, then
\begin{equation*}
H^{i-1}_{\{x\}}(\bar X_x)
=\widetilde H^{,i}(\bar V)
\xrightarrow{\ \sim\ }
\widetilde H^{,i}(\bar V_{\bar\eta})
=\phi^i
\qquad\text{for }i\leqslant m.
\end{equation*}
The $\phi^i$ are indeed the cohomology sheaves of $\bar V$ with support
in $V_s$; under the hypotheses of \SGARefLink{sga7.sga7i.ei.n4-5}{4.5}, condition (a) makes (a)
applicable with $m=\infty$ (SGA 4 XV 2.1).

}{Let $X_x$ be the strict localization of $X$ at $x$, and put
$V=X_x-\{x\}$. For every trait $S'$ finite over $S$, put similarly
\[
X'_x=X_x\times_S S',\qquad V'=V\times_S S'.
\]
Let $\bar S$ be the normalization of $S$ in $\bar\eta$, and put
\[
\bar X_x=X_x\times_S\bar S,
\qquad \bar V=V\times_S\bar S.
\]
We have:

(1) If $\phi^i=0$ on $(X_x)_s$ (i.e. at every generization of $x$ in
$X_s$) for $i\leqslant m$, then
\begin{equation*}
H^{i-1}_{\{x\}}(\bar X_x)
=\widetilde H^{,i}(\bar V)
\xrightarrow{\ \sim\ }
\widetilde H^{,i}(\bar V_{\bar\eta})
=\phi^i
\qquad\text{for }i\leqslant m.
\end{equation*}
The $\phi^i$ are indeed the cohomology sheaves of $\bar V$ with support
in $V_s$; under the hypotheses of \SGARefLink{sga7.sga7i.ei.n4-5}{4.5}, condition (a) makes (1)
applicable with $m=\infty$ (SGA 4 XV 2.1).

}{E7I-I-013}(2) Hypothesis (b) of \SGARefLink{sga7.sga7i.ei.n4-5}{4.5} is stable under base change. By SGA 2 XIV
5.6, it implies
\begin{equation*}
\widetilde H^{,i}(V')
=H^{i-1}_{\{x\}}(X'_x)=0
\qquad\text{for }i<n.
\end{equation*}
We conclude by observing that $\widetilde H^{,i}(\bar V)$ is the
filtered colimit of the $\widetilde H^{,i}(V')$.

Combining \SGARefLink{sga7.sga7i.ei.n4-2}{4.2} and \SGARefLink{sga7.sga7i.ei.n4-5}{4.5} gives:

\underline{Corollary} \SGARefTarget{sga7.sga7i.ei.n4-6}{4.6}. \underline{Let $S$ be a henselian trait of
characteristic $0$, let $X/S$ be a relative complete intersection of
relative dimension $n$, and let $x\in X_s$ be an isolated non-smooth
point in its fibre}. \underline{Then $\phi^i_x=0$ for $i\ne n$.}

\underline{Corollary} \SGARefTarget{sga7.sga7i.ei.n4-7}{4.7}. \underline{Under the hypotheses of \SGARefLink{sga7.sga7i.ei.n4-6}{4.6},
$\phi^n_x$ is a flat }$\Lambda$-\underline{module}.

\SGAvar{Indeed, if $\Lambda'$ is a finite $\Lambda$-algebra, a general nonsense
(SGA 4 XVII 5.2.11) gives, for the $\phi^i_{\Lambda}$,


}{Indeed, if $\Lambda'$ is a finite $\Lambda$-algebra, a general nonsense
(SGA 4 XVII 5.2.11) gives, for the $\phi^i_{\Lambda'}$,


}{E7I-I-014}\begin{equation*}
\phi^{n-i}_{\Lambda'}
=\Tor_i^\Lambda(\Lambda',\phi^n_\Lambda),
\end{equation*}
and we apply \SGARefLink{sga7.sga7i.ei.n4-6}{4.6} to $\Lambda'$.

\SGAvar{\underline{Variant} \SGARefTarget{sga7.sga7i.ei.n4-8}{4.8}. In \SGARefLink{sga7.sga7i.ei.n4-5}{4.5}, omit condition (a), and let $Z$ be the
non- locus of $f$. One may still conclude that
\begin{equation*}
\phi^i_x=0\qquad\text{for }i<n-\dim Z_s.
\end{equation*}
We proceed by descending induction on $\dim\overline{\{x\}}$. The
induction hypothesis permits us to apply \SGARefLink{sga7.sga7i.ei.n4-5}{4.5} (1) with
$m=n-\dim Z_s$, and the proof is completed as before.

}{\underline{Variant} \SGARefTarget{sga7.sga7i.ei.n4-8}{4.8}. In \SGARefLink{sga7.sga7i.ei.n4-5}{4.5}, omit condition (a), and let $Z$ be the
non-smooth locus of $f$. One may still conclude that
\begin{equation*}
\phi^i_x=0\qquad\text{for }i<n-\dim Z_s.
\end{equation*}
We proceed by descending induction on $\dim\overline{\{x\}}$. The
induction hypothesis permits us to apply \SGARefLink{sga7.sga7i.ei.n4-5}{4.5} (1) with
$m=n-\dim Z_s$, and the proof is completed as before.

}{E7I-I-015}\SGApage{I}{16}
\SGARefTarget{sga7.sga7i.ei.n4-9}{4.9}. It is possible to obtain a vanishing criterion using only
information about the geometric fibres of $f$. For complete
intersections, however, the result is one degree weaker than \SGARefLink{sga7.sga7i.ei.n4-5}{4.5}.

\underline{Proposition} \SGARefTarget{sga7.sga7i.ei.n4-10}{4.10}. \underline{Under the general hypotheses
of \SGARefLink{sga7.sga7i.ei.n4-5}{4.5}, suppose that:}

(a) \underline{$x$ is an isolated non-smooth point in its fibre;}

\SGAvar{(b) \underline{for every point $y$ of $X_{\bar\eta}$ whose closure in
$X_{\bar\eta}$ has dimension $d(y)$ and which is a generization of $x$,
the \'{e}tale depth satisfies}
\begin{equation*}
\operatorname{depth}_{\mathrm{\acute{e}t},y}(X_{\bar\eta})
\geqslant n-d(g),
\quad\text{i.e.}\quad
H^i_{\{y\}}(X_{\bar\eta})=0
\quad\underline{\text{for }}i\leqslant n-d(g);
\end{equation*}

}{(b) \underline{for every point $y$ of $X_{\bar\eta}$ whose closure in
$X_{\bar\eta}$ has dimension $d(y)$ and which is a generization of $x$,
the \'{e}tale depth satisfies}
\begin{equation*}
\operatorname{depth}_{\mathrm{\acute{e}t},y}(X_{\bar\eta})
\geqslant n-d(y),
\quad\text{i.e.}\quad
H^i_{\{y\}}(X_{\bar\eta})=0
\quad\underline{\text{for }}i\leqslant n-d(y);
\end{equation*}

}{E7I-I-016}(c) \underline{$\operatorname{depth}_{\mathrm{\acute{e}t},x}(X_s)
\geqslant n$.}

\underline{Then $\phi^i_x=0$ for $i<n-1$.}

Retain the notation of \SGARefLink{sga7.sga7i.ei.n4-5}{4.5}. We may apply \SGARefLink{sga7.sga7i.ei.n4-5}{4.5} (1) with $m=\infty$.
By the local Lefschetz theorem (SGA 2 XIV), hypothesis (b) gives
\begin{center}
\begin{tikzcd}[row sep=large, column sep=huge]
H^i(V') \arrow[r, "\sim"] \& H^i(V_s) \& {(\ i<n\text{-}1\ )} \\
H^i(V') \arrow[r, hook] \& H^i(V_s) \& {(\ i=n\text{-}1\ )}
\end{tikzcd}
\end{center}
Condition (c) then gives
\[
H^i(\bar V)=\varinjlim H^i(V')=0
\qquad\text{for }i\leqslant n-1,
\]
which proves \SGARefLink{sga7.sga7i.ei.n4-10}{4.10}.

\SGARefTarget{sga7.sga7i.ei.n4-11}{4.11}. As in \SGARefLink{sga7.sga7i.ei.n4-8}{4.8}, one can give a variant of \SGARefLink{sga7.sga7i.ei.n4-10}{4.10} in which (a) is not
assumed.




\SGApage{I}{\(\approx\)17}
5. \underline{Action of monodromy on $\pi_1$}

In this section (Grothendieck's expos\'{e} of 19 November 1968), we
shall use the semistable reduction theorem for curves. Artin and Winters
[2] gave a direct proof of this theorem; it can also be deduced [3] from
the analogous theorem for abelian varieties (IX \S~3, using \SGARefLink{sga7.sga7i.ei.n3-6}{3.6}). We
shall also use Schlessinger's theory (Rim's Expos\'{e} VI).

\SGARefTarget{sga7.sga7i.ei.n5-1}{5.1}. Let $k$ be an algebraically closed field, let $C$ be a projective
curve over $k$, and let $x_1,\ldots,x_n$ be a finite set of closed points
of $C$. The pointed curve
\[
X=(C;x_1,\ldots,x_n)
\]
is called \underline{stable} if:

(a) $C$ is reduced, and its only singularities are double points with
distinct tangents;

(b) the $x_i$ are distinct smooth points of $C$;

\SGAvar{(c) if $C_1$ is an irreducible component of $C$, and $E$ is the set of
points of the normalization $C'_1$ of $C_1$ lying above a singular point
of $C_1$ or above one of the $x_i$, then, if $C'_1$ has genus $0$
(respectively $1$), the set $E$ has at least $3$ (respectively $1$)
elements.

}{(c) if $C_1$ is an irreducible component of $C$, and $E$ is the set of
points of the normalization $C'_1$ of $C_1$ lying above a singular point
of $C$ or above one of the $x_i$, then, if $C'_1$ has genus $0$
(respectively $1$), the set $E$ has at least $3$ (respectively $1$)
elements.

}{E7I-I-017}Condition (c) is equivalent to each of the following:

(c') $X$ has no nontrivial infinitesimal automorphisms.

(c'') If $\omega$ is the invertible dualizing sheaf, the invertible
sheaf
\[
\omega'=\omega\Bigl(\sum_i x_i\Bigr)
\]
is ample.

An $S$-scheme $C$ equipped with sections $x_1,\ldots,x_n$ is a
\underline{stable curve} over $S$ if its geometric fibres are stable
curves. In [3], \S\S~1--2, only the case $n=0$ is considered. The
extension of the results of loc. cit. to the general case is immediate:

(1) If $X/k$ is stable, then
$H^1(C,\omega'^{\otimes n})=0$ for $n\geqslant2$, and
$\omega'^{\otimes n}$ is very ample for $n\geqslant3$ (cf. [3] 1.2).

(2) The formal moduli scheme $M$ of $X$ is smooth, and $C$ remains
singular over a normal-crossings divisor of $M$.


\SGApage{I}{\(\approx\)18}
(3) If $X$ and $Y$ are stable over $S$, then $\Isom_S(X,Y)$ is finite
and unramified over $S$ (isomorphisms are required to respect the marked
points).

(4) If $X/k$ is stable, the map
\[
\Aut_k(X)\longrightarrow\Aut_k\Pic^\circ(C)
\]
is injective.

Finally, we have:

\SGAvar{\underline{Proposition} \SGARefTarget{sga7.sga7i.ei.n5-1-1}{5.1.1}. \underline{Let
$X_\eta=(C_\eta;x_1,\ldots,x_n)$ be stable over the generic point
$\eta$ of a trait $S$, with $C_\eta$ smooth. There exist a finite
separable extension $\eta'$ of $\eta$ and a stable curve $X'$ over the
normalization $S'$ of $S$ in $\eta$ such that}
\begin{equation*}
X'\otimes_{S'}\eta'=X_\eta\otimes_\eta\eta'.
\end{equation*}

This follows from the usual semistable reduction theorem, which permits
one to choose $\eta'$ so that the geometric special fibre of the minimal
model $C''$ of $C_{\eta'}$ over $S'$ satisfies (a). One then blows up,
iteratively if necessary, \underline{smooth} points of $C''$ so that
(b) is satisfied. Next one contracts the chains of smooth rational
components of self-intersection two which contain none of the
$\bar x_i$ in the special fibre; the result is $C'$ (cf. [3], proof of
2.3, p.~88).

}{\underline{Proposition} \SGARefTarget{sga7.sga7i.ei.n5-1-1}{5.1.1}. \underline{Let
$X_\eta=(C_\eta;x_1,\ldots,x_n)$ be stable over the generic point
$\eta$ of a trait $S$, with $C_\eta$ smooth. There exist a finite
separable extension $\eta'$ of $\eta$ and a stable curve $X'$ over the
normalization $S'$ of $S$ in $\eta'$ such that}
\begin{equation*}
X'\otimes_{S'}\eta'=X_\eta\otimes_\eta\eta'.
\end{equation*}

This follows from the usual semistable reduction theorem, which permits
one to choose $\eta'$ so that the geometric special fibre of the minimal
model $C''$ of $C_{\eta'}$ over $S'$ satisfies (a). One then blows up,
iteratively if necessary, \underline{smooth} points of $C''$ so that
(b) is satisfied. Next one contracts the chains of smooth rational
components of self-intersection $-2$ which contain none of the
$\bar x_i$ in the special fibre; the result is $C'$ (cf. [3], proof of
2.3, p.~88).

}{E7I-I-018}As in loc. cit., one can check that $C'$ exists as soon as the Jacobian
of $C_\eta$ has semistable reduction. The smoothness hypothesis on
$C_\eta$ is moreover unnecessary.

\SGAvar{\SGARefTarget{sga7.sga7i.ei.n5-2}{5.2}. Let $X_\eta$ be a geometrically connected scheme of finite type
over the fraction field of a strictly local trait $S$. For a geometric
point $\xi$ of $X_{\bar\eta}$, there is an exact sequence
\begin{equation}\tag{\SGARefTarget{sga7.sga7i.ei.n5-2-1}{5.2.1}}
\begin{tikzcd}[column sep=large, row sep=tiny]
0 \arrow[r, dash] \& \pi_1(X_{\bar\eta},\xi) \arrow[r]
  \& \pi_1(X,\xi) \arrow[r]
  \& \Gal(\bar\eta/\eta) \arrow[r] \& 0 .
\end{tikzcd}
\end{equation}
It induces a homomorphism from $\Gal(\bar\eta/\eta)$ to the group of
outer automorphisms of $\pi_1(X_{\bar\eta})$. Passing to the largest
prime-to-$p$ quotient $\pi_1^{(p)}(X_{\bar\eta})$ of
$\pi_1(X_{\bar\eta})$, one obtains
\begin{equation}\tag{\SGARefTarget{sga7.sga7i.ei.n5-2-2}{5.2.2}}
\Gal(\bar\eta/\eta)
\longrightarrow
\Aut_{\mathrm{ext}}\bigl(\pi_1^{(p)}(X_{\bar\eta})\bigr).
\end{equation}

Suppose that $X_\eta$ has a rational point $x$, and let $\bar x$ be the
geometric point
\[
\bar\eta\longrightarrow\eta\xrightarrow{\ x\ }X_\eta.
\]
For $\xi=\bar x$, the exact sequence (\SGARefLink{sga7.sga7i.ei.n5-2-1}{5.2.2}) splits canonically, giving


}{\SGARefTarget{sga7.sga7i.ei.n5-2}{5.2}. Let $X_\eta$ be a geometrically connected scheme of finite type
over the fraction field of a strictly local trait $S$. For a geometric
point $\xi$ of $X_{\bar\eta}$, there is an exact sequence
\begin{equation}\tag{\SGARefTarget{sga7.sga7i.ei.n5-2-1}{5.2.1}}
\begin{tikzcd}[column sep=large, row sep=tiny]
0 \arrow[r] \& \pi_1(X_{\bar\eta},\xi) \arrow[r]
  \& \pi_1(X,\xi) \arrow[r]
  \& \Gal(\bar\eta/\eta) \arrow[r] \& 0 .
\end{tikzcd}
\end{equation}
It induces a homomorphism from $\Gal(\bar\eta/\eta)$ to the group of
outer automorphisms of $\pi_1(X_{\bar\eta})$. Passing to the largest
prime-to-$p$ quotient $\pi_1^{(p)}(X_{\bar\eta})$ of
$\pi_1(X_{\bar\eta})$, one obtains
\begin{equation}\tag{\SGARefTarget{sga7.sga7i.ei.n5-2-2}{5.2.2}}
\Gal(\bar\eta/\eta)
\longrightarrow
\Aut_{\mathrm{ext}}\bigl(\pi_1^{(p)}(X_{\bar\eta})\bigr).
\end{equation}

Suppose that $X_\eta$ has a rational point $x$, and let $\bar x$ be the
geometric point
\[
\bar\eta\longrightarrow\eta\xrightarrow{\ x\ }X_\eta.
\]
For $\xi=\bar x$, the exact sequence (\SGARefLink{sga7.sga7i.ei.n5-2-1}{5.2.1}) splits canonically, giving


}{E7I-I-019}\begin{equation}\tag{\SGARefTarget{sga7.sga7i.ei.n5-2-3}{5.2.3}}
[x]:\Gal(\bar\eta/\eta)
\longrightarrow\Aut\bigl(\pi_1(X_{\bar\eta},\bar x)\bigr)
\longrightarrow\Aut\bigl(\pi_1^{(p)}(X_{\bar\eta},\bar x)\bigr).
\end{equation}

\SGApage{I}{19}
\underline{Theorem} \SGARefTarget{sga7.sga7i.ei.n5-3}{5.3}. \underline{The image of $P$ under (\SGARefLink{sga7.sga7i.ei.n5-2-2}{5.2.2})
(or under (\SGARefLink{sga7.sga7i.ei.n5-2-3}{5.2.3}), which amounts to the same thing) is finite.}

\SGARefTarget{sga7.sga7i.ei.n5-3-1}{5.3.1}. \underline{Reduction to the case in which $X_\eta$ is
geometrically normal and integral.}

Let $X'_\eta$ be the normalization of $X_\eta$, let
$(X'_{i\eta})_i$ be the family of connected components of $X'_\eta$,
put
\[
X''_\eta=X'_\eta\times_{X_\eta}X'_\eta,
\]
and let $(X''_{j\eta})_j$ be the family of connected components of
$X''_\eta$. After a finite extension of $k(\eta)$, we may suppose that
the $X'_{i\eta}$ are geometrically normal and integral and have rational
points $x_i$, and that the $X''_{j\eta}$ are geometrically connected and
have rational points $y_j$.

For $k=1$ or $2$, and whenever $p_k(y_j)\in X'_{i\eta}$, choose a class
of paths $\ell_{kij}$ from $\operatorname{pr}_k(\bar y_j)$ to $\bar x_i$.
It is an element of a torsor under
$\pi_1(X'_{i\bar\eta},\bar x_i)$. Let $\bar\ell_{kij}$ be the
corresponding prime-to-$p$ path (the same path modulo
$\Ker(\pi\longrightarrow\pi^{(p)})$). We may choose
$\bar\ell_{kij}$ to be invariant under $P$: the set of prime-to-$p$
paths is an inverse limit of finite sets of cardinality prime to $p$ on
which the pro-$p$ group $P$ acts.

\SGAvar{Consequently, if a finite-index subgroup $P'$ of $P$ acts trivially on
the groups $\pi_1^{(p)}(X'_{i\bar\eta},\bar x_i)$, the Van Kampen
formulas---equivalently, the fact that the normalization map
$X'_{i\bar\eta}\longrightarrow X_{i\bar\eta}$ is effective descent for
\'{e}tale coverings---show that $P'$ acts trivially on
$\pi_1^{(p)}(X_{\bar\eta},\bar x)$.

}{Consequently, if a finite-index subgroup $P'$ of $P$ acts trivially on
the groups $\pi_1^{(p)}(X'_{i\bar\eta},\bar x_i)$, the Van Kampen
formulas---equivalently, the fact that the normalization map
$X'_{\bar\eta}\longrightarrow X_{\bar\eta}$ is effective descent for
\'{e}tale coverings---show that $P'$ acts trivially on
$\pi_1^{(p)}(X_{\bar\eta},\bar x)$.

}{E7I-I-020}\SGARefTarget{sga7.sga7i.ei.n5-3-2}{5.3.2}. If $X_{\bar\eta}$ is normal and $U$ is an open subset, then
$\pi_1(U)$ maps onto $\pi_1(X_{\bar\eta})$; we may therefore suppose
that $X_{\bar\eta}$ is smooth and affine. By cutting with
\underline{generic} hyperplane sections (which replaces $S$ by its
strict localization at the generic point of the special fibre of an
$\mathbb A^n_S$) and applying Bertini, we reduce to the case in which
$X_\eta$ is \underline{smooth, geometrically connected, and of
dimension one}.

After a finite extension of $k(\eta)$, we may suppose that $X_\eta$ is
the complement in a smooth projective curve $\bar X_\eta$ of a finite
set of rational points $x_i$, whose number we are free to increase. By
the semistable reduction theorem (\SGARefLink{sga7.sga7i.ei.n5-1-1}{5.1.1}), we reduce to the case in which
there is a stable curve $X=(C;x_1,\ldots,x_n)$ over $S$ such that
\[
X_\eta=C_\eta-\{x_1,\ldots,x_n\}.
\]
We then apply the following result.


\SGApage{I}{20}
\underline{Theorem} \SGARefTarget{sga7.sga7i.ei.n5-4}{5.4}. \underline{Let $S$ be a strictly local scheme,
let $f:\bar X\longrightarrow S$ be a proper flat morphism whose special
fibre is a curve with ordinary double points as its only non-smooth
points, and let $Y$ be a subscheme of $\bar X$ which is the union of a
finite number of disjoint sections contained in the smooth locus. Let
$x$ be a rational point of $X=\bar X-Y$. Let $U$ be the open subset of
$S$ over which $X$ is smooth, and let $\xi$ be a geometric point of $U$.
Then the image of $\pi_1(U,\xi)$ in
$\Aut(\pi_1^{(p)}(X_\xi,x\xi))$ is abelian and prime to $p$.}

Observe that $\pi_1(U,\xi)$ can act on
$\pi_1^{(p)}(X_\xi,x\xi)$ only because the groups $\pi_1^{(p)}$
satisfy the specialization theorem.

We reduce first to the case in which $S$ is complete, and then to the
universal case in which $S$ is a formal moduli scheme for the special
fibre equipped with its sections. Then $S$ is the spectrum of a formal
power-series ring over a Cohen ring:
\[
S=\Spec\bigl(W[[t_1,\ldots,t_n]]\bigr),
\]
and $U$ is the complement of a normal-crossings divisor with equation
$t_1\cdots t_m=0$. In this universal case, Abhyankar's lemma shows that
$\pi_1(U,\xi)$ is abelian and prime to $p$, which proves the theorem.

\underline{Remark} \SGARefTarget{sga7.sga7i.ei.n5-5}{5.5}. One can prove a result analogous to \SGARefLink{sga7.sga7i.ei.n5-4}{5.4} in
arbitrary relative dimension by taking a generic plane section of
dimension one and applying Bertini.



6. \underline{Appendix}: \underline{arithmetic proof of the stable reduction theorem}

(by P. Deligne)

\SGAvar{Let $S$ be a trait and let $A_\eta$ be an abelian variety over $k(\eta)$,
of dimension $d\neq 0$. For a local morphism of traits
$S'\longrightarrow S$, denote by $A_{\eta'}$ the abelian variety over
$k(\eta')$ obtained from $A_\eta$ by extension of scalars. Let
$\mathcal A_{S'}$ be the N\'eron model of $A_{\eta'}$ over $S'$, let
$\mathcal A_{S's'}$ be its special fibre, and let
$\mathcal A^\circ_{S's'}$ be the identity component of
$\mathcal A_{S's'}$. The stable reduction theorem is as follows.


\underline{Theorem \SGARefTarget{sga7.sga7i.ei.n6-1}{6.1}. For a suitable $S'$, $A_{\eta'}$ has stable
reduction, i.e. $\mathcal A^\circ_{S's'}$ is an extension of an abelian
variety by a torus.}

}{Let $S$ be a trait and let $A_\eta$ be an abelian variety over $k(\eta)$,
of dimension $d\neq 0$. For a local morphism of traits
$S'\longrightarrow S$, denote by $A_{\eta'}$ the abelian variety over
$k(\eta')$ obtained from $A_\eta$ by extension of scalars. Let
$\mathcal A_{S'}$ be the N\'eron model of $A_{\eta'}$ over $S'$, let
$\mathcal A_{S',s'}$ be its special fibre, and let
$\mathcal A^\circ_{S',s'}$ be the identity component of
$\mathcal A_{S',s'}$. The stable reduction theorem is as follows.


\underline{Theorem \SGARefTarget{sga7.sga7i.ei.n6-1}{6.1}. For a suitable $S'$, $A_{\eta'}$ has stable
reduction, i.e. $\mathcal A^\circ_{S',s'}$ is an extension of an abelian
variety by a torus.}

}{E7I-I-021}\SGApage{I}{\(\approx\)21}
It follows fairly readily from \SGARefLink{sga7.sga7i.ei.n6-1}{6.1} that one may choose $S'$ such that
$k(\eta')$ is a finite separable extension of $k(\eta)$. Suppose first
that the residue field of $S$ is of finite type over the prime field (the
same argument would work if $k(s)$ were radicial over such a field).

Let $\ell$ be a prime different from the residual characteristic exponent
$p$, let $T_\ell(A)$ be the Tate module, let
$V_\ell(A)=T_\ell(A)\otimes\mathbf Q_\ell$, and let
\[
 \rho:\Gal(\bar\eta/\eta)\longrightarrow
 \mathrm{GL}(T_\ell(A))
\]
be the natural representation. After a preliminary extension of scalars,
we are reduced to the case in which the inertia group $I$ acts
unipotently (\SGARefLink{sga7.sga7i.ei.n1-2}{1.2}), and hence through its largest pro-$\ell$ quotient
$\mathbf Z_\ell(1)$ (\SGARefLink{sga7.sga7i.ei.n0-3}{0.3}). Let $T$ be the image in
$\mathrm{GL}(T_\ell(A))$ of a generator of $\mathbf Z_\ell(1)$, and let
$r$ be the largest integer $\geq 0$ such that $(T-1)^r\neq 0$. We have:

\SGAvar{\underline{Lemma \SGARefTarget{sga7.sga7i.ei.n6-2}{6.2}. The map}
\[
 (\sigma,v)\in I\times V_\ell(A)
 \longmapsto (\rho(\sigma)-1)^r(v)
\]
\underline{induces the morphism and the isomorphism in the commutative
diagram}
\begin{equation}\tag{\SGARefTarget{sga7.sga7i.ei.n6-2-1}{6.2.1}}
\begin{tikzcd}[column sep=huge, row sep=huge]
V_\ell(A)_I\otimes\mathbf Q_\ell(r)
  \arrow[r,"M"] \arrow[d]
\& V_\ell(A)^I \\
V_\ell(A)/\Ker(T-1)^r\otimes\mathbf Q_\ell(r)
  \arrow[r,"\sim"]
\& \operatorname{Im}(T-1)^r \arrow[u]
\end{tikzcd}\quad .
\end{equation}
$V_\ell(A)_I$ \underline{and} $V_\ell(A)^I$ \underline{are
$\Gal(\bar s/s)$-modules, and $M$ is Galois-equivariant.}

}{\underline{Lemma \SGARefTarget{sga7.sga7i.ei.n6-2}{6.2}. The map}
\[
 (\sigma,v)\in I\times V_\ell(A)
 \longmapsto (\rho(\sigma)-1)^r(v)
\]
\underline{induces the morphism and the isomorphism in the commutative
diagram}
\begin{equation}\tag{\SGARefTarget{sga7.sga7i.ei.n6-2-1}{6.2.1}}
\begin{tikzcd}[column sep=huge, row sep=huge]
V_\ell(A)_I\otimes\mathbf Q_\ell(r)
  \arrow[r,"M"] \arrow[d]
\& V_\ell(A)^I \\
\bigl(V_\ell(A)/\Ker(T-1)^r\bigr)\otimes\mathbf Q_\ell(r)
  \arrow[r,"\sim"]
\& \operatorname{Im}(T-1)^r \arrow[u]
\end{tikzcd}\quad .
\end{equation}
$V_\ell(A)_I$ \underline{and} $V_\ell(A)^I$ \underline{are
$\Gal(\bar s/s)$-modules, and $M$ is Galois-equivariant.}

}{E7I-I-022}\SGARefTarget{sga7.sga7i.ei.n6-3}{6.3}. An $\ell$-adic representation
\[
 \tau:\Gal(\bar s/s)\longrightarrow\mathrm{GL}(V)
\]
is said to be of weight $n$ ($n\in\mathbf Z$) if the following condition
is satisfied:

(*) For a suitable model $X$ of $k(s)$ (that is, an irreducible variety
of finite type over the prime field such that $k(s)=k(X)$), $\tau$
factors through $\pi_1(X,\bar s)$ and, for every closed point $x$ of
$X$, the eigenvalues of $\tau(F_x)$ are algebraic numbers all of whose
complex conjugates have absolute value $|k(x)|^{n/2}$ (here $F_x$
denotes the geometric Frobenius $\phi_x^{-1}$).


\SGApage{I}{\(\approx\)22}
If $V$ (respectively $W$) has weight $n$ (respectively $m$) and
$n\neq m$, then $\Hom_{\Gal}(V,W)=0$.

\SGARefTarget{sga7.sga7i.ei.n6-4}{6.4}. It follows from the universal property of the N\'eron model that
$V_\ell(A)^I=V_\ell(\mathcal A^\circ_{S,s})$. If $a$ is the dimension
of the largest abelian quotient of $\mathcal A^\circ_{S,s}$ and $m$ is
the dimension of its multiplicative part, then, by Weil,
$V_\ell(A)^I$ is an extension of a representation of dimension $2a$ and
weight $-1$ by a representation of dimension $m$ and weight $-2$.

\SGAvar{\SGARefTarget{sga7.sga7i.ei.n6-5}{6.5}. Let $A^*$ be the abelian variety dual to $A$. Recall that
$V_\ell(A^*)$ and $V_\ell(A)$ (and hence $V_\ell(A^*)^I$ and
$V_\ell(A)_I$ are in duality with values in $\mathbf Q_\ell(1)$. If
$a^*$ and $m^*$ are defined for $A^*$ as $a$ and $m$ are for $A$, it
follows from \SGARefLink{sga7.sga7i.ei.n6-4}{6.4} that $V_\ell(A)_I$ is an extension of a representation
of dimension $m^*$ and weight $0$ by a representation of dimension
$2a^*$ and weight $-1$.

}{\SGARefTarget{sga7.sga7i.ei.n6-5}{6.5}. Let $A^*$ be the abelian variety dual to $A$. Recall that
$V_\ell(A^*)$ and $V_\ell(A)$ (and hence $V_\ell(A^*)^I$ and
$V_\ell(A)_I$) are in duality with values in $\mathbf Q_\ell(1)$. If
$a^*$ and $m^*$ are defined for $A^*$ as $a$ and $m$ are for $A$, it
follows from \SGARefLink{sga7.sga7i.ei.n6-4}{6.4} that $V_\ell(A)_I$ is an extension of a representation
of dimension $m^*$ and weight $0$ by a representation of dimension
$2a^*$ and weight $-1$.

}{E7I-I-023}\SGARefTarget{sga7.sga7i.ei.n6-6}{6.6}. The Galois module $\mathbf Q_\ell(r)$ has weight $-2r$. Since
$\operatorname{Im}(T-1)^r\neq 0$, it follows from \SGARefLink{sga7.sga7i.ei.n6-4}{6.4}, \SGARefLink{sga7.sga7i.ei.n6-5}{6.5}, and the isomorphism in \SGARefLink{sga7.sga7i.ei.n6-2}{6.2}
that $r\leq 1$. If $r=0$, then
\[
\left\{
\begin{array}{l}
a+m\leq d,\\[1ex]
2a+m=2d,
\end{array}
\right.
\]
so $m=0$, $a=d$, and $A_\eta$ has good reduction. Suppose that $r=1$.
It then follows from \SGARefLink{sga7.sga7i.ei.n6-2}{6.2} that $\operatorname{Im}(T-1)$ has weight $-2$ and that
$V_\ell(A)/V_\ell(A)^I$ has weight $0$; these spaces have the same
dimension $m_1\leq m$.

We then have
\[
 \dim V_\ell(A)^I=2d-m_1=2a+m,
\]
and
\[
 2\dim\mathcal A_{S,s}\geq 2(a+m)\geq
 2a+m+m_1=2d=2\dim\mathcal A_{S,s},
\]
so $A_\eta$ has stable reduction
($\dim\mathcal A_{S,s}=a+m$). At the same time, we have obtained
$m_1=m$; that is,
\[
 \operatorname{Im}(T-1)\subset V_\ell(A)^I\simeq
 V_\ell(\mathcal A^\circ_{S,s})
\]
is the $V_\ell$ of the multiplicative part of
$\mathcal A^\circ_{S,s}$.

\SGARefTarget{sga7.sga7i.ei.n6-7}{6.7}. To treat the general case, we shall use (\SGARefLink{sga7.sga7i.ei.n0-5}{0.5}) (cf. \SGARefLink{sga7.sga7i.ei.n1-3}{1.3}). We reduce
to the case in which:


\SGApage{I}{\(\approx\)23}
(a) the conditions of (\SGARefLink{sga7.sga7i.ei.n0-5-2}{0.5.2}) or (\SGARefLink{sga7.sga7i.ei.n0-5-3}{0.5.3}) are satisfied;

(b) $\Gal(\bar\eta/\eta)$ acts trivially on $A_\ell$;

(c) the action of $I$ on $T_\ell(A)$ is unipotent.

Under these hypotheses, we shall prove that $A_\eta$ has stable
reduction. Hypothesis (a) allows us to apply (\SGARefLink{sga7.sga7i.ei.n0-5-1}{0.5.1}) and write
$S=\varprojlim S_i$ as in \SGARefLink{sga7.sga7i.ei.n1-3}{1.3}, whose notation we retain. Let $s_i$ be
the closed point of $S_i$.

\SGAvar{For sufficiently large $i$, the identity component $\mathcal A^\circ$
of the N\'eron model $\mathcal A_S$ descends to a smooth group scheme
$\mathcal A_i$ with connected fibres over $S_i$, and
$(\mathcal A_i)_\ell$ is constant over $S_i-D$. The group
$\pi_1(S_i-D_i,\bar\eta)$ acts on $T_\ell(A)$ and acts trivially on
$T_\ell(A)/\ell T_\ell(A)\simeq A_\ell$. As in \SGARefLink{sga7.sga7i.ei.n1-3}{1.3}, one sees that
$P_i$ acts trivially. We then have, by (\SGARefLink{sga7.sga7i.ei.n1-3-1}{1.3.1}),
\[
 T_\ell(A)_I=T_\ell(A)_{I_i},\qquad
 T_\ell(\mathcal A_{i,s_i})=T_\ell(\mathcal A^\circ_s)
 =T_\ell(A)^I=T_\ell(A^{I_i}),
\]
and (\SGARefLink{sga7.sga7i.ei.n6-2-1}{6.2.1}) is a diagram of $\Gal(\bar s_i/s_i)$-modules to which the
arguments of \SGARefLink{sga7.sga7i.ei.n6-4}{6.4}--\SGARefLink{sga7.sga7i.ei.n6-6}{6.6} may be applied.


}{For sufficiently large $i$, the identity component $\mathcal A^\circ$
of the N\'eron model $\mathcal A_S$ descends to a smooth group scheme
$\mathcal A_i$ with connected fibres over $S_i$, and
$(\mathcal A_i)_\ell$ is constant over $S_i-D_i$. The group
$\pi_1(S_i-D_i,\bar\eta)$ acts on $T_\ell(A)$ and acts trivially on
$T_\ell(A)/\ell T_\ell(A)\simeq A_\ell$. As in \SGARefLink{sga7.sga7i.ei.n1-3}{1.3}, one sees that
$P_i$ acts trivially. We then have, by (\SGARefLink{sga7.sga7i.ei.n1-3-1}{1.3.1}),
\[
 T_\ell(A)_I=T_\ell(A)_{I_i},\qquad
 T_\ell(\mathcal A_{i,s_i})=T_\ell(\mathcal A^\circ_s)
 =T_\ell(A)^I=T_\ell(A)^{I_i},
\]
and (\SGARefLink{sga7.sga7i.ei.n6-2-1}{6.2.1}) is a diagram of $\Gal(\bar s_i/s_i)$-modules to which the
arguments of \SGARefLink{sga7.sga7i.ei.n6-4}{6.4}--\SGARefLink{sga7.sga7i.ei.n6-6}{6.6} may be applied.


}{E7I-I-024}\begin{center}
\underline{BIBLIOGRAPHY}
\end{center}

\SGAvar{\begin{enumerate}[label={[\arabic*]},leftmargin=3em]
\item M. Artin, ``Algebraic approximation of structures over complete
local rings,'' \emph{Publ. Math. I.H.E.S.} \underline{36} (1969),
pp. 23--58.

\item M. Artin and G. Winters, ``Degenerate fibres and reduction of
curves.''

\item P. Deligne and D. Mumford, ``The irreductibility of the space of
curves of a given genus,'' \emph{Publ. Math. I.H.E.S.} \underline{36}
(1969), pp. 75--110.

\item J.-P. Serre, \emph{Corps locaux}, Publ. Math. Univ.
Nancago--Hermann, 1962.

\item J.-P. Serre and J. Tate, ``Good reduction of abelian varieties,''
\emph{Ann. of Math.} \underline{88}, no. 3 (1968), pp. 492--517.
\end{enumerate}


}{\begin{enumerate}[label={[\arabic*]},leftmargin=3em]
\item M. Artin, ``Algebraic approximation of structures over complete
local rings,'' \emph{Publ. Math. I.H.E.S.} \underline{36} (1969),
pp. 23--58.

\item M. Artin and G. Winters, ``Degenerate fibres and reduction of
curves.''

\item P. Deligne and D. Mumford, ``The irreducibility of the space of
curves of a given genus,'' \emph{Publ. Math. I.H.E.S.} \underline{36}
(1969), pp. 75--110.

\item J.-P. Serre, \emph{Corps locaux}, Publ. Math. Univ.
Nancago--Hermann, 1962.

\item J.-P. Serre and J. Tate, ``Good reduction of abelian varieties,''
\emph{Ann. of Math.} \underline{88}, no. 3 (1968), pp. 492--517.
\end{enumerate}


}{E7I-I-025}\markboth{Exposé I}{Exposé I}

\ifnum\SGAform=1
\clearpage\section*{Critical apparatus}\begingroup\small\raggedright\sloppy
Readings are given in the exact TeX notation of the two English witnesses. Decimal suffixes identify the individual differences in each marked passage.

\par\medskip\noindent\hypertarget{app-E7I-I-001.1}{}\textbf{E7I-I-001.1}\quad printed folio 2.\par
\noindent Printed reading: {\ttfamily\footnotesize Hironoka\allowbreak{}}\par
\noindent Corrected reading: {\ttfamily\footnotesize Hironaka\allowbreak{}}\par
\noindent Source: {\ttfamily\footnotesize 7I:0;\allowbreak{} fixes/FIXLOG\_7I\_CURR\allowbreak{}ENT.jsonl\#607;\allowbreak{} fixes/MISPRINTS\_7I.j\allowbreak{}son\#0;\allowbreak{} port/decisions/decis\allowbreak{}ions\_en7I\_r2\_001.jso\allowbreak{}n\#C606;\allowbreak{} fixes/FIXLOG\_7I\_CURR\allowbreak{}ENT.jsonl\#C606;\allowbreak{} 7I:0\allowbreak{}}\par

\par\medskip\noindent\hypertarget{app-E7I-I-002.1}{}\textbf{E7I-I-002.1}\quad printed folio 4.\par
\noindent Printed reading: {\ttfamily\footnotesize u\allowbreak{}}\par
\noindent Corrected reading: {\ttfamily\footnotesize \textbackslash{}sigma\allowbreak{}}\par
\noindent Source: {\ttfamily\footnotesize ED-EN-FORMULA-G001;\allowbreak{} fr/sga7i/\{diplomatic\allowbreak{},current\}/expose\_I\_b\allowbreak{}ody.tex;\allowbreak{} folio\allowbreak{} 4\allowbreak{}}\par

\par\medskip\noindent\hypertarget{app-E7I-I-003.1}{}\textbf{E7I-I-003.1}\quad printed folio 4.\par
\noindent Printed reading: \emph{(empty)}\par
\noindent Corrected reading: {\ttfamily\footnotesize \textbackslash{}bar\allowbreak{}}\par
\noindent Source: {\ttfamily\footnotesize fr/sga7i/\{diplomatic\allowbreak{},current\}/expose\_I\_b\allowbreak{}ody.tex;\allowbreak{} folio\allowbreak{} 4;\allowbreak{} formula\allowbreak{} 124;\allowbreak{} port/decisions/decis\allowbreak{}ions\_en7I\_r2\_001.jso\allowbreak{}n\#C607;\allowbreak{} fixes/FIXLOG\_7I\_CURR\allowbreak{}ENT.jsonl\#C607;\allowbreak{} 7I:3\allowbreak{}}\par

\par\medskip\noindent\hypertarget{app-E7I-I-004.1}{}\textbf{E7I-I-004.1}\quad printed folio 6.\par
\noindent Printed reading: {\ttfamily\footnotesize 1\allowbreak{}}\par
\noindent Corrected reading: {\ttfamily\footnotesize 0\allowbreak{}}\par
\noindent Source: {\ttfamily\footnotesize ADOPT-F7I-I-007.1;\allowbreak{} ADOPT-F7I-I-007.2;\allowbreak{} fr/sga7i/diplomatic/\allowbreak{}expose\_I\_body.tex;\allowbreak{} fr/sga7i/current/exp\allowbreak{}ose\_I\_body.tex\allowbreak{}}\par

\par\medskip\noindent\hypertarget{app-E7I-I-004.2}{}\textbf{E7I-I-004.2}\quad printed folio 6.\par
\noindent Printed reading: {\ttfamily\footnotesize 1\allowbreak{}}\par
\noindent Corrected reading: {\ttfamily\footnotesize 0\allowbreak{}}\par
\noindent Source: {\ttfamily\footnotesize ADOPT-F7I-I-007.1;\allowbreak{} ADOPT-F7I-I-007.2;\allowbreak{} fr/sga7i/diplomatic/\allowbreak{}expose\_I\_body.tex;\allowbreak{} fr/sga7i/current/exp\allowbreak{}ose\_I\_body.tex\allowbreak{}}\par

\par\medskip\noindent\hypertarget{app-E7I-I-005.1}{}\textbf{E7I-I-005.1}\quad printed folio 7.\par
\noindent Printed reading: \emph{(empty)}\par
\noindent Corrected reading: {\ttfamily\footnotesize \_o\allowbreak{}}\par
\noindent Source: {\ttfamily\footnotesize fr/sga7i/\{diplomatic\allowbreak{},current\}/expose\_I\_b\allowbreak{}ody.tex;\allowbreak{} folio\allowbreak{} 7\allowbreak{}}\par

\par\medskip\noindent\hypertarget{app-E7I-I-005.2}{}\textbf{E7I-I-005.2}\quad printed folio 7.\par
\noindent Printed reading: \emph{(empty)}\par
\noindent Corrected reading: {\ttfamily\footnotesize \_o\allowbreak{}}\par
\noindent Source: {\ttfamily\footnotesize fr/sga7i/\{diplomatic\allowbreak{},current\}/expose\_I\_b\allowbreak{}ody.tex;\allowbreak{} folio\allowbreak{} 7;\allowbreak{} formula\allowbreak{} 255\allowbreak{}}\par

\par\medskip\noindent\hypertarget{app-E7I-I-005.3}{}\textbf{E7I-I-005.3}\quad printed folio 7.\par
\noindent Printed reading: \emph{(empty)}\par
\noindent Corrected reading: {\ttfamily\footnotesize \_o\allowbreak{}}\par
\noindent Source: {\ttfamily\footnotesize fr/sga7i/\{diplomatic\allowbreak{},current\}/expose\_I\_b\allowbreak{}ody.tex;\allowbreak{} folio\allowbreak{} 7;\allowbreak{} formula\allowbreak{} 257\allowbreak{}}\par

\par\medskip\noindent\hypertarget{app-E7I-I-005.4}{}\textbf{E7I-I-005.4}\quad printed folio 7.\par
\noindent Printed reading: {\ttfamily\footnotesize Gal\allowbreak{}}\par
\noindent Corrected reading: {\ttfamily\footnotesize GL\allowbreak{}}\par
\noindent Source: {\ttfamily\footnotesize fr/sga7i/\{diplomatic\allowbreak{},current\}/expose\_I\_b\allowbreak{}ody.tex;\allowbreak{} folio\allowbreak{} 7;\allowbreak{} formula\allowbreak{} 262\allowbreak{}}\par

\par\medskip\noindent\hypertarget{app-E7I-I-006.1}{}\textbf{E7I-I-006.1}\quad printed folio 8.\par
\noindent Printed reading: {\ttfamily\footnotesize \_\allowbreak{}}\par
\noindent Corrected reading: {\ttfamily\footnotesize \textasciicircum{}\allowbreak{}}\par
\noindent Source: {\ttfamily\footnotesize fr/sga7i/\{diplomatic\allowbreak{},current\}/expose\_I\_b\allowbreak{}ody.tex;\allowbreak{} folio\allowbreak{} 8;\allowbreak{} formula\allowbreak{} 311\allowbreak{}}\par

\par\medskip\noindent\hypertarget{app-E7I-I-007.1}{}\textbf{E7I-I-007.1}\quad printed folio 8.\par
\noindent Printed reading: \emph{(empty)}\par
\noindent Corrected reading: {\ttfamily\footnotesize \textbackslash{}bar\allowbreak{}}\par
\noindent Source: {\ttfamily\footnotesize fr/sga7i/\{diplomatic\allowbreak{},current\}/expose\_I\_b\allowbreak{}ody.tex;\allowbreak{} folio\allowbreak{} 8;\allowbreak{} formula\allowbreak{} 315\allowbreak{}}\par

\par\medskip\noindent\hypertarget{app-E7I-I-008.1}{}\textbf{E7I-I-008.1}\quad printed folio 9.\par
\noindent Printed reading: \emph{(empty)}\par
\noindent Corrected reading: {\ttfamily\footnotesize +1\allowbreak{}}\par
\noindent Source: {\ttfamily\footnotesize fr/sga7i/\{diplomatic\allowbreak{},current\}/expose\_I\_b\allowbreak{}ody.tex;\allowbreak{} folio\allowbreak{} 9;\allowbreak{} formula\allowbreak{} 348\allowbreak{}}\par

\par\medskip\noindent\hypertarget{app-E7I-I-009.1}{}\textbf{E7I-I-009.1}\quad printed folio 10.\par
\noindent Printed reading: \emph{(empty)}\par
\noindent Corrected reading: {\ttfamily\footnotesize 0.\allowbreak{}}\par
\noindent Source: {\ttfamily\footnotesize ADOPT-F7I-I-013.1;\allowbreak{} fr/sga7i/diplomatic/\allowbreak{}expose\_I\_body.tex;\allowbreak{} fr/sga7i/current/exp\allowbreak{}ose\_I\_body.tex\allowbreak{}}\par

\par\medskip\noindent\hypertarget{app-E7I-I-009.2}{}\textbf{E7I-I-009.2}\quad printed folio 10.\par
\noindent Printed reading: {\ttfamily\footnotesize X\allowbreak{}}\par
\noindent Corrected reading: {\ttfamily\footnotesize x\allowbreak{}}\par
\noindent Source: {\ttfamily\footnotesize fr/sga7i/\{diplomatic\allowbreak{},current\}/expose\_I\_b\allowbreak{}ody.tex;\allowbreak{} folio\allowbreak{} 10;\allowbreak{} formula\allowbreak{} 384\allowbreak{}}\par

\par\medskip\noindent\hypertarget{app-E7I-I-010.1}{}\textbf{E7I-I-010.1}\quad printed folio 11.\par
\noindent Printed reading: \emph{(empty)}\par
\noindent Corrected reading: {\ttfamily\footnotesize \textbackslash{}bar\allowbreak{}}\par
\noindent Source: {\ttfamily\footnotesize fr/sga7i/\{diplomatic\allowbreak{},current\}/expose\_I\_b\allowbreak{}ody.tex;\allowbreak{} folio\allowbreak{} 11;\allowbreak{} formula\allowbreak{} 445\allowbreak{}}\par

\par\medskip\noindent\hypertarget{app-E7I-I-010.2}{}\textbf{E7I-I-010.2}\quad printed folio 11.\par
\noindent Printed reading: {\ttfamily\footnotesize xd\allowbreak{}}\par
\noindent Corrected reading: {\ttfamily\footnotesize nd\allowbreak{}}\par
\noindent Source: {\ttfamily\footnotesize fr/sga7i/\{diplomatic\allowbreak{},current\}/expose\_I\_b\allowbreak{}ody.tex;\allowbreak{} folio\allowbreak{} 11;\allowbreak{} formula\allowbreak{} 451\allowbreak{}}\par

\par\medskip\noindent\hypertarget{app-E7I-I-011.1}{}\textbf{E7I-I-011.1}\quad printed folio 12.\par
\noindent Printed reading: \emph{(empty)}\par
\noindent Corrected reading: {\ttfamily\footnotesize +1\allowbreak{}}\par
\noindent Source: {\ttfamily\footnotesize fr/sga7i/\{diplomatic\allowbreak{},current\}/expose\_I\_b\allowbreak{}ody.tex;\allowbreak{} folio\allowbreak{} 12;\allowbreak{} formula\allowbreak{} 482\allowbreak{}}\par

\par\medskip\noindent\hypertarget{app-E7I-I-012.1}{}\textbf{E7I-I-012.1}\quad printed folio 13.\par
\noindent Printed reading: \emph{(empty)}\par
\noindent Corrected reading: {\ttfamily\footnotesize \textbackslash{}bar\allowbreak{}}\par
\noindent Source: {\ttfamily\footnotesize fr/sga7i/\{diplomatic\allowbreak{},current\}/expose\_I\_b\allowbreak{}ody.tex;\allowbreak{} folio\allowbreak{} 13;\allowbreak{} formula\allowbreak{} 527\allowbreak{}}\par

\par\medskip\noindent\hypertarget{app-E7I-I-012.2}{}\textbf{E7I-I-012.2}\quad printed folio 13.\par
\noindent Printed reading: \emph{(empty)}\par
\noindent Corrected reading: {\ttfamily\footnotesize \_\textbackslash{}ell\allowbreak{}}\par
\noindent Source: {\ttfamily\footnotesize fr/sga7i/\{diplomatic\allowbreak{},current\}/expose\_I\_b\allowbreak{}ody.tex;\allowbreak{} folio\allowbreak{} 13;\allowbreak{} formula\allowbreak{} 530\allowbreak{}}\par

\par\medskip\noindent\hypertarget{app-E7I-I-013.1}{}\textbf{E7I-I-013.1}\quad printed folio 15.\par
\noindent Printed reading: \emph{(empty)}\par
\noindent Corrected reading: {\ttfamily\footnotesize \textbackslash{}bar\allowbreak{}}\par
\noindent Source: {\ttfamily\footnotesize fr/sga7i/\{diplomatic\allowbreak{},current\}/expose\_I\_b\allowbreak{}ody.tex;\allowbreak{} folio\allowbreak{} 15;\allowbreak{} formula\allowbreak{} 629;\allowbreak{} port/decisions/decis\allowbreak{}ions\_en7I\_r2\_001.jso\allowbreak{}n\#C614;\allowbreak{} fixes/FIXLOG\_7I\_CURR\allowbreak{}ENT.jsonl\#C614;\allowbreak{} 7I:25\allowbreak{}}\par

\par\medskip\noindent\hypertarget{app-E7I-I-013.2}{}\textbf{E7I-I-013.2}\quad printed folio 15.\par
\noindent Printed reading: {\ttfamily\footnotesize a\allowbreak{}}\par
\noindent Corrected reading: {\ttfamily\footnotesize 1\allowbreak{}}\par
\noindent Source: {\ttfamily\footnotesize port/decisions/decis\allowbreak{}ions\_en7I\_r2\_001.jso\allowbreak{}n\#C615;\allowbreak{} fixes/FIXLOG\_7I\_CURR\allowbreak{}ENT.jsonl\#C615;\allowbreak{} 7I:26\allowbreak{}}\par

\par\medskip\noindent\hypertarget{app-E7I-I-014.1}{}\textbf{E7I-I-014.1}\quad printed folio 15.\par
\noindent Printed reading: \emph{(empty)}\par
\noindent Corrected reading: {\ttfamily\footnotesize '\allowbreak{}}\par
\noindent Source: {\ttfamily\footnotesize fr/sga7i/\{diplomatic\allowbreak{},current\}/expose\_I\_b\allowbreak{}ody.tex;\allowbreak{} folio\allowbreak{} 15;\allowbreak{} formula\allowbreak{} 664;\allowbreak{} port/decisions/decis\allowbreak{}ions\_en7I\_r2\_001.jso\allowbreak{}n\#C616;\allowbreak{} fixes/FIXLOG\_7I\_CURR\allowbreak{}ENT.jsonl\#C616;\allowbreak{} 7I:27\allowbreak{}}\par

\par\medskip\noindent\hypertarget{app-E7I-I-015.1}{}\textbf{E7I-I-015.1}\quad printed folio 16.\par
\noindent Printed reading: \emph{(empty)}\par
\noindent Corrected reading: {\ttfamily\footnotesize smooth\allowbreak{}}\par
\noindent Source: {\ttfamily\footnotesize WITNESS-F7I-I-022.1;\allowbreak{} fr/sga7i/diplomatic/\allowbreak{}expose\_I\_body.tex;\allowbreak{} fr/sga7i/current/exp\allowbreak{}ose\_I\_body.tex\allowbreak{}}\par

\par\medskip\noindent\hypertarget{app-E7I-I-016.1}{}\textbf{E7I-I-016.1}\quad printed folio 16.\par
\noindent Printed reading: {\ttfamily\footnotesize g\allowbreak{}}\par
\noindent Corrected reading: {\ttfamily\footnotesize y\allowbreak{}}\par
\noindent Source: {\ttfamily\footnotesize ED-EN-FORMULA-G003;\allowbreak{} fr/sga7i/\{diplomatic\allowbreak{},current\}/expose\_I\_b\allowbreak{}ody.tex;\allowbreak{} folio\allowbreak{} 16\allowbreak{}}\par

\par\medskip\noindent\hypertarget{app-E7I-I-016.2}{}\textbf{E7I-I-016.2}\quad printed folio 16.\par
\noindent Printed reading: {\ttfamily\footnotesize g\allowbreak{}}\par
\noindent Corrected reading: {\ttfamily\footnotesize y\allowbreak{}}\par
\noindent Source: {\ttfamily\footnotesize ED-EN-FORMULA-G003;\allowbreak{} fr/sga7i/\{diplomatic\allowbreak{},current\}/expose\_I\_b\allowbreak{}ody.tex;\allowbreak{} folio\allowbreak{} 16\allowbreak{}}\par

\par\medskip\noindent\hypertarget{app-E7I-I-017.1}{}\textbf{E7I-I-017.1}\quad printed folio 17.\par
\noindent Printed reading: {\ttfamily\footnotesize \_1\allowbreak{}}\par
\noindent Corrected reading: \emph{(empty)}\par
\noindent Source: {\ttfamily\footnotesize fr/sga7i/\{diplomatic\allowbreak{},current\}/expose\_I\_b\allowbreak{}ody.tex;\allowbreak{} folio\allowbreak{} 17;\allowbreak{} formula\allowbreak{} 714\allowbreak{}}\par

\par\medskip\noindent\hypertarget{app-E7I-I-018.1}{}\textbf{E7I-I-018.1}\quad printed folio 18.\par
\noindent Printed reading: \emph{(empty)}\par
\noindent Corrected reading: {\ttfamily\footnotesize '\allowbreak{}}\par
\noindent Source: {\ttfamily\footnotesize fr/sga7i/\{diplomatic\allowbreak{},current\}/expose\_I\_b\allowbreak{}ody.tex;\allowbreak{} folio\allowbreak{} 18;\allowbreak{} formula\allowbreak{} 756\allowbreak{}}\par

\par\medskip\noindent\hypertarget{app-E7I-I-018.2}{}\textbf{E7I-I-018.2}\quad printed folio 18.\par
\noindent Printed reading: {\ttfamily\footnotesize two\allowbreak{}}\par
\noindent Corrected reading: {\ttfamily\footnotesize \$-2\$\allowbreak{}}\par
\noindent Source: {\ttfamily\footnotesize fr/sga7i/\{diplomatic\allowbreak{},current\}/expose\_I\_b\allowbreak{}ody.tex;\allowbreak{} folio\allowbreak{} 18\allowbreak{}}\par

\par\medskip\noindent\hypertarget{app-E7I-I-019.1}{}\textbf{E7I-I-019.1}\quad printed folio 18.\par
\noindent Printed reading: {\ttfamily\footnotesize ,\allowbreak{} dash\allowbreak{}}\par
\noindent Corrected reading: \emph{(empty)}\par
\noindent Source: {\ttfamily\footnotesize fr/sga7i/\{diplomatic\allowbreak{},current\}/expose\_I\_b\allowbreak{}ody.tex;\allowbreak{} folio\allowbreak{} 18;\allowbreak{} formula\allowbreak{} 773\allowbreak{}}\par

\par\medskip\noindent\hypertarget{app-E7I-I-019.2}{}\textbf{E7I-I-019.2}\quad printed folio 18.\par
\noindent Printed reading: {\ttfamily\footnotesize 2\allowbreak{}}\par
\noindent Corrected reading: {\ttfamily\footnotesize 1\allowbreak{}}\par
\noindent Source: {\ttfamily\footnotesize 7I:35;\allowbreak{} fixes/MISPRINTS\_7I.j\allowbreak{}son\#35\allowbreak{}}\par

\par\medskip\noindent\hypertarget{app-E7I-I-020.1}{}\textbf{E7I-I-020.1}\quad printed folio 19.\par
\noindent Printed reading: {\ttfamily\footnotesize i\allowbreak{}}\par
\noindent Corrected reading: \emph{(empty)}\par
\noindent Source: {\ttfamily\footnotesize fr/sga7i/\{diplomatic\allowbreak{},current\}/expose\_I\_b\allowbreak{}ody.tex;\allowbreak{} folio\allowbreak{} 19;\allowbreak{} formula\allowbreak{} 820\allowbreak{}}\par

\par\medskip\noindent\hypertarget{app-E7I-I-020.2}{}\textbf{E7I-I-020.2}\quad printed folio 19.\par
\noindent Printed reading: {\ttfamily\footnotesize i\allowbreak{}}\par
\noindent Corrected reading: \emph{(empty)}\par
\noindent Source: {\ttfamily\footnotesize fr/sga7i/\{diplomatic\allowbreak{},current\}/expose\_I\_b\allowbreak{}ody.tex;\allowbreak{} folio\allowbreak{} 19;\allowbreak{} formula\allowbreak{} 820\allowbreak{}}\par

\par\medskip\noindent\hypertarget{app-E7I-I-021.1}{}\textbf{E7I-I-021.1}\quad printed folio 20.\par
\noindent Printed reading: \emph{(empty)}\par
\noindent Corrected reading: {\ttfamily\footnotesize ,\allowbreak{}}\par
\noindent Source: {\ttfamily\footnotesize fr/sga7i/\{diplomatic\allowbreak{},current\}/expose\_I\_b\allowbreak{}ody.tex;\allowbreak{} folio\allowbreak{} 20;\allowbreak{} formula\allowbreak{} 875\allowbreak{}}\par

\par\medskip\noindent\hypertarget{app-E7I-I-021.2}{}\textbf{E7I-I-021.2}\quad printed folio 20.\par
\noindent Printed reading: \emph{(empty)}\par
\noindent Corrected reading: {\ttfamily\footnotesize ,\allowbreak{}}\par
\noindent Source: {\ttfamily\footnotesize fr/sga7i/\{diplomatic\allowbreak{},current\}/expose\_I\_b\allowbreak{}ody.tex;\allowbreak{} folio\allowbreak{} 20;\allowbreak{} formula\allowbreak{} 876\allowbreak{}}\par

\par\medskip\noindent\hypertarget{app-E7I-I-021.3}{}\textbf{E7I-I-021.3}\quad printed folio 20.\par
\noindent Printed reading: \emph{(empty)}\par
\noindent Corrected reading: {\ttfamily\footnotesize ,\allowbreak{}}\par
\noindent Source: {\ttfamily\footnotesize fr/sga7i/\{diplomatic\allowbreak{},current\}/expose\_I\_b\allowbreak{}ody.tex;\allowbreak{} folio\allowbreak{} 20;\allowbreak{} formula\allowbreak{} 877\allowbreak{}}\par

\par\medskip\noindent\hypertarget{app-E7I-I-021.4}{}\textbf{E7I-I-021.4}\quad printed folio 21.\par
\noindent Printed reading: \emph{(empty)}\par
\noindent Corrected reading: {\ttfamily\footnotesize ,\allowbreak{}}\par
\noindent Source: {\ttfamily\footnotesize fr/sga7i/\{diplomatic\allowbreak{},current\}/expose\_I\_b\allowbreak{}ody.tex;\allowbreak{} folio\allowbreak{} 21;\allowbreak{} formula\allowbreak{} 880\allowbreak{}}\par

\par\medskip\noindent\hypertarget{app-E7I-I-022.1}{}\textbf{E7I-I-022.1}\quad printed folio 21.\par
\noindent Printed reading: \emph{(empty)}\par
\noindent Corrected reading: {\ttfamily\footnotesize \textbackslash{}bigl(\allowbreak{}}\par
\noindent Source: {\ttfamily\footnotesize fr/sga7i/\{diplomatic\allowbreak{},current\}/expose\_I\_b\allowbreak{}ody.tex;\allowbreak{} folio\allowbreak{} 21;\allowbreak{} formula\allowbreak{} 902\allowbreak{}}\par

\par\medskip\noindent\hypertarget{app-E7I-I-022.2}{}\textbf{E7I-I-022.2}\quad printed folio 21.\par
\noindent Printed reading: \emph{(empty)}\par
\noindent Corrected reading: {\ttfamily\footnotesize \textbackslash{}bigr)\allowbreak{}}\par
\noindent Source: {\ttfamily\footnotesize fr/sga7i/\{diplomatic\allowbreak{},current\}/expose\_I\_b\allowbreak{}ody.tex;\allowbreak{} folio\allowbreak{} 21;\allowbreak{} formula\allowbreak{} 902\allowbreak{}}\par

\par\medskip\noindent\hypertarget{app-E7I-I-023.1}{}\textbf{E7I-I-023.1}\quad printed folio 22.\par
\noindent Printed reading: \emph{(empty)}\par
\noindent Corrected reading: {\ttfamily\footnotesize )\allowbreak{}}\par
\noindent Source: {\ttfamily\footnotesize port/decisions/decis\allowbreak{}ions\_en7I\_r2\_001.jso\allowbreak{}n\#C618;\allowbreak{} fixes/FIXLOG\_7I\_CURR\allowbreak{}ENT.jsonl\#C618;\allowbreak{} 7I:40\allowbreak{}}\par

\par\medskip\noindent\hypertarget{app-E7I-I-024.1}{}\textbf{E7I-I-024.1}\quad printed folio 23.\par
\noindent Printed reading: \emph{(empty)}\par
\noindent Corrected reading: {\ttfamily\footnotesize \_i\allowbreak{}}\par
\noindent Source: {\ttfamily\footnotesize fr/sga7i/\{diplomatic\allowbreak{},current\}/expose\_I\_b\allowbreak{}ody.tex;\allowbreak{} folio\allowbreak{} 23;\allowbreak{} formula\allowbreak{} 999\allowbreak{}}\par

\par\medskip\noindent\hypertarget{app-E7I-I-024.2}{}\textbf{E7I-I-024.2}\quad printed folio 23.\par
\noindent Printed reading: \emph{(empty)}\par
\noindent Corrected reading: {\ttfamily\footnotesize )\allowbreak{}}\par
\noindent Source: {\ttfamily\footnotesize fr/sga7i/\{diplomatic\allowbreak{},current\}/expose\_I\_b\allowbreak{}ody.tex;\allowbreak{} folio\allowbreak{} 23;\allowbreak{} formula\allowbreak{} 1004\allowbreak{}}\par

\par\medskip\noindent\hypertarget{app-E7I-I-024.3}{}\textbf{E7I-I-024.3}\quad printed folio 23.\par
\noindent Printed reading: {\ttfamily\footnotesize )\allowbreak{}}\par
\noindent Corrected reading: \emph{(empty)}\par
\noindent Source: {\ttfamily\footnotesize fr/sga7i/\{diplomatic\allowbreak{},current\}/expose\_I\_b\allowbreak{}ody.tex;\allowbreak{} folio\allowbreak{} 23;\allowbreak{} formula\allowbreak{} 1004\allowbreak{}}\par

\par\medskip\noindent\hypertarget{app-E7I-I-025.1}{}\textbf{E7I-I-025.1}\quad printed folio 24.\par
\noindent Printed reading: {\ttfamily\footnotesize irreductibility\allowbreak{}}\par
\noindent Corrected reading: {\ttfamily\footnotesize irreducibility\allowbreak{}}\par
\noindent Source: {\ttfamily\footnotesize 7I:43;\allowbreak{} fixes/MISPRINTS\_7I.j\allowbreak{}son\#43\allowbreak{}}\par

\par\endgroup\normalsize\fi
\clearpage\SGAtoc{Exposé II}
\selectlanguage{english}
\SGApage{I}{25}

\begin{center}
\underline{\SGARefAliasTarget{sga7.sga7i.eii.nroot}Expos\'e II}

\vspace{1em}
\underline{\underline{Finiteness properties of the fundamental group}}

\vspace{1em}
\underline{by Mich\`ele Raynaud}
\end{center}

Let $k$ be a separably closed field of characteristic $p\geq 0$, let
$X$ be a connected scheme of finite type over $k$, and let
$\Pi_1^{(p)}(X)$ be the largest ``prime-to-$p$'' quotient of $\Pi_1(X)$.
In this expos\'e we prove that, when $X$ is a surface,
$\Pi_1^{(p)}(X)$ is topologically finitely presented, and that the same
holds for arbitrary $X$ if resolution of singularities is assumed.

For a surface, the proof method, due to J.-P. Murre, consists in reducing
to a fibration over $P^1$ having properties (i), (ii), and (iii) of the
following proposition.

\underline{Proposition \SGARefTarget{sga7.sga7i.eii.n2-1}{2.1}. Let $k$ be an algebraically closed field and
let $S$ be a projective, irreducible, smooth surface over $k$. Then there
exists a blow-up}
\[
 g:S'\longrightarrow S,
\]
\underline{where $S'$ is a regular surface equipped with a fibration}
\[
 f:S'\longrightarrow P^1,
\]
\underline{such that $f$ has a section $\sigma$ and the following
conditions hold:}

(i) \underline{$f$ is smooth at the points of $\sigma(P^1)$.}

(ii) \underline{The fibres of $f$ are geometrically integral curves
having only ordinary singularities.}

(iii) \underline{There exists a nonempty open subset $U$ of $P^1$ such
that the fibres of $f$ over the points of $U$ are smooth.}

The proposition follows from the existence of a projective embedding
$S\longrightarrow P^r$ for which a Lefschetz pencil exists (XVII 2.5).
Indeed, let $D=\{H_t\}_{t\in P^1}$ be such a pencil of hyperplanes. Recall
that it has the following properties (loc. cit. \SGARefLink{sga7.sga7i.eii.n2-2}{2.2}):


\SGApage{I}{\(\approx\)26}
(a) The axis $\Delta$ of $D$ meets $S$ transversely; the closed
subscheme $\Delta\cap S$ is therefore reduced and consists of finitely
many closed points $x_1,\ldots,x_n$ of $S$.

(b) There exists a nonempty open subset $U$ of $P^1$ such that, for
every $t\in U$, $H_t$ meets $S$ transversely.

(c) For every $t\in P^1-U$, $H_t$ meets $S$ transversely except at one
point, which is an ordinary quadratic singularity.

Through every point $x$ of $S$ not belonging to $\Delta$ there passes a
unique hyperplane $H_t$. The map assigning $t\in P^1$ to $x$ is a
rational map $a:S\dashrightarrow P^1$ with domain of definition
$S-(\Delta\cap S)$.

Consider the blow-up $g:S'\longrightarrow S$ of the closed subscheme
$S\cap\Delta$. The scheme $S'$ is regular, and the fibre $S'_{x_i}$ over
$x_i$ is the projective bundle defined by the tangent space of $S$ at
$x_i$, hence is isomorphic to $P^1$ (EGA IV 19.4.2, 19.4.4). Moreover,
the rational map $f=ag$ is everywhere defined (XVIII 3.1). Finally, to
each $x_i$ one can associate the section of $f$ that sends
$t\in P^1$ to the tangent vector at $x_i$ of $H_t\cap S$.

Let us verify conditions (i), (ii), and (iii). For each $t\in P^1$, the
curves $S'_t$ and $S_t=H_t\cap S$ are isomorphic. Indeed, the morphism
$S'_t\longrightarrow S_t$ is an isomorphism except possibly at the
points $x_i$, where $S_t$ is regular; moreover, since $S'_t$ is a
complete intersection in $S'$, it is integral, and is therefore
isomorphic to $S_t$.

\SGAvar{Conditions (i) and (ii) are precisely conditions (a) and (b). Finally,
$f$ is smooth at the points of $\sigma$, since $f$ is flat and the
fibres $S'_t$ are geometrically regular at $\sigma(t)$.


}{Conditions (ii) and (iii) are precisely conditions (b) and (c). Finally,
$f$ is smooth at the points of $\sigma$, since $f$ is flat and the
fibres $S'_t$ are geometrically regular at $\sigma(t)$.


}{E7I-II-001}\SGApage{I}{27}
\underline{Corollary \SGARefTarget{sga7.sga7i.eii.n2-2}{2.2}. With the notation of \SGARefLink{sga7.sga7i.eii.n2-1}{2.1}, let $Y$ be a
normal-crossings divisor in $S$, and let $Y'=Y\times_S S'$. Then $f$,
$g$, and $U$ may be chosen so that $Y'\times_{P^1}U$ is a divisor with
normal crossings relative to $U$.}

\SGAvar{It is known that, if a Lefschetz pencil exists for an embedding
$S\longrightarrow P^r$, the Lefschetz pencils form an open subset of the
variety of lines in $P^r$ (XVII 3.2.1). The same is true of the set of
Lefschetz pencils whose axis does not meet $Y$ and for which there is a
nonempty open subset $U_1$ of $P^1$ such that $Y\cap H_t$ is smooth for
$t\in U_1$. It is then enough to construct $f$ and $g$ from a pencil
satisfying these two conditions and replace by $U\cap U_1$.

}{It is known that, if a Lefschetz pencil exists for an embedding
$S\longrightarrow P^r$, the Lefschetz pencils form an open subset of the
variety of lines in $P^r$ (XVII 3.2.1). The same is true of the set of
Lefschetz pencils whose axis does not meet $Y$ and for which there is a
nonempty open subset $U_1$ of $P^1$ such that $Y\cap H_t$ is smooth for
$t\in U_1$. It is then enough to construct $f$ and $g$ from a pencil
satisfying these two conditions and replace $U$ by $U\cap U_1$.

}{E7I-II-002}\SGARefTarget{sga7.sga7i.eii.n2-3-0}{2.3.0}. Let $G$ be a profinite group. Recall that $G$ is said to be
topologically finitely generated if there is an integer $n$ such that
$G$ is isomorphic to a quotient of the free pro-group $F(n)$ on $n$
generators. Put $K=\Ker(F(n)\longrightarrow G)$. The group $G$ is said
to be topologically finitely presented if, in addition, there are
finitely many elements $k_1,\ldots,k_r$ of $K$ such that the smallest
closed normal subgroup of $K$ containing $k_1,\ldots,k_r$ is $K$.

\SGAvar{\underline{Theorem \SGARefTarget{sga7.sga7i.eii.n2-3-1}{2.3.1}. Let $k$ be a separably closed field of
characteristic $p\geq 0$, and let $X$ be a connected scheme of finite
type over $k$. Suppose that schemes of finite type and dimension at most
$\dim X$ over an algebraic closure of $k$ admit strong resolution of
singularities (SGA 5 I 3.1.5). This condition is satisfied when
$\dim X=2$, by [1], or when $k=0$, by [2]. Then the
fundamental group $\Pi_1^{(p)}(X)$ is topologically finitely presented.}

}{\underline{Theorem \SGARefTarget{sga7.sga7i.eii.n2-3-1}{2.3.1}. Let $k$ be a separably closed field of
characteristic $p\geq 0$, and let $X$ be a connected scheme of finite
type over $k$. Suppose that schemes of finite type and dimension at most
$\dim X$ over an algebraic closure of $k$ admit strong resolution of
singularities (SGA 5 I 3.1.5). This condition is satisfied when
$\dim X=2$, by [1], or when $\operatorname{char}k=0$, by [2]. Then the
fundamental group $\Pi_1^{(p)}(X)$ is topologically finitely presented.}

}{E7I-II-003}By SGA 1 IX 4.10, we may suppose that $k$ is algebraically closed.

1) \underline{Reduction to the case in which $X$ is a regular surface.}

\SGAvar{If $f:X'\longrightarrow X$ is a morphism of finite type, put
$X''=X'\times_X S'$, and let $(X'_a)_{a\in A}$ and
$(X''_b)_{b\in B}$ be the families of connected components of $X'$ and
$X''$, respectively. When $f$ is an effective descent morphism for the
category of finite \'etale coverings, the fundamental group $\Pi_1(X)$
can be computed explicitly


}{If $f:X'\longrightarrow X$ is a morphism of finite type, put
$X''=X'\times_X X'$, and let $(X'_a)_{a\in A}$ and
$(X''_b)_{b\in B}$ be the families of connected components of $X'$ and
$X''$, respectively. When $f$ is an effective descent morphism for the
category of finite \'etale coverings, the fundamental group $\Pi_1(X)$
can be computed explicitly


}{E7I-II-004}\SGApage{I}{\(\approx\)28}
in terms of the fundamental groups $\Pi_1(X'_a)$ and $\Pi_1(X''_b)$.
It follows (SGA 1 IX 5.2, \SGARefLink{sga7.sga7i.eix.n5-3}{5.3}) that if the $\Pi_1(X'_a)$ are
topologically finitely generated, then so is $\Pi_1(X)$; and if the
$\Pi_1(X'_a)$ are topologically finitely presented while the
$\Pi_1(X''_b)$ are topologically finitely generated, then $\Pi_1(X)$ is
topologically finitely presented.

Since $X$ admits resolution of singularities, there is a regular scheme
$X'$ and a proper birational morphism $f:X'\longrightarrow X$. Since
$f$ is an effective descent morphism for finite \'etale coverings, it
suffices to prove the theorem under the additional assumption that $X$
is regular. One first deduces that $\Pi_1(X)$ is topologically finitely
generated; knowing that $\Pi_1(X')$ is topologically finitely presented
and that the $\Pi_1(X''_b)$ are topologically finitely generated then
shows that $\Pi_1(X)$ is topologically finitely presented.

Now let $(U_i)_{i\in I}$ be a finite affine open covering of $X$, and
let $f:\coprod U_i\longrightarrow X$ be the canonical morphism. Since
$f$ is an effective descent morphism for finite \'etale coverings, it is
enough to prove the theorem for the $U_i$. We are thus reduced to the
case in which $X$ is affine and smooth.

Assume henceforth that $X$ is affine and smooth over $k$, with
$\dim X>2$. The Lefschetz theorem for quasi-projective schemes
([3] V 7.4) applies: if $Y$ is a sufficiently general hyperplane section
of $X$, then
\[
 \Pi_1(Y)\simeq\Pi_1(X).
\]
It therefore suffices to prove the theorem for $Y$, reducing us to the
case $\dim X=2$.

2) \underline{The case in which $X$ is a smooth surface over $k$.}

By the strong resolution theorem for surfaces, there exist a smooth,
integral, projective surface $S$ and an open immersion
$i:X\longrightarrow S$ such that $Y=S-X$ is a normal-crossings divisor.
We apply Corollary \SGARefLink{sga7.sga7i.eii.n2-2}{2.2} and retain its notation.


\SGApage{I}{\(\approx\)29}
\SGAvar{We show that it is enough to prove that the fundamental group of
$X'=X\times_S S'$ is topologically finitely presented. Since $X$ and
$X'$ are regular and
\[
 \operatorname{codim}\!\left(X-\bigcup_{1\leq i\leq n}\{x_i\},X\right)
 \geq 2,
\]
every \'etale covering of $X'$ induces on
$X-\bigcup_{1\leq i\leq n}\{x_i\}$ an \'etale covering which extends
uniquely to an \'etale covering of $X$ (SGA 2 1.11). This gives an
isomorphism
\[
 \Pi_1(X')\simeq\Pi_1(X).
\]

Consider the cartesian diagram
\[
\begin{tikzcd}[column sep=large,row sep=large]
X' \arrow[d,"h"'] \& X'_U \arrow[l] \arrow[d,"h_U"] \\
P^1 \& U \arrow[l]
\end{tikzcd}
\]
where $h=f|_{X'}$. Let $L$ be the set of primes different from $p$, and
let $\bar\eta$ be a geometric point over the generic point $\eta$ of
$P^1$. We verify the hypotheses of SGA 1 XIII 4.8. The morphism $h$ has
a section $\sigma$, because $f$ does. To prove condition (a) of
SGA 1 XIII 4.7, note first that, since the fibres of $h$ are
geometrically integral, $h$ is $0$-acyclic and locally $0$-acyclic
(SGA 4 XV 4.1, 1.16). It remains to show that, for every locally constant
sheaf of sets $F$ on $X'$, $(F,h)$ is cohomologically proper in
dimensions $\leq 0$. If $\xi$ is a closed point of $P^1$ and $Z$ denotes
the integral closure of $\Spec\mathcal O_{P^1,\xi}$ in $\bar\eta$, then
$X'_Z=X'\times_{P^1}Z$ satisfies condition $(S_2)$ at its closed points
and is regular at every other point, hence is normal. Consequently, an
\'etale covering of $X'$ whose restriction to $X'_{\bar\eta}$ is
connected is necessarily connected. This proves the assertion, and
hence condition (a).

}{We show that it is enough to prove that the fundamental group of
$X'=X\times_S S'$ is topologically finitely presented. Since $X$ and
$X'$ are regular and
\[
 \operatorname{codim}\!\left(\bigcup_{1\leq i\leq n}\{x_i\},X\right)
 \geq 2,
\]
every \'etale covering of $X'$ induces on
$X-\bigcup_{1\leq i\leq n}\{x_i\}$ an \'etale covering which extends
uniquely to an \'etale covering of $X$ (SGA 2 1.11). This gives an
isomorphism
\[
 \Pi_1(X')\simeq\Pi_1(X).
\]

Consider the cartesian diagram
\[
\begin{tikzcd}[column sep=large,row sep=large]
X' \arrow[d,"h"'] \& X'_U \arrow[l] \arrow[d,"h_U"] \\
P^1 \& U \arrow[l]
\end{tikzcd}
\]
where $h=f|_{X'}$. Let $L$ be the set of primes different from $p$, and
let $\bar\eta$ be a geometric point over the generic point $\eta$ of
$P^1$. We verify the hypotheses of SGA 1 XIII 4.8. The morphism $h$ has
a section $\sigma$, because $f$ does. To prove condition (a) of
SGA 1 XIII 4.7, note first that, since the fibres of $h$ are
geometrically integral, $h$ is $0$-acyclic and locally $0$-acyclic
(SGA 4 XV 4.1, 1.16). It remains to show that, for every locally constant
sheaf of sets $F$ on $X'$, $(F,h)$ is cohomologically proper in
dimensions $\leq 0$. If $\xi$ is a closed point of $P^1$ and $Z$ denotes
the integral closure of $\Spec\mathcal O_{P^1,\xi}$ in $\bar\eta$, then
$X'_Z=X'\times_{P^1}Z$ satisfies condition $(S_2)$ at its closed points
and is regular at every other point, hence is normal. Consequently, an
\'etale covering of $X'_Z$ whose restriction to $X'_{\bar\eta}$ is
connected is necessarily connected. This proves the assertion, and
hence condition (a).

}{E7I-II-005}Since $X'_U$ is the complement in $S'_U$ of a divisor with normal
crossings relative to $U$, $h_U$ is locally $1$-aspherical for $L$
(SGA 4 XV 2.1); and, for every finite constant sheaf of $L$-groups $F$
on $X'_U$, $(F,h_U)$ is cohomologically proper in dimensions $\leq 1$
(SGA 1 XIII 2.4), which gives condition (b). Since $P^1$ is normal and
$T=P^1-U$ has codimension $1$, we have


\SGApage{I}{\(\approx\)30}
\SGAvar{\[
 \operatorname{prof\,\acute et}_T(S)\geq 2,
\]
which gives condition (c). The condition
\[
 \operatorname{prof\,hop}_x(X')\geq 3
\]
holds at every closed point of $X'$ by the purity theorem, which gives
condition (e). Finally, for every $t\in P^1$, $h$ is smooth at
$\sigma(t)$, and hence is locally $1$-aspherical for $L$ at $\sigma(t)$,
which gives condition (d).

By SGA 1 XIII 4.7, if $a$ is a geometric point of
$X'_{\bar\eta}$, there is an exact sequence
\[
 1\longrightarrow\Pi_1^{(p)}(X'_{\bar\eta},a)
 \longrightarrow\Pi'_1(X'_U,a)
 \longrightarrow\Pi_1(U,a)\longrightarrow 1,
\]
and the choice of the section $\sigma$ defines a splitting of this exact
sequence. Let $K$ be the image of $\Pi_1(U,a)$ in
$\Aut(\Pi_1^{(p)}(X'_{\bar\eta},a))$. The group of coinvariants under
$K$,
\[
 \Pi_1^{(p)}(X'_{\bar\eta},a)_K,
\]
is the quotient of $\Pi_1^{(p)}(X'_{\bar\eta},a)$ by the closed normal
subgroup generated by the elements $x^{-1}qx$, where
$x\in\Pi_1^{(p)}(X'_{\bar\eta},a)$ and $q\in K$. By loc. cit. (4.7.4),
there is an isomorphism
\[
 \Pi_1^{(p)}(X',a)\simeq
 \Pi_1^{(p)}(X'_{\bar\eta},a).
\]

}{\[
 \operatorname{prof\,\acute et}_T(P^1)\geq 2,
\]
which gives condition (c). The condition
\[
 \operatorname{prof\,hop}_x(X')\geq 3
\]
holds at every closed point of $X'$ by the purity theorem, which gives
condition (e). Finally, for every $t\in P^1$, $h$ is smooth at
$\sigma(t)$, and hence is locally $1$-aspherical for $L$ at $\sigma(t)$,
which gives condition (d).

By SGA 1 XIII 4.7, if $a$ is a geometric point of
$X'_{\bar\eta}$, there is an exact sequence
\[
 1\longrightarrow\Pi_1^{(p)}(X'_{\bar\eta},a)
 \longrightarrow\Pi'_1(X'_U,a)
 \longrightarrow\Pi_1(U,a)\longrightarrow 1,
\]
and the choice of the section $\sigma$ defines a splitting of this exact
sequence. Let $K$ be the image of $\Pi_1(U,a)$ in
$\Aut(\Pi_1^{(p)}(X'_{\bar\eta},a))$. The group of coinvariants under
$K$,
\[
 \Pi_1^{(p)}(X'_{\bar\eta},a)_K,
\]
is the quotient of $\Pi_1^{(p)}(X'_{\bar\eta},a)$ by the closed normal
subgroup generated by the elements $x^{-1}q(x)$, where
$x\in\Pi_1^{(p)}(X'_{\bar\eta},a)$ and $q\in K$. By loc. cit. (4.7.4),
there is an isomorphism
\[
 \Pi_1^{(p)}(X',a)\simeq
 \Pi_1^{(p)}(X'_{\bar\eta},a)_K.
\]

}{E7I-II-006}Write $P^1-U=\coprod_{1\leq i\leq n}\{t_i\}$. For each $i$, let
$\bar T_i$ be the strict localization of $P^1$ at a geometric point
$\bar t_i$ over $t_i$, and let $s_i$ be the generic point of $\bar T_i$.
An \'etale covering of $U$ whose restriction to every $s_i$ is trivial
extends to an \'etale covering of $P^1$, and is therefore trivial. For
each $i$, let $\overline{k(s_i)}$ be an algebraic closure of $k(s_i)$.
The preceding assertion says that $\Pi_1(U,a)$ is the smallest closed
normal subgroup containing the images of the Galois groups
$\Gal(\overline{k(s_i)}/k(s_i))$ in $\Pi_1(U,a)$. By I \SGARefLink{sga7.sga7i.ei.n5-3}{5.3}, the image
$K_i$ of $\Gal(\overline{k(s_i)}/k(s_i))$ in
$\Aut(\Pi_1^{(p)}(X'_{\bar\eta},a))$ is finite. Since $K$ is the
smallest closed normal subgroup containing all the $K_i$, and since
$\Pi_1^{(p)}(X'_{\bar\eta},a)$ is topologically finitely presented
(SGA 1 XIII 3.2), this proves that
$\Pi_1^{(p)}(X'_{\bar\eta},a)_K$, and hence also the isomorphic group
$\Pi_1^{(p)}(X)$, is topologically finitely presented.


\SGApage{I}{31}
\underline{Remark \SGARefTarget{sga7.sga7i.eii.n2-3-2}{2.3.2}.} Resolution of singularities is not needed to
prove Theorem \SGARefLink{sga7.sga7i.eii.n2-3-1}{2.3.1} when $X$ is the open complement of a
normal-crossings divisor in a projective scheme smooth over $k$.

\begin{center}\rule{1.5cm}{0.4pt}\end{center}

\begin{center}\underline{BIBLIOGRAPHY}\end{center}

\SGAvar{\begin{enumerate}[label={[\arabic*]},leftmargin=3em]
\item S. Abhyankar, ``Local uniformization of algebraic surfaces over
ground field of characteristic $p=0$,'' \emph{Amer. J. of Math.},
vol. 63, 1956.

\item H. Hironaka, ``Resolution of singularities of an algebraic variety
over a field of characteristic zero,'' \emph{Ann. of Math.}, vol. 79,
1964.

\item M. Raynaud, \emph{Th\'eor\`eme de Lefschetz en cohomologie des
faisceaux coh\'erents et en cohomologie \'etale}, thesis, forthcoming.
\end{enumerate}


}{\begin{enumerate}[label={[\arabic*]},leftmargin=3em]
\item S. Abhyankar, ``Local uniformization on algebraic surfaces over
ground fields of characteristic $p\neq0$,'' \emph{Ann. of Math.},
vol. 63, no. 3, 1956, pp. 491--526.

\item H. Hironaka, ``Resolution of singularities of an algebraic variety
over a field of characteristic zero,'' \emph{Ann. of Math.}, vol. 79,
1964.

\item M. Raynaud, \emph{Th\'eor\`eme de Lefschetz en cohomologie des
faisceaux coh\'erents et en cohomologie \'etale}, thesis, forthcoming.
\end{enumerate}


}{E7I-II-007}\markboth{Exposé II}{Exposé II}

\ifnum\SGAform=1
\clearpage\section*{Critical apparatus}\begingroup\small\raggedright\sloppy
Readings are given in the exact TeX notation of the two English witnesses. Decimal suffixes identify the individual differences in each marked passage.

\par\medskip\noindent\hypertarget{app-E7I-II-001.1}{}\textbf{E7I-II-001.1}\quad printed folio 26.\par
\noindent Printed reading: {\ttfamily\footnotesize i\allowbreak{}}\par
\noindent Corrected reading: {\ttfamily\footnotesize ii\allowbreak{}}\par
\noindent Source: {\ttfamily\footnotesize ED-EN-NONMATH-I-II26\allowbreak{};\allowbreak{} 7I:44;\allowbreak{} WITNESS-F7I-II-001.3\allowbreak{};\allowbreak{} WITNESS-F7I-II-001.4\allowbreak{};\allowbreak{} fixes/MISPRINTS\_7I.j\allowbreak{}son\#44;\allowbreak{} fr/sga7i/diplomatic/\allowbreak{}expose\_II\_body.tex;\allowbreak{} fr/sga7i/current/exp\allowbreak{}ose\_II\_body.tex;\allowbreak{} F7I-II-001.1,\allowbreak{} F7I-II-001.2,\allowbreak{} F7I-II-001.3,\allowbreak{} F7I-II-001.4\allowbreak{}}\par

\par\medskip\noindent\hypertarget{app-E7I-II-001.2}{}\textbf{E7I-II-001.2}\quad printed folio 26.\par
\noindent Printed reading: {\ttfamily\footnotesize ii\allowbreak{}}\par
\noindent Corrected reading: {\ttfamily\footnotesize iii\allowbreak{}}\par
\noindent Source: {\ttfamily\footnotesize ED-EN-NONMATH-I-II26\allowbreak{};\allowbreak{} 7I:44;\allowbreak{} WITNESS-F7I-II-001.3\allowbreak{};\allowbreak{} WITNESS-F7I-II-001.4\allowbreak{};\allowbreak{} fixes/MISPRINTS\_7I.j\allowbreak{}son\#44;\allowbreak{} fr/sga7i/diplomatic/\allowbreak{}expose\_II\_body.tex;\allowbreak{} fr/sga7i/current/exp\allowbreak{}ose\_II\_body.tex;\allowbreak{} F7I-II-001.1,\allowbreak{} F7I-II-001.2,\allowbreak{} F7I-II-001.3,\allowbreak{} F7I-II-001.4\allowbreak{}}\par

\par\medskip\noindent\hypertarget{app-E7I-II-001.3}{}\textbf{E7I-II-001.3}\quad printed folio 26.\par
\noindent Printed reading: {\ttfamily\footnotesize a\allowbreak{}}\par
\noindent Corrected reading: {\ttfamily\footnotesize b\allowbreak{}}\par
\noindent Source: {\ttfamily\footnotesize ED-EN-NONMATH-I-II26\allowbreak{};\allowbreak{} 7I:44;\allowbreak{} WITNESS-F7I-II-001.3\allowbreak{};\allowbreak{} WITNESS-F7I-II-001.4\allowbreak{};\allowbreak{} fixes/MISPRINTS\_7I.j\allowbreak{}son\#44;\allowbreak{} fr/sga7i/diplomatic/\allowbreak{}expose\_II\_body.tex;\allowbreak{} fr/sga7i/current/exp\allowbreak{}ose\_II\_body.tex;\allowbreak{} F7I-II-001.1,\allowbreak{} F7I-II-001.2,\allowbreak{} F7I-II-001.3,\allowbreak{} F7I-II-001.4\allowbreak{}}\par

\par\medskip\noindent\hypertarget{app-E7I-II-001.4}{}\textbf{E7I-II-001.4}\quad printed folio 26.\par
\noindent Printed reading: {\ttfamily\footnotesize b\allowbreak{}}\par
\noindent Corrected reading: {\ttfamily\footnotesize c\allowbreak{}}\par
\noindent Source: {\ttfamily\footnotesize ED-EN-NONMATH-I-II26\allowbreak{};\allowbreak{} 7I:44;\allowbreak{} WITNESS-F7I-II-001.3\allowbreak{};\allowbreak{} WITNESS-F7I-II-001.4\allowbreak{};\allowbreak{} fixes/MISPRINTS\_7I.j\allowbreak{}son\#44;\allowbreak{} fr/sga7i/diplomatic/\allowbreak{}expose\_II\_body.tex;\allowbreak{} fr/sga7i/current/exp\allowbreak{}ose\_II\_body.tex;\allowbreak{} F7I-II-001.1,\allowbreak{} F7I-II-001.2,\allowbreak{} F7I-II-001.3,\allowbreak{} F7I-II-001.4\allowbreak{}}\par

\par\medskip\noindent\hypertarget{app-E7I-II-002.1}{}\textbf{E7I-II-002.1}\quad printed folio 27.\par
\noindent Printed reading: \emph{(empty)}\par
\noindent Corrected reading: {\ttfamily\footnotesize \$U\$\allowbreak{} }\par
\noindent Source: {\ttfamily\footnotesize ED-EN-FORMULA-G005;\allowbreak{} fr/sga7i/\{diplomatic\allowbreak{},current\}/expose\_II\_\allowbreak{}body.tex;\allowbreak{} folio\allowbreak{} 27\allowbreak{}}\par

\par\medskip\noindent\hypertarget{app-E7I-II-003.1}{}\textbf{E7I-II-003.1}\quad printed folio 27.\par
\noindent Printed reading: \emph{(empty)}\par
\noindent Corrected reading: {\ttfamily\footnotesize \textbackslash{}operatorname\{char\}\allowbreak{}}\par
\noindent Source: {\ttfamily\footnotesize fr/sga7i/\{diplomatic\allowbreak{},current\}/expose\_II\_\allowbreak{}body.tex;\allowbreak{} folio\allowbreak{} 27;\allowbreak{} formula\allowbreak{} 123\allowbreak{}}\par

\par\medskip\noindent\hypertarget{app-E7I-II-004.1}{}\textbf{E7I-II-004.1}\quad printed folio 27.\par
\noindent Printed reading: {\ttfamily\footnotesize S\allowbreak{}}\par
\noindent Corrected reading: {\ttfamily\footnotesize X\allowbreak{}}\par
\noindent Source: {\ttfamily\footnotesize fr/sga7i/\{diplomatic\allowbreak{},current\}/expose\_II\_\allowbreak{}body.tex;\allowbreak{} folio\allowbreak{} 27;\allowbreak{} formula\allowbreak{} 128\allowbreak{}}\par

\par\medskip\noindent\hypertarget{app-E7I-II-005.1}{}\textbf{E7I-II-005.1}\quad printed folio 29.\par
\noindent Printed reading: {\ttfamily\footnotesize X-\allowbreak{}}\par
\noindent Corrected reading: \emph{(empty)}\par
\noindent Source: {\ttfamily\footnotesize fr/sga7i/\{diplomatic\allowbreak{},current\}/expose\_II\_\allowbreak{}body.tex;\allowbreak{} folio\allowbreak{} 29;\allowbreak{} formula\allowbreak{} 173\allowbreak{}}\par

\par\medskip\noindent\hypertarget{app-E7I-II-005.2}{}\textbf{E7I-II-005.2}\quad printed folio 29.\par
\noindent Printed reading: \emph{(empty)}\par
\noindent Corrected reading: {\ttfamily\footnotesize \_Z\allowbreak{}}\par
\noindent Source: {\ttfamily\footnotesize fr/sga7i/\{diplomatic\allowbreak{},current\}/expose\_II\_\allowbreak{}body.tex;\allowbreak{} folio\allowbreak{} 29;\allowbreak{} formula\allowbreak{} 215\allowbreak{}}\par

\par\medskip\noindent\hypertarget{app-E7I-II-006.1}{}\textbf{E7I-II-006.1}\quad printed folio 30.\par
\noindent Printed reading: {\ttfamily\footnotesize S\allowbreak{}}\par
\noindent Corrected reading: {\ttfamily\footnotesize P\textasciicircum{}1\allowbreak{}}\par
\noindent Source: {\ttfamily\footnotesize fr/sga7i/\{diplomatic\allowbreak{},current\}/expose\_II\_\allowbreak{}body.tex;\allowbreak{} folio\allowbreak{} 30;\allowbreak{} formula\allowbreak{} 230\allowbreak{}}\par

\par\medskip\noindent\hypertarget{app-E7I-II-006.2}{}\textbf{E7I-II-006.2}\quad printed folio 30.\par
\noindent Printed reading: {\ttfamily\footnotesize qx\allowbreak{}}\par
\noindent Corrected reading: {\ttfamily\footnotesize q(x)\allowbreak{}}\par
\noindent Source: {\ttfamily\footnotesize fr/sga7i/\{diplomatic\allowbreak{},current\}/expose\_II\_\allowbreak{}body.tex;\allowbreak{} folio\allowbreak{} 30;\allowbreak{} formula\allowbreak{} 248\allowbreak{}}\par

\par\medskip\noindent\hypertarget{app-E7I-II-006.3}{}\textbf{E7I-II-006.3}\quad printed folio 30.\par
\noindent Printed reading: \emph{(empty)}\par
\noindent Corrected reading: {\ttfamily\footnotesize \_K\allowbreak{}}\par
\noindent Source: {\ttfamily\footnotesize fr/sga7i/\{diplomatic\allowbreak{},current\}/expose\_II\_\allowbreak{}body.tex;\allowbreak{} folio\allowbreak{} 30;\allowbreak{} formula\allowbreak{} 251\allowbreak{}}\par

\par\medskip\noindent\hypertarget{app-E7I-II-007.1}{}\textbf{E7I-II-007.1}\quad printed folio 31.\par
\noindent Printed reading: {\ttfamily\footnotesize of\allowbreak{}}\par
\noindent Corrected reading: {\ttfamily\footnotesize on\allowbreak{}}\par
\noindent Source: {\ttfamily\footnotesize 7I:52;\allowbreak{} ADOPT-F7I-II-008.5;\allowbreak{} ADOPT-F7I-II-008.6;\allowbreak{} ADOPT-F7I-II-008.7;\allowbreak{} fixes/MISPRINTS\_7I.j\allowbreak{}son\#52;\allowbreak{} fr/sga7i/diplomatic/\allowbreak{}expose\_II\_body.tex;\allowbreak{} fr/sga7i/current/exp\allowbreak{}ose\_II\_body.tex\allowbreak{}}\par

\par\medskip\noindent\hypertarget{app-E7I-II-007.2}{}\textbf{E7I-II-007.2}\quad printed folio 31.\par
\noindent Printed reading: {\ttfamily\footnotesize field\allowbreak{}}\par
\noindent Corrected reading: {\ttfamily\footnotesize fields\allowbreak{}}\par
\noindent Source: {\ttfamily\footnotesize 7I:52;\allowbreak{} ADOPT-F7I-II-008.5;\allowbreak{} ADOPT-F7I-II-008.6;\allowbreak{} ADOPT-F7I-II-008.7;\allowbreak{} fixes/MISPRINTS\_7I.j\allowbreak{}son\#52;\allowbreak{} fr/sga7i/diplomatic/\allowbreak{}expose\_II\_body.tex;\allowbreak{} fr/sga7i/current/exp\allowbreak{}ose\_II\_body.tex\allowbreak{}}\par

\par\medskip\noindent\hypertarget{app-E7I-II-007.3}{}\textbf{E7I-II-007.3}\quad printed folio 31.\par
\noindent Printed reading: {\ttfamily\footnotesize =\allowbreak{}}\par
\noindent Corrected reading: {\ttfamily\footnotesize \textbackslash{}neq\allowbreak{}}\par
\noindent Source: {\ttfamily\footnotesize 7I:52;\allowbreak{} ADOPT-F7I-II-008.5;\allowbreak{} ADOPT-F7I-II-008.6;\allowbreak{} ADOPT-F7I-II-008.7;\allowbreak{} fixes/MISPRINTS\_7I.j\allowbreak{}son\#52;\allowbreak{} fr/sga7i/diplomatic/\allowbreak{}expose\_II\_body.tex;\allowbreak{} fr/sga7i/current/exp\allowbreak{}ose\_II\_body.tex;\allowbreak{} fr/sga7i/\{diplomatic\allowbreak{},current\}/expose\_II\_\allowbreak{}body.tex;\allowbreak{} folio\allowbreak{} 31;\allowbreak{} formula\allowbreak{} 279\allowbreak{}}\par

\par\medskip\noindent\hypertarget{app-E7I-II-007.4}{}\textbf{E7I-II-007.4}\quad printed folio 31.\par
\noindent Printed reading: {\ttfamily\footnotesize Amer.\allowbreak{} J\allowbreak{}}\par
\noindent Corrected reading: {\ttfamily\footnotesize Ann\allowbreak{}}\par
\noindent Source: {\ttfamily\footnotesize 7I:52;\allowbreak{} ADOPT-F7I-II-008.5;\allowbreak{} ADOPT-F7I-II-008.6;\allowbreak{} ADOPT-F7I-II-008.7;\allowbreak{} fixes/MISPRINTS\_7I.j\allowbreak{}son\#52;\allowbreak{} fr/sga7i/diplomatic/\allowbreak{}expose\_II\_body.tex;\allowbreak{} fr/sga7i/current/exp\allowbreak{}ose\_II\_body.tex\allowbreak{}}\par

\par\medskip\noindent\hypertarget{app-E7I-II-007.5}{}\textbf{E7I-II-007.5}\quad printed folio 31.\par
\noindent Printed reading: \emph{(empty)}\par
\noindent Corrected reading: {\ttfamily\footnotesize no.\allowbreak{} 3,\allowbreak{} }\par
\noindent Source: {\ttfamily\footnotesize 7I:52;\allowbreak{} ADOPT-F7I-II-008.5;\allowbreak{} ADOPT-F7I-II-008.6;\allowbreak{} ADOPT-F7I-II-008.7;\allowbreak{} fixes/MISPRINTS\_7I.j\allowbreak{}son\#52;\allowbreak{} fr/sga7i/diplomatic/\allowbreak{}expose\_II\_body.tex;\allowbreak{} fr/sga7i/current/exp\allowbreak{}ose\_II\_body.tex\allowbreak{}}\par

\par\medskip\noindent\hypertarget{app-E7I-II-007.6}{}\textbf{E7I-II-007.6}\quad printed folio 31.\par
\noindent Printed reading: \emph{(empty)}\par
\noindent Corrected reading: {\ttfamily\footnotesize ,\allowbreak{} pp.\allowbreak{} 491--526\allowbreak{}}\par
\noindent Source: {\ttfamily\footnotesize 7I:52;\allowbreak{} ADOPT-F7I-II-008.5;\allowbreak{} ADOPT-F7I-II-008.6;\allowbreak{} ADOPT-F7I-II-008.7;\allowbreak{} fixes/MISPRINTS\_7I.j\allowbreak{}son\#52;\allowbreak{} fr/sga7i/diplomatic/\allowbreak{}expose\_II\_body.tex;\allowbreak{} fr/sga7i/current/exp\allowbreak{}ose\_II\_body.tex\allowbreak{}}\par

\par\endgroup\normalsize\fi
\clearpage\SGAtoc{Exposé VI}
\selectlanguage{english}

\SGApage{I}{32}


S G A  7


\underline{Exposé VI}


\underline{\underline{FORMAL DEFORMATION THEORY}}


\underline{by  D. S. RIM}


1. \underline{Formal existence theorem}

We recall the notion of cofibered category (S.G.A.1.IV). Let  $\underline{C}$  be a fixed category. A cofibered category  $\underline{A}$  over  $\underline{C}$  is, by definition, a category  $\underline{A}$  together with a functor  $P : \underline{A} \to \underline{C}$  such that

\underline{Ax.1.}  For any map  $f$  in  $\underline{C}$  and any object  $a$  in  $\underline{A}$  with  $P(a) =$ the source of  $f$  , there exists an  $f$-morphism  $\alpha : a \to b$  which is cocartesian (i.e. for any  $f$-morphism  $\alpha' : a \to b'$  there exists a unique  $T$-morphism  $\tau : b \to b'$  in  $\underline{A}$  such that  $\alpha' = \tau\cdot\alpha$  where  $T =$ the target of  $f$ ).

\SGAvar{\underline{Ax.2.}  A composit of cocartesian morphism in  $\underline{A}$  is again cocartesian.

}{\underline{Ax.2.}  A composite of cocartesian morphisms in  $\underline{A}$  is again cocartesian.

}{F7I-VI-001}The following properties of a cofibered category  $P : \underline{A} \to \underline{C}$  follow immediately from the definitions:

(cancellation)  Let  $a \overset{\alpha}{\underset{\beta}{\rightrightarrows}} a'$  be any two cocartesian maps such that  $P(\alpha) = P(\beta)$  . If there exists a cocartesian map  $b \overset{\gamma}{\to} a$  such that  $\alpha\gamma = \beta\gamma$ , then  $\alpha = \beta$  .

(factorization)  Given cocartesian maps  $a \overset{\alpha}{\to} a'$  ,  $a \overset{\beta}{\to} a''$  , there exists a cocartesian map  $a' \overset{\gamma}{\to} a''$  such that  $\gamma\alpha = \beta$  if and only if there exists  $f : P(a') \to P(a'')$  in  $\underline{C}$  such that  $f\cdot P(\alpha) = P(\beta)$  . Furthermore  $\gamma$  is unique.

Let  $\underline{A}$  be a cofibered category over  $\underline{C}$  . For each object  $S$  in  $\underline{C}$  we denote by  $\underline{A}(S)$  the subcategory of  $\underline{A}$  whose objects are the objects  $a$  in  $\underline{A}$  such that  $P(a) = S$  , and  $\mathrm{Hom}_{\underline{A}(S)}(a,a') = \{\alpha \in \mathrm{Hom}_{\underline{A}}(a,a') \mid P(\alpha) = \mathrm{id}_S\}$  . A cofibered groupoid over  $\underline{C}$  is a cofibered category  $\underline{A}$  over  $\underline{C}$  such that  $\underline{A}(S)$
\SGApage{I}{33}
is a groupoid for every object  $S$  in  $\underline{C}$ . We recall that a groupoid is a category in which every morphism is an isomorphism.  Thus, a cofibered groupoid over $\underline{C}$  is a cofibered category  $P : \underline{A} \to \underline{C}$  such that every map in  $\underline{A}$  is cocartesian. Furthermore, for every map  $\alpha \in F\ell(\underline{A})$  ,  $\alpha$  is an isomorphism in  $\underline{A}$  if and only if  $P(\alpha)$  is an isomorphism.

\SGAvar{Throughout the rest of this section, we fix a complete noetherian local ring  $\Lambda$  with the residue field  $k$  , and we set  $\underline{C}_{\Lambda}$  to be the category of artinian local  $\Lambda$-algebras with the same residue field  $k$  , where the maps are local homomorphisms over  $\Lambda$  . From 1.2 on, a cofibered groupoid  $\underline{A}$  over  $\underline{C}_{\Lambda}$ is always assumed to have the property that  $\underline{A}(k) = \{e\}$  , where  $\{e\}$  denotes the category in which  $\mathrm{Hom}_{\{e\}}(a,\ b)$  consists of one and only one element for any  two objects  $a,\ b$  in  $\{e\}$  . A map  $a' \xrightarrow{\ \alpha\ } a$  in  $\underline{A}$  will be called simply a $\underline{C}_{\Lambda}$ -map if  $p(a') = p(a)$  and  $p(\alpha) = 1_{p(a)}$  . A map  $a' \xrightarrow{\ \alpha\ } a$  on  $\underline{A}$  will be called surjective, by abuse of language, if  $p(\alpha) : p(a') \longrightarrow p(a)$  is surjective in  $\underline{C}_{\Lambda}$  .

}{Throughout the rest of this section, we fix a complete noetherian local ring  $\Lambda$  with the residue field  $k$  , and we set  $\underline{C}_{\Lambda}$  to be the category of artinian local  $\Lambda$-algebras with the same residue field  $k$  , where the maps are local homomorphisms over  $\Lambda$  . From 1.2 on, a cofibered groupoid  $\underline{A}$  over  $\underline{C}_{\Lambda}$ is always assumed to have the property that  $\underline{A}(k) = \{e\}$  , where  $\{e\}$  denotes the category in which  $\mathrm{Hom}_{\{e\}}(a,\ b)$  consists of one and only one element for any  two objects  $a,\ b$  in  $\{e\}$  . A map  $a' \xrightarrow{\ \alpha\ } a$  in  $\underline{A}$  will be called simply a $\underline{C}_{\Lambda}$ -map if  $p(a') = p(a)$  and  $p(\alpha) = 1_{p(a)}$  . A map  $a' \xrightarrow{\ \alpha\ } a$  in  $\underline{A}$  will be called surjective, by abuse of language, if  $p(\alpha) : p(a') \longrightarrow p(a)$  is surjective in  $\underline{C}_{\Lambda}$  .

}{F7I-VI-002}\underline{Ex.1.1.(a)}  Let  $F : \underline{C} \longrightarrow$(categories) be a functor. Define the category $\underline{F}$  as follows:  an object in  $\underline{F}$  is a pair  $(S,a)$  where  $S \in \mathrm{ob}(\underline{C})$  and $a \in \mathrm{ob}F(S)$  , and a map  $(S,a) \to (S',a')$  is a pair  $(f,u')$  where  $f : S \to S'$  is a map in  $\underline{C}$  and  $u' : F(f)a \to a'$  is a map in  $F(S')$  . Composition of two maps $(S,a) \xrightarrow{\ (f,u')\ } (S',a') \xrightarrow{\ (g,u'')\ } (S'',a'')$  is given by $(gf,\ u''\cdot(F(g)u')) : (S,a) \to (S'',a'')$  . With respect to the functor  $P : \underline{F} \to \underline{C}$ given by  $(S,a) \mapsto S$  ,  $\underline{F}$  becomes a cofibered category over  $\underline{C}$  . If  $F(S)$  is a groupoid for every  $S \in \mathrm{ob}(\underline{C})$  , then  $\underline{F}$  is a cofibered groupoid over  $\underline{C}$  . In particular, a functor  $F : \underline{C} \to$ (groups) yields a cofibered group, and a functor $F : \underline{C} \to$ (Ens)  yields a cofibered discrete groupoid. We also note that a natural transformation  $\varphi : F \to G$  , where  $F,G : \underline{C} \to$ (categories) are functors, induces a cocartesian functor  $\varphi : \underline{F} \to \underline{G}$  over  $\underline{C}$  .

\underline{Ex.1.1.(b)}  Let  $X$  ,  $Y$  be schemes over  $\Lambda$  together with a fixed
\SGApage{I}{34}
isomorphism  $X \underset{\Lambda}{\otimes} k \xleftrightarrow{\ \tau\ } Y \underset{\Lambda}{\otimes} k$  . Then the functor  $\mathrm{Isom}_{X,Y} : \underline{C}_{\Lambda} \to$ (groupoids) given by  $\mathrm{Isom}_{X,Y}(S) =$ the set of isomorphisms  $X \underset{\Lambda}{\otimes} S \xleftrightarrow{\ \sigma\ } Y \underset{\Lambda}{\otimes} S$  such that $\sigma \underset{\Lambda}{\otimes} k = \tau$  yields a cofibered groupoid  $P : \underline{\mathrm{Isom}}_{X,Y} \to \underline{C}_{\Lambda}$  .

\underline{Ex.1.1.(c)}  Let  $X_0$  be a fixed scheme over  $k$  , and let  $\underline{M}_{X_0}$  be the category consisting of deformations of  $X_0$  over  $\underline{C}_{\Lambda}$  (see §4). Thus an object in $\underline{M}_{X_0}$  is a pair  $(R,X)$  where  $R$  is in  $\underline{C}_{\Lambda}$  and  $X$  is a deformation of  $X_0$  to $R$  . Then the functor  $P : \underline{M}_{X_0} \to \underline{C}_{\Lambda}$  given by  $(R,X) \mathrel{\mapstochar\rightsquigarrow} R$  defines a cofibered groupoid over  $\underline{C}_{\Lambda}$  .

\underline{Ex.1.1.(d)}  For any groupoid  $X$  we denote by  $[X]$  the discrete groupoid consisting of the isomorphism classes of objects in  $X$  . If  $P : \underline{F} \to \underline{C}_{\Lambda}$  is a cofibered groupoid, we obtain a functor  $[\underline{F}] : \underline{C}_{\Lambda} \to$ (Ens)  given by $[\underline{F}](S) = [\underline{F}(S)]$  , and in turn a cofibered discrete groupoid, which is, by abuse of notation, denoted by  $[\underline{F}]$  .

\underline{Definition} 1.2.  \underline{A cofibered groupoid}  $\underline{A}$  \underline{over}  $\underline{C}_{\Lambda}$  \underline{is called semi-homogeneous if the following properties hold}:

(1) \underline{every diagram}

\begin{center}
\begin{tikzcd}
a'' \arrow[r, dashed] \arrow[d, dotted] \& a' \arrow[d] \\
a_1 \arrow[r] \& a
\end{tikzcd}
\end{center}

\underline{of solid arrows in}  $\underline{A}$  \underline{where}  $a' \to a$  \underline{is surjective can be completed into a commutative diagram including dotted arrows}.

(2) \underline{Let}

\begin{center}
\begin{tikzcd}[column sep=large]
\& R \\
R \underset{k}{\times} k[\varepsilon] \arrow[ur, "\pi_1"] \arrow[dr, "\pi_2"'] \& \\
\& k[\varepsilon]
\end{tikzcd}
\end{center}

\underline{be the projection maps, where}  $k[\varepsilon]$  \underline{is the ring of dual numbers over}  $k$  , \underline{and let}

\begin{center}
\begin{tikzcd}[column sep=huge]
a' \arrow[dr, "\alpha"] \& \\
\& a \\
b' \arrow[ur, "\beta"'] \&
\end{tikzcd}
\end{center}

\underline{be}  $\pi_1$\underline{-maps in}  $\underline{A}$  . \underline{Then there exists a}  $\underline{C}_{\Lambda}$  \underline{-map}  $\gamma : a' \longrightarrow b'$  \underline{with}  $\alpha = \beta\cdot\gamma$ \underline{if and only if}  $(\pi_2)_{*}(a') \approx (\pi_2)_{*}(b')$  .

\SGAvar{\underline{Remark} 1.3.(a)  A reformulation of the notion in terms of  2-categories will be given later in 2.6(b). We note that if  $\underline{A}$  is semi-homogeneous. Then so is  $[\underline{A}]$  , and in particular 1.2(2) entails the canonical map
\SGApage{I}{35}
$[\underline{A}(R \times_{k} k[\epsilon])] \longrightarrow [\underline{A}(R)] \times [\underline{A}(k[\epsilon])]$
to be bijective.

}{\underline{Remark} 1.3.(a)  A reformulation of the notion in terms of  2-categories will be given later in 2.6(b). We note that if  $\underline{A}$  is semi-homogeneous, then so is  $[\underline{A}]$  , and in particular 1.2(2) entails the canonical map
\SGApage{I}{35}
$[\underline{A}(R \times_{k} k[\epsilon])] \longrightarrow [\underline{A}(R)] \times [\underline{A}(k[\epsilon])]$
to be bijective.

}{F7I-VI-003}\underline{Remark} 1.3.(b)    Consider a commutative diagram

\begin{center}
\begin{tikzcd}[row sep=large, column sep=large]
 \& a' \arrow[dl, "v_{1}"'] \arrow[dr, "v_{2}"] \& \\
a_{1} \arrow[dr, "u_{1}"'] \& \& a_{2} \arrow[dl, "u_{2}"] \\
 \& a \&
\end{tikzcd}
\end{center}









in a cofibered category  $\underline{A}$  over  $\underline{C}_{\Lambda}$ . Set  $R = P(a)$ , $R' = P(a')$  and
$R_{i} = P(a_{i})$  for  $i = 1, 2$ . If we set  $a''$  to be the  "direct image"  under the
map  $(P(v_{1}),P(v_{2})) : R' \to R_{1} \times_{R} R_{2}$ , we obtain a commutative diagram

\begin{center}
\begin{tikzcd}[row sep=large, column sep=large]
 \& a' \arrow[d] \arrow[ddl, "v_{1}"'] \arrow[ddr, "v_{2}"] \& \\
 \& a'' \arrow[dl] \arrow[dr] \& \\
a_{1} \arrow[dr, "u_{1}"'] \& \& a_{2} \arrow[dl, "u_{2}"] \\
 \& a \&
\end{tikzcd}
\end{center}













Thus 1.2(1) can be replaced by

(1)':  every diagram

\begin{tikzcd}[cramped, column sep=small, row sep=small]
a_{1} \arrow[dr] \& \& a_{2} \arrow[dl] \\
 \& a \&
\end{tikzcd}  in  $\underline{A}$  above  \begin{tikzcd}[cramped, column sep=small, row sep=small]
R_{1} \arrow[dr] \& \& R_{2} \arrow[dl] \\
 \& R \&
\end{tikzcd}  in  $\underline{C}_{\Lambda}$ , where









$R_{2} \to R$  is surjective can be completed to a commutative diagram

\begin{center}
\begin{tikzcd}[row sep=large, column sep=large]
 \& a' \arrow[dl] \arrow[dr] \& \\
a_{1} \arrow[dr] \& \& a_{2} \arrow[dl] \\
 \& a \&
\end{tikzcd}  above  \begin{tikzcd}[row sep=large, column sep=large]
 \& R_{1} \times_{R} R_{2} \arrow[dl, "\pi_{1}"'] \arrow[dr, "\pi_{2}"] \& \\
R_{1} \arrow[dr] \& \& R_{2} \arrow[dl] \\
 \& R \&
\end{tikzcd}
\end{center}












where  $\pi_{i}$  ($i = 1,2$)  are projections.

\SGAvar{\underline{Remark} 1.3.(c)  If  $\underline{A}$  is a semi-homogeneous cofibered groupoid over  $\underline{C}_{\Lambda}$ , then
 $\underline{A}(k [\epsilon])]$ is canonically provided with a vector space structure over  $k$ , where
$k [\epsilon]$ is the ring of dual numbers over  $k$ . Indeed, let  $\underline{\underline{V}}$  be the category of
\SGApage{I}{36}
finite-dimensional vector spaces over  $k$ . Then any functor  $H : \underline{\underline{V}} \to (Ens)$
preserving the product is canonically factored through

}{\underline{Remark} 1.3.(c)  If  $\underline{A}$  is a semi-homogeneous cofibered groupoid over  $\underline{C}_{\Lambda}$ , then
 $[\underline{A}(k [\epsilon])]$ is canonically provided with a vector space structure over  $k$ , where
$k [\epsilon]$ is the ring of dual numbers over  $k$ . Indeed, let  $\underline{\underline{V}}$  be the category of
\SGApage{I}{36}
finite-dimensional vector spaces over  $k$ . Then any functor  $H : \underline{\underline{V}} \to (Ens)$
preserving the product is canonically factored through

}{F7I-VI-004}\begin{center}
\begin{tikzcd}[row sep=large, column sep=large]
\underline{\underline{V}} \arrow[rr, "H"] \arrow[dr, "\overline{H}"'] \& \& (Ens) \\
 \& \underline{\underline{V}} \arrow[ur] \&
\end{tikzcd}
\end{center}

where  $\underline{\underline{V}} \to (Ens)$  is the forgetful functor. On the other hand, the functor
$G : \underline{\underline{V}} \longrightarrow \underline{C}_{\wedge}$  given by  $W \longmapsto k[W] = k \oplus W$ , with  $W^{2} = 0$ , transforms
products in  $\underline{\underline{V}}$  into fibre products over  $k$  in  $\underline{C}_{\wedge}$ . If  $\underline{A}$  is a semi-homogeneous cofibered groupoid over  $\underline{C}_{\wedge}$ , then  $[\underline{A}(k[V] \times_{k} k[W])] \longrightarrow [\underline{A}(k[V])] \times [\underline{A}(k[W])]$
is bijective by 1.3(a), and the functor  $V \longmapsto [\underline{A}(k[V])]$  commutes with products,
hence the result.

\underline{Definition} 1.4. (1) \underline{Let}  $R \in ob(\underline{C}_{\wedge})$ . \underline{A small extension of}  $R$  \underline{in}  $\underline{C}_{\wedge}$  \underline{is a}
\underline{surjective map}  $R' \xrightarrow{\ f\ } R$  \underline{in}  $\underline{C}_{\wedge}$  \underline{such that}  $(\mathrm{Ker}\ f)^{2} = 0$  \underline{and}  $\mathrm{Ker}\ f \simeq k$ .
\underline{In other words,}  $R' \xrightarrow{\ f\ } R$  \underline{in}  $\underline{C}_{\wedge}$  \underline{is a small extension of}  $R$  \underline{if}

$0 \longrightarrow k \xrightarrow{\ t\ } R' \xrightarrow{\ f\ } R \longrightarrow 0$
\underline{is exact for some}  $t \in R'$  \underline{with}  $t^{2} = 0$ .

        (2)  \underline{Let}  $\underline{A}$  \underline{be a cofibered groupoid over}  $\underline{C}_{\wedge}$ , \underline{and let}  $a$  \underline{be an object}
\underline{in}  $\underline{A}$ . \underline{A small prolongation of}  $a$  \underline{is an arrow}  $a' \xrightarrow{\ \alpha\ } a$  \underline{in}  $\underline{A}$  \underline{such that}
$P(\alpha)$  \underline{is a small extension in}  $\underline{C}_{\wedge}$ . \underline{An arrow}  $a' \xrightarrow{\ \alpha\ } a$  \underline{in}  $\underline{A}$  \underline{will be called}
"\underline{versal for small prolongations of}  $a$  "  \underline{if, for every small prolongation}
$a'_{1} \longrightarrow a$  \underline{in}  $\underline{A}$ , \underline{there exists a commutative diagram}

\begin{center}
\begin{tikzcd}[row sep=large, column sep=large]
a' \arrow[rr] \arrow[dr, "\alpha"'] \& \& a'_{1} \arrow[dl] \\
 \& a \&
\end{tikzcd}
\end{center}

\underline{Proposition} 1.5.    \underline{Let}  $\underline{A}$  \underline{be a semi-homogeneous cofibered groupoid over}  $\underline{C}_{\wedge}$ .

(1) \underline{Let}  \begin{tikzcd}[cramped, row sep=tiny, column sep=small]
 \& R \\
R \times_{k} k[\epsilon] \arrow[ur, "\pi_{1}"] \arrow[dr, "\pi_{2}"'] \& \\
 \& k[\epsilon]
\end{tikzcd}  \underline{be projection maps in}  $\underline{C}_{\wedge}$ , \underline{and let}  $a' \xrightarrow{\ \alpha\ } a$








\SGApage{I}{37}
\underline{be a}  $\pi_1$-\underline{map.} \underline{Then there exists a map}  $a \xrightarrow{\ \beta\ } a'$  \underline{such that}  $\alpha\,\beta = 1_a$ \underline{if and only if there exists a map}  $a \longrightarrow (\pi_2)_* a'$  .

(2)  \underline{Let}  $R' \xrightarrow{\ f\ } R$  \underline{be a small extension in}  $\underline{C}_{\Lambda}$  , \underline{and consider a commutative} \underline{diagram}

\SGAvar{\begin{center}
\begin{tikzcd}[row sep=large, column sep=small]
 \& a'' \arrow[dl, "\alpha'"'] \arrow[dr, "\alpha'_1"] \& \& \& \& \& R'\underset{R}{\times}R' \arrow[dl] \arrow[dr] \& \\
a' \arrow[dr, "\alpha"'] \& \& a' \arrow[dl, "\alpha_1"] \& \& \underline{above} \& R' \arrow[dr] \& \& R' \arrow[dl] \\
 \& a \& \& \& \& \& R \&
\end{tikzcd}
\end{center}
















\underline{Define the canonical map}  $\tilde{f} : R' \underset{R}{\times} R' \longrightarrow k[\mathrm{Ker}\ f]$  by  $\tilde{f}(r'_1,\ r'_2) = r'_1(0)+(r'_1-r'_2)$ . \underline{If there exists a map}  $a' \longrightarrow (\tilde{f})_*(a'')$  , \underline{then there exists a map}  $\gamma: a' \longrightarrow a$ \underline{such that}  $\alpha = \alpha_1 \cdot \gamma$  .

\underline{Proof}  (1) The part  "only if"  is clear. Now assume that there exists a map $\varphi : a \longrightarrow (\pi_2)_* a'$  . If we set  $h : R \longrightarrow R \underset{k}{\times} k[\epsilon]$  by  $h = 1 \times P(\varphi)$  , then the composite map  $R \xrightarrow{\ h\ } R \underset{h}{\times} k[\epsilon] \xrightarrow{\ \pi_1\ } R$  is the identity, and therefore we get maps  $a \overset{\gamma}{\underset{\alpha_1}{\rightleftarrows}} f_* a$  such that  $\alpha_1 \cdot \gamma = 1_a$  where  $P(\gamma) = h$  and  $P(\alpha_1) = \pi_1$  . On the other hand,  $\pi_2 \cdot h = P(\varphi)$  and hence there exists a  $\pi_2$ -map $h_* a \longrightarrow (\pi_2)_* a'$  . This means that  $(\pi_2)_*(h_* a) = (\pi_2)_* a'$  , and therefore by 1.2(2) we have a  $\underline{C}_{\Lambda}$ -map  $u : h_* a \longrightarrow a'$  such that  $\alpha \cdot u = \alpha_1$  . Define $\beta : a \longrightarrow a'$  by  $\beta = u \cdot \gamma$  . Then  $\alpha \cdot \beta = \alpha \cdot u \cdot \gamma = \alpha_1 \cdot \gamma = 1_a$  .

(2) The map  $1 \times \tilde{f} : R' \underset{R}{\times} h' \longrightarrow R' \underset{k}{\times} k[\mathrm{Ker}\ f]$  is an isomorphism compatible with the projections on the first factor, and  $\pi_2 \cdot (1 \times \tilde{f}) = \tilde{f}$  where $\pi_2 : R' \underset{k}{\times} k[\mathrm{Ker}\ f] \longrightarrow k[\mathrm{Ker}\ f]$  is the projection map. We note that $k[\mathrm{Ker}\ f] \simeq k[\epsilon]$  since  $R' \xrightarrow{\ f\ } R$  is a small extension. Therefore, if there exists a map  $a' \longrightarrow (\tilde{f})_* a'' = (\pi_2)_* [(1\ \tilde{f})_* a'']$  , then by (1) above there exists a map  $\gamma': a' \longrightarrow a''$  such that  $\alpha' \cdot \gamma' = 1_{a'}$  . Define  $\gamma : a' \longrightarrow a$ by  $\gamma = \alpha'_1 \cdot \gamma'$  . We then have  $\alpha_1 \cdot \gamma = \alpha$  .


}{\begin{center}
\begin{tikzcd}[row sep=large, column sep=small]
 \& a'' \arrow[dl, "\alpha'"'] \arrow[dr, "\alpha'_1"] \& \& \& \& \& R'\underset{R}{\times}R' \arrow[dl] \arrow[dr] \& \\
a' \arrow[dr, "\alpha"'] \& \& a'_1 \arrow[dl, "\alpha_1"] \& \& \underline{above} \& R' \arrow[dr] \& \& R' \arrow[dl] \\
 \& a \& \& \& \& \& R \&
\end{tikzcd}
\end{center}
















\underline{Define the canonical map}  $\tilde{f} : R' \underset{R}{\times} R' \longrightarrow k[\mathrm{Ker}\ f]$  by  $\tilde{f}(r'_1,\ r'_2) = r'_1(0)+(r'_1-r'_2)$ . \underline{If there exists a map}  $a' \longrightarrow (\tilde{f})_*(a'')$  , \underline{then there exists a map}  $\gamma: a' \longrightarrow a'_1$ \underline{such that}  $\alpha = \alpha_1 \cdot \gamma$  .

\underline{Proof}  (1) The part  "only if"  is clear. Now assume that there exists a map $\varphi : a \longrightarrow (\pi_2)_* a'$  . If we set  $h : R \longrightarrow R \underset{k}{\times} k[\epsilon]$  by  $h = 1 \times P(\varphi)$  , then the composite map  $R \xrightarrow{\ h\ } R \underset{k}{\times} k[\epsilon] \xrightarrow{\ \pi_1\ } R$  is the identity, and therefore we get maps  $a \overset{\gamma}{\underset{\alpha_1}{\rightleftarrows}} h_* a$  such that  $\alpha_1 \cdot \gamma = 1_a$  where  $P(\gamma) = h$  and  $P(\alpha_1) = \pi_1$  . On the other hand,  $\pi_2 \cdot h = P(\varphi)$  and hence there exists a  $\pi_2$ -map $h_* a \longrightarrow (\pi_2)_* a'$  . This means that  $(\pi_2)_*(h_* a) = (\pi_2)_* a'$  , and therefore by 1.2(2) we have a  $\underline{C}_{\Lambda}$ -map  $u : h_* a \longrightarrow a'$  such that  $\alpha \cdot u = \alpha_1$  . Define $\beta : a \longrightarrow a'$  by  $\beta = u \cdot \gamma$  . Then  $\alpha \cdot \beta = \alpha \cdot u \cdot \gamma = \alpha_1 \cdot \gamma = 1_a$  .

(2) The map  $1 \times \tilde{f} : R' \underset{R}{\times} R' \longrightarrow R' \underset{k}{\times} k[\mathrm{Ker}\ f]$  is an isomorphism compatible with the projections on the first factor, and  $\pi_2 \cdot (1 \times \tilde{f}) = \tilde{f}$  where $\pi_2 : R' \underset{k}{\times} k[\mathrm{Ker}\ f] \longrightarrow k[\mathrm{Ker}\ f]$  is the projection map. We note that $k[\mathrm{Ker}\ f] \simeq k[\epsilon]$  since  $R' \xrightarrow{\ f\ } R$  is a small extension. Therefore, if there exists a map  $a' \longrightarrow (\tilde{f})_* a'' = (\pi_2)_* [(1 \times \tilde{f})_* a'']$  , then by (1) above there exists a map  $\gamma': a' \longrightarrow a''$  such that  $\alpha' \cdot \gamma' = 1_{a'}$  . Define  $\gamma : a' \longrightarrow a'_1$ by  $\gamma = \alpha'_1 \cdot \gamma'$  . We then have  $\alpha_1 \cdot \gamma = \alpha$  .


}{F7I-VI-005}\SGApage{I}{38}

\underline{Corollary} 1.6.   \underline{Let}  $\underline{A}$  \underline{be a semi-homogeneous cofibered groupoid over}  $\underline{C}_{\Lambda}$ \underline{and consider a diagram}

\begin{center}
\begin{tikzcd}[column sep=small]
a' \arrow[dr, "\alpha"'] \& \& a_1 \arrow[dl, "\alpha_1"] \\
 \& a \&
\end{tikzcd}
\end{center}

\SGAvar{\underline{in}  $\underline{A}$  \underline{where}  $\alpha_1$  \underline{is a small prolongation.} \underline{Assume that}  $a$  \underline{is versal for small} \underline{prolongations of}  $e$  . \underline{Then there exists a map}  $a' \overset{\ \rho\ }{\relbar\joinrel\relbar\joinrel\relbar} a_1$  \underline{such that}  $\alpha = \alpha_1 \rho$ \underline{in}  $\underline{A}$  \underline{if and only if there exists a map}  $f : P(a') \rightarrow P(a_1)$  \underline{such that} $P(\alpha) = P(\alpha_1)\cdot f$  \underline{in}  $\underline{C}_{\Lambda}$  .

}{\underline{in}  $\underline{A}$  \underline{where}  $\alpha_1$  \underline{is a small prolongation.} \underline{Assume that}  $a$  \underline{is versal for small} \underline{prolongations of}  $e$  . \underline{Then there exists a map}  $a' \xrightarrow{\ \rho\ } a_1$  \underline{such that}  $\alpha = \alpha_1 \rho$ \underline{in}  $\underline{A}$  \underline{if and only if there exists a map}  $f : P(a') \rightarrow P(a_1)$  \underline{such that} $P(\alpha) = P(\alpha_1)\cdot f$  \underline{in}  $\underline{C}_{\Lambda}$  .

}{F7I-VI-006}\underline{Proof}:  Assume that  $P(\alpha) = P(\alpha_1)\cdot f$  where  $f : P(a') \rightarrow P(a_1)$  . Replacing  $a'$ by its direct image under  $f$  , we may and shall assume that  $P(\alpha) = P(\alpha_1)$  . Set  $R' = P(a_1)$  ,  $R = P(a)$  . Since  $\underline{A}$  is semi-homogeneous, we get a commutative diagram

\begin{center}
\begin{tikzcd}[row sep=large, column sep=small]
 \& a'' \arrow[dl, "\beta"'] \arrow[dr, "\beta_1"] \& \& \& \& \& R'\underset{R}{\times}R' \arrow[dl] \arrow[dr] \& \\
a' \arrow[dr, "\alpha"'] \& \& a_1 \arrow[dl, "\alpha_1"] \& \& \text{above} \& R' \arrow[dr] \& \& R' \arrow[dl] \\
 \& a \& \& \& \& \& R \&
\end{tikzcd}
\end{center}
















Since  $a$  is versal for small prolongations of  $e$  , so is  $a'$ and therefore our statement follows from 1.5(2).

\SGAvar{\underline{Definition} 1.7.   \underline{For each object}  $R$  \underline{in}  $\underline{C}_{\Lambda}$ \underline{we put}  $\overline{\Omega}_R = M/(m_{\Lambda}R+M^2)$  \underline{where}  $\underline{m}_{\Lambda}$ $M$  \underline{are the maximal ideals of}  $\Lambda$ ,  $R$  \underline{respectively.} \underline{It is a finite-dimensional vector space over}  $k$  , \underline{and}  $R \longmapsto \overline{\Omega}_R$  \underline{defines a functor}  $\overline{\Omega} : \underline{C}_{\Lambda} \rightarrow \underline{\underline{V}}$  \underline{where} $\underline{\underline{V}}$  \underline{is the category of finite-dimensional vector spaces over}  $k$  . \underline{If}  $\underline{A}$  \underline{is a} \underline{cofibered category over}  $\underline{C}_{\Lambda}$  , \underline{then the composite functor}  $\underline{A} \xrightarrow{\ P\ } \underline{C}_{\Lambda} \rightarrow \underline{\underline{V}}$  \underline{is, by} \underline{abuse of notation denoted by}  $\overline{\Omega}$  \underline{again.}

}{\underline{Definition} 1.7.   \underline{For each object}  $R$  \underline{in}  $\underline{C}_{\Lambda}$ \underline{we put}  $\overline{\Omega}_R = M/(m_{\Lambda}R+M^2)$  \underline{where}  $\underline{m}_{\Lambda}$ , $M$  \underline{are the maximal ideals of}  $\Lambda$ ,  $R$  \underline{respectively.} \underline{It is a finite-dimensional vector space over}  $k$  , \underline{and}  $R \longmapsto \overline{\Omega}_R$  \underline{defines a functor}  $\overline{\Omega} : \underline{C}_{\Lambda} \rightarrow \underline{\underline{V}}$  \underline{where} $\underline{\underline{V}}$  \underline{is the category of finite-dimensional vector spaces over}  $k$  . \underline{If}  $\underline{A}$  \underline{is a} \underline{cofibered category over}  $\underline{C}_{\Lambda}$  , \underline{then the composite functor}  $\underline{A} \xrightarrow{\ P\ } \underline{C}_{\Lambda} \rightarrow \underline{\underline{V}}$  \underline{is, by} \underline{abuse of notation denoted by}  $\overline{\Omega}$  \underline{again.}

}{F7I-VI-007}\underline{Proposition} 1.8.   \underline{Let}  $\underline{A}$  \underline{be a semi-homogeneous cofibered groupoid over}  $\underline{C}_{\Lambda}$  . \underline{Then, given an object}  $a$  \underline{in}  $\underline{A}$  \underline{versal for small prolongations of}  $e$  , \underline{there} \underline{exists}  $a' \rightarrow a$  \underline{such that}
\SGApage{I}{39}
(i)  $a'$  \underline{is versal for small prolongations of}  $a$

(ii)  $p(a')$ $\longrightarrow$ $p(a)$  \underline{is surjective; and the kernel is annihilated by the maximal ideal of}  $p(a)$  .

(iii)  $\overline{\Omega}_{a'} \longrightarrow \overline{\Omega}_{a}$  \underline{is an isomorphism}.

\underline{Proof}:  Set  $R = P(a)$  and consider an exact sequence  $0 \to J \to S \to R \to 0$  where $S$  is a formal power-series ring over  $\wedge$  in a finite number of variables, and $J \subset M^2$  where  $M$  is the maximal ideal of  $S$  . Let  $\Sigma$  be the set of ideals  $J_1$ in  $S$  such that (i)  $J \supset J_1 \supset MJ$  and (ii) there exists  $\pi_1$-morphism  $a_1 \to a$  in  $\underline{A}$  where  $\pi_1 : S/J_1 \to R$  is the canonical map. The property 1.2(1) of  $\underline{A}$  together with the fact that  $J/MJ$  is a finite dimensional vector space over  $k$  entails that there exists a smallest ideal in the family  $\Sigma$  . Let it be  $J'$  and choose a  $\pi'$-morphism  $a' \overset{g}{\to} a$  where  $\pi' : S/J' \to R$  is the canonical map. Note that $\overline{\Omega}_{a'} \to \overline{\Omega}_{a}$  is an isomorphism since  $J' \subset J \subset M^2$  . We claim that  $a' \overset{g}{\to} a$  is versal for small prolongations of  $a$  . Indeed, let  $a_1 \to a$  be a small prolongation. If we put  $R_1 = P(a_1)$  , we obtain an exact commutative diagram

\begin{center}
\begin{tikzcd}[column sep=small]
0 \arrow[r] \& J/MJ \arrow[r] \arrow[d, "h"] \& S/MJ \arrow[r] \arrow[d, "h_1"] \& R \arrow[r] \arrow[d, equal] \& 0 \\
0 \arrow[r] \& k \arrow[r] \& R_1 \arrow[r] \& R \arrow[r] \& 0
\end{tikzcd}
\end{center}









\SGAvar{If  $h = 0$  , then  $\pi_1 : R_1 \to R$  admits a cross-section so that  $h_1$  induces $f : S/J'$ $\to R_1$  . If  $h \neq 0$  , then  $h_1 : S/MJ \to R_1$  is surjective and hence  $h_1$ induces  $f : S/J'$ $\to R_1$  by the definiton of the ideal  $J'$  . Therefore in either cases we obtain a commutative diagram

\begin{center}
\begin{tikzcd}
S/J' \arrow[rr, "f", no head] \arrow[dr, "\pi'"'] \& \& R_1 \arrow[dl, "\pi_1"] \\
\& R \&
\end{tikzcd}
\end{center}

}{If  $h = 0$  , then  $\pi_1 : R_1 \to R$  admits a cross-section so that  $h_1$  induces $f : S/J'$ $\to R_1$  . If  $h \neq 0$  , then  $h_1 : S/MJ \to R_1$  is surjective and hence  $h_1$ induces  $f : S/J'$ $\to R_1$  by the definition of the ideal  $J'$  . Therefore in either case we obtain a commutative diagram

\begin{center}
\begin{tikzcd}
S/J' \arrow[rr, "f"] \arrow[dr, "\pi'"'] \& \& R_1 \arrow[dl, "\pi_1"] \\
\& R \&
\end{tikzcd}
\end{center}

}{F7I-VI-008}Since  $a$  is versal for small prolongations of  $e$  , it follows from 1.6 that we obtain a commutative diagram
\SGApage{I}{40}
\begin{tikzcd}
a' \arrow[rr, "\rho"] \arrow[dr] \& \& a_1 \arrow[dl] \\
\& a \&
\end{tikzcd}

Let  $P : \underline{A} \to \underline{C}_\wedge$  be a cofibered category (in groupoids). It is then clear that we obtain a cofibered category  $\mathrm{pro}(P) : \mathrm{pro}(\underline{A}) \to \mathrm{pro}(\underline{C}_\wedge)$  (in groupoids) (pour la définition de la catégorie de pro-objets  $\mathrm{pro}(\underline{A})$  , voir A. Grothendieck, Sém. Bourbaki 195).

\begin{center}
\begin{tikzcd}
\underline{A} \arrow[r] \arrow[d, "P"'] \& \mathrm{pro}(\underline{A}) \arrow[d, "\mathbf{pro}(P)"] \\
\underline{C}_\wedge \arrow[r] \& \mathrm{pro}(\underline{C}_\wedge)
\end{tikzcd}
\end{center}









Now let  $\widehat{\underline{C}}_\wedge$  be the full subcategory of  pro-$\underline{C}_\wedge$  in which $\mathrm{ob}(\widehat{\underline{C}}_\wedge) = \{(R_i)_{i \in \underline{N}}$  ,  $R_i \in \mathrm{ob}(\underline{C}_\wedge) \mid \varphi_i : R_i \to R_{i-1}$  is surjective for all  $i$  and induces isomorphisms  $\underline{m}_i/\underline{m}_\wedge + \underline{m}_i^2 \to \underline{m}_{i-1}/\underline{m}_\wedge + \underline{m}_{i-1}^2$  for all sufficiently large  $i$  $\}$ where  $\underline{N}$  stands for the ordered category of natural numbers, and  $\underline{m}_i$  is the maximal ideal of  $R_i$  . We set  $\widehat{\underline{A}}$  to be the full subcategory of pro-$\underline{A}$  in which

\begin{equation*}
\mathrm{ob}(\widehat{\underline{A}}) = \{a \in \mathrm{ob}(\mathrm{pro}(\underline{A})) \mid \mathrm{pro}(P)\ (a) \in \mathrm{ob}(\widehat{\underline{C}}_\wedge)\}
\end{equation*}
We then obtain a cofibered category  $\widehat{P} : \widehat{\underline{A}} \to \widehat{\underline{C}}_\wedge$  (in groupoids) with a commutative diagram

\begin{center}
\begin{tikzcd}
\underline{A} \arrow[r, hook] \arrow[d, "P"'] \& \widehat{\underline{A}} \arrow[d, "\widehat{P}"] \\
\underline{C}_\wedge \arrow[r, hook] \& \widehat{\underline{C}}_\wedge
\end{tikzcd}
\end{center}










We note that the category  $\widehat{\underline{C}}_\wedge$  is canonically equivalent to the category of complete local $\wedge$-algebras of formally finite type over  $\wedge$  with the fixed residue field  $k$  . Henceforth we shall make such identification.

\underline{Definition} 1.9.   \underline{Let}  $\underline{A}$  \underline{be a cofibered groupoid over}  $\underline{C}_\wedge$  . \underline{An object}  $\xi$  \underline{in} $\widehat{\underline{A}}$  \underline{is called a quasi-initial object of}  $\widehat{\underline{A}}$  \underline{if for every object}  $a$  \underline{in}  $\widehat{\underline{A}}$  \underline{there exists a map}  $\xi \to a$  \underline{in}  $\widehat{\underline{A}}$  . \underline{A quasi-initial object}  $\xi$  \underline{of}  $\widehat{\underline{A}}$  \underline{is called minimal}
\SGApage{I}{41}
\underline{if the morphism}  $\xi : h_R \to [\hat{\underline{A}}]$  \underline{induces a bijection} $\xi(k[\epsilon])$: $\mathrm{Hom}_{\wedge\text{-alg}}(R,k[\epsilon]) \to [\underline{A}(k[\epsilon])]$ , \underline{where} $R = P(\xi)$ . \underline{An object} $\xi$ \underline{in} $\hat{\underline{A}}$ \underline{is called a versal formal object of} $\underline{A}$ \underline{if every surjective map} $\alpha : a' \to a$ \underline{in} $\underline{A}$ \underline{induces a surjective map} $\mathrm{Hom}_{\hat{\underline{A}}}(\xi,a') \to \mathrm{Hom}_{\hat{\underline{A}}}(\xi,a)$ .

\underline{A versal formal object} $\xi$ \underline{of} $\underline{A}$ \underline{is called minimal if the morphism} $\xi : h_R \to [\hat{\underline{A}}]$ \underline{induces a bijection} $\xi(k[\epsilon]) : \mathrm{Hom}_{\wedge\text{-alg}}(R,k[\epsilon]) \to [\underline{A}(k[\epsilon])]$ \underline{where} $R = P(\xi)$ .

\underline{Proposition} 1.10.  \underline{Let} $\underline{A}$ \underline{be a cofibered groupoid over} $\underline{C}_{\wedge}$ .

  (i)  \underline{every (minimally) versal formal object of} $\underline{A}$ \underline{is a quasi-initial (minimal) object of} $\hat{\underline{A}}$ .

\SGAvar{  (ii) \underline{If} $\xi$ \underline{is an object in} $\hat{A}$ \underline{such that} $\widetilde{\xi}(k[\epsilon])$: $\mathrm{Hom}_{\wedge\text{-alg}}(R,k[\epsilon])\to [\underline{A}(k[\epsilon])]$ \underline{is injective, where} $P(\xi) = R$ , \underline{then every endomorphism of} $\xi$ \underline{in} $\hat{A}$ \underline{is an automorphism. In particular a quasi-initial minimal object of} $\hat{\underline{A}}$ , \underline{if it exists, is unique up to (non-canonical) isomorphism.}

\underline{Proof}: (i) let $\xi$ be a versal formal object of $\underline{A}$ . We must show that, for any object $a$ in $\hat{\underline{A}}$ , we have a map $\xi \to a$ . By the definition of $\hat{\underline{A}}$ , there exists a sequence of surjective maps in $\underline{A}$ : $--- \to a_n \xrightarrow{\ \alpha_n\ } a_{n-1} \to --- \to a_1 \overset{\alpha_1}{\relbar\joinrel\relbar} a_0$ such that $a = \varprojlim_n a_n$ . Assume that there exists maps $\varphi_i : \xi \to a_i$ such that $\alpha_i\varphi_i = \varphi_{i-1}$ for all $i < n$ . Since $\xi$ is a versal formal object of $\underline{A}$ and $\alpha_n : a_n \to a_{n-1}$ is surjective in $\underline{A}$ , we obtain a map $\xi \xrightarrow{\ \varphi_n\ } a_n$ such that $\alpha_n\varphi_n = \varphi_{n-1}$ . If we set $\varphi = \varprojlim_n \varphi_n$ , we have $\varphi : \xi \to a$ .

}{  (ii) \underline{If} $\xi$ \underline{is an object in} $\hat{\underline{A}}$ \underline{such that} $\widetilde{\xi}(k[\epsilon])$: $\mathrm{Hom}_{\wedge\text{-alg}}(R,k[\epsilon])\to [\underline{A}(k[\epsilon])]$ \underline{is injective, where} $P(\xi) = R$ , \underline{then every endomorphism of} $\xi$ \underline{in} $\hat{\underline{A}}$ \underline{is an automorphism. In particular a quasi-initial minimal object of} $\hat{\underline{A}}$ , \underline{if it exists, is unique up to (non-canonical) isomorphism.}

\underline{Proof}: (i) let $\xi$ be a versal formal object of $\underline{A}$ . We must show that, for any object $a$ in $\hat{\underline{A}}$ , we have a map $\xi \to a$ . By the definition of $\hat{\underline{A}}$ , there exists a sequence of surjective maps in $\underline{A}$ : $--- \to a_n \xrightarrow{\ \alpha_n\ } a_{n-1} \to --- \to a_1 \overset{\alpha_1}{\relbar\joinrel\relbar} a_0$ such that $a = \varprojlim_n a_n$ . Assume that there exist maps $\varphi_i : \xi \to a_i$ such that $\alpha_i\varphi_i = \varphi_{i-1}$ for all $i < n$ . Since $\xi$ is a versal formal object of $\underline{A}$ and $\alpha_n : a_n \to a_{n-1}$ is surjective in $\underline{A}$ , we obtain a map $\xi \xrightarrow{\ \varphi_n\ } a_n$ such that $\alpha_n\varphi_n = \varphi_{n-1}$ . If we set $\varphi = \varprojlim_n \varphi_n$ , we have $\varphi : \xi \to a$ .

}{F7I-VI-009}(ii) Let $\xi \xrightarrow{\ \alpha\ } \xi$ be a map in $\hat{A}$ . Then for every $\lambda \in \mathrm{Hom}_{\wedge\text{-alg}}(R,k[\epsilon])$ we have $\lambda_*\xi = \lambda_*\alpha\xi \simeq (\lambda\cdot P(\alpha))_*\xi$ and hence by the hypothesis on $\xi$ we have that $\lambda = \lambda\cdot P(\alpha)$ for all $\lambda \in \mathrm{Hom}_{\wedge\text{-alg}}(R,k[\epsilon])$ , and therefore $P(\alpha) : R \to R$ is surjective. However it is trivial to see that every surjective endomorphism of a noetherian ring is always an automorphism and hence $P(\alpha)$ is an automorphism of $R$ and therefore $\alpha$ is an automorphism of $\xi$ .

\underline{Theorem} 1.11. (Schlessinger) \underline{Let} $\underline{A}$ \underline{be a cofibered groupoid over} $\underline{C}_{\Lambda}$ . \underline{Then} $\underline{A}$
\SGApage{I}{42}
\underline{admits a minimally-versal formal object if and only if}

  (i)  $\underline{A}$  \underline{is semi-homogeneous}

  (ii)  $[\underline{A}(k[\epsilon])]$  \underline{is a finite-dimensional vector space over}  $k$  (cf. 1.3(c)).

\underline{Proof}; (\underline{Necessity}) Let $\xi$ in $\hat{\underline{A}}$ be a minimally-versal formal object of $\underline{A}$ , and consider a diagram

\begin{center}
\begin{tikzcd}[column sep=small, row sep=large]
 a_1 \arrow[dr] \& \& a_2 \arrow[dl, "\alpha"] \\
 \& a \&
\end{tikzcd}
\end{center}

\SGAvar{in $\underline{A}$ where $\alpha$ is surjective. Since $\xi$ is a quasi-initial object of $\hat{\underline{A}}$ , there exists a map $\xi \to a_1$ . Since $\alpha$ is surjective, the composit map $\xi \to a_1 \to a$ can be lifted to $a_2$ i.e., we obtain a commutative diagram

}{in $\underline{A}$ where $\alpha$ is surjective. Since $\xi$ is a quasi-initial object of $\hat{\underline{A}}$ , there exists a map $\xi \to a_1$ . Since $\alpha$ is surjective, the composite map $\xi \to a_1 \to a$ can be lifted to $a_2$ i.e., we obtain a commutative diagram

}{F7I-VI-010}\begin{center}
\begin{tikzcd}[column sep=small, row sep=large]
 \& \xi \arrow[dl] \arrow[dr] \& \\
 a_1 \arrow[dr] \& \& a_2 \arrow[dl, "\alpha"] \\
 \& a \&
\end{tikzcd}
\end{center}

which proves the condition 1.2(1). We now prove the condition 1.2(2). Let

\SGAvar{\begin{center}
\begin{tikzcd}[column sep=large]
 \& S \\
 S{\times}k[\epsilon] \arrow[ur, "\pi_1"] \arrow[dr, "\pi_2"'] \& \\
 \& k[\epsilon]
\end{tikzcd}
\end{center}

}{\begin{center}
\begin{tikzcd}[column sep=large]
 \& S \\
 S{\times}_k k[\epsilon] \arrow[ur, "\pi_1"] \arrow[dr, "\pi_2"'] \& \\
 \& k[\epsilon]
\end{tikzcd}
\end{center}

}{F7I-VI-011}be the projections,

and let

\begin{center}
\begin{tikzcd}[column sep=small, row sep=large]
 a' \arrow[dr, "\alpha"'] \& \& b' \arrow[dl, "\beta"] \\
 \& a \&
\end{tikzcd}
\end{center}

be $\pi_1$-maps in $\underline{A}$ . By the definition of $\xi$ , we obtain a commutative diagram

\begin{center}
\begin{tikzcd}[column sep=small, row sep=large]
 \& \xi \arrow[dl, "u"'] \arrow[dr, "v"] \& \\
 a' \arrow[dr, "\alpha"'] \& \& b' \arrow[dl, "\beta"] \\
 \& a \&
\end{tikzcd}
\end{center}

\SGAvar{and in particular we have $\pi_1\cdot P(u) = \pi_1\cdot P(v)$ . Assume now that $(\pi_2)_*a' \approx (\pi_2)_*b'$ . We then have $\pi_2\cdot P(u) = \pi_2\cdot P(v)$ because of the minimality of
\SGApage{I}{43}
$\xi$ , and hence  $P(u) = \pi_1 P(u)\times \pi_2 P(u) = \pi_1 P(v)\times \pi_2 P(v) = P(v)$  . Since  $\underline{A}$  is a cofibered groupoid, we obtain a unique isomorphism  $a' \overset{\sigma}{\to} b'$  in $\underline{A}(S\times k[\epsilon])$  such that  $\sigma u = v$  . Then  $\beta\sigma u = \alpha u$  entails that  $\beta\sigma= \alpha$  since  $\underline{A}$ is a cofibered category.

(\underline{Sufficiency}) Let  $c_1,c_2,\dots,c_n$  be objects in  $\underline{A}(k[\epsilon])$  which form a $k$-basis of the vector space  $[\underline{A}(k[\epsilon])]$  , and choose  $a_1$  in  $\underline{A}(k[\epsilon]\times_k\dots\times_k k[\epsilon])$ such that  $(\pi_j)_* a_1 \approx c_j$  . Then  $\tilde{a}(k[\epsilon]) : \mathrm{Hom}_{\Lambda\text{-alg}}(R_1,k[\epsilon]) \to [\underline{A}(k[\epsilon])]$  is bijective, where  $R_1 = P(a_1)$  . In particular  $a_1 \to e$  is versal for small prolongations of  $e$  . Choose  $a_2 \to a_1$  which is versal for small prolongations of  $a_1$ such that  $\overline{\Omega}_{a_2} \to \overline{\Omega}_{a_1}$  is bijective. We then choose  $a_3 \to a_2$  , versal for small prolongations of  $a_2$  such that  $\overline{\Omega}_{a_3} \to \overline{\Omega}_{a_2}$  , is bijective etc., and set $\xi = \underset{n}{\underleftarrow{\mathrm{Lim}}}\, a_n$  . We claim that  $\xi$  is a minimally-versal formal object of  $\underline{A}$  . Indeed, since  $\overline{\Omega}_{a_n} \to \overline{\Omega}_{a_{n-1}}$  is bijective, for all  $n$  ,  $\overline{\Omega}_{\xi} \to \overline{\Omega}_{a_1}$  is bijective so that $\xi(k[\epsilon]) : \mathrm{Hom}_{\Lambda\text{-alg}}(R,k[\epsilon] \to [\underline{A}(k[\epsilon])]$  is bijective. Now consider a diagram

}{and in particular we have $\pi_1\cdot P(u) = \pi_1\cdot P(v)$ . Assume now that $(\pi_2)_*a' \approx (\pi_2)_*b'$ . We then have $\pi_2\cdot P(u) = \pi_2\cdot P(v)$ because of the minimality of
\SGApage{I}{43}
$\xi$ , and hence  $P(u) = \pi_1 P(u)\times \pi_2 P(u) = \pi_1 P(v)\times \pi_2 P(v) = P(v)$  . Since  $\underline{A}$  is a cofibered groupoid, we obtain a unique isomorphism  $a' \overset{\sigma}{\to} b'$  in $\underline{A}(S\times_k k[\epsilon])$  such that  $\sigma u = v$  . Then  $\beta\sigma u = \alpha u$  entails that  $\beta\sigma= \alpha$  since  $\underline{A}$ is a cofibered category.

(\underline{Sufficiency}) Let  $c_1,c_2,\dots,c_n$  be objects in  $\underline{A}(k[\epsilon])$  which form a $k$-basis of the vector space  $[\underline{A}(k[\epsilon])]$  , and choose  $a_1$  in  $\underline{A}(k[\epsilon]\times_k\dots\times_k k[\epsilon])$ such that  $(\pi_j)_* a_1 \approx c_j$  . Then  $\tilde{a}_1(k[\epsilon]) : \mathrm{Hom}_{\Lambda\text{-alg}}(R_1,k[\epsilon]) \to [\underline{A}(k[\epsilon])]$  is bijective, where  $R_1 = P(a_1)$  . In particular  $a_1 \to e$  is versal for small prolongations of  $e$  . Choose  $a_2 \to a_1$  which is versal for small prolongations of  $a_1$ such that  $\overline{\Omega}_{a_2} \to \overline{\Omega}_{a_1}$  is bijective. We then choose  $a_3 \to a_2$  , versal for small prolongations of  $a_2$  such that  $\overline{\Omega}_{a_3} \to \overline{\Omega}_{a_2}$  is bijective etc., and set $\xi = \underset{n}{\underleftarrow{\mathrm{Lim}}}\, a_n$  . We claim that  $\xi$  is a minimally-versal formal object of  $\underline{A}$  . Indeed, since  $\overline{\Omega}_{a_n} \to \overline{\Omega}_{a_{n-1}}$  is bijective, for all  $n$  ,  $\overline{\Omega}_{\xi} \to \overline{\Omega}_{a_1}$  is bijective so that $\xi(k[\epsilon]) : \mathrm{Hom}_{\Lambda\text{-alg}}(R,k[\epsilon]) \to [\underline{A}(k[\epsilon])]$  is bijective. Now consider a diagram

}{F7I-VI-012}\begin{center}
\begin{tikzcd}
 \& a' \arrow[d, "\alpha"] \\
\xi \arrow[r, "h"] \& a
\end{tikzcd}
\end{center}

\SGAvar{where  $\alpha$  is surjective in  $\underline{A}$  . We must show that  $h$  can be lifted to  $a'$  . For this, one may assume, by induction, that  $a'$  is a small prolongation of  $a$  . Since  $P(a)$  is an artinian ring,  $h : \xi \to a$  is factored as a composit map $\xi \to a_q \longrightarrow a$  for some  $q$  . Since  $\underline{A}$  is semi-homogeneous, we get a diagram

}{where  $\alpha$  is surjective in  $\underline{A}$  . We must show that  $h$  can be lifted to  $a'$  . For this, one may assume, by induction, that  $a'$  is a small prolongation of  $a$  . Since  $P(a)$  is an artinian ring,  $h : \xi \to a$  is factored as a composite map $\xi \to a_q \longrightarrow a$  for some  $q$  . Since  $\underline{A}$  is semi-homogeneous, we get a diagram

}{F7I-VI-013}\begin{center}
\begin{tikzcd}
 \& a'' \arrow[r] \arrow[d, "\alpha'"] \& a' \arrow[d, "\alpha"] \\
\xi \arrow[r] \& a_q \arrow[r] \& a
\end{tikzcd}
\end{center}

where  $\alpha'$  is also a small prolongation. We then obtain a commutative diagram

\begin{center}
\begin{tikzcd}
 \& \& a'' \arrow[d, "\alpha'"] \\
\xi \arrow[r] \& a_{q+1} \arrow[r] \arrow[ur] \& a_q
\end{tikzcd}
\end{center}

and therefore we are done.

\underline{Remark} 1.12.  Let
\SGApage{I}{44}
\begin{tikzcd}
\dots\dots \arrow[r] \& a_n \arrow[r] \& a_{n-1} \arrow[r] \& \dots \arrow[r] \& a_2 \arrow[r] \& a_1 \arrow[r] \& a_o = e
\end{tikzcd}

\SGAvar{be an infinite sequences of arrows in  $\underline{A}$  such that

}{be an infinite sequence of arrows in  $\underline{A}$  such that

}{F7I-VI-014}\begin{equation*}
\left\{\begin{array}{l}
a_i \ \text{ is versal for small prolongations of } \ a_{i-1} \\[2pt]
P(a_i) \longrightarrow P(a_{i-1}) \ \text{ is surjective} \\[2pt]
\overline{\Omega}_{a_i} \dashrightarrow \overline{\Omega}_{a_{i-1}} \ \text{ is bijective for } \ i >> 0
\end{array}\right.
\end{equation*}
Then the above proof shows that  $\underleftarrow{\mathrm{Lim}}\, a_i$  is a versal formal object of  $\underline{A}$  .

\underline{Remark} 1.13.  Let  $\underline{A}$  be a cofibered groupoid over  $\underline{C}_\wedge$  . If  $\underline{A}$  has the property 1.11(i) and (ii), then so does  $[\underline{A}]$  , and  $\xi \in \mathrm{ob}(\hat{\underline{A}})$  is a minimally-versal formal object of  $\underline{A}$  if and only if it is so for  $[\underline{A}]$  by 1.10(ii).

\underline{Remark} 1.14.  Let  $\underline{A}$  be a cofibered groupoid over  $\underline{C}_\wedge$  , admitting a minimally-versal formal object  $\xi$  . Then an invariant of  $R = P(\xi)$  , which is a complete local  $\wedge$-algebra of formally finite type over  $\wedge$  , is automatically an invariant of  $\underline{A}$  over  $\underline{C}_\wedge$  . Therefore, for example, we shall say that  $\underline{A}$  is (formally) smooth over  $\underline{C}_\wedge$  if  $R$  is (formally) smooth over  $\wedge$  i.e., a formal power-series algebra over  $\wedge$  , or we say that  $\underline{A}$  is irreducible over  $\underline{C}_\wedge$  if  $\mathrm{Spec}(R)$  is irreducible etc. In the next section we introduce the notion of  "smoothness"  of a functor  $\underline{A} \to \underline{B}$  over  $\underline{C}_\wedge$  , where  $\underline{A}, \underline{B}$  are cofibered categories in groupoids over  $\underline{C}_\wedge$  . It turns out (cf. 2.4) that  "$\underline{A}$  is (formally) smooth over  $\underline{C}_\wedge$"  is equivalent to  " $P : \underline{A} \to \underline{C}_\wedge$  is a smooth functor"  , and therefore no confusion arises. We also define the dimension of  $\underline{A}$  over  $\underline{C}_\wedge$  by  $\dim_{\underline{C}_\wedge} \underline{A} = \dim\, k \otimes_\wedge R$  . For this we note the following fact:  Let  $\underline{C}_k$  be the category of artinian local $k$-algebras with the same residue field  $k$  . Then  $\underline{C}_k \to \underline{C}_\wedge$  is a fully-faithful imbedding, whose (left) adjoint functor  $\underline{C}_\wedge \to \underline{C}_k$  is given by  $R \to k \otimes_\wedge R$  . Let $\underline{A}|\underline{C}_k$  be the pull-back of  $P : \underline{A} \to \underline{C}_\wedge$  under  $\underline{C}_k \to \underline{C}_\wedge$  . Then it is trivial to see that if  $\xi$  is a minimally-versal formal object for  $\underline{A}$  over  $\underline{C}_\wedge$  , then  $k \otimes_\wedge \xi$ is a minimally-versal formal object for  $\underline{A}|\underline{C}_k$  over  $\underline{C}_k$  , where  $\xi \to k \otimes_\wedge \xi$  is a cocartesian arrow above  $R \to k \otimes_\wedge R$  . It follows that  $\dim_{\underline{C}_\wedge} \underline{A} = \dim_{\underline{C}_k} (\underline{A}|\underline{C}_k)$  .


\SGApage{I}{45}

\SGAvar{1.15.  In the remaining of this paragraph, we biefly explain how to modify the preceeding theory in order to allow for residue field extensions.

}{1.15.  In the remaining of this paragraph, we briefly explain how to modify the preceding theory in order to allow for residue field extensions.

}{F7I-VI-015}Let $\Lambda$ and $K$ be complete noetherian local rings, with $K$ a $\Lambda$-algebra. We assume that

(a)  the natural map from $\Lambda$ to $K$ is a local homomorphism;

(b)  the residue field $K'$ of $K$ is an extension of finite type of the residue field $k$ of $\Lambda$ .

The most interesting case is when $K = K'$ and $K$ is a finite inseparable extension of $k$ .

We denote by $\hat{C}_{\Lambda,K}$ the category whose objects are the complete noetherian local rings $R$ , provided with a structure of $\Lambda$-algebra and with a $\Lambda$-epimorphism $s : R \longrightarrow K$ .

We denote by $C_{\Lambda,K}$ the subcategory of $\hat{C}_{\Lambda,K}$ consisting of the $(R, s)$ such that $\mathrm{Ker}(s)$ is a $R$-module of finite length. The category $\hat{C}_{\Lambda,K}$ can be identified with a full subcategory of $\mathrm{Pro}(C_{\Lambda,K})$ .

If $K$ is artinian, so is any object in $C_{\Lambda,K}$ .

For any finite dimensional vector space $V$ over $K'$ we denote by $D_V(K)$ the algebra of dual numbers over $K$ defined by $V$ ; its elements are of the form $k + v\,\varepsilon$ (with $k \in K$ , $v \in V$ and $\varepsilon^2 = 0$ : $(k + v\,\varepsilon)\,(k' + v'\,\varepsilon) = (k\,k') + (k\,v' + k'\,v)\varepsilon$  ). The map $s : D_V(K) \to K$ : $s(k + v\,\varepsilon) = k$ , turns $D_V(K)$ into an object of $C_{\Lambda,K}$ . The role played before by $k[\varepsilon]$ will be played here by $D_{K'}(K)$ .

\SGAvar{As explained in 1.8, any cofibered groupoid $\underline{A}$ over $C_{\Lambda,K}$ extends in a natural way to a cofibered category $\hat{\underline{A}}$ over $\hat{C}_{\Lambda,K}$ . A map in $\underline{A}$ or $\hat{\underline{A}}$ will be called \underline{surjective} if the corresponding homomorphism of algebras in $C_{\Lambda,k}$ or in $\hat{C}_{\Lambda,K}$ is.

}{As explained in 1.8, any cofibered groupoid $\underline{A}$ over $C_{\Lambda,K}$ extends in a natural way to a cofibered category $\hat{\underline{A}}$ over $\hat{C}_{\Lambda,K}$ . A map in $\underline{A}$ or $\hat{\underline{A}}$ will be called \underline{surjective} if the corresponding homomorphism of algebras in $C_{\Lambda,K}$ or in $\hat{C}_{\Lambda,K}$ is.

}{F7I-VI-016}\SGApage{I}{46}

\underline{We assume that} $\underline{A}(K) = \{e\}$ .

\underline{Definition} 1.16.  \underline{A cofibered groupoïd} $\underline{A}$ \underline{over} $C_{\Lambda,K}$ \underline{is called} \underline{semi-homogeneous} \underline{if the following properties hold:}

(I)  \underline{every diagram}

\begin{center}
\begin{tikzcd}[row sep=large, column sep=large]
a'' \arrow[r, dashed] \arrow[d, dashed] \& a' \arrow[d] \\
a_1 \arrow[r] \& a
\end{tikzcd}
\end{center}











\underline{of solid arrows in} $\underline{A}$ , \underline{where} $a' \longrightarrow a$ \underline{is surjective, can be completed into} \underline{a commutative diagram including dotted arrows;}

(II)  \underline{For any} $R$ , \underline{the natural map}

\begin{equation*}
[\underline{A}(R \times_K D_{K'}(K))] \to [\underline{A}(R) \times \underline{A}\,(D_{K'}(K))] \simeq [\underline{A}(R)] \times [\underline{A}(D_{K'}(K))]
\end{equation*}
\underline{is bijective.}

If $\underline{A}$ is semi-homogeneous, so is $[\underline{A}]$ .

1.17.  Let us first consider the restriction of $\underline{A}$ to the full subcategory of $C_{\Lambda,K}$ whose objects are the $D_V(K)$ . Let $\Omega$ be the vector space $\Omega^1_{K/\Lambda} \otimes_K K'$ , and $d : K \to \Omega$ the derivative. A map $f : D_V(K) \to D_W(K)$ can be written:

\SGAvar{\begin{equation*}
f(k + v\,\varepsilon) = k + (u(v) + D\,d\,k)
\end{equation*}
with $u \in \mathrm{Hom}_{K'}(V, W)$ and $D \in \mathrm{Hom}_{K'}(\Omega, W)$ . If we identify $f$ with $(u, D)$ , the composition law is

}{\begin{equation*}
f(k + v\,\varepsilon) = k + (u(v) + D\,d\,k)\,\varepsilon
\end{equation*}
with $u \in \mathrm{Hom}_{K'}(V, W)$ and $D \in \mathrm{Hom}_{K'}(\Omega, W)$ . If we identify $f$ with $(u, D)$ , the composition law is

}{F7I-VI-017}\begin{equation*}
(u, D) \circ (u', D') = (uu', D + u\,D')\ .
\end{equation*}
The condition (II) entails

(II')  the functor from $K'$-vector spaces to sets

\begin{equation*}
V \longmapsto [\underline{A}(D_V(K))]
\end{equation*}
commutes with products.

\underline{Lemma} 1.18.  \underline{Assume condition} (II') \underline{is satisfied by} $\underline{A}$ . \underline{Then}
\SGApage{I}{47}
(1)  $[\underline{A}(D_{K'}(K))]$  \underline{carries a unique vector space structure, such that}  $[\underline{A}(D_{V}(K))] \simeq V \otimes_{K'} [\underline{A}(D_{K'}(K))]$  \underline{as a functor in}  $V$ .

(2)   \underline{There is a linear map}  $\gamma : \operatorname{Hom}(\Omega,K') \rightarrow [\underline{A}(D_{K'}(K))]$  \underline{such that, for any}  $D \in \operatorname{Hom}(\Omega, K')$  ,

\SGAvar{\begin{equation*}
[\underline{A}((k, D))] : [\underline{A}(D_{K'}(K))] \rightarrow [\underline{A}(D_{K'}(K))] \quad \text{is} \quad x \longmapsto x + \gamma(D) \quad .
\end{equation*}
(3)   \underline{For any}  $(u, D) : D_{V}(K) \longrightarrow D_{W}(K)$  , \underline{the map}

}{\begin{equation*}
[\underline{A}((1, D))] : [\underline{A}(D_{K'}(K))] \rightarrow [\underline{A}(D_{K'}(K))] \quad \text{is} \quad x \longmapsto x + \gamma(D) \quad .
\end{equation*}
(3)   \underline{For any}  $(u, D) : D_{V}(K) \longrightarrow D_{W}(K)$  , \underline{the map}

}{F7I-VI-018}\begin{equation*}
(u, D) : [\underline{A}(D_{V}(K))] \rightarrow [\underline{A}(D_{W}(K))] \quad \text{i.e.}
\end{equation*}
\begin{equation*}
(u, D) : V \otimes [\underline{A}(D_{K'}(K))] \rightarrow W \otimes [\underline{A}(D_{K'}(K))] \quad \underline{\text{is}}
\end{equation*}
\begin{equation*}
x \longmapsto u(x) + \gamma(D) \quad .
\end{equation*}
Let us introduce the notations

\begin{equation*}
\begin{aligned}
F(V) \quad &= \quad [\underline{A}(D_{V}(K))] \\
A \quad &= \quad F(K') \\
F(u,D) : F(V) \rightarrow F(W) \quad &= \ \text{the map induced by} \ \ (u,D) : D_{V}(K) \rightarrow D_{W}(K) \\
F(D) \quad : F(V) \rightarrow F(V) \quad &: \quad F(D) = F(1, D) \quad .
\end{aligned}
\end{equation*}
By (II') and the fact that  $K'$  is a  $K'$-vector space object in the category of $K'$-vector spaces,  $F(K')$  is a  $K'$-vector space in (Ens), hence a  $K'$-vector space, and (1) is clear. Let us now identify  $F(V)$  and  $V \otimes A$  .

We deduce from the identity

\begin{equation*}
(u, D) = (1, D) \circ (u, 0) \quad \text{that}
\end{equation*}
\begin{equation*}
F(u, D) = F(D) \circ (u \otimes 1_{A}) \quad .
\end{equation*}
The functoriality of  $F$  amounts to

\begin{equation*}
\begin{aligned}
F(0) &= \mathrm{Id} \\
F(D) \circ F(D') &= F(D + D') \\
(u \otimes 1) \circ F(D) \ &= F(u\,D) \circ (u \otimes 1_{A})
\end{aligned}
\end{equation*}
For any  $D : \Omega \rightarrow V$  and any  $n \in N$  , if  $i$  and  $p$  are the maps

\SGAvar{\begin{equation*}
i \ ; \ V \lhook\joinrel\longrightarrow V \oplus K'^{\,n} \qquad\qquad p : V \oplus K'^{\,n} \longrightarrow K^{n} \quad , 
\end{equation*}

}{\begin{equation*}
i \ ; \ V \lhook\joinrel\longrightarrow V \oplus K'^{\,n} \qquad\qquad p : V \oplus K'^{\,n} \longrightarrow K'^{\,n} \quad , 
\end{equation*}

}{F7I-VI-019}\SGApage{I}{48}

\SGAvar{ then  $p i D = 0$  and the following diagramm is commutative

\begin{center}
\begin{tikzcd}[column sep=huge, row sep=large]
V \otimes A \arrow[r] \arrow[d, "F(D)"] \& (V \oplus K'^{\,n}) \otimes A \arrow[r] \arrow[d, "{F(i,D)}"] \& K'^{\,n}\!\otimes A \arrow[d, equal] \\
V \otimes A \arrow[r] \& (V \oplus K'^{\,n}) \otimes A \arrow[r] \& K'^{\,n}\!\otimes A
\end{tikzcd}
\end{center}













Let  $x : K^{n} \rightarrow V$  be a map,  $x_{i} = x(e_{i})$  , and let  $a_{i}$  be a family of elements of  $A$  . We have  $(1_{V} + x) \circ i \circ D = D$  . Hence, if

\begin{equation*}
F(i\,D)\ (\Sigma\ e_{i}\ a_{i}) = \Sigma\ e_{i}\ a_{i} + a \qquad (a \in V \otimes H) \quad , \ \text{then}
\end{equation*}
\begin{equation*}
F(D)\ (\Sigma\ x_{i}\ a_{i}) \ = \Sigma\ x_{i}\ a_{i} + a
\end{equation*}
and  $a$  does not depend on the  $x_{i}$  . Hence,

\begin{equation*}
\begin{aligned}
F(D)\ (y) \ &= y + G(D) &&\qquad , \ \text{with} \\
G(D + D') &= G(D) + G(D') & \\
u\ G(D) \quad &= F(u\,D) &&\qquad (\text{for any} \ \ u : V \rightarrow W) \quad .
\end{aligned}
\end{equation*}
One easily checks that any such  $G$  is defined, for some  $\gamma \in \Omega \otimes A$  , by the formula

}{ then  $p i D = 0$  and the following diagram is commutative

\begin{center}
\begin{tikzcd}[column sep=huge, row sep=large]
V \otimes A \arrow[r] \arrow[d, "F(D)"] \& (V \oplus K'^{\,n}) \otimes A \arrow[r] \arrow[d, "{F(i\,D)}"] \& K'^{\,n}\!\otimes A \arrow[d, equal] \\
V \otimes A \arrow[r] \& (V \oplus K'^{\,n}) \otimes A \arrow[r] \& K'^{\,n}\!\otimes A
\end{tikzcd}
\end{center}













Let  $x : K'^{\,n} \rightarrow V$  be a map,  $x_{i} = x(e_{i})$  , and let  $a_{i}$  be a family of elements of  $A$  . We have  $(1_{V} + x) \circ i \circ D = D$  . Hence, if

\begin{equation*}
F(i\,D)\ (\Sigma\ e_{i}\ a_{i}) = \Sigma\ e_{i}\ a_{i} + a \qquad (a \in V \otimes A) \quad , \ \text{then}
\end{equation*}
\begin{equation*}
F(D)\ (\Sigma\ x_{i}\ a_{i}) \ = \Sigma\ x_{i}\ a_{i} + a
\end{equation*}
and  $a$  does not depend on the  $x_{i}$  . Hence,

\begin{equation*}
\begin{aligned}
F(D)\ (y) \ &= y + G(D) &&\qquad , \ \text{with} \\
G(D + D') &= G(D) + G(D') & \\
u\ G(D) \quad &= G(u\,D) &&\qquad (\text{for any} \ \ u : V \rightarrow W) \quad .
\end{aligned}
\end{equation*}
One easily checks that any such  $G$  is defined, for some  $\gamma \in \Omega \otimes A$  , by the formula

}{F7I-VI-020}\begin{equation*}
G(D) = D \otimes 1_{A}\ (\gamma) \quad ,
\end{equation*}
as promised in (2) (3).

\underline{Definition} 1.19.   (1) \underline{An object}  $\xi$  \underline{in}  $\widehat{\underline{A}}$  \underline{is called}   \underline{versal formal object if} \underline{every surjective map}  $\alpha : a' \longrightarrow a$  \underline{in}  $\underline{A}$  \underline{induces a surjective map}

\SGAvar{\begin{equation*}
\operatorname{Hom}_{\widehat{\underline{A}}}(\xi, a') \longrightarrow \operatorname{Hom}_{\widehat{\underline{A}}}\ (\xi, a) \quad .
\end{equation*}
(2)   \underline{A formal versal object}  $\xi$  \underline{in}  $\widehat{A}$  , \underline{with image}  $(R, s)$  \underline{in}  $\widehat{C}_{\Lambda,K}$  , \underline{is} \underline{called  minimal  if for any two maps}  $\varphi_{1}$ , $\varphi_{2} : R \rightrightarrows D_{K'}(K)$  , \underline{if}  $\varphi_{1}(\xi)$  \underline{is iso-} \underline{morphic to}  $\varphi_{2}(\xi)$  , \underline{then there exists a derivation}  $D$  \underline{of}  $K$  \underline{into}  $K'$  \underline{such that}  $(1, D)\ (\varphi_{1}) = \varphi_{2}$  , i.e.  $\varphi_{2}\text{-}\varphi_{2} = D\,s$  .

}{\begin{equation*}
\operatorname{Hom}_{\widehat{\underline{A}}}(\xi, a') \longrightarrow \operatorname{Hom}_{\widehat{\underline{A}}}\ (\xi, a) \quad .
\end{equation*}
(2)   \underline{A formal versal object}  $\xi$  \underline{in}  $\widehat{A}$  , \underline{with image}  $(R, s)$  \underline{in}  $\widehat{C}_{\Lambda,K}$  , \underline{is} \underline{called  minimal  if for any two maps}  $\varphi_{1}$ , $\varphi_{2} : R \rightrightarrows D_{K'}(K)$  , \underline{if}  $\varphi_{1}(\xi)$  \underline{is iso-} \underline{morphic to}  $\varphi_{2}(\xi)$  , \underline{then there exists a derivation}  $D$  \underline{of}  $K$  \underline{into}  $K'$  \underline{such that}  $(1, D)\ (\varphi_{1}) = \varphi_{2}$  , i.e.  $\varphi_{2}\text{-}\varphi_{1} = D\,s$  .

}{F7I-VI-021}\underline{Theorem} 1.20.   \underline{If}  $\underline{A}$  \underline{is semi-homogeneous and if}  ,

\begin{equation*}
\dim_{K'}\ ([A(D_{K'}(K))]) < \infty \quad ,
\end{equation*}
\underline{then}  $\underline{A}$  \underline{admits a minimally versal formal object, and any two of them are}
\SGApage{I}{49}
\underline{isomorphic}.

With the notations of 1.18, let  $H^*$  be a subspace of  $[\underline{A}(D_{K'}(K))]$ , supplementary to the image of  $\gamma$ , and let  $H$  be the dual of  $H^*$ . We define $R_o = D_H(K)$ . Let  $\xi_o$  be an object in  $\underline{A}(R_o)$ , whose image in $[\underline{A}(R_o)] \simeq H \otimes [\underline{A}(D_{K'}(K))] \simeq \operatorname{Hom}(H^*,[\underline{A}(D_{K'}(K))])$  is the inclusion of  $H^*$  in $[\underline{A}(D_{K'}(K))]$ . Then

(1)  The maps  $\operatorname{Hom}_{C_{\wedge,K}}(D_H(K),\ D_V(K)) \longrightarrow [\underline{A}(D_V(K))]$  are surjective .

(2)  The minimality condition 1.19(2)  is valid for  $(R_o,\ \xi_o)$  .

Choose now some  $S \in \hat{C}_{\wedge,K}$ , formally smooth over  $\wedge$ , and an epimorphism $t : S \longrightarrow R_o$  such that

\begin{equation*}
\Omega_{R_o/\wedge} \otimes K' \overset{\sim}{\longrightarrow} \Omega_{S/\wedge} \otimes K' \quad .
\end{equation*}
If  $M$  is the maximal ideal of  $S$ , we inductively define, for  $n \geq 1$ , $R_n = S/I_n$  and  $\xi_n \in \operatorname{Ob}(\underline{A}(R_n))$  by the properties that

(a)  $I_{n+1}$  is the intersection of the ideals  $I : I_n M \subset I \subset I_n$ , such that  $\xi_n$ extends to an object  $\xi$  in  $\underline{A}(S/I)$  .

(b)  $\xi_{n+1}$  is an extension of  $\xi_n$  .

If  $R = \varprojlim R_n$  and  $\xi = \varprojlim \xi_n$ , then one checks that  $(R,\ \xi)$  is a minimally-versal formal object, and the only one up to isomorphism.


2. \underline{Prorepresentable cofibered groupoids.}

\SGAvar{Let  $\underline{A}$  be a cofibered groupoid over  $\underline{C}_\wedge$ , admitting a minimally-versal formal object  $\xi$ . In general we can not recover  $\underline{A}$  (up to  $\underline{C}_\wedge$-equivalence) out of  $\xi$  alone..One has to know when two maps  $R \overset{f}{\underset{g}{\rightrightarrows}} S$  yield a  $\underline{C}_\wedge$-isomorphism $\xi(f) \longrightarrow \xi(g)$ , where  $R = P(\xi)$ . As an application of 1.11  and 1.12 we now consider the question of  "recovering  $\underline{A}$ "  up to  $\underline{C}_\wedge$-equivalence. Indeed, we
\SGApage{I}{50}
shall have a quite satisfactory answer applicable to a wide class of cofibered groupoids.

}{Let  $\underline{A}$  be a cofibered groupoid over  $\underline{C}_\wedge$ , admitting a minimally-versal formal object  $\xi$ . In general we can not recover  $\underline{A}$  (up to  $\underline{C}_\wedge$-equivalence) out of  $\xi$  alone. One has to know when two maps  $R \overset{f}{\underset{g}{\rightrightarrows}} S$  yield a  $\underline{C}_\wedge$-isomorphism $\xi(f) \longrightarrow \xi(g)$ , where  $R = P(\xi)$ . As an application of 1.11  and 1.12 we now consider the question of  "recovering  $\underline{A}$ "  up to  $\underline{C}_\wedge$-equivalence. Indeed, we
\SGApage{I}{50}
shall have a quite satisfactory answer applicable to a wide class of cofibered groupoids.

}{F7I-VI-022}A number of cofibered groupoids which usually occur in practice enjoy a slightly stronger property than those of semi-homogeneous ones, as it will be defined explicitly in 2.5. To investigate them it is convenient to rephrase what we have done so far in terms of  "smooth functors" . All the notations and terminologies in the previous section remain in force throughout in this section.

\underline{Definition} 2.1.  \underline{Let}  $\varphi$ : $\underline{A} \longrightarrow \underline{B}$  \underline{be a functor of cofibered groupoids over}  $\underline{C}_\wedge$ . \underline{We say that}  $\varphi$  \underline{is smooth if every surjective map}  $\beta : b' \longrightarrow \varphi(a)$  \underline{in}  $\underline{B}$  \underline{can be completed to a commutative diagram}

\begin{equation}\tag{*}
\begin{tikzcd}[row sep=large, column sep=huge]
\varphi(a') \arrow[r] \arrow[dr, "\varphi(\alpha)"'] \& b' \arrow[d, "\beta"] \\
 \& \varphi(a)
\end{tikzcd}
\end{equation}










\underline{for some}  $\alpha : a' \longrightarrow a$  \underline{in}  $\underline{A}$ . \underline{We say that}  $\varphi$  \underline{is minimally smooth if it is smooth and}  $\varphi(k[\epsilon]) : [\underline{A}(k[\epsilon])] \longrightarrow [\underline{B}(k[\epsilon])]$  \underline{is bijective}.

\underline{Remark} 2.2(a)  Let  $\varphi$ : $\underline{A} \longrightarrow \underline{B}$  be a functor of cofibered groupoids over  $\underline{C}_\wedge$ . We note that  " $\varphi$  is smooth" $\Leftrightarrow$ "every surjective map  $\beta : b' \longrightarrow \varphi(a)$  in  $\underline{B}$ can be completed into a commutative diagram (*) such that  $\varphi(a') \longrightarrow b'$  is a $\underline{C}_\wedge$-isomorphism"  $\Leftrightarrow$  "every small prolongation  $\beta : b' \longrightarrow \varphi(a)$  in  $\underline{B}$  can be completed into a commutative diagram (*) ".

\underline{Remark} 2.2(b)  For any cofibered groupoid  $\underline{A}$  over  $\underline{C}_\wedge$ , the natural functor $\underline{A} \longrightarrow [\underline{A}]$  is minimally smooth.

\underline{Remark} 2.2(c)  For each object  $R$  in  $\hat{\underline{C}}_\wedge$ , we set  $h_R : \underline{C}_\wedge \longrightarrow$ (Ens)  to be the functor defined by  $h_R(S) = \operatorname{Hom}_{\wedge\text{-alg}}(R,\ S)$  for all  $S$  in  $\underline{C}_\wedge$ . It associates a cofibered discrete groupoid  $\underline{h}_R$  over  $\underline{C}_\wedge$ , and a map  $R \longrightarrow S$  in  $\underline{C}_\wedge$
\SGApage{I}{51}
induces a functor  $\underline{h}_S \longrightarrow \underline{h}_R$  of cofibered groupoids. The functor  $\underline{h}_S \longrightarrow \underline{h}_R$  induced by a map  $R \longrightarrow S$  in  $\underline{C}_\wedge$  is smooth if and only if  $S$  is a formal power-series ring over  $R$ .

\SGAvar{\underline{Remark} 2.2(d)  Given an object  $R$  in  $\widehat{\underline{C}}_\wedge$ , let us denote by  $(R,\ \underline{C}_\wedge)$  the category of arrows in  $\widehat{\underline{C}}_\wedge$  with the source  $R$  and a target in  $\underline{C}_\wedge$ . Let  $\underline{A}$  be a cofibered groupoid over  $\underline{C}_\wedge$ . Given  $\xi \in \mathrm{ob}\ \widehat{\underline{A}}$  with  $P(\xi) = R$ , let  $\underline{A}' \longrightarrow (R,\ \underline{C}_\wedge)$  be the pull-back of  $\underline{A} \longrightarrow \underline{C}_\wedge$  under the canonical functor  $(R,\ \underline{C}_\wedge) \longrightarrow \underline{C}_\wedge$  and choose a section  $\xi^{\#}$  such that  $\xi^{\#}(1_R) = \xi$ . (This simply means that we choose, for each map  $R \xrightarrow{\ f\ } S$ , a $f$-map  $\xi \longrightarrow \xi(f)$ ). We then obtain a functor  $\xi^{\#} : \underline{h}_R \longrightarrow \underline{A}$  over  $\underline{C}_\wedge$  by sitting  $(R \xrightarrow{\ f\ } S) \longmapsto \xi(f)$ . Of course this functor does not depend on the choice of  $\xi^{\#}$  up to natural equivalence of functors. Now  " $\xi^{\#} : \underline{h}_R \longrightarrow \underline{A}$  is (minimally) smooth"  means that every surjective map  $a' \xrightarrow{\ \alpha\ } \xi(f)$  in  $\underline{A}$  can be completed into a commutative diagram

}{\underline{Remark} 2.2(d)  Given an object  $R$  in  $\widehat{\underline{C}}_\wedge$ , let us denote by  $(R,\ \underline{C}_\wedge)$  the category of arrows in  $\widehat{\underline{C}}_\wedge$  with the source  $R$  and a target in  $\underline{C}_\wedge$ . Let  $\underline{A}$  be a cofibered groupoid over  $\underline{C}_\wedge$ . Given  $\xi \in \mathrm{ob}\ \widehat{\underline{A}}$  with  $P(\xi) = R$ , let  $\underline{A}' \longrightarrow (R,\ \underline{C}_\wedge)$  be the pull-back of  $\underline{A} \longrightarrow \underline{C}_\wedge$  under the canonical functor  $(R,\ \underline{C}_\wedge) \longrightarrow \underline{C}_\wedge$  and choose a section  $\xi^{\#}$  such that  $\xi^{\#}(1_R) = \xi$ . (This simply means that we choose, for each map  $R \xrightarrow{\ f\ } S$ , a $f$-map  $\xi \longrightarrow \xi(f)$ ). We then obtain a functor  $\xi^{\#} : \underline{h}_R \longrightarrow \underline{A}$  over  $\underline{C}_\wedge$  by setting  $(R \xrightarrow{\ f\ } S) \longmapsto \xi(f)$ . Of course this functor does not depend on the choice of  $\xi^{\#}$  up to natural equivalence of functors. Now  " $\xi^{\#} : \underline{h}_R \longrightarrow \underline{A}$  is (minimally) smooth"  means that every surjective map  $a' \xrightarrow{\ \alpha\ } \xi(f)$  in  $\underline{A}$  can be completed into a commutative diagram

}{F7I-VI-023}\begin{center}
\begin{tikzcd}[row sep=large, column sep=huge]
\xi(f') \arrow[r] \arrow[dr, "\xi(\ \ )"'] \& a' \arrow[d, "\alpha"] \\
 \& \xi(f)
\end{tikzcd}
\end{center}











(and  $\underline{h}_R(k[\epsilon]) \longrightarrow \underline{A}(k[\epsilon])$   is bijective), which means precisely that  $\xi$  is a (minimally) versal formal object of  $\underline{A}$  over  $\underline{C}_\wedge$ . Thus we conclude that the existence of a (minimally) versal formal object of  $\underline{A}$  is equivalent to the existence of a (minimally) smooth functor  $\underline{h}_R \longrightarrow \underline{A}$  over  $\underline{C}_\wedge$  for some  $R \in \mathrm{ob}(\widehat{\underline{C}}_\wedge)$ .

The following general properties on smoothness are trivial to verify.

\SGAvar{\underline{Proposition} 2.3.  \underline{Let}  $\underline{A} \xrightarrow{\ \varphi\ } \underline{B} \xrightarrow{\ \varphi\ } \underline{C}$  \underline{be functors of cofibered groupoids over}  $\underline{C}_\wedge$ .

}{\underline{Proposition} 2.3.  \underline{Let}  $\underline{A} \xrightarrow{\ \varphi\ } \underline{B} \xrightarrow{\ \psi\ } \underline{C}$  \underline{be functors of cofibered groupoids over}  $\underline{C}_\wedge$ .

}{F7I-VI-024}(a)  \underline{If}  $\varphi, \psi$  \underline{are both smooth, then so is}  $\psi.\varphi$ .

(b)  \underline{If}  $\psi.\varphi$  \underline{is smooth and}  $\varphi$  \underline{is surjective on isomorphism classes of objects,}
\SGApage{I}{52}
\underline{then}  $\psi$  \underline{is smooth.}

\SGAvar{\underline{Remark} 2.4.  Let  $\underline{A}$  be a cofibered groupoids over  $\underline{C}_\wedge$ , admitting a versal formal object  $\xi$ . This means, by 2.2(d), that  $\tilde{\xi} : \underline{h}_R \to \underline{A}$  is smooth, where  $R = P(\xi)$ . Therefore we have,  " $\underline{A}$  is (formally) smooth over  $\underline{C}_\wedge$  (c.f. 1.13) i.e.,  $R$  is a formal power-series algebra over  $\wedge$ "  $\Leftrightarrow$  "the composite functor  $\underline{h}_R \xrightarrow{\ \tilde{\xi}\ } \underline{A} \xrightarrow{\ P\ } \underline{C}_\wedge$   $(= \underline{h}_\wedge)$   is smooth"  $\Leftrightarrow$  " $\underline{A} \xrightarrow{\ P\ } \underline{C}_\wedge$  is a smooth functor"  by 2.3(b).

}{\underline{Remark} 2.4.  Let  $\underline{A}$  be a cofibered groupoid over  $\underline{C}_\wedge$ , admitting a versal formal object  $\xi$ . This means, by 2.2(d), that  $\tilde{\xi} : \underline{h}_R \to \underline{A}$  is smooth, where  $R = P(\xi)$ . Therefore we have,  " $\underline{A}$  is (formally) smooth over  $\underline{C}_\wedge$  (c.f. 1.14) i.e.,  $R$  is a formal power-series algebra over  $\wedge$ "  $\Leftrightarrow$  "the composite functor  $\underline{h}_R \xrightarrow{\ \tilde{\xi}\ } \underline{A} \xrightarrow{\ P\ } \underline{C}_\wedge$   $(= \underline{h}_\wedge)$   is smooth"  $\Leftrightarrow$  " $\underline{A} \xrightarrow{\ P\ } \underline{C}_\wedge$  is a smooth functor"  by 2.3(b).

}{F7I-VI-025}We now consider a stronger property than being semi-homogeneous.

\underline{Definition} 2.5.  \underline{Let}  $\underline{A}$  \underline{be a cofibered groupoid over}  $\underline{C}_\wedge$ . \underline{It is called homogeneous if every diagram}

\begin{center}
\begin{tikzcd}[row sep=large, column sep=huge]
 \& a' \arrow[d] \\
b' \arrow[r] \& a
\end{tikzcd}
\end{center}

\underline{where}  $a' \to a$  \underline{is surjective admits a fibre-product}  $a' \underset{a}{\times} b'$  \underline{in}  $\underline{A}$  \underline{such that}

\begin{equation*}
P(a' \times b') = P(a') \underset{P(a)}{\times} P(b') \ .
\end{equation*}
\underline{Remark} 2.6(a)  Let  $\underline{A}$  be a cofibered groupoid over  $\underline{C}_\wedge$ , and consider a diagram

(*) \qquad
\begin{tikzcd}[row sep=small, column sep=small]
 \& a' \arrow[dl] \arrow[dr] \& \\
a_1 \arrow[dr] \& \& a_2 \arrow[dl] \\
 \& a \&
\end{tikzcd}
\qquad \text{in } $\underline{A}$ \text{ above the diagram} \qquad
\begin{tikzcd}[row sep=small, column sep=small]
 \& R_1 \underset{R}{\times} R_2 \arrow[dl, "\pi_1"'] \arrow[dr, "\pi_2"] \& \\
R_1 \arrow[dr] \& \& R_2 \arrow[dl] \\
 \& R \&
\end{tikzcd}
















in  $\underline{C}_\wedge$  where  $\pi_i$  are the projections. If a fibre-product  $a_1 \underset{a}{\times} a_2$  exists in  $\underline{A}$ , we obtain a unique map  $\varphi : a' \to a_1 \underset{a}{\times} a_2$  compatible with the projections. If  $P(a_1 \underset{a}{\times} a_2) = R_1 \underset{R}{\times} R_2$ , then  $P(\varphi)$  is the identity map on  $R_1 \underset{R}{\times} R_2$  and therefore  $\varphi$  is a  $\underline{C}_\wedge$-isomorphism i.e.,  $a' = a_1 \underset{a}{\times} a_2$ . In other words, if  $\underline{A}$  is a homogeneous cofibered groupoid over  $\underline{C}_\wedge$ , every  diagram (*) defines a fibre-product  $a_1 \underset{a}{\times} a_2$  provided  $R_1 \to R$  is surjective in  $\underline{C}_\wedge$ . This fact entails immediately
\SGApage{I}{53}
that a homogeneous cofibered groupoid is, a fortiori, semi-homogeneous.


\underline{Remark} 2.6(b)   Consider a diagram

\begin{center}
\begin{tikzcd}[row sep=large, column sep=large]
X \arrow[dr, "\varphi"'] \& \& Y \arrow[dl, "\psi"] \\
 \& Z \&
\end{tikzcd}
\end{center}

of categories and functors. The ordinary fibre-product  $X \times_Z Y$  is a category in which an object is a pair  $(x,y)$  such that  $\varphi(x) = \psi(y)$ , and a map $(x,y) \to (x', y')$  consists of a pair  $\alpha : x \to x'$ , $\beta : y \to y'$  such that $\varphi(\alpha) = \psi(\beta)$ . The fibre-product  $X \times^2_Z Y$  in the sense of  2-categories is a category in which an object is a triple  $(x,y; \gamma)$ , where  $\varphi(x) \xrightarrow{\gamma} \psi(y)$  is an isomorphism, and a map  $(x,y; \gamma) \to (x',y'; \gamma')$  consists of a pair  $x \xrightarrow{\alpha} x'$   $y \xrightarrow{\beta} y'$  with a commutative diagram

\begin{center}
\begin{tikzcd}[row sep=huge, column sep=huge]
\varphi(x) \arrow[r, "\varphi(\alpha)"] \arrow[d, "\gamma"'] \& \varphi(x') \arrow[d, "\gamma'"] \\
\psi(y) \arrow[r, "\psi(\beta)"'] \& \psi(y')
\end{tikzcd}
\end{center}









One can rephrase the notions of (semi)homogeneous cofibered groupoid or smooth functor in terms of fibre-product of categories (in the sense of 2-category). Now let  $P : \underline{A} \to \underline{C}_\wedge$  be a cofibered category, and consider an object $R_1 \times_R R_2$  in  $\underline{C}_\wedge$ . If we choose a direct image functor  $(\pi_i)_*: \underline{A}(R_1 \times_R R_2) \to \underline{A}(R_i)$ , where  $\pi_i : R_1 \times_R R_2 \to R_i$  is the projection for  $i = 1,2$,  we obtain a functor $\underline{A}(R_1 \times_R R_2) \to \underline{A}(R_1) \times^2_{\underline{A}(R)} \underline{A}(R_2)$ . It is clear from the definition of cofibered category that this functor, up to canonical natural equivalence, does not depend on the choice of direct image functors. It is now clear from the definition that a cofibered groupoid  $\underline{A}$  over  $\underline{C}_\wedge$  is homogeneous over  $\underline{C}_\wedge$  if and only if $\underline{A}(R_1 \times_R R_2) \to \underline{A}(R_1) \times^2_{\underline{A}(R)} \underline{A}(R_2)$  is an equivalence of categories whenever  $R_1 \to R$  is surjective in  $\underline{C}_\wedge$ , and  $\underline{A}$  is semi-homogeneous if and only if

\SGAvar{\begin{equation*}
[\underline{A}(R \times_R R_2)] \longrightarrow [\underline{A}(R_1) \times^2_{\underline{A}(R)} \underline{A}(R_2)] 
\end{equation*}

}{\begin{equation*}
[\underline{A}(R_1 \times_R R_2)] \longrightarrow [\underline{A}(R_1) \times^2_{\underline{A}(R)} \underline{A}(R_2)] 
\end{equation*}

}{F7I-VI-026}\SGApage{I}{54}

 is surjective whenever  $R_1 \longrightarrow R$  is surjective in  $\underline{C}_\wedge$ , and is bijective whenever  $R = k$  and  $R_1 = k[\epsilon]$ . Now let  $\varphi : \underline{A} \to \underline{B}$  be a functor of cofibered groupoids over  $\underline{C}_\wedge$ . For each map  $R' \to R$  in  $\underline{C}$  , $\varphi$  induces a functor $\underline{A}(R') \to \underline{A}(R) \times^2_{\underline{B}(R)} \underline{B}(R')$  uniquely determined up to a (canonical) natural equivalence. It is again trivial to see that  $\varphi$  is smooth if and only if the functor $A(R') \to \underline{A}(R) \times^2_{\underline{B}(R)} \underline{B}(R')$  is surjective on isomorphism classes of objects, whenever $R' \to R$  is surjective.


\underline{Remark} 2.6(c). If  $\underline{A}$  is a homogeneous cofibered groupoid over  $\underline{C}_\wedge$ , the functor $A(k[V \times W]) \to \underline{A}(k[V])\times \underline{A}(k[W])$         is an equivalence of categories for any finite-dimensional vector spaces  $V, W$  over  $k$  (We recall our convention $\underline{A}(k) = \{e\}$  ). In particular, we have

\begin{equation*}
\mathrm{Aut}_{\underline{A}(k[V\times W])}(1_{V\times W}) \xrightarrow{\ \sim\ } \mathrm{Aut}_{\underline{A}(k[V])}(1_V)\times\mathrm{Aut}_{\underline{A}(k[W])}(1_W)
\end{equation*}
where  $1_V$  is the direct image of  $e$  under the unique map  $k \to k[V]$  in  $\underline{C}_\wedge$  . It follows that if  $\underline{A}$  is homogeneous over  $\underline{C}_\wedge$ , $\mathrm{Aut}(1_\epsilon)$  as well as  $[\underline{A}(k[\epsilon])]$ are canonically provided with vector space structures over  $k$  .


\underline{Remark} 2.6(d). For any object  $R$  in  $\hat{\underline{C}}_\wedge$ , the cofibered (discrete) groupoid $\underline{h}_R$  over  $\underline{C}_\wedge$  is homogeneous.

If  $\underline{A}$  is semi-homogeneous, then so is  $[\underline{A}]$ . However, the homogeneity of  $\underline{A}$  does not imply in general the homogeneity of  $[\underline{A}]$ . Indeed we have


\underline{Proposition} 2.7.  \underline{Let}  $\underline{A}$  \underline{be a homogeneous cofibered groupoid over}  $\underline{C}_\wedge$ . \underline{The following statements are equivalent.}

(i)  $[\underline{A}]$  \underline{is homogeneous i.e.}  $[\underline{A}(R' \times_R S)] \longrightarrow [\underline{A}(R')] \times_{[\underline{A}(R)]} [\underline{A}(S)]$  \underline{is bijective} \underline{for every small extension}  $R' \to R$  .

(ii) \underline{For every small prolongation}  $a' \to a$  \underline{in}  $\underline{A}$ , \underline{the map}

\begin{equation*}
\mathrm{Aut}_{\underline{A}(R')}(a') \longrightarrow \mathrm{Aut}_{\underline{A}(R)}(a) 
\end{equation*}

\SGApage{I}{55}

 \underline{is surjective, where}  $R' = P(a')$  ,  $R = P(a)$  .

\underline{Proof}.  (i) $\Rightarrow$ (ii)  Let  $R' \xrightarrow{\ f\ } R$  be a small extension in  $\underline{C}_{\wedge}$  and let  $a' \to a$  be a f-map in  $\underline{A}$ . Given  $\tau \in \mathrm{Aut}_{\underline{A}(R)}(a)$  , we consider  $a' \times_{a} a'$  and  $a' \times_{a} a'_{\tau}$  where  $a'_{\tau} \longrightarrow a$  denotes the composite map  $a' \longrightarrow a \xrightarrow{\ \tau\ } a$  . Since  $[A(R' \times_{R} R')] \longrightarrow [\underline{A}(R')] \times_{[\underline{A}(R)]} [\underline{A}(R')]$  is bijective, we conclude the existence of an isomorphism  $a' \times_{a} a' \xrightarrow{\ u\ } a' \times_{a} a'_{\tau}$  , and  $u$  induces, via the projection on the second factor, an isomorphism  $a' \xrightarrow{\ \tau'\ } a'$  which is a lifting of  $\tau$  .

(ii) $\Rightarrow$ (i)  Let  $a''_{i}$  ($i=1,2$)  be objects in  $\underline{A}(R' \times_{R} S)$  which have the same image under  $[\underline{A}(R' \times_{R} S)] \to [\underline{A}(R')] \times_{[\underline{A}(R)]} [\underline{A}(S)]$  . We then obtain a diagram of solid arrows

\begin{center}
\begin{tikzcd}[column sep=small, row sep=small]
\&\& a''_{2} \arrow[ddll] \arrow[ddrr] \&\& \\
\&\& a''_{1} \arrow[u, dashed] \arrow[dl] \arrow[dr] \&\& \\
a'_{2} \arrow[ddrr] \& a'_{1} \arrow[l, "\sim"] \arrow[dr] \&\& b_{1} \arrow[r, "\tau"] \arrow[dl] \& b_{2} \arrow[ddll] \\
\&\& a_{1} \arrow[d, "\bar{\tau}"] \&\& \\
\&\& a_{2} \&\&
\end{tikzcd}
\qquad \text{above} \qquad
\begin{tikzcd}[column sep=small, row sep=small]
\& R' \times_{R} S \arrow[dl] \arrow[dr] \& \\
R' \arrow[dr] \&\& S \arrow[dl] \\
\& R \&
\end{tikzcd}
\end{center}



































where  $\bar{\tau}$ ; $a_{1} \to a_{2}$  is the induced isomorphism from  $\tau$  so that the square box on the lower right corner is commutative. If the condition (ii) is verified, one can always replace the given isomorphism  $a'_{1} \overset{\sim}{\to} a'_{2}$  so that the square box on the left corner is also commutative. Since  $a''_{i} = a'_{i} \times_{a_{i}} b_{i}$  for  $i = 1, 2$  by  2.6(a), we obtain the isomorphism  $a''_{1} \overset{\sim}{\to} a''_{2}$  . This proves that  $[\underline{A}(R' \times_{R} S)] \longrightarrow [\underline{A}(R')] \times_{[\underline{A}(R)]} [\underline{A}(S)]$  is injective, and hence is bijective since  $\underline{A}$  is homogeneous by hypothesis.

Now let  $\underline{A}$  be a homogeneous cofibered groupoid over  $\underline{C}_{\Lambda}$  , and let  $R' \xrightarrow{\ f\ } R$  be a small extension in  $\underline{C}_{\Lambda}$  . We have a canonical map
\SGApage{I}{56}
$\tilde{f} : R' \times_{R} R' \longrightarrow k[\mathrm{Ker}\ f]$  given by  $(r'_{1}\ ,\ r'_{2}) \longrightarrow r'_{1}(0) + (r'_{1} - r'_{2})$  , and  $1 \times \tilde{f} : R' \times_{R} R' \longrightarrow R' \times_{k} k[\mathrm{Ker}\ f]$  is an isomorphism. Given any two f-maps  $a_{i} \to a$  ($i=1,2$) in  $\underline{A}$  ,  $a_{1} \times_{a} a_{2}$  is an object over  $R' \times_{R} R'$  and in turn we may consider  $\tau(a_{1},\ a_{2}) = (\tilde{f})_{*}(a_{1} \times_{a} a_{2})$  which is an object in  $\underline{A}(k[\mathrm{Ker}\ f])$  . Since  $1 \times \tilde{f} : R' \times_{R} R' \longrightarrow R' \times_{k} k[\mathrm{Ker}\ f]$  is an isomorphism, we have the isomorphism  $a_{1} \times_{a} a_{2} \simeq a_{1} \times \tau(a_{1},\ a_{2})$  compatible with the projections on the first factor. The statement 1.5(2) takes a more definite form for a homogeneous groupoid.

\underline{Lemma} 2.8.  \underline{Let}  $\underline{A}$  \underline{be a homogeneous cofibered groupoid over}  $\underline{C}_{\Lambda}$  , \underline{and}  $a_{i} \xrightarrow{\ \alpha_{i}\ } a$  ($i = 1, 2$)  f-\underline{maps in}  $\underline{A}$  . \underline{Then}

(1)  \underline{There exists}  $\rho$  : $a_{1} \to a_{2}$  \underline{in}  $\underline{A}$  \underline{with}  $\alpha_{2} \cdot \rho = \alpha_{1}$  \underline{if and only if} \underline{there exists an arrow}  $a_{1} \to \tau(a_{1},\ a_{2})$  \underline{in}  $\underline{A}$  .

(2)  \underline{There exists}  $\rho : a_{1} \to a_{2}$  \underline{with}  $\alpha_{2} \cdot \rho = \alpha_{1}$  \underline{such that}  $P(\rho) = \mathrm{id}_{R'}$ \underline{if and only if}  $\tau(a_{1},\ a_{2}) = 0$  \underline{in}  $[\underline{A}(k[\mathrm{Ker}\ f])]$  .

\SGAvar{\underline{Proof}:  We prove both at the same time. There exists  $\rho : a_{1} \to a_{2}$  in  $\underline{A}$  with  $\alpha_{2} \cdot \rho = \alpha_{1}$  (such that  $P(\rho) = \mathrm{id}_{R'}$ )  $\Leftrightarrow$  there exists  $h : a_{1} \to a_{1} \times_{a} a_{2}$  in  $\underline{A}$  with  $\mathrm{pr}_{1} \circ h = \mathrm{id}_{a_{1}}$  (such that  $P(h)$ = diagonal map)  $\Leftrightarrow$ there exists  $g : a_{1} \to a_{1} \times \tau(a_{1},\ a_{2})$  in  $\underline{A}$  with  $\mathrm{pr}_{1} \circ g = \mathrm{id}_{a_{1}}$  (such that  $P(g) : R' \to R' \times_{k} k[\mathrm{Ker}\ f]$  is given by  $r' \to (r',\ r'(0))$  $\Leftrightarrow$ there exists an arrow  $\gamma : a_{1} \to \tau(a_{1},\ a_{2})$  in  $\underline{A}$  (such that  $P(\gamma) : R' \to k[\mathrm{Ker}\ f]$  is given by  $r' \to r'(0))$  . However, there exists an arrow  $\gamma : a_{1} \to \tau(a_{1},\ a_{2})$  in  $\underline{A}$  such that  $P(\gamma) : R' \to k[\mathrm{Ker}\ f]$  is factored through  $R' \to k \to k[\mathrm{Ker}\ f]$  if and only if  $\tau(a_{1},\ a_{2}) = 0$  .

}{\underline{Proof}:  We prove both at the same time. There exists  $\rho : a_{1} \to a_{2}$  in  $\underline{A}$  with  $\alpha_{2} \cdot \rho = \alpha_{1}$  (such that  $P(\rho) = \mathrm{id}_{R'}$ )  $\Leftrightarrow$  there exists  $h : a_{1} \to a_{1} \times_{a} a_{2}$  in  $\underline{A}$  with  $\mathrm{pr}_{1} \circ h = \mathrm{id}_{a_{1}}$  (such that  $P(h)$ = diagonal map)  $\Leftrightarrow$ there exists  $g : a_{1} \to a_{1} \times \tau(a_{1},\ a_{2})$  in  $\underline{A}$  with  $\mathrm{pr}_{1} \circ g = \mathrm{id}_{a_{1}}$  (such that  $P(g) : R' \to R' \times_{k} k[\mathrm{Ker}\ f]$  is given by  $r' \to (r',\ r'(0)))$  $\Leftrightarrow$ there exists an arrow  $\gamma : a_{1} \to \tau(a_{1},\ a_{2})$  in  $\underline{A}$  (such that  $P(\gamma) : R' \to k[\mathrm{Ker}\ f]$  is given by  $r' \to r'(0))$  . However, there exists an arrow  $\gamma : a_{1} \to \tau(a_{1},\ a_{2})$  in  $\underline{A}$  such that  $P(\gamma) : R' \to k[\mathrm{Ker}\ f]$  is factored through  $R' \to k \to k[\mathrm{Ker}\ f]$  if and only if  $\tau(a_{1},\ a_{2}) = 0$  .

}{F7I-VI-027}\underline{Proposition} 2.9.  (a)  \underline{Let}  $\underline{A} \xrightarrow{\ \varphi\ } \underline{B} \xrightarrow{\ \psi\ } \underline{C}$  \underline{be functors of cofibered groupoids} \underline{over}  $\underline{C}_{\Lambda}$  . \underline{Assume that}

(i)  $\underline{B}$  \underline{is semi-homogeneous and}  $\underline{C}$  \underline{is homogeneous}

(ii)  $\psi(k[\epsilon]) : [\underline{B}(k[\epsilon])] \longrightarrow [\underline{C}(k[\epsilon])]$  \underline{is injective.}

\underline{Then the composite functor}  $\psi \circ \varphi$  \underline{is smooth if and only if}  $\varphi$  \underline{and}  $\psi$  \underline{are both}
\SGApage{I}{57}
\underline{smooth.}

(b)  \underline{Let} $\varphi : \underline{A} \longrightarrow \underline{B}$ \underline{be a smooth functor of cofibered groupoids over} $\underline{C}_{\Lambda}$ . \underline{If} $\underline{A}$ \underline{is semi-homogeneous and} $\underline{B}$ \underline{is homogeneous, then} $\widehat{\varphi} : \widehat{\underline{A}} \to \widehat{\underline{B}}$ \underline{is} \underline{smooth} (i.e., \underline{for any surjective map} $b' \xrightarrow{\ \beta\ } \widehat{\varphi}(a)$ \underline{in} $\widehat{\underline{B}}$ , \underline{there exists} $\alpha : a' \to a$ \underline{in} $\widehat{\underline{A}}$ \underline{and a} $\widehat{\underline{C}}_{\Lambda}$\underline{-isomorphism} $\rho : \widehat{\varphi}(a') \xrightarrow{\ \sim\ } b'$ \underline{such that} $\beta \cdot \rho = \varphi(\alpha))$.

\SGAvar{\underline{Proof}: (a)  Assume that the composite functor $\psi \circ \varphi$ is smooth. If $\varphi$ is smooth, then so is $\psi$ by 2.3(b), and therefore it suffices to show that $\varphi$ is smooth. Let a small prolongation $b' \xrightarrow{\ \beta\ } \varphi(a)$ in $\underline{B}$ be given. Since the composite functor $\varphi \circ \psi$ is smooth, there exists an arrow $a' \xrightarrow{\ \alpha\ } a$ in $\underline{A}$ and a $\underline{C}_{\Lambda}$-isomorphism $\rho : \psi(b') \longrightarrow \psi\ \varphi(a')$ with a commutative diagram

}{\underline{Proof}: (a)  Assume that the composite functor $\psi \circ \varphi$ is smooth. If $\varphi$ is smooth, then so is $\psi$ by 2.3(b), and therefore it suffices to show that $\varphi$ is smooth. Let a small prolongation $b' \xrightarrow{\ \beta\ } \varphi(a)$ in $\underline{B}$ be given. Since the composite functor $\psi \circ \varphi$ is smooth, there exists an arrow $a' \xrightarrow{\ \alpha\ } a$ in $\underline{A}$ and a $\underline{C}_{\Lambda}$-isomorphism $\rho : \psi(b') \longrightarrow \psi\ \varphi(a')$ with a commutative diagram

}{F7I-VI-028}\begin{center}
\begin{tikzcd}[column sep=huge, row sep=huge]
\psi(b') \ar[rr, "\rho"] \ar[dr, "\psi(\beta)"'] \& \& \psi\ \varphi(a') \ar[dl, "\psi\ \varphi(\alpha)"] \\
 \& \psi\ \varphi(a) \&
\end{tikzcd}
\end{center}

Since $\underline{B}$ is semi-homogeneous, we can choose a commutative diagram

\begin{center}
\begin{tikzcd}[column sep=small, row sep=small]
 \& b'' \ar[dl] \ar[dr] \& \& \& \& R' \underset{R}{\times} R' \ar[dl] \ar[dr] \& \\
b' \ar[dr, "\beta"'] \& \& \varphi(a') \ar[dl, "\varphi(\alpha)"] \& \text{above} \& R' \ar[dr] \& \& R' \ar[dl] \\
 \& \varphi(a) \& \& \& \& R \&
\end{tikzcd}
\end{center}













\SGAvar{where  $R = P(a)$ , $R' = P(b') = P(a')$ . Since $\underline{C}$ is homogeneous, $\psi(b'') = \psi(b') \underset{\psi\ \varphi(a)}{\times} \psi\varphi(a')$ by 2.6(a), and it follows from 2.8(2) that $0 = \tau(\psi\ \varphi(a'), \psi(b')) = (\widetilde{f})_{*}\psi(b'')$ , and hence $(\widetilde{f})_{*}(b'') = 0$ be the hypothesis (ii). It follows from 1.5(2) that there exists a map $\gamma : b' \longrightarrow \varphi(a')$ such that $\beta = \varphi(\alpha)\cdot\gamma$ .

}{where  $R = P(a)$ , $R' = P(b') = P(a')$ . Since $\underline{C}$ is homogeneous, $\psi(b'') = \psi(b') \underset{\psi\ \varphi(a)}{\times} \psi\varphi(a')$ by 2.6(a), and it follows from 2.8(2) that $0 = \tau(\psi\ \varphi(a'), \psi(b')) = (\widetilde{f})_{*}\psi(b'')$ , and hence $(\widetilde{f})_{*}(b'') = 0$ by the hypothesis (ii). It follows from 1.5(2) that there exists a map $\gamma : b' \longrightarrow \varphi(a')$ such that $\beta = \varphi(\alpha)\cdot\gamma$ .

}{F7I-VI-029}(b)  Let a surjective map $\beta : b' \to \varphi(a)$ in $\widehat{\underline{B}}$ be given. This means that we are given a sequence of commutative diagrams
\SGApage{I}{58}
\begin{tikzcd}[column sep=large]
{} \ar[r, dashed] \& b_3' \ar[r, "\beta_2'"] \ar[d, "\beta_3"] \& b_2' \ar[r, "\beta_1'"] \ar[d, "\beta_2"] \& b_1' \ar[d, "\beta_1"] \\
{} \ar[r, dashed] \& \varphi(a_3) \ar[r, "\varphi(\alpha_2)"'] \& \varphi(a_2) \ar[r, "\varphi(\alpha_1)"'] \& \varphi(a_1)
\end{tikzcd}















in $\underline{B}$ , where $\alpha_{v}$ , $\beta_{v}$ , $\beta_{v}'$ are surjective for all $v$ , such that $\beta = \varprojlim_{v} \beta_{v}$ . To show that $\widehat{\varphi}$ is smooth, it suffices to prove that a commutative diagram below of solid arrows in which $\alpha_{n}$ , $\beta_{n}$ , $\beta_{n+1}$ , $\beta_{n}'$ are surjective can be completed into a commutative diagram including dotted arrows.

\begin{center}
\begin{tikzcd}[column sep=large, row sep=large]
\varphi(a_{n+1}') \ar[rrrr, dashed, "\varphi(\ \ )"] \ar[dr, dashed, leftrightarrow] \ar[ddr, dashed, "\varphi(\ \ )"'] \& \& \& \& \varphi(a_n') \ar[dl, leftrightarrow] \ar[ddl, "\varphi(\alpha_n')"] \\
 \& b_{n+1}' \ar[rr, "\beta_n'"] \ar[d, "\beta_{n+1}"] \& \& b_n' \ar[d, "\beta_n"'] \& \\
 \& \varphi(a_{n+1}) \ar[rr, "\varphi(\alpha_n)"'] \& \& \varphi(a_n) \&
\end{tikzcd}
\end{center}


















\SGAvar{Now $\underline{A}$ is, by hypotheses, semi-homogeneous and $\alpha_{n} : a_{n}' \to a_{n}$ is surjective, and therefore there exists a commutative diagram

}{Now $\underline{A}$ is, by hypotheses, semi-homogeneous and $\alpha'_{n} : a_{n}' \to a_{n}$ is surjective, and therefore there exists a commutative diagram

}{F7I-VI-030}\begin{center}
\begin{tikzcd}[column sep=small, row sep=small]
 \& a' \ar[dl, "\gamma"'] \ar[dr, "\gamma'"] \& \& \& \& R_{n+1} \underset{R_n}{\times} R_n' \ar[dl] \ar[dr] \& \& \\
a_{n+1} \ar[dr, "\alpha_n"'] \& \& a_n' \ar[dl, "\alpha_n'"] \& \text{in } \underline{A} \text{ over} \& R_{n+1} \ar[dr] \& \& R_n' \ar[dl] \& \text{in } \underline{C}_{\Lambda} . \\
 \& a_n \& \& \& \& R_n \& \&
\end{tikzcd}
\end{center}













The hypothesis that $\underline{B}$ is homogeneous entails that

\begin{tikzcd}[column sep=huge, row sep=huge]
 \& \varphi(a') \ar[dl, "\varphi(\gamma)"'] \ar[dr, "\varphi(\gamma')"] \& \\
\varphi(a_{n+1}) \ar[dr, "\varphi(\alpha_n)"'] \& \& \varphi(a_n') \ar[dl, "\varphi(\alpha_n')"] \\
 \& \varphi(a_n) \&
\end{tikzcd}









\SGApage{I}{59}
is a diagram of a fibre-product in $\underline{B}$ , and therefore we obtain a unique map $\delta : b'_{n+1} \to \varphi(a')$ in $\underline{B}$ with an appropriate commutativity. Since $\beta_{n+1}$ , $\beta'_n$ are both surjective, so is $\delta$ and therefore, since $\varphi$ is smooth, there exists $a'_{n+1} \xrightarrow{\ \alpha'\ } a'$ in $\underline{A}$ and a $\underline{C}_{\Lambda}$-isomorphism $\varphi(a'_{n+1}) \xrightarrow[\sim]{\rho} b'_{n+1}$ such that $\delta \cdot \rho = \varphi(\alpha')$ . We then obtain the commutative diagram

\SGAvar{\begin{center}
\begin{tikzcd}[row sep=huge, column sep=huge]
\varphi(a'_{n+1})
  \arrow[rrr, "\varphi(\gamma'\cdot\alpha')"]
  \arrow[dr, "\rho"]
  \arrow[ddr, "\varphi(\gamma\cdot\alpha')"']
\& \& \& \varphi(a'_n)
  \arrow[dl, "\sim"]
  \arrow[ddl, "\varphi(\alpha'_n)"] \\
\& b'_{n+1} \arrow[r] \arrow[d] \& b'_n \arrow[d] \& \\
\& \varphi(a_{n+1} \arrow[r] \& \varphi(a_n) \&
\end{tikzcd}
\end{center}


















}{\begin{center}
\begin{tikzcd}[row sep=huge, column sep=huge]
\varphi(a'_{n+1})
  \arrow[rrr, "\varphi(\gamma'\cdot\alpha')"]
  \arrow[dr, "\rho"]
  \arrow[ddr, "\varphi(\gamma\cdot\alpha')"']
\& \& \& \varphi(a'_n)
  \arrow[dl, "\sim"]
  \arrow[ddl, "\varphi(\alpha'_n)"] \\
\& b'_{n+1} \arrow[r] \arrow[d] \& b'_n \arrow[d] \& \\
\& \varphi(a_{n+1}) \arrow[r] \& \varphi(a_n) \&
\end{tikzcd}
\end{center}


















}{F7I-VI-031}\underline{Corollary} 2.10.  (a) \underline{Let} $\varphi : \underline{A} \to \underline{B}$ \underline{be a smooth functor of cofibered groupoids over} $\underline{C}_{\Lambda}$ . \underline{Let} $\xi$ \underline{be a versal formal object of} $\underline{A}$ , \underline{and} $\eta$ \underline{a minimally-versal formal object of} $\underline{B}$ . \underline{If} $\underline{B}$ \underline{is homogeneous, then for any map} $\beta : \eta \to \varphi(\xi)$ \underline{in} $\underline{\hat{B}}$ , $P(\xi)$ \underline{is a formal power-series ring over} $P(\eta)$ \underline{through} $P(\beta) : P(\eta) \to P(\varphi(\xi)) = P(\xi)$ .

(b)  \underline{Let} $\underline{A}$ \underline{be a homogeneous cofibered groupoid over} $\underline{C}_{\Lambda}$ , \underline{and let} $\xi$ , $\eta$ \underline{be versal formal objects of} $\underline{A}$ . \underline{If} $\eta$ \underline{is minimal, then for any map} $\alpha : \eta \to \xi$ \underline{in} $\hat{A}$ , $P(\xi)$ \underline{is a formal power-series ring over} $P(\eta)$ \underline{through} $P(\alpha) : P(\eta) \to P(\xi)$ .

\underline{Proof}:  (a)  A map $\beta : \eta \to \varphi(\xi)$ induces a commutative diagram

\begin{center}
\begin{tikzcd}[row sep=huge, column sep=huge]
\underline{h}_{P(\xi)} \arrow[r, "\widetilde{\beta}"] \arrow[d, "\widetilde{\xi}"'] \& \underline{h}_{P(\eta)} \arrow[d, "\widetilde{\eta}"] \\
\underline{A} \arrow[r, "\varphi"'] \& \underline{B}
\end{tikzcd}
\end{center}









Since $\widetilde{\xi}$ , $\varphi$ are smooth, $\widetilde{\eta}\cdot\widetilde{\beta} = \varphi\cdot\widetilde{\xi}$ is smooth. On the other hand,
\SGApage{I}{60}
$\widetilde{\eta}(k[\epsilon]) : \underline{h}_{P(\eta)}(k[\epsilon]) \to [\underline{B}(k[\epsilon])]$ is an isomorphism and therefore $\widetilde{\beta}$ is smooth by 2.9(a), which means that $P(\xi)$ is a formal power-series ring over $P(\eta)$ .

(b)  follows from (a) by setting $\underline{A} = \underline{B}$ .

\SGAvar{Finally we consider the question of  "representability"  of a cofibered groupoid. Let $\underline{C}$ be a fixed category. Each object $R$ in $\underline{C}$ defines a functor $h_R : \underline{C} \to$ (Ens)  given by  $h_R(S) = \mathrm{Hom}_{\underline{C}}(R,S)$ , and in turn we obtain a functor $h : \underline{C}^0 \to \underline{\mathrm{Hom}}(\underline{C},(\mathrm{Ens}))$ which is fully-faithful. A  "cocartegory in $\underline{C}$ "  is a diagram

}{Finally we consider the question of  "representability"  of a cofibered groupoid. Let $\underline{C}$ be a fixed category. Each object $R$ in $\underline{C}$ defines a functor $h_R : \underline{C} \to$ (Ens)  given by  $h_R(S) = \mathrm{Hom}_{\underline{C}}(R,S)$ , and in turn we obtain a functor $h : \underline{C}^0 \to \underline{\mathrm{Hom}}(\underline{C},(\mathrm{Ens}))$ which is fully-faithful. A  "cocategory in $\underline{C}$ "  is a diagram

}{F7I-VI-032}\begin{tikzcd}[column sep=huge]
h_{R_1} \times_{h_{R_0}} h_{R_1} \arrow[r] \& h_{R_1} \arrow[r, shift left=3, "s_1"] \arrow[r, shift right=3, "s_2"'] \& h_{R_0} \arrow[l, "\epsilon" description]
\end{tikzcd}
$\qquad\text{with } s_i \cdot \epsilon = \text{identity for } i = 1,2$










\begin{equation*}
(\text{where } \ h_{R_1} \times_{h_{R_0}} h_{R_1} \ = \ (h_{R_1})_{s_1} \times_{h_{R_0}} {}_{s_2}(h_{R_1})) \ \text{ such that}
\end{equation*}
(i)  $c$  is associative and is compatible with  $s_i$  for  $i = 1,2$

\begin{multline*}
(\mathrm{ii}) \quad h_{R_1} \xrightarrow{\ \sim\ } h_{R_0} \times_{h_{R_0}} h_{R_1} \xrightarrow{\ \epsilon \times 1\ } h_{R_1} \times_{h_{R_0}} h_{R_1} \xrightarrow{\ c\ } h_{R_1}
\\[2ex]
\phantom{(\mathrm{ii})} \quad h_{R_1} \xrightarrow{\ \sim\ } h_{R_1} \times_{h_{R_0}} h_{R_0} \xrightarrow{\ 1 \times \epsilon\ } h_{R_1} \times_{h_{R_0}} h_{R_1} \xrightarrow{\ c\ } h_{R_1}
\\[1ex]
\text{are both identities.}
\end{multline*}
If  $(R_0, R_1, s_1, s_2, \epsilon, c)$  is a cocategory in $\underline{C}$ , then for an object $S$ in $\underline{C}$ , we can define a category  $h_{(R_0,\ R_1,\ \ldots)}(S)$  by setting  $\mathrm{ob}\ h_{(R_0,\ R_1,\ldots)}(S) = \mathrm{Hom}_{\underline{C}}(R_0,\ S)$ , and  $\mathrm{Fl}\ h_{(R_0,\ R_1,\ldots)}(S) = \mathrm{Hom}_{\underline{C}}(R_1,S)$  in which the composition rule of arrows is given by  $c : h_{R_1} \times_{h_{R_0}} h_{R_1} \to h_{R_1}$ . Then $h_{(R_0,\ R_1,\ldots)} : \underline{C} \to$ (categories)  is a functor and in turn defines a cofibered category $\underline{h}_{(R_0,\ R_1,\ldots)}$ over $\underline{C}$ . For instance, any object $R$ in $\underline{C}$ defines a discrete cocategory $\underline{h}_R = \underline{h}_{(R,\ R,\ s_1,\ s_2,\ \epsilon,\ c)}$ by setting $s_1$, $s_2$, $\epsilon$, $c$ to be all equal to the identity map. A homomorphism  $(R_0,\ R_1,\ldots) \to (R_0{}',\ R_1{}',\ldots)$
\SGApage{I}{61}
of cocategories in $\underline{C}$ is a pair of maps $h_{R_0} \to h_{R'_0}$ , $h_{R_1} \to h_{R'_1}$ in $\underline{C}$ such that the induced map

\begin{equation*}
h_{(R_0,\ R_1,\ldots)}(S) \to h_{(R_0',\ R_1',\ldots)}(S)
\end{equation*}
is a functor for every object $S$ in $\underline{C}$ . For instance, each cocategory $(R_0,\ R_1,\ldots)$ is provided with a homomorphism $\epsilon : \underline{h}_{R_0} \to \underline{h}_{(R_0,\ R_1,\ldots)}$ of cocategories. A cogroupoid in $\underline{C}$ is by definition a cocategory $(R_0,\ R_1,\ldots)$ in $\underline{C}$ such that $h_{(R_0,\ R_1,\ldots)}(S)$ is a groupoid for every object $S$ in $\underline{C}$ . In other words, a cogroupoid in $\underline{C}$ is a cocategory $(R_0,\ R_1,\ldots)$ in $\underline{C}$ such that $\underline{h}_{(R_0,\ R_1,\ldots)}$ is a cofibered category in groupoids over $\underline{C}$ .

In practice, a cocategory or cogroupoid occur in a following fashion: Let $P : \underline{F} \to \underline{C}$ be a cofibered category, and let $\xi$ be a fixed object in $\underline{F}$ . Set $R_0 = P(\xi)$ and let $(R_0,\ \underline{C})$ be the category of arrows in $\underline{C}$ with the source $R_0$ . Let $P : \underline{F}' \to (R_0,\ \underline{C})$ be the pull-back of $\underline{F}$ under the canonical functor $(R_0,\ \underline{C}) \to \underline{C}$ , and choose a cocartesian section $\xi^{\#}$ of $\underline{F}'$ over $(R_0,\ \underline{C})$ such that $\xi^{\#}(\mathrm{id}_{R_0}) = \xi$ . (This simply means that we choose, for each arrow $R_0 \xrightarrow{f} S$ in $\underline{C}$ , a cocartesian $f$-map $\xi \to \xi^{\#}(f)$ ). Define the functor

\SGAvar{\begin{equation*}
\mathrm{Hom}_{\xi^{\#}} : (R_0 \amalg R_0,\ \underline{C}) \to (\mathrm{Ens}) \qquad \text{given by}
\end{equation*}
\begin{equation*}
X = (f_1,\ f_2) \longmapsto \mathrm{Hom}_{\underline{F}(S)}(\xi^{\#}(f_1),\ \xi^{\#}(f_2)) \qquad \text{where}
\end{equation*}
$S$ $=$ the target of $f_1$ ( $=$ the target of $f_2$ ). We note that this functor does not depend on the choice of $\xi^{\#}$ up to (canonical) natural equivalence. Henceforth, by abuse of notation, we shall simply write $\xi$ for $\xi^{\#}$ . If this functor is representable by an object $X_0 = (R_0 \ \overset{s_1}{\underset{s_2}{=\!=}}\ R_1)$ in $(R_0 \qquad R_0,\ \underline{C})$ i.e. if there exists $R_0$-map $\varphi : \xi(s_1) \to \xi(s_2)$ in $\underline{F}$ such that the induced map $\widetilde{\varphi} : \mathrm{Hom}_{(R_0 \amalg R_0,\ \underline{C})}(X_0,\ X) \to \mathrm{Hom}_{\xi}(X)$ is bijective for all $X$ in $(R_0 \amalg R_0,\ \underline{C})$ , then the object $X_0$ defines a cocategory $(R_0,\ R_1,\ldots)$ in $\underline{C}$ in a natural way. Indeed, a triple of arrows
\SGApage{I}{62}
\begin{tikzcd}[column sep=huge]
R_0 \arrow[r, shift left=2.5, "u"] \arrow[r, "v"] \arrow[r, shift right=2.5, "w"'] \& S
\end{tikzcd}

}{\begin{equation*}
\mathrm{Hom}_{\xi^{\#}} : (R_0 \amalg R_0,\ \underline{C}) \to (\mathrm{Ens}) \qquad \text{given by}
\end{equation*}
\begin{equation*}
X = (f_1,\ f_2) \longmapsto \mathrm{Hom}_{\underline{F}(S)}(\xi^{\#}(f_1),\ \xi^{\#}(f_2)) \qquad \text{where}
\end{equation*}
$S$ $=$ the target of $f_1$ ( $=$ the target of $f_2$ ). We note that this functor does not depend on the choice of $\xi^{\#}$ up to (canonical) natural equivalence. Henceforth, by abuse of notation, we shall simply write $\xi$ for $\xi^{\#}$ . If this functor is representable by an object $X_0 = (R_0 \ \overset{s_1}{\underset{s_2}{\rightrightarrows}}\ R_1)$ in $(R_0 \amalg R_0,\ \underline{C})$ i.e. if there exists $R_0$-map $\varphi : \xi(s_1) \to \xi(s_2)$ in $\underline{F}$ such that the induced map $\widetilde{\varphi} : \mathrm{Hom}_{(R_0 \amalg R_0,\ \underline{C})}(X_0,\ X) \to \mathrm{Hom}_{\xi}(X)$ is bijective for all $X$ in $(R_0 \amalg R_0,\ \underline{C})$ , then the object $X_0$ defines a cocategory $(R_0,\ R_1,\ldots)$ in $\underline{C}$ in a natural way. Indeed, a triple of arrows
\SGApage{I}{62}
\begin{tikzcd}[column sep=huge]
R_0 \arrow[r, shift left=2.5, "u"] \arrow[r, "v"] \arrow[r, shift right=2.5, "w"'] \& S
\end{tikzcd}

}{F7I-VI-033}in $\underline{C}$ defines a composition

\begin{equation*}
\mathrm{Hom}_{\underline{F}(S)}(\xi(v),\ \xi(w)) \times \mathrm{Hom}_{\underline{F}(S)}(\xi(u),\ \xi(v)) \to \mathrm{Hom}_{\underline{F}(S)}(\xi(u),\ \xi(w))
\end{equation*}
and in turn yields the composition rule

\begin{equation*}
h_{R_1}(S) \underset{h_{R_0}(S)}{\times} h_{R_1}(S) \to h_{R_1}(S)
\end{equation*}
depending functorially on $S$ . Consequently the functor $\mathrm{Hom}_{\xi}$ defines the cocategory $(R_0,\ R_1,\ldots)$ in $\underline{C}$ and in turn defines a cofibered category $\underline{h}_{(R_0,\ R_1,\ldots)}$ in $\underline{C}$ . By definition, ob $\underline{h}_{(R_0,\ R_1,\ldots)} = \coprod_{S} \mathrm{Hom}_{\underline{C}}(R_0,\ S)$ (disjoint union for $S \in \mathrm{ob}(\underline{C})$ ), and an arrow $(R_0 \xrightarrow{f} S) \to (R_0 \xrightarrow{f'} S')$ is a pair $S \xrightarrow{h} S'$ and $R_1 \xrightarrow{u} S'$ in $\underline{C}$ such that $u \cdot s_1 = h \cdot f$ and $u \cdot s_2 = f'$ , and the structural functor $P : \underline{h}_{(R_0,\ R_1,\ldots)} \to \underline{C}$ is given by $(R_0 \xrightarrow{f} S) \longmapsto S$ . Now the pair $(\xi,\ \varphi)$ defines the functor

\SGAvar{\begin{equation*}
(\xi,\ \varphi) : \underline{h}_{(R_0,\ R_1,\ldots)} \to \underline{F} \quad \text{over} \quad \underline{C} \quad \text{by}
\end{equation*}
\begin{equation*}
(R_0 \xrightarrow{f} S) \ \longmapsto\ \xi(f)
\end{equation*}
\begin{equation*}
[(R_0 \xrightarrow{f} S) \xrightarrow{\ (h,u)\ } (R_0 \to S')] \longmapsto \ \text{the composit map}\ \ \xi(f) \xrightarrow{\ h_{*}\ } \xi(hf) \xrightarrow{\ u(\varphi)\ } \xi(f') \ .
\end{equation*}
We note that, for each object $S$ in $\underline{C}$ , the functor $(\xi,\varphi)(S)$   $\underline{h}_{(R_0,R_1,\ldots)}(S) \to \underline{F}(S)$ is fully-faithful by virtue of the definition of $(R_0 \rightrightarrows R_1)$ coming from the representability of the functor $\mathrm{Hom}_{\xi}$ . Consequently the functor $(\xi,\varphi) : \underline{h}_{(R_0,R_1,\ldots)} \to \underline{F}$ is an $\underline{C}$-equivalence if and only if $\xi$ is a quasi-initial object of $\underline{F}$ over $\underline{C}$ i.e., for every object $a$ in $\underline{F}$ there exists a cocartesian arrow $\xi \to a$ . Instead of this functor $\mathrm{Hom}_{\xi}$ one may also consider the functor

}{\begin{equation*}
(\xi,\ \varphi) : \underline{h}_{(R_0,\ R_1,\ldots)} \to \underline{F} \quad \text{over} \quad \underline{C} \quad \text{by}
\end{equation*}
\begin{equation*}
(R_0 \xrightarrow{f} S) \ \longmapsto\ \xi(f)
\end{equation*}
\begin{equation*}
[(R_0 \xrightarrow{f} S) \xrightarrow{\ (h,u)\ } (R_0 \to S')] \longmapsto \ \text{the composite map}\ \ \xi(f) \xrightarrow{\ h_{*}\ } \xi(hf) \xrightarrow{\ u(\varphi)\ } \xi(f') \ .
\end{equation*}
We note that, for each object $S$ in $\underline{C}$ , the functor $(\xi,\varphi)(S) :$   $\underline{h}_{(R_0,R_1,\ldots)}(S) \to \underline{F}(S)$ is fully-faithful by virtue of the definition of $(R_0 \rightrightarrows R_1)$ coming from the representability of the functor $\mathrm{Hom}_{\xi}$ . Consequently the functor $(\xi,\varphi) : \underline{h}_{(R_0,R_1,\ldots)} \to \underline{F}$ is a $\underline{C}$-equivalence if and only if $\xi$ is a quasi-initial object of $\underline{F}$ over $\underline{C}$ i.e., for every object $a$ in $\underline{F}$ there exists a cocartesian arrow $\xi \to a$ . Instead of this functor $\mathrm{Hom}_{\xi}$ one may also consider the functor

}{F7I-VI-034}\begin{equation*}
\mathrm{Isom}_{\xi} : (R_0 \amalg R_0,\ \underline{C}) \to (\mathrm{Ens}) 
\end{equation*}

\SGApage{I}{63}

 given by  $X = (f_1,f_2) \longmapsto \mathrm{Isom}_{\underline{F}(S)}(\xi(f_1),\xi(f_2))$  where  $S$  is the target of  $f_1$ ( $=$ the target of  $f_2$ ). If this functor is representable by an object  $(R_0 \rightrightarrows R_1)$ in  $(R_0 \amalg R_0, \underline{C})$  , then it defines  a cogroupoid  $(R_0,R_1,\ldots)$  in  $\underline{C}$  together with the functor  $(\xi,\varphi) : \underline{h}_{(R_0,R_1,\ldots)} \to \underline{F}$ , which is, when  $\underline{F}$  is a cofibered groupoid, a  $\underline{C}$-equivalence if and only if  $\xi$  is a quasi-initial object of  $\underline{F}$  over  $\underline{C}$ .


\underline{Definition and proposition} 2.11.  \underline{A cogroupoid}  $(R_0,R_1,\ldots)$  \underline{in}  $\hat{\underline{C}}_\Lambda$  \underline{is called smooth if any one of the following equivalent conditions holds:}

 (i)  $R_1$  \underline{is a formal power-series ring over}  $R_0$  \underline{via} $s_1$  (\underline{and}  $s_2$ )

 (ii)  $\epsilon : \underline{h}_{R_0} \longrightarrow \underline{h}_{(R_0,R_1,\ldots)}$  \underline{is smooth over}  $\underline{C}_\Lambda$ .

\underline{A minimally smooth cogroupoid or a normalized smooth cogroupoid in}  $\hat{\underline{C}}_\Lambda$  \underline{is a smooth cogroupoid}  $(R_0,R_1,\ldots)$  \underline{in}  $\hat{\underline{C}}_\Lambda$  \underline{such that}  $\underline{h}_{(R_0,R_1,\ldots)}(k[\epsilon])$  \underline{is totally disconnected} i.e.,  $s_1(k[\epsilon]) = s_2(k[\epsilon])$  .

\SGAvar{\underline{Proof}:  $\epsilon : \underline{h}_{R_0} \to \underline{h}_{(R_0,R_1,\ldots)}$  is smooth over  $\underline{C}_\Lambda \Leftrightarrow$ given any surjective map  $(S',f') \xrightarrow{\ (h,u)\ } (S,\ f)$  in  $\underline{h}_{(R_0,R_1,\ldots)}$  , there exists a lifting  $f''$  of  $f$ to  $S'$  together with an isomorphism  $(S',f'') \xrightarrow{\ (1,u')\ } (S',f') \Leftrightarrow$ given  $u \in h_{R_1}(S)$ and a lifting  $f'$  or  $s_1(u)$  to  $S'$  , there exists a lifting  $u'$  of  $u$  to  $S'$ such that  $s_1(u') = f'$  , where  $S' \to S$  is an arbitrary surjective map in $\underline{C}_\Lambda \Leftrightarrow s_1 : \underline{h}_{R_1} \to \underline{h}_{R_0}$  is smooth over  $\underline{C}_\Lambda$ .


}{\underline{Proof}:  $\epsilon : \underline{h}_{R_0} \to \underline{h}_{(R_0,R_1,\ldots)}$  is smooth over  $\underline{C}_\Lambda \Leftrightarrow$ given any surjective map  $(S',f') \xrightarrow{\ (h,u)\ } (S,\ f)$  in  $\underline{h}_{(R_0,R_1,\ldots)}$  , there exists a lifting  $f''$  of  $f$ to  $S'$  together with an isomorphism  $(S',f'') \xrightarrow{\ (1,u')\ } (S',f') \Leftrightarrow$ given  $u \in h_{R_1}(S)$ and a lifting  $f'$  of  $s_1(u)$  to  $S'$  , there exists a lifting  $u'$  of  $u$  to  $S'$ such that  $s_1(u') = f'$  , where  $S' \to S$  is an arbitrary surjective map in $\underline{C}_\Lambda \Leftrightarrow s_1 : \underline{h}_{R_1} \to \underline{h}_{R_0}$  is smooth over  $\underline{C}_\Lambda$ .


}{F7I-VI-035}\underline{Proposition} 2.12.  \underline{Let}  $(R_0,R_1,\ldots)$  \underline{be a cogroupoid in}  $\hat{\underline{C}}_\Lambda$ .

(a)  \underline{If}  $(R_0,R_1,\ldots)$  \underline{is smooth, then}  $\underline{h}_{(R_0,R_1,\ldots)}$  \underline{is homogeneous.}

(b)  \underline{If}  $s_1(k[\epsilon]) = s_2(k[\epsilon])$  , \underline{then the following statements are equivalent}

          (i)  $(R_0,R_1,\ldots)$  \underline{is smooth}

         (ii)  $\underline{h}_{(R_0,R_1,\ldots)}$  \underline{is homogeneous}

        (iii)  $\underline{h}_{(R_0,R_1,\ldots)}$  \underline{is semi-homogeneous}


\SGApage{I}{64}

(c)  \underline{If}  $\underline{h}_{(R_0,R_1,\ldots)}$  \underline{is semi-homogeneous, then it is equivalent to}  $\underline{h}_{(S_0,S_1,\ldots)}$  \underline{for some minimally smooth cogroupoid}  $S_*$  .

\underline{Proof}:   (a)  Consider a diagram

\begin{center}
\begin{tikzcd}[column sep=huge, row sep=large]
(S_1,f_1) \arrow[dr, "{(h_1,u_1)}"'] \& \& (S_2,f_2) \arrow[dl, "{(h_2,u_2)}"] \\
\& (S,f) \&
\end{tikzcd}
\end{center}

in  $\underline{h}_{(R_0,R_1,\ldots)}$  where  $h_2 : S_2 \longrightarrow S$  is surjective in  $\underline{C}_\Lambda$ . Since $s_1 : \underline{h}_{R_1} \longrightarrow \underline{h}_{R_0}$  is smooth by hypothesis,  $u_1^{-1}\, u_2$  can be lifted to some $v \in h_{R_1}(S_2)$  such that  $s_1(v) = f_2$ . Since  $h_2 \cdot s_2(v) = h_1 \cdot f_1$ , $f_1 \times s_2(v) \in h_{R_0}(S_1 \underset{S}{\times} S_2)$  is defined and we obtain the commutative diagram

\begin{center}
\begin{tikzcd}[column sep=huge, row sep=large]
\& (S_1 \underset{S}{\times} S_2,\ f_1 \times s_2(v)) \arrow[dl, "\pi_1"'] \arrow[dr, "{(\pi_2,v^{-1})}"] \& \\
(S_1,f_1) \arrow[dr, "{(h_1,u_1)}"'] \& \& (S_2,f_2) \arrow[dl, "{(h_2,u_2)}"] \\
\& (S,f) \&
\end{tikzcd}
\end{center}










We claim that the above diagram defines a fibre-product  $(S_1,f_1) \underset{(S,f)}{\times} (S_2,f_2)$  in $\underline{h}_{(R_0,R_1,\ldots)}$ . Indeed, for any object  $(T,\ g)$  in  $\underline{h}_{(R_0,R_1,\ldots)}$  we have

{\small \begin{multline*}
\mathrm{Hom}((T,g),(S_1\underset{S}{\times}S_2,f_1\times s_2(v)))=\{(h_1'\times h_2',w_1\times w_2)\in h_T(S_1\underset{S}{\times}S_2)\times h_{R_1}(S_1\underset{S}{\times}S_2)\,|\,s_1(w_1\times w_2)=(h_1'\times h_2')\ g\ ,\\
s_2(w_1\times w_2)=f_1\times s_2(v)\} = \{(h_1'\times h_2',w_1,w_2) \in h_T(S_1\underset{S}{\times}S_2)\ \times h_{R_1}(S_1)\times h_{R_1}(S_2)\,|\,h_1w_1 = h_2w_2,\\
s_1(w_1) = h_1'\ g,\ s_1(w_2) = h_2'\ g,\ s_2(w_1) = f_1,\ s_2(w_2) = s_2(v)\} =\\
\{(h_1'\times h_2',w_1,w_2) \in h_T(S_1\underset{S}{\times}S_2)\times h_{R_1}(S_1)\times h_{R_1}(S_2)\,|\ u_1^{-1}\ u_2\ h_2(v^{-1}w_2)=h_1\ w_1\ ,\ s_1(w_1) = h_1'\ \ g, 
\end{multline*}}

\SGApage{I}{65}

\SGAvar{{\footnotesize \begin{multline*}s_2(w_1)=f_1, \quad s_1(v^{-1}w_2)=h_2'g, \quad s_2(v^{-1}w_2)=f_2\} \;\overset{\sim}{\longrightarrow}\; \{(h_1'\times h_2', w_1, w_2') \in h_T(S_1\times_S S_2)\times h_{R_1}(S_1)\times h_{R_1}(S_2)\;\big| \\ u_2\, h_2\, w_2' \;=\; u_1\, h_1\, w_1, \quad s_1(w_1) \;=\; h_1'g, \quad s_2(w_1)=\; f_1, \quad s_1(w_2') \;=\; h_2'g, \quad s_2(w_2') \;=\; f_2\} \;= \\ \mathrm{Hom}((T,\ g),\ (S_1,\ f_1) \;\times_{\mathrm{Hom}((T,g),\ (S,f))}\; \mathrm{Hom}((T,\ g),\ (S_2,\ f_2))\ .\end{multline*}}
(b)  It suffices to show (iii) $\Rightarrow$ (i) . Since $\underline{h}_{(R_0,R_1,\ldots)}$ is semi-homogeneous and $\dim[\underline{h}_{(R_0,R_1,\ldots)}(k[\epsilon])] = \dim h_{R_0}(k[\epsilon]) < \infty$ , $\underline{h}_{(R_0,R_1,\ldots)}$ admits a minimally versal formal object. On the other hand, $(R_0,\ 1_{R_0})$ is a quasi-initial minimal object of $\widehat{\underline{h}}_{(R_0,R_1,\ldots)}$ since $\underline{h}_{R_0}(k[\epsilon]) \to [\underline{h}_{(R_0,R_1,\ldots,)}(k[\epsilon])]$ is bijective, and therefore $(R_0,\ 1_{R_0})$ is a minimally versal formal object of $\underline{h}_{(R_0,R_1,\ldots)}$ by 1.10(ii), which means that $\epsilon : \underline{h}_{R_0} \longrightarrow \underline{h}_{(R_0,R_1,\ldots)}$ is minimally smooth.

(c)  Let $(S_0,\ \xi_0)$ be a minimally versal formal object for $\underline{h}_{(R_0,R_1,\ldots)}$ (1.11). For any object $\xi : R_0 \to T$ in $\underline{h}_{(R_0,\ R_1,\ldots)}(T)$ , the functor $\mathrm{Isom}(\xi,\xi)$ from $C_\Lambda/T \otimes_\Lambda T$ to (sets) which to any $R$ associates the sets of isomorphisms between the two images of $\xi$ on $R : R_0 \overset{\xi}{\longrightarrow} T \rightrightarrows T \otimes_\Lambda T \longrightarrow R$ , is represented by $R_1 \otimes_{R_0 \hat{\otimes}_\Lambda R_0} (T \otimes_\Lambda T)$ . By a passage to limit, one get that $\mathrm{Isom}(\xi_0,\ \xi_0)$ is pro-represented by some $S_1$ , and $\underline{h}_{(R_0,R_1,\ldots)}$ is equivalent to $\underline{h}_{(S_0,S_1,\ldots)}$ . By (b), $S_*$ is minimally smooth.

}{{\footnotesize \begin{multline*}s_2(w_1)=f_1, \quad s_1(v^{-1}w_2)=h_2'g, \quad s_2(v^{-1}w_2)=f_2\} \;\overset{\sim}{\longrightarrow}\; \{(h_1'\times h_2', w_1, w_2') \in h_T(S_1\times_S S_2)\times h_{R_1}(S_1)\times h_{R_1}(S_2)\;\big| \\ u_2\, h_2\, w_2' \;=\; u_1\, h_1\, w_1, \quad s_1(w_1) \;=\; h_1'g, \quad s_2(w_1)=\; f_1, \quad s_1(w_2') \;=\; h_2'g, \quad s_2(w_2') \;=\; f_2\} \;= \\ \mathrm{Hom}((T,\ g),\ (S_1,\ f_1)) \;\times_{\mathrm{Hom}((T,g),\ (S,f))}\; \mathrm{Hom}((T,\ g),\ (S_2,\ f_2))\ .\end{multline*}}
(b)  It suffices to show (iii) $\Rightarrow$ (i) . Since $\underline{h}_{(R_0,R_1,\ldots)}$ is semi-homogeneous and $\dim[\underline{h}_{(R_0,R_1,\ldots)}(k[\epsilon])] = \dim h_{R_0}(k[\epsilon]) < \infty$ , $\underline{h}_{(R_0,R_1,\ldots)}$ admits a minimally versal formal object. On the other hand, $(R_0,\ 1_{R_0})$ is a quasi-initial minimal object of $\widehat{\underline{h}}_{(R_0,R_1,\ldots)}$ since $\underline{h}_{R_0}(k[\epsilon]) \to [\underline{h}_{(R_0,R_1,\ldots)}(k[\epsilon])]$ is bijective, and therefore $(R_0,\ 1_{R_0})$ is a minimally versal formal object of $\underline{h}_{(R_0,R_1,\ldots)}$ by 1.10(ii), which means that $\epsilon : \underline{h}_{R_0} \longrightarrow \underline{h}_{(R_0,R_1,\ldots)}$ is minimally smooth.

(c)  Let $(S_0,\ \xi_0)$ be a minimally versal formal object for $\underline{h}_{(R_0,R_1,\ldots)}$ (1.11). For any object $\xi : R_0 \to T$ in $\underline{h}_{(R_0,\ R_1,\ldots)}(T)$ , the functor $\mathrm{Isom}(\xi,\xi)$ from $\underline{C}_\Lambda/T \hat{\otimes}_\Lambda T$ to (sets) which to any $R$ associates the set of isomorphisms between the two images of $\xi$ on $R : R_0 \overset{\xi}{\longrightarrow} T \rightrightarrows T \hat{\otimes}_\Lambda T \longrightarrow R$ , is represented by $R_1 \otimes_{R_0 \hat{\otimes}_\Lambda R_0} (T \otimes_\Lambda T)$ . By a passage to limit, one gets that $\mathrm{Isom}(\xi_0,\ \xi_0)$ is pro-represented by some $S_1$ , and $\underline{h}_{(R_0,R_1,\ldots)}$ is equivalent to $\underline{h}_{(S_0,S_1,\ldots)}$ . By (b), $S_*$ is minimally smooth.

}{F7I-VI-036}\underline{Definition} 2.13.  \underline{A cofibered groupoid} $\underline{A}$ \underline{over} $\underline{C}_\Lambda$ \underline{is called (minimally) pro-representable if it is} $\underline{C}_\Lambda$\underline{-equivalent to} $\underline{h}_{(R_0,R_1,\ldots)}$ \underline{for some (normalized) cogroupoid} $(R_0,R_1,\ldots)$ \underline{in} $\widehat{\underline{C}}_\Lambda$ .

\underline{Proposition} 2.14.  \underline{If} $(R_0,\ R_1,\ldots)$ \underline{is a normalized cogroupoid in} $\widehat{\underline{C}}_\Lambda$ , \underline{then} \underline{every homomorphism} $(\alpha_0,\alpha_1) : (R_0,R_1,\ldots) \to (R_0,R_1,\ldots)$ \underline{which induces} $\underline{C}_\Lambda$\underline{-equivalence} $\underline{h}_{(\alpha_0,\alpha_1)} : \underline{h}_{(R_0,R_1,\ldots)} \longrightarrow \underline{h}_{(R_0,\ R_1,\ldots)}$ \underline{is an isomorphism. In}
\SGApage{I}{66}
\underline{particular, for any cofibered groupoid} $\underline{A}$ \underline{over} $\underline{C}_\Lambda$ , \underline{a normalized prorepresentation, if it exists, is unique up to an isomorphism.}

\underline{Proof}:  Consider the commutative diagram

\begin{center}
\begin{tikzcd}[column sep=huge, row sep=huge]
\underline{h}_{(R_0,R_1,\ldots)} \arrow[r, "{\underline{h}_{(\alpha_0,\alpha_1)}}"] \& \underline{h}_{(R_0,R_1,\ldots)} \\
\underline{h}_{R_0} \arrow[u, "\epsilon"] \arrow[r, "\underline{h}_{\alpha_0}"'] \& \underline{h}_{R_0} \arrow[u, "\epsilon"']
\end{tikzcd}
\end{center}













\SGAvar{where $\underline{h}_{(\alpha_0,\alpha_1)}$ is a $\underline{C}_\Lambda$-equivalence. Since $\underline{h}_{(R_0,R_1,\ldots)}(k[\epsilon])$ is totally disconnected, $[\underline{h}_{(\alpha_0,\alpha_1)}(k[\epsilon])$ and $[\epsilon(k[\epsilon])]$ are both bijections, and hence so is $\underline{h}_{\alpha_0}(k[\epsilon])$ .

}{where $\underline{h}_{(\alpha_0,\alpha_1)}$ is a $\underline{C}_\Lambda$-equivalence. Since $\underline{h}_{(R_0,R_1,\ldots)}(k[\epsilon])$ is totally disconnected, $[\underline{h}_{(\alpha_0,\alpha_1)}(k[\epsilon])]$ and $[\epsilon(k[\epsilon])]$ are both bijections, and hence so is $\underline{h}_{\alpha_0}(k[\epsilon])$ .

}{F7I-VI-037}Consequently, $\alpha_0 : R_0 \to R_0$ is a surjective endomorphism and hence is an automorphism. It then follows that $\alpha_1$ is also an automorphism of $R_1$ .

\underline{Lemma} 2.15.  \underline{Let} $A,\ B$ \underline{be discrete groupoids over} $\underline{C}_\wedge$ \underline{and} $\varphi : A \to B$ \underline{a minimally smooth functor}. \underline{If} $A,\ B$ \underline{are homogeneous}, \underline{then} $\varphi$ \underline{is a} $\underline{C}_\wedge$\underline{-equivalence.}

\underline{Proof}:  We must show that $\varphi(S) : A(S) \longrightarrow B(S)$ is bijective for all $S \in \mathrm{ob}(\underline{C}_\wedge)$ . We note that $\varphi(S)$ is surjective for all $S \in \mathrm{ob}(\underline{C}_\wedge)$ and bijective for $S = k[\epsilon]$ . We argue by induction on the length of $S$ . Let $S' \overset{f}{\longrightarrow} S$ be a small extension in $\underline{C}_\wedge$ . Since $A,\ B$ are homogeneous discrete groupoids, the canonical isomorphism $1 \times \tilde{f} : S' \times_S S' \longrightarrow S' \times_k k[\mathrm{Ker}\ f]$ induces the commutative diagram
\SGApage{I}{67}
\begin{tikzcd}[column sep=huge, row sep=huge]
A(S') \underset{A(S)}{\times} A(S') \arrow[r, "\sim"] \arrow[d] \& A(S') \times A(k[\mathrm{Ker}\ f]) \arrow[d] \\
B(S') \underset{B(S)}{\times} B(S') \arrow[r, "\sim"] \& B(S') \times B(k[\mathrm{Ker}\ f])
\end{tikzcd}













Therefore  $A(S')$ , $B(S')$  are principal homogeneous sets over  $A(S)$ , $B(S)$  with the structural group  $A(k[\mathrm{Ker}\ f]) \xrightarrow{\ \sim\ } B(k[\mathrm{Ker}\ f])$  . Consequently if  $\varphi(S)$  is bijective, then  $\varphi(S')$  is injective and hence is bijective.

\SGAvar{\underline{Corollary} 2.16.  \underline{Let}  $\underline{A}$  \underline{be a homogeneous groupoid over}  $\underline{C}_{\wedge}$  \underline{such that} $\mathrm{Aut}_{\underline{A}(k[\epsilon])}(1_{\epsilon})$  \underline{is of finite-dimensional. Then for any}  $\xi \in \mathrm{ob}(\hat{\underline{A}})$  , \underline{the functor} $\mathrm{Isom}_{\xi} : \underline{C}_{R \hat{\otimes}_{\wedge} R} \longrightarrow (\mathrm{Ens})$  \underline{is prorepresentable, where}  $R = P(\xi)$  .

}{\underline{Corollary} 2.16.  \underline{Let}  $\underline{A}$  \underline{be a homogeneous groupoid over}  $\underline{C}_{\wedge}$  \underline{such that} $\mathrm{Aut}_{\underline{A}(k[\epsilon])}(1_{\epsilon})$  \underline{is finite-dimensional. Then for any}  $\xi \in \mathrm{ob}(\hat{\underline{A}})$  , \underline{the functor} $\mathrm{Isom}_{\xi} : \underline{C}_{R \hat{\otimes}_{\wedge} R} \longrightarrow (\mathrm{Ens})$  \underline{is prorepresentable, where}  $R = P(\xi)$  .

}{F7I-VI-038}\underline{Proof}:  Assume that  $\mathrm{Isom}_{\xi}$  is homogeneous. We then have, by 1.11, a minimally smooth functor  $\Phi : \underline{h}_{R_{1}} \to \underline{\mathrm{Isom}}_{\xi}$  , which is a  $\underline{C}_{R \hat{\otimes}_{\wedge} R}$-equivalence by 2.15 i.e., $\mathrm{Isom}_{\xi}$  is prorepresentable. Therefore it suffices to show that  $\mathrm{Isom}_{\xi}$  is homogeneous. Consider a diagram

\begin{center}
\begin{tikzcd}[column sep=large, row sep=large]
T_1 \arrow[dr] \& \& T_2 \arrow[dl] \\
 \& T \&
\end{tikzcd}
\end{center}

in  $\underline{C}_{R \hat{\otimes}_{\wedge} R}$  where  $T_{2} \longrightarrow T$  is surjective. Since  $\underline{A}$  is homogeneous, $\xi_{i}(T_{1}) \underset{\xi_{i}(T)}{\times} \xi_{i}(T_{2})$  exists in  $\underline{A}$  , and we have a canonical isomorphism

$\xi_{i}(T_{1} \underset{T}{\times} T_{2}) \simeq \xi_{i}(T_{1}) \underset{\xi_{i}(T)}{\times} \xi_{i}(T_{2})$  for  $i = 1,2$  (cf. 2.6(a)). It follows that

$\mathrm{Isom}_{\xi}(T_{1} \underset{T}{\times} T_{2}) \longrightarrow \mathrm{Isom}_{\xi}(T_{1}) \underset{\mathrm{Isom}_{\xi}(T)}{\times} \mathrm{Isom}_{\xi}(T_{2})$  is an isomorphism.

\underline{Theorem} 2.17.  \underline{Let}  $\underline{A}$  \underline{be a cofibered groupoid over}  $\underline{C}_{\wedge}$  . \underline{Then it is prorepresentable by a (normalized) smooth cogroupoid if and only if}
\SGApage{I}{68}
(1)  $\underline{A}$  \underline{is homogeneous}

(2)  $[A(k[\epsilon])]$  \underline{and}  $[\mathrm{Aut}(1_{\epsilon})]$  \underline{are both finite-dimensional vector spaces over}  $k$  .

\underline{Proof}:  Assume that  $\underline{A}$  is prorepresentable by a smooth cogroupoid i.e.,  $\underline{A}$  is  $\underline{C}_{\wedge}$-equivalent to  $\underline{h} = \underline{h}_{(R_{0},R_{1},\ldots)}$  for some smooth cogroupoid  $(R_{0},R_{1},\ldots)$  in  $\hat{\underline{C}}_{\wedge}$  . Since  $\underline{h}$  is homogeneous by 2.12(a), and

\SGAvar{\begin{equation*}
\begin{aligned}
[\underline{h}(k[\epsilon])] &= \mathrm{Coker}\ (\mathrm{Hom}_{\wedge\text{-alg}}(R_{1},k[\epsilon] \ \underset{s_{2}}{\overset{s_{1}}{\rightrightarrows}}\ \mathrm{Hom}_{\wedge\text{-alg}}(R_{0},k[\epsilon])) \qquad \text{as well as} \\[2ex]
\mathrm{Aut}_{\underline{h}}(1_{\epsilon}) &= \mathrm{Ker}(\mathrm{Hom}_{\wedge\text{-alg}}(R_{1},k[\epsilon]) \xrightarrow{\ s_{1}\ } \mathrm{Hom}_{\wedge\text{-alg}}(R_{0},k[\epsilon])) \\[2ex]
&= \mathrm{Ker}(\mathrm{Hom}_{\wedge\text{-alg}}(R_{1},k[\epsilon]) \xrightarrow{\ s_{2}\ } \mathrm{Hom}_{\wedge\text{-alg}}(R_{0},k[\epsilon])) \qquad \text{are both finite-dimensional}
\end{aligned}
\end{equation*}
vector spaces over  $k$  , we must have the same properties on  $\underline{A}$  . Conversely assume (1) and (2) on  $\underline{A}$  . We claim that  $\underline{A}$  is prorepresentable by a normalized smooth cogroupoid. Since  $\underline{A}$  is homogeneous and  $[\underline{A}(k[\epsilon])]$  is a finite-dimensional vector space over  $k$  ,  $\underline{A}$  admits a minimally versal formal object  $\xi$  . Set  $P(\xi) = R_{0}$  and consider the functor  $\mathrm{Isom}_{\xi} : \underline{C}_{R_{0} \hat{\otimes}_{\wedge} R_{0}} \to (\mathrm{Ens})$  given by

\begin{equation*}
T \longmapsto \mathrm{Isom}_{\underline{A}(T)}(\xi_{1}(T),\xi_{2}(T))
\end{equation*}
where  $\xi_{i}(T) = \xi(f_{i})$  and  $f_{i} : R_{0} \to T$  denotes the composit maps

}{\begin{equation*}
\begin{aligned}
[\underline{h}(k[\epsilon])] &= \mathrm{Coker}\ (\mathrm{Hom}_{\wedge\text{-alg}}(R_{1},k[\epsilon]) \ \underset{s_{2}}{\overset{s_{1}}{\rightrightarrows}}\ \mathrm{Hom}_{\wedge\text{-alg}}(R_{0},k[\epsilon])) \qquad \text{as well as} \\[2ex]
\mathrm{Aut}_{\underline{h}}(1_{\epsilon}) &= \mathrm{Ker}(\mathrm{Hom}_{\wedge\text{-alg}}(R_{1},k[\epsilon]) \xrightarrow{\ s_{1}\ } \mathrm{Hom}_{\wedge\text{-alg}}(R_{0},k[\epsilon])) \\[2ex]
&= \mathrm{Ker}(\mathrm{Hom}_{\wedge\text{-alg}}(R_{1},k[\epsilon]) \xrightarrow{\ s_{2}\ } \mathrm{Hom}_{\wedge\text{-alg}}(R_{0},k[\epsilon])) \qquad \text{are both finite-dimensional}
\end{aligned}
\end{equation*}
vector spaces over  $k$  , we must have the same properties on  $\underline{A}$  . Conversely assume (1) and (2) on  $\underline{A}$  . We claim that  $\underline{A}$  is prorepresentable by a normalized smooth cogroupoid. Since  $\underline{A}$  is homogeneous and  $[\underline{A}(k[\epsilon])]$  is a finite-dimensional vector space over  $k$  ,  $\underline{A}$  admits a minimally versal formal object  $\xi$  . Set  $P(\xi) = R_{0}$  and consider the functor  $\mathrm{Isom}_{\xi} : \underline{C}_{R_{0} \hat{\otimes}_{\wedge} R_{0}} \to (\mathrm{Ens})$  given by

\begin{equation*}
T \longmapsto \mathrm{Isom}_{\underline{A}(T)}(\xi_{1}(T),\xi_{2}(T))
\end{equation*}
where  $\xi_{i}(T) = \xi(f_{i})$  and  $f_{i} : R_{0} \to T$  denotes the composite maps

}{F7I-VI-039}\begin{center}
\begin{tikzcd}[column sep=large]
R_0 \arrow[r, shift left=1.2ex, "s_1"] \arrow[r, shift right=1.2ex, "s_2"'] \& R_0 \otimes_{\wedge} R_0 \arrow[r] \& T \ .
\end{tikzcd}
\end{center}









\SGAvar{This functor is pro-representable by 2.16, i.e. there exists  $R_{1} \in \mathrm{ob}(\hat{\underline{C}}_{R_{0} \hat{\otimes}_{\wedge} R_{0}})$  and an isomorphism  $\xi_{1}(R_{1}) \xrightarrow{\ \varphi\ } \xi_{2}(R_{1})$  such that  $\Phi : h_{R_{1}} \to \mathrm{Isom}_{\xi}$  is an isomorphism over  $\underline{C}_{R_{0} \hat{\otimes}_{\wedge} R_{0}}$  . Then  $(R_{0},R_{1},\ldots)$  defines a cogroupoid in  $\hat{\underline{C}}_{\wedge}$  and the functor $(\xi,\varphi) : \underline{h}_{(R_{0},R_{1},\ldots)} \to \underline{A}$  is a  $\underline{C}_{\wedge}$-equivalence. Since  $\xi$  is a minimally versal formal object of  $\underline{A}$  , the composit functor  $\underline{h}_{(R_{0},R_{0},\ldots)} \xrightarrow{\ \epsilon\ } \underline{h}_{(R_{0},R_{1},\ldots)} \xrightarrow{\ (\xi,\varphi)\ } \underline{A}$ is, by definition, minimally smooth and therefore so is  $\epsilon$  since  $(\xi,\varphi)$  is a $\underline{C}_{\wedge}$-equivalence i.e.,  $(R_{0},R_{1},\ldots)$  is a normalized smooth cogroupoid in  $\hat{\underline{C}}_{\wedge}$  .


}{This functor is pro-representable by 2.16, i.e. there exists  $R_{1} \in \mathrm{ob}(\hat{\underline{C}}_{R_{0} \hat{\otimes}_{\wedge} R_{0}})$  and an isomorphism  $\xi_{1}(R_{1}) \xrightarrow{\ \varphi\ } \xi_{2}(R_{1})$  such that  $\Phi : h_{R_{1}} \to \mathrm{Isom}_{\xi}$  is an isomorphism over  $\underline{C}_{R_{0} \hat{\otimes}_{\wedge} R_{0}}$  . Then  $(R_{0},R_{1},\ldots)$  defines a cogroupoid in  $\hat{\underline{C}}_{\wedge}$  and the functor $(\xi,\varphi) : \underline{h}_{(R_{0},R_{1},\ldots)} \to \underline{A}$  is a  $\underline{C}_{\wedge}$-equivalence. Since  $\xi$  is a minimally versal formal object of  $\underline{A}$  , the composite functor  $\underline{h}_{(R_{0},R_{0},\ldots)} \xrightarrow{\ \epsilon\ } \underline{h}_{(R_{0},R_{1},\ldots)} \xrightarrow{\ (\xi,\varphi)\ } \underline{A}$ is, by definition, minimally smooth and therefore so is  $\epsilon$  since  $(\xi,\varphi)$  is a $\underline{C}_{\wedge}$-equivalence i.e.,  $(R_{0},R_{1},\ldots)$  is a normalized smooth cogroupoid in  $\hat{\underline{C}}_{\wedge}$  .


}{F7I-VI-040}\SGApage{I}{69}

\underline{Corollary} 2.18.  \underline{A functor}  $F : \underline{C}_\Lambda \longrightarrow$ (Ens)  \underline{is pro-representable if and only if it is homogeneous and}  $\dim F(k[\epsilon]) < \infty$  .

\underline{Corollary} 2.19.  \underline{Every smooth cogroupoid in}  $\hat{\underline{C}}_\Lambda$  \underline{can be normalized up to isomorphism.}

Given a cogroupoid  $(R_0,R_1,\dots)$  in  $\hat{\underline{C}}_\Lambda$ , we set  $\bar{R}_1 = R_1/J$  where  $J$  is the ideal in  $R_1$  generated by  $\{s_1(r) - s_2(r) \mid r \in R_0\}$ . Then  $(R_0,\bar{R}_1,\dots)$  defines a formal group scheme over  $R_0$  with respect to the induced composition rule. Now let  $\underline{A}$  be a cofibered groupoid over  $\underline{C}_\Lambda$  prorepresented by a normalized smooth cogroupoid  $(R_0,R_1,\dots)$  via  $(\xi,\varphi)$ . Let  $\underline{C}_k \hookrightarrow \underline{C}_{R_0} \hookrightarrow \underline{C}_{R_0 \hat{\otimes}_\Lambda R_0}$  be the embeddings corresponding to the surjective maps  $R_0 \hat{\otimes}_\Lambda R_0 \longrightarrow R_0 \longrightarrow k$ , and let  $\mathrm{Aut}_\xi$ , $\mathrm{Aut}^o$  be the restriction functors of  $\mathrm{Isom}_\xi$  on  $\underline{C}_{R_0}$ , $\underline{C}_k$  respectively. In other words,  $\mathrm{Aut}_\xi(S)$ $=$ the automorphism group of  $\xi(S)$  for every  $S \in \mathrm{ob}(\underline{C}_{R_0})$ , and  $\mathrm{Aut}^o(S)$ $=$ the automorphism group of  $1_S$  for every  $S \in \mathrm{ob}(\underline{C}_k)$ , where  $1_S$ $=$ the direct image of  $1 \in \underline{A}(k)$  under the map  $k \to S$ . Since  $R_1$  prorepresents the functor  $\mathrm{Isom}_\xi$  on  $\underline{C}_{R_0 \hat{\otimes}_\Lambda R_0}$ , it follows that  $\bar{R}_1$  prorepresents the functor  $\mathrm{Aut}_\xi$  on  $\underline{C}_{R_0}$ , and  $k \otimes_{R_0} \bar{R}_1$  prorepresents  $\mathrm{Aut}^o$  on  $\underline{C}_k$ . Though  $\underline{A}$  is prorepresented by a smooth cogroupoid  $(R_0,R_1,\dots)$  , neither  $(R_0,\bar{R}_1,\dots)$  nor  $(k,\ k \otimes_{R_0} \bar{R}_1,\dots)$  are in general smooth. Indeed we have

\SGAvar{\underline{Corollary} 2.20.  \underline{Let}  $\underline{A}$  \underline{be a homogeneous cogroupoid over}  $\underline{C}_\Lambda$  \underline{prorepresented by a normalized smooth cogroupoid}  $(R_0,R_1,\dots)$  \underline{via}  $(\xi,\varphi)$  .

}{\underline{Corollary} 2.20.  \underline{Let}  $\underline{A}$  \underline{be a homogeneous cofibered groupoid over}  $\underline{C}_\Lambda$  \underline{prorepresented by a normalized smooth cogroupoid}  $(R_0,R_1,\dots)$  \underline{via}  $(\xi,\varphi)$  .

}{F7I-VI-041}(a)  $(R_0,\bar{R}_1,\dots)$  \underline{and}  $(k,\ k \otimes_{R_0} \bar{R}_1,\dots)$  \underline{prorepresent the functors}  $\mathrm{Aut}_\xi$  \underline{and}  $\mathrm{Aut}^o$  \underline{respectively.}

(b)  \underline{The following statements are equivalent}

(i)   \underline{the formal group scheme}  $\underline{h}_{(R_0,\bar{R}_1,\dots)}$  \underline{is formally smooth over}  $R_0$  ;

(ii)  \underline{for every surjective arrow}  $a' \to a$  \underline{in}  $\underline{A}$  ,
$\mathrm{Aut}_{P(a')}(a') \longrightarrow \mathrm{Aut}_{P(a)}(a)$  \underline{is surjective},

(iii)  $R_0$  \underline{prorepresents the functor}  $[\underline{A}]$  \underline{over}  $\underline{C}_\Lambda$  .


\SGApage{I}{70}

\underline{Proof}:  We have already noted (a). As for (b),  (i) $\Leftrightarrow$ (ii)  is clear whereas (ii) $\Leftrightarrow$ (iii)  follows from 2.7 and 2.15, since  $\tilde{\xi} : h_{R_0} \longrightarrow [\underline{A}]$  is minimally smooth.

\SGAvar{\underline{Remark} 2.21.  Let  $\underline{A}$  be a cofibered groupoid over  $\underline{C}_\Lambda$  , and let  $\underline{A}|\underline{C}_k$  be the pull-back of  $p : \underline{A} \longrightarrow \underline{C}_\Lambda$  under  $\underline{C}_k \hookrightarrow \underline{C}_\Lambda$ . Then it is trivial to see that if  $\underline{A}$  is prorepresented by  $(R_0,R_1,\dots)$  via  $(\xi,\varphi)$  , then  $\underline{A}|\underline{C}_k$  is prorepresented by  $(k \otimes_\Lambda R_0\ ,\ k \otimes_\Lambda R_1,\dots)$  via  $(k \otimes_\Lambda \xi\ ,\ k \otimes_\Lambda \varphi)$  where  $k \otimes_\Lambda \xi$ $=$ its direct image of  $\xi$  under the map  $R_0 \longrightarrow k \otimes_\Lambda R_0$ ,  $k \otimes_\Lambda \varphi$ $=$ the image of  $\varphi$  under  $\mathrm{Isom}_{\underline{A}(R_0)}(\xi(s_1),\ \xi(s_2)) \longrightarrow \mathrm{Isom}_{\underline{A}(k \otimes_\Lambda R_1)}(k \otimes_\Lambda \xi(s_1),\ k \otimes_\Lambda \xi(s_2))$ . If  $(R_0,R_1,\dots)$  is normalized, then so is  $(k \otimes_\Lambda R_0,\ k \otimes_\Lambda R_1,\dots)$  , and the relative dimension of  $R_1$  over  $R_0$  is equal to the relative dimension of  $k \otimes_\Lambda R_1$  over  $k \otimes_\Lambda R_0$  since  $R_1$  is smooth over  $R_0$  .

\underline{Remark} 2.22.  Let  $\underline{A}$  be prorepresented by smooth cogroupoid  $(R_0,R_1,\dots)$  and  $(R'_0,\ R'_1,\dots)$  where we assume that  $(R_0,R_1,\dots)$  is normalized. We then obtain a homomorphism  $(i_0,i_1) : (R_0,R_1,\dots) \longrightarrow (R'_0,\ R'_1,\dots)$  . Now  $R'_0$  is a formal power-series ring over  $R_0$  by 2.10(b), and the homomorphism  $(i_0,i_1)$  induces an isomorphism  $R'_0 \otimes_{R_0} (R_0,R_1,\dots) \overset{\sim}{\longrightarrow} (R'_0,\ R'_1,\dots)$  .


}{\underline{Remark} 2.21.  Let  $\underline{A}$  be a cofibered groupoid over  $\underline{C}_\Lambda$  , and let  $\underline{A}|\underline{C}_k$  be the pull-back of  $p : \underline{A} \longrightarrow \underline{C}_\Lambda$  under  $\underline{C}_k \hookrightarrow \underline{C}_\Lambda$ . Then it is trivial to see that if  $\underline{A}$  is prorepresented by  $(R_0,R_1,\dots)$  via  $(\xi,\varphi)$  , then  $\underline{A}|\underline{C}_k$  is prorepresented by  $(k \otimes_\Lambda R_0\ ,\ k \otimes_\Lambda R_1,\dots)$  via  $(k \otimes_\Lambda \xi\ ,\ k \otimes_\Lambda \varphi)$  where  $k \otimes_\Lambda \xi$ $=$ the direct image of  $\xi$  under the map  $R_0 \longrightarrow k \otimes_\Lambda R_0$ ,  $k \otimes_\Lambda \varphi$ $=$ the image of  $\varphi$  under  $\mathrm{Isom}_{\underline{A}(R_1)}(\xi(s_1),\ \xi(s_2)) \longrightarrow \mathrm{Isom}_{\underline{A}(k \otimes_\Lambda R_1)}(k \otimes_\Lambda \xi(s_1),\ k \otimes_\Lambda \xi(s_2))$ . If  $(R_0,R_1,\dots)$  is normalized, then so is  $(k \otimes_\Lambda R_0,\ k \otimes_\Lambda R_1,\dots)$  , and the relative dimension of  $R_1$  over  $R_0$  is equal to the relative dimension of  $k \otimes_\Lambda R_1$  over  $k \otimes_\Lambda R_0$  since  $R_1$  is smooth over  $R_0$  .

\underline{Remark} 2.22.  Let  $\underline{A}$  be prorepresented by smooth cogroupoids  $(R_0,R_1,\dots)$  and  $(R'_0,\ R'_1,\dots)$  where we assume that  $(R_0,R_1,\dots)$  is normalized. We then obtain a homomorphism  $(i_0,i_1) : (R_0,R_1,\dots) \longrightarrow (R'_0,\ R'_1,\dots)$  . Now  $R'_0$  is a formal power-series ring over  $R_0$  by 2.10(b), and the homomorphism  $(i_0,i_1)$  induces an isomorphism  $R'_0 \otimes_{R_0} (R_0,R_1,\dots) \overset{\sim}{\longrightarrow} (R'_0,\ R'_1,\dots)$  .


}{F7I-VI-042}3.  \underline{The coherent sheaves}  $\underline{D}^i_X$

One of the nice functors which is usually not pro-representable but nevertheless admits a formal versal object is the deformation functor. To this end we recall a few basic facts on commutative ring extensions and its obstruction theory. This section contains nothing new (possibly except in its systematic approach via Grothendieck cohomology on site). A purpose of this section is to make this expose as much self-contained. (Indeed we prove here more than we need later). The basic facts are stated in theorem 3.12. Throughout rings are meant to be commutative rings with unit.


\SGApage{I}{71}

Given a ring  $R$  we denote by  $\underline{C}_R$  the category of $R$-algebras. For each object  $A$  in  $\underline{C}_R$  we set

Cov($A$) $=$ the set of all surjective arrows  $\varphi : A' \to A$  in  $\underline{C}_R$  with target  $A$ .

It clearly defines a pretopology and in turn we obtain a site  $\underline{C}_R$ . For each object  $A$  in  $\underline{C}_R$ , we shall denote by  $\underline{C}_{A|R}$  the category of arrows in  $\underline{C}_R$  with target  $A$ . We may and shall provide the category  $\underline{C}_{A|R}$  with the topology induced from the forgetful functor  $j_A : \underline{C}_{A|R} \to \underline{C}_R$ . We note that the category  $\underline{C}_{A|R}$  has a final object  $(1_A : A \to A)$  and an initial object  $(i_A : R \to A)$ , and admits a fibre-product as well as coproduct. Henceforth we shall denote the object  $(B \to A)$  in  $\underline{C}_{A|R}$ , by abuse of notation, simply by  $B$  in case when there is no possible confusion. Given an abelian category  $\underline{C}$ , a  $\underline{C}$-valued sheaf on  $\underline{C}_{A|R}$  is, by definition, a functor  $F : \underline{C}^{O}_{A|R} \to \underline{C}$  such that, for any surjective arrow  $B' \to B$  in  $\underline{C}_{A|R}$ ,

\begin{equation*}
0 \to F(B) \to F(B') \rightrightarrows F(B' \underset{B}{\times} B')
\end{equation*}
is exact. The  \underline{additive}  presheaves are those which preserve coproduct i.e.,  $F(B_1 \underset{R}{\otimes} B_2) = F(B_1) \oplus F(B_2)$ . For example, an $A$-module  $M$  defines the functor

\begin{equation*}
\mathrm{Der}_R(-,M): \underline{C}^{O}_{A|R} \to (A\text{-mod})
\end{equation*}
given by  $B \to \mathrm{Der}_R(B,M) =$ $R$-linear derivations of  $B$  with values in  $M$  regarded as $B$-module via the structural map  $B \to A$ , and it is clearly an additive sheaf on  $\underline{C}^{O}_{A|R}$ , where ($A$-mod) denotes the category of $A$-modules. A presheaf  $F$  on  $\underline{C}_{A|R}$  with values in  ($A$-mod)  is called "finite type" if  $F(B)$  is an $A$-module of finite type whenever  $B$  is an $R$-algebra of finite type. For example, if  $A$  is noetherian and  $M$  is an $A$-module of finite type, then  $\mathrm{Der}_R(-,M)$  is an additive sheaf of finite type on the site  $\underline{C}_{A|R}$ .

A few more words on notations and basic facts in homological algebra:

3.1(a)  Throughout the rest of this section, we shall denote by  $\widetilde{\underline{C}}_{A|R}$  (or  $\overset{\vee}{\underline{C}}_{A|R}$ )  the category of sheaves (or presheaves) with values in ($A$-mod). We have the (left exact) injection functor  $i_{A|R}: \widetilde{\underline{C}}_{A|R} \to \overset{\vee}{\underline{C}}_{A|R}$ , and the associated-sheaf (exact)
\SGApage{I}{72}
functor  $G_{A|R}: \overset{\vee}{\underline{C}}_{A|R} \to \widetilde{\underline{C}}_{A|R}$  which is adjoint to the functor  $i_{A|R}$ . An object  $B \in ob(\underline{C}_{A|R})$  induces a continuous functor  $j : \underline{C}_{B|R} \to \underline{C}_{A|R}$  and in turn a commutative diagram

\begin{equation}\tag{*}
\begin{tikzcd}[column sep=huge, row sep=huge]
\widetilde{\underline{C}}_{A|R} \arrow[r, "\widetilde{j}"] \arrow[d, "i_{A|R}"'] \& \widetilde{\underline{C}}_{B|R} \arrow[d] \\
\overset{\vee}{\underline{C}}_{A|R} \arrow[r, "j^{\vee}"'] \& \overset{\vee}{\underline{C}}_{B|R}
\end{tikzcd}
\end{equation}
















where the restriction functor  $\widetilde{j}$  (as well as  $j^{\vee}$ )  is an exact functor, since the site  $\underline{C}_{B|R}$  is nothing but the category of arrows in  $\underline{C}_{A|R}$  having target  $B$ , with the induced topology. A ring homomorphism  $S \to R$  induces a continuous functor  $\beta : \underline{C}_{A|R} \to \underline{C}_{A|S}$ , and in turn a commutative diagram

\begin{equation}\tag{**}
\begin{tikzcd}[column sep=huge, row sep=huge]
\widetilde{\underline{C}}_{A|S} \arrow[r, "\widetilde{\beta}"] \arrow[d, "i_{A|S}"'] \& \widetilde{\underline{C}}_{A|R} \arrow[d] \\
\overset{\vee}{\underline{C}}_{A|S} \arrow[r, "\beta^{\vee}"'] \& \overset{\vee}{\underline{C}}_{A|R}
\end{tikzcd}
\end{equation}















Note that  $\widetilde{\beta}$  is not necessarily an exact functor.

3.1(b)  The right derived functor of the injection functor  $i_{A|R}: \widetilde{\underline{C}}_{A|R} \to \overset{\vee}{\underline{C}}_{A|R}$  will be denoted by  $\underline{H}^{\circ}_{A|R}$ . For each sheaf  $F$  on  $\underline{C}_{A|R}$ , we set

\SGAvar{\begin{equation*}
H^{\circ}(A|R,F) = [H^{\circ}_{A|R}(F)](A)
\end{equation*}
For any  $B' \to B$  in  Cov($B$) , let us define the functor  $H^{\circ}((B' \to B),-) : \overset{\vee}{\underline{C}}_{A|R} \to (A\text{-mod})$  by  $H^{\circ}((B'\to B),F) = Ker(F(B') \rightrightarrows F(B'\underset{B}{\times} B'))$ . The right derived functors of  $H^{\circ}((B' \to B),-)$  will be denoted by  $H^{\circ}((B'\to B),-)$ . One knows that the  $H^{\circ}((B' \to B),F)$  are the cohomology groups of the (Čech)
\SGApage{I}{73}
semi-simplicial module

}{\begin{equation*}
H^{\circ}(A|R,F) = [H^{\circ}_{A|R}(F)](A)
\end{equation*}
For any  $B' \to B$  in  Cov($B$) , let us define the functor  $H^{0}((B' \to B),-) : \overset{\vee}{\underline{C}}_{A|R} \to (A\text{-mod})$  by  $H^{0}((B'\to B),F) = Ker(F(B') \rightrightarrows F(B'\underset{B}{\times} B'))$ . The right derived functors of  $H^{0}((B' \to B),-)$  will be denoted by  $H^{\circ}((B'\to B),-)$ . One knows that the  $H^{\circ}((B' \to B),F)$  are the cohomology groups of the (Čech)
\SGApage{I}{73}
semi-simplicial module

}{F7I-VI-043}\begin{center}
\begin{tikzcd}[column sep=large]
F(B') \arrow[r, shift left=2] \arrow[r, shift right=2] \& F(B'\times_B B') \arrow[r, shift left=3] \arrow[r] \arrow[r, shift right=3] \& F(B'\ \times_B B'\ \times_B B') \arrow[r, shift left=4.5] \arrow[r, shift left=1.5] \arrow[r, shift right=1.5] \arrow[r, shift right=4.5] \& \cdots
\end{tikzcd}
\end{center}








\SGAvar{The right derived functor of  $\check{\underline{H}}^{0}_{A|R} : \underline{C}_{A|R} \to \underline{C}_{A|R}$  will be denoted by  $\check{\underline{H}}^{\bullet}_{A|R}$ ,

}{The right derived functor of  $\check{\underline{H}}^{0}_{A|R} : \overset{\vee}{\underline{C}}_{A|R} \to \overset{\vee}{\underline{C}}_{A|R}$  will be denoted by  $\check{\underline{H}}^{\bullet}_{A|R}$ ,

}{F7I-VI-044}where

\begin{equation*}
[\check{\underline{H}}^{0}_{A|R}(F)](B) \;=\; \lim_{\{B' \to B\}\, \in\, \mathrm{Cov}(B)} H^{0}(\{B' \to B\},F) \ .
\end{equation*}
For each presheaf  $F$  on  $\underline{C}_{A|R}$  , we set

\begin{equation*}
\check{H}^{\bullet}(A|R,F) \;=\; [\check{H}^{\bullet}_{A|R}(F)](A) \;=\; \lim_{\{A' \to A\}\, \in\, \mathrm{Cov}(A)} H^{\bullet}(\{A' \to A\},F)
\end{equation*}
for any sheaf  $F$  on  $\underline{C}_{A|R}$  , we have the usual spectral sequence

\begin{equation*}
\check{H}^{p}(A|R,\underline{H}^{q}_{A|R}(F)) \;\Rightarrow\; H^{p+q}(A|R,F)
\end{equation*}
where we note that  $\check{H}^{0}(A|R,\underline{H}^{q}(F)) = 0$  for all  $q > 0$ .

In particular

\SGAvar{\begin{equation*}
\left\{
\begin{array}{l}
H^{1}(A|R,F) \;=\; H^{1}(A|R,F) \\[2ex]
0 \to \check{H}^{2}(A|R,F) \to H^{2}(A|R,F) \to \check{H}^{1}(A|R,\underline{H}^{1}_{A\ R}(F)) \to \check{H}^{3}(A|R,F) \quad \text{is exact.}
\end{array}
\right.
\end{equation*}
3.1(c)   If  $B$  is an object in  $\underline{C}_{A|R}$  , the commutative diagram (*) in 2.1(a)

}{\begin{equation*}
\left\{
\begin{array}{l}
\check{H}^{1}(A|R,F) \;=\; H^{1}(A|R,F) \\[2ex]
0 \to \check{H}^{2}(A|R,F) \to H^{2}(A|R,F) \to \check{H}^{1}(A|R,\underline{H}^{1}_{A|R}(F)) \to \check{H}^{3}(A|R,F) \quad \text{is exact.}
\end{array}
\right.
\end{equation*}
3.1(c)   If  $B$  is an object in  $\underline{C}_{A|R}$  , the commutative diagram (*) in 3.1(a)

}{F7I-VI-045}together with the fact that  $\tilde{j}$  as well as  $j^{\vee}$  are exact functors yields the

canonical isomorphism

\begin{equation*}
H^{\bullet}(B|R,\tilde{j}F) \;=\; [\underline{H}^{\bullet}_{A|R}(F)](B)
\end{equation*}
On the other hand, a ring homomorphism  $S → R$  induces the commutative diagram (**)

\SGAvar{in 2.1(a). Now  $\tilde{\beta}$  is not necessarily an exact functor. However,  $β$  preserves a

}{in 3.1(a). Now  $\tilde{\beta}$  is not necessarily an exact functor. However,  $β$  preserves a

}{F7I-VI-046}fibre-product, and thus we obtain a spectral sequence

\begin{equation*}
\underline{H}^{p}_{A|R}(R^{q}\tilde{\beta}F) \;\Rightarrow\; \check{\beta}\underline{H}^{p+q}_{A|S}(F) \qquad \text{i.e.,} \qquad H^{p}(A|R,R^{q}\tilde{\beta}F) \;\Rightarrow\; H^{p+q}(A|S,F)
\end{equation*}
Since we have the canonical isomorphism  $\tilde{\beta} = G_{A|R}\circ\check{\beta}\circ i_{A|S}$  and  $\beta^{\#} = G_{A|R}\circ\check{\beta}$  is an
\SGApage{I}{74}
exact functor, we have  $R^{\bullet}\tilde{\beta}F = \beta^{\#}\underline{H}^{\bullet}_{A|S}(F)$  . Thus the spectral sequence above can

now be written as

\begin{equation*}
H^{p}(A|R,\beta^{\#}\underline{H}^{q}_{A|S}(F)) \;\Rightarrow\; H^{p+q}(A|S,F) \ .
\end{equation*}
\underline{Lemma} 3.2.  \underline{Given}  $B \in ob(\underline{C}_{A|R})$  , \underline{choose a surjective arrow}  $P → B$  in  $\underline{C}_{B|R}$  ,

\underline{where}  $P$  \underline{is a polynomial algebra over}  $R$  . \underline{Then for any}  $F \in ob(\check{\underline{C}}_{A|R})$  , \underline{the}

\underline{natural map}  $H^{\bullet}(\{P \to B\},F) \to \check{\underline{H}}^{\bullet}(F)\ (B)$  \underline{is an isomorphism.}

\SGAvar{\underline{Proof}:  It suffice to see that the natural map

}{\underline{Proof}:  It suffices to see that the natural map

}{F7I-VI-047}\begin{equation*}
H^{0}(\{P \to B\},F) \to \check{\underline{H}}^{0}(F)\ (B) \;=\; \varinjlim_{\{B' \to B\}\, \in\, \mathrm{Cov}(B)} H^{0}(\{B' \to B\},F)
\end{equation*}
is an isomorphism, which is clear since, for any  $\{B' \to B\} \in \mathrm{Cov}(B)$  , there exists

an arrow  $\{P \to B\} \to \{B' \to B\}$  .


\underline{Corollary} 3.3.

(a)  \underline{Assume that}  $A$  \underline{is noetherian}. \underline{If}  $F \in ob(\widetilde{\underline{C}}_{A|R})$  \underline{is of finite type,}

\underline{then}  $\underline{H}^{n}(F)$  \underline{is of finite type for all}  $n$  .

(b)  \underline{For any polynomial algebra}  $P$  \underline{over}  $R$  , \underline{we have}  $H^{n}(P|R,F) = 0$

\underline{for all}  $n > 0$  \underline{and for all}  $F \in ob(\widetilde{\underline{C}}_{P|R})$  .

(c)  \underline{A diagram}  $0 \to F' \to F \to F'' \to 0$  \underline{in}  $\widetilde{\underline{C}}_{A|R}$  \underline{is exact if and only if}

$0 \to F'(P) \to F(P) \to F''(P) \to 0$  \underline{is exact for every}  $P \in ob(\underline{C}_{A|R})$  \underline{which is a polynomial algebra over}  $R$  .

\underline{Proof}:

(a)  In view of the spectral sequence

\begin{equation*}
\check{H}^{p}(B|R,\underline{H}^{q}(F)) \;\Rightarrow\; H^{p+q}(B|R,F)
\end{equation*}
together with the fact that  $\check{H}^{0}(\underline{H}^{q}(F)) = 0$  for all  $q > 0$  , it suffices to

\SGAvar{show that  $\underline{H}^{n}(G)$  is of finite type for all  $n$  and for all presheaf  $G$  of finite

}{show that  $\underline{H}^{n}(G)$  is of finite type for all  $n$  and for all presheaves  $G$  of finite

}{F7I-VI-048}type. However  $\check{H}^{\bullet}(B|R,G) = H^{\bullet}(\{P \to B\},G)$  , where  $P$  is a polynomial algebra over

$R$  and  $\{P \to B\} \in \mathrm{Cov}(B)$  , and furthermore  $H^{\bullet}(\{P \to B\},G)$  is simply the
\SGApage{I}{75}
cohomology gotten from the complex

\begin{multline*}
\dots\dots\dots\dots\dots G(P \underset{B}{\times} P \underset{B}{\times} P) \;\substack{\longleftarrow\\ \longleftarrow\\ \longleftarrow}\; G(P \underset{B}{\times} P) \substack{\longleftarrow\\ \longleftarrow}\; G(P) \ .
\end{multline*}
If  $B$  is an R-algebra of finite type, then we may take  $P$  to be of finite type over  $R$  so that  $P\underset{B}{\times}\ ....\ \underset{B}{\times}P$  are all R-algebras of finite type. If  $G$  is of finite type, then  $G(P\underset{B}{\times}\ ....\ \underset{B}{\times}P)$  are all A-modules of finite type, so that  $H^{\bullet}(\{P \to B\},G)$  is an A-module of finite type since  $A$  is noetherian.

\SGAvar{(b)  Let  $P$  be a polynomial algebra over  $R$ . If  $B \to P$  is any surjective R-algebra map, then it amits a section and hence  $H^{i}(\{B \to P\},G) = 0$  for all  $i > 0$ , so that  $\check{H}^{i}(P|R,\ G) = 0$  for all  $i > 0$  and for all presheaves  $G$  on  $\underline{C}_{P|R}$ . It follows that the spectral sequence  $\check{H}^{p}(P|R,\ \underline{H}^{q}_{P|R}(F)) \Rightarrow H^{p+q}(P|R,\ F)$  degenerates and hence  $\check{H}^{0}(P|R,\ \underline{H}^{n}(F)) \simeq H^{n}(P|R,\ F)$  for all  $n$ . However  $\check{H}^{0}(P|R,\ \underline{H}^{n}(F)) = 0$  for all  $n > 0$  whenever  $F$  is a sheaf, and hence  $H^{n}(P|R,\ F) = 0$  for all  $n > 0$  and for all  $F \in ob(\widetilde{\underline{C}}_{P|R})$ .

}{(b)  Let  $P$  be a polynomial algebra over  $R$ . If  $B \to P$  is any surjective R-algebra map, then it admits a section and hence  $H^{i}(\{B \to P\},G) = 0$  for all  $i > 0$ , so that  $\check{H}^{i}(P|R,\ G) = 0$  for all  $i > 0$  and for all presheaves  $G$  on  $\underline{C}_{P|R}$ . It follows that the spectral sequence  $\check{H}^{p}(P|R,\ \underline{H}^{q}_{P|R}(F)) \Rightarrow H^{p+q}(P|R,\ F)$  degenerates and hence  $\check{H}^{0}(P|R,\ \underline{H}^{n}(F)) \simeq H^{n}(P|R,\ F)$  for all  $n$ . However  $\check{H}^{0}(P|R,\ \underline{H}^{n}(F)) = 0$  for all  $n > 0$  whenever  $F$  is a sheaf, and hence  $H^{n}(P|R,\ F) = 0$  for all  $n > 0$  and for all  $F \in ob(\widetilde{\underline{C}}_{P|R})$ .

}{F7I-VI-049}(c)  If  $0 \to F' \to F \to F'' \to 0$  is an exact sequence in  $\widetilde{\underline{C}}_{A|R}$ , then

\SGAvar{\begin{equation*}
0 \to F'(B) \to F(B) \to F''(B) \to H'(B|R,\ F') \to H'(B|R,\ F) \to \ ....
\end{equation*}
is exact. If  $P \in ob(\underline{C}_{A|R})$  is a polynomial algebra over  $R$ , then  $H^{1}(P|R,\ F') = 0$  by 3.3(b) and hence

\begin{equation*}
0 \to F' \to F \to F'' \to 0
\end{equation*}
is exact. As for the other direction, let  $\underline{P}_{A|R}$  be the full subcategory of  $\underline{C}_{A|R}$  in which

}{\begin{equation*}
0 \to F'(B) \to F(B) \to F''(B) \to H^{1}(B|R,\ F') \to H^{1}(B|R,\ F) \to \ ....
\end{equation*}
is exact. If  $P \in ob(\underline{C}_{A|R})$  is a polynomial algebra over  $R$ , then  $H^{1}(P|R,\ F') = 0$  by 3.3(b) and hence

\begin{equation*}
0 \to F'(P) \to F(P) \to F''(P) \to 0
\end{equation*}
is exact. As for the other direction, let  $\underline{P}_{A|R}$  be the full subcategory of  $\underline{C}_{A|R}$  in which

}{F7I-VI-050}\begin{equation*}
ob(\underline{P}_{A|R}) = \{P \in ob(\underline{C}_{A|R})|\ \ P \text{ is a polynomial algebra over } R\} \ ,
\end{equation*}
provided with the induced topology, and let  $p : \underline{P}_{A|R} \to \underline{C}_{A|R}$  be the inclusion functor. Since every object in  $\underline{C}_{A|R}$  can be covered by an object in  $\underline{P}_{A|R}$ , it follows from a comparison lemma (          ) that  $\widetilde{p} : \widetilde{\underline{C}}_{A|R} \to \widetilde{\underline{P}}_{A|R}$  is an equivalence of categories. Consequently, if  $0 \to F'(P) \to F(P) \to F''(P) \to 0$  is exact for every  $P \in ob(\underline{P}_{A|R})$ , then  $0 \to \widetilde{p}F' \to \widetilde{p}F \to \widetilde{p}F'' \to 0$  is exact in  $\widetilde{\underline{P}}_{A|R}$  and hence  $0 \to F' \to F \to F'' \to 0$  is exact in  $\widetilde{\underline{C}}_{A|R}$  .


\SGApage{I}{76}

\SGAvar{\underline{Remark} 3.4.  Every surjective arrow in  $\underline{P}_{A|R}$  admits a section. It follows immediately that every presheaf on  $\underline{P}_{A|R}$  is actually a sheaf i.e.,  $\widetilde{\underline{P}}_{A|R} = \underline{P}_{A|R}$ . Consequently we find that the restriction functor  $\widetilde{\underline{C}}_{A|R} \to \check{\underline{P}}_{A|R}$  is an equivalence of categories.

}{\underline{Remark} 3.4.  Every surjective arrow in  $\underline{P}_{A|R}$  admits a section. It follows immediately that every presheaf on  $\underline{P}_{A|R}$  is actually a sheaf i.e.,  $\widetilde{\underline{P}}_{A|R} = \check{\underline{P}}_{A|R}$ . Consequently we find that the restriction functor  $\widetilde{\underline{C}}_{A|R} \to \check{\underline{P}}_{A|R}$  is an equivalence of categories.

}{F7I-VI-051}\underline{Lemma} 3.5.  \underline{Consider a diagram of R-algebras}

\begin{center}
\begin{tikzcd}
 \& C \arrow[d] \\
B' \arrow[r, "\text{surjective}"'] \& B
\end{tikzcd}
\end{center}

\underline{and let}  $S$  \underline{be an R-algebra.}

(i)  \underline{If}  $\mathrm{Tor}^{R}_{1}(B,S) = 0$  , \underline{then the canonical map}

\begin{equation*}
(B'\underset{B}{\times} C) \underset{R}{\otimes} S \to (B'\underset{R}{\otimes} S) \underset{(B \underset{R}{\otimes} S)}{\times} (C \underset{R}{\otimes} S)
\end{equation*}
\underline{is an isomorphism.}

\SGAvar{(ii)  \underline{Assume that}  $B'$  \underline{is R-flat.} \underline{If}  $\mathrm{Tor}^{R}_{i}(C,S) = 0$  \underline{for all}  $0 < i < n$  \underline{and}  $\mathrm{Tor}^{R}_{i}(B,S) = 0$  \underline{for all}  $0 < i < n$  , \underline{then}  $\mathrm{Tor}^{R}_{i}(B'\underset{B}{\times}C,\ S) = 0$  \underline{for all}  $0 < i < n$  .

}{(ii)  \underline{Assume that}  $B'$  \underline{is R-flat.} \underline{If}  $\mathrm{Tor}^{R}_{i}(C,S) = 0$  \underline{for all}  $0 < i < n$  \underline{and}  $\mathrm{Tor}^{R}_{i}(B,S) = 0$  \underline{for all}  $0 < i \leq n$  , \underline{then}  $\mathrm{Tor}^{R}_{i}(B'\underset{B}{\times}C,\ S) = 0$  \underline{for all}  $0 < i < n$  .

}{F7I-VI-052}\underline{Proof}:  Let  $0 \to J \to B' \to B \to 0$  be exact.

\SGAvar{(i)  If  $\mathrm{Tor}^{R}_{i}(B,\ S) = 0$  , then

}{(i)  If  $\mathrm{Tor}^{R}_{1}(B,\ S) = 0$  , then

}{F7I-VI-053}\begin{equation*}
0 \to J\underset{R}{\otimes}S \to B'\underset{R}{\otimes}S \to B\underset{R}{\otimes}S \to 0
\end{equation*}
is exact so that

\begin{equation*}
0 \to J\underset{R}{\otimes}S \to (B'\underset{R}{\otimes}S) \underset{(B\underset{R}{\otimes}S)}{\times} (C\underset{R}{\otimes}S) \to C\underset{R}{\otimes}S \to 0
\end{equation*}
is exact. We thus obtain an exact commutative diagram

\begin{tikzcd}
 \& J \underset{R}{\otimes} S \arrow[r] \arrow[d, equal] \& (B'\underset{B}{\times}C)\underset{R}{\otimes}S \arrow[r] \arrow[d] \& C \underset{R}{\otimes} S \arrow[r] \arrow[d, equal] \& 0 \\
0 \arrow[r] \& J \underset{R}{\otimes} S \arrow[r] \& (B'\underset{R}{\otimes}S) \underset{(B\underset{R}{\otimes}S)}{\times} (C \underset{R}{\otimes}S) \arrow[r] \& C \underset{R}{\otimes} S \arrow[r] \& 0
\end{tikzcd}















\SGAvar{\SGApage{I}{77}
and therefore $(B'\otimes_B C)\otimes_R S \to (B'\otimes_R S)\ \times_{(B\otimes_R S)}\ (C \otimes_R S)$  is an isomorphism.

}{\SGApage{I}{77}
and therefore $(B'\times_B C)\otimes_R S \to (B'\otimes_R S)\ \times_{(B\otimes_R S)}\ (C \otimes_R S)$  is an isomorphism.

}{F7I-VI-054}(ii)  Since  $B'$  is R-flat and  $\mathrm{Tor}^R_i(B,\ S) = 0$  for all  $0 < i \leq n$   by hypothesis, we have  $\mathrm{Tor}^R_i(J,\ S) = 0$  for all  $0 < i < n$  . Then the exact sequence  $0 \to J \to B'\times_B C \to C \to 0$  together with the hypothesis that  $\mathrm{Tor}^R_i(C,\ S) = 0$  for all  $0 < i < n$  entails that  $\mathrm{Tor}^R_i(B'\times_B C,\ S) = 0$  for all  $0 < i < n$  .


\underline{Corollary} 3.6.

\SGAvar{(a)  \underline{Let}  $A, S$  \underline{be} R-\underline{algebras. If}  $\mathrm{Tor}^R_1(A,\ S) = 0$  , \underline{then the canonical map}  $\check{H}^{\cdot}(A\otimes_R S|S,\ G) \to \check{H}^{\cdot}(A|R,\ SG)$  \underline{is an isomorphism for any}  $G \in \mathrm{ob}(\check{\underline{C}}_{A\otimes_R S|S})$  \underline{where the functor}  $S : \underline{C}_{A|R} \to \underline{C}_{A\otimes_R S|S}$  \underline{is given by}  $X \longmapsto X\otimes_R S$  . \underline{In particular if}  $S$  \underline{is} R-\underline{flat, then the canonical map}  $\check{S}\check{H}^{\cdot}(G) \to \check{H}(\check{S}G)$  \underline{is an isomorphism for every}  $G \in \mathrm{ob}(\check{\underline{C}}_{A\otimes_R S|S})$  .

}{(a)  \underline{Let}  $A, S$  \underline{be} R-\underline{algebras. If}  $\mathrm{Tor}^R_1(A,\ S) = 0$  , \underline{then the canonical map}  $\check{H}^{\cdot}(A\otimes_R S|S,\ G) \to \check{H}^{\cdot}(A|R,\ \check{S}G)$  \underline{is an isomorphism for any}  $G \in \mathrm{ob}(\check{\underline{C}}_{A\otimes_R S|S})$  \underline{where the functor}  $S : \underline{C}_{A|R} \to \underline{C}_{A\otimes_R S|S}$  \underline{is given by}  $X \longmapsto X\otimes_R S$  . \underline{In particular if}  $S$  \underline{is} R-\underline{flat, then the canonical map}  $\check{S}\underline{\check{H}}^{\cdot}(G) \to \underline{\check{H}}^{\cdot}(\check{S}G)$  \underline{is an isomorphism for every}  $G \in \mathrm{ob}(\check{\underline{C}}_{A\otimes_R S|S})$  .

}{F7I-VI-055}(b)  \underline{Let}  $P \to B$  \underline{be a surjective} R-\underline{algebra map, where}  $P$  \underline{is a polynomial algebra over}  $R$  . \underline{If}  $\mathrm{Tor}^R_i(B,\ S) = 0$  \underline{for all}  $0 < i \leq n$  , \underline{then} $\mathrm{Tor}^R_i(P\times_B P\times_B \ldots \times_B P,\ S) = 0$  \underline{for all}  $0 < i < n$  .


\underline{Proof}:

(a)  Given  $B \in \mathrm{ob}(\underline{C}_{A|R})$  , choose  $\{P \to B\} \in \mathrm{Cov}(B)$  where  $P$  is a polynomial algebra over  $R$  . If  $\mathrm{Tor}^R_1(B,\ S) = 0$  , the natural map

\begin{equation*}
(\overbrace{P\times_B \ \ldots\ \times_B P}^{\,n\,})\otimes_R S \to (P\otimes_R S)\times_{(B\otimes_R S)} \ldots \times_{(B\otimes_R S)} (P\otimes_R S)
\end{equation*}
is an isomorphism for all  $n$  by 3.5(i), and hence  $\check{H}^{\cdot}(B|R,\ \check{S}G) = H^{\cdot}(\{P \to B\},\ \check{S}G) = H^{\cdot}(\{P\otimes_R S \to B\otimes_R S\},\ G) = \check{H}^{\cdot}(B\otimes_R S|S,\ G)$  . If  $S$  is R-flat, then  $\mathrm{Tor}^R_1(X,\ S) = 0$  for every  $X \in \mathrm{ob}(\underline{C}_{A|R})$  and hence $\check{H}^{\cdot}(X\otimes_R S|S,\ G) \to \check{H}^{\cdot}(X|R,\ \check{S}G)$  is an isomorphism for every  $X \in \mathrm{ob}(\underline{C}_{A|R})$  i.e. $\check{S}\underline{\check{H}}^{\cdot}(G) \to \underline{\check{H}}^{\cdot}(\check{S}G)$  is an isomorphism.

(b)  Set  $P^{(m)} = \overbrace{P\times_B \ \ldots\ \times_B P}^{\,m+1\,}$  . If  $m = 0$  there is nothing to prove since  $P$  is R-flat. Assume that  $\mathrm{Tor}^R_i(P^{(m-1)},\ S) = 0$  for all  $0 < i < n$  . Since  $P \to B$  is surjective and  $P^{(m)} = P\times_B P^{(m-1)}$  , it follows from 3.5(ii) that $\mathrm{Tor}^R_i(P^{(m)},\ S) = 0$  for all  $0 < i < n$  .


\SGApage{I}{78}


\underline{Proposition} 3.7.

(a)  \underline{Let}  $A, S$  \underline{be} R-\underline{algebras. If}  $\mathrm{Tor}^R_i(A,\ S) = 0$  \underline{for all}  $0 < i \leq n$  , \underline{then for any sheaf}  $F$  \underline{on}  $\underline{C}_{A\otimes_R S|S}$  , \underline{the canonical map}  $H^i(A\otimes_R S|S,\ F) \to H^i(A|R,\ \widetilde{S}F)$ \underline{is an isomorphism for all}  $i \leq n$  , \underline{where}  $S : \underline{C}_{A|R} \to \underline{C}_{A\otimes_R S|S}$  \underline{is given by} $X \longmapsto X\otimes_R S$  .

(b)  \underline{Let}  $P$  \underline{be a polynomial algebra over}  $R$  . \underline{Then for any additive sheaf}  $F$  \underline{on}  $\underline{C}_{P\otimes_R A|R}$  , \underline{the canonical map}  $H^n(P\otimes_R A|R,\ F) \to H^n(A|R,\ F)$  \underline{is an isomorphism for all}  $n > 0$  .

(c)  \underline{Let}  $S \to R \to A$  \underline{be ring homomorphisms. Then for any additive sheaf}  $F$  \underline{on}  $\underline{C}_{A|S}$  , \underline{we have an exact sequence}

{\small \begin{multline*}
0 \to H^0(A|R,\ F_R) \to H^0(A|S,\ F) \to H^0(R|S,\ F) \to H^1(A|R,\ F_R) \to H^1(A|S,\ F) \to H^1(R|S,\ F) \to \\ \ldots\ldots\ldots \to H^n(A|R,\ F_R) \to H^n(A|S,\ F) \to H^n(R|S,\ F) \to H^{n+1}(A|R,\ F_R) \to \ldots.
\end{multline*}}
\underline{where}  $F_R$  \underline{is a sheaf on}  $\underline{C}_{A|R}$ \underline{defined by}

\begin{equation*}
F_R(X) = \mathrm{Ker}(F(X) \to F(R))
\end{equation*}
\underline{for all}  $X \in \mathrm{ob}(\underline{C}_{A|R})$  \underline{via the forgetful functor}  $\underline{C}_{A|R} \to \underline{C}_{A|S}$  .


\underline{Proof}:

(a)  Consider the continuous functor

\SGAvar{\begin{equation*}
S : \underline{C}_{A|R} \to \underline{C}_{A\otimes_R S|S}
\end{equation*}
where  $S(X) = X\otimes_R S$  . Since the functor  $S$  sends a polynomial algebras over  $R$  to polynomial algebras over  $S$  ,  $\widetilde{S} : \widetilde{\underline{C}}_{A\otimes_R S|S} \to \widetilde{\underline{C}}_{A|R}$  is an exact functor by 3.3(c) and in turn we obtain the canonical map  $\check{S}\underline{H}^{\cdot}_{A\otimes_R S|S}(F) \to \underline{H}^{\cdot}_{A|R}(\widetilde{S}F)$  , and in turn we obtain the commutative diagram of spectral sequences:

}{\begin{equation*}
S : \underline{C}_{A|R} \to \underline{C}_{A\otimes_R S|S}
\end{equation*}
where  $S(X) = X\otimes_R S$  . Since the functor  $S$  sends polynomial algebras over  $R$  to polynomial algebras over  $S$  ,  $\widetilde{S} : \widetilde{\underline{C}}_{A\otimes_R S|S} \to \widetilde{\underline{C}}_{A|R}$  is an exact functor by 3.3(c) and in turn we obtain the canonical map  $\check{S}\underline{H}^{\cdot}_{A\otimes_R S|S}(F) \to \underline{H}^{\cdot}_{A|R}(\widetilde{S}F)$  , and in turn we obtain the commutative diagram of spectral sequences:

}{F7I-VI-056}\begin{center}
\begin{tikzcd}[column sep=huge, row sep=large]
E^{p,q}_S(S(X)) = \check{H}^p(X\otimes_R S|S,\ \underline{H}^q_{A\otimes_R S|S}(F)) \arrow[r, Rightarrow] \arrow[d] \& H^{p+q}(X\otimes_R S|S,\ F) \arrow[dd] \\
\check{H}^p(X|R,\ \check{S}\underline{H}^q_{A\otimes_R S|S}(F)) \arrow[d] \& \\
E^{p,q}_R(X) =\ \ \check{H}^p\ (X|R,\ \underline{H}^q_{A|R}(\widetilde{S}F)) \arrow[r, Rightarrow] \& H^{p+q}(X|R,\ \widetilde{S}F)
\end{tikzcd}
\end{center}




















\SGApage{I}{79}

\SGAvar{We note that  $E^{p,q}_{S}(S(X)) \to \check{H}^{p}(X|R,\ \check{S}\underline{H}^{q}_{A\underset{R}{\otimes}S|S}(F))$  is an isomorphism whenever $\operatorname{Tor}^{R}_{1}(X,\ S) = 0$  by 3.6(a). Assume now that $[\check{S}\underline{H}^{q}_{A\underset{R}{\otimes}S|S}(F)](X) \to [\underline{H}^{q}_{A|R}(\widetilde{S}F)](X)$  is an isomorphism for all  $q < n$  and for all $X \in \mathrm{ob}(\underline{C}_{A|R})$  such that  $\operatorname{Tor}^{R}_{i}(X,\ S) = 0$  for all  $0\ \ \ i\ \ \ n$ . It follows immediately from 3.6(b) that  $E^{p,q}_{S}(S(X)) \to E^{p,q}_{R}(X)$  is an isomorphism for all  $q < n$ and for all  $p$ , provided  $\operatorname{Tor}^{R}_{i}(X,\ S) = 0$  for all  $0 < i \leq n$ . Since $E^{0,n}_{S}(S(X)) \to E^{0,n}_{R}(X)$  is an isomorphism (since both sides are zero), we get that $H^{n}(X\underset{R}{\phantom{\otimes}}S|S,\ F) \to H^{n}(X|R,\ \widetilde{S}F)$  is an isomorphism whenever  $\operatorname{Tor}^{R}_{i}(X,\ S) = 0$  for all $0 < i \leq n$  i.e.,  $[\check{S}\underline{H}^{i}_{A\underset{R}{\otimes}S|S}(F)](X) \to [\underline{H}^{i}_{A|R}(\widetilde{S}F)](X)$  is an isomorphism for all $i \leq n$  whenever  $\operatorname{Tor}^{R}_{i}(X,\ S) = 0$  for all  $0 < i \leq n$ .

(b)  Let  $P$  be a polynomial algebra over  $R$  and  let $f : \underline{C}_{A|R} \to \underline{C}_{A\underset{R}{\otimes}P|R}$  be the continuous functor given by  $f(X) = X\underset{R}{\otimes}P$ . Let  $P' \to X$ be a surjective arrow in  $\underline{C}_{A|R}$  where  $P'$  is a polynomial algebra over  $R$ . Then $P'\underset{R}{\otimes}P \to X\underset{R}{\otimes}P$  is a surjective arrow in  $\underline{C}_{A\underset{R}{\otimes}P|R}$ , and  $P'\underset{R}{\otimes}P$  is again a polynomial algebra over  $R$ . Since  $P$  is flat over  $R$ , it follows from 3.5(i) that $\check{H}^{\cdot}(X|R,\ \check{f}G) = H^{\cdot}(\{P' \to X\},\ \check{f}G) = H\ (\{P'\underset{R}{\otimes}P \to X\underset{R}{\otimes}P\},\ G) = \check{H}^{\cdot}(X\underset{R}{\otimes}P|R,\ G)$ for every  $G \in \mathrm{ob}(\underline{\check{C}}_{A\underset{R}{\otimes}P|R})$ . It follows that  $\widetilde{f} : \underline{\widetilde{C}}_{A\underset{R}{\otimes}P|R} \to \underline{C}_{A|R}$  sends acyclic ones to acyclic ones. On the other hand,  $f$  sends polynomial algebras over  $R$  to polynomial algebras over  $R$ , and hence  $\widetilde{f} : \underline{\widetilde{C}}_{A\underset{R}{\otimes}P|R} \to\!\underline{\widetilde{C}}_{A|R}$  is an exact functor by 3.2(c). It follows that  $H^{\cdot}(A\underset{R}{\otimes}P|R,\ F) \to H^{\cdot}(A|R,\ \widetilde{f}F)$  is an isomorphism for every $F \in \mathrm{ob}(\underline{\widetilde{C}}_{A\underset{R}{\otimes}P|R})$ . Now for each  $X \in \mathrm{ob}(\underline{C}_{A|R})$ , we have canonical maps

}{We note that  $E^{p,q}_{S}(S(X)) \to \check{H}^{p}(X|R,\ \check{S}\underline{H}^{q}_{A\underset{R}{\otimes}S|S}(F))$  is an isomorphism whenever $\operatorname{Tor}^{R}_{1}(X,\ S) = 0$  by 3.6(a). Assume now that $[\check{S}\underline{H}^{q}_{A\underset{R}{\otimes}S|S}(F)](X) \to [\underline{H}^{q}_{A|R}(\widetilde{S}F)](X)$  is an isomorphism for all  $q < n$  and for all $X \in \mathrm{ob}(\underline{C}_{A|R})$  such that  $\operatorname{Tor}^{R}_{i}(X,\ S) = 0$  for all  $0 < i < n$ . It follows immediately from 3.6(b) that  $E^{p,q}_{S}(S(X)) \to E^{p,q}_{R}(X)$  is an isomorphism for all  $q < n$ and for all  $p$ , provided  $\operatorname{Tor}^{R}_{i}(X,\ S) = 0$  for all  $0 < i \leq n$ . Since $E^{0,n}_{S}(S(X)) \to E^{0,n}_{R}(X)$  is an isomorphism (since both sides are zero), we get that $H^{n}(X\underset{R}{\otimes}S|S,\ F) \to H^{n}(X|R,\ \widetilde{S}F)$  is an isomorphism whenever  $\operatorname{Tor}^{R}_{i}(X,\ S) = 0$  for all $0 < i \leq n$  i.e.,  $[\check{S}\underline{H}^{i}_{A\underset{R}{\otimes}S|S}(F)](X) \to [\underline{H}^{i}_{A|R}(\widetilde{S}F)](X)$  is an isomorphism for all $i \leq n$  whenever  $\operatorname{Tor}^{R}_{i}(X,\ S) = 0$  for all  $0 < i \leq n$ .

(b)  Let  $P$  be a polynomial algebra over  $R$  and  let $f : \underline{C}_{A|R} \to \underline{C}_{A\underset{R}{\otimes}P|R}$  be the continuous functor given by  $f(X) = X\underset{R}{\otimes}P$ . Let  $P' \to X$ be a surjective arrow in  $\underline{C}_{A|R}$  where  $P'$  is a polynomial algebra over  $R$ . Then $P'\underset{R}{\otimes}P \to X\underset{R}{\otimes}P$  is a surjective arrow in  $\underline{C}_{A\underset{R}{\otimes}P|R}$ , and  $P'\underset{R}{\otimes}P$  is again a polynomial algebra over  $R$ . Since  $P$  is flat over  $R$ , it follows from 3.5(i) that $\check{H}^{\cdot}(X|R,\ \check{f}G) = H^{\cdot}(\{P' \to X\},\ \check{f}G) = H^{\cdot}\ (\{P'\underset{R}{\otimes}P \to X\underset{R}{\otimes}P\},\ G) = \check{H}^{\cdot}(X\underset{R}{\otimes}P|R,\ G)$ for every  $G \in \mathrm{ob}(\underline{\check{C}}_{A\underset{R}{\otimes}P|R})$ . It follows that  $\widetilde{f} : \underline{\widetilde{C}}_{A\underset{R}{\otimes}P|R} \to \underline{\widetilde{C}}_{A|R}$  sends acyclic ones to acyclic ones. On the other hand,  $f$  sends polynomial algebras over  $R$  to polynomial algebras over  $R$ , and hence  $\widetilde{f} : \underline{\widetilde{C}}_{A\underset{R}{\otimes}P|R} \to\!\underline{\widetilde{C}}_{A|R}$  is an exact functor by 3.3(c). It follows that  $H^{\cdot}(A\underset{R}{\otimes}P|R,\ F) \to H^{\cdot}(A|R,\ \widetilde{f}F)$  is an isomorphism for every $F \in \mathrm{ob}(\underline{\widetilde{C}}_{A\underset{R}{\otimes}P|R})$ . Now for each  $X \in \mathrm{ob}(\underline{C}_{A|R})$ , we have canonical maps

}{F7I-VI-057}\begin{equation*}
(\widetilde{f}F)\ (X) = F(X\underset{R}{\otimes}P) \to F(X)\oplus F(P)\ \ ,
\end{equation*}
and hence a canonical map  $\widetilde{f}F \to F\oplus F(P)$  where  $F(P)$  denotes a constant sheaf (and  $F$  on the right hand side denotes, strictly speaking, the restriction of  $F$ on  $\underline{C}_{A|R}$ ) . If  $F$  is an additive sheaf, then  $(\widetilde{f}F)\ (X) = F(X\underset{R}{\otimes}P) \overset{\sim}{\to} F(X)\oplus F(P)$ is an isomorphism for every  $X \in \mathrm{ob}(\underline{C}_{A|R})$  and hence  $\widetilde{f}F \overset{\sim}{\to} F\oplus F(P)$  is an isomorphism. Therefore 


\SGApage{I}{80}

\begin{equation*}
 H^{n}(A|R,\ \widetilde{f}F) \overset{\sim}{\to} H^{n}(A|R,\ F\oplus F(P)) = H^{n}(A|R,\ F)
\end{equation*}
for all  $n > 0$  since  $F(P)$ , being a constant sheaf, is flask. It follows that $H^{n}(A\underset{R}{\otimes}P|R,\ F) \to H^{n}(A|R,\ F)$  is an isomorphism for all  $n > 0$ , provided  $F$  is an additive sheaf on  $\underline{C}_{A\underset{R}{\otimes}P|R}$ .

(c)  (communicated from Rinehart) Consider the forgetful functor $\beta : \underline{C}_{A|R} \to \underline{C}_{A|S}$ . We have a spectral sequence

\SGAvar{\begin{equation*}
E^{p,q} = H^{p}(A|R,\ \beta^{\#}\underline{H}^{q}_{A|S}(F))\ \Longrightarrow\ H^{p+q}(A|S,\ F)
\end{equation*}
where  $\beta^{\#} = G_{A|R}\circ\check{\beta}$ ,  $G_{A|R} : \underline{\check{C}}_{A|R} \to \underline{\widetilde{C}}_{A|R}$  is "associated-sheaf"  functor. Now let  $F$  be an additive sheaf. We claim that  $\beta^{\#}\underline{H}^{q}_{A|S}(F)$  is a constant sheaf for all  $q > 0$ . Indeed, if  $P$  is a polynomial algebra over  $R$ , then  $P = R\underset{S}{\otimes}P'$ for a polynomial algebra  $P'$  over  $S$  and hence

}{\begin{equation*}
E^{p,q} = H^{p}(A|R,\ \beta^{\#}\underline{H}^{q}_{A|S}(F))\ \Longrightarrow\ H^{p+q}(A|S,\ F)
\end{equation*}
where  $\beta^{\#} = G_{A|R}\circ\check{\beta}$ ,  $G_{A|R} : \underline{\check{C}}_{A|R} \to \underline{\widetilde{C}}_{A|R}$  is the "associated-sheaf"  functor. Now let  $F$  be an additive sheaf. We claim that  $\beta^{\#}\underline{H}^{q}_{A|S}(F)$  is a constant sheaf for all  $q > 0$ . Indeed, if  $P$  is a polynomial algebra over  $R$ , then  $P = R\underset{S}{\otimes}P'$ for a polynomial algebra  $P'$  over  $S$  and hence

}{F7I-VI-058}\begin{equation*}
[\beta^{\#}\underline{H}^{q}_{A|S}(F)](P) = H^{q}(P|S,\ F) = H^{q}(R\underset{S}{\otimes}P'|S,\ F) \to H^{q}(R|S,\ F)
\end{equation*}
is an isomorphism for all  $q > 0$  by 3.7(b), and surjective for  $q = 0$ . It follows from 3.3(c) that the canonical map  $\beta^{\#}\underline{H}^{q}_{A|S}(F) \to H^{q}(R|S,\ F) =$ constant sheaf is an isomorphism for all  $q > 0$ , and is surjective for  $q = 0$ . It follows that $E^{p,q} = H^{p}(A|R,\ \beta^{\#}\underline{H}^{q}_{A|S}(F)) = 0$  for all  $p,\ q > 0$ , and hence the above spectral sequence yields the exact sequence

\begin{equation*}
(*)\quad 0 \to E^{1,0} \to H^{1} \to E^{0,1}\!\to\ E^{2,0} \to H^{2} \to E^{0,2} \to E^{3,0}\!\to\ H^{3} \to\ ....
\end{equation*}
where we note that  $E^{0,q} = [\beta^{\#}\underline{H}^{q}_{A|S}(F)](A) \overset{\sim}{\to} H^{q}(R|S,\ F)$  for all  $q > 0$  and $E^{p,0} = H^{p}(A|R,\ \widetilde{\beta}F)$ . On the other hand, the exact sequence  $0 \to F_{R} \to \widetilde{\beta}F \to F(R) \to 0$ gives us the exact sequence

{\footnotesize (**)\quad\begin{tikzcd}[column sep=tiny, row sep=large]
0 \arrow[r] \& H^{0}(A|R,F_{R}) \arrow[r] \& H^{0}(A|R,\widetilde{\beta}F) \arrow[r] \arrow[d, equal] \& H^{0}(A|R,F(R)) \arrow[r] \arrow[d, equal] \& H^{1}(A|R,F_{R}) \arrow[r] \& H^{1}(A|R,\widetilde{\beta}F) \arrow[r] \& 0 \\
 \& \& H^{0}(A|S,F) \& H^{0}(R|S,F) \& \& \&
\end{tikzcd}}










\SGAvar{The exact sequences (*) and (**) combined gives the desired exact sequence.

}{The exact sequences (*) and (**) combined give the desired exact sequence.

}{F7I-VI-059}\SGApage{I}{81}


\underline{Corollary} 3.8.

(a)  \underline{If}  $A, B$  \underline{are} $R$-\underline{algebras such that}  $\mathrm{Tor}^{R}_{i}(A, B) = 0$  \underline{for all} $0 < i \leq n$  , \underline{then for any additive sheaf}  $F \in \mathrm{ob}(\widetilde{\underline{C}}_{A\underset{R}{\otimes}B|R})$  , \underline{the canonical map}

\begin{equation*}
H^{i}(A\underset{R}{\otimes}B|R,\ F) \rightarrow H^{i}(A|R,\ F) \oplus H^{i}(B|R,\ F)
\end{equation*}
\underline{is an isomorphism for all}  $i \leq n$  .

(b)  \underline{Let}  $A$  \underline{be an} $R$-\underline{algebra such that}  $\mathrm{Tor}^{R}_{i}(A, A) = 0$  \underline{for all}  $i > 0$  . \underline{If}  $A\underset{R}{\otimes}A \rightarrow A$  \underline{is an isomorphism, then for any additive sheaf}  $F \in \mathrm{ob}(\widetilde{\underline{C}}_{A|R})$  , \underline{we have}  $H^{i}(A|R,\ F) = 0$  \underline{for all}  $i$  .


\underline{Proof:}

\SGAvar{(a)  Let  $F \in \mathrm{ob}(\widetilde{\underline{C}}_{A\underset{R}{\otimes}B|R})$  be additive. By 3.7(c), the ring homomorphisms $R \rightarrow A \rightarrow A\underset{R}{\otimes}B$  gives us the exact sequence

\begin{equation}\tag{*}
... \rightarrow H^{i}(A\underset{R}{\otimes}B\ A,\ F_{A}) \rightarrow H^{i}(A\underset{R}{\otimes}B\ R,\ F) \rightarrow H^{i}(A|R,\ F) \rightarrow ....
\end{equation}
consider now the commutative diagram

}{(a)  Let  $F \in \mathrm{ob}(\widetilde{\underline{C}}_{A\underset{R}{\otimes}B|R})$  be additive. By 3.7(c), the ring homomorphisms $R \rightarrow A \rightarrow A\underset{R}{\otimes}B$  give us the exact sequence

\begin{equation}\tag{*}
... \rightarrow H^{i}(A\underset{R}{\otimes}B|A,\ F_{A}) \rightarrow H^{i}(A\underset{R}{\otimes}B|R,\ F) \rightarrow H^{i}(A|R,\ F) \rightarrow ....
\end{equation}
consider now the commutative diagram

}{F7I-VI-060}\begin{center}
\begin{tikzcd}[column sep=huge, row sep=huge]
H^{i}(A\underset{R}{\otimes}B|A,\ F_{A}) \arrow[r] \arrow[d] \& H^{i}(A\underset{R}{\otimes}B|R,\ F) \arrow[d] \\
H^{i}(B|R,\ \widetilde{A}F_{A}) \arrow[r, "\approx"] \& H^{i}(B|R,\ F)
\end{tikzcd}
\end{center}









where the functor  $A : \underline{C}_{B|R} \rightarrow \underline{C}_{A\underset{R}{\otimes}B|A}$  is given by  $X \longmapsto A\underset{R}{\otimes}X$  , and the map $H^{i}(B|R,\ \widetilde{A}F_{A}) \rightarrow H^{i}(B|R,\ F)$  is induced from the map  $\widetilde{A}F_{A} \rightarrow F$  , which is an isomorphism since  $F$  is additive.  Note that the commutativity of the above diagram, because of its functorial nature, is reduced to the case  $i = 0$  .

If  $\mathrm{Tor}^{R}_{i}(A, B) = 0$  for all  $0 < i \leq n$  , then  $H^{i}(A\underset{R}{\otimes}B|A,\ F_{A}) \rightarrow H^{i}(B|R,\ \widetilde{A}F_{A})$  is an isomorphism for all  $i \leq n$  by 3.7(a), and hence the composite map

\begin{equation*}
H^{i}(A\underset{R}{\otimes}B|A,\ F_{A}) \rightarrow H^{i}(A\underset{R}{\otimes}B|R,\ F) \rightarrow H^{i}(B|R,\ F)
\end{equation*}
and (because of the symmetry) 


\SGApage{I}{82}

\begin{equation*}
 H^{i}(A\underset{R}{\otimes}B|B,\ F_{B}) \rightarrow H^{i}(A\underset{R}{\otimes}B|R,\ F) \rightarrow H^{i}(A|R,\ F)
\end{equation*}
are isomorphisms for all  $i \leq n$  . This implies that the exact sequence (*) breaks up to short split-exact sequences

\begin{equation*}
0 \rightarrow H^{i}(A\underset{R}{\otimes}B|A,\ F_{A}) \rightarrow H^{i}(A\underset{R}{\otimes}B|R,\ F) \rightarrow H^{i}(A|R,\ F) \rightarrow 0
\end{equation*}
for all  $i \leq n$  , and that

\begin{equation*}
H^{i}(A\underset{R}{\otimes}B|R,\ F) \rightarrow H^{i}(A|R,\ F) \oplus H^{i}(B|R,\ F)
\end{equation*}
is an isomorphism for all  $i \leq n$  .

(b)  Since  $\mathrm{Tor}^{R}_{i}(A,A) = 0$  for all  $i > 0$  , it follows from 3.8(a) that

\begin{equation*}
H^{i}(A\underset{R}{\otimes}A|R,\ F) \rightarrow H^{i}(A|R,\ F) \oplus H^{i}(A|R,\ F)
\end{equation*}
is an isomorphism for all  $i$  . On the other hand,  $A\underset{R}{\otimes}A \rightarrow A$  is an isomorphism by hypothesis and hence  $H^{i}(A|R,\ F) \rightarrow H^{i}(A\underset{R}{\otimes}A|R,\ F)$  is an isomorphism for all  $i$  . This means that the diagonal map  $H^{i}(A|R,\ F) \rightarrow H^{i}(A|R,\ F) \oplus H^{i}(A|R,\ F)$  is an isomorphism for all  $i$  i.e.,  $H^{i}(A|R,\ F) = 0$  for all  $i$  .

\underline{Corollary} 3.9:  Let  $S$  be a multiplicative subset of an $R$-algebra  $A$  . If  $F \in \mathrm{ob}\!\left(\widetilde{\underline{C}}_{S^{-1}A|R}\right)$  is additive such that  $F_{A}$  is also additive, then the canonical map

\begin{equation*}
H^{i}(S^{-1}A|R,\ F) \rightarrow H^{i}(A|R,\ F)
\end{equation*}
is an isomorphism for all  $i$  , where  $F_{A} \in \mathrm{ob}\!\left(\widetilde{\underline{C}}_{S^{-1}A|A}\right)$  is defined by $F_{A}(X) = \mathrm{Ker}(F(X) \rightarrow F(A))$  for all  $X \in \mathrm{ob}\!\left(\underline{C}_{S^{-1}A|A}\right)$  .

\SGAvar{\underline{Proof}:  Since  $F$  is additive, the ring homomorphisms  $R \rightarrow A \rightarrow S^{-1}A$  gives us the exact sequence

}{\underline{Proof}:  Since  $F$  is additive, the ring homomorphisms  $R \rightarrow A \rightarrow S^{-1}A$  give us the exact sequence

}{F7I-VI-061}\begin{equation*}
.... \rightarrow H^{n}(S^{-1}A|A,\ F_{A}) \rightarrow H^{n}(S^{-1}A|R,\ F) \rightarrow H^{n}(A|R,\ F) \rightarrow ....
\end{equation*}
On the other hand, since  $S^{-1}A\underset{A}{\otimes}S^{-1}A \rightarrow S^{-1}A$  is an isomorphism and  $F_{A}$  is additive by hypothesis, it follows from 3.8(b) that  $H^{n}(S^{-1}A|A,\ F_{A}) = 0$  for all  $n$  , and hence  $H^{n}(S^{-1}A|R,\ F) \rightarrow H^{n}(A|R,\ F)$  is an isomorphism for all  $n$  .


\SGApage{I}{83}

Now let  $A$  be an $R$-algebra. An $A$-module  $M$  defines an additive sheaf $\mathrm{Der}_R(-,\ M) \in \mathrm{ob}(\widetilde{\underline{C}}_{A|R})$  by  $X \to \mathrm{Der}_R(X,\ M)$  . We now make the following definition:


\underline{Definition} 3.10:  \underline{Let}  $A$  \underline{be an} $R$-\underline{algebra. For each} $A$-\underline{module}  $M$  \underline{we set}

\begin{equation*}
D^i(A|R,\ M) = H^i(A|R,\ \mathrm{Der}_R(-,\ M)) \ \ .
\end{equation*}
We want to make an explicit computation of low-dimensional ones i.e., $D^i(A|R,\ M)$  for  $i = 1,\ 2$  : Choose a surjective arrow  $P \to A$  in  $\underline{C}_{A|R}$  where  $P$  is a polynomial algebra over  $R$  , and let  $0 \to I \to P \to A \to 0$  be exact. Then


\underline{Lemma} 3.10:

\begin{equation*}
\begin{aligned}
D^1(A|R,\ M) &= \mathrm{Coker}(\mathrm{Der}_R(P,\ M) \to \mathrm{Hom}_P(I,\ M)) \\
D^2(A|R,\ M) &= H^1(A|R,\ \underline{H}^1_{A|R}(\mathrm{Der}_R(-,\ M)))
\end{aligned}
\end{equation*}
\underline{Proof}:  Consider the simplicial algebra

\begin{multline}\tag{*}
........\ \ P_A^{(2)} \ \mathrel{\substack{\to\\ \to\\ \to}} \ P_A^{(1)} \ \mathrel{\substack{\to\\ \to}} \ P_A^{(0)}
\end{multline}
where  $P_A^{(n)} = \overbrace{P\times_A \ldots \times_A P}^{\,n+1}$  . For each  $n$  we have the exact sequence

\begin{equation*}
0 \to I^{(n)} \to P_A^{(n)} \to A \to 0
\end{equation*}
where  $I^{(n)} = \overbrace{I\times I\times \ldots \times I}^{\,n+1}$  , and it is easy to see that

(a):

\begin{equation*}
0 \to \mathrm{Hom}_A(\overline{I}^{(n)}/\Delta(\overline{I}),\ M) \to \mathrm{Der}_R(P_A^{(n)},\ M) \xrightarrow{\ \mathrm{Der}(\Delta,\ M)\ } \mathrm{Der}_R(P,\ M) \to 0
\end{equation*}
is exact for all  $n$  , where  $\overline{I} = I/I^2$  and  $\Delta$  stands for the diagonal map. We also note the exact sequence

(b):

\SGAvar{\begin{equation*}
0 \to \mathrm{Hom}_A(\overline{I}^{(n)}/\Delta(\overline{I}),\ M) \to \mathrm{Hom}_A(\overline{I}^{(n)},\ M) \xrightarrow{\ \mathrm{Hom}(\Delta,\ M)\ } \mathrm{Hom}_A(I,\ M) \to 0
\end{equation*}
The exact sequences (a) and (b) yield the respective exact sequences of simplicial $A$-modules. Since  $\{I \to 0\}$ , $\{I \to I\}$  admits sections, we have

}{\begin{equation*}
0 \to \mathrm{Hom}_A(\overline{I}^{(n)}/\Delta(\overline{I}),\ M) \to \mathrm{Hom}_A(\overline{I}^{(n)},\ M) \xrightarrow{\ \mathrm{Hom}(\Delta,\ M)\ } \mathrm{Hom}_A(\overline{I},\ M) \to 0
\end{equation*}
The exact sequences (a) and (b) yield the respective exact sequences of simplicial $A$-modules. Since  $\{I \to 0\}$ , $\{I \to I\}$  admit sections, we have

}{F7I-VI-062}\begin{equation*}
H^i(\{I \to 0\},\ \mathrm{Hom}_A(-,\ M)) = H^i(\{I \to I\},\ \mathrm{Hom}_A(-,\ M)) = 0 
\end{equation*}

\SGApage{I}{84}

 for all  $i > 0$  , and it follows readily that

\begin{equation*}
\left\{\begin{array}{l}
0 \to \mathrm{Der}_R(A,\ M) \to \mathrm{Der}_R(P,\ M) \to \mathrm{Hom}_A(\overline{I},\ M) \to H^1(\{P \to A\},\ \mathrm{Der}_R(-,\ M)) \to 0 \\[4pt]
\hspace{9cm}\text{is exact} \\[4pt]
H^i(\{P \to A\},\ \mathrm{Der}_R(-,\ M)) = 0 \ \ \text{for all} \ \ i > 1
\end{array}\right.
\end{equation*}
It follows from 3.1(b) that

\begin{equation*}
\begin{aligned}
D^1(A|R,\ M) &= \check{H}^1(A|R,\ \mathrm{Der}_R(-,\ M)) = H^1(\{P \to A\},\ \mathrm{Der}_R(-,\ M)) \\
&= \mathrm{Coker}(\mathrm{Der}_R(P,\ M) \to \mathrm{Hom}_A(\overline{I},\ M)) \\
&= \mathrm{Coker}(\mathrm{Der}_R(P,\ M) \to \mathrm{Hom}_P(I,\ M)) \ \ , \\
D^2(A|R,\ M) &= H^1(A|R,\ \underline{H}^1_{A|R}(\mathrm{Der}_R(-,\ M))) \ \ .
\end{aligned}
\end{equation*}
Now choose a free $P$-module  $F$  and a $P$-linear map  $F \xrightarrow{\ \gamma\ } P$  such that  $\gamma(F) = I$  . Associated to $P$-linear map  $\gamma : F \to P$  , there is a Koszul complex, which is simply denoted, by abuse of notation, by  $K_I$  .


\SGAvar{\underline{Lemma} 3.11:   $D^2(A|R,\ M) \simeq \mathrm{Coker}(H^1(K_I,\ M) \to \mathrm{Hom}_A(H_1(K^I),\ M))$ .

}{\underline{Lemma} 3.11:   $D^2(A|R,\ M) \simeq \mathrm{Coker}(H^1(K_I,\ M) \to \mathrm{Hom}_A(H_1(K_I),\ M))$ .

}{F7I-VI-063}\underline{Proof}:  We have

\begin{equation*}
\begin{aligned}
D^2(A|R,\ M) &= H^1(A|R,\ H^1_{A|R}(\mathrm{Der}_R(-,\ M))) \\
&= \mathrm{Ker}(H^1(P_A^{(1)}|R,\ M) \ \mathrel{\substack{\to\\ \to\\ \to}} \ H^1(P_A^{(2)}|R,\ M))
\end{aligned}
\end{equation*}
since  $[\underline{H}^1_{A|R}(\mathrm{Der}_R(-,\ M))](P_A^0) = [\underline{H}^1_{A|R}(\mathrm{Der}_R(-,\ M))](P) = 0$  .

Let  $0 \to K \to F \xrightarrow{\ \gamma\ } I \to 0$  be exact. If we set  $P[F]$  to be the symmetric algebra over  $P$  generated by  $F$  , then  $P[F]$  and  $P[F]\otimes P[F]$  are polynomial algebras over  $R$  and we obtain the surjective arrows

\begin{equation*}
\begin{aligned}
P[F] \to P_A^{(1)} &= P\times_A P &&\text{given by} && X \to (0,\ \gamma(X)) \\
P[F]\otimes_P P[F] \to P_A^{(2)} &= P\times_A P\times_A P && && X\otimes 1 \to (0,\ \gamma(X),\ 0),\ 1\otimes X \to (0,\ 0,\ \gamma(X))
\end{aligned}
\end{equation*}
where  $X \in F$  . If we set  $J = \mathrm{Ker}(P[F] \to P_A^{(1)})$ and
$J' = \mathrm{Ker}(P[F]\otimes_P P[F] \to P_A^{(2)})$  , then 


\SGApage{I}{85}

 
\begin{align*}
J &= (F) \cap (F_\gamma) = (K) + (F\cdot F_\gamma)\\
J' &= (F\otimes 1 + 1\otimes F) \cap (F_\gamma\otimes 1 + 1\otimes F) \cap (F\otimes 1 + 1\otimes F_\gamma)\\
&= (K\otimes 1 + 1\otimes K) + (F\cdot F_\gamma\otimes 1 + 1\otimes F\cdot F_\gamma) + (F\otimes F)
\end{align*}

where  $F_\gamma$ $=$ the ideal generated by  $\{X - \gamma(X)|X \in F\}$  . Now  $\pi_i : P_A^{(2)} \to P_A^{(1)}$  $i = 0, 1, 2$  can be lifted to  $\hat{\pi}_i : P[F]\underset{P}{\otimes}P[F] \to P[F]$  where, for each  $X \in F$  ,

$$\pi_0 = \begin{cases} X\otimes 1 \to \gamma(X) - X\\[3ex] 1\otimes X \to X\end{cases}\qquad\qquad \pi_1 = \begin{cases} X\otimes 1 \to 0\\[3ex] 1\otimes X \to X\end{cases}\qquad\qquad \pi_2 = \begin{cases} X\otimes 1 \to X\\[3ex] 1\otimes X \to 0\end{cases}\qquad .$$
A trivial computation shows that  $\pi_0 - \pi_1 + \pi_2$  sends  $K\otimes 1 + 1\otimes K + 1\otimes F\cdot F_\gamma$  to zero, whereas it sends  $(F\cdot F_\gamma\otimes 1 + F\otimes F)$  onto  $(F\cdot F_\gamma)$  . Now for any $D \in \mathrm{Der}_R(P[F]\underset{P}{\otimes}P[F], M)$  ,  $D(F\cdot F_\gamma\otimes 1 + F\otimes F) = 0$  since  $I\cdot M = 0$  , and consequently

\begin{align*}
D^2(A|R, M) &= \mathrm{Ker}(D^1(P_A^{(1)}|R, M) \mathrel{\substack{\to\\ \to\\ \to}} D^1(P_A^{(2)}|R, M))\\
&= \{[f] \in D^1(P_A^{(1)}|R, M)|f(F\cdot F_\gamma) = 0\}\\
&= \mathrm{Coker}(\mathrm{Der}_R(P[F], M) \to \mathrm{Hom}_{P[F]}((F) \cap (F_\gamma)/(F\cdot F_\gamma), M))\\
&= \mathrm{Coker}(\mathrm{Der}_P(P[F], M) \to \mathrm{Hom}_{P[F]}((F) \cap (F_\gamma)/(F)\cdot(F_\gamma), M))\quad .
\end{align*}

\SGAvar{Now
$$\begin{aligned}
0 \to (F) \to P[F] &\overset{\;\;P[0]\;\;}{\relbar\joinrel\relbar\joinrel\relbar} P \to 0\\
0 \to (F_\gamma) \to P[F] &\overset{\;\;P[\gamma]\;\;}{\relbar\joinrel\relbar\joinrel\relbar} P \to 0 \quad .
\end{aligned}$$
are exact and hence  $(F) \cap (F_\gamma)/(F)\cdot(F_\gamma) = \mathrm{Tor}_1^{P[F]}(P, P_\gamma) = H_1(K_I)$  where  $K_I$  is the Koszul complex associated  to the $P$-linear map  $F \xrightarrow{\gamma} I(\subset P)$  . We thus have

\begin{align*}
D^2(A|R, M) &= \mathrm{Coker}(\mathrm{Hom}_P(F, M) \to \mathrm{Hom}_P(H_1(K_I), M))\\
&= \mathrm{Coker}(H^1(K_I, M) \to \mathrm{Hom}_A(H^1(K_I), M))\quad .
\end{align*}

}{Now
$$\begin{aligned}
0 \to (F) \to P[F] &\xrightarrow{\;\;P[0]\;\;} P \to 0\\
0 \to (F_\gamma) \to P[F] &\xrightarrow{\;\;P[\gamma]\;\;} P \to 0 \quad .
\end{aligned}$$
are exact and hence  $(F) \cap (F_\gamma)/(F)\cdot(F_\gamma) = \mathrm{Tor}_1^{P[F]}(P, P_\gamma) = H_1(K_I)$  where  $K_I$  is the Koszul complex associated  to the $P$-linear map  $F \xrightarrow{\gamma} I(\subset P)$  . We thus have

\begin{align*}
D^2(A|R, M) &= \mathrm{Coker}(\mathrm{Hom}_P(F, M) \to \mathrm{Hom}_P(H_1(K_I), M))\\
&= \mathrm{Coker}(H^1(K_I, M) \to \mathrm{Hom}_A(H_1(K_I), M))\quad .
\end{align*}

}{F7I-VI-064}We can now state all the basic facts needed for our purpose, summarizing the various results we have already established.


\SGApage{I}{86}


\underline{Theorem} 3.12.

(1)  \underline{If}  $0 \to M' \to M \to M'' \to 0$  \underline{is an exact sequence of} $A$-\underline{modules, then}

\begin{gather*}
0 \to D^0(A|R, M') \to D^0(A|R, M) \to D^0(A|R, M'') \to D^1(A|R, M') \to\\
D^1(A|R, M)\to\;\; D^1(A|R, M'') \to\; ........ \;\to D^n(A|R, M') \to\\
D^n(A|R, M) \to D^n(A|R, M'') \to \;...
\end{gather*}

\underline{is exact}.

(2)  \underline{Let}  $S \to R \to A$  \underline{be ring homomorphisms. Then for any $A$-module}  $M$  \underline{we have the exact sequence}

\SGAvar{\begin{gather*}
0 \to D^0(A\;R, M) \to D^0(A\;S, M) \to D^0(R\;S, M) \to D^1(A\;R, M) \to\\
D^1(A\;S, M) \to D^1(R\;S, M) \to\; ........ \;\to D^n(A\;R, M) \to\\
D^n(A\;S, M) \to D^n(R\;S, M) \to \;....
\end{gather*}

}{\begin{gather*}
0 \to D^0(A|R, M) \to D^0(A|S, M) \to D^0(R|S, M) \to D^1(A|R, M) \to\\
D^1(A|S, M) \to D^1(R|S, M) \to\; ........ \;\to D^n(A|R, M) \to\\
D^n(A|S, M) \to D^n(R|S, M) \to \;....
\end{gather*}

}{F7I-VI-065}(3)  \underline{Let}  $A, S$  \underline{be $R$-algebras such that}  $\mathrm{Tor}_i^R(A, S) = 0$  \underline{for all}  $i > 0$  . \underline{Then for any}  $A\underset{R}{\otimes}S$-\underline{module}  $M$  , \underline{the canonical map}  $D^i(A\underset{R}{\otimes}S|S, M) \to D^i(A|R, M)$  \underline{is an} \underline{isomorphism for all}  $i$  . \underline{In particular, if}  $A$  \underline{is} $R$-\underline{flat we have}

$$D^i(A\otimes S|S, M) \xrightarrow{\sim} D^i(A|R, M)$$
\underline{for all}  $i$  \underline{and for all} $R$-\underline{algebras}  $S$  .

(4)  \underline{Let}  $S$  \underline{be a multiplicative subset of}  $A$  . \underline{Then for any}  $S^{-1}A$-\underline{module} $N$  , \underline{the canonical map}

$$D^i(S^{-1}A|R, N) \to D^i(A|R, N)$$
\underline{is an isomorphism for all}  $i$  . \underline{Consequently, for any} $A$-\underline{module}  $M$  \underline{we have the} \underline{canonical isomorphism}

$$D^i(S^{-1}A|R, S^{-1}M) \xrightarrow{\sim} S^{-1}D^i(A|R, M)$$
\underline{for all}  $i$  .

(5)  \underline{If}  $A, B$  \underline{are} $R$-\underline{algebras such that}  $\mathrm{Tor}_i^R(A, B) = 0$  \underline{for all} $0 < i \leq n$  , \underline{then for any}  $A\underset{R}{\otimes}B$-\underline{module}  $M$  , \underline{the canonical map}

$$D^i(A\underset{R}{\otimes}B|R, M) \to D^i(A|R, M) \oplus D^i(B|R, M)$$
\underline{is an isomorphism for all}  $i \leq n$  . \underline{In particular, if}  $\mathrm{Tor}_i^R(A, A) = 0$  \underline{for all}
\SGApage{I}{87}
$i > 0$  \underline{and}

\begin{equation*}
\Delta \ : \ \mathrm{Spec}(A) \ \to \ \mathrm{Spec}(A) \underset{\mathrm{Spec}(R)}{\times} \mathrm{Spec}(A)
\end{equation*}
\underline{is an open immersion} (\underline{for instance}  $A/R$  \underline{etale}), \underline{then}  $D^i(A|R, M) = 0$  \underline{for all} $i > 0$  \underline{and for all} $A$-\underline{modules}  $M$ .

(6) \underline{Let}  $A$  \underline{be an} $R$-\underline{algebra of essentially finite type over}  $R$ , \underline{where} $R$  \underline{is noetherian, and let}  $E$  \underline{be a flat} $A$-\underline{module. Then we have a canonical isomorphism}

\begin{equation*}
E \underset{A}{\otimes} D^i(A|R, \ M) \ \overset{\sim}{\longrightarrow} \ D^i(A|R, \ E \underset{A}{\otimes} M)
\end{equation*}
\underline{for all}  $i \geq 0$  \underline{and for all} $A$-\underline{modules}  $M$ .

\SGAvar{(7) \underline{If}  $R$  \underline{is noetherian and}  $A$  \underline{is an} $R$-\underline{algebra of essentially finite} \underline{type, then}  $D^i(A|R, M)$  \underline{is} $A$-\underline{module of finite type for all}  $i$ , \underline{provided}  $M$  \underline{is} \underline{an} $A$-\underline{module of finite type}.

}{(7) \underline{If}  $R$  \underline{is noetherian and}  $A$  \underline{is an} $R$-\underline{algebra of essentially finite} \underline{type, then}  $D^i(A|R, M)$  \underline{is an} $A$-\underline{module of finite type for all}  $i$ , \underline{provided}  $M$  \underline{is} \underline{an} $A$-\underline{module of finite type}.

}{F7I-VI-066}(8) \underline{Let}  $P$  \underline{be a polynomial algebra over}  $R$  \underline{and let}  $0 \to I \to P \to A \to 0$ \underline{be exact. Then}  $D^0(A|R, M) = \mathrm{Der}_R(A, M)$  ,

\SGAvar{\begin{equation*}
\begin{aligned}
D^1(A|R, \ M) \ &= \ \mathrm{Coker}(\mathrm{Der}_R(P, \ M) \ \to \ \mathrm{Hom}_A(I/I^2, \ M)\\
&= \ \underline{\text{the set of isomorphism classes of}}\\
&\qquad\ \ \text{R-}\underline{\text{algebra extensions of}} \ \ A \ \ \underline{\text{by}} \ \ M \ ,\\
D^2(A|R, \ M) \ &= \ \mathrm{Coker}(H^1(K_I, \ M) \ \to \ \mathrm{Hom}_A(H_1(K_I), \ M)) \ .
\end{aligned}
\end{equation*}

}{\begin{equation*}
\begin{aligned}
D^1(A|R, \ M) \ &= \ \mathrm{Coker}(\mathrm{Der}_R(P, \ M) \ \to \ \mathrm{Hom}_A(I/I^2, \ M))\\
&= \ \underline{\text{the set of isomorphism classes of}}\\
&\qquad\ \ \text{R-}\underline{\text{algebra extensions of}} \ \ A \ \ \underline{\text{by}} \ \ M \ ,\\
D^2(A|R, \ M) \ &= \ \mathrm{Coker}(H^1(K_I, \ M) \ \to \ \mathrm{Hom}_A(H_1(K_I), \ M)) \ .
\end{aligned}
\end{equation*}

}{F7I-VI-067}\underline{Proof}:

(1) If  $0 \to M' \to M \to M'' \to 0$  is an exact sequence of $A$-modules, then for any polynomial algebra  $P$  over  $R$  ,  $0 \to \mathrm{Der}_R(P, M') \to \mathrm{Der}_R(P, M) \to \mathrm{Der}_R(P, M'') \to 0$ is exact and hence  $0 \to \mathrm{Der}_R(-, M') \to \mathrm{Der}_R(-, M) \to \mathrm{Der}_R(-, M'') \to 0$  is an exact sequence of sheaves by 3.3(c).

(2) $\mathrm{Der}_S(-, M)$  is an additive sheaf on  $\underline{C}_{A|S}$  , and $\mathrm{Der}_R(B, M) = \mathrm{Ker}(\mathrm{Der}_S(B, M) \to \mathrm{Der}_S(R, M))$  for every  $B \in \mathrm{ob}(\underline{C}_{A|R})$  . Therefore our assertion follows from 3.7(c).

(3) Let  $S : \underline{C}_{A|R} \to \underline{C}_{A \underset{R}{\otimes} S|S}$  be the functor given by  $X \longmapsto X \underset{R}{\otimes} S$  . Then $\widetilde{S}\mathrm{Der}_S(-, M) = \mathrm{Der}_R(-, M)$  and hence our assertion follows from 3.7(a).

\SGApage{I}{88}

(4)  Follows from 3.9. Indeed,  $F = \mathrm{Der}_R(-, N)$  is an additive sheaf on $\underline{C}_{S^{-1}A|R}$  , and  $F_A \in \mathrm{ob}\left(\underline{C}^{\sim}_{S^{-1}A|A}\right)$  is also additive since $F_A(X) = \mathrm{Ker}(F(X) \to F(A)) = \mathrm{Ker}(\mathrm{Der}_R(X, N) \to \mathrm{Der}_R(A, N)) = \mathrm{Der}_A(X, N)$  i.e., $F_A = \mathrm{Der}_A(-, N)$  .

(5)  The first assertion follows from 3.8(a). Now assume that $\mathrm{Tor}^R_i(A, A) = 0$  for all  $i > 0$  and that

\begin{equation*}
\Delta \ : \ \mathrm{Spec}(A) \ \to \ \mathrm{Spec}(A) \underset{\mathrm{Spec}(R)}{\times} \mathrm{Spec}(A)
\end{equation*}
is an open immersion i.e.,  $(A \underset{R}{\otimes} A)_{\varphi^{-1}(\underline{m})}$  $\to A_{\underline{m}}$  is an isomorphism for all maximal ideals  $\underline{m}$  of  $A$  , where  $\varphi : A \underset{R}{\otimes} A \to A$  is the multiplication map. Then the first assertion combined with 3.12(4) entails that the diagonal map

\begin{equation*}
\Delta_{\underline{m}} \ : \ D^i(A|R, \ M)_{\underline{m}} \ \to \ D^i(A|R, \ M)_{\underline{m}} \ \oplus \ D^i(A|R, \ M)_{\underline{m}}
\end{equation*}
is an isomorphism for every maximal ideal  $\underline{m}$  of  $A$  and hence

\begin{equation*}
\Delta \ : \ D^i(A|R, \ M) \ \to \ D^i(A|R, \ M) \ \oplus \ D^i(A|R, \ M)
\end{equation*}
is an isomorphism, which entails that  $D^i(A|R, M) = 0$  for all  $i$  .

(6)  In view of 3.12(4), we may assume that  $A$  is an R-algebra of finite type. Now consider the map

\begin{equation*}
E \underset{A}{\otimes} \mathrm{Der}_R(-, \ M) \ \to \ \mathrm{Der}_R(-, \ E \underset{A}{\otimes} M)
\end{equation*}
in  $\underline{C}^{\sim}_{A|R}$  . Since  $E$  is A-flat, we have an isomorphism

\begin{equation*}
E \underset{A}{\otimes} \mathrm{Der}_R(B, \ M) \ \overset{\sim}{\longrightarrow} \ \mathrm{Der}_R(B, \ E \underset{A}{\otimes} M)
\end{equation*}
whenever  $B \in \mathrm{ob}(\underline{C}_{A|R})$  is of finite type over  $R$  . It follows from 3.3(b) that

\begin{equation*}
D^i(A|R, \ E \underset{A}{\otimes} \mathrm{Der}_R(-, \ M)) \ \to \ D^i(A|R, \ E \underset{A}{\otimes} M)
\end{equation*}
is an isomorphism. On the other hand, the functor  $E \underset{A}{\otimes} -$  :  $(A\text{-mod}) \to (A\text{-mod})$  is exact, and therefore we have

\begin{equation*}
E \underset{A}{\otimes} D^i(A|R, \ M) \ = \ D^i(A|R, \ E \underset{A}{\otimes} \mathrm{Der}_R(-, \ M)) \ \ ,
\end{equation*}
and hence our assertion.

(7)  follows from 3.3(a).


\SGApage{I}{89}

(8)  In view of 3.10 and 3.11, it suffices to show that $\mathrm{Coker}(\mathrm{Der}_R(P, M) \to \mathrm{Hom}_A(I/I^2, M)) = \mathrm{Exalcomm}(A|R, M)$ , where  $\mathrm{Exalcomm}(A|R, M)$ denotes the set of isomorphism classes of R-algebra extensions of  $A$  by  $M$ . Let $0 \to M \to A' \to A \to 0$   be an R-algebra extension of  $A$  by  $M$ . Since  $P$  is a polynomial algebra over  $R$ , we obtain a commutative diagram

\begin{center}
\begin{tikzcd}[column sep=large]
0 \arrow[r] \& I \arrow[r] \arrow[d] \& P \arrow[r] \arrow[d] \& A \arrow[r] \arrow[d, equal] \& 0 \\
0 \arrow[r] \& M \arrow[r] \& A' \arrow[r] \& A \arrow[r] \& 0
\end{tikzcd}
\end{center}











\SGAvar{where  $I \to M$  is uniquely determined by the extensions  $A'$  up to  $\mathrm{Der}_R(P, M)$  . We thus obtain a map

}{where  $I \to M$  is uniquely determined by the extension  $A'$  up to  $\mathrm{Der}_R(P, M)$  . We thus obtain a map

}{F7I-VI-068}\begin{equation*}
\Phi \ : \ \mathrm{Exalcomm}(A|R, M) \to \mathrm{Coker}(\mathrm{Der}_R(P, M) \to \mathrm{Hom}_A(I/I^2, M))  \ .
\end{equation*}
Conversely a P-linear map  $I \to M$  yields an extension by push-out

\begin{center}
\begin{tikzcd}[column sep=large]
0 \arrow[r] \& I \arrow[r] \arrow[d] \& P \arrow[r] \arrow[d] \& A \arrow[r] \arrow[d, equal] \& 0 \\
0 \arrow[r] \& M \arrow[r] \& A' \arrow[r] \& A \arrow[r] \& 0
\end{tikzcd}
\end{center}










\SGAvar{Two P-linear maps  $I \rightrightarrows M$  which differ by  $\mathrm{Der}_R(P, M)$  is trivially seen to yield isomorphic extensions and therefore we obtain a map

\begin{equation*}
\Psi \ : \ \mathrm{Coker}(\mathrm{Der}_R(P, M) \to \mathrm{Hom}_A(I/I^2, M)) \to \mathrm{Exalcomm}(A|R, M)  \ .
\end{equation*}
It is straightforward to verify that  $\Phi$, $\Psi$  are inverse process to each other.

\underline{Corollary} 3.13.  (a)  R-\underline{algebra}  $A$  \underline{is formally smooth over}  $R$  \underline{if and only if} $D^1(A|R, M) = 0$  \underline{for every} A-\underline{module}  $M$  .

}{Two P-linear maps  $I \rightrightarrows M$  which differ by  $\mathrm{Der}_R(P, M)$  are trivially seen to yield isomorphic extensions and therefore we obtain a map

\begin{equation*}
\Psi \ : \ \mathrm{Coker}(\mathrm{Der}_R(P, M) \to \mathrm{Hom}_A(I/I^2, M)) \to \mathrm{Exalcomm}(A|R, M)  \ .
\end{equation*}
It is straightforward to verify that  $\Phi$, $\Psi$  are inverse processes to each other.

\underline{Corollary} 3.13.  (a)  The R-\underline{algebra}  $A$  \underline{is formally smooth over}  $R$  \underline{if and only if} $D^1(A|R, M) = 0$  \underline{for every} A-\underline{module}  $M$  .

}{F7I-VI-069}(b)  \underline{Let}  $R$  \underline{be noetherian and}  $A$  \underline{an} R-\underline{algebra of finite type.} \underline{Then}  $A$ \underline{is a relative complete-intersection over}  $R$  \underline{if and only if}  $D^2(A|R, M) = 0$  \underline{for every} A-\underline{module}  $M$  .

(c)  \underline{Let}  $R$  \underline{be noetherian and}  $A$  \underline{an} R-\underline{algebra of finite type.} \underline{Then for} \underline{any multiplicative subset}  $S$  \underline{in}  $A$  , \underline{we have a canonical isomorphism}

\begin{equation*}
D^i(S^{-1}A|R, \ S^{-1}M) \xrightarrow{\ \sim\ } S^{-1}D^i(A|R, M)  \ \underline{\text{for all}} \ \ i \geq 0  \ .
\end{equation*}

\SGApage{I}{90}

(d)  \underline{Let}  $R$  \underline{be noetherian and}  $A$  \underline{a local} R-\underline{algebra of essentially finite} \underline{type. We denote by}  $\hat{A}$  \underline{the} \underline{m}-\underline{adic completion of}  $A$  \underline{where}  $\underline{m}$ = \underline{the maximal ideal of} $A$  .  \underline{Then for any} A-\underline{module}  $E$  \underline{of finite type we have a canonical isomorphism}

\begin{equation*}
D^1(\hat{A}|R, \ \hat{A} \underset{A}{\otimes} E) \xrightarrow{\ \sim\ } \hat{A} \underset{A}{\otimes} D^1(A|R, E)  \ .
\end{equation*}
\underline{Proof}:  (a)  follows immediately from  $D^1(A|R, M) = \mathrm{Exalcomm}(A|R, M)$ .

(b)  Let  $P$  be a polynomial algebra over  $R$  in a finite number of variables and let  $0 \to I \to P \to A \to 0$  be exact. Localizing if necessary, we may assume that  $I$  is generated by a P-regular sequence which would imply that  $H_1(K_I) = 0$ and therefore  $D^2(A|R, M) = \mathrm{Coker}(H^1(K_I, M) \to \mathrm{Hom}_A(H_1(K_I), M)) = 0$  . Conversely assume that  $D^2(A|R, M) = 0$  for every A-module  $M$  .

Let  $0 \to I \to P \to A \to 0$  be exact where  $P$  is a polynomial algebra over $R$  in a finite number of variables. We have

\SGAvar{\begin{equation*}
0 = D^2(A|R, M) = \mathrm{Coker}(H^1(K_I, M) \to \mathrm{Hom}_A(H_1(K_I), M))
\end{equation*}
i.e.,  $H^1(K_I, M) \to \mathrm{Hom}_A(H_1(K_I), M)$  is surjective for all A-module  $M$  . Choose an exact sequence  $0 \to K \to F \to I \to 0$  where  $F$  is a free P-module of finite type. We note that  $H_1(K_I) = K/K_0$  where  $K_0 = \mathrm{Im}(\overset{2}{\wedge} F \to \overset{1}{\wedge} F)$  . The fact that $D^2(A|R, H_1(K_I)) = 0$  i.e.,

\begin{equation*}
H^1(K_I, \ H_1(K_I)) \to \mathrm{Hom}_A(H_1(K_I), \ H_1(K_I))  \ \text{ is surjective}
\end{equation*}
entails that  $0 \to K/K_0 \to F/IF \to I/I^2 \to 0$  is exact and splits as A-modules and in particular  $I/I^2$  is projective as A-module. Therefore, localizing  $A$  if necessary, we may assume that  $I/I^2$  is  A-free, and that  $F/IF \to I/I^2$  is an isomorphism. Then the fact that  $0 \to K/K_0 \to F/IF \to I/I^2 \to 0$  is exact and splits as A-modules entails that  $H_1(K_I) = K/K_0 = 0$  . This means that  $I$  is generated by a P-regular sequence i.e.,  $A$  is a relative complete intersection over  $R$  .

}{\begin{equation*}
0 = D^2(A|R, M) = \mathrm{Coker}(H^1(K_I, M) \to \mathrm{Hom}_A(H_1(K_I), M))
\end{equation*}
i.e.,  $H^1(K_I, M) \to \mathrm{Hom}_A(H_1(K_I), M)$  is surjective for every A-module  $M$  . Choose an exact sequence  $0 \to K \to F \to I \to 0$  where  $F$  is a free P-module of finite type. We note that  $H_1(K_I) = K/K_0$  where  $K_0 = \mathrm{Im}(\overset{2}{\wedge} F \to \overset{1}{\wedge} F)$  . The fact that $D^2(A|R, H_1(K_I)) = 0$  i.e.,

\begin{equation*}
H^1(K_I, \ H_1(K_I)) \to \mathrm{Hom}_A(H_1(K_I), \ H_1(K_I))  \ \text{ is surjective}
\end{equation*}
entails that  $0 \to K/K_0 \to F/IF \to I/I^2 \to 0$  is exact and splits as A-modules and in particular  $I/I^2$  is projective as an A-module. Therefore, localizing  $A$  if necessary, we may assume that  $I/I^2$  is  A-free, and that  $F/IF \to I/I^2$  is an isomorphism. Then the fact that  $0 \to K/K_0 \to F/IF \to I/I^2 \to 0$  is exact and splits as A-modules entails that  $H_1(K_I) = K/K_0 = 0$  . This means that  $I$  is generated by a P-regular sequence i.e.,  $A$  is a relative complete intersection over  $R$  .

}{F7I-VI-070}(c)  follows immediately from 3.12(4) and (6).

(d)  Firstly it follows from 3.12(6) that

\SGAvar{\begin{equation*}
D^i(A\;R, \ \hat{A}\underset{A}{\otimes}E) \xrightarrow{\ \sim\ } \hat{A}\underset{A}{\otimes}D^i(A|R, E)
\end{equation*}
is an isomorphism for all  $i$  , and therefore it suffices to show that
\SGApage{I}{91}
$D^1(\hat{A}|R,\ \hat{A}\otimes_A E) \to D^1(A|R,\ \hat{A}\otimes_A E)$  is an isomorphism. Our proof is based on the description of  $D^1$  in terms of algebra extensions. Let  $0 \to \hat{A}\otimes_A E \to A' \xrightarrow{\ \varphi\ } A \to 0$  be an R-algebra extension. The fact that  $\hat{A}\otimes_A E$  is an $\hat{A}$-module of finite type entails readily that  $A'$  is a local R-algebra, complete with respect to $\underline{m}'$-adic topology where  $\underline{m}'$ = the maximal ideal of  $A'$ . Therefore if there is an R-algebra map  $\psi : A \to A'$  such that  $\varphi \circ \psi = \mathrm{id}_A$ , then  $\psi$  extends uniquely to  $\hat{\psi} : \hat{A} \to A'$  such that  $\varphi \circ \hat{\psi} = \mathrm{id}_{\hat{A}}$ . This proves that

\begin{equation*}
D^1(\hat{A}|R,\ \hat{A}\otimes_A E) \to D^1(A|R,\ \hat{A}\otimes_A E)
\end{equation*}
is injective. Now let  $0 \to \hat{A}\otimes_A E \to B \xrightarrow{\ \varphi\ } A \to 0$  be an R-algebra extension. Firstly we claim that the topology on  $\hat{E} = \hat{A}\otimes_A E$  induced from  $\underline{m}_B$-adic topology on  $B$  coincides with the topology on  $\hat{E}$  as $\hat{A}$-module. Indeed, since  $A$  is essentially of finite type over  $R$ , one can find a subalgebra  $B_0$  of  $B$  such that  $B_0$  is a local R-algebra of essentially finite type over  $R$  and that  $\varphi(B_0) = A$ . Set  $E_0 = E \cap B_0$  and denote by  $\underline{m}_B$ , $\underline{m}_0$ , $\underline{m}$  the maximal ideals of  $B$ , $B_0$ , $A$  respectively. Then  $\underline{m}_B = \underline{m}_0 + \hat{E}$  and hence  $\underline{m}_B^{n+1} = \underline{m}_0^{n+1} + \underline{m}^n\hat{E}$  for all  $n$ . Since  $B_0$  is noetherian, there exists an integer  $r > 0$  by Artin-Rees Theorem such that

\begin{equation*}
\begin{aligned}
\hat{E} \cap \underline{m}_B^{n+1} &= \underline{m}^n\hat{E} + \hat{E} \cap \underline{m}_0^{n+1} = \underline{m}^n\hat{E} + E_0 \cap \underline{m}_0^{n+1}\\
&= \underline{m}^n\hat{E} + \underline{m}^{n+1-r}(E_0 \cap \underline{m}_0^{\,r}) \subset \underline{m}^{n+1-r}\hat{E}
\end{aligned}
\end{equation*}
for all  $n \geq r$ , which proves our assertion about the two topologies on  $E$  . In particular  $\hat{E} = \hat{A}\otimes_A E$  is complete with respect to the topology induced from $\underline{m}_B$-adic topology on  $B$ , and hence we obtain the exact commutative diagram

}{\begin{equation*}
D^i(A|R, \ \hat{A}\underset{A}{\otimes}E) \xrightarrow{\ \sim\ } \hat{A}\underset{A}{\otimes}D^i(A|R, E)
\end{equation*}
is an isomorphism for all  $i$  , and therefore it suffices to show that
\SGApage{I}{91}
$D^1(\hat{A}|R,\ \hat{A}\otimes_A E) \to D^1(A|R,\ \hat{A}\otimes_A E)$  is an isomorphism. Our proof is based on the description of  $D^1$  in terms of algebra extensions. Let  $0 \to \hat{A}\otimes_A E \to A' \xrightarrow{\ \varphi\ } \hat{A} \to 0$  be an R-algebra extension whose restriction to $A$ splits. The fact that  $\hat{A}\otimes_A E$  is an $\hat{A}$-module of finite type entails readily that  $A'$  is a local R-algebra, complete with respect to $\underline{m}'$-adic topology where  $\underline{m}'$ = the maximal ideal of  $A'$ . Therefore if there is an R-algebra map  $\psi : A \to A'$  lifting the canonical map $A \to \hat{A}$ , then  $\psi$  extends uniquely to  $\hat{\psi} : \hat{A} \to A'$  such that  $\varphi \circ \hat{\psi} = \mathrm{id}_{\hat{A}}$ . This proves that

\begin{equation*}
D^1(\hat{A}|R,\ \hat{A}\otimes_A E) \to D^1(A|R,\ \hat{A}\otimes_A E)
\end{equation*}
is injective. Now let  $0 \to \hat{A}\otimes_A E \to B \xrightarrow{\ \varphi\ } A \to 0$  be an R-algebra extension. Firstly we claim that the topology on  $\hat{E} = \hat{A}\otimes_A E$  induced from  $\underline{m}_B$-adic topology on  $B$  coincides with the topology on  $\hat{E}$  as $\hat{A}$-module. Indeed, since  $A$  is essentially of finite type over  $R$ , one can find a subalgebra  $B_0$  of  $B$  such that  $B_0$  is a local R-algebra of essentially finite type over  $R$  and that  $\varphi(B_0) = A$ . Set  $E_0 = \hat{E} \cap B_0$  and denote by  $\underline{m}_B$ , $\underline{m}_0$ , $\underline{m}$  the maximal ideals of  $B$ , $B_0$ , $A$  respectively. Then  $\underline{m}_B = \underline{m}_0 + \hat{E}$  and hence  $\underline{m}_B^{n+1} = \underline{m}_0^{n+1} + \underline{m}^n\hat{E}$  for all  $n$ . Since  $B_0$  is noetherian, there exists an integer  $r > 0$  by Artin-Rees Theorem such that

\begin{equation*}
\begin{aligned}
\hat{E} \cap \underline{m}_B^{n+1} &= \underline{m}^n\hat{E} + \hat{E} \cap \underline{m}_0^{n+1} = \underline{m}^n\hat{E} + E_0 \cap \underline{m}_0^{n+1}\\
&= \underline{m}^n\hat{E} + \underline{m}^{n+1-r}(E_0 \cap \underline{m}_0^{\,r}) \subset \underline{m}^{n+1-r}\hat{E}
\end{aligned}
\end{equation*}
for all  $n \geq r$ , which proves our assertion about the two topologies on  $\hat{E}$  . In particular  $\hat{E} = \hat{A}\otimes_A E$  is complete with respect to the topology induced from $\underline{m}_B$-adic topology on  $B$ , and hence we obtain the exact commutative diagram

}{F7I-VI-071}\begin{center}
\begin{tikzcd}
0 \arrow[r] \& \hat{E} \arrow[r] \arrow[d, equal] \& B \arrow[r] \arrow[d] \& A \arrow[r] \arrow[d] \& 0 \\
0 \arrow[r] \& \hat{E} \arrow[r] \& \hat{B} \arrow[r] \& \hat{A} \arrow[r] \& 0
\end{tikzcd}
\end{center}












\SGAvar{where  $\hat{B}$  stands for  $\underline{m}_B$-adic completion of  $B$  . This proves that  $D^1$
$D^1(\hat{A}|R,\ \hat{E}) \to D^1(A|R,\ \hat{E})$  is surjective and hence we are done.


}{where  $\hat{B}$  stands for  $\underline{m}_B$-adic completion of  $B$  . This proves that
$D^1(\hat{A}|R,\ \hat{E}) \to D^1(A|R,\ \hat{E})$  is surjective and hence we are done.


}{F7I-VI-072}\SGApage{I}{92}

\underline{Corollary} 3.14. \underline{Let}  $R$  \underline{be noetherian and}  $A$  \underline{an R-algebra of finite type, smooth everywhere except at one closed point}  $x \in \mathrm{Spec}(A)$  . \underline{We then have isomorphisms}

\begin{equation*}
D^1(A|R,\ A) \simeq D^1(A_x|R,\ A_x) \simeq D^1(\hat{A}_x|R,\ \hat{A}_x)\ .
\end{equation*}
\underline{Proof}: Since  $A$  is smooth over  $R$  except at  $x \in \mathrm{Spec}(A)$  ,  $D^1(A|R,\ A)$  is an A-module of finite type with its support on the set  $\{x\}$  , and therefore we have

\begin{equation*}
D^1(A|R,\ A) = D^1(A|R,\ A)_x = \widehat{D^1(A|R,\ A)}_x\ .
\end{equation*}
Therefore our assertion follows from 3.13.

\SGAvar{\underline{Remark} 3.15.  Let  $R$  be topological ring and  $A$  a topological R-algebra. We then define, for any A-module  $M$  ,

\begin{equation*}
D^i_{\mathrm{top}}(A|R,\ E) = \varinjlim_{\alpha} D^i(A_\alpha|R_\alpha,\ A_\alpha \otimes_\alpha E)
\end{equation*}
where  $A_\alpha = A/J_\alpha$ , $R_\alpha = R/I_\alpha$  , and  $J_\alpha$ , $I_\alpha$  are both open in  $A, R$  respectively. Then  $A$  is formally smooth topological R-algebra means that  $D^1_{\mathrm{top}}(A|R,\ E) = 0$  for every A-module  $E$  .

}{\underline{Remark} 3.15.  Let  $R$  be a topological ring and  $A$  a topological R-algebra. We then define, for any A-module  $E$  ,

\begin{equation*}
D^i_{\mathrm{top}}(A|R,\ E) = \varinjlim_{\alpha} D^i(A_\alpha|R_\alpha,\ A_\alpha \otimes_A E)
\end{equation*}
where  $A_\alpha = A/J_\alpha$ , $R_\alpha = R/I_\alpha$  , and  $J_\alpha$ , $I_\alpha$  are both open in  $A, R$  respectively. Then to say that  $A$  is formally smooth as a topological R-algebra means that  $D^1_{\mathrm{top}}(A|R,\ E) = 0$  for every A-module  $E$  .

}{F7I-VI-073}        Now let  $X$  be a scheme over a ring  $R$  and  $\underline{E}$  a quasi-coherent  X-module. We then define the sheaves  $\underline{D}^i_{X|R}(\underline{E}) = \underline{D}^i(X|R,\ \underline{E})$  on  $X$  via presheaves on  $X$  given by

\SGAvar{\begin{equation*}
[\underline{D}^i_{X|R}(\underline{E})](U) = D^i(\Gamma(U,\ \underline{O}_X)|R,\ \Gamma(U,\ \underline{E}))
\end{equation*}
for every affine open set  $U \subset X$  . If  $R$  is noetherian and  $X$  is of finite type over  $R$  , then the property 3.13(a) shows that  $\underline{D}^i_{X|R}(\underline{E})$  can actually be constructed canonically, as a quasi-coherent X-module.

        In case when there is no possible confusion we shall abreviate by setting  $\underline{D}^i_{X|R} = \underline{D}^i_{X|R}(\underline{O}_X)$  . All the statements of 3.13 can now be translated in terms of these quasi-coherent sheaves  $\underline{D}^i_{X|R}(\underline{E})$  , which we leave to the readers. We list below a few properties which is most relevant to our purpose.

}{\begin{equation*}
[\underline{D}^i_{X|R}(\underline{E})](U) = D^i(\Gamma(U,\ \underline{O}_X)|R,\ \Gamma(U,\ \underline{E}))
\end{equation*}
for every affine open set  $U \subset X$  . If  $R$  is noetherian and  $X$  is of finite type over  $R$  , then the property 3.13(c) shows that  $\underline{D}^i_{X|R}(\underline{E})$  can actually be constructed canonically, as a quasi-coherent X-module.

        In case when there is no possible confusion we shall abbreviate by setting  $\underline{D}^i_{X|R} = \underline{D}^i_{X|R}(\underline{O}_X)$  . All the statements of 3.13 can now be translated in terms of these quasi-coherent sheaves  $\underline{D}^i_{X|R}(\underline{E})$  , which we leave to the readers. We list below a few properties which are most relevant to our purpose.

}{F7I-VI-074}\underline{Corollary} 3.15. \underline{Let}  $R$  \underline{be noetherian and}  $X$  \underline{an} R-\underline{scheme of finite type. Then}
\SGApage{I}{93}
(a)  \underline{For any coherent} X-\underline{module}  $\underline{E}$ , $\underline{D}^i_{X|R}(\underline{E})$  \underline{is again a coherent} X-\underline{module} \underline{for all}  $i$ , \underline{and furthermore}  $\operatorname{Supp} \underline{D}^i_{X|R}(\underline{E}) \subset X^{\mathrm{Sing}}$  \underline{for all}  $i > 0$  \underline{where} $X^{\mathrm{Sing}} = \{x \in X| X$ \underline{is non-smooth over}  $R$  \underline{at the point}  $x \in X\}$ .

(b)  \underline{If}  $X$  \underline{is flat over}  $R$ , \underline{then for any} R-\underline{algebra}  $S$ , \underline{we have the} \underline{canonical isomorphism}

\begin{equation*}
\underline{D}^i_{X\underset{R}{\otimes}S\,|\,S}(\underline{E}) \overset{\sim}{\to} \underline{D}^i_{X|R}(\underline{E})
\end{equation*}
\underline{for all}  $i$ .

(c)  \underline{If}  $X$  \underline{is a relative complete intersection over}  $R$ , \underline{then for every} \underline{surjective ring homomorphism}  $S \to R$  \underline{we have the exact sequence}

\begin{equation*}
0 \to \underline{D}^1_{X|R}(\underline{E}) \to \underline{D}^1_{X|S}(\underline{E}) \to \underline{\mathrm{Hom}}_{\underline{O}_X}(I/I^2\underset{R}{\otimes}\underline{O}_X, \underline{E}) \to 0
\end{equation*}
\underline{where}  $I = \mathrm{Ker}(S \to R)$ . \underline{If, furthermore,}  $R$  \underline{is local with the residue field}  $k$ \underline{and}  $X$  \underline{is} R-\underline{flat,} \underline{then we have the exact sequence}

\SGAvar{\begin{equation*}
0 \to \underline{D}^1_{X_0|h}(\underline{O}_{X_0}) \to \underline{D}^1_{X|S}(\underline{O}_{X_0}) \to D^1(R|S,\ k)\underset{k}{\otimes}\underline{O}_{X_0} \to 0
\end{equation*}
\underline{where}  $X_0 = X\underset{R}{\otimes}k$ .

\underline{Proof}:  (a)  Follows from 3.12(6) and (8) whereas (b) follows from 3.12(3).

(c)  Since  $X$  is a relative complete-intersection over  $R$ , we have $\underline{D}^i_{X|R}(\underline{E}) = 0$  for all  $i \geq 2$  by 3.12(10) and hence

}{\begin{equation*}
0 \to \underline{D}^1_{X_0|k}(\underline{O}_{X_0}) \to \underline{D}^1_{X|S}(\underline{O}_{X_0}) \to D^1(R|S,\ k)\underset{k}{\otimes}\underline{O}_{X_0} \to 0
\end{equation*}
\underline{where}  $X_0 = X\underset{R}{\otimes}k$ .

\underline{Proof}:  (a)  Follows from 3.12(7) and the vanishing for smooth algebras (cf. 3.13(a)), whereas (b) follows from 3.12(3).

(c)  Since  $X$  is a relative complete-intersection over  $R$ , we have $\underline{D}^i_{X|R}(\underline{E}) = 0$  for all  $i \geq 2$  by the complete-intersection vanishing (cf. 3.13(b)) and hence

}{F7I-VI-075}\begin{equation*}
0 \to \underline{D}^0_{X|R}(\underline{E}) \to \underline{D}^0_{X|S}(\underline{E}) \to \underline{D}^0_{R|S}(\underline{E}) \to \underline{D}^1_{X|R}(\underline{E}) \to \underline{D}^1_{X|S}(\underline{E}) \to \underline{D}^1_{R|S}(\underline{E}) \to 0
\end{equation*}
is exact by 3.12(2), where  $\underline{D}^i_{R|S}(\underline{E})$  are quasi-coherent sheaves of  $X$  defined by

\begin{equation*}
[\underline{D}^i_{R|S}(\underline{E})](U) = D^i(R|S,\ \underline{E}(U))
\end{equation*}
for every affine open subset  $U \subset X$ . In particular,  $\underline{D}^0_{R|S}(\underline{E}) = 0$  since  $S \to R$ is surjective, and

\SGAvar{\begin{equation*}
\underline{D}^1_{R|S}(\underline{E}) = \underline{\mathrm{Hom}}_{\underline{O}_X}(I/I^2\underset{R}{\otimes}\underline{O}_X, \underline{E})
\end{equation*}
by 3.12(7), and hence the first exact sequence. If  $X$  is R-flat, we have
\SGApage{I}{94}
$\underline{D}^i_{X_0|k}(\underline{E}) \overset{\sim}{\to} \underline{D}^i_{X|R}(\underline{E})$  for all  $i$  and for all quasi-coherent $X_0$-module  $\underline{E}$  and in particular

}{\begin{equation*}
\underline{D}^1_{R|S}(\underline{E}) = \underline{\mathrm{Hom}}_{\underline{O}_X}(I/I^2\underset{R}{\otimes}\underline{O}_X, \underline{E})
\end{equation*}
by 3.12(8), and hence the first exact sequence. If  $X$  is R-flat, we have
\SGApage{I}{94}
$\underline{D}^i_{X_0|k}(\underline{E}) \overset{\sim}{\to} \underline{D}^i_{X|R}(\underline{E})$  for all  $i$  and for every quasi-coherent $X_0$-module  $\underline{E}$  and in particular

}{F7I-VI-076}\begin{equation*}
\underline{D}^1_{X_0|k}(\underline{O}_{X_0}) \overset{\sim}{\to} \underline{D}^1_{X|R}(\underline{O}_{X_0})
\end{equation*}
from which the second exact sequence follows.

\SGAvar{\underline{Remark} 3.16:  The quasi-coherent sheaves  $\underline{D}^i_{X|R}(\underline{E})$  we have defined should appear as  $E_2^{0,i}$  of a spectral sequence relating the local ones and the global ones. L. Illusie has obtained such a theory (Lecture Notes in Mathematics 239). His method is via  "homotopical algebra"  generalizing the methods of M. Andre and D. Quillen rather than  "cohomological algebra" . It seems that the approach we have taken here can be globalized via Grothendieck cohomology on the category of schemes with an appropriate topology generalizing the one on the affine category.


}{\underline{Remark} 3.16:  The quasi-coherent sheaves  $\underline{D}^i_{X|R}(\underline{E})$  we have defined should appear as  $E_2^{0,i}$  of a spectral sequence relating the local ones and the global ones. L. Illusie has obtained such a theory (Lecture Notes in Mathematics 239). His method is via  "homotopical algebra"  generalizing the methods of M. André and D. Quillen rather than  "cohomological algebra" . It seems that the approach we have taken here can be globalized via Grothendieck cohomology on the category of schemes with an appropriate topology generalizing the one on the affine category.


}{F7I-VI-077}4.  \underline{Formal moduli of deformations}

One amongst the nice cofibered categories in which the isomorphism classes of objects are usually not prorepresentable, but nevertheless admit a formal versal object is the cofibered category of deformations. We now come back to the situations and notations of paragraph 1. Thus, throughout the section,  $\Lambda$  is a fixed complete local noetherian ring with the residue field  $k$ , and  $\underline{C}_{\Lambda}$  is the category of artinian local $\Lambda$-algebras with the same residue field  $k$ .

Let  $X_0$  be a fixed scheme over the field  $k$ . By an (infinitesimal) deformation of  $X_0$ , we mean a pair  $(R, X)$  where  $R \in \mathrm{ob}(\underline{C}_{\Lambda})$  and  $X$  is a flat R-scheme together with a commutative diagram

\begin{tikzcd}[column sep=huge, row sep=huge]
X_0 \arrow[r, "i_X"] \arrow[d] \& X \arrow[d, "\mathrm{flat}"] \\
\mathrm{Spec}(k) \arrow[r] \& \mathrm{Spec}(R)
\end{tikzcd}












\SGAvar{\SGApage{I}{95}
such that the induced morphism $X_0 \to X \times_{\mathrm{spec}(R)} \mathrm{Spec}(k)$ is an isomorphism. We note that, since $R$ is an artinian local ring, $i_X : X_0 \to X$ is necessarily a closed immersion which is a homeomorphism on underlying topological spaces. If $(R, X)$ , $(R', X')$ are deformations of $X_0$ to $R, R'$ respectively, we define an arrow $(R', X') \to (R, X)$ to consist of a map $R' \xrightarrow{\ \varphi\ } R$ in $\underline{C}_\Lambda$ and a morphism $\Phi : X \to X'$ of schemes with a commutative diagram

}{\SGApage{I}{95}
such that the induced morphism $X_0 \to X \times_{\mathrm{Spec}(R)} \mathrm{Spec}(k)$ is an isomorphism. We note that, since $R$ is an artinian local ring, $i_X : X_0 \to X$ is necessarily a closed immersion which is a homeomorphism on underlying topological spaces. If $(R, X)$ , $(R', X')$ are deformations of $X_0$ to $R, R'$ respectively, we define an arrow $(R', X') \to (R, X)$ to consist of a map $R' \xrightarrow{\ \varphi\ } R$ in $\underline{C}_\Lambda$ and a morphism $\Phi : X \to X'$ of schemes with a commutative diagram

}{F7I-VI-078}\begin{center}
\begin{tikzcd}[column sep=huge, row sep=large]
X \arrow[rr, "\Phi"] \arrow[ddd] \& \& X' \arrow[ddd] \\
 \& X_0 \arrow[ul, "i_X"'] \arrow[ur, "i_{X'}"] \arrow[d] \& \\
 \& \mathrm{Spec}(k) \arrow[dl] \arrow[dr] \& \\
\mathrm{Spec}(R) \arrow[rr, "\mathrm{Spec}(\varphi)"'] \& \& \mathrm{Spec}(R')
\end{tikzcd}
\end{center}












\SGAvar{We thus obtain the category $\underline{M}_{X_0}$ of (infinitesimal) deformations of $X_0$ in which the objects and the arrows are defined as above. The category $\underline{M}_{X_0}$ is provided with a functor $p : \underline{M}_{X_0} \to \underline{C}_\Lambda$ defined by $(R, X) \to R$ . Now let $(R, X) \in \mathrm{ob}(\underline{M}_{X_0})$ and a map $R \xrightarrow{\ \varphi\ } S$ in $\underline{C}_\Lambda$ be given. We then obtain a commutative diagram

}{We thus obtain the category $\underline{M}_{X_0}$ of (infinitesimal) deformations of $X_0$ in which the objects and the arrows are defined as above. The category $\underline{M}_{X_0}$ is provided with a functor $P : \underline{M}_{X_0} \to \underline{C}_\Lambda$ defined by $(R, X) \to R$ . Now let $(R, X) \in \mathrm{ob}(\underline{M}_{X_0})$ and a map $R \xrightarrow{\ \varphi\ } S$ in $\underline{C}_\Lambda$ be given. We then obtain a commutative diagram

}{F7I-VI-079}\begin{center}
\begin{tikzcd}[column sep=huge, row sep=huge]
X \times_{\mathrm{Spec}(R)} \mathrm{Spec}(S) \arrow[r, "p_1"] \arrow[d, "\mathrm{flat}"'] \& X \arrow[d, "\mathrm{flat}"] \\
\mathrm{Spec}(S) \arrow[r, "\mathrm{Spec}(\varphi)"'] \& \mathrm{Spec}(R)
\end{tikzcd}
\end{center}









and $\varphi$-arrow $(\varphi, p_1) : (R, X) \to (S, X \times_{\mathrm{Spec}(R)} \mathrm{Spec}(S))$ is obviously cocartesian in $\underline{M}_{X_0}$ . It follows that the category $\underline{M}_{X_0}$ is a cofibered category over $\underline{C}_\Lambda$ via $P : \underline{M}_{X_0} \to \underline{C}_\Lambda$ . We observe the following facts:


\SGApage{I}{96}

4.1(a): Let $(1_R, \Phi) : (R, X_1) \to (R, X_2)$ be an arrow in $\underline{M}_{X_0}$ i.e., $\Phi : X_2 \to X_1$ is a morphism over $R$ such that $X_2 \otimes_R k \to X_1 \otimes_R k$ is an isomorphism. Since $R$ is artinian and $X_i$ is $R$-flat, it follows immediately that $\Phi$ is an isomorphism so that $(1_R, \Phi) : (R, X_1) \to (R, X_2)$ is an isomorphism in $\underline{M}_{X_0}$ . Consequently, $\underline{M}_{X_0}$ is a cofibered category in groupoids over $\underline{C}_\Lambda$ .

4.1(b): Suppose that $X_0$ is affine, say $X_0 = \mathrm{Spec}(A_0)$ . Then for any object $(R, X) \in \mathrm{ob}(\underline{M}_{X_0})$ , $X$ is also affine. Consequently, for any $R \in \mathrm{ob}(\underline{C}_\Lambda)$ , the category $\underline{M}_{X_0}(R)$ is canonically equivalent, via the sections, to the category in which an object is a flat algebra $A$ over $R$ together with an $R$-algebra map $j_A : A \to A_0$ inducing an isomorphism $A \otimes_R k \xrightarrow{\ \sim\ } A_0$ , and a map $\Phi : A_1 \to A_2$ is an $R$-algebra map with a commutative diagram

\begin{center}
\begin{tikzcd}[column sep=huge, row sep=huge]
A_1 \arrow[r, "\Phi"] \arrow[d, "j_{A_1}"'] \& A_2 \arrow[d, "j_{A_2}"] \\
A_0 \arrow[r, equal] \& A_0
\end{tikzcd}
\end{center}

4.1(c): Consider the cofibered category $\widehat{\underline{M}}_{X_0}$ in groupoids over $\widehat{\underline{C}}_\Lambda$ . By definition, the category $\widehat{\underline{M}}_{X_0}(R)$ for any $R \in \mathrm{ob}(\widehat{\underline{C}}_\Lambda)$ is canonically equivalent to the category of formal schemes $\mathfrak{X}$ , adic over the formal spectrum $\mathrm{Spf}(R)$ , such that $\mathfrak{X}_n = \mathfrak{X} \otimes_R R/\underline{m}^{n+1}$ is flat over $R/\underline{m}^{n+1}$ for every $n \geq 0$ (or what amounts to the same, if $X_0$ is noetherian, such that $\mathfrak{X}$ is flat over $\mathrm{Spf}(R)$) . Consequently, if $X_0$ is noetherian, $\widehat{\underline{M}}_{X_0}$ is canonically $\widehat{\underline{C}}_\Lambda$-equivalent to the category of deformations of $X_0$ over the objects in $\widehat{\underline{C}}_\Lambda$ , in the sense of formal schemes.

4.1(d): Let $(R, X) \in \mathrm{ob}(\underline{M}_{X_0})$ i.e., we have a commutative diagram
\SGApage{I}{97}
\[
\begin{tikzcd}
X_0 \arrow[r, "i_X"] \arrow[d] \& X \arrow[d, "\text{flat}"] \\
\mathrm{Spec}(k) \arrow[r] \& \mathrm{Spec}(R)
\end{tikzcd}
\]

such that  $X_0 \to X \times_{\mathrm{Spec}(R)} \mathrm{Spec}(k)$  is an isomorphism. Now let  $U_0 \subset X_0$  be a fixed open subscheme of  $X_0$ . Since  $i_X : X_0 \to X$  is a closed immersion which is a homeomorphism on the underlying topological spaces, we obtain a commutative diagram

\begin{center}
\begin{tikzcd}
U_0 \arrow[r, "i_X|U_0"] \arrow[d] \& X|U_0 \arrow[d, "\text{flat}"] \\
\mathrm{Spec}(k) \arrow[r] \& \mathrm{Spec}(R)
\end{tikzcd}
\end{center}









such that  $U_0 \to (X|U_0) \times_{\mathrm{Spec}(R)} \mathrm{Spec}(k) = X \otimes_R k|U_0$  is an isomorphism, so that $(R, X|U_0)$  is a deformation of  $U_0$  to  $R$ . Furthermore, for any arrow $(\varphi, \Phi) : (R, X) \to (R', X')$  in  $\underline{M}_{X_0}$ ,  $(\varphi, \Phi|U_0) : (R, X|U_0) \to (R', X'|U_0)$  is an arrow in  $\underline{M}_{U_0}$ . Consequently we have the restriction functor  $\underline{M}_{X_0} \to \underline{M}_{U_0}$  over $\underline{C}_\Lambda$ . One may also consider a localization as well as a formalization at a point. Namely let  $x_0$  be a fixed point in  $X_0$ . If  $(R, X) \in \mathrm{ob}(\underline{M}_{X_0})$ , then $\mathrm{Spec}(\underline{O}_{X,x_0})$  is a deformation of  $\mathrm{Spec}(\underline{O}_{X_0,x_0})$  to  $R$  and  $\mathrm{Spec}(\hat{\underline{O}}_{X,x_0})$  is a deformation of  $\mathrm{Spec}(\hat{\underline{O}}_{X_0,x_0})$  over  $R$ , where  $\hat{\ }$  stands for the completion with respect to the linear topology given by the powers of the maximal ideal. We thus obtain functors  $\underline{M}_{X_0} \to \underline{M}_{\mathrm{Spec}(\underline{O}_{X_0,x_0})} \to \underline{M}_{\mathrm{Spec}(\hat{\underline{O}}_{X_0,x_0})}$ .

Firstly we establish that  $\underline{M}_{X_0}$  is a homogeneous cofibered category over $\underline{C}_\Lambda$ .

\underline{Lemma} 4.2.  \underline{Consider a diagram}

\begin{center}
\begin{tikzcd}
 \& R' \arrow[d] \\
R \arrow[r, "\text{surjective}"'] \& \bar{R}
\end{tikzcd}
\end{center}

\underline{in}  $\underline{C}_\Lambda$ , \underline{and let}  $E, E'$  \underline{be flat modules over}  $R, R'$  \underline{respectively. Then for any} $R'$-\underline{homomorphism}  $E' \to \bar{E} = E \otimes_R \bar{R}$ , \underline{the fibre-product}  $E \times_{\bar{E}} E'$  \underline{is flat as} $R \times_{\bar{R}} R'$-\underline{module}


\SGApage{I}{98}

\underline{Proof}:  Let  $0 \to I \to R \to \bar{R} \to 0$  be exact so that  $0 \to I \to R \times_{\bar{R}} R' \to R' \to 0$  is exact. Since  $E$  is $R$-flat,  $0 \to I \otimes_R E \to E \to \bar{E} \to 0$  is exact and hence we obtain the exact sequence

\SGAvar{\begin{equation}\tag{*}
0 \to I \otimes_R E \to E \times_{\bar{E}} E' \to E' \to 0 \ .
\end{equation}
Set  $S = R \times_{\bar{R}} R$ . Since  $I$  is nilpotent,  $E \times_{\bar{E}} E'$  is flat as an $S$-module if and only if  $(E \times_{\bar{E}} E') \otimes_S R'$  is flat as an $R'$-module and  $\mathrm{Tor}_1^S (E \times_{\bar{E}} E', R') = 0$ . However, we note that  $I \otimes_R E = I \otimes_S (E \times_{\bar{E}} E')$  and  $E' = R' \otimes_S (E \times_{\bar{E}} E')$  and therefore the exact sequence (*) can be rewritten as

}{\begin{equation}\tag{*}
0 \to I \otimes_R E \to E \times_{\bar{E}} E' \to E' \to 0 \ .
\end{equation}
Set  $S = R \times_{\bar{R}} R'$ . Since  $I$  is nilpotent,  $E \times_{\bar{E}} E'$  is flat as an $S$-module if and only if  $(E \times_{\bar{E}} E') \otimes_S R'$  is flat as an $R'$-module and  $\mathrm{Tor}_1^S (E \times_{\bar{E}} E', R') = 0$ . However, we note that  $I \otimes_R E = I \otimes_S (E \times_{\bar{E}} E')$  and  $E' = R' \otimes_S (E \times_{\bar{E}} E')$  and therefore the exact sequence (*) can be rewritten as

}{F7I-VI-080}\begin{equation*}
0 \ \to \ I \otimes_S (E \times_{\bar{E}} E') \ \to \ E \times_{\bar{E}} E' \ \to \ R' \otimes_S (E \times_{\bar{E}} E') \ \to \ 0 \ .
\end{equation*}
This means that  $\mathrm{Tor}_1^S (R', E \times_{\bar{E}} E') = 0$  and  $R' \otimes_S (E \times_{\bar{E}} E') = E'$  is flat as $R'$-module, and therefore  $E \times_{\bar{E}} E'$  is $S$-flat.

\underline{Corollary} 4.3.  \underline{Given a diagram}

\begin{center}
\begin{tikzcd}[column sep=large]
(R,\ X) \arrow[dr, "{(\varphi,\ \Phi)}"'] \& \& (R',\ X') \arrow[dl, "{(\varphi',\ \Phi')}"] \\
 \& (\bar{R},\ \bar{X}) \&
\end{tikzcd}
\end{center}








\underline{in}  $\underline{M}_{X_0}$ , \underline{the amalgamated sum}  $X \underset{\bar{X}}{\amalg} X'$  \underline{(in the sense of schemes) is a flat scheme over}  $R \times_{\bar{R}} R'$ , \underline{provided}  $\varphi : R \to \bar{R}$  \underline{is surjective. Consequently,}  $\underline{M}_{X_0}$ \underline{is a homogeneous cofibered category in groupoids over}  $\underline{C}_\Lambda$ .

\SGAvar{\underline{Proof}:  For each affine open  $U_0 \subset X_0$ , we have by definition that $X \underset{\bar{X}}{\amalg} X'|U_0 \simeq \mathrm{Spec}(\Gamma(U_0,\ X) \times_{\Gamma(U_0,\bar{X})} \Gamma(U_0,\ X'))$ . Consequently, if  $\varphi : R \to \bar{R}$ is surjective,  $X \underset{\bar{X}}{\amalg} X'$  is a flat scheme over  $R \times_{\bar{R}} R'$  by 4.2, and the morphism $i_X \underset{i_{\bar{X}}}{\amalg} i_{X'} : X_0 \to X \underset{\bar{X}}{\phantom{\amalg}} X'$  induces an isomorphism  $X_0 \to (X \underset{\bar{X}}{\amalg} X') \otimes_S k$ , where $S = R \times_{\bar{R}} R'$ . Therefore  $(R \times_{\bar{R}} R' ,\ X \underset{\bar{X}}{\amalg} X') \in \mathrm{ob}(\underline{M}_{X_0})$  and we obviously have that  $(R \times_{\bar{R}} R',\ X \underset{\bar{X}}{\amalg} X') = (R,\ X) \times_{(\bar{R},\bar{X})} (R',\ X')$  in the category  $\underline{M}_{X_0}$ , which
\SGApage{I}{99}
means that $\underline{M}_{X_0}$ is homogeneous over $\underline{C}_\Lambda$ .

}{\underline{Proof}:  For each affine open  $U_0 \subset X_0$ , we have by definition that $X \underset{\bar{X}}{\amalg} X'|U_0 \simeq \mathrm{Spec}(\Gamma(U_0,\ X) \times_{\Gamma(U_0,\bar{X})} \Gamma(U_0,\ X'))$ . Consequently, if  $\varphi : R \to \bar{R}$ is surjective,  $X \underset{\bar{X}}{\amalg} X'$  is a flat scheme over  $R \times_{\bar{R}} R'$  by 4.2, and the morphism $i_X \underset{i_{\bar{X}}}{\amalg} i_{X'} : X_0 \to X \underset{\bar{X}}{\amalg} X'$  induces an isomorphism  $X_0 \to (X \underset{\bar{X}}{\amalg} X') \otimes_S k$ , where $S = R \times_{\bar{R}} R'$ . Therefore  $(R \times_{\bar{R}} R' ,\ X \underset{\bar{X}}{\amalg} X') \in \mathrm{ob}(\underline{M}_{X_0})$  and we obviously have that  $(R \times_{\bar{R}} R',\ X \underset{\bar{X}}{\amalg} X') = (R,\ X) \times_{(\bar{R},\bar{X})} (R',\ X')$  in the category  $\underline{M}_{X_0}$ , which
\SGApage{I}{99}
means that $\underline{M}_{X_0}$ is homogeneous over $\underline{C}_\Lambda$ .

}{F7I-VI-081}\underline{Lemma} 4.4.  \underline{For any scheme} $X_0$ \underline{over} $k$ , \underline{we have a canonical exact sequence}

\begin{equation*}
0 \;\to\; H^{1}(X_{0},\ \underline{D}^{0}_{X_{0}}) \;\to\; [\underline{M}_{X_{0}}(k[\varepsilon])] \;\to\; H^{0}(X_{0},\ \underline{D}^{1}_{X_{0}})
\end{equation*}
\underline{where} $\underline{D}^{i}_{X_{0}}$ \underline{are as defined in paragraph 3.}

\underline{Proof}:  (Special case)  Assume that $X_{0}$ is affine and let $\Gamma(X_{0},\ \underline{O}_{X_{0}}) = A_{0}$ . Then as it was noted in 4.1(b), the category $\underline{M}_{X_{0}}(k[\varepsilon])$ is canonically equivalent to the category $\underline{E}$ , where

\SGAvar{\begin{equation*}
\mathrm{ob}(\underline{E}) : \text{object in } \underline{E} \text{ is a flat algebra } A \text{ over } k[\varepsilon]
\end{equation*}
together with a  $k[\varepsilon]$ -algebra map  $j_{A} : A \to A_{0}$  inducing the isomorphism $k \otimes_{k[\varepsilon]} j_{A} : k \otimes_{k[\varepsilon]} A \overset{\sim}{\relbar\joinrel\relbar\joinrel\relbar} A_{0}$ . Alternatively, an object in $\underline{E}$ is a $k$-algebra $A$  together with an exact sequence  $0 \to A_{0} \overset{\varepsilon}{\relbar\joinrel\relbar\joinrel\relbar} A \overset{j_{A}}{\relbar\joinrel\relbar\joinrel\relbar} A_{0} \to 0$  in which $\varepsilon(A_{0})^{2} = 0$  as an ideal in  $A$ .

}{\begin{equation*}
\mathrm{ob}(\underline{E}) : \text{object in } \underline{E} \text{ is a flat algebra } A \text{ over } k[\varepsilon]
\end{equation*}
together with a  $k[\varepsilon]$ -algebra map  $j_{A} : A \to A_{0}$  inducing the isomorphism $k \otimes_{k[\varepsilon]} j_{A} : k \otimes_{k[\varepsilon]} A \overset{\sim}{\longrightarrow} A_{0}$ . Alternatively, an object in $\underline{E}$ is a $k$-algebra $A$  together with an exact sequence  $0 \to A_{0} \overset{\varepsilon}{\longrightarrow} A \overset{j_{A}}{\longrightarrow} A_{0} \to 0$  in which $\varepsilon(A_{0})^{2} = 0$  as an ideal in  $A$ .

}{F7I-VI-082}F$\ell$($\underline{E}$) :  An arrow  $A_{1} \to A_{2}$  in  $\underline{E}$  with respect to first description of the object in  $\underline{E}$  is a  $k[\varepsilon]$-algebra map  $\Phi : A_{1} \to A_{2}$  such that  $j_{A_{1}} = j_{A_{2}} \circ \Phi$ . Therefore, an arrow  $A_{1} \to A_{2}$  in  $\underline{E}$  with respect to the alternative description of objects in  $\underline{E}$  is a $k$-algebra map  $\Phi : A_{1} \to A_{2}$  yielding a commutative diagram

\begin{center}
\begin{tikzcd}[column sep=large]
0 \arrow[r] \& A_{0} \arrow[r, "\varepsilon"] \arrow[d, equal] \& A_{1} \arrow[r, "j_{A_{1}}"] \arrow[d, "\Phi"] \& A_{0} \arrow[r] \arrow[d, equal] \& 0 \\
0 \arrow[r] \& A_{0} \arrow[r] \& A_{2} \arrow[r, "j_{A_{2}}"] \& A_{0} \arrow[r] \& 0
\end{tikzcd}\quad .
\end{center}
















\SGAvar{It follows that  $[\underline{M}_{X_{0}}(k[\varepsilon])]$  is canonically isomorphic, via  $\Gamma(X_{0},\ -)$ , to the set of isomorphism classes of $k$-algebra extensions of  $A_{0}$  by  $A_{0}$ . In other words, we have a canonical bijection  $[\underline{M}_{X_{0}}(k[\varepsilon])] \overset{\sim}{\relbar\joinrel\relbar\joinrel\relbar} H^{0}(X_{0},\ \underline{D}^{1}_{X_{0}})$  via $\Gamma(X_{0},\ -)$, which is trivially seen to respect the vector space structures on both sides.


}{It follows that  $[\underline{M}_{X_{0}}(k[\varepsilon])]$  is canonically isomorphic, via  $\Gamma(X_{0},\ -)$ , to the set of isomorphism classes of $k$-algebra extensions of  $A_{0}$  by  $A_{0}$ . In other words, we have a canonical bijection  $[\underline{M}_{X_{0}}(k[\varepsilon])] \overset{\sim}{\longrightarrow} H^{0}(X_{0},\ \underline{D}^{1}_{X_{0}})$  via $\Gamma(X_{0},\ -)$, which is trivially seen to respect the vector space structures on both sides.


}{F7I-VI-083}\SGApage{I}{100}

\SGAvar{(General case)  For each affine open  $U_{0} \subset X_{0}$ , we get a canonical map $[\underline{M}_{X_{0}}(k[\varepsilon])] \to H^{0}(U_{0},\ \underline{D}^{1}_{X_{0}})$  defined as the composite map $[\underline{M}_{X_{0}}(k[\varepsilon])] \to [\underline{M}_{U_{0}}(k[\varepsilon])] \overset{\sim}{\relbar\joinrel\relbar\joinrel\relbar} H^{0}(U_{0},\ \underline{D}^{1}_{X_{0}})$  and hence we obtain a canonical map $[\underline{M}_{X_{0}}(k[\varepsilon])] \to H^{0}(X_{0},\ \underline{D}^{1}_{X_{0}})$ . Its kernel is the set of isomorphism classes of locally-trivial deformations of  $X_{0}$  to  $k[\varepsilon]$ , i.e., the set of isomorphism classes of deformations of  $X_{0}$  to  $k[\varepsilon]$ , locally (on $X$ ) isomorphic to $X_{0} \otimes_{k} k[\varepsilon]$ . Since  $\underline{D}^{0}_{X_{0}} = \underline{\mathrm{Der}}_{k}\,(\underline{O}_{X_{0}},\ \underline{O}_{X_{0}}) = \underline{\mathrm{Aut}}(X_{0} \otimes_{k} k[\varepsilon]|X_{0})$  $=$ the sheaf of germs of automorphisms of  $X_{0} \otimes_{k} k[\varepsilon]$  inducing the identity on the subscheme $X_{0}$ , it is canonically isomorphic to  $H^{1}(X_{0},\ \underline{D}^{0}_{X_{0}})$ , so that we obtain the desired exact sequence.

}{(General case)  For each affine open  $U_{0} \subset X_{0}$ , we get a canonical map $[\underline{M}_{X_{0}}(k[\varepsilon])] \to H^{0}(U_{0},\ \underline{D}^{1}_{X_{0}})$  defined as the composite map $[\underline{M}_{X_{0}}(k[\varepsilon])] \to [\underline{M}_{U_{0}}(k[\varepsilon])] \overset{\sim}{\longrightarrow} H^{0}(U_{0},\ \underline{D}^{1}_{X_{0}})$  and hence we obtain a canonical map $[\underline{M}_{X_{0}}(k[\varepsilon])] \to H^{0}(X_{0},\ \underline{D}^{1}_{X_{0}})$ . Its kernel is the set of isomorphism classes of locally-trivial deformations of  $X_{0}$  to  $k[\varepsilon]$ , i.e., the set of isomorphism classes of deformations of  $X_{0}$  to  $k[\varepsilon]$ , locally (on $X$ ) isomorphic to $X_{0} \otimes_{k} k[\varepsilon]$ . Since  $\underline{D}^{0}_{X_{0}} = \underline{\mathrm{Der}}_{k}\,(\underline{O}_{X_{0}},\ \underline{O}_{X_{0}}) = \underline{\mathrm{Aut}}(X_{0} \otimes_{k} k[\varepsilon]|X_{0})$  $=$ the sheaf of germs of automorphisms of  $X_{0} \otimes_{k} k[\varepsilon]$  inducing the identity on the subscheme $X_{0}$ , it is canonically isomorphic to  $H^{1}(X_{0},\ \underline{D}^{0}_{X_{0}})$ , so that we obtain the desired exact sequence.

}{F7I-VI-084}\underline{Theorem} 4.5.  \underline{Let} $X_{0}$ \underline{be a scheme of finite type over} $k$ \underline{for which either} (i) $X_{0}$ \underline{is proper over} $k$ \underline{or} (ii) $X_{0}$ \underline{is affine but} $X_{0}^{\mathrm{sing}} = \{x \in X_{0}|X_{0}$ \underline{is non-smooth over} $k$ \underline{at the point} $x \}$ \underline{is a finite set. Then}

(a)  $\underline{M}_{X_{0}}$  \underline{admits a formal versal object} (\underline{called a formal versal deformation of}  $X_{0}$ ) \underline{over}  $\underline{C}_{\Lambda}$ . \underline{If}  $(R,\ \mathfrak{X})$  \underline{is a formal versal deformation of}  $X_{0}$ , \underline{then for any} $(S,\ \mathfrak{Y}) \in \mathrm{ob}(\widehat{\underline{M}}_{X_{0}})$ , \underline{there is a map}  $(\varphi,\Phi) : (R,\mathfrak{X}) \to (S,\mathfrak{Y})$  \underline{in}  $\widehat{\underline{M}}_{X_{0}}$  \underline{such that the diagram below is cartesian}

\begin{center}
\begin{tikzcd}[column sep=huge, row sep=large]
\mathfrak{Y} \arrow[r, "\Phi"] \arrow[d] \& \mathfrak{X} \arrow[d] \\
\mathrm{Spf}(S) \arrow[r, "\mathrm{Spf}(\varphi)"'] \& \mathrm{Spf}(R)
\end{tikzcd}\quad .
\end{center}









\underline{Furthermore a formal versal deformation}  $(R,\mathfrak{X})$  of  $X_{0}$  \underline{prorepresents the functor} $[\underline{M}_{X_{0}}]$  \underline{if and only if, for every surjective map}  $R' \to R$  \underline{in}  $\underline{C}_{\Lambda}$  \underline{and a deformation} $(R',\ X')$  \underline{of}  $X_{0}$  \underline{to}  $R'$ , \underline{the homomorphism}

\begin{equation*}
\mathrm{Aut}_{R'}(X'|X_{0}) \to \mathrm{Aut}_{R}(X' \underset{R'}{\otimes} R|X_{0})
\end{equation*}
\underline{is surjective.}

(b)  \underline{If}  $X_{0}$  \underline{is proper over}  $k$ , \underline{then}  $\underline{M}_{X_{0}}$  \underline{is smoothly pseudo-representable}
\SGApage{I}{101}
\underline{over} $\underline{C}_{\Lambda}$ .

\SGAvar{\underline{Proof}:  (a)  Since  $\underline{M}_{X_{0}}$  is a homogeneous cofibered category in groupoids over $\underline{C}_{\Lambda}$ by 4.3, it suffices to see that  $[\underline{M}(k[\epsilon])]$  is a finite-dimensional vector space over  $k$  by 1.11 and 2.9(b).

}{\underline{Proof}:  (a)  Since  $\underline{M}_{X_{0}}$  is a homogeneous cofibered category in groupoids over $\underline{C}_{\Lambda}$ by 4.3, it suffices to see that  $[\underline{M}_{X_{0}}(k[\epsilon])]$  is a finite-dimensional vector space over  $k$  by 1.11 and 2.9(b).

}{F7I-VI-085}Now consider the exact sequence

\SGAvar{\begin{equation}\tag{*}
0 \to H^{1}(X_{0}, \underline{D}^{0}_{X_{0}}) \to [\underline{M}_{X_{0}}(k[\epsilon])] \to H^{0}(X_{0}, \underline{D}^{1}_{X_{0}}) \ .
\end{equation}
If  $X_{0}$  is proper over  $k$  , then  $H^{1}(X_{0}, \underline{D}^{0}_{X_{0}})$  and  $H^{0}(X_{0}, \underline{D}^{1}_{X_{0}})$  are both finite dimensional since  $\underline{D}^{i}_{X_{0}}$  are coherent for all  $i \geq 0$  by 2.13(a). Now assume that $X_{0}$  is affine such that  $X_{0}^{\mathrm{sing}}$  is a finite set. Since  $\underline{D}^{1}_{X_{0}}$  is a coherent sheaf on  $X_{0}$  such that  $\mathrm{Supp}\ \underline{D}^{1}_{X_{0}} \subset X_{0}^{\mathrm{sing}}$  by 2.13(a), we find that  $H^{0}(X_{0}, \underline{D}^{1}_{X_{0}})$  is a finite-dimensional vector space over  $k$  , whereas  $H^{1}(X_{0}, \underline{D}^{0}_{X_{0}}) = 0$  since  $X_{0}$ is affine. Therefore, in either case, it follows from the exact sequence (*) that $[\underline{M}_{X_{0}}(k[\epsilon])]$  is a finite dimensional vector space over  $k$  . The second statement of (a) follows from 2.7.

}{\begin{equation}\tag{*}
0 \to H^{1}(X_{0}, \underline{D}^{0}_{X_{0}}) \to [\underline{M}_{X_{0}}(k[\epsilon])] \to H^{0}(X_{0}, \underline{D}^{1}_{X_{0}}) \ .
\end{equation}
If  $X_{0}$  is proper over  $k$  , then  $H^{1}(X_{0}, \underline{D}^{0}_{X_{0}})$  and  $H^{0}(X_{0}, \underline{D}^{1}_{X_{0}})$  are both finite dimensional since  $\underline{D}^{i}_{X_{0}}$  are coherent for all  $i \geq 0$  by 3.15(a). Now assume that $X_{0}$  is affine such that  $X_{0}^{\mathrm{sing}}$  is a finite set. Since  $\underline{D}^{1}_{X_{0}}$  is a coherent sheaf on  $X_{0}$  such that  $\mathrm{Supp}\ \underline{D}^{1}_{X_{0}} \subset X_{0}^{\mathrm{sing}}$  by 3.15(a), we find that  $H^{0}(X_{0}, \underline{D}^{1}_{X_{0}})$  is a finite-dimensional vector space over  $k$  , whereas  $H^{1}(X_{0}, \underline{D}^{0}_{X_{0}}) = 0$  since  $X_{0}$ is affine. Therefore, in either case, it follows from the exact sequence (*) that $[\underline{M}_{X_{0}}(k[\epsilon])]$  is a finite dimensional vector space over  $k$  . The second statement of (a) follows from 2.7.

}{F7I-VI-086}          (b)  If  $X_{0}$  is proper,  $\mathrm{Aut}\ (X_{0} \otimes_{k} k[\epsilon]|X_{0}) = H^{0}(X_{0}, \underline{D}^{0}_{X_{0}})$  is a finite dimensional vector space and hence our statement follows from 2.13.

\underline{Remark 4.6(a)}:  Let  $X_{0}$  be affine. Then a formal versal deformation of  $X_{0}$ exists if and only if  $\underline{D}^{1}_{X_{0}}$  has its support on a finite number of points. However this does not imply that non-smooth points of  $X_{0}$  are all isolated. There are affine schemes  $X_{0}$  over $k$  having a non-isolated non-smooth point but  $\mathrm{supp}\ \underline{D}^{1}_{X_{0}}$ is a finite set.

\underline{Remark 4.6(b)}:  Henceforth a deformation of  $X_{0}$  is meant by an object in  $\underline{M}_{X_{0}}$  . In case when there is a possible confusion, an object in  $\underline{M}_{X_{0}}$  will be referred as an infinitesimal deformation of  $X_{0}$  .

\underline{Remark 4.6(c)}:  Let  $X_{0}$  be a scheme over  $k$  . We have fixed a complete local noetherian ring  $\Lambda$  with the residue field  $k$  , and considered the deformations
\SGApage{I}{102}
of  $X_{0}$  over the category  $\underline{C}_{\Lambda}$  . (In practice one usually takes  $\Lambda = k$  or  $W(k)$ if  $k$  is a perfect field of positive characteristic). Now  $\underline{C}_{k} \subset \underline{C}_{\Lambda}$  is a fully faithful imbedding, and that  $\underline{M}_{X_{0}}|\underline{C}_{k}$  is simply the category of deformations of $X_{0}$  over  $S$  where  $S \in \mathrm{ob}(\underline{C}_{k})$  . Therefore if  $(R, \mathfrak{X})$  is a formal versal deformation of  $X_{0}$  over  $\underline{C}_{\Lambda}$  , then  $(k \otimes_{\Lambda} R, k \otimes_{\Lambda} \mathfrak{X})$  is a formal versal deformation of $X_{0}$  over  $\underline{C}_{k}$  . It follows, for example, that  $\dim_{\underline{C}_{\Lambda}} \underline{M}_{X_{0}} = \dim k \otimes_{\Lambda} R$  does not depend on the choice of  $\Lambda$  but depends only on  $X_{0}$  .

          One can generalize the above theorem in many ways. Let us for instance consider a pair  $Y_{0} \to X_{0}$  where  $X_{0}$  is a scheme of finite type over  $k$  and  $Y_{0}$ a fixed closed subscheme of  $X_{0}$  . For each  $R \in \mathrm{ob}(\underline{C}_{\Lambda})$  , a deformation of $(X_{0}, Y_{0})$  to  $R$  is meant to be a commutative diagram

\SGAvar{\begin{tikzcd}[row sep=large, column sep=huge]
Y_0 \arrow[rrr] \arrow[dr, hook] \arrow[ddr] \& \& \& Y \arrow[dl, hook] \arrow[ddl, "\text{flat}"] \\
 \& X_0 \arrow[r] \arrow[d] \& X \arrow[d, no head, "\text{flat}"] \& \\
 \& \mathrm{Spec}(k) \arrow[r] \& \mathrm{Spec}(R) \&
\end{tikzcd}


















}{\begin{tikzcd}[row sep=large, column sep=huge]
Y_0 \arrow[rrr] \arrow[dr, hook] \arrow[ddr] \& \& \& Y \arrow[dl, hook] \arrow[ddl, "\text{flat}"] \\
 \& X_0 \arrow[r] \arrow[d] \& X \arrow[d, "\text{flat}"] \& \\
 \& \mathrm{Spec}(k) \arrow[r] \& \mathrm{Spec}(R) \&
\end{tikzcd}


















}{F7I-VI-087}such that  $(X_{0}, Y_{0}) \to (X \times_{\mathrm{Spec}(R)} \mathrm{Spec}(k), Y \times_{\mathrm{Spec}(R)} \mathrm{Spec}(k))$  is an isomorphism. Defining a morphism  $(R'; X', Y') \to (R; X, Y)$  of deformations of  $(X_{0}, Y_{0})$  to be the obvious ones, we obtain the cofibered category  $\underline{M}_{X_{0},Y_{0}}$  in groupoids over $\underline{C}_{\Lambda}$  . (This is a process often used in rigidification of automorphisms in moduli problems of deformations.) One sees immediately, by the same argument as in 4.3 based on 4.2, that  $\underline{M}_{X_{0},Y_{0}}$  is a homogeneous cofibered category in groupoids over $\underline{C}_{\Lambda}$  .

\underline{Lemma} 4.7.  $0 \to H^{0}(Y_{0}, \underline{\mathrm{Hom}}_{\underline{O}_{Y_{0}}} (\underline{I}/\underline{I}^{2}, \underline{O}_{Y_{0}})) \to [\underline{M}_{X_{0},Y_{0}}(k[\epsilon])] \to [\underline{M}_{X_{0}}(k[\epsilon])]$ \underline{is exact, where}  $\underline{I}$  \underline{is the ideal sheaf of}  $\underline{O}_{X_{0}}$  \underline{defining}  $Y_{0}$  i.e. 


\SGApage{I}{103}

\SGAvar{\begin{equation*}
 0 \to \underline{I} \to \underline{O}_{X_0} \to \underline{O}_{Y_0} \to 0 \qquad \underline{\text{is exact,}}
\end{equation*}
\underline{Proof}: Let $\underline{B}$ be the full subcategory of $\underline{M}_{X_0,Y_0}(k[\epsilon])$ whose objects are of the form $(k[\epsilon]; X \times_{\mathrm{Spec}(k)} \mathrm{Spec}(k[\epsilon]), Y)$ . We then have $[B] = \mathrm{Ker}\,([\underline{M}_{X_0,Y_0}(k[\epsilon])] \to [\underline{M}_{X_0}(k[\epsilon])]\,)$ . The category $\underline{B}$ is the groupoid of deformations of $Y_0$ to $k[\epsilon]$ "inside $X_0 \times_{\mathrm{Spec}(k)} \mathrm{Spec}(k[\epsilon])$ ", i.e., an object in $\underline{B}$ is a commutative diagram

}{\begin{equation*}
 0 \to \underline{I} \to \underline{O}_{X_0} \to \underline{O}_{Y_0} \to 0 \qquad \underline{\text{is exact,}}
\end{equation*}
\underline{Proof}: Let $\underline{B}$ be the full subcategory of $\underline{M}_{X_0,Y_0}(k[\epsilon])$ whose objects are of the form $(k[\epsilon]; X_0 \times_{\mathrm{Spec}(k)} \mathrm{Spec}(k[\epsilon]), Y)$ . We then have $[B] = \mathrm{Ker}\,([\underline{M}_{X_0,Y_0}(k[\epsilon])] \to [\underline{M}_{X_0}(k[\epsilon])]\,)$ . The category $\underline{B}$ is the groupoid of deformations of $Y_0$ to $k[\epsilon]$ "inside $X_0 \times_{\mathrm{Spec}(k)} \mathrm{Spec}(k[\epsilon])$ ", i.e., an object in $\underline{B}$ is a commutative diagram

}{F7I-VI-088}\begin{center}
\begin{tikzcd}[column sep=large, row sep=large]
Y \arrow[rr, "i_Y"] \arrow[dr, "\text{flat}"'] \& \& X_0 \times_{\mathrm{Spec}(k)} \mathrm{Spec}(k[\epsilon]) \arrow[dl] \\
 \& \mathrm{Spec}(k[\epsilon]) \&
\end{tikzcd}
\end{center}










such that $\mathrm{Spec}(k) \times_{\mathrm{Spec}(k[\epsilon])} i_Y = i_{Y_0}$ , and an arrow $(Y_1, i_{Y_1}) \to (Y_2, i_{Y_2})$ is a $k[\epsilon]$-morphism $\Phi : Y_1 \to Y_2$ (necessarily an isomorphism) such that $\mathrm{Spec}(k) \times_{\mathrm{Spec}(k[\epsilon])} \Phi = \mathrm{id}_{Y_0}$ and $i_{Y_2} \circ \Phi = i_{Y_1}$ . Alternatively, an object in $\underline{B}$ is a pair $(Y, j_Y)$ , where $Y$ is a deformation of $Y_0$ to $k[\epsilon]$ and $j_Y : Y \to X_0$ is a morphism over $k$ such that the composite morphism $Y_0 \to Y \xrightarrow{\;j_Y\;} X_0$ is $i_{Y_0}$ , and an arrow $(Y_1, j_{Y_1}) \to (Y_2, j_{Y_2})$ is a morphism $\Phi : Y_1 \to Y_2$ such that $k \otimes_{k[\epsilon]} \Phi = \mathrm{id}_{Y_0}$ and $j_{Y_2} \circ \Phi = j_{Y_1}$ . Therefore the association

\SGAvar{$(Y, j_Y) \;\longmapsto\;$
\begin{tikzcd}[column sep=large]
0 \arrow[r] \& \underline{I} \arrow[r] \arrow[d] \& \underline{O}_{X_0} \arrow[r] \arrow[d, "\mathfrak{J}_Y"] \& \underline{O}_{Y_0} \arrow[r] \arrow[d, equal] \& 0 \\
0 \arrow[r] \& \underline{O}_{Y_0} \arrow[r, "\epsilon"] \& \underline{O}_Y \arrow[r] \& \underline{O}_{Y_0} \arrow[r] \& 0
\end{tikzcd}















}{$(Y, j_Y) \;\longmapsto\;$
\begin{tikzcd}[column sep=large]
0 \arrow[r] \& \underline{I} \arrow[r] \arrow[d] \& \underline{O}_{X_0} \arrow[r] \arrow[d, "j_Y"] \& \underline{O}_{Y_0} \arrow[r] \arrow[d, equal] \& 0 \\
0 \arrow[r] \& \underline{O}_{Y_0} \arrow[r, "\epsilon"] \& \underline{O}_Y \arrow[r] \& \underline{O}_{Y_0} \arrow[r] \& 0
\end{tikzcd}















}{F7I-VI-089}gives rise to an isomorphism

\begin{equation*}
[B] \;\xrightarrow{\;j|\underline{I}\;}\; \mathrm{Hom}_{\underline{O}_{X_0}}(\underline{I},\, \underline{O}_{Y_0}) = \mathrm{Hom}_{\underline{O}_{Y_0}}(\underline{I}/\underline{I}^2,\, \underline{O}_{Y_0}) = H^0(Y_0,\, \underline{\mathrm{Hom}}_{\underline{O}_{Y_0}}(\underline{I}/\underline{I}^2,\, \underline{O}_{Y_0}))\;.
\end{equation*}

\SGApage{I}{104}

\underline{Remark} 4.8.  For a more general formula implying the above lemma, see SGA3 III 4.4. One may also note the following fact. One can define for any $S$-scheme $X$ and a quasi-coherent $X$-module $\underline{E}$ a quasi-coherent $X$-module $\underline{D}^i_{X|S}(\underline{E})$ for all $i \geq 0$ , generalizing the case when $S$ is affine. As usual we simply write $\underline{D}^i_{X|S}(\underline{O}_X) = \underline{D}^i_{X|S}$ . We then have

\SGAvar{\begin{equation*}
H^0(Y_0,\, \underline{\mathrm{Hom}}_{Y_0}(\underline{I}/\underline{I}^2,\, \underline{O}_{Y_0})) = H^0(Y_0,\, \underline{D}^1_{Y_0|X})
\end{equation*}
where $0 \to \underline{I} \to \underline{O}_{X_0} \to \underline{O}_{Y_0} \to 0$ is exact.

}{\begin{equation*}
H^0(Y_0,\, \underline{\mathrm{Hom}}_{Y_0}(\underline{I}/\underline{I}^2,\, \underline{O}_{Y_0})) = H^0(Y_0,\, \underline{D}^1_{Y_0|X_0})
\end{equation*}
where $0 \to \underline{I} \to \underline{O}_{X_0} \to \underline{O}_{Y_0} \to 0$ is exact.

}{F7I-VI-090}\underline{Theorem} 4.9.  \underline{Let} $X_0$ \underline{be} $k$-\underline{scheme of finite type and} $Y_0$ \underline{a closed subscheme of} $X_0$ . \underline{Assume that either} (i) $X_0$ \underline{is proper over} $k$ \underline{or} (ii) $X_0$ \underline{is affine and smooth over} $k$ \underline{outside a finite number of points and} $Y_0$ \underline{is proper over} $k$ . \underline{We then have}

(a) $\underline{M}_{X_0,\,Y_0}$ \underline{admits a formal versal deformation over} $\underline{C}_\wedge$ , \underline{and it pro-represents the functor} $[\underline{M}_{X_0,\,Y_0}]$ \underline{if and only if, for every surjective map} $R' \to R$ \underline{in} $\underline{C}_\wedge$ \underline{and a deformation} $(R'; X', Y')$ ,

\SGAvar{\begin{equation*}
\mathrm{Aut}_{R'}((X', Y')|X_0) \to \mathrm{Aut}_R(X' \otimes_{R'} R,\, Y',\, \otimes_{R'} R)|X_0)
\end{equation*}
\underline{is surjective.}

}{\begin{equation*}
\mathrm{Aut}_{R'}((X', Y')|X_0) \to \mathrm{Aut}_R((X' \otimes_{R'} R,\, Y' \otimes_{R'} R)|X_0)
\end{equation*}
\underline{is surjective.}

}{F7I-VI-091}(b) \underline{If furthermore} $H^0(X_0,\, \underline{D}^0_{Y_0|k} \times_{\underline{D}^0_{X_0|k}(\underline{O}_{Y_0})} \underline{D}^0_{X_0|k})$ \underline{is finite dimensional,} \underline{then} $\underline{M}_{X_0,\,Y_0}$ \underline{is smoothly pseudo-representable.}

\underline{Proof}: It is obvious from 4.4 and 4.6 together with 1.11, 2.7, and 2.13.

\SGAvar{We now consider the passage to localization and completion in the deformation theory     of schemes. Firstly on a notation: Let $X_0$ be a $k$-scheme of finite type. Given a fixed point $x_0 \in X_0$ , we have natural functors

}{We now consider the passage to localization and completion in the deformation theory     of schemes. First, some notation: Let $X_0$ be a $k$-scheme of finite type. Given a fixed point $x_0 \in X_0$ , we have natural functors

}{F7I-VI-092}\begin{equation*}
\underline{M}_{X_0} \to \underline{M}_{\mathrm{Spec}(\underline{O}_{X_0,\, x_0})} \to \underline{M}_{\mathrm{Spec}(\widehat{\underline{O}}_{X_0,\, x_0})} 
\end{equation*}

\SGApage{I}{105}

 and in turn we obtain the functors

\SGAvar{\begin{equation*}
\widehat{\underline{M}}_{X_0} \to \widehat{\underline{M}}_{\mathrm{Spec}(\underline{O}_{X_0,\ x_0})} \to \widehat{\underline{M}}_{\mathrm{Spec}(\widehat{\underline{O}}_{X_0,\ x_0})}\ .
\end{equation*}
As it was remarked in 4.1(c), $\underline{M}_{X_0}$  is canonically equivalent to the category in which an object is a pair  $(R,\ \mathfrak{X})$  where  $R \in \mathrm{ob}(\underline{C}_{\Lambda})$  and  $\mathfrak{X}$  is a flat $R$-adic formal scheme together with a fixed isomorphism  $\mathfrak{X}_0 \longleftrightarrow X_0$ . Henceforth we shall write, when there is no possible confusion, the images of  $(R,\ \mathfrak{X})$  under $\underline{M}_{X_0} \to \underline{M}_{\mathrm{Spec}(\underline{O}_{X_0,\ x_0})}$  and  $\underline{M}_{X_0} \to \underline{M}_{\mathrm{Spec}(\underline{O}_{X_0,\ x_0})}$  by  $(R,\ \mathfrak{X}_{(x_0)})$  and  $(R,\ \mathfrak{X}_{(x_0)})$ respectively. Therefore if  $\mathfrak{X} = \varinjlim\limits_{\nu} \mathfrak{X}_{\nu}$  where  $\mathfrak{X}_{\nu} = \mathfrak{X} \underset{\mathrm{Spf}(R)}{\times} \mathrm{Spec}(R/\underline{m}^{\nu+1})$    then

\begin{equation*}
\mathfrak{X}_{(x_0)} = \varinjlim\limits_{\nu} \mathrm{Spec}(\underline{O}_{\mathfrak{X}_{\nu},\ x_0})
\end{equation*}
\begin{equation*}
\mathfrak{X}_{(\hat{X}_0)} = \varinjlim\limits_{\nu} \mathrm{Spec}(\widehat{\underline{O}}_{\mathfrak{X}_{\nu},\ x_0})\ .
\end{equation*}
\underline{Theorem} 4.10. \underline{Let}  $X_0$  \underline{be an affine scheme of finite type over}  $k$  .

}{\begin{equation*}
\widehat{\underline{M}}_{X_0} \to \widehat{\underline{M}}_{\mathrm{Spec}(\underline{O}_{X_0,\ x_0})} \to \widehat{\underline{M}}_{\mathrm{Spec}(\widehat{\underline{O}}_{X_0,\ x_0})}\ .
\end{equation*}
As was remarked in 4.1(c), $\widehat{\underline{M}}_{X_0}$  is canonically equivalent to the category in which an object is a pair  $(R,\ \mathfrak{X})$  where  $R \in \mathrm{ob}(\widehat{\underline{C}}_{\Lambda})$  and  $\mathfrak{X}$  is a flat $R$-adic formal scheme together with a fixed isomorphism  $\mathfrak{X}_0 \longleftrightarrow X_0$ . Henceforth we shall denote, when there is no possible confusion, the images of  $(R,\ \mathfrak{X})$  under $\widehat{\underline{M}}_{X_0} \to \widehat{\underline{M}}_{\mathrm{Spec}(\underline{O}_{X_0,\ x_0})}$  and  $\widehat{\underline{M}}_{X_0} \to \widehat{\underline{M}}_{\mathrm{Spec}(\widehat{\underline{O}}_{X_0,\ x_0})}$  by  $(R,\ \mathfrak{X}_{(x_0)})$  and  $(R,\ \mathfrak{X}_{(\hat{x}_0)})$ respectively. Therefore if  $\mathfrak{X} = \varinjlim\limits_{\nu} \mathfrak{X}_{\nu}$  where  $\mathfrak{X}_{\nu} = \mathfrak{X} \underset{\mathrm{Spf}(R)}{\times} \mathrm{Spec}(R/\underline{m}^{\nu+1})$    then

\begin{equation*}
\mathfrak{X}_{(x_0)} = \varinjlim\limits_{\nu} \mathrm{Spec}(\underline{O}_{\mathfrak{X}_{\nu},\ x_0})
\end{equation*}
\begin{equation*}
\mathfrak{X}_{(\hat{x}_0)} = \varinjlim\limits_{\nu} \mathrm{Spec}(\widehat{\underline{O}}_{\mathfrak{X}_{\nu},\ x_0})\ .
\end{equation*}
\underline{Theorem} 4.10. \underline{Let}  $X_0$  \underline{be an affine scheme of finite type over}  $k$  .

}{F7I-VI-093}(a)  \underline{If}  $X_0$  \underline{is locally a complete-intersection, then}  $\underline{M}_{X_0}$  \underline{is smooth over}  $\underline{C}_{\Lambda}$  .

(b)  \underline{Assume that, for a fixed closed point}  $x \in X_0$  ,  $X_0 - \{x\}$  \underline{is smooth over}  $k$  . \underline{Then}

\begin{equation*}
\underline{M}_{X_0} \to \underline{M}_{\mathrm{Spec}(\underline{O}_{X_0,\ x})} \to \underline{M}_{\mathrm{Spec}(\widehat{\underline{O}}_{X_0,\ x})}
\end{equation*}
\underline{are both minimally smooth functors. In particular, if}  $(R,\ \mathfrak{X})$  \underline{is a formal versal deformation of}  $X_0$  , \underline{then}  $(R,\ \mathfrak{X}_{(x)})$  \underline{and}  $(R,\ \mathfrak{X}_{(\hat{x})})$  \underline{are formal versal deformations of}  $\mathrm{Spec}(\underline{O}_{X_0,\ x})$  \underline{and}  $\mathrm{Spec}(\widehat{\underline{O}}_{X_0,\ x})$  \underline{respectively}.

\SGAvar{\underline{Proof}:  In this proof we go back to the notations of the cohomology of commutative algebras defined in the paragraph 3. We set  $X_0 = \mathrm{Spec}(A_0)$


}{\underline{Proof}:  In this proof we go back to the notations of the cohomology of commutative algebras defined in paragraph 3. We set  $X_0 = \mathrm{Spec}(A_0)$


}{F7I-VI-094}\SGApage{I}{106}

(a)  It suffices to show that  $P : [\underline{M}_{X_0}] \to \underline{C}_{\Lambda}$  is smooth.

\SGAvar{Let  $(R,\ \mathrm{Spec}(A))$  be a deformation of  $\mathrm{Spec}(A_0)$  to  $R$  , where  $R \in \mathrm{ob}(\underline{C}_{\Lambda})$  . Given a small extension  $R' \to R$  in  $\underline{C}_{\Lambda}$  (i.e.,  $0 \to k \to R' \to R \to 0$  is exact), we must show that  $(R,\ \mathrm{Spec}(A))$  can be lifted to  $R'$  . The exact sequence of  $D^i$  with coefficient  $A_0$  applied to the ring homomorphisms  $R' \to R \to A$  gives us the exact sequence

\begin{equation*}
\ldots \to D^1(A|R,\ A_0) \to D^1(A|R',\ A_0) \to D^1(R|R',\ k) \underset{k}{\otimes} A_0 \to D^2(A|R,\ A_0) \to \ldots
\end{equation*}
Since  $R' \to R$  is a small extension, $D^1(R'|R,\ k)$  is 1-dimensional vector space over  $k$  generated by the element (denoted by  $[R']$) represented by the extension $0 \to k \to R' \to R \to 0$  . On the other hand,  $A$  is $R$-flat and hence $D^i(A|R,\ A_0) = D^i(A_0|k,A_0)$  for all  $i > 0$  and in particular  $D^2(A|R,\ A_0) = 0$ by 3.12(11) since  $A_0$  is locally a complete-intersection. Consequently we get the exact sequence

}{Let  $(R,\ \mathrm{Spec}(A))$  be a deformation of  $\mathrm{Spec}(A_0)$  to  $R$  , where  $R \in \mathrm{ob}(\underline{C}_{\Lambda})$  . Given a small extension  $R' \to R$  in  $\underline{C}_{\Lambda}$  (i.e.,  $0 \to k \to R' \to R \to 0$  is exact), we must show that  $(R,\ \mathrm{Spec}(A))$  can be lifted to  $R'$  . The exact sequence for  $D^i$  with coefficients in  $A_0$  applied to the ring homomorphisms  $R' \to R \to A$  gives us the exact sequence

\begin{equation*}
\ldots \to D^1(A|R,\ A_0) \to D^1(A|R',\ A_0) \to D^1(R|R',\ k) \underset{k}{\otimes} A_0 \to D^2(A|R,\ A_0) \to \ldots
\end{equation*}
Since  $R' \to R$  is a small extension, $D^1(R|R',\ k)$  is a 1-dimensional vector space over  $k$  generated by the element (denoted by  $[R']$) represented by the extension $0 \to k \to R' \to R \to 0$  . On the other hand,  $A$  is $R$-flat and hence $D^i(A|R,\ A_0) = D^i(A_0|k,A_0)$  for all  $i > 0$  and in particular  $D^2(A|R,\ A_0) = 0$ by 3.12(8) since  $A_0$  is locally a complete-intersection. Consequently we get the exact sequence

}{F7I-VI-095}\begin{equation*}
\to D^1(A_0|k,\ A_0) \to D^1(A|R',\ A_0) \to D^1(R|R',\ k) \underset{k}{\otimes} A_0 \to 0
\end{equation*}
This means that there exists an $R'$-algebra extension of  $A$  by  $A_0$  , denoted by $A'$  , such that  $[A'] \to [R'] \otimes 1$  under the map  $D^1(A|R',\ A_0) \to D^1(R|R',\ k) \underset{k}{\otimes} A_0$  , which means that the deformation  $(R,\ \mathrm{Spec}(A))$  can be lifted to a deformation $(R',\ \mathrm{Spec}(A'))$  .

(b)  By 3.14 we have the isomorphisms

\SGAvar{\begin{equation*}
D^1(A_0|k,\ A_0) \simeq D^1((A_0)_x|k,\ (A_0)_x) \simeq D^1(\widehat{(A_0)}_x|k,\ \widehat{(A_0)}_x)
\end{equation*}
i.e.,   $[\underline{M}_{X_0}(k[\epsilon])] \simeq [\underline{M}_{\mathrm{Spec}(\underline{O}_{X_0,\ x})}(k[\epsilon])] \simeq [\underline{M}_{\mathrm{Spec}(\widehat{\underline{O}}_{X_0,\ x})}(k[\epsilon])]$ .

}{\begin{equation*}
D^1(A_0|k,\ A_0) \simeq D^1((A_0)_x|k,\ (A_0)_x) \simeq D^1(\widehat{(A_0)_x}|k,\ \widehat{(A_0)_x})
\end{equation*}
i.e.,   $[\underline{M}_{X_0}(k[\epsilon])] \simeq [\underline{M}_{\mathrm{Spec}(\underline{O}_{X_0,\ x})}(k[\epsilon])] \simeq [\underline{M}_{\mathrm{Spec}(\widehat{\underline{O}}_{X_0,\ x})}(k[\epsilon])]$ .

}{F7I-VI-096}It suffices to show, by 2.7(a), that

\SGAvar{\begin{equation*}
[\underline{M}_{X_0}] \to [\underline{M}_{\mathrm{Spec}(\underline{O}_{X_0,\ x})}] \to [\underline{M}_{\mathrm{Spec}(\widehat{\underline{O}}_{X_0,\ x})}]
\end{equation*}
are both smooth functors. Let  $(R,\ \mathrm{Spec}(A))$  be a deformation of  $X_0 = \mathrm{Spec}(A_0)$  , and let  $0 \to k \to R' \to R \to 0$  be a small extension in  $\underline{C}_{\Lambda}$  . As above, the ring homomorphisms  $R' \to R \to A$  gives us the exact commutative diagram
\SGApage{I}{107}
{\small
\[
\begin{tikzcd}[column sep=small, row sep=large]
{} \arrow[r] \& D^1(A_0|k,A_0) \arrow[r] \arrow[d, "\approx"'] \& D^1(A|R',A_0) \arrow[r, "p"] \arrow[d] \& D^1(R|R',k)\underset{k}{\otimes}A_0 \arrow[r] \arrow[d] \& D^2(A_0|k,A_0) \arrow[r] \arrow[d, "\approx"] \& \dots \\
{} \arrow[r] \& D^1(A_0|k,(A_0)_x) \arrow[r] \& D^1(A|R',(A_0)_x) \arrow[r] \& D^1(R|R',k)\underset{k}{\otimes}(A_0)_x \arrow[r] \& D^2(A_0|k,(A_0)_x) \arrow[r] \& \dots \\
{} \arrow[r] \& D^1((A_0)_x|k,(A_0)_x) \arrow[r] \arrow[u] \& D^1(A_x|R',(A_0)_x) \arrow[r, "p_x"] \arrow[u] \& D^1(R|R',k)\underset{k}{\otimes}(A_0)_x \arrow[r] \arrow[u] \& D^2((A_0)_x|k,(A_0)_x) \arrow[r] \arrow[u] \& \dots
\end{tikzcd}
\]}























where the indicated isomorphism comes from the fact that  $A_0$  is smooth over  $k$  except at one closed point  $x$  so that

$$D^i(A_0|k,\ (A_0)_x) = D^i(A_0|k,\ A_0)_x = D^i(A_0|k,\ A_0)$$
for all  $i > 0$ . Now suppose that we are given an arrow $(R',\ \mathrm{Spec}(B)) \to (R,\ \mathrm{Spec}(A_x))$  in  $[\underline{M}_{\mathrm{Spec}((A_0)_x)}]$ . It means that we are given a class $[B] \in D^1(A_x|R',\ (A_0)_x)$  such that  $p_x([B]) = [R']$ . Then a trivial diagram chasing shows that there exists $[A'] \in D^1(A|R',\ A_0)$  such that $p([A']) = [R']$  and  $A'_x = B$  i.e.,  $(R',\ \mathrm{Spec}(A'))$  is a deformation of  $\mathrm{Spec}(A_0)$ to  $R'$  such that  $\mathrm{Spec}(A')_x = \mathrm{Spec}(B)$ . This proves that $\underline{M}_{X_0} \to \underline{M}_{\mathrm{Spec}(\underline{O}_{X_0,\ x})}$ is smooth. The smoothness of the functor  $\underline{M}_{\mathrm{Spec}((A_0)_x)} \to \underline{M}_{\mathrm{Spec}(\widehat{(A_0)_x})}$ comes from the similar diagram together with the fact that

$$D^i((A_0)_x|k,\ \widehat{(A_0)_x}) = \widehat{D^i}((A_0)_x|k,\ (A_0)_x) = D^i((A_0)_x|k,\ (A_0)_x)$$
for all  $i > 0$  since  $x$  is an only possible non-smooth point.

}{\begin{equation*}
[\underline{M}_{X_0}] \to [\underline{M}_{\mathrm{Spec}(\underline{O}_{X_0,\ x})}] \to [\underline{M}_{\mathrm{Spec}(\widehat{\underline{O}}_{X_0,\ x})}]
\end{equation*}
are both smooth functors. Let  $(R,\ \mathrm{Spec}(A))$  be a deformation of  $X_0 = \mathrm{Spec}(A_0)$  , and let  $0 \to k \to R' \to R \to 0$  be a small extension in  $\underline{C}_{\Lambda}$  . As above, the ring homomorphisms  $R' \to R \to A$  give us the exact commutative diagram
\SGApage{I}{107}
{\small
\[
\begin{tikzcd}[column sep=small, row sep=large]
{} \arrow[r] \& D^1(A_0|k,A_0) \arrow[r] \arrow[d, "\approx"'] \& D^1(A|R',A_0) \arrow[r, "p"] \arrow[d] \& D^1(R|R',k)\underset{k}{\otimes}A_0 \arrow[r] \arrow[d] \& D^2(A_0|k,A_0) \arrow[r] \arrow[d, "\approx"] \& \dots \\
{} \arrow[r] \& D^1(A_0|k,(A_0)_x) \arrow[r] \& D^1(A|R',(A_0)_x) \arrow[r] \& D^1(R|R',k)\underset{k}{\otimes}(A_0)_x \arrow[r] \& D^2(A_0|k,(A_0)_x) \arrow[r] \& \dots \\
{} \arrow[r] \& D^1((A_0)_x|k,(A_0)_x) \arrow[r] \arrow[u] \& D^1(A_x|R',(A_0)_x) \arrow[r, "p_x"] \arrow[u] \& D^1(R|R',k)\underset{k}{\otimes}(A_0)_x \arrow[r] \arrow[u] \& D^2((A_0)_x|k,(A_0)_x) \arrow[r] \arrow[u] \& \dots
\end{tikzcd}
\]}























where the indicated isomorphisms come from the fact that  $A_0$  is smooth over  $k$  except at one closed point  $x$  so that

$$D^i(A_0|k,\ (A_0)_x) = D^i(A_0|k,\ A_0)_x = D^i(A_0|k,\ A_0)$$
for all  $i > 0$ . Now suppose that we are given an arrow $(R',\ \mathrm{Spec}(B)) \to (R,\ \mathrm{Spec}(A_x))$  in  $[\underline{M}_{\mathrm{Spec}((A_0)_x)}]$ . It means that we are given a class $[B] \in D^1(A_x|R',\ (A_0)_x)$  such that  $p_x([B]) = [R']$ . Then a diagram chase shows that there exists $[A'] \in D^1(A|R',\ A_0)$  such that $p([A']) = [R']$  and  $A'_x = B$  i.e.,  $(R',\ \mathrm{Spec}(A'))$  is a deformation of  $\mathrm{Spec}(A_0)$ to  $R'$  such that  $\mathrm{Spec}(A')_x = \mathrm{Spec}(B)$ . This proves that $\underline{M}_{X_0} \to \underline{M}_{\mathrm{Spec}(\underline{O}_{X_0,\ x})}$ is smooth. The smoothness of the functor  $\underline{M}_{\mathrm{Spec}((A_0)_x)} \to \underline{M}_{\mathrm{Spec}(\widehat{(A_0)_x})}$ comes from the similar diagram together with the fact that

$$D^i((A_0)_x|k,\ \widehat{(A_0)_x}) = \widehat{D^i((A_0)_x|k,\ (A_0)_x)} = D^i((A_0)_x|k,\ (A_0)_x)$$
for all  $i > 0$  since  $x$  is the only possible non-smooth point.

}{F7I-VI-097}\underline{Theorem} 4.11.  \underline{Let}  $X_0$  \underline{be a} (\underline{separated}) $k$-\underline{scheme of finite type, smooth outside a finite closed set}  $X_0^{\mathrm{Sing}} = \{x_1,\dots,x_r\}$  \underline{of}  $X_0$ . \underline{Let}  $U_0^1,\dots,U_0^r$  \underline{be affine open neighbourhoods of}  $x_1,\dots,x_r$  \underline{such that}  $U_0^i \cap X_0^{\mathrm{Sing}} = \{x_i\}$ . \underline{If} $H^2(X_0,\ \underline{D}^0_{X_0}) = 0$ , \underline{then the natural functor}  $\underline{M}_{X_0} \to \underline{M}_{U_0^1} \times\dots\times \underline{M}_{U_0^r}$   \underline{is smooth.}

\underline{Proof}: We can complete  $U_0^1, \dots ,U_0^r$  into an affine open covering of  $X_0$  by choosing $U_0^{r+1},\dots,U_0^N$  such that   $U_0^j \cap X_0^{\mathrm{Sing}} = \emptyset$  for all  $j > r$ . Now let $R' \to R$  be a small extension in  $\underline{C}_{\Lambda}$ , and let  $(R',\ U'^1),\dots,(R',\ U'^r)$  be deformations of  $U_0^1,\dots,U_0^r$  to  $R'$  such that  $U'^i \underset{R'}{\otimes} R = X|U_0^i$ ,  $1 \le i \le r$ , for
\SGApage{I}{108}
some deformation  $(R,\ X)$  of  $X_0$  to  $R$ . We must show that there exists a deformation  $(R',\ X')$  of  $X_0$  to  $R'$  such that  $X' \underset{R'}{\otimes} R = X$  and  $X'|U_0^i = U'^i$  for $1 \le i \le r$ . Since  $U_0^j$  for  $j > r$  is smooth over  $k$ ,  $X|U_0^j$  can be lifted, by 4.9(a), to a deformation  $(R',\ U'^j)$  i.e., there exist deformations  $(R',\ U'^j)$ of  $U_0^j$  to  $R'$  such that  $U'^j \underset{R'}{\otimes} R = X|U_0^j$  for all  $j > r$ . Since, for any pair $i \ne j$ ,  $U_0^i \cap U_0^j$  is affine and smooth over  $k$ , we obtain an isomorphism

\SGAvar{$$\tau_{ij} : U'^i|U_0^i \cap U_0^j \xrightarrow{\ \sim\ } U'^j|U_0^i \cap U_0^j$$
such that  $\tau_{ij} \underset{R'}{\otimes} R = \mathrm{id}_{X|U_0^i \cap U_0^j}$ , and in turn defines a cohomology class in $H^2(X_0\ \underline{D}^0_{X_0})$  given by  $U_0^i \cap U_0^j \cap U_0^k \to \tau_{ij}\tau_{jk}\tau_{ik}^{\ -1}$ . Since  $H^2(X_0,\ \underline{D}^0_{X_0}) = 0$  by hypothesis, the obstruction for glueing these to obtain a global deformation vanishes, and hence we get a deformation  $(R',\ X')$  of  $X_0$  to  $R'$  such that $X' \underset{R'}{\otimes} R = X$  and  $X'|U_0^j \simeq U'^j$ . This completes our proof.

\underline{Remark} 4.12.  The previous proof assumes  $X_0$  to be separated, which is unnecessary. In fact, for a fixed deformation  $(R,\ X)$  of  $X_0$  and deformations  $(R',\ U'^j)$  of $U_0^i$ , the deformations  $(R',\ V')$  of open set  $V$  of  $X$  to  $R'$  compatible with the deformation  $(R',\ U'^j)$  already given, form a  "gerbe"  (in Giraud's terminology) whose structure sheaf is  $\underline{D}^0_{X_0}$ . Hence the obstruction to the existence of a section of this  "gerbe"  (i.e., a global deformation  $(R',\ X')$  of  $X$  compatible with the given datum  $(R',\ U'^j))$  is in $H^2(X_0,\ \underline{D}^0_{X_0})$  which is zero by assumption.

}{$$\tau_{ij} : U'^i|U_0^i \cap U_0^j \xrightarrow{\ \sim\ } U'^j|U_0^i \cap U_0^j$$
such that  $\tau_{ij} \underset{R'}{\otimes} R = \mathrm{id}_{X|U_0^i \cap U_0^j}$ , and in turn defines a cohomology class in $H^2(X_0,\ \underline{D}^0_{X_0})$  given by  $U_0^i \cap U_0^j \cap U_0^k \to \tau_{ij}\tau_{jk}\tau_{ik}^{\ -1}$ . Since  $H^2(X_0,\ \underline{D}^0_{X_0}) = 0$  by hypothesis, the obstruction for glueing these to obtain a global deformation vanishes, and hence we get a deformation  $(R',\ X')$  of  $X_0$  to  $R'$  such that $X' \underset{R'}{\otimes} R = X$  and  $X'|U_0^j \simeq U'^j$ . This completes our proof.

\underline{Remark} 4.12.  The previous proof assumes  $X_0$  to be separated, which is unnecessary. In fact, for a fixed deformation  $(R,\ X)$  of  $X_0$  and deformations  $(R',\ U'^j)$  of $U_0^j$ , the deformations  $(R',\ V')$  of open set  $V$  of  $X$  to  $R'$  compatible with the deformation  $(R',\ U'^j)$  already given, form a  "gerbe"  (in Giraud's terminology) whose structure sheaf is  $\underline{D}^0_{X_0}$ . Hence the obstruction to the existence of a section of this  "gerbe"  (i.e., a global deformation  $(R',\ X')$  of  $X$  compatible with the given datum  $(R',\ U'^j))$  is in $H^2(X_0,\ \underline{D}^0_{X_0})$  which is zero by assumption.

}{F7I-VI-098}\underline{Corollary} 4.13.  \underline{Let}  $X_0$  \underline{be a} $k$-\underline{scheme of finite type, smooth outside a finite subset}  $X_0^{\mathrm{Sing}} = \{x_1,\dots,x_r\}$  of  $X_0$  \underline{such that}  $H^2(X_0,\ \underline{D}^0_{X_0}) = 0$ . \underline{Then}

\SGAvar{$$\text{(a)}\quad \underline{M}_{X_0} \to \underline{M}_{\mathrm{Spec}(\widehat{\underline{O}}_{X_0,x_1})} \times\dots\times \underline{M}_{\mathrm{Spec}(\widehat{\underline{O}}_{X_0,x_r})}\quad \underline{\text{is smooth.}}$$
$$\underline{\text{In particular,}}\ \ \dim \underline{M}_{X_0} = \dim H^1(X_0,\ \underline{D}^0_{X_0}) + \sum_{c=1}^{r} \ \dim \underline{M}_{\mathrm{Spec}(\widehat{\underline{O}}_{X_0,x_0})}\ .$$
(b)  \underline{If furthermore}  $X_0$  \underline{is locally a complete-intersection, then}  $\underline{M}_{X_0}$ \underline{is smooth.}

}{$$\text{(a)}\quad \underline{M}_{X_0} \to \underline{M}_{\mathrm{Spec}(\widehat{\underline{O}}_{X_0,x_1})} \times\dots\times \underline{M}_{\mathrm{Spec}(\widehat{\underline{O}}_{X_0,x_r})}\quad \underline{\text{is smooth.}}$$
$$\underline{\text{In particular,}}\ \ \dim \underline{M}_{X_0} = \dim H^1(X_0,\ \underline{D}^0_{X_0}) + \sum_{c=1}^{r} \ \dim \underline{M}_{\mathrm{Spec}(\widehat{\underline{O}}_{X_0,x_c})}\ .$$
(b)  \underline{If furthermore}  $X_0$  \underline{is locally a complete-intersection, then}  $\underline{M}_{X_0}$ \underline{is smooth.}

}{F7I-VI-099}\SGApage{I}{109}

\underline{Proof}:

(a)  The first statement follows from 4.10 and 4.9(b). Now let  $(R, \mathfrak{X})$  and  $(R_1, \mathfrak{X}_1),...,(R_r, \mathfrak{X}_r)$  be formal versal deformations of  $X_0$  and  $\mathrm{Spec}(\underline{O}_{X_0,x_1}),...,\mathrm{Spec}(\underline{O}_{X_0,x_r})$  respectively. We then obtain a commutative diagram

\begin{center}
\begin{tikzcd}[column sep=huge, row sep=huge]
\underline{h}_R \arrow[r] \arrow[d] \& \underline{h}_{R_1 \widehat{\otimes}_\Lambda ... \widehat{\otimes}_\Lambda R_r} \arrow[d] \\
\underline{M}_{X_0} \arrow[r] \& \underline{M}_{\mathrm{Spec}(\widehat{\underline{O}}_{X_0,x_1})} \times ... \times \underline{M}_{\mathrm{Spec}(\widehat{\underline{O}}_{X_0,x_r})}
\end{tikzcd}
\end{center}















Since the composite functor

\begin{equation*}
\underline{h}_R \to \underline{M}_{X_0} \to \underline{M}_{\mathrm{Spec}(\widehat{\underline{O}}_{X_0,x_1})} \times ... \times \underline{M}_{\mathrm{Spec}(\widehat{\underline{O}}_{X_0,x_r})}
\end{equation*}
is smooth whereas

\begin{equation*}
\underline{h}_{R_1 \widehat{\otimes}_\Lambda ... \widehat{\otimes}_\Lambda R_r} \to \underline{M}_{\mathrm{Spec}(\widehat{\underline{O}}_{X_0,x_1})} \times .... \times \underline{M}_{\mathrm{Spec}(\widehat{\underline{O}}_{X_0,x_r})}
\end{equation*}
is minimally smooth, it follows from 2.7(a) that

$\underline{h}_R \to \underline{h}_{R_1 \widehat{\otimes}_\Lambda ... \widehat{\otimes}_\Lambda R_r}$   is smooth i.e.,  $R$  is a formal power-series ring over  $R_1 \widehat{\otimes}_\Lambda ... \widehat{\otimes}_\Lambda R_r$ . Since

\SGAvar{\begin{equation*}
0 \to H^1(X_0,\underline{D}^0_{X_0}) \to [\underline{M}_{X_0}(k[\epsilon])] \to [\underline{M}_{\mathrm{Spec}(\widehat{\underline{O}}_{X_0,x_1})}(k[\epsilon])] ... [\underline{M}_{\mathrm{Spec}(\widehat{\underline{O}}_{X_0,x_r})}(k[\epsilon])]
\end{equation*}
is exact by 4.4 i.e.,

}{\begin{equation*}
0 \to H^1(X_0,\underline{D}^0_{X_0}) \to [\underline{M}_{X_0}(k[\epsilon])] \to [\underline{M}_{\mathrm{Spec}(\widehat{\underline{O}}_{X_0,x_1})}(k[\epsilon])] \times ... \times [\underline{M}_{\mathrm{Spec}(\widehat{\underline{O}}_{X_0,x_r})}(k[\epsilon])]
\end{equation*}
is exact by 4.4 i.e.,

}{F7I-VI-100}\begin{equation*}
0 \to H^1(X_0, \underline{D}^0_{X_0}) \to \mathrm{Hom}_{\Lambda-\mathrm{alg}}(R, k[\epsilon]) \to \mathrm{Hom}_{\Lambda-\mathrm{alg}}(R_1 \widehat{\otimes}_\Lambda ... \widehat{\otimes}_\Lambda R_r, k[\epsilon])
\end{equation*}
is exact, it follows that  $R$  is a formal power-series ring over    $R_1 \widehat{\otimes}_\Lambda ... \widehat{\otimes}_\Lambda R_r$ in $d$-variables where  $d = \dim_k H^1(X_0, \underline{D}^0_{X_0})$  .

Therefore

\begin{equation*}
\begin{aligned}
\dim \underline{M}_{X_0} &= \dim k \otimes_\Lambda R = d + \dim k \otimes_\Lambda (R_1 \widehat{\otimes}_\Lambda ... \widehat{\otimes}_\Lambda R_r) \\
&= d + \dim(k \otimes_\Lambda R_1) \widehat{\otimes}_k ... \widehat{\otimes}_k (k \otimes_\Lambda R_r) \\
&= d + \sum_{i=1}^{r} \dim k \otimes_\Lambda R_i = d + \sum_{i=1}^{r} \dim \underline{M}_{\mathrm{Spec}(\widehat{\underline{O}}_{X_0,x_i})}
\end{aligned}
\end{equation*}

\SGApage{I}{110}

(b)  follows from 4.10 and 4.9(a).

\SGAvar{4.14.    We close this paragraph with an explicit construction of a formal versal deformation in a simple situation. In fact, if  $X_0$  is affine and is a complete-intersection, it is trivial to write down explicitely. Thus let

\begin{equation*}
A_0 = k[x_1,...,x_n]/(f_1,...,f_r) \qquad S/J
\end{equation*}
where  $S = k[x_1,...,x_n]$  ,  $J = (f_1,...,f_r)$  , and  $f_1,...,f_r$  is an $S$-regular sequence. We then obtain the exact sequence

\begin{equation*}
... \to \mathrm{Der}_k(S, A_0) \to \mathrm{Hom}_{A_0}(J/J^2, A_0) \to D^1(A_0|k, A_0) \to 0
\end{equation*}
We note that  $J/J^2$  is  $A_0$-free module of rank  $r$  generated by $f_1(\mathrm{mod}\ J^2),...,f_r(\mathrm{mod}\ J^2)$  , since  $f_1,...,f_r$  is an $S$-regular sequence. We now assume that  $D^1(A_0|k, A_0)$  is a finite-dimensional vector space over  $k$  , say  $A$ has isolated non-smooth points only. Let  $d = \dim_k D^1(A_0|k, A_0)$  and choose $p_j : J/J^2 \to A_0$   $(j = 1, 2,...,d)$  which form a $k$-basis of  $D^1(A_0|k, A_0)$  . Now choose  $P_{jk} \in \Lambda[x_1,...,x_n]$  such that  $p_j(f_i) = P_{ji}$  (mod $\underline{m}_\Lambda + J$)  where  $\underline{m}_\Lambda$  is the maximal ideal of  $\Lambda$  . We then consider

\begin{equation*}
\begin{aligned}
R &= \Lambda[[t_1,...,t_d]] \\
B &= R[x_1,...,x_n]/(G_1,...,G_r)
\end{aligned}
\end{equation*}
where  $G_j = F_j - \sum_{i=1}^{d} t_i P_{ji}$  and  $F_j$  is a polynomial with coefficients in  $\Lambda$ gotten from  $f_j$  by lifting the coefficients from  $k$  to  $\Lambda$  . Since $R[x_1,...,x_n]$  is $R$-flat and  $G_1,...,G_r$  modulo  $\underline{m}_R$  (which becomes  $f_1,...,f_r$) is a regular sequence in  $k \otimes_R R[x_1,...,x_n]$  , it follows that  $B$  is $R$-flat, and $k \otimes_R B = A_0$  . Furthermore,

}{4.14.    We close this paragraph with an explicit construction of a formal versal deformation in a simple situation. In fact, if  $X_0$  is affine and is a complete-intersection, it is trivial to write down explicitly. Thus let

\begin{equation*}
A_0 = k[x_1,...,x_n]/(f_1,...,f_r) = S/J
\end{equation*}
where  $S = k[x_1,...,x_n]$  ,  $J = (f_1,...,f_r)$  , and  $f_1,...,f_r$  is an $S$-regular sequence. We then obtain the exact sequence

\begin{equation*}
... \to \mathrm{Der}_k(S, A_0) \to \mathrm{Hom}_{A_0}(J/J^2, A_0) \to D^1(A_0|k, A_0) \to 0
\end{equation*}
We note that  $J/J^2$  is  $A_0$-free module of rank  $r$  generated by $f_1(\mathrm{mod}\ J^2),...,f_r(\mathrm{mod}\ J^2)$  , since  $f_1,...,f_r$  is an $S$-regular sequence. We now assume that  $D^1(A_0|k, A_0)$  is a finite-dimensional vector space over  $k$  , say  $A_0$ has isolated non-smooth points only. Let  $d = \dim_k D^1(A_0|k, A_0)$  and choose $p_j : J/J^2 \to A_0$   $(j = 1, 2,...,d)$  which form a $k$-basis of  $D^1(A_0|k, A_0)$  . Now choose  $P_{jk} \in \Lambda[x_1,...,x_n]$  such that  $p_j(f_i) = P_{ji}$  (mod $\underline{m}_\Lambda + J$)  where  $\underline{m}_\Lambda$  is the maximal ideal of  $\Lambda$  . We then consider

\begin{equation*}
\begin{aligned}
R &= \Lambda[[t_1,...,t_d]] \\
B &= R[x_1,...,x_n]/(G_1,...,G_r)
\end{aligned}
\end{equation*}
where  $G_j = F_j - \sum_{i=1}^{d} t_i P_{ij}$  and  $F_j$  is a polynomial with coefficients in  $\Lambda$ gotten from  $f_j$  by lifting the coefficients from  $k$  to  $\Lambda$  . Since $R[x_1,...,x_n]$  is $R$-flat and  $G_1,...,G_r$  modulo  $\underline{m}_R$  (which becomes  $f_1,...,f_r$) is a regular sequence in  $k \otimes_R R[x_1,...,x_n]$  , it follows that  $B$  is $R$-flat, and $k \otimes_R B = A_0$  . Furthermore,

}{F7I-VI-101}\begin{equation*}
\mathrm{Hom}_{\Lambda-\mathrm{alg}}(R, k[\epsilon]) \to D^1(A_0|k, A_0) = [\underline{M}_{\mathrm{Spec}(A_0)}(k[\epsilon])]
\end{equation*}
is bijective by our construction of  $B$  . Therefore, if we set 


\SGApage{I}{111}

\begin{equation*}
 \mathrm{Spf}(B) \;=\; \varinjlim_{v}\; \mathrm{Spec}(R/\underline{m}^{v+1}\ \underset{R}{\otimes}\ B)
\end{equation*}
then  $(R,\ \mathrm{Spf}(B))$  is an object in  $\widehat{\underline{M}}_{\mathrm{Spec}(A_0)}$  and the functor  $\underline{h}_R \to \underline{M}_{\mathrm{Spec}(A_0)}$  induced by  $(R,\ \mathrm{Spf}(B))$  is smooth by 1.6, since  $R$  is a formal power-series ring over  $\Lambda$  . In other words,  $(R,\ \mathrm{Spf}(B))$  is a formal versal object of  $\underline{M}_{\mathrm{Spec}(A_0)}$  .

We note, by 4.10(b), that if we set

\begin{equation*}
B = R[[X_1,\ldots,X_n]]/(G_1,\ldots,G_r)\ \ ,
\end{equation*}
then  $(R,\ \mathrm{Spf}(B))$  is a formal versal deformation of  $\mathrm{Spec}(A_0)$  where  $A_0 = k[[X_1,\ldots,X_n]]/(f_1,\ldots,f_r)$  .

\underline{Remark} 4.15.  Let  $X_0$  be a $k$-scheme of finite type with isolated non-smooth points only, and let  $(R,\ \mathfrak{X})$  be a formal versal deformation of  $X_0$  . The above construction shows that if  $X_0$  is affine and locally a complete-intersection, then  $\mathfrak{X}$  is a formalization of an $R$-scheme of finite type. This fact can easily be generalized to the case when  $X_0$  is not necessarily affine but  $H^2(X_0,\ \underline{D}^0_{X_0}) = 0$  .

\SGAvar{\underline{Remark} 4.16.  The \underline{global}  theory of the relative cotangent complex (the latter theory being reduced to the truncated complex of length 1 , which is sufficient however for many purposes of deformation theory) allows to generalize and simplify certain results, such as 4.10. Namely let  $S$  be a scheme,  $S_0 \subset S' \subset S$  two subschemes of  $S$  defined by ideals  $J \supset I$  such that  $JI = 0$  ,  $X'$  a scheme flat over  $S'$  ,  $L. = L^{X_0/S_0}$  the relative cotangent complex of  $X_0$  over  $S_0$  , which is an element of the derived category  $D(X_0,\ \underline{O}_X)$  . We consider the category  $\underline{F}$  of all flat schemes  $X$  over  $S$  which extend  $X'/S'$  , i.e. together with an  $S'$-isomorphism  $X \underset{S}{\times} S' \simeq X'$  . Then

(a)  For an object  $X'$  of  $\underline{F}$  , the group of automorphisms is canonically isomorphic to  $\mathrm{Ext}^0(X_0;\ L.,I_{X_0})$  (where  $I_{X_0}$  is the inverse image of  $I$  , viewed as a module on  $S_0$  , on  $X_0)$  .


}{\underline{Remark} 4.16.  The \underline{global}  theory of the relative cotangent complex (the latter theory being reduced to the truncated complex of length 1 , which is sufficient however for many purposes of deformation theory) allows to generalize and simplify certain results, such as 4.10. Namely let  $S$  be a scheme,  $S_0 \subset S' \subset S$  two subschemes of  $S$  defined by ideals  $J \supset I$  such that  $JI = 0$  ,  $X'$  a scheme flat over  $S'$  ,  $L. = L^{X_0/S_0}$  the relative cotangent complex of  $X_0$  over  $S_0$  , which is an element of the derived category  $D(X_0,\ \underline{O}_{X_0})$  . We consider the category  $\underline{F}$  of all flat schemes  $X$  over  $S$  which extend  $X'/S'$  , i.e. together with an  $S'$-isomorphism  $X \underset{S}{\times} S' \simeq X'$  . Then

(a)  For an object  $X$  of  $\underline{F}$  , the group of automorphisms is canonically isomorphic to  $\mathrm{Ext}^0(X_0;\ L.,I_{X_0})$  (where  $I_{X_0}$  is the inverse image of  $I$  , viewed as a module on  $S_0$  , on  $X_0)$  .


}{F7I-VI-102}\SGApage{I}{112}

(b)  The set of isomorphism classes of objects of  $\underline{F}$  is empty or a torsor (= a principal homogeneous space) under the group  $\mathrm{Ext}^1(X_0;\ L.,I_{X_0})$  .

(c)  One can define a canonical element  $c \in \mathrm{Ext}^2(X_0;\ L.,I_{X_0})$  whose vanishing is necessary and sufficient for  $\underline{F}$  to be non empty.

In order to study the  $\mathrm{Ext}^i(X_0;\ L.,I_{X_0})$  , it is often convenient to write

\SGAvar{\begin{equation*}
\mathrm{Ext}^i(X_0;L.,I_{X_0}) = H^i(X_0,\ \underline{RHom}(L.,I_{X_0}))\ \ .
\end{equation*}
and to notice that the complex  $\underline{RHom}(L.,I_{X_0})$  has as cohomology sheaves the sheaves  $\underline{D}^i_{X_0}(I_{X_0})$  introduced in §3. For  $i = 0, 1, 2$,  these sheaves may be viewed respectively as the sheaf of automorphisms of any object  $X$  of  $\underline{F}$  (understood:  automorphisms inducing the identity on  $X')$  , the sheaf of indeterminacy of isomorphism classes of local solutions to the flat deformation problem, and the sheaf of obstructions to the existence of local solutions. On the other hand, the standart spectral sequence gives us

\begin{equation*}
\mathrm{Ext}^*(X_0 \quad L.,I_{X_0}) \leftarrow E_2^{p,q} = H^p(X_0,\ \underline{D}^q_{X_0}(I_{X_0}))\ \ ,
\end{equation*}
which yields the well-known isomorphism

}{\begin{equation*}
\mathrm{Ext}^i(X_0;L.,I_{X_0}) = H^i(X_0,\ \underline{RHom}(L.,I_{X_0}))\ \ .
\end{equation*}
and to notice that the complex  $\underline{RHom}(L.,I_{X_0})$  has as cohomology sheaves the sheaves  $\underline{D}^i_{X_0}(I_{X_0})$  introduced in §3. For  $i = 0, 1, 2$,  these sheaves may be viewed respectively as the sheaf of automorphisms of any object  $X$  of  $\underline{F}$  (understood:  automorphisms inducing the identity on  $X')$  , the sheaf of indeterminacy of isomorphism classes of local solutions to the flat deformation problem, and the sheaf of obstructions to the existence of local solutions. On the other hand, the standard spectral sequence gives us

\begin{equation*}
\mathrm{Ext}^*(X_0;\ L.,I_{X_0}) \leftarrow E_2^{p,q} = H^p(X_0,\ \underline{D}^q_{X_0}(I_{X_0}))\ \ ,
\end{equation*}
which yields the well-known isomorphism

}{F7I-VI-103}\begin{equation*}
\mathrm{Aut}(X/X') \simeq H^0(X_0,\underline{D}^0_{X_0}(I_{X_0})) \simeq H^0(X_0,\underline{Hom}(\underline{\Omega}^1_{X_0/S_0},I_{X_0})) = \mathrm{Hom}(\underline{\Omega}^1_{X_0/S_0},I_{X_0})\ \ ,
\end{equation*}
and the exact sequence in low degrees

\SGAvar{\begin{equation*}
0 - H^1(X_0,\underline{D}^0_{X_0}) - \mathrm{Ext}^1 - H^0(X_0,\underline{D}^1_{X_0}) - H^2(X_0,\underline{D}^0_{X_0}) - \mathrm{Ext}^2
\end{equation*}
This shows in particular that if  $H^2(X_0,\underline{D}^0_{X_0}) = 0$  , and if there exists any local flat deformation  $X$  of  $X'/S'$  , then there exists a flat global deformation which corresponds to given local deformations. More precisely and generally, for a given local deformations (i.e. a section of the  "sheaf of isomorphism classes of local flat deformations" ), the argument 4.10 shows that the obstruction to finding a global flat deformation compatible with these is in  $H^2(X_0,\ \underline{D}^0_{X_0})$  , hence such a global flat deformation exists if  $H^2(X_0,\ \underline{D}^0_{X_0}) = 0$  : this is essentially 4.10.


}{\begin{equation*}
0 - H^1(X_0,\underline{D}^0_{X_0}) - \mathrm{Ext}^1 - H^0(X_0,\underline{D}^1_{X_0}) - H^2(X_0,\underline{D}^0_{X_0}) - \mathrm{Ext}^2
\end{equation*}
This shows in particular that if  $H^2(X_0,\underline{D}^0_{X_0}) = 0$  , and if there exists any local flat deformation  $X$  of  $X'/S'$  , then there exists a flat global deformation which corresponds to given local deformations. More precisely and generally, for given local deformations (i.e. a section of the  "sheaf of isomorphism classes of local flat deformations" ), the argument 4.10 shows that the obstruction to finding a global flat deformation compatible with these is in  $H^2(X_0,\ \underline{D}^0_{X_0})$  , hence such a global flat deformation exists if  $H^2(X_0,\ \underline{D}^0_{X_0}) = 0$  : this is essentially 4.10.


}{F7I-VI-104}\SGApage{I}{113}

On the other hand, the spectral sequence shows that the  $\mathrm{Ext}^2$  is zero whenever we have the relations

\SGAvar{\begin{equation}\tag{*}
H^0(X_0,\ \underline{D}^2_{X_0}\ ) = 0 \quad ,\quad H^1(X_0,\ \underline{D}^1_{X_0}\ ) = 0 \quad ,\quad H^2(X_0,\ \underline{D}^0_{X_0}\ ) = 0 \quad .
\end{equation}
The first relation is always satisfied if  $X_0$  is a relative complete intersection over  $S_0$  (as then  $\underline{D}^2_{X_0} = 0$  for any module  $M$  on  $X_0$ ), the second is satisfied provided the support of  $\underline{D}^1_{X_0}$  is discrete, as is the case if  $X_0$  is smooth over  $S_0$  except on a discrete subset of  $X_0$  over  $S_0$ . When the previous three relations $(*)$ are all satisfied, then by (c) above,  $\underline{F}$  is non empty, which generalizes 4.12(b).

More geometrically, and without using the global theory of the relative cotangentcomplex (but only the globalizing of the functors  $\underline{D}^i_{X_0}$  ), one can get a simple geometric interpretation of the three successive obstructions to deforming flatly  $X'$  to  $X$ , lying in the three left hand terms of $(*)$. Namely by the local theory, the first obstruction  $c_1$  in  $H^0(X_0,\ \underline{D}^2_{X_0})$  can be defined as the obstruction to finding local solutions, when this obstruction vanishes, the sheaf  $\underline{F}$  of isomorphism classes of local solutions of the problem is (by the local theory) in a canonical way a torsor under  $\underline{D}^1_{X_0}$ , and the second obstruction  $c_2$ , lying in  $H^1(X_0,\ \underline{D}^1_{X_0})$ , is the obstruction to finding a section of this torsor, i.e. of finding a  "coherent system of isomorphism classes of local solutions" ; if this second obstruction also vanishes, then, choosing arbitrarily a section of  $\underline{F}$ , one gets, as noted above, an obstruction  $c_3$  in  $H^2(X_0,\ \underline{D}^0_{X_0})$  to finding a global solution corresponding to the given system of (isomorphism classes of) local solutions. Using the fact that the section of  $F$  is determined up to adding a section of  $\underline{D}^1_{X_0}$ , we get a map  $d : H^0(X_0,\ \underline{D}^1_{X_0}) \to H^2(X_0,\ \underline{D}^0_{X_0})$ , and a solution to the deformation problem exists if and only if  $c_1 = 0$ ,  $c_2 = 0$ , and  $c_3$  is zero as an element of Coker $d$  (each  $c_i$  being defined when the previous ones are zero). Of course, the groups  $H^0(X_0,\ \underline{D}^2_{X_0})$ ,  $H^1(X_0,\ \underline{D}^1_{X_0})$  and
\SGApage{I}{114}
$H^2(X_0,\ \underline{D}^0_{X_0})/\mathrm{Im}\ d$  are just the initial terms  $E_2^{0,2}$ ,  $E_2^{1,1}$ , and  $E_3^{2,0}$  of the spectral sequence above,  $d$  being the differential (up to sign ?)

}{\begin{equation}\tag{*}
H^0(X_0,\ \underline{D}^2_{X_0}\ ) = 0 \quad ,\quad H^1(X_0,\ \underline{D}^1_{X_0}\ ) = 0 \quad ,\quad H^2(X_0,\ \underline{D}^0_{X_0}\ ) = 0 \quad .
\end{equation}
The first relation is always satisfied if  $X_0$  is a relative complete intersection over  $S_0$  (as then  $\underline{D}^2_{X_0} = 0$  for any module  $M$  on  $X_0$ ), the second is satisfied provided the support of  $\underline{D}^1_{X_0}$  is discrete, as is the case if  $X_0$  is smooth over  $S_0$  except on a discrete subset of  $X_0$  over  $S_0$ . When the previous three relations $(*)$ are all satisfied, then by (c) above,  $\underline{F}$  is non empty, which generalizes 4.13(b).

More geometrically, and without using the global theory of the relative cotangent complex (but only the globalizing of the functors  $\underline{D}^i_{X_0}$  ), one can get a simple geometric interpretation of the three successive obstructions to deforming flatly  $X'$  to  $X$ , lying in the three left hand terms of $(*)$. Namely by the local theory, the first obstruction  $c_1$  in  $H^0(X_0,\ \underline{D}^2_{X_0})$  can be defined as the obstruction to finding local solutions, when this obstruction vanishes, the sheaf  $\underline{F}$  of isomorphism classes of local solutions of the problem is (by the local theory) in a canonical way a torsor under  $\underline{D}^1_{X_0}$ , and the second obstruction  $c_2$ , lying in  $H^1(X_0,\ \underline{D}^1_{X_0})$ , is the obstruction to finding a section of this torsor, i.e. of finding a  "coherent system of isomorphism classes of local solutions" ; if this second obstruction also vanishes, then, choosing arbitrarily a section of  $\underline{F}$ , one gets, as noted above, an obstruction  $c_3$  in  $H^2(X_0,\ \underline{D}^0_{X_0})$  to finding a global solution corresponding to the given system of (isomorphism classes of) local solutions. Using the fact that the section of  $\underline{F}$  is determined up to adding a section of  $\underline{D}^1_{X_0}$ , we get a map  $d : H^0(X_0,\ \underline{D}^1_{X_0}) \to H^2(X_0,\ \underline{D}^0_{X_0})$ , and a solution to the deformation problem exists if and only if  $c_1 = 0$ ,  $c_2 = 0$ , and  $c_3$  is zero as an element of Coker $d$  (each  $c_i$  being defined when the previous ones are zero). Of course, the groups  $H^0(X_0,\ \underline{D}^2_{X_0})$ ,  $H^1(X_0,\ \underline{D}^1_{X_0})$  and
\SGApage{I}{114}
$H^2(X_0,\ \underline{D}^0_{X_0})/\mathrm{Im}\ d$  are just the initial terms  $E_2^{0,2}$ ,  $E_2^{1,1}$ , and  $E_3^{2,0}$  of the spectral sequence above,  $d$  being the differential (up to sign ?)

}{F7I-VI-105}\begin{equation*}
d_2 : E_2^{0,1} \to E_2^{2,0} \quad .
\end{equation*}

5.  \underline{Formal Jacobian subschemes}

Let  $X_0$  be a scheme of finite type over a field  $k$ , and let  $(R,\ \mathfrak{X})$  be a formal versal deformation of  $X_0$ . Then  $\mathfrak{X}$  is not an ordinary scheme over  $R$  but an $R$-adic formal scheme. The purpose of this section is to clarify  "the non-smooth locus of  $\mathfrak{X}$  over  $R$ ".

\SGAvar{We fistly recall the definition of Fitting ideals associated to a module of finite type: Let  $A$  be a commutative ring and  $E$  an $A$-module of finite type. Choose a presentation

}{We first recall the definition of Fitting ideals associated to a module of finite type: Let  $A$  be a commutative ring and  $E$  an $A$-module of finite type. Choose a presentation

}{F7I-VI-106}\begin{equation*}
K \to F \to E \to 0
\end{equation*}
where  $F$  is locally-free of constant rank  $n$ . We define

\SGAvar{\begin{equation*}
I^P_A(E) = \mathrm{Im}(\overset{n-p}{\wedge} K \otimes \overset{n-p}{\wedge} F^* \to A) \qquad \text{for}\quad p = 0,\ 1,\ 2,\ \ldots
\end{equation*}
The ideals  $I^P_A(E)$  thus defined does not depend on the choice of a presentation and therefore are invariants of $A$-module  $E$ , called the Fitting ideals of $A$-module  $E$ . The following properties readily follow from the definition.

}{\begin{equation*}
I^p_A(E) = \mathrm{Im}(\overset{n-p}{\wedge} K \otimes \overset{n-p}{\wedge} F^* \to A) \qquad \text{for}\quad p = 0,\ 1,\ 2,\ \ldots
\end{equation*}
The ideals  $I^p_A(E)$  thus defined do not depend on the choice of a presentation and therefore are invariants of the $A$-module  $E$ , called the Fitting ideals of the $A$-module  $E$ . The following properties readily follow from the definition.

}{F7I-VI-107}5.1.    (a) For any ring homomorphism  $A \to B$ , the canonical map

\SGAvar{\begin{equation*}
B \underset{A}{\otimes} A/I^p_A(E) \overset{\sim}{\to} B/I^p_B(B \underset{A}{\otimes} E)
\end{equation*}
is an isomorphism for all  $P$ . In particular we have

}{\begin{equation*}
B \underset{A}{\otimes} A/I^p_A(E) \overset{\sim}{\to} B/I^p_B(B \underset{A}{\otimes} E)
\end{equation*}
is an isomorphism for all  $p$ . In particular we have

}{F7I-VI-108}\begin{equation*}
I^p_{S^{-1}A}(S^{-1}E) = S^{-1}I^p_A(E)
\end{equation*}
for any multiplicative subset  $S$  in  $A$ .


\SGApage{I}{115}

(b)  We have  $I_A^0(E) \subset I_A^1(E) \subset I_A^2(E) \subset \ \dots$ , and, for each point $x \in \operatorname{Spec}(A)$ , $I_A^p(E)_x = A_x$  for all  $p \geq \dim_{\varkappa(x)} E(x)$ , where  $E(x) = \varkappa(x) \underset{A}{\otimes} E$ .

(c)  $x \in \operatorname{Spec}(A) - \operatorname{Supp} A/I_A^p(E)$  if and only if, for any (or for some) presentation

\begin{equation*}
K \to F \to E \to 0
\end{equation*}
where  $F \approx A^n$ , $K_x$  contains  $A_x^{n-p}$  which is isomorphic to a direct summand of $F_x$ . In particular, we have that

\begin{equation*}
\operatorname{Supp} A/I_A^0(E) = \operatorname{Supp} E \ .
\end{equation*}
(d)

\begin{equation*}
I_A^r(E_1 \oplus E_2) = \sum_{p+q=r} \ I_A^p(E_1) \cdot I_A^q(E_2)
\end{equation*}
(e)  If  $A$  is a noetherian adic ring, then

\SGAvar{\begin{equation*}
A/I_A^p(E) = \varprojlim_{v} A_v \big/ I_{A_v}^p (E_v)
\end{equation*}
where  $A = \varprojlim_{v} A_v$  and  $E_v = E \underset{A}{\otimes} A_v$ .

}{\begin{equation*}
A/I_A^p(E) = \varprojlim_{\nu} A_\nu \big/ I_{A_\nu}^p (E_\nu)
\end{equation*}
where  $A = \varprojlim_{\nu} A_\nu$  and  $E_\nu = E \underset{A}{\otimes} A_\nu$ .

}{F7I-VI-109}In view of the property 5.1(a), our definition of Fitting ideals is readily globalized. Thus let  $X$  be any scheme, and  $\underline{E}$  a quasi-coherent $X$-module of finite presentation. We then obtain a sequence of quasi-coherent ideal sheaves $\underline{I}_X^p(\underline{E})$  for  $p \geq 0$  defined by

\SGAvar{\begin{equation*}
\Gamma(U, \ \underline{I}_X^p(\underline{E})) \xrightarrow{\sim} I^p_{\Gamma(U, \ \underline{O}_x)}(\Gamma(U, \ \underline{E}))
\end{equation*}
for every affine open subset  $U \subset X$ .

Now let  $X \to Y$  be a morphism of schemes. The relative Kähler-differentials $\underline{\Omega}_{X|Y}(= \underline{\Omega}^1_{X|Y})$  is a quasi-coherent $X$-module. If the morphism  $X \to Y$  is of essentially finite presentation, then  $\underline{\Omega}_{X|Y}$  is a quasi-coherent $X$-module of finite presentation, and in turn we may consider the ideal sheaves  $\underline{I}_X^p(\underline{\Omega}_{X|Y})$  ..

}{\begin{equation*}
\Gamma(U, \ \underline{I}_X^p(\underline{E})) \xrightarrow{\sim} I^p_{\Gamma(U, \ \underline{O}_X)}(\Gamma(U, \ \underline{E}))
\end{equation*}
for every affine open subset  $U \subset X$ .

Now let  $X \to Y$  be a morphism of schemes. The relative module of Kähler differentials $\underline{\Omega}_{X|Y}(= \underline{\Omega}^1_{X|Y})$  is a quasi-coherent $X$-module. If the morphism  $X \to Y$  is of essentially finite presentation, then  $\underline{\Omega}_{X|Y}$  is a quasi-coherent $X$-module of finite presentation, and in turn we may consider the ideal sheaves  $\underline{I}_X^p(\underline{\Omega}_{X|Y})$  .

}{F7I-VI-110}\underline{Definition} 5.2.  \underline{Let}  $X \to Y$  \underline{be a morphism of essentially finite presentation.} \underline{For each integer}  $p \geq 0$ , \underline{we define}  $J_X^p(X|Y)$  \underline{to be the closed subscheme of}  $X$
\SGApage{I}{116}
\underline{defined by the ideal sheaf}  $\underline{I}_X^p(X|Y) = \underline{I}_X^p(\underline{\Omega}_{X|Y})$ , \underline{called the Jacobian subschemes} \underline{of}  $X$  \underline{over}  $Y$ . \underline{For each point}  $x \in X$ , \underline{we set}  $X_x = \operatorname{Spec}(\underline{O}_{X,x})$ .

\underline{Theorem} 5.3.  \underline{Let}  $X \to Y$  \underline{be a morphism of essentially finite presentation. Then}

(1)  \underline{For any base change}  $Y \leftarrow Y'$ , \underline{the canonical map} $J^p_{X'}(X'|Y') \xrightarrow{\sim} J^p_X(X|Y) \underset{Y}{\times} Y'$  \underline{is an isomorphism for all}  $p$  \underline{where}  $X' = X \underset{Y}{\times} Y'$ .

(2)  \underline{For each point}  $x \in X$ , \underline{we have the canonical isomorphism} $J^p_{X_x}(X_x|Y) \xrightarrow{\sim} J^p_X(X|Y)_x$ .

(3)  $J^0_X(X|Y) \supset J^1_X(X|Y) \supset J^2_X(X|Y) \supset \dots,$  \underline{and}  $x \notin J^p_X(X|Y)$  \underline{for all} $p \geq \dim_{\varkappa(x)} \Omega_{X|Y}(x)$ .

(4)  $x \notin J^p_X(X|Y)$  \underline{if and only if there exists an open neighbourhood}  $U$ \underline{of}  $x$  \underline{which admits a closed immersion into a}  $Y$-\underline{scheme}  $U'$  \underline{such that}  $U'$  \underline{is} \underline{smooth over}  $Y$  at  $x'$  \underline{and}  $\dim_{x'} U'_{f(x)} = p$  \underline{where}  $x'$  \underline{is the image of}  $x$  \underline{under} $U \to U'$ .

\begin{equation*}
(5) \quad \underline{I}^r_{X_1 \underset{Y}{\times} X_2}(X_1 \underset{Y}{\times} X_2 | Y) = \sum_{p+q=r} \ \pi_1^* \underline{I}^p_{X_1}(X_1|Y) \cdot \pi_2^* \underline{I}^q_{X_2}(X_2|Y)
\end{equation*}
\underline{where}  $\pi_i : X_1 \underset{Y}{\times} X_2 \to X_i$  \underline{is the projection}.

\SGAvar{\underline{Proof}:  (3) follows from 5.1(b) whereas (1) follows from 5.1(a) together with the fact that  $P^*\underline{\Omega}_{X|Y} \xrightarrow{\sim} \underline{\Omega}_{X'|Y'}$  is an isomorphism where  $P : X \underset{Y}{\times} Y' \to X$  is the projection. (2) also follows from 5.1(a) since  $\underline{O}_{X,x} \underset{\underline{O}_X}{\otimes} \underline{\Omega}_{X|Y} \xrightarrow{\sim} \underline{\Omega}_{X_x|Y}$  is an isomorphism (5) follows from 5.1(d) and the isomorphism

}{\underline{Proof}:  (3) follows from 5.1(b) whereas (1) follows from 5.1(a) together with the fact that  $P^*\underline{\Omega}_{X|Y} \xrightarrow{\sim} \underline{\Omega}_{X'|Y'}$  is an isomorphism where  $P : X \underset{Y}{\times} Y' \to X$  is the projection. (2) also follows from 5.1(a) since  $\underline{O}_{X,x} \underset{\underline{O}_X}{\otimes} \underline{\Omega}_{X|Y} \xrightarrow{\sim} \underline{\Omega}_{X_x|Y}$  is an isomorphism. (5) follows from 5.1(d) and the isomorphism

}{F7I-VI-111}\begin{equation*}
\pi_1^*\underline{\Omega}_{X_1|Y} \oplus \pi_2^*\underline{\Omega}_{X_2|Y} \xrightarrow{\sim} \underline{\Omega}_{X_1 \underset{Y}{\times} X_2|Y}
\end{equation*}
(4)  Assume the existence of a diagram

\SGAvar{\begin{tikzcd}[column sep=huge, row sep=huge]
U \arrow[d, "f"'] \& U' \arrow[l, "i"] \arrow[dl] \\
Y \&
\end{tikzcd}











}{\begin{tikzcd}[column sep=huge, row sep=huge]
U \arrow[r, "i"] \arrow[d, "f"'] \& U' \arrow[dl] \\
Y \&
\end{tikzcd}











}{F7I-VI-112}\SGApage{I}{117}
where $U'$ is smooth over $Y$ at $x' = i(x)$ and $\dim_{x'}U'_{f(x)} = p$ . Then $\dim_{\varkappa(x')}\underline{\Omega}_{U'|Y}(x') = p$ and $i^{*}\underline{\Omega}_{U'|Y} \to \underline{\Omega}_{U|Y} \to 0$ is exact, so that $\dim_{\varkappa(x)}\underline{\Omega}_{U|Y}(x) \leq p$ which implies that $x \notin J^{p}_{U}(U|Y)$ by 5.3(3). Since $J^{p}_{U}(U|Y) = J^{p}_{X}(X|Y)|U$ , this means that $x \notin J^{p}_{X}(X|Y)$ . Now assume conversely that $x \notin J^{p}_{X}(X|Y)$ .

The question being local, we may assume that $X, Y$ are both affine, say $X = \mathrm{Spec}(B)$ , $Y = \mathrm{Spec}(A)$ . Choose a presentation over $A$

\SGAvar{\begin{equation*}
0 \to J \to P \to B \to 0
\end{equation*}
where $P = A[x_1,\ldots,x_n]$ . Then $J/J^{2} \overset{d}{\to} B \underset{P}{\otimes} \Omega_{P|A} \to \Omega_{B|A} \to 0$ is exact and $B \underset{P}{\otimes} \Omega_{P|A}$ is $B$-free of rank $n$ , and therefore $I^{p}_{B}(\Omega_{B|A}) = \mathrm{Im}(\overset{n-p}{\wedge}J/J^{2} \underset{B}{\otimes} (B \underset{P}{\otimes} \overset{n-p}{\wedge}\Omega_{P|A})^{*} \to B)$ . Therefore, $x \notin J^{p}_{X}(X|Y)$ $\Leftrightarrow$ there exists $f_1,\ldots,f_{n-p}$ in $J$ such that $\left[\det\left(\dfrac{f_i}{x_j}\right)\right](x) \neq 0$ for some choice of $n-p$ variables among $x_1, x_2,\ldots,x_n$   there exists $f_1, f_2,\ldots,f_{n-p}$ in $J$ such that $\mathrm{Spec}(P/(f_1,\ldots,f_{n-p})) \to \mathrm{Spec}(A)$ is smooth at $x'$ where $x'$ is the image of $x$ under $\mathrm{Spec}(B) \to \mathrm{Spec}(P/(f_1,\ldots,f_{n-p}))$ . Thus it suffices to take $U' = \mathrm{Spec}(P/(f_1,\ldots,f_{n-p}))$ .

}{\begin{equation*}
0 \to J \to P \to B \to 0
\end{equation*}
where $P = A[x_1,\ldots,x_n]$ . Then $J/J^{2} \overset{d}{\to} B \underset{P}{\otimes} \Omega_{P|A} \to \Omega_{B|A} \to 0$ is exact and $B \underset{P}{\otimes} \Omega_{P|A}$ is $B$-free of rank $n$ , and therefore $I^{p}_{B}(\Omega_{B|A}) = \mathrm{Im}(\overset{n-p}{\wedge}J/J^{2} \underset{B}{\otimes} (B \underset{P}{\otimes} \overset{n-p}{\wedge}\Omega_{P|A})^{*} \to B)$ . Therefore, $x \notin J^{p}_{X}(X|Y)$ $\Leftrightarrow$ there exists $f_1,\ldots,f_{n-p}$ in $J$ such that $\left[\det\left(\dfrac{\partial f_i}{\partial x_j}\right)\right](x) \neq 0$ for some choice of $n-p$ variables among $x_1, x_2,\ldots,x_n$   there exists $f_1, f_2,\ldots,f_{n-p}$ in $J$ such that $\mathrm{Spec}(P/(f_1,\ldots,f_{n-p})) \to \mathrm{Spec}(A)$ is smooth at $x'$ where $x'$ is the image of $x$ under $\mathrm{Spec}(B) \to \mathrm{Spec}(P/(f_1,\ldots,f_{n-p}))$ . Thus it suffices to take $U' = \mathrm{Spec}(P/(f_1,\ldots,f_{n-p}))$ .

}{F7I-VI-113}The following is a generalization of the usual Jacobian criteria for relative situation.

\underline{Theorem} 5.4.   \underline{Let} $X \overset{f}{\to} Y$ \underline{be a morphism of essentially finite presentation. Assume that there exists a fixed integer} $n \geq 0$ \underline{with either conditions:}

(1)  $f$  \underline{is flat and} $\dim_{x}X_{f(x)} = n$ \underline{for every point} $x \in X$ .

(2)  $f$  \underline{is a relative complete-intersection with virtual dimension} $n$ (SGA6 VIII 1.9).

\underline{Then} $x \notin J^{n}_{X}(X|Y)$ \underline{if and only if} $f$ \underline{is smooth at} $x$ .

\underline{Proof:}  We have, by 5.3(4) that $x \notin J^{n}_{X}(X|Y)$ if and only if there exists a diagram
\SGApage{I}{118}
\begin{tikzcd}
U \arrow[r, hook, "i"] \arrow[d, "f"'] \& U' \arrow[dl, "f'"] \\
Y \&
\end{tikzcd}








such that $f'$ is smooth at $x'$ and $\dim_{x'}U'_{f(x)} = n$ , where $i$ is a closed immersion and $x' = i(x)$ . Thus it suffices to show that either conditions (1) or (2) entails the isomorphism of $i$ at the point $x$ . Set $A = \underline{O}_{Y,f(x)}$ , $B' = O_{U',x'}$ , $B = O_{U,x}$ and let

\begin{equation*}
0 \to I \to B' \to B \to 0
\end{equation*}
be exact. We must show that $I = 0$ . Assume (1). Since $B$ is $A$-flat,

\begin{equation*}
0 \to I/\underline{m}I \to B'/\underline{m}B' \to B/\underline{m}B \to 0
\end{equation*}
is exact where $\underline{m}$ is the maximal ideal of $A$ . By hypothesis, $\dim_{x}U_{f(x)} = \dim_{x'}U'_{f(x)}$ so that $\dim B'/\underline{m}B' = \dim B/\underline{m}B$ . However $B'/\underline{m}B'$ is a regular local ring and therefore it entails that $I/\underline{m}I = 0$ so that $I = 0$ . Now assume (2), and consider the exact sequence

\begin{equation*}
I/I^{2} \overset{d}{\to} B \underset{B'}{\otimes} \Omega_{B'|A} \to \Omega_{B|A} \to 0
\end{equation*}
Since $X$ is a relative complete-intersection over $Y$ and $U' \to Y$ is smooth at $x$ , it follows that $d : I/I^{2} \to B \underset{B'}{\otimes} \Omega_{B'|A}$ is injective at $x$ and $I/I^{2}$ and $B \underset{B'}{\otimes} \Omega_{B'|A}$ are both free at $x$ . Therefore by the definition of virtual dimension of relative complete-intersection we must have that

\begin{equation*}
n = \dim_{\varkappa(x)}\Omega_{B'|A}(x') - \dim_{\varkappa(x)}I/I^{2}(x) = n - \dim_{\varkappa(x)}I/I^{2}(x)
\end{equation*}
i.e., $I/I^{2} = 0$ i.e., $I = 0$ .

\underline{Definition} 5.5.   \underline{Let} $X \overset{f}{\to} Y$ \underline{be a morphism of essentially finite presentation.} \underline{In case when there exists a fixed integer} $n$ \underline{with the property} (1) \underline{or} (2) \underline{in} 5.4,
\SGApage{I}{119}
\underline{we set}  $\underline{I}_X(X|Y) = \underline{I}^n_X(X|Y)$  \underline{and}  $J_X(X|Y) = J^n_X(X|Y)$ .

We can extend this notion to formal schemes in an obvious manner.

\SGAvar{Thus let $\mathfrak{X}$ be a noetherian formal scheme and $\underline{E}$ a coherent $\mathfrak{X}$-module. Then we can define the ideal sheaf $\underline{I}^P_{\mathfrak{X}}(\underline{E})$ by setting

\begin{equation*}
\underline{I}^P_{\mathfrak{X}}(\underline{E})(U) \;=\; I^P_{\mathfrak{X}(U)}(\underline{E}(U))
\end{equation*}
for every affine open $U \subset \mathfrak{X}$ . If $\mathfrak{X} = \varinjlim_{v} X_v$ , then $\underline{E} = \varinjlim_{v} \underline{E}_v$ where $\underline{E}_v = \underline{E} \otimes_{\underline{O}_{\mathfrak{X}}} \underline{O}_{X_v}$ and it follows from the fact that $\mathfrak{X}$ is noetherian that

\begin{equation*}
\underline{I}^P_{\mathfrak{X}}(\underline{E}) \;=\; \varprojlim_{v} \underline{I}^P_{X_v}(\underline{E}_v)
\end{equation*}
Now let $S$ be a noetherian formal scheme and $\mathfrak{X}$ an $S$-adic formal scheme of essentially finite type. Let $\underline{I}_{\mathfrak{X}}$ be the ideal sheaf defining the closed immersion $\Delta : \mathfrak{X} \to \mathfrak{X} \underset{S}{\times} \mathfrak{X}$ (= the fibre-product in the category of formal schemes) i.e.,

}{Thus let $\mathfrak{X}$ be a noetherian formal scheme and $\underline{E}$ a coherent $\mathfrak{X}$-module. Then we can define the ideal sheaf $\underline{I}^p_{\mathfrak{X}}(\underline{E})$ by setting

\begin{equation*}
\underline{I}^p_{\mathfrak{X}}(\underline{E})(U) \;=\; I^p_{\mathfrak{X}(U)}(\underline{E}(U))
\end{equation*}
for every affine open $U \subset \mathfrak{X}$ . If $\mathfrak{X} = \varinjlim_{v} X_v$ , then $\underline{E} = \varprojlim_{v} \underline{E}_v$ where $\underline{E}_v = \underline{E} \otimes_{\underline{O}_{\mathfrak{X}}} \underline{O}_{X_v}$ and it follows from the fact that $\mathfrak{X}$ is noetherian that

\begin{equation*}
\underline{I}^p_{\mathfrak{X}}(\underline{E}) \;=\; \varprojlim_{v} \underline{I}^p_{X_v}(\underline{E}_v)
\end{equation*}
Now let $S$ be a noetherian formal scheme and $\mathfrak{X}$ an $S$-adic formal scheme of essentially finite type. Let $\underline{I}_{\mathfrak{X}}$ be the ideal sheaf defining the closed immersion $\Delta : \mathfrak{X} \to \mathfrak{X} \underset{S}{\times} \mathfrak{X}$ (= the fibre-product in the category of formal schemes) i.e.,

}{F7I-VI-114}\begin{equation*}
0 \to \underline{I}_{\mathfrak{X}} \to \underline{O}_{\mathfrak{X} \underset{S}{\times} \mathfrak{X}} \to \underline{O}_{\mathfrak{X}} \to 0
\end{equation*}
is exact. We define $\underline{\Omega}_{\mathfrak{X}|S} = \underline{I}_{\mathfrak{X}}/I^2_{\mathfrak{X}}$ . Since $\mathfrak{X}$ is an $S$-adic formal scheme of essentially finite type, it follows that $\mathfrak{X} = \varinjlim_{v} X_v$ where $X_v = \mathfrak{X} \underset{S}{\times} S_v$ , and hence $\underline{\Omega}_{\mathfrak{X}|S} \simeq \varprojlim_{v} \underline{\Omega}_{X_v|S_v}$ is a coherent $\mathfrak{X}$-module. We note that for any formal $S$-scheme $S'$ of finite type, we have a canonical isomorphism

\begin{equation*}
\underline{\Omega}_{\mathfrak{X} \underset{S}{\times} S'|S'} \;\simeq\; \underline{O}_{S'} \otimes_{\underline{O}_S} \underline{\Omega}_{\mathfrak{X}|S} \quad .
\end{equation*}
\underline{Definition} 5.6. \underline{Let} $S$ \underline{be a noetherian formal scheme and} $\mathfrak{X}$ \underline{an} $S$-\underline{adic formal scheme of essentially finite type. We define the ideal sheaf} $\underline{I}^p_{\mathfrak{X}}(\mathfrak{X}|S)$ \underline{by}

\begin{equation*}
\underline{I}^p_{\mathfrak{X}}(\mathfrak{X}|S) \;=\; \underline{I}^p_{\mathfrak{X}}(\underline{\Omega}_{\mathfrak{X}|S})
\end{equation*}
\underline{and we set} $J^p_{\mathfrak{X}}(\mathfrak{X}|S)$ \underline{the formal closed subscheme of} $\mathfrak{X}$ \underline{defined by} $\underline{I}^p_{\mathfrak{X}}(\mathfrak{X}|S)$ .

Since $\Omega_{\mathfrak{X}|S} = \varprojlim_{n} \Omega_{X_n|S_n}$ we have 


\SGApage{I}{120}

\begin{equation*}
 I^p_{\mathfrak{X}}(\mathfrak{X}|S) \;=\; \varprojlim_{n} I^p_{X_n}(X_n|S_n)
\end{equation*}
\begin{equation*}
J^p_{\mathfrak{X}}(\mathfrak{X}|S) \;=\; \varinjlim_{n} J^p_{X_n}(X_n|S_n) \quad .
\end{equation*}
$\mathfrak{X}_{(x)}$ is an $S$-adic formal scheme of essentially finite type and we have a commutative diagram

\begin{center}
\begin{tikzcd}[column sep=huge, row sep=large]
\mathfrak{X}_{(\widehat{x})} \arrow[rr, "i_{(\widehat{x})}"] \arrow[dr] \& \& \mathfrak{X} \\
\& \mathfrak{X}_{(x)} \arrow[ur, "i_{(x)}"'] \&
\end{tikzcd}
\end{center}












We note that if $x$ is an isolated point of $\mathfrak{X}$ , then $i_{(x)}$ , $i_{(\widehat{x})}$ are both isomorphisms to the connected component of $\mathfrak{X}$ containing the point $x$ .

In particular, if $|\mathfrak{X}| = \{x_1,\dots,x_m\}$ is a finite discrete space, then we have a canonical identification $\mathfrak{X} = \coprod_{1 \leq i \leq m} \mathfrak{X}_{(x_i)}$ (disjoint union).

\underline{Corollary} 5.7. \underline{Let} $S$ \underline{be a noetherian formal scheme, and} $\mathfrak{X}$ \underline{an} $S$-\underline{adic formal scheme of essentially finite type. We then have}

\SGAvar{(1) $J^p_{\mathfrak{X}}(\mathfrak{X}|S)$ \underline{commutes with base change} i.e., $J^p_{\mathfrak{X}'}(\mathfrak{X}'|S') = J^p_{\mathfrak{X}}(\mathfrak{X}|S) \underset{S}{\times} S'$ \underline{where} $\mathfrak{X} = \mathfrak{X} \underset{S}{\times} S'$ .

(2) \underline{The connected components of} $J^p_{\mathfrak{X}}(\mathfrak{X}|S)$ \underline{is in bijective correspondence with that of} $J^p_{X_0}(X_0|S_0)$ . \underline{For each point} $x \in |J^p_{X_0}(X_0|S_0)|$ \underline{we have a canonical isomorphism} $J^p_{\mathfrak{X}}(\mathfrak{X}|S)_{(x)} \simeq J^p_{\mathfrak{X}_{(x)}}(\mathfrak{X}_{(x)}|S)$ . \underline{In particular if}

}{(1) $J^p_{\mathfrak{X}}(\mathfrak{X}|S)$ \underline{commutes with base change} i.e., $J^p_{\mathfrak{X}'}(\mathfrak{X}'|S') = J^p_{\mathfrak{X}}(\mathfrak{X}|S) \underset{S}{\times} S'$ \underline{where} $\mathfrak{X}' = \mathfrak{X} \underset{S}{\times} S'$ .

(2) \underline{The connected components of} $J^p_{\mathfrak{X}}(\mathfrak{X}|S)$ \underline{are in bijective correspondence with those of} $J^p_{X_0}(X_0|S_0)$ . \underline{For each point} $x \in |J^p_{X_0}(X_0|S_0)|$ \underline{we have a canonical isomorphism} $J^p_{\mathfrak{X}}(\mathfrak{X}|S)_{(x)} \simeq J^p_{\mathfrak{X}_{(x)}}(\mathfrak{X}_{(x)}|S)$ . \underline{In particular if}

}{F7I-VI-115}\begin{equation*}
|J^p_{X_0}(X_0|S_0)| \;=\; \{x_1,\dots,x_m\}
\end{equation*}
\underline{is discrete, then}

\SGAvar{\begin{equation*}
J^p_{\mathfrak{X}}(\mathfrak{X}|S) \;-\; \coprod_{1 \leq i \leq m} J^p_{\mathfrak{X}_{(x_i)}}(\mathfrak{X}_{(x_i)}|S)
\end{equation*}
\underline{(disjoint union)}.


}{\begin{equation*}
J^p_{\mathfrak{X}}(\mathfrak{X}|S) \;=\; \coprod_{1 \leq i \leq m} J^p_{\mathfrak{X}_{(x_i)}}(\mathfrak{X}_{(x_i)}|S)
\end{equation*}
\underline{(disjoint union)}.


}{F7I-VI-116}\SGApage{I}{121}

(3) \underline{Let} $S$ \underline{be an ordinary noetherian scheme and} $X$ \underline{an} $S$-\underline{scheme of essentially finite type. Given a closed subscheme} $S_0 \subset S$ , \underline{let} $\widehat{S}$ \underline{be the formal completion of} $S$ \underline{along} $S_0$ .

\underline{If we set}

\begin{equation*}
\widehat{X} \;=\; X \underset{S}{\times} \widehat{S} \;=\; \underline{\text{the formal completion of}}\;\; X \;\; \underline{\text{along}} \;\; X_0 \;=\; X \underset{S}{\times} S_0 \;\;,\;\; \underline{\text{then}}
\end{equation*}
\begin{equation*}
J^{p}_{\widehat{X}}(\widehat{X}|\widehat{S}) \;=\; J^{p}_{X}(X|S) \underset{S}{\times} \widehat{S} \;=\; \underline{\text{the formal completion of}} \;\; J^{p}_{X}(X|S) \;\; \underline{\text{along}} \;\; J^{p}_{X_0}(X_0|S_0) \;\;.
\end{equation*}
(4) \underline{If there exists a fixed integer} $n \geq 0$ \underline{with the property} 5.4(1) \underline{or} (2) \underline{for} $X_\nu \to S_\nu$ \underline{for all} $\nu$ , \underline{then} $x \notin J^{n}_{\mathfrak{X}}(\mathfrak{X}|S)$ \underline{if and only if} $\mathfrak{X}$ \underline{is formally smooth over} $S$ \underline{at} $x$ .

\underline{Proof}: (1) follows from $\underline{\Omega}_{\mathfrak{X}'|S'} \simeq \underline{O}_{S'} \widehat{\otimes}_{\underline{O}_S} \underline{\Omega}_{\mathfrak{X}|S}$ whereas (2) follows from

\SGAvar{\begin{equation*}
\underline{O}_{\mathfrak{X}_{(x)}} \;\widehat{\otimes}_{\underline{O}_{\mathfrak{X}}}\; \underline{\Omega}_{\mathfrak{X}|S} \;\simeq\; \underline{\Omega}_{\mathfrak{X}_{(x)}|S} \;\;.
\end{equation*}
\begin{equation*}
(3)\quad J^{p}_{\widehat{X}}(\widehat{X}|\widehat{S}) \;=\; \varinjlim_{\nu} \; J^{p}_{X \underset{S}{\times} S_\nu}\!\left(X \underset{S}{\times} S_\nu \,\middle|\, S_\nu\right) \;=\; \varinjlim_{\nu} \; J^{p}_{X}(X|S) \underset{S}{\times} S_\nu \;=\; J^{p}_{X}(X|S) \underset{S}{\times} \widehat{S} \;\;.
\end{equation*}
\begin{multline*}
(4)\quad \text{Since } \mathfrak{X} \text{ is an S-adic formal scheme, } \mathfrak{X} \to S \text{ is formally smooth at}\\ x \Leftrightarrow X_\nu \to S_\nu \text{ is formally smooth at } x \text{ for all } \nu \;\Leftrightarrow\; x \notin J^{n}_{X_\nu}(X_\nu|S_\nu) \text{ for all } \nu\\ \Leftrightarrow\; x \notin \varinjlim_{\nu} J^{n}_{X_\nu}(X_\nu|S_\nu) \;=\; J^{n}_{\mathfrak{X}}(\mathfrak{X}|S) \;\;.
\end{multline*}
\underline{Remark} 5.8. Let $\mathfrak{X}$ be an $S$-adic formal scheme of essentially finite type. Given a point $x \in \mathfrak{X}$ we also consider the $S$-adic formal scheme $\mathfrak{X}_{(\hat{X})} = \varinjlim_{\nu} \mathrm{Spec}\big(\widehat{\underline{O}_{X_\nu,x}}\big)$ . We thus have the natural morphisms $\mathfrak{X}_{(\hat{X})} \to \mathfrak{X}_{(x)} \to \mathfrak{X}$ . Now $\widehat{\underline{O}_{X_0,x}}$ is in general not essentially finite type over $S_{(y)}$ where $y$ = the image of $x$ under $\mathfrak{X} \to S$ . For simplicity, let us set $A_\nu = \underline{O}_{X_\nu,x}$ , $R_\nu = \underline{O}_{S_\nu,y}$ . Now the natural map $\Omega_{A_\nu|R_\nu} \to \Omega_{\widehat{A}_\nu|R_\nu}$ induces an isomorphism $\widehat{\Omega_{A_\nu|R_\nu}} \to \widehat{\Omega_{A_\nu|R_\nu}}$ , and therefore $\varprojlim_{\nu} \widehat{\Omega_{A_\nu|R_\nu}} \;\xrightarrow{\;\sim\;}\; \lim_{\nu} \widehat{\Omega_{\widehat{A}_\nu|R_\nu}}$ where $\widehat{\phantom{xx}}$ stands for the separable completion of a module with respect to the topology induced from the ring. Therefore if we define $\widehat{\Omega_{\mathfrak{X}_{(\hat{x})}|S}} = \varprojlim_{\nu} \widehat{\Omega_{\widehat{\underline{O}}_{X_\nu,x}|\underline{O}_{S_\nu}}}$ , then $\widehat{\Omega_{\mathfrak{X}_{(\hat{x})}|S}}$ is a coherent $\mathfrak{X}_{(\hat{x})}$-module. Thus if
\SGApage{I}{122}
we set $J^{p}_{\mathfrak{X}_{(\hat{x})}}(\mathfrak{X}_{(\hat{x})}|S)$ to be the formal closed subscheme of $\mathfrak{X}_{(\hat{x})}$ defined by the ideal sheaf $I^{p}_{\mathfrak{X}_{(\hat{x})}}(\mathfrak{X}_{(\hat{x})}|S) = I^{p}_{\mathfrak{X}_{(\hat{x})}}(\Omega_{\mathfrak{X}_{(\hat{x})}|S})$ , then we have

}{\begin{equation*}
\underline{O}_{\mathfrak{X}_{(x)}} \;\widehat{\otimes}_{\underline{O}_{\mathfrak{X}}}\; \underline{\Omega}_{\mathfrak{X}|S} \;\simeq\; \underline{\Omega}_{\mathfrak{X}_{(x)}|S} \;\;.
\end{equation*}
\begin{equation*}
(3)\quad J^{p}_{\widehat{X}}(\widehat{X}|\widehat{S}) \;=\; \varinjlim_{\nu} \; J^{p}_{X \underset{S}{\times} S_\nu}\!\left(X \underset{S}{\times} S_\nu \,\middle|\, S_\nu\right) \;=\; \varinjlim_{\nu} \; J^{p}_{X}(X|S) \underset{S}{\times} S_\nu \;=\; J^{p}_{X}(X|S) \underset{S}{\times} \widehat{S} \;\;.
\end{equation*}
\begin{multline*}
(4)\quad \text{Since } \mathfrak{X} \text{ is an S-adic formal scheme, } \mathfrak{X} \to S \text{ is formally smooth at}\\ x \Leftrightarrow X_\nu \to S_\nu \text{ is formally smooth at } x \text{ for all } \nu \;\Leftrightarrow\; x \notin J^{n}_{X_\nu}(X_\nu|S_\nu) \text{ for all } \nu\\ \Leftrightarrow\; x \notin \varinjlim_{\nu} J^{n}_{X_\nu}(X_\nu|S_\nu) \;=\; J^{n}_{\mathfrak{X}}(\mathfrak{X}|S) \;\;.
\end{multline*}
\underline{Remark} 5.8. Let $\mathfrak{X}$ be an $S$-adic formal scheme of essentially finite type. Given a point $x \in \mathfrak{X}$ we also consider the $S$-adic formal scheme $\mathfrak{X}_{(\hat{x})} = \varinjlim_{\nu} \mathrm{Spec}\big(\widehat{\underline{O}_{X_\nu,x}}\big)$ . We thus have the natural morphisms $\mathfrak{X}_{(\hat{x})} \to \mathfrak{X}_{(x)} \to \mathfrak{X}$ . Now $\widehat{\underline{O}_{X_0,x}}$ is in general not essentially finite type over $S_{(y)}$ where $y$ = the image of $x$ under $\mathfrak{X} \to S$ . For simplicity, let us set $A_\nu = \underline{O}_{X_\nu,x}$ , $R_\nu = \underline{O}_{S_\nu,y}$ . Now the natural map $\Omega_{A_\nu|R_\nu} \to \Omega_{\widehat{A}_\nu|R_\nu}$ induces an isomorphism $\widehat{\Omega_{A_\nu|R_\nu}} \to \widehat{\Omega_{\widehat{A}_\nu|R_\nu}}$ , and therefore $\varprojlim_{\nu} \widehat{\Omega_{A_\nu|R_\nu}} \;\xrightarrow{\;\sim\;}\; \varprojlim_{\nu} \widehat{\Omega_{\widehat{A}_\nu|R_\nu}}$ where $\widehat{\phantom{xx}}$ stands for the separated completion of a module with respect to the topology induced from the ring. Therefore if we define $\widehat{\Omega_{\mathfrak{X}_{(\hat{x})}|S}} = \varprojlim_{\nu} \widehat{\Omega_{\widehat{\underline{O}}_{X_\nu,x}|\underline{O}_{S_\nu,y}}}$ , then $\widehat{\Omega_{\mathfrak{X}_{(\hat{x})}|S}}$ is a coherent $\mathfrak{X}_{(\hat{x})}$-module. Thus if
\SGApage{I}{122}
we set $J^{p}_{\mathfrak{X}_{(\hat{x})}}(\mathfrak{X}_{(\hat{x})}|S)$ to be the formal closed subscheme of $\mathfrak{X}_{(\hat{x})}$ defined by the ideal sheaf $\underline{I}^{p}_{\mathfrak{X}_{(\hat{x})}}(\mathfrak{X}_{(\hat{x})}|S) = \underline{I}^{p}_{\mathfrak{X}_{(\hat{x})}}(\widehat{\Omega}_{\mathfrak{X}_{(\hat{x})}|S})$ , then we have

}{F7I-VI-117}\begin{equation*}
J^{p}_{\mathfrak{X}_{(x)}}(\mathfrak{X}_{(x)}|S)_{(\hat{x})} \;\simeq\; J^{p}_{\mathfrak{X}_{(\hat{x})}}(\mathfrak{X}_{(\hat{x})}|S) \;\;.
\end{equation*}
We note that if $x$ is an isolated point of $J^{p}_{\mathfrak{X}}(\mathfrak{X}|S)$ , then $J^{p}_{\mathfrak{X}_{(\hat{x})}}(\mathfrak{X}_{(\hat{x})}|S) \to J^{p}_{\mathfrak{X}_{(x)}}(\mathfrak{X}_{(x)}|S)$ is an isomorphism.

\SGAvar{To consider the "image" of $J^{p}_{\mathfrak{X}}(\mathfrak{X}|S)$ under the morphism $\mathfrak{X} \xrightarrow{\,f\,} S$ , one may simply take the formal subscheme of $S$ defined by the ideal sheaves $\mathrm{Ker}(\underline{O}_S \to \underline{O}_{\mathfrak{X}}/I^{p}_{\mathfrak{X}}(\mathfrak{X}|S))$ . However, it is more convenient for our purpose to use the Fitting ideal again.

}{To consider the "image" of $J^{p}_{\mathfrak{X}}(\mathfrak{X}|S)$ under the morphism $\mathfrak{X} \xrightarrow{\,f\,} S$ , one may simply take the formal subscheme of $S$ defined by the ideal sheaf $\mathrm{Ker}(\underline{O}_S \to \underline{O}_{\mathfrak{X}}/\underline{I}^{p}_{\mathfrak{X}}(\mathfrak{X}|S))$ . However, it is more convenient for our purpose to use the Fitting ideal again.

}{F7I-VI-118}\underline{Definition} 5.9. \underline{Let} $S$ \underline{be a noetherian formal scheme and} $f : \mathfrak{X} \to S$ \underline{an} $S$-\underline{adic formal scheme of essentially finite type. Assume that} $J^{p}_{\mathfrak{X}}(\mathfrak{X}|S)$ \underline{is proper over} $S$ , \underline{so that} $f_{*}\underline{O}_{J^{p}_{\mathfrak{X}}(\mathfrak{X}|S)}$ \underline{is a coherent} $\underline{O}_S$-\underline{module. We then define the ideal sheaf} $\underline{I}^{p}_{S}(\mathfrak{X}|S)$ \underline{of} $\underline{O}_S$ \underline{by}

\begin{equation*}
\underline{I}^{p}_{S}(\mathfrak{X}|S) \;=\; \underline{I}^{0}_{S}\big(f_{*}\underline{O}_{J^{p}_{\mathfrak{X}}(\mathfrak{X}|S)}\big)
\end{equation*}
\underline{and set} $J^{p}_{S}(\mathfrak{X}|S)$ \underline{to be the closed formal subscheme of} $S$ \underline{defined by the ideal sheaf} $\underline{I}^{p}_{S}(\mathfrak{X}|S)$ . \underline{In case when there exists a fixed integer} $n$ \underline{having the property} (1) \underline{or} (2) \underline{in} 5.4 \underline{for} $X_\nu \to S_\nu$ \underline{for all} $\nu$ , \underline{we set}

\begin{equation*}
\underline{I}_{\mathfrak{X}}(\mathfrak{X}|S) \;=\; \underline{I}^{n}_{\mathfrak{X}}(\mathfrak{X}|S) \;\;,\qquad J_{\mathfrak{X}}(\mathfrak{X}|S) \;=\; J^{n}_{\mathfrak{X}}(\mathfrak{X}|S)
\end{equation*}
\begin{equation*}
\underline{I}_{S}(\mathfrak{X}|S) \;=\; \underline{I}^{n}_{S}(\mathfrak{X}|S) \;\;,\qquad J_{S}(\mathfrak{X}|S) \;=\; J^{n}_{S}(\mathfrak{X}|S)
\end{equation*}
\underline{Corollary} 5.10. \underline{Let} $S$ \underline{be a noetherian formal scheme and} $\mathfrak{X} \xrightarrow{\,f\,} S$ \underline{an} $S$-\underline{adic formal scheme, essentially of finite type such that} $J^{p}_{\mathfrak{X}}(\mathfrak{X}|S)$ \underline{is proper over} $S$ . \underline{Then}

\begin{equation*}
(1)\quad \big|J^{p}_{S}(\mathfrak{X}|S)\big| \;=\; \underline{\text{the image of}} \;\; \big|J^{p}_{\mathfrak{X}}(\mathfrak{X}|S)\big| \;\; \underline{\text{under}} \;\; f : \mathfrak{X} \to S \;\;.
\end{equation*}

\SGApage{I}{123}

\SGAvar{(2)  $J^p_S(\mathfrak{X}|S)$  \underline{commute with base change}  i.e.,  $J^p_{S'}(\mathfrak{X}'|S') = J^p_S(\mathfrak{X}|S) \times_S S'$  .

}{(2)  $J^p_S(\mathfrak{X}|S)$  \underline{commutes with base change}  i.e.,  $J^p_{S'}(\mathfrak{X}'|S') = J^p_S(\mathfrak{X}|S) \times_S S'$  .

}{F7I-VI-119}(3)  If  $J^p_{\mathfrak{X}}(\mathfrak{X}|S) = \coprod_{1\leq i\leq m} \mathfrak{Y}^{(i)}$  \underline{(disjoint union), then}

$\underline{I}^p_S(\mathfrak{X}|S) = \prod_{1\leq i\leq m} \underline{I}^0_S(f_*\underline{O}_{\mathfrak{Y}(i)})$  . \underline{In particular if}

\begin{equation*}
|J^p_{X_0}(X_0|S_0)| = \{x_1,\ldots,x_m\}
\end{equation*}
\underline{is discrete, then}

\begin{equation*}
\underline{I}^p_S(\mathfrak{X}|S) = \prod_{1\leq i\leq m} \underline{I}^p_S(\mathfrak{X}_{(x_i)}|S) = \prod_{1\leq i\leq m} \underline{I}^p_S(\mathfrak{X}_{(\hat{x}_i)}|S)
\end{equation*}
(4)  \underline{Let}  $S$  \underline{be an ordinary noetherian scheme and}  $X$  \underline{an} $S$-\underline{scheme of essentially finite type. Given a closed subscheme}  $S_0 \to S$  , \underline{let}  $\hat{S}$  \underline{be the formal completion of}  $S$  \underline{along}  $S_0$  . \underline{If we set}  $\hat{X} = X \times_S \hat{S}$ $=$ \underline{the formal completion of}  $X$  \underline{along}  $X_0 = X \times_S S_0$  , \underline{then}  $J^p_{\hat{S}}(\hat{X}|\hat{S}) = J^p_S(X|S) \times_S \hat{S}$ $=$ \underline{the formal completion of}  $J^p_S(X|S)$  \underline{along}  $J^p_{S_0}(X_0|S_0)$  .

\underline{Proof}:  Supp $f_*\underline{O}_{J^p_{\mathfrak{X}}(\mathfrak{X}|S)}$ $=$ $|J^p_S(\mathfrak{X}|S)|$  by 5.1(c), and hence (1) follows. (2) and (4) are clear from 5.1(a), and (3) follows from 5.1(d), and the isomorphism $J^p_{\mathfrak{X}_{(\hat{x})}}(\mathfrak{X}_{(\hat{x})}|S) \xrightarrow{\sim} J^p_{\mathfrak{X}_{(x)}}(\mathfrak{X}_{(x)}|S)$  for an isolated point  $x$  in  $J^p_{\mathfrak{X}}(\mathfrak{X}|S)$  .

\SGAvar{\underline{Remark} 5.11.  Our definition of Jacobian ideal, adopted from the classical one over base field, commutes with base change and gives a smoothness criteria under a mild condition of 5.4. However one can define, probably in many different ways, a canonical subscheme which gives the non-smooth locus in a general setting. We shall outline one here:  Given an $A$-algebra  $B$  , define  $I_{B|A} = \bigcap_E \mathrm{Ann}\ D^1(B|A,E)$ where  $E$  runs through all  $B$-modules of finite type. This ideal  $I_{B|A}$  commutes with localizations by 3.13(c), and therefore, given a morphism  $X \xrightarrow{\ f\ } Y$  , we obtain the ideal sheaf  $\underline{I}_{X|Y}$  . If  $f$  is a morphism of essentially finite type, then  $f$  is smooth at  $x$  if and only if  $x \notin \mathrm{Supp}\ \underline{O}_X/\underline{O}_{X|Y}$  by 3.13(a). This notion is in an obvious way carried over to adic formal schemes of essentially
\SGApage{I}{124}
finite type and thus it is equally suitable for deformation theory.


}{\underline{Remark} 5.11.  Our definition of Jacobian ideal, adopted from the classical one over base field, commutes with base change and gives a smoothness criteria under a mild condition of 5.4. However one can define, probably in many different ways, a canonical subscheme which gives the non-smooth locus in a general setting. We shall outline one here:  Given an $A$-algebra  $B$  , define  $I_{B|A} = \bigcap_E \mathrm{Ann}\ D^1(B|A,E)$ where  $E$  runs through all  $B$-modules of finite type. This ideal  $I_{B|A}$  commutes with localizations by 3.13(c), and therefore, given a morphism  $X \xrightarrow{\ f\ } Y$  , we obtain the ideal sheaf  $\underline{I}_{X|Y}$  . If  $f$  is a morphism of essentially finite type, then  $f$  is smooth at  $x$  if and only if  $x \notin \mathrm{Supp}\ \underline{O}_X/\underline{I}_{X|Y}$  by 3.13(a). This notion is in an obvious way carried over to adic formal schemes of essentially
\SGApage{I}{124}
finite type and thus it is equally suitable for deformation theory.


}{F7I-VI-120}6.  \underline{Non-degenerate quadratic singularities}

The objective in this section is to investigate the deformations of a $k$-scheme having non-degenerate quadratic singularities only. In view of 4.13, it becomes purely a local problem in an appropriate situation.

We go back to the notations of  §4. Thus  $\Lambda$  is a fixed complete noetherian local ring with the residue field  $k$  , and  $\underline{C}_\Lambda$  stands for the category of local artinian $\Lambda$-algebras with the same residue field  $k$  .

\SGAvar{Let  $X$  be a $k$-scheme of finite type, and let  $x$  be a closed point of $X$  . We recall that the point  $x$  is said to be a non-degenerate quadratic singularity rational over  $k$  if  $\hat{O}_{X,x}$ $-$ $k$ $[[x_1,\ldots,x_n]]/(f)$  with $(\partial f/\partial x_1,\ldots,\partial f/\partial x_n) = (x_1,\ldots,x_n)$  . One  can always assume  $f$  to be a quadratic form. Indeed we have

}{Let  $X$  be a $k$-scheme of finite type, and let  $x$  be a closed point of $X$  . We recall that the point  $x$  is said to be a non-degenerate quadratic singularity rational over  $k$  if  $\hat{O}_{X,x}$ $\simeq$ $k$ $[[x_1,\ldots,x_n]]/(f)$  with $(\partial f/\partial x_1,\ldots,\partial f/\partial x_n) = (x_1,\ldots,x_n)$  . One  can always assume  $f$  to be a quadratic form. Indeed we have

}{F7I-VI-121}\underline{Lemma} 6.1.  \underline{Let}  $f \in k[[x_1,\ldots,x_n]]$  . \underline{If}  $(\partial f/\partial x_1,\ldots,\partial f/\partial x_n) = M$  $(= (x_1,\ldots,x_n))$, \underline{then one can find}  $x_1',\ldots,x_n'$  \underline{with}  $x_i = x_i'$  (mod $M^2$)  \underline{such that} $f = q(x_1',\ldots,x_n')$  \underline{where}  $q$  \underline{is a quadratic form.}

\SGAvar{\underline{Proof}:  We do this by a successive approximation. Write  $f = f_2 + h_3$  where  $f_2$ is a quadratic form and  $h_3 \in M^3$  . Since  $(\partial f/\partial x_1,\ldots,\partial f/\partial x_n) = M$  , there exists $\epsilon_1^{(1)},\ldots,\epsilon_n^{(1)}$  in  $M^2$  such that  $\Sigma\ \epsilon_i^{(1)}\ \partial f/\partial x_i \equiv h_3$  (mod $M^4$)  . If we set $X_i^{(1)} = X_i - \epsilon_i^{(1)}$  , then  $f(x_1^{(1)},\ldots,x_n^{(1)}) =$ $f(x_1,\ldots,x_n) - \Sigma\ \epsilon_i^{(1)}\ \partial f/\partial x_i + \ldots \equiv f_2(x_1,\ldots,x_n)$  (mod $M^4$)  i.e., $f(x_1^{(1)},\ldots,x_n^{(1)}) \equiv f_2(x_1,\ldots,x_n) + h_4$  with  $h_4 \in M^4$  . Now $(\partial f/\partial x_1^{(1)},\ldots,\partial f/\partial x_n^{(1)}) = (\partial f/\partial x_1,\ldots,\partial f/\partial x_n) = M$  and therefore we can choose
\SGApage{I}{125}
$\varepsilon_1^{(2)},\dots,\varepsilon_n^{(2)}$  in  $M^3$  such that  $\Sigma\ \varepsilon_i^{(2)}\ \partial f/\partial x_u^{(1)} \equiv h_4\ (\mathrm{mod}\ M^5)$ . Then $f(x_1^{(1)} - \varepsilon_1^{(2)},\dots,x_n^{(1)} - \varepsilon_n^{(2)}) = f(x_1^{(1)},\dots,x_n^{(1)}) - \Sigma\ \varepsilon_i^{(2)}\partial f/\partial x_i^{(1)} + \dots \equiv f_2(x_1,\dots,x_n)\ (\mathrm{mod}\ M^5)$ . Continuing this process we can find $\varepsilon_i^{(1)},\ \varepsilon_i^{(2)},\ \varepsilon_i^{(3)},\dots$  $(i = 1,2,\dots,n)$  with  $\varepsilon_i^{(v)} \in M^{v+1}$  such that if we set $x_i^{(v)} = x_i - (\varepsilon_i^{(1)} + \dots + \varepsilon_i^{(v)})$  then  $f(x_1^{(v)},\dots,x_n^{(v)}) \equiv f_2(x_1,\dots,x_n)$ $(\mathrm{mod}\ M^{v+3})$ . Thus if we set  $\varepsilon_i = \sum_{v} \varepsilon_i^{(v)}$  then  $f(x_1-\varepsilon_1,\dots,x_n-\varepsilon_n) =$ $= f_2(x_1,\dots,x_n)$  i.e.,  $f(x_1,\dots,x_n) = f_2(x_1+\varepsilon_1,\dots,x_n+\varepsilon_n)$ .

}{\underline{Proof}:  We do this by a successive approximation. Write  $f = f_2 + h_3$  where  $f_2$ is a quadratic form and  $h_3 \in M^3$  . Since  $(\partial f/\partial x_1,\ldots,\partial f/\partial x_n) = M$  , there exists $\epsilon_1^{(1)},\ldots,\epsilon_n^{(1)}$  in  $M^2$  such that  $\Sigma\ \epsilon_i^{(1)}\ \partial f/\partial x_i \equiv h_3$  (mod $M^4$)  . If we set $x_i^{(1)} = x_i - \epsilon_i^{(1)}$  , then  $f(x_1^{(1)},\ldots,x_n^{(1)}) =$ $f(x_1,\ldots,x_n) - \Sigma\ \epsilon_i^{(1)}\ \partial f/\partial x_i + \ldots \equiv f_2(x_1,\ldots,x_n)$  (mod $M^4$)  i.e., $f(x_1^{(1)},\ldots,x_n^{(1)}) \equiv f_2(x_1,\ldots,x_n) + h_4$  with  $h_4 \in M^4$  . Now $(\partial f/\partial x_1^{(1)},\ldots,\partial f/\partial x_n^{(1)}) = (\partial f/\partial x_1,\ldots,\partial f/\partial x_n) = M$  and therefore we can choose
\SGApage{I}{125}
$\varepsilon_1^{(2)},\dots,\varepsilon_n^{(2)}$  in  $M^3$  such that  $\Sigma\ \varepsilon_i^{(2)}\ \partial f/\partial x_i^{(1)} \equiv h_4\ (\mathrm{mod}\ M^5)$ . Then $f(x_1^{(1)} - \varepsilon_1^{(2)},\dots,x_n^{(1)} - \varepsilon_n^{(2)}) = f(x_1^{(1)},\dots,x_n^{(1)}) - \Sigma\ \varepsilon_i^{(2)}\partial f/\partial x_i^{(1)} + \dots \equiv f_2(x_1,\dots,x_n)\ (\mathrm{mod}\ M^5)$ . Continuing this process we can find $\varepsilon_i^{(1)},\ \varepsilon_i^{(2)},\ \varepsilon_i^{(3)},\dots$  $(i = 1,2,\dots,n)$  with  $\varepsilon_i^{(v)} \in M^{v+1}$  such that if we set $x_i^{(v)} = x_i - (\varepsilon_i^{(1)} + \dots + \varepsilon_i^{(v)})$  then  $f(x_1^{(v)},\dots,x_n^{(v)}) \equiv f_2(x_1,\dots,x_n)$ $(\mathrm{mod}\ M^{v+3})$ . Thus if we set  $\varepsilon_i = \sum_{v} \varepsilon_i^{(v)}$  then  $f(x_1-\varepsilon_1,\dots,x_n-\varepsilon_n) =$ $= f_2(x_1,\dots,x_n)$  i.e.,  $f(x_1,\dots,x_n) = f_2(x_1+\varepsilon_1,\dots,x_n+\varepsilon_n)$ .

}{F7I-VI-122}\underline{Lemma} 6.2.   \underline{Let}  $A_0$  \underline{be a local} $k$-\underline{algebra, essentially of finite type such that}

\begin{equation*}
\hat{A}_0 \simeq k[[x_1,\dots,x_n]]/(f)
\end{equation*}
\underline{where}  $f$  \underline{is a non-degenerate quadratic form, and let}  $(R,\ \mathfrak{X})$  \underline{be a formal versal deformation of}  $X_0 = \mathrm{Spec}(A_0)$  . \underline{Then}

(1)  $J_{\mathfrak{X}}(\mathfrak{X}|R) \rightarrow J_R(\mathfrak{X}|R) \rightarrow \mathrm{Spec}(\Lambda)$  \underline{are both isomorphisms.}

\SGAvar{(2)  $J_R(\mathfrak{X}|R)$  \underline{is a principal ideal in}  $R$  \underline{generated by a free formal generator of}  $R$  \underline{over}  $\Lambda$  i.e.  $I_R(\mathfrak{X}|R) = (t)$  \underline{and}  $R = \Lambda[[t]]$  .

}{(2)  $I_R(\mathfrak{X}|R)$  \underline{is a principal ideal in}  $R$  \underline{generated by a free formal generator of}  $R$  \underline{over}  $\Lambda$  i.e.  $I_R(\mathfrak{X}|R) = (t)$  \underline{and}  $R = \Lambda[[t]]$  .

}{F7I-VI-123}(3)  $\mathfrak{X} \rightarrow \mathrm{Spec}(\Lambda)$  \underline{is formally smooth.}

\SGAvar{\underline{Proof}:  We set  $\mathfrak{X} = \mathrm{Spf}(A)$  . By 4.10(b),  $(R,\ \mathrm{Spf}(\hat{A}))$  is a formal versal deformation of  $\hat{X}_0 = \mathrm{Spec}(\hat{A}_0)$  where  $\hat{A}_0 = \varprojlim_{v} A \otimes_R R_v$  . Now  $\hat{A}_0 \simeq k[[x_1,\dots,x_n]]/(f)$  and  $(\partial f/\partial x_1,\dots,\partial f/\partial x_n) = (x_1,\dots,x_n)$  in $k[[x_1,\dots,x_n]]$  . It follows readily that  $[\underline{M}_{\hat{X}_0}(k[\varepsilon])] \simeq k$  and, by 4.14, we get

}{\underline{Proof}:  We set  $\mathfrak{X} = \mathrm{Spf}(A)$  . By 4.10(b),  $(R,\ \mathrm{Spf}(\hat{A}))$  is a formal versal deformation of  $\hat{X}_0 = \mathrm{Spec}(\hat{A}_0)$  where  $\hat{A} = \varprojlim_{v} A \otimes_R R_v$  . Now  $\hat{A}_0 \simeq k[[x_1,\dots,x_n]]/(f)$  and  $(\partial f/\partial x_1,\dots,\partial f/\partial x_n) = (x_1,\dots,x_n)$  in $k[[x_1,\dots,x_n]]$  . It follows readily that  $[\underline{M}_{\hat{X}_0}(k[\varepsilon])] \simeq k$  and, by 4.14, we get

}{F7I-VI-124}\begin{equation*}
R \simeq \Lambda[[t]] \quad\text{and}\quad \hat{A} \simeq R[[x_1,\dots,x_n]]/(F - t)
\end{equation*}
where  $F$  is a quadratic polynomial in  $\Lambda[x_1,\dots,x_n]$  $(\subset R[[x_1,\dots,x_n]])$ obtained from  $f$  by lifting the coefficients to  $\Lambda$  . Now

\SGAvar{\begin{equation*}
0 \rightarrow \hat{A} \ \overset{df}{\relbar\joinrel\relbar\joinrel\relbar}\ \hat{A}dx_1 \oplus \dots \oplus \hat{A}dx_n \ \relbar\joinrel\relbar\joinrel\relbar\joinrel\relbar\joinrel\relbar\ \hat{\Omega}_{\hat{A}|R} \rightarrow 0 
\end{equation*}

}{\begin{equation*}
0 \rightarrow \hat{A} \ \overset{dF}{\longrightarrow}\ \hat{A}dx_1 \oplus \dots \oplus \hat{A}dx_n \ \xrightarrow{\hspace{3em}}\ \hat{\Omega}_{\hat{A}|R} \rightarrow 0 
\end{equation*}

}{F7I-VI-125}\SGApage{I}{126}

\SGAvar{ is exact where  $df : \hat{A} \rightarrow \hat{A}dx_1 \oplus \dots \oplus \hat{A}dx_n$  is given by $1 \rightarrow (\partial f/\partial x_1)dx_1 + \dots + (\partial f/\partial x_n)dx_n$  , and therefore $I_{\hat{A}}(\hat{A}|R) = (\partial F/\partial x_1,\dots,\partial F/\partial x_n) = (x_1,\dots,x_n)$  , and hence $\hat{A}/I_{\hat{A}}(\hat{A}|R) = \hat{A}/(x_1,\dots,x_n) = R[[x_1,\dots,x_n]]/(x_1,\dots,x_n,F - t) \simeq \Lambda$ i.e., the canonical map  $\Lambda \rightarrow \hat{A}/I_{\hat{A}}(\hat{A}|R)$  is an isomorphism. It follows that $\hat{A} \rightarrow A/I_A(A|R)$  is an isomorphism, and consequently  $\Lambda \rightarrow R/I_R(A|R)$  is also an isomorphism. This proves (1) and (2). On the other hand, (3) follows from $\hat{A} \simeq R[[x_1,\dots,x_n]]/(F - t) \simeq \Lambda[[t,x_1,\dots,x_n]]/(F - t) \simeq \Lambda[[x_1,\dots,x_n]]$   as $\Lambda$-algebras.

}{ is exact where  $dF : \hat{A} \rightarrow \hat{A}dx_1 \oplus \dots \oplus \hat{A}dx_n$  is given by $1 \rightarrow (\partial F/\partial x_1)dx_1 + \dots + (\partial F/\partial x_n)dx_n$  , and therefore $I_{\hat{A}}(\hat{A}|R) = (\partial F/\partial x_1,\dots,\partial F/\partial x_n) = (x_1,\dots,x_n)$  , and hence $\hat{A}/I_{\hat{A}}(\hat{A}|R) = \hat{A}/(x_1,\dots,x_n) = R[[x_1,\dots,x_n]]/(x_1,\dots,x_n,F - t) \simeq \Lambda$ i.e., the canonical map  $\Lambda \rightarrow \hat{A}/I_{\hat{A}}(\hat{A}|R)$  is an isomorphism. It follows that $\hat{A}/I_{\hat{A}}(\hat{A}|R) \rightarrow A/I_A(A|R)$  is an isomorphism, and consequently  $\Lambda \rightarrow R/I_R(A|R)$  is also an isomorphism. This proves (1) and (2). On the other hand, (3) follows from $\hat{A} \simeq R[[x_1,\dots,x_n]]/(F - t) \simeq \Lambda[[t,x_1,\dots,x_n]]/(F - t) \simeq \Lambda[[x_1,\dots,x_n]]$   as $\Lambda$-algebras.

}{F7I-VI-126}\underline{Corollary} 6.3.  \underline{Let}  $X_0$  \underline{be a k-scheme of finite type, and let}  $x$  \underline{in}  $X_0$  \underline{be a non-degenerate quadratic singularity rational over}  $k$  .

(1)  \underline{Let}  $(R,\ \mathfrak{X})$ ,  $(R,\ \mathfrak{Y})$  \underline{be any two deformations of}  $X_0$  \underline{where} $R \in \mathrm{ob}(\hat{\underline{C}}_{\Lambda})$  . \underline{Assume that}  $k = k^2$  . \underline{Then}  $\mathfrak{X}_{(\widehat{x})} \simeq \mathfrak{Y}_{(\widehat{x})}$  \underline{as} $R$-\underline{adic formal schemes if and only if}  $I_R(\mathfrak{X}_{(\mathrm{x})}|R) = I_R(\mathfrak{Y}_{(\mathrm{x})}|R)$  .

(2)  \underline{Let}  $R \in \mathrm{ob}(\hat{\underline{C}}_{\Lambda})$  \underline{be a regular local ring, and let}  $(R,\ \mathfrak{X})$  \underline{be a deformation of}  $X_0$  . \underline{Then}  $O_{\mathfrak{X},x}$  \underline{is a regular local ring if and only if} $I_R(\mathfrak{X}_{(\mathrm{x})}|R) \not\subset \underline{m}_R^{\,2}$  i.e., \underline{if and only if}  $I_R(\mathfrak{X}_{(\mathrm{x})}|R)$  \underline{is a regular divisor in} $\mathrm{Spf}(R)$  .

\underline{Proof:}  Let  $(R_0,\ Z)$  be a formal versal deformation of  $\mathrm{Spec}(\hat{\underline{O}}_{X_0,x})$  . Since  $x$ in  $X_0$  is a non-degenerate quadratic singularity rational over  $k$  , we may assume that  $R_0 \simeq \Lambda[[t]]$  and  $Z \simeq \mathrm{Spf}(R_0[[x_1,\dots,x_n]]/(F - t))$  as in 6.2.

(1)  $(R,\ \mathfrak{X}_{(\widehat{x})})$ ,  $(R,\ \mathfrak{Y}_{(\widehat{x})})$  are deformations of  $\mathrm{Spec}(\hat{\underline{O}}_{X_0,x})$  , and therefore they are induced from  $(R_0,\ Z)$  by continuous local $\Lambda$-algebra maps $\varphi,\ \psi : R_0 \rightarrow R$  respectively, i.e., if we set  $a = \varphi(t)$  ,  $b = \psi(t)$  , then

\SGAvar{\begin{equation*}
\mathfrak{X}_{(\widehat{x})} \simeq \mathrm{Spf}(R[[X_1,\dots,X_n]]/(R - a)) \quad ,
\end{equation*}
\begin{equation*}
\mathfrak{Y}_{(\widehat{x})} \simeq \mathrm{Spf}(R[[X_1,\dots,X_n]]/(F - b))
\end{equation*}
as $R$-adic formal schemes. Assume that  $I_R(\mathfrak{X}_{(\mathrm{x})}|R) = I_R(\mathfrak{Y}_{(\mathrm{x})}|R)$  i.e.,
\SGApage{I}{127}
$(a) = (b)$ in $R$ . Then $a = ab$ for some unit $u$ in $R$ . Since $k = k^2$ and $R$ is complete, we can write $u = v^2$ for some unit $v$ in $R$ , i.e., $a = v^2b$ . Then $(F - b) = (v^2F - a)$ , so that $X_i \to vX_i$ establishes an $R$-algebra isomorphism $R[[X_1,\ldots,X_n]]/(F - a) \;\overset{\sim}{\relbar\relbar\relbar}\; R[[X_1,\ldots,X_n]]/(F - b)$ since $F$ is a quadratic form. The other direction is trivial.

}{\begin{equation*}
\mathfrak{X}_{(\widehat{x})} \simeq \mathrm{Spf}(R[[X_1,\dots,X_n]]/(F - a)) \quad ,
\end{equation*}
\begin{equation*}
\mathfrak{Y}_{(\widehat{x})} \simeq \mathrm{Spf}(R[[X_1,\dots,X_n]]/(F - b))
\end{equation*}
as $R$-adic formal schemes. Assume that  $I_R(\mathfrak{X}_{(\mathrm{x})}|R) = I_R(\mathfrak{Y}_{(\mathrm{x})}|R)$  i.e.,
\SGApage{I}{127}
$(a) = (b)$ in $R$ . Then $a = ub$ for some unit $u$ in $R$ . We construct a matrix $T\in\operatorname{GL}_n(R)$ such that $F(TX)=uF(X)$.\footnote{Editorial proof repair: the printed inference that every unit of $R$ has a square root need not hold in characteristic $2$. The following argument retains the original hypotheses.} Write $\mathfrak m=\mathfrak m_R$, and interpret polynomial congruences coefficientwise. Since $k=k^2$, choose $c\in R^\times$ whose residue satisfies $\bar c^2=\bar u$, and put $T_1=cI_n$. Then $F(T_1X)\equiv uF(X)\pmod{\mathfrak m}$. By the nondegeneracy condition in Lemma 6.2, the coefficient matrix of the linear forms $\partial F/\partial X_i$ is invertible modulo $\mathfrak m$, hence over $R$.  Suppose $T_\nu\in\operatorname{GL}_n(R)$ satisfies $F(T_\nu X)\equiv uF(X)\pmod{\mathfrak m^\nu}$, with $\nu\geq1$. The forms $(\partial F/\partial X_i)(T_\nu X)$ form an $R$-basis of the linear forms. Consequently, there are linear forms $L_i$ with coefficients in $\mathfrak m^\nu$ such that \[ E_\nu(X):=uF(X)-F(T_\nu X)  =\sum_i L_i(X)\frac{\partial F}{\partial X_i}(T_\nu X). \] Indeed, every quadratic monomial is a product of two linear forms, and expressing one factor in this basis preserves its coefficient ideal. Define $H_\nu$ by $H_\nu X=(L_i(X))_i$, and set $T_{\nu+1}=T_\nu+H_\nu$. The quadratic identity gives \[ F(T_{\nu+1}X)=F(T_\nu X)+E_\nu(X)+F(H_\nu X)  \equiv uF(X)\pmod{\mathfrak m^{\nu+1}}, \] since the coefficients of $F(H_\nu X)$ lie in $\mathfrak m^{2\nu}\subseteq\mathfrak m^{\nu+1}$. Moreover, $T_{\nu+1}\equiv T_\nu\pmod{\mathfrak m^\nu}$. Completeness and separatedness therefore give a limit $T\equiv cI_n\pmod{\mathfrak m}$ with $F(TX)=uF(X)$; in particular, $T$ is invertible. The substitution $X\mapsto TX$ sends $F-a$ to $u(F-b)$ and induces the required isomorphism \[ R[[X_1,\ldots,X_n]]/(F-a)  \xrightarrow{\ \sim\ } R[[X_1,\ldots,X_n]]/(F-b). \] The converse follows from invariance of $I_R$ under $R$-isomorphisms.

}{F7I-VI-127}(2)  There exists a continuous local $\Lambda$-algebra map $\varphi : R_0 = \Lambda[[t]] \to R$ such that $\mathfrak{X}_{(\widehat{x})} \simeq \mathrm{Spf}(R[[X_1,\ldots,X_n]]/(F - \varphi(t)))$ . We note that $I_R(\mathfrak{X}_{(x)}|R) = I_R(\mathfrak{X}_{(\widehat{x})}|R) = (\varphi(t))$ , and that $R[[X_1,\ldots,X_n]]/(F - \varphi(t)) = \hat{\underline{O}}_{\mathfrak{X},x}$ is regular if and only if $\underline{O}_{\mathfrak{X},x}$ is regular. Now $R$ is regular and hence $R[[X_1,\ldots,X_n]]$ is also a regular local ring. Thus if we set $M$ to be the maximal ideal of $R[[X_1,\ldots,X_n]]$ , then $R[[X_1,\ldots,X_n]]/(F - \varphi(t))$ is regular if and only if $F - \varphi(t)$ is not in $M^2$ , i.e. if and only if $\varphi(t) \notin \underline{m}_R^{\,2}$ , i.e., if and only if $I_R(\mathfrak{X}_{(x)}|R) \not\subset \underline{m}_R^{\,2}$ .

\underline{Theorem} 6.4.  \underline{Let} $X_0$ \underline{be a} $k$-\underline{scheme of finite type with isolated non-smooth points only, and let} $(R, \mathfrak{X})$ \underline{be a formal versal deformation of} $X_0$ . \underline{Denote by} $\{z_1,\ldots,z_m\}$ \underline{the set of non-smooth points in} $X_0$ . \underline{Assume that} $H^2(X_0, \underline{D}^0_{X_0}) = 0$ , \underline{and all} $z_i$'s \underline{are non-degenerate quadratic singularity rational over} $k$ . \underline{Then}

(1)  $R$ \underline{is formally smooth over} $\Lambda$ i.e., $R$ \underline{is a formal power-series ring over} $\Lambda$ ;

(2)  $J_{\mathfrak{X}}(\mathfrak{X}|R) = \coprod_{1\leq i\leq m} J_{\mathfrak{X}_{(z_i)}}(\mathfrak{X}_{(z_i)}|R)$  \underline{(disjoint union), and}

$J_{\mathfrak{X}_{(z_i)}}(\mathfrak{X}_{(z_i)}|R) \;\overset{\sim}{\relbar\relbar\relbar}\; J_R(\mathfrak{X}_{(z_i)}|R)$  \underline{is an isomorphism};

(3)  $I_R(\mathfrak{X}_{(z_i)}|R)$ \underline{is a principal ideal and in turn so is}

$I_R(\mathfrak{X}|R) = \prod_{1\leq\ i\leq\ m} I_R(\mathfrak{X}_{(z_i)}|R)$ . \underline{Furthermore if we set} $I_R(\mathfrak{X}_{(z_i)}|R) = (t_i)$ , \underline{then} $t_1,\ldots,t_m$ \underline{can be extended to a system of free formal generators of} $R$ \underline{over} $\Lambda$ (\underline{as formal power-series ring over} $\Lambda$ ).

(4)  $\mathfrak{X} \to \mathrm{Spec}(\Lambda)$ \underline{is smooth}, i.e. $J_{\mathfrak{X}}(\mathfrak{X}|\Lambda) = \varphi$ .


\SGApage{I}{128}

(5)  \underline{If} $X_0$ \underline{is a complete irreducible curve, then} $\dim \underline{M}_{X_0} = 3g - 3 + a$ \underline{where} $g = \dim_k H^1(X_0, \underline{O}_{X_0})$ \underline{and} $a = \dim_k H^0(X_0, \underline{D}^0_{X_0})$ .

\SGAvar{\underline{Proof:}  (1) follows from (2) and (3): The first statement of (2) follows from 5.7(2). Now let $(R_i, \mathfrak{Y}_i)$ be a formal versal deformation of $\mathrm{Spec}(\hat{\underline{O}}_{X_0,z_i})$ . Since $(R, \mathfrak{X}_{(\widehat{z}_i)})$ is a deformation of $\mathrm{Spec}(\hat{\underline{O}}_{X_0,z_i})$ , there exists a continuous local $\Lambda$-algebra maps $\varphi_i = R_i \to R$ with a cartesian diagram

}{\underline{Proof:}  (1) follows from (2) and (3): The first statement of (2) follows from 5.7(2). Now let $(R_i, \mathfrak{Y}_i)$ be a formal versal deformation of $\mathrm{Spec}(\hat{\underline{O}}_{X_0,z_i})$ . Since $(R, \mathfrak{X}_{(\widehat{z}_i)})$ is a deformation of $\mathrm{Spec}(\hat{\underline{O}}_{X_0,z_i})$ , there exists a continuous local $\Lambda$-algebra maps $\varphi_i : R_i \to R$ with a cartesian diagram

}{F7I-VI-128}\begin{center}
\begin{tikzcd}[column sep=huge, row sep=huge]
\mathfrak{X}_{(\widehat{z}_i)} \arrow[r] \arrow[d] \& \mathfrak{Y}_i \arrow[d] \\
\mathrm{Spf}(R) \arrow[r, "{\mathrm{Spf}(\varphi_i)}"'] \& \mathrm{Spf}(R_i)
\end{tikzcd}
\end{center}













Therefore we have

\begin{equation*}
J_{\mathfrak{X}_{(z_i)}}(\mathfrak{X}_{(z_i)}|R) = J_{\mathfrak{X}_{(\widehat{z}_i)}}(\mathfrak{X}_{(\widehat{z}_i)}|R) \simeq J_{\mathfrak{Y}_i}(\mathfrak{Y}_i|R_i) \underset{\mathrm{Spf}(R_i)}{\times} \mathrm{Spf}(R)
\end{equation*}
\begin{equation*}
J_R(\mathfrak{X}_{(z_i)}|R) = J_R(\mathfrak{X}_{(\widehat{z}_i)}|R) \simeq J_{R_i}(\mathfrak{Y}_i|R_i) \underset{\mathrm{Spf}(R_i)}{\times} \mathrm{Spf}(R) \quad .
\end{equation*}
Since $J_{\mathfrak{Y}_i}(\mathfrak{Y}_i|R_i) \to J_{R_i}(\mathfrak{Y}_i|R_i)$ is an isomorphism by 6.2(1), we have that the canonical morphism $J_{\mathfrak{X}_{(z_i)}}(\mathfrak{X}_{(z_i)}|R) \to J_R(\mathfrak{X}_{(z_i)}|R)$ is an isomorphism. Now $I_{R_i}(\mathfrak{Y}_i|R_i) = (u_i) = (u_i)$ with $R_i = \Lambda[[u_i]]$ by 6.2(2), and therefore

\SGAvar{\begin{equation*}
I_R(\mathfrak{X}_{(z_i)}|R) = I_R(\mathfrak{X}_{(\widehat{z}_i)}|R) = (t_i)
\end{equation*}
where $t_i = \varphi_i(u_i)$ , and therefore $I_R(\mathfrak{X}|R) = \prod_{1\leq\ i\leq\ m} I_R(\mathfrak{X}_{(z_i)}|R) = (\ \prod_{i=1}^{n} t_i\ )$ .

}{\begin{equation*}
I_R(\mathfrak{X}_{(z_i)}|R) = I_R(\mathfrak{X}_{(\widehat{z}_i)}|R) = (t_i)
\end{equation*}
where $t_i = \varphi_i(u_i)$ , and therefore $I_R(\mathfrak{X}|R) = \prod_{1\leq\ i\leq\ m} I_R(\mathfrak{X}_{(z_i)}|R) = (\ \prod_{i=1}^{m} t_i\ )$ .

}{F7I-VI-129}Now $\underline{M}_{X_0} \to \underline{M}_{X_0,(\widehat{z}_1)} \times \ldots \times \underline{M}_{X_0,(\widehat{z}_m)}$ is smooth by and therefore $R_1 \hat{\otimes}_{\Lambda} \ldots \hat{\otimes}_{\Lambda} R_m \to R$ is formally smooth by 4.13, which implies that $t_1,\ldots,t_m$ is extended to a system of free formal generators of $R$ over $\Lambda$ .


\SGApage{I}{129}

\SGAvar{(4)  The projection morphism  $\mathfrak{X}_{(\widehat{z}_i)} \overset{\sim}{\rightarrow} \mathfrak{Y}_i \times_{\mathrm{Spf}(R_i)} \mathrm{Spf}(R) \rightarrow \mathfrak{Y}_i$  is formally smooth since  $\mathrm{Spf}(R) \overset{\ \mathrm{Spf}(\varphi_i)\ }{\relbar\joinrel\relbar\joinrel\relbar\joinrel\relbar} \mathrm{Spf}(R_i)$  is formally smooth. On the other hand, $\mathfrak{Y}_i \rightarrow \mathrm{Spec}(\Lambda)$  is formally smooth by 6.2(3)  and therefore  $\mathfrak{X}_{(\widehat{z}_i)} \rightarrow \mathrm{Spec}(\Lambda)$  is formally smooth. Consequently  $\mathfrak{X} \rightarrow \mathrm{Spec}(\Lambda)$  is formally smooth since  $\{z_1,\dots,z_m\}$  is the set of non-smooth points of  $X_0$ .

}{(4)  The projection morphism  $\mathfrak{X}_{(\widehat{z}_i)} \overset{\sim}{\rightarrow} \mathfrak{Y}_i \times_{\mathrm{Spf}(R_i)} \mathrm{Spf}(R) \rightarrow \mathfrak{Y}_i$  is formally smooth since  $\mathrm{Spf}(R) \xrightarrow{\ \mathrm{Spf}(\varphi_i)\ } \mathrm{Spf}(R_i)$  is formally smooth. On the other hand, $\mathfrak{Y}_i \rightarrow \mathrm{Spec}(\Lambda)$  is formally smooth by 6.2(3)  and therefore  $\mathfrak{X}_{(\widehat{z}_i)} \rightarrow \mathrm{Spec}(\Lambda)$  is formally smooth. Consequently  $\mathfrak{X} \rightarrow \mathrm{Spec}(\Lambda)$  is formally smooth since  $\{z_1,\dots,z_m\}$  is the set of non-smooth points of  $X_0$ .

}{F7I-VI-130}(5)  Since  $R$  is formally smooth over  $\Lambda$  and

\begin{equation*}
0 \rightarrow H^1(X_0,\ \underline{D}^0_{X_0})\ \rightarrow [\underline{M}_{X_0}(k[\varepsilon])] \rightarrow H^0(X_0,\ \underline{D}^1_{X_0}) \rightarrow 0
\end{equation*}
is exact by 4.4 and 4.11, we have

\begin{equation*}
\begin{aligned}
\dim \underline{M}_{X_0} \ &= \ \dim k \otimes R \ = \ \dim_k H^1(X_0,\ \underline{D}^0_{X_0}) + \dim H^0(X_0,\ \underline{D}^1_{X_0})\\
&= \ \dim_k H^1(X_0,\ \underline{D}^0_{X_0}) + m \ .
\end{aligned}
\end{equation*}
Then it suffices to show that

\begin{equation*}
\dim_k H^1(X_0,\ \underline{D}^0_{X_0}) + m = 3g - 3 + \dim_k H^0(X_0,\ \underline{D}^0_{X_0})
\end{equation*}
i.e.,  $\chi(\underline{D}^0_{X_0}) = 3 - 3g + m$

\SGAvar{where we set  $\chi(\underline{E}) = \dim_k H^0(X_0,\ \underline{E}) - \dim_k H^1(X_0,\ \underline{E})$  for any coherent  $X_0$-module  $\underline{E}$ . Since  $\lambda(\underline{D}^0_{X_0})$  does not change when one replaces the base field  $k$  by its algebraic closure, we may assume that  $k$  is algebraically closed. Set  $A = \underline{O}_{X_0,z}$ , where  $z$  is a non-smooth point of  $X_0$ , and let  $A'$  be its normalization. By hypothesis, we have  $A \simeq k[[x,y]]/(f)$  where  $f$  is a non-degenerate quadratic form. Since  $k$  is algebraically closed, we may assume that  $f = xy$ , i.e., $A \simeq k[[x,y]]/(xy)$ . From this we conclude immediately that  $\ell(A'/A) = 1$  and  $\mathrm{Der}_k(A,\ \underline{m}_A) = \mathrm{Der}_k(A,\ A)$ . Let  $C$  be the conductor of  $A$  in  $A'$ . Since $A$ is a complete-intersection we must have that  $\ell(A/C) = \ell(A'/A) = 1$  so that  $C = \underline{m}_A$ . We thus obtain that  $A' = \{a' \in K |\ a'\underline{m}_A \subset \underline{m}_A\}$  where  $K$  is the fraction field of  $A$ . We now claim that  $\mathrm{Der}_k(A,\ A) \subset \mathrm{Der}_k(A',\ A')$  as $A$-modules contained in  $\mathrm{Der}_k(K,\ K)$ . Indeed, let  $X \in \mathrm{Der}_k(A,\ A)$ . Then for any  $a' \in A'$  and
\SGApage{I}{130}
$x \in \underline{m}_A$  we have  $xa' \in \underline{m}_A$  so that  $X(xa') \in \underline{m}_A$  and  $a'X(x) \in \underline{m}_A$  since  $\mathrm{Der}_k(A,\ \underline{m}_A) = \mathrm{Der}_k(A,\ A)$ . Consequently  $X(xa') = a'X(x) + xX(a')$  entails that  $xX(a') \in \underline{m}_A$  for all  $x \in \underline{m}_A$ , and therefore  $X(a') \in A'$ . Furthermore, the exact sequence

}{where we set  $\chi(\underline{E}) = \dim_k H^0(X_0,\ \underline{E}) - \dim_k H^1(X_0,\ \underline{E})$  for any coherent  $X_0$-module  $\underline{E}$ . Since  $\chi(\underline{D}^0_{X_0})$  does not change when one replaces the base field  $k$  by its algebraic closure, we may assume that  $k$  is algebraically closed. Set  $A = \underline{O}_{X_0,z}$ , where  $z$  is a non-smooth point of  $X_0$ , and let  $A'$  be its normalization. By hypothesis, we have  $A \simeq k[[x,y]]/(f)$  where  $f$  is a non-degenerate quadratic form. Since  $k$  is algebraically closed, we may assume that  $f = xy$ , i.e., $A \simeq k[[x,y]]/(xy)$ . From this we conclude immediately that  $\ell(A'/A) = 1$  and  $\mathrm{Der}_k(A,\ \underline{m}_A) = \mathrm{Der}_k(A,\ A)$ . Let  $C$  be the conductor of  $A$  in  $A'$ . Since $A$ is a complete-intersection we must have that  $\ell(A/C) = \ell(A'/A) = 1$  so that  $C = \underline{m}_A$ . We thus obtain that  $A' = \{a' \in K |\ a'\underline{m}_A \subset \underline{m}_A\}$  where  $K$  is the fraction field of  $A$ . We now claim that  $\mathrm{Der}_k(A,\ A) \subset \mathrm{Der}_k(A',\ A')$  as $A$-modules contained in  $\mathrm{Der}_k(K,\ K)$ . Indeed, let  $X \in \mathrm{Der}_k(A,\ A)$ . Then for any  $a' \in A'$  and
\SGApage{I}{130}
$x \in \underline{m}_A$  we have  $xa' \in \underline{m}_A$  so that  $X(xa') \in \underline{m}_A$  and  $a'X(x) \in \underline{m}_A$  since  $\mathrm{Der}_k(A,\ \underline{m}_A) = \mathrm{Der}_k(A,\ A)$ . Consequently  $X(xa') = a'X(x) + xX(a')$  entails that  $xX(a') \in \underline{m}_A$  for all  $x \in \underline{m}_A$ , and therefore  $X(a') \in A'$ . Furthermore, the exact sequence

}{F7I-VI-131}\begin{equation*}
0 \ \rightarrow\ \mathrm{Der}_k(A,\ A) \rightarrow \mathrm{Der}_k(A',\ A') \rightarrow A'/\underline{m}_A A' \rightarrow\ 0
\end{equation*}
entails the exact sequence

\begin{equation*}
0 \ \rightarrow\ \mathrm{Der}_k(A,\ A) \rightarrow \mathrm{Der}_k(A',\ A') \rightarrow A'/C \rightarrow 0
\end{equation*}
Therefore if we set  $X'_0$  to be the normalization of  $X_0$ , then

\begin{equation*}
0 \rightarrow \underline{D}^0_{X_0} \rightarrow \underline{D}^0_{X'_0}\ \rightarrow \underline{O}_{X'_0}/C \rightarrow 0
\end{equation*}
is exact. Since  $X_0$  is locally a complete-intersection, we have  $\ell(\underline{O}_{X'_0}/C) = 2\ell(\underline{O}_{X'_0}/\underline{O}_{X_0}) = 2m$ , whereas  $\chi(\underline{D}^0_{X'_0}) = 3 - 3(g-m)$  by Riemann-Roch theorem. Consequently

\SGAvar{\begin{equation*}
\chi(\underline{D}^0_{X^0_0}) = \chi(\underline{D}^0_{X'_0}) - \ell(\underline{O}_{X'_0}/C) = 3 - 3(g-m) - 2m = 3 - 3g + m \ .
\end{equation*}
This completes our proof.

}{\begin{equation*}
\chi(\underline{D}^0_{X_0}) = \chi(\underline{D}^0_{X'_0}) - \ell(\underline{O}_{X'_0}/C) = 3 - 3(g-m) - 2m = 3 - 3g + m \ .
\end{equation*}
This completes our proof.

}{F7I-VI-132}\underline{Remark} 6.5.  In the formula 6.4(5), the number  $a = H^0(X_0,\ \underline{D}^0_{X_0})$  can be easily computed in terms of Riemann-Roch theorem. Thus

\begin{equation*}
a = \left\{ \begin{array}{lll} 0 & \text{if} & g \geq 2 \\[4pt] 1 & \text{if} & g = 1 \\[4pt] 3 & \text{if} & g = 0 \end{array} \right.
\end{equation*}
Finally we add a remark on characteristic 2. If the characteristic of the base field  $k$  is  $2$ , there is no non-degenerate quadratic form on  $n$  variables if  $n$  is odd. In other words there does not exist a non-degenerate quadratic singularity of even dimension if  $\mathrm{char}\, k = 2$ . In this case one may consider a quadratic form of maximum non-degeneracy: 


\SGApage{I}{131}

\underline{Definition} 6.6.  \underline{Let}  $X$  \underline{be a} k-\underline{scheme of finite type.} \underline{A closed point}  $z \in X$  \underline{is called an ordinary quadratic singularity rational over}  $k$  \underline{if the following holds:}

\SGAvar{(i)  \underline{If either}  $char k \neq 2$  \underline{or}  $char k = 2$  \underline{and}  $dim O_{X,z}$  \underline{is odd, then} $a$  \underline{is a non-degenerate quadratic singularity rational over}  $k$  .

}{(i)  \underline{If either}  $char k \neq 2$  \underline{or}  $char k = 2$  \underline{and}  $dim O_{X,z}$  \underline{is odd, then} $z$  \underline{is a non-degenerate quadratic singularity rational over}  $k$  .

}{F7I-VI-133}(ii)  \underline{If}  $char k = 2$  \underline{and}  $dim O_{X,z}$  \underline{is even, then} $\bar{k} \underset{k}{\otimes} \hat{O}_{X,z} \simeq \bar{k}[[X_0, X_1, \ldots, X_{2n}]]/(f)$  \underline{with}

\begin{equation*}
f = X_0^2 + X_1X_2 + X_3X_4 + \ldots + X_{2n-1}X_{2n}
\end{equation*}
\underline{where}  $\bar{k}$  \underline{is the algebraic closure of}  $k$  .

\underline{Lemma} 6.7.  \underline{Let}  $k$  \underline{be a field of characteristic}  $2$  \underline{and}  \underline{set} $A_0 = k[[X_0, X_1, \ldots ,X_{2n}]]/(f)$  \underline{where}

\begin{equation*}
f = X_0^2 + X_1X_2 + X_3X_4 + \ldots + X_{2n-1}X_{2n} \quad ,
\end{equation*}
\underline{and let}  $(R, \mathfrak{X})$  \underline{be a formal versal deformation of}  $X_0 = \mathrm{Spec}(A_0)$  . \underline{Then} $R = \Lambda[[t, t_1]]$  , \underline{where}  $t, t_1$  \underline{are free formal generators of}  $R$  \underline{over}  $\Lambda$  , \underline{such} \underline{that}  $I_R(\mathfrak{X}|R) = (t^2)$  , \underline{and}  $\mathfrak{X} \to \mathrm{Spec}(\Lambda)$  \underline{is formally smooth.}

\underline{Proof:}  We set  $\mathfrak{X} = \mathrm{Spf}(A)$  . Since

\begin{equation*}
(f, \frac{\partial f}{\partial X_0} , \ldots , \frac{\partial f}{\partial X_{2n}} ) = (X_0^2, X_1, \ldots , X_{2n}) \quad ,
\end{equation*}
it follows that  $D^1(A_0|k, A_0) \simeq  k[X_0]/(X_0^2)$  and therefore  $[\underline{M}_{X_0}(k[\epsilon])]$  is 2-dimensional vector space over  $k$  . Furthermore, it follows from 4.14 that

\begin{equation*}
R \simeq \Lambda[[t, t_1]] \quad \text{and} \quad A \simeq R[[X_0, X_1, \ldots ,X_{2n}]]/(F)
\end{equation*}
where  $F = f + t_1 + tX_0$  . Now the exact sequence

\begin{center}
\begin{tikzcd}[column sep=large]
0 \arrow[r] \& A \arrow[r, "dF"] \& AdX_0 \oplus \ldots \oplus AdX_{2n} \arrow[r] \& \hat{\Omega}_{A|R} \arrow[r] \& 0
\end{tikzcd}
\end{center}










where

\SGAvar{\begin{equation*}
dF = 1 \to \sum_{0 \leq i \leq 2n} \frac{\partial f}{\partial X_i} dX_i + dX_0 \quad , 
\end{equation*}

}{\begin{equation*}
dF = 1 \to \sum_{0 \leq i \leq 2n} \frac{\partial f}{\partial X_i} dX_i + t\,dX_0 \quad , 
\end{equation*}

}{F7I-VI-134}\SGApage{I}{132}

 entails that

\begin{equation*}
I_A(A|R) = (t, x_1, x_2, \ldots, x_{2n}) = (x_0^2 + t_1, t, x_1, \ldots, x_{2n})
\end{equation*}
so that

\begin{equation*}
A/I_A(A|R) \simeq A/(X_0^2 + t_1, t, x_1, \ldots, x_{2n}) \simeq \Lambda[[t_1, X_0]]/(X_0^2 - t_1)  \ .
\end{equation*}
Since  $\Lambda[[t_1, X_0]]/(X_0^2 - t_1) \simeq R/(t) \oplus R/(t)$  as R-modules, it follows that $I_R(A|R) = (t^2)$  , which proves our first assertion. On the other hand,

\begin{equation*}
\begin{aligned} A &\simeq \Lambda[[t, t_1, x_o, x_1, \ldots, x_{2n}]]/(X_0^2 + X_1X_2 + \ldots + X_{2n-1}X_{2n} + tX_0 + t_1) \\ &\simeq \Lambda[[t, X_0, X_1,\ldots,X_{2n}]] \end{aligned}
\end{equation*}
and therefore  $A$  is formally smooth over  $\Lambda$  .

The above statement 6.7 yields a similar result as 6.4, proof of which we leave to the readers.

\underline{Corollary} 6.8.  \underline{Let}  $X_0$  \underline{be a} k-\underline{scheme of finite type with isolated non-smooth points only such that}  $H^2(X_0, \underline{D}^0_{X_0}) = 0$  , \underline{and let}  $(R, \mathfrak{X})$  \underline{be a formal versal deformation of}  $X_0$  . \underline{Denote by}  $\{z_1, \ldots, z_m\}$  \underline{the set of non-smooth points in} $X_0$  , \underline{and assume that}  $char k = 2$  , $dim X_0$  \underline{is even, and all}  $z_i$'s  \underline{are ordinary quadratic singularities. Then}

(1)  $R$  \underline{is formally smooth over}  $\Lambda$

(2)  $I_R(\mathfrak{X}_{(z_i)}|R)$  \underline{is a principal ideal and in turn so is}

\SGAvar{$I_R(\mathfrak{X}|R) = \mathop{\hphantom{\bigcap}}_{1 \leq i \leq m} I_R(\mathfrak{X}_{(z_i)}|R)$  . \underline{Furthermore,}  $I_R(\mathfrak{X}_{(z_i)}|R) = (t_i^2)$  \underline{and} $t_1, \ldots , t_m$  \underline{can be extended to a system of free formal generators of}  $R$  \underline{over} $\Lambda$  .

}{$I_R(\mathfrak{X}|R) = \prod_{1 \leq i \leq m} I_R(\mathfrak{X}_{(z_i)}|R)$  . \underline{Furthermore,}  $I_R(\mathfrak{X}_{(z_i)}|R) = (t_i^2)$  \underline{and} $t_1, \ldots , t_m$  \underline{can be extended to a system of free formal generators of}  $R$  \underline{over} $\Lambda$  .

}{F7I-VI-135}(3)  $\mathfrak{X} \to \mathrm{Spec}(\Lambda)$  \underline{is formally smooth}

(4)  \underline{If}  $k$  \underline{is algebraically closed and}  $X_0$  \underline{is a complete irreducible curve, then}  $dim \underline{M}_{X_0} = 3g - 3 + a$  \underline{where}  $g = dim_k H^1(X_0, \underline{O}_{X_0})$  \underline{and} $a = dim_k H^0(X_0, \underline{D}^0_{X_0})$  .
\markboth{Exposé VI}{Exposé VI}

\ifnum\SGAform=1
\clearpage\section*{Critical apparatus}\begingroup\small\raggedright\sloppy
Readings are given in the exact TeX notation of the two English witnesses. Decimal suffixes identify the individual differences in each marked passage.

\par\medskip\noindent\hypertarget{app-F7I-VI-001.1}{}\textbf{F7I-VI-001.1}\quad printed folio 32.\par
\noindent Printed reading: {\ttfamily\footnotesize composit\allowbreak{}}\par
\noindent Corrected reading: {\ttfamily\footnotesize composite\allowbreak{}}\par
\noindent Source: {\ttfamily\footnotesize 7I:53\allowbreak{}}\par

\par\medskip\noindent\hypertarget{app-F7I-VI-001.2}{}\textbf{F7I-VI-001.2}\quad printed folio 32.\par
\noindent Printed reading: {\ttfamily\footnotesize morphism\allowbreak{}}\par
\noindent Corrected reading: {\ttfamily\footnotesize morphisms\allowbreak{}}\par
\noindent Source: {\ttfamily\footnotesize 7I:53\allowbreak{}}\par

\par\medskip\noindent\hypertarget{app-F7I-VI-002.1}{}\textbf{F7I-VI-002.1}\quad printed folio 33.\par
\noindent Printed reading: {\ttfamily\footnotesize on\allowbreak{}}\par
\noindent Corrected reading: {\ttfamily\footnotesize in\allowbreak{}}\par
\noindent Source: {\ttfamily\footnotesize 7I:55\allowbreak{}}\par

\par\medskip\noindent\hypertarget{app-F7I-VI-003.1}{}\textbf{F7I-VI-003.1}\quad printed folio 34.\par
\noindent Printed reading: {\ttfamily\footnotesize .\allowbreak{}}\par
\noindent Corrected reading: {\ttfamily\footnotesize ,\allowbreak{}}\par
\noindent Source: {\ttfamily\footnotesize 7I:56\allowbreak{}}\par

\par\medskip\noindent\hypertarget{app-F7I-VI-003.2}{}\textbf{F7I-VI-003.2}\quad printed folio 34.\par
\noindent Printed reading: {\ttfamily\footnotesize Then\allowbreak{}}\par
\noindent Corrected reading: {\ttfamily\footnotesize then\allowbreak{}}\par
\noindent Source: {\ttfamily\footnotesize 7I:56\allowbreak{}}\par

\par\medskip\noindent\hypertarget{app-F7I-VI-004.1}{}\textbf{F7I-VI-004.1}\quad printed folio 35.\par
\noindent Printed reading: \emph{(empty)}\par
\noindent Corrected reading: {\ttfamily\footnotesize [\allowbreak{}}\par
\noindent Source: {\ttfamily\footnotesize 7I:57\allowbreak{}}\par

\par\medskip\noindent\hypertarget{app-F7I-VI-005.1}{}\textbf{F7I-VI-005.1}\quad printed folio 37.\par
\noindent Printed reading: \emph{(empty)}\par
\noindent Corrected reading: {\ttfamily\footnotesize \_1\allowbreak{}}\par
\noindent Source: {\ttfamily\footnotesize 7I:60\allowbreak{}}\par

\par\medskip\noindent\hypertarget{app-F7I-VI-005.2}{}\textbf{F7I-VI-005.2}\quad printed folio 37.\par
\noindent Printed reading: \emph{(empty)}\par
\noindent Corrected reading: {\ttfamily\footnotesize '\_1\allowbreak{}}\par
\noindent Source: {\ttfamily\footnotesize 7I:61\allowbreak{}}\par

\par\medskip\noindent\hypertarget{app-F7I-VI-005.3}{}\textbf{F7I-VI-005.3}\quad printed folio 37.\par
\noindent Printed reading: {\ttfamily\footnotesize h\allowbreak{}}\par
\noindent Corrected reading: {\ttfamily\footnotesize k\allowbreak{}}\par
\noindent Source: {\ttfamily\footnotesize FIX-7I-D-0025\allowbreak{}}\par

\par\medskip\noindent\hypertarget{app-F7I-VI-005.4}{}\textbf{F7I-VI-005.4}\quad printed folio 37.\par
\noindent Printed reading: {\ttfamily\footnotesize f\allowbreak{}}\par
\noindent Corrected reading: {\ttfamily\footnotesize h\allowbreak{}}\par
\noindent Source: {\ttfamily\footnotesize 7I:59\allowbreak{}}\par

\par\medskip\noindent\hypertarget{app-F7I-VI-005.5}{}\textbf{F7I-VI-005.5}\quad printed folio 37.\par
\noindent Printed reading: {\ttfamily\footnotesize h\allowbreak{}}\par
\noindent Corrected reading: {\ttfamily\footnotesize R\allowbreak{}}\par
\noindent Source: {\ttfamily\footnotesize 7I:62\allowbreak{}}\par

\par\medskip\noindent\hypertarget{app-F7I-VI-005.6}{}\textbf{F7I-VI-005.6}\quad printed folio 37.\par
\noindent Printed reading: {\ttfamily\footnotesize \textbackslash{}\allowbreak{} }\par
\noindent Corrected reading: {\ttfamily\footnotesize  \textbackslash{}times\allowbreak{} }\par
\noindent Source: {\ttfamily\footnotesize 7I:63\allowbreak{}}\par

\par\medskip\noindent\hypertarget{app-F7I-VI-005.7}{}\textbf{F7I-VI-005.7}\quad printed folio 37.\par
\noindent Printed reading: \emph{(empty)}\par
\noindent Corrected reading: {\ttfamily\footnotesize '\_1\allowbreak{}}\par
\noindent Source: {\ttfamily\footnotesize 7I:64\allowbreak{}}\par

\par\medskip\noindent\hypertarget{app-F7I-VI-006.1}{}\textbf{F7I-VI-006.1}\quad printed folio 38.\par
\noindent Printed reading: {\ttfamily\footnotesize \textbackslash{}overset\allowbreak{}}\par
\noindent Corrected reading: {\ttfamily\footnotesize \textbackslash{}xrightarrow\allowbreak{}}\par
\noindent Source: {\ttfamily\footnotesize FIX-7I-D-0027\allowbreak{}}\par

\par\medskip\noindent\hypertarget{app-F7I-VI-006.2}{}\textbf{F7I-VI-006.2}\quad printed folio 38.\par
\noindent Printed reading: {\ttfamily\footnotesize \}\{\textbackslash{}relbar\textbackslash{}joinrel\textbackslash{}re\allowbreak{}lbar\textbackslash{}joinrel\textbackslash{}relbar\allowbreak{}}\par
\noindent Corrected reading: \emph{(empty)}\par
\noindent Source: {\ttfamily\footnotesize WITNESS-F7I-VI-006.2\allowbreak{}}\par

\par\medskip\noindent\hypertarget{app-F7I-VI-007.1}{}\textbf{F7I-VI-007.1}\quad printed folio 38.\par
\noindent Printed reading: \emph{(empty)}\par
\noindent Corrected reading: {\ttfamily\footnotesize  ,\allowbreak{}}\par
\noindent Source: {\ttfamily\footnotesize 7I:66\allowbreak{}}\par

\par\medskip\noindent\hypertarget{app-F7I-VI-008.1}{}\textbf{F7I-VI-008.1}\quad printed folio 39.\par
\noindent Printed reading: {\ttfamily\footnotesize definiton\allowbreak{}}\par
\noindent Corrected reading: {\ttfamily\footnotesize definition\allowbreak{}}\par
\noindent Source: {\ttfamily\footnotesize 7I:68\allowbreak{}}\par

\par\medskip\noindent\hypertarget{app-F7I-VI-008.2}{}\textbf{F7I-VI-008.2}\quad printed folio 39.\par
\noindent Printed reading: {\ttfamily\footnotesize cases\allowbreak{}}\par
\noindent Corrected reading: {\ttfamily\footnotesize case\allowbreak{}}\par
\noindent Source: {\ttfamily\footnotesize 7I:69\allowbreak{}}\par

\par\medskip\noindent\hypertarget{app-F7I-VI-008.3}{}\textbf{F7I-VI-008.3}\quad printed folio 39.\par
\noindent Printed reading: {\ttfamily\footnotesize ,\allowbreak{} no\allowbreak{} head\allowbreak{}}\par
\noindent Corrected reading: \emph{(empty)}\par
\noindent Source: {\ttfamily\footnotesize FIX-7I-D-0029\allowbreak{}}\par

\par\medskip\noindent\hypertarget{app-F7I-VI-009.1}{}\textbf{F7I-VI-009.1}\quad printed folio 41.\par
\noindent Printed reading: \emph{(empty)}\par
\noindent Corrected reading: {\ttfamily\footnotesize \textbackslash{}underline\{\allowbreak{}}\par
\noindent Source: {\ttfamily\footnotesize FIX-7I-D-0031\allowbreak{}}\par

\par\medskip\noindent\hypertarget{app-F7I-VI-009.2}{}\textbf{F7I-VI-009.2}\quad printed folio 41.\par
\noindent Printed reading: \emph{(empty)}\par
\noindent Corrected reading: {\ttfamily\footnotesize \}\allowbreak{}}\par
\noindent Source: {\ttfamily\footnotesize FIX-7I-D-0031\allowbreak{}}\par

\par\medskip\noindent\hypertarget{app-F7I-VI-009.3}{}\textbf{F7I-VI-009.3}\quad printed folio 41.\par
\noindent Printed reading: \emph{(empty)}\par
\noindent Corrected reading: {\ttfamily\footnotesize \textbackslash{}underline\{\allowbreak{}}\par
\noindent Source: {\ttfamily\footnotesize FIX-7I-D-0032\allowbreak{}}\par

\par\medskip\noindent\hypertarget{app-F7I-VI-009.4}{}\textbf{F7I-VI-009.4}\quad printed folio 41.\par
\noindent Printed reading: \emph{(empty)}\par
\noindent Corrected reading: {\ttfamily\footnotesize \}\allowbreak{}}\par
\noindent Source: {\ttfamily\footnotesize FIX-7I-D-0032\allowbreak{}}\par

\par\medskip\noindent\hypertarget{app-F7I-VI-009.5}{}\textbf{F7I-VI-009.5}\quad printed folio 41.\par
\noindent Printed reading: {\ttfamily\footnotesize exists\allowbreak{}}\par
\noindent Corrected reading: {\ttfamily\footnotesize exist\allowbreak{}}\par
\noindent Source: {\ttfamily\footnotesize 7I:75\allowbreak{}}\par

\par\medskip\noindent\hypertarget{app-F7I-VI-010.1}{}\textbf{F7I-VI-010.1}\quad printed folio 42.\par
\noindent Printed reading: {\ttfamily\footnotesize composit\allowbreak{}}\par
\noindent Corrected reading: {\ttfamily\footnotesize composite\allowbreak{}}\par
\noindent Source: {\ttfamily\footnotesize 7I:77\allowbreak{}}\par

\par\medskip\noindent\hypertarget{app-F7I-VI-011.1}{}\textbf{F7I-VI-011.1}\quad printed folio 42.\par
\noindent Printed reading: \emph{(empty)}\par
\noindent Corrected reading: {\ttfamily\footnotesize \_k\allowbreak{} }\par
\noindent Source: {\ttfamily\footnotesize ADOPT-F7I-VI-011.1\allowbreak{}}\par

\par\medskip\noindent\hypertarget{app-F7I-VI-012.1}{}\textbf{F7I-VI-012.1}\quad printed folio 43.\par
\noindent Printed reading: \emph{(empty)}\par
\noindent Corrected reading: {\ttfamily\footnotesize \_k\allowbreak{}}\par
\noindent Source: {\ttfamily\footnotesize ADOPT-F7I-VI-012.1\allowbreak{}}\par

\par\medskip\noindent\hypertarget{app-F7I-VI-012.2}{}\textbf{F7I-VI-012.2}\quad printed folio 43.\par
\noindent Printed reading: \emph{(empty)}\par
\noindent Corrected reading: {\ttfamily\footnotesize \_1\allowbreak{}}\par
\noindent Source: {\ttfamily\footnotesize 7I:78\allowbreak{}}\par

\par\medskip\noindent\hypertarget{app-F7I-VI-012.3}{}\textbf{F7I-VI-012.3}\quad printed folio 43.\par
\noindent Printed reading: {\ttfamily\footnotesize ,\allowbreak{} }\par
\noindent Corrected reading: \emph{(empty)}\par
\noindent Source: {\ttfamily\footnotesize 7I:79\allowbreak{}}\par

\par\medskip\noindent\hypertarget{app-F7I-VI-012.4}{}\textbf{F7I-VI-012.4}\quad printed folio 43.\par
\noindent Printed reading: \emph{(empty)}\par
\noindent Corrected reading: {\ttfamily\footnotesize )\allowbreak{}}\par
\noindent Source: {\ttfamily\footnotesize 7I:80\allowbreak{}}\par

\par\medskip\noindent\hypertarget{app-F7I-VI-013.1}{}\textbf{F7I-VI-013.1}\quad printed folio 43.\par
\noindent Printed reading: {\ttfamily\footnotesize composit\allowbreak{}}\par
\noindent Corrected reading: {\ttfamily\footnotesize composite\allowbreak{}}\par
\noindent Source: {\ttfamily\footnotesize 7I:81\allowbreak{}}\par

\par\medskip\noindent\hypertarget{app-F7I-VI-014.1}{}\textbf{F7I-VI-014.1}\quad printed folio 44.\par
\noindent Printed reading: {\ttfamily\footnotesize sequences\allowbreak{}}\par
\noindent Corrected reading: {\ttfamily\footnotesize sequence\allowbreak{}}\par
\noindent Source: {\ttfamily\footnotesize 7I:82\allowbreak{}}\par

\par\medskip\noindent\hypertarget{app-F7I-VI-015.1}{}\textbf{F7I-VI-015.1}\quad printed folio 45.\par
\noindent Printed reading: {\ttfamily\footnotesize biefly\allowbreak{}}\par
\noindent Corrected reading: {\ttfamily\footnotesize briefly\allowbreak{}}\par
\noindent Source: {\ttfamily\footnotesize 7I:83\allowbreak{}}\par

\par\medskip\noindent\hypertarget{app-F7I-VI-015.2}{}\textbf{F7I-VI-015.2}\quad printed folio 45.\par
\noindent Printed reading: {\ttfamily\footnotesize preceeding\allowbreak{}}\par
\noindent Corrected reading: {\ttfamily\footnotesize preceding\allowbreak{}}\par
\noindent Source: {\ttfamily\footnotesize 7I:84\allowbreak{}}\par

\par\medskip\noindent\hypertarget{app-F7I-VI-016.1}{}\textbf{F7I-VI-016.1}\quad printed folio 45.\par
\noindent Printed reading: {\ttfamily\footnotesize k\allowbreak{}}\par
\noindent Corrected reading: {\ttfamily\footnotesize K\allowbreak{}}\par
\noindent Source: {\ttfamily\footnotesize 7I:85\allowbreak{}}\par

\par\medskip\noindent\hypertarget{app-F7I-VI-017.1}{}\textbf{F7I-VI-017.1}\quad printed folio 46.\par
\noindent Printed reading: \emph{(empty)}\par
\noindent Corrected reading: {\ttfamily\footnotesize \textbackslash{},\textbackslash{}varepsilon\allowbreak{}}\par
\noindent Source: {\ttfamily\footnotesize 7I:86\allowbreak{}}\par

\par\medskip\noindent\hypertarget{app-F7I-VI-018.1}{}\textbf{F7I-VI-018.1}\quad printed folio 47.\par
\noindent Printed reading: {\ttfamily\footnotesize k\allowbreak{}}\par
\noindent Corrected reading: {\ttfamily\footnotesize 1\allowbreak{}}\par
\noindent Source: {\ttfamily\footnotesize 7I:87\allowbreak{}}\par

\par\medskip\noindent\hypertarget{app-F7I-VI-019.1}{}\textbf{F7I-VI-019.1}\quad printed folio 47.\par
\noindent Printed reading: \emph{(empty)}\par
\noindent Corrected reading: {\ttfamily\footnotesize '\allowbreak{}}\par
\noindent Source: {\ttfamily\footnotesize 7I:89\allowbreak{}}\par

\par\medskip\noindent\hypertarget{app-F7I-VI-019.2}{}\textbf{F7I-VI-019.2}\quad printed folio 47.\par
\noindent Printed reading: \emph{(empty)}\par
\noindent Corrected reading: {\ttfamily\footnotesize \textbackslash{},\allowbreak{}}\par
\noindent Source: {\ttfamily\footnotesize 7I:89\allowbreak{}}\par

\par\medskip\noindent\hypertarget{app-F7I-VI-020.1}{}\textbf{F7I-VI-020.1}\quad printed folio 48.\par
\noindent Printed reading: {\ttfamily\footnotesize diagramm\allowbreak{}}\par
\noindent Corrected reading: {\ttfamily\footnotesize diagram\allowbreak{}}\par
\noindent Source: {\ttfamily\footnotesize 7I:90\allowbreak{}}\par

\par\medskip\noindent\hypertarget{app-F7I-VI-020.2}{}\textbf{F7I-VI-020.2}\quad printed folio 48.\par
\noindent Printed reading: {\ttfamily\footnotesize ,\allowbreak{}}\par
\noindent Corrected reading: {\ttfamily\footnotesize \textbackslash{},\allowbreak{}}\par
\noindent Source: {\ttfamily\footnotesize 7I:91\allowbreak{}}\par

\par\medskip\noindent\hypertarget{app-F7I-VI-020.3}{}\textbf{F7I-VI-020.3}\quad printed folio 48.\par
\noindent Printed reading: \emph{(empty)}\par
\noindent Corrected reading: {\ttfamily\footnotesize '\allowbreak{}}\par
\noindent Source: {\ttfamily\footnotesize 7I:92\allowbreak{}}\par

\par\medskip\noindent\hypertarget{app-F7I-VI-020.4}{}\textbf{F7I-VI-020.4}\quad printed folio 48.\par
\noindent Printed reading: \emph{(empty)}\par
\noindent Corrected reading: {\ttfamily\footnotesize \textbackslash{},\allowbreak{}}\par
\noindent Source: {\ttfamily\footnotesize 7I:92\allowbreak{}}\par

\par\medskip\noindent\hypertarget{app-F7I-VI-020.5}{}\textbf{F7I-VI-020.5}\quad printed folio 48.\par
\noindent Printed reading: {\ttfamily\footnotesize H\allowbreak{}}\par
\noindent Corrected reading: {\ttfamily\footnotesize A\allowbreak{}}\par
\noindent Source: {\ttfamily\footnotesize 7I:93\allowbreak{}}\par

\par\medskip\noindent\hypertarget{app-F7I-VI-020.6}{}\textbf{F7I-VI-020.6}\quad printed folio 48.\par
\noindent Printed reading: {\ttfamily\footnotesize F\allowbreak{}}\par
\noindent Corrected reading: {\ttfamily\footnotesize G\allowbreak{}}\par
\noindent Source: {\ttfamily\footnotesize 7I:94\allowbreak{}}\par

\par\medskip\noindent\hypertarget{app-F7I-VI-021.1}{}\textbf{F7I-VI-021.1}\quad printed folio 48.\par
\noindent Printed reading: {\ttfamily\footnotesize 2\allowbreak{}}\par
\noindent Corrected reading: {\ttfamily\footnotesize 1\allowbreak{}}\par
\noindent Source: {\ttfamily\footnotesize 7I:96\allowbreak{}}\par

\par\medskip\noindent\hypertarget{app-F7I-VI-022.1}{}\textbf{F7I-VI-022.1}\quad printed folio 49.\par
\noindent Printed reading: {\ttfamily\footnotesize .\allowbreak{}}\par
\noindent Corrected reading: \emph{(1 espace(s), 0 retour(s) de ligne dans le témoin TeX)}\par
\noindent Source: {\ttfamily\footnotesize 7I:98\allowbreak{}}\par

\par\medskip\noindent\hypertarget{app-F7I-VI-023.1}{}\textbf{F7I-VI-023.1}\quad printed folio 51.\par
\noindent Printed reading: {\ttfamily\footnotesize sitting\allowbreak{}}\par
\noindent Corrected reading: {\ttfamily\footnotesize setting\allowbreak{}}\par
\noindent Source: {\ttfamily\footnotesize 7I:99\allowbreak{}}\par

\par\medskip\noindent\hypertarget{app-F7I-VI-024.1}{}\textbf{F7I-VI-024.1}\quad printed folio 51.\par
\noindent Printed reading: {\ttfamily\footnotesize \textbackslash{}varphi\allowbreak{}}\par
\noindent Corrected reading: {\ttfamily\footnotesize \textbackslash{}psi\allowbreak{}}\par
\noindent Source: {\ttfamily\footnotesize 7I:100\allowbreak{}}\par

\par\medskip\noindent\hypertarget{app-F7I-VI-025.1}{}\textbf{F7I-VI-025.1}\quad printed folio 52.\par
\noindent Printed reading: {\ttfamily\footnotesize groupoids\allowbreak{}}\par
\noindent Corrected reading: {\ttfamily\footnotesize groupoid\allowbreak{}}\par
\noindent Source: {\ttfamily\footnotesize 7I:101\allowbreak{}}\par

\par\medskip\noindent\hypertarget{app-F7I-VI-025.2}{}\textbf{F7I-VI-025.2}\quad printed folio 52.\par
\noindent Printed reading: {\ttfamily\footnotesize 13\allowbreak{}}\par
\noindent Corrected reading: {\ttfamily\footnotesize 14\allowbreak{}}\par
\noindent Source: {\ttfamily\footnotesize 7I:102\allowbreak{}}\par

\par\medskip\noindent\hypertarget{app-F7I-VI-026.1}{}\textbf{F7I-VI-026.1}\quad printed folio 53.\par
\noindent Printed reading: \emph{(empty)}\par
\noindent Corrected reading: {\ttfamily\footnotesize \_1\allowbreak{}}\par
\noindent Source: {\ttfamily\footnotesize 7I:104\allowbreak{}}\par

\par\medskip\noindent\hypertarget{app-F7I-VI-027.1}{}\textbf{F7I-VI-027.1}\quad printed folio 56.\par
\noindent Printed reading: \emph{(empty)}\par
\noindent Corrected reading: {\ttfamily\footnotesize )\allowbreak{}}\par
\noindent Source: {\ttfamily\footnotesize 7I:110\allowbreak{}}\par

\par\medskip\noindent\hypertarget{app-F7I-VI-028.1}{}\textbf{F7I-VI-028.1}\quad printed folio 57.\par
\noindent Printed reading: {\ttfamily\footnotesize \textbackslash{}varphi\allowbreak{}}\par
\noindent Corrected reading: {\ttfamily\footnotesize \textbackslash{}psi\allowbreak{}}\par
\noindent Source: {\ttfamily\footnotesize 7I:111\allowbreak{}}\par

\par\medskip\noindent\hypertarget{app-F7I-VI-028.2}{}\textbf{F7I-VI-028.2}\quad printed folio 57.\par
\noindent Printed reading: {\ttfamily\footnotesize \textbackslash{}psi\allowbreak{}}\par
\noindent Corrected reading: {\ttfamily\footnotesize \textbackslash{}varphi\allowbreak{}}\par
\noindent Source: {\ttfamily\footnotesize 7I:111\allowbreak{}}\par

\par\medskip\noindent\hypertarget{app-F7I-VI-029.1}{}\textbf{F7I-VI-029.1}\quad printed folio 57.\par
\noindent Printed reading: {\ttfamily\footnotesize be\allowbreak{}}\par
\noindent Corrected reading: {\ttfamily\footnotesize by\allowbreak{}}\par
\noindent Source: {\ttfamily\footnotesize 7I:112\allowbreak{}}\par

\par\medskip\noindent\hypertarget{app-F7I-VI-030.1}{}\textbf{F7I-VI-030.1}\quad printed folio 58.\par
\noindent Printed reading: \emph{(empty)}\par
\noindent Corrected reading: {\ttfamily\footnotesize '\allowbreak{}}\par
\noindent Source: {\ttfamily\footnotesize WITNESS-F7I-VI-030.1\allowbreak{}}\par

\par\medskip\noindent\hypertarget{app-F7I-VI-031.1}{}\textbf{F7I-VI-031.1}\quad printed folio 59.\par
\noindent Printed reading: \emph{(empty)}\par
\noindent Corrected reading: {\ttfamily\footnotesize )\allowbreak{}}\par
\noindent Source: {\ttfamily\footnotesize 7I:114\allowbreak{}}\par

\par\medskip\noindent\hypertarget{app-F7I-VI-032.1}{}\textbf{F7I-VI-032.1}\quad printed folio 60.\par
\noindent Printed reading: {\ttfamily\footnotesize cocartegory\allowbreak{}}\par
\noindent Corrected reading: {\ttfamily\footnotesize cocategory\allowbreak{}}\par
\noindent Source: {\ttfamily\footnotesize 7I:116\allowbreak{}}\par

\par\medskip\noindent\hypertarget{app-F7I-VI-033.1}{}\textbf{F7I-VI-033.1}\quad printed folio 61.\par
\noindent Printed reading: {\ttfamily\footnotesize =\textbackslash{}!=\allowbreak{}}\par
\noindent Corrected reading: {\ttfamily\footnotesize \textbackslash{}rightrightarrows\allowbreak{}}\par
\noindent Source: {\ttfamily\footnotesize FIX-7I-D-0041\allowbreak{}}\par

\par\medskip\noindent\hypertarget{app-F7I-VI-033.2}{}\textbf{F7I-VI-033.2}\quad printed folio 61.\par
\noindent Printed reading: {\ttfamily\footnotesize \textbackslash{}qquad\allowbreak{}}\par
\noindent Corrected reading: {\ttfamily\footnotesize \textbackslash{}amalg\allowbreak{}}\par
\noindent Source: {\ttfamily\footnotesize 7I:118\allowbreak{}}\par

\par\medskip\noindent\hypertarget{app-F7I-VI-034.1}{}\textbf{F7I-VI-034.1}\quad printed folio 62.\par
\noindent Printed reading: {\ttfamily\footnotesize composit\allowbreak{}}\par
\noindent Corrected reading: {\ttfamily\footnotesize composite\allowbreak{}}\par
\noindent Source: {\ttfamily\footnotesize 7I:120\allowbreak{}}\par

\par\medskip\noindent\hypertarget{app-F7I-VI-034.2}{}\textbf{F7I-VI-034.2}\quad printed folio 62.\par
\noindent Printed reading: \emph{(empty)}\par
\noindent Corrected reading: {\ttfamily\footnotesize  :\allowbreak{}}\par
\noindent Source: {\ttfamily\footnotesize 7I:121\allowbreak{}}\par

\par\medskip\noindent\hypertarget{app-F7I-VI-034.3}{}\textbf{F7I-VI-034.3}\quad printed folio 62.\par
\noindent Printed reading: {\ttfamily\footnotesize an\allowbreak{}}\par
\noindent Corrected reading: {\ttfamily\footnotesize a\allowbreak{}}\par
\noindent Source: {\ttfamily\footnotesize 7I:122\allowbreak{}}\par

\par\medskip\noindent\hypertarget{app-F7I-VI-035.1}{}\textbf{F7I-VI-035.1}\quad printed folio 63.\par
\noindent Printed reading: {\ttfamily\footnotesize or\allowbreak{}}\par
\noindent Corrected reading: {\ttfamily\footnotesize of\allowbreak{}}\par
\noindent Source: {\ttfamily\footnotesize 7I:123\allowbreak{}}\par

\par\medskip\noindent\hypertarget{app-F7I-VI-036.1}{}\textbf{F7I-VI-036.1}\quad printed folio 65.\par
\noindent Printed reading: \emph{(empty)}\par
\noindent Corrected reading: {\ttfamily\footnotesize )\allowbreak{}}\par
\noindent Source: {\ttfamily\footnotesize 7I:124\allowbreak{}}\par

\par\medskip\noindent\hypertarget{app-F7I-VI-036.2}{}\textbf{F7I-VI-036.2}\quad printed folio 65.\par
\noindent Printed reading: {\ttfamily\footnotesize ,\allowbreak{}}\par
\noindent Corrected reading: \emph{(empty)}\par
\noindent Source: {\ttfamily\footnotesize 7I:125\allowbreak{}}\par

\par\medskip\noindent\hypertarget{app-F7I-VI-036.3}{}\textbf{F7I-VI-036.3}\quad printed folio 65.\par
\noindent Printed reading: \emph{(empty)}\par
\noindent Corrected reading: {\ttfamily\footnotesize \textbackslash{}underline\{\allowbreak{}}\par
\noindent Source: {\ttfamily\footnotesize 7I:126\allowbreak{}}\par

\par\medskip\noindent\hypertarget{app-F7I-VI-036.4}{}\textbf{F7I-VI-036.4}\quad printed folio 65.\par
\noindent Printed reading: \emph{(empty)}\par
\noindent Corrected reading: {\ttfamily\footnotesize \}\allowbreak{}}\par
\noindent Source: {\ttfamily\footnotesize WITNESS-F7I-VI-036.4\allowbreak{}}\par

\par\medskip\noindent\hypertarget{app-F7I-VI-036.5}{}\textbf{F7I-VI-036.5}\quad printed folio 65.\par
\noindent Printed reading: \emph{(empty)}\par
\noindent Corrected reading: {\ttfamily\footnotesize \textbackslash{}hat\{\allowbreak{}}\par
\noindent Source: {\ttfamily\footnotesize WITNESS-F7I-VI-036.5\allowbreak{}}\par

\par\medskip\noindent\hypertarget{app-F7I-VI-036.6}{}\textbf{F7I-VI-036.6}\quad printed folio 65.\par
\noindent Printed reading: \emph{(empty)}\par
\noindent Corrected reading: {\ttfamily\footnotesize \}\allowbreak{}}\par
\noindent Source: {\ttfamily\footnotesize WITNESS-F7I-VI-036.6\allowbreak{}}\par

\par\medskip\noindent\hypertarget{app-F7I-VI-036.7}{}\textbf{F7I-VI-036.7}\quad printed folio 65.\par
\noindent Printed reading: {\ttfamily\footnotesize sets\allowbreak{}}\par
\noindent Corrected reading: {\ttfamily\footnotesize set\allowbreak{}}\par
\noindent Source: {\ttfamily\footnotesize WITNESS-F7I-VI-036.7\allowbreak{}}\par

\par\medskip\noindent\hypertarget{app-F7I-VI-036.8}{}\textbf{F7I-VI-036.8}\quad printed folio 65.\par
\noindent Printed reading: \emph{(empty)}\par
\noindent Corrected reading: {\ttfamily\footnotesize \textbackslash{}hat\{\allowbreak{}}\par
\noindent Source: {\ttfamily\footnotesize ADOPT-F7I-VI-036.8\allowbreak{}}\par

\par\medskip\noindent\hypertarget{app-F7I-VI-036.9}{}\textbf{F7I-VI-036.9}\quad printed folio 65.\par
\noindent Printed reading: \emph{(empty)}\par
\noindent Corrected reading: {\ttfamily\footnotesize \}\allowbreak{}}\par
\noindent Source: {\ttfamily\footnotesize ADOPT-F7I-VI-036.9\allowbreak{}}\par

\par\medskip\noindent\hypertarget{app-F7I-VI-036.10}{}\textbf{F7I-VI-036.10}\quad printed folio 65.\par
\noindent Printed reading: {\ttfamily\footnotesize get\allowbreak{}}\par
\noindent Corrected reading: {\ttfamily\footnotesize gets\allowbreak{}}\par
\noindent Source: {\ttfamily\footnotesize 7I:127\allowbreak{}}\par

\par\medskip\noindent\hypertarget{app-F7I-VI-037.1}{}\textbf{F7I-VI-037.1}\quad printed folio 66.\par
\noindent Printed reading: \emph{(empty)}\par
\noindent Corrected reading: {\ttfamily\footnotesize ]\allowbreak{}}\par
\noindent Source: {\ttfamily\footnotesize 7I:128\allowbreak{}}\par

\par\medskip\noindent\hypertarget{app-F7I-VI-038.1}{}\textbf{F7I-VI-038.1}\quad printed folio 67.\par
\noindent Printed reading: {\ttfamily\footnotesize  of\allowbreak{}}\par
\noindent Corrected reading: \emph{(empty)}\par
\noindent Source: {\ttfamily\footnotesize 7I:129\allowbreak{}}\par

\par\medskip\noindent\hypertarget{app-F7I-VI-039.1}{}\textbf{F7I-VI-039.1}\quad printed folio 68.\par
\noindent Printed reading: \emph{(empty)}\par
\noindent Corrected reading: {\ttfamily\footnotesize )\allowbreak{}}\par
\noindent Source: {\ttfamily\footnotesize 7I:131\allowbreak{}}\par

\par\medskip\noindent\hypertarget{app-F7I-VI-039.2}{}\textbf{F7I-VI-039.2}\quad printed folio 68.\par
\noindent Printed reading: {\ttfamily\footnotesize composit\allowbreak{}}\par
\noindent Corrected reading: {\ttfamily\footnotesize composite\allowbreak{}}\par
\noindent Source: {\ttfamily\footnotesize 7I:132\allowbreak{}}\par

\par\medskip\noindent\hypertarget{app-F7I-VI-040.1}{}\textbf{F7I-VI-040.1}\quad printed folio 68.\par
\noindent Printed reading: {\ttfamily\footnotesize composit\allowbreak{}}\par
\noindent Corrected reading: {\ttfamily\footnotesize composite\allowbreak{}}\par
\noindent Source: {\ttfamily\footnotesize 7I:133\allowbreak{}}\par

\par\medskip\noindent\hypertarget{app-F7I-VI-041.1}{}\textbf{F7I-VI-041.1}\quad printed folio 69.\par
\noindent Printed reading: {\ttfamily\footnotesize cogroupoid\allowbreak{}}\par
\noindent Corrected reading: {\ttfamily\footnotesize cofibered\allowbreak{} groupoid\allowbreak{}}\par
\noindent Source: {\ttfamily\footnotesize WITNESS-F7I-VI-041.1\allowbreak{}}\par

\par\medskip\noindent\hypertarget{app-F7I-VI-042.1}{}\textbf{F7I-VI-042.1}\quad printed folio 70.\par
\noindent Printed reading: {\ttfamily\footnotesize its\allowbreak{}}\par
\noindent Corrected reading: {\ttfamily\footnotesize the\allowbreak{}}\par
\noindent Source: {\ttfamily\footnotesize 7I:136\allowbreak{}}\par

\par\medskip\noindent\hypertarget{app-F7I-VI-042.2}{}\textbf{F7I-VI-042.2}\quad printed folio 70.\par
\noindent Printed reading: {\ttfamily\footnotesize 0\allowbreak{}}\par
\noindent Corrected reading: {\ttfamily\footnotesize 1\allowbreak{}}\par
\noindent Source: {\ttfamily\footnotesize 7I:135\allowbreak{}}\par

\par\medskip\noindent\hypertarget{app-F7I-VI-042.3}{}\textbf{F7I-VI-042.3}\quad printed folio 70.\par
\noindent Printed reading: {\ttfamily\footnotesize cogroupoid\allowbreak{}}\par
\noindent Corrected reading: {\ttfamily\footnotesize cogroupoids\allowbreak{}}\par
\noindent Source: {\ttfamily\footnotesize ADOPT-F7I-VI-042.3\allowbreak{}}\par

\par\medskip\noindent\hypertarget{app-F7I-VI-043.1}{}\textbf{F7I-VI-043.1}\quad printed folio 72.\par
\noindent Printed reading: {\ttfamily\footnotesize \textbackslash{}circ\allowbreak{}}\par
\noindent Corrected reading: {\ttfamily\footnotesize 0\allowbreak{}}\par
\noindent Source: {\ttfamily\footnotesize 7I:137\allowbreak{}}\par

\par\medskip\noindent\hypertarget{app-F7I-VI-043.2}{}\textbf{F7I-VI-043.2}\quad printed folio 72.\par
\noindent Printed reading: {\ttfamily\footnotesize \textbackslash{}circ\allowbreak{}}\par
\noindent Corrected reading: {\ttfamily\footnotesize 0\allowbreak{}}\par
\noindent Source: {\ttfamily\footnotesize 7I:138\allowbreak{}}\par

\par\medskip\noindent\hypertarget{app-F7I-VI-043.3}{}\textbf{F7I-VI-043.3}\quad printed folio 72.\par
\noindent Printed reading: {\ttfamily\footnotesize \textbackslash{}circ\allowbreak{}}\par
\noindent Corrected reading: {\ttfamily\footnotesize 0\allowbreak{}}\par
\noindent Source: {\ttfamily\footnotesize 7I:139\allowbreak{}}\par

\par\medskip\noindent\hypertarget{app-F7I-VI-044.1}{}\textbf{F7I-VI-044.1}\quad printed folio 73.\par
\noindent Printed reading: \emph{(empty)}\par
\noindent Corrected reading: {\ttfamily\footnotesize \textbackslash{}overset\{\textbackslash{}vee\}\{\allowbreak{}}\par
\noindent Source: {\ttfamily\footnotesize 7I:140\allowbreak{}}\par

\par\medskip\noindent\hypertarget{app-F7I-VI-044.2}{}\textbf{F7I-VI-044.2}\quad printed folio 73.\par
\noindent Printed reading: \emph{(empty)}\par
\noindent Corrected reading: {\ttfamily\footnotesize \}\allowbreak{}}\par
\noindent Source: {\ttfamily\footnotesize 7I:140\allowbreak{}}\par

\par\medskip\noindent\hypertarget{app-F7I-VI-044.3}{}\textbf{F7I-VI-044.3}\quad printed folio 73.\par
\noindent Printed reading: \emph{(empty)}\par
\noindent Corrected reading: {\ttfamily\footnotesize \textbackslash{}overset\{\textbackslash{}vee\}\{\allowbreak{}}\par
\noindent Source: {\ttfamily\footnotesize 7I:140\allowbreak{}}\par

\par\medskip\noindent\hypertarget{app-F7I-VI-044.4}{}\textbf{F7I-VI-044.4}\quad printed folio 73.\par
\noindent Printed reading: \emph{(empty)}\par
\noindent Corrected reading: {\ttfamily\footnotesize \}\allowbreak{}}\par
\noindent Source: {\ttfamily\footnotesize 7I:140\allowbreak{}}\par

\par\medskip\noindent\hypertarget{app-F7I-VI-045.1}{}\textbf{F7I-VI-045.1}\quad printed folio 73.\par
\noindent Printed reading: \emph{(empty)}\par
\noindent Corrected reading: {\ttfamily\footnotesize \textbackslash{}check\{\allowbreak{}}\par
\noindent Source: {\ttfamily\footnotesize 7I:141\allowbreak{}}\par

\par\medskip\noindent\hypertarget{app-F7I-VI-045.2}{}\textbf{F7I-VI-045.2}\quad printed folio 73.\par
\noindent Printed reading: \emph{(empty)}\par
\noindent Corrected reading: {\ttfamily\footnotesize \}\allowbreak{}}\par
\noindent Source: {\ttfamily\footnotesize 7I:141\allowbreak{}}\par

\par\medskip\noindent\hypertarget{app-F7I-VI-045.3}{}\textbf{F7I-VI-045.3}\quad printed folio 73.\par
\noindent Printed reading: {\ttfamily\footnotesize \textbackslash{}\allowbreak{} }\par
\noindent Corrected reading: {\ttfamily\footnotesize |\allowbreak{}}\par
\noindent Source: {\ttfamily\footnotesize 7I:142\allowbreak{}}\par

\par\medskip\noindent\hypertarget{app-F7I-VI-045.4}{}\textbf{F7I-VI-045.4}\quad printed folio 73.\par
\noindent Printed reading: {\ttfamily\footnotesize 2\allowbreak{}}\par
\noindent Corrected reading: {\ttfamily\footnotesize 3\allowbreak{}}\par
\noindent Source: {\ttfamily\footnotesize 7I:143\allowbreak{}}\par

\par\medskip\noindent\hypertarget{app-F7I-VI-046.1}{}\textbf{F7I-VI-046.1}\quad printed folio 73.\par
\noindent Printed reading: {\ttfamily\footnotesize 2\allowbreak{}}\par
\noindent Corrected reading: {\ttfamily\footnotesize 3\allowbreak{}}\par
\noindent Source: {\ttfamily\footnotesize 7I:144\allowbreak{}}\par

\par\medskip\noindent\hypertarget{app-F7I-VI-047.1}{}\textbf{F7I-VI-047.1}\quad printed folio 74.\par
\noindent Printed reading: {\ttfamily\footnotesize suffice\allowbreak{}}\par
\noindent Corrected reading: {\ttfamily\footnotesize suffices\allowbreak{}}\par
\noindent Source: {\ttfamily\footnotesize 7I:145\allowbreak{}}\par

\par\medskip\noindent\hypertarget{app-F7I-VI-048.1}{}\textbf{F7I-VI-048.1}\quad printed folio 74.\par
\noindent Printed reading: {\ttfamily\footnotesize presheaf\allowbreak{}}\par
\noindent Corrected reading: {\ttfamily\footnotesize presheaves\allowbreak{}}\par
\noindent Source: {\ttfamily\footnotesize 7I:146\allowbreak{}}\par

\par\medskip\noindent\hypertarget{app-F7I-VI-049.1}{}\textbf{F7I-VI-049.1}\quad printed folio 75.\par
\noindent Printed reading: {\ttfamily\footnotesize amits\allowbreak{}}\par
\noindent Corrected reading: {\ttfamily\footnotesize admits\allowbreak{}}\par
\noindent Source: {\ttfamily\footnotesize 7I:147\allowbreak{}}\par

\par\medskip\noindent\hypertarget{app-F7I-VI-050.1}{}\textbf{F7I-VI-050.1}\quad printed folio 75.\par
\noindent Printed reading: {\ttfamily\footnotesize '\allowbreak{}}\par
\noindent Corrected reading: {\ttfamily\footnotesize \textasciicircum{}\{1\}\allowbreak{}}\par
\noindent Source: {\ttfamily\footnotesize 7I:148\allowbreak{}}\par

\par\medskip\noindent\hypertarget{app-F7I-VI-050.2}{}\textbf{F7I-VI-050.2}\quad printed folio 75.\par
\noindent Printed reading: {\ttfamily\footnotesize '\allowbreak{}}\par
\noindent Corrected reading: {\ttfamily\footnotesize \textasciicircum{}\{1\}\allowbreak{}}\par
\noindent Source: {\ttfamily\footnotesize 7I:148\allowbreak{}}\par

\par\medskip\noindent\hypertarget{app-F7I-VI-050.3}{}\textbf{F7I-VI-050.3}\quad printed folio 75.\par
\noindent Printed reading: \emph{(empty)}\par
\noindent Corrected reading: {\ttfamily\footnotesize (P)\allowbreak{}}\par
\noindent Source: {\ttfamily\footnotesize 7I:149\allowbreak{}}\par

\par\medskip\noindent\hypertarget{app-F7I-VI-050.4}{}\textbf{F7I-VI-050.4}\quad printed folio 75.\par
\noindent Printed reading: \emph{(empty)}\par
\noindent Corrected reading: {\ttfamily\footnotesize (P)\allowbreak{}}\par
\noindent Source: {\ttfamily\footnotesize 7I:149\allowbreak{}}\par

\par\medskip\noindent\hypertarget{app-F7I-VI-050.5}{}\textbf{F7I-VI-050.5}\quad printed folio 75.\par
\noindent Printed reading: \emph{(empty)}\par
\noindent Corrected reading: {\ttfamily\footnotesize (P)\allowbreak{}}\par
\noindent Source: {\ttfamily\footnotesize 7I:149\allowbreak{}}\par

\par\medskip\noindent\hypertarget{app-F7I-VI-051.1}{}\textbf{F7I-VI-051.1}\quad printed folio 76.\par
\noindent Printed reading: \emph{(empty)}\par
\noindent Corrected reading: {\ttfamily\footnotesize \textbackslash{}check\{\allowbreak{}}\par
\noindent Source: {\ttfamily\footnotesize 7I:150\allowbreak{}}\par

\par\medskip\noindent\hypertarget{app-F7I-VI-051.2}{}\textbf{F7I-VI-051.2}\quad printed folio 76.\par
\noindent Printed reading: \emph{(empty)}\par
\noindent Corrected reading: {\ttfamily\footnotesize \}\allowbreak{}}\par
\noindent Source: {\ttfamily\footnotesize 7I:150\allowbreak{}}\par

\par\medskip\noindent\hypertarget{app-F7I-VI-052.1}{}\textbf{F7I-VI-052.1}\quad printed folio 76.\par
\noindent Printed reading: {\ttfamily\footnotesize <\allowbreak{}}\par
\noindent Corrected reading: {\ttfamily\footnotesize \textbackslash{}leq\allowbreak{}}\par
\noindent Source: {\ttfamily\footnotesize 7I:151\allowbreak{}}\par

\par\medskip\noindent\hypertarget{app-F7I-VI-053.1}{}\textbf{F7I-VI-053.1}\quad printed folio 76.\par
\noindent Printed reading: {\ttfamily\footnotesize i\allowbreak{}}\par
\noindent Corrected reading: {\ttfamily\footnotesize 1\allowbreak{}}\par
\noindent Source: {\ttfamily\footnotesize FIX-7I-D-0066\allowbreak{}}\par

\par\medskip\noindent\hypertarget{app-F7I-VI-054.1}{}\textbf{F7I-VI-054.1}\quad printed folio 77.\par
\noindent Printed reading: {\ttfamily\footnotesize \textbackslash{}otimes\allowbreak{}}\par
\noindent Corrected reading: {\ttfamily\footnotesize \textbackslash{}times\allowbreak{}}\par
\noindent Source: {\ttfamily\footnotesize FIX-7I-D-0067\allowbreak{}}\par

\par\medskip\noindent\hypertarget{app-F7I-VI-055.1}{}\textbf{F7I-VI-055.1}\quad printed folio 77.\par
\noindent Printed reading: {\ttfamily\footnotesize SG\allowbreak{}}\par
\noindent Corrected reading: {\ttfamily\footnotesize \textbackslash{}check\{S\}G\allowbreak{}}\par
\noindent Source: {\ttfamily\footnotesize 7I:154\allowbreak{}}\par

\par\medskip\noindent\hypertarget{app-F7I-VI-055.2}{}\textbf{F7I-VI-055.2}\quad printed folio 77.\par
\noindent Printed reading: \emph{(empty)}\par
\noindent Corrected reading: {\ttfamily\footnotesize \textbackslash{}underline\{\allowbreak{}}\par
\noindent Source: {\ttfamily\footnotesize 7I:155\allowbreak{}}\par

\par\medskip\noindent\hypertarget{app-F7I-VI-055.3}{}\textbf{F7I-VI-055.3}\quad printed folio 77.\par
\noindent Printed reading: \emph{(empty)}\par
\noindent Corrected reading: {\ttfamily\footnotesize \}\allowbreak{}}\par
\noindent Source: {\ttfamily\footnotesize 7I:155\allowbreak{}}\par

\par\medskip\noindent\hypertarget{app-F7I-VI-055.4}{}\textbf{F7I-VI-055.4}\quad printed folio 77.\par
\noindent Printed reading: \emph{(empty)}\par
\noindent Corrected reading: {\ttfamily\footnotesize \textbackslash{}underline\{\allowbreak{}}\par
\noindent Source: {\ttfamily\footnotesize 7I:155\allowbreak{}}\par

\par\medskip\noindent\hypertarget{app-F7I-VI-055.5}{}\textbf{F7I-VI-055.5}\quad printed folio 77.\par
\noindent Printed reading: \emph{(empty)}\par
\noindent Corrected reading: {\ttfamily\footnotesize \}\}\textasciicircum{}\{\textbackslash{}cdot\allowbreak{}}\par
\noindent Source: {\ttfamily\footnotesize 7I:155\allowbreak{}}\par

\par\medskip\noindent\hypertarget{app-F7I-VI-056.1}{}\textbf{F7I-VI-056.1}\quad printed folio 78.\par
\noindent Printed reading: {\ttfamily\footnotesize  a\allowbreak{}}\par
\noindent Corrected reading: \emph{(empty)}\par
\noindent Source: {\ttfamily\footnotesize 7I:156\allowbreak{}}\par

\par\medskip\noindent\hypertarget{app-F7I-VI-057.1}{}\textbf{F7I-VI-057.1}\quad printed folio 79.\par
\noindent Printed reading: {\ttfamily\footnotesize \textbackslash{}\allowbreak{} \textbackslash{}\allowbreak{} \textbackslash{}\allowbreak{} }\par
\noindent Corrected reading: {\ttfamily\footnotesize  <\allowbreak{} }\par
\noindent Source: {\ttfamily\footnotesize 7I:157\allowbreak{}}\par

\par\medskip\noindent\hypertarget{app-F7I-VI-057.2}{}\textbf{F7I-VI-057.2}\quad printed folio 79.\par
\noindent Printed reading: {\ttfamily\footnotesize \textbackslash{}\allowbreak{} \textbackslash{}\allowbreak{} \textbackslash{}\allowbreak{} }\par
\noindent Corrected reading: {\ttfamily\footnotesize  <\allowbreak{} }\par
\noindent Source: {\ttfamily\footnotesize 7I:157\allowbreak{}}\par

\par\medskip\noindent\hypertarget{app-F7I-VI-057.3}{}\textbf{F7I-VI-057.3}\quad printed folio 79.\par
\noindent Printed reading: {\ttfamily\footnotesize \textbackslash{}phantom\{\allowbreak{}}\par
\noindent Corrected reading: \emph{(empty)}\par
\noindent Source: {\ttfamily\footnotesize 7I:158\allowbreak{}}\par

\par\medskip\noindent\hypertarget{app-F7I-VI-057.4}{}\textbf{F7I-VI-057.4}\quad printed folio 79.\par
\noindent Printed reading: {\ttfamily\footnotesize \}\allowbreak{}}\par
\noindent Corrected reading: \emph{(empty)}\par
\noindent Source: {\ttfamily\footnotesize WITNESS-F7I-VI-057.4\allowbreak{}}\par

\par\medskip\noindent\hypertarget{app-F7I-VI-057.5}{}\textbf{F7I-VI-057.5}\quad printed folio 79.\par
\noindent Printed reading: \emph{(empty)}\par
\noindent Corrected reading: {\ttfamily\footnotesize \textasciicircum{}\{\textbackslash{}cdot\}\allowbreak{}}\par
\noindent Source: {\ttfamily\footnotesize WITNESS-F7I-VI-057.5\allowbreak{}}\par

\par\medskip\noindent\hypertarget{app-F7I-VI-057.6}{}\textbf{F7I-VI-057.6}\quad printed folio 79.\par
\noindent Printed reading: \emph{(empty)}\par
\noindent Corrected reading: {\ttfamily\footnotesize \textbackslash{}widetilde\{\allowbreak{}}\par
\noindent Source: {\ttfamily\footnotesize 7I:160\allowbreak{}}\par

\par\medskip\noindent\hypertarget{app-F7I-VI-057.7}{}\textbf{F7I-VI-057.7}\quad printed folio 79.\par
\noindent Printed reading: \emph{(empty)}\par
\noindent Corrected reading: {\ttfamily\footnotesize \}\allowbreak{}}\par
\noindent Source: {\ttfamily\footnotesize 7I:160\allowbreak{}}\par

\par\medskip\noindent\hypertarget{app-F7I-VI-057.8}{}\textbf{F7I-VI-057.8}\quad printed folio 79.\par
\noindent Printed reading: {\ttfamily\footnotesize 2\allowbreak{}}\par
\noindent Corrected reading: {\ttfamily\footnotesize 3\allowbreak{}}\par
\noindent Source: {\ttfamily\footnotesize 7I:161\allowbreak{}}\par

\par\medskip\noindent\hypertarget{app-F7I-VI-058.1}{}\textbf{F7I-VI-058.1}\quad printed folio 80.\par
\noindent Printed reading: \emph{(empty)}\par
\noindent Corrected reading: {\ttfamily\footnotesize  the\allowbreak{}}\par
\noindent Source: {\ttfamily\footnotesize 7I:162\allowbreak{}}\par

\par\medskip\noindent\hypertarget{app-F7I-VI-059.1}{}\textbf{F7I-VI-059.1}\quad printed folio 80.\par
\noindent Printed reading: {\ttfamily\footnotesize gives\allowbreak{}}\par
\noindent Corrected reading: {\ttfamily\footnotesize give\allowbreak{}}\par
\noindent Source: {\ttfamily\footnotesize 7I:163\allowbreak{}}\par

\par\medskip\noindent\hypertarget{app-F7I-VI-060.1}{}\textbf{F7I-VI-060.1}\quad printed folio 81.\par
\noindent Printed reading: {\ttfamily\footnotesize gives\allowbreak{}}\par
\noindent Corrected reading: {\ttfamily\footnotesize give\allowbreak{}}\par
\noindent Source: {\ttfamily\footnotesize 7I:165\allowbreak{}}\par

\par\medskip\noindent\hypertarget{app-F7I-VI-060.2}{}\textbf{F7I-VI-060.2}\quad printed folio 81.\par
\noindent Printed reading: {\ttfamily\footnotesize \textbackslash{}\allowbreak{} }\par
\noindent Corrected reading: {\ttfamily\footnotesize |\allowbreak{}}\par
\noindent Source: {\ttfamily\footnotesize 7I:164\allowbreak{}}\par

\par\medskip\noindent\hypertarget{app-F7I-VI-060.3}{}\textbf{F7I-VI-060.3}\quad printed folio 81.\par
\noindent Printed reading: {\ttfamily\footnotesize \textbackslash{}\allowbreak{} }\par
\noindent Corrected reading: {\ttfamily\footnotesize |\allowbreak{}}\par
\noindent Source: {\ttfamily\footnotesize 7I:164\allowbreak{}}\par

\par\medskip\noindent\hypertarget{app-F7I-VI-061.1}{}\textbf{F7I-VI-061.1}\quad printed folio 82.\par
\noindent Printed reading: {\ttfamily\footnotesize gives\allowbreak{}}\par
\noindent Corrected reading: {\ttfamily\footnotesize give\allowbreak{}}\par
\noindent Source: {\ttfamily\footnotesize 7I:166\allowbreak{}}\par

\par\medskip\noindent\hypertarget{app-F7I-VI-062.1}{}\textbf{F7I-VI-062.1}\quad printed folio 83.\par
\noindent Printed reading: \emph{(empty)}\par
\noindent Corrected reading: {\ttfamily\footnotesize \textbackslash{}overline\{\allowbreak{}}\par
\noindent Source: {\ttfamily\footnotesize 7I:167\allowbreak{}}\par

\par\medskip\noindent\hypertarget{app-F7I-VI-062.2}{}\textbf{F7I-VI-062.2}\quad printed folio 83.\par
\noindent Printed reading: \emph{(empty)}\par
\noindent Corrected reading: {\ttfamily\footnotesize \}\allowbreak{}}\par
\noindent Source: {\ttfamily\footnotesize 7I:167\allowbreak{}}\par

\par\medskip\noindent\hypertarget{app-F7I-VI-062.3}{}\textbf{F7I-VI-062.3}\quad printed folio 83.\par
\noindent Printed reading: {\ttfamily\footnotesize admits\allowbreak{}}\par
\noindent Corrected reading: {\ttfamily\footnotesize admit\allowbreak{}}\par
\noindent Source: {\ttfamily\footnotesize 7I:168\allowbreak{}}\par

\par\medskip\noindent\hypertarget{app-F7I-VI-063.1}{}\textbf{F7I-VI-063.1}\quad printed folio 84.\par
\noindent Printed reading: {\ttfamily\footnotesize \textasciicircum{}\allowbreak{}}\par
\noindent Corrected reading: {\ttfamily\footnotesize \_\allowbreak{}}\par
\noindent Source: {\ttfamily\footnotesize 7I:169\allowbreak{}}\par

\par\medskip\noindent\hypertarget{app-F7I-VI-064.1}{}\textbf{F7I-VI-064.1}\quad printed folio 85.\par
\noindent Printed reading: {\ttfamily\footnotesize \textbackslash{}overset\allowbreak{}}\par
\noindent Corrected reading: {\ttfamily\footnotesize \textbackslash{}xrightarrow\allowbreak{}}\par
\noindent Source: {\ttfamily\footnotesize FIX-7I-D-0078\allowbreak{}}\par

\par\medskip\noindent\hypertarget{app-F7I-VI-064.2}{}\textbf{F7I-VI-064.2}\quad printed folio 85.\par
\noindent Printed reading: {\ttfamily\footnotesize \}\{\textbackslash{}relbar\textbackslash{}joinrel\textbackslash{}re\allowbreak{}lbar\textbackslash{}joinrel\textbackslash{}relbar\allowbreak{}}\par
\noindent Corrected reading: \emph{(empty)}\par
\noindent Source: {\ttfamily\footnotesize WITNESS-F7I-VI-064.2\allowbreak{}}\par

\par\medskip\noindent\hypertarget{app-F7I-VI-064.3}{}\textbf{F7I-VI-064.3}\quad printed folio 85.\par
\noindent Printed reading: {\ttfamily\footnotesize \textbackslash{}overset\allowbreak{}}\par
\noindent Corrected reading: {\ttfamily\footnotesize \textbackslash{}xrightarrow\allowbreak{}}\par
\noindent Source: {\ttfamily\footnotesize FIX-7I-D-0079\allowbreak{}}\par

\par\medskip\noindent\hypertarget{app-F7I-VI-064.4}{}\textbf{F7I-VI-064.4}\quad printed folio 85.\par
\noindent Printed reading: {\ttfamily\footnotesize \}\{\textbackslash{}relbar\textbackslash{}joinrel\textbackslash{}re\allowbreak{}lbar\textbackslash{}joinrel\textbackslash{}relbar\allowbreak{}}\par
\noindent Corrected reading: \emph{(empty)}\par
\noindent Source: {\ttfamily\footnotesize WITNESS-F7I-VI-064.4\allowbreak{}}\par

\par\medskip\noindent\hypertarget{app-F7I-VI-064.5}{}\textbf{F7I-VI-064.5}\quad printed folio 85.\par
\noindent Printed reading: {\ttfamily\footnotesize \textasciicircum{}\allowbreak{}}\par
\noindent Corrected reading: {\ttfamily\footnotesize \_\allowbreak{}}\par
\noindent Source: {\ttfamily\footnotesize 7I:174\allowbreak{}}\par

\par\medskip\noindent\hypertarget{app-F7I-VI-065.1}{}\textbf{F7I-VI-065.1}\quad printed folio 86.\par
\noindent Printed reading: {\ttfamily\footnotesize \textbackslash{};\allowbreak{}}\par
\noindent Corrected reading: {\ttfamily\footnotesize |\allowbreak{}}\par
\noindent Source: {\ttfamily\footnotesize 7I:175\allowbreak{}}\par

\par\medskip\noindent\hypertarget{app-F7I-VI-065.2}{}\textbf{F7I-VI-065.2}\quad printed folio 86.\par
\noindent Printed reading: {\ttfamily\footnotesize \textbackslash{};\allowbreak{}}\par
\noindent Corrected reading: {\ttfamily\footnotesize |\allowbreak{}}\par
\noindent Source: {\ttfamily\footnotesize 7I:175\allowbreak{}}\par

\par\medskip\noindent\hypertarget{app-F7I-VI-065.3}{}\textbf{F7I-VI-065.3}\quad printed folio 86.\par
\noindent Printed reading: {\ttfamily\footnotesize \textbackslash{};\allowbreak{}}\par
\noindent Corrected reading: {\ttfamily\footnotesize |\allowbreak{}}\par
\noindent Source: {\ttfamily\footnotesize 7I:175\allowbreak{}}\par

\par\medskip\noindent\hypertarget{app-F7I-VI-065.4}{}\textbf{F7I-VI-065.4}\quad printed folio 86.\par
\noindent Printed reading: {\ttfamily\footnotesize \textbackslash{};\allowbreak{}}\par
\noindent Corrected reading: {\ttfamily\footnotesize |\allowbreak{}}\par
\noindent Source: {\ttfamily\footnotesize 7I:175\allowbreak{}}\par

\par\medskip\noindent\hypertarget{app-F7I-VI-065.5}{}\textbf{F7I-VI-065.5}\quad printed folio 86.\par
\noindent Printed reading: {\ttfamily\footnotesize \textbackslash{};\allowbreak{}}\par
\noindent Corrected reading: {\ttfamily\footnotesize |\allowbreak{}}\par
\noindent Source: {\ttfamily\footnotesize 7I:175\allowbreak{}}\par

\par\medskip\noindent\hypertarget{app-F7I-VI-065.6}{}\textbf{F7I-VI-065.6}\quad printed folio 86.\par
\noindent Printed reading: {\ttfamily\footnotesize \textbackslash{};\allowbreak{}}\par
\noindent Corrected reading: {\ttfamily\footnotesize |\allowbreak{}}\par
\noindent Source: {\ttfamily\footnotesize 7I:175\allowbreak{}}\par

\par\medskip\noindent\hypertarget{app-F7I-VI-065.7}{}\textbf{F7I-VI-065.7}\quad printed folio 86.\par
\noindent Printed reading: {\ttfamily\footnotesize \textbackslash{};\allowbreak{}}\par
\noindent Corrected reading: {\ttfamily\footnotesize |\allowbreak{}}\par
\noindent Source: {\ttfamily\footnotesize 7I:175\allowbreak{}}\par

\par\medskip\noindent\hypertarget{app-F7I-VI-065.8}{}\textbf{F7I-VI-065.8}\quad printed folio 86.\par
\noindent Printed reading: {\ttfamily\footnotesize \textbackslash{};\allowbreak{}}\par
\noindent Corrected reading: {\ttfamily\footnotesize |\allowbreak{}}\par
\noindent Source: {\ttfamily\footnotesize 7I:175\allowbreak{}}\par

\par\medskip\noindent\hypertarget{app-F7I-VI-065.9}{}\textbf{F7I-VI-065.9}\quad printed folio 86.\par
\noindent Printed reading: {\ttfamily\footnotesize \textbackslash{};\allowbreak{}}\par
\noindent Corrected reading: {\ttfamily\footnotesize |\allowbreak{}}\par
\noindent Source: {\ttfamily\footnotesize 7I:175\allowbreak{}}\par

\par\medskip\noindent\hypertarget{app-F7I-VI-066.1}{}\textbf{F7I-VI-066.1}\quad printed folio 87.\par
\noindent Printed reading: \emph{(empty)}\par
\noindent Corrected reading: {\ttfamily\footnotesize  an\allowbreak{}}\par
\noindent Source: {\ttfamily\footnotesize 7I:176\allowbreak{}}\par

\par\medskip\noindent\hypertarget{app-F7I-VI-067.1}{}\textbf{F7I-VI-067.1}\quad printed folio 87.\par
\noindent Printed reading: \emph{(empty)}\par
\noindent Corrected reading: {\ttfamily\footnotesize )\allowbreak{}}\par
\noindent Source: {\ttfamily\footnotesize FIX-7I-D-0081\allowbreak{}}\par

\par\medskip\noindent\hypertarget{app-F7I-VI-068.1}{}\textbf{F7I-VI-068.1}\quad printed folio 89.\par
\noindent Printed reading: {\ttfamily\footnotesize extensions\allowbreak{}}\par
\noindent Corrected reading: {\ttfamily\footnotesize extension\allowbreak{}}\par
\noindent Source: {\ttfamily\footnotesize 7I:178\allowbreak{}}\par

\par\medskip\noindent\hypertarget{app-F7I-VI-069.1}{}\textbf{F7I-VI-069.1}\quad printed folio 89.\par
\noindent Printed reading: {\ttfamily\footnotesize is\allowbreak{}}\par
\noindent Corrected reading: {\ttfamily\footnotesize are\allowbreak{}}\par
\noindent Source: {\ttfamily\footnotesize 7I:179\allowbreak{}}\par

\par\medskip\noindent\hypertarget{app-F7I-VI-069.2}{}\textbf{F7I-VI-069.2}\quad printed folio 89.\par
\noindent Printed reading: {\ttfamily\footnotesize process\allowbreak{}}\par
\noindent Corrected reading: {\ttfamily\footnotesize processes\allowbreak{}}\par
\noindent Source: {\ttfamily\footnotesize 7I:180\allowbreak{}}\par

\par\medskip\noindent\hypertarget{app-F7I-VI-069.3}{}\textbf{F7I-VI-069.3}\quad printed folio 89.\par
\noindent Printed reading: \emph{(empty)}\par
\noindent Corrected reading: {\ttfamily\footnotesize The\allowbreak{} }\par
\noindent Source: {\ttfamily\footnotesize WITNESS-F7I-VI-069.3\allowbreak{}}\par

\par\medskip\noindent\hypertarget{app-F7I-VI-070.1}{}\textbf{F7I-VI-070.1}\quad printed folio 90.\par
\noindent Printed reading: {\ttfamily\footnotesize all\allowbreak{}}\par
\noindent Corrected reading: {\ttfamily\footnotesize every\allowbreak{}}\par
\noindent Source: {\ttfamily\footnotesize WITNESS-F7I-VI-070.1\allowbreak{}}\par

\par\medskip\noindent\hypertarget{app-F7I-VI-070.2}{}\textbf{F7I-VI-070.2}\quad printed folio 90.\par
\noindent Printed reading: \emph{(empty)}\par
\noindent Corrected reading: {\ttfamily\footnotesize an\allowbreak{} }\par
\noindent Source: {\ttfamily\footnotesize ADOPT-F7I-VI-070.2\allowbreak{}}\par

\par\medskip\noindent\hypertarget{app-F7I-VI-071.1}{}\textbf{F7I-VI-071.1}\quad printed folio 90.\par
\noindent Printed reading: {\ttfamily\footnotesize \textbackslash{};\allowbreak{}}\par
\noindent Corrected reading: {\ttfamily\footnotesize |\allowbreak{}}\par
\noindent Source: {\ttfamily\footnotesize 7I:183\allowbreak{}}\par

\par\medskip\noindent\hypertarget{app-F7I-VI-071.2}{}\textbf{F7I-VI-071.2}\quad printed folio 91.\par
\noindent Printed reading: \emph{(empty)}\par
\noindent Corrected reading: {\ttfamily\footnotesize \textbackslash{}hat\{\allowbreak{}}\par
\noindent Source: {\ttfamily\footnotesize ADOPT-F7I-VI-071.2\allowbreak{}}\par

\par\medskip\noindent\hypertarget{app-F7I-VI-071.3}{}\textbf{F7I-VI-071.3}\quad printed folio 91.\par
\noindent Printed reading: \emph{(empty)}\par
\noindent Corrected reading: {\ttfamily\footnotesize \}\allowbreak{}}\par
\noindent Source: {\ttfamily\footnotesize ADOPT-F7I-VI-071.3\allowbreak{}}\par

\par\medskip\noindent\hypertarget{app-F7I-VI-071.4}{}\textbf{F7I-VI-071.4}\quad printed folio 91.\par
\noindent Printed reading: \emph{(empty)}\par
\noindent Corrected reading: {\ttfamily\footnotesize  whose\allowbreak{} restriction\allowbreak{} to\allowbreak{} \$A\$\allowbreak{} splits\allowbreak{}}\par
\noindent Source: {\ttfamily\footnotesize ADOPT-F7I-VI-071.4\allowbreak{}}\par

\par\medskip\noindent\hypertarget{app-F7I-VI-071.5}{}\textbf{F7I-VI-071.5}\quad printed folio 91.\par
\noindent Printed reading: {\ttfamily\footnotesize such\allowbreak{}}\par
\noindent Corrected reading: {\ttfamily\footnotesize lifting\allowbreak{}}\par
\noindent Source: {\ttfamily\footnotesize ADOPT-F7I-VI-071.5\allowbreak{}}\par

\par\medskip\noindent\hypertarget{app-F7I-VI-071.6}{}\textbf{F7I-VI-071.6}\quad printed folio 91.\par
\noindent Printed reading: {\ttfamily\footnotesize that\allowbreak{} }\par
\noindent Corrected reading: {\ttfamily\footnotesize the\allowbreak{} canonical\allowbreak{} map\allowbreak{} }\par
\noindent Source: {\ttfamily\footnotesize ADOPT-F7I-VI-071.6\allowbreak{}}\par

\par\medskip\noindent\hypertarget{app-F7I-VI-071.7}{}\textbf{F7I-VI-071.7}\quad printed folio 91.\par
\noindent Printed reading: {\ttfamily\footnotesize \textbackslash{}varphi\allowbreak{}}\par
\noindent Corrected reading: {\ttfamily\footnotesize A\allowbreak{}}\par
\noindent Source: {\ttfamily\footnotesize ADOPT-F7I-VI-071.7\allowbreak{}}\par

\par\medskip\noindent\hypertarget{app-F7I-VI-071.8}{}\textbf{F7I-VI-071.8}\quad printed folio 91.\par
\noindent Printed reading: {\ttfamily\footnotesize \textbackslash{}circ\allowbreak{}}\par
\noindent Corrected reading: {\ttfamily\footnotesize \textbackslash{}to\allowbreak{}}\par
\noindent Source: {\ttfamily\footnotesize ADOPT-F7I-VI-071.8\allowbreak{}}\par

\par\medskip\noindent\hypertarget{app-F7I-VI-071.9}{}\textbf{F7I-VI-071.9}\quad printed folio 91.\par
\noindent Printed reading: {\ttfamily\footnotesize \textbackslash{}psi\allowbreak{} =\allowbreak{} \textbackslash{}mathrm\allowbreak{}}\par
\noindent Corrected reading: {\ttfamily\footnotesize \textbackslash{}hat\allowbreak{}}\par
\noindent Source: {\ttfamily\footnotesize ADOPT-F7I-VI-071.9\allowbreak{}}\par

\par\medskip\noindent\hypertarget{app-F7I-VI-071.10}{}\textbf{F7I-VI-071.10}\quad printed folio 91.\par
\noindent Printed reading: {\ttfamily\footnotesize id\allowbreak{}}\par
\noindent Corrected reading: {\ttfamily\footnotesize A\allowbreak{}}\par
\noindent Source: {\ttfamily\footnotesize ADOPT-F7I-VI-071.10\allowbreak{}}\par

\par\medskip\noindent\hypertarget{app-F7I-VI-071.11}{}\textbf{F7I-VI-071.11}\quad printed folio 91.\par
\noindent Printed reading: {\ttfamily\footnotesize \_A\allowbreak{}}\par
\noindent Corrected reading: \emph{(empty)}\par
\noindent Source: {\ttfamily\footnotesize ADOPT-F7I-VI-071.11\allowbreak{}}\par

\par\medskip\noindent\hypertarget{app-F7I-VI-071.12}{}\textbf{F7I-VI-071.12}\quad printed folio 91.\par
\noindent Printed reading: \emph{(empty)}\par
\noindent Corrected reading: {\ttfamily\footnotesize \textbackslash{}hat\{\allowbreak{}}\par
\noindent Source: {\ttfamily\footnotesize 7I:190\allowbreak{}}\par

\par\medskip\noindent\hypertarget{app-F7I-VI-071.13}{}\textbf{F7I-VI-071.13}\quad printed folio 91.\par
\noindent Printed reading: \emph{(empty)}\par
\noindent Corrected reading: {\ttfamily\footnotesize \}\allowbreak{}}\par
\noindent Source: {\ttfamily\footnotesize 7I:190\allowbreak{}}\par

\par\medskip\noindent\hypertarget{app-F7I-VI-071.14}{}\textbf{F7I-VI-071.14}\quad printed folio 91.\par
\noindent Printed reading: \emph{(empty)}\par
\noindent Corrected reading: {\ttfamily\footnotesize \textbackslash{}hat\{\allowbreak{}}\par
\noindent Source: {\ttfamily\footnotesize 7I:191\allowbreak{}}\par

\par\medskip\noindent\hypertarget{app-F7I-VI-071.15}{}\textbf{F7I-VI-071.15}\quad printed folio 91.\par
\noindent Printed reading: \emph{(empty)}\par
\noindent Corrected reading: {\ttfamily\footnotesize \}\allowbreak{}}\par
\noindent Source: {\ttfamily\footnotesize 7I:191\allowbreak{}}\par

\par\medskip\noindent\hypertarget{app-F7I-VI-072.1}{}\textbf{F7I-VI-072.1}\quad printed folio 91.\par
\noindent Printed reading: {\ttfamily\footnotesize  \$D\textasciicircum{}1\$\allowbreak{}}\par
\noindent Corrected reading: \emph{(empty)}\par
\noindent Source: {\ttfamily\footnotesize WITNESS-F7I-VI-072.1\allowbreak{}}\par

\par\medskip\noindent\hypertarget{app-F7I-VI-073.1}{}\textbf{F7I-VI-073.1}\quad printed folio 92.\par
\noindent Printed reading: \emph{(empty)}\par
\noindent Corrected reading: {\ttfamily\footnotesize  a\allowbreak{}}\par
\noindent Source: {\ttfamily\footnotesize WITNESS-F7I-VI-073.1\allowbreak{}}\par

\par\medskip\noindent\hypertarget{app-F7I-VI-073.2}{}\textbf{F7I-VI-073.2}\quad printed folio 92.\par
\noindent Printed reading: {\ttfamily\footnotesize M\allowbreak{}}\par
\noindent Corrected reading: {\ttfamily\footnotesize E\allowbreak{}}\par
\noindent Source: {\ttfamily\footnotesize 7I:193\allowbreak{}}\par

\par\medskip\noindent\hypertarget{app-F7I-VI-073.3}{}\textbf{F7I-VI-073.3}\quad printed folio 92.\par
\noindent Printed reading: {\ttfamily\footnotesize \textbackslash{}alpha\allowbreak{}}\par
\noindent Corrected reading: {\ttfamily\footnotesize A\allowbreak{}}\par
\noindent Source: {\ttfamily\footnotesize 7I:194\allowbreak{}}\par

\par\medskip\noindent\hypertarget{app-F7I-VI-073.4}{}\textbf{F7I-VI-073.4}\quad printed folio 92.\par
\noindent Printed reading: \emph{(empty)}\par
\noindent Corrected reading: {\ttfamily\footnotesize  to\allowbreak{} say\allowbreak{} that\allowbreak{}}\par
\noindent Source: {\ttfamily\footnotesize ADOPT-F7I-VI-073.4\allowbreak{}}\par

\par\medskip\noindent\hypertarget{app-F7I-VI-073.5}{}\textbf{F7I-VI-073.5}\quad printed folio 92.\par
\noindent Printed reading: \emph{(empty)}\par
\noindent Corrected reading: {\ttfamily\footnotesize  as\allowbreak{} a\allowbreak{}}\par
\noindent Source: {\ttfamily\footnotesize ADOPT-F7I-VI-073.5\allowbreak{}}\par

\par\medskip\noindent\hypertarget{app-F7I-VI-074.1}{}\textbf{F7I-VI-074.1}\quad printed folio 92.\par
\noindent Printed reading: {\ttfamily\footnotesize a\allowbreak{}}\par
\noindent Corrected reading: {\ttfamily\footnotesize c\allowbreak{}}\par
\noindent Source: {\ttfamily\footnotesize 7I:195\allowbreak{}}\par

\par\medskip\noindent\hypertarget{app-F7I-VI-074.2}{}\textbf{F7I-VI-074.2}\quad printed folio 92.\par
\noindent Printed reading: {\ttfamily\footnotesize abreviate\allowbreak{}}\par
\noindent Corrected reading: {\ttfamily\footnotesize abbreviate\allowbreak{}}\par
\noindent Source: {\ttfamily\footnotesize 7I:196\allowbreak{}}\par

\par\medskip\noindent\hypertarget{app-F7I-VI-074.3}{}\textbf{F7I-VI-074.3}\quad printed folio 92.\par
\noindent Printed reading: {\ttfamily\footnotesize is\allowbreak{}}\par
\noindent Corrected reading: {\ttfamily\footnotesize are\allowbreak{}}\par
\noindent Source: {\ttfamily\footnotesize 7I:197\allowbreak{}}\par

\par\medskip\noindent\hypertarget{app-F7I-VI-075.1}{}\textbf{F7I-VI-075.1}\quad printed folio 93.\par
\noindent Printed reading: {\ttfamily\footnotesize h\allowbreak{}}\par
\noindent Corrected reading: {\ttfamily\footnotesize k\allowbreak{}}\par
\noindent Source: {\ttfamily\footnotesize 7I:198\allowbreak{}}\par

\par\medskip\noindent\hypertarget{app-F7I-VI-075.2}{}\textbf{F7I-VI-075.2}\quad printed folio 93.\par
\noindent Printed reading: {\ttfamily\footnotesize 6\allowbreak{}}\par
\noindent Corrected reading: {\ttfamily\footnotesize 7\allowbreak{}}\par
\noindent Source: {\ttfamily\footnotesize WITNESS-F7I-VI-075.2\allowbreak{}}\par

\par\medskip\noindent\hypertarget{app-F7I-VI-075.3}{}\textbf{F7I-VI-075.3}\quad printed folio 93.\par
\noindent Printed reading: \emph{(empty)}\par
\noindent Corrected reading: {\ttfamily\footnotesize the\allowbreak{} vanishing\allowbreak{} for\allowbreak{} smooth\allowbreak{} algebras\allowbreak{} }\par
\noindent Source: {\ttfamily\footnotesize WITNESS-F7I-VI-075.3\allowbreak{}}\par

\par\medskip\noindent\hypertarget{app-F7I-VI-075.4}{}\textbf{F7I-VI-075.4}\quad printed folio 93.\par
\noindent Printed reading: {\ttfamily\footnotesize 8\allowbreak{}}\par
\noindent Corrected reading: {\ttfamily\footnotesize cf.\allowbreak{} 3.13(a\allowbreak{}}\par
\noindent Source: {\ttfamily\footnotesize WITNESS-F7I-VI-075.4\allowbreak{}}\par

\par\medskip\noindent\hypertarget{app-F7I-VI-075.5}{}\textbf{F7I-VI-075.5}\quad printed folio 93.\par
\noindent Printed reading: \emph{(empty)}\par
\noindent Corrected reading: {\ttfamily\footnotesize ),\allowbreak{}}\par
\noindent Source: {\ttfamily\footnotesize WITNESS-F7I-VI-075.5\allowbreak{}}\par

\par\medskip\noindent\hypertarget{app-F7I-VI-075.6}{}\textbf{F7I-VI-075.6}\quad printed folio 93.\par
\noindent Printed reading: \emph{(empty)}\par
\noindent Corrected reading: {\ttfamily\footnotesize the\allowbreak{} complete-intersectio\allowbreak{}n\allowbreak{} vanishing\allowbreak{} (cf.\allowbreak{} }\par
\noindent Source: {\ttfamily\footnotesize WITNESS-F7I-VI-075.6\allowbreak{}}\par

\par\medskip\noindent\hypertarget{app-F7I-VI-075.7}{}\textbf{F7I-VI-075.7}\quad printed folio 93.\par
\noindent Printed reading: {\ttfamily\footnotesize 12\allowbreak{}}\par
\noindent Corrected reading: {\ttfamily\footnotesize 13\allowbreak{}}\par
\noindent Source: {\ttfamily\footnotesize 7I:200\allowbreak{}}\par

\par\medskip\noindent\hypertarget{app-F7I-VI-075.8}{}\textbf{F7I-VI-075.8}\quad printed folio 93.\par
\noindent Printed reading: {\ttfamily\footnotesize 10\allowbreak{}}\par
\noindent Corrected reading: {\ttfamily\footnotesize b)\allowbreak{}}\par
\noindent Source: {\ttfamily\footnotesize WITNESS-F7I-VI-075.8\allowbreak{}}\par

\par\medskip\noindent\hypertarget{app-F7I-VI-076.1}{}\textbf{F7I-VI-076.1}\quad printed folio 93.\par
\noindent Printed reading: {\ttfamily\footnotesize 7\allowbreak{}}\par
\noindent Corrected reading: {\ttfamily\footnotesize 8\allowbreak{}}\par
\noindent Source: {\ttfamily\footnotesize 7I:201\allowbreak{}}\par

\par\medskip\noindent\hypertarget{app-F7I-VI-076.2}{}\textbf{F7I-VI-076.2}\quad printed folio 94.\par
\noindent Printed reading: {\ttfamily\footnotesize all\allowbreak{}}\par
\noindent Corrected reading: {\ttfamily\footnotesize every\allowbreak{}}\par
\noindent Source: {\ttfamily\footnotesize WITNESS-F7I-VI-076.2\allowbreak{}}\par

\par\medskip\noindent\hypertarget{app-F7I-VI-077.1}{}\textbf{F7I-VI-077.1}\quad printed folio 94.\par
\noindent Printed reading: {\ttfamily\footnotesize Andre\allowbreak{}}\par
\noindent Corrected reading: {\ttfamily\footnotesize André\allowbreak{}}\par
\noindent Source: {\ttfamily\footnotesize ADOPT-F7I-VI-077.1\allowbreak{}}\par

\par\medskip\noindent\hypertarget{app-F7I-VI-078.1}{}\textbf{F7I-VI-078.1}\quad printed folio 95.\par
\noindent Printed reading: {\ttfamily\footnotesize spec\allowbreak{}}\par
\noindent Corrected reading: {\ttfamily\footnotesize Spec\allowbreak{}}\par
\noindent Source: {\ttfamily\footnotesize FIX-7I-D-0103\allowbreak{}}\par

\par\medskip\noindent\hypertarget{app-F7I-VI-079.1}{}\textbf{F7I-VI-079.1}\quad printed folio 95.\par
\noindent Printed reading: {\ttfamily\footnotesize p\allowbreak{}}\par
\noindent Corrected reading: {\ttfamily\footnotesize P\allowbreak{}}\par
\noindent Source: {\ttfamily\footnotesize ADOPT-F7I-VI-079.1\allowbreak{}}\par

\par\medskip\noindent\hypertarget{app-F7I-VI-080.1}{}\textbf{F7I-VI-080.1}\quad printed folio 98.\par
\noindent Printed reading: \emph{(empty)}\par
\noindent Corrected reading: {\ttfamily\footnotesize '\allowbreak{}}\par
\noindent Source: {\ttfamily\footnotesize 7I:184\allowbreak{}}\par

\par\medskip\noindent\hypertarget{app-F7I-VI-081.1}{}\textbf{F7I-VI-081.1}\quad printed folio 98.\par
\noindent Printed reading: {\ttfamily\footnotesize \textbackslash{}phantom\{\allowbreak{}}\par
\noindent Corrected reading: \emph{(empty)}\par
\noindent Source: {\ttfamily\footnotesize 7I:185\allowbreak{}}\par

\par\medskip\noindent\hypertarget{app-F7I-VI-081.2}{}\textbf{F7I-VI-081.2}\quad printed folio 98.\par
\noindent Printed reading: {\ttfamily\footnotesize \}\allowbreak{}}\par
\noindent Corrected reading: \emph{(empty)}\par
\noindent Source: {\ttfamily\footnotesize 7I:185\allowbreak{}}\par

\par\medskip\noindent\hypertarget{app-F7I-VI-082.1}{}\textbf{F7I-VI-082.1}\quad printed folio 99.\par
\noindent Printed reading: {\ttfamily\footnotesize \textbackslash{}relbar\textbackslash{}joinrel\textbackslash{}relb\allowbreak{}ar\textbackslash{}joinrel\textbackslash{}relbar\allowbreak{}}\par
\noindent Corrected reading: {\ttfamily\footnotesize \textbackslash{}longrightarrow\allowbreak{}}\par
\noindent Source: {\ttfamily\footnotesize FIX-7I-D-0084\allowbreak{}}\par

\par\medskip\noindent\hypertarget{app-F7I-VI-082.2}{}\textbf{F7I-VI-082.2}\quad printed folio 99.\par
\noindent Printed reading: {\ttfamily\footnotesize \textbackslash{}relbar\textbackslash{}joinrel\textbackslash{}relb\allowbreak{}ar\textbackslash{}joinrel\textbackslash{}relbar\allowbreak{}}\par
\noindent Corrected reading: {\ttfamily\footnotesize \textbackslash{}longrightarrow\allowbreak{}}\par
\noindent Source: {\ttfamily\footnotesize FIX-7I-D-0085\allowbreak{}}\par

\par\medskip\noindent\hypertarget{app-F7I-VI-082.3}{}\textbf{F7I-VI-082.3}\quad printed folio 99.\par
\noindent Printed reading: {\ttfamily\footnotesize \textbackslash{}relbar\textbackslash{}joinrel\textbackslash{}relb\allowbreak{}ar\textbackslash{}joinrel\textbackslash{}relbar\allowbreak{}}\par
\noindent Corrected reading: {\ttfamily\footnotesize \textbackslash{}longrightarrow\allowbreak{}}\par
\noindent Source: {\ttfamily\footnotesize FIX-7I-D-0085\allowbreak{}}\par

\par\medskip\noindent\hypertarget{app-F7I-VI-083.1}{}\textbf{F7I-VI-083.1}\quad printed folio 99.\par
\noindent Printed reading: {\ttfamily\footnotesize \textbackslash{}relbar\textbackslash{}joinrel\textbackslash{}relb\allowbreak{}ar\textbackslash{}joinrel\textbackslash{}relbar\allowbreak{}}\par
\noindent Corrected reading: {\ttfamily\footnotesize \textbackslash{}longrightarrow\allowbreak{}}\par
\noindent Source: {\ttfamily\footnotesize FIX-7I-D-0087\allowbreak{}}\par

\par\medskip\noindent\hypertarget{app-F7I-VI-084.1}{}\textbf{F7I-VI-084.1}\quad printed folio 100.\par
\noindent Printed reading: {\ttfamily\footnotesize \textbackslash{}relbar\textbackslash{}joinrel\textbackslash{}relb\allowbreak{}ar\textbackslash{}joinrel\textbackslash{}relbar\allowbreak{}}\par
\noindent Corrected reading: {\ttfamily\footnotesize \textbackslash{}longrightarrow\allowbreak{}}\par
\noindent Source: {\ttfamily\footnotesize FIX-7I-D-0088\allowbreak{}}\par

\par\medskip\noindent\hypertarget{app-F7I-VI-085.1}{}\textbf{F7I-VI-085.1}\quad printed folio 101.\par
\noindent Printed reading: \emph{(empty)}\par
\noindent Corrected reading: {\ttfamily\footnotesize \_\{X\_\{0\}\}\allowbreak{}}\par
\noindent Source: {\ttfamily\footnotesize 7I:204\allowbreak{}}\par

\par\medskip\noindent\hypertarget{app-F7I-VI-086.1}{}\textbf{F7I-VI-086.1}\quad printed folio 101.\par
\noindent Printed reading: {\ttfamily\footnotesize 2\allowbreak{}}\par
\noindent Corrected reading: {\ttfamily\footnotesize 3\allowbreak{}}\par
\noindent Source: {\ttfamily\footnotesize ADOPT-F7I-VI-086.1\allowbreak{}}\par

\par\medskip\noindent\hypertarget{app-F7I-VI-086.2}{}\textbf{F7I-VI-086.2}\quad printed folio 101.\par
\noindent Printed reading: {\ttfamily\footnotesize 13\allowbreak{}}\par
\noindent Corrected reading: {\ttfamily\footnotesize 15\allowbreak{}}\par
\noindent Source: {\ttfamily\footnotesize ADOPT-F7I-VI-086.2\allowbreak{}}\par

\par\medskip\noindent\hypertarget{app-F7I-VI-086.3}{}\textbf{F7I-VI-086.3}\quad printed folio 101.\par
\noindent Printed reading: {\ttfamily\footnotesize 2\allowbreak{}}\par
\noindent Corrected reading: {\ttfamily\footnotesize 3\allowbreak{}}\par
\noindent Source: {\ttfamily\footnotesize ADOPT-F7I-VI-086.3\allowbreak{}}\par

\par\medskip\noindent\hypertarget{app-F7I-VI-086.4}{}\textbf{F7I-VI-086.4}\quad printed folio 101.\par
\noindent Printed reading: {\ttfamily\footnotesize 13\allowbreak{}}\par
\noindent Corrected reading: {\ttfamily\footnotesize 15\allowbreak{}}\par
\noindent Source: {\ttfamily\footnotesize ADOPT-F7I-VI-086.4\allowbreak{}}\par

\par\medskip\noindent\hypertarget{app-F7I-VI-087.1}{}\textbf{F7I-VI-087.1}\quad printed folio 102.\par
\noindent Printed reading: {\ttfamily\footnotesize no\allowbreak{} head,\allowbreak{} }\par
\noindent Corrected reading: \emph{(empty)}\par
\noindent Source: {\ttfamily\footnotesize WITNESS-F7I-VI-087.1\allowbreak{}}\par

\par\medskip\noindent\hypertarget{app-F7I-VI-088.1}{}\textbf{F7I-VI-088.1}\quad printed folio 103.\par
\noindent Printed reading: \emph{(empty)}\par
\noindent Corrected reading: {\ttfamily\footnotesize \_0\allowbreak{}}\par
\noindent Source: {\ttfamily\footnotesize 7I:206\allowbreak{}}\par

\par\medskip\noindent\hypertarget{app-F7I-VI-089.1}{}\textbf{F7I-VI-089.1}\quad printed folio 103.\par
\noindent Printed reading: {\ttfamily\footnotesize \textbackslash{}mathfrak\{J\}\allowbreak{}}\par
\noindent Corrected reading: {\ttfamily\footnotesize j\allowbreak{}}\par
\noindent Source: {\ttfamily\footnotesize FIX-7I-D-0107\allowbreak{}}\par

\par\medskip\noindent\hypertarget{app-F7I-VI-090.1}{}\textbf{F7I-VI-090.1}\quad printed folio 104.\par
\noindent Printed reading: \emph{(empty)}\par
\noindent Corrected reading: {\ttfamily\footnotesize \_0\allowbreak{}}\par
\noindent Source: {\ttfamily\footnotesize 7I:208\allowbreak{}}\par

\par\medskip\noindent\hypertarget{app-F7I-VI-091.1}{}\textbf{F7I-VI-091.1}\quad printed folio 104.\par
\noindent Printed reading: \emph{(empty)}\par
\noindent Corrected reading: {\ttfamily\footnotesize (\allowbreak{}}\par
\noindent Source: {\ttfamily\footnotesize WITNESS-F7I-VI-091.1\allowbreak{}}\par

\par\medskip\noindent\hypertarget{app-F7I-VI-091.2}{}\textbf{F7I-VI-091.2}\quad printed folio 104.\par
\noindent Printed reading: {\ttfamily\footnotesize ,\textbackslash{},\allowbreak{}}\par
\noindent Corrected reading: \emph{(empty)}\par
\noindent Source: {\ttfamily\footnotesize 7I:209\allowbreak{}}\par

\par\medskip\noindent\hypertarget{app-F7I-VI-092.1}{}\textbf{F7I-VI-092.1}\quad printed folio 104.\par
\noindent Printed reading: {\ttfamily\footnotesize Firstly\allowbreak{}}\par
\noindent Corrected reading: {\ttfamily\footnotesize First,\allowbreak{}}\par
\noindent Source: {\ttfamily\footnotesize ADOPT-F7I-VI-092.1\allowbreak{}}\par

\par\medskip\noindent\hypertarget{app-F7I-VI-092.2}{}\textbf{F7I-VI-092.2}\quad printed folio 104.\par
\noindent Printed reading: {\ttfamily\footnotesize on\allowbreak{} a\allowbreak{}}\par
\noindent Corrected reading: {\ttfamily\footnotesize some\allowbreak{}}\par
\noindent Source: {\ttfamily\footnotesize ADOPT-F7I-VI-092.2\allowbreak{}}\par

\par\medskip\noindent\hypertarget{app-F7I-VI-093.1}{}\textbf{F7I-VI-093.1}\quad printed folio 105.\par
\noindent Printed reading: {\ttfamily\footnotesize it\allowbreak{} }\par
\noindent Corrected reading: \emph{(empty)}\par
\noindent Source: {\ttfamily\footnotesize ADOPT-F7I-VI-093.1\allowbreak{}}\par

\par\medskip\noindent\hypertarget{app-F7I-VI-093.2}{}\textbf{F7I-VI-093.2}\quad printed folio 105.\par
\noindent Printed reading: \emph{(empty)}\par
\noindent Corrected reading: {\ttfamily\footnotesize \textbackslash{}widehat\{\allowbreak{}}\par
\noindent Source: {\ttfamily\footnotesize ADOPT-F7I-VI-093.2\allowbreak{}}\par

\par\medskip\noindent\hypertarget{app-F7I-VI-093.3}{}\textbf{F7I-VI-093.3}\quad printed folio 105.\par
\noindent Printed reading: \emph{(empty)}\par
\noindent Corrected reading: {\ttfamily\footnotesize \}\allowbreak{}}\par
\noindent Source: {\ttfamily\footnotesize ADOPT-F7I-VI-093.3\allowbreak{}}\par

\par\medskip\noindent\hypertarget{app-F7I-VI-093.4}{}\textbf{F7I-VI-093.4}\quad printed folio 105.\par
\noindent Printed reading: \emph{(empty)}\par
\noindent Corrected reading: {\ttfamily\footnotesize \textbackslash{}widehat\{\allowbreak{}}\par
\noindent Source: {\ttfamily\footnotesize ADOPT-F7I-VI-093.4\allowbreak{}}\par

\par\medskip\noindent\hypertarget{app-F7I-VI-093.5}{}\textbf{F7I-VI-093.5}\quad printed folio 105.\par
\noindent Printed reading: \emph{(empty)}\par
\noindent Corrected reading: {\ttfamily\footnotesize \}\allowbreak{}}\par
\noindent Source: {\ttfamily\footnotesize ADOPT-F7I-VI-093.5\allowbreak{}}\par

\par\medskip\noindent\hypertarget{app-F7I-VI-093.6}{}\textbf{F7I-VI-093.6}\quad printed folio 105.\par
\noindent Printed reading: {\ttfamily\footnotesize write\allowbreak{}}\par
\noindent Corrected reading: {\ttfamily\footnotesize denote\allowbreak{}}\par
\noindent Source: {\ttfamily\footnotesize ADOPT-F7I-VI-093.6\allowbreak{}}\par

\par\medskip\noindent\hypertarget{app-F7I-VI-093.7}{}\textbf{F7I-VI-093.7}\quad printed folio 105.\par
\noindent Printed reading: \emph{(empty)}\par
\noindent Corrected reading: {\ttfamily\footnotesize \textbackslash{}widehat\{\allowbreak{}}\par
\noindent Source: {\ttfamily\footnotesize ADOPT-F7I-VI-093.7\allowbreak{}}\par

\par\medskip\noindent\hypertarget{app-F7I-VI-093.8}{}\textbf{F7I-VI-093.8}\quad printed folio 105.\par
\noindent Printed reading: \emph{(empty)}\par
\noindent Corrected reading: {\ttfamily\footnotesize \}\allowbreak{}}\par
\noindent Source: {\ttfamily\footnotesize ADOPT-F7I-VI-093.8\allowbreak{}}\par

\par\medskip\noindent\hypertarget{app-F7I-VI-093.9}{}\textbf{F7I-VI-093.9}\quad printed folio 105.\par
\noindent Printed reading: \emph{(empty)}\par
\noindent Corrected reading: {\ttfamily\footnotesize \textbackslash{}widehat\{\allowbreak{}}\par
\noindent Source: {\ttfamily\footnotesize ADOPT-F7I-VI-093.9\allowbreak{}}\par

\par\medskip\noindent\hypertarget{app-F7I-VI-093.10}{}\textbf{F7I-VI-093.10}\quad printed folio 105.\par
\noindent Printed reading: \emph{(empty)}\par
\noindent Corrected reading: {\ttfamily\footnotesize \}\allowbreak{}}\par
\noindent Source: {\ttfamily\footnotesize ADOPT-F7I-VI-093.10\allowbreak{}}\par

\par\medskip\noindent\hypertarget{app-F7I-VI-093.11}{}\textbf{F7I-VI-093.11}\quad printed folio 105.\par
\noindent Printed reading: \emph{(empty)}\par
\noindent Corrected reading: {\ttfamily\footnotesize \textbackslash{}widehat\{\allowbreak{}}\par
\noindent Source: {\ttfamily\footnotesize WITNESS-F7I-VI-093.1\allowbreak{}1\allowbreak{}}\par

\par\medskip\noindent\hypertarget{app-F7I-VI-093.12}{}\textbf{F7I-VI-093.12}\quad printed folio 105.\par
\noindent Printed reading: \emph{(empty)}\par
\noindent Corrected reading: {\ttfamily\footnotesize \}\allowbreak{}}\par
\noindent Source: {\ttfamily\footnotesize WITNESS-F7I-VI-093.1\allowbreak{}2\allowbreak{}}\par

\par\medskip\noindent\hypertarget{app-F7I-VI-093.13}{}\textbf{F7I-VI-093.13}\quad printed folio 105.\par
\noindent Printed reading: \emph{(empty)}\par
\noindent Corrected reading: {\ttfamily\footnotesize \textbackslash{}widehat\{\allowbreak{}}\par
\noindent Source: {\ttfamily\footnotesize WITNESS-F7I-VI-093.1\allowbreak{}3\allowbreak{}}\par

\par\medskip\noindent\hypertarget{app-F7I-VI-093.14}{}\textbf{F7I-VI-093.14}\quad printed folio 105.\par
\noindent Printed reading: \emph{(empty)}\par
\noindent Corrected reading: {\ttfamily\footnotesize \}\allowbreak{}}\par
\noindent Source: {\ttfamily\footnotesize WITNESS-F7I-VI-093.1\allowbreak{}4\allowbreak{}}\par

\par\medskip\noindent\hypertarget{app-F7I-VI-093.15}{}\textbf{F7I-VI-093.15}\quad printed folio 105.\par
\noindent Printed reading: \emph{(empty)}\par
\noindent Corrected reading: {\ttfamily\footnotesize \textbackslash{}widehat\{\allowbreak{}}\par
\noindent Source: {\ttfamily\footnotesize WITNESS-F7I-VI-093.1\allowbreak{}5\allowbreak{}}\par

\par\medskip\noindent\hypertarget{app-F7I-VI-093.16}{}\textbf{F7I-VI-093.16}\quad printed folio 105.\par
\noindent Printed reading: \emph{(empty)}\par
\noindent Corrected reading: {\ttfamily\footnotesize \}\allowbreak{}}\par
\noindent Source: {\ttfamily\footnotesize WITNESS-F7I-VI-093.1\allowbreak{}6\allowbreak{}}\par

\par\medskip\noindent\hypertarget{app-F7I-VI-093.17}{}\textbf{F7I-VI-093.17}\quad printed folio 105.\par
\noindent Printed reading: \emph{(empty)}\par
\noindent Corrected reading: {\ttfamily\footnotesize \textbackslash{}hat\{\allowbreak{}}\par
\noindent Source: {\ttfamily\footnotesize 7I:211\allowbreak{}}\par

\par\medskip\noindent\hypertarget{app-F7I-VI-093.18}{}\textbf{F7I-VI-093.18}\quad printed folio 105.\par
\noindent Printed reading: \emph{(empty)}\par
\noindent Corrected reading: {\ttfamily\footnotesize \}\allowbreak{}}\par
\noindent Source: {\ttfamily\footnotesize 7I:211\allowbreak{}}\par

\par\medskip\noindent\hypertarget{app-F7I-VI-093.19}{}\textbf{F7I-VI-093.19}\quad printed folio 105.\par
\noindent Printed reading: {\ttfamily\footnotesize X\allowbreak{}}\par
\noindent Corrected reading: {\ttfamily\footnotesize x\allowbreak{}}\par
\noindent Source: {\ttfamily\footnotesize FIX-7I-D-0112\allowbreak{}}\par

\par\medskip\noindent\hypertarget{app-F7I-VI-094.1}{}\textbf{F7I-VI-094.1}\quad printed folio 105.\par
\noindent Printed reading: {\ttfamily\footnotesize the\allowbreak{} }\par
\noindent Corrected reading: \emph{(empty)}\par
\noindent Source: {\ttfamily\footnotesize ADOPT-F7I-VI-094.1\allowbreak{}}\par

\par\medskip\noindent\hypertarget{app-F7I-VI-095.1}{}\textbf{F7I-VI-095.1}\quad printed folio 106.\par
\noindent Printed reading: {\ttfamily\footnotesize of\allowbreak{}}\par
\noindent Corrected reading: {\ttfamily\footnotesize for\allowbreak{}}\par
\noindent Source: {\ttfamily\footnotesize ADOPT-F7I-VI-095.1\allowbreak{}}\par

\par\medskip\noindent\hypertarget{app-F7I-VI-095.2}{}\textbf{F7I-VI-095.2}\quad printed folio 106.\par
\noindent Printed reading: {\ttfamily\footnotesize coefficient\allowbreak{}}\par
\noindent Corrected reading: {\ttfamily\footnotesize coefficients\allowbreak{} in\allowbreak{}}\par
\noindent Source: {\ttfamily\footnotesize ADOPT-F7I-VI-095.2\allowbreak{}}\par

\par\medskip\noindent\hypertarget{app-F7I-VI-095.3}{}\textbf{F7I-VI-095.3}\quad printed folio 106.\par
\noindent Printed reading: {\ttfamily\footnotesize '\allowbreak{}}\par
\noindent Corrected reading: \emph{(empty)}\par
\noindent Source: {\ttfamily\footnotesize 7I:229\allowbreak{}}\par

\par\medskip\noindent\hypertarget{app-F7I-VI-095.4}{}\textbf{F7I-VI-095.4}\quad printed folio 106.\par
\noindent Printed reading: \emph{(empty)}\par
\noindent Corrected reading: {\ttfamily\footnotesize '\allowbreak{}}\par
\noindent Source: {\ttfamily\footnotesize 7I:229\allowbreak{}}\par

\par\medskip\noindent\hypertarget{app-F7I-VI-095.5}{}\textbf{F7I-VI-095.5}\quad printed folio 106.\par
\noindent Printed reading: \emph{(empty)}\par
\noindent Corrected reading: {\ttfamily\footnotesize  a\allowbreak{}}\par
\noindent Source: {\ttfamily\footnotesize WITNESS-F7I-VI-095.5\allowbreak{}}\par

\par\medskip\noindent\hypertarget{app-F7I-VI-095.6}{}\textbf{F7I-VI-095.6}\quad printed folio 106.\par
\noindent Printed reading: {\ttfamily\footnotesize 11\allowbreak{}}\par
\noindent Corrected reading: {\ttfamily\footnotesize 8\allowbreak{}}\par
\noindent Source: {\ttfamily\footnotesize ADOPT-F7I-VI-095.6\allowbreak{}}\par

\par\medskip\noindent\hypertarget{app-F7I-VI-096.1}{}\textbf{F7I-VI-096.1}\quad printed folio 106.\par
\noindent Printed reading: {\ttfamily\footnotesize \}\allowbreak{}}\par
\noindent Corrected reading: \emph{(empty)}\par
\noindent Source: {\ttfamily\footnotesize ADOPT-F7I-VI-096.1\allowbreak{}}\par

\par\medskip\noindent\hypertarget{app-F7I-VI-096.2}{}\textbf{F7I-VI-096.2}\quad printed folio 106.\par
\noindent Printed reading: \emph{(empty)}\par
\noindent Corrected reading: {\ttfamily\footnotesize \}\allowbreak{}}\par
\noindent Source: {\ttfamily\footnotesize ADOPT-F7I-VI-096.2\allowbreak{}}\par

\par\medskip\noindent\hypertarget{app-F7I-VI-096.3}{}\textbf{F7I-VI-096.3}\quad printed folio 106.\par
\noindent Printed reading: {\ttfamily\footnotesize \}\allowbreak{}}\par
\noindent Corrected reading: \emph{(empty)}\par
\noindent Source: {\ttfamily\footnotesize ADOPT-F7I-VI-096.3\allowbreak{}}\par

\par\medskip\noindent\hypertarget{app-F7I-VI-096.4}{}\textbf{F7I-VI-096.4}\quad printed folio 106.\par
\noindent Printed reading: \emph{(empty)}\par
\noindent Corrected reading: {\ttfamily\footnotesize \}\allowbreak{}}\par
\noindent Source: {\ttfamily\footnotesize ADOPT-F7I-VI-096.4\allowbreak{}}\par

\par\medskip\noindent\hypertarget{app-F7I-VI-097.1}{}\textbf{F7I-VI-097.1}\quad printed folio 106.\par
\noindent Printed reading: {\ttfamily\footnotesize gives\allowbreak{}}\par
\noindent Corrected reading: {\ttfamily\footnotesize give\allowbreak{}}\par
\noindent Source: {\ttfamily\footnotesize ADOPT-F7I-VI-097.1\allowbreak{}}\par

\par\medskip\noindent\hypertarget{app-F7I-VI-097.2}{}\textbf{F7I-VI-097.2}\quad printed folio 107.\par
\noindent Printed reading: {\ttfamily\footnotesize isomorphism\allowbreak{}}\par
\noindent Corrected reading: {\ttfamily\footnotesize isomorphisms\allowbreak{}}\par
\noindent Source: {\ttfamily\footnotesize ADOPT-F7I-VI-097.2\allowbreak{}}\par

\par\medskip\noindent\hypertarget{app-F7I-VI-097.3}{}\textbf{F7I-VI-097.3}\quad printed folio 107.\par
\noindent Printed reading: {\ttfamily\footnotesize comes\allowbreak{}}\par
\noindent Corrected reading: {\ttfamily\footnotesize come\allowbreak{}}\par
\noindent Source: {\ttfamily\footnotesize ADOPT-F7I-VI-097.3\allowbreak{}}\par

\par\medskip\noindent\hypertarget{app-F7I-VI-097.4}{}\textbf{F7I-VI-097.4}\quad printed folio 107.\par
\noindent Printed reading: {\ttfamily\footnotesize trivial\allowbreak{} }\par
\noindent Corrected reading: \emph{(empty)}\par
\noindent Source: {\ttfamily\footnotesize ADOPT-F7I-VI-097.4\allowbreak{}}\par

\par\medskip\noindent\hypertarget{app-F7I-VI-097.5}{}\textbf{F7I-VI-097.5}\quad printed folio 107.\par
\noindent Printed reading: {\ttfamily\footnotesize chasing\allowbreak{}}\par
\noindent Corrected reading: {\ttfamily\footnotesize chase\allowbreak{}}\par
\noindent Source: {\ttfamily\footnotesize ADOPT-F7I-VI-097.5\allowbreak{}}\par

\par\medskip\noindent\hypertarget{app-F7I-VI-097.6}{}\textbf{F7I-VI-097.6}\quad printed folio 107.\par
\noindent Printed reading: {\ttfamily\footnotesize \}\allowbreak{}}\par
\noindent Corrected reading: \emph{(empty)}\par
\noindent Source: {\ttfamily\footnotesize ADOPT-F7I-VI-097.6\allowbreak{}}\par

\par\medskip\noindent\hypertarget{app-F7I-VI-097.7}{}\textbf{F7I-VI-097.7}\quad printed folio 107.\par
\noindent Printed reading: \emph{(empty)}\par
\noindent Corrected reading: {\ttfamily\footnotesize \}\allowbreak{}}\par
\noindent Source: {\ttfamily\footnotesize ADOPT-F7I-VI-097.7\allowbreak{}}\par

\par\medskip\noindent\hypertarget{app-F7I-VI-097.8}{}\textbf{F7I-VI-097.8}\quad printed folio 107.\par
\noindent Printed reading: {\ttfamily\footnotesize an\allowbreak{}}\par
\noindent Corrected reading: {\ttfamily\footnotesize the\allowbreak{}}\par
\noindent Source: {\ttfamily\footnotesize ADOPT-F7I-VI-097.8\allowbreak{}}\par

\par\medskip\noindent\hypertarget{app-F7I-VI-098.1}{}\textbf{F7I-VI-098.1}\quad printed folio 108.\par
\noindent Printed reading: \emph{(empty)}\par
\noindent Corrected reading: {\ttfamily\footnotesize ,\allowbreak{}}\par
\noindent Source: {\ttfamily\footnotesize 7I:230\allowbreak{}}\par

\par\medskip\noindent\hypertarget{app-F7I-VI-098.2}{}\textbf{F7I-VI-098.2}\quad printed folio 108.\par
\noindent Printed reading: {\ttfamily\footnotesize i\allowbreak{}}\par
\noindent Corrected reading: {\ttfamily\footnotesize j\allowbreak{}}\par
\noindent Source: {\ttfamily\footnotesize 7I:231\allowbreak{}}\par

\par\medskip\noindent\hypertarget{app-F7I-VI-099.1}{}\textbf{F7I-VI-099.1}\quad printed folio 108.\par
\noindent Printed reading: {\ttfamily\footnotesize 0\allowbreak{}}\par
\noindent Corrected reading: {\ttfamily\footnotesize c\allowbreak{}}\par
\noindent Source: {\ttfamily\footnotesize 7I:232\allowbreak{}}\par

\par\medskip\noindent\hypertarget{app-F7I-VI-100.1}{}\textbf{F7I-VI-100.1}\quad printed folio 109.\par
\noindent Printed reading: \emph{(empty)}\par
\noindent Corrected reading: {\ttfamily\footnotesize \textbackslash{}times\allowbreak{} }\par
\noindent Source: {\ttfamily\footnotesize 7I:233\allowbreak{}}\par

\par\medskip\noindent\hypertarget{app-F7I-VI-100.2}{}\textbf{F7I-VI-100.2}\quad printed folio 109.\par
\noindent Printed reading: \emph{(empty)}\par
\noindent Corrected reading: {\ttfamily\footnotesize  \textbackslash{}times\allowbreak{}}\par
\noindent Source: {\ttfamily\footnotesize 7I:233\allowbreak{}}\par

\par\medskip\noindent\hypertarget{app-F7I-VI-101.1}{}\textbf{F7I-VI-101.1}\quad printed folio 110.\par
\noindent Printed reading: {\ttfamily\footnotesize explicitely\allowbreak{}}\par
\noindent Corrected reading: {\ttfamily\footnotesize explicitly\allowbreak{}}\par
\noindent Source: {\ttfamily\footnotesize 7I:234\allowbreak{}}\par

\par\medskip\noindent\hypertarget{app-F7I-VI-101.2}{}\textbf{F7I-VI-101.2}\quad printed folio 110.\par
\noindent Printed reading: {\ttfamily\footnotesize \textbackslash{}qquad\allowbreak{}}\par
\noindent Corrected reading: {\ttfamily\footnotesize =\allowbreak{}}\par
\noindent Source: {\ttfamily\footnotesize 7I:235\allowbreak{}}\par

\par\medskip\noindent\hypertarget{app-F7I-VI-101.3}{}\textbf{F7I-VI-101.3}\quad printed folio 110.\par
\noindent Printed reading: \emph{(empty)}\par
\noindent Corrected reading: {\ttfamily\footnotesize \_0\allowbreak{}}\par
\noindent Source: {\ttfamily\footnotesize 7I:236\allowbreak{}}\par

\par\medskip\noindent\hypertarget{app-F7I-VI-101.4}{}\textbf{F7I-VI-101.4}\quad printed folio 110.\par
\noindent Printed reading: {\ttfamily\footnotesize ji\allowbreak{}}\par
\noindent Corrected reading: {\ttfamily\footnotesize ij\allowbreak{}}\par
\noindent Source: {\ttfamily\footnotesize 7I:237\allowbreak{}}\par

\par\medskip\noindent\hypertarget{app-F7I-VI-102.1}{}\textbf{F7I-VI-102.1}\quad printed folio 111.\par
\noindent Printed reading: \emph{(empty)}\par
\noindent Corrected reading: {\ttfamily\footnotesize \{\allowbreak{}}\par
\noindent Source: {\ttfamily\footnotesize 7I:213\allowbreak{}}\par

\par\medskip\noindent\hypertarget{app-F7I-VI-102.2}{}\textbf{F7I-VI-102.2}\quad printed folio 111.\par
\noindent Printed reading: \emph{(empty)}\par
\noindent Corrected reading: {\ttfamily\footnotesize \_0\}\allowbreak{}}\par
\noindent Source: {\ttfamily\footnotesize 7I:213\allowbreak{}}\par

\par\medskip\noindent\hypertarget{app-F7I-VI-102.3}{}\textbf{F7I-VI-102.3}\quad printed folio 111.\par
\noindent Printed reading: {\ttfamily\footnotesize '\allowbreak{}}\par
\noindent Corrected reading: \emph{(empty)}\par
\noindent Source: {\ttfamily\footnotesize 7I:214\allowbreak{}}\par

\par\medskip\noindent\hypertarget{app-F7I-VI-103.1}{}\textbf{F7I-VI-103.1}\quad printed folio 112.\par
\noindent Printed reading: {\ttfamily\footnotesize standart\allowbreak{}}\par
\noindent Corrected reading: {\ttfamily\footnotesize standard\allowbreak{}}\par
\noindent Source: {\ttfamily\footnotesize 7I:216\allowbreak{}}\par

\par\medskip\noindent\hypertarget{app-F7I-VI-103.2}{}\textbf{F7I-VI-103.2}\quad printed folio 112.\par
\noindent Printed reading: {\ttfamily\footnotesize  \textbackslash{}quad\allowbreak{} }\par
\noindent Corrected reading: {\ttfamily\footnotesize ;\textbackslash{}\allowbreak{} }\par
\noindent Source: {\ttfamily\footnotesize 7I:215\allowbreak{}}\par

\par\medskip\noindent\hypertarget{app-F7I-VI-104.1}{}\textbf{F7I-VI-104.1}\quad printed folio 112.\par
\noindent Printed reading: {\ttfamily\footnotesize a\allowbreak{} }\par
\noindent Corrected reading: \emph{(empty)}\par
\noindent Source: {\ttfamily\footnotesize 7I:218\allowbreak{}}\par

\par\medskip\noindent\hypertarget{app-F7I-VI-105.1}{}\textbf{F7I-VI-105.1}\quad printed folio 113.\par
\noindent Printed reading: {\ttfamily\footnotesize 12\allowbreak{}}\par
\noindent Corrected reading: {\ttfamily\footnotesize 13\allowbreak{}}\par
\noindent Source: {\ttfamily\footnotesize WITNESS-F7I-VI-105.1\allowbreak{}}\par

\par\medskip\noindent\hypertarget{app-F7I-VI-105.2}{}\textbf{F7I-VI-105.2}\quad printed folio 113.\par
\noindent Printed reading: {\ttfamily\footnotesize cotangentcomplex\allowbreak{}}\par
\noindent Corrected reading: {\ttfamily\footnotesize cotangent\allowbreak{} complex\allowbreak{}}\par
\noindent Source: {\ttfamily\footnotesize 7I:219\allowbreak{}}\par

\par\medskip\noindent\hypertarget{app-F7I-VI-105.3}{}\textbf{F7I-VI-105.3}\quad printed folio 113.\par
\noindent Printed reading: \emph{(empty)}\par
\noindent Corrected reading: {\ttfamily\footnotesize \textbackslash{}underline\{\allowbreak{}}\par
\noindent Source: {\ttfamily\footnotesize 7I:220\allowbreak{}}\par

\par\medskip\noindent\hypertarget{app-F7I-VI-105.4}{}\textbf{F7I-VI-105.4}\quad printed folio 113.\par
\noindent Printed reading: \emph{(empty)}\par
\noindent Corrected reading: {\ttfamily\footnotesize \}\allowbreak{}}\par
\noindent Source: {\ttfamily\footnotesize 7I:220\allowbreak{}}\par

\par\medskip\noindent\hypertarget{app-F7I-VI-106.1}{}\textbf{F7I-VI-106.1}\quad printed folio 114.\par
\noindent Printed reading: {\ttfamily\footnotesize fistly\allowbreak{}}\par
\noindent Corrected reading: {\ttfamily\footnotesize first\allowbreak{}}\par
\noindent Source: {\ttfamily\footnotesize WITNESS-F7I-VI-106.1\allowbreak{}}\par

\par\medskip\noindent\hypertarget{app-F7I-VI-107.1}{}\textbf{F7I-VI-107.1}\quad printed folio 114.\par
\noindent Printed reading: {\ttfamily\footnotesize P\allowbreak{}}\par
\noindent Corrected reading: {\ttfamily\footnotesize p\allowbreak{}}\par
\noindent Source: {\ttfamily\footnotesize 7I:222\allowbreak{}}\par

\par\medskip\noindent\hypertarget{app-F7I-VI-107.2}{}\textbf{F7I-VI-107.2}\quad printed folio 114.\par
\noindent Printed reading: {\ttfamily\footnotesize P\allowbreak{}}\par
\noindent Corrected reading: {\ttfamily\footnotesize p\allowbreak{}}\par
\noindent Source: {\ttfamily\footnotesize 7I:223\allowbreak{}}\par

\par\medskip\noindent\hypertarget{app-F7I-VI-107.3}{}\textbf{F7I-VI-107.3}\quad printed folio 114.\par
\noindent Printed reading: {\ttfamily\footnotesize does\allowbreak{}}\par
\noindent Corrected reading: {\ttfamily\footnotesize do\allowbreak{}}\par
\noindent Source: {\ttfamily\footnotesize 7I:224\allowbreak{}}\par

\par\medskip\noindent\hypertarget{app-F7I-VI-107.4}{}\textbf{F7I-VI-107.4}\quad printed folio 114.\par
\noindent Printed reading: \emph{(empty)}\par
\noindent Corrected reading: {\ttfamily\footnotesize the\allowbreak{} }\par
\noindent Source: {\ttfamily\footnotesize WITNESS-F7I-VI-107.4\allowbreak{}}\par

\par\medskip\noindent\hypertarget{app-F7I-VI-107.5}{}\textbf{F7I-VI-107.5}\quad printed folio 114.\par
\noindent Printed reading: \emph{(empty)}\par
\noindent Corrected reading: {\ttfamily\footnotesize  the\allowbreak{}}\par
\noindent Source: {\ttfamily\footnotesize WITNESS-F7I-VI-107.5\allowbreak{}}\par

\par\medskip\noindent\hypertarget{app-F7I-VI-108.1}{}\textbf{F7I-VI-108.1}\quad printed folio 114.\par
\noindent Printed reading: {\ttfamily\footnotesize P\allowbreak{}}\par
\noindent Corrected reading: {\ttfamily\footnotesize p\allowbreak{}}\par
\noindent Source: {\ttfamily\footnotesize 7I:226\allowbreak{}}\par

\par\medskip\noindent\hypertarget{app-F7I-VI-109.1}{}\textbf{F7I-VI-109.1}\quad printed folio 115.\par
\noindent Printed reading: {\ttfamily\footnotesize v\allowbreak{}}\par
\noindent Corrected reading: {\ttfamily\footnotesize \textbackslash{}nu\allowbreak{}}\par
\noindent Source: {\ttfamily\footnotesize ADOPT-F7I-VI-109.1\allowbreak{}}\par

\par\medskip\noindent\hypertarget{app-F7I-VI-109.2}{}\textbf{F7I-VI-109.2}\quad printed folio 115.\par
\noindent Printed reading: {\ttfamily\footnotesize v\allowbreak{}}\par
\noindent Corrected reading: {\ttfamily\footnotesize \textbackslash{}nu\allowbreak{}}\par
\noindent Source: {\ttfamily\footnotesize ADOPT-F7I-VI-109.2\allowbreak{}}\par

\par\medskip\noindent\hypertarget{app-F7I-VI-109.3}{}\textbf{F7I-VI-109.3}\quad printed folio 115.\par
\noindent Printed reading: {\ttfamily\footnotesize v\allowbreak{}}\par
\noindent Corrected reading: {\ttfamily\footnotesize \textbackslash{}nu\allowbreak{}}\par
\noindent Source: {\ttfamily\footnotesize ADOPT-F7I-VI-109.3\allowbreak{}}\par

\par\medskip\noindent\hypertarget{app-F7I-VI-109.4}{}\textbf{F7I-VI-109.4}\quad printed folio 115.\par
\noindent Printed reading: {\ttfamily\footnotesize v\allowbreak{}}\par
\noindent Corrected reading: {\ttfamily\footnotesize \textbackslash{}nu\allowbreak{}}\par
\noindent Source: {\ttfamily\footnotesize ADOPT-F7I-VI-109.4\allowbreak{}}\par

\par\medskip\noindent\hypertarget{app-F7I-VI-109.5}{}\textbf{F7I-VI-109.5}\quad printed folio 115.\par
\noindent Printed reading: {\ttfamily\footnotesize v\allowbreak{}}\par
\noindent Corrected reading: {\ttfamily\footnotesize \textbackslash{}nu\allowbreak{}}\par
\noindent Source: {\ttfamily\footnotesize ADOPT-F7I-VI-109.5\allowbreak{}}\par

\par\medskip\noindent\hypertarget{app-F7I-VI-109.6}{}\textbf{F7I-VI-109.6}\quad printed folio 115.\par
\noindent Printed reading: {\ttfamily\footnotesize v\allowbreak{}}\par
\noindent Corrected reading: {\ttfamily\footnotesize \textbackslash{}nu\allowbreak{}}\par
\noindent Source: {\ttfamily\footnotesize ADOPT-F7I-VI-109.6\allowbreak{}}\par

\par\medskip\noindent\hypertarget{app-F7I-VI-109.7}{}\textbf{F7I-VI-109.7}\quad printed folio 115.\par
\noindent Printed reading: {\ttfamily\footnotesize v\allowbreak{}}\par
\noindent Corrected reading: {\ttfamily\footnotesize \textbackslash{}nu\allowbreak{}}\par
\noindent Source: {\ttfamily\footnotesize ADOPT-F7I-VI-109.7\allowbreak{}}\par

\par\medskip\noindent\hypertarget{app-F7I-VI-109.8}{}\textbf{F7I-VI-109.8}\quad printed folio 115.\par
\noindent Printed reading: {\ttfamily\footnotesize v\allowbreak{}}\par
\noindent Corrected reading: {\ttfamily\footnotesize \textbackslash{}nu\allowbreak{}}\par
\noindent Source: {\ttfamily\footnotesize ADOPT-F7I-VI-109.8\allowbreak{}}\par

\par\medskip\noindent\hypertarget{app-F7I-VI-110.1}{}\textbf{F7I-VI-110.1}\quad printed folio 115.\par
\noindent Printed reading: {\ttfamily\footnotesize x\allowbreak{}}\par
\noindent Corrected reading: {\ttfamily\footnotesize X\allowbreak{}}\par
\noindent Source: {\ttfamily\footnotesize 7I:227\allowbreak{}}\par

\par\medskip\noindent\hypertarget{app-F7I-VI-110.2}{}\textbf{F7I-VI-110.2}\quad printed folio 115.\par
\noindent Printed reading: \emph{(empty)}\par
\noindent Corrected reading: {\ttfamily\footnotesize module\allowbreak{} of\allowbreak{} }\par
\noindent Source: {\ttfamily\footnotesize ADOPT-F7I-VI-110.2\allowbreak{}}\par

\par\medskip\noindent\hypertarget{app-F7I-VI-110.3}{}\textbf{F7I-VI-110.3}\quad printed folio 115.\par
\noindent Printed reading: {\ttfamily\footnotesize -\allowbreak{}}\par
\noindent Corrected reading: \emph{(1 espace(s), 0 retour(s) de ligne dans le témoin TeX)}\par
\noindent Source: {\ttfamily\footnotesize ADOPT-F7I-VI-110.3\allowbreak{}}\par

\par\medskip\noindent\hypertarget{app-F7I-VI-110.4}{}\textbf{F7I-VI-110.4}\quad printed folio 115.\par
\noindent Printed reading: {\ttfamily\footnotesize .\allowbreak{}}\par
\noindent Corrected reading: \emph{(empty)}\par
\noindent Source: {\ttfamily\footnotesize 7I:228\allowbreak{}}\par

\par\medskip\noindent\hypertarget{app-F7I-VI-111.1}{}\textbf{F7I-VI-111.1}\quad printed folio 116.\par
\noindent Printed reading: \emph{(empty)}\par
\noindent Corrected reading: {\ttfamily\footnotesize .\allowbreak{}}\par
\noindent Source: {\ttfamily\footnotesize 7I:239\allowbreak{}}\par

\par\medskip\noindent\hypertarget{app-F7I-VI-112.1}{}\textbf{F7I-VI-112.1}\quad printed folio 116.\par
\noindent Printed reading: \emph{(empty)}\par
\noindent Corrected reading: {\ttfamily\footnotesize r,\allowbreak{} "i"]\allowbreak{} \textbackslash{}arrow[\allowbreak{}}\par
\noindent Source: {\ttfamily\footnotesize WITNESS-F7I-VI-112.1\allowbreak{}}\par

\par\medskip\noindent\hypertarget{app-F7I-VI-112.2}{}\textbf{F7I-VI-112.2}\quad printed folio 116.\par
\noindent Printed reading: {\ttfamily\footnotesize  \textbackslash{}arrow[l,\allowbreak{} "i"]\allowbreak{}}\par
\noindent Corrected reading: \emph{(empty)}\par
\noindent Source: {\ttfamily\footnotesize WITNESS-F7I-VI-112.2\allowbreak{}}\par

\par\medskip\noindent\hypertarget{app-F7I-VI-113.1}{}\textbf{F7I-VI-113.1}\quad printed folio 117.\par
\noindent Printed reading: \emph{(empty)}\par
\noindent Corrected reading: {\ttfamily\footnotesize \textbackslash{}partial\allowbreak{} }\par
\noindent Source: {\ttfamily\footnotesize 7I:240\allowbreak{}}\par

\par\medskip\noindent\hypertarget{app-F7I-VI-113.2}{}\textbf{F7I-VI-113.2}\quad printed folio 117.\par
\noindent Printed reading: \emph{(empty)}\par
\noindent Corrected reading: {\ttfamily\footnotesize \textbackslash{}partial\allowbreak{} }\par
\noindent Source: {\ttfamily\footnotesize 7I:240\allowbreak{}}\par

\par\medskip\noindent\hypertarget{app-F7I-VI-114.1}{}\textbf{F7I-VI-114.1}\quad printed folio 119.\par
\noindent Printed reading: {\ttfamily\footnotesize P\allowbreak{}}\par
\noindent Corrected reading: {\ttfamily\footnotesize p\allowbreak{}}\par
\noindent Source: {\ttfamily\footnotesize FIX-7I-D-0131\allowbreak{}}\par

\par\medskip\noindent\hypertarget{app-F7I-VI-114.2}{}\textbf{F7I-VI-114.2}\quad printed folio 119.\par
\noindent Printed reading: {\ttfamily\footnotesize P\allowbreak{}}\par
\noindent Corrected reading: {\ttfamily\footnotesize p\allowbreak{}}\par
\noindent Source: {\ttfamily\footnotesize FIX-7I-D-0132\allowbreak{}}\par

\par\medskip\noindent\hypertarget{app-F7I-VI-114.3}{}\textbf{F7I-VI-114.3}\quad printed folio 119.\par
\noindent Printed reading: {\ttfamily\footnotesize P\allowbreak{}}\par
\noindent Corrected reading: {\ttfamily\footnotesize p\allowbreak{}}\par
\noindent Source: {\ttfamily\footnotesize FIX-7I-D-0132\allowbreak{}}\par

\par\medskip\noindent\hypertarget{app-F7I-VI-114.4}{}\textbf{F7I-VI-114.4}\quad printed folio 119.\par
\noindent Printed reading: {\ttfamily\footnotesize \textbackslash{}varinjlim\allowbreak{}}\par
\noindent Corrected reading: {\ttfamily\footnotesize \textbackslash{}varprojlim\allowbreak{}}\par
\noindent Source: {\ttfamily\footnotesize 7I:245\allowbreak{}}\par

\par\medskip\noindent\hypertarget{app-F7I-VI-114.5}{}\textbf{F7I-VI-114.5}\quad printed folio 119.\par
\noindent Printed reading: {\ttfamily\footnotesize P\allowbreak{}}\par
\noindent Corrected reading: {\ttfamily\footnotesize p\allowbreak{}}\par
\noindent Source: {\ttfamily\footnotesize FIX-7I-D-0133\allowbreak{}}\par

\par\medskip\noindent\hypertarget{app-F7I-VI-114.6}{}\textbf{F7I-VI-114.6}\quad printed folio 119.\par
\noindent Printed reading: {\ttfamily\footnotesize P\allowbreak{}}\par
\noindent Corrected reading: {\ttfamily\footnotesize p\allowbreak{}}\par
\noindent Source: {\ttfamily\footnotesize FIX-7I-D-0133\allowbreak{}}\par

\par\medskip\noindent\hypertarget{app-F7I-VI-115.1}{}\textbf{F7I-VI-115.1}\quad printed folio 120.\par
\noindent Printed reading: \emph{(empty)}\par
\noindent Corrected reading: {\ttfamily\footnotesize '\allowbreak{}}\par
\noindent Source: {\ttfamily\footnotesize 7I:246\allowbreak{}}\par

\par\medskip\noindent\hypertarget{app-F7I-VI-115.2}{}\textbf{F7I-VI-115.2}\quad printed folio 120.\par
\noindent Printed reading: {\ttfamily\footnotesize is\allowbreak{}}\par
\noindent Corrected reading: {\ttfamily\footnotesize are\allowbreak{}}\par
\noindent Source: {\ttfamily\footnotesize ADOPT-F7I-VI-115.2\allowbreak{}}\par

\par\medskip\noindent\hypertarget{app-F7I-VI-115.3}{}\textbf{F7I-VI-115.3}\quad printed folio 120.\par
\noindent Printed reading: {\ttfamily\footnotesize that\allowbreak{}}\par
\noindent Corrected reading: {\ttfamily\footnotesize those\allowbreak{}}\par
\noindent Source: {\ttfamily\footnotesize ADOPT-F7I-VI-115.3\allowbreak{}}\par

\par\medskip\noindent\hypertarget{app-F7I-VI-116.1}{}\textbf{F7I-VI-116.1}\quad printed folio 120.\par
\noindent Printed reading: {\ttfamily\footnotesize -\allowbreak{}}\par
\noindent Corrected reading: {\ttfamily\footnotesize =\allowbreak{}}\par
\noindent Source: {\ttfamily\footnotesize 7I:247\allowbreak{}}\par

\par\medskip\noindent\hypertarget{app-F7I-VI-117.1}{}\textbf{F7I-VI-117.1}\quad printed folio 121.\par
\noindent Printed reading: {\ttfamily\footnotesize X\allowbreak{}}\par
\noindent Corrected reading: {\ttfamily\footnotesize x\allowbreak{}}\par
\noindent Source: {\ttfamily\footnotesize FIX-7I-D-0142\allowbreak{}}\par

\par\medskip\noindent\hypertarget{app-F7I-VI-117.2}{}\textbf{F7I-VI-117.2}\quad printed folio 121.\par
\noindent Printed reading: {\ttfamily\footnotesize X\allowbreak{}}\par
\noindent Corrected reading: {\ttfamily\footnotesize x\allowbreak{}}\par
\noindent Source: {\ttfamily\footnotesize FIX-7I-D-0143\allowbreak{}}\par

\par\medskip\noindent\hypertarget{app-F7I-VI-117.3}{}\textbf{F7I-VI-117.3}\quad printed folio 121.\par
\noindent Printed reading: \emph{(empty)}\par
\noindent Corrected reading: {\ttfamily\footnotesize \textbackslash{}widehat\{\allowbreak{}}\par
\noindent Source: {\ttfamily\footnotesize 7I:250\allowbreak{}}\par

\par\medskip\noindent\hypertarget{app-F7I-VI-117.4}{}\textbf{F7I-VI-117.4}\quad printed folio 121.\par
\noindent Printed reading: \emph{(empty)}\par
\noindent Corrected reading: {\ttfamily\footnotesize \}\allowbreak{}}\par
\noindent Source: {\ttfamily\footnotesize 7I:250\allowbreak{}}\par

\par\medskip\noindent\hypertarget{app-F7I-VI-117.5}{}\textbf{F7I-VI-117.5}\quad printed folio 121.\par
\noindent Printed reading: {\ttfamily\footnotesize \textbackslash{}lim\allowbreak{}}\par
\noindent Corrected reading: {\ttfamily\footnotesize \textbackslash{}varprojlim\allowbreak{}}\par
\noindent Source: {\ttfamily\footnotesize 7I:251\allowbreak{}}\par

\par\medskip\noindent\hypertarget{app-F7I-VI-117.6}{}\textbf{F7I-VI-117.6}\quad printed folio 121.\par
\noindent Printed reading: {\ttfamily\footnotesize separable\allowbreak{}}\par
\noindent Corrected reading: {\ttfamily\footnotesize separated\allowbreak{}}\par
\noindent Source: {\ttfamily\footnotesize WITNESS-F7I-VI-117.6\allowbreak{}}\par

\par\medskip\noindent\hypertarget{app-F7I-VI-117.7}{}\textbf{F7I-VI-117.7}\quad printed folio 121.\par
\noindent Printed reading: \emph{(empty)}\par
\noindent Corrected reading: {\ttfamily\footnotesize ,y\allowbreak{}}\par
\noindent Source: {\ttfamily\footnotesize WITNESS-F7I-VI-117.7\allowbreak{}}\par

\par\medskip\noindent\hypertarget{app-F7I-VI-117.8}{}\textbf{F7I-VI-117.8}\quad printed folio 122.\par
\noindent Printed reading: \emph{(empty)}\par
\noindent Corrected reading: {\ttfamily\footnotesize \textbackslash{}underline\{\allowbreak{}}\par
\noindent Source: {\ttfamily\footnotesize ADOPT-F7I-VI-117.8\allowbreak{}}\par

\par\medskip\noindent\hypertarget{app-F7I-VI-117.9}{}\textbf{F7I-VI-117.9}\quad printed folio 122.\par
\noindent Printed reading: \emph{(empty)}\par
\noindent Corrected reading: {\ttfamily\footnotesize \}\allowbreak{}}\par
\noindent Source: {\ttfamily\footnotesize ADOPT-F7I-VI-117.9\allowbreak{}}\par

\par\medskip\noindent\hypertarget{app-F7I-VI-117.10}{}\textbf{F7I-VI-117.10}\quad printed folio 122.\par
\noindent Printed reading: \emph{(empty)}\par
\noindent Corrected reading: {\ttfamily\footnotesize \textbackslash{}underline\{\allowbreak{}}\par
\noindent Source: {\ttfamily\footnotesize WITNESS-F7I-VI-117.1\allowbreak{}0\allowbreak{}}\par

\par\medskip\noindent\hypertarget{app-F7I-VI-117.11}{}\textbf{F7I-VI-117.11}\quad printed folio 122.\par
\noindent Printed reading: \emph{(empty)}\par
\noindent Corrected reading: {\ttfamily\footnotesize \}\allowbreak{}}\par
\noindent Source: {\ttfamily\footnotesize WITNESS-F7I-VI-117.1\allowbreak{}1\allowbreak{}}\par

\par\medskip\noindent\hypertarget{app-F7I-VI-117.12}{}\textbf{F7I-VI-117.12}\quad printed folio 122.\par
\noindent Printed reading: \emph{(empty)}\par
\noindent Corrected reading: {\ttfamily\footnotesize \textbackslash{}widehat\{\allowbreak{}}\par
\noindent Source: {\ttfamily\footnotesize 7I:253\allowbreak{}}\par

\par\medskip\noindent\hypertarget{app-F7I-VI-117.13}{}\textbf{F7I-VI-117.13}\quad printed folio 122.\par
\noindent Printed reading: \emph{(empty)}\par
\noindent Corrected reading: {\ttfamily\footnotesize \}\allowbreak{}}\par
\noindent Source: {\ttfamily\footnotesize WITNESS-F7I-VI-117.1\allowbreak{}3\allowbreak{}}\par

\par\medskip\noindent\hypertarget{app-F7I-VI-118.1}{}\textbf{F7I-VI-118.1}\quad printed folio 122.\par
\noindent Printed reading: {\ttfamily\footnotesize sheaves\allowbreak{}}\par
\noindent Corrected reading: {\ttfamily\footnotesize sheaf\allowbreak{}}\par
\noindent Source: {\ttfamily\footnotesize ADOPT-F7I-VI-118.1\allowbreak{}}\par

\par\medskip\noindent\hypertarget{app-F7I-VI-118.2}{}\textbf{F7I-VI-118.2}\quad printed folio 122.\par
\noindent Printed reading: \emph{(empty)}\par
\noindent Corrected reading: {\ttfamily\footnotesize \textbackslash{}underline\{\allowbreak{}}\par
\noindent Source: {\ttfamily\footnotesize ADOPT-F7I-VI-118.2\allowbreak{}}\par

\par\medskip\noindent\hypertarget{app-F7I-VI-118.3}{}\textbf{F7I-VI-118.3}\quad printed folio 122.\par
\noindent Printed reading: \emph{(empty)}\par
\noindent Corrected reading: {\ttfamily\footnotesize \}\allowbreak{}}\par
\noindent Source: {\ttfamily\footnotesize ADOPT-F7I-VI-118.3\allowbreak{}}\par

\par\medskip\noindent\hypertarget{app-F7I-VI-119.1}{}\textbf{F7I-VI-119.1}\quad printed folio 123.\par
\noindent Printed reading: {\ttfamily\footnotesize commute\allowbreak{}}\par
\noindent Corrected reading: {\ttfamily\footnotesize commutes\allowbreak{}}\par
\noindent Source: {\ttfamily\footnotesize 7I:254\allowbreak{}}\par

\par\medskip\noindent\hypertarget{app-F7I-VI-120.1}{}\textbf{F7I-VI-120.1}\quad printed folio 123.\par
\noindent Printed reading: {\ttfamily\footnotesize O\allowbreak{}}\par
\noindent Corrected reading: {\ttfamily\footnotesize I\allowbreak{}}\par
\noindent Source: {\ttfamily\footnotesize 7I:255\allowbreak{}}\par

\par\medskip\noindent\hypertarget{app-F7I-VI-121.1}{}\textbf{F7I-VI-121.1}\quad printed folio 124.\par
\noindent Printed reading: {\ttfamily\footnotesize -\allowbreak{}}\par
\noindent Corrected reading: {\ttfamily\footnotesize \textbackslash{}simeq\allowbreak{}}\par
\noindent Source: {\ttfamily\footnotesize 7I:256\allowbreak{}}\par

\par\medskip\noindent\hypertarget{app-F7I-VI-122.1}{}\textbf{F7I-VI-122.1}\quad printed folio 124.\par
\noindent Printed reading: {\ttfamily\footnotesize X\allowbreak{}}\par
\noindent Corrected reading: {\ttfamily\footnotesize x\allowbreak{}}\par
\noindent Source: {\ttfamily\footnotesize 7I:257\allowbreak{}}\par

\par\medskip\noindent\hypertarget{app-F7I-VI-122.2}{}\textbf{F7I-VI-122.2}\quad printed folio 124.\par
\noindent Printed reading: {\ttfamily\footnotesize X\allowbreak{}}\par
\noindent Corrected reading: {\ttfamily\footnotesize x\allowbreak{}}\par
\noindent Source: {\ttfamily\footnotesize 7I:257\allowbreak{}}\par

\par\medskip\noindent\hypertarget{app-F7I-VI-122.3}{}\textbf{F7I-VI-122.3}\quad printed folio 125.\par
\noindent Printed reading: {\ttfamily\footnotesize u\allowbreak{}}\par
\noindent Corrected reading: {\ttfamily\footnotesize i\allowbreak{}}\par
\noindent Source: {\ttfamily\footnotesize 7I:258\allowbreak{}}\par

\par\medskip\noindent\hypertarget{app-F7I-VI-123.1}{}\textbf{F7I-VI-123.1}\quad printed folio 125.\par
\noindent Printed reading: {\ttfamily\footnotesize J\allowbreak{}}\par
\noindent Corrected reading: {\ttfamily\footnotesize I\allowbreak{}}\par
\noindent Source: {\ttfamily\footnotesize 7I:259\allowbreak{}}\par

\par\medskip\noindent\hypertarget{app-F7I-VI-124.1}{}\textbf{F7I-VI-124.1}\quad printed folio 125.\par
\noindent Printed reading: {\ttfamily\footnotesize \_0\allowbreak{}}\par
\noindent Corrected reading: \emph{(empty)}\par
\noindent Source: {\ttfamily\footnotesize ADOPT-F7I-VI-124.1\allowbreak{}}\par

\par\medskip\noindent\hypertarget{app-F7I-VI-125.1}{}\textbf{F7I-VI-125.1}\quad printed folio 125.\par
\noindent Printed reading: {\ttfamily\footnotesize df\allowbreak{}}\par
\noindent Corrected reading: {\ttfamily\footnotesize dF\allowbreak{}}\par
\noindent Source: {\ttfamily\footnotesize WITNESS-F7I-VI-125.1\allowbreak{}}\par

\par\medskip\noindent\hypertarget{app-F7I-VI-125.2}{}\textbf{F7I-VI-125.2}\quad printed folio 125.\par
\noindent Printed reading: {\ttfamily\footnotesize \textbackslash{}relbar\textbackslash{}joinrel\textbackslash{}relb\allowbreak{}ar\textbackslash{}joinrel\textbackslash{}relbar\allowbreak{}}\par
\noindent Corrected reading: {\ttfamily\footnotesize \textbackslash{}longrightarrow\allowbreak{}}\par
\noindent Source: {\ttfamily\footnotesize FIX-7I-D-0163\allowbreak{}}\par

\par\medskip\noindent\hypertarget{app-F7I-VI-125.3}{}\textbf{F7I-VI-125.3}\quad printed folio 125.\par
\noindent Printed reading: {\ttfamily\footnotesize \textbackslash{}relbar\textbackslash{}joinrel\textbackslash{}relb\allowbreak{}ar\textbackslash{}joinrel\textbackslash{}relbar\textbackslash{}jo\allowbreak{}inrel\textbackslash{}relbar\textbackslash{}joinrel\allowbreak{}\textbackslash{}relbar\allowbreak{}}\par
\noindent Corrected reading: {\ttfamily\footnotesize \textbackslash{}xrightarrow\{\textbackslash{}hspace\allowbreak{}\{3em\}\}\allowbreak{}}\par
\noindent Source: {\ttfamily\footnotesize FIX-7I-D-0164\allowbreak{}}\par

\par\medskip\noindent\hypertarget{app-F7I-VI-126.1}{}\textbf{F7I-VI-126.1}\quad printed folio 126.\par
\noindent Printed reading: {\ttfamily\footnotesize df\allowbreak{}}\par
\noindent Corrected reading: {\ttfamily\footnotesize dF\allowbreak{}}\par
\noindent Source: {\ttfamily\footnotesize ADOPT-F7I-VI-126.1\allowbreak{}}\par

\par\medskip\noindent\hypertarget{app-F7I-VI-126.2}{}\textbf{F7I-VI-126.2}\quad printed folio 126.\par
\noindent Printed reading: {\ttfamily\footnotesize f\allowbreak{}}\par
\noindent Corrected reading: {\ttfamily\footnotesize F\allowbreak{}}\par
\noindent Source: {\ttfamily\footnotesize ADOPT-F7I-VI-126.2\allowbreak{}}\par

\par\medskip\noindent\hypertarget{app-F7I-VI-126.3}{}\textbf{F7I-VI-126.3}\quad printed folio 126.\par
\noindent Printed reading: {\ttfamily\footnotesize f\allowbreak{}}\par
\noindent Corrected reading: {\ttfamily\footnotesize F\allowbreak{}}\par
\noindent Source: {\ttfamily\footnotesize ADOPT-F7I-VI-126.3\allowbreak{}}\par

\par\medskip\noindent\hypertarget{app-F7I-VI-126.4}{}\textbf{F7I-VI-126.4}\quad printed folio 126.\par
\noindent Printed reading: \emph{(empty)}\par
\noindent Corrected reading: {\ttfamily\footnotesize /I\_\{\textbackslash{}hat\{A\}\}(\textbackslash{}hat\{A\}\allowbreak{}|R)\allowbreak{}}\par
\noindent Source: {\ttfamily\footnotesize WITNESS-F7I-VI-126.4\allowbreak{}}\par

\par\medskip\noindent\hypertarget{app-F7I-VI-127.1}{}\textbf{F7I-VI-127.1}\quad printed folio 126.\par
\noindent Printed reading: {\ttfamily\footnotesize R\allowbreak{}}\par
\noindent Corrected reading: {\ttfamily\footnotesize F\allowbreak{}}\par
\noindent Source: {\ttfamily\footnotesize 7I:268\allowbreak{}}\par

\par\medskip\noindent\hypertarget{app-F7I-VI-127.2}{}\textbf{F7I-VI-127.2}\quad printed folio 127.\par
\noindent Printed reading: {\ttfamily\footnotesize ab\allowbreak{}}\par
\noindent Corrected reading: {\ttfamily\footnotesize ub\allowbreak{}}\par
\noindent Source: {\ttfamily\footnotesize 7I:269\allowbreak{}}\par

\par\medskip\noindent\hypertarget{app-F7I-VI-127.3}{}\textbf{F7I-VI-127.3}\quad printed folio 127.\par
\noindent Printed reading: {\ttfamily\footnotesize Since\allowbreak{}}\par
\noindent Corrected reading: {\ttfamily\footnotesize We\allowbreak{} construct\allowbreak{} a\allowbreak{} matrix\allowbreak{}}\par
\noindent Source: {\ttfamily\footnotesize WITNESS-F7I-VI-127.3\allowbreak{}}\par

\par\medskip\noindent\hypertarget{app-F7I-VI-127.4}{}\textbf{F7I-VI-127.4}\quad printed folio 127.\par
\noindent Printed reading: {\ttfamily\footnotesize k\allowbreak{} =\allowbreak{} k\textasciicircum{}2\allowbreak{}}\par
\noindent Corrected reading: {\ttfamily\footnotesize T\textbackslash{}in\textbackslash{}operatorname\{GL\allowbreak{}\}\_n(R)\allowbreak{}}\par
\noindent Source: {\ttfamily\footnotesize WITNESS-F7I-VI-127.4\allowbreak{}}\par

\par\medskip\noindent\hypertarget{app-F7I-VI-127.5}{}\textbf{F7I-VI-127.5}\quad printed folio 127.\par
\noindent Printed reading: {\ttfamily\footnotesize and\allowbreak{}}\par
\noindent Corrected reading: {\ttfamily\footnotesize such\allowbreak{} that\allowbreak{} \$F(TX)=uF(X)\$.\textbackslash{}footn\allowbreak{}ote\{Editorial\allowbreak{} proof\allowbreak{} repair:\allowbreak{} the\allowbreak{} printed\allowbreak{} inference\allowbreak{} that\allowbreak{} every\allowbreak{} unit\allowbreak{} of\allowbreak{}}\par
\noindent Source: {\ttfamily\footnotesize WITNESS-F7I-VI-127.5\allowbreak{}}\par

\par\medskip\noindent\hypertarget{app-F7I-VI-127.6}{}\textbf{F7I-VI-127.6}\quad printed folio 127.\par
\noindent Printed reading: {\ttfamily\footnotesize is\allowbreak{}}\par
\noindent Corrected reading: {\ttfamily\footnotesize has\allowbreak{}}\par
\noindent Source: {\ttfamily\footnotesize ADOPT-F7I-VI-127.6\allowbreak{}}\par

\par\medskip\noindent\hypertarget{app-F7I-VI-127.7}{}\textbf{F7I-VI-127.7}\quad printed folio 127.\par
\noindent Printed reading: {\ttfamily\footnotesize complete\allowbreak{}}\par
\noindent Corrected reading: {\ttfamily\footnotesize a\allowbreak{} square\allowbreak{} root\allowbreak{} need\allowbreak{} not\allowbreak{} hold\allowbreak{} in\allowbreak{} characteristic\allowbreak{} \$2\$.\allowbreak{} The\allowbreak{} following\allowbreak{} argument\allowbreak{} retains\allowbreak{} the\allowbreak{} original\allowbreak{} hypotheses.\}\allowbreak{} Write\allowbreak{} \$\textbackslash{}mathfrak\allowbreak{} m=\textbackslash{}mathfrak\allowbreak{} m\_R\$\allowbreak{}}\par
\noindent Source: {\ttfamily\footnotesize ADOPT-F7I-VI-127.7\allowbreak{}}\par

\par\medskip\noindent\hypertarget{app-F7I-VI-127.8}{}\textbf{F7I-VI-127.8}\quad printed folio 127.\par
\noindent Printed reading: {\ttfamily\footnotesize we\allowbreak{}}\par
\noindent Corrected reading: {\ttfamily\footnotesize and\allowbreak{}}\par
\noindent Source: {\ttfamily\footnotesize ADOPT-F7I-VI-127.8\allowbreak{}}\par

\par\medskip\noindent\hypertarget{app-F7I-VI-127.9}{}\textbf{F7I-VI-127.9}\quad printed folio 127.\par
\noindent Printed reading: {\ttfamily\footnotesize can\allowbreak{}}\par
\noindent Corrected reading: {\ttfamily\footnotesize interpret\allowbreak{}}\par
\noindent Source: {\ttfamily\footnotesize ADOPT-F7I-VI-127.9\allowbreak{}}\par

\par\medskip\noindent\hypertarget{app-F7I-VI-127.10}{}\textbf{F7I-VI-127.10}\quad printed folio 127.\par
\noindent Printed reading: {\ttfamily\footnotesize write\allowbreak{}}\par
\noindent Corrected reading: {\ttfamily\footnotesize polynomial\allowbreak{} congruences\allowbreak{} coefficientwise.\allowbreak{} Since\allowbreak{}}\par
\noindent Source: {\ttfamily\footnotesize ADOPT-F7I-VI-127.10\allowbreak{}}\par

\par\medskip\noindent\hypertarget{app-F7I-VI-127.11}{}\textbf{F7I-VI-127.11}\quad printed folio 127.\par
\noindent Printed reading: {\ttfamily\footnotesize u\allowbreak{} }\par
\noindent Corrected reading: {\ttfamily\footnotesize k\allowbreak{}}\par
\noindent Source: {\ttfamily\footnotesize WITNESS-F7I-VI-127.1\allowbreak{}1\allowbreak{}}\par

\par\medskip\noindent\hypertarget{app-F7I-VI-127.12}{}\textbf{F7I-VI-127.12}\quad printed folio 127.\par
\noindent Printed reading: {\ttfamily\footnotesize  v\allowbreak{}}\par
\noindent Corrected reading: {\ttfamily\footnotesize k\allowbreak{}}\par
\noindent Source: {\ttfamily\footnotesize WITNESS-F7I-VI-127.1\allowbreak{}2\allowbreak{}}\par

\par\medskip\noindent\hypertarget{app-F7I-VI-127.13}{}\textbf{F7I-VI-127.13}\quad printed folio 127.\par
\noindent Printed reading: \emph{(empty)}\par
\noindent Corrected reading: {\ttfamily\footnotesize ,\allowbreak{}}\par
\noindent Source: {\ttfamily\footnotesize ADOPT-F7I-VI-127.13\allowbreak{}}\par

\par\medskip\noindent\hypertarget{app-F7I-VI-127.14}{}\textbf{F7I-VI-127.14}\quad printed folio 127.\par
\noindent Printed reading: {\ttfamily\footnotesize for\allowbreak{} some\allowbreak{} unit\allowbreak{}}\par
\noindent Corrected reading: {\ttfamily\footnotesize choose\allowbreak{}}\par
\noindent Source: {\ttfamily\footnotesize ADOPT-F7I-VI-127.14\allowbreak{}}\par

\par\medskip\noindent\hypertarget{app-F7I-VI-127.15}{}\textbf{F7I-VI-127.15}\quad printed folio 127.\par
\noindent Printed reading: {\ttfamily\footnotesize v\allowbreak{}}\par
\noindent Corrected reading: {\ttfamily\footnotesize c\textbackslash{}in\allowbreak{} R\textasciicircum{}\textbackslash{}times\allowbreak{}}\par
\noindent Source: {\ttfamily\footnotesize ADOPT-F7I-VI-127.15\allowbreak{}}\par

\par\medskip\noindent\hypertarget{app-F7I-VI-127.16}{}\textbf{F7I-VI-127.16}\quad printed folio 127.\par
\noindent Printed reading: {\ttfamily\footnotesize in\allowbreak{}}\par
\noindent Corrected reading: {\ttfamily\footnotesize whose\allowbreak{} residue\allowbreak{} satisfies\allowbreak{}}\par
\noindent Source: {\ttfamily\footnotesize ADOPT-F7I-VI-127.16\allowbreak{}}\par

\par\medskip\noindent\hypertarget{app-F7I-VI-127.17}{}\textbf{F7I-VI-127.17}\quad printed folio 127.\par
\noindent Printed reading: {\ttfamily\footnotesize R\allowbreak{}}\par
\noindent Corrected reading: {\ttfamily\footnotesize \textbackslash{}bar\allowbreak{} c\textasciicircum{}2=\textbackslash{}bar\allowbreak{} u\allowbreak{}}\par
\noindent Source: {\ttfamily\footnotesize ADOPT-F7I-VI-127.17\allowbreak{}}\par

\par\medskip\noindent\hypertarget{app-F7I-VI-127.18}{}\textbf{F7I-VI-127.18}\quad printed folio 127.\par
\noindent Printed reading: \emph{(1 espace(s), 0 retour(s) de ligne dans le témoin TeX)}\par
\noindent Corrected reading: \emph{(empty)}\par
\noindent Source: {\ttfamily\footnotesize ADOPT-F7I-VI-127.18\allowbreak{}}\par

\par\medskip\noindent\hypertarget{app-F7I-VI-127.19}{}\textbf{F7I-VI-127.19}\quad printed folio 127.\par
\noindent Printed reading: {\ttfamily\footnotesize i.e.,\allowbreak{}}\par
\noindent Corrected reading: {\ttfamily\footnotesize and\allowbreak{} put\allowbreak{}}\par
\noindent Source: {\ttfamily\footnotesize ADOPT-F7I-VI-127.19\allowbreak{}}\par

\par\medskip\noindent\hypertarget{app-F7I-VI-127.20}{}\textbf{F7I-VI-127.20}\quad printed folio 127.\par
\noindent Printed reading: {\ttfamily\footnotesize a\allowbreak{} }\par
\noindent Corrected reading: {\ttfamily\footnotesize T\_1\allowbreak{}}\par
\noindent Source: {\ttfamily\footnotesize ADOPT-F7I-VI-127.20\allowbreak{}}\par

\par\medskip\noindent\hypertarget{app-F7I-VI-127.21}{}\textbf{F7I-VI-127.21}\quad printed folio 127.\par
\noindent Printed reading: {\ttfamily\footnotesize  v\textasciicircum{}2b\allowbreak{}}\par
\noindent Corrected reading: {\ttfamily\footnotesize cI\_n\allowbreak{}}\par
\noindent Source: {\ttfamily\footnotesize ADOPT-F7I-VI-127.21\allowbreak{}}\par

\par\medskip\noindent\hypertarget{app-F7I-VI-127.22}{}\textbf{F7I-VI-127.22}\quad printed folio 127.\par
\noindent Printed reading: \emph{(1 espace(s), 0 retour(s) de ligne dans le témoin TeX)}\par
\noindent Corrected reading: \emph{(empty)}\par
\noindent Source: {\ttfamily\footnotesize ADOPT-F7I-VI-127.22\allowbreak{}}\par

\par\medskip\noindent\hypertarget{app-F7I-VI-127.23}{}\textbf{F7I-VI-127.23}\quad printed folio 127.\par
\noindent Printed reading: \emph{(empty)}\par
\noindent Corrected reading: {\ttfamily\footnotesize F\allowbreak{}}\par
\noindent Source: {\ttfamily\footnotesize ADOPT-F7I-VI-127.23\allowbreak{}}\par

\par\medskip\noindent\hypertarget{app-F7I-VI-127.24}{}\textbf{F7I-VI-127.24}\quad printed folio 127.\par
\noindent Printed reading: {\ttfamily\footnotesize F\allowbreak{}}\par
\noindent Corrected reading: {\ttfamily\footnotesize T\_1X)\textbackslash{}equiv\allowbreak{}}\par
\noindent Source: {\ttfamily\footnotesize ADOPT-F7I-VI-127.24\allowbreak{}}\par

\par\medskip\noindent\hypertarget{app-F7I-VI-127.25}{}\textbf{F7I-VI-127.25}\quad printed folio 127.\par
\noindent Printed reading: {\ttfamily\footnotesize -\allowbreak{}}\par
\noindent Corrected reading: {\ttfamily\footnotesize uF(X)\textbackslash{}pmod\{\textbackslash{}mathfrak\allowbreak{}}\par
\noindent Source: {\ttfamily\footnotesize ADOPT-F7I-VI-127.25\allowbreak{}}\par

\par\medskip\noindent\hypertarget{app-F7I-VI-127.26}{}\textbf{F7I-VI-127.26}\quad printed folio 127.\par
\noindent Printed reading: {\ttfamily\footnotesize b)\allowbreak{}}\par
\noindent Corrected reading: {\ttfamily\footnotesize m\}\$.\allowbreak{}}\par
\noindent Source: {\ttfamily\footnotesize ADOPT-F7I-VI-127.26\allowbreak{}}\par

\par\medskip\noindent\hypertarget{app-F7I-VI-127.27}{}\textbf{F7I-VI-127.27}\quad printed folio 127.\par
\noindent Printed reading: {\ttfamily\footnotesize =\allowbreak{}}\par
\noindent Corrected reading: {\ttfamily\footnotesize By\allowbreak{}}\par
\noindent Source: {\ttfamily\footnotesize ADOPT-F7I-VI-127.27\allowbreak{}}\par

\par\medskip\noindent\hypertarget{app-F7I-VI-127.28}{}\textbf{F7I-VI-127.28}\quad printed folio 127.\par
\noindent Printed reading: {\ttfamily\footnotesize (v\textasciicircum{}2F\allowbreak{}}\par
\noindent Corrected reading: {\ttfamily\footnotesize the\allowbreak{}}\par
\noindent Source: {\ttfamily\footnotesize ADOPT-F7I-VI-127.28\allowbreak{}}\par

\par\medskip\noindent\hypertarget{app-F7I-VI-127.29}{}\textbf{F7I-VI-127.29}\quad printed folio 127.\par
\noindent Printed reading: {\ttfamily\footnotesize -\allowbreak{}}\par
\noindent Corrected reading: {\ttfamily\footnotesize nondegeneracy\allowbreak{}}\par
\noindent Source: {\ttfamily\footnotesize ADOPT-F7I-VI-127.29\allowbreak{}}\par

\par\medskip\noindent\hypertarget{app-F7I-VI-127.30}{}\textbf{F7I-VI-127.30}\quad printed folio 127.\par
\noindent Printed reading: {\ttfamily\footnotesize a)\$\allowbreak{}}\par
\noindent Corrected reading: {\ttfamily\footnotesize condition\allowbreak{}}\par
\noindent Source: {\ttfamily\footnotesize ADOPT-F7I-VI-127.30\allowbreak{}}\par

\par\medskip\noindent\hypertarget{app-F7I-VI-127.31}{}\textbf{F7I-VI-127.31}\quad printed folio 127.\par
\noindent Printed reading: \emph{(empty)}\par
\noindent Corrected reading: {\ttfamily\footnotesize in\allowbreak{} Lemma\allowbreak{} 6.2\allowbreak{}}\par
\noindent Source: {\ttfamily\footnotesize ADOPT-F7I-VI-127.31\allowbreak{}}\par

\par\medskip\noindent\hypertarget{app-F7I-VI-127.32}{}\textbf{F7I-VI-127.32}\quad printed folio 127.\par
\noindent Printed reading: {\ttfamily\footnotesize so\allowbreak{}}\par
\noindent Corrected reading: {\ttfamily\footnotesize the\allowbreak{}}\par
\noindent Source: {\ttfamily\footnotesize ADOPT-F7I-VI-127.32\allowbreak{}}\par

\par\medskip\noindent\hypertarget{app-F7I-VI-127.33}{}\textbf{F7I-VI-127.33}\quad printed folio 127.\par
\noindent Printed reading: {\ttfamily\footnotesize that\allowbreak{}}\par
\noindent Corrected reading: {\ttfamily\footnotesize coefficient\allowbreak{} matrix\allowbreak{} of\allowbreak{} the\allowbreak{} linear\allowbreak{} forms\allowbreak{}}\par
\noindent Source: {\ttfamily\footnotesize ADOPT-F7I-VI-127.33\allowbreak{}}\par

\par\medskip\noindent\hypertarget{app-F7I-VI-127.34}{}\textbf{F7I-VI-127.34}\quad printed folio 127.\par
\noindent Printed reading: \emph{(empty)}\par
\noindent Corrected reading: {\ttfamily\footnotesize \textbackslash{}partial\allowbreak{} F/\textbackslash{}partial\allowbreak{} }\par
\noindent Source: {\ttfamily\footnotesize ADOPT-F7I-VI-127.34\allowbreak{}}\par

\par\medskip\noindent\hypertarget{app-F7I-VI-127.35}{}\textbf{F7I-VI-127.35}\quad printed folio 127.\par
\noindent Printed reading: {\ttfamily\footnotesize \_i\allowbreak{} \textbackslash{}to\allowbreak{} vX\allowbreak{}}\par
\noindent Corrected reading: \emph{(empty)}\par
\noindent Source: {\ttfamily\footnotesize ADOPT-F7I-VI-127.35\allowbreak{}}\par

\par\medskip\noindent\hypertarget{app-F7I-VI-127.36}{}\textbf{F7I-VI-127.36}\quad printed folio 127.\par
\noindent Printed reading: {\ttfamily\footnotesize establishes\allowbreak{}}\par
\noindent Corrected reading: {\ttfamily\footnotesize is\allowbreak{} invertible\allowbreak{} modulo\allowbreak{} \$\textbackslash{}mathfrak\allowbreak{} m\$,\allowbreak{} hence\allowbreak{} over\allowbreak{} \$R\$.\allowbreak{} Suppose\allowbreak{} \$T\_\textbackslash{}nu\textbackslash{}in\textbackslash{}operatorna\allowbreak{}me\{GL\}\_n(R)\$\allowbreak{} satisfies\allowbreak{} \$F(T\_\textbackslash{}nu\allowbreak{} X)\textbackslash{}equiv\allowbreak{} uF(X)\textbackslash{}pmod\{\textbackslash{}mathfrak\allowbreak{} m\textasciicircum{}\textbackslash{}nu\}\$,\allowbreak{} with\allowbreak{} \$\textbackslash{}nu\textbackslash{}geq1\$.\allowbreak{} The\allowbreak{} forms\allowbreak{} \$(\textbackslash{}partial\allowbreak{} F/\textbackslash{}partial\allowbreak{} X\_i)(T\_\textbackslash{}nu\allowbreak{} X)\$\allowbreak{} form\allowbreak{}}\par
\noindent Source: {\ttfamily\footnotesize ADOPT-F7I-VI-127.36\allowbreak{}}\par

\par\medskip\noindent\hypertarget{app-F7I-VI-127.37}{}\textbf{F7I-VI-127.37}\quad printed folio 127.\par
\noindent Printed reading: {\ttfamily\footnotesize algebra\allowbreak{}}\par
\noindent Corrected reading: {\ttfamily\footnotesize basis\allowbreak{} of\allowbreak{} the\allowbreak{} linear\allowbreak{} forms.\allowbreak{} Consequently,\allowbreak{} there\allowbreak{} are\allowbreak{} linear\allowbreak{} forms\allowbreak{} \$L\_i\$\allowbreak{} with\allowbreak{} coefficients\allowbreak{} in\allowbreak{} \$\textbackslash{}mathfrak\allowbreak{} m\textasciicircum{}\textbackslash{}nu\$\allowbreak{} such\allowbreak{} that\allowbreak{} \textbackslash{}[\allowbreak{} E\_\textbackslash{}nu(X):=uF(X)-F(T\_\allowbreak{}\textbackslash{}nu\allowbreak{} X)\allowbreak{} =\textbackslash{}sum\_i\allowbreak{} L\_i(X)\textbackslash{}frac\{\textbackslash{}partial\allowbreak{} F\}\{\textbackslash{}partial\allowbreak{} X\_i\}(T\_\textbackslash{}nu\allowbreak{} X).\allowbreak{} \textbackslash{}]\allowbreak{} Indeed,\allowbreak{} every\allowbreak{} quadratic\allowbreak{} monomial\allowbreak{} is\allowbreak{} a\allowbreak{} product\allowbreak{} of\allowbreak{} two\allowbreak{} linear\allowbreak{} forms,\allowbreak{} and\allowbreak{} expressing\allowbreak{} one\allowbreak{} factor\allowbreak{} in\allowbreak{} this\allowbreak{} basis\allowbreak{} preserves\allowbreak{} its\allowbreak{} coefficient\allowbreak{} ideal.\allowbreak{} Define\allowbreak{} \$H\_\textbackslash{}nu\$\allowbreak{} by\allowbreak{} \$H\_\textbackslash{}nu\allowbreak{} X=(L\_i(X))\_i\$,\allowbreak{} and\allowbreak{} set\allowbreak{} \$T\_\{\textbackslash{}nu+1\}=T\_\textbackslash{}nu+H\_\textbackslash{}\allowbreak{}nu\$.\allowbreak{} The\allowbreak{} quadratic\allowbreak{} identity\allowbreak{} gives\allowbreak{} \textbackslash{}[\allowbreak{} F(T\_\{\textbackslash{}nu+1\}X)=F(T\_\textbackslash{}n\allowbreak{}u\allowbreak{} X)+E\_\textbackslash{}nu(X)+F(H\_\textbackslash{}nu\allowbreak{} X)\allowbreak{} \textbackslash{}equiv\allowbreak{} uF(X)\textbackslash{}pmod\{\textbackslash{}mathfrak\allowbreak{} m\textasciicircum{}\{\textbackslash{}nu+1\}\},\allowbreak{} \textbackslash{}]\allowbreak{} since\allowbreak{} the\allowbreak{} coefficients\allowbreak{} of\allowbreak{} \$F(H\_\textbackslash{}nu\allowbreak{} X)\$\allowbreak{} lie\allowbreak{} in\allowbreak{} \$\textbackslash{}mathfrak\allowbreak{} m\textasciicircum{}\{2\textbackslash{}nu\}\textbackslash{}subseteq\textbackslash{}ma\allowbreak{}thfrak\allowbreak{} m\textasciicircum{}\{\textbackslash{}nu+1\}\$.\allowbreak{} Moreover,\allowbreak{} \$T\_\{\textbackslash{}nu+1\}\textbackslash{}equiv\allowbreak{} T\_\textbackslash{}nu\textbackslash{}pmod\{\textbackslash{}mathfrak\allowbreak{} m\textasciicircum{}\textbackslash{}nu\}\$.\allowbreak{} Completeness\allowbreak{} and\allowbreak{} separatedness\allowbreak{} therefore\allowbreak{} give\allowbreak{} a\allowbreak{} limit\allowbreak{} \$T\textbackslash{}equiv\allowbreak{} cI\_n\textbackslash{}pmod\{\textbackslash{}mathfrak\allowbreak{} m\}\$\allowbreak{} with\allowbreak{} \$F(TX)=uF(X)\$;\allowbreak{} in\allowbreak{} particular,\allowbreak{} \$T\$\allowbreak{} is\allowbreak{} invertible.\allowbreak{} The\allowbreak{} substitution\allowbreak{} \$X\textbackslash{}mapsto\allowbreak{} TX\$\allowbreak{} sends\allowbreak{} \$F-a\$\allowbreak{} to\allowbreak{} \$u(F-b)\$\allowbreak{} and\allowbreak{} induces\allowbreak{} the\allowbreak{} required\allowbreak{}}\par
\noindent Source: {\ttfamily\footnotesize WITNESS-F7I-VI-127.3\allowbreak{}7\allowbreak{}}\par

\par\medskip\noindent\hypertarget{app-F7I-VI-127.38}{}\textbf{F7I-VI-127.38}\quad printed folio 127.\par
\noindent Printed reading: {\ttfamily\footnotesize \$R[[X\_1,\textbackslash{}ldots,X\_n]]\allowbreak{}/(F\allowbreak{} -\allowbreak{} a)\allowbreak{} \textbackslash{};\textbackslash{}overset\{\textbackslash{}sim\}\{\textbackslash{}re\allowbreak{}lbar\textbackslash{}relbar\textbackslash{}relbar\}\textbackslash{}\allowbreak{};\allowbreak{}}\par
\noindent Corrected reading: {\ttfamily\footnotesize \textbackslash{}[\allowbreak{}}\par
\noindent Source: {\ttfamily\footnotesize WITNESS-F7I-VI-127.3\allowbreak{}8\allowbreak{}}\par

\par\medskip\noindent\hypertarget{app-F7I-VI-127.39}{}\textbf{F7I-VI-127.39}\quad printed folio 127.\par
\noindent Printed reading: \emph{(empty)}\par
\noindent Corrected reading: {\ttfamily\footnotesize -a)\allowbreak{} \textbackslash{}xrightarrow\{\textbackslash{}\allowbreak{} \textbackslash{}sim\textbackslash{}\allowbreak{} \}\allowbreak{}}\par
\noindent Source: {\ttfamily\footnotesize WITNESS-F7I-VI-127.3\allowbreak{}9\allowbreak{}}\par

\par\medskip\noindent\hypertarget{app-F7I-VI-127.40}{}\textbf{F7I-VI-127.40}\quad printed folio 127.\par
\noindent Printed reading: \emph{(empty)}\par
\noindent Corrected reading: {\ttfamily\footnotesize R[[X\_1,\textbackslash{}ldots,X\_n]]/\allowbreak{}(F\allowbreak{}}\par
\noindent Source: {\ttfamily\footnotesize WITNESS-F7I-VI-127.4\allowbreak{}0\allowbreak{}}\par

\par\medskip\noindent\hypertarget{app-F7I-VI-127.41}{}\textbf{F7I-VI-127.41}\quad printed folio 127.\par
\noindent Printed reading: {\ttfamily\footnotesize \$\allowbreak{}}\par
\noindent Corrected reading: {\ttfamily\footnotesize .\allowbreak{}}\par
\noindent Source: {\ttfamily\footnotesize WITNESS-F7I-VI-127.4\allowbreak{}1\allowbreak{}}\par

\par\medskip\noindent\hypertarget{app-F7I-VI-127.42}{}\textbf{F7I-VI-127.42}\quad printed folio 127.\par
\noindent Printed reading: {\ttfamily\footnotesize since\allowbreak{} \$F\$\allowbreak{} is\allowbreak{} a\allowbreak{} quadratic\allowbreak{} form.\allowbreak{}}\par
\noindent Corrected reading: {\ttfamily\footnotesize \textbackslash{}]\allowbreak{}}\par
\noindent Source: {\ttfamily\footnotesize WITNESS-F7I-VI-127.4\allowbreak{}2\allowbreak{}}\par

\par\medskip\noindent\hypertarget{app-F7I-VI-127.43}{}\textbf{F7I-VI-127.43}\quad printed folio 127.\par
\noindent Printed reading: {\ttfamily\footnotesize other\allowbreak{}}\par
\noindent Corrected reading: {\ttfamily\footnotesize converse\allowbreak{}}\par
\noindent Source: {\ttfamily\footnotesize ADOPT-F7I-VI-127.43\allowbreak{}}\par

\par\medskip\noindent\hypertarget{app-F7I-VI-127.44}{}\textbf{F7I-VI-127.44}\quad printed folio 127.\par
\noindent Printed reading: {\ttfamily\footnotesize direction\allowbreak{}}\par
\noindent Corrected reading: {\ttfamily\footnotesize follows\allowbreak{}}\par
\noindent Source: {\ttfamily\footnotesize ADOPT-F7I-VI-127.44\allowbreak{}}\par

\par\medskip\noindent\hypertarget{app-F7I-VI-127.45}{}\textbf{F7I-VI-127.45}\quad printed folio 127.\par
\noindent Printed reading: {\ttfamily\footnotesize is\allowbreak{}}\par
\noindent Corrected reading: {\ttfamily\footnotesize from\allowbreak{}}\par
\noindent Source: {\ttfamily\footnotesize ADOPT-F7I-VI-127.45\allowbreak{}}\par

\par\medskip\noindent\hypertarget{app-F7I-VI-127.46}{}\textbf{F7I-VI-127.46}\quad printed folio 127.\par
\noindent Printed reading: {\ttfamily\footnotesize trivial\allowbreak{}}\par
\noindent Corrected reading: {\ttfamily\footnotesize invariance\allowbreak{} of\allowbreak{} \$I\_R\$\allowbreak{} under\allowbreak{} \$R\$-isomorphisms\allowbreak{}}\par
\noindent Source: {\ttfamily\footnotesize ADOPT-F7I-VI-127.46\allowbreak{}}\par

\par\medskip\noindent\hypertarget{app-F7I-VI-128.1}{}\textbf{F7I-VI-128.1}\quad printed folio 128.\par
\noindent Printed reading: {\ttfamily\footnotesize =\allowbreak{}}\par
\noindent Corrected reading: {\ttfamily\footnotesize :\allowbreak{}}\par
\noindent Source: {\ttfamily\footnotesize 7I:273\allowbreak{}}\par

\par\medskip\noindent\hypertarget{app-F7I-VI-129.1}{}\textbf{F7I-VI-129.1}\quad printed folio 128.\par
\noindent Printed reading: {\ttfamily\footnotesize n\allowbreak{}}\par
\noindent Corrected reading: {\ttfamily\footnotesize m\allowbreak{}}\par
\noindent Source: {\ttfamily\footnotesize 7I:274\allowbreak{}}\par

\par\medskip\noindent\hypertarget{app-F7I-VI-130.1}{}\textbf{F7I-VI-130.1}\quad printed folio 129.\par
\noindent Printed reading: {\ttfamily\footnotesize \textbackslash{}overset\allowbreak{}}\par
\noindent Corrected reading: {\ttfamily\footnotesize \textbackslash{}xrightarrow\allowbreak{}}\par
\noindent Source: {\ttfamily\footnotesize FIX-7I-D-0195\allowbreak{}}\par

\par\medskip\noindent\hypertarget{app-F7I-VI-130.2}{}\textbf{F7I-VI-130.2}\quad printed folio 129.\par
\noindent Printed reading: {\ttfamily\footnotesize \}\{\textbackslash{}relbar\textbackslash{}joinrel\textbackslash{}re\allowbreak{}lbar\textbackslash{}joinrel\textbackslash{}relbar\textbackslash{}\allowbreak{}joinrel\textbackslash{}relbar\allowbreak{}}\par
\noindent Corrected reading: \emph{(empty)}\par
\noindent Source: {\ttfamily\footnotesize WITNESS-F7I-VI-130.2\allowbreak{}}\par

\par\medskip\noindent\hypertarget{app-F7I-VI-131.1}{}\textbf{F7I-VI-131.1}\quad printed folio 129.\par
\noindent Printed reading: {\ttfamily\footnotesize \textbackslash{}lambda\allowbreak{}}\par
\noindent Corrected reading: {\ttfamily\footnotesize \textbackslash{}chi\allowbreak{}}\par
\noindent Source: {\ttfamily\footnotesize 7I:277\allowbreak{}}\par

\par\medskip\noindent\hypertarget{app-F7I-VI-132.1}{}\textbf{F7I-VI-132.1}\quad printed folio 130.\par
\noindent Printed reading: {\ttfamily\footnotesize \textasciicircum{}0\allowbreak{}}\par
\noindent Corrected reading: \emph{(empty)}\par
\noindent Source: {\ttfamily\footnotesize FIX-7I-D-0200\allowbreak{}}\par

\par\medskip\noindent\hypertarget{app-F7I-VI-133.1}{}\textbf{F7I-VI-133.1}\quad printed folio 131.\par
\noindent Printed reading: {\ttfamily\footnotesize a\allowbreak{}}\par
\noindent Corrected reading: {\ttfamily\footnotesize z\allowbreak{}}\par
\noindent Source: {\ttfamily\footnotesize 7I:262\allowbreak{}}\par

\par\medskip\noindent\hypertarget{app-F7I-VI-134.1}{}\textbf{F7I-VI-134.1}\quad printed folio 131.\par
\noindent Printed reading: \emph{(empty)}\par
\noindent Corrected reading: {\ttfamily\footnotesize t\textbackslash{},\allowbreak{}}\par
\noindent Source: {\ttfamily\footnotesize ADOPT-F7I-VI-134.1\allowbreak{}}\par

\par\medskip\noindent\hypertarget{app-F7I-VI-135.1}{}\textbf{F7I-VI-135.1}\quad printed folio 132.\par
\noindent Printed reading: {\ttfamily\footnotesize \textbackslash{}mathop\{\textbackslash{}hphantom\{\textbackslash{}b\allowbreak{}igcap\}\}\allowbreak{}}\par
\noindent Corrected reading: {\ttfamily\footnotesize \textbackslash{}prod\allowbreak{}}\par
\noindent Source: {\ttfamily\footnotesize 7I:264\allowbreak{}}\par

\par\endgroup\normalsize\fi
\clearpage\SGAtoc{Exposé VII}
\selectlanguage{english}
\SGApage{I}{133}


\begin{center}
\SGARefAliasTarget{sga7.sga7i.evii.nroot}SGA 7

\vspace{1em}
\underline{Expos\'e VII}

\vspace{1em}
\underline{\underline{BIEXTENSIONS OF SHEAVES OF GROUPS}}

\vspace{1em}
\underline{by A. Grothendieck}
\end{center}

0. \underline{Introduction}

The notion of a \underline{biextension} was introduced by D.~Mumford
[6], in the context of formal groups, as a convenient way of expressing
the relations between the formal groups associated with an abelian scheme
and the dual abelian scheme. In the present expos\'e we undertake a
systematic study of this notion in the general context of sheaves of
groups on a given site. In particular, we show that it fits remarkably
simply into the classical cohomological formalism of
$\overset{\mathrm L}{\otimes}$ and $\mathrm R\underline{\mathrm{Hom}}$
(\SGARefLink{sga7.sga7i.evii.n3-6-4}{3.6.4} and \SGARefLink{sga7.sga7i.evii.n3-6-5}{3.6.5}). In the following expos\'e we shall make explicit the
most interesting special cases, notably that of biextensions of group
schemes by the multiplicative group $\underline G_m$ (working with such
sites as the Zariski site, the fpqc site, or an intermediate site). We
shall see, in particular, that if $A$ and $B$ are two abelian schemes over
a scheme $S$, biextensions of $(A,B)$ by $\underline G_m$ admit a
remarkable interpretation as, essentially, the classical divisorial
correspondences on $A,B$. This gives a particularly convenient explicit
cohomological meaning to the notion of divisorial correspondence (already
implicit in Tate duality for abelian varieties over $p$-adic fields).
This interpretation of divisorial correspondences is also convenient for
following the fate of such a correspondence when $A$ and $B$ (assumed to
be defined over the fraction field of a discrete valuation ring) are made
to degenerate by means of N\'eron theory.


\SGApage{I}{\(\approx\)134}
The results of Expos\'e VIII will show, for example, that one obtains a
biextension of the pair of identity components of the N\'eron models,
hence a biextension of their connected special fibres, which plays the
role of a specialization of the original divisorial correspondence.
These results will be used in Expos\'e IX, in conjunction with the
monodromy theorem (I~3), to prove certain results about the N\'eron models
of abelian varieties, the most important being the ``semistable reduction
theorem'' for the N\'eron model.

Throughout what follows, unless expressly stated otherwise, we work in a
fixed $\underline U$-topos denoted by $\underline T$; that is (SGA~4,
Expos\'e~IV), $\underline T$ is a category equivalent to the category of
sheaves on a suitable $\underline U$-site (which may, if desired, be taken
to be $\underline T$ itself, endowed with its canonical topology). In our
terminology we shall often identify the objects of $\underline T$ with
the sheaves of sets on $\underline T$. All groups, extensions of groups,
and so forth considered below are relative to this topos $\underline T$.

1. \underline{Complements on extensions of groups}

\SGARefTarget{sga7.sga7i.evii.n1-1}{1.1}. Let $P$ and $G$ be two groups of $\underline T$, and let

\begin{equation}\tag{\SGARefTarget{sga7.sga7i.evii.n1-1-1}{1.1.1}}
\begin{tikzcd}[column sep=large]
1 \arrow[r] \& G \arrow[r, "i"] \& E \arrow[r, "j"] \& P \arrow[r] \& 1
\end{tikzcd}
\end{equation}

be an extension of $P$ by $G$. Thus $E$ is a group,
$j:E\longrightarrow P$ is an epimorphism of groups, and $i$ induces an
isomorphism of $G$ with $\ker j$. We shall reinterpret the datum of such
an extension in different terms, which will be convenient for the
definition and study of biextensions.


\SGApage{I}{135}
\SGAvar{\SGARefTarget{sga7.sga7i.evii.n1-1-2}{1.1.2}. First, the homomorphism $i$ makes $G$ act on $E$ in two ways, on
the left and on the right, by setting
\[
 gxg'=i(g)xi(g')
\]
for $g,g'\in G(S)$ and $x\in E(S)$, where $S$ is any object of
$\underline T$. If, by means of $j$, we regard $E$ as an object ``over
$P$'', that is, as an object of the induced topos $E_{/P}$,
this amounts to saying that the group $G_P=G\times P$ induced by $G$ in
$\underline T_{/P}$ acts on $E$ on the left and on the right. Moreover,
it is well known that each of these actions makes $E$ a
\underline{torsor} (a left torsor and a right torsor, respectively) over
$P$, with group $G_P$; more precisely still, it makes $E$ a
\underline{bitorsor} under $(G_P,G_P)$. Recall [2] that this means that
$E$ is a right torsor $E_d$ under the right action of $G_P$, and that the
left action induces an isomorphism of $G_P$ with the group of
automorphisms of $E_d$ (the sheaf of groups of automorphisms of $E$ that
commute with the right action of $G_P$ on $E$). This condition is
equivalent to the symmetric condition: the left action of $G_P$ on $E$
makes $E$ a left torsor $E_s$, and the right action induces an
isomorphism of the opposite group $G_P^{\mathrm o}$ with the group of
automorphisms of $E_s$.

\SGARefTarget{sga7.sga7i.evii.n1-1-2-1}{1.1.2.1}. When $\underline T$ is the ``point topos'', that is, the
category of sets (equivalent to the category of sheaves on the one-point
topological space), so that $P$ and $G$ are ordinary groups, the datum of
a bitorsor $E$ under $(G_P,G_P)$ is plainly equivalent to the datum, for
each $p\in P$, of a bitorsor $E_p$ under $(G,G)$, namely the fibre of $E$
at $p$. An analogous interpretation is available in the general case.
Indeed, giving a bitorsor $E$ under $(G_P,G_P)$ is equivalent to giving,
for every ``point'' $p\in P(S)$, where $S$ is any object of
$\underline T$, a bitorsor $E_p$ under $(G_S,G_S)$ over $S$, in such a
way that the formation of $E_p$ is ``compatible with base change along
every morphism $S'\longrightarrow S$''. Thus, if
$p':S'\longrightarrow P$ is the composite morphism, one is also given an
isomorphism of bitorsors from $P_{p}$ to the inverse image of $P_p$ under
$S'\longrightarrow S$, and these isomorphisms satisfy the usual
transitivity condition for a composite
$S''\longrightarrow S'\longrightarrow S$.


When stated in canonical language (SGA~1, Expos\'e~VI), this reduces to
the following tautology: giving a bitorsor under $(G_P,G_P)$ over $P$ is
equivalent to giving a cartesian section, over the category
$\underline T_{/P}$, of the fibred category of bitorsors under
$(G_S,G_S)$ over a variable base $S\in\operatorname{ob}\underline
T_{/P}$. More precisely, the ``value at $P$'' functor induces an
equivalence between the category of the cartesian sections in question
and the category of bitorsors under $(G_S,G_P)$. This follows trivially
from the fact that $P$ is a final object of the base category
$\underline T_{/P}$ of our fibred category: for \underline{every}
category fibred over $\underline T_{/P}$, the ``value at $S$'' functor is
an equivalence between the category of cartesian sections of that fibred
category and its fibre category over $P$. Despite its tautological
nature, however, this interpretation will be useful for expressing
certain compatibilities in language close to the set-theoretic case,
using points with values in an object $S$ in place of ordinary points.




}{\SGARefTarget{sga7.sga7i.evii.n1-1-2}{1.1.2}. First, the homomorphism $i$ makes $G$ act on $E$ in two ways, on
the left and on the right, by setting
\[
 gxg'=i(g)xi(g')
\]
for $g,g'\in G(S)$ and $x\in E(S)$, where $S$ is any object of
$\underline T$. If, by means of $j$, we regard $E$ as an object ``over
$P$'', that is, as an object of the induced topos $\underline T_{/P}$,
this amounts to saying that the group $G_P=G\times P$ induced by $G$ in
$\underline T_{/P}$ acts on $E$ on the left and on the right. Moreover,
it is well known that each of these actions makes $E$ a
\underline{torsor} (a left torsor and a right torsor, respectively) over
$P$, with group $G_P$; more precisely still, it makes $E$ a
\underline{bitorsor} under $(G_P,G_P)$. Recall [2] that this means that
$E$ is a right torsor $E_d$ under the right action of $G_P$, and that the
left action induces an isomorphism of $G_P$ with the group of
automorphisms of $E_d$ (the sheaf of groups of automorphisms of $E$ that
commute with the right action of $G_P$ on $E$). This condition is
equivalent to the symmetric condition: the left action of $G_P$ on $E$
makes $E$ a left torsor $E_s$, and the right action induces an
isomorphism of the opposite group $G_P^{\mathrm o}$ with the group of
automorphisms of $E_s$.

\SGARefTarget{sga7.sga7i.evii.n1-1-2-1}{1.1.2.1}. When $\underline T$ is the ``point topos'', that is, the
category of sets (equivalent to the category of sheaves on the one-point
topological space), so that $P$ and $G$ are ordinary groups, the datum of
a bitorsor $E$ under $(G_P,G_P)$ is plainly equivalent to the datum, for
each $p\in P$, of a bitorsor $E_p$ under $(G,G)$, namely the fibre of $E$
at $p$. An analogous interpretation is available in the general case.
Indeed, giving a bitorsor $E$ under $(G_P,G_P)$ is equivalent to giving,
for every ``point'' $p\in P(S)$, where $S$ is any object of
$\underline T$, a bitorsor $E_p$ under $(G_S,G_S)$ over $S$, in such a
way that the formation of $E_p$ is ``compatible with base change along
every morphism $S'\longrightarrow S$''. Thus, if
$p':S'\longrightarrow P$ is the composite morphism, one is also given an
isomorphism of bitorsors from $E_{p'}$ to the inverse image of $E_p$ under
$S'\longrightarrow S$, and these isomorphisms satisfy the usual
transitivity condition for a composite
$S''\longrightarrow S'\longrightarrow S$.


When stated in canonical language (SGA~1, Expos\'e~VI), this reduces to
the following tautology: giving a bitorsor under $(G_P,G_P)$ over $P$ is
equivalent to giving a cartesian section, over the category
$\underline T_{/P}$, of the fibred category of bitorsors under
$(G_S,G_S)$ over a variable base $S\in\operatorname{ob}\underline
T_{/P}$. More precisely, the ``value at $P$'' functor induces an
equivalence between the category of the cartesian sections in question
and the category of bitorsors under $(G_P,G_P)$. This follows trivially
from the fact that $P$ is a final object of the base category
$\underline T_{/P}$ of our fibred category: for \underline{every}
category fibred over $\underline T_{/P}$, the ``value at $P$'' functor is
an equivalence between the category of cartesian sections of that fibred
category and its fibre category over $P$. Despite its tautological
nature, however, this interpretation will be useful for expressing
certain compatibilities in language close to the set-theoretic case,
using points with values in an object $S$ in place of ordinary points.




}{E7I-VII-001}\SGApage{I}{136}
\SGARefTarget{sga7.sga7i.evii.n1-1-3}{1.1.3}. In terms of the extension structure (\SGARefLink{sga7.sga7i.evii.n1-1-1}{1.1.1}), the bitorsor
structure of \SGARefLink{sga7.sga7i.evii.n1-1-2}{1.1.2} on $E$ under $(G_P,G_P)$ used only part of the
multiplicative structure of $E$: namely, multiplication of pairs of
``points'' of which at least one comes from a point of $G$ (via $i$). We
shall now interpret the multiplicative structure of $E$ as an additional
structure on the bitorsor under consideration.


\SGApage{I}{\(\approx\)137}
\SGAvar{For this purpose, first consider two sections $p,p'$ of $P$. With the
notation of \SGARefLink{sga7.sga7i.evii.n1-1-2-1}{1.1.2.1}, these give two bitorsors $E_p,E_{p'}$ under $(G,G)$,
namely the inverse images in $E$ of the sections
$p,p':e\longrightarrow P$ (where $e$ denotes the final object of
$\underline T$). It is then clear that multiplication in $E$ induces a
morphism of objects of $\underline T$
\[
 \varphi:P_p\times P_{p'}\longrightarrow P_{pp'},
\]
which, in terms of the left and right structures of the three bitorsors
involved (and, as usual, using points with values in an arbitrary object
$S\in\operatorname{ob}\underline T$), satisfies
\begin{equation}\tag{\SGARefTarget{sga7.sga7i.evii.n1-1-3-0}{1.1.3.0}}
\left\{
\begin{aligned}
 \varphi(gx,x')&=g\varphi(x,x'),
 &\qquad \varphi(x,x'g')&=\varphi(x,x')g',\\
 \varphi(xg,x')&=\varphi(x,gx').
\end{aligned}
\right.
\end{equation}
Here $g,g'\in G(S)$, $x\in P_p(S)$, and $x'\in P_{p'}(S)$. In terms of
the contracted product [2]
\[
 P_p\mathbin{\overset{G}{\times}}P_{p'},
\]
the third relation says that $\varphi$ factors as
\[
 P_p\times P_{p'}\longrightarrow
 P_p\mathbin{\overset{G}{\times}}P_{p'}
 \xrightarrow{\ \varphi_{p,p'}\ }P_{pp'},
\]
where the first arrow is the canonical quotient morphism. The first two
relations, on the other hand, say that the morphism thus obtained,
\begin{equation}\tag{\SGARefTarget{sga7.sga7i.evii.n1-1-3-1}{1.1.3.1}}
 \varphi_{p,p'}:
 P_p\mathbin{\overset{G}{\times}}P_{p'}\longrightarrow P_{pp'},
\end{equation}
is a homomorphism of \underline{bitorsors}, where the bitorsor structure
on the source is the one for which the left (respectively, right)
operations of $G$ come from the left (respectively, right) operations of
$G$ on the first factor $P_p$ (respectively, on the second factor
$P_{p'}$).


}{For this purpose, first consider two sections $p,p'$ of $P$. With the
notation of \SGARefLink{sga7.sga7i.evii.n1-1-2-1}{1.1.2.1}, these give two bitorsors $E_p,E_{p'}$ under $(G,G)$,
namely the inverse images in $E$ of the sections
$p,p':e\longrightarrow P$ (where $e$ denotes the final object of
$\underline T$). It is then clear that multiplication in $E$ induces a
morphism of objects of $\underline T$
\[
 \varphi:E_p\times E_{p'}\longrightarrow E_{pp'},
\]
which, in terms of the left and right structures of the three bitorsors
involved (and, as usual, using points with values in an arbitrary object
$S\in\operatorname{ob}\underline T$), satisfies
\begin{equation}\tag{\SGARefTarget{sga7.sga7i.evii.n1-1-3-0}{1.1.3.0}}
\left\{
\begin{aligned}
 \varphi(gx,x')&=g\varphi(x,x'),
 &\qquad \varphi(x,x'g')&=\varphi(x,x')g',\\
 \varphi(xg,x')&=\varphi(x,gx').
\end{aligned}
\right.
\end{equation}
Here $g,g'\in G(S)$, $x\in E_p(S)$, and $x'\in E_{p'}(S)$. In terms of
the contracted product [2]
\[
 E_p\mathbin{\overset{G}{\times}}E_{p'},
\]
the third relation says that $\varphi$ factors as
\[
 E_p\times E_{p'}\longrightarrow
 E_p\mathbin{\overset{G}{\times}}E_{p'}
 \xrightarrow{\ \varphi_{p,p'}\ }E_{pp'},
\]
where the first arrow is the canonical quotient morphism. The first two
relations, on the other hand, say that the morphism thus obtained,
\begin{equation}\tag{\SGARefTarget{sga7.sga7i.evii.n1-1-3-1}{1.1.3.1}}
 \varphi_{p,p'}:
 E_p\mathbin{\overset{G}{\times}}E_{p'}\longrightarrow E_{pp'},
\end{equation}
is a homomorphism of \underline{bitorsors}, where the bitorsor structure
on the source is the one for which the left (respectively, right)
operations of $G$ come from the left (respectively, right) operations of
$G$ on the first factor $E_p$ (respectively, on the second factor
$E_{p'}$).


}{E7I-VII-002}\SGApage{I}{\(\approx\)138}
Notice that, like every homomorphism of bitorsors under a fixed pair of
groups, $\varphi_{p,p'}$ is necessarily an \underline{isomorphism}.

The definition of the preceding isomorphism (\SGARefLink{sga7.sga7i.evii.n1-1-3-1}{1.1.3.1}) immediately extends
to the case in which $p,p'$ are points of $P$ with values in an arbitrary
object $S$ of $\underline T$, rather than sections of $P$: one simply
applies the preceding discussion to the induced topos
$\underline T_{/S}$ and to the extension of groups over $S$ obtained from
(\SGARefLink{sga7.sga7i.evii.n1-1-1}{1.1.1}) by base change. The formation of the isomorphisms
$\varphi_{p,p'}$ is then compatible with every base change
$S'\longrightarrow S$. In other words, it corresponds to a homomorphism
of cartesian sections of the fibred category over the category
$\underline T_{/P\times P}$ of pairs of arrows
$p,p':S\rightrightarrows P$ with target $P$.

Giving a system of $\varphi_{p,p'}$ satisfying this condition is, of
course, equivalent to giving $\varphi_{p,p'}$ in the ``universal'' case
$S=P\times P$, where $p,p'$ are the two projections
$\operatorname{pr}_i$ ($i=1,2$) from $P\times P$ to $P$. It is then an
isomorphism of bitorsors under
$(G_{P\times P},G_{P\times P})$:
\begin{equation}\tag{\SGARefTarget{sga7.sga7i.evii.n1-1-3-2}{1.1.3.2}}
 \varphi:\operatorname{pr}_1^*(E)\,\operatorname{pr}_2^*(E)
 \overset{\sim}{\longrightarrow}\pi^*(E),
\end{equation}
where
\[
 \pi:P\times P\longrightarrow P
\]
is the morphism of objects of $\underline T$ defining the multiplicative
structure of the group $P$. The compatibility properties below are most
conveniently stated, however, in terms of the data (\SGARefLink{sga7.sga7i.evii.n1-1-3-1}{1.1.3.1}), avoiding
diagrams that are less close to intuition.


\SGApage{I}{139}
\SGAvar{\SGARefTarget{sga7.sga7i.evii.n1-1-4}{1.1.4}. In writing the relations (\SGARefLink{sga7.sga7i.evii.n1-1-3-0}{1.1.3.0}), which enabled us to express
the given composition law on $E$ in terms of isomorphisms of bitorsors
(\SGARefLink{sga7.sga7i.evii.n1-1-3-1}{1.1.3.1}), we used associativity of this composition law only for a
product of three ``points'' of $E$ of which at least one comes from a
point of $G$. To express full associativity in terms of the
$\varphi_{p,p'}$, one must require, for three points $p,p',p"$ of $P$
with values in an object $S$ of $\underline T$, that the following
diagram be commutative:
\begin{equation}\tag{\SGARefTarget{sga7.sga7i.evii.n1-1-4-1}{1.1.4.1}}
\begin{tikzcd}[column sep=huge, row sep=large]
 \& E_{pp'}E_{p''} \arrow[dr, "{\varphi_{pp',p''}}"] \& \\
 E_pE_{p'}E_{p''}
   \arrow[ur, "{\varphi_{p,p'}\wedge\id}"]
   \arrow[dr, "{\id\wedge\varphi_{p',p''}}"']
 \&\& E_{pp'p''} \\
 \& E_pE_{p'p''} \arrow[ur, "{\varphi_{p,p'p''}}"'] \&
\end{tikzcd}
\end{equation}
In this diagram we have written the contracted product of bitorsors
without the signs $\mathbin{\overset{G_S}{\times}}$. It is understood
that the contracted product $Q.R$ of two bitorsors is formed using, in
passing to the quotient in $Q\times R$, the action of the structural
group $G$ \underline{on the right of} $Q$ \underline{and on the left of}
$R$. This contracted-product operation is associative up to a canonical
isomorphism [2], and we make the identification
\[
 (E_pE_{p'})E_{p''}\simeq E_p(E_{p'}E_{p''}),
\]
writing both sides simply as $E_pE_{p'}E_{p''}$. Obvious conventions of
this kind will be used throughout the rest of this expos\'e without being
expressly noted each time.

\SGARefTarget{sga7.sga7i.evii.n1-1-4-2}{1.1.4.2}. Just as the isomorphisms $\varphi_{p,p'}$ of (\SGARefLink{sga7.sga7i.evii.n1-1-3-1}{1.1.3.1}) are
determined by the ``universal'' case $S=P\times P$,
$p=\operatorname{pr}_1$, $p'=\operatorname{pr}_2$ considered in
(\SGARefLink{sga7.sga7i.evii.n1-1-3-2}{1.1.3.2}), so the associativity condition (\SGARefLink{sga7.sga7i.evii.n1-1-4-1}{1.1.4.1}) may likewise be
checked only in the universal case
$S=P\times P\times P$, $p=\operatorname{pr}_1$,
$p'=\operatorname{pr}_2$, $p"=\operatorname{pr}_3$. We leave it to the
reader to write the diagram in whatever notation is preferred.




\SGARefTarget{sga7.sga7i.evii.n1-1-5}{1.1.5}. Conversely, suppose that the groups $P$ and $G$ of
$\underline T$ are given, together with a bitorsor
\begin{equation}\tag{\SGARefTarget{sga7.sga7i.evii.n1-1-5-1}{1.1.5.1}}
 j:E\longrightarrow P
\end{equation}
under $(G_P,G_P)$ (the two actions of $G_P$ on $E$, or equivalently of
$G$ on $P$, being understood in the notation), and a system of
isomorphisms of bitorsors over a variable base $S$,
\begin{equation}\tag{\SGARefTarget{sga7.sga7i.evii.n1-1-5-2}{1.1.5.2}}
 \varphi_{p,p'}:E_pE_{p'}\overset{\sim}{\longrightarrow}E_{pp'}
 \qquad (p,p'\in P(S)),
\end{equation}
compatible with base change $S'\longrightarrow S$; in other words, an
isomorphism of bitorsors (\SGARefLink{sga7.sga7i.evii.n1-1-3-2}{1.1.3.2}). Suppose that the corresponding
diagrams (\SGARefLink{sga7.sga7i.evii.n1-1-4-1}{1.1.4.1}) commute. I claim that \underline{there is a unique
extension structure on $E$ of $P$ by $G$ such that the preceding data
are obtained from it by the constructions of \SGARefLink{sga7.sga7i.evii.n1-1-2}{1.1.2} and \SGARefLink{sga7.sga7i.evii.n1-1-3}{1.1.3}}.

}{\SGARefTarget{sga7.sga7i.evii.n1-1-4}{1.1.4}. In writing the relations (\SGARefLink{sga7.sga7i.evii.n1-1-3-0}{1.1.3.0}), which enabled us to express
the given composition law on $E$ in terms of isomorphisms of bitorsors
(\SGARefLink{sga7.sga7i.evii.n1-1-3-1}{1.1.3.1}), we used associativity of this composition law only for a
product of three ``points'' of $E$ of which at least one comes from a
point of $G$. To express full associativity in terms of the
$\varphi_{p,p'}$, one must require, for three points $p,p',p''$ of $P$
with values in an object $S$ of $\underline T$, that the following
diagram be commutative:
\begin{equation}\tag{\SGARefTarget{sga7.sga7i.evii.n1-1-4-1}{1.1.4.1}}
\begin{tikzcd}[column sep=huge, row sep=large]
 \& E_{pp'}E_{p''} \arrow[dr, "{\varphi_{pp',p''}}"] \& \\
 E_pE_{p'}E_{p''}
   \arrow[ur, "{\varphi_{p,p'}\wedge\id}"]
   \arrow[dr, "{\id\wedge\varphi_{p',p''}}"']
 \&\& E_{pp'p''} \\
 \& E_pE_{p'p''} \arrow[ur, "{\varphi_{p,p'p''}}"'] \&
\end{tikzcd}
\end{equation}
In this diagram we have written the contracted product of bitorsors
without the signs $\mathbin{\overset{G_S}{\times}}$. It is understood
that the contracted product $Q.R$ of two bitorsors is formed using, in
passing to the quotient in $Q\times R$, the action of the structural
group $G$ \underline{on the right of} $Q$ \underline{and on the left of}
$R$. This contracted-product operation is associative up to a canonical
isomorphism [2], and we make the identification
\[
 (E_pE_{p'})E_{p''}\simeq E_p(E_{p'}E_{p''}),
\]
writing both sides simply as $E_pE_{p'}E_{p''}$. Obvious conventions of
this kind will be used throughout the rest of this expos\'e without being
expressly noted each time.

\SGARefTarget{sga7.sga7i.evii.n1-1-4-2}{1.1.4.2}. Just as the isomorphisms $\varphi_{p,p'}$ of (\SGARefLink{sga7.sga7i.evii.n1-1-3-1}{1.1.3.1}) are
determined by the ``universal'' case $S=P\times P$,
$p=\operatorname{pr}_1$, $p'=\operatorname{pr}_2$ considered in
(\SGARefLink{sga7.sga7i.evii.n1-1-3-2}{1.1.3.2}), so the associativity condition (\SGARefLink{sga7.sga7i.evii.n1-1-4-1}{1.1.4.1}) may likewise be
checked only in the universal case
$S=P\times P\times P$, $p=\operatorname{pr}_1$,
$p'=\operatorname{pr}_2$, $p''=\operatorname{pr}_3$. We leave it to the
reader to write the diagram in whatever notation is preferred.




\SGARefTarget{sga7.sga7i.evii.n1-1-5}{1.1.5}. Conversely, suppose that the groups $P$ and $G$ of
$\underline T$ are given, together with a bitorsor
\begin{equation}\tag{\SGARefTarget{sga7.sga7i.evii.n1-1-5-1}{1.1.5.1}}
 j:E\longrightarrow P
\end{equation}
under $(G_P,G_P)$ (the two actions of $G_P$ on $E$, or equivalently of
$G$ on $E$, being understood in the notation), and a system of
isomorphisms of bitorsors over a variable base $S$,
\begin{equation}\tag{\SGARefTarget{sga7.sga7i.evii.n1-1-5-2}{1.1.5.2}}
 \varphi_{p,p'}:E_pE_{p'}\overset{\sim}{\longrightarrow}E_{pp'}
 \qquad (p,p'\in P(S)),
\end{equation}
compatible with base change $S'\longrightarrow S$; in other words, an
isomorphism of bitorsors (\SGARefLink{sga7.sga7i.evii.n1-1-3-2}{1.1.3.2}). Suppose that the corresponding
diagrams (\SGARefLink{sga7.sga7i.evii.n1-1-4-1}{1.1.4.1}) commute. I claim that \underline{there is a unique
extension structure on $E$ of $P$ by $G$ such that the preceding data
are obtained from it by the constructions of \SGARefLink{sga7.sga7i.evii.n1-1-2}{1.1.2} and \SGARefLink{sga7.sga7i.evii.n1-1-3}{1.1.3}}.

}{E7I-VII-003}\SGApage{I}{\(\approx\)140}
First observe that, on taking $p,p'$ in (\SGARefLink{sga7.sga7i.evii.n1-1-5-2}{1.1.5.2}) to be the identity
section $1$ of $P$, one obtains an isomorphism
\begin{equation}\tag{\SGARefTarget{sga7.sga7i.evii.n1-1-5-3}{1.1.5.3}}
 \varphi_{1,1}:E_1\overset{\sim}{\longleftarrow}
 E_1\mathbin{\overset{G}{\times}}E_1.
\end{equation}
Taking the contracted product, on the left say, of both sides with the
inverse bitorsor $E_1^{-1}$ of $E_1$ [2] therefore gives a canonical
isomorphism
\begin{equation}\tag{\SGARefTarget{sga7.sga7i.evii.n1-1-5-4}{1.1.5.4}}
 E_1\simeq\text{the unit bitorsor }{}_sG_d
\end{equation}
(that is, $G$ acting on itself by left and right translations), in view
of the canonical isomorphism


\SGApage{I}{\(\approx\)141}
\[
 E_1^{-1}.E_1\simeq\text{the unit bitorsor }{}_sG_d
\]
of loc.~cit., valid for every bitorsor. Equivalently, one has a section
$1_E$ of $E_1$, hence of $E$,
\begin{equation}\tag{\SGARefTarget{sga7.sga7i.evii.n1-1-5-5}{1.1.5.5}}
 1_E\in\Gamma(E_1)=\operatorname{Hom}(e,E_1)\subset\Gamma(E),
\end{equation}
namely the image under (\SGARefLink{sga7.sga7i.evii.n1-1-5-4}{1.1.5.4}) of the distinguished section of
${}_sG_d$, satisfying
\[
 g.1_E=1_E.g\qquad(g\in G(S)).
\]
By construction this section is characterized by
\begin{equation}\tag{\SGARefTarget{sga7.sga7i.evii.n1-1-5-6}{1.1.5.6}}
 \varphi_{1,1}(1_E,1_E)=1_E.
\end{equation}

On the other hand, giving the morphism of torsors $\varphi$ of
(\SGARefLink{sga7.sga7i.evii.n1-1-3-2}{1.1.3.2}) is plainly equivalent to giving a morphism
\begin{equation}\tag{\SGARefTarget{sga7.sga7i.evii.n1-1-5-7}{1.1.5.7}}
 \pi_E:E\times E\longrightarrow E,
 \qquad\text{also denoted }(x,y)\longmapsto xy,
\end{equation}
satisfying
\begin{equation}\tag{\SGARefTarget{sga7.sga7i.evii.n1-1-5-8}{1.1.5.8}}
 j(xy)=j(x)j(y)\qquad(x,y\in E(S)),
\end{equation}
so that $j$ is compatible with the multiplicative structures $\pi_E$
on $E$ and $\pi$ on $P$, together with the three relations (\SGARefLink{sga7.sga7i.evii.n1-1-3-0}{1.1.3.0})
(preferably beginning with the last of these). Finally define a morphism
of objects of $\underline T$,
\[
 i:G\longrightarrow E,
\]
by the formula
\begin{equation}\tag{\SGARefTarget{sga7.sga7i.evii.n1-1-5-9}{1.1.5.9}}
 i(g)=g.1_E=1_E.g.
\end{equation}


\SGApage{I}{\(\approx\)142}
The equality of the last two terms was already observed in constructing
$1_E$. I now claim that \underline{the multiplication law $\pi_E$ on
$E$ makes $E$ a group with identity section $1_E$, and that
$j:E\longrightarrow P$ and $i:G\longrightarrow E$ are homomorphisms of
groups making $E$ an extension of $P$ by $G$}.

\SGAvar{Indeed, associativity of the composition law on $E$ is precisely the
commutativity of the diagrams (\SGARefLink{sga7.sga7i.evii.n1-1-4-1}{1.1.4.1}). The construction of $\pi_E$ by
means of the isomorphisms $\varphi_{p,p'}$ of (\SGARefLink{sga7.sga7i.evii.n1-1-3-1}{1.1.3.1}) also shows that,
for every $x\in E(S)$, left and right translation by $x$ in $E(S)$ are
injective. Since $1_E1_E=1_E$ by (\SGARefLink{sga7.sga7i.evii.n1-1-5-6}{1.1.5.6}), it follows at once that
$1_E$ is a two-sided identity of $E$. We leave to the reader the
verification that an $x\in E(S)$ lying over $p\in P(S)$ has an inverse;
one uses the canonical isomorphism deduced from $\varphi_{p,p^{-1}}$,
\begin{equation}\tag{\SGARefTarget{sga7.sga7i.evii.n1-1-5-10}{1.1.5.10}}
 P_{p^{-1}}\simeq(P_p)^{\mathrm o}.
\end{equation}
Thus $E$ is indeed a \underline{group}. We have already observed that
$j$ is compatible with the composition laws (\SGARefLink{sga7.sga7i.evii.n1-1-5-8}{1.1.5.8}), and the same is
true of $i$, since, for $g,g'\in G(S)$,
\begin{align*}
 i(g)i(g')
 &= (g.1_E)(g'.1_E)\\
 &= (g.1_E.g')1_E &&\text{(the last relation of (\SGARefLink{sga7.sga7i.evii.n1-1-3-0}{1.1.3.0}))}\\
 &= g.1_E.g' &&\text{(since $1_E$ is an identity)}\\
 &= g.(1_E.g')=g.(g'.1_E)\\
 &= (gg').1_E &&\text{by (\SGARefLink{sga7.sga7i.evii.n1-1-5-9}{1.1.5.9})}\\
 &= i(gg').
\end{align*}
Since $j$ is the structural projection of a torsor, it is an
epimorphism of objects of $\underline T$. It remains only to see that
$i$ induces an isomorphism of $G$ with $\ker j=E_1$, which is clear from
the definition (\SGARefLink{sga7.sga7i.evii.n1-1-5-9}{1.1.5.9}) and the construction of $1_E$.


}{Indeed, associativity of the composition law on $E$ is precisely the
commutativity of the diagrams (\SGARefLink{sga7.sga7i.evii.n1-1-4-1}{1.1.4.1}). The construction of $\pi_E$ by
means of the isomorphisms $\varphi_{p,p'}$ of (\SGARefLink{sga7.sga7i.evii.n1-1-3-1}{1.1.3.1}) also shows that,
for every $x\in E(S)$, left and right translation by $x$ in $E(S)$ are
injective. Since $1_E1_E=1_E$ by (\SGARefLink{sga7.sga7i.evii.n1-1-5-6}{1.1.5.6}), it follows at once that
$1_E$ is a two-sided identity of $E$. We leave to the reader the
verification that an $x\in E(S)$ lying over $p\in P(S)$ has an inverse;
one uses the canonical isomorphism deduced from $\varphi_{p,p^{-1}}$,
\begin{equation}\tag{\SGARefTarget{sga7.sga7i.evii.n1-1-5-10}{1.1.5.10}}
 E_{p^{-1}}\simeq(E_p)^{\mathrm o}.
\end{equation}
Thus $E$ is indeed a \underline{group}. We have already observed that
$j$ is compatible with the composition laws (\SGARefLink{sga7.sga7i.evii.n1-1-5-8}{1.1.5.8}), and the same is
true of $i$, since, for $g,g'\in G(S)$,
\begin{align*}
 i(g)i(g')
 &= (g.1_E)(g'.1_E)\\
 &= (g.1_E.g')1_E &&\text{(the last relation of (\SGARefLink{sga7.sga7i.evii.n1-1-3-0}{1.1.3.0}))}\\
 &= g.1_E.g' &&\text{(since $1_E$ is an identity)}\\
 &= g.(1_E.g')=g.(g'.1_E)\\
 &= (gg').1_E &&\text{by (\SGARefLink{sga7.sga7i.evii.n1-1-5-9}{1.1.5.9})}\\
 &= i(gg').
\end{align*}
Since $j$ is the structural projection of a torsor, it is an
epimorphism of objects of $\underline T$. It remains only to see that
$i$ induces an isomorphism of $G$ with $\ker j=E_1$, which is clear from
the definition (\SGARefLink{sga7.sga7i.evii.n1-1-5-9}{1.1.5.9}) and the construction of $1_E$.


}{E7I-VII-004}\SGApage{I}{143}
\SGARefTarget{sga7.sga7i.evii.n1-1-6}{1.1.6}. We shall accept that the preceding arguments prove the following.
For a fixed topos $\underline T$, the preceding constructions give an
\underline{equivalence between the category of extensions of groups of
$\underline T$} as in (\SGARefLink{sga7.sga7i.evii.n1-1-1}{1.1.1}), where $P,G,E$ are all variable,
\underline{and the category of systems $(P,G,E,\varphi)$ in which $P$
and $G$ are groups of $\underline T$, $E$ is a bitorsor under
$(G_P,G_P)$ over $P$, and $\varphi$ is an isomorphism of bitorsors}
(\SGARefLink{sga7.sga7i.evii.n1-1-3-2}{1.1.3.2}) (or, equivalently, a system of base-change-compatible
isomorphisms of bitorsors (\SGARefLink{sga7.sga7i.evii.n1-1-3-1}{1.1.3.1})), \underline{satisfying the
associativity condition expressed by commutativity of diagram (\SGARefLink{sga7.sga7i.evii.n1-1-4-2}{1.1.4.2})}
(or, equivalently, of the diagrams (\SGARefLink{sga7.sga7i.evii.n1-1-4-1}{1.1.4.1})).

We leave it to the reader to spell out, using the preceding dictionary,
the evident notion of inverse image (respectively, direct image) of an
extension (\SGARefLink{sga7.sga7i.evii.n1-1-1}{1.1.1}) under a homomorphism of groups
\[
 P'\longrightarrow P
 \qquad\bigl(\text{respectively, }G\longrightarrow G'\bigr),
\]
or under both operations simultaneously. We likewise leave to the reader
the compatibility of the preceding equivalence of categories with the
inverse-image functor $f^*$ associated with a morphism of topoi
\[
 f:\underline T'\longrightarrow\underline T.
\]
For this one should use the datum $\varphi$ in the form (\SGARefLink{sga7.sga7i.evii.n1-1-3-2}{1.1.3.2}), in
preference to the $\varphi_{p,p'}$ of (\SGARefLink{sga7.sga7i.evii.n1-1-3-1}{1.1.3.1}); it then defines the
corresponding datum $f^*(\varphi)$ on the topos $\underline T'$.




\SGARefTarget{sga7.sga7i.evii.n1-2}{1.2}. \underline{The commutative case}. From \SGARefLink{sga7.sga7i.evii.n1-4}{1.4} onward we shall be
concerned only with \underline{commutative} extensions of groups. We
therefore assume that $P$ and $G$ are commutative, and spell out here,
in terms of the bitorsor description $(P,G,E,\varphi)$ of an extension,
the condition that $E$ be a commutative group.


\SGApage{I}{\(\approx\)144}
\SGAvar{An obviously necessary first condition---which in fact expresses that
$i(G)$ be a \underline{central} subgroup of $E$---is that the bitorsors
$E_p$, for $p\in P(S)$, be ordinary torsors; that is, the left and right
actions of $G$ on $P_p$ must coincide. Equivalently, $E$ over $P$ is an
ordinary torsor under $G_P$. (Thus there will henceforth be no need to
speak of bitorsors, and we shall consider only ordinary torsors.) Assuming
this preliminary condition, commutativity of $E$ is expressed by
commutativity of the following diagrams, for $p,p'\in P(S)$:
\begin{equation}\tag{\SGARefTarget{sga7.sga7i.evii.n1-2-1}{1.2.1}}
\begin{tikzcd}[column sep=large, row sep=large]
 E_pE_{p'} \arrow[r, "{\varphi_{p,p'}}"] \arrow[d, "\mathrm{sym}"']
   \& E_{pp'} \arrow[d, "\id"] \\
 E_{p'}E_p \arrow[r, "{\varphi_{p',p}}"]
   \& E_{p'p}.
\end{tikzcd}
\end{equation}
Here the first vertical arrow is the canonical commutativity isomorphism
(defined because $G$ is now commutative); moreover $pp'=p'p$ because
$P$ is commutative. This condition is plainly equivalent to the
corresponding condition in the universal situation
$S=P\times P$, $p=\operatorname{pr}_1$,
$p'=\operatorname{pr}_2$, which we need not write out again.




}{An obviously necessary first condition---which in fact expresses that
$i(G)$ be a \underline{central} subgroup of $E$---is that the bitorsors
$E_p$, for $p\in P(S)$, be ordinary torsors; that is, the left and right
actions of $G$ on $E_p$ must coincide. Equivalently, $E$ over $P$ is an
ordinary torsor under $G_P$. (Thus there will henceforth be no need to
speak of bitorsors, and we shall consider only ordinary torsors.) Assuming
this preliminary condition, commutativity of $E$ is expressed by
commutativity of the following diagrams, for $p,p'\in P(S)$:
\begin{equation}\tag{\SGARefTarget{sga7.sga7i.evii.n1-2-1}{1.2.1}}
\begin{tikzcd}[column sep=large, row sep=large]
 E_pE_{p'} \arrow[r, "{\varphi_{p,p'}}"] \arrow[d, "\mathrm{sym}"']
   \& E_{pp'} \arrow[d, "\id"] \\
 E_{p'}E_p \arrow[r, "{\varphi_{p',p}}"]
   \& E_{p'p}.
\end{tikzcd}
\end{equation}
Here the first vertical arrow is the canonical commutativity isomorphism
(defined because $G$ is now commutative); moreover $pp'=p'p$ because
$P$ is commutative. This condition is plainly equivalent to the
corresponding condition in the universal situation
$S=P\times P$, $p=\operatorname{pr}_1$,
$p'=\operatorname{pr}_2$, which we need not write out again.




}{E7I-VII-005}\SGARefTarget{sga7.sga7i.evii.n1-3}{1.3}. \underline{Relations with rigidified or birigidified torsors}

\SGARefTarget{sga7.sga7i.evii.n1-3-1-0}{\SGARefAliasTarget{sga7.sga7i.evii.n1-3-1}1.3.1.0}. Let $P$ be an object of $\underline T$ equipped with a section
$e_P$, let $G$ be a group of $\underline T$, and let $E$ be a torsor
under $G_P$. A \underline{rigidification of $E$}, relative to the given
section $e_P$ of $P$, is a trivialization of the torsor
$E_1=e_P^*(E)$ under $G$;


\SGApage{I}{\(\approx\)145}
that is, a section of this torsor. This terminology is inspired by the
case where $G$ is commutative and
\begin{equation}\tag{\SGARefTarget{sga7.sga7i.evii.n1-3-1-1}{1.3.1.1}}
 \operatorname{Hom}_{\mathrm{pt}}(P,G)\simeq(e),
\end{equation}
where the left-hand side denotes the set of morphisms carrying $e_P$ to
the identity section $e_G$ of $G$. This relation is equivalent to saying
that the natural homomorphism
\begin{equation}\tag{\SGARefTarget{sga7.sga7i.evii.n1-3-1-2}{1.3.1.2}}
 \Gamma(G)\xrightarrow{\ \sim\ }\operatorname{Hom}(P,G),
\end{equation}
which associates with a section $g$ of $G$ the constant morphism
$P\longrightarrow G$ with value $g$, is an isomorphism. It is immediate
that these equivalent conditions are also equivalent to saying that
every automorphism of the torsor $E$ preserving the given
rigidification is the identity automorphism (such automorphisms
correspond precisely to the elements of the left-hand side of
(\SGARefLink{sga7.sga7i.evii.n1-3-1-1}{1.3.1.1})). In particular, under condition (\SGARefLink{sga7.sga7i.evii.n1-3-1-1}{1.3.1.1}), the category of
torsors under $G_P$ rigidified relative to $e_P$ is ``rigid''.

\SGARefTarget{sga7.sga7i.evii.n1-3-2}{1.3.2}. Let, in addition, $Q$ be an object of $\underline T$ equipped
with a section $e_Q$, and now let $E$ denote a torsor under
$G_{P\times Q}$. Since $e_Q$ defines a section
$e_{P\times Q/P}$ of $P\times Q$ over $P$, one may consider
rigidifications $\alpha$ of $E$ relative to this section, working in the
induced topos $\underline T_{/P}$; explicitly, these are sections of the
torsor induced by $E$ on $P\times e_Q$. Interchanging the roles of $P$
and $Q$, one may likewise consider rigidifications $\beta$ of $E$
relative to the section $e_{P\times Q/Q}$ of $P\times Q$ over $Q$
defined by $e_P$. Two such partial rigidifications are said to be
\underline{compatible} if the sections $\alpha$ and $\beta$ agree on
$e_P\times e_Q$, that is, if they define the same section of the torsor
$E_{1,1}$ under $G$, the inverse image


\SGApage{I}{\(\approx\)146}
of $E$ under the section $(e_P,e_Q)$ of $P\times Q$. A
\underline{birigidification} of the torsor $E$ relative to the pair of
sections $(e_P,e_Q)$ is a pair of compatible partial rigidifications.

If $G$ is commutative, then $E$ admits a birigidification
$(\alpha,\beta)$ if and only if it admits partial rigidifications
$\alpha$ and $\beta$; moreover one may find a birigidification in which
$\alpha$ (or $\beta$) is prescribed in advance. Indeed, it suffices to
correct $\beta$ by the morphism $Q\longrightarrow G$ obtained as the
composite
\[
 Q\longrightarrow e\xrightarrow{\lambda}G,
\]
where $\lambda$ is defined by
$\beta e_Q=(\alpha e_P)\lambda$.

Still in the commutative case, if $E,E'$ are two birigidified torsors
under $G_{P\times Q}$, saying that they are isomorphic merely as
$G_{P\times Q}$-torsors is equivalent to saying that they are
isomorphic as birigidified $G_{P\times Q}$-torsors. In other words, if
there is an isomorphism $u:E\longrightarrow E'$ of
$G_{P\times Q}$-torsors, then there is one
$v:E\longrightarrow E'$ that preserves the birigidifications. Indeed,
measuring the discrepancy of $u$ from the given birigidifications gives
morphisms $f:P\longrightarrow G$ and $g:Q\longrightarrow G$ such that
$f(e_P)=g(e_Q)=a\in G(e)$. There is then a morphism
$P\times Q\longrightarrow G$ whose restrictions along
$e_{P\times Q/P}$ and $e_{P\times Q/Q}$ are respectively $f$ and $g$,
namely
\[
 h(p,q)=f(p)g(q)a^{-1},
\]
and one may take $v=uh$.

\SGAvar{\SGARefTarget{sga7.sga7i.evii.n1-3-3}{1.3.3}. The terminology of \SGARefLink{sga7.sga7i.evii.n1-3-1}{1.3.1} extends immediately to the case of a
\underline{bitorsor} $E$ under $(G_P,G_P)$. A rigidification of the
bitorsor $E$ is then required to be an isomorphism from $e_P(E)$ to
the trivial bitorsor ${}_sG_d$; equivalently, it is defined by a
\underline{central} section $e_E$ of $e_P(E)$, one satisfying
\[
 g.e_E=e_E.g\qquad(g\in G(S)).
\]
The terminology of birigidifications extends in the same way to
bitorsors.




}{\SGARefTarget{sga7.sga7i.evii.n1-3-3}{1.3.3}. The terminology of \SGARefLink{sga7.sga7i.evii.n1-3-1}{1.3.1} extends immediately to the case of a
\underline{bitorsor} $E$ under $(G_P,G_P)$. A rigidification of the
bitorsor $E$ is then required to be an isomorphism from $e_P^*(E)$ to
the trivial bitorsor ${}_sG_d$; equivalently, it is defined by a
\underline{central} section $e_E$ of $e_P^*(E)$, one satisfying
\[
 g.e_E=e_E.g\qquad(g\in G(S)).
\]
The terminology of birigidifications extends in the same way to
bitorsors.




}{E7I-VII-006}\SGApage{I}{147}
\SGARefTarget{sga7.sga7i.evii.n1-3-4}{1.3.4}. Return to the case in which two groups $P$ and $G$ are given,
and choose the identity section $e_P$ in order to define
rigidifications of bitorsors over $P$ and birigidifications of bitorsors
over $P\times P$. There are natural functors
\begin{equation}\tag{\SGARefTarget{sga7.sga7i.evii.n1-3-4-1}{1.3.4.1}}
 \mathbf{EXT}(P;G)\xrightarrow{\ i\ }
 \mathrm{BITORSRIG}(P,G)\xrightarrow{\ \delta\ }
 \mathbf{BITORSBIRIG}(P,P;G),
\end{equation}
where the three expressions denote, respectively, the category of
extensions of the group $P$ by $G$, the category of bitorsors under
$(G_P,G_P)$ rigidified at $e_P$, and the category of bitorsors under
$(G_{P\times P},G_{P\times P})$ birigidified at $(e_P,e_P)$. The functor
$i$ associates with an extension $E$ of $P$ by $G$ the corresponding
bitorsor over $P$ (\SGARefLink{sga7.sga7i.evii.n1-1-6}{1.1.6}), rigidified by the identity section of $E$.
The functor $\delta$ is defined by
\begin{equation}\tag{\SGARefTarget{sga7.sga7i.evii.n1-3-4-2}{1.3.4.2}}
 \delta(E)=\pi^*(E)\,\pr_1^*(E)^{-1}\,\pr_2^*(E)^{-1},
\end{equation}
where $\pi:P\times P\longrightarrow P$ is the composition law of $P$,
and $\pr_1,\pr_2$ are the two projections. There is a
\underline{canonical birigidification of $\delta(E)$}. Indeed, the
restriction of $\delta(E)$ to $P\times e_P$ is canonically isomorphic
to
\[
 \id_P^*(E)\,\id_P^*(E)^{-1}\,e^*(E)^{-1},
\]
where $e=\pr_2\operatorname{inj}_1$ is the constant morphism
$P\longrightarrow P$ with value $e_P$. The given rigidification of $E$
provides an isomorphism from $e^*(E)$ to the trivial bitorsor, and hence
an isomorphism of the restriction of $\delta(E)$ to $P\times e_P$ with
the trivial bitorsor. Symmetrically one defines a partial
rigidification along $e_P\times P$, and the two partial
rigidifications are immediately seen to be compatible.


\SGApage{I}{148}
\SGARefTarget{sga7.sga7i.evii.n1-3-4-3}{1.3.4.3}. Observe that the composite $\delta i$ of the functors in
(\SGARefLink{sga7.sga7i.evii.n1-3-4-1}{1.3.4.1}) is canonically isomorphic to the trivial functor, that is, to
the constant functor whose value is the trivial birigidified bitorsor
over $P\times P$. Indeed, for every extension $E$ of $P$ by $G$, the
isomorphism $\varphi$ of (\SGARefLink{sga7.sga7i.evii.n1-1-3-2}{1.1.3.2}) plainly defines a trivialization of
$\delta i(E)$ compatible with its canonical birigidification.




\SGAvar{\underline{Proposition \SGARefTarget{sga7.sga7i.evii.n1-3-5}{1.3.5}. Let $P$ and $Q$ be two groups of
$\underline T$.}

}{\underline{Proposition \SGARefTarget{sga7.sga7i.evii.n1-3-5}{1.3.5}. Let $P$ and $G$ be two groups of
$\underline T$.}

}{E7I-VII-007}a) \underline{Suppose that every morphism of sheaves of sets
$P\times P\longrightarrow G$ which is trivial on $P\times e_P$ and on
$e_P\times P$ is trivial. Then the functor $i$ of (\SGARefLink{sga7.sga7i.evii.n1-3-4-1}{1.3.4.1}) is fully
faithful. In particular, on a bitorsor $E$ under $(G_P,G_P)$ rigidified
relative to $e_P$, that is, equipped with a central section $e_E$ of
$E_1=e_P^*(E)$, there is at most one extension structure of $P$ by $G$
compatible with its bitorsor structure and having $e_E$ as identity
section.}

\SGAvar{b) \underline{Suppose in addition that every morphism of sheaves of
sets $P\times P\times P\longrightarrow G$ which is trivial on the
partial products $P\times P\times e_P$, $P\times e_P\times P$, and
$e_P\times P\times P$ is trivial. Let $E$ be a torsor under
$G_P$. In order that $E$ admit an extension structure of $P$ by
$G$ compatible with its bitorsor structure, it is necessary and
sufficient that the bitorsor under
$(G_{P\times P},G_{P\times P})$ underlying}
\[
 \delta(E)=\pi^*(E)\,\pr_1^*(E)^{-1}\,\pr_2^*(E)^{-1}
\]
\underline{be trivial. When this is so, for every central section
$e_E$ of $E_1=e_P^*(E)$, that is, every rigidification of $E$ relative
to $e_P$, there is a unique extension structure on $E$ compatible with
its bitorsor structure and having $e_E$ as identity section. In other
words, $E$, regarded as a bitorsor rigidified relative to $e_P$, lies
in the essential image of $i$. Moreover, the birigidified bitorsor
$\delta(E)$ has one and only one trivialization.}


}{b) \underline{Suppose in addition that every morphism of sheaves of
sets $P\times P\times P\longrightarrow G$ which is trivial on the
partial products $P\times P\times e_P$, $P\times e_P\times P$, and
$e_P\times P\times P$ is trivial. Let $E$ be a bitorsor under
$(G_P,G_P)$. In order that $E$ admit an extension structure of $P$ by
$G$ compatible with its bitorsor structure, it is necessary and
sufficient that the bitorsor under
$(G_{P\times P},G_{P\times P})$ underlying}
\[
 \delta(E)=\pi^*(E)\,\pr_1^*(E)^{-1}\,\pr_2^*(E)^{-1}
\]
\underline{be trivial. When this is so, for every central section
$e_E$ of $E_1=e_P^*(E)$, that is, every rigidification of $E$ relative
to $e_P$, there is a unique extension structure on $E$ compatible with
its bitorsor structure and having $e_E$ as identity section. In other
words, $E$, regarded as a bitorsor rigidified relative to $e_P$, lies
in the essential image of $i$. Moreover, the birigidified bitorsor
$\delta(E)$ has one and only one trivialization.}


}{E7I-VII-008}\SGApage{I}{\(\approx\)149}
Thus, in terms of diagram (\SGARefLink{sga7.sga7i.evii.n1-3-4-1}{1.3.4.1}), the essential image of the fully
faithful functor $i$ is ``the kernel'' of $\delta$. More picturesquely,
diagram (\SGARefLink{sga7.sga7i.evii.n1-3-4-1}{1.3.4.1}) is an ``exact sequence of pointed categories,'' and
even of categories endowed with pseudo-group structures (in the
terminology introduced in \SGARefLink{sga7.sga7i.evii.n2-5-1}{2.5.1} below).

\SGAvar{Let us prove \SGARefLink{sga7.sga7i.evii.n1-3-5}{1.3.5}. For a), let $E,E'$ be two extensions of $P$ by $G$,
and let $f:E\longrightarrow E'$ be a morphism of the underlying
rigidified bitorsors. We show that $f$ is a morphism of extensions, or,
what plainly amounts to the same thing, a homomorphism of groups. We
must prove $f(xy)=f(x)f(y)$. If $j':E'\longrightarrow P$ denotes the
structural morphism of the torsor $E'$ over $P$, then
$j'f(xy)=j'(f(x)f(y))$: the two sides are respectively $j(f(f(xy)))$ and
$j(f(f(x)f(y)))$, which are equal because $j$ is multiplicative. Hence
\begin{equation*}
 f(xy)=u(j(x),j(y))f(x)f(y),
\end{equation*}
where
\begin{equation*}
 u:P\times P\longrightarrow G
\end{equation*}
is a uniquely determined morphism. Setting first $x=1$ and then $y=1$
in the preceding relation, and using $f(1)=1$, gives
\begin{equation*}
 u(p,1)=u(1,p)=1.
\end{equation*}
Thus $u$ is trivial on $P\times e_P$ and on $e_P\times P$; by
hypothesis it is trivial, and therefore $f(xy)=f(x)f(y)$, as required.

For b), necessity of the stated condition for $E$ to admit an
extension structure compatible with its bitorsor structure is immediate
from \SGARefLink{sga7.sga7i.evii.n1-3-4-3}{1.3.4.3}. It remains to prove sufficiency for the first assertion.
Pulling a trivialization $\varphi$ of $\delta(E)$ back along the
identity section of $P\times P$ gives a trivialization of
$E_1=e_P(E)$ (compare 1.5.1), that is, a rigidification of $E$
relative to $e_P$. Fix this rigidification. A given trivialization
$\varphi$ of $\delta(E)$ need not be compatible with the corresponding
birigidification. Relative to the natural partial rigidifications of
$\delta(E)$, its restrictions to $P\times e_P$ and $e_P\times P$ are
described by central morphisms
\begin{equation*}
 f:P\longrightarrow G,qquad g:P\longrightarrow G,
\end{equation*}
whose values at the identity section of $P$ coincide:
$f(1)=g(1)=a\in\Gamma G$. The central morphism
\begin{equation*}
 h:P\times P\longrightarrow G,qquad h(p,q)=f(p)g(q)a^{-1},
\end{equation*}
has restrictions $f$ and $g$ to $P\times e_P$ and $e_P\times P$.
Correcting $\varphi$ by $h$, that is, replacing it by
$\varphi(p,q)h(p,q)^{-1}$, therefore gives a trivialization of
$\delta(E)$ compatible with its birigidification. Such a trivialization
is unique: two of them differ by a central morphism
$P\times P\longrightarrow G$ trivial on $P\times e_P$ and
$e_P\times P$, hence trivial by hypothesis. This proves the last
assertion of \SGARefLink{sga7.sga7i.evii.n1-3-5}{1.3.5} b).

}{Let us prove \SGARefLink{sga7.sga7i.evii.n1-3-5}{1.3.5}. For a), let $E,E'$ be two extensions of $P$ by $G$,
and let $f:E\longrightarrow E'$ be a morphism of the underlying
rigidified bitorsors. We show that $f$ is a morphism of extensions, or,
what plainly amounts to the same thing, a homomorphism of groups. We
must prove $f(xy)=f(x)f(y)$. If $j':E'\longrightarrow P$ denotes the
structural morphism of the torsor $E'$ over $P$, then
$j'f(xy)=j'(f(x)f(y))$: the two sides are respectively $j(xy)$ and
$j(x)j(y)$, which are equal because $j$ is multiplicative. Hence
\begin{equation*}
 f(xy)=u(j(x),j(y))f(x)f(y),
\end{equation*}
where
\begin{equation*}
 u:P\times P\longrightarrow G
\end{equation*}
is a uniquely determined morphism. Setting first $x=1$ and then $y=1$
in the preceding relation, and using $f(1)=1$, gives
\begin{equation*}
 u(p,1)=u(1,p)=1.
\end{equation*}
Thus $u$ is trivial on $P\times e_P$ and on $e_P\times P$; by
hypothesis it is trivial, and therefore $f(xy)=f(x)f(y)$, as required.

For b), necessity of the stated condition for $E$ to admit an
extension structure compatible with its bitorsor structure is immediate
from \SGARefLink{sga7.sga7i.evii.n1-3-4-3}{1.3.4.3}. It remains to prove sufficiency for the first assertion.
Pulling a trivialization $\varphi$ of $\delta(E)$ back along the
identity section of $P\times P$ gives a trivialization of
$E_1=e_P^*(E)$ (compare 1.5.1), that is, a rigidification of $E$
relative to $e_P$. Fix this rigidification. A given trivialization
$\varphi$ of $\delta(E)$ need not be compatible with the corresponding
birigidification. Relative to the natural partial rigidifications of
$\delta(E)$, its restrictions to $P\times e_P$ and $e_P\times P$ are
described by central morphisms
\begin{equation*}
 f:P\longrightarrow G,qquad g:P\longrightarrow G,
\end{equation*}
whose values at the identity section of $P$ coincide:
$f(1)=g(1)=a\in\Gamma G$. The central morphism
\begin{equation*}
 h:P\times P\longrightarrow G,qquad h(p,q)=f(p)g(q)a^{-1},
\end{equation*}
has restrictions $f$ and $g$ to $P\times e_P$ and $e_P\times P$.
Correcting $\varphi$ by $h$, that is, replacing it by
$\varphi(p,q)h(p,q)^{-1}$, therefore gives a trivialization of
$\delta(E)$ compatible with its birigidification. Such a trivialization
is unique: two of them differ by a central morphism
$P\times P\longrightarrow G$ trivial on $P\times e_P$ and
$e_P\times P$, hence trivial by hypothesis. This proves the last
assertion of \SGARefLink{sga7.sga7i.evii.n1-3-5}{1.3.5} b).

}{E7I-VII-009}\SGApage{I}{\(\approx\)150}
It remains to prove that $\varphi$, viewed as an isomorphism of the form
(\SGARefLink{sga7.sga7i.evii.n1-1-3-2}{1.1.3.2}), satisfies the associativity condition (\SGARefLink{sga7.sga7i.evii.n1-1-4-1}{1.1.4.1}); by our
choice of $\varphi$, $e_E$ is already a two-sided identity for the
composition law defined by $\varphi$. The two composites in (\SGARefLink{sga7.sga7i.evii.n1-1-4-1}{1.1.4.1})
are isomorphisms of bitorsors under
$(G_{P\times P\times P},G_{P\times P\times P})$, so they differ by a
central morphism
\begin{equation*}
 u:P\times P\times P\longrightarrow G
\end{equation*}
through the formula
\begin{equation*}
 (xy)z=u(jx,jy,jz)\,x(yz).
\end{equation*}


\SGApage{I}{\(\approx\)151}
\SGAvar{Setting successively $y=z=1$, $z=x=1$, and $x=y=1$, and using the fact that
$1$ is a two-sided identity for the composition law, shows that $u$ is
trivial on the three partial products $ P\times e_P\times e_P$,
$e_P\times P\times e_P$, and $e_P\times e_P\times P$. The hypothesis then
implies that $u$ is trivial, so $(xy)z=x(yz)$. This completes the proof.




}{Setting successively $z=1$, $x=1$, and $y=1$, and using the fact that
$1$ is a two-sided identity for the composition law, shows that $u$ is
trivial on the three partial products $P\times P\times e_P$,
$e_P\times P\times P$, and $P\times e_P\times P$. The hypothesis then
implies that $u$ is trivial, so $(xy)z=x(yz)$. This completes the proof.




}{E7I-VII-010}\underline{Corollary \SGARefTarget{sga7.sga7i.evii.n1-3-6}{1.3.6}. Suppose that every morphism of sheaves of
sets $P\times G\longrightarrow G$ which is trivial on $P\times e_G$
and on $e_P\times G$ is trivial. Then, if $G$ is commutative, $G$ is
central in every extension of $P$ by $G$. If $P$ is also commutative
and the hypothesis of \SGARefLink{sga7.sga7i.evii.n1-3-5}{1.3.5} a) holds, then every extension $E$ of $P$
by $G$ is commutative.}

\SGAvar{For the first assertion, we must prove that the action of $P$ on $G$
defined by the extension $G$ is trivial. This action is expressed by a
morphism $u:P\times G\longrightarrow G$. The restrictions to
$P\times e_G$ and $e_P\times G$ of the corresponding morphism
\[
 (p,g)\longmapsto u(p,g)g^{-1}
\]
are plainly trivial. Hence this last morphism is trivial, so
$u(p,g)=g$, as required. Now, if $E$ is an extension of $P$ by $G$, its
opposite group $E^{\mathrm o}$ is an extension of $P^{\mathrm o}$ by
$G^{\mathrm o}$, and the bitorsor structure underlying
$i(E^{\mathrm o})$ is obtained from $i(E)$ by interchanging the left
and right actions. When $E$ is a central extension, one therefore has
$i(E)=i(E^{\mathrm o})$; if, in addition, $P$ is commutative, so that
$P=P^{\mathrm o}$, the uniqueness assertion \SGARefLink{sga7.sga7i.evii.n1-3-5}{1.3.5} a) gives
$E=E^{\mathrm o}$. Thus $E$ is commutative, proving the corollary.


}{For the first assertion, we must prove that the action of $P$ on $G$
defined by the extension $E$ is trivial. This action is expressed by a
morphism $u:P\times G\longrightarrow G$. The restrictions to
$P\times e_G$ and $e_P\times G$ of the corresponding morphism
\[
 (p,g)\longmapsto u(p,g)g^{-1}
\]
are plainly trivial. Hence this last morphism is trivial, so
$u(p,g)=g$, as required. Now, if $E$ is an extension of $P$ by $G$, its
opposite group $E^{\mathrm o}$ is an extension of $P^{\mathrm o}$ by
$G^{\mathrm o}$, and the bitorsor structure underlying
$i(E^{\mathrm o})$ is obtained from $i(E)$ by interchanging the left
and right actions. When $E$ is a central extension, one therefore has
$i(E)=i(E^{\mathrm o})$; if, in addition, $P$ is commutative, so that
$P=P^{\mathrm o}$, the uniqueness assertion \SGARefLink{sga7.sga7i.evii.n1-3-5}{1.3.5} a) gives
$E=E^{\mathrm o}$. Thus $E$ is commutative, proving the corollary.


}{E7I-VII-011}\SGApage{I}{152}
\SGARefTarget{sga7.sga7i.evii.n1-3-7}{1.3.7}. One interesting special case in which all the hypotheses of
\SGARefLink{sga7.sga7i.evii.n1-3-5}{1.3.5} and \SGARefLink{sga7.sga7i.evii.n1-3-6}{1.3.6} hold is
\begin{equation}\tag{\SGARefTarget{sga7.sga7i.evii.n1-3-7-1}{1.3.7.1}}
 \underline{\operatorname{Hom}}_{\mathrm{pt}}(P,G)=e
 \quad\text{(the final sheaf)},
\end{equation}
where the left-hand side denotes the sheaf of morphisms of pointed
sheaves of sets from $P$ to $G$. Independently of the marked sections,
this condition may also be written in the form, analogous to (\SGARefLink{sga7.sga7i.evii.n1-3-1-2}{1.3.1.2}),
\begin{equation}\tag{\SGARefTarget{sga7.sga7i.evii.n1-3-7-2}{1.3.7.2}}
 G\simeq\underline{\operatorname{Hom}}_{\mathrm{sets}}(P,G),
\end{equation}
where the second member denotes the sheaf of morphisms between the
underlying sheaves of sets of $P$ and $G$. These conditions also mean
that, for every object $Q$ of $\underline T$, the map
$u\longmapsto u\circ\pr_2$
\begin{equation}\tag{\SGARefTarget{sga7.sga7i.evii.n1-3-7-3}{1.3.7.3}}
 \operatorname{Hom}_{\mathrm{sets}}(Q,G)
 \xrightarrow{\ \sim\ }
 \operatorname{Hom}_{\mathrm{sets}}(P\times Q,G)
\end{equation}
is bijective. This property of $P$ (with $G$ fixed) is stable under
Cartesian products; in particular it implies the same property for
$P\times P$, $P\times P\times P$, and so forth. It follows immediately
that it implies the hypotheses made in \SGARefLink{sga7.sga7i.evii.n1-3-5}{1.3.5} and \SGARefLink{sga7.sga7i.evii.n1-3-6}{1.3.6}. It also implies
that all the categories occurring in (\SGARefLink{sga7.sga7i.evii.n1-3-4-1}{1.3.4.1}) are \underline{rigid},
that is, every automorphism of an object of any one of them is the
identity.

Moreover, since condition (\SGARefLink{sga7.sga7i.evii.n1-3-7-1}{1.3.7.1}) on $P,G$ is stable under every base
change $S\longrightarrow e$, descent together with the uniqueness
result \SGARefLink{sga7.sga7i.evii.n1-3-5}{1.3.5} a) shows that a rigidified bitorsor $E$ over $P$ comes from
an extension of $P$ by $G$ as soon as it does so \underline{locally over
$e$}. By \SGARefLink{sga7.sga7i.evii.n1-3-5}{1.3.5} b), this may also be expressed by the vanishing of the
section
\begin{equation*}
 \underline{\delta}(E)\in\Gamma R^1h_*(G_{P\times P})
\end{equation*}
defined by $\delta(E)$, where
\begin{equation*}
 h:P\times P\longrightarrow e
\end{equation*}
is the structural morphism.




\SGApage{I}{153}
\SGARefTarget{sga7.sga7i.evii.n1-3-8}{1.3.8}. Let $p:P\longrightarrow S$ be a morphism of schemes which is
\underline{coherent} (that is, quasi-compact and quasi-separated) and
\underline{flat}, and suppose that
\begin{equation}\tag{\SGARefTarget{sga7.sga7i.evii.n1-3-8-1}{1.3.8.1}}
 p_*(\mathcal O_P)\xleftarrow{\ \sim\ }\mathcal O_S.
\end{equation}
Since the formation of $p_*(\mathcal O_P)$ commutes with every flat base
change (by coherence of $p$), the stated conditions on the relative
scheme $P/S$ are stable under Cartesian products and hence hold, in
particular, for the Cartesian powers $P\times P$,
$P\times P\times P$, and so forth. They are also stable under every flat
base change.

\SGAvar{Now let $G$ be an \underline{affine} scheme over $S$. Relation
(\SGARefLink{sga7.sga7i.evii.n1-3-7-2}{1.3.7.2}) follows at once from (\SGARefLink{sga7.sga7i.evii.n1-3-8-1}{1.3.8.1}), since
\[
 \operatorname{Hom}_S(P,G)
 =\operatorname{Hom}_{\mathrm{Alg}}(\underline A,p_*(\mathcal O_P))
 \simeq\operatorname{Hom}_{\mathrm{Alg}}(\underline A,\mathcal O_S),
\]
where $\underline A$ is the quasi-coherent algebra over $S$ such that
$G=Spec( A)$. By the preceding observation,
the same conclusion applies to $P\times P$, $P\times P\times P$, and
so on. Likewise, if $G$ is flat over $S$, the same conclusion applies
to the relative scheme $P_G$ over $G$ with respect to the relative
scheme $G_G$ over $G$.

}{Now let $G$ be an \underline{affine} scheme over $S$. Relation
(\SGARefLink{sga7.sga7i.evii.n1-3-7-2}{1.3.7.2}) follows at once from (\SGARefLink{sga7.sga7i.evii.n1-3-8-1}{1.3.8.1}), since
\[
 \operatorname{Hom}_S(P,G)
 =\operatorname{Hom}_{\mathrm{Alg}}(\underline A,p_*(\mathcal O_P))
 \simeq\operatorname{Hom}_{\mathrm{Alg}}(\underline A,\mathcal O_S),
\]
where $\underline A$ is the quasi-coherent algebra over $S$ such that
$G=\operatorname{Spec}(\underline A)$. By the preceding observation,
the same conclusion applies to $P\times P$, $P\times P\times P$, and
so on. Likewise, if $G$ is flat over $S$, the same conclusion applies
to the relative scheme $P_G$ over $G$ with respect to the relative
scheme $G_G$ over $G$.

}{E7I-VII-012}When $P$ and $G$ are group schemes over $S$, and one works with a
topology on $(\mathrm{Sch})_{/S}$ coarser than the canonical topology,
so that $P$ and $G$ may be identified with sheaves on this site, every
morphism from $P$, $P\times P$, $P\times P\times P$, and so on to $G$
which is compatible with the marked sections is trivial. Similarly, if
$G$ is flat, every $S$-morphism $P\times G\longrightarrow G$ compatible
with the marked sections is trivial, as one sees by interpreting it as
a $G$-morphism from $P_G$ to $G_G$. Thus the hypotheses of \SGARefLink{sga7.sga7i.evii.n1-3-5}{1.3.5} hold,
as do those of \SGARefLink{sga7.sga7i.evii.n1-3-6}{1.3.6} when $G$ is flat. Moreover, the three categories
occurring in (\SGARefLink{sga7.sga7i.evii.n1-3-4-1}{1.3.4.1}) are rigid.


\SGApage{I}{154}
\SGAvar{\SGARefTarget{sga7.sga7i.evii.n1-3-8-2}{1.3.8.2}. The most interesting case in which (\SGARefLink{sga7.sga7i.evii.n1-3-8-1}{1.3.8.1}) holds, and
continues to hold after every extension of the base, is that in which
$P$ is an \underline{abelian scheme} over $S$; then the conditions of
\SGARefLink{sga7.sga7i.evii.n1-3-7}{1.3.7} themselves hold. In this case the conclusions of \SGARefLink{sga7.sga7i.evii.n1-3-5}{3.5} and \SGARefLink{sga7.sga7i.evii.n1-3-6}{3.6}
are therefore valid for every affine group $G$ over $S$. This is the
case originally considered by J.~P. Serre, and the preceding
developments were plainly inspired by it. An interesting case in which
(\SGARefLink{sga7.sga7i.evii.n1-3-8-1}{1.3.8.1}) holds without holding universally is that in which $P$ is the
\underline{N\'eron model} of an abelian scheme over the fraction field
of a discrete valuation ring with spectrum $S$.




}{\SGARefTarget{sga7.sga7i.evii.n1-3-8-2}{1.3.8.2}. The most interesting case in which (\SGARefLink{sga7.sga7i.evii.n1-3-8-1}{1.3.8.1}) holds, and
continues to hold after every extension of the base, is that in which
$P$ is an \underline{abelian scheme} over $S$; then the conditions of
\SGARefLink{sga7.sga7i.evii.n1-3-7}{1.3.7} themselves hold. In this case the conclusions of \SGARefLink{sga7.sga7i.evii.n1-3-5}{1.3.5} and \SGARefLink{sga7.sga7i.evii.n1-3-6}{1.3.6}
are therefore valid for every affine group $G$ over $S$. This is the
case originally considered by J.~P. Serre, and the preceding
developments were plainly inspired by it. An interesting case in which
(\SGARefLink{sga7.sga7i.evii.n1-3-8-1}{1.3.8.1}) holds without holding universally is that in which $P$ is the
\underline{N\'eron model} of an abelian scheme over the fraction field
of a discrete valuation ring with spectrum $S$.




}{E7I-VII-013}\SGARefTarget{sga7.sga7i.evii.n1-4}{1.4}. \underline{Extensions of a free commutative group}

Let $I$ be an object of $\underline T$, and put
\begin{equation}\tag{\SGARefTarget{sga7.sga7i.evii.n1-4-1}{1.4.1}}
 P=\mathbb Z[I],
\end{equation}
the free $\mathbb Z$-module generated by $I$ (SGA 4, IV, 11). With
every extension $E$ of $P$ by a given $\mathbb Z$-module $G$ (from now
on only commutative group extensions will occur), associate the
corresponding torsor under $G_P$, and then restrict it to $I$ along the
canonical morphism
\[
 I\longrightarrow P=\mathbb Z[I].
\]


\SGApage{I}{\(\approx\)155}
This gives a canonical functor
\begin{equation}\tag{\SGARefTarget{sga7.sga7i.evii.n1-4-2}{1.4.2}}
 \operatorname{EXT}(P,G)\longrightarrow
 \operatorname{TORS}(I,G_I).
\end{equation}
I claim that this functor is an \underline{equivalence of categories}.
Equivalently, for every $G$, the maps
\begin{equation}\tag{\SGARefTarget{sga7.sga7i.evii.n1-4-3}{1.4.3}}
 \operatorname{Hom}_{\mathbb Z}(P,G)\longrightarrow
 H^0(I,G_I)=G(I)
\end{equation}
and
\begin{equation}\tag{\SGARefTarget{sga7.sga7i.evii.n1-4-4}{1.4.4}}
 \operatorname{Ext}^1_{\mathbb Z}(P,G)\longrightarrow H^1(I,G_I)
\end{equation}
are isomorphisms. More generally, if $A$ is a ring (not necessarily
commutative) of $\underline T$ and $P=A[I]$ is the free $A$-module
generated by $I$, there are canonical isomorphisms, functorial in the
variable $A$-module $G$,
\begin{equation}\tag{\SGARefTarget{sga7.sga7i.evii.n1-4-5}{1.4.5}}
 \operatorname{Ext}^i_A(A[I],G)\xrightarrow{\ \sim\ }H^i(I,G_I).
\end{equation}
For $A=\mathbb Z$ and $i=0,1$, these coincide with (\SGARefLink{sga7.sga7i.evii.n1-4-3}{1.4.3}) and
(\SGARefLink{sga7.sga7i.evii.n1-4-4}{1.4.4}). To define (\SGARefLink{sga7.sga7i.evii.n1-4-5}{1.4.5}), recall that its left-hand side is the
derived functor of
\[
 G\longmapsto\operatorname{Hom}_A(A[I],G),
\]
which, by the definition of $A[I]$, is isomorphic to
$G\longmapsto G(I)=H^0(I,G_I)$. Since $G\longmapsto G_I$ carries
injective modules to injective modules of $\underline T_{/I}$ (SGA 4,
V), the assertion follows. We leave to the reader the verification of
the stated compatibility with (\SGARefLink{sga7.sga7i.evii.n1-4-3}{1.4.3}) and (\SGARefLink{sga7.sga7i.evii.n1-4-4}{1.4.4}).

Finally, the functor (\SGARefLink{sga7.sga7i.evii.n1-4-2}{1.4.2}) is compatible with inverse images under a
morphism of topoi
\[
 f:\underline T'\longrightarrow\underline T.
\]




\SGApage{I}{156}
2. \underline{The notion of a biextension of abelian sheaves}

\SGARefTarget{sga7.sga7i.evii.n2-0}{2.0}. Throughout this section, $P,Q,G$ denote commutative groups of the
topos $\underline T$. We shall frequently use the Cartesian diagram
\begin{equation}\tag{\SGARefTarget{sga7.sga7i.evii.n2-0-1}{2.0.1}}
\begin{tikzcd}[column sep=4em,row sep=2.6em]
 P \arrow[d] \& P\times Q \arrow[l] \arrow[d] \\
 e \& Q \arrow[l]
\end{tikzcd}
\end{equation}
and the groups $G_P$ (respectively $G_Q$, $G_{P\times Q}$) over $P$
(respectively $Q$, $P\times Q$), obtained by inverse image from the
group $G$ over $e$. We identify $G_{P\times Q}$ either with the inverse
image $(G_P)_{P\times Q}$ under $\pr_1$, or with the inverse image
$(G_Q)_{P\times Q}$ under $\pr_2$. In this notation the group
structures of $P,Q,P\times Q$ play no role; indeed, we shall almost
never use the product-group structure on $P\times Q$. We shall instead
use on $P\times Q$ the structure of a group over $P$ obtained by inverse
image from $Q$ over $e$,
\begin{equation}\tag{\SGARefTarget{sga7.sga7i.evii.n2-0-2}{2.0.2}}
 P\times Q=Q_P,
\end{equation}
and, similarly, the structure of a group over $Q$ obtained by inverse
image from $P$ over $e$,
\begin{equation}\tag{\SGARefTarget{sga7.sga7i.evii.n2-0-3}{2.0.3}}
 P\times Q=P_Q.
\end{equation}
Thus over $P$ we have the two groups $G_P,Q_P$, and over $Q$ the two
groups $G_Q,P_Q$.

If $E$ is a torsor over $P\times Q$ under $G_{P\times Q}$, it may be
viewed either as a torsor over $Q_P$ under
$G_{P\times Q}=(G_P)_{Q_P}$, or as a torsor over $P_Q$ under
$G_{P\times Q}=(G_Q)_{P_Q}$. It therefore makes sense to equip $E$ with
a commutative extension structure in either variable. Applied in the
induced topoi $\underline T_{/P}$ and $\underline T_{/Q}$, \SGARefLink{sga7.sga7i.evii.n1-1-6}{1.1.6} gives
a convenient explicit description, using the canonical isomorphisms
\[
 \underline T_{/P\times Q}\simeq
 (\underline T_{/P})_{/P\times Q},\qquad
 \underline T_{/P\times Q}\simeq
 (\underline T_{/Q})_{/P\times Q}.
\]

\SGApage{I}{\(\approx\)157}
\SGAvar{An extension structure in the $P$-direction---that is, of $Q_P$ by
$G_P$---is equivalent to giving, for every object $S$ of $\underline T$
and elements $p,p'\in P(S)$, $q\in Q(S)$, an isomorphism of torsors
under $G_S$
\begin{equation}\tag{\SGARefTarget{sga7.sga7i.evii.n2-0-4}{2.0.4}}
 \varphi_{p,p';q}:E_{p,q}E_{p',q}
 \xrightarrow{\ \sim\ }E_{pp',q}.
\end{equation}
As usual (cf. \SGARefLink{sga7.sga7i.evii.n1-1-2-1}{1.1.2.1}), for $(p,q)\in(P\times Q)(S)$, $E_{p,q}$ denotes
the torsor under $G_S$ obtained from $E$ by inverse image along
$(p,q):S\longrightarrow P\times Q$. The data (\SGARefLink{sga7.sga7i.evii.n2-0-4}{2.0.4}) must be
\underline{functorial in} $S$, associative, and commutative. The
associativity condition is the commutativity of
\begin{equation}\tag{\SGARefTarget{sga7.sga7i.evii.n2-0-5}{2.0.5}}
\begin{tikzcd}[column sep=5.1em,row sep=3.4em]
 \& E_{pp',q}E_{p'',q}
   \arrow[rd,"{\varphi_{pp',p'';q}}"] \& \\
 E_{p,q}E_{p',q}E_{p'',q}
   \arrow[ru,"{\varphi_{p,p;q}\wedge\mathrm{id}}"]
   \arrow[rd,"{\mathrm{id}\wedge\varphi_{p',p'';q}}"']
 \&\& E_{pp'p'',q} \\
 \& E_{p,q}E_{p'p'',q}
   \arrow[ru,"{\varphi_{p,p'p'';q}}"'] \&
\end{tikzcd}
\end{equation}
for $p,p',p''\in P(S)$ and $q\in Q(S)$. Commutativity is expressed by
\begin{equation}\tag{\SGARefTarget{sga7.sga7i.evii.n2-0-6}{2.0.6}}
\begin{tikzcd}[column sep=5em,row sep=3.1em]
 E_{p,q}E_{p',q} \arrow[r,"{\varphi_{p,p';q}}"]
   \arrow[d,"\mathrm{sym}"']
 \& E_{pp',q} \arrow[d,"\mathrm{id}"] \\
 E_{p',q}E_{p,q} \arrow[r,"{\varphi_{p',p;q}}"']
 \& E_{p'p,q}.
\end{tikzcd}
\end{equation}
Here $p,p'\in P(S)$ and $q\in Q(S)$.




Symmetrically, an extension structure in the $Q$-direction---that is,
of $P_Q$ by $G_Q$---is described by isomorphisms of torsors under $G_S$
\begin{equation}\tag{\SGARefTarget{sga7.sga7i.evii.n2-0-7}{2.0.7}}
 \psi_{p;q,q'}:E_{p,q}E_{p,q'}
 \xrightarrow{\ \sim\ }E_{pp',q},
\end{equation}
for $p\in P(S)$ and $q,q'\in Q(S)$. These are functorial in $S$ and
satisfy the symmetric associativity and commutativity conditions
\begin{equation}\tag{\SGARefTarget{sga7.sga7i.evii.n2-0-8}{2.0.8}}
\begin{tikzcd}[column sep=5.1em,row sep=3.4em]
 \& E_{p,qq'}E_{p,q''}
   \arrow[rd,"{\psi_{p;qq',q''}}"] \& \\
 E_{p,q}E_{p,q'}E_{p,q''}
   \arrow[ru,"{\psi_{p;q,q'}\wedge\mathrm{id}}"]
   \arrow[rd,"{\mathrm{id}\wedge\psi_{p;q',q''}}"']
 \&\& E_{p,qq'q''} \\
 \& E_{p,q}E_{p,q'q''}
   \arrow[ru,"{\psi_{p;q;q'q''}}"'] \&
\end{tikzcd}
\end{equation}
for $p\in P(S)$ and $q,q',q''\in Q(S)$, and
\begin{equation}\tag{\SGARefTarget{sga7.sga7i.evii.n2-0-9}{2.0.9}}
\begin{tikzcd}[column sep=5em,row sep=3.1em]
 E_{p,q}E_{p,q'} \arrow[r,"{\psi_{p;q,q'}}"]
   \arrow[d,"\mathrm{sym}"']
 \& E_{p,qq'} \arrow[d,"\mathrm{id}"] \\
 E_{p,q'}E_{p,q} \arrow[r,"{\psi_{p;q',q}}"']
 \& E_{p,q'q}.
\end{tikzcd}
\end{equation}
Here $p\in P(S)$ and $q,q'\in Q(S)$.


\underline{Definition \SGARefTarget{sga7.sga7i.evii.n2-1}{2.1}. Let $P,Q,G$ be abelian groups of
$\underline T$, and let $E$ be a torsor under $G_{P\times Q}$ equipped
with an extension structure of $Q_P$ by $G_P$, defined by the
isomorphisms (\SGARefLink{sga7.sga7i.evii.n2-0-4}{2.0.4}), and an extension structure of $P_Q$ by $G_Q$,
defined by the isomorphisms (\SGARefLink{sga7.sga7i.evii.n2-0-7}{2.0.7}), each functorial in $S$ and subject
to the associativity and commutativity conditions above. The two
structures are called compatible if, for every object $S$ of
$\underline T$ and $p,p'\in P(S)$, $q,q'\in Q(S)$, the following
diagram commutes:}
\begin{equation}\tag{\SGARefTarget{sga7.sga7i.evii.n2-1-1}{2.1.1}}
\begin{tikzcd}[column sep=2.7em,row sep=3.2em,font=\scriptsize]
 \&\& E_{p,qq'}E_{p',qq'}
    \arrow[rrdd,"{\varphi_{p,p';qq'}}"] \&\& \\
 E_{p,q}E_{p,q'}E_{p',q}E_{p',q'}
    \arrow[rru,"{\psi_{p;q,q'}\wedge\psi_{p';q,q'}}"]
    \arrow[d,"{}"'] \&\&\&\& \\
 E_{p,q}E_{p',q}E_{p,q'}E_{p',q'}
    \arrow[rrd,"{\varphi_{p,p';q}\wedge\varphi_{p,p';q'}}"']
 \&\&\&\& E_{pp',qq'} \\
 \&\& E_{pp',q}E_{pp',q'}
    \arrow[rru,"{\psi_{pp';q,q'}}"'] \&\&
\end{tikzcd}
\end{equation}
\underline{where the vertical arrow is induced by the symmetry
isomorphism}
\[
 E_{p,q'}E_{p',q}\xrightarrow{\ \sim\ }E_{p',q}E_{p,q'}.
\]
\underline{A \emph{biextension} of $(P,Q)$ by $G$ is a torsor $E$
under $G_{P\times Q}$ equipped with an extension structure of $Q_P$ by $G_P$
and an extension structure of $G_Q$ by $P_Q$, these two structures being
compatible in the preceding sense.}


}{An extension structure in the $P$-direction---that is, of $P_Q$ by
$G_Q$---is equivalent to giving, for every object $S$ of $\underline T$
and elements $p,p'\in P(S)$, $q\in Q(S)$, an isomorphism of torsors
under $G_S$
\begin{equation}\tag{\SGARefTarget{sga7.sga7i.evii.n2-0-4}{2.0.4}}
 \varphi_{p,p';q}:E_{p,q}E_{p',q}
 \xrightarrow{\ \sim\ }E_{pp',q}.
\end{equation}
As usual (cf. \SGARefLink{sga7.sga7i.evii.n1-1-2-1}{1.1.2.1}), for $(p,q)\in(P\times Q)(S)$, $E_{p,q}$ denotes
the torsor under $G_S$ obtained from $E$ by inverse image along
$(p,q):S\longrightarrow P\times Q$. The data (\SGARefLink{sga7.sga7i.evii.n2-0-4}{2.0.4}) must be
\underline{functorial in} $S$, associative, and commutative. The
associativity condition is the commutativity of
\begin{equation}\tag{\SGARefTarget{sga7.sga7i.evii.n2-0-5}{2.0.5}}
\begin{tikzcd}[column sep=5.1em,row sep=3.4em]
 \& E_{pp',q}E_{p'',q}
   \arrow[rd,"{\varphi_{pp',p'';q}}"] \& \\
 E_{p,q}E_{p',q}E_{p'',q}
   \arrow[ru,"{\varphi_{p,p';q}\wedge\mathrm{id}}"]
   \arrow[rd,"{\mathrm{id}\wedge\varphi_{p',p'';q}}"']
 \&\& E_{pp'p'',q} \\
 \& E_{p,q}E_{p'p'',q}
   \arrow[ru,"{\varphi_{p,p'p'';q}}"'] \&
\end{tikzcd}
\end{equation}
for $p,p',p''\in P(S)$ and $q\in Q(S)$. Commutativity is expressed by
\begin{equation}\tag{\SGARefTarget{sga7.sga7i.evii.n2-0-6}{2.0.6}}
\begin{tikzcd}[column sep=5em,row sep=3.1em]
 E_{p,q}E_{p',q} \arrow[r,"{\varphi_{p,p';q}}"]
   \arrow[d,"\mathrm{sym}"']
 \& E_{pp',q} \arrow[d,"\mathrm{id}"] \\
 E_{p',q}E_{p,q} \arrow[r,"{\varphi_{p',p;q}}"']
 \& E_{p'p,q}.
\end{tikzcd}
\end{equation}
Here $p,p'\in P(S)$ and $q\in Q(S)$.




Symmetrically, an extension structure in the $Q$-direction---that is,
of $Q_P$ by $G_P$---is described by isomorphisms of torsors under $G_S$
\begin{equation}\tag{\SGARefTarget{sga7.sga7i.evii.n2-0-7}{2.0.7}}
 \psi_{p;q,q'}:E_{p,q}E_{p,q'}
 \xrightarrow{\ \sim\ }E_{p,qq'},
\end{equation}
for $p\in P(S)$ and $q,q'\in Q(S)$. These are functorial in $S$ and
satisfy the symmetric associativity and commutativity conditions
\begin{equation}\tag{\SGARefTarget{sga7.sga7i.evii.n2-0-8}{2.0.8}}
\begin{tikzcd}[column sep=5.1em,row sep=3.4em]
 \& E_{p,qq'}E_{p,q''}
   \arrow[rd,"{\psi_{p;qq',q''}}"] \& \\
 E_{p,q}E_{p,q'}E_{p,q''}
   \arrow[ru,"{\psi_{p;q,q'}\wedge\mathrm{id}}"]
   \arrow[rd,"{\mathrm{id}\wedge\psi_{p;q',q''}}"']
 \&\& E_{p,qq'q''} \\
 \& E_{p,q}E_{p,q'q''}
   \arrow[ru,"{\psi_{p;q,q'q''}}"'] \&
\end{tikzcd}
\end{equation}
for $p\in P(S)$ and $q,q',q''\in Q(S)$, and
\begin{equation}\tag{\SGARefTarget{sga7.sga7i.evii.n2-0-9}{2.0.9}}
\begin{tikzcd}[column sep=5em,row sep=3.1em]
 E_{p,q}E_{p,q'} \arrow[r,"{\psi_{p;q,q'}}"]
   \arrow[d,"\mathrm{sym}"']
 \& E_{p,qq'} \arrow[d,"\mathrm{id}"] \\
 E_{p,q'}E_{p,q} \arrow[r,"{\psi_{p;q',q}}"']
 \& E_{p,q'q}.
\end{tikzcd}
\end{equation}
Here $p\in P(S)$ and $q,q'\in Q(S)$.


\underline{Definition \SGARefTarget{sga7.sga7i.evii.n2-1}{2.1}. Let $P,Q,G$ be abelian groups of
$\underline T$, and let $E$ be a torsor under $G_{P\times Q}$ equipped
with an extension structure in the $P$-direction, defined by the
isomorphisms (\SGARefLink{sga7.sga7i.evii.n2-0-4}{2.0.4}), and an extension structure in the $Q$-direction,
defined by the isomorphisms (\SGARefLink{sga7.sga7i.evii.n2-0-7}{2.0.7}), each functorial in $S$ and subject
to the associativity and commutativity conditions above. The two
structures are called compatible if, for every object $S$ of
$\underline T$ and $p,p'\in P(S)$, $q,q'\in Q(S)$, the following
diagram commutes:}
\begin{equation}\tag{\SGARefTarget{sga7.sga7i.evii.n2-1-1}{2.1.1}}
\begin{tikzcd}[column sep=2.7em,row sep=3.2em,font=\scriptsize]
 \&\& E_{p,qq'}E_{p',qq'}
    \arrow[rrdd,"{\varphi_{p,p';qq'}}"] \&\& \\
 E_{p,q}E_{p,q'}E_{p',q}E_{p',q'}
    \arrow[rru,"{\psi_{p;q,q'}\wedge\psi_{p';q,q'}}"]
    \arrow[d,"{}"'] \&\&\&\& \\
 E_{p,q}E_{p',q}E_{p,q'}E_{p',q'}
    \arrow[rrd,"{\varphi_{p,p';q}\wedge\varphi_{p,p';q'}}"']
 \&\&\&\& E_{pp',qq'} \\
 \&\& E_{pp',q}E_{pp',q'}
    \arrow[rru,"{\psi_{pp';q,q'}}"'] \&\&
\end{tikzcd}
\end{equation}
\underline{where the vertical arrow is induced by the symmetry
isomorphism}
\[
 E_{p,q'}E_{p',q}\xrightarrow{\ \sim\ }E_{p',q}E_{p,q'}.
\]
\underline{A \emph{biextension} of $(P,Q)$ by $G$ is a torsor $E$
under $G_{P\times Q}$ equipped with compatible extension structures in
the $P$- and $Q$-directions in the preceding sense.}


}{E7I-VII-014}\SGApage{I}{160}
\SGARefLink{sga7.sga7i.evii.n2-1-1}{2.1.1}. Equivalently, a biextension structure on the
$G_{P\times Q}$-torsor $E$ is given by two partially defined composition
laws $\varphi(x,y)$ and $\psi(x,y)$ on $E$. The first is defined when
$\pr_2j(x)=\pr_2j(y)$, and the second when
$\pr_1j(x)=\pr_1j(y)$, where $j:E\longrightarrow P\times Q$ is the
structural morphism of the torsor. Besides associativity relations of
the type
\begin{equation}\tag{\SGARefTarget{sga7.sga7i.evii.n2-1-1-1}{2.1.1.1}}
 \varphi(gx,y)=\varphi(x,gy)=g\varphi(x,y)
 \qquad(x,y\in E(S),\ g\in G(S))
\end{equation}
and the analogous relations for $\psi$, the axioms say that these
partial composition laws are associative and commutative and satisfy
\begin{equation}\tag{\SGARefTarget{sga7.sga7i.evii.n2-1-1-2}{2.1.1.2}}
 \varphi(\psi(x,y),\psi(z,t))
 =\psi(\varphi(x,z),\varphi(y,t)),
\end{equation}
where $x,y,z,t$ lie respectively in
$E_{p,q}(S),E_{p,q'}(S),E_{p',q}(S),E_{p',q'}(S)$.




\SGAvar{\SGARefTarget{sga7.sga7i.evii.n2-2}{2.2}. \underline{Identity sections}. Consider first the identity section
of $E$ regarded as an extension of $Q_P$ by $G_P$,
\begin{equation}\tag{\SGARefTarget{sga7.sga7i.evii.n2-2-1}{2.2.1}}
 e_{E/P}\in\Gamma(E/P)=\operatorname{Hom}_P(P,E),
\end{equation}
and, symmetrically, the identity section of $E$ regarded as an extension
of $P_G$ by $G_Q$,
\begin{equation}\tag{\SGARefTarget{sga7.sga7i.evii.n2-2-2}{2.2.2}}
 e_{E/Q}\in\Gamma(E/Q)=\operatorname{Hom}_Q(Q,E).
\end{equation}
In particular, \underline{the torsor $E$ over $P\times Q$ under
$G_{P\times Q}$ is canonically trivialized over $1_P\times Q$ by
$e_{E/P}$, and over $P\times1_Q$ by $e_{E/Q}$. These trivializations
coincide over the identity section $1_P\times1_Q$ of $P\times Q$}, and
hence define the same section
\begin{equation}\tag{\SGARefTarget{sga7.sga7i.evii.n2-2-3}{2.2.3}}
 e_E\in\Gamma(E),
\end{equation}
which, by abuse of language, will be called the \underline{identity
section} (or \underline{neutral section}) of $E$.


}{\SGARefTarget{sga7.sga7i.evii.n2-2}{2.2}. \underline{Identity sections}. Consider first the identity section
of $E$ regarded as an extension of $Q_P$ by $G_P$,
\begin{equation}\tag{\SGARefTarget{sga7.sga7i.evii.n2-2-1}{2.2.1}}
 e_{E/P}\in\Gamma(E/P)=\operatorname{Hom}_P(P,E),
\end{equation}
and, symmetrically, the identity section of $E$ regarded as an extension
of $P_Q$ by $G_Q$,
\begin{equation}\tag{\SGARefTarget{sga7.sga7i.evii.n2-2-2}{2.2.2}}
 e_{E/Q}\in\Gamma(E/Q)=\operatorname{Hom}_Q(Q,E).
\end{equation}
In particular, \underline{the torsor $E$ over $P\times Q$ under
$G_{P\times Q}$ is canonically trivialized over $P\times1_Q$ by
$e_{E/P}$, and over $1_P\times Q$ by $e_{E/Q}$. These trivializations
coincide over the identity section $1_P\times1_Q$ of $P\times Q$}, and
hence define the same section
\begin{equation}\tag{\SGARefTarget{sga7.sga7i.evii.n2-2-3}{2.2.3}}
 e_E\in\Gamma(E),
\end{equation}
which, by abuse of language, will be called the \underline{identity
section} (or \underline{neutral section}) of $E$.


}{E7I-VII-015}\SGApage{I}{\(\approx\)161}
To prove the asserted equality, let $e,e'$ be the two sections of $E$
deduced respectively from $e_{E/P}$ and $e_{E/Q}$. They are
characterized as the identity sections in the inverse image
$E_{1_P,1_Q}$ of $E$ along the identity section of $P\times Q$, for the
two composition laws $\varphi$ and $\psi$, respectively. Applying
(\SGARefLink{sga7.sga7i.evii.n2-1-1-2}{2.1.1.2}) to the four sections $e,e',e',e$ of $E_{1,1}$ gives $e$ on
the left and $e'$ on the right; hence $e=e'$. It follows moreover that
\underline{the two composition laws $\varphi$ and $\psi$ on $E_{1,1}$
are identical}. Indeed, under transport of structure by the morphisms
$g\mapsto g.e$ and $g\mapsto g.e'$, respectively, both correspond to
the group law of $G$, and $e=e'$.

\SGAvar{The extension of $Q_P$ by $G_P$ over the base $P$, restricted over the
identity section $1_P$ of $P$, gives an extension of $Q$ by $G$ on the
inverse image of $1_P\times Q$ in $G$. The homomorphism $e_{E/Q}$ is
immediately seen to be multiplicative for this extension structure, so
the extension is canonically split. Symmetrically, the extension of $P$
by $G$ obtained by restricting the extension of $P_Q$ by $G_Q$ over
$Q$ to the identity section $1_Q$ is canonically split by $e_{E/P}$.




}{The extension of $Q_P$ by $G_P$ over the base $P$, restricted over the
identity section $1_P$ of $P$, gives an extension of $Q$ by $G$ on the
inverse image of $1_P\times Q$ in $E$. The homomorphism $e_{E/Q}$ is
immediately seen to be multiplicative for this extension structure, so
the extension is canonically split. Symmetrically, the extension of $P$
by $G$ obtained by restricting the extension of $P_Q$ by $G_Q$ over
$Q$ to the identity section $1_Q$ is canonically split by $e_{E/P}$.




}{E7I-VII-016}\SGApage{I}{162}
\SGARefTarget{sga7.sga7i.evii.n2-3}{2.3}. \underline{Morphisms of biextensions}. Let $E$ be a biextension
of $(P,Q)$ by $G$, and let $E'$ likewise be a biextension of
$(P',Q')$ by $G'$, where $P',Q',G'$ are again commutative groups of
$\underline T$. A \underline{morphism of biextensions} from $E$ to
$E'$ is a system $(u,v,w,f)$, where
\begin{equation}\tag{\SGARefTarget{sga7.sga7i.evii.n2-3-1}{2.3.1}}
 u:P\longrightarrow P',\qquad
 v:Q\longrightarrow Q',\qquad
 w:G\longrightarrow G'
\end{equation}
are morphisms of groups, and
\begin{equation}\tag{\SGARefTarget{sga7.sga7i.evii.n2-3-2}{2.3.2}}
 f:E\longrightarrow E'
\end{equation}
is a morphism of the sheaves of sets underlying $E$ and $E'$, such
that $f$ is simultaneously a homomorphism of group extensions over
$P$ and $P'$, respectively, associated with the relative group
homomorphisms
\[
 u\times v:Q_P\longrightarrow Q'_{P'},
 \qquad
 u\times w:G_P\longrightarrow G'_{P'},
\]
and a homomorphism of group extensions over $Q$ and $Q'$, respectively,
associated with the relative group homomorphisms
\[
 u\times v:P_Q\longrightarrow P'_{Q'},
 \qquad
 v\times w:G_Q\longrightarrow G'_{Q'}.
\]
Notice moreover that $u,v,w$ are determined by $f$ alone. Indeed,
passing to the quotient recovers $u\times v$ from $f$, and then
recovers $u$ and $v$ by restriction to $P\times1_Q$ and
$1_P\times Q$, respectively; restricting $f$ to $G.e_E$ recovers
$w$. Thus no ambiguity arises from identifying the morphism
$(u,v,w,f)$ with $f$.

After taking the usual precaution of adjoining their sources and
targets to morphisms (the canonical genuflection), one sees that
\underline{the biextensions of the topos $\underline T$ are the
objects of a category}, whose objects and morphisms are those defined
in \SGARefLink{sga7.sga7i.evii.n2-1}{2.1} and \SGARefLink{sga7.sga7i.evii.n2-3}{2.3} and whose composition of morphisms is defined in the
evident way. Moreover, $(P,Q,G,E)\mapsto(P,Q,G)$ gives a canonical
functor
\begin{equation}\tag{\SGARefTarget{sga7.sga7i.evii.n2-3-3}{2.3.3}}
 \operatorname{BIEXT}(\underline T)
 \longrightarrow
 \operatorname{GRAB}(\underline T)\times
 \operatorname{GRAB}(\underline T)\times
 \operatorname{GRAB}(\underline T),
\end{equation}
where $\operatorname{BIEXT}(\underline T)$ denotes the category of
biextensions over $\underline T$, and
$\operatorname{GRAB}(\underline T)$ the category of abelian groups
over $\underline T$. It is useful to decompose this canonical functor
into the two functors $E\mapsto(P,Q)$,
\begin{equation}\tag{\SGARefTarget{sga7.sga7i.evii.n2-3-4}{2.3.4}}
 \operatorname{BIEXT}(\underline T)
 \longrightarrow
 \operatorname{GRAB}(\underline T)\times
 \operatorname{GRAB}(\underline T),
\end{equation}
and $E\mapsto G$,
\begin{equation}\tag{\SGARefTarget{sga7.sga7i.evii.n2-3-5}{2.3.5}}
 \operatorname{BIEXT}(\underline T)
 \longrightarrow
 \operatorname{GRAB}(\underline T).
\end{equation}




\SGApage{I}{163}
\SGARefTarget{sga7.sga7i.evii.n2-4}{2.4}. \underline{The category $\operatorname{BIEXT}(P,Q;G)$ for
variable $P,Q,G$}.

I claim that \underline{the canonical functor (\SGARefLink{sga7.sga7i.evii.n2-3-4}{2.3.4}) is fibred}, and
that \underline{the canonical functor (\SGARefLink{sga7.sga7i.evii.n2-3-5}{2.3.5}) is cofibred}. The first
assertion means, in particular, that if $E$ is a biextension of
$(P,Q)$ by $G$, and if
\[
 u:P'\longrightarrow P,\qquad v:Q'\longrightarrow Q
\]
are morphisms of abelian groups, then there is a biextension $E'$ of
$(P',Q')$ by $G$ and a morphism of biextensions
\[
 f:E'\longrightarrow E
\]
compatible with $u,v,\mathrm{id}_G$, having the usual universal
property for morphisms of biextensions into $E$ compatible with
$u,v,\mathrm{id}_G$.

\SGAvar{The second assertion is expressed similarly. In particular, if $E$ is
a biextension of $(P,G)$ by $G$, and if
\[
 w:G\longrightarrow G'
\]
is a morphism of abelian groups, then there is a biextension $E'$ of
$(P,Q)$ by $G'$ and a morphism of biextensions
\[
 g:E\longrightarrow E'
\]
compatible with $\mathrm{id}_P,\mathrm{id}_Q,w$, having the usual
universal property for morphisms of biextensions from $E$ compatible
with $\mathrm{id}_P,\mathrm{id}_Q,w$. We leave it to the reader to
spell out the two constructions whose existence has just been
asserted; they are obtained easily by paraphrasing the familiar
analogous constructions for extensions of ordinary groups.

Thus, for three abelian groups $P,Q,G$ of $\underline T$, write
\begin{equation}\tag{\SGARefTarget{sga7.sga7i.evii.n2-4-1}{2.4.1}}
 \operatorname{BIEXT}(P,Q;G)
\end{equation}
for the category of biextensions of $(P;Q)$ by $G$, defined as the
fibre category of (\SGARefLink{sga7.sga7i.evii.n2-3-3}{2.3.3}) over $(P,Q,G)$ (so its morphisms are the
morphisms of biextensions compatible with the identity morphisms of
$P,Q,G$). Then \underline{this category depends covariantly on $G$ and
contravariantly on $P,Q$}. Moreover, with these variances,
\underline{the category $\operatorname{BIEXT}(P,Q;G)$ depends
additively on each of the three arguments $P,Q,G$}. For example, for
the argument $G$, this means precisely:

\begin{enumerate}[label=\alph*)]
\item $\operatorname{BIEXT}(P,Q;0)$ is equivalent to the one-point
category;
\item for two abelian groups $G,G'$ of $\underline T$, the functor
\begin{equation}\tag{\SGARefTarget{sga7.sga7i.evii.n2-4-2}{2.4.2}}
 \operatorname{BIEXT}(P,G;G\times G')
 \longrightarrow
 \operatorname{BIEXT}(P,Q;G)\times
 \operatorname{BIEXT}(P,Q;G')
\end{equation}
deduced from the two projections $G\times G'\longrightarrow G$ and
$G\times G'\longrightarrow G'$ is an
\underline{equivalence of categories}.
\end{enumerate}
In the language of [4], one also says that, for fixed $P,Q$, the
category of biextensions of $(P,Q)$ by variable groups $G$ is, via the
functor $E\mapsto G$ induced by (\SGARefLink{sga7.sga7i.evii.n2-3-5}{2.3.5}), an
\underline{additive cofibred category} over
$\operatorname{GRAB}(\underline T)$.

}{The second assertion is expressed similarly. In particular, if $E$ is
a biextension of $(P,Q)$ by $G$, and if
\[
 w:G\longrightarrow G'
\]
is a morphism of abelian groups, then there is a biextension $E'$ of
$(P,Q)$ by $G'$ and a morphism of biextensions
\[
 g:E\longrightarrow E'
\]
compatible with $\mathrm{id}_P,\mathrm{id}_Q,w$, having the usual
universal property for morphisms of biextensions from $E$ compatible
with $\mathrm{id}_P,\mathrm{id}_Q,w$. We leave it to the reader to
spell out the two constructions whose existence has just been
asserted; they are obtained easily by paraphrasing the familiar
analogous constructions for extensions of ordinary groups.

Thus, for three abelian groups $P,Q,G$ of $\underline T$, write
\begin{equation}\tag{\SGARefTarget{sga7.sga7i.evii.n2-4-1}{2.4.1}}
 \operatorname{BIEXT}(P,Q;G)
\end{equation}
for the category of biextensions of $(P,Q)$ by $G$, defined as the
fibre category of (\SGARefLink{sga7.sga7i.evii.n2-3-3}{2.3.3}) over $(P,Q,G)$ (so its morphisms are the
morphisms of biextensions compatible with the identity morphisms of
$P,Q,G$). Then \underline{this category depends covariantly on $G$ and
contravariantly on $P,Q$}. Moreover, with these variances,
\underline{the category $\operatorname{BIEXT}(P,Q;G)$ depends
additively on each of the three arguments $P,Q,G$}. For example, for
the argument $G$, this means precisely:

\begin{enumerate}[label=\alph*)]
\item $\operatorname{BIEXT}(P,Q;0)$ is equivalent to the one-point
category;
\item for two abelian groups $G,G'$ of $\underline T$, the functor
\begin{equation}\tag{\SGARefTarget{sga7.sga7i.evii.n2-4-2}{2.4.2}}
 \operatorname{BIEXT}(P,Q;G\times G')
 \longrightarrow
 \operatorname{BIEXT}(P,Q;G)\times
 \operatorname{BIEXT}(P,Q;G')
\end{equation}
deduced from the two projections $G\times G'\longrightarrow G$ and
$G\times G'\longrightarrow G'$ is an
\underline{equivalence of categories}.
\end{enumerate}
In the language of [4], one also says that, for fixed $P,Q$, the
category of biextensions of $(P,Q)$ by variable groups $G$ is, via the
functor $E\mapsto G$ induced by (\SGARefLink{sga7.sga7i.evii.n2-3-5}{2.3.5}), an
\underline{additive cofibred category} over
$\operatorname{GRAB}(\underline T)$.

}{E7I-VII-017}\SGApage{I}{165}
\underline{Remark \SGARefTarget{sga7.sga7i.evii.n2-4-3}{2.4.3}.} In the terminology of \emph{loc.\ cit.}, the
preceding cofibred category is not merely additive but even
\underline{left exact} [4, \SGARefLink{sga7.sga7i.evii.n2-2}{2.2}]. We shall not verify this property here,
since it will reappear below in \SGARefLink{sga7.sga7i.evii.n3-6}{3.6} as a consequence of more precise
facts. We only point out that the general theory of \emph{loc.\ cit.}
would then describe the structure of this cofibred category, up to
equivalence of fibred categories, in terms of a chain complex
$L_\bullet$ in $\operatorname{GRAB}(\underline T)$ depending only on
$P$ and $Q$. We shall also give below a direct construction of this
complex as
\[
 P\overset{\mathbb L}{\otimes}Q,
\]
the tensor product taken in a derived category.




\SGARefTarget{sga7.sga7i.evii.n2-5}{2.5}. \underline{Structure of the category
$\operatorname{BIEXT}(P,Q;G)$}.

The general formal considerations of [4, no.~1] give several useful
consequences of the additive dependence of the categories
$\operatorname{BIEXT}(P,Q;G)$ on $G$.

\SGARefTarget{sga7.sga7i.evii.n2-5-1}{2.5.1}. First, for fixed $G$, the category
$\operatorname{BIEXT}(P,Q;G)$ has a composition law which assigns
functorially to two biextensions $E,E'$ of $(P,Q)$ by $G$ a
\underline{product biextension}, denoted $E.E'$. This construction is
simply the analogue of the familiar product of two extensions of one
commutative group by another. It is compatible with that construction
whether $E,E'$ are viewed as group extensions of $Q_P$ by $G_P$, or
as group extensions of $P_Q$ by $G_Q$. In particular, as a torsor
under $G_{P\times Q}$, $EE'$ is the product of the torsors under
$G_{P\times Q}$ defined by $E$ and $E'$.

\SGApage{I}{\(\approx\)166}
The formation of $E.E'$ is moreover \underline{associative and
commutative} up to canonical isomorphisms, and this composition law on
$\operatorname{BIEXT}(P,Q;G)$ has a two-sided neutral object, called
the \underline{trivial biextension of $(P,Q)$ by $G$}. It may be
constructed, as in \emph{loc.\ cit.}, from the unique biextension
$E_0$ of $(P,Q)$ by the zero group $0$, up to isomorphism, by means of
the homomorphism $0\longrightarrow G$. From this description one
immediately obtains an alternative description of the trivial
biextension as the trivial torsor
\[
 E=P\times Q\times G.
\]
Its two group structures over $P$ and over $Q$ are obtained,
respectively, from the product group structures on $Q\times G$ and
$P\times G$ by the base changes $P\longrightarrow e$ and
$Q\longrightarrow e$.

Finally, for every biextension $E$ of $(P,Q)$ by $G$, one can define
the \underline{inverse biextension}, denoted $E^{-1}$, together with
an isomorphism from $EE^{-1}$ to the trivial biextension. This
determines $E^{-1}$ up to unique isomorphism and makes its dependence
on $E$ functorial.

In short, the category $\operatorname{BIEXT}(P,Q;G)$ behaves like a
group object in the category $\mathbf{Cat}$ of categories, except that
associativity, commutativity, unit, and inverse relations hold only up
to specified isomorphisms. One may call it a
\underline{pseudo-group} in the 2-category $\mathbf{Cat}$.

\SGApage{I}{167}
\SGARefTarget{sga7.sga7i.evii.n2-5-2}{2.5.2}. For fixed $P,Q,G$, the preceding properties imply the
following facts. For every biextension $E$ of $(P,Q)$ by $G$, the
monoid of endomorphisms of $E$ is canonically isomorphic to the monoid
of endomorphisms of the trivial biextension $\underline 1$: an
endomorphism $u$ of $\underline 1$ is sent to
$u.\mathrm{id}_E$ on $\underline 1.E\simeq E$. The endomorphism monoid
of $\underline 1$ is in fact a \underline{commutative group}, denoted
\begin{equation}\tag{\SGARefTarget{sga7.sga7i.evii.n2-5-2-1}{2.5.2.1}}
 \operatorname{Biext}^{0}(P,Q;G).
\end{equation}
Its composition law is also given by
$(u,v)\mapsto u.v\in
\operatorname{End}(\underline 1.\underline 1)
\simeq\operatorname{End}(\underline 1)$.

Likewise, the composition law $(E,E')\mapsto EE'$ equips the set
\begin{equation}\tag{\SGARefTarget{sga7.sga7i.evii.n2-5-2-2}{2.5.2.2}}
 \operatorname{Biext}^{1}(P,Q;G)
\end{equation}
of isomorphism classes of biextensions of $(P,Q)$ by $G$ with the
structure of a \underline{commutative group}; the inverse of the class
of $E$ is the class of $E^{-1}$.

\SGARefTarget{sga7.sga7i.evii.n2-5-3}{2.5.3}. Given a group homomorphism $G\longrightarrow G'$, the
corresponding functor
\begin{equation}\tag{\SGARefTarget{sga7.sga7i.evii.n2-5-3-1}{2.5.3.1}}
 \operatorname{BIEXT}(P,Q;G)
 \longrightarrow
 \operatorname{BIEXT}(P,Q;G')
\end{equation}
is compatible, up to canonical isomorphisms, with the operations
$E.E'$ and $E^{-1}$ of \SGARefLink{sga7.sga7i.evii.n2-5-1}{2.5.1}. Consequently, the two evident maps
\begin{equation}\tag{\SGARefTarget{sga7.sga7i.evii.n2-5-3-2}{2.5.3.2}}
 \operatorname{Biext}^{i}(P,Q;G)
 \longrightarrow
 \operatorname{Biext}^{i}(P,Q;G')
 \qquad(i=0,1)
\end{equation}
are \underline{group homomorphisms}. Moreover, if $E$ is a
biextension of $(P,Q)$ by $G$ and $E'$ the corresponding
biextension by $G'$, the natural homomorphism
\begin{equation}\tag{\SGARefTarget{sga7.sga7i.evii.n2-5-3-3}{2.5.3.3}}
 \operatorname{Hom}(E,E)=\operatorname{Aut}(E)
 \longrightarrow
 \operatorname{Hom}(E',E')=\operatorname{Aut}(E')
\end{equation}
is precisely (\SGARefLink{sga7.sga7i.evii.n2-5-3-2}{2.5.3.2}) for $i=0$.

\SGApage{I}{168}
\SGARefTarget{sga7.sga7i.evii.n2-5-4}{2.5.4}. The group $\operatorname{Biext}^{0}(P,Q;G)$ admits a simple
explicit description. If $E$ is a biextension of $(P,Q)$ by $G$, an
automorphism $u$ of $E$ is first of all an automorphism of the
underlying torsor under $G_{P\times Q}$, and hence is given by a
morphism
\[
 u_0:P\times Q\longrightarrow G
\]
through the formula
\[
 u(x)=u_0(j(x))x.
\]
Compatibility of $u$ with the first, respectively the second,
extension structure on $E$ is equivalent to additivity of $u_0$ in
the second, respectively the first, argument. Thus
\underline{$u$ is an automorphism of the biextension structure of $E$
if and only if $u_0$ is bilinear}. This gives a
\underline{canonical bijection}
\begin{equation}\tag{\SGARefTarget{sga7.sga7i.evii.n2-5-4-1}{2.5.4.1}}
 \operatorname{Biext}^{0}(P,Q;G)
 \simeq
 \operatorname{Hom}(P\otimes Q,G).
\end{equation}
It is readily checked that this bijection is independent of $E$,
respects the group structures, and is functorial in $G$ for the
functoriality on the left described in (\SGARefLink{sga7.sga7i.evii.n2-5-3-2}{2.5.3.2}).

Below we shall determine $\operatorname{Biext}^{1}(P,Q;G)$ and find
that, canonically and functorially in $G$, it is isomorphic to
\[
 \operatorname{Ext}^{1}
 \bigl(P\overset{\mathbb L}{\otimes}Q,G\bigr).
\]




\SGARefTarget{sga7.sga7i.evii.n2-6}{2.6}. In \SGARefLink{sga7.sga7i.evii.n2-5}{2.5} only $G$ was varied, while $P$ and $Q$ were held fixed.
We now say a few words about variation in $P,Q$. The essential fact,
from which the others follow formally, is that the contravariant
dependence of $\operatorname{BIEXT}(P,Q;G)$ on $P,Q$
\underline{commutes} with its covariant dependence on $G$. Thus the
categories in question may be viewed as the fibre categories of a
\underline{cofibred category}
\begin{equation}\tag{\SGARefTarget{sga7.sga7i.evii.n2-6-1}{2.6.1}}
 \operatorname{BIEXT}'(\underline T)
 \longrightarrow
 \operatorname{GRAB}(\underline T)^{\mathrm{op}}\times
 \operatorname{GRAB}(\underline T)^{\mathrm{op}}\times
 \operatorname{GRAB}(\underline T),
\end{equation}
where the superscript $\mathrm{op}$ denotes the opposite category and
the structural functor is $E\mapsto(P,Q,G)$. This cofibred category is
additive in each of its three arguments, or, as we shall say,
\underline{tri-additive}.

\SGApage{I}{\(\approx\)169}
\SGAvar{By paraphrasing standard arguments, which we shall not reproduce, it
follows that the composition $E.E'$ in the fibre category
$\operatorname{BIEXT}(P,P';G)$ may be described using the additive
structure of any one of the three arguments: the three resulting
pseudo-group structures are canonically isomorphic. Consequently, the
abelian groups
\begin{equation}\tag{\SGARefTarget{sga7.sga7i.evii.n2-6-2}{2.6.2}}
 \operatorname{Biext}^{i}(P,Q;G),\qquad i=0,1,
\end{equation}
may be regarded as two trilinear functors in $P,Q,G$, respectively in
$\operatorname{GRAB}(\underline T)^{\mathrm{op}}$,
$\operatorname{GRAB}(\underline T)^{\mathrm{op}}$, and
$\operatorname{GRAB}(\underline T)$. Equivalently, they are
\underline{trilinear functors} of variable $P,Q,G$ in
$\operatorname{GRAB}(\underline T)$, contravariant in $P,Q$ and
covariant in $G$, with values in the category of abelian groups.
For this variance, the isomorphism (\SGARefLink{sga7.sga7i.evii.n2-5-4-1}{2.5.4.1}) is functorial in all of
$P,Q,G$, and not merely in $G$. The analogous assertion will hold
below for the description of $\operatorname{Biext}^{1}$ by an
$\operatorname{Ext}^{1}$.

}{By paraphrasing standard arguments, which we shall not reproduce, it
follows that the composition $E.E'$ in the fibre category
$\operatorname{BIEXT}(P,Q;G)$ may be described using the additive
structure of any one of the three arguments: the three resulting
pseudo-group structures are canonically isomorphic. Consequently, the
abelian groups
\begin{equation}\tag{\SGARefTarget{sga7.sga7i.evii.n2-6-2}{2.6.2}}
 \operatorname{Biext}^{i}(P,Q;G),\qquad i=0,1,
\end{equation}
may be regarded as two trilinear functors in $P,Q,G$, respectively in
$\operatorname{GRAB}(\underline T)^{\mathrm{op}}$,
$\operatorname{GRAB}(\underline T)^{\mathrm{op}}$, and
$\operatorname{GRAB}(\underline T)$. Equivalently, they are
\underline{trilinear functors} of variable $P,Q,G$ in
$\operatorname{GRAB}(\underline T)$, contravariant in $P,Q$ and
covariant in $G$, with values in the category of abelian groups.
For this variance, the isomorphism (\SGARefLink{sga7.sga7i.evii.n2-5-4-1}{2.5.4.1}) is functorial in all of
$P,Q,G$, and not merely in $G$. The analogous assertion will hold
below for the description of $\operatorname{Biext}^{1}$ by an
$\operatorname{Ext}^{1}$.

}{E7I-VII-018}\SGARefTarget{sga7.sga7i.evii.n2-7}{2.7}. \underline{The symmetric biextension of a given biextension}.

If $E$ is a biextension of $(P,Q)$ by $G$, define its
\underline{symmetric biextension} $\,{}^{s}E$, which is a biextension
of $(Q,P)$ by $G$. The underlying sheaf of sets of $\,{}^{s}E$ is that
of $E$, and the actions of $G$ on $E$ and $\,{}^{s}E$ are the same.
The structural morphism
\[
 {}^{s}E\longrightarrow Q\times P
\]
is the composite of $E\longrightarrow P\times Q$ with the symmetry
$P\times Q\longrightarrow Q\times P$; finally, the left and right
extension structures of $\,{}^{s}E$ are, respectively, the right and
left extension structures of $E$. Evidently $E\mapsto{}^{s}E$ is an
automorphism of the category $\operatorname{BIEXT}(\underline T)$ of
biextensions, and indeed an \underline{involutive automorphism}:
\[
 {}^{s}({}^{s}E)=E.
\]
Every statement about bitorsors $E,\ldots$ therefore gives a symmetric
statement upon being applied to $\,{}^{s}E,\ldots$. We may call it the
statement obtained by \emph{exchanging left and right}. In general we
shall spell out only one of the two symmetric statements, choosing
between left and right, for example, according to our political
preferences.




\SGApage{I}{170}
\SGARefTarget{sga7.sga7i.evii.n2-8}{2.8}. \underline{Change of topos. Localization.}

Let
\[
 f:\underline T'\longrightarrow\underline T
\]
be a morphism of topoi. The inverse-image functor $f^*$ transforms a
biextension of $(P,Q)$ by $G$ into a biextension $f^*(E)$ of
$(f^*P,f^*Q)$ by $f^*G$, as always for a structure expressible in
terms of finite projective limits and arbitrary inductive limits. To
prove compatibility of the two extension structures on $f^*(E)$
induced by those on $E$, one expresses compatibility by commutativity
of (\SGARefLink{sga7.sga7i.evii.n2-1-1}{2.1.1}) in the universal case
\[
 S=P\times P\times Q\times Q,
\]
where $p,p',q,q'$ are the four projections. As expected,
$E\mapsto f^*(E)$ gives a canonical functor
\[
 f^*:\operatorname{BIEXT}(\underline T)
 \longrightarrow
 \operatorname{BIEXT}(\underline T').
\]

\SGApage{I}{\(\approx\)171}
In what follows we shall use without further mention the fact that
every construction made with biextensions is compatible with inverse
images under morphisms of topoi. This applies in particular to the
product $EE'$ and inverse $E^{-1}$ considered in \SGARefLink{sga7.sga7i.evii.n2-5}{2.5}.

One may take for $\underline T'$ an \underline{induced topos}
$\underline T_{/S'}$, and for $f$ the canonical morphism, called the
localization morphism (SGA 4 IV 5),
\[
 f=i_{S'}:\underline T_{/S'}\longrightarrow\underline T.
\]
Then $f^*(E)$ is called the \underline{biextension induced} by $E$ on
$S'$, and is also denoted, as usual, by $E|_{S'}$ or $E_{S'}$.

Since restriction has the usual transitivity properties, the
categories $\operatorname{BIEXT}(\underline T_{/S})$, for variable
$S\in\underline T$, form the fibres of a \underline{fibred category}
over $\underline T$ (SGA 1 VI 8). In fact, this category is even a
\underline{stack} over $\underline T$ in the sense of Giraud [2]:
morphisms and objects in its fibres can be glued. For example, to give
a biextension over $S$ up to canonical isomorphism, it suffices to give
a covering family (that is, an epimorphic family) $S_i\longrightarrow
S$, a biextension $E_i$ over each $S_i$ (an object of
$\underline T_{/S_i}$), and, for every pair $(i,j)$, an isomorphism
between $E_i|_{S_j}$ and $E_j|_{S_i}$, subject to the usual cocycle
condition.

We call this the \underline{stack of biextensions over $\underline T$}.
Replacing $\underline T$ in (\SGARefLink{sga7.sga7i.evii.n2-3-3}{2.3.3}), and hence also in (\SGARefLink{sga7.sga7i.evii.n2-3-4}{2.3.4}) and
(\SGARefLink{sga7.sga7i.evii.n2-3-5}{2.3.5}), by $\underline T_{/S}$ for variable $S$ gives morphisms of
stacks over $\underline T$. The same applies to the operations $E.E'$
and $E^{-1}$ of (\SGARefLink{sga7.sga7i.evii.n2-5-1}{2.5.1}) in the categories
$\operatorname{BIEXT}(P_S,Q_S;G_S)$ for variable $S$, where $P,Q,G$
are fixed commutative groups over $\underline T$.




\SGApage{I}{172}
\SGARefTarget{sga7.sga7i.evii.n2-9}{2.9}. \underline{Relations with birigidified torsors}.

In the terminology of \SGARefLink{sga7.sga7i.evii.n1-3-2}{1.3.2}, we saw in \SGARefLink{sga7.sga7i.evii.n2-2}{2.2} that, for every
biextension $E$ of $(P,Q)$ by $G$, the underlying torsor structure
under $G_{P\times Q}$ is canonically birigidified relative to the pair
of identity sections of $P$ and $Q$. This gives a canonical functor
\begin{equation}\tag{\SGARefTarget{sga7.sga7i.evii.n2-9-1}{2.9.1}}
 \operatorname{BIEXT}(P,Q;G)
 \xrightarrow{\ i\ }
 \operatorname{TORSBIRIG}(P,Q;G),
\end{equation}
where the category on the right consists of birigidified torsors over
$P\times Q$ under $G_{P\times Q}$.

Using the identity sections of $P$ and $Q$, for every product of $n$
objects each isomorphic to $P$ or $Q$ one obtains a canonical
morphism $\operatorname{inj}_i$ from the $i$-th factor into that
product. Thus, for a morphism from such a product to $G$, one may speak
of its restriction to the $i$-th factor. Consider the products
\begin{equation}\tag{\SGARefTarget{sga7.sga7i.evii.n2-9-2}{2.9.2}}
 \left\{
 \begin{array}{l}
 P\times Q,\\[1mm]
 P\times Q\times Q,\quad P\times Q\times Q\times Q,\\[1mm]
 P\times P\times Q,\quad P\times P\times P\times Q,\\[1mm]
 P\times P\times Q\times Q.
 \end{array}
 \right.
\end{equation}

\underline{Proposition \SGARefTarget{sga7.sga7i.evii.n2-9-3}{2.9.3}. Suppose that every morphism from one of
the products (\SGARefLink{sga7.sga7i.evii.n2-9-2}{2.9.2}) to $G$ which is trivial on each factor of that
product} (\underline{embedded by the identity sections as just explained}) \underline{is
trivial. Then:}

\medskip
\noindent
a) \underline{The canonical functor (\SGARefLink{sga7.sga7i.evii.n2-9-1}{2.9.1}) is fully faithful, and
both sides of (\SGARefLink{sga7.sga7i.evii.n2-9-1}{2.9.1}) are rigid categories. In particular, if $E$ is a
torsor under $G_{P\times Q}$ birigidified relative to $(e_P,e_Q)$,
there is at most one biextension structure on $E$ compatible with its
birigidified torsor structure.}

\SGApage{I}{\(\approx\)173}
\medskip
\noindent
b) \underline{Such a biextension structure exists on $E$ if and only
if $E$ admits an extension structure of $Q_P$ by $G_P$ compatible with
its torsor structure under $G_P$, and an extension structure of $P_Q$
by $G_Q$ compatible with its torsor structure under $G_Q$
(cf.\ \SGARefLink{sga7.sga7i.evii.n1-3-5}{1.3.5}(b)). In that case there is a unique extension structure of
$Q_P$ by $G_P$ compatible with the torsor structure under $G_P$ and
the given rigidification over $P\times e_Q$, and a unique extension
structure of $P_Q$ by $G_Q$ compatible with the torsor structure under
$G_Q$ and the given rigidification over $e_P\times Q$; these two
extension structures are compatible in the sense of \SGARefLink{sga7.sga7i.evii.n2-1}{2.1}.}

\underline{Proof.} a) If $E$ is a birigidified torsor over
$P\times Q$, its automorphism group identifies with the group of
morphisms $P\times Q\longrightarrow G$ trivial on the two partial
factors. By the hypothesis applied to $P\times Q$, this group contains
only the identity automorphism, proving rigidity. Now let $E,E'$ be
biextensions of $(P,Q)$ by $G$. To see that every morphism
$f:E\longrightarrow E'$ of birigidified torsors is a morphism of
biextensions, apply \SGARefLink{sga7.sga7i.evii.n1-3-5}{1.3.5}(a) to the two extension structures defined
by $E,E'$ over $P$ and over $Q$, respectively. This uses the
hypothesis for $P\times Q\times Q$ and $P\times P\times Q$ in
(\SGARefLink{sga7.sga7i.evii.n2-9-2}{2.9.2}).

\SGAvar{\noindent
b) Necessity is clear. For sufficiency, \SGARefLink{sga7.sga7i.evii.n1-3-5}{1.3.5}(b) first shows that the
stated condition gives unique extension structures of $Q_P$ by $G_P$
and of $P_Q$ by $G_P$, compatible with the torsor structure and the
rigidifications. Here the hypothesis is used for
$P\times Q\times Q\times Q$ and $P\times P\times P\times Q$.
It remains to prove that the two extension structures are compatible,
that is, that diagram (\SGARefLink{sga7.sga7i.evii.n2-1-1}{2.1.1}) commutes. Its two composites are
isomorphisms of torsors over $P\times P\times Q\times Q$, and hence
differ by a morphism
\[
 u:P\times P\times Q\times Q\longrightarrow G.
\]
Because the partial trivializations of $E$ are identity sections for
the laws $\varphi$ and $\psi$, respectively, $u$ is trivial on each of
the four factors. The hypothesis for
$P\times P\times Q\times Q$ therefore makes $u$ trivial. This proves
the proposition.




\underline{Corollary \SGARefTarget{sga7.sga7i.evii.n2-9-4}{2.9.4}. Let $P,Q,G$ be commutative groups
satisfying}
\begin{equation}\tag{\SGARefTarget{sga7.sga7i.evii.n2-9-4-1}{2.9.4.1}}
 \underline{\operatorname{Hom}}_{\mathrm{pt}}(P,G)=e,
 \qquad
 \underline{\operatorname{Hom}}_{\mathrm{pt}}(Q,G)=e
\end{equation}
\underline{(cf.\ \SGARefLink{sga7.sga7i.evii.n1-3-7}{1.3.7}). Then the preliminary conditions of 2.9.2 are
satisfied, and conclusions (a) and (b) of that proposition hold.
Moreover, if $E$ is a torsor under $G_{P\times Q}$ birigidified
relative to $(e_P,e_Q)$, it admits a biextension structure of $(P,Q)$
by $G$ compatible with its given structures if and only if the two
morphisms of pointed sheaves which it defines,}
\begin{equation}\tag{\SGARefTarget{sga7.sga7i.evii.n2-9-4-2}{2.9.4.2}}
 Q\longrightarrow R^1p_*(G_P),
 \qquad
 P\longrightarrow R^1q_*(G_Q),
\end{equation}
\underline{where $p:P\longrightarrow e$ and $q:Q\longrightarrow e$
are the structural morphisms, are morphisms of groups.}

The first assertion was proved in \SGARefLink{sga7.sga7i.evii.n1-3-7}{1.3.7}. For the second, necessity is
clear. To prove sufficiency, apply Proposition \SGARefLink{sga7.sga7i.evii.n2-9-3}{2.9.2}(b) and \SGARefLink{sga7.sga7i.evii.n1-3-5}{1.3.5}(b).
Multiplicativity of the first morphism in (\SGARefLink{sga7.sga7i.evii.n2-9-4-2}{2.9.4.2}) says that the
rigidified torsor $E$ under $G_P$ satisfies the condition of \SGARefLink{sga7.sga7i.evii.n1-3-5}{1.3.5}(b)
\underline{locally over $P$}. As observed in \SGARefLink{sga7.sga7i.evii.n1-3-7}{1.3.7}, that condition is
then true globally over $P$, so $E$ admits an extension structure of
$Q_P$ by $G_P$ compatible with its rigidified torsor structure.
Symmetrically, $E$ admits an extension structure of $P_Q$ by $G_Q$
compatible with its rigidified torsor structure under $G_Q$. The
conclusion follows from Proposition \SGARefLink{sga7.sga7i.evii.n2-9-3}{2.9.2}(b).

\underline{Example \SGARefTarget{sga7.sga7i.evii.n2-9-5}{2.9.5}.} Let $S$ be a scheme, let $P$ and $Q$ be
abelian schemes over $S$, and let $G=\underline G_{mS}$ be the
multiplicative group scheme over $S$. We noted in \SGARefLink{sga7.sga7i.evii.n1-3-8}{1.3.8} that
(\SGARefLink{sga7.sga7i.evii.n2-9-4-1}{2.9.4.1}) then holds. A birigidified torsor $E$ over $P\times Q$ is
equivalently a birigidified invertible sheaf over $P\times Q$. Recall
that a \underline{divisorial correspondence on $P,Q$} is the
isomorphism class of such a birigidified invertible sheaf (see, for
example, [1]). In the fppf topology, the morphisms (\SGARefLink{sga7.sga7i.evii.n2-9-4-2}{2.9.4.2}) are
\begin{equation}\tag{\SGARefTarget{sga7.sga7i.evii.n2-9-5-1}{2.9.5.1}}
 Q\longrightarrow\underline{\operatorname{Pic}}_{P/S},
 \qquad
 P\longrightarrow\underline{\operatorname{Pic}}_{Q/S},
\end{equation}
defined by the divisorial correspondence. The theorem of the square
(\emph{loc.\ cit.}) says that these morphisms necessarily respect the
group structures. Thus \SGARefLink{sga7.sga7i.evii.n2-9-4}{2.9.4} shows that
\underline{the category of biextensions of $(P,Q)$ by
$G=\underline G_{m,S}$ is equivalent to the category of birigidified
invertible sheaves on $P\times Q$}; it is therefore a rigid category,
and its set of isomorphism classes identifies with the set of
divisorial correspondences on $(P,Q)$. This identification plainly
respects the natural group laws on the two sets, the first being that
of (\SGARefLink{sga7.sga7i.evii.n2-5-2-2}{2.5.2.2}).

\underline{Remark \SGARefTarget{sga7.sga7i.evii.n2-9-6}{2.9.6}.} Additivity of the first morphism in
(\SGARefLink{sga7.sga7i.evii.n2-9-4-2}{2.9.4.2}) is equivalent to factorization of the second through the
subsheaf $\underline{\operatorname{Ext}}^1(Q,G)$; note that, by \SGARefLink{sga7.sga7i.evii.n1-3-5}{1.3.5},
the canonical morphism
\[
 \underline{\operatorname{Ext}}^1(Q,G)
 \longrightarrow R^1q_*(G_Q)
\]
is injective. Likewise, additivity of the second morphism in
(\SGARefLink{sga7.sga7i.evii.n2-9-4-2}{2.9.4.2}) means that the first factors through
$\underline{\operatorname{Ext}}^1(P,G)$. This gives group
isomorphisms
\begin{equation}\tag{\SGARefTarget{sga7.sga7i.evii.n2-9-6-1}{2.9.6.1}}
 \operatorname{Biext}(P,Q;G)
 \simeq
 \operatorname{Hom}\bigl(Q,\underline{\operatorname{Ext}}^1(P,G)\bigr)
 \simeq
 \operatorname{Hom}\bigl(P,\underline{\operatorname{Ext}}^1(Q,G)\bigr).
\end{equation}
In the situation of \SGARefLink{sga7.sga7i.evii.n2-9-5}{2.9.5} these reduce to the familiar isomorphisms
\begin{equation}\tag{\SGARefTarget{sga7.sga7i.evii.n2-9-6-2}{2.9.6.2}}
 \operatorname{Corr}(P,Q)
 \simeq \operatorname{Hom}(P,Q')
 \simeq \operatorname{Hom}(Q,P'),
\end{equation}
where
\[
 P'=\underline{\operatorname{Ext}}^1(P,\underline G_m),
 \qquad
 Q'=\underline{\operatorname{Ext}}^1(Q,\underline G_m)
\]
are the dual abelian schemes of $P$ and $Q$, identified respectively
with $\underline{\operatorname{Pic}}^0_{P/S}$ and
$\underline{\operatorname{Pic}}^0_{Q/S}$. In both formulas,
$\operatorname{Hom}$ denotes the group of \underline{group
homomorphisms}.




}{\noindent
b) Necessity is clear. For sufficiency, \SGARefLink{sga7.sga7i.evii.n1-3-5}{1.3.5}(b) first shows that the
stated condition gives unique extension structures of $Q_P$ by $G_P$
and of $P_Q$ by $G_Q$, compatible with the torsor structure and the
rigidifications. Here the hypothesis is used for
$P\times Q\times Q\times Q$ and $P\times P\times P\times Q$.
It remains to prove that the two extension structures are compatible,
that is, that diagram (\SGARefLink{sga7.sga7i.evii.n2-1-1}{2.1.1}) commutes. Its two composites are
isomorphisms of torsors over $P\times P\times Q\times Q$, and hence
differ by a morphism
\[
 u:P\times P\times Q\times Q\longrightarrow G.
\]
Because the partial trivializations of $E$ are identity sections for
the laws $\varphi$ and $\psi$, respectively, $u$ is trivial on each of
the four factors. The hypothesis for
$P\times P\times Q\times Q$ therefore makes $u$ trivial. This proves
the proposition.




\underline{Corollary \SGARefTarget{sga7.sga7i.evii.n2-9-4}{2.9.4}. Let $P,Q,G$ be commutative groups
satisfying}
\begin{equation}\tag{\SGARefTarget{sga7.sga7i.evii.n2-9-4-1}{2.9.4.1}}
 \underline{\operatorname{Hom}}_{\mathrm{pt}}(P,G)=e,
 \qquad
 \underline{\operatorname{Hom}}_{\mathrm{pt}}(Q,G)=e
\end{equation}
\underline{(cf.\ \SGARefLink{sga7.sga7i.evii.n1-3-7}{1.3.7}). Then the hypothesis of Proposition 2.9.3 is
satisfied, and conclusions (a) and (b) of that proposition hold.
Moreover, if $E$ is a torsor under $G_{P\times Q}$ birigidified
relative to $(e_P,e_Q)$, it admits a biextension structure of $(P,Q)$
by $G$ compatible with its given structures if and only if the two
morphisms of pointed sheaves which it defines,}
\begin{equation}\tag{\SGARefTarget{sga7.sga7i.evii.n2-9-4-2}{2.9.4.2}}
 Q\longrightarrow R^1p_*(G_P),
 \qquad
 P\longrightarrow R^1q_*(G_Q),
\end{equation}
\underline{where $p:P\longrightarrow e$ and $q:Q\longrightarrow e$
are the structural morphisms, are morphisms of groups.}

The first assertion was proved in \SGARefLink{sga7.sga7i.evii.n1-3-7}{1.3.7}. For the second, necessity is
clear. To prove sufficiency, apply Proposition \SGARefLink{sga7.sga7i.evii.n2-9-3}{2.9.3}(b) and \SGARefLink{sga7.sga7i.evii.n1-3-5}{1.3.5}(b).
Multiplicativity of the first morphism in (\SGARefLink{sga7.sga7i.evii.n2-9-4-2}{2.9.4.2}) says that the
rigidified torsor $E$ under $G_P$ satisfies the condition of \SGARefLink{sga7.sga7i.evii.n1-3-5}{1.3.5}(b)
\underline{locally over $P$}. As observed in \SGARefLink{sga7.sga7i.evii.n1-3-7}{1.3.7}, that condition is
then true globally over $P$, so $E$ admits an extension structure of
$Q_P$ by $G_P$ compatible with its rigidified torsor structure.
Symmetrically, $E$ admits an extension structure of $P_Q$ by $G_Q$
compatible with its rigidified torsor structure under $G_Q$. The
conclusion follows from Proposition \SGARefLink{sga7.sga7i.evii.n2-9-3}{2.9.3}(b).

\underline{Example \SGARefTarget{sga7.sga7i.evii.n2-9-5}{2.9.5}.} Let $S$ be a scheme, let $P$ and $Q$ be
abelian schemes over $S$, and let $G=\underline G_{m,S}$ be the
multiplicative group scheme over $S$. We noted in \SGARefLink{sga7.sga7i.evii.n1-3-8}{1.3.8} that
(\SGARefLink{sga7.sga7i.evii.n2-9-4-1}{2.9.4.1}) then holds. A birigidified torsor $E$ over $P\times Q$ is
equivalently a birigidified invertible sheaf over $P\times Q$. Recall
that a \underline{divisorial correspondence on $P,Q$} is the
isomorphism class of such a birigidified invertible sheaf (see, for
example, [1]). In the fppf topology, the morphisms (\SGARefLink{sga7.sga7i.evii.n2-9-4-2}{2.9.4.2}) are
\begin{equation}\tag{\SGARefTarget{sga7.sga7i.evii.n2-9-5-1}{2.9.5.1}}
 Q\longrightarrow\underline{\operatorname{Pic}}_{P/S},
 \qquad
 P\longrightarrow\underline{\operatorname{Pic}}_{Q/S},
\end{equation}
defined by the divisorial correspondence. The theorem of the square
(\emph{loc.\ cit.}) says that these morphisms necessarily respect the
group structures. Thus \SGARefLink{sga7.sga7i.evii.n2-9-4}{2.9.4} shows that
\underline{the category of biextensions of $(P,Q)$ by
$G=\underline G_{m,S}$ is equivalent to the category of birigidified
invertible sheaves on $P\times Q$}; it is therefore a rigid category,
and its set of isomorphism classes identifies with the set of
divisorial correspondences on $(P,Q)$. This identification plainly
respects the natural group laws on the two sets, the first being that
of (\SGARefLink{sga7.sga7i.evii.n2-5-2-2}{2.5.2.2}).

\underline{Remark \SGARefTarget{sga7.sga7i.evii.n2-9-6}{2.9.6}.} Additivity of the first morphism in
(\SGARefLink{sga7.sga7i.evii.n2-9-4-2}{2.9.4.2}) is equivalent to factorization of the second through the
subsheaf $\underline{\operatorname{Ext}}^1(Q,G)$; note that, by \SGARefLink{sga7.sga7i.evii.n1-3-5}{1.3.5},
the canonical morphism
\[
 \underline{\operatorname{Ext}}^1(Q,G)
 \longrightarrow R^1q_*(G_Q)
\]
is injective. Likewise, additivity of the second morphism in
(\SGARefLink{sga7.sga7i.evii.n2-9-4-2}{2.9.4.2}) means that the first factors through
$\underline{\operatorname{Ext}}^1(P,G)$. This gives group
isomorphisms
\begin{equation}\tag{\SGARefTarget{sga7.sga7i.evii.n2-9-6-1}{2.9.6.1}}
 \operatorname{Biext}^{1}(P,Q;G)
 \simeq
 \operatorname{Hom}\bigl(Q,\underline{\operatorname{Ext}}^1(P,G)\bigr)
 \simeq
 \operatorname{Hom}\bigl(P,\underline{\operatorname{Ext}}^1(Q,G)\bigr).
\end{equation}
In the situation of \SGARefLink{sga7.sga7i.evii.n2-9-5}{2.9.5} these reduce to the familiar isomorphisms
\begin{equation}\tag{\SGARefTarget{sga7.sga7i.evii.n2-9-6-2}{2.9.6.2}}
 \operatorname{Corr}(P,Q)
 \simeq \operatorname{Hom}(P,Q')
 \simeq \operatorname{Hom}(Q,P'),
\end{equation}
where
\[
 P'=\underline{\operatorname{Ext}}^1(P,\underline G_m),
 \qquad
 Q'=\underline{\operatorname{Ext}}^1(Q,\underline G_m)
\]
are the dual abelian schemes of $P$ and $Q$, identified respectively
with $\underline{\operatorname{Pic}}^0_{P/S}$ and
$\underline{\operatorname{Pic}}^0_{Q/S}$. In both formulas,
$\operatorname{Hom}$ denotes the group of \underline{group
homomorphisms}.




}{E7I-VII-019}\SGApage{I}{176}
\SGARefTarget{sga7.sga7i.evii.n2-10}{2.10}. \underline{Various generalizations of the notion of
biextension}. We mention them for completeness, although we shall not
use them.

\SGARefTarget{sga7.sga7i.evii.n2-10-1}{2.10.1}. One can define a biextension of $(P,Q)$ by $G$ without
assuming $P$ or $Q$ commutative, only $G$ being commutative. It is a
torsor under $G_{P\times Q}$ equipped with central extension
structures of $Q_P$ by $G_P$ and of $P_Q$ by $G_Q$, each compatible
with the torsor structure and compatible with one another, meaning
that they make (\SGARefLink{sga7.sga7i.evii.n2-1-1}{2.1.1}) commute. Most of the preceding developments
remain valid, \emph{mutatis mutandis}, for this more general notion.
Even when $P,Q$ are commutative, it is more general than the notion of
\SGARefLink{sga7.sga7i.evii.n2-1}{2.1}, which is the case in which the two underlying relative group
structures on $E$ are \underline{commutative}. In the noncommutative
case, the category of biextensions of $(P,Q)$ by $G$ still depends
additively on $G$, but not on $P$ and $Q$.

\SGApage{I}{177}
\SGARefTarget{sga7.sga7i.evii.n2-10-2}{2.10.2}. Paraphrasing \SGARefLink{sga7.sga7i.evii.n2-1}{2.1}, for a fixed integer $n\geq1$ one defines an
\underline{$n$-extension of a sequence of $n$ commutative groups}
$(P_1,P_2,\ldots,P_n)$ by a commutative group $G$. It consists of a
torsor $E$ under $G_{P_1\times\cdots\times P_n}$, equipped with $n$
partial extension structures over the partial products of $n-1$
factors, with a compatibility relation of type (\SGARefLink{sga7.sga7i.evii.n2-1-1}{2.1.1}) for every pair
of distinct indices between $1$ and $n$.

\SGAvar{If $n\geq3$, the $P_i$ are abelian schemes over a scheme $S$, and
$G$ is the multiplicative group over $\underline G_m$
(compare \SGARefLink{sga7.sga7i.evii.n2-9-3}{2.9.2}), every $n$-extension of
$(P_1,\ldots,P_n)$ by $G$ is \underline{trivial}.

}{If $n\geq3$, the $P_i$ are abelian schemes over a scheme $S$, and
$G=\underline G_{m,S}$ is the multiplicative group over $S$
(compare \SGARefLink{sga7.sga7i.evii.n2-9-3}{2.9.3}), every $n$-extension of
$(P_1,\ldots,P_n)$ by $G$ is \underline{trivial}.

}{E7I-VII-020}\SGARefTarget{sga7.sga7i.evii.n2-10-3}{2.10.3}. The notion of biextension in \SGARefLink{sga7.sga7i.evii.n2-1}{2.1} may be viewed as a variant of
the notion of a commutative extension of commutative groups, that is,
of an extension of $\mathbb Z$-modules. Let $A$ be a ring of
$\underline T$, assumed commutative for simplicity, and let $P,Q,G$
be three $A$-modules. More generally one can define an
\underline{$A$-module biextension} of $(P,Q)$ by $G$, a variant of an
extension of $A$-modules.

\SGApage{I}{\(\approx\)178}
Such a biextension is a biextension in the sense of \SGARefLink{sga7.sga7i.evii.n2-1}{2.1} whose
underlying partial extension structures are enriched to extension
structures of $A_P$-modules and $A_Q$-modules, respectively,
compatible with the relative module structures already present on
$Q_P,G_P$ and $P_Q,G_Q$. The two relative module structures on $E$
may be interpreted as two actions of the underlying multiplicative
monoid of $A$ on the sheaf of sets underlying $E$, and these two
actions are required to commute. This condition is added to the
compatibility expressed by (\SGARefLink{sga7.sga7i.evii.n2-1-1}{2.1.1}), and is itself expressed by an
analogous diagram, which we leave to the reader. Finally, the first
action of $A$ on $E$ is required to be distributive with respect not
only to the first composition law on $E$ but also to the second; the
same is required of the second action.

\SGAvar{The dictionary of \SGARefLink{sga7.sga7i.evii.n1-1-6}{1.1.6} also has an $A$-linear variant. An extension
of the $A$-module $P$ by the $A$-module $G$ amounts to a family of
torsors $E_p$ under $G_p$, parametrized by points $p\in P(S)$ for
variable $S\in\operatorname{ob}\underline T$, with $E_p$ varying
\underline{$A$-linearly}, rather than merely additively, in $p$.
Besides the isomorphisms $\varphi_{p,p'}$ of (\SGARefLink{sga7.sga7i.evii.n1-1-3-1}{1.1.3.1}), this is
expressed by isomorphisms
\[
 \mu_{a,p}:\,{}^{a}P_p\xrightarrow{\ \sim\ }P_{ap}
 \qquad(p\in P(S),\ a\in A(S)),
\]
where $\,{}^{a}E_p$ denotes the image of the $G_S$-torsor $P_p$ under
extension of the structure group along
\[
 g\longmapsto ag:G_S\longrightarrow G_S.
\]
We leave to the reader the conditions on $\varphi$ and $\mu$ which
make these data correspond to an extension of $A$-modules $E$ of $P$
by $G$ (cf.\ \SGARefLink{sga7.sga7i.evii.n3-5-4}{3.5.4}(b)).




}{The dictionary of \SGARefLink{sga7.sga7i.evii.n1-1-6}{1.1.6} also has an $A$-linear variant. An extension
of the $A$-module $P$ by the $A$-module $G$ amounts to a family of
torsors $E_p$ under $G_p$, parametrized by points $p\in P(S)$ for
variable $S\in\operatorname{ob}\underline T$, with $E_p$ varying
\underline{$A$-linearly}, rather than merely additively, in $p$.
Besides the isomorphisms $\varphi_{p,p'}$ of (\SGARefLink{sga7.sga7i.evii.n1-1-3-1}{1.1.3.1}), this is
expressed by isomorphisms
\[
 \mu_{a,p}:\,{}^{a}E_p\xrightarrow{\ \sim\ }E_{ap}
 \qquad(p\in P(S),\ a\in A(S)),
\]
where $\,{}^{a}E_p$ denotes the image of the $G_S$-torsor $E_p$ under
extension of the structure group along
\[
 g\longmapsto ag:G_S\longrightarrow G_S.
\]
We leave to the reader the conditions on $\varphi$ and $\mu$ which
make these data correspond to an extension of $A$-modules $E$ of $P$
by $G$ (cf.\ \SGARefLink{sga7.sga7i.evii.n3-5-4}{3.5.4}(b)).




}{E7I-VII-021}\SGApage{I}{179}
\section*{\SGARefTarget{sga7.sga7i.evii.n3}{3}. Supplements on homological algebra}

The principal purpose of this section is to define the isomorphism
\[
 \Biext^1(P,Q;G)\simeq
 \Ext^1(P\overset{\mathbb L}{\otimes}Q,G)
\]
announced above. For this we shall need certain developments in
homological algebra which are of independent interest and will occupy
sections \SGARefLink{sga7.sga7i.evii.n3-1}{3.1} through \SGARefLink{sga7.sga7i.evii.n3-4}{3.4}. The developments in \SGARefLink{sga7.sga7i.evii.n3-1}{3.1} and \SGARefLink{sga7.sga7i.evii.n3-3}{3.3} may be
regarded as supplements to [4], especially to section 6.9 therein.

\subsection*{\SGARefTarget{sga7.sga7i.evii.n3-1}{3.1}. Construction of an additive cofibred category
$\underline\Psi_{L_\bullet}$ in terms of a chain complex $L_\bullet$
in an abelian category $\underline C$}

Let $\underline C$ be an abelian category and let $L_\bullet$ be a
chain complex in $\underline C$,
\begin{equation}\tag{\SGARefTarget{sga7.sga7i.evii.n3-1-1}{3.1.1}}
 L_2\longrightarrow L_1\longrightarrow L_0\longrightarrow0,
\end{equation}
written with lower-degree conventions, so that the differential has
degree $-1$. Thus $L_i=0$ for $i<0$. We shall define a cofibred
category
\[
 \underline\Psi=\underline\Psi_{L_\bullet}
\]
over $\underline C$; up to canonical equivalence of cofibred
categories it is the category of the same name in [4, 6.9]. The
objects of $\underline\Psi_{L_\bullet}$ are pairs
\[
 (X,\alpha),
\]
where $X$ is an extension of $L_0$ by an object $A$ of
$\underline C$, depending on $X$,
\begin{equation}\tag{\SGARefTarget{sga7.sga7i.evii.n3-1-2}{3.1.2}}
 0\longrightarrow A\xrightarrow{i}X\xrightarrow{j}L_0
 \longrightarrow0,
\end{equation}
and $\alpha$ is a trivialization of the inverse image $d_0^*(X)$ of
this extension by $d_0:L_1\longrightarrow L_0$. It is required that
the corresponding trivialization of the inverse image
$d_1^*d_0^*(X)$ by $d_1:L_2\longrightarrow L_1$ be the evident one,
obtained from the transitivity isomorphism
\[
 d_1^*d_0^*(X)\simeq(d_0d_1)^*(X)
\]
and the relation $d_0d_1=0$.

\SGApage{I}{\(\approx\)180}
Equivalently, $\alpha$ is a lifting of $d_0:L_1\longrightarrow L_0$
to a morphism $\alpha:L_1\longrightarrow X$, subject to
$\alpha d_1=0$:
\begin{equation}\tag{\SGARefTarget{sga7.sga7i.evii.n3-1-3}{3.1.3}}
\begin{tikzcd}[row sep=large,column sep=large]
 \&\&0\arrow[d]\\
 \&\&A\arrow[d,"i"]\\
 \&\&X\arrow[d,"j"]\\
 L_2\arrow[r,"d_1"']\&L_1\arrow[r,"d_0"']\arrow[ru,"\alpha"]
   \&L_0\arrow[d]\\
 \&\&0
\end{tikzcd}
\qquad,\qquad \alpha d_1=0.\quad(*)
\end{equation}
Here is a simpler description, pointed out by Deligne:
$\underline\Psi_{L_\bullet}(A)$ is the category of extensions of the
complex $L_\bullet$ by $A$, regarded as a complex concentrated in
degree zero.

In this form one sees that the description of the objects
$(X,\alpha)$ depends only on the chain complex $[L_\bullet]$
\underline{truncated in degree 1}, obtained from $L_\bullet$ by
``killing $H_i(L_\bullet)$ for $i\geq2$'': its terms are zero in
degrees at least $2$, agree with those of $L_\bullet$ in degrees at
most $0$, and its degree-$1$ term is the cokernel of
$d_1:L_2\longrightarrow L_1$.

A homomorphism from $(X,\alpha)$ to $(X',\alpha')$ is a homomorphism
of extensions $u:X\longrightarrow X'$, compatible with the identity
on $L_0$, such that $u\alpha=\alpha'$. Homomorphisms are composed by
composing the corresponding homomorphisms of extensions. This gives
the announced category $\underline\Psi_{L_\bullet}$. Up to canonical
isomorphism, it depends only on the degree-$1$ truncation
$[L_\bullet]$ of $L_\bullet$.

\SGApage{I}{\(\approx\)181}
There is an evident functor $(X,\alpha)\longmapsto A$,
\begin{equation}\tag{\SGARefTarget{sga7.sga7i.evii.n3-1-4}{3.1.4}}
 \pi:\underline\Psi_{L_\bullet}\longrightarrow\underline C.
\end{equation}
\underline{This functor is cofibred}. Indeed, let
$u:A\longrightarrow B$ be a morphism of $\underline C$, and let
$(X,\alpha)$ be an object of the fibre
$\underline\Psi_{L_\bullet}(A)$. Thus $X$ is an extension of $L_0$ by
$A$. Form the extension $u_*(X)$ of $L_0$ by $B$, the direct image of
$X$ under $u$, and equip it with the partial trivialization
$\beta=f\alpha$, where $f:X\longrightarrow u_*(X)$ is the canonical
morphism. It follows immediately that
\[
 f:(X,\alpha)\longrightarrow(u_*(X),\beta)
\]
is cocartesian for (\SGARefLink{sga7.sga7i.evii.n3-1-4}{3.1.4}): it has the usual universal property for
$u$-morphisms in $\underline\Psi_{L_\bullet}$ with source
$(X,\alpha)$. This is just the familiar and elementary corresponding
universal property of $f:X\longrightarrow u_*(X)$.

\SGAvar{Finally, the cofibred category $\underline\Psi_{L_\bullet}$ is
\underline{additive}; that is, it satisfies
\begin{equation}\tag{\SGARefTarget{sga7.sga7i.evii.n3-1-5}{3.1.5}}
 \underline\Psi_{L_\bullet}(0)
 \quad\text{is equivalent to the one-point category},
\end{equation}
and, for arbitrary $A,B$ in $\underline C$,
\begin{equation}\tag{\SGARefTarget{sga7.sga7i.evii.n3-1-6}{3.1.6}}
 \underline\Psi_{L_\bullet}(A\times B)
 \longrightarrow
 \underline\Psi_L(A)\times
 \underline\Psi_{L_\bullet}(B)
 \quad\text{is an equivalence}.
\end{equation}
This follows immediately from the corresponding properties of the
cofibred category of extensions of $L_0$ by variable objects of
$\underline C$. We shall freely use the consequences of the
additivity of $\underline\Psi_{L_\bullet}$ established in [4, 1] and
recalled above in \SGARefLink{sga7.sga7i.evii.n2-5}{2.5}.

}{Finally, the cofibred category $\underline\Psi_{L_\bullet}$ is
\underline{additive}; that is, it satisfies
\begin{equation}\tag{\SGARefTarget{sga7.sga7i.evii.n3-1-5}{3.1.5}}
 \underline\Psi_{L_\bullet}(0)
 \quad\text{is equivalent to the one-point category},
\end{equation}
and, for arbitrary $A,B$ in $\underline C$,
\begin{equation}\tag{\SGARefTarget{sga7.sga7i.evii.n3-1-6}{3.1.6}}
 \underline\Psi_{L_\bullet}(A\times B)
 \longrightarrow
 \underline\Psi_{L_\bullet}(A)\times
 \underline\Psi_{L_\bullet}(B)
 \quad\text{is an equivalence}.
\end{equation}
This follows immediately from the corresponding properties of the
cofibred category of extensions of $L_0$ by variable objects of
$\underline C$. We shall freely use the consequences of the
additivity of $\underline\Psi_{L_\bullet}$ established in [4, 1] and
recalled above in \SGARefLink{sga7.sga7i.evii.n2-5}{2.5}.

}{E7I-VII-022}\SGApage{I}{182}
\SGARefTarget{sga7.sga7i.evii.n3-1-7}{3.1.7}. Among these consequences, for every object $A$ of
$\underline C$ the category $\underline\Psi_{L_\bullet}(A)$ has a
\underline{zero}, or \underline{trivial}, object (\SGARefLink{sga7.sga7i.evii.n2-5-1}{2.5.1}). It is the
trivial extension $X=L_0+A$ of $L_0$ by $A$, equipped with the
partial trivialization $\alpha:L_1\longrightarrow X$ whose components
are $d_0$ and $0$.

\SGAvar{Moreover, by \SGARefLink{sga7.sga7i.evii.n2-5-2}{2.5.2}, for every $(X,\alpha)$ in
$\underline\Psi_{L_\bullet}(A)$ the automorphism group of this object
is canonically isomorphic to the automorphism group of the trivial
object. Denote the latter commutative group by
$\Psi^0_{L_\bullet}(A)$. The explicit description above gives the
canonical isomorphisms
\begin{multline}\tag{\SGARefTarget{sga7.sga7i.evii.n3-1-8}{3.1.8}}
 \Psi^0_{L_\bullet}(A)
 =\Ker\bigl(\Hom(L_0,A)\longrightarrow\Hom(L_1,A)\bigr)
 \simeq\Hom(H_0(L_\bullet),A)\\
 \simeq\Ext^0(L_\bullet,A),
\end{multline}
induced by the usual identification of the automorphism group of the
extension $X$ with $\Hom(L_0,A)$. The isomorphism (\SGARefLink{sga7.sga7i.evii.n3-1-8}{3.1.7.1}) is plainly
functorial in $A$.

}{Moreover, by \SGARefLink{sga7.sga7i.evii.n2-5-2}{2.5.2}, for every $(X,\alpha)$ in
$\underline\Psi_{L_\bullet}(A)$ the automorphism group of this object
is canonically isomorphic to the automorphism group of the trivial
object. Denote the latter commutative group by
$\Psi^0_{L_\bullet}(A)$. The explicit description above gives the
canonical isomorphisms
\begin{multline}\tag{\SGARefTarget{sga7.sga7i.evii.n3-1-8}{3.1.8}}
 \Psi^0_{L_\bullet}(A)
 =\Ker\bigl(\Hom(L_0,A)\longrightarrow\Hom(L_1,A)\bigr)
 \simeq\Hom(H_0(L_\bullet),A)\\
 \simeq\Ext^0(L_\bullet,A),
\end{multline}
induced by the usual identification of the automorphism group of the
extension $X$ with $\Hom(L_0,A)$. The isomorphism (\SGARefLink{sga7.sga7i.evii.n3-1-8}{3.1.8}) is plainly
functorial in $A$.

}{E7I-VII-023}\SGARefTarget{sga7.sga7i.evii.n3-1-9}{3.1.9}. Finally, recall from \SGARefLink{sga7.sga7i.evii.n2-5-2}{2.5.2} that the natural multiplicative
structure on the fibre category $\underline\Psi_{L_\bullet}(A)$
endows its set of isomorphism classes with a natural group structure.
We denote this group by $\Psi^1_{L_\bullet}(A)$; it is an additive
functor of $A$.




\SGApage{I}{183}
\subsection*{\SGARefTarget{sga7.sga7i.evii.n3-2}{3.2}. Determination of $\Psi^1_{L_\bullet}(A)$}

Consider an object $(X,\alpha)$ (\SGARefLink{sga7.sga7i.evii.n3-1-3}{3.1.3}) of
$\underline\Psi_{L_\bullet}(A)$. Consider the cochain complex of
length one
\begin{equation}\tag{\SGARefTarget{sga7.sga7i.evii.n3-2-1}{3.2.1}}
 [X\longrightarrow L_0]:\quad
 0\longrightarrow X\longrightarrow L_0\longrightarrow0\longrightarrow
 \cdots
\end{equation}
defined by the morphism $j:X\longrightarrow L_0$, with $X$ as its
degree-zero component. It is a resolution of $A$, and hence may be
identified with $A$ as an object of the derived category
$D(\underline C)$ [5]. The diagram (\SGARefLink{sga7.sga7i.evii.n3-1-3}{3.1.3}) then gives a degree-one
morphism of complexes
\begin{equation}\tag{\SGARefTarget{sga7.sga7i.evii.n3-2-2}{3.2.2}}
 L_\bullet\longrightarrow[X\longrightarrow L_0],
\end{equation}
made explicit by the following diagram with anticommutative squares:
\begin{equation}\tag{\SGARefTarget{sga7.sga7i.evii.n3-2-3}{3.2.3}}
\begin{tikzcd}[column sep=large,row sep=huge]
 \cdots \& L_2\arrow[r,"d_1"]
 \& L_1\arrow[r,"d_0"]\arrow[d,"u^{-1}=-\alpha"']
 \& L_0\arrow[r]\arrow[d,"u^0=\mathrm{id}"'] \& 0 \& \cdots\\
 \cdots \& 0\arrow[r] \& X\arrow[r,"j"] \& L_0\arrow[r] \& 0 \& \cdots
\end{tikzcd}
\end{equation}
(We use the usual convention which associates with the chain complex
$L_\bullet$ a complex whose differential has degree $1$ and whose
degree-$i$ component is $L_{-i}$.) We thus obtain a degree-one
morphism in the derived category
\begin{equation}\tag{\SGARefTarget{sga7.sga7i.evii.n3-2-3.occ-2}{3.2.3}}
 c(X,\alpha):L_\bullet\longrightarrow A
 \qquad
 \bigl(\text{a degree-one morphism in }D(\underline C)\bigr),
\end{equation}
which plainly depends only on the isomorphism class of $(X,\alpha)$.
We have thereby defined a canonical map
\begin{equation}\tag{\SGARefTarget{sga7.sga7i.evii.n3-2-4}{3.2.4}}
 \Psi^1_{L_\bullet}(A)\longrightarrow\Ext^1(L_\bullet,A)
 \overset{\mathrm{def}}{=}
 \Hom_{D(\underline C)}(L_\bullet,A[1]).
\end{equation}

\SGApage{I}{184}
\underline{Theorem} \SGARefTarget{sga7.sga7i.evii.n3-2-5}{3.2.5}. \underline{The map} (\SGARefLink{sga7.sga7i.evii.n3-2-4}{3.2.4})
\underline{is an isomorphism of groups, functorial in the object}
$A$ \underline{of} $\underline C$.

We leave to the reader the verification that (\SGARefLink{sga7.sga7i.evii.n3-2-4}{3.2.4}) is additive and
functorial in $A$, and shall confine ourselves to proving that it is
bijective. For injectivity, let us first show that if the morphism of
complexes $u$ in (\SGARefLink{sga7.sga7i.evii.n3-2-3}{3.2.3}) is zero in $D(\underline C)$, then it is
already zero in the homotopy category $K(\underline C)$; that is,
there exists a morphism
\[
 s:L_0\longrightarrow X
\]
such that
\begin{equation}\tag{\SGARefTarget{sga7.sga7i.evii.n3-2-5-1}{3.2.5.1}}
 js=u^0=\mathrm{id}_{L_0},
 \qquad sd_0=-u^{-1}=\alpha.
\end{equation}
Indeed, the hypothesis says that the conclusion becomes true after
the morphism of complexes under consideration is composed with a
suitable quasi-isomorphism $v$
\begin{equation}\tag{*}
\begin{tikzcd}[column sep=large,row sep=huge]
 0\arrow[r] \& X\arrow[r,"j"]\arrow[d]
 \& L_0\arrow[r]\arrow[d] \& 0\\
 0\arrow[r] \& C^0\arrow[r,"j'"] \& C^1\arrow[r] \& C^2\arrow[r] \& \cdots
\end{tikzcd}
\end{equation}
(and, after replacing $C^i$ by $0$ for $i\geq2$ and $C^1$ by
$\operatorname{Im}(C^0\longrightarrow C^1)
=\operatorname{Ker}(C^1\longrightarrow C^2)$ if necessary, we may
suppose that $C^i=0$ for $i\geq2$). But the morphism (*) induces the
identity on $A$, and therefore identifies $X$ with the fibre product
of $L_0$ and $C^0$ over $C^1$. A morphism
$s':L_0\longrightarrow C^0$ defining a homotopy $vu\sim0$, and hence
satisfying $j's'=v^1u^0$, consequently factors through a morphism
$s:L_0\longrightarrow X$ satisfying
$js=u^0=\mathrm{id}_{L_0}$. Thus $s$ splits the extension $X$ of
$L_0$ by $A$, so the complex $X\longrightarrow L_0$ is isomorphic
\underline{in} $K(\underline C)$ to the complex $K^\bullet$
concentrated in $A$. For the latter complex it is clear that
\[
 \Hom_{K(\underline C)}(L_\bullet,K^\bullet)
 \longrightarrow
 \Hom_{D(\underline C)}(L_\bullet,K^\bullet)
\]
is an isomorphism, proving the assertion.

\SGApage{I}{\(\approx\)185}
\SGAvar{Now the choice of a morphism $j:L_0\longrightarrow X$ satisfying
(\SGARefLink{sga7.sga7i.evii.n3-2-5-1}{3.2.5.1}) is precisely equivalent to a trivialization of the
extension $X$ of $L_0$ by $A$ compatible with the partial
trivialization $\alpha$. This proves the injectivity of (\SGARefLink{sga7.sga7i.evii.n3-2-4}{3.2.4}).

}{Now the choice of a morphism $s:L_0\longrightarrow X$ satisfying
(\SGARefLink{sga7.sga7i.evii.n3-2-5-1}{3.2.5.1}) is precisely equivalent to a trivialization of the
extension $X$ of $L_0$ by $A$ compatible with the partial
trivialization $\alpha$. This proves the injectivity of (\SGARefLink{sga7.sga7i.evii.n3-2-4}{3.2.4}).

}{E7I-VII-024}For surjectivity, start with an element of the right-hand side. By
definition it may be described by a degree-one morphism of complexes
\begin{equation}\tag{\SGARefTarget{sga7.sga7i.evii.n3-2-5-2}{3.2.5.2}}
\begin{tikzcd}[column sep=large,row sep=huge]
 \ldots \& L_2\arrow[r,"d_1"]
 \& L_1\arrow[r,"d_0"]\arrow[d,"u^{-1}"']
 \& L_0\arrow[r]\arrow[d,"u^0"'] \& 0\\
 \& 0\arrow[r] \& C^0\arrow[r] \& C^1\arrow[r] \& \ldots
\end{tikzcd}
\end{equation}
with anticommutative squares, where $C^\bullet$ is a resolution of
$A$ and, as above, we may suppose that $C^i=0$ for $i\geq2$. Thus
$C^0$ is an extension of $C^1$ by $A$, and we may form the extension
$X$ obtained by pulling it back along $L_0\longrightarrow C^1$.
Anticommutativity in (\SGARefLink{sga7.sga7i.evii.n3-2-5-2}{3.2.5.2}) shows that $-u^{-1}$ factors through a
morphism $\alpha:L_1\longrightarrow X$ which is a partial splitting,
that is, $j\alpha=d_0$ (with $j$ as in (\SGARefLink{sga7.sga7i.evii.n3-1-2}{3.1.2})). Moreover, the
relation $u^{-1}d_1=0$ in (\SGARefLink{sga7.sga7i.evii.n3-2-5-2}{3.2.5.2}) implies
$\alpha d_1=0$. Hence $(X,\alpha)$ is indeed an object of
$\underline\Psi_{L_\bullet}(A)$. By construction, the morphism
(\SGARefLink{sga7.sga7i.evii.n3-2-5-2}{3.2.5.2}) is the composite of the morphism (\SGARefLink{sga7.sga7i.evii.n3-2-3.occ-2}{3.2.3}) associated with
$(X,\alpha)$ and the morphism
\[
 [X\longrightarrow L_0]\longrightarrow C^\bullet
 =[C^0\longrightarrow C^1]
\]
whose components are $(f,u^0)$, where $f:X\longrightarrow C^0$ is
the canonical morphism obtained from the definition of $X$ as the
pullback of the extension $C^0$.

\SGApage{I}{\(\approx\)186}
The preceding morphism is a morphism of resolutions of $A$ and
induces the identity on $A$. It follows, by the definition
$\Ext^1=\Hom^1_{D(\underline C)}$, that the given element of
$\Ext^1(L_\bullet,A)$ described by (\SGARefLink{sga7.sga7i.evii.n3-2-5-2}{3.2.5.2}) is precisely
$c(X,\alpha)$. This proves the surjectivity of (\SGARefLink{sga7.sga7i.evii.n3-2-4}{3.2.4}) and completes
the proof of \SGARefLink{sga7.sga7i.evii.n3-2-5}{3.2.5}.




\subsection*{\SGARefTarget{sga7.sga7i.evii.n3-3}{3.3}. Properties of the cofibred category
$\underline\Psi_{L_\bullet}$}

In this section, which is not essential for understanding what
follows, we freely use the notions of [4], our purpose being to
compare the preceding results with those of that work.

\underline{Theorem} \SGARefTarget{sga7.sga7i.evii.n3-3-1}{3.3.1}. \underline{The additive cofibred category}
$\underline\Psi_{L_\bullet}$ \underline{over} $\underline C$
\underline{is left exact} ([4, \SGARefLink{sga7.sga7i.evii.n2-2}{2.2}]), \underline{admits a
quasi-maximal object} ([4, 4.4]), \underline{and, for every object}
$E=(X,\alpha)$ \underline{of this category, there exists a typical
complex} $L^E_\bullet$ \underline{for} $E$ ([4, \SGARefLink{sga7.sga7i.evii.n3-1}{3.1}]).
\underline{Moreover, the typical pro-complex of}
$\underline\Psi_{L_\bullet}$ ([4, \SGARefLink{sga7.sga7i.evii.n3-7}{3.7} and 4.4])
\underline{is canonically isomorphic in} $D(\underline C)$
\underline{to the truncated complex} $[L_\bullet]$
\underline{considered in} \SGARefLink{sga7.sga7i.evii.n3-1}{3.1}.

We now spell out the meaning of the several assertions in \SGARefLink{sga7.sga7i.evii.n3-3-1}{3.3.1}.

\SGARefTarget{sga7.sga7i.evii.n3-3-2}{3.3.2}. To say that the cofibred category
$\underline\Psi_{L_\bullet}$ is \underline{left exact} means that it
is additive and also satisfies the following condition:

\SGAvar{\SGARefTarget{sga7.sga7i.evii.n3-3-2-1}{3.3.2.1}. For every exact sequence
\[
 0\longrightarrow A'\xrightarrow{u}A\xrightarrow{v}A''
 \longrightarrow0
\]
in $\underline C$, the natural functor (in the notation of [4, 1.8])
\[
 \underline\Psi_{L_\bullet}(A)
 \longrightarrow
 \underline\Psi_{L_\bullet}([A\longrightarrow A''])
\]
is an equivalence of categories.

Recall that if $\underline\Psi$ is a cofibred category over
$\underline C$ and $v:A\longrightarrow A''$ is a morphism in
$\underline C$, then $\underline\Psi([A\longrightarrow A''])$ is the
category of pairs $(E,\rho)$ in which $E$ is an object of
$\underline\Psi(A)$ and $\rho$ is a ``trivialization'' of the
corresponding object $v_*(E)$ of $\underline\Psi(A'')$; that is, an
isomorphism from it to the zero object (\SGARefLink{sga7.sga7i.evii.n2-5-1}{2.5.1}) of
$\underline\Psi(A'')$. The definition of the functor
\[
 \underline\Psi(A')\longrightarrow
 \underline\Psi([A\longrightarrow A''])
\]
in \SGARefLink{sga7.sga7i.evii.n3-3-2-1}{3.3.2.1} is evident. When
$\underline\Psi=\underline\Psi_{L_\bullet}$, the objects of
$\underline\Psi_{L_\bullet}([A\longrightarrow A''])$ are explicitly
the triples $(X,\alpha,\rho)$, where $(X,\alpha)$ is as in (\SGARefLink{sga7.sga7i.evii.n3-1-3}{3.1.3})
and $\rho:X\longrightarrow A''$ is a morphism in $\underline C$
satisfying
\begin{equation}\tag{*}
 \rho i=v,\qquad \rho\alpha=0.
\end{equation}
If $v:A\longrightarrow A''$ belongs to a short exact sequence as in
\SGARefLink{sga7.sga7i.evii.n3-3-2-1}{3.2.3}, then, as is well known, giving a morphism satisfying the
first relation (*) is equivalent, by the left exactness of the
cofibred category of extensions of $L_0$ by variable objects of
$\underline C$, to giving an extension $X'$ of $L_0$ by
$A'\simeq\Ker v$ which is a subextension of $X$: one takes
$X'=\Ker\rho$. The second relation (*) says that $\alpha$ factors
through a morphism $\alpha':L_1\longrightarrow X'$. Up to the
compatibilities whose verification is immediate, this proves the
left exactness stated in \SGARefLink{sga7.sga7i.evii.n3-3-2-1}{3.2.3}.

\SGARefTarget{sga7.sga7i.evii.n3-3-3}{3.3.3}. Recall that an object $E_0$ of a cofibred category
$\underline\Psi$ over $\underline C$ is called \underline{quasi-maximal}
if, for every object $E$ of $\underline\Psi$ lying over an object
$A$ of $\underline C$, there is a morphism $E\longrightarrow E'$ in
$\underline\Psi$ inducing a \underline{monomorphism}
$A\longrightarrow A'$ in $\underline C$, such that $E$ is
\underline{dominated by} $E_0$, meaning that there is a morphism
$E_0\longrightarrow E'$. If $\underline C$ has enough injective
objects, this is equivalent to saying that every object of
$\underline\Psi$ over an injective object of $\underline C$ is
dominated by $E_0$.

For $\underline\Psi=\underline\Psi_{L_\bullet}$ there is the
following canonical quasi-maximal object. Put
\[
 L'_1=\Coker(d_1:L_2\longrightarrow L_1)
\]
and let $E_0$ denote the object $(X_0,\alpha_0)$ of
$\underline\Psi_{L_\bullet}(L'_1)$, where $X_0=L_0\oplus L'_1$ is the trivial extension of $L_0$ by $L'_1$, and where $\alpha_0$ has components
$(d'_1,\mathrm{id}_{L'_1})$, with $d'_1:L_1\longrightarrow L_0$ induced
by $d_1:L_1\longrightarrow L_0$ by passage to the quotient.

}{\SGARefTarget{sga7.sga7i.evii.n3-3-2-1}{3.3.2.1}. For every exact sequence
\[
 0\longrightarrow A'\xrightarrow{u}A\xrightarrow{v}A''
 \longrightarrow0
\]
in $\underline C$, the natural functor (in the notation of [4, 1.8])
\[
 \underline\Psi_{L_\bullet}(A')
 \longrightarrow
 \underline\Psi_{L_\bullet}([A\longrightarrow A''])
\]
is an equivalence of categories.

Recall that if $\underline\Psi$ is a cofibred category over
$\underline C$ and $v:A\longrightarrow A''$ is a morphism in
$\underline C$, then $\underline\Psi([A\longrightarrow A''])$ is the
category of pairs $(E,\rho)$ in which $E$ is an object of
$\underline\Psi(A)$ and $\rho$ is a ``trivialization'' of the
corresponding object $v_*(E)$ of $\underline\Psi(A'')$; that is, an
isomorphism from it to the zero object (\SGARefLink{sga7.sga7i.evii.n2-5-1}{2.5.1}) of
$\underline\Psi(A'')$. The definition of the functor
\[
 \underline\Psi(A')\longrightarrow
 \underline\Psi([A\longrightarrow A''])
\]
in \SGARefLink{sga7.sga7i.evii.n3-3-2-1}{3.3.2.1} is evident. When
$\underline\Psi=\underline\Psi_{L_\bullet}$, the objects of
$\underline\Psi_{L_\bullet}([A\longrightarrow A''])$ are explicitly
the triples $(X,\alpha,\rho)$, where $(X,\alpha)$ is as in (\SGARefLink{sga7.sga7i.evii.n3-1-3}{3.1.3})
and $\rho:X\longrightarrow A''$ is a morphism in $\underline C$
satisfying
\begin{equation}\tag{*}
 \rho i=v,\qquad \rho\alpha=0.
\end{equation}
If $v:A\longrightarrow A''$ belongs to a short exact sequence as in
\SGARefLink{sga7.sga7i.evii.n3-3-2-1}{3.3.2.1}, then, as is well known, giving a morphism satisfying the
first relation (*) is equivalent, by the left exactness of the
cofibred category of extensions of $L_0$ by variable objects of
$\underline C$, to giving an extension $X'$ of $L_0$ by
$A'\simeq\Ker v$ which is a subextension of $X$: one takes
$X'=\Ker\rho$. The second relation (*) says that $\alpha$ factors
through a morphism $\alpha':L_1\longrightarrow X'$. Up to the
compatibilities whose verification is immediate, this proves the
left exactness stated in \SGARefLink{sga7.sga7i.evii.n3-3-2-1}{3.3.2.1}.

\SGARefTarget{sga7.sga7i.evii.n3-3-3}{3.3.3}. Recall that an object $E_0$ of a cofibred category
$\underline\Psi$ over $\underline C$ is called \underline{quasi-maximal}
if, for every object $E$ of $\underline\Psi$ lying over an object
$A$ of $\underline C$, there is a morphism $E\longrightarrow E'$ in
$\underline\Psi$ inducing a \underline{monomorphism}
$A\longrightarrow A'$ in $\underline C$, such that $E'$ is
\underline{dominated by} $E_0$, meaning that there is a morphism
$E_0\longrightarrow E'$. If $\underline C$ has enough injective
objects, this is equivalent to saying that every object of
$\underline\Psi$ over an injective object of $\underline C$ is
dominated by $E_0$.

For $\underline\Psi=\underline\Psi_{L_\bullet}$ there is the
following canonical quasi-maximal object. Put
\[
 L'_1=\Coker(d_1:L_2\longrightarrow L_1)
\]
and let $\bar d_0:L'_1\longrightarrow L_0$ be the morphism induced
by $d_0:L_1\longrightarrow L_0$. Let $E_0=(X_0,\alpha_0)$ be the
object of $\underline\Psi_{L_\bullet}(L'_1)$ for which
$X_0=L_0\oplus L'_1$ is the trivial extension of $L_0$ by $L'_1$
and $\alpha_0$ is the composite
\[
 L_1\twoheadrightarrow L'_1
 \xrightarrow{(\bar d_0,\,\mathrm{id}_{L'_1})}
 L_0\oplus L'_1.
\]

}{E7I-VII-025}\SGApage{I}{\(\approx\)188}
To see that $E_0$ is quasi-maximal, observe that for every object
$(X,\alpha)$ of $\underline\Psi_{L_\bullet}$ one can find a morphism
from it to an object $(X',\alpha')$, inducing a
\underline{monomorphism} $A\longrightarrow A'$ in $\underline C$,
such that the extension $X'$ of $L_0$ by $A'$ is \underline{split}.
Indeed, it suffices to take $A'=X$ and
$(X',\alpha')=i_*(X,\alpha)$, with the notation of (\SGARefLink{sga7.sga7i.evii.n3-1-2}{3.1.2}). On the
other hand, it is immediate that any object $(X',\alpha')$ of
$\underline\Psi$ for which the extension $X'$ is split is dominated
by $E_0$.

\SGARefTarget{sga7.sga7i.evii.n3-3-4}{3.3.4}. Let $E$ be an object of an additive cofibred category
$\underline\Psi$ over $\underline C$, lying above an object $A$ of
$\underline C$, and consider the following functor of the variable
object $B$ of $\underline C$:
\begin{equation}\tag{\SGARefTarget{sga7.sga7i.evii.n3-3-4-1}{3.3.4.1}}
 B\longmapsto\Hom(E,\theta_B),
\end{equation}
where $\theta_B$ denotes the trivial object of $\underline\Psi(B)$.
We say that $E$ \underline{admits a typical complex} if this functor
is representable by an object $\Omega_E$ of $\underline C$. Since
the functor in question maps canonically to the functor
$B\longmapsto\Hom(A,B)$ represented by $A$, one obtains a canonical
morphism
\begin{equation}\tag{\SGARefTarget{sga7.sga7i.evii.n3-3-4-2}{3.3.4.2}}
 A\longrightarrow\Omega_E,
\end{equation}
and the corresponding length-one \underline{chain complex}
$A\longrightarrow\Omega_E$ in $\underline C$ is called the
\underline{typical complex of} $E$.

\SGApage{I}{\(\approx\)189}
\SGAvar{When $\underline\Psi=\underline\Psi_{L_\bullet}$ and
$E=(X,\alpha)$, one checks immediately that the functor (\SGARefLink{sga7.sga7i.evii.n3-3-4-1}{3.3.4.1}) is
represented by $\Coker\alpha$, and that the canonical morphism
(\SGARefLink{sga7.sga7i.evii.n3-3-4-2}{3.3.4.2}) is simply the composite
\[
 A\xrightarrow{j}X\xrightarrow{\mathrm{can}}\Coker\alpha.
\]

}{When $\underline\Psi=\underline\Psi_{L_\bullet}$ and
$E=(X,\alpha)$, one checks immediately that the functor (\SGARefLink{sga7.sga7i.evii.n3-3-4-1}{3.3.4.1}) is
represented by $\Coker\alpha$, and that the canonical morphism
(\SGARefLink{sga7.sga7i.evii.n3-3-4-2}{3.3.4.2}) is simply the composite
\[
 A\xrightarrow{i}X\xrightarrow{\mathrm{can}}\Coker\alpha.
\]

}{E7I-VII-026}\SGARefTarget{sga7.sga7i.evii.n3-3-5}{3.3.5}. If an additive cofibred category $\underline\Psi$ over
$\underline C$ admits a quasi-maximal object (\SGARefLink{sga7.sga7i.evii.n3-3-3}{3.3.3}), and the typical
complexes of its objects exist, then the typical complex $T_\bullet$
of a quasi-maximal object $E_0$, regarded as an object of the
\underline{derived category} $D(\underline C)$, is independent, up
to canonical isomorphism, of the choice of $E_0$ ([4, 4.4]). It is
called the \underline{typical complex of the additive cofibred
category} $\underline\Psi$. If, in addition, $\underline\Psi$ is
left exact (\SGARefLink{sga7.sga7i.evii.n3-3-2}{3.3.2}), then $\underline\Psi$ can be reconstructed from
the typical complex $T_\bullet$, up to an equivalence of fibred
categories over $\underline C$ unique up to unique isomorphism
([4, 6.10]). There are also canonical isomorphisms, functorial in
$A\in\underline C$ ([4, 5.4]),
\begin{equation}\tag{\SGARefTarget{sga7.sga7i.evii.n3-3-5-1}{3.3.5.1}}
 \Psi^0(A)\xrightarrow{\sim}\Ext^0(T_\bullet,A),
 \qquad
 \Psi^1(A)\xrightarrow{\sim}\Ext^1(T_\bullet,A).
\end{equation}

To complete the comparison with [4], we shall prove that when
$\underline\Psi=\underline\Psi_{L_\bullet}$ its typical complex is
the truncation $[L_\bullet]$. This gives another definition of the
isomorphisms (\SGARefLink{sga7.sga7i.evii.n3-1-8}{3.1.8}) and (\SGARefLink{sga7.sga7i.evii.n3-2-4}{3.2.4}) in terms of (\SGARefLink{sga7.sga7i.evii.n3-3-5-1}{3.3.5.1}), and one
checks that it is compatible with the definitions in \SGARefLink{sga7.sga7i.evii.n3-1}{3.1} and \SGARefLink{sga7.sga7i.evii.n3-2}{3.2}.

\SGApage{I}{\(\approx\)190}
\SGAvar{For this assertion we may plainly restrict to the case where
$L_\bullet$ is already truncated in degree $1$, since \SGARefLink{sga7.sga7i.evii.n3-1}{3.1} noted that
truncation does not change $\underline\Psi_{L_\bullet}$ up to an
isomorphism of cofibred categories. The explicit construction in
\SGARefLink{sga7.sga7i.evii.n3-3-4}{3.3.4} then shows immediately that the typical complex of the
quasi-maximal object described in \SGARefLink{sga7.sga7i.evii.n3-3-3}{3.3.3)} is canonically isomorphic to
$L_\bullet$.

\underline{Remark} \SGARefTarget{sga7.sga7i.evii.n3-3-6}{3.3.6}. Let
\[
 0\longrightarrow A'\longrightarrow A\longrightarrow A''
 \longrightarrow0
\]
be an exact sequence in $\underline C$. It gives an exact sequence
of $\Ext^i$ groups
\[
 0\longrightarrow\Ext^0(L_\bullet,A')
 \longrightarrow\Ext^0(L_\bullet,A)
 \longrightarrow\cdots\longrightarrow\Ext^1(L_\bullet,A)
 \longrightarrow\Ext^1(L_\bullet,A''),
\]
which could of course be continued to the right by introducing
$\Ext^i$ for $i\geq2$. Using the isomorphisms (\SGARefLink{sga7.sga7i.evii.n3-1-8}{3.1.8}) and (\SGARefLink{sga7.sga7i.evii.n3-2-4}{3.2.4}),
we obtain the exact sequence
\begin{multline}\tag{\SGARefTarget{sga7.sga7i.evii.n3-3-6-1}{3.3.6.1}}
 0\longrightarrow\Psi^0_{L_\bullet}(A')
 \longrightarrow\Psi^0_{L_\bullet}(A)
 \longrightarrow\Psi^0_{L_\bullet}(A'')
 \xrightarrow{o}\Psi^1_{L_\bullet}(A')\\
 \longrightarrow\Psi^1_{L_\bullet}(A)
 \longrightarrow\Psi^1_{L_\bullet}(A'').
\end{multline}
All arrows other than $\partial$ are defined simply by functoriality
of $\Psi^0_{L_\bullet}$ and $\Psi^1_{L_\bullet}$. On the other hand,
[4, \SGARefLink{sga7.sga7i.evii.n2-5-2}{2.5.2}] defines an exact sequence (3.3.6.1), for every left-exact
cofibred category $\underline\Psi$ over $\underline C$, through a
morphism
\begin{equation}\tag{\SGARefTarget{sga7.sga7i.evii.n3-3-6-2}{3.3.6.2}}
 \partial:\Psi^0(A'')\longrightarrow\Psi^1(A').
\end{equation}
It associates with an element $u$ of the left-hand side the class,
in the right-hand side, of the object of $\underline\Psi(A')$,
determined up to canonical isomorphism, whose image in
$\underline\Psi([A\longrightarrow A''])$ (notation of 3.3.2) is the
pair $(X_0,u)$, where $X_0$ is the trivial object of
$\underline\Psi(A)$.

}{For this assertion we may plainly restrict to the case where
$L_\bullet$ is already truncated in degree $1$, since \SGARefLink{sga7.sga7i.evii.n3-1}{3.1} noted that
truncation does not change $\underline\Psi_{L_\bullet}$ up to an
isomorphism of cofibred categories. The explicit construction in
\SGARefLink{sga7.sga7i.evii.n3-3-4}{3.3.4} then shows immediately that the typical complex of the
quasi-maximal object described in \SGARefLink{sga7.sga7i.evii.n3-3-3}{3.3.3} is canonically isomorphic to
$L_\bullet$.

\underline{Remark} \SGARefTarget{sga7.sga7i.evii.n3-3-6}{3.3.6}. Let
\[
 0\longrightarrow A'\longrightarrow A\longrightarrow A''
 \longrightarrow0
\]
be an exact sequence in $\underline C$. It gives an exact sequence
of $\Ext^i$ groups
\[
 0\longrightarrow\Ext^0(L_\bullet,A')
 \longrightarrow\Ext^0(L_\bullet,A)
 \longrightarrow\cdots\longrightarrow\Ext^1(L_\bullet,A)
 \longrightarrow\Ext^1(L_\bullet,A''),
\]
which could of course be continued to the right by introducing
$\Ext^i$ for $i\geq2$. Using the isomorphisms (\SGARefLink{sga7.sga7i.evii.n3-1-8}{3.1.8}) and (\SGARefLink{sga7.sga7i.evii.n3-2-4}{3.2.4}),
we obtain the exact sequence
\begin{multline}\tag{\SGARefTarget{sga7.sga7i.evii.n3-3-6-1}{3.3.6.1}}
 0\longrightarrow\Psi^0_{L_\bullet}(A')
 \longrightarrow\Psi^0_{L_\bullet}(A)
 \longrightarrow\Psi^0_{L_\bullet}(A'')
 \xrightarrow{\partial}\Psi^1_{L_\bullet}(A')\\
 \longrightarrow\Psi^1_{L_\bullet}(A)
 \longrightarrow\Psi^1_{L_\bullet}(A'').
\end{multline}
All arrows other than $\partial$ are defined simply by functoriality
of $\Psi^0_{L_\bullet}$ and $\Psi^1_{L_\bullet}$. On the other hand,
[4, \SGARefLink{sga7.sga7i.evii.n2-5-2}{2.5.2}] defines an exact sequence (3.3.6.1), for every left-exact
cofibred category $\underline\Psi$ over $\underline C$, through a
morphism
\begin{equation}\tag{\SGARefTarget{sga7.sga7i.evii.n3-3-6-2}{3.3.6.2}}
 \partial:\Psi^0(A'')\longrightarrow\Psi^1(A').
\end{equation}
It associates with an element $u$ of the left-hand side the class,
in the right-hand side, of the object of $\underline\Psi(A')$,
determined up to canonical isomorphism, whose image in
$\underline\Psi([A\longrightarrow A''])$ (notation of 3.3.2) is the
pair $(X_0,u)$, where $X_0$ is the trivial object of
$\underline\Psi(A)$.

}{E7I-VII-027}\SGApage{I}{\(\approx\)191}
For $\underline\Psi=\underline\Psi_{L_\bullet}$, the morphism
(\SGARefLink{sga7.sga7i.evii.n3-3-6-2}{3.3.6.2}) defined by this \underline{``geometric'' construction is
the negative of the boundary morphism} $\partial$ \underline{in}
(\SGARefLink{sga7.sga7i.evii.n3-3-6-1}{3.3.6.1}). Thus, after the identifications above and a change of
sign in $\partial$, the latter exact sequence is precisely that of
[4, \SGARefLink{sga7.sga7i.evii.n2-5-2}{2.5.2}]. This follows from the final assertion of [4, 5.4], in
view of the compatibilities noted in \SGARefLink{sga7.sga7i.evii.n3-3-5}{3.3.5}.

\underline{Remark} \SGARefTarget{sga7.sga7i.evii.n3-3-6-3}{3.3.6.3}. We take this opportunity to clarify a
point in [4, no.~3.11] concerning the fundamental isomorphisms
(3.11.3) and (3.11.4), which play an essential role in the proof of
[4, 5.4]. The diagram on page 35 of [4] is commutative, not
anticommutative, and therefore does not as it stands define a
morphism
\[
 u:[L^X_1\longrightarrow L^X_0]\longrightarrow K^\bullet[1].
\]
It must be made anticommutative by replacing $u^1$ by $-u^1$. One
could instead adopt the opposite convention and replace $u^0$ by
$-u^0$; however, the stated convention, in which $u^1$ becomes
$-u^1$, is the one for which (3.11.3) generalizes the isomorphism
(\SGARefLink{sga7.sga7i.evii.n3-5-3}{3.5.3}) of [4], as claimed there. With this clarification, the claim
made in [4, 5.4] must also be corrected: the two boundary morphisms
considered there \underline{differ by a sign}.

\SGApage{I}{192}
\SGAvar{\SGARefTarget{sga7.sga7i.evii.n3-3-7}{3.3.7}. In [4, 6.13] the variance of the cofibred category
$\underline\Psi_{L_\bullet}$ over $\underline C$ as $L_\bullet$
varies was also made explicit, in particular by reconstructing the
category of cocartesian functors from $\underline\Psi_{L_\bullet}$
to $\underline\Psi_{L'_\bullet}$ in terms of the chain complexes
$L_\bullet$ and $L'_\bullet$. The set of isomorphism classes of such
cocartesian functors is canonically in bijection with the set of isomorphisms
\[
 \Hom_{D(\underline C)}(L'_\bullet,L_\bullet).
\]
We shall use only a small part of these results, which we develop in
the next section using the simple description of
$\underline\Psi_{L_\bullet}$ given in \SGARefLink{sga7.sga7i.evii.n3-1}{3.1}.



}{\SGARefTarget{sga7.sga7i.evii.n3-3-7}{3.3.7}. In [4, 6.13] the variance of the cofibred category
$\underline\Psi_{L_\bullet}$ over $\underline C$ as $L_\bullet$
varies was also made explicit, in particular by reconstructing the
category of cocartesian functors from $\underline\Psi_{L_\bullet}$
to $\underline\Psi_{L'_\bullet}$ in terms of the chain complexes
$L_\bullet$ and $L'_\bullet$. The set of isomorphism classes of such
cocartesian functors is canonically in bijection with
\[
 \Hom_{D(\underline C)}(L'_\bullet,L_\bullet).
\]
We shall use only a small part of these results, which we develop in
the next section using the simple description of
$\underline\Psi_{L_\bullet}$ given in \SGARefLink{sga7.sga7i.evii.n3-1}{3.1}.



}{E7I-VII-028}\subsection*{\SGARefTarget{sga7.sga7i.evii.n3-4}{3.4}. Variance of $\underline\Psi_{L_\bullet}$ with
$L_\bullet$}

Consider a morphism of complexes
\begin{equation}\tag{\SGARefTarget{sga7.sga7i.evii.n3-4-1}{3.4.1}}
 f:L'_\bullet\longrightarrow L_\bullet
\end{equation}
given by a commutative diagram
\begin{equation}\tag{\SGARefTarget{sga7.sga7i.evii.n3-4-2}{3.4.2}}
\begin{tikzcd}[column sep=large,row sep=large]
 {}\arrow[r]\&L'_2\arrow[r]\arrow[d,"f_2"']
 \&L'_1\arrow[r]\arrow[d,"f_1"']
 \&L'_0\arrow[r]\arrow[d,"f_0"']\&0\\
 {}\arrow[r]\&L_2\arrow[r]\&L_1\arrow[r]\&L_0\arrow[r]\&0.
\end{tikzcd}
\end{equation}
We associate with it a canonical functor
\begin{equation}\tag{\SGARefTarget{sga7.sga7i.evii.n3-4-3}{3.4.3}}
 f^*:\underline\Psi_{L_\bullet}\longrightarrow
 \underline\Psi_{L'_\bullet}.
\end{equation}
It is in fact a cocartesian functor of cofibred categories over
$\underline C$: it is compatible with the structural functors to
$\underline C$ and carries cocartesian morphisms to cocartesian
morphisms; see SGA 1, Exposé VI.

\SGAvar{If $(X,\alpha)$ is an object of $\underline\Psi_L$, define
$f^*(X,\alpha)=(X',\alpha')$ as follows. The extension $X'$ of
$L'_0$ by $A$ is the pullback, along $f_0$, of the given extension
$X$ of $L_0$ by $A$, and $\alpha'$ is the partial trivialization
deduced from $\alpha$ by transitivity in the second square of
(\SGARefLink{sga7.sga7i.evii.n3-4-2}{3.4.2}). The compatibility of $\alpha'$ with the relation
$d'_0d'_1=0$ follows immediately from the analogous compatibility
for $\alpha$. By construction, (\SGARefLink{sga7.sga7i.evii.n3-4-3}{3.4.3}) commutes with the structural
functors to $\underline C$. Its cocartesian character follows at
once from the fact that, in the category of extensions of objects of
$\underline C$, direct images of extensions commute with inverse
images of extensions.

The functor (\SGARefLink{sga7.sga7i.evii.n3-4-3}{3.4.3}) induces morphisms of functors
\begin{equation}\tag{\SGARefTarget{sga7.sga7i.evii.n3-4-4}{3.4.4}}
 f^i:\Psi^i_{L_\bullet}\longrightarrow\Psi^i_{L'_\bullet},
 \qquad i=0,1
\end{equation}
(compare \SGARefLink{sga7.sga7i.evii.n2-6}{2.6}). The morphism (\SGARefLink{sga7.sga7i.evii.n3-4-1}{3.4.1}), on the other hand, induces
morphisms
\begin{equation}\tag{\SGARefTarget{sga7.sga7i.evii.n3-4-5}{3.4.5}}
 f^i:\Ext^i(L_\bullet,-)\qquad\Ext^i(L'_\bullet,-),
 \qquad i\in\mathbb Z.
\end{equation}
For $i=0,1$, the morphisms (\SGARefLink{sga7.sga7i.evii.n3-4-4}{3.4.4}) and (\SGARefLink{sga7.sga7i.evii.n3-4-5}{3.4.5}) are compatible with
the isomorphisms (\SGARefLink{sga7.sga7i.evii.n3-1-8}{3.1.8}) and (\SGARefLink{sga7.sga7i.evii.n3-2-4}{3.2.4}). In other words, for
$A\in\operatorname{Ob}\underline C$ and $i=0,1$, the following
diagram is commutative:
\begin{equation}\tag{\SGARefTarget{sga7.sga7i.evii.n3-4-6}{3.4.6}}
\begin{tikzcd}[column sep=large,row sep=large]
 \Psi^i_{L_\bullet}(A)\arrow[r,"\sim"]\arrow[d]
 \&\Ext^i(L_\bullet,A)\arrow[d]\\
 \Psi^i_{L'_\bullet}(A)\arrow[r,"\sim"]
 \&\Ext^i(L'_\bullet,A).
\end{tikzcd}
\end{equation}
The verification of this compatibility is essentially immediate and
is left to the reader.

}{If $(X,\alpha)$ is an object of $\underline\Psi_{L_\bullet}$, define
$f^*(X,\alpha)=(X',\alpha')$ as follows. The extension $X'$ of
$L'_0$ by $A$ is the pullback, along $f_0$, of the given extension
$X$ of $L_0$ by $A$, and $\alpha'$ is the partial trivialization
deduced from $\alpha$ by transitivity in the second square of
(\SGARefLink{sga7.sga7i.evii.n3-4-2}{3.4.2}). The compatibility of $\alpha'$ with the relation
$d'_0d'_1=0$ follows immediately from the analogous compatibility
for $\alpha$. By construction, (\SGARefLink{sga7.sga7i.evii.n3-4-3}{3.4.3}) commutes with the structural
functors to $\underline C$. Its cocartesian character follows at
once from the fact that, in the category of extensions of objects of
$\underline C$, direct images of extensions commute with inverse
images of extensions.

The functor (\SGARefLink{sga7.sga7i.evii.n3-4-3}{3.4.3}) induces morphisms of functors
\begin{equation}\tag{\SGARefTarget{sga7.sga7i.evii.n3-4-4}{3.4.4}}
 f^i:\Psi^i_{L_\bullet}\longrightarrow\Psi^i_{L'_\bullet},
 \qquad i=0,1
\end{equation}
(compare \SGARefLink{sga7.sga7i.evii.n2-6}{2.6}). The morphism (\SGARefLink{sga7.sga7i.evii.n3-4-1}{3.4.1}), on the other hand, induces
morphisms
\begin{equation}\tag{\SGARefTarget{sga7.sga7i.evii.n3-4-5}{3.4.5}}
 f^i:\Ext^i(L_\bullet,-)\longrightarrow\Ext^i(L'_\bullet,-),
 \qquad i\in\mathbb Z.
\end{equation}
For $i=0,1$, the morphisms (\SGARefLink{sga7.sga7i.evii.n3-4-4}{3.4.4}) and (\SGARefLink{sga7.sga7i.evii.n3-4-5}{3.4.5}) are compatible with
the isomorphisms (\SGARefLink{sga7.sga7i.evii.n3-1-8}{3.1.8}) and (\SGARefLink{sga7.sga7i.evii.n3-2-4}{3.2.4}). In other words, for
$A\in\operatorname{Ob}\underline C$ and $i=0,1$, the following
diagram is commutative:
\begin{equation}\tag{\SGARefTarget{sga7.sga7i.evii.n3-4-6}{3.4.6}}
\begin{tikzcd}[column sep=large,row sep=large]
 \Psi^i_{L_\bullet}(A)\arrow[r,"\sim"]\arrow[d]
 \&\Ext^i(L_\bullet,A)\arrow[d]\\
 \Psi^i_{L'_\bullet}(A)\arrow[r,"\sim"]
 \&\Ext^i(L'_\bullet,A).
\end{tikzcd}
\end{equation}
The verification of this compatibility is essentially immediate and
is left to the reader.

}{E7I-VII-029}\SGApage{I}{194}
\underline{Proposition} \SGARefTarget{sga7.sga7i.evii.n3-4-7}{3.4.7}. \underline{The functor} (\SGARefLink{sga7.sga7i.evii.n3-4-3}{3.4.3})
\underline{is faithful} (\underline{respectively fully faithful,
respectively an equivalence of categories}) \underline{if and only
if}
\[
 H_0(f):H_0(L'_\bullet)\longrightarrow H_0(L_\bullet)
\]
\underline{is an epimorphism} (\underline{respectively, if}
$H_0(f)$ \underline{is an isomorphism and} $H_1(f)$
\underline{is an epimorphism; respectively, if} $H_0(f)$
\underline{and} $H_1(f)$ \underline{are isomorphisms}).

\SGAvar{Indeed, with $f^i$ denoting the functorial morphism (\SGARefLink{sga7.sga7i.evii.n3-4-4}{3.4.4}), the
required condition is that $f^0$ be a monomorphism (respectively,
that $f^0$ be an isomorphism and $f$ a monomorphism; respectively,
that $f^0$ and $f^1$ be isomorphisms). By the commutativity of
(\SGARefLink{sga7.sga7i.evii.n3-4-6}{3.4.6}), these conditions may be expressed by the analogous
conditions on the functors $\Ext^i$, $i=0,1$. The proposition then
reduces to a straightforward statement of homological algebra,
which we leave to the reader as an exercise.




\SGARefTarget{sga7.sga7i.evii.n3-4-8}{3.4.8}. Consider now an exact sequence of chain complexes
\begin{equation}\tag{\SGARefTarget{sga7.sga7i.evii.n3-4-8-1}{3.4.8.1}}
 0\longrightarrow L_\bullet\xrightarrow{f}L'_\bullet
 \xrightarrow{g}L''_\bullet\longrightarrow0.
\end{equation}
It gives functors of the form (\SGARefLink{sga7.sga7i.evii.n3-4-3}{3.4.3}),
\[
 \underline\Psi_{L''_\bullet}\xrightarrow{g^*}
 \underline\Psi_{L'_\bullet}\xrightarrow{f^*}
 \underline\Psi_{L_\bullet}.
\]
The transitivity property implicit in functors of the form (\SGARefLink{sga7.sga7i.evii.n3-4-3}{3.4.3})
gives an isomorphism of functors
\[
 f^*g^*\simeq(gf)^*=(0_{L_\bullet\to L''_\bullet})^*.
\]
The last member is evidently canonically isomorphic to the
cocartesian zero functor from $\underline\Psi_{L''_\bullet}$ to
$\underline\Psi_{L_\bullet}$.  The preceding isomorphism therefore
defines a canonical functor
\begin{equation}\tag{\SGARefTarget{sga7.sga7i.evii.n3-4-8-2}{3.4.8.2}}
 \underline\Psi_{L''_\bullet}\longrightarrow
 \underline\Psi_{[L_\bullet\longrightarrow L'_\bullet]}.
\end{equation}
Here $\underline\Psi_{[L_\bullet\longrightarrow L'_\bullet]}$
denotes the category of pairs $(X',t)$ in which $X'$ is an object of
$\underline\Psi_{L'_\bullet}$ and $t$ is a trivialization of its
inverse image $f^*(X')$, that is, an isomorphism
\[
 f^*(X')\simeq0_{L_\bullet,A'},
\]
where $0_{L_\bullet,A}$ denotes the zero object of
$\underline\Psi_{L_\bullet}(A)$ and $A$ is the object of
$\underline C$ underlying $X'$.  We claim that \underline{the
functor} (\SGARefLink{sga7.sga7i.evii.n3-4-8-2}{3.4.8.2}) \underline{is an equivalence of categories}.
Together with the left exactness in $A$ established in \SGARefLink{sga7.sga7i.evii.n3-3-1}{3.3.1}, this
completes the exactness property in $L_\bullet$ of the expression
$\underline\Psi_{L_\bullet}(A)$.

Let us sketch the proof.  To give an object of the target of
(\SGARefLink{sga7.sga7i.evii.n3-4-8-2}{3.4.8.2}) is to give an extension
\[
 0\longrightarrow A\xrightarrow{i'}X'\xrightarrow{j'}L'_0
 \longrightarrow0
\]
of $L'_0$ by an object $A$ of $\underline C$, together with a
morphism $\alpha':L'_1\longrightarrow X'$ satisfying
$j\alpha'=d'_0$ and $\alpha'd'_1=0$, and a morphism
$t:L_0\longrightarrow X'$ satisfying $jt=f_0$ and
$td_0=\alpha'f_1$.  These data are displayed in the following
diagram:
\begin{center}
\begin{tikzpicture}[
  x=1cm,y=1cm,>=Latex,
  every node/.style={font=\small,inner sep=1.2pt}
]
  \node (z00) at (0,0) {$0$};
  \node (l0)  at (1.3,0) {$L_0$};
  \node (lp0) at (2.8,0) {$L'_0$};
  \node (lpp0) at (4.3,0) {$L''_0$};
  \node (z05) at (5.5,0) {$0$};
  \node (z10) at (0,-1.35) {$0$};
  \node (l1)  at (1.3,-1.35) {$L_1$};
  \node (lp1) at (2.8,-1.35) {$L'_1$};
  \node (lpp1) at (4.3,-1.35) {$L''_1$};
  \node (z15) at (5.5,-1.35) {$0$};
  \node (z20) at (0,-2.70) {$0$};
  \node (l2)  at (1.3,-2.70) {$L_2$};
  \node (lp2) at (2.8,-2.70) {$L'_2$};
  \node (lpp2) at (4.3,-2.70) {$L''_2$};
  \node (z25) at (5.5,-2.70) {$0$};
  \node (xp) at (3.25,1.25) {$X'$};
  \node (ap) at (4.0,2.25) {$A$};
  \node (zap) at (4.75,3.25) {$0$};
  \node (xpp) at (5.85,1.25) {$X''$};
  \node (app) at (6.60,2.25) {$A$};
  \node (zapp) at (7.35,3.25) {$0$};
  \draw[->] (z00)--(l0);
  \draw[->] (l0)--node[above]{$f_0$}(lp0);
  \draw[->] (lp0)--node[above]{$g_0$}(lpp0);
  \draw[->] (lpp0)--(z05);
  \draw[->] (z10)--(l1);
  \draw[->] (l1)--node[above]{$f_1$}(lp1);
  \draw[->] (lp1)--node[above]{$g_1$}(lpp1);
  \draw[->] (lpp1)--(z15);
  \draw[->] (z20)--(l2);
  \draw[->] (l2)--(lp2);
  \draw[->] (lp2)--(lpp2);
  \draw[->] (lpp2)--(z25);
  \draw[->] (l1)--node[right]{$d_0$}(l0);
  \draw[->] (lp1)--node[right]{$d'_0$}(lp0);
  \draw[->] (lpp1)--node[right]{$d''_0$}(lpp0);
  \draw[->] (l2)--(l1);
  \draw[->] (lp2)--node[right]{$d'_1$}(lp1);
  \draw[->] (lpp2)--node[right]{$d''_1$}(lpp1);
  \draw[->] (zap)--(ap);
  \draw[->] (ap)--node[left]{$i'$}(xp);
  \draw[->,dashed] (zapp)--(app);
  \draw[->,dashed] (app)--node[left]{$i''$}(xpp);
  \draw[->] (l0) to[bend left=22] node[above]{$t$} (xp);
  \draw[->] (xp)--node[left]{$j'$}(lp0);
  \draw[->] (lp1)--node[pos=.82,right]{$\alpha'$}(xp);
  \draw[->,dashed] (xpp)--node[above]{$j''$}(lpp0);
  \draw[->,dashed] (lpp1) to[bend right=25]
    node[pos=.82,right]{$\alpha''$} (xpp);
  \node[anchor=west] at (7.55,-0.70) {$
    \left\{
    \begin{aligned}
      j'\alpha'\&=d'_0,\\
      \alpha'd'_1\&=0,\\
      j't\&=f_0,\\
      td_0\&=\alpha'f_1.
    \end{aligned}
    \right.$};
\end{tikzpicture}
\end{center}

The familiar left exactness, in the first argument, of extensions
of an object $L$ by an object $A$, applied to the first row of the
diagram, shows that a morphism $t$ satisfying $jt=f_0$ allows one
to interpret $X'$ as the inverse image of an extension $X''$ of
$L''_0$ by $A$.  Under this interpretation, $t$ corresponds to the
natural trivialization of the inverse image of that extension along
$g_0f_0=0$.

The morphism $\alpha'$ satisfying $j\alpha'=d'_0$ and
$\alpha'f_1=td_0$ is interpreted by essentially the same dictionary,
now applied to the second row of the diagram, as arising from a
morphism $\alpha'':L''_1\longrightarrow X''$ satisfying
$j''\alpha''=d''_0$.  Here $\alpha'$ and $\alpha''$ are regarded as
trivializations of the inverse images of $X'$ and $X''$ along
$d'_0$ and $d''_0$, respectively.  Finally, the relation
$\alpha'd'_1=0$ is then equivalent to $\alpha''d''_1=0$.  Thus the
data $(X',\alpha',t)$, an object of
$\underline\Psi_{[L_\bullet\longrightarrow L'_\bullet]}$, are
equivalent to the data $(X'',\alpha'')$, an object of
$\underline\Psi_{L_\bullet}$.  This proves, in substance, that
(\SGARefLink{sga7.sga7i.evii.n3-4-8-2}{3.4.8.2}) is indeed an equivalence of categories.

\SGARefTarget{sga7.sga7i.evii.n3-4-9}{3.4.9}. Continue to consider the exact sequence (\SGARefLink{sga7.sga7i.evii.n3-4-8-1}{3.4.8.1}), and the
exact sequence of $\Ext^i$ groups that it defines:
\begin{multline}\tag{\SGARefTarget{sga7.sga7i.evii.n3-4-9-1}{3.4.9.1}}
 0\longrightarrow\Ext^0(L''_\bullet,A)
 \xrightarrow{g^0}\Ext^0(L'_\bullet A)
 \xrightarrow{f^0}\Ext^0(L_\bullet,A)
 \xrightarrow{\partial^0}\Ext^1(L''_\bullet,A)\\
 \xrightarrow{g^1}\Ext^1(L'_\bullet,A)
 \xrightarrow{f^1}\Ext^1(L_\bullet,A)\longrightarrow\cdots.
\end{multline}
The canonical isomorphisms of the forms (\SGARefLink{sga7.sga7i.evii.n3-1-8}{3.1.8}) and (\SGARefLink{sga7.sga7i.evii.n3-2-4}{3.2.4}) give a
geometric interpretation of the first six terms of this exact
sequence in terms of the categories $\underline\Psi_{L_\bullet}$,
$\underline\Psi_{L'_\bullet}$, and
$\underline\Psi_{L''_\bullet}$.  Compatibility (\SGARefLink{sga7.sga7i.evii.n3-4-6}{3.4.6}) likewise
gives a geometric interpretation of the arrows
$f^0,g^0,f^1,g^1$.  We shall now give such an interpretation of
$\partial^0$, using the result just proved that (\SGARefLink{sga7.sga7i.evii.n3-4-8-2}{3.4.8.2}) is an
equivalence of categories.

Let
\[
 u\in\Ext^0(L_\bullet,A),
 \qquad\text{i.e.}\qquad
 u:L_0\longrightarrow A,\quad ud_0=0.
\]
Consider the object $(0,\overline u)$ of
$\underline\Psi_{[L_\bullet\to L']}$.  Here $0$ is the
zero object of $\underline\Psi_{L'_\bullet}$, whose inverse image
along $f_0$ is therefore identified with the zero object of
$\underline\Psi_{L_\bullet}$, and $\overline u$ is the automorphism
of the latter defined by $u$ through the canonical isomorphism
(\SGARefLink{sga7.sga7i.evii.n3-1-8}{3.1.8}).  Since (\SGARefLink{sga7.sga7i.evii.n3-4-8-2}{3.4.8.2}) is an equivalence, $(0,\overline u)$ comes
from an object of $\underline\Psi_{L''_\bullet}(A)$, determined up
to unique isomorphism.  Its isomorphism class is an element of
\[
 \Psi^1_{L''_\bullet}(A)\simeq\Ext^1(L''_\bullet,A).
\]
\underline{We claim that this element is precisely}
$\partial^0(u)$.

}{Indeed, with $f^i$ denoting the functorial morphism (\SGARefLink{sga7.sga7i.evii.n3-4-4}{3.4.4}), the
required condition is that $f^0$ be a monomorphism (respectively,
that $f^0$ be an isomorphism and $f^1$ a monomorphism; respectively,
that $f^0$ and $f^1$ be isomorphisms). By the commutativity of
(\SGARefLink{sga7.sga7i.evii.n3-4-6}{3.4.6}), these conditions may be expressed by the analogous
conditions on the functors $\Ext^i$, $i=0,1$. The proposition then
reduces to a straightforward statement of homological algebra,
which we leave to the reader as an exercise.




\SGARefTarget{sga7.sga7i.evii.n3-4-8}{3.4.8}. Consider now an exact sequence of chain complexes
\begin{equation}\tag{\SGARefTarget{sga7.sga7i.evii.n3-4-8-1}{3.4.8.1}}
 0\longrightarrow L_\bullet\xrightarrow{f}L'_\bullet
 \xrightarrow{g}L''_\bullet\longrightarrow0.
\end{equation}
It gives functors of the form (\SGARefLink{sga7.sga7i.evii.n3-4-3}{3.4.3}),
\[
 \underline\Psi_{L''_\bullet}\xrightarrow{g^*}
 \underline\Psi_{L'_\bullet}\xrightarrow{f^*}
 \underline\Psi_{L_\bullet}.
\]
The transitivity property implicit in functors of the form (\SGARefLink{sga7.sga7i.evii.n3-4-3}{3.4.3})
gives an isomorphism of functors
\[
 f^*g^*\simeq(gf)^*=(0_{L_\bullet\to L''_\bullet})^*.
\]
The last member is evidently canonically isomorphic to the
cocartesian zero functor from $\underline\Psi_{L''_\bullet}$ to
$\underline\Psi_{L_\bullet}$.  The preceding isomorphism therefore
defines a canonical functor
\begin{equation}\tag{\SGARefTarget{sga7.sga7i.evii.n3-4-8-2}{3.4.8.2}}
 \underline\Psi_{L''_\bullet}\longrightarrow
 \underline\Psi_{[L_\bullet\longrightarrow L'_\bullet]}.
\end{equation}
Here $\underline\Psi_{[L_\bullet\longrightarrow L'_\bullet]}$
denotes the category of pairs $(X',t)$ in which $X'$ is an object of
$\underline\Psi_{L'_\bullet}$ and $t$ is a trivialization of its
inverse image $f^*(X')$, that is, an isomorphism
\[
 f^*(X')\simeq0_{L_\bullet,A},
\]
where $0_{L_\bullet,A}$ denotes the zero object of
$\underline\Psi_{L_\bullet}(A)$ and $A$ is the object of
$\underline C$ underlying $X'$.  We claim that \underline{the
functor} (\SGARefLink{sga7.sga7i.evii.n3-4-8-2}{3.4.8.2}) \underline{is an equivalence of categories}.
Together with the left exactness in $A$ established in \SGARefLink{sga7.sga7i.evii.n3-3-1}{3.3.1}, this
completes the exactness property in $L_\bullet$ of the expression
$\underline\Psi_{L_\bullet}(A)$.

Let us sketch the proof.  To give an object of the target of
(\SGARefLink{sga7.sga7i.evii.n3-4-8-2}{3.4.8.2}) is to give an extension
\[
 0\longrightarrow A\xrightarrow{i'}X'\xrightarrow{j'}L'_0
 \longrightarrow0
\]
of $L'_0$ by an object $A$ of $\underline C$, together with a
morphism $\alpha':L'_1\longrightarrow X'$ satisfying
$j'\alpha'=d'_0$ and $\alpha'd'_1=0$, and a morphism
$t:L_0\longrightarrow X'$ satisfying $j't=f_0$ and
$td_0=\alpha'f_1$.  These data are displayed in the following
diagram:
\begin{center}
\begin{tikzpicture}[
  x=1cm,y=1cm,>=Latex,
  every node/.style={font=\small,inner sep=1.2pt}
]
  \node (z00) at (0,0) {$0$};
  \node (l0)  at (1.3,0) {$L_0$};
  \node (lp0) at (2.8,0) {$L'_0$};
  \node (lpp0) at (4.3,0) {$L''_0$};
  \node (z05) at (5.5,0) {$0$};
  \node (z10) at (0,-1.35) {$0$};
  \node (l1)  at (1.3,-1.35) {$L_1$};
  \node (lp1) at (2.8,-1.35) {$L'_1$};
  \node (lpp1) at (4.3,-1.35) {$L''_1$};
  \node (z15) at (5.5,-1.35) {$0$};
  \node (z20) at (0,-2.70) {$0$};
  \node (l2)  at (1.3,-2.70) {$L_2$};
  \node (lp2) at (2.8,-2.70) {$L'_2$};
  \node (lpp2) at (4.3,-2.70) {$L''_2$};
  \node (z25) at (5.5,-2.70) {$0$};
  \node (xp) at (3.25,1.25) {$X'$};
  \node (ap) at (4.0,2.25) {$A$};
  \node (zap) at (4.75,3.25) {$0$};
  \node (xpp) at (5.85,1.25) {$X''$};
  \node (app) at (6.60,2.25) {$A$};
  \node (zapp) at (7.35,3.25) {$0$};
  \draw[->] (z00)--(l0);
  \draw[->] (l0)--node[above]{$f_0$}(lp0);
  \draw[->] (lp0)--node[above]{$g_0$}(lpp0);
  \draw[->] (lpp0)--(z05);
  \draw[->] (z10)--(l1);
  \draw[->] (l1)--node[above]{$f_1$}(lp1);
  \draw[->] (lp1)--node[above]{$g_1$}(lpp1);
  \draw[->] (lpp1)--(z15);
  \draw[->] (z20)--(l2);
  \draw[->] (l2)--(lp2);
  \draw[->] (lp2)--(lpp2);
  \draw[->] (lpp2)--(z25);
  \draw[->] (l1)--node[right]{$d_0$}(l0);
  \draw[->] (lp1)--node[right]{$d'_0$}(lp0);
  \draw[->] (lpp1)--node[right]{$d''_0$}(lpp0);
  \draw[->] (l2)--(l1);
  \draw[->] (lp2)--node[right]{$d'_1$}(lp1);
  \draw[->] (lpp2)--node[right]{$d''_1$}(lpp1);
  \draw[->] (zap)--(ap);
  \draw[->] (ap)--node[left]{$i'$}(xp);
  \draw[->,dashed] (zapp)--(app);
  \draw[->,dashed] (app)--node[left]{$i''$}(xpp);
  \draw[->] (l0) to[bend left=22] node[above]{$t$} (xp);
  \draw[->] (xp)--node[left]{$j'$}(lp0);
  \draw[->] (lp1)--node[pos=.82,right]{$\alpha'$}(xp);
  \draw[->,dashed] (xpp)--node[above]{$j''$}(lpp0);
  \draw[->,dashed] (lpp1) to[bend right=25]
    node[pos=.82,right]{$\alpha''$} (xpp);
  \node[anchor=west] at (7.55,-0.70) {$
    \left\{
    \begin{aligned}
      j'\alpha'\&=d'_0,\\
      \alpha'd'_1\&=0,\\
      j't\&=f_0,\\
      td_0\&=\alpha'f_1.
    \end{aligned}
    \right.$};
\end{tikzpicture}
\end{center}

The familiar left exactness, in the first argument, of extensions
of an object $L$ by an object $A$, applied to the first row of the
diagram, shows that a morphism $t$ satisfying $j't=f_0$ allows one
to interpret $X'$ as the inverse image of an extension $X''$ of
$L''_0$ by $A$.  Under this interpretation, $t$ corresponds to the
natural trivialization of the inverse image of that extension along
$g_0f_0=0$.

The morphism $\alpha'$ satisfying $j'\alpha'=d'_0$ and
$\alpha'f_1=td_0$ is interpreted by essentially the same dictionary,
now applied to the second row of the diagram, as arising from a
morphism $\alpha'':L''_1\longrightarrow X''$ satisfying
$j''\alpha''=d''_0$.  Here $\alpha'$ and $\alpha''$ are regarded as
trivializations of the inverse images of $X'$ and $X''$ along
$d'_0$ and $d''_0$, respectively.  Finally, the relation
$\alpha'd'_1=0$ is then equivalent to $\alpha''d''_1=0$.  Thus the
data $(X',\alpha',t)$, an object of
$\underline\Psi_{[L_\bullet\longrightarrow L'_\bullet]}$, are
equivalent to the data $(X'',\alpha'')$, an object of
$\underline\Psi_{L''_\bullet}$.  This proves, in substance, that
(\SGARefLink{sga7.sga7i.evii.n3-4-8-2}{3.4.8.2}) is indeed an equivalence of categories.

\SGARefTarget{sga7.sga7i.evii.n3-4-9}{3.4.9}. Continue to consider the exact sequence (\SGARefLink{sga7.sga7i.evii.n3-4-8-1}{3.4.8.1}), and the
exact sequence of $\Ext^i$ groups that it defines:
\begin{multline}\tag{\SGARefTarget{sga7.sga7i.evii.n3-4-9-1}{3.4.9.1}}
 0\longrightarrow\Ext^0(L''_\bullet,A)
 \xrightarrow{g^0}\Ext^0(L'_\bullet,A)
 \xrightarrow{f^0}\Ext^0(L_\bullet,A)
 \xrightarrow{\partial^0}\Ext^1(L''_\bullet,A)\\
 \xrightarrow{g^1}\Ext^1(L'_\bullet,A)
 \xrightarrow{f^1}\Ext^1(L_\bullet,A)\longrightarrow\cdots.
\end{multline}
The canonical isomorphisms of the forms (\SGARefLink{sga7.sga7i.evii.n3-1-8}{3.1.8}) and (\SGARefLink{sga7.sga7i.evii.n3-2-4}{3.2.4}) give a
geometric interpretation of the first six terms of this exact
sequence in terms of the categories $\underline\Psi_{L_\bullet}$,
$\underline\Psi_{L'_\bullet}$, and
$\underline\Psi_{L''_\bullet}$.  Compatibility (\SGARefLink{sga7.sga7i.evii.n3-4-6}{3.4.6}) likewise
gives a geometric interpretation of the arrows
$f^0,g^0,f^1,g^1$.  We shall now give such an interpretation of
$\partial^0$, using the result just proved that (\SGARefLink{sga7.sga7i.evii.n3-4-8-2}{3.4.8.2}) is an
equivalence of categories.

Let
\[
 u\in\Ext^0(L_\bullet,A),
 \qquad\text{i.e.}\qquad
 u:L_0\longrightarrow A,\quad ud_0=0.
\]
Consider the object $(0,\overline u)$ of
$\underline\Psi_{[L_\bullet\to L'_\bullet]}$.  Here $0$ is the
zero object of $\underline\Psi_{L'_\bullet}$, whose inverse image
along $f_0$ is therefore identified with the zero object of
$\underline\Psi_{L_\bullet}$, and $\overline u$ is the automorphism
of the latter defined by $u$ through the canonical isomorphism
(\SGARefLink{sga7.sga7i.evii.n3-1-8}{3.1.8}).  Since (\SGARefLink{sga7.sga7i.evii.n3-4-8-2}{3.4.8.2}) is an equivalence, $(0,\overline u)$ comes
from an object of $\underline\Psi_{L''_\bullet}(A)$, determined up
to unique isomorphism.  Its isomorphism class is an element of
\[
 \Psi^1_{L''_\bullet}(A)\simeq\Ext^1(L''_\bullet,A).
\]
\underline{We claim that this element is precisely}
$\partial^0(u)$.

}{E7I-VII-030}\SGApage{I}{\(\approx\)197}
We prove this compatibility by reducing it to the analogous one in
[4, Cor.~5.5], where the sign must be corrected as indicated in
3.3.6.3.  The argument also gives another proof of the equivalence
(\SGARefLink{sga7.sga7i.evii.n3-4-8-2}{3.4.8.2}), by reducing it to [4, Prop.~\SGARefLink{sga7.sga7i.evii.n2-7}{2.7}].  Consider the opposite
category $\underline C^{\mathrm o}$ and interpret the construction
of $\underline\Psi_{L_\bullet}(A)$, for fixed $A$, together with
the constructions of \SGARefLink{sga7.sga7i.evii.n3-4-8}{3.4.8} and the present section, in terms of
$\underline C^{\mathrm o}$ and the category $\underline E$ cofibred
over $\underline C^{\mathrm o}$ whose objects are extensions of the
fixed object $A$ of $\underline C^{\mathrm o}$ by a variable object
$L$ of $\underline C^{\mathrm o}$.  This is an additive cofibred
category over $\underline C^{\mathrm o}$, and is even left exact
there (in the sense of [4, \SGARefLink{sga7.sga7i.evii.n2-2}{2.2}]).

A chain complex $L_\bullet$ of $\underline C$ may be interpreted as
a cochain complex of $\underline C^{\mathrm o}$.  Under this
interpretation, the construction of $\underline\Psi_{L_\bullet}(A)$
coincides with the opposite of the category
$\underline{\widehat E}(L_\bullet)$ in the sense of [4, 1.8].  The
functor (\SGARefLink{sga7.sga7i.evii.n3-4-8-2}{3.4.8.2}) is then identified with the opposite of the functor
of the form [4, (2.1.2)], with $E$ replaced by its extension
$\underline{\widehat E}$ to the category of cochain complexes.
Consequently, the assertion that this functor is an equivalence is
the special case [4, \SGARefLink{sga7.sga7i.evii.n2-7}{2.7}, a)$\Rightarrow$d)].

\SGApage{I}{\(\approx\)198}
On the other hand, the exact sequence (\SGARefLink{sga7.sga7i.evii.n3-4-9-1}{3.4.9.1}) is identified with
the exact sequence of the groups $\Ext^i(A,L''_\bullet)$,
$\Ext^i(A,L'_\bullet)$, and $\Ext^i(A,L_\bullet)$, apart from a
factor $(-1)^i$ on $\partial^i$ (Verdier dixit).  Similarly, the
isomorphisms (\SGARefLink{sga7.sga7i.evii.n3-1-8}{3.1.8}) and (\SGARefLink{sga7.sga7i.evii.n3-2-4}{3.2.4}) are identified with the analogous
isomorphisms (3.11.5) and (3.11.4).  The asserted compatibility is
therefore the special case of [4, Cor.~5.5] corrected as above.

\underline{Remark} \SGARefTarget{sga7.sga7i.evii.n3-4-9-2}{3.4.9.2}. The editor, being very afraid of signs,
confesses that he does not have complete confidence in the preceding
proof: he fears that signs may have been omitted in passing from
$C$ to $C^{\mathrm o}$.  He recommends that a courageous reader
verify the asserted compatibility directly (or, respectively, its
opposite).  Fortunately, these sign questions will play no role in
the following expos\'es.



\subsection*{\SGARefTarget{sga7.sga7i.evii.n3-5}{3.5}. A canonical partial flat resolution of an
arbitrary $\mathbb Z$-module}

\SGAvar{We apply the preceding results when $\underline C$ is the category
of commutative groups of the topos $\underline T$, that is, the
category of $\mathbb Z$-modules of $\underline T$.  Let $P$ be an
object of $\underline C$.  We shall define, functorially in $P$, a
chain complex $L_\bullet(P)$ augmented to $P$:
\begin{equation}\tag{\SGARefTarget{sga7.sga7i.evii.n3-5-1}{3.5.1}}
\begin{tikzcd}[column sep=small,row sep=large]
 0\arrow[r]\&\overbrace{\mathbb Z[P\times P]\times
   \mathbb Z[P\times P\times P]}^{L_2}
   \arrow[r,"d_1"]
 \&\overbrace{\mathbb Z[P\times P]}^{L_1}
   \arrow[r,"d_0"]
 \&\overbrace{\mathbb Z[P]}^{L_0}
   \arrow[r]\arrow[d,"\varepsilon"']\&0\\
 \&\&\&P\&{}
\end{tikzcd}
\end{equation}
with $L_i(P)=0$ for $i\geq3$.  The components are displayed in
(\SGARefLink{sga7.sga7i.evii.n3-5-1}{3.5.1}), using the notation of \SGARefLink{sga7.sga7i.evii.n1-4}{1.4}.  To describe the differentials
$d_0,d_1$ and the augmentation $\varepsilon$, write $[p]$ for the
section of $\mathbb Z[P](S)$ defined by $p\in P(S)$, and similarly
write $[p,q]$ and $[p,q,r]$ for the corresponding sections of
$\mathbb Z[P\times P](S)$ and
$\mathbb Z[P\times P\times P](S)$.  Then set
\begin{equation}\tag{\SGARefTarget{sga7.sga7i.evii.n3-5-2}{3.5.2}}
\left\{
\begin{aligned}
 \varepsilon[p]&=p,\\
 d_0[p,q]&=[p+q]-[p]-[q],\\
 d_1[p,q]&=[p,q]-[q,p],\\
 d_1[p,q,r]&=[p+q,r]-[p,q+r]+[p,q]-[q,r].
\end{aligned}
\right.
\end{equation}
These homomorphisms of commutative groups do define a complex
$L_\bullet(P)$ augmented to $P$:
\[
 \varepsilon d_0=0,\qquad d_1d_0=0.
\]
The first relation is immediate.  The second splits into two
relations according to the two factors of $L_2(P)$; they express,
respectively, commutativity and associativity of the composition law
on $P$.

}{We apply the preceding results when $\underline C$ is the category
of commutative groups of the topos $\underline T$, that is, the
category of $\mathbb Z$-modules of $\underline T$.  Let $P$ be an
object of $\underline C$.  We shall define, functorially in $P$, a
chain complex $L_\bullet(P)$ augmented to $P$:
\begin{equation}\tag{\SGARefTarget{sga7.sga7i.evii.n3-5-1}{3.5.1}}
\begin{tikzcd}[column sep=small,row sep=large]
 0\arrow[r]\&\overbrace{\mathbb Z[P\times P]\times
   \mathbb Z[P\times P\times P]}^{L_2}
   \arrow[r,"d_1"]
 \&\overbrace{\mathbb Z[P\times P]}^{L_1}
   \arrow[r,"d_0"]
 \&\overbrace{\mathbb Z[P]}^{L_0}
   \arrow[r]\arrow[d,"\varepsilon"']\&0\\
 \&\&\&P\&{}
\end{tikzcd}
\end{equation}
with $L_i(P)=0$ for $i\geq3$.  The components are displayed in
(\SGARefLink{sga7.sga7i.evii.n3-5-1}{3.5.1}), using the notation of \SGARefLink{sga7.sga7i.evii.n1-4}{1.4}.  To describe the differentials
$d_0,d_1$ and the augmentation $\varepsilon$, write $[p]$ for the
section of $\mathbb Z[P](S)$ defined by $p\in P(S)$, and similarly
write $[p,q]$ and $[p,q,r]$ for the corresponding sections of
$\mathbb Z[P\times P](S)$ and
$\mathbb Z[P\times P\times P](S)$.  Then set
\begin{equation}\tag{\SGARefTarget{sga7.sga7i.evii.n3-5-2}{3.5.2}}
\left\{
\begin{aligned}
 \varepsilon[p]&=p,\\
 d_0[p,q]&=[p+q]-[p]-[q],\\
 d_1[p,q]&=[p,q]-[q,p],\\
 d_1[p,q,r]&=[p+q,r]-[p,q+r]+[p,q]-[q,r].
\end{aligned}
\right.
\end{equation}
These homomorphisms of commutative groups do define a complex
$L_\bullet(P)$ augmented to $P$:
\[
 \varepsilon d_0=0,\qquad d_0d_1=0.
\]
The first relation is immediate.  The second splits into two
relations according to the two factors of $L_2(P)$; they express,
respectively, commutativity and associativity of the composition law
on $P$.

}{E7I-VII-031}\SGApage{I}{\(\approx\)200}
The augmented complex depends functorially on $P$.  A morphism
$u:P\longrightarrow P'$ gives a commutative diagram
\[
\begin{tikzcd}[column sep=huge,row sep=large]
 L_\bullet(P)\arrow[r,"L_\bullet(u)"]\arrow[d]
 \&L_\bullet(P')\arrow[d]\\
 P\arrow[r,"u"]\&P'.
\end{tikzcd}
\]
Its components are flat, as is every free $\mathbb Z$-module.

\underline{Proposition} \SGARefTarget{sga7.sga7i.evii.n3-5-3}{3.5.3}. \underline{One has}
$H_1(L_\bullet(P))=0$, \underline{and the augmentation}
$\varepsilon$ \underline{induces an isomorphism}
\[
 H_0(L_\bullet(P))\xrightarrow{\sim}P.
\]

For the proof one can make the standard reduction to the case in
which $\underline T$ is the category of sets and $P$ is an ordinary
$\mathbb Z$-module, where (\SGARefLink{sga7.sga7i.evii.n3-5-1}{3.5.1}) is likely familiar.  We instead
give a direct proof that does not use this reduction, applying \SGARefLink{sga7.sga7i.evii.n3-4-7}{3.4.7}
to the augmentation
\[
 \varepsilon:L_\bullet(P)\longrightarrow P,
\]
where $P$ is regarded, as usual, as a chain complex concentrated in
degree zero.  It remains to prove that, for every object $G$ of
$\underline C$, the functor
\begin{equation}\tag{*}
 \varepsilon^*:\underline\Psi_P(G)\longrightarrow
 \underline\Psi_{L_\bullet(P)}(G)
\end{equation}
is an equivalence of categories.

The source of (*) is plainly isomorphic to the category of
extensions of the $\mathbb Z$-module $P$ by $G$.  To describe its
target, use \SGARefLink{sga7.sga7i.evii.n1-4}{1.4} to identify the categories of extensions of
$L_0,L_1,L_2$ by $G$, respectively, with the category of torsors
under $G_P$, the category of torsors under $G_{P\times P}$, and the
category of pairs consisting of a torsor under $G_{P\times P}$ and
a torsor under $G_{P\times P\times P}$.

\SGApage{I}{\(\approx\)201}
\SGAvar{Thus an object of the target of (*) is a pair $(E,\varphi)$, where
$E$ is a torsor under $G_P$ and $\varphi$ is a trivialization of the
torsor $E'$ on $P\times P$ corresponding to the inverse image along
$d_0$ of the extension of $\mathbb Z[P]$ by $G$ defined by $E$.
The trivialization is subject to the compatibility with
$d_0d_1=0$ described below.  Formula (\SGARefLink{sga7.sga7i.evii.n3-5-2}{3.5.2}) immediately gives
\[
 E'=\pi^*E\,\pr_1^*(E)^{-1}\,\pr_2^*(E)^{-1},
\]
and a trivialization of this torsor is the same as an isomorphism,
also denoted by $\varphi$,
\[
 \varphi:\pr_1^*E\,\pr_2^*E\xrightarrow{\sim}\pi^*E.
\]
Expanding the compatibility of $\varphi$ with $d_1d_0=0$ gives the
commutativity condition (\SGARefLink{sga7.sga7i.evii.n1-2-1}{1.2.1}) and the associativity condition
(\SGARefLink{sga7.sga7i.evii.n1-1-4-1}{1.1.4.1}), corresponding to the two factors of $L_2$.  The assertion
that (*) is an equivalence is therefore precisely the dictionary of
\SGARefLink{sga7.sga7i.evii.n1-1-6}{1.1.6}.  This proves \SGARefLink{sga7.sga7i.evii.n3-5-3}{3.5.3}.

}{Thus an object of the target of (*) is a pair $(E,\varphi)$, where
$E$ is a torsor under $G_P$ and $\varphi$ is a trivialization of the
torsor $E'$ on $P\times P$ corresponding to the inverse image along
$d_0$ of the extension of $\mathbb Z[P]$ by $G$ defined by $E$.
The trivialization is subject to the compatibility with
$d_0d_1=0$ described below.  Formula (\SGARefLink{sga7.sga7i.evii.n3-5-2}{3.5.2}) immediately gives
\[
 E'=\pi^*E\,\pr_1^*(E)^{-1}\,\pr_2^*(E)^{-1},
\]
and a trivialization of this torsor is the same as an isomorphism,
also denoted by $\varphi$,
\[
 \varphi:\pr_1^*E\,\pr_2^*E\xrightarrow{\sim}\pi^*E.
\]
Expanding the compatibility of $\varphi$ with $d_0d_1=0$ gives the
commutativity condition (\SGARefLink{sga7.sga7i.evii.n1-2-1}{1.2.1}) and the associativity condition
(\SGARefLink{sga7.sga7i.evii.n1-1-4-1}{1.1.4.1}), corresponding to the two factors of $L_2$.  The assertion
that (*) is an equivalence is therefore precisely the dictionary of
\SGARefLink{sga7.sga7i.evii.n1-1-6}{1.1.6}.  This proves \SGARefLink{sga7.sga7i.evii.n3-5-3}{3.5.3}.

}{E7I-VII-032}\underline{Remark} \SGARefTarget{sga7.sga7i.evii.n3-5-4}{3.5.4}. a) It would be interesting to strengthen
\SGARefLink{sga7.sga7i.evii.n3-5-3}{3.5.3} by constructing a resolution $L_\bullet(P)$ of $P$, functorial
in $P$, whose components are functorially isomorphic to sums of free
$\mathbb Z$-modules generated by cartesian powers of $P$.  In view
of isomorphisms of the form (\SGARefLink{sga7.sga7i.evii.n1-4-5}{1.4.5}), this would give a method for
partially computing $\Ext^i(P,A)$ as the abutment of a spectral
sequence whose initial term is expressed through groups
\[
 H^j(P\times P\times\cdots\times P,A).
\]
Here we obtain only very partial information about $\Ext^1$, and no
information about the higher $\Ext^i$.  There is, in any case, a
natural way to continue the partial resolution $L_\bullet(P)$ by one
degree.\footnote{L. Breen constructed such a resolution (editor's
note).}

\SGApage{I}{\(\approx\)202}
b) There is a variant of (\SGARefLink{sga7.sga7i.evii.n3-5-1}{3.5.1}) and \SGARefLink{sga7.sga7i.evii.n3-5-3}{3.5.3} in the category of
$A$-modules, where $A$ is a fixed ring of $\underline T$.  The form
of $L_\bullet(P)$ is more complicated because the description \SGARefLink{sga7.sga7i.evii.n1-1-6}{1.1.6}
must be enlarged (compare \SGARefLink{sga7.sga7i.evii.n2-10-3}{2.10.3}).  One obtains
\begin{equation}\tag{\SGARefTarget{sga7.sga7i.evii.n3-5-5}{3.5.5}}
\left\{
\begin{aligned}
 L_0={}&A[P],\\
 L_1={}&A[P\times P]+A[A\times P],\\
 L_2={}&A[P\times P\times P]+A[P\times P]
       +A[A\times A\times P]+A[A\times P\times P].
\end{aligned}
\right.
\end{equation}
The augmented-complex structure is given by (\SGARefLink{sga7.sga7i.evii.n3-5-2}{3.5.2}), supplemented by
\begin{equation}\tag{\SGARefTarget{sga7.sga7i.evii.n3-5-6}{3.5.6}}
\left\{
\begin{aligned}
 d_0[\lambda,p]={}&[\lambda p]-\lambda[p],\\
 d_1[\lambda,\mu,p]={}&[\lambda,\mu p]-[\lambda\mu,p]
                         +\lambda[\mu,p],\\
 d_1[\lambda,p,q]={}&[\lambda,p+q]-[\lambda,p]-[\lambda,q]
                       +\lambda[p,q]-[\lambda p,\lambda q].
\end{aligned}
\right.
\end{equation}
The relation $d_0d_1=0$ follows from the analogous relation on the
components appearing in (\SGARefLink{sga7.sga7i.evii.n3-5-2}{3.5.2}), which, as we saw, is equivalent to
associativity and commutativity of addition on $P$, together with
the relations
\[
 d_0d_1[\lambda,\mu,p]=0,
 \qquad d_0d_1[\lambda,p,q]=0.
\]
The latter express, respectively, associativity of the given action
of $A$ on $P$,
\[
 \lambda(\mu p)=(\lambda\mu)p,
\]
and its distributivity over addition,
\[
 \lambda(p+q)=\lambda p+\lambda q.
\]

\SGApage{I}{\(\approx\)203}
For a torsor $E$ under $G_P$, a trivialization of $d_0^*E$ amounts
to an internal composition law on $E$ (satisfying the compatibility
conditions already described for the internal law on $E$ and its
torsor structure), together with an external composition law of $A$
on $E$, compatible with the action of $A$ on $P$ and on $G$.  The
compatibility of this trivialization with the evident trivialization
of $d_1^*d_0^*E$ is equivalent to the conditions already described
above, which say that the internal law on $E$ is a commutative-group
law extending $P$ by $G$, together with associativity of the external
law and distributivity of the latter with respect to the internal
law.  These conditions say exactly that $E$, with its two
composition laws and its torsor structure, is an $A$-module extension
of the $A$-modules $P$ and $G$.

It follows, by essentially the same method as in \SGARefLink{sga7.sga7i.evii.n3-5-3}{3.5.3}, that
$L_\bullet$ is a partial resolution of the module $P$, that is,
$H_1(L_\bullet)=0$.  It would again be interesting to extend this
partial resolution to an infinite resolution of the same kind, with
components of the form
\[
 A[A\times A\times\cdots\times P\times P\times\cdots].
\]
Compare the analogous question in a) for commutative groups.



\SGApage{I}{204}
\subsection*{\SGARefTarget{sga7.sga7i.evii.n3-6}{3.6}. Application to the description of the cofibred
category of biextensions of a given pair $(P,Q)$ of commutative
groups}

\SGAvar{Let $P$ and $Q$ be two $\mathbb Z$-modules of the topos
$\underline T$.  Choose a flat resolution $L'_\bullet(P)$ of $P$
that agrees in degrees at most $2$ with the augmented complex
$L_\bullet(P)$ of \SGARefLink{sga7.sga7i.evii.n3-5}{3.5}.  This is possible by \SGARefLink{sga7.sga7i.evii.n3-5-3}{3.5.3} and because every
$\mathbb Z$-module on $\underline T$ is a quotient of a free, hence
flat, $\mathbb Z$-module.  Similarly choose a flat resolution
$L'_\bullet(Q)$ of $Q$ agreeing with $L_\bullet(Q)$ in degrees at
most $2$.  We have canonical homomorphisms
\begin{equation}\tag{\SGARefTarget{sga7.sga7i.evii.n3-6-1}{3.6.1}}
 L_\bullet(P)\longrightarrow L'_\bullet(P),
 \qquad
 L_\bullet(Q)\longrightarrow L'_\bullet(Q),
\end{equation}
and hence a canonical homomorphism
\begin{equation}\tag{\SGARefTarget{sga7.sga7i.evii.n3-6-2}{3.6.2}}
\begin{split}
 L_\bullet(P,Q)
 &\overset{\mathrm{def}}{=}L_\bullet(P)\otimes L_\bullet(Q)
 \longrightarrow L'_\bullet(P,Q)\\
 &\overset{\mathrm{def}}{=}L'_\bullet(P)\otimes L'_\bullet(Q)
 \simeq P\mathbin{\overset{\mathbb L}{\otimes}}Q.
\end{split}
\end{equation}
It is an isomorphism on components in degrees at most $2$.
Consequently, by definition,
\begin{equation}\tag{\SGARefTarget{sga7.sga7i.evii.n3-6-3}{3.6.3}}
 \underline\Psi_{L_\bullet(P,Q)}
 =
 \underline\Psi_{L_\bullet(P',Q')}.
\end{equation}

}{Let $P$ and $Q$ be two $\mathbb Z$-modules of the topos
$\underline T$.  Choose a flat resolution $L'_\bullet(P)$ of $P$
that agrees in degrees at most $2$ with the augmented complex
$L_\bullet(P)$ of \SGARefLink{sga7.sga7i.evii.n3-5}{3.5}.  This is possible by \SGARefLink{sga7.sga7i.evii.n3-5-3}{3.5.3} and because every
$\mathbb Z$-module on $\underline T$ is a quotient of a free, hence
flat, $\mathbb Z$-module.  Similarly choose a flat resolution
$L'_\bullet(Q)$ of $Q$ agreeing with $L_\bullet(Q)$ in degrees at
most $2$.  We have canonical homomorphisms
\begin{equation}\tag{\SGARefTarget{sga7.sga7i.evii.n3-6-1}{3.6.1}}
 L_\bullet(P)\longrightarrow L'_\bullet(P),
 \qquad
 L_\bullet(Q)\longrightarrow L'_\bullet(Q),
\end{equation}
and hence a canonical homomorphism
\begin{equation}\tag{\SGARefTarget{sga7.sga7i.evii.n3-6-2}{3.6.2}}
\begin{split}
 L_\bullet(P,Q)
 &\overset{\mathrm{def}}{=}L_\bullet(P)\otimes L_\bullet(Q)
 \longrightarrow L'_\bullet(P,Q)\\
 &\overset{\mathrm{def}}{=}L'_\bullet(P)\otimes L'_\bullet(Q)
 \simeq P\mathbin{\overset{\mathbb L}{\otimes}}Q.
\end{split}
\end{equation}
It is an isomorphism on components in degrees at most $2$.
Consequently, by definition,
\begin{equation}\tag{\SGARefTarget{sga7.sga7i.evii.n3-6-3}{3.6.3}}
 \underline\Psi_{L_\bullet(P,Q)}
 =
 \underline\Psi_{L'_\bullet(P,Q)}.
\end{equation}

}{E7I-VII-033}\underline{Theorem} \SGARefTarget{sga7.sga7i.evii.n3-6-4}{3.6.4}. \underline{The additive cofibred category
over the category}
\[
 \underline C=\operatorname{GRAB}(\underline T)
\]
\underline{of commutative groups of $\underline T$, consisting of
biextensions of $(P,Q)$ by a variable abelian group $G$} (\SGARefLink{sga7.sga7i.evii.n2-4}{2.4}),
\underline{is canonically equivalent to the additive cofibred
category}
\[
 \underline\Psi_{L_\bullet(P,Q)}
 =
 \underline\Psi_{L'_\bullet(P,Q)}
\]
\underline{defined from}
$L_\bullet(P,Q)=L_\bullet(P)\otimes L_\bullet(Q)$
\underline{as in \SGARefLink{sga7.sga7i.evii.n3-1}{3.1}}.

This and \SGARefLink{sga7.sga7i.evii.n3-2-5}{3.2.5} give the principal result of the present section,
announced above:

\SGApage{I}{205}
\underline{Corollary} \SGARefTarget{sga7.sga7i.evii.n3-6-5}{3.6.5}. \underline{For fixed $P$ and $Q$ there
is a canonical isomorphism, functorial in $G$,}
\begin{equation}\tag{\SGARefTarget{sga7.sga7i.evii.n3-6-5-1}{3.6.5.1}}
 \Biext^1(P,Q;G)
 \xrightarrow{\sim}
 \Ext^1(P\mathbin{\overset{\mathbb L}{\otimes}}Q,G).
\end{equation}

To prove \SGARefLink{sga7.sga7i.evii.n3-6-5}{3.6.5}, first make explicit the components of degrees
$0,1,2$ of $L_\bullet(P,Q)$:
\begin{equation}\tag{\SGARefTarget{sga7.sga7i.evii.n3-6-6}{3.6.6}}
\left\{
\begin{aligned}
 L_0(P,Q)
  ={}&L_0(P)\otimes L_0(Q)
     =\mathbb Z[P\times Q],\\
 L_1(P,Q)
  ={}&L_0(P)\otimes L_1(Q)+L_1(P)\otimes L_0(Q)\\
  ={}&\mathbb Z[P\times Q\times Q]
     +\mathbb Z[P\times P\times Q],\\
 L_2(P,Q)
  ={}&L_0(P)\otimes L_2(Q)+L_1(P)\otimes L_1(Q)
     +L_2(P)\otimes L_0(Q)\\
  ={}&\mathbb Z[P\times Q\times Q]
     +\mathbb Z[P\times Q\times Q\times Q]\\
   &+\mathbb Z[P\times P\times Q\times Q]\\
   &+\mathbb Z[P\times P\times Q]
     +\mathbb Z[P\times P\times P\times Q].
\end{aligned}
\right.
\end{equation}

\SGAvar{To describe the cofibred category
$\underline\Psi_{L_\bullet(P,Q)}$ through its fibres
$\underline\Psi_{L_\bullet(P,Q)}(G)$, use \SGARefLink{sga7.sga7i.evii.n1-4}{1.4} to express extensions
of $L_i(P,Q)$, for $i=0,1,2$, by $G$ in terms of torsors.  An
extension of $L_0(P\times Q)$ by $G$ is a torsor under $G_{P\times Q}$.
An extension of $L_1(P,Q)$ by $G$ is a pair consisting of a torsor
under $G_{P\times Q\times Q}$ and one under
$G_{P\times P\times Q}$.  An extension of $L_2(P,Q)$ by $G$ is
similarly a system of five torsors, under the groups obtained from
$G$ by base change to
\[
\begin{gathered}
 P\times Q\times Q,\quad
 P\times Q\times Q\times Q,\quad
 P\times P\times Q\times Q,\\
 P\times P\times Q,\quad
 P\times P\times P\times Q.
\end{gathered}
\]

Under this dictionary, a pair $(X,\alpha)$ consisting of an
extension of $L_0(P,Q)$ by $G$ and a trivialization of
$d_0^*X$ on $L_1(P,Q)$ becomes a torsor $E$ on $P\times Q$ equipped
with two trivializations $(\varphi,\psi)$ of the torsors
$(E_1,E_2)$ on $P\times Q\times Q$ and
$P\times P\times Q$, respectively, obtained from $E$ by inverse
image along
\[
 d_0:L_1(P,Q)\longrightarrow L_0(P,Q).
\]
Expanding $d_0$, as in the proof of \SGARefLink{sga7.sga7i.evii.n3-5-3}{3.5.3}, identifies
$\varphi$ and $\psi$ with isomorphisms
\[
 \varphi:q_1^*E\,q_2^*E\xrightarrow{\sim}\pi_Q^*E,
 \qquad
 \psi:p_1^*E\,p_2^*E\xrightarrow{\sim}\pi_P^*E.
\]
Here $\pi_Q,q_1q_2:P\times Q\times Q\longrightarrow  Q$
are induced by the composition law and the first and second
projections $Q\times Q\longrightarrow Q$; the maps
\[
 \pi_P,p_1,p_2:P\times P\times Q\longrightarrow P\times Q
\]
are defined symmetrically.  Thus $\varphi$ and $\psi$ are data of
the forms (\SGARefLink{sga7.sga7i.evii.n2-0-4}{2.0.4}) and (\SGARefLink{sga7.sga7i.evii.n2-0-7}{2.0.7}) appearing in \SGARefLink{sga7.sga7i.evii.n2-1}{2.1}.

Finally, the condition that $(X,\alpha)$ be an object of
$\underline\Psi_{L_\bullet(P,Q)}(G)$, namely compatibility of the
partial trivialization $\alpha$ with $d_0d_1=0$, becomes five
compatibility relations for the structure $(E,\varphi,\psi)$ on
torsors over, respectively,
\[
\begin{gathered}
 P\times Q\times Q,\quad
 P\times Q\times Q\times Q,\quad
 P\times P\times Q\times Q,\\
 P\times P\times Q,\quad
 P\times P\times P\times Q.
\end{gathered}
\]
They correspond to the five terms in the expression (\SGARefLink{sga7.sga7i.evii.n3-6-6}{3.6.6}) for
$L_2(P,Q)$.  As in the proof of \SGARefLink{sga7.sga7i.evii.n3-5-3}{3.5.3}, the first two are precisely
the commutativity relation (\SGARefLink{sga7.sga7i.evii.n2-0-6}{2.0.6}) and the associativity relation
(\SGARefLink{sga7.sga7i.evii.n2-0-5}{2.0.5}); together they say that $\varphi$ gives $E$ the structure
of a commutative extension of $Q_P$ by $G_P$.  Symmetrically, the
last two say that $\psi$ gives $E$ the structure of a commutative
extension of $P_Q$ by $G_P$.  The middle relation, on torsors over
$P\times P\times Q\times Q$, is exactly the commutativity of
diagram (\SGARefLink{sga7.sga7i.evii.n2-1-1}{2.1.1}).  It says that $\varphi$ and $\psi$ are compatible,
hence, in view of the other conditions, define on $E$ a
biextension structure of $(P,Q)$ by $G$.

}{To describe the cofibred category
$\underline\Psi_{L_\bullet(P,Q)}$ through its fibres
$\underline\Psi_{L_\bullet(P,Q)}(G)$, use \SGARefLink{sga7.sga7i.evii.n1-4}{1.4} to express extensions
of $L_i(P,Q)$, for $i=0,1,2$, by $G$ in terms of torsors.  An
extension of $L_0(P,Q)$ by $G$ is a torsor under $G_{P\times Q}$.
An extension of $L_1(P,Q)$ by $G$ is a pair consisting of a torsor
under $G_{P\times Q\times Q}$ and one under
$G_{P\times P\times Q}$.  An extension of $L_2(P,Q)$ by $G$ is
similarly a system of five torsors, under the groups obtained from
$G$ by base change to
\[
\begin{gathered}
 P\times Q\times Q,\quad
 P\times Q\times Q\times Q,\quad
 P\times P\times Q\times Q,\\
 P\times P\times Q,\quad
 P\times P\times P\times Q.
\end{gathered}
\]

Under this dictionary, a pair $(X,\alpha)$ consisting of an
extension of $L_0(P,Q)$ by $G$ and a trivialization of
$d_0^*X$ on $L_1(P,Q)$ becomes a torsor $E$ on $P\times Q$ equipped
with two trivializations $(\varphi,\psi)$ of the torsors
$(E_1,E_2)$ on $P\times Q\times Q$ and
$P\times P\times Q$, respectively, obtained from $E$ by inverse
image along
\[
 d_0:L_1(P,Q)\longrightarrow L_0(P,Q).
\]
Expanding $d_0$, as in the proof of \SGARefLink{sga7.sga7i.evii.n3-5-3}{3.5.3}, identifies
$\varphi$ and $\psi$ with isomorphisms
\[
 \varphi:q_1^*E\,q_2^*E\xrightarrow{\sim}\pi_Q^*E,
 \qquad
 \psi:p_1^*E\,p_2^*E\xrightarrow{\sim}\pi_P^*E.
\]
Here $\pi_Q,q_1,q_2:P\times Q\times Q\longrightarrow P\times Q$
are induced by the composition law and the first and second
projections $Q\times Q\longrightarrow Q$; the maps
\[
 \pi_P,p_1,p_2:P\times P\times Q\longrightarrow P\times Q
\]
are defined symmetrically.  Thus $\varphi$ and $\psi$ are data of
the forms (\SGARefLink{sga7.sga7i.evii.n2-0-4}{2.0.4}) and (\SGARefLink{sga7.sga7i.evii.n2-0-7}{2.0.7}) appearing in \SGARefLink{sga7.sga7i.evii.n2-1}{2.1}.

Finally, the condition that $(X,\alpha)$ be an object of
$\underline\Psi_{L_\bullet(P,Q)}(G)$, namely compatibility of the
partial trivialization $\alpha$ with $d_0d_1=0$, becomes five
compatibility relations for the structure $(E,\varphi,\psi)$ on
torsors over, respectively,
\[
\begin{gathered}
 P\times Q\times Q,\quad
 P\times Q\times Q\times Q,\quad
 P\times P\times Q\times Q,\\
 P\times P\times Q,\quad
 P\times P\times P\times Q.
\end{gathered}
\]
They correspond to the five terms in the expression (\SGARefLink{sga7.sga7i.evii.n3-6-6}{3.6.6}) for
$L_2(P,Q)$.  As in the proof of \SGARefLink{sga7.sga7i.evii.n3-5-3}{3.5.3}, the first two are precisely
the commutativity relation (\SGARefLink{sga7.sga7i.evii.n2-0-6}{2.0.6}) and the associativity relation
(\SGARefLink{sga7.sga7i.evii.n2-0-5}{2.0.5}); together they say that $\varphi$ gives $E$ the structure
of a commutative extension of $Q_P$ by $G_P$.  Symmetrically, the
last two say that $\psi$ gives $E$ the structure of a commutative
extension of $P_Q$ by $G_Q$.  The middle relation, on torsors over
$P\times P\times Q\times Q$, is exactly the commutativity of
diagram (\SGARefLink{sga7.sga7i.evii.n2-1-1}{2.1.1}).  It says that $\varphi$ and $\psi$ are compatible,
hence, in view of the other conditions, define on $E$ a
biextension structure of $(P,Q)$ by $G$.

}{E7I-VII-034}\SGApage{I}{\(\approx\)206}
The verification of these interpretations is immediate and is left
to the reader.  The preceding discussion indeed defines an
equivalence of categories cofibred over
$\underline C=\operatorname{GRAB}(\underline T)$, as asserted in
\SGARefLink{sga7.sga7i.evii.n3-6-4}{3.6.4}.

\SGApage{I}{207}
\underline{Remark} \SGARefTarget{sga7.sga7i.evii.n3-6-7}{3.6.7}. Each of the three generalizations of
biextensions mentioned in \SGARefLink{sga7.sga7i.evii.n2-10}{2.10} has a corresponding variant of \SGARefLink{sga7.sga7i.evii.n3-6-4}{3.6.4}
and hence of \SGARefLink{sga7.sga7i.evii.n3-6-5}{3.6.5}.  The generalization to any number of factors
(\SGARefLink{sga7.sga7i.evii.n2-10-2}{2.10.2}) is immediate, as is the generalization to biextensions of
$A$-modules (\SGARefLink{sga7.sga7i.evii.n2-10-3}{2.10.3}), once the variant of $L_\bullet(P)$ described
in \SGARefLink{sga7.sga7i.evii.n3-5-4}{3.5.4} b) has been constructed.  From the viewpoint of
homological algebra, the case in which $P$ and $Q$ are not assumed
commutative (\SGARefLink{sga7.sga7i.evii.n2-10-1}{2.10.1}) is more interesting.

Associate with a not necessarily commutative group $P$ of
$\underline T$ a complex $K_\bullet(P)$ truncated in degree $2$.
It is described as $L_\bullet(P)$ was in \SGARefLink{sga7.sga7i.evii.n3-5}{3.5}, except that
\[
 K_2(P)=\mathbb Z[P\times P\times P],
\]
the first factor occurring in $L_2(P)$ being omitted; $d_0$ and
$d_1$ are defined as in (\SGARefLink{sga7.sga7i.evii.n3-5-2}{3.5.2}).  Then
\[
 H_0(K_\bullet(P))=P/[P,P]
 \simeq\underline H_1(P,\mathbb Z),
\]
and similarly
\[
 H_1(K_\bullet(P))\simeq\underline H_2(P,\mathbb Z).
\]
Indeed, $K_\bullet(P)$ is precisely the truncation in degree $2$ of
the chain complex $C_\bullet(P,\mathbb Z)$ of $P$ with coefficients
in $\mathbb Z$.

The argument proving \SGARefLink{sga7.sga7i.evii.n3-5-3}{3.5.3} also shows that the cofibred category
$\underline\Psi_{K_\bullet(P)}$ over
$\underline C=\operatorname{GRAB}(\underline T)$ is canonically
equivalent to the category of central extensions of $P$ by variable
commutative groups $G$.  Likewise, for two groups $P,Q$, the
cofibred category of biextensions of $(P,Q)$ in the sense of \SGARefLink{sga7.sga7i.evii.n2-10-1}{2.10.1}
is equivalent to $\underline\Psi_{K_\bullet(P,Q)}$, where
\[
 K_\bullet(P,Q)
 =
 K_\bullet(P)\mathbin{\overset{\mathbb L}{\otimes}}K_\bullet(Q).
\]

\SGApage{I}{\(\approx\)208}
These considerations suggest that, when $P$ is commutative, the
canonical resolution sought in \SGARefLink{sga7.sga7i.evii.n3-5-4}{3.5.4} a) might be obtained by a
suitable modification of $C_\bullet(P,\mathbb Z)$.



\subsection*{\SGARefTarget{sga7.sga7i.evii.n3-7}{3.7}. Various compatibilities}

In this section we examine the behaviour, for variable $P$ and $Q$,
of the equivalence of cofibred categories in \SGARefLink{sga7.sga7i.evii.n3-6-4}{3.6.4}.  The variance of
the cofibred category of biextensions of $(P,Q)$ by a variable $G$
with respect to variable $P,Q$ is defined in principle in \SGARefLink{sga7.sga7i.evii.n2-4}{2.4}; it may
also be expressed by saying that one has a cofibred functor (\SGARefLink{sga7.sga7i.evii.n2-6-1}{2.6.1}).
Notice that the variance in $P,Q$ is \underline{contravariant}.  The
variance of the cofibred category
$\underline\Psi_{L_\bullet(P,Q)}$ with respect to variable $P,Q$ is
also clear, since
$L_\bullet(P,Q)=L_\bullet(P)\otimes L_\bullet(Q)$ depends
functorially on the pair $(P,Q)$, while $\underline\Psi_{L_\bullet}$
depends contravariantly on $L_\bullet$ (\SGARefLink{sga7.sga7i.evii.n3-4}{3.4}).  With this understood,
a first compatibility is expressed, for every triple of morphisms
\begin{equation}\tag{\SGARefTarget{sga7.sga7i.evii.n3-7-1}{3.7.1}}
 P'\longrightarrow P,\qquad
 Q'\longrightarrow Q,\qquad
 G\longrightarrow G',
\end{equation}
by the fact that the diagram of functors
\begin{equation}\tag{\SGARefTarget{sga7.sga7i.evii.n3-7-2}{3.7.2}}
\begin{tikzcd}[column sep=huge,row sep=large]
 \underline\Psi_{L_\bullet(P,Q)}(G)
   \arrow[r]\arrow[d]
 \&\underline\Psi_{L_\bullet(P',Q')}(G')
   \arrow[d]\\
 \operatorname{BIEXT}(P,Q;G)
   \arrow[r]
 \&\operatorname{BIEXT}(P',Q';G')
\end{tikzcd}
\end{equation}
is commutative up to a canonical isomorphism.  We shall admit this
compatibility: making the commutativity isomorphism explicit could
only be as obvious as it is tedious (as, indeed, could the complete
explication it presupposes of the equivalence of cofibred categories
in \SGARefLink{sga7.sga7i.evii.n3-6-4}{3.6.4}).

\SGApage{I}{209}
\SGARefTarget{sga7.sga7i.evii.n3-7-3}{3.7.3}. It follows in particular from the preceding compatibility
that the isomorphism (\SGARefLink{sga7.sga7i.evii.n3-6-5-1}{3.6.5.1}) is functorial not only in $G$, but also
in $P$ and $Q$.  One may state, similarly and more trivially, the
analogous functoriality of the isomorphism defined through (\SGARefLink{sga7.sga7i.evii.n3-1-8}{3.1.8}),
just as (\SGARefLink{sga7.sga7i.evii.n3-6-5-1}{3.6.5.1}) was defined through (\SGARefLink{sga7.sga7i.evii.n3-2-4}{3.2.4}):
\begin{equation}\tag{\SGARefTarget{sga7.sga7i.evii.n3-7-4}{3.7.4}}
 \Biext^0(P,Q;G)
 \xrightarrow{\sim}
 \Ext^0(P\mathbin{\overset{\mathbb L}{\otimes}}Q,G)
 \simeq\Hom(P\otimes Q,G).
\end{equation}
But such an isomorphism was already obtained in (\SGARefLink{sga7.sga7i.evii.n2-5-4-1}{2.5.4.1}), and was
observed in \SGARefLink{sga7.sga7i.evii.n2-6}{2.6} to be functorial in its three arguments.  It is
immediate to verify that the two isomorphisms (\SGARefLink{sga7.sga7i.evii.n3-7-4}{3.7.4}) and (\SGARefLink{sga7.sga7i.evii.n2-5-4-1}{2.5.4.1})
are equal.

\SGARefTarget{sga7.sga7i.evii.n3-7-5}{3.7.5}. The value of the isomorphism in \SGARefLink{sga7.sga7i.evii.n3-6-5}{3.6.5} is that it expresses the
remarkable group of geometric origin $\Biext^1(P,Q;G)$ in terms of
the very simple homological-algebra invariant
$\Ext^1(P\mathbin{\overset{\mathbb L}{\otimes}}Q,G)$.  The latter
belongs to the sequence of invariants
$\Ext^i(P\mathbin{\overset{\mathbb L}{\otimes}}Q,G)$, to which the
usual formalism of homological algebra applies.  In particular, the
sequence of these $\Ext^i$ is a \underline{cohomological functor} in
the sense of [3] (or a \underline{connected sequence of functors} in
the terminology of [5]) in each of the three arguments $P,Q,G$.
Thus an exact sequence
\begin{equation}\tag{\SGARefTarget{sga7.sga7i.evii.n3-7-5-1}{3.7.5.1}}
 0\longrightarrow G'\longrightarrow G\longrightarrow G''
 \longrightarrow0
\end{equation}
(respectively,
\begin{equation}\tag{\SGARefTarget{sga7.sga7i.evii.n3-7-5-2}{3.7.5.2}}
 0\longrightarrow P'\longrightarrow P\longrightarrow P''
 \longrightarrow0,\qquad\text{or}\qquad
 0\longrightarrow Q'\longrightarrow Q\longrightarrow Q''
 \longrightarrow0)
\end{equation}
gives rise to an infinite exact sequence of the $\Ext^i$.  Its first
six terms, which involve only $\Ext^0$ and $\Ext^1$, may be
interpreted in terms of biextensions by \SGARefLink{sga7.sga7i.evii.n3-6-5}{3.6.5} and \SGARefLink{sga7.sga7i.evii.n3-7-4}{3.7.4}.  We have
already indicated in \SGARefLink{sga7.sga7i.evii.n3-3-6}{3.3.6} a geometric interpretation of the exact
sequence deduced from (\SGARefLink{sga7.sga7i.evii.n3-7-5-1}{3.7.5.1}), and especially of the boundary map
\begin{equation}\tag{\SGARefTarget{sga7.sga7i.evii.n3-7-5-3}{3.7.5.3}}
 \partial:\Biext^0(P,Q;G'')\longrightarrow
 \Biext^1(P,Q;G'),
\end{equation}
within the framework of a \underline{left exact cofibred category}.
One may seek an analogous interpretation of the six-term exact
sequence associated with either short exact sequence in (\SGARefLink{sga7.sga7i.evii.n3-7-5-2}{3.7.5.2}),
say the first; equivalently, in view of \SGARefLink{sga7.sga7i.evii.n3-7-3}{3.7.3}, one may seek a
geometric interpretation of the boundary homomorphism
\begin{equation}\tag{\SGARefTarget{sga7.sga7i.evii.n3-7-5-4}{3.7.5.4}}
 \partial:\Biext^0(P',Q;G)\longrightarrow
 \Biext^1(P'',Q;G).
\end{equation}
Indeed, we have the following result.

\SGApage{I}{210}
\underline{Proposition} \SGARefTarget{sga7.sga7i.evii.n3-7-6}{3.7.6}. \underline{The cofibred category over}
\[
 \operatorname{GRAB}(\underline T)^{\mathrm{op}}
 \times\operatorname{GRAB}(\underline T)^{\mathrm{op}}
 \times\operatorname{GRAB}(\underline T)
\]
\underline{of biextensions of $(P,Q)$ by $G$, with $P,Q,G$
variable} (cf. (\SGARefLink{sga7.sga7i.evii.n2-6-1}{2.6.1})), \underline{is left exact not only in the
argument $G$, but also in each of the arguments $P$ and $Q$.
Moreover, the boundary homomorphism (\SGARefLink{sga7.sga7i.evii.n3-7-5-4}{3.7.5.4}) deduced from the exact
sequence of the $\Ext^i$ associated with (\SGARefLink{sga7.sga7i.evii.n3-7-5-2}{3.7.5.2}), by means of the
isomorphisms \SGARefLink{sga7.sga7i.evii.n3-6-5}{3.6.5} and \SGARefLink{sga7.sga7i.evii.n3-7-4}{3.7.4}, is the boundary homomorphism in the
theory of left exact cofibred categories over}
$\underline C=\operatorname{GRAB}(\underline T)^{\mathrm{op}}$
\underline{recalled in \SGARefLink{sga7.sga7i.evii.n3-3-6}{3.3.6}, applied to the cofibred category of
biextensions of $(P,Q)$ by $G$, where $Q$ and $G$ are fixed and $P$
varies} (subject to a possible sign error in \SGARefLink{sga7.sga7i.evii.n3-4-9}{3.4.9}; cf. \SGARefLink{sga7.sga7i.evii.n3-4-9-2}{3.4.9.2}).

\SGApage{I}{\(\approx\)211}
We sketch the proof of \SGARefLink{sga7.sga7i.evii.n3-7-6}{3.7.6}.  By \SGARefLink{sga7.sga7i.evii.n3-6-4}{3.6.4} and compatibility (\SGARefLink{sga7.sga7i.evii.n3-7-2}{3.7.2}),
we interpret biextensions of $(P,Q)$ by $G$ as objects of
$\underline\Psi_{L_\bullet(P,Q)}(G)$.  Consider the homomorphism
\[
 L_\bullet(P,Q)
 =L_\bullet(P)\otimes L_\bullet(Q)
 \longrightarrow P\otimes L_\bullet(Q)
\]
defined by the augmentation $L_\bullet(P)\longrightarrow P$.
It follows from \SGARefLink{sga7.sga7i.evii.n3-5-3}{3.5.3} and the flatness of $L_\bullet(Q)$ that this
homomorphism induces isomorphisms on $H_0$ and $H_1$, and hence by
\SGARefLink{sga7.sga7i.evii.n3-4-7}{3.4.7} an equivalence
\[
 \underline\Psi_{P\otimes L_\bullet(Q)}
 \longrightarrow
 \underline\Psi_{L_\bullet(P,Q)}.
\]
We may therefore interpret biextensions of $(P,Q)$ by $G$ as objects
of $\underline\Psi_{P\otimes L_\bullet(Q)}(G)$.  The exact sequence
(\SGARefLink{sga7.sga7i.evii.n3-7-5-2}{3.7.5.2}) now gives an exact sequence of chain complexes
\begin{equation}\tag{\SGARefTarget{sga7.sga7i.evii.n3-7-6-1}{3.7.6.1}}
 0\longrightarrow P'\otimes L_\bullet(Q)
 \longrightarrow P\otimes L_\bullet(Q)
 \longrightarrow P''\otimes L_\bullet(Q)
 \longrightarrow0.
\end{equation}
The asserted left exactness follows from \SGARefLink{sga7.sga7i.evii.n3-4-8}{3.4.8} applied to this exact
sequence of complexes.  (A direct verification would in any event
be immediate.)  Moreover, the homomorphism (\SGARefLink{sga7.sga7i.evii.n3-7-5-4}{3.7.5.4}) is precisely the
boundary homomorphism $\Ext^0\longrightarrow\Ext^1$ deduced from this
exact sequence by taking the $\Ext^i$ with values in $G$.  The last
assertion of \SGARefLink{sga7.sga7i.evii.n3-7-6}{3.7.6} is consequently a special case of \SGARefLink{sga7.sga7i.evii.n3-4-9}{3.4.9}.



\SGApage{I}{212}
\subsection*{\SGARefTarget{sga7.sga7i.evii.n3-8}{3.8}. Biextensions of $(P,Q)$ trivialized over extensions of $P$ and $Q$}

\SGAvar{First consider two exact sequences of abelian groups
\begin{equation}\tag{\SGARefTarget{sga7.sga7i.evii.n3-8-1}{3.8.1}}
 0\longrightarrow L_1\xrightarrow{\ i\ }L_0\xrightarrow{\ u\ }P
 \longrightarrow0,\qquad
 0\longrightarrow M_1\xrightarrow{\ j\ }M_0\xrightarrow{\ v\ }Q
 \longrightarrow0.
\end{equation}
We shall describe the category of biextensions $E$ of $(P,Q)$ by an
abelian group $G$ for which the inverse image $(u,v)^*(E)$ is trivial;
that is, those for which there exists a homomorphism $s$ making the
following triangle commutative:
\begin{equation}\tag{\SGARefTarget{sga7.sga7i.evii.n3-8-2}{3.8.2}}
\begin{tikzcd}[column sep=huge,row sep=huge]
 E\arrow[d]\arrow[dr,dash,"s"]\&\\
 P\times Q\&L_0\times M_0\arrow[l,"u\times v"']
\end{tikzcd}
\end{equation}
and satisfying the additivity condition required in each of its two
arguments.  The choice of such an $s$ therefore identifies
$(u,v)^*(E)$ with the trivial biextension
$E_0=L_0\times M_0\times G$ of $(L_0,M_0)$ by $G$.  The left
exactness in $P,Q$ of the cofibred category of biextensions of
$(P,Q)$ by fixed $G$ (\SGARefLink{sga7.sga7i.evii.n3-7-6}{3.7.6}) then shows that to give the biextension
$E$ of $(P,Q)$, together with the trivialization $s$ of its inverse
image, is essentially equivalent to giving two partial
trivializations of the trivial biextension $E_0$ of $(L_0,M_0)$ by
$G$, over $(L_1,M_0)$ and $(L_0,M_1)$ respectively, which agree over
$(L_1,M_1)$.  Since the biextensions induced by $E_0$ over
$(L_1,M_0)$, $(L_0,M_1)$, and $(L_1,M_1)$ identify with the trivial
biextensions, a partial trivialization over $(L_1,M_0)$ (respectively
over $(L_0,M_1)$) is equivalent to giving a homomorphism (or pairing)
\begin{equation}\tag{\SGARefTarget{sga7.sga7i.evii.n3-8-3}{3.8.3}}
 f:L_1\otimes M_0\longrightarrow G
 \qquad
 \bigl(\text{respectively }g:L_0\otimes M_1\longrightarrow G\bigr).
\end{equation}
The compatibility in question simply means that the two
homomorphisms agree over $L_1\otimes M_1$, namely
\begin{equation}\tag{\SGARefTarget{sga7.sga7i.evii.n3-8-4}{3.8.4}}
 f(x_1\otimes y_1)=g(x_1\otimes y_1),
\end{equation}
for sections $x_1,y_1$ of $L_1,M_1$ over an object of
$\underline T$.  In terms of the homomorphism $s$ of (\SGARefLink{sga7.sga7i.evii.n3-8-2}{3.8.2}), the
forms $f$ and $g$ are characterized by
\begin{equation}\tag{\SGARefTarget{sga7.sga7i.evii.n3-8-5}{3.8.5}}
 s(x_1,y_0)=f(x_1\otimes y_0)e_{E/Q}(y_0),\qquad
 s(x_0,y_1)=g(x_0\otimes y_1)e_{E/P}(x_0),
\end{equation}
which again makes clear the necessity of condition (\SGARefLink{sga7.sga7i.evii.n3-8-4}{3.8.4}).

}{First consider two exact sequences of abelian groups
\begin{equation}\tag{\SGARefTarget{sga7.sga7i.evii.n3-8-1}{3.8.1}}
 0\longrightarrow L_1\xrightarrow{\ i\ }L_0\xrightarrow{\ u\ }P
 \longrightarrow0,\qquad
 0\longrightarrow M_1\xrightarrow{\ j\ }M_0\xrightarrow{\ v\ }Q
 \longrightarrow0.
\end{equation}
We shall describe the category of biextensions $E$ of $(P,Q)$ by an
abelian group $G$ for which the inverse image $(u,v)^*(E)$ is trivial;
that is, those for which there exists a homomorphism $s$ making the
following triangle commutative:
\begin{equation}\tag{\SGARefTarget{sga7.sga7i.evii.n3-8-2}{3.8.2}}
\begin{tikzcd}[column sep=huge,row sep=huge]
 E\arrow[d]\&\\
 P\times Q\&L_0\times M_0\arrow[l,"u\times v"']\arrow[ul,"s"]
\end{tikzcd}
\end{equation}
and satisfying the additivity condition required in each of its two
arguments.  The choice of such an $s$ therefore identifies
$(u,v)^*(E)$ with the trivial biextension
$E_0=L_0\times M_0\times G$ of $(L_0,M_0)$ by $G$.  The left
exactness in $P,Q$ of the cofibred category of biextensions of
$(P,Q)$ by fixed $G$ (\SGARefLink{sga7.sga7i.evii.n3-7-6}{3.7.6}) then shows that to give the biextension
$E$ of $(P,Q)$, together with the trivialization $s$ of its inverse
image, is essentially equivalent to giving two partial
trivializations of the trivial biextension $E_0$ of $(L_0,M_0)$ by
$G$, over $(L_1,M_0)$ and $(L_0,M_1)$ respectively, which agree over
$(L_1,M_1)$.  Since the biextensions induced by $E_0$ over
$(L_1,M_0)$, $(L_0,M_1)$, and $(L_1,M_1)$ identify with the trivial
biextensions, a partial trivialization over $(L_1,M_0)$ (respectively
over $(L_0,M_1)$) is equivalent to giving a homomorphism (or pairing)
\begin{equation}\tag{\SGARefTarget{sga7.sga7i.evii.n3-8-3}{3.8.3}}
 f:L_1\otimes M_0\longrightarrow G
 \qquad
 \bigl(\text{respectively }g:L_0\otimes M_1\longrightarrow G\bigr).
\end{equation}
The compatibility in question simply means that the two
homomorphisms agree over $L_1\otimes M_1$, namely
\begin{equation}\tag{\SGARefTarget{sga7.sga7i.evii.n3-8-4}{3.8.4}}
 f(x_1\otimes y_1)=g(x_1\otimes y_1),
\end{equation}
for sections $x_1,y_1$ of $L_1,M_1$ over an object of
$\underline T$.  In terms of the homomorphism $s$ of (\SGARefLink{sga7.sga7i.evii.n3-8-2}{3.8.2}), the
forms $f$ and $g$ are characterized by
\begin{equation}\tag{\SGARefTarget{sga7.sga7i.evii.n3-8-5}{3.8.5}}
 s(x_1,y_0)=f(x_1\otimes y_0)e_{E/Q}(y_0),\qquad
 s(x_0,y_1)=g(x_0\otimes y_1)e_{E/P}(x_0),
\end{equation}
which again makes clear the necessity of condition (\SGARefLink{sga7.sga7i.evii.n3-8-4}{3.8.4}).

}{E7I-VII-035}\SGApage{I}{\(\approx\)213}
Once the trivialization $s$ has been chosen, the other
trivializations are the morphisms
$s':L_0\times M_0\longrightarrow E$ of the form
\begin{equation}\tag{\SGARefTarget{sga7.sga7i.evii.n3-8-6}{3.8.6}}
 s'(x_0,y_0)=s(x_0,y_0)+h(x_0\otimes y_0),
\end{equation}
where
\[
 h:L_0\otimes M_0\longrightarrow G
\]
is an arbitrary form.  It is then clear from (\SGARefLink{sga7.sga7i.evii.n3-8-5}{3.8.5}) that the forms
$f,g$ defined by $s$ are changed into
\begin{equation}\tag{\SGARefTarget{sga7.sga7i.evii.n3-8-7}{3.8.7}}
\begin{aligned}
 f'(x_1\otimes y_0)&=f(x_1\otimes y_0)+h(x_1\otimes y_0),\\
 g'(x_0\otimes y_1)&=g(x_0\otimes y_1)+h(x_0\otimes y_1).
\end{aligned}
\end{equation}
Thus we obtain:

\underline{Proposition} \SGARefTarget{sga7.sga7i.evii.n3-8-8}{3.8.8}.
\underline{The subgroup of} $\Biext^1(P,Q;G)$ \underline{formed by the classes
of biextensions whose inverse image under $(u,v)$ is a trivial
biextension of $(L_0,M_0)$} (that is, the kernel of
$\Biext^1(P,Q;G)\longrightarrow\Biext^1(L_0,M_0;G)$)
\underline{is canonically isomorphic, by the preceding description,
to the quotient of the group of pairs of forms $(f,g)$ in (\SGARefLink{sga7.sga7i.evii.n3-8-3}{3.8.3})
satisfying compatibility condition (\SGARefLink{sga7.sga7i.evii.n3-8-4}{3.8.4}), by the pairs of forms
which arise by restriction from a form}
$h:L_0\otimes M_0\longrightarrow G$.

\SGApage{I}{\(\approx\)214}
We now interpret this isomorphism in terms of homological algebra,
taking into account the canonical isomorphism
\[
 \Biext^1(P,Q;G)\simeq
 \Ext^1(P\mathbin{\overset{\mathbb L}{\otimes}}Q,G)
\]
of \SGARefLink{sga7.sga7i.evii.n3-6-5}{3.6.5}.  To this end, interpret the data (\SGARefLink{sga7.sga7i.evii.n3-8-1}{3.8.1}) as two complexes
resolving $P$ and $Q$, respectively:
\begin{equation}\tag{\SGARefTarget{sga7.sga7i.evii.n3-8-9}{3.8.9}}
 L_\bullet:\ 0\longrightarrow L_1\longrightarrow L_0
 \longrightarrow0,\qquad
 M_\bullet:\ 0\longrightarrow M_1\longrightarrow M_0
 \longrightarrow0,
\end{equation}
and consider the tensor-product complex $L_\bullet\otimes M_\bullet$
\[
 0\longrightarrow L_1\otimes M_1
 \xrightarrow{d_2}
 L_1\otimes M_0+L_0\otimes M_1
 \xrightarrow{d_1}
 L_0\otimes M_0\longrightarrow0.
\]
There is, moreover, a canonical homomorphism in the derived category
\[
 P\mathbin{\overset{\mathbb L}{\otimes}}Q
 \longrightarrow L_\bullet\otimes M_\bullet,
\]
and hence an equally canonical homomorphism
\begin{equation}\tag{\SGARefTarget{sga7.sga7i.evii.n3-8-10}{3.8.10}}
 H^1\!\left(\Hom^{\cdot}(L_\bullet\otimes M_\bullet,G)\right)
 \longrightarrow
 \Ext^1(P\mathbin{\overset{\mathbb L}{\otimes}}Q,G).
\end{equation}
The group of $1$-cocycles of
$\Hom^{\cdot}(L_\bullet\otimes M_\bullet,G)$ identifies with the
group of pairs $(f,g)$ as in (\SGARefLink{sga7.sga7i.evii.n3-8-3}{3.8.3}) which satisfy compatibility
condition (\SGARefLink{sga7.sga7i.evii.n3-8-4}{3.8.4}): to such a pair one associates the homomorphism
$F:L_1\otimes M_0+L_0\otimes M_1\longrightarrow G$ with components
$(f,g)$,
\begin{equation}\tag{\SGARefTarget{sga7.sga7i.evii.n3-8-11}{3.8.11}}
 F(x_1\otimes y_0+x_0\otimes y_1)
 =f(x_1\otimes y_0)+g(x_0\otimes y_1).
\end{equation}
Moreover, the coboundary of the $0$-cochain
$h:L_0\otimes M_0\longrightarrow G$ is precisely the pair of its
restrictions to $L_1\otimes M_0$ and $L_0\otimes M_1$.  Consequently,
the group of pairs $(f,g)$ modulo the pairs $(h,h)$ described in
\SGARefLink{sga7.sga7i.evii.n3-8-8}{3.8.8} is canonically isomorphic to the first term of (\SGARefLink{sga7.sga7i.evii.n3-8-10}{3.8.10}).
With this understood, one has the following compatibility.

\SGApage{I}{215}
\underline{Proposition} \SGARefTarget{sga7.sga7i.evii.n3-8-12}{3.8.12}.
\underline{If}
\[
 \xi\in\Ker\!\left(
 \Biext^1(P,Q;G)\longrightarrow\Biext^1(L_0,M_0;G)\right)
\]
\underline{is described by the pair $(f,g)$ as in \SGARefLink{sga7.sga7i.evii.n3-8-8}{3.8.8}, then the
image of $\xi$ in}
$\Ext^1(P\mathbin{\overset{\mathbb L}{\otimes}}Q,G)$
\underline{under the isomorphism (\SGARefLink{sga7.sga7i.evii.n3-6-5-1}{3.6.5.1}) is the image under
(\SGARefLink{sga7.sga7i.evii.n3-8-10}{3.8.10}) of the element of its first term defined by the cocycle
(\SGARefLink{sga7.sga7i.evii.n3-8-11}{3.8.11})}.

The verification of this compatibility is left to the reader.

\SGARefTarget{sga7.sga7i.evii.n3-8-13}{3.8.13}. We shall apply this especially in the case
\[
 L_0=\mathbb Z[P],\qquad M_0=\mathbb Z[Q],
\]
where $L_1$ and $M_1$ are respectively the kernels of the canonical
homomorphisms $\mathbb Z[P]\longrightarrow P$ and
$\mathbb Z[Q]\longrightarrow Q$.  Notice the canonical isomorphisms
\[
 L_0\mathbin{\overset{\mathbb L}{\otimes}}M_0
 \simeq L_0\otimes M_0
 \simeq\mathbb Z[P\times Q].
\]
Since $L_0$ and $M_0$ are flat, (\SGARefLink{sga7.sga7i.evii.n1-4-5}{1.4.5}) gives canonical isomorphisms
\[
 \Ext^i(L_0\mathbin{\overset{\mathbb L}{\otimes}}M_0,G)
 \simeq H^i(P\times Q,G_{P\times Q}).
\]
It follows from this and \SGARefLink{sga7.sga7i.evii.n3-6-5}{3.6.5}, or directly from the definitions and
\SGARefLink{sga7.sga7i.evii.n1-4}{1.4}, that the functor which assigns to each biextension of
$(L_0,M_0)$ by $G$ the $G_{P\times Q}$-torsor over $P\times Q$ that
it induces is an equivalence of categories.  (This remains valid
without using an additive structure on $P,Q$.)  Consequently, if
$E$ is a biextension of $(P,Q)$ by $G$, giving a trivialization of
its inverse image over $(L_0,M_0)$ is equivalent to giving a section
of $E$ considered only as an object over $P\times Q$, that is, a
trivialization of $E$ considered as a torsor under $G_{P\times Q}$.

\SGApage{I}{\(\approx\)216}
Biextensions equipped with this additional structure are therefore
described by the pairs $(f,g)$ of (\SGARefLink{sga7.sga7i.evii.n3-8-3}{3.8.3}) satisfying compatibility
(\SGARefLink{sga7.sga7i.evii.n3-8-4}{3.8.4}), and their classification up to isomorphism is the one given
in \SGARefLink{sga7.sga7i.evii.n3-8-8}{3.8.8}.

With the notation of \SGARefLink{sga7.sga7i.evii.n3-5}{3.5}, one may also regard $L_1$ as the cokernel
of $L_2(P)\longrightarrow L_1(P)$, and similarly for $M_1$.  This
makes it possible to express the data $(f,g)$ in terms of a system
of homomorphisms
\[
\left\{
\begin{aligned}
 f'&:L_1(P)\otimes L_0(Q)
     =\mathbb Z(P\times P\times Q)\longrightarrow G,\\
 g'&:L_0(P)\otimes L_1(Q)
     =\mathbb Z(P\times Q\times Q)\longrightarrow G,
\end{aligned}
\right.
\]
that is, homomorphisms of sheaves of \underline{sets}
\[
 P\times P\times Q\longrightarrow G,\qquad
 P\times Q\times Q\longrightarrow G,
\]
satisfying five compatibility conditions.  These express that $f'$
and $g'$ factor through $f$ and $g$---which gives $2\times2$
equations, in view of the expression of $L_2(P)$ and $L_2(Q)$ as
sums of two terms---together with condition (\SGARefLink{sga7.sga7i.evii.n3-8-4}{3.8.4}).  One then sees
that $f'$ and $g'$ are precisely the partial multiplication data
defining a biextension, and that the five conditions in question are
exactly the five axioms for biextensions considered in \SGARefLink{sga7.sga7i.evii.n2-1}{2.1} (cf. the
proof of \SGARefLink{sga7.sga7i.evii.n3-6-4}{3.6.4}).

\subsection*{Bibliography to Expos\'e VII}

\SGAvar{\begin{enumerate}[label={[\arabic*]}]
\item M. Demazure, J. Giraud, and M. Raynaud,
\emph{Schémas abéliens}, Séminaire de la Faculté des Sciences
d'Orsay 1967/68 (to appear).
\item J. Giraud, \emph{Cohomologie non abélienne}, thesis, Paris,
1967 (to appear).
\item A. Grothendieck, \emph{Sur quelques points d'Algèbre
homologique}, Tohoku Mathematical Journal 9 (1956), 119--221.
\item A. Grothendieck, \emph{Catégories cofibrées additives et
Complexe cotangent relatif}, Lectures Notes in Math. 79,
Springer, 1968.
\item J.-L. Verdier, \emph{Catégories dérivées des catégories
abéliennes} (to appear).
\item D. Mumford, \emph{Biextensions of formal groups}, Tata
Institute of Fundamental Research, 1968 (to appear).
\end{enumerate}


}{\begin{enumerate}[label={[\arabic*]}]
\item M. Demazure, J. Giraud, and M. Raynaud,
\emph{Schémas abéliens}, Séminaire de la Faculté des Sciences
d'Orsay 1967/68 (to appear).
\item J. Giraud, \emph{Cohomologie non abélienne}, thesis, Paris,
1967 (to appear).
\item A. Grothendieck, \emph{Sur quelques points d'Algèbre
homologique}, Tohoku Mathematical Journal 9 (1957), 119--221.
\item A. Grothendieck, \emph{Catégories cofibrées additives et
Complexe cotangent relatif}, Lecture Notes in Mathematics 79,
Springer, 1968.
\item J.-L. Verdier, \emph{Catégories dérivées des catégories
abéliennes} (to appear).
\item D. Mumford, \emph{Biextensions of formal groups}, Tata
Institute of Fundamental Research, 1968 (to appear).
\end{enumerate}


}{E7I-VII-036}\markboth{Exposé VII}{Exposé VII}

\ifnum\SGAform=1
\clearpage\section*{Critical apparatus}\begingroup\small\raggedright\sloppy
Readings are given in the exact TeX notation of the two English witnesses. Decimal suffixes identify the individual differences in each marked passage.

\par\medskip\noindent\hypertarget{app-E7I-VII-001.1}{}\textbf{E7I-VII-001.1}\quad printed folio 135.\par
\noindent Printed reading: {\ttfamily\footnotesize E\allowbreak{}}\par
\noindent Corrected reading: {\ttfamily\footnotesize \textbackslash{}underline\allowbreak{} T\allowbreak{}}\par
\noindent Source: {\ttfamily\footnotesize fr/sga7i/\{diplomatic\allowbreak{},current\}/expose\_VII\allowbreak{}\_body.tex;\allowbreak{} folio\allowbreak{} 135;\allowbreak{} formula\allowbreak{} 43\allowbreak{}}\par

\par\medskip\noindent\hypertarget{app-E7I-VII-001.2}{}\textbf{E7I-VII-001.2}\quad printed folio 136.\par
\noindent Printed reading: {\ttfamily\footnotesize P\allowbreak{}}\par
\noindent Corrected reading: {\ttfamily\footnotesize E\allowbreak{}}\par
\noindent Source: {\ttfamily\footnotesize fr/sga7i/\{diplomatic\allowbreak{},current\}/expose\_VII\allowbreak{}\_body.tex;\allowbreak{} folio\allowbreak{} 136;\allowbreak{} formula\allowbreak{} 88\allowbreak{}}\par

\par\medskip\noindent\hypertarget{app-E7I-VII-001.3}{}\textbf{E7I-VII-001.3}\quad printed folio 136.\par
\noindent Printed reading: \emph{(empty)}\par
\noindent Corrected reading: {\ttfamily\footnotesize '\allowbreak{}}\par
\noindent Source: {\ttfamily\footnotesize fr/sga7i/\{diplomatic\allowbreak{},current\}/expose\_VII\allowbreak{}\_body.tex;\allowbreak{} folio\allowbreak{} 136;\allowbreak{} formula\allowbreak{} 88\allowbreak{}}\par

\par\medskip\noindent\hypertarget{app-E7I-VII-001.4}{}\textbf{E7I-VII-001.4}\quad printed folio 136.\par
\noindent Printed reading: {\ttfamily\footnotesize P\allowbreak{}}\par
\noindent Corrected reading: {\ttfamily\footnotesize E\allowbreak{}}\par
\noindent Source: {\ttfamily\footnotesize fr/sga7i/\{diplomatic\allowbreak{},current\}/expose\_VII\allowbreak{}\_body.tex;\allowbreak{} folio\allowbreak{} 136;\allowbreak{} formula\allowbreak{} 89\allowbreak{}}\par

\par\medskip\noindent\hypertarget{app-E7I-VII-001.5}{}\textbf{E7I-VII-001.5}\quad printed folio 136.\par
\noindent Printed reading: {\ttfamily\footnotesize S\allowbreak{}}\par
\noindent Corrected reading: {\ttfamily\footnotesize P\allowbreak{}}\par
\noindent Source: {\ttfamily\footnotesize port/decisions/decis\allowbreak{}ions\_en7I\_r2\_002.jso\allowbreak{}n\#C724;\allowbreak{} fixes/FIXLOG\_7I\_CURR\allowbreak{}ENT.jsonl\#C724;\allowbreak{} 7I:286\allowbreak{}}\par

\par\medskip\noindent\hypertarget{app-E7I-VII-001.6}{}\textbf{E7I-VII-001.6}\quad printed folio 136.\par
\noindent Printed reading: {\ttfamily\footnotesize S\allowbreak{}}\par
\noindent Corrected reading: {\ttfamily\footnotesize P\allowbreak{}}\par
\noindent Source: {\ttfamily\footnotesize fr/sga7i/\{diplomatic\allowbreak{},current\}/expose\_VII\allowbreak{}\_body.tex;\allowbreak{} folio\allowbreak{} 136;\allowbreak{} formula\allowbreak{} 102\allowbreak{}}\par

\par\medskip\noindent\hypertarget{app-E7I-VII-002.1}{}\textbf{E7I-VII-002.1}\quad printed folio 137.\par
\noindent Printed reading: {\ttfamily\footnotesize P\allowbreak{}}\par
\noindent Corrected reading: {\ttfamily\footnotesize E\allowbreak{}}\par
\noindent Source: {\ttfamily\footnotesize fr/sga7i/\{diplomatic\allowbreak{},current\}/expose\_VII\allowbreak{}\_body.tex;\allowbreak{} folio\allowbreak{} 137;\allowbreak{} formula\allowbreak{} 123\allowbreak{}}\par

\par\medskip\noindent\hypertarget{app-E7I-VII-002.2}{}\textbf{E7I-VII-002.2}\quad printed folio 137.\par
\noindent Printed reading: {\ttfamily\footnotesize P\allowbreak{}}\par
\noindent Corrected reading: {\ttfamily\footnotesize E\allowbreak{}}\par
\noindent Source: {\ttfamily\footnotesize fr/sga7i/\{diplomatic\allowbreak{},current\}/expose\_VII\allowbreak{}\_body.tex;\allowbreak{} folio\allowbreak{} 137;\allowbreak{} formula\allowbreak{} 123\allowbreak{}}\par

\par\medskip\noindent\hypertarget{app-E7I-VII-002.3}{}\textbf{E7I-VII-002.3}\quad printed folio 137.\par
\noindent Printed reading: {\ttfamily\footnotesize P\allowbreak{}}\par
\noindent Corrected reading: {\ttfamily\footnotesize E\allowbreak{}}\par
\noindent Source: {\ttfamily\footnotesize fr/sga7i/\{diplomatic\allowbreak{},current\}/expose\_VII\allowbreak{}\_body.tex;\allowbreak{} folio\allowbreak{} 137;\allowbreak{} formula\allowbreak{} 123\allowbreak{}}\par

\par\medskip\noindent\hypertarget{app-E7I-VII-002.4}{}\textbf{E7I-VII-002.4}\quad printed folio 137.\par
\noindent Printed reading: {\ttfamily\footnotesize P\allowbreak{}}\par
\noindent Corrected reading: {\ttfamily\footnotesize E\allowbreak{}}\par
\noindent Source: {\ttfamily\footnotesize fr/sga7i/\{diplomatic\allowbreak{},current\}/expose\_VII\allowbreak{}\_body.tex;\allowbreak{} folio\allowbreak{} 137;\allowbreak{} formula\allowbreak{} 127\allowbreak{}}\par

\par\medskip\noindent\hypertarget{app-E7I-VII-002.5}{}\textbf{E7I-VII-002.5}\quad printed folio 137.\par
\noindent Printed reading: {\ttfamily\footnotesize P\allowbreak{}}\par
\noindent Corrected reading: {\ttfamily\footnotesize E\allowbreak{}}\par
\noindent Source: {\ttfamily\footnotesize fr/sga7i/\{diplomatic\allowbreak{},current\}/expose\_VII\allowbreak{}\_body.tex;\allowbreak{} folio\allowbreak{} 137;\allowbreak{} formula\allowbreak{} 128\allowbreak{}}\par

\par\medskip\noindent\hypertarget{app-E7I-VII-002.6}{}\textbf{E7I-VII-002.6}\quad printed folio 137.\par
\noindent Printed reading: {\ttfamily\footnotesize P\allowbreak{}}\par
\noindent Corrected reading: {\ttfamily\footnotesize E\allowbreak{}}\par
\noindent Source: {\ttfamily\footnotesize fr/sga7i/\{diplomatic\allowbreak{},current\}/expose\_VII\allowbreak{}\_body.tex;\allowbreak{} folio\allowbreak{} 137;\allowbreak{} formula\allowbreak{} 129\allowbreak{}}\par

\par\medskip\noindent\hypertarget{app-E7I-VII-002.7}{}\textbf{E7I-VII-002.7}\quad printed folio 137.\par
\noindent Printed reading: {\ttfamily\footnotesize P\allowbreak{}}\par
\noindent Corrected reading: {\ttfamily\footnotesize E\allowbreak{}}\par
\noindent Source: {\ttfamily\footnotesize fr/sga7i/\{diplomatic\allowbreak{},current\}/expose\_VII\allowbreak{}\_body.tex;\allowbreak{} folio\allowbreak{} 137;\allowbreak{} formula\allowbreak{} 129\allowbreak{}}\par

\par\medskip\noindent\hypertarget{app-E7I-VII-002.8}{}\textbf{E7I-VII-002.8}\quad printed folio 137.\par
\noindent Printed reading: {\ttfamily\footnotesize P\allowbreak{}}\par
\noindent Corrected reading: {\ttfamily\footnotesize E\allowbreak{}}\par
\noindent Source: {\ttfamily\footnotesize fr/sga7i/\{diplomatic\allowbreak{},current\}/expose\_VII\allowbreak{}\_body.tex;\allowbreak{} folio\allowbreak{} 137;\allowbreak{} formula\allowbreak{} 131\allowbreak{}}\par

\par\medskip\noindent\hypertarget{app-E7I-VII-002.9}{}\textbf{E7I-VII-002.9}\quad printed folio 137.\par
\noindent Printed reading: {\ttfamily\footnotesize P\allowbreak{}}\par
\noindent Corrected reading: {\ttfamily\footnotesize E\allowbreak{}}\par
\noindent Source: {\ttfamily\footnotesize fr/sga7i/\{diplomatic\allowbreak{},current\}/expose\_VII\allowbreak{}\_body.tex;\allowbreak{} folio\allowbreak{} 137;\allowbreak{} formula\allowbreak{} 131\allowbreak{}}\par

\par\medskip\noindent\hypertarget{app-E7I-VII-002.10}{}\textbf{E7I-VII-002.10}\quad printed folio 137.\par
\noindent Printed reading: {\ttfamily\footnotesize P\allowbreak{}}\par
\noindent Corrected reading: {\ttfamily\footnotesize E\allowbreak{}}\par
\noindent Source: {\ttfamily\footnotesize fr/sga7i/\{diplomatic\allowbreak{},current\}/expose\_VII\allowbreak{}\_body.tex;\allowbreak{} folio\allowbreak{} 137;\allowbreak{} formula\allowbreak{} 131\allowbreak{}}\par

\par\medskip\noindent\hypertarget{app-E7I-VII-002.11}{}\textbf{E7I-VII-002.11}\quad printed folio 137.\par
\noindent Printed reading: {\ttfamily\footnotesize P\allowbreak{}}\par
\noindent Corrected reading: {\ttfamily\footnotesize E\allowbreak{}}\par
\noindent Source: {\ttfamily\footnotesize fr/sga7i/\{diplomatic\allowbreak{},current\}/expose\_VII\allowbreak{}\_body.tex;\allowbreak{} folio\allowbreak{} 137;\allowbreak{} formula\allowbreak{} 131\allowbreak{}}\par

\par\medskip\noindent\hypertarget{app-E7I-VII-002.12}{}\textbf{E7I-VII-002.12}\quad printed folio 137.\par
\noindent Printed reading: {\ttfamily\footnotesize P\allowbreak{}}\par
\noindent Corrected reading: {\ttfamily\footnotesize E\allowbreak{}}\par
\noindent Source: {\ttfamily\footnotesize fr/sga7i/\{diplomatic\allowbreak{},current\}/expose\_VII\allowbreak{}\_body.tex;\allowbreak{} folio\allowbreak{} 137;\allowbreak{} formula\allowbreak{} 131\allowbreak{}}\par

\par\medskip\noindent\hypertarget{app-E7I-VII-002.13}{}\textbf{E7I-VII-002.13}\quad printed folio 137.\par
\noindent Printed reading: {\ttfamily\footnotesize P\allowbreak{}}\par
\noindent Corrected reading: {\ttfamily\footnotesize E\allowbreak{}}\par
\noindent Source: {\ttfamily\footnotesize fr/sga7i/\{diplomatic\allowbreak{},current\}/expose\_VII\allowbreak{}\_body.tex;\allowbreak{} folio\allowbreak{} 137\allowbreak{}}\par

\par\medskip\noindent\hypertarget{app-E7I-VII-002.14}{}\textbf{E7I-VII-002.14}\quad printed folio 137.\par
\noindent Printed reading: {\ttfamily\footnotesize P\allowbreak{}}\par
\noindent Corrected reading: {\ttfamily\footnotesize E\allowbreak{}}\par
\noindent Source: {\ttfamily\footnotesize fr/sga7i/\{diplomatic\allowbreak{},current\}/expose\_VII\allowbreak{}\_body.tex;\allowbreak{} folio\allowbreak{} 137\allowbreak{}}\par

\par\medskip\noindent\hypertarget{app-E7I-VII-002.15}{}\textbf{E7I-VII-002.15}\quad printed folio 137.\par
\noindent Printed reading: {\ttfamily\footnotesize P\allowbreak{}}\par
\noindent Corrected reading: {\ttfamily\footnotesize E\allowbreak{}}\par
\noindent Source: {\ttfamily\footnotesize fr/sga7i/\{diplomatic\allowbreak{},current\}/expose\_VII\allowbreak{}\_body.tex;\allowbreak{} folio\allowbreak{} 137\allowbreak{}}\par

\par\medskip\noindent\hypertarget{app-E7I-VII-002.16}{}\textbf{E7I-VII-002.16}\quad printed folio 138.\par
\noindent Printed reading: {\ttfamily\footnotesize P\allowbreak{}}\par
\noindent Corrected reading: {\ttfamily\footnotesize E\allowbreak{}}\par
\noindent Source: {\ttfamily\footnotesize fr/sga7i/\{diplomatic\allowbreak{},current\}/expose\_VII\allowbreak{}\_body.tex;\allowbreak{} folio\allowbreak{} 138;\allowbreak{} formula\allowbreak{} 135\allowbreak{}}\par

\par\medskip\noindent\hypertarget{app-E7I-VII-002.17}{}\textbf{E7I-VII-002.17}\quad printed folio 138.\par
\noindent Printed reading: {\ttfamily\footnotesize P\allowbreak{}}\par
\noindent Corrected reading: {\ttfamily\footnotesize E\allowbreak{}}\par
\noindent Source: {\ttfamily\footnotesize fr/sga7i/\{diplomatic\allowbreak{},current\}/expose\_VII\allowbreak{}\_body.tex;\allowbreak{} folio\allowbreak{} 138;\allowbreak{} formula\allowbreak{} 136\allowbreak{}}\par

\par\medskip\noindent\hypertarget{app-E7I-VII-003.1}{}\textbf{E7I-VII-003.1}\quad printed folio 139.\par
\noindent Printed reading: {\ttfamily\footnotesize "\allowbreak{}}\par
\noindent Corrected reading: {\ttfamily\footnotesize ''\allowbreak{}}\par
\noindent Source: {\ttfamily\footnotesize fr/sga7i/\{diplomatic\allowbreak{},current\}/expose\_VII\allowbreak{}\_body.tex;\allowbreak{} folio\allowbreak{} 139;\allowbreak{} formula\allowbreak{} 167\allowbreak{}}\par

\par\medskip\noindent\hypertarget{app-E7I-VII-003.2}{}\textbf{E7I-VII-003.2}\quad printed folio 140.\par
\noindent Printed reading: {\ttfamily\footnotesize "\allowbreak{}}\par
\noindent Corrected reading: {\ttfamily\footnotesize ''\allowbreak{}}\par
\noindent Source: {\ttfamily\footnotesize ED-EN-PORT-WITNESS-7\allowbreak{}I-C1025;\allowbreak{} fr/sga7i/diplomatic/\allowbreak{}expose\_VII\_body.tex;\allowbreak{} folio\allowbreak{} 140;\allowbreak{} fr/sga7i/\{diplomatic\allowbreak{},current\}/expose\_VII\allowbreak{}\_body.tex;\allowbreak{} folio\allowbreak{} 140;\allowbreak{} formula\allowbreak{} 190\allowbreak{}}\par

\par\medskip\noindent\hypertarget{app-E7I-VII-003.3}{}\textbf{E7I-VII-003.3}\quad printed folio 140.\par
\noindent Printed reading: {\ttfamily\footnotesize P\allowbreak{}}\par
\noindent Corrected reading: {\ttfamily\footnotesize E\allowbreak{}}\par
\noindent Source: {\ttfamily\footnotesize fr/sga7i/\{diplomatic\allowbreak{},current\}/expose\_VII\allowbreak{}\_body.tex;\allowbreak{} folio\allowbreak{} 140;\allowbreak{} formula\allowbreak{} 199\allowbreak{}}\par

\par\medskip\noindent\hypertarget{app-E7I-VII-004.1}{}\textbf{E7I-VII-004.1}\quad printed folio 142.\par
\noindent Printed reading: {\ttfamily\footnotesize P\allowbreak{}}\par
\noindent Corrected reading: {\ttfamily\footnotesize E\allowbreak{}}\par
\noindent Source: {\ttfamily\footnotesize fr/sga7i/\{diplomatic\allowbreak{},current\}/expose\_VII\allowbreak{}\_body.tex;\allowbreak{} folio\allowbreak{} 142;\allowbreak{} formula\allowbreak{} 255\allowbreak{}}\par

\par\medskip\noindent\hypertarget{app-E7I-VII-004.2}{}\textbf{E7I-VII-004.2}\quad printed folio 142.\par
\noindent Printed reading: {\ttfamily\footnotesize P\allowbreak{}}\par
\noindent Corrected reading: {\ttfamily\footnotesize E\allowbreak{}}\par
\noindent Source: {\ttfamily\footnotesize fr/sga7i/\{diplomatic\allowbreak{},current\}/expose\_VII\allowbreak{}\_body.tex;\allowbreak{} folio\allowbreak{} 142;\allowbreak{} formula\allowbreak{} 255\allowbreak{}}\par

\par\medskip\noindent\hypertarget{app-E7I-VII-005.1}{}\textbf{E7I-VII-005.1}\quad printed folio 144.\par
\noindent Printed reading: {\ttfamily\footnotesize P\allowbreak{}}\par
\noindent Corrected reading: {\ttfamily\footnotesize E\allowbreak{}}\par
\noindent Source: {\ttfamily\footnotesize fr/sga7i/\{diplomatic\allowbreak{},current\}/expose\_VII\allowbreak{}\_body.tex;\allowbreak{} folio\allowbreak{} 144;\allowbreak{} formula\allowbreak{} 296\allowbreak{}}\par

\par\medskip\noindent\hypertarget{app-E7I-VII-006.1}{}\textbf{E7I-VII-006.1}\quad printed folio 146.\par
\noindent Printed reading: \emph{(empty)}\par
\noindent Corrected reading: {\ttfamily\footnotesize \textasciicircum{}*\allowbreak{}}\par
\noindent Source: {\ttfamily\footnotesize fr/sga7i/\{diplomatic\allowbreak{},current\}/expose\_VII\allowbreak{}\_body.tex;\allowbreak{} folio\allowbreak{} 146;\allowbreak{} formula\allowbreak{} 417\allowbreak{}}\par

\par\medskip\noindent\hypertarget{app-E7I-VII-006.2}{}\textbf{E7I-VII-006.2}\quad printed folio 146.\par
\noindent Printed reading: \emph{(empty)}\par
\noindent Corrected reading: {\ttfamily\footnotesize \textasciicircum{}*\allowbreak{}}\par
\noindent Source: {\ttfamily\footnotesize fr/sga7i/\{diplomatic\allowbreak{},current\}/expose\_VII\allowbreak{}\_body.tex;\allowbreak{} folio\allowbreak{} 146;\allowbreak{} formula\allowbreak{} 420\allowbreak{}}\par

\par\medskip\noindent\hypertarget{app-E7I-VII-007.1}{}\textbf{E7I-VII-007.1}\quad printed folio 148.\par
\noindent Printed reading: {\ttfamily\footnotesize Q\allowbreak{}}\par
\noindent Corrected reading: {\ttfamily\footnotesize G\allowbreak{}}\par
\noindent Source: {\ttfamily\footnotesize fr/sga7i/\{diplomatic\allowbreak{},current\}/expose\_VII\allowbreak{}\_body.tex;\allowbreak{} folio\allowbreak{} 148;\allowbreak{} formula\allowbreak{} 470\allowbreak{}}\par

\par\medskip\noindent\hypertarget{app-E7I-VII-008.1}{}\textbf{E7I-VII-008.1}\quad printed folio 148.\par
\noindent Printed reading: {\ttfamily\footnotesize torsor\allowbreak{}}\par
\noindent Corrected reading: {\ttfamily\footnotesize bitorsor\allowbreak{}}\par
\noindent Source: {\ttfamily\footnotesize 7I:300;\allowbreak{} fixes/MISPRINTS\_7I.j\allowbreak{}son\#300\allowbreak{}}\par

\par\medskip\noindent\hypertarget{app-E7I-VII-008.2}{}\textbf{E7I-VII-008.2}\quad printed folio 148.\par
\noindent Printed reading: \emph{(empty)}\par
\noindent Corrected reading: {\ttfamily\footnotesize (\allowbreak{}}\par
\noindent Source: {\ttfamily\footnotesize 7I:300;\allowbreak{} fixes/MISPRINTS\_7I.j\allowbreak{}son\#300;\allowbreak{} fr/sga7i/\{diplomatic\allowbreak{},current\}/expose\_VII\allowbreak{}\_body.tex;\allowbreak{} folio\allowbreak{} 148;\allowbreak{} formula\allowbreak{} 491\allowbreak{}}\par

\par\medskip\noindent\hypertarget{app-E7I-VII-008.3}{}\textbf{E7I-VII-008.3}\quad printed folio 148.\par
\noindent Printed reading: \emph{(empty)}\par
\noindent Corrected reading: {\ttfamily\footnotesize ,G\_P)\allowbreak{}}\par
\noindent Source: {\ttfamily\footnotesize 7I:300;\allowbreak{} fixes/MISPRINTS\_7I.j\allowbreak{}son\#300;\allowbreak{} fr/sga7i/\{diplomatic\allowbreak{},current\}/expose\_VII\allowbreak{}\_body.tex;\allowbreak{} folio\allowbreak{} 148;\allowbreak{} formula\allowbreak{} 491\allowbreak{}}\par

\par\medskip\noindent\hypertarget{app-E7I-VII-009.1}{}\textbf{E7I-VII-009.1}\quad printed folio 149.\par
\noindent Printed reading: {\ttfamily\footnotesize f(f(\allowbreak{}}\par
\noindent Corrected reading: \emph{(empty)}\par
\noindent Source: {\ttfamily\footnotesize port/decisions/decis\allowbreak{}ions\_en7I\_r2\_002.jso\allowbreak{}n\#C731;\allowbreak{} 7I:301\allowbreak{}}\par

\par\medskip\noindent\hypertarget{app-E7I-VII-009.2}{}\textbf{E7I-VII-009.2}\quad printed folio 149.\par
\noindent Printed reading: {\ttfamily\footnotesize ))\allowbreak{}}\par
\noindent Corrected reading: \emph{(empty)}\par
\noindent Source: {\ttfamily\footnotesize port/decisions/decis\allowbreak{}ions\_en7I\_r2\_002.jso\allowbreak{}n\#C731;\allowbreak{} 7I:301\allowbreak{}}\par

\par\medskip\noindent\hypertarget{app-E7I-VII-009.3}{}\textbf{E7I-VII-009.3}\quad printed folio 149.\par
\noindent Printed reading: {\ttfamily\footnotesize f(f(\allowbreak{}}\par
\noindent Corrected reading: \emph{(empty)}\par
\noindent Source: {\ttfamily\footnotesize port/decisions/decis\allowbreak{}ions\_en7I\_r2\_002.jso\allowbreak{}n\#C731;\allowbreak{} 7I:301\allowbreak{}}\par

\par\medskip\noindent\hypertarget{app-E7I-VII-009.4}{}\textbf{E7I-VII-009.4}\quad printed folio 149.\par
\noindent Printed reading: {\ttfamily\footnotesize f\allowbreak{}}\par
\noindent Corrected reading: {\ttfamily\footnotesize j\allowbreak{}}\par
\noindent Source: {\ttfamily\footnotesize port/decisions/decis\allowbreak{}ions\_en7I\_r2\_002.jso\allowbreak{}n\#C731;\allowbreak{} 7I:301;\allowbreak{} fr/sga7i/\{diplomatic\allowbreak{},current\}/expose\_VII\allowbreak{}\_body.tex;\allowbreak{} folio\allowbreak{} 149;\allowbreak{} formula\allowbreak{} 520\allowbreak{}}\par

\par\medskip\noindent\hypertarget{app-E7I-VII-009.5}{}\textbf{E7I-VII-009.5}\quad printed folio 149.\par
\noindent Printed reading: {\ttfamily\footnotesize ))\allowbreak{}}\par
\noindent Corrected reading: \emph{(empty)}\par
\noindent Source: {\ttfamily\footnotesize port/decisions/decis\allowbreak{}ions\_en7I\_r2\_002.jso\allowbreak{}n\#C731;\allowbreak{} 7I:301\allowbreak{}}\par

\par\medskip\noindent\hypertarget{app-E7I-VII-009.6}{}\textbf{E7I-VII-009.6}\quad printed folio 150.\par
\noindent Printed reading: \emph{(empty)}\par
\noindent Corrected reading: {\ttfamily\footnotesize \textasciicircum{}*\allowbreak{}}\par
\noindent Source: {\ttfamily\footnotesize fr/sga7i/\{diplomatic\allowbreak{},current\}/expose\_VII\allowbreak{}\_body.tex;\allowbreak{} folio\allowbreak{} 150;\allowbreak{} formula\allowbreak{} 537\allowbreak{}}\par

\par\medskip\noindent\hypertarget{app-E7I-VII-010.1}{}\textbf{E7I-VII-010.1}\quad printed folio 151.\par
\noindent Printed reading: {\ttfamily\footnotesize y=\allowbreak{}}\par
\noindent Corrected reading: \emph{(empty)}\par
\noindent Source: {\ttfamily\footnotesize fr/sga7i/\{diplomatic\allowbreak{},current\}/expose\_VII\allowbreak{}\_body.tex;\allowbreak{} folio\allowbreak{} 151;\allowbreak{} formula\allowbreak{} 568\allowbreak{}}\par

\par\medskip\noindent\hypertarget{app-E7I-VII-010.2}{}\textbf{E7I-VII-010.2}\quad printed folio 151.\par
\noindent Printed reading: {\ttfamily\footnotesize z=\allowbreak{}}\par
\noindent Corrected reading: \emph{(empty)}\par
\noindent Source: {\ttfamily\footnotesize fr/sga7i/\{diplomatic\allowbreak{},current\}/expose\_VII\allowbreak{}\_body.tex;\allowbreak{} folio\allowbreak{} 151;\allowbreak{} formula\allowbreak{} 569\allowbreak{}}\par

\par\medskip\noindent\hypertarget{app-E7I-VII-010.3}{}\textbf{E7I-VII-010.3}\quad printed folio 151.\par
\noindent Printed reading: {\ttfamily\footnotesize x=\allowbreak{}}\par
\noindent Corrected reading: \emph{(empty)}\par
\noindent Source: {\ttfamily\footnotesize fr/sga7i/\{diplomatic\allowbreak{},current\}/expose\_VII\allowbreak{}\_body.tex;\allowbreak{} folio\allowbreak{} 151;\allowbreak{} formula\allowbreak{} 570\allowbreak{}}\par

\par\medskip\noindent\hypertarget{app-E7I-VII-010.4}{}\textbf{E7I-VII-010.4}\quad printed folio 151.\par
\noindent Printed reading: {\ttfamily\footnotesize e\_\allowbreak{}}\par
\noindent Corrected reading: \emph{(empty)}\par
\noindent Source: {\ttfamily\footnotesize fr/sga7i/\{diplomatic\allowbreak{},current\}/expose\_VII\allowbreak{}\_body.tex;\allowbreak{} folio\allowbreak{} 151;\allowbreak{} formula\allowbreak{} 573\allowbreak{}}\par

\par\medskip\noindent\hypertarget{app-E7I-VII-010.5}{}\textbf{E7I-VII-010.5}\quad printed folio 151.\par
\noindent Printed reading: {\ttfamily\footnotesize e\_\allowbreak{}}\par
\noindent Corrected reading: \emph{(empty)}\par
\noindent Source: {\ttfamily\footnotesize fr/sga7i/\{diplomatic\allowbreak{},current\}/expose\_VII\allowbreak{}\_body.tex;\allowbreak{} folio\allowbreak{} 151;\allowbreak{} formula\allowbreak{} 574\allowbreak{}}\par

\par\medskip\noindent\hypertarget{app-E7I-VII-010.6}{}\textbf{E7I-VII-010.6}\quad printed folio 151.\par
\noindent Printed reading: {\ttfamily\footnotesize e\_\allowbreak{}}\par
\noindent Corrected reading: \emph{(empty)}\par
\noindent Source: {\ttfamily\footnotesize fr/sga7i/\{diplomatic\allowbreak{},current\}/expose\_VII\allowbreak{}\_body.tex;\allowbreak{} folio\allowbreak{} 151;\allowbreak{} formula\allowbreak{} 575\allowbreak{}}\par

\par\medskip\noindent\hypertarget{app-E7I-VII-011.1}{}\textbf{E7I-VII-011.1}\quad printed folio 151.\par
\noindent Printed reading: {\ttfamily\footnotesize G\allowbreak{}}\par
\noindent Corrected reading: {\ttfamily\footnotesize E\allowbreak{}}\par
\noindent Source: {\ttfamily\footnotesize fr/sga7i/\{diplomatic\allowbreak{},current\}/expose\_VII\allowbreak{}\_body.tex;\allowbreak{} folio\allowbreak{} 151;\allowbreak{} formula\allowbreak{} 590\allowbreak{}}\par

\par\medskip\noindent\hypertarget{app-E7I-VII-012.1}{}\textbf{E7I-VII-012.1}\quad printed folio 153.\par
\noindent Printed reading: \emph{(empty)}\par
\noindent Corrected reading: {\ttfamily\footnotesize \textbackslash{}operatorname\{\allowbreak{}}\par
\noindent Source: {\ttfamily\footnotesize fr/sga7i/\{diplomatic\allowbreak{},current\}/expose\_VII\allowbreak{}\_body.tex;\allowbreak{} folio\allowbreak{} 153;\allowbreak{} formula\allowbreak{} 648\allowbreak{}}\par

\par\medskip\noindent\hypertarget{app-E7I-VII-012.2}{}\textbf{E7I-VII-012.2}\quad printed folio 153.\par
\noindent Printed reading: \emph{(empty)}\par
\noindent Corrected reading: {\ttfamily\footnotesize \}\allowbreak{}}\par
\noindent Source: {\ttfamily\footnotesize fr/sga7i/\{diplomatic\allowbreak{},current\}/expose\_VII\allowbreak{}\_body.tex;\allowbreak{} folio\allowbreak{} 153;\allowbreak{} formula\allowbreak{} 648\allowbreak{}}\par

\par\medskip\noindent\hypertarget{app-E7I-VII-012.3}{}\textbf{E7I-VII-012.3}\quad printed folio 153.\par
\noindent Printed reading: \emph{(empty)}\par
\noindent Corrected reading: {\ttfamily\footnotesize \textbackslash{}underline\allowbreak{}}\par
\noindent Source: {\ttfamily\footnotesize fr/sga7i/\{diplomatic\allowbreak{},current\}/expose\_VII\allowbreak{}\_body.tex;\allowbreak{} folio\allowbreak{} 153;\allowbreak{} formula\allowbreak{} 648\allowbreak{}}\par

\par\medskip\noindent\hypertarget{app-E7I-VII-013.1}{}\textbf{E7I-VII-013.1}\quad printed folio 154.\par
\noindent Printed reading: \emph{(empty)}\par
\noindent Corrected reading: {\ttfamily\footnotesize 1.\allowbreak{}}\par
\noindent Source: {\ttfamily\footnotesize 7I:305;\allowbreak{} fixes/MISPRINTS\_7I.j\allowbreak{}son\#305\allowbreak{}}\par

\par\medskip\noindent\hypertarget{app-E7I-VII-013.2}{}\textbf{E7I-VII-013.2}\quad printed folio 154.\par
\noindent Printed reading: \emph{(empty)}\par
\noindent Corrected reading: {\ttfamily\footnotesize 1.\allowbreak{}}\par
\noindent Source: {\ttfamily\footnotesize 7I:305;\allowbreak{} fixes/MISPRINTS\_7I.j\allowbreak{}son\#305\allowbreak{}}\par

\par\medskip\noindent\hypertarget{app-E7I-VII-014.1}{}\textbf{E7I-VII-014.1}\quad printed folio 157.\par
\noindent Printed reading: \emph{(empty)}\par
\noindent Corrected reading: {\ttfamily\footnotesize P\_\allowbreak{}}\par
\noindent Source: {\ttfamily\footnotesize ED-EN-FORMULA-G019;\allowbreak{} fr/sga7i/\{diplomatic\allowbreak{},current\}/expose\_VII\allowbreak{}\_body.tex;\allowbreak{} folio\allowbreak{} 157\allowbreak{}}\par

\par\medskip\noindent\hypertarget{app-E7I-VII-014.2}{}\textbf{E7I-VII-014.2}\quad printed folio 157.\par
\noindent Printed reading: {\ttfamily\footnotesize \_P\allowbreak{}}\par
\noindent Corrected reading: \emph{(empty)}\par
\noindent Source: {\ttfamily\footnotesize ED-EN-FORMULA-G019;\allowbreak{} fr/sga7i/\{diplomatic\allowbreak{},current\}/expose\_VII\allowbreak{}\_body.tex;\allowbreak{} folio\allowbreak{} 157\allowbreak{}}\par

\par\medskip\noindent\hypertarget{app-E7I-VII-014.3}{}\textbf{E7I-VII-014.3}\quad printed folio 157.\par
\noindent Printed reading: {\ttfamily\footnotesize P\allowbreak{}}\par
\noindent Corrected reading: {\ttfamily\footnotesize Q\allowbreak{}}\par
\noindent Source: {\ttfamily\footnotesize fr/sga7i/\{diplomatic\allowbreak{},current\}/expose\_VII\allowbreak{}\_body.tex;\allowbreak{} folio\allowbreak{} 157;\allowbreak{} formula\allowbreak{} 765\allowbreak{}}\par

\par\medskip\noindent\hypertarget{app-E7I-VII-014.4}{}\textbf{E7I-VII-014.4}\quad printed folio 157.\par
\noindent Printed reading: \emph{(empty)}\par
\noindent Corrected reading: {\ttfamily\footnotesize '\allowbreak{}}\par
\noindent Source: {\ttfamily\footnotesize fr/sga7i/\{diplomatic\allowbreak{},current\}/expose\_VII\allowbreak{}\_body.tex;\allowbreak{} folio\allowbreak{} 157;\allowbreak{} formula\allowbreak{} 782\allowbreak{}}\par

\par\medskip\noindent\hypertarget{app-E7I-VII-014.5}{}\textbf{E7I-VII-014.5}\quad printed folio 158.\par
\noindent Printed reading: \emph{(empty)}\par
\noindent Corrected reading: {\ttfamily\footnotesize Q\_\allowbreak{}}\par
\noindent Source: {\ttfamily\footnotesize fr/sga7i/\{diplomatic\allowbreak{},current\}/expose\_VII\allowbreak{}\_body.tex;\allowbreak{} folio\allowbreak{} 158;\allowbreak{} formula\allowbreak{} 787\allowbreak{}}\par

\par\medskip\noindent\hypertarget{app-E7I-VII-014.6}{}\textbf{E7I-VII-014.6}\quad printed folio 158.\par
\noindent Printed reading: {\ttfamily\footnotesize \_Q\allowbreak{}}\par
\noindent Corrected reading: \emph{(empty)}\par
\noindent Source: {\ttfamily\footnotesize fr/sga7i/\{diplomatic\allowbreak{},current\}/expose\_VII\allowbreak{}\_body.tex;\allowbreak{} folio\allowbreak{} 158;\allowbreak{} formula\allowbreak{} 787\allowbreak{}}\par

\par\medskip\noindent\hypertarget{app-E7I-VII-014.7}{}\textbf{E7I-VII-014.7}\quad printed folio 158.\par
\noindent Printed reading: {\ttfamily\footnotesize Q\allowbreak{}}\par
\noindent Corrected reading: {\ttfamily\footnotesize P\allowbreak{}}\par
\noindent Source: {\ttfamily\footnotesize fr/sga7i/\{diplomatic\allowbreak{},current\}/expose\_VII\allowbreak{}\_body.tex;\allowbreak{} folio\allowbreak{} 158;\allowbreak{} formula\allowbreak{} 788\allowbreak{}}\par

\par\medskip\noindent\hypertarget{app-E7I-VII-014.8}{}\textbf{E7I-VII-014.8}\quad printed folio 158.\par
\noindent Printed reading: {\ttfamily\footnotesize pp\allowbreak{}}\par
\noindent Corrected reading: {\ttfamily\footnotesize p,qq\allowbreak{}}\par
\noindent Source: {\ttfamily\footnotesize fr/sga7i/\{diplomatic\allowbreak{},current\}/expose\_VII\allowbreak{}\_body.tex;\allowbreak{} folio\allowbreak{} 158;\allowbreak{} formula\allowbreak{} 791\allowbreak{}}\par

\par\medskip\noindent\hypertarget{app-E7I-VII-014.9}{}\textbf{E7I-VII-014.9}\quad printed folio 158.\par
\noindent Printed reading: {\ttfamily\footnotesize ,q\allowbreak{}}\par
\noindent Corrected reading: \emph{(empty)}\par
\noindent Source: {\ttfamily\footnotesize fr/sga7i/\{diplomatic\allowbreak{},current\}/expose\_VII\allowbreak{}\_body.tex;\allowbreak{} folio\allowbreak{} 158;\allowbreak{} formula\allowbreak{} 791\allowbreak{}}\par

\par\medskip\noindent\hypertarget{app-E7I-VII-014.10}{}\textbf{E7I-VII-014.10}\quad printed folio 158.\par
\noindent Printed reading: {\ttfamily\footnotesize ;\allowbreak{}}\par
\noindent Corrected reading: {\ttfamily\footnotesize ,\allowbreak{}}\par
\noindent Source: {\ttfamily\footnotesize fr/sga7i/\{diplomatic\allowbreak{},current\}/expose\_VII\allowbreak{}\_body.tex;\allowbreak{} folio\allowbreak{} 158;\allowbreak{} formula\allowbreak{} 796\allowbreak{}}\par

\par\medskip\noindent\hypertarget{app-E7I-VII-014.11}{}\textbf{E7I-VII-014.11}\quad printed folio 159.\par
\noindent Printed reading: {\ttfamily\footnotesize of\allowbreak{}}\par
\noindent Corrected reading: {\ttfamily\footnotesize in\allowbreak{} the\allowbreak{}}\par
\noindent Source: {\ttfamily\footnotesize fr/sga7i/\{diplomatic\allowbreak{},current\}/expose\_VII\allowbreak{}\_body.tex;\allowbreak{} folio\allowbreak{} 159\allowbreak{}}\par

\par\medskip\noindent\hypertarget{app-E7I-VII-014.12}{}\textbf{E7I-VII-014.12}\quad printed folio 159.\par
\noindent Printed reading: {\ttfamily\footnotesize Q\_\allowbreak{}}\par
\noindent Corrected reading: \emph{(empty)}\par
\noindent Source: {\ttfamily\footnotesize fr/sga7i/\{diplomatic\allowbreak{},current\}/expose\_VII\allowbreak{}\_body.tex;\allowbreak{} folio\allowbreak{} 159\allowbreak{}}\par

\par\medskip\noindent\hypertarget{app-E7I-VII-014.13}{}\textbf{E7I-VII-014.13}\quad printed folio 159.\par
\noindent Printed reading: {\ttfamily\footnotesize  by\allowbreak{} \$G\_P\$\allowbreak{}}\par
\noindent Corrected reading: {\ttfamily\footnotesize -direction\allowbreak{}}\par
\noindent Source: {\ttfamily\footnotesize fr/sga7i/\{diplomatic\allowbreak{},current\}/expose\_VII\allowbreak{}\_body.tex;\allowbreak{} folio\allowbreak{} 159\allowbreak{}}\par

\par\medskip\noindent\hypertarget{app-E7I-VII-014.14}{}\textbf{E7I-VII-014.14}\quad printed folio 159.\par
\noindent Printed reading: {\ttfamily\footnotesize of\allowbreak{}}\par
\noindent Corrected reading: {\ttfamily\footnotesize in\allowbreak{} the\allowbreak{}}\par
\noindent Source: {\ttfamily\footnotesize fr/sga7i/\{diplomatic\allowbreak{},current\}/expose\_VII\allowbreak{}\_body.tex;\allowbreak{} folio\allowbreak{} 159\allowbreak{}}\par

\par\medskip\noindent\hypertarget{app-E7I-VII-014.15}{}\textbf{E7I-VII-014.15}\quad printed folio 159.\par
\noindent Printed reading: {\ttfamily\footnotesize P\_\allowbreak{}}\par
\noindent Corrected reading: \emph{(empty)}\par
\noindent Source: {\ttfamily\footnotesize fr/sga7i/\{diplomatic\allowbreak{},current\}/expose\_VII\allowbreak{}\_body.tex;\allowbreak{} folio\allowbreak{} 159\allowbreak{}}\par

\par\medskip\noindent\hypertarget{app-E7I-VII-014.16}{}\textbf{E7I-VII-014.16}\quad printed folio 159.\par
\noindent Printed reading: {\ttfamily\footnotesize  by\allowbreak{} \$G\_Q\$\allowbreak{}}\par
\noindent Corrected reading: {\ttfamily\footnotesize -direction\allowbreak{}}\par
\noindent Source: {\ttfamily\footnotesize fr/sga7i/\{diplomatic\allowbreak{},current\}/expose\_VII\allowbreak{}\_body.tex;\allowbreak{} folio\allowbreak{} 159\allowbreak{}}\par

\par\medskip\noindent\hypertarget{app-E7I-VII-014.17}{}\textbf{E7I-VII-014.17}\quad printed folio 159.\par
\noindent Printed reading: {\ttfamily\footnotesize an\allowbreak{}}\par
\noindent Corrected reading: {\ttfamily\footnotesize compatible\allowbreak{}}\par
\noindent Source: {\ttfamily\footnotesize fr/sga7i/\{diplomatic\allowbreak{},current\}/expose\_VII\allowbreak{}\_body.tex;\allowbreak{} folio\allowbreak{} 159\allowbreak{}}\par

\par\medskip\noindent\hypertarget{app-E7I-VII-014.18}{}\textbf{E7I-VII-014.18}\quad printed folio 159.\par
\noindent Printed reading: {\ttfamily\footnotesize structure\allowbreak{}}\par
\noindent Corrected reading: {\ttfamily\footnotesize structures\allowbreak{}}\par
\noindent Source: {\ttfamily\footnotesize fr/sga7i/\{diplomatic\allowbreak{},current\}/expose\_VII\allowbreak{}\_body.tex;\allowbreak{} folio\allowbreak{} 159\allowbreak{}}\par

\par\medskip\noindent\hypertarget{app-E7I-VII-014.19}{}\textbf{E7I-VII-014.19}\quad printed folio 159.\par
\noindent Printed reading: {\ttfamily\footnotesize of\allowbreak{}}\par
\noindent Corrected reading: {\ttfamily\footnotesize in\allowbreak{} the\allowbreak{} \$P\$-\allowbreak{} and\allowbreak{}}\par
\noindent Source: {\ttfamily\footnotesize fr/sga7i/\{diplomatic\allowbreak{},current\}/expose\_VII\allowbreak{}\_body.tex;\allowbreak{} folio\allowbreak{} 159\allowbreak{}}\par

\par\medskip\noindent\hypertarget{app-E7I-VII-014.20}{}\textbf{E7I-VII-014.20}\quad printed folio 159.\par
\noindent Printed reading: {\ttfamily\footnotesize \_P\allowbreak{}}\par
\noindent Corrected reading: \emph{(empty)}\par
\noindent Source: {\ttfamily\footnotesize fr/sga7i/\{diplomatic\allowbreak{},current\}/expose\_VII\allowbreak{}\_body.tex;\allowbreak{} folio\allowbreak{} 159\allowbreak{}}\par

\par\medskip\noindent\hypertarget{app-E7I-VII-014.21}{}\textbf{E7I-VII-014.21}\quad printed folio 159.\par
\noindent Printed reading: {\ttfamily\footnotesize  by\allowbreak{} \$G\_P\$\allowbreak{} and\allowbreak{} an\allowbreak{} extension\allowbreak{} structure\allowbreak{} of\allowbreak{} \$G\_Q\$\allowbreak{} by\allowbreak{} \$P\_Q\$,\allowbreak{} these\allowbreak{} two\allowbreak{} structures\allowbreak{} being\allowbreak{} compatible\allowbreak{}}\par
\noindent Corrected reading: {\ttfamily\footnotesize -directions\allowbreak{}}\par
\noindent Source: {\ttfamily\footnotesize fr/sga7i/\{diplomatic\allowbreak{},current\}/expose\_VII\allowbreak{}\_body.tex;\allowbreak{} folio\allowbreak{} 159\allowbreak{}}\par

\par\medskip\noindent\hypertarget{app-E7I-VII-015.1}{}\textbf{E7I-VII-015.1}\quad printed folio 160.\par
\noindent Printed reading: {\ttfamily\footnotesize G\allowbreak{}}\par
\noindent Corrected reading: {\ttfamily\footnotesize Q\allowbreak{}}\par
\noindent Source: {\ttfamily\footnotesize fr/sga7i/\{diplomatic\allowbreak{},current\}/expose\_VII\allowbreak{}\_body.tex;\allowbreak{} folio\allowbreak{} 160;\allowbreak{} formula\allowbreak{} 844\allowbreak{}}\par

\par\medskip\noindent\hypertarget{app-E7I-VII-015.2}{}\textbf{E7I-VII-015.2}\quad printed folio 160.\par
\noindent Printed reading: \emph{(empty)}\par
\noindent Corrected reading: {\ttfamily\footnotesize P\textbackslash{}times\allowbreak{}}\par
\noindent Source: {\ttfamily\footnotesize ED-EN-FORMULA-G011;\allowbreak{} fr/sga7i/\{diplomatic\allowbreak{},current\}/expose\_VII\allowbreak{}\_body.tex;\allowbreak{} folio\allowbreak{} 160\allowbreak{}}\par

\par\medskip\noindent\hypertarget{app-E7I-VII-015.3}{}\textbf{E7I-VII-015.3}\quad printed folio 160.\par
\noindent Printed reading: {\ttfamily\footnotesize P\textbackslash{}times\allowbreak{} }\par
\noindent Corrected reading: \emph{(empty)}\par
\noindent Source: {\ttfamily\footnotesize ED-EN-FORMULA-G011;\allowbreak{} fr/sga7i/\{diplomatic\allowbreak{},current\}/expose\_VII\allowbreak{}\_body.tex;\allowbreak{} folio\allowbreak{} 160\allowbreak{}}\par

\par\medskip\noindent\hypertarget{app-E7I-VII-015.4}{}\textbf{E7I-VII-015.4}\quad printed folio 160.\par
\noindent Printed reading: \emph{(empty)}\par
\noindent Corrected reading: {\ttfamily\footnotesize 1\_\allowbreak{}}\par
\noindent Source: {\ttfamily\footnotesize ED-EN-FORMULA-G011;\allowbreak{} fr/sga7i/\{diplomatic\allowbreak{},current\}/expose\_VII\allowbreak{}\_body.tex;\allowbreak{} folio\allowbreak{} 160\allowbreak{}}\par

\par\medskip\noindent\hypertarget{app-E7I-VII-015.5}{}\textbf{E7I-VII-015.5}\quad printed folio 160.\par
\noindent Printed reading: {\ttfamily\footnotesize 1\_\allowbreak{}}\par
\noindent Corrected reading: \emph{(1 espace(s), 0 retour(s) de ligne dans le témoin TeX)}\par
\noindent Source: {\ttfamily\footnotesize ED-EN-FORMULA-G011;\allowbreak{} fr/sga7i/\{diplomatic\allowbreak{},current\}/expose\_VII\allowbreak{}\_body.tex;\allowbreak{} folio\allowbreak{} 160\allowbreak{}}\par

\par\medskip\noindent\hypertarget{app-E7I-VII-016.1}{}\textbf{E7I-VII-016.1}\quad printed folio 161.\par
\noindent Printed reading: {\ttfamily\footnotesize G\allowbreak{}}\par
\noindent Corrected reading: {\ttfamily\footnotesize E\allowbreak{}}\par
\noindent Source: {\ttfamily\footnotesize fr/sga7i/\{diplomatic\allowbreak{},current\}/expose\_VII\allowbreak{}\_body.tex;\allowbreak{} folio\allowbreak{} 161;\allowbreak{} formula\allowbreak{} 888\allowbreak{}}\par

\par\medskip\noindent\hypertarget{app-E7I-VII-017.1}{}\textbf{E7I-VII-017.1}\quad printed folio 163.\par
\noindent Printed reading: {\ttfamily\footnotesize G\allowbreak{}}\par
\noindent Corrected reading: {\ttfamily\footnotesize Q\allowbreak{}}\par
\noindent Source: {\ttfamily\footnotesize fr/sga7i/\{diplomatic\allowbreak{},current\}/expose\_VII\allowbreak{}\_body.tex;\allowbreak{} folio\allowbreak{} 163;\allowbreak{} formula\allowbreak{} 958\allowbreak{}}\par

\par\medskip\noindent\hypertarget{app-E7I-VII-017.2}{}\textbf{E7I-VII-017.2}\quad printed folio 164.\par
\noindent Printed reading: {\ttfamily\footnotesize ;\allowbreak{}}\par
\noindent Corrected reading: {\ttfamily\footnotesize ,\allowbreak{}}\par
\noindent Source: {\ttfamily\footnotesize fr/sga7i/\{diplomatic\allowbreak{},current\}/expose\_VII\allowbreak{}\_body.tex;\allowbreak{} folio\allowbreak{} 164;\allowbreak{} formula\allowbreak{} 973\allowbreak{}}\par

\par\medskip\noindent\hypertarget{app-E7I-VII-017.3}{}\textbf{E7I-VII-017.3}\quad printed folio 164.\par
\noindent Printed reading: {\ttfamily\footnotesize G\allowbreak{}}\par
\noindent Corrected reading: {\ttfamily\footnotesize Q\allowbreak{}}\par
\noindent Source: {\ttfamily\footnotesize fr/sga7i/\{diplomatic\allowbreak{},current\}/expose\_VII\allowbreak{}\_body.tex;\allowbreak{} folio\allowbreak{} 164;\allowbreak{} formula\allowbreak{} 985\allowbreak{}}\par

\par\medskip\noindent\hypertarget{app-E7I-VII-018.1}{}\textbf{E7I-VII-018.1}\quad printed folio 169.\par
\noindent Printed reading: {\ttfamily\footnotesize P'\allowbreak{}}\par
\noindent Corrected reading: {\ttfamily\footnotesize Q\allowbreak{}}\par
\noindent Source: {\ttfamily\footnotesize fr/sga7i/\{diplomatic\allowbreak{},current\}/expose\_VII\allowbreak{}\_body.tex;\allowbreak{} folio\allowbreak{} 169\allowbreak{}}\par

\par\medskip\noindent\hypertarget{app-E7I-VII-019.1}{}\textbf{E7I-VII-019.1}\quad printed folio 173.\par
\noindent Printed reading: {\ttfamily\footnotesize P\allowbreak{}}\par
\noindent Corrected reading: {\ttfamily\footnotesize Q\allowbreak{}}\par
\noindent Source: {\ttfamily\footnotesize fr/sga7i/\{diplomatic\allowbreak{},current\}/expose\_VII\allowbreak{}\_body.tex;\allowbreak{} folio\allowbreak{} 173;\allowbreak{} formula\allowbreak{} 1281\allowbreak{}}\par

\par\medskip\noindent\hypertarget{app-E7I-VII-019.2}{}\textbf{E7I-VII-019.2}\quad printed folio 174.\par
\noindent Printed reading: {\ttfamily\footnotesize preliminary\allowbreak{} conditions\allowbreak{}}\par
\noindent Corrected reading: {\ttfamily\footnotesize hypothesis\allowbreak{}}\par
\noindent Source: {\ttfamily\footnotesize 7I:343;\allowbreak{} WITNESS-F7I-VII-040.\allowbreak{}1;\allowbreak{} WITNESS-F7I-VII-040.\allowbreak{}2;\allowbreak{} fixes/MISPRINTS\_7I.j\allowbreak{}son\#343;\allowbreak{} fr/sga7i/diplomatic/\allowbreak{}expose\_VII\_body.tex;\allowbreak{} fr/sga7i/current/exp\allowbreak{}ose\_VII\_body.tex\allowbreak{}}\par

\par\medskip\noindent\hypertarget{app-E7I-VII-019.3}{}\textbf{E7I-VII-019.3}\quad printed folio 174.\par
\noindent Printed reading: \emph{(empty)}\par
\noindent Corrected reading: {\ttfamily\footnotesize  Proposition\allowbreak{}}\par
\noindent Source: {\ttfamily\footnotesize 7I:343;\allowbreak{} WITNESS-F7I-VII-040.\allowbreak{}1;\allowbreak{} WITNESS-F7I-VII-040.\allowbreak{}2;\allowbreak{} fixes/MISPRINTS\_7I.j\allowbreak{}son\#343;\allowbreak{} fr/sga7i/diplomatic/\allowbreak{}expose\_VII\_body.tex;\allowbreak{} fr/sga7i/current/exp\allowbreak{}ose\_VII\_body.tex\allowbreak{}}\par

\par\medskip\noindent\hypertarget{app-E7I-VII-019.4}{}\textbf{E7I-VII-019.4}\quad printed folio 174.\par
\noindent Printed reading: {\ttfamily\footnotesize 2\allowbreak{}}\par
\noindent Corrected reading: {\ttfamily\footnotesize 3\allowbreak{}}\par
\noindent Source: {\ttfamily\footnotesize 7I:343;\allowbreak{} WITNESS-F7I-VII-040.\allowbreak{}1;\allowbreak{} WITNESS-F7I-VII-040.\allowbreak{}2;\allowbreak{} fixes/MISPRINTS\_7I.j\allowbreak{}son\#343;\allowbreak{} fr/sga7i/diplomatic/\allowbreak{}expose\_VII\_body.tex;\allowbreak{} fr/sga7i/current/exp\allowbreak{}ose\_VII\_body.tex\allowbreak{}}\par

\par\medskip\noindent\hypertarget{app-E7I-VII-019.5}{}\textbf{E7I-VII-019.5}\quad printed folio 174.\par
\noindent Printed reading: {\ttfamily\footnotesize are\allowbreak{}}\par
\noindent Corrected reading: {\ttfamily\footnotesize is\allowbreak{}}\par
\noindent Source: {\ttfamily\footnotesize 7I:343;\allowbreak{} WITNESS-F7I-VII-040.\allowbreak{}1;\allowbreak{} WITNESS-F7I-VII-040.\allowbreak{}2;\allowbreak{} fixes/MISPRINTS\_7I.j\allowbreak{}son\#343;\allowbreak{} fr/sga7i/diplomatic/\allowbreak{}expose\_VII\_body.tex;\allowbreak{} fr/sga7i/current/exp\allowbreak{}ose\_VII\_body.tex\allowbreak{}}\par

\par\medskip\noindent\hypertarget{app-E7I-VII-019.6}{}\textbf{E7I-VII-019.6}\quad printed folio 175.\par
\noindent Printed reading: {\ttfamily\footnotesize 2\allowbreak{}}\par
\noindent Corrected reading: {\ttfamily\footnotesize 3\allowbreak{}}\par
\noindent Source: {\ttfamily\footnotesize 7I:345;\allowbreak{} fixes/MISPRINTS\_7I.j\allowbreak{}son\#345\allowbreak{}}\par

\par\medskip\noindent\hypertarget{app-E7I-VII-019.7}{}\textbf{E7I-VII-019.7}\quad printed folio 175.\par
\noindent Printed reading: {\ttfamily\footnotesize 2\allowbreak{}}\par
\noindent Corrected reading: {\ttfamily\footnotesize 3\allowbreak{}}\par
\noindent Source: {\ttfamily\footnotesize 7I:346;\allowbreak{} fixes/MISPRINTS\_7I.j\allowbreak{}son\#346\allowbreak{}}\par

\par\medskip\noindent\hypertarget{app-E7I-VII-019.8}{}\textbf{E7I-VII-019.8}\quad printed folio 175.\par
\noindent Printed reading: {\ttfamily\footnotesize mS\allowbreak{}}\par
\noindent Corrected reading: {\ttfamily\footnotesize m,S\allowbreak{}}\par
\noindent Source: {\ttfamily\footnotesize fr/sga7i/\{diplomatic\allowbreak{},current\}/expose\_VII\allowbreak{}\_body.tex;\allowbreak{} folio\allowbreak{} 175;\allowbreak{} formula\allowbreak{} 1320\allowbreak{}}\par

\par\medskip\noindent\hypertarget{app-E7I-VII-019.9}{}\textbf{E7I-VII-019.9}\quad printed folio 176.\par
\noindent Printed reading: \emph{(empty)}\par
\noindent Corrected reading: {\ttfamily\footnotesize \textasciicircum{}\{1\}\allowbreak{}}\par
\noindent Source: {\ttfamily\footnotesize ED-EN-FORMULA-G021;\allowbreak{} fr/sga7i/\{diplomatic\allowbreak{},current\}/expose\_VII\allowbreak{}\_body.tex;\allowbreak{} folio\allowbreak{} 176\allowbreak{}}\par

\par\medskip\noindent\hypertarget{app-E7I-VII-020.1}{}\textbf{E7I-VII-020.1}\quad printed folio 177.\par
\noindent Printed reading: \emph{(empty)}\par
\noindent Corrected reading: {\ttfamily\footnotesize =\textbackslash{}underline\allowbreak{} G\_\{m,S\}\allowbreak{}}\par
\noindent Source: {\ttfamily\footnotesize ED-EN-FORMULA-G023;\allowbreak{} fr/sga7i/\{diplomatic\allowbreak{},current\}/expose\_VII\allowbreak{}\_body.tex;\allowbreak{} folio\allowbreak{} 177;\allowbreak{} ED-EN-NONMATH-I-VII1\allowbreak{}77;\allowbreak{} WITNESS-F7I-VII-044.\allowbreak{}1;\allowbreak{} WITNESS-F7I-VII-044.\allowbreak{}2;\allowbreak{} WITNESS-F7I-VII-044.\allowbreak{}3;\allowbreak{} fr/sga7i/diplomatic/\allowbreak{}expose\_VII\_body.tex;\allowbreak{} fr/sga7i/current/exp\allowbreak{}ose\_VII\_body.tex;\allowbreak{} F7I-VII-044.1,\allowbreak{} F7I-VII-044.2,\allowbreak{} F7I-VII-044.3\allowbreak{}}\par

\par\medskip\noindent\hypertarget{app-E7I-VII-020.2}{}\textbf{E7I-VII-020.2}\quad printed folio 177.\par
\noindent Printed reading: {\ttfamily\footnotesize \textbackslash{}underline\allowbreak{} G\_m\allowbreak{}}\par
\noindent Corrected reading: {\ttfamily\footnotesize S\allowbreak{}}\par
\noindent Source: {\ttfamily\footnotesize ED-EN-NONMATH-I-VII1\allowbreak{}77;\allowbreak{} WITNESS-F7I-VII-044.\allowbreak{}1;\allowbreak{} WITNESS-F7I-VII-044.\allowbreak{}2;\allowbreak{} WITNESS-F7I-VII-044.\allowbreak{}3;\allowbreak{} fr/sga7i/diplomatic/\allowbreak{}expose\_VII\_body.tex;\allowbreak{} fr/sga7i/current/exp\allowbreak{}ose\_VII\_body.tex;\allowbreak{} F7I-VII-044.1,\allowbreak{} F7I-VII-044.2,\allowbreak{} F7I-VII-044.3\allowbreak{}}\par

\par\medskip\noindent\hypertarget{app-E7I-VII-020.3}{}\textbf{E7I-VII-020.3}\quad printed folio 177.\par
\noindent Printed reading: {\ttfamily\footnotesize 2\allowbreak{}}\par
\noindent Corrected reading: {\ttfamily\footnotesize 3\allowbreak{}}\par
\noindent Source: {\ttfamily\footnotesize ED-EN-NONMATH-I-VII1\allowbreak{}77;\allowbreak{} WITNESS-F7I-VII-044.\allowbreak{}1;\allowbreak{} WITNESS-F7I-VII-044.\allowbreak{}2;\allowbreak{} WITNESS-F7I-VII-044.\allowbreak{}3;\allowbreak{} fr/sga7i/diplomatic/\allowbreak{}expose\_VII\_body.tex;\allowbreak{} fr/sga7i/current/exp\allowbreak{}ose\_VII\_body.tex;\allowbreak{} F7I-VII-044.1,\allowbreak{} F7I-VII-044.2,\allowbreak{} F7I-VII-044.3\allowbreak{}}\par

\par\medskip\noindent\hypertarget{app-E7I-VII-021.1}{}\textbf{E7I-VII-021.1}\quad printed folio 178.\par
\noindent Printed reading: {\ttfamily\footnotesize P\allowbreak{}}\par
\noindent Corrected reading: {\ttfamily\footnotesize E\allowbreak{}}\par
\noindent Source: {\ttfamily\footnotesize fr/sga7i/\{diplomatic\allowbreak{},current\}/expose\_VII\allowbreak{}\_body.tex;\allowbreak{} folio\allowbreak{} 178;\allowbreak{} formula\allowbreak{} 1412\allowbreak{}}\par

\par\medskip\noindent\hypertarget{app-E7I-VII-021.2}{}\textbf{E7I-VII-021.2}\quad printed folio 178.\par
\noindent Printed reading: {\ttfamily\footnotesize P\allowbreak{}}\par
\noindent Corrected reading: {\ttfamily\footnotesize E\allowbreak{}}\par
\noindent Source: {\ttfamily\footnotesize fr/sga7i/\{diplomatic\allowbreak{},current\}/expose\_VII\allowbreak{}\_body.tex;\allowbreak{} folio\allowbreak{} 178;\allowbreak{} formula\allowbreak{} 1412\allowbreak{}}\par

\par\medskip\noindent\hypertarget{app-E7I-VII-021.3}{}\textbf{E7I-VII-021.3}\quad printed folio 178.\par
\noindent Printed reading: {\ttfamily\footnotesize P\allowbreak{}}\par
\noindent Corrected reading: {\ttfamily\footnotesize E\allowbreak{}}\par
\noindent Source: {\ttfamily\footnotesize fr/sga7i/\{diplomatic\allowbreak{},current\}/expose\_VII\allowbreak{}\_body.tex;\allowbreak{} folio\allowbreak{} 178;\allowbreak{} formula\allowbreak{} 1414\allowbreak{}}\par

\par\medskip\noindent\hypertarget{app-E7I-VII-022.1}{}\textbf{E7I-VII-022.1}\quad printed folio 181.\par
\noindent Printed reading: \emph{(empty)}\par
\noindent Corrected reading: {\ttfamily\footnotesize \{\allowbreak{}}\par
\noindent Source: {\ttfamily\footnotesize port/decisions/decis\allowbreak{}ions\_en7I\_r2\_002.jso\allowbreak{}n\#C756;\allowbreak{} 7I:359\allowbreak{}}\par

\par\medskip\noindent\hypertarget{app-E7I-VII-022.2}{}\textbf{E7I-VII-022.2}\quad printed folio 181.\par
\noindent Printed reading: \emph{(empty)}\par
\noindent Corrected reading: {\ttfamily\footnotesize \_\textbackslash{}bullet\}\allowbreak{}}\par
\noindent Source: {\ttfamily\footnotesize port/decisions/decis\allowbreak{}ions\_en7I\_r2\_002.jso\allowbreak{}n\#C756;\allowbreak{} 7I:359\allowbreak{}}\par

\par\medskip\noindent\hypertarget{app-E7I-VII-023.1}{}\textbf{E7I-VII-023.1}\quad printed folio 182.\par
\noindent Printed reading: {\ttfamily\footnotesize 7.1\allowbreak{}}\par
\noindent Corrected reading: {\ttfamily\footnotesize 8\allowbreak{}}\par
\noindent Source: {\ttfamily\footnotesize 7I:360;\allowbreak{} fixes/MISPRINTS\_7I.j\allowbreak{}son\#360\allowbreak{}}\par

\par\medskip\noindent\hypertarget{app-E7I-VII-024.1}{}\textbf{E7I-VII-024.1}\quad printed folio 185.\par
\noindent Printed reading: {\ttfamily\footnotesize j\allowbreak{}}\par
\noindent Corrected reading: {\ttfamily\footnotesize s\allowbreak{}}\par
\noindent Source: {\ttfamily\footnotesize ED-EN-FORMULA-G025;\allowbreak{} fr/sga7i/\{diplomatic\allowbreak{},current\}/expose\_VII\allowbreak{}\_body.tex;\allowbreak{} folio\allowbreak{} 185\allowbreak{}}\par

\par\medskip\noindent\hypertarget{app-E7I-VII-025.1}{}\textbf{E7I-VII-025.1}\quad printed folio 187.\par
\noindent Printed reading: \emph{(empty)}\par
\noindent Corrected reading: {\ttfamily\footnotesize '\allowbreak{}}\par
\noindent Source: {\ttfamily\footnotesize fr/sga7i/\{diplomatic\allowbreak{},current\}/expose\_VII\allowbreak{}\_body.tex;\allowbreak{} folio\allowbreak{} 187;\allowbreak{} formula\allowbreak{} 1644;\allowbreak{} port/PORT\_EN7I\_LOG.j\allowbreak{}sonl\#C762;\allowbreak{} fixes/FIXLOG\_7I\_CURR\allowbreak{}ENT.jsonl\#C762;\allowbreak{} 7I:368\allowbreak{}}\par

\par\medskip\noindent\hypertarget{app-E7I-VII-025.2}{}\textbf{E7I-VII-025.2}\quad printed folio 187.\par
\noindent Printed reading: \emph{(empty)}\par
\noindent Corrected reading: {\ttfamily\footnotesize 3.\allowbreak{}}\par
\noindent Source: {\ttfamily\footnotesize 7I:369;\allowbreak{} WITNESS-F7I-VII-055.\allowbreak{}2;\allowbreak{} fixes/MISPRINTS\_7I.j\allowbreak{}son\#369;\allowbreak{} fr/sga7i/diplomatic/\allowbreak{}expose\_VII\_body.tex;\allowbreak{} fr/sga7i/current/exp\allowbreak{}ose\_VII\_body.tex\allowbreak{}}\par

\par\medskip\noindent\hypertarget{app-E7I-VII-025.3}{}\textbf{E7I-VII-025.3}\quad printed folio 187.\par
\noindent Printed reading: {\ttfamily\footnotesize 3\allowbreak{}}\par
\noindent Corrected reading: {\ttfamily\footnotesize 1\allowbreak{}}\par
\noindent Source: {\ttfamily\footnotesize 7I:369;\allowbreak{} WITNESS-F7I-VII-055.\allowbreak{}2;\allowbreak{} fixes/MISPRINTS\_7I.j\allowbreak{}son\#369;\allowbreak{} fr/sga7i/diplomatic/\allowbreak{}expose\_VII\_body.tex;\allowbreak{} fr/sga7i/current/exp\allowbreak{}ose\_VII\_body.tex\allowbreak{}}\par

\par\medskip\noindent\hypertarget{app-E7I-VII-025.4}{}\textbf{E7I-VII-025.4}\quad printed folio 187.\par
\noindent Printed reading: \emph{(empty)}\par
\noindent Corrected reading: {\ttfamily\footnotesize 3.\allowbreak{}}\par
\noindent Source: {\ttfamily\footnotesize WITNESS-F7I-VII-055.\allowbreak{}3;\allowbreak{} WITNESS-F7I-VII-055.\allowbreak{}4;\allowbreak{} fr/sga7i/diplomatic/\allowbreak{}expose\_VII\_body.tex;\allowbreak{} fr/sga7i/current/exp\allowbreak{}ose\_VII\_body.tex\allowbreak{}}\par

\par\medskip\noindent\hypertarget{app-E7I-VII-025.5}{}\textbf{E7I-VII-025.5}\quad printed folio 187.\par
\noindent Printed reading: {\ttfamily\footnotesize 3\allowbreak{}}\par
\noindent Corrected reading: {\ttfamily\footnotesize 1\allowbreak{}}\par
\noindent Source: {\ttfamily\footnotesize WITNESS-F7I-VII-055.\allowbreak{}3;\allowbreak{} WITNESS-F7I-VII-055.\allowbreak{}4;\allowbreak{} fr/sga7i/diplomatic/\allowbreak{}expose\_VII\_body.tex;\allowbreak{} fr/sga7i/current/exp\allowbreak{}ose\_VII\_body.tex\allowbreak{}}\par

\par\medskip\noindent\hypertarget{app-E7I-VII-025.6}{}\textbf{E7I-VII-025.6}\quad printed folio 188.\par
\noindent Printed reading: \emph{(empty)}\par
\noindent Corrected reading: {\ttfamily\footnotesize '\allowbreak{}}\par
\noindent Source: {\ttfamily\footnotesize fr/sga7i/\{diplomatic\allowbreak{},current\}/expose\_VII\allowbreak{}\_body.tex;\allowbreak{} folio\allowbreak{} 188;\allowbreak{} formula\allowbreak{} 1686;\allowbreak{} port/PORT\_EN7I\_LOG.j\allowbreak{}sonl\#C763;\allowbreak{} fixes/FIXLOG\_7I\_CURR\allowbreak{}ENT.jsonl\#C763;\allowbreak{} 7I:371\allowbreak{}}\par

\par\medskip\noindent\hypertarget{app-E7I-VII-025.7}{}\textbf{E7I-VII-025.7}\quad printed folio 188.\par
\noindent Printed reading: {\ttfamily\footnotesize E\allowbreak{}}\par
\noindent Corrected reading: {\ttfamily\footnotesize \textbackslash{}bar\allowbreak{} d\_0:L'\_1\textbackslash{}longrightar\allowbreak{}row\allowbreak{} L\allowbreak{}}\par
\noindent Source: {\ttfamily\footnotesize fr/sga7i/\{diplomatic\allowbreak{},current\}/expose\_VII\allowbreak{}\_body.tex;\allowbreak{} folio\allowbreak{} 188\allowbreak{}}\par

\par\medskip\noindent\hypertarget{app-E7I-VII-025.8}{}\textbf{E7I-VII-025.8}\quad printed folio 188.\par
\noindent Printed reading: {\ttfamily\footnotesize denote\allowbreak{}}\par
\noindent Corrected reading: {\ttfamily\footnotesize be\allowbreak{}}\par
\noindent Source: {\ttfamily\footnotesize fr/sga7i/\{diplomatic\allowbreak{},current\}/expose\_VII\allowbreak{}\_body.tex;\allowbreak{} folio\allowbreak{} 188\allowbreak{}}\par

\par\medskip\noindent\hypertarget{app-E7I-VII-025.9}{}\textbf{E7I-VII-025.9}\quad printed folio 188.\par
\noindent Printed reading: {\ttfamily\footnotesize object\allowbreak{}}\par
\noindent Corrected reading: {\ttfamily\footnotesize morphism\allowbreak{} induced\allowbreak{} by\allowbreak{}}\par
\noindent Source: {\ttfamily\footnotesize fr/sga7i/\{diplomatic\allowbreak{},current\}/expose\_VII\allowbreak{}\_body.tex;\allowbreak{} folio\allowbreak{} 188\allowbreak{}}\par

\par\medskip\noindent\hypertarget{app-E7I-VII-025.10}{}\textbf{E7I-VII-025.10}\quad printed folio 188.\par
\noindent Printed reading: \emph{(empty)}\par
\noindent Corrected reading: {\ttfamily\footnotesize d\_0:L\_1\textbackslash{}longrightarr\allowbreak{}ow\allowbreak{} L\_0\$.\allowbreak{} Let\allowbreak{} \$E\_0=\allowbreak{}}\par
\noindent Source: {\ttfamily\footnotesize fr/sga7i/\{diplomatic\allowbreak{},current\}/expose\_VII\allowbreak{}\_body.tex;\allowbreak{} folio\allowbreak{} 188\allowbreak{}}\par

\par\medskip\noindent\hypertarget{app-E7I-VII-025.11}{}\textbf{E7I-VII-025.11}\quad printed folio 188.\par
\noindent Printed reading: \emph{(empty)}\par
\noindent Corrected reading: {\ttfamily\footnotesize be\allowbreak{} the\allowbreak{} object\allowbreak{} }\par
\noindent Source: {\ttfamily\footnotesize fr/sga7i/\{diplomatic\allowbreak{},current\}/expose\_VII\allowbreak{}\_body.tex;\allowbreak{} folio\allowbreak{} 188\allowbreak{}}\par

\par\medskip\noindent\hypertarget{app-E7I-VII-025.12}{}\textbf{E7I-VII-025.12}\quad printed folio 188.\par
\noindent Printed reading: {\ttfamily\footnotesize ,\allowbreak{}}\par
\noindent Corrected reading: \emph{(empty)}\par
\noindent Source: {\ttfamily\footnotesize fr/sga7i/\{diplomatic\allowbreak{},current\}/expose\_VII\allowbreak{}\_body.tex;\allowbreak{} folio\allowbreak{} 188\allowbreak{}}\par

\par\medskip\noindent\hypertarget{app-E7I-VII-025.13}{}\textbf{E7I-VII-025.13}\quad printed folio 188.\par
\noindent Printed reading: {\ttfamily\footnotesize where\allowbreak{}}\par
\noindent Corrected reading: {\ttfamily\footnotesize for\allowbreak{}}\par
\noindent Source: {\ttfamily\footnotesize fr/sga7i/\{diplomatic\allowbreak{},current\}/expose\_VII\allowbreak{}\_body.tex;\allowbreak{} folio\allowbreak{} 188\allowbreak{}}\par

\par\medskip\noindent\hypertarget{app-E7I-VII-025.14}{}\textbf{E7I-VII-025.14}\quad printed folio 188.\par
\noindent Printed reading: \emph{(empty)}\par
\noindent Corrected reading: {\ttfamily\footnotesize which\allowbreak{} }\par
\noindent Source: {\ttfamily\footnotesize fr/sga7i/\{diplomatic\allowbreak{},current\}/expose\_VII\allowbreak{}\_body.tex;\allowbreak{} folio\allowbreak{} 188\allowbreak{}}\par

\par\medskip\noindent\hypertarget{app-E7I-VII-025.15}{}\textbf{E7I-VII-025.15}\quad printed folio 188.\par
\noindent Printed reading: {\ttfamily\footnotesize ,\allowbreak{} }\par
\noindent Corrected reading: \emph{(0 espace(s), 1 retour(s) de ligne dans le témoin TeX)}\par
\noindent Source: {\ttfamily\footnotesize fr/sga7i/\{diplomatic\allowbreak{},current\}/expose\_VII\allowbreak{}\_body.tex;\allowbreak{} folio\allowbreak{} 188\allowbreak{}}\par

\par\medskip\noindent\hypertarget{app-E7I-VII-025.16}{}\textbf{E7I-VII-025.16}\quad printed folio 188.\par
\noindent Printed reading: {\ttfamily\footnotesize  where\allowbreak{}}\par
\noindent Corrected reading: \emph{(empty)}\par
\noindent Source: {\ttfamily\footnotesize fr/sga7i/\{diplomatic\allowbreak{},current\}/expose\_VII\allowbreak{}\_body.tex;\allowbreak{} folio\allowbreak{} 188\allowbreak{}}\par

\par\medskip\noindent\hypertarget{app-E7I-VII-025.17}{}\textbf{E7I-VII-025.17}\quad printed folio 188.\par
\noindent Printed reading: {\ttfamily\footnotesize has\allowbreak{}}\par
\noindent Corrected reading: {\ttfamily\footnotesize is\allowbreak{}}\par
\noindent Source: {\ttfamily\footnotesize fr/sga7i/\{diplomatic\allowbreak{},current\}/expose\_VII\allowbreak{}\_body.tex;\allowbreak{} folio\allowbreak{} 188\allowbreak{}}\par

\par\medskip\noindent\hypertarget{app-E7I-VII-025.18}{}\textbf{E7I-VII-025.18}\quad printed folio 188.\par
\noindent Printed reading: {\ttfamily\footnotesize components\allowbreak{}}\par
\noindent Corrected reading: {\ttfamily\footnotesize the\allowbreak{} composite\allowbreak{}}\par
\noindent Source: {\ttfamily\footnotesize fr/sga7i/\{diplomatic\allowbreak{},current\}/expose\_VII\allowbreak{}\_body.tex;\allowbreak{} folio\allowbreak{} 188\allowbreak{}}\par

\par\medskip\noindent\hypertarget{app-E7I-VII-025.19}{}\textbf{E7I-VII-025.19}\quad printed folio 188.\par
\noindent Printed reading: {\ttfamily\footnotesize \$(d\allowbreak{}}\par
\noindent Corrected reading: {\ttfamily\footnotesize \textbackslash{}[\allowbreak{} L\_1\textbackslash{}twoheadrightarro\allowbreak{}w\allowbreak{} L\allowbreak{}}\par
\noindent Source: {\ttfamily\footnotesize fr/sga7i/\{diplomatic\allowbreak{},current\}/expose\_VII\allowbreak{}\_body.tex;\allowbreak{} folio\allowbreak{} 188\allowbreak{}}\par

\par\medskip\noindent\hypertarget{app-E7I-VII-025.20}{}\textbf{E7I-VII-025.20}\quad printed folio 188.\par
\noindent Printed reading: \emph{(empty)}\par
\noindent Corrected reading: {\ttfamily\footnotesize  \textbackslash{}xrightarrow\{(\textbackslash{}bar\allowbreak{} d\_0\allowbreak{}}\par
\noindent Source: {\ttfamily\footnotesize fr/sga7i/\{diplomatic\allowbreak{},current\}/expose\_VII\allowbreak{}\_body.tex;\allowbreak{} folio\allowbreak{} 188\allowbreak{}}\par

\par\medskip\noindent\hypertarget{app-E7I-VII-025.21}{}\textbf{E7I-VII-025.21}\quad printed folio 188.\par
\noindent Printed reading: \emph{(empty)}\par
\noindent Corrected reading: {\ttfamily\footnotesize \textbackslash{},\allowbreak{}}\par
\noindent Source: {\ttfamily\footnotesize fr/sga7i/\{diplomatic\allowbreak{},current\}/expose\_VII\allowbreak{}\_body.tex;\allowbreak{} folio\allowbreak{} 188\allowbreak{}}\par

\par\medskip\noindent\hypertarget{app-E7I-VII-025.22}{}\textbf{E7I-VII-025.22}\quad printed folio 188.\par
\noindent Printed reading: {\ttfamily\footnotesize \$,\allowbreak{}}\par
\noindent Corrected reading: {\ttfamily\footnotesize \}\allowbreak{} L\_0\textbackslash{}oplus\allowbreak{}}\par
\noindent Source: {\ttfamily\footnotesize fr/sga7i/\{diplomatic\allowbreak{},current\}/expose\_VII\allowbreak{}\_body.tex;\allowbreak{} folio\allowbreak{} 188\allowbreak{}}\par

\par\medskip\noindent\hypertarget{app-E7I-VII-025.23}{}\textbf{E7I-VII-025.23}\quad printed folio 188.\par
\noindent Printed reading: {\ttfamily\footnotesize with\allowbreak{} \$d\allowbreak{}}\par
\noindent Corrected reading: {\ttfamily\footnotesize L\allowbreak{}}\par
\noindent Source: {\ttfamily\footnotesize fr/sga7i/\{diplomatic\allowbreak{},current\}/expose\_VII\allowbreak{}\_body.tex;\allowbreak{} folio\allowbreak{} 188\allowbreak{}}\par

\par\medskip\noindent\hypertarget{app-E7I-VII-025.24}{}\textbf{E7I-VII-025.24}\quad printed folio 188.\par
\noindent Printed reading: {\ttfamily\footnotesize :L\_1\textbackslash{}longrightarrow\allowbreak{} L\_0\$\allowbreak{} induced\allowbreak{}}\par
\noindent Corrected reading: {\ttfamily\footnotesize .\allowbreak{}}\par
\noindent Source: {\ttfamily\footnotesize fr/sga7i/\{diplomatic\allowbreak{},current\}/expose\_VII\allowbreak{}\_body.tex;\allowbreak{} folio\allowbreak{} 188\allowbreak{}}\par

\par\medskip\noindent\hypertarget{app-E7I-VII-025.25}{}\textbf{E7I-VII-025.25}\quad printed folio 188.\par
\noindent Printed reading: {\ttfamily\footnotesize by\allowbreak{} \$d\_1:L\_1\textbackslash{}longrightar\allowbreak{}row\allowbreak{} L\_0\$\allowbreak{} by\allowbreak{} passage\allowbreak{} to\allowbreak{} the\allowbreak{} quotient.\allowbreak{}}\par
\noindent Corrected reading: {\ttfamily\footnotesize \textbackslash{}]\allowbreak{}}\par
\noindent Source: {\ttfamily\footnotesize fr/sga7i/\{diplomatic\allowbreak{},current\}/expose\_VII\allowbreak{}\_body.tex;\allowbreak{} folio\allowbreak{} 188\allowbreak{}}\par

\par\medskip\noindent\hypertarget{app-E7I-VII-026.1}{}\textbf{E7I-VII-026.1}\quad printed folio 189.\par
\noindent Printed reading: {\ttfamily\footnotesize j\allowbreak{}}\par
\noindent Corrected reading: {\ttfamily\footnotesize i\allowbreak{}}\par
\noindent Source: {\ttfamily\footnotesize fr/sga7i/\{diplomatic\allowbreak{},current\}/expose\_VII\allowbreak{}\_body.tex;\allowbreak{} folio\allowbreak{} 189;\allowbreak{} formula\allowbreak{} 1744\allowbreak{}}\par

\par\medskip\noindent\hypertarget{app-E7I-VII-027.1}{}\textbf{E7I-VII-027.1}\quad printed folio 190.\par
\noindent Printed reading: {\ttfamily\footnotesize )\allowbreak{}}\par
\noindent Corrected reading: \emph{(empty)}\par
\noindent Source: {\ttfamily\footnotesize port/decisions/decis\allowbreak{}ions\_en7I\_r2\_002.jso\allowbreak{}n\#C766;\allowbreak{} fixes/FIXLOG\_7I\_CURR\allowbreak{}ENT.jsonl\#C766;\allowbreak{} 7I:379\allowbreak{}}\par

\par\medskip\noindent\hypertarget{app-E7I-VII-027.2}{}\textbf{E7I-VII-027.2}\quad printed folio 190.\par
\noindent Printed reading: {\ttfamily\footnotesize o\allowbreak{}}\par
\noindent Corrected reading: {\ttfamily\footnotesize \textbackslash{}partial\allowbreak{}}\par
\noindent Source: {\ttfamily\footnotesize fr/sga7i/\{diplomatic\allowbreak{},current\}/expose\_VII\allowbreak{}\_body.tex;\allowbreak{} folio\allowbreak{} 190\allowbreak{}}\par

\par\medskip\noindent\hypertarget{app-E7I-VII-028.1}{}\textbf{E7I-VII-028.1}\quad printed folio 192.\par
\noindent Printed reading: {\ttfamily\footnotesize  the\allowbreak{} set\allowbreak{} of\allowbreak{} isomorphisms\allowbreak{}}\par
\noindent Corrected reading: \emph{(empty)}\par
\noindent Source: {\ttfamily\footnotesize 7I:384;\allowbreak{} fixes/FIXLOG\_7I\_CURR\allowbreak{}ENT.jsonl\#770\allowbreak{}}\par

\par\medskip\noindent\hypertarget{app-E7I-VII-029.1}{}\textbf{E7I-VII-029.1}\quad printed folio 192.\par
\noindent Printed reading: \emph{(empty)}\par
\noindent Corrected reading: {\ttfamily\footnotesize \{\allowbreak{}}\par
\noindent Source: {\ttfamily\footnotesize ED-EN-FORMULA-G027;\allowbreak{} fr/sga7i/\{diplomatic\allowbreak{},current\}/expose\_VII\allowbreak{}\_body.tex;\allowbreak{} folio\allowbreak{} 192;\allowbreak{} port/decisions/decis\allowbreak{}ions\_en7I\_r2\_002.jso\allowbreak{}n\#C770;\allowbreak{} 7I:385\allowbreak{}}\par

\par\medskip\noindent\hypertarget{app-E7I-VII-029.2}{}\textbf{E7I-VII-029.2}\quad printed folio 192.\par
\noindent Printed reading: \emph{(empty)}\par
\noindent Corrected reading: {\ttfamily\footnotesize \_\textbackslash{}bullet\}\allowbreak{}}\par
\noindent Source: {\ttfamily\footnotesize ED-EN-FORMULA-G027;\allowbreak{} fr/sga7i/\{diplomatic\allowbreak{},current\}/expose\_VII\allowbreak{}\_body.tex;\allowbreak{} folio\allowbreak{} 192;\allowbreak{} port/decisions/decis\allowbreak{}ions\_en7I\_r2\_002.jso\allowbreak{}n\#C770;\allowbreak{} 7I:385\allowbreak{}}\par

\par\medskip\noindent\hypertarget{app-E7I-VII-029.3}{}\textbf{E7I-VII-029.3}\quad printed folio 193.\par
\noindent Printed reading: {\ttfamily\footnotesize \textbackslash{}qquad\allowbreak{}}\par
\noindent Corrected reading: {\ttfamily\footnotesize \textbackslash{}longrightarrow\allowbreak{}}\par
\noindent Source: {\ttfamily\footnotesize fr/sga7i/\{diplomatic\allowbreak{},current\}/expose\_VII\allowbreak{}\_body.tex;\allowbreak{} folio\allowbreak{} 193\allowbreak{}}\par

\par\medskip\noindent\hypertarget{app-E7I-VII-030.1}{}\textbf{E7I-VII-030.1}\quad printed folio 194.\par
\noindent Printed reading: \emph{(empty)}\par
\noindent Corrected reading: {\ttfamily\footnotesize \textasciicircum{}1\allowbreak{}}\par
\noindent Source: {\ttfamily\footnotesize fr/sga7i/\{diplomatic\allowbreak{},current\}/expose\_VII\allowbreak{}\_body.tex;\allowbreak{} folio\allowbreak{} 194;\allowbreak{} formula\allowbreak{} 1850\allowbreak{}}\par

\par\medskip\noindent\hypertarget{app-E7I-VII-030.2}{}\textbf{E7I-VII-030.2}\quad printed folio 195.\par
\noindent Printed reading: {\ttfamily\footnotesize '\allowbreak{}}\par
\noindent Corrected reading: \emph{(empty)}\par
\noindent Source: {\ttfamily\footnotesize fr/sga7i/\{diplomatic\allowbreak{},current\}/expose\_VII\allowbreak{}\_body.tex;\allowbreak{} folio\allowbreak{} 195;\allowbreak{} formula\allowbreak{} 1871\allowbreak{}}\par

\par\medskip\noindent\hypertarget{app-E7I-VII-030.3}{}\textbf{E7I-VII-030.3}\quad printed folio 195.\par
\noindent Printed reading: \emph{(empty)}\par
\noindent Corrected reading: {\ttfamily\footnotesize '\allowbreak{}}\par
\noindent Source: {\ttfamily\footnotesize fr/sga7i/\{diplomatic\allowbreak{},current\}/expose\_VII\allowbreak{}\_body.tex;\allowbreak{} folio\allowbreak{} 195;\allowbreak{} formula\allowbreak{} 1885\allowbreak{}}\par

\par\medskip\noindent\hypertarget{app-E7I-VII-030.4}{}\textbf{E7I-VII-030.4}\quad printed folio 195.\par
\noindent Printed reading: {\ttfamily\footnotesize jt\allowbreak{}}\par
\noindent Corrected reading: {\ttfamily\footnotesize j't\allowbreak{}}\par
\noindent Source: {\ttfamily\footnotesize ED-EN-FORMULA-G029;\allowbreak{} fr/sga7i/\{diplomatic\allowbreak{},current\}/expose\_VII\allowbreak{}\_body.tex;\allowbreak{} folio\allowbreak{} 195\allowbreak{}}\par

\par\medskip\noindent\hypertarget{app-E7I-VII-030.5}{}\textbf{E7I-VII-030.5}\quad printed folio 196.\par
\noindent Printed reading: {\ttfamily\footnotesize jt\allowbreak{}}\par
\noindent Corrected reading: {\ttfamily\footnotesize j't\allowbreak{}}\par
\noindent Source: {\ttfamily\footnotesize fr/sga7i/\{diplomatic\allowbreak{},current\}/expose\_VII\allowbreak{}\_body.tex;\allowbreak{} folio\allowbreak{} 196\allowbreak{}}\par

\par\medskip\noindent\hypertarget{app-E7I-VII-030.6}{}\textbf{E7I-VII-030.6}\quad printed folio 196.\par
\noindent Printed reading: \emph{(empty)}\par
\noindent Corrected reading: {\ttfamily\footnotesize '\allowbreak{}}\par
\noindent Source: {\ttfamily\footnotesize fr/sga7i/\{diplomatic\allowbreak{},current\}/expose\_VII\allowbreak{}\_body.tex;\allowbreak{} folio\allowbreak{} 196;\allowbreak{} formula\allowbreak{} 1902\allowbreak{}}\par

\par\medskip\noindent\hypertarget{app-E7I-VII-030.7}{}\textbf{E7I-VII-030.7}\quad printed folio 196.\par
\noindent Printed reading: \emph{(empty)}\par
\noindent Corrected reading: {\ttfamily\footnotesize ''\allowbreak{}}\par
\noindent Source: {\ttfamily\footnotesize ED-EN-FORMULA-G031;\allowbreak{} fr/sga7i/\{diplomatic\allowbreak{},current\}/expose\_VII\allowbreak{}\_body.tex;\allowbreak{} folio\allowbreak{} 196\allowbreak{}}\par

\par\medskip\noindent\hypertarget{app-E7I-VII-030.8}{}\textbf{E7I-VII-030.8}\quad printed folio 196.\par
\noindent Printed reading: \emph{(1 espace(s), 0 retour(s) de ligne dans le témoin TeX)}\par
\noindent Corrected reading: {\ttfamily\footnotesize ,\allowbreak{}}\par
\noindent Source: {\ttfamily\footnotesize fr/sga7i/\{diplomatic\allowbreak{},current\}/expose\_VII\allowbreak{}\_body.tex;\allowbreak{} folio\allowbreak{} 196\allowbreak{}}\par

\par\medskip\noindent\hypertarget{app-E7I-VII-030.9}{}\textbf{E7I-VII-030.9}\quad printed folio 197.\par
\noindent Printed reading: \emph{(empty)}\par
\noindent Corrected reading: {\ttfamily\footnotesize \_\textbackslash{}bullet\allowbreak{}}\par
\noindent Source: {\ttfamily\footnotesize ED-EN-FORMULA-G033;\allowbreak{} fr/sga7i/\{diplomatic\allowbreak{},current\}/expose\_VII\allowbreak{}\_body.tex;\allowbreak{} folio\allowbreak{} 197;\allowbreak{} port/decisions/decis\allowbreak{}ions\_en7I\_r2\_002.jso\allowbreak{}n\#C772;\allowbreak{} 7I:399\allowbreak{}}\par

\par\medskip\noindent\hypertarget{app-E7I-VII-031.1}{}\textbf{E7I-VII-031.1}\quad printed folio 199.\par
\noindent Printed reading: {\ttfamily\footnotesize 1d\allowbreak{}}\par
\noindent Corrected reading: {\ttfamily\footnotesize 0d\allowbreak{}}\par
\noindent Source: {\ttfamily\footnotesize fr/sga7i/\{diplomatic\allowbreak{},current\}/expose\_VII\allowbreak{}\_body.tex;\allowbreak{} folio\allowbreak{} 199\allowbreak{}}\par

\par\medskip\noindent\hypertarget{app-E7I-VII-031.2}{}\textbf{E7I-VII-031.2}\quad printed folio 199.\par
\noindent Printed reading: {\ttfamily\footnotesize 0\allowbreak{}}\par
\noindent Corrected reading: {\ttfamily\footnotesize 1\allowbreak{}}\par
\noindent Source: {\ttfamily\footnotesize fr/sga7i/\{diplomatic\allowbreak{},current\}/expose\_VII\allowbreak{}\_body.tex;\allowbreak{} folio\allowbreak{} 199\allowbreak{}}\par

\par\medskip\noindent\hypertarget{app-E7I-VII-032.1}{}\textbf{E7I-VII-032.1}\quad printed folio 201.\par
\noindent Printed reading: {\ttfamily\footnotesize 1d\allowbreak{}}\par
\noindent Corrected reading: {\ttfamily\footnotesize 0d\allowbreak{}}\par
\noindent Source: {\ttfamily\footnotesize fr/sga7i/\{diplomatic\allowbreak{},current\}/expose\_VII\allowbreak{}\_body.tex;\allowbreak{} folio\allowbreak{} 201;\allowbreak{} formula\allowbreak{} 2067\allowbreak{}}\par

\par\medskip\noindent\hypertarget{app-E7I-VII-032.2}{}\textbf{E7I-VII-032.2}\quad printed folio 201.\par
\noindent Printed reading: {\ttfamily\footnotesize 0\allowbreak{}}\par
\noindent Corrected reading: {\ttfamily\footnotesize 1\allowbreak{}}\par
\noindent Source: {\ttfamily\footnotesize fr/sga7i/\{diplomatic\allowbreak{},current\}/expose\_VII\allowbreak{}\_body.tex;\allowbreak{} folio\allowbreak{} 201;\allowbreak{} formula\allowbreak{} 2067\allowbreak{}}\par

\par\medskip\noindent\hypertarget{app-E7I-VII-033.1}{}\textbf{E7I-VII-033.1}\quad printed folio 204.\par
\noindent Printed reading: \emph{(empty)}\par
\noindent Corrected reading: {\ttfamily\footnotesize '\allowbreak{}}\par
\noindent Source: {\ttfamily\footnotesize ED-EN-FORMULA-G035;\allowbreak{} fr/sga7i/\{diplomatic\allowbreak{},current\}/expose\_VII\allowbreak{}\_body.tex;\allowbreak{} folio\allowbreak{} 204\allowbreak{}}\par

\par\medskip\noindent\hypertarget{app-E7I-VII-033.2}{}\textbf{E7I-VII-033.2}\quad printed folio 204.\par
\noindent Printed reading: {\ttfamily\footnotesize '\allowbreak{}}\par
\noindent Corrected reading: \emph{(empty)}\par
\noindent Source: {\ttfamily\footnotesize ED-EN-FORMULA-G035;\allowbreak{} fr/sga7i/\{diplomatic\allowbreak{},current\}/expose\_VII\allowbreak{}\_body.tex;\allowbreak{} folio\allowbreak{} 204\allowbreak{}}\par

\par\medskip\noindent\hypertarget{app-E7I-VII-033.3}{}\textbf{E7I-VII-033.3}\quad printed folio 204.\par
\noindent Printed reading: {\ttfamily\footnotesize '\allowbreak{}}\par
\noindent Corrected reading: \emph{(empty)}\par
\noindent Source: {\ttfamily\footnotesize ED-EN-FORMULA-G035;\allowbreak{} fr/sga7i/\{diplomatic\allowbreak{},current\}/expose\_VII\allowbreak{}\_body.tex;\allowbreak{} folio\allowbreak{} 204\allowbreak{}}\par

\par\medskip\noindent\hypertarget{app-E7I-VII-034.1}{}\textbf{E7I-VII-034.1}\quad printed folio 205.\par
\noindent Printed reading: {\ttfamily\footnotesize \textbackslash{}times\allowbreak{} }\par
\noindent Corrected reading: {\ttfamily\footnotesize ,\allowbreak{}}\par
\noindent Source: {\ttfamily\footnotesize fr/sga7i/\{diplomatic\allowbreak{},current\}/expose\_VII\allowbreak{}\_body.tex;\allowbreak{} folio\allowbreak{} 205;\allowbreak{} formula\allowbreak{} 2158\allowbreak{}}\par

\par\medskip\noindent\hypertarget{app-E7I-VII-034.2}{}\textbf{E7I-VII-034.2}\quad printed folio 206.\par
\noindent Printed reading: {\ttfamily\footnotesize 1q\allowbreak{}}\par
\noindent Corrected reading: {\ttfamily\footnotesize 1,q\allowbreak{}}\par
\noindent Source: {\ttfamily\footnotesize fr/sga7i/\{diplomatic\allowbreak{},current\}/expose\_VII\allowbreak{}\_body.tex;\allowbreak{} folio\allowbreak{} 206;\allowbreak{} formula\allowbreak{} 2187\allowbreak{}}\par

\par\medskip\noindent\hypertarget{app-E7I-VII-034.3}{}\textbf{E7I-VII-034.3}\quad printed folio 206.\par
\noindent Printed reading: \emph{(2 espace(s), 0 retour(s) de ligne dans le témoin TeX)}\par
\noindent Corrected reading: {\ttfamily\footnotesize  P\textbackslash{}times\allowbreak{} }\par
\noindent Source: {\ttfamily\footnotesize fr/sga7i/\{diplomatic\allowbreak{},current\}/expose\_VII\allowbreak{}\_body.tex;\allowbreak{} folio\allowbreak{} 206;\allowbreak{} formula\allowbreak{} 2188\allowbreak{}}\par

\par\medskip\noindent\hypertarget{app-E7I-VII-034.4}{}\textbf{E7I-VII-034.4}\quad printed folio 206.\par
\noindent Printed reading: {\ttfamily\footnotesize P\allowbreak{}}\par
\noindent Corrected reading: {\ttfamily\footnotesize Q\allowbreak{}}\par
\noindent Source: {\ttfamily\footnotesize fr/sga7i/\{diplomatic\allowbreak{},current\}/expose\_VII\allowbreak{}\_body.tex;\allowbreak{} folio\allowbreak{} 206;\allowbreak{} formula\allowbreak{} 2211\allowbreak{}}\par

\par\medskip\noindent\hypertarget{app-E7I-VII-035.1}{}\textbf{E7I-VII-035.1}\quad printed folio 212.\par
\noindent Printed reading: {\ttfamily\footnotesize \textbackslash{}arrow[dr,dash,"s"]\allowbreak{}}\par
\noindent Corrected reading: \emph{(empty)}\par
\noindent Source: {\ttfamily\footnotesize fr/sga7i/\{diplomatic\allowbreak{},current\}/expose\_VII\allowbreak{}\_body.tex;\allowbreak{} folio\allowbreak{} 212\allowbreak{}}\par

\par\medskip\noindent\hypertarget{app-E7I-VII-035.2}{}\textbf{E7I-VII-035.2}\quad printed folio 212.\par
\noindent Printed reading: \emph{(empty)}\par
\noindent Corrected reading: {\ttfamily\footnotesize ]\textbackslash{}arrow[ul,"s"\allowbreak{}}\par
\noindent Source: {\ttfamily\footnotesize fr/sga7i/\{diplomatic\allowbreak{},current\}/expose\_VII\allowbreak{}\_body.tex;\allowbreak{} folio\allowbreak{} 212\allowbreak{}}\par

\par\medskip\noindent\hypertarget{app-E7I-VII-036.1}{}\textbf{E7I-VII-036.1}\quad printed folio 217.\par
\noindent Printed reading: {\ttfamily\footnotesize 1956\allowbreak{}}\par
\noindent Corrected reading: {\ttfamily\footnotesize 1957\allowbreak{}}\par
\noindent Source: {\ttfamily\footnotesize port/PORT\_EN7I\_LOG.j\allowbreak{}sonl\#C779;\allowbreak{} fixes/FIXLOG\_7I\_CURR\allowbreak{}ENT.jsonl\#C779;\allowbreak{} 7I:418\allowbreak{}}\par

\par\medskip\noindent\hypertarget{app-E7I-VII-036.2}{}\textbf{E7I-VII-036.2}\quad printed folio 217.\par
\noindent Printed reading: {\ttfamily\footnotesize Lectures\allowbreak{}}\par
\noindent Corrected reading: {\ttfamily\footnotesize Lecture\allowbreak{}}\par
\noindent Source: {\ttfamily\footnotesize 7I:419;\allowbreak{} WITNESS-F7I-VII-078.\allowbreak{}3;\allowbreak{} fixes/MISPRINTS\_7I.j\allowbreak{}son\#419;\allowbreak{} fr/sga7i/diplomatic/\allowbreak{}expose\_VII\_body.tex;\allowbreak{} fr/sga7i/current/exp\allowbreak{}ose\_VII\_body.tex\allowbreak{}}\par

\par\medskip\noindent\hypertarget{app-E7I-VII-036.3}{}\textbf{E7I-VII-036.3}\quad printed folio 217.\par
\noindent Printed reading: {\ttfamily\footnotesize Math.\allowbreak{}}\par
\noindent Corrected reading: {\ttfamily\footnotesize Mathematics\allowbreak{}}\par
\noindent Source: {\ttfamily\footnotesize 7I:419;\allowbreak{} WITNESS-F7I-VII-078.\allowbreak{}3;\allowbreak{} fixes/MISPRINTS\_7I.j\allowbreak{}son\#419;\allowbreak{} fr/sga7i/diplomatic/\allowbreak{}expose\_VII\_body.tex;\allowbreak{} fr/sga7i/current/exp\allowbreak{}ose\_VII\_body.tex\allowbreak{}}\par

\par\endgroup\normalsize\fi
\clearpage\SGAtoc{Exposé VIII}
\selectlanguage{english}
\SGApage{I}{218}

\begin{center}
\SGARefAliasTarget{sga7.sga7i.eviii.nroot}SGA 7

\vspace{1em}
\underline{Expos\'e VIII}

\vspace{1em}
\underline{\underline{COMPLEMENTS ON BIEXTENSIONS. GENERAL PROPERTIES}}

\underline{\underline{OF BIEXTENSIONS OF GROUP SCHEMES}}

\vspace{1em}
\underline{by A. Grothendieck}
\end{center}

\subsection*{Contents}

\begin{enumerate}[label=\arabic*.]
\setcounter{enumi}{-1}
\item Introduction
\item Various special cases over an arbitrary topos
\item Pairings defined by a biextension
\item Biextensions of group schemes $(P,Q)$ by $\underline G_m$: generalities
\item Extensions and biextensions of smooth connected group schemes over a field
\item Extensions and biextensions by constant truncated torsion-free group schemes
\item Extensions and biextensions by the N\'eron model of $\underline G_m$
\item Canonical prolongations of extensions and biextensions by $\underline G_m$
\end{enumerate}

0. \underline{Introduction}

\SGARefTarget{sga7.sga7i.eviii.n0-1}{0.1}. We continue here the general study of biextensions begun in the
preceding expos\'e.  First we give some complements in the case of an
arbitrary base topos, developing in particular, with the detail it
deserves, the formalism of the pairings associated with biextensions
(Section 2), whose importance in the theory of abelian schemes is well
known.  We then begin the study of biextensions of general group
schemes, the case of greatest interest for us being that of
biextensions by the multiplicative group $\underline G_m$.  We shall
nevertheless be obliged, as a technical intermediate step, also to
study biextensions by other types of group schemes (Sections 5 and 6).

\SGApage{I}{\(\approx\)219}
The most interesting result for what follows is given in \SGARefLink{sga7.sga7i.eviii.n7-1}{7.1}, which
gives conditions under which, for $P$ and $Q$ smooth over a trait with
generic point $\eta$, a biextension of $(P_\eta,Q_\eta)$ by
$\underline G_{m\eta}$ can be prolonged to a biextension of $(P,Q)$ by
$\underline G_m$, determined up to a unique isomorphism.  This result
will be very convenient in the following expos\'e, devoted to the
theory of the N\'eron model of abelian varieties.  It is likely that
one could manage there at less expense, without the imposing
paperweight of Expos\'es VII and VIII (which will no longer be used in
the remainder of the Seminar); it nevertheless seems likely to us that
the developments given in these expos\'es are destined to be useful
elsewhere as well.

\SGARefTarget{sga7.sga7i.eviii.n0-2}{0.2}. When we work with group schemes over a scheme $S$, it will be
understood that we work with the fppf topos of $S$.  This is the topos
associated with the $\underline U$-site of schemes locally of finite
presentation over $S$ (and belonging to $\underline U$), endowed with
the fppf topology (SGA 3 IV 6.3).  The reader will readily verify that
most of the statements we shall give would also be valid (and
essentially equivalent to the fppf statements) for other topologies on
$S$, whose underlying site contains at least the schemes locally of
finite presentation over $S$, such as the Zariski, \'etale, and fpqc
topologies.  One reason is that, by descent theory, the category of
extensions of a group scheme $P$ locally of finite presentation over
$S$ by the group scheme $\underline G_m$ does not depend, up to
equivalence, on the chosen topology; the same statement holds for the
notion of biextension.

The technical reason for preferring the fppf topology (or a finer
topology, such as fpqc) to coarser topologies such as the Zariski and
\'etale topologies is that, for every integer $n>0$, the homomorphism
$n\cdot\id$ from $\underline G_m$ to itself is an \underline{fppf
epimorphism} (because it is faithfully flat and of finite
presentation), but in general is not a Zariski or an \'etale
epimorphism.  Notice that it is an \'etale epimorphism if and only if
$n$ is prime to the residue characteristics of $S$, and a Zariski
epimorphism if and only if $S$ is empty.

\SGApage{I}{220}
\SGARefTarget{sga7.sga7i.eviii.n0-3}{0.3}. Finally, most of the statements concerning extensions or
biextensions by $\underline G_m$ would remain valid if that group were
replaced by a \underline{torus} $G$.  Since such a $G$ is locally
isomorphic (for the fppf, and even for the \'etale topology) to a group
$\underline G_m^r$, this is an essentially trivial generalization.



\subsection*{\SGARefTarget{sga7.sga7i.eviii.n1}{1}. Various special cases over an arbitrary topos}

\SGARefTarget{sga7.sga7i.eviii.n1-1}{1.1}. In this section and the next, as in Expos\'e VII, $T$ denotes a
topos, and $P,Q,G$ denote three abelian groups of $T$.  Thus there is a
canonical isomorphism (VII \SGARefLink{sga7.sga7i.evii.n3-6-5}{3.6.5})
\begin{equation}\tag{\SGARefTarget{sga7.sga7i.eviii.n1-1-1}{1.1.1}}
 \Biext^1(P,Q;G)
 \simeq
 \Ext^1\!\left(P\mathbin{\overset{\mathbb L}{\otimes}}Q,G\right).
\end{equation}
Recall the canonical isomorphisms (called the Cartan isomorphisms)
\begin{equation}\tag{\SGARefTarget{sga7.sga7i.eviii.n1-1-2}{1.1.2}}
\begin{split}
 \Ext^1\!\left(P\mathbin{\overset{\mathbb L}{\otimes}}Q,G\right)
 &\simeq \Ext^1\!\left(P,R\uHom(Q,G)\right)\\
 &\simeq \Ext^1\!\left(Q,R\uHom(P,G)\right).
\end{split}
\end{equation}
All these isomorphisms are functorial in the three arguments $P,Q,G$.

\SGAvar{From (\SGARefLink{sga7.sga7i.eviii.n1-1-1}{1.1.1}) one deduces the
following infinite exact sequence,
taking into account that the homology objects of
$P\mathbin{\overset{\mathbb L}{\otimes}}Q$ are the $\Tor_i(P,Q)$,
which vanish for $i\geqslant1$:
\begin{equation}\tag{\SGARefTarget{sga7.sga7i.eviii.n1-1-3}{1.1.3}}
\begin{aligned}
0\longrightarrow{}&
 \Ext^1(P\otimes Q,G)
 \longrightarrow \Biext^1(P,Q;G)
 \longrightarrow \Hom(\Tor_1(P,Q),G)\\
 \longrightarrow{}&
 \Ext^2(P\otimes Q,G)
 \longrightarrow
 \Ext^2\!\left(P\mathbin{\overset{\mathbb L}{\otimes}}Q,G\right)
 \longrightarrow\cdots .
\end{aligned}
\end{equation}
On the other hand, since the cohomology objects of $R\uHom(Q,G)$ are
the $\uExt^i(Q,G)$, the first equality in (\SGARefLink{sga7.sga7i.eviii.n1-1-2}{1.1.2}) gives rise to the
five-term exact sequence
\begin{equation}\tag{\SGARefTarget{sga7.sga7i.eviii.n1-1-4}{1.1.4}}
\begin{aligned}
0\longrightarrow{}&
 \Ext^1(P,\uHom(Q,G))
 \longrightarrow \Biext^1(P,Q;G)
 \longrightarrow \Hom(P,\uExt^1(Q,G))\\
 {}&
 \Ext^2(P,\uHom(Q,P))
 \longrightarrow \Ext^2(P,R\uHom(Q,G)),
\end{aligned}
\end{equation}
associated with the usual spectral sequence
\[
 \Ext^*(P,R\uHom(Q,G))
 \Longleftarrow
 \Ext^p(P,\uExt^q(Q,G)).
\]
We leave it to the reader, should the need arise, to make the first
arrow of (\SGARefLink{sga7.sga7i.eviii.n1-1-3}{1.1.3}), respectively of (\SGARefLink{sga7.sga7i.eviii.n1-1-4}{1.1.4}), explicit by means of the
\underline{natural fully faithful functor}
\begin{equation}\tag{\SGARefTarget{sga7.sga7i.eviii.n1-1-5}{1.1.5}}
 \operatorname{EXT}(P\otimes Q,G)
 \longrightarrow
 \operatorname{BIEXT}(P,Q;G),
\end{equation}
respectively
\begin{equation}\tag{\SGARefTarget{sga7.sga7i.eviii.n1-1-6}{1.1.6}}
 \operatorname{EXT}(P,\uHom(Q,G))
 \longrightarrow
 \operatorname{BIEXT}(P,Q;G),
\end{equation}
where $\operatorname{EXT}$ and $\operatorname{BIEXT}$ denote the
evident categories of extensions and biextensions.

}{Taking the standard truncation of the derived tensor product, whose
homology objects are the $\Tor_i(P,Q)$, one obtains from (\SGARefLink{sga7.sga7i.eviii.n1-1-1}{1.1.1}) the
following infinite exact sequence:
\begin{equation}\tag{\SGARefTarget{sga7.sga7i.eviii.n1-1-3}{1.1.3}}
\begin{aligned}
0\longrightarrow{}&
 \Ext^1(P\otimes Q,G)
 \longrightarrow \Biext^1(P,Q;G)
 \longrightarrow \Hom(\Tor_1(P,Q),G)\\
 \longrightarrow{}&
 \Ext^2(P\otimes Q,G)
 \longrightarrow
 \Ext^2\!\left(P\mathbin{\overset{\mathbb L}{\otimes}}Q,G\right)
 \longrightarrow\cdots .
\end{aligned}
\end{equation}
On the other hand, since the cohomology objects of $R\uHom(Q,G)$ are
the $\uExt^i(Q,G)$, the first equality in (\SGARefLink{sga7.sga7i.eviii.n1-1-2}{1.1.2}) gives rise to the
five-term exact sequence
\begin{equation}\tag{\SGARefTarget{sga7.sga7i.eviii.n1-1-4}{1.1.4}}
\begin{aligned}
0\longrightarrow{}&
 \Ext^1(P,\uHom(Q,G))
 \longrightarrow \Biext^1(P,Q;G)
 \longrightarrow \Hom(P,\uExt^1(Q,G))\\
 \longrightarrow{}&
 \Ext^2(P,\uHom(Q,G))
 \longrightarrow \Ext^2(P,R\uHom(Q,G)),
\end{aligned}
\end{equation}
associated with the usual spectral sequence
\[
 \Ext^*(P,R\uHom(Q,G))
 \Longleftarrow
 \Ext^p(P,\uExt^q(Q,G)).
\]
We leave it to the reader, should the need arise, to make the first
arrow of (\SGARefLink{sga7.sga7i.eviii.n1-1-3}{1.1.3}), respectively of (\SGARefLink{sga7.sga7i.eviii.n1-1-4}{1.1.4}), explicit by means of the
\underline{natural fully faithful functor}
\begin{equation}\tag{\SGARefTarget{sga7.sga7i.eviii.n1-1-5}{1.1.5}}
 \operatorname{EXT}(P\otimes Q,G)
 \longrightarrow
 \operatorname{BIEXT}(P,Q;G),
\end{equation}
respectively
\begin{equation}\tag{\SGARefTarget{sga7.sga7i.eviii.n1-1-6}{1.1.6}}
 \operatorname{EXT}(P,\uHom(Q,G))
 \longrightarrow
 \operatorname{BIEXT}(P,Q;G),
\end{equation}
where $\operatorname{EXT}$ and $\operatorname{BIEXT}$ denote the
evident categories of extensions and biextensions.

}{E7I-VIII-001}\SGApage{I}{\(\approx\)221}
Of course, the isomorphism
\[
 \Biext^1(P,Q;G)\simeq\Ext^1(Q,R\uHom(P,G))
\]
gives rise to an exact sequence analogous to (\SGARefLink{sga7.sga7i.eviii.n1-1-3}{1.1.3}), with the roles
of $P$ and $Q$ interchanged.  Its first arrow can be made explicit by
means of the fully faithful functor
\begin{equation}\tag{\SGARefTarget{sga7.sga7i.eviii.n1-1-7}{1.1.7}}
 \operatorname{EXT}(Q,\uHom(P,G))
 \longrightarrow
 \operatorname{BIEXT}(P,Q;G)
\end{equation}
analogous to (\SGARefLink{sga7.sga7i.eviii.n1-1-6}{1.1.6}).

\SGARefTarget{sga7.sga7i.eviii.n1-2}{1.2}. Suppose that
\begin{equation}\tag{\SGARefTarget{sga7.sga7i.eviii.n1-2-1}{1.2.1}}
 \Tor_1(P,Q)=0.
\end{equation}
Then (\SGARefLink{sga7.sga7i.eviii.n1-1-3}{1.1.3}) supplies a canonical isomorphism
\begin{equation}\tag{\SGARefTarget{sga7.sga7i.eviii.n1-2-2}{1.2.2}}
 \Biext^1(P,Q;G)\simeq\Ext^1(P\otimes Q,G).
\end{equation}
This holds, for example, if $P$ or $Q$ is a \underline{flat}
$\mathbb Z$-module.  More precisely, under condition (\SGARefLink{sga7.sga7i.eviii.n1-2-1}{1.2.1}), the
functor (\SGARefLink{sga7.sga7i.eviii.n1-1-5}{1.1.5}) is an equivalence of categories.

\SGApage{I}{222}
\SGAvar{\SGARefTarget{sga7.sga7i.eviii.n1-3}{1.3}. Suppose that
\begin{equation}\tag{\SGARefTarget{sga7.sga7i.eviii.n1-3-1}{1.3.1}}
 P\otimes Q=0.
\end{equation}
Then (\SGARefLink{sga7.sga7i.eviii.n1-1-3}{1.1.3}) supplies a canonical isomorphism
\begin{equation}\tag{\SGARefTarget{sga7.sga7i.eviii.n1-3-2}{1.3.2}}
 \Biext^1(P,Q;G)\simeq\Hom(\Tor_1(P,Q),G).
\end{equation}
The case of greatest interest to us in which (\SGARefLink{sga7.sga7i.eviii.n1-3-1}{1.3.1}) holds is that in
which there is an integer $n>0$ such that
\begin{equation}\tag{\SGARefTarget{sga7.sga7i.eviii.n1-3-3}{1.3.3}}
 nP=P,\qquad
 Q=\varinjlim_i {}_{n^i}Q;
\end{equation}
that is, $P$ is divisible by $n$, and $Q$ is the limit of subgroups
annihilated by powers of $n$.  Put, as usual,
\[
 T_n(P)=({}_{n^i}P)_{i\geq0}\quad\text{(a projective system)},\qquad
 {}_{n^\infty}P=\varinjlim_i{}_{n^i}P.
\]
From (\SGARefLink{sga7.sga7i.eviii.n1-3-3}{1.3.3}) one obtains the following formulas, more precise than
(\SGARefLink{sga7.sga7i.eviii.n1-3-1}{1.3.1}):
\begin{equation}\tag{\SGARefTarget{sga7.sga7i.eviii.n1-3-4}{1.3.4}}
\begin{split}
 P\otimes Q&=0,\\
 \Tor_1(P,Q)
 &\simeq\varinjlim_i
   \bigl({}_{n^i}P\otimes{}_{n^i}Q\bigr)
 \simeq T_n(P)\otimes Q\\
 &\simeq{}_{n^\infty}P\otimes T_n(Q).
\end{split}
\end{equation}
Thus the isomorphism (\SGARefLink{sga7.sga7i.eviii.n1-3-2}{1.3.2}) here assumes the more concrete form
\begin{equation}\tag{\SGARefTarget{sga7.sga7i.eviii.n1-3-5}{1.3.5}}
\begin{split}
 \Biext^1(P,Q;G)
 &\simeq\Hom(T_n(P),D(Q))\\
 &\simeq\Hom(Q,D(T_n(P))\\
 &\simeq\Hom(T_n(P)\otimes T_n(Q),T_n(G)).
\end{split}
\end{equation}
Let us specify the meaning of these formulas.

\begin{enumerate}[label=\alph*)]
\item In the inductive system occurring in (\SGARefLink{sga7.sga7i.eviii.n1-3-4}{1.3.4}), the transition
homomorphism
\[
 {}_{n^i}P\otimes{}_{n^i}Q
 \longrightarrow
 {}_{n^j}P\otimes{}_{n^j}Q
 \qquad(i\leq j)
\]
is described by identifying the first member with
$\,{}_{n^j}P\otimes{}_{n^i}Q$ by means of
$n^{j-i}\otimes\id$, and then using the homomorphism
\[
 \id\otimes\operatorname{inj}:
 {}_{n^j}P\otimes{}_{n^i}Q
 \longrightarrow
 {}_{n^j}P\otimes{}_{n^j}Q,
\]
where
$\operatorname{inj}:{}_{n^i}Q\longrightarrow{}_{n^j}Q$
denotes the inclusion.

\item By definition, put
\[
 T_n(P)\otimes Q
 \overset{\mathrm{def}}{=}
 \varinjlim_i\bigl(T_n(P)\otimes{}_{n^i}Q\bigr),
\]
where
\[
 T_n(P)\otimes M
 \overset{\mathrm{def}}{=}
 {}_{n^i}P\otimes M
 \qquad\text{if }n^iM=0.
\]
The meaning of the third term written in (\SGARefLink{sga7.sga7i.eviii.n1-3-4}{1.3.4}) is specified in the
same way when $nQ=Q$ is also assumed, so that $T_n(Q)$, like
$T_n(P)$, is a strict projective system.  The last term of (\SGARefLink{sga7.sga7i.eviii.n1-3-5}{1.3.4}) is
a Hom between projective systems of groups.

\item The letter $D$ denotes $\uHom(-,G)$.  The object $D(Q)$ is
defined as the projective system of groups
\[
 D(Q)=\bigl(D({}_{n^i}Q)\bigr)_{i\geq0}
 =\bigl(\uHom({}_{n^i}Q,G)\bigr)_{i\geq0},
\]
the second member of (\SGARefLink{sga7.sga7i.eviii.n1-3-5}{1.3.5}) being the group of homomorphisms of
projective systems of groups.  The following definition gives a
simpler interpretation of the last member:
\[
 D(T_n(P))=\varinjlim_iD({}_{n^i}P),
\]
where the transition morphisms of this inductive system are obtained
by transposition from the transition morphisms
\[
 n^{j-1}:{}_{n^j}P\longrightarrow{}_{n^i}P
\]
of $T_n(P)$.
\end{enumerate}

}{\SGARefTarget{sga7.sga7i.eviii.n1-3}{1.3}. Suppose that
\begin{equation}\tag{\SGARefTarget{sga7.sga7i.eviii.n1-3-1}{1.3.1}}
 P\otimes Q=0.
\end{equation}
Then (\SGARefLink{sga7.sga7i.eviii.n1-1-3}{1.1.3}) supplies a canonical isomorphism
\begin{equation}\tag{\SGARefTarget{sga7.sga7i.eviii.n1-3-2}{1.3.2}}
 \Biext^1(P,Q;G)\simeq\Hom(\Tor_1(P,Q),G).
\end{equation}
The case of greatest interest to us in which (\SGARefLink{sga7.sga7i.eviii.n1-3-1}{1.3.1}) holds is that in
which there is an integer $n>0$ such that
\begin{equation}\tag{\SGARefTarget{sga7.sga7i.eviii.n1-3-3}{1.3.3}}
 nP=P,\qquad
 Q=\varinjlim_i {}_{n^i}Q;
\end{equation}
that is, $P$ is divisible by $n$, and $Q$ is the limit of subgroups
annihilated by powers of $n$.  Put, as usual,
\[
 T_n(P)=({}_{n^i}P)_{i\geq0}\quad\text{(a projective system)},\qquad
 {}_{n^\infty}P=\varinjlim_i{}_{n^i}P.
\]
From (\SGARefLink{sga7.sga7i.eviii.n1-3-3}{1.3.3}) one obtains the following formulas, more precise than
(\SGARefLink{sga7.sga7i.eviii.n1-3-1}{1.3.1}):
\begin{equation}\tag{\SGARefTarget{sga7.sga7i.eviii.n1-3-4}{1.3.4}}
\begin{split}
 P\otimes Q&=0,\\
 \Tor_1(P,Q)
 &\simeq\varinjlim_i
   \bigl({}_{n^i}P\otimes{}_{n^i}Q\bigr)
 \simeq T_n(P)\otimes Q\\
 &\simeq{}_{n^\infty}P\otimes T_n(Q).
\end{split}
\end{equation}
Thus the isomorphism (\SGARefLink{sga7.sga7i.eviii.n1-3-2}{1.3.2}) here assumes the more concrete form
\begin{equation}\tag{\SGARefTarget{sga7.sga7i.eviii.n1-3-5}{1.3.5}}
\begin{split}
 \Biext^1(P,Q;G)
 &\simeq\Hom(T_n(P),D(Q))\\
 &\simeq\Hom(Q,D(T_n(P)))\\
 &\simeq\Hom(T_n(P)\otimes T_n(Q),T_n(G)).
\end{split}
\end{equation}
Let us specify the meaning of these formulas.

\begin{enumerate}[label=\alph*)]
\item In the inductive system occurring in (\SGARefLink{sga7.sga7i.eviii.n1-3-4}{1.3.4}), the transition
homomorphism
\[
 {}_{n^i}P\otimes{}_{n^i}Q
 \longrightarrow
 {}_{n^j}P\otimes{}_{n^j}Q
 \qquad(i\leq j)
\]
is described by identifying the first member with
$\,{}_{n^j}P\otimes{}_{n^i}Q$ by means of
$n^{j-i}\otimes\id$, and then using the homomorphism
\[
 \id\otimes\operatorname{inj}:
 {}_{n^j}P\otimes{}_{n^i}Q
 \longrightarrow
 {}_{n^j}P\otimes{}_{n^j}Q,
\]
where
$\operatorname{inj}:{}_{n^i}Q\longrightarrow{}_{n^j}Q$
denotes the inclusion.

\item By definition, put
\[
 T_n(P)\otimes Q
 \overset{\mathrm{def}}{=}
 \varinjlim_i\bigl(T_n(P)\otimes{}_{n^i}Q\bigr),
\]
where
\[
 T_n(P)\otimes M
 \overset{\mathrm{def}}{=}
 {}_{n^i}P\otimes M
 \qquad\text{if }n^iM=0.
\]
The meaning of the third term written in (\SGARefLink{sga7.sga7i.eviii.n1-3-4}{1.3.4}) is specified in the
same way when $nQ=Q$ is also assumed, so that $T_n(Q)$, like
$T_n(P)$, is a strict projective system.  The last term of (\SGARefLink{sga7.sga7i.eviii.n1-3-5}{1.3.5}) is
a Hom between projective systems of groups.

\item The letter $D$ denotes $\uHom(-,G)$.  The object $D(Q)$ is
defined as the projective system of groups
\[
 D(Q)=\bigl(D({}_{n^i}Q)\bigr)_{i\geq0}
 =\bigl(\uHom({}_{n^i}Q,G)\bigr)_{i\geq0},
\]
the second member of (\SGARefLink{sga7.sga7i.eviii.n1-3-5}{1.3.5}) being the group of homomorphisms of
projective systems of groups.  The following definition gives a
simpler interpretation of the next member:
\[
 D(T_n(P))=\varinjlim_iD({}_{n^i}P),
\]
where the transition morphisms of this inductive system are obtained
by transposition from the transition morphisms
\[
 n^{j-i}:{}_{n^j}P\longrightarrow{}_{n^i}P
\]
of $T_n(P)$.
\end{enumerate}

}{E7I-VIII-002}\SGApage{I}{\(\approx\)224}
We leave it to the reader to make explicit the calculation of
$\Tor_1(P,Q)$ indicated in (\SGARefLink{sga7.sga7i.eviii.n1-3-4}{1.3.4}), and to derive (\SGARefLink{sga7.sga7i.eviii.n1-3-5}{1.3.5}) from it.
We note only that one uses the relations, deduced from the second
hypothesis in (\SGARefLink{sga7.sga7i.eviii.n1-3-3}{1.3.3}),
\[
 \Tor_p(P,Q)\simeq
 \varinjlim_i\Tor_p(P,{}_{n^i}Q).
\]
For $p=0$ this immediately gives $P\otimes Q=0$, and it reduces the
calculation of $\Tor_1(P,Q)$ to that of the
$\Tor_1(P,{}_{n^i}Q)$ and their transition morphisms.  Thus, when
$Q$ is annihilated by an integer $m=n^i$ such that $mP=P$, it remains
to make explicit an isomorphism
\begin{equation}\tag{\SGARefTarget{sga7.sga7i.eviii.n1-3-6}{1.3.6}}
 \Tor_1(P,Q)
 \xleftarrow{\ \sim\ }
 {}_mP\otimes{}_mQ
 \qquad(mP=P,\ mQ=0).
\end{equation}
There are several ways to do this.  One consists in using the exact
Tor sequence associated with
\[
 0\longrightarrow{}_mP\longrightarrow P
 \xrightarrow{\ m\ }P\longrightarrow0.
\]
One may also use the homomorphism (\SGARefLink{sga7.sga7i.eviii.n2-1-1}{2.1.1}) below; the fact that the
transition morphisms are those specified in part (a) is then precisely
\SGARefLink{sga7.sga7i.eviii.n2-1-11}{2.1.11}.  In any case, the choice of (\SGARefLink{sga7.sga7i.eviii.n1-3-6}{1.3.6}) amounts to the choice of a
sign; see \SGARefLink{sga7.sga7i.eviii.n2-1-10}{2.1.10}.

\SGAvar{\SGARefTarget{sga7.sga7i.eviii.n1-4}{1.4}. Suppose that
\begin{equation}\tag{\SGARefTarget{sga7.sga7i.eviii.n1-4-1}{1.4.1}}
 \uHom(Q,G)=0.
\end{equation}
Then (\SGARefLink{sga7.sga7i.eviii.n1-1-4}{1.1.4}) supplies a canonical isomorphism
\begin{equation}\tag{\SGARefTarget{sga7.sga7i.eviii.n1-4-2}{1.4.2}}
 \Biext^1(P,Q;G)
 \xrightarrow{\ \sim\ }
 \Hom(P,\uExt^1(Q,G)).
\end{equation}
This case occurs in particular when $Q$ is an abelian scheme and
$G=\underline G_m$; see VII \SGARefLink{sga7.sga7i.evii.n1-3-8}{1.3.8}.  Hence the functor in $P$,
\[
 P\longmapsto\Biext^1(P,Q;G),
\]
is \underline{representable} by $Q^*=\uExt^1(Q,G)$.  In particular,
there is a canonical biextension $E^Q$ of $(Q^*,Q)$ by $G$, determined
up to a unique isomorphism, since its automorphism group is
\[
 \Hom(P,\uHom(Q,G))=0.
\]
This biextension, or its symmetric biextension $\,{}^sE^Q$ of
$(Q,Q^*)$ by $G$, may be called the \underline{Weil biextension} of
$(Q^*,Q)$, respectively $(Q,Q^*)$, by $G$.

}{\SGARefTarget{sga7.sga7i.eviii.n1-4}{1.4}. Suppose that
\begin{equation}\tag{\SGARefTarget{sga7.sga7i.eviii.n1-4-1}{1.4.1}}
 \uHom(Q,G)=0.
\end{equation}
Then (\SGARefLink{sga7.sga7i.eviii.n1-1-4}{1.1.4}) supplies a canonical isomorphism
\begin{equation}\tag{\SGARefTarget{sga7.sga7i.eviii.n1-4-2}{1.4.2}}
 \Biext^1(P,Q;G)
 \xrightarrow{\ \sim\ }
 \Hom(P,\uExt^1(Q,G)).
\end{equation}
This case occurs in particular when $Q$ is an abelian scheme and
$G=\underline G_m$; see VII \SGARefLink{sga7.sga7i.evii.n1-3-8}{1.3.8}.  Hence the functor in $P$,
\[
 P\longmapsto\Biext^1(P,Q;G),
\]
is \underline{representable} by $Q^*=\uExt^1(Q,G)$.  In particular,
there is a canonical biextension $E^Q$ of $(Q^*,Q)$ by $G$, determined
up to a unique isomorphism, since its automorphism group is
\[
 \Hom(Q^*,\uHom(Q,G))=0.
\]
This biextension, or its symmetric biextension $\,{}^sE^Q$ of
$(Q,Q^*)$ by $G$, may be called the \underline{Weil biextension} of
$(Q^*,Q)$, respectively $(Q,Q^*)$, by $G$.

}{E7I-VIII-003}\SGApage{I}{225}
\SGARefTarget{sga7.sga7i.eviii.n1-5}{1.5}. Suppose that
\begin{equation}\tag{\SGARefTarget{sga7.sga7i.eviii.n1-5-1}{1.5.1}}
 \Hom(P,\uExt^1(Q,G))=0.
\end{equation}
This holds in particular if $\uExt^1(Q,G)=0$; for example, if $Q$ is
a finite locally free group scheme or a group of multiplicative type
(SGA 3 IX), and $G=\underline G_m$; see \SGARefLink{sga7.sga7i.eviii.n3-4}{3.4} below.  Then (\SGARefLink{sga7.sga7i.eviii.n1-1-4}{1.1.4})
supplies a canonical isomorphism
\begin{equation}\tag{\SGARefTarget{sga7.sga7i.eviii.n1-5-2}{1.5.2}}
 \Biext^1(P,Q;G)
 \xleftarrow{\ \sim\ }
 \Ext^1(P,\uHom(Q,G)),
\end{equation}
where $\uHom(Q,G)$ plays the role of a "dual group" of $Q$ (relative
to $G$).  More precisely, it follows from (\SGARefLink{sga7.sga7i.eviii.n1-5-2}{1.5.2}) that the fully
faithful functor (\SGARefLink{sga7.sga7i.eviii.n1-1-6}{1.1.6}) is an \underline{equivalence of categories}.

\SGAvar{\SGARefTarget{sga7.sga7i.eviii.n1-6}{1.6}. Suppose that the two relations
\begin{equation}\tag{\SGARefTarget{sga7.sga7i.eviii.n1-6-1}{1.6.1}}
 \uHom(P,G)=0,\qquad
 \Hom(P,\uExt^1(Q,G))=0
\end{equation}
hold simultaneously.  By \SGARefLink{sga7.sga7i.eviii.n1-4}{1.4}, these two
relations give isomorphisms
\[
 \Biext^1(P,Q;G)
 \xrightarrow{\ \sim\ }
 \Hom(Q,\uExt^1(P,G)),
\]
\[
 \Biext^1(P,Q;G)
 \xleftarrow{\ \sim\ }
 \Ext^1(P,\uHom(Q,G)).
\]
Their composite is the canonical isomorphism
\begin{equation}\tag{\SGARefTarget{sga7.sga7i.eviii.n1-6-2}{1.6.2}}
 \Ext^1(P,\uHom(Q,G))
 \xrightarrow{\ \sim\ }
 \Hom(Q,\uExt^1(P,G))
 \quad\bigl(\simeq\Biext^1(P,G;G)\bigr),
\end{equation}
which plays the role of a duality formula.

Notice also that the first relation in (\SGARefLink{sga7.sga7i.eviii.n1-6-1}{1.6.1}) plainly implies
$\Hom(P\otimes Q,G)=0$.  Hence the category
$\operatorname{EXT}(P,\uHom(Q,G))$ of commutative extensions of $P$
by $\uHom(Q,G)$ is rigid (every automorphism of an object is the
identity), as is $\operatorname{BIEXT}(P,Q;G)$.  The structure of
these categories is therefore known once the group of isomorphism
classes of objects is known; by (\SGARefLink{sga7.sga7i.eviii.n1-6-2}{1.6.2}), this group is canonically
isomorphic to
\[
 \Hom(Q,\uExt^1(P,Q)).
\]

}{\SGARefTarget{sga7.sga7i.eviii.n1-6}{1.6}. Suppose that the two relations
\begin{equation}\tag{\SGARefTarget{sga7.sga7i.eviii.n1-6-1}{1.6.1}}
 \uHom(P,G)=0,\qquad
 \Hom(P,\uExt^1(Q,G))=0
\end{equation}
hold simultaneously.  By \SGARefLink{sga7.sga7i.eviii.n1-4}{1.4} and \SGARefLink{sga7.sga7i.eviii.n1-5}{1.5}, respectively, these two
relations give isomorphisms
\[
 \Biext^1(P,Q;G)
 \xrightarrow{\ \sim\ }
 \Hom(Q,\uExt^1(P,G)),
\]
\[
 \Biext^1(P,Q;G)
 \xleftarrow{\ \sim\ }
 \Ext^1(P,\uHom(Q,G)).
\]
Their composite is the canonical isomorphism
\begin{equation}\tag{\SGARefTarget{sga7.sga7i.eviii.n1-6-2}{1.6.2}}
 \Ext^1(P,\uHom(Q,G))
 \xrightarrow{\ \sim\ }
 \Hom(Q,\uExt^1(P,G))
 \quad\bigl(\simeq\Biext^1(P,Q;G)\bigr),
\end{equation}
which plays the role of a duality formula.

Notice also that the first relation in (\SGARefLink{sga7.sga7i.eviii.n1-6-1}{1.6.1}) plainly implies
$\Hom(P\otimes Q,G)=0$.  Hence the category
$\operatorname{EXT}(P,\uHom(Q,G))$ of commutative extensions of $P$
by $\uHom(Q,G)$ is rigid (every automorphism of an object is the
identity), as is $\operatorname{BIEXT}(P,Q;G)$.  The structure of
these categories is therefore known once the group of isomorphism
classes of objects is known; by (\SGARefLink{sga7.sga7i.eviii.n1-6-2}{1.6.2}), this group is canonically
isomorphic to
\[
 \Hom(Q,\uExt^1(P,G)).
\]

}{E7I-VIII-004}\SGApage{I}{\(\approx\)226}
The map (\SGARefLink{sga7.sga7i.eviii.n1-6-2}{1.6.2}) is in fact defined in all cases, without assuming
(\SGARefLink{sga7.sga7i.eviii.n1-6-1}{1.6.1}), and can also be described directly, without passing through
$\Biext^1(P,Q;G)$, in terms of the natural pairing
\[
 Q\times Q'\longrightarrow G,
\qquad Q'=\uHom(Q,G).
\]
Interpreting it as a homomorphism
$Q\longrightarrow\uHom(Q',G)$ gives a composite
\begin{equation}\tag{\SGARefTarget{sga7.sga7i.eviii.n1-6-3}{1.6.3}}
\begin{split}
 \Ext^1(P,Q')\times Q
 &\longrightarrow
 \Ext^1(P,Q')\times\Hom(Q',G)\\
 &\longrightarrow
 \Ext^1(P,G),
\end{split}
\end{equation}
which recovers (\SGARefLink{sga7.sga7i.eviii.n1-6-2}{1.6.2}).

For a useful case in which the considerations of \SGARefLink{sga7.sga7i.eviii.n1-6}{1.6} apply, see \SGARefLink{sga7.sga7i.eviii.n3-7}{3.7}
below.



\subsection*{\SGARefTarget{sga7.sga7i.eviii.n2}{2}. Pairings defined by a biextension}

\SGARefTarget{sga7.sga7i.eviii.n2-1}{2.1}. Let $P$ and $Q$ be two abelian groups of $T$, and let $n>0$ be
an integer.  Recall that $\,{}_nP$ denotes the kernel of multiplication
by $n\cdot\id_P$.  We shall define a canonical homomorphism
\begin{equation}\tag{\SGARefTarget{sga7.sga7i.eviii.n2-1-1}{2.1.1}}
 \alpha=\alpha^n_{P,Q}:
 {}_nP\otimes{}_nQ\longrightarrow\uTor_1(P,Q),
\end{equation}
which will be the composite
\begin{equation}\tag{\SGARefTarget{sga7.sga7i.eviii.n2-1-2}{2.1.2}}
 {}_nP\otimes{}_nQ
 \longrightarrow \uTor_1({}_nP,{}_nQ)
 \longrightarrow \uTor_1(P,Q),
\end{equation}
where the first arrow is $\alpha^n_{{}_nP,{}_nQ}$ and the second is
induced by the inclusions $\,{}_nP\longrightarrow P$ and
$\,{}_nQ\longrightarrow Q$.  We are thus reduced to defining (\SGARefLink{sga7.sga7i.eviii.n2-1-1}{2.1.1})
when $nP=0$ and $nQ=0$, that is, when $P$ and $Q$ are
$\mathbb Z/n\mathbb Z$-modules.

\SGApage{I}{\(\approx\)227}
\SGAvar{Suppose only that $nQ=0$.  Taking a flat resolution $L_\bullet$ of
$P$ (as a $\mathbb Z$-module), one obtains canonical isomorphisms
\[
 P\mathbin{\overset{\mathbb L}{\otimes}_{\mathbb Z}}Q
 \simeq L_\bullet\otimes_{\mathbb Z}Q
 \simeq
 (L_\bullet\otimes_{\mathbb Z}\mathbb Z/n\mathbb Z)
 \otimes_{\mathbb Z/n\mathbb Z}Q,
\]
which may be written as an isomorphism in the derived category
\begin{equation}\tag{\SGARefTarget{sga7.sga7i.eviii.n2-1-3}{2.1.3}}
 P\mathbin{\overset{\mathbb L}{\otimes}_{\mathbb Z}}Q
 \simeq
 \left(P\mathbin{\overset{\mathbb L}{\otimes}_{\mathbb Z}}
       \mathbb Z/n\mathbb Z\right)
 \mathbin{\overset{\mathbb L}{\otimes}_{\mathbb Z/n\mathbb Z}}Q.
\end{equation}
This gives rise to the homological K\"unneth-type spectral sequence
(cf.\ SGA 1 IV 5.2)
\begin{equation}\tag{\SGARefTarget{sga7.sga7i.eviii.n2-1-4}{2.1.4}}
 \uTor^{\mathbb Z}_*(P,Q)
 \Longleftarrow
 E^2_{p,q}
 =
 \uTor^{\mathbb Z/n\mathbb Z}_p
 \bigl(\uTor^{\mathbb Z}_q(P,\mathbb Z/n\mathbb Z),Q\bigr).
\end{equation}
Recall also the canonical isomorphisms
\[
 \uTor^{\mathbb Z}_0(P,\mathbb Z/n\mathbb Z)
 =P_n\mathrel{\mathop:}=P/nP,
 \qquad
 \uTor^{\mathbb Z}_1(P,\mathbb Z/n\mathbb Z)\simeq{}_nP,
\]
and
\[
 \uTor^{\mathbb Z}_i(P,\mathbb Z/n\mathbb Z)=0
 \qquad\text{for }i\geq2.
\]
They follow by computing these objects from the canonical flat
resolution of $(\mathbb Z/n\mathbb Z)_T$:
\[
\begin{tikzcd}[column sep=large,row sep=large]
 0\arrow[r]\&
 \mathbb Z_T\arrow[r,"n"]\arrow[d,equal]\&
 \mathbb Z_T\arrow[r]\arrow[d,equal]\&
 0\\
 \&L_1\&L_0\&
\end{tikzcd}
\]
whose augmentation
$L_0\longrightarrow(\mathbb Z/n\mathbb Z)_T$ sends $1$ to $1$.
Consequently, the spectral sequence (\SGARefLink{sga7.sga7i.eviii.n2-1-4}{2.1.4}) reduces essentially to
the exact sequence
\begin{equation}\tag{\SGARefTarget{sga7.sga7i.eviii.n2-1-5}{2.1.5}}
\begin{split}
0\longrightarrow{}&
 \uTor^{\mathbb Z/n\mathbb Z}_2(P_n,Q)
 \longrightarrow{}_nP\otimes Q
 \xrightarrow{\alpha}\uTor^{\mathbb Z}_1(P,Q)\\
\longrightarrow{}&
 \uTor^{\mathbb Z/n\mathbb Z}_1(P_n,Q)
 \longrightarrow0,
\end{split}
\end{equation}
where $P_n=P\otimes_{\mathbb Z}\mathbb Z/n\mathbb Z=P/nP$.
The desired homomorphism $\alpha$ is the edge homomorphism
$E^2_{0,1}\longrightarrow\uTor^{\mathbb Z}_1$ in this exact sequence.
Since the sequence is plainly functorial in $P$ and $Q$, functoriality
in $P$ shows that the homomorphism in question is indeed the composite
\[
 {}_nP\otimes Q
 \longrightarrow\uTor_1({}_nP,Q)
 \longrightarrow\uTor_1(P,Q),
\]
where the first homomorphism is the $\alpha$ just defined, with
$(P,Q)$ replaced by $({}_nP,Q)$, and the second is induced by the
inclusion $\,{}_nP\longrightarrow Q$.

}{Suppose only that $nQ=0$.  Taking a flat resolution $L_\bullet$ of
$P$ (as a $\mathbb Z$-module), one obtains canonical isomorphisms
\[
 P\mathbin{\overset{\mathbb L}{\otimes}_{\mathbb Z}}Q
 \simeq L_\bullet\otimes_{\mathbb Z}Q
 \simeq
 (L_\bullet\otimes_{\mathbb Z}\mathbb Z/n\mathbb Z)
 \otimes_{\mathbb Z/n\mathbb Z}Q,
\]
which may be written as an isomorphism in the derived category
\begin{equation}\tag{\SGARefTarget{sga7.sga7i.eviii.n2-1-3}{2.1.3}}
 P\mathbin{\overset{\mathbb L}{\otimes}_{\mathbb Z}}Q
 \simeq
 \left(P\mathbin{\overset{\mathbb L}{\otimes}_{\mathbb Z}}
       \mathbb Z/n\mathbb Z\right)
 \mathbin{\overset{\mathbb L}{\otimes}_{\mathbb Z/n\mathbb Z}}Q.
\end{equation}
This gives rise to the homological K\"unneth-type spectral sequence
(cf.\ SGA 1 IV 5.2)
\begin{equation}\tag{\SGARefTarget{sga7.sga7i.eviii.n2-1-4}{2.1.4}}
 \uTor^{\mathbb Z}_*(P,Q)
 \Longleftarrow
 E^2_{p,q}
 =
 \uTor^{\mathbb Z/n\mathbb Z}_p
 \bigl(\uTor^{\mathbb Z}_q(P,\mathbb Z/n\mathbb Z),Q\bigr).
\end{equation}
Recall also the canonical isomorphisms
\[
 \uTor^{\mathbb Z}_0(P,\mathbb Z/n\mathbb Z)
 =P_n\mathrel{\mathop:}=P/nP,
 \qquad
 \uTor^{\mathbb Z}_1(P,\mathbb Z/n\mathbb Z)\simeq{}_nP,
\]
and
\[
 \uTor^{\mathbb Z}_i(P,\mathbb Z/n\mathbb Z)=0
 \qquad\text{for }i\geq2.
\]
They follow by computing these objects from the canonical flat
resolution of $(\mathbb Z/n\mathbb Z)_T$:
\[
\begin{tikzcd}[column sep=large,row sep=large]
 0\arrow[r]\&
 \mathbb Z_T\arrow[r,"n"]\arrow[d,equal]\&
 \mathbb Z_T\arrow[r]\arrow[d,equal]\&
 0\\
 \&L_1\&L_0\&
\end{tikzcd}
\]
whose augmentation
$L_0\longrightarrow(\mathbb Z/n\mathbb Z)_T$ sends $1$ to $1$.
Consequently, the spectral sequence (\SGARefLink{sga7.sga7i.eviii.n2-1-4}{2.1.4}) reduces essentially to
the exact sequence
\begin{equation}\tag{\SGARefTarget{sga7.sga7i.eviii.n2-1-5}{2.1.5}}
\begin{split}
0\longrightarrow{}&
 \uTor^{\mathbb Z/n\mathbb Z}_2(P_n,Q)
 \longrightarrow{}_nP\otimes Q
 \xrightarrow{\alpha}\uTor^{\mathbb Z}_1(P,Q)\\
\longrightarrow{}&
 \uTor^{\mathbb Z/n\mathbb Z}_1(P_n,Q)
 \longrightarrow0,
\end{split}
\end{equation}
where $P_n=P\otimes_{\mathbb Z}\mathbb Z/n\mathbb Z=P/nP$.
The desired homomorphism $\alpha$ is the edge homomorphism
$E^2_{0,1}\longrightarrow\uTor^{\mathbb Z}_1$ in this exact sequence.
Since the sequence is plainly functorial in $P$ and $Q$, functoriality
in $P$ shows that the homomorphism in question is indeed the composite
\[
 {}_nP\otimes Q
 \longrightarrow\uTor_1({}_nP,Q)
 \longrightarrow\uTor_1(P,Q),
\]
where the first homomorphism is the $\alpha$ just defined, with
$(P,Q)$ replaced by $({}_nP,Q)$, and the second is induced by the
inclusion $\,{}_nP\longrightarrow P$.

}{E7I-VIII-005}\SGApage{I}{228}
\SGARefTarget{sga7.sga7i.eviii.n2-1-6}{2.1.6}. This therefore defines, for arbitrary $P$ and $Q$, a
homomorphism (\SGARefLink{sga7.sga7i.eviii.n2-1-1}{2.1.1}) satisfying the factorization condition (\SGARefLink{sga7.sga7i.eviii.n2-1-2}{2.1.2}).
Moreover, when $nQ=0$, the exact sequence (\SGARefLink{sga7.sga7i.eviii.n2-1-5}{2.1.5}) shows that $\alpha$
is an isomorphism whenever $P_n$ is flat over
$\mathbb Z/n\mathbb Z$.  This holds in particular if $P_n=0$, that
is, if $nP=P$ and $P$ is $n$-divisible, or if $P$ itself is a flat
$\mathbb Z/n\mathbb Z$-module.  It also holds if $Q$ is a flat
$\mathbb Z/n\mathbb Z$-module.

\SGARefTarget{sga7.sga7i.eviii.n2-1-7}{2.1.7}. Let us make the calculation of the homomorphism $\alpha$ in
(\SGARefLink{sga7.sga7i.eviii.n2-1-5}{2.1.5}) explicit.  To determine it, it suffices, for every section
$q$ of $Q$ (over an object $S$ of $T$; take the final object for
simplicity), to determine the corresponding homomorphism
\[
 {}_nP\longrightarrow\uTor^{\mathbb Z}_1(P,Q).
\]
The section $q$ may be regarded as the image of the section $1$ of
$(\mathbb Z/n\mathbb Z)_T$ under the uniquely determined
homomorphism
\[
 u_q:(\mathbb Z/n\mathbb Z)_T\longrightarrow Q.
\]
Functoriality of $\alpha$ in $Q$ therefore reduces us to the case
$Q=(\mathbb Z/n\mathbb Z)_T$, with $q$ its unit section.  In this
case, identify
\[
 \uTor^{\mathbb Z}_1(P,Q)
 =
 \uTor^{\mathbb Z}_1(P,\mathbb Z/n\mathbb Z)
\]
with $\,{}_nP$ as explained after (\SGARefLink{sga7.sga7i.eviii.n2-1-4}{2.1.4}).  The homomorphism in question
is then the identity.  Equivalently, the edge homomorphism
\[
 E^2_{0,q}
 =
 \uTor^{\mathbb Z}_q(P,\mathbb Z/n\mathbb Z)
 \otimes_{\mathbb Z/n\mathbb Z}\mathbb Z/n\mathbb Z
 \longrightarrow
 \uTor^{\mathbb Z}_q(P,\mathbb Z/n\mathbb Z)
\]
is the identity.  This follows immediately from the definitions.
The analogous assertion remains valid when $\mathbb Z$ and its
quotient $\mathbb Z/n\mathbb Z$ are replaced by any ring $A$ on $T$
and any quotient ring $A/I$.

\SGApage{I}{229}
\SGAvar{\SGARefTarget{sga7.sga7i.eviii.n2-1-8}{2.1.8}. The homomorphism $\alpha$ of (\SGARefLink{sga7.sga7i.eviii.n2-1-1}{2.1.1}) may also be determined by
describing, for every pair of sections $p$ of $\,{}_nP$ and $q$ of
$\,{}_nQ$ (over an object $S$ of $T$, again taken to be final to
simplify the notation), the section of
$\uTor^{\mathbb Z}_1(P,Q)$ determined by $(p,q)$.  By functoriality in
$P$ and $Q$, as above, we are reduced to the case
\[
 P=Q=(\mathbb Z/n\mathbb Z)_T,
\]
with $p$ and $q$ the unit section.  Thus we need only determine
\[
 \alpha:
 (\mathbb Z/n\mathbb Z)_T\otimes(\mathbb Z/n\mathbb Z)_T
 \longrightarrow
 \uTor^{\mathbb Z}_1
 \bigl(\mathbb Z/n\mathbb Z)_T,(\mathbb Z/n\mathbb Z)_T\bigr).
\]
The description in \SGARefLink{sga7.sga7i.eviii.n2-1-7}{2.1.7} shows that this is the isomorphism
\begin{equation}\tag{\SGARefTarget{sga7.sga7i.eviii.n2-1-8-1}{2.1.8.1}}
 (\mathbb Z/n\mathbb Z)_T
 \xrightarrow{\sim}
 \uTor^{\mathbb Z}_1
 (\mathbb Z/n\mathbb Z)_T,(\mathbb Z/n\mathbb Z)_T\bigr)
\end{equation}
deduced from the canonical flat resolution
\begin{equation}\tag{\SGARefTarget{sga7.sga7i.eviii.n2-1-8-2}{2.1.8.2}}
 0\longrightarrow\mathbb Z_T
 \xrightarrow{n}\mathbb Z_T
 \longrightarrow0
\end{equation}
of the \underline{first} factor
$P=(\mathbb Z/n\mathbb Z)_T$ occurring in $\uTor_1$.



}{\SGARefTarget{sga7.sga7i.eviii.n2-1-8}{2.1.8}. The homomorphism $\alpha$ of (\SGARefLink{sga7.sga7i.eviii.n2-1-1}{2.1.1}) may also be determined by
describing, for every pair of sections $p$ of $\,{}_nP$ and $q$ of
$\,{}_nQ$ (over an object $S$ of $T$, again taken to be final to
simplify the notation), the section of
$\uTor^{\mathbb Z}_1(P,Q)$ determined by $(p,q)$.  By functoriality in
$P$ and $Q$, as above, we are reduced to the case
\[
 P=Q=(\mathbb Z/n\mathbb Z)_T,
\]
with $p$ and $q$ the unit section.  Thus we need only determine
\[
 \alpha:
 (\mathbb Z/n\mathbb Z)_T\otimes(\mathbb Z/n\mathbb Z)_T
 \longrightarrow
 \uTor^{\mathbb Z}_1
 \bigl((\mathbb Z/n\mathbb Z)_T,(\mathbb Z/n\mathbb Z)_T\bigr).
\]
The description in \SGARefLink{sga7.sga7i.eviii.n2-1-7}{2.1.7} shows that this is the isomorphism
\begin{equation}\tag{\SGARefTarget{sga7.sga7i.eviii.n2-1-8-1}{2.1.8.1}}
 (\mathbb Z/n\mathbb Z)_T
 \xrightarrow{\sim}
 \uTor^{\mathbb Z}_1
 \bigl((\mathbb Z/n\mathbb Z)_T,(\mathbb Z/n\mathbb Z)_T\bigr)
\end{equation}
deduced from the canonical flat resolution
\begin{equation}\tag{\SGARefTarget{sga7.sga7i.eviii.n2-1-8-2}{2.1.8.2}}
 0\longrightarrow\mathbb Z_T
 \xrightarrow{n}\mathbb Z_T
 \longrightarrow0
\end{equation}
of the \underline{first} factor
$P=(\mathbb Z/n\mathbb Z)_T$ occurring in $\uTor_1$.



}{E7I-VIII-006}\SGApage{I}{230}
\SGARefTarget{sga7.sga7i.eviii.n2-1-9}{2.1.9}. The definition given in (\SGARefLink{sga7.sga7i.eviii.n2-1-1}{2.1.1}) is not symmetric.  By
interchanging the roles of $P$ and $Q$ in the preceding construction,
one may define an analogous homomorphism
\begin{equation}\tag{\SGARefTarget{sga7.sga7i.eviii.n2-1-9-1}{2.1.9.1}}
 \beta^n_{P,Q}:{}_nP\otimes{}_nQ\longrightarrow\uTor_1(P,Q),
\end{equation}
first reducing to the case $nP=0$ and then using a flat resolution
\underline{of} $Q$.  Equivalently, $\beta$ may be defined in terms of
$\alpha$ by requiring the following diagram to commute:
\begin{equation}\tag{\SGARefTarget{sga7.sga7i.eviii.n2-1-9-2}{2.1.9.2}}
\begin{tikzcd}[column sep=huge,row sep=huge]
 {}_nP\otimes{}_nQ
   \arrow[r,"\sim"]
   \arrow[d,"{\beta^n_{P,Q}}"']\&
 {}_nQ\otimes{}_nP
   \arrow[d,"{\alpha^n_{Q,P}}"]\\
 \uTor_1(P,Q)\arrow[r,"\sim"]\&
 \uTor_1(Q,P).
\end{tikzcd}
\end{equation}
Here the horizontal arrows are the symmetry isomorphisms.  With this
convention, the relation between $\alpha$ and $\beta$ is
\begin{equation}\tag{\SGARefTarget{sga7.sga7i.eviii.n2-1-9-3}{2.1.9.3}}
 \beta^n_{P,Q}=-\alpha^n_{P,Q}.
\end{equation}
In other words:

\medskip
\noindent\underline{Lemma \SGARefTarget{sga7.sga7i.eviii.n2-1-10}{2.1.10}.}
\underline{Let $P$ and $Q$ be two commutative groups and let $n>0$ be
an integer.  Then the following diagram is anticommutative, where
$\alpha$ is defined in (\SGARefLink{sga7.sga7i.eviii.n2-1-1}{2.1.1}):}
\begin{equation}\tag{\SGARefTarget{sga7.sga7i.eviii.n2-1-10-1}{2.1.10.1}}
\begin{tikzcd}[column sep=huge,row sep=huge]
 {}_nP\otimes{}_nQ
   \arrow[r,"\sim"]
   \arrow[d,"{\alpha^n_{P,Q}}"']\&
 {}_nQ\otimes{}_nP
   \arrow[d,"{\alpha^n_{Q,P}}"]\\
 \uTor_1(P,Q)
   \arrow[r,"\sim"]
   \arrow[ur,phantom,"-" description]\&
 \uTor_1(Q,P).
\end{tikzcd}
\end{equation}
\underline{The horizontal arrows are the symmetry isomorphisms.}

\medskip
\noindent\underline{Proof.}
The procedure of \SGARefLink{sga7.sga7i.eviii.n2-1-8}{2.1.8} immediately reduces us to the case
$P=Q=(\mathbb Z/n\mathbb Z)_T$.  We must then prove that the two
isomorphisms (\SGARefLink{sga7.sga7i.eviii.n2-1-8-1}{2.1.8.1}), obtained by resolving respectively the first
and second arguments of $\uTor_1$ by (\SGARefLink{sga7.sga7i.eviii.n2-1-8-2}{2.1.8.2}), differ by a sign
(equivalently, that the symmetry isomorphism on
$\uTor^{\mathbb Z}_1(\mathbb Z/n\mathbb Z,\mathbb Z/n\mathbb Z)$ is
$-\id$).  We may plainly suppose that $T$ is the punctual topos, so
that we work with ordinary commutative groups.  Let $L_\bullet$ be the
complex
\[
 L_\bullet=\bigl(0\longrightarrow\mathbb Z
 \xrightarrow{n}\mathbb Z\longrightarrow0\bigr),
\]
regarded as a flat resolution of
$P=Q=\mathbb Z/n\mathbb Z$.  Use
$L_\bullet\otimes L_\bullet$ to compute
\[
 \Tor_1(P,Q)=\Tor_1(\mathbb Z/n\mathbb Z,\mathbb Z/n\mathbb Z).
\]
The isomorphisms with
$H_1(L_\bullet\otimes\mathbb Z/n\mathbb Z)$, which defines $\alpha$,
and with
$H_1(\mathbb Z/n\mathbb Z\otimes L_\bullet)$, which defines $\beta$,
are induced respectively by the augmentations
$L_\bullet\longrightarrow\mathbb Z/n\mathbb Z$ of the second and
first factors of $L_\bullet\otimes L_\bullet$.  It follows at once
that $\alpha(1,1)$ is represented by the cycle
\[
 -1_0\otimes1_1+1_1\otimes1_0
 \in(L_\bullet\otimes L_\bullet)_1
 =L_0\otimes L_1+L_1\otimes L_0,
\]
where $1_0$ and $1_1$ denote the element $1$ of $L_0$ and $L_1$
respectively, whereas $\beta(1,1)$ is represented by
$1_0\otimes1_1-1_1\otimes1_0$.  This proves \SGARefLink{sga7.sga7i.eviii.n2-1-10}{2.1.10}.

\SGApage{I}{\(\approx\)231}
To finish these generalities on the pairings $\alpha$, we must still
specify the relations among the homomorphisms $\alpha^n_{P,Q}$ as
$n$ varies.  The main result is the following.

\SGAvar{\medskip
\noindent\underline{Lemma \SGARefTarget{sga7.sga7i.eviii.n2-1-11}{2.1.11}.}
\underline{Let $P,Q$ be two commutative groups of $T$, and let $n,n'$
be positive integers such that $n'$ is a multiple of $n$, say
$n'=mn$.  Then the following diagram is commutative:}
\begin{equation}\tag{\SGARefTarget{sga7.sga7i.eviii.n2-1-11-1}{2.1.11.1}}
\begin{tikzcd}[column sep=large,row sep=large]
 \&{}_{n'}P\otimes{}_nQ
   \arrow[ld,"P_m\otimes\id_{{}_nQ}"']
   \arrow[rd,"\id_{{}_{n'}P}\otimes i_m"]\&\\
 {}_nP\otimes{}_nQ
   \arrow[rd,"{\alpha^n_{P,Q}}"']\&\&
 {}_{n'}P\otimes{}_{n'}Q
   \arrow[ld,"{\alpha^{n'}_{P,Q}}"]\\
 \&\uTor_1(P,Q)\&
\end{tikzcd}
\end{equation}
\underline{Here $P_m:{}_{n'}P\longrightarrow{}_nP$ is the homomorphism
induced by $m\cdot\id_P$, and
$i_m:{}_nQ\longrightarrow{}_{n'}Q$ is the inclusion.  Likewise, the
diagram (\SGARefTarget{sga7.sga7i.eviii.n2-1-11-2}{2.11.2}) symmetric to (\SGARefLink{sga7.sga7i.eviii.n2-1-11-1}{2.11.1}) is commutative; it has the
same vertices as (\SGARefLink{sga7.sga7i.eviii.n2-1-11-1}{2.11.1}), except that the upper vertex
$\,{}_{n'}P\otimes{}_nQ$ is replaced by
$\,{}_nP\otimes{}_{n'}Q$.}

}{\medskip
\noindent\underline{Lemma \SGARefTarget{sga7.sga7i.eviii.n2-1-11}{2.1.11}.}
\underline{Let $P,Q$ be two commutative groups of $T$, and let $n,n'$
be positive integers such that $n'$ is a multiple of $n$, say
$n'=mn$.  Then the following diagram is commutative:}
\begin{equation}\tag{\SGARefTarget{sga7.sga7i.eviii.n2-1-11-1}{2.1.11.1}}
\begin{tikzcd}[column sep=large,row sep=large]
 \&{}_{n'}P\otimes{}_nQ
   \arrow[ld,"P_m\otimes\id_{{}_nQ}"']
   \arrow[rd,"\id_{{}_{n'}P}\otimes i_m"]\&\\
 {}_nP\otimes{}_nQ
   \arrow[rd,"{\alpha^n_{P,Q}}"']\&\&
 {}_{n'}P\otimes{}_{n'}Q
   \arrow[ld,"{\alpha^{n'}_{P,Q}}"]\\
 \&\uTor_1(P,Q)\&
\end{tikzcd}
\end{equation}
\underline{Here $P_m:{}_{n'}P\longrightarrow{}_nP$ is the homomorphism
induced by $m\cdot\id_P$, and
$i_m:{}_nQ\longrightarrow{}_{n'}Q$ is the inclusion.  Likewise, the
diagram (\SGARefTarget{sga7.sga7i.eviii.n2-1-11-2}{2.1.11.2}) symmetric to (\SGARefLink{sga7.sga7i.eviii.n2-1-11-1}{2.1.11.1}) is commutative; it has the
same vertices as (\SGARefLink{sga7.sga7i.eviii.n2-1-11-1}{2.1.11.1}), except that the upper vertex
$\,{}_{n'}P\otimes{}_nQ$ is replaced by
$\,{}_nP\otimes{}_{n'}Q$.}

}{E7I-VIII-007}\SGApage{I}{\(\approx\)232}
\medskip
It suffices to prove the first assertion; the second follows by
symmetry in view of \SGARefLink{sga7.sga7i.eviii.n2-1-10}{2.1.10}.  By the procedure of \SGARefLink{sga7.sga7i.eviii.n2-1-8}{2.1.8}, we are reduced
to the punctual topos, where $P$ and $Q$ are ordinary commutative
groups and
\[
 P=\mathbb Z/n'\mathbb Z,\qquad Q=\mathbb Z/n'\mathbb Z.
\]
In fact, for $P$ we use only the hypothesis $n'P=0$.  Using the
resolution
\[
 0\longrightarrow\mathbb Z\xrightarrow{n'}\mathbb Z
 \longrightarrow0
\]
of $Q$, identify $\Tor_1(P,Q)$ with $\,{}_{n'}P=P$.  Likewise, using
\[
 0\longrightarrow\mathbb Z\xrightarrow n\mathbb Z
 \longrightarrow0
\]
to resolve
$\,{}_nQ\simeq\mathbb Z/n\mathbb Z$, identify
$\Tor_1({}_nP,{}_nQ)$ with $\,{}_nP$.  The augmentation in the latter
resolution sends $1_0$ to $m\cdot1_Q$, where $1_Q$ denotes the element
$1$ of $Q=\mathbb Z/n'\mathbb Z$.  Under these identifications the
homomorphism
\[
\begin{tikzcd}[column sep=large,row sep=large]
 \Tor_1({}_nP,{}_nQ)\arrow[r]\arrow[d,no head,"\wr"']\&
 \Tor_1(P,Q)\arrow[d,no head,"\wr"']\\
 {}_nP\&{}_{n'}P=P
\end{tikzcd}
\]
induced by the inclusions
$\,{}_nP\hookrightarrow P$ and $\,{}_nQ\hookrightarrow Q$ is simply
the natural inclusion $\,{}_nP\hookrightarrow{}_{n'}P$.

Consequently, identify the upper vertex of (\SGARefLink{sga7.sga7i.eviii.n2-1-11-1}{2.1.11.1}) with $P/nP$
using the generator $m\cdot1_Q$ of $\,{}_nQ$, which identifies
$\,{}_nQ$ with $\mathbb Z/n\mathbb Z$.  The diagram becomes
\[
\begin{tikzcd}[column sep=large,row sep=large]
 \&P/nP
   \arrow[ld,"x\longmapsto mx"']
   \arrow[rd,"x\longmapsto mx"]\&\\
 {}_nP\arrow[rd,"x\longmapsto x"']\&\&
 P\arrow[ld,"x\longmapsto x"]\\
 \&{}_{n'}P\&
\end{tikzcd}
\]
which is plainly commutative.

\SGApage{I}{233}
\medskip
\noindent\underline{Corollary \SGARefTarget{sga7.sga7i.eviii.n2-1-12}{2.1.12}.}
\underline{With the notation of \SGARefLink{sga7.sga7i.eviii.n2-1-11}{2.1.11}, the following diagram is
commutative:}
\begin{equation}\tag{\SGARefTarget{sga7.sga7i.eviii.n2-1-12-1}{2.1.12.1}}
\begin{tikzcd}[column sep=huge,row sep=huge]
 {}_{n'}P\otimes{}_{n'}Q
   \arrow[r,"m\otimes m"]
   \arrow[d,"{\alpha^{n'}_{P,Q}}"']\&
 {}_nP\otimes{}_nQ
   \arrow[d,"{\alpha^n_{P,Q}}"]\\
 {}_{n'}\uTor_1(P,Q)
   \arrow[r,"m"]\&
 {}_n\uTor_1(P,Q).
\end{tikzcd}
\end{equation}
\underline{In the description of the horizontal arrows, $m$ denotes
the homomorphisms $\,{}_{n'}E\longrightarrow{}_nE$ induced by
$m\cdot\id_E$, for $E=P,Q$, and $\uTor_1(P,Q)$ respectively.}

\medskip
Indeed, by \SGARefLink{sga7.sga7i.eviii.n2-1-11}{2.1.11}, for
$x\in{}_{n'}P(S)$ and $y\in{}_{n'}Q(S)$ one has
\[
 \alpha^n(mx\otimes my)=\alpha^{n'}(x\otimes my).
\]
The right-hand side equals
$m\alpha^{n'}(x\otimes y)$, since
$x\otimes my=m(x\otimes y)$.  This proves the commutativity of
(\SGARefLink{sga7.sga7i.eviii.n2-1-12-1}{2.1.12.1}).



\SGApage{I}{234}
\SGARefTarget{sga7.sga7i.eviii.n2-1-13}{2.1.13}. Let $\ell$ be a prime number.  For every commutative group
$E$, denote by $T_\ell(E)$ the inverse system
\begin{equation}\tag{\SGARefTarget{sga7.sga7i.eviii.n2-1-13-1}{2.1.13.1}}
 T_\ell(E)=({}_{\ell^\nu}E)_{\nu\geq0},
\end{equation}
which plainly depends functorially on $E$.  The compatibilities of
\SGARefLink{sga7.sga7i.eviii.n2-1-12}{2.1.12}, for $n,n'$ of the form $\ell^\nu,\ell^{\nu'}$ with
$\nu'\geq\nu$, show that the homomorphisms
$\alpha^{\ell^\nu}_{P,Q}$, as $\nu$ varies, define a homomorphism of
inverse systems, functorial in $P,Q$:
\begin{equation}\tag{\SGARefTarget{sga7.sga7i.eviii.n2-1-13-2}{2.1.13.2}}
 \alpha^{(\ell)}_{P,Q}:
 T_\ell(P)\otimes T_\ell(Q)
 \longrightarrow T_\ell\bigl(\uTor_1(P,Q)\bigr).
\end{equation}
Knowledge of these homomorphisms for every prime $\ell$ is plainly
equivalent to knowledge of the $\alpha^n_{P,Q}$ for every integer
$n$.  The antisymmetry relations of \SGARefLink{sga7.sga7i.eviii.n2-1-10}{2.1.10} then become the analogous
$\ell$-adic antisymmetry relations, for every $\ell$, expressed by
the anticommutativity of the diagrams
\begin{equation}\tag{\SGARefTarget{sga7.sga7i.eviii.n2-1-13-3}{2.1.13.3}}
\begin{tikzcd}[column sep=huge,row sep=huge]
 T_\ell(P)\otimes T_\ell(Q)
   \arrow[r]
   \arrow[d,"{\alpha^{(\ell)}_{P,Q}}"']
   \arrow[rd,phantom,"-" description]\&
 T_\ell(Q)\otimes T_\ell(P)
   \arrow[d,"{\alpha^{(\ell)}_{Q,P}}"]\\
 T_\ell\bigl(\uTor_1(P,Q)\bigr)
   \arrow[r]\&
 T_\ell\bigl(\uTor_1(Q,P)\bigr).
\end{tikzcd}
\end{equation}
Here the horizontal arrows are the symmetry arrows.

\medskip
\noindent\underline{Remark \SGARefTarget{sga7.sga7i.eviii.n2-1-14}{2.1.14}.}
The construction of the homomorphism (\SGARefLink{sga7.sga7i.eviii.n2-1-1}{2.1.1}), and the results \SGARefLink{sga7.sga7i.eviii.n2-1-10}{2.1.10}
and \SGARefLink{sga7.sga7i.eviii.n2-1-11}{2.1.11} concerning it, may be generalized to the situation in
which one is given a commutative ring $\underline A$ on $T$ and two
ideals $\underline a$ and $\underline b$ of $\underline A$.  In the
case treated above, both ideals are the ideal
$n\mathbb Z_T$ of the constant ring $\mathbb Z_T$ on $T$.  One then
defines a homomorphism, functorial in the $\underline A$-modules
$P,Q$,
\begin{equation}\tag{\SGARefTarget{sga7.sga7i.eviii.n2-1-14-1}{2.1.14.1}}
 {}_{\underline a}P\otimes{}_{\underline b}Q
 \xrightarrow{\ \alpha\ }
 \uHom_{\underline A}
 \left(
   \frac{\underline a\cap\underline b}
        {\underline a\mathbin{\cdot}\underline b},
   \uTor^{\underline A}_1(P,Q)
 \right),
\end{equation}
as the composite of the homomorphisms
\[
\begin{split}
 {}_{\underline a}P\otimes{}_{\underline b}Q
 &\longrightarrow
 \uHom(\underline A/\underline a,P)
 \otimes
 \uHom(\underline A/\underline b,Q)\\
 &\longrightarrow
 \uHom\!\left(
   \uTor_1(\underline A/\underline a,\underline A/\underline b),
   \uTor_1(P,Q)
 \right),
\end{split}
\]
using the well-known canonical isomorphism
\begin{equation}\tag{\SGARefTarget{sga7.sga7i.eviii.n2-1-14-2}{2.1.14.2}}
 \uTor_1(\underline A/\underline a,\underline A/\underline b)
 \simeq
 \frac{\underline a\cap\underline b}
      {\underline a\mathbin{\cdot}\underline b}.
\end{equation}
In fact, the sign in \SGARefLink{sga7.sga7i.eviii.n2-1-10}{2.1.10} is ``hidden'' in the choice of this
isomorphism, according as one works with the resolution
$0\longrightarrow\underline a\longrightarrow\underline A
\longrightarrow0$ of $\underline A/\underline a$, or with the
resolution $0\longrightarrow\underline b\longrightarrow\underline A
\longrightarrow0$ of $\underline A/\underline b$.  Here it becomes
the anticommutativity of the diagram
\begin{equation}\tag{\SGARefTarget{sga7.sga7i.eviii.n2-1-14-3}{2.1.14.3}}
\begin{tikzcd}[column sep=huge,row sep=huge]
 {}_{\underline a}P\otimes{}_{\underline b}Q
   \arrow[r]
   \arrow[d]
   \arrow[rd,phantom,"-1" description]\&
 {}_{\underline b}Q\otimes{}_{\underline a}P
   \arrow[d]\\
 \uHom\!\left(
   \dfrac{\underline a\cap\underline b}
         {\underline a\mathbin{\cdot}\underline b},
   \uTor_1(P,Q)
 \right)
   \arrow[r]\&
 \uHom\!\left(
   \dfrac{\underline a\cap\underline b}
         {\underline a\mathbin{\cdot}\underline b},
   \uTor_1(Q,P)
 \right).
\end{tikzcd}
\end{equation}
The horizontal arrows are the symmetry arrows.  When
$\underline a=\underline b$, this anticommutativity may also be
expressed by saying that the symmetry automorphism on
$\uTor_1(\underline A/\underline a,\underline A/\underline a)$ is
$-\id$.

\SGApage{I}{\(\approx\)235}
\SGAvar{In the special case where
$\underline a=\underline b=n\underline A$, with $n$ a section of
$\underline A$, multiplication by $n$ induces an epimorphism
\[
 \underline A/\underline a
 \longrightarrow
 \frac{\underline a\cap\underline a}
      {\underline a\mathbin{\cdot}\underline a}
 =\underline a/\underline a^2,
\]
and hence a functorial injective isomorphism, in the
$\underline A$-module $T$,
\[
 \uHom(\underline a/\underline a^2,T)
 \longrightarrow{}_{\underline a}T.
\]
Thus in this case the homomorphism (\SGARefLink{sga7.sga7i.eviii.n2-1-14-1}{2.1.14.1}) may be interpreted as a
homomorphism
\begin{equation}\tag{\SGARefTarget{sga7.sga7i.eviii.n2-1-14-4}{2.1.14.4}}
 {}_nP\otimes{}_nQ
 \longrightarrow{}_n\uTor_1(P,Q),
\end{equation}
which is closer in form to (\SGARefLink{sga7.sga7i.eviii.n2-1-1}{2.1.1}).

}{In the special case where
$\underline a=\underline b=n\underline A$, with $n$ a section of
$\underline A$, multiplication by $n$ induces an epimorphism
\[
 \underline A/\underline a
 \longrightarrow
 \frac{\underline a\cap\underline a}
      {\underline a\mathbin{\cdot}\underline a}
 =\underline a/\underline a^2,
\]
and hence a functorial injective homomorphism, in the
$\underline A$-module $T$,
\[
 \uHom(\underline a/\underline a^2,T)
 \longrightarrow{}_{\underline a}T.
\]
Thus in this case the homomorphism (\SGARefLink{sga7.sga7i.eviii.n2-1-14-1}{2.1.14.1}) may be interpreted as a
homomorphism
\begin{equation}\tag{\SGARefTarget{sga7.sga7i.eviii.n2-1-14-4}{2.1.14.4}}
 {}_nP\otimes{}_nQ
 \longrightarrow{}_n\uTor_1(P,Q),
\end{equation}
which is closer in form to (\SGARefLink{sga7.sga7i.eviii.n2-1-1}{2.1.1}).

}{E7I-VIII-008}\SGApage{I}{\(\approx\)236}
Continue to suppose that
$\underline a=\underline b=n\underline A$, and that $n$ is a regular
section of $\underline A$.  The spectral sequence (\SGARefLink{sga7.sga7i.eviii.n2-1-4}{2.1.4}) still makes
sense in the present setting and yields the following exact sequence
of low-degree terms, generalizing (\SGARefLink{sga7.sga7i.eviii.n2-1-5}{2.1.5}), where $Q$ is now an
$\underline A$-module annihilated by $\underline a$:
\begin{equation}\tag{\SGARefTarget{sga7.sga7i.eviii.n2-1-14-5}{2.1.14.5}}
\begin{split}
 \uTor^{\underline A}_2(P,Q)
 &\longrightarrow
 \uTor^{\underline A/n\underline A}_2(P_n,Q)
 \longrightarrow{}_nP\otimes Q
 \xrightarrow{\ \alpha\ }
 \uTor^{\underline A}_1(P,Q)\\
 &\longrightarrow
 \uTor^{\underline A/n\underline A}_1(P_n,Q)
 \longrightarrow0.
\end{split}
\end{equation}

\SGAvar{Return now to arbitrary $\underline a$ and $\underline b$, and let
$p$ and $q$ be two sections of $\underline A$.  We then obtain the
commutative diagram below; (\SGARefLink{sga7.sga7i.eviii.n2-1-11-1}{2.1.11.1}) is essentially the special case
obtained by taking $q=1$:
\begin{equation}\tag{\SGARefTarget{sga7.sga7i.eviii.n2-1-14-6}{2.1.14.6}}
{\footnotesize
\begin{tikzcd}[row sep=large,column sep=tiny]
 \&{}_{p\underline a}P\otimes{}_{q\underline b}Q
   \arrow[dl,"p\otimes q"']
   \arrow[dr,"i\otimes j"]\&\\
 {}_{\underline a}P\otimes{}_{\underline b}Q
   \arrow[d,"\alpha"']\&\&
 {}_{pq\underline a}P\otimes{}_{pq\underline b}Q
   \arrow[d,"\alpha"]\\
 \uHom\!\left(
   \dfrac{\underline a\cap\underline b}
         {\underline a\mathbin{\cdot}\underline b},
   \uTor_1(P,Q)
 \right)
   \arrow[dr,hook,"\mathrm{incl.}"']\&\&
 \uHom\!\left(
   \dfrac{pq\underline a\cap pq\underline b}
         {pq\underline a\mathbin{\cdot}pq\underline b},
   \uTor_1(P,Q)
 \right)
   \arrow[dl,hook,"(pq)^*"]\\
 \&\uHom\!\left(
   \dfrac{\underline a\cap\underline b}
         {pq\underline a\mathbin{\cdot}\underline b},
   \uTor_1(P,Q)
 \right)\&
\end{tikzcd}}
\end{equation}
Here
\[
 i:{}_{p\underline a}P\longrightarrow{}_{pq\underline a}P,
 \qquad
 j:{}_{q\underline b}Q\longrightarrow{}_{pq\underline b}Q
\]
are the canonical inclusions, and $(pq)^*$ denotes the injective
homomorphism obtained, by passing to $\uHom$, from the surjective
homomorphism
\[
 \frac{\underline a\cap\underline b}
      {pq\underline a\mathbin{\cdot}\underline b}
 \longrightarrow
 \frac{pq\underline a\cap\underline b}
      {pq\underline a\mathbin{\cdot}pq\underline b}
\]
induced by multiplication by $pq$.  When
$\underline a=\underline b=n\underline A$, this diagram may be
rewritten in a simpler form, given in (\SGARefLink{sga7.sga7i.eviii.n2-1-11-1}{2.1.11.1}) when $q=1$.



}{Return now to arbitrary $\underline a$ and $\underline b$, and let
$p$ and $q$ be two sections of $\underline A$.  We then obtain the
commutative diagram below; (\SGARefLink{sga7.sga7i.eviii.n2-1-11-1}{2.1.11.1}) is essentially the special case
obtained by taking $q=1$:
\begin{equation}\tag{\SGARefTarget{sga7.sga7i.eviii.n2-1-14-6}{2.1.14.6}}
{\footnotesize
\begin{tikzcd}[row sep=large,column sep=tiny]
 \&{}_{p\underline a}P\otimes{}_{q\underline b}Q
   \arrow[dl,"p\otimes q"']
   \arrow[dr,"i\otimes j"]\&\\
 {}_{\underline a}P\otimes{}_{\underline b}Q
   \arrow[d,"\alpha"']\&\&
 {}_{pq\underline a}P\otimes{}_{pq\underline b}Q
   \arrow[d,"\alpha"]\\
 \uHom\!\left(
   \dfrac{\underline a\cap\underline b}
         {\underline a\mathbin{\cdot}\underline b},
   \uTor_1(P,Q)
 \right)
   \arrow[dr,hook,"\mathrm{incl.}"']\&\&
 \uHom\!\left(
   \dfrac{pq\underline a\cap pq\underline b}
         {pq\underline a\mathbin{\cdot}pq\underline b},
   \uTor_1(P,Q)
 \right)
   \arrow[dl,hook,"(pq)^*"]\\
 \&\uHom\!\left(
   \dfrac{\underline a\cap\underline b}
         {pq\underline a\mathbin{\cdot}\underline b},
   \uTor_1(P,Q)
 \right)\&
\end{tikzcd}}
\end{equation}
Here
\[
 i:{}_{p\underline a}P\longrightarrow{}_{pq\underline a}P,
 \qquad
 j:{}_{q\underline b}Q\longrightarrow{}_{pq\underline b}Q
\]
are the canonical inclusions, and $(pq)^*$ denotes the injective
homomorphism obtained, by passing to $\uHom$, from the surjective
homomorphism
\[
 \frac{\underline a\cap\underline b}
      {pq\underline a\mathbin{\cdot}\underline b}
 \longrightarrow
 \frac{pq\underline a\cap pq\underline b}
      {pq\underline a\mathbin{\cdot}pq\underline b}
\]
induced by multiplication by $pq$.  When
$\underline a=\underline b=n\underline A$, this diagram may be
rewritten in a simpler form, given in (\SGARefLink{sga7.sga7i.eviii.n2-1-11-1}{2.1.11.1}) when $q=1$.



}{E7I-VIII-009}\SGApage{I}{237}
\SGARefTarget{sga7.sga7i.eviii.n2-2}{2.2}. Consider now, in addition to $P$ and $Q$, a commutative group
$G$.  In (\SGARefLink{sga7.sga7i.eviii.n1-1-3}{1.1.3}) we defined a canonical homomorphism, functorial in
$P,Q$, and $G$,
\begin{equation}\tag{\SGARefTarget{sga7.sga7i.eviii.n2-2-1}{2.2.1}}
 \mathrm{Biext}^1(P,Q;G)
 \longrightarrow
 \operatorname{Hom}\bigl(\operatorname{Tor}_1(P,Q),G\bigr).
\end{equation}
On the other hand, if $n$ is a positive integer (respectively, if
$\ell$ is a prime number), the homomorphism $\alpha^n_{P,Q}$ of
(\SGARefLink{sga7.sga7i.eviii.n2-1-1}{2.1.1}) (respectively, the homomorphism $\alpha^{(\ell)}_{P,Q}$ of
(\SGARefLink{sga7.sga7i.eviii.n2-1-13-2}{2.1.13.2})) gives a homomorphism
\[
 \operatorname{Hom}\bigl(\operatorname{Tor}_1(P,Q),G\bigr)
 \longrightarrow
 \operatorname{Hom}({}_nP\otimes{}_nQ,G)
\]
(respectively, a homomorphism
\[
 \operatorname{Hom}\bigl(\operatorname{Tor}_1(P,Q),G\bigr)
 \longrightarrow
 \operatorname{Hom}\bigl(
   T_\ell(P)\otimes T_\ell(Q),T_\ell(G)
 \bigr)).
\]
These last homomorphisms are also functorial in $P,Q,G$.
Composing them with (\SGARefLink{sga7.sga7i.eviii.n2-2-1}{2.2.1}), we obtain canonical homomorphisms,
functorial in $P,Q,G$,
\begin{equation}\tag{\SGARefTarget{sga7.sga7i.eviii.n2-2-2}{2.2.2}}
 \varphi^n=\varphi^n_{P,Q}:
 \mathrm{Biext}^1(P,Q;G)
 \longrightarrow
 \operatorname{Hom}({}_nP\otimes{}_nQ,G),
\end{equation}
\begin{equation}\tag{\SGARefTarget{sga7.sga7i.eviii.n2-2-3}{2.2.3}}
 \varphi^{(\ell)}=\varphi^{(\ell)}_{P,Q}:
 \mathrm{Biext}^1(P,Q;G)
 \longrightarrow
 \operatorname{Hom}\bigl(
   T_\ell(P)\otimes T_\ell(Q),T_\ell(G)
 \bigr).
\end{equation}
For
\begin{equation}\tag{\SGARefTarget{sga7.sga7i.eviii.n2-2-4}{2.2.4}}
 \xi\in\mathrm{Biext}^1(P,Q;G),
\end{equation}
denote its image under (\SGARefLink{sga7.sga7i.eviii.n2-2-2}{2.2.2}) (respectively, under (\SGARefLink{sga7.sga7i.eviii.n2-2-3}{2.2.3})) by
$\varphi^n_\xi$ (respectively, by $\varphi^{(\ell)}_\xi$):
\begin{equation}\tag{\SGARefTarget{sga7.sga7i.eviii.n2-2-5}{2.2.5}}
 \varphi^n_\xi:{}_nP\otimes{}_nQ
 \longrightarrow G
 \qquad\text{(or, equivalently, }{}_nG\text{)},
\end{equation}
\begin{equation}\tag{\SGARefTarget{sga7.sga7i.eviii.n2-2-6}{2.2.6}}
 \varphi^{(\ell)}_\xi:
 T_\ell(P)\otimes T_\ell(Q)
 \longrightarrow T_\ell(G).
\end{equation}

\SGApage{I}{238}
\SGAvar{\SGARefTarget{sga7.sga7i.eviii.n2-2-7}{2.2.7}. Plainly, knowledge of $\varphi^{(\ell)}_\xi$ is equivalent
to knowledge of the $\varphi^n_\xi$ for all $n=\ell^\nu$, with
$\nu\geq0$.  For fixed $n$, knowledge of $\varphi^n_\xi$ is
equivalent to knowledge, on the image of
\[
 \varphi^n_{P,Q}:{}_nP\otimes{}_nQ\longrightarrow\uTor_1(P,Q;G),
\]
of the homomorphism
\begin{equation}\tag{\SGARefTarget{sga7.sga7i.eviii.n2-2-8}{2.2.8}}
 \uTor_1(P,Q)\longrightarrow G
\end{equation}
which is the image of $\xi$ under (\SGARefLink{sga7.sga7i.eviii.n2-2-1}{2.2.1}).

}{\SGARefTarget{sga7.sga7i.eviii.n2-2-7}{2.2.7}. Plainly, knowledge of $\varphi^{(\ell)}_\xi$ is equivalent
to knowledge of the $\varphi^n_\xi$ for all $n=\ell^\nu$, with
$\nu\geq0$.  For fixed $n$, knowledge of $\varphi^n_\xi$ is
equivalent to knowledge, on the image of
\[
 \alpha^n_{P,Q}:{}_nP\otimes{}_nQ\longrightarrow\uTor_1(P,Q),
\]
of the homomorphism
\begin{equation}\tag{\SGARefTarget{sga7.sga7i.eviii.n2-2-8}{2.2.8}}
 \uTor_1(P,Q)\longrightarrow G
\end{equation}
which is the image of $\xi$ under (\SGARefLink{sga7.sga7i.eviii.n2-2-1}{2.2.1}).

}{E7I-VIII-010}Observe that if $P=nP$ and $Q=nQ$, then the preceding image is
$\,{}_n\uTor_1(P,Q)$.  This follows readily by considering the exact
sequence (\SGARefLink{sga7.sga7i.eviii.n1-3-5}{1.3.5}) and the analogous exact sequence for $Q$.
Furthermore, if $P$ and $Q$ are divisible, that is, if $nP=P$ and
$nQ=Q$ for every positive integer $n$, one readily checks that
$\uTor_1(P,Q)$ is torsion, that is,
\[
 \uTor_1(P,Q)=\varinjlim_n{}_{n}\uTor_1(P,Q).
\]
Thus, in this case, knowledge of the $\varphi^n_\xi$ for all $n$ is
equivalent to knowledge of (\SGARefLink{sga7.sga7i.eviii.n2-2-8}{2.2.8}), whereas knowledge of
$\varphi^{(\ell)}_\xi$ for a given prime $\ell$ is equivalent to
knowledge of the restriction of (\SGARefLink{sga7.sga7i.eviii.n2-2-8}{2.2.8}) to the $\ell$-primary
component of $\uTor_1(P,Q)$.

\SGAvar{We now make explicit the symmetry properties of the pairings (\SGARefLink{sga7.sga7i.eviii.n2-2-5}{2.2.4})
and (\SGARefLink{sga7.sga7i.eviii.n2-2-6}{2.2.6}) associated with biextensions.  First observe that there is
a commutative diagram
\begin{equation}\tag{\SGARefTarget{sga7.sga7i.eviii.n2-2-9}{2.2.9}}
\begin{tikzcd}[column sep=large,row sep=huge]
 \mathrm{Biext}^1(P,Q;G)
   \arrow[r,"\sim"]
   \arrow[d]\&
 \mathrm{Biext}^1(Q,P;G)
   \arrow[d]\\
 \operatorname{Hom}\bigl(\uTor_1(P,Q),G\bigr)
   \arrow[r,"\sim"]\&
 \operatorname{Hom}\bigl(\uTor_1(Q,P),G\bigr),
\end{tikzcd}
\end{equation}
where the horizontal arrows are the symmetry arrows (the first of
which was defined in VII \SGARefLink{sga7.sga7i.evii.n2-7}{2.7}), and the vertical arrows are (\SGARefLink{sga7.sga7i.eviii.n2-2-1}{2.2.1}).
The commutativity follows from that of the two constituent squares in
the fuller diagram
\begin{equation}\tag{\SGARefTarget{sga7.sga7i.eviii.n2-2-10}{2.2.10}}
\begin{tikzcd}[column sep=large,row sep=large]
 \mathrm{Biext}^1(P,Q;G)
   \arrow[r,"\sim"]
   \arrow[d]\&
 \mathrm{Biext}^1(Q,P;G)
   \arrow[d]\\
 \operatorname{Ext}^1\!\left(
   P\mathbin{\overset{\mathbb L}{\otimes}}Q,G
 \right)
   \arrow[r,"\sim"]
   \arrow[d]\&
 \operatorname{Ext}^1\!\left(
   Q\mathbin{\overset{\mathbb L}{\otimes}}P,G
 \right)
   \arrow[d]\\
 \operatorname{Hom}\bigl(\uTor_1(P,Q),G\bigr)
   \arrow[r,"\sim"]\&
 \operatorname{Hom}\bigl(\uTor_1(Q,P),G\bigr).
\end{tikzcd}
\end{equation}
The first vertical arrow in each column is the canonical isomorphism
of VII \SGARefLink{sga7.sga7i.evii.n3-6-5}{3.6.5}, and the second is the evident homomorphism.  The second
square in (\SGARefLink{sga7.sga7i.eviii.n2-2-10}{2.2.10}) is trivially commutative: by definition, the
symmetry isomorphism
\[
 P\mathbin{\overset{\mathbb L}{\otimes}}Q
 \xrightarrow{\ \sim\ }
 Q\mathbin{\overset{\mathbb L}{\otimes}}P
\]
induces the symmetry isomorphism
$\uTor_1(P,Q)\simeq\uTor_1(Q,P)$, and, for a variable complex
$L_\bullet$, the canonical isomorphism
\[
 \operatorname{Ext}^1(L_\bullet,G)
 \simeq
 \operatorname{Hom}\bigl(H_1(L_\bullet),G\bigr)
\]
is functorial in $L_\bullet$.

}{We now make explicit the symmetry properties of the pairings (\SGARefLink{sga7.sga7i.eviii.n2-2-5}{2.2.5})
and (\SGARefLink{sga7.sga7i.eviii.n2-2-6}{2.2.6}) associated with biextensions.  First observe that there is
a commutative diagram
\begin{equation}\tag{\SGARefTarget{sga7.sga7i.eviii.n2-2-9}{2.2.9}}
\begin{tikzcd}[column sep=large,row sep=huge]
 \mathrm{Biext}^1(P,Q;G)
   \arrow[r,"\sim"]
   \arrow[d]\&
 \mathrm{Biext}^1(Q,P;G)
   \arrow[d]\\
 \operatorname{Hom}\bigl(\uTor_1(P,Q),G\bigr)
   \arrow[r,"\sim"]\&
 \operatorname{Hom}\bigl(\uTor_1(Q,P),G\bigr),
\end{tikzcd}
\end{equation}
where the horizontal arrows are the symmetry arrows (the first of
which was defined in VII \SGARefLink{sga7.sga7i.evii.n2-7}{2.7}), and the vertical arrows are (\SGARefLink{sga7.sga7i.eviii.n2-2-1}{2.2.1}).
The commutativity follows from that of the two constituent squares in
the fuller diagram
\begin{equation}\tag{\SGARefTarget{sga7.sga7i.eviii.n2-2-10}{2.2.10}}
\begin{tikzcd}[column sep=large,row sep=large]
 \mathrm{Biext}^1(P,Q;G)
   \arrow[r,"\sim"]
   \arrow[d]\&
 \mathrm{Biext}^1(Q,P;G)
   \arrow[d]\\
 \operatorname{Ext}^1\!\left(
   P\mathbin{\overset{\mathbb L}{\otimes}}Q,G
 \right)
   \arrow[r,"\sim"]
   \arrow[d]\&
 \operatorname{Ext}^1\!\left(
   Q\mathbin{\overset{\mathbb L}{\otimes}}P,G
 \right)
   \arrow[d]\\
 \operatorname{Hom}\bigl(\uTor_1(P,Q),G\bigr)
   \arrow[r,"\sim"]\&
 \operatorname{Hom}\bigl(\uTor_1(Q,P),G\bigr).
\end{tikzcd}
\end{equation}
The first vertical arrow in each column is the canonical isomorphism
of VII \SGARefLink{sga7.sga7i.evii.n3-6-5}{3.6.5}, and the second is the evident homomorphism.  The second
square in (\SGARefLink{sga7.sga7i.eviii.n2-2-10}{2.2.10}) is trivially commutative: by definition, the
symmetry isomorphism
\[
 P\mathbin{\overset{\mathbb L}{\otimes}}Q
 \xrightarrow{\ \sim\ }
 Q\mathbin{\overset{\mathbb L}{\otimes}}P
\]
induces the symmetry isomorphism
$\uTor_1(P,Q)\simeq\uTor_1(Q,P)$, and, for a variable complex
$L_\bullet$, the canonical isomorphism
\[
 \operatorname{Ext}^1(L_\bullet,G)
 \simeq
 \operatorname{Hom}\bigl(H_1(L_\bullet),G\bigr)
\]
is functorial in $L_\bullet$.

}{E7I-VIII-011}\SGApage{I}{\(\approx\)239}
As for the commutativity of the first square of (\SGARefLink{sga7.sga7i.eviii.n2-2-10}{2.2.10}), it follows
by making explicit simultaneously the symmetry isomorphisms and the
vertical isomorphisms defined in VII \SGARefLink{sga7.sga7i.evii.n3-6-5}{3.6.5}.  We leave the details of
the verification to the reader.  One uses the fact that the symmetry
isomorphism
\[
 L_\bullet(P)\otimes L_\bullet(Q)
 \simeq
 L_\bullet(Q)\otimes L_\bullet(P)
\]
induces in degree $1$ the isomorphism
\[
 L_0(P)\otimes L_1(Q)+L_1(P)\otimes L_0(Q)
 \simeq
 L_0(Q)\otimes L_1(P)+L_1(Q)\otimes L_0(P),
\]
which, on corresponding summands, is given by the ordinary symmetry
isomorphism of the underlying ungraded modules (and hence introduces
no sign).

In view of the commutativity of (\SGARefLink{sga7.sga7i.eviii.n2-2-9}{2.2.9}), the anticommutativity of
(\SGARefLink{sga7.sga7i.eviii.n2-1-10-1}{2.1.10.1}), in terms of the $\varphi^n$ or
$\varphi^{(\ell)}$, becomes the following relation.

\SGApage{I}{240}
\SGAvar{\medskip
\noindent\underline{Proposition \SGARefTarget{sga7.sga7i.eviii.n2-2-11}{2.2.11}.}
\underline{The following diagrams are anticommutative:}
\begin{equation}\tag{\SGARefTarget{sga7.sga7i.eviii.n2-2-11-1}{2.2.11.1}}
\begin{tikzcd}[row sep=large,column sep=huge]
 \mathrm{Biext}^1(P,Q;G)
   \arrow[rr,"\sim"]
   \arrow[dd,"{\varphi^n_{P,Q}}"']\&\&
 \mathrm{Biext}^1(Q,P;G)
   \arrow[dd,"{\varphi^n_{Q,P}}"]\\
 \&- \&\\
 \operatorname{Hom}({}_nP\otimes{}_nQ,G)
   \arrow[rr,"\sim"]\&\&
 \operatorname{Hom}({}_nQ\otimes{}_nP,G),
\end{tikzcd}
\end{equation}
\begin{equation}\tag{\SGARefTarget{sga7.sga7i.eviii.n2-2-11-2}{2.2.11.2}}
\begin{tikzcd}[row sep=large,column sep=huge]
 \mathrm{Biext}^1(P,Q;G)
   \arrow[rr,"\sim"]
   \arrow[dd,"{\varphi^{(\ell)}_{P,Q}}"']\&\&
 \mathrm{Biext}^1(Q,P;G)
   \arrow[dd,"{\varphi^{(\ell)}_{Q,P}}"]\\
 \&- \&\\
 \operatorname{Hom}\bigl(
   T_\ell(P)\otimes T_\ell(Q),T_\ell(G)
   \arrow[rr,"\sim"]\&\&
 \operatorname{Hom}\bigl(
   T_\ell(Q)\otimes T_\ell(P),T_\ell(G)
 \bigr),
\end{tikzcd}
\end{equation}
\underline{where the horizontal arrows are the symmetry arrows.}

}{\medskip
\noindent\underline{Proposition \SGARefTarget{sga7.sga7i.eviii.n2-2-11}{2.2.11}.}
\underline{The following diagrams are anticommutative:}
\begin{equation}\tag{\SGARefTarget{sga7.sga7i.eviii.n2-2-11-1}{2.2.11.1}}
\begin{tikzcd}[row sep=large,column sep=huge]
 \mathrm{Biext}^1(P,Q;G)
   \arrow[rr,"\sim"]
   \arrow[dd,"{\varphi^n_{P,Q}}"']\&\&
 \mathrm{Biext}^1(Q,P;G)
   \arrow[dd,"{\varphi^n_{Q,P}}"]\\
 \&- \&\\
 \operatorname{Hom}({}_nP\otimes{}_nQ,G)
   \arrow[rr,"\sim"]\&\&
 \operatorname{Hom}({}_nQ\otimes{}_nP,G),
\end{tikzcd}
\end{equation}
\begin{equation}\tag{\SGARefTarget{sga7.sga7i.eviii.n2-2-11-2}{2.2.11.2}}
\begin{tikzcd}[row sep=large,column sep=huge]
 \mathrm{Biext}^1(P,Q;G)
   \arrow[rr,"\sim"]
   \arrow[dd,"{\varphi^{(\ell)}_{P,Q}}"']\&\&
 \mathrm{Biext}^1(Q,P;G)
   \arrow[dd,"{\varphi^{(\ell)}_{Q,P}}"]\\
 \&- \&\\
 \operatorname{Hom}\bigl(
   T_\ell(P)\otimes T_\ell(Q),T_\ell(G)
 \bigr)
   \arrow[rr,"\sim"]\&\&
 \operatorname{Hom}\bigl(
   T_\ell(Q)\otimes T_\ell(P),T_\ell(G)
 \bigr),
\end{tikzcd}
\end{equation}
\underline{where the horizontal arrows are the symmetry arrows.}

}{E7I-VIII-012}We record the important corollary.

\medskip
\noindent\underline{Corollary \SGARefTarget{sga7.sga7i.eviii.n2-2-12}{2.2.12}.}
\underline{Suppose that $P=Q$, and suppose that $\xi$ is the class of
a symmetric (respectively, antisymmetric) biextension of $(P,P)$ by
$G$; that is, suppose ${}^s\xi=\xi$ (respectively,
${}^s\xi=-\xi$), where ${}^s\xi$ denotes the image of $\xi$ under
symmetry.  Then the bilinear maps $\varphi^n_\xi$, on
${}_nP\times{}_nP$, and $\varphi^{(\ell)}_\xi$, on
$T_\ell(P)\times T_\ell(P)$, are antisymmetric (respectively,
symmetric).}

\SGApage{I}{241}
\medskip
\noindent\underline{Remark \SGARefTarget{sga7.sga7i.eviii.n2-2-13}{2.2.13}.}
Suppose that $\xi$ is symmetric.  If $n$ (respectively, $\ell$) is
odd, then \SGARefLink{sga7.sga7i.eviii.n2-2-12}{2.2.12} shows that $\varphi^n_\xi$ (respectively,
$\varphi^{(\ell)}_\xi$) is in fact an alternating form.  This
conclusion is not valid in general without the restriction that $n$
(respectively, $\ell$) be odd, as is already seen when $P$ is the
ordinary group $\mathbb Z/2\mathbb Z$ and
$G=\mathbb Z/2\mathbb Z$.

Suppose, however, that $2P=P$.  This is the case in particular when
$P$ is an abelian scheme, working as usual in the fppf or fpqc
topology.  The preceding forms are then still alternating.  It is
enough to prove that $\varphi^{2^\nu}_\xi(x,x)=0$.  Since $2P=P$, a
given element $x\in{}_{2^\nu}P(S)$ may be written as $x=2y$, with
$y\in{}_{2^{\nu+1}}P(S)$.  Hence
\[
 \varphi^{2^\nu}_\xi(x,x)
 =2\varphi^{2^{\nu+1}}_\xi(y,y)=0,
\]
because $\varphi^{2^{\nu+1}}_\xi$ is antisymmetric.



\SGARefTarget{sga7.sga7i.eviii.n2-3}{2.3}. \underline{Relation with ``Cartier duality''.}
The description given in \SGARefLink{sga7.sga7i.eviii.n2-2}{2.2} of the pairings defined by a
biextension, through the isomorphism theorem VII \SGARefLink{sga7.sga7i.evii.n3-6-5}{3.6.5} and the
homomorphism (\SGARefLink{sga7.sga7i.eviii.n2-2-1}{2.2.1}), is plainly not the definition formerly used for
abelian schemes, at a time when the very notion of biextension did not
yet exist.  That definition proceeds through Cartier duality, and we
now give its relation to the present one.

Consider again the homomorphism
\begin{equation}\tag{\SGARefTarget{sga7.sga7i.eviii.n2-3-1}{2.3.1}}
 \mathrm{Biext}^1(P,Q;G)
 \longrightarrow
 \operatorname{Hom}\bigl(P,\uExt^1(Q,G)\bigr)
\end{equation}
arising from the exact sequence (\SGARefLink{sga7.sga7i.eviii.n1-4}{1.4}).  It is in fact an isomorphism
when
\begin{equation}\tag{\SGARefTarget{sga7.sga7i.eviii.n2-3-2}{2.3.2}}
 \uHom(Q,G)=0.
\end{equation}
This condition holds in particular when $Q$ is an abelian scheme over
a scheme $S$ and $G$ is the multiplicative group $\underline G_m$
(VII \SGARefLink{sga7.sga7i.evii.n1-3-8}{1.3.8}).

\SGApage{I}{\(\approx\)242}
There is also an evident ``restriction to $nQ$'' homomorphism
\[
 \uExt^1(Q,G)\longrightarrow\uExt^1(nQ,G),
\]
and hence a composite homomorphism
\begin{equation}\tag{\SGARefTarget{sga7.sga7i.eviii.n2-3-4}{2.3.4}}
\begin{split}
 \operatorname{Hom}\bigl(P,\uExt^1(Q,G)\bigr)
 &\longrightarrow
 \operatorname{Hom}\bigl(P,\uExt^1(nQ,G)\bigr)\\
 &\longrightarrow
 \operatorname{Hom}\bigl({}_nP,{}_n\uExt^1(nQ,G)\bigr),
\end{split}
\end{equation}
where the second arrow is restriction to $\,{}_nP$.  Composing with
(\SGARefLink{sga7.sga7i.eviii.n2-3-1}{2.3.1}), we obtain a homomorphism
\begin{equation}\tag{\SGARefTarget{sga7.sga7i.eviii.n2-3-5}{2.3.5}}
 \mathrm{Biext}^1(P,Q;G)
 \longrightarrow
 \operatorname{Hom}\bigl({}_nP,{}_n\uExt^1(nQ,G)\bigr).
\end{equation}

On the other hand, the exact sequence
\begin{equation}\tag{\SGARefTarget{sga7.sga7i.eviii.n2-3-6}{2.3.6}}
 0\longrightarrow{}_nQ\longrightarrow Q
 \xrightarrow{\ n\ }nQ\longrightarrow0
\end{equation}
gives rise to a long exact sequence of the $\uExt^i(-,G)$, and hence
to the short exact sequence
\begin{equation}\tag{\SGARefTarget{sga7.sga7i.eviii.n2-3-7}{2.3.7}}
 0\longrightarrow\uHom(Q,G)_n
 \longrightarrow\uHom({}_nQ,G)
 \xrightarrow{\ \partial\ }
 {}_n\uExt^1(nQ,G)
 \longrightarrow0,
\end{equation}
where, as usual, $E_n$ denotes $E/nE$.  We retain from this sequence
the canonical surjective homomorphism
\begin{equation}\tag{\SGARefTarget{sga7.sga7i.eviii.n2-3-8}{2.3.8}}
 \uHom({}_nQ,G)
 \xrightarrow{\hspace{3em}\partial\hspace{3em}}
 {}_n\uExt^1(nQ,G),
\end{equation}
which is an isomorphism if $\uHom(Q,G)_n=0$, and a fortiori if (\SGARefLink{sga7.sga7i.eviii.n2-3-2}{2.3.2})
holds.  From (\SGARefLink{sga7.sga7i.eviii.n2-3-8}{2.3.8}) we deduce a homomorphism
\begin{equation}\tag{\SGARefTarget{sga7.sga7i.eviii.n2-3-9}{2.3.9}}
\begin{split}
 \operatorname{Hom}\bigl({}_nP,\uHom({}_nQ,G)\bigr)
 &\simeq
 \operatorname{Hom}({}_nP\otimes{}_nQ,G)\\
 &\longrightarrow
 \operatorname{Hom}\bigl({}_nP,{}_n\uExt^1(nQ,G)\bigr),
\end{split}
\end{equation}
which is an isomorphism when (\SGARefLink{sga7.sga7i.eviii.n2-3-2}{2.3.2}) holds.

\SGAvar{All the homomorphisms under consideration are functorial in $P,Q,G$.
The homomorphisms (\SGARefLink{sga7.sga7i.eviii.n2-3-5}{2.3.5}), (\SGARefLink{sga7.sga7i.eviii.n2-3-9}{2.3.9}), and (\SGARefLink{sga7.sga7i.eviii.n2-2-2}{2.2.2}) form the diagram
\begin{equation}\tag{\SGARefTarget{sga7.sga7i.eviii.n2-3-10}{2.3.10}}
\begin{tikzcd}[column sep=large,row sep=large]
 \&\mathrm{Biext}^1(P,Q;G)
   \arrow[ddl,"\varphi^n"']
   \arrow[ddr,"(2.3.5)"]\&\\
 \&-\&\\
 \operatorname{Hom}({}_nP\otimes{}_nQ,G)
   \arrow[rr,"(2.3.9)"]\&\&
 \operatorname{Hom}\bigl({}_nP,{}_n\uExt^1(nQnG)\bigr),
\end{tikzcd}
\end{equation}
where the horizontal arrow (\SGARefLink{sga7.sga7i.eviii.n2-3-9}{2.3.9}) is an isomorphism if (\SGARefLink{sga7.sga7i.eviii.n2-3-2}{2.3.2})
holds.

}{All the homomorphisms under consideration are functorial in $P,Q,G$.
The homomorphisms (\SGARefLink{sga7.sga7i.eviii.n2-3-5}{2.3.5}), (\SGARefLink{sga7.sga7i.eviii.n2-3-9}{2.3.9}), and (\SGARefLink{sga7.sga7i.eviii.n2-2-2}{2.2.2}) form the diagram
\begin{equation}\tag{\SGARefTarget{sga7.sga7i.eviii.n2-3-10}{2.3.10}}
\begin{tikzcd}[column sep=large,row sep=large]
 \&\mathrm{Biext}^1(P,Q;G)
   \arrow[ddl,"\varphi^n"']
   \arrow[ddr,"(2.3.5)"]\&\\
 \&-\&\\
 \operatorname{Hom}({}_nP\otimes{}_nQ,G)
   \arrow[rr,"(2.3.9)"]\&\&
 \operatorname{Hom}\bigl({}_nP,{}_n\uExt^1(nQ,G)\bigr),
\end{tikzcd}
\end{equation}
where the horizontal arrow (\SGARefLink{sga7.sga7i.eviii.n2-3-9}{2.3.9}) is an isomorphism if (\SGARefLink{sga7.sga7i.eviii.n2-3-2}{2.3.2})
holds.

}{E7I-VIII-013}\SGApage{I}{243}
\medskip
\noindent\underline{Proposition \SGARefTarget{sga7.sga7i.eviii.n2-3-11}{2.3.11}.}
\underline{The diagram (\SGARefLink{sga7.sga7i.eviii.n2-3-10}{2.3.10}) is anticommutative.}

When $\uHom(Q,G)=0$, this statement therefore recovers a
characterization of the homomorphism $\varphi^n$ of (\SGARefLink{sga7.sga7i.eviii.n2-2-2}{2.2.2}) in terms
of the homomorphisms (\SGARefLink{sga7.sga7i.eviii.n2-3-5}{2.3.5}) and (\SGARefLink{sga7.sga7i.eviii.n2-3-9}{2.3.9}).  This is the
characterization formerly used in the theory of abelian schemes,
where moreover $nQ=Q$.



\SGAvar{\medskip
\noindent\underline{Proof of \SGARefLink{sga7.sga7i.eviii.n2-3-11}{2.5.11}.}\footnote{For a more elegant
proof, see \SGARefLink{sga7.sga7i.eviii.n2-3-19}{2.3.19} below.}
Let $\xi$ be an element of $\mathrm{Biext}^1(P,Q;G)$, let
\[
 \varphi_\xi:{}_nP\otimes{}_nQ\longrightarrow G
\]
be the associated pairing, and let
\[
 u_\xi:{}_nP\longrightarrow{}_n\uExt^1(nQ,G)
\]
be the image of $\xi$ under (\SGARefLink{sga7.sga7i.eviii.n2-3-5}{2.3.5}).  We must prove that $-u_\xi$ is
obtained from $\varphi_\xi$ by the arrow (\SGARefLink{sga7.sga7i.eviii.n2-3-9}{2.3.9}), that is, that it is
the composite
\begin{equation}\tag{*}
 {}_nP
 \xrightarrow{\ \varphi_\xi\ }
 \uHom({}_nQ,G)
 \xrightarrow{\ (\SGARefLink{sga7.sga7i.eviii.n2-3-8}{2.3.8})\ }
 {}_n\uExt^1(nQ,G).
\end{equation}
To prove this, we may take a section $x$ of $\,{}_nP$ and prove that
its image under $-u_\xi$ equals its image under the composite (*).
If desired, this would allow us to reduce to the case
$P=(\mathbb Z/n\mathbb Z)_T$, but we shall not need that reduction.

}{\medskip
\noindent\underline{Proof of \SGARefLink{sga7.sga7i.eviii.n2-3-11}{2.3.11}.}\footnote{For a more elegant
proof, see \SGARefLink{sga7.sga7i.eviii.n2-3-19}{2.3.19} below.}
Let $\xi$ be an element of $\mathrm{Biext}^1(P,Q;G)$, let
\[
 \varphi_\xi:{}_nP\otimes{}_nQ\longrightarrow G
\]
be the associated pairing, and let
\[
 u_\xi:{}_nP\longrightarrow{}_n\uExt^1(nQ,G)
\]
be the image of $\xi$ under (\SGARefLink{sga7.sga7i.eviii.n2-3-5}{2.3.5}).  We must prove that $-u_\xi$ is
obtained from $\varphi_\xi$ by the arrow (\SGARefLink{sga7.sga7i.eviii.n2-3-9}{2.3.9}), that is, that it is
the composite
\begin{equation}\tag{*}
 {}_nP
 \xrightarrow{\ \varphi_\xi\ }
 \uHom({}_nQ,G)
 \xrightarrow{\ (\SGARefLink{sga7.sga7i.eviii.n2-3-8}{2.3.8})\ }
 {}_n\uExt^1(nQ,G).
\end{equation}
To prove this, we may take a section $x$ of $\,{}_nP$ and prove that
its image under $-u_\xi$ equals its image under the composite (*).
If desired, this would allow us to reduce to the case
$P=(\mathbb Z/n\mathbb Z)_T$, but we shall not need that reduction.

}{E7I-VIII-014}\medskip
\noindent a) Choose flat resolutions $L_\bullet$ and $M_\bullet$,
\[
 0\longrightarrow L_1\longrightarrow L_0\longrightarrow0,
 \qquad
 0\longrightarrow M_1\longrightarrow M_0\longrightarrow0,
\]
of $P$ and $Q$, respectively.  A flat resolution of $P$, for example,
is obtained by writing $P$ as a quotient of a flat
$\mathbb Z_T$-module $L_0$ and taking $L_1$ to be the kernel of
$L_0\longrightarrow P$; this kernel is flat because it is
torsion-free.  We obtain the complex $L_\bullet\otimes M_\bullet$
\[
 0\longrightarrow L_1\otimes M_1
 \xrightarrow{d_2}
 L_1\otimes M_0+L_0\otimes M_1
 \xrightarrow{d_1}
 L_0\otimes M_0
 \longrightarrow0,
\]
whose sheaf $\underline H_1$ will serve as $\uTor_1(P,Q)$.

\SGApage{I}{\(\approx\)244}
Let $L'_0$ (respectively, $M'_0$) be the subsheaf of $L_0$
(respectively, $M_0$) which is the inverse image of $\,{}_nP$
(respectively, $\,{}_nQ$).  Thus
\[
 0\longrightarrow L_1\longrightarrow L'_0\longrightarrow0
 \qquad
 \left(\text{respectively, }
 0\longrightarrow M_1\longrightarrow M'_0\longrightarrow0\right)
\]
is a flat resolution of $\,{}_nP$ (respectively, of $\,{}_nQ$).  The
homomorphism
\[
 {}_nP\otimes{}_nQ
 \longrightarrow
 \uTor_1(P,Q)=\underline H_1(L_\bullet\otimes M_\bullet)
\]
may then be described, by passage to the quotient, from the
homomorphism
\[
 L'_0\otimes M'_0
 \longrightarrow
 \underline H_1(L_\bullet\otimes M_\bullet)
\]
given by
\[
 x_0\otimes y_0
 \longmapsto
 \text{the class of the $1$-cycle }
 [nx_0]\otimes y_0-x_0\otimes[ny_0].
\]
Here the brackets indicate that the element in question is regarded
as a section of the subsheaf $L_1$ of $L_0$ (respectively, $M_1$ of
$M_0$).  This description follows from the calculation made in the
proof of \SGARefLink{sga7.sga7i.eviii.n2-1-10}{2.1.10}.

\SGAvar{\medskip
\noindent b) Take an injective resolution
\[
 C^\bullet=\bigl(0\longrightarrow C^0
 \xrightarrow{d^0}C^1\longrightarrow\cdots\bigr)
\]
of $G$, and identify
\[
 \mathrm{Biext}^1(P,Q;G)
 \simeq
 \operatorname{Ext}^1\!\left(
   P\mathbin{\overset{\mathbb L}{\otimes}}Q,G
 \right)
 =H^1\!\left(
   \operatorname{Hom}^{\boldsymbol\cdot}
   (L_\bullet\otimes M_\bullet,C^{\boldsymbol\cdot})
 \right).
\]
The element $\xi$ may therefore be represented by a pair of
homomorphisms $f_1,f_0$, giving the anticommutative diagram
\begin{equation}\tag{\SGARefTarget{sga7.sga7i.eviii.n2-3-12}{2.3.12}}
\begin{tikzcd}[column sep=huge,row sep=huge]
 L_1\otimes M_0+L_0\otimes M_1
   \arrow[r,"d_1"]
   \arrow[d,"f_1"']
   \arrow[rd,phantom,"-" description]\&
 L_0\otimes M_0
   \arrow[d,"f_0"']\\
 C^0\arrow[r,"d^0"]\&C^1.
\end{tikzcd}
\end{equation}
By definition, the pairing
\[
 \varphi_\xi:{}_nP\otimes{}_nQ\longrightarrow G
\]
is obtained, by passage to the quotient, from the homomorphism
\[
 \varphi:L'_0\otimes M'_0\longrightarrow G
\]
given by
\begin{equation}\tag{\SGARefTarget{sga7.sga7i.eviii.n2-3-13}{2.3.13}}
 \varphi(x_0\otimes y_0)
 =f_1\bigl([nx_0]\otimes y_0-x_0\otimes[ny_0]\bigr)
 \in G=\ker(C^0\longrightarrow C^1).
\end{equation}
Since assertion \SGARefLink{sga7.sga7i.eviii.n2-3-11}{2.3.11} is local on $T$, we may suppose that the given
section $x$ of $\,{}_nP$ lifts to a section $x_0$ of $L'_0$.
Regarding $y_0$ as variable, the preceding formula then defines a
homomorphism
\begin{equation}\tag{\SGARefTarget{sga7.sga7i.eviii.n2-3-14}{2.3.14}}
 \varphi':M'_0\longrightarrow G,
 \qquad
 y_0\longmapsto
 f_1\bigl([nx_0]\otimes y_0\to x_0\otimes[ny_0]\bigr),
\end{equation}
which passes to the quotient to give
\[
 {}_nQ\longrightarrow G,
 \qquad
 y\longmapsto\varphi_\xi(x,y),
\]
the homomorphism associated with $\varphi_\xi$ and $x$.



}{\medskip
\noindent b) Take an injective resolution
\[
 C^\bullet=\bigl(0\longrightarrow C^0
 \xrightarrow{d^0}C^1\longrightarrow\cdots\bigr)
\]
of $G$, and identify
\[
 \mathrm{Biext}^1(P,Q;G)
 \simeq
 \operatorname{Ext}^1\!\left(
   P\mathbin{\overset{\mathbb L}{\otimes}}Q,G
 \right)
 =H^1\!\left(
   \operatorname{Hom}^{\boldsymbol\cdot}
   (L_\bullet\otimes M_\bullet,C^{\boldsymbol\cdot})
 \right).
\]
The element $\xi$ may therefore be represented by a pair of
homomorphisms $f_1,f_0$, giving the anticommutative diagram
\begin{equation}\tag{\SGARefTarget{sga7.sga7i.eviii.n2-3-12}{2.3.12}}
\begin{tikzcd}[column sep=huge,row sep=huge]
 L_1\otimes M_0+L_0\otimes M_1
   \arrow[r,"d_1"]
   \arrow[d,"f_1"']
   \arrow[rd,phantom,"-" description]\&
 L_0\otimes M_0
   \arrow[d,"f_0"']\\
 C^0\arrow[r,"d^0"]\&C^1.
\end{tikzcd}
\end{equation}
By definition, the pairing
\[
 \varphi_\xi:{}_nP\otimes{}_nQ\longrightarrow G
\]
is obtained, by passage to the quotient, from the homomorphism
\[
 \varphi:L'_0\otimes M'_0\longrightarrow G
\]
given by
\begin{equation}\tag{\SGARefTarget{sga7.sga7i.eviii.n2-3-13}{2.3.13}}
 \varphi(x_0\otimes y_0)
 =f_1\bigl([nx_0]\otimes y_0-x_0\otimes[ny_0]\bigr)
 \in G=\ker(C^0\longrightarrow C^1).
\end{equation}
Since assertion \SGARefLink{sga7.sga7i.eviii.n2-3-11}{2.3.11} is local on $T$, we may suppose that the given
section $x$ of $\,{}_nP$ lifts to a section $x_0$ of $L'_0$.
Regarding $y_0$ as variable, the preceding formula then defines a
homomorphism
\begin{equation}\tag{\SGARefTarget{sga7.sga7i.eviii.n2-3-14}{2.3.14}}
 \varphi':M'_0\longrightarrow G,
 \qquad
 y_0\longmapsto
 f_1\bigl([nx_0]\otimes y_0-x_0\otimes[ny_0]\bigr),
\end{equation}
which passes to the quotient to give
\[
 {}_nQ\longrightarrow G,
 \qquad
 y\longmapsto\varphi_\xi(x,y),
\]
the homomorphism associated with $\varphi_\xi$ and $x$.



}{E7I-VIII-015}\SGApage{I}{\(\approx\)245}
\medskip
\noindent c) We must take the image of this homomorphism under the
connecting operator $\partial$ of (\SGARefLink{sga7.sga7i.eviii.n2-3-7}{2.3.7}), associated with the exact
sequence (\SGARefLink{sga7.sga7i.eviii.n2-3-6}{2.3.6}).  Now $nQ\simeq Q/{}_nQ$ is plainly also isomorphic
to $M_0/M'_0$, so that we have a homomorphism of exact sequences
\[
\begin{tikzcd}[column sep=large,row sep=large]
 0\arrow[r]\&M'_0\arrow[r]\arrow[d]\&M_0\arrow[r]\arrow[d]
   \&M_0/M'_0\arrow[r]\arrow[d,"\wr"']\&0\\
 0\arrow[r]\&{}_nQ\arrow[r]\&Q\arrow[r]\&nQ\arrow[r]\&0\&{,}
\end{tikzcd}
\]
where the last vertical arrow is an isomorphism.  This shows that the
desired section of $\uExt^1(nQ,G)$ may also be computed as the image
of $\varphi'$ in (\SGARefLink{sga7.sga7i.eviii.n2-3-14}{2.3.14}) under the connecting operator
\[
 \uHom(M'_0,G)\xrightarrow{\ \partial\ }
 \uExt^1(M_0/M'_0,G),
\]
or, equivalently, as the section of the second term defined by the
element of global $\mathrm{Ext}^1$ which is the image of $\varphi'$
under the connecting operator
\[
 \mathrm{Hom}(M'_0,G)\longrightarrow
 \mathrm{Ext}^1(M_0/M'_0,G).
\]
The usual construction of the connecting operator gives the following
recipe.  Extend the composite (\SGARefLink{sga7.sga7i.eviii.n2-3-13}{2.3.13})
\[
 M'_0\xrightarrow{\ \varphi'\ }G\longrightarrow C^0
\]
to a homomorphism
\[
 \psi:M_0\longrightarrow C^0,
\]
which is possible since $C^0$ is injective, and consider the unique
homomorphism
\begin{equation}\tag{\SGARefTarget{sga7.sga7i.eviii.n2-3-15}{2.3.15}}
 \lambda:M_0/M'_0\longrightarrow C^1
\end{equation}
which makes the diagram
\[
\begin{tikzcd}[column sep=large,row sep=large]
 0\arrow[r]\&M'_0\arrow[r]\arrow[d,"\varphi'"']
   \&M_0\arrow[r]\arrow[d,"\psi"']
   \&M_0/M'_0\arrow[r]\arrow[d,"\lambda"']\&0\\
 0\arrow[r]\&G\arrow[r]\&C^0\arrow[r]\&C^1\&\&{.}
\end{tikzcd}
\]
commutative.  It is clear that (\SGARefLink{sga7.sga7i.eviii.n2-3-15}{2.3.15}) takes values in the
$1$-cocycles---that is, it is a $1$-cocycle---and that this cocycle
represents the desired element of
\[
 \mathrm{Ext}^1(nQ,G)=\mathrm{Ext}^1(M_0/M'_0,G).
\]

\SGApage{I}{\(\approx\)247}
\medskip
\noindent d) It remains to describe explicitly the section of
$\uExt^1(nQ,G)$ obtained from $\xi$ and $x$ by means of arrow (\SGARefLink{sga7.sga7i.eviii.n2-3-5}{2.3.5})
in diagram (\SGARefLink{sga7.sga7i.eviii.n2-3-10}{2.3.10}).  For this purpose we may disregard the fact that
the section $x$ of $P$ is annihilated by $n$.  In terms of the chosen
lift $x_0$ of $x$ to a section of $L_0$, the description is as follows.
This lift defines an augmentation homomorphism
\[
 u:M_\bullet\longrightarrow L_\bullet\otimes M_\bullet,
 \qquad
 y_0\longmapsto x_0\otimes y_0,
 \quad
 y_1\longmapsto x_0\otimes y_1.
\]
The section of $\uExt^1(Q,G)$ obtained from (\SGARefLink{sga7.sga7i.eviii.n2-3-1}{2.3.1}) is therefore the
one obtained from the element $\xi$ of
$\mathrm{Ext}^1(L_\bullet\otimes M_\bullet,G)$ by composition with
the preceding homomorphism.  Thus it is represented by the
degree-one homomorphism
\[
\begin{tikzcd}[column sep=large,row sep=huge]
 0\arrow[r]\&M_1\arrow[r,"d_1"]
   \arrow[d,"{y_1\mapsto f_1(x_0\otimes y_1)}"']
   \arrow[rd,phantom,"-" description]
   \&M_0\arrow[r]\arrow[d,"{y_0\mapsto f_0(x_0\otimes y_0)}"]\&0\\
 0\arrow[r]\&C^0\arrow[r,"d^0"]\&C^1\&\ldots\&{.}
\end{tikzcd}
\]
To restrict this class to $nQ\hookrightarrow Q$, use the isomorphism
$nQ\simeq M_0/M'_0$, which gives a flat resolution
\[
 0\longrightarrow M'_0\longrightarrow M_0\longrightarrow0
\]
of $nQ$.  The inclusion $nQ\hookrightarrow Q$ lifts to the following
homomorphism of resolutions:
\[
\begin{tikzcd}[column sep=large,row sep=huge]
 0\arrow[r]\&M'_0\arrow[r]
   \arrow[d,"{y'_0\mapsto[ny'_0]}"']
   \arrow[rd,phantom,"+" description]
   \&M_0\arrow[r]\arrow[d,"{y_0\mapsto ny_0}"]\&0\\
 0\arrow[r]\&M_1\arrow[r]\&M_0\arrow[r]\&0\&{.}
\end{tikzcd}
\]
By composition, the desired section of $\uExt^1(nQ,G)$ is therefore
associated with the element of $\mathrm{Ext}^1(nQ,G)$ represented by
the degree-one homomorphism of complexes
\[
\begin{tikzcd}[column sep=large,row sep=large]
 0\arrow[r,dash]\&M'_0\arrow[r,dash,"d_1"]\arrow[d,"g_1"']
   \&M_0\arrow[r,dash]\arrow[d,"g_0"]\&0\&{}\\
 0\arrow[r,dash]\&C^0\arrow[r,dash]\&C^1\arrow[r,dash]\&\dots\&{,}
\end{tikzcd}
\]
where
\[
 g_1(y'_0)=f_1\bigl(x_0\otimes[ny'_0]\bigr),
 \qquad
 g_0(y_0)=n f_0(x_0\otimes y_0).
\]

\SGApage{I}{\(\approx\)248}
\medskip
\noindent e) Using now the fact that $x$ is a section of $\,{}_nP$,
we show that the preceding element of $\mathrm{Ext}^1(nQ,G)$ is the
negative of the element defined by (\SGARefLink{sga7.sga7i.eviii.n2-3-15}{2.3.15}).  For any homomorphism
\[
 h:M_0\longrightarrow C^0,
\]
the element of $\mathrm{Ext}^1$ considered in d) is unchanged if
$g_1,g_0$ are replaced, respectively, by
\[
 g'_1=g_1+h d_1,
 \qquad
 g'_0=g_0-d^0h.
\]
We have at our disposal the homomorphism
$\psi:M_0\longrightarrow C^0$ constructed in c), which by definition
satisfies
\[
 \psi(y'_0)=
 f_1\bigl([nx_0]\otimes y'_0-x_0\otimes[ny'_0]\bigr)
\]
for every section $y'_0$ of $M'_0$.  Equivalently,
\[
 f_1\bigl(x_0\otimes[ny'_0]\bigr)
 \bigl(=g_1(y'_0)\bigr)
 =f_1\bigl([nx_0]\otimes y'_0\bigr)-\psi(y'_0).
\]
Define $h:M_0\longrightarrow C^0$ by
\[
 h(y_0)=-f_1\bigl([nx_0]\otimes y_0\bigr)+\psi(y_0).
\]
By construction, then,
\[
 g'_1=0,
 \qquad
 g'_0(y_0)=n f_0(x_0\otimes y_0)
 +d^0f_1\bigl([nx_0]\otimes y_0\bigr)-d^0\psi(y_0).
\]
Since $d^0f_1=-f_0d_1$, we have
\[
 d^0f_1\bigl([nx_0]\otimes y_0\bigr)
 =-f_0(nx_0\otimes y_0),
\]
and hence
\[
 g'_1=0,
 \qquad
 g'_0=-d^0\psi.
\]
Thus the element in question is the one defined by the homomorphism
\[
 M_0/M'_0\simeq nQ\longrightarrow C^1
\]
induced by $-d^0\psi$ upon passage to the quotient.  We therefore
obtain the negative of the element defined by (\SGARefLink{sga7.sga7i.eviii.n2-3-15}{2.3.15}).

\SGApage{I}{\(\approx\)249}
\SGAvar{\medskip
\noindent\underline{Remarks \SGARefTarget{sga7.sga7i.eviii.n2-3-16}{2.3.16}.}  The preceding proof actually
gives a slightly more precise compatibility than the one stated in
\SGARefLink{sga7.sga7i.eviii.n2-3-11}{2.3.11}, obtained by choosing a section $x$ of $\,{}_nP$.  This section
then defines a homomorphism
\[
 {}_nQ\longrightarrow G,
 \qquad
 y\longmapsto\varphi_\xi(x,y),
\]
and hence, by the connecting homomorphism associated with (\SGARefLink{sga7.sga7i.eviii.n2-3-6}{2.3.6}), an
element of $\mathrm{Ext}^1(nQ,G)$.  On the other hand, if we forget
that $x$ is annihilated by $n$, we consider the homomorphism
$\mathbb Z_T\longrightarrow P$ defined by $X$.  Its inverse image is a
biextension of $(\mathbb Z_T,Q)$ by $G$, whose class may be interpreted
as an element of $\mathrm{Ext}^1(Q,G)$ and therefore induces an element
of $\mathrm{Ext}^1(nQ,G)$.  With these conventions, the two elements
obtained in $\mathrm{Ext}^1(nQ,G)$ are \underline{opposite}.



}{\medskip
\noindent\underline{Remarks \SGARefTarget{sga7.sga7i.eviii.n2-3-16}{2.3.16}.}  The preceding proof actually
gives a slightly more precise compatibility than the one stated in
\SGARefLink{sga7.sga7i.eviii.n2-3-11}{2.3.11}, obtained by choosing a section $x$ of $\,{}_nP$.  This section
then defines a homomorphism
\[
 {}_nQ\longrightarrow G,
 \qquad
 y\longmapsto\varphi_\xi(x,y),
\]
and hence, by the connecting homomorphism associated with (\SGARefLink{sga7.sga7i.eviii.n2-3-6}{2.3.6}), an
element of $\mathrm{Ext}^1(nQ,G)$.  On the other hand, if we forget
that $x$ is annihilated by $n$, we consider the homomorphism
$\mathbb Z_T\longrightarrow P$ defined by $x$.  Its inverse image is a
biextension of $(\mathbb Z_T,Q)$ by $G$, whose class may be interpreted
as an element of $\mathrm{Ext}^1(Q,G)$ and therefore induces an element
of $\mathrm{Ext}^1(nQ,G)$.  With these conventions, the two elements
obtained in $\mathrm{Ext}^1(nQ,G)$ are \underline{opposite}.



}{E7I-VIII-016}\medskip
\noindent \SGARefTarget{sga7.sga7i.eviii.n2-3-17}{2.3.17}.  Let us use the explicit cocycle description of the
pairing
\[
 \varphi_\xi^n:{}_nP\otimes{}_nQ\longrightarrow G
\]
given in (\SGARefLink{sga7.sga7i.eviii.n2-3-12}{2.3.12}) to obtain another expression for it.  First suppose
that the biextension $E$ of class $\xi$ admits a section $s$ over
$P\times Q$.  As in VII \SGARefLink{sga7.sga7i.evii.n3-8-13}{3.8.13}, take for the flat resolution of $P$
the one with $L_0=\mathbb Z[P]$ and $L_1$ the kernel of the canonical
homomorphism $\mathbb Z[P]\longrightarrow P$, and define
$M_\bullet$ similarly.  There is then a unique biextension
trivialization of the inverse image of $E$ on $(L_0,M_0)$ whose
restriction to $P\times Q$ is given by $s$; denote this trivialization
again by $s$ (cf. VII \SGARefLink{sga7.sga7i.evii.n3-8-2}{3.8.2}).  By VII \SGARefLink{sga7.sga7i.evii.n3-8-12}{3.8.12}, the class of $\xi$ in
\[
 \mathrm{Ext}^1\!\left(
 P\mathbin{\overset{\mathbb L}{\otimes}}Q,G\right)
\]
is represented by the cocycle with components $(f,g)$ characterized by
\[
 s(x_1,y_0)=f(x_1\otimes y_0)e_{E/Q}(y_0),
 \qquad
 s(x_0,y_1)=g(x_0\otimes y_1)e_{E/P}(x_0).
\]
If $x$ and $y$ are sections of $\,{}_nP$ and $\,{}_nQ$, lift them to
the corresponding sections $x_0=[x]$ and $y_0=[y]$ of $L_0$ and
$M_0$.  Applying formula (\SGARefLink{sga7.sga7i.eviii.n2-3-13}{2.3.13}), and putting $z=s(x,y)$, gives
\begin{equation}\tag{\SGARefTarget{sga7.sga7i.eviii.n2-3-18}{2.3.18}}
 \varphi_\xi^n(x,y)
 ={}^n z/e_{E/Q}(y)-z^n/e_{E/P}(x).
\end{equation}
Here the exponent $n$ written on the left, respectively on the right,
of $z$ denotes the $n$th power for the left, respectively right,
multiplicative structure on $E$.  It therefore gives a section of $E$
over $(0,y)$, respectively $(x,0)$; the meaning of the ``quotient'' of
such an element by $e_{E/Q}(y)$, respectively $e_{E/P}(x)$, is clear.

\SGApage{I}{\(\approx\)250}
The preceding formula (\SGARefLink{sga7.sga7i.eviii.n2-3-18}{2.3.18}) still makes sense whenever $z$ is a
global section of $E$ lifting the sections $x,y$ of
$\,{}_nP,\,{}_nQ$.  I claim that (\SGARefLink{sga7.sga7i.eviii.n2-3-18}{2.3.18}) is valid for every such
section, without assuming that a section of $E$ over $P\times Q$
exists.  The section $(x,y)$ of $P\times Q$ always lifts locally to a
section of $E$, since $E\longrightarrow P\times Q$ is an epimorphism;
thus the formula computes $\varphi_\xi^n(x,y)$ in every case.
Moreover, if $z$ is replaced by another section $z'$ of $E$ lifting
$(x,y)$, then $z'=z+g$ for a section $g$ of $G$.  The two terms on the
right-hand side of (\SGARefLink{sga7.sga7i.eviii.n2-3-18}{2.3.18}) corresponding to $z'$ are the terms
corresponding to $z$, each increased by $ng$.  Their difference is
therefore independent of $z$, showing that (\SGARefLink{sga7.sga7i.eviii.n2-3-18}{2.3.18}) defines a
well-determined homomorphism
\[
 {}_nP\otimes{}_nQ\longrightarrow G.
\]

\SGApage{I}{\(\approx\)251}
\SGAvar{It remains to show that this homomorphism is indeed $\varphi^n$.
For this we reduce at once to the case
\[
 P=Q=(\mathbb Z/n\mathbb Z)_T.
\]
Then $P\times Q$ is of the form $I_T$, where
$I=\mathbb Z/n\mathbb Z\times\mathbb Z/n\mathbb Z$ is a finite set.
In this case the torsor $E$ under $G_{P\times Q}$---which identifies
with a family of torsors $E_i$ under $G$, for $i\in I$---plainly
becomes trivial locally on $T$.  We are thus reduced to the already
treated case in which $E$ has a global section over $P\times Q$.



}{It remains to show that this homomorphism is indeed $\varphi_\xi^n$.
For this we reduce at once to the case
\[
 P=Q=(\mathbb Z/n\mathbb Z)_T.
\]
Then $P\times Q$ is of the form $I_T$, where
$I=\mathbb Z/n\mathbb Z\times\mathbb Z/n\mathbb Z$ is a finite set.
In this case the torsor $E$ under $G_{P\times Q}$---which identifies
with a family of torsors $E_i$ under $G$, for $i\in I$---plainly
becomes trivial locally on $T$.  We are thus reduced to the already
treated case in which $E$ has a global section over $P\times Q$.



}{E7I-VIII-017}\medskip
\noindent \SGARefTarget{sga7.sga7i.eviii.n2-3-19}{2.3.19}.  \underline{Another proof of \SGARefLink{sga7.sga7i.eviii.n2-3-11}{2.3.11}} (by
Deligne).

\medskip
\noindent a) Reduction to $P=\mathbb Z/n\mathbb Z$ and to
$T=(\mathrm{Sets})$.

\medskip
\noindent b) Let $\partial$ be the following arrow from $nQ$ to
$\,{}_nQ[1]$:
\[
 nQ\xleftarrow[\text{``$n$''}]{\ \sim\ }
 [\,{}_nQ\longrightarrow Q\,]
 \longrightarrow{}_nQ[1].
\]
The following diagram is anticommutative:
\[
\begin{tikzcd}[row sep=huge,column sep=huge]
 nQ[0]\arrow[r,"1"]\arrow[d,"\partial"']
   \&{[Q\xrightarrow{n}Q]}\\
 {}_nQ[1]\arrow[ru,"1"']\&{.}
\end{tikzcd}
\]
Indeed, the arrows
\[
\begin{tikzcd}[column sep=small,row sep=large]
 {[{}_nQ}\arrow[r]\arrow[dd,"0"']\&{Q]}\arrow[dd,"n"]
   \&\&{[{}_nQ}\arrow[r]\arrow[dd,"1"']\&{Q]}\arrow[dd,"0"']
   \&\text{are opposite}
   \&{[{}_nQ}\arrow[r]\&{Q]}\arrow[ddl,"1"]\\
 \&\&\text{and}\&\&\&\text{by means of}\quad H\,:\&\&\\
 {[Q}\arrow[r,"n"]\&{Q]}\&\&{[Q}\arrow[r,"n"]\&{Q]}
   \&\&{[Q}\arrow[r,"n"]\&{Q]}
\end{tikzcd}
\]

\SGApage{I}{\(\approx\)252}
\medskip
\noindent c) Diagram \SGARefLink{sga7.sga7i.eviii.n2-3-10}{2.3.10} is
\[
\begin{tikzcd}[column sep=small,row sep=small]
 \&\&\&H^1\mathrm{RHom}
   (\mathbb Z\mathbin{\overset{\mathbb L}{\otimes}}Q,G)
   \arrow[ddd]\\
 \&\&\&+\\
 H^1\mathrm{RHom}
   (\mathbb Z/n\mathbb Z\mathbin{\overset{\mathbb L}{\otimes}}Q,G)
   \arrow[rr]\arrow[dd]\arrow[uurrr]
 \&\&H^1\mathrm{RHom}
   (\mathbb Z/n\mathbb Z,\mathrm{RHom}(Q,G))
   \arrow[dd]\&\\
 \&\&\&\mathrm{Ext}^1(Q,G)\arrow[ddl]\\
 \mathrm{Hom}(\mathrm{Tor}_1(\mathbb Z/n\mathbb Z,Q),G)\arrow[d]
 \&|[draw,circle,inner sep=1pt]|{-\,?}
 \&\mathrm{Hom}(\mathbb Z/n\mathbb Z,\mathrm{Ext}^1(Q,G))
   \arrow[d]\&\\
 \mathrm{Hom}({}_nQ,G)\arrow[rr]
 \&\&\mathrm{Ext}^1(nQ,G)\&
\end{tikzcd}
\]
and:

\smallskip
\noindent c1) We have an isomorphism
\[
 \mathbb Z/n\mathbb Z
 \mathbin{\overset{\mathbb L}{\otimes}}Q
 \simeq[Q\xrightarrow{n}Q],
\]
and, under this isomorphism, the first composite vertical arrow is
induced by
\[
 {}_nQ[1]\longrightarrow[Q\xrightarrow{n}Q].
\]

\smallskip
\noindent c2) The second horizontal arrow is the transpose of the
arrow $\partial$ displayed above in b).

\smallskip
\noindent c3) If the diagram is traversed clockwise, it is clear that
one obtains the arrow induced by
\[
 nQ[0]\longrightarrow[Q\xrightarrow{n}Q].
\]
This proves \SGARefLink{sga7.sga7i.eviii.n2-3-11}{2.3.11}.



\SGApage{I}{253}
\medskip
\noindent \SGARefTarget{sga7.sga7i.eviii.n2-4}{2.4}.  Let
\[
 f:\underline T'\longrightarrow\underline T
\]
be a morphism of topoi.  There is then a canonical homomorphism
(VII \SGARefLink{sga7.sga7i.evii.n2-8}{2.8})
\[
 \Biext^1(P,Q;G)\longrightarrow\Biext^1(P',Q';G'),
\]
where primes denote inverse-image groups.  Since, for every
commutative group $E$ of $\underline T$, there is plainly a canonical
isomorphism
\[
 ({}_nE)'\simeq{}_n(E'),
\]
the homomorphisms (\SGARefLink{sga7.sga7i.eviii.n2-2-2}{2.2.2}) over $\underline T$ and $\underline T'$ give
a canonical diagram
\begin{equation}\tag{\SGARefTarget{sga7.sga7i.eviii.n2-4-1}{2.4.1}}
\begin{tikzcd}[column sep=huge,row sep=huge]
 \Biext^1(P,Q;G)\arrow[r]\arrow[d,"{\varphi^n_{P,Q}}"']
   \&\Biext^1(P',Q';G')\arrow[d,"{\varphi^n_{P',Q'}}"]\\
 \Hom({}_nP\otimes{}_nQ,G)\arrow[r]
   \&\Hom({}_nP'\otimes{}_nQ',G').
\end{tikzcd}
\end{equation}
Similarly, the homomorphisms (\SGARefLink{sga7.sga7i.eviii.n2-2-3}{2.2.3}) give an analogous diagram in
terms of the $T_\ell$.  I claim that these diagrams are
\underline{commutative}.  The verification is essentially immediate
from the definitions and is left to the reader.  More generally,
notice that \underline{all} the canonical homomorphisms considered in
this section and the preceding one are \underline{compatible with
change of topos}, in the evident sense.  We shall henceforth use
compatibilities of this kind without further mention.




\SGApage{I}{\(\approx\)254}
\subsection*{\SGARefTarget{sga7.sga7i.eviii.n3}{3}. Biextensions of group schemes $(P,Q)$ by
$\underline G_m$: generalities}

Throughout this section we fix a base scheme $S$ and, as announced in
Section 0, work in the fppf topos of $S$.

\medskip
\noindent \SGARefTarget{sga7.sga7i.eviii.n3-1}{3.1}.  Let us first recall the most important example of a
group $\Biext^1(P,Q;\underline G_m)$: the case where $P$ and $Q$ are
two abelian schemes over $S$, and where $\Biext^1$ is interpreted as
the group of divisorial correspondences on $(P,Q)$ (VII \SGARefLink{sga7.sga7i.evii.n2-9-4}{2.9.4}).  The
$\ell$-adic pairings of the preceding section are then the familiar
pairings [1] in the theory of abelian schemes.

\medskip
\noindent \SGARefTarget{sga7.sga7i.eviii.n3-2}{3.2}.  More generally, suppose only that $Q$ is an abelian
scheme.  It is well known (VII \SGARefLink{sga7.sga7i.evii.n1-3-8}{1.3.8}) that
\begin{equation}\tag{\SGARefTarget{sga7.sga7i.eviii.n3-2-1}{3.2.1}}
 \uHom(Q,\underline G_m)=0,
\end{equation}
and hence (\SGARefLink{sga7.sga7i.eviii.n1-4}{1.4}) there is a canonical isomorphism
\begin{equation}\tag{\SGARefTarget{sga7.sga7i.eviii.n3-2-2}{3.2.2}}
 \Biext^1(P,Q;\underline G_m)\simeq\Hom(P,Q^*),
\end{equation}
where
\begin{equation}\tag{\SGARefTarget{sga7.sga7i.eviii.n3-2-3}{3.2.3}}
 Q^*=\uExt^1(Q,\underline G_m)
\end{equation}
is the dual abelian scheme of $Q$.  (For the representability of
$Q^*$, due to M. Raynaud in the general case, see [1].)

When $P$ is also an abelian scheme, symmetrically to (\SGARefLink{sga7.sga7i.eviii.n3-2-2}{3.2.2}), there is
a canonical isomorphism
\begin{equation}\tag{\SGARefTarget{sga7.sga7i.eviii.n3-2-4}{3.2.4}}
 \Biext^1(P,Q;\underline G_m)\simeq\Hom(Q,P^*).
\end{equation}
Thus the datum of a biextension $\xi$ of $(P,Q)$ by $\underline G_m$
may be interpreted equivalently as the datum of a homomorphism
\[
 u:P\longrightarrow Q^*
\]
or of a homomorphism
\[
 v:Q\longrightarrow P^*.
\]

\SGApage{I}{\(\approx\)255}
Recall that a homomorphism of abelian schemes $w:A\longrightarrow B$
is an isomorphism if and only if, for every integer $n>0$, the induced
homomorphism ${}_nw:{}_nA\longrightarrow{}_nB$ is an isomorphism.  In
the case of $u$ and $v$ above, ${}_nu:{}_nP\longrightarrow{}_nQ^*$ is
identified with the homomorphism
${}_nP\longrightarrow D({}_nQ)$ defined by the pairing
\[
 \varphi_\xi^n:{}_nP\times{}_nQ\longrightarrow\underline G_m
\]
(see \SGARefLink{sga7.sga7i.eviii.n2-3-11}{2.3.11}), whereas ${}_nv:{}_nQ\longrightarrow{}_nP^*$ is
identified with the homomorphism
${}_nQ\longrightarrow D({}_nP)$ defined by the pairing
\[
 \psi_\xi^n:{}_nP\times{}_nQ\longrightarrow\underline G_m.
\]
Now $\varphi_\xi^n=-\psi_\xi^n$ (\SGARefLink{sga7.sga7i.eviii.n2-1-10}{2.1.10}), so, up to sign,
${}_nu$ and ${}_nv$ are transposes of one another.  Since ${}_nP$ and
${}_nQ$ are finite and locally free over $S$, and therefore satisfy
Cartier duality, it follows that ${}_nu$ is an isomorphism if and only
if ${}_nv$ is an isomorphism.  Consequently $u$ is an isomorphism if
and only if $v$ is an isomorphism.

An equivalent, though slightly less symmetric, formulation is the
following.  Consider the canonical biextension (the Weil
biextension) of $(P,P^*)$ by $\underline G_m$, corresponding to the
identity homomorphism $v:P^*\longrightarrow P^*$.  Then the
corresponding canonical homomorphism
\[
 u:P\longrightarrow P^{**}
\]
is an \underline{isomorphism}.  This is the \underline{biduality
theorem} of Barsotti and Cartier.




\SGApage{I}{256}
\medskip
\noindent \SGARefTarget{sga7.sga7i.eviii.n3-3}{3.3}.  We now turn to the case where $Q$ is finite locally
free over $S$, or is a group scheme of multiplicative type over $S$.
The key lemma in this case is the following well-known result.

\medskip
\noindent\underline{Proposition \SGARefTarget{sga7.sga7i.eviii.n3-3-1}{3.3.1}.}  \underline{Let $Q$ be a
commutative group scheme over $S$, and suppose that $Q$ is finite
locally free over $S$, or that $Q$ is of finite type and of
multiplicative type over $S$ (SGA 3 IX).  Then}
\[
 \uExt^1(Q,\underline G_m)=0.
\]

\noindent\underline{Proof.}  We distinguish the two cases.

\medskip
\noindent a) Suppose that $Q$ is finite locally free over $S$.  We
must prove that every extension $E$ of $Q$ by $\underline G_m$ splits
locally for the fppf topology.  We may plainly suppose $S$ affine;
there is then an integer $n>0$ such that $nQ=0$ (SGA 3 IX 6.1).
Moreover, $n\cdot\id$ is an epimorphism of $\underline G_m$ for the
fppf topology, and hence the Kummer exact sequence
\[
 0\longrightarrow\mu_n\longrightarrow\underline G_m
 \xrightarrow{\ n\ }\underline G_m\longrightarrow0
\]
gives an exact sequence of the $\Ext^i(Q,-)$ and, in particular, an
\underline{epimorphism}
\[
 \Ext^1(Q,\mu_n)\longrightarrow
 {}_n\Ext^1(Q,\underline G_m)=\Ext^1(Q,\underline G_m).
\]
The last equality follows because multiplication by $n$, being zero
on $Q$, is also zero on $\Ext^1(Q,\underline G_m)$.  Thus the
extension $E$ comes from an extension $F$ of $Q$ by $\mu_n$.

The group $F$ is represented by a finite locally free scheme over $S$
(fpqc descent theory, SGA 1 VII).  By the known properties of Cartier
duality for finite locally free groups (SGA 3 VII 3),
\[
 D(F)=\uHom(F,\underline G_m)
\]
is an extension of $D(\mu_n)$ by $D(Q)$.  Here one uses the following
fact: if $A\longrightarrow B$ is a monomorphism of commutative finite
locally free groups over $S$, then
$D(B)\longrightarrow D(A)$ is faithfully flat.  Since $D(B)$ and
$D(A)$ are themselves finite locally free over $S$, the fibrewise
criterion for faithful flatness reduces us to the case where $S$ is
the spectrum of a field $k$.  In that case the assertion says that
the homomorphism of affine coordinate rings
\[
 k(D(A))\longrightarrow k(D(B))
\]
is injective (SGA 3 $\mathrm{VI}_B$ 11.14).  This homomorphism is the
transpose of $k(B)\longrightarrow k(A)$, which is surjective because
$A\longrightarrow B$ is a closed immersion (SGA 3 $\mathrm{VI}_B$
1.4.2).  Hence $k(D(A))\longrightarrow k(D(B))$ is indeed injective.

\SGApage{I}{\(\approx\)257}
The inclusion $\mu_n\longrightarrow\underline G_m$ identifies with a
section of
\[
 D(\mu_n)=\uHom(\mu_n,\underline G_m),
\]
and it remains only to prove that, locally for the fppf topology on
$S$, this inclusion extends to a homomorphism
$F\longrightarrow\underline G_m$; equivalently, that the section of
$D(\mu_n)$ just considered lifts fppf-locally to a section of $D(F)$.
This is clear, since $D(F)\longrightarrow D(\mu_n)$ is faithfully flat
and of finite presentation.

\SGAvar{\medskip
\noindent b) By SGA 3 X 4.5, after replacing $S$ by an \`etale cover,
we may suppose that $Q$ is diagonalizable, say $Q=D(M)$, where $M$ is
a finitely generated $\mathbb Z$-module.  Let $T$ be the torsion
submodule of $M$, and put $L=M/T$, so that $L\simeq\mathbb Z^r$.  The
exact sequence
\[
 0\longrightarrow T\longrightarrow M\longrightarrow L
 \longrightarrow0
\]
then gives an exact sequence (SGA 3 IX 3.1)
\[
 0\longrightarrow D(L)\longrightarrow D(M)\longrightarrow D(L)
 \longrightarrow0.
\]
It is therefore enough to prove separately that the $\uExt^1$ of
$D(L)$ and of $D(T)$ by $\underline G_m$ vanish.  The case of $D(T)$
follows from a), and the case of $D(L)$ reduces to $L=\mathbb Z$.
Thus it remains to check that
\[
 \uExt^1(\underline G_m,\underline G_m)=0,
\]
that is, that every extension $E$ of $\underline G_m$ by
$\underline G_m$ is fppf-locally trivial.

}{\medskip
\noindent b) By SGA 3 X 4.5, after replacing $S$ by an \`etale cover,
we may suppose that $Q$ is diagonalizable, say $Q=D(M)$, where $M$ is
a finitely generated $\mathbb Z$-module.  Let $T$ be the torsion
submodule of $M$, and put $L=M/T$, so that $L\simeq\mathbb Z^r$.  The
exact sequence
\[
 0\longrightarrow T\longrightarrow M\longrightarrow L
 \longrightarrow0
\]
then gives an exact sequence (SGA 3 IX 3.1)
\[
 0\longrightarrow D(L)\longrightarrow D(M)\longrightarrow D(T)
 \longrightarrow0.
\]
It is therefore enough to prove separately that the $\uExt^1$ of
$D(L)$ and of $D(T)$ by $\underline G_m$ vanish.  The case of $D(T)$
follows from a), and the case of $D(L)$ reduces to $L=\mathbb Z$.
Thus it remains to check that
\[
 \uExt^1(\underline G_m,\underline G_m)=0,
\]
that is, that every extension $E$ of $\underline G_m$ by
$\underline G_m$ is fppf-locally trivial.

}{E7I-VIII-018}\SGApage{I}{\(\approx\)258}
Such an extension is represented by an affine flat group scheme over
$S$ (fpqc descent theory, SGA 1 VII) whose fibres are tori; it is
therefore a torus (SGA 3 X 4.9).  By SGA 3 X 4.5 we may suppose that
$E$ is a split torus
\[
 E=D(R)\simeq\underline G_m^{,2}.
\]
After a Zariski localization, we may also suppose that the
monomorphism $\underline G_m\hookrightarrow E=D(R)$ comes from a
constant morphism $R\longrightarrow\mathbb Z_S$ (SGA 3 VIII 1.4),
which is necessarily an epimorphism (SGA 3 VIII 3.2 b)).  Choosing
$r\in R$ above $1\in\mathbb Z$ gives a homomorphism
$E\longrightarrow\underline G_m$ splitting the extension, and proves
the assertion.

\medskip
\noindent\underline{Corollary \SGARefTarget{sga7.sga7i.eviii.n3-3-2}{3.3.2}.}  \underline{Let $Q$ be as in
\SGARefLink{sga7.sga7i.eviii.n3-3-1}{3.3.1}.  Then, for every $P$, there is a canonical isomorphism}
\begin{equation}\tag{\SGARefTarget{sga7.sga7i.eviii.n3-3-2-1}{3.3.2.1}}
 \Biext^1(P,Q;\underline G_m)
 \xleftarrow{\ \sim\ }\Ext^1(P,D(Q)),
\end{equation}
\underline{where $D(Q)$ denotes the Cartier dual
$\uHom(Q,\underline G_m)$ of $Q$.  More precisely, the canonical
functor (\SGARefLink{sga7.sga7i.eviii.n1-1-6}{1.1.6}) is an equivalence from the category of extensions of
$P$ by $D(Q)$ to the category of biextensions of $(P,Q)$ by
$\underline G_m$.}

This is the special case of \SGARefLink{sga7.sga7i.eviii.n1-5}{1.5} obtained by applying \SGARefLink{sga7.sga7i.eviii.n3-3-1}{3.3.1}.

\SGApage{I}{\(\approx\)259}
\medskip
\noindent\underline{Remarks \SGARefTarget{sga7.sga7i.eviii.n3-3-3}{3.3.3}.}  Under the hypotheses of \SGARefLink{sga7.sga7i.eviii.n3-3-1}{3.3.1},
$D(Q)$ is represented by a separated group scheme, locally of finite
presentation and locally quasi-finite over $S$.  More precisely, if
$Q$ is finite locally free over $S$, then so is $D(Q)$, as is well
known (SGA 3 VII 3).  If $Q$ is of multiplicative type, then $D(Q)$
is a \underline{finitely generated twisted constant group} (SGA 3 X
5.1), as was seen in SGA 3 IX 4.5 and X 5.7.  If $P$ is a group
scheme, every extension of $P$ by $D(Q)$ is therefore representable,
by a familiar lemma from fpqc descent theory (SGA 3 X 5.4).

\medskip
\noindent \SGARefTarget{sga7.sga7i.eviii.n3-3-4}{3.3.4}.  It is false that
$\uExt^1(Q,\underline G_m)=0$ for every affine flat group scheme $Q$
of finite presentation over $S$.  This already fails for
\[
 Q=\underline G_{a_S},\qquad
 S=\Spec k[t]/(t^2),
\]
where $k$ is a field of characteristic $p>0$.




\medskip
\noindent\underline{Proposition \SGARefTarget{sga7.sga7i.eviii.n3-4}{3.4}.}  \underline{Let $P$ be a smooth
group scheme of finite presentation over $S$, with connected fibres,
and let $Q$ be a torus over $S$.  Then the category of biextensions
of $(P,Q)$ by $\underline G_m$ is equivalent to the one-point
category; that is,}
\begin{equation}\tag{\SGARefTarget{sga7.sga7i.eviii.n3-4-1}{3.4.1}}
 \Biext^0(P,Q;\underline G_m)=
 \Biext^1(P,Q;\underline G_m)=0.
\end{equation}

\noindent\underline{Proof.}  Put $M=D(Q)$.  This is a twisted
constant group over $S$ whose fibres are isomorphic to groups
$\mathbb Z^r$ (\SGARefLink{sga7.sga7i.eviii.n3-3-3}{3.3.3}).  By \SGARefLink{sga7.sga7i.eviii.n3-3-2}{3.3.2}, the category in question is
equivalent to the category of extensions of $P$ by $M$.  This
category is \underline{rigid}: the automorphism group of any object
is the trivial group.  This is the first equality in (\SGARefLink{sga7.sga7i.eviii.n3-4-1}{3.4.1}), which
may also be written
\[
 \Hom(P,M)=0.
\]
It follows at once from the fact that $M$ is etale over $S$, whereas
$P$ has connected fibres.  It remains to prove
\[
 \Ext^1(P,M)=0,
\]
that is, that every extension $E$ of $P$ by $M$ is trivial.

\SGApage{I}{\(\approx\)260}
By the uniqueness of a splitting when one exists, it is enough to
prove that $E$ is locally trivial.  We may therefore reduce to the
case $M=\mathbb Z_S^r$, and then to $M=\mathbb Z_S$.  When $S$ is the
spectrum of a field, the conclusion follows from \SGARefLink{sga7.sga7i.eviii.n5-5}{5.5} (i) below.
(Notice that we shall not use \SGARefLink{sga7.sga7i.eviii.n3-4}{3.4} before IX 6, so in particular there
is no vicious circle.)  We now reduce the general case to this
special case.

\SGAvar{It is enough to find a section $f$ of $P$ over $P$ such that
$f(0)=0$, for such a section is necessarily multiplicative.  Indeed,
\[
 f(x+y)-f(x)-f(y)=0:
\]
the left-hand side defines a morphism
$P\times_S P\longrightarrow\mathbb Z_S$ which vanishes on the zero
section of $P\times_S P$.  It therefore vanishes identically, since
$P\times_S P$ has connected fibres (the geometric fibres of $P$
being connected).  The required conclusion now follows from the
following lemma.

}{It is enough to find a section $f$ of $E$ over $P$ such that
$f(0)=0$, for such a section is necessarily multiplicative.  Indeed,
\[
 f(x+y)-f(x)-f(y)=0:
\]
the left-hand side defines a morphism
$P\times_S P\longrightarrow\mathbb Z_S$ which vanishes on the zero
section of $P\times_S P$.  It therefore vanishes identically, since
$P\times_S P$ has connected fibres (the geometric fibres of $P$
being connected).  The required conclusion now follows from the
following lemma.

}{E7I-VIII-019}\medskip
\noindent\underline{Lemma \SGARefTarget{sga7.sga7i.eviii.n3-4-2}{3.4.2}.}  \underline{Let $S$ be a scheme,
$P$ a smooth $S$-scheme of finite presentation with geometrically
connected fibres, $Z$ a group, and $E$ a $Z_P$-torsor (for the fppf,
etale, or fpqc topology; these notions are equivalent here, cf. SGA
3 X 5.4).  Let $e$ be a section of $P$ over $S$, and let $e'$ be a
section of $E$ over $S$ above $e$ (that is, a section of $e^{-1}(E)$
over $S$).}

\SGApage{I}{\(\approx\)261}
\underline{The torsor $E$ is trivial if and only if, for every
$s\in S$, the induced torsor $E_s$ over $P_s$ is trivial.  In that
case there is a unique section $f$ of $E$ over $P$ such that
$f\circ e=e'$.}

\noindent\underline{Proof.}  Uniqueness is clear.  Two such sections
$f$ and $g$ ``differ'' by a morphism
\[
 h:P\longrightarrow Z_S,
 \qquad h\circ e=1,
\]
and such a morphism is constantly equal to $1$, since $P$ has
connected fibres.  Moreover, if a section $f$ of $E$ over $P$
exists, it can always be corrected by a morphism
$P\longrightarrow Z_S$ coming from a morphism $z:S\longrightarrow
Z_S$, so as to obtain a section $f'$ satisfying $f'\circ e=e'$.
One need only choose $z$ so that
\[
 f(e(t))=e'(t)z(t).
\]

Applying this observation fibrewise, for every $s\in S$ we obtain a
well-determined section $f_s$ of $E_s$ over $P_s$ taking $e_s$ to
$e'_s$.  It remains to prove that the $f_s$ arise from a section $f$
of $E$ over $P$.  Equivalently (EGA IV 17.9.3), we must prove that
\[
 U=\bigcup_{s\in S}f_s(P_s)\subset E
\]
is open in $E$.

Suppose first that $S$ is noetherian.  Then $U$ is constructible.  To
see this, for each $s\in S$ let $T$ be the closure of $s$, endowed
with its reduced induced structure.  There is an open neighbourhood
$T'$ of $s$ in $T$ and a section of $E_{T'}$ over $P_{T'}$ extending
$f_s$ and compatible with the marked sections $e$ and $e'$ (as
follows directly from EGA IV 8.8.2).  By the criterion of EGA IV
1.10.1, it remains to show that $U$ is stable under generization.
The usual argument using EGA II 7.1.9 reduces this to the case in
which $S$ is the spectrum of a discrete valuation ring.

Since $P$ is smooth over $S$, it is regular and hence normal; and
since it has connected fibres, it is connected.  The given section
of $E$ over the generic fibre of $P$ over $S$ therefore extends
uniquely to a section $f$ of $E$ over $P$ (EGA IV 18.10.19 and
18.10.20), necessarily compatible with $e$ and $e'$.  Thus in this
case $U=f(P)$ is open.  This proves \SGARefLink{sga7.sga7i.eviii.n3-4}{3.4} when $S$ is noetherian.  The
general case reduces readily to this one by the usual limit
arguments of EGA IV, 8 and 9; we leave the details to the reader.

\SGApage{I}{262}
\medskip
\noindent\underline{Remark \SGARefTarget{sga7.sga7i.eviii.n3-4-3}{3.4.3}.}  Lemma 3.4.2 can be ``crushed''
under learned references, namely SGA 4 XV 1.15 (the case $n=1$) and
\SGARefLink{sga7.sga7i.eviii.n4-1}{4.1}.




\medskip
\noindent\underline{Corollary \SGARefTarget{sga7.sga7i.eviii.n3-5}{3.5}.}  \underline{Let $P$ be a smooth
group scheme of finite presentation over $S$, with connected fibres,
and let}
\[
 0\longrightarrow Q'\longrightarrow Q\longrightarrow Q''
 \longrightarrow0
\]
\underline{be an exact sequence of commutative groups over $S$, with
$Q'$ a torus.  Then the inverse-image functor on biextensions}
\begin{equation}\tag{\SGARefTarget{sga7.sga7i.eviii.n3-5-1}{3.5.1}}
 \operatorname{BIEXT}(P,Q'';\underline G_m)
 \xrightarrow{\ \sim\ }
 \operatorname{BIEXT}(P,Q;\underline G_m)
\end{equation}
\underline{is an equivalence of categories.  In particular, the
following natural homomorphisms are isomorphisms:}
\begin{equation}\tag{\SGARefTarget{sga7.sga7i.eviii.n3-5-2}{3.5.2}}
 \Biext^i(P,Q'';\underline G_m)
 \xrightarrow{\ \sim\ }
 \Biext^i(P,Q;\underline G_m),
 \qquad i=0,1.
\end{equation}

Indeed, up to equivalence, $\operatorname{BIEXT}(P,Q'';\underline
G_m)$ identifies with the category of biextensions $E$ of $(P,Q)$ by
$\underline G_m$ \underline{equipped} with a trivialization of the
induced biextension $E'$ of $(P,Q')$ by $\underline G_m$ (VII
\SGARefLink{sga7.sga7i.evii.n3-7-6}{3.7.6}).  By \SGARefLink{sga7.sga7i.eviii.n3-4}{3.4} there is exactly one such trivialization of $E'$ for
each given biextension $E$.  One can of course also prove directly
that (\SGARefLink{sga7.sga7i.eviii.n3-5-2}{3.5.2}) is an isomorphism for $i=0,1$ - which is equivalent to
the conclusion of \SGARefLink{sga7.sga7i.eviii.n3-5}{3.5} - by using the exact sequence of the
$\Biext^i$ in VII \SGARefLink{sga7.sga7i.evii.n3-7-5}{3.7.5} associated with
\[
 0\longrightarrow Q'\longrightarrow Q\longrightarrow Q''
 \longrightarrow0.
\]

\SGApage{I}{263}
\SGAvar{\medskip
\noindent\underline{Remark \SGARefTarget{sga7.sga7i.eviii.n3-6}{3.6}.}  The symmetric statement obtained
from \SGARefLink{sga7.sga7i.eviii.n3-5}{3.5} by interchanging left and right is of course also valid; it
also follows formally from \SGARefLink{sga7.sga7i.eviii.n3-5}{3.5} by passing to symmetric biextensions
(VII \SGARefLink{sga7.sga7i.evii.n2-7}{2.7}).  Suppose, in addition to the data and hypotheses of \SGARefLink{sga7.sga7i.eviii.n3-5}{3.5},
that the symmetric data and hypotheses are given: namely, that $Q$
has connected fibres (equivalently, that $Q''$ has connected fibres,
since $Q'$ has connected fibres), and that there is an exact sequence
\[
 0\longrightarrow P'\longrightarrow P\longrightarrow P''
 \longrightarrow0
\]
with $P'$ a torus.  Successive application of \SGARefLink{sga7.sga7i.eviii.n3-5}{3.5} and of its
symmetric form, applied to this last sequence and to $Q''$, shows
that the functor
\begin{equation}\tag{\SGARefTarget{sga7.sga7i.eviii.n3-6-1}{3.6.1}}
 \operatorname{BIEXT}(P'',Q'';\underline G_m)
 \xrightarrow{\ \sim\ }
 \operatorname{BIEXT}(P,Q;\underline G_m)
\end{equation}
is an equivalence of categories.  In particular, the following
homomorphism is an isomorphism:
\begin{equation}\tag{\SGARefTarget{sga7.sga7i.eviii.n3-6-2}{3.6.2}}
 \Biext^i(P'',Q'';\underline G_m)
 \xrightarrow{\ \sim\ }
 \Biext^1(P,Q;\underline G_m),
 \qquad i=0,1.
\end{equation}

}{\medskip
\noindent\underline{Remark \SGARefTarget{sga7.sga7i.eviii.n3-6}{3.6}.}  The symmetric statement obtained
from \SGARefLink{sga7.sga7i.eviii.n3-5}{3.5} by interchanging left and right is of course also valid; it
also follows formally from \SGARefLink{sga7.sga7i.eviii.n3-5}{3.5} by passing to symmetric biextensions
(VII \SGARefLink{sga7.sga7i.evii.n2-7}{2.7}).  Suppose, in addition to the data and hypotheses of \SGARefLink{sga7.sga7i.eviii.n3-5}{3.5},
that the symmetric data and hypotheses are given: namely, that $Q$
has connected fibres (equivalently, that $Q''$ has connected fibres,
since $Q'$ has connected fibres), and that there is an exact sequence
\[
 0\longrightarrow P'\longrightarrow P\longrightarrow P''
 \longrightarrow0
\]
with $P'$ a torus.  Successive application of \SGARefLink{sga7.sga7i.eviii.n3-5}{3.5} and of its
symmetric form, applied to this last sequence and to $Q''$, shows
that the functor
\begin{equation}\tag{\SGARefTarget{sga7.sga7i.eviii.n3-6-1}{3.6.1}}
 \operatorname{BIEXT}(P'',Q'';\underline G_m)
 \xrightarrow{\ \sim\ }
 \operatorname{BIEXT}(P,Q;\underline G_m)
\end{equation}
is an equivalence of categories.  In particular, the following
homomorphism is an isomorphism:
\begin{equation}\tag{\SGARefTarget{sga7.sga7i.eviii.n3-6-2}{3.6.2}}
 \Biext^i(P'',Q'';\underline G_m)
 \xrightarrow{\ \sim\ }
 \Biext^i(P,Q;\underline G_m),
 \qquad i=0,1.
\end{equation}

}{E7I-VIII-020}\medskip
\noindent\underline{Proposition \SGARefTarget{sga7.sga7i.eviii.n3-7}{3.7}.}  \underline{Let $A$ be an
abelian scheme over $S$, let $T$ be a group scheme of multiplicative
type (SGA 3 IX) and of finite type over $S$, and let}
\[
 \underline M=\uHom(T,\underline G_m)
\]
\underline{be the Cartier dual of $T$ (a finitely generated twisted
constant group over $S$, by \SGARefLink{sga7.sga7i.eviii.n3-3-3}{3.3.3}).  Then the categories
$\operatorname{EXT}(A,T)$ and
$\operatorname{BIEXT}(A,\underline M;\underline G_m)$, formed
respectively by commutative extensions of $A$ by $T$ and by
biextensions of $(A,\underline M)$ by $\underline G_m$, are rigid and
canonically equivalent.  Their groups of isomorphism classes are
canonically isomorphic to $\Hom(\underline M,A^*)$, where}
\[
 A^*=\uExt^1(A,\underline G_m)
\]
\underline{is the dual abelian scheme of $A$ (\SGARefLink{sga7.sga7i.eviii.n3-2}{3.2}).  In particular,
there is a canonical isomorphism}
\begin{equation}\tag{\SGARefTarget{sga7.sga7i.eviii.n3-7-1}{3.7.1}}
 \Ext^1(A,T)\xrightarrow{\ \sim\ }
 \Hom(\underline M,A^*).
\end{equation}

\SGApage{I}{\(\approx\)264}
\SGAvar{This is the special case of \SGARefLink{sga7.sga7i.eviii.n1-6}{1.6} obtained by taking $P=A$ and
$Q=\underline M$.  By the biduality theorem (SGA 3 X 5.7), we do
indeed have
\[
 T\simeq\uHom(\underline M,\underline G_m),
\]
and it remains only to verify the conditions (\SGARefLink{sga7.sga7i.eviii.n1-6-1}{1.6.1}), which here read
\[
 \uHom(A,\underline G_m)=0,
 \qquad
 \Hom\bigl(A,\uExt^1(\underline M,\underline G_m)\bigr)=0.
\]
The first equality was recalled in (\SGARefLink{sga7.sga7i.eviii.n3-2-1}{3.2.1}); the second may be
strengthened to
\[
 \uExt^1(\underline M,\underline G_m)=0.
\]
This assertion is local on $S$, so we may suppose that $\underline
M$ is constant, $\underline M=M_S$.  Decomposing $M$ as a direct sum
of cyclic $\mathbb Z$-modules reduces us to $M=\mathbb Z$, which is
trivial, and to $\mathbb Z=\mathbb Z/n\mathbb Z$, which follows from the
$n$-divisibility of $\underline G_m$ (or, more learnedly, from
\SGARefLink{sga7.sga7i.eviii.n3-3-1}{3.3.1}).

}{This is the special case of \SGARefLink{sga7.sga7i.eviii.n1-6}{1.6} obtained by taking $P=A$ and
$Q=\underline M$.  By the biduality theorem (SGA 3 X 5.7), we do
indeed have
\[
 T\simeq\uHom(\underline M,\underline G_m),
\]
and it remains only to verify the conditions (\SGARefLink{sga7.sga7i.eviii.n1-6-1}{1.6.1}), which here read
\[
 \uHom(A,\underline G_m)=0,
 \qquad
 \Hom\bigl(A,\uExt^1(\underline M,\underline G_m)\bigr)=0.
\]
The first equality was recalled in (\SGARefLink{sga7.sga7i.eviii.n3-2-1}{3.2.1}); the second may be
strengthened to
\[
 \uExt^1(\underline M,\underline G_m)=0.
\]
This assertion is local on $S$, so we may suppose that $\underline
M$ is constant, $\underline M=M_S$.  Decomposing $M$ as a direct sum
of cyclic $\mathbb Z$-modules reduces us to $M=\mathbb Z$, which is
trivial, and to $M=\mathbb Z/n\mathbb Z$, which follows from the
$n$-divisibility of $\underline G_m$ (or, more learnedly, from
\SGARefLink{sga7.sga7i.eviii.n3-3-1}{3.3.1}).

}{E7I-VIII-021}When $T=\underline G_m$, and hence $\underline M=\mathbb Z_S$,
(\SGARefLink{sga7.sga7i.eviii.n3-7-1}{3.7.1}) reduces to the identity isomorphism (\SGARefLink{sga7.sga7i.eviii.n3-2-3}{3.2.3}).  Of course, a
direct proof of (\SGARefLink{sga7.sga7i.eviii.n3-7-1}{3.7.1}) would be essentially trivial, without
introducing the group
\[
 \Ext^1\bigl(A\mathbin{\mathop{\otimes}^{\mathbb L}}\underline M,
              \underline G_m\bigr)
 =\Biext^1(A,\underline M;\underline G_m).
\]




\SGApage{I}{\(\approx\)265}
\subsection*{\SGARefTarget{sga7.sga7i.eviii.n4}{4}. Extensions and biextensions of smooth connected group
schemes over a field}

We recall the following lemma due to Rosenlicht [3 bis, Proposition
3].

\medskip
\noindent\underline{Lemma \SGARefTarget{sga7.sga7i.eviii.n4-1}{4.1}.}  \underline{Let $k$ be a field, and
let $X$ and $Y$ be two separable, geometrically connected $k$-schemes
of finite type, equipped respectively with $k$-rational points $e_X$
and $e_Y$.  Then every morphism}
\[
 X\times Y\longrightarrow\underline G_m
\]
\underline{which is trivial on the two subschemes $X\times e_Y$ and
$e_X\times Y$ is trivial.}

We shall use this result only when $X$ and $Y$ are group schemes over
$k$; in that case the hypotheses imply that $X$ and $Y$ are in fact
smooth over $k$.  By induction, \SGARefLink{sga7.sga7i.eviii.n4-1}{4.1} also implies the analogous result
for any number of factors.  We may therefore apply the results of VII
\SGARefLink{sga7.sga7i.evii.n1-3}{1.3} and VII \SGARefLink{sga7.sga7i.evii.n2-9}{2.9}, and obtain the following two propositions.

\SGAvar{\medskip
\noindent\underline{Proposition \SGARefTarget{sga7.sga7i.eviii.n4-2}{4.2}.}  \underline{Let $P$ be a
smooth connected group over the field $k$.  The natural functor from
the category of group extensions of $P$ by $\underline G_m$ to the
category of line bundles on $P$ rigidified at the identity element
$e_P$ is fully faithful}.  \underline{Its essential image consists of the
rigidified line bundles $\mathbb L$ on $P$ for which}
\[
 \pi^*P\simeq
 \pr_1^*P\pr_2^*P,
\]
\underline{where $\pi:P\times P\longrightarrow P$ is the group law
of $P$, and $\pr_1,\pr_2$ are the two projections.}

}{\medskip
\noindent\underline{Proposition \SGARefTarget{sga7.sga7i.eviii.n4-2}{4.2}.}  \underline{Let $P$ be a
smooth connected group over the field $k$.  The natural functor from
the category of group extensions of $P$ by $\underline G_m$ to the
category of line bundles on $P$ rigidified at the identity element
$e_P$ is fully faithful}.  \underline{Its essential image consists of the
rigidified line bundles $\mathbb L$ on $P$ for which}
\[
 \pi^*\mathbb L\simeq
 \pr_1^*\mathbb L\otimes\pr_2^*\mathbb L,
\]
\underline{where $\pi:P\times P\longrightarrow P$ is the group law
of $P$, and $\pr_1,\pr_2$ are the two projections.}

}{E7I-VIII-022}\medskip
\noindent\underline{Proposition \SGARefTarget{sga7.sga7i.eviii.n4-3}{4.3}.}  \underline{Let $P$ and $Q$ be
two smooth connected algebraic groups over the field $k$.  The
natural functor from the category of biextensions of $(P,Q)$ by
$\underline G_m$ to the category of line bundles on $P\times Q$
birigidified with respect to the identity sections $e_P$ and $e_Q$ of
$P$ and $Q$ is fully faithful; moreover, both categories are rigid.
A birigidified line bundle $\mathbb L$ on $P\times Q$ comes from a
biextension if and only if its pullbacks to $P\times P\times Q$ and
$P\times Q\times Q$ by $\pi_P\times\id_Q$ and
$\id_P\times\pi_Q$ satisfy the ``theorem of the cube'' ([1] or
[2]).}

\SGApage{I}{\(\approx\)266}
This last condition is simply the translation of the conditions in
VII \SGARefLink{sga7.sga7i.evii.n2-9-3}{2.9.3} b), in the explicit form given in VII \SGARefLink{sga7.sga7i.evii.n3-5}{3.5} b).

\medskip
\noindent \SGARefTarget{sga7.sga7i.eviii.n4-4}{4.4}.  I do not know whether the cube condition in \SGARefLink{sga7.sga7i.eviii.n4-3}{4.3} is
automatically satisfied.  At least it is known [2, IV 2.6] that, for
a given $\mathbb L$, there is always an integer $n\geq1$, a power of
the characteristic exponent of $k$, such that
$\mathbb L^{\otimes n}$ satisfies that condition and hence comes from
a biextension of $(P,Q)$ by $\underline G_m$.  It is also always
possible to find a purely inseparable extension $k'$ of $k$ such that
the required condition becomes true after the base extension
$k\longrightarrow k'$; after this extension, therefore, $\mathbb L$
is a birigidified line bundle arising from a biextension of $(P,Q)$
by $\underline G_m$.

In particular, when $k$ is perfect, the condition in \SGARefLink{sga7.sga7i.eviii.n4-3}{4.3} is always
satisfied.  This last fact, which formally implies the preceding
statements by a trace argument, can also be seen by applying \SGARefLink{sga7.sga7i.eviii.n4-7}{4.7}
below; cf. \SGARefLink{sga7.sga7i.eviii.n4-9}{4.9}.




We recall, in the form of a lemma, the following result from EGA
$\mathrm{Err}_{\mathrm{IV}}$ 53, 21.4.13.

\SGApage{I}{267}
\medskip
\noindent\underline{Lemma \SGARefTarget{sga7.sga7i.eviii.n4-5}{4.5}.}  \underline{Let
$f:X\longrightarrow Y$ be a smooth surjective morphism with integral
fibres, where $Y$ is normal and integral with generic point $y$.  Let
$\mathbb L$ be a line bundle on $X$ whose restriction to $X_y$ is
trivial.  Then there is a line bundle $\mathbb M$ on $Y$ such that}
\[
 \mathbb L\simeq f^*\mathbb M.
\]

\medskip
\noindent\underline{Proposition \SGARefTarget{sga7.sga7i.eviii.n4-6}{4.6}.}  \underline{Let $P,Q$ be two
smooth connected algebraic groups over the field $k$, with $P$
affine.  Suppose in addition either that $P$ admits a composition
series each of whose factors is a torus or is isomorphic to the
additive group $\underline G_a$ (a condition satisfied in particular
when $k$ is perfect [SGA 3, XVII 7.2.1, 4.1.1 (i)$\Longleftrightarrow$
(ii), and 4.1.5]), or that $Q$ is an extension of an abelian variety
$B$ by a smooth connected affine algebraic group satisfying the
preceding condition.  Under either hypothesis, every birigidified
line bundle $\mathbb L$ on $P\times Q$ is trivial; a fortiori, by
\SGARefLink{sga7.sga7i.eviii.n4-3}{4.3}, every biextension of $(P,Q)$ by $\underline G_m$ is trivial.}

\noindent\underline{Proof.}  First observe that every automorphism of
a birigidified line bundle is the identity (\SGARefLink{sga7.sga7i.eviii.n4-3}{4.3}).  Descent theory
therefore shows that it is enough to prove that $\mathbb L$ becomes
trivial after a finite separable extension of $k$.  Passing to the
limit immediately reduces us to the case where $k$ is separably
closed.

Every torus over such a field is split.  Hence, if $P$ admits a
composition series as in \SGARefLink{sga7.sga7i.eviii.n4-6}{4.6}, it also admits one whose quotients are
isomorphic to $\underline G_a$ or $\underline G_m$.  Repeated
application of \SGARefLink{sga7.sga7i.eviii.n4-5}{4.5}, together with
\[
 \Pic(\underline G_{a_K})=\Pic(\underline G_{m_K})=0
\]
for every field $K$, gives
\[
 \Pic(P)=0.
\]
This equality remains true after every extension of the base field
$k$.  Applying \SGARefLink{sga7.sga7i.eviii.n4-5}{4.5} to
\[
 \pr_2:P\times Q\longrightarrow Q
\]
and to $\mathbb L$, we find a line bundle $\mathbb M$ on $Q$ such
that $\mathbb L\simeq\pr_2^*\mathbb M$.  Since the restriction of
$\mathbb L$ to $e_P\times Q$ is trivial, $\mathbb M$ is trivial, and
so is $\mathbb L$.  This proves the first case.

\SGApage{I}{\(\approx\)268}
\SGAvar{In the second case, let $Q$ be an extension of an abelian scheme $B$
by a group $M$ admitting a composition series with quotients
isomorphic to $\underline G_m$ or $\underline G_a$.  Apply \SGARefLink{sga7.sga7i.eviii.n4-5}{4.5} to
\[
 f:X=P\times Q\longrightarrow Y=P\times B.
\]
The fibre over the generic point $y$ of $P\times B$ is an
$M_y$-torsor.  It is necessarily trivial by the particular structure
of $M$ and the classical equalities
\[
 H^1(K,\underline G_a)=0,
 \qquad
 H^1(K,\underline G_m)=0,
\]
valid for every field $K$, here taken to be $k(y)$.  Since
$\Pic(M)=0$, we have $\Pic(X_y)=0$; thus \SGARefLink{sga7.sga7i.eviii.n4-5}{4.5} applies and shows that
$\mathbb L$ is the pullback of a line bundle $\mathbb N$ on
$Y=P\times A$.

}{In the second case, let $Q$ be an extension of an abelian scheme $B$
by a group $M$ admitting a composition series with quotients
isomorphic to $\underline G_m$ or $\underline G_a$.  Apply \SGARefLink{sga7.sga7i.eviii.n4-5}{4.5} to
\[
 f:X=P\times Q\longrightarrow Y=P\times B.
\]
The fibre over the generic point $y$ of $P\times B$ is an
$M_y$-torsor.  It is necessarily trivial by the particular structure
of $M$ and the classical equalities
\[
 H^1(K,\underline G_a)=0,
 \qquad
 H^1(K,\underline G_m)=0,
\]
valid for every field $K$, here taken to be $k(y)$.  Since
$\Pic(M)=0$, we have $\Pic(X_y)=0$; thus \SGARefLink{sga7.sga7i.eviii.n4-5}{4.5} applies and shows that
$\mathbb L$ is the pullback of a line bundle $\mathbb N$ on
$Y=P\times B$.

}{E7I-VIII-023}Let $\mathbb N_1$ and $\mathbb N_2$ be the restrictions of $\mathbb
N$ to $P\times e_B$ and $e_P\times B$, respectively.  Replacing
$\mathbb N$ by
\[
 \mathbb N\otimes\pr_1^*(\mathbb N_1)^{-1}
 \otimes\pr_2^*(\mathbb N_2)^{-1},
\]
we may suppose that $\mathbb N$ is birigidified and that $\mathbb L$
is its birigidified pullback.  We are thereby reduced to the case in
which $Q$ itself is an abelian scheme.

Then $\mathbb L$ defines a canonical morphism
\begin{equation}\tag{*}
 P\longrightarrow\underline{\Pic}_{Q/k}
\end{equation}
taking the identity to the identity.  Since $P$ is connected, this
morphism factors through the dual abelian scheme
\[
 Q'=\underline{\Pic}^{,0}_{Q/k}.
\]
It is well known that every morphism of schemes from a smooth
connected affine algebraic group to an abelian scheme is constant.
This is a geometric assertion, so we may suppose $k$ algebraically
closed; a devissage then reduces it immediately to $P=\underline
G_a$ or $P=\underline G_m$, and every rational map from
$\mathbb P^1_k$ to an abelian scheme over $k$ is constant.

\SGApage{I}{\(\approx\)269}
\SGAvar{The constancy, hence vanishing, of (*) means that $\mathbb L$ is the
pullback of a line bundle $\mathbb M$ on $P$.  Since $\mathbb L$ is
trivial on $P\times e_A$, $\mathbb M$ is trivial, and therefore so is
$\mathbb L$.  This completes the proof.




\medskip
\noindent\underline{Proposition \SGARefTarget{sga7.sga7i.eviii.n4-7}{4.7}.}  \underline{Let $P$ and $Q$
be two smooth connected algebraic groups over the field $k$.
Suppose that $P$ is an extension of an abelian scheme $A$ by a smooth
connected affine algebraic group $L$ admitting a composition series
each of whose factors is isomorphic to a torus or to $\underline
G_a$} (\underline{a condition always satisfied when $k$ is perfect}).  \underline{Then the
functor considered in \SGARefLink{sga7.sga7i.eviii.n1-3}{1.3} is an equivalence of categories; that is,
every birigidified line bundle $\mathbb L$ on $P\times Q$ comes from
a biextension of $(P,Q)$.  Moreover, the functor between the rigid
categories}
\[
 \operatorname{BIEXT}(A,Q;\underline G_m)
 \longrightarrow
 \operatorname{BIEXT}(P,Q;\underline G_m)
\]
\underline{is an equivalence of categories; in other words, it
induces an isomorphism of groups}
\begin{equation}\tag{\SGARefTarget{sga7.sga7i.eviii.n4-7-1}{4.7.1}}
 \Biext^1(A,Q;\underline G_m)
 \xrightarrow{\ \sim\ }
 \Biext^1(P,Q;\underline G_m).
\end{equation}

}{The constancy, hence vanishing, of (*) means that $\mathbb L$ is the
pullback of a line bundle $\mathbb M$ on $P$.  Since $\mathbb L$ is
trivial on $P\times e_Q$, $\mathbb M$ is trivial, and therefore so is
$\mathbb L$.  This completes the proof.




\medskip
\noindent\underline{Proposition \SGARefTarget{sga7.sga7i.eviii.n4-7}{4.7}.}  \underline{Let $P$ and $Q$
be two smooth connected algebraic groups over the field $k$.
Suppose that $P$ is an extension of an abelian scheme $A$ by a smooth
connected affine algebraic group $L$ admitting a composition series
each of whose factors is isomorphic to a torus or to $\underline
G_a$} (\underline{a condition always satisfied when $k$ is perfect}).  \underline{Then the
functor considered in \SGARefLink{sga7.sga7i.eviii.n1-3}{4.3} is an equivalence of categories; that is,
every birigidified line bundle $\mathbb L$ on $P\times Q$ comes from
a biextension of $(P,Q)$.  Moreover, the functor between the rigid
categories}
\[
 \operatorname{BIEXT}(A,Q;\underline G_m)
 \longrightarrow
 \operatorname{BIEXT}(P,Q;\underline G_m)
\]
\underline{is an equivalence of categories; in other words, it
induces an isomorphism of groups}
\begin{equation}\tag{\SGARefTarget{sga7.sga7i.eviii.n4-7-1}{4.7.1}}
 \Biext^1(A,Q;\underline G_m)
 \xrightarrow{\ \sim\ }
 \Biext^1(P,Q;\underline G_m).
\end{equation}

}{E7I-VIII-024}For the first assertion, descent and \SGARefLink{sga7.sga7i.eviii.n4-3}{4.3} again reduce us to the case
where $k$ is separably closed.  Proceeding as in the proof of \SGARefLink{sga7.sga7i.eviii.n4-6}{4.6},
we deduce from \SGARefLink{sga7.sga7i.eviii.n4-5}{4.5} and the hypothesis on $L$ that $\mathbb L$ is
isomorphic to the pullback of a birigidified line bundle $\mathbb M$
on $A\times Q$.  For the first assertion of \SGARefLink{sga7.sga7i.eviii.n4-7}{4.7} we are thus reduced
to the case $A=P$.  The line bundle $\mathbb L$ then defines a
morphism of pointed schemes
\[
 f:Q\longrightarrow\underline{\Pic}_{P/k},
\]
which necessarily factors through the dual abelian scheme $P'$ of
$P$.  Let $\mathbb L'$ be the pullback of the Weil sheaf on
$P\times P'$ along
\[
 \id_P\times f:P\times Q\longrightarrow P\times P'.
\]
This is a birigidified line bundle which, by the definition of $f$,
differs from $\mathbb L$ only by the pullback of a line bundle
$\mathbb M$ on $P$.  Since both $\mathbb L'$ and $\mathbb L$ have
trivial restriction to $P\times e_Q$, it follows that $\mathbb M$ is
trivial and hence that $\mathbb L\simeq\mathbb L'$.  We noted in VII
\SGARefLink{sga7.sga7i.evii.n2-9-4}{2.9.4} that divisorial correspondences on a product of two abelian
schemes arise from biextensions.  Consequently $\mathbb L'$, and
therefore $\mathbb L$, comes from a biextension.

\SGApage{I}{\(\approx\)270}
Returning to the general case, where $P$ is an extension of $A$ by
$L$, it remains to prove that (\SGARefLink{sga7.sga7i.eviii.n4-7-1}{4.7.1}) is an isomorphism.  This follows
from the exact sequence of the $\Biext^i$, for $i=0,1$, associated
with
\[
 0\longrightarrow L\longrightarrow P\longrightarrow A
 \longrightarrow0,
\]
and from the relations
\[
 \Biext^0(L,Q;\underline G_m)=0,
 \qquad
 \Biext^1(L,Q;\underline G_m)=0.
\]
The first is a special case of the rigidity assertion in \SGARefLink{sga7.sga7i.eviii.n4-3}{4.3}, and
the second is precisely the first case of \SGARefLink{sga7.sga7i.eviii.n4-6}{4.6}.

\SGApage{I}{271}
\medskip
\noindent\underline{Corollary \SGARefTarget{sga7.sga7i.eviii.n4-8}{4.8}.}  \underline{Under the hypotheses
of \SGARefLink{sga7.sga7i.eviii.n4-7}{4.7}, suppose that $Q$ is an extension of an algebraic group $B$}
(\underline{necessarily smooth and connected}) \underline{by a smooth connected affine
algebraic group $M$.  Then the functor between the rigid categories}
\[
 \operatorname{BIEXT}(A,B;\underline G_m)
 \longrightarrow
 \operatorname{BIEXT}(P,Q;\underline G_m)
\]
\underline{is an equivalence of categories; in other words, it
induces an isomorphism of groups}
\begin{equation}\tag{\SGARefTarget{sga7.sga7i.eviii.n4-8-1}{4.8.1}}
 \Biext^1(A,B;\underline G_m)
 \xrightarrow{\ \sim\ }
 \Biext^1(P,Q;\underline G_m).
\end{equation}

In view of (\SGARefLink{sga7.sga7i.eviii.n4-7-1}{4.7.1}), it remains to prove that
\[
 \Biext^1(A,B;\underline G_m)
 \xrightarrow{\ \sim\ }
 \Biext^1(A,Q;\underline G_m)
\]
is an isomorphism.  As in the proof of (\SGARefLink{sga7.sga7i.eviii.n4-7-1}{4.7.1}), this follows from the
exact sequence of the $\Biext^i$ associated with
\[
 0\longrightarrow M\longrightarrow Q\longrightarrow B
 \longrightarrow0,
\]
together with the relations
\[
 \Biext^0(P,M;\underline G_m)=0,
 \qquad
 \Biext^1(P,M;\underline G_m)=0.
\]
Indeed, (\SGARefLink{sga7.sga7i.eviii.n4-7-1}{4.7.1}), applied with $M$ in place of $Q$, identifies these
groups with the corresponding groups having $A$ in the first
variable, as required by the preceding exact sequence.  The second
vanishing displayed above is a special case of the second case of
\SGARefLink{sga7.sga7i.eviii.n4-6}{4.6}, after interchanging the two factors.

\SGApage{I}{272}
\medskip
\noindent\underline{Remark \SGARefTarget{sga7.sga7i.eviii.n4-9}{4.9}.}  Let $P$ and $Q$ be two smooth
connected algebraic groups over a \underline{perfect} field $k$.
It is then known [3] that each is an extension of an abelian scheme
by a smooth connected affine group, the latter being the product of
a torus with an iterated extension of copies of $\underline G_a$
(cf. the reference given in \SGARefLink{sga7.sga7i.eviii.n4-6}{4.6}):
\[
 0\longrightarrow L\longrightarrow P\longrightarrow A
 \longrightarrow0,
 \qquad
 0\longrightarrow M\longrightarrow Q\longrightarrow B
 \longrightarrow0.
\]
Thus \SGARefLink{sga7.sga7i.eviii.n4-8}{4.8} applies and shows that the theory of biextensions of $(P,Q)$
by $\underline G_m$ reduces to the analogous (and much more
classical!) theory for the pair $(A,B)$ of abelian varieties.

\medskip
\noindent\underline{Corollary \SGARefTarget{sga7.sga7i.eviii.n4-10}{4.10}.}  \underline{Let $P$ and $Q$ be
two smooth connected algebraic groups over $k$, let $L$ be a smooth
connected affine algebraic subgroup of $P$, and let $E$ be a
biextension of $(P,Q)$ by $\underline G_m$.  By \SGARefLink{sga7.sga7i.eviii.n2-2}{2.2}, $E$ induces a
pairing}
\[
 T_\ell(P)\times T_\ell(Q)\longrightarrow T_\ell(\underline G_m).
\]
\underline{Then the induced pairing}
\[
 T_\ell(L)\times T_\ell(Q)\longrightarrow T_\ell(\underline G_m)
\]
\underline{is zero.}

Indeed, we may plainly suppose that $k$ is algebraically closed.  We
are then under the hypotheses of \SGARefLink{sga7.sga7i.eviii.n4-6}{4.6}, which imply that the
restriction of $E$ to $(L,Q)$ is the trivial biextension.  The
asserted conclusion follows, since the induced pairing in question
is the pairing defined by this restriction.



\SGApage{I}{273}
\medskip
\noindent \SGARefTarget{sga7.sga7i.eviii.n4-11}{4.11}.  Let $P$ and $Q$ be algebraic groups smooth over
$k$, let $E$ be a biextension of $(P,Q)$ by $\underline G_m$, and
let $T$ be the largest subtorus of $P$.  We shall determine the left
annihilator of the pairing
\begin{equation}\tag{\SGARefTarget{sga7.sga7i.eviii.n4-11-1}{4.11.1}}
 T_\ell(P)\times T_\ell(Q)
 \longrightarrow T_\ell(\underline G_{m,k}),
\end{equation}
where $\ell$ is again a fixed prime number.  By \SGARefLink{sga7.sga7i.eviii.n4-10}{4.10} this annihilator
always contains the subgroup $T_\ell(T)$ of $T_\ell(P)$.

Choose an extension $k'$ of $k$ over which $P$ and $Q$ can each be
written as an extension of an abelian scheme by an affine group
scheme that is an iterated extension of tori and copies of
$\underline G_a$.  If $P$ and $Q$ have unipotent rank zero, we may
take $k'=k$; in general, we may take for $k'$ the perfect closure of
$k$ (or a finite radicial extension).  Let $A'$ and $B'$ be the
abelian parts of $P_{k'}$ and $Q_{k'}$, respectively.  By \SGARefLink{sga7.sga7i.eviii.n4-7}{4.7}, the
biextension $E_{k'}$ of $(P_{k'},Q_{k'})$ comes from a biextension
$F'$ of $(A',B')$, uniquely determined up to unique isomorphism.
Consequently the pairing
\begin{equation}\tag{\SGARefTarget{sga7.sga7i.eviii.n4-11-2}{4.11.2}}
 T_\ell(P_{k'})\times T_\ell(Q_{k'})
 \longrightarrow T_\ell(\underline G_{m,k'}),
\end{equation}
obtained from (\SGARefLink{sga7.sga7i.eviii.n4-11-1}{4.11.1}) by the base-field extension $k\to k'$, and
also defined by the biextension $E_{k'}$, is obtained from the
pairing
\begin{equation}\tag{\SGARefTarget{sga7.sga7i.eviii.n4-11-3}{4.11.3}}
 T_\ell(A')\times T_\ell(B')
 \longrightarrow T_\ell(\underline G_{m,k'})
\end{equation}
defined by $F'$, by composing with the canonical homomorphisms
(which are in fact epimorphisms if $P$ and $Q$ are
$\ell$-divisible)
\begin{equation}\tag{\SGARefTarget{sga7.sga7i.eviii.n4-11-4}{4.11.4}}
 T_\ell(P_{k'})\longrightarrow T_\ell(A'),
 \qquad
 T_\ell(Q_{k'})\longrightarrow T_\ell(B').
\end{equation}

\SGAvar{It follows that the left annihilator of (\SGARefLink{sga7.sga7i.eviii.n4-11-2}{4.11.2}), which is obtained
from the left annihilator of (\SGARefLink{sga7.sga7i.eviii.n4-11-1}{4.11.1}) by base change, is the inverse
image under the first homomorphism in (\SGARefLink{sga7.sga7i.eviii.n4-11-4}{4.11.4}) of the left
annihilator of (\SGARefLink{sga7.sga7i.eviii.n4-11-3}{4.11.3}).  Let $C'$ be the largest abelian subscheme
of $A'$ such that the restriction of $F'$ to $(C',B')$ is the
trivial biextension.  Equivalently, $C'$ is the reduced identity
component of the kernel of the homomorphism
\begin{equation}\tag{\SGARefTarget{sga7.sga7i.eviii.n4-11-5}{4.11.5}}
 A'\longrightarrow (B')^*
\end{equation}
defined by $F'$, where $(B')^*$ is the dual abelian scheme of $B'$.
The left annihilator of (\SGARefLink{sga7.sga7i.eviii.n4-11-3}{4.11.3}) identifies with $T_\ell(C')$;
therefore the left annihilator of (\SGARefLink{sga7.sga7i.eviii.n4-11-1}{1.11.1}) is precisely the inverse
image of
\[
 T_\ell(C')\lhook\joinrel\longrightarrow T_\ell(A')
\]
under the first morphism in (\SGARefLink{sga7.sga7i.eviii.n4-11-4}{4.11.4}).

}{It follows that the left annihilator of (\SGARefLink{sga7.sga7i.eviii.n4-11-2}{4.11.2}), which is obtained
from the left annihilator of (\SGARefLink{sga7.sga7i.eviii.n4-11-1}{4.11.1}) by base change, is the inverse
image under the first homomorphism in (\SGARefLink{sga7.sga7i.eviii.n4-11-4}{4.11.4}) of the left
annihilator of (\SGARefLink{sga7.sga7i.eviii.n4-11-3}{4.11.3}).  Let $C'$ be the largest abelian subscheme
of $A'$ such that the restriction of $F'$ to $(C',B')$ is the
trivial biextension.  Equivalently, $C'$ is the reduced identity
component of the kernel of the homomorphism
\begin{equation}\tag{\SGARefTarget{sga7.sga7i.eviii.n4-11-5}{4.11.5}}
 A'\longrightarrow (B')^*
\end{equation}
defined by $F'$, where $(B')^*$ is the dual abelian scheme of $B'$.
The left annihilator of (\SGARefLink{sga7.sga7i.eviii.n4-11-3}{4.11.3}) identifies with $T_\ell(C')$;
therefore the left annihilator of (\SGARefLink{sga7.sga7i.eviii.n4-11-1}{4.11.1}) is precisely the inverse
image of
\[
 T_\ell(C')\lhook\joinrel\longrightarrow T_\ell(A')
\]
under the first morphism in (\SGARefLink{sga7.sga7i.eviii.n4-11-4}{4.11.4}).

}{E7I-VIII-025}\SGApage{I}{\(\approx\)274}
In particular, to say that the left annihilator of (\SGARefLink{sga7.sga7i.eviii.n4-11-1}{4.11.1}) is
exactly $T_\ell(T)$ -- equivalently, that the left annihilator of
(\SGARefLink{sga7.sga7i.eviii.n4-11-2}{4.11.2}) is the kernel $T_\ell(T')$ of the first morphism in
(\SGARefLink{sga7.sga7i.eviii.n4-11-4}{4.11.4}) -- is to say that (\SGARefLink{sga7.sga7i.eviii.n4-11-3}{4.11.3}) is nondegenerate on the left,
that is, that $C'=0$, or equivalently that (\SGARefLink{sga7.sga7i.eviii.n4-11-5}{4.11.5}) has finite
kernel.  In this last form the condition is visibly independent of
$\ell$, and in any of its forms it is plainly independent of the
choice of $k'$.  When it holds, we shall simply say that the given
biextension $E$ of $(P,Q)$ by $\underline G_m$ is
\underline{nondegenerate on the left}.  By duality, this also means
that the homomorphism
\begin{equation}\tag{\SGARefTarget{sga7.sga7i.eviii.n4-11-6}{4.11.6}}
 B'\longrightarrow (A')^*
\end{equation}
defined by $F'$ is an epimorphism, that is, is surjective.

The notion of a biextension \underline{nondegenerate on the right}
is defined symmetrically: (\SGARefLink{sga7.sga7i.eviii.n4-11-6}{4.11.6}) must have finite kernel, or,
equivalently, (\SGARefLink{sga7.sga7i.eviii.n4-11-5}{4.11.5}) must be an epimorphism.  The biextension $E$
of $(P,Q)$ by $\underline G_{m,k}$ is called
\underline{nondegenerate} if it is nondegenerate on both the left
and the right.  This is equivalent to saying that (\SGARefLink{sga7.sga7i.eviii.n4-11-5}{4.11.5}), or
equivalently (\SGARefLink{sga7.sga7i.eviii.n4-11-6}{4.11.6}), is an isogeny.

\SGApage{I}{\(\approx\)275}
It is immediate from the definition that these notions are
invariant under every extension of the base field $k$.



\medskip
\noindent \SGARefTarget{sga7.sga7i.eviii.n4-12}{4.12}.  Let $P$ be a smooth connected group scheme, let
$\mathbb L$ be a line bundle on $P$, and put
\begin{equation}\tag{\SGARefTarget{sga7.sga7i.eviii.n4-12-1}{4.12.1}}
 \delta(\mathbb L)=
 \pi^*\mathbb L\otimes
 \pr_1^*(\mathbb L)^{-1}\otimes
 \pr_2^*(\mathbb L)^{-1},
\end{equation}
with the notation of \SGARefLink{sga7.sga7i.eviii.n4-2}{4.2}.  As already observed in VII \SGARefLink{sga7.sga7i.evii.n1-3-4}{1.3.4}, this is
a birigidified line bundle on $P\times P$.  By \SGARefLink{sga7.sga7i.eviii.n4-3}{4.3}, if
$\delta(\mathbb L)$ comes from a biextension, that biextension is
determined up to unique isomorphism.  In particular, by \SGARefLink{sga7.sga7i.eviii.n4-7}{4.7}, it is
defined when $P$ is an extension of an abelian scheme $A$ by an
iterated extension of tori and copies of $\underline G_a$.  Whenever
it is defined, this biextension is plainly symmetric.

\SGAvar{In the preceding favourable case, and assuming that the maximal
torus of $P$ is split, the construction of $\delta(\mathbb L)$ can
be reduced to the case where $P$ is an abelian scheme.  Indeed, the
construction is functorial in $P$, while the homomorphism
\begin{equation}\tag{\SGARefTarget{sga7.sga7i.eviii.n4-12-2}{4.12.2}}
 \Pic(A)\longrightarrow\Pic(P)
\end{equation}
induced by the canonical quotient homomorphism
$f:A \longrightarrow P$ is surjective by \SGARefLink{sga7.sga7i.eviii.n4-5}{4.5} (cf. the beginning of
the proof of \SGARefLink{sga7.sga7i.eviii.n4-6}{4.6}).  We can therefore find an isomorphism
\[
 \mathbb L^*\simeq f^*\mathbb M,
\]
where $\mathbb M$ is a line bundle on $\mathbb M$.  Then
$\delta(\mathbb L)$ is the pullback of the biextension
$\delta(\mathbb M)$ of $(A,A)$ by $\underline G_{m,k}$.

}{In the preceding favourable case, and assuming that the maximal
torus of $P$ is split, the construction of $\delta(\mathbb L)$ can
be reduced to the case where $P$ is an abelian scheme.  Indeed, the
construction is functorial in $P$, while the homomorphism
\begin{equation}\tag{\SGARefTarget{sga7.sga7i.eviii.n4-12-2}{4.12.2}}
 \Pic(A)\longrightarrow\Pic(P)
\end{equation}
induced by the canonical quotient homomorphism
$f:P\longrightarrow A$ is surjective by \SGARefLink{sga7.sga7i.eviii.n4-5}{4.5} (cf. the beginning of
the proof of \SGARefLink{sga7.sga7i.eviii.n4-6}{4.6}).  We can therefore find an isomorphism
\[
 \mathbb L\simeq f^*\mathbb M,
\]
where $\mathbb M$ is a line bundle on $A$.  Then
$\delta(\mathbb L)$ is the pullback of the biextension
$\delta(\mathbb M)$ of $(A,A)$ by $\underline G_{m,k}$.

}{E7I-VIII-026}\SGApage{I}{276}
\medskip
\noindent\underline{Proposition \SGARefTarget{sga7.sga7i.eviii.n4-13}{4.13}.}  \underline{Under the
hypotheses of \SGARefLink{sga7.sga7i.eviii.n4-12}{4.12}, suppose that $\mathbb L$ is an ample line bundle
on $P$.  If the biextension $\delta(\mathbb L)$ of $(P,P)$ by
$\underline G_{m,k}$ is defined, then it is nondegenerate in the
sense of \SGARefLink{sga7.sga7i.eviii.n4-11}{4.11}.  If $P$ is an extension of an abelian scheme $A$ by a
smooth connected affine group $L$ satisfying the condition of \SGARefLink{sga7.sga7i.eviii.n4-7}{4.7},
then the biextension of $(A,A)$ from which $\delta(\mathbb L)$ comes
by \SGARefLink{sga7.sga7i.eviii.n4-8}{4.8} is of the form $\delta(\mathbb M)$, where $\mathbb M$ is an
ample line bundle on $A$.}

We reduce to the case where $k$ is algebraically closed.  By the
definition, it is then enough, with the notation of \SGARefLink{sga7.sga7i.eviii.n4-12}{4.12}, to prove
that the biextension $\delta(\mathbb M)$ of $(A,A)$ is
nondegenerate.  M. Raynaud proves in his thesis [2, XI 1.11] that if
$\mathbb L=f^*\mathbb M$ is ample, then $\mathbb M$ is ample as
well.  The theory of A. Weil then gives the familiar conclusion that
$\delta(\mathbb M)$ is nondegenerate [1, 2], that is, that the
homomorphism from $A$ to its dual abelian scheme is an isogeny.

\medskip
\noindent\underline{Remark \SGARefTarget{sga7.sga7i.eviii.n4-14}{4.14}.}  Suppose we are in the favourable
case where $P$ is an extension of an abelian scheme $A$ by a group
$L$ which is an iterated extension of tori and copies of
$\underline G_a$, and consider the canonical composite
\begin{equation}\tag{\SGARefTarget{sga7.sga7i.eviii.n4-14-1}{4.14.1}}
 \operatorname{NS}(A)
 \xrightarrow{\ \delta\ }
 \Biext^1(A,A;\underline G_m)
 \xrightarrow{\ \sim\ }
 \Biext^1(P,P;\underline G_m).
\end{equation}
By \SGARefLink{sga7.sga7i.eviii.n4-8}{4.8} the second arrow is bijective.  On the other hand, the theory
of abelian varieties shows that the first arrow identifies the
Neron--Severi group $\operatorname{NS}(A)$ of $A$ with the subgroup
of
\[
 \Biext^1(A,A;\underline G_m)=\operatorname{Corr}(A,A)
\]
consisting of the classes of symmetric biextensions, or equivalently
symmetric divisorial correspondences.  Thus (\SGARefLink{sga7.sga7i.eviii.n4-14-1}{4.14.1}) identifies
$\operatorname{NS}(A)$ with the group of symmetric biextensions of
$(P,P)$ by $\underline G_m$.  By \SGARefLink{sga7.sga7i.eviii.n4-7}{4.7} this group can equally be
described as the group of classes of symmetric divisorial
correspondences on $(P,P)$.

\SGApage{I}{\(\approx\)277}
It is natural to call a divisorial correspondence on $(P,P)$
\underline{ample}, and likewise to call the corresponding
biextension of $(P,P)$ by $\underline G_m$ \underline{ample}, when
it comes from an ample element of $\operatorname{NS}(A)$.  This
condition is plainly invariant under extension of the base field;
the notion therefore still makes sense without the preceding
restrictive hypothesis on $P$, assuming only that $P$ is smooth and
connected.

Suppose further that $L$ is ``split solvable over $k$'', in the
sense that it admits a devissage by copies of $\underline G_a$ and
$\underline G_m$.  The preceding observation and the surjectivity of
(\SGARefLink{sga7.sga7i.eviii.n4-12-2}{4.12.2}) show that the image of
\begin{equation}\tag{\SGARefTarget{sga7.sga7i.eviii.n4-14-2}{4.14.2}}
 \Pic(P)\xrightarrow{\ \delta\ }
 \Biext^1(P,P;\underline G_m)
\end{equation}
consists of the symmetric biextensions, and that its kernel is the
subgroup $\Pic^0(P)$ of $\Pic(P)$ which is the image of $\Pic^0(A)$.
Consequently, if we put
\[
 \operatorname{NS}(P)=\Pic(P)/\Pic^0(P)
 \simeq\operatorname{NS}(A),
\]
we may also interpret (\SGARefLink{sga7.sga7i.eviii.n4-14-1}{4.14.1}) as a canonical isomorphism
\begin{equation}\tag{\SGARefTarget{sga7.sga7i.eviii.n4-14-3}{4.14.3}}
 \operatorname{NS}(P)
 \xrightarrow{\ \sim\ }
 \Biext^1(P,P;\underline G_m)^{\mathrm{sym}}.
\end{equation}

\SGApage{I}{\(\approx\)278}
\SGAvar{Finally, the theorem of Raynaud cited in \SGARefLink{sga7.sga7i.eviii.n4-13}{4.13} says that an element
$\xi$ of $\Pic(P)$ is ample if and only if its image
$\delta(\xi)$ under (\SGARefLink{sga7.sga7i.eviii.n4-14-2}{4.14.1}) is the class of an ample biextension.



}{Finally, the theorem of Raynaud cited in \SGARefLink{sga7.sga7i.eviii.n4-13}{4.13} says that an element
$\xi$ of $\Pic(P)$ is ample if and only if its image
$\delta(\xi)$ under (\SGARefLink{sga7.sga7i.eviii.n4-14-2}{4.14.2}) is the class of an ample biextension.



}{E7I-VIII-027}\medskip
\noindent\underline{5.  Extensions and biextensions by torsion-free
truncated constant group schemes.}

In this section and the next we give some auxiliary results that will
be used in section 7 in the case $Z=\mathbb Z$.

\medskip
\noindent\underline{Proposition \SGARefTarget{sga7.sga7i.eviii.n5-1}{5.1}.}  \underline{Let $S$ be an
irreducible geometrically unibranch scheme with generic point $\eta$,
and let $Z$ be a torsion-free group, hence giving a constant group
scheme $Z_S$ over $S$.  Then every torsor $E$ under $Z_S$ is trivial,
and the functor}
\[
 E\longmapsto E(\eta)
\]
\underline{from the category of torsors under $Z_S$ to the category of
torsors under $Z=Z_S(\eta)$ is an equivalence of categories.}

\SGAvar{A familiar lemma of descent theory (SGA 3 X 5.4) shows that every
torsor $E$ under $G_S$ is representable.  In fact, that result applies
to every group scheme over $S$ which is separated, locally of finite
presentation, and locally quasi-finite over $S$.  Since $Z_S$ is
etale over $S$, fpqc descent shows that the same is true of every
torsor $E$ under $Z_S$; such a torsor is therefore locally trivial for
the \underline{etale topology} (and not merely for the fppf topology).
To prove that $E$ is trivial, it is consequently enough to prove
\[
 H^1(S_{\mathrm{\acute{e}t}},Z_S)=e.
\]
Likewise, the full-faithfulness assertion of \SGARefLink{sga7.sga7i.eviii.n5-1}{5.1} reduces to the
relation
\[
 Z\xrightarrow{\ \sim\ }H^0(S_{\mathrm{\acute{e}t}},Z_S).
\]
The latter is immediate.  To prove the former, consider the inclusion
\[
 i:\eta\longrightarrow S.
\]
Because $S$ is geometrically unibranch, we have
\[
 Z_S\xrightarrow{\ \sim\ }i_*(Z_\eta)
\]
(SGA 4 IX 2.14.1).  The noncommutative low-degree exact sequence that
replaces the Leray spectral sequence for $i:\eta\to S$ (SGA 4 XII
3.2) therefore begins
\[
 e\longrightarrow H^1(S,i_*(Z_\eta))
   \longrightarrow H^1(\eta,Z_\eta)\longrightarrow\cdots,
\]
and shows that
\[
 H^1(S,Z_S)\longrightarrow H^1(\eta,Z_\eta))
\]
is injective.  We are thus reduced to the case in which $S$ is the
spectrum of a field $k$.  If $\pi$ is the Galois group of a separable
closure of $k$, then
\[
 H^1(k,Z)=\Hom_{\mathrm{cont}}(\pi,Z)/Z=e,
\]
because $E$ is torsion-free.  This completes the proof.

\medskip
\noindent\underline{Corollary \SGARefTarget{sga7.sga7i.eviii.n5-2}{5.2}.}  \underline{With the notation of
\SGARefLink{sga7.sga7i.eviii.n5-1}{5.1}, let $P$ be a smooth scheme over $S$.  Then every torsor $E$ under
$Z_P$ is trivial, and the restriction functor}
\[
 E\longmapsto E_\eta
\]
\underline{from the category of torsors on $P_\eta$ under $Z_{P_\eta}$ to the
category of torsors on $P_\eta$ under $Z_{P_\eta}$ is an equivalence of
categories.}

Indeed, $P$ is geometrically unibranch (EGA IV 17.5.7).  Moreover,
every irreducible component of $P$ dominates $S$, because $P$ is flat
over $S$, and the set of irreducible components of $P$ is locally
finite.  Thus $P$ is the disjoint sum of irreducible geometrically
unibranch schemes $P_i$.  The assertion now follows immediately from
\SGARefLink{sga7.sga7i.eviii.n5-1}{5.1} applied to the various $P_i$ and their generic fibres $P_{i\eta}$,
noting that the generic point $p_i$ of $P_i$ is also the generic point
of $P_{i}$.

}{A familiar lemma of descent theory (SGA 3 X 5.4) shows that every
torsor $E$ under $Z_S$ is representable.  In fact, that result applies
to every group scheme over $S$ which is separated, locally of finite
presentation, and locally quasi-finite over $S$.  Since $Z_S$ is
etale over $S$, fpqc descent shows that the same is true of every
torsor $E$ under $Z_S$; such a torsor is therefore locally trivial for
the \underline{etale topology} (and not merely for the fppf topology).
To prove that $E$ is trivial, it is consequently enough to prove
\[
 H^1(S_{\mathrm{\acute{e}t}},Z_S)=e.
\]
Likewise, the full-faithfulness assertion of \SGARefLink{sga7.sga7i.eviii.n5-1}{5.1} reduces to the
relation
\[
 Z\xrightarrow{\ \sim\ }H^0(S_{\mathrm{\acute{e}t}},Z_S).
\]
The latter is immediate.  To prove the former, consider the inclusion
\[
 i:\eta\longrightarrow S.
\]
Because $S$ is geometrically unibranch, we have
\[
 Z_S\xrightarrow{\ \sim\ }i_*(Z_\eta)
\]
(SGA 4 IX 2.14.1).  The noncommutative low-degree exact sequence that
replaces the Leray spectral sequence for $i:\eta\to S$ (SGA 4 XII
3.2) therefore begins
\[
 e\longrightarrow H^1(S,i_*(Z_\eta))
   \longrightarrow H^1(\eta,Z_\eta)\longrightarrow\cdots,
\]
and shows that
\[
 H^1(S,Z_S)\longrightarrow H^1(\eta,Z_\eta)
\]
is injective.  We are thus reduced to the case in which $S$ is the
spectrum of a field $k$.  If $\pi$ is the Galois group of a separable
closure of $k$, then
\[
 H^1(k,Z)=\Hom_{\mathrm{cont}}(\pi,Z)/Z=e,
\]
because $Z$ is torsion-free.  This completes the proof.

\medskip
\noindent\underline{Corollary \SGARefTarget{sga7.sga7i.eviii.n5-2}{5.2}.}  \underline{With the notation of
\SGARefLink{sga7.sga7i.eviii.n5-1}{5.1}, let $P$ be a smooth scheme over $S$.  Then every torsor $E$ under
$Z_P$ is trivial, and the restriction functor}
\[
 E\longmapsto E_\eta
\]
\underline{from the category of torsors on $P$ under $Z_P$ to the
category of torsors on $P_\eta$ under $Z_{P_\eta}$ is an equivalence of
categories.}

Indeed, $P$ is geometrically unibranch (EGA IV 17.5.7).  Moreover,
every irreducible component of $P$ dominates $S$, because $P$ is flat
over $S$, and the set of irreducible components of $P$ is locally
finite.  Thus $P$ is the disjoint sum of irreducible geometrically
unibranch schemes $P_i$.  The assertion now follows immediately from
\SGARefLink{sga7.sga7i.eviii.n5-1}{5.1} applied to the various $P_i$ and their generic fibres $P_{i\eta}$,
noting that the generic point $p_i$ of $P_i$ is also the generic point
of $P_{i\eta}$.

}{E7I-VIII-028}\SGApage{I}{280}
\medskip
\noindent \SGARefTarget{sga7.sga7i.eviii.n5-2-1}{5.2.1}.  The preceding corollary immediately gives the
analogous statement for bitorsors on $P$ under $(Z_P,Z_P)$: restriction
to the generic fibre induces an equivalence from their category to the
category of bitorsors on $P_\eta$ under
$(Z_{P_\eta},Z_{P_\eta})$.  Indeed, it is enough to express the bitorsor
structure on $E$ in terms of its structure as a right torsor under
$Z_P$, together with an isomorphism from $Z_P$ to the commutant of
$Z_P$ for its action on $E$.  Since $E$ is necessarily trivial, that
commutant is itself isomorphic to $Z_P$.  We are thus reduced, on each
irreducible component $X$ of $P$, to the following immediate
observation: if $X$ is a nonempty connected scheme, the functor
$Z\mapsto Z_X$ from sets to $X$-schemes is fully faithful.

\medskip
\noindent\underline{Proposition \SGARefTarget{sga7.sga7i.eviii.n5-3}{5.3}.}  \underline{Let $S$ and $Z$ be
as in \SGARefLink{sga7.sga7i.eviii.n5-1}{5.1}, and let $P$ be a smooth group scheme over $S$.  Consider
extensions $E$ of group schemes of $P$ by $Z_S$.  Then restriction to
the generic fibre}
\[
 E\longmapsto E_\eta
\]
\underline{is an equivalence from the category of these extensions to
the category of extensions of $P_\eta$ by $Z_\eta$.  If $P$ and $Z$
are commutative, it also induces an equivalence between the
corresponding categories of commutative extensions.}

\SGApage{I}{\(\approx\)281}
First observe that every extension of sheaves of groups of $P$ by
$Z_S$ (and similarly of $P_\eta$ by $Z_\eta$) gives, by \SGARefLink{sga7.sga7i.eviii.n5-2}{5.2}, a trivial
torsor, hence in particular a representable one.  The categories in
\SGARefLink{sga7.sga7i.eviii.n5-3}{5.3} may therefore equally well be interpreted as categories of
extensions of sheaves of groups.  The description in VII \SGARefLink{sga7.sga7i.evii.n1-1-6}{1.1.6} and
\SGARefLink{sga7.sga7i.evii.n1-2}{1.2} expresses these in terms of bitorsors on
$P$, $P\times P$, and $P\times P\times P$ (respectively on
$P_\eta$, $P_\eta\times P_\eta$, and
$P_\eta\times P_\eta\times P_\eta$), under the pairs of groups
obtained from $(Z,Z)$ by base change.  Applying \SGARefLink{sga7.sga7i.eviii.n5-2}{5.2} in the commutative
case and \SGARefLink{sga7.sga7i.eviii.n5-2-1}{5.2.1} in the general case to the smooth $S$-schemes
$P$, $P\times P$, and $P\times P\times P$ gives the asserted
conclusion.  The same argument gives the following result.

\medskip
\noindent\underline{Corollary \SGARefTarget{sga7.sga7i.eviii.n5-4}{5.4}.}  \underline{Let $S,Z$ be as in
\SGARefLink{sga7.sga7i.eviii.n5-1}{5.1}, and let $P,Q$ be two smooth group schemes over $S$.  Suppose that
$P,Q,$ and $Z$ are commutative, and consider biextensions $E$ of
$(P,Q)$ by $Z_S$.  Then restriction to the generic fibre}
\[
 E\longmapsto E_\eta
\]
\underline{is an equivalence from the category of these biextensions
to the category of biextensions of $(P_\eta,Q_\eta)$ by $Z_\eta$.}

The preceding results reduce the determination of categories of
extensions or biextensions by a torsion-free constant group to the
case in which the base scheme is the spectrum of a field.  In that
case we obtain the following.

\SGApage{I}{282}
\medskip
\noindent\underline{Proposition \SGARefTarget{sga7.sga7i.eviii.n5-5}{5.5}.}  \underline{Let $k$ be a field,
$Z$ a torsion-free group, $P$ a $k$-group scheme locally of finite
type, $P^\circ$ its identity component, and
$\Phi=P/P^\circ$ its group of connected components (SGA 3
$\mathrm{VI}_A$ 4).  Then:}

\smallskip
\noindent (i) \underline{The natural functor from the category of group
schemes which are extensions (respectively commutative extensions) of
$\Phi$ by $Z_k$ to the category of group schemes which are extensions
(respectively commutative extensions) of $P$ by $Z_k$ is an
equivalence.  In particular, if $P$ is connected, every extension of
$P$ by $Z_k$ is trivial, and the category of such extensions is
rigid.}

\SGAvar{\smallskip
\noindent (ii) \underline{Suppose that $\Phi$ is a torsion group and
that $P$ and $Z$ are commutative.  Then the category of extensions of
$P$ by $Z_k$ is rigid: every automorphism of an object of this
category is trivial.  If, in addition, there is an integer $n>0$ such
that $n\id_\Phi=0$---for example, if $\Phi$ is finite, as happens when
$Z$ is of finite type---then the group of classes of commutative
extensions of $P$ by $G_k$ is given by the canonical isomorphism}
\begin{equation}\tag{\SGARefTarget{sga7.sga7i.eviii.n5-5-1}{5.5.1}}
 \Ext^1(P,Z_k)
 \xleftarrow{\ \sim\ }
 \Hom\bigl(\Phi,(Z\otimes_{\mathbb Z}\mathbb Q)/Z\bigr),
\end{equation}
\underline{defined by the boundary operator in the exact sequence of
$\Ext$ groups associated with the exact sequence}
\begin{equation}\tag{\SGARefTarget{sga7.sga7i.eviii.n5-5-2}{5.5.2}}
 0\longrightarrow Z\longrightarrow Z\otimes_{\mathbb Z}\mathbb Q
 \longrightarrow (Z\otimes_{\mathbb Z}\mathbb Q)/Z
 \longrightarrow 0.
\end{equation}

The canonical morphism
\[
 P\longrightarrow\Phi
\]
plainly induces an isomorphism on the sets of irreducible components,
since it is flat with irreducible fibres.  By \SGARefLink{sga7.sga7i.eviii.n5-1}{5.1}, the induced functor
from torsors under $Z_\Phi$ to torsors under $Z_P$ is therefore an
equivalence.  (Recall that a group scheme locally of finite type over
$k$ is geometrically unibranch.)  The same conclusion applies to
bitorsors and to the morphisms
\[
 P\times P\longrightarrow\Phi\times\Phi,
 \qquad
 P\times P\times P\longrightarrow\Phi\times\Phi\times\Phi.
\]
The description in VII \SGARefLink{sga7.sga7i.evii.n1-1-6}{1.1.6} of extensions of $P$ by $Z_k$ and of
$\Phi$ by $Z_k$ (and VII \SGARefLink{sga7.sga7i.evii.n1-3}{1.3} in the commutative case) now gives the
assertions in (i).

We restrict attention to commutative extensions.  The automorphism
group of an extension of $P$ by $G_k$ is
\[
 \Hom(P,G_k)\simeq\Hom(\Phi,G_k),
\]
which is plainly zero when $\Phi$ is a torsion group.  Likewise, if
$V=G\otimes_{\mathbb Z}\mathbb Q$ and $\Phi$ is annihilated by the
integer $n>0$, then
\[
 \Ext^i(\Phi,V_k)=0\qquad\text{for every }i\in\mathbb Z,
\]
because the groups in question are $\mathbb Q$-vector spaces
annihilated by $n$.  The isomorphism (\SGARefLink{sga7.sga7i.eviii.n5-5-1}{5.5.1}) follows at once.

}{\smallskip
\noindent (ii) \underline{Suppose that $\Phi$ is a torsion group and
that $P$ and $Z$ are commutative.  Then the category of extensions of
$P$ by $Z_k$ is rigid: every automorphism of an object of this
category is trivial.  If, in addition, there is an integer $n>0$ such
that $n\id_\Phi=0$---for example, if $\Phi$ is finite, as happens when
$P$ is of finite type---then the group of classes of commutative
extensions of $P$ by $Z_k$ is given by the canonical isomorphism}
\begin{equation}\tag{\SGARefTarget{sga7.sga7i.eviii.n5-5-1}{5.5.1}}
 \Ext^1(P,Z_k)
 \xleftarrow{\ \sim\ }
 \Hom\bigl(\Phi,(Z\otimes_{\mathbb Z}\mathbb Q)/Z\bigr),
\end{equation}
\underline{defined by the boundary operator in the exact sequence of
$\Ext$ groups associated with the exact sequence}
\begin{equation}\tag{\SGARefTarget{sga7.sga7i.eviii.n5-5-2}{5.5.2}}
 0\longrightarrow Z\longrightarrow Z\otimes_{\mathbb Z}\mathbb Q
 \longrightarrow (Z\otimes_{\mathbb Z}\mathbb Q)/Z
 \longrightarrow 0.
\end{equation}

The canonical morphism
\[
 P\longrightarrow\Phi
\]
plainly induces an isomorphism on the sets of irreducible components,
since it is flat with irreducible fibres.  By \SGARefLink{sga7.sga7i.eviii.n5-1}{5.1}, the induced functor
from torsors under $Z_\Phi$ to torsors under $Z_P$ is therefore an
equivalence.  (Recall that a group scheme locally of finite type over
$k$ is geometrically unibranch.)  The same conclusion applies to
bitorsors and to the morphisms
\[
 P\times P\longrightarrow\Phi\times\Phi,
 \qquad
 P\times P\times P\longrightarrow\Phi\times\Phi\times\Phi.
\]
The description in VII \SGARefLink{sga7.sga7i.evii.n1-1-6}{1.1.6} of extensions of $P$ by $Z_k$ and of
$\Phi$ by $Z_k$ (and VII \SGARefLink{sga7.sga7i.evii.n1-3}{1.2} in the commutative case) now gives the
assertions in (i).

We restrict attention to commutative extensions.  The automorphism
group of an extension of $P$ by $Z_k$ is
\[
 \Hom(P,Z_k)\simeq\Hom(\Phi,Z_k),
\]
which is plainly zero when $\Phi$ is a torsion group.  Likewise, if
$V=Z\otimes_{\mathbb Z}\mathbb Q$ and $\Phi$ is annihilated by the
integer $n>0$, then
\[
 \Ext^i(\Phi,V_k)=0\qquad\text{for every }i\in\mathbb Z,
\]
because the groups in question are $\mathbb Q$-vector spaces
annihilated by $n$.  The isomorphism (\SGARefLink{sga7.sga7i.eviii.n5-5-1}{5.5.1}) follows at once.

}{E7I-VIII-029}\SGApage{I}{\(\approx\)283}
The same method proves, essentially, the following.

\SGAvar{\medskip
\noindent\underline{Corollary \SGARefTarget{sga7.sga7i.eviii.n5-6}{5.6}.}  \underline{Let $k$ be a field,
$Z$ a torsion-free commutative group, $P$ and $Q$ two commutative
group schemes locally of finite type over $k$, and let
$\Phi=P/P^\circ$ and $\Psi=Q/Q^\circ$ be their groups of neutral
components.  Then:}

}{\medskip
\noindent\underline{Corollary \SGARefTarget{sga7.sga7i.eviii.n5-6}{5.6}.}  \underline{Let $k$ be a field,
$Z$ a torsion-free commutative group, $P$ and $Q$ two commutative
group schemes locally of finite type over $k$, and let
$\Phi=P/P^\circ$ and $\Psi=Q/Q^\circ$ be their groups of connected
components.  Then:}

}{E7I-VIII-030}\smallskip
\noindent (i) \underline{The natural functor from the category of
biextensions of $(\Phi,\Psi)$ by $Z_k$ to the category of
biextensions of $(P,Q)$ by $Z_k$ is an equivalence.  In particular,
if $P$ or $Q$ is connected, the category of biextensions of $(P,Q)$
by $Z_k$ is equivalent to the terminal category.}

\SGApage{I}{\(\approx\)284}
\SGAvar{\smallskip
\noindent (ii) \underline{Suppose that $\Phi$ or $\Psi$ is a torsion
group.  Then the category of biextensions of $(P,Q)$ by $Z_k$ is
rigid.  If, in addition, there is an integer $n>0$ such that
$n\id_\Phi=0$ or $n\id_\Psi=0$, then the group of classes of
biextensions of $(P,Q)$ by $Z_k$ is given by the canonical
isomorphism}
\begin{equation}\tag{\SGARefTarget{sga7.sga7i.eviii.n5-6-1}{5.6.1}}
 \Biext^1(P,Q;G_k)
 \xleftarrow{\ \sim\ }
 \Hom\bigl(\Phi\otimes\Psi,(Z\otimes_{\mathbb Z}\mathbb Q)/Z\bigr),
\end{equation}
\underline{defined by the boundary operator in the exact sequence of
$\Biext^i$ groups $(i=0,1)$ associated with (\SGARefLink{sga7.sga7i.eviii.n5-5-2}{5.5.2}).}




\medskip
\noindent \SGARefTarget{sga7.sga7i.eviii.n5-7}{5.7}.  Let $S$ be a scheme, $T$ a closed subscheme, and $Z$
a set.  The \underline{constant scheme over $S$, truncated along $T$,
with value $Z$}, denoted $Z_{S,T}$, is the scheme obtained by taking
the sum of copies $S_z$ of $S$, one for each $z\in Z$, and gluing all
these copies along the open subscheme $U=S\setminus T$.  Thus
\begin{equation}\tag{\SGARefTarget{sga7.sga7i.eviii.n5-7-1}{5.7.1}}
 Z_{S,T}\times_S U\simeq U,
 \qquad
 Z_{S,T}\times_S T\simeq Z_T,
\end{equation}
and $Z_{S,T}$ is etale over $S$.  For a variable $S$-scheme $S'$
there is an isomorphism, functorial in $S'$,
\begin{equation}\tag{\SGARefTarget{sga7.sga7i.eviii.n5-7-2}{5.7.2}}
 \Hom_S(S',Z_{S,T})\simeq\Hom_T(S'_T,Z_T).
\end{equation}
The map is restriction to $S'_T$, using the second isomorphism in
(\SGARefLink{sga7.sga7i.eviii.n5-7-1}{5.7.1}); we leave it to the reader to check that the map is bijective.
If
\[
 i:T\longrightarrow S
\]
denotes the inclusion, the functorial isomorphism (\SGARefLink{sga7.sga7i.eviii.n5-7-2}{5.7.2}) may also be
written as an isomorphism of sheaves
\begin{equation}\tag{\SGARefTarget{sga7.sga7i.eviii.n5-7-3}{5.7.3}}
 Z_{S,T}\simeq i_*(Z_T).
\end{equation}
This formula, or the definition directly, shows that as $Z$ varies,
$Z_{S,T}$ depends functorially on $Z$, and that this functor commutes
with finite inverse limits.  In particular, it takes groups to group
schemes and commutative groups to commutative group schemes, as is also
clear from the form (\SGARefLink{sga7.sga7i.eviii.n5-7-2}{2.7.2}) of the functor represented by $Z_{S,T}$.

}{\smallskip
\noindent (ii) \underline{Suppose that $\Phi$ or $\Psi$ is a torsion
group.  Then the category of biextensions of $(P,Q)$ by $Z_k$ is
rigid.  If, in addition, there is an integer $n>0$ such that
$n\id_\Phi=0$ or $n\id_\Psi=0$, then the group of classes of
biextensions of $(P,Q)$ by $Z_k$ is given by the canonical
isomorphism}
\begin{equation}\tag{\SGARefTarget{sga7.sga7i.eviii.n5-6-1}{5.6.1}}
 \Biext^1(P,Q;Z_k)
 \xleftarrow{\ \sim\ }
 \Hom\bigl(\Phi\otimes\Psi,(Z\otimes_{\mathbb Z}\mathbb Q)/Z\bigr),
\end{equation}
\underline{defined by the boundary operator in the exact sequence of
$\Biext^i$ groups $(i=0,1)$ associated with (\SGARefLink{sga7.sga7i.eviii.n5-5-2}{5.5.2}).}




\medskip
\noindent \SGARefTarget{sga7.sga7i.eviii.n5-7}{5.7}.  Let $S$ be a scheme, $T$ a closed subscheme, and $Z$
a set.  The \underline{constant scheme over $S$, truncated along $T$,
with value $Z$}, denoted $Z_{S,T}$, is the scheme obtained by taking
the sum of copies $S_z$ of $S$, one for each $z\in Z$, and gluing all
these copies along the open subscheme $U=S\setminus T$.  Thus
\begin{equation}\tag{\SGARefTarget{sga7.sga7i.eviii.n5-7-1}{5.7.1}}
 Z_{S,T}\times_S U\simeq U,
 \qquad
 Z_{S,T}\times_S T\simeq Z_T,
\end{equation}
and $Z_{S,T}$ is etale over $S$.  For a variable $S$-scheme $S'$
there is an isomorphism, functorial in $S'$,
\begin{equation}\tag{\SGARefTarget{sga7.sga7i.eviii.n5-7-2}{5.7.2}}
 \Hom_S(S',Z_{S,T})\simeq\Hom_T(S'_T,Z_T).
\end{equation}
The map is restriction to $S'_T$, using the second isomorphism in
(\SGARefLink{sga7.sga7i.eviii.n5-7-1}{5.7.1}); we leave it to the reader to check that the map is bijective.
If
\[
 i:T\longrightarrow S
\]
denotes the inclusion, the functorial isomorphism (\SGARefLink{sga7.sga7i.eviii.n5-7-2}{5.7.2}) may also be
written as an isomorphism of sheaves
\begin{equation}\tag{\SGARefTarget{sga7.sga7i.eviii.n5-7-3}{5.7.3}}
 Z_{S,T}\simeq i_*(Z_T).
\end{equation}
This formula, or the definition directly, shows that as $Z$ varies,
$Z_{S,T}$ depends functorially on $Z$, and that this functor commutes
with finite inverse limits.  In particular, it takes groups to group
schemes and commutative groups to commutative group schemes, as is also
clear from the form (\SGARefLink{sga7.sga7i.eviii.n5-7-2}{5.7.2}) of the functor represented by $Z_{S,T}$.

}{E7I-VIII-031}\SGApage{I}{\(\approx\)285}
It is also immediate from the definition, or from (\SGARefLink{sga7.sga7i.eviii.n5-7-3}{5.7.3}), that the
formation of $Z_{S,T}$ commutes with every base change
$f:S'\to S$.  This gives an isomorphism, functorial in $Z$,
\begin{equation}\tag{\SGARefTarget{sga7.sga7i.eviii.n5-7-4}{5.7.4}}
 Z_{S,T}\times_S S'\simeq Z_{S',T'},
 \qquad T'=T\times_S S'=T_{S'}.
\end{equation}

\SGAvar{\medskip
\noindent \SGARefTarget{sga7.sga7i.eviii.n5-8}{5.8}.  We shall study torsors $E$ under a group scheme of the
form $Z_{S,T}$, where $Z$ is a group, assumed commutative here for
simplicity.  There is a natural functor
\begin{equation}\tag{\SGARefTarget{sga7.sga7i.eviii.n5-8-0}{5.8.0}}
 E\longmapsto E_T
\end{equation}
from the category of torsors under $Z_{S,T}$ to the category of
torsors under $Z_T$.  We claim that this functor is an
\underline{equivalence of categories}.  Its full faithfulness means
that for two torsors $E,F$ under $G_{S,T}$, the map
\[
 \Hom(E,F)\longrightarrow\Hom(E_T,F_T)
\]
is bijective.  By fpqc descent this reduces immediately to the case in
which $E$ and $F$ are trivial, hence to $E=F=G_{S,T}$, where the
assertion is precisely the bijectivity of (\SGARefLink{sga7.sga7i.eviii.n5-7-2}{5.7.2}) for $S'=S$.

}{\medskip
\noindent \SGARefTarget{sga7.sga7i.eviii.n5-8}{5.8}.  We shall study torsors $E$ under a group scheme of the
form $Z_{S,T}$, where $Z$ is a group, assumed commutative here for
simplicity.  There is a natural functor
\begin{equation}\tag{\SGARefTarget{sga7.sga7i.eviii.n5-8-0}{5.8.0}}
 E\longmapsto E_T
\end{equation}
from the category of torsors under $Z_{S,T}$ to the category of
torsors under $Z_T$.  We claim that this functor is an
\underline{equivalence of categories}.  Its full faithfulness means
that for two torsors $E,F$ under $Z_{S,T}$, the map
\[
 \Hom(E,F)\longrightarrow\Hom(E_T,F_T)
\]
is bijective.  By fpqc descent this reduces immediately to the case in
which $E$ and $F$ are trivial, hence to $E=F=Z_{S,T}$, where the
assertion is precisely the bijectivity of (\SGARefLink{sga7.sga7i.eviii.n5-7-2}{5.7.2}) for $S'=S$.

}{E7I-VIII-032}\SGApage{I}{\(\approx\)286}
It remains to prove that (\SGARefLink{sga7.sga7i.eviii.n5-8-0}{5.8.0}) is essentially surjective.  Let $E_T$
be a torsor under $Z_T$.  We already observed in the proof of \SGARefLink{sga7.sga7i.eviii.n5-1}{5.1} that
$E_T$ is representable and locally trivial for the etale topology.  It
therefore suffices to prove that
\[
 H^1(S_{\mathrm{\acute{e}t}},Z_{S,T})
 \longrightarrow H^1(T_{\mathrm{\acute{e}t}},Z_T)
\]
is bijective.  This follows from (\SGARefLink{sga7.sga7i.eviii.n5-7-3}{5.7.3}) and the relation
\[
 R^1i_{\mathrm{\acute{e}t},*}(Z_T)=0,
\]
which is the special case SGA 4 VIII 5.6.

\medskip
\noindent \SGARefTarget{sga7.sga7i.eviii.n5-8-1}{5.8.1}.  Let now $P$ be a commutative group scheme over $S$.
We wish to determine the category of extensions, as sheaves of
commutative groups, of $P$ by $Z_{S,T}$.  For this we use the
description in VII \SGARefLink{sga7.sga7i.evii.n1-1-6}{1.1.6} and \SGARefLink{sga7.sga7i.evii.n1-3}{1.3} in terms of torsors on
$P$, $P\times P$, and $P\times P\times P$, under the inverse images
of $Z_{S,T}$.  By (\SGARefLink{sga7.sga7i.eviii.n5-7-4}{5.7.4}), these inverse images are again of the same
form as $Z_{S,T}$.  Applying the preceding result to those torsors
shows that they correspond to torsors on
$P_T$, $P_T\times_TP_T$, and
$P_T\times_TP_T\times_TP_T$ under the inverse images of $Z_T$.
We therefore obtain the following.

\SGAvar{\medskip
\noindent\underline{Proposition \SGARefTarget{sga7.sga7i.eviii.n5-9}{5.9}.}  \underline{Let $S$ be a scheme,
$T$ a closed subscheme, $G$ a commutative group, $G_{S,T}$ the
corresponding truncated constant group scheme over $S$ defined in
\SGARefLink{sga7.sga7i.eviii.n5-7}{5.7}, and $P$ a commutative group scheme over $S$.  Then restriction to
$T$ induces an equivalence between the category of sheaves of
commutative groups which are extensions of $P$ by $Z_{S,T}$ and the
category of sheaves of commutative groups which are extensions of
$P_T$ by $Z_T$.}

}{\medskip
\noindent\underline{Proposition \SGARefTarget{sga7.sga7i.eviii.n5-9}{5.9}.}  \underline{Let $S$ be a scheme,
$T$ a closed subscheme, $Z$ a commutative group, $Z_{S,T}$ the
corresponding truncated constant group scheme over $S$ defined in
\SGARefLink{sga7.sga7i.eviii.n5-7}{5.7}, and $P$ a commutative group scheme over $S$.  Then restriction to
$T$ induces an equivalence between the category of sheaves of
commutative groups which are extensions of $P$ by $Z_{S,T}$ and the
category of sheaves of commutative groups which are extensions of
$P_T$ by $Z_T$.}

}{E7I-VIII-033}\SGApage{I}{\(\approx\)287}
An essentially identical argument gives the following.

\medskip
\noindent\underline{Corollary \SGARefTarget{sga7.sga7i.eviii.n5-10}{5.10}.}  \underline{With the notation of
\SGARefLink{sga7.sga7i.eviii.n5-9}{5.9}, let $Q$ be another commutative group scheme over $S$.  Then
restriction to $T$ induces an equivalence between the category of
biextensions of $(P,Q)$ by $Z_{S,T}$ and the category of biextensions
of $(P_T,Q_T)$ by $Z_T$.}




\medskip
\noindent\underline{6.  Extensions and biextensions by the Neron model of}
$\underline G_m$.

\medskip
\noindent \SGARefTarget{sga7.sga7i.eviii.n6-1}{6.1}.  Let $S$ be a scheme, $T$ a closed subscheme, and $Z$ a set,
giving a truncated constant scheme $Z_{S,T}$ over $S$ (\SGARefLink{sga7.sga7i.eviii.n5-7}{5.7}).  There is a
canonical homomorphism, functorial in $Z$, of etale schemes over $S$,
\begin{equation}\tag{\SGARefTarget{sga7.sga7i.eviii.n6-1-1}{6.1.1}}
 Z_S\longrightarrow Z_{S,T},
\end{equation}
which is a \underline{surjective} morphism, as follows at once from (\SGARefLink{sga7.sga7i.eviii.n5-7-1}{5.7.1}).
Its functoriality in $Z$ shows that when $Z$ is a group, (\SGARefLink{sga7.sga7i.eviii.n6-1-1}{6.1.1}) is a
homomorphism of etale group schemes over $S$.  Its kernel is therefore an
etale group scheme over $S$, denoted $Z_{S,U}$, and there is a canonical
exact sequence, functorial in $Z$,
\begin{equation}\tag{\SGARefTarget{sga7.sga7i.eviii.n6-1-2}{6.1.2}}
 0\longrightarrow Z_{S,U}\longrightarrow Z_S\longrightarrow Z_{S,T}
 \longrightarrow 0.
\end{equation}
It follows immediately from (\SGARefLink{sga7.sga7i.eviii.n5-7-1}{5.7.1}) that $Z_{S,U}$ is the sum of a family of
$S$-schemes $S_n$, $n\in Z$, where $S_n=S$ for $n=0$ and $S_n=U$ for
$n\ne0$.  This description of $Z_{S,U}$ readily shows that for every group
$F$ over $S$, restriction to $U$ induces a bijection
\begin{equation}\tag{\SGARefTarget{sga7.sga7i.eviii.n6-1-3}{6.1.3}}
 \mathrm{Hom}_{\mathrm{gr}}(Z_{S,U},F)
 \xrightarrow{\ \sim\ }
 \mathrm{Hom}_{\mathrm{gr}}(Z_U,F|U).
\end{equation}

\SGApage{I}{288}
\medskip
\noindent \SGARefTarget{sga7.sga7i.eviii.n6-1-4}{6.1.4}.  Henceforth we suppose that $Z$ is commutative, and we wish
to examine extensions of commutative groups of $Z_{S,T}$ by a given
commutative group $G$ over $S$.  By the exact sequence (\SGARefLink{sga7.sga7i.eviii.n6-1-2}{6.1.2}), the category
of these extensions is equivalent to the category of extensions $F$ of
$Z_S$ by $G$, furnished with a trivialization of the inverse image of the
extension along $Z_{S,U}\to Z_S$.  By (\SGARefLink{sga7.sga7i.eviii.n6-1-3}{6.1.3}), such a trivialization amounts
to a trivialization of the extension $F|U$ of $Z_U$ by $G|U$.

\SGAvar{For what follows we take $Z$ to be the group $\mathbb Z$ of integers.  Thus
$\mathbb Z_S$ is the free group on one generator, and VII \SGARefLink{sga7.sga7i.evii.n1-4}{1.4} gives an
equivalence between the category of extensions of $\mathbb Z_Z$ by $G$ and
the category of torsors under $G$.  Applying the same result to
$\mathbb Z_U$ and $G|U$ gives the following.

\medskip
\noindent\underline{Lemma \SGARefTarget{sga7.sga7i.eviii.n6-2}{6.2}.}  \underline{Let $S$ be a scheme, $T$ a
closed subscheme, $U=S\setminus T$, and $G$ a commutative group over $S$.
The category of commutative groups over $S$ which are extensions of
$\mathbb Z_{S,T}$ (defined in (\SGARefLink{sga7.sga7i.eviii.n5-7}{5.7})) by $G$ is equivalent to the category of
pairs $(F,s)$ consisting of a torsor $F$ under $G$ and a section $s$ of
$F|U$.  If $\alpha:S\to\mathbb Z_{S,T}$ is the section of
$\mathbb Z_{S,T}$ which is the image of the section $1$ of $\mathbb Z_S$
over ${}_{\overline G}S$, then the torsor corresponding to an extension $\overline G$ of
$\mathbb Z_{S,T}$ by $G$ is $F=\alpha^*(\overline G)$ under $G$, furnished
with the section of $F|U$ given by restricting the unit section of
$\overline G$ to $U$.}




}{For what follows we take $Z$ to be the group $\mathbb Z$ of integers.  Thus
$\mathbb Z_S$ is the free group on one generator, and VII \SGARefLink{sga7.sga7i.evii.n1-4}{1.4} gives an
equivalence between the category of extensions of $\mathbb Z_S$ by $G$ and
the category of torsors under $G$.  Applying the same result to
$\mathbb Z_U$ and $G|U$ gives the following.

\medskip
\noindent\underline{Lemma \SGARefTarget{sga7.sga7i.eviii.n6-2}{6.2}.}  \underline{Let $S$ be a scheme, $T$ a
closed subscheme, $U=S\setminus T$, and $G$ a commutative group over $S$.
The category of commutative groups over $S$ which are extensions of
$\mathbb Z_{S,T}$ (defined in (\SGARefLink{sga7.sga7i.eviii.n5-7}{5.7})) by $G$ is equivalent to the category of
pairs $(F,s)$ consisting of a torsor $F$ under $G$ and a section $s$ of
$F|U$.  If $\alpha:S\to\mathbb Z_{S,T}$ is the section of
$\mathbb Z_{S,T}$ which is the image of the section $1$ of $\mathbb Z_S$
over $S$, then the torsor corresponding to an extension $\overline G$ of
$\mathbb Z_{S,T}$ by $G$ is $F=\alpha^*(\overline G)$ under $G$, furnished
with the section of $F|U$ given by restricting the unit section of
$\overline G$ to $U$.}




}{E7I-VIII-034}\SGApage{I}{289}
\medskip
\noindent \SGARefTarget{sga7.sga7i.eviii.n6-2-1}{6.2.1}.  If $\overline G$ corresponds to the pair $(F,s)$, then the
inverse image $f_n^*(\overline G)$ in $\overline G$ of the section $f_n$ of
$\mathbb Z_{S,T}$ defined by the constant section of $\mathbb Z_S$ with
value $n\in\mathbb Z$ is canonically identified, as a torsor under $G$, with
the $n$-th power $F^{(n)}$ of the $G$-torsor $F$.  The canonical isomorphism
of $G_U$-torsors
\[
 f_n^*(\overline G)|U\simeq G_U,
\]
defined by the first relation in (\SGARefLink{sga7.sga7i.eviii.n5-7-1}{5.7.1}), is determined by the section
$s^{(n)}$ of $F_U^{(n)}$, the $n$-th power of the given section $s$ of
$F_U$.  This makes the following description of the sections of
$\overline G$ over $S$ explicit.  A section $\overline g$ of $\overline G$
over $S$ is determined by:

\smallskip
\noindent a) a section $\underline n$ of $\mathbb Z_{S,T}$ over $S$, which
may be interpreted as a section of $\mathbb Z_T$, that is, as a locally
constant function on $T$ with values in $\mathbb Z$, or equivalently as a
partition
\[
 T=\coprod_{i\in\mathbb Z}T_i
\]
of $T$ into a disjoint sum of open subsets indexed by $i\in\mathbb Z$;

\SGAvar{\smallskip
\noindent b) a family of sections
\[
 \begin{gathered}
  \overline g_i\in\Gamma(F^{(i)}/S_i'),\qquad i\in\mathbb Z,\\
  S_i'=S\setminus T_i',\qquad
  T_i'=\displaystyle\bigcup_{\substack{j\in\mathbb Z\\j\ne i}}T_j,
 \end{gathered}
\]
subject to the following compatibility condition.  For $i,j\in\mathbb Z$,
$ g_i$ and $ g_j$ agree on $S_i'\cap S_j'=U$ after
$F^{(i)}|U$ and $F^{(j)}|U$ have been identified with $G_U$ by means of the
sections $s^{(i)}$ and $s^{(j)}$, respectively.  In other words, there must
exist a necessarily unique section $h_{ij}\in\Gamma(G/U)$ such that
\[
 \overline g_i=h_{ij}s^{(i)},\qquad
 \overline g_j=h_{ij}s^{(j)}.
\]
Since the description in \SGARefLink{sga7.sga7i.eviii.n6-2}{3.2} of an extension $\overline G$ of
$\mathbb Z_{S,T}$ by $G$ is plainly compatible with every base change
$S'\to S$, the preceding gives a complete description of the sheaf
$\overline G$ as a functor of the argument
$S'\in\operatorname{Ob}(\mathrm{Sch}_{/S})$.

}{\smallskip
\noindent b) a family of sections
\[
 \begin{gathered}
  \overline g_i\in\Gamma(F^{(i)}/S_i'),\qquad i\in\mathbb Z,\\
  S_i'=S\setminus T_i',\qquad
  T_i'=\displaystyle\bigcup_{\substack{j\in\mathbb Z\\j\ne i}}T_j,
 \end{gathered}
\]
subject to the following compatibility condition.  For $i,j\in\mathbb Z$,
$\overline g_i$ and $\overline g_j$ agree on $S_i'\cap S_j'=U$ after
$F^{(i)}|U$ and $F^{(j)}|U$ have been identified with $G_U$ by means of the
sections $s^{(i)}$ and $s^{(j)}$, respectively.  In other words, there must
exist a necessarily unique section $h_{ij}\in\Gamma(G/U)$ such that
\[
 \overline g_i=h_{ij}s^{(i)},\qquad
 \overline g_j=h_{ij}s^{(j)}.
\]
Since the description in \SGARefLink{sga7.sga7i.eviii.n6-2}{6.2} of an extension $\overline G$ of
$\mathbb Z_{S,T}$ by $G$ is plainly compatible with every base change
$S'\to S$, the preceding gives a complete description of the sheaf
$\overline G$ as a functor of the argument
$S'\in\operatorname{Ob}(\mathrm{Sch}_{/S})$.

}{E7I-VIII-035}\SGApage{I}{\(\approx\)290}
\medskip
\noindent \SGARefTarget{sga7.sga7i.eviii.n6-2-2}{6.2.2}.  If $T$ is a scheme which is a sum of nonempty connected
schemes $T_i$, $i\in I$, the description in \SGARefLink{sga7.sga7i.eviii.n6-2-1}{6.2.1} of the sections of
$\overline G$ over $S$ may be simplified.  Such a section is determined by
a family of integers $(n_i)_{i\in I}$ indexed by $I$, and a family of
sections
\[
 \overline g_i\in\Gamma(F^{(n_i)}/S_i'),\qquad i\in I,
\]
where
\[
 S_i'=S\setminus T_i',
 \qquad
 T_i'=\coprod_{\substack{j\in I\\j\ne i}}T_j.
\]
The $\overline g_i$ satisfy the following compatibility condition: for every
pair $i,j\in I$, the restrictions of $\overline g_i$ and $\overline g_j$ to
$S_i'\cap S_j'=U$ agree under the identifications of $F^{(n_i)}|U$ and
$F^{(n_j)}|U$ deduced from the respective sections $s^{(n_i)}$ and
$s^{(n_j)}$ of these torsors.

\medskip
\noindent \SGARefTarget{sga7.sga7i.eviii.n6-2-3}{6.2.3}.  Finally, if $T$ is assumed nonempty and connected, giving a
section $\overline g$ of $\overline G$ over $S$ amounts to giving an integer
$n$ and a section $\overline g$ of $F^{(n)}$ over $S$.  There is no longer
any compatibility condition to impose, so the section $s$ of $F|U$ does not
occur in the description of the group of sections of $\overline G$ over
$S$.  (Of course, $s$ occurs again as soon as one seeks to describe the
points of $\overline G$ with values in a general $S$-scheme $S'$.)

\medskip
\noindent \SGARefTarget{sga7.sga7i.eviii.n6-2-4}{6.2.4}.  Since $\overline G$ is the union of the ``open
subschemes'' $f_n^*(\overline G)$, $n\in\mathbb Z$, described in \SGARefLink{sga7.sga7i.eviii.n6-2-1}{6.2.1}, it
follows that $\overline G$ is representable if and only if the torsors
$F^{(n)}$ are representable.  For example, it suffices that $G$ be
representable and affine over $S$, by the descent theory of SGA 1 VIII 2.1.



\SGApage{I}{291}
\SGAvar{\medskip
\noindent \SGARefTarget{sga7.sga7i.eviii.n6-3}{6.3}.  In what follows we suppose that
$G=\underline G_{mS}$.  Every extension $\underline G$ of
$\mathbb Z_{S,T}$ by $G=\underline G_{mS}$ is then representable (\SGARefLink{sga7.sga7i.eviii.n6-2-4}{6.2.4}),
and is of course smooth over $S$, being an extension of the etale group
scheme $\mathbb Z_{S,T}$ by the smooth group scheme $\underline G_{mS}$.
Moreover, giving a torsor $F$ under $\underline G_m$ is equivalent to giving
an invertible Module $\mathbb L$ over $S$ (by associating to the latter the
torsor of invertible sections of $\mathbb L$).  Thus $\overline G$ is
described by a pair $(\mathbb L,s)$, where $\mathbb L$ is an invertible
Module over $S$ and $s$ is a trivialization of $\mathbb L$ over
$U=S\setminus T$, that is, an invertible section of $\mathbb L|U$.

When $U$ is schematically dense in $S$, and $S$ is locally noetherian, or is
reduced with a locally finite set of irreducible components, giving such a
pair is equivalent to giving a \underline{divisor $D$ on $S$ such that}
\begin{equation}\tag{\SGARefTarget{sga7.sga7i.eviii.n6-3-1}{6.3.1}}
 \operatorname{supp}D\subset T,
\end{equation}
by associating to such a divisor the pair
\begin{equation}\tag{\SGARefTarget{sga7.sga7i.eviii.n6-3-2}{6.3.2}}
 \bigl(\underline O_S(-D),s_{D}|U\bigr)
\end{equation}
formed by the invertible Module $\underline O_S(-D)$ associated with $-D$
(EGA IV 21.2.8.1) and the restriction to $U$ of its canonical rational
section $s_{-D}$ (EGA IV 21.2.11 and 20.2.11(ii)).  In any event, without
any hypothesis on $S$, a divisor $D$ on $S$ satisfying (\SGARefLink{sga7.sga7i.eviii.n6-3-1}{6.3.1}) defines by
(\SGARefLink{sga7.sga7i.eviii.n6-3-2}{6.3.2}) a pair $(\mathbb L,s)$ and hence an extension $ G$ of
$\mathbb Z_{S,T}$ by $\underline G_m$.  We call it the
\underline{Neron--Raynaud model of $\underline G_m$} associated with the
divisor $D$ and the closed subscheme $T$ of $S$ satisfying (\SGARefLink{sga7.sga7i.eviii.n6-3-1}{6.3.1}).  If
$T=\operatorname{supp}D$, we omit $T$ from the terminology.

}{\medskip
\noindent \SGARefTarget{sga7.sga7i.eviii.n6-3}{6.3}.  In what follows we suppose that
$G=\underline G_{mS}$.  Every extension $\overline G$ of
$\mathbb Z_{S,T}$ by $G=\underline G_{mS}$ is then representable (\SGARefLink{sga7.sga7i.eviii.n6-2-4}{6.2.4}),
and is of course smooth over $S$, being an extension of the etale group
scheme $\mathbb Z_{S,T}$ by the smooth group scheme $\underline G_{mS}$.
Moreover, giving a torsor $F$ under $\underline G_m$ is equivalent to giving
an invertible Module $\mathbb L$ over $S$ (by associating to the latter the
torsor of invertible sections of $\mathbb L$).  Thus $\overline G$ is
described by a pair $(\mathbb L,s)$, where $\mathbb L$ is an invertible
Module over $S$ and $s$ is a trivialization of $\mathbb L$ over
$U=S\setminus T$, that is, an invertible section of $\mathbb L|U$.

When $U$ is schematically dense in $S$, and $S$ is locally noetherian, or is
reduced with a locally finite set of irreducible components, giving such a
pair is equivalent to giving a \underline{divisor $D$ on $S$ such that}
\begin{equation}\tag{\SGARefTarget{sga7.sga7i.eviii.n6-3-1}{6.3.1}}
 \operatorname{supp}D\subset T,
\end{equation}
by associating to such a divisor the pair
\begin{equation}\tag{\SGARefTarget{sga7.sga7i.eviii.n6-3-2}{6.3.2}}
 \bigl(\underline O_S(-D),s_{-D}|U\bigr)
\end{equation}
formed by the invertible Module $\underline O_S(-D)$ associated with $-D$
(EGA IV 21.2.8.1) and the restriction to $U$ of its canonical rational
section $s_{-D}$ (EGA IV 21.2.11 and 20.2.11(ii)).  In any event, without
any hypothesis on $S$, a divisor $D$ on $S$ satisfying (\SGARefLink{sga7.sga7i.eviii.n6-3-1}{6.3.1}) defines by
(\SGARefLink{sga7.sga7i.eviii.n6-3-2}{6.3.2}) a pair $(\mathbb L,s)$ and hence an extension $\overline G$ of
$\mathbb Z_{S,T}$ by $\underline G_m$.  We call it the
\underline{Neron--Raynaud model of $\underline G_m$} associated with the
divisor $D$ and the closed subscheme $T$ of $S$ satisfying (\SGARefLink{sga7.sga7i.eviii.n6-3-1}{6.3.1}).  If
$T=\operatorname{supp}D$, we omit $T$ from the terminology.

}{E7I-VIII-036}\SGApage{I}{\(\approx\)292}
It is immediate that the formation of this extension $\overline G$ commutes
with every base change $f:S'\to S$ for which the inverse-image divisor
$f^*(D)$ is defined (EGA IV 21.4.2), and in particular with every
\underline{flat} base change.

\medskip
\noindent \SGARefTarget{sga7.sga7i.eviii.n6-3-3}{6.3.3}.  The general description in \SGARefLink{sga7.sga7i.eviii.n6-2-1}{6.2.1} of the sections of
$\overline G$ immediately shows that they may be identified with the
sections $f$ of the sheaf $\mathcal M_S^*$ of regular meromorphic sections
on $S$ (EGA IV 20.1.8) whose divisor $\operatorname{div}(f)$ is locally, for
the Zariski topology, an integral multiple of $D$.  If $V$ is an open subset
of $S$ on which the restriction of $\operatorname{div}(f)$ equals that of
$nD$, where $n\in\mathbb Z$, then the section of $\mathbb Z_{S,T}$ defined
by $f$ is simply the image in (\SGARefLink{sga7.sga7i.eviii.n6-1-2}{6.1.2}) of the constant section of
$\mathbb Z_S$ with value $n$.  The preceding description of the group of
sections of $\overline G$ over $S$ remains valid over every open subset of
$S$, and more generally for the group of points of $\overline G$ with
values in any flat $S$-scheme $S'$.  Since $\overline G$ itself is flat
(indeed smooth) over $S$, this description again gives a complete
characterization of the Neron--Raynaud model of $\underline G_m$ defined by
$D$ and $T\supset\operatorname{supp}D$.

\medskip
\noindent \SGARefTarget{sga7.sga7i.eviii.n6-4}{6.4}.  Suppose now that $S$ is \underline{normal}, that its set of
irreducible components is locally finite (so that $S$ is a sum of normal
integral schemes), and that $D$ is a positive divisor with all
multiplicities equal to $1$ (EGA IV 21.6), that $T=\operatorname{supp}D$,
and that $T$ is locally irreducible, that is, its irreducible components are
disjoint.  Then for every rational function $f$ on $S$ which is regular on
$U=S\setminus T$, the divisor $\operatorname{div}(f)$ is locally a multiple
of $D$.  Thus in this case there is a canonical isomorphism of groups
\begin{equation}\tag{\SGARefTarget{sga7.sga7i.eviii.n6-4-1}{6.4.1}}
 \Gamma(S,\overline G)\xrightarrow{\ \sim\ }
 \Gamma(U,\underline G_{mU}).
\end{equation}
Suppose further that $T$ is geometrically unibranch and locally integral, a
hypothesis stable under smooth base change (EGA IV 17.5.7).  Since all the
hypotheses are then stable under base change, (\SGARefLink{sga7.sga7i.eviii.n6-4-1}{6.4.1}) implies that
\underline{for every $S'$ smooth over $S$ there is a canonical group
isomorphism, functorial in $S'$,}
\begin{equation}\tag{\SGARefTarget{sga7.sga7i.eviii.n6-4-2}{6.4.2}}
 \mathrm{Hom}_S(S',\overline G)\xrightarrow{\ \sim\ }
 \mathrm{Hom}_U(S'_U,\underline G_{mU})
 \simeq\Gamma(S'_U,\underline O_{S'_U}^*).
\end{equation}
Since $\overline G$ is itself smooth over $S$, this isomorphism characterizes
the group scheme $\overline G$ over $S$ up to unique isomorphism.  It also
shows that $\overline G$ plays the role of a ``Neron model'' for the group
scheme $\underline G_{mU}$ over the open subset $U$ of $S$.  Together with
the fact that it was first introduced and used by M.~Raynaud, this justifies
the terminology adopted in \SGARefLink{sga7.sga7i.eviii.n6-3}{6.3} for $\overline G$.




\SGApage{I}{293}
\medskip
\noindent\underline{Theorem \SGARefTarget{sga7.sga7i.eviii.n6-5}{6.5}.}  \underline{Let $S$ be a normal locally
integral scheme and let $D$ be a positive divisor on $S$, without multiple
components, whose support is geometrically unibranch and locally integral.
Let $\overline G$ be the Neron--Raynaud model of $\underline G_m$ over $S$
relative to $D$ (\SGARefLink{sga7.sga7i.eviii.n6-3}{6.3}), and consider the category of torsors $E$ over $S$
under the group scheme $\overline G$.  The restriction functor}
\[
 E\longmapsto E_U=E|U
\]
\underline{from this category to the category of torsors over $U$ under the
group scheme $\underline G_{mU}$ is fully faithful.  If $S$ is regular (more
generally, if its local rings are factorial), this functor is an equivalence
of categories.}

Let us prove that the functor in question is fully faithful, that is, that
for two torsors $E,F$ under $\overline G$, the restriction map
\begin{equation}\tag{*}
 \mathrm{Hom}(E,F)\longrightarrow\mathrm{Hom}(E_U,F_U)
\end{equation}
is bijective, where $\mathrm{Hom}$ denotes the set of homomorphisms of
\underline{torsors}.  First observe that this question is plainly
\underline{local on $S$ for the Zariski topology}.  Indeed, suppose that
$S$ is the union of open subsets $S_i$ such that for every open subset $S'$
of $S$ contained in some $S_i$ (and hence in particular for $S'=S_i$ or
$S_i\cap S_j$), the map
\[
 \mathrm{Hom}(E_{S'},F_{S'})
 \longrightarrow
 \mathrm{Hom}(E_{S'\cap U},F_{S'\cap U})
\]
is bijective.  Then (*) is bijective.  On the other hand, it follows at once
from \SGARefLink{sga7.sga7i.eviii.n5-8}{5.8} applied to $S,T,\mathbb Z$, and from \SGARefLink{sga7.sga7i.eviii.n5-1}{5.1} applied to $T,\mathbb Z$,
that every torsor under $\overline G$ is locally trivial for the Zariski
topology.  These two observations reduce us to proving the bijectivity of
(*) when $E=F=\overline G$.  In that case the map is precisely (\SGARefLink{sga7.sga7i.eviii.n6-4-1}{6.4.1}),
whose bijectivity has already been established.

\SGApage{I}{\(\approx\)294}
\SGAvar{We now prove that the functor $E\mapsto E_U$ is essentially surjective when
the local rings of $X$ are factorial.  Since $\underline G_{mS}$ is
identified with an open subgroup scheme of $\overline G$, it is enough to
show that every torsor under $\underline G_{mU}$ is isomorphic to the
restriction to $U$ of a torsor under $\underline G_{mS}$.  Equivalently,
every invertible Module over $U$ must be isomorphic to the restriction of an
invertible Module over $S$.  This follows from the hypothesis on $S$
(EGA IV 21.6.11, where only the locally noetherian case is treated; that is
the only case needed below).  This proves \SGARefLink{sga7.sga7i.eviii.n6-5}{6.5}.

}{We now prove that the functor $E\mapsto E_U$ is essentially surjective when
the local rings of $S$ are factorial.  Since $\underline G_{mS}$ is
identified with an open subgroup scheme of $\overline G$, it is enough to
show that every torsor under $\underline G_{mU}$ is isomorphic to the
restriction to $U$ of a torsor under $\underline G_{mS}$.  Equivalently,
every invertible Module over $U$ must be isomorphic to the restriction of an
invertible Module over $S$.  This follows from the hypothesis on $S$
(EGA IV 21.6.11, where only the locally noetherian case is treated; that is
the only case needed below).  This proves \SGARefLink{sga7.sga7i.eviii.n6-5}{6.5}.

}{E7I-VIII-037}\SGApage{I}{295}
\medskip
\noindent \SGARefTarget{sga7.sga7i.eviii.n6-5-1}{6.5.1}.  Suppose now, in addition, that $S$ is locally noetherian
and \underline{regular}.  Since $S$ is then locally factorial, giving $D$
is equivalent to giving a reduced, purely $1$-codimensional closed
subscheme, again denoted $D$, which is geometrically unibranch.  Notice also
that the hypotheses on $S$ and $D$ are stable under every smooth base change
$S'\to S$, so that the conclusion of \SGARefLink{sga7.sga7i.eviii.n6-5}{6.5} applies to $S'$ with the subscheme
$D'=D\times_S S'$.  The description in VII \SGARefLink{sga7.sga7i.evii.n1-1-6}{1.1.6} of extension groups now
immediately gives the following consequence.

\medskip
\noindent\underline{Corollary \SGARefTarget{sga7.sga7i.eviii.n6-6}{6.6}.}  \underline{Let $S$ be a locally
noetherian regular scheme, let $D$ be a reduced, geometrically unibranch,
purely} $1$-\underline{codimensional closed subscheme (thus defining a divisor on $S$),
let $\overline G$ be the Neron--Raynaud model of $\underline G_m$ relative
to $D$, and let $P$ be a smooth commutative group scheme over $S$.  Consider
the category of commutative extensions, as sheaves of groups, of $P$ by
$\overline G$.  Restriction to $U$, $E\mapsto E|U$, induces an equivalence
between this category and the category of extensions of $P_U$ by
$\underline G_{mU}$.}

\medskip
\noindent \SGARefTarget{sga7.sga7i.eviii.n6-6-1}{6.6.1}.  The same method gives the analogous statement for
extensions which are not necessarily commutative (with $P$ no longer
assumed commutative).  We leave the details of the proof to the reader; they
are slightly more involved because bitorsors must be considered in place of
torsors.

By essentially the same proof as for \SGARefLink{sga7.sga7i.eviii.n6-6}{6.6}, one likewise obtains:

\SGApage{I}{296}
\medskip
\noindent\underline{Corollary \SGARefTarget{sga7.sga7i.eviii.n6-7}{6.7}.}  \underline{Let $S,D,\overline G$ be as
in \SGARefLink{sga7.sga7i.eviii.n6-6}{6.6}, and let $P,Q$ be two smooth commutative group schemes over $S$.
Consider the category of biextensions of $(P,Q)$ by $\overline G$.  Then the
restriction functor to $U$, $E\mapsto E|U$, from this category to the
category of biextensions of $(P_U,Q_U)$ by $\underline G_{mU}$ is an
equivalence of categories.}




\medskip
\noindent\underline{Remarks \SGARefTarget{sga7.sga7i.eviii.n6-8}{6.8}.}\quad a) We shall apply \SGARefLink{sga7.sga7i.eviii.n6-6}{6.6} and \SGARefLink{sga7.sga7i.eviii.n6-7}{6.7} when
$S$ is a trait and $D$ is the divisor defined by the inclusion of the closed
point $s$ of $S$ into $S$.  In this case $U$ is the spectrum of the fraction
field $K$ of $S$, and $\overline G$ is the usual Neron model of the
multiplicative group over $K$, introduced by Raynaud.  Notice that even in
this special case, the proof of \SGARefLink{sga7.sga7i.eviii.n6-6}{6.6} and \SGARefLink{sga7.sga7i.eviii.n6-7}{6.7} uses \SGARefLink{sga7.sga7i.eviii.n6-5}{6.5} over bases more general
than traits (such as $P$ and its Cartesian powers).  Nevertheless, in this
case it is enough to know the \underline{existence} of the Neron model
$\overline G$ in the special situation under consideration: \SGARefLink{sga7.sga7i.eviii.n6-5}{6.5} for base
schemes smooth over $S$, and hence \SGARefLink{sga7.sga7i.eviii.n6-6}{6.6} and \SGARefLink{sga7.sga7i.eviii.n6-7}{6.7}, then follow formally.  For
us, the introduction of $\overline G$ will chiefly serve as a convenient
intermediate step in studying, in the next section, the problem of extending
extensions and biextensions by the multiplicative group.

\smallskip
\noindent b) A Neron--Raynaud model $\overline G$ of
$\underline G_{mU}$, satisfying the fundamental condition (\SGARefLink{sga7.sga7i.eviii.n6-4-2}{6.4.2}) for every
$S'$ smooth over $S$, can be defined under much more general conditions than
those considered above, with $U$ merely a schematically dense open subset
of a normal locally noetherian scheme $S$.  In this case $\overline G$ is an
fppf \underline{sheaf}, not necessarily representable.  Let $\underline D$
denote the sheaf, on the small etale site of $S$, of divisors on $S$ whose
support is contained in $T=S\setminus U$, and, by abuse of language, also
denote by the same letter its canonical extension to a sheaf on
$(\mathrm{Sch})_{/S}$.  Its sections over a relative scheme
$f:S'\to S$ are the sections of $f_{\mathrm{\acute{e}t}}^*(\underline D)$
over $S'$.  The desired model is an extension of groups
\begin{equation}\tag{\SGARefTarget{sga7.sga7i.eviii.n6-8-1}{6.8.1}}
 0\longrightarrow\underline G_{mS}\longrightarrow\overline G
 \longrightarrow\underline D\longrightarrow0.
\end{equation}
The construction of a canonical extension (\SGARefLink{sga7.sga7i.eviii.n6-8-1}{6.8.1}), without any normality
hypothesis on $S$, presents no difficulty when one uses the general
description VII \SGARefLink{sga7.sga7i.evii.n1-1-6}{1.1.6}; at the same time one obtains (\SGARefLink{sga7.sga7i.eviii.n6-4-2}{6.4.2}) for $S'$ etale
over $S$.  For the same relation to hold for every $S'$ smooth over $S$, it
is equivalent to verify that the formation of the sheaf $\underline D$
commutes with every smooth base change $f:S'\to S$.  For normal $S$, this
follows from EGA $\mathrm{Err}_{\mathrm{IV}}$ 53, 21.4.9.

\SGApage{I}{\(\approx\)297}
\medskip
\noindent \SGARefTarget{sga7.sga7i.eviii.n6-9}{6.9}.\quad ``\underline{Self-criticism}'' (observations of
P.~Deligne).\footnote{The author of the present Expose having admitted
defeat, the reader is invited, if necessary, to rewrite sections 5 and 6
according to the accompanying directions.}

Sections 5 and 6 are really rather uninspiring.  One periodically has the
impression that the formalism of $i_*,i^*$ is being redone ``by hand'',
without seeing what is geometric and what is ``general nonsense''.

I propose extracting the ``geometric'' statements a), b), c):

\SGAvar{\smallskip
\noindent a) (cf. \SGARefLink{sga7.sga7i.eviii.n6-8}{6.8} b)) Let $j:V\hookrightarrow S$ be a schematically
dense open subset of a normal noetherian scheme $S$.  The etale sheaf
$\underline G_{mS}$ is then identified with a subsheaf of the etale sheaf
$j_*\underline G_{mU}$.  Let $\underline D$ be the quotient etale sheaf in
the exact sequence
\[
 0\longrightarrow\underline G_{mS}\longrightarrow
 j_*\underline G_{mS}\longrightarrow\underline D\longrightarrow0.
\]
\underline{Then} the formation of $\underline D$ commutes with every smooth
base change (EGA $\mathrm{Err}_{\mathrm{IV}}$ 53, 21.4.9); \underline{moreover},
$R^1j_*\underline G_m=0$ (Theorem 90) if the strictly local rings of $S$ are
factorial.

}{\smallskip
\noindent a) (cf. \SGARefLink{sga7.sga7i.eviii.n6-8}{6.8} b)) Let $j:U\hookrightarrow S$ be a schematically
dense open subset of a normal noetherian scheme $S$.  The etale sheaf
$\underline G_{mS}$ is then identified with a subsheaf of the etale sheaf
$j_*\underline G_{mU}$.  Let $\underline D$ be the quotient etale sheaf in
the exact sequence
\[
 0\longrightarrow\underline G_{mS}\longrightarrow
 j_*\underline G_{mU}\longrightarrow\underline D\longrightarrow0.
\]
\underline{Then} the formation of $\underline D$ commutes with every smooth
base change (EGA $\mathrm{Err}_{\mathrm{IV}}$ 53, 21.4.9); \underline{moreover},
$R^1j_*\underline G_m=0$ (Theorem 90) if the strictly local rings of $S$ are
factorial.

}{E7I-VIII-038}\SGApage{I}{\(\approx\)298}
\smallskip
\noindent b) If $U=S\setminus Y$, where $Y$ is a geometrically unibranch
divisor, and if $i$ is the inclusion of $Y$ into $S$, then
\[
 \underline D=i_*\mathbb Z_Y.
\]

\SGAvar{\smallskip
\noindent c) If $Y$ is a closed subscheme of a scheme $X$, and if $E$ is a set,
then the etale sheaf $i_*\underline E_Y$ is represented by an etale scheme
over $X$, whose formation commutes with every base change.

}{\smallskip
\noindent c) If $Y$ is a closed subscheme of a scheme $X$, with closed immersion
$i:Y\hookrightarrow X$, and if $E$ is a set,
then the etale sheaf $i_*\underline E_Y$ is represented by an etale scheme
over $X$, whose formation commutes with every base change.

}{E7I-VIII-039}\medskip
\noindent\underline{Corollary.}  \underline{Let $S$ be a normal noetherian
scheme, let $Y$ be a geometrically unibranch divisor, and let
$U=S\setminus Y$:}
\begin{center}
\begin{tikzcd}[column sep=huge]
 U \arrow[r, hook, "j"'] \& S \& Y \arrow[l, "i"']
\end{tikzcd}
\qquad .
\end{center}
\underline{There exists a unique smooth group scheme $\mathcal G_m$ over
$S$ such that after every smooth base change one has
$\mathcal G_m|S_{\mathrm{\acute{e}t}}=j_*(\underline G_m)_U$.}

Indeed, the preceding requirements determine $\mathcal G_m$ as a sheaf, for
the etale topology, on the category of schemes smooth over $S$.  By a) and
b), the sequence
\[
 0\longrightarrow\underline G_m\longrightarrow\mathcal G_m
 \xrightarrow{\ q\ }i_*\mathbb Z_Y\longrightarrow0
\]
is exact.  By c), $i_*\mathbb Z_Y$ is representable.  Finally, the morphism
$q$ is representable.

Use of the preceding ``small smooth site'' should make it possible to deduce
the statements of section 6 formally from a).



\SGApage{I}{\(\approx\)299}
\medskip
\noindent\underline{7.  Canonical extensions of extensions and
biextensions by} $\underline G_m$.

\medskip
\noindent \SGARefTarget{sga7.sga7i.eviii.n7-0}{7.0}.  We now reap the fruit of our preparations in sections 5 and
6.  First let the base scheme $S$ be a \underline{trait}, that is, the
spectrum of a discrete valuation ring; let $\eta$ be its generic point,
$s$ its closed point, and $k=k(s)$.  Let $P$ be a smooth commutative group
scheme over $S$, and write $P_0=P_s$ for its special fibre.  In the present
section we consider only commutative groups; in particular, all group
extensions under consideration are commutative.

\medskip
\noindent\underline{Theorem \SGARefTarget{sga7.sga7i.eviii.n7-1}{7.1}.}  \underline{The notation is that recalled
in \SGARefLink{sga7.sga7i.eviii.n7-0}{7.0}.}

\smallskip
\noindent a) \underline{Suppose that the group of connected components of
the special fibre of $P$,}
\[
 \Phi=P_0/P_0^{\circ},
\]
\underline{is a torsion group---as is the case, in particular, if $P$ is of
finite type over $S$.  Then the functor $E\mapsto E_\eta$ from the category
of extensions of $P$ by $\underline G_{mS}$ to the category of extensions
of $P_\eta$ by $\underline G_{m\eta}$ is fully faithful.  Suppose that there
is an integer $n>0$ such that $n\Phi=0$.  Then every extension $E_\eta$ of
$P_\eta$ by $\underline G_{m\eta}$ has an associated canonical
homomorphism}
\begin{equation}\tag{\SGARefTarget{sga7.sga7i.eviii.n7-1-1}{7.1.1}}
 d(E_\eta):\Phi\longrightarrow(\mathbb Q/\mathbb Z)_k,
\end{equation}
\underline{whose vanishing is necessary and sufficient for $E_\eta$ to be
isomorphic to the generic fibre of an extension of $P$ by
$\underline G_{mS}$.  In particular, if $P_0$ is connected, the preceding
functor $E\mapsto E_\eta$ is an equivalence of categories.}

\SGAvar{\smallskip
\noindent b) \underline{Let $P$ and $Q$ be two smooth group schemes over
$S$, put}
\[
 \Phi=P_0/P_0^{\circ},\qquad \Psi=Q_0/Q_0^{\circ},
\]
\underline{and suppose that either $\Phi$ or $\Psi$ is a torsion group} (\underline{for
example, that either $P$ or $Q$ is of finite type over $S$}).  \underline{Then the
functor $E\mapsto E_\eta$ from the category of biextensions of $(P,Q)$ by
$\underline G_{mS}$ to the category of biextensions of
$(P_\eta,Q_\eta)$ by $\underline G_{m\eta}$ is fully faithful.  Suppose that
there is an integer $n>0$ such that $n\Phi=0$ or $n\Psi=0$.  Then every
biextension $E_\eta$ of $(P_\eta,Q_\eta)$ by $\underline G_{m\eta}$ has an
associated canonical pairing}
\begin{equation}\tag{\SGARefTarget{sga7.sga7i.eviii.n7-1-2}{7.1.2}}
 d(E_\eta):\Phi\otimes_{\mathbb Z}\Psi
 \longrightarrow(\mathbb Q/\mathbb Z)_k,
\end{equation}
\underline{whose vanishing is necessary and sufficient for $E$ to be
isomorphic to the generic fibre of a biextension $E$ of $(P,Q)$ by
$\underline G_{mS}$.  In particular, if either $P_0$ or $Q_0$ is connected,
the preceding functor $E\mapsto E_\eta$ is an equivalence of categories.}




\medskip
\noindent\underline{Proof.}\quad a) Let $\overline G$ be the
Neron--Raynaud model of $\underline G_m$.  The functor under consideration
is the composite of the functor ``extension of the kernel
$\underline G_{mS}\to\overline G$'' with the functor ``restriction to
$\eta$'' defined on the category of biextensions by $\overline G$.  The latter
functor is an equivalence of categories by \SGARefLink{sga7.sga7i.eviii.n6-6}{6.6}, so it remains to study the
first.  For this purpose we use the exact sequence of group schemes
\begin{equation}\tag{\SGARefTarget{sga7.sga7i.eviii.n7-1-3}{7.1.3}}
 0\longrightarrow\underline G_{mS}\longrightarrow\overline G
 \longrightarrow\mathbb Z_{S,s}\longrightarrow0,
\end{equation}
where the notation $\mathbb Z_{S,s}$ is that of \SGARefLink{sga7.sga7i.eviii.n5-7}{5.7}.  The exact sequence of
$\mathrm{Ext}^i$, $i=0,1$, associated with (\SGARefLink{sga7.sga7i.eviii.n7-1-3}{7.1.3}) shows that the functor in
question is fully faithful if and only if
\begin{equation}\tag{\SGARefTarget{sga7.sga7i.eviii.n7-1-4}{7.1.4}}
 \mathrm{Hom}(P,\mathbb Z_{S,s})=0,
\end{equation}
and that the obstruction for an extension $\overline E$ of $P$ by
$\overline G$ to come from an extension $E$ of $P$ by
$\underline G_{mS}$ lies in
\begin{equation}\tag{\SGARefTarget{sga7.sga7i.eviii.n7-1-5}{7.1.5}}
 \mathrm{Ext}^1(P,\mathbb Z_{S,s}).
\end{equation}
By (\SGARefLink{sga7.sga7i.eviii.n5-7-2}{5.7.2}), the first group in (\SGARefLink{sga7.sga7i.eviii.n7-1-4}{7.1.4}) is isomorphic to
$\mathrm{Hom}(P_0,\mathbb Z_k)$, which in turn is isomorphic to
$\mathrm{Hom}(\Phi,\mathbb Z_k)$.  Thus (\SGARefLink{sga7.sga7i.eviii.n7-1-4}{7.1.4}) is equivalent to
\begin{equation}\tag{7.1.4 bis}
 \mathrm{Hom}(\Phi,\mathbb Z_k)=0.
\end{equation}
This relation certainly holds if $\Phi$ is a torsion group.  Moreover, if
$n\Phi=0$, then \SGARefLink{sga7.sga7i.eviii.n5-9}{5.9} and \SGARefLink{sga7.sga7i.eviii.n5-5}{5.5} give canonical isomorphisms
\begin{equation}\tag{\SGARefTarget{sga7.sga7i.eviii.n7-1-6}{7.1.6}}
 \mathrm{Ext}^1(P,\mathbb Z_{S,s})
 \simeq\mathrm{Ext}^1(P_0,\mathbb Z_k)
 \simeq\mathrm{Hom}\bigl(\Phi,(\mathbb Q/\mathbb Z)_k\bigr),
\end{equation}
which completes the proof of \SGARefLink{sga7.sga7i.eviii.n7-1}{7.1}(a).

\smallskip
\noindent b) The proof is essentially the same.  By \SGARefLink{sga7.sga7i.eviii.n6-7}{6.7}, we reduce to
studying the functor ``extension of the kernel
$\underline G_{mS}\to\overline G$'' from the category of biextensions of
$(P,Q)$ by $\underline G_{mS}$ to the category of biextensions of $(P,Q)$ by
$\overline G$.  We use the exact sequence (\SGARefLink{sga7.sga7i.eviii.n7-1-3}{7.1.3}) and the associated exact
sequence of $\mathrm{Biext}^i$, $i=0,1$ (VII \SGARefLink{sga7.sga7i.evii.n3-7-5}{3.7.5}).  Full faithfulness of
the functor is equivalent to
\begin{equation}\tag{\SGARefTarget{sga7.sga7i.eviii.n7-1-7}{7.1.7}}
 \mathrm{Hom}(P\otimes Q,\mathbb Z_{S,s})=0.
\end{equation}
By (\SGARefLink{sga7.sga7i.eviii.n5-7-2}{5.7.2}), the left side is isomorphic to
$\mathrm{Hom}(P_0\otimes Q_0,\mathbb Z_k)$ and hence to
$\mathrm{Hom}(\Phi\otimes\Psi,\mathbb Z_k)$; this group is indeed zero if
either $\Phi$ or $\Psi$ is a torsion group.  On the other hand, for a given
biextension $\overline E$ of $(P,Q)$ by $\overline G$, the obstruction to
its coming from a biextension $E$ of $(P,Q)$ by $\underline G_{mS}$ lies in
\[
 \mathrm{Biext}^1(P,Q;\mathbb Z_{S,s}).
\]
By \SGARefLink{sga7.sga7i.eviii.n5-10}{5.10} this is isomorphic to
$\mathrm{Biext}^1(P_0,Q_0;\mathbb Z_k)$, which, when $n\Phi=0$ or
$n\Psi=0$, is itself isomorphic by \SGARefLink{sga7.sga7i.eviii.n5-6}{5.6} to
\[
 \mathrm{Hom}\bigl(\Phi\otimes\Psi,Q/\mathbb Z\bigr).
\]
This completes the proof of \SGARefLink{sga7.sga7i.eviii.n7-1}{7.1}.




}{\smallskip
\noindent b) \underline{Let $P$ and $Q$ be two smooth group schemes over
$S$, put}
\[
 \Phi=P_0/P_0^{\circ},\qquad \Psi=Q_0/Q_0^{\circ},
\]
\underline{and suppose that either $\Phi$ or $\Psi$ is a torsion group} (\underline{for
example, that either $P$ or $Q$ is of finite type over $S$}).  \underline{Then the
functor $E\mapsto E_\eta$ from the category of biextensions of $(P,Q)$ by
$\underline G_{mS}$ to the category of biextensions of
$(P_\eta,Q_\eta)$ by $\underline G_{m\eta}$ is fully faithful.  Suppose that
there is an integer $n>0$ such that $n\Phi=0$ or $n\Psi=0$.  Then every
biextension $E_\eta$ of $(P_\eta,Q_\eta)$ by $\underline G_{m\eta}$ has an
associated canonical pairing}
\begin{equation}\tag{\SGARefTarget{sga7.sga7i.eviii.n7-1-2}{7.1.2}}
 d(E_\eta):\Phi\otimes_{\mathbb Z}\Psi
 \longrightarrow(\mathbb Q/\mathbb Z)_k,
\end{equation}
\underline{whose vanishing is necessary and sufficient for $E_\eta$ to be
isomorphic to the generic fibre of a biextension $E$ of $(P,Q)$ by
$\underline G_{mS}$.  In particular, if either $P_0$ or $Q_0$ is connected,
the preceding functor $E\mapsto E_\eta$ is an equivalence of categories.}




\medskip
\noindent\underline{Proof.}\quad a) Let $\overline G$ be the
Neron--Raynaud model of $\underline G_m$.  The functor under consideration
is the composite of the functor ``extension of the kernel
$\underline G_{mS}\to\overline G$'' with the functor ``restriction to
$\eta$'' defined on the category of extensions by $\overline G$.  The latter
functor is an equivalence of categories by \SGARefLink{sga7.sga7i.eviii.n6-6}{6.6}, so it remains to study the
first.  For this purpose we use the exact sequence of group schemes
\begin{equation}\tag{\SGARefTarget{sga7.sga7i.eviii.n7-1-3}{7.1.3}}
 0\longrightarrow\underline G_{mS}\longrightarrow\overline G
 \longrightarrow\mathbb Z_{S,s}\longrightarrow0,
\end{equation}
where the notation $\mathbb Z_{S,s}$ is that of \SGARefLink{sga7.sga7i.eviii.n5-7}{5.7}.  The exact sequence of
$\mathrm{Ext}^i$, $i=0,1$, associated with (\SGARefLink{sga7.sga7i.eviii.n7-1-3}{7.1.3}) shows that the functor in
question is fully faithful if and only if
\begin{equation}\tag{\SGARefTarget{sga7.sga7i.eviii.n7-1-4}{7.1.4}}
 \mathrm{Hom}(P,\mathbb Z_{S,s})=0,
\end{equation}
and that the obstruction for an extension $\overline E$ of $P$ by
$\overline G$ to come from an extension $E$ of $P$ by
$\underline G_{mS}$ lies in
\begin{equation}\tag{\SGARefTarget{sga7.sga7i.eviii.n7-1-5}{7.1.5}}
 \mathrm{Ext}^1(P,\mathbb Z_{S,s}).
\end{equation}
By (\SGARefLink{sga7.sga7i.eviii.n5-7-2}{5.7.2}), the first group in (\SGARefLink{sga7.sga7i.eviii.n7-1-4}{7.1.4}) is isomorphic to
$\mathrm{Hom}(P_0,\mathbb Z_k)$, which in turn is isomorphic to
$\mathrm{Hom}(\Phi,\mathbb Z_k)$.  Thus (\SGARefLink{sga7.sga7i.eviii.n7-1-4}{7.1.4}) is equivalent to
\begin{equation}\tag{7.1.4 bis}
 \mathrm{Hom}(\Phi,\mathbb Z_k)=0.
\end{equation}
This relation certainly holds if $\Phi$ is a torsion group.  Moreover, if
$n\Phi=0$, then \SGARefLink{sga7.sga7i.eviii.n5-9}{5.9} and \SGARefLink{sga7.sga7i.eviii.n5-5}{5.5} give canonical isomorphisms
\begin{equation}\tag{\SGARefTarget{sga7.sga7i.eviii.n7-1-6}{7.1.6}}
 \mathrm{Ext}^1(P,\mathbb Z_{S,s})
 \simeq\mathrm{Ext}^1(P_0,\mathbb Z_k)
 \simeq\mathrm{Hom}\bigl(\Phi,(\mathbb Q/\mathbb Z)_k\bigr),
\end{equation}
which completes the proof of \SGARefLink{sga7.sga7i.eviii.n7-1}{7.1}(a).

\smallskip
\noindent b) The proof is essentially the same.  By \SGARefLink{sga7.sga7i.eviii.n6-7}{6.7}, we reduce to
studying the functor ``extension of the kernel
$\underline G_{mS}\to\overline G$'' from the category of biextensions of
$(P,Q)$ by $\underline G_{mS}$ to the category of biextensions of $(P,Q)$ by
$\overline G$.  We use the exact sequence (\SGARefLink{sga7.sga7i.eviii.n7-1-3}{7.1.3}) and the associated exact
sequence of $\mathrm{Biext}^i$, $i=0,1$ (VII \SGARefLink{sga7.sga7i.evii.n3-7-5}{3.7.5}).  Full faithfulness of
the functor is equivalent to
\begin{equation}\tag{\SGARefTarget{sga7.sga7i.eviii.n7-1-7}{7.1.7}}
 \mathrm{Hom}(P\otimes Q,\mathbb Z_{S,s})=0.
\end{equation}
By (\SGARefLink{sga7.sga7i.eviii.n5-7-2}{5.7.2}), the left side is isomorphic to
$\mathrm{Hom}(P_0\otimes Q_0,\mathbb Z_k)$ and hence to
$\mathrm{Hom}(\Phi\otimes\Psi,\mathbb Z_k)$; this group is indeed zero if
either $\Phi$ or $\Psi$ is a torsion group.  On the other hand, for a given
biextension $\overline E$ of $(P,Q)$ by $\overline G$, the obstruction to
its coming from a biextension $E$ of $(P,Q)$ by $\underline G_{mS}$ lies in
\[
 \mathrm{Biext}^1(P,Q;\mathbb Z_{S,s}).
\]
By \SGARefLink{sga7.sga7i.eviii.n5-10}{5.10} this is isomorphic to
$\mathrm{Biext}^1(P_0,Q_0;\mathbb Z_k)$, which, when $n\Phi=0$ or
$n\Psi=0$, is itself isomorphic by \SGARefLink{sga7.sga7i.eviii.n5-6}{5.6} to
\[
 \mathrm{Hom}\bigl(\Phi\otimes\Psi,(\mathbb Q/\mathbb Z)_k\bigr).
\]
This completes the proof of \SGARefLink{sga7.sga7i.eviii.n7-1}{7.1}.




}{E7I-VIII-040}\SGApage{I}{302}
\medskip
\noindent\underline{Remark \SGARefTarget{sga7.sga7i.eviii.n7-2}{7.2}.}  Given a smooth group scheme $P$ over a
connected regular locally noetherian scheme $S$ with generic point $\eta$,
one may ask to study the functor $E\mapsto E_\eta$ from the category of
extensions of $P$ by $\underline G_{mS}$ to the category of extensions of
$P_\eta$ by $\underline G_{m\eta}$.  This problem is readily reduced to the
one treated in \SGARefLink{sga7.sga7i.eviii.n7-1}{7.1}(a).  For simplicity, suppose that $P$ is of finite type
over $S$.  It follows easily from \SGARefLink{sga7.sga7i.eviii.n7-1}{7.1} that this functor is fully
faithful.\footnote{For this it is enough that $S$ be normal rather than
regular.}  One readily deduces that if $E_\eta$ is an extension of $P_\eta$
by $\underline G_{m\eta}$, then it is isomorphic to the generic fibre of an
extension $E$ of $P$ by $\underline G_{mS}$ if and only if, for every
$s\in S$ such that
\[
 \dim\operatorname{Spec}(\underline O_{S,s})=1,
\]
the analogous extension problem obtained by base change
$\operatorname{Spec}(\underline O_{S,s})\to S$ has a solution.  By \SGARefLink{sga7.sga7i.eviii.n7-1}{7.1}(a),
this is expressed by the vanishing of a certain homomorphism
\begin{equation}\tag{\SGARefTarget{sga7.sga7i.eviii.n7-2-1}{7.2.1}}
 d(E_\eta,s):P_s/P_s^{\circ}
 \longrightarrow(\mathbb Q/\mathbb Z)_{k(s)}.
\end{equation}

The preceding assertion plainly reduces to the following one.  Let $U$ be
an open subset of $S$ whose complement $T$ satisfies
\[
 \operatorname{codim}(T,S)\geq2.
\]
Then the functor $E\mapsto E|U$ from the category of extensions of $P$ by
$\underline G_{mS}$ to the category of extensions of $P_U$ by
$\underline G_{mU}$ is an equivalence of categories.  In terms of the
description VII \SGARefLink{sga7.sga7i.evii.n1-1-6}{1.1.6} this assertion is essentially trivial: for every
smooth $S$-scheme $X$ (such as $P$, $P\times P$,
$P\times P\times P$, and so on), $X$ is locally factorial and
$X_T=X\setminus X_U$ has codimension at least $2$ in $X$.  Consequently,
the functor $F\mapsto F|X_U$ from the category of torsors under
$\underline G_{mX}$ to the category of torsors under
$\underline G_{mX_U}$ is an equivalence of categories.

\SGApage{I}{\(\approx\)303}
There are entirely analogous results for the problem of extending
biextensions; we leave their formulation to the reader.




\medskip
\noindent \SGARefTarget{sga7.sga7i.eviii.n7-3}{7.3}.  We shall make explicit some functorial properties of the
obstructions defined in (\SGARefLink{sga7.sga7i.eviii.n7-1-1}{7.1.1}) and (\SGARefLink{sga7.sga7i.eviii.n7-1-2}{7.1.2}), with $P,Q,S$ variable, and
some relations among them.  In what follows, we shall suppose that the
conditions allowing these obstructions to be defined are satisfied.  Since
the verifications present no difficulty, we leave them to the reader.

\SGAvar{\medskip
\noindent \SGARefTarget{sga7.sga7i.eviii.n7-3-1}{7.3.1}.  Let
\[
 u:P'\longrightarrow P
\]
be a morphism of smooth group schemes over $S$, and let $E_\eta$ be an
extension of $P_\eta$ by $\underline G_{m\eta}$.  Pullback by
$u_\eta:P'_\eta\to P_\eta$ gives an extension $E'_\eta$ of $P'_\eta$ by
$\underline G_{m}$.  With this notation, the following diagram is
commutative:
\begin{equation}\tag{\SGARefTarget{sga7.sga7i.eviii.n7-3-1-1}{7.3.1.1}}
\begin{tikzcd}[row sep=large,column sep=huge]
 P'_0/{P'_0}^{\circ}
   \arrow[dd,"u_0"']
   \arrow[rd,"d(E'_\eta)"] \& \\
 \& (\mathbb Q/\mathbb Z)_k \\
 P_0/P_0^{\circ}
   \arrow[ru,"d(E_\eta)"'] \&
\end{tikzcd}
\qquad .
\end{equation}

}{\medskip
\noindent \SGARefTarget{sga7.sga7i.eviii.n7-3-1}{7.3.1}.  Let
\[
 u:P'\longrightarrow P
\]
be a morphism of smooth group schemes over $S$, and let $E_\eta$ be an
extension of $P_\eta$ by $\underline G_{m\eta}$.  Pullback by
$u_\eta:P'_\eta\to P_\eta$ gives an extension $E'_\eta$ of $P'_\eta$ by
$\underline G_{m\eta}$.  With this notation, the following diagram is
commutative:
\begin{equation}\tag{\SGARefTarget{sga7.sga7i.eviii.n7-3-1-1}{7.3.1.1}}
\begin{tikzcd}[row sep=large,column sep=huge]
 P'_0/{P'_0}^{\circ}
   \arrow[dd,"u_0"']
   \arrow[rd,"d(E'_\eta)"] \& \\
 \& (\mathbb Q/\mathbb Z)_k \\
 P_0/P_0^{\circ}
   \arrow[ru,"d(E_\eta)"'] \&
\end{tikzcd}
\qquad .
\end{equation}

}{E7I-VIII-041}\SGApage{I}{\(\approx\)304}
Similarly, if we have morphisms of smooth group schemes over $S$
\[
 u:P'\longrightarrow P,\qquad v:Q'\longrightarrow Q,
\]
and a biextension $E_\eta$ of $(P_\eta,Q_\eta)$ by
$\underline G_{m\eta}$, whose pullback by $(u_\eta,v_\eta)$ is a
biextension $E'_\eta$ of $(P'_\eta,Q'_\eta)$ by
$\underline G_{m\eta}$, then the following diagram is commutative:
\begin{equation}\tag{\SGARefTarget{sga7.sga7i.eviii.n7-3-1-2}{7.3.1.2}}
\begin{tikzcd}[row sep=large,column sep=huge]
 P'_0/{P'_0}^{\circ}\otimes Q'_0/{Q'_0}^{\circ}
   \arrow[dd,"u_0\otimes v_0"']
   \arrow[rd,"d(E'_\eta)"] \& \\
 \& (\mathbb Q/\mathbb Z)_k \\
 P_0/P_0^{\circ}\otimes Q_0/Q_0^{\circ}
   \arrow[ru,"d(E_\eta)"'] \&
\end{tikzcd}
\end{equation}
The essentially mechanical verification of these compatibilities is left
to the reader.



\medskip
\noindent \SGARefTarget{sga7.sga7i.eviii.n7-3-2}{7.3.2}.  It follows from the compatibility (\SGARefLink{sga7.sga7i.eviii.n7-3-1-1}{7.3.1.1}) that, if
$E_\eta$ is an extension of $P_\eta$ by $\underline G_{m\eta}$, there is a
largest group subscheme $V$ of $P$ such that the restriction of $E_\eta$ to
$V_\eta$ extends to an extension of $V$ by $\underline G_{mS}$.  It is the
open subgroup $V$ of $P$ containing $P_\eta$ and defined by the condition
\begin{equation}\tag{\SGARefTarget{sga7.sga7i.eviii.n7-3-2-1}{7.3.2.1}}
 V_0/V_0^{\circ}=\ker d(E_\eta).
\end{equation}
If $u:P'\to P$ is then a morphism of smooth group schemes over $S$, the
extension $E'_\eta$, obtained by pulling $E_\eta$ back along $u_\eta$,
extends to an extension of $P'$ by $\underline G_{mS}$ if and only if
\begin{equation}\tag{\SGARefTarget{sga7.sga7i.eviii.n7-3-2-2}{7.3.2.2}}
 u(P')\subset V.
\end{equation}

\SGApage{I}{\(\approx\)305}
\SGAvar{In the case of a biextension $E$ of $(P_\eta,Q_\eta)$ by
$\underline G_{m\eta}$, there is likewise a largest group subscheme $V$ of
$P$ such that the restriction of $E$ to $(V_\eta,Q_\eta)$ extends to
a biextension of $(V,Q)$ by $\underline G_{mS}$.  It is the open subgroup
of $P$ containing the generic fibre $P_\eta$ and characterized by the fact
that $V_0/V_0^{\circ}$ is the left annihilator in $\Phi$ of the pairing
(\SGARefLink{sga7.sga7i.eviii.n7-1-2}{7.1.2}), that is, the kernel of the corresponding homomorphism from $\Phi$
to $\underline{\mathrm{Hom}}(\Psi,(\mathbb Q/\mathbb Z)_k)$:
\begin{equation}\tag{\SGARefTarget{sga7.sga7i.eviii.n7-3-2-3}{7.3.2.3}}
 V_0/V_0^{\circ}
 =\ker\!\left(
   \Phi\xrightarrow{\ d(E_\eta)\ }
   \underline{\mathrm{Hom}}(\Psi,\mathbb Q/\mathbb Z)_k)
 \right).
\end{equation}
If $u:P'\to P$ is a morphism of smooth group schemes over $S$, the
biextension of $(P'_\eta,Q_\eta)$ by $\underline G_{m\eta}$ obtained by
pulling $E_\eta$ back along $u_\eta$ extends to a biextension of $(P',Q)$
by $\underline G_{mS}$ if and only if (\SGARefLink{sga7.sga7i.eviii.n7-3-2-2}{7.3.2.2}) holds.  Symmetric results
hold upon interchanging the roles of $P$ and $Q$.

}{In the case of a biextension $E_\eta$ of $(P_\eta,Q_\eta)$ by
$\underline G_{m\eta}$, there is likewise a largest group subscheme $V$ of
$P$ such that the restriction of $E_\eta$ to $(V_\eta,Q_\eta)$ extends to
a biextension of $(V,Q)$ by $\underline G_{mS}$.  It is the open subgroup
of $P$ containing the generic fibre $P_\eta$ and characterized by the fact
that $V_0/V_0^{\circ}$ is the left annihilator in $\Phi$ of the pairing
(\SGARefLink{sga7.sga7i.eviii.n7-1-2}{7.1.2}), that is, the kernel of the corresponding homomorphism from $\Phi$
to $\underline{\mathrm{Hom}}(\Psi,(\mathbb Q/\mathbb Z)_k)$:
\begin{equation}\tag{\SGARefTarget{sga7.sga7i.eviii.n7-3-2-3}{7.3.2.3}}
 V_0/V_0^{\circ}
 =\ker\!\left(
   \Phi\xrightarrow{\ d(E_\eta)\ }
   \underline{\mathrm{Hom}}(\Psi,(\mathbb Q/\mathbb Z)_k)
 \right).
\end{equation}
If $u:P'\to P$ is a morphism of smooth group schemes over $S$, the
biextension of $(P'_\eta,Q_\eta)$ by $\underline G_{m\eta}$ obtained by
pulling $E_\eta$ back along $u_\eta$ extends to a biextension of $(P',Q)$
by $\underline G_{mS}$ if and only if (\SGARefLink{sga7.sga7i.eviii.n7-3-2-2}{7.3.2.2}) holds.  Symmetric results
hold upon interchanging the roles of $P$ and $Q$.

}{E7I-VIII-042}Of course, in general there is no largest \underline{pair} of subgroups
$P'\to P$, $Q'\to Q$ for which the restriction $E'_\eta$ of $E_\eta$ to a
biextension of $(P'_\eta,Q'_\eta)$ by $\underline G_{m\eta}$ extends to a
biextension of $(P',Q')$ by $\underline G_{mS}$.  One can only say that,
for a given pair $(P',Q')$, or more generally for a given pair of
morphisms $P'\to P$, $Q'\to Q$ of smooth group schemes, the desired
extension exists if and only if the restriction of the form $d(E_\eta)$
to
\[
 P'_0/{P'_0}^{\circ}\otimes Q'_0/{Q'_0}^{\circ}
\]
is zero.

\SGApage{I}{306}
\medskip
\noindent \SGARefTarget{sga7.sga7i.eviii.n7-3-3}{7.3.3}.  \underline{Symmetry}.  Let $E_\eta$ be a biextension of
$(P_\eta,Q_\eta)$ by $\underline G_{m\eta}$, and consider the transposed
biextension ${}^{s}E_\eta$ (VII \SGARefLink{sga7.sga7i.evii.n2-7}{2.7}), which is a biextension of
$(Q_\eta,P_\eta)$ by $\underline G_{m\eta}$.  The following diagram is
commutative:
\begin{equation}\tag{\SGARefTarget{sga7.sga7i.eviii.n7-3-3-1}{7.3.3.1}}
\begin{tikzcd}[row sep=large,column sep=huge]
 \Phi\otimes\Psi
   \arrow[dd,"\mathrm{sym}"']
   \arrow[rd,"d(E_\eta)"] \& \\
 \& (\mathbb Q/\mathbb Z)_k \\
 \Psi\otimes\Phi
   \arrow[ru,"d({}^{s}E_\eta)"'] \&
\end{tikzcd}
\qquad ,
\end{equation}
where the vertical arrow is the canonical symmetry.  (Compare this
statement with the \underline{antisymmetry} statement of \SGARefLink{sga7.sga7i.eviii.n2-2-10}{2.2.10}.)  It
follows in particular that, if $P=Q$ and $E_\eta$ is a
\underline{symmetric} biextension of $(P_\eta,Q_\eta)$ by
$\underline G_{m\eta}$, that is, if $E_\eta$ is isomorphic to
${}^{s}E_\eta$, then the bilinear form $d(E_\eta)$ on $\Phi\times\Phi$
with values in $(\mathbb Q/\mathbb Z)_k$ is \underline{symmetric} (and not
alternating, as in \SGARefLink{sga7.sga7i.eviii.n2-2-12}{2.2.12}!).  If it is zero, that is, if $E_\eta$ extends
to a biextension $E$ of $(P,P)$ by $\underline G_{mS}$, the latter is
necessarily symmetric as well: by \SGARefLink{sga7.sga7i.eviii.n7-1}{7.1}(a), every isomorphism
${}^{s}E_\eta\to E_\eta$ extends uniquely to an isomorphism
${}^{s}E\to E$.



\SGApage{I}{307}
\medskip
\noindent \SGARefTarget{sga7.sga7i.eviii.n7-3-4}{7.3.4}.  \underline{Reduction of obstructions of type (\SGARefLink{sga7.sga7i.eviii.n7-1-2}{7.1.2}) to
those of type (\SGARefLink{sga7.sga7i.eviii.n7-1-1}{7.1.1}).}

\SGAvar{Suppose, for definiteness, that there is an integer $n>0$ such that
$n\Phi=0$, and let
\[
 q\in Q(S)
\]
be a section of $Q$.  By an abuse of notation, write
\[
 q_0\in\Psi(k)
\]
for the image in $\Psi(k)$ of the value of $q$ at $s$.  Regarding $E_\eta$
as an extension of $P_\eta Q_\eta$ by
$\underline G_{mQ_\eta}$ over $Q_\eta$, pullback gives an extension
\[
 q_\eta^*(E_\eta),\qquad
 \text{an extension of $P_\eta$ by $\underline G_{m\eta}$},
\]
and hence an obstruction
\begin{equation}\tag{\SGARefTarget{sga7.sga7i.eviii.n7-3-4-1}{7.3.4.1}}
 d_q\overset{\mathrm{dfn}}{=}
 d\bigl(q_\eta^*(E_*)\bigr):
 \Phi\longrightarrow(\mathbb Q/\mathbb Z)_k.
\end{equation}
With $d=d(E_\eta)$, one has the relation
\begin{equation}\tag{\SGARefTarget{sga7.sga7i.eviii.n7-3-4-2}{7.3.4.2}}
 d_q(x)=d(x,q_0),
\end{equation}
valid for every point $x$ of $\Phi$ with values in a $k$-scheme $T$.

}{Suppose, for definiteness, that there is an integer $n>0$ such that
$n\Phi=0$, and let
\[
 q\in Q(S)
\]
be a section of $Q$.  By an abuse of notation, write
\[
 q_0\in\Psi(k)
\]
for the image in $\Psi(k)$ of the value of $q$ at $s$.  Regarding $E_\eta$
as an extension of $P_\eta\times_\eta Q_\eta$ by
$\underline G_{mQ_\eta}$ over $Q_\eta$, pullback gives an extension
\[
 q_\eta^*(E_\eta),\qquad
 \text{an extension of $P_\eta$ by $\underline G_{m\eta}$},
\]
and hence an obstruction
\begin{equation}\tag{\SGARefTarget{sga7.sga7i.eviii.n7-3-4-1}{7.3.4.1}}
 d_q\overset{\mathrm{dfn}}{=}
 d\bigl(q_\eta^*(E_\eta)\bigr):
 \Phi\longrightarrow(\mathbb Q/\mathbb Z)_k.
\end{equation}
With $d=d(E_\eta)$, one has the relation
\begin{equation}\tag{\SGARefTarget{sga7.sga7i.eviii.n7-3-4-2}{7.3.4.2}}
 d_q(x)=d(x,q_0),
\end{equation}
valid for every point $x$ of $\Phi$ with values in a $k$-scheme $T$.

}{E7I-VIII-043}\smallskip
\noindent \SGARefTarget{sga7.sga7i.eviii.n7-3-4-3}{7.3.4.3}.  For every $q\in Q(S)$, this formula determines $d_q$
in terms of the pairing $d$; conversely, it plainly recovers $d$ from the
invariants $d_q$, at least when every point of $\Psi$ comes from a section
of $Q$.  By Hensel's lemma, this condition holds automatically when $S$ is
strictly local.  The pairing $d$ is known once it is known over the
separable closure of $k$, and its formation commutes with etale localization
on $S$ by \SGARefLink{sga7.sga7i.eviii.n7-3-5-3}{7.3.5.3} below.  Thus, after replacing $S$ by its strict
localization if necessary, the determination of $d$ reduces to that of the
invariants $d_q$.

\SGApage{I}{\(\approx\)308}
In particular, when every point of $\Psi$ comes from a section $q$ of $Q$,
the pairing $d$ is zero---that is, $E_\eta$ extends to a biextension $E$ of
$(P,Q)$ by $\underline G_m$---if and only if $d_q$ is zero for every
section $q$ of $Q$, that is, if and only if $q_\eta^*(E_\eta)$ extends to
an extension of $P$ by $\underline G_m$.  With no hypothesis on $\Psi$,
the analogous criterion remains valid provided that the $q$ are allowed to
be points with values in local etale $S$-schemes $S'$.

\smallskip
\noindent \SGARefTarget{sga7.sga7i.eviii.n7-3-4-4}{7.3.4.4}.  Conversely, starting from the situation of \SGARefLink{sga7.sga7i.eviii.n7-1}{7.1}(a), the
obstruction (\SGARefLink{sga7.sga7i.eviii.n7-1-1}{7.1.1}) may be interpreted as an obstruction of type (\SGARefLink{sga7.sga7i.eviii.n7-1-2}{7.1.2}):
take $Q=\mathbb Z_S$ and the biextension of $(P_\eta,\mathbb Z_\eta)$
deduced from $E_\eta$.  Indeed, it suffices to apply the preceding
construction with $q$ equal to the section $1$ of $\mathbb Z_S$.

\SGAvar{\medskip
\noindent \SGARefTarget{sga7.sga7i.eviii.n7-3-5}{7.3.5}.  \underline{Effect of a base change $S'\to S$.}  Suppose
given a dominant local morphism of traits
\[
 f:S'\longrightarrow S,
\]
and denote by $\eta'$ the generic point of $S'$, by $s'$ its closed point,
and put $k'=k(s')$.  If $S=\operatorname{Spec}(V)$ and
$S'=\operatorname{Spec}(V')$, then $V'$ is a discrete valuation ring
dominating the discrete valuation ring $V$.  Put
\begin{equation}\tag{\SGARefTarget{sga7.sga7i.eviii.n7-3-5-1}{7.3.5.1}}
 e(S'/S)=e(V'/V)=\operatorname{length}_{V'}(V'/\mathfrak m V),
\end{equation}
where $\mathfrak m$ is the maximal ideal of $V$.

}{\medskip
\noindent \SGARefTarget{sga7.sga7i.eviii.n7-3-5}{7.3.5}.  \underline{Effect of a base change $S'\to S$.}  Suppose
given a dominant local morphism of traits
\[
 f:S'\longrightarrow S,
\]
and denote by $\eta'$ the generic point of $S'$, by $s'$ its closed point,
and put $k'=k(s')$.  If $S=\operatorname{Spec}(V)$ and
$S'=\operatorname{Spec}(V')$, then $V'$ is a discrete valuation ring
dominating the discrete valuation ring $V$.  Put
\begin{equation}\tag{\SGARefTarget{sga7.sga7i.eviii.n7-3-5-1}{7.3.5.1}}
 e(S'/S)=e(V'/V)=\operatorname{length}_{V'}(V'/\mathfrak m V'),
\end{equation}
where $\mathfrak m$ is the maximal ideal of $V$.

}{E7I-VIII-044}Let now $P$ be as in \SGARefLink{sga7.sga7i.eviii.n7-1}{7.1}(a), let $E_\eta$ be an extension of $P_\eta$ by
$\underline G_{m\eta}$, put $P'=P\times_S S'$, and let $E'_{\eta'}$ be
the extension of $P'_{\eta'}$ by $\underline G_{m\eta'}$ obtained by
pulling back $E_\eta$.  With the notation of \SGARefLink{sga7.sga7i.eviii.n4-1}{4.1}(a), there is a canonical
isomorphism
\[
 \Phi'\overset{\mathrm{dfn}}{=}P'_0/{P'_0}^{\circ}
 \simeq\Phi\otimes_k k'\simeq\Phi_{k'},
\]
and hence one may consider
\[
 d(E_\eta)_{k'}:\Phi'\longrightarrow(\mathbb Q/\mathbb Z)_{k'}.
\]
Then
\begin{equation}\tag{\SGARefTarget{sga7.sga7i.eviii.n7-3-5-2}{7.3.5.2}}
 d(E'_{\eta'})=e(S'/S)\,d(E_\eta)_{k'}.
\end{equation}

\SGApage{I}{\(\approx\)309}
Likewise, under the hypotheses of \SGARefLink{sga7.sga7i.eviii.n7-1}{7.1}(b) and with the evident notation,
(\SGARefLink{sga7.sga7i.eviii.n7-3-5-2}{7.3.5.2}) holds with $d(E_\eta)$ now denoting the pairing (\SGARefLink{sga7.sga7i.eviii.n7-1-2}{7.1.2}),
$d(E_\eta)_{k'}$ the pairing
\[
 \Phi'\otimes\Psi'\longrightarrow(\mathbb Q/\mathbb Z)_{k'}
\]
obtained from it by base change $k\to k'$, and $E'_{\eta'}$ the
biextension of $(P'_{\eta'},Q'_{\eta'})$ by
$\underline G_{m\eta'}$ obtained from the biextension $E_\eta$ of
$(P_\eta,Q_\eta)$ by $\underline G_{m\eta}$ through the base change
$\eta'\to\eta$.

\smallskip
\noindent \SGARefTarget{sga7.sga7i.eviii.n7-3-5-3}{7.3.5.3}.  An interesting case is $e(S'/S)=1$, that is,
\[
 \mathfrak m'=\mathfrak m V',
\]
where $\mathfrak m'$ denotes the maximal ideal of $V'$.  This holds in
particular if $V'$ is etale over $V$, or if $V'$ is the henselization,
strict henselization, or completion of $V$.  In this case (\SGARefLink{sga7.sga7i.eviii.n7-3-5-2}{7.3.5.2}) simply
says that \underline{the formation of obstruction (\SGARefLink{sga7.sga7i.eviii.n7-1-1}{7.1.1}), respectively
(\SGARefLink{sga7.sga7i.eviii.n7-1-2}{7.1.2}), commutes with the base change under consideration $S'\to S$}.

\smallskip
\noindent \SGARefTarget{sga7.sga7i.eviii.n7-3-5-4}{7.3.5.4}.  On the other hand, (\SGARefLink{sga7.sga7i.eviii.n7-3-5-2}{7.3.5.2}) in the general case shows
that, for every extension $E_\eta$ as in \SGARefLink{sga7.sga7i.eviii.n7-1}{7.1}(a) (respectively every
biextension $E_\eta$ as in \SGARefLink{sga7.sga7i.eviii.n7-1}{7.1}(b)), there is a \underline{finite} morphism
of traits $S'\to S$ such that $k(\eta')$ is a separable extension of
$k(\eta)$ and such that the extension (respectively biextension)
$E'_{\eta'}$ obtained from $E_\eta$ by base change $\eta'\to\eta$ extends
to an extension (respectively biextension) $E'$ over all of $S'$.  Indeed,
it is enough to choose $S'/S$ so that $e(S'/S)$ is a multiple of the
integer $n>0$ occurring in \SGARefLink{sga7.sga7i.eviii.n7-1}{7.1}, and such an extension can always be found
with $k(\eta')$ separable over $k(\eta)$.



\SGApage{I}{310}
\medskip
\noindent\underline{Proposition \SGARefTarget{sga7.sga7i.eviii.n7-4}{7.4}.  Let $S$ be a reduced scheme.}

\smallskip
\noindent a) \underline{Let $P$ be a smooth group scheme over $S$ whose
fibres at the maximal points of $S$ are connected.  Then the natural
functor}
\begin{equation}\tag{\SGARefTarget{sga7.sga7i.eviii.n7-4-1}{7.4.1}}
 \mathrm{EXT}(P,\underline G_{mS})
 \longrightarrow
 \mathrm{TORSRIG}(P,e_P;\underline G_{mS})
\end{equation}
\underline{from the category of extensions of $P$ by $\underline G_{mS}$
to the category of torsors under $\underline G_{mP}$ rigidified along the
identity section of $P$ is fully faithful.  If $S$ is normal and $P$ has
connected fibres, an object $E$ of the target of (\SGARefLink{sga7.sga7i.eviii.n7-4-1}{7.4.1}) belongs to the
essential image of this functor if and only if, for every maximal point
$\eta$ of $S$, the rigidified torsor $E_\eta$ on $P_\eta$ comes from an
extension of $P_\eta$ by $\underline G_{m\eta}$.}

\smallskip
\noindent b) \underline{Let $P,Q$ be two smooth group schemes over $S$
whose fibres at the maximal points of $S$ are connected.  Then the natural
functor}
\begin{equation}\tag{\SGARefTarget{sga7.sga7i.eviii.n7-4-2}{7.4.2}}
 \mathrm{BIEXT}(P,Q;\underline G_{mS})
 \longrightarrow
 \mathrm{TORSBIRIG}(P,Q;\underline G_{mS})
\end{equation}
\underline{from the category of biextensions of $(P,Q)$ by
$\underline G_{mS}$ to the category of torsors under
$\underline G_{m,P\times_S Q}$ birigidified along the pair of identity
sections of $P$ and $Q$ is fully faithful.  If $S$ is normal and either $P$
or $Q$ has connected fibres, an object $E$ of the target of (\SGARefLink{sga7.sga7i.eviii.n7-4-2}{7.4.2}) belongs
to the essential image of the functor if and only if, for every maximal
point $\eta$ of $S$, $E_\eta$ comes from a biextension of
$(P_\eta,Q_\eta)$ by $\underline G_{m\eta}$.}

\SGApage{I}{\(\approx\)311}
We prove (a); the essentially identical proof of (b) is left to the reader.
For the full-faithfulness assertion, note that \SGARefLink{sga7.sga7i.eviii.n4-1}{4.1} remains valid if $X$ and
$Y$ are flat schemes over a reduced scheme $S$ whose fibres at the maximal
points satisfy the conditions stated in \SGARefLink{sga7.sga7i.eviii.n4-1}{4.1}.  Indeed, the family of maximal
fibres of $X\times_S Y$ (which is flat over the reduced scheme $S$) is
schematically dense in $X\times_S Y$.  Hence a morphism from
$X\times_S Y$ to the separated $S$-scheme $\underline G_{mS}$ is trivial
as soon as it is trivial on the maximal fibres.  Thus the full faithfulness
of (\SGARefLink{sga7.sga7i.eviii.n7-4-1}{7.4.1}) is proved as in \SGARefLink{sga7.sga7i.eviii.n4-2}{4.2}, and the characterization of its essential
image given there remains valid.

It remains to prove that, if the rigidified torsor $E$ on $P$ is such that
the $E_\eta$ come from extensions, then so does $E$, assuming that $S$ is
normal and $P$ has connected fibres.  The latter condition implies that
$P$ is of finite presentation over $S$ (SGA 3 $\mathrm{VI}_B$ 5.5).
Full faithfulness also shows that the question is local on $S$, so that we
may suppose $S$ affine and local.  The usual limit argument, using
Nagata's theorem on the finiteness of the integral closure of a reduced
finitely generated $\mathbb Z$-algebra (EGA IV 7.8), reduces us further to
the case in which $S$ is noetherian.

Then \SGARefLink{sga7.sga7i.eviii.n7-1}{7.1}(a) shows that there is an open subset $U\subset S$ whose complement
has codimension at least $2$ and such that $E_U$ comes from an extension of
$P_U$ by $\underline G_{mU}$.  In other words, the desired isomorphism
\[
 \pi^*(E)\simeq
 \operatorname{pr}_1^*(E)\operatorname{pr}_2^*(E)
\]
can be constructed over $(P\times_S P)_U$.  This is an open subset of the
normal scheme $P\times_S P$ whose complement has codimension at least $2$.
The isomorphism on $(P\times_S P)_U$ therefore extends over all of
$P\times_S P$, completing the proof.

\SGApage{I}{312}
\medskip
\noindent\underline{Corollary \SGARefTarget{sga7.sga7i.eviii.n7-5}{7.5}.  Let $S$ be a normal scheme and let $P$
and $Q$ be two smooth group schemes over $S$, with connected fibres at the
maximal points, either $P$ or $Q$ having connected fibres.  Suppose in
addition that, for every maximal point $\eta$ of $S$, either $P_\eta$ or
$Q_\eta$ is an extension of an abelian scheme by a torus (or, more
generally, admits a composition series whose factors are abelian schemes,
tori, or groups $\underline G_a$).  Then the functor (\SGARefLink{sga7.sga7i.eviii.n7-4-2}{7.4.2}) is an
equivalence of categories.}

This is an immediate consequence of \SGARefLink{sga7.sga7i.eviii.n7-4}{7.4}(b) and \SGARefLink{sga7.sga7i.eviii.n4-7}{4.7}.

\bigskip
\noindent\underline{BIBLIOGRAPHY}

\begin{enumerate}
\item[{[1]}] M. Demazure, J. Giraud, M. Raynaud,
  \emph{Sch\'emas ab\'eliens}, S\'eminaire Orsay 1967/68, forthcoming.
\item[{[2]}] M. Raynaud, \emph{Faisceaux amples sur les sch\'emas en groupes
  et les espaces homog\`enes}, Paris thesis, 1968 (forthcoming in Lecture
  Notes, Springer).
\item[{[3]}] M. Rosenlicht, \emph{Some basic theorems on algebraic groups},
  Amer. Jour. Math. \textbf{78}, no.~2 (1956), pp.~401--443.
\item[{[3bis]}] M. Rosenlicht, \emph{Some rationality questions on algebraic
  groups}, Annali di Matematica, Series 4, no.~43 (1957), pp.~25--50.
\item[{[4]}] J.-P. Serre, \emph{Corps locaux}, Act. Sci. Ind. 1296 (1962),
  Hermann, Paris.
\end{enumerate}


\markboth{Exposé VIII}{Exposé VIII}

\ifnum\SGAform=1
\clearpage\section*{Critical apparatus}\begingroup\small\raggedright\sloppy
Readings are given in the exact TeX notation of the two English witnesses. Decimal suffixes identify the individual differences in each marked passage.

\par\medskip\noindent\hypertarget{app-E7I-VIII-001.1}{}\textbf{E7I-VIII-001.1}\quad printed folio 220.\par
\noindent Printed reading: {\ttfamily\footnotesize From\allowbreak{}}\par
\noindent Corrected reading: {\ttfamily\footnotesize Taking\allowbreak{} the\allowbreak{} standard\allowbreak{} truncation\allowbreak{} of\allowbreak{} the\allowbreak{} derived\allowbreak{} tensor\allowbreak{} product,\allowbreak{} whose\allowbreak{} homology\allowbreak{} objects\allowbreak{} are\allowbreak{} the\allowbreak{} \$\textbackslash{}Tor\_i(P,Q)\$,\allowbreak{} one\allowbreak{} obtains\allowbreak{} from\allowbreak{}}\par
\noindent Source: {\ttfamily\footnotesize ADOPT-F7I-VIII-001.1\allowbreak{};\allowbreak{} ADOPT-F7I-VIII-001.2\allowbreak{};\allowbreak{} ADOPT-F7I-VIII-001.3\allowbreak{};\allowbreak{} WITNESS-F7I-VIII-001\allowbreak{}.4;\allowbreak{} WITNESS-F7I-VIII-001\allowbreak{}.5;\allowbreak{} WITNESS-F7I-VIII-001\allowbreak{}.6;\allowbreak{} WITNESS-F7I-VIII-001\allowbreak{}.7;\allowbreak{} WITNESS-F7I-VIII-001\allowbreak{}.8;\allowbreak{} WITNESS-F7I-VIII-001\allowbreak{}.9;\allowbreak{} fr/sga7i/diplomatic/\allowbreak{}expose\_VIII\_body.tex\allowbreak{};\allowbreak{} fr/sga7i/current/exp\allowbreak{}ose\_VIII\_body.tex\allowbreak{}}\par

\par\medskip\noindent\hypertarget{app-E7I-VIII-001.2}{}\textbf{E7I-VIII-001.2}\quad printed folio 220.\par
\noindent Printed reading: {\ttfamily\footnotesize one\allowbreak{} deduces\allowbreak{} }\par
\noindent Corrected reading: \emph{(empty)}\par
\noindent Source: {\ttfamily\footnotesize ADOPT-F7I-VIII-001.1\allowbreak{};\allowbreak{} ADOPT-F7I-VIII-001.2\allowbreak{};\allowbreak{} ADOPT-F7I-VIII-001.3\allowbreak{};\allowbreak{} WITNESS-F7I-VIII-001\allowbreak{}.4;\allowbreak{} WITNESS-F7I-VIII-001\allowbreak{}.5;\allowbreak{} WITNESS-F7I-VIII-001\allowbreak{}.6;\allowbreak{} WITNESS-F7I-VIII-001\allowbreak{}.7;\allowbreak{} WITNESS-F7I-VIII-001\allowbreak{}.8;\allowbreak{} WITNESS-F7I-VIII-001\allowbreak{}.9;\allowbreak{} fr/sga7i/diplomatic/\allowbreak{}expose\_VIII\_body.tex\allowbreak{};\allowbreak{} fr/sga7i/current/exp\allowbreak{}ose\_VIII\_body.tex\allowbreak{}}\par

\par\medskip\noindent\hypertarget{app-E7I-VIII-001.3}{}\textbf{E7I-VIII-001.3}\quad printed folio 220.\par
\noindent Printed reading: {\ttfamily\footnotesize ,\allowbreak{} taking\allowbreak{} into\allowbreak{} account\allowbreak{} that\allowbreak{} the\allowbreak{} homology\allowbreak{} objects\allowbreak{} of\allowbreak{} \$P\textbackslash{}mathbin\{\textbackslash{}overset\{\allowbreak{}\textbackslash{}mathbb\allowbreak{} L\}\{\textbackslash{}otimes\}\}Q\$\allowbreak{} are\allowbreak{} the\allowbreak{} \$\textbackslash{}Tor\_i(P,Q)\$,\allowbreak{} which\allowbreak{} vanish\allowbreak{} for\allowbreak{} \$i\textbackslash{}geqslant1\$\allowbreak{}}\par
\noindent Corrected reading: \emph{(empty)}\par
\noindent Source: {\ttfamily\footnotesize ADOPT-F7I-VIII-001.1\allowbreak{};\allowbreak{} ADOPT-F7I-VIII-001.2\allowbreak{};\allowbreak{} ADOPT-F7I-VIII-001.3\allowbreak{};\allowbreak{} WITNESS-F7I-VIII-001\allowbreak{}.4;\allowbreak{} WITNESS-F7I-VIII-001\allowbreak{}.5;\allowbreak{} WITNESS-F7I-VIII-001\allowbreak{}.6;\allowbreak{} WITNESS-F7I-VIII-001\allowbreak{}.7;\allowbreak{} WITNESS-F7I-VIII-001\allowbreak{}.8;\allowbreak{} WITNESS-F7I-VIII-001\allowbreak{}.9;\allowbreak{} fr/sga7i/diplomatic/\allowbreak{}expose\_VIII\_body.tex\allowbreak{};\allowbreak{} fr/sga7i/current/exp\allowbreak{}ose\_VIII\_body.tex\allowbreak{}}\par

\par\medskip\noindent\hypertarget{app-E7I-VIII-001.4}{}\textbf{E7I-VIII-001.4}\quad printed folio 220.\par
\noindent Printed reading: \emph{(empty)}\par
\noindent Corrected reading: {\ttfamily\footnotesize \textbackslash{}longrightarrow\allowbreak{}}\par
\noindent Source: {\ttfamily\footnotesize fr/sga7i/\{diplomatic\allowbreak{},current\}/expose\_VII\allowbreak{}I\_body.tex;\allowbreak{} folio\allowbreak{} 220;\allowbreak{} formula\allowbreak{} 47;\allowbreak{} port/decisions/decis\allowbreak{}ions\_en7I\_r2\_002.jso\allowbreak{}n\#C780;\allowbreak{} fixes/FIXLOG\_7I\_CURR\allowbreak{}ENT.jsonl\#C780;\allowbreak{} 7I:421\allowbreak{}}\par

\par\medskip\noindent\hypertarget{app-E7I-VIII-001.5}{}\textbf{E7I-VIII-001.5}\quad printed folio 220.\par
\noindent Printed reading: {\ttfamily\footnotesize P\allowbreak{}}\par
\noindent Corrected reading: {\ttfamily\footnotesize G\allowbreak{}}\par
\noindent Source: {\ttfamily\footnotesize fr/sga7i/\{diplomatic\allowbreak{},current\}/expose\_VII\allowbreak{}I\_body.tex;\allowbreak{} folio\allowbreak{} 220;\allowbreak{} formula\allowbreak{} 47\allowbreak{}}\par

\par\medskip\noindent\hypertarget{app-E7I-VIII-002.1}{}\textbf{E7I-VIII-002.1}\quad printed folio 222.\par
\noindent Printed reading: \emph{(empty)}\par
\noindent Corrected reading: {\ttfamily\footnotesize )\allowbreak{}}\par
\noindent Source: {\ttfamily\footnotesize ED-EN-FORMULA-G039;\allowbreak{} fr/sga7i/\{diplomatic\allowbreak{},current\}/expose\_VII\allowbreak{}I\_body.tex;\allowbreak{} folio\allowbreak{} 222\allowbreak{}}\par

\par\medskip\noindent\hypertarget{app-E7I-VIII-002.2}{}\textbf{E7I-VIII-002.2}\quad printed folio 223.\par
\noindent Printed reading: {\ttfamily\footnotesize 4\allowbreak{}}\par
\noindent Corrected reading: {\ttfamily\footnotesize 5\allowbreak{}}\par
\noindent Source: {\ttfamily\footnotesize 7I:424;\allowbreak{} fixes/MISPRINTS\_7I.j\allowbreak{}son\#424\allowbreak{}}\par

\par\medskip\noindent\hypertarget{app-E7I-VIII-002.3}{}\textbf{E7I-VIII-002.3}\quad printed folio 223.\par
\noindent Printed reading: {\ttfamily\footnotesize last\allowbreak{}}\par
\noindent Corrected reading: {\ttfamily\footnotesize next\allowbreak{}}\par
\noindent Source: {\ttfamily\footnotesize 7I:426;\allowbreak{} fixes/MISPRINTS\_7I.j\allowbreak{}son\#426\allowbreak{}}\par

\par\medskip\noindent\hypertarget{app-E7I-VIII-002.4}{}\textbf{E7I-VIII-002.4}\quad printed folio 224.\par
\noindent Printed reading: {\ttfamily\footnotesize 1\allowbreak{}}\par
\noindent Corrected reading: {\ttfamily\footnotesize i\allowbreak{}}\par
\noindent Source: {\ttfamily\footnotesize fr/sga7i/\{diplomatic\allowbreak{},current\}/expose\_VII\allowbreak{}I\_body.tex;\allowbreak{} folio\allowbreak{} 224;\allowbreak{} formula\allowbreak{} 87\allowbreak{}}\par

\par\medskip\noindent\hypertarget{app-E7I-VIII-003.1}{}\textbf{E7I-VIII-003.1}\quad printed folio 225.\par
\noindent Printed reading: {\ttfamily\footnotesize P\allowbreak{}}\par
\noindent Corrected reading: {\ttfamily\footnotesize Q\textasciicircum{}*\allowbreak{}}\par
\noindent Source: {\ttfamily\footnotesize fr/sga7i/\{diplomatic\allowbreak{},current\}/expose\_VII\allowbreak{}I\_body.tex;\allowbreak{} folio\allowbreak{} 225;\allowbreak{} formula\allowbreak{} 113\allowbreak{}}\par

\par\medskip\noindent\hypertarget{app-E7I-VIII-004.1}{}\textbf{E7I-VIII-004.1}\quad printed folio 225.\par
\noindent Printed reading: \emph{(empty)}\par
\noindent Corrected reading: {\ttfamily\footnotesize  and\allowbreak{} \textbackslash{}SGARefLink\{sga7.sga\allowbreak{}7i.eviii.n1-5\}\{1.5\},\allowbreak{} respectively\allowbreak{}}\par
\noindent Source: {\ttfamily\footnotesize port/decisions/decis\allowbreak{}ions\_en7I\_r2\_002.jso\allowbreak{}n\#C783;\allowbreak{} 7I:433\allowbreak{}}\par

\par\medskip\noindent\hypertarget{app-E7I-VIII-004.2}{}\textbf{E7I-VIII-004.2}\quad printed folio 226.\par
\noindent Printed reading: {\ttfamily\footnotesize G\allowbreak{}}\par
\noindent Corrected reading: {\ttfamily\footnotesize Q\allowbreak{}}\par
\noindent Source: {\ttfamily\footnotesize fr/sga7i/\{diplomatic\allowbreak{},current\}/expose\_VII\allowbreak{}I\_body.tex;\allowbreak{} folio\allowbreak{} 226;\allowbreak{} formula\allowbreak{} 131\allowbreak{}}\par

\par\medskip\noindent\hypertarget{app-E7I-VIII-004.3}{}\textbf{E7I-VIII-004.3}\quad printed folio 226.\par
\noindent Printed reading: {\ttfamily\footnotesize Q\allowbreak{}}\par
\noindent Corrected reading: {\ttfamily\footnotesize G\allowbreak{}}\par
\noindent Source: {\ttfamily\footnotesize fr/sga7i/\{diplomatic\allowbreak{},current\}/expose\_VII\allowbreak{}I\_body.tex;\allowbreak{} folio\allowbreak{} 226\allowbreak{}}\par

\par\medskip\noindent\hypertarget{app-E7I-VIII-005.1}{}\textbf{E7I-VIII-005.1}\quad printed folio 228.\par
\noindent Printed reading: {\ttfamily\footnotesize Q\allowbreak{}}\par
\noindent Corrected reading: {\ttfamily\footnotesize P\allowbreak{}}\par
\noindent Source: {\ttfamily\footnotesize fr/sga7i/\{diplomatic\allowbreak{},current\}/expose\_VII\allowbreak{}I\_body.tex;\allowbreak{} folio\allowbreak{} 228;\allowbreak{} formula\allowbreak{} 184\allowbreak{}}\par

\par\medskip\noindent\hypertarget{app-E7I-VIII-006.1}{}\textbf{E7I-VIII-006.1}\quad printed folio 229.\par
\noindent Printed reading: \emph{(empty)}\par
\noindent Corrected reading: {\ttfamily\footnotesize (\allowbreak{}}\par
\noindent Source: {\ttfamily\footnotesize fr/sga7i/\{diplomatic\allowbreak{},current\}/expose\_VII\allowbreak{}I\_body.tex;\allowbreak{} folio\allowbreak{} 229;\allowbreak{} formula\allowbreak{} 240\allowbreak{}}\par

\par\medskip\noindent\hypertarget{app-E7I-VIII-006.2}{}\textbf{E7I-VIII-006.2}\quad printed folio 229.\par
\noindent Printed reading: \emph{(empty)}\par
\noindent Corrected reading: {\ttfamily\footnotesize \textbackslash{}bigl(\allowbreak{}}\par
\noindent Source: {\ttfamily\footnotesize ED-EN-FORMULA-G041;\allowbreak{} fr/sga7i/\{diplomatic\allowbreak{},current\}/expose\_VII\allowbreak{}I\_body.tex;\allowbreak{} folio\allowbreak{} 229\allowbreak{}}\par

\par\medskip\noindent\hypertarget{app-E7I-VIII-007.1}{}\textbf{E7I-VIII-007.1}\quad printed folio 232.\par
\noindent Printed reading: \emph{(empty)}\par
\noindent Corrected reading: {\ttfamily\footnotesize 1.\allowbreak{}}\par
\noindent Source: {\ttfamily\footnotesize ED-EN-NONMATH-I-VIII\allowbreak{}232;\allowbreak{} 7I:443;\allowbreak{} 7I:444;\allowbreak{} 7I:445;\allowbreak{} fixes/MISPRINTS\_7I.j\allowbreak{}son\#443;\allowbreak{} fixes/MISPRINTS\_7I.j\allowbreak{}son\#444;\allowbreak{} fixes/MISPRINTS\_7I.j\allowbreak{}son\#445;\allowbreak{} F7I-VIII-013.1,\allowbreak{} F7I-VIII-013.2,\allowbreak{} F7I-VIII-013.3\allowbreak{}}\par

\par\medskip\noindent\hypertarget{app-E7I-VIII-007.2}{}\textbf{E7I-VIII-007.2}\quad printed folio 232.\par
\noindent Printed reading: \emph{(empty)}\par
\noindent Corrected reading: {\ttfamily\footnotesize .1\allowbreak{}}\par
\noindent Source: {\ttfamily\footnotesize ED-EN-NONMATH-I-VIII\allowbreak{}232;\allowbreak{} 7I:443;\allowbreak{} 7I:444;\allowbreak{} 7I:445;\allowbreak{} fixes/MISPRINTS\_7I.j\allowbreak{}son\#443;\allowbreak{} fixes/MISPRINTS\_7I.j\allowbreak{}son\#444;\allowbreak{} fixes/MISPRINTS\_7I.j\allowbreak{}son\#445;\allowbreak{} F7I-VIII-013.1,\allowbreak{} F7I-VIII-013.2,\allowbreak{} F7I-VIII-013.3\allowbreak{}}\par

\par\medskip\noindent\hypertarget{app-E7I-VIII-007.3}{}\textbf{E7I-VIII-007.3}\quad printed folio 232.\par
\noindent Printed reading: \emph{(empty)}\par
\noindent Corrected reading: {\ttfamily\footnotesize .1\allowbreak{}}\par
\noindent Source: {\ttfamily\footnotesize ED-EN-NONMATH-I-VIII\allowbreak{}232;\allowbreak{} 7I:443;\allowbreak{} 7I:444;\allowbreak{} 7I:445;\allowbreak{} fixes/MISPRINTS\_7I.j\allowbreak{}son\#443;\allowbreak{} fixes/MISPRINTS\_7I.j\allowbreak{}son\#444;\allowbreak{} fixes/MISPRINTS\_7I.j\allowbreak{}son\#445;\allowbreak{} F7I-VIII-013.1,\allowbreak{} F7I-VIII-013.2,\allowbreak{} F7I-VIII-013.3\allowbreak{}}\par

\par\medskip\noindent\hypertarget{app-E7I-VIII-008.1}{}\textbf{E7I-VIII-008.1}\quad printed folio 235.\par
\noindent Printed reading: {\ttfamily\footnotesize isomorphism\allowbreak{}}\par
\noindent Corrected reading: {\ttfamily\footnotesize homomorphism\allowbreak{}}\par
\noindent Source: {\ttfamily\footnotesize 7I:448;\allowbreak{} fixes/MISPRINTS\_7I.j\allowbreak{}son\#448\allowbreak{}}\par

\par\medskip\noindent\hypertarget{app-E7I-VIII-009.1}{}\textbf{E7I-VIII-009.1}\quad printed folio 237.\par
\noindent Printed reading: \emph{(empty)}\par
\noindent Corrected reading: {\ttfamily\footnotesize  pq\allowbreak{}}\par
\noindent Source: {\ttfamily\footnotesize fr/sga7i/\{diplomatic\allowbreak{},current\}/expose\_VII\allowbreak{}I\_body.tex;\allowbreak{} folio\allowbreak{} 237\allowbreak{}}\par

\par\medskip\noindent\hypertarget{app-E7I-VIII-010.1}{}\textbf{E7I-VIII-010.1}\quad printed folio 238.\par
\noindent Printed reading: {\ttfamily\footnotesize \textbackslash{}varphi\allowbreak{}}\par
\noindent Corrected reading: {\ttfamily\footnotesize \textbackslash{}alpha\allowbreak{}}\par
\noindent Source: {\ttfamily\footnotesize ED-EN-FORMULA-G043;\allowbreak{} fr/sga7i/\{diplomatic\allowbreak{},current\}/expose\_VII\allowbreak{}I\_body.tex;\allowbreak{} folio\allowbreak{} 238\allowbreak{}}\par

\par\medskip\noindent\hypertarget{app-E7I-VIII-010.2}{}\textbf{E7I-VIII-010.2}\quad printed folio 238.\par
\noindent Printed reading: {\ttfamily\footnotesize ;G\allowbreak{}}\par
\noindent Corrected reading: \emph{(empty)}\par
\noindent Source: {\ttfamily\footnotesize ED-EN-FORMULA-G043;\allowbreak{} fr/sga7i/\{diplomatic\allowbreak{},current\}/expose\_VII\allowbreak{}I\_body.tex;\allowbreak{} folio\allowbreak{} 238\allowbreak{}}\par

\par\medskip\noindent\hypertarget{app-E7I-VIII-011.1}{}\textbf{E7I-VIII-011.1}\quad printed folio 238.\par
\noindent Printed reading: {\ttfamily\footnotesize 4\allowbreak{}}\par
\noindent Corrected reading: {\ttfamily\footnotesize 5\allowbreak{}}\par
\noindent Source: {\ttfamily\footnotesize 7I:452;\allowbreak{} fixes/MISPRINTS\_7I.j\allowbreak{}son\#452\allowbreak{}}\par

\par\medskip\noindent\hypertarget{app-E7I-VIII-012.1}{}\textbf{E7I-VIII-012.1}\quad printed folio 240.\par
\noindent Printed reading: \emph{(4 espace(s), 2 retour(s) de ligne dans le témoin TeX)}\par
\noindent Corrected reading: {\ttfamily\footnotesize  \textbackslash{}bigr)\allowbreak{} }\par
\noindent Source: {\ttfamily\footnotesize fr/sga7i/\{diplomatic\allowbreak{},current\}/expose\_VII\allowbreak{}I\_body.tex;\allowbreak{} folio\allowbreak{} 240;\allowbreak{} formula\allowbreak{} 528\allowbreak{}}\par

\par\medskip\noindent\hypertarget{app-E7I-VIII-013.1}{}\textbf{E7I-VIII-013.1}\quad printed folio 243.\par
\noindent Printed reading: {\ttfamily\footnotesize nQnG\allowbreak{}}\par
\noindent Corrected reading: {\ttfamily\footnotesize nQ,G\allowbreak{}}\par
\noindent Source: {\ttfamily\footnotesize fr/sga7i/\{diplomatic\allowbreak{},current\}/expose\_VII\allowbreak{}I\_body.tex;\allowbreak{} folio\allowbreak{} 243;\allowbreak{} formula\allowbreak{} 582\allowbreak{}}\par

\par\medskip\noindent\hypertarget{app-E7I-VIII-014.1}{}\textbf{E7I-VIII-014.1}\quad printed folio 243.\par
\noindent Printed reading: {\ttfamily\footnotesize 5\allowbreak{}}\par
\noindent Corrected reading: {\ttfamily\footnotesize 3\allowbreak{}}\par
\noindent Source: {\ttfamily\footnotesize 7I:455;\allowbreak{} fixes/MISPRINTS\_7I.j\allowbreak{}son\#455\allowbreak{}}\par

\par\medskip\noindent\hypertarget{app-E7I-VIII-015.1}{}\textbf{E7I-VIII-015.1}\quad printed folio 245.\par
\noindent Printed reading: {\ttfamily\footnotesize \textbackslash{}to\allowbreak{} }\par
\noindent Corrected reading: {\ttfamily\footnotesize -\allowbreak{}}\par
\noindent Source: {\ttfamily\footnotesize fr/sga7i/\{diplomatic\allowbreak{},current\}/expose\_VII\allowbreak{}I\_body.tex;\allowbreak{} folio\allowbreak{} 245\allowbreak{}}\par

\par\medskip\noindent\hypertarget{app-E7I-VIII-016.1}{}\textbf{E7I-VIII-016.1}\quad printed folio 249.\par
\noindent Printed reading: {\ttfamily\footnotesize X\allowbreak{}}\par
\noindent Corrected reading: {\ttfamily\footnotesize x\allowbreak{}}\par
\noindent Source: {\ttfamily\footnotesize fr/sga7i/\{diplomatic\allowbreak{},current\}/expose\_VII\allowbreak{}I\_body.tex;\allowbreak{} folio\allowbreak{} 249;\allowbreak{} formula\allowbreak{} 735\allowbreak{}}\par

\par\medskip\noindent\hypertarget{app-E7I-VIII-017.1}{}\textbf{E7I-VIII-017.1}\quad printed folio 251.\par
\noindent Printed reading: \emph{(empty)}\par
\noindent Corrected reading: {\ttfamily\footnotesize \_\textbackslash{}xi\allowbreak{}}\par
\noindent Source: {\ttfamily\footnotesize fr/sga7i/\{diplomatic\allowbreak{},current\}/expose\_VII\allowbreak{}I\_body.tex;\allowbreak{} folio\allowbreak{} 251;\allowbreak{} formula\allowbreak{} 802\allowbreak{}}\par

\par\medskip\noindent\hypertarget{app-E7I-VIII-018.1}{}\textbf{E7I-VIII-018.1}\quad printed folio 258.\par
\noindent Printed reading: {\ttfamily\footnotesize L\allowbreak{}}\par
\noindent Corrected reading: {\ttfamily\footnotesize T\allowbreak{}}\par
\noindent Source: {\ttfamily\footnotesize fr/sga7i/\{diplomatic\allowbreak{},current\}/expose\_VII\allowbreak{}I\_body.tex;\allowbreak{} folio\allowbreak{} 258;\allowbreak{} formula\allowbreak{} 950\allowbreak{}}\par

\par\medskip\noindent\hypertarget{app-E7I-VIII-019.1}{}\textbf{E7I-VIII-019.1}\quad printed folio 260.\par
\noindent Printed reading: {\ttfamily\footnotesize P\allowbreak{}}\par
\noindent Corrected reading: {\ttfamily\footnotesize E\allowbreak{}}\par
\noindent Source: {\ttfamily\footnotesize fr/sga7i/\{diplomatic\allowbreak{},current\}/expose\_VII\allowbreak{}I\_body.tex;\allowbreak{} folio\allowbreak{} 260;\allowbreak{} formula\allowbreak{} 1024\allowbreak{}}\par

\par\medskip\noindent\hypertarget{app-E7I-VIII-020.1}{}\textbf{E7I-VIII-020.1}\quad printed folio 263.\par
\noindent Printed reading: {\ttfamily\footnotesize 1\allowbreak{}}\par
\noindent Corrected reading: {\ttfamily\footnotesize i\allowbreak{}}\par
\noindent Source: {\ttfamily\footnotesize fr/sga7i/\{diplomatic\allowbreak{},current\}/expose\_VII\allowbreak{}I\_body.tex;\allowbreak{} folio\allowbreak{} 263;\allowbreak{} formula\allowbreak{} 1138\allowbreak{}}\par

\par\medskip\noindent\hypertarget{app-E7I-VIII-021.1}{}\textbf{E7I-VIII-021.1}\quad printed folio 264.\par
\noindent Printed reading: {\ttfamily\footnotesize \textbackslash{}mathbb\allowbreak{} Z\allowbreak{}}\par
\noindent Corrected reading: {\ttfamily\footnotesize M\allowbreak{}}\par
\noindent Source: {\ttfamily\footnotesize ED-EN-FONT-G012;\allowbreak{} F7I-VIII-033.2;\allowbreak{} fr/sga7i/\{diplomatic\allowbreak{},current\}/expose\_VII\allowbreak{}I\_body.tex;\allowbreak{} folio\allowbreak{} 264\allowbreak{}}\par

\par\medskip\noindent\hypertarget{app-E7I-VIII-022.1}{}\textbf{E7I-VIII-022.1}\quad printed folio 265.\par
\noindent Printed reading: {\ttfamily\footnotesize P\allowbreak{}}\par
\noindent Corrected reading: {\ttfamily\footnotesize \textbackslash{}mathbb\allowbreak{} L\allowbreak{}}\par
\noindent Source: {\ttfamily\footnotesize ED-EN-FORMULA-G045;\allowbreak{} fr/sga7i/\{diplomatic\allowbreak{},current\}/expose\_VII\allowbreak{}I\_body.tex;\allowbreak{} folio\allowbreak{} 265\allowbreak{}}\par

\par\medskip\noindent\hypertarget{app-E7I-VIII-022.2}{}\textbf{E7I-VIII-022.2}\quad printed folio 265.\par
\noindent Printed reading: {\ttfamily\footnotesize P\allowbreak{}}\par
\noindent Corrected reading: {\ttfamily\footnotesize \textbackslash{}mathbb\allowbreak{} L\textbackslash{}otimes\allowbreak{}}\par
\noindent Source: {\ttfamily\footnotesize ED-EN-FORMULA-G045;\allowbreak{} fr/sga7i/\{diplomatic\allowbreak{},current\}/expose\_VII\allowbreak{}I\_body.tex;\allowbreak{} folio\allowbreak{} 265\allowbreak{}}\par

\par\medskip\noindent\hypertarget{app-E7I-VIII-022.3}{}\textbf{E7I-VIII-022.3}\quad printed folio 265.\par
\noindent Printed reading: {\ttfamily\footnotesize P\allowbreak{}}\par
\noindent Corrected reading: {\ttfamily\footnotesize \textbackslash{}mathbb\allowbreak{} L\allowbreak{}}\par
\noindent Source: {\ttfamily\footnotesize ED-EN-FORMULA-G045;\allowbreak{} fr/sga7i/\{diplomatic\allowbreak{},current\}/expose\_VII\allowbreak{}I\_body.tex;\allowbreak{} folio\allowbreak{} 265\allowbreak{}}\par

\par\medskip\noindent\hypertarget{app-E7I-VIII-023.1}{}\textbf{E7I-VIII-023.1}\quad printed folio 268.\par
\noindent Printed reading: {\ttfamily\footnotesize A\allowbreak{}}\par
\noindent Corrected reading: {\ttfamily\footnotesize B\allowbreak{}}\par
\noindent Source: {\ttfamily\footnotesize fr/sga7i/\{diplomatic\allowbreak{},current\}/expose\_VII\allowbreak{}I\_body.tex;\allowbreak{} folio\allowbreak{} 268;\allowbreak{} formula\allowbreak{} 1299\allowbreak{}}\par

\par\medskip\noindent\hypertarget{app-E7I-VIII-024.1}{}\textbf{E7I-VIII-024.1}\quad printed folio 269.\par
\noindent Printed reading: {\ttfamily\footnotesize A\allowbreak{}}\par
\noindent Corrected reading: {\ttfamily\footnotesize Q\allowbreak{}}\par
\noindent Source: {\ttfamily\footnotesize fr/sga7i/\{diplomatic\allowbreak{},current\}/expose\_VII\allowbreak{}I\_body.tex;\allowbreak{} folio\allowbreak{} 269;\allowbreak{} formula\allowbreak{} 1322\allowbreak{}}\par

\par\medskip\noindent\hypertarget{app-E7I-VIII-024.2}{}\textbf{E7I-VIII-024.2}\quad printed folio 269.\par
\noindent Printed reading: {\ttfamily\footnotesize 1\allowbreak{}}\par
\noindent Corrected reading: {\ttfamily\footnotesize 4\allowbreak{}}\par
\noindent Source: {\ttfamily\footnotesize port/PORT\_EN7I\_LOG.j\allowbreak{}sonl\#C798;\allowbreak{} fixes/FIXLOG\_7I\_CURR\allowbreak{}ENT.jsonl\#C798;\allowbreak{} 7I:817\allowbreak{}}\par

\par\medskip\noindent\hypertarget{app-E7I-VIII-025.1}{}\textbf{E7I-VIII-025.1}\quad printed folio 274.\par
\noindent Printed reading: {\ttfamily\footnotesize 1\allowbreak{}}\par
\noindent Corrected reading: {\ttfamily\footnotesize 4\allowbreak{}}\par
\noindent Source: {\ttfamily\footnotesize 7I:483;\allowbreak{} fixes/MISPRINTS\_7I.j\allowbreak{}son\#483\allowbreak{}}\par

\par\medskip\noindent\hypertarget{app-E7I-VIII-026.1}{}\textbf{E7I-VIII-026.1}\quad printed folio 276.\par
\noindent Printed reading: {\ttfamily\footnotesize A\allowbreak{} }\par
\noindent Corrected reading: {\ttfamily\footnotesize P\allowbreak{}}\par
\noindent Source: {\ttfamily\footnotesize fr/sga7i/\{diplomatic\allowbreak{},current\}/expose\_VII\allowbreak{}I\_body.tex;\allowbreak{} folio\allowbreak{} 276;\allowbreak{} formula\allowbreak{} 1476\allowbreak{}}\par

\par\medskip\noindent\hypertarget{app-E7I-VIII-026.2}{}\textbf{E7I-VIII-026.2}\quad printed folio 276.\par
\noindent Printed reading: {\ttfamily\footnotesize P\allowbreak{}}\par
\noindent Corrected reading: {\ttfamily\footnotesize A\allowbreak{}}\par
\noindent Source: {\ttfamily\footnotesize fr/sga7i/\{diplomatic\allowbreak{},current\}/expose\_VII\allowbreak{}I\_body.tex;\allowbreak{} folio\allowbreak{} 276;\allowbreak{} formula\allowbreak{} 1476\allowbreak{}}\par

\par\medskip\noindent\hypertarget{app-E7I-VIII-026.3}{}\textbf{E7I-VIII-026.3}\quad printed folio 276.\par
\noindent Printed reading: {\ttfamily\footnotesize \textasciicircum{}*\allowbreak{}}\par
\noindent Corrected reading: \emph{(empty)}\par
\noindent Source: {\ttfamily\footnotesize fr/sga7i/\{diplomatic\allowbreak{},current\}/expose\_VII\allowbreak{}I\_body.tex;\allowbreak{} folio\allowbreak{} 276\allowbreak{}}\par

\par\medskip\noindent\hypertarget{app-E7I-VIII-026.4}{}\textbf{E7I-VIII-026.4}\quad printed folio 276.\par
\noindent Printed reading: {\ttfamily\footnotesize \textbackslash{}mathbb\allowbreak{} M\allowbreak{}}\par
\noindent Corrected reading: {\ttfamily\footnotesize A\allowbreak{}}\par
\noindent Source: {\ttfamily\footnotesize ED-EN-FORMULA-G047;\allowbreak{} fr/sga7i/\{diplomatic\allowbreak{},current\}/expose\_VII\allowbreak{}I\_body.tex;\allowbreak{} folio\allowbreak{} 276\allowbreak{}}\par

\par\medskip\noindent\hypertarget{app-E7I-VIII-027.1}{}\textbf{E7I-VIII-027.1}\quad printed folio 278.\par
\noindent Printed reading: {\ttfamily\footnotesize 1\allowbreak{}}\par
\noindent Corrected reading: {\ttfamily\footnotesize 2\allowbreak{}}\par
\noindent Source: {\ttfamily\footnotesize ADOPT-F7I-VIII-043.1\allowbreak{};\allowbreak{} fr/sga7i/diplomatic/\allowbreak{}expose\_VIII\_body.tex\allowbreak{};\allowbreak{} fr/sga7i/current/exp\allowbreak{}ose\_VIII\_body.tex\allowbreak{}}\par

\par\medskip\noindent\hypertarget{app-E7I-VIII-028.1}{}\textbf{E7I-VIII-028.1}\quad printed folio 278.\par
\noindent Printed reading: {\ttfamily\footnotesize G\allowbreak{}}\par
\noindent Corrected reading: {\ttfamily\footnotesize Z\allowbreak{}}\par
\noindent Source: {\ttfamily\footnotesize fr/sga7i/\{diplomatic\allowbreak{},current\}/expose\_VII\allowbreak{}I\_body.tex;\allowbreak{} folio\allowbreak{} 278;\allowbreak{} formula\allowbreak{} 1548\allowbreak{}}\par

\par\medskip\noindent\hypertarget{app-E7I-VIII-028.2}{}\textbf{E7I-VIII-028.2}\quad printed folio 279.\par
\noindent Printed reading: {\ttfamily\footnotesize )\allowbreak{}}\par
\noindent Corrected reading: \emph{(empty)}\par
\noindent Source: {\ttfamily\footnotesize fr/sga7i/\{diplomatic\allowbreak{},current\}/expose\_VII\allowbreak{}I\_body.tex;\allowbreak{} folio\allowbreak{} 279;\allowbreak{} formula\allowbreak{} 1563;\allowbreak{} port/decisions/decis\allowbreak{}ions\_en7I\_r2\_002.jso\allowbreak{}n\#C804;\allowbreak{} fixes/FIXLOG\_7I\_CURR\allowbreak{}ENT.jsonl\#C804;\allowbreak{} 7I:491\allowbreak{}}\par

\par\medskip\noindent\hypertarget{app-E7I-VIII-028.3}{}\textbf{E7I-VIII-028.3}\quad printed folio 279.\par
\noindent Printed reading: {\ttfamily\footnotesize E\allowbreak{}}\par
\noindent Corrected reading: {\ttfamily\footnotesize Z\allowbreak{}}\par
\noindent Source: {\ttfamily\footnotesize fr/sga7i/\{diplomatic\allowbreak{},current\}/expose\_VII\allowbreak{}I\_body.tex;\allowbreak{} folio\allowbreak{} 279;\allowbreak{} formula\allowbreak{} 1569\allowbreak{}}\par

\par\medskip\noindent\hypertarget{app-E7I-VIII-028.4}{}\textbf{E7I-VIII-028.4}\quad printed folio 279.\par
\noindent Printed reading: {\ttfamily\footnotesize \_\textbackslash{}eta\allowbreak{}}\par
\noindent Corrected reading: \emph{(empty)}\par
\noindent Source: {\ttfamily\footnotesize fr/sga7i/\{diplomatic\allowbreak{},current\}/expose\_VII\allowbreak{}I\_body.tex;\allowbreak{} folio\allowbreak{} 279;\allowbreak{} formula\allowbreak{} 1575\allowbreak{}}\par

\par\medskip\noindent\hypertarget{app-E7I-VIII-028.5}{}\textbf{E7I-VIII-028.5}\quad printed folio 279.\par
\noindent Printed reading: {\ttfamily\footnotesize \{\allowbreak{}}\par
\noindent Corrected reading: \emph{(empty)}\par
\noindent Source: {\ttfamily\footnotesize fr/sga7i/\{diplomatic\allowbreak{},current\}/expose\_VII\allowbreak{}I\_body.tex;\allowbreak{} folio\allowbreak{} 279;\allowbreak{} formula\allowbreak{} 1576\allowbreak{}}\par

\par\medskip\noindent\hypertarget{app-E7I-VIII-028.6}{}\textbf{E7I-VIII-028.6}\quad printed folio 279.\par
\noindent Printed reading: {\ttfamily\footnotesize \_\textbackslash{}eta\}\allowbreak{}}\par
\noindent Corrected reading: \emph{(empty)}\par
\noindent Source: {\ttfamily\footnotesize fr/sga7i/\{diplomatic\allowbreak{},current\}/expose\_VII\allowbreak{}I\_body.tex;\allowbreak{} folio\allowbreak{} 279;\allowbreak{} formula\allowbreak{} 1576\allowbreak{}}\par

\par\medskip\noindent\hypertarget{app-E7I-VIII-028.7}{}\textbf{E7I-VIII-028.7}\quad printed folio 280.\par
\noindent Printed reading: \emph{(empty)}\par
\noindent Corrected reading: {\ttfamily\footnotesize \textbackslash{}eta\allowbreak{}}\par
\noindent Source: {\ttfamily\footnotesize fr/sga7i/\{diplomatic\allowbreak{},current\}/expose\_VII\allowbreak{}I\_body.tex;\allowbreak{} folio\allowbreak{} 280;\allowbreak{} formula\allowbreak{} 1591\allowbreak{}}\par

\par\medskip\noindent\hypertarget{app-E7I-VIII-029.1}{}\textbf{E7I-VIII-029.1}\quad printed folio 282.\par
\noindent Printed reading: {\ttfamily\footnotesize Z\allowbreak{}}\par
\noindent Corrected reading: {\ttfamily\footnotesize P\allowbreak{}}\par
\noindent Source: {\ttfamily\footnotesize ED-EN-FORMULA-G037;\allowbreak{} fr/sga7i/\{diplomatic\allowbreak{},current\}/expose\_VII\allowbreak{}I\_body.tex;\allowbreak{} folio\allowbreak{} 282\allowbreak{}}\par

\par\medskip\noindent\hypertarget{app-E7I-VIII-029.2}{}\textbf{E7I-VIII-029.2}\quad printed folio 282.\par
\noindent Printed reading: {\ttfamily\footnotesize G\allowbreak{}}\par
\noindent Corrected reading: {\ttfamily\footnotesize Z\allowbreak{}}\par
\noindent Source: {\ttfamily\footnotesize ED-EN-FORMULA-G037;\allowbreak{} fr/sga7i/\{diplomatic\allowbreak{},current\}/expose\_VII\allowbreak{}I\_body.tex;\allowbreak{} folio\allowbreak{} 282\allowbreak{}}\par

\par\medskip\noindent\hypertarget{app-E7I-VIII-029.3}{}\textbf{E7I-VIII-029.3}\quad printed folio 283.\par
\noindent Printed reading: {\ttfamily\footnotesize 3\allowbreak{}}\par
\noindent Corrected reading: {\ttfamily\footnotesize 2\allowbreak{}}\par
\noindent Source: {\ttfamily\footnotesize port/PORT\_EN7I\_LOG.j\allowbreak{}sonl\#C1026;\allowbreak{} fixes/FIXLOG\_7I\_CURR\allowbreak{}ENT.jsonl\#C1026;\allowbreak{} 7I:1066\allowbreak{}}\par

\par\medskip\noindent\hypertarget{app-E7I-VIII-029.4}{}\textbf{E7I-VIII-029.4}\quad printed folio 283.\par
\noindent Printed reading: {\ttfamily\footnotesize G\allowbreak{}}\par
\noindent Corrected reading: {\ttfamily\footnotesize Z\allowbreak{}}\par
\noindent Source: {\ttfamily\footnotesize fr/sga7i/\{diplomatic\allowbreak{},current\}/expose\_VII\allowbreak{}I\_body.tex;\allowbreak{} folio\allowbreak{} 283;\allowbreak{} formula\allowbreak{} 1685\allowbreak{}}\par

\par\medskip\noindent\hypertarget{app-E7I-VIII-029.5}{}\textbf{E7I-VIII-029.5}\quad printed folio 283.\par
\noindent Printed reading: {\ttfamily\footnotesize G\allowbreak{}}\par
\noindent Corrected reading: {\ttfamily\footnotesize Z\allowbreak{}}\par
\noindent Source: {\ttfamily\footnotesize fr/sga7i/\{diplomatic\allowbreak{},current\}/expose\_VII\allowbreak{}I\_body.tex;\allowbreak{} folio\allowbreak{} 283;\allowbreak{} formula\allowbreak{} 1686\allowbreak{}}\par

\par\medskip\noindent\hypertarget{app-E7I-VIII-029.6}{}\textbf{E7I-VIII-029.6}\quad printed folio 283.\par
\noindent Printed reading: {\ttfamily\footnotesize G\allowbreak{}}\par
\noindent Corrected reading: {\ttfamily\footnotesize Z\allowbreak{}}\par
\noindent Source: {\ttfamily\footnotesize fr/sga7i/\{diplomatic\allowbreak{},current\}/expose\_VII\allowbreak{}I\_body.tex;\allowbreak{} folio\allowbreak{} 283;\allowbreak{} formula\allowbreak{} 1686\allowbreak{}}\par

\par\medskip\noindent\hypertarget{app-E7I-VIII-029.7}{}\textbf{E7I-VIII-029.7}\quad printed folio 283.\par
\noindent Printed reading: {\ttfamily\footnotesize G\allowbreak{}}\par
\noindent Corrected reading: {\ttfamily\footnotesize Z\allowbreak{}}\par
\noindent Source: {\ttfamily\footnotesize fr/sga7i/\{diplomatic\allowbreak{},current\}/expose\_VII\allowbreak{}I\_body.tex;\allowbreak{} folio\allowbreak{} 283;\allowbreak{} formula\allowbreak{} 1688\allowbreak{}}\par

\par\medskip\noindent\hypertarget{app-E7I-VIII-030.1}{}\textbf{E7I-VIII-030.1}\quad printed folio 283.\par
\noindent Printed reading: {\ttfamily\footnotesize neutral\allowbreak{}}\par
\noindent Corrected reading: {\ttfamily\footnotesize connected\allowbreak{}}\par
\noindent Source: {\ttfamily\footnotesize 7I:502;\allowbreak{} 7I:1072;\allowbreak{} fixes/MISPRINTS\_7I.j\allowbreak{}son\#502;\allowbreak{} fixes/MISPRINTS\_7I.j\allowbreak{}son\#1072\allowbreak{}}\par

\par\medskip\noindent\hypertarget{app-E7I-VIII-031.1}{}\textbf{E7I-VIII-031.1}\quad printed folio 284.\par
\noindent Printed reading: {\ttfamily\footnotesize G\allowbreak{}}\par
\noindent Corrected reading: {\ttfamily\footnotesize Z\allowbreak{}}\par
\noindent Source: {\ttfamily\footnotesize fr/sga7i/\{diplomatic\allowbreak{},current\}/expose\_VII\allowbreak{}I\_body.tex;\allowbreak{} folio\allowbreak{} 284;\allowbreak{} formula\allowbreak{} 1717\allowbreak{}}\par

\par\medskip\noindent\hypertarget{app-E7I-VIII-031.2}{}\textbf{E7I-VIII-031.2}\quad printed folio 285.\par
\noindent Printed reading: {\ttfamily\footnotesize 2\allowbreak{}}\par
\noindent Corrected reading: {\ttfamily\footnotesize 5\allowbreak{}}\par
\noindent Source: {\ttfamily\footnotesize 7I:1074;\allowbreak{} fixes/MISPRINTS\_7I.j\allowbreak{}son\#1074\allowbreak{}}\par

\par\medskip\noindent\hypertarget{app-E7I-VIII-032.1}{}\textbf{E7I-VIII-032.1}\quad printed folio 285.\par
\noindent Printed reading: {\ttfamily\footnotesize G\allowbreak{}}\par
\noindent Corrected reading: {\ttfamily\footnotesize Z\allowbreak{}}\par
\noindent Source: {\ttfamily\footnotesize fr/sga7i/\{diplomatic\allowbreak{},current\}/expose\_VII\allowbreak{}I\_body.tex;\allowbreak{} folio\allowbreak{} 285;\allowbreak{} formula\allowbreak{} 1758\allowbreak{}}\par

\par\medskip\noindent\hypertarget{app-E7I-VIII-032.2}{}\textbf{E7I-VIII-032.2}\quad printed folio 286.\par
\noindent Printed reading: {\ttfamily\footnotesize G\allowbreak{}}\par
\noindent Corrected reading: {\ttfamily\footnotesize Z\allowbreak{}}\par
\noindent Source: {\ttfamily\footnotesize fr/sga7i/\{diplomatic\allowbreak{},current\}/expose\_VII\allowbreak{}I\_body.tex;\allowbreak{} folio\allowbreak{} 286;\allowbreak{} formula\allowbreak{} 1762\allowbreak{}}\par

\par\medskip\noindent\hypertarget{app-E7I-VIII-033.1}{}\textbf{E7I-VIII-033.1}\quad printed folio 286.\par
\noindent Printed reading: {\ttfamily\footnotesize G\allowbreak{}}\par
\noindent Corrected reading: {\ttfamily\footnotesize Z\allowbreak{}}\par
\noindent Source: {\ttfamily\footnotesize fr/sga7i/\{diplomatic\allowbreak{},current\}/expose\_VII\allowbreak{}I\_body.tex;\allowbreak{} folio\allowbreak{} 286;\allowbreak{} formula\allowbreak{} 1784\allowbreak{}}\par

\par\medskip\noindent\hypertarget{app-E7I-VIII-033.2}{}\textbf{E7I-VIII-033.2}\quad printed folio 286.\par
\noindent Printed reading: {\ttfamily\footnotesize G\allowbreak{}}\par
\noindent Corrected reading: {\ttfamily\footnotesize Z\allowbreak{}}\par
\noindent Source: {\ttfamily\footnotesize fr/sga7i/\{diplomatic\allowbreak{},current\}/expose\_VII\allowbreak{}I\_body.tex;\allowbreak{} folio\allowbreak{} 286;\allowbreak{} formula\allowbreak{} 1785\allowbreak{}}\par

\par\medskip\noindent\hypertarget{app-E7I-VIII-034.1}{}\textbf{E7I-VIII-034.1}\quad printed folio 288.\par
\noindent Printed reading: {\ttfamily\footnotesize Z\allowbreak{}}\par
\noindent Corrected reading: {\ttfamily\footnotesize S\allowbreak{}}\par
\noindent Source: {\ttfamily\footnotesize fr/sga7i/\{diplomatic\allowbreak{},current\}/expose\_VII\allowbreak{}I\_body.tex;\allowbreak{} folio\allowbreak{} 288;\allowbreak{} formula\allowbreak{} 1844\allowbreak{}}\par

\par\medskip\noindent\hypertarget{app-E7I-VIII-034.2}{}\textbf{E7I-VIII-034.2}\quad printed folio 288.\par
\noindent Printed reading: {\ttfamily\footnotesize \{\}\_\{\textbackslash{}overline\allowbreak{} G\}\allowbreak{}}\par
\noindent Corrected reading: \emph{(empty)}\par
\noindent Source: {\ttfamily\footnotesize fr/sga7i/\{diplomatic\allowbreak{},current\}/expose\_VII\allowbreak{}I\_body.tex;\allowbreak{} folio\allowbreak{} 288\allowbreak{}}\par

\par\medskip\noindent\hypertarget{app-E7I-VIII-035.1}{}\textbf{E7I-VIII-035.1}\quad printed folio 289.\par
\noindent Printed reading: \emph{(empty)}\par
\noindent Corrected reading: {\ttfamily\footnotesize \textbackslash{}overline\allowbreak{}}\par
\noindent Source: {\ttfamily\footnotesize fr/sga7i/\{diplomatic\allowbreak{},current\}/expose\_VII\allowbreak{}I\_body.tex;\allowbreak{} folio\allowbreak{} 289;\allowbreak{} formula\allowbreak{} 1913;\allowbreak{} port/decisions/decis\allowbreak{}ions\_en7I\_r2\_002.jso\allowbreak{}n\#C807;\allowbreak{} fixes/FIXLOG\_7I\_CURR\allowbreak{}ENT.jsonl\#C807;\allowbreak{} 7I:508\allowbreak{}}\par

\par\medskip\noindent\hypertarget{app-E7I-VIII-035.2}{}\textbf{E7I-VIII-035.2}\quad printed folio 289.\par
\noindent Printed reading: \emph{(empty)}\par
\noindent Corrected reading: {\ttfamily\footnotesize \textbackslash{}overline\allowbreak{}}\par
\noindent Source: {\ttfamily\footnotesize fr/sga7i/\{diplomatic\allowbreak{},current\}/expose\_VII\allowbreak{}I\_body.tex;\allowbreak{} folio\allowbreak{} 289;\allowbreak{} formula\allowbreak{} 1914;\allowbreak{} port/decisions/decis\allowbreak{}ions\_en7I\_r2\_002.jso\allowbreak{}n\#C807;\allowbreak{} fixes/FIXLOG\_7I\_CURR\allowbreak{}ENT.jsonl\#C807;\allowbreak{} 7I:508\allowbreak{}}\par

\par\medskip\noindent\hypertarget{app-E7I-VIII-035.3}{}\textbf{E7I-VIII-035.3}\quad printed folio 289.\par
\noindent Printed reading: {\ttfamily\footnotesize 3\allowbreak{}}\par
\noindent Corrected reading: {\ttfamily\footnotesize 6\allowbreak{}}\par
\noindent Source: {\ttfamily\footnotesize 7I:1061;\allowbreak{} fixes/MISPRINTS\_7I.j\allowbreak{}son\#1061\allowbreak{}}\par

\par\medskip\noindent\hypertarget{app-E7I-VIII-036.1}{}\textbf{E7I-VIII-036.1}\quad printed folio 291.\par
\noindent Printed reading: {\ttfamily\footnotesize \textbackslash{}underline\allowbreak{}}\par
\noindent Corrected reading: {\ttfamily\footnotesize \textbackslash{}overline\allowbreak{}}\par
\noindent Source: {\ttfamily\footnotesize fr/sga7i/\{diplomatic\allowbreak{},current\}/expose\_VII\allowbreak{}I\_body.tex;\allowbreak{} folio\allowbreak{} 291;\allowbreak{} formula\allowbreak{} 1973\allowbreak{}}\par

\par\medskip\noindent\hypertarget{app-E7I-VIII-036.2}{}\textbf{E7I-VIII-036.2}\quad printed folio 291.\par
\noindent Printed reading: \emph{(empty)}\par
\noindent Corrected reading: {\ttfamily\footnotesize -\allowbreak{}}\par
\noindent Source: {\ttfamily\footnotesize fr/sga7i/\{diplomatic\allowbreak{},current\}/expose\_VII\allowbreak{}I\_body.tex;\allowbreak{} folio\allowbreak{} 291;\allowbreak{} formula\allowbreak{} 1998\allowbreak{}}\par

\par\medskip\noindent\hypertarget{app-E7I-VIII-036.3}{}\textbf{E7I-VIII-036.3}\quad printed folio 291.\par
\noindent Printed reading: \emph{(empty)}\par
\noindent Corrected reading: {\ttfamily\footnotesize \textbackslash{}overline\allowbreak{}}\par
\noindent Source: {\ttfamily\footnotesize fr/sga7i/\{diplomatic\allowbreak{},current\}/expose\_VII\allowbreak{}I\_body.tex;\allowbreak{} folio\allowbreak{} 291;\allowbreak{} formula\allowbreak{} 2007\allowbreak{}}\par

\par\medskip\noindent\hypertarget{app-E7I-VIII-037.1}{}\textbf{E7I-VIII-037.1}\quad printed folio 294.\par
\noindent Printed reading: {\ttfamily\footnotesize X\allowbreak{}}\par
\noindent Corrected reading: {\ttfamily\footnotesize S\allowbreak{}}\par
\noindent Source: {\ttfamily\footnotesize fr/sga7i/\{diplomatic\allowbreak{},current\}/expose\_VII\allowbreak{}I\_body.tex;\allowbreak{} folio\allowbreak{} 294;\allowbreak{} formula\allowbreak{} 2104\allowbreak{}}\par

\par\medskip\noindent\hypertarget{app-E7I-VIII-038.1}{}\textbf{E7I-VIII-038.1}\quad printed folio 297.\par
\noindent Printed reading: {\ttfamily\footnotesize V\allowbreak{}}\par
\noindent Corrected reading: {\ttfamily\footnotesize U\allowbreak{}}\par
\noindent Source: {\ttfamily\footnotesize fr/sga7i/\{diplomatic\allowbreak{},current\}/expose\_VII\allowbreak{}I\_body.tex;\allowbreak{} folio\allowbreak{} 297;\allowbreak{} formula\allowbreak{} 2188\allowbreak{}}\par

\par\medskip\noindent\hypertarget{app-E7I-VIII-038.2}{}\textbf{E7I-VIII-038.2}\quad printed folio 297.\par
\noindent Printed reading: {\ttfamily\footnotesize mS\allowbreak{}}\par
\noindent Corrected reading: {\ttfamily\footnotesize mU\allowbreak{}}\par
\noindent Source: {\ttfamily\footnotesize fr/sga7i/\{diplomatic\allowbreak{},current\}/expose\_VII\allowbreak{}I\_body.tex;\allowbreak{} folio\allowbreak{} 297;\allowbreak{} formula\allowbreak{} 2193\allowbreak{}}\par

\par\medskip\noindent\hypertarget{app-E7I-VIII-039.1}{}\textbf{E7I-VIII-039.1}\quad printed folio 298.\par
\noindent Printed reading: \emph{(empty)}\par
\noindent Corrected reading: {\ttfamily\footnotesize X\$,\allowbreak{} with\allowbreak{} closed\allowbreak{} immersion\allowbreak{} \$i:Y\textbackslash{}hookrightarrow\allowbreak{} }\par
\noindent Source: {\ttfamily\footnotesize fr/sga7i/\{diplomatic\allowbreak{},current\}/expose\_VII\allowbreak{}I\_body.tex;\allowbreak{} folio\allowbreak{} 298\allowbreak{}}\par

\par\medskip\noindent\hypertarget{app-E7I-VIII-040.1}{}\textbf{E7I-VIII-040.1}\quad printed folio 300.\par
\noindent Printed reading: \emph{(empty)}\par
\noindent Corrected reading: {\ttfamily\footnotesize \_\textbackslash{}eta\allowbreak{}}\par
\noindent Source: {\ttfamily\footnotesize fr/sga7i/\{diplomatic\allowbreak{},current\}/expose\_VII\allowbreak{}I\_body.tex;\allowbreak{} folio\allowbreak{} 300;\allowbreak{} formula\allowbreak{} 2271\allowbreak{}}\par

\par\medskip\noindent\hypertarget{app-E7I-VIII-040.2}{}\textbf{E7I-VIII-040.2}\quad printed folio 300.\par
\noindent Printed reading: {\ttfamily\footnotesize biextensions\allowbreak{}}\par
\noindent Corrected reading: {\ttfamily\footnotesize extensions\allowbreak{}}\par
\noindent Source: {\ttfamily\footnotesize 7I:524;\allowbreak{} 7I:1094;\allowbreak{} fixes/MISPRINTS\_7I.j\allowbreak{}son\#524;\allowbreak{} fixes/MISPRINTS\_7I.j\allowbreak{}son\#1094\allowbreak{}}\par

\par\medskip\noindent\hypertarget{app-E7I-VIII-040.3}{}\textbf{E7I-VIII-040.3}\quad printed folio 302.\par
\noindent Printed reading: \emph{(empty)}\par
\noindent Corrected reading: {\ttfamily\footnotesize (\textbackslash{}mathbb\allowbreak{} }\par
\noindent Source: {\ttfamily\footnotesize ED-EN-FONT-G013;\allowbreak{} F7I-VIII-067.2;\allowbreak{} fr/sga7i/\{diplomatic\allowbreak{},current\}/expose\_VII\allowbreak{}I\_body.tex;\allowbreak{} folio\allowbreak{} 302\allowbreak{}}\par

\par\medskip\noindent\hypertarget{app-E7I-VIII-040.4}{}\textbf{E7I-VIII-040.4}\quad printed folio 302.\par
\noindent Printed reading: \emph{(empty)}\par
\noindent Corrected reading: {\ttfamily\footnotesize )\_k\allowbreak{}}\par
\noindent Source: {\ttfamily\footnotesize ED-EN-FONT-G013;\allowbreak{} F7I-VIII-067.2;\allowbreak{} fr/sga7i/\{diplomatic\allowbreak{},current\}/expose\_VII\allowbreak{}I\_body.tex;\allowbreak{} folio\allowbreak{} 302\allowbreak{}}\par

\par\medskip\noindent\hypertarget{app-E7I-VIII-041.1}{}\textbf{E7I-VIII-041.1}\quad printed folio 303.\par
\noindent Printed reading: \emph{(empty)}\par
\noindent Corrected reading: {\ttfamily\footnotesize \textbackslash{}eta\allowbreak{}}\par
\noindent Source: {\ttfamily\footnotesize fr/sga7i/\{diplomatic\allowbreak{},current\}/expose\_VII\allowbreak{}I\_body.tex;\allowbreak{} folio\allowbreak{} 303;\allowbreak{} formula\allowbreak{} 2373\allowbreak{}}\par

\par\medskip\noindent\hypertarget{app-E7I-VIII-042.1}{}\textbf{E7I-VIII-042.1}\quad printed folio 305.\par
\noindent Printed reading: \emph{(empty)}\par
\noindent Corrected reading: {\ttfamily\footnotesize \_\textbackslash{}eta\allowbreak{}}\par
\noindent Source: {\ttfamily\footnotesize fr/sga7i/\{diplomatic\allowbreak{},current\}/expose\_VII\allowbreak{}I\_body.tex;\allowbreak{} folio\allowbreak{} 305;\allowbreak{} formula\allowbreak{} 2406\allowbreak{}}\par

\par\medskip\noindent\hypertarget{app-E7I-VIII-042.2}{}\textbf{E7I-VIII-042.2}\quad printed folio 305.\par
\noindent Printed reading: \emph{(empty)}\par
\noindent Corrected reading: {\ttfamily\footnotesize \_\textbackslash{}eta\allowbreak{}}\par
\noindent Source: {\ttfamily\footnotesize fr/sga7i/\{diplomatic\allowbreak{},current\}/expose\_VII\allowbreak{}I\_body.tex;\allowbreak{} folio\allowbreak{} 305;\allowbreak{} formula\allowbreak{} 2411\allowbreak{}}\par

\par\medskip\noindent\hypertarget{app-E7I-VIII-042.3}{}\textbf{E7I-VIII-042.3}\quad printed folio 305.\par
\noindent Printed reading: \emph{(empty)}\par
\noindent Corrected reading: {\ttfamily\footnotesize (\allowbreak{}}\par
\noindent Source: {\ttfamily\footnotesize ED-EN-FORMULA-G049;\allowbreak{} fr/sga7i/\{diplomatic\allowbreak{},current\}/expose\_VII\allowbreak{}I\_body.tex;\allowbreak{} folio\allowbreak{} 305;\allowbreak{} port/decisions/decis\allowbreak{}ions\_en7I\_r2\_003.jso\allowbreak{}n\#C817;\allowbreak{} fixes/FIXLOG\_7I\_CURR\allowbreak{}ENT.jsonl\#C817;\allowbreak{} 7I:529\allowbreak{}}\par

\par\medskip\noindent\hypertarget{app-E7I-VIII-043.1}{}\textbf{E7I-VIII-043.1}\quad printed folio 307.\par
\noindent Printed reading: \emph{(empty)}\par
\noindent Corrected reading: {\ttfamily\footnotesize \_\textbackslash{}eta\textbackslash{}times\allowbreak{}}\par
\noindent Source: {\ttfamily\footnotesize fr/sga7i/\{diplomatic\allowbreak{},current\}/expose\_VII\allowbreak{}I\_body.tex;\allowbreak{} folio\allowbreak{} 307;\allowbreak{} formula\allowbreak{} 2489\allowbreak{}}\par

\par\medskip\noindent\hypertarget{app-E7I-VIII-043.2}{}\textbf{E7I-VIII-043.2}\quad printed folio 307.\par
\noindent Printed reading: {\ttfamily\footnotesize *\allowbreak{}}\par
\noindent Corrected reading: {\ttfamily\footnotesize \textbackslash{}eta\allowbreak{}}\par
\noindent Source: {\ttfamily\footnotesize fr/sga7i/\{diplomatic\allowbreak{},current\}/expose\_VII\allowbreak{}I\_body.tex;\allowbreak{} folio\allowbreak{} 307\allowbreak{}}\par

\par\medskip\noindent\hypertarget{app-E7I-VIII-044.1}{}\textbf{E7I-VIII-044.1}\quad printed folio 308.\par
\noindent Printed reading: \emph{(empty)}\par
\noindent Corrected reading: {\ttfamily\footnotesize '\allowbreak{}}\par
\noindent Source: {\ttfamily\footnotesize ED-EN-FORMULA-G051;\allowbreak{} fr/sga7i/\{diplomatic\allowbreak{},current\}/expose\_VII\allowbreak{}I\_body.tex;\allowbreak{} folio\allowbreak{} 308;\allowbreak{} port/decisions/decis\allowbreak{}ions\_en7I\_r4\_001.jso\allowbreak{}n\#C1031;\allowbreak{} fixes/FIXLOG\_7I\_CURR\allowbreak{}ENT.jsonl\#C1031;\allowbreak{} 7I:1098\allowbreak{}}\par

\par\endgroup\normalsize\fi
\clearpage\SGAtoc{Exposé IX}
\selectlanguage{english}
\SGApage{I}{313}

\begin{center}
\SGARefAliasTarget{sga7.sga7i.eix.nroot}SGA 7

\vspace{1em}
\underline{Expos\'e IX}

\vspace{1em}
\underline{\underline{N\'ERON MODELS AND MONODROMY}}

\vspace{1em}
\underline{by A. Grothendieck}

(with an appendix by M. Raynaud)
\end{center}

\subsection*{Contents}

\begin{enumerate}[label=\arabic*.]
\setcounter{enumi}{-1}
\item Introduction
\item N\'eron model of abelian schemes: notation; the canonical pairing
  $\Phi_o\times\Phi'_o\longrightarrow(\mathbb Q/\mathbb Z)_k$, and the
  canonical biextension $W^o$ of $(A^o,{A'}^o)$
\item Fixed part and toric part of $T_\ell(A_K)$. Criteria for good
  reduction. Orthogonality theorem for $\ell\ne p$
\item The semistable-reduction case. The semistable reduction theorem
\item Application to a conjecture of Serre--Tate and to the conductor
\item The orthogonality theorem for $\ell=p$ (semistable case). Duality
  of the abelian schemes $B_o$ and $B'_o$. Characterization of the fixed
  part. Criteria for good reduction, essentially good reduction, and
  semistable reduction
\item Remarks on the construction of the toric part and the fixed part
  of $T_p(A)$ in the non-semistable-reduction case
\item The Raynaud extension $G^\natural$ over $S$ attached to the
  abelian scheme $A_K$ with semistable reduction. Duality of the abelian
  schemes $B,B'$ over $S$
\item The Raynaud extension $A_K^{\natural\,o}$ in the semistable case,
  and the ind-group $A_K^\natural$
\item Definition of the monodromy pairing
  $\underline M_\ell\otimes\underline M'_\ell\longrightarrow
  \mathbb Z_{\ell S}$
\item Properties of the monodromy pairing. Integrality and positivity
  theorem
\item Connected components of the N\'eron model: duality relation,
  asymptotic behavior
\item N\'eron model of a Jacobian and the Picard--Lefschetz formula
\item Relations with transcendental theory: the complex-analytic case
  \begin{enumerate}[label=\arabic{enumi}.\arabic*.]
  \item Review of analytic groups over $S$
  \item Formal completion along a fiber
  \item The analytic Raynaud extension
  \item Interpretation in terms of Deligne's ``mixed Hodge structures''
  \item The relatively algebraic case
  \end{enumerate}
\item Relations with transcendental theory: the rigid-analytic case.
  The canonical homomorphism
  $\underline M_K\longrightarrow A_K^{\natural\,o}$
\end{enumerate}




\SGApage{I}{\(\approx\)314}
0. \underline{Introduction}

\SGARefTarget{sga7.sga7i.eix.n0-1}{0.1}. In the present expos\'e we apply the results of the preceding
expos\'es to the study of the \underline{N\'eron model} (\SGARefLink{sga7.sga7i.eix.n1-1}{1.1}) $A$ of an
abelian scheme $A_K$ defined over the function field $K$ of the
Henselian trait $S$. We shall see that the most important properties of
the N\'eron model can be interpreted in terms of the action of the local
monodromy group $I$ on the Tate module
$T_\ell(A_K(\overline K))$ (where $\ell$ is a prime number distinct
from $p$). Since the dual of the latter can be interpreted as the
$\ell$-adic cohomology group $H^1(A_{\overline K},\mathbb Z_\ell)$,
and since, for every smooth projective scheme $X$ over $K$,
$H^1(X_{\overline K},\mathbb Z_\ell)$ is identified with
$H^1(A_{\overline K},\mathbb Z_\ell)$, where $A_{\overline K}$ is the
Albanese variety of $X_{\overline K}$, the present expos\'e may also be
viewed as a detailed study of \underline{local monodromy phenomena for}
$\ell$-\underline{adic} $H^1$ (or, better still, for ``motivic'' $H^1$)
\underline{of smooth projective varieties over} $K$. From this point of
view, it seems clear that the principal results of the present expos\'e
are destined to be subsumed within a ``N\'eron theory'' for motives of
arbitrary weight, that is, for $H^i$ ($\ell$-adic, de Rham, Hodge, and
so forth) with arbitrary $i$, of which at present we are only beginning
to catch a glimpse. (See [9] in this connection, and more particularly
Deligne's Conjectures 9.8 through 9.13 in the cited report.)

\SGApage{I}{315}
\SGARefTarget{sga7.sga7i.eix.n0-2}{0.2}. The most important result of the present expos\'e is the
``\underline{semistable reduction theorem}'' \SGARefLink{sga7.sga7i.eix.n3-6}{3.6}, which asserts that,
after a finite separable extension $K'$ of $K$, the N\'eron model of
$A_{K'}$ has a special fiber whose identity component is an extension
of an abelian scheme by a torus. This result was conjectured by
J.-P. Serre in 1964, and the principle of the proof given here
\footnote{It is developed in a letter from the lecturer to J.-P. Serre
dated 30 October 1964.}
goes back to the same year. It appears in \SGARefLink{sga7.sga7i.eix.n3-5}{3.5} that this property of
``semistable reduction'' is equivalent to the action of the monodromy
group $I$ on $T_\ell(A_K(\overline K))$ being ``essentially unipotent
of step 2'' (where $\ell$ is a fixed prime distinct from the residual
characteristic $p$ of $S$). In this form it was proved, in terms of
$\ell$-adic $H^1$, in Expos\'e III as a special case of the
``\underline{monodromy theorem}'', using resolution of singularities
for excellent schemes of dimension $2$ (due to S. Abhyankar). To prove
the equivalence of the two properties, we shall also need an
``orthogonality theorem'' \SGARefLink{sga7.sga7i.eix.n2-4}{2.4}, due to Igusa [11] in the special case in
which $A_K$ is the Jacobian of a smooth curve $X_K$ over $K$, the
generic fiber of a regular connected projective scheme $X$ over $S$
whose geometric special fiber has no singular points other than
ordinary quadratic singularities. This is indeed a semistable-reduction
case, but it would not suffice for our purposes; it is chiefly in order
to establish the orthogonality theorem in its general form that we use
the convenient tool supplied by the theory of biextensions developed
in the two preceding expos\'es. Let us also point out that D. Mumford
has found a very different proof of the semistable reduction theorem,
one which does not use resolution of singularities and instead uses
his theory of theta functions (unpublished); unfortunately, he must
restrict himself to residual characteristic different from $2$.

\SGApage{I}{316}
\SGARefTarget{sga7.sga7i.eix.n0-3}{0.3}. The semistable reduction theorem is the principal tool used in the
subsequent sections to study N\'eron models, which we undertake here
from the \underline{``algebraic'' point of view}. It also seems
indispensable for studying abelian varieties over $K$ from the
\underline{``rigid-analytic'' point of view}, or from the point of view
of ``theta functions'' (these terms apparently being synonymous when
abelian varieties over complete valued fields are in question), as
developed recently by D. Mumford and M. Raynaud following work of Tate
and Morikawa. (In Sections 13 and 14 we shall say a few words about the
relations between the algebraic and transcendental points of view.)
Finally, the same theorem is a key ingredient in the proof of the
``\underline{semistable reduction theorem for algebraic curves}'' of
Artin--Deligne--Mumford [3],\footnote{P. Deligne gave the proof of this
theorem in the oral Seminar, but it was not considered useful to
reproduce it in the Seminar texts.} which we have already had occasion
to use in the study of the action of monodromy on the fundamental group
(the ``essentially tame action theorem'' of Expos\'e V), and which is
used again in the proof of the integrality and positivity theorem \SGARefLink{sga7.sga7i.eix.n10-4}{10.4}
in the present expos\'e.\footnote{There is another proof of the latter
theorem, by way of rigid-analytic theory, which does not use the
semistable reduction theorem.}

\SGARefTarget{sga7.sga7i.eix.n0-4}{0.4}. The present expos\'e differs from the other expos\'es of the
Seminar in that, in studying the Tate module $T_\ell(A_K)$, we have
systematically taken into account also the case in which $\ell$ equals
the residual characteristic $p$. It turns out that the language of
\underline{Barsotti--Tate groups} [33]\footnote{Or ``$p$-divisible
groups'' in Tate's terminology.} makes it possible to extend to this
case all the most important phenomena established for $\ell\ne p$,
subject on occasion to having available Tate's fundamental theorem
[33, Theorem 4], which at present has been established only when $K$
has characteristic zero. Of course, in the latter case the pro-$p$
Barsotti--Tate group $T_p(A_K)$ may still be interpreted in terms of
the action of the monodromy group $I$ on the $p$-adic module
$T_p(A_K(\overline K))$, but one should bear in mind that most of the
statements suggested by this superficial analogy with the case
$\ell\ne p$ are grossly false. (Thus it is no longer true that the
monodromy action is essentially unipotent when $\ell=p$.) By using more
judicious translations, however (such as \SGARefLink{sga7.sga7i.eix.n5-13}{5.13} below), it is already
clear that the study of pro-$p$ Barsotti--Tate groups defined
geometrically over the function field of a trait of residual
characteristic $p>0$, whether or not the trait has mixed
characteristic, must be regarded as part of the theory of ``local
monodromy'', in which these groups play the role of \underline{$p$-adic
local systems} on $\operatorname{Spec}(K)$. In fact, these are
essentially the ``local systems'' that enter into the \underline{study,
from the ``$p$-adic'' point of view, of $H^1$ of smooth projective
varieties over} $K$. For the still-to-be-developed study of higher
$H^i$ from the $p$-adic point of view, it appears that the category of
pro-$p$ Barsotti--Tate groups over $K$, respectively over $S$, will
have to be suitably enlarged, taking inspiration from the viewpoints
provided by ``crystalline cohomology''.

\SGApage{I}{318}
\SGARefTarget{sga7.sga7i.eix.n0-5}{0.5}. The expos\'e given here follows the oral lectures very closely up
to Section 4, but is markedly more detailed and complete from the
following sections onward, which had been only partially sketched. It
will not be used in the subsequent expos\'es. Some results of the
present expos\'e (\SGARefLink{sga7.sga7i.eix.n10-4}{10.4} and \SGARefLink{sga7.sga7i.eix.n12-5}{12.5}), on the other hand, use the
\underline{Picard--Lefschetz formula}, which will be established only
later.

\begin{center}
\underline{LEITFADEN}

\vspace{0.5em}
\begin{tikzpicture}[line width=0.35pt,every node/.style={fill=white,inner sep=1.2pt}]
  \coordinate (p1)  at ( 0.00, 0.00);
  \coordinate (p2)  at ( 0.00,-0.76);
  \coordinate (p3)  at ( 0.00,-1.53);
  \coordinate (p4)  at (-1.60,-1.92);
  \coordinate (p5)  at ( 0.00,-2.30);
  \coordinate (p6)  at (-1.62,-2.69);
  \coordinate (p7)  at ( 1.32,-2.87);
  \coordinate (p8)  at ( 0.04,-3.35);
  \coordinate (p9)  at ( 3.02,-2.69);
  \coordinate (p10) at ( 3.10,-3.35);
  \coordinate (p11) at ( 2.66,-4.21);
  \coordinate (p12) at ( 3.76,-4.21);
  \coordinate (p13) at ( 0.12,-4.37);
  \coordinate (p14) at ( 1.04,-4.37);
  \draw (p1)--(p2) (p2)--(p3) (p3)--(p4) (p3)--(p5)
        (p5)--(p6) (p5)--(p7) (p5)--(p8) (p5)--(p9)
        (p7)--(p8) (p7)--(p13) (p7)--(p14)
        (p8)--(p13) (p8)--(p14) (p9)--(p10)
        (p10)--(p11) (p10)--(p12) (p10)--(p13) (p10)--(p14);
  \node at (p1)  {1};  \node at (p2)  {2};
  \node at (p3)  {3};  \node at (p4)  {4};
  \node at (p5)  {5};  \node at (p6)  {6};
  \node at (p7)  {7};  \node at (p8)  {8};
  \node at (p9)  {9};  \node at (p10) {10};
  \node at (p11) {11}; \node at (p12) {12};
  \node at (p13) {13}; \node at (p14) {14};
\end{tikzpicture}
\end{center}




\medskip
\noindent
\underline{1. Neron models of abelian schemes: notation; the canonical pairing}
\[
 \Phi_o\times\Phi'_o\longrightarrow(\mathbb Q/\mathbb Z)_k,
\]
\underline{and the canonical biextension $W^o$ of $(A^o,{A'}^o)$.}

\SGApage{I}{319}
\SGAvar{\medskip
\noindent
\SGARefTarget{sga7.sga7i.eix.n1-0}{1.0}. Throughout this expos\'e we fix a trait $S$, for which we use the
usual notation $s$, $\eta$, $k$, $X$, and so forth, of I~1.  Let $A_K$
be an abelian scheme over $K$, and denote its dual abelian scheme [5] by
$A'_K$.  The duality between $A_K$ and $A'_K$ is given by a divisorial
correspondence $W_K$ on $(A_K,A'_K)$, that is, an invertible sheaf on
$A_K\times A'_K$, birigidified with respect to the identity sections of
$A_K$ and $A'_K$.  It is determined up to unique isomorphism by its
class (cf. VIII~\SGARefLink{sga7.sga7i.eviii.n3-1}{3.1})
\begin{equation}\tag{\SGARefTarget{sga7.sga7i.eix.n1-0-1}{1.0.1}}
 w_K\in\mathrm{Biext}^1(A_K,A'_K;\mathbb G_{mK}).
\end{equation}
An important role will be played by the corresponding pairing
(VIII~\SGARefLink{sga7.sga7i.eviii.n2-2}{2.2})
\begin{equation}\tag{\SGARefTarget{sga7.sga7i.eix.n1-0-2}{1.0.2}}
 \varphi:T_\ell(A_K)\times T_\ell(A'_K)
 \longrightarrow\mathbb Z_\ell(1)_K
 \overset{\mathrm{def}}{=}T_\ell(\mathbb G_{mK}),
\end{equation}
where $\ell$ is an arbitrary prime number.  When
$\ell\ne p=\operatorname{char}k$, giving $\varphi$ is equivalent to
giving a pairing of finite free $\mathbb Z_\ell$-modules
\begin{equation}\tag{\SGARefTarget{sga7.sga7i.eix.n1-0-3}{1.0.3}}
 \varphi:T_\ell(A_K(\overline K))\times
 T_\ell(A'_K(\overline K))\longrightarrow T_\ell(\overline K^*),
\end{equation}
compatible with the action of
\[
 \pi=\operatorname{Gal}(\overline K/K)
\]
on the three modules in question, where $\overline K$ is a separable
closure of $K$.  It is known (VIII~\SGARefLink{sga7.sga7i.eviii.n3-2}{3.2}) that this pairing is a perfect
duality: it identifies each of the modules
$T_\ell(A_K(\overline K))$ and $T_\ell(A'_K(\overline K))$ with the
dual of the other, with values in the invertible module
$\mathbb Z_\ell(1)=T_\ell(\overline K^*)$.

}{\medskip
\noindent
\SGARefTarget{sga7.sga7i.eix.n1-0}{1.0}. Throughout this expos\'e we fix a trait $S$, for which we use the
usual notation $s$, $\eta$, $k$, $K$, and so forth, of I~1.  Let $A_K$
be an abelian scheme over $K$, and denote its dual abelian scheme [5] by
$A'_K$.  The duality between $A_K$ and $A'_K$ is given by a divisorial
correspondence $W_K$ on $(A_K,A'_K)$, that is, an invertible sheaf on
$A_K\times A'_K$, birigidified with respect to the identity sections of
$A_K$ and $A'_K$.  It is determined up to unique isomorphism by its
class (cf. VIII~\SGARefLink{sga7.sga7i.eviii.n3-1}{3.1})
\begin{equation}\tag{\SGARefTarget{sga7.sga7i.eix.n1-0-1}{1.0.1}}
 w_K\in\mathrm{Biext}^1(A_K,A'_K;\mathbb G_{mK}).
\end{equation}
An important role will be played by the corresponding pairing
(VIII~\SGARefLink{sga7.sga7i.eviii.n2-2}{2.2})
\begin{equation}\tag{\SGARefTarget{sga7.sga7i.eix.n1-0-2}{1.0.2}}
 \varphi:T_\ell(A_K)\times T_\ell(A'_K)
 \longrightarrow\mathbb Z_\ell(1)_K
 \overset{\mathrm{def}}{=}T_\ell(\mathbb G_{mK}),
\end{equation}
where $\ell$ is an arbitrary prime number.  When
$\ell\ne p=\operatorname{char}k$, giving $\varphi$ is equivalent to
giving a pairing of finite free $\mathbb Z_\ell$-modules
\begin{equation}\tag{\SGARefTarget{sga7.sga7i.eix.n1-0-3}{1.0.3}}
 \varphi:T_\ell(A_K(\overline K))\times
 T_\ell(A'_K(\overline K))\longrightarrow T_\ell(\overline K^*),
\end{equation}
compatible with the action of
\[
 \pi=\operatorname{Gal}(\overline K/K)
\]
on the three modules in question, where $\overline K$ is a separable
closure of $K$.  It is known (VIII~\SGARefLink{sga7.sga7i.eviii.n3-2}{3.2}) that this pairing is a perfect
duality: it identifies each of the modules
$T_\ell(A_K(\overline K))$ and $T_\ell(A'_K(\overline K))$ with the
dual of the other, with values in the invertible module
$\mathbb Z_\ell(1)=T_\ell(\overline K^*)$.

}{E7I-IX-001}\SGApage{I}{\(\approx\)320}
At times we shall suppose that a polarization of $A_K$ has been chosen;
thus we are given an element
\begin{equation}\tag{\SGARefTarget{sga7.sga7i.eix.n1-0-4}{1.0.4}}
 \xi\in\operatorname{NS}(A_K/K)
 =\underline{\operatorname{NS}}_{A_K/K}(K),
\end{equation}
where
\[
 \underline{\operatorname{NS}}_{A_K/K}
 \overset{\mathrm{def}}{=}
 \underline{\operatorname{Pic}}_{A_K/K}/
 \bigl(\underline{\operatorname{Pic}}{}^o_{A_K/K}=A'_K\bigr).
\]
In the usual way it defines a homomorphism
\begin{equation}\tag{\SGARefTarget{sga7.sga7i.eix.n1-0-5}{1.0.5}}
 \psi_\xi:A_K\longrightarrow A'_K,
\end{equation}
hence a homomorphism
\begin{equation}\tag{\SGARefTarget{sga7.sga7i.eix.n1-0-6}{1.0.6}}
 T_\ell(\psi_\xi),\quad\text{or simply}\quad
 \psi_\xi:T_\ell(A_K)\longrightarrow T_\ell(A'_K),
\end{equation}
and a bilinear form $\varphi_\xi$
\begin{equation}\tag{\SGARefTarget{sga7.sga7i.eix.n1-0-7}{1.0.7}}
 \begin{split}
 \varphi_\xi(x,y)&=\varphi(x,\psi_\xi(y)),\\
 \varphi_\xi:T_\ell(A_K)\times T_\ell(A_K)
 &\longrightarrow\mathbb Z_\ell(1)_K,
 \end{split}
\end{equation}
which, for $\ell\ne p$, may also be interpreted as a pairing of
finite free $\mathbb Z_\ell$-modules with operators:
\begin{equation}\tag{\SGARefTarget{sga7.sga7i.eix.n1-0-8}{1.0.8}}
 \varphi_\xi:T_\ell(A_K(\overline K))\times
 T_\ell(A_K(\overline K))\longrightarrow T_\ell(\overline K^*).
\end{equation}
These constructions make sense for every element of the
Neron--Severi group of $A_K$, and yield an \underline{alternating} form
$\varphi_\xi$, in either interpretation (\SGARefLink{sga7.sga7i.eix.n1-0-7}{1.0.7}) or (\SGARefLink{sga7.sga7i.eix.n1-0-8}{1.0.8})
(VIII~\SGARefLink{sga7.sga7i.eviii.n2-2-11}{2.2.11}).  The fact that $\xi$ is a \underline{polarization}
implies that $\psi_\xi$ in (\SGARefLink{sga7.sga7i.eix.n1-0-5}{1.0.5}) is an isogeny; equivalently, since
$\dim A_K=\dim A'_K$, it is an epimorphism of groups.  This is also
equivalent to the injectivity of $T_\ell(\psi_\xi)$ in (\SGARefLink{sga7.sga7i.eix.n1-0-6}{1.0.6}) (or to
its having finite cokernel), and finally to the form $\varphi_\xi$ in
(\SGARefLink{sga7.sga7i.eix.n1-0-7}{1.0.7}) being \underline{separating}.

\SGApage{I}{321}
\medskip
\noindent
\SGARefTarget{sga7.sga7i.eix.n1-1}{1.1}. Recall that there exists a \underline{smooth scheme} $A$ over
$S$, called the \underline{Neron model of $A_K$}, which represents the
functor
\begin{equation}\tag{\SGARefTarget{sga7.sga7i.eix.n1-1-1}{1.1.1}}
 S'\longmapsto\operatorname{Hom}(S'_\eta,A_K)
\end{equation}
on the category of smooth $S$-schemes.  Thus, by definition, there is
an isomorphism, functorial in the smooth relative $S$-scheme $S'$,
\begin{equation}\tag{\SGARefTarget{sga7.sga7i.eix.n1-1-2}{1.1.2}}
 \operatorname{Hom}_S(S',A)\simeq
 \operatorname{Hom}_K(S'_\eta,A_K).
\end{equation}
Neron [19] constructed $A$ when the residue field $k$ of $S$ is
perfect.  Raynaud extended the construction to the general case in an
unwritten seminar held at IHES in 1966--67 (see [21] for the statement
of Raynaud's principal results, which complete those of Neron in
loc. cit.).  We shall admit this existence result here, together with
the fact that $A$ is a quasi-projective \underline{scheme} over $S$.
Applying (\SGARefLink{sga7.sga7i.eix.n1-1-2}{1.1.2}) when $S'$ is smooth \underline{over} $\eta$, one also
finds that the following canonical morphism is an isomorphism:
\begin{equation}\tag{\SGARefTarget{sga7.sga7i.eix.n1-1-3}{1.1.3}}
 A\times_S\operatorname{Spec}(K)\xrightarrow{\ \sim\ }A_K.
\end{equation}
This justifies the paired notation $A,A_K$; when convenient we shall
also write $A_\eta$ in place of $A_K$.  Using (\SGARefLink{sga7.sga7i.eix.n1-1-3}{1.1.3}) as an identifying
isomorphism, the canonical bijection (\SGARefLink{sga7.sga7i.eix.n1-1-2}{1.1.2}) is simply the map given by
\underline{restriction to the generic fibers}.

Since the functor (\SGARefLink{sga7.sga7i.eix.n1-1-1}{1.1.1}) plainly has the structure of a commutative
group, the Neron model $A$ representing it itself has the structure of
a \underline{smooth commutative group scheme over $S$}.  The remainder
of this expos\'e is devoted to the study of this group scheme.

\SGApage{I}{\(\approx\)322}
An important invariant arising from the Neron model is, of course, the
smooth commutative group scheme of finite type over $k$
\begin{equation}\tag{\SGARefTarget{sga7.sga7i.eix.n1-1-4}{1.1.4}}
 A_o=A\times_Ss=A_k.
\end{equation}
Following the general notation, we denote its identity component by
$A_o^o$ and also consider the group scheme
\begin{equation}\tag{\SGARefTarget{sga7.sga7i.eix.n1-1-5}{1.1.5}}
 \Phi_o=A_o/A_o^o
\end{equation}
of connected components of $A_o$, which is finite and etale over $k$.
In terms of a separable closure $\overline k$ of $k$, it is therefore
determined by the ordinary finite group
\begin{equation}\tag{\SGARefTarget{sga7.sga7i.eix.n1-1-6}{1.1.6}}
 \Phi_o(\overline k),
\end{equation}
together with the natural action of
\[
 \pi_o=\operatorname{Gal}(\overline k/k)
\]
on this group.

\SGAvar{Along with the Neron model $A$ of $A_K$, we shall study
\begin{equation}\tag{\SGARefTarget{sga7.sga7i.eix.n1-1-7}{1.1.7}}
 A'=\text{the Neron model of }A'_K,
\end{equation}
which is likewise a smooth commutative group scheme of finite type over
$S$.  We shall denote by
\begin{equation}\tag{\SGARefTarget{sga7.sga7i.eix.n1-1-8}{1.1.8}}
 \Phi'_o=A'_o/{A'}_o^o
\end{equation}
the group of connected components of its special fiber $A_o$.



}{Along with the Neron model $A$ of $A_K$, we shall study
\begin{equation}\tag{\SGARefTarget{sga7.sga7i.eix.n1-1-7}{1.1.7}}
 A'=\text{the Neron model of }A'_K,
\end{equation}
which is likewise a smooth commutative group scheme of finite type over
$S$.  We shall denote by
\begin{equation}\tag{\SGARefTarget{sga7.sga7i.eix.n1-1-8}{1.1.8}}
 \Phi'_o=A'_o/{A'}_o^o
\end{equation}
the group of connected components of its special fiber $A'_o$.



}{E7I-IX-002}\SGApage{I}{323}
\medskip
\noindent
\SGARefTarget{sga7.sga7i.eix.n1-2}{1.2}. Apply VIII~\SGARefLink{sga7.sga7i.eviii.n7-1}{7.1}(b) to the pair $(A,A')$ and to the biextension
$W_K$ of $(A_K,A'_K)$ by $\mathbb G_{mK}$ in (\SGARefLink{sga7.sga7i.eix.n1-0-1}{1.0.1}).  We thereby
obtain a \underline{canonical pairing}
\begin{equation}\tag{\SGARefTarget{sga7.sga7i.eix.n1-2-1}{1.2.1}}
 \Phi_o\times\Phi'_o\longrightarrow(\mathbb Q/\mathbb Z)_k,
\end{equation}
that is, a pairing of $\pi_o$-modules
\begin{equation}\tag{\SGARefTarget{sga7.sga7i.eix.n1-2-2}{1.2.2}}
 \Phi_o(\overline k)\times\Phi'_o(\overline k)
 \longrightarrow\mathbb Q/\mathbb Z.
\end{equation}
Here the pairing appears as the obstruction to extending $W_K$ to a
biextension $W$ of $(A,A')$ by $\mathbb G_{mS}$; the meaning of this
point of view is explained in VIII~\SGARefLink{sga7.sga7i.eviii.n7-3-2}{7.3.2}.  We record at once the
following conjecture concerning this pairing.

\medskip
\noindent
\underline{Conjecture \SGARefTarget{sga7.sga7i.eix.n1-3}{1.3}.} \underline{The pairing} (\SGARefLink{sga7.sga7i.eix.n1-2-1}{1.2.1})
\underline{is a perfect duality}: \underline{it identifies each of the
component groups} $\Phi_o$ \underline{and} $\Phi'_o$
\underline{with the dual of the other, with values in}
$(\mathbb Q/\mathbb Z)_k$.

\medskip
\noindent
\SGARefTarget{sga7.sga7i.eix.n1-3-1}{1.3.1}. The formation of (\SGARefLink{sga7.sga7i.eix.n1-2-1}{1.2.1}) commutes with formally etale base
change (VIII~\SGARefLink{sga7.sga7i.eviii.n7-3-5-3}{7.3.5.3}), and the formation of Neron models has the
analogous property.  Consequently, to prove Conjecture~\SGARefLink{sga7.sga7i.eix.n1-3}{1.3} it would
suffice to treat the case in which $S$ is complete and has separably
closed residue field.  Artin and Mazur proved the conjecture when
$A_K$ is the Jacobian of a proper smooth curve over $K$, using their
general, unpublished theory of the self-duality of the relative
Jacobian (in the sense of derived categories) of a proper curve scheme
over $S$.  They also proved Conjecture~\SGARefLink{sga7.sga7i.eix.n1-3}{1.3} when the residue field is
\underline{finite}, using Tate duality [32].  We shall likewise see in
\SGARefLink{sga7.sga7i.eix.n11-4}{11.4} and \SGARefLink{sga7.sga7i.eix.n11-5}{11.5} (subject to compatibility checks) that \SGARefLink{sga7.sga7i.eix.n1-3}{1.3} holds on the
$\ell$-primary components for $\ell\ne p$, and that it holds in the
case of ``semistable reduction''.  These results make \SGARefLink{sga7.sga7i.eix.n1-3}{1.3} extremely
plausible.  Moreover, \SGARefLink{sga7.sga7i.eix.n1-3}{1.3} may also be regarded as a by-product of a
general duality theory for group schemes over fraction fields of
complete discrete valuation rings, formally analogous to SGA~5~I.
Such a theory remains heuristic at present; it would also contain
local class field theory in the form given to it by Serre [27].

\SGApage{I}{324}
\medskip
\noindent
\SGARefTarget{sga7.sga7i.eix.n1-3-2}{1.3.2}. For a given abelian scheme $A_K$, the validity of \SGARefLink{sga7.sga7i.eix.n1-3}{1.3} means that
the pairing (\SGARefLink{sga7.sga7i.eix.n1-2-1}{1.2.1}) is \underline{separating} on both the left and the
right.  By VIII~\SGARefLink{sga7.sga7i.eviii.n7-3-2}{7.3.2}, separation on the left says precisely that the
only open subgroup scheme $V$ of $A$ for which $W_K$ can be extended to
a biextension of $(V,A')$ by $\mathbb G_{mS}$ is the identity component
$A^o$ of $A$.  There is a symmetric interpretation of separation on
the right.

\SGAvar{\medskip
\noindent
\SGARefTarget{sga7.sga7i.eix.n1-4}{1.4}. More generally, suppose that
\begin{equation}\tag{\SGARefTarget{sga7.sga7i.eix.n1-4-1}{1.4.1}}
 U\subset A\qquad\text{and}\qquad U'\subset A'
\end{equation}
are open subgroup schemes such that $W_K$ extends to a biextension of
$(U,U')$ by $\mathbb G_{mS}$.  By VIII~\SGARefLink{sga7.sga7i.eviii.n7-3-2}{7.3.2}, this says explicitly that
the subgroups
\begin{equation}\tag{\SGARefTarget{sga7.sga7i.eix.n1-4-2}{1.4.2}}
 \Phi'_o=U_o/U_o^o\subset\Phi_o,
 \qquad
 \Psi'_o=U'_o/{U'}_o^o\subset\Psi_o
\end{equation}
defined by $U$ and $U'$ are orthogonal under (\SGARefLink{sga7.sga7i.eix.n1-2-1}{1.2.1}).
We denote this extension (unique up to unique isomorphism) of $W_K$ by
$W$.  It is determined up to unique isomorphism by its class
\begin{equation}\tag{\SGARefTarget{sga7.sga7i.eix.n1-4-3}{1.4.3}}
 w\in\mathrm{Biext}^1(U,U';\mathbb G_{mS}),
\end{equation}
which is itself characterized by the condition that its image in
$\mathrm{Biext}^1(A_K,A'_K;\mathbb G_{mK})$ be the class $w_K$ of
(\SGARefLink{sga7.sga7i.eix.n1-0-1}{1.0.1}).

}{\medskip
\noindent
\SGARefTarget{sga7.sga7i.eix.n1-4}{1.4}. More generally, suppose that
\begin{equation}\tag{\SGARefTarget{sga7.sga7i.eix.n1-4-1}{1.4.1}}
 U\subset A\qquad\text{and}\qquad U'\subset A'
\end{equation}
are open subgroup schemes such that $W_K$ extends to a biextension of
$(U,U')$ by $\mathbb G_{mS}$.  By VIII~\SGARefLink{sga7.sga7i.eviii.n7-3-2}{7.3.2}, this says explicitly that
the subgroups
\begin{equation}\tag{\SGARefTarget{sga7.sga7i.eix.n1-4-2}{1.4.2}}
 F_o=U_o/U_o^o\subset\Phi_o,
 \qquad
 F'_o=U'_o/{U'}_o^o\subset\Phi'_o
\end{equation}
defined by $U$ and $U'$ are orthogonal under (\SGARefLink{sga7.sga7i.eix.n1-2-1}{1.2.1}).
We denote this extension (unique up to unique isomorphism) of $W_K$ by
$W$.  It is determined up to unique isomorphism by its class
\begin{equation}\tag{\SGARefTarget{sga7.sga7i.eix.n1-4-3}{1.4.3}}
 w\in\mathrm{Biext}^1(U,U';\mathbb G_{mS}),
\end{equation}
which is itself characterized by the condition that its image in
$\mathrm{Biext}^1(A_K,A'_K;\mathbb G_{mK})$ be the class $w_K$ of
(\SGARefLink{sga7.sga7i.eix.n1-0-1}{1.0.1}).

}{E7I-IX-003}\SGApage{I}{\(\approx\)325}
For example, one may take for $U$ and $U'$ the identity components
$A^o$ and $A'{}^o$ of $A$ and $A'$.  One thereby obtains a
\underline{canonical biextension $W^o$, or $W$, of $(A^o,A'{}^o)$ by
$\mathbb G_{mS}$}, giving rise to a class
\begin{equation}\tag{\SGARefTarget{sga7.sga7i.eix.n1-4-4}{1.4.4}}
 w=w^o\in\mathrm{Biext}^1(A^o,A'{}^o;\mathbb G_{mS}),
\end{equation}
with which we shall work chiefly.  Evidently, the restrictions to
$(A^o,A'{}^o)$ of the classes $w$ in (\SGARefLink{sga7.sga7i.eix.n1-4-3}{1.4.3}), for the various pairs
$(U,U')$, are all equal to $w^o$ in (\SGARefLink{sga7.sga7i.eix.n1-4-4}{1.4.4}).

\medskip
\noindent
\underline{Proposition \SGARefTarget{sga7.sga7i.eix.n1-5}{1.5}.} \underline{Consider the special fiber
$W_s$ of the canonical biextension} $W$ \underline{of}
$(A^o,A'{}^o)$ \underline{by} $\mathbb G_{mS}$ \underline{in (\SGARefLink{sga7.sga7i.eix.n1-4-4}{1.4.4}),
so that $W_s$ is a biextension of $(A_o^o,{A'}_o^o)$ by}
$\mathbb G_{mk}$.  \underline{This biextension is separating}
(VIII~\SGARefLink{sga7.sga7i.eviii.n4-11}{4.11}).

Indeed, choose an ample invertible sheaf $\mathbb L$ on $A^o$.  Then
$\mathbb L_\eta$ defines a homomorphism (cf. \SGARefLink{sga7.sga7i.eix.n1-0-5}{1.0.5})
\[
 A_K\longrightarrow A'_K,
\]
which extends to a homomorphism $A\longrightarrow A'$, and hence gives
\[
 \psi:A^o\longrightarrow A'{}^o.
\]
We may consider the biextension $W_{\mathbb L}$ of $(A^o,A^o)$ by
$\mathbb G_{mS}$ obtained by pulling back $W$ along
$(\operatorname{id}_{A^o},\psi)$.  Its restriction to the generic
fibers is precisely the biextension $\delta(\mathbb L_\eta)$ of
VIII~\SGARefLink{sga7.sga7i.eviii.n4-12}{4.12}.  Hence, by VIII~\SGARefLink{sga7.sga7i.eviii.n7-4}{7.4}, $W_{\mathbb L}$ is associated with the
birigidified invertible sheaf $\delta(\mathbb L)$, and consequently
$(W_{\mathbb L})_s$ is defined by $\delta(\mathbb L_s)$.  On the other
hand, it may also be described as the pullback of $W_s$ by the
morphisms
\[
 \operatorname{id}_{A_o^o}:A_o^o\longrightarrow A_o^o,
 \qquad
 \psi_o:A_o^o\longrightarrow {A'}_o^o.
\]
Since $\mathbb L$, and hence $\mathbb L_s$, is ample, VIII~\SGARefLink{sga7.sga7i.eviii.n4-13}{4.13} shows
that $W_{\mathbb L_s}$ is separating, in particular separating on the
left.  A fortiori $W_s$ is therefore separating on the left.  By
symmetry it is also separating on the right, which proves \SGARefLink{sga7.sga7i.eix.n1-5}{1.5}.

\SGApage{I}{\(\approx\)326}
We shall substantially sharpen \SGARefLink{sga7.sga7i.eix.n1-5}{1.5} when $A_o^o$ has unipotent rank
zero (\SGARefLink{sga7.sga7i.eix.n5-4}{5.4}).



\medskip
\noindent
\underline{2. Fixed and toric parts of $T_\ell(A_K)$.  Criteria for
good reduction.  Orthogonality theorem for $\ell\ne p$.}

\medskip
\noindent
\SGARefTarget{sga7.sga7i.eix.n2-1}{2.1}. Like every commutative algebraic group over $k$, $A_o$ has a
largest subtorus, whose formation commutes with every extension of the
base field [SGA~3 XVII~7.2.1, XII~1.12]:
\begin{equation}\tag{\SGARefTarget{sga7.sga7i.eix.n2-1-1}{2.1.1}}
 T_o\subset A_o.
\end{equation}
When $k$ is perfect, $T_o$ is characterized by the condition that
$A_o^o/T_o$ be an extension of an abelian scheme $B_o$ over $k$ by a
smooth connected unipotent group [25].  The latter necessarily has a
composition series all of whose quotients are isomorphic to
$\mathbb G_{ak}$ (SGA~3 XVII~4.1.3), and one has an exact sequence
\begin{equation}\tag{\SGARefTarget{sga7.sga7i.eix.n2-1-2}{2.1.2}}
 0\longrightarrow U_o\longrightarrow A_o^o/T_o
 \longrightarrow B_o\longrightarrow0.
\end{equation}
In any case, whether or not $k$ is perfect, \underline{if} $A_o^o$ is
an extension of an abelian scheme $B_o$ by a smooth connected affine
scheme $L_o$,
\begin{equation}\tag{\SGARefTarget{sga7.sga7i.eix.n2-1-3}{2.1.3}}
 0\longrightarrow L_o\longrightarrow A_o^o
 \longrightarrow B_o\longrightarrow0,
\end{equation}
this extension structure is known to be unique [25].  Moreover,
$T_o\subset L_o$, and $T_o$ is the largest subtorus of $L_o$.  Thus
one has the exact sequence (\SGARefLink{sga7.sga7i.eix.n2-1-2}{2.1.2}), with
\begin{equation}\tag{\SGARefTarget{sga7.sga7i.eix.n2-1-4}{2.1.4}}
 U_o=L_o/T_o.
\end{equation}

\SGApage{I}{327}
\medskip
\noindent
\SGARefTarget{sga7.sga7i.eix.n2-1-5}{2.1.5}. Let us note that $A_o^o$ has \underline{unipotent rank zero}
(that is, after passage to an algebraic closure of $k$, it has no
subgroup isomorphic to $\mathbb G_a$) if and only if $A_o^o/T_o$ is an
abelian scheme $B_o$.  In that case the preceding hypotheses apply
with $U_o=0$.

\SGAvar{Let $\ell$ be a prime different from the characteristic of $k$, so that
$A_o^o$ is $\ell$-divisible.  The inclusion (\SGARefLink{sga7.sga7i.eix.n2-1-1}{2.1.1}) then gives an exact
sequence of pro-group schemes:
\begin{equation}\tag{\SGARefTarget{sga7.sga7i.eix.n2-1-6}{2.1.6}}
\begin{tikzcd}[row sep=large,column sep=large]
0 \arrow[r] \& T_\ell(T_o) \arrow[r]
 \& T_\ell(A_o^o) \arrow[r] \arrow[d,"\wr"]
 \& T_\ell(A_o^o/T_o) \arrow[r] \arrow[d,"\wr"] \& 0\\
 \& \& T_\ell(A_o) \& T_\ell(A_o/T_o) \& {,}
\end{tikzcd}
\end{equation}
whose terms may also be interpreted as twisted constant free
$\ell$-adic sheaves.  When an exact sequence (\SGARefLink{sga7.sga7i.eix.n2-1-2}{2.1.2}) is available---for
example, when $k$ is perfect or when $A_o^o$ has unipotent rank
zero---one has
\begin{equation}\tag{\SGARefTarget{sga7.sga7i.eix.n2-1-7}{2.1.7}}
 T_\ell(A_o^o/T_o)\xrightarrow{\ \sim\ }T_\ell(B_o),
\end{equation}
so that the exact sequence (\SGARefLink{sga7.sga7i.eix.n2-1-6}{2.1.5}) becomes
\begin{equation}\tag{\SGARefTarget{sga7.sga7i.eix.n2-1-8}{2.1.8}}
 0\longrightarrow T_\ell(T_o)\longrightarrow T_\ell(A_o^o)
 \longrightarrow T_\ell(B_o)\longrightarrow0.
\end{equation}
Put
\begin{equation}\tag{\SGARefTarget{sga7.sga7i.eix.n2-1-9}{2.1.9}}
 n=\dim A_K=\dim A_o,\qquad
 \mu=\dim T_o,\qquad
 \lambda=\dim U_o,\qquad
 \alpha=\dim B_o.
\end{equation}
Then
\begin{equation}\tag{\SGARefTarget{sga7.sga7i.eix.n2-1-10}{2.1.10}}
 n=\mu+\lambda+\alpha,
\end{equation}
and
\begin{equation}\tag{\SGARefTarget{sga7.sga7i.eix.n2-1-11}{2.1.11}}
 \left\{
 \begin{aligned}
 \operatorname{rank}T_\ell(T_o)&=\mu,&
 \operatorname{rank}T_\ell(A_o^o/T_o)
 &=\operatorname{rank}T_\ell(B_o)=2\alpha,\\
 \operatorname{rank}T_\ell(A_o^o)&=\mu+2\alpha.
 \end{aligned}
 \right.
\end{equation}
The second relation is a well-known theorem of Weil [5, 15].

}{Let $\ell$ be a prime different from the characteristic of $k$, so that
$A_o^o$ is $\ell$-divisible.  The inclusion (\SGARefLink{sga7.sga7i.eix.n2-1-1}{2.1.1}) then gives an exact
sequence of pro-group schemes:
\begin{equation}\tag{\SGARefTarget{sga7.sga7i.eix.n2-1-6}{2.1.6}}
\begin{tikzcd}[row sep=large,column sep=large]
0 \arrow[r] \& T_\ell(T_o) \arrow[r]
 \& T_\ell(A_o^o) \arrow[r] \arrow[d,"\wr"]
 \& T_\ell(A_o^o/T_o) \arrow[r] \arrow[d,"\wr"] \& 0\\
 \& \& T_\ell(A_o) \& T_\ell(A_o/T_o) \& {,}
\end{tikzcd}
\end{equation}
whose terms may also be interpreted as twisted constant free
$\ell$-adic sheaves.  When an exact sequence (\SGARefLink{sga7.sga7i.eix.n2-1-2}{2.1.2}) is available---for
example, when $k$ is perfect or when $A_o^o$ has unipotent rank
zero---one has
\begin{equation}\tag{\SGARefTarget{sga7.sga7i.eix.n2-1-7}{2.1.7}}
 T_\ell(A_o^o/T_o)\xrightarrow{\ \sim\ }T_\ell(B_o),
\end{equation}
so that the exact sequence (\SGARefLink{sga7.sga7i.eix.n2-1-6}{2.1.6}) becomes
\begin{equation}\tag{\SGARefTarget{sga7.sga7i.eix.n2-1-8}{2.1.8}}
 0\longrightarrow T_\ell(T_o)\longrightarrow T_\ell(A_o^o)
 \longrightarrow T_\ell(B_o)\longrightarrow0.
\end{equation}
Put
\begin{equation}\tag{\SGARefTarget{sga7.sga7i.eix.n2-1-9}{2.1.9}}
 n=\dim A_K=\dim A_o,\qquad
 \mu=\dim T_o,\qquad
 \lambda=\dim U_o,\qquad
 \alpha=\dim B_o.
\end{equation}
Then
\begin{equation}\tag{\SGARefTarget{sga7.sga7i.eix.n2-1-10}{2.1.10}}
 n=\mu+\lambda+\alpha,
\end{equation}
and
\begin{equation}\tag{\SGARefTarget{sga7.sga7i.eix.n2-1-11}{2.1.11}}
 \left\{
 \begin{aligned}
 \operatorname{rank}T_\ell(T_o)&=\mu,&
 \operatorname{rank}T_\ell(A_o^o/T_o)
 &=\operatorname{rank}T_\ell(B_o)=2\alpha,\\
 \operatorname{rank}T_\ell(A_o^o)&=\mu+2\alpha.
 \end{aligned}
 \right.
\end{equation}
The second relation is a well-known theorem of Weil [5, 15].

}{E7I-IX-004}\SGApage{I}{328}
\medskip
\noindent
\SGARefTarget{sga7.sga7i.eix.n2-1-12}{2.1.12}. The integers $n,\mu,\lambda,\alpha$ are plainly defined, after
passing to an algebraic closure of $k$, even without assuming the
existence of an exact sequence (\SGARefLink{sga7.sga7i.eix.n2-1-2}{2.1.2}).  They are respectively called
the dimension of $A_o$, the \underline{reductive rank of $A_o$}
(equal to the dimension of $T_o$), the \underline{unipotent rank of
$A_o$}, and the \underline{abelian rank of $A_o$}.  They are still
related by (\SGARefLink{sga7.sga7i.eix.n2-1-10}{2.1.10}) and (\SGARefLink{sga7.sga7i.eix.n2-1-11}{2.1.11}), with the term
$\operatorname{rank}T_\ell(B_o)$ omitted from the latter.

\SGAvar{\medskip
\noindent
\SGARefTarget{sga7.sga7i.eix.n2-1-13}{2.1.13}. Suppose that $k$ has characteristic $p>0$ and that $A_o^o$ is
$p$-divisible, or equivalently that it has unipotent rank zero
(cf. \SGARefLink{sga7.sga7i.eix.n2-1-5}{1.5.5}).  One should then also consider the pro-group $T_p(A_o^o)$,
which is a \underline{Barsotti--Tate group} (or a
\underline{$p$-divisible group}, in the terminology of [33]).  The
exact sequence (\SGARefLink{sga7.sga7i.eix.n2-1-8}{2.1.8}) remains valid for $\ell=p$ and is then an exact
sequence of Barsotti--Tate groups.  Here $\lambda=0$, and hence
\begin{equation}\tag{\SGARefTarget{sga7.sga7i.eix.n2-1-13-1}{2.1.13.1}}
 n=\mu+\alpha.
\end{equation}
The formulas (\SGARefLink{sga7.sga7i.eix.n2-1-11}{2.1.11}) also remain valid, provided that rank is
interpreted as the rank (or ``height'') of the Barsotti--Tate groups in
question.

}{\medskip
\noindent
\SGARefTarget{sga7.sga7i.eix.n2-1-13}{2.1.13}. Suppose that $k$ has characteristic $p>0$ and that $A_o^o$ is
$p$-divisible, or equivalently that it has unipotent rank zero
(cf. \SGARefLink{sga7.sga7i.eix.n2-1-5}{2.1.5}).  One should then also consider the pro-group $T_p(A_o^o)$,
which is a \underline{Barsotti--Tate group} (or a
\underline{$p$-divisible group}, in the terminology of [33]).  The
exact sequence (\SGARefLink{sga7.sga7i.eix.n2-1-8}{2.1.8}) remains valid for $\ell=p$ and is then an exact
sequence of Barsotti--Tate groups.  Here $\lambda=0$, and hence
\begin{equation}\tag{\SGARefTarget{sga7.sga7i.eix.n2-1-13-1}{2.1.13.1}}
 n=\mu+\alpha.
\end{equation}
The formulas (\SGARefLink{sga7.sga7i.eix.n2-1-11}{2.1.11}) also remain valid, provided that rank is
interpreted as the rank (or ``height'') of the Barsotti--Tate groups in
question.

}{E7I-IX-005}\SGApage{I}{329}
\medskip
\noindent
\SGARefTarget{sga7.sga7i.eix.n2-1-14}{2.1.14}. For ${A'}_o^o$ we likewise introduce the invariants
\[
 T'_o\subset {A'}_o^o,\qquad U'_o,\qquad B'_o,
\]
where $T'_o$ is the toric part of ${A'}_o^o$, $B'_o$ its abelian part
(defined in particular when $k$ is perfect or when ${A'}_o^o$ has
unipotent rank zero), and $U'_o$ its unipotent part (defined, for
example, when $k$ is perfect, or trivially when ${A'}_o^o$ has
unipotent rank zero).  There is a corresponding exact sequence of the
form (\SGARefLink{sga7.sga7i.eix.n2-1-6}{2.1.6}), or (\SGARefLink{sga7.sga7i.eix.n2-1-8}{2.1.8}), which we shall not write again.  One could
introduce integers $\mu',\lambda',\alpha'$ for ${A'}_o^o$ as in
(\SGARefLink{sga7.sga7i.eix.n2-1-9}{2.1.9}), but we shall see below (\SGARefLink{sga7.sga7i.eix.n2-2-7}{2.2.7}) that they are the same as those
for $A_o^o$; this spares us from introducing new notation for them.



\medskip
\noindent
\SGARefTarget{sga7.sga7i.eix.n2-2}{2.2}. We propose to make more precise the relations between
$T_\ell(A_K)$ and $T_\ell(A_o)=T_\ell(A_o^o)$ by bringing
$T_\ell(A^o)$ into play.  First recall the following.

\medskip
\noindent
\underline{Lemma} \SGARefTarget{sga7.sga7i.eix.n2-2-1}{2.2.1}. \underline{Let $m>0$ be an integer.  Then the
following conditions are equivalent:}

\smallskip
\noindent
(i) $m\,\mathrm{id}_{A^o}$ \underline{is flat}.

\noindent
(ii) $m\,\mathrm{id}_{A^o}$ \underline{is surjective}.

\noindent
(iii) $m\,\mathrm{id}_{A^o}$ \underline{is quasi-finite}.

\noindent
(iv) $m$ \underline{is prime to the characteristic of $k$}, or $A_o^o$
\underline{has unipotent rank zero} (cf. \SGARefLink{sga7.sga7i.eix.n2-1-5}{2.1.5}).

\smallskip
\noindent
\underline{Moreover, if these conditions hold for $m$, they hold for
$m^\nu$ for every $\nu>0$; the groups}
\[
 {}_{m^\nu}A^o\underset{\mathrm{dfn}}{=}
 \ker(m^\nu\,\mathrm{id}_{A^o})
\]
\underline{are quasi-finite, separated, and flat over $S$, and, as
$\nu$ varies, they form an inverse system of group schemes whose
transition morphisms}
\[
 x\longmapsto m^{\nu'-\nu}x:
 {}_{m^{\nu'}}A^o\longrightarrow{}_{m^\nu}A^o
 \qquad(\nu'\geq\nu)
\]
\underline{are faithfully flat}.

\SGApage{I}{\(\approx\)330}
Indeed, the fiberwise criterion for flatness (EGA~IV~11.3.5) shows that
(i), like (ii) and (iii), is a condition on the fibers; these three
conditions are equivalent by SGA~3 $\mathrm{VI}_A$~5.6, since the
fibers in question are smooth and connected.  In the form (ii), a
classical theorem of Weil [5, 15] shows that the conditions hold on the
generic fiber, which is an abelian scheme; it is therefore enough to
state them on the special fiber.  Their equivalence with (iv) is then
also well known from the structure theory.  In any of these forms,
stability of the condition under $m\mapsto m^\nu$ is immediate.  It
follows that the ${}_{m^\nu}A^o$ are quasi-finite flat group schemes
over $S$, and they are plainly separated over $S$ because $A^o$ is.
Finally, faithful flatness of the transition morphisms---equivalently,
the fact that they are epimorphisms for the fppf topology---follows
formally from the fact that $m\,\mathrm{id}_{A^o}$ is faithfully flat,
hence an epimorphism for that topology.

\medskip
\noindent
\SGARefTarget{sga7.sga7i.eix.n2-2-2}{2.2.2}. Consequently, we shall take $m$ to be a prime $\ell$ satisfying
the conditions of \SGARefLink{sga7.sga7i.eix.n2-2-1}{2.2.1}.  We then denote by
\begin{equation}\tag{\SGARefTarget{sga7.sga7i.eix.n2-2-2-1}{2.2.2.1}}
 T_\ell(A^o)=({}_{\ell^\nu}A^o)_{\nu>0}
\end{equation}
the inverse system described in \SGARefLink{sga7.sga7i.eix.n2-2-1}{2.2.1}.  Under the two base changes
$\eta\longrightarrow S$ and $s\longrightarrow S$, it recovers the
inverse systems $T_\ell(A_K)$ and
$T_\ell(A_o^o)=T_\ell(A_o)$ considered in \SGARefLink{sga7.sga7i.eix.n1-0}{1.0} and \SGARefLink{sga7.sga7i.eix.n2-1}{2.1} respectively.

\SGApage{I}{331}
\medskip
\noindent
\SGARefTarget{sga7.sga7i.eix.n2-2-3}{2.2.3}. Suppose now that $S$ is henselian.  Every quasi-finite separated
scheme $X$ over $S$ then has a canonical decomposition as a disjoint
union
\begin{equation}\tag{\SGARefTarget{sga7.sga7i.eix.n2-2-3-1}{2.2.3.1}}
 X=X^f\coprod X',
\end{equation}
where $X^f$ is \underline{finite over $S$} and $X'_o=\varnothing$, that
is, $X'$ is confined to its generic fiber $X'_\eta$; thus $X^f$ has the
same special fiber as $X$.  This decomposition is plainly functorial in
$X$.  In particular, if $X$ is a group scheme over $S$, then $X^f$ is
a subgroup scheme.  If $X$ is flat, respectively \'{e}tale, over $S$,
then so is $X^f$.

We apply this to the components ${}_{\ell^\nu}A^o$ of $T_\ell(A^o)$.
This gives a projective subsystem of groups
\begin{equation}\tag{\SGARefTarget{sga7.sga7i.eix.n2-2-3-2}{2.2.3.2}}
 \bigl(({}_{\ell^\nu}A^o)^f\bigr)_{\nu>0}
 \underset{\mathrm{dfn}}{=}T_\ell(A^o)^f
 \subset T_\ell(A^o),
\end{equation}
called the \underline{fixed part} of the Tate module $T_\ell(A^o)$.
By construction, there are canonical isomorphisms, functorial in $A$,
\begin{equation}\tag{\SGARefTarget{sga7.sga7i.eix.n2-2-3-3}{2.2.3.3}}
 T_\ell(A^o)^f\times_S s\xrightarrow{\ \sim\ }
 T_\ell(A^o)\times_Ss\xrightarrow{\ \sim\ }T_\ell(A_o^o).
\end{equation}
On the other hand, the generic fiber of $T_\ell(A^o)^f$ defines a
canonical sub-pro-group of $T_\ell(A_K)$, again called the
``\underline{fixed part}'' of this Tate module:
\begin{equation}\tag{\SGARefTarget{sga7.sga7i.eix.n2-2-3-4}{2.2.3.4}}
 T_\ell(A_K)^f\subset T_\ell(A_K).
\end{equation}

\SGApage{I}{332}
\medskip
\noindent
\SGARefTarget{sga7.sga7i.eix.n2-2-4}{2.2.4}. When $\ell\ne p=\operatorname{char}k$, the fixed part
$T_\ell(A^o)^f$ may be interpreted as a twisted constant $\ell$-adic
sheaf on $S$ (the $({}_{\ell^\nu}A^o)^f$ being finite \'{e}tale over
$S$).  Since the functor
\[
 X\longmapsto X_o=X\times_Ss
\]
from finite \'{e}tale schemes over $S$ to finite \'{e}tale schemes
over $k$ is an equivalence of categories when $S$ is henselian, there
is an analogous equivalence for the categories of twisted constant
$\ell$-adic sheaves over $S$ and over $k$.  Hence (\SGARefLink{sga7.sga7i.eix.n2-2-3-3}{2.2.3.3}), which
gives the special fiber of $T_\ell(A^o)^f$, characterizes this fixed
part up to unique isomorphism.  On the other hand, as an $\ell$-adic
subsheaf of $T_\ell(A^o)$ it is plainly characterized by its generic
fiber (\SGARefLink{sga7.sga7i.eix.n2-2-3-4}{2.2.3.4}), or equivalently by the corresponding
$\mathbb Z_\ell$-submodule, stable under
$\pi=\operatorname{Gal}(\overline K/K)$:
\begin{equation}\tag{\SGARefTarget{sga7.sga7i.eix.n2-2-4-1}{2.2.4.1}}
 T_\ell(A_K)^f(\overline K)
 \underset{\mathrm{dfn}}{=}T_\ell(A(\overline K))^f
 \subset T_\ell(A(\overline K)).
\end{equation}
In fact, one has the following result.

\medskip
\noindent
\underline{Proposition} \SGARefTarget{sga7.sga7i.eix.n2-2-5}{2.2.5}. \underline{Suppose, as above, that
$\ell\ne p$ and that $S$ is henselian.  Then the submodule (\SGARefLink{sga7.sga7i.eix.n2-2-4-1}{2.2.4.1})
of $T_\ell(A(\overline K))=T_\ell(A_K)(\overline K)$ corresponding to
the fixed part $T_\ell(A_K)^f$ in (\SGARefLink{sga7.sga7i.eix.n2-2-3-4}{2.2.3.4}) is the submodule
$T_\ell(A(\overline K))^I$ of invariants under the inertia group
$I\subset\pi$.}

\SGAvar{This follows, by passage to the inverse limit, from the analogous
relation valid for every integer $m>0$ prime to
$p=\operatorname{char}k$:
\begin{equation}\tag{\SGARefTarget{sga7.sga7i.eix.n2-2-5-1}{2.2.5.1}}
 ({}_mA)^f(\overline K)=({}_mA(\overline K))^I.
\end{equation}
To prove this relation, introduce the strict localization $\widetilde S$
of $S$, whose fraction field $\widetilde K$ is the maximal unramified
subextension of $\overline K/K$.  Since $({}_mA)^f$ is finite
\'{e}tale over $S$, its points with values in $\overline K$ already
have values in $\widetilde K$.  Thus the left-hand side of (\SGARefLink{sga7.sga7i.eix.n2-2-5-1}{2.2.5.1})
is
$\operatorname{Hom}_S(\widetilde S,({}_mA)^f)$, whereas the right-hand
side is, by definition of $\widetilde K$, equal to ${}_mA(\widetilde K)$.
By the definition of $({}_mA)^f$, and since $\widetilde S$ is an
inverse limit of finite \'{e}tale $S$-schemes $S_i$, the former is also
equal to
\[
 \operatorname{Hom}_S(\widetilde S,{}_mA)
 ={}_m\operatorname{Hom}_S(\widetilde S,A).
\]
By the defining property of the N\'{e}ron model,
$\operatorname{Hom}_S(\widetilde S,A)=A(\widetilde K)$, which proves
(\SGARefLink{sga7.sga7i.eix.n2-2-5-1}{6.2.5.1}).

}{This follows, by passage to the inverse limit, from the analogous
relation valid for every integer $m>0$ prime to
$p=\operatorname{char}k$:
\begin{equation}\tag{\SGARefTarget{sga7.sga7i.eix.n2-2-5-1}{2.2.5.1}}
 ({}_mA)^f(\overline K)=({}_mA(\overline K))^I.
\end{equation}
To prove this relation, introduce the strict localization $\widetilde S$
of $S$, whose fraction field $\widetilde K$ is the maximal unramified
subextension of $\overline K/K$.  Since $({}_mA)^f$ is finite
\'{e}tale over $S$, its points with values in $\overline K$ already
have values in $\widetilde K$.  Thus the left-hand side of (\SGARefLink{sga7.sga7i.eix.n2-2-5-1}{2.2.5.1})
is
$\operatorname{Hom}_S(\widetilde S,({}_mA)^f)$, whereas the right-hand
side is, by definition of $\widetilde K$, equal to ${}_mA(\widetilde K)$.
By the definition of $({}_mA)^f$, and since $\widetilde S$ is an
inverse limit of finite \'{e}tale $S$-schemes $S_i$, the former is also
equal to
\[
 \operatorname{Hom}_S(\widetilde S,{}_mA)
 ={}_m\operatorname{Hom}_S(\widetilde S,A).
\]
By the defining property of the N\'{e}ron model,
$\operatorname{Hom}_S(\widetilde S,A)=A(\widetilde K)$, which proves
(\SGARefLink{sga7.sga7i.eix.n2-2-5-1}{2.2.5.1}).

}{E7I-IX-006}\SGApage{I}{\(\approx\)333}
To deduce \SGARefLink{sga7.sga7i.eix.n2-2-5}{2.2.5}, observe that (\SGARefLink{sga7.sga7i.eix.n2-2-5-1}{2.2.5.1}) gives
\begin{equation}\tag{*}
 T_\ell(A(\overline K))^I
 =\varprojlim_\nu({}_{\ell^\nu}A)^f(\overline K)
 \supset\varprojlim_\nu({}_{\ell^\nu}A^o)^f(\overline K).
\end{equation}
If $(x_\nu)$ belongs to the rightmost term, then for every
$\nu\geq0$, $x_\nu$ is a section of $A$ that is infinitely divisible
by $\ell$.  Its value at $s$ is therefore an element of $A_o(k)$
infinitely divisible by $\ell$, and consequently lies in $A_o^o(k)$.
This proves equality of the last two terms of (*) and completes the
proof of \SGARefLink{sga7.sga7i.eix.n2-2-5}{2.2.5}.

\medskip
\noindent
\underline{Proposition} \SGARefTarget{sga7.sga7i.eix.n2-2-6}{2.2.6}. \underline{Let $C_K$ be a second
abelian scheme over $K$, with N\'{e}ron model $C$, and consider a
homomorphism $u_K:C_K\longrightarrow A_K$, which therefore extends to
a homomorphism $u:C\longrightarrow A$.  Suppose that $u_K$ is an
isogeny.  Then, for every prime $\ell$,
$T_\ell(u_K):T_\ell(C_K)\longrightarrow T_\ell(A_K)$ induces an
isogeny of fixed parts}
\begin{equation}\tag{\SGARefTarget{sga7.sga7i.eix.n2-2-6-1}{2.2.6.1}}
 T_\ell(C_K)^f\longrightarrow T_\ell(A_K)^f,
\end{equation}
\underline{and likewise the map induced by
$u_o:C_o\longrightarrow A_o$}
\begin{equation}\tag{\SGARefTarget{sga7.sga7i.eix.n2-2-6-2}{2.2.6.2}}
 T_\ell(u_o^o):T_\ell(C_o^o)\longrightarrow T_\ell(A_o^o)
\end{equation}
\underline{is an isogeny.  Finally, $u_o:C_o\longrightarrow A_o$
induces an isogeny from the largest subtorus of $C_o$ to the largest
subtorus $T_o$ of $A_o$; and over an extension $k'/k$ on which the}
``\underline{abelian parts}'' \underline{of $C_o$ and $A_o$ are defined (for example, over
the perfect closure of $k$), $u_o$ induces an isogeny between the
abelian parts of $C_{ok'}$ and $A_{ok'}$.}

\SGApage{I}{\(\approx\)334}
This statement follows formally from the fact that
$T_\ell(A_K)^f$, $T_\ell(A_o^o)$, the largest subtorus of $A_o$, and
the ``abelian part'' of $A_{ok'}$ are additive functors of the varying
$A_K$, and from the fact that $u_K$ is an isogeny precisely in the
sense that there is a morphism $v_K:A_K\longrightarrow C_K$ such that
\[
 u_Kv_K=n\,\mathrm{id},\qquad v_Ku_K=n\,\mathrm{id}
\]
for some integer $n>0$.

\medskip
\noindent
\underline{Corollary} \SGARefTarget{sga7.sga7i.eix.n2-2-7}{2.2.7}. \underline{If $A_K$ and $C_K$ are
isogenous abelian schemes over $K$, then the corresponding invariants
$\mu,\lambda,\alpha$ in (\SGARefLink{sga7.sga7i.eix.n2-1-11}{2.1.11}) are the same for $A_K$ and $C_K$.
In particular, they are the same for $A_K$ and its dual $A'_K$.}

\medskip
\noindent
\underline{Remarks} \SGARefTarget{sga7.sga7i.eix.n2-2-8}{2.2.8}. In the two preceding statements, \SGARefLink{sga7.sga7i.eix.n2-2-6}{2.2.6} and
\SGARefLink{sga7.sga7i.eix.n2-2-7}{2.2.7}, we assumed that $S$ is henselian.  It is immediate, however,
that \SGARefLink{sga7.sga7i.eix.n2-2-7}{2.2.7}, as well as the assertions of \SGARefLink{sga7.sga7i.eix.n2-2-6}{2.2.6} concerning the
properties of the induced homomorphism
$u_o:C_o\longrightarrow A_o$, holds without this hypothesis.  Indeed,
one reduces to the case already treated by the base change
$S^h\longrightarrow S$, where $S^h$ is the henselization of $S$, using
the fact---immediate from the definitions---that formation of N\'{e}ron
models commutes with this base change.

\SGApage{I}{335}
\medskip
\noindent
\underline{Corollary} \SGARefTarget{sga7.sga7i.eix.n2-2-9}{2.2.9}. (\underline{Good reduction.})
\underline{Let $S$ be a trait, let $A_K$ be an abelian scheme over its
fraction field $K$, let $A$ be the N\'{e}ron model of $A_K$, and let
$\ell$ be a prime different from $\operatorname{char}k$.  The following
conditions are equivalent.}

\smallskip
\noindent
(i) \underline{There exists an abelian scheme over $S$ whose generic
fiber is isomorphic to $A_K$.}

\noindent
(ii) $A$ \underline{is an abelian scheme}.

\noindent
(ii$^o$) $A^o$ \underline{is an abelian scheme}.

\noindent
(iii) $A_o$ \underline{is an abelian scheme}.

\noindent
(iii$^o$) $A_o^o$ \underline{is an abelian scheme}, \underline{that is,
with the notation of (\SGARefLink{sga7.sga7i.eix.n2-1-11}{2.1.11}), $\lambda=\mu=0$}.

\noindent
(iv) $A$ \underline{is proper over $S$}.

\noindent
(iv$^o$) $A^o$ \underline{is proper over $S$}.

\noindent
(v) $T_\ell(A_K)$ \underline{is} ``\underline{unramified over $S$}'', that is,
\underline{it extends to a twisted constant $\ell$-adic sheaf over
$S$}, or equivalently \underline{the inertia group $I$ acts trivially
on $T_\ell(A_K(\overline K))$}.

\noindent
(v bis) \underline{For every integer $\nu\geq0$,
${}_{\ell^\nu}A_K$ is} ``\underline{unramified over $S$}'', that is,
\underline{it extends to an \'{e}tale covering of $S$}, or equivalently
\underline{the inertia group $I$ acts trivially on
${}_{\ell^\nu}A_K(\overline K)$}.

\smallskip
\noindent
\underline{Finally, an abelian scheme over $S$ as in (i) is a N\'{e}ron
model of $A_K$, and in particular is determined up to unique
isomorphism.}

\medskip
\noindent
\SGARefTarget{sga7.sga7i.eix.n2-2-9-1}{2.2.9.1}. We say that $A_K$ has \underline{good reduction} over $S$ if
it satisfies condition (i) above.  Criterion (v), or (v bis), is known
as the \underline{N\'{e}ron--Ogg--Shafarevich criterion} for good
reduction.

Let us prove \SGARefLink{sga7.sga7i.eix.n2-2-9}{2.2.9}.  It is well known (cf., for example, [5]) that an
abelian scheme over $S$ satisfies the universal property of a
N\'{e}ron model; this reduces to the fact that a rational section of an
abelian scheme over a regular base is defined everywhere.  This proves
the equivalence
$(i)\Longleftrightarrow(ii)\Longleftrightarrow(ii^o)$, as well as the
last assertion of \SGARefLink{sga7.sga7i.eix.n2-2-9}{2.2.9}.  These conditions plainly imply (iii), (iv),
and (iv$^o$), each of which implies (iii$^o$).  Let us prove that
(iii$^o$)$\Longrightarrow$(ii$^o$), which will establish the
equivalence of (i) through (iv$^o$).  This follows from the fact that a
smooth group scheme over $S$ whose fibers are abelian schemes is proper
over $S$ (and hence is an abelian scheme), by EGA~IV~15.7.10.

\SGApage{I}{\(\approx\)336}
On the other hand, in view of \SGARefLink{sga7.sga7i.eix.n2-2-5}{2.2.5} and the fact that
\[
 T_\ell(A_K(\overline K))/T_\ell(A_K(\overline K))^I
\]
is plainly a torsion-free $\mathbb Z_\ell$-module, condition (v) says
that $T_\ell(A_K(\overline K))$ and its fixed part have the same rank.
By (\SGARefLink{sga7.sga7i.eix.n2-1-11}{2.1.11}) and (\SGARefLink{sga7.sga7i.eix.n2-2-3-3}{2.2.3.3}), this is the condition
\[
 \mu+2\alpha=2(\lambda+\mu+\alpha),
\]
that is, $2\lambda+\mu=0$, or $\lambda=\mu=0$, which is precisely
(iii$^o$).  Finally, the equivalence of (v) and (v bis) is clear from
the definitions.  This completes the proof of \SGARefLink{sga7.sga7i.eix.n2-2-9}{2.2.9}.



\SGAvar{\medskip
\noindent
\SGARefTarget{sga7.sga7i.eix.n2-3}{2.3}. Suppose once again that $S$ is henselian and that
$\ell\ne p=\operatorname{char}k$.  We now use the inclusion
\begin{equation}\tag{\SGARefTarget{sga7.sga7i.eix.n2-3-1}{2.3.1}}
 T_\ell(T_o)\hookrightarrow T_\ell(A_o^o)
\end{equation}
of (\SGARefLink{sga7.sga7i.eix.n2-1-6}{2.1.6}), together with the equivalence of categories recalled in
\SGARefLink{sga7.sga7i.eix.n2-2-4}{2.2.4}.  We obtain a canonical twisted constant $\ell$-adic subsheaf
of $T_\ell(A^o)^f$, called the \underline{toric part of}
$T_\ell(A^o)$,
\begin{equation}\tag{\SGARefTarget{sga7.sga7i.eix.n2-3-2}{2.3.2}}
 T_\ell(A^o)^t\hookrightarrow T_\ell(A^o)^f
 \hookrightarrow T_\ell(A^o),
\end{equation}
characterized by the condition that its special fiber is the inclusion
(\SGARefLink{sga7.sga7i.eix.n2-3-1}{2.3.1}):
\begin{equation}\tag{\SGARefTarget{sga7.sga7i.eix.n2-3-3}{2.3.3}}
 T_\ell(A^o)^t\times_Ss=T_\ell(T_o)
 \hookrightarrow T_\ell(A_o^o).
\end{equation}
Later we shall give an interpretation of the toric part of
$T_\ell(A^o)$ as the Tate module $T_\ell(T)$ of a torus $T$ over $S$
lifting the torus $T_o$ over $k$.  That interpretation will allow the
definition of the toric part to be extended to the case $\ell=p$.
For the moment, and to fix ideas, we confine ourselves to the case
$\ell\ne p$.  In this case the toric part of $T_\ell(A^o)$ may be
described by the corresponding $\mathbb Z_\ell$-submodule of
$T_\ell(A(\overline K))$:
\begin{equation}\tag{\SGARefTarget{sga7.sga7i.eix.n2-3-4}{2.3.4}}
 T_\ell(A(\overline K))^t\hookrightarrow
 T_\ell(A(\overline K))^f\hookrightarrow
 T_\ell(A(\overline K)).
\end{equation}
It is stable under the action of $\pi$; its first term is defined as
$T_\ell(A^o)^t(\overline K)$ and is, as one would expect, again called
the \underline{toric submodule} of the Tate module
$T_\ell(A(\overline K))$.  Thus this Tate module carries the canonical
three-step filtration (\SGARefLink{sga7.sga7i.eix.n2-3-4}{2.3.4}).  Applying the same considerations to
the dual abelian scheme $A'_K$, we obtain the analogous three-step
filtration
\begin{equation}\tag{\SGARefTarget{sga7.sga7i.eix.n2-3-5}{2.3.5}}
 T_\ell(A'(\overline K))^t\hookrightarrow
 T_\ell(A'(\overline K))^f\hookrightarrow
 T_\ell(A'(\overline K)).
\end{equation}
The middle terms in (\SGARefLink{sga7.sga7i.eix.n2-3-4}{2.3.4}) and (\SGARefLink{sga7.sga7i.eix.n2-3-5}{2.3.5}), the ``fixed parts,'' are
described in terms of the $\pi$-module structures of the Tate modules
of $A(\overline K)$ and $(A'(\overline K)$ as their submodules of
invariants under the inertia group $I$ (\SGARefLink{sga7.sga7i.eix.n2-2-5}{1.6.5}).  The toric parts are
then completely described in terms of the fixed parts and the pairing
$\varphi$ of (\SGARefLink{sga7.sga7i.eix.n1-0-2}{1.0.2}), by the following theorem (a special case of which
already appears in [11]).

}{\medskip
\noindent
\SGARefTarget{sga7.sga7i.eix.n2-3}{2.3}. Suppose once again that $S$ is henselian and that
$\ell\ne p=\operatorname{char}k$.  We now use the inclusion
\begin{equation}\tag{\SGARefTarget{sga7.sga7i.eix.n2-3-1}{2.3.1}}
 T_\ell(T_o)\hookrightarrow T_\ell(A_o^o)
\end{equation}
of (\SGARefLink{sga7.sga7i.eix.n2-1-6}{2.1.6}), together with the equivalence of categories recalled in
\SGARefLink{sga7.sga7i.eix.n2-2-4}{2.2.4}.  We obtain a canonical twisted constant $\ell$-adic subsheaf
of $T_\ell(A^o)^f$, called the \underline{toric part of}
$T_\ell(A^o)$,
\begin{equation}\tag{\SGARefTarget{sga7.sga7i.eix.n2-3-2}{2.3.2}}
 T_\ell(A^o)^t\hookrightarrow T_\ell(A^o)^f
 \hookrightarrow T_\ell(A^o),
\end{equation}
characterized by the condition that its special fiber is the inclusion
(\SGARefLink{sga7.sga7i.eix.n2-3-1}{2.3.1}):
\begin{equation}\tag{\SGARefTarget{sga7.sga7i.eix.n2-3-3}{2.3.3}}
 T_\ell(A^o)^t\times_Ss=T_\ell(T_o)
 \hookrightarrow T_\ell(A_o^o).
\end{equation}
Later we shall give an interpretation of the toric part of
$T_\ell(A^o)$ as the Tate module $T_\ell(T)$ of a torus $T$ over $S$
lifting the torus $T_o$ over $k$.  That interpretation will allow the
definition of the toric part to be extended to the case $\ell=p$.
For the moment, and to fix ideas, we confine ourselves to the case
$\ell\ne p$.  In this case the toric part of $T_\ell(A^o)$ may be
described by the corresponding $\mathbb Z_\ell$-submodule of
$T_\ell(A(\overline K))$:
\begin{equation}\tag{\SGARefTarget{sga7.sga7i.eix.n2-3-4}{2.3.4}}
 T_\ell(A(\overline K))^t\hookrightarrow
 T_\ell(A(\overline K))^f\hookrightarrow
 T_\ell(A(\overline K)).
\end{equation}
It is stable under the action of $\pi$; its first term is defined as
$T_\ell(A^o)^t(\overline K)$ and is, as one would expect, again called
the \underline{toric submodule} of the Tate module
$T_\ell(A(\overline K))$.  Thus this Tate module carries the canonical
three-step filtration (\SGARefLink{sga7.sga7i.eix.n2-3-4}{2.3.4}).  Applying the same considerations to
the dual abelian scheme $A'_K$, we obtain the analogous three-step
filtration
\begin{equation}\tag{\SGARefTarget{sga7.sga7i.eix.n2-3-5}{2.3.5}}
 T_\ell(A'(\overline K))^t\hookrightarrow
 T_\ell(A'(\overline K))^f\hookrightarrow
 T_\ell(A'(\overline K)).
\end{equation}
The middle terms in (\SGARefLink{sga7.sga7i.eix.n2-3-4}{2.3.4}) and (\SGARefLink{sga7.sga7i.eix.n2-3-5}{2.3.5}), the ``fixed parts,'' are
described in terms of the $\pi$-module structures of the Tate modules
of $A(\overline K)$ and $A'(\overline K)$ as their submodules of
invariants under the inertia group $I$ (\SGARefLink{sga7.sga7i.eix.n2-2-5}{2.2.5}).  The toric parts are
then completely described in terms of the fixed parts and the pairing
$\varphi$ of (\SGARefLink{sga7.sga7i.eix.n1-0-2}{1.0.2}), by the following theorem (a special case of which
already appears in [11]).

}{E7I-IX-007}\SGApage{I}{338}
\medskip
\noindent
\underline{Theorem} \SGARefTarget{sga7.sga7i.eix.n2-4}{2.4}. (\underline{Orthogonality theorem}.)
\underline{Let $S$ be a henselian trait, let $A_K$ be an abelian scheme
over its fraction field $K$, let $A'_K$ be the dual abelian scheme, and
let $\ell$ be a prime different from the residual characteristic of
$S$.  The Tate modules $T_\ell(A(\overline K))$ and
$T_\ell(A'(\overline K))$ then carry the canonical filtrations (\SGARefLink{sga7.sga7i.eix.n2-3-4}{2.3.4})
and (\SGARefLink{sga7.sga7i.eix.n2-3-5}{2.3.5}), respectively, by the ``fixed part'' and the ``toric
part.''  The toric part $T_\ell(A(\overline K))^t$ is the intersection
of the fixed part $T_\ell(A(\overline K))^f$ with the orthogonal of the
fixed part
$T_\ell(A'(\overline K))^f=T_\ell(A'(\overline K))^I$ with respect to
the canonical pairing $\varphi$ of (\SGARefLink{sga7.sga7i.eix.n1-0-3}{1.0.3}).  Symmetrically, the toric
part $T_\ell(A'(\overline K))^t$ is the intersection of
$T_\ell(A'(\overline K))^f$ with the orthogonal of the fixed part
$T_\ell(A(\overline K))^f=T_\ell(A(\overline K))^I$.  (Here the
superscript $I$ denotes the submodule of invariants under the action of
the inertia group $I$.)}

\SGAvar{\smallskip
\noindent
\underline{Proof.}\quad a) Let us prove that
$T_\ell(A(\overline K))^t$ is orthogonal to
$T_\ell(A'(\overline K))^f$.  To do so, we interpret the restriction of
$\varphi$ to the fixed parts of the Tate modules using the canonical
biextension $W^o=W$ of (\SGARefLink{sga7.sga7i.eix.n1-4-4}{1.4.4}).  This is the canonical extension of the
biextension $W_K$ of $(A_K;A'_K)$ by $\mathbb G_{mK}$ to a biextension
of $(A^o,A'{}^o)$ by $\mathbb G_{mS}$.  By VIII~\SGARefLink{sga7.sga7i.eviii.n2-2}{2.2}, $W$ gives rise to
a pairing
\begin{equation}\tag{\SGARefTarget{sga7.sga7i.eix.n2-4-1}{2.4.1}}
 \varphi_w:T_\ell(A^o)\times T_\ell(A'{}^o)
 \longrightarrow T_\ell(\mathbb G_{mS}),
\end{equation}
from which the pairing $\varphi$ of (\SGARefLink{sga7.sga7i.eix.n1-0-3}{1.0.3}) is obtained by passing to
$\overline K$-valued points.  It follows that the restriction of
$\varphi$ to the fixed parts considered in \SGARefLink{sga7.sga7i.eix.n2-4}{2.4} is obtained, by passing
to $\overline K$-valued points, from the pairing
\begin{equation}\tag{\SGARefTarget{sga7.sga7i.eix.n2-4-2}{2.4.2}}
 T_\ell(A^o)^f\times T_\ell(A'{}^o)^f
 \longrightarrow T_\ell(\mathbb G_{mS}),
\end{equation}
the restriction of (\SGARefLink{sga7.sga7i.eix.n2-4-1}{2.4.1}).  Since the two factors on the left of
(\SGARefLink{sga7.sga7i.eix.n2-4-2}{2.4.2}) are now twisted constant $\ell$-adic sheaves, this pairing may
be studied on the special fibers.  The fibers of the two factors are
precisely the Tate modules $T_\ell(A_o^o)$ and
$T_\ell(A'{}_o^o)$ (\SGARefLink{sga7.sga7i.eix.n2-2-3-3}{2.2.3.3}), and the pairing induced on special fibers
by (\SGARefLink{sga7.sga7i.eix.n2-4-2}{2.4.2}) is
\begin{equation}\tag{\SGARefTarget{sga7.sga7i.eix.n2-4-3}{2.4.3}}
 T_\ell(A_o^o)\times T_\ell(A'{}_o^o)
 \longrightarrow T_\ell(\mathbb G_{mk}),
\end{equation}
obtained from the biextension $W_o$ of $(A_o^o,A'{}_o^o)$ by
$\mathbb G_{mk}$, the restriction of $W$ to the special fibers.

Recall from \SGARefLink{sga7.sga7i.eix.n2-3}{2.3} that the toric part of
$T_\ell(A(\overline K))$ is obtained by passing to
$\overline K$-valued points in the twisted constant $\ell$-adic
subsheaf $T_\ell(A^o)^t$ of $T_\ell(A^o)^f$.  Thus the orthogonality
relation to be proved says that the pairing induced by (\SGARefLink{sga7.sga7i.eix.n2-4-2}{2.4.2}),
\begin{equation}\tag{\SGARefTarget{sga7.sga7i.eix.n2-4-4}{2.4.4}}
 T_\ell(A^o)^t\times T_\ell(A'{}^o)^f
 \longrightarrow T_\ell(\mathbb G_{mS}),
\end{equation}
is identically zero.  It is enough to verify this for the corresponding
pairing on the special fibers, which by the definition of the toric
part is
\begin{equation}\tag{\SGARefTarget{sga7.sga7i.eix.n2-4-5}{2.4.5}}
 T_\ell(T_o)\times T_\ell(A{}_o^o)
 \longrightarrow T_\ell(\mathbb G_{mk}),
\end{equation}
induced by (\SGARefLink{sga7.sga7i.eix.n2-4-3}{2.4.3}).  By VIII~\SGARefLink{sga7.sga7i.eviii.n4-10}{4.10} this pairing is necessarily zero,
which proves the asserted orthogonality.

\smallskip
\noindent
b) Let us prove that the inclusion
\begin{equation}\tag{\SGARefTarget{sga7.sga7i.eix.n2-4-6}{2.4.6}}
 T_\ell(A(\overline K))^t\subset
 T_\ell(A(\overline K))^f\cap
 \bigl(T_\ell(A'(\overline K))^f\bigr)^\perp
\end{equation}
is an equality, where $\perp$ denotes the orthogonal.  Equivalently, in
the pairing (\SGARefLink{sga7.sga7i.eix.n2-4-2}{2.4.2}), the annihilator of the second factor---which, as
we have just seen, contains the toric part $T_\ell(A^o)^t$---is equal
to that toric part.  Since these are again relations between twisted
constant $\ell$-adic sheaves, it is enough to check the assertion on
the special fibers: namely, that the inclusion already obtained (\SGARefLink{sga7.sga7i.eix.n2-4-7}{1.10})
\begin{equation}\tag{\SGARefTarget{sga7.sga7i.eix.n2-4-7}{2.4.7}}
 T_\ell(T_o)\subset T_\ell(A'{}_o^o)^\perp
\end{equation}
is an equality.  In other words, we must verify that the biextension
$W_o$ of $(A_o^o,A'{}_o^o)$ is \underline{nondegenerate on the left}
(VIII~\SGARefLink{sga7.sga7i.eviii.n4-11}{4.11}).  But this is precisely what we already observed in \SGARefLink{sga7.sga7i.eix.n1-5}{1.5},
as a consequence of a result from Raynaud's thesis.



\medskip
\noindent
\SGARefTarget{sga7.sga7i.eix.n2-5}{2.5}. The orthogonality theorem has an equivalent formulation using a
polarization $\xi$ of $A_K$ and the alternating bilinear form (\SGARefLink{sga7.sga7i.eix.n1-0-8}{1.0.8}),
compatible with the operations of $\pi$,
\begin{equation}\tag{\SGARefTarget{sga7.sga7i.eix.n2-5-1}{2.5.1}}
 \varphi_\xi:T_\ell(A(\overline K))\times T_\ell(A(\overline K))
 \longrightarrow T_\ell(\overline K^*),
\end{equation}
obtained from the form $\varphi$ of (\SGARefLink{sga7.sga7i.eix.n1-0-3}{1.0.3}) by transporting structure in
the second argument along the isogeny of $\mathbb Z_\ell$-modules
\[
 T_\ell(A(\overline K))\longrightarrow T_\ell(A'(\overline K))
\]
induced by the isogeny $\psi_\xi:A_K\longrightarrow A'_K$ defined by
$\xi$.  The equality in (\SGARefLink{sga7.sga7i.eix.n2-4-6}{2.4.6}) then takes the simple form
\begin{equation}\tag{\SGARefTarget{sga7.sga7i.eix.n2-5-2}{2.5.2}}
 W=V\cap V^\perp,
\end{equation}
where we have set
\begin{equation}\tag{\SGARefTarget{sga7.sga7i.eix.n2-5-3}{2.5.3}}
 \left\{
 \begin{array}{l}
 U=T_\ell(A(\overline K)),\qquad
 V=U^I=\text{the ``fixed part'' of }T_\ell(A(\overline K)),\\[2mm]
 \hspace{4em}W=\text{the toric part of }T_\ell(A(\overline K)),
 \end{array}
 \right.
\end{equation}
and $\perp$ denotes the orthogonal with respect to the polarization form
(\SGARefLink{sga7.sga7i.eix.n2-5-1}{2.5.1}).  These modules give the following diagram of inclusions:
\begin{equation}\tag{\SGARefTarget{sga7.sga7i.eix.n2-5-4}{2.5.4}}
\begin{tikzcd}[column sep=large,row sep=large]
 \&\& V\arrow[dl,hook',"2\lambda"'] \&\& \\
 U \& W^\perp=V+V^\perp\arrow[l,hook',"\mu"']
 \&\& W=V\cap V^\perp
    \arrow[ul,hook',"2\alpha"']
    \arrow[dl,hook,"2\lambda"]
 \& 0\arrow[l,hook',"\mu"'] \\
 \&\& V^\perp\arrow[ul,hook,"2\alpha"] \&\&
\end{tikzcd}
\qquad ,
\end{equation}
where the integers $2\alpha$, $2\lambda$, and $\mu$ are the ranks of
the quotients of consecutive modules in the diagram, with the notation
of (\SGARefLink{sga7.sga7i.eix.n2-1-9}{2.1.9}).  (That $W$ has rank $\mu$ is the definition of $\mu$; that
$V/W$ has rank $2\alpha$ follows from the second relation in (\SGARefLink{sga7.sga7i.eix.n2-1-11}{2.1.11})
together with the exact sequence (\SGARefLink{sga7.sga7i.eix.n2-1-6}{2.1.6}).  Hence, by orthogonality,
$\operatorname{rank}(U/W^\perp)=\mu$ and
$\operatorname{rank}(W^\perp/V^\perp)=2\alpha$.  Finally, the formulas
\begin{equation}\tag{\SGARefTarget{sga7.sga7i.eix.n2-5-5}{2.5.5}}
 \operatorname{rank}\bigl(V/(V\cap V^\perp)\bigr)
 =\operatorname{rank}(V/W)=2\lambda
\end{equation}
and the dual formula follow from the preceding ones, in view of
(\SGARefLink{sga7.sga7i.eix.n2-1-10}{2.1.10}) and $\operatorname{rank}U=n$.)

}{\smallskip
\noindent
\underline{Proof.}\quad a) Let us prove that
$T_\ell(A(\overline K))^t$ is orthogonal to
$T_\ell(A'(\overline K))^f$.  To do so, we interpret the restriction of
$\varphi$ to the fixed parts of the Tate modules using the canonical
biextension $W^o=W$ of (\SGARefLink{sga7.sga7i.eix.n1-4-4}{1.4.4}).  This is the canonical extension of the
biextension $W_K$ of $(A_K,A'_K)$ by $\mathbb G_{mK}$ to a biextension
of $(A^o,A'{}^o)$ by $\mathbb G_{mS}$.  By VIII~\SGARefLink{sga7.sga7i.eviii.n2-2}{2.2}, $W$ gives rise to
a pairing
\begin{equation}\tag{\SGARefTarget{sga7.sga7i.eix.n2-4-1}{2.4.1}}
 \varphi_W:T_\ell(A^o)\times T_\ell(A'{}^o)
 \longrightarrow T_\ell(\mathbb G_{mS}),
\end{equation}
from which the pairing $\varphi$ of (\SGARefLink{sga7.sga7i.eix.n1-0-3}{1.0.3}) is obtained by passing to
$\overline K$-valued points.  It follows that the restriction of
$\varphi$ to the fixed parts considered in \SGARefLink{sga7.sga7i.eix.n2-4}{2.4} is obtained, by passing
to $\overline K$-valued points, from the pairing
\begin{equation}\tag{\SGARefTarget{sga7.sga7i.eix.n2-4-2}{2.4.2}}
 T_\ell(A^o)^f\times T_\ell(A'{}^o)^f
 \longrightarrow T_\ell(\mathbb G_{mS}),
\end{equation}
the restriction of (\SGARefLink{sga7.sga7i.eix.n2-4-1}{2.4.1}).  Since the two factors on the left of
(\SGARefLink{sga7.sga7i.eix.n2-4-2}{2.4.2}) are now twisted constant $\ell$-adic sheaves, this pairing may
be studied on the special fibers.  The fibers of the two factors are
precisely the Tate modules $T_\ell(A_o^o)$ and
$T_\ell(A'{}_o^o)$ (\SGARefLink{sga7.sga7i.eix.n2-2-3-3}{2.2.3.3}), and the pairing induced on special fibers
by (\SGARefLink{sga7.sga7i.eix.n2-4-2}{2.4.2}) is
\begin{equation}\tag{\SGARefTarget{sga7.sga7i.eix.n2-4-3}{2.4.3}}
 T_\ell(A_o^o)\times T_\ell(A'{}_o^o)
 \longrightarrow T_\ell(\mathbb G_{mk}),
\end{equation}
obtained from the biextension $W_o$ of $(A_o^o,A'{}_o^o)$ by
$\mathbb G_{mk}$, the restriction of $W$ to the special fibers.

Recall from \SGARefLink{sga7.sga7i.eix.n2-3}{2.3} that the toric part of
$T_\ell(A(\overline K))$ is obtained by passing to
$\overline K$-valued points in the twisted constant $\ell$-adic
subsheaf $T_\ell(A^o)^t$ of $T_\ell(A^o)^f$.  Thus the orthogonality
relation to be proved says that the pairing induced by (\SGARefLink{sga7.sga7i.eix.n2-4-2}{2.4.2}),
\begin{equation}\tag{\SGARefTarget{sga7.sga7i.eix.n2-4-4}{2.4.4}}
 T_\ell(A^o)^t\times T_\ell(A'{}^o)^f
 \longrightarrow T_\ell(\mathbb G_{mS}),
\end{equation}
is identically zero.  It is enough to verify this for the corresponding
pairing on the special fibers, which by the definition of the toric
part is
\begin{equation}\tag{\SGARefTarget{sga7.sga7i.eix.n2-4-5}{2.4.5}}
 T_\ell(T_o)\times T_\ell(A'{}_o^o)
 \longrightarrow T_\ell(\mathbb G_{mk}),
\end{equation}
induced by (\SGARefLink{sga7.sga7i.eix.n2-4-3}{2.4.3}).  By VIII~\SGARefLink{sga7.sga7i.eviii.n4-10}{4.10} this pairing is necessarily zero,
which proves the asserted orthogonality.

\smallskip
\noindent
b) Let us prove that the inclusion
\begin{equation}\tag{\SGARefTarget{sga7.sga7i.eix.n2-4-6}{2.4.6}}
 T_\ell(A(\overline K))^t\subset
 T_\ell(A(\overline K))^f\cap
 \bigl(T_\ell(A'(\overline K))^f\bigr)^\perp
\end{equation}
is an equality, where $\perp$ denotes the orthogonal.  Equivalently, in
the pairing (\SGARefLink{sga7.sga7i.eix.n2-4-2}{2.4.2}), the annihilator of the second factor---which, as
we have just seen, contains the toric part $T_\ell(A^o)^t$---is equal
to that toric part.  Since these are again relations between twisted
constant $\ell$-adic sheaves, it is enough to check the assertion on
the special fibers: namely, that the inclusion already obtained (\SGARefLink{sga7.sga7i.eix.n2-4-7}{2.4.7})
\begin{equation}\tag{\SGARefTarget{sga7.sga7i.eix.n2-4-7}{2.4.7}}
 T_\ell(T_o)\subset T_\ell(A'{}_o^o)^\perp
\end{equation}
is an equality.  In other words, we must verify that the biextension
$W_o$ of $(A_o^o,A'{}_o^o)$ is \underline{nondegenerate on the left}
(VIII~\SGARefLink{sga7.sga7i.eviii.n4-11}{4.11}).  But this is precisely what we already observed in \SGARefLink{sga7.sga7i.eix.n1-5}{1.5},
as a consequence of a result from Raynaud's thesis.



\medskip
\noindent
\SGARefTarget{sga7.sga7i.eix.n2-5}{2.5}. The orthogonality theorem has an equivalent formulation using a
polarization $\xi$ of $A_K$ and the alternating bilinear form (\SGARefLink{sga7.sga7i.eix.n1-0-8}{1.0.8}),
compatible with the operations of $\pi$,
\begin{equation}\tag{\SGARefTarget{sga7.sga7i.eix.n2-5-1}{2.5.1}}
 \varphi_\xi:T_\ell(A(\overline K))\times T_\ell(A(\overline K))
 \longrightarrow T_\ell(\overline K^*),
\end{equation}
obtained from the form $\varphi$ of (\SGARefLink{sga7.sga7i.eix.n1-0-3}{1.0.3}) by transporting structure in
the second argument along the isogeny of $\mathbb Z_\ell$-modules
\[
 T_\ell(A(\overline K))\longrightarrow T_\ell(A'(\overline K))
\]
induced by the isogeny $\psi_\xi:A_K\longrightarrow A'_K$ defined by
$\xi$.  The equality in (\SGARefLink{sga7.sga7i.eix.n2-4-6}{2.4.6}) then takes the simple form
\begin{equation}\tag{\SGARefTarget{sga7.sga7i.eix.n2-5-2}{2.5.2}}
 W=V\cap V^\perp,
\end{equation}
where we have set
\begin{equation}\tag{\SGARefTarget{sga7.sga7i.eix.n2-5-3}{2.5.3}}
 \left\{
 \begin{array}{l}
 U=T_\ell(A(\overline K)),\qquad
 V=U^I=\text{the ``fixed part'' of }T_\ell(A(\overline K)),\\[2mm]
 \hspace{4em}W=\text{the toric part of }T_\ell(A(\overline K)),
 \end{array}
 \right.
\end{equation}
and $\perp$ denotes the orthogonal with respect to the polarization form
(\SGARefLink{sga7.sga7i.eix.n2-5-1}{2.5.1}).  These modules give the following diagram of inclusions:
\begin{equation}\tag{\SGARefTarget{sga7.sga7i.eix.n2-5-4}{2.5.4}}
\begin{tikzcd}[column sep=large,row sep=large]
 \&\& V\arrow[dl,hook',"2\lambda"'] \&\& \\
 U \& W^\perp=V+V^\perp\arrow[l,hook',"\mu"']
 \&\& W=V\cap V^\perp
    \arrow[ul,hook',"2\alpha"']
    \arrow[dl,hook,"2\lambda"]
 \& 0\arrow[l,hook',"\mu"'] \\
 \&\& V^\perp\arrow[ul,hook,"2\alpha"] \&\&
\end{tikzcd}
\qquad ,
\end{equation}
where the integers $2\alpha$, $2\lambda$, and $\mu$ are the ranks of
the quotients of consecutive modules in the diagram, with the notation
of (\SGARefLink{sga7.sga7i.eix.n2-1-9}{2.1.9}).  (That $W$ has rank $\mu$ is the definition of $\mu$; that
$V/W$ has rank $2\alpha$ follows from the second relation in (\SGARefLink{sga7.sga7i.eix.n2-1-11}{2.1.11})
together with the exact sequence (\SGARefLink{sga7.sga7i.eix.n2-1-6}{2.1.6}).  Hence, by orthogonality,
$\operatorname{rank}(U/W^\perp)=\mu$ and
$\operatorname{rank}(W^\perp/V^\perp)=2\alpha$.  Finally, the formulas
\begin{equation}\tag{\SGARefTarget{sga7.sga7i.eix.n2-5-5}{2.5.5}}
 \operatorname{rank}\bigl(V^\perp/(V\cap V^\perp)\bigr)
 =\operatorname{rank}(V^\perp/W)=2\lambda
\end{equation}
and the dual formula follow from the preceding ones, in view of
(\SGARefLink{sga7.sga7i.eix.n2-1-10}{2.1.10}) and $\operatorname{rank}U=2n$.)

}{E7I-IX-008}\SGApage{I}{\(\approx\)341}
\medskip
\noindent
\underline{Remarks} \SGARefTarget{sga7.sga7i.eix.n2-6}{2.6}.\quad a) Here we obtained the orthogonality
theorem \SGARefLink{sga7.sga7i.eix.n2-4}{2.4} as an easy consequence of the theory of biextensions
developed in the preceding sections.  The \underline{inclusion}
(\SGARefLink{sga7.sga7i.eix.n2-4-6}{2.4.6}) can be proved without using the theory of biextensions, by an
arithmetic argument in the style of the proof of the ``monodromy
theorem'' given in Expos\'e I.  One reduces to a situation in which $S$
has finite residue field, where the orthogonality result follows from a
simple consideration of \underline{weights} for the Frobenius operation
(using A.~Weil's result on the eigenvalues of that Frobenius).  There is
no need to invoke resolution of singularities here: one may state the
desired result for group schemes $A$ over schemes $S$ under sufficiently
general conditions.  Unfortunately, it does not appear possible to
obtain by this method equality in the preceding inclusion.  That
equality will apparently be used in an essential way in the proof of
the ``semistable reduction theorem'' in the next section.

\SGApage{I}{\(\approx\)342}
\smallskip
\noindent
b) The orthogonality theorem \SGARefLink{sga7.sga7i.eix.n2-4}{2.4} allows the most important properties
of the N\'{e}ron model of an abelian variety $A_K$ to be expressed in
terms of the representation of the inertia group $I$ on the Tate module
$T_\ell(A(\overline K))$, where $\ell$ is a fixed prime different from
the residual characteristic.  (The case $\ell=p$ will be considered
separately later, in section 5.)  As a first statement of this kind,
though one that does not use the orthogonality theorem, recall the
good-reduction criterion \SGARefLink{sga7.sga7i.eix.n2-2-9}{2.2.9}.  We shall give many further examples
below, beginning in the next section (\SGARefLink{sga7.sga7i.eix.n3-5}{3.5}).



\SGApage{I}{\(\approx\)343}
\medskip
\noindent
3. \underline{The case of semistable reduction.  The semistable reduction
theorem.}

We first give an auxiliary result in the style of SGA 3.

\SGAvar{\medskip
\noindent
\underline{Proposition} \SGARefTarget{sga7.sga7i.eix.n3-1}{3.1}. \underline{Let $S$ be a scheme, let $G$ be a
commutative $S$-group scheme that is flat, separated, and of finite
presentation over $S$, let $H$ be a quasi-separated group scheme locally of
finite type over $S$, let $u:G\longrightarrow H$ be a homomorphism of
$S$-group schemes, and put $K=\ker u$.  For every $s\in S$, let $\bar s$
denote the spectrum of an algebraic closure of $k(s)$, and denote by
$\mu(s)$ (respectively, $\alpha(s)$) the reductive rank (respectively, the
abelian rank) of $K_s$, that is, the dimension of the largest subtorus
(respectively, of the largest abelian-variety quotient) of
$(K_{\bar s})^o_{\mathrm{red}}$.  Let $s\in S$, let $t\in S$ be a
generalization of $S$, and let $p$ be the characteristic exponent of
$k(s)$.}

}{\medskip
\noindent
\underline{Proposition} \SGARefTarget{sga7.sga7i.eix.n3-1}{3.1}. \underline{Let $S$ be a scheme, let $G$ be a
commutative $S$-group scheme that is flat, separated, and of finite
presentation over $S$, let $H$ be a quasi-separated group scheme locally of
finite type over $S$, let $u:G\longrightarrow H$ be a homomorphism of
$S$-group schemes, and put $K=\ker u$.  For every $s\in S$, let $\bar s$
denote the spectrum of an algebraic closure of $k(s)$, and denote by
$\mu(s)$ (respectively, $\alpha(s)$) the reductive rank (respectively, the
abelian rank) of $K_s$, that is, the dimension of the largest subtorus
(respectively, of the largest abelian-variety quotient) of
$(K_{\bar s})^o_{\mathrm{red}}$.  Let $s\in S$, let $t\in S$ be a
generalization of $s$, and let $p$ be the characteristic exponent of
$k(s)$.}

}{E7I-IX-009}\smallskip
\noindent
a) \underline{For every integer $n>0$ prime to $p$, ${}_nK$ is etale, of
finite type, and separated over $S$ in a neighborhood of $s$, and one has
the divisibility relation}
\begin{equation}\tag{\SGARefTarget{sga7.sga7i.eix.n3-1-1}{3.1.1}}
 \operatorname{rank}({}_nK_s)\mid \operatorname{rank}({}_nK_t).
\end{equation}

\noindent
b) \underline{One has}
\begin{equation}\tag{\SGARefTarget{sga7.sga7i.eix.n3-1-2}{3.1.2}}
 \mu(s)+2\alpha(s)\leq \mu(t)+2\alpha(t).
\end{equation}

\noindent
c) \underline{If $(K_{\bar t})^o_{\mathrm{red}}$ is unipotent, then so is
$(K_{\bar s})^o_{\mathrm{red}}$.}

\smallskip
\noindent
d) \underline{Suppose that $G_s$ has unipotent rank zero, that $H$ is smooth
over $S$, and that $u_t$ is an isogeny.  Then $u$ is flat and quasi-finite
over an open neighborhood $U$ of $s$; a fortiori, $K|_U$ is flat and
quasi-finite.}

\smallskip
\noindent
e) \underline{Under the hypotheses of d), suppose in addition that $u_t$ is
a monomorphism (and hence an isomorphism from $G_t$ onto an open subgroup of
$H_t$).  Then $u_s$ is an isomorphism from $G_s$ onto an open subgroup of
$H_s$.  If, moreover, $S$ is irreducible with generic point $t$, then there
is an open neighborhood $U$ of $s$ such that $u$ induces an isomorphism from
$G|_U$ onto an open subgroup of $H|_U$.}

\SGApage{I}{\(\approx\)344}
\medskip
\noindent
\underline{Proof.}

\smallskip
\noindent
a) We may suppose that $n$ is invertible on $S$.  Then $n\,\mathrm{id}_G$
is etale on every fibre (SGA 3, XV, 1.3); since $G$ is flat and of finite
presentation over $S$, the fibrewise criterion for etaleness (EGA IV,
17.8.2) shows that $n\,\mathrm{id}_G$ is etale.  Its kernel ${}_nG$ is
therefore etale; it is also separated and of finite presentation over $S$,
since $G$ is.  Similarly, ${}_nH$ is locally of finite type,
quasi-separated, and formally unramified over $S$,
\footnote{At least when $H$ is commutative.  In the general case, ${}_nH$
is still unramified over $S$ along its unit section, which is all that is
needed here.}
so its unit section is an open immersion (EGA IV, 17.4.2 b')).  Consequently,
the kernel of
\[
 {}_nu:{}_nG\longrightarrow {}_nH
\]
is an open subscheme of ${}_nG$, hence is etale over $S$; it is
quasi-compact in ${}_nG$ and therefore of finite presentation over $S$, and
it is separated over $S$ because $G$ is.  This kernel is plainly ${}_nK$.

Formula (\SGARefLink{sga7.sga7i.eix.n3-1-1}{3.1.1}) follows from these properties of ${}_nK$.  Indeed, after a
base change we may reduce to the case where $S$ is henselian local with
closed point $s$, and then use the decomposition of ${}_nK$ into
${}_nK^f$ and a scheme with empty special fibre (cf. \SGARefLink{sga7.sga7i.eix.n2-2-3}{2.2.3}).  More precisely,
the same argument gives a monomorphism on geometric fibres
\[
 {}_nK(\bar s)\lhook\joinrel\longrightarrow {}_nK(\bar t).
\]

\SGApage{I}{\(\approx\)345}
\smallskip
\noindent
b) Choose an integer $n>0$ prime to the characteristic exponent of $k(s)$
and to the ranks of the quotients of the geometric fibres of $K$ at $s$ and
$t$ by their reduced identity components.  The two sides of (\SGARefLink{sga7.sga7i.eix.n3-1-1}{3.1.1}) are
then obtained by raising $n$, respectively, to the powers whose exponents
are the two sides of (\SGARefLink{sga7.sga7i.eix.n3-1-2}{3.1.2}), as follows from the known structure of
commutative algebraic groups.  Thus b) follows from a).

\SGAvar{\smallskip
\noindent
c) The hypothesis that $(K_{\bar t})^o_{\mathrm{s}}$ is unipotent is
equivalent to the equalities $\mu(t)=\alpha(t)=0$, and likewise for
$K_{\bar s}$.  Assertion c) therefore follows formally from b).

}{\smallskip
\noindent
c) The hypothesis that $(K_{\bar t})^o_{\mathrm{red}}$ is unipotent is
equivalent to the equalities $\mu(t)=\alpha(t)=0$, and likewise for
$K_{\bar s}$.  Assertion c) therefore follows formally from b).

}{E7I-IX-010}\smallskip
\noindent
d) Since $u_t$ is an isogeny, $K_t$ is finite.  We are therefore under the
hypotheses of c), and $(K_{\bar s})^o_{\mathrm{red}}$ is unipotent.  As
$G_s$ has unipotent rank zero, it follows that $K_s$ is finite.  Moreover,
the fact that $u_t$ is an isogeny gives
$\dim G_t=\dim H_t$.  Since $G$ and $H$ are flat and locally of finite
presentation, one has
\[
 \dim G_s=\dim G_t,\qquad \dim H_s=\dim H_t
 \quad\text{(SGA 3, VI$_B$, 4.3)},
\]
and hence $\dim G_s=\dim H_s$.  Since $\dim\ker u_s=0$, it follows that
$\dim H_s=\dim u_s(G_s)$ (SGA 3, VI$_B$, 1.2), which, because $H_s$ is
smooth, is equivalent to $u_s$ being flat (SGA 3, VI$_A$, 5.6).  There is
therefore an open neighborhood $U$ of $s$ such that $u$ is flat and
quasi-finite at the points of the unit section over $U$ (SGA 3, VI$_B$,
2.2(i) and \SGARefLink{sga7.sga7i.eix.n2-5}{2.5}(ii)).  As for $u_s$, it follows that $u_{s'}$ is flat for
every $s'\in U$; hence $u|_{G_U}$ is flat and quasi-finite, since it has
these properties on every fibre.

\smallskip
\noindent
e) For the first assertion, a base change immediately reduces us to the
case where $S$ is integral with generic point $t$, and hence to proving the
second assertion of e).  That assertion follows at once from d) and the
following lemma.

\SGApage{I}{346}
\SGAvar{\medskip
\noindent
\underline{Lemma} \SGARefTarget{sga7.sga7i.eix.n3-1-3}{3.1.3}. \underline{Let $S$ be a scheme, let $K$ be a
quasi-finite $S$-scheme of finite presentation that is separated and flat
over $S$, let $s\in S$, and let $t\in T$ be a generalization of $s$.  Then}
\[
 \operatorname{rank}K_s\leq \operatorname{rank}K_t.
\]
\underline{If $\operatorname{rank}$} $K_\eta\leq1$ \underline{for every maximal point
$\eta$ of} $S$, \underline{then $K\longrightarrow S$ is an open immersion (and hence
an isomorphism if it admits a section).}

}{\medskip
\noindent
\underline{Lemma} \SGARefTarget{sga7.sga7i.eix.n3-1-3}{3.1.3}. \underline{Let $S$ be a scheme, let $K$ be a
quasi-finite $S$-scheme of finite presentation that is separated and flat
over $S$, let $s\in S$, and let $t\in S$ be a generalization of $s$.  Then}
\[
 \operatorname{rank}K_s\leq \operatorname{rank}K_t.
\]
\underline{If $\operatorname{rank}$} $K_\eta\leq1$ \underline{for every maximal point
$\eta$ of} $S$, \underline{then $K\longrightarrow S$ is an open immersion (and hence
an isomorphism if it admits a section).}

}{E7I-IX-011}For the first assertion, strict localization reduces us to the case where
$S$ is strictly local.  Decomposing $K$ as the disjoint union of a finite
scheme $K^f$ and a scheme with empty special fibre, we obtain
\[
 \operatorname{rank}K_s
 =\operatorname{rank}K^f_s
 =\operatorname{rank}K^f_t
 \leq\operatorname{rank}K_t.
\]
The second equality follows because $K^f$ is flat, finite, and of finite
presentation over $S$, and therefore locally free over $S$.

For the second assertion, the preceding argument already shows that
$\operatorname{rank}K_s\leq1$ for every $s\in S$.  Thus $K_s$ is empty or
isomorphic to $\operatorname{Spec}k(s)$.  Consequently,
$K\longrightarrow S$ is a monomorphism (SGA 3, VI$_B$, 2.11), and, being
flat and of finite presentation, it is an open immersion (EGA IV, 17.9.1).
This proves the lemma.

\SGAvar{\medskip
\noindent
\underline{Remark} \SGARefTarget{sga7.sga7i.eix.n3-1-4}{3.1.4}.  Part \SGARefLink{sga7.sga7i.eix.n3-1}{3.1} a), and the remainder of Proposition
\SGARefLink{sga7.sga7i.eix.n3-1}{3.1}, give no information about ${}_nK_s$ when $n$ is a power of $p$ (for
$p>1$).  At least when $G$ is smooth (and $S$ is locally noetherian or $H$
is of finite presentation over $S$), essentially the same method as in the
proof of a) gives such information.  Indeed, reducing to the case where
$S$ is complete local noetherian and using SGA 3, XV, 1.6 c), X, 3.2, IX,
\SGARefLink{sga7.sga7i.eix.n3-6}{3.6}, and IX, \SGARefLink{sga7.sga7i.eix.n5-2}{5.2}, one obtains
\begin{equation}\tag{\SGARefTarget{sga7.sga7i.eix.n3-1-4-1}{3.1.4.1}}
 \operatorname{rank}_{\mathrm{mult}}({}_nK_s)
 \mid
 \operatorname{rank}_{\mathrm{mult}}({}_nK_t),
\end{equation}
where $\operatorname{rank}_{\mathrm{mult}}$ denotes the rank of the largest
subgroup of multiplicative type of the commutative algebraic group under
consideration (SGA 3, XVII, 7.2.1).  Notice that a torsion algebraic group
of multiplicative type is finite.  Formula (\SGARefLink{sga7.sga7i.eix.n3-1-4-1}{3.1.4.1}) is valid for every
integer $n\geq1$.  It follows from SGA , XVII, 4.6.1(vi) that, in c), if
$K_t$ is unipotent (SGA 3, VII, 1.1, then $K_s$ is unipotent as well.



}{\medskip
\noindent
\underline{Remark} \SGARefTarget{sga7.sga7i.eix.n3-1-4}{3.1.4}.  Part \SGARefLink{sga7.sga7i.eix.n3-1}{3.1} a), and the remainder of Proposition
\SGARefLink{sga7.sga7i.eix.n3-1}{3.1}, give no information about ${}_nK_s$ when $n$ is a power of $p$ (for
$p>1$).  At least when $G$ is smooth (and $S$ is locally noetherian or $H$
is of finite presentation over $S$), essentially the same method as in the
proof of a) gives such information.  Indeed, reducing to the case where
$S$ is complete local noetherian and using SGA 3, XV, 1.6 c), X, 3.2, IX,
\SGARefLink{sga7.sga7i.eix.n3-6}{3.6}, and IX, \SGARefLink{sga7.sga7i.eix.n5-2}{5.2}, one obtains
\begin{equation}\tag{\SGARefTarget{sga7.sga7i.eix.n3-1-4-1}{3.1.4.1}}
 \operatorname{rank}_{\mathrm{mult}}({}_nK_s)
 \mid
 \operatorname{rank}_{\mathrm{mult}}({}_nK_t),
\end{equation}
where $\operatorname{rank}_{\mathrm{mult}}$ denotes the rank of the largest
subgroup of multiplicative type of the commutative algebraic group under
consideration (SGA 3, XVII, 7.2.1).  Notice that a torsion algebraic group
of multiplicative type is finite.  Formula (\SGARefLink{sga7.sga7i.eix.n3-1-4-1}{3.1.4.1}) is valid for every
integer $n\geq1$.  It follows from SGA 3, XVII, 4.6.1(vi) that, in c), if
$K_t$ is unipotent (SGA 3, XVII, 1.1), then $K_s$ is unipotent as well.



}{E7I-IX-012}\SGApage{I}{347}
\medskip
\noindent
\underline{Proposition} \SGARefTarget{sga7.sga7i.eix.n3-2}{3.2}. \underline{Let $S$ be a regular integral
locally noetherian scheme of dimension} $1$, \underline{let $K$ be its field of rational
functions, let $A_K$ be an abelian scheme over $K$, and let $A$ be its Neron
model over $S$, which is a smooth separated group scheme over $S$} (\underline{and is
quasi-projective over $S$ if $S$ is noetherian [23, VIII, 2, 3\textdegree]).
The following conditions are equivalent:}

\smallskip
\noindent
(i) \underline{For every $s\in S$, $A_s^o$ has unipotent rank zero; that is,
it is an extension of an abelian scheme by a torus (\SGARefLink{sga7.sga7i.eix.n2-1-5}{2.1.5}).}

\smallskip
\noindent
(ii) \underline{There exists a smooth separated group scheme $G$ of finite
type over $S$, extending $A_K$, whose fibres have unipotent rank zero.}

\smallskip
\noindent
\underline{Moreover, if $G$ is as in (ii), then the unique morphism
$G\longrightarrow A$ inducing the given isomorphism on the generic fibres
is an isomorphism from $G$ onto an open subgroup of $A$; it therefore induces
an isomorphism of identity components}
\[
 G^o\xrightarrow{\ \sim\ }A^o.
\]

\SGApage{I}{348}
\medskip
\noindent
\SGARefTarget{sga7.sga7i.eix.n3-2-1}{3.2.1}.  Let us note at once that the definition of the Neron model $A$ of
$A_K$ is the same as in the local case recalled in \SGARefLink{sga7.sga7i.eix.n1-1}{1.1}.  The existence of
this Neron model, and the fact that it is separated and of finite type over
$S$, are well known and in fact reduce immediately to the local case.  Of
course, for every closed point $s\in S$,
\[
 S'=\operatorname{Spec}(\mathcal O_{S,s})
\]
is a trait, and $A\times_S S'$ is precisely the Neron model of $A_K$ over
$S'$ in the sense considered in \SGARefLink{sga7.sga7i.eix.n1-1}{1.1}.

Let us prove \SGARefLink{sga7.sga7i.eix.n3-2}{3.2}.  It is trivial that (i) implies (ii), since it suffices to
take $G=A$.  It remains to prove that if $G$ is as in (ii), then
$G\longrightarrow A$ is an open immersion.  This is indeed a special case
of \SGARefLink{sga7.sga7i.eix.n3-1}{3.1} e).

\medskip
\noindent
\underline{Corollary} \SGARefTarget{sga7.sga7i.eix.n3-3}{3.3}. \underline{With the notation of \SGARefLink{sga7.sga7i.eix.n3-2}{3.2}, suppose
that the identity components of the fibres of the Neron model have
unipotent rank zero.  Let $S'$ be a scheme satisfying the same hypotheses
as $S$, let $K'$ be its field of rational functions, let
$f:S'\longrightarrow S$ be a dominant morphism, let}
\[
 A_{K'}\simeq A_K\otimes_KK'
 \simeq \text{the generic fibre of }A\times_SS',
\]
\underline{let $A'$ be the Neron model of $A_{K'}$, and consider the unique
morphism}
\[
 u:A\times_SS'\longrightarrow A'
\]
\underline{inducing the identity on the generic fibres.  Then $u$ is an
open immersion and therefore induces an isomorphism on identity
components}
\begin{equation}\tag{\SGARefTarget{sga7.sga7i.eix.n3-3-1}{3.3.1}}
 A^o\times_SS'\xrightarrow{\ \sim\ }{A'}^o.
\end{equation}
Indeed, it suffices to apply \SGARefLink{sga7.sga7i.eix.n3-2}{3.2} over $S'$, taking
$G\simeq A\times_SS'$.

\SGApage{I}{349}
\medskip
\noindent
\SGARefTarget{sga7.sga7i.eix.n3-3-2}{3.3.2}.  Thus the equivalent properties of \SGARefLink{sga7.sga7i.eix.n3-2}{3.2} imply a partial
``stability'' property for Neron models.  One should bear in mind that when
$A_K$ does not have good reduction, that is, when $A$ is not an abelian
scheme, it is not true that, for every $S'$ as in \SGARefLink{sga7.sga7i.eix.n3-3}{3.3} (even for $S'$ finite
over $S$ and $K'/K$ separable), the morphism
\[
 A\times_SS'\longrightarrow A'
\]
itself is an isomorphism, or that it induces isomorphisms on the groups of
connected components of the fibres (see \SGARefLink{sga7.sga7i.eix.n11-10}{11.10} below for the behavior of
these groups of connected components under extension of the base).  These
observations, reinforced by \SGARefLink{sga7.sga7i.eix.n3-9}{3.9} below, justify the following terminology,
introduced by Mumford (and which also appears to agree with the notions of
the same name studied in his book [17]).

\medskip
\noindent
\underline{Definition} \SGARefTarget{sga7.sga7i.eix.n3-4}{3.4}. \underline{With the notation of \SGARefLink{sga7.sga7i.eix.n3-2}{3.2}, the
abelian scheme $A_K$ over $K$ is said to have semistable reduction over $S$
if the equivalent conditions (i) and (ii) of \SGARefLink{sga7.sga7i.eix.n3-2}{3.2} are satisfied.}

\medskip
\noindent
\SGARefTarget{sga7.sga7i.eix.n3-4-0}{3.4.0}.  Notice that in the form (ii), this condition does not invoke the
notion of a Neron model.  Moreover, by localization it reduces immediately
to the case where $S$ is a trait, as follows for example from \SGARefLink{sga7.sga7i.eix.n3-2-1}{3.2.1}.
Furthermore, since formation of the Neron model commutes with base change
from $S$ to its henselization or its strict henselization, the notion even
reduces to the case where $S$ is a strictly local trait.



\SGApage{I}{350}
\medskip
\noindent
\underline{Proposition} \SGARefTarget{sga7.sga7i.eix.n3-5}{3.5}. (\underline{Galois criterion for semistable
reduction}).  \underline{Suppose that $S$ is a trait.  The following conditions on the
abelian scheme $A_K$ are equivalent} (the notation being that of (\SGARefLink{sga7.sga7i.eix.n2-5-3}{2.5.3}),
applied to the situation deduced from $(A_K,S)$ by passing to the
henselization of $S$):

\smallskip
\noindent
(i) \underline{$A_K$ has semistable reduction over $S$} (\SGARefLink{sga7.sga7i.eix.n3-4}{3.4}).

\SGAvar{\smallskip
\noindent
(ii) \underline{With the notation of (\SGARefLink{sga7.sga7i.eix.n2-5-3}{2.5.3}) and (\SGARefLink{sga7.sga7i.eix.n2-5-4}{2.5.4}), one has}
\[
 V \subset V^\perp.
\]

}{\smallskip
\noindent
(ii) \underline{With the notation of (\SGARefLink{sga7.sga7i.eix.n2-5-3}{2.5.3}) and (\SGARefLink{sga7.sga7i.eix.n2-5-4}{2.5.4}), one has}
\[
 V^\perp\subset V.
\]

}{E7I-IX-013}\noindent
(iii) \underline{For the toric part $W$ of the Tate module one has}
\[
 W=V^\perp.
\]

\noindent
(iv) \underline{The action of the inertia group $I$ on
$U=T_\ell(A(\overline K))$ is unipotent of level $2$} (that is, there is an
$I$-stable submodule $U'$ of $U$ such that $I$ acts trivially on $U'$ and
on $U/U'$).

\SGAvar{We may suppose that $S$ is henselian (\SGARefLink{sga7.sga7i.eix.n3-4-0}{3.4.0}).  With the notation of \SGARefLink{sga7.sga7i.eix.n2-1-9}{2.1.9},
(i) means that
\[
 \lambda=0.
\]
By the rank formulas following (\SGARefLink{sga7.sga7i.eix.n2-5-4}{2.5.4}), this is equivalent to
$V=V\cap V^\perp$, that is, to (ii).  Since
$W=V\cap V^\perp$ by the orthogonality theorem \SGARefLink{sga7.sga7i.eix.n2-4}{2.4}, conditions (ii) and
(iii) are equivalent.  Moreover, (iv) also says that $I$ acts trivially
on $U^I=V$.  By duality this means that it acts trivially on
$V^\perp$, that is, that $V^\perp\subset U^I=V$, which is (ii).  This
completes the proof.

}{We may suppose that $S$ is henselian (\SGARefLink{sga7.sga7i.eix.n3-4-0}{3.4.0}).  With the notation of \SGARefLink{sga7.sga7i.eix.n2-1-9}{2.1.9},
(i) means that
\[
 \lambda=0.
\]
By the rank formulas following (\SGARefLink{sga7.sga7i.eix.n2-5-4}{2.5.4}), this is equivalent to
$V^\perp=V\cap V^\perp$, that is, to (ii).  Since
$W=V\cap V^\perp$ by the orthogonality theorem \SGARefLink{sga7.sga7i.eix.n2-4}{2.4}, conditions (ii) and
(iii) are equivalent.  Moreover, (iv) also says that $I$ acts trivially
on $U/U^I=U/V$.  By duality this means that it acts trivially on
$V^\perp$, that is, that $V^\perp\subset U^I=V$, which is (ii).  This
completes the proof.

}{E7I-IX-014}\medskip
\noindent
\underline{Remarks} \SGARefTarget{sga7.sga7i.eix.n3-5-1}{3.5.1}.  Conditions (ii) and (iv) plainly depend only
on the action of $I$ on the $\mathbb Q_\ell$-vector space
$U\otimes_{\mathbb Z_\ell}\mathbb Q_\ell$, and could have been formulated
in terms of that action.  Let us also note that these conditions are
manifestly stable under every base change by a morphism of traits
$S'\longrightarrow S$.  Finally, if (iv) holds for $A_K$, it plainly holds
for every abelian subscheme and every quotient abelian scheme of $A_K$,
and in particular for every abelian scheme isogenous to $A_K$ (for
example, for $A'_K$).

\SGApage{I}{351}
\medskip
\noindent
\underline{Corollary} \SGARefTarget{sga7.sga7i.eix.n3-5-2}{3.5.2}. \underline{Let $g$ be a topological generator
of $\mathbb Z_\ell(1)=T_\ell(\overline k^{\,*})$.  The equivalent conditions
of \SGARefLink{sga7.sga7i.eix.n3-5}{3.5} are also equivalent to the following condition:}

\SGAvar{\smallskip
\noindent
(v) \underline{The action of the inertia group $I$ on
$U=T_\ell(A(\overline K))$ factors through the maximal pro-$\ell$ quotient
$I(\ell)$ of $I$, which is canonically isomorphic to $\mathbb Z_\ell(1)$ $(I\qquad)$,
and the action $g_U$ of the generator $g$ on $U$ satisfies}
\begin{equation}\tag{\SGARefTarget{sga7.sga7i.eix.n3-5-3}{3.5.3}}
 (1-g_U)^2=0.
\end{equation}
Indeed, for every continuous unipotent action of a profinite group $I$ on
an $\ell$-adic vector space, one sees at once that the image of $I$ is a
pro-$\ell$ group.  It follows that (v) is equivalent to (iv).

}{\smallskip
\noindent
(v) \underline{The action of the inertia group $I$ on
$U=T_\ell(A(\overline K))$ factors through the maximal pro-$\ell$ quotient
$I(\ell)$ of $I$, which is canonically isomorphic to $\mathbb Z_\ell(1)$,
and the action $g_U$ of the generator $g$ on $U$ satisfies}
\begin{equation}\tag{\SGARefTarget{sga7.sga7i.eix.n3-5-3}{3.5.3}}
 (1-g_U)^2=0.
\end{equation}
Indeed, for every continuous unipotent action of a profinite group $I$ on
an $\ell$-adic vector space, one sees at once that the image of $I$ is a
pro-$\ell$ group.  It follows that (v) is equivalent to (iv).

}{E7I-IX-015}We can now prove the central result of the present expose:

\medskip
\noindent
\underline{Theorem} \SGARefTarget{sga7.sga7i.eix.n3-6}{3.6}. \underline{(Semistable reduction theorem).  Let
$S$ be a connected regular noetherian scheme of dimension $1$, let $K$ be
the field of rational functions of $S$, and let $A_K$ be an abelian scheme
over $K$.  Then there is a finite Galois extension $K'$ of $K$ such that
the abelian scheme $A_{K'}=A_K\otimes_KK'$ has semistable reduction} (\SGARefLink{sga7.sga7i.eix.n3-4}{3.4})
\underline{over the normalization $S'$ of $S$ in $K'$.}

\SGApage{I}{\(\approx\)352}
Notice that $S'$ is connected, regular, noetherian, and of dimension $1$,
so that the stated conclusion makes sense.  To prove \SGARefLink{sga7.sga7i.eix.n3-6}{3.6}, choose a
separable closure $\overline K$ of $K$ and write $\overline K$ as the
filtered union of its finite Galois subextensions $K_i$.  Let $S_i$ be the
normalization of $S$ in $K_i$, and let $U$ be the largest open subset of
$S$ over which $A_K$ has good reduction.  Since $S$ is connected,
noetherian, and one-dimensional, and $U\ne\varnothing$, the complement
$T=S-U$ is a finite set of closed points of $S$, and
$\mathcal O_{S,t}$ is a discrete valuation ring for every $t\in T$.  It is
enough to find, for every $t\in T$, an index $i=i(t)$ such that $A_{K_i}$
has semistable reduction at the points of $S_i$ above $t$: one then takes
for $K'$ the compositum of the extensions $K_{i(t)}$, for $t\in T$.
Thus we are reduced to the case where $S$ is a trait.

Moreover, since the Galois group permutes the points of the normalization
lying above the closed point $s$ of $S$, it is enough to find $K'$ such
that $A_{K'}$ has semistable reduction at \emph{one} point of the
normalization $S'$.  Here it is unnecessary at first to require that
$K'/K$ be Galois: it is enough to find a finite etale extension $K'/K$
satisfying the preceding condition, since its Galois closure satisfies the
same condition.  Using criterion \SGARefLink{sga7.sga7i.eix.n3-5}{3.5}(iv), we see that Theorem \SGARefLink{sga7.sga7i.eix.n3-6}{3.6} reduces
to the following assertion.

\medskip
\noindent
\underline{Corollary} \SGARefTarget{sga7.sga7i.eix.n3-7}{3.7}. \underline{Suppose again that $S$ is a trait.
Then there is an open subgroup $\pi'$ of $\pi$ such that the action of
$I'=I\cap\pi'$ on $U=T_\ell(A(\overline K))$ is unipotent of level $2$;
equivalently, there is an open subgroup $I'$ of $I$ whose action on $U$ is
unipotent of level $2$.}

\SGApage{I}{\(\approx\)353}
There is a canonical isomorphism, compatible with the actions of $\pi$,
\begin{equation}\tag{*}
 H^1(A_{\overline K},\mathbb Z_\ell)
 \simeq
 \operatorname{Hom}\bigl(T_\ell(A(\overline K)),\mathbb Z_\ell\bigr),
\end{equation}
which shows that assertion \SGARefLink{sga7.sga7i.eix.n3-7}{3.7} is equivalent to the same assertion for
the actions of $I$ on the left-hand side of (*).  But that assertion was
proved in Expose III (using resolution of singularities for excellent
schemes of dimension $2$).

\medskip
\noindent
\underline{Corollary} \SGARefTarget{sga7.sga7i.eix.n3-8}{3.8}. \underline{The equivalent conditions (i)--(iv)
of \SGARefLink{sga7.sga7i.eix.n3-5}{3.5} are also equivalent to the following condition:}

\smallskip
\noindent
(vi) \quad \underline{The action of $I$ on
$U=T_\ell(A(\overline K))$ is unipotent.}

\SGAvar{Indeed, to prove that (v) implies (vi), observe that if a group $I$
acting unipotently on a finite-dimensional vector space $E$ over a field
of characteristic zero (here $\mathbb Q_\ell$) has a subgroup $I'$ of
finite index acting with level $m$ (here $m=2$), then $I$ itself acts with
level $m$.  In fact, since $I'$ has finite index in $I$, the algebraic
group $G'\subset\operatorname{GL}(E)$ that it defines has finite index in
the group $G$ defined by $I$; but $G$ is unipotent and hence connected (in
characteristic zero), so $G=G'$.  Thus $I$ and $I'$ have the same
``level of unipotence.''

}{Indeed, to prove that (vi) implies (iv), observe that if a group $I$
acting unipotently on a finite-dimensional vector space $E$ over a field
of characteristic zero (here $\mathbb Q_\ell$) has a subgroup $I'$ of
finite index acting with level $m$ (here $m=2$), then $I$ itself acts with
level $m$.  In fact, since $I'$ has finite index in $I$, the algebraic
group $G'\subset\operatorname{GL}(E)$ that it defines has finite index in
the group $G$ defined by $I$; but $G$ is unipotent and hence connected (in
characteristic zero), so $G=G'$.  Thus $I$ and $I'$ have the same
``level of unipotence.''

}{E7I-IX-016}We also obtain a converse to \SGARefLink{sga7.sga7i.eix.n3-3}{3.3}:

\medskip
\noindent
\underline{Corollary} \SGARefTarget{sga7.sga7i.eix.n3-9}{3.9}. \underline{The notation is that of \SGARefLink{sga7.sga7i.eix.n3-2}{3.2}.  The
conditions (i) and (ii) of \SGARefLink{sga7.sga7i.eix.n3-2}{3.2} are also equivalent to the following
condition:}

\smallskip
\noindent
(iii) \quad \underline{For every finite etale extension $K'$ of $K$, let
$S'$ denote the normalization of $S$ in $K'$, and let $A'$ denote the
Neron model of
$A_{K'}=A\otimes_SK'=A_K\otimes_KK'$.  Then the canonical morphism}
\[
 A\times_SS'\longrightarrow A'
\]
\underline{induces an isomorphism on identity components.}

\SGApage{I}{\(\approx\)354}
The necessity of the condition was proved, in a stronger form, in \SGARefLink{sga7.sga7i.eix.n3-3}{3.3}.
For sufficiency, choose $K'$ as in \SGARefLink{sga7.sga7i.eix.n3-6}{3.6}.  Since the fibres of ${A'}^o$ have
unipotent rank zero, the same is true of the fibres of
$(A\times_SS')^o=A^o\times_SS'$, and hence also of the fibres of $A^o$,
from which the former are obtained by extension of the base field.



\medskip
\noindent
\underline{4. Application to a conjecture of Serre--Tate and to the
conductor}

\SGAvar{\medskip
\noindent
\SGARefTarget{sga7.sga7i.eix.n4-1}{4.1}. We suppose that $S$ is henselian, and retain the notation of (\SGARefLink{sga7.sga7i.eix.n2-5-3}{2.5.3}),
with $\ell\ne p$.
\footnote{The case $\ell\ne p$ will be studied in \S\S5 and 6 below.}
To study the action of $\pi=\operatorname{Gal}(\overline K/K)$ and of its
inertia subgroup $I$ on $U=T_\ell(A(\overline K))$, we shall use the fact
that there exists an open subgroup $I'$ of $I$, which we may moreover take
to be normal in $\pi$, such that the action of $I'$ on $U$ is unipotent of
level two; that is,
\[
 U^{I'}\supset (U^{I'})^\perp.
\]
It is then clear that $U^{I'}$ does not depend on the particular choice of
$I'$.  We might call it the \underline{essentially fixed submodule} of
$U$, and shall denote it by $U^{\mathrm{ef}}$.  Its orthogonal will be
called the \underline{essentially toric submodule} of $U$, and will be
denoted by $U^{\mathrm{et}}$.  We thus have a three-step filtration
\begin{equation}\tag{\SGARefTarget{sga7.sga7i.eix.n4-1-1}{4.1.1}}
 U\supset U^{\mathrm{ef}}\supset U^{\mathrm{et}}\supset0,
\end{equation}
which is plainly invariant under the action of $\pi$.

Let us also note that there is a largest subgroup of $I$ acting
unipotently (or, equivalently here, unipotently of level two) on $U$: it
is the subgroup $I'$ consisting of the elements acting unipotently.  It is
also the kernel of the representation of $I$ on the graded module
associated with the filtration (\SGARefLink{sga7.sga7i.eix.n4-1-1}{4.1.1}), or, equivalently in the present
case, the kernel of the induced representation of $I$ on
$U^{\mathrm{ef}}$.  By the Galois criterion for semistable reduction
(\SGARefLink{sga7.sga7i.eix.n3-5}{3.5}), this subgroup $I'$ corresponds to \underline{the smallest Galois
extension} $\widetilde K'$ of $\widetilde K$ (the maximal unramified
extension of $K$) \underline{such that} $A_{\widetilde K'}$ \underline{has
semistable reduction}.  This description also shows that the subgroup
$I'$ of $I$ is independent of the choice of the prime $\ell\ne p$.

}{\medskip
\noindent
\SGARefTarget{sga7.sga7i.eix.n4-1}{4.1}. We suppose that $S$ is henselian, and retain the notation of (\SGARefLink{sga7.sga7i.eix.n2-5-3}{2.5.3}),
with $\ell\ne p$.
\footnote{The case $\ell=p$ will be studied in \S\S5 and 6 below.}
To study the action of $\pi=\operatorname{Gal}(\overline K/K)$ and of its
inertia subgroup $I$ on $U=T_\ell(A(\overline K))$, we shall use the fact
that there exists an open subgroup $I'$ of $I$, which we may moreover take
to be normal in $\pi$, such that the action of $I'$ on $U$ is unipotent of
level two; that is,
\[
 U^{I'}\supset (U^{I'})^\perp.
\]
It is then clear that $U^{I'}$ does not depend on the particular choice of
$I'$.  We might call it the \underline{essentially fixed submodule} of
$U$, and shall denote it by $U^{\mathrm{ef}}$.  Its orthogonal will be
called the \underline{essentially toric submodule} of $U$, and will be
denoted by $U^{\mathrm{et}}$.  We thus have a three-step filtration
\begin{equation}\tag{\SGARefTarget{sga7.sga7i.eix.n4-1-1}{4.1.1}}
 U\supset U^{\mathrm{ef}}\supset U^{\mathrm{et}}\supset0,
\end{equation}
which is plainly invariant under the action of $\pi$.

Let us also note that there is a largest subgroup of $I$ acting
unipotently (or, equivalently here, unipotently of level two) on $U$: it
is the subgroup $I'$ consisting of the elements acting unipotently.  It is
also the kernel of the representation of $I$ on the graded module
associated with the filtration (\SGARefLink{sga7.sga7i.eix.n4-1-1}{4.1.1}), or, equivalently in the present
case, the kernel of the induced representation of $I$ on
$U^{\mathrm{ef}}$.  By the Galois criterion for semistable reduction
(\SGARefLink{sga7.sga7i.eix.n3-5}{3.5}), this subgroup $I$ corresponds to \underline{the smallest Galois
extension} $\widetilde K'$ of $\widetilde K$ (the maximal unramified
extension of $K$) \underline{such that} $A_{\widetilde K'}$ \underline{has
semistable reduction}.  This description also shows that the subgroup
$I'$ of $I$ is independent of the choice of the prime $\ell\ne p$.

}{E7I-IX-017}\SGApage{I}{\(\approx\)355}
Notice that, once the representation of $\pi$ on $U^{\mathrm{ef}}$, and
hence on $U^{\mathrm{et}}$, is known, the representation on
$U/U^{\mathrm{ef}}$ is known up to isogeny, by means of the isogeny
induced by the polarization pairing (\SGARefLink{sga7.sga7i.eix.n2-5-1}{2.5.1}):
\begin{equation}\tag{\SGARefTarget{sga7.sga7i.eix.n4-1-2}{4.1.2}}
 U/U^{\mathrm{ef}}\longrightarrow (U^{\mathrm{et}})^\vee(1)
 \qquad\text{an isogeny},
\end{equation}
where $(1)$ denotes the Tate twist by
$\mathbb Z_\ell(1)=T_\ell(\overline K^*)$; this homomorphism is plainly
compatible with the operations of $\pi$.  More precisely, there is a
canonical \underline{isomorphism}
\[
 U/U^{\mathrm{ef}}\simeq ((U')^{\mathrm{et}})^\vee(1),
 \qquad U'=T_\ell(A'(\overline K)).
\]
It therefore remains to study more closely the action of $\pi$ on
$U^{\mathrm{ef}}$.

By construction this action is trivial on the normal subgroup $I'$ of
$I$, and hence factors through a representation of the quotient group
\begin{equation}\tag{\SGARefTarget{sga7.sga7i.eix.n4-1-3}{4.1.3}}
 \pi'=\pi/I',
\end{equation}
an action of which we shall now give a description ``independent of
$\ell$.''

\SGApage{I}{356}
\medskip
\noindent
\SGARefTarget{sga7.sga7i.eix.n4-2}{4.2}. Let $K'$ be a Galois extension of $K$ corresponding to an open
normal subgroup $V$ of $\pi$ such that $V\cap I=I'$.  The Galois group
$G$ of $K'/K$ then occurs in an exact sequence
\begin{equation}\tag{\SGARefTarget{sga7.sga7i.eix.n4-2-1}{4.2.1}}
 1\longrightarrow\Gamma\longrightarrow G\longrightarrow G_o
 \longrightarrow1,
\end{equation}
where
\begin{equation}\tag{\SGARefTarget{sga7.sga7i.eix.n4-2-2}{4.2.2}}
 \Gamma=I/I'
\end{equation}
is the inertia subgroup, and $G_o$ is the Galois group of the largest
unramified subextension $K''$ of $K'/K$.  Thus (\SGARefLink{sga7.sga7i.eix.n4-2-1}{4.2.1}) is the image of the
analogous exact sequence
\begin{equation}\tag{\SGARefTarget{sga7.sga7i.eix.n4-2-3}{4.2.3}}
 1\longrightarrow\Gamma\longrightarrow\pi'\longrightarrow\pi_o
 \longrightarrow1,
\end{equation}
where $\pi'$ is the group (\SGARefLink{sga7.sga7i.eix.n4-1-3}{4.1.3}) acting on $U^{\mathrm{ef}}$.  We
therefore have the following diagram of field inclusions:
\begin{equation}\tag{\SGARefTarget{sga7.sga7i.eix.n4-2-4}{4.2.4}}
\begin{tikzpicture}[baseline=(current bounding box.center),
  every node/.style={inner sep=2pt}]
  \node (K) at (0,0) {$K$};
  \node (Ktwo) at (4.25,0) {$K''=\widetilde K\cap K'$};
  \node (Kt) at (8.5,0) {$\widetilde K$};
  \node (Kp) at (4.25,-1.8) {$K'$};
  \node (Ktp) at (8.5,-1.8) {$\widetilde K'$};
  \node (Kb) at (8.5,-3.6) {$\overline K$};
  \draw[->] (K) -- node[above] {$G_o$} (Ktwo);
  \draw[->] (Ktwo) -- (Kt);
  \draw[->] (Ktwo) -- node[right] {$\Gamma$} (Kp);
  \draw[->] (Kp) -- (Ktp);
  \draw[->] (Kt) -- node[right] {$\Gamma$} (Ktp);
  \draw[->] (Ktp) -- node[right] {$I'$} (Kb);
  \draw[-] (K) -- node[below left,pos=.58] {$G$} (Kp);
  \draw[-] (K) to[out=65,in=115,looseness=.85]
    node[above,pos=.78] {$\pi_o$} (Kt);
  \draw[-] (Kt) to[out=12,in=12,looseness=1.25]
    node[right,pos=.52] {$I$} (Kb);
\end{tikzpicture}
\end{equation}

\SGAvar{Let $S'$ be the normalization of $S$ in $K'$, and let $A^*$ be the
Neron model of $A_{K'}$.  Then, by the Galois criterion for semistable
reduction (\SGARefLink{sga7.sga7i.eix.n3-5}{4.5}), $A_{K'}$ \underline{has semistable reduction over
$S'$}; that is, $A^{*o}_o$ \underline{is an extension of an abelian
scheme by a torus}.  On the other hand, by transport of structure, $G$
\underline{acts on $A^*$, and hence on $A^{*o}$, by automorphisms in a
manner compatible with the given action of $G$ on the base scheme $S'$
of $A^*$}.  It is then clear that \underline{the action of $\pi'$ on}
\[
 U^{\mathrm{ef}}=T_\ell(A^{*o})^f(\overline K)
\]
\underline{is obtained by transport of structure from the action of
$\pi'$ on $(A^{*o}/K',\widetilde K')$}, where $\pi'$ acts on $A^*$
through $G$.  Taking into account the canonical isomorphism
\begin{equation}\tag{\SGARefTarget{sga7.sga7i.eix.n4-2-6}{4.2.6}}
 U^{\mathrm{ef}}\simeq T_\ell(A^{*o})^f(\overline K)
 \simeq T_\ell(A^{*o}_o)(\overline k)
\end{equation}
(valid for every twisted constant $\ell$-adic sheaf over $S'$; cf. I\hspace{2.5em}), we obtain the following
description of the action of $\pi'$ on $U^{\mathrm{ef}}$.

}{Let $S'$ be the normalization of $S$ in $K'$, and let $A^*$ be the
Neron model of $A_{K'}$.  Then, by the Galois criterion for semistable
reduction (\SGARefLink{sga7.sga7i.eix.n3-5}{3.5}), $A_{K'}$ \underline{has semistable reduction over
$S'$}; that is, $A^{*o}_o$ \underline{is an extension of an abelian
scheme by a torus}.  On the other hand, by transport of structure, $G$
\underline{acts on $A^*$, and hence on $A^{*o}$, by automorphisms in a
manner compatible with the given action of $G$ on the base scheme $S'$
of $A^*$}.  It is then clear that \underline{the action of $\pi'$ on}
\[
 U^{\mathrm{ef}}=T_\ell(A^{*o})^f(\overline K)
\]
\underline{is obtained by transport of structure from the action of
$\pi'$ on $(A^{*o}/K',\widetilde K')$}, where $\pi'$ acts on $A^*$
through $G$.  Taking into account the canonical isomorphism
\begin{equation}\tag{\SGARefTarget{sga7.sga7i.eix.n4-2-6}{4.2.6}}
 U^{\mathrm{ef}}\simeq T_\ell(A^{*o})^f(\overline K)
 \simeq T_\ell(A^{*o}_o)(\overline k)
\end{equation}
(valid for every twisted constant $\ell$-adic sheaf over $S'$; cf.
Expos\'e I [reference omitted in the source]), we obtain the following
description of the action of $\pi'$ on $U^{\mathrm{ef}}$.

}{E7I-IX-018}\SGApage{I}{\(\approx\)357}
Consider the smooth connected group scheme $A^{*o}_o$ over the residue
field $k'\hookrightarrow\overline k$ of $S'$, which is a quasi-Galois
extension of $k$ with Galois group $G_o$.  The quotient group $G$ of
$\pi'$ acts on $A^{*o}_o$ compatibly with its action on $k'$ through
$G_o$.  Moreover, $\pi'$ acts on $\overline k$ through its quotient
$\pi_o$, and hence acts on the situation
$(A^{*o}_o/k',\overline k)$.  Thus, by transport of structure, $\pi'$
acts on
\[
 U^{\mathrm{ef}}=T_\ell(A^{*o}_o)(\overline k),
\]
and this is the desired action of $\pi'$.  In particular, the subgroup
$\Gamma$ of $G$, which acts trivially on $k'$, acts on $A^{*o}_o$ by
$k'$-automorphisms, and the corresponding representation of
$\Gamma=I/I'$ on $T_\ell(A^{*o}_o)(\overline k)$ is the representation
deduced from the natural representation of the inertia group $I$ on
$U^{\mathrm{ef}}$.

\SGAvar{On the other hand, the essentially toric part $U^{\mathrm{et}}$ of $U$
corresponds simply to the submodule
\begin{equation}\tag{\SGARefTarget{sga7.sga7i.eix.n4-2-7}{4.2.7}}
 U^{\mathrm{et}}\simeq T_\ell(T^*_o)(\overline k)
 \subset U^{\mathrm{ef}}\simeq T_\ell(A^{*o}_o)(\overline k),
\end{equation}
where $T^*_o$ is the largest subtorus of $A^{*o}_o$.  Likewise, put
\begin{equation}\tag{\SGARefTarget{sga7.sga7i.eix.n4-2-8}{4.2.8}}
 B^*_o=A^{*o}_o/T^*_o
\end{equation}
for the abelian part of $A^{*o}_o$.  The action of $G$ on $A^{*o}_o$
induces separately actions of $G$ on $T^*_o$ and on $B^*_o$, and the
action of $\pi'$ on $U^{\mathrm{et}}$ and on
$U^{\mathrm{ef}}/U^{\mathrm{et}}$ is obtained from these actions by
means of the isomorphism (\SGARefLink{sga7.sga7i.eix.n4-2-7}{4.2.7}) and the corresponding isomorphism
\begin{equation}\tag{\SGARefTarget{sga7.sga7i.eix.n4-2-9}{4.2.9}}
 U^{\mathrm{ef}}/U^{\mathrm{et}}
 \simeq T_\ell(B^*_o)(\overline k).
\end{equation}
It follows from this and the isogeny (\SGARefLink{sga7.sga7i.eix.n4-1-2}{4.1.2}) that the character
$g\longmapsto\operatorname{Tr}g_U$ of the action of $\pi'$ (or of
$\pi$) on $U$ is given by
\begin{equation}\tag{\SGARefTarget{sga7.sga7i.eix.n4-2-10}{4.2.10}}
 \operatorname{Tr}g_{U_\ell}
 =\operatorname{Tr}g_{B^*_o,\ell}
 +\operatorname{Tr}g_{T^*_o,\ell}
 +\chi_\ell(g)\operatorname{Tr}g^{-1}_{T^*_o,\ell}.
\end{equation}
Here $U_\ell$ has been written for $U$ on the left to recall the chosen
prime $\ell\ne p$; $g_{B^*_o,\ell}$ and $g_{T^*_o,\ell}$ denote,
respectively, the actions of $g$ on the Tate modules
$T_\ell(B^*_o)(\overline k)$ and $T_\ell(T^*_o)(\overline k)$ induced by
its actions on $B^*_o$ and $T^*_o$ just described; and
\begin{equation}\tag{\SGARefTarget{sga7.sga7i.eix.n4-2-11}{4.2.11}}
 \chi_\ell:\pi'\longrightarrow\pi_o\longrightarrow\mathbb Z_\ell^*
\end{equation}
is the character of $\pi'$ (or of $\pi_o$) induced by the action on
$\mathbb Z_\ell(1)=T_\ell(\overline k^*)$---a character that is trivial
on $\Gamma^{o}$.

}{On the other hand, the essentially toric part $U^{\mathrm{et}}$ of $U$
corresponds simply to the submodule
\begin{equation}\tag{\SGARefTarget{sga7.sga7i.eix.n4-2-7}{4.2.7}}
 U^{\mathrm{et}}\simeq T_\ell(T^*_o)(\overline k)
 \subset U^{\mathrm{ef}}\simeq T_\ell(A^{*o}_o)(\overline k),
\end{equation}
where $T^*_o$ is the largest subtorus of $A^{*o}_o$.  Likewise, put
\begin{equation}\tag{\SGARefTarget{sga7.sga7i.eix.n4-2-8}{4.2.8}}
 B^*_o=A^{*o}_o/T^*_o
\end{equation}
for the abelian part of $A^{*o}_o$.  The action of $G$ on $A^{*o}_o$
induces separately actions of $G$ on $T^*_o$ and on $B^*_o$, and the
action of $\pi'$ on $U^{\mathrm{et}}$ and on
$U^{\mathrm{ef}}/U^{\mathrm{et}}$ is obtained from these actions by
means of the isomorphism (\SGARefLink{sga7.sga7i.eix.n4-2-7}{4.2.7}) and the corresponding isomorphism
\begin{equation}\tag{\SGARefTarget{sga7.sga7i.eix.n4-2-9}{4.2.9}}
 U^{\mathrm{ef}}/U^{\mathrm{et}}
 \simeq T_\ell(B^*_o)(\overline k).
\end{equation}
It follows from this and the isogeny (\SGARefLink{sga7.sga7i.eix.n4-1-2}{4.1.2}) that the character
$g\longmapsto\operatorname{Tr}g_U$ of the action of $\pi'$ (or of
$\pi$) on $U$ is given by
\begin{equation}\tag{\SGARefTarget{sga7.sga7i.eix.n4-2-10}{4.2.10}}
 \operatorname{Tr}g_{U_\ell}
 =\operatorname{Tr}g_{B^*_o,\ell}
 +\operatorname{Tr}g_{T^*_o,\ell}
 +\chi_\ell(g)\operatorname{Tr}g^{-1}_{T^*_o,\ell}.
\end{equation}
Here $U_\ell$ has been written for $U$ on the left to recall the chosen
prime $\ell\ne p$; $g_{B^*_o,\ell}$ and $g_{T^*_o,\ell}$ denote,
respectively, the actions of $g$ on the Tate modules
$T_\ell(B^*_o)(\overline k)$ and $T_\ell(T^*_o)(\overline k)$ induced by
its actions on $B^*_o$ and $T^*_o$ just described; and
\begin{equation}\tag{\SGARefTarget{sga7.sga7i.eix.n4-2-11}{4.2.11}}
 \chi_\ell:\pi'\longrightarrow\pi_o\longrightarrow\mathbb Z_\ell^*
\end{equation}
is the character of $\pi'$ (or of $\pi_o$) induced by the action on
$\mathbb Z_\ell(1)=T_\ell(\overline k^*)$---a character that is trivial
on $\Gamma$.

}{E7I-IX-019}\SGApage{I}{\(\approx\)358}
We are now in a position to prove a special case of a conjecture of
Serre--Tate [29, Appendix].
\footnote{This is essentially the case $i=1$ of those conjectures for
the $H^i$.}



\SGApage{I}{359}
\medskip
\noindent
\underline{Theorem} \SGARefTarget{sga7.sga7i.eix.n4-3}{4.3}. \underline{Let $S$ be a henselian trait with
residue field $k$ and function field $K$, and let $A_K$ be an abelian
scheme over $K$.  Consider the representation of}
$\pi=\operatorname{Gal}(\overline K/K)$ \underline{on the Tate module}
$U_\ell=T_\ell(A(\overline K))$, \underline{where $\ell$ is a prime
distinct from the residual characteristic.}

\smallskip
\noindent
a) \underline{There exists an open subgroup $I'$ of the inertia group
that acts unipotently on $U_\ell$, so that the character of the
representation of $\pi$ on $U_\ell$ is in fact a function on $\pi/I'$.}
\underline{Moreover, on $\Gamma=I/I'$ this function takes integral
values and is independent of $\ell$.  More generally, for every $g\in I$,
the coefficients of the characteristic polynomial of $g_{U_\ell}$ are
rational integers independent of $\ell$.}

\smallskip
\noindent
b) \underline{Suppose that $k$ is a finite field with $q$ elements, so
that}
\[
 \pi_o=\pi/I\simeq\operatorname{Gal}(\overline k/k)
\]
\underline{is identified with the profinite completion $\widehat{\mathbb
Z}$ of $\mathbb Z$, the generator $1$ of $\mathbb Z$ corresponding to
the Frobenius automorphism}
\[
 f_q:x\longmapsto x^q
\]
\underline{of $\overline k/k$.  Then, for every $g\in\pi$ whose image in
$\pi_o$ lies in $\mathbb Z$, the trace} (\underline{and, more generally, every
coefficient of the characteristic polynomial}) \underline{of $g_{U_\ell}$ is a
rational number independent of $\ell$.  It is even an integer if the
image of $g$ in $\pi_o$ is of the form $f_q^{,n}$, with $n\geq0$.}

Let us prove, for example, the assertions concerning traces.  The
assertions concerning the coefficients of the characteristic polynomial
could be proved in the same way.  One may also note that they follow
formally from the assertions about traces, by expressing these
coefficients in terms of the traces of the iterates (although those
formulas introduce denominators and therefore do not give the asserted
integrality), together with a lemma of Dwork--Fatou, for which a pleasant
proof is given by Kleiman in [14, 2.8].

\SGApage{I}{\(\approx\)360}
\smallskip
\noindent
a) We use formula (\SGARefLink{sga7.sga7i.eix.n4-2-10}{4.2.10}), which here becomes
\begin{equation}\tag{\SGARefTarget{sga7.sga7i.eix.n4-3-0}{4.3.0}}
 \operatorname{Tr}g_{U_\ell}
 =\operatorname{Tr}g_{B^*_o,\ell}
 +\operatorname{Tr}g_{T^*_o,\ell}
 +\operatorname{Tr}g^{-1}_{T^*_o,\ell}.
\end{equation}
Here $g_{B^*_o}$ and $g_{T^*_o}$ are now $k'$-automorphisms of $B^*_o$
and $T^*_o$, respectively.  It therefore suffices to show that the two
corresponding traces are rational integers independent of $\ell$, which
is a general ``geometric'' assertion about algebraic groups.  For
$g_{B^*_o}$ this is a well-known result of Weil in the theory of abelian
varieties [5] [15].  For $g_{T^*_o}$ it is practically immediate: one
describes the representation of $\Gamma$ on $T=T^*_o$ by means of the
corresponding representation on the dual group $M'\simeq\mathbb Z^\mu$.
The required trace is that of the contragredient of the latter
representation, by the canonical isomorphism
\[
 T_\ell(T)\simeq\check M_\ell(1)
 \mathrel{\mathop{=}^{\mathrm{def}}}
 \check M'\otimes_{\mathbb Z}\mathbb Z_\ell(1).
\]

\smallskip
\noindent
b) If the image of $g$ in $\pi_o$ is $f_q^{,n}$, formula (\SGARefLink{sga7.sga7i.eix.n4-2-10}{4.2.10}) now
becomes
\begin{equation}\tag{\SGARefTarget{sga7.sga7i.eix.n4-3-1}{4.3.1}}
 \operatorname{Tr}g_{U,\ell}
 =\operatorname{Tr}g_{B^*_o,\ell}
 +\operatorname{Tr}g_{T^*_o,\ell}
 +q^n\operatorname{Tr}g^{-1}_{T^*_o,\ell}.
\end{equation}
We are thus reduced to proving the following lemma, in which $k$ is to be
replaced by an algebraic closure of $k'$, and $A$ by
$A^{*o}_o\otimes_{k'}\overline k$.

\medskip
\noindent
\underline{Lemma} \SGARefTarget{sga7.sga7i.eix.n4-3-2}{4.3.2}. \underline{Let $k$ be the algebraic closure of
a finite field, let $A$ be a commutative group scheme of finite type over
$k$, and let $g$ be an automorphism of this group scheme compatible with
the automorphism $f$ of $\operatorname{Spec}(k)$ obtained by transport of
structure from an iterated Frobenius automorphism}
\[
 x\longmapsto x^q,
 \qquad q=p^r,\quad r\geq0,
\]
\underline{where $q$ is a power of the characteristic.  Thus the action
of $g$ on functions $\varphi$ satisfies}
\[
 g^*(\lambda\varphi)=f^*(\lambda)g^*(\varphi)
 =\lambda^{q^{-1}}g^*(\varphi)
 =\lambda^{p^{-r}}g^*(\varphi)
 \qquad(\lambda\in k).
\]
\underline{Let $g_\ell$ be the automorphism induced by $g$ on
$T_\ell(A(k))$ by ``transport of structure,'' where $\ell$ is a prime
different from $p$.  Then the trace of $g_\ell$ (and, more generally,
every coefficient of its characteristic polynomial) is a rational
integer independent of $\ell$, and its eigenvalues are algebraic integers
of absolute value $q^{1/2}$ or $q$.  If $A$ is an abelian scheme, their
absolute values are $q^{1/2}$; if $A$ is affine, the eigenvalues are of
the form $q\varepsilon_i$, where the $\varepsilon_i$ are roots of unity.}

\SGApage{I}{\(\approx\)361}
Indeed, consider the ``$q$-th power'' endomorphism $f_q^A$ of the
underlying absolute scheme of $A$, defined as the identity on points and
as the $q$-th power map on the structure sheaf.  It is an endomorphism of
the group scheme $A$ compatible with the endomorphism
$f_q^{\operatorname{Spec}(k)}=f^{-1}$ of the base
$\operatorname{Spec}(k)$.  Consequently the composite
\begin{equation}\tag{\SGARefTarget{sga7.sga7i.eix.n4-3-3}{4.3.3}}
 F=f_q^A g
\end{equation}
is an endomorphism of the group scheme $A$ compatible with the identity
on $\operatorname{Spec}(k)$; in other words, it is a $k$-endomorphism of
the group scheme $A$.  I claim that the automorphism $\Gamma_g$ of $A(k)$
defined by $g$ through transport of structure, namely
\begin{equation}\tag{\SGARefTarget{sga7.sga7i.eix.n4-3-4}{4.3.4}}
 \Gamma_g(u)=g\circ u\circ f^{-1}
 \qquad\text{for }u\in A(k),
\end{equation}
may also be described by
\begin{equation}\tag{\SGARefTarget{sga7.sga7i.eix.n4-3-5}{4.3.5}}
 \Gamma_g(u)=F(u)
 \qquad\text{for }u\in A(k).
\end{equation}
To verify this, it is enough to observe that the underlying set maps of
the morphisms $\operatorname{Spec}(k)\longrightarrow A$ represented by
the two sides are the same.  This reduces us to the equality obtained
from (\SGARefLink{sga7.sga7i.eix.n4-3-5}{4.3.5}) by composing its left-hand member on the left with $f_q^A$
and its right-hand member on the right with $f^{-1}$, since $f_q^A$ and
$f^{-1}$ induce the identity on the underlying sets of the source
schemes.  The resulting formula follows immediately from (\SGARefLink{sga7.sga7i.eix.n4-3-3}{4.3.3}) and
(\SGARefLink{sga7.sga7i.eix.n4-3-4}{4.3.4}).

\SGApage{I}{\(\approx\)362}
Applying (\SGARefLink{sga7.sga7i.eix.n4-3-5}{4.3.5}) to $u\in{}_{\ell^\nu}A(k)$ and passing to the inverse
limit over $\nu$, we find that $g_\ell$ may also be interpreted as the
endomorphism $T_\ell(F)$ of $T_\ell(A(k))$ associated with the
endomorphism $F$ of $A$.  It is known, on the other hand, that
$T_\ell(F)$ has a characteristic polynomial with integral coefficients
independent of $\ell$.  By devissage, this reduces to the case in which
$A$ is either an abelian scheme or a torus, where the required result was
already used in part a) of the proof.  The same argument also proves the
assertions about the eigenvalues, in view, in the abelian-scheme case, of
the classical results of Weil [15].

This proves \SGARefLink{sga7.sga7i.eix.n4-3-2}{4.3.2} and completes the proof of \SGARefLink{sga7.sga7i.eix.n4-3}{4.3}.  At the same time it
proves the following result.

\medskip
\noindent
\underline{Corollary} \SGARefTarget{sga7.sga7i.eix.n4-4}{4.4}. \underline{Suppose that the hypotheses of
\SGARefLink{sga7.sga7i.eix.n4-3}{4.3} b) hold.  Then, for every $g\in\pi'$ whose image in $\pi_o$ is of the
form $f_q^n$ with $n\geq0$, the eigenvalues of the induced automorphisms
of}
\[
 U^{\mathrm{et}},\qquad
 U^{\mathrm{ef}}/U^{\mathrm{et}},\qquad
 U/U^{\mathrm{ef}}
\]
\underline{are algebraic integers whose absolute values are,
respectively, $q^n$, $q^{n/2}$, and} $1$.  \underline{For the two extreme pieces one
may say more precisely that the corresponding eigenvalues are,
respectively,}
\[
 (q^n\varepsilon_i)_{1\leq i\leq\mu}
 \qquad\text{and}\qquad
 (\varepsilon_i)_{1\leq i\leq\mu},
\]
\underline{where the $\varepsilon_i$ $(1\leq i\leq\mu)$ are roots of
unity.}

\SGApage{I}{\(\approx\)363}
Indeed, formula (\SGARefLink{sga7.sga7i.eix.n4-1-2}{4.1.2}) gives for the last piece the eigenvalues
\[
 \frac{q^n}{q^n\varepsilon_i}=\varepsilon_i^{-1}.
\]
But since $\varepsilon_i$ is a root of unity, one has
$\varepsilon_i^{-1}=\overline{\varepsilon_i}$ (complex conjugate), and
since the family of the $\varepsilon_i$ is stable under complex
conjugation (being the family of roots of an equation with integral
coefficients), the family of the $\varepsilon_i^{-1}$ is also the family
of the $\varepsilon_i$.  This proves the last assertion of \SGARefLink{sga7.sga7i.eix.n4-4}{4.4}.



\medskip
\noindent
\SGARefTarget{sga7.sga7i.eix.n4-5}{4.5}. We shall give an application of \SGARefLink{sga7.sga7i.eix.n4-3}{4.3} a) to the theory of the local
conductor of abelian schemes.  Suppose that $k$ is algebraically
closed.\footnote{It would in fact suffice for $k$ to be perfect, but,
after passing to the strict henselization of $S$, one is immediately
reduced to the case considered here.}  Recall (SGA 5 X 6.2, or [20])
that, for every commutative etale $\ell$-group scheme $F_K$ over $K$
($\ell\neq p$), one defines an integer
\begin{equation}\tag{\SGARefTarget{sga7.sga7i.eix.n4-5-1}{4.5.1}}
 \alpha_s(F_K)=\alpha_s(i_*F_K),
\end{equation}
where $i:\eta\longrightarrow S$ denotes the inclusion, as follows.  One
chooses a finite Galois extension $K'$ of $K$, with group $\Gamma$,
that splits $F_K$, so that $F_K$ is represented by a finite
$\ell$-group $M$ on which $\Gamma$ acts.  On the other hand, one
associates with the extension $K'/K$ a certain projective
$\mathbb Z_\ell[\Gamma]$-module, defined up to isomorphism,
\begin{equation}\tag{\SGARefTarget{sga7.sga7i.eix.n4-5-2}{4.5.2}}
 \operatorname{Sw}=\operatorname{Sw}(K'/K,\mathbb Z_\ell),
\end{equation}
and sets
\begin{equation}\tag{\SGARefTarget{sga7.sga7i.eix.n4-5-3}{4.5.3}}
 \alpha_s(F_K)
 =\operatorname{length}
   \operatorname{Hom}_{\mathbb Z_\ell(\Gamma)}(\operatorname{Sw},M).
\end{equation}
One checks (loc.\ cit.) that this expression is independent of the
choice of the extension $K'/K$.  Moreover, the projectivity of
$\operatorname{Sw}$ as a $\mathbb Z_\ell[\Gamma]$-module implies that
$\alpha_s(F_K)$ is additive in exact sequences.

\SGApage{I}{\(\approx\)364}
We apply this definition to the sheaf ${}_\ell A_K$, where $A_K$ is our
abelian scheme over $K$.  In this case one therefore obtains an integer
\begin{equation}\tag{\SGARefTarget{sga7.sga7i.eix.n4-5-4}{4.5.4}}
 \alpha_s(A_K;\ell)
 =\operatorname{length}
   \operatorname{Hom}_{\mathbb Z_\ell(\Gamma)}
   \bigl(\operatorname{Sw},
         U_\ell\otimes_{\mathbb Z_\ell}\mathbb F_\ell\bigr),
 \qquad
 U_\ell=T_\ell(A(\overline K)),
\end{equation}
where $\Gamma$ is chosen as above (a choice that a priori depends on
$\ell$).  The integer (\SGARefLink{sga7.sga7i.eix.n4-5-4}{4.5.4}) is called the \emph{wild conductor
exponent} of the abelian scheme $A_K$ relative to the discrete valuation
ring $V\subset K$.\footnote{One should also define the \emph{tame
conductor exponent} of $A_K$ relative to $V$ as the difference between
the lengths over $\mathbb Z/\ell\mathbb Z$ of the geometric generic and
special fibres of ${}_\ell A^o$.  It is independent of $\ell$ and equals
$2\lambda+\mu$ in the notation of (1.1.11).  \underline{The} (total) \underline{conductor
exponent} of $A_K$ is the sum of the tame and wild exponents; it is the
local term occurring in the Euler--Poincare formula for ${}_\ell A^o$ in
SGA 5 X 6.2 and [20].}  After passage to the henselization it is of
course defined without assuming that $V$ is henselian, and one is
immediately reduced to the case in which $S$ is strictly local.  We
intend to prove the following result.

\medskip
\noindent
\underline{Corollary} \SGARefTarget{sga7.sga7i.eix.n4-6}{4.6}. \underline{The wild conductor exponent
(\SGARefLink{sga7.sga7i.eix.n4-5-4}{4.5.4}) is independent of the choice of the prime $\ell\neq p$.}

\SGApage{I}{\(\approx\)365}
We may suppose that $S$ is strictly local.  We compute (\SGARefLink{sga7.sga7i.eix.n4-5-4}{4.5.4}), replacing
$\operatorname{Sw}$ by any other projective
$\mathbb Z_\ell(\Gamma)$-module $L$.  Observe that the filtration
(\SGARefLink{sga7.sga7i.eix.n4-1-1}{4.1.1}) of $U_\ell$ consists of $\mathbb Z_\ell$-submodules that are
manifestly direct summands.  It therefore induces a corresponding
filtration of
\[
 U(\ell)=U_\ell\otimes_{\mathbb Z_\ell}\mathbb F_\ell,
\]
whose successive factors are obtained, by tensoring with
$\mathbb F_\ell$, from the factors $V_i$ $(i=1,2,3)$
\[
 U^{\mathrm{et}},\qquad
 U^{\mathrm{ef}}/U^{\mathrm{et}},\qquad
 U/U^{\mathrm{ef}}
\]
appearing in (\SGARefLink{sga7.sga7i.eix.n4-1-1}{4.1.1}).  This filtration is, of course, stable under
$\Gamma$.  Additivity thus decomposes
\[
 \operatorname{length}
 \operatorname{Hom}_{\mathbb Z_\ell(\Gamma)}(L,U(\ell))
\]
as a sum of three terms, the lengths of the groups
\begin{equation}\tag{*}
 \operatorname{Hom}_{\mathbb Z_\ell(\Gamma)}
   (L,V_i\otimes_{\mathbb Z_\ell}\mathbb F_\ell)
 \simeq
 \operatorname{Hom}_{\mathbb Z_\ell(\Gamma)}(L,V_i)
   \otimes_{\mathbb Z_\ell}\mathbb F_\ell.
\end{equation}
The isomorphism in (*) again follows from the assumed projectivity of
$L$; it is only necessary to choose $K'$ large enough for $\Gamma$ to
act on the $V_i$.  Now
$\operatorname{Hom}_{\mathbb Z_\ell(\Gamma)}(L,V_i)$ is plainly a
torsion-free $\mathbb Z_\ell$-module, so its rank is the length of the
right-hand side of (*).  If $\varphi_i$ is the character of the
representation of $\Gamma$ on $V_i$, and $\psi$ the character of the
representation of $I$ on $\check L$, we obtain
\begin{equation}\tag{\SGARefTarget{sga7.sga7i.eix.n4-6-1}{4.6.1}}
 \operatorname{length}
 \operatorname{Hom}_{\mathbb Z_\ell(\Gamma)}(L,U(\ell))
 =
 \sum_{1\leq i\leq3}N^{-1}
 \sum_{g\in\Gamma}\psi(g)\varphi_i(g)
\end{equation}
\[
 =N^{-1}\sum_{g\in\Gamma}\psi(g)\varphi(g),
 \qquad N=\operatorname{card}\Gamma,
\]
where $\varphi(g)=\sum_i\varphi_i(g)$ is also the function on $\Gamma$
obtained, by passage to the quotient, from the character of the
monodromy representation of the inertia group $I$ on $U$.  By \SGARefLink{sga7.sga7i.eix.n4-3}{4.3} a),
this character is an integral-valued rational function independent of
$\ell$.  Taking $L=\operatorname{Sw}$, the same is true of the character
$\psi(g)$ of $L$, which is
\[
 a(g)-N\delta(g)+1,
\]
where $\delta$ is the Dirac function on $\Gamma$ and $a$ is the Artin
conductor of $K'/K$ [26].  This proves \SGARefLink{sga7.sga7i.eix.n4-6}{4.6} and at the same time gives the
following expression, independent of $\ell$, for the wild conductor
exponent:
\begin{equation}\tag{\SGARefTarget{sga7.sga7i.eix.n4-6-2}{4.6.2}}
 \alpha_s(A_K)
 =N^{-1}\sum_{g\in\Gamma}
   \bigl(a(g)-N\delta(g)+1\bigr)\varphi(g).
\end{equation}
Here $\Gamma$ is the Galois group of a finite Galois extension $K'$ of
$K$ that splits $U^{\mathrm{ef}}$ (for example the smallest such
extension; cf.\ \SGARefLink{sga7.sga7i.eix.n4-1}{4.1}), $a(g)$ is the Artin conductor on $\Gamma$, and
$\varphi(g)$ is the trace function on $\Gamma$ described in (\SGARefLink{sga7.sga7i.eix.n4-3-0}{4.3.0}).
In the semistable-reduction case one may take $K'=K$, hence
$\Gamma=\{e\}$, and obtains
\begin{equation}\tag{\SGARefTarget{sga7.sga7i.eix.n4-6-3}{4.6.3}}
 \alpha_s(A_K)=0
 \qquad\text{(semistable-reduction case).}
\end{equation}

\SGApage{I}{\(\approx\)366}
To conclude, let us state explicitly a criterion for semistable
reduction that follows from the general considerations of \SGARefLink{sga7.sga7i.eix.n4-2}{4.2}.

\medskip
\noindent
\underline{Proposition} \SGARefTarget{sga7.sga7i.eix.n4-7}{4.7}. \underline{(M.\ Raynaud's criterion for
semistable reduction.)  Let $S$ be a trait, let $A_K$ be an abelian
scheme over its function field $K$, and let $n>1$ be an integer prime
to the residual characteristic $p$.  Suppose that ${}_nA_K$ is
unramified relative to $S$, that is, that the inertia group acts
trivially on ${}_nA(\overline K)$.  Then either $A_K$ has semistable
reduction over $S$, or $n=2$; in the latter case, if $S$ is strictly
local, $A_K$ acquires semistable reduction over a finite Galois extension
$K'$ of $K$ whose Galois group is isomorphic to
$(\mathbb Z/2\mathbb Z)^r$.}\footnote{As follows from an example
communicated to me by J.\ Tate, \SGARefLink{sga7.sga7i.eix.n4-7}{4.7} is no longer valid when $n$ equals
the residual characteristic $p$, even if $K$ has characteristic zero
and $A_K$ is an elliptic curve with complex multiplication.}

We may plainly suppose that $S$ is strictly local.  With the notation
of \SGARefLink{sga7.sga7i.eix.n4-1}{4.1} and \SGARefLink{sga7.sga7i.eix.n4-2}{4.2}, take $K'$ as small as possible; in other words, suppose
that $\Gamma$ acts faithfully on $A^{*o}_o$.  By hypothesis $\Gamma$
acts trivially on ${}_nA(\overline K)$, and hence a fortiori on
\[
 {}_nA(\overline K)^{\mathrm{ef}}
 \simeq{}_nA^*_o(k)\supset{}_nA^{*o}_o(k).
\]
The conclusion follows from the next two lemmas.

\SGApage{I}{367}
\medskip
\noindent
\underline{Lemma} \SGARefTarget{sga7.sga7i.eix.n4-7-1}{4.7.1}. \underline{(J.-P.\ Serre.)  Let $G$ be a
commutative group scheme over a field $k$, which is an extension of an
abelian scheme by a torus, and let $n\geq2$ be an integer.  Every
finite-order automorphism of $G$ that induces the identity on ${}_nG$ is
the identity if $n\neq2$, and is either the identity or has order two if
$n=2$.}

\medskip
\noindent
\underline{Lemma} \SGARefTarget{sga7.sga7i.eix.n4-7-2}{4.7.2}. \underline{A finite group in which every
element distinct from $1$ has order $2$ is isomorphic to
$(\mathbb Z/2\mathbb Z)^r$ for a suitable $r\geq0$.}

\smallskip
\noindent
\SGARefTarget{sga7.sga7i.eix.n4-7-3}{4.7.3}. Let us recall the well-known proof of \SGARefLink{sga7.sga7i.eix.n4-7-1}{4.7.1}.  Since
$n\,\operatorname{id}_G$ is an epimorphism, the hypothesis that $u$
induces the identity on ${}_nG$, or equivalently that $1-u$ induces zero
on ${}_nG$, is expressed by the condition
\[
 1-u=nv
\]
in the endomorphism ring of $G$.  To prove the asserted conclusion, we
reduce, by decomposing the finite cyclic group generated by $u$ into its
$p$-primary components for the various primes $p$ dividing the order of
$u$, to the case in which $u$ has order $p^N$.

\SGAvar{Let $E$ be the subring of $\operatorname{End}(G)$ generated by $v$.
Since $u=1-nv$ and $u^{p^N}=1$, it is a quotient of
\[
 \mathbb Z[T]\big/\bigl(T^{p^N}-1\bigr);
\]
it is therefore of finite type over $\mathbb Z$.  It is moreover
torsion-free (a formal consequence of the divisibility of $G$), and thus
embeds in
\[
 E_{\mathbb Q}=E\otimes_{\mathbb Z}\mathbb Q.
\]
The latter is a quotient, after the change of variable $u=1-nT$, of
$\mathbb Q[u]/(u^{p^N}-1)$, which is a direct product of finite field
extensions of $\mathbb Q$; the same is therefore true of $E_{\mathbb Q}$.
Considering the images of $E$, $u$, and $v$ in the component fields
$K_i$ of $E_{\mathbb Q}$, it remains to prove the following assertion.
Let $K=\mathbb Q(u)$ be generated by a primitive $p^n$-th root of unity
$u$, and let $n\geq2$ be an integer such that
\[
 u-1\equiv0\pmod n
\]
in the ring of integers of $K$.  Then $u=1$, or $u=-1$ and $n=2$.

}{Let $E$ be the subring of $\operatorname{End}(G)$ generated by $v$.
Since $u=1-nv$ and $u^{p^N}=1$, it is a quotient of
\[
 \mathbb Z[T]\big/\bigl((1-nT)^{p^N}-1\bigr);
\]
it is therefore of finite type over $\mathbb Z$.  It is moreover
torsion-free (a formal consequence of the divisibility of $G$), and thus
embeds in
\[
 E_{\mathbb Q}=E\otimes_{\mathbb Z}\mathbb Q.
\]
The latter is a quotient, after the change of variable $u=1-nT$, of
$\mathbb Q[u]/(u^{p^N}-1)$, which is a direct product of finite field
extensions of $\mathbb Q$; the same is therefore true of $E_{\mathbb Q}$.
Considering the images of $E$, $u$, and $v$ in the component fields
$K_i$ of $E_{\mathbb Q}$, it remains to prove the following assertion.
Let $K=\mathbb Q(u)$ be generated by a primitive $p^m$-th root of unity
$u$, and let $n\geq2$ be an integer such that
\[
 u-1\equiv0\pmod n
\]
in the ring of integers of $K$.  Then $u=1$, or $u=-1$ and $n=2$.

}{E7I-IX-020}\SGApage{I}{\(\approx\)368}
This follows at once from the familiar fact that, if $u\neq1$, that is,
$m\geq1$, the ideal $(u-1)$ is prime and $(p)$ is a power of it, with
exponent greater than $1$ except when $p=2$ and $m=1$, that is, when
$u=-1$.  Since $(u-1)$ is prime and divides $p$, an integer $n\geq2$ can
divide $u-1$ only if $n=p$; and $p$ can divide $u-1$ only if
$(u-1)=(p)$, that is, if $u=-1$.  Hence $n=2$.

\smallskip
\noindent
\SGARefTarget{sga7.sga7i.eix.n4-7-4}{4.7.4}. Let us prove \SGARefLink{sga7.sga7i.eix.n4-7-2}{4.7.2}.  It is enough to prove that $\Gamma$ is
commutative.  If $x,y\in\Gamma$, then $(xy)(xy)=1$; since $x^2=1$ and
$y^2=1$, this may also be written
\[
 xy=(xy)^{-1}=y^{-1}x^{-1}=yx,
\]
as required.

\medskip
\noindent
\underline{Remark} \SGARefTarget{sga7.sga7i.eix.n4-7-5}{4.7.5}. The proof given in \SGARefLink{sga7.sga7i.eix.n4-7-3}{4.7.3}, together with
\SGARefLink{sga7.sga7i.eix.n4-7-2}{4.7.2}, shows that if $G$ is a divisible commutative group in a topos,
$n\geq2$ is an integer, and $\Gamma$ is a finite group acting faithfully
on $G$ while inducing the identity on ${}_nG$, then $\Gamma$ is the
trivial group if $n\geq3$, and is isomorphic to a group
$(\mathbb Z/2\mathbb Z)^r$ if $n=2$.



\SGApage{I}{\(\approx\)369}
\medskip
\noindent
\underline{5. The orthogonality theorem in the case $\ell=p$}
(\underline{semistable case}).  \underline{Duality of the abelian schemes $B_o$ and $B'_o$.
Characterization of the fixed part.  Criteria for good reduction,
essentially good reduction, and semistable reduction.}

\medskip
\noindent
\SGARefTarget{sga7.sga7i.eix.n5-1}{5.1}. Let us place ourselves in the situation of \SGARefLink{sga7.sga7i.eix.n2-2-3}{2.2.3}, where $S$ is
henselian and $A^o$ is $\ell$-divisible, so that the \underline{fixed
part}
\[
 T_\ell(A^o)^f\lhook\joinrel\longrightarrow T_\ell(A)
\]
is defined as the pro-$\ell$ Barsotti--Tate group formed by the
$({}_{\ell^\nu}A^o)^f$.  We shall now give a definition of the toric
part (\SGARefLink{sga7.sga7i.eix.n2-3-2}{2.3.2}) of $T_\ell(A^o)$ which is independent of the hypothesis
$\ell\neq p$ made in \SGARefLink{sga7.sga7i.eix.n2-3}{2.3}, but which instead assumes that $S$ is
\underline{complete}.\footnote{This hypothesis is in fact not
essential, as is shown in \SGARefLink{sga7.sga7i.eix.n6-1}{6.1} below.}  For this purpose, for every
integer $n\geq0$ introduce the local Artinian scheme
\begin{equation}\tag{\SGARefTarget{sga7.sga7i.eix.n5-1-1}{5.1.1}}
 S_n=\operatorname{Spec}(V/\mathfrak m^{n+1}),
\end{equation}
the $n$-th infinitesimal neighbourhood of $s$ in $S$, and, for every
$S$-scheme $X$, put $X_n=X\times_S S_n$.  Since
$(A^o)_n=(A_n)^o$ is smooth over $S_n$, it is known (SGA 3 IX 3.6 bis)
that $T_o$ is the special fibre of a unique subtorus
\begin{equation}\tag{\SGARefTarget{sga7.sga7i.eix.n5-1-2}{5.1.2}}
 T_n\subset A_n^o.
\end{equation}
By the uniqueness of the $T_n$, they are obtained from one another by
the base changes $S_{n'}\longrightarrow S_n$.  In other words, they
define a \underline{formal subtorus} $\widehat T$ of the group
$\widehat A^o$ (in the category of formal schemes over $S$), the formal
completion of $A^o$ along the special fibre:
\begin{equation}\tag{\SGARefTarget{sga7.sga7i.eix.n5-1-3}{5.1.3}}
 \widehat T\subset\widehat A^o.
\end{equation}

Since the functor $X\longmapsto\widehat X$ (formal completion along the
special fibre) induces an equivalence between the category of finite
schemes over $S$ and that of finite formal schemes over $S$ (EGA III
4.8), we shall simply identify a finite scheme over $S$ with the
corresponding formal scheme.  With this identification, the definition
of the \emph{fixed parts} immediately gives, for every $\nu\geq0$,
\begin{equation}\tag{\SGARefTarget{sga7.sga7i.eix.n5-1-4}{5.1.4}}
 ({}_{\ell^\nu}A^o)^f={}_{\ell^\nu}(\widehat A^o),
\end{equation}
which may also be written as an identity of projective systems:
\begin{equation}\tag{\SGARefTarget{sga7.sga7i.eix.n5-1-5}{5.1.5}}
 T_\ell(A^o)^f=T_\ell(\widehat A^o).
\end{equation}
We therefore obtain a Barsotti--Tate pro-subgroup of the fixed part of
$T_\ell(A^o)$ by taking
$T_\ell(\widehat T)\subset T_\ell(\widehat A^o)$.  We accordingly
define the \underline{toric part} of $T_\ell(A^o)$ by
\begin{equation}\tag{\SGARefTarget{sga7.sga7i.eix.n5-1-6}{5.1.6}}
 T_\ell(A^o)^t
 \stackrel{\mathrm{def}}{=}
 T_\ell(\widehat T)
 \subset T_\ell(\widehat A^o)
 =T_\ell(A^o)^f.
\end{equation}
Thus, by definition,
\begin{equation}\tag{\SGARefTarget{sga7.sga7i.eix.n5-1-7}{5.1.7}}
 T_\ell(A^o)^t=({}_{\ell^\nu}\widehat T)_{\nu\geq0},
\end{equation}
the finite flat $S$-group ${}_{\ell^\nu}\widehat T$ itself being
described as the inductive limit of the finite flat $S_n$-groups
\begin{equation}\tag{\SGARefTarget{sga7.sga7i.eix.n5-1-8}{5.1.8}}
 {}_{\ell^\nu}(T_n),\qquad n\ \text{variable};
\end{equation}
that is, its affine algebra is described as the projective limit of the
affine algebras of the groups (\SGARefLink{sga7.sga7i.eix.n5-1-8}{5.1.8}).  Of course,
\begin{equation}\tag{\SGARefTarget{sga7.sga7i.eix.n5-1-9}{5.1.9}}
 \bigl(T_\ell(A^o)^t\bigr)_n=T_\ell(T_n).
\end{equation}
For $n=0$ this is precisely formula (\SGARefLink{sga7.sga7i.eix.n2-3-3}{2.3.3}), by which we characterized
the toric part (as a profinite etale group) in the case $\ell\neq p$.
This shows that definition (\SGARefLink{sga7.sga7i.eix.n5-1-6}{5.1.6}) is compatible with that given in
\SGARefLink{sga7.sga7i.eix.n2-3}{2.3} when $\ell\neq p$.

\SGApage{I}{\(\approx\)371}
Passing to generic fibres, the fixed and toric parts of $T_\ell(A^o)$
define fixed and toric parts in $T_\ell(A_K)$:
\begin{equation}\tag{\SGARefTarget{sga7.sga7i.eix.n5-1-10}{5.1.10}}
 T_\ell(A_K)^t\subset T_\ell(A_K)^f\subset T_\ell(A_K).
\end{equation}
These are inclusions of pro-$\ell$ Barsotti--Tate groups.  Consider the
analogous filtration corresponding to the dual abelian scheme $A'_K$:
\begin{equation}\tag{\SGARefTarget{sga7.sga7i.eix.n5-1-11}{5.1.11}}
 T_\ell(A'_K)^t\subset T_\ell(A'_K)^f\subset T_\ell(A'_K).
\end{equation}
We then have the following result, which generalizes \SGARefLink{sga7.sga7i.eix.n2-4}{2.4} (at least in
the case where $S$ is complete, to which one could easily have reduced
in \SGARefLink{sga7.sga7i.eix.n2-4}{2.4}).




\SGAvar{\medskip
\noindent
\underline{Theorem} \SGARefTarget{sga7.sga7i.eix.n5-2}{5.2}. \underline{Suppose that $S$ is a complete
trait, and let $A_K$ be an abelian scheme over $K$, $A'_K$ the dual
abelian scheme, and $\ell$ a prime such that the connected Neron model
$A^o$ is $\ell$-divisible (cf.\ \SGARefLink{sga7.sga7i.eix.n2-2-1}{2.2.1}), so that one then has the
filtrations (\SGARefLink{sga7.sga7i.eix.n5-1-10}{5.1.10}) and (\SGARefLink{sga7.sga7i.eix.n5-1-11}{5.1.11}) of the pro-groups
$T_\ell(A_K)$ and $T_\ell(A_K')$.  Then the toric parts are
recovered from the fixed parts by the formula}
\begin{equation}\tag{\SGARefTarget{sga7.sga7i.eix.n5-2-1}{5.2.1}}
 T_\ell(A_K)^t
 =
 T_\ell(A_K)^f\cap\bigl(T_\ell(A'_K)^f\bigr)^\perp,
\end{equation}
\underline{and by the analogous formula obtained by interchanging the
roles of $A_K$ and $A'_K$, where $\perp$ denotes the orthogonal
complement with respect to the canonical pairing (\SGARefLink{sga7.sga7i.eix.n1-0-2}{1.0.2}).}

}{\medskip
\noindent
\underline{Theorem} \SGARefTarget{sga7.sga7i.eix.n5-2}{5.2}. \underline{Suppose that $S$ is a complete
trait, and let $A_K$ be an abelian scheme over $K$, $A'_K$ the dual
abelian scheme, and $\ell$ a prime such that the connected Neron model
$A^o$ is $\ell$-divisible (cf.\ \SGARefLink{sga7.sga7i.eix.n2-2-1}{2.2.1}), so that one then has the
filtrations (\SGARefLink{sga7.sga7i.eix.n5-1-10}{5.1.10}) and (\SGARefLink{sga7.sga7i.eix.n5-1-11}{5.1.11}) of the pro-groups
$T_\ell(A_K)$ and $T_\ell(A'_K)$.  Then the toric parts are
recovered from the fixed parts by the formula}
\begin{equation}\tag{\SGARefTarget{sga7.sga7i.eix.n5-2-1}{5.2.1}}
 T_\ell(A_K)^t
 =
 T_\ell(A_K)^f\cap\bigl(T_\ell(A'_K)^f\bigr)^\perp,
\end{equation}
\underline{and by the analogous formula obtained by interchanging the
roles of $A_K$ and $A'_K$, where $\perp$ denotes the orthogonal
complement with respect to the canonical pairing (\SGARefLink{sga7.sga7i.eix.n1-0-2}{1.0.2}).}

}{E7I-IX-021}\SGApage{I}{\(\approx\)372}
\underline{Finally, when $A_K$ has semistable reduction over $S$
(which follows from the other hypotheses in the case $\ell=p$), one
even has the formula}
\begin{equation}\tag{\SGARefTarget{sga7.sga7i.eix.n5-2-2}{5.2.2}}
 T_\ell(A_K)^t=\bigl(T_\ell(A'_K)^f\bigr)^\perp.
\end{equation}

\smallskip
\noindent
\SGARefTarget{sga7.sga7i.eix.n5-2-3}{5.2.3}. The meaning of formula (\SGARefLink{sga7.sga7i.eix.n5-2-1}{5.2.1}) is quite clear when $\ell$ is
prime to the characteristic of $K$: as in \SGARefLink{sga7.sga7i.eix.n2-4}{2.4}, it can then be
interpreted by passing to the ordinary $\mathbb Z_\ell$-modules, with
an action of
$\pi=\operatorname{Gal}(\overline K/K)$, formed by the
$\overline K$-points of the pro-groups involved.  Let us spell it out
in the general case, using the notion of the dual
\begin{equation}\tag{\SGARefTarget{sga7.sga7i.eix.n5-2-3-1}{5.2.3.1}}
 D(G)=\operatorname{Hom}\bigl(G,T_\ell(\mathbb G_m)\bigr)
\end{equation}
of a pro-$\ell$ Barsotti--Tate group $G=(G(\nu))$.  By definition,
this is the projective system of Cartier duals $D(G(\nu))$
(SGA 3 VII 3.3.1 and VIII 1), whose transition morphisms are defined in
the evident way (for example, by transposition from the inclusion
morphisms among the $G(\nu)$).  The duality theory of abelian
varieties tells us that the pairing mentioned in \SGARefLink{sga7.sga7i.eix.n5-2}{5.2} actually defines
an \underline{isomorphism}
\begin{equation}\tag{\SGARefTarget{sga7.sga7i.eix.n5-2-3-2}{5.2.3.2}}
 T_\ell(A'_K)\simeq D\bigl(T_\ell(A_K)\bigr),
\end{equation}
and hence, by composition, a morphism
\begin{equation}\tag{\SGARefTarget{sga7.sga7i.eix.n5-2-3-3}{5.2.3.3}}
 T_\ell(A_K)^f\longrightarrow D\bigl(T_\ell(A'_K)^f\bigr).
\end{equation}
Formula (\SGARefLink{sga7.sga7i.eix.n5-2-1}{5.2.1}) says that the kernel of the preceding homomorphism
(which a priori is a pro-subgroup of its first term) is precisely the
toric part $T_\ell(A_K)^t$.  Notice that this is an
\underline{isomorphism of pro-objects, and not of projective systems}:
we do not assert (and it would in general be false) that, for every
integer $\nu$, $T_\ell(A_K)^t(\nu)$ equals the kernel of
\[
 T_\ell(A_K)^f(\nu)
 \longrightarrow
 D\bigl(T_\ell(A'_K)^f\bigr)(\nu);
\]
the latter can a priori be larger.  We shall see below, however, that
it is not larger if $A_K$ has semistable reduction; by the hypotheses
made in \SGARefLink{sga7.sga7i.eix.n5-2}{5.2}, this is in particular the case when $\ell=p$.

\SGApage{I}{373}
\smallskip
\noindent
\SGARefTarget{sga7.sga7i.eix.n5-2-4}{5.2.4}. The proof of (\SGARefLink{sga7.sga7i.eix.n5-2-1}{5.2.1}) merely repeats, with practically no change,
the proof of \SGARefLink{sga7.sga7i.eix.n2-4}{2.4}.  Observe that the induced pairing
\begin{equation}\tag{\SGARefTarget{sga7.sga7i.eix.n5-2-4-1}{5.2.4.1}}
 T_\ell(A_K)^f\times T_\ell(A'_K)^f
 \longrightarrow T_\ell(\mathbb G_{mK})
\end{equation}
which gives rise to (\SGARefLink{sga7.sga7i.eix.n5-2-3-3}{5.2.3.3}) is induced by a pairing
\begin{equation}\tag{\SGARefTarget{sga7.sga7i.eix.n5-2-4-2}{5.2.4.2}}
 T_\ell(A^o)^f\times T_\ell({A'}^o)^f
 \longrightarrow T_\ell(\mathbb G_{mS})
\end{equation}
of pro-Barsotti--Tate groups over $S$, itself induced by the pairing
of pro-groups
\begin{equation}\tag{\SGARefTarget{sga7.sga7i.eix.n5-2-4-3}{5.2.4.3}}
 T_\ell(A^o)\times T_\ell({A'}^o)
 \longrightarrow T_\ell(\mathbb G_{mS})
\end{equation}
deduced from the biextension $W$ of $(A^o,{A'}^o)$ by
$\mathbb G_{mS}$.  Consequently, the pairing
\begin{equation}\tag{\SGARefTarget{sga7.sga7i.eix.n5-2-4-4}{5.2.4.4}}
 T_\ell(A_o^o)\times T_\ell({A'}_o^o)
 \longrightarrow T_\ell(\mathbb G_{mk}),
\end{equation}
obtained from (\SGARefLink{sga7.sga7i.eix.n5-2-4-2}{5.2.4.2}) by passing to the special fibres, is precisely
the pairing deduced from the biextension $W_o$ of
$(A_o^o,{A'}_o^o)$ by $\mathbb G_{mk}$.  By VIII \SGARefLink{sga7.sga7i.eviii.n4-11}{4.11} and \SGARefLink{sga7.sga7i.eviii.n1-5}{1.5},
the left annihilator of the pairing (\SGARefLink{sga7.sga7i.eix.n5-2-4-4}{5.2.4.4}) is therefore
$T_\ell(T_o)$, which is precisely the special fibre of the toric part
$T_\ell(A^o)^t$.  Returning to the definition of left annihilators
given in \SGARefLink{sga7.sga7i.eix.n5-2-3}{5.2.3}, we conclude by means of the following lemma.

\SGApage{I}{374}
\medskip
\noindent
\underline{Lemma} \SGARefTarget{sga7.sga7i.eix.n5-2-5}{5.2.5}. \underline{Let $u:G\longrightarrow H$ be a
homomorphism of pro}-$\ell$ \underline{Barsotti--Tate groups over $S$.  If
$u_o:G_o\longrightarrow H_o$ is zero, then $u$ is zero.  If
$u_o:G_o\longrightarrow H_o$ is an isogeny (respectively an
isomorphism), then $u$ is an isogeny (respectively an isomorphism).}

Indeed, apply the first assertion of \SGARefLink{sga7.sga7i.eix.n5-2-5}{5.2.5} to the composite morphism
\[
 T_\ell(A^o)^t
 \longrightarrow T_\ell(A^o)^f
 \longrightarrow D\bigl(T_\ell({A'}^o)^f\bigr),
\]
where the second morphism is deduced from the pairing (\SGARefLink{sga7.sga7i.eix.n5-2-4-2}{5.2.4.2}).  The
composite is zero, so the pairing under consideration vanishes on the
toric part on the left.  By symmetry it also vanishes on the toric part
on the right, and hence descends to a pairing on the quotients, or
\emph{abelian pieces},
\[
 T_\ell(A^o)^{ab}\times T_\ell({A'}^o)^{ab}
 \longrightarrow T_\ell(\mathbb G_m).
\]
These quotients are again pro-Barsotti--Tate groups, as one checks
readily.  It remains to prove that the morphism corresponding to this
pairing,
\[
 T_\ell(A^o)^{ab}
 \longrightarrow D\bigl(T_\ell({A'}^o)^{ab}\bigr),
\]
is an isogeny, which again follows from \SGARefLink{sga7.sga7i.eix.n5-2-5}{5.2.5}.

We leave to the reader the easy proof of \SGARefLink{sga7.sga7i.eix.n5-2-5}{5.2.5}, a result valid over any
connected Noetherian base.  Moreover, below in \SGARefLink{sga7.sga7i.eix.n7-4-7}{7.4.7}, by a somewhat
more precise analysis of the biextensions $W_n$ of
$(A_n^o,{A'}_n^o)$, we shall find a quite evident direct proof of
the two special cases in which we had to use \SGARefLink{sga7.sga7i.eix.n5-2-5}{5.2.5}.

\SGApage{I}{375}
\smallskip
\noindent
\SGARefTarget{sga7.sga7i.eix.n5-2-6}{5.2.6}. Suppose now that $A_K$ has semistable reduction over $S$,
so that $\ell$ may be any prime whatsoever (equal or unequal to
$p$).  Moreover, by the Galois criterion for semistable reduction
\SGARefLink{sga7.sga7i.eix.n3-5}{3.5}, for $\ell\neq p$ one has
\begin{equation}\tag{\SGARefTarget{sga7.sga7i.eix.n5-2-7}{5.2.7}}
 \bigl(T_\ell(A'_K)^f\bigr)^\perp
 \ \bigl(=T_\ell(A_K)^t\bigr)
 \subset T_\ell(A_K)^f.
\end{equation}
To finish the proof of \SGARefLink{sga7.sga7i.eix.n5-2}{5.2}, it remains to show that this inclusion is
valid without restriction on $\ell$.

Indeed, the left-hand term is canonically isomorphic to
\[
 D\bigl(T_\ell(A'_K)/T_\ell(A'_K)^f\bigr)
\]
and therefore has \emph{rank} (or \emph{height})
$2n-2\alpha-\mu=\mu$, in the notation of (\SGARefLink{sga7.sga7i.eix.n2-1-9}{2.1.9}), since by
hypothesis $\lambda=0$.  This is also the rank of
$T_\ell(A_K)^t$.  Since the inclusion
\begin{equation}\tag{*}
 T_\ell(A_K)^t\subset T_\ell(A_K)
\end{equation}
is \underline{strict}, in the sense that it already induces a
monomorphism on every term of the projective systems under
consideration, it follows readily that the left-hand term of (*) is
maximal among the pro-Barsotti--Tate subgroups of $T_\ell(A_K)$
having the same rank.  Thus equality holds in (\SGARefLink{sga7.sga7i.eix.n5-2-7}{5.2.7}); by the
orthogonality theorem this is equivalent to the inclusion between the
two extreme terms of (\SGARefLink{sga7.sga7i.eix.n5-2-7}{5.2.7}).




\medskip
\noindent
\SGARefTarget{sga7.sga7i.eix.n5-3}{5.3}. Using the orthogonality correspondence between strict
pro-Barsotti--Tate subgroups (in the sense specified above) of the dual
pro-groups $T_\ell(A_K)$ and $T_\ell(A'_K)$, one formally deduces
from (\SGARefLink{sga7.sga7i.eix.n5-2-2}{5.2.2}) that the pairing (\SGARefLink{sga7.sga7i.eix.n1-0-2}{1.0.2}) between these pro-groups induces a
pairing
\begin{equation}\tag{\SGARefTarget{sga7.sga7i.eix.n5-3-2}{5.3.2}}
 T_\ell(A_K)^{ab}\times T_\ell(A'_K)^{ab}
 \longrightarrow T_\ell(\mathbb G_{mK})
\end{equation}
which is a \underline{perfect duality}, that is, it identifies each of
the two factors on the left with the dual of the other (in the sense of
\SGARefLink{sga7.sga7i.eix.n5-2-3}{5.2.3}).\footnote{The superscripts $ab$ denote the \emph{abelian
parts}, the quotient of the fixed part by the toric part.}  It is also
clear that (\SGARefLink{sga7.sga7i.eix.n5-3-2}{5.3.2}) is induced by the pairing
\begin{equation}\tag{\SGARefTarget{sga7.sga7i.eix.n5-3-3}{5.3.3}}
 T_\ell(A^o)^{ab}\times T_\ell(A^{\prime o})^{ab}
 \longrightarrow T_\ell(\mathbb G_{mS})
\end{equation}
obtained from the pairing (\SGARefLink{sga7.sga7i.eix.n5-2-4-2}{5.2.4.2}) by passing to the quotients.  Since
this is a perfect duality on the generic fibres, that is, since the
corresponding homomorphism
\begin{equation}\tag{\SGARefTarget{sga7.sga7i.eix.n5-3-4}{5.3.4}}
 T_\ell(A^o)^{ab}
 \longrightarrow D\bigl(T_\ell(A^{\prime o})^{ab}\bigr)
\end{equation}
is an isomorphism on the generic fibres, one naturally wishes to deduce
that it is so over the whole of $S$, or equivalently on the special
fibres.  This deduction is trivially valid if $\ell\neq p$.  It is
also valid, at least in unequal characteristics when $\ell=p$, by a
deep theorem of Tate [33, Th.\ 4].  We shall nevertheless establish it
directly below in \SGARefLink{sga7.sga7i.eix.n7-4-6}{7.4.6}, without using Tate's theorem.

\SGApage{I}{\(\approx\)376}
To make precise the geometric meaning of the result we wish to prove,
observe that the pairing obtained from (\SGARefLink{sga7.sga7i.eix.n5-3-3}{5.3.3}) by passage to the
special fibres is precisely the pairing
\begin{equation}\tag{\SGARefTarget{sga7.sga7i.eix.n5-3-5}{5.3.5}}
 T_\ell(B_o)\times T_\ell(B'_o)
 \longrightarrow T_\ell(\mathbb G_{mk})
\end{equation}
deduced from the biextension $W_o$ of
$(A_o^o,A_o^{\prime o})$ by passage to the quotient, where
$B_o$ and $B'_o$ denote respectively the abelian parts of
$A_o^o$ and $A_o^{\prime o}$.  By VIII \SGARefLink{sga7.sga7i.eviii.n4-8}{4.8}, the biextension
$W_o$ comes from a biextension $W_o^{ab}$ of
$(B_o,B'_o)$ by $\mathbb G_{mk}$, unique up to a unique
isomorphism.  Thus the pairing (\SGARefLink{sga7.sga7i.eix.n5-3-5}{5.3.5}) may also be interpreted as the
one deduced from the divisorial correspondence $W_o^{ab}$, and the
corresponding homomorphism
\begin{equation}\tag{\SGARefTarget{sga7.sga7i.eix.n5-3-6}{5.3.6}}
 T_\ell(B_o)
 \longrightarrow D\bigl(T_\ell(B'_o)\bigr)
 \simeq T_\ell\bigl(D(B'_o)\bigr)
\end{equation}
(where, in the last term, $D$ denotes the dual abelian scheme) is
precisely the morphism obtained by applying $T_\ell$ functorially to
the homomorphism
\begin{equation}\tag{\SGARefTarget{sga7.sga7i.eix.n5-3-7}{5.3.7}}
 B_o\longrightarrow D(B'_o)
\end{equation}
defined by $W_o^{ab}$.  Consequently, by the standard dictionary
(cf.\ [5]), saying that (5.3.5) is a perfect duality is also saying that
the kernel of (\SGARefLink{sga7.sga7i.eix.n5-3-7}{5.3.7}), which is a finite group over $k$ by \SGARefLink{sga7.sga7i.eix.n1-5}{1.5}, has
zero $\ell$-primary component.  Saying that this holds for every
$\ell$ therefore means that the isogeny (\SGARefLink{sga7.sga7i.eix.n5-3-7}{5.3.7}) is an isomorphism,
that is, that the divisorial correspondence $W_o^{ab}$ is a perfect
duality between $B_o$ and $B'_o$.  The preceding considerations
show only that the kernel of (\SGARefLink{sga7.sga7i.eix.n5-3-7}{5.3.7}) is a $p$-group.  We therefore
state the more complete result.

\SGApage{I}{377}
\SGAvar{\medskip
\noindent
\underline{Theorem} \SGARefTarget{sga7.sga7i.eix.n5-4}{5.4}. \underline{Let $S$ be a trait, and let
$A_K$ be an abelian scheme over $S$ having semistable reduction
over $S$.  Let $B_o$ and $B'_o$ denote respectively the
abelian parts of the Neron models of $A_K$ and of the dual abelian
scheme $A'_K$.  Then the canonical divisorial correspondence
$W_o^{ab}$ on $(B_o,B'_o)$, deduced by VIII \SGARefLink{sga7.sga7i.eviii.n4-8}{4.8} from the
canonical biextension $W_o$ of $(A_o^o,A_o^{\prime o})$ (the
special fibre of the biextension $W$ of \SGARefLink{sga7.sga7i.eix.n1-4}{1.4}), is a perfect duality
between $B_o$ and $B'_o$; that is, it defines an isomorphism
of each of these abelian schemes with the dual of the other.}

}{\medskip
\noindent
\underline{Theorem} \SGARefTarget{sga7.sga7i.eix.n5-4}{5.4}. \underline{Let $S$ be a trait, and let
$A_K$ be an abelian scheme over $K$ having semistable reduction
over $S$.  Let $B_o$ and $B'_o$ denote respectively the
abelian parts of the Neron models of $A_K$ and of the dual abelian
scheme $A'_K$.  Then the canonical divisorial correspondence
$W_o^{ab}$ on $(B_o,B'_o)$, deduced by VIII \SGARefLink{sga7.sga7i.eviii.n4-8}{4.8} from the
canonical biextension $W_o$ of $(A_o^o,A_o^{\prime o})$ (the
special fibre of the biextension $W$ of \SGARefLink{sga7.sga7i.eix.n1-4}{1.4}), is a perfect duality
between $B_o$ and $B'_o$; that is, it defines an isomorphism
of each of these abelian schemes with the dual of the other.}

}{E7I-IX-022}\SGApage{I}{\(\approx\)378}
The proof of \SGARefLink{sga7.sga7i.eix.n5-4}{5.4} in the general case will be given in section 7, as a
consequence of a very natural and important construction due to
Raynaud.  Let us merely observe here that, in proving \SGARefLink{sga7.sga7i.eix.n5-4}{5.4}, one may
assume that $S$ is complete, by \SGARefLink{sga7.sga7i.eix.n3-3}{3.3}.


\SGAvar{\par\medskip
\noindent
\SGARefTarget{sga7.sga7i.eix.n5-5}{5.5}. To express conveniently the extreme factors in the filtrations
(\SGARefLink{sga7.sga7i.eix.n5-1-10}{5.1.10}) and (\SGARefLink{sga7.sga7i.eix.n5-1-11}{5.1.11}), in the case of semistable reduction, let us
introduce the twisted constant group schemes over $S$
\begin{equation}\tag{\SGARefTarget{sga7.sga7i.eix.n5-5-1}{5.5.1}}
 \underline M'_o=D(T_o),\qquad \underline M_o=D(T'_o),
\end{equation}
dual to the maximal tori $T_o,T'_o$ of the special fibres of $A,A'$
(SGA 3, X 5.7). Since $S$ is complete local, these group schemes may
be regarded as the special fibres of twisted constant schemes
\begin{equation}\tag{\SGARefTarget{sga7.sga7i.eix.n5-5-2}{5.5.2}}
 \underline M',\qquad \underline M\quad\hbox{over }S.
\end{equation}
Thus, if the torus $T_o$ is split (for example, if $k$ is separably
closed), that is, if $\underline M'_o$ is the constant group defined
by the ordinary group
\begin{equation}\tag{\SGARefTarget{sga7.sga7i.eix.n5-5-3}{5.5.3}}
 M'=\operatorname{Hom}_{k\text{-}\mathrm{gr}}(T_o,\mathbb G_{mk}),
\end{equation}
then $\underline M'$ is the constant group defined by $M'$:
\begin{equation*}
 \underline M'=M'_S.
\end{equation*}
One may moreover observe that, since $T_o$ and $T'_o$ are isogenous
(\SGARefLink{sga7.sga7i.eix.n2-2-6}{2.2.6}), $T'_o$ is split whenever $T_o$ is. In this case one should
therefore also introduce
\begin{equation}\tag{\SGARefTarget{sga7.sga7i.eix.n5-5-4}{5.5.4}}
 M=\operatorname{Hom}_{k\text{-}\mathrm{gr}}(T'_o,\mathbb G_{mk}),
\end{equation}
and then
\begin{equation*}
 \underline M=M_S.
\end{equation*}
With this notation, for every integer $n$ there are evidently
isomorphisms
\begin{equation}\tag{\SGARefTarget{sga7.sga7i.eix.n5-5-5}{5.5.5}}
 \underline M'_n\simeq D(T_n),\qquad
 \underline M_n\simeq D(T'_n),
\end{equation}
where $T'_n$ is defined in terms of $A'$ as $T_n$ is in terms of $A$
(\SGARefLink{sga7.sga7i.eix.n5-1-2}{5.1.2}). Hence there follow canonical isomorphisms (where
$\mathbb Z_\ell(1)=T_\ell(\mathbb G_{mS})$)
\begin{equation}\tag{\SGARefTarget{sga7.sga7i.eix.n5-5-6}{5.5.6}}
 T_\ell(T_n)\simeq
 \check{\underline M}'_n\otimes_{\mathbb Z}\mathbb Z_\ell(1),
 \qquad
 T_\ell(T'_n)\simeq
 \check{\underline M}_n\otimes_{\mathbb Z}\mathbb Z_\ell(1),
\end{equation}
and, on passing to the limit over $n$, canonical isomorphisms
\begin{equation}\tag{\SGARefTarget{sga7.sga7i.eix.n5-5-7}{5.5.7}}
 \left\{
 \begin{aligned}
  T_\ell(A^o)^t=T_\ell(\widehat T)
   &\simeq \check{\underline M}'\otimes\mathbb Z_\ell(1),\\[1ex]
  T_\ell(A^{\prime o})^t=T_\ell(\widehat T')
   &\simeq \check{\underline M}\otimes\mathbb Z_\ell(1).
 \end{aligned}
 \right.
\end{equation}
In the case of semistable reduction, duality, together with the
orthogonality theorem \SGARefLink{sga7.sga7i.eix.n5-2}{5.2} and relation (\SGARefLink{sga7.sga7i.eix.n5-2-2}{5.2.2}), then gives canonical
isomorphisms
\begin{equation}\tag{\SGARefTarget{sga7.sga7i.eix.n5-5-8}{5.5.8}}
 \left\{
 \begin{aligned}
  T_\ell(A_K)/T_\ell(A_K)^f
   &\simeq \underline M_K\otimes_{\mathbb Z}\mathbb Z_\ell,\\[1ex]
  T_\ell(A')/T_\ell(A'_K)^f
   &\simeq \underline M'_K\otimes_{\mathbb Z}\mathbb Z_\ell.
 \end{aligned}
 \right.
\end{equation}
Likewise, the same formula (\SGARefLink{sga7.sga7i.eix.n5-2-2}{5.2.2}) gives canonical isomorphisms
\begin{equation}\tag{\SGARefTarget{sga7.sga7i.eix.n5-5-9}{5.5.9}}
 \left\{
 \begin{aligned}
  T_\ell(A_K)/T_\ell(A_K)^t
   &\simeq D\bigl(T_\ell(A'_K)^f\bigr),\\[1ex]
  T_\ell(A'_K)/T_\ell(A'_K)^t
   &\simeq D\bigl(T_\ell(A_K)^f\bigr),
 \end{aligned}
 \right.
\end{equation}
where one should recall that the groups $T_\ell(A_K)^f$ and
$T_\ell(A'_K)^f$ occurring on the right are the restrictions to
$\eta$ of the pro-$\ell$ Barsotti--Tate groups $T_\ell(A^o)^f$ and
$T_\ell(A^{\prime o})^f$, respectively, defined over $S$. Here $D$
denotes duality in the sense of pro-Barsotti--Tate groups, namely
$\underline{\operatorname{Hom}}(-,\mathbb Z_\ell(1))$. Thus we have:

}{\par\medskip
\noindent
\SGARefTarget{sga7.sga7i.eix.n5-5}{5.5}. To express conveniently the extreme factors in the filtrations
(\SGARefLink{sga7.sga7i.eix.n5-1-10}{5.1.10}) and (\SGARefLink{sga7.sga7i.eix.n5-1-11}{5.1.11}), in the case of semistable reduction, let us
introduce the twisted constant group schemes over $S$
\begin{equation}\tag{\SGARefTarget{sga7.sga7i.eix.n5-5-1}{5.5.1}}
 \underline M'_o=D(T_o),\qquad \underline M_o=D(T'_o),
\end{equation}
dual to the maximal tori $T_o,T'_o$ of the special fibres of $A,A'$
(SGA 3, X 5.7). Since $S$ is complete local, these group schemes may
be regarded as the special fibres of twisted constant schemes
\begin{equation}\tag{\SGARefTarget{sga7.sga7i.eix.n5-5-2}{5.5.2}}
 \underline M',\qquad \underline M\quad\hbox{over }S.
\end{equation}
Thus, if the torus $T_o$ is split (for example, if $k$ is separably
closed), that is, if $\underline M'_o$ is the constant group defined
by the ordinary group
\begin{equation}\tag{\SGARefTarget{sga7.sga7i.eix.n5-5-3}{5.5.3}}
 M'=\operatorname{Hom}_{k\text{-}\mathrm{gr}}(T_o,\mathbb G_{mk}),
\end{equation}
then $\underline M'$ is the constant group defined by $M'$:
\begin{equation*}
 \underline M'=M'_S.
\end{equation*}
One may moreover observe that, since $T_o$ and $T'_o$ are isogenous
(\SGARefLink{sga7.sga7i.eix.n2-2-6}{2.2.6}), $T'_o$ is split whenever $T_o$ is. In this case one should
therefore also introduce
\begin{equation}\tag{\SGARefTarget{sga7.sga7i.eix.n5-5-4}{5.5.4}}
 M=\operatorname{Hom}_{k\text{-}\mathrm{gr}}(T'_o,\mathbb G_{mk}),
\end{equation}
and then
\begin{equation*}
 \underline M=M_S.
\end{equation*}
With this notation, for every integer $n$ there are evidently
isomorphisms
\begin{equation}\tag{\SGARefTarget{sga7.sga7i.eix.n5-5-5}{5.5.5}}
 \underline M'_n\simeq D(T_n),\qquad
 \underline M_n\simeq D(T'_n),
\end{equation}
where $T'_n$ is defined in terms of $A'$ as $T_n$ is in terms of $A$
(\SGARefLink{sga7.sga7i.eix.n5-1-2}{5.1.2}). Hence there follow canonical isomorphisms (where
$\mathbb Z_\ell(1)=T_\ell(\mathbb G_{mS})$)
\begin{equation}\tag{\SGARefTarget{sga7.sga7i.eix.n5-5-6}{5.5.6}}
 T_\ell(T_n)\simeq
 \check{\underline M}'_n\otimes_{\mathbb Z}\mathbb Z_\ell(1),
 \qquad
 T_\ell(T'_n)\simeq
 \check{\underline M}_n\otimes_{\mathbb Z}\mathbb Z_\ell(1),
\end{equation}
and, on passing to the limit over $n$, canonical isomorphisms
\begin{equation}\tag{\SGARefTarget{sga7.sga7i.eix.n5-5-7}{5.5.7}}
 \left\{
 \begin{aligned}
  T_\ell(A^o)^t=T_\ell(\widehat T)
   &\simeq \check{\underline M}'\otimes\mathbb Z_\ell(1),\\[1ex]
  T_\ell(A^{\prime o})^t=T_\ell(\widehat T')
   &\simeq \check{\underline M}\otimes\mathbb Z_\ell(1).
 \end{aligned}
 \right.
\end{equation}
In the case of semistable reduction, duality, together with the
orthogonality theorem \SGARefLink{sga7.sga7i.eix.n5-2}{5.2} and relation (\SGARefLink{sga7.sga7i.eix.n5-2-2}{5.2.2}), then gives canonical
isomorphisms
\begin{equation}\tag{\SGARefTarget{sga7.sga7i.eix.n5-5-8}{5.5.8}}
 \left\{
 \begin{aligned}
  T_\ell(A_K)/T_\ell(A_K)^f
   &\simeq \underline M_K\otimes_{\mathbb Z}\mathbb Z_\ell,\\[1ex]
  T_\ell(A'_K)/T_\ell(A'_K)^f
   &\simeq \underline M'_K\otimes_{\mathbb Z}\mathbb Z_\ell.
 \end{aligned}
 \right.
\end{equation}
Likewise, the same formula (\SGARefLink{sga7.sga7i.eix.n5-2-2}{5.2.2}) gives canonical isomorphisms
\begin{equation}\tag{\SGARefTarget{sga7.sga7i.eix.n5-5-9}{5.5.9}}
 \left\{
 \begin{aligned}
  T_\ell(A_K)/T_\ell(A_K)^t
   &\simeq D\bigl(T_\ell(A'_K)^f\bigr),\\[1ex]
  T_\ell(A'_K)/T_\ell(A'_K)^t
   &\simeq D\bigl(T_\ell(A_K)^f\bigr),
 \end{aligned}
 \right.
\end{equation}
where one should recall that the groups $T_\ell(A_K)^f$ and
$T_\ell(A'_K)^f$ occurring on the right are the restrictions to
$\eta$ of the pro-$\ell$ Barsotti--Tate groups $T_\ell(A^o)^f$ and
$T_\ell(A^{\prime o})^f$, respectively, defined over $S$. Here $D$
denotes duality in the sense of pro-Barsotti--Tate groups, namely
$\underline{\operatorname{Hom}}(-,\mathbb Z_\ell(1))$. Thus we have:

}{E7I-IX-023}\SGApage{I}{380}
\medskip
\noindent
\underline{Proposition} \SGARefTarget{sga7.sga7i.eix.n5-6}{5.6}. \underline{Suppose that $S$ is complete
and that $A_K$ has semistable reduction over $S$. Then, for every
prime number $\ell$} (\underline{which may equal the residue characteristic of
$S$}), \underline{the pro}-$\ell$ \underline{Barsotti--Tate group $T_\ell(A_K)^t$ is the
generic fibre of a pro-Barsotti--Tate group over $S$ ``of toric
type'', given by (\SGARefLink{sga7.sga7i.eix.n5-5-7}{5.5.7}), whereas
$T_\ell(A_K)/T_\ell(A_K)^f$ is the generic fibre of a
pro-Barsotti--Tate group over $S$ of pro-\'{e}tale type, given by
(\SGARefLink{sga7.sga7i.eix.n5-5-8}{5.5.8}). Finally, the groups $T_\ell(A_K)^f$ and
$T_\ell(A_K)/T_\ell(A_K)^t$ are the generic fibres of
pro-Barsotti--Tate groups defined over $S$, equal respectively to
$T_\ell(A^o)^f$ and to
$D(T_\ell(A^{\prime o})^f)=
\underline{\operatorname{Hom}}(T_\ell(A^{\prime o})^f,
\mathbb Z_\ell(1))$.}

\SGAvar{\medskip
\noindent
\SGARefTarget{sga7.sga7i.eix.n5-6-1}{5.6.1}. Observe that, for an \'{e}tale covering $S'$ of $S$, splitting
$T_\ell(A_K)^t$ is equivalent to splitting
$T_\ell(A_K)/T_\ell(A_K)^f$: these conditions say respectively that
$S'$ splits the twisted constant groups $\underline M$ and
$\underline M'$, and these are isogenous because $T_o$ and $T'_o$
are. At the same time one sees that there is a smallest \'{e}tale
extension $S'$ of $S$ which does the job.

}{\medskip
\noindent
\SGARefTarget{sga7.sga7i.eix.n5-6-1}{5.6.1}. Observe that, for an \'{e}tale covering $S'$ of $S$, splitting
$T_\ell(A_K)^t$ is equivalent to splitting
$T_\ell(A_K)/T_\ell(A_K)^f$: these conditions say respectively that
$S'$ splits the twisted constant groups $\underline M'$ and
$\underline M$, and these are isogenous because $T_o$ and $T'_o$
are. At the same time one sees that there is a smallest \'{e}tale
extension $S'$ of $S$ which does the job.

}{E7I-IX-024}We also record the following easy consequence of the definitions in
\SGARefLink{sga7.sga7i.eix.n2-2-3}{2.2.3}.


\SGApage{I}{381}
\par\medskip
\noindent
\underline{Proposition} \SGARefTarget{sga7.sga7i.eix.n5-7}{5.7}. \underline{Suppose that $A^o$ is
$\ell$-divisible and that $S$ is henselian. Let $F_K$ be a
constructible $\ell$-adic sheaf over $K$ which is ``unramified'', that
is, which extends to a twisted constant $\ell$-adic sheaf $F$ over
$S$, and let $u:F_K\longrightarrow T_\ell(A_K)$ be a homomorphism of
pro-groups. Then $u$ factors through a homomorphism from $F$ to
$T_\ell(A_K)^f$.}

Indeed, write
\begin{equation}\tag{\SGARefTarget{sga7.sga7i.eix.n5-7-1}{5.7.1}}
 u=(u(\nu))_{\nu\geqslant0},\qquad
 u(\nu):F_K(\nu)\longrightarrow T_\ell(A_K)(\nu)
 ={}_{\ell^\nu}A_K.
\end{equation}
By the universal property of $A$, since $F(\nu)$ is \'{e}tale and
hence smooth over $S$, the morphism $u(\nu)$ comes from a
homomorphism
\[
 v(\nu):F(\nu)\longrightarrow{}_{\ell^\nu}A.
\]
By definition of the fixed part, this homomorphism factors as
\[
 v(\nu):F(\nu)\longrightarrow({}_{\ell^\nu}A)^f.
\]
Since the image of the induced morphism on special fibres lies in the
infinitely $\ell$-divisible part of $A_o$, it follows that $v(\nu)$
in fact factors through $({}_{\ell^\nu}A^o)^f$. Passing to the limit
over $\nu$ gives a homomorphism
\[
 v:F\longrightarrow T_\ell(A^o)^f,
\]
and the induced morphism on generic fibres plainly gives the required
factorization.

\medskip
\noindent
\underline{Remarks} \SGARefTarget{sga7.sga7i.eix.n5-7-2}{5.7.2}. When $\ell\neq p$, property \SGARefLink{sga7.sga7i.eix.n5-7}{5.7} uniquely
characterizes the sub-pro-group $T_\ell(A_K)^f$ of $T_\ell(A_K)$ as
its largest unramified $\ell$-adic subsheaf. When $\ell=p$, the
statement of \SGARefLink{sga7.sga7i.eix.n5-7}{5.7} ceases to be very satisfactory. We shall give a
considerably stronger property, which can serve as an intrinsic
characterization of the fixed part of $T_p(A_K)$ (in terms of the
pro-Barsotti--Tate group structure of $T_p(A_K)$ alone), at least in
unequal characteristics.

\SGApage{I}{382}
\medskip
\noindent
\underline{Theorem} \SGARefTarget{sga7.sga7i.eix.n5-8}{5.8} (M. Raynaud). \underline{Let $A_K$ be an
abelian scheme with semistable reduction over the function field $K$
of a henselian trait $S$. Let $\ell$ be a prime number (the
interesting case being $\ell=p$), and consider the pro-$\ell$
Barsotti--Tate group $X_\eta=T_\ell(A_K)$ and its ``fixed part''
$X_\eta^f$, which is the generic fibre of the pro-$\ell$
Barsotti--Tate group $X^f=T_\ell(A^o)^f$. If $\ell=p$, suppose that
$\operatorname{char}K=0$ (so that Tate's theorem [33, Th.\ 4] is
available over $S$). Let $F$ be a pro-$\ell$ Barsotti--Tate group over
$S$. Then every homomorphism $F_\eta\longrightarrow X_\eta$ factors
through $X_\eta^f$.}

Passing to the ind-$\ell$ groups associated with the pro-$\ell$
Barsotti--Tate groups under consideration, the preceding statement
becomes equivalent, by the cited theorem of Tate, to the following.

\medskip
\noindent
\underline{Corollary} \SGARefTarget{sga7.sga7i.eix.n5-9}{5.9}. \underline{If $F$ is an ind-$\ell$
Barsotti--Tate group, then every homomorphism
$u_\eta:F_\eta\longrightarrow A_\eta$ extends uniquely to a
homomorphism $u:F\longrightarrow A$.}

In other words, the characteristic Neron property (\SGARefLink{sga7.sga7i.eix.n1-1-2}{1.1.2}) of the
Neron model remains valid for homomorphisms from ind-group schemes of
Barsotti--Tate type, with no smoothness hypothesis on them. (Only the
case of ind-\'{e}tale groups would follow directly from the
definition.)

\SGApage{I}{383}
\medskip
\noindent
\underline{Remark} \SGARefTarget{sga7.sga7i.eix.n5-9-1}{5.9.1}. In fact, M. Raynaud proves \SGARefLink{sga7.sga7i.eix.n5-9}{5.9} without
assuming that $A_K$ has semistable reduction, and even for every Neron
group scheme $A$ over $S$ (whose generic fibre need not be an abelian
scheme). Of course, the assumption that $S$ is henselian is also
unnecessary in the formulation of \SGARefLink{sga7.sga7i.eix.n5-9}{5.9}. Raynaud's proof of this more
general case reduces it to the special case considered in \SGARefLink{sga7.sga7i.eix.n5-8}{5.8}, using
an unpublished property in the construction of Neron models. Here we
shall give only Raynaud's proof of \SGARefLink{sga7.sga7i.eix.n5-9}{5.9}.

\SGAvar{\medskip
\noindent
\underline{Proof of \SGARefLink{sga7.sga7i.eix.n5-9}{5.9}.} It is known [33] that $F$ is an extension
of an ind-\'{e}tale Barsotti--Tate group $F'$ by a Barsotti--Tate group
$F''$ with connected special fibre, so that for every integer $n$
there is an exact sequence
\[
 0\longrightarrow F''(n)\longrightarrow F(n)
 \longrightarrow F'(n)\longrightarrow0,
\]
with $F'(n)$ connected and $F''(n)_o$ purely infinitesimal.
The given homomorphism $u_\eta$ is defined by a sequence of
homomorphisms
\[
 u_\eta(n):F(n)_\eta\longrightarrow A_\eta,
\]
and it is enough to extend these separately to homomorphisms
$u(n):F(n)\longrightarrow A$. Lemma \SGARefLink{sga7.sga7i.eix.n5-9-2}{5.9.2} below then reduces us to
proving that the homomorphism
$u''_\eta:F''_\eta\longrightarrow A_\eta$ induced by $u_\eta$
extends to $u'':F''\longrightarrow A$. Consequently, \underline{we
may suppose that $F$ has connected special fibre}.

}{\medskip
\noindent
\underline{Proof of \SGARefLink{sga7.sga7i.eix.n5-9}{5.9}.} It is known [33] that $F$ is an extension
of an ind-\'{e}tale Barsotti--Tate group $F'$ by a Barsotti--Tate group
$F''$ with connected special fibre, so that for every integer $n$
there is an exact sequence
\[
 0\longrightarrow F''(n)\longrightarrow F(n)
 \longrightarrow F'(n)\longrightarrow0,
\]
with $F'(n)$ \'{e}tale and $F''(n)_o$ purely infinitesimal.
The given homomorphism $u_\eta$ is defined by a sequence of
homomorphisms
\[
 u_\eta(n):F(n)_\eta\longrightarrow A_\eta,
\]
and it is enough to extend these separately to homomorphisms
$u(n):F(n)\longrightarrow A$. Lemma \SGARefLink{sga7.sga7i.eix.n5-9-2}{5.9.2} below then reduces us to
proving that the homomorphism
$u''_\eta:F''_\eta\longrightarrow A_\eta$ induced by $u_\eta$
extends to $u'':F''\longrightarrow A$. Consequently, \underline{we
may suppose that $F$ has connected special fibre}.

}{E7I-IX-025}We may regard $u_\eta$ as a homomorphism from $F_\eta$ to the
ind-$\ell$ Barsotti--Tate group $Y_\eta$ associated with
$X_\eta=T_\ell(A_K)$, and then consider the composite
\begin{equation}\tag{*}
 v_\eta:F_\eta\longrightarrow Y_\eta
 \longrightarrow Y_\eta/Y_\eta^f.
\end{equation}
By \SGARefLink{sga7.sga7i.eix.n5-6}{5.6}, the latter group extends to an ind-\'{e}tale
Barsotti--Tate group $Z$ over $S$, namely
\[
 \underline M\otimes_{\mathbb Z}(\mathbb Q_\ell/\mathbb Z_\ell).
\]
By Tate's theorem cited in \SGARefLink{sga7.sga7i.eix.n5-8}{5.8}, the homomorphism (*) then extends to a
homomorphism
\[
 v:F\longrightarrow Z.
\]
Since $F$ has connected special fibre and $Z$ is ind-\'{e}tale, this
morphism is zero. Hence $v_\eta$ is zero as well, and $u_\eta$
therefore factors through $Y_\eta^f$, as required.

\SGApage{I}{\(\approx\)384}
It remains to prove the following lemma.

\medskip
\noindent
\underline{Lemma} \SGARefTarget{sga7.sga7i.eix.n5-9-2}{5.9.2}. Let $S$ be a trait, let $A$ be
a Neron group scheme over $S$ (\SGARefLink{sga7.sga7i.eix.n1-1-2}{1.1.2}), let
\[
 1\longrightarrow G'\longrightarrow G\longrightarrow G''
 \longrightarrow1
\]
be an extension of group schemes over $S$, with $G''$ smooth, let
$u:G'\longrightarrow A$ be a group homomorphism, and let
$v_\eta:G_\eta\longrightarrow A_\eta$ be a group homomorphism
extending $u_\eta$. Suppose that the extension of $G''$ by $A$
obtained from the extension $G$ by pushout along $u$ is representable
\footnote{Unpublished results of M. Raynaud imply, moreover, that this
condition is always satisfied.}---as is the case, for example, if
$G\longrightarrow G''$ admits a section locally for the ``finite
locally free'' topology (cf.\ SGA 3, IV 6.3) and if $A$ is
quasi-projective over $S$. Then there exists a unique group
homomorphism $v:G\longrightarrow A$ extending both $v_\eta$ and $u$.

Indeed, the last hypothesis reduces us to the case in which $G'=A$,
$u$ is the identity, and $v_\eta:G_\eta\longrightarrow A_\eta$ is a
projection which we seek to extend to a projection
$v:G\longrightarrow A$. But $G$ is smooth, being an extension of two
smooth group schemes, and the conclusion follows from the definition
of the Neron property of $A$.

\SGApage{I}{385}
\medskip
\noindent
\underline{Corollary} \SGARefTarget{sga7.sga7i.eix.n5-10}{5.10}. \underline{Let $A_K$ be an abelian
scheme over the function field} $K$ \underline{of a trait $S$, and let $\ell$ be
a prime number. When $\ell=p$, suppose that
$\operatorname{char}K=0$, so that Tate's theorem already invoked in
\SGARefLink{sga7.sga7i.eix.n5-8}{5.8} is available. Then $A_K$ has good reduction over $S$ if and only
if the pro-$\ell$ Barsotti--Tate group $T_\ell(A_K)$ ``has good
reduction over $S$'', that is, extends to a Barsotti--Tate group over
$S$.}

(This is the result conjectured in [8, 4.9].)

Necessity being clear, it remains to prove sufficiency; observe that
the case $\ell\neq p$ has already been treated in \SGARefLink{sga7.sga7i.eix.n2-2-9}{2.2.9}. We may
suppose that $S$ is henselian (\SGARefLink{sga7.sga7i.eix.n2-2-9}{2.2.9}). First suppose that $A_K$ has
semistable reduction over $S$. Applying \SGARefLink{sga7.sga7i.eix.n5-8}{5.8} with $F=X$ gives
\[
 T_\ell(A_K)=T_\ell(A_K)^f.
\]
This also says that the maximal torus of $A_o$ is zero, that is, that
$A_o^o$ is an abelian scheme; the conclusion then follows from
\SGARefLink{sga7.sga7i.eix.n2-2-9}{2.2.9} (iii${}^o$).

We now treat the general case. By \SGARefLink{sga7.sga7i.eix.n2-2-9}{2.2.9} we may plainly assume that
$S$ is strictly local. Let $K'/K$ be a finite Galois extension such
that $A_{K'}$ has semistable reduction over the normalization $S'$ of
$S$ in $K'$ (\SGARefLink{sga7.sga7i.eix.n3-6}{3.6}). If $A^*$ denotes the corresponding Neron model,
$G=\operatorname{Gal}(K'/K)$ acts on the situation $(A^*,S')$, and
hence on $A_o^{*o}$. By \SGARefLink{sga7.sga7i.eix.n4-1}{4.1} it is enough to prove that $G$ acts
trivially on the latter; this indeed implies that $A_K$ has
semistable reduction, reducing us to the case already treated.

\SGApage{I}{\(\approx\)386}
\SGAvar{Let $X$ be a Barsotti--Tate group over $S$ extending
$X_\eta=T_\ell(A_K)$. By Tate's uniqueness theorem already invoked,
there is a unique isomorphism
\[
 X\times_S S'\simeq T_\ell(A^{*o})
\]
which induces the identity on generic fibres. By transport of
structure, this isomorphism is compatible with the actions of $G$.
Now $X$ plainly acts trivially on the special fibre of
$X\times_S S'$, and hence on $T_\ell(A_o^{*o})$. It follows at once
that $G$ acts trivially on $A_o^{*o}$, completing the proof.


}{Let $X$ be a Barsotti--Tate group over $S$ extending
$X_\eta=T_\ell(A_K)$. By Tate's uniqueness theorem already invoked,
there is a unique isomorphism
\[
 X\times_S S'\simeq T_\ell(A^{*o})
\]
which induces the identity on generic fibres. By transport of
structure, this isomorphism is compatible with the actions of $G$.
Now $G$ plainly acts trivially on the special fibre of
$X\times_S S'$, and hence on $T_\ell(A_o^{*o})$. It follows at once
that $G$ acts trivially on $A_o^{*o}$, completing the proof.


}{E7I-IX-026}\medskip
\noindent
\underline{Remark} \SGARefTarget{sga7.sga7i.eix.n5-11}{5.11}. Once Tate's theorem has been proved without any
restrictive hypothesis on $\operatorname{char}K$, the preceding results
\SGARefLink{sga7.sga7i.eix.n5-8}{5.8}, \SGARefLink{sga7.sga7i.eix.n5-9}{5.9}, and \SGARefLink{sga7.sga7i.eix.n5-10}{5.10}, as well as \SGARefLink{sga7.sga7i.eix.n5-13}{5.13} below, will plainly hold without
such a restriction.

We record a definition which will also be useful elsewhere.

\medskip
\noindent
\underline{Definition} \SGARefTarget{sga7.sga7i.eix.n5-12}{5.12}. Let $S$ be a henselian trait, let $\ell$ be
a prime number, and let $X_\eta$ be a pro-$\ell$ Barsotti--Tate group
over the generic point $\eta$ of $S$. We say that $X_\eta$ has
\emph{good reduction over $S$} if it extends to a Barsotti--Tate group
over $S$ (cf. \SGARefLink{sga7.sga7i.eix.n5-10}{5.10}). We say that it has \emph{virtual good reduction
over $S$} if there exists a filtration
\[
 \bigl(\operatorname{Fil}^i(X_\eta)\bigr)_{i\geq0}
\]
by strict Barsotti--Tate subgroups\footnote{That is, subobjects as
projective systems, or equivalently subobjects term by term.} such that
the Barsotti--Tate groups
\[
 \operatorname{Gr}^i(X_\eta)=
 \operatorname{Fil}^i(X_\eta)/\operatorname{Fil}^{i+1}(X_\eta)
\]
have good reduction over $S$. If the filtration can be chosen to have
length at most $n$ (that is, if at most $n$ of the
$\operatorname{Gr}^i$ are nonzero), we say that $X_\eta$ has
\emph{virtual good reduction of level $n$ over $S$}.

\SGApage{I}{\(\approx\)387}
Finally, we say that $X_\eta$ has \emph{essential good reduction over
$S$} (respectively, \emph{essential virtual good reduction},
respectively, \emph{essential virtual good reduction of level $n$}) if
there exists a finite extension $K'/K$ such that $X_{K'}$ has good
reduction (respectively, virtual good reduction, respectively, virtual
good reduction of level $n$) over the normalization $S'$ of $S$ in
$K'$.

\medskip
\noindent
\SGARefTarget{sga7.sga7i.eix.n5-12-1}{5.12.1}. When $\ell\neq p$, the categories of pro-$\ell$
Barsotti--Tate groups over $\eta$ and over $S$ may be interpreted in
terms of $\ell$-adic representations of
\[
 \pi=\pi_1(\eta), \qquad \pi_o=\pi_1(S),
\]
respectively. The six notions just introduced can then be expressed in
terms of the natural representation of the inertia group $I$ on
$X_\eta(\overline K)$: the representation is, respectively, trivial,
unipotent, unipotent of level $n$, essentially trivial, essentially
unipotent, or essentially unipotent of level $n$.

\medskip
\noindent
\SGARefTarget{sga7.sga7i.eix.n5-12-2}{5.12.2}. Take $X_\eta=T_\ell(A_K)$, where $A_K$ is an abelian scheme
over $K$. We saw in \SGARefLink{sga7.sga7i.eix.n5-6}{5.6} that, if $A_K$ has semistable reduction over
$S$, then $X_\eta$ has virtual good reduction of level 2 (and we shall
see the converse in a moment). Consequently, by the semistable
reduction theorem \SGARefLink{sga7.sga7i.eix.n3-6}{3.6}, $X_\eta$ has essential virtual good reduction of
level 2. This statement should be regarded as the ``right'' formulation
of the geometric monodromy theorem in Expos\'e III for $H^1$ of a proper
smooth scheme over $K$, without excluding the case $\ell=p$ (which was
not considered there). The geometric monodromy theorem---or, for the
present, conjecture---for the higher $H^i$ given in Expos\'e III should
likewise be reformulated so as to include ``$p$-adic'' coefficients, as
was observed in \SGARefLink{sga7.sga7i.eix.n0-4}{0.4}.

\SGApage{I}{\(\approx\)388}
We now summarize the relations between the reduction properties of
abelian schemes over $K$ and those of the associated Barsotti--Tate
groups.

\medskip
\noindent
\underline{Proposition} \SGARefTarget{sga7.sga7i.eix.n5-13}{5.13}. Let $S$ be a henselian trait, let $A_K$
be an abelian scheme over its fraction field, and let $\ell$ be a prime
number. When $\ell=p$, assume that $K$ has characteristic zero
(cf. \SGARefLink{sga7.sga7i.eix.n5-11}{5.11}).

\SGAvar{\smallskip
\noindent
a) The abelian scheme $A$ has good reduction over $S$ (\SGARefLink{sga7.sga7i.eix.n2-2-9}{2.2.9}) if and
only if $T_\ell(A_K)$ has good reduction over $S$ (\SGARefLink{sga7.sga7i.eix.n5-12}{5.12}).

}{\smallskip
\noindent
a) The abelian scheme $A_K$ has good reduction over $S$ (\SGARefLink{sga7.sga7i.eix.n2-2-9}{2.2.9}) if and
only if $T_\ell(A_K)$ has good reduction over $S$ (\SGARefLink{sga7.sga7i.eix.n5-12}{5.12}).

}{E7I-IX-027}\smallskip
\noindent
b) The abelian scheme $A_K$ has essential good reduction over $S$---that
is, there exists a finite extension $K'/K$ such that $A_{K'}$ has good
reduction over the normalization $S'$ of $S$ in $K'$---if and only if
$T_\ell(A_K)$ has essential good reduction over $S$ (\SGARefLink{sga7.sga7i.eix.n5-12}{5.12}). In that
case $K'$ may be chosen to be a finite separable extension of $K$.

\SGAvar{\smallskip
\noindent
c) The following conditions are equivalent:
\begin{enumerate}
\item[1)] $A$ has semistable reduction over $S$ (3.4);
\item[2)] $T_\ell(A_K)$ has virtual good reduction of level 2 over $S$
  (\SGARefLink{sga7.sga7i.eix.n5-12}{5.12});
\item[3)] $T_\ell(A_K)$ has virtual good reduction over $S$.
\end{enumerate}

}{\smallskip
\noindent
c) The following conditions are equivalent:
\begin{enumerate}
\item[1)] $A_K$ has semistable reduction over $S$ (3.4);
\item[2)] $T_\ell(A_K)$ has virtual good reduction of level 2 over $S$
  (\SGARefLink{sga7.sga7i.eix.n5-12}{5.12});
\item[3)] $T_\ell(A_K)$ has virtual good reduction over $S$.
\end{enumerate}

}{E7I-IX-028}\noindent
\underline{Proof}. Part a) is just \SGARefLink{sga7.sga7i.eix.n5-10}{5.10} and has been included for
reference.

\SGAvar{For b), the first assertion follows immediately from a) and the
definitions. To prove the final assertion, we may suppose that
$\ell\neq p$. In view of a), it then suffices to take for $K'$ the
finite separable extension of $K$ corresponding to an open subgroup
$\pi'$ of
\[
 \pi=\operatorname{Gal}(\overline K/\overline K)
\]
whose intersection with the inertia group $I$ is contained in the
kernel of the representation of $I$ on
$T_\ell(A_K(\overline K))$; this kernel is open by hypothesis
(cf. \SGARefLink{sga7.sga7i.eix.n5-12-1}{5.12.1}).

}{For b), the first assertion follows immediately from a) and the
definitions. To prove the final assertion, we may suppose that
$\ell\neq p$. In view of a), it then suffices to take for $K'$ the
finite separable extension of $K$ corresponding to an open subgroup
$\pi'$ of
\[
 \pi=\operatorname{Gal}(\overline K/K)
\]
whose intersection with the inertia group $I$ is contained in the
kernel of the representation of $I$ on
$T_\ell(A_K(\overline K))$; this kernel is open by hypothesis
(cf. \SGARefLink{sga7.sga7i.eix.n5-12-1}{5.12.1}).

}{E7I-IX-029}\SGApage{I}{\(\approx\)389}
For c), we have already noted that 1) implies 2) (\SGARefLink{sga7.sga7i.eix.n5-12-2}{5.12.2}), and it is
immediate that 2) implies 3). It remains to prove that 3) implies 1).
For $\ell\neq p$, this is precisely the Galois criterion for
semistable reduction \SGARefLink{sga7.sga7i.eix.n3-5}{3.5}. For $\ell=p$, argue as follows. Choose a
prime $\ell'\neq p$ and apply the considerations of \SGARefLink{sga7.sga7i.eix.n4-1}{4.1} to it. We may
suppose that $S$ is strictly local (\SGARefLink{sga7.sga7i.eix.n3-4-0}{3.4.0}), and choose a finite Galois
extension $K'/K$, with group $G$, such that $A_{K'}$ has semistable
reduction. We must prove that the representation of $G$ on each of the
three quotients occurring in (\SGARefLink{sga7.sga7i.eix.n4-1-1}{4.1.1}) is trivial. Equivalently, their
direct-sum representation is trivial, or again its trace is the
``trivial'' trace, that is, is constant. We saw in \SGARefLink{sga7.sga7i.eix.n4-3}{4.3}(a) that this
trace function is independent of the choice of $\ell'\neq p$, and in
\SGARefLink{sga7.sga7i.eix.n6-4}{6.4} below we indicate a method of calculating it in terms of
$T_p(A_K)$ as well. That calculation plainly gives a constant function
when $T_p(A_K)$ has virtual good reduction, and this completes the
proof.

\medskip
\noindent
\underline{Remark} \SGARefTarget{sga7.sga7i.eix.n5-13-1}{5.13.1}. It seems plausible that, if $X_\eta$ is a
pro-$\ell$ Barsotti--Tate group over $\eta$ having virtual good
reduction over $S$ and essential virtual good reduction of level $n$,
then it has virtual good reduction of level $n$. (This is true when
$\ell\neq p$.)

The following result, an easy consequence of \SGARefLink{sga7.sga7i.eix.n5-7}{5.7} and of the
semistable reduction theorem \SGARefLink{sga7.sga7i.eix.n3-6}{3.6}, was pointed out to me by
J.-P. Serre.

\SGApage{I}{390}
\medskip
\noindent
\underline{Proposition} \SGARefTarget{sga7.sga7i.eix.n5-14}{5.14} (Criterion for essential good reduction).
Let $S$ be a trait with function field $K$, let $A_K$ be an abelian
scheme over $K$, and let $\ell$ be a prime number distinct from
$\operatorname{char}K$.

\smallskip
\noindent
a) Suppose that $\ell\neq p=\operatorname{char}k$. Then the following
conditions are equivalent:
\begin{enumerate}
\item[(i)] $A_K$ has essential good reduction over $S$; that is, there
  exists a finite separable extension $K'/K$ such that $A_{K'}$ has
  good reduction over the normalization $S'$ of $S$ in $K'$;
\item[(ii)] the action of the inertia group $I$ on
  $U_\ell=T_\ell(A(\overline K))$ is essentially trivial; that is, it
  is trivial on an open subgroup of $I$, or equivalently factors
  through a finite quotient $\Gamma$ of $I$;
\item[(iii)] the action of $I$ on
  $U_\ell\otimes_{\mathbb Z_\ell}\mathbb Q_\ell$ is semisimple.
\end{enumerate}

\SGAvar{\noindent
b) Suppose that $\ell=p$. If $n=\dim A_K\neq0$, then the image of $I$
in the automorphism group of $U_\ell$ is not finite; more precisely,
its image in the group $\mathbb Z_\ell$ of automorphisms of
$\bigwedge^{2n}U_\ell$ has finite index. A sufficient (but by no means
necessary) condition for $A_K$ to have essential good reduction is
that the action of $I$ on
$U_\ell\otimes_{\mathbb Z_\ell}\mathbb Q_\ell$ be semisimple.

}{\noindent
b) Suppose that $\ell=p$. If $n=\dim A_K\neq0$, then the image of $I$
in the automorphism group of $U_\ell$ is not finite; more precisely,
its image in the group $\mathbb Z_\ell^\times$ of automorphisms of
$\bigwedge^{2n}U_\ell$ has finite index. A sufficient (but by no means
necessary) condition for $A_K$ to have essential good reduction is
that the action of $I$ on
$U_\ell\otimes_{\mathbb Z_\ell}\mathbb Q_\ell$ be semisimple.

}{E7I-IX-030}\noindent
\underline{Proof}. When $\ell\neq p$, the equivalence of (i) and (ii)
follows immediately from the good-reduction criterion \SGARefLink{sga7.sga7i.eix.n2-2-9}{2.2.9}, and (ii)
implies (iii) because $\mathbb Q_\ell$ has characteristic zero. The
homomorphism
\[
 \pi=\operatorname{Gal}(\overline K/K)
 \longrightarrow \mathbb Z_\ell^\times
 \simeq\operatorname{Aut}\!\left(\bigwedge^{2n}U_\ell\right)
\]
is $\chi^n$, where $\chi$ is induced by the character
\[
 \mathbb Z_\ell^\times=\operatorname{Aut}(\mathbb Z_\ell(1))
\]
arising from the action on the $\ell^\nu$-th roots of unity. It is
known that, when $\ell=p$, the image of $I$ under this character has
finite index [33]. The fact that the action of $I$ on
$U_p\otimes\mathbb Q_p$ need not be semisimple even when $A_K$ is an
elliptic curve with good reduction is exhibited in [28, IV A 2.4].

\SGApage{I}{\(\approx\)391}
It remains to prove that (iii) implies (i), for every $\ell$. We use
the semistable reduction theorem \SGARefLink{sga7.sga7i.eix.n3-6}{3.6}, which immediately reduces us to
the case where $A_K$ has semistable reduction over $S$: hypothesis
(iii) is preserved upon replacing $K$ by a finite separable extension
$K'$, or equivalently $I$ by an open subgroup $I'$. When $\ell\neq p$,
the action of $I$ on $U_\ell\otimes\mathbb Q_\ell$ is unipotent
(\SGARefLink{sga7.sga7i.eix.n3-5}{3.5}(v)) and semisimple, hence trivial, and the good-reduction
criterion \SGARefLink{sga7.sga7i.eix.n2-2-9}{2.2.9} gives the result. We now give an argument valid with no
restriction on $\ell$, which yields a slightly sharper statement.

\medskip
\noindent
\underline{Corollary} \SGARefTarget{sga7.sga7i.eix.n5-15}{5.15}. With the notation of \SGARefLink{sga7.sga7i.eix.n5-14}{5.14}, suppose that
$A_K$ has semistable reduction over $S$, and that the fixed part
\[
 (U_{\ell\mathbb Q_\ell})^f
 \stackrel{\mathrm{def}}{=}(U_\ell^f)_{\mathbb Q_\ell}
\]
admits an $I$-stable complement in $U_{\ell\mathbb Q_\ell}$. Then
$A_K$ has good reduction.

It is enough to prove under these hypotheses that $U_\ell=U_\ell^f$.
By \SGARefLink{sga7.sga7i.eix.n5-6}{5.6} this says precisely that the toric part $T_o$ of $A_o^o$ is
zero, so that $A_o^o$ is an abelian scheme; hence $A$ is an abelian
scheme by \SGARefLink{sga7.sga7i.eix.n2-2-9}{2.2.9}. We may plainly suppose that $k$ is separably closed.
The intersection of the asserted complement in \SGARefLink{sga7.sga7i.eix.n5-15}{5.15} with $U_\ell$
then defines a pro-$\ell$ subgroup $F_K$ of $T_\ell(A_K)$ which is
isogenous to the cotoric part
\[
 T_\ell(A_K)/T_\ell(A_K)^f.
\]
The latter is pro-\'{e}tale and unramified by \SGARefLink{sga7.sga7i.eix.n5-6}{5.6}, and the same is
therefore true of $F_K$. It follows from \SGARefLink{sga7.sga7i.eix.n5-7}{5.7} that $F_K$ is contained
in $T_\ell(A_K)^f$; thus $F_K=0$, proving that $U_\ell$ is equal to
its fixed part.



\SGApage{I}{\(\approx\)392}
\medskip
\noindent
\underline{6. Remarks on the construction of the toric part and the fixed
part of} $T_p(A)$ \underline{in the case of non-semistable reduction}.

\medskip
\noindent
\SGARefTarget{sga7.sga7i.eix.n6-1}{6.1}. The construction of the toric part of $T_\ell(A_K)$ given in \SGARefLink{sga7.sga7i.eix.n5-1}{5.1}
still makes sense as soon as $S$ is complete, without assuming that
$A^o$ is $\ell$-divisible; that is (when $p=\ell$), without assuming
semistable reduction.  That hypothesis had been used only in order to
formulate the orthogonality theorem \SGARefLink{sga7.sga7i.eix.n5-2}{5.2}.  Moreover, one can construct
$T_p(A_K)^t$ while assuming only that $S$ is henselian: one shows that
the toric part $T_p(A_{\widehat K})^t$ over the completion, constructed
by the method of \SGARefLink{sga7.sga7i.eix.n5-1}{5.1}, comes from a Barsotti--Tate pro-subgroup of
$T_p(A_K)$, necessarily determined uniquely.

More precisely, one can construct a pro-$p$ Barsotti--Tate group
$T_p(A)^t$ \underline{over} $S$ and a homomorphism
\[
 T_p(A)^t\longrightarrow T_p(A)
\]
of projective systems which is a monomorphism term by term and whose
pullback to $\widehat S$ gives the analogous objects constructed by the
method of \SGARefLink{sga7.sga7i.eix.n5-1}{5.1}.  This follows by paraphrasing the construction of \SGARefLink{sga7.sga7i.eix.n5-1}{5.1},
except that, instead of lifting the whole maximal torus $T_o$ of $A_o$,
for each integer $\nu\geq0$ one merely lifts the subgroup
${}_{p^\nu}T_o$ to a multiplicative-type subgroup
$({}_{p^\nu}A^o)^t$ of ${}_{p^\nu}A^o$.

\SGApage{I}{\(\approx\)393}
For the existence and uniqueness of such a lifting, note that the
functor
\[
 (\mathrm{Sch})^o_{/S}\longrightarrow\mathrm{Sets},
 \qquad
 S'\longmapsto
 \{\text{finite multiplicative-type subgroups of }A^o_{S'}\}
\]
is represented by an $S$-scheme $X$ of finite type (cf. the proof of
SGA 3 XI 3.12(a)), which is \'etale over $S$ by SGA 3 IX 3.6 bis.
Consequently, there is exactly one section of $X$ over $S$ extending a
given section of $X_s$ over $s$, which proves the assertion.

The same argument shows more generally that, for every
\underline{finite multiplicative-type subgroup} $H_o$ of $A_o$, there
is a unique finite multiplicative-type subgroup $H$ of $A$ extending
it.  Consequently, for every \underline{toric} Barsotti--Tate
$p$-subgroup $H_o$ of $A_o$ (that is, one whose components are of
multiplicative type), there is a unique Barsotti--Tate $p$-subgroup $H$
of $A$ extending it.  This in turn gives a Barsotti--Tate $p$-subgroup
of ${}_{p^\infty}A_K$, or, equivalently, a pro-$p$ Barsotti--Tate
subgroup of $T_p(A_K)$.

In particular, this procedure applies when $H_o$ is the largest toric
Barsotti--Tate $p$-subgroup of $A_o$.  It is obtained by taking the
largest $p$-divisible subgroup $\overline A_o$ of $A_o^o$
($\overline A_o=p^N A_o^o$ for $N$ large) and the ``toric part'' of
${}_{p^\infty}\overline A_o$ (defined, for example, by SGA 3 XVII
7.2.1).  We thus obtain a Barsotti--Tate subgroup of
${}_{p^\infty}A_K$ which might be called the \underline{toric part} of
the latter.

\medskip
\noindent
\SGARefTarget{sga7.sga7i.eix.n6-2}{6.2}. Again let $\overline A_o$ be the largest $p$-divisible subgroup
scheme of $A_o^o$, and consider the Barsotti--Tate $p$-group
\[
 F_o={}_{p^\infty}(\overline A_o).
\]
It is immediate that $F_o$ may be characterized as the largest
Barsotti--Tate $p$-subgroup of $A_o$ (or of $A_o^o$).  It is natural to
ask whether this Barsotti--Tate group can be extended to a
Barsotti--Tate subgroup $F$ of $A$, which would then play the role of
the \underline{fixed part} $T_\ell(A)^f$ of (\SGARefLink{sga7.sga7i.eix.n2-2-3-2}{2.2.3.2}) in the case in
which $A^o$ is $\ell$-divisible.

\SGApage{I}{\(\approx\)394}
In the semistable case, where $\overline A_o=A_o^o$, the method of \SGARefLink{sga7.sga7i.eix.n2-2}{2.2}
answers the question affirmatively.  In the general case, whenever the
problem has a solution, it is readily seen to be unique.  The
corresponding pro-subgroups of $T_p(A_K)$ and $T_p(A'_K)$ (which again
deserve to be called the \underline{fixed parts} of the pro-groups in
question) still satisfy the orthogonality theorem \SGARefLink{sga7.sga7i.eix.n5-2}{5.2}, and the fixed
part $T_p(A_K)^f$ may still be characterized as in \SGARefLink{sga7.sga7i.eix.n5-8}{5.8}.

\medskip
\noindent
\SGARefTarget{sga7.sga7i.eix.n6-2-1}{6.2.1}. We shall nevertheless show that in general the question has no
affirmative answer; that is, \underline{the fixed part $T_p(A)^f$
cannot in general be defined} in the case of non-semistable reduction.
Indeed, by a lifting theorem of Serre--Tate [30], when $S$ is complete,
the existence of the fixed part is equivalent to the existence of
infinitesimal liftings $\overline A_n$ of $\overline A_o$ in the $A_n$,
or again to the existence of a smooth group scheme $\overline A$ over
$S$, an extension of an abelian scheme by a torus $T$, having special
fibre $\overline A_o$, together with a formal homomorphism
\[
 \widehat{\overline A}\longrightarrow\widehat A
\]
extending the inclusion on the special fibres.

Thus, when $\overline A_o$ is an abelian scheme, that is, when $T_o=0$,
this is equivalent to the existence of an abelian scheme $\overline A$
extending $\overline A_o$ and a homomorphism
$\overline A\longrightarrow A$ extending the inclusion of the special
fibres.  It will therefore suffice to give an example in which $T_o=0$,
$A_o^o$ is neither unipotent nor proper, and $A_K$ is simple over $K$.

\SGApage{I}{395}
\medskip
\noindent
\SGARefTarget{sga7.sga7i.eix.n6-2-2}{6.2.2}. Here is an example, due to M. Raynaud.  It is based on the
observation that, if $S$ is a trait, $S'$ is the normalization of $S$
in a finite separable extension $K'$ of $K$, and $A'$ is an abelian
scheme over $S'$, then the group scheme
\[
 A=\prod_{S'/S} A'/S'
\]
over $S$ is a N\'eron model: it identifies with the N\'eron model of its
generic fibre, which is indeed an abelian scheme.  This follows
immediately from the definitions and from the fact that $A'$ itself is
a N\'eron model over $S'$.  Then
\[
 A_o=\prod_{S'_s/s}(A'\times_{S'}S'_s)/S'_s
\]
is an extension of an abelian scheme, nonzero if $A'\neq0$, by a
connected unipotent scheme, nonzero if $A'\neq0$ and $S'$ is ramified
over $S$.

\SGAvar{It therefore suffices to give an example in which $A'\neq0$, $S'$ is
ramified over $S$ (a complete local trait), and
\[
 A_K=\prod_{K'/K}A'_{K'}/K'
\]
is simple over $K'$.  Take $K'/S$ to be a ramified quadratic extension
(such an extension always exists), and take $A'$ to be an elliptic
curve over $S'$.  Then $A$ has relative dimension two, and saying that
$A_K$ is simple over $K$ means that there is no nonzero homomorphism
\[
 B_K\longrightarrow A_K,
\]
where $B_K$ is an abelian scheme of dimension one over $K$.  Giving
such a homomorphism is equivalent to giving a nonzero homomorphism
\[
 B_{K'}\longrightarrow A'_{K'},
\]
which would necessarily be an isogeny.  Thus $A_K$ is simple over $K$
precisely when $A'_{K'}$ is not isogenous to an elliptic curve coming
from $K$.

}{It therefore suffices to give an example in which $A'\neq0$, $S'$ is
ramified over $S$ (a complete local trait), and
\[
 A_K=\prod_{K'/K}A'_{K'}/K'
\]
is simple over $K$.  Take $K'/K$ to be a ramified quadratic extension
(such an extension always exists), and take $A'$ to be an elliptic
curve over $S'$.  Then $A$ has relative dimension two, and saying that
$A_K$ is simple over $K$ means that there is no nonzero homomorphism
\[
 B_K\longrightarrow A_K,
\]
where $B_K$ is an abelian scheme of dimension one over $K$.  Giving
such a homomorphism is equivalent to giving a nonzero homomorphism
\[
 B_{K'}\longrightarrow A'_{K'},
\]
which would necessarily be an isogeny.  Thus $A_K$ is simple over $K$
precisely when $A'_{K'}$ is not isogenous to an elliptic curve coming
from $K$.

}{E7I-IX-031}If the contrary held, the $j$-invariant of $A'$ would be algebraically
dependent (over the prime field) on its conjugate $\overline j$ under
the nontrivial $K$-automorphism of $K'$.  If $j$ is chosen so that
\[
 a=j+\overline j=\operatorname{Tr}_{K'/K}j
\]
is transcendental over the prime field, this is equivalent to saying
that $j$ is algebraically dependent on $a$.  By cardinality
considerations, however, there exists an element $a\in V$
transcendental over the prime subfield $P$ of $K$, and then an element
$b\in V'$ (the normalization of $V$ in $K'$) which is transcendental
over $P(a)$ and satisfies
\[
 \operatorname{Tr}_{V'/V}b=0.
\]
It is enough to take $j=a+b\in V'$.  We may moreover suppose that
$b\in\mathfrak r(V')$, and that the image of $a$ in
$k=V/\mathfrak m$ is a prescribed element of $k$, chosen distinct from
the exceptional values $j_o=0$ and $j_o=12^3$.  This is possible even
when $k$ has only two elements, because then these two contemplated
values of $j_o$ coincide.  Thus the image $j_o$ of $j$ in $k(s')$ is
different from $0$ and $12^3$, and it is known that there then exists
an elliptic curve over $S'$ with $j$-invariant $j$.

\SGApage{I}{396}
\medskip
\noindent
\SGARefTarget{sga7.sga7i.eix.n6-3}{6.3}. Return to the case in which $S$ is merely assumed henselian.  In
analogy with the definitions of \SGARefLink{sga7.sga7i.eix.n4-1}{4.1} for $\ell\neq p$, one can then
define a filtration
\begin{equation}\tag{\SGARefTarget{sga7.sga7i.eix.n6-3-1}{6.3.1}}
 T_p(A_K)\supset T_p(A_K)^{\mathrm{ef}}
 \supset T_p(A_K)^{\mathrm{et}}\supset0
\end{equation}
by strict Barsotti--Tate pro-subgroups, called respectively the
\underline{essentially fixed part} and the
\underline{essentially toric part} of $T_p(A_K)$.

\SGAvar{To do so, choose a finite Galois extension $K'$ of $K$ such that
$A_{K'}$ has semistable reduction over the normalization $S'$ of $S$
in $K'$.  The fixed part $T_p(A_{K'})^f$ of
$T_p(A_{K'})=(T_p(A_K))_{K'}$ is then defined (\SGARefLink{sga7.sga7i.eix.n2-2-3}{2.2.3}), as is the toric
part (\SGARefLink{sga7.sga7i.eix.n6-1}{6.1}).  The latter is moreover the orthogonal complement of the
fixed part $T_p(A'_{K})^f$ in $T_p(A'_{K})$ (\SGARefLink{sga7.sga7i.eix.n5-2}{5.2}).  Plainly
\[
 T_p(A_{K'})^t\subset T_p(A_{K'})^f
\]
(this follows from \SGARefLink{sga7.sga7i.eix.n5-3}{5.3}, or directly from the construction of the toric
part in \SGARefLink{sga7.sga7i.eix.n6-1}{6.1}), and the Galois group of $K'/K$ preserves these two strict
pro-subgroups of $(T_p(A_K))_{K'}$.  Descent therefore gives the desired
filtration (\SGARefLink{sga7.sga7i.eix.n6-3-1}{6.3.1}).  It is independent of the choice of $K'$, by the
semistable-reduction property \SGARefLink{sga7.sga7i.eix.n3-3}{3.3}.  Descent also plainly gives the
orthogonality property
\begin{equation}\tag{\SGARefTarget{sga7.sga7i.eix.n6-3-2}{6.3.2}}
 \begin{split}
 T_p(A_K)^{\mathrm{et}}
   &=\bigl(T_p(A'_k)^{\mathrm{ef}}\bigr)^\perp,\\
 T_p(A_K)^{\mathrm{ef}}
   &=\bigl(T_p(A'_k)^{\mathrm{et}}\bigr)^\perp.
 \end{split}
\end{equation}

One could give an intrinsic characterization of the filtration (\SGARefLink{sga7.sga7i.eix.n6-3-1}{6.3.1})
in terms of the Barsotti--Tate pro-group $X=T_p(A_K)$ over the fraction
field $K$ of $S$, in the spirit of \SGARefLink{sga7.sga7i.eix.n5-8}{5.8}---at least, for the moment, in
unequal characteristics.  The essentially fixed part is the largest
Barsotti--Tate pro-subgroup having ``essentially good reduction over
$S$'': after a finite separable extension $K'$ of $K$, it gives a
Barsotti--Tate pro-group $X_{K'}$ having good reduction over the
normalization $S'$ of $S$ in $K$, that is, extending to a
Barsotti--Tate pro-group over $S'$.  Such an extension is necessarily
unique up to a unique isomorphism, by the already-cited theorem of Tate
[33, Th. 4].  The essentially toric part is then defined in terms of
the essentially fixed part by orthogonality (\SGARefLink{sga7.sga7i.eix.n6-3-2}{6.3.2}).

Let us also record the following immediate consequence of \SGARefLink{sga7.sga7i.eix.n5-8}{5.8} and the
definitions; for $\ell=p$ it reduces to \SGARefLink{sga7.sga7i.eix.n5-12}{5.12}(a).

\medskip
\noindent
\underline{Proposition} \SGARefTarget{sga7.sga7i.eix.n6-3-3}{6.3.3}. \underline{Let $\ell$ be a prime
number.  If $\ell=p$, assume that $\operatorname{char}K=0$.}
\footnote{This hypothesis is probably unnecessary; cf. \SGARefLink{sga7.sga7i.eix.n5-11}{5.11}.}
\underline{Then $A_K$ has essential good reduction over $S$ (\SGARefLink{sga7.sga7i.eix.n5-12}{5.12}(a))
if and only if its essentially fixed part}
$T_\ell(A_K)^{\mathrm{ef}}$
\underline{(\SGARefLink{sga7.sga7i.eix.n4-1-1}{4.1.1} and \SGARefLink{sga7.sga7i.eix.n6-3-1}{6.3.1}) is all of $T_\ell(A_K)$.}

}{To do so, choose a finite Galois extension $K'$ of $K$ such that
$A_{K'}$ has semistable reduction over the normalization $S'$ of $S$
in $K'$.  The fixed part $T_p(A_{K'})^f$ of
$T_p(A_{K'})=(T_p(A_K))_{K'}$ is then defined (\SGARefLink{sga7.sga7i.eix.n2-2-3}{2.2.3}), as is the toric
part (\SGARefLink{sga7.sga7i.eix.n6-1}{6.1}).  The latter is moreover the orthogonal complement of the
fixed part $T_p(A'_{K'})^f$ in $T_p(A'_{K'})$ (\SGARefLink{sga7.sga7i.eix.n5-2}{5.2}).  Plainly
\[
 T_p(A_{K'})^t\subset T_p(A_{K'})^f
\]
(this follows from \SGARefLink{sga7.sga7i.eix.n5-3}{5.3}, or directly from the construction of the toric
part in \SGARefLink{sga7.sga7i.eix.n6-1}{6.1}), and the Galois group of $K'/K$ preserves these two strict
pro-subgroups of $(T_p(A_K))_{K'}$.  Descent therefore gives the desired
filtration (\SGARefLink{sga7.sga7i.eix.n6-3-1}{6.3.1}).  It is independent of the choice of $K'$, by the
semistable-reduction property \SGARefLink{sga7.sga7i.eix.n3-3}{3.3}.  Descent also plainly gives the
orthogonality property
\begin{equation}\tag{\SGARefTarget{sga7.sga7i.eix.n6-3-2}{6.3.2}}
 \begin{split}
 T_p(A_K)^{\mathrm{et}}
   &=\bigl(T_p(A'_K)^{\mathrm{ef}}\bigr)^\perp,\\
 T_p(A_K)^{\mathrm{ef}}
   &=\bigl(T_p(A'_K)^{\mathrm{et}}\bigr)^\perp.
 \end{split}
\end{equation}

One could give an intrinsic characterization of the filtration (\SGARefLink{sga7.sga7i.eix.n6-3-1}{6.3.1})
in terms of the Barsotti--Tate pro-group $X=T_p(A_K)$ over the fraction
field $K$ of $S$, in the spirit of \SGARefLink{sga7.sga7i.eix.n5-8}{5.8}---at least, for the moment, in
unequal characteristics.  The essentially fixed part is the largest
Barsotti--Tate pro-subgroup having ``essentially good reduction over
$S$'': after a finite separable extension $K'$ of $K$, it gives a
Barsotti--Tate pro-group $X_{K'}$ having good reduction over the
normalization $S'$ of $S$ in $K'$, that is, extending to a
Barsotti--Tate pro-group over $S'$.  Such an extension is necessarily
unique up to a unique isomorphism, by the already-cited theorem of Tate
[33, Th. 4].  The essentially toric part is then defined in terms of
the essentially fixed part by orthogonality (\SGARefLink{sga7.sga7i.eix.n6-3-2}{6.3.2}).

Let us also record the following immediate consequence of \SGARefLink{sga7.sga7i.eix.n5-8}{5.8} and the
definitions; for $\ell=p$ it reduces to \SGARefLink{sga7.sga7i.eix.n5-12}{5.13}(b).

\medskip
\noindent
\underline{Proposition} \SGARefTarget{sga7.sga7i.eix.n6-3-3}{6.3.3}. \underline{Let $\ell$ be a prime
number.  If $\ell=p$, assume that $\operatorname{char}K=0$.}
\footnote{This hypothesis is probably unnecessary; cf. \SGARefLink{sga7.sga7i.eix.n5-11}{5.11}.}
\underline{Then $A_K$ has essential good reduction over $S$ (\SGARefLink{sga7.sga7i.eix.n5-12}{5.13}(b))
if and only if its essentially fixed part}
$T_\ell(A_K)^{\mathrm{ef}}$
\underline{(\SGARefLink{sga7.sga7i.eix.n4-1-1}{4.1.1} and \SGARefLink{sga7.sga7i.eix.n6-3-1}{6.3.1}) is all of $T_\ell(A_K)$.}

}{E7I-IX-032}\SGApage{I}{398}
\medskip
\noindent
\SGARefTarget{sga7.sga7i.eix.n6-4}{6.4}. Recall that, when $\ell\neq p$, in \SGARefLink{sga7.sga7i.eix.n4-3}{4.3}(a) we exhibited a
remarkable integer-valued function on the inertia group $I$ of $S$,
independent of the choice of $\ell$.  When the residue field is perfect,
knowledge of this function suffices to reconstruct the conductor
exponent of $A$ (\SGARefLink{sga7.sga7i.eix.n4-5}{4.5}): one takes, on a sufficiently large quotient
$\Gamma=I/I'$ of $I$, the inner product of this function with the Artin
conductor on $\Gamma$ (cf. (\SGARefLink{sga7.sga7i.eix.n4-6-2}{4.6.2})).  It is useful to observe that this
integer-valued function on $I$ can also be recovered from the knowledge
of $T_\ell(A_K)$ for $\ell=p$---at least, for the time being, in unequal
characteristics.

\SGAvar{When the residue field $k$ is finite, an analogous observation applies
to the $\mathbb Q$-valued function considered in \SGARefLink{sga7.sga7i.eix.n4-3}{4.3}(b), defined on the
subgroup $\pi'$ of $\pi$ consisting of the elements whose image in
$\pi_o$ has the form $f_q^{\,n}$, $n\in\mathbb Z$.  We indicate the
principle of the verification.  Using the filtration (\SGARefLink{sga7.sga7i.eix.n6-3-1}{6.3.1}), one is
reduced to attaching a suitable function (on $I$, respectively on
$\pi'$) to a pro-$p$ Barsotti--Tate group $Y$ over $K$ having
``essentially good reduction over $S$'' in the sense of \SGARefLink{sga7.sga7i.eix.n5-12}{6.5}.  Here $Y$
will be one of the three quotients arising from the filtration (\SGARefLink{sga7.sga7i.eix.n6-3-1}{6.3.1}).

}{When the residue field $k$ is finite, an analogous observation applies
to the $\mathbb Q$-valued function considered in \SGARefLink{sga7.sga7i.eix.n4-3}{4.3}(b), defined on the
subgroup $\pi'$ of $\pi$ consisting of the elements whose image in
$\pi_o$ has the form $f_q^{\,n}$, $n\in\mathbb Z$.  We indicate the
principle of the verification.  Using the filtration (\SGARefLink{sga7.sga7i.eix.n6-3-1}{6.3.1}), one is
reduced to attaching a suitable function (on $I$, respectively on
$\pi'$) to a pro-$p$ Barsotti--Tate group $Y$ over $K$ having
``essentially good reduction over $S$'' in the sense of \SGARefLink{sga7.sga7i.eix.n5-12}{5.12}.  Here $Y$
will be one of the three quotients arising from the filtration (\SGARefLink{sga7.sga7i.eix.n6-3-1}{6.3.1}).

}{E7I-IX-033}In the case of \SGARefLink{sga7.sga7i.eix.n4-3}{4.3}(a), a suitable residue-field extension within $S$
(EGA $0_{\mathrm{III}}$ 10.3.1) reduces us to the case of a perfect
residue field, or even an algebraically closed one if desired.  Choose
a finite Galois extension $K'$ of $K$ such that $Y_{K'}$ has good
reduction over $S'$.  We then have a Barsotti--Tate group $Y_{S'}$ over
$S'$ and an action of $G$ on $Y_{S'}$ compatible with its action on
$S'$.  If $k'$ is the residue field of $S'$, then $G$ acts, not
necessarily faithfully, on $k'$ (whose invariant field is $k$), and on
$Y_{k'}$ compatibly with its action on $k'$.

\SGApage{I}{\(\approx\)399}
By Dieudonn\'e theory [16 bis], the Barsotti--Tate pro-$p$ group $Y_{k'}$
has a canonically associated finite free module $M$ (of rank equal to
the rank of $Y_{k'}$) over the ring $W(k')$ of Witt vectors of $k'$.
The group $G$ plainly acts on $M$ compatibly with its actions on
$W(k')$, through its actions on $k'$.  Hence the inertia subgroup
$\Gamma$ of $G$ acts $W(k')$-linearly on $M$; since $\Gamma$ is a
quotient of the inertia group $I$, this gives a trace function on
$\Gamma$, and hence on $I$, which is the desired function.

A priori this function takes values in $W(k')$,
\footnote{In fact it is immediate that this function always takes values
in $W(k)$, since it is invariant under $\operatorname{Gal}(k'/k)$.}
but it agrees with the analogous function defined in terms of the
$T_\ell$ for $\ell\neq p$ whenever the action of $G$ on $Y_{S'}$ is
obtained from an action of $G$ on a smooth $p$-divisible group scheme
$C'$ over $S'$ by taking $Y=T_p(C')$ (cf. [25 bis, Th. 7, p. 737]).
It is therefore integer-valued in this case.  The case of \SGARefLink{sga7.sga7i.eix.n4-3}{4.3}(b) is
handled similarly: to define a trace function on $\pi'$, one uses the
method of the proof of \SGARefLink{sga7.sga7i.eix.n4-3-2}{4.3.2} to ``linearize'' on $M$ the action defined
by an element $g\in\pi$ lying above an $f_q^n$, with $n\geq0$.



\medskip
\noindent
7. \underline{The Raynaud extension $G^\natural$ over $S$ associated
with the abelian scheme $A_K$ having semistable reduction.  Duality of
the abelian schemes $B,B'$ over $S$.}

\medskip
\noindent
\SGARefTarget{sga7.sga7i.eix.n7-1}{7.1}. Let $S$ be a Noetherian scheme which is the spectrum of a ring
$V$, separated and complete for the topology defined by an ideal
$\underline m$ of $V$.  Let $A$ be a smooth group scheme over $S$, let
$A^o$ be its identity component, and suppose that $A_o/A_o^o$ is
finite (hence finite etale) over $S_o$, and that $A_o^o$ is an
extension of an abelian scheme by a torus.  As usual, we have put
$A_n=A\times_S S_n$, where
$S_n=\operatorname{Spec}(V/\underline m^{n+1})$.  We shall be
interested chiefly in the case where $S$ is a trait, the generic fiber
$A_K$ is an abelian scheme having semistable reduction over $S$, and
$A$ is the Neron model of $A_K$.

\SGApage{I}{\(\approx\)400}
Let $T_o$ be the maximal torus of $A_o$, so that
$B_o=A_o^o/T_o$ is an abelian scheme.  The construction of \SGARefLink{sga7.sga7i.eix.n5-1}{5.1} then
applies and gives a formal subtorus $\widehat T$ of the formal
completion $\widehat A$.  For every $n$, the quotient
\begin{equation}\tag{\SGARefTarget{sga7.sga7i.eix.n7-1-1}{7.1.1}}
 B_n=A_n^o/T_n
\end{equation}
is representable (SGA~3 $\mathrm{VI}_A$~3.2); more precisely, it is a
flat scheme of finite type over $S_n$ whose pullback to $S_o$ is the
abelian scheme $B_o$.  Consequently $B_n$ is an \underline{abelian
scheme} over $S_n$.  As $n$ varies, the $B_n$ are obtained from one
another by base change and therefore define a \underline{formal
abelian scheme}
\begin{equation}\tag{\SGARefTarget{sga7.sga7i.eix.n7-1-2}{7.1.2}}
 \widehat B=(B_n)_{n\geqslant0}\,.
\end{equation}
We shall give conditions implying that this formal abelian scheme is
\underline{algebraizable} (EGA~III~5.4.1), more precisely, that it is
the formal completion of a projective abelian scheme over $S$, which
we shall then denote by $B$.  By loc. cit., $B$ is determined up to a
unique isomorphism (as an $S$-scheme; the group structure of
$\widehat B$ then comes uniquely from a group structure on $B$, again
by loc. cit.).  This abelian scheme $B$ will then be called the
\underline{abelian part} of the scheme $A$, and will sometimes be
denoted by $A^{\mathrm{ab}}$.

\SGApage{I}{401}
\SGAvar{\medskip
\noindent
\SGARefTarget{sga7.sga7i.eix.n7-1-3}{7.1.3}. By EGA~III~5.4.5, projectivity of $\widehat B$ means that one
can find an invertible sheaf $\widehat{\mathbb M}$ on $\widehat B$,
that is, a sequence of invertible sheaves $\mathbb M_n$ on the $B_n$
which glue as $n$ varies, and which is relatively ample over $S$, that
is, such that $\mathbb M_o$ is ample.  The natural idea would be to
start with an ample invertible sheaf $\mathbb M$ on $G$ and show that,
for every $n$, its restriction $\mathbb L_n$ to $G_n$ is the pullback
of an invertible sheaf $\mathbb M_n$ on $B_n$.  I do not know whether such
$\mathbb M_n$ necessarily exist (and strongly doubt it).  We shall
follow a slightly different method, constructing a
\underline{polarization} of $ B$, that is, a projective system
of polarizations $\rho_n$ of the $B_n$ which glue.  It is known [5]
that the doubles of these polarizations canonically define
(symmetric) invertible sheaves $\mathbb M_n$ on the $B_n$.  By their
canonical nature these sheaves necessarily glue and furnish the
desired $\mathbb M$.

}{\medskip
\noindent
\SGARefTarget{sga7.sga7i.eix.n7-1-3}{7.1.3}. By EGA~III~5.4.5, projectivity of $\widehat B$ means that one
can find an invertible sheaf $\widehat{\mathbb M}$ on $\widehat B$,
that is, a sequence of invertible sheaves $\mathbb M_n$ on the $B_n$
which glue as $n$ varies, and which is relatively ample over $S$, that
is, such that $\mathbb M_o$ is ample.  The natural idea would be to
start with an ample invertible sheaf $\mathbb L$ on $A$ and show that,
for every $n$, its restriction $\mathbb L_n$ to $A_n$ is the pullback
of an invertible sheaf $\mathbb M_n$ on $B_n$.  I do not know whether such
$\mathbb M_n$ necessarily exist (and strongly doubt it).  We shall
follow a slightly different method, constructing a
\underline{polarization} of $\widehat B$, that is, a projective system
of polarizations $\rho_n$ of the $B_n$ which glue.  It is known [5]
that the doubles of these polarizations canonically define
(symmetric) invertible sheaves $\mathbb M_n$ on the $B_n$.  By their
canonical nature these sheaves necessarily glue and furnish the
desired $\widehat{\mathbb M}$.

}{E7I-IX-034}We shall interpret an element of $\operatorname{NS}(B_n/S_n)$ as a
symmetric divisorial correspondence on $(B_n,B_n)$ [5], that is, as a
biextension of $(B_n,B_n)$ by $\mathbb G_{mS_n}$.  By VIII~\SGARefLink{sga7.sga7i.eviii.n3-5}{3.5}, these
biextensions correspond bijectively to biextensions of $(A_n,A_n)$ by
$\mathbb G_{mS_n}$.  To obtain a coherent system of such biextensions
$W_n$, it suffices to start with
\begin{equation}\tag{\SGARefTarget{sga7.sga7i.eix.n7-1-4}{7.1.4}}
 W,\quad\text{a symmetric biextension of $(A,A)$ by $\mathbb G_{mS}$},
\end{equation}
and take the $W_n$ obtained from $W$ by the base changes
$S_n\longrightarrow S$.  We must in addition suppose that the element
of $\operatorname{NS}(B_o)$ defined by $W_o$ is a polarization of
$B_o$; we thereby obtain a polarization of $\widehat B$.

\SGApage{I}{402}
\medskip
\noindent
\underline{Proposition \SGARefTarget{sga7.sga7i.eix.n7-1-5}{7.1.5}.} \underline{With the preceding notation
for $S,A,\widehat B$, suppose that $S$ is normal.  Then $\widehat B$ is
algebraizable, and the abelian scheme $B$ over $S$ from which it comes
is projective over $S$.}

We may plainly suppose that $A=A^o$ and that $S$ is integral.  By
Raynaud [23, XI~1.13], the hypothesis that $S$ is normal implies that
$A$ is quasi-projective over $S$.  Choose a relatively ample
invertible sheaf $\mathbb L$ on $A$, and form the birigidified
invertible sheaf
\[
 \delta(\mathbb L)=
 \pi^*(\mathbb L)\operatorname{pr}_1^*(\mathbb L)^{-1}
 \operatorname{pr}_2^*(\mathbb L)^{-1}
\]
on $A\times_S A$, as in VIII~\SGARefLink{sga7.sga7i.eviii.n4-12}{4.12}.  Since the fibers of $A_o$ have
unipotent rank zero, the same is true of the fibers of $A$ (this
follows easily from \SGARefLink{sga7.sga7i.eix.n3-1}{3.1}).  Therefore, by VIII~\SGARefLink{sga7.sga7i.eviii.n7-5}{7.5},
$\delta(\mathbb L)$ defines a biextension $W$ of $(A,A)$, plainly
symmetric, which by VIII~\SGARefLink{sga7.sga7i.eviii.n4-13}{4.13} (that is, essentially by Raynaud
[23, XI~1.11]) satisfies the required ampleness condition at the
points of $S_o$.

\medskip
\noindent
\underline{Remark \SGARefTarget{sga7.sga7i.eix.n7-1-6}{7.1.6}.} Without assuming that $S$ is normal, one
can prove that $\widehat B$ is always algebraizable by reducing to the
normal case as in \SGARefLink{sga7.sga7i.eix.n13-5-3}{13.5.3}(a).



\medskip
\noindent
\SGARefTarget{sga7.sga7i.eix.n7-2}{7.2}. We now record the following lemma.

\SGAvar{\medskip
\noindent
\underline{Lemma \SGARefTarget{sga7.sga7i.eix.n7-2-1}{7.2.1}.} Let $S$ be the spectrum of a Noetherian ring
$V$, separated and complete for the topology defined by an ideal
$\underline m$ of $V$.  Let $C$ be the category of group schemes
$\Lambda$ over $S$, smooth and of finite type over $S$, such that
$\Lambda/\Lambda^o$ is a finite etale group scheme over $S$ and
$\Lambda^o$ is an extension of an abelian scheme $B$ by an isotrivial
torus $T$ (SGA~3 IX~1.1).  Let $C'$ be the category of group objects
in the category of formal schemes of finite type (EGA~I~10.3.3) over
$\widehat S=\operatorname{Spf}(A)$, and consider the functor
\[
 A\longmapsto\widehat A
\]
from $C$ to $C'$ which takes the formal completion along
$A_o=A\times_S S_o$, where
$S_o=\operatorname{Spec}(V/\underline m)$. Then:

}{\medskip
\noindent
\underline{Lemma \SGARefTarget{sga7.sga7i.eix.n7-2-1}{7.2.1}.} Let $S$ be the spectrum of a Noetherian ring
$V$, separated and complete for the topology defined by an ideal
$\underline m$ of $V$.  Let $C$ be the category of group schemes
$\Lambda$ over $S$, smooth and of finite type over $S$, such that
$\Lambda/\Lambda^o$ is a finite etale group scheme over $S$ and
$\Lambda^o$ is an extension of an abelian scheme $B$ by an isotrivial
torus $T$ (SGA~3 IX~1.1).  Let $C'$ be the category of group objects
in the category of formal schemes of finite type (EGA~I~10.13.3) over
$\widehat S=\operatorname{Spf}(V)$, and consider the functor
\[
 A\longmapsto\widehat A
\]
from $C$ to $C'$ which takes the formal completion along
$A_o=A\times_S S_o$, where
$S_o=\operatorname{Spec}(V/\underline m)$. Then:

}{E7I-IX-035}\SGApage{I}{\(\approx\)403}
\smallskip
\noindent
\underline{(a) The preceding functor is fully faithful.}

\SGAvar{\smallskip
\noindent
\underline{(b) Let $\mathcal A$ be an object of $C'$, smooth over
$\widehat S$}; that is, suppose that the ordinary schemes
\[
 A_n=A\times_{\widehat S}S_n,
 \qquad
 S_n=\operatorname{Spec}(A/\underline m^{n+1}),
\]
are smooth over $S_n$.  Suppose that $A_o/A_o^o$ is a finite etale
group scheme over $S_o$, and that $A_o^o$ is an extension of an
abelian scheme $B_o$ by an isotrivial torus $T_o$.  Put
\[
 \mathcal A^o=(A_n^o)_{n\geq0},
 \qquad
 \mathcal A/\mathcal A^o=(A_n/A_n^o)_{n\geq0}.
\]
\underline{Then $\mathcal A/\mathcal A^o$ is finite etale over
$\widehat S$, and $\mathcal A^o$ is, in an essentially unique manner
and functorially in $\mathcal A$, an extension of a formal abelian
scheme $\mathcal B$ over $\widehat S$ by a formal torus $\mathcal T$
over $\widehat S$, reducing respectively to $B_o$ and $T_o$.  The
formal group scheme $\mathcal A$ belongs to the essential image of
$C\longrightarrow C'$---or, as we shall also say, is
``algebraizable''---if and only if $\mathcal B$ is algebraizable.}

}{\smallskip
\noindent
\underline{(b) Let $\mathcal A$ be an object of $C'$, smooth over
$\widehat S$}; that is, suppose that the ordinary schemes
\[
 A_n=\mathcal A\times_{\widehat S}S_n,
 \qquad
 S_n=\operatorname{Spec}(V/\underline m^{n+1}),
\]
are smooth over $S_n$.  Suppose that $A_o/A_o^o$ is a finite etale
group scheme over $S_o$, and that $A_o^o$ is an extension of an
abelian scheme $B_o$ by an isotrivial torus $T_o$.  Put
\[
 \mathcal A^o=(A_n^o)_{n\geq0},
 \qquad
 \mathcal A/\mathcal A^o=(A_n/A_n^o)_{n\geq0}.
\]
\underline{Then $\mathcal A/\mathcal A^o$ is finite etale over
$\widehat S$, and $\mathcal A^o$ is, in an essentially unique manner
and functorially in $\mathcal A$, an extension of a formal abelian
scheme $\mathcal B$ over $\widehat S$ by a formal torus $\mathcal T$
over $\widehat S$, reducing respectively to $B_o$ and $T_o$.  The
formal group scheme $\mathcal A$ belongs to the essential image of
$C\longrightarrow C'$---or, as we shall also say, is
``algebraizable''---if and only if $\mathcal B$ is algebraizable.}

}{E7I-IX-036}\medskip
\noindent
\SGARefTarget{sga7.sga7i.eix.n7-2-1-1}{7.2.1.1}. First observe, over an arbitrary scheme $S$, that if $A$ is
a smooth group scheme such that $A^o$ is an extension of an abelian
scheme by a torus, this extension structure is essentially unique and
is functorial in $A$ in the evident sense.  Indeed, the assertion
reduces at once to the fact that every morphism from a torus $T$ to an
abelian scheme $B$ is zero.  This is well known on geometric fibers
[15], and SGA~3 IX~5.3 then gives the conclusion.  Applying the result
to the compatible systems of smooth group schemes $A_n$ over the
$S_n$ occurring in \SGARefLink{sga7.sga7i.eix.n7-2-1}{7.2.1} gives the analogous result for formal group
schemes over $S$.  Moreover, we already observed in \SGARefLink{sga7.sga7i.eix.n7-1}{7.1} that, if
$A_o^o$ is an extension of an abelian scheme by a torus, the same is
true of $A_n^o$.  This proves the first assertion of (b).

\SGApage{I}{404}
\medskip
\noindent
\SGARefTarget{sga7.sga7i.eix.n7-2-1-2}{7.2.1.2}. Suppose next that the formal group scheme $\mathcal A$ over
$S$ is of the form $\widehat A$, with $A\in\operatorname{ob}C$.  If
$A^o$ is an extension of an abelian scheme $B$ by a torus $T$, then
plainly $\mathcal B\simeq\widehat B$.  Thus algebraizability of
$\mathcal A$ implies algebraizability of $\mathcal B$.

\SGAvar{Conversely, suppose that $\mathcal B$ is algebraizable, that is, that
it comes from an abelian scheme $B$, and let us prove that $\mathcal A$
is algebraizable.  We first prove that
$\mathcal A^o=(A_n^o)$ is algebraizable.  By VIII~\SGARefLink{sga7.sga7i.eviii.n3-7}{3.7}, the extensions
$A_n^o$ of $B_n$ by $T_n$ may be interpreted as homomorphisms
\begin{equation}\tag{*}
 \underline M_n\longrightarrow B'_n,
\end{equation}
where $\underline M_n$ is the twisted constant group over $S_n$ dual
to $T_n$ (SGA~3 X~5.7), and $B'_n$ is the abelian scheme dual to
$B_n$.  Restriction to $S_o$, or to any $S_n$, induces an equivalence
from the category of isotrivial tori over $S$ to the corresponding
category over $S_o$ (SGA~3 X~3.2).  Hence the $T_n$ come from a torus
$T$ over $S$, determined up to a unique isomorphism.  Denote its dual
group by $M$; its restriction to $S_n$ is canonically
isomorphic to $\underline M_n$.

}{Conversely, suppose that $\mathcal B$ is algebraizable, that is, that
it comes from an abelian scheme $B$, and let us prove that $\mathcal A$
is algebraizable.  We first prove that
$\mathcal A^o=(A_n^o)$ is algebraizable.  By VIII~\SGARefLink{sga7.sga7i.eviii.n3-7}{3.7}, the extensions
$A_n^o$ of $B_n$ by $T_n$ may be interpreted as homomorphisms
\begin{equation}\tag{*}
 \underline M_n\longrightarrow B'_n,
\end{equation}
where $\underline M_n$ is the twisted constant group over $S_n$ dual
to $T_n$ (SGA~3 X~5.7), and $B'_n$ is the abelian scheme dual to
$B_n$.  Restriction to $S_o$, or to any $S_n$, induces an equivalence
from the category of isotrivial tori over $S$ to the corresponding
category over $S_o$ (SGA~3 X~3.2).  Hence the $T_n$ come from a torus
$T$ over $S$, determined up to a unique isomorphism.  Denote its dual
group by $\underline M$; its restriction to $S_n$ is canonically
isomorphic to $\underline M_n$.

}{E7I-IX-037}The isotriviality of $T$ means that $\underline M$ is a disjoint union
of $S$-schemes $\underline M^i$ which are finite over $S$
(SGA~3 X~5.11).  It follows at once that the homomorphisms $(*)$ come
from a unique homomorphism
\begin{equation}\tag{**}
 \underline M\longrightarrow B',
\end{equation}
where $B'$ is the abelian scheme dual to $B$.  In turn, this
homomorphism is interpreted as an extension of $B$ by $T$, say $A^o$,
which plainly supplies an algebraization of $\mathcal A^o$.

\SGApage{I}{\(\approx\)405}
\SGAvar{\medskip
\noindent
\SGARefTarget{sga7.sga7i.eix.n7-2-1-3}{7.2.1.3}. On the other hand, by hypothesis $A_o/A_o^o$ is a finite
group scheme $\Phi_o$ over $S$, and the $A_n/A_n^o$ are likewise
finite group schemes $\Phi_n$ over $S$; they are necessarily etale
over $S$.  There is therefore an
etale group scheme $\Phi$ over $S$, determined up to a unique
isomorphism, whose formal completion is the formal group scheme
$\mathcal A/\mathcal A^o$ defined by the $\Phi_n$.

}{\medskip
\noindent
\SGARefTarget{sga7.sga7i.eix.n7-2-1-3}{7.2.1.3}. On the other hand, by hypothesis $A_o/A_o^o$ is a finite
group scheme $\Phi_o$ over $S_o$, and the $A_n/A_n^o$ are likewise
finite group schemes $\Phi_n$ over $S_n$; they are necessarily etale
over $S_n$.  There is therefore an
etale group scheme $\Phi$ over $S$, determined up to a unique
isomorphism, whose formal completion is the formal group scheme
$\mathcal A/\mathcal A^o$ defined by the $\Phi_n$.

}{E7I-IX-038}Since $G=A^o$ is divisible, the category of extensions, as fppf
sheaves, of a finite locally constant group over $S$ such as $\Phi$
by the group scheme $A^o$ is canonically equivalent to the category
of torsors under the sheaf
\[
 H=\underline{\operatorname{Hom}}_{\mathrm{gr}}(\Phi,G),
\]
and this equivalence is compatible with every base change.  Applying
the same description to extensions of $\Phi_n$ by $G_n$ in terms of
torsors under
\[
 H_n\simeq
 \underline{\operatorname{Hom}}_{\mathrm{gr}}(\Phi_n,G_n),
\]
the $A_n$ define a coherent system of torsors $P_n$ under the $H_n$.
Because $G$ is an extension of an abelian scheme by a torus, it follows
readily that $H$ is a group scheme \underline{finite and locally free
over $S$}.  Consequently, the functor
\[
 P\longmapsto\widehat P=(P_n)
\]
is an equivalence between the category of torsors over $S$ under $H$
and the category of ``formal'' torsors $\widehat P$ under
$\widehat H$, that is, coherent systems of torsors $P_n$ under the
$H_n$.

\SGApage{I}{\(\approx\)406}
Thus the $P_n$ defined by the extensions $A_n$ come from a single
torsor $P$ over $S$ under $H$.  We obtain an extension $A$ of $\Phi$
by $G=A^o$ which algebraizes $\mathcal A$, except that its
representability remains to be checked.  Regard $A$ as a torsor over
$\Phi$ under $A^o_\Phi$.  This torsor becomes trivial after base
change by the finite, surjective, locally free morphism
\[
 P\times_S\Phi\longrightarrow\Phi,
\]
since the base change $P\longrightarrow S$ splits the extension $A$
of $\Phi$ by $A^o$.  One may therefore conclude by one of Raynaud's
many criteria for effectiveness of descent [23, XI~3.1(6)].

\medskip
\noindent
\SGARefTarget{sga7.sga7i.eix.n7-2-1-4}{7.2.1.4}. This completes the proof of (b), and only the full
faithfulness assertion (a) remains.  We leave it to the reader to
verify that this assertion is essentially contained in the method of
proof of (b), or, if preferred, to find an independent direct proof.

\medskip
\noindent
\SGARefTarget{sga7.sga7i.eix.n7-2-1-5}{7.2.1.5}. Notice that the appeal to Raynaud's delicate descent result
becomes unnecessary if $A^o$ is assumed quasi-projective over $S$, or,
equivalently, if $B$ is projective over $S$.  This condition holds in
particular when one starts from the formal group scheme $\widehat A$
over $S$, with $A$ as in \SGARefLink{sga7.sga7i.eix.n7-1-5}{7.1.5}.  This is the only case in which we
shall use \SGARefLink{sga7.sga7i.eix.n7-2-1}{7.2.1}.



\medskip
\noindent
\SGARefTarget{sga7.sga7i.eix.n7-2-2}{7.2.2}. We apply \SGARefLink{sga7.sga7i.eix.n7-2-1}{7.2.1} to the formal group scheme $\widehat A$
considered in \SGARefLink{sga7.sga7i.eix.n7-1}{7.1}, in the favorable case treated there in which the
formal abelian scheme $\widehat B$ is algebraizable (cf. \SGARefLink{sga7.sga7i.eix.n7-1-5}{7.1.5}), and
we suppose in addition that the torus $T_o$ over $S_o$ is isotrivial.
We thereby obtain a group scheme $A^\natural$ over $S$, called the
\underline{Raynaud group} of $A$ over $S$, characterized up to
isomorphism by a canonical isomorphism of formal group schemes over
$\operatorname{Spf}(V)=\widehat S$:
\begin{equation}\tag{\SGARefTarget{sga7.sga7i.eix.n7-2-3}{7.2.3}}
 (A^\natural)^\wedge\simeq\widehat A.
\end{equation}
It follows from \SGARefLink{sga7.sga7i.eix.n7-2-1}{7.2.1} that
\begin{equation}\tag{\SGARefTarget{sga7.sga7i.eix.n7-2-4}{7.2.4}}
 \Phi=A^\natural/A^{\natural o}
\end{equation}
is a finite etale group over $S$, and that the identity component
\[
 A^{\natural o}=A^{o\natural}
\]
is endowed with the structure of an extension
\begin{equation}\tag{\SGARefTarget{sga7.sga7i.eix.n7-2-5}{7.2.5}}
 0\longrightarrow T\longrightarrow A^{\natural o}
 \longrightarrow B\longrightarrow0,
\end{equation}
where $B$ is the abelian scheme over $S$ whose formal completion is
the formal abelian scheme $\widehat B$ constructed in \SGARefLink{sga7.sga7i.eix.n7-1}{7.1}, and $T$ is
characterized, up to a unique isomorphism, as the isotrivial torus
over $S$ lifting $T_o$.  On restriction to the $S_n$, the extension
structure (\SGARefLink{sga7.sga7i.eix.n7-2-5}{7.2.5}) gives the extension structures on
$A_n^{\natural o}=A_n^o$ defined by (\SGARefLink{sga7.sga7i.eix.n7-1-1}{7.1.1}).

\SGApage{I}{407}
\medskip
\noindent
\SGARefTarget{sga7.sga7i.eix.n7-2-6}{7.2.6}. It also follows immediately from \SGARefLink{sga7.sga7i.eix.n7-2-1}{7.2.1} that the Raynaud group
$A^\natural$, as well as the exact sequence (\SGARefLink{sga7.sga7i.eix.n7-2-5}{7.2.5}), depends
functorially on $A$ (satisfying the conditions under consideration),
and that its formation commutes with every base change
$S'\longrightarrow S$ associated with a homomorphism of adic
topological rings $V\longrightarrow V'$.



\medskip
\noindent
\SGARefTarget{sga7.sga7i.eix.n7-3}{7.3}. Let us now consider, for a fixed prime number $\ell$, the
pro-$\ell$ group $T_{\ell}(A^{o})$. Identifying a finite scheme over
$S$ with the corresponding formal scheme, and defining the
``\underline{fixed part}'' $T_{\ell}(A^{o})^{f}$ as usual, we obtain
an isomorphism as in (\SGARefLink{sga7.sga7i.eix.n5-1-5}{5.1.5}):

\SGApage{I}{\(\approx\)408}
\begin{equation}\tag{\SGARefTarget{sga7.sga7i.eix.n7-3-1}{7.3.1}}
T_{\ell}(A^{o})^{f} \simeq T_{\ell}(\widehat{A}^{o}) \quad ,
\end{equation}
and, defining the ``\underline{toric part}'' as in (\SGARefLink{sga7.sga7i.eix.n5-1-6}{5.1.6}),

\begin{equation}\tag{\SGARefTarget{sga7.sga7i.eix.n7-3-2}{7.3.2}}
T_{\ell}(A^{o})^{t} = T_{\ell}(\widehat{T}) \simeq T_{\ell}(T)
 \subset T_{\ell}(A^{o})^{f} \simeq T_{\ell}(\widehat{A}^{o}) \quad ,
\end{equation}
we immediately obtain, for the ``\underline{abelian part},'' an
isomorphism

\begin{equation}\tag{\SGARefTarget{sga7.sga7i.eix.n7-3-3}{7.3.3}}
T_{\ell}(A^{o})^{ab} \stackrel{\mathrm{def}}{=}
T_{\ell}(A^{o})^{f}/T_{\ell}(A^{o})^{t}
 \simeq T_{\ell}(\widehat{B}) \quad .
\end{equation}
We have not yet used here the hypothesis that $\widehat B$ is
algebraizable, the hypothesis that allowed us to construct the
abelian scheme $B$ and the Raynaud extension of $B$ by $T$. Since
there is, of course, a canonical isomorphism (the special case of
(\SGARefLink{sga7.sga7i.eix.n7-3-1}{7.3.1}) applied to the abelian scheme $B$!)
$T_{\ell}(\widehat B)\simeq T_{\ell}(B)$, the isomorphism (\SGARefLink{sga7.sga7i.eix.n7-3-3}{7.3.3})
also gives us

\begin{equation}\tag{\SGARefTarget{sga7.sga7i.eix.n7-3-4}{7.3.4}}
T_{\ell}(A^{o})^{ab} \simeq T_{\ell}(B) \quad .
\end{equation}
Let us also note in passing that (\SGARefLink{sga7.sga7i.eix.n7-2-3}{7.2.3}) and (\SGARefLink{sga7.sga7i.eix.n7-3-1}{7.3.1}) imply a canonical
isomorphism

\begin{equation}\tag{\SGARefTarget{sga7.sga7i.eix.n7-3-5}{7.3.5}}
T_{\ell}(A^{o})^{f} \simeq T_{\ell}(A^{\natural o}) \quad ,
\end{equation}
since the right-hand side is already a profinite group, that is,
identical to its fixed part. The canonical exact sequence

\begin{equation}\tag{\SGARefTarget{sga7.sga7i.eix.n7-3-6}{7.3.6}}
0 \longrightarrow T_{\ell}(A^{o})^{t}
 \longrightarrow T_{\ell}(A^{o})^{f}
 \longrightarrow T_{\ell}(B) \longrightarrow 0 \quad ,
\end{equation}
which can in any case (without the algebraizability hypothesis on
$\widehat B$) be interpreted as the exact sequence obtained by
applying the functor $T_{\ell}$ to the exact sequence

\begin{equation*}
0 \longrightarrow \widehat{T} \longrightarrow \widehat{A}^{o}
 \longrightarrow \widehat{B} \longrightarrow 0 \quad ,
\end{equation*}
can also be regarded as obtained by applying the functor $T_{\ell}$
to the canonical exact sequence of ordinary group schemes

\SGApage{I}{\(\approx\)409}
\begin{equation*}
0 \longrightarrow T \longrightarrow A^{\natural o}
 \longrightarrow B \longrightarrow 0 \qquad .
\end{equation*}



\medskip
\noindent
\SGARefTarget{sga7.sga7i.eix.n7-4}{7.4}. Suppose now that a second group scheme $A'$ over $S$ is given,
satisfying the same hypotheses as $A$, and suppose that $A$ and $A'$
have connected fibres. Let $W$ be a biextension of $(A,A')$ by
$\mathbb G_{mS}$. We thereby obtain a pairing (VIII \SGARefLink{sga7.sga7i.eviii.n2-2}{2.2})

\begin{equation}\tag{\SGARefTarget{sga7.sga7i.eix.n7-4-1}{7.4.1}}
T_{\ell}(A) \times T_{\ell}(A') \longrightarrow
T_{\ell}(\mathbb G_{mS}) = \mathbb Z_{\ell}(1)
\end{equation}
and, by restriction to the fixed parts, a pairing

\begin{equation}\tag{\SGARefTarget{sga7.sga7i.eix.n7-4-2}{7.4.2}}
T_{\ell}(A)^{f} \times T_{\ell}(A')^{f}
 = T_{\ell}(\widehat A) \times T_{\ell}(\widehat {A'})
 \longrightarrow \mathbb Z_{\ell}(1) \qquad .
\end{equation}
This can also be interpreted as the pairing defined by the
biextension $\widehat W$ of $(\widehat A,\widehat {A'})$ by
$\widehat{\mathbb G}_{mS}$ obtained from $W$ by completion (here, for
simplicity, defining a biextension of formal group schemes as a
coherent family of biextensions $W_n$ of $(A_n,A'_n)$ by
$\mathbb G_{mS_n}$). Now we know (VIII \SGARefLink{sga7.sga7i.eviii.n3-5}{3.5}) that $W_n$ comes from a
biextension $W_n^{ab}$ of $(B_n,B'_n)$, where $B'_n$ is the abelian
part of $A'_n$, so that (\SGARefLink{sga7.sga7i.eix.n7-4-2}{7.4.2}) factors through

\begin{equation}\tag{\SGARefTarget{sga7.sga7i.eix.n7-4-3}{7.4.3}}
T_{\ell}(A)^{ab} \times T_{\ell}(A')^{ab}
 = T_{\ell}(\widehat B) \times T_{\ell}(\widehat {B'})
 \longrightarrow \mathbb Z_{\ell}(1) \qquad .
\end{equation}
This is also the pairing defined by the biextension

\SGAvar{\begin{equation*}
\widehat W^{ab}=(W_n^{ab})
\end{equation*}
of $(\widehat B,\widehat {B'})$ by $\widehat{\mathbb G}_{mS}$. When
$\widehat B$ and $\widehat {B'}$ are algebraizable, so that abelian
schemes $B,B'$ over $S$ are available, EGA IV 5.1.4 shows that
$\widehat W^{ab}$ comes from a biextension, defined up to a unique
isomorphism,

}{\begin{equation*}
\widehat W^{ab}=(W_n^{ab})
\end{equation*}
of $(\widehat B,\widehat {B'})$ by $\widehat{\mathbb G}_{mS}$. When
$\widehat B$ and $\widehat {B'}$ are algebraizable, so that abelian
schemes $B,B'$ over $S$ are available, EGA III 5.1.4 shows that
$\widehat W^{ab}$ comes from a biextension, defined up to a unique
isomorphism,

}{E7I-IX-039}\SGApage{I}{\(\approx\)410}
\begin{equation*}
W^{ab}\quad\text{of}\quad(B,B')\quad\text{by}\quad\mathbb G_{mS},
\end{equation*}
and the corresponding pairing

\begin{equation}\tag{\SGARefTarget{sga7.sga7i.eix.n7-4-4}{7.4.4}}
T_{\ell}(B) \times T_{\ell}(B')
 \longrightarrow T_{\ell}(\mathbb G_{mS})
\end{equation}
is canonically isomorphic to (\SGARefLink{sga7.sga7i.eix.n7-4-3}{7.4.3}).

Writing $D(B')$ for the abelian scheme dual to $B'$, $W^{ab}$ defines
a homomorphism

\begin{equation}\tag{\SGARefTarget{sga7.sga7i.eix.n7-4-5}{7.4.5}}
u:B\longrightarrow D(B') \qquad ,
\end{equation}
whose properties reflect those of the pairings (\SGARefLink{sga7.sga7i.eix.n7-4-4}{7.4.4}), that is,
(\SGARefLink{sga7.sga7i.eix.n7-4-3}{7.4.3}). For example, the pairing (\SGARefLink{sga7.sga7i.eix.n7-4-3}{7.4.3}) is a perfect duality at a
point $t\in S$ if and only if
$u_t:B_t\longrightarrow D(B')_t$ is an isogeny whose kernel has rank
prime to $\ell$. This implies that $u$ itself is an isogeny (since
$S$ is local, hence connected), and $\ker u$ is flat and finite over
$S$, so that $\ker u$ has the same rank at every point. Thus the
condition just considered is independent of the chosen point $t$.

\medskip
\noindent
\SGARefTarget{sga7.sga7i.eix.n7-4-6}{7.4.6}. This completes in particular the proof of \SGARefLink{sga7.sga7i.eix.n5-4}{5.4}. Indeed, under
the conditions of that theorem (where $S$ may be assumed complete),
applying the preceding discussion to $A^o$ and ${A'}^o$ shows that,
for every $\ell$, the pairing (\SGARefLink{sga7.sga7i.eix.n7-4-3}{7.4.3}) is a perfect duality at the
generic point $\eta$, and hence also at the point $s$.

\SGApage{I}{411}
\SGAvar{\medskip
\noindent
\SGARefTarget{sga7.sga7i.eix.n7-4-7}{7.4.7}. The preceding considerations likewise make the orthogonality
theorem \SGARefLink{sga7.sga7i.eix.n5-2}{5.2} clear, without recourse to Lemma \SGARefLink{sga7.sga7i.eix.n5-2-4}{5.2.4}. Indeed, along the
way in \SGARefLink{sga7.sga7i.eix.n7-4}{7.4} we established that the pairing
$(\SGARefLink{sga7.sga7i.eix.n5-2-3-1}{5.2.3.1})=(\SGARefLink{sga7.sga7i.eix.n7-4-2}{7.4.2})$ vanishes on the toric parts on both sides, and
therefore factors through the pairing (\SGARefLink{sga7.sga7i.eix.n7-4-3}{7.4.3}). It remains only to
establish that this pairing is nondegenerate on the generic fibres of
the pro-groups under consideration, knowing that it is so on the
special fibres. This follows from the fact that if (\SGARefLink{sga7.sga7i.eix.n7-4-5}{7.4.5}) induces an
isogeny on the special fibres, then it is an isogeny.

}{\medskip
\noindent
\SGARefTarget{sga7.sga7i.eix.n7-4-7}{7.4.7}. The preceding considerations likewise make the orthogonality
theorem \SGARefLink{sga7.sga7i.eix.n5-2}{5.2} clear, without recourse to Lemma \SGARefLink{sga7.sga7i.eix.n5-2-4}{5.2.5}. Indeed, along the
way in \SGARefLink{sga7.sga7i.eix.n7-4}{7.4} we established that the pairing
$(\SGARefLink{sga7.sga7i.eix.n5-2-3-1}{5.2.4.2})=(\SGARefLink{sga7.sga7i.eix.n7-4-2}{7.4.2})$ vanishes on the toric parts on both sides, and
therefore factors through the pairing (\SGARefLink{sga7.sga7i.eix.n7-4-3}{7.4.3}). It remains only to
establish that this pairing is nondegenerate on the generic fibres of
the pro-groups under consideration, knowing that it is so on the
special fibres. This follows from the fact that if (\SGARefLink{sga7.sga7i.eix.n7-4-5}{7.4.5}) induces an
isogeny on the special fibres, then it is an isogeny.

}{E7I-IX-040}\medskip
\noindent
\SGARefTarget{sga7.sga7i.eix.n7-4-8}{7.4.8}. Return to the general conditions of \SGARefLink{sga7.sga7i.eix.n7-4}{7.4}, and suppose that an
open subset $U\subset S-S_o$ is given such that $A_U$ and $A'_U$ are
\underline{abelian schemes}, dual to one another through the
biextension $W_U$. Restricting (\SGARefLink{sga7.sga7i.eix.n7-3-6}{7.3.6}), and the analogous exact
sequence for $A'$, to $U$, we obtain canonical filtrations on the
Barsotti--Tate pro-$\ell$ groups $T_\ell(A_U)$ and $T_\ell(A'_U)$:

\begin{equation}\tag{\SGARefTarget{sga7.sga7i.eix.n7-4-9}{7.4.9}}
\left\{
\begin{aligned}
T_\ell(A_U)&\supset T_\ell(A_U)^f\supset T_\ell(A_U)^t\supset0,\\[4ex]
T_\ell(A'_U)&\supset T_\ell(A'_U)^f\supset T_\ell(A'_U)^t\supset0
\qquad .
\end{aligned}
\right.
\end{equation}
I claim that \underline{these two filtrations are dual to one another}
with respect to the pairing (\SGARefLink{sga7.sga7i.eix.n7-4-1}{7.4.1}), which here is a perfect duality
(VIII \SGARefLink{sga7.sga7i.eviii.n3-2}{3.2}); that is, relation (\SGARefLink{sga7.sga7i.eix.n5-2-2}{5.2.2}) and its symmetric counterpart
still hold, with $A_K,A'_K$ replaced by $A_U,A'_U$. We saw above that
this orthogonality assertion is also equivalent to saying that the
biextension $W^{ab}$ of $(B,B')$ by $\mathbb G_{mS}$ is a perfect
duality over $U$ (or, equivalently when $U$ is dense in $S$, over
$S$). It is enough to verify this property at every point $\eta$ of
$U$; the usual argument then reduces us to the already known case in
which $S$ is a trait with generic point $\eta$.

\SGApage{I}{\(\approx\)412}
The orthogonality relation also allows us to extend formulas (\SGARefLink{sga7.sga7i.eix.n5-5-7}{5.5.7}),
(\SGARefLink{sga7.sga7i.eix.n5-5-8}{5.5.8}), and (\SGARefLink{sga7.sga7i.eix.n5-5-9}{5.5.9}) to the present case, replacing restriction to the
generic point $\eta$ by restriction to the open subset $U$.
Consequently, \SGARefLink{sga7.sga7i.eix.n5-6}{5.6} likewise extends to the present case. We refrain
from rewriting the preceding formulas and statement here in this
context.



Let us summarize, from the point of view of N\'{e}ron models, the
results obtained concerning the abelian parts and their pairings.

\medskip
\noindent
\underline{Proposition} \SGARefTarget{sga7.sga7i.eix.n7-5}{7.5}. \underline{Let $S$ be a complete trait
with function field $K$. To every abelian scheme $A_K$ over $K$
having semistable reduction over $S$, the procedure of \SGARefLink{sga7.sga7i.eix.n7-1}{7.1}, applied
to the identity component $A^o$ of the N\'{e}ron model $A$ of $A_K$,
associates its abelian part $B=A^{ab}$, which is an abelian scheme
over $S$ depending functorially on $A_K$. For every other abelian
scheme $A'_K$ with semistable reduction, one can associate to every
divisorial correspondence $W_K$ on $(A_K,A'_K)$ a unique divisorial
correspondence $W^{ab}$ on $(A^{ab},{A'}^{ab})$ such that, for every
integer $n\geq0$, the pullback of $W_n^{ab}$ along
$A_n^o\longrightarrow B_n=A_n^{oab}$ and
${A'}_n^o\longrightarrow B'_n={A'}_n^{oab}$ is the biextension $W_n$
of $(A_n^o,{A'}_n^o)$ obtained from $W$ by reduction, where $W$ is
the canonical extension (\SGARefLink{sga7.sga7i.eix.n1-4}{1.4}) of $W_K$ to $(A^o,{A'}^o)$.

The formation of} $W^{ab}$ \underline{is compatible with the formation of
pullbacks of correspondences by morphisms of abelian schemes}
$\overline A_K\longrightarrow A_K$ \underline{and}
$\overline {A'}_K\longrightarrow A'_K$. \underline{If} $W_K$ \underline{is a perfect
duality}, \underline{then so is} $W^{ab}$. \underline{If} $W_K$ \underline{is the divisorial
correspondence on} $(A_K,A_K)$ \underline{associated with a polarization of}
$A_K$, \underline{then} $W^{ab}$ \underline{is associated with a polarization of} $A^{ab}$.

\SGApage{I}{\(\approx\)413}
The compatibility of the formation of the $W^{ab}$ with pullbacks is
immediate from the indicated characterization of the $W^{ab}$. The
additional assertion concerning polarizations was established in
\SGARefLink{sga7.sga7i.eix.n7-1-5}{7.1.5} at the same time as the algebraizability of $\widehat B$.

\medskip
\noindent
\SGARefTarget{sga7.sga7i.eix.n7-6}{7.6}. Continue to suppose that $S$ is a complete trait, with $A$ the
N\'{e}ron model of an abelian scheme $A_K$ over $K$. Set

\begin{equation}\tag{\SGARefTarget{sga7.sga7i.eix.n7-6-1}{7.6.1}}
A_K^{\natural}=(A^{\natural})_K \quad ,
\end{equation}
where $A^{\natural}$ is the Raynaud group scheme (\SGARefLink{sga7.sga7i.eix.n7-2-2}{7.2.2}): it is an
extension of the finite etale group $\Phi$ by
$A^{\natural o}$, while $A^{\natural o}$ is an extension of the
abelian scheme $B$ by the torus $T$. Thus $A_K^{\natural}$ is itself
an extension of $\Phi_K$ by
$(A_K^{\natural})^o=(A^{\natural o})_K=(A^{o\natural})_K$, which is an
extension of the abelian scheme $B_K$ by the torus $T_K$. For every
prime number $\ell$, (\SGARefLink{sga7.sga7i.eix.n7-3-5}{7.3.5}) and (\SGARefLink{sga7.sga7i.eix.n7-3-2}{7.3.2}) give

\begin{equation}\tag{\SGARefTarget{sga7.sga7i.eix.n7-6-2}{7.6.2}}
T_{\ell}(A_K^{\natural}) \simeq T_{\ell}(A_K^{\natural o})
 \simeq T_{\ell}(A_K)^f \; , \qquad
T_{\ell}(T_K) \simeq T_{\ell}(A_K)^t \quad .
\end{equation}

\medskip
\noindent
\SGARefTarget{sga7.sga7i.eix.n7-6-3}{7.6.3}. By \SGARefLink{sga7.sga7i.eix.n3-3}{3.3}, the formation of $A_K^{\natural o}$ commutes with
every base change associated with a morphism of complete traits
$S'\longrightarrow S$.



\SGApage{I}{\(\approx\)414}
\medskip
\noindent
8. \underline{The Raynaud extension} $A_K^{\natural o}$ \underline{in the
non-semistable case}, \underline{and the ind-group} $A_K^{\flat}$

\medskip
\noindent
\SGARefTarget{sga7.sga7i.eix.n8-1}{8.1}. Continue to suppose that the trait $S$ is complete, but no
longer suppose necessarily that $A_K$ has semistable reduction. As
usual, let $K'$ be a finite Galois extension of $K$ such that
$A_{K'}$ has semistable reduction relative to the normalization $S'$
of $S$ in $K'$. Consider the Raynaud extension
$A_{K'}^{\natural}$ associated with $A_{K'}$ (\SGARefLink{sga7.sga7i.eix.n7-6-1}{7.6.1}). It is clear
that the Galois group of $K'/K$ acts by transport of structure on
this group scheme; hence, by Galois descent, it gives a group scheme
over $K$, denoted $A_{K,K'}^{\natural}$, which depends on the choice
of the extension $K'$. Its identity component, however, is
independent of this choice up to canonical isomorphism, by \SGARefLink{sga7.sga7i.eix.n7-6-3}{7.6.3},
and will simply be denoted $A_K^{\natural o}$. Descent shows that
$A_K^{\natural o}$ is an extension

\SGAvar{\begin{equation}\tag{\SGARefTarget{sga7.sga7i.eix.n8-1-1}{8.1.1}}
0 \longrightarrow T_K \longrightarrow A_K^{\natural o}
 \longrightarrow B_K \longrightarrow 0 \quad ,
\end{equation}
where $T_K$ is a torus and $B_K$ an abelian scheme over $K$.
Likewise, by the definition of the essentially fixed and essentially
toric parts of $T_{\ell}(A_K)$ (\SGARefLink{sga7.sga7i.eix.n4-1}{4.1} and \SGARefLink{sga7.sga7i.eix.n6-3}{6.5}), the isomorphisms
(\SGARefLink{sga7.sga7i.eix.n7-6-2}{7.6.2}) give canonical isomorphisms

}{\begin{equation}\tag{\SGARefTarget{sga7.sga7i.eix.n8-1-1}{8.1.1}}
0 \longrightarrow T_K \longrightarrow A_K^{\natural o}
 \longrightarrow B_K \longrightarrow 0 \quad ,
\end{equation}
where $T_K$ is a torus and $B_K$ an abelian scheme over $K$.
Likewise, by the definition of the essentially fixed and essentially
toric parts of $T_{\ell}(A_K)$ (\SGARefLink{sga7.sga7i.eix.n4-1}{4.1} and \SGARefLink{sga7.sga7i.eix.n6-3}{6.3}), the isomorphisms
(\SGARefLink{sga7.sga7i.eix.n7-6-2}{7.6.2}) give canonical isomorphisms

}{E7I-IX-041}\begin{equation}\tag{\SGARefTarget{sga7.sga7i.eix.n8-1-2}{8.1.2}}
T_{\ell}(A_K^{\natural}) \simeq T_{\ell}(A_K)^{ef} \quad ,
\qquad T_{\ell}(T_K) \simeq T_{\ell}(A_K)^{et} \quad ,
\end{equation}
and hence, on passing to the quotient,

\begin{equation}\tag{\SGARefTarget{sga7.sga7i.eix.n8-1-3}{8.1.3}}
T_{\ell}(B_K) \simeq
T_{\ell}(A_K)^{ef}/T_{\ell}(A_K)^{et} \quad .
\end{equation}

\SGApage{I}{\(\approx\)415}
This therefore determines two of the quotients occurring in the
filtration

\SGAvar{\begin{equation}\tag{\SGARefTarget{sga7.sga7i.eix.n8-1-4}{8.1.4}}
T_\ell(A_K) \supset T_\ell(A_K)^{ef}
 \supset T_\ell(A_K)^{et} \supset 0
\end{equation}
Thus the extension of group schemes (\SGARefLink{sga7.sga7i.eix.n8-1-1}{8.1.1}) describes two quotients
in this filtration. For the third quotient
$T_\ell(A_K)/T_\ell(A_K)^{ef}$, (\SGARefLink{sga7.sga7i.eix.n6-3-2}{6.5.2}) gives a canonical
isomorphism

}{\begin{equation}\tag{\SGARefTarget{sga7.sga7i.eix.n8-1-4}{8.1.4}}
T_\ell(A_K) \supset T_\ell(A_K)^{ef}
 \supset T_\ell(A_K)^{et} \supset 0
\end{equation}
Thus the extension of group schemes (\SGARefLink{sga7.sga7i.eix.n8-1-1}{8.1.1}) describes two quotients
in this filtration. For the third quotient
$T_\ell(A_K)/T_\ell(A_K)^{ef}$, (\SGARefLink{sga7.sga7i.eix.n6-3-2}{6.3.2}) gives a canonical
isomorphism

}{E7I-IX-042}\begin{equation}\tag{\SGARefTarget{sga7.sga7i.eix.n8-1-5}{8.1.5}}
T_\ell(A_K)/T_\ell(A_K)^{ef}
 \simeq D(T_\ell(A'_K)^{et}) \quad ,
\end{equation}
where $D=\underline{\operatorname{Hom}}(-,T_\ell(\mathbb{G}_m))$
denotes Cartier duality (in the sense of Barsotti--Tate
pro-groups), and $A'_K$ denotes the abelian scheme dual to $A_K$.
Consider then the exact sequence of type (\SGARefLink{sga7.sga7i.eix.n8-1-1}{8.1.1})

\begin{equation}\tag{\SGARefTarget{sga7.sga7i.eix.n8-1-6}{8.1.6}}
0 \longrightarrow T'_K \longrightarrow {A'_K}^{\natural o}
 \longrightarrow B'_K \longrightarrow 0
\end{equation}
associated with the abelian scheme $A'_K$ over $K$, together with
the corresponding isomorphisms of type (\SGARefLink{sga7.sga7i.eix.n8-1-2}{8.1.2}). Formula (\SGARefLink{sga7.sga7i.eix.n8-1-5}{8.1.5})
then gives

\begin{equation}\tag{\SGARefTarget{sga7.sga7i.eix.n8-1-7}{8.1.7}}
T_\ell(A_K)/T_\ell(A_K)^{ef}
 \simeq D(T_\ell(T'_K))
 \simeq \underline{M}_K\otimes_{\mathbb{Z}}\mathbb{Z}_\ell \quad ,
\end{equation}
where, generalizing the notation introduced in (\SGARefLink{sga7.sga7i.eix.n5-5-2}{5.5.2}), we have set

\begin{equation}\tag{\SGARefTarget{sga7.sga7i.eix.n8-1-8}{8.1.8}}
\underline{M}_K=D(T'_K) \quad , \qquad
\underline{M}'_K=D(T_K) \quad ,
\end{equation}
and where $D$ now denotes Cartier duality
$\underline{\operatorname{Hom}}(-,\mathbb{G}_{mK})$ in the sense of
ordinary group schemes. Thus $\underline{M}_K$ and
$\underline{M}'_K$ are twisted constant groups over $K$ which become
constant over $K'$ and isomorphic to groups of the form
$\mathbb{Z}^{r}$, where $r=\dim T=\dim T'$.



\SGApage{I}{416}
\SGAvar{\medskip
\noindent
\SGARefTarget{sga7.sga7i.eix.n8-2}{8.2}. The hypotheses are those of \SGARefLink{sga7.sga7i.eix.n8-1}{8.1}. It is clear that the extension
(\SGARefLink{sga7.sga7i.eix.n8-1-1}{8.1.1}) depends functorially on the abelian scheme $A_K$, and that
its formation commutes with every morphism of traits
$\widetilde S\longrightarrow S$. Choose a polarization of $A_K$,
hence an isogeny $A_K\longrightarrow A_K'$. It induces a
homomorphism from (\SGARefLink{sga7.sga7i.eix.n8-1-1}{8.1.1}) to (\SGARefLink{sga7.sga7i.eix.n8-1-6}{8.1.6}), which is immediately seen, by
looking at the $T_\ell$, to be an isogeny term by term. Likewise,
using \SGARefLink{sga7.sga7i.eix.n7-5}{7.5} and descent, one sees that every biextension by
$\mathbb{G}_{mK}$ of abelian schemes $A_K$ and $\overline A_K$ over
$K$ defines a biextension of $(B_K,\overline B_K)$ by
$\mathbb{G}_{mK}$, where $B_K$ and $\overline B_K$ denote the
abelian quotients of the Raynaud extensions corresponding to $A_K$
and $\overline A_K$, respectively. When $W$ is a perfect duality, so
is the corresponding biextension of $(B_K,\overline B_K)$. In
particular, with the notation of \SGARefLink{sga7.sga7i.eix.n7-6}{7.6}, the abelian schemes $B_K$ and
$B'_K$ are canonically dual to one another.

}{\medskip
\noindent
\SGARefTarget{sga7.sga7i.eix.n8-2}{8.2}. The hypotheses are those of \SGARefLink{sga7.sga7i.eix.n8-1}{8.1}. It is clear that the extension
(\SGARefLink{sga7.sga7i.eix.n8-1-1}{8.1.1}) depends functorially on the abelian scheme $A_K$, and that
its formation commutes with every morphism of traits
$\widetilde S\longrightarrow S$. Choose a polarization of $A_K$,
hence an isogeny $A_K\longrightarrow A'_K$. It induces a
homomorphism from (\SGARefLink{sga7.sga7i.eix.n8-1-1}{8.1.1}) to (\SGARefLink{sga7.sga7i.eix.n8-1-6}{8.1.6}), which is immediately seen, by
looking at the $T_\ell$, to be an isogeny term by term. Likewise,
using \SGARefLink{sga7.sga7i.eix.n7-5}{7.5} and descent, one sees that every biextension by
$\mathbb{G}_{mK}$ of abelian schemes $A_K$ and $\overline A_K$ over
$K$ defines a biextension of $(B_K,\overline B_K)$ by
$\mathbb{G}_{mK}$, where $B_K$ and $\overline B_K$ denote the
abelian quotients of the Raynaud extensions corresponding to $A_K$
and $\overline A_K$, respectively. When $W$ is a perfect duality, so
is the corresponding biextension of $(B_K,\overline B_K)$. In
particular, with the notation of \SGARefLink{sga7.sga7i.eix.n7-6}{7.6}, the abelian schemes $B_K$ and
$B'_K$ are canonically dual to one another.

}{E7I-IX-043}The correspondence between biextensions on $(A_K,\overline A_K)$
and biextensions on $(B_K,\overline B_K)$ has evident functoriality
properties as $A_K$ and $\overline A_K$ vary. These follow by descent
from the analogous properties stated in \SGARefLink{sga7.sga7i.eix.n7-5}{7.5}, and we leave their
formulation to the reader. We also leave to the reader the
compatibility of the preceding construction with every change of
traits $\widetilde S\longrightarrow S$.

\SGAvar{\medskip
\noindent
\SGARefTarget{sga7.sga7i.eix.n8-3}{8.3}. Return to the group $A_{K,K'}^{\natural}$ and study its
variation with $K'$. It is clear that two isomorphic extensions
$K'$ and $K'_1$ of $K'$ give canonically isomorphic groups, so it is
enough to consider the subextensions $K'$ of the separable closure
$\overline K$ of $K$. If $K'\hookrightarrow K''$ are two such finite
Galois subextensions of $\overline K$, it follows immediately from
\SGARefLink{sga7.sga7i.eix.n3-3}{3.3} that there is a canonical \underline{open immersion}

\begin{equation}\tag{\SGARefTarget{sga7.sga7i.eix.n8-3-1}{8.3.1}}
A^{\natural}_{K,K'} \hookrightarrow A^{\natural}_{K,K''}
 \qquad ,
\end{equation}
which identifies the first group with an open subgroup of the
second. As $K'$ varies, the $A_{K,K'}$ thus form an
inductive system of smooth groups of finite type over $K$, whose
transition morphisms are open immersions. If we set

}{\medskip
\noindent
\SGARefTarget{sga7.sga7i.eix.n8-3}{8.3}. Return to the group $A_{K,K'}^{\natural}$ and study its
variation with $K'$. It is clear that two isomorphic extensions
$K'$ and $K'_1$ of $K$ give canonically isomorphic groups, so it is
enough to consider the subextensions $K'$ of the separable closure
$\overline K$ of $K$. If $K'\hookrightarrow K''$ are two such finite
Galois subextensions of $\overline K$, it follows immediately from
\SGARefLink{sga7.sga7i.eix.n3-3}{3.3} that there is a canonical \underline{open immersion}

\begin{equation}\tag{\SGARefTarget{sga7.sga7i.eix.n8-3-1}{8.3.1}}
A^{\natural}_{K,K'} \hookrightarrow A^{\natural}_{K,K''}
 \qquad ,
\end{equation}
which identifies the first group with an open subgroup of the
second. As $K'$ varies, the $A^{\natural}_{K,K'}$ thus form an
inductive system of smooth groups of finite type over $K$, whose
transition morphisms are open immersions. If we set

}{E7I-IX-044}\SGApage{I}{\(\approx\)417}
\begin{equation}\tag{\SGARefTarget{sga7.sga7i.eix.n8-3-2}{8.3.2}}
\Phi_{K,K'}=A^{\natural}_{K,K'}
 /{A^{\natural}_{K,K'}}^{o} \qquad ,
\end{equation}
the $\Phi_{K,K'}$ themselves form an inductive system of finite
etale groups over $K$, which will be studied in greater detail
below.

We shall set

\begin{equation}\tag{\SGARefTarget{sga7.sga7i.eix.n8-3-3}{8.3.3}}
A^{\flat}_{K}=\varinjlim_{K'} A^{\natural}_{K,K'} \qquad ,
\end{equation}
where the limit is taken in the category of fppf sheaves over $K$.
Since the transition morphisms in this inductive system are open
immersions, it is clear moreover that this limit is represented by
a group scheme over $K$ (in general non-quasi-compact), which we
shall naturally denote by the same symbol $A^{\flat}_{K}$. For this
group scheme one has

\begin{equation}\tag{\SGARefTarget{sga7.sga7i.eix.n8-3-4}{8.3.4}}
{A_{K}^{\flat}}^{o}=A_{K}^{\natural o} \quad , \qquad
A^{\flat}_{K}/{A_{K}^{\flat}}^{o}
 \simeq \varinjlim_{K'}\Phi_{K,K'} \qquad .
\end{equation}
The latter group will be determined in \SGARefLink{sga7.sga7i.eix.n11-10}{11.10}: it is canonically
isomorphic to $\underline{M}_{K}\otimes\mathbb{Q}/\mathbb{Z}$, with
the notation of (\SGARefLink{sga7.sga7i.eix.n8-1-8}{8.1.8}).



\SGApage{I}{\(\approx\)418}
\medskip
\noindent
9. \underline{Definition of the monodromy pairing}
$\underline{M}_{\ell}\otimes\underline{M}'_{\ell}
 \longrightarrow \mathbb{Z}_{\ell S}$.

\SGAvar{\medskip
\noindent
\SGARefTarget{sga7.sga7i.eix.n9-1}{9.1}. In this section, suppose that $S$ is henselian and that $A_K$
has semistable reduction over $S$, or, equivalently, that $A'_K$
has semistable reduction over $S$ (\SGARefLink{sga7.sga7i.eix.n3-5-1}{3.5.1}). We then have over $S$ two
twisted constant $\mathbb{Z}$-modules $\underline{M}$ and
${M}'$ (\SGARefLink{sga7.sga7i.eix.n5-5-2}{5.5.2}). Over a suitable finite etale covering of
$S$ (and hence over $S$ itself if $S$ is strictly local, that is, if
$k$ is separably closed), they split into constant groups with
values $M$ and $M'$, respectively, noncanonically isomorphic to
$\mathbb{Z}^{r}$, where $r$ is the common dimension of the maximal
tori of $A_o$ and $A'_o$. For every prime number $\ell$, we shall
consider the corresponding $\ell$-adic sheaves

}{\medskip
\noindent
\SGARefTarget{sga7.sga7i.eix.n9-1}{9.1}. In this section, suppose that $S$ is henselian and that $A_K$
has semistable reduction over $S$, or, equivalently, that $A'_K$
has semistable reduction over $S$ (\SGARefLink{sga7.sga7i.eix.n3-5-1}{3.5.1}). We then have over $S$ two
twisted constant $\mathbb{Z}$-modules $\underline{M}$ and
$\underline{M}'$ (\SGARefLink{sga7.sga7i.eix.n5-5-2}{5.5.2}). Over a suitable finite etale covering of
$S$ (and hence over $S$ itself if $S$ is strictly local, that is, if
$k$ is separably closed), they split into constant groups with
values $M$ and $M'$, respectively, noncanonically isomorphic to
$\mathbb{Z}^{r}$, where $r$ is the common dimension of the maximal
tori of $A_o$ and $A'_o$. For every prime number $\ell$, we shall
consider the corresponding $\ell$-adic sheaves

}{E7I-IX-045}\begin{equation}\tag{\SGARefTarget{sga7.sga7i.eix.n9-1-1}{9.1.1}}
\underline{M}_{\ell}=\underline{M}\otimes_{\mathbb{Z}}\mathbb{Z}_{\ell}
 \quad , \qquad
\underline{M}'_{\ell}=\underline{M}'\otimes_{\mathbb{Z}}\mathbb{Z}_{\ell}
 \qquad ,
\end{equation}
whose geometric meaning is made explicit by formulas (\SGARefLink{sga7.sga7i.eix.n5-5-7}{5.5.7}) and
(\SGARefLink{sga7.sga7i.eix.n5-5-8}{5.5.8}). In this section we shall define a canonical pairing, called
the \underline{monodromy pairing},

\begin{equation}\tag{\SGARefTarget{sga7.sga7i.eix.n9-1-2}{9.1.2}}
u_{\ell}:\underline{M}_{\ell}\otimes\underline{M}'_{\ell}
 \longrightarrow\mathbb{Z}_{\ell} \qquad ,
\end{equation}
which will be defined in terms of the Barsotti--Tate pro-$\ell$-groups
$T_{\ell}(A^{o})$ and $T_{\ell}({A'}^{o})$, and even in terms of the
generic fiber $T_{\ell}(A_K)$ of the first of these groups (subject,
when $\ell=p$, to the availability of Tate's theorem [33, Th.~4],
that is, for the time being, provided we are then in unequal
characteristic). When $\ell\neq p$, this is equivalent to saying
that the pairing (\SGARefLink{sga7.sga7i.eix.n9-1-2}{9.1.2}) is described in terms of the Galois
$\pi$-module

\begin{equation}\tag{\SGARefTarget{sga7.sga7i.eix.n9-1-3}{9.1.3}}
U=T_{\ell}(A_K)(\overline K)=T_{\ell}(A_{\overline K})
\end{equation}
considered in (\SGARefLink{sga7.sga7i.eix.n2-5-3}{2.5.3}). More precisely, we shall see that knowing the
pairing is essentially equivalent to knowing the action of the
monodromy group $I$ on $U$.



\SGApage{I}{419}
\SGAvar{\medskip
\noindent
\SGARefTarget{sga7.sga7i.eix.n9-2}{9.2}. \underline{The case $\ell\neq p$}. Begin with the simpler case
$\ell\neq p$, using the notation introduced in (\SGARefLink{sga7.sga7i.eix.n2-5-3}{2.5.3}). Note that,
up to torsion, the relation $V=U^I$ also means that $V^{\perp}$ is
the module generated by the elements $gx-x$, where $u\in U$ and
$g\in I$, and in any case one has

}{\medskip
\noindent
\SGARefTarget{sga7.sga7i.eix.n9-2}{9.2}. \underline{The case $\ell\neq p$}. Begin with the simpler case
$\ell\neq p$, using the notation introduced in (\SGARefLink{sga7.sga7i.eix.n2-5-3}{2.5.3}). Note that,
up to torsion, the relation $V=U^I$ also means that $V^{\perp}$ is
the module generated by the elements $gx-x$, where $x\in U$ and
$g\in I$, and in any case one has

}{E7I-IX-046}\begin{equation}\tag{\SGARefTarget{sga7.sga7i.eix.n9-2-1}{9.2.1}}
gx-x\in V^{\perp}\quad\text{for }x\in U,\ g\in I \qquad .
\end{equation}
Thus, for $g\in I$,

\SGAvar{\begin{equation}\tag{\SGARefTarget{sga7.sga7i.eix.n9-2-2}{9.2.2}}
gx=x+u(g)\mathbin{.}x \qquad ,
\end{equation}
where $g \longmapsto u(g)$ is a homomorphism from $U$ to
$V^{\perp}$ which vanishes on $V=U^I$, or equivalently

}{\begin{equation}\tag{\SGARefTarget{sga7.sga7i.eix.n9-2-2}{9.2.2}}
gx=x+u(g)\mathbin{.}x \qquad ,
\end{equation}
where $u(g)$ is a homomorphism from $U$ to
$V^{\perp}$ which vanishes on $V=U^I$, or equivalently

}{E7I-IX-047}\begin{equation}\tag{\SGARefTarget{sga7.sga7i.eix.n9-2-3}{9.2.3}}
u\in\operatorname{Hom}(U/V,V^{\perp})
 \simeq (U/V)^{\vee}\otimes V^{\perp} \qquad .
\end{equation}
In the case of semistable reduction under consideration, that is,
when $V^{\perp}\subset V$, formula (\SGARefLink{sga7.sga7i.eix.n9-2-2}{9.2.2}) says that the action of
$g$ on $U$ is a \underline{transvection} parallel to
$V^{\perp}=W$, and composition of these transvections is simply
addition of the corresponding homomorphisms $u$ in (\SGARefLink{sga7.sga7i.eix.n9-2-3}{9.2.3}).
Consequently, the action of the inertia group $I$ on $U$ is
completely described by a group homomorphism

\begin{equation}\tag{\SGARefTarget{sga7.sga7i.eix.n9-2-4}{9.2.4}}
I\longrightarrow (U/V)^{\vee}\otimes V^{\perp} \qquad .
\end{equation}

\SGApage{I}{\(\approx\)420}
The orthogonality theorem, applied both to $A_K$ and to $A'_K$,
gives

\begin{equation*}
V^{\perp}=W \quad , \qquad
(U/V)^{\vee}\simeq W'(-1) \qquad ,
\end{equation*}
where $W'$ is the toric part of
$U'=T_{\ell}(A'(\overline K))$, and $(-1)$ denotes the Tate twist.
Moreover, the homomorphism (\SGARefLink{sga7.sga7i.eix.n9-2-4}{9.2.4}) factors through
$I(\ell)\simeq\mathbb{Z}_{\ell}(1)$, so that (\SGARefLink{sga7.sga7i.eix.n9-2-4}{9.2.4}) may be
interpreted as a canonical homomorphism

\begin{equation}\tag{\SGARefTarget{sga7.sga7i.eix.n9-2-5}{9.2.5}}
u:\mathbb{Z}_{\ell}(1)\longrightarrow W'\otimes W(-1)
 \qquad ,
\end{equation}
or, equivalently, as an element

\begin{equation}\tag{\SGARefTarget{sga7.sga7i.eix.n9-2-6}{9.2.6}}
u\in W'\otimes W(-2) \qquad .
\end{equation}
Recall that $W$ and $W'$ are the groups of $\overline K$-valued
points of the two twisted constant $\ell$-adic sheaves
$T_{\ell}(A^o)^t$ and $T_{\ell}({A'}^o)^t$, whose special fibers are
$T_{\ell}(T_o)$ and $T_{\ell}(T'_o)$, respectively. We obtain the
canonical isomorphisms

\SGAvar{\begin{equation}\tag{\SGARefTarget{sga7.sga7i.eix.n9-2-7}{9.2.7}}
\left\{
\begin{array}{l}
W\simeq T_{\ell}(T_o)(\overline k))
 \simeq\check M\otimes\mathbb{Z}_{\ell}(1),\\[3ex]
W'\simeq T_{\ell}(T_o(\overline k))
 \simeq M'\otimes\mathbb{Z}_{\ell}(1),
\end{array}
\right.
\end{equation}
where $M$ and $M'$ denote the character groups over $\overline k$
of $T_o$ and $T'_o$, respectively:

}{\begin{equation}\tag{\SGARefTarget{sga7.sga7i.eix.n9-2-7}{9.2.7}}
\left\{
\begin{array}{l}
W\simeq T_{\ell}(T_o)(\overline k)
 \simeq\check M\otimes\mathbb{Z}_{\ell}(1),\\[3ex]
W'\simeq T_{\ell}(T'_o(\overline k))
 \simeq\check M'\otimes\mathbb{Z}_{\ell}(1),
\end{array}
\right.
\end{equation}
where $M$ and $M'$ denote the character groups over $\overline k$
of $T_o$ and $T'_o$, respectively:

}{E7I-IX-048}\begin{equation}\tag{\SGARefTarget{sga7.sga7i.eix.n9-2-8}{9.2.8}}
T_o=D(M) \quad , \qquad T'_o=D(M') \qquad .
\end{equation}
It follows that the element $u$ of (\SGARefLink{sga7.sga7i.eix.n9-2-6}{9.2.6}), which describes the
action of the inertia group $I$ on the Tate module $U$ of $A_K$, is
interpreted as a canonical element of
$(\check M'\otimes_{\mathbb{Z}}\check M)
 \otimes_{\mathbb{Z}}\mathbb{Z}_{\ell}$. It is moreover invariant
under $\pi_o=\pi/I=\operatorname{Gal}(\overline k/k)$, since by
definition it is invariant under transport of structure:

\SGApage{I}{\(\approx\)421}
\SGAvar{\begin{equation}\tag{\SGARefTarget{sga7.sga7i.eix.n9-2-9}{9.2.9}}
u\in
(\check M'\otimes_{\mathbb{Z}}\check M)
 \otimes_{\mathbb{Z}}\mathbb{Z}_{\ell}\bigr)^{I}
\simeq
\operatorname{Hom}_{\mathbb{Z}}
 (\underline M\otimes\underline M',\mathbb{Z}_{\ell})
\simeq
\operatorname{Hom}_{\mathbb{Z}_{\ell}}
 (\underline M_{\ell}\otimes\underline M'_{\ell},
  \mathbb{Z}_{\ell}) \quad .
\end{equation}
Up to writing $u$ in place of $u_{\ell}$, this is the element
(\SGARefLink{sga7.sga7i.eix.n9-1-2}{9.1.2}) that we wished to define.

}{\begin{equation}\tag{\SGARefTarget{sga7.sga7i.eix.n9-2-9}{9.2.9}}
u\in
\bigl((\check M'\otimes_{\mathbb{Z}}\check M)
 \otimes_{\mathbb{Z}}\mathbb{Z}_{\ell}\bigr)^{\pi_o}
\simeq
\operatorname{Hom}_{\mathbb{Z}}
 (\underline M\otimes\underline M',\mathbb{Z}_{\ell})
\simeq
\operatorname{Hom}_{\mathbb{Z}_{\ell}}
 (\underline M_{\ell}\otimes\underline M'_{\ell},
  \mathbb{Z}_{\ell}) \quad .
\end{equation}
Up to writing $u$ in place of $u_{\ell}$, this is the element
(\SGARefLink{sga7.sga7i.eix.n9-1-2}{9.1.2}) that we wished to define.

}{E7I-IX-049}To define (\SGARefLink{sga7.sga7i.eix.n9-1-2}{9.1.2}) in the general case, we shall need a few
preliminaries, which will be developed in the next two sections.



\medskip
\noindent
\SGARefTarget{sga7.sga7i.eix.n9-3}{9.3}. \underline{Mixed extensions.}

Let $\underline C$ be an abelian category, and let

\begin{equation}\tag{\SGARefTarget{sga7.sga7i.eix.n9-3-1}{9.3.1}}
P,\ R,\ Q
\end{equation}
be three objects of $\underline C$. We are interested in classifying the
objects $X$ of $\underline C$ endowed with a three-step filtration

\begin{equation}\tag{*}
X^0=X\supset X^1\supset X^2\supset X^3=0
\end{equation}
and with given isomorphisms

\SGAvar{\begin{equation*}
X^0/X^1\xrightarrow{\ \sim\ }P,
\qquad X^1/X^2\xrightarrow{\ \sim\ }R,
\qquad X^2/X^3=X^2\xrightarrow{\ \sim\ }Q.
\end{equation*}
Putting

}{\begin{equation*}
X^0/X^1\xrightarrow{\ \sim\ }Q,
\qquad X^1/X^2\xrightarrow{\ \sim\ }R,
\qquad X^2/X^3=X^2\xrightarrow{\ \sim\ }P.
\end{equation*}
Putting

}{E7I-IX-050}\begin{equation*}
E=X^0/X^2,
\qquad F=X^1/X^3=X^1,
\end{equation*}
we see that $X$ gives rise to two ordinary extensions

\begin{equation}\tag{\SGARefTarget{sga7.sga7i.eix.n9-3-2}{9.3.2}}
0\longrightarrow P\longrightarrow F\longrightarrow R\longrightarrow0,
\qquad
0\longrightarrow R\longrightarrow E\longrightarrow Q\longrightarrow0,
\end{equation}
of $R$ by $P$ and of $Q$ by $R$, respectively. Our point of view here will
be to start, in addition to the objects (\SGARefLink{sga7.sga7i.eix.n9-3-1}{9.3.1}), with the extension
structures (\SGARefLink{sga7.sga7i.eix.n9-3-2}{9.3.2}), which are likewise assumed to be given, and to consider
the corresponding three-step filtered structures $X$.

\SGApage{I}{\(\approx\)422}
More precisely, a \underline{mixed extension} of the extension $E$ by the
extension $F$ will mean an object $X$ of $\underline C$, endowed with a
three-step filtration (*) and with given isomorphisms

\begin{equation}\tag{\SGARefTarget{sga7.sga7i.eix.n9-3-3}{9.3.3}}
E\simeq X^0/X^2,
\qquad F\simeq X^1/X^3=X^1,
\end{equation}
compatible with the filtrations: they carry $X^1/X^2$ to
$\operatorname{Im}(R\longrightarrow E)$ and $X^2$ to
$\operatorname{Im}(P\longrightarrow F)$, respectively, and the
isomorphism $R\simeq F/P$ (and hence $R\simeq X^1/X^2$) deduced from
(\SGARefLink{sga7.sga7i.eix.n9-3-3}{9.3.3}) is required to equal the one deduced from the first exact sequence
in (\SGARefLink{sga7.sga7i.eix.n9-3-2}{9.3.2}).

For fixed $P,Q,R$ and fixed extensions $E$ and $F$, the mixed extensions of
$E$ by $F$ form a category in the evident way; it is in fact a groupoid
(every arrow is invertible). The automorphism group of a mixed extension
$X$ of $E$ by $F$ is canonically isomorphic to $\operatorname{Hom}(Q,P)$
through

\begin{equation}\tag{\SGARefTarget{sga7.sga7i.eix.n9-3-4}{9.3.4}}
\operatorname{Hom}(Q,P)
 \xrightarrow{\ \sim\ }
\operatorname{Aut}_{\mathrm{extpan}}(X),
\qquad
u\longmapsto\operatorname{id}_X+\bar u,
\end{equation}
where $\bar u$ is the endomorphism of the object of $\underline C$
underlying $X$ obtained by composing
$X\longrightarrow Q\xrightarrow{u}P\longrightarrow X$; the outer arrows
are respectively the canonical epimorphism and the canonical monomorphism
coming from the given mixed-extension structure on $X$.

One could make the structure of the category
$\operatorname{EXTPAN}(E,F)$ of mixed extensions of $E$ by $F$ precise by
regarding it as a ``torsor'' under the ``group-category'' (VII \SGARefLink{sga7.sga7i.evii.n2-5}{2.5})
$\operatorname{EXT}(Q,P)$ of extensions of $Q$ by $P$, using a natural
functor analogous to the Brauer composition of ordinary extensions:

\SGApage{I}{\(\approx\)423}
\begin{equation}\tag{\SGARefTarget{sga7.sga7i.eix.n9-3-5}{9.3.5}}
\begin{aligned}
\omega:
\operatorname{EXT}(Q,P)\times\operatorname{EXTPAN}(E,F)
 &\longrightarrow \operatorname{EXTPAN}(E,F),\\
\omega(G,X)&=G\wedge X.
\end{aligned}
\end{equation}
For every object $X$ of $\operatorname{EXTPAN}(E,F)$, the partial functor
$G\longmapsto G\wedge X$ is an equivalence from
$\operatorname{EXT}(Q,P)$ to $\operatorname{EXTPAN}(E,F)$.

We shall give only a brief description of $\omega$. Regard a mixed
extension $X$ of $E$ by $F$ as defining an ordinary extension, also denoted
$X$, of $Q$ by $F$ (viewed as an object of $\underline C$). For every
extension $G$ of $Q$ by $P$, let $\bar G$ denote the extension of $Q$ by
$F$ induced by the inclusion $P\longrightarrow F$. Set

\begin{equation}\tag{\SGARefTarget{sga7.sga7i.eix.n9-3-6}{9.3.6}}
G\wedge X=\omega(G,X)
 \overset{\mathrm{dfn}}{=}\bar G\wedge X,
\end{equation}
where $\wedge$ denotes Brauer composition in the category of extensions of
$Q$ by $F$. The three-step filtration of the right-hand side $X'$ of
(\SGARefLink{sga7.sga7i.eix.n9-3-6}{9.3.6}) is clear, as is the second isomorphism of type (\SGARefLink{sga7.sga7i.eix.n9-3-3}{9.3.3}) attached to
$X'$.

To define the first isomorphism of type (\SGARefLink{sga7.sga7i.eix.n9-3-3}{9.3.3}) for $X'$, observe that the
functor

\SGAvar{\begin{equation*}
Y\longmapsto\widetilde Y,
\qquad
\operatorname{EXT}(Q,F)\longrightarrow\operatorname{EXT}(Q,R),
\end{equation*}
induced by the given morphism $F\longrightarrow R$ in (\SGARefLink{sga7.sga7i.eix.n9-3-2}{9.3.2}) is
compatible with Brauer composition. It therefore sends $X'$ to an
extension $\widetilde X'$ canonically isomorphic to
$\widetilde{\bar G}\wedge\widetilde X$. The first isomorphism (\SGARefLink{sga7.sga7i.eix.n9-3-3}{9.3.3})
gives an isomorphism $\widetilde X\simeq F$, while
$\widetilde{\bar G}$---the extension induced from $G$ by the composite
$P\longrightarrow F\longrightarrow R$---is canonically isomorphic to the
trivial extension, since that composite $P\longrightarrow R$ is zero.
Consequently there is a canonical isomorphism
$\widetilde X'\simeq F$, which specifies on the right-hand side of (\SGARefLink{sga7.sga7i.eix.n9-3-6}{9.3.6})
its mixed-extension structure of $E$ by $F$.


}{\begin{equation*}
Y\longmapsto\widetilde Y,
\qquad
\operatorname{EXT}(Q,F)\longrightarrow\operatorname{EXT}(Q,R),
\end{equation*}
induced by the given morphism $F\longrightarrow R$ in (\SGARefLink{sga7.sga7i.eix.n9-3-2}{9.3.2}) is
compatible with Brauer composition. It therefore sends $X'$ to an
extension $\widetilde X'$ canonically isomorphic to
$\widetilde{\bar G}\wedge\widetilde X$. The first isomorphism (\SGARefLink{sga7.sga7i.eix.n9-3-3}{9.3.3})
gives an isomorphism $\widetilde X\simeq E$, while
$\widetilde{\bar G}$---the extension induced from $G$ by the composite
$P\longrightarrow F\longrightarrow R$---is canonically isomorphic to the
trivial extension, since that composite $P\longrightarrow R$ is zero.
Consequently there is a canonical isomorphism
$\widetilde X'\simeq E$, which specifies on the right-hand side of (\SGARefLink{sga7.sga7i.eix.n9-3-6}{9.3.6})
its mixed-extension structure of $E$ by $F$.


}{E7I-IX-051}\SGApage{I}{\(\approx\)424}
We leave to the reader the definition of $\omega$ on the arrows of
$\operatorname{EXT}(Q,P)\times\operatorname{EXTPAN}(E,F)$, the
verification that (\SGARefLink{sga7.sga7i.eix.n9-3-5}{9.3.5}) has the asserted property---namely, for every
\underline{fixed} mixed extension $X$, the functor
$G\longmapsto G\wedge X$ is an equivalence from
$\operatorname{EXT}(Q,P)$ to $\operatorname{EXTPAN}(E,F)$---and the proof
of the associativity formula

\begin{equation}\tag{\SGARefTarget{sga7.sga7i.eix.n9-3-7}{9.3.7}}
G_1\wedge(G_2\wedge X)
 \simeq (G_1\wedge G_2)\wedge X,
\end{equation}
where $G_1\wedge G_2$ denotes Brauer composition in
$\operatorname{EXT}(Q,P)$.

The preceding discussion immediately gives parts (a) and (b) of the
following result.

\medskip
\noindent
\underline{Proposition \SGARefTarget{sga7.sga7i.eix.n9-3-8}{9.3.8}. Given two extensions} $E$ \underline{and} $F$ \underline{as in
(\SGARefLink{sga7.sga7i.eix.n9-3-2}{9.3.2}) in the abelian category $\mathbf C$, consider the category
$\operatorname{EXTPAN}(E,F)$ of mixed extensions of} $E$ \underline{by} $F$. \underline{Then:}

\smallskip
\noindent
a) \underline{$\operatorname{EXTPAN}(E,F)$ is a groupoid. For every one
of its objects the automorphism group is canonically isomorphic to}
$\operatorname{Ext}^0(Q,P)=\operatorname{Hom}(Q,P)$.

\smallskip
\noindent
b) \underline{The set $\operatorname{Extpan}(E,F)$ of isomorphism classes
of objects of $\operatorname{EXTPAN}(E,F)$ is empty, or is naturally a
torsor under the group} $\operatorname{Ext}^1(Q,P)$.

\smallskip
\noindent
c) \underline{One can canonically define a class}

\begin{equation}\tag{\SGARefTarget{sga7.sga7i.eix.n9-3-9}{9.3.9}}
c(E,F)\in\operatorname{Ext}^2(Q,P),
\end{equation}
\underline{whose vanishing is necessary and sufficient for}
$\operatorname{Extpan}(E,F)\neq\varnothing$, \underline{that is, for the
existence of a mixed extension of} $E$ \underline{by} $F$.

\SGAvar{It remains to prove (d). A mixed
extension of $E$ by $F$ may also be described as an extension $X$ of $E$
by $P$ together with an isomorphism of extensions of $R$ by $P$

}{It remains to prove (c). A mixed
extension of $E$ by $F$ may also be described as an extension $X$ of $E$
by $P$ together with an isomorphism of extensions of $R$ by $P$

}{E7I-IX-052}\SGApage{I}{\(\approx\)425}
\begin{equation*}
X^1\simeq F,
\end{equation*}
where $X^1$ denotes the image of the extension $X$ under the morphism
$R\longrightarrow E$ in (\SGARefLink{sga7.sga7i.eix.n9-3-2}{9.3.2}). Such a mixed extension therefore exists
if and only if there is an element of $\operatorname{Ext}^1(E,P)$ whose
image in $\operatorname{Ext}^1(R,P)$ is the class $\xi$ of $F$. Consider
the exact sequence of the $\operatorname{Ext}^i(-,P)$ associated with the
second exact sequence in (\SGARefLink{sga7.sga7i.eix.n9-3-2}{9.3.2}):

\SGAvar{\begin{equation*}
\cdots\longrightarrow
\operatorname{Ext}^1(E,P)\longrightarrow
\operatorname{Ext}^1(R,P)
\xrightarrow{\partial}
\operatorname{Ext}^2(R,Q)\longrightarrow\cdots .
\end{equation*}
It is enough to define $c(E,F)$ by

}{\begin{equation*}
\cdots\longrightarrow
\operatorname{Ext}^1(E,P)\longrightarrow
\operatorname{Ext}^1(R,P)
\xrightarrow{\partial}
\operatorname{Ext}^2(Q,P)\longrightarrow\cdots .
\end{equation*}
It is enough to define $c(E,F)$ by

}{E7I-IX-053}\begin{equation}\tag{\SGARefTarget{sga7.sga7i.eix.n9-3-10}{9.3.10}}
c(E,F)=\partial\xi,
\qquad
\xi=\operatorname{cl}(F)\in\operatorname{Ext}^1(R,P).
\end{equation}

Now suppose that we are given an exact functor

\begin{equation}\tag{\SGARefTarget{sga7.sga7i.eix.n9-3-11}{9.3.11}}
r:\underline C\longrightarrow\underline C',
\end{equation}
where $\underline C'$ is another abelian category. Write
$P',Q',R',E',F'$ for the objects of $\underline C'$ obtained from
$P,Q,R,E,F$ by applying $r$, so that $E'$ is an extension of $Q'$ by
$R'$ and $F'$ is an extension of $R'$ by $P'$. Up to canonical
isomorphism, the following diagram of functors is commutative:

\SGAvar{\begin{equation}\tag{\SGARefTarget{sga7.sga7i.eix.n9-3-12}{9.3.12}}
\begin{tikzcd}[column sep=large,row sep=large]
\operatorname{EXT}(Q,P)\times\operatorname{EXTPAN}(E,F)
  \arrow[r,"\omega"]\arrow[d]
\& \operatorname{EXTPAN}(E,F)\arrow[d]\\
\operatorname{EXT}(Q',P')\times\operatorname{EXTPAN}(E',F')
  \arrow[r,"\omega"']
\& \operatorname{EXTPAN}(E',F') \quad ,
\end{tikzcd}
\end{equation}
The vertical arrows are the evident functors induced by the exact functor
$r$, and $\omega'$ is the arrow analogous to $\omega$ in (\SGARefLink{sga7.sga7i.eix.n9-3-5}{9.3.5}), with
$\underline C$ replaced by $\underline C'$. The lower arrow in the
printed diagram is labelled $\omega$ without a prime; the explanatory
sentence identifies it as $\omega'$, which is the notation used here.
We obtain analogous compatibilities, with the functor $r$, of the
structures considered in \SGARefLink{sga7.sga7i.eix.n9-3-8}{9.3.8}; their explicit formulation is left to the
reader.

}{\begin{equation}\tag{\SGARefTarget{sga7.sga7i.eix.n9-3-12}{9.3.12}}
\begin{tikzcd}[column sep=large,row sep=large]
\operatorname{EXT}(Q,P)\times\operatorname{EXTPAN}(E,F)
  \arrow[r,"\omega"]\arrow[d]
\& \operatorname{EXTPAN}(E,F)\arrow[d]\\
\operatorname{EXT}(Q',P')\times\operatorname{EXTPAN}(E',F')
  \arrow[r,"\omega'"']
\& \operatorname{EXTPAN}(E',F') \quad ,
\end{tikzcd}
\end{equation}
The vertical arrows are the evident functors induced by the exact functor
$r$, and $\omega'$ is the arrow analogous to $\omega$ in (\SGARefLink{sga7.sga7i.eix.n9-3-5}{9.3.5}), with
$\underline C$ replaced by $\underline C'$. The lower arrow in the
printed diagram is labelled $\omega$ without a prime; the explanatory
sentence identifies it as $\omega'$, which is the notation used here.
We obtain analogous compatibilities, with the functor $r$, of the
structures considered in \SGARefLink{sga7.sga7i.eix.n9-3-8}{9.3.8}; their explicit formulation is left to the
reader.

}{E7I-IX-054}\SGApage{I}{\(\approx\)426}
We shall use the following consequence of this commutativity. Suppose that
$\operatorname{EXTPAN}(E,F)$ is nonempty and choose an object $X$ in it.
Then, by means of $G'\longmapsto G'\wedge X'$, where $X'=r(X)$, the
category $\operatorname{EXTPAN}(E',F')$ is equivalent to
$\operatorname{EXT}(Q',P')$. Thus an object $Y'$ of
$\operatorname{EXTPAN}(E',F')$ is known up to isomorphism when the
isomorphism class of the corresponding object of
$\operatorname{EXT}(Q',P')$ is known, that is, when the corresponding
element

\begin{equation*}
c_X(Y')\in\operatorname{Ext}^1(Q',P')
\end{equation*}
is known.

If $X$ is replaced by another object $X_1=H\wedge X$, where $H$ is an
extension of $Q$ by $P$, then one immediately obtains


\SGAvar{\begin{equation*}
c_{X_1}(Y')=c_X(Y)-\operatorname{cl}(H')
\end{equation*}
from the associativity formula (\SGARefLink{sga7.sga7i.eix.n9-3-7}{9.3.7}). Hence, for fixed $Y'$ and varying
$X$, the elements $c_X(Y')$ run through exactly one coset

}{\begin{equation*}
c_{X_1}(Y')=c_X(Y')-\operatorname{cl}(H')
\end{equation*}
from the associativity formula (\SGARefLink{sga7.sga7i.eix.n9-3-7}{9.3.7}). Hence, for fixed $Y'$ and varying
$X$, the elements $c_X(Y')$ run through exactly one coset

}{E7I-IX-055}\begin{equation}\tag{\SGARefTarget{sga7.sga7i.eix.n9-3-13}{9.3.13}}
c(Y')\in
\operatorname{Ext}^1(Q',P')/
\operatorname{Im}\operatorname{Ext}^1(Q,P),
\end{equation}
which is therefore canonically associated with the mixed extension $Y'$ of
$E'$ by $F'$. By construction, its vanishing is necessary and sufficient
for $Y'$ to be isomorphic to the transform under $r$ of a mixed extension
$X$ of $E$ by $F$.



\SGApage{I}{427}
\medskip
\noindent
\SGARefTarget{sga7.sga7i.eix.n9-4}{9.4}. Let $S$ be a henselian trait with generic point $\eta$, let $n$ be an
integer $>0$, put $\Lambda=\mathbb Z/n\mathbb Z$, and let $\underline C$ be
the category of $\Lambda$-modules on $S$ (for the fppf site on $S$; cf.
VIII \SGARefLink{sga7.sga7i.eviii.n0-2}{0.2}) and $\underline C'$ the category of $\Lambda$-modules on $\eta$.
Let $P,R,Q$ be objects of $\underline C$, and let $E,F$ be extensions of
$Q$ by $R$ and of $R$ by $P$, respectively (\SGARefLink{sga7.sga7i.eix.n9-3-2}{9.3.2}). Suppose that $Q$ is a
locally free $\Lambda_S$-module of finite type and that $P$ is represented
by a finite group of multiplicative type over $S$ (SGA 3 IX), so that it is
of the form

\begin{equation*}
P=D(Q^*),
\end{equation*}
where $D=\underline{\operatorname{Hom}}(-,\mathbb G_{mS})$ denotes Cartier
duality. Let $P_\eta,R_\eta,Q_\eta,E_\eta,F_\eta$ denote the restrictions of
$P,R,Q,E,F$ to $\eta$, and suppose that a mixed extension $X_\eta$ of
$E_\eta$ by $F_\eta$ is given. Under these conditions we shall canonically
define a pairing

\begin{equation}\tag{\SGARefTarget{sga7.sga7i.eix.n9-4-1}{9.4.1}}
c(X_\eta):Q\otimes_\Lambda Q^*\longrightarrow\Lambda_S
\qquad(\Lambda=\mathbb Z/n\mathbb Z).
\end{equation}

For this purpose, first \underline{suppose} that $S$ is \underline{strictly local}. Then $Q$
and $Q^*$ are constant, say $Q=Q_{oS}$ and $Q^*=Q^*_{oS}$, and giving
(\SGARefLink{sga7.sga7i.eix.n9-4-1}{9.4.1}) amounts to giving an element

\begin{equation}\tag{\SGARefTarget{sga7.sga7i.eix.n9-4-2}{9.4.2}}
c(X_\eta)\in\check Q_o\otimes\check Q_o^*,
\end{equation}
where the check denotes the ordinary dual $\operatorname{Hom}(-,\Lambda)$.
We apply the considerations of \SGARefLink{sga7.sga7i.eix.n9-3}{9.3} with the restriction functor
$r:\underline C\longrightarrow\underline C'$. First observe that a mixed
extension of $E$ by $F$ exists. Indeed, the obstruction to its existence
lies in $\operatorname{Ext}^2(Q,P)$ (\SGARefLink{sga7.sga7i.eix.n9-3-8}{9.3.8}(c)), and it is enough to prove

\begin{equation*}
\operatorname{Ext}^2(Q,P)=0.
\end{equation*}

This is a special case of the following lemma.

\SGApage{I}{428}
\medskip
\noindent
\underline{Lemma \SGARefTarget{sga7.sga7i.eix.n9-4-3}{9.4.3}. Under the preceding hypotheses, with $S$ strictly
local, one has:}

\smallskip
\noindent
a) \underline{canonical isomorphisms}

\begin{equation*}
\operatorname{Ext}^i_\Lambda(Q,P)\simeq
H^i(S,\check Q\otimes P),
\qquad
\operatorname{Ext}^i_\Lambda(Q_\eta,P_\eta)\simeq
H^i(\eta,\check Q\otimes P);
\end{equation*}

\noindent
b) \underline{a canonical isomorphism}

\begin{equation}\tag{b)}
\check Q\otimes P\simeq D(Q\otimes Q^*);
\end{equation}

\noindent
c) \underline{for every finite constant group $H=H_{oS}$ over $S$,}

\SGAvar{\begin{equation}\tag{c)}
H^i(S,D(H)=H^i(\eta,D(H)_\eta)=0
\qquad\underline{\text{for }}i\geq2;
\end{equation}

}{\begin{equation}\tag{c)}
H^i(S,D(H))=H^i(\eta,D(H)_\eta)=0
\qquad\underline{\text{for }}i\geq2;
\end{equation}

}{E7I-IX-056}\noindent
d) \underline{canonical isomorphisms}

\SGAvar{\begin{equation}\tag{d)}
\begin{aligned}
H^1(S,D(H))&\simeq\mathbb G_m(S)\otimes\check H^o
              =V^*\otimes\check H_o,\\
H^1(\eta,D(H)_\eta)&\simeq\mathbb G_m(\eta)\otimes H_o
              =K^*\otimes\check H_o,
\end{aligned}
\end{equation}
\underline{where the check in (d) denotes $\operatorname{Hom}(-,\mathbb
Q/\mathbb Z)$.}

}{\begin{equation}\tag{d)}
\begin{aligned}
H^1(S,D(H))&\simeq\mathbb G_m(S)\otimes\check H_o
              =V^*\otimes\check H_o,\\
H^1(\eta,D(H)_\eta)&\simeq\mathbb G_m(\eta)\otimes\check H_o
              =K^*\otimes\check H_o,
\end{aligned}
\end{equation}
\underline{where the check in (d) denotes $\operatorname{Hom}(-,\mathbb
Q/\mathbb Z)$.}

}{E7I-IX-057}The formulas in (a) and (b) follow immediately from the fact that $Q$ is a
locally free $\Lambda_S$-module of finite type. To prove (c), we plainly
reduce to the case $H_o=\mathbb Z/m\mathbb Z$, that is,
$D(H)=\mu_m$. The cohomology sequence associated with the Kummer exact
sequence

\begin{equation*}
0\longrightarrow\mu_m\longrightarrow\mathbb G_m
 \xrightarrow{\ m\cdot\operatorname{id}\ }\mathbb G_m
 \longrightarrow0
\end{equation*}
then reduces us to proving

\begin{equation*}
H^i(S,\mathbb G_m)=H^i(\eta,\mathbb G_m)=0
\qquad\text{for }i\geq1.
\end{equation*}
These groups are the groups $H^i_{\mathrm{\acute et}}$ computed on the
étale site [6, Theorem \SGARefLink{sga7.sga7i.eix.n11-7}{11.7}], and therefore vanish for $S$, since $S$ is
strictly local. Their vanishing for $\eta$ is likewise well known
[6, Theorem \SGARefLink{sga7.sga7i.eix.n1-1}{1.1}].

\SGApage{I}{\(\approx\)429}
It remains to describe the isomorphisms in (d). Over an arbitrary base $U$
there is a canonical homomorphism

\begin{equation}\tag{\SGARefTarget{sga7.sga7i.eix.n9-4-3-1}{9.4.3.1}}
\mathbb G_m(U)\otimes\check H_o
\longrightarrow H^1(U,D(H_{oU})).
\end{equation}
If $H_o$ is annihilated by an integer $m>0$, this homomorphism is described
using the boundary map
$\mathbb G_m(U)\xrightarrow{\partial}H^1(U,\mu_m)$ associated with the
Kummer sequence: it is the composite

\SGAvar{\begin{equation*}
\begin{aligned}
\mathbb G_m(U)\otimes\check H_o
&\longrightarrow H^1(U,\mu_m)\otimes\check H^o\\
&\longrightarrow H^1(U,\mu_m\otimes\check H_o)
 \longrightarrow H^1(U,D(H_{oS})).
\end{aligned}
\end{equation*}
Here the second arrow is deduced from the evident homomorphism
$\check H_o\longrightarrow
\operatorname{Hom}(\mu_m,\mu_m\otimes\check H_o)$, and the third is a
special case of (b), with $m$ in place of $n$. We leave it to the reader to
check that (\SGARefLink{sga7.sga7i.eix.n9-4-3-1}{9.4.3.1}) is independent of the choice of the integer $m$ such
that $mH_o=0$. We also claim that (\SGARefLink{sga7.sga7i.eix.n9-4-3-1}{9.4.3.1}) is always injective and that its
cokernel vanishes if $\operatorname{Pic}(U)=0$. (In fact, the cokernel is
always isomorphic to $\operatorname{Hom}(H_o,\operatorname{Pic}(U))$, but
this is immaterial here.) Indeed, we reduce to the case
$H_o=\mathbb Z/n\mathbb Z$, and the assertion then follows immediately from
the cohomology sequence of the Kummer sequence above. This proves (d), and
immediately yields:

}{\begin{equation*}
\begin{aligned}
\mathbb G_m(U)\otimes\check H_o
&\longrightarrow H^1(U,\mu_m)\otimes\check H_o\\
&\longrightarrow H^1(U,\mu_m\otimes\check H_o)
 \longrightarrow H^1(U,D(H_{oU})).
\end{aligned}
\end{equation*}
Here the second arrow is deduced from the evident homomorphism
$\check H_o\longrightarrow
\operatorname{Hom}(\mu_m,\mu_m\otimes\check H_o)$, and the third is a
special case of (b), with $m$ in place of $n$. We leave it to the reader to
check that (\SGARefLink{sga7.sga7i.eix.n9-4-3-1}{9.4.3.1}) is independent of the choice of the integer $m$ such
that $mH_o=0$. We also claim that (\SGARefLink{sga7.sga7i.eix.n9-4-3-1}{9.4.3.1}) is always injective and that its
cokernel vanishes if $\operatorname{Pic}(U)=0$. (In fact, the cokernel is
always isomorphic to $\operatorname{Hom}(H_o,\operatorname{Pic}(U))$, but
this is immaterial here.) Indeed, we reduce to the case
$H_o=\mathbb Z/n\mathbb Z$, and the assertion then follows immediately from
the cohomology sequence of the Kummer sequence above. This proves (d), and
immediately yields:

}{E7I-IX-058}\medskip
\noindent
\underline{Corollary \SGARefTarget{sga7.sga7i.eix.n9-4-4}{9.4.4}. With the notation of \SGARefLink{sga7.sga7i.eix.n9-4-3}{9.4.3}, there is a canonical
isomorphism}

\SGAvar{\noindent
e)
\begin{equation*}
H^1(\eta,D(H)_\eta)/\operatorname{Im}H^1(S,D(H)
\simeq\check H_o,
\end{equation*}
\underline{and hence, on combining} (a) \underline{and} (b), \underline{a canonical isomorphism}

}{\noindent
e)
\begin{equation*}
H^1(\eta,D(H)_\eta)/\operatorname{Im}H^1(S,D(H))
\simeq\check H_o,
\end{equation*}
\underline{and hence, on combining} (a) \underline{and} (b), \underline{a canonical isomorphism}

}{E7I-IX-059}\begin{equation}\tag{\SGARefTarget{sga7.sga7i.eix.n9-4-5}{9.4.5}}
\operatorname{Ext}^1(Q_\eta,P_\eta)/
\operatorname{Im}\operatorname{Ext}^1(Q,P)
\simeq\check Q_o\otimes\check Q_o^*.
\end{equation}

\SGApage{I}{\(\approx\)430}
Returning to the definition of (\SGARefLink{sga7.sga7i.eix.n9-4-2}{9.4.2}) when $S$ is strictly local, the
relation $\operatorname{Ext}^2(Q,P)=0$ makes it possible to define the
canonical element (\SGARefLink{sga7.sga7i.eix.n9-3-13}{9.3.13}), taking for $Y'$ the given mixed extension
$X_\eta$ of $E_\eta$ by $F_\eta$. Formula (\SGARefLink{sga7.sga7i.eix.n9-4-5}{9.4.5}) identifies this element
$c(X_\eta)$ with an element of $\check Q_o\otimes\check Q_o^*$, thereby
completing the definition of (\SGARefLink{sga7.sga7i.eix.n9-4-2}{9.4.2}). The general considerations of \SGARefLink{sga7.sga7i.eix.n9-3}{9.3}
also show that this element is the obstruction to extending $X_\eta$ to a
mixed extension $X$ of $E$ by $F$.

\medskip
\noindent
\SGARefTarget{sga7.sga7i.eix.n9-4-6}{9.4.6}. The definition of (\SGARefLink{sga7.sga7i.eix.n9-4-1}{9.4.1}) when $S$ is henselian now reduces
immediately to the strictly local case by Galois descent from the strict
henselization of $S$.

\medskip
\noindent
\underline{Proposition \SGARefTarget{sga7.sga7i.eix.n9-4-7}{9.4.7}. Let $P,Q,R,E,F$ be as above, except that the
trait $S$ need not be henselian. Consider the functor}

\begin{equation}\tag{\SGARefTarget{sga7.sga7i.eix.n9-4-7-1}{9.4.7.1}}
X\longmapsto X_\eta
\end{equation}
\underline{from the category of mixed extensions of $E$ by $F$ to the
category of mixed extensions of $E_\eta$ by $F_\eta$. For every mixed
extension $X_\eta$ of $E_\eta$ by $F_\eta$, consider the pairing}

\begin{equation}\tag{\SGARefTarget{sga7.sga7i.eix.n9-4-7-2}{9.4.7.2}}
c_o(X_\eta):Q_o\otimes Q_o^*\longrightarrow\Lambda_s
\end{equation}
\underline{induced by the pairing of type (\SGARefLink{sga7.sga7i.eix.n9-4-1}{9.4.1}) defined over the
henselization $S^h$ of} $S$. \underline{Then the functor (\SGARefLink{sga7.sga7i.eix.n9-4-7-1}{9.4.7.1}) is fully faithful,
and a mixed extension $X_\eta$ of $E_\eta$ by $F_\eta$ belongs to its
essential image if and only if the pairing (\SGARefLink{sga7.sga7i.eix.n9-4-7-2}{9.4.7.2}) is zero.}



\SGApage{I}{431}
\medskip
\noindent
\SGARefTarget{sga7.sga7i.eix.n9-4-7-3}{9.4.7.3}. This statement still makes sense and remains valid if one assumes
only that $S$ is a connected regular noetherian scheme of dimension $1$.
Instead of the single pairing (\SGARefLink{sga7.sga7i.eix.n9-4-7-2}{9.4.7.2}), one must then consider a pairing
$Q_s\otimes Q_s^*\longrightarrow\Lambda_s$ for each closed point $s$ of
$S$. It will be most convenient to prove \SGARefLink{sga7.sga7i.eix.n9-4-7}{9.4.7} in this more general form.

\smallskip
\noindent
a) For every mixed extension $X$, the map

\SGAvar{\begin{equation*}
\operatorname{End}(X)\longrightarrow\operatorname{End}(X_\eta)
\end{equation*}
is bijective. The two sides are respectively the sections over $S$ and
$\eta$ of
$\underline{\operatorname{Hom}}(P,Q)=D(Q\otimes Q^*)$; hence the assertion
follows from the normality of $S$ and the finiteness of
$D(Q\otimes Q^*)$ over $S$. It follows at once that if $X$
and $Y$ are two mixed extensions already known to be isomorphic, then

}{\begin{equation*}
\operatorname{End}(X)\longrightarrow\operatorname{End}(X_\eta)
\end{equation*}
is bijective. The two sides are respectively the sections over $S$ and
$\eta$ of
$\underline{\operatorname{Hom}}(Q,P)=D(Q\otimes Q^*)$; hence the assertion
follows from the normality of $S$ and the finiteness of
$D(Q\otimes Q^*)$ over $S$. It follows at once that if $X$
and $Y$ are two mixed extensions already known to be isomorphic, then

}{E7I-IX-060}\begin{equation*}
\operatorname{Hom}(X,Y)\longrightarrow
\operatorname{Hom}(X_\eta,Y_\eta)
\end{equation*}
is bijective.

\smallskip
\noindent
b) If $X$ and $Y$ are two mixed extensions, every isomorphism
$X_\eta\xrightarrow{\sim}Y_\eta$ comes from an isomorphism
$X\longrightarrow Y$. By (a), the question is local for the étale
topology, so we may reduce to the case in which $S$ is strictly local. In
that case \SGARefLink{sga7.sga7i.eix.n9-4-3}{9.4.3} shows that the map

\SGAvar{\begin{equation*}
\operatorname{Ext}^1_\Lambda(P,Q)
\longrightarrow\operatorname{Ext}^1_\Lambda(P_\eta\ ,Q_\eta)
\end{equation*}
identifies with

}{\begin{equation*}
\operatorname{Ext}^1_\Lambda(Q,P)
\longrightarrow\operatorname{Ext}^1_\Lambda(Q_\eta,P_\eta)
\end{equation*}
identifies with

}{E7I-IX-061}\begin{equation*}
V^*\otimes\check H_o\longrightarrow K^*\otimes\check H_o,
\qquad H_o=Q_o\otimes Q_o^*.
\end{equation*}
The latter is injective, since the exact sequence
$0\longrightarrow V^*\longrightarrow K^*\longrightarrow\mathbb Z
\longrightarrow0$ is split. By the discussion of \SGARefLink{sga7.sga7i.eix.n9-3}{9.3} this implies that if
$X_\eta$ and $Y_\eta$ are isomorphic, then so are $X$ and $Y$; the
conclusion now follows from (a).

\SGApage{I}{\(\approx\)432}
Thus (a) and (b) show that $X\longmapsto X_\eta$ is \underline{fully} faithful.

\SGAvar{\smallskip
\noindent
c) Let $X_\eta$ be a mixed extension over $\eta$. We prove that it extends
to $S$ if and only if the pairings (\SGARefLink{sga7.sga7i.eix.n9-4-7-2}{9.4.6.2}) vanish. Here again, the
uniqueness already proved shows that the question is local for the étale
topology, reducing us once more to the strictly local case. The conclusion
then holds by construction, as observed immediately before \SGARefLink{sga7.sga7i.eix.n9-4-6}{9.4.6}.


}{\smallskip
\noindent
c) Let $X_\eta$ be a mixed extension over $\eta$. We prove that it extends
to $S$ if and only if the pairings (\SGARefLink{sga7.sga7i.eix.n9-4-7-2}{9.4.7.2}) vanish. Here again, the
uniqueness already proved shows that the question is local for the étale
topology, reducing us once more to the strictly local case. The conclusion
then holds by construction, as observed immediately before \SGARefLink{sga7.sga7i.eix.n9-4-6}{9.4.6}.


}{E7I-IX-062}This completes the proof of \SGARefLink{sga7.sga7i.eix.n9-4-7}{9.4.7}.



\medskip
\noindent
\SGARefTarget{sga7.sga7i.eix.n9-5}{9.5}. Fix a prime number $\ell$, pro-$\ell$ Barsotti--Tate groups over $S$

\begin{equation}\tag{\SGARefTarget{sga7.sga7i.eix.n9-5-1}{9.5.1}}
P=(P(n))_{n\geq0},
\qquad R=(R(n))_{n\geq0},
\qquad Q=(Q(n))_{n\geq0},
\end{equation}
and extensions of pro-$\ell$ Barsotti--Tate groups

\begin{equation}\tag{\SGARefTarget{sga7.sga7i.eix.n9-5-2}{9.5.2}}
0\longrightarrow P\longrightarrow F\longrightarrow R\longrightarrow0,
\qquad
0\longrightarrow R\longrightarrow E\longrightarrow Q\longrightarrow0.
\end{equation}
Thus, for every integer $n\geq0$, these define extensions of
$\Lambda_n$-modules, where
$\Lambda_n=\mathbb Z/\ell^{\,n+1}\mathbb Z$:

\begin{equation}\tag{\SGARefTarget{sga7.sga7i.eix.n9-5-3}{9.5.3}}
\begin{gathered}
0\longrightarrow P(n)\longrightarrow F(n)\longrightarrow R(n)
 \longrightarrow0,\\
0\longrightarrow R(n)\longrightarrow E(n)\longrightarrow Q(n)
 \longrightarrow0.
\end{gathered}
\end{equation}
We assume that $Q$ is of étale type and $P$ of multiplicative type.
Finally, suppose that a pro-$\ell$ Barsotti--Tate group $X_\eta$ is given
over $\eta$ which is a \underline{mixed extension} of $E_\eta$ by $F_\eta$. By this
we mean that it is endowed, as a pro-$\ell$ Barsotti--Tate group, with a
three-step filtration whose partial extensions are supplied with
isomorphisms to $E_\eta$ and $F_\eta$, respectively. It follows that, for
every $n$, $X_\eta(n)$ is endowed with the structure of a mixed extension
of $E(n)_\eta$ by $F(n)_\eta$. We may therefore associate with it a pairing
of type (\SGARefLink{sga7.sga7i.eix.n9-4-1}{9.4.1}):

\SGApage{I}{\(\approx\)433}
\begin{equation}\tag{\SGARefTarget{sga7.sga7i.eix.n9-5-4}{9.5.4}}
c(X(n)_\eta):Q(n)\otimes Q^*(n)
\longrightarrow\Lambda_n=\mathbb Z/\ell^{\,n+1}\mathbb Z,
\end{equation}
where the $Q^*(n)$ are the finite étale groups over $S$ defined by

\begin{equation*}
P(n)\simeq D(Q^*(n)),
\qquad\text{i.e.}\qquad Q^*(n)=D(P(n)),
\end{equation*}
and $D$ again denotes Cartier duality. The $Q^*(n)$ form a pro-$\ell$ group
in the evident way, denoted by $Q^*$. We leave it to the reader to verify
that, as $n$ varies, the homomorphisms (\SGARefLink{sga7.sga7i.eix.n9-5-4}{9.5.4}) define a homomorphism of
projective systems, that is, a homomorphism of twisted constant
$\ell$-adic sheaves over $S$:

\begin{equation}\tag{\SGARefTarget{sga7.sga7i.eix.n9-5-5}{9.5.5}}
c(X_\eta):Q\otimes Q^*\longrightarrow\mathbb Z_{\ell S}.
\end{equation}

We mention in passing the following result, which will not be used below.

\medskip
\noindent
\underline{Proposition \SGARefTarget{sga7.sga7i.eix.n9-5-6}{9.5.6}. Let $P,Q,R,E,F$ be as above, except that the
trait $S$ need not be henselian. Consider the functor}

\begin{equation}\tag{\SGARefTarget{sga7.sga7i.eix.n9-5-6-1}{9.5.6.1}}
X\longmapsto X_\eta
\end{equation}
\underline{from the category of Barsotti--Tate mixed extensions of $E$ by
$F$ to the category of mixed extensions of $E_\eta$ by $F_\eta$. For every
Barsotti--Tate group $X_\eta$ which is a mixed extension of $E_\eta$ by
$F_\eta$, consider the corresponding pairing}

\SGApage{I}{\(\approx\)434}
\begin{equation}\tag{\SGARefTarget{sga7.sga7i.eix.n9-5-6-2}{9.5.6.2}}
Q_o\otimes Q_o^*\longrightarrow\mathbb Z_{\ell s}
\end{equation}
\underline{induced by the pairing of type (\SGARefLink{sga7.sga7i.eix.n9-5-5}{9.5.5}) defined over the
henselization of $S$. Then the functor (\SGARefLink{sga7.sga7i.eix.n9-5-6-1}{9.5.6.1}) is fully faithful, and its
essential image consists of those Barsotti--Tate groups $X_\eta$ which are
mixed extensions of $E_\eta$ by $F_\eta$ and for which the corresponding
pairing (\SGARefLink{sga7.sga7i.eix.n9-5-6-2}{9.5.6.2}) is zero}. \underline{Moreover, when $\operatorname{char}K=0$} (\underline{so that
Tate's theorem [33, Theorem 4] applies}), \underline{$X_\eta$ belongs to the essential
image if and only if it extends to a Barsotti--Tate group over $S$.}

\medskip
\noindent
\SGARefTarget{sga7.sga7i.eix.n9-5-6-3}{9.5.6.3}. We leave it to the reader to extend this statement to the case in
which $S$ is assumed only to be a connected regular noetherian scheme of
dimension $1$.

\SGAvar{The first assertion of \SGARefLink{sga7.sga7i.eix.n9-5-6}{9.5.6} follows immediately from \SGARefLink{sga7.sga7i.eix.n9-4-7}{9.4.6}, applied to the
components of the Barsotti--Tate groups under consideration.

It remains to prove that if a Barsotti--Tate group $X_\eta$, which is a
mixed extension of $E_\eta$ by $F_\eta$, extends to a Barsotti--Tate group
$X_\eta$ over $S$, then it also extends as a mixed extension.

}{The first assertion of \SGARefLink{sga7.sga7i.eix.n9-5-6}{9.5.6} follows immediately from \SGARefLink{sga7.sga7i.eix.n9-4-7}{9.4.7}, applied to the
components of the Barsotti--Tate groups under consideration.

It remains to prove that if a Barsotti--Tate group $X_\eta$, which is a
mixed extension of $E_\eta$ by $F_\eta$, extends to a Barsotti--Tate group
$X$ over $S$, then it also extends as a mixed extension.

}{E7I-IX-063}By Tate's theorem [33], the homomorphism
$X_\eta\longrightarrow Q_\eta$ extends uniquely to a homomorphism

\SGAvar{\begin{equation*}
X\longrightarrow Q.
\end{equation*}
For every integer $n$, the corresponding homomorphism
$X(n)\longrightarrow Q(n)$ is surjective on the generic fibres and hence is
surjective. Since $Q(n)$ is étale over $n$, the map
$X(n)\longrightarrow Q(n)$ is flat, and consequently it is an epimorphism.
Thus $X\longrightarrow Q$ is an \underline{epimorphism component by component}.

}{\begin{equation*}
X\longrightarrow Q.
\end{equation*}
For every integer $n$, the corresponding homomorphism
$X(n)\longrightarrow Q(n)$ is surjective on the generic fibres and hence is
surjective. Since $Q(n)$ is étale over $S$, the map
$X(n)\longrightarrow Q(n)$ is flat, and consequently it is an epimorphism.
Thus $X\longrightarrow Q$ is an \underline{epimorphism component by component}.

}{E7I-IX-064}Let $\overline F$ be the kernel of $X\longrightarrow Q$. Then
$\overline F$ is a Barsotti--Tate group and $X$ is an extension of $Q$ by
$\overline F$. The isomorphism
$F_\eta\xrightarrow{\sim}\overline F_\eta$ extends, by Tate's theorem, to
an isomorphism $F\simeq\overline F$. Hence $X$ is an extension of $Q$ by
$F$. Consider now the quotient $\overline E=X/P$. A further application of
Tate's theorem gives an isomorphism $E\simeq\overline E$. It follows that
$X$ is a mixed extension of $E$ by $F$ extending the mixed extension
$X_\eta$, as required.

\SGApage{I}{\(\approx\)435}
\medskip
\noindent
\underline{Remarks \SGARefTarget{sga7.sga7i.eix.n9-5-6-4}{9.5.6.4}.}

\smallskip
\noindent
a) Of course, the last assertion of \SGARefLink{sga7.sga7i.eix.n9-5-6}{9.5.6} should hold without restriction
on $K$, since the same should be true of the invoked theorem of Tate.

\SGAvar{\smallskip
\noindent
b) One can give another proof of the good-reduction criterion \SGARefLink{sga7.sga7i.eix.n5-10}{5.10}. As in
\SGARefLink{sga7.sga7i.eix.n5-10}{5.10}, reduce to the case where $A_K$ has semistable reduction, and then
apply \SGARefLink{sga7.sga7i.eix.n9-5-6}{9.5.6} to $T_\ell(A_K)$ regarded as a mixed extension (cf. \SGARefLink{sga7.sga7i.eix.n9-6}{9.6}). This
implies $u_\ell=0$; using \SGARefLink{sga7.sga7i.eix.n11-6-2}{11.6.2}(a) below then gives
$T_\ell(A_K)^t=0$.

\medskip
\noindent
\SGARefTarget{sga7.sga7i.eix.n9-5-8}{9.5.8}. When $\ell\neq p$, the Barsotti--Tate groups under consideration
over $S$ and over $K$, respectively, may simply be interpreted in terms of
twisted constant $\ell$-adic sheaves over $S$, or equivalently in terms of
representations of
$\pi_o=\operatorname{Gal}(\overline k/k)$ and
$\pi=\operatorname{Gal}(\overline K/K)$, respectively, on free
$\mathbb Z_\ell$-modules of finite type. One then sees that the homomorphism
(\SGARefLink{sga7.sga7i.eix.n9-5-5}{9.5.6}) is the canonical homomorphism expressing the action of the inertia
subgroup $I\subset\pi$ on $X_\eta(\overline K)$, as defined by the method
of \SGARefLink{sga7.sga7i.eix.n9-2}{9.2}.

}{\smallskip
\noindent
b) One can give another proof of the good-reduction criterion \SGARefLink{sga7.sga7i.eix.n5-10}{5.10}. As in
\SGARefLink{sga7.sga7i.eix.n5-10}{5.10}, reduce to the case where $A_K$ has semistable reduction, and then
apply \SGARefLink{sga7.sga7i.eix.n9-5-6}{9.5.6} to $T_\ell(A_K)$ regarded as a mixed extension (cf. \SGARefLink{sga7.sga7i.eix.n9-6}{9.6}). This
implies $u_\ell=0$; using \SGARefLink{sga7.sga7i.eix.n11-6-2}{11.5.2}(a) below then gives
$T_\ell(A_K)^t=0$.

\medskip
\noindent
\SGARefTarget{sga7.sga7i.eix.n9-5-8}{9.5.8}. When $\ell\neq p$, the Barsotti--Tate groups under consideration
over $S$ and over $K$, respectively, may simply be interpreted in terms of
twisted constant $\ell$-adic sheaves over $S$, or equivalently in terms of
representations of
$\pi_o=\operatorname{Gal}(\overline k/k)$ and
$\pi=\operatorname{Gal}(\overline K/K)$, respectively, on free
$\mathbb Z_\ell$-modules of finite type. One then sees that the homomorphism
(\SGARefLink{sga7.sga7i.eix.n9-5-5}{9.5.5}) is the canonical homomorphism expressing the action of the inertia
subgroup $I\subset\pi$ on $X_\eta(\overline K)$, as defined by the method
of \SGARefLink{sga7.sga7i.eix.n9-2}{9.2}.

}{E7I-IX-065}To verify this, one reduces immediately to the situation of \SGARefLink{sga7.sga7i.eix.n9-4}{9.4}, taking
$n$ there to be a power of $\ell$ and $R$ to be a locally constant
$\Lambda$-module (as are $Q$ and $P$); the same is then true of $E,F$, and
$X_\eta$. Verification that the two definitions agree is left to the
reader.



\SGApage{I}{436}
\medskip
\noindent
\SGARefTarget{sga7.sga7i.eix.n9-6}{9.6}. We are now in a position to define the announced canonical pairing
(\SGARefLink{sga7.sga7i.eix.n9-1-2}{9.1.2}), without any restriction on $\ell$, in such a way that, for
$\ell\neq p$, this definition agrees with the one given in \SGARefLink{sga7.sga7i.eix.n9-2}{9.2}. For this
purpose, we apply the considerations of \SGARefLink{sga7.sga7i.eix.n9-5}{9.5} to the following situation:

\begin{equation}\tag{\SGARefTarget{sga7.sga7i.eix.n9-6-1}{9.6.1}}
\left\{
\begin{aligned}
P&=T_\ell(A^o)^t=\check{\underline M}{}'_\ell(1),
& Q&=\underline M_\ell,
& R&=T_\ell(A^o)^f/T_\ell(A^o)^t,\\[1.25ex]
F&=T_\ell(A^o)^f,
& E&=D(F')=D\bigl(T_\ell({A'}^o)^f\bigr).
\end{aligned}
\right.
\end{equation}
In the last formula, $D$ denotes Cartier duality in the sense of
pro-$\ell$ Barsotti--Tate groups; that is, $D(Y)$ is the
\underline{projective} system of the $D(Y(n))$, with the usual transition
morphisms. The structure of $F$ as an extension of $R$ by $P$ is clear.
Likewise,

\begin{equation*}
F'=T_\ell({A'}^o)^f
\end{equation*}
has the structure of an extension of $R'$ by
$P'=T_\ell({A'}^o)^t=\check{\underline M}_\ell(1)$. Applying Cartier
duality, we conclude that $E=D(F')$ is endowed with the structure of an
extension of $D(P')=\underline M_\ell$ by $D(R')$. Now $D(R')$ is
canonically isomorphic to $R$ by \SGARefLink{sga7.sga7i.eix.n5-5}{5.5} (proved in \SGARefLink{sga7.sga7i.eix.n7-4-6}{7.4.6}); this specifies the
structure of $E$ as an extension of $Q$ by $R$.

\SGApage{I}{\(\approx\)437}
As the mixed extension $X_\eta$ of $E_\eta$ by $F_\eta$, we take

\begin{equation}\tag{\SGARefTarget{sga7.sga7i.eix.n9-6-2}{9.6.2}}
X_\eta=T_\ell(A_K)=X^0_\eta
 \supset X^1_\eta=T_\ell(A_K)^f
 \supset X^2_\eta=T_\ell(A_K)^t
 \supset0.
\end{equation}
The isomorphism of extensions of $R_\eta$ by $P_\eta$

\begin{equation*}
X^1_\eta\simeq F_\eta
\end{equation*}
is clear. The isomorphism of extensions of $Q_\eta$ by $R_\eta$

\SGAvar{\begin{equation*}
X_\eta/X^2_\eta\simeq F_\eta=D(F'_\eta)
\end{equation*}
is clear from the orthogonality theorem \SGARefLink{sga7.sga7i.eix.n5-2}{5.2}.


We therefore have the data required to define a pairing (\SGARefLink{sga7.sga7i.eix.n9-5-5}{9.5.6}), which
here, since $Q^*=\underline M'_\ell$, takes the form (\SGARefLink{sga7.sga7i.eix.n9-1-2}{9.1.2}).

Its compatibility with the definition of \SGARefLink{sga7.sga7i.eix.n9-2}{9.2} follows from \SGARefLink{sga7.sga7i.eix.n9-5-8}{9.5.8}. The
considerations of \SGARefLink{sga7.sga7i.eix.n5-8}{5.8} also show, subject to the validity of Tate's theorem
[33] (hence at least when $\operatorname{char}K=0$), that the pairing in
question can be described in terms of the single pro-$\ell$
Barsotti--Tate group $X_\eta=T_\ell(A_K)$ over $S$: its three-step
filtration, and the extensions $F$ and $E$ extending $X^1_\eta$ and
$X_\eta/X^2_\eta$, can be defined canonically in terms of $X_\eta$ alone.


}{\begin{equation*}
X_\eta/X^2_\eta\simeq E_\eta=D(F'_\eta)
\end{equation*}
is clear from the orthogonality theorem \SGARefLink{sga7.sga7i.eix.n5-2}{5.2}.


We therefore have the data required to define a pairing (\SGARefLink{sga7.sga7i.eix.n9-5-5}{9.5.5}), which
here, since $Q^*=\underline M'_\ell$, takes the form (\SGARefLink{sga7.sga7i.eix.n9-1-2}{9.1.2}).

Its compatibility with the definition of \SGARefLink{sga7.sga7i.eix.n9-2}{9.2} follows from \SGARefLink{sga7.sga7i.eix.n9-5-8}{9.5.8}. The
considerations of \SGARefLink{sga7.sga7i.eix.n5-8}{5.8} also show, subject to the validity of Tate's theorem
[33] (hence at least when $\operatorname{char}K=0$), that the pairing in
question can be described in terms of the single pro-$\ell$
Barsotti--Tate group $X_\eta=T_\ell(A_K)$ over $\eta$: its three-step
filtration, and the extensions $F$ and $E$ extending $X^1_\eta$ and
$X_\eta/X^2_\eta$, can be defined canonically in terms of $X_\eta$ alone.


}{E7I-IX-066}\medskip
\noindent
\underline{Remark \SGARefTarget{sga7.sga7i.eix.n9-7}{9.7}.} We no longer necessarily assume that $A_K$ has
semistable reduction over the henselian trait $S$. Let $K'$ be a finite
Galois extension of $K$ such that $A_{K'}$ has semistable reduction over
the normalization $S'$ of $S$ in $K'$ (\SGARefLink{sga7.sga7i.eix.n3-6}{3.6}). With the notation of (\SGARefLink{sga7.sga7i.eix.n8-1-8}{8.1.8}),
Galois descent from the monodromy pairing for $A_{K'}$ then defines a
pairing

\SGApage{I}{\(\approx\)438}
\begin{equation}\tag{*}
v_\ell^{K'}:\underline M_K\otimes\underline M'_K
 \longrightarrow\mathbb Z_{\ell K}.
\end{equation}
A priori this pairing depends on the choice of $K'$, but it follows from
(10.4.3) below that the pairing

\begin{equation}\tag{\SGARefTarget{sga7.sga7i.eix.n9-7-1}{9.7.1}}
u_\ell:\underline M_K\otimes\underline M'_K
 \longrightarrow\mathbb Q_{\ell K}
\end{equation}
given by

\begin{equation*}
u_\ell=\frac{1}{e(S'/S)}\,v_\ell^{K'}
\end{equation*}
is independent of it, where

\begin{equation*}
e(S'/S)=\operatorname{length}(V'/\mathfrak mV').
\end{equation*}
We may again call (\SGARefLink{sga7.sga7i.eix.n9-7-1}{9.7.1}) the \underline{monodromy pairing} for the
abelian scheme $A_K$ over $K$ and the prime number $\ell$ under
consideration: when $A_K$ itself has semistable reduction, it plainly
coincides with the restriction to $\eta$ of the monodromy pairing (\SGARefLink{sga7.sga7i.eix.n9-1-2}{9.1.2}).
It follows from \SGARefLink{sga7.sga7i.eix.n10-4}{10.4} below that this pairing in fact always has values in
$\mathbb Q_K$ and is independent of $\ell$.

One may ask whether it has values in $\mathbb Z_K$, which is equivalent to
saying that (\SGARefLink{sga7.sga7i.eix.n9-7-1}{9.7.1}) has values in $\mathbb Z_{\ell K}$ for every prime
number $\ell$. This is so at least when $A_K$ is an elliptic curve. The
only case requiring verification is then the one in which
$\underline M_K\simeq\underline M'_K$ has rank $1$, so that
$\underline M_K\otimes\underline M'_K$ is canonically isomorphic to
$\mathbb Z_K$ and the monodromy pairing (\SGARefLink{sga7.sga7i.eix.n9-7-1}{9.7.1}) can simply be interpreted
as an ordinary rational number. It follows from the results of Néron [19]
that this number is $v(j)$, where $j$ is the absolute invariant of $A_K$
and $v$ is the normalized valuation of $K$.




\SGApage{I}{\(\approx\)439}
\medskip
\noindent
10. \underline{Properties of the monodromy pairing}. \underline{Integrality
and positivity theorem}

In this section we shall first give certain essentially trivial formal
properties of the monodromy pairings, leaving the (somewhat tedious)
details of the verifications to the reader. We shall also state a deeper
theorem (\SGARefLink{sga7.sga7i.eix.n10-4}{10.4}) concerning these pairings and indicate how its proof may be
reduced to a special case of Jacobians, to be treated (by transcendental
methods) in section 12, which is in fact logically independent of section
11.

\medskip
\noindent
\SGARefTarget{sga7.sga7i.eix.n10-1}{10.1}. \underline{Transport of structure}. It is clear from the definition
of the monodromy pairing (\SGARefLink{sga7.sga7i.eix.n9-1-2}{9.1.2}) that this pairing is compatible with
transport of structure in the situation $(S,A_K,A'_K)$.

\SGAvar{\medskip
\noindent
\SGARefTarget{sga7.sga7i.eix.n10-2}{10.2}. \underline{Symmetry}. Starting with the dual scheme $A'_K$ of $A_K$
and applying the construction of section 9 to it, we must introduce the
dual scheme $(A'_K)'$ of $A_K$, which is canonically isomorphic to $A_K$
by the biduality theorem ([5] or VIII 3.2). Denote by
$u_\ell^{A_K}$ the monodromy pairing (\SGARefLink{sga7.sga7i.eix.n9-1-2}{9.1.2}) relative to $A_K$. In the same
way we obtain a pairing

}{\medskip
\noindent
\SGARefTarget{sga7.sga7i.eix.n10-2}{10.2}. \underline{Symmetry}. Starting with the dual scheme $A'_K$ of $A_K$
and applying the construction of section 9 to it, we must introduce the
dual scheme $(A'_K)'$ of $A'_K$, which is canonically isomorphic to $A_K$
by the biduality theorem ([5] or VIII 3.2). Denote by
$u_\ell^{A_K}$ the monodromy pairing (\SGARefLink{sga7.sga7i.eix.n9-1-2}{9.1.2}) relative to $A_K$. In the same
way we obtain a pairing

}{E7I-IX-067}\begin{equation}\tag{\SGARefTarget{sga7.sga7i.eix.n10-2-1}{10.2.1}}
u_\ell^{A'_K}:\underline M'_\ell\otimes\underline M_\ell
 \longrightarrow\mathbb Z_\ell.
\end{equation}
I claim that the latter is obtained from $u_\ell^{A_K}$ by the symmetry

\begin{equation*}
s:\underline M'_\ell\otimes\underline M_\ell
 \xrightarrow{\ \sim\ }
 \underline M_\ell\otimes\underline M'_\ell;
\end{equation*}
that is,

\begin{equation}\tag{\SGARefTarget{sga7.sga7i.eix.n10-2-2}{10.2.2}}
u_\ell^{A'_K}=u_\ell^{A_K}\circ s.
\end{equation}

\SGApage{I}{\(\approx\)440}
Taking account of the known anticommutativity property of the pairings
$\varphi_\xi$ ([5] or VIII 2.2.11), the symmetry formula (\SGARefLink{sga7.sga7i.eix.n10-2-2}{10.2.2}) may be
reformulated as follows in terms of a fixed polarization $\xi$ of $A_K$.
It defines a homomorphism

\begin{equation*}
\widetilde\xi:A_K\longrightarrow A'_K
\end{equation*}
and hence a corresponding homomorphism ($t$ standing for ``toric'')

\begin{equation}\tag{\SGARefTarget{sga7.sga7i.eix.n10-2-3}{10.2.3}}
\xi_t:\underline M\longrightarrow\underline M'.
\end{equation}
We may therefore form the pairing

\begin{equation}\tag{\SGARefTarget{sga7.sga7i.eix.n10-2-4}{10.2.4}}
\begin{aligned}
\varphi^t_{\xi,\ell}:\underline M_\ell\otimes\underline M_\ell
 &\longrightarrow\mathbb Z_\ell,\\
\varphi^t_{\xi,\ell}(x,y)&=u_\ell^{A_K}\bigl(x,\xi_t(y)\bigr).
\end{aligned}
\end{equation}
The preceding form $\varphi^t_{\xi,\ell}$ is \underline{symmetric}:

\begin{equation}\tag{\SGARefTarget{sga7.sga7i.eix.n10-2-5}{10.2.5}}
\varphi^t_{\xi,\ell}(x,y)=\varphi^t_{\xi,\ell}(y,x).
\end{equation}
In fact, this assertion is valid for every element $\xi$ of the
Néron--Severi group of $A_K$, not necessarily a polarization. But only
when $\xi$ is nondegenerate---that is, defines an \underline{isogeny}
$A_K\longrightarrow A'_K$, and hence also an isogeny (\SGARefLink{sga7.sga7i.eix.n10-2-3}{10.2.3})---is the
symmetry relation (\SGARefLink{sga7.sga7i.eix.n10-2-5}{10.2.5}) equivalent to (\SGARefLink{sga7.sga7i.eix.n10-2-2}{10.2.2}). When $\xi$ is a
polarization, we shall moreover sharpen (\SGARefLink{sga7.sga7i.eix.n10-2-5}{10.2.5}) considerably in \SGARefLink{sga7.sga7i.eix.n10-4}{10.4}(b)
below.

Let us finally mention a third interpretation of the symmetry relation.
Consider a biextension (that is, a divisorial correspondence) $\xi$ on a
pair $(A_K,\overline A_K)$ of abelian schemes with semistable reduction
over $K$. It defines a homomorphism

\begin{equation*}
A_K\longrightarrow\overline A'_K.
\end{equation*}
Passing to the Néron models, this defines a homomorphism on the maximal
tori of the special fibres, and therefore, in the reverse direction, a
homomorphism on their character groups:

\SGApage{I}{\(\approx\)441}
\begin{equation}\tag{\SGARefTarget{sga7.sga7i.eix.n10-2-6}{10.2.6}}
\xi_t:\overline{\underline M}\longrightarrow\underline M',
\end{equation}
where $\overline{\underline M}$ is defined from $\overline A_K$ as
$\underline M$ is defined from $A_K$. We can therefore form the pairing

\SGAvar{\begin{equation}\tag{\SGARefTarget{sga7.sga7i.eix.n10-2-7}{10.2.7}}
\begin{aligned}
\varphi^t_\xi:\underline M_\ell\otimes
 \overline{\underline M}_\ell&\longrightarrow\mathbb Z_\ell,\\
\varphi^t_\xi(x,y)&=\varphi^t\bigl(x,\xi_t(y)\bigr).
\end{aligned}
\end{equation}

}{\begin{equation}\tag{\SGARefTarget{sga7.sga7i.eix.n10-2-7}{10.2.7}}
\begin{aligned}
\varphi^t_\xi:\underline M_\ell\otimes
 \overline{\underline M}_\ell&\longrightarrow\mathbb Z_\ell,\\
\varphi^t_\xi(x,y)&=u_\ell^{A_K}\bigl(x,\xi_t(y)\bigr).
\end{aligned}
\end{equation}

}{E7I-IX-068}Notice in passing that, for a variable abelian scheme $A_K$ with semistable
reduction, the twisted constant group $\underline M$ is a covariant
functor of $A_K$. Formation of the pairings (\SGARefLink{sga7.sga7i.eix.n10-2-7}{10.2.7}) associated with
biextensions of abelian schemes with semistable reduction over $(K,S)$ is
then compatible, in the evident sense, with taking inverse images of
biextensions.

\SGAvar{When $\xi$ is the canonical (``Weil'') biextension of $(A_K,A'_K)$, the
pairing (\SGARefLink{sga7.sga7i.eix.n10-2-7}{10.2.7}) is of course precisely the monodromy pairing (\SGARefLink{sga7.sga7i.eix.n9-1-2}{9.2.1}).

}{When $\xi$ is the canonical (``Weil'') biextension of $(A_K,A'_K)$, the
pairing (\SGARefLink{sga7.sga7i.eix.n10-2-7}{10.2.7}) is of course precisely the monodromy pairing (\SGARefLink{sga7.sga7i.eix.n9-1-2}{9.1.2}).

}{E7I-IX-069}Now consider the transpose ${}^s\xi$ of $\xi$, a biextension of
$(\overline A_K,A_K)$, and hence the pairing

\begin{equation*}
\varphi^t_{{}^s\xi}:\overline{\underline M}_\ell\otimes
 \underline M_\ell\longrightarrow\mathbb Z_\ell.
\end{equation*}
The symmetry relation is

\begin{equation}\tag{\SGARefTarget{sga7.sga7i.eix.n10-2-8}{10.2.8}}
\varphi^t_{{}^s\xi}=\varphi^t_\xi\circ s,
\qquad
s:\overline{\underline M}_\ell\otimes\underline M_\ell
 \xrightarrow{\ \sim\ }
 \underline M_\ell\otimes\overline{\underline M}_\ell,
\end{equation}
where $s$ is again the symmetry isomorphism. Compare this relation with
the antisymmetry relation of VIII \SGARefLink{sga7.sga7i.eviii.n2-2-11}{2.2.11}.




\SGApage{I}{442}
\medskip
\noindent
\SGARefTarget{sga7.sga7i.eix.n10-3}{10.3}. \underline{Change of trait}. Let

\begin{equation}\tag{\SGARefTarget{sga7.sga7i.eix.n10-3-1}{10.3.1}}
g:S'\longrightarrow S
\end{equation}
be a dominant morphism of henselian traits, associated with an injective
local homomorphism of discrete valuation rings

\begin{equation*}
V\longrightarrow V'.
\end{equation*}
As usual, put

\begin{equation}\tag{\SGARefTarget{sga7.sga7i.eix.n10-3-2}{10.3.2}}
e(S'/S)=\operatorname{length}_V\bigl(V'/\underline mV'\bigr),
\end{equation}
where $\underline m$ is the maximal ideal of $V$. It is then known that
the homomorphism induced by (\SGARefLink{sga7.sga7i.eix.n10-3-1}{10.3.1}) on the ``tame'' parts of the inertia
groups,

\begin{equation}\tag{\SGARefTarget{sga7.sga7i.eix.n10-3-3}{10.3.3}}
g_*:I'_t\simeq\prod_{\ell\ne p}\mathbb Z_\ell(1)(k')
 \longrightarrow
 I_t=\prod_{\ell\ne p}\mathbb Z_\ell(1)(k),
\end{equation}
is given by

\begin{equation}\tag{\SGARefTarget{sga7.sga7i.eix.n10-3-4}{10.3.4}}
g_*(x)=e(S'/S)x,
\end{equation}
after the two sides of (\SGARefLink{sga7.sga7i.eix.n10-3-3}{10.3.3}) have been identified by means of the
evident isomorphisms
$\mathbb Z_\ell(1)(k)\xrightarrow{\ \sim\ }\mathbb Z_\ell(1)(k')$.


\SGAvar{On the other hand, if $A_K$ is an abelian scheme over $K$ with
semistable reduction over $S$, formation of the twisted lattices
$\underline M$ and $\underline M'$ commutes with every base change
$S'\longrightarrow S$ as above, as follows immediately from the
definitions (especially \SGARefLink{sga7.sga7i.eix.n3-4}{3.4}). We may therefore regard the pairing
$u_\ell^{A_{K}}$ defined by $A_{K'}$, where $K'$ is the fraction field
of $V'$, as a pairing

\begin{equation*}
u_\ell^{A_{K'}}:
 \underline M_{\ell S'}\otimes\underline M'_{\ell S'}
 \longrightarrow\mathbb Z_{\ell S'},
\end{equation*}
and compare it with the pairing $u_{\ell S'}^{A_K}$ obtained from
$u_\ell^{A_K}$ by base change. When $\ell\ne p$, the construction of the
pairings $u_\ell$ given in \SGARefLink{sga7.sga7i.eix.n9-2}{9.2}, together with formula (\SGARefLink{sga7.sga7i.eix.n10-3-4}{10.2.3}), shows at
once that

}{On the other hand, if $A_K$ is an abelian scheme over $K$ with
semistable reduction over $S$, formation of the twisted lattices
$\underline M$ and $\underline M'$ commutes with every base change
$S'\longrightarrow S$ as above, as follows immediately from the
definitions (especially \SGARefLink{sga7.sga7i.eix.n3-4}{3.4}). We may therefore regard the pairing
$u_\ell^{A_{K'}}$ defined by $A_{K'}$, where $K'$ is the fraction field
of $V'$, as a pairing

\begin{equation*}
u_\ell^{A_{K'}}:
 \underline M_{\ell S'}\otimes\underline M'_{\ell S'}
 \longrightarrow\mathbb Z_{\ell S'},
\end{equation*}
and compare it with the pairing $u_{\ell S'}^{A_K}$ obtained from
$u_\ell^{A_K}$ by base change. When $\ell\ne p$, the construction of the
pairings $u_\ell$ given in \SGARefLink{sga7.sga7i.eix.n9-2}{9.2}, together with formula (\SGARefLink{sga7.sga7i.eix.n10-3-4}{10.3.4}), shows at
once that

}{E7I-IX-070}\SGApage{I}{\(\approx\)443}
\begin{equation}\tag{\SGARefTarget{sga7.sga7i.eix.n10-3-5}{10.3.5}}
u_\ell^{A_{K'}}=e(S'/S)u_\ell^{A_K}.
\end{equation}
In fact this formula remains valid for $\ell=p$. To see this, return to
the general construction of \SGARefLink{sga7.sga7i.eix.n9-3}{9.3}--\SGARefLink{sga7.sga7i.eix.n9-6}{9.6} and reduce to the analogous
compatibility assertion for the isomorphism \SGARefLink{sga7.sga7i.eix.n9-4-4}{9.4.4}(e), under a change of
strictly local traits $S'\longrightarrow S$. Notice that the isomorphism
in question is induced by the normalized valuation
$K^*/V^*\simeq\mathbb Z$, and therefore gives the commutative diagram


\begin{center}
\begin{tikzcd}
K^*/V^* \arrow[r, "\sim"] \arrow[d]
  \& \mathbb Z \arrow[d, "e(S'/S)"] \\
K^{\prime *}/V^{\prime *} \arrow[r, "\sim"]
  \& \mathbb Z .
\end{tikzcd}
\end{center}

Here is a deeper result concerning the monodromy pairings. We shall prove
it by transcendental methods in section 12, relying on results that will
be proved later in the Seminar. (We shall not, in fact, use \SGARefLink{sga7.sga7i.eix.n10-4}{10.4} in the
remainder of the Seminar, any more than we shall use the other results of
the present expos\'e.)

\SGApage{I}{444}
\medskip
\noindent
\underline{Theorem} \SGARefTarget{sga7.sga7i.eix.n10-4}{10.4}. \underline{Let $S$ be a henselian trait, and
let $A_K$ be an abelian scheme over $K$ with semistable reduction over
$S$. Consider the lattices $\underline M$ and $\underline M'$ associated
as in \SGARefLink{sga7.sga7i.eix.n9-1}{9.1} (cf. \SGARefLink{sga7.sga7i.eix.n5-5-2}{5.5.2}). Then the following assertions hold:}

\smallskip
\noindent
a) \underline{There exists a unique pairing}

\begin{equation}\tag{\SGARefTarget{sga7.sga7i.eix.n10-4-1}{10.4.1}}
u:\underline M\otimes\underline M'\longrightarrow\mathbb Z_S
\end{equation}
\underline{such that, for every prime number $\ell$,
$u\otimes_{\mathbb Z}\mathbb Z_\ell$ is the monodromy pairing $u_\ell$
of (\SGARefLink{sga7.sga7i.eix.n9-1-2}{9.1.2}). This pairing is separating.}

\smallskip
\noindent
b) \underline{Let $\xi$ be a polarization of $A_K$, thus defining a
homomorphism (\SGARefLink{sga7.sga7i.eix.n10-2-3}{10.2.3})}

\begin{equation*}
\xi_t:\underline M\longrightarrow\underline M',
\end{equation*}
\underline{and consider the form}

\begin{equation}\tag{\SGARefTarget{sga7.sga7i.eix.n10-4-2}{10.4.2}}
\begin{aligned}
\varphi_\xi^t:\underline M\otimes\underline M
  &\longrightarrow\mathbb Z_S,\\
\varphi_\xi^t(x,y)&=u\bigl(x,\xi_t(y)\bigr).
\end{aligned}
\end{equation}
\underline{This form is symmetric, positive, and nondegenerate.}

Of course, in terms of the geometric fibre $M_o$ of $\underline M$ at a
given geometric point of $S$, the preceding assertion says that the
corresponding form

\begin{equation*}
M_o\otimes M_o\longrightarrow\mathbb Z
\end{equation*}
is symmetric, positive, and nondegenerate. The pairing (\SGARefLink{sga7.sga7i.eix.n10-4-1}{10.4.1}) will
again be called the \underline{monodromy pairing}.



\SGApage{I}{445}
\medskip
\noindent
\SGARefTarget{sga7.sga7i.eix.n10-5}{10.5}. \underline{Reduction of \SGARefLink{sga7.sga7i.eix.n10-4}{10.4} to a particular case}

\smallskip
\noindent
\SGARefTarget{sga7.sga7i.eix.n10-5-1}{10.5.1}. The uniqueness of a $u$ satisfying the stated condition (even
for a single $\ell$) is clear. On the other hand, existence simply means
that every $u_\ell$ in fact maps $\underline M\otimes\underline M'$ into
$\mathbb Z_S$, and that the homomorphism thereby obtained

\begin{equation*}
\underline M\otimes\underline M'\longrightarrow\mathbb Z_S
\end{equation*}
is independent of $\ell$. In fact it even suffices to prove the analogous
assertion for the $u_\ell\otimes_{\mathbb Z_\ell}\mathbb Q_\ell$ and for
a homomorphism

\begin{equation*}
u_{\mathbb Q}:\underline M_{\mathbb Q}\otimes
 \underline M'_{\mathbb Q}\longrightarrow\mathbb Q_S.
\end{equation*}

Suppose, then, that existence has been established after a base change
$S'\longrightarrow S$ as in \SGARefLink{sga7.sga7i.eix.n9-3}{9.3}. It follows immediately from (\SGARefLink{sga7.sga7i.eix.n10-3-5}{10.3.5})
that there exists a homomorphism

\SGAvar{\begin{equation*}
u:\underline M\otimes\underline M'
 \longrightarrow\frac{1}{e(S'/S)}\mathbb Z_S\subset Q_S
\end{equation*}
compatible with the $u_\ell\otimes_{\mathbb Z_\ell}\mathbb Q_\ell$;
this proves a). As for b), it is then plain that, if it holds for
$A_{K'}$, it also holds for $A_K$.

}{\begin{equation*}
u:\underline M\otimes\underline M'
 \longrightarrow\frac{1}{e(S'/S)}\mathbb Z_S\subset\mathbb Q_S
\end{equation*}
compatible with the $u_\ell\otimes_{\mathbb Z_\ell}\mathbb Q_\ell$;
this proves a). As for b), it is then plain that, if it holds for
$A_{K'}$, it also holds for $A_K$.

}{E7I-IX-071}Consequently, in proving \SGARefLink{sga7.sga7i.eix.n10-4}{10.4} we may make any dominant change of traits
$S'\longrightarrow S$. In particular, this allows us to suppose that
$S$ is complete with algebraically closed residue field. In that case
$\underline M$ and $\underline M'$ are constant group schemes, with
values $M$ and $M'$, finitely generated free $\mathbb Z$-modules, and a)
and b) may be reformulated more conveniently in terms of the latter.

\SGAvar{\smallskip
\noindent
\SGARefTarget{sga7.sga7i.eix.n10-5-2}{10.5.2}. Suppose now that $A_K$ is isomorphic to a direct factor of an
abelian scheme $\overline A_K$ over $k$ with semistable reduction, and
that \SGARefLink{sga7.sga7i.eix.n10-4}{10.4} has been proved for the latter. It follows at once that \SGARefLink{sga7.sga7i.eix.n10-4}{10.4}
also holds for $A_K$, in view of the evident functorial behaviour of the
lattices $\underline M,\underline M'$ and of the monodromy pairings with
respect to direct sums of abelian schemes with semistable reduction.


}{\smallskip
\noindent
\SGARefTarget{sga7.sga7i.eix.n10-5-2}{10.5.2}. Suppose now that $A_K$ is isomorphic to a direct factor of an
abelian scheme $\overline A_K$ over $K$ with semistable reduction, and
that \SGARefLink{sga7.sga7i.eix.n10-4}{10.4} has been proved for the latter. It follows at once that \SGARefLink{sga7.sga7i.eix.n10-4}{10.4}
also holds for $A_K$, in view of the evident functorial behaviour of the
lattices $\underline M,\underline M'$ and of the monodromy pairings with
respect to direct sums of abelian schemes with semistable reduction.


}{E7I-IX-072}\SGApage{I}{\(\approx\)446}
It is equally immediate from \SGARefLink{sga7.sga7i.eix.n10-5-1}{10.5.1} that the validity of \SGARefLink{sga7.sga7i.eix.n10-4}{10.4} depends
only on the isogeny class of $A_K$. Together with the preceding remark
and the ``Poincar\'e reducibility theorem'' [15], this shows that, if
$A_K$ is a quotient of an abelian scheme $\overline A_K$ for which \SGARefLink{sga7.sga7i.eix.n10-4}{10.4}
holds, then \SGARefLink{sga7.sga7i.eix.n10-4}{10.4} also holds for $A_K$.


\smallskip
\noindent
\SGARefTarget{sga7.sga7i.eix.n10-5-3}{10.5.3}. We now use the well-known fact [31, Th.~11 (p.~20)] that, after
making a finite extension of $K$ if necessary (which is permissible by
\SGARefLink{sga7.sga7i.eix.n10-5-1}{10.5.1}), every abelian scheme $A_K$ is a quotient of the Jacobian

\SGAvar{\begin{equation*}
\underline{\operatorname{Pic}}^0_{X_K/K}
\end{equation*}
of a proper, smooth, geometrically connected curve $X_K$ over $X$.
Thus, by \SGARefLink{sga7.sga7i.eix.n10-5-2}{10.5.2}, we are reduced to proving \SGARefLink{sga7.sga7i.eix.n10-4}{10.4} for such a Jacobian.


}{\begin{equation*}
\underline{\operatorname{Pic}}^0_{X_K/K}
\end{equation*}
of a proper, smooth, geometrically connected curve $X_K$ over $K$.
Thus, by \SGARefLink{sga7.sga7i.eix.n10-5-2}{10.5.2}, we are reduced to proving \SGARefLink{sga7.sga7i.eix.n10-4}{10.4} for such a Jacobian.


}{E7I-IX-073}\smallskip
\noindent
\SGARefTarget{sga7.sga7i.eix.n10-5-4}{10.5.4}. By a theorem of Artin--Mumford [3] (already cited and used in
Expos\'e V), after making a further finite extension $K'$ of $K$ if
necessary, we may suppose that $X_K$ extends to a projective and flat
scheme over $S$ which is regular and whose geometric special fibre
$X_{\bar s}$ has no singularities other than ordinary quadratic
singularities (that is, here, normal crossings of two branches).

In this case we shall prove \SGARefLink{sga7.sga7i.eix.n10-4}{10.4} a) in \SGARefLink{sga7.sga7i.eix.n12-5}{12.5} below, by describing the
monodromy pairing (\SGARefLink{sga7.sga7i.eix.n10-4-1}{10.4.1}) completely in terms of the geometric data of
the situation; at the same time we shall verify \SGARefLink{sga7.sga7i.eix.n10-4}{10.4} b) for the canonical
polarization of the Jacobian.

\SGApage{I}{447}
\smallskip
\noindent
\SGARefTarget{sga7.sga7i.eix.n10-5-5}{10.5.5}. It remains to see that, if for a given $A_K$ with semistable
reduction assertion \SGARefLink{sga7.sga7i.eix.n10-4}{10.4} a) holds and assertion \SGARefLink{sga7.sga7i.eix.n10-4}{10.4} b) holds for one
given polarization $\xi_o$ of $A_K$, then the same conclusion holds for
every polarization $\xi$ of $A_K$. Indeed, let

\begin{equation*}
\widetilde\xi_o,\widetilde\xi:A_K\longrightarrow A'_K
\end{equation*}
be the homomorphisms defined by $\xi_o$ and $\xi$; these are isogenies.
Then

\begin{equation*}
\widetilde\xi=\widetilde\xi_o\alpha
 ={}^t(\alpha^*)\widetilde\xi_o,
 \qquad
\alpha\in\operatorname{End}(A_K)\otimes_{\mathbb Z}\mathbb Q,
\end{equation*}
where $\alpha\longmapsto\alpha^*$ denotes the canonical involution of
the algebra $\operatorname{End}(A_K)\otimes_{\mathbb Z}\mathbb Q$
defined by the polarization $\xi_o$, and the superscript $t$ in
${}^t(\alpha^*)$ denotes transpose (associating with an endomorphism of
$A_K$ an endomorphism of $A'_K$). On the other hand, it is known
[15, Chap.~V, Th.~3] that

\begin{equation*}
\alpha=\sum_i\beta_i^*\beta_i,
 \qquad
\beta_i\in\operatorname{End}(A_K)\otimes_{\mathbb Z}\mathbb Q,
\end{equation*}
and hence

\begin{equation*}
\widetilde\xi=\sum_i{}^t\beta_i\,\widetilde\xi_o\,\beta_i.
\end{equation*}
Equivalently, in terms of the biextensions (that is, divisorial
correspondences) associated with the polarizations under consideration
(VIII \SGARefLink{sga7.sga7i.eviii.n4-14}{4.14}),

\begin{equation}\tag{\SGARefTarget{sga7.sga7i.eix.n10-5-5-1}{10.5.5.1}}
\delta(\xi)=\sum_i\delta(\xi_o)\circ(\beta_i,\beta_i).
\end{equation}
By what was said after (\SGARefLink{sga7.sga7i.eix.n10-2-7}{10.2.7}), this formula implies (with the notation
of (\SGARefLink{sga7.sga7i.eix.n10-4-2}{10.4.2}))

\begin{equation}\tag{\SGARefTarget{sga7.sga7i.eix.n10-5-5-2}{10.5.5.2}}
\varphi^t_\xi
 =\sum_i\varphi^t_{\xi_o}(\beta_i,\beta_i)
 \in\operatorname{Hom}
 \bigl(M_{\mathbb Q}\otimes_{\mathbb Q}M_{\mathbb Q},\mathbb Q\bigr).
\end{equation}
It indeed follows that, if $\varphi^t_{\xi_o}$ is symmetric and positive,
then so is $\varphi^t_\xi$.

\SGApage{I}{\(\approx\)448}
Finally, the nondegeneracy of $\varphi^t_{\xi_o}$ is plainly equivalent
to the assertion that the pairing (\SGARefLink{sga7.sga7i.eix.n10-4-1}{10.4.1}) is separating (since
$\xi_{ot}:M\otimes_{\mathbb Z}\mathbb Q
\longrightarrow M'\otimes_{\mathbb Z}\mathbb Q$ is an isomorphism,
$\widetilde\xi_o$ being an isogeny), and is therefore likewise
independent of the chosen polarization. One may also observe that, in
order to verify that (\SGARefLink{sga7.sga7i.eix.n10-4-1}{10.4.1}) is separating, it suffices to verify the
same assertion for the monodromy pairing $u_\ell$ of (\SGARefLink{sga7.sga7i.eix.n9-1-2}{9.1.2}) for one
prime number $\ell$. Taking $\ell\ne p$, one may use the description of
$u_\ell$ in \SGARefLink{sga7.sga7i.eix.n9-2}{9.2}; the fact that $u_\ell$ is separating on the left (and
hence on both sides, for reasons of rank) is then immediate.



\medskip
\noindent
11. \underline{Connected components of the N\'eron model: duality
relation and asymptotic behaviour}

\smallskip
\noindent
\SGARefTarget{sga7.sga7i.eix.n11-0}{11.0}. The purpose of the present section is to study the finite \'etale
$k$-group

\begin{equation}\tag{\SGARefTarget{sga7.sga7i.eix.n11-0-1}{11.0.1}}
\Phi_o=A_o/A_o^{,o}
\end{equation}
already considered in section 1. For every prime number $\ell$, denote
by

\SGAvar{\begin{equation*}
\Phi_o(\ell)
\end{equation*}
its $\ell$-primary component. Thus knowledge of $\Phi_o$ is equivalent
to knowledge of the $\Phi_o(1)$, by the formula

}{\begin{equation*}
\Phi_o(\ell)
\end{equation*}
its $\ell$-primary component. Thus knowledge of $\Phi_o$ is equivalent
to knowledge of the $\Phi_o(\ell)$, by the formula

}{E7I-IX-074}\begin{equation}\tag{\SGARefTarget{sga7.sga7i.eix.n11-0-2}{11.0.2}}
\Phi_o=\coprod_\ell\Phi_o(\ell).
\end{equation}

\SGApage{I}{\(\approx\)449}
Notice that the groups $\Phi_o$ and $\Phi_o(\ell)$ are determined,
respectively, by the ordinary finite groups $\Phi_o(\bar k)$ and
$\Phi_o(\ell)(\bar k)$ together with the natural action of

\begin{equation*}
\pi_o=\operatorname{Gal}(\bar k/k)
\end{equation*}
on them. On the other hand, it is clear that $\Phi_o(\bar k)$ is also
the group of connected components of the N\'eron model of
$A_{\widetilde K}$, where $\widetilde K$ is the fraction field of the
strict henselization of $S$ relative to the chosen separable closure
$\bar k$ of $k$. Thus, in studying $\Phi_o$, we may reduce to the case
where $S$ is strictly local; the actions of

\begin{equation*}
\pi_o\simeq\operatorname{Gal}(\widetilde K/K)
\end{equation*}
on $\Phi_o(\bar k)$ are then recovered automatically by transport of
structure. We could likewise reduce further to the case where $S$ is
complete, using the base change $\widehat S\longrightarrow S$ and the
fact that formation of the N\'eron model commutes with it
[21, Th.~\SGARefLink{sga7.sga7i.eix.n3-4}{3.4} b)].

\smallskip
\noindent
\SGARefTarget{sga7.sga7i.eix.n11-1}{11.1}. Suppose, then, that $S$ is strictly local, so that $\Phi_o$ may be
identified with the ordinary finite group $\Phi_o(k)$. We obtain

\begin{equation}\tag{\SGARefTarget{sga7.sga7i.eix.n11-1-1}{11.1.1}}
\Phi_o(k)\simeq A_o(k)/A_o^o(k)
 \simeq A(S)/A^o(S)
 \simeq A(K)/A(K)^o,
\end{equation}
where the first equality follows from the fact that $A_o$ is smooth over
$k$, the second from Hensel's lemma since $A$ is smooth over $S$ and
$S$ is henselian, and the third from the relation
$A(S)\xrightarrow{\sim}A(K)$ and from the convention

\begin{equation}\tag{\SGARefTarget{sga7.sga7i.eix.n11-1-2}{11.1.2}}
A(K)^o\mathrel{\mathop:}=A^o(S).
\end{equation}
Thus $A(K)^o$ appears as a natural subgroup of finite index in $A(K)$
(and is not to be confused with $A^o(K)=A(K)$). The isomorphisms
(\SGARefLink{sga7.sga7i.eix.n11-1-1}{11.1.1}) are plainly functorial in $A$ and compatible with transport of
structure in $(S,A_K)$, and in particular with the action of any Galois
groups.

\SGApage{I}{\(\approx\)450}
Now let $\ell$ be a prime number different from $p$. Then $A^o$ is
$\ell$-divisible; more precisely, $\ell\operatorname{id}_{A^o}$ is an
\'etale surjective morphism. It consequently induces a surjective
endomorphism of $A^o(S)=A(K)^o$, which is therefore an
\underline{$\ell$-divisible group}. It follows at once that, for every
integer $n\geq 0$ prime to $p=\operatorname{char}(k)$,

\begin{equation}\tag{\SGARefTarget{sga7.sga7i.eix.n11-1-3}{11.1.3}}
{}_n\Phi_o(k)\simeq{}_nA(K)/{}_nA(K)^o,
 \qquad (n,p)=1.
\end{equation}
By the definition (\SGARefLink{sga7.sga7i.eix.n2-2-3}{2.2.3}) of the fixed part of ${}_nA_K$, we have

\begin{equation}\tag{\SGARefTarget{sga7.sga7i.eix.n11-1-4}{11.1.4}}
{}_nA(K)^o=({}_nA_K)^f(\bar K)
 \simeq T_n(A_K)^f(\bar K)
 \otimes_{\mathbb Z_n}\mathbb Z/n\mathbb Z.
\end{equation}
Taking $n$ to be a power of $\ell$, and putting as in (\SGARefLink{sga7.sga7i.eix.n2-5-3}{2.5.3})

\begin{equation}\tag{\SGARefTarget{sga7.sga7i.eix.n11-1-5}{11.1.5}}
U=T_\ell\bigl(A(\bar K)\bigr),
 \qquad
T_\ell\bigl(A(\bar K)\bigr)^f=U^I,
\end{equation}
where $U^I$ denotes the submodule of $U$ consisting of invariants under
the inertia group $I$, we may write (\SGARefLink{sga7.sga7i.eix.n11-1-3}{11.1.3}) in the form

\begin{equation}\tag{\SGARefTarget{sga7.sga7i.eix.n11-1-6}{11.1.6}}
{}_n\Phi_o(k)\simeq (U_n)^I/(U^I)_n,
 \qquad n=\ell^{r+1},
\end{equation}
where the subscript $n$ denotes reduction modulo $n$, that is, tensor
product over $\mathbb Z_\ell$ with
$\mathbb Z_\ell/n\mathbb Z_\ell$. The transition morphisms between the
groups occurring on the two sides of (\SGARefLink{sga7.sga7i.eix.n11-1-6}{11.1.6}) being evident, passage to
the inductive limit over the powers of $\ell$ gives the following
result.

\SGApage{I}{451}
\medskip
\noindent
\underline{Proposition \SGARefTarget{sga7.sga7i.eix.n11-2}{11.2}.} \underline{Suppose that $S$ is strictly
local, and let $\ell$ be a prime number different from $p$. Then there is
a canonical isomorphism}

\begin{equation}\tag{\SGARefTarget{sga7.sga7i.eix.n11-2-1}{11.2.1}}
\Phi_o(\ell)\simeq
 \frac{(U\otimes D_\ell)^I}{U^I\otimes D_\ell},
\end{equation}
\underline{where $U$ is defined in (\SGARefLink{sga7.sga7i.eix.n11-1-5}{11.1.5}) and
$D_\ell=\mathbb Q_\ell/\mathbb Z_\ell$.}

Thus the group $\Phi_o(\ell)$ can be described solely in terms of the
Tate module $U=T_\ell(A(\bar K))$, regarded as a Galois module under the
action of the inertia group $I$.



\medskip
\noindent
\SGARefTarget{sga7.sga7i.eix.n11-3}{11.3}. We shall use the explicit expression (\SGARefLink{sga7.sga7i.eix.n11-2-1}{11.2.1}) to sketch the
definition of a canonical pairing

\begin{equation}\tag{\SGARefTarget{sga7.sga7i.eix.n11-3-1}{11.3.1}}
\Phi_o(\ell)\otimes\Phi'_o(\ell)\longrightarrow\mathbb Q/\mathbb Z
\end{equation}
which is a perfect duality; here $\Phi'_o$ denotes, as in section 1, the
group of connected components associated with the N\'eron model of the
\underline{dual} $A'_K$ of $A_K$. There is little doubt (although the
author has dispensed with the verification) that this pairing is none
other than, possibly up to sign, the one defined in (\SGARefLink{sga7.sga7i.eix.n1-2-1}{1.2.1}), and the
reader is invited, as an exercise, to establish the required
compatibility (with the correct sign!).

Put

\begin{equation*}
U'=T_\ell\bigl(A'(\bar K)\bigr),
 \qquad
T=T_\ell(1)(\bar K)=T_\ell(\bar K^*).
\end{equation*}
There is a \underline{perfect pairing} of finitely generated free
$\mathbb Z_\ell$-modules, compatible with the actions of $I$,

\begin{equation}\tag{\SGARefTarget{sga7.sga7i.eix.n11-3-2}{11.3.2}}
\varphi:U\otimes U'\longrightarrow T,
 \qquad (T\simeq\mathbb Z_\ell\text{ noncanonically}),
\end{equation}
and, by the known structure of $I$, an exact sequence


\SGApage{I}{\(\approx\)452}
\begin{equation}\tag{\SGARefTarget{sga7.sga7i.eix.n11-3-3}{11.3.3}}
1\longrightarrow R\longrightarrow I\longrightarrow T\longrightarrow1,
\end{equation}
where $R$ is a profinite group \underline{prime to} $\ell$. Starting now
from the data (\SGARefLink{sga7.sga7i.eix.n11-3-2}{11.3.2}) and (\SGARefLink{sga7.sga7i.eix.n11-3-3}{11.3.3}), introduce the groups

\SGAvar{\begin{equation}\tag{\SGARefTarget{sga7.sga7i.eix.n11-3-4}{11,3.4}}
\begin{aligned}
\Phi&=\frac{(U\otimes D_\ell)^I}{U^I\otimes D_\ell},\qquad
\Phi'=\frac{(U'\otimes D_\ell)^I}{{U'}^I\otimes D_\ell},\qquad
D_\ell=\mathbb Q_\ell/\mathbb Z_\ell.
\end{aligned}
\end{equation}
We shall define a perfect duality

}{\begin{equation}\tag{\SGARefTarget{sga7.sga7i.eix.n11-3-4}{11.3.4}}
\begin{aligned}
\Phi&=\frac{(U\otimes D_\ell)^I}{U^I\otimes D_\ell},\qquad
\Phi'=\frac{(U'\otimes D_\ell)^I}{{U'}^I\otimes D_\ell},\qquad
D_\ell=\mathbb Q_\ell/\mathbb Z_\ell.
\end{aligned}
\end{equation}
We shall define a perfect duality

}{E7I-IX-075}\begin{equation}\tag{\SGARefTarget{sga7.sga7i.eix.n11-3-5}{11.3.5}}
\Phi\otimes\Phi'\longrightarrow\mathbb Q/\mathbb Z,
\end{equation}
which will in particular give the announced pairing (\SGARefLink{sga7.sga7i.eix.n11-3-1}{11.3.1}). Consider
the exact sequence of $I$-modules

\begin{equation}\tag{\SGARefTarget{sga7.sga7i.eix.n11-3-6}{11.3.6}}
0\longrightarrow U\xrightarrow{\ i\ }
 U\otimes_{\mathbb Z_\ell}\mathbb Q_\ell
 \xrightarrow{\ j\ }U\otimes D_\ell\longrightarrow0
\end{equation}
and the associated exact cohomology sequence, where the cohomology of
$I$ is defined by \underline{continuous cochains} with coefficients in
the topological groups under consideration:

\SGAvar{\begin{equation}\tag{\SGARefTarget{sga7.sga7i.eix.n11-3-7}{11.3.7}}
\resizebox{\textwidth}{!}{
\begin{tikzcd}[ampersand replacement=\&, column sep=small, row sep=large]
0 \arrow[r]
 \& H^o(I,U) \arrow[r, "i^o"] \arrow[d, no head, "\wr"]
 \& H^o(I,U\otimes\mathbb Q_\ell) \arrow[r, "j^o"]
      \arrow[d, no head, "\wr"]
 \& H^o(I,U\otimes D_\ell) \arrow[r, "\partial"]
      \arrow[d, no head, "\wr"]
 \& H^1(I,U) \arrow[r, "i^1"]
 \& H^1(I,U\otimes\mathbb Q_\ell) \arrow[d, no head, "\wr"] \\
{}
 \& U^I
 \& (U\otimes Q_\ell)^I\simeq U^I\otimes_{\mathbb Z}\mathbb Q
 \& (U\otimes D_\ell)^I
 \& {}
 \& H^1(I,U)\otimes_{\mathbb Z}\mathbb Q .
\end{tikzcd}}
\end{equation}
It is immediate that the group $\Phi$ defined by (\SGARefLink{sga7.sga7i.eix.n11-3-4}{11.3.4}) is precisely
$\operatorname{coker}j^o$, and is therefore canonically isomorphic to
$\ker i^1$. Thus we obtain a canonical exact sequence

}{\begin{equation}\tag{\SGARefTarget{sga7.sga7i.eix.n11-3-7}{11.3.7}}
\resizebox{\textwidth}{!}{
\begin{tikzcd}[ampersand replacement=\&, column sep=small, row sep=large]
0 \arrow[r]
 \& H^o(I,U) \arrow[r, "i^o"] \arrow[d, no head, "\wr"]
 \& H^o(I,U\otimes\mathbb Q_\ell) \arrow[r, "j^o"]
      \arrow[d, no head, "\wr"]
 \& H^o(I,U\otimes D_\ell) \arrow[r, "\partial"]
      \arrow[d, no head, "\wr"]
 \& H^1(I,U) \arrow[r, "i^1"]
 \& H^1(I,U\otimes\mathbb Q_\ell) \arrow[d, no head, "\wr"] \\
{}
 \& U^I
 \& (U\otimes\mathbb Q_\ell)^I\simeq U^I\otimes_{\mathbb Z}\mathbb Q
 \& (U\otimes D_\ell)^I
 \& {}
 \& H^1(I,U)\otimes_{\mathbb Z}\mathbb Q .
\end{tikzcd}}
\end{equation}
It is immediate that the group $\Phi$ defined by (\SGARefLink{sga7.sga7i.eix.n11-3-4}{11.3.4}) is precisely
$\operatorname{coker}j^o$, and is therefore canonically isomorphic to
$\ker i^1$. Thus we obtain a canonical exact sequence

}{E7I-IX-076}\begin{equation}\tag{\SGARefTarget{sga7.sga7i.eix.n11-3-8}{11.3.8}}
0\longrightarrow\Phi\longrightarrow H^1(I,U)
 \longrightarrow H^1(I,U)\otimes_{\mathbb Z}\mathbb Q,
\end{equation}
which shows that $\Phi$ \underline{is canonically isomorphic to the
torsion subgroup of $H^1(I,U)$}.

\SGApage{I}{\(\approx\)453}
On the other hand, from the exact sequence (\SGARefLink{sga7.sga7i.eix.n11-3-3}{11.3.3}) we deduce, for every
$n=\ell^r$, an isomorphism

\begin{equation}\tag{\SGARefTarget{sga7.sga7i.eix.n11-3-9}{11.3.9}}
H^1(I,T_n)\simeq\operatorname{Hom}(I,T_n)
 \xleftarrow{\sim}\operatorname{Hom}(T,T_n)
 \xleftarrow{\sim}\mathbb Z/n\mathbb Z.
\end{equation}
(Notice that $I$ acts trivially on $T$, and hence on $T_n=T/nT$.) It
follows that, for every finite $I$-$\mathbb Z_\ell$-module $M$, if

\begin{equation}\tag{\SGARefTarget{sga7.sga7i.eix.n11-3-10}{11.3.10}}
M'=\operatorname{Hom}(M,T\otimes D_\ell)
\end{equation}
denotes its ``Cartier dual,'' there is a canonical pairing

\begin{equation}\tag{\SGARefTarget{sga7.sga7i.eix.n11-3-11}{11.3.11}}
H^i(I,M)\otimes H^{1-i}(I,M')
 \longrightarrow H^1(I,T\otimes D_\ell)\simeq D_\ell,
\end{equation}
where the last isomorphism is deduced from the isomorphisms (\SGARefLink{sga7.sga7i.eix.n11-3-9}{11.3.9}) by
passage to the inductive limit over the powers $n$ of $\ell$. It is well
known that these pairings are \underline{perfect dualities} (see, for
example, SGA 5 I 5). Taking $M$ in these pairings to be of the form
$U_n$, so that $M'=U'_n$, and passing to the limit, we obtain a perfect
pairing

\begin{equation}\tag{\SGARefTarget{sga7.sga7i.eix.n11-3-12}{11.3.12}}
H^o(I,U\otimes D_\ell)\times H^1(I,U')\longrightarrow D_\ell,
\end{equation}
where the second factor on the left is a finitely generated
$\mathbb Z_\ell$-module. A further passage to the limit in the pairings
(\SGARefLink{sga7.sga7i.eix.n11-3-12}{11.3.12}), for transition morphisms of the form
$\ell^r\operatorname{id}$, gives a second perfect pairing

\begin{equation}\tag{\SGARefTarget{sga7.sga7i.eix.n11-3-13}{11.3.13}}
H^o(I,U\otimes\mathbb Q_\ell)
 \times H^1(I,U'\otimes\mathbb Q_\ell)
 \longrightarrow D_\ell
\end{equation}
of finite-dimensional vector spaces over $\mathbb Q_\ell$. Let $i'$ and
$j'$ denote the morphisms defined as in (\SGARefLink{sga7.sga7i.eix.n11-3-6}{11.3.6}), with $U$ replaced by
$U'$. It is then easy to see that the following homomorphisms

\SGApage{I}{\(\approx\)454}
\begin{center}
\begin{tikzcd}[column sep=6em, row sep=5em]
H^o(I,U\otimes\mathbb Q_\ell) \arrow[r, "j^o"]
  \& H^o(I,U\otimes D_\ell) \\
H^1(I,U'\otimes\mathbb Q_\ell)
  \& H^1(I,U') \arrow[l, "{i'}^1"']
\end{tikzcd}
\end{center}
are \underline{transposes} of one another under the pairings (\SGARefLink{sga7.sga7i.eix.n11-3-12}{11.3.12})
and (\SGARefLink{sga7.sga7i.eix.n11-3-13}{11.3.13}). It follows immediately that
$\operatorname{coker}j^o$ and $\ker {i'}^1$ are dual to one another by a
canonical pairing with values in $D_\ell$; this is the announced pairing
(\SGARefLink{sga7.sga7i.eix.n11-3-5}{11.3.5}).



\medskip
\noindent
\SGARefTarget{sga7.sga7i.eix.n11-4}{11.4}. The preceding description of a pairing (\SGARefLink{sga7.sga7i.eix.n11-3-1}{11.3.1}), by means of the
expression (\SGARefLink{sga7.sga7i.eix.n11-2-1}{11.2.1}) for the $\ell$-primary component of the group of
connected components, is valid only when $\ell\ne p$. Suppose now that
$A_K$ has \underline{semistable reduction} over $S$. We shall again
define pairings that are perfect dualities

\SGAvar{\begin{equation}\tag{\SGARefTarget{sga7.sga7i.eix.n11-4-1}{11.4.1}}
\Phi_o(\ell)\otimes\Phi'_o(\ell)\longrightarrow\mathbb Q/\mathbb Z
\end{equation}
without any restriction on $\ell$. As in \SGARefLink{sga7.sga7i.eix.n11-3}{11.4}, one should establish that this recovers (up to a sign, perhaps,
which ought to be specified) the pairing defined in (\SGARefLink{sga7.sga7i.eix.n1-2-1}{1.2.1}). This would
therefore establish Conjecture \SGARefLink{sga7.sga7i.eix.n1-3}{1.3} in the case of semistable reduction,
just as it does when one restricts to the $\ell$-primary components for
$\ell\ne p$.

}{\begin{equation}\tag{\SGARefTarget{sga7.sga7i.eix.n11-4-1}{11.4.1}}
\Phi_o(\ell)\otimes\Phi'_o(\ell)\longrightarrow\mathbb Q/\mathbb Z
\end{equation}
without any restriction on $\ell$. As in \SGARefLink{sga7.sga7i.eix.n11-3}{11.3}, one should establish that this recovers (up to a sign, perhaps,
which ought to be specified) the pairing defined in (\SGARefLink{sga7.sga7i.eix.n1-2-1}{1.2.1}). This would
therefore establish Conjecture \SGARefLink{sga7.sga7i.eix.n1-3}{1.3} in the case of semistable reduction,
just as it does when one restricts to the $\ell$-primary components for
$\ell\ne p$.

}{E7I-IX-077}To define (\SGARefLink{sga7.sga7i.eix.n11-4-1}{11.4.1}), we shall give, essentially, an expression for
$\Phi_o(\ell)$ in terms of the Barsotti--Tate group $T_\ell(A_K)$ over
$K$; in the present case this will play the same role as \SGARefLink{sga7.sga7i.eix.n11-2}{11.2} did in
the case treated above.

\SGApage{I}{455}
\medskip
\noindent
\underline{Theorem} \SGARefTarget{sga7.sga7i.eix.n11-5}{11.5}. \underline{Suppose that $S$ is henselian and
that $A_K$ has semistable reduction over $S$. Let $\ell$ be a prime and
consider the monodromy pairing} (\SGARefLink{sga7.sga7i.eix.n9-1-2}{9.1.2}), \underline{and hence the
canonical homomorphism}

\SGAvar{\begin{equation*}
\widetilde u_\ell;\underline M_\ell
 \longrightarrow\underline{\check M}'_\ell.
\end{equation*}
\underline{Put $\Phi(\ell)=\operatorname{coker}\widetilde u_\ell$.
Then there is a canonical isomorphism}

}{\begin{equation*}
\widetilde u_\ell:\underline M_\ell
 \longrightarrow\underline{\check M}'_\ell.
\end{equation*}
\underline{Put $\Phi(\ell)=\operatorname{coker}\widetilde u_\ell$.
Then there is a canonical isomorphism}

}{E7I-IX-078}\begin{equation}\tag{\SGARefTarget{sga7.sga7i.eix.n11-5-1}{11.5.1}}
\Phi_o(\ell)\simeq\Phi(\ell)\times_S s.
\end{equation}

This statement immediately gives the definition of a perfect pairing
(\SGARefLink{sga7.sga7i.eix.n11-4-1}{11.4.1}). Indeed, applying \SGARefLink{sga7.sga7i.eix.n11-5}{11.5} to $A'_K$ in place of $A_K$ and using
(\SGARefLink{sga7.sga7i.eix.n10-2-2}{10.2.2}), we also obtain a canonical isomorphism

\begin{equation*}
\Phi'_o(\ell)\simeq\Phi'(\ell)\times_S s,
\qquad
\Phi'(\ell)=\operatorname{coker}\bigl(
 {}^t\!\widetilde u_\ell:\underline M'_\ell
 \longrightarrow\underline{\check M}_\ell\bigr).
\end{equation*}

\medskip
\noindent
\underline{Remarks} \SGARefTarget{sga7.sga7i.eix.n11-5-2}{11.5.2}. a) Theorem 11.5 implies in particular that
$\widetilde u_\ell$ is an isogeny, that is, that $u_\ell$ is
separating, including in the case $\ell=p$, where this assertion is not
immediate from the definition in \SGARefLink{sga7.sga7i.eix.n9-6}{9.6}. Notice that we do not use here
the (considerably deeper) Theorem \SGARefLink{sga7.sga7i.eix.n10-4}{10.4}, which will be proved in the next
section.

\smallskip
\noindent
b) Using \SGARefLink{sga7.sga7i.eix.n10-4}{10.4}, which supplies us with a monodromy pairing

\begin{equation*}
u:\underline M\otimes\underline M'\longrightarrow\mathbb Z_S,
\end{equation*}
and denoting by $\Phi$ the cokernel of the corresponding homomorphism

\begin{equation*}
\widetilde u:\underline M\longrightarrow\underline{\check M}',
\end{equation*}
we obtain an equivalent formulation of \SGARefLink{sga7.sga7i.eix.n11-5}{11.5}, independent of the choice
of $\ell$, in the canonical isomorphism

\SGApage{I}{\(\approx\)456}
\begin{equation}\tag{\SGARefTarget{sga7.sga7i.eix.n11-5-3}{11.5.3}}
\Phi_o\simeq\Phi\times_S s.
\end{equation}



\medskip
\noindent
\SGARefTarget{sga7.sga7i.eix.n11-6}{11.6}. To define the isomorphism (\SGARefLink{sga7.sga7i.eix.n11-5-1}{11.5.1}), we may restrict ourselves to
the case in which $S$ is strictly local (\SGARefLink{sga7.sga7i.eix.n11-0}{11.0}), so that $\underline M$
and $\underline M'$ are constant, with values

\begin{equation}\tag{\SGARefTarget{sga7.sga7i.eix.n11-6-1}{11.6.1}}
M=D(T'_o),\qquad M'=D(T_o),
\end{equation}
where $T_o$ and $T'_o$ are respectively the maximal tori of $A_o$ and
$A'_o$. Likewise identifying $\Phi_o$ with an ordinary finite group,
we must define a canonical isomorphism between its component
$\Phi_o(\ell)$ and

\begin{equation*}
\Phi(\ell)=\operatorname{coker}\bigl(
 \widetilde u_\ell:M_\ell\longrightarrow\check M'_\ell\bigr).
\end{equation*}
In fact, we shall directly define a canonical isomorphism

\begin{equation}\tag{\SGARefTarget{sga7.sga7i.eix.n11-6-2}{11.6.2}}
\Phi_o(\ell)\simeq
 \ker\bigl(\widetilde u_\ell\otimes D_\ell:
 M_\ell\otimes D_\ell\longrightarrow
 \check M'_\ell\otimes D_\ell\bigr).
\end{equation}
Since $\Phi_o(\ell)$ is finite and $M_\ell$ and $\check M'_\ell$ have
the same rank, it will first follow that
$\widetilde u_\ell:M_\ell\to\check M'_\ell$ is an isogeny. It then
follows that the right-hand side of (\SGARefLink{sga7.sga7i.eix.n11-6-2}{11.6.2}) is canonically isomorphic
to $\Phi(\ell)$:

\begin{equation}\tag{\SGARefTarget{sga7.sga7i.eix.n11-6-3}{11.6.3}}
\Phi(\ell)=\operatorname{coker}\widetilde u_\ell
 \simeq
 \ker(\widetilde u_\ell\otimes D_\ell).
\end{equation}
Indeed, identifying $M_\ell$ by means of $\widetilde u_\ell$ with a
finite-index submodule of $\check M'_\ell$, we may write the source and
target of $\widetilde u_\ell\otimes D_\ell$ as

\SGAvar{\begin{equation*}
\Phi(\ell)=V/M_\ell,\qquad
\Phi'(\ell)=V/\check M'_\ell,
\end{equation*}
where

}{\begin{equation*}
M_\ell\otimes D_\ell=V/M_\ell,\qquad
\check M'_\ell\otimes D_\ell=V/\check M'_\ell,
\end{equation*}
where

}{E7I-IX-079}\begin{equation*}
V=M_\ell\otimes\mathbb Q_\ell
 \simeq\check M'_\ell\otimes\mathbb Q_\ell.
\end{equation*}
Composing (\SGARefLink{sga7.sga7i.eix.n11-6-2}{11.6.2}) with (\SGARefLink{sga7.sga7i.eix.n11-6-3}{11.6.3}) now immediately gives the isomorphism
(\SGARefLink{sga7.sga7i.eix.n11-5-1}{11.5.1}) to be defined, so everything reduces to defining the canonical
isomorphism (\SGARefLink{sga7.sga7i.eix.n11-6-2}{11.6.2}).

\SGApage{I}{\(\approx\)457}
\SGAvar{Writing $M_\ell\otimes D_\ell$ (respectively,
$\check M_\ell\otimes D_\ell$) as the inductive limit of the $M_n$
(respectively, the $\check M'_n$), where $n$ ranges over the powers of
$\ell$, it is equivalent to define canonical isomorphisms

}{Writing $M_\ell\otimes D_\ell$ (respectively,
$\check M'_\ell\otimes D_\ell$) as the inductive limit of the $M_n$
(respectively, the $\check M'_n$), where $n$ ranges over the powers of
$\ell$, it is equivalent to define canonical isomorphisms

}{E7I-IX-080}\begin{equation}\tag{\SGARefTarget{sga7.sga7i.eix.n11-6-4}{11.6.4}}
{}_n\Phi_o\xrightarrow{\ \sim\ }
 \ker\bigl(\widetilde u_n:M_n\longrightarrow\check M'_n\bigr),
 \qquad n=\ell^{r+1},
\end{equation}
compatible with the evident transition morphisms as $n$ varies. Here,
by an abuse of notation, $\widetilde u_n$ now denotes the homomorphism
obtained by reducing
$\widetilde u_\ell:M_\ell\to\check M'_\ell$ modulo $n$.

For this purpose, consider the group scheme ${}_nA$ over $S$. By the
semistable-reduction hypothesis it is flat and quasi-finite over $S$
(\SGARefLink{sga7.sga7i.eix.n2-2-1}{2.2.1}). There is a canonical filtration

\begin{equation}\tag{\SGARefTarget{sga7.sga7i.eix.n11-6-5}{11.6.5}}
{}_nA\supset({}_nA^o)^f\supset({}_nA^o)^t\supset0,
\end{equation}
where the superscript $f$ (respectively, $t$) denotes the ``fixed part''
(\SGARefLink{sga7.sga7i.eix.n2-2-3}{2.2.3}) (respectively, the ``toric part'' (\SGARefLink{sga7.sga7i.eix.n5-1}{5.1})) of the group scheme to
which it is applied. In the definition of the toric part in \SGARefLink{sga7.sga7i.eix.n5-1}{5.1}, we had
assumed $S$ complete; this is permissible here (\SGARefLink{sga7.sga7i.eix.n11-0}{11.0}). (We saw moreover
in \SGARefLink{sga7.sga7i.eix.n6-1}{6.1} that the construction generalizes to the henselian case, if one
insists.) Put

\begin{equation}\tag{\SGARefTarget{sga7.sga7i.eix.n11-6-6}{11.6.6}}
{}_n\Psi={}_nA/({}_nA^o)^f.
\end{equation}
This is a flat, separated, quasi-finite group scheme over $S$
(SGA 3 $\mathrm{VI}_A$ 3.2 and $\mathrm{VI}_B$ \SGARefLink{sga7.sga7i.eix.n9-2}{9.2}), with generic fibre

\begin{equation*}
{}_n\Psi_K=({}_nA_K)/({}_nA_K)^f,
\end{equation*}
where, by definition,
$({}_nA_K)^f=({}_nA^o)^f_K$ denotes the $r$-th constituent
$(n=\ell^{r+1})$ of the Barsotti--Tate pro-group
$T_\ell(A_K)^f$. The first of relations (\SGARefLink{sga7.sga7i.eix.n5-5-8}{5.5.8}) therefore gives a
canonical isomorphism

\SGApage{I}{\(\approx\)458}
\begin{equation}\tag{\SGARefTarget{sga7.sga7i.eix.n11-6-7}{11.6.7}}
{}_n\Psi_K\xrightarrow{\ \sim\ }\underline M_{nK},
\qquad
\underline M_n\mathrel{\mathop:}=\underline M
 \otimes_{\mathbb Z}(\mathbb Z/n\mathbb Z).
\end{equation}
In fact, we have:

\medskip
\noindent
\underline{Lemma} \SGARefTarget{sga7.sga7i.eix.n11-6-8}{11.6.8}. \underline{The group scheme ${}_n\Psi$ over
$S$ defined in (\SGARefLink{sga7.sga7i.eix.n11-6-6}{11.6.6}) is \emph{\'etale} over $S$, and the isomorphism
(\SGARefLink{sga7.sga7i.eix.n11-6-7}{11.6.7}) extends uniquely to a morphism}

\begin{equation}\tag{\SGARefTarget{sga7.sga7i.eix.n11-6-8-1}{11.6.8.1}}
{}_n\Psi\lhook\joinrel\longrightarrow\underline M_n
\end{equation}
\underline{which is an open immersion.}

To see that ${}_n\Psi$ is \'etale, in view of (\SGARefLink{sga7.sga7i.eix.n11-6-7}{11.6.7}) it remains only
to verify this at the points of its special fibre, or equivalently to
verify that this fibre is \'etale. This is immediate, since the fibre in
question is

\begin{equation*}
{}_nA_o/{}_nA_o^o\subset A_o/A_o^o=\Phi_o.
\end{equation*}
It follows that ${}_n\Psi$ is normal; and since $\underline M_n$ is
finite over $S$, the existence of the extension (\SGARefLink{sga7.sga7i.eix.n11-6-8-1}{11.6.8.1}) follows at
once. This is a morphism of \'etale group schemes over $S$, so its
kernel is \'etale over $S$. Since ${}_n\Psi$, and hence this kernel, is
separated, while its generic fibre is the unit group by (\SGARefLink{sga7.sga7i.eix.n11-6-7}{11.6.7}), the
kernel is the unit group (\SGARefLink{sga7.sga7i.eix.n3-1-3}{3.1.3}). Thus (\SGARefLink{sga7.sga7i.eix.n11-6-8-1}{11.6.8.1}) is an \'etale
monomorphism, hence an open immersion (EGA IV 17.9.1), proving the lemma.

Now put

\begin{equation}\tag{\SGARefTarget{sga7.sga7i.eix.n11-6-9}{11.6.9}}
{}_n\Phi=({}_nA)^f/({}_nA^o)^f\subset{}_n\Psi.
\end{equation}
Thus ${}_n\Phi$ is a finite open subscheme of ${}_n\Psi$, since
$({}_nA)^f$ is a finite open subgroup scheme of ${}_nA$. More
precisely, it is immediate that

\SGApage{I}{\(\approx\)459}
\begin{equation}\tag{\SGARefTarget{sga7.sga7i.eix.n11-6-10}{11.6.10}}
{}_n\Phi\xrightarrow{\ \sim\ }({}_n\Psi)^f.
\end{equation}
Composing (\SGARefLink{sga7.sga7i.eix.n11-6-9}{11.6.9}) with (\SGARefLink{sga7.sga7i.eix.n11-6-8-1}{11.6.8.1}), we obtain an open immersion

\begin{equation}\tag{\SGARefTarget{sga7.sga7i.eix.n11-6-11}{11.6.11}}
{}_n\Phi\lhook\joinrel\longrightarrow\underline M_n.
\end{equation}
On the other hand,

\begin{equation*}
{}_n\Phi\times_Ss={}_n\Psi\times_Ss={}_nA_o/{}_nA_o^o,
\end{equation*}
and, since $A_o^o$ is $n$-divisible (\SGARefLink{sga7.sga7i.eix.n2-2-1}{2.2.1}), the last member is
immediately identified with ${}_n\Phi_o$. We therefore obtain an
isomorphism

\begin{equation}\tag{\SGARefTarget{sga7.sga7i.eix.n11-6-12}{11.6.12}}
{}_n\Phi\times_Ss\xrightarrow{\ \sim\ }{}_n\Phi_o,
\end{equation}
which shows that ${}_n\Phi$ is characterized, up to unique isomorphism,
as the finite \'etale scheme over $S$ whose special fibre is
${}_n\Phi_o$. Thus, to define the desired isomorphism (\SGARefLink{sga7.sga7i.eix.n11-6-4}{11.6.4}) by means
of the open immersion (\SGARefLink{sga7.sga7i.eix.n11-6-11}{11.6.11}), it is enough to prove the following.

\medskip
\noindent
\underline{Lemma} \SGARefTarget{sga7.sga7i.eix.n11-6-13}{11.6.13}. \underline{The image of the canonical
homomorphism (\SGARefLink{sga7.sga7i.eix.n11-6-11}{11.6.11}) is equal to the kernel of}

\begin{equation*}
\widetilde u_n:\underline M_n\longrightarrow\underline{\check M}'_n.
\end{equation*}

\SGApage{I}{455}
The required compatibility of the homomorphisms (\SGARefLink{sga7.sga7i.eix.n11-6-11}{11.6.11}) with the
transition homomorphisms as $n$ varies is immediate from the
definitions. Consequently Lemma \SGARefLink{sga7.sga7i.eix.n11-6-13}{11.6.13} will complete the proof of
Theorem \SGARefLink{sga7.sga7i.eix.n11-5}{11.5}.



\SGApage{I}{\(\approx\)460}
To prove the lemma, note that

\begin{equation*}
K_n=\ker\widetilde u_n
\end{equation*}
is also the left annihilator of the canonical pairing

\begin{equation*}
u_n:\underline M_n\times\underline M'_n
 \longrightarrow(\mathbb Z/n\mathbb Z)_S,
\end{equation*}
whose explicit definition is given in \SGARefLink{sga7.sga7i.eix.n9-4}{9.4} and \SGARefLink{sga7.sga7i.eix.n9-6}{9.6} in terms of the
extensions

\SGAvar{\begin{equation}\tag{\SGARefTarget{sga7.sga7i.eix.n11-6-14}{11.6.14}}
\resizebox{\textwidth}{!}{
\begin{tikzcd}[ampersand replacement=\&, column sep=small, row sep=large]
0 \arrow[r]
 \& {\ensuremath{\check{\underline M}'_n(1)}} \arrow[r]
 \& F_n \arrow[r] \arrow[d, equal]
 \& R_n \arrow[r]
 \& 0
 \& {,}
 \& 0 \arrow[r]
 \& {\ensuremath{\check{\underline M}_n(1)}} \arrow[r]
 \& F'_n \arrow[r] \arrow[d, equal]
 \& R'_n \arrow[r]
 \& 0 \\
{}
 \& {}
 \& ({}_nA^o)^f
 \& {}
 \& {}
 \& {}
 \& {}
 \& {}
 \& ({}_nA'^o)^f
 \& {}
 \& {}
\end{tikzcd}}
\end{equation}
(where $R_n$ and $R'_n$ are Cartier dual to one another), and in terms
of the mixed extension $X_{n\eta}$ of $E_{n\eta}$ by $F_{n\eta}$,
which is the generic fibre of (\SGARefLink{sga7.sga7i.eix.n11-6-5}{11.7.5}), where we put

\begin{equation}\tag{\SGARefTarget{sga7.sga7i.eix.n11-6-15}{11.6.15}}
\begin{tikzcd}[column sep=large, row sep=large]
E_n=D(F_n),
 \& 0 \arrow[r]
 \& R_n \arrow[r] \arrow[d, no head, "\wr"']
 \& E_n \arrow[r]
 \& \underline M_n \arrow[r]
 \& 0
 \& {.} \\
{}
 \& {}
 \& D(R'_n)
 \& {}
 \& {}
 \& {}
 \& {}
\end{tikzcd}
\end{equation}
It follows immediately from \SGARefLink{sga7.sga7i.eix.n9-4-7}{9.4.6} that $K_n$ is the largest among the finite \'etale subgroups
$K'_n$ of $\underline M_n$ for which the inverse image
$\overline X_{n\eta}$ of $K'_n$ in the mixed extension $X_{n\eta}$,
regarded as a mixed extension of $\overline E_{n\eta}$ by $F_{n\eta}$,
extends to a mixed extension of $\overline E_n$ by $F_n$. Here
$\overline E_n$ is the evident subgroup extension of $K_n$ by
$R_n$.

}{\begin{equation}\tag{\SGARefTarget{sga7.sga7i.eix.n11-6-14}{11.6.14}}
\resizebox{\textwidth}{!}{
\begin{tikzcd}[ampersand replacement=\&, column sep=small, row sep=large]
0 \arrow[r]
 \& {\ensuremath{\check{\underline M}'_n(1)}} \arrow[r]
 \& F_n \arrow[r] \arrow[d, equal]
 \& R_n \arrow[r]
 \& 0
 \& {,}
 \& 0 \arrow[r]
 \& {\ensuremath{\check{\underline M}_n(1)}} \arrow[r]
 \& F'_n \arrow[r] \arrow[d, equal]
 \& R'_n \arrow[r]
 \& 0 \\
{}
 \& {}
 \& ({}_nA^o)^f
 \& {}
 \& {}
 \& {}
 \& {}
 \& {}
 \& ({}_nA'^o)^f
 \& {}
 \& {}
\end{tikzcd}}
\end{equation}
(where $R_n$ and $R'_n$ are Cartier dual to one another), and in terms
of the mixed extension $X_{n\eta}$ of $E_{n\eta}$ by $F_{n\eta}$,
which is the generic fibre of (\SGARefLink{sga7.sga7i.eix.n11-6-5}{11.6.5}), where we put

\begin{equation}\tag{\SGARefTarget{sga7.sga7i.eix.n11-6-15}{11.6.15}}
\begin{tikzcd}[column sep=large, row sep=large]
E_n=D(F'_n),
 \& 0 \arrow[r]
 \& R_n \arrow[r] \arrow[d, no head, "\wr"']
 \& E_n \arrow[r]
 \& \underline M_n \arrow[r]
 \& 0
 \& {.} \\
{}
 \& {}
 \& D(R'_n)
 \& {}
 \& {}
 \& {}
 \& {}
\end{tikzcd}
\end{equation}
It follows immediately from Proposition \SGARefLink{sga7.sga7i.eix.n9-4-7}{9.4.7} that $K_n$ is the largest among the finite \'etale subgroups
$K'_n$ of $\underline M_n$ for which the inverse image
$\overline X_{n\eta}$ of $K'_n$ in the mixed extension $X_{n\eta}$,
regarded as a mixed extension of $\overline E_{n\eta}$ by $F_{n\eta}$,
extends to a mixed extension of $\overline E_n$ by $F_n$. Here
$\overline E_n$ is the evident subgroup extension of $K'_n$ by
$R_n$.

}{E7I-IX-081}This condition is plainly satisfied for the subgroup ${}_n\Phi$ of
$\underline M_n$ in (\SGARefLink{sga7.sga7i.eix.n11-6-11}{11.6.11}), since the corresponding mixed extension
$({}_nA)^f$ over $S$ is available. Hence ${}_n\Phi\subset K_n$. To
prove the opposite inclusion, consider the mixed extension
$\overline X_n$ of $\overline E_n$ by $F_n$ extending
$\overline X_{n\eta}$, defined as above by taking $K'_n=K_n$. Thus it
is an extension

\SGApage{I}{\(\approx\)461}
\begin{equation*}
0\longrightarrow F_n\longrightarrow\overline X_n
 \longrightarrow K_n\longrightarrow0,
\end{equation*}
and there is a canonical homomorphism

\begin{equation}\tag{*}
F_n=({}_nA^o)^f\lhook\joinrel\longrightarrow A,
\end{equation}
which on the generic fibre is induced by the canonical homomorphism

\begin{equation}\tag{**}
\overline X_{n\eta}\lhook\joinrel\longrightarrow X_{n\eta}
 ={}_nA_K\lhook\joinrel\longrightarrow A_K.
\end{equation}
It follows immediately from (\SGARefLink{sga7.sga7i.eix.n5-9-2}{5.9.2}) that the homomorphism $(**)$ extends
to a unique homomorphism

\begin{equation}\tag{***}
\overline X_n\longrightarrow A
\end{equation}
inducing $(*)$. Since this last homomorphism maps the generic fibre into
${}_nA$, which is a closed subscheme of $A$, it factors through
$\overline X_n\to{}_nA$; and since $\overline X_n$ is finite, it factors
through $({}_nA)^f$. Looking at the induced homomorphism on generic
fibres, we conclude that

\begin{equation*}
\overline X_{n\eta}\subset({}_nA)^f_\eta,
\quad\text{i.e.}\quad
K_{n\eta}\subset{}_n\Phi_\eta,
\quad\text{i.e.}\quad
K_n\subset{}_n\Phi.
\end{equation*}
This proves the lemma.

\SGAvar{\medskip
\noindent
\underline{Remarks} \SGARefTarget{sga7.sga7i.eix.n11-7}{11.7}. The construction made in \SGARefLink{sga7.sga7i.eix.n11-6}{11.7}
shows that, from the extensions (\SGARefLink{sga7.sga7i.eix.n11-6-14}{11.6.14}) over $S$ and the mixed
extension $X_{n\eta}={}_nA_K$ of $E_n$ by $F_n$, one can reconstruct,
up to unique isomorphism, the mixed extension (\SGARefLink{sga7.sga7i.eix.n11-6-5}{11.6.5}), and hence also
the subgroup

}{\medskip
\noindent
\underline{Remarks} \SGARefTarget{sga7.sga7i.eix.n11-7}{11.7}. The construction made in \SGARefLink{sga7.sga7i.eix.n11-6}{11.6}
shows that, from the extensions (\SGARefLink{sga7.sga7i.eix.n11-6-14}{11.6.14}) over $S$ and the mixed
extension $X_{n\eta}={}_nA_K$ of $E_n$ by $F_n$, one can reconstruct,
up to unique isomorphism, the mixed extension (\SGARefLink{sga7.sga7i.eix.n11-6-5}{11.6.5}), and hence also
the subgroup

}{E7I-IX-082}\begin{equation*}
({}_nA)^f\supset({}_nA^o)^f
\end{equation*}
of ${}_nA$. Indeed, by \SGARefLink{sga7.sga7i.eix.n9-4}{9.4} these data determine the pairing

\SGAvar{\begin{equation*}
u_n:\underline M_n\times\underline M'_n
 \longrightarrow(\mathbb Z/n\mathbb Z)_S,
\end{equation*}
and the preceding construction gives $({}_nA)^f$ as the mixed extension
$\overline X_n$ of $\overline E_n$ by $F_n$, extending the mixed
extension $\overline X_{n\eta}$ induced by $X_{n\eta}$. Here
$\overline E_n$ is the inverse image of $\ker\widetilde u_n$ under
$E_n\to\underline M_n$. By (9.4.8), this mixed extension
$\overline X_n$ is thus determined \underline{up to unique
isomorphism}. We then recover ${}_nA$ by gluing $\overline X_n$ and
$X_{n\eta}$ along their common open subscheme $\overline X_{n\eta}$.

}{\begin{equation*}
u_n:\underline M_n\times\underline M'_n
 \longrightarrow(\mathbb Z/n\mathbb Z)_S,
\end{equation*}
and the preceding construction gives $({}_nA)^f$ as the mixed extension
$\overline X_n$ of $\overline E_n$ by $F_n$, extending the mixed
extension $\overline X_{n\eta}$ induced by $X_{n\eta}$. Here
$\overline E_n$ is the inverse image of $\ker\widetilde u_n$ under
$E_n\to\underline M_n$. By (9.4.7), this mixed extension
$\overline X_n$ is thus determined \underline{up to unique
isomorphism}. We then recover ${}_nA$ by gluing $\overline X_n$ and
$X_{n\eta}$ along their common open subscheme $\overline X_{n\eta}$.

}{E7I-IX-083}\SGApage{I}{\(\approx\)462}
Thus, if the structural data (\SGARefLink{sga7.sga7i.eix.n9-6-1}{9.6.1}) and the mixed extension

\begin{equation*}
X_\eta=T_\ell(A_K)
\end{equation*}
of $E_\eta$ by $F_\eta$ are known---equivalently, if the extensions
(\SGARefLink{sga7.sga7i.eix.n11-6-14}{11.6.14}) and $X_{n\eta}$ are known for every $n=\ell^{r+1}$,
$r\geq0$---then one can reconstruct the inductive (or, if preferred,
projective) system of mixed extensions (\SGARefLink{sga7.sga7i.eix.n11-6-5}{11.6.5}).



\SGARefTarget{sga7.sga7i.eix.n11-8}{11.8}. As in \SGARefLink{sga7.sga7i.eix.n8-3}{8.3}, regard the separable closure $\overline K$ of $K$ as the
inductive limit of its finite subextensions $K'$. For every such $K'$, let
$A[K']$ denote the N\'{e}ron model of

\begin{equation*}
A_{K'}=A_K\otimes_K K'.
\end{equation*}
It gives a finite \'{e}tale group scheme

\SGAvar{\begin{equation}\tag{\SGARefTarget{sga7.sga7i.eix.n11-8-1}{11.8.1}}
\Phi_o[K']=A[K']_o,A[K']_o^o
\end{equation}
of connected components, defined over the residue field of the normalization
$S[K']$ of $S$ in $K'$. For simplicity, suppose that the residue field of
$S$ is separably closed, so that the same is true of the residue fields of
the $S[K']$. The groups (\SGARefLink{sga7.sga7i.eix.n11-8-1}{11.8.1}) may then be identified with ordinary finite
groups $\Phi[K']$, and it is immediate that, as $K'$ varies, these form an
inductive system, which we now propose to study. We shall be interested in
particular in the group

}{\begin{equation}\tag{\SGARefTarget{sga7.sga7i.eix.n11-8-1}{11.8.1}}
\Phi_o[K']=A[K']_o/A[K']_o^o
\end{equation}
of connected components, defined over the residue field of the normalization
$S[K']$ of $S$ in $K'$. For simplicity, suppose that the residue field of
$S$ is separably closed, so that the same is true of the residue fields of
the $S[K']$. The groups (\SGARefLink{sga7.sga7i.eix.n11-8-1}{11.8.1}) may then be identified with ordinary finite
groups $\Phi[K']$, and it is immediate that, as $K'$ varies, these form an
inductive system, which we now propose to study. We shall be interested in
particular in the group

}{E7I-IX-084}\SGApage{I}{\(\approx\)463}
\begin{equation}\tag{\SGARefTarget{sga7.sga7i.eix.n11-8-2}{11.8.2}}
\Psi=\varinjlim_{K'}\Phi[K'].
\end{equation}
We know from \SGARefLink{sga7.sga7i.eix.n3-6}{3.6} that, for $K'$ sufficiently large, say
$K'\supset K'_o$, $A[K']$ has semistable reduction. Hence for
$K'\supset K'_o$ the transition morphisms are injective (\SGARefLink{sga7.sga7i.eix.n3-3}{3.3}). Moreover,
with the notation of \SGARefLink{sga7.sga7i.eix.n5-5}{5.5}, retained in the present section, there is for
$K'\supset K'_o$ a canonical isomorphism

\SGAvar{\begin{equation}\tag{\SGARefTarget{sga7.sga7i.eix.n11-8-3}{11.8.3}}
\Psi[K'](\ell)\simeq
\ker\bigl(\widetilde u[K']_\ell\otimes D_\ell:
 M_\ell\otimes D_\ell\longrightarrow
 \check M'_\ell\otimes D_\ell\bigr),
\end{equation}
coming from (\SGARefLink{sga7.sga7i.eix.n11-6-2}{11.6.2}) applied to $A[K']$. One also verifies immediately that
the corresponding immersions

}{\begin{equation}\tag{\SGARefTarget{sga7.sga7i.eix.n11-8-3}{11.8.3}}
\Phi[K'](\ell)\simeq
\ker\bigl(\widetilde u[K']_\ell\otimes D_\ell:
 M_\ell\otimes D_\ell\longrightarrow
 \check M'_\ell\otimes D_\ell\bigr),
\end{equation}
coming from (\SGARefLink{sga7.sga7i.eix.n11-6-2}{11.6.2}) applied to $A[K']$. One also verifies immediately that
the corresponding immersions

}{E7I-IX-085}\begin{equation*}
\Phi[K'](\ell)\lhook\joinrel\longrightarrow M_\ell\otimes D_\ell
\end{equation*}
are compatible with the transition morphisms as $K'$ varies. Passing to the
limit therefore gives a canonical monomorphism

\begin{equation}\tag{\SGARefTarget{sga7.sga7i.eix.n11-8-4}{11.8.4}}
\Psi(\ell)\mathrel{\mathop{=}^{\mathrm{def}}}
 \varinjlim_{K'}\Phi[K'](\ell)
 \lhook\joinrel\longrightarrow M_\ell\otimes D_\ell
\end{equation}
from the $\ell$-primary component of $\Psi$; summing over all $\ell$ gives

\begin{equation}\tag{\SGARefTarget{sga7.sga7i.eix.n11-8-5}{11.8.5}}
\Psi\mathrel{\mathop{=}^{\mathrm{def}}}
 \varinjlim_{K'}\Phi[K']
 \longrightarrow M\otimes\mathbb Q/\mathbb Z.
\end{equation}

\medskip
\noindent
\underline{Theorem} \SGARefTarget{sga7.sga7i.eix.n11-9}{11.9}. \underline{The canonical homomorphism (\SGARefLink{sga7.sga7i.eix.n11-8-5}{11.8.5})
is an isomorphism}.

We already know that it is a monomorphism. To prove that it is an
isomorphism, we must use the characterization (\SGARefLink{sga7.sga7i.eix.n11-8-3}{11.8.3}) of the image of
$\Phi[K'](\ell)$ and the behaviour of $\widetilde u[K']$ as $K'$ varies.
That behaviour is given by (\SGARefLink{sga7.sga7i.eix.n10-3-4}{10.3.4}).

\SGApage{I}{\(\approx\)464}
\SGAvar{We may reduce to the case in which $S$ is strictly local. To prove \SGARefLink{sga7.sga7i.eix.n11-9}{11.9} we
are then reduced to the following familiar assertion, applied to $S[K'_o]$:
if $S$ is a strictly local trait with fraction field $K$, then for every
integer $n\geq1$ there is a separable extension $K'$ of $K$ whose degree is
divisible by $n$ (cf. I\hspace{3em}).


}{We may reduce to the case in which $S$ is strictly local. To prove \SGARefLink{sga7.sga7i.eix.n11-9}{11.9} we
are then reduced to the following familiar assertion, applied to $S[K'_o]$:
if $S$ is a strictly local trait with fraction field $K$, then for every
integer $n\geq1$ there is a separable extension $K'$ of $K$ whose degree is
divisible by $n$.


}{E7I-IX-086}\medskip
\noindent
\underline{Corollary} \SGARefTarget{sga7.sga7i.eix.n11-10}{11.10}. \underline{Under the hypotheses ($S$
complete) and with the notation of \SGARefLink{sga7.sga7i.eix.n8-3}{8.3}, the ind-finite group scheme over
$K$ of connected components of the group scheme $A_K^{\flat}$,
namely $A_K^{\flat}/{A_K^{\flat}}^o$, is canonically isomorphic to}
$\underline M_K\otimes\mathbb Q/\mathbb Z$, \underline{where
$\underline M_K$ is the twisted constant group over $K$ defined in
(\SGARefLink{sga7.sga7i.eix.n8-1-8}{8.1.8})}.

\SGAvar{This follows immediately from (8.2.4) and \SGARefLink{sga7.sga7i.eix.n11-9}{11.9}, in view of the definition
of $\underline M_K$.

}{This follows immediately from (8.3.4) and \SGARefLink{sga7.sga7i.eix.n11-9}{11.9}, in view of the definition
of $\underline M_K$.

}{E7I-IX-087}\medskip
\noindent
\underline{Remarks} \SGARefTarget{sga7.sga7i.eix.n11-11}{11.11}. a) If one restricts to the $\ell$-primary
components for $\ell\neq p$, Theorem \SGARefLink{sga7.sga7i.eix.n11-9}{11.9} is immediate in terms of \SGARefLink{sga7.sga7i.eix.n11-2}{11.2} and
does not require the more delicate considerations of \SGARefLink{sga7.sga7i.eix.n11-6}{11.6} (via \SGARefLink{sga7.sga7i.eix.n9-4}{9.4}).
Likewise, if $K$ is assumed to be \underline{of characteristic zero}, \SGARefLink{sga7.sga7i.eix.n11-9}{11.9}
can be proved without using (\SGARefLink{sga7.sga7i.eix.n11-8-3}{11.8.3}): it is immediate that the image of
${}_n\Phi[K'](\ell)$ in

\begin{equation*}
M_n={}_n(M_\ell\otimes D_\ell)
\end{equation*}
contains the image of the group ${}_nA(K')$ of $K'$-valued points of
${}_nA$. For $K'$ sufficiently large this group is equal to
${}_nA(\overline K)$, and therefore maps onto $M_n$. If
$\operatorname{char}K=p$, on the other hand, it appears that there is no
``trivial'' proof---that is, no proof avoiding \SGARefLink{sga7.sga7i.eix.n9-4}{9.4}---of \SGARefLink{sga7.sga7i.eix.n11-10}{11.10} for the
$p$-primary component.

b) The existence of an isomorphism

\begin{equation*}
\Psi\simeq(\mathbb Q/\mathbb Z)^r,
\end{equation*}
together with the correct value of $r$, was conjectured in 1964 by
J.-P. Serre.


\SGApage{I}{\(\approx\)465}
\medskip
\noindent
12. \underline{N\'eron model of a Jacobian and the Picard--Lefschetz
formula}

Let us first recall the following general results of M. Raynaud
[22, Ths.~5 and 7; 22 bis].

\medskip
\noindent
\underline{Theorem} \SGARefTarget{sga7.sga7i.eix.n12-1}{12.1} (M. Raynaud). \underline{Let $S$ be a strictly
local trait, and let $f:X\to S$ be a scheme proper and flat over $S$, with
one-dimensional fibres, such that}

\begin{equation}\tag{\SGARefTarget{sga7.sga7i.eix.n12-1-1}{12.1.1}}
 f_*\underline{\mathcal O}_X\simeq\underline{\mathcal O}_S,
\end{equation}

\noindent
\underline{and such that the local rings of $X$ are factorial. For every
maximal point $x_i$ of $X_o$, put}

\SGAvar{\begin{equation}\tag{\SGARefTarget{sga7.sga7i.eix.n12-1-2}{12.1.2}}
 d_i=\operatorname{length}\underline{\mathcal O}_{X_o,x},
 \qquad d_i^t=\mu_i d_i,
\end{equation}

}{\begin{equation}\tag{\SGARefTarget{sga7.sga7i.eix.n12-1-2}{12.1.2}}
 d_i=\operatorname{length}\underline{\mathcal O}_{X_o,x_i},
 \qquad d_i^t=\mu_i d_i,
\end{equation}

}{E7I-IX-088}\noindent
\underline{where $\mu_i$ is the radicial multiplicity of $k(x_i)$ over
$k$} (EGA IV, 4.7.4) (\underline{thus a power of the characteristic
exponent of $k$}). \underline{Let $d$ (respectively $d^t$) be the gcd of the integers
$d_i$ (respectively $d_i^t$); hence $d=d^t$ if $k$ is perfect. Put}

\begin{equation}\tag{\SGARefTarget{sga7.sga7i.eix.n12-1-3}{12.1.3}}
 P=\underline{\operatorname{Pic}}_{X/S}:
   (\operatorname{Sch})_{/S}^{\,o}\longrightarrow(\operatorname{Ab}).
\end{equation}

\noindent
a) \underline{Consider the following conditions}:

\smallskip
\noindent
1) $d^t=1$.

\smallskip
\noindent
2) $X$ \underline{has a zero-cycle of degree $1$}.

\smallskip
\noindent
3) \underline{There exists an invertible sheaf $L$ on $X\times P^o$ such
that the corresponding homomorphism}
$P^o_\eta\to\underline{\operatorname{Pic}}_{X_\eta/\eta}=P_\eta$
\underline{is the canonical inclusion}. (It is known, by J.-P. Murre,
that $P_\eta$ is representable by a group scheme locally of finite type
over $\eta$; see, for example, SGA 6, XII, 1.5.)

\SGApage{I}{\(\approx\)466}
\smallskip
\noindent
4) $X$ \underline{is cohomologically flat over $S$} (EGA III, 7.8.1)
(equivalently, cohomologically flat in dimension zero, or equivalently
$k\xrightarrow{\sim}H^o(X_o,\underline{\mathcal O}_{X_o})$),
\underline{and $d=1$}.

\smallskip
\noindent
5) $P$ \underline{is representable} (equivalently, representable by a
group scheme smooth over $S$).

\smallskip
\noindent
6) $P^o$ \underline{is representable} (equivalently, representable by a
group scheme smooth over $S$). (For the definition of $P^o$, see SGA 3,
$\mathrm{VI}_B$, \SGARefLink{sga7.sga7i.eix.n3-1}{3.1}.)

\smallskip
\noindent
7) $P^o$ \underline{is separated over $S$}, that is, it satisfies the
valuative criterion of separatedness of EGA II, 7.2.3 (equivalently,
every section $s$ of $P^o$ over $S$ for which $s(\eta)=0$ is zero).

\smallskip
\noindent
8) $d=1$.

\smallskip
\noindent
\underline{Then one has the implications}

\begin{equation}\tag{\SGARefTarget{sga7.sga7i.eix.n12-1-4}{12.1.4}}
 1)\Longleftrightarrow2)\Longrightarrow3)\Longrightarrow4)
 \Longleftrightarrow5)\Longleftrightarrow6)\Longleftrightarrow7)
 \Longrightarrow8).
\end{equation}

\noindent
\underline{In particular, if $d=d^t$, for example if $k$ is perfect, the
conditions 1) through 8) are equivalent}.

\medskip
\noindent
b) \underline{Suppose that $d^t=1$, so that $P$ is representable by a
group scheme smooth over $S$. Then the schematic closure of the identity
section of $P$ is a group subscheme $E$ which is \`etale over $S$. If $I$
is the set of reduced irreducible components $X_o^i$ of $X_o$, one obtains
an isomorphism}

\begin{equation}\tag{\SGARefTarget{sga7.sga7i.eix.n12-1-5}{12.1.5}}
 \mathbb Z^I\simeq\Gamma(S,E)
\end{equation}

\SGApage{I}{\(\approx\)467}
\SGAvar{\noindent
\underline{by associating to every $i\in I$ the section $S_i$ of $P$ over
$S$ (zero over $\eta$) defined by the invertible sheaf
$\underline{\mathcal O}_X(X_o^i)$ on $S$}. (Because the local rings of $X$
are factorial by hypothesis, the cycles $X_o^i$ are indeed divisors; thus
$\underline{\mathcal O}_X(X_i^o)$ has its stated meaning as the invertible
sheaf associated with that divisor; see EGA IV, 21.2.8.1.) \underline{One
has}

}{\noindent
\underline{by associating to every $i\in I$ the section $S_i$ of $P$ over
$S$ (zero over $\eta$) defined by the invertible sheaf
$\underline{\mathcal O}_X(X_o^i)$ on $X$}. (Because the local rings of $X$
are factorial by hypothesis, the cycles $X_o^i$ are indeed divisors; thus
$\underline{\mathcal O}_X(X_o^i)$ has its stated meaning as the invertible
sheaf associated with that divisor; see EGA IV, 21.2.8.1.) \underline{One
has}

}{E7I-IX-089}\begin{equation}\tag{\SGARefTarget{sga7.sga7i.eix.n12-1-6}{12.1.6}}
 P^o\cap E=\underline{\text{the zero section of }}P \quad .
\end{equation}

\noindent
c) \underline{Under the conditions of b), the quotient}

\begin{equation}\tag{\SGARefTarget{sga7.sga7i.eix.n12-1-7}{12.1.7}}
 Q=P/E
\end{equation}

\SGAvar{\noindent
\underline{is representable by a smooth separated group scheme over $S$,
and this group scheme is N\'eron, that is, it satisfies the universal
property (\SGARefLink{sga7.sga7i.eix.n1-1-2}{1.1.2}) in terms of its special fibre}.

}{\noindent
\underline{is representable by a smooth separated group scheme over $S$,
and this group scheme is N\'eron, that is, it satisfies the universal
property (\SGARefLink{sga7.sga7i.eix.n1-1-2}{1.1.2}) in terms of its generic fibre}.

}{E7I-IX-090}\medskip
\noindent
d) \underline{Under the conditions of c), suppose in addition that the
geometric generic fibre of $X$ is integral. Then there is a canonical
exact sequence}

\begin{equation}\tag{\SGARefTarget{sga7.sga7i.eix.n12-1-8}{12.1.8}}
 0\longrightarrow A\longrightarrow Q
   \xrightarrow{\ \deg\ }\mathbb Z_S\longrightarrow0,
\end{equation}

\noindent
\underline{where $\deg$ is induced by the ``degree'' homomorphism
$P\to\mathbb Z_S$, and where $A$ is a smooth separated group scheme of
finite type over $S$, which is likewise N\'eron, that is, a N\'eron model
of}

\begin{equation}\tag{\SGARefTarget{sga7.sga7i.eix.n12-1-9}{12.1.9}}
 A_\eta=P_\eta^o=
 \underline{\operatorname{Pic}}^o_{X_\eta/\eta} \quad .
\end{equation}

\noindent
\underline{Moreover, the composite homomorphism $P^o\to P\to Q$ induces
an isomorphism}

\begin{equation}\tag{\SGARefTarget{sga7.sga7i.eix.n12-1-10}{12.1.10}}
 P^o\xrightarrow{\ \sim\ }A^o.
\end{equation}

\SGApage{I}{468}
\medskip
\noindent
\SGARefTarget{sga7.sga7i.eix.n12-1-11}{12.1.11}. The case of interest to us is that in which $X_\eta$ is smooth
over $\eta$, so that, in view of (\SGARefLink{sga7.sga7i.eix.n12-1-1}{12.1.1}), the geometric generic fibre of
$X$ is smooth and connected (hence integral), and in which $X_\eta$ has a
zero-cycle of degree one. All the conditions considered in \SGARefLink{sga7.sga7i.eix.n12-1}{12.1} (and in
particular conditions 1) through 8) of a)) then hold. Moreover,
$A_\eta=P_\eta^o$ is then an abelian scheme, and d) gives us a very
explicit construction of its N\'eron model $A$. In particular, (\SGARefLink{sga7.sga7i.eix.n12-1-10}{12.1.10})
gives an isomorphism

\begin{equation}\tag{\SGARefTarget{sga7.sga7i.eix.n12-1-12}{12.1.12}}
 A_o^{\,o}\simeq\underline{\operatorname{Pic}}^o_{X_o/k},
\end{equation}

\noindent
which shows in particular that $A$ \underline{has semistable reduction
over $S$ if and only if}
$\underline{\operatorname{Pic}}^o_{X_o/k}$ \underline{has unipotent rank
zero}. When $X$ is regular and $X_o$ is separable over $k$, this is known
to hold if and only if the geometric special fibre $\overline X_o$ has no
singular points other than ordinary double points (that is, points whose
completed local ring is the same as that of the union of the two
coordinate axes of affine $2$-space at the origin). Verification of this
fact, together with the ``well-known'' description of the Picard scheme
of a proper separable curve over a field $k$, using EGA IV, 21.8.5, is left
to the reader as an exercise. Notice also that the implication
``$\overline X_o$ has ordinary double points'' $\Longrightarrow$
``$A_\eta=\underline{\operatorname{Pic}}^o_{X_\eta/\eta}$ has semistable
reduction'' follows immediately from the ``geometric monodromy theorem''
III, in view of the Galois
criterion for semistable reduction in \SGARefLink{sga7.sga7i.eix.n3-5}{3.5}. This argument, however, does
not give the important isomorphism (\SGARefLink{sga7.sga7i.eix.n12-1-10}{12.1.10}) and its consequence
(\SGARefLink{sga7.sga7i.eix.n12-1-12}{12.1.12}).

\SGApage{I}{469}
\SGAvar{\medskip
\noindent
\SGARefTarget{sga7.sga7i.eix.n12-2}{12.2}. For the proof of \SGARefLink{sga7.sga7i.eix.n12-1}{12.1}, we refer to [22 bis]; in fact we shall need
only the case where $X_\eta$ is smooth over $\eta$. Let us also note that
[22 bis] proves relations between Picard schemes and N\'eron models under
conditions substantially more general than those of \SGARefLink{sga7.sga7i.eix.n12-1}{12.1} c) and d). As
their statements are more technical, we have chosen not to include them
in this review. We mention only that b), c), and d) are easy consequences
of a), and that in a) the only delicate points are the implications
$3)\Longrightarrow4)\Longrightarrow5)$ and $7)\Longrightarrow4)$, the
first being the most difficult. This point is trivial, however, if one
assumes that $X_o$ is separable over $k$ (which implies that $X$ has a
section over $S$, since $S$ is strictly local, and hence that condition
1), $d^t=1$, holds). We advise the reader to give the proof in this
particularly simple case, assuming in addition, if necessary, that
$X_\eta$ is smooth, so that $P_\eta^o$ is an abelian scheme. One may then
use its N\'eron model $A$, and $A^o$ represents
$\underline{\operatorname{Pic}}^o_{X/S}$. We shall use \SGARefLink{sga7.sga7i.eix.n12-1}{12.1} only in this
particular case, and under the further assumption that $\overline X_o$
has no singularities other than ordinary double points, that is, by
(\SGARefLink{sga7.sga7i.eix.n12-1-11}{12.1.11}), that $A_\eta$ has semistable reduction. This is the case
considered in 12.5.4, to which we restrict ourselves from now on.


}{\medskip
\noindent
\SGARefTarget{sga7.sga7i.eix.n12-2}{12.2}. For the proof of \SGARefLink{sga7.sga7i.eix.n12-1}{12.1}, we refer to [22 bis]; in fact we shall need
only the case where $X_\eta$ is smooth over $\eta$. Let us also note that
[22 bis] proves relations between Picard schemes and N\'eron models under
conditions substantially more general than those of \SGARefLink{sga7.sga7i.eix.n12-1}{12.1} c) and d). As
their statements are more technical, we have chosen not to include them
in this review. We mention only that b), c), and d) are easy consequences
of a), and that in a) the only delicate points are the implications
$3)\Longrightarrow4)\Longrightarrow5)$ and $7)\Longrightarrow4)$, the
first being the most difficult. This point is trivial, however, if one
assumes that $X_o$ is separable over $k$ (which implies that $X$ has a
section over $S$, since $S$ is strictly local, and hence that condition
1), $d^t=1$, holds). We advise the reader to give the proof in this
particularly simple case, assuming in addition, if necessary, that
$X_\eta$ is smooth, so that $P_\eta^o$ is an abelian scheme. One may then
use its N\'eron model $A$, and $A^o$ represents
$\underline{\operatorname{Pic}}^o_{X/S}$. We shall use \SGARefLink{sga7.sga7i.eix.n12-1}{12.1} only in this
particular case, and under the further assumption that $\overline X_o$
has no singularities other than ordinary double points, that is, by
(\SGARefLink{sga7.sga7i.eix.n12-1-11}{12.1.11}), that $A_\eta$ has semistable reduction. This is the case
considered in 10.5.4, to which we restrict ourselves from now on.


}{E7I-IX-091}\medskip
\noindent
\SGARefTarget{sga7.sga7i.eix.n12-3}{12.3}. Let $C$ be a proper separable curve over a field $k$, which for
simplicity we suppose separably closed (or even algebraically closed, if
the reader prefers). If $\widetilde C$ is the normalization of $C$, it is
known that $\underline{\operatorname{Pic}}^o_{\widetilde C/k}$ has
``reductive rank'' zero (that is, its maximal torus is zero). Consequently
the maximal torus of $\underline{\operatorname{Pic}}^o_{C/k}$ is contained
in

\SGApage{I}{\(\approx\)470}
\begin{equation*}
 N=\ker\bigl(\underline{\operatorname{Pic}}_{C/k}
       \longrightarrow\underline{\operatorname{Pic}}_{\widetilde C/k}\bigr),
\end{equation*}

\SGAvar{\noindent
and is therefore canonically isomorphic to the maximal torus of this last
group. The structure of $N$ is easily made explicit by EGA IV, 21.8.5:
$N$ is canonically an extension of a torus $T$, which we shall describe,
by a connected unipotent group $U$ (which itself has, over $k$, a
composition series whose quotients are isomorphic to the additive group
$\mathbb G_a$). Moreover, if $T'$ is the maximal torus of $N$, then
$T'\cap U$ is both unipotent and of multiplicative type, and is therefore
zero (SGA 3, XVII, 2.4), while the image of $T'$ in $T$ is all of $T$
(SGA 3, XII, 7.1 e)). Thus $T'\to T$ is an isomorphism. Hence $T$ is
canonically isomorphic to the maximal torus of $N$, and therefore to that
of $\underline{\operatorname{Pic}}^o_{X/k}$.


}{\noindent
and is therefore canonically isomorphic to the maximal torus of this last
group. The structure of $N$ is easily made explicit by EGA IV, 21.8.5:
$N$ is canonically an extension of a torus $T$, which we shall describe,
by a connected unipotent group $U$ (which itself has, over $k$, a
composition series whose quotients are isomorphic to the additive group
$\mathbb G_a$). Moreover, if $T'$ is the maximal torus of $N$, then
$T'\cap U$ is both unipotent and of multiplicative type, and is therefore
zero (SGA 3, XVII, 2.4), while the image of $T'$ in $T$ is all of $T$
(SGA 3, XII, 7.1 e)). Thus $T'\to T$ is an isomorphism. Hence $T$ is
canonically isomorphic to the maximal torus of $N$, and therefore to that
of $\underline{\operatorname{Pic}}^o_{C/k}$.


}{E7I-IX-092}To make the structure of $N$ explicit, it is convenient to suppose that
$k$ is separably closed. The irreducible components $C_i$ of $C$ are then
geometrically irreducible, and the nonnormal points of $C$ are radicial
over $k$. Let $J$ be the set of irreducible components of $\widetilde C$,
and for every $x\in C$ let $J(x)$ be the set of ``branches'' of $C$ through
$x$ (that is, the points of $\widetilde C$ above $x$). Finally put

\begin{equation}\tag{\SGARefTarget{sga7.sga7i.eix.n12-3-1}{12.3.1}}
 R(x)=\mathbb Z^{J(x)}/\operatorname{Im}\mathbb Z,
\end{equation}

\noindent
where $\mathbb Z\to\mathbb Z^{J(x)}$ is the diagonal map, and

\begin{equation}\tag{\SGARefTarget{sga7.sga7i.eix.n12-3-2}{12.3.2}}
 R=\left(\bigoplus_x R(x)\right)/\operatorname{Im}\mathbb Z^J.
\end{equation}

\noindent
Clearly it is enough to take the sum over the finite set of points of $C$
which belong to at least two branches of $C$. Thus $R$ is a finitely
generated $\mathbb Z$-module. There is then a canonical isomorphism

\SGApage{I}{\(\approx\)471}
\SGAvar{\begin{equation}\tag{\SGARefTarget{sga7.sga7i.eix.n12-3-3}{12.2.3}}
 T\ \bigl(\simeq\text{ the maximal torus of }
       \underline{\operatorname{Pic}}^o_{C/k}\bigr)
 \simeq R\otimes\mathbb G_{m,k}.
\end{equation}

}{\begin{equation}\tag{\SGARefTarget{sga7.sga7i.eix.n12-3-3}{12.3.3}}
 T\ \bigl(\simeq\text{ the maximal torus of }
       \underline{\operatorname{Pic}}^o_{C/k}\bigr)
 \simeq R\otimes\mathbb G_{m,k}.
\end{equation}

}{E7I-IX-093}\noindent
In other words, the character group $M$ of $T$ is

\SGAvar{\begin{equation}\tag{\SGARefTarget{sga7.sga7i.eix.n12-3-4}{12.2.4}}
 M\mathrel{\mathop{=}^{\mathrm{def}}}D(T)=\check R
 \simeq\ker\left(\bigoplus_x\check R(x)\longrightarrow\mathbb Z^J\right),
\end{equation}

}{\begin{equation}\tag{\SGARefTarget{sga7.sga7i.eix.n12-3-4}{12.3.4}}
 M\mathrel{\mathop{=}^{\mathrm{def}}}D(T)=\check R
 \simeq\ker\left(\bigoplus_x\check R(x)\longrightarrow\mathbb Z^J\right),
\end{equation}

}{E7I-IX-094}\noindent
where $\check R(x)$ may also be interpreted as the set of
integer-valued functions on $J(x)$ whose values have sum zero.


\medskip
\noindent
\SGARefTarget{sga7.sga7i.eix.n12-3-5}{12.3.5}. We shall need only the case in which at most two branches of $C$
pass through any point $x$ (as happens, in particular, for ordinary double
points). If $\operatorname{card}J(x)=2$, then $R(x)$ is isomorphic to
$\mathbb Z$, but the isomorphism depends on an ordering of the two branches
$(C',C'')$ through $x$, or equivalently on the choice of one branch $C'$:
to $C'$ one associates the element of $R(x)$ which is the canonical image
of $0_{C'}+1_{C''}$. One then obtains an isomorphism

\begin{equation}\tag{\SGARefTarget{sga7.sga7i.eix.n12-3-6}{12.3.6}}
 R=\mathbb Z^I/\operatorname{Im}\mathbb Z^J,
 \qquad
 T\simeq\mathbb G_m^I/\operatorname{Im}\mathbb G_m^J,
 \qquad
 M=\ker\bigl(\mathbb Z^I\longrightarrow\mathbb Z^J\bigr),
\end{equation}

\noindent
where $I$ is the set of points $x\in C$ which belong to two branches of
$C$, after a choice of signs as just explained.

\medskip
\noindent
\underline{Remarks} \SGARefTarget{sga7.sga7i.eix.n12-3-7}{12.3.7}. A sometimes convenient way to describe the
character group $M$ of the maximal torus $T$ of
$\underline{\operatorname{Pic}}^o_{C/k}$ is the following. Consider the
``unoriented one-dimensional graph'' $\Gamma(C)$ (N. Bourbaki,
\emph{Groupes et alg\`ebres de Lie}, Chapter IV, Appendix), whose vertices
are the irreducible components of $\widetilde C$. Two vertices
$\widetilde C_i$ and $\widetilde C_j$ are joined by a set of edges in
bijective correspondence with the pairs of \underline{distinct} branches
of $C$, contained respectively in $\widetilde C_i$ and
$\widetilde C_j$, which pass through the same point of $C$. This set is
$\widetilde C_i\times_C\widetilde C_j$ when $i\neq j$; when $i=j$, it is

\SGApage{I}{\(\approx\)472}
\begin{equation*}
 \bigl(\widetilde C_i\times_C\widetilde C_i-
       \Delta(\widetilde C_i)\bigr)/(\text{symmetry}),
\end{equation*}

\noindent
where $\Delta:\widetilde C\to\widetilde C\times_C\widetilde C$ is the
diagonal. (It is therefore empty if $C_i$ is geometrically unibranch.) One
then has canonical isomorphisms

\begin{equation}\tag{\SGARefTarget{sga7.sga7i.eix.n12-3-7.occ-2}{12.3.7}}
 M=D(T)\xleftarrow{\ \sim\ }H_1(\Gamma(C),\mathbb Z),
 \qquad
 \check M=R\simeq H^1(\Gamma(C),\mathbb Z).
\end{equation}

We shall not need this combinatorial interpretation of $M$, and shall
confine ourselves to defining the canonical homomorphism

\SGAvar{\begin{equation}\notag\SGARefAliasTarget{sga7.sga7i.eix.n12-3-8}
 H_1(\Gamma(C),\mathbb Z)\longrightarrow M,
\end{equation}

}{\begin{equation}\tag{\SGARefTarget{sga7.sga7i.eix.n12-3-8}{12.3.8}}
 H_1(\Gamma(C),\mathbb Z)\longrightarrow M,
\end{equation}

}{E7I-IX-095}\noindent
leaving it to the reader to prove that this is an isomorphism. Let $I$ be
the set of points of $C$ which belong to at least two branches of $C$,
let $J$, as above, be the set of irreducible components of $\widetilde C$
(or of $C$, which amounts to the same thing), and let $K$ be the set of
``branches'' of $C$ at the points of $I$, that is, the inverse image of
$I$ in $\widetilde C$. There is a canonical map

\begin{equation}\tag{\SGARefTarget{sga7.sga7i.eix.n12-3-9}{12.3.9}}
 h:K\longrightarrow I\times J,
 \qquad h=(f,g),
\end{equation}

\noindent
and the description (\SGARefLink{sga7.sga7i.eix.n12-3-4}{12.3.4}) of $M$ is expressed in terms of $h$ as

\begin{equation}\tag{\SGARefTarget{sga7.sga7i.eix.n12-3-10}{12.3.10}}
 M\simeq\ker\left(
  \mathbb Z^{(h)}:\mathbb Z^{(K)}\longrightarrow
  \mathbb Z^{(I\times J)}\right).
\end{equation}

\SGAvar{Indeed, the expression (\SGARefLink{sga7.sga7i.eix.n12-3-3}{12.3.3}) for $R=\check M$ may plainly be
interpreted as the cokernel of the homomorphism
$\mathbb Z^{I\times J}\to\mathbb Z^K$ induced by $h$, which is the
transpose of $\mathbb Z^{(h)}$. The graph $\Gamma(C)$ is likewise described
in terms of $h$ by

}{Indeed, the expression (\SGARefLink{sga7.sga7i.eix.n12-3-3}{12.3.2}) for $R=\check M$ may plainly be
interpreted as the cokernel of the homomorphism
$\mathbb Z^{I\times J}\to\mathbb Z^K$ induced by $h$, which is the
transpose of $\mathbb Z^{(h)}$. The graph $\Gamma(C)$ is likewise described
in terms of $h$ by

}{E7I-IX-096}\SGApage{I}{\(\approx\)473}
\begin{equation*}
 \Gamma_0=J,
 \qquad
 \Gamma_1=(K\times_I K-\Delta(K))/(\text{symmetry}).
\end{equation*}

Choose an orientation on each edge of $\Gamma(C)$, and denote by
$\Gamma_1$ the set of these oriented edges $(x,C',C'')$. Then

\begin{equation}\tag{\SGARefTarget{sga7.sga7i.eix.n12-3-11}{12.3.11}}
 H_1(\Gamma(C),\mathbb Z)=
 \ker\left(d:\mathbb Z^{(\Gamma_1)}\longrightarrow\mathbb Z^{(J)}\right),
\end{equation}

\noindent
where the differential is

\begin{equation}\tag{\SGARefTarget{sga7.sga7i.eix.n12-3-12}{12.3.12}}
 d(x,C',C'')=g(C'')-g(C').
\end{equation}

The homomorphism (\SGARefLink{sga7.sga7i.eix.n12-3-8}{12.3.8}) is induced by the following homomorphism $v$ of
chain complexes:

\SGAvar{\begin{equation}\tag{\SGARefTarget{sga7.sga7i.eix.n12-3-13}{12.3.13}}
 \begin{split}
 v=(v_0,v_1):C_\bullet(\Gamma(C))
   &=\left[\mathbb Z^{(\Gamma)}\longrightarrow\mathbb Z^{(J)}\right]\\
   &\longrightarrow
     \left[\mathbb Z^{(K)}\longrightarrow\mathbb Z^{(I\times J)}\right],
 \end{split}
\end{equation}

}{\begin{equation}\tag{\SGARefTarget{sga7.sga7i.eix.n12-3-13}{12.3.13}}
 \begin{split}
 v=(v_0,v_1):C_\bullet(\Gamma(C))
   &=\left[\mathbb Z^{(\Gamma_1)}\longrightarrow\mathbb Z^{(J)}\right]\\
   &\longrightarrow
     \left[\mathbb Z^{(K)}\longrightarrow\mathbb Z^{(I\times J)}\right],
 \end{split}
\end{equation}

}{E7I-IX-097}\noindent
given by

\begin{equation}\tag{\SGARefTarget{sga7.sga7i.eix.n12-3-14}{12.3.14}}
 v_1(x,C',C'')=(x,C'')-(x,C'),
 \qquad
 v_0(\widetilde C_j)=(0,\widetilde C_j).
\end{equation}


\SGAvar{\medskip
\noindent
\SGARefTarget{sga7.sga7i.eix.n12-4}{12.4}. We now place ourselves in the case considered at the end of \SGARefLink{sga7.sga7i.eix.n12-2}{12.2}.
By (\SGARefLink{sga7.sga7i.eix.n12-1-12}{12.1.12}) and (\SGARefLink{sga7.sga7i.eix.n12-3-6}{12.3.6}), applied with $C=X_o$, the character group
$M$ of the maximal torus $T_o$ of the special fibre $A_o$ of the
N\'eron model $A$ of
$A=\underline{\operatorname{Pic}}^o_{X_\eta/\eta}$ is given by

}{\medskip
\noindent
\SGARefTarget{sga7.sga7i.eix.n12-4}{12.4}. We now place ourselves in the case considered at the end of \SGARefLink{sga7.sga7i.eix.n12-2}{12.2}.
By (\SGARefLink{sga7.sga7i.eix.n12-1-12}{12.1.12}) and (\SGARefLink{sga7.sga7i.eix.n12-3-6}{12.3.6}), applied with $C=X_o$, the character group
$M$ of the maximal torus $T_o$ of the special fibre $A_o$ of the
N\'eron model $A$ of
$A_\eta=\underline{\operatorname{Pic}}^o_{X_\eta/\eta}$ is given by

}{E7I-IX-098}\begin{equation}\tag{\SGARefTarget{sga7.sga7i.eix.n12-4-1}{12.4.1}}
 M=\ker\left(\mathbb Z^I\xrightarrow{\ d\ }\mathbb Z^J\right),
\end{equation}

\noindent
where $I$ is the set of singular points of $X_o$, and $J$ is the set of
its irreducible components, the homomorphism
$\mathbb Z^I\to\mathbb Z^J$ being given by

\begin{equation}\tag{\SGARefTarget{sga7.sga7i.eix.n12-4-2}{12.4.2}}
 d(x_i)=C'_i-C''_i.
\end{equation}

\SGApage{I}{\(\approx\)474}
Here $C'_i$ and $C''_i$ are the two irreducible components of
$\widetilde C$ corresponding to the branches through the double point
$x_i$. The operator $d$ depends on the choice of an order
$(C'_i,C''_i)$ on these two branches, and the isomorphism (\SGARefLink{sga7.sga7i.eix.n12-4-1}{12.4.1})
plainly depends on this system of choices. Nevertheless, if, for
$i\in I$, we denote by

\begin{equation}\tag{\SGARefTarget{sga7.sga7i.eix.n12-4-3}{12.4.3}}
 \varphi_{x_i}:M\longrightarrow\mathbb Z
\end{equation}

\noindent
the composite

\begin{equation*}
 M\longrightarrow\mathbb Z^I
   \xrightarrow{\ \operatorname{pr}_i\ }\mathbb Z,
\end{equation*}

\noindent
where the first homomorphism comes from (\SGARefLink{sga7.sga7i.eix.n12-4-1}{12.4.1}), then, on returning to
the expression (\SGARefLink{sga7.sga7i.eix.n12-3-4}{12.3.4}), one sees at once that $\varphi_{x_i}$ is
well-defined up to sign. More precisely, it depends only on the order
chosen between the two branches $C'_i,C''_i$ through $x_i$, and changes
sign when that order is reversed. It follows in particular that the
bilinear form

\begin{equation}\tag{\SGARefTarget{sga7.sga7i.eix.n12-4-4}{12.4.4}}
 \varphi_{x_i}\otimes\varphi_{x_i}
 \in \check M\otimes_{\mathbb Z}\check M
 \simeq\operatorname{Hom}(M\otimes M,\mathbb Z)
\end{equation}

\noindent
is determined canonically by the singular point $x_i$, without any sign
ambiguity. We may therefore put

\begin{equation}\tag{\SGARefTarget{sga7.sga7i.eix.n12-4-5}{12.4.5}}
 u=\sum_i\varphi_{x_i}\otimes\varphi_{x_i}:
 M\otimes_{\mathbb Z}M\longrightarrow\mathbb Z.
\end{equation}

This gives a well-defined pairing, independent of every choice of
signs. It is plainly \underline{positive definite}, since it is the
restriction to $M\hookrightarrow\mathbb Z^I$ of the standard quadratic
form $\sum_i X_i^2$ on $\mathbb Z^I$.

\SGApage{I}{\(\approx\)475}
The notation $M$ for the character group of $T_o$ is not, a priori,
consistent with that of the preceding sections, where this group was
written $M'$, while $M$ denoted the analogous group associated with the
dual abelian scheme $A'_K$ of $A_K$. But since $A_K$ is a Jacobian, it
has a canonical polarization $\xi$ [5], which induces a canonical
isomorphism

\begin{equation}\tag{\SGARefTarget{sga7.sga7i.eix.n12-4-6}{12.4.6}}
 \psi_\xi:A_K\xrightarrow{\ \sim\ }A'_K.
\end{equation}

\noindent
By means of this isomorphism we shall identify $A_K$ and $A'_K$, so that
there is no longer any need to distinguish $M$ from $M'$. In particular,
for every prime number $\ell$, the monodromy pairing (\SGARefLink{sga7.sga7i.eix.n9-1-2}{9.1.2}) may be
regarded as a pairing

\begin{equation}\tag{\SGARefTarget{sga7.sga7i.eix.n12-4-7}{12.4.7}}
 u_\ell:M_\ell\otimes_{\mathbb Z_\ell}M_\ell
 \longrightarrow\mathbb Z_\ell.
\end{equation}

With this understood, we can state the principal result of the present
section:

\SGAvar{\medskip
\noindent
\underline{Theorem} \SGARefTarget{sga7.sga7i.eix.n12-5}{12,5} (\underline{Picard--Lefschetz formula}).
\underline{Let $S$ be a strictly local trait, and let $X$ be a regular
scheme, projective and flat over $S$, with one-dimensional fibres, whose
geometric generic fibre is smooth and connected and whose geometric
special fibre has no singularities other than ordinary quadratic
singularities. With the notation introduced in \SGARefLink{sga7.sga7i.eix.n12-4}{12.4}, for every prime
number $\ell$, the monodromy pairing (\SGARefLink{sga7.sga7i.eix.n12-4-7}{12.4.7}) is equal to the pairing
obtained from (\SGARefLink{sga7.sga7i.eix.n12-4-5}{12.4.5}) by extension of scalars from $\mathbb Z$ to
$\mathbb Z_\ell$.}


}{\medskip
\noindent
\underline{Theorem} \SGARefTarget{sga7.sga7i.eix.n12-5}{12.5} (\underline{Picard--Lefschetz formula}).
\underline{Let $S$ be a strictly local trait, and let $X$ be a regular
scheme, projective and flat over $S$, with one-dimensional fibres, whose
geometric generic fibre is smooth and connected and whose geometric
special fibre has no singularities other than ordinary quadratic
singularities. With the notation introduced in \SGARefLink{sga7.sga7i.eix.n12-4}{12.4}, for every prime
number $\ell$, the monodromy pairing (\SGARefLink{sga7.sga7i.eix.n12-4-7}{12.4.7}) is equal to the pairing
obtained from (\SGARefLink{sga7.sga7i.eix.n12-4-5}{12.4.5}) by extension of scalars from $\mathbb Z$ to
$\mathbb Z_\ell$.}


}{E7I-IX-099}\SGApage{I}{476}
\medskip
\noindent
\SGARefTarget{sga7.sga7i.eix.n12-6}{12.6}. Notice that this statement proves \SGARefLink{sga7.sga7i.eix.n10-4}{10.4} in the particular case
under consideration when, in \SGARefLink{sga7.sga7i.eix.n10-4}{10.4} b), one takes for $\xi$ the canonical
polarization considered there. As we showed in \SGARefLink{sga7.sga7i.eix.n10-5}{10.5}, this suffices to
imply \SGARefLink{sga7.sga7i.eix.n10-4}{10.4} in the general case.

\medskip
\noindent
\SGARefTarget{sga7.sga7i.eix.n12-7}{12.7}. \underline{Proof of \SGARefLink{sga7.sga7i.eix.n12-5}{12.5}. The case $\ell\ne p$.} It relies
essentially on the \underline{cohomological Picard--Lefschetz formula}
in the situation considered in \SGARefLink{sga7.sga7i.eix.n12-5}{12.5}, a formula giving the action of the
monodromy group $I$ on the following group:

\begin{equation}\tag{\SGARefTarget{sga7.sga7i.eix.n12-7-1}{12.7.1}}
 H^1(X_{\bar\eta},\mathbb Z_\ell)
 \simeq H^1(A_{\bar\eta},\mathbb Z_\ell)
 \simeq\operatorname{Hom}
   \bigl(T_\ell(A_{\bar\eta}),\mathbb Z_\ell\bigr)
\end{equation}

\noindent
when $\ell\ne p$. This formula, which will be proved in a later expos\'e
(Expos\'e XV), and indeed in a substantially more general case, is
essentially identical to assertion \SGARefLink{sga7.sga7i.eix.n12-5}{12.5} for the prime $\ell$ in
question, in view of the definition in \SGARefLink{sga7.sga7i.eix.n9-2}{9.2} of $u_\ell$ for $\ell\ne p$.
It therefore remains to treat the case $\ell=p$. We shall do this by a
method which reduces it to characteristic zero (and, moreover, to the
case where $X_o$ has only one singular point), which we now explain.

\SGAvar{\medskip
\noindent
\underline{Remarks \SGARefTarget{sga7.sga7i.eix.n12-7-2}{12.7.2}.} Notice also that, once the announced
reduction has been made, the usual arguments would allow us to reduce to
the case in which the base field is the field $\mathbb C$ of complex
numbers, and then to invoke the Picard--Lefschetz formula in its classical
form---except that there does not seem to be a satisfactory reference
for that formula. The proof of the Picard--Lefschetz formula given in our
Seminar will proceed similarly, by reduction to the transcendental case,
which is amenable to properly transcendental methods. One may therefore
say that, in any event, the proof of \SGARefLink{sga7.sga7i.eix.n12-5}{12.5} (and hence of \SGARefLink{sga7.sga7i.eix.n10-4}{10.4} proposed
here is transcendental in nature. Let us also point out that the only
other known proof of \SGARefLink{sga7.sga7i.eix.n10-4}{10.4} comes from rigid-analytic theory [24], and must
for that reason likewise be regarded as essentially ``transcendental.''
It would be interesting to find purely algebraic proofs of these results
(although it is by no means evident that this is possible!).


}{\medskip
\noindent
\underline{Remarks \SGARefTarget{sga7.sga7i.eix.n12-7-2}{12.7.2}.} Notice also that, once the announced
reduction has been made, the usual arguments would allow us to reduce to
the case in which the base field is the field $\mathbb C$ of complex
numbers, and then to invoke the Picard--Lefschetz formula in its classical
form---except that there does not seem to be a satisfactory reference
for that formula. The proof of the Picard--Lefschetz formula given in our
Seminar will proceed similarly, by reduction to the transcendental case,
which is amenable to properly transcendental methods. One may therefore
say that, in any event, the proof of \SGARefLink{sga7.sga7i.eix.n12-5}{12.5} (and hence of \SGARefLink{sga7.sga7i.eix.n10-4}{10.4}) proposed
here is transcendental in nature. Let us also point out that the only
other known proof of \SGARefLink{sga7.sga7i.eix.n10-4}{10.4} comes from rigid-analytic theory [24], and must
for that reason likewise be regarded as essentially ``transcendental.''
It would be interesting to find purely algebraic proofs of these results
(although it is by no means evident that this is possible!).


}{E7I-IX-100}\SGApage{I}{477}
\medskip
\noindent
\SGARefTarget{sga7.sga7i.eix.n12-8}{12.8}. \underline{Proof of \SGARefLink{sga7.sga7i.eix.n12-5}{12.5}: reduction to the case in which $S$ has
characteristic zero.} We shall use an argument based on the same principle
as the proof of the ``theorem on the essentially tame action of monodromy
on $\pi_1$'' in Expos\'e V. Notice that
we may assume $S$ complete, by \SGARefLink{sga7.sga7i.eix.n10-3}{10.3}. Let $W$ be a Cohen $p$-ring with
residue field $k$, and choose a homomorphism $W\longrightarrow V$ inducing
the identity on residue fields (EGA $0_{\mathrm{IV}}$, 19.8.6). Let
$(x_\lambda)_{\lambda\in L}$ be the set of singular points of $X_o$, and
consider the formal moduli scheme $\underline S$ of $X_o$ over $W$, together
with the ``universal curve''
\[
 \underline f:\underline X\longrightarrow\underline S
\]
whose special fibre is $X_o$ (Expos\'e VI). We saw in
the cited place that $\underline S$ is isomorphic to the formal spectrum of
a power-series ring $W[[T_1,\ldots,T_N]]$, and that in $\underline S$ one
can find a normal-crossings divisor

\begin{equation}\tag{\SGARefTarget{sga7.sga7i.eix.n12-8-1}{12.8.1}}
 D=\sum_{\lambda\in L}D_\lambda,
\end{equation}

\noindent
whose irreducible components $D_\lambda$ are indexed by $L$, each
$D_\lambda$ being described by an equation of the form $T_{i_\lambda}=0$.
(If desired, one may suppose $i_\lambda=\lambda$ after indexing $L$ by a
set of integers $[1,r]$.) The component $D_\lambda$ is the image under
$\underline f:\underline X\longrightarrow\underline S$ of the component
$D'_\lambda$ through the singular point $x_\lambda$ of $X_o$ in the
non-smoothness locus
\[
 D'=\bigcup_{\lambda\in L}D'_\lambda
\]
of $\underline f$. By the construction itself,
$f:X\longrightarrow S$ is isomorphic to the pullback of
$\underline f:\underline X\longrightarrow\underline S$ by a suitable
$W$-morphism $g:S\longrightarrow\underline S$:

\SGApage{I}{\(\approx\)478}
\begin{equation}\tag{\SGARefTarget{sga7.sga7i.eix.n12-8-2}{12.8.2}}
\begin{tikzcd}
 \underline X \arrow[d,"\underline f"'] \&
 X \arrow[l] \arrow[d,"f"] \\
 \underline S \& S \arrow[l,"g"']
\end{tikzcd}
\qquad .
\end{equation}

\noindent
Set

\begin{equation}\tag{\SGARefTarget{sga7.sga7i.eix.n12-8-3}{12.8.3}}
 P=\underline{\operatorname{Pic}}^{o}_{\underline X/\underline S}.
\end{equation}

\SGAvar{\noindent
Since the special fibre of $\underline f$ is geometrically reduced of
dimension $1$, it follows that $\underline F$ is cohomologically flat in
dimension zero, by EGA III, 7.8.1, and that $P$ is formally smooth over
$ S$ (the latter fact follows from
$H^2(X_o,\underline{\mathcal O}_{X_o})=0$, since $X_o$ has dimension $1$).
For simplicity, we shall assume that $P$ \underline{is representable by a
group scheme over $ S$}; it will necessarily be a smooth group
scheme with connected fibres. This is a special case of an unpublished
result of Mumford cited in [22] (and considerably generalized by Raynaud in
the note cited there, when the base is a trait). Since no proof of either
Mumford's or Raynaud's result has been published, the reader may prefer to
use the general results of M. Artin [1]. They give the representability of
$P$---under the sole assumption that $\underline f$ is proper, flat, and
cohomologically flat in dimension zero---by an ``algebraic space'' in
Artin's sense, which would suffice for our purpose. Put

}{\noindent
Since the special fibre of $\underline f$ is geometrically reduced of
dimension $1$, it follows that $\underline f$ is cohomologically flat in
dimension zero, by EGA III, 7.8.1, and that $P$ is formally smooth over
$\underline S$ (the latter fact follows from
$H^2(X_o,\underline{\mathcal O}_{X_o})=0$, since $X_o$ has dimension $1$).
For simplicity, we shall assume that $P$ \underline{is representable by a
group scheme over $\underline S$}; it will necessarily be a smooth group
scheme with connected fibres. This is a special case of an unpublished
result of Mumford cited in [22] (and considerably generalized by Raynaud in
the note cited there, when the base is a trait). Since no proof of either
Mumford's or Raynaud's result has been published, the reader may prefer to
use the general results of M. Artin [1]. They give the representability of
$P$---under the sole assumption that $\underline f$ is proper, flat, and
cohomologically flat in dimension zero---by an ``algebraic space'' in
Artin's sense, which would suffice for our purpose. Put

}{E7I-IX-101}\SGApage{I}{\(\approx\)479}
\begin{equation}\tag{\SGARefTarget{sga7.sga7i.eix.n12-8-4}{12.8.4}}
 U=\underline S-\operatorname{supp}D.
\end{equation}

\noindent
Thus $U$ is the largest open subset of $\underline S$ over which
$\underline X$ is smooth, and hence
\[
 P|_U=\underline{\operatorname{Pic}}^{o}_{\underline X_U/U}
\]
is an abelian scheme over $U$. Its pullback to $\eta$ by
$g:\eta\longrightarrow U$ is the abelian scheme $A_\eta$ over $\eta$ that
we wish to study. More precisely, there is a canonical isomorphism over
$S$ (\SGARefLink{sga7.sga7i.eix.n12-1-10}{12.1.10}):

\begin{equation}\tag{\SGARefTarget{sga7.sga7i.eix.n12-8-5}{12.8.5}}
 P\times_{\underline S}S\simeq A^{o}.
\end{equation}

\noindent
Since the special fibre $P_o=A_o^{o}$ is $\ell$-divisible for every
$\ell$, the same holds for $P$ itself; all its fibres are therefore
extensions of an abelian variety by a torus. Consider the pro-$\ell$ group
$T_\ell(P)$ over $\underline S$, whose pullback by $g$ is precisely
$T_\ell(A^{o})$. Applying the constructions of \SGARefLink{sga7.sga7i.eix.n7-3}{7.3}, one defines a
filtration of this pro-group by a ``fixed part'' and a ``toric part'':

\begin{equation}\tag{\SGARefTarget{sga7.sga7i.eix.n12-8-6}{12.8.6}}
 T_\ell(P)\supset T_\ell(P)^f=T_\ell(\widehat P)
 \supset T_\ell(P)^t=T_\ell(T)\supset0.
\end{equation}

\noindent
Its pullback by $g$ is the analogous filtration of $T_\ell(A^{o})$
studied in sections 2 and 5.

Consider the restriction of this filtration over the open subset $U$ of
$\underline S$. There we have the canonical self-duality of $P_U$, as the
Jacobian of a smooth curve scheme over $U$, and hence a corresponding
pairing

\begin{equation}\tag{\SGARefTarget{sga7.sga7i.eix.n12-8-7}{12.8.7}}
 \varphi:T_\ell(P_U)\otimes T_\ell(P_U)
 \longrightarrow T_\ell(\mathbb G_{mU}).
\end{equation}

\SGApage{I}{\(\approx\)480}
\SGAvar{\noindent
This is a \underline{self-duality}, in the sense of Barsotti--Tate groups,
of $T_\ell(B_U)$. We saw in \SGARefLink{sga7.sga7i.eix.n7-4-8}{7.4.8} that, for this pairing,
$T_\ell(P_U)^f$ \underline{and} $T_\ell(P_U)^t$ \underline{are each
other's orthogonal complements}. It follows that
\[
 E_U=T_\ell(P_U)/T_\ell(P_U)^t
\]
is canonically isomorphic to the Barsotti--Tate pro-group dual to
$T_\ell(P_U)^f=F_U$. It therefore extends canonically to a Barsotti--Tate
group $E$ over $ S$, namely

}{\noindent
This is a \underline{self-duality}, in the sense of Barsotti--Tate groups,
of $T_\ell(P_U)$. We saw in \SGARefLink{sga7.sga7i.eix.n7-4-8}{7.4.8} that, for this pairing,
$T_\ell(P_U)^f$ \underline{and} $T_\ell(P_U)^t$ \underline{are each
other's orthogonal complements}. It follows that
\[
 E_U=T_\ell(P_U)/T_\ell(P_U)^t
\]
is canonically isomorphic to the Barsotti--Tate pro-group dual to
$T_\ell(P_U)^f=F_U$. It therefore extends canonically to a Barsotti--Tate
group $E$ over $\underline S$, namely

}{E7I-IX-102}\begin{equation}\tag{\SGARefTarget{sga7.sga7i.eix.n12-8-8}{12.8.8}}
 E=D(F),\qquad\text{where}\qquad F=T_\ell(P)^f.
\end{equation}

\noindent
We are now in a typical situation of mixed extensions (\SGARefLink{sga7.sga7i.eix.n9-3}{9.3}), and can repeat
the constructions of \SGARefLink{sga7.sga7i.eix.n9-4}{9.4}--\SGARefLink{sga7.sga7i.eix.n9-6}{9.6} (for whose having been written in the stale
and piss-poor setting of a base scheme reduced to a trait we apologize!):
$T_\ell(P_U)$ is a mixed extension of $E_U$ by $F_U$, where $E$ and $F$
are Barsotti--Tate groups over $\underline S$ which are extensions, dual to
one another:


\SGAvar{\begin{equation}\tag{\SGARefTarget{sga7.sga7i.eix.n12-8-9}{12.8.9}}
\begin{tikzcd}[row sep=small,column sep=small]
0 \arrow[r] \& P \arrow[rr] \arrow[d,dash,"\wr"]
 \& \& F \arrow[r] \arrow[d,dash,"\wr"] \& R \arrow[r] \& 0 \& \\
\& T_\ell(P)^t\simeq
  \underline{\operatorname{Hom}}
    \bigl(M,T_\ell(\underline{\mathbb G}_{m\underline S})
\& \& T_\ell(P)^f \& \& \& \\
\& \& \& \& \& \& \\
0 \arrow[r] \& R \arrow[rr] \arrow[d,dash,"\wr"]
\& \& E \arrow[r] \& Q \arrow[r] \arrow[d,dash,"\wr"] \& 0 \& \\
\& D(R) \& \& \& M\otimes\mathbb Z_{\ell\,\underline S} \& \& .
\end{tikzcd}
\end{equation}

\noindent
 Notice that $R$ is
canonically self-dual, as follows from the self-duality (\SGARefLink{sga7.sga7i.eix.n12-8-7}{12.8.7}) and the
orthogonality statement noted above. The only difference from the cited
construction is that the pairing obtained on $M_\ell\times M_\ell$ takes
values in the group
\[
 \bigl(\mathbb G_m(U)/\mathbb G_m(\underline S)
 \otimes_{\mathbb Z}\mathbb Z_\ell,
\]
which is no longer canonically isomorphic to $\mathbb Z_\ell$, but rather
to $\mathbb Z_\ell^{\,L}$, by means of the isomorphism

}{\begin{equation}\tag{\SGARefTarget{sga7.sga7i.eix.n12-8-9}{12.8.9}}
\begin{tikzcd}[row sep=small,column sep=small]
0 \arrow[r] \& P \arrow[rr] \arrow[d,dash,"\wr"]
 \& \& F \arrow[r] \arrow[d,dash,"\wr"] \& R \arrow[r] \& 0 \& \\
\& T_\ell(P)^t\simeq
  \underline{\operatorname{Hom}}
    \bigl(M,T_\ell(\mathbb G_{m\underline S})\bigr)
\& \& T_\ell(P)^f \& \& \& \\
\& \& \& \& \& \& \\
0 \arrow[r] \& R \arrow[rr] \arrow[d,dash,"\wr"]
\& \& E \arrow[r] \& Q \arrow[r] \arrow[d,dash,"\wr"] \& 0 \& \\
\& D(R) \& \& \& M\otimes\mathbb Z_{\ell\,\underline S} \& \& .
\end{tikzcd}
\end{equation}

\noindent
 Notice that $R$ is
canonically self-dual, as follows from the self-duality (\SGARefLink{sga7.sga7i.eix.n12-8-7}{12.8.7}) and the
orthogonality statement noted above. The only difference from the cited
construction is that the pairing obtained on $M_\ell\times M_\ell$ takes
values in the group
\[
 \bigl(\mathbb G_m(U)/\mathbb G_m(\underline S)\bigr)
 \otimes_{\mathbb Z}\mathbb Z_\ell,
\]
which is no longer canonically isomorphic to $\mathbb Z_\ell$, but rather
to $\mathbb Z_\ell^{\,L}$, by means of the isomorphism

}{E7I-IX-103}\SGApage{I}{\(\approx\)481}
\begin{equation}\tag{\SGARefTarget{sga7.sga7i.eix.n12-8-10}{12.8.10}}
 \mathbb Z^L\simeq
 \mathbb G_m(U)/\mathbb G_m(\underline S).
\end{equation}

\noindent
It maps the element indexed by $\lambda$ in the canonical basis of
$\mathbb Z^L$ to the class of $T_{i_\lambda}$ in the right-hand side,
where $T_{i_\lambda}$ is an equation for the divisor $D_\lambda$. We thus
obtain a canonical pairing

\begin{equation}\tag{\SGARefTarget{sga7.sga7i.eix.n12-8-11}{12.8.11}}
 \underline{u}_\ell:M_\ell\otimes M_\ell
 \longrightarrow\mathbb Z_\ell^{\,L}\qquad .
\end{equation}

\SGAvar{\noindent
It is essentially immediate that the monodromy pairing $u_\ell$ (\SGARefLink{sga7.sga7i.eix.n12-4-7}{12.4.7})
over $S$ which we propose to study is the composite of
${u}_\ell$ with the homomorphism
$\mathbb Z_\ell^{\,L}\longrightarrow\mathbb Z_\ell$ obtained, by
tensoring with $\mathbb Z_\ell$ over $\mathbb Z$, from the homomorphism

}{\noindent
It is essentially immediate that the monodromy pairing $u_\ell$ (\SGARefLink{sga7.sga7i.eix.n12-4-7}{12.4.7})
over $S$ which we propose to study is the composite of
$\underline{u}_\ell$ with the homomorphism
$\mathbb Z_\ell^{\,L}\longrightarrow\mathbb Z_\ell$ obtained, by
tensoring with $\mathbb Z_\ell$ over $\mathbb Z$, from the homomorphism

}{E7I-IX-104}\begin{center}
\begin{tikzcd}
 \mathbb G_m(U)/\mathbb G_m(\underline S)
   \arrow[r] \arrow[d,dash,"\wr"'] \&
 \mathbb G_m(\eta)/\mathbb G_m(S)
   \arrow[d,dash,"\wr"'] \\
 \mathbb Z^L \& \mathbb Z
\end{tikzcd}
\end{center}

\noindent
induced by $g:S\longrightarrow\underline S$ (which sends $\eta$ into
$U$). This homomorphism is given by an element

\begin{equation}\tag{\SGARefTarget{sga7.sga7i.eix.n12-8-12}{12.8.12}}
 (d_\lambda)_{\lambda\in L}\in\mathbb Z^L,\qquad
 d_\lambda\geqslant1\quad\text{for every }\lambda\in L,
\end{equation}

\SGApage{I}{\(\approx\)482}
\SGAvar{\noindent
where, for every $\lambda\in L$, $d_\lambda$ is the multiplicity of the
inverse image in $S$ of the divisor $D_\lambda$ on $\underline S$. It in
fact depends only on the completion
$\widehat{\underline{\mathcal O}}_{X,x_\lambda}$ as a $V$-algebra
(Expos\'e VI). Moreover, it is
known (Expos\'e VI) that, for a given $\lambda$, one has $d_\lambda=1$ if
and only if $x$ is \underline{regular} at $x_\lambda$. Since in \SGARefLink{sga7.sga7i.eix.n12-5}{12.5} we assumed $X$ regular, that is, all the
$d_\lambda$ equal to $1$, assertion \SGARefLink{sga7.sga7i.eix.n12-5}{12.5} will therefore follow from the
following lemma, in which $X,S,g$ no longer occur.

}{\noindent
where, for every $\lambda\in L$, $d_\lambda$ is the multiplicity of the
inverse image in $S$ of the divisor $D_\lambda$ on $\underline S$. It in
fact depends only on the completion
$\widehat{\underline{\mathcal O}}_{X,x_\lambda}$ as a $V$-algebra
(Expos\'e VI). Moreover, it is
known (Expos\'e VI) that, for a given $\lambda$, one has $d_\lambda=1$ if
and only if $X$ is \underline{regular} at $x_\lambda$. Since in \SGARefLink{sga7.sga7i.eix.n12-5}{12.5} we assumed $X$ regular, that is, all the
$d_\lambda$ equal to $1$, assertion \SGARefLink{sga7.sga7i.eix.n12-5}{12.5} will therefore follow from the
following lemma, in which $X,S,g$ no longer occur.

}{E7I-IX-105}\medskip
\noindent
\underline{Lemma} \SGARefTarget{sga7.sga7i.eix.n12-8-13}{12.8.13}. \underline{If $\lambda\in L$, the composite of
(\SGARefLink{sga7.sga7i.eix.n12-8-11}{12.8.11}) with $\operatorname{pr}_\lambda$ is equal to the form
$\varphi_{x_\lambda}\otimes\varphi_{x_\lambda}$ of (\SGARefLink{sga7.sga7i.eix.n12-4-4}{12.4.4}).}

\SGAvar{\medskip
To prove this last assertion, consider the completed strict localization
$\underline S^\lambda$ of $\underline S$ at the generic point of
$D$, and put
\[
 \underline X^\lambda=
  X\times_{ S}\underline S^\lambda.
\]
The special fibre $ X_o^\lambda$ then has only one singular
point $y_\lambda$, which is again an ordinary quadratic singularity. In
the corresponding monodromy pairing

}{\medskip
To prove this last assertion, consider the completed strict localization
$\underline S^\lambda$ of $\underline S$ at the generic point of
$D_\lambda$, and put
\[
 \underline X^\lambda=
 \underline X\times_{\underline S}\underline S^\lambda.
\]
The special fibre $\underline X_o^\lambda$ then has only one singular
point $y_\lambda$, which is again an ordinary quadratic singularity. In
the corresponding monodromy pairing

}{E7I-IX-106}\begin{equation}\tag{\SGARefTarget{sga7.sga7i.eix.n12-8-14}{12.8.14}}
 u_\ell^\lambda:M_\ell^\lambda\otimes M_\ell^\lambda
 \longrightarrow\mathbb Z_\ell\qquad ,
\end{equation}

\noindent
$M^\lambda$ consequently has rank $1$ or $0$. We leave to the reader the
verification of the following easy facts, which make explicit, in the
special case in which we need them, the behaviour under localization of
the preceding constructions.

\medskip
\noindent
a) There is a unique subtorus $T_\lambda$ of
$P\times_{\underline S}D_\lambda$ whose restriction to the generic point
$s_\lambda$ of $D_\lambda$ is the maximal torus of $P_{s_\lambda}$ (which
has dimension $0$ or $1$).

\SGApage{I}{\(\approx\)483}
\medskip
\noindent
b) Consider the transpose homomorphism $M\longrightarrow M^\lambda$ on
character groups of the inclusion
$T_{\lambda o}\hookrightarrow T_o$, where $T_o$ is the maximal torus of
$P_o$, and the form
$\varphi_{y_\lambda}:M^\lambda\longrightarrow\mathbb Z$, defined up to
sign as in \SGARefLink{sga7.sga7i.eix.n12-4-3}{12.4.3} (with $X$ replaced by $\underline X^\lambda$). Then,
up to sign, the following diagram is commutative:

\SGAvar{\begin{equation}\tag{\SGARefTarget{sga7.sga7i.eix.n12-8-15}{12.8.15}}
\begin{tikzcd}[column sep=huge,row sep=large]
 M \arrow[rr] \arrow[dr,"\varphi_{x_\lambda}"']
 \& \& M^\lambda \arrow[dl,"\varphi_{\lambda_\lambda}"] \\
 \& \mathbb Z \& \qquad .
\end{tikzcd}
\end{equation}

}{\begin{equation}\tag{\SGARefTarget{sga7.sga7i.eix.n12-8-15}{12.8.15}}
\begin{tikzcd}[column sep=huge,row sep=large]
 M \arrow[rr] \arrow[dr,"\varphi_{x_\lambda}"']
 \& \& M^\lambda \arrow[dl,"\varphi_{y_\lambda}"] \\
 \& \mathbb Z \& \qquad .
\end{tikzcd}
\end{equation}

}{E7I-IX-107}\noindent
 (In fact, choosing an ordering of the two branches
of $X_o$ through $x_\lambda$ determines an analogous choice at the point
$y_\lambda$ of $\underline X_o^\lambda$, and for choices corresponding in
this way the preceding diagram is commutative.)

\medskip
\noindent
c) The following diagram is commutative:

\begin{center}
\begin{tikzcd}[column sep=huge,row sep=large]
 M_\ell\otimes M_\ell
   \arrow[rr] \arrow[dr,"\operatorname{pr}_\lambda\underline{u}_\ell"']
 \& \& M_\ell^\lambda\otimes M_\ell^\lambda
   \arrow[dl,"u_\ell^\lambda"] \\
 \& \mathbb Z_\ell \& \qquad ,
\end{tikzcd}
\end{center}

\SGAvar{\noindent
where the notation is that of \SGARefLink{sga7.sga7i.eix.n12-8-13}{12.8.13} and (\SGARefLink{sga7.sga7i.eix.n12-8-14}{12.8.14}). These relations show
at once that, to prove \SGARefLink{sga7.sga7i.eix.n12-8-13}{12.8.13}, it is enough to prove the analogous
relation for $\underline X^\lambda$:
\[
 u_\ell^\lambda=
 (\varphi_{y_\lambda}\otimes\varphi_{y_\lambda})
 \otimes_{\mathbb Z}\mathbb Z_\ell,
\]
which is nothing other than \SGARefLink{sga7.sga7i.eix.n12-5}{12.5} for $\underline X^\lambda$ over
$\underline S^\lambda$. Now $\underline S$ is a trait of
\underline{characteristic zero}. This completes the
announced reduction.

}{\noindent
where the notation is that of \SGARefLink{sga7.sga7i.eix.n12-8-13}{12.8.13} and (\SGARefLink{sga7.sga7i.eix.n12-8-14}{12.8.14}). These relations show
at once that, to prove \SGARefLink{sga7.sga7i.eix.n12-8-13}{12.8.13}, it is enough to prove the analogous
relation for $\underline X^\lambda$:
\[
 u_\ell^\lambda=
 (\varphi_{y_\lambda}\otimes\varphi_{y_\lambda})
 \otimes_{\mathbb Z}\mathbb Z_\ell,
\]
which is nothing other than \SGARefLink{sga7.sga7i.eix.n12-5}{12.5} for $\underline X^\lambda$ over
$\underline S^\lambda$. Now $\underline S^\lambda$ is a trait of
\underline{characteristic zero}. This completes the
announced reduction.

}{E7I-IX-108}\SGApage{I}{\(\approx\)484}
\medskip
\noindent
\underline{Remarks \SGARefTarget{sga7.sga7i.eix.n12-9}{12.9}.} Notice that this reduction simultaneously
reduces us to the case in which $X_o$ has only one singular point, the case
classically considered in the Picard--Lefschetz formula. Notice also that,
in this case, $X_o$ has one or two irreducible components, corresponding
respectively to $\operatorname{rank}M=1$ and $\operatorname{rank}M=0$ (the
latter being the case in which the Jacobian of the generic fibre has good
reduction).

\medskip
\noindent
\SGARefTarget{sga7.sga7i.eix.n12-10}{12.10}. The proof given in \SGARefLink{sga7.sga7i.eix.n12-8}{12.8} yields a result more general than \SGARefLink{sga7.sga7i.eix.n12-5}{12.5},
valid without assuming that $X$ is regular, by introducing the integers
$d_\lambda\geqslant1$ of (\SGARefLink{sga7.sga7i.eix.n12-8-12}{12.8.12}): all the monodromy pairings are then
deduced from the pairing

\begin{equation}\tag{\SGARefTarget{sga7.sga7i.eix.n12-10-1}{12.10.1}}
 u=\sum_{\lambda\in L}d_\lambda\,
   \varphi_{x_\lambda}\otimes\varphi_{x_\lambda}:
   M\otimes M\longrightarrow\mathbb Z\qquad .
\end{equation}


\medskip
\noindent
13. \underline{Links with transcendental theory: the complex analytic
case}.

For the general notions concerning analytic spaces, we refer to [2].

\medskip
\noindent
\SGARefTarget{sga7.sga7i.eix.n13-1}{13.1}. \underline{Review of analytic groups over $S$.} Let $S$ be an
analytic space, and let $A$ be an analytic space over $S$, endowed with
the structure of a commutative group over $S$ (that is, $A$ is a
commutative group in the category of analytic spaces over $S$). We assume
that $A$ is \underline{smooth} over $S$, and denote by

\SGApage{I}{\(\approx\)485}
\begin{equation}\tag{\SGARefTarget{sga7.sga7i.eix.n13-1-1}{13.1.1}}
 \underline E=\underline{\operatorname{Lie}}(A/S)
\end{equation}

\noindent
the locally free module on $S$ obtained by pulling the relative tangent
sheaf of $A$ over $S$ back along the zero section. We shall admit here the
definition of the exponential homomorphism

\begin{equation}\tag{\SGARefTarget{sga7.sga7i.eix.n13-1-2}{13.1.2}}
 \exp:E\longrightarrow A\qquad ,
\end{equation}

\noindent
where $E$ is the analytic vector bundle on $S$ whose module of
holomorphic sections is $\underline E$. In a neighbourhood of a point
$s\in S$, this homomorphism is first defined near the zero section of $E$
by the familiar methods of the theory of formal groups in characteristic
zero, and is then extended in the evident way, by multiplicativity, to
all of $E$. It is a homomorphism of relative analytic groups over $S$;
it is \underline{\'etale}, and its image is the open subspace $A^o$ of
$A$ which is the union of the identity components of the fibres of $A$
over $S$. If $R$ is the analytic subspace which is the kernel of $\exp$,
we therefore have an exact sequence of relative analytic groups

\begin{equation}\tag{\SGARefTarget{sga7.sga7i.eix.n13-1-3}{13.1.3}}
 0\longrightarrow R\longrightarrow E\longrightarrow A^o
 \longrightarrow0\qquad .
\end{equation}

\noindent
Here the morphism $E\longrightarrow A^o$ is surjective and \'etale, so
$R$ is a relative analytic group \underline{\'etale} over $S$. It is an
analytic subspace of $E$, and is closed in $E$ if and only if $A^o$ is
separated over $S$, that is, if and only if the zero section of $A^o$ is
closed in $A^o$.

\SGApage{I}{486}
\medskip
\noindent
\SGARefTarget{sga7.sga7i.eix.n13-1-4}{13.1.4}. Of course, the exact sequence (\SGARefLink{sga7.sga7i.eix.n13-1-3}{13.1.3}) is
\underline{functorial in $A$}, and commutes with every base change
$S'\longrightarrow S$. In this way one obtains an equivalence of
categories, compatible with base changes, between the category of
relative analytic groups over $S$ which are smooth and have connected
fibres, and the category of pairs $(E,R)$ consisting of a locally free
vector bundle $E$ over $S$ and an analytic subgroup $R$ of $E$ which is
\'etale over $S$.

\medskip
\noindent
\SGARefTarget{sga7.sga7i.eix.n13-1-5}{13.1.5}. When $A$ is an ``\underline{abeloid}'' group over $S$, by which
we mean that it is smooth, proper, and separated over $S$, with connected
fibres---so that in particular $A=A^o$---then $R$ is a locally constant
group over $S$, whose rank at each point $s$ (as a $\mathbb Z$-module) is
equal to $2\operatorname{rank}_s(E)$, and conversely. In this way one
obtains the familiar interpretation of an abeloid space over $S$ in
terms of a ``\underline{local system}'' $R$ over $S$ and a
\underline{complex} analytic vector-bundle structure on the real
analytic bundle $R\otimes_{\mathbb Z}\mathbb R$ defined by $R$.

\SGAvar{\medskip
\noindent
\SGARefTarget{sga7.sga7i.eix.n13-2}{13.2}. \underline{Formal completion along a fibre.} We now fix a point
$s\in S$, and shall work locally in a neighbourhood of $s$. For every
integer $n\geqslant0$, let $S_n$ denote the $n$-th infinitesimal
neighbourhood of $s$ in $S$, and, for every analytic space $X$ over $S$,
put likewise $X_n=X\times_S S_n$. As $n$ varies, the system of the $X_n$
then forms an inductive system of analytic spaces over $S$, which we
denote by $\widehat X$ and call the \underline{formal completion of $X$
along $X_s$}. It clearly depends functorially on $X$, and its formation
commutes with finite projective limits; hence $X\longmapsto\widehat X$
takes groups to groups. Notice that $\widehat X$ in fact depends only on
the restriction of $X$ to an arbitrary open neighbourhood $U$ of $S$ in
$S$.

}{\medskip
\noindent
\SGARefTarget{sga7.sga7i.eix.n13-2}{13.2}. \underline{Formal completion along a fibre.} We now fix a point
$s\in S$, and shall work locally in a neighbourhood of $s$. For every
integer $n\geqslant0$, let $S_n$ denote the $n$-th infinitesimal
neighbourhood of $s$ in $S$, and, for every analytic space $X$ over $S$,
put likewise $X_n=X\times_S S_n$. As $n$ varies, the system of the $X_n$
then forms an inductive system of analytic spaces over $S$, which we
denote by $\widehat X$ and call the \underline{formal completion of $X$
along $X_s$}. It clearly depends functorially on $X$, and its formation
commutes with finite projective limits; hence $X\longmapsto\widehat X$
takes groups to groups. Notice that $\widehat X$ in fact depends only on
the restriction of $X$ to an arbitrary open neighbourhood $U$ of $s$ in
$S$.

}{E7I-IX-109}\SGApage{I}{487}
\medskip
\noindent
\underline{Lemma \SGARefTarget{sga7.sga7i.eix.n13-2-1}{13.2.1}.} \underline{The data are those of \SGARefLink{sga7.sga7i.eix.n13-1}{13.1}, and
suppose that the following conditions hold:} a) $A=A^o$;
b) $R$ \underline{is constant in a neighbourhood of $s$}; and
c) $R_s$ \underline{spans the complex vector space $E_s$}.
\underline{Let $A'$ be a second smooth analytic group over $S$, and
consider the map}

\begin{equation*}
 u\longmapsto\widehat u:
 \operatorname{Hom}_{S\text{-}\mathrm{gr}}(A,A')
 \longrightarrow
 \operatorname{Hom}_{\widehat S\text{-}\mathrm{gr}}
   (\widehat A,\widehat {A}')\quad ,
\end{equation*}

\noindent
\underline{where the notation is that of \SGARefLink{sga7.sga7i.eix.n13-2}{13.2}, and the analogous map
obtained by restricting over a variable open neighbourhood $U$ of $s$
and passing to the inductive limit:}

\begin{equation}\tag{\SGARefTarget{sga7.sga7i.eix.n13-2-1-1}{13.2.1.1}}
 u\longmapsto\widehat u:
 \operatorname{Hom}_{(S,s)\text{-}\mathrm{gr}}(A,A')
 \longrightarrow
 \operatorname{Hom}_{\widehat S\text{-}\mathrm{gr}}
   (\widehat A,\widehat {A}')\quad .
\end{equation}

\noindent
(\underline{Here $(S,s)$ denotes the germ of the analytic space $S$ at
$s$.}) \underline{Then the latter map is bijective.}

Indeed, consider the exact sequence

\SGAvar{\begin{equation*}
 0\longrightarrow R'\longrightarrow E'\longrightarrow {A'}^o
 \longrightarrow C
\end{equation*}

}{\begin{equation*}
 0\longrightarrow R'\longrightarrow E'\longrightarrow {A'}^o
 \longrightarrow0
\end{equation*}

}{E7I-IX-110}\noindent
analogous to (\SGARefLink{sga7.sga7i.eix.n13-1-3}{13.1.3}) and defined by $A'$. Since the homomorphisms
$A\longrightarrow A'$ necessarily factor through ${A'}^o$, and the same
holds after the base changes $U\longrightarrow S$ and
$S_n\longrightarrow S$, we may assume that $A'={A'}^o$. By the
dictionary of \SGARefLink{sga7.sga7i.eix.n13-1-4}{13.1.4}, it remains to prove that giving a germ of a vector
bundle homomorphism $E\longrightarrow E'$ which takes $R$ into $R'$ is
equivalent to giving a coherent system of homomorphisms
$E_n\longrightarrow E'_n$ which take $R_n$ into $R'_n$.

\SGApage{I}{\(\approx\)488}
If $\underline{\mathcal O}_s$ denotes the local ring of $S$ at $s$,
giving a coherent system of homomorphisms
$E_n\longrightarrow E'_n$ is equivalent to giving a homomorphism of
$\widehat{\underline{\mathcal O}}_s$-modules

\begin{equation*}
 \widehat F:\widehat{\underline E}_s
 \longrightarrow\widehat{\underline E}'_s,
\end{equation*}

\SGAvar{\noindent
where $\underline E_s$ and $ E'_s$ are the free
$\underline{\mathcal O}_s$-modules obtained at $s$ from the
$\underline{\mathcal O}_s$-modules $\underline E$ and $\underline E'$,
corresponding to $E$ and $E'$. On the other hand, a germ of a vector
bundle homomorphism $E\longrightarrow E'$ is equivalent to giving a
homomorphism of $\underline{\mathcal O}_s$-modules

}{\noindent
where $\underline E_s$ and $\underline E'_s$ are the free
$\underline{\mathcal O}_s$-modules obtained at $s$ from the
$\underline{\mathcal O}_S$-modules $\underline E$ and $\underline E'$,
corresponding to $E$ and $E'$. On the other hand, a germ of a vector
bundle homomorphism $E\longrightarrow E'$ is equivalent to giving a
homomorphism of $\underline{\mathcal O}_s$-modules

}{E7I-IX-111}\begin{equation*}
 F:\underline E_s\longrightarrow\underline E'_s\quad .
\end{equation*}

\noindent
This already shows that (\SGARefLink{sga7.sga7i.eix.n13-2-1-1}{13.2.1.1}) is injective. For surjectivity, it is
enough to prove that if $\widehat F$ is a homomorphism as above such that,
for every $n$, the induced homomorphism
$F_n:E_n\longrightarrow E'_n$ takes $R_n$ into $R'_n$, then
$\widehat F$ comes from an $F$---that is, it takes $\underline E_s$ into
$\underline E'_s$---and the corresponding germ of homomorphisms
$E\longrightarrow E'$ takes $R$ into $R'$.

\SGAvar{By hypothesis c), $R_s=R_o$ spans $E_s$; it follows from Nakayama's
lemma that the group $R_o$ (identified either with the group of sections
of $\widehat R$ over $\widehat S$, or with the group of germs of sections
of $R$ over $S$) generates $\widehat{ E}_s$. Thus, to prove that
$\widehat F(\underline E_s)\subset\underline E'_s$, it is enough to prove
that $\widehat F(R_o)\subset\underline E'_s$. But
$\widehat F(R_o)\subset R'_o$, where $R'_o$ is itself interpreted as the
group of germs of sections of $R'\subset E'$ in a neighbourhood of $s$.
This gives the asserted inclusion, and hence the existence of $F$.
Moreover, we see that $F(R_o)\subset R'_o$, where $R_o$ and $R'_o$ are
regarded as groups of germs of sections of $E$ and $E'$, respectively.
By b), $R_o$ generates $R$ in a neighbourhood of $s$; it follows that
$F(R)\subset R'$ in a neighbourhood of $s$, which completes the proof.

\medskip
\noindent
\underline{Corollary \SGARefTarget{sga7.sga7i.eix.n13-2-2}{13.2.2}.} \underline{If $A$ satisfies conditions
a), b), and c) of \SGARefLink{sga7.sga7i.eix.n13-2-1}{13.2.2}, it is determined up to unique isomorphism in a
neighbourhood of $s$ by the knowledge of the $\widehat S$-group
$\widehat A$.}


}{By hypothesis c), $R_s=R_o$ spans $E_s$; it follows from Nakayama's
lemma that the group $R_o$ (identified either with the group of sections
of $\widehat R$ over $\widehat S$, or with the group of germs of sections
of $R$ over $S$) generates $\widehat{\underline E}_s$. Thus, to prove that
$\widehat F(\underline E_s)\subset\underline E'_s$, it is enough to prove
that $\widehat F(R_o)\subset\underline E'_s$. But
$\widehat F(R_o)\subset R'_o$, where $R'_o$ is itself interpreted as the
group of germs of sections of $R'\subset E'$ in a neighbourhood of $s$.
This gives the asserted inclusion, and hence the existence of $F$.
Moreover, we see that $F(R_o)\subset R'_o$, where $R_o$ and $R'_o$ are
regarded as groups of germs of sections of $E$ and $E'$, respectively.
By b), $R_o$ generates $R$ in a neighbourhood of $s$; it follows that
$F(R)\subset R'$ in a neighbourhood of $s$, which completes the proof.

\medskip
\noindent
\underline{Corollary \SGARefTarget{sga7.sga7i.eix.n13-2-2}{13.2.2}.} \underline{If $A$ satisfies conditions
a), b), and c) of \SGARefLink{sga7.sga7i.eix.n13-2-1}{13.2.1}, it is determined up to unique isomorphism in a
neighbourhood of $s$ by the knowledge of the $\widehat S$-group
$\widehat A$.}


}{E7I-IX-112}\SGApage{I}{\(\approx\)489}
\medskip
\noindent
\SGARefTarget{sga7.sga7i.eix.n13-3}{13.3}. \underline{The analytic Raynaud extension}.

\SGAvar{We are now in a position to paraphrase, in analytic geometry, the
construction of the Raynaud extension given in 7. Suppose that $A$ is
as in \SGARefLink{sga7.sga7i.eix.n13-1}{13.1}, with $A=A_o$,

and consider $R_o=R_3$, which is a free $\mathbb Z$-module of finite
type. After restricting $S$ to a sufficiently small open neighbourhood
of $s$, and interpreting $R_o$ as the group of germs of sections of $R$
near $s$, one obtains a homomorphism

}{We are now in a position to paraphrase, in analytic geometry, the
construction of the Raynaud extension given in 7. Suppose that $A$ is
as in \SGARefLink{sga7.sga7i.eix.n13-1}{13.1}, with $A=A^o$,

and consider $R_o=R_s$, which is a free $\mathbb Z$-module of finite
type. After restricting $S$ to a sufficiently small open neighbourhood
of $s$, and interpreting $R_o$ as the group of germs of sections of $R$
near $s$, one obtains a homomorphism

}{E7I-IX-113}\begin{equation}\tag{\SGARefTarget{sga7.sga7i.eix.n13-3-1}{13.3.1}}
 R_{oS}\longrightarrow R\subset E,
\end{equation}

\noindent
where $R_{oS}$ denotes the constant analytic group over $S$ with value
$R_o$. On choosing the neighbourhood of $s$ small enough that every
section of $\underline{\mathcal O}_S$ whose germ at $s$ is zero is itself
zero, one sees that the preceding homomorphism is injective. Its image is
therefore a constant analytic subgroup of $R$; it is plainly the largest
constant subgroup of $R$, and is called the \underline{fixed part} of
$R$ and denoted by $R^f$:

\begin{equation}\tag{\SGARefTarget{sga7.sga7i.eix.n13-3-2}{13.3.2}}
 R_{oS}\simeq R^f\hookrightarrow R,
 \qquad\text{inducing an isomorphism }R_s^f\simeq R_s.
\end{equation}

\noindent
We may then consider the quotient group

\begin{equation}\tag{\SGARefTarget{sga7.sga7i.eix.n13-3-3}{13.3.3}}
 A^\natural=E/R^f,
\end{equation}

\SGApage{I}{\(\approx\)490}
\noindent
and, in view of (\SGARefLink{sga7.sga7i.eix.n13-1-3}{13.1.3}), obtain a canonical exact sequence

\begin{equation}\tag{\SGARefTarget{sga7.sga7i.eix.n13-3-4}{13.3.4}}
 0\longrightarrow R/R^f\stackrel{i}{\longrightarrow}A^\natural
 \longrightarrow A\longrightarrow0.
\end{equation}

\noindent
Notice that, after restricting $S$ to a neighbourhood of $s$, $R^f$ is
always closed in $E$, and hence $A^\natural$ is separated over $S$.
Notice also that the homomorphism $A^\natural\longrightarrow A$ is
\'etale, and therefore induces an isomorphism

\begin{equation}\tag{\SGARefTarget{sga7.sga7i.eix.n13-3-5}{13.3.5}}
 \underline{\operatorname{Lie}}(A^\natural)
 \stackrel{\sim}{\longrightarrow}
 \underline{\operatorname{Lie}}(A)=\underline E.
\end{equation}

\noindent
Moreover, for every $n\geqslant0$, the canonical homomorphism
$A^\natural\longrightarrow A$ induces an isomorphism
$A_n^\natural\simeq A_n$ (since $R_n^f\stackrel{\sim}{\longrightarrow}
R_n$), and consequently an isomorphism

\begin{equation}\tag{\SGARefTarget{sga7.sga7i.eix.n13-3-6}{13.3.6}}
 \widehat{A^\natural}\stackrel{\sim}{\longrightarrow}\widehat A.
\end{equation}

\noindent
This last fact shows that $A^\natural$ does indeed play the role of the
group, likewise denoted by $A^\natural$, that we considered in 7. Here
it has been constructed without any hypothesis on the fibre
$A_s=A_o$ of the form ``an extension of an abeloid group by a torus''.
Such a hypothesis will nevertheless be necessary if we wish to recover
on $A^\natural$ the important extension structure (\SGARefLink{sga7.sga7i.eix.n7-2-3}{7.2.3}). By contrast,
the homomorphism $A^\natural\longrightarrow A$ in (\SGARefLink{sga7.sga7i.eix.n13-3-4}{13.3.4}) is essentially
transcendental in nature and has no counterpart in the purely algebraic
theory of \S~7. In \S~14 we shall return to the question of a purely
algebraic interpretation of the homomorphism
$i:R/R^f\longrightarrow A^\natural$ in (\SGARefLink{sga7.sga7i.eix.n13-3-4}{13.3.4}).

Suppose first that we are given a complex subtorus

\begin{equation}\tag{\SGARefTarget{sga7.sga7i.eix.n13-3-7}{13.3.7}}
 T_o\subset A_o,
\end{equation}

\noindent
where by a ``\underline{complex torus}'' we mean an analytic group
isomorphic to a Cartesian power of $\mathbb C^*$. Let

\SGApage{I}{\(\approx\)491}
\begin{equation}\tag{\SGARefTarget{sga7.sga7i.eix.n13-3-8}{13.3.8}}
 T=T_o\times S
\end{equation}

\noindent
be the constant torus over $S$ with value $T_o$. Proceeding as in \SGARefLink{sga7.sga7i.eix.n5-1}{5.1},
one obtains a unique homomorphism (indeed, an immersion) of inductive
systems of analytic groups

\begin{equation*}
 \widehat T\longrightarrow\widehat{A^\natural},
\end{equation*}

\noindent
which over $S_o$ reduces to (\SGARefLink{sga7.sga7i.eix.n13-3-7}{13.3.7}). By \SGARefLink{sga7.sga7i.eix.n13-2-2}{13.2.2}, this homomorphism comes
from a unique germ of homomorphisms

\begin{equation}\tag{\SGARefTarget{sga7.sga7i.eix.n13-3-9}{13.3.9}}
 T\longrightarrow A^\natural,
\end{equation}

\noindent
and, after taking $S$ small enough around $s$, we may suppose that it is
defined over all of $S$. The corresponding homomorphism

\begin{equation*}
 \underline{\operatorname{Lie}}(T)\longrightarrow
 \underline{\operatorname{Lie}}(A^\natural)\simeq\underline E
\end{equation*}

\SGAvar{\noindent
is injective over each $S_n$, hence over $\widehat S$, and therefore over
$\underline{\mathcal O}_s$ by flat descent. Thus there is a neighbourhood
of $S$ in $S$, which we may take to be all of $S$, on which
$\underline{\operatorname{Lie}}(T)$ identifies with a locally direct
summand $\underline E^t$ of $\underline E$; denote the corresponding
vector subbundle of $E$ by $E^t$:

}{\noindent
is injective over each $S_n$, hence over $\widehat S$, and therefore over
$\underline{\mathcal O}_s$ by flat descent. Thus there is a neighbourhood
of $s$ in $S$, which we may take to be all of $S$, on which
$\underline{\operatorname{Lie}}(T)$ identifies with a locally direct
summand $\underline E^t$ of $\underline E$; denote the corresponding
vector subbundle of $E$ by $E^t$:

}{E7I-IX-114}\begin{equation}\tag{\SGARefTarget{sga7.sga7i.eix.n13-3-10}{13.3.10}}
 \underline{\operatorname{Lie}}(T)\xrightarrow{\ \sim\ }
 \underline E^t\subset\underline E\simeq
 \underline{\operatorname{Lie}}(A^\natural)\simeq
 \underline{\operatorname{Lie}}(A),
 \qquad E^t\subset E.
\end{equation}

\noindent
Likewise, let

\begin{equation}\tag{\SGARefTarget{sga7.sga7i.eix.n13-3-11}{13.3.11}}
 R_o^t\subset R_o,
 \qquad R^t=R_{oS}^t\subset R_{oS}\simeq R^f
\end{equation}

\noindent
denote respectively the sublattice of $R_o$ defined by the subtorus
$T_o$ of $A_o$, and the corresponding constant analytic group over $S$.
Then the homomorphism (\SGARefLink{sga7.sga7i.eix.n13-3-9}{13.3.9}) corresponds simply to the homomorphism

\SGApage{I}{\(\approx\)492}
\begin{equation*}
 E^t/R_{oS}^t=E^t/R^t\longrightarrow E/R_{oS}=E/R^f
\end{equation*}

\noindent
induced by the inclusions

\begin{equation*}
 E^t\hookrightarrow E,
 \qquad R^t\hookrightarrow R^f.
\end{equation*}

\noindent
Using the fact that (\SGARefLink{sga7.sga7i.eix.n13-3-9}{13.3.9}) induces an immersion on the fibre over $s$,
we conclude that the homomorphism

\begin{equation*}
 R_o/R_o^t\longrightarrow E_o/E_o^t=(E/E^t)_o
\end{equation*}

\noindent
is injective. It follows at once that

\begin{equation*}
 R_o/R_o^t\longrightarrow(E/E^t)_{s'}
\end{equation*}

\noindent
is still injective for $s'$ near $s$. Thus, after shrinking $S$ around
$s$, we may suppose that (\SGARefLink{sga7.sga7i.eix.n13-3-9}{13.3.9}) \underline{is a closed immersion}.

We are therefore led to introduce

\begin{equation}\tag{\SGARefTarget{sga7.sga7i.eix.n13-3-12}{13.3.12}}
 B=A^\natural/T\simeq(E/E^t)/(R^f/R^t),
\end{equation}

\noindent
so that $A^\natural$ acquires an extension structure

\begin{equation}\tag{\SGARefTarget{sga7.sga7i.eix.n13-3-13}{13.3.13}}
 0\longrightarrow T\longrightarrow A^\natural
 \longrightarrow B\longrightarrow0,
\end{equation}

\noindent
where $B$ is a relative analytic group over $S$, smooth and separated
over $S$, with connected fibres. By construction we therefore have

\begin{equation}\tag{\SGARefTarget{sga7.sga7i.eix.n13-3-14}{13.3.14}}
 \underline{\operatorname{Lie}}(B)\simeq
 \underline E/\underline E^t,
\end{equation}

\noindent
and the exponential exact sequence for $B$ is

\SGApage{I}{\(\approx\)493}
\begin{equation}\tag{\SGARefTarget{sga7.sga7i.eix.n13-3-15}{13.3.15}}
 0\longrightarrow R^f/R^t\longrightarrow E/E^t
 \longrightarrow B\longrightarrow0.
\end{equation}

\noindent
Notice that

\begin{equation}\tag{\SGARefTarget{sga7.sga7i.eix.n13-3-16}{13.3.16}}
 R^f/R^t=(R_o/R_o^t)_S
\end{equation}

\noindent
is a constant analytic group over $S$. It follows from \SGARefLink{sga7.sga7i.eix.n13-1-5}{13.1.5} that $B$
is an abeloid group over $S$ in a neighbourhood of $s$ if and only if
its fibre at $s$,

\begin{equation*}
 B_o=A_o/T_o,
\end{equation*}

\noindent
is an abeloid analytic group; in that case $A_o$ appears as an extension
of such a group $B_o$ by a torus $T_o$. \underline{Every way of writing
$A_o$ as an extension of an abeloid group by a complex torus} $(*)$
\underline{therefore gives rise on $A$, in a neighbourhood of} $s$, \underline{to a unique
extension structure} (\SGARefLink{sga7.sga7i.eix.n13-3-13}{13.3.13}) \underline{of an abeloid group $B$ by a
torus $T$, inducing the given extension structure at $s$}.


\medskip
\noindent
\SGARefTarget{sga7.sga7i.eix.n13-4}{13.4}. \underline{Interpretation in terms of Deligne's} ``\underline{mixed Hodge
structures}''.

\SGAvar{Let us first return to the dictionary of \SGARefLink{sga7.sga7i.eix.n13-1-4}{13.1.4}, which we shall
reformulate in a slightly different way, better suited to introducing
the ``Hodge structures'' relevant to the present situation.  For this
purpose, let $\mathbb R$ be the sheaf of sections of $R$ over $S$, and
introduce the $\underline{\mathcal O}_S$-module

}{Let us first return to the dictionary of \SGARefLink{sga7.sga7i.eix.n13-1-4}{13.1.4}, which we shall
reformulate in a slightly different way, better suited to introducing
the ``Hodge structures'' relevant to the present situation.  For this
purpose, let $\underline R$ be the sheaf of sections of $R$ over $S$, and
introduce the $\underline{\mathcal O}_S$-module

}{E7I-IX-115}\begin{equation}\tag{\SGARefTarget{sga7.sga7i.eix.n13-4-1}{13.4.1}}
 \mathcal E(A)=\mathcal E=
 \underline R\otimes_{\mathbb Z}\underline{\mathcal O}_S\qquad .
\end{equation}

\noindent
It is locally free of finite type when $R$ is locally constant.  (Note
that $\mathcal E$ need not be coherent when $R$ is not locally
constant.)  The homomorphism (\SGARefLink{sga7.sga7i.eix.n13-1-3}{13.1.3})
$\underline R\longrightarrow\underline E$ then defines a canonical
homomorphism

\SGApage{I}{\(\approx\)494}
\begin{equation}\tag{\SGARefTarget{sga7.sga7i.eix.n13-4-2}{13.4.2}}
 \mathcal E=\underline R\otimes_{\mathbb Z}
 \underline{\mathcal O}_S\longrightarrow\underline E\qquad .
\end{equation}

\noindent
We denote its kernel by $\underline F(A)$, or simply by $\underline F$;
thus there is an exact sequence

\begin{equation}\tag{\SGARefTarget{sga7.sga7i.eix.n13-4-3}{13.4.3}}
 0\longrightarrow\underline F\longrightarrow\mathcal E
 \longrightarrow\underline E\qquad .
\end{equation}

\noindent
The most interesting case is that in which (\SGARefLink{sga7.sga7i.eix.n13-4-2}{13.4.2}) is an epimorphism,
that is, in which $R_s$ spans the complex vector space $E_s$ for every
$s\in S$.  We then have an exact sequence

\begin{equation}\tag{\SGARefTarget{sga7.sga7i.eix.n13-4-4}{13.4.4}}
 0\longrightarrow\underline F\longrightarrow\mathcal E
 \longrightarrow\underline E\longrightarrow0\qquad .
\end{equation}

\noindent
If $R$ is locally constant, this is a short exact sequence of locally
free modules over $S$.  Knowledge of the pair $(E,R)$ is then equivalent
to knowledge of the pair $(R,\underline F)$, where $R$ is a locally
constant analytic group over $S$ whose fibres are free $\mathbb Z$-modules
of finite type, and $\underline F$ is a locally direct-summand submodule
of $\mathcal E=\underline R\otimes_{\mathbb Z}
\underline{\mathcal O}_S$, subject in addition to

\begin{equation}\tag{\SGARefTarget{sga7.sga7i.eix.n13-4-5}{13.4.5}}
 \underline F\cap\underline R=0\qquad .
\end{equation}

\noindent
From the standpoint of Hodge structures, $\mathcal E$ should be regarded
as filtered by the ``\underline{Hodge filtration}''

\SGAvar{\begin{equation}\tag{\SGARefTarget{sga7.sga7i.eix.n13-4-6}{13.4.6}}
 \operatorname{Fil}^{1}(\mathcal E)=\mathcal E,
 \qquad \operatorname{Fil}^{0}(\mathcal E)=\underline F,
 \qquad \operatorname{Fil}^{-1}(\mathcal E)=0\qquad .
\end{equation}

}{\begin{equation}\tag{\SGARefTarget{sga7.sga7i.eix.n13-4-6}{13.4.6}}
 \operatorname{Fil}^{-1}(\mathcal E)=\mathcal E,
 \qquad \operatorname{Fil}^{0}(\mathcal E)=\underline F,
 \qquad \operatorname{Fil}^{1}(\mathcal E)=0\qquad .
\end{equation}

}{E7I-IX-116}\noindent
Notice that the exact sequence (\SGARefLink{sga7.sga7i.eix.n13-4-3}{13.4.3}) still depends functorially on the
analytic group $A$ over $S$, and commutes with arbitrary base change.

\SGApage{I}{\(\approx\)495}
Returning now to the conditions of \SGARefLink{sga7.sga7i.eix.n13-3}{13.3}, put

\begin{equation}\tag{\SGARefTarget{sga7.sga7i.eix.n13-4-7}{13.4.7}}
 \mathcal E^f=\underline R^f\otimes_{\mathbb Z}
 \underline{\mathcal O}_S,
 \qquad
 \mathcal E^t=\underline R^t\otimes_{\mathbb Z}
 \underline{\mathcal O}_S\qquad .
\end{equation}

\noindent
These modules are defined on a suitable neighbourhood of $s$, which we
may suppose to be all of $S$, and define a filtration of $\mathcal E$,
called the \underline{weight filtration},

\begin{equation}\tag{\SGARefTarget{sga7.sga7i.eix.n13-4-8}{13.4.8}}
 \mathcal E\supset\mathcal E^f\supset\mathcal E^t\supset0\qquad .
\end{equation}

\SGAvar{\noindent
Here $\mathcal E^t$, $\mathcal E^f$, and $\mathcal E$ are to be regarded
as having weights $-2$, $-1$, and $0$, respectively; thus this is an
increasing filtration.  Like $R^t$, the module $\mathcal E^t$ depends on
the choice of a subtorus $T_o$ of $A_o=\Lambda_s$.  Since $R^f$ and $R^t$ are
locally constant, \underline{$\mathcal E^f$ and $\mathcal E^t$ are
locally free}; they may also be interpreted as

}{\noindent
Here $\mathcal E^t$, $\mathcal E^f$, and $\mathcal E$ are to be regarded
as having weights $-2$, $-1$, and $0$, respectively; thus this is an
increasing filtration.  Like $R^t$, the module $\mathcal E^t$ depends on
the choice of a subtorus $T_o$ of $A_o=A_s$.  Since $R^f$ and $R^t$ are
locally constant, \underline{$\mathcal E^f$ and $\mathcal E^t$ are
locally free}; they may also be interpreted as

}{E7I-IX-117}\begin{equation}\tag{\SGARefTarget{sga7.sga7i.eix.n13-4-9}{13.4.9}}
 \mathcal E^f=\mathcal E(A^\natural),
 \qquad \mathcal E^t=\mathcal E(T)\qquad .
\end{equation}

\noindent
Applying the functoriality of (\SGARefLink{sga7.sga7i.eix.n13-4-3}{13.4.3}), as $A$ varies, to the
homomorphisms

\begin{equation*}
 T\longrightarrow A^\natural\longrightarrow A
\end{equation*}

\noindent
shows at once that $\underline F(A^\natural)$ and $\underline F(T)$
identify respectively with

\begin{equation}\tag{\SGARefTarget{sga7.sga7i.eix.n13-4-10}{13.4.10}}
 \underline F(A^\natural)\simeq
 \underline F^f\stackrel{\mathrm{def}}{=}
 \underline F\cap\mathcal E^f,
 \qquad
 \underline F(T)=\underline F^t=
 \underline F\cap\mathcal E^t=0\qquad .
\end{equation}

\noindent
The last equality follows immediately from the fact that
$\underline F(T)=0$ for a torus $T$: in this case the exact sequence of
type (\SGARefLink{sga7.sga7i.eix.n13-4-3}{13.4.3}) is
$0\longrightarrow0\longrightarrow\mathcal E(T)
\longrightarrow\underline{\operatorname{Lie}}(T)\longrightarrow0$.
At the same time we see that the composite homomorphism

\begin{equation*}
 \mathcal E^t\longrightarrow\mathcal E\longrightarrow\underline E
\end{equation*}

\noindent
induces an isomorphism

\SGApage{I}{\(\approx\)496}
\begin{equation}\tag{\SGARefTarget{sga7.sga7i.eix.n13-4-11}{13.4.11}}
 \mathcal E^t\stackrel{\sim}{\longrightarrow}\underline E^t
 \simeq\underline{\operatorname{Lie}}(T)\qquad .
\end{equation}

\noindent
\underline{Suppose now that $R_s$ spans the complex vector space $E_s$},
as happens in particular in the favourable case \SGARefLink{sga7.sga7i.eix.n13-3-16}{13.3.16}, where $A_s$ is
an extension of an abeloid group by a complex torus $T_o$.  After
restricting $S$ to a suitable neighbourhood of $s$, the canonical
homomorphism of type (\SGARefLink{sga7.sga7i.eix.n13-4-2}{13.4.2}), with $A$ replaced by $A^\natural$,

\begin{equation*}
 \mathcal E^f=\mathcal E(A^\natural)
 \longrightarrow\underline{\operatorname{Lie}}(A^\natural)
 \simeq\underline E
\end{equation*}

\noindent
is an epimorphism; a fortiori, so is (\SGARefLink{sga7.sga7i.eix.n13-4-2}{13.4.2}).  We therefore have the
exact sequences (\SGARefLink{sga7.sga7i.eix.n13-4-4}{13.4.4}) and

\begin{equation}\tag{\SGARefTarget{sga7.sga7i.eix.n13-4-12}{13.4.12}}
 0\longrightarrow\underline F^f\longrightarrow\mathcal E^f
 \longrightarrow\underline E\longrightarrow0\qquad .
\end{equation}

\noindent
In particular,

\begin{equation}\tag{\SGARefTarget{sga7.sga7i.eix.n13-4-13}{13.4.13}}
 \underline F+\mathcal E^f=\mathcal E\qquad ,
\end{equation}

\SGAvar{\noindent
since $\mathcal E/\underline F\simeq\underline E$.  Moreover, since
$B=A^\natural/T$ is a quotient of $A$,

}{\noindent
since $\mathcal E/\underline F\simeq\underline E$.  Moreover, since
$B=A^\natural/T$ is a quotient of $A^\natural$,

}{E7I-IX-118}it satisfies the analogous surjectivity condition, and hence gives an
exact sequence of type (\SGARefLink{sga7.sga7i.eix.n13-4-4}{13.4.4}):

\SGAvar{\begin{equation}\tag{\SGARefTarget{sga7.sga7i.eix.n13-4-14}{13.4.14}}
 0\longrightarrow\underline F(B)\longrightarrow(B)
 \longrightarrow\underline{\operatorname{Lie}}(B)
 \longrightarrow0\qquad .
\end{equation}

}{\begin{equation}\tag{\SGARefTarget{sga7.sga7i.eix.n13-4-14}{13.4.14}}
 0\longrightarrow\underline F(B)\longrightarrow\mathcal E(B)
 \longrightarrow\underline{\operatorname{Lie}}(B)
 \longrightarrow0\qquad .
\end{equation}

}{E7I-IX-119}\noindent
Its interpretation in terms of the modules already introduced is
immediate.  Indeed, by (\SGARefLink{sga7.sga7i.eix.n13-3-12}{13.3.12}),

\begin{equation}\tag{\SGARefTarget{sga7.sga7i.eix.n13-4-15}{13.4.15}}
 \left\{
 \begin{array}{l}
  \mathcal E(B)\simeq
  (\underline R^f/\underline R^t)\otimes_{\mathbb Z}
  \underline{\mathcal O}_S
  \simeq\mathcal E^f/\mathcal E^t,\\[2ex]
  \underline{\operatorname{Lie}}(B)
  \simeq\underline E/\underline E^t.
 \end{array}
 \right.
\end{equation}

\SGApage{I}{\(\approx\)497}
\noindent
It follows from (\SGARefLink{sga7.sga7i.eix.n13-4-11}{13.4.11}) and the relation
$\underline F\cap\mathcal E^t=0$ in (\SGARefLink{sga7.sga7i.eix.n13-4-10}{13.4.10}) that
$\underline F(B)=\ker(\mathcal E^f/\mathcal E^t
\longrightarrow\underline E/\underline E^t)$ is isomorphic to
$\underline F^f$ through the homomorphism

\begin{equation}\tag{\SGARefTarget{sga7.sga7i.eix.n13-4-16}{13.4.16}}
 \underline F^f\simeq\underline F(A^\natural)
 \stackrel{\sim}{\longrightarrow}\underline F(B)\qquad .
\end{equation}

\noindent
Let us summarize the properties just obtained for the weight filtration
(\SGARefLink{sga7.sga7i.eix.n13-4-8}{13.4.8}), and its relation with the Hodge filtration (\SGARefLink{sga7.sga7i.eix.n13-4-6}{13.4.6}) of
$\mathcal E$ defined by the submodule $\underline F$, in the favourable
case \SGARefLink{sga7.sga7i.eix.n13-3-16}{13.3.16} where $A_s$ is an extension of an abeloid analytic group
$B_o$ by a torus $T_o$.  We use this torus to define $R^t$, so that $B$
is an abeloid analytic group over $S$.

\medskip
\noindent
\underline{Proposition \SGARefTarget{sga7.sga7i.eix.n13-4-17}{13.4.17}.} \underline{Under the preceding
conditions, the Hodge and weight filtrations on}
$\mathcal E=\underline R\otimes_{\mathbb Z}\underline{\mathcal O}_S$,

\begin{equation}\tag{\SGARefTarget{sga7.sga7i.eix.n13-4-17-1}{13.4.17.1}}
 \mathcal E\supset\underline F,
 \qquad
 \mathcal E\supset\mathcal E^f\supset\mathcal E^t\supset0,
\end{equation}

\noindent
\underline{satisfy the following conditions:}

\smallskip
\noindent
a) \underline{The weight filtration is obtained, by tensoring with
$\underline{\mathcal O}_S$ over $\mathbb Z$, from the analogous
filtration on $\underline R$}

\begin{equation*}
 \underline R\supset\underline R^f\supset\underline R^t\supset0,
\end{equation*}

\noindent
\underline{where $\underline R^f$ and $\underline R^t$ are locally
constant analytic groups over} $S$, \underline{with values free $\mathbb Z$-modules
of finite type}, namely
$R_o^f=R_s=\pi_1(A_s)$ and $R_o^t=\pi_1(T_o)$.

\smallskip
\noindent
b) \underline{One has the relations}

\begin{equation}\tag{*}
 \underline F\cap\mathcal E^t=0,
 \qquad \underline F+\mathcal E^f=\mathcal E\qquad .
\end{equation}

\SGApage{I}{\(\approx\)498}
\noindent
\underline{In other words, the filtration on $\mathcal E^t$ induced by
the Hodge filtration}, with the degree conventions of (\SGARefLink{sga7.sga7i.eix.n13-4-6}{13.4.6}),
\underline{is concentrated in degree $-1$}, whereas the filtration on
$\mathcal E/\mathcal E^f$ induced by the Hodge filtration is
\underline{concentrated in degree $0$}.

\smallskip
\noindent
c) \underline{Consider on}

\begin{equation*}
 \mathcal E^{\mathrm{ab}}
 \stackrel{\mathrm{def}}{=}\mathcal E^f/\mathcal E^t
 \simeq\underline R^{\mathrm{ab}}\otimes_{\mathbb Z}
 \underline{\mathcal O}_S,
 \qquad
 \underline R^{\mathrm{ab}}=\underline R^f/\underline R^t
 \quad\text{(cf. a))},
\end{equation*}

\noindent
\underline{the filtration induced by the Hodge filtration on
$\mathcal E$}.  It is therefore defined by the submodule
$\underline F^{\mathrm{ab}}$ of $\mathcal E^{\mathrm{ab}}$ which is the
image of

\begin{equation*}
 \underline F^f\stackrel{\mathrm{def}}{=}
 \underline F\cap\mathcal E^f\longrightarrow\mathcal E^{\mathrm{ab}}.
\end{equation*}

\noindent
The arrow is a monomorphism by the first relation $(*)$.
\underline{Then the pair}

\begin{equation*}
 \bigl(\underline R^{\mathrm{ab}},
 \underline F^{\mathrm{ab}}\subset
 \underline R^{\mathrm{ab}}\otimes_{\mathbb Z}
 \underline{\mathcal O}_S=\mathcal E^{\mathrm{ab}}\bigr)
\end{equation*}

\noindent
\underline{is isomorphic to the pair}
$(\underline R(B),\underline F(B))$
\underline{associated with an abeloid group over $S$}, as explained
after (\SGARefLink{sga7.sga7i.eix.n13-4-4}{13.4.4}).  Thus

\begin{equation*}
 B\simeq\mathcal E^{\mathrm{ab}}/
 (F^{\mathrm{ab}}+R^{\mathrm{ab}}),
\end{equation*}

\noindent
where the symbols without underlining denote the corresponding analytic
groups over $S$.

\medskip
\noindent
\underline{Remark \SGARefTarget{sga7.sga7i.eix.n13-4-18}{13.4.18}.} In Deligne's terminology [4], the conclusion
of the preceding proposition, after restriction to an open subset $U$
of $S$ on which $R$ is locally constant, says that $R$, endowed with the
two filtrations (\SGARefLink{sga7.sga7i.eix.n13-4-17-1}{13.4.17.1}) on
$\mathcal E=\underline R\otimes_{\mathbb Z}\underline{\mathcal O}_S$,
is an ``analytic family of mixed Hodge structures over $U$''.  The
situation just described served as the model for Deligne's conjectures
on N\'eron theory in higher dimensions, stated in [9, 9.8--9.13].

There is, however, an important feature of the algebraic theory studied
in \S~7 which is missing from the description in \SGARefLink{sga7.sga7i.eix.n13-4-17}{13.4.17}.  Suppose that
$A$ is an abeloid group over an open subset $U$ complementary to a
nowhere-dense analytic subset $Z$ of $S$---say $Z=\{s\}$ when
$\dim S=1$---so that, by (\SGARefLink{sga7.sga7i.eix.n13-1-5}{13.1.5}), $R|U$ is a local system over $U$.
Then the local system $R/R^t$ over $U$ ought to be \underline{constant};
more precisely, it ought to be dual to the local system ${R'}^f$, where
$R'$ is the analytic group, \`etale over $S$, which would correspond
(as $R$ corresponds to $A$) to a suitable analytic group $A'$ over $S$
such that $A'|U$ is the abeloid group dual to $A|U$.  In the notation of
\S~5, the algebraic analogue is the duality between
$T_\ell(A_K)/T_\ell(A_K)^t$ and $T_\ell(A'_K)^f$ in (\SGARefLink{sga7.sga7i.eix.n5-5-9}{5.5.9}), which
follows from the duality theorem \SGARefLink{sga7.sga7i.eix.n5-2}{5.2}.

\SGApage{I}{\(\approx\)499}
I do not know whether such a statement holds under only the hypotheses
of \SGARefLink{sga7.sga7i.eix.n13-4-17}{13.4.17}, or even whether the action of $\pi_1(U,t)$ on $R_t$ is
unipotent for $t\in U$.  Nor do the naive methods used in the present
paragraph seem capable of proving such a result, even under the
additional hypothesis that $A|U$ is endowed with a ``polarization'',
that is, with an invertible sheaf ample on every fibre.  This is,
however, something that the ``semi-simple group people'' apparently know
how to prove, at least when $S$ is nonsingular of dimension $1$; see the
report [9].  In the next subsection we shall show that everything works
under relative algebraicity hypotheses on $A$, and at the same time
relate the result to the construction of \S~7.


\SGApage{I}{\(\approx\)500}
\medskip
\noindent
\SGARefTarget{sga7.sga7i.eix.n13-5}{13.5}. \underline{The case of a relatively algebraic situation.}
We now suppose that $A$ comes from a relative group scheme over $S$
which is smooth and separated, and whose fibres are extensions of
abelian schemes by tori.  We denote this relative scheme by the same
letter $A$.  (We allow ourselves this abuse of notation because the
functor $G\longmapsto G^{\mathrm{an}}$, from the category of relative
group schemes of the type under consideration to the category of
analytic groups over $S$, is fully faithful, as would be easy to prove
using results of relative GAGA type [10].)  For everything concerning
the notion of a relative scheme and the relations between relative
schemes over $S$ and analytic spaces over $S$, we refer to [10] (see
also [2, Exp. 15]).

If $\underline O$ denotes the local ring of $S$ at $s$, then the
relative scheme $A$ is determined near $s$ by the ordinary scheme
$A_{\underline O}$ over $\operatorname{Spec}(\underline O)$ obtained
from $A$ by localization at $s$.  It is then clear that the formal
completion of $A_{\underline O}$ along its special fibre (isomorphic to
$A_s=A_o$) is isomorphic to the formal completion of the analytic group
$A$ along $A_s$ (\SGARefLink{sga7.sga7i.eix.n13-2}{13.2}), where, by the abuse of notation introduced
above, the $A_n=A\times_S S_n$ are identified with algebraic schemes
over $S$.  In the favourable case in which $\hat A$, viewed as a formal
group scheme over $\hat{\underline O}$, is algebraizable (\SGARefLink{sga7.sga7i.eix.n7-2-1}{7.2.1}), as is
the case in particular when $S$ is normal at $s$ and $A$ is
quasi-projective over $S$ (\SGARefLink{sga7.sga7i.eix.n7-1-5}{7.1.5}), we therefore obtain a group scheme
$A^{\natural}_{\hat{\underline O}}$ over
$\operatorname{Spec}(\hat{\underline O})$, which is an extension of an
abelian scheme $B_{\hat{\underline O}}$ by a torus.  For simplicity, we
suppose that $B_{\hat{\underline O}}$ is projective over
$\hat{\underline O}$ (as is indeed the case, we know, if
$\underline O$, and hence $\hat{\underline O}$, is normal).  We shall
then prove:

\SGApage{I}{501}
\medskip
\noindent
\underline{Proposition \SGARefTarget{sga7.sga7i.eix.n13-5-1}{13.5.1}.} \underline{Under the preceding
conditions, the analytic group $A^{\natural}$ of (\SGARefLink{sga7.sga7i.eix.n13-3-3}{13.3.3}) comes, in a
neighbourhood of $s$, from a relative group scheme determined up to a
unique isomorphism, which is an extension of a projective relative
abelian scheme over $S$ by a torus.  If
$A^{\natural}_{\underline O}$ denotes the group scheme it defines over
$\underline O=\underline O_{S,s}$, then the group scheme}

\begin{equation*}
 A^{\natural}_{\underline O}\otimes_{\underline O}
 \hat{\underline O}
\end{equation*}

\noindent
\underline{is canonically isomorphic, with its extension structure, to
the Raynaud extension $A^{\natural}_{\hat{\underline O}}$ introduced
above.}

The latter assertion follows immediately from the first, in view of
\SGARefLink{sga7.sga7i.eix.n7-2-1}{7.2.1}, since, by (\SGARefLink{sga7.sga7i.eix.n13-3-6}{13.3.6}), the formal completion of
$A^{\natural}_{\underline O}\otimes_{\underline O}\hat{\underline O}$
along its special fibre is $\hat A_{\underline O}$.  To prove the
algebraicity of $A^{\natural}$ near $s$, proceeding as in \SGARefLink{sga7.sga7i.eix.n7-2-1}{7.2.1}, we are
immediately reduced to proving that of the abelian quotient $B$ of
$A^{\natural}$, which follows from the next lemma.

\SGAvar{\medskip
\noindent
\underline{Lemma \SGARefTarget{sga7.sga7i.eix.n13-5-2}{13.5.2}.} \underline{Let $B$ be an abeloid analytic
group over the analytic space $S$, and let $s$ be a point of $S$ such
that the completion $\hat B$ along the fibre $B_s$ is algebraizable;
that is, the $B_n=B\times_S S_n$ are algebraizable and $\hat B$, now
viewed as a formal scheme over $\hat{\underline C}_{S,s}$, is
algebraizable, or equivalently comes from an abelian scheme
$B_{\hat{\underline O}}$ over $\hat{\underline C}$.  Suppose moreover
that $B_{\hat{\underline O}}$ is projective over
$\hat{\underline O}$.  Then $B$ is projective over $S$ in a
neighbourhood of $s$.}

}{\medskip
\noindent
\underline{Lemma \SGARefTarget{sga7.sga7i.eix.n13-5-2}{13.5.2}.} \underline{Let $B$ be an abeloid analytic
group over the analytic space $S$, and let $s$ be a point of $S$ such
that the completion $\hat B$ along the fibre $B_s$ is algebraizable;
that is, the $B_n=B\times_S S_n$ are algebraizable and $\hat B$, now
viewed as a formal scheme over $\hat{\underline O}_{S,s}$, is
algebraizable, or equivalently comes from an abelian scheme
$B_{\hat{\underline O}}$ over $\hat{\underline O}$.  Suppose moreover
that $B_{\hat{\underline O}}$ is projective over
$\hat{\underline O}$.  Then $B$ is projective over $S$ in a
neighbourhood of $s$.}

}{E7I-IX-120}The hypothesis means that there is a polarization of the formal
analytic space $\hat B$, that is, a sequence of polarizations $p_n$ of
the $B_n$ which ``glue together''.  The conclusion means that there is
a polarization $p$ of $B$ over a neighbourhood of $s$.  Let $P$ be the
analytic space over $S$ of relative polarizations of $B$ over $S$,
which represents the functor

\SGApage{I}{\(\approx\)502}
\begin{equation*}
 S'\longmapsto
 \{\text{polarizations of $B_{S'}$ over $S'$}\}
\end{equation*}

\noindent
on the category of analytic spaces $S'$ over $S$.  This is an open
subfunctor of the N\'eron--Severi functor $\underline{\mathrm{NS}}_{B/S}$
over $S$.  The simplest standard arguments, which we omit here, show
that $\underline{\mathrm{NS}}_{B/S}$, and hence also $P$, is represented
by an \underline{unramified} analytic space over $S$.  The hypothesis
then says that there is a \underline{formal} section $p$ of $P$ over
$S$ at $s$.  Since $P$ is unramified and therefore locally quasi-finite
over $S$, it follows from [2, Exp. 18, no.~1] that every formal section
of $P$ over $S$ at $s$ comes from an analytic section near $s$, that is,
from a polarization of $B$ over a neighbourhood of $s$.  This proves
the assertion.

\medskip
\noindent
\underline{Remarks \SGARefTarget{sga7.sga7i.eix.n13-5-3}{13.5.3}.}

\SGAvar{\smallskip
\noindent
a) If in \SGARefLink{sga7.sga7i.eix.n13-5-2}{13.5.2} we omit the hypothesis that
$B_{\hat{\underline O}}$ is projective over $S$,

}{\smallskip
\noindent
a) If in \SGARefLink{sga7.sga7i.eix.n13-5-2}{13.5.2} we omit the hypothesis that
$B_{\hat{\underline O}}$ is projective over $\hat{\underline O}$,

}{E7I-IX-121}we can still conclude that $B$ is algebraizable near $s$.  For this one
uses an unpublished result (due independently to M.~Artin and
P.~Deligne, using the theory of the sheaves $\underline T^i$ of VI 2)
which asserts that, if $X$ is proper over $S$ and $X_{\mathrm{red}}$ is
algebraizable, then $X$ is algebraizable.  This reduces us to the case
where $S$ is reduced.  On the other hand, when $S$ is normal, so that
$\hat{\underline O}$ is normal, every abelian scheme over
$\hat{\underline O}$ is known to be projective over
$\hat{\underline O}$ (see, for example, [23]), and 13.5.2 applies.  To
obtain the case of an arbitrary reduced $S$, one uses non-flat descent
from the normalization $S'$ of $S$, by means of the following
(delicate) result, which we shall not prove here: a proper morphism of
noetherian schemes $f:S'\longrightarrow S$ which is an epimorphism
(that is, such that
$\underline O_S\longrightarrow f_*(\underline O_{S'})$ is injective)
is an \underline{effective descent morphism} for the fibred category of
relative abelian schemes.

\SGApage{I}{\(\approx\)503}
\smallskip
\noindent
b) It appears that the proof of \SGARefLink{sga7.sga7i.eix.n13-5-2}{13.5.2} should also prove the analogous
result for every analytic space $B$ proper over $S$.  Probably more
delicate is the question whether an analytic space $X$ proper over $S$
whose formal completion $\hat X$ along $X_s$ is algebraizable is itself
algebraizable over $S$ near $s$.  We merely point out that here the
notion of relative algebraizability of $X$ over $S$ should be taken in
a broader sense than in [10]: near $s$, $X$ should be required to come
from an \underline{algebraic space} over $\underline O_{S,s}$ in the
sense of M.~Artin [1].

\SGAvar{\medskip
\noindent
\SGARefTarget{sga7.sga7i.eix.n13-5-4}{13.5.4}. Let us place ourselves under the conditions of \SGARefLink{sga7.sga7i.eix.n13-5-1}{13.5.1}.  For
every integer $n>C$, choose also the standard isomorphism of groups over
$\mathbb C$

}{\medskip
\noindent
\SGARefTarget{sga7.sga7i.eix.n13-5-4}{13.5.4}. Let us place ourselves under the conditions of \SGARefLink{sga7.sga7i.eix.n13-5-1}{13.5.1}.  For
every integer $n>0$, choose also the standard isomorphism of groups over
$\mathbb C$

}{E7I-IX-122}\begin{equation*}
 \mu_{n\mathbb C}\simeq(\mathbb Z/n\mathbb Z)_{\mathbb C},
 \qquad m\longmapsto\exp(2\pi i m/n)\qquad .
\end{equation*}

\noindent
By passing to the inverse limit over the powers of a fixed prime
$\ell$, this gives a canonical isomorphism of $\ell$-adic sheaves over
$\mathbb C$,

\begin{equation*}
 \mathbb Z_\ell(1)_{\mathbb C}\simeq\mathbb Z_{\ell\mathbb C}
 \qquad ,
\end{equation*}

\noindent
which we shall use to identify the two sides.

Applying the exact sequence of
$\underline{\operatorname{Ext}}^i_{\mathbb Z}
(\mathbb Z/n\mathbb Z,-)$ to the exact sequence (\SGARefLink{sga7.sga7i.eix.n13-1-3}{13.1.3}), we obtain a
canonical isomorphism

\SGApage{I}{\(\approx\)504}
\begin{equation}\tag{\SGARefTarget{sga7.sga7i.eix.n13-5-5}{13.5.5}}
 {}_nA^{\mathrm o}\simeq
 \underline R\otimes_{\mathbb Z}\mathbb Z/n\mathbb Z
 \qquad ,
\end{equation}

\noindent
which, as $n$ runs through the powers of a given prime $\ell$, gives an
isomorphism of inverse systems of groups over $S$,

\begin{equation}\tag{\SGARefTarget{sga7.sga7i.eix.n13-5-6}{13.5.6}}
 T_\ell(A^{\mathrm o})\simeq
 \underline R\otimes_{\mathbb Z}\mathbb Z_\ell
 \qquad ,
\end{equation}

\SGAvar{\noindent
where $\mathbb Z_\ell$ is regarded as the inverse system of the
$\mathbb Z/\ell^{n+1}\mathbb Z$.  In the situation considered here, in
which $A$ comes from a separated group scheme over $S$, it is then
immediately clear that the
$\underline R\otimes\mathbb Z/n\mathbb Z$ may be regarded as relative
schemes.  This gives a meaning to the symbol
$(\underline R\otimes_{\mathbb Z}\mathbb Z_\ell)_{\underline O}$
(whereas $ R_{\underline O}$ would in general have none!), and yields a
canonical isomorphism

}{\noindent
where $\mathbb Z_\ell$ is regarded as the inverse system of the
$\mathbb Z/\ell^{n+1}\mathbb Z$.  In the situation considered here, in
which $A$ comes from a separated group scheme over $S$, it is then
immediately clear that the
$\underline R\otimes\mathbb Z/n\mathbb Z$ may be regarded as relative
schemes.  This gives a meaning to the symbol
$(\underline R\otimes_{\mathbb Z}\mathbb Z_\ell)_{\underline O}$
(whereas $\underline R_{\underline O}$ would in general have none!), and yields a
canonical isomorphism

}{E7I-IX-123}\begin{equation}\tag{\SGARefTarget{sga7.sga7i.eix.n13-5-7}{13.5.7}}
 T_\ell(A_{\underline O})\simeq
 (\underline R\otimes_{\mathbb Z}\mathbb Z_\ell)_{\underline O}
 \qquad .
\end{equation}

\noindent
It follows immediately from the constructions of paragraph 2 and \SGARefLink{sga7.sga7i.eix.n13-3}{13.3}
that the preceding isomorphism carries the fixed part to the fixed part
and the toric part to the toric part:

\begin{equation}\tag{\SGARefTarget{sga7.sga7i.eix.n13-5-8}{13.5.8}}
 T_\ell(A_{\underline O})^f\simeq
 (\underline R^f\otimes_{\mathbb Z}\mathbb Z_\ell)_{\underline O},
 \qquad
 T_\ell(A_{\underline O})^t\simeq
 (\underline R^t\otimes_{\mathbb Z}\mathbb Z_\ell)_{\underline O}
 \qquad .
\end{equation}

\noindent
If $A'$ is a relative analytic group over $S$ satisfying the same
hypotheses as $A$, and if $W_{\underline O}$ is a biextension of
$(A_{\underline O},A'_{\underline O})$ by
$\mathbb G_{m\underline O}$, hence defines over a neighbourhood of $s$
a biextension of $(A,A')$ by $\mathbb G_{mS}$ (now using the notion in
the complex-analytic sense, obtained by working in the site of analytic
spaces over $S$ for a suitable topology), then, by VIII \SGARefLink{sga7.sga7i.eviii.n2-2}{2.2}, $W$
defines pairings

\SGApage{I}{\(\approx\)505}
\begin{equation}\tag{\SGARefTarget{sga7.sga7i.eix.n13-5-9}{13.5.9}}
 (\underline R\otimes\mathbb Z_\ell)\times
 (\underline R'\otimes_{\mathbb Z}\mathbb Z_\ell)
 \simeq T_\ell(A)\times T_\ell(A')
 \longrightarrow\mathbb Z_\ell(1)_S
 \simeq\mathbb Z_{\ell S}\qquad ,
\end{equation}

\noindent
and it is again clear that the corresponding pairing

\begin{equation}\tag{\SGARefTarget{sga7.sga7i.eix.n13-5-10}{13.5.10}}
 T_\ell(A_{\underline O})\times T_\ell(A'_{\underline O})
 \longrightarrow\mathbb Z_\ell(1)_S
\end{equation}

\noindent
is precisely the one associated, by loc. cit., with the biextension
$W_{\underline O}$ of $(A_{\underline O},A'_{\underline O})$.  It is
also not difficult to see that, after restricting to a neighbourhood
of $s$, the pairings (\SGARefLink{sga7.sga7i.eix.n13-5-9}{13.5.9}) arise from a unique pairing (no longer
depending on $\ell$)

\begin{equation}\tag{\SGARefTarget{sga7.sga7i.eix.n13-5-11}{13.5.11}}
 \underline R\otimes\underline R'\longrightarrow\mathbb Z
 \qquad .
\end{equation}

\SGAvar{\noindent
With these preliminaries, the algebraic results concerning the
$T_\ell(A_{\underline O})$ can be used to draw conclusions about the
sheaf $\underline R$.  Thus the pairing (\SGARefLink{sga7.sga7i.eix.n13-5-11}{13.5.11}) is zero on
$\underline R^t\times\underline R'{}^f$ and on
$\underline R^f\times\underline R'{}^t$ (\SGARefLink{sga7.sga7i.eix.n7-4-8}{7.4.8}).  Moreover, if $S$ is
nonsingular of dimension $1$ and, over $U=S-\{s\}$, $A$ and $A'$ are
relative abelian schemes dual to one another through $W$, then
$\underline R^t_U$ and $\underline R'{}^f_U$, as well as
$\underline R^f_U$ and $ R'{}^t_U$, are exact annihilators of
one another over a sufficiently small neighbourhood of $s$ (\SGARefLink{sga7.sga7i.eix.n7-4-8}{7.4.8}).
Since (\SGARefLink{sga7.sga7i.eix.n13-5-11}{13.5.11}) is a perfect duality over $U$, it follows that there is
a duality over $U$ between $\underline R/\underline R^t$ and
$\underline R'{}^f$.  In particular,
$\underline R/\underline R^t$ is constant, and a fortiori the action of
$\pi_1(U,t)$ on $\underline R_t$ is unipotent of level $2$.  Similarly,
over $U$ there is a perfect duality between
$\underline R/\underline R^f$ and $\underline R'{}^t$.  Since the
latter is also isomorphic to the dual torus of the lattice $M_S$ over
$S$, where $M$ is the group defined in (\SGARefLink{sga7.sga7i.eix.n5-5-4}{5.5.4}), we obtain a canonical
isomorphism

}{\noindent
With these preliminaries, the algebraic results concerning the
$T_\ell(A_{\underline O})$ can be used to draw conclusions about the
sheaf $\underline R$.  Thus the pairing (\SGARefLink{sga7.sga7i.eix.n13-5-11}{13.5.11}) is zero on
$\underline R^t\times\underline R'{}^f$ and on
$\underline R^f\times\underline R'{}^t$ (\SGARefLink{sga7.sga7i.eix.n7-4-8}{7.4.8}).  Moreover, if $S$ is
nonsingular of dimension $1$ and, over $U=S-\{s\}$, $A$ and $A'$ are
relative abelian schemes dual to one another through $W$, then
$\underline R^t_U$ and $\underline R'{}^f_U$, as well as
$\underline R^f_U$ and $\underline R'{}^t_U$, are exact annihilators of
one another over a sufficiently small neighbourhood of $s$ (\SGARefLink{sga7.sga7i.eix.n7-4-8}{7.4.8}).
Since (\SGARefLink{sga7.sga7i.eix.n13-5-11}{13.5.11}) is a perfect duality over $U$, it follows that there is
a duality over $U$ between $\underline R/\underline R^t$ and
$\underline R'{}^f$.  In particular,
$\underline R/\underline R^t$ is constant, and a fortiori the action of
$\pi_1(U,t)$ on $\underline R_t$ is unipotent of level $2$.  Similarly,
over $U$ there is a perfect duality between
$\underline R/\underline R^f$ and $\underline R'{}^t$.  Since the
latter is also isomorphic to the dual torus of the lattice $M_S$ over
$S$, where $M$ is the group defined in (\SGARefLink{sga7.sga7i.eix.n5-5-4}{5.5.4}), we obtain a canonical
isomorphism

}{E7I-IX-124}\SGApage{I}{\(\approx\)506}
\begin{equation}\tag{\SGARefTarget{sga7.sga7i.eix.n13-5-12}{13.5.12}}
 \underline R_U/\underline R^f_U\simeq M_U
 \qquad ,
\end{equation}

\noindent
and the homomorphism $i$ of (\SGARefLink{sga7.sga7i.eix.n13-3-4}{13.3.4}) may be written as a homomorphism

\begin{equation}\tag{\SGARefTarget{sga7.sga7i.eix.n13-5-13}{13.5.13}}
 i:M_U\longrightarrow A^\natural
 \qquad ,
\end{equation}

\noindent
which we shall encounter again in a slightly different context in the
next paragraph.


\medskip
\noindent
14. \underline{Links with transcendental theory: the rigid-analytic
case}.  \underline{The canonical homomorphism}
$\underline M_K\longrightarrow A_K^{\natural\mathrm o}$.

\medskip
\noindent
\SGARefTarget{sga7.sga7i.eix.n14-1}{14.1}. We assume here the foundations of the theory of rigid-analytic
spaces, for which one may consult [12], [13], and [34].

Let $S$ be a complete trait, and let $A_K$ be an abelian scheme over its
field of fractions.  In \SGARefLink{sga7.sga7i.eix.n8-1}{8.1}, by descent from the Raynaud extension in
the semistable case, we constructed an extension
$G_K=A_K^{\natural\mathrm o}$ (\SGARefLink{sga7.sga7i.eix.n8-1-1}{8.1.1})

\begin{equation}\tag{\SGARefTarget{sga7.sga7i.eix.n14-1-1}{14.1.1}}
 0\longrightarrow T_K\longrightarrow G_K\longrightarrow B_K
 \longrightarrow0
\end{equation}

\noindent
of an abelian scheme $B_K$ by a torus $T_K$.  Considering the dual
abelian scheme $A'_K$, we likewise have an extension
$G'_K=A'{}^{\natural\mathrm o}_K$

\begin{equation}\tag{\SGARefTarget{sga7.sga7i.eix.n14-1-2}{14.1.2}}
 0\longrightarrow T'_K\longrightarrow G'_K\longrightarrow B'_K
 \longrightarrow0\qquad ,
\end{equation}

\noindent
where $B'_K$ is moreover identified with the abelian scheme dual to
$B_K$ (\SGARefLink{sga7.sga7i.eix.n8-2}{8.2}).  As in (\SGARefLink{sga7.sga7i.eix.n8-1-8}{8.1.8}), let $\underline M_K$ and
$\underline M'_K$ denote the lattices over $K$ dual to $T'_K$ and
$T_K$, respectively:

\begin{equation}\tag{\SGARefTarget{sga7.sga7i.eix.n14-1-3}{14.1.3}}
 \underline M_K=D(T'_K),\qquad \underline M'_K=D(T_K)\qquad .
\end{equation}

\noindent
Thus the extension structures (\SGARefLink{sga7.sga7i.eix.n14-1-2}{14.1.2}) and (\SGARefLink{sga7.sga7i.eix.n14-1-1}{14.1.1}) may respectively be
interpreted (VIII \SGARefLink{sga7.sga7i.eviii.n3-7}{3.7}) as equivalent to the data of homomorphisms

\SGApage{I}{\(\approx\)507}
\begin{equation}\tag{\SGARefTarget{sga7.sga7i.eix.n14-1-4}{14.1.4}}
 \underline M_K\xrightarrow{\ p_K\ }B_K,\qquad
 \underline M'_K\xrightarrow{\ p'_K\ }B'_K\qquad .
\end{equation}

Let $W_K^{\natural}$ be the biextension of $(B_K,B'_K)$ by
$\mathbb G_{mK}$ which defines the canonical duality between these
abelian schemes (\SGARefLink{sga7.sga7i.eix.n8-2}{8.2}), and let $E_K$ be the biextension of
$(\underline M_K,\underline M'_K)$ by $\mathbb G_{mK}$ obtained by
pulling $W_K^{\natural}$ back along the homomorphisms (\SGARefLink{sga7.sga7i.eix.n14-1-4}{14.1.4}).  One
checks immediately that trivializations of $E_K$ correspond
canonically and bijectively to homomorphisms

\begin{equation}\tag{\SGARefTarget{sga7.sga7i.eix.n14-1-5}{14.1.5}}
 i:\underline M_K\longrightarrow G_K
\end{equation}

\noindent
lifting $p_K:\underline M_K\longrightarrow B_K$, or, symmetrically, to
homomorphisms

\begin{equation}\tag{\SGARefTarget{sga7.sga7i.eix.n14-1-6}{14.1.6}}
 i':\underline M'_K\longrightarrow G'_K
\end{equation}

\SGAvar{\noindent
lifting $P'_K:\underline M'_K\longrightarrow B'_K$.

}{\noindent
lifting $p'_K:\underline M'_K\longrightarrow B'_K$.

}{E7I-IX-125}With this understood, write a superscript ${}^{\mathrm{an}}$ for the
rigid-analytic groups over $K$ associated with the algebraic groups
under consideration [13].  In [24] one defines a canonical exact
sequence

\begin{equation}\tag{\SGARefTarget{sga7.sga7i.eix.n14-1-7}{14.1.7}}
 0\longrightarrow\underline M_K^{\mathrm{an}}
 \xrightarrow{\ i^{\mathrm{an}}\ }G_K^{\mathrm{an}}
 \xrightarrow{\ e\ }A_K^{\mathrm{an}}\longrightarrow0\qquad .
\end{equation}

\noindent
Here $i^{\mathrm{an}}$, which may also be interpreted as a homomorphism
$i:\underline M_K\longrightarrow G_K$ of ordinary group schemes,
\underline{lifts the homomorphism} $p$ in (\SGARefLink{sga7.sga7i.eix.n14-1-4}{14.1.4}), that is, it is of
the form considered in (\SGARefLink{sga7.sga7i.eix.n14-1-5}{14.1.5}).  This exact sequence is the
rigid-analytic analogue of the exact sequence (\SGARefLink{sga7.sga7i.eix.n13-3-4}{13.3.4}).  It depends
functorially on $A_K$, and its formation is compatible with every
change of traits $S'\longrightarrow S$.

\SGApage{I}{\(\approx\)508}
The latter fact also shows that, in order to define and study (\SGARefLink{sga7.sga7i.eix.n14-1-7}{14.1.7}),
finite descent reduces us to the case where $A_K$ has semistable
reduction.  In this case the group schemes
$T_K,G_K,B_K,\underline M_K,\ldots$ are the generic fibres of group
schemes $T,G,B,M,\ldots$ over $S$, with $G$ an extension of the abelian
scheme $B$ by the torus $T$ (paragraph 7).  We may then consider the
formal completion $\hat G$ of $G$ along its special fibre, and by
construction there is a canonical isomorphism (\SGARefLink{sga7.sga7i.eix.n7-2-6}{7.2.6})

\begin{equation}\tag{\SGARefTarget{sga7.sga7i.eix.n14-1-8}{14.1.8}}
 \hat G\xrightarrow{\ \sim\ }\widehat{A^{\mathrm o}}
 \qquad .
\end{equation}

\noindent
Passing to the ``generic fibres'', which are rigid-analytic groups over
$K$, gives a canonical isomorphism

\begin{equation}\tag{\SGARefTarget{sga7.sga7i.eix.n14-1-9}{14.1.9}}
 \hat G^{\mathrm{an}}\xrightarrow{\ \sim\ }
 \widehat{A^{\mathrm o}}\,{}^{\mathrm{an}}
 \qquad .
\end{equation}

\noindent
On the other hand, there are canonical homomorphisms (rigid-analytic
open immersions, but not in general isomorphisms)

\begin{equation}\tag{\SGARefTarget{sga7.sga7i.eix.n14-1-10}{14.1.10}}
 \hat G^{\mathrm{an}}\lhook\joinrel\longrightarrow G_K^{\mathrm{an}},
 \qquad
 \widehat{A^{\mathrm o}}\,{}^{\mathrm{an}}
 \lhook\joinrel\longrightarrow A_K^{\mathrm{an}}
 \qquad .
\end{equation}

\noindent
The homomorphism $e$ in (\SGARefLink{sga7.sga7i.eix.n14-1-7}{14.1.7}) is uniquely characterized by the
commutativity of the diagram

\begin{equation}\tag{\SGARefTarget{sga7.sga7i.eix.n14-1-11}{14.1.11}}
\begin{tikzcd}
 \hat G^{\mathrm{an}} \arrow[r] \arrow[d]
   \& \widehat{A^{\mathrm o}}\,{}^{\mathrm{an}} \arrow[d] \\
 G_K^{\mathrm{an}} \arrow[r] \& A_K^{\mathrm{an}}
\end{tikzcd} \quad ,
\end{equation}

\noindent
where the unnamed arrows are those of (\SGARefLink{sga7.sga7i.eix.n14-1-9}{14.1.9}) and (\SGARefLink{sga7.sga7i.eix.n14-1-10}{14.1.10}).

\medskip
\noindent
\underline{Remarks \SGARefTarget{sga7.sga7i.eix.n14-1-12}{14.1.12}.}

\smallskip
\noindent
a) Notice that the uniqueness of $e$ making the preceding diagram
commute is immediate, since the rigid-analytic space
$G_K^{\mathrm{an}}$ is connected.  The existence of $e$, on the other
hand, is by no means trivial, and on this point we merely refer to the
cited note of M.~Raynaud.  Moreover, the definition of
$i^{\mathrm{an}}$ given in loc. cit. is likewise transcendental in
nature, although it may be interpreted as an object of an essentially
algebraic nature (\SGARefLink{sga7.sga7i.eix.n14-1-5}{14.1.5}), once the construction of $G$ itself, and
hence of $G_K$, has been obtained through the constructions of
\underline{formal geometry} in paragraph 7.  I know of no purely
algebro-geometric description of it, but below (\SGARefLink{sga7.sga7i.eix.n14-4}{14.4}) I shall propose
an ``arithmetic'' characterization, obtained by reducing to a situation
of finite type over the ring $\mathbb Z$.  Thus, from the standpoint of
algebraic geometry, the homomorphism (\SGARefLink{sga7.sga7i.eix.n14-1-5}{14.1.5}) seems to be a more hidden
piece of structure than those encountered in the preceding paragraphs.

\SGApage{I}{\(\approx\)509}
\smallskip
\noindent
b) As explained above, the homomorphism $i$ is canonically associated
with a well-determined trivialization of the biextension $E_K$ of
$(\underline M_K,\underline M'_K)$, or equivalently with a
well-determined homomorphism (\SGARefLink{sga7.sga7i.eix.n14-1-6}{14.1.6}) lifting $p'_K$.  One then proves
that the latter is precisely the homomorphism occurring in the exact
sequence

\begin{equation}\tag{\SGARefTarget{sga7.sga7i.eix.n14-1-13}{14.1.13}}
 0\longrightarrow\underline M'{}^{\mathrm{an}}_K
 \xrightarrow{\ i'{}^{\mathrm{an}}\ }G'{}^{\mathrm{an}}_K
 \longrightarrow A'{}^{\mathrm{an}}_K\longrightarrow0
\end{equation}

\noindent
analogous to (\SGARefLink{sga7.sga7i.eix.n14-1-7}{14.1.7}), associated with the abelian scheme $A'_K$ dual
to $A_K$ (unless it is the opposite???).


\SGApage{I}{510}
\par\medskip
\noindent
\SGARefTarget{sga7.sga7i.eix.n14-2}{14.2}. \underline{Relations of the homomorphism}
$i:\underline M_K\longrightarrow G_K$ \underline{with monodromy.}

Let $n$ be an integer $>0$.  From the exact sequence (\SGARefLink{sga7.sga7i.eix.n14-1-7}{14.1.7}), by
means of the exact sequence of the
$\underline{\operatorname{Ext}}^i(\mathbb Z/n\mathbb Z,-)$ on the
site of rigid-analytic spaces over $K$, one then obtains a canonical
exact sequence (taking account of the fact that
$\underline M_K^{\mathrm{an}}$ is torsion-free and
$G_K^{\mathrm{an}}$ is $n$-divisible):

\begin{equation*}
 0\longrightarrow{}_nG_K^{\mathrm{an}}
 \longrightarrow{}_nA_K^{\mathrm{an}}
 \longrightarrow\underline M_{K,n}^{\mathrm{an}}
 \longrightarrow0\qquad ,
\end{equation*}

\noindent
where
$\underline M_{K,n}\overset{\mathrm{dfn}}{=}
\underline M_K/n\underline M_K$.  Since all the terms occurring in
this exact sequence come from finite group schemes over $K$, it may be
rewritten as arising from an exact sequence

\begin{equation}\tag{\SGARefTarget{sga7.sga7i.eix.n14-2-1}{14.2.1}}
 0\longrightarrow{}_nG_K\longrightarrow{}_nA_K
 \longrightarrow\underline M_{K,n}\longrightarrow0\qquad .
\end{equation}

\noindent
As $n$ varies, these exact sequences behave in the expected way, and,
on passing to the inverse limit over the powers of a prime number
$\ell$, one obtains the analogous exact sequence

\begin{equation}\tag{\SGARefTarget{sga7.sga7i.eix.n14-2-2}{14.2.2}}
 0\longrightarrow T_\ell(G_K)\longrightarrow T_\ell(A_K)
 \longrightarrow\underline M_K\otimes_{\mathbb Z}\mathbb Z_\ell
 \longrightarrow0\qquad .
\end{equation}

\noindent
But one already has such an exact sequence by purely algebraic means,
thanks to the first formula (\SGARefLink{sga7.sga7i.eix.n8-1-2}{8.1.2}) and to (\SGARefLink{sga7.sga7i.eix.n8-1-7}{8.1.7}), which give

\begin{equation}\tag{\SGARefTarget{sga7.sga7i.eix.n14-2-3}{14.2.3}}
 T_\ell(G_K)\xrightarrow{\ \sim\ }T_\ell(A_K)^{\mathrm{ef}}
 \subset T_\ell(A_K),\qquad
 T_\ell(A_K)/T_\ell(A_K)^{\mathrm{ef}}
 \simeq\underline M_K\otimes_{\mathbb Z}\mathbb Z_\ell\qquad .
\end{equation}

\noindent
With this understood, I claim that \underline{the homomorphisms}
(\SGARefLink{sga7.sga7i.eix.n14-2-3}{14.2.3}) \underline{arising from the algebraic theory are precisely
those deduced from the exact sequence} (\SGARefLink{sga7.sga7i.eix.n14-2-2}{14.2.2}) \underline{arising
from the transcendental theory}.  For the first homomorphism, which
corresponds to the first homomorphism of (\SGARefLink{sga7.sga7i.eix.n14-2-2}{14.2.2}), our assertion follows
at once, after reduction to the semistable situation, from the
commutativity of (\SGARefLink{sga7.sga7i.eix.n14-1-11}{14.1.11}), since the vertical arrows there induce
isomorphisms on the $T_\ell$, while applying $T_\ell$ to the horizontal
arrows gives, respectively, the arrow under consideration in (\SGARefLink{sga7.sga7i.eix.n14-2-3}{14.2.3})
and that in (\SGARefLink{sga7.sga7i.eix.n14-2-2}{14.2.2}).  We shall admit the asserted compatibility for
the two isomorphisms obtained between
$T_\ell(A_K)/T_\ell(A_K)^{\mathrm{ef}}$ and
$\underline M_K\otimes_{\mathbb Z}\mathbb Z_\ell$.  Notice, moreover,
that \underline{this compatibility} (for a given $\ell$)
\underline{already suffices to characterize uniquely, by
transcendental means, the homomorphism} (\SGARefLink{sga7.sga7i.eix.n14-1-5}{14.1.5}) \underline{occurring
in} (\SGARefLink{sga7.sga7i.eix.n14-1-7}{14.1.7}), interpreted as an isomorphism from
$\underline M_K^{\mathrm{an}}$ onto $\operatorname{Ker}e$.

\SGApage{I}{\(\approx\)511}
One may also say that the homomorphism $e$ of (\SGARefLink{sga7.sga7i.eix.n14-1-7}{14.1.7}) identifies the
$T_\ell(A_K)/T_\ell(A_K)^{\mathrm{ef}}$ with the $\ell$-adic
completions $\underline R_K\otimes\mathbb Z_\ell$ of a single lattice
$\underline R_K=\operatorname{Ker}e$ over $K$, and that this
``integral structure'' on the preceding system of pro-$\ell$ groups is
compatible, in the evident sense, with the integral structure arising
from formulas (\SGARefLink{sga7.sga7i.eix.n8-1-7}{8.1.7}) (orthogonality theorem \SGARefLink{sga7.sga7i.eix.n5-2}{5.2}).  This then defines
canonically an isomorphism

\begin{equation*}
 i:\underline M_K\xrightarrow{\ \sim\ }\operatorname{Ker}e
 \qquad .
\end{equation*}

Since the homomorphism $i$ (\SGARefLink{sga7.sga7i.eix.n14-1-5}{14.1.5}) enables one to reconstruct the
usual extension structure on $T_\ell(A_K)$, namely (\SGARefLink{sga7.sga7i.eix.n14-2-2}{14.2.2}), and hence
also its structure as a \underline{blended extension}, described in
\SGARefLink{sga7.sga7i.eix.n9-6}{9.6}, by means of the exact sequence deduced from (\SGARefLink{sga7.sga7i.eix.n14-1-1}{14.1.1}),

\begin{equation*}
\begin{array}{ccccccccc}
0&\longrightarrow&T_\ell(T_K)&\longrightarrow&T_\ell(G_K)
 &\longrightarrow&T_\ell(B_K)&\longrightarrow&0\\[5ex]
 &&\underline M'_K\otimes\mathbb Z_\ell(1)&&&&&&
\end{array}
\end{equation*}

\SGApage{I}{\(\approx\)512}
\noindent
it also enables one, in the semistable case, to reconstruct the
\underline{monodromy pairings} (\SGARefLink{sga7.sga7i.eix.n9-1-2}{9.1.2})

\begin{equation}\tag{\SGARefTarget{sga7.sga7i.eix.n14-2-4}{14.2.4}}
 u_\ell:\underline M\otimes\underline M'
 \longrightarrow\mathbb Z_\ell\qquad ,
\end{equation}

\SGApage{I}{436}
\noindent
these having been defined precisely in terms of blended extensions in
\SGARefLink{sga7.sga7i.eix.n9-6}{9.6}.  Let us merely indicate the result of translating the construction
of \SGARefLink{sga7.sga7i.eix.n9-6}{9.6}.  For this purpose, observe that in the semistable case $p_K$
and $p'_K$ arise from homomorphisms

\SGApage{I}{\(\approx\)512}
\SGAvar{\begin{equation*}
 p:\underline M\longrightarrow B,\qquad
 \underline M'\longrightarrow B'\qquad ,
\end{equation*}

\noindent
where $B,B'$ are abelian schemes over $S$, dual to one another.  Pulling
back by $(p,p')$ therefore gives a biextension $E$ of
$(\underline M, M')$ by $\mathbb G_{mS}$ whose generic fibre
is the biextension $E_K$ already considered in \SGARefLink{sga7.sga7i.eix.n14-1}{4.1}.

}{\begin{equation*}
 p:\underline M\longrightarrow B,\qquad
 p':\underline M'\longrightarrow B'\qquad ,
\end{equation*}

\noindent
where $B,B'$ are abelian schemes over $S$, dual to one another.  Pulling
back by $(p,p')$ therefore gives a biextension $E$ of
$(\underline M,\underline M')$ by $\mathbb G_{mS}$ whose generic fibre
is the biextension $E_K$ already considered in \SGARefLink{sga7.sga7i.eix.n14-1}{14.1}.

}{E7I-IX-126}With this said, suppose for simplicity that $\underline M$ and
$\underline M'$ are constant, with values $M$ and $M'$.  For every
pair $(m,m')\in M\times M'$, $E_{m,m'}$ is a torsor under
$\mathbb G_{mS}$, and hence corresponds to an invertible module
$L_{m,m'}$ over the ring $V$ of $S$.  It therefore admits a basis
$u_0$, determined up to multiplication by a unit of $V$.  Consequently
one obtains a canonical ``valuation'' map

\begin{equation*}
 \underline v:E_{m,m'}(K)\longrightarrow\mathbb Z
\end{equation*}

\noindent
which associates with every basis $u$ of
$L_{m,m'}\otimes_V K$, written $u=\lambda u_0$ ($\lambda\in K^*$),
the element

\begin{equation*}
 \underline v(u)=\underline v(\lambda)\in\mathbb Z\qquad ,
\end{equation*}

\noindent
where $\underline v$ on the right denotes the valuation associated
with $V$ (normalized in the usual way, so that its value group is
$\mathbb Z$); the resulting integer is plainly independent of the
choice of $u_0$.  If $\varphi$ is now a trivialization of $E$, one may
consider the corresponding map

\SGApage{I}{\(\approx\)513}
\begin{equation}\tag{\SGARefTarget{sga7.sga7i.eix.n14-2-5}{14.2.5}}
 (m,m')\longmapsto\underline v\bigl(\varphi(m,m')\bigr):
 M\times M'\longrightarrow\mathbb Z\qquad .
\end{equation}

\noindent
With this understood, one finds that the \underline{monodromy
homomorphism} (\SGARefLink{sga7.sga7i.eix.n14-2-4}{14.2.4}) \underline{takes its values in} $\mathbb Z$
\underline{and is precisely the homomorphism} (\SGARefLink{sga7.sga7i.eix.n14-2-5}{14.2.5}), \underline{in
which} $\varphi$ \underline{is the splitting of} $E_K$
\underline{associated with the homomorphism} $i$ (\SGARefLink{sga7.sga7i.eix.n14-1-5}{14.1.5}).

\medskip
\noindent
\underline{Remarks \SGARefTarget{sga7.sga7i.eix.n14-2-6}{14.2.6}.}  In particular, one recovers (by means of
the transcendental rigid-analytic theory) the fact that the monodromy
pairing is integral and independent of $\ell$ (\SGARefLink{sga7.sga7i.eix.n10-4}{10.4} a)), proved by a
different (and likewise transcendental) method in \SGARefLink{sga7.sga7i.eix.n12-5}{12.5}.  According to
M.~Raynaud, rigid-analytic theory also recovers the positivity assertion
\SGARefLink{sga7.sga7i.eix.n10-4}{10.4} b) (which is likewise contained in \SGARefLink{sga7.sga7i.eix.n12-5}{12.5}).

\medskip
\noindent
\underline{Remarks \SGARefTarget{sga7.sga7i.eix.n14-2-7}{14.2.7}.}  Regard
$\underline M_K\otimes_{\mathbb Z}\mathbb Q$ as an extension of
$\underline M_K\otimes_{\mathbb Z}\mathbb Q/\mathbb Z$ by
$\underline M_K$.  Using the homomorphism (\SGARefLink{sga7.sga7i.eix.n14-1-5}{14.1.5})
$i:\underline M_K\longrightarrow G_K$, one deduces an extension
$A_K^{\flat}$

\begin{equation}\tag{\SGARefTarget{sga7.sga7i.eix.n14-2-8}{14.2.8}}
 0\longrightarrow G_K\longrightarrow A_K^{\flat}
 \longrightarrow\underline M_K\otimes\mathbb Q/\mathbb Z
 \longrightarrow0\qquad .
\end{equation}

\noindent
Observe moreover that, since $G_K$ is a divisible group, the datum of
such an extension is equivalent to that of a family of extensions of
the form (\SGARefLink{sga7.sga7i.eix.n14-2-2}{14.2.2}), namely those deduced from (\SGARefLink{sga7.sga7i.eix.n14-2-8}{14.2.8}) by applying the
$T_\ell$.  Since $i$ is a monomorphism, the extensions in question
which define (\SGARefLink{sga7.sga7i.eix.n14-2-8}{14.2.8}) are precisely the extensions (\SGARefLink{sga7.sga7i.eix.n14-2-2}{14.2.2}) arising
from the exact sequence (\SGARefLink{sga7.sga7i.eix.n14-2-1}{14.2.1}).  As these latter extensions are
defined by purely algebraic means, the same is true of the extension
(\SGARefLink{sga7.sga7i.eix.n14-2-8}{14.2.8}).  In fact, the characterization just given shows at once that
this extension is canonically isomorphic to the group, likewise denoted
$A_K^{\flat}$, introduced in (8.2.3), in view of the isomorphism \SGARefLink{sga7.sga7i.eix.n11-11}{11.11} from
the group of connected components of the latter to
$\underline M_K\otimes\mathbb Q/\mathbb Z$.


\SGApage{I}{514}
\par\medskip
\noindent
\SGARefTarget{sga7.sga7i.eix.n14-3}{14.3}. \underline{Application to an arithmetic characterization of}
$i:\underline M_K\longrightarrow G_K$.

The characterization we have in mind rests on the observation that,
for every prime number $\ell$, the extension (\SGARefLink{sga7.sga7i.eix.n14-2-2}{14.2.2}) of
$\underline M_K\otimes\mathbb Z_\ell$ by $T_\ell(G_K)$ deduced from
$i$ is known by purely algebraic means.  In particular, the image of
$i$ under the canonical homomorphism

\begin{equation}\tag{\SGARefTarget{sga7.sga7i.eix.n14-3-1}{14.3.1}}
 \operatorname{Hom}(\underline M_K,G_K)\longrightarrow
 \prod_\ell\operatorname{Ext}^1
 \bigl(\underline M_K\otimes\mathbb Z_\ell,T_\ell(G_K)\bigr)
 \qquad .
\end{equation}

\noindent
is known.  In the case (to which one may always reduce by descent) in
which $\underline M_K$ is constant, with value $M$, the preceding
homomorphism may be interpreted as

\SGAvar{\begin{equation}\tag{\SGARefTarget{sga7.sga7i.eix.n14-3-2}{14.3.2}}
 \operatorname{Hom}\bigl(M,G_K(K)\bigr)\longrightarrow
 \prod_\ell\operatorname{Hom}
 \bigl(M,H^1(K,T_\ell(G_K)),
\end{equation}

}{\begin{equation}\tag{\SGARefTarget{sga7.sga7i.eix.n14-3-2}{14.3.2}}
 \operatorname{Hom}\bigl(M,G_K(K)\bigr)\longrightarrow
 \prod_\ell\operatorname{Hom}
 \bigl(M,H^1(K,T_\ell(G_K))\bigr),
\end{equation}

}{E7I-IX-127}\noindent
deduced from the canonical homomorphism

\begin{equation}\tag{\SGARefTarget{sga7.sga7i.eix.n14-3-3}{14.3.3}}
 G_K(K)\longrightarrow\prod_\ell H^1(K,T_\ell(G_K))
 \qquad ,
\end{equation}

\noindent
which comes from the connecting homomorphisms in the Kummer exact
sequences

\begin{equation*}
 0\longrightarrow{}_nG_K\longrightarrow G_K
 \xrightarrow{\ n\ }G_K\longrightarrow0\qquad .
\end{equation*}

\SGApage{I}{\(\approx\)515}
(The homomorphism (\SGARefLink{sga7.sga7i.eix.n14-3-1}{14.3.1}) could in fact always be interpreted as the
connecting homomorphism analogous to (\SGARefLink{sga7.sga7i.eix.n14-3-3}{14.3.3}) for the twisted group
$\check{\underline M}_K\otimes G_K=
\underline{\operatorname{Hom}}(\underline M_K,G_K)$.)  The kernel of
(\SGARefLink{sga7.sga7i.eix.n14-3-3}{14.3.3}) plainly consists of the elements of $G_K(K)$ which are
\underline{divisible by} $n$ \underline{for every integer}
$n\geq1$, or, as we shall say, which are \underline{infinitely
divisible}.  Thus this already determines $i$ modulo an element of
$\operatorname{Hom}(M,g)$, where $g$ is the subgroup of $G_K(K)$
formed by its infinitely divisible elements:

\begin{equation*}
 g=\bigcap_{n\geq1}nG_K(K)\qquad .
\end{equation*}

We now have the following result.

\medskip
\noindent
\underline{Lemma \SGARefTarget{sga7.sga7i.eix.n14-3-4}{14.3.4}.}  \underline{Let} $V$ \underline{be a
discrete valuation ring whose residue field} $k$ \underline{has
characteristic} $p>0$ \underline{and is of finite type over its prime
field, and let} $G_K$ \underline{be a group scheme of finite type over}
$K$, \underline{of unipotent rank zero if} $\operatorname{char}K=0$.
\underline{Then the group} $g$ \underline{of infinitely divisible
elements of} $G_K(K)$ \underline{is zero; that is,} (\underline{if}
$G_K$ \underline{is divisible}) \underline{the homomorphism} (\SGARefLink{sga7.sga7i.eix.n14-3-3}{14.3.3})
\underline{is injective}.

\SGAvar{After making a finite extension $K'/K$ and replacing $V$ by its
normalization in $K'$, we may suppose that $G_K$ admits a composition
series whose successive quotients are of one of the following types:
a constant finite group, an abelian variety, $\mathbb G_{mK}$, or
$\mathbb G_{aK}$.  We are therefore immediately reduced to the case
where $\mathbb G_K$ itself is of one of these types.  The case of a constant
finite group, or of $\mathbb G_{aK}$, is trivial.  In the abelian case,
$G_K$ may be regarded as the generic fibre of its N\'{e}ron model $G$,
so we are reduced to showing that, if $G$ is a scheme of finite type
over $S$, the subgroup of $G(S)$ formed by the infinitely divisible
points is zero.

This is true of the group $G(k)$, as follows at once from the usual
d\'{e}vissage of $G_o$ and from the \underline{Mordell--Weil theorem}
[16] (which says that $G_o(k)$ is a finitely generated $\mathbb Z$-module
if $G_o$ is an abelian scheme over the field of finite type $k$).
Consequently $g$ is contained in the kernel $G(S)^{(1)}$ of
$G(S)\longrightarrow G(k)$.  The kernel of
$G(S)^{(1)}\longrightarrow G(V/\underline m^{,n+1})$ has a filtration
whose quotients are isomorphic to vector groups over $k$ (after
supposing, as we may, that $S$ is complete), and these are killed by
$p$.  It follows that $g$ is contained in the kernel of the preceding
homomorphism for every $n$, and hence that it is zero.  The case
$G_K=\mathbb G_{mK}$ reduces to the preceding one, since
$G_K(K)\longrightarrow K^*$ is an extension of $\mathbb Z$ (which has no nonzero
infinitely divisible elements) by
$V^*=\mathbb G_m(S)$, and the latter satisfies the same condition by
what we have just proved.

}{After making a finite extension $K'/K$ and replacing $V$ by its
normalization in $K'$, we may suppose that $G_K$ admits a composition
series whose successive quotients are of one of the following types:
a constant finite group, an abelian variety, $\mathbb G_{mK}$, or
$\mathbb G_{aK}$.  We are therefore immediately reduced to the case
where $G_K$ itself is of one of these types.  The case of a constant
finite group, or of $\mathbb G_{aK}$, is trivial.  In the abelian case,
$G_K$ may be regarded as the generic fibre of its N\'{e}ron model $G$,
so we are reduced to showing that, if $G$ is a scheme of finite type
over $S$, the subgroup of $G(S)$ formed by the infinitely divisible
points is zero.

This is true of the group $G(k)$, as follows at once from the usual
d\'{e}vissage of $G_o$ and from the \underline{Mordell--Weil theorem}
[16] (which says that $G_o(k)$ is a finitely generated $\mathbb Z$-module
if $G_o$ is an abelian scheme over the field of finite type $k$).
Consequently $g$ is contained in the kernel $G(S)^{(1)}$ of
$G(S)\longrightarrow G(k)$.  The kernel of
$G(S)^{(1)}\longrightarrow G(V/\underline m^{,n+1})$ has a filtration
whose quotients are isomorphic to vector groups over $k$ (after
supposing, as we may, that $S$ is complete), and these are killed by
$p$.  It follows that $g$ is contained in the kernel of the preceding
homomorphism for every $n$, and hence that it is zero.  The case
$G_K=\mathbb G_{mK}$ reduces to the preceding one, since
$G_K(K)=K^*$ is an extension of $\mathbb Z$ (which has no nonzero
infinitely divisible elements) by
$V^*=\mathbb G_m(S)$, and the latter satisfies the same condition by
what we have just proved.

}{E7I-IX-128}\SGApage{I}{\(\approx\)516}
\medskip
\noindent
\underline{Remarks \SGARefTarget{sga7.sga7i.eix.n14-3-4-1}{14.3.4.1}.}  Notice that in \SGARefLink{sga7.sga7i.eix.n14-3-4}{14.3.4} it is essential
to require divisibility by \underline{all} integers $\geq1$, and not
only by the powers of a fixed prime $\ell$; equivalently, one must
consider the kernel of (\SGARefLink{sga7.sga7i.eix.n14-3-3}{14.3.3}), rather than restricting the right-hand
side to a single factor $H^1(K,T_\ell(G_K))$.  Indeed, if $G_K$ comes
from a smooth group scheme $G$ over $S$ and $S$ is henselian, every
element of the kernel of $G(S)\longrightarrow G(k)$ is infinitely
$\ell$-divisible for every $\ell\neq p$!  The argument proving \SGARefLink{sga7.sga7i.eix.n14-3-4}{14.3.4}
shows that the subgroup of $G_K(K)$ formed by the infinitely
$p$-divisible elements (that is, the kernel of
$G_K(K)\longrightarrow H^1(K,T_p(G_K))$ when $G_K$ is $p$-divisible)
is a \underline{finite} group.

\SGApage{I}{517}
\medskip
\noindent
\SGARefTarget{sga7.sga7i.eix.n14-3-5}{14.3.5}. The preceding considerations show that, when the residue
field $k$ has characteristic $p>0$ and is of finite type in the
absolute sense, the homomorphism $i$ (\SGARefLink{sga7.sga7i.eix.n14-1-5}{14.1.5}) is completely determined
by its image under (\SGARefLink{sga7.sga7i.eix.n14-3-1}{14.3.1}), that is, by the extensions (\SGARefLink{sga7.sga7i.eix.n14-2-2}{14.2.2}).  To
deduce from this a principle characterizing $i$ in the general case,
with no restriction on $k$, observe that in paragraph 7 we were able to
construct the scheme $G$ over $S$ under conditions considerably more
general than those in which $S$ is a complete trait and one has an
abelian variety with semistable reduction over its field of fractions.

\SGAvar{More generally, let $S$ denote the spectrum of a normal noetherian ring
$V$, separated and complete for an $\underline m$-adic topology, and
consider a smooth quasi-projective group scheme $A$ over $S$, with
connected fibres.  Suppose that
$A_o=A\times_S S_o$, where
$S_o=\operatorname{Spec}(V/\underline m)$, is an extension of an
abelian scheme $B_o$ by an isotrivial torus $T_o$, and that $A_U$ is an
abelian scheme, where $U\subset S-S_o$ is a nonempty open subset.  In
\SGARefLink{sga7.sga7i.eix.n7-2-2}{7.22} we then constructed a scheme $G$ over $S$, an extension of an
abelian scheme $B$ by a torus $T$, characterized up to unique
isomorphism by the isomorphism (\SGARefLink{sga7.sga7i.eix.n14-1-8}{14.1.8})

}{More generally, let $S$ denote the spectrum of a normal noetherian ring
$V$, separated and complete for an $\underline m$-adic topology, and
consider a smooth quasi-projective group scheme $A$ over $S$, with
connected fibres.  Suppose that
$A_o=A\times_S S_o$, where
$S_o=\operatorname{Spec}(V/\underline m)$, is an extension of an
abelian scheme $B_o$ by an isotrivial torus $T_o$, and that $A_U$ is an
abelian scheme, where $U\subset S-S_o$ is a nonempty open subset.  In
\SGARefLink{sga7.sga7i.eix.n7-2-2}{7.2.2} we then constructed a scheme $G$ over $S$, an extension of an
abelian scheme $B$ by a torus $T$, characterized up to unique
isomorphism by the isomorphism (\SGARefLink{sga7.sga7i.eix.n14-1-8}{14.1.8})

}{E7I-IX-129}\begin{equation}\tag{\SGARefTarget{sga7.sga7i.eix.n14-3-5-1}{14.3.5.1}}
 \hat G\xrightarrow{\ \sim\ }\hat A
\end{equation}

\noindent
on the formal completions along $G_o$ and $A_o$.  For every prime
number $\ell$ one deduces an isomorphism (\SGARefLink{sga7.sga7i.eix.n7-3-5}{7.3.5})

\SGApage{I}{\(\approx\)518}
\begin{equation}\tag{\SGARefTarget{sga7.sga7i.eix.n14-3-5-2}{14.3.5.2}}
 T_\ell(G)\xrightarrow{\ \sim\ }T_\ell(A)^f\qquad ,
\end{equation}

\noindent
and hence, on restricting to $U$, the structure of an extension of
pro-$\ell$ Barsotti--Tate groups

\begin{equation}\tag{\SGARefTarget{sga7.sga7i.eix.n14-3-5-3}{14.3.5.3}}
\begin{tikzcd}[row sep=large,column sep=large]
0 \arrow[r] \& T_\ell(G_U) \arrow[r] \arrow[d,dash,"\wr"']
 \& T_\ell(A_U) \arrow[r]
 \& T_\ell(A_U)/T_\ell(A_U)^f \arrow[r] \& 0\quad .\\
\& T_\ell(A_U)^f
\end{tikzcd}
\end{equation}

\SGAvar{Suppose now that a second group scheme $G'$ over $S$ is given,
satisfying the same hypotheses as $A$, together with a biextension $W$
of $(A,A')$ by $\mathbb G_m$ which over $U$ induces a duality between
the abelian schemes $A_U$ and $A'_U$.  (In fact, by VIII \SGARefLink{sga7.sga7i.eviii.n7-1}{7.1}, giving
$W$ is equivalent to giving such a duality between $A_U$ and $A'_U$.)
Then $W$ defines pairings

}{Suppose now that a second group scheme $A'$ over $S$ is given,
satisfying the same hypotheses as $A$, together with a biextension $W$
of $(A,A')$ by $\mathbb G_m$ which over $U$ induces a duality between
the abelian schemes $A_U$ and $A'_U$.  (In fact, by VIII \SGARefLink{sga7.sga7i.eviii.n7-1}{7.1}, giving
$W$ is equivalent to giving such a duality between $A_U$ and $A'_U$.)
Then $W$ defines pairings

}{E7I-IX-130}\begin{equation}\tag{\SGARefTarget{sga7.sga7i.eix.n14-3-5-4}{14.3.5.4}}
 T_\ell(A_U)\times T_\ell(A'_U)\longrightarrow\mathbb Z_\ell(1)
 \qquad ,
\end{equation}

\SGAvar{\noindent
which give rise to the orthogonality theorem (\SGARefLink{sga7.sga7i.eix.n7-4-8}{7.4.8}):
$T_\ell(A_U)^f$ is the orthogonal of
$T_\ell(A'_U)^t=T_\ell(T'_U)$, where $T$ is the toric part of the
Raynaud extension $G'={A'}^{\natural}$ associated with $A'$.  Thus one
again obtains an isomorphism

}{\noindent
which give rise to the orthogonality theorem (\SGARefLink{sga7.sga7i.eix.n7-4-8}{7.4.8}):
$T_\ell(A_U)^f$ is the orthogonal of
$T_\ell(A'_U)^t=T_\ell(T'_U)$, where $T'$ is the toric part of the
Raynaud extension $G'={A'}^{\natural}$ associated with $A'$.  Thus one
again obtains an isomorphism

}{E7I-IX-131}\begin{equation}\tag{\SGARefTarget{sga7.sga7i.eix.n14-3-5-5}{14.3.5.5}}
 T_\ell(A_U)/T_\ell(A_U)^f
 \simeq\underline M_U\otimes\mathbb Z_\ell\qquad ,
\end{equation}

\SGAvar{\noindent
where $\underline M=D(T')$ is the lattice over $S$ dual to the torus
$T'$.  Consequently (\SGARefLink{sga7.sga7i.eix.n14-3-5-3}{14.3.8}) and (\SGARefLink{sga7.sga7i.eix.n14-3-5-5}{14.3.10}) give an extension

}{\noindent
where $\underline M=D(T')$ is the lattice over $S$ dual to the torus
$T'$.  Consequently (\SGARefLink{sga7.sga7i.eix.n14-3-5-3}{14.3.5.3}) and (\SGARefLink{sga7.sga7i.eix.n14-3-5-5}{14.3.5.5}) give an extension

}{E7I-IX-132}\begin{equation}\tag{\SGARefTarget{sga7.sga7i.eix.n14-3-5-6}{14.3.5.6}}
 0\longrightarrow T_\ell(G_U)\longrightarrow T_\ell(A_U)
 \longrightarrow\underline M_U\otimes\mathbb Z_\ell
 \longrightarrow0\qquad .
\end{equation}

With this understood, we may state the following conjecture.


\SGApage{I}{\(\approx\)519}
\par\medskip
\noindent
\underline{Conjecture \SGARefTarget{sga7.sga7i.eix.n14-4}{14.4}.}  The hypotheses being those just set out in
\SGARefLink{sga7.sga7i.eix.n14-3-5}{14.3.5}, there exists a homomorphism

\begin{equation}\tag{\SGARefTarget{sga7.sga7i.eix.n14-4-1}{14.4.1}}
 i:\underline M_U\longrightarrow G_U
\end{equation}

\noindent
such that, for every prime number $\ell$, the extension (\SGARefLink{sga7.sga7i.eix.n14-3-5-6}{14.3.5.6}) is
isomorphic to the extension deduced from $i$ as explained for (\SGARefLink{sga7.sga7i.eix.n14-2-2}{14.2.2}).
More precisely, there is one and only one way of associating with every
situation of the type considered in \SGARefLink{sga7.sga7i.eix.n14-3-5}{14.3.5} a homomorphism (\SGARefLink{sga7.sga7i.eix.n14-4-1}{14.4.1})
satisfying the preceding condition, in such a way that the formation of
$i$ is compatible with every base change $S'\longrightarrow S$
associated with a continuous homomorphism $V\longrightarrow V'$ of
noetherian adic rings.

\medskip
\noindent
\SGARefTarget{sga7.sga7i.eix.n14-4-2}{14.4.2}. It seems that the existence of $i$, compatibly with base change,
follows from M. Raynaud's rigid-analytic theory, in view of the fact that
he is developing it in a framework more general than that stated in
\SGARefLink{sga7.sga7i.eix.n14-1}{14.1}, working, as here, over general noetherian adic rings.  Moreover,
to prove the existence and uniqueness asserted in \SGARefLink{sga7.sga7i.eix.n14-4}{14.4}, one is
immediately reduced to considering only those $S$ for which $S_o$ is a
scheme of finite type over $\mathbb Z$.  Indeed, $V$ is the filtered
inductive limit of subrings $V_i$ of finite type over $\mathbb Z$; these
may be supposed normal by Nagata's theorem (EGA IV 7.8), and, for $i$
sufficiently large, the situation $(A,A',W)$ comes from an analogous
situation over $V_i$, or alternatively over the separated completion of
$V_i$ for the topology defined by $m_i=\underline m V_i$.  This gives
the asserted reduction.

\SGApage{I}{520}
\SGAvar{\medskip
\noindent
\SGARefTarget{sga7.sga7i.eix.n14-4-3}{14.4.3}. Observe now that $S-S_o$ is a Jacobson scheme (EGA IV 10.5.8).
Thus, to see that two solutions $i,i'$ of the problem, for a given
situation $(A,A',W)$, are equal, it is enough to prove that these
homomorphisms are equal at every \underline{closed} point $\eta$ of
$U$.  Let $S'$ be the normalization of the closure of $\eta$ in $S$.
Then $S'$ is a complete trait, and the uniqueness result of \SGARefLink{sga7.sga7i.eix.n14-3-4}{4.3} may be
applied to the pullback to $S'$ of the situation over $S$, since the
residue field of $S'$ is \underline{finite}.  Once uniqueness is thus
established, it implies compatibility of the formation of $i$ with base
change, subject to existence.

}{\medskip
\noindent
\SGARefTarget{sga7.sga7i.eix.n14-4-3}{14.4.3}. Observe now that $S-S_o$ is a Jacobson scheme (EGA IV 10.5.8).
Thus, to see that two solutions $i,i'$ of the problem, for a given
situation $(A,A',W)$, are equal, it is enough to prove that these
homomorphisms are equal at every \underline{closed} point $\eta$ of
$U$.  Let $S'$ be the normalization of the closure of $\eta$ in $S$.
Then $S'$ is a complete trait, and the uniqueness result of \SGARefLink{sga7.sga7i.eix.n14-3-4}{14.3.4} may be
applied to the pullback to $S'$ of the situation over $S$, since the
residue field of $S'$ is \underline{finite}.  Once uniqueness is thus
established, it implies compatibility of the formation of $i$ with base
change, subject to existence.

}{E7I-IX-133}\medskip
\noindent
\underline{Remarks \SGARefTarget{sga7.sga7i.eix.n14-5}{14.5}.}  It would remain to formulate with the desired
generality, and to establish, a compatibility between the homomorphism
$i$ introduced in \SGARefLink{sga7.sga7i.eix.n14-4}{14.4} by a rigid-analytic method and the analogous
homomorphism (\SGARefLink{sga7.sga7i.eix.n13-5-13}{13.5.13}) of complex-analytic theory---or, better still, to
find a common generalization of them!



\bigskip
\noindent\underline{BIBLIOGRAPHY}

\SGAvar{\begin{enumerate}
\item[{[1]}] M. Artin, \emph{Algebraization of formal moduli I}
  (forthcoming).
\item[{[2]}] H. Cartan, \emph{Familles d'espaces complexes et fondements
  de la g\'{e}om\'{e}trie analytique}, S\'{e}minaire ENS 1960/61.
\item[{[3]}] P. Deligne, D. Mumford, \emph{The irreducibility of the space
  of curves of given genus}, Publ. Math. no.~36 (1969).
\item[{[4]}] P. Deligne, work in preparation on Hodge theory (forthcoming
  in Publ. Math., 1970).
\item[{[5]}] M. Demazure, J. Giraud, M. Raynaud, \emph{Sch\'{e}mas
  ab\'{e}liens}, S\'{e}minaire Orsay 1967/68, forthcoming.
\item[{[6]}] A. Grothendieck, \emph{Le Groupe de Brauer III}, in
  \emph{Dix Expos\'{e}s sur la cohomologie des sch\'{e}mas}, North-Holland
  Publishing Co.
\item[{[7]}] A. Grothendieck, \emph{Crystals and the De Rham cohomology of
  schémes} (notesby I. Coates and O. Jussila), in \emph{Dix expos\'{e}s
  sur la cohomologie des sch\'{e}mas}, Noth-Holland Publishing Co.
\item[{[8]}] A. Grothendieck, \emph{Un th\'{e}or\`{e}me sur les
  homomorphismes de sch\'{e}mas ab\'{e}liens}, Invent. Math. \textbf{2}
  (1966), pp.~59--78.
\item[{[9]}] P. A. Griffiths, \emph{Report on variation of Hodge
  structures} (forthcoming).
\item[{[10]}] M. Hakim, \emph{Topos annel\'{e}s et sch\'{e}mas relatifs},
  forthcoming in Grundlehren der Mathematik, Springer (1970).
\item[{[11]}] J. I. Igusa, \emph{Abstract vanishing cycle theory},
  Proceedings of the Japan Academy \textbf{34}, no.~9 (1958),
  pp.~589--593.
\item[{[12]}] R. Kiehl, \emph{Der Endlichkeitssatzf\"{u}r eigentliche
  Abbildnngen}, Invent. Math. \textbf{2} (1967), pp.~191--214.
\item[{[13]}] R. Kiehl.
\item[{[14]}] S. L. Kleiman, \emph{Algebraic cycles and the Weil
  conjectures}, in \emph{Dix expos\'{e}s sur la cohomologie des sch\'{e}mas},
  North-Holland Publishing Co.
\item[{[15]}] S. Lang, \emph{Abelian varieties}, Interscience Publishers,
  no.~7 (1959).
\item[{[16]}] S. Lang and A. N\'{e}ron, \emph{Rational points of abelian
  varieties over function fields}, Amer. J. Math. \textbf{81}, no.~1
  (1959), pp.~95--118.
\item[{[16bis]}] U. I. Manin, \emph{Th\'{e}orie des groupes formels
  commutatifs}, Uspekhi Mat. Nauk \textbf{XVIII} (1963), pp.~4--90.
\item[{[17]}] D. Mumford, \emph{Geometric invariant theory}, Ergebnisse
  der Mathematik 34, Springer (1965).
\item[{[18]}] D. Mumford, \emph{Abstract Theta functions}, Advanced
  Science Seminar in Algebraic Geometry, Bowdoin College, 1967 (notes by
  H. Pittie).
\item[{[19]}] A. N\'{e}ron, \emph{Mod\`{e}les minimaux des vari\'{e}t\'{e}s
  ab\'{e}liennes sur les corps locaux et globaux}, Publ. Math. 21 (1964),
  pp.~5--128.
\item[{[20]}] M. Raynaud, \emph{Caract\'{e}ristique d'Euler--Poincar\'{e}
  d'un faisceau et et cohomologie des vari\'{e}t\'{e}s ab\'{e}liennes}, in
  \emph{Dix expos\'{e}s sur la cohomologie des sch\'{e}mas}, North-Holland
  Publishing Co.
\item[{[21]}] M. Raynaud, \emph{Mod\`{e}les de N\'{e}ron}, C. R. Acad.
  Sci. Paris \textbf{262} (February 1966), pp.~413--416.
\item[{[22]}] M. Raynaud, \emph{Sp\'{e}cialisation du foncteur de Picard},
  C. R. Acad. Sci. \textbf{264}, pp.~941--943 and 1001--1004.
\item[{[22bis]}] M. Raynaud, \emph{Sp\'{e}cialisation du foncteur de
  Picard}, forthcoming in Publ. Math.
\item[{[23]}] M. Raynaud, \emph{Faisceaux amples sur les sch\'{e}mas en
  groupes et les espaces homog\`{e}nes}, Paris thesis, 1968 (forthcoming
  in Lecture Notes, Springer, 1970).
\item[{[24]}] M. Raynaud.
\item[{[25]}] M. Rosenlicht, \emph{Some basic theorems on algebraic
  groups}, Amer. J. Math. \textbf{78}, no.~2 (1956), pp.~401--443.
\item[{[25bis]}] J.-P. Serre, \emph{Quelques propri\'{e}t\'{e}s des
  vari\'{e}t\'{e}s ab\'{e}liennes en caract\'{e}ristique $p$}, Amer. J.
  Math. \textbf{80} (1958), pp.~715--739.
\item[{[26]}] J.-P. Serre, \emph{Corps locaux}, Act. Sci. Ind. 1296
  (1962), Herman, Paris.
\item[{[27]}] J.-P. Serre, \emph{Sur les corps locaux \`{a} corps
  r\'{e}siduel alg\'{e}briquement clos}, Bull. Soc. Math. France
  \textbf{89} (1961), pp.~105--154.
\item[{[28]}] J.-P. Serre, \emph{Abelian $\ell$-adic representations and
  elliptic curves}, Benjamin Inc. (1968).
\item[{[29]}] J.-P. Serre and J. Tate, \emph{Good reduction of abelian
  varieties}, Ann. of Math. \textbf{88}, no.~3 (1968), pp.~492--517.
\item[{[30]}] J.-P. Serre and J. Tate, \emph{Elliptic curves and formal
  groups} (notes of the Summer Institute on Algebraic Geometry, Woods
  Hole, 1964).
\item[{[31]}] J.-P. Serre, \emph{Morphismes universels et vari\'{e}t\'{e}
  d'Albanese}, S\'{e}minaire Chevalley 1958/59, Expos\'{e} 10.
\item[{[32]}] J. Tate, \emph{WC-groups over $p$-adic fields}, S\'{e}minaire
  Bourbaki no.~156 (1957).
\item[{[33]}] J. Tate, \emph{$p$-divisible groups}, Proceedings of a
  Conference on Local Fields, Driebergen (1966), Springer.
\item[{[34]}] J. Tate, \emph{Rigid analytic spaces} (mimeographed by
  IHES).
\end{enumerate}



}{\begin{enumerate}
\item[{[1]}] M. Artin, \emph{Algebraization of formal moduli I}
  (forthcoming).
\item[{[2]}] H. Cartan, \emph{Familles d'espaces complexes et fondements
  de la g\'{e}om\'{e}trie analytique}, S\'{e}minaire ENS 1960/61.
\item[{[3]}] P. Deligne, D. Mumford, \emph{The irreducibility of the space
  of curves of given genus}, Publ. Math. no.~36 (1969).
\item[{[4]}] P. Deligne, work in preparation on Hodge theory (forthcoming
  in Publ. Math., 1970).
\item[{[5]}] M. Demazure, J. Giraud, M. Raynaud, \emph{Sch\'{e}mas
  ab\'{e}liens}, S\'{e}minaire Orsay 1967/68, forthcoming.
\item[{[6]}] A. Grothendieck, \emph{Le Groupe de Brauer III}, in
  \emph{Dix Expos\'{e}s sur la cohomologie des sch\'{e}mas}, North-Holland
  Publishing Co.
\item[{[7]}] A. Grothendieck, \emph{Crystals and the De Rham cohomology of
  schemes} (notes by I. Coates and O. Jussila), in \emph{Dix expos\'{e}s
  sur la cohomologie des sch\'{e}mas}, North-Holland Publishing Co.
\item[{[8]}] A. Grothendieck, \emph{Un th\'{e}or\`{e}me sur les
  homomorphismes de sch\'{e}mas ab\'{e}liens}, Invent. Math. \textbf{2}
  (1966), pp.~59--78.
\item[{[9]}] P. A. Griffiths, \emph{Report on variation of Hodge
  structures} (forthcoming).
\item[{[10]}] M. Hakim, \emph{Topos annel\'{e}s et sch\'{e}mas relatifs},
  forthcoming in Grundlehren der Mathematik, Springer (1970).
\item[{[11]}] J. I. Igusa, \emph{Abstract vanishing cycle theory},
  Proceedings of the Japan Academy \textbf{34}, no.~9 (1958),
  pp.~589--593.
\item[{[12]}] R. Kiehl, \emph{Der Endlichkeitssatz f\"{u}r eigentliche
  Abbildungen}, Invent. Math. \textbf{2} (1967), pp.~191--214.
\item[{[13]}] R. Kiehl.
\item[{[14]}] S. L. Kleiman, \emph{Algebraic cycles and the Weil
  conjectures}, in \emph{Dix expos\'{e}s sur la cohomologie des sch\'{e}mas},
  North-Holland Publishing Co.
\item[{[15]}] S. Lang, \emph{Abelian varieties}, Interscience Publishers,
  no.~7 (1959).
\item[{[16]}] S. Lang and A. N\'{e}ron, \emph{Rational points of abelian
  varieties over function fields}, Amer. J. Math. \textbf{81}, no.~1
  (1959), pp.~95--118.
\item[{[16bis]}] U. I. Manin, \emph{Th\'{e}orie des groupes formels
  commutatifs}, Uspekhi Mat. Nauk \textbf{XVIII} (1963), pp.~4--90.
\item[{[17]}] D. Mumford, \emph{Geometric invariant theory}, Ergebnisse
  der Mathematik 34, Springer (1965).
\item[{[18]}] D. Mumford, \emph{Abstract Theta functions}, Advanced
  Science Seminar in Algebraic Geometry, Bowdoin College, 1967 (notes by
  H. Pittie).
\item[{[19]}] A. N\'{e}ron, \emph{Mod\`{e}les minimaux des vari\'{e}t\'{e}s
  ab\'{e}liennes sur les corps locaux et globaux}, Publ. Math. 21 (1964),
  pp.~5--128.
\item[{[20]}] M. Raynaud, \emph{Caract\'{e}ristique d'Euler--Poincar\'{e}
  d'un faisceau et cohomologie des vari\'{e}t\'{e}s ab\'{e}liennes}, in
  \emph{Dix expos\'{e}s sur la cohomologie des sch\'{e}mas}, North-Holland
  Publishing Co.
\item[{[21]}] M. Raynaud, \emph{Mod\`{e}les de N\'{e}ron}, C. R. Acad.
  Sci. Paris \textbf{262} (February 1966), pp.~413--416.
\item[{[22]}] M. Raynaud, \emph{Sp\'{e}cialisation du foncteur de Picard},
  C. R. Acad. Sci. \textbf{264}, pp.~941--943 and 1001--1004.
\item[{[22bis]}] M. Raynaud, \emph{Sp\'{e}cialisation du foncteur de
  Picard}, forthcoming in Publ. Math.
\item[{[23]}] M. Raynaud, \emph{Faisceaux amples sur les sch\'{e}mas en
  groupes et les espaces homog\`{e}nes}, Paris thesis, 1968 (forthcoming
  in Lecture Notes, Springer, 1970).
\item[{[24]}] M. Raynaud.
\item[{[25]}] M. Rosenlicht, \emph{Some basic theorems on algebraic
  groups}, Amer. J. Math. \textbf{78}, no.~2 (1956), pp.~401--443.
\item[{[25bis]}] J.-P. Serre, \emph{Quelques propri\'{e}t\'{e}s des
  vari\'{e}t\'{e}s ab\'{e}liennes en caract\'{e}ristique $p$}, Amer. J.
  Math. \textbf{80} (1958), pp.~715--739.
\item[{[26]}] J.-P. Serre, \emph{Corps locaux}, Act. Sci. Ind. 1296
  (1962), Hermann, Paris.
\item[{[27]}] J.-P. Serre, \emph{Sur les corps locaux \`{a} corps
  r\'{e}siduel alg\'{e}briquement clos}, Bull. Soc. Math. France
  \textbf{89} (1961), pp.~105--154.
\item[{[28]}] J.-P. Serre, \emph{Abelian $\ell$-adic representations and
  elliptic curves}, Benjamin Inc. (1968).
\item[{[29]}] J.-P. Serre and J. Tate, \emph{Good reduction of abelian
  varieties}, Ann. of Math. \textbf{88}, no.~3 (1968), pp.~492--517.
\item[{[30]}] J.-P. Serre and J. Tate, \emph{Elliptic curves and formal
  groups} (notes of the Summer Institute on Algebraic Geometry, Woods
  Hole, 1964).
\item[{[31]}] J.-P. Serre, \emph{Morphismes universels et vari\'{e}t\'{e}
  d'Albanese}, S\'{e}minaire Chevalley 1958/59, Expos\'{e} 10.
\item[{[32]}] J. Tate, \emph{WC-groups over $p$-adic fields}, S\'{e}minaire
  Bourbaki no.~156 (1957).
\item[{[33]}] J. Tate, \emph{$p$-divisible groups}, Proceedings of a
  Conference on Local Fields, Driebergen (1966), Springer.
\item[{[34]}] J. Tate, \emph{Rigid analytic spaces} (mimeographed by
  IHES).
\end{enumerate}



}{E7I-IX-134}\markboth{Exposé IX}{Exposé IX}

\ifnum\SGAform=1
\clearpage\section*{Critical apparatus}\begingroup\small\raggedright\sloppy
Readings are given in the exact TeX notation of the two English witnesses. Decimal suffixes identify the individual differences in each marked passage.

\par\medskip\noindent\hypertarget{app-E7I-IX-001.1}{}\textbf{E7I-IX-001.1}\quad printed folio 319.\par
\noindent Printed reading: {\ttfamily\footnotesize X\allowbreak{}}\par
\noindent Corrected reading: {\ttfamily\footnotesize K\allowbreak{}}\par
\noindent Source: {\ttfamily\footnotesize fr/sga7i/\{diplomatic\allowbreak{},current\}/expose\_IX\_\allowbreak{}body.tex;\allowbreak{} folio\allowbreak{} 319;\allowbreak{} formula\allowbreak{} 94;\allowbreak{} port/decisions/decis\allowbreak{}ions\_en7I\_r4\_001.jso\allowbreak{}n\#C1033;\allowbreak{} fixes/FIXLOG\_7I\_CURR\allowbreak{}ENT.jsonl\#C1033;\allowbreak{} 7I:1117\allowbreak{}}\par

\par\medskip\noindent\hypertarget{app-E7I-IX-002.1}{}\textbf{E7I-IX-002.1}\quad printed folio 322.\par
\noindent Printed reading: \emph{(empty)}\par
\noindent Corrected reading: {\ttfamily\footnotesize '\allowbreak{}}\par
\noindent Source: {\ttfamily\footnotesize fr/sga7i/\{diplomatic\allowbreak{},current\}/expose\_IX\_\allowbreak{}body.tex;\allowbreak{} folio\allowbreak{} 322;\allowbreak{} formula\allowbreak{} 170\allowbreak{}}\par

\par\medskip\noindent\hypertarget{app-E7I-IX-003.1}{}\textbf{E7I-IX-003.1}\quad printed folio 324.\par
\noindent Printed reading: {\ttfamily\footnotesize \textbackslash{}Phi'\allowbreak{}}\par
\noindent Corrected reading: {\ttfamily\footnotesize F\allowbreak{}}\par
\noindent Source: {\ttfamily\footnotesize fr/sga7i/\{diplomatic\allowbreak{},current\}/expose\_IX\_\allowbreak{}body.tex;\allowbreak{} folio\allowbreak{} 324\allowbreak{}}\par

\par\medskip\noindent\hypertarget{app-E7I-IX-003.2}{}\textbf{E7I-IX-003.2}\quad printed folio 324.\par
\noindent Printed reading: {\ttfamily\footnotesize \textbackslash{}Psi\allowbreak{}}\par
\noindent Corrected reading: {\ttfamily\footnotesize F\allowbreak{}}\par
\noindent Source: {\ttfamily\footnotesize fr/sga7i/\{diplomatic\allowbreak{},current\}/expose\_IX\_\allowbreak{}body.tex;\allowbreak{} folio\allowbreak{} 324\allowbreak{}}\par

\par\medskip\noindent\hypertarget{app-E7I-IX-003.3}{}\textbf{E7I-IX-003.3}\quad printed folio 324.\par
\noindent Printed reading: {\ttfamily\footnotesize \textbackslash{}Psi\allowbreak{}}\par
\noindent Corrected reading: {\ttfamily\footnotesize \textbackslash{}Phi'\allowbreak{}}\par
\noindent Source: {\ttfamily\footnotesize fr/sga7i/\{diplomatic\allowbreak{},current\}/expose\_IX\_\allowbreak{}body.tex;\allowbreak{} folio\allowbreak{} 324;\allowbreak{} fr/sga7i/\{diplomatic\allowbreak{},current\}/expose\_IX\_\allowbreak{}body.tex;\allowbreak{} folio\allowbreak{} 324;\allowbreak{} formula\allowbreak{} 205\allowbreak{}}\par

\par\medskip\noindent\hypertarget{app-E7I-IX-004.1}{}\textbf{E7I-IX-004.1}\quad printed folio 328.\par
\noindent Printed reading: {\ttfamily\footnotesize 5\allowbreak{}}\par
\noindent Corrected reading: {\ttfamily\footnotesize 6\allowbreak{}}\par
\noindent Source: {\ttfamily\footnotesize port/decisions/decis\allowbreak{}ions\_en7I\_r2\_003.jso\allowbreak{}n\#C831;\allowbreak{} fixes/FIXLOG\_7I\_CURR\allowbreak{}ENT.jsonl\#C831;\allowbreak{} 7I:547\allowbreak{}}\par

\par\medskip\noindent\hypertarget{app-E7I-IX-005.1}{}\textbf{E7I-IX-005.1}\quad printed folio 329.\par
\noindent Printed reading: \emph{(empty)}\par
\noindent Corrected reading: {\ttfamily\footnotesize 2.\allowbreak{}}\par
\noindent Source: {\ttfamily\footnotesize 7I:549;\allowbreak{} 7I:1111;\allowbreak{} WITNESS-F7I-IX-014.2\allowbreak{};\allowbreak{} fixes/MISPRINTS\_7I.j\allowbreak{}son\#549;\allowbreak{} fixes/MISPRINTS\_7I.j\allowbreak{}son\#1111;\allowbreak{} fr/sga7i/diplomatic/\allowbreak{}expose\_IX\_body.tex;\allowbreak{} fr/sga7i/current/exp\allowbreak{}ose\_IX\_body.tex\allowbreak{}}\par

\par\medskip\noindent\hypertarget{app-E7I-IX-005.2}{}\textbf{E7I-IX-005.2}\quad printed folio 329.\par
\noindent Printed reading: {\ttfamily\footnotesize .5\allowbreak{}}\par
\noindent Corrected reading: \emph{(empty)}\par
\noindent Source: {\ttfamily\footnotesize 7I:549;\allowbreak{} 7I:1111;\allowbreak{} WITNESS-F7I-IX-014.2\allowbreak{};\allowbreak{} fixes/MISPRINTS\_7I.j\allowbreak{}son\#549;\allowbreak{} fixes/MISPRINTS\_7I.j\allowbreak{}son\#1111;\allowbreak{} fr/sga7i/diplomatic/\allowbreak{}expose\_IX\_body.tex;\allowbreak{} fr/sga7i/current/exp\allowbreak{}ose\_IX\_body.tex\allowbreak{}}\par

\par\medskip\noindent\hypertarget{app-E7I-IX-006.1}{}\textbf{E7I-IX-006.1}\quad printed folio 333.\par
\noindent Printed reading: {\ttfamily\footnotesize 6\allowbreak{}}\par
\noindent Corrected reading: {\ttfamily\footnotesize 2\allowbreak{}}\par
\noindent Source: {\ttfamily\footnotesize 7I:551;\allowbreak{} 7I:1119;\allowbreak{} fixes/MISPRINTS\_7I.j\allowbreak{}son\#551;\allowbreak{} fixes/MISPRINTS\_7I.j\allowbreak{}son\#1119\allowbreak{}}\par

\par\medskip\noindent\hypertarget{app-E7I-IX-007.1}{}\textbf{E7I-IX-007.1}\quad printed folio 337.\par
\noindent Printed reading: {\ttfamily\footnotesize (\allowbreak{}}\par
\noindent Corrected reading: \emph{(empty)}\par
\noindent Source: {\ttfamily\footnotesize fr/sga7i/\{diplomatic\allowbreak{},current\}/expose\_IX\_\allowbreak{}body.tex;\allowbreak{} folio\allowbreak{} 337;\allowbreak{} formula\allowbreak{} 600;\allowbreak{} port/decisions/decis\allowbreak{}ions\_en7I\_r2\_003.jso\allowbreak{}n\#C838;\allowbreak{} fixes/FIXLOG\_7I\_CURR\allowbreak{}ENT.jsonl\#C838;\allowbreak{} 7I:556\allowbreak{}}\par

\par\medskip\noindent\hypertarget{app-E7I-IX-007.2}{}\textbf{E7I-IX-007.2}\quad printed folio 337.\par
\noindent Printed reading: {\ttfamily\footnotesize 1\allowbreak{}}\par
\noindent Corrected reading: {\ttfamily\footnotesize 2\allowbreak{}}\par
\noindent Source: {\ttfamily\footnotesize 7I:1133;\allowbreak{} fixes/MISPRINTS\_7I.j\allowbreak{}son\#1133\allowbreak{}}\par

\par\medskip\noindent\hypertarget{app-E7I-IX-007.3}{}\textbf{E7I-IX-007.3}\quad printed folio 337.\par
\noindent Printed reading: {\ttfamily\footnotesize 6\allowbreak{}}\par
\noindent Corrected reading: {\ttfamily\footnotesize 2\allowbreak{}}\par
\noindent Source: {\ttfamily\footnotesize 7I:1133;\allowbreak{} fixes/MISPRINTS\_7I.j\allowbreak{}son\#1133\allowbreak{}}\par

\par\medskip\noindent\hypertarget{app-E7I-IX-008.1}{}\textbf{E7I-IX-008.1}\quad printed folio 338.\par
\noindent Printed reading: {\ttfamily\footnotesize ;\allowbreak{}}\par
\noindent Corrected reading: {\ttfamily\footnotesize ,\allowbreak{}}\par
\noindent Source: {\ttfamily\footnotesize fr/sga7i/\{diplomatic\allowbreak{},current\}/expose\_IX\_\allowbreak{}body.tex;\allowbreak{} folio\allowbreak{} 338;\allowbreak{} formula\allowbreak{} 625\allowbreak{}}\par

\par\medskip\noindent\hypertarget{app-E7I-IX-008.2}{}\textbf{E7I-IX-008.2}\quad printed folio 338.\par
\noindent Printed reading: {\ttfamily\footnotesize w\allowbreak{}}\par
\noindent Corrected reading: {\ttfamily\footnotesize W\allowbreak{}}\par
\noindent Source: {\ttfamily\footnotesize fr/sga7i/\{diplomatic\allowbreak{},current\}/expose\_IX\_\allowbreak{}body.tex;\allowbreak{} folio\allowbreak{} 338;\allowbreak{} formula\allowbreak{} 630\allowbreak{}}\par

\par\medskip\noindent\hypertarget{app-E7I-IX-008.3}{}\textbf{E7I-IX-008.3}\quad printed folio 339.\par
\noindent Printed reading: \emph{(empty)}\par
\noindent Corrected reading: {\ttfamily\footnotesize '\allowbreak{}}\par
\noindent Source: {\ttfamily\footnotesize ED-EN-FORMULA-G055;\allowbreak{} fr/sga7i/\{diplomatic\allowbreak{},current\}/expose\_IX\_\allowbreak{}body.tex;\allowbreak{} folio\allowbreak{} 339\allowbreak{}}\par

\par\medskip\noindent\hypertarget{app-E7I-IX-008.4}{}\textbf{E7I-IX-008.4}\quad printed folio 340.\par
\noindent Printed reading: {\ttfamily\footnotesize 1\allowbreak{}}\par
\noindent Corrected reading: {\ttfamily\footnotesize 2\allowbreak{}}\par
\noindent Source: {\ttfamily\footnotesize ED-EN-NONMATH-I-IX34\allowbreak{}0;\allowbreak{} WITNESS-F7I-IX-023.2\allowbreak{};\allowbreak{} WITNESS-F7I-IX-023.3\allowbreak{};\allowbreak{} fr/sga7i/diplomatic/\allowbreak{}expose\_IX\_body.tex;\allowbreak{} fr/sga7i/current/exp\allowbreak{}ose\_IX\_body.tex;\allowbreak{} F7I-IX-023.2,\allowbreak{} F7I-IX-023.3\allowbreak{}}\par

\par\medskip\noindent\hypertarget{app-E7I-IX-008.5}{}\textbf{E7I-IX-008.5}\quad printed folio 340.\par
\noindent Printed reading: {\ttfamily\footnotesize 10\allowbreak{}}\par
\noindent Corrected reading: {\ttfamily\footnotesize 4.7\allowbreak{}}\par
\noindent Source: {\ttfamily\footnotesize ED-EN-NONMATH-I-IX34\allowbreak{}0;\allowbreak{} WITNESS-F7I-IX-023.2\allowbreak{};\allowbreak{} WITNESS-F7I-IX-023.3\allowbreak{};\allowbreak{} fr/sga7i/diplomatic/\allowbreak{}expose\_IX\_body.tex;\allowbreak{} fr/sga7i/current/exp\allowbreak{}ose\_IX\_body.tex;\allowbreak{} F7I-IX-023.2,\allowbreak{} F7I-IX-023.3\allowbreak{}}\par

\par\medskip\noindent\hypertarget{app-E7I-IX-008.6}{}\textbf{E7I-IX-008.6}\quad printed folio 341.\par
\noindent Printed reading: \emph{(empty)}\par
\noindent Corrected reading: {\ttfamily\footnotesize \textasciicircum{}\textbackslash{}perp\allowbreak{}}\par
\noindent Source: {\ttfamily\footnotesize ED-EN-PORT-WITNESS-7\allowbreak{}I-C844;\allowbreak{} fr/sga7i/diplomatic/\allowbreak{}expose\_IX\_body.tex;\allowbreak{} folio\allowbreak{} 341;\allowbreak{} port/PORT\_EN7I\_LOG.j\allowbreak{}sonl\#C844\#C844;\allowbreak{} 7I:565\allowbreak{}}\par

\par\medskip\noindent\hypertarget{app-E7I-IX-008.7}{}\textbf{E7I-IX-008.7}\quad printed folio 341.\par
\noindent Printed reading: \emph{(empty)}\par
\noindent Corrected reading: {\ttfamily\footnotesize \textasciicircum{}\textbackslash{}perp\allowbreak{}}\par
\noindent Source: {\ttfamily\footnotesize ED-EN-PORT-WITNESS-7\allowbreak{}I-C844;\allowbreak{} fr/sga7i/diplomatic/\allowbreak{}expose\_IX\_body.tex;\allowbreak{} folio\allowbreak{} 341;\allowbreak{} port/PORT\_EN7I\_LOG.j\allowbreak{}sonl\#C844\#C844;\allowbreak{} 7I:565\allowbreak{}}\par

\par\medskip\noindent\hypertarget{app-E7I-IX-008.8}{}\textbf{E7I-IX-008.8}\quad printed folio 341.\par
\noindent Printed reading: {\ttfamily\footnotesize n\allowbreak{}}\par
\noindent Corrected reading: {\ttfamily\footnotesize 2n\allowbreak{}}\par
\noindent Source: {\ttfamily\footnotesize fr/sga7i/\{diplomatic\allowbreak{},current\}/expose\_IX\_\allowbreak{}body.tex;\allowbreak{} folio\allowbreak{} 341;\allowbreak{} formula\allowbreak{} 681;\allowbreak{} port/PORT\_EN7I\_LOG.j\allowbreak{}sonl\#C1034;\allowbreak{} fixes/FIXLOG\_7I\_CURR\allowbreak{}ENT.jsonl\#C1034;\allowbreak{} 7I:1124\allowbreak{}}\par

\par\medskip\noindent\hypertarget{app-E7I-IX-009.1}{}\textbf{E7I-IX-009.1}\quad printed folio 343.\par
\noindent Printed reading: {\ttfamily\footnotesize S\allowbreak{}}\par
\noindent Corrected reading: {\ttfamily\footnotesize s\allowbreak{}}\par
\noindent Source: {\ttfamily\footnotesize fr/sga7i/\{diplomatic\allowbreak{},current\}/expose\_IX\_\allowbreak{}body.tex;\allowbreak{} folio\allowbreak{} 343;\allowbreak{} formula\allowbreak{} 708;\allowbreak{} port/decisions/decis\allowbreak{}ions\_en7I\_r2\_003.jso\allowbreak{}n\#C845;\allowbreak{} fixes/FIXLOG\_7I\_CURR\allowbreak{}ENT.jsonl\#C845;\allowbreak{} 7I:566\allowbreak{}}\par

\par\medskip\noindent\hypertarget{app-E7I-IX-010.1}{}\textbf{E7I-IX-010.1}\quad printed folio 345.\par
\noindent Printed reading: {\ttfamily\footnotesize s\allowbreak{}}\par
\noindent Corrected reading: {\ttfamily\footnotesize red\allowbreak{}}\par
\noindent Source: {\ttfamily\footnotesize fr/sga7i/\{diplomatic\allowbreak{},current\}/expose\_IX\_\allowbreak{}body.tex;\allowbreak{} folio\allowbreak{} 345;\allowbreak{} formula\allowbreak{} 775;\allowbreak{} port/decisions/decis\allowbreak{}ions\_en7I\_r2\_003.jso\allowbreak{}n\#C848;\allowbreak{} fixes/FIXLOG\_7I\_CURR\allowbreak{}ENT.jsonl\#C848;\allowbreak{} 7I:569\allowbreak{}}\par

\par\medskip\noindent\hypertarget{app-E7I-IX-011.1}{}\textbf{E7I-IX-011.1}\quad printed folio 346.\par
\noindent Printed reading: {\ttfamily\footnotesize T\allowbreak{}}\par
\noindent Corrected reading: {\ttfamily\footnotesize S\allowbreak{}}\par
\noindent Source: {\ttfamily\footnotesize fr/sga7i/\{diplomatic\allowbreak{},current\}/expose\_IX\_\allowbreak{}body.tex;\allowbreak{} folio\allowbreak{} 346;\allowbreak{} formula\allowbreak{} 812\allowbreak{}}\par

\par\medskip\noindent\hypertarget{app-E7I-IX-012.1}{}\textbf{E7I-IX-012.1}\quad printed folio 347.\par
\noindent Printed reading: \emph{(empty)}\par
\noindent Corrected reading: {\ttfamily\footnotesize 3\allowbreak{}}\par
\noindent Source: {\ttfamily\footnotesize port/decisions/decis\allowbreak{}ions\_en7I\_r2\_003.jso\allowbreak{}n\#C850;\allowbreak{} fixes/FIXLOG\_7I\_CURR\allowbreak{}ENT.jsonl\#C850;\allowbreak{} 7I:572\allowbreak{}}\par

\par\medskip\noindent\hypertarget{app-E7I-IX-012.2}{}\textbf{E7I-IX-012.2}\quad printed folio 347.\par
\noindent Printed reading: {\ttfamily\footnotesize VII\allowbreak{}}\par
\noindent Corrected reading: {\ttfamily\footnotesize XVII\allowbreak{}}\par
\noindent Source: {\ttfamily\footnotesize ED-EN-PORT-WITNESS-7\allowbreak{}I-C1035;\allowbreak{} fr/sga7i/diplomatic/\allowbreak{}expose\_IX\_body.tex;\allowbreak{} folio\allowbreak{} 347;\allowbreak{} port/decisions/decis\allowbreak{}ions\_en7I\_r2\_003.jso\allowbreak{}n\#C851;\allowbreak{} 7I:573\allowbreak{}}\par

\par\medskip\noindent\hypertarget{app-E7I-IX-012.3}{}\textbf{E7I-IX-012.3}\quad printed folio 347.\par
\noindent Printed reading: \emph{(empty)}\par
\noindent Corrected reading: {\ttfamily\footnotesize )\allowbreak{}}\par
\noindent Source: {\ttfamily\footnotesize ED-EN-PORT-WITNESS-7\allowbreak{}I-C1035;\allowbreak{} fr/sga7i/diplomatic/\allowbreak{}expose\_IX\_body.tex;\allowbreak{} folio\allowbreak{} 347;\allowbreak{} port/decisions/decis\allowbreak{}ions\_en7I\_r2\_003.jso\allowbreak{}n\#C851;\allowbreak{} 7I:573\allowbreak{}}\par

\par\medskip\noindent\hypertarget{app-E7I-IX-013.1}{}\textbf{E7I-IX-013.1}\quad printed folio 350.\par
\noindent Printed reading: \emph{(1 espace(s), 0 retour(s) de ligne dans le témoin TeX)}\par
\noindent Corrected reading: {\ttfamily\footnotesize \textasciicircum{}\textbackslash{}perp\allowbreak{}}\par
\noindent Source: {\ttfamily\footnotesize fr/sga7i/\{diplomatic\allowbreak{},current\}/expose\_IX\_\allowbreak{}body.tex;\allowbreak{} folio\allowbreak{} 350;\allowbreak{} formula\allowbreak{} 906\allowbreak{}}\par

\par\medskip\noindent\hypertarget{app-E7I-IX-013.2}{}\textbf{E7I-IX-013.2}\quad printed folio 350.\par
\noindent Printed reading: {\ttfamily\footnotesize \textasciicircum{}\textbackslash{}perp\allowbreak{}}\par
\noindent Corrected reading: \emph{(empty)}\par
\noindent Source: {\ttfamily\footnotesize fr/sga7i/\{diplomatic\allowbreak{},current\}/expose\_IX\_\allowbreak{}body.tex;\allowbreak{} folio\allowbreak{} 350;\allowbreak{} formula\allowbreak{} 906\allowbreak{}}\par

\par\medskip\noindent\hypertarget{app-E7I-IX-014.1}{}\textbf{E7I-IX-014.1}\quad printed folio 350.\par
\noindent Printed reading: \emph{(empty)}\par
\noindent Corrected reading: {\ttfamily\footnotesize \textasciicircum{}\textbackslash{}perp\allowbreak{}}\par
\noindent Source: {\ttfamily\footnotesize fr/sga7i/\{diplomatic\allowbreak{},current\}/expose\_IX\_\allowbreak{}body.tex;\allowbreak{} folio\allowbreak{} 350;\allowbreak{} formula\allowbreak{} 919;\allowbreak{} port/decisions/decis\allowbreak{}ions\_en7I\_r2\_003.jso\allowbreak{}n\#C853;\allowbreak{} fixes/FIXLOG\_7I\_CURR\allowbreak{}ENT.jsonl\#C853;\allowbreak{} 7I:576\allowbreak{}}\par

\par\medskip\noindent\hypertarget{app-E7I-IX-014.2}{}\textbf{E7I-IX-014.2}\quad printed folio 350.\par
\noindent Printed reading: \emph{(empty)}\par
\noindent Corrected reading: {\ttfamily\footnotesize /U\allowbreak{}}\par
\noindent Source: {\ttfamily\footnotesize fr/sga7i/\{diplomatic\allowbreak{},current\}/expose\_IX\_\allowbreak{}body.tex;\allowbreak{} folio\allowbreak{} 350;\allowbreak{} formula\allowbreak{} 923\allowbreak{}}\par

\par\medskip\noindent\hypertarget{app-E7I-IX-014.3}{}\textbf{E7I-IX-014.3}\quad printed folio 350.\par
\noindent Printed reading: \emph{(empty)}\par
\noindent Corrected reading: {\ttfamily\footnotesize U/\allowbreak{}}\par
\noindent Source: {\ttfamily\footnotesize fr/sga7i/\{diplomatic\allowbreak{},current\}/expose\_IX\_\allowbreak{}body.tex;\allowbreak{} folio\allowbreak{} 350;\allowbreak{} formula\allowbreak{} 923\allowbreak{}}\par

\par\medskip\noindent\hypertarget{app-E7I-IX-015.1}{}\textbf{E7I-IX-015.1}\quad printed folio 351.\par
\noindent Printed reading: {\ttfamily\footnotesize )\$\allowbreak{} \$(I\textbackslash{}qquad\allowbreak{}}\par
\noindent Corrected reading: \emph{(empty)}\par
\noindent Source: {\ttfamily\footnotesize ADOPT-F7I-IX-034.1;\allowbreak{} fr/sga7i/diplomatic/\allowbreak{}expose\_IX\_body.tex;\allowbreak{} fr/sga7i/current/exp\allowbreak{}ose\_IX\_body.tex\allowbreak{}}\par

\par\medskip\noindent\hypertarget{app-E7I-IX-016.1}{}\textbf{E7I-IX-016.1}\quad printed folio 353.\par
\noindent Printed reading: {\ttfamily\footnotesize v\allowbreak{}}\par
\noindent Corrected reading: {\ttfamily\footnotesize vi\allowbreak{}}\par
\noindent Source: {\ttfamily\footnotesize 7I:581;\allowbreak{} WITNESS-F7I-IX-036.2\allowbreak{};\allowbreak{} fixes/MISPRINTS\_7I.j\allowbreak{}son\#581;\allowbreak{} fr/sga7i/diplomatic/\allowbreak{}expose\_IX\_body.tex;\allowbreak{} fr/sga7i/current/exp\allowbreak{}ose\_IX\_body.tex\allowbreak{}}\par

\par\medskip\noindent\hypertarget{app-E7I-IX-016.2}{}\textbf{E7I-IX-016.2}\quad printed folio 353.\par
\noindent Printed reading: {\ttfamily\footnotesize vi\allowbreak{}}\par
\noindent Corrected reading: {\ttfamily\footnotesize iv\allowbreak{}}\par
\noindent Source: {\ttfamily\footnotesize 7I:581;\allowbreak{} WITNESS-F7I-IX-036.2\allowbreak{};\allowbreak{} fixes/MISPRINTS\_7I.j\allowbreak{}son\#581;\allowbreak{} fr/sga7i/diplomatic/\allowbreak{}expose\_IX\_body.tex;\allowbreak{} fr/sga7i/current/exp\allowbreak{}ose\_IX\_body.tex\allowbreak{}}\par

\par\medskip\noindent\hypertarget{app-E7I-IX-017.1}{}\textbf{E7I-IX-017.1}\quad printed folio 354.\par
\noindent Printed reading: {\ttfamily\footnotesize \textbackslash{}ne\allowbreak{} }\par
\noindent Corrected reading: {\ttfamily\footnotesize =\allowbreak{}}\par
\noindent Source: {\ttfamily\footnotesize ED-EN-FORMULA-G057;\allowbreak{} fr/sga7i/\{diplomatic\allowbreak{},current\}/expose\_IX\_\allowbreak{}body.tex;\allowbreak{} folio\allowbreak{} 354\allowbreak{}}\par

\par\medskip\noindent\hypertarget{app-E7I-IX-017.2}{}\textbf{E7I-IX-017.2}\quad printed folio 355.\par
\noindent Printed reading: {\ttfamily\footnotesize '\allowbreak{}}\par
\noindent Corrected reading: \emph{(empty)}\par
\noindent Source: {\ttfamily\footnotesize EN-REPAIR-C857;\allowbreak{} fixes/MISPRINTS\_7I.j\allowbreak{}son;\allowbreak{} fixes/FIXLOG\_7I\_CURR\allowbreak{}ENT.jsonl\#C857\allowbreak{}}\par

\par\medskip\noindent\hypertarget{app-E7I-IX-018.1}{}\textbf{E7I-IX-018.1}\quad printed folio 356.\par
\noindent Printed reading: {\ttfamily\footnotesize 4\allowbreak{}}\par
\noindent Corrected reading: {\ttfamily\footnotesize 3\allowbreak{}}\par
\noindent Source: {\ttfamily\footnotesize 7I:584;\allowbreak{} fixes/MISPRINTS\_7I.j\allowbreak{}son\#584\allowbreak{}}\par

\par\medskip\noindent\hypertarget{app-E7I-IX-018.2}{}\textbf{E7I-IX-018.2}\quad printed folio 357.\par
\noindent Printed reading: \emph{(empty)}\par
\noindent Corrected reading: {\ttfamily\footnotesize  Expos\textbackslash{}'e\allowbreak{}}\par
\noindent Source: {\ttfamily\footnotesize ADOPT-F7I-IX-040.3;\allowbreak{} ADOPT-F7I-IX-040.4;\allowbreak{} fr/sga7i/diplomatic/\allowbreak{}expose\_IX\_body.tex;\allowbreak{} fr/sga7i/current/exp\allowbreak{}ose\_IX\_body.tex\allowbreak{}}\par

\par\medskip\noindent\hypertarget{app-E7I-IX-018.3}{}\textbf{E7I-IX-018.3}\quad printed folio 357.\par
\noindent Printed reading: {\ttfamily\footnotesize \textbackslash{}hspace\{2.5em\}\allowbreak{}}\par
\noindent Corrected reading: {\ttfamily\footnotesize  [reference\allowbreak{} omitted\allowbreak{} in\allowbreak{} the\allowbreak{} source]\allowbreak{}}\par
\noindent Source: {\ttfamily\footnotesize ADOPT-F7I-IX-040.3;\allowbreak{} ADOPT-F7I-IX-040.4;\allowbreak{} fr/sga7i/diplomatic/\allowbreak{}expose\_IX\_body.tex;\allowbreak{} fr/sga7i/current/exp\allowbreak{}ose\_IX\_body.tex\allowbreak{}}\par

\par\medskip\noindent\hypertarget{app-E7I-IX-019.1}{}\textbf{E7I-IX-019.1}\quad printed folio 358.\par
\noindent Printed reading: {\ttfamily\footnotesize \textasciicircum{}\{o\}\allowbreak{}}\par
\noindent Corrected reading: \emph{(empty)}\par
\noindent Source: {\ttfamily\footnotesize ED-EN-FORMULA-G059;\allowbreak{} fr/sga7i/\{diplomatic\allowbreak{},current\}/expose\_IX\_\allowbreak{}body.tex;\allowbreak{} folio\allowbreak{} 358\allowbreak{}}\par

\par\medskip\noindent\hypertarget{app-E7I-IX-020.1}{}\textbf{E7I-IX-020.1}\quad printed folio 367.\par
\noindent Printed reading: {\ttfamily\footnotesize T\allowbreak{}}\par
\noindent Corrected reading: {\ttfamily\footnotesize (1-nT)\allowbreak{}}\par
\noindent Source: {\ttfamily\footnotesize fr/sga7i/\{diplomatic\allowbreak{},current\}/expose\_IX\_\allowbreak{}body.tex;\allowbreak{} folio\allowbreak{} 367;\allowbreak{} formula\allowbreak{} 1512\allowbreak{}}\par

\par\medskip\noindent\hypertarget{app-E7I-IX-020.2}{}\textbf{E7I-IX-020.2}\quad printed folio 368.\par
\noindent Printed reading: {\ttfamily\footnotesize n\allowbreak{}}\par
\noindent Corrected reading: {\ttfamily\footnotesize m\allowbreak{}}\par
\noindent Source: {\ttfamily\footnotesize fr/sga7i/\{diplomatic\allowbreak{},current\}/expose\_IX\_\allowbreak{}body.tex;\allowbreak{} folio\allowbreak{} 368;\allowbreak{} formula\allowbreak{} 1523\allowbreak{}}\par

\par\medskip\noindent\hypertarget{app-E7I-IX-021.1}{}\textbf{E7I-IX-021.1}\quad printed folio 371.\par
\noindent Printed reading: \emph{(empty)}\par
\noindent Corrected reading: {\ttfamily\footnotesize '\allowbreak{}}\par
\noindent Source: {\ttfamily\footnotesize fr/sga7i/\{diplomatic\allowbreak{},current\}/expose\_IX\_\allowbreak{}body.tex;\allowbreak{} folio\allowbreak{} 371;\allowbreak{} formula\allowbreak{} 1630\allowbreak{}}\par

\par\medskip\noindent\hypertarget{app-E7I-IX-021.2}{}\textbf{E7I-IX-021.2}\quad printed folio 371.\par
\noindent Printed reading: {\ttfamily\footnotesize '\allowbreak{}}\par
\noindent Corrected reading: \emph{(empty)}\par
\noindent Source: {\ttfamily\footnotesize fr/sga7i/\{diplomatic\allowbreak{},current\}/expose\_IX\_\allowbreak{}body.tex;\allowbreak{} folio\allowbreak{} 371;\allowbreak{} formula\allowbreak{} 1630\allowbreak{}}\par

\par\medskip\noindent\hypertarget{app-E7I-IX-022.1}{}\textbf{E7I-IX-022.1}\quad printed folio 377.\par
\noindent Printed reading: {\ttfamily\footnotesize S\allowbreak{}}\par
\noindent Corrected reading: {\ttfamily\footnotesize K\allowbreak{}}\par
\noindent Source: {\ttfamily\footnotesize fr/sga7i/\{diplomatic\allowbreak{},current\}/expose\_IX\_\allowbreak{}body.tex;\allowbreak{} folio\allowbreak{} 377;\allowbreak{} formula\allowbreak{} 1730\allowbreak{}}\par

\par\medskip\noindent\hypertarget{app-E7I-IX-023.1}{}\textbf{E7I-IX-023.1}\quad printed folio 379.\par
\noindent Printed reading: \emph{(empty)}\par
\noindent Corrected reading: {\ttfamily\footnotesize \_K\allowbreak{}}\par
\noindent Source: {\ttfamily\footnotesize fr/sga7i/\{diplomatic\allowbreak{},current\}/expose\_IX\_\allowbreak{}body.tex;\allowbreak{} folio\allowbreak{} 379;\allowbreak{} formula\allowbreak{} 1774;\allowbreak{} ED-en7I-C874;\allowbreak{} OPEN\_ITEMS\_DECIDED.t\allowbreak{}sv\allowbreak{}}\par

\par\medskip\noindent\hypertarget{app-E7I-IX-024.1}{}\textbf{E7I-IX-024.1}\quad printed folio 380.\par
\noindent Printed reading: \emph{(empty)}\par
\noindent Corrected reading: {\ttfamily\footnotesize '\allowbreak{}}\par
\noindent Source: {\ttfamily\footnotesize fr/sga7i/\{diplomatic\allowbreak{},current\}/expose\_IX\_\allowbreak{}body.tex;\allowbreak{} folio\allowbreak{} 380;\allowbreak{} port/PORT\_EN7I\_LOG.j\allowbreak{}sonl\#C877;\allowbreak{} fixes/FIXLOG\_7I\_CURR\allowbreak{}ENT.jsonl\#C877;\allowbreak{} 7I:612\allowbreak{}}\par

\par\medskip\noindent\hypertarget{app-E7I-IX-024.2}{}\textbf{E7I-IX-024.2}\quad printed folio 380.\par
\noindent Printed reading: {\ttfamily\footnotesize '\allowbreak{}}\par
\noindent Corrected reading: \emph{(empty)}\par
\noindent Source: {\ttfamily\footnotesize fr/sga7i/\{diplomatic\allowbreak{},current\}/expose\_IX\_\allowbreak{}body.tex;\allowbreak{} folio\allowbreak{} 380;\allowbreak{} port/PORT\_EN7I\_LOG.j\allowbreak{}sonl\#C877;\allowbreak{} fixes/FIXLOG\_7I\_CURR\allowbreak{}ENT.jsonl\#C877;\allowbreak{} 7I:612\allowbreak{}}\par

\par\medskip\noindent\hypertarget{app-E7I-IX-025.1}{}\textbf{E7I-IX-025.1}\quad printed folio 383.\par
\noindent Printed reading: {\ttfamily\footnotesize connected\allowbreak{}}\par
\noindent Corrected reading: {\ttfamily\footnotesize \textbackslash{}'\{e\}tale\allowbreak{}}\par
\noindent Source: {\ttfamily\footnotesize port/decisions/decis\allowbreak{}ions\_en7I\_r2\_003.jso\allowbreak{}n\#C878;\allowbreak{} 7I:613\allowbreak{}}\par

\par\medskip\noindent\hypertarget{app-E7I-IX-026.1}{}\textbf{E7I-IX-026.1}\quad printed folio 386.\par
\noindent Printed reading: {\ttfamily\footnotesize X\allowbreak{}}\par
\noindent Corrected reading: {\ttfamily\footnotesize G\allowbreak{}}\par
\noindent Source: {\ttfamily\footnotesize port/decisions/decis\allowbreak{}ions\_en7I\_r2\_003.jso\allowbreak{}n\#C880;\allowbreak{} 7I:615;\allowbreak{} fr/sga7i/\{diplomatic\allowbreak{},current\}/expose\_IX\_\allowbreak{}body.tex;\allowbreak{} folio\allowbreak{} 386;\allowbreak{} formula\allowbreak{} 1964\allowbreak{}}\par

\par\medskip\noindent\hypertarget{app-E7I-IX-027.1}{}\textbf{E7I-IX-027.1}\quad printed folio 388.\par
\noindent Printed reading: \emph{(empty)}\par
\noindent Corrected reading: {\ttfamily\footnotesize \_K\allowbreak{}}\par
\noindent Source: {\ttfamily\footnotesize fr/sga7i/\{diplomatic\allowbreak{},current\}/expose\_IX\_\allowbreak{}body.tex;\allowbreak{} folio\allowbreak{} 388;\allowbreak{} formula\allowbreak{} 2028\allowbreak{}}\par

\par\medskip\noindent\hypertarget{app-E7I-IX-028.1}{}\textbf{E7I-IX-028.1}\quad printed folio 388.\par
\noindent Printed reading: \emph{(empty)}\par
\noindent Corrected reading: {\ttfamily\footnotesize \_K\allowbreak{}}\par
\noindent Source: {\ttfamily\footnotesize fr/sga7i/\{diplomatic\allowbreak{},current\}/expose\_IX\_\allowbreak{}body.tex;\allowbreak{} folio\allowbreak{} 388;\allowbreak{} formula\allowbreak{} 2044\allowbreak{}}\par

\par\medskip\noindent\hypertarget{app-E7I-IX-029.1}{}\textbf{E7I-IX-029.1}\quad printed folio 389.\par
\noindent Printed reading: {\ttfamily\footnotesize \textbackslash{}overline\allowbreak{} }\par
\noindent Corrected reading: \emph{(empty)}\par
\noindent Source: {\ttfamily\footnotesize fr/sga7i/\{diplomatic\allowbreak{},current\}/expose\_IX\_\allowbreak{}body.tex;\allowbreak{} folio\allowbreak{} 389;\allowbreak{} formula\allowbreak{} 2054\allowbreak{}}\par

\par\medskip\noindent\hypertarget{app-E7I-IX-030.1}{}\textbf{E7I-IX-030.1}\quad printed folio 390.\par
\noindent Printed reading: \emph{(empty)}\par
\noindent Corrected reading: {\ttfamily\footnotesize \textasciicircum{}\textbackslash{}times\allowbreak{}}\par
\noindent Source: {\ttfamily\footnotesize fr/sga7i/\{diplomatic\allowbreak{},current\}/expose\_IX\_\allowbreak{}body.tex;\allowbreak{} folio\allowbreak{} 390;\allowbreak{} formula\allowbreak{} 2106\allowbreak{}}\par

\par\medskip\noindent\hypertarget{app-E7I-IX-031.1}{}\textbf{E7I-IX-031.1}\quad printed folio 395.\par
\noindent Printed reading: {\ttfamily\footnotesize '\allowbreak{}}\par
\noindent Corrected reading: \emph{(empty)}\par
\noindent Source: {\ttfamily\footnotesize ED-EN-FORMULA-G053;\allowbreak{} fr/sga7i/\{diplomatic\allowbreak{},current\}/expose\_IX\_\allowbreak{}body.tex;\allowbreak{} folio\allowbreak{} 395\allowbreak{}}\par

\par\medskip\noindent\hypertarget{app-E7I-IX-031.2}{}\textbf{E7I-IX-031.2}\quad printed folio 395.\par
\noindent Printed reading: {\ttfamily\footnotesize S\allowbreak{}}\par
\noindent Corrected reading: {\ttfamily\footnotesize K\allowbreak{}}\par
\noindent Source: {\ttfamily\footnotesize ED-EN-FORMULA-G053;\allowbreak{} fr/sga7i/\{diplomatic\allowbreak{},current\}/expose\_IX\_\allowbreak{}body.tex;\allowbreak{} folio\allowbreak{} 395\allowbreak{}}\par

\par\medskip\noindent\hypertarget{app-E7I-IX-032.1}{}\textbf{E7I-IX-032.1}\quad printed folio 396.\par
\noindent Printed reading: \emph{(empty)}\par
\noindent Corrected reading: {\ttfamily\footnotesize '\allowbreak{}}\par
\noindent Source: {\ttfamily\footnotesize fr/sga7i/\{diplomatic\allowbreak{},current\}/expose\_IX\_\allowbreak{}body.tex;\allowbreak{} folio\allowbreak{} 396;\allowbreak{} formula\allowbreak{} 2331;\allowbreak{} port/decisions/decis\allowbreak{}ions\_en7I\_r2\_003.jso\allowbreak{}n\#C884;\allowbreak{} fixes/FIXLOG\_7I\_CURR\allowbreak{}ENT.jsonl\#C884;\allowbreak{} 7I:625\allowbreak{}}\par

\par\medskip\noindent\hypertarget{app-E7I-IX-032.2}{}\textbf{E7I-IX-032.2}\quad printed folio 396.\par
\noindent Printed reading: \emph{(empty)}\par
\noindent Corrected reading: {\ttfamily\footnotesize '\allowbreak{}}\par
\noindent Source: {\ttfamily\footnotesize fr/sga7i/\{diplomatic\allowbreak{},current\}/expose\_IX\_\allowbreak{}body.tex;\allowbreak{} folio\allowbreak{} 396;\allowbreak{} formula\allowbreak{} 2332;\allowbreak{} port/decisions/decis\allowbreak{}ions\_en7I\_r2\_003.jso\allowbreak{}n\#C884;\allowbreak{} fixes/FIXLOG\_7I\_CURR\allowbreak{}ENT.jsonl\#C884;\allowbreak{} 7I:625\allowbreak{}}\par

\par\medskip\noindent\hypertarget{app-E7I-IX-032.3}{}\textbf{E7I-IX-032.3}\quad printed folio 397.\par
\noindent Printed reading: {\ttfamily\footnotesize k\allowbreak{}}\par
\noindent Corrected reading: {\ttfamily\footnotesize K\allowbreak{}}\par
\noindent Source: {\ttfamily\footnotesize ED-EN-PORT-WITNESS-7\allowbreak{}I-C527;\allowbreak{} fr/sga7i/diplomatic/\allowbreak{}expose\_IX\_body.tex;\allowbreak{} folio\allowbreak{} 397;\allowbreak{} port/PORT\_EN7I\_LOG.j\allowbreak{}sonl\#C527\#C527;\allowbreak{} }\par

\par\medskip\noindent\hypertarget{app-E7I-IX-032.4}{}\textbf{E7I-IX-032.4}\quad printed folio 397.\par
\noindent Printed reading: {\ttfamily\footnotesize k\allowbreak{}}\par
\noindent Corrected reading: {\ttfamily\footnotesize K\allowbreak{}}\par
\noindent Source: {\ttfamily\footnotesize ED-EN-PORT-WITNESS-7\allowbreak{}I-C527;\allowbreak{} fr/sga7i/diplomatic/\allowbreak{}expose\_IX\_body.tex;\allowbreak{} folio\allowbreak{} 397;\allowbreak{} port/PORT\_EN7I\_LOG.j\allowbreak{}sonl\#C527\#C527;\allowbreak{} }\par

\par\medskip\noindent\hypertarget{app-E7I-IX-032.5}{}\textbf{E7I-IX-032.5}\quad printed folio 397.\par
\noindent Printed reading: \emph{(empty)}\par
\noindent Corrected reading: {\ttfamily\footnotesize '\allowbreak{}}\par
\noindent Source: {\ttfamily\footnotesize fr/sga7i/\{diplomatic\allowbreak{},current\}/expose\_IX\_\allowbreak{}body.tex;\allowbreak{} folio\allowbreak{} 397;\allowbreak{} formula\allowbreak{} 2347\allowbreak{}}\par

\par\medskip\noindent\hypertarget{app-E7I-IX-032.6}{}\textbf{E7I-IX-032.6}\quad printed folio 397.\par
\noindent Printed reading: {\ttfamily\footnotesize 12\allowbreak{}}\par
\noindent Corrected reading: {\ttfamily\footnotesize 13\allowbreak{}}\par
\noindent Source: {\ttfamily\footnotesize port/PORT\_EN7I\_LOG.j\allowbreak{}sonl\#C885;\allowbreak{} fixes/FIXLOG\_7I\_CURR\allowbreak{}ENT.jsonl\#C885;\allowbreak{} 7I:628\allowbreak{}}\par

\par\medskip\noindent\hypertarget{app-E7I-IX-032.7}{}\textbf{E7I-IX-032.7}\quad printed folio 397.\par
\noindent Printed reading: {\ttfamily\footnotesize a\allowbreak{}}\par
\noindent Corrected reading: {\ttfamily\footnotesize b\allowbreak{}}\par
\noindent Source: {\ttfamily\footnotesize port/PORT\_EN7I\_LOG.j\allowbreak{}sonl\#C885;\allowbreak{} fixes/FIXLOG\_7I\_CURR\allowbreak{}ENT.jsonl\#C885;\allowbreak{} 7I:628\allowbreak{}}\par

\par\medskip\noindent\hypertarget{app-E7I-IX-032.8}{}\textbf{E7I-IX-032.8}\quad printed folio 397.\par
\noindent Printed reading: {\ttfamily\footnotesize 12\allowbreak{}}\par
\noindent Corrected reading: {\ttfamily\footnotesize 13\allowbreak{}}\par
\noindent Source: {\ttfamily\footnotesize port/PORT\_EN7I\_LOG.j\allowbreak{}sonl\#C886;\allowbreak{} fixes/FIXLOG\_7I\_CURR\allowbreak{}ENT.jsonl\#C886;\allowbreak{} 7I:629\allowbreak{}}\par

\par\medskip\noindent\hypertarget{app-E7I-IX-032.9}{}\textbf{E7I-IX-032.9}\quad printed folio 397.\par
\noindent Printed reading: {\ttfamily\footnotesize a\allowbreak{}}\par
\noindent Corrected reading: {\ttfamily\footnotesize b\allowbreak{}}\par
\noindent Source: {\ttfamily\footnotesize port/PORT\_EN7I\_LOG.j\allowbreak{}sonl\#C886;\allowbreak{} fixes/FIXLOG\_7I\_CURR\allowbreak{}ENT.jsonl\#C886;\allowbreak{} 7I:629\allowbreak{}}\par

\par\medskip\noindent\hypertarget{app-E7I-IX-033.1}{}\textbf{E7I-IX-033.1}\quad printed folio 398.\par
\noindent Printed reading: {\ttfamily\footnotesize 6\allowbreak{}}\par
\noindent Corrected reading: {\ttfamily\footnotesize 5\allowbreak{}}\par
\noindent Source: {\ttfamily\footnotesize WITNESS-F7I-IX-074.2\allowbreak{};\allowbreak{} WITNESS-F7I-IX-074.3\allowbreak{};\allowbreak{} fr/sga7i/diplomatic/\allowbreak{}expose\_IX\_body.tex;\allowbreak{} fr/sga7i/current/exp\allowbreak{}ose\_IX\_body.tex\allowbreak{}}\par

\par\medskip\noindent\hypertarget{app-E7I-IX-033.2}{}\textbf{E7I-IX-033.2}\quad printed folio 398.\par
\noindent Printed reading: {\ttfamily\footnotesize 5\allowbreak{}}\par
\noindent Corrected reading: {\ttfamily\footnotesize 12\allowbreak{}}\par
\noindent Source: {\ttfamily\footnotesize WITNESS-F7I-IX-074.2\allowbreak{};\allowbreak{} WITNESS-F7I-IX-074.3\allowbreak{};\allowbreak{} fr/sga7i/diplomatic/\allowbreak{}expose\_IX\_body.tex;\allowbreak{} fr/sga7i/current/exp\allowbreak{}ose\_IX\_body.tex\allowbreak{}}\par

\par\medskip\noindent\hypertarget{app-E7I-IX-034.1}{}\textbf{E7I-IX-034.1}\quad printed folio 401.\par
\noindent Printed reading: {\ttfamily\footnotesize M\allowbreak{}}\par
\noindent Corrected reading: {\ttfamily\footnotesize L\allowbreak{}}\par
\noindent Source: {\ttfamily\footnotesize ED-EN-FORMULA-G061;\allowbreak{} fr/sga7i/\{diplomatic\allowbreak{},current\}/expose\_IX\_\allowbreak{}body.tex;\allowbreak{} folio\allowbreak{} 401\allowbreak{}}\par

\par\medskip\noindent\hypertarget{app-E7I-IX-034.2}{}\textbf{E7I-IX-034.2}\quad printed folio 401.\par
\noindent Printed reading: {\ttfamily\footnotesize G\allowbreak{}}\par
\noindent Corrected reading: {\ttfamily\footnotesize A\allowbreak{}}\par
\noindent Source: {\ttfamily\footnotesize fr/sga7i/\{diplomatic\allowbreak{},current\}/expose\_IX\_\allowbreak{}body.tex;\allowbreak{} folio\allowbreak{} 401;\allowbreak{} formula\allowbreak{} 2490\allowbreak{}}\par

\par\medskip\noindent\hypertarget{app-E7I-IX-034.3}{}\textbf{E7I-IX-034.3}\quad printed folio 401.\par
\noindent Printed reading: {\ttfamily\footnotesize G\allowbreak{}}\par
\noindent Corrected reading: {\ttfamily\footnotesize A\allowbreak{}}\par
\noindent Source: {\ttfamily\footnotesize fr/sga7i/\{diplomatic\allowbreak{},current\}/expose\_IX\_\allowbreak{}body.tex;\allowbreak{} folio\allowbreak{} 401;\allowbreak{} formula\allowbreak{} 2494\allowbreak{}}\par

\par\medskip\noindent\hypertarget{app-E7I-IX-034.4}{}\textbf{E7I-IX-034.4}\quad printed folio 401.\par
\noindent Printed reading: \emph{(empty)}\par
\noindent Corrected reading: {\ttfamily\footnotesize \textbackslash{}widehat\allowbreak{}}\par
\noindent Source: {\ttfamily\footnotesize fr/sga7i/\{diplomatic\allowbreak{},current\}/expose\_IX\_\allowbreak{}body.tex;\allowbreak{} folio\allowbreak{} 401;\allowbreak{} formula\allowbreak{} 2498;\allowbreak{} port/decisions/decis\allowbreak{}ions\_en7I\_r2\_003.jso\allowbreak{}n\#C890;\allowbreak{} fixes/FIXLOG\_7I\_CURR\allowbreak{}ENT.jsonl\#C890;\allowbreak{} 7I:827\allowbreak{}}\par

\par\medskip\noindent\hypertarget{app-E7I-IX-034.5}{}\textbf{E7I-IX-034.5}\quad printed folio 401.\par
\noindent Printed reading: \emph{(empty)}\par
\noindent Corrected reading: {\ttfamily\footnotesize \textbackslash{}widehat\{\allowbreak{}}\par
\noindent Source: {\ttfamily\footnotesize fr/sga7i/\{diplomatic\allowbreak{},current\}/expose\_IX\_\allowbreak{}body.tex;\allowbreak{} folio\allowbreak{} 401\allowbreak{}}\par

\par\medskip\noindent\hypertarget{app-E7I-IX-034.6}{}\textbf{E7I-IX-034.6}\quad printed folio 401.\par
\noindent Printed reading: \emph{(empty)}\par
\noindent Corrected reading: {\ttfamily\footnotesize \}\allowbreak{}}\par
\noindent Source: {\ttfamily\footnotesize fr/sga7i/\{diplomatic\allowbreak{},current\}/expose\_IX\_\allowbreak{}body.tex;\allowbreak{} folio\allowbreak{} 401\allowbreak{}}\par

\par\medskip\noindent\hypertarget{app-E7I-IX-035.1}{}\textbf{E7I-IX-035.1}\quad printed folio 403.\par
\noindent Printed reading: {\ttfamily\footnotesize 3\allowbreak{}}\par
\noindent Corrected reading: {\ttfamily\footnotesize 13\allowbreak{}}\par
\noindent Source: {\ttfamily\footnotesize port/PORT\_EN7I\_LOG.j\allowbreak{}sonl\#C893;\allowbreak{} fixes/FIXLOG\_7I\_CURR\allowbreak{}ENT.jsonl\#C893;\allowbreak{} 7I:833\allowbreak{}}\par

\par\medskip\noindent\hypertarget{app-E7I-IX-035.2}{}\textbf{E7I-IX-035.2}\quad printed folio 403.\par
\noindent Printed reading: {\ttfamily\footnotesize A\allowbreak{}}\par
\noindent Corrected reading: {\ttfamily\footnotesize V\allowbreak{}}\par
\noindent Source: {\ttfamily\footnotesize fr/sga7i/\{diplomatic\allowbreak{},current\}/expose\_IX\_\allowbreak{}body.tex;\allowbreak{} folio\allowbreak{} 403;\allowbreak{} ED-EN-FORMULA-G063;\allowbreak{} fr/sga7i/\{diplomatic\allowbreak{},current\}/expose\_IX\_\allowbreak{}body.tex;\allowbreak{} folio\allowbreak{} 403\allowbreak{}}\par

\par\medskip\noindent\hypertarget{app-E7I-IX-036.1}{}\textbf{E7I-IX-036.1}\quad printed folio 403.\par
\noindent Printed reading: \emph{(empty)}\par
\noindent Corrected reading: {\ttfamily\footnotesize \textbackslash{}mathcal\allowbreak{} }\par
\noindent Source: {\ttfamily\footnotesize ED-EN-FONT-G000;\allowbreak{} F7I-IX-078.2;\allowbreak{} fr/sga7i/\{diplomatic\allowbreak{},current\}/expose\_IX\_\allowbreak{}body.tex;\allowbreak{} folio\allowbreak{} 403;\allowbreak{} port/decisions/decis\allowbreak{}ions\_en7I\_r2\_003.jso\allowbreak{}n\#C892;\allowbreak{} 7I:642\allowbreak{}}\par

\par\medskip\noindent\hypertarget{app-E7I-IX-036.2}{}\textbf{E7I-IX-036.2}\quad printed folio 403.\par
\noindent Printed reading: {\ttfamily\footnotesize A\allowbreak{}}\par
\noindent Corrected reading: {\ttfamily\footnotesize V\allowbreak{}}\par
\noindent Source: {\ttfamily\footnotesize fr/sga7i/\{diplomatic\allowbreak{},current\}/expose\_IX\_\allowbreak{}body.tex;\allowbreak{} folio\allowbreak{} 403;\allowbreak{} formula\allowbreak{} 2567\allowbreak{}}\par

\par\medskip\noindent\hypertarget{app-E7I-IX-037.1}{}\textbf{E7I-IX-037.1}\quad printed folio 404.\par
\noindent Printed reading: \emph{(empty)}\par
\noindent Corrected reading: {\ttfamily\footnotesize \textbackslash{}underline\allowbreak{} }\par
\noindent Source: {\ttfamily\footnotesize port/decisions/decis\allowbreak{}ions\_en7I\_r2\_003.jso\allowbreak{}n\#C894;\allowbreak{} 7I:643\allowbreak{}}\par

\par\medskip\noindent\hypertarget{app-E7I-IX-038.1}{}\textbf{E7I-IX-038.1}\quad printed folio 405.\par
\noindent Printed reading: \emph{(empty)}\par
\noindent Corrected reading: {\ttfamily\footnotesize \_o\allowbreak{}}\par
\noindent Source: {\ttfamily\footnotesize fr/sga7i/\{diplomatic\allowbreak{},current\}/expose\_IX\_\allowbreak{}body.tex;\allowbreak{} folio\allowbreak{} 405;\allowbreak{} formula\allowbreak{} 2646\allowbreak{}}\par

\par\medskip\noindent\hypertarget{app-E7I-IX-038.2}{}\textbf{E7I-IX-038.2}\quad printed folio 405.\par
\noindent Printed reading: \emph{(empty)}\par
\noindent Corrected reading: {\ttfamily\footnotesize \_n\allowbreak{}}\par
\noindent Source: {\ttfamily\footnotesize fr/sga7i/\{diplomatic\allowbreak{},current\}/expose\_IX\_\allowbreak{}body.tex;\allowbreak{} folio\allowbreak{} 405;\allowbreak{} formula\allowbreak{} 2649\allowbreak{}}\par

\par\medskip\noindent\hypertarget{app-E7I-IX-038.3}{}\textbf{E7I-IX-038.3}\quad printed folio 405.\par
\noindent Printed reading: \emph{(empty)}\par
\noindent Corrected reading: {\ttfamily\footnotesize \_n\allowbreak{}}\par
\noindent Source: {\ttfamily\footnotesize fr/sga7i/\{diplomatic\allowbreak{},current\}/expose\_IX\_\allowbreak{}body.tex;\allowbreak{} folio\allowbreak{} 405;\allowbreak{} formula\allowbreak{} 2650\allowbreak{}}\par

\par\medskip\noindent\hypertarget{app-E7I-IX-039.1}{}\textbf{E7I-IX-039.1}\quad printed folio 410.\par
\noindent Printed reading: {\ttfamily\footnotesize IV\allowbreak{}}\par
\noindent Corrected reading: {\ttfamily\footnotesize III\allowbreak{}}\par
\noindent Source: {\ttfamily\footnotesize port/PORT\_EN7I\_LOG.j\allowbreak{}sonl\#C900;\allowbreak{} fixes/FIXLOG\_7I\_CURR\allowbreak{}ENT.jsonl\#C900;\allowbreak{} 7I:649\allowbreak{}}\par

\par\medskip\noindent\hypertarget{app-E7I-IX-040.1}{}\textbf{E7I-IX-040.1}\quad printed folio 411.\par
\noindent Printed reading: {\ttfamily\footnotesize 4\allowbreak{}}\par
\noindent Corrected reading: {\ttfamily\footnotesize 5\allowbreak{}}\par
\noindent Source: {\ttfamily\footnotesize port/PORT\_EN7I\_LOG.j\allowbreak{}sonl\#C901;\allowbreak{} fixes/FIXLOG\_7I\_CURR\allowbreak{}ENT.jsonl\#C901;\allowbreak{} 7I:840\allowbreak{}}\par

\par\medskip\noindent\hypertarget{app-E7I-IX-040.2}{}\textbf{E7I-IX-040.2}\quad printed folio 411.\par
\noindent Printed reading: {\ttfamily\footnotesize 3\allowbreak{}}\par
\noindent Corrected reading: {\ttfamily\footnotesize 4\allowbreak{}}\par
\noindent Source: {\ttfamily\footnotesize fr/sga7i/\{diplomatic\allowbreak{},current\}/expose\_IX\_\allowbreak{}body.tex;\allowbreak{} folio\allowbreak{} 411;\allowbreak{} formula\allowbreak{} 2800;\allowbreak{} port/PORT\_EN7I\_LOG.j\allowbreak{}sonl\#C902;\allowbreak{} fixes/FIXLOG\_7I\_CURR\allowbreak{}ENT.jsonl\#C902;\allowbreak{} 7I:841\allowbreak{}}\par

\par\medskip\noindent\hypertarget{app-E7I-IX-040.3}{}\textbf{E7I-IX-040.3}\quad printed folio 411.\par
\noindent Printed reading: {\ttfamily\footnotesize 1\allowbreak{}}\par
\noindent Corrected reading: {\ttfamily\footnotesize 2\allowbreak{}}\par
\noindent Source: {\ttfamily\footnotesize fr/sga7i/\{diplomatic\allowbreak{},current\}/expose\_IX\_\allowbreak{}body.tex;\allowbreak{} folio\allowbreak{} 411;\allowbreak{} formula\allowbreak{} 2800;\allowbreak{} port/PORT\_EN7I\_LOG.j\allowbreak{}sonl\#C902;\allowbreak{} fixes/FIXLOG\_7I\_CURR\allowbreak{}ENT.jsonl\#C902;\allowbreak{} 7I:841\allowbreak{}}\par

\par\medskip\noindent\hypertarget{app-E7I-IX-041.1}{}\textbf{E7I-IX-041.1}\quad printed folio 414.\par
\noindent Printed reading: {\ttfamily\footnotesize 5\allowbreak{}}\par
\noindent Corrected reading: {\ttfamily\footnotesize 3\allowbreak{}}\par
\noindent Source: {\ttfamily\footnotesize 7I:842;\allowbreak{} fixes/MISPRINTS\_7I.j\allowbreak{}son\#842\allowbreak{}}\par

\par\medskip\noindent\hypertarget{app-E7I-IX-042.1}{}\textbf{E7I-IX-042.1}\quad printed folio 415.\par
\noindent Printed reading: {\ttfamily\footnotesize 5\allowbreak{}}\par
\noindent Corrected reading: {\ttfamily\footnotesize 3\allowbreak{}}\par
\noindent Source: {\ttfamily\footnotesize 7I:844;\allowbreak{} fixes/MISPRINTS\_7I.j\allowbreak{}son\#844\allowbreak{}}\par

\par\medskip\noindent\hypertarget{app-E7I-IX-043.1}{}\textbf{E7I-IX-043.1}\quad printed folio 416.\par
\noindent Printed reading: \emph{(empty)}\par
\noindent Corrected reading: {\ttfamily\footnotesize '\allowbreak{}}\par
\noindent Source: {\ttfamily\footnotesize fr/sga7i/\{diplomatic\allowbreak{},current\}/expose\_IX\_\allowbreak{}body.tex;\allowbreak{} folio\allowbreak{} 416;\allowbreak{} formula\allowbreak{} 2930\allowbreak{}}\par

\par\medskip\noindent\hypertarget{app-E7I-IX-043.2}{}\textbf{E7I-IX-043.2}\quad printed folio 416.\par
\noindent Printed reading: {\ttfamily\footnotesize '\allowbreak{}}\par
\noindent Corrected reading: \emph{(empty)}\par
\noindent Source: {\ttfamily\footnotesize fr/sga7i/\{diplomatic\allowbreak{},current\}/expose\_IX\_\allowbreak{}body.tex;\allowbreak{} folio\allowbreak{} 416;\allowbreak{} formula\allowbreak{} 2930\allowbreak{}}\par

\par\medskip\noindent\hypertarget{app-E7I-IX-044.1}{}\textbf{E7I-IX-044.1}\quad printed folio 416.\par
\noindent Printed reading: {\ttfamily\footnotesize '\allowbreak{}}\par
\noindent Corrected reading: \emph{(empty)}\par
\noindent Source: {\ttfamily\footnotesize fr/sga7i/\{diplomatic\allowbreak{},current\}/expose\_IX\_\allowbreak{}body.tex;\allowbreak{} folio\allowbreak{} 416;\allowbreak{} formula\allowbreak{} 2954\allowbreak{}}\par

\par\medskip\noindent\hypertarget{app-E7I-IX-044.2}{}\textbf{E7I-IX-044.2}\quad printed folio 417.\par
\noindent Printed reading: \emph{(empty)}\par
\noindent Corrected reading: {\ttfamily\footnotesize \textasciicircum{}\{\textbackslash{}natural\}\allowbreak{}}\par
\noindent Source: {\ttfamily\footnotesize ED-EN-FORMULA-G065;\allowbreak{} fr/sga7i/\{diplomatic\allowbreak{},current\}/expose\_IX\_\allowbreak{}body.tex;\allowbreak{} folio\allowbreak{} 417;\allowbreak{} port/decisions/decis\allowbreak{}ions\_en7I\_r2\_003.jso\allowbreak{}n\#C904;\allowbreak{} 7I:653\allowbreak{}}\par

\par\medskip\noindent\hypertarget{app-E7I-IX-045.1}{}\textbf{E7I-IX-045.1}\quad printed folio 418.\par
\noindent Printed reading: \emph{(empty)}\par
\noindent Corrected reading: {\ttfamily\footnotesize \textbackslash{}underline\allowbreak{}}\par
\noindent Source: {\ttfamily\footnotesize fr/sga7i/\{diplomatic\allowbreak{},current\}/expose\_IX\_\allowbreak{}body.tex;\allowbreak{} folio\allowbreak{} 418;\allowbreak{} formula\allowbreak{} 2982;\allowbreak{} port/decisions/decis\allowbreak{}ions\_en7I\_r2\_003.jso\allowbreak{}n\#C906;\allowbreak{} fixes/FIXLOG\_7I\_CURR\allowbreak{}ENT.jsonl\#C906;\allowbreak{} 7I:655\allowbreak{}}\par

\par\medskip\noindent\hypertarget{app-E7I-IX-046.1}{}\textbf{E7I-IX-046.1}\quad printed folio 419.\par
\noindent Printed reading: {\ttfamily\footnotesize u\allowbreak{}}\par
\noindent Corrected reading: {\ttfamily\footnotesize x\allowbreak{}}\par
\noindent Source: {\ttfamily\footnotesize fr/sga7i/\{diplomatic\allowbreak{},current\}/expose\_IX\_\allowbreak{}body.tex;\allowbreak{} folio\allowbreak{} 419;\allowbreak{} formula\allowbreak{} 3014\allowbreak{}}\par

\par\medskip\noindent\hypertarget{app-E7I-IX-047.1}{}\textbf{E7I-IX-047.1}\quad printed folio 419.\par
\noindent Printed reading: {\ttfamily\footnotesize g\allowbreak{} \textbackslash{}longmapsto\allowbreak{} }\par
\noindent Corrected reading: \emph{(empty)}\par
\noindent Source: {\ttfamily\footnotesize fr/sga7i/\{diplomatic\allowbreak{},current\}/expose\_IX\_\allowbreak{}body.tex;\allowbreak{} folio\allowbreak{} 419;\allowbreak{} formula\allowbreak{} 3019\allowbreak{}}\par

\par\medskip\noindent\hypertarget{app-E7I-IX-048.1}{}\textbf{E7I-IX-048.1}\quad printed folio 420.\par
\noindent Printed reading: {\ttfamily\footnotesize )\allowbreak{}}\par
\noindent Corrected reading: \emph{(empty)}\par
\noindent Source: {\ttfamily\footnotesize fr/sga7i/\{diplomatic\allowbreak{},current\}/expose\_IX\_\allowbreak{}body.tex;\allowbreak{} folio\allowbreak{} 420;\allowbreak{} formula\allowbreak{} 3053\allowbreak{}}\par

\par\medskip\noindent\hypertarget{app-E7I-IX-048.2}{}\textbf{E7I-IX-048.2}\quad printed folio 420.\par
\noindent Printed reading: \emph{(empty)}\par
\noindent Corrected reading: {\ttfamily\footnotesize '\allowbreak{}}\par
\noindent Source: {\ttfamily\footnotesize fr/sga7i/\{diplomatic\allowbreak{},current\}/expose\_IX\_\allowbreak{}body.tex;\allowbreak{} folio\allowbreak{} 420;\allowbreak{} formula\allowbreak{} 3053\allowbreak{}}\par

\par\medskip\noindent\hypertarget{app-E7I-IX-048.3}{}\textbf{E7I-IX-048.3}\quad printed folio 420.\par
\noindent Printed reading: \emph{(empty)}\par
\noindent Corrected reading: {\ttfamily\footnotesize \textbackslash{}check\allowbreak{}}\par
\noindent Source: {\ttfamily\footnotesize fr/sga7i/\{diplomatic\allowbreak{},current\}/expose\_IX\_\allowbreak{}body.tex;\allowbreak{} folio\allowbreak{} 420;\allowbreak{} formula\allowbreak{} 3053;\allowbreak{} port/decisions/decis\allowbreak{}ions\_en7I\_r2\_003.jso\allowbreak{}n\#C908;\allowbreak{} fixes/FIXLOG\_7I\_CURR\allowbreak{}ENT.jsonl\#C908;\allowbreak{} 7I:660\allowbreak{}}\par

\par\medskip\noindent\hypertarget{app-E7I-IX-049.1}{}\textbf{E7I-IX-049.1}\quad printed folio 421.\par
\noindent Printed reading: \emph{(empty)}\par
\noindent Corrected reading: {\ttfamily\footnotesize \textbackslash{}bigl(\allowbreak{}}\par
\noindent Source: {\ttfamily\footnotesize fr/sga7i/\{diplomatic\allowbreak{},current\}/expose\_IX\_\allowbreak{}body.tex;\allowbreak{} folio\allowbreak{} 421;\allowbreak{} formula\allowbreak{} 3067\allowbreak{}}\par

\par\medskip\noindent\hypertarget{app-E7I-IX-049.2}{}\textbf{E7I-IX-049.2}\quad printed folio 421.\par
\noindent Printed reading: {\ttfamily\footnotesize I\allowbreak{}}\par
\noindent Corrected reading: {\ttfamily\footnotesize \textbackslash{}pi\_o\allowbreak{}}\par
\noindent Source: {\ttfamily\footnotesize fr/sga7i/\{diplomatic\allowbreak{},current\}/expose\_IX\_\allowbreak{}body.tex;\allowbreak{} folio\allowbreak{} 421;\allowbreak{} formula\allowbreak{} 3067\allowbreak{}}\par

\par\medskip\noindent\hypertarget{app-E7I-IX-050.1}{}\textbf{E7I-IX-050.1}\quad printed folio 421.\par
\noindent Printed reading: {\ttfamily\footnotesize P\allowbreak{}}\par
\noindent Corrected reading: {\ttfamily\footnotesize Q\allowbreak{}}\par
\noindent Source: {\ttfamily\footnotesize fr/sga7i/\{diplomatic\allowbreak{},current\}/expose\_IX\_\allowbreak{}body.tex;\allowbreak{} folio\allowbreak{} 421;\allowbreak{} formula\allowbreak{} 3076\allowbreak{}}\par

\par\medskip\noindent\hypertarget{app-E7I-IX-050.2}{}\textbf{E7I-IX-050.2}\quad printed folio 421.\par
\noindent Printed reading: {\ttfamily\footnotesize Q\allowbreak{}}\par
\noindent Corrected reading: {\ttfamily\footnotesize P\allowbreak{}}\par
\noindent Source: {\ttfamily\footnotesize fr/sga7i/\{diplomatic\allowbreak{},current\}/expose\_IX\_\allowbreak{}body.tex;\allowbreak{} folio\allowbreak{} 421;\allowbreak{} formula\allowbreak{} 3076\allowbreak{}}\par

\par\medskip\noindent\hypertarget{app-E7I-IX-051.1}{}\textbf{E7I-IX-051.1}\quad printed folio 423.\par
\noindent Printed reading: {\ttfamily\footnotesize F\allowbreak{}}\par
\noindent Corrected reading: {\ttfamily\footnotesize E\allowbreak{}}\par
\noindent Source: {\ttfamily\footnotesize fr/sga7i/\{diplomatic\allowbreak{},current\}/expose\_IX\_\allowbreak{}body.tex;\allowbreak{} folio\allowbreak{} 423\allowbreak{}}\par

\par\medskip\noindent\hypertarget{app-E7I-IX-051.2}{}\textbf{E7I-IX-051.2}\quad printed folio 423.\par
\noindent Printed reading: {\ttfamily\footnotesize F\allowbreak{}}\par
\noindent Corrected reading: {\ttfamily\footnotesize E\allowbreak{}}\par
\noindent Source: {\ttfamily\footnotesize fr/sga7i/\{diplomatic\allowbreak{},current\}/expose\_IX\_\allowbreak{}body.tex;\allowbreak{} folio\allowbreak{} 423;\allowbreak{} formula\allowbreak{} 3156\allowbreak{}}\par

\par\medskip\noindent\hypertarget{app-E7I-IX-052.1}{}\textbf{E7I-IX-052.1}\quad printed folio 425.\par
\noindent Printed reading: {\ttfamily\footnotesize d\allowbreak{}}\par
\noindent Corrected reading: {\ttfamily\footnotesize c\allowbreak{}}\par
\noindent Source: {\ttfamily\footnotesize 7I:667;\allowbreak{} 7I:863;\allowbreak{} fixes/MISPRINTS\_7I.j\allowbreak{}son\#667;\allowbreak{} fixes/MISPRINTS\_7I.j\allowbreak{}son\#863\allowbreak{}}\par

\par\medskip\noindent\hypertarget{app-E7I-IX-053.1}{}\textbf{E7I-IX-053.1}\quad printed folio 425.\par
\noindent Printed reading: {\ttfamily\footnotesize R\allowbreak{}}\par
\noindent Corrected reading: {\ttfamily\footnotesize Q\allowbreak{}}\par
\noindent Source: {\ttfamily\footnotesize ED-EN-FORMULA-G067;\allowbreak{} fr/sga7i/\{diplomatic\allowbreak{},current\}/expose\_IX\_\allowbreak{}body.tex;\allowbreak{} folio\allowbreak{} 425\allowbreak{}}\par

\par\medskip\noindent\hypertarget{app-E7I-IX-053.2}{}\textbf{E7I-IX-053.2}\quad printed folio 425.\par
\noindent Printed reading: {\ttfamily\footnotesize Q\allowbreak{}}\par
\noindent Corrected reading: {\ttfamily\footnotesize P\allowbreak{}}\par
\noindent Source: {\ttfamily\footnotesize ED-EN-FORMULA-G067;\allowbreak{} fr/sga7i/\{diplomatic\allowbreak{},current\}/expose\_IX\_\allowbreak{}body.tex;\allowbreak{} folio\allowbreak{} 425\allowbreak{}}\par

\par\medskip\noindent\hypertarget{app-E7I-IX-054.1}{}\textbf{E7I-IX-054.1}\quad printed folio 425.\par
\noindent Printed reading: \emph{(empty)}\par
\noindent Corrected reading: {\ttfamily\footnotesize '\allowbreak{}}\par
\noindent Source: {\ttfamily\footnotesize fr/sga7i/\{diplomatic\allowbreak{},current\}/expose\_IX\_\allowbreak{}body.tex;\allowbreak{} folio\allowbreak{} 425\allowbreak{}}\par

\par\medskip\noindent\hypertarget{app-E7I-IX-055.1}{}\textbf{E7I-IX-055.1}\quad printed folio 426.\par
\noindent Printed reading: \emph{(empty)}\par
\noindent Corrected reading: {\ttfamily\footnotesize '\allowbreak{}}\par
\noindent Source: {\ttfamily\footnotesize fr/sga7i/\{diplomatic\allowbreak{},current\}/expose\_IX\_\allowbreak{}body.tex;\allowbreak{} folio\allowbreak{} 426;\allowbreak{} formula\allowbreak{} 3238\allowbreak{}}\par

\par\medskip\noindent\hypertarget{app-E7I-IX-056.1}{}\textbf{E7I-IX-056.1}\quad printed folio 428.\par
\noindent Printed reading: \emph{(empty)}\par
\noindent Corrected reading: {\ttfamily\footnotesize )\allowbreak{}}\par
\noindent Source: {\ttfamily\footnotesize fr/sga7i/\{diplomatic\allowbreak{},current\}/expose\_IX\_\allowbreak{}body.tex;\allowbreak{} folio\allowbreak{} 428;\allowbreak{} formula\allowbreak{} 3302;\allowbreak{} port/decisions/decis\allowbreak{}ions\_en7I\_r2\_003.jso\allowbreak{}n\#C914;\allowbreak{} fixes/FIXLOG\_7I\_CURR\allowbreak{}ENT.jsonl\#C914;\allowbreak{} 7I:674\allowbreak{}}\par

\par\medskip\noindent\hypertarget{app-E7I-IX-057.1}{}\textbf{E7I-IX-057.1}\quad printed folio 429.\par
\noindent Printed reading: {\ttfamily\footnotesize \textasciicircum{}\allowbreak{}}\par
\noindent Corrected reading: {\ttfamily\footnotesize \_\allowbreak{}}\par
\noindent Source: {\ttfamily\footnotesize port/decisions/decis\allowbreak{}ions\_en7I\_r2\_003.jso\allowbreak{}n\#C917;\allowbreak{} fixes/FIXLOG\_7I\_CURR\allowbreak{}ENT.jsonl\#C917;\allowbreak{} 7I:677\allowbreak{}}\par

\par\medskip\noindent\hypertarget{app-E7I-IX-057.2}{}\textbf{E7I-IX-057.2}\quad printed folio 428.\par
\noindent Printed reading: \emph{(empty)}\par
\noindent Corrected reading: {\ttfamily\footnotesize \textbackslash{}check\allowbreak{}}\par
\noindent Source: {\ttfamily\footnotesize fr/sga7i/\{diplomatic\allowbreak{},current\}/expose\_IX\_\allowbreak{}body.tex;\allowbreak{} folio\allowbreak{} 428;\allowbreak{} formula\allowbreak{} 3303;\allowbreak{} port/decisions/decis\allowbreak{}ions\_en7I\_r2\_003.jso\allowbreak{}n\#C915;\allowbreak{} fixes/FIXLOG\_7I\_CURR\allowbreak{}ENT.jsonl\#C915;\allowbreak{} 7I:675\allowbreak{}}\par

\par\medskip\noindent\hypertarget{app-E7I-IX-058.1}{}\textbf{E7I-IX-058.1}\quad printed folio 429.\par
\noindent Printed reading: {\ttfamily\footnotesize \textasciicircum{}\allowbreak{}}\par
\noindent Corrected reading: {\ttfamily\footnotesize \_\allowbreak{}}\par
\noindent Source: {\ttfamily\footnotesize ED-EN-FORMULA-G069;\allowbreak{} fr/sga7i/\{diplomatic\allowbreak{},current\}/expose\_IX\_\allowbreak{}body.tex;\allowbreak{} folio\allowbreak{} 429\allowbreak{}}\par

\par\medskip\noindent\hypertarget{app-E7I-IX-058.2}{}\textbf{E7I-IX-058.2}\quad printed folio 429.\par
\noindent Printed reading: {\ttfamily\footnotesize oS\allowbreak{}}\par
\noindent Corrected reading: {\ttfamily\footnotesize oU\allowbreak{}}\par
\noindent Source: {\ttfamily\footnotesize ED-EN-FORMULA-G069;\allowbreak{} fr/sga7i/\{diplomatic\allowbreak{},current\}/expose\_IX\_\allowbreak{}body.tex;\allowbreak{} folio\allowbreak{} 429;\allowbreak{} port/decisions/decis\allowbreak{}ions\_en7I\_r2\_003.jso\allowbreak{}n\#C918;\allowbreak{} fixes/FIXLOG\_7I\_CURR\allowbreak{}ENT.jsonl\#C918;\allowbreak{} 7I:678\allowbreak{}}\par

\par\medskip\noindent\hypertarget{app-E7I-IX-059.1}{}\textbf{E7I-IX-059.1}\quad printed folio 429.\par
\noindent Printed reading: \emph{(empty)}\par
\noindent Corrected reading: {\ttfamily\footnotesize )\allowbreak{}}\par
\noindent Source: {\ttfamily\footnotesize fr/sga7i/\{diplomatic\allowbreak{},current\}/expose\_IX\_\allowbreak{}body.tex;\allowbreak{} folio\allowbreak{} 429;\allowbreak{} formula\allowbreak{} 3332;\allowbreak{} port/decisions/decis\allowbreak{}ions\_en7I\_r2\_003.jso\allowbreak{}n\#C919;\allowbreak{} fixes/FIXLOG\_7I\_CURR\allowbreak{}ENT.jsonl\#C919;\allowbreak{} 7I:679\allowbreak{}}\par

\par\medskip\noindent\hypertarget{app-E7I-IX-060.1}{}\textbf{E7I-IX-060.1}\quad printed folio 431.\par
\noindent Printed reading: \emph{(empty)}\par
\noindent Corrected reading: {\ttfamily\footnotesize Q,\allowbreak{}}\par
\noindent Source: {\ttfamily\footnotesize fr/sga7i/\{diplomatic\allowbreak{},current\}/expose\_IX\_\allowbreak{}body.tex;\allowbreak{} folio\allowbreak{} 431\allowbreak{}}\par

\par\medskip\noindent\hypertarget{app-E7I-IX-060.2}{}\textbf{E7I-IX-060.2}\quad printed folio 431.\par
\noindent Printed reading: {\ttfamily\footnotesize ,Q\allowbreak{}}\par
\noindent Corrected reading: \emph{(empty)}\par
\noindent Source: {\ttfamily\footnotesize fr/sga7i/\{diplomatic\allowbreak{},current\}/expose\_IX\_\allowbreak{}body.tex;\allowbreak{} folio\allowbreak{} 431\allowbreak{}}\par

\par\medskip\noindent\hypertarget{app-E7I-IX-061.1}{}\textbf{E7I-IX-061.1}\quad printed folio 431.\par
\noindent Printed reading: \emph{(empty)}\par
\noindent Corrected reading: {\ttfamily\footnotesize Q,\allowbreak{}}\par
\noindent Source: {\ttfamily\footnotesize fr/sga7i/\{diplomatic\allowbreak{},current\}/expose\_IX\_\allowbreak{}body.tex;\allowbreak{} folio\allowbreak{} 431;\allowbreak{} formula\allowbreak{} 3385\allowbreak{}}\par

\par\medskip\noindent\hypertarget{app-E7I-IX-061.2}{}\textbf{E7I-IX-061.2}\quad printed folio 431.\par
\noindent Printed reading: {\ttfamily\footnotesize ,Q\allowbreak{}}\par
\noindent Corrected reading: \emph{(empty)}\par
\noindent Source: {\ttfamily\footnotesize fr/sga7i/\{diplomatic\allowbreak{},current\}/expose\_IX\_\allowbreak{}body.tex;\allowbreak{} folio\allowbreak{} 431;\allowbreak{} formula\allowbreak{} 3385\allowbreak{}}\par

\par\medskip\noindent\hypertarget{app-E7I-IX-061.3}{}\textbf{E7I-IX-061.3}\quad printed folio 431.\par
\noindent Printed reading: {\ttfamily\footnotesize P\allowbreak{}}\par
\noindent Corrected reading: {\ttfamily\footnotesize Q\allowbreak{}}\par
\noindent Source: {\ttfamily\footnotesize fr/sga7i/\{diplomatic\allowbreak{},current\}/expose\_IX\_\allowbreak{}body.tex;\allowbreak{} folio\allowbreak{} 431;\allowbreak{} formula\allowbreak{} 3385\allowbreak{}}\par

\par\medskip\noindent\hypertarget{app-E7I-IX-061.4}{}\textbf{E7I-IX-061.4}\quad printed folio 431.\par
\noindent Printed reading: {\ttfamily\footnotesize \textbackslash{}\allowbreak{} }\par
\noindent Corrected reading: \emph{(empty)}\par
\noindent Source: {\ttfamily\footnotesize fr/sga7i/\{diplomatic\allowbreak{},current\}/expose\_IX\_\allowbreak{}body.tex;\allowbreak{} folio\allowbreak{} 431;\allowbreak{} formula\allowbreak{} 3385\allowbreak{}}\par

\par\medskip\noindent\hypertarget{app-E7I-IX-061.5}{}\textbf{E7I-IX-061.5}\quad printed folio 431.\par
\noindent Printed reading: {\ttfamily\footnotesize Q\allowbreak{}}\par
\noindent Corrected reading: {\ttfamily\footnotesize P\allowbreak{}}\par
\noindent Source: {\ttfamily\footnotesize fr/sga7i/\{diplomatic\allowbreak{},current\}/expose\_IX\_\allowbreak{}body.tex;\allowbreak{} folio\allowbreak{} 431;\allowbreak{} formula\allowbreak{} 3385\allowbreak{}}\par

\par\medskip\noindent\hypertarget{app-E7I-IX-062.1}{}\textbf{E7I-IX-062.1}\quad printed folio 432.\par
\noindent Printed reading: {\ttfamily\footnotesize 6\allowbreak{}}\par
\noindent Corrected reading: {\ttfamily\footnotesize 7\allowbreak{}}\par
\noindent Source: {\ttfamily\footnotesize 7I:883;\allowbreak{} fixes/MISPRINTS\_7I.j\allowbreak{}son\#883\allowbreak{}}\par

\par\medskip\noindent\hypertarget{app-E7I-IX-063.1}{}\textbf{E7I-IX-063.1}\quad printed folio 434.\par
\noindent Printed reading: {\ttfamily\footnotesize 6\allowbreak{}}\par
\noindent Corrected reading: {\ttfamily\footnotesize 7\allowbreak{}}\par
\noindent Source: {\ttfamily\footnotesize 7I:885;\allowbreak{} fixes/MISPRINTS\_7I.j\allowbreak{}son\#885\allowbreak{}}\par

\par\medskip\noindent\hypertarget{app-E7I-IX-063.2}{}\textbf{E7I-IX-063.2}\quad printed folio 434.\par
\noindent Printed reading: {\ttfamily\footnotesize \_\textbackslash{}eta\allowbreak{}}\par
\noindent Corrected reading: \emph{(empty)}\par
\noindent Source: {\ttfamily\footnotesize fr/sga7i/\{diplomatic\allowbreak{},current\}/expose\_IX\_\allowbreak{}body.tex;\allowbreak{} folio\allowbreak{} 434;\allowbreak{} formula\allowbreak{} 3456\allowbreak{}}\par

\par\medskip\noindent\hypertarget{app-E7I-IX-064.1}{}\textbf{E7I-IX-064.1}\quad printed folio 435.\par
\noindent Printed reading: {\ttfamily\footnotesize n\allowbreak{}}\par
\noindent Corrected reading: {\ttfamily\footnotesize S\allowbreak{}}\par
\noindent Source: {\ttfamily\footnotesize fr/sga7i/\{diplomatic\allowbreak{},current\}/expose\_IX\_\allowbreak{}body.tex;\allowbreak{} folio\allowbreak{} 435;\allowbreak{} formula\allowbreak{} 3463\allowbreak{}}\par

\par\medskip\noindent\hypertarget{app-E7I-IX-065.1}{}\textbf{E7I-IX-065.1}\quad printed folio 435.\par
\noindent Printed reading: {\ttfamily\footnotesize 6\allowbreak{}}\par
\noindent Corrected reading: {\ttfamily\footnotesize 5\allowbreak{}}\par
\noindent Source: {\ttfamily\footnotesize port/PORT\_EN7I\_LOG.j\allowbreak{}sonl\#C921;\allowbreak{} fixes/FIXLOG\_7I\_CURR\allowbreak{}ENT.jsonl\#C921;\allowbreak{} 7I:890\allowbreak{}}\par

\par\medskip\noindent\hypertarget{app-E7I-IX-065.2}{}\textbf{E7I-IX-065.2}\quad printed folio 436.\par
\noindent Printed reading: {\ttfamily\footnotesize 6\allowbreak{}}\par
\noindent Corrected reading: {\ttfamily\footnotesize 5\allowbreak{}}\par
\noindent Source: {\ttfamily\footnotesize WITNESS-F7I-IX-116.1\allowbreak{};\allowbreak{} fr/sga7i/diplomatic/\allowbreak{}expose\_IX\_body.tex;\allowbreak{} fr/sga7i/current/exp\allowbreak{}ose\_IX\_body.tex\allowbreak{}}\par

\par\medskip\noindent\hypertarget{app-E7I-IX-066.1}{}\textbf{E7I-IX-066.1}\quad printed folio 437.\par
\noindent Printed reading: {\ttfamily\footnotesize F\allowbreak{}}\par
\noindent Corrected reading: {\ttfamily\footnotesize E\allowbreak{}}\par
\noindent Source: {\ttfamily\footnotesize fr/sga7i/\{diplomatic\allowbreak{},current\}/expose\_IX\_\allowbreak{}body.tex;\allowbreak{} folio\allowbreak{} 437;\allowbreak{} formula\allowbreak{} 3537\allowbreak{}}\par

\par\medskip\noindent\hypertarget{app-E7I-IX-066.2}{}\textbf{E7I-IX-066.2}\quad printed folio 437.\par
\noindent Printed reading: {\ttfamily\footnotesize 6\allowbreak{}}\par
\noindent Corrected reading: {\ttfamily\footnotesize 5\allowbreak{}}\par
\noindent Source: {\ttfamily\footnotesize ADOPT-F7I-IX-117.2;\allowbreak{} fr/sga7i/diplomatic/\allowbreak{}expose\_IX\_body.tex;\allowbreak{} fr/sga7i/current/exp\allowbreak{}ose\_IX\_body.tex\allowbreak{}}\par

\par\medskip\noindent\hypertarget{app-E7I-IX-066.3}{}\textbf{E7I-IX-066.3}\quad printed folio 437.\par
\noindent Printed reading: {\ttfamily\footnotesize S\allowbreak{}}\par
\noindent Corrected reading: {\ttfamily\footnotesize \textbackslash{}eta\allowbreak{}}\par
\noindent Source: {\ttfamily\footnotesize fr/sga7i/\{diplomatic\allowbreak{},current\}/expose\_IX\_\allowbreak{}body.tex;\allowbreak{} folio\allowbreak{} 437;\allowbreak{} formula\allowbreak{} 3542\allowbreak{}}\par

\par\medskip\noindent\hypertarget{app-E7I-IX-067.1}{}\textbf{E7I-IX-067.1}\quad printed folio 439.\par
\noindent Printed reading: \emph{(empty)}\par
\noindent Corrected reading: {\ttfamily\footnotesize '\allowbreak{}}\par
\noindent Source: {\ttfamily\footnotesize ED-EN-FORMULA-G071;\allowbreak{} fr/sga7i/\{diplomatic\allowbreak{},current\}/expose\_IX\_\allowbreak{}body.tex;\allowbreak{} folio\allowbreak{} 439;\allowbreak{} port/decisions/decis\allowbreak{}ions\_en7I\_r2\_003.jso\allowbreak{}n\#C924;\allowbreak{} 7I:688\allowbreak{}}\par

\par\medskip\noindent\hypertarget{app-E7I-IX-068.1}{}\textbf{E7I-IX-068.1}\quad printed folio 441.\par
\noindent Printed reading: {\ttfamily\footnotesize \textbackslash{}varphi\allowbreak{}}\par
\noindent Corrected reading: {\ttfamily\footnotesize u\_\textbackslash{}ell\allowbreak{}}\par
\noindent Source: {\ttfamily\footnotesize port/decisions/decis\allowbreak{}ions\_en7I\_r2\_003.jso\allowbreak{}n\#C927;\allowbreak{} 7I:691\allowbreak{}}\par

\par\medskip\noindent\hypertarget{app-E7I-IX-068.2}{}\textbf{E7I-IX-068.2}\quad printed folio 441.\par
\noindent Printed reading: {\ttfamily\footnotesize t\allowbreak{}}\par
\noindent Corrected reading: {\ttfamily\footnotesize \{A\_K\}\allowbreak{}}\par
\noindent Source: {\ttfamily\footnotesize port/decisions/decis\allowbreak{}ions\_en7I\_r2\_003.jso\allowbreak{}n\#C927;\allowbreak{} 7I:691\allowbreak{}}\par

\par\medskip\noindent\hypertarget{app-E7I-IX-069.1}{}\textbf{E7I-IX-069.1}\quad printed folio 441.\par
\noindent Printed reading: \emph{(empty)}\par
\noindent Corrected reading: {\ttfamily\footnotesize 1.\allowbreak{}}\par
\noindent Source: {\ttfamily\footnotesize 7I:893;\allowbreak{} fixes/MISPRINTS\_7I.j\allowbreak{}son\#893\allowbreak{}}\par

\par\medskip\noindent\hypertarget{app-E7I-IX-069.2}{}\textbf{E7I-IX-069.2}\quad printed folio 441.\par
\noindent Printed reading: {\ttfamily\footnotesize .1\allowbreak{}}\par
\noindent Corrected reading: \emph{(empty)}\par
\noindent Source: {\ttfamily\footnotesize 7I:893;\allowbreak{} fixes/MISPRINTS\_7I.j\allowbreak{}son\#893\allowbreak{}}\par

\par\medskip\noindent\hypertarget{app-E7I-IX-070.1}{}\textbf{E7I-IX-070.1}\quad printed folio 442.\par
\noindent Printed reading: \emph{(empty)}\par
\noindent Corrected reading: {\ttfamily\footnotesize '\allowbreak{}}\par
\noindent Source: {\ttfamily\footnotesize fr/sga7i/\{diplomatic\allowbreak{},current\}/expose\_IX\_\allowbreak{}body.tex;\allowbreak{} folio\allowbreak{} 442;\allowbreak{} formula\allowbreak{} 3644;\allowbreak{} port/decisions/decis\allowbreak{}ions\_en7I\_r2\_004.jso\allowbreak{}n\#C928;\allowbreak{} fixes/FIXLOG\_7I\_CURR\allowbreak{}ENT.jsonl\#C928;\allowbreak{} 7I:692\allowbreak{}}\par

\par\medskip\noindent\hypertarget{app-E7I-IX-070.2}{}\textbf{E7I-IX-070.2}\quad printed folio 443.\par
\noindent Printed reading: {\ttfamily\footnotesize 2\allowbreak{}}\par
\noindent Corrected reading: {\ttfamily\footnotesize 3\allowbreak{}}\par
\noindent Source: {\ttfamily\footnotesize port/decisions/decis\allowbreak{}ions\_en7I\_r2\_004.jso\allowbreak{}n\#C930;\allowbreak{} fixes/FIXLOG\_7I\_CURR\allowbreak{}ENT.jsonl\#C930;\allowbreak{} 7I:694\allowbreak{}}\par

\par\medskip\noindent\hypertarget{app-E7I-IX-070.3}{}\textbf{E7I-IX-070.3}\quad printed folio 443.\par
\noindent Printed reading: {\ttfamily\footnotesize 3\allowbreak{}}\par
\noindent Corrected reading: {\ttfamily\footnotesize 4\allowbreak{}}\par
\noindent Source: {\ttfamily\footnotesize port/decisions/decis\allowbreak{}ions\_en7I\_r2\_004.jso\allowbreak{}n\#C930;\allowbreak{} fixes/FIXLOG\_7I\_CURR\allowbreak{}ENT.jsonl\#C930;\allowbreak{} 7I:694\allowbreak{}}\par

\par\medskip\noindent\hypertarget{app-E7I-IX-071.1}{}\textbf{E7I-IX-071.1}\quad printed folio 445.\par
\noindent Printed reading: \emph{(empty)}\par
\noindent Corrected reading: {\ttfamily\footnotesize \textbackslash{}mathbb\allowbreak{}}\par
\noindent Source: {\ttfamily\footnotesize ED-EN-FONT-G001;\allowbreak{} F7I-IX-125.1;\allowbreak{} fr/sga7i/\{diplomatic\allowbreak{},current\}/expose\_IX\_\allowbreak{}body.tex;\allowbreak{} folio\allowbreak{} 445\allowbreak{}}\par

\par\medskip\noindent\hypertarget{app-E7I-IX-072.1}{}\textbf{E7I-IX-072.1}\quad printed folio 445.\par
\noindent Printed reading: {\ttfamily\footnotesize k\allowbreak{}}\par
\noindent Corrected reading: {\ttfamily\footnotesize K\allowbreak{}}\par
\noindent Source: {\ttfamily\footnotesize fr/sga7i/\{diplomatic\allowbreak{},current\}/expose\_IX\_\allowbreak{}body.tex;\allowbreak{} folio\allowbreak{} 445;\allowbreak{} formula\allowbreak{} 3697\allowbreak{}}\par

\par\medskip\noindent\hypertarget{app-E7I-IX-073.1}{}\textbf{E7I-IX-073.1}\quad printed folio 446.\par
\noindent Printed reading: {\ttfamily\footnotesize X\allowbreak{}}\par
\noindent Corrected reading: {\ttfamily\footnotesize K\allowbreak{}}\par
\noindent Source: {\ttfamily\footnotesize fr/sga7i/\{diplomatic\allowbreak{},current\}/expose\_IX\_\allowbreak{}body.tex;\allowbreak{} folio\allowbreak{} 446;\allowbreak{} formula\allowbreak{} 3708\allowbreak{}}\par

\par\medskip\noindent\hypertarget{app-E7I-IX-074.1}{}\textbf{E7I-IX-074.1}\quad printed folio 448.\par
\noindent Printed reading: {\ttfamily\footnotesize 1\allowbreak{}}\par
\noindent Corrected reading: {\ttfamily\footnotesize \textbackslash{}ell\allowbreak{}}\par
\noindent Source: {\ttfamily\footnotesize fr/sga7i/\{diplomatic\allowbreak{},current\}/expose\_IX\_\allowbreak{}body.tex;\allowbreak{} folio\allowbreak{} 448;\allowbreak{} formula\allowbreak{} 3752\allowbreak{}}\par

\par\medskip\noindent\hypertarget{app-E7I-IX-075.1}{}\textbf{E7I-IX-075.1}\quad printed folio 452.\par
\noindent Printed reading: {\ttfamily\footnotesize ,\allowbreak{}}\par
\noindent Corrected reading: {\ttfamily\footnotesize .\allowbreak{}}\par
\noindent Source: {\ttfamily\footnotesize ED-EN-NORM-G026;\allowbreak{} complete\allowbreak{} normalized-equal\allowbreak{} formula\allowbreak{} audit;\allowbreak{} fr/sga7i/\{diplomatic\allowbreak{},current\}/expose\_IX\_\allowbreak{}body.tex;\allowbreak{} folio\allowbreak{} 452\allowbreak{}}\par

\par\medskip\noindent\hypertarget{app-E7I-IX-076.1}{}\textbf{E7I-IX-076.1}\quad printed folio 452.\par
\noindent Printed reading: \emph{(empty)}\par
\noindent Corrected reading: {\ttfamily\footnotesize \textbackslash{}mathbb\allowbreak{}}\par
\noindent Source: {\ttfamily\footnotesize ED-EN-FONT-G002;\allowbreak{} F7I-IX-131.1;\allowbreak{} fr/sga7i/\{diplomatic\allowbreak{},current\}/expose\_IX\_\allowbreak{}body.tex;\allowbreak{} folio\allowbreak{} 452\allowbreak{}}\par

\par\medskip\noindent\hypertarget{app-E7I-IX-077.1}{}\textbf{E7I-IX-077.1}\quad printed folio 454.\par
\noindent Printed reading: {\ttfamily\footnotesize 4\allowbreak{}}\par
\noindent Corrected reading: {\ttfamily\footnotesize 3\allowbreak{}}\par
\noindent Source: {\ttfamily\footnotesize 7I:907;\allowbreak{} fixes/MISPRINTS\_7I.j\allowbreak{}son\#907\allowbreak{}}\par

\par\medskip\noindent\hypertarget{app-E7I-IX-078.1}{}\textbf{E7I-IX-078.1}\quad printed folio 455.\par
\noindent Printed reading: {\ttfamily\footnotesize ;\allowbreak{}}\par
\noindent Corrected reading: {\ttfamily\footnotesize :\allowbreak{}}\par
\noindent Source: {\ttfamily\footnotesize fr/sga7i/\{diplomatic\allowbreak{},current\}/expose\_IX\_\allowbreak{}body.tex;\allowbreak{} folio\allowbreak{} 455;\allowbreak{} formula\allowbreak{} 3893\allowbreak{}}\par

\par\medskip\noindent\hypertarget{app-E7I-IX-079.1}{}\textbf{E7I-IX-079.1}\quad printed folio 456.\par
\noindent Printed reading: {\ttfamily\footnotesize \textbackslash{}Phi(\allowbreak{}}\par
\noindent Corrected reading: {\ttfamily\footnotesize M\_\allowbreak{}}\par
\noindent Source: {\ttfamily\footnotesize fr/sga7i/\{diplomatic\allowbreak{},current\}/expose\_IX\_\allowbreak{}body.tex;\allowbreak{} folio\allowbreak{} 456;\allowbreak{} formula\allowbreak{} 3928\allowbreak{}}\par

\par\medskip\noindent\hypertarget{app-E7I-IX-079.2}{}\textbf{E7I-IX-079.2}\quad printed folio 456.\par
\noindent Printed reading: {\ttfamily\footnotesize )\allowbreak{}}\par
\noindent Corrected reading: {\ttfamily\footnotesize \textbackslash{}otimes\allowbreak{} D\_\textbackslash{}ell\allowbreak{}}\par
\noindent Source: {\ttfamily\footnotesize fr/sga7i/\{diplomatic\allowbreak{},current\}/expose\_IX\_\allowbreak{}body.tex;\allowbreak{} folio\allowbreak{} 456;\allowbreak{} formula\allowbreak{} 3928\allowbreak{}}\par

\par\medskip\noindent\hypertarget{app-E7I-IX-079.3}{}\textbf{E7I-IX-079.3}\quad printed folio 456.\par
\noindent Printed reading: {\ttfamily\footnotesize \textbackslash{}Phi\allowbreak{}}\par
\noindent Corrected reading: {\ttfamily\footnotesize \textbackslash{}check\allowbreak{} M\allowbreak{}}\par
\noindent Source: {\ttfamily\footnotesize fr/sga7i/\{diplomatic\allowbreak{},current\}/expose\_IX\_\allowbreak{}body.tex;\allowbreak{} folio\allowbreak{} 456;\allowbreak{} formula\allowbreak{} 3928\allowbreak{}}\par

\par\medskip\noindent\hypertarget{app-E7I-IX-079.4}{}\textbf{E7I-IX-079.4}\quad printed folio 456.\par
\noindent Printed reading: {\ttfamily\footnotesize (\allowbreak{}}\par
\noindent Corrected reading: {\ttfamily\footnotesize \_\allowbreak{}}\par
\noindent Source: {\ttfamily\footnotesize fr/sga7i/\{diplomatic\allowbreak{},current\}/expose\_IX\_\allowbreak{}body.tex;\allowbreak{} folio\allowbreak{} 456;\allowbreak{} formula\allowbreak{} 3928\allowbreak{}}\par

\par\medskip\noindent\hypertarget{app-E7I-IX-079.5}{}\textbf{E7I-IX-079.5}\quad printed folio 456.\par
\noindent Printed reading: {\ttfamily\footnotesize )\allowbreak{}}\par
\noindent Corrected reading: {\ttfamily\footnotesize \textbackslash{}otimes\allowbreak{} D\_\textbackslash{}ell\allowbreak{}}\par
\noindent Source: {\ttfamily\footnotesize fr/sga7i/\{diplomatic\allowbreak{},current\}/expose\_IX\_\allowbreak{}body.tex;\allowbreak{} folio\allowbreak{} 456;\allowbreak{} formula\allowbreak{} 3928\allowbreak{}}\par

\par\medskip\noindent\hypertarget{app-E7I-IX-080.1}{}\textbf{E7I-IX-080.1}\quad printed folio 457.\par
\noindent Printed reading: \emph{(empty)}\par
\noindent Corrected reading: {\ttfamily\footnotesize '\allowbreak{}}\par
\noindent Source: {\ttfamily\footnotesize fr/sga7i/\{diplomatic\allowbreak{},current\}/expose\_IX\_\allowbreak{}body.tex;\allowbreak{} folio\allowbreak{} 457;\allowbreak{} formula\allowbreak{} 3931\allowbreak{}}\par

\par\medskip\noindent\hypertarget{app-E7I-IX-081.1}{}\textbf{E7I-IX-081.1}\quad printed folio 460.\par
\noindent Printed reading: {\ttfamily\footnotesize 7\allowbreak{}}\par
\noindent Corrected reading: {\ttfamily\footnotesize 6\allowbreak{}}\par
\noindent Source: {\ttfamily\footnotesize 7I:709;\allowbreak{} 7I:916;\allowbreak{} fixes/MISPRINTS\_7I.j\allowbreak{}son\#709;\allowbreak{} fixes/MISPRINTS\_7I.j\allowbreak{}son\#916\allowbreak{}}\par

\par\medskip\noindent\hypertarget{app-E7I-IX-081.2}{}\textbf{E7I-IX-081.2}\quad printed folio 460.\par
\noindent Printed reading: \emph{(empty)}\par
\noindent Corrected reading: {\ttfamily\footnotesize '\allowbreak{}}\par
\noindent Source: {\ttfamily\footnotesize fr/sga7i/\{diplomatic\allowbreak{},current\}/expose\_IX\_\allowbreak{}body.tex;\allowbreak{} folio\allowbreak{} 460;\allowbreak{} formula\allowbreak{} 3995;\allowbreak{} port/PORT\_EN7I\_LOG.j\allowbreak{}sonl\#C938;\allowbreak{} fixes/FIXLOG\_7I\_CURR\allowbreak{}ENT.jsonl\#C938;\allowbreak{} 7I:711\allowbreak{}}\par

\par\medskip\noindent\hypertarget{app-E7I-IX-081.3}{}\textbf{E7I-IX-081.3}\quad printed folio 460.\par
\noindent Printed reading: \emph{(empty)}\par
\noindent Corrected reading: {\ttfamily\footnotesize Proposition\allowbreak{} }\par
\noindent Source: {\ttfamily\footnotesize ADOPT-F7I-IX-137.3;\allowbreak{} fr/sga7i/diplomatic/\allowbreak{}expose\_IX\_body.tex;\allowbreak{} fr/sga7i/current/exp\allowbreak{}ose\_IX\_body.tex\allowbreak{}}\par

\par\medskip\noindent\hypertarget{app-E7I-IX-081.4}{}\textbf{E7I-IX-081.4}\quad printed folio 460.\par
\noindent Printed reading: {\ttfamily\footnotesize 6\allowbreak{}}\par
\noindent Corrected reading: {\ttfamily\footnotesize 7\allowbreak{}}\par
\noindent Source: {\ttfamily\footnotesize ADOPT-F7I-IX-137.3;\allowbreak{} fr/sga7i/diplomatic/\allowbreak{}expose\_IX\_body.tex;\allowbreak{} fr/sga7i/current/exp\allowbreak{}ose\_IX\_body.tex\allowbreak{}}\par

\par\medskip\noindent\hypertarget{app-E7I-IX-081.5}{}\textbf{E7I-IX-081.5}\quad printed folio 460.\par
\noindent Printed reading: \emph{(empty)}\par
\noindent Corrected reading: {\ttfamily\footnotesize '\allowbreak{}}\par
\noindent Source: {\ttfamily\footnotesize fr/sga7i/\{diplomatic\allowbreak{},current\}/expose\_IX\_\allowbreak{}body.tex;\allowbreak{} folio\allowbreak{} 460;\allowbreak{} formula\allowbreak{} 4005\allowbreak{}}\par

\par\medskip\noindent\hypertarget{app-E7I-IX-082.1}{}\textbf{E7I-IX-082.1}\quad printed folio 461.\par
\noindent Printed reading: {\ttfamily\footnotesize 7\allowbreak{}}\par
\noindent Corrected reading: {\ttfamily\footnotesize 6\allowbreak{}}\par
\noindent Source: {\ttfamily\footnotesize 7I:712;\allowbreak{} 7I:919;\allowbreak{} fixes/MISPRINTS\_7I.j\allowbreak{}son\#712;\allowbreak{} fixes/MISPRINTS\_7I.j\allowbreak{}son\#919\allowbreak{}}\par

\par\medskip\noindent\hypertarget{app-E7I-IX-083.1}{}\textbf{E7I-IX-083.1}\quad printed folio 462.\par
\noindent Printed reading: {\ttfamily\footnotesize 8\allowbreak{}}\par
\noindent Corrected reading: {\ttfamily\footnotesize 7\allowbreak{}}\par
\noindent Source: {\ttfamily\footnotesize port/PORT\_EN7I\_LOG.j\allowbreak{}sonl\#C939;\allowbreak{} fixes/FIXLOG\_7I\_CURR\allowbreak{}ENT.jsonl\#C939;\allowbreak{} 7I:922\allowbreak{}}\par

\par\medskip\noindent\hypertarget{app-E7I-IX-084.1}{}\textbf{E7I-IX-084.1}\quad printed folio 462.\par
\noindent Printed reading: {\ttfamily\footnotesize ,\allowbreak{}}\par
\noindent Corrected reading: {\ttfamily\footnotesize /\allowbreak{}}\par
\noindent Source: {\ttfamily\footnotesize fr/sga7i/\{diplomatic\allowbreak{},current\}/expose\_IX\_\allowbreak{}body.tex;\allowbreak{} folio\allowbreak{} 462;\allowbreak{} formula\allowbreak{} 4065\allowbreak{}}\par

\par\medskip\noindent\hypertarget{app-E7I-IX-085.1}{}\textbf{E7I-IX-085.1}\quad printed folio 463.\par
\noindent Printed reading: {\ttfamily\footnotesize \textbackslash{}Psi\allowbreak{}}\par
\noindent Corrected reading: {\ttfamily\footnotesize \textbackslash{}Phi\allowbreak{}}\par
\noindent Source: {\ttfamily\footnotesize fr/sga7i/\{diplomatic\allowbreak{},current\}/expose\_IX\_\allowbreak{}body.tex;\allowbreak{} folio\allowbreak{} 463;\allowbreak{} formula\allowbreak{} 4079\allowbreak{}}\par

\par\medskip\noindent\hypertarget{app-E7I-IX-086.1}{}\textbf{E7I-IX-086.1}\quad printed folio 464.\par
\noindent Printed reading: {\ttfamily\footnotesize  (cf.\allowbreak{} I\textbackslash{}hspace\{3em\})\allowbreak{}}\par
\noindent Corrected reading: \emph{(empty)}\par
\noindent Source: {\ttfamily\footnotesize ADOPT-F7I-IX-142.1;\allowbreak{} fr/sga7i/diplomatic/\allowbreak{}expose\_IX\_body.tex;\allowbreak{} fr/sga7i/current/exp\allowbreak{}ose\_IX\_body.tex\allowbreak{}}\par

\par\medskip\noindent\hypertarget{app-E7I-IX-087.1}{}\textbf{E7I-IX-087.1}\quad printed folio 464.\par
\noindent Printed reading: {\ttfamily\footnotesize 2\allowbreak{}}\par
\noindent Corrected reading: {\ttfamily\footnotesize 3\allowbreak{}}\par
\noindent Source: {\ttfamily\footnotesize port/PORT\_EN7I\_LOG.j\allowbreak{}sonl\#C940;\allowbreak{} fixes/FIXLOG\_7I\_CURR\allowbreak{}ENT.jsonl\#C940;\allowbreak{} 7I:924\allowbreak{}}\par

\par\medskip\noindent\hypertarget{app-E7I-IX-088.1}{}\textbf{E7I-IX-088.1}\quad printed folio 465.\par
\noindent Printed reading: \emph{(empty)}\par
\noindent Corrected reading: {\ttfamily\footnotesize \_i\allowbreak{}}\par
\noindent Source: {\ttfamily\footnotesize ED-EN-FORMULA-G073;\allowbreak{} fr/sga7i/\{diplomatic\allowbreak{},current\}/expose\_IX\_\allowbreak{}body.tex;\allowbreak{} folio\allowbreak{} 465;\allowbreak{} port/decisions/decis\allowbreak{}ions\_en7I\_r2\_004.jso\allowbreak{}n\#C941;\allowbreak{} fixes/FIXLOG\_7I\_CURR\allowbreak{}ENT.jsonl\#C941;\allowbreak{} 7I:716\allowbreak{}}\par

\par\medskip\noindent\hypertarget{app-E7I-IX-089.1}{}\textbf{E7I-IX-089.1}\quad printed folio 467.\par
\noindent Printed reading: {\ttfamily\footnotesize S\allowbreak{}}\par
\noindent Corrected reading: {\ttfamily\footnotesize X\allowbreak{}}\par
\noindent Source: {\ttfamily\footnotesize fr/sga7i/\{diplomatic\allowbreak{},current\}/expose\_IX\_\allowbreak{}body.tex;\allowbreak{} folio\allowbreak{} 467;\allowbreak{} formula\allowbreak{} 4186\allowbreak{}}\par

\par\medskip\noindent\hypertarget{app-E7I-IX-089.2}{}\textbf{E7I-IX-089.2}\quad printed folio 467.\par
\noindent Printed reading: \emph{(empty)}\par
\noindent Corrected reading: {\ttfamily\footnotesize o\textasciicircum{}\allowbreak{}}\par
\noindent Source: {\ttfamily\footnotesize fr/sga7i/\{diplomatic\allowbreak{},current\}/expose\_IX\_\allowbreak{}body.tex;\allowbreak{} folio\allowbreak{} 467;\allowbreak{} formula\allowbreak{} 4189;\allowbreak{} port/decisions/decis\allowbreak{}ions\_en7I\_r2\_004.jso\allowbreak{}n\#C942;\allowbreak{} fixes/FIXLOG\_7I\_CURR\allowbreak{}ENT.jsonl\#C942;\allowbreak{} 7I:717\allowbreak{}}\par

\par\medskip\noindent\hypertarget{app-E7I-IX-089.3}{}\textbf{E7I-IX-089.3}\quad printed folio 467.\par
\noindent Printed reading: {\ttfamily\footnotesize \textasciicircum{}o\allowbreak{}}\par
\noindent Corrected reading: \emph{(empty)}\par
\noindent Source: {\ttfamily\footnotesize fr/sga7i/\{diplomatic\allowbreak{},current\}/expose\_IX\_\allowbreak{}body.tex;\allowbreak{} folio\allowbreak{} 467;\allowbreak{} formula\allowbreak{} 4189;\allowbreak{} port/decisions/decis\allowbreak{}ions\_en7I\_r2\_004.jso\allowbreak{}n\#C942;\allowbreak{} fixes/FIXLOG\_7I\_CURR\allowbreak{}ENT.jsonl\#C942;\allowbreak{} 7I:717\allowbreak{}}\par

\par\medskip\noindent\hypertarget{app-E7I-IX-090.1}{}\textbf{E7I-IX-090.1}\quad printed folio 467.\par
\noindent Printed reading: {\ttfamily\footnotesize special\allowbreak{}}\par
\noindent Corrected reading: {\ttfamily\footnotesize generic\allowbreak{}}\par
\noindent Source: {\ttfamily\footnotesize port/PORT\_EN7I\_LOG.j\allowbreak{}sonl\#C943;\allowbreak{} fixes/FIXLOG\_7I\_CURR\allowbreak{}ENT.jsonl\#C943;\allowbreak{} 7I:718\allowbreak{}}\par

\par\medskip\noindent\hypertarget{app-E7I-IX-091.1}{}\textbf{E7I-IX-091.1}\quad printed folio 469.\par
\noindent Printed reading: {\ttfamily\footnotesize 12\allowbreak{}}\par
\noindent Corrected reading: {\ttfamily\footnotesize 10\allowbreak{}}\par
\noindent Source: {\ttfamily\footnotesize port/PORT\_EN7I\_LOG.j\allowbreak{}sonl\#C944;\allowbreak{} fixes/FIXLOG\_7I\_CURR\allowbreak{}ENT.jsonl\#C944;\allowbreak{} 7I:928\allowbreak{}}\par

\par\medskip\noindent\hypertarget{app-E7I-IX-092.1}{}\textbf{E7I-IX-092.1}\quad printed folio 470.\par
\noindent Printed reading: {\ttfamily\footnotesize X\allowbreak{}}\par
\noindent Corrected reading: {\ttfamily\footnotesize C\allowbreak{}}\par
\noindent Source: {\ttfamily\footnotesize fr/sga7i/\{diplomatic\allowbreak{},current\}/expose\_IX\_\allowbreak{}body.tex;\allowbreak{} folio\allowbreak{} 470;\allowbreak{} formula\allowbreak{} 4261\allowbreak{}}\par

\par\medskip\noindent\hypertarget{app-E7I-IX-093.1}{}\textbf{E7I-IX-093.1}\quad printed folio 471.\par
\noindent Printed reading: {\ttfamily\footnotesize 2\allowbreak{}}\par
\noindent Corrected reading: {\ttfamily\footnotesize 3\allowbreak{}}\par
\noindent Source: {\ttfamily\footnotesize 7I:719;\allowbreak{} 7I:929;\allowbreak{} fixes/MISPRINTS\_7I.j\allowbreak{}son\#719;\allowbreak{} fixes/MISPRINTS\_7I.j\allowbreak{}son\#929;\allowbreak{} ED-EN-NORM-G027;\allowbreak{} complete\allowbreak{} normalized-equal\allowbreak{} formula\allowbreak{} audit;\allowbreak{} fr/sga7i/\{diplomatic\allowbreak{},current\}/expose\_IX\_\allowbreak{}body.tex;\allowbreak{} folio\allowbreak{} 471\allowbreak{}}\par

\par\medskip\noindent\hypertarget{app-E7I-IX-094.1}{}\textbf{E7I-IX-094.1}\quad printed folio 471.\par
\noindent Printed reading: {\ttfamily\footnotesize 2\allowbreak{}}\par
\noindent Corrected reading: {\ttfamily\footnotesize 3\allowbreak{}}\par
\noindent Source: {\ttfamily\footnotesize 7I:720;\allowbreak{} 7I:930;\allowbreak{} fixes/MISPRINTS\_7I.j\allowbreak{}son\#720;\allowbreak{} fixes/MISPRINTS\_7I.j\allowbreak{}son\#930;\allowbreak{} ED-EN-NORM-G028;\allowbreak{} complete\allowbreak{} normalized-equal\allowbreak{} formula\allowbreak{} audit;\allowbreak{} fr/sga7i/\{diplomatic\allowbreak{},current\}/expose\_IX\_\allowbreak{}body.tex;\allowbreak{} folio\allowbreak{} 471\allowbreak{}}\par

\par\medskip\noindent\hypertarget{app-E7I-IX-095.1}{}\textbf{E7I-IX-095.1}\quad printed folio 472.\par
\noindent Printed reading: {\ttfamily\footnotesize \textbackslash{}notag\textbackslash{}SGARefAliasTa\allowbreak{}rget\allowbreak{}}\par
\noindent Corrected reading: {\ttfamily\footnotesize \textbackslash{}tag\{\textbackslash{}SGARefTarget\allowbreak{}}\par
\noindent Source: {\ttfamily\footnotesize ED-EN-NORM-G029;\allowbreak{} complete\allowbreak{} normalized-equal\allowbreak{} formula\allowbreak{} audit;\allowbreak{} fr/sga7i/\{diplomatic\allowbreak{},current\}/expose\_IX\_\allowbreak{}body.tex;\allowbreak{} folio\allowbreak{} 472\allowbreak{}}\par

\par\medskip\noindent\hypertarget{app-E7I-IX-095.2}{}\textbf{E7I-IX-095.2}\quad printed folio 472.\par
\noindent Printed reading: \emph{(empty)}\par
\noindent Corrected reading: {\ttfamily\footnotesize \}\{12.3.8\}\allowbreak{}}\par
\noindent Source: {\ttfamily\footnotesize ED-EN-NORM-G029;\allowbreak{} complete\allowbreak{} normalized-equal\allowbreak{} formula\allowbreak{} audit;\allowbreak{} fr/sga7i/\{diplomatic\allowbreak{},current\}/expose\_IX\_\allowbreak{}body.tex;\allowbreak{} folio\allowbreak{} 472\allowbreak{}}\par

\par\medskip\noindent\hypertarget{app-E7I-IX-096.1}{}\textbf{E7I-IX-096.1}\quad printed folio 472.\par
\noindent Printed reading: {\ttfamily\footnotesize 3\allowbreak{}}\par
\noindent Corrected reading: {\ttfamily\footnotesize 2\allowbreak{}}\par
\noindent Source: {\ttfamily\footnotesize port/PORT\_EN7I\_LOG.j\allowbreak{}sonl\#C948;\allowbreak{} fixes/FIXLOG\_7I\_CURR\allowbreak{}ENT.jsonl\#C948;\allowbreak{} 7I:936\allowbreak{}}\par

\par\medskip\noindent\hypertarget{app-E7I-IX-097.1}{}\textbf{E7I-IX-097.1}\quad printed folio 473.\par
\noindent Printed reading: \emph{(empty)}\par
\noindent Corrected reading: {\ttfamily\footnotesize \_1\allowbreak{}}\par
\noindent Source: {\ttfamily\footnotesize fr/sga7i/\{diplomatic\allowbreak{},current\}/expose\_IX\_\allowbreak{}body.tex;\allowbreak{} folio\allowbreak{} 473\allowbreak{}}\par

\par\medskip\noindent\hypertarget{app-E7I-IX-098.1}{}\textbf{E7I-IX-098.1}\quad printed folio 473.\par
\noindent Printed reading: \emph{(empty)}\par
\noindent Corrected reading: {\ttfamily\footnotesize \_\textbackslash{}eta\allowbreak{}}\par
\noindent Source: {\ttfamily\footnotesize fr/sga7i/\{diplomatic\allowbreak{},current\}/expose\_IX\_\allowbreak{}body.tex;\allowbreak{} folio\allowbreak{} 473;\allowbreak{} formula\allowbreak{} 4358;\allowbreak{} port/PORT\_EN7I\_LOG.j\allowbreak{}sonl\#C951;\allowbreak{} fixes/FIXLOG\_7I\_CURR\allowbreak{}ENT.jsonl\#C951;\allowbreak{} 7I:941\allowbreak{}}\par

\par\medskip\noindent\hypertarget{app-E7I-IX-099.1}{}\textbf{E7I-IX-099.1}\quad printed folio 475.\par
\noindent Printed reading: {\ttfamily\footnotesize ,\allowbreak{}}\par
\noindent Corrected reading: {\ttfamily\footnotesize .\allowbreak{}}\par
\noindent Source: {\ttfamily\footnotesize port/decisions/decis\allowbreak{}ions\_en7I\_r2\_004.jso\allowbreak{}n\#C955;\allowbreak{} fixes/FIXLOG\_7I\_CURR\allowbreak{}ENT.jsonl\#C955;\allowbreak{} 7I:732\allowbreak{}}\par

\par\medskip\noindent\hypertarget{app-E7I-IX-100.1}{}\textbf{E7I-IX-100.1}\quad printed folio 477.\par
\noindent Printed reading: \emph{(empty)}\par
\noindent Corrected reading: {\ttfamily\footnotesize )\allowbreak{}}\par
\noindent Source: {\ttfamily\footnotesize port/decisions/decis\allowbreak{}ions\_en7I\_r2\_004.jso\allowbreak{}n\#C958;\allowbreak{} fixes/FIXLOG\_7I\_CURR\allowbreak{}ENT.jsonl\#C958;\allowbreak{} 7I:735\allowbreak{}}\par

\par\medskip\noindent\hypertarget{app-E7I-IX-101.1}{}\textbf{E7I-IX-101.1}\quad printed folio 478.\par
\noindent Printed reading: {\ttfamily\footnotesize F\allowbreak{}}\par
\noindent Corrected reading: {\ttfamily\footnotesize f\allowbreak{}}\par
\noindent Source: {\ttfamily\footnotesize fr/sga7i/\{diplomatic\allowbreak{},current\}/expose\_IX\_\allowbreak{}body.tex;\allowbreak{} folio\allowbreak{} 478;\allowbreak{} formula\allowbreak{} 4472\allowbreak{}}\par

\par\medskip\noindent\hypertarget{app-E7I-IX-101.2}{}\textbf{E7I-IX-101.2}\quad printed folio 478.\par
\noindent Printed reading: \emph{(empty)}\par
\noindent Corrected reading: {\ttfamily\footnotesize \textbackslash{}underline\allowbreak{}}\par
\noindent Source: {\ttfamily\footnotesize fr/sga7i/\{diplomatic\allowbreak{},current\}/expose\_IX\_\allowbreak{}body.tex;\allowbreak{} folio\allowbreak{} 478;\allowbreak{} formula\allowbreak{} 4474;\allowbreak{} port/decisions/decis\allowbreak{}ions\_en7I\_r2\_004.jso\allowbreak{}n\#C959;\allowbreak{} fixes/FIXLOG\_7I\_CURR\allowbreak{}ENT.jsonl\#C959;\allowbreak{} 7I:961\allowbreak{}}\par

\par\medskip\noindent\hypertarget{app-E7I-IX-101.3}{}\textbf{E7I-IX-101.3}\quad printed folio 478.\par
\noindent Printed reading: \emph{(empty)}\par
\noindent Corrected reading: {\ttfamily\footnotesize \textbackslash{}underline\allowbreak{}}\par
\noindent Source: {\ttfamily\footnotesize fr/sga7i/\{diplomatic\allowbreak{},current\}/expose\_IX\_\allowbreak{}body.tex;\allowbreak{} folio\allowbreak{} 478;\allowbreak{} formula\allowbreak{} 4478;\allowbreak{} port/decisions/decis\allowbreak{}ions\_en7I\_r2\_004.jso\allowbreak{}n\#C960;\allowbreak{} fixes/FIXLOG\_7I\_CURR\allowbreak{}ENT.jsonl\#C960;\allowbreak{} 7I:962\allowbreak{}}\par

\par\medskip\noindent\hypertarget{app-E7I-IX-102.1}{}\textbf{E7I-IX-102.1}\quad printed folio 480.\par
\noindent Printed reading: {\ttfamily\footnotesize B\allowbreak{}}\par
\noindent Corrected reading: {\ttfamily\footnotesize P\allowbreak{}}\par
\noindent Source: {\ttfamily\footnotesize fr/sga7i/\{diplomatic\allowbreak{},current\}/expose\_IX\_\allowbreak{}body.tex;\allowbreak{} folio\allowbreak{} 480;\allowbreak{} formula\allowbreak{} 4510\allowbreak{}}\par

\par\medskip\noindent\hypertarget{app-E7I-IX-102.2}{}\textbf{E7I-IX-102.2}\quad printed folio 480.\par
\noindent Printed reading: \emph{(empty)}\par
\noindent Corrected reading: {\ttfamily\footnotesize \textbackslash{}underline\allowbreak{}}\par
\noindent Source: {\ttfamily\footnotesize fr/sga7i/\{diplomatic\allowbreak{},current\}/expose\_IX\_\allowbreak{}body.tex;\allowbreak{} folio\allowbreak{} 480;\allowbreak{} formula\allowbreak{} 4516;\allowbreak{} port/decisions/decis\allowbreak{}ions\_en7I\_r2\_004.jso\allowbreak{}n\#C961;\allowbreak{} fixes/FIXLOG\_7I\_CURR\allowbreak{}ENT.jsonl\#C961;\allowbreak{} 7I:964\allowbreak{}}\par

\par\medskip\noindent\hypertarget{app-E7I-IX-103.1}{}\textbf{E7I-IX-103.1}\quad printed folio 480.\par
\noindent Printed reading: {\ttfamily\footnotesize \textbackslash{}underline\{\allowbreak{}}\par
\noindent Corrected reading: \emph{(empty)}\par
\noindent Source: {\ttfamily\footnotesize fr/sga7i/\{diplomatic\allowbreak{},current\}/expose\_IX\_\allowbreak{}body.tex;\allowbreak{} folio\allowbreak{} 480;\allowbreak{} formula\allowbreak{} 4524\allowbreak{}}\par

\par\medskip\noindent\hypertarget{app-E7I-IX-103.2}{}\textbf{E7I-IX-103.2}\quad printed folio 480.\par
\noindent Printed reading: {\ttfamily\footnotesize \}\allowbreak{}}\par
\noindent Corrected reading: \emph{(empty)}\par
\noindent Source: {\ttfamily\footnotesize fr/sga7i/\{diplomatic\allowbreak{},current\}/expose\_IX\_\allowbreak{}body.tex;\allowbreak{} folio\allowbreak{} 480;\allowbreak{} formula\allowbreak{} 4524\allowbreak{}}\par

\par\medskip\noindent\hypertarget{app-E7I-IX-103.3}{}\textbf{E7I-IX-103.3}\quad printed folio 480.\par
\noindent Printed reading: \emph{(empty)}\par
\noindent Corrected reading: {\ttfamily\footnotesize )\textbackslash{}bigr\allowbreak{}}\par
\noindent Source: {\ttfamily\footnotesize fr/sga7i/\{diplomatic\allowbreak{},current\}/expose\_IX\_\allowbreak{}body.tex;\allowbreak{} folio\allowbreak{} 480;\allowbreak{} formula\allowbreak{} 4524\allowbreak{}}\par

\par\medskip\noindent\hypertarget{app-E7I-IX-103.4}{}\textbf{E7I-IX-103.4}\quad printed folio 481.\par
\noindent Printed reading: \emph{(empty)}\par
\noindent Corrected reading: {\ttfamily\footnotesize \textbackslash{}bigr)\allowbreak{}}\par
\noindent Source: {\ttfamily\footnotesize fr/sga7i/\{diplomatic\allowbreak{},current\}/expose\_IX\_\allowbreak{}body.tex;\allowbreak{} folio\allowbreak{} 481;\allowbreak{} formula\allowbreak{} 4532;\allowbreak{} port/decisions/decis\allowbreak{}ions\_en7I\_r2\_004.jso\allowbreak{}n\#C962;\allowbreak{} fixes/FIXLOG\_7I\_CURR\allowbreak{}ENT.jsonl\#C962;\allowbreak{} 7I:740\allowbreak{}}\par

\par\medskip\noindent\hypertarget{app-E7I-IX-104.1}{}\textbf{E7I-IX-104.1}\quad printed folio 481.\par
\noindent Printed reading: \emph{(empty)}\par
\noindent Corrected reading: {\ttfamily\footnotesize \textbackslash{}underline\allowbreak{}}\par
\noindent Source: {\ttfamily\footnotesize fr/sga7i/\{diplomatic\allowbreak{},current\}/expose\_IX\_\allowbreak{}body.tex;\allowbreak{} folio\allowbreak{} 481;\allowbreak{} formula\allowbreak{} 4544;\allowbreak{} port/decisions/decis\allowbreak{}ions\_en7I\_r2\_004.jso\allowbreak{}n\#C963;\allowbreak{} fixes/FIXLOG\_7I\_CURR\allowbreak{}ENT.jsonl\#C963;\allowbreak{} 7I:741\allowbreak{}}\par

\par\medskip\noindent\hypertarget{app-E7I-IX-105.1}{}\textbf{E7I-IX-105.1}\quad printed folio 482.\par
\noindent Printed reading: {\ttfamily\footnotesize x\allowbreak{}}\par
\noindent Corrected reading: {\ttfamily\footnotesize X\allowbreak{}}\par
\noindent Source: {\ttfamily\footnotesize fr/sga7i/\{diplomatic\allowbreak{},current\}/expose\_IX\_\allowbreak{}body.tex;\allowbreak{} folio\allowbreak{} 482;\allowbreak{} formula\allowbreak{} 4567\allowbreak{}}\par

\par\medskip\noindent\hypertarget{app-E7I-IX-106.1}{}\textbf{E7I-IX-106.1}\quad printed folio 482.\par
\noindent Printed reading: \emph{(empty)}\par
\noindent Corrected reading: {\ttfamily\footnotesize \_\textbackslash{}lambda\allowbreak{}}\par
\noindent Source: {\ttfamily\footnotesize fr/sga7i/\{diplomatic\allowbreak{},current\}/expose\_IX\_\allowbreak{}body.tex;\allowbreak{} folio\allowbreak{} 482;\allowbreak{} formula\allowbreak{} 4577\allowbreak{}}\par

\par\medskip\noindent\hypertarget{app-E7I-IX-106.2}{}\textbf{E7I-IX-106.2}\quad printed folio 482.\par
\noindent Printed reading: \emph{(2 espace(s), 1 retour(s) de ligne dans le témoin TeX)}\par
\noindent Corrected reading: {\ttfamily\footnotesize  \textbackslash{}underline\allowbreak{} }\par
\noindent Source: {\ttfamily\footnotesize fr/sga7i/\{diplomatic\allowbreak{},current\}/expose\_IX\_\allowbreak{}body.tex;\allowbreak{} folio\allowbreak{} 482;\allowbreak{} formula\allowbreak{} 4578\allowbreak{}}\par

\par\medskip\noindent\hypertarget{app-E7I-IX-106.3}{}\textbf{E7I-IX-106.3}\quad printed folio 482.\par
\noindent Printed reading: \emph{(empty)}\par
\noindent Corrected reading: {\ttfamily\footnotesize \textbackslash{}underline\allowbreak{}}\par
\noindent Source: {\ttfamily\footnotesize fr/sga7i/\{diplomatic\allowbreak{},current\}/expose\_IX\_\allowbreak{}body.tex;\allowbreak{} folio\allowbreak{} 482;\allowbreak{} formula\allowbreak{} 4578\allowbreak{}}\par

\par\medskip\noindent\hypertarget{app-E7I-IX-106.4}{}\textbf{E7I-IX-106.4}\quad printed folio 482.\par
\noindent Printed reading: \emph{(empty)}\par
\noindent Corrected reading: {\ttfamily\footnotesize \textbackslash{}underline\allowbreak{}}\par
\noindent Source: {\ttfamily\footnotesize fr/sga7i/\{diplomatic\allowbreak{},current\}/expose\_IX\_\allowbreak{}body.tex;\allowbreak{} folio\allowbreak{} 482;\allowbreak{} formula\allowbreak{} 4579;\allowbreak{} port/decisions/decis\allowbreak{}ions\_en7I\_r2\_004.jso\allowbreak{}n\#C964;\allowbreak{} fixes/FIXLOG\_7I\_CURR\allowbreak{}ENT.jsonl\#C964;\allowbreak{} 7I:952\allowbreak{}}\par

\par\medskip\noindent\hypertarget{app-E7I-IX-107.1}{}\textbf{E7I-IX-107.1}\quad printed folio 483.\par
\noindent Printed reading: {\ttfamily\footnotesize \textbackslash{}lambda\allowbreak{}}\par
\noindent Corrected reading: {\ttfamily\footnotesize y\allowbreak{}}\par
\noindent Source: {\ttfamily\footnotesize fr/sga7i/\{diplomatic\allowbreak{},current\}/expose\_IX\_\allowbreak{}body.tex;\allowbreak{} folio\allowbreak{} 483;\allowbreak{} formula\allowbreak{} 4594\allowbreak{}}\par

\par\medskip\noindent\hypertarget{app-E7I-IX-108.1}{}\textbf{E7I-IX-108.1}\quad printed folio 484.\par
\noindent Printed reading: \emph{(empty)}\par
\noindent Corrected reading: {\ttfamily\footnotesize \textasciicircum{}\textbackslash{}lambda\allowbreak{}}\par
\noindent Source: {\ttfamily\footnotesize fr/sga7i/\{diplomatic\allowbreak{},current\}/expose\_IX\_\allowbreak{}body.tex;\allowbreak{} folio\allowbreak{} 484;\allowbreak{} formula\allowbreak{} 4630\allowbreak{}}\par

\par\medskip\noindent\hypertarget{app-E7I-IX-109.1}{}\textbf{E7I-IX-109.1}\quad printed folio 487.\par
\noindent Printed reading: {\ttfamily\footnotesize S\allowbreak{}}\par
\noindent Corrected reading: {\ttfamily\footnotesize s\allowbreak{}}\par
\noindent Source: {\ttfamily\footnotesize fr/sga7i/\{diplomatic\allowbreak{},current\}/expose\_IX\_\allowbreak{}body.tex;\allowbreak{} folio\allowbreak{} 487;\allowbreak{} formula\allowbreak{} 4719;\allowbreak{} port/decisions/decis\allowbreak{}ions\_en7I\_r2\_004.jso\allowbreak{}n\#C968;\allowbreak{} fixes/FIXLOG\_7I\_CURR\allowbreak{}ENT.jsonl\#C968;\allowbreak{} 7I:749\allowbreak{}}\par

\par\medskip\noindent\hypertarget{app-E7I-IX-110.1}{}\textbf{E7I-IX-110.1}\quad printed folio 487.\par
\noindent Printed reading: {\ttfamily\footnotesize  C\allowbreak{}}\par
\noindent Corrected reading: {\ttfamily\footnotesize 0\allowbreak{}}\par
\noindent Source: {\ttfamily\footnotesize fr/sga7i/\{diplomatic\allowbreak{},current\}/expose\_IX\_\allowbreak{}body.tex;\allowbreak{} folio\allowbreak{} 487;\allowbreak{} formula\allowbreak{} 4734;\allowbreak{} port/PORT\_EN7I\_LOG.j\allowbreak{}sonl\#C969;\allowbreak{} fixes/FIXLOG\_7I\_CURR\allowbreak{}ENT.jsonl\#C969;\allowbreak{} 7I:750\allowbreak{}}\par

\par\medskip\noindent\hypertarget{app-E7I-IX-111.1}{}\textbf{E7I-IX-111.1}\quad printed folio 488.\par
\noindent Printed reading: \emph{(empty)}\par
\noindent Corrected reading: {\ttfamily\footnotesize \textbackslash{}underline\allowbreak{}}\par
\noindent Source: {\ttfamily\footnotesize fr/sga7i/\{diplomatic\allowbreak{},current\}/expose\_IX\_\allowbreak{}body.tex;\allowbreak{} folio\allowbreak{} 488;\allowbreak{} formula\allowbreak{} 4754;\allowbreak{} port/decisions/decis\allowbreak{}ions\_en7I\_r2\_004.jso\allowbreak{}n\#C970;\allowbreak{} fixes/FIXLOG\_7I\_CURR\allowbreak{}ENT.jsonl\#C970;\allowbreak{} 7I:751\allowbreak{}}\par

\par\medskip\noindent\hypertarget{app-E7I-IX-111.2}{}\textbf{E7I-IX-111.2}\quad printed folio 488.\par
\noindent Printed reading: {\ttfamily\footnotesize s\allowbreak{}}\par
\noindent Corrected reading: {\ttfamily\footnotesize S\allowbreak{}}\par
\noindent Source: {\ttfamily\footnotesize fr/sga7i/\{diplomatic\allowbreak{},current\}/expose\_IX\_\allowbreak{}body.tex;\allowbreak{} folio\allowbreak{} 488;\allowbreak{} formula\allowbreak{} 4757;\allowbreak{} port/PORT\_EN7I\_LOG.j\allowbreak{}sonl\#C971;\allowbreak{} fixes/FIXLOG\_7I\_CURR\allowbreak{}ENT.jsonl\#C971;\allowbreak{} 7I:752\allowbreak{}}\par

\par\medskip\noindent\hypertarget{app-E7I-IX-112.1}{}\textbf{E7I-IX-112.1}\quad printed folio 488.\par
\noindent Printed reading: \emph{(empty)}\par
\noindent Corrected reading: {\ttfamily\footnotesize \textbackslash{}underline\allowbreak{}}\par
\noindent Source: {\ttfamily\footnotesize fr/sga7i/\{diplomatic\allowbreak{},current\}/expose\_IX\_\allowbreak{}body.tex;\allowbreak{} folio\allowbreak{} 488;\allowbreak{} formula\allowbreak{} 4783\allowbreak{}}\par

\par\medskip\noindent\hypertarget{app-E7I-IX-112.2}{}\textbf{E7I-IX-112.2}\quad printed folio 489.\par
\noindent Printed reading: {\ttfamily\footnotesize 2\allowbreak{}}\par
\noindent Corrected reading: {\ttfamily\footnotesize 1\allowbreak{}}\par
\noindent Source: {\ttfamily\footnotesize 7I:972;\allowbreak{} fixes/MISPRINTS\_7I.j\allowbreak{}son\#972\allowbreak{}}\par

\par\medskip\noindent\hypertarget{app-E7I-IX-113.1}{}\textbf{E7I-IX-113.1}\quad printed folio 489.\par
\noindent Printed reading: {\ttfamily\footnotesize \_\allowbreak{}}\par
\noindent Corrected reading: {\ttfamily\footnotesize \textasciicircum{}\allowbreak{}}\par
\noindent Source: {\ttfamily\footnotesize fr/sga7i/\{diplomatic\allowbreak{},current\}/expose\_IX\_\allowbreak{}body.tex;\allowbreak{} folio\allowbreak{} 489;\allowbreak{} formula\allowbreak{} 4806\allowbreak{}}\par

\par\medskip\noindent\hypertarget{app-E7I-IX-113.2}{}\textbf{E7I-IX-113.2}\quad printed folio 489.\par
\noindent Printed reading: {\ttfamily\footnotesize 3\allowbreak{}}\par
\noindent Corrected reading: {\ttfamily\footnotesize s\allowbreak{}}\par
\noindent Source: {\ttfamily\footnotesize fr/sga7i/\{diplomatic\allowbreak{},current\}/expose\_IX\_\allowbreak{}body.tex;\allowbreak{} folio\allowbreak{} 489;\allowbreak{} formula\allowbreak{} 4807;\allowbreak{} port/decisions/decis\allowbreak{}ions\_en7I\_r2\_004.jso\allowbreak{}n\#C973;\allowbreak{} fixes/FIXLOG\_7I\_CURR\allowbreak{}ENT.jsonl\#C973;\allowbreak{} 7I:755\allowbreak{}}\par

\par\medskip\noindent\hypertarget{app-E7I-IX-114.1}{}\textbf{E7I-IX-114.1}\quad printed folio 491.\par
\noindent Printed reading: {\ttfamily\footnotesize S\allowbreak{}}\par
\noindent Corrected reading: {\ttfamily\footnotesize s\allowbreak{}}\par
\noindent Source: {\ttfamily\footnotesize fr/sga7i/\{diplomatic\allowbreak{},current\}/expose\_IX\_\allowbreak{}body.tex;\allowbreak{} folio\allowbreak{} 491;\allowbreak{} formula\allowbreak{} 4862;\allowbreak{} port/decisions/decis\allowbreak{}ions\_en7I\_r2\_004.jso\allowbreak{}n\#C975;\allowbreak{} fixes/FIXLOG\_7I\_CURR\allowbreak{}ENT.jsonl\#C975;\allowbreak{} 7I:757\allowbreak{}}\par

\par\medskip\noindent\hypertarget{app-E7I-IX-115.1}{}\textbf{E7I-IX-115.1}\quad printed folio 493.\par
\noindent Printed reading: {\ttfamily\footnotesize \textbackslash{}mathbb\allowbreak{}}\par
\noindent Corrected reading: {\ttfamily\footnotesize \textbackslash{}underline\allowbreak{}}\par
\noindent Source: {\ttfamily\footnotesize ED-EN-FONT-G014;\allowbreak{} F7I-IX-176.1;\allowbreak{} fr/sga7i/\{diplomatic\allowbreak{},current\}/expose\_IX\_\allowbreak{}body.tex;\allowbreak{} folio\allowbreak{} 493\allowbreak{}}\par

\par\medskip\noindent\hypertarget{app-E7I-IX-116.1}{}\textbf{E7I-IX-116.1}\quad printed folio 494.\par
\noindent Printed reading: \emph{(empty)}\par
\noindent Corrected reading: {\ttfamily\footnotesize -\allowbreak{}}\par
\noindent Source: {\ttfamily\footnotesize fr/sga7i/\{diplomatic\allowbreak{},current\}/expose\_IX\_\allowbreak{}body.tex;\allowbreak{} folio\allowbreak{} 494;\allowbreak{} formula\allowbreak{} 4937;\allowbreak{} port/PORT\_EN7I\_LOG.j\allowbreak{}sonl\#C594;\allowbreak{} fixes/FIXLOG\_7I\_CURR\allowbreak{}ENT.jsonl\#C594;\allowbreak{} }\par

\par\medskip\noindent\hypertarget{app-E7I-IX-116.2}{}\textbf{E7I-IX-116.2}\quad printed folio 494.\par
\noindent Printed reading: {\ttfamily\footnotesize -\allowbreak{}}\par
\noindent Corrected reading: \emph{(empty)}\par
\noindent Source: {\ttfamily\footnotesize fr/sga7i/\{diplomatic\allowbreak{},current\}/expose\_IX\_\allowbreak{}body.tex;\allowbreak{} folio\allowbreak{} 494;\allowbreak{} formula\allowbreak{} 4937;\allowbreak{} port/PORT\_EN7I\_LOG.j\allowbreak{}sonl\#C594;\allowbreak{} fixes/FIXLOG\_7I\_CURR\allowbreak{}ENT.jsonl\#C594;\allowbreak{} }\par

\par\medskip\noindent\hypertarget{app-E7I-IX-117.1}{}\textbf{E7I-IX-117.1}\quad printed folio 495.\par
\noindent Printed reading: {\ttfamily\footnotesize \textbackslash{}Lambda\allowbreak{}}\par
\noindent Corrected reading: {\ttfamily\footnotesize A\allowbreak{}}\par
\noindent Source: {\ttfamily\footnotesize fr/sga7i/\{diplomatic\allowbreak{},current\}/expose\_IX\_\allowbreak{}body.tex;\allowbreak{} folio\allowbreak{} 495;\allowbreak{} formula\allowbreak{} 4953;\allowbreak{} port/decisions/decis\allowbreak{}ions\_en7I\_r2\_004.jso\allowbreak{}n\#C978;\allowbreak{} fixes/FIXLOG\_7I\_CURR\allowbreak{}ENT.jsonl\#C978;\allowbreak{} 7I:760\allowbreak{}}\par

\par\medskip\noindent\hypertarget{app-E7I-IX-118.1}{}\textbf{E7I-IX-118.1}\quad printed folio 496.\par
\noindent Printed reading: \emph{(empty)}\par
\noindent Corrected reading: {\ttfamily\footnotesize \textasciicircum{}\textbackslash{}natural\allowbreak{}}\par
\noindent Source: {\ttfamily\footnotesize fr/sga7i/\{diplomatic\allowbreak{},current\}/expose\_IX\_\allowbreak{}body.tex;\allowbreak{} folio\allowbreak{} 496;\allowbreak{} formula\allowbreak{} 4983\allowbreak{}}\par

\par\medskip\noindent\hypertarget{app-E7I-IX-119.1}{}\textbf{E7I-IX-119.1}\quad printed folio 496.\par
\noindent Printed reading: \emph{(empty)}\par
\noindent Corrected reading: {\ttfamily\footnotesize \textbackslash{}mathcal\allowbreak{} E\allowbreak{}}\par
\noindent Source: {\ttfamily\footnotesize ED-EN-FORMULA-G075;\allowbreak{} fr/sga7i/\{diplomatic\allowbreak{},current\}/expose\_IX\_\allowbreak{}body.tex;\allowbreak{} folio\allowbreak{} 496;\allowbreak{} port/decisions/decis\allowbreak{}ions\_en7I\_r2\_004.jso\allowbreak{}n\#C979;\allowbreak{} 7I:761\allowbreak{}}\par

\par\medskip\noindent\hypertarget{app-E7I-IX-120.1}{}\textbf{E7I-IX-120.1}\quad printed folio 501.\par
\noindent Printed reading: {\ttfamily\footnotesize C\allowbreak{}}\par
\noindent Corrected reading: {\ttfamily\footnotesize O\allowbreak{}}\par
\noindent Source: {\ttfamily\footnotesize fr/sga7i/\{diplomatic\allowbreak{},current\}/expose\_IX\_\allowbreak{}body.tex;\allowbreak{} folio\allowbreak{} 501;\allowbreak{} formula\allowbreak{} 5117;\allowbreak{} port/decisions/decis\allowbreak{}ions\_en7I\_r2\_004.jso\allowbreak{}n\#C982;\allowbreak{} fixes/FIXLOG\_7I\_CURR\allowbreak{}ENT.jsonl\#C982;\allowbreak{} 7I:764\allowbreak{}}\par

\par\medskip\noindent\hypertarget{app-E7I-IX-120.2}{}\textbf{E7I-IX-120.2}\quad printed folio 501.\par
\noindent Printed reading: {\ttfamily\footnotesize C\allowbreak{}}\par
\noindent Corrected reading: {\ttfamily\footnotesize O\allowbreak{}}\par
\noindent Source: {\ttfamily\footnotesize fr/sga7i/\{diplomatic\allowbreak{},current\}/expose\_IX\_\allowbreak{}body.tex;\allowbreak{} folio\allowbreak{} 501;\allowbreak{} formula\allowbreak{} 5119;\allowbreak{} port/decisions/decis\allowbreak{}ions\_en7I\_r2\_004.jso\allowbreak{}n\#C983;\allowbreak{} fixes/FIXLOG\_7I\_CURR\allowbreak{}ENT.jsonl\#C983;\allowbreak{} 7I:765\allowbreak{}}\par

\par\medskip\noindent\hypertarget{app-E7I-IX-121.1}{}\textbf{E7I-IX-121.1}\quad printed folio 502.\par
\noindent Printed reading: {\ttfamily\footnotesize S\allowbreak{}}\par
\noindent Corrected reading: {\ttfamily\footnotesize \textbackslash{}hat\{\textbackslash{}underline\allowbreak{} O\}\allowbreak{}}\par
\noindent Source: {\ttfamily\footnotesize fr/sga7i/\{diplomatic\allowbreak{},current\}/expose\_IX\_\allowbreak{}body.tex;\allowbreak{} folio\allowbreak{} 502\allowbreak{}}\par

\par\medskip\noindent\hypertarget{app-E7I-IX-122.1}{}\textbf{E7I-IX-122.1}\quad printed folio 503.\par
\noindent Printed reading: {\ttfamily\footnotesize C\allowbreak{}}\par
\noindent Corrected reading: {\ttfamily\footnotesize 0\allowbreak{}}\par
\noindent Source: {\ttfamily\footnotesize fr/sga7i/\{diplomatic\allowbreak{},current\}/expose\_IX\_\allowbreak{}body.tex;\allowbreak{} folio\allowbreak{} 503;\allowbreak{} formula\allowbreak{} 5189\allowbreak{}}\par

\par\medskip\noindent\hypertarget{app-E7I-IX-123.1}{}\textbf{E7I-IX-123.1}\quad printed folio 504.\par
\noindent Printed reading: \emph{(empty)}\par
\noindent Corrected reading: {\ttfamily\footnotesize \textbackslash{}underline\allowbreak{}}\par
\noindent Source: {\ttfamily\footnotesize fr/sga7i/\{diplomatic\allowbreak{},current\}/expose\_IX\_\allowbreak{}body.tex;\allowbreak{} folio\allowbreak{} 504;\allowbreak{} formula\allowbreak{} 5208\allowbreak{}}\par

\par\medskip\noindent\hypertarget{app-E7I-IX-124.1}{}\textbf{E7I-IX-124.1}\quad printed folio 505.\par
\noindent Printed reading: \emph{(empty)}\par
\noindent Corrected reading: {\ttfamily\footnotesize \textbackslash{}underline\allowbreak{}}\par
\noindent Source: {\ttfamily\footnotesize fr/sga7i/\{diplomatic\allowbreak{},current\}/expose\_IX\_\allowbreak{}body.tex;\allowbreak{} folio\allowbreak{} 505;\allowbreak{} formula\allowbreak{} 5241;\allowbreak{} port/decisions/decis\allowbreak{}ions\_en7I\_r2\_004.jso\allowbreak{}n\#C985;\allowbreak{} fixes/FIXLOG\_7I\_CURR\allowbreak{}ENT.jsonl\#C985;\allowbreak{} 7I:767\allowbreak{}}\par

\par\medskip\noindent\hypertarget{app-E7I-IX-125.1}{}\textbf{E7I-IX-125.1}\quad printed folio 507.\par
\noindent Printed reading: {\ttfamily\footnotesize P\allowbreak{}}\par
\noindent Corrected reading: {\ttfamily\footnotesize p\allowbreak{}}\par
\noindent Source: {\ttfamily\footnotesize fr/sga7i/\{diplomatic\allowbreak{},current\}/expose\_IX\_\allowbreak{}body.tex;\allowbreak{} folio\allowbreak{} 507;\allowbreak{} formula\allowbreak{} 5289;\allowbreak{} port/decisions/decis\allowbreak{}ions\_en7I\_r2\_004.jso\allowbreak{}n\#C987;\allowbreak{} fixes/FIXLOG\_7I\_CURR\allowbreak{}ENT.jsonl\#C987;\allowbreak{} 7I:769\allowbreak{}}\par

\par\medskip\noindent\hypertarget{app-E7I-IX-126.1}{}\textbf{E7I-IX-126.1}\quad printed folio 512.\par
\noindent Printed reading: \emph{(empty)}\par
\noindent Corrected reading: {\ttfamily\footnotesize p':\allowbreak{}}\par
\noindent Source: {\ttfamily\footnotesize fr/sga7i/\{diplomatic\allowbreak{},current\}/expose\_IX\_\allowbreak{}body.tex;\allowbreak{} folio\allowbreak{} 512;\allowbreak{} formula\allowbreak{} 5373\allowbreak{}}\par

\par\medskip\noindent\hypertarget{app-E7I-IX-126.2}{}\textbf{E7I-IX-126.2}\quad printed folio 512.\par
\noindent Printed reading: \emph{(empty)}\par
\noindent Corrected reading: {\ttfamily\footnotesize \textbackslash{}underline\allowbreak{}}\par
\noindent Source: {\ttfamily\footnotesize fr/sga7i/\{diplomatic\allowbreak{},current\}/expose\_IX\_\allowbreak{}body.tex;\allowbreak{} folio\allowbreak{} 512;\allowbreak{} formula\allowbreak{} 5378\allowbreak{}}\par

\par\medskip\noindent\hypertarget{app-E7I-IX-126.3}{}\textbf{E7I-IX-126.3}\quad printed folio 512.\par
\noindent Printed reading: {\ttfamily\footnotesize 4\allowbreak{}}\par
\noindent Corrected reading: {\ttfamily\footnotesize 14\allowbreak{}}\par
\noindent Source: {\ttfamily\footnotesize ED-EN-NONMATH-I-IX51\allowbreak{}2;\allowbreak{} WITNESS-F7I-IX-194.4\allowbreak{};\allowbreak{} fr/sga7i/diplomatic/\allowbreak{}expose\_IX\_body.tex;\allowbreak{} fr/sga7i/current/exp\allowbreak{}ose\_IX\_body.tex;\allowbreak{} F7I-IX-194.4\allowbreak{}}\par

\par\medskip\noindent\hypertarget{app-E7I-IX-127.1}{}\textbf{E7I-IX-127.1}\quad printed folio 514.\par
\noindent Printed reading: \emph{(empty)}\par
\noindent Corrected reading: {\ttfamily\footnotesize \textbackslash{}bigr)\allowbreak{}}\par
\noindent Source: {\ttfamily\footnotesize fr/sga7i/\{diplomatic\allowbreak{},current\}/expose\_IX\_\allowbreak{}body.tex;\allowbreak{} folio\allowbreak{} 514;\allowbreak{} formula\allowbreak{} 5429\allowbreak{}}\par

\par\medskip\noindent\hypertarget{app-E7I-IX-128.1}{}\textbf{E7I-IX-128.1}\quad printed folio 515.\par
\noindent Printed reading: {\ttfamily\footnotesize \textbackslash{}mathbb\allowbreak{} }\par
\noindent Corrected reading: \emph{(empty)}\par
\noindent Source: {\ttfamily\footnotesize ED-EN-FONT-G003;\allowbreak{} F7I-IX-197.2;\allowbreak{} fr/sga7i/\{diplomatic\allowbreak{},current\}/expose\_IX\_\allowbreak{}body.tex;\allowbreak{} folio\allowbreak{} 515\allowbreak{}}\par

\par\medskip\noindent\hypertarget{app-E7I-IX-128.2}{}\textbf{E7I-IX-128.2}\quad printed folio 516.\par
\noindent Printed reading: {\ttfamily\footnotesize \textbackslash{}longrightarrow\allowbreak{} }\par
\noindent Corrected reading: {\ttfamily\footnotesize =\allowbreak{}}\par
\noindent Source: {\ttfamily\footnotesize fr/sga7i/\{diplomatic\allowbreak{},current\}/expose\_IX\_\allowbreak{}body.tex;\allowbreak{} folio\allowbreak{} 516;\allowbreak{} formula\allowbreak{} 5480\allowbreak{}}\par

\par\medskip\noindent\hypertarget{app-E7I-IX-129.1}{}\textbf{E7I-IX-129.1}\quad printed folio 517.\par
\noindent Printed reading: {\ttfamily\footnotesize 22\allowbreak{}}\par
\noindent Corrected reading: {\ttfamily\footnotesize 2.2\allowbreak{}}\par
\noindent Source: {\ttfamily\footnotesize 7I:1015;\allowbreak{} fixes/MISPRINTS\_7I.j\allowbreak{}son\#1015\allowbreak{}}\par

\par\medskip\noindent\hypertarget{app-E7I-IX-130.1}{}\textbf{E7I-IX-130.1}\quad printed folio 518.\par
\noindent Printed reading: {\ttfamily\footnotesize G\allowbreak{}}\par
\noindent Corrected reading: {\ttfamily\footnotesize A\allowbreak{}}\par
\noindent Source: {\ttfamily\footnotesize fr/sga7i/\{diplomatic\allowbreak{},current\}/expose\_IX\_\allowbreak{}body.tex;\allowbreak{} folio\allowbreak{} 518;\allowbreak{} formula\allowbreak{} 5530;\allowbreak{} port/decisions/decis\allowbreak{}ions\_en7I\_r2\_004.jso\allowbreak{}n\#C997;\allowbreak{} fixes/FIXLOG\_7I\_CURR\allowbreak{}ENT.jsonl\#C997;\allowbreak{} 7I:784\allowbreak{}}\par

\par\medskip\noindent\hypertarget{app-E7I-IX-131.1}{}\textbf{E7I-IX-131.1}\quad printed folio 518.\par
\noindent Printed reading: \emph{(empty)}\par
\noindent Corrected reading: {\ttfamily\footnotesize '\allowbreak{}}\par
\noindent Source: {\ttfamily\footnotesize fr/sga7i/\{diplomatic\allowbreak{},current\}/expose\_IX\_\allowbreak{}body.tex;\allowbreak{} folio\allowbreak{} 518;\allowbreak{} formula\allowbreak{} 5546;\allowbreak{} port/decisions/decis\allowbreak{}ions\_en7I\_r2\_004.jso\allowbreak{}n\#C998;\allowbreak{} fixes/FIXLOG\_7I\_CURR\allowbreak{}ENT.jsonl\#C998;\allowbreak{} 7I:785\allowbreak{}}\par

\par\medskip\noindent\hypertarget{app-E7I-IX-132.1}{}\textbf{E7I-IX-132.1}\quad printed folio 518.\par
\noindent Printed reading: {\ttfamily\footnotesize 8\allowbreak{}}\par
\noindent Corrected reading: {\ttfamily\footnotesize 5.3\allowbreak{}}\par
\noindent Source: {\ttfamily\footnotesize 7I:1018;\allowbreak{} fixes/MISPRINTS\_7I.j\allowbreak{}son\#1018\allowbreak{}}\par

\par\medskip\noindent\hypertarget{app-E7I-IX-132.2}{}\textbf{E7I-IX-132.2}\quad printed folio 518.\par
\noindent Printed reading: {\ttfamily\footnotesize 10\allowbreak{}}\par
\noindent Corrected reading: {\ttfamily\footnotesize 5.5\allowbreak{}}\par
\noindent Source: {\ttfamily\footnotesize 7I:1018;\allowbreak{} fixes/MISPRINTS\_7I.j\allowbreak{}son\#1018\allowbreak{}}\par

\par\medskip\noindent\hypertarget{app-E7I-IX-133.1}{}\textbf{E7I-IX-133.1}\quad printed folio 520.\par
\noindent Printed reading: {\ttfamily\footnotesize 4\allowbreak{}}\par
\noindent Corrected reading: {\ttfamily\footnotesize 14\allowbreak{}}\par
\noindent Source: {\ttfamily\footnotesize 7I:789;\allowbreak{} WITNESS-F7I-IX-202.2\allowbreak{};\allowbreak{} fixes/MISPRINTS\_7I.j\allowbreak{}son\#789;\allowbreak{} fr/sga7i/diplomatic/\allowbreak{}expose\_IX\_body.tex;\allowbreak{} fr/sga7i/current/exp\allowbreak{}ose\_IX\_body.tex\allowbreak{}}\par

\par\medskip\noindent\hypertarget{app-E7I-IX-133.2}{}\textbf{E7I-IX-133.2}\quad printed folio 520.\par
\noindent Printed reading: \emph{(empty)}\par
\noindent Corrected reading: {\ttfamily\footnotesize .4\allowbreak{}}\par
\noindent Source: {\ttfamily\footnotesize 7I:789;\allowbreak{} WITNESS-F7I-IX-202.2\allowbreak{};\allowbreak{} fixes/MISPRINTS\_7I.j\allowbreak{}son\#789;\allowbreak{} fr/sga7i/diplomatic/\allowbreak{}expose\_IX\_body.tex;\allowbreak{} fr/sga7i/current/exp\allowbreak{}ose\_IX\_body.tex\allowbreak{}}\par

\par\medskip\noindent\hypertarget{app-E7I-IX-134.1}{}\textbf{E7I-IX-134.1}\quad printed folio 521.\par
\noindent Printed reading: {\ttfamily\footnotesize schémes\allowbreak{}}\par
\noindent Corrected reading: {\ttfamily\footnotesize schemes\allowbreak{}}\par
\noindent Source: {\ttfamily\footnotesize 7I:790;\allowbreak{} 7I:1022;\allowbreak{} 7I:791;\allowbreak{} 7I:1024;\allowbreak{} FIX-7I-D-0701;\allowbreak{} fixes/FIXLOG\_7I.json\allowbreak{}l\#701;\allowbreak{} fixes/MISPRINTS\_7I.j\allowbreak{}son\#790;\allowbreak{} fixes/MISPRINTS\_7I.j\allowbreak{}son\#1022;\allowbreak{} fixes/MISPRINTS\_7I.j\allowbreak{}son\#791;\allowbreak{} fixes/MISPRINTS\_7I.j\allowbreak{}son\#1024\allowbreak{}}\par

\par\medskip\noindent\hypertarget{app-E7I-IX-134.2}{}\textbf{E7I-IX-134.2}\quad printed folio 521.\par
\noindent Printed reading: {\ttfamily\footnotesize Noth\allowbreak{}}\par
\noindent Corrected reading: {\ttfamily\footnotesize North\allowbreak{}}\par
\noindent Source: {\ttfamily\footnotesize 7I:790;\allowbreak{} 7I:1022;\allowbreak{} 7I:791;\allowbreak{} 7I:1024;\allowbreak{} FIX-7I-D-0701;\allowbreak{} fixes/FIXLOG\_7I.json\allowbreak{}l\#701;\allowbreak{} fixes/MISPRINTS\_7I.j\allowbreak{}son\#790;\allowbreak{} fixes/MISPRINTS\_7I.j\allowbreak{}son\#1022;\allowbreak{} fixes/MISPRINTS\_7I.j\allowbreak{}son\#791;\allowbreak{} fixes/MISPRINTS\_7I.j\allowbreak{}son\#1024\allowbreak{}}\par

\par\medskip\noindent\hypertarget{app-E7I-IX-134.3}{}\textbf{E7I-IX-134.3}\quad printed folio 521.\par
\noindent Printed reading: {\ttfamily\footnotesize Abbildnngen\allowbreak{}}\par
\noindent Corrected reading: {\ttfamily\footnotesize Abbildungen\allowbreak{}}\par
\noindent Source: {\ttfamily\footnotesize 7I:1026;\allowbreak{} FIX-7I-D-0702;\allowbreak{} fixes/FIXLOG\_7I.json\allowbreak{}l\#702;\allowbreak{} fixes/MISPRINTS\_7I.j\allowbreak{}son\#1026\allowbreak{}}\par

\par\medskip\noindent\hypertarget{app-E7I-IX-134.4}{}\textbf{E7I-IX-134.4}\quad printed folio 522.\par
\noindent Printed reading: {\ttfamily\footnotesize  et\allowbreak{}}\par
\noindent Corrected reading: \emph{(empty)}\par
\noindent Source: {\ttfamily\footnotesize 7I:793;\allowbreak{} 7I:1027;\allowbreak{} fixes/MISPRINTS\_7I.j\allowbreak{}son\#793;\allowbreak{} fixes/MISPRINTS\_7I.j\allowbreak{}son\#1027\allowbreak{}}\par

\par\medskip\noindent\hypertarget{app-E7I-IX-134.5}{}\textbf{E7I-IX-134.5}\quad printed folio 523.\par
\noindent Printed reading: {\ttfamily\footnotesize Herman\allowbreak{}}\par
\noindent Corrected reading: {\ttfamily\footnotesize Hermann\allowbreak{}}\par
\noindent Source: {\ttfamily\footnotesize 7I:794;\allowbreak{} 7I:1028;\allowbreak{} fixes/MISPRINTS\_7I.j\allowbreak{}son\#794;\allowbreak{} fixes/MISPRINTS\_7I.j\allowbreak{}son\#1028\allowbreak{}}\par

\par\endgroup\normalsize\fi
\clearpage\SGAtoc{Publisher catalogue}
\selectlanguage{english}
\clearpage
\SGApage{I}{525}



\begingroup
\small
\sloppy
Lecture Notes in Mathematics


Comprehensive leaflet on request

Vol. 111: K. H. Mayer, Relationen zwischen charakteristischen Zahlen. III, 99 Seiten. 1969. DM 16,--

Vol. 112: Colloquium on Methods of Optimization. Edited by N. N. Moiseev. IV, 293 pages. 1970. DM 18,--

Vol. 113: R. Wille, Kongruenzklassengeometrien. III, 99 Seiten. 1970. DM 16,--

Vol. 114: H. Jacquet and R. P. Langlands, Automorphic Forms on GL (2). VII, 548 pages. 1970. DM 24,--

Vol. 115: K. H. Roggenkamp and V. Huber-Dyson, Lattices over Orders I. XIX, 290 pages. 1970. DM 18,--

\SGAvar{Vol. 116: Sèminaire Pierre Lelong (Analyse) Année 1969. IV, 195 pages. 1970. DM 16,--

}{Vol. 116: Séminaire Pierre Lelong (Analyse) Année 1969. IV, 195 pages. 1970. DM 16,--

}{E7I-Catalogue-001}Vol. 117 : Y. Meyer, Nombres de Pisot, Nombres de Salem et Analyse Harmonique. 63 pages. 1970. DM 16,--

Vol. 118: Proceedings of the 15th Scandinavian Congress, Oslo 1968. Edited by K. E. Aubert and W. Ljunggren. IV, 162 pages. 1970. DM 16,--

Vol. 119: M. Raynaud, Faisceaux amples sur les schémas en groupes et les espaces homogènes. III, 219 pages. 1970. DM 16,--

Vol. 120: D. Siefkes, Büchi's Monadic Second Order Successor Arithmetic. XII, 130 Seiten. 1970. DM 16,--

Vol. 121: H. S. Bear, Lectures on Gleason Parts. III, 47 pages. 1970. DM 16,--

Vol. 122: H. Zieschang, E. Vogt und H.-D. Coldewey, Flächen und ebene diskontinuierliche Gruppen. VIII, 203 Seiten. 1970. DM 16,--

Vol. 123: A. V. Jategaonkar, Left Principal Ideal Rings.VI, 145 pages. 1970. DM 16,--

\SGAvar{Vol. 124: Séminare de Probabilités IV. Edited by P. A. Meyer. IV, 282 pages. 1970. DM 20,--

}{Vol. 124: Séminaire de Probabilités IV. Edited by P. A. Meyer. IV, 282 pages. 1970. DM 20,--

}{E7I-Catalogue-002}Vol. 125: Symposium on Automatic Demonstration. V, 310 pages.1970. DM 20,--

Vol. 126: P. Schapira, Théorie des Hyperfonctions. XI,157 pages.1970. DM 16,--

Vol. 127: I. Stewart, Lie Algebras. IV, 97 pages. 1970. DM 16,--

Vol. 128: M. Takesaki, Tomita's Theory of Modular Hilbert Algebras and its Applications. II, 123 pages. 1970. DM 16,--

Vol. 129: K. H. Hofmann, The Duality of Compact Semigroups and C$^*$- Bigebras. XII, 142 pages. 1970. DM 16,--

Vol. 130: F. Lorenz, Quadratische Formen über Körpern. II, 77 Seiten. 1970. DM 16,--

Vol. 131: A Borel et al., Seminar on Algebraic Groups and Related Finite Groups. VII, 321 pages. 1970. DM 22,--

Vol. 132: Symposium on Optimization. III, 348 pages. 1970. DM 22,--

Vol. 133: F. Topsøe, Topology and Measure. XIV, 79 pages. 1970. DM 16,--

Vol. 134: L. Smith, Lectures on the Eilenberg-Moore Spectral Sequence. VII, 142 pages. 1970. DM 16,--

Vol. 135: W. Stoll, Value Distribution of Holomorphic Maps into Compact Complex Manifolds. II, 267 pages. 1970. DM 18,--

\SGAvar{Vol. 136 : M. Karoubi et al., Séminaire Heidelberg-Saarbrücken-Strasbuorg sur la K-Théorie. IV, 264 pages. 1970. DM 18,--

}{Vol. 136 : M. Karoubi et al., Séminaire Heidelberg-Saarbrücken-Strasbourg sur la K-Théorie. IV, 264 pages. 1970. DM 18,--

}{E7I-Catalogue-003}Vol. 137 : Reports of the Midwest Category Seminar IV. Edited by S. MacLane. III, 139 pages. 1970. DM 16,--

Vol. 138: D. Foata et M. Schützenberger, Théorie Géométrique des Polynômes Eulériens. V, 94 pages. 1970. DM 16,--

Vol. 139: A. Badrikian, Séminaire sur les Fonctions Aléatoires Linéaires et les Mesures Cylindriques. VII, 221 pages. 1970. DM 18,--

Vol. 140: Lectures in Modern Analysis and Applications II. Edited by C. T. Taam. VI, 119 pages. 1970. DM 16,--

Vol. 141: G. Jameson, Ordered Linear Spaces. XV, 194 pages. 1970. DM 16,--

Vol. 142: K. W. Roggenkamp, Lattices over Orders II. V, 388 pages. 1970. DM 22,--

Vol. 143: K. W. Gruenberg, Cohomological Topics in Group Theory. XIV, 275 pages. 1970. DM 20,--

Vol. 144: Seminar on Differential Equations and Dynamical Systems, II. Edited by J. A. Yorke. VIII, 268 pages. 1970. DM 20,--

Vol. 145: E. J. Dubuc, Kan Extensions in Enriched Category Theory. XVI, 173 pages. 1970. DM 16,--

Vol. 146: A. B. Altman and S. Kleiman, Introduction to Grothendieck Duality Theory. II, 192 pages. 1970. DM 18,--

Vol. 147: D. E. Dobbs, Cech Cohomological Dimensions for Commutative Rings. VI, 176 pages. 1970. DM 16,--

Vol. 148: R. Azencott, Espaces de Poisson des Groupes Localement Compacts. IX, 141 pages. 1970. DM 16,--

Vol. 149: R. G. Swan and E. G. Evans, K-Theory of Finite Groups and Orders. IV, 237 pages. 1970. DM 20,--

Vol. 150: Heyer, Dualität lokalkompakter Gruppen. XIII, 372 Seiten. 1970. DM 20,--

Vol. 151: M. Demazure et A. Grothendieck, Schémas en Groupes I. (SGA 3). XV, 562 pages. 1970. DM 24,--

Vol. 152: M. Demazure et A. Grothendieck, Schémas en Groupes II. (SGA 3). IX, 654 pages. 1970. DM 24,--

Vol. 153: M. Demazure et A. Grothendieck, Schémas en Groupes III. (SGA 3). VIII, 529 pages. 1970. DM 24,--

Vol. 154: A. Lascoux et M. Berger, Variétés Kähleriennes Compactes. VII, 83 pages. 1970. DM 16,--

Vol. 155: Several Complex Variables I, Maryland 1970. Edited by J. Horváth. IV, 214 pages. 1970. DM 18,--

Vol. 156: R. Hartshorne, Ample Subvarieties of Algebraic Varieties. XIV, 256 pages. 1970. DM 20,--

Vol. 157: T. tom Dieck, K. H. Kamps und D. Puppe, Homotopietheorie. VI, 265 Seiten. 1970. DM 20,--

Vol. 158: T. G. Ostrom, Finite Translation Planes. IV. 112 pages. 1970. DM 16,--

Vol. 159: R. Ansorge und R. Hass. Konvergenz von Differenzenverfahren für lineare und nichtlineare Anfangswertaufgaben. VIII, 145 Seiten. 1970. DM 16,--

\SGAvar{Vol. 160: L. Sucheston, Constributions to Ergodic Theory and Probability. VII, 277 pages. 1970. DM 20,--

}{Vol. 160: L. Sucheston, Contributions to Ergodic Theory and Probability. VII, 277 pages. 1970. DM 20,--

}{E7I-Catalogue-004}Vol. 161: J. Stasheff, H-Spaces from a Homotopy Point of View. VI, 95 pages. 1970. DM 16,--

Vol. 162: Harish-Chandra and van Dijk, Harmonic Analysis on Reductive p-adic Groups. IV, 125 pages. 1970. DM 16,--

\SGAvar{Vol. 163: P. Deligne, Equations Différentielles à Points Singuliers Reguliers. III, 133 pages. 1970. DM 16,--

Vol. 164: J. P. Ferrier, Seminaire sur les Algebres Complètes. II, 69 pages. 1970. DM 16,--

}{Vol. 163: P. Deligne, Equations Différentielles à Points Singuliers Réguliers. III, 133 pages. 1970. DM 16,--

Vol. 164: J. P. Ferrier, Séminaire sur les Algèbres Complètes. II, 69 pages. 1970. DM 16,--

}{E7I-Catalogue-005}Vol. 165: J. M. Cohen, Stable Homotopy. V, 194 pages. 1970. DM 16,--

Vol. 166: A. J. Silberger, PGL$_2$ over the p-adics: its Representations, Spherical Functions, and Fourier Analysis. VII, 202 pages. 1970. DM 18,--

Vol. 167: Lavrentiev, Romanov and Vasiliev, Multidimensional Inverse Problems for Differential Equations. V, 59 pages. 1970. DM 16,--

Vol. 168: F. P. Peterson, The Steenrod Algebra and its Applications: A conference to Celebrate N. E. Steenrod's Sixtieth Birthday. VII, 317 pages. 1970. DM 22,--

Vol. 169: M. Raynaud, Anneaux Locaux Henséliens. V, 129 pages. 1970. DM 16,--

Vol. 170: Lectures in Modern Analysis and Applications III. Edited by C. T. Taam. VI, 213 pages. 1970. DM 18,--

Vol. 171: Set-Valued Mappings, Selections and Topological Properties of $2^X$. Edited by W. M. Fleischman. X, 110 pages. 1970. DM 16,--

Vol. 172: Y.-T. Siu and G. Trautmann, Gap-Sheaves and Extension of Coherent Analytic Subsheaves. V, 172 pages. 1971. DM 16,--

Vol. 173: J. N. Mordeson and B. Vinograde, Structure of Arbitrary Purely Inseparable Extension Fields. IV, 138 pages. 1970. DM 16,--

Vol. 174: B. Iversen, Linear Determinants with Applications to the Picard Scheme of a Family of Algebraic Curves. VI, 69 pages. 1970. DM 16,--

Vol. 175: M. Brelot, On Topologies and Boundaries in Potential Theory. VI, 176 pages. 1971. DM 18,--

Vol. 176: H. Popp, Fundamentalgruppen algebraischer Mannigfaltigkeiten. IV, 154 Seiten. 1970. DM 16,--

Vol. 177: J. Lambek, Torsion Theories, Additive Semantics and Rings of Quotients. VI, 94 pages. 1971. DM 16,--

Please turn over



\noindent\rule{\linewidth}{0.4pt}

\SGApage{I}{526}
Vol. 178: Th. Bröcker und T. tom Dieck, Kobordismentheorie. XVI, 191 Seiten. 1970. DM 18,--

\SGAvar{Vol. 179: Seminaire Bourbaki -- vol. 1968/69. Exposés 347-363. IV. 295 pages. 1971. DM 22,--

}{Vol. 179: Séminaire Bourbaki -- vol. 1968/69. Exposés 347-363. IV. 295 pages. 1971. DM 22,--

}{E7I-Catalogue-006}Vol. 180: Séminaire Bourbaki -- vol. 1969/70. Exposés 364-381. IV, 310 pages. 1971. DM 22,--

Vol. 181: F. DeMeyer and E. Ingraham, Separable Algebras over Commutative Rings. V, 157 pages. 1971. DM 16.--

Vol. 182: L. D. Baumert. Cyclic Difference Sets. VI, 166 pages. 1971. DM 16,--

Vol. 183: Analytic Theory of Differential Equations. Edited by P. F. Hsieh and A. W. J. Stoddart. VI, 225 pages. 1971. DM 20,--

Vol. 184: Symposium on Several Complex Variables, Park City, Utah, 1970. Edited by R. M. Brooks. V, 234 pages. 1971. DM 20,--

Vol. 185: Several Complex Variables II, Maryland 1970. Edited by J. Horváth. III, 287 pages. 1971. DM 24,--

Vol. 186: Recent Trends in Graph Theory. Edited by M. Capobianco/J. B. Frechen/M. Krolik. VI, 219 pages. 1971. DM 18.--

Vol. 187: H. S. Shapiro, Topics in Approximation Theory. VIII, 275 pages. 1971. DM 22,--

Vol. 188: Symposium on Semantics of Algorithmic Languages. Edited by E. Engeler. VI, 372 pages. 1971. DM 26,--

Vol. 189: A. Weil, Dirichlet Series and Automorphic Forms. V, 164 pages. 1971. DM 16,--

Vol. 190: Martingales. A Report on a Meeting at Oberwolfach, May 17-23, 1970. Edited by H. Dinges. V, 75 pages. 1971. DM 16,--

Vol. 191: Séminaire de Probabilités V. Edited by P. A. Meyer. IV, 372 pages. 1971. DM 26,--

Vol. 192: Proceedings of Liverpool Singularities -- Symposium I. Edited by C. T. C. Wall. V, 319 pages. 1971. DM 24,--

Vol. 193: Symposium on the Theory of Numerical Analysis. Edited by J. Ll. Morris. VI, 152 pages. 1971. DM 16,--

Vol. 194: M. Berger, P. Gauduchon et E. Mazet. Le Spectre d'une Variété Riemannienne. VII, 251 pages. 1971. DM 22,--

Vol. 195: Reports of the Midwest Category Seminar V. Edited by J.W. Gray and S. Mac Lane.III, 255 pages. 1971. DM 22,--

Vol. 196: H-spaces -- Neuchâtel (Suisse)- Août 1970. Edited by F. Sigrist, V, 156 pages. 1971. DM 16,--

Vol. 197: Manifolds -- Amsterdam 1970. Edited by N. H. Kuiper. V, 231 pages. 1971. DM 20,--

Vol. 198: M. Hervé, Analytic and Plurisubharmonic Functions in Finite and Infinite Dimensional Spaces. VI, 90 pages. 1971. DM 16.--

Vol. 199: Ch. J. Mozzochi, On the Pointwise Convergence of Fourier Series. VII, 87 pages. 1971. DM 16,--

Vol. 2 00: U. Neri, Singular Integrals. VII, 272 pages. 1971. DM 22,--

Vol. 201 : J. H. van Lint, Coding Theory. VII, 136 pages. 1971. DM 16,--

Vol. 202: J. Benedetto, Harmonic Analysis on Totally Disconnected Sets. VIII, 261 pages. 1971. DM 22,--

Vol. 203: D. Knutson, Algebraic Spaces. VI, 261 pages. 1971. DM 22,--

Vol. 204: A. Zygmund, Intégrales Singulières. IV, 53 pages. 1971. DM 16,--

Vol. 205: Séminaire Pierre Lelong (Analyse) Année 1970. VI, 243 pages. 1971. DM 20,--

Vol. 206: Symposium on Differential Equations and Dynamical Systems. Edited by D. Chillingworth. XI, 173 pages. 1971. DM 16,--

Vol. 207: L. Bernstein, The Jacobi-Perron Algorithm -- Its Theory and Application. IV, 161 pages. 1971. DM 16,--

Vol. 208: A. Grothendieck and J. P. Murre, The Tame Fundamental Group of a Formal Neighbourhood of a Divisor with Normal Crossings on a Scheme. VIII, 133 pages. 1971. DM 16,--

Vol. 209: Proceedings of Liverpool Singularities \quad Symposium II. Edited by C. T. C. Wall. V, 280 pages. 1971. DM 22,--

Vol. 210: M. Eichler, Projective Varieties and Modular Forms. III, 118 pages. 1971. DM 16,--

Vol. 211: Théorie des Matroïdes. Edité par C. P. Bruter. III, 108 pages. 1971. DM 16,--

Vol. 212: B. Scarpellini, Proof Theory and Intuitionistic Systems. VII, 291 pages. 1971. DM 24,--

Vol. 213: H. Hogbe-Nlend, Théorie des Bornologies et Applications. V, 168 pages. 1971. DM 18,--

Vol. 214: M. Smorodinsky, Ergodic Theory, Entropy. V, 64 pages. 1971. DM 16,--

Vol. 215: P. Antonelli, D. Burghelea and P. J. Kahn, The Concordance-Homotopy Groups of Geometric Automorphism Groups. X, 140 pages. 1971. DM 16,--

Vol. 216: H. Maaß, Siegel's Modular Forms and Dirichlet Series. VII, 328 pages. 1971. DM 20,--

Vol. 217: T. J. Jech, Lectures in Set Theory with Particular Emphasis on the Method of Forcing. V, 137 pages. 1971. DM 16,--

Vol. 218: C. P. Schnorr, Zufälligkeit und Wahrscheinlichkeit. IV, 212.Seiten 1971. DM 20,--

Vol. 219: N. L. Alling and N. Greenleaf, Foundations of the Theory of Klein Surfaces. IX, 117 pages. 1971. DM 16,--

Vol. 220: W. A. Coppel, Disconjugacy. V, 148 pages. 1971. DM 16,--

Vol. 221: P. Gabriel und F. Ulmer, Lokal präsentierbare Kategorien. V, 200 Seiten. 1971. DM 18,--

Vol. 222: C. Meghea, Compactification des Espaces Harmoniques. III, 108 pages. 1971. DM 16,--

Vol. 223: U. Felgner, Models of ZF-Set Theory. VI, 173 pages. 1971. DM 16,--

\SGAvar{Vol. 224: Revètements Etales et Groupe Fondamental. (SGA 1). Dirigé par A. Grothendieck XXII, 447 pages. 1971. DM 30,--

}{Vol. 224: Revêtements Etales et Groupe Fondamental. (SGA 1). Dirigé par A. Grothendieck XXII, 447 pages. 1971. DM 30,--

}{E7I-Catalogue-007}Vol. 225: Théorie des Intersections et Théorème de Riemann-Roch. (SGA 6). Dirigé par P. Berthelot, A. Grothendieck et L. Illusie. XII, 700 pages. 1971. DM 40,--

Vol. 226: Seminar on Potential Theory, II. Edited by H. Bauer. IV, 170 pages. 1971. DM 18,--

Vol. 227: H. L. Montgomery, Topics in Multiplicative Number Theory. IX, 178 pages. 1971. DM 18,--

Vol. 228: Conference on Applications of Numerical Analysis. Edited by J. Ll. Morris. X, 358 pages. 1971. DM 26,--

Vol. 229: J. Väisälä, Lectures on n-Dimensional Quasiconformal Mappings. XIV, 144 pages. 1971. DM 16,--

Vol. 230: L. Waelbroeck, Topological Vector Spaces and Algebras. VII, 158 pages. 1971. DM 16,--

Vol. 231: H. Reiter, L$^1$-Algebras and Segal Algebras. XI, 113 pages. 1971. DM 16,--

Vol. 232: T. H. Ganelius, Tauberian Remainder Theorems. VI, 75 pages. 1971. DM 16,--

Vol. 233: C. P. Tsokos and W. J. Padgett. Random Integral Equations with Applications to Stochastic Systems. VII, 174 pages. 1971. DM 18,--

Vol. 234: A. Andreotti and W. Stoll. Analytic and Algebraic Dependence of Meromorphic Functions. III, 390 pages. 1971. DM 26,--

Vol. 235: Global Differentiable Dynamics. Edited by O. Hájek, A. J. Lohwater, and R. McCann. X, 140 pages. 1971. DM 16,--

Vol. 236 : M. Barr, P. A. Grillet, and D. H. van Osdol. Exact Categories and Categories of Sheaves. VII, 239 pages. 1971, DM 20,--

Vol. 237: B. Stenström. Rings and Modules of Quotients. VII, 136 pages. 1971. DM 16,--

Vol. 238: Der kanonische Modul eines Cohen-Macaulay-Rings. Herausgegeben von Jürgen Herzog und Ernst Kunz. VI, 103 Seiten. 1971. DM 16,--

Vol. 239: L. Illusie, Complexe Cotangent et Déformations I. XV, 355 pages. 1971. DM 26,--

Vol. 240: A. Kerber, Representations of Permutation Groups I. VII, 192 pages. 1971. DM 18,--

Vol. 241: S. Kaneyuki, Homogeneous Bounded Domains and Siegel Domains. V, 89 pages. 1971. DM 16,--

Vol. 242: R. R. Coifman et G. Weiss, Analyse Harmonique Non-Commutative sur Certains Espaces. V, 160 pages. 1971. DM 16,--

Vol. 243: Japan-United States Seminar on Ordinary Differential and Functional Equations. Edited by M. Urabe. VIII, 332 pages. 1971. DM 26,--

Vol. 244: Séminaire Bourbaki -- vol. 1970/71. Exposés 382--399. IV, 356 pages. 1971. DM 26,--

Vol. 245: D. E. Cohen, Groups of Cohomological Dimension One. V, 99 pages. 1972. DM 16,--

\noindent\rule{\linewidth}{0.4pt}



\noindent\rule{\linewidth}{0.4pt}

\SGApage{I}{527}
Vol. 246: Lectures on Rings and Modules. Tulane University Ring and Operator Theory Year, 1970--1971. Volume I. X, 661 pages. 1972. DM 40,--

Vol. 247: Lectures on Operator Algebras. Tulane University Ring and Operator Theory Year, 1970--1971. Volume II. XI, 786 pages. 1972. DM 40,--

Vol. 248: Lectures on the Applications of Sheaves to Ring Theory. Tulane University Ring and Operator Theory Year, 1970--1971. Volume III. VIII, 315 pages. 1971. DM 26,--

Vol. 249: Symposium on Algebraic Topology. Edited by P. J. Hilton. VII, 111 pages. 1971. DM 16,--

Vol. 250: B. Jónsson, Topics in Universal Algebra. VI, 220 pages. 1972. DM 20,--

Vol. 251: The Theory of Arithmetic Functions. Edited by A. A. Gioia and D. L. Goldsmith VI, 287 pages. 1972. DM 24,--

Vol. 252: D. A. Stone, Stratified Polyhedra. IX, 193 pages. 1972. DM 18,--

Vol. 253: V. Komkov, Optimal Control Theory for the Damping of Vibrations of Simple Elastic Systems. V, 240 pages. 1972. DM 20,--

Vol. 254: C. U. Jensen, Les Foncteurs Dérivés de $\varprojlim$ et leurs Applications en Théorie des Modules. V, 103 pages. 1972. DM 16,--

Vol. 255: Conference in Mathematical Logic -- London '70. Edited by W. Hodges. VIII, 351 pages. 1972. DM 26,--

Vol. 256: C. A. Berenstein and M. A. Dostal, Analytically Uniform Spaces and their Applications to Convolution Equations. VII, 130 pages. 1972. DM 16,--

Vol. 257: R. B. Holmes, A Course on Optimization and Best Approximation. VIII, 233 pages. 1972. DM 20,--

Vol. 258: Séminaire de Probabilités VI. Edited by P. A. Meyer. VI, 253 pages. 1972. DM 22,--

Vol. 259: N. Moulis, Structures de Fredholm sur les Variétés Hilbertiennes. V, 123 pages. 1972. DM 16,--

Vol. 260: R. Godement and H. Jacquet, Zeta Functions of Simple Algebras. IX, 188 pages. 1972. DM 18,--

Vol. 261: A. Guichardet, Symmetric Hilbert Spaces and Related Topics. V, 197 pages. 1972. DM 18,--

Vol. 262: H. G. Zimmer, Computational Problems, Methods, and Results in Algebraic Number Theory. V, 103 pages. 1972. DM 16,--

Vol. 263: T. Parthasarathy, Selection Theorems and their Applications. VII, 101 pages. 1972. DM 16,--

Vol. 264: W. Messing, The Crystals Associated to Barsotti-Tate Groups: with Applications to Abelian Schemes. III, 190 pages. 1972. DM 18,--

Vol. 265: N. Saavedra Rivano, Catégories Tannakiennes. II, 418 pages. 1972. DM 26,--

Vol. 266: Conference on Harmonic Analysis. Edited by D. Gulick and R. L. Lipsman. VI, 323 pages. 1972. DM 24,--

Vol. 267: Numerische Lösung nichtlinearer partieller Differential- und Integro-Differentialgleichungen. Herausgegeben von R. Ansorge und W. Törnig, VI, 339 Seiten. 1972. DM 26,--

Vol. 268: C. G. Simader, On Dirichlet's Boundary Value Problem. IV, 238 pages. 1972. DM 20,--

Vol. 269: Théorie des Topos et Cohomologie Etale des Schémas. (SGA 4). Dirigé par M. Artin, A. Grothendieck et J. L. Verdier. XIX, 525 pages. 1972. DM 50,--

Vol. 270: Théorie des Topos et Cohomologie Etale des Schémas. Tome 2. (SGA 4). Dirigé par M. Artin, A. Grothendieck et J. L. Verdier. V, 314 pages. 1972. DM 50,--

Vol. 271: J. P. May, The Geometry of Iterated Loop Spaces. IX, 175 pages. 1972. DM 18,--

Vol. 272: K. R. Parthasarathy and K. Schmidt, Positive Definite Kernels, Continuous Tensor Products, and Central Limit Theorems of Probability Theory. VI, 107 pages. 1972. DM 16,--

Vol. 273: U. Seip, Kompakt erzeugte Vektorräume und Analysis. IX, 119 Seiten. 1972. DM 16,--

\SGAvar{Vol. 274: Toposes, Algebraic Geometry and Logic. Edited by. F. W. Lawvere. VI, 189 pages. 1972. DM 18,--

}{Vol. 274: Toposes, Algebraic Geometry and Logic. Edited by F. W. Lawvere. VI, 189 pages. 1972. DM 18,--

}{E7I-Catalogue-008}Vol. 275: Séminaire Pierre Lelong (Analyse) Année 1970--1971. VI, 181 pages. 1972. DM 18,--

Vol. 276: A. Borel, Représentations de Groupes Localement Compacts. V, 98 pages. 1972. DM 16,--

Vol. 277: Séminaire Banach. Edité par C. Houzel. VII, 229 pages. 1972. DM 20,--

Vol. 278: H. Jacquet, Automorphic Forms on GL(2). Part II. XIII, 142 pages. 1972. DM 16,--

Vol. 279: R. Bott, S. Gitler and I. M. James, Lectures on Algebraic and Differential Topology. V, 174 pages. 1972. DM 18,--

Vol. 280: Conference on the Theory of Ordinary and Partial Differential Equations. Edited by W. N. Everitt and B. D. Sleeman. XV, 367 pages. 1972. DM 26,--

Vol. 281: Coherence in Categories. Edited by S. Mac Lane. VII, 235 pages. 1972. DM 20,--

Vol. 282: W. Klingenberg und P. Flaschel, Riemannsche Hilbertmannigfaltigkeiten. Periodische Geodätische. VII, 211 Seiten. 1972. DM 20,--

Vol. 283: L. Illusie, Complexe Cotangent et Déformations II. VII, 304 pages. 1972. DM 24,--

Vol. 284: P. A. Meyer, Martingales and Stochastic Integrals I. VI, 89 pages. 1972. DM 16,--

Vol. 285: P. de la Harpe, Classical Banach-Lie Algebras and Banach-Lie Groups of Operators in Hilbert Space. III, 160 pages. 1972. DM 16,--

Vol. 286: S. Murakami, On Automorphisms of Siegel Domains. V, 95 pages. 1972. DM 16,--

Vol. 288: Groupes de Monodromie en Géométrie Algébrique. (SGA 7 I). Dirigé par A. Grothendieck. IX, 523 pages. 1972. DM 50,--

\noindent\rule{\linewidth}{0.4pt}



\SGApage{I}{528}
ISBN 3-540-05987-3

ISBN 0-387-05987-3
\endgroup



\markboth{Publisher catalogue}{Publisher catalogue}

\ifnum\SGAform=1
\clearpage\section*{Critical apparatus}\begingroup\small\raggedright\sloppy
Readings are given in the exact TeX notation of the two English witnesses. Decimal suffixes identify the individual differences in each marked passage.

\par\medskip\noindent\hypertarget{app-E7I-Catalogue-001.1}{}\textbf{E7I-Catalogue-001.1}\quad printed folio 525.\par
\noindent Printed reading: {\ttfamily\footnotesize Sèminaire\allowbreak{}}\par
\noindent Corrected reading: {\ttfamily\footnotesize Séminaire\allowbreak{}}\par
\noindent Source: {\ttfamily\footnotesize port/PORT\_EN7I\_LOG.j\allowbreak{}sonl\#C1000;\allowbreak{} fixes/FIXLOG\_7I\_CURR\allowbreak{}ENT.jsonl\#C1000;\allowbreak{} 7I:796\allowbreak{}}\par

\par\medskip\noindent\hypertarget{app-E7I-Catalogue-002.1}{}\textbf{E7I-Catalogue-002.1}\quad printed folio 525.\par
\noindent Printed reading: {\ttfamily\footnotesize Séminare\allowbreak{}}\par
\noindent Corrected reading: {\ttfamily\footnotesize Séminaire\allowbreak{}}\par
\noindent Source: {\ttfamily\footnotesize port/PORT\_EN7I\_LOG.j\allowbreak{}sonl\#C1001;\allowbreak{} fixes/FIXLOG\_7I\_CURR\allowbreak{}ENT.jsonl\#C1001;\allowbreak{} 7I:797\allowbreak{}}\par

\par\medskip\noindent\hypertarget{app-E7I-Catalogue-003.1}{}\textbf{E7I-Catalogue-003.1}\quad printed folio 525.\par
\noindent Printed reading: {\ttfamily\footnotesize Strasbuorg\allowbreak{}}\par
\noindent Corrected reading: {\ttfamily\footnotesize Strasbourg\allowbreak{}}\par
\noindent Source: {\ttfamily\footnotesize port/PORT\_EN7I\_LOG.j\allowbreak{}sonl\#C1002;\allowbreak{} fixes/FIXLOG\_7I\_CURR\allowbreak{}ENT.jsonl\#C1002;\allowbreak{} 7I:798\allowbreak{}}\par

\par\medskip\noindent\hypertarget{app-E7I-Catalogue-004.1}{}\textbf{E7I-Catalogue-004.1}\quad printed folio 525.\par
\noindent Printed reading: {\ttfamily\footnotesize Constributions\allowbreak{}}\par
\noindent Corrected reading: {\ttfamily\footnotesize Contributions\allowbreak{}}\par
\noindent Source: {\ttfamily\footnotesize port/PORT\_EN7I\_LOG.j\allowbreak{}sonl\#C1003;\allowbreak{} fixes/FIXLOG\_7I\_CURR\allowbreak{}ENT.jsonl\#C1003;\allowbreak{} 7I:799\allowbreak{}}\par

\par\medskip\noindent\hypertarget{app-E7I-Catalogue-005.1}{}\textbf{E7I-Catalogue-005.1}\quad printed folio 525.\par
\noindent Printed reading: {\ttfamily\footnotesize Reguliers\allowbreak{}}\par
\noindent Corrected reading: {\ttfamily\footnotesize Réguliers\allowbreak{}}\par
\noindent Source: {\ttfamily\footnotesize port/PORT\_EN7I\_LOG.j\allowbreak{}sonl\#C1004;\allowbreak{} fixes/FIXLOG\_7I\_CURR\allowbreak{}ENT.jsonl\#C1004;\allowbreak{} 7I:800\allowbreak{}}\par

\par\medskip\noindent\hypertarget{app-E7I-Catalogue-005.2}{}\textbf{E7I-Catalogue-005.2}\quad printed folio 525.\par
\noindent Printed reading: {\ttfamily\footnotesize Seminaire\allowbreak{}}\par
\noindent Corrected reading: {\ttfamily\footnotesize Séminaire\allowbreak{}}\par
\noindent Source: {\ttfamily\footnotesize port/PORT\_EN7I\_LOG.j\allowbreak{}sonl\#C1005;\allowbreak{} fixes/FIXLOG\_7I\_CURR\allowbreak{}ENT.jsonl\#C1005;\allowbreak{} 7I:801\allowbreak{}}\par

\par\medskip\noindent\hypertarget{app-E7I-Catalogue-005.3}{}\textbf{E7I-Catalogue-005.3}\quad printed folio 525.\par
\noindent Printed reading: {\ttfamily\footnotesize Algebres\allowbreak{}}\par
\noindent Corrected reading: {\ttfamily\footnotesize Algèbres\allowbreak{}}\par
\noindent Source: {\ttfamily\footnotesize port/PORT\_EN7I\_LOG.j\allowbreak{}sonl\#C1005;\allowbreak{} fixes/FIXLOG\_7I\_CURR\allowbreak{}ENT.jsonl\#C1005;\allowbreak{} 7I:801\allowbreak{}}\par

\par\medskip\noindent\hypertarget{app-E7I-Catalogue-006.1}{}\textbf{E7I-Catalogue-006.1}\quad printed folio 526.\par
\noindent Printed reading: {\ttfamily\footnotesize Seminaire\allowbreak{}}\par
\noindent Corrected reading: {\ttfamily\footnotesize Séminaire\allowbreak{}}\par
\noindent Source: {\ttfamily\footnotesize port/PORT\_EN7I\_LOG.j\allowbreak{}sonl\#C1012;\allowbreak{} fixes/FIXLOG\_7I\_CURR\allowbreak{}ENT.jsonl\#C1012;\allowbreak{} 7I:1040\allowbreak{}}\par

\par\medskip\noindent\hypertarget{app-E7I-Catalogue-007.1}{}\textbf{E7I-Catalogue-007.1}\quad printed folio 526.\par
\noindent Printed reading: {\ttfamily\footnotesize Revètements\allowbreak{}}\par
\noindent Corrected reading: {\ttfamily\footnotesize Revêtements\allowbreak{}}\par
\noindent Source: {\ttfamily\footnotesize port/PORT\_EN7I\_LOG.j\allowbreak{}sonl\#C1013;\allowbreak{} fixes/FIXLOG\_7I\_CURR\allowbreak{}ENT.jsonl\#C1013;\allowbreak{} 7I:1044\allowbreak{}}\par

\par\medskip\noindent\hypertarget{app-E7I-Catalogue-008.1}{}\textbf{E7I-Catalogue-008.1}\quad printed folio 527.\par
\noindent Printed reading: {\ttfamily\footnotesize .\allowbreak{}}\par
\noindent Corrected reading: \emph{(empty)}\par
\noindent Source: {\ttfamily\footnotesize port/PORT\_EN7I\_LOG.j\allowbreak{}sonl\#C1014;\allowbreak{} fixes/FIXLOG\_7I\_CURR\allowbreak{}ENT.jsonl\#C1014;\allowbreak{} 7I:1045\allowbreak{}}\par

\par\endgroup\normalsize\fi
\part{SGA 7 II}
\clearpage\SGAtoc{Exposé X}
\selectlanguage{english}
\SGApage{II}{1}

\begin{center}
\underline{\underline{INTERSECTIONS ON REGULAR SURFACES}}

\medskip
\underline{by P. Deligne}
\end{center}

In the first part of this exposé, after some preliminaries from intersection theory
(contained in $[5]\cup[7]$), we give Artin's proof of a special case of a theorem of
Raynaud [9] (1.13), and an a priori proof of a theorem which Néron [8] had verified
case by case (1.15).  The theorems of [9] used in [2], and hence in IX, follow fairly
easily from (1.13).

In the second part, continuing in the same vein, I explain how the ``Italians'' [1],
[3] calculated intersection numbers on surfaces.  The main results, 2.9 and 2.11,
may also be found in Northcott [12] and [13].  To conclude, in 2.27, 2.30, and 2.34
I give a conjectural relation between the dimension of the versal formal moduli
scheme of a singularity of a reduced curve, the invariant $\delta$ of the singularity,
and an invariant measuring its symmetry.

This exposé will not be used in the remainder of the seminar.

\medskip
1. \underline{Intersection numbers on arithmetic surfaces}

1.0. Throughout this section, let $X$ be a regular scheme of dimension 2, let $Y$ be
a reduced divisor on $X$, and let $k$ be a subfield of
$H^0(Y,\underline{O}_Y)$.  We assume that $Y$ is a curve (EGA II 7.4.1) proper over
$k$.  Here are three examples of this situation.

\underline{Example} 1.1. Let $(S,\eta,s)$ be a trait and let
$f:X\longrightarrow S$ be a regular, proper, and flat $S$-scheme whose fibres are
purely one-dimensional.  We take $Y=(f^{-1}(s))_{\mathrm{red}}$ and $k=k(s)$.  From
1.10 onward we shall restrict ourselves to this case.


\SGApage{II}{2}
\underline{Example} 1.2. Let $X_o$ be the spectrum of a normal local ring of
dimension 2, let $s$ be the closed point of $X_o$, and let
$f:X\longrightarrow X_o$ be a regular connected $X_o$-scheme, proper over $X_o$,
distinct from $X_o$ but agreeing with $X_o$ over $X_o\setminus\{s\}$.  We take
$Y=(f^{-1}(s))_{\mathrm{red}}$ and $k=k(s)$.

\underline{Example} 1.3. We take $X$ to be a complete regular surface over a field
$k$, and $Y$ to be a reduced curve on $X$.

Recall the following lemma (SGA 6 IV, Ex.~1.9).

\SGAvar{\underline{Lemma} 1.4. \underline{Let} $Z_o$ \underline{be a closed subscheme of a
Noetherian scheme} $Z$. \underline{Then the inclusion of} $Z_o$ \underline{in} $Z$
\underline{induces an isomorphism from the Grothendieck group of the category of
coherent} $O_{\mathbb Z_o}$\underline{-modules to the Grothendieck group of the
category of coherent sheaves on} $Z$ \underline{with support in} $Z_o$.

}{\underline{Lemma} 1.4. \underline{Let} $Z_o$ \underline{be a closed subscheme of a
Noetherian scheme} $Z$. \underline{Then the inclusion of} $Z_o$ \underline{in} $Z$
\underline{induces an isomorphism from the Grothendieck group of the category of
coherent} $\underline{O}_{Z_o}$\underline{-modules to the Grothendieck group of the
category of coherent sheaves on} $Z$ \underline{with support in} $Z_o$.

}{E7II-X-001}This lemma allows us to define the Euler--Poincaré characteristic $\chi(\mathcal F)$
of a coherent sheaf $\mathcal F$ on $X$ supported on $Y$ so that:

(i) If $\mathcal F$ is a coherent sheaf of $\underline{O}_Y$-modules, then
\begin{equation}\tag{1.4.1}
\chi(\mathcal F)=\dim_k H^0(\mathcal F)-\dim_k H^1(\mathcal F).
\end{equation}

(ii) $\chi$ is additive with respect to short exact sequences.

If the cohomology sheaves of a complex $K$ of $\underline{O}_X$-modules are coherent
sheaves supported on $Y$, we set
\begin{equation}\tag{1.4.2}
\chi(K)=\sum_i(-1)^i\chi(\underline{H}^i(K)).
\end{equation}
If $\mathcal F$ is supported at a point $y$ of $Y$, then
\begin{equation}\tag{1.4.3}
\chi(\mathcal F)=\operatorname{length}_{\underline{O}_y}(\mathcal F_y)[k(y):k].
\end{equation}


\SGApage{II}{\(\approx\)3}
Under the hypotheses of (1.1), (1.2), or (1.3), one has
\begin{equation}\tag{1.4.4}
\chi(\mathcal F)=\operatorname{length}f_*\mathcal F-
\operatorname{length}R^1f_*\mathcal F.
\end{equation}
Let $D$ be a divisor $\geq0$ supported on $Y$, and let $\mathcal L$ be an invertible
sheaf on $D$.  There is then a unique way to define the degree
$\deg_D(\mathcal L)$ such that
\begin{equation}\tag{1.4.5}
\left\{
\begin{array}{l}
\deg_D(\mathcal L'\otimes\mathcal L'')=
\deg_D(\mathcal L')+\deg_D(\mathcal L''),\\[3ex]
\text{for }\mathcal L\subset\underline O,
\text{ one has }\deg_D(\mathcal L)=-\chi(\underline O/\mathcal L).
\end{array}
\right.
\end{equation}
The Riemann--Roch formula on $D$ is then
\begin{equation}\tag{1.4.6}
\chi(\mathcal L)=\chi(\underline{O}_D)+\deg_D(\mathcal L).
\end{equation}
If $D=\sum n_iD_i$ is the decomposition of $D$ into reduced irreducible divisors,
then
\begin{equation}\tag{1.4.7}
\deg_D(\mathcal L)=\sum n_i\deg_{D_i}(\mathcal L).
\end{equation}

\underline{Definition} 1.5. \underline{Let} $D$ \underline{and} $E$
\underline{be two divisors} $\geq0$ \underline{on} $X$, \underline{and suppose that}
$D$ \underline{is supported on} $Y$. \underline{The intersection number of} $D$
\underline{and} $E$ \underline{(relative to the field} $k$ (1.0)) \underline{is
defined by}
\begin{equation*}
(D,E)=\chi(\underline{O}_D\overset{\mathbb L}{\otimes}\underline{O}_E)
=\chi(\underline{O}_D\otimes\underline{O}_E)
-\chi(\operatorname{Tor}^1(\underline{O}_D,\underline{O}_E)).
\end{equation*}
Indeed, under the hypotheses made here,
$\operatorname{Tor}^2(\underline{O}_D,\underline{O}_E)=0$.

\underline{Proposition} 1.6. \underline{Under the hypotheses of} 1.5,
\underline{one has}:

(i) \quad \underline{The intersection number} $(D,E)$ \underline{is additive in the
divisors} $D$ \underline{and} $E$; \underline{this allows us, by linearity, to define
the intersection number of two divisors which are not necessarily positive.}


\SGApage{II}{\(\approx\)4}
(ii) \quad $(D,E)=\deg_D(\underline O(E))$.

(iii) \quad \underline{If} $D$ \underline{and} $E$ \underline{meet at only finitely
many points, then} $(D,E)=\chi(\underline{O}_D\otimes\underline{O}_E)\geq0$,
\underline{and equality can occur only if} $D$ \underline{and} $E$
\underline{are disjoint.}

(iv) \quad \underline{If} $D$ \underline{and} $E$ \underline{are supported on} $Y$,
\underline{then} $(D,E)=(E,D)$.

(v) \quad \underline{If} $N$ \underline{is the conormal sheaf of} $D$,
\underline{then}
\begin{equation*}
D^2=(D,D)=-\deg_D(N).
\end{equation*}

(vi) \quad \underline{If} $E$, \underline{not necessarily positive, is the divisor
of a meromorphic function on} $X$, \underline{then} $(D,E)=0$.

\underline{Proof}. The exact sequence
\begin{equation*}
0\longrightarrow\underline O(-E)\longrightarrow\underline O_X
\longrightarrow\underline O_E\longrightarrow0
\end{equation*}
gives, by (1.4.6),
\begin{equation*}
\begin{array}{rl}
(D,E)\overset{\mathrm{def}}{=}&
\chi(\underline{O}_D\overset{\mathbb L}{\otimes}\underline{O}_E)\\
=&\chi(\underline{O}_D\overset{\mathbb L}{\otimes}\underline{O}_X)
-\chi(\underline{O}_D\overset{\mathbb L}{\otimes}\underline O(-E))\\[1ex]
=&\chi(\underline{O}_D)-\chi(\underline O(-E)|D)
=-\deg_D(\underline O(-E))=\deg_D(\underline O(E)).
\end{array}
\end{equation*}
This proves (ii), (vi), and biadditivity, by (1.4.5) and (1.4.7).  Under the
hypotheses of (iii), one has
$\operatorname{Tor}^1(\underline{O}_D,\underline{O}_E)=0$, which gives the
conclusion.  Under the hypotheses of (v), one has
\begin{equation*}
\underline{O}_D\otimes\underline{O}_D=\underline{O}_D
\quad\text{and}\quad
\operatorname{Tor}^1(\underline{O}_D,\underline{O}_D)\simeq N
\quad(\text{SGA 6 VII 2.5}),
\end{equation*}
hence
\begin{equation*}
(D,D)=\chi(\underline{O}_D)-\chi(N)=-\deg_D(N).
\end{equation*}
Finally, (iv) follows from the symmetry of the tensor product.

Notice that (ii), (iv), and (vi) remain true without assuming the divisors positive.


\SGApage{II}{5}
\underline{Proposition} 1.7. \underline{Let} $(Y_i)_{1\leq i\leq k}$
\underline{be the irreducible components of} $Y$, $f$ \underline{a meromorphic
function on} $X$, \underline{and}
\begin{equation*}
\operatorname{div}(f)=\sum_i n_iY_i+Z
\end{equation*}
\underline{its divisor, with} $|Z|\cap Y$ \underline{finite. Then, putting}
$D_i=n_iY_i$, \underline{one has, for} $(a_i)\in\mathbb Q^k$:
\begin{equation*}
\left(\sum_i a_iD_i\right)^2
=-\sum_{i<j}(a_i-a_j)^2(D_i,D_j)-\sum_i a_i^2(D_i,Z).
\end{equation*}

\SGAvar{\underline{Proof} (after [7]). By 1.6 (v), (vi),
\begin{equation*}
\begin{aligned}
\left(\sum_i a_iD_i\right)^2
&=\sum_i a_i\left(D_i,\sum_j a_jD_j\right)
=\sum_i a_i\left(D_i,\sum_j a_jD_j-a_i\operatorname{div}f\right)\\
&=\sum_i a_i\left(D_i,\sum_{j\ne i}(a_j-a_i)D_j-a_iZ\right)\\
&=\sum_{i\ne j}a_i(a_j-a_i)(D_i,D_j)-\sum_i a_i^2(D_i,Z)\\
&=-\sum_{i<j}(a_i-a_j)^2(D_iD_j)-\sum_i a_i^2(D_i,Z).
\end{aligned}
\end{equation*}

}{\underline{Proof} (after [7]). By 1.6 (v), (vi),
\begin{equation*}
\begin{aligned}
\left(\sum_i a_iD_i\right)^2
&=\sum_i a_i\left(D_i,\sum_j a_jD_j\right)
=\sum_i a_i\left(D_i,\sum_j a_jD_j-a_i\operatorname{div}f\right)\\
&=\sum_i a_i\left(D_i,\sum_{j\ne i}(a_j-a_i)D_j-a_iZ\right)\\
&=\sum_{i\ne j}a_i(a_j-a_i)(D_i,D_j)-\sum_i a_i^2(D_i,Z)\\
&=-\sum_{i<j}(a_i-a_j)^2(D_i,D_j)-\sum_i a_i^2(D_i,Z).
\end{aligned}
\end{equation*}

}{E7II-X-002}\underline{Corollary} 1.8. \underline{Under the hypotheses of Example} 1.1,
\underline{take for} $f$ \underline{a uniformizer} $\pi$ \underline{of} $S$.
\underline{Here} $Z=0$ \underline{and}
\begin{equation*}
\left(\sum_i k_iY_i\right)^2
=-\sum_{i<j}\frac{1}{n_in_j}(k_in_j-k_jn_i)^2(Y_i,Y_j);
\end{equation*}
\underline{thus the intersection form is negative, and if} $Y$
\underline{is connected, its kernel consists precisely of the multiples of}
$\operatorname{div}(\pi)$.

\underline{Corollary} 1.9 (Mumford [7]). \underline{Under the hypotheses of Example}
1.2, \underline{take for} $f$ \underline{a nonzero element of the maximal ideal of
the local ring} $\underline O_s$. \underline{Here} $Z>0$, \underline{so the
intersection form is negative definite.}



\SGApage{II}{6}
1.10. \underline{For the remainder of this section we restrict ourselves to the
case considered in Example} 1.1.

\SGAvar{Lichtenbaum [5] proved that, in this case, $X$ is automatically projective over
$S$, so that there is no difficulty in defining the complex
$Rf^{!}\underline{O}_S$ ([4] III 8.7, p.~190).  Since $f$ has relative dimension
1 and $X$ and $S$ are regular (Gorenstein would suffice), there exists ([4] V 8.3
and V 9.1) an invertible sheaf $\omega_{X/S}$ on $S$, unique up to unique
isomorphism, such that

}{Lichtenbaum [5] proved that, in this case, $X$ is automatically projective over
$S$, so that there is no difficulty in defining the complex
$Rf^{!}\underline{O}_S$ ([4] III 8.7, p.~190).  Since $f$ has relative dimension
1 and $X$ and $S$ are regular (Gorenstein would suffice), there exists ([4] V 8.3
and V 9.1) an invertible sheaf $\omega_{X/S}$ on $X$, unique up to unique
isomorphism, such that

}{E7II-X-003}\begin{equation}\tag{1.10.1}
Rf^{!}\underline{O}_S \;\simeq\; \omega_{X/S}[1] \quad .
\end{equation}
We denote by $K$ any divisor (effective or not) such that
$\omega_{X/S}\simeq\underline{O}(K)$.

Let $D$ be a divisor $\geq0$ supported on $Y$, and let $N$ be its conormal sheaf:

\begin{center}
\begin{tikzcd}[column sep=large, row sep=large]
D \arrow[rr, "i"] \arrow[dr, swap, "g"] \& {} \& X \arrow[dl, "f"] \\
{} \& S \& {}
\end{tikzcd}
\end{center}
One has

\begin{equation}\tag{1.10.2}
Rg^{!}\underline{O}_S
=Ri^{!}Rf^{!}\underline{O}_S
=Ri^{!}\omega_{X/S}[1]
=\omega_{X/S}\big|_D\otimes N^{\vee}[0].
\end{equation}
By the duality theorem ([4] V 11.1, p.~210),

\begin{equation}\tag{1.10.3}
Rg_*Rg^{!}\underline{O}_S
\;\simeq\;
\operatorname{RHom}(Rg_*\underline{O}_D,\underline{O}_S) \quad .
\end{equation}
Moreover, for every $\underline{O}_S$-module $M$ supported at the closed point,

\begin{equation}\tag{1.10.4}
\chi\bigl(\operatorname{RHom}(M,\underline{O}_S)\bigr)=-\chi(M),
\end{equation}
so that here, using (1.4.4) and (1.10.3), we obtain

\begin{equation*}
\chi(Rg^{!}\underline{O}_S)=-\chi(\underline{O}_D) \quad .
\end{equation*}
Since, by (1.4.6),

\begin{equation*}
\chi(Rg^{!}\underline{O}_S)
=\chi(\underline{O}_D)+\deg_D(Rg^{!}\underline{O}_S),
\end{equation*}


\SGApage{II}{\(\approx\)7}
we conclude from (1.10.2) that

\SGAvar{\begin{equation*}
-2\chi(\underline{O}_D)
=\deg_D(Rg^{!}\underline{O}_D)
=\deg_D(\omega_{X/S}|_D\otimes N^{\vee})
=\deg_D(\omega_{X/S})-\deg_D(N),
\end{equation*}
which, by 1.6(ii) and (v), may be rewritten as follows.

}{\begin{equation*}
-2\chi(\underline{O}_D)
=\deg_D(Rg^{!}\underline{O}_S)
=\deg_D(\omega_{X/S}|_D\otimes N^{\vee})
=\deg_D(\omega_{X/S})-\deg_D(N),
\end{equation*}
which, by 1.6(ii) and (v), may be rewritten as follows.

}{E7II-X-004}\underline{Proposition} 1.11. \underline{Under the hypotheses of} (1.1),
\underline{for every divisor} $D\geq0$ \underline{supported on} $Y$,
\underline{one has}

\begin{equation*}
\chi(\underline{O}_D)=-\frac12(D,D+K) \quad .
\end{equation*}

1.12. We denote by $d$ the greatest common divisor of the geometric
multiplicities of the irreducible components of $X_s=f^{-1}(s)$; that is, the
greatest common divisor of the lengths of the local rings at the generic points
of the geometric special fibre of $f$.  If $X_\eta$ has a rational point, or
equivalently if $f:X\longrightarrow S$ has a section $x$, then $d=1$; more
precisely, $f$ is smooth at $x(s)$.  Indeed, $x$ is a closed immersion of regular
schemes, hence a regular immersion, and EGA IV 17.12.1 applies.

Here is Artin's proof of a special case of a theorem of Raynaud [9], which only
assumes that $d$ is prime to the characteristic $p$ of $k$.

\underline{Theorem} 1.13. \underline{With the preceding notation, if} $Y$
\underline{is geometrically connected and} $d=1$, \underline{then} $X$
\underline{is cohomologically flat over} $S$ \underline{in dimension} $0$
(EGA III 7.8.1); \underline{more precisely} (EGA III 7.6.9(ii)),
\underline{one has}

\begin{equation*}
k\xrightarrow{\;\sim\;}H^0(X_s,\underline{O}_{X_s}).
\end{equation*}
It is easy to reduce to the case in which $S$ is strictly local.  Since $Y$ is
geometrically connected, $H^0(Y,\underline{O}_Y)$ is a purely inseparable
extension of $k$; and since $(d,p)=1$, at least one irreducible component of $Y$
is


\SGApage{II}{\(\approx\)8}
geometrically reduced.  It follows readily that
$H^0(Y,\underline{O}_Y)=k$.

Let $D$ be a divisor on $X$ such that

\begin{equation*}
Y\leq D<X_s \quad .
\end{equation*}
Since $d=1$, $D$ is not a rational multiple of $X_s$; hence, by 1.8,

\begin{equation*}
(D,X_s-D)=-(D,D)>0 \quad .
\end{equation*}
The divisor $Y$ therefore has an irreducible component $E$ such that

\begin{equation*}
0<E\leq X_s-D \quad\text{and}\quad (D,E)>0 \quad .
\end{equation*}
Consider the diagram

\begin{center}
\begin{tikzcd}[column sep=large, row sep=large]
0 \arrow[r] \& \underline{O}(-D-E) \arrow[r] \arrow[d, "\varphi"] \&
\underline{O} \arrow[r] \arrow[d, equal] \&
\underline{O}_{D+E} \arrow[r] \arrow[d, "\psi"] \& 0 \\
0 \arrow[r] \& \underline{O}(-D) \arrow[r] \&
\underline{O} \arrow[r] \& \underline{O}_D \arrow[r] \& 0
\end{tikzcd}
\end{center}
The snake lemma gives an isomorphism between $\operatorname{Ker}(\psi)$ and
$\operatorname{Coker}(\varphi)$.  This cokernel is precisely the restriction of
$\underline{O}(-D)$ to $E$, whence an exact sequence

\SGAvar{\begin{equation*}
0\longrightarrow\underline{O}(-D)|_E
\longrightarrow\underline{O}_{D+E}
\longrightarrow\underline{O}_D\longrightarrow0 \quad .
\end{equation*}
By construction,
$\deg_E(\underline{O}(-D)=-(D,E)<0$, so that
$H^0(\underline{O}(-D)|_E)=0$ and the restriction map

}{\begin{equation*}
0\longrightarrow\underline{O}(-D)|_E
\longrightarrow\underline{O}_{D+E}
\longrightarrow\underline{O}_D\longrightarrow0 \quad .
\end{equation*}
By construction,
$\deg_E(\underline{O}(-D)|_E)=-(D,E)<0$, so that
$H^0(\underline{O}(-D)|_E)=0$ and the restriction map

}{E7II-X-005}\begin{equation*}
0\longrightarrow H^0(\underline{O}_{D+E})
\longrightarrow H^0(\underline{O}_D)
\end{equation*}
is injective.  This allows us, by induction, to construct a sequence of divisors

\begin{equation*}
Y=D_0\leq D_1\leq\cdots\leq D_n=X_s
\end{equation*}
such that $H^0(\underline{O}_{D_i})$ injects into
$H^0(\underline{O}_{D_{i-1}})$.  Thus
$H^0(\underline{O}_{X_s})=H^0(\underline{O}_Y)=k$, which completes the proof of
1.13.


\SGApage{II}{\(\approx\)9}
The following fact is proved in [5] or [10].

\underline{Theorem} 1.14 (Castelnuovo, Enriques, \ldots).
\underline{Under the hypotheses of} (1.1), \underline{let} $D$
\underline{be a reduced irreducible divisor contained in} $Y$.
\underline{Suppose that} $H^1(D,\underline{O}_D)=0$ \underline{and that, on putting}
$k'=H^0(D,\underline{O}_D)$, \underline{one has}

\SGAvar{\begin{equation*}
(D,D)=-[k':k] \quad .
\end{equation*}
\underline{Then} $D$ \underline{is a projective line over} $k'$, \underline{and
there exist a regular} $S$\underline{-scheme} $X_1$, \underline{a closed point}
$x_1$ \underline{of} $X_1$ \underline{with residue field} $k'$, \underline{and an
isomorphism from} $X$ \underline{to the} $S$\underline{-scheme obtained from}
$X_1$ \underline{by blowing up} $x_1$, \underline{under which} $D$
\underline{is identified with the fibre over} $x_1$ \underline{in} $X_1$.

}{\begin{equation*}
(D,D)=-[k':k] \quad .
\end{equation*}
\underline{Then} $D$ \underline{is a projective line over} $k'$, \underline{and
there exist a regular} $S$\underline{-scheme} $X_1$, \underline{a closed point}
$x_1$ \underline{of} $X_1$ \underline{with residue field} $k'$, \underline{and an
isomorphism from} $X$ \underline{to the} $S$\underline{-scheme obtained from}
$X_1$ \underline{by blowing up} $x_1$, \underline{under which} $D$
\underline{is identified with the fibre over} $x_1$.

}{E7II-X-006}A divisor $D$ having these properties is called an \underline{exceptional curve
of the first kind}.

I cannot resist the pleasure of proving the following proposition of Néron [8],
Proposition 1$'$, p.~90.

\underline{Proposition} 1.15. \underline{Suppose that} $X_\eta$
\underline{is an elliptic curve (in particular, that it has a section), and that}
$X$ \underline{contains no exceptional curve of the first kind.  Then, with the
notation of} (1.10.1), \underline{the canonical map}

\begin{equation*}
f^*f_*\omega_{X/S}\xrightarrow{\;\sim\;}\omega_{X/S}
\end{equation*}
\underline{is universally an isomorphism}.

Since $d=1$ (1.12), by (1.13) it is enough to prove that $\omega_{X/S}$ is
isomorphic to $\underline{O}_X$.  Since $\omega_{X/S}$ is trivial on the elliptic
curve $X_\eta$, we know a priori that

\begin{equation*}
\omega_{X/S}\simeq\underline{O}\Bigl(\sum_i n_iY_i\Bigr),
\end{equation*}


\SGApage{II}{\(\approx\)10}
where the $Y_i$ are the irreducible components of $Y$; and since
$\operatorname{div}(\pi)\sim0$, we may assume the $n_i$ positive.  If
$\sum_i n_iY_i$ is proportional to $\operatorname{div}(\pi)$, it is an integral
multiple of $\operatorname{div}(\pi)$ and we are done.  Otherwise,

\begin{equation*}
\Bigl(\sum_i n_iY_i\Bigr)^2<0,
\end{equation*}
there is an $i$ such that

\begin{equation*}
\Bigl(Y_i,\sum_j n_jY_j\Bigr)<0,
\end{equation*}
and the following lemma applies.

\underline{Lemma} 1.16. \underline{Under the hypotheses of} 1.1,
\underline{suppose that} $Y$ \underline{is connected and that either} $Y$
\underline{is reducible or} $\chi(\underline{O}_Y)\leq0$.
\underline{Then an irreducible component} $Y_i$ \underline{of} $Y$
\underline{is an exceptional curve of the first kind if and only if}
$(Y_i,K)<0$.

Let $k'=H^0(Y_i,\underline{O}_{Y_i})$.  The Riemann--Roch formula 1.11 gives

\begin{equation*}
\chi(\underline{O}_{Y_i})
=[k':k]\bigl(1-\dim_{k'}H^1(Y_i,\underline{O}_{Y_i})\bigr)
=-\frac12\bigl((Y_i,Y_i)+(Y_i,K)\bigr) \quad .
\end{equation*}
If $Y_i$ is exceptional of the first kind, it follows that
$(Y_i,K)=-[k':k]<0$.  Conversely, if $(Y_i,K)<0$, then
$\chi(\underline{O}_{Y_i})>0$, which implies
$\chi(\underline{O}_{Y_i})=[k':k]$ and that $Y$ is reducible.  Hence, by 1.8,
$(Y_i,Y_i)<0$.  The negative integers $(Y_i,Y_i)$ and $(Y_i,K)$ are divisible by
$[k':k]$, and therefore necessarily

\begin{equation*}
(Y_i,Y_i)=(Y_i,K)=-[k':k] \quad .
\end{equation*}

\underline{Exercise}. (a) Under the hypotheses of (1.1), suppose that $X_\eta$ is
a principal homogeneous space under an elliptic curve and that $X$ contains no
exceptional curve of the first kind.  Show that on $X$ the Riemann--Roch formula
(1.11) takes the form

\begin{equation*}
\chi(\underline{O}_D)=-\frac12(D,D).
\end{equation*}


\SGApage{II}{\(\approx\)11}
(b) Suppose in addition that the residue field $k(s)$ is algebraically closed.
Show that the diagram of irreducible components of $X_s$ is a multiple of one of
the diagrams considered by Néron [8], pp.~124--125, fourth column (cases a through
c8).  The diagram determines the genera of the reduced irreducible components,
their multiplicities in $X_s$, and the nature and disposition of the singularities
of $(X_s)_{\mathrm{red}}$.

First treat the case in which $X_s$ is irreducible.  Once this case has been
excluded, prove that every reduced irreducible component is a projective line of
self-intersection $-2$.  If $X_s=\sum_i n_iD_i$ is the decomposition of $X_s$
into reduced irreducible divisors, rewrite this condition as

\SGAvar{\begin{equation}\tag{*}
n_i=\frac12\sum_{j\ne i}n_i(D_i\,D_j) \quad .
\end{equation}
On the other hand, $X_s$ is connected, so that:

}{\begin{equation}\tag{*}
n_i=\frac12\sum_{j\ne i}n_j(D_i\,D_j) \quad .
\end{equation}
On the other hand, $X_s$ is connected, so that:

}{E7II-X-007}(**) For any $i$ and $j$, there exists a chain
$i=i_1,i_2,\ldots,i_n=j$ with
$(D_{i_k}\,D_{i_{k+1}})\ne0$.

Now find all solutions in integers of $(*)+(**)$.  First show that, except in
Néron's cases b2 and c2, the intersection numbers $(D_i\,D_j)$ are either zero or
one.



2. \underline{Italian intersections}

Let $S'$ and $S$ be two smooth projective surfaces over a field $k$, and let
$f:S'\longrightarrow S$ be a birational morphism.  If $D$ and $E$ are two
divisors on $S$, then

\begin{equation*}
(D,E)=(f^*D,f^*E),
\end{equation*}


as one sees, for example, by moving $D$ and $E$.

\SGApage{II}{\(\approx\)12}
It is therefore natural to seek a ``birational'' expression for the intersection
number.  I propose to explain how the Italians proceeded.  My principal source
of information is [1], which contains a long bibliography.  See also Manin [6]
and Northcott [12], [13].

(2.0) Let $S$ be the spectrum of a regular local ring of dimension 2, let $s$ be
the closed point of $S$, and let $k$ be a subfield of $k(s)$ over which $k(s)$ is
finite.  Intersection numbers and Euler--Poincaré characteristics will be
computed over $k$.

In this section, an $S$-scheme $p:S'\longrightarrow S$ will be called a
\underline{birational transform of} $S$ if it is regular and connected, and if
$p$ is proper and induces an isomorphism
$S'\setminus p^{-1}(s)\longrightarrow S\setminus\{s\}$.  Everything said about
$S$ will also apply to the localizations of its birational transforms at a
closed point.

(2.1) Let $f:(S'',p'')\longrightarrow(S',p')$ be a morphism of $S$-schemes
between birational transforms of $S$.  It is known [5] that $f$ has a unique
factorization

\begin{equation*}
f=f_n\cdots f_1,\qquad f_i:S_{i+1}\longrightarrow S_i,
\end{equation*}
such that $S_{i+1}$ is obtained from $S_i$ by blowing up a finite set $X_i$ of
closed points, and $f_i(X_{i+1})\subset X_i$.  The set $X_i$ is the set of points
of $S_i$ at which the projection of $S''$ onto $S_i$ is not an isomorphism.

In particular, $f$ induces an isomorphism over the complement of a finite subset
of $S'$ containing the generic points of the irreducible components of
$p'^{-1}(s)$.  Hence $f^{-1}$ defines an injection


\SGApage{II}{\(\approx\)13}
from the set $\operatorname{Irr}(p'^{-1}(s))$ of irreducible components of
$p'^{-1}(s)$ into $\operatorname{Irr}(p''^{-1}(s))$.

The birational transforms of $S$ have no $S$-automorphisms.  We write $I$ for the
set of $S$-isomorphism classes of birational transforms of $S$, ordered by
domination.  The set $I$ is filtered to the right and has least element $S$.

\underline{Definition} 2.2. \underline{The set} $S^*$ \underline{of points of}
$S$ \underline{infinitely near to} $s$ \underline{is the direct-limit set}

\begin{equation*}
S^*=\varinjlim_I\operatorname{Irr}(p^{-1}(s)) \quad .
\end{equation*}
If an element $t\in S^*$ is represented on a birational transform $S'$ of $S$ by
a complete reduced irreducible curve $T$, the field
$H^0(T,\underline{O}_T)$ depends only on $t$, not on the choice of $S'$.  It is
called the \underline{residue field} $k(t)$ of $t$.

Let $S_1$ be the blow-up of $S$ along $s$.  The projective line $s^*$, the inverse
image of $\{s\}$ in $S_1$, is called the \underline{curvetta} of $s$ and defines
an element $s^*$ of $S^*$ with residue field $k(s)$.

Let $S'$ be a birational transform of $S$, let $t$ be a closed point of $S'$, and
let $S'_t$ be the spectrum of the local ring of $S'$ at $t$.  There is an evident
injection $(S'_t)^*\hookrightarrow S^*$; in particular, $t$ defines an element
$t^*$ of $S^*$ with the same residue field.  We write $(S')^*$ for the union in
$S^*$ of the $(S'_t)^*$ as $t$ ranges over the closed points of $S'$.

The following proposition follows from the description 2.1 of the birational
transforms of $S$.


\SGApage{II}{14}
\underline{Proposition} 2.3. \underline{For every} $t\in S^*$,
\underline{there is a unique sequence of birational transforms}
$(S_i)_{0\leq i\leq n+1}$ \underline{of} $S$ \underline{and closed points}
$t_i\in S_i$ $(0\leq i\leq n)$ \underline{such that}

(i)\quad $S_0=S$, \underline{and} $S_{i+1}$ \underline{is the blow-up of} $S_i$
\underline{along} $\{t_i\}$;

(ii)\quad $t_{i+1}\in t_i^*$ $(0\leq i<n)$, \underline{and} $t=t_n^*$.

\underline{Moreover, for every birational transform} $(S',p')$ \underline{of}
$S$, \underline{the element} $t$ \underline{belongs to the image of}
$\operatorname{Irr}(p'^{-1}(s))$ \underline{in} $S^*$ \underline{if and only if}
$S'$ \underline{dominates} $S_{n+1}$.

The integer $n$ is called the order of $t$.  We write $S(t)$ (respectively
$S_+(t)$) for the birational transform $S_n$ (respectively $S_{n+1}$) of $S$.
By abuse of notation, we also write $t$ for the closed point $t_n$ of $S_n$, and
$t^*$ for its inverse image $t_n^*$ in $S_{n+1}$.

We equip $S^*$ with the following order:

\SGAvar{\begin{equation*}
t_1\prec t_2
\quad\Longleftrightarrow\quad
S_+(t_2)\ \text{dominates}\ S_+(t_1) \quad .
\end{equation*}
If $t_1\prec t_2$, then $k(t_2)$ is an extension of $k(t_1)$.  The ordered set
$S^*$ has least element $s^*$, and every bounded subset of $S^*$ is finite and
totally ordered ($S^*$ is a ``tree rooted at $s^*$'').  Moreover, the order $n$ of
$t\in S^*$ is the distance in the tree $S^*$ from $s^*$ to $t$ ($n = \#[s^*,t]$).

}{\begin{equation*}
t_1\prec t_2
\quad\Longleftrightarrow\quad
S_+(t_2)\ \text{dominates}\ S_+(t_1) \quad .
\end{equation*}
If $t_1\prec t_2$, then $k(t_2)$ is an extension of $k(t_1)$.  The ordered set
$S^*$ has least element $s^*$, and every bounded subset of $S^*$ is finite and
totally ordered ($S^*$ is a ``tree rooted at $s^*$'').  Moreover, the order $n$ of
$t\in S^*$ is the distance in the tree $S^*$ from $s^*$ to $t$ ($n = \#]s^*,t]$).

}{E7II-X-008}The following proposition makes the description 2.1 of the birational
transforms of $S$ explicit.

\underline{Proposition} 2.4. \underline{The function assigning to a birational
transform} $p:S'\to S$ \underline{the subset}
$\operatorname{Irr}(p^{-1}(s))$ \underline{of} $S^*$ \underline{induces a
bijection between} $I$ (2.1) \underline{and the finite subsets of} $S^*$
\underline{which, together with each element,}


\SGApage{II}{\(\approx\)15}
\underline{contain all its predecessors}.

Let $D$ be a divisor $\geq0$ on $S$, and let $t\in S^*$.  We say that $t$
\underline{belongs to} $D$ if the closed point $t$ of $S(t)$ (2.3) belongs to the
strict transform (the schematic closure of $D-\{s\}$) of $D$ on $S(t)$.

\underline{Definition} 2.5. \underline{Let} $t,t'\in S^*$.
\underline{We say that} $t'$ \underline{is proximate to} $t$ \underline{if}
$t'\in S_+(t)^*$ (2.2) \underline{and belongs to the curvetta} $t^*$ \underline{of}
$t$ (2.3).

\SGAvar{If $t'$ is an immediate successor of $t$, then $t'$ is proximate to $t$.  If in a
sequence $(t_i)_{0\leq i\leq n}$ every $t_{i+1}$ is an immediate successor of
$t_i$, and if $t_n$ is proximate to $t_o$, then every $t_i$ with $0<i\leq n$ is
proximate to $t_o$, while $t_n$ is not proximate to $t_i$ for
$0<i\leq n-1$.

}{If $t'$ is an immediate successor of $t$, then $t'$ is proximate to $t$.  If in a
sequence $(t_i)_{0\leq i\leq n}$ every $t_{i+1}$ is an immediate successor of
$t_i$, and if $t_n$ is proximate to $t_0$, then every $t_i$ with $0<i\leq n$ is
proximate to $t_0$, while $t_n$ is not proximate to $t_i$ for
$0<i<n-1$.

}{E7II-X-009}(2.6) Let $A$ be a Noetherian local ring of dimension $n$, let $T$ be its
spectrum, let $m$ be its maximal ideal, let $k=A/m$ be its residue field, and let
$P=\sum_{0\leq i<n}a_iX^i$ be its Hilbert--Samuel polynomial:

\SGAvar{\begin{equation*}
P(n)=\dim_k(m^n/m^{n+1})\qquad\text{for sufficiently large }n.
\end{equation*}
Recall that the multiplicity $\mu_s$ of the closed point $s$ in $T$ is defined by

}{\begin{equation*}
P(r)=\dim_k(m^r/m^{r+1})\qquad\text{for sufficiently large }r.
\end{equation*}
Recall that the multiplicity $\mu_s$ of the closed point $s$ in $T$ is defined by

}{E7II-X-010}\begin{equation*}
a_{n-1}=\frac{\mu_s}{(n-1)!} \quad .
\end{equation*}
When $n=1$, this simply says

\SGAvar{\begin{equation*}
\mu_s=\dim_k(m^n/m^{n+1})\qquad\text{for sufficiently large }n.
\end{equation*}
Geometrically, $\mu_s$ may be defined as the degree, in the projective space
$\mathbf P_k(m/m^2)$, of the special fibre of the blow-up of $T$ along the closed
subscheme $\{s\}$.


}{\begin{equation*}
\mu_s=\dim_k(m^r/m^{r+1})\qquad\text{for sufficiently large }r.
\end{equation*}
Geometrically, $\mu_s$ may be defined as the degree, in the projective space
$\mathbf P_k(m/m^2)$, of the special fibre of the blow-up of $T$ along the closed
subscheme $\{s\}$.


}{E7II-X-011}\SGApage{II}{\(\approx\)16}
Recall that if $Y$ and $Z$ are two closed subschemes of a scheme $X$, then the
strict transform of $Y$ in the blow-up $\widetilde X$ of $X$ along $Z$---the
schematic closure of $Y\setminus Z$ in $\widetilde X$---is precisely the blow-up
of $Y$ along $Y\cap Z$.

Suppose that $T$ is a closed subscheme of a regular scheme $U$.  Let
$p:U_1\longrightarrow U$ be the blow-up of $U$ along $\{s\}$, let $s^*$ be the
inverse image of $s$ in $U_1$, and let $T'$ be the strict transform of $T$ in
$U_1$.  The preceding discussion shows that the multiplicity of $T$ at $s$ is
the degree, in the projective space $s^*$, of the subscheme $s^*\cap T'$.

Suppose further that $T$ is a divisor on $U$.  Let $T^*$ be its total transform
in $U_1$ (the inverse-image divisor), and define $\nu$ by
$T^*=\nu s^*+T'$.  If $X$ is a line in $s^*$, and Euler--Poincaré
characteristics are calculated as alternating sums of dimensions over $k$, then

\begin{equation}\tag{2.6.1}
(X,s^*)=
\chi(\underline{O}_X\overset{\mathbb L}{\otimes}\underline{O}_{s^*})
=\deg_X\underline{O}(s^*)=-1,
\end{equation}
and we have seen that

\begin{equation}\tag{2.6.2}
(X,T')=
\chi(\underline{O}_X\overset{\mathbb L}{\otimes}\underline{O}_{T'})
=\deg_{s^*}(s^*\cap T')=\mu_s \quad .
\end{equation}
The projection formula

\begin{equation*}
Rp_*(\underline{O}_X\overset{\mathbb L}{\otimes}\underline{O}_{T^*})
=Rp_*(\underline{O}_X\overset{\mathbb L}{\otimes}Lp^*\underline{O}_T)
=Rp_*\underline{O}_X\overset{\mathbb L}{\otimes}\underline{O}_T
\end{equation*}
and its corollary

\begin{equation*}
(X,T^*)=\nu(X,s^*)+(X,T')=0
\end{equation*}
give, by (2.6.1) and (2.6.2), the identity

\begin{equation}\tag{2.6.3}
\mu_s=\nu \quad .
\end{equation}


\SGApage{II}{17}
\underline{Proposition} 2.7. \underline{Let} $T$ \underline{be a divisor}
$\geq0$ \underline{on a regular scheme} $U$, $s$ \underline{a closed point of}
$U$, $U_1$ \underline{the blow-up of} $U$ \underline{along} $\{s\}$, $s^*$
\underline{the inverse image of} $s$ \underline{in} $U_1$, $T'$ \underline{the
strict transform of} $T$ \underline{on} $U_1$, \underline{and} $T^*$
\underline{its total transform}.

\underline{The following three quantities are equal}:

(i)\quad \underline{the multiplicity, in Samuel's sense, of the point} $s$
\underline{on} $T$;

(ii)\quad \underline{the degree, in the projective space} $s^*$,
\underline{of} $s^*\cap T'$;

(iii)\quad \underline{the multiplicity of} $s^*$ \underline{in the divisor}
$T^*$.

2.8. Return to the hypotheses of 2.0.  Let $D$ be a divisor $\geq0$ on $S$ and
let $t\in S^*$.  The \underline{multiplicity} $\mu_t(D)$ \underline{of} $D$
\underline{at} $t$ is, by definition, the multiplicity at $t$ of the strict
transform of $D$ on $S(t)$ (2.3).  The support of the function
$t\mapsto\mu_t(D)$ is the set of points infinitely near to $s$ that belong to
$D$ (2.4).  When $D$ is regular, the numbers
$\mu_t(D)[k(t):k(s)]$ are zero or one: indeed, the strict transforms $D'$ of $D$
are isomorphic to $D$, and one applies the definition (2.6) of multiplicity.

Intersection numbers can be calculated in terms of multiplicities.

\underline{Theorem} 2.9. \underline{If the divisors} $D$ \underline{and} $E$
\underline{meet only at} $s$, \underline{then}

\begin{equation*}
(D,E)=\sum_{t\in S^*}\mu_t(D)\mu_t(E)[k(t):k],
\end{equation*}
\underline{where the sum is over the points infinitely near to} $s$
\underline{which belong to both} $D$ \underline{and} $E$.


\underline{Proof}. Let $p:S'\longrightarrow S$ be a birational transform of $S$.
We denote the total transforms of $D$ and $E$ by $D^*$ and $E^*$, and their
strict transforms by $D'$ and $E'$.


\SGApage{II}{18}
\underline{Lemma} 2.9.1. (i) \underline{The total transform} $D^*$
\underline{is characterized by}:

\SGAvar{(a)\quad $D^*-D'$ \underline{is supported on} $f^{-1}(s)$;

(b)\quad \underline{if} $X$ \underline{is supported on} $f^{-1}(s)$,
\underline{then} $(D^*,X)=0$.

}{(a)\quad $D^*-D'$ \underline{is supported on} $p^{-1}(s)$;

(b)\quad \underline{if} $X$ \underline{is supported on} $p^{-1}(s)$,
\underline{then} $(D^*,X)=0$.

}{E7II-X-012}(ii)\quad $(D,E)=(D^*,E^*)=(D^*,E')=(D',E^*)$.

\SGAvar{Assertion (i)(a) is immediate, and (i)(b) follows from 1.6(vi).  That (a) and (b)
determine $D^*$ follows from 1.9.  By construction, the cokernel of the map
$\underline{O}_E\longrightarrow f_*\underline{O}_{E}$ is finite; hence

\begin{equation*}
\begin{aligned}
(D^*,E')
&=\chi(f^*\underline{O}_D\otimes\underline{O}_{E'})\\
&=\chi(\underline{O}_D\otimes f_*\underline{O}_{E'})
=\chi(\underline{O}_D\otimes\underline{O}_E)
=(D,E),
\end{aligned}
\end{equation*}
which, by (i), implies (ii).

}{Assertion (i)(a) is immediate, and (i)(b) follows from 1.6(vi).  That (a) and (b)
determine $D^*$ follows from 1.9.  By construction, the cokernel of the map
$\underline{O}_E\longrightarrow p_*\underline{O}_{E'}$ is finite; hence

\begin{equation*}
\begin{aligned}
(D^*,E')
&=\chi(Lp^*\underline{O}_D\otimes\underline{O}_{E'})\\
&=\chi(\underline{O}_D\otimes Rp_*\underline{O}_{E'})
=\chi(\underline{O}_D\otimes\underline{O}_E)
=(D,E),
\end{aligned}
\end{equation*}
which, by (i), implies (ii).

}{E7II-X-013}Take for $S'$ the blow-up of $S$ at $s$.  By (2.9.1), (2.6.3), and 2.7,

\SGAvar{\begin{equation}\tag{2.9.2}
\begin{aligned}
(D,E)
&=(D^*,E')
=(D'+\mu_s(S)s^*,E')\\
&=(D',E')+\mu_s(D)(s^*,E')\\
&=(D',E')+\mu_s(D)\mu_s(E)[k(s):k].
\end{aligned}
\end{equation}
Theorem 2.9 now follows by induction on $(D,E)$, through repeated application of
(2.9.2).

}{\begin{equation}\tag{2.9.2}
\begin{aligned}
(D,E)
&=(D^*,E')
=(D'+\mu_s(D)s^*,E')\\
&=(D',E')+\mu_s(D)(s^*,E')\\
&=(D',E')+\mu_s(D)\mu_s(E)[k(s):k].
\end{aligned}
\end{equation}
Theorem 2.9 now follows by induction on $(D,E)$, through repeated application of
(2.9.2).

}{E7II-X-014}(2.10) Let $A$ be a reduced Noetherian local ring of dimension one, let $s$ be the
closed point of $\operatorname{Spec}(A)$, let $\widetilde A$ be the normalization
of $A$, and let $k$ be a field over which $k(s)$ is finite.  Suppose that
$\widetilde A$ is finite over $A$ (equivalently, that the completion $\widehat A$
is reduced).  The \underline{$\delta$-invariant of the singularity of}
$\operatorname{Spec}(A)$ \underline{at} $s$ is the length over $k$ of the
$A$-module $\widetilde A/A$.


\SGApage{II}{19}
\underline{Theorem} 2.11. \underline{Let} $D$ \underline{be a reduced curve on}
$S$. \underline{Then}:

(i)\quad \underline{for every} $t\in S^*$, \underline{with the notation of} 2.8,
\underline{one has}

\begin{equation*}
\mu_t(D)
=\sum_{\substack{t'\ \mathrm{proximate}\\\mathrm{to}\ t}}
\mu_{t'}(D)[k(t'):k(t)] \quad .
\end{equation*}

(ii)\quad \underline{Suppose that the normalization} $\widetilde D$
\underline{of} $D$ \underline{is finite over} $D$, \underline{and let}
$\delta=\chi(\underline{O}_{\widetilde D}/\underline{O}_D)$
\underline{be the invariant of the singularity of} $D$ \underline{at} $s$.
\underline{For all but finitely many} $t\in S^*$, $\mu_t(D)$ \underline{is zero
or one, and}

\begin{equation*}
\delta
=\sum_{t\in S^*}\frac12\mu_t(D)\bigl(\mu_t(D)-1\bigr)[k(t):k] \quad .
\end{equation*}

Let $p:S_1\longrightarrow S$ be the blow-up of $S$ along $\{s\}$, let $D^*$ be
the total transform of $D$, and let $D'$ be its strict transform.  By (2.7),

\begin{equation}\tag{2.11.1}
D^*=D'+\mu_s(D)s^*
\end{equation}
and
\begin{equation}\tag{2.11.2}
\mu_s(D)=(s^*,D')\,[k(s):k]^{-1} \quad .
\end{equation}
The curvetta $s^*$ is nonsingular.  Thus, by 2.8 and Definition 2.5, formula 2.9,
applied to the divisors $s^*$ and $D'$ on $S_1$, becomes

\SGAvar{\begin{equation}\tag{2.11.3}
(s^*,D)
=\sum_{\substack{t'\ \mathrm{proximate}\\\mathrm{to}\ s}}
\mu_{t'}(D)[k(t'):k] \quad .
\end{equation}
Assertion (i) for $s=t$ follows from (2.11.2) and (2.11.3); the general case is
obtained by applying this result to the localization of $S(t)$ at $t$ (2.3).

}{\begin{equation}\tag{2.11.3}
(s^*,D')
=\sum_{\substack{t'\ \mathrm{proximate}\\\mathrm{to}\ s}}
\mu_{t'}(D)[k(t'):k] \quad .
\end{equation}
Assertion (i) for $s=t$ follows from (2.11.2) and (2.11.3); the general case is
obtained by applying this result to the localization of $S(t)$ at $t$ (2.3).

}{E7II-X-015}Let $m$ be the maximal ideal defining $\{s\}$, and let $n\geq0$.  It is known
(SGA VII 3.5) that

\begin{equation}\tag{2.11.4}
p_*\underline{O}_{S_1}(-ns^*)=m^n
\quad\text{and}\quad
R^ip_*\underline{O}_{S_1}(-ns^*)=0\quad(i>0).
\end{equation}


\SGApage{II}{\(\approx\)20}
It follows from this and from the exact sequence

\begin{equation*}
0\longrightarrow\underline{O}_{S_1}(-ns^*)
\longrightarrow\underline{O}_{S_1}
\longrightarrow\underline{O}_{ns^*}
\longrightarrow0
\end{equation*}
that

\begin{equation}\tag{2.11.5}
p_*\underline{O}_{ns^*}=\underline{O}_S/m^n
\quad\text{and}\quad
R^ip_*\underline{O}_{ns^*}=0\quad(i>0),
\end{equation}
and that

\begin{equation}\tag{2.11.6}
\chi(\underline{O}_{ns^*})
=\chi(\underline{O}_S/m^n)
=\frac{n(n+1)}2[k(s):k] \quad .
\end{equation}
Put $n=0$ in (2.11.4).  For every complex $K$ of coherent sheaves on $S$, the
projection formula gives

\SGAvar{\begin{equation*}
K\xrightarrow{\;\sim\;}Rp_*Lp^*K
\qquad
\bigl(\simeq K\overset{\mathbb L}{\otimes}Rp_*\underline{O}_{S_1}\bigr).
\end{equation*}
In particular, since
$Lf^*\underline{O}_D=f^*\underline{O}_D=\underline{O}_{D^*}$, one has

}{\begin{equation*}
K\xrightarrow{\;\sim\;}Rp_*Lp^*K
\qquad
\bigl(\simeq K\overset{\mathbb L}{\otimes}Rp_*\underline{O}_{S_1}\bigr).
\end{equation*}
In particular, since
$Lp^*\underline{O}_D=p^*\underline{O}_D=\underline{O}_{D^*}$, one has

}{E7II-X-016}\begin{equation}\tag{2.11.7}
p_*\underline{O}_{D^*}=\underline{O}_D
\quad\text{and}\quad
R^ip_*\underline{O}_{D^*}=0\quad(i>0).
\end{equation}
Put $E=\mu_s(D)s^*$.  By (2.11.1), (2.9.1)(i)(b), and (2.6.1),

\begin{equation}\tag{2.11.8}
\begin{aligned}
\chi(\underline{O}_{D'\cap E})
&=(D',E)=(D^*-E,E)=-(E,E)\\
&=-\mu_s(D)^2(s^*,s^*)
=\mu_s(D)^2[k(s):k].
\end{aligned}
\end{equation}
The exact sequence

\begin{equation*}
0\longrightarrow\underline{O}_{D^*}
\longrightarrow\underline{O}_{D'}\oplus\underline{O}_E
\longrightarrow\underline{O}_{D'\cap E}
\longrightarrow0
\end{equation*}
gives, by direct image and (2.11.7), the exact sequence

\SGAvar{\begin{equation*}
0\longrightarrow\underline{O}_D
\longrightarrow f_*\underline{O}_{D'}\oplus f_*\underline{O}_E
\longrightarrow f_*\underline{O}_{D'\cap E}
\longrightarrow0,
\end{equation*}
and hence

\begin{equation*}
0\longrightarrow f_*\underline{O}_{D'}/\underline{O}_D
\longrightarrow f_*\underline{O}_E \longrightarrow f_*\ \underline{O}_{D'\cap E} \longrightarrow0 \quad .
\end{equation*}


}{\begin{equation*}
0\longrightarrow\underline{O}_D
\longrightarrow p_*\underline{O}_{D'}\oplus p_*\underline{O}_E
\longrightarrow p_*\underline{O}_{D'\cap E}
\longrightarrow0,
\end{equation*}
and hence

\begin{equation*}
0\longrightarrow p_*\underline{O}_E
\longrightarrow p_*\underline{O}_{D'\cap E}
\longrightarrow p_*\underline{O}_{D'}/\underline{O}_D
\longrightarrow0 \quad .
\end{equation*}


}{E7II-X-017}\SGApage{II}{\(\approx\)21}
Consequently, by (2.11.6) and (2.11.8),

\SGAvar{\begin{equation}\tag{2.11.9}
\begin{aligned}
\chi(f_*\underline{O}_{D'}/\underline{O}_D)
&=\frac12\mu_s(D)\bigl(\mu_s(D)+1\bigr)[k(s): k] - \mu_s(D)^2[k(s):k]\\
&=\frac12\mu_s(D)\bigl(\mu_s(D)-1\bigr)[k(s):k] \quad .
\end{aligned}
\end{equation}
The schemes $D$ and $D'$ are reduced, differ only over $s$, and $D'$ is finite
over $D$.  They therefore have the same normalization, and the $\delta$-invariant
of the singularity of $D$ at $s$ is the sum of
$\chi(f_*\underline{O}_{D'}/\underline{O}_D)$ and the invariants $S$ of the
singularities of $D'$ at the points of $D'\cap s^*$.  We prove 2.11(ii) by
induction on $\delta$.  If $\mu_s(D)>1$, the induction hypothesis applies to
$D'$, and 2.11(ii) follows from (2.11.9).  If $\mu_s(D)=1$, then
$(D',s^*)=1$ and $D'\cap s^*$ is a regular divisor on $D'$; by EGA IV 19.1.1,
$D'$ is regular.  Formula (2.11.9) then shows that $D$ is itself regular, and
2.11(ii) follows from 2.8.



}{\begin{equation}\tag{2.11.9}
\begin{aligned}
\chi(p_*\underline{O}_{D'}/\underline{O}_D)
&=\mu_s(D)^2[k(s):k]
-\frac12\mu_s(D)\bigl(\mu_s(D)+1\bigr)[k(s):k]\\
&=\frac12\mu_s(D)\bigl(\mu_s(D)-1\bigr)[k(s):k] \quad .
\end{aligned}
\end{equation}
The schemes $D$ and $D'$ are reduced, differ only over $s$, and $D'$ is finite
over $D$.  They therefore have the same normalization, and the $\delta$-invariant
of the singularity of $D$ at $s$ is the sum of
$\chi(p_*\underline{O}_{D'}/\underline{O}_D)$ and the invariants $\delta$ of the
singularities of $D'$ at the points of $D'\cap s^*$.  We prove 2.11(ii) by
induction on $\delta$.  If $\mu_s(D)>1$, the induction hypothesis applies to
$D'$, and 2.11(ii) follows from (2.11.9).  If $\mu_s(D)=1$, then
$(D',s^*)=1$ and $D'\cap s^*$ is a regular divisor on $D'$; by EGA IV 19.1.1,
$D'$ is regular.  Formula (2.11.9) then shows that $D$ is itself regular, and
2.11(ii) follows from 2.8.



}{E7II-X-018}2.12. Let $t\in S^*$, and let $S'$ be a birational transform of $S$ which
dominates $S_+(t)$ (2.3).  The \underline{total transform} of $t$ on $S'$,
again denoted by $t^*$, is the total transform on $S'$ of the divisor $t^*$ on
$S_+(t)$.

Let $\nu:S^*\longrightarrow\mathbf N$ be a function of finite support, and
suppose that, for every $t\in S^*$,

\begin{equation}\tag{2.12.1}
\nu(t)\geq
\sum_{\substack{t'\ \mathrm{proximate}\\\mathrm{to}\ t}}
\nu(t')[k(t'):k(t)]
\end{equation}
(the proximate-point inequality).  Let $f:S'\longrightarrow S$ be a birational
transform of $S$ which dominates $S_+(t)$ for each $t$ in the support of $\nu$,
and let
\[
\nu^*=\sum_{t\in S^*}\nu(t)t^*
\]
be the corresponding divisor on $S'$.


\SGApage{II}{\(\approx\)22}
We write $I(\nu)$ for the ideal of $\underline{O}_S$ obtained by direct image
under $f$ of the ideal of $\underline{O}_{S'}$ defining $\nu^*$:

\begin{equation*}
I(\nu)=f_*\underline{O}_{S'}(-\nu^*) \quad .
\end{equation*}
By (2.11.4), this ideal does not depend on the choice of $S'$.  The ideals of
$S$ obtained in this way are those which Zariski [11] calls ``complete.''

\underline{Theorem} 2.13. \underline{Under the preceding hypotheses, one has}

\begin{equation*}
\chi(\underline{O}_S/I(\nu))
=\sum_{t\in S^*}\frac{\nu(t)(\nu(t)+1)}2[k(t):k],
\end{equation*}
\underline{and}

\begin{equation*}
R^if_*\underline{O}_{S'}(-\nu^*)=0
\quad\text{for }i>0.
\end{equation*}

In old French, this is expressed briefly by saying: ``requiring a plane curve to
pass $n$ times through a point $t$, whether proper or infinitely near, imposes
$n(n+1)/2$ conditions, generally independent, on the coefficients of its
equation.''

We first prove the following lemma.

\SGAvar{\underline{Lemma} 2.14. \underline{Let} $p:S_1\longrightarrow S$
\underline{be the blow-up of} $S$ \underline{along} $\{s\}$; \underline{let}
$T$ \underline{be a closed subscheme of} $S_1$ \underline{with ideal} $q$;
\underline{let} $D$ \underline{be a divisor}; \underline{let} $n$ \underline{be an integer}; \underline{let}
$J=\underline{O}_{S_1}(-s^*)$ \underline{be the ideal defining the divisor}
$s^*$; \underline{put}
$Y=\operatorname{Spec}(\underline{O}_{S_1}/J^nq)$; \underline{and let} $D$
\underline{be a regular divisor on} $S$, \underline{with strict transform} $D'$
\underline{on} $S_1$. \underline{Suppose that} $T\cap s^*$ \underline{and}
$T\cap D'$ \underline{are finite}.

}{\underline{Lemma} 2.14. \underline{Let} $p:S_1\longrightarrow S$
\underline{be the blow-up of} $S$ \underline{along} $\{s\}$; \underline{let}
$T$ \underline{be a closed subscheme of} $S_1$ \underline{with ideal} $q$;
\underline{let} $n$ \underline{be an integer}; \underline{let}
$J=\underline{O}_{S_1}(-s^*)$ \underline{be the ideal defining the divisor}
$s^*$; \underline{put}
$Y=\operatorname{Spec}(\underline{O}_{S_1}/J^nq)$; \underline{and let} $D$
\underline{be a regular divisor on} $S$, \underline{with strict transform} $D'$
\underline{on} $S_1$. \underline{Suppose that} $T\cap s^*$ \underline{and}
$T\cap D'$ \underline{are finite}.

}{E7II-X-019}(i)\quad \underline{If}
$\operatorname{length}_{k(s)}(T\cap s^*)\leq n+1$, \underline{then}

\begin{equation*}
R^ip_*(qJ^n)=0\quad\text{for }i>0,
\end{equation*}
\underline{and, if} $T$ \underline{is Artinian,}

\begin{equation*}
\operatorname{length}_{k(s)}
\bigl(\underline{O}_S/p_*(qJ^n)\bigr)
=\frac{n(n+1)}2+\operatorname{length}_{k(s)}(T).
\end{equation*}


(ii)\quad \underline{If}
$\operatorname{length}_{k(s)}(T\cap s^*)\leq n$, \underline{then}

\SGApage{II}{\(\approx\)23}
\begin{equation*}
\operatorname{length}_{k(s)}
\Bigl(D\cap
\operatorname{Spec}\bigl(\underline{O}_S/p_*(qJ^n)\bigr)\Bigr)
=n+\operatorname{length}_{k(s)}(D'\cap T) \quad .
\end{equation*}

Consider the exact sequence on $S_1$

\SGAvar{\begin{equation}\tag{2.14.1}
0\longrightarrow J^nq
\longrightarrow J^n
\xrightarrow{\,u\,}J^n\otimes\underline{O}_T
\longrightarrow0 \quad .
\end{equation}
To prove that $R^ip_*(J^nq)=0$ for $i>0$, it is enough, by (2.11.4), to show
that $p_*(u)$ is surjective.  By the theorem on holomorphic functions, it suffices to
prove that the maps

}{\begin{equation}\tag{2.14.1}
0\longrightarrow J^nq
\longrightarrow J^n
\xrightarrow{\,u\,}J^n\otimes\underline{O}_T
\longrightarrow0 \quad .
\end{equation}
To prove that $R^ip_*(J^nq)=0$ for $i>0$, it is enough, by (2.11.4), to show
that $p_*(u)$ is surjective.  By the theorem on formal functions, it suffices to
prove that the maps

}{E7II-X-020}\begin{equation*}
u_k:p_*(J^n\otimes\underline{O}/J^k)
\longrightarrow
p_*(J^n\otimes\underline{O}_T\otimes\underline{O}/J^k)
\end{equation*}
are surjective.

Write $\operatorname{Gr}_J$ for the associated graded object of the $J$-adic
filtration.  The module $J^m/J^{m+1}$ is isomorphic to
$\underline{O}_{s^*}(m)$; hence $R^1p_*(J^m/J^{m+1})=0$, and the map
\[
p_*(J^n/J^{n+k+1})\longrightarrow p_*(J^n/J^{n+k})
\]
is surjective.  Therefore, to prove by induction on $k$ that the maps $u_k$ are
surjective, it suffices to prove the surjectivity of

\begin{equation*}
v_k:p_*\bigl(J^n\otimes\operatorname{Gr}_J^k(\underline{O})\bigr)
\longrightarrow
p_*\bigl(J^n\otimes\operatorname{Gr}_J^k(\underline{O}_T)\bigr)
\qquad(k\geq0).
\end{equation*}
The module $\operatorname{Gr}_J^k(\underline{O}_T)$ is a quotient of
$J^k/J^{k+1}\otimes\underline{O}_T$, and thus has length at most $n+1$ over
$k(s)$.  The kernel $N$ of the surjective map

\begin{equation*}
v'_k:J^n\otimes\operatorname{Gr}_J^k(\underline{O})
\simeq\underline{O}_{s^*}(n+k)
\longrightarrow
J^n\otimes\operatorname{Gr}_J^k(\underline{O}_T)
\end{equation*}
is therefore an invertible sheaf on the projective line $s^*$ of degree
$\geq-1$.  Hence $R^1p_*N=0$, and $v_k=p_*(v'_k)$ is surjective.

If $T$ is Artinian, (2.14.1) now gives

\begin{equation*}
\operatorname{length}_{k(s)}
\bigl(p_*J^n/p_*(J^nq)\bigr)
=\operatorname{length}_{k(s)}(\underline{O}_T),
\end{equation*}


\SGApage{II}{\(\approx\)24}
and 2.14(i) follows from (2.11.4) and (2.11.6).

\SGAvar{Let $I_{D'\cap T}$ be the (principal ideal of $\underline{O}_{D'}$
defining $D'\cap T$. Under the hypotheses (ii), the exact sequence
\begin{equation*}
0\longrightarrow I_{D'\cap T}\longrightarrow\underline{O}_{D'\cap T}
\longrightarrow O_T\longrightarrow0
\end{equation*}
shows that $D'\cap T$ satisfies hypothesis (i).

}{Let $I_{D'\cap T}$ be the (principal) ideal of $\underline{O}_{D'}$
defining $D'\cap T$. Under the hypotheses (ii), the exact sequence
\begin{equation*}
0\longrightarrow I_{D'\cap T}\longrightarrow\underline{O}_{D'\cup T}
\longrightarrow O_T\longrightarrow0
\end{equation*}
shows that $D'\cup T$ satisfies hypothesis (i).

}{E7II-X-021}The projection $p$ induces an isomorphism between $D$ and $D'$.  Under this
isomorphism, the ideal of $\underline{O}_D$ defining the subscheme
\[
D\cap\operatorname{Spec}\bigl(\underline{O}_S/p_*(qJ^n)\bigr)
\]
corresponds to the ideal of $D'$ generated by the global sections of $qJ^n$ on
$S_1$.  This ideal is contained in the ideal of $D'$ defining $D'\cap Y$.
It is in fact equal to it, because every section of
$\underline{O}_{D'\cup Y}$ extends to a section of $\underline{O}_{S_1}$:
if $I_{D'\cup Y}$ is the ideal defining $D'\cup Y$, the obstruction to
extension lies in $R^1p_*(I_{D'\cup Y})$.  One has

\begin{equation*}
\begin{aligned}
I_{D'\cup Y}
&=\underline{O}(-D')\cap qJ^n\\
&=(\underline{O}(-D')\cap J^n)\cap qJ^n\\
&=\underline{O}(-D')J^n\cap qJ^n\\
&=(\underline{O}(-D')\cap q)J^n,
\end{aligned}
\end{equation*}
and by (i), applied to $D'\cup T$, there is no obstruction.

Thus

\begin{equation}\tag{2.14.2}
\operatorname{length}_{k(s)}
\Bigl(D\cap
\operatorname{Spec}\bigl(\underline{O}_S/p_*(qJ^n)\bigr)\Bigr)
=\operatorname{length}_{k(s)}(D'\cap Y) \quad .
\end{equation}
The isomorphism

\begin{equation*}
J^n\otimes
\bigl(\underline{O}_{D'}/q\underline{O}_{D'}\bigr)
\xrightarrow{\;\sim\;}
\underline{O}_{D'}J^n/
\underline{O}_{D'}qJ^n
\end{equation*}
gives the exact sequence

\SGAvar{\begin{equation*}
0\longrightarrow
J^n\otimes_{D'\cap T}
\longrightarrow\underline{O}_{D'\cap Y}
\longrightarrow\underline{O}_{D'\cap ns^*}
\longrightarrow0,
\end{equation*}
and hence


}{\begin{equation*}
0\longrightarrow
J^n\otimes\underline{O}_{D'\cap T}
\longrightarrow\underline{O}_{D'\cap Y}
\longrightarrow\underline{O}_{D'\cap ns^*}
\longrightarrow0,
\end{equation*}
and hence


}{E7II-X-022}\SGApage{II}{\(\approx\)25}
\begin{equation}\tag{2.14.3}
\operatorname{length}_{k(s)}(D'\cap Y)
=n(D',s^*)+\operatorname{length}_{k(s)}(D'\cap T).
\end{equation}
Since $(D',s^*)=1$, assertion 2.14(ii) follows from (2.14.2) and (2.14.3).

We now prove, by induction on the number of elements in the support of $\nu$,
the assertions of 2.13 and the following assertion:

\SGAvar{\begin{equation}\tag{2.14.4}
\text{if $D$ is a regular divisor on $S$, then}\qquad
\chi\Bigl(D\cap\operatorname{Spec}(\underline{O}_S/I(\nu))
=\sum_{\substack{t\in S^*\\t\ \mathrm{on}\ D}}
\nu(t)[k(t):k].
\end{equation}
If $\nu=0$, then $\nu^*=0$, $I(\nu)=\underline{O}_S$, and 2.13 and (2.14.4)
are immediate (the second assertion of 2.13 follows by repeated application of
(2.11.4) with $n=0$).  Otherwise $\nu(s)>0$, and there is an $S$-morphism
$g:S'\longrightarrow S_1$.

}{\begin{equation}\tag{2.14.4}
\text{if $D$ is a regular divisor on $S$, then}\qquad
\chi\Bigl(D\cap\operatorname{Spec}(\underline{O}_S/I(\nu))\Bigr)
=\sum_{\substack{t\in S^*\\t\ \mathrm{on}\ D}}
\nu(t)[k(t):k].
\end{equation}
If $\nu=0$, then $\nu^*=0$, $I(\nu)=\underline{O}_S$, and 2.13 and (2.14.4)
are immediate (the second assertion of 2.13 follows by repeated application of
(2.11.4) with $n=0$).  Otherwise $\nu(s)>0$, and there is an $S$-morphism
$g:S'\longrightarrow S_1$.

}{E7II-X-023}For each closed point $a$ of $S_1$, let $S_{1a}$ be the localization of $S_1$
at $a$, let $\nu_a$ be the restriction of $\nu$ to $S_{1a}^*$, let
$S'_a=S'\times_{S_1}S_{1a}$, let
$I(\nu_a)=g_*\underline{O}_{S'_a}(-\nu_a^*)$, and put
\[
T_a=\operatorname{Spec}(\underline{O}_{S_1}/I(\nu_a)),\qquad
T=\coprod_aT_a,
\]
where the union is finite and disjoint.  Let $q$ be the ideal of $T$.  With the
notation of 2.14, one has

\begin{equation}\tag{2.14.5}
g_*\underline{O}_{S'}(-\nu^*)=qJ^{\nu(s)}.
\end{equation}
The induction hypothesis applies to $(S_{1a},S'_a,\nu_a)$.  Hence

\begin{equation}\tag{2.14.6}
R^ig_*\underline{O}_{S'}(-\nu^*)=0\quad\text{for }i>0,
\end{equation}
\begin{equation}\tag{2.14.7}
\chi(T)=
\sum_{t\in S^*\setminus\{s\}}
\frac{\nu(t)(\nu(t)+1)}2[k(t):k],
\end{equation}
and, by (2.14.4) applied to the divisor $s^*$ on $S_1$ and by (2.12.1),

\begin{equation*}
\operatorname{length}_{k(s)}(T\cap s^*)
=\sum_{\substack{t\ \mathrm{proximate}\\\mathrm{to}\ s}}
\nu(t)[k(t):k(s)]
\leq\nu(s).
\end{equation*}


\SGApage{II}{\(\approx\)26}
This inequality allows us to apply 2.14.  From (2.14.5) we obtain

\begin{equation*}
R^ip_*g_*\underline{O}_{S'}(-\nu^*)=0
\quad\text{for }i>0
\end{equation*}
and
\begin{equation*}
\operatorname{length}_{k(s)}
\bigl(\underline{O}_S/p_*g_*\underline{O}_{S'}(-\nu^*)\bigr)
=\frac{\nu(s)(\nu(s)+1)}2+\operatorname{length}_{k(s)}(T).
\end{equation*}
The assertions of 2.13 therefore follow from (2.14.6) and (2.14.7).  Finally,
2.14(ii) gives

\SGAvar{\begin{equation*}
\operatorname{length}_{k(s)}
\Bigl(D\cap\operatorname{Spec}(\underline{O}_S/I(\nu))
=\nu(s)+\operatorname{length}_{k(s)}(D'\cap T).
\end{equation*}
Applying (2.14.4) to $D'$ gives

\begin{equation*}
\begin{aligned}
\chi\Bigl(D\cap\operatorname{Spec}(\underline{O}_S/I(\nu))
&=\nu(s)[k(s):k]
+\sum_{{t\ \mathrm{on}\ D}}
\nu(t)[k(t):k]\\
&=\sum_{\substack{t\in S^*\\t\ \mathrm{on}\ D}}
\nu(t)[k(t):k].
\end{aligned}
\end{equation*}
This completes the proof of 2.13.


}{\begin{equation*}
\operatorname{length}_{k(s)}
\Bigl(D\cap\operatorname{Spec}(\underline{O}_S/I(\nu))\Bigr)
=\nu(s)+\operatorname{length}_{k(s)}(D'\cap T).
\end{equation*}
Applying (2.14.4) to $D'$ gives

\begin{equation*}
\begin{aligned}
\chi\Bigl(D\cap\operatorname{Spec}(\underline{O}_S/I(\nu))\Bigr)
&=\nu(s)[k(s):k]
+\sum_{\substack{t\ \mathrm{on}\ D\\t\ne s}}
\nu(t)[k(t):k]\\
&=\sum_{\substack{t\in S^*\\t\ \mathrm{on}\ D}}
\nu(t)[k(t):k].
\end{aligned}
\end{equation*}
This completes the proof of 2.13.


}{E7II-X-024}(2.15) Fragments of the preceding discussion generalize to higher dimension.  Let
$S$ be the spectrum of a regular local ring of dimension $n$, with closed point
$s$, where $k(s)$ is a finite extension of $k$.  We may again consider the set
$I$ of schemes $f:S'\longrightarrow S$ obtained from $S$ by successive blow-ups
of closed points, and put, as in 2.1,

\SGAvar{\begin{equation*}
S^*=\varinjlim_{s'\in I}\operatorname{Irr}\bigl(f^{-1}(s).
\end{equation*}
The set $S^*$ will again be called the set of \underline{points of $S$
infinitely near to $s$}.

}{\begin{equation*}
S^*=\varinjlim_{S'\in I}\operatorname{Irr}\bigl(f^{-1}(s)\bigr).
\end{equation*}
The set $S^*$ will again be called the set of \underline{points of $S$
infinitely near to $s$}.

}{E7II-X-025}If $X$ is a closed subscheme of $S$ and $t\in S^*$, the multiplicity of $X$ at
$t$ is defined as in 2.8.

Let $f:S'\longrightarrow S$ be a composite of blow-up morphisms at closed
points, let $D$ be a divisor on $S$, and let $D^*=f^*D$ be its total transform
on $S'$.


\SGApage{II}{\(\approx\)27}
Let $E$ be a curve on $S$, let $E'$ be its strict transform on $S'$, and suppose
that $D\cap E$ is supported at $s$.  The projection formula again gives

\begin{equation}\tag{2.16}
(D,E)=(D^*,E').
\end{equation}
Proceeding as in 2.9, one proves with the aid of (2.16) that

\begin{equation}\tag{2.17}
(D,E)=\sum_{t\in S^*}\mu_t(D)\mu_t(E)[k(t):k] \quad .
\end{equation}
Define proximate points as in 2.5.  Proceeding as in 2.11, one verifies that

\SGAvar{\begin{equation}\tag{2.18}
\mu_t(E)=
\sum_{\substack{t'\ \mathrm{proximate}\\\mathrm{to}\ t}}
\mu_{t'}(E)[k(t'):k] \quad .
\end{equation}
On the other hand, formula 2.11(ii) is false in general (see the examples in
2.25).

}{\begin{equation}\tag{2.18}
\mu_t(E)=
\sum_{\substack{t'\ \mathrm{proximate}\\\mathrm{to}\ t}}
\mu_{t'}(E)[k(t'):k(t)] \quad .
\end{equation}
On the other hand, formula 2.11(ii) is false in general (see the examples in
2.25).

}{E7II-X-026}2.19. Let $S$ be a regular projective surface over a field $k$.  Write $S^*$
for the disjoint sum, taken over the closed points of $S$,

\begin{equation*}
S^*=\coprod_{s\in S}\operatorname{Spec}(\underline{O}_{s,S})^*,
\end{equation*}
and call $S^*$ the set of \underline{points infinitely near to the points of}
$S$.

Let $I$ be the set of surfaces obtained from $S$ by successive blow-ups of
closed points.  The set $I$ is cofinal in the set of regular models of the
function field of $S$, ordered by domination.  For $i\in I$, let $S_i$ denote
the corresponding surface, equipped with $p_i:S_i\longrightarrow S$.  We write
$\operatorname{Div}^{\pm}(S_i)$ for the group of divisors, effective or not, on
$S_i$.

\underline{Definition} 2.20. (i) \underline{The group of ``curves'' on $S$ is
the inductive limit}

\begin{equation*}
\varinjlim_{i\in I}\operatorname{Div}^{\pm}(S_i).
\end{equation*}

(ii) \underline{The category of ``invertible sheaves'' on $S$ is the inductive
limit}


\SGApage{II}{\(\approx\)28}
(SGA 4 VI 6) \underline{of the categories of invertible sheaves on the}
$(S_i)_{i\in I}$.

These objects are birational invariants of $S$.

\underline{Example} 2.21. (i) Every ``curve'' $C$ on $S$ defines an
``invertible sheaf'' $\underline{O}(C)$ on $S$.

(ii) Every divisor $C$ on $S$ defines a ``curve'' on $S$, denoted by $C^*$.

(iii) If $t\in S^*$, define the surface $S_+(t)$ and the divisor $t^*$ on
$S_+(t)$ as in 2.3; the ``curve'' defined by this divisor will again be denoted
by $t^*$.

(2.22) One verifies easily that a ``curve'' $C$ (respectively an ``invertible
sheaf'' $\mathcal L$) can be written in one and only one way in the form

\begin{equation*}
C=C_0^*-\sum_{t\in S^*}i_t t^*
\qquad\text{respectively}\qquad
\mathcal L=\mathcal L_0\otimes\underline{O}
\left(-\sum_{t\in S^*}i_t t^*\right),
\end{equation*}
where $C_0$ (respectively $\mathcal L_0$) is a divisor (respectively an
invertible sheaf) on $S$, and where the integers $i_t$ are almost all zero.  We
say that $C$ (respectively $\mathcal L$) is the curve $C_0$ (respectively the
invertible sheaf $\mathcal L_0$) endowed with the virtual multiplicities $i_t$.

The completion
$\operatorname{Div}^{\pm}(S)\times\mathbb Z^{S^*}$ of the group
$\operatorname{Div}^{\pm}(S)\times\mathbb Z^{(S^*)}$ of ``curves'' on $S$ is
again a birational invariant of $S$; it is called the group of ``generalized
curves'' on $S$.  On $S$, a ``generalized curve'' $C$ is represented by a
divisor $C_0$ endowed with a possibly infinite collection of virtual
multiplicities.

``Generalized invertible sheaves'' on $S$ are defined in the same way.


\SGApage{II}{29}
\underline{Example} 2.23. Suppose that $k$ is perfect.  Then the invertible
sheaf $\Omega^2_{S/k}$, endowed with virtual multiplicity $-1$ at every point
$t\in S^*$, is a birational invariant of $S$; denote it by
$\Omega^2_{S/k}{}^*$.

\underline{Example} 2.24. Let $f:S_1\longrightarrow S_2$ be a dominant
morphism of regular projective surfaces over $k$.  The inverse image of a
``curve'' or an ``invertible sheaf'' on $S_2$ is defined without difficulty by
passing to a model of $S_2$ on which this ``curve'' or ``invertible sheaf'' is
an ordinary divisor or invertible sheaf.  By passage to the limit (I have not
checked whether this passage to the limit is possible), one similarly defines
the inverse image of a ``generalized curve'' or a ``generalized invertible
sheaf.''  Suppose that $k$ is perfect and that $f$ is generically \'{e}tale.
The \underline{different} $\mathfrak d(f)$ of $f$ is then the effective divisor
on $S_1$ characterized by

\begin{equation*}
f^*\Omega^2_{S_2/k}\bigl(\mathfrak d(f)\bigr)
\overset{\sim}{\longrightarrow}\Omega^2_{S_1/k} \quad .
\end{equation*}
This different is not a birational invariant of $f$; to obtain a birational
invariant one must consider the generalized coincidence curve

\begin{equation*}
\mathfrak d^*(f)=\mathfrak d(f)+\sum_{t\in S_1^*}t^*
-f^*\left(\sum_{t\in S_2^*}t^*\right),
\end{equation*}
or, briefly,

\SGAvar{\begin{equation*}
\mathfrak d^*(f)=\Omega^2_{S_1/k}{}^*
-f^*\Omega^2_{S_2/k}{}^* \quad .
\end{equation*}
This ``generalized curve'' is supported on the closed subset of $S_2$ where
$f$ is not \'{e}tale.


}{\begin{equation*}
\mathfrak d^*(f)=\Omega^2_{S_1/k}{}^*
-f^*\Omega^2_{S_2/k}{}^* \quad .
\end{equation*}
This ``generalized curve'' is supported on the closed subset of $S_1$ where
$f$ is not \'{e}tale.


}{E7II-X-027}\SGApage{II}{\(\approx\)30}
(2.25) \underline{Example}.

Let $k$ be an algebraically closed field and let $E_n$ be affine $n$-space
over $k$.  Let $C$ be a reduced curve in $E_n$ having a singularity at the
origin.  With this singularity we associate its \underline{Enriques diagram}
(cf. [3]): its vertices are the points of $C$ infinitely near the origin
(2.4 and 2.15), and each vertex is joined by a line to the vertices proximate
to it (2.5 and 2.15).  This diagram determines the multiplicities of these
points on $C$ (2.11 and 2.18) and, when $n=2$, the invariant $\delta$ (2.10
and 2.11) of the singularity of $C$ at the origin.

For the simplest plane singularities one obtains the following table.  A
number next to a vertex is its multiplicity.

\begin{center}
\small
\setlength{\tabcolsep}{5pt}
\renewcommand{\arraystretch}{1.35}
\begin{tabular}{c c c c}
Sketch & Equation & $\delta$ & Enriques diagram \\

\begin{tikzpicture}[baseline=-.5ex,x=.55cm,y=.55cm]
  \draw[line width=.45pt] (-.8,0)--(.8,0);
\end{tikzpicture}
& $y=0$ & $0$ &
\begin{tikzpicture}[baseline=-.5ex,x=.56cm,y=.56cm]
  \foreach \x in {0,...,6}{
    \fill (\x,0) circle (1.05pt);
    \node[above=1pt,scale=.72] at (\x,0) {$1$};
  }
  \draw (0,0)--(6,0);
  \node[right=2pt,scale=.85] at (6,0) {$\cdots$};
\end{tikzpicture}
\\

\begin{tikzpicture}[baseline=-.5ex,x=.55cm,y=.55cm]
  \draw[line width=.45pt] (-.75,-.42)--(.75,.42);
  \draw[line width=.45pt] (-.75,.42)--(.75,-.42);
\end{tikzpicture}
& $xy=0$ & $1$ &
\begin{tikzpicture}[baseline=-.5ex,x=.72cm,y=.62cm]
  \fill (0,0) circle (1.05pt);
  \node[left=1pt,scale=.72] at (0,0) {$2$};
  \foreach \x in {1,...,5}{
    \fill (\x,.48) circle (1.05pt);
    \fill (\x,-.48) circle (1.05pt);
    \node[above=1pt,scale=.72] at (\x,.48) {$1$};
    \node[above=1pt,scale=.72] at (\x,-.48) {$1$};
  }
  \draw (0,0)--(1,.48)--(5,.48);
  \draw (0,0)--(1,-.48)--(5,-.48);
  \node[right=2pt,scale=.85] at (5,.48) {$\cdots$};
  \node[right=2pt,scale=.85] at (5,-.48) {$\cdots$};
\end{tikzpicture}
\\

\begin{tikzpicture}[baseline=-.5ex,x=.55cm,y=.55cm]
  \draw[line width=.45pt] (-.85,0)--(.85,0);
  \draw[line width=.45pt,domain=-.72:.72,samples=40]
    plot (\x,{.72*\x*\x});
\end{tikzpicture}
& $y(y-x^2)=0$ & $2$ &
\begin{tikzpicture}[baseline=-.5ex,x=.68cm,y=.62cm]
  \fill (0,0) circle (1.05pt);
  \fill (.9,0) circle (1.05pt);
  \node[above=1pt,scale=.72] at (0,0) {$2$};
  \node[above=1pt,scale=.72] at (.9,0) {$2$};
  \foreach \x/\lab in {1.8/1,2.7/1,3.6/1,4.5/0}{
    \fill (\x,.48) circle (1.05pt);
    \fill (\x,-.48) circle (1.05pt);
    \ifnum\lab=1
      \node[above=1pt,scale=.72] at (\x,.48) {$1$};
      \node[below=1pt,scale=.72] at (\x,-.48) {$1$};
    \fi
  }
  \draw (0,0)--(.9,0)--(1.8,.48)--(4.5,.48);
  \draw (.9,0)--(1.8,-.48)--(4.5,-.48);
  \node[right=2pt,scale=.72] at (4.5,-.48) {$\cdot\ \cdot$};
\end{tikzpicture}
\\

\begin{tikzpicture}[baseline=-.5ex,x=.55cm,y=.55cm]
  \draw[line width=.45pt]
    (0,0) .. controls (.28,.05) and (.55,.35) .. (.72,.78);
  \draw[line width=.45pt]
    (0,0) .. controls (.28,-.05) and (.55,-.35) .. (.72,-.78);
\end{tikzpicture}
& $y^2=x^3$ & $1$ &
\begin{tikzpicture}[baseline=-.5ex,x=.72cm,y=.65cm]
  \coordinate (A) at (0,.42);
  \coordinate (B) at (.9,.42);
  \coordinate (C) at (1.7,-.25);
  \coordinate (D) at (2.6,-.25);
  \coordinate (E) at (3.5,-.25);
  \foreach \p in {A,B,C,D,E}{\fill (\p) circle (1.05pt);}
  \node[above=1pt,scale=.72] at (A) {$2$};
  \foreach \p in {B,C,D,E}{\node[above=1pt,scale=.72] at (\p) {$1$};}
  \draw (A)--(B)--(C)--(D)--(E);
  \draw (A)--(C);
  \node[right=2pt,scale=.85] at (E) {$\cdots$};
\end{tikzpicture}
\end{tabular}
\end{center}

If $n>2$, the invariant $\delta$ is no longer determined by this diagram.
For example, the union of three concurrent coplanar lines (respectively,
noncoplanar lines) has invariant $\delta=3$ (respectively, $\delta=2$),
although the Enriques diagram is the same.  I do not know whether this
phenomenon also occurs for unibranch curves.

Here are two examples of diagrams in the case $n=3$ which cannot occur for
plane curves (cf. 2.5).  We give, in order, the equations of the curve, the
equations of the pure transform after blowing up the origin, parametric
equations of the curve, the invariant $\delta$, and the Enriques diagram.


\SGApage{II}{\(\approx\)31}
\SGAvar{\begin{equation}\tag{2.25.1}
\left\{
\begin{gathered}
  \left\{\begin{aligned}y^2&=xz^4,\\ z^3&=x^2,
  \end{aligned}\right.\qquad
  \left\{\begin{aligned}(y/z)^2&=(x/z)^3,\\ z&=(x/z)^2,
  \end{aligned}\right.\\[3pt]
  \left\{\begin{aligned}x&=u^6,\\y&=u^7,\\z&=u^4,
  \end{aligned}\right.\qquad \delta=5
\end{gathered}
\right.
\end{equation}

}{\begin{equation}\tag{2.25.1}
\left\{
\begin{gathered}
  \left\{\begin{aligned}y^2&=xz^2,\\ z^3&=x^2,
  \end{aligned}\right.\qquad
  \left\{\begin{aligned}(y/z)^2&=(x/z)^3,\\ z&=(x/z)^2,
  \end{aligned}\right.\\[3pt]
  \left\{\begin{aligned}x&=u^6,\\y&=u^7,\\z&=u^4,
  \end{aligned}\right.\qquad \delta=5
\end{gathered}
\right.
\end{equation}

}{E7II-X-028}\begin{center}
\begin{tikzpicture}[x=.78cm,y=.72cm]
  \coordinate (v1) at (0,0);
  \coordinate (v2) at (1,0);
  \coordinate (v3) at (1.75,-.85);
  \coordinate (v4) at (1.75,-1.70);
  \coordinate (v5) at (2.65,-2.35);
  \coordinate (v6) at (3.65,-2.62);
  \foreach \p in {v1,v2,v3,v4,v5,v6}{\fill (\p) circle (1.15pt);}
  \node[above left=1pt] at (v1) {$4$};
  \node[above=1pt] at (v2) {$2$};
  \node[right=2pt] at (v3) {$1$};
  \node[right=2pt] at (v4) {$1$};
  \node[above right=1pt] at (v5) {$1$};
  \node[above right=1pt] at (v6) {$1$};
  \draw (v1)--(v2) (v1)--(v3) (v1)--(v4);
  \draw (v2)--(v3) (v2)--(v4) (v3)--(v4);
  \draw (v4)--(v5)--(v6);
  \fill (4.08,-2.86) circle (.7pt);
  \fill (4.38,-3.05) circle (.7pt);
  \fill (4.68,-3.24) circle (.7pt);
\end{tikzpicture}
\end{center}

\begin{equation}\tag{2.25.2}
\left\{
\begin{gathered}
  \left\{\begin{aligned}y^2z&=x^3,\\ z^3&=y^2,
  \end{aligned}\right.\qquad
  \left\{\begin{aligned}(y/z)^2&=(x/z)^3,\\ z&=(y/z)^2,
  \end{aligned}\right.\\[3pt]
  \left\{\begin{aligned}x&=u^8,\\y&=u^9,\\z&=u^6,
  \end{aligned}\right.\qquad \delta=10
\end{gathered}
\right.
\end{equation}

\begin{center}
\begin{tikzpicture}[x=.72cm,y=.65cm]
  \coordinate (w1) at (0,0);
  \coordinate (w2) at (1,0);
  \coordinate (w3) at (1.75,-.72);
  \coordinate (w4) at (1.75,-1.44);
  \coordinate (w5) at (1.75,-2.16);
  \coordinate (w6) at (1.75,-2.88);
  \coordinate (w7) at (2.55,-3.70);
  \coordinate (w8) at (3.35,-4.52);
  \foreach \p in {w1,w2,w3,w4,w5,w6,w7,w8}{\fill (\p) circle (1.15pt);}
  \node[above left=1pt] at (w1) {$6$};
  \node[above=1pt] at (w2) {$2$};
  \foreach \p in {w3,w4,w5,w6}{\node[right=2pt] at (\p) {$1$};}
  \node[above right=1pt] at (w7) {$1$};
  \node[above right=1pt] at (w8) {$1$};
  \draw (w1)--(w2);
  \foreach \p in {w3,w4,w5,w6}{\draw (w1)--(\p);}
  \draw (w2)--(w3) (w2)--(w4);
  \draw (w3)--(w4)--(w5)--(w6)--(w7)--(w8);
  \fill (3.72,-4.86) circle (.7pt);
  \fill (4.02,-5.10) circle (.7pt);
  \fill (4.32,-5.34) circle (.7pt);
\end{tikzpicture}
\end{center}


2.26. Theorem 2.27 below was stated as a false conjecture in the first version
of this exposé (circulated by the IHES).  It owes much to conversations with
D. S. Rim and includes some of his results (D. S. Rim, ``Torsion
Differentials and Deformations,'' Theorem 2.7(i)).


Let $C$ be a reduced projective curve over an algebraically closed field $k$,
and let $s\in C(k)$.  We shall write

\SGApage{II}{\(\approx\)32}
\begin{description}[leftmargin=2.7cm,style=nextline]
\item[$A_s$] for the local ring of $C$ at $s$ (or its completion, which would
make no difference below);
\item[$\widetilde A_s$] for the normalization of $A_s$;
\item[$\delta(s)$] for $\dim_k(\widetilde A_s/A_s)$;
\item[$m_0(s)$] for the dimension of the formal versal scheme $T$ of
deformations of $\operatorname{Spec}(A_s)$.
\end{description}

For an irreducible component $E$ of $T$, put likewise
$m_0^E(s)=\dim(E)$.

\SGAvar{The module $\widetilde\tau$ of derivations of $\widetilde A_s$ and the module
$\tau$ of derivations of $A_s$ are torsion-free.  If $A_s$ is integral, they
embed in the module of derivations of the fraction field of $A$; in general,
they embed in the module of derivations of the total quotient ring of $A_s$
(a product of fields).  Moreover, through these embeddings, $\tau$ and
$\widetilde\tau$ are commensurable: both
$\tau/(\tau\cap\widetilde\tau)$ and
$\widetilde\tau/(\tau\cap\widetilde\tau)$ are finite-dimensional over $k$.
Set

}{The module $\widetilde\tau$ of derivations of $\widetilde A_s$ and the module
$\tau$ of derivations of $A_s$ are torsion-free.  If $A_s$ is integral, they
embed in the module of derivations of the fraction field of $A_s$; in general,
they embed in the module of derivations of the total quotient ring of $A_s$
(a product of fields).  Moreover, through these embeddings, $\tau$ and
$\widetilde\tau$ are commensurable: both
$\tau/(\tau\cap\widetilde\tau)$ and
$\widetilde\tau/(\tau\cap\widetilde\tau)$ are finite-dimensional over $k$.
Set

}{E7II-X-029}\begin{equation*}
m_1(s)=[\widetilde\tau:\tau]
\mathrel{\mathop=\limits_{\mathrm{def}}}
\dim_k\frac{\widetilde\tau}{\tau\cap\widetilde\tau}
-\dim_k\frac{\tau}{\tau\cap\widetilde\tau}.
\end{equation*}

\underline{Theorem} 2.27. \underline{Let $E$ be an irreducible component of
the formal versal scheme $T$ of deformations of $A_s$. If the deformation of
$A_s$ over the generic point $e$ of $E$ is smooth---that is, if the image of
the formal non-smoothness subscheme (defined by a Jacobian ideal: VI 5), which
is finite over $T$, does not contain $e$---then}

\begin{equation*}
3\delta(s)=m_0^E(s)+m_1(s).
\end{equation*}


\SGApage{II}{\(\approx\)33}
The validity of 2.27 depends only on the completion of $A_s$.  Replacing
$(C,s)$ by a pair $(C',s)$ formally isomorphic to it at $s$, we may suppose
that the following conditions hold:

\begin{enumerate}[label=(\alph*),leftmargin=1.2cm]
\item $s$ is the only singular point of $C$;
\item $C$ has no infinitesimal automorphisms.
\end{enumerate}

Condition (b) ensures that the deformation functor of $C$ is pro-representable.
It is even effectively pro-representable, since there is no obstruction to
lifting a prescribed ample invertible sheaf on $C$.  Let
$f:X_0\longrightarrow T_0$ be the universal deformation of $C$, and let $t_0$
be the closed point of $T_0$, so that $C\overset\sim\longrightarrow X_{0t_0}$.

\underline{Lemma} 2.28. \underline{There exist a scheme $T_1$ of finite type
over $k$, a point $t_1\in T_1(k)$, and a proper flat morphism
$f:X_1\longrightarrow T_1$, whose fiber at $t_1$ is $C$, such that for every
$u\in T_1(k)$ the completion of $T_1$ at $u$ is a formal moduli scheme for the
fiber $X_{1u}$.}

\SGAvar{Here is a direct proof of this corollary of general results of Artin.  Let
$\mathcal L$ be a very ample invertible sheaf on $C$ such that
$H^1(C,\mathcal L)=0$, and put
$\mathbb P=\mathbb P(H^o(C,\mathcal L))$.  We have an embedding
$v:C\hookrightarrow\mathbb P$.  Let $T_2$ be the Hilbert scheme parametrizing
the subvarieties of $\mathbb P$ having the same Hilbert polynomial as
$v(C)\subset\mathbb P$, and let $t_2\in T_2$ correspond to $v(C)$.  The
completion $\widehat T_2$ of $T_2$ at $t_2$ is smooth over $S_0$.  The vertical
tangent bundle of $\widehat T_2/T_0$ is induced by the subbundle $L$ of the
tangent bundle of $T_2$ which is the image of
$\operatorname{Lie}(\underline{\operatorname{Aut}}(\mathbb P))$ under the
action of $\underline{\operatorname{Aut}}(\mathbb P)$ on $T_2$.

In a neighborhood of $t_2$, let $f_1\ldots f_n$, with
$n=\dim\underline{\operatorname{Aut}}(\mathbb P)$, be equations defining a
subscheme $T'_1$ of $T_2$ which passes through $t_2$ and whose completion at
$t_2$ is étale over $T_0$.  Take for $T_1$ a neighborhood of $t$ in $T'_1$
on which the tangent bundle is transverse to $L$.


}{Here is a direct proof of this corollary of general results of Artin.  Let
$\mathcal L$ be a very ample invertible sheaf on $C$ such that
$H^1(C,\mathcal L)=0$, and put
$\mathbb P=\mathbb P(H^0(C,\mathcal L))$.  We have an embedding
$v:C\hookrightarrow\mathbb P$.  Let $T_2$ be the Hilbert scheme parametrizing
the subvarieties of $\mathbb P$ having the same Hilbert polynomial as
$v(C)\subset\mathbb P$, and let $t_2\in T_2$ correspond to $v(C)$.  The
completion $\widehat T_2$ of $T_2$ at $t_2$ is smooth over $T_0$.  The vertical
tangent bundle of $\widehat T_2/T_0$ is induced by the subbundle $L$ of the
tangent bundle of $T_2$ which is the image of
$\operatorname{Lie}(\underline{\operatorname{Aut}}(\mathbb P))$ under the
action of $\underline{\operatorname{Aut}}(\mathbb P)$ on $T_2$.

In a neighborhood of $t_2$, let $f_1,\ldots,f_n$, with
$n=\dim\underline{\operatorname{Aut}}(\mathbb P)$, be equations defining a
subscheme $T'_1$ of $T_2$ which passes through $t_2$ and whose completion at
$t_2$ is étale over $T_0$.  Take for $T_1$ a neighborhood of $t_2$ in $T'_1$
on which the tangent bundle is transverse to $L$.


}{E7II-X-030}\SGApage{II}{34}
2.29. The morphism of functors

\begin{equation*}
(\text{deformations of }C)\longrightarrow(\text{deformations of }A_s)
\end{equation*}

corresponds to a morphism $p:T_0\longrightarrow T$.  By VI 2.10, $p$ is
smooth.  Let $t$ be the closed point of $T$.  The Zariski tangent space of
$p^{-1}(t)$ at $t_0$ classifies the locally trivial deformations of $C$ over
$k[\varepsilon]$, where $\varepsilon^2=0$.  Putting $E_0=p^{-1}(E)$, we obtain

\begin{equation}\tag{2.29.1}
\dim(E_0)=\dim(E)+\dim H^1(C,T^1_{C/k}).
\end{equation}

Since $X_0/T_0$ is flat, all its fibers have the same arithmetic genus $p$,
where $1-p=\chi(\underline O)$.  By hypothesis, the fiber at the generic point
of $E_0$ is smooth.  The formal moduli scheme of a smooth proper curve being
smooth of dimension $3p-3$, Lemma 2.28 gives

\SGAvar{\begin{equation}\tag{2.29.2}
\dim(E)=3p-3.
\end{equation}

}{\begin{equation}\tag{2.29.2}
\dim(E_0)=3p-3.
\end{equation}

}{E7II-X-031}Let $\widetilde C$ be the normalization of $C$, of arithmetic genus
$\widetilde p=p-\delta(s)$.  We have

\SGAvar{\begin{equation}\tag{2.29.3}
\begin{aligned}
\dim H^1(C,T^1_{C/k})
&=-\chi(C,T^1_{C/k})\\
&=-\chi(\widetilde C,T^1_{\widetilde C/k})+m_1(s)\\
&=-[3(p-\delta(s))-3]+m_1(s).
\end{aligned}
\end{equation}

}{\begin{equation}\tag{2.29.3}
\begin{aligned}
\dim H^1(C,T^1_{C/k})
&=-\chi(C,T^1_{C/k})\\
&=-\chi(\widetilde C,T^1_{\widetilde C/k})+m_1(s)\\
&=[3(p-\delta(s))-3]+m_1(s).
\end{aligned}
\end{equation}

}{E7II-X-032}Formulas (2.29.1), (2.29.2), and (2.29.3) prove 2.27.

\underline{Conjecture} 2.30. \underline{Every reduced proper curve has a
smooth deformation.}

This conjecture implies, by 2.27 and 2.28, that the formal versal scheme of
deformations of $A_s$ is always equidimensional, of dimension
$3\delta(s)-m_1(s)$.



\SGApage{II}{35}
2.31. Let $C$ be a reduced proper curve over $k$, of arithmetic genus $p$.
We shall write

\begin{description}[leftmargin=2.7cm,style=nextline]
\item[$M_0$] for the dimension of the formal versal scheme $T$ of deformations
of $C$;
\item[$M_1$] for the dimension of the Lie algebra of
$\underline{\operatorname{Aut}}(C)$.
\end{description}

For an irreducible component $E$ of $T$, write $M_0^E$ for the dimension of
$E$.  The global analogue of 2.27 is the following.

\underline{Proposition} 2.32. \underline{If the deformation of $C$ over the
generic point of $E$ is smooth, then}

\begin{equation*}
3p-3=M_0^E-M_1.
\end{equation*}

Let $X\subset C(k)$ be a finite set of smooth points of $C$ such that every
vector field on $C$ which vanishes at every point $x\in X$ is zero.  The
deformation functor of $(C,X)$ is pro-representable and satisfies a variant of
2.28.  The formal moduli scheme of $(C,X)$ is smooth over $T$, of relative
dimension $\#X-M_1$, and the analogue of 2.28 gives

\SGAvar{\begin{equation*}
M_0+(\#X-M_1)=3p-3+\#X,
\end{equation*}

}{\begin{equation*}
M_0^E+(\#X-M_1)=3p-3+\#X,
\end{equation*}

}{E7II-X-033}which proves 2.32.

Suppose now that $k$ has \underline{characteristic zero}.  Then
$M_1=\dim\underline{\operatorname{Aut}}(C)$.  Here is an analogous
interpretation of $m_1(s)$.


\SGApage{II}{36}
\underline{Lemma} 2.33. \underline{If $k$ has characteristic} zero, \underline{every
derivation of $A_s$ extends to a derivation of $\widetilde A_s$.}

If $D$ is a derivation of $A_s$, then

\SGAvar{\begin{equation*}
\exp(Dt)=\sum\frac{D^nt^n}{n!}
\end{equation*}

}{\begin{equation*}
\exp(Dt)=\sum_{n\geq0}\frac{D^nt^n}{n!}
\end{equation*}

}{E7II-X-034}is an automorphism of the completion at $(s,0)$ of
$A_s\otimes_k k[t]$.  By transport of structure, this automorphism acts on the
normalization $(\widetilde A_s\otimes_k k[t])^{\wedge}$ of that ring.  Reducing
modulo $t^2$, we obtain an automorphism $1+Dt$ of
$\widetilde A_s^{\wedge}\otimes_k k[t]/(t^2)$, and $D$ thereby extends to a
derivation of $\widetilde A_s^{\wedge}$.  Since it also extends to the fraction
field of $A_s$ (to the total quotient ring when $A_s$ is not integral), it
extends to $\widetilde A_s$.

Suppose that $s$ is singular.  A similar argument then shows that the vector
field on $\operatorname{Spec}(\widetilde A_s)$ defined by $D$ vanishes at every
point above $s$.

Let $c$ be an ideal of $\widetilde A_s$ contained in the conductor of
$\widetilde A_s/A_s$.  Put

\begin{equation*}
\widetilde G_c=\underline{\operatorname{Aut}}(\widetilde A_s/c),
\end{equation*}

and let $G_c$ be the subgroup of $\widetilde G_c$ which stabilizes the
subalgebra $A_s/c$.  If $d\subset c$, then $\widetilde G_d$ stabilizes $c/d$
and maps onto $\widetilde G_c$.  Notice that
$g\in\widetilde G_c$ is determined by $g(t)$ for a uniformizer $t$, and that
$g(t)$ may be any uniformizer.  Consequently the pointed homogeneous space
$\widetilde G_c/G_c$ is independent of $c$; denote it by $\widetilde G/G$.

Let $\widetilde\tau_0$ be the $\widetilde A_s$-module of vector fields on
$\operatorname{Spec}(\widetilde A_s)$ which vanish at the points above $s$.
When $s$ is singular, one has

\begin{equation*}
\dim_k(\widetilde\tau/\widetilde\tau_0)
=\text{the number of branches},
\end{equation*}

and

\begin{equation*}
\widetilde\tau_0/\tau\simeq
\text{the tangent space at the origin of }\widetilde G/G.
\end{equation*}


\SGApage{II}{37}
\underline{Proposition} 2.34. \underline{In characteristic zero, when $s$ is
singular, the integer $m_1(s)$ is the sum of the number of branches and the
dimension of $\widetilde G/G$.}

\bigskip
\begin{center}
\underline{\textsc{Bibliography}}
\end{center}

\SGAvar{\begin{description}[leftmargin=1.1cm,labelwidth=.8cm,labelsep=.3cm]
\item[{[1]}] P. Defrise, \textit{Étude locale des correspondances rationnelles
entre surfaces algébriques}, thèse d'agrégation, Brussels, 1949.

\item[{[2]}] P. Deligne and D. Mumford, ``The irreducibility of the space of
curves of a given genus,'' \textit{Publ. Math. IHES} \underline{36} (1969),
75--110.

\item[{[3]}] Enriques--Chisini, \textit{Lezioni sulla teoria geometrica delle
equazioni e delle funzioni algebriche}, Zanichelli, Bologna, 1924.

\item[{[4]}] R. Hartshorne, \textit{Residues and Duality}, Lecture Notes in
Mathematics \underline{20}, Springer, 1966.

\item[{[5]}] S. Lichtenbaum, ``Curves over discrete valuation rings,''
\textit{Amer. J. Math.} 15, no. 2 (April 1968), 380--405.

\item[{[6]}] I. Manin, ``Rational surfaces over perfect fields,''
\textit{Publ. Math. IHES}, no. 30, 55--114.

\item[{[7]}] D. Mumford, ``The topology of normal singularities of an
algebraic surface and a criterion for simplicity,'' \textit{Publ. Math. IHES},
no. 9, 5--22.

\item[{[8]}] A. Néron, \textit{Modèles minimaux des variétés abéliennes sur
les corps locaux et globaux}, Publ. Math., no. 21.


\item[{[9]}] M. Raynaud, ``Spécialisation du foncteur de Picard. Critères
numériques de représentabilité,'' \textit{C. R. Acad. Sci. Paris} 264
(June 1967), 1001--1004.

\item[{[10]}] Saint-Donat, \textit{Sur un théorème de G. Castelnuovo et
F. Enriques}, thèse de troisième cycle, Lyon.

\item[{[11]}] O. Zariski, ``On polynomial ideals defined by infinitely near
base points,'' \textit{Amer. J. Math.} 60 (1938), 151.

\item[{[12]}] D. G. Northcott, ``On the genus formula for plane curves,''
\textit{J. London Math. Soc.} 30 (1955), 376--382.

\item[{[13]}] D. G. Northcott, ``Abstract dilatations and infinitely near
points,'' \textit{Proc. Cambridge Philos. Soc.} 52 (1956), 178--197.

\item[] The notion of an ``infinitely near point'' is discussed in Lipman's
appendix to Chapter II of the following work.

\item[{[14]}] O. Zariski, \textit{Algebraic Surfaces}, Ergebnisse
\underline{61}, Springer-Verlag.
\end{description}


}{\begin{description}[leftmargin=1.1cm,labelwidth=.8cm,labelsep=.3cm]
\item[{[1]}] P. Defrise, \textit{Étude locale des correspondances rationnelles
entre surfaces algébriques}, thèse d'agrégation, Brussels, 1949.

\item[{[2]}] P. Deligne and D. Mumford, ``The irreducibility of the space of
curves of a given genus,'' \textit{Publ. Math. IHES} \underline{36} (1969),
75--110.

\item[{[3]}] Enriques--Chisini, \textit{Lezioni sulla teoria geometrica delle
equazioni e delle funzioni algebriche}, Zanichelli, Bologna, 1924.

\item[{[4]}] R. Hartshorne, \textit{Residues and Duality}, Lecture Notes in
Mathematics \underline{20}, Springer, 1966.

\item[{[5]}] S. Lichtenbaum, ``Curves over discrete valuation rings,''
\textit{Amer. J. Math.} 90, no. 2 (April 1968), 380--405.

\item[{[6]}] I. Manin, ``Rational surfaces over perfect fields,''
\textit{Publ. Math. IHES}, no. 30, 55--114.

\item[{[7]}] D. Mumford, ``The topology of normal singularities of an
algebraic surface and a criterion for simplicity,'' \textit{Publ. Math. IHES},
no. 9, 5--22.

\item[{[8]}] A. Néron, \textit{Modèles minimaux des variétés abéliennes sur
les corps locaux et globaux}, Publ. Math., no. 21.


\item[{[9]}] M. Raynaud, ``Spécialisation du foncteur de Picard. Critères
numériques de représentabilité,'' \textit{C. R. Acad. Sci. Paris} 264
(June 1967), 1001--1004.

\item[{[10]}] Saint-Donat, \textit{Sur un théorème de G. Castelnuovo et
F. Enriques}, thèse de troisième cycle, Lyon.

\item[{[11]}] O. Zariski, ``On polynomial ideals defined by infinitely near
base points,'' \textit{Amer. J. Math.} 60 (1938), 151.

\item[{[12]}] D. G. Northcott, ``On the genus formula for plane curves,''
\textit{J. London Math. Soc.} 30 (1955), 376--382.

\item[{[13]}] D. G. Northcott, ``Abstract dilatations and infinitely near
points,'' \textit{Proc. Cambridge Philos. Soc.} 52 (1956), 178--197.

\item[] The notion of an ``infinitely near point'' is discussed in Lipman's
appendix to Chapter II of the following work.

\item[{[14]}] O. Zariski, \textit{Algebraic Surfaces}, Ergebnisse
\underline{61}, Springer-Verlag.
\end{description}


}{E7II-X-035}\markboth{Exposé X}{Exposé X}

\ifnum\SGAform=1
\clearpage\section*{Critical apparatus}\begingroup\small\raggedright\sloppy
Readings are given in the exact TeX notation of the two English witnesses. Decimal suffixes identify the individual differences in each marked passage.

\par\medskip\noindent\hypertarget{app-E7II-X-001.1}{}\textbf{E7II-X-001.1}\quad printed folio 2.\par
\noindent Printed reading: \emph{(empty)}\par
\noindent Corrected reading: {\ttfamily\footnotesize \textbackslash{}underline\{\allowbreak{}}\par
\noindent Source: {\ttfamily\footnotesize ED-EN-FONT-G015;\allowbreak{} F7II-X-002.2\allowbreak{} F7II-X-002.4;\allowbreak{} fr/sga7ii/\{diplomati\allowbreak{}c,current\}/expose\_X\_\allowbreak{}body.tex;\allowbreak{} folio\allowbreak{} 2\allowbreak{}}\par

\par\medskip\noindent\hypertarget{app-E7II-X-001.2}{}\textbf{E7II-X-001.2}\quad printed folio 2.\par
\noindent Printed reading: \emph{(empty)}\par
\noindent Corrected reading: {\ttfamily\footnotesize \}\allowbreak{}}\par
\noindent Source: {\ttfamily\footnotesize ED-EN-FONT-G015;\allowbreak{} F7II-X-002.2\allowbreak{} F7II-X-002.4;\allowbreak{} fr/sga7ii/\{diplomati\allowbreak{}c,current\}/expose\_X\_\allowbreak{}body.tex;\allowbreak{} folio\allowbreak{} 2\allowbreak{}}\par

\par\medskip\noindent\hypertarget{app-E7II-X-001.3}{}\textbf{E7II-X-001.3}\quad printed folio 2.\par
\noindent Printed reading: {\ttfamily\footnotesize \textbackslash{}mathbb\allowbreak{} }\par
\noindent Corrected reading: \emph{(empty)}\par
\noindent Source: {\ttfamily\footnotesize ED-EN-FONT-G015;\allowbreak{} F7II-X-002.2\allowbreak{} F7II-X-002.4;\allowbreak{} fr/sga7ii/\{diplomati\allowbreak{}c,current\}/expose\_X\_\allowbreak{}body.tex;\allowbreak{} folio\allowbreak{} 2\allowbreak{}}\par

\par\medskip\noindent\hypertarget{app-E7II-X-002.1}{}\textbf{E7II-X-002.1}\quad printed folio 5.\par
\noindent Printed reading: {\ttfamily\footnotesize iD\allowbreak{}}\par
\noindent Corrected reading: {\ttfamily\footnotesize i,D\allowbreak{}}\par
\noindent Source: {\ttfamily\footnotesize fr/sga7ii/\{diplomati\allowbreak{}c,current\}/expose\_X\_\allowbreak{}body.tex;\allowbreak{} folio\allowbreak{} 5;\allowbreak{} formula\allowbreak{} 112\allowbreak{}}\par

\par\medskip\noindent\hypertarget{app-E7II-X-003.1}{}\textbf{E7II-X-003.1}\quad printed folio 6.\par
\noindent Printed reading: {\ttfamily\footnotesize S\allowbreak{}}\par
\noindent Corrected reading: {\ttfamily\footnotesize X\allowbreak{}}\par
\noindent Source: {\ttfamily\footnotesize fr/sga7ii/\{diplomati\allowbreak{}c,current\}/expose\_X\_\allowbreak{}body.tex;\allowbreak{} folio\allowbreak{} 6;\allowbreak{} formula\allowbreak{} 130\allowbreak{}}\par

\par\medskip\noindent\hypertarget{app-E7II-X-004.1}{}\textbf{E7II-X-004.1}\quad printed folio 7.\par
\noindent Printed reading: {\ttfamily\footnotesize D\allowbreak{}}\par
\noindent Corrected reading: {\ttfamily\footnotesize S\allowbreak{}}\par
\noindent Source: {\ttfamily\footnotesize fr/sga7ii/\{diplomati\allowbreak{}c,current\}/expose\_X\_\allowbreak{}body.tex;\allowbreak{} folio\allowbreak{} 7;\allowbreak{} formula\allowbreak{} 146\allowbreak{}}\par

\par\medskip\noindent\hypertarget{app-E7II-X-005.1}{}\textbf{E7II-X-005.1}\quad printed folio 8.\par
\noindent Printed reading: \emph{(empty)}\par
\noindent Corrected reading: {\ttfamily\footnotesize |\_E)\allowbreak{}}\par
\noindent Source: {\ttfamily\footnotesize fr/sga7ii/\{diplomati\allowbreak{}c,current\}/expose\_X\_\allowbreak{}body.tex;\allowbreak{} folio\allowbreak{} 8\allowbreak{}}\par

\par\medskip\noindent\hypertarget{app-E7II-X-006.1}{}\textbf{E7II-X-006.1}\quad printed folio 9.\par
\noindent Printed reading: {\ttfamily\footnotesize  \textbackslash{}underline\{in\}\allowbreak{} \$X\_1\$\allowbreak{}}\par
\noindent Corrected reading: \emph{(empty)}\par
\noindent Source: {\ttfamily\footnotesize ED-EN-FORMULA-G079;\allowbreak{} fr/sga7ii/\{diplomati\allowbreak{}c,current\}/expose\_X\_\allowbreak{}body.tex;\allowbreak{} folio\allowbreak{} 9\allowbreak{}}\par

\par\medskip\noindent\hypertarget{app-E7II-X-007.1}{}\textbf{E7II-X-007.1}\quad printed folio 11.\par
\noindent Printed reading: {\ttfamily\footnotesize i\allowbreak{}}\par
\noindent Corrected reading: {\ttfamily\footnotesize j\allowbreak{}}\par
\noindent Source: {\ttfamily\footnotesize fr/sga7ii/\{diplomati\allowbreak{}c,current\}/expose\_X\_\allowbreak{}body.tex;\allowbreak{} folio\allowbreak{} 11;\allowbreak{} formula\allowbreak{} 267\allowbreak{}}\par

\par\medskip\noindent\hypertarget{app-E7II-X-008.1}{}\textbf{E7II-X-008.1}\quad printed folio 14.\par
\noindent Printed reading: {\ttfamily\footnotesize [\allowbreak{}}\par
\noindent Corrected reading: {\ttfamily\footnotesize ]\allowbreak{}}\par
\noindent Source: {\ttfamily\footnotesize ED-en7II-C508;\allowbreak{} OPEN\_ITEMS\_DECIDED.t\allowbreak{}sv\allowbreak{}}\par

\par\medskip\noindent\hypertarget{app-E7II-X-009.1}{}\textbf{E7II-X-009.1}\quad printed folio 15.\par
\noindent Printed reading: {\ttfamily\footnotesize o\allowbreak{}}\par
\noindent Corrected reading: {\ttfamily\footnotesize 0\allowbreak{}}\par
\noindent Source: {\ttfamily\footnotesize fr/sga7ii/\{diplomati\allowbreak{}c,current\}/expose\_X\_\allowbreak{}body.tex;\allowbreak{} folio\allowbreak{} 15;\allowbreak{} formula\allowbreak{} 448\allowbreak{}}\par

\par\medskip\noindent\hypertarget{app-E7II-X-009.2}{}\textbf{E7II-X-009.2}\quad printed folio 15.\par
\noindent Printed reading: {\ttfamily\footnotesize o\allowbreak{}}\par
\noindent Corrected reading: {\ttfamily\footnotesize 0\allowbreak{}}\par
\noindent Source: {\ttfamily\footnotesize fr/sga7ii/\{diplomati\allowbreak{}c,current\}/expose\_X\_\allowbreak{}body.tex;\allowbreak{} folio\allowbreak{} 15;\allowbreak{} formula\allowbreak{} 450\allowbreak{}}\par

\par\medskip\noindent\hypertarget{app-E7II-X-009.3}{}\textbf{E7II-X-009.3}\quad printed folio 15.\par
\noindent Printed reading: {\ttfamily\footnotesize \textbackslash{}leq\allowbreak{} }\par
\noindent Corrected reading: {\ttfamily\footnotesize <\allowbreak{}}\par
\noindent Source: {\ttfamily\footnotesize fr/sga7ii/\{diplomati\allowbreak{}c,current\}/expose\_X\_\allowbreak{}body.tex;\allowbreak{} folio\allowbreak{} 15;\allowbreak{} formula\allowbreak{} 454;\allowbreak{} port/PORT\_EN7II\_LOG.\allowbreak{}jsonl\#C509;\allowbreak{} fixes/FIXLOG\_7II\_CUR\allowbreak{}RENT.jsonl\#C509;\allowbreak{} 7II:14\allowbreak{}}\par

\par\medskip\noindent\hypertarget{app-E7II-X-010.1}{}\textbf{E7II-X-010.1}\quad printed folio 15.\par
\noindent Printed reading: {\ttfamily\footnotesize n\allowbreak{}}\par
\noindent Corrected reading: {\ttfamily\footnotesize r\allowbreak{}}\par
\noindent Source: {\ttfamily\footnotesize ED-EN-FORMULA-G081;\allowbreak{} fr/sga7ii/\{diplomati\allowbreak{}c,current\}/expose\_X\_\allowbreak{}body.tex;\allowbreak{} folio\allowbreak{} 15\allowbreak{}}\par

\par\medskip\noindent\hypertarget{app-E7II-X-010.2}{}\textbf{E7II-X-010.2}\quad printed folio 15.\par
\noindent Printed reading: {\ttfamily\footnotesize n\allowbreak{}}\par
\noindent Corrected reading: {\ttfamily\footnotesize r\allowbreak{}}\par
\noindent Source: {\ttfamily\footnotesize ED-EN-FORMULA-G081;\allowbreak{} fr/sga7ii/\{diplomati\allowbreak{}c,current\}/expose\_X\_\allowbreak{}body.tex;\allowbreak{} folio\allowbreak{} 15\allowbreak{}}\par

\par\medskip\noindent\hypertarget{app-E7II-X-010.3}{}\textbf{E7II-X-010.3}\quad printed folio 15.\par
\noindent Printed reading: {\ttfamily\footnotesize n\allowbreak{}}\par
\noindent Corrected reading: {\ttfamily\footnotesize r\allowbreak{}}\par
\noindent Source: {\ttfamily\footnotesize ED-EN-FORMULA-G081;\allowbreak{} fr/sga7ii/\{diplomati\allowbreak{}c,current\}/expose\_X\_\allowbreak{}body.tex;\allowbreak{} folio\allowbreak{} 15\allowbreak{}}\par

\par\medskip\noindent\hypertarget{app-E7II-X-010.4}{}\textbf{E7II-X-010.4}\quad printed folio 15.\par
\noindent Printed reading: {\ttfamily\footnotesize n\allowbreak{}}\par
\noindent Corrected reading: {\ttfamily\footnotesize r\allowbreak{}}\par
\noindent Source: {\ttfamily\footnotesize ED-EN-FORMULA-G081;\allowbreak{} fr/sga7ii/\{diplomati\allowbreak{}c,current\}/expose\_X\_\allowbreak{}body.tex;\allowbreak{} folio\allowbreak{} 15\allowbreak{}}\par

\par\medskip\noindent\hypertarget{app-E7II-X-011.1}{}\textbf{E7II-X-011.1}\quad printed folio 15.\par
\noindent Printed reading: {\ttfamily\footnotesize n\allowbreak{}}\par
\noindent Corrected reading: {\ttfamily\footnotesize r\allowbreak{}}\par
\noindent Source: {\ttfamily\footnotesize fr/sga7ii/\{diplomati\allowbreak{}c,current\}/expose\_X\_\allowbreak{}body.tex;\allowbreak{} folio\allowbreak{} 15\allowbreak{}}\par

\par\medskip\noindent\hypertarget{app-E7II-X-011.2}{}\textbf{E7II-X-011.2}\quad printed folio 15.\par
\noindent Printed reading: {\ttfamily\footnotesize n\allowbreak{}}\par
\noindent Corrected reading: {\ttfamily\footnotesize r\allowbreak{}}\par
\noindent Source: {\ttfamily\footnotesize fr/sga7ii/\{diplomati\allowbreak{}c,current\}/expose\_X\_\allowbreak{}body.tex;\allowbreak{} folio\allowbreak{} 15\allowbreak{}}\par

\par\medskip\noindent\hypertarget{app-E7II-X-011.3}{}\textbf{E7II-X-011.3}\quad printed folio 15.\par
\noindent Printed reading: {\ttfamily\footnotesize n\allowbreak{}}\par
\noindent Corrected reading: {\ttfamily\footnotesize r\allowbreak{}}\par
\noindent Source: {\ttfamily\footnotesize fr/sga7ii/\{diplomati\allowbreak{}c,current\}/expose\_X\_\allowbreak{}body.tex;\allowbreak{} folio\allowbreak{} 15\allowbreak{}}\par

\par\medskip\noindent\hypertarget{app-E7II-X-012.1}{}\textbf{E7II-X-012.1}\quad printed folio 18.\par
\noindent Printed reading: {\ttfamily\footnotesize f\allowbreak{}}\par
\noindent Corrected reading: {\ttfamily\footnotesize p\allowbreak{}}\par
\noindent Source: {\ttfamily\footnotesize fr/sga7ii/\{diplomati\allowbreak{}c,current\}/expose\_X\_\allowbreak{}body.tex;\allowbreak{} folio\allowbreak{} 18;\allowbreak{} formula\allowbreak{} 570\allowbreak{}}\par

\par\medskip\noindent\hypertarget{app-E7II-X-012.2}{}\textbf{E7II-X-012.2}\quad printed folio 18.\par
\noindent Printed reading: {\ttfamily\footnotesize f\allowbreak{}}\par
\noindent Corrected reading: {\ttfamily\footnotesize p\allowbreak{}}\par
\noindent Source: {\ttfamily\footnotesize fr/sga7ii/\{diplomati\allowbreak{}c,current\}/expose\_X\_\allowbreak{}body.tex;\allowbreak{} folio\allowbreak{} 18;\allowbreak{} formula\allowbreak{} 572\allowbreak{}}\par

\par\medskip\noindent\hypertarget{app-E7II-X-013.1}{}\textbf{E7II-X-013.1}\quad printed folio 18.\par
\noindent Printed reading: {\ttfamily\footnotesize f\allowbreak{}}\par
\noindent Corrected reading: {\ttfamily\footnotesize p\allowbreak{}}\par
\noindent Source: {\ttfamily\footnotesize ED-EN-FORMULA-G083;\allowbreak{} fr/sga7ii/\{diplomati\allowbreak{}c,current\}/expose\_X\_\allowbreak{}body.tex;\allowbreak{} folio\allowbreak{} 18\allowbreak{}}\par

\par\medskip\noindent\hypertarget{app-E7II-X-013.2}{}\textbf{E7II-X-013.2}\quad printed folio 18.\par
\noindent Printed reading: \emph{(empty)}\par
\noindent Corrected reading: {\ttfamily\footnotesize '\allowbreak{}}\par
\noindent Source: {\ttfamily\footnotesize ED-EN-FORMULA-G083;\allowbreak{} fr/sga7ii/\{diplomati\allowbreak{}c,current\}/expose\_X\_\allowbreak{}body.tex;\allowbreak{} folio\allowbreak{} 18\allowbreak{}}\par

\par\medskip\noindent\hypertarget{app-E7II-X-013.3}{}\textbf{E7II-X-013.3}\quad printed folio 18.\par
\noindent Printed reading: {\ttfamily\footnotesize f\allowbreak{}}\par
\noindent Corrected reading: {\ttfamily\footnotesize Lp\allowbreak{}}\par
\noindent Source: {\ttfamily\footnotesize fr/sga7ii/\{diplomati\allowbreak{}c,current\}/expose\_X\_\allowbreak{}body.tex;\allowbreak{} folio\allowbreak{} 18\allowbreak{}}\par

\par\medskip\noindent\hypertarget{app-E7II-X-013.4}{}\textbf{E7II-X-013.4}\quad printed folio 18.\par
\noindent Printed reading: {\ttfamily\footnotesize f\allowbreak{}}\par
\noindent Corrected reading: {\ttfamily\footnotesize Rp\allowbreak{}}\par
\noindent Source: {\ttfamily\footnotesize fr/sga7ii/\{diplomati\allowbreak{}c,current\}/expose\_X\_\allowbreak{}body.tex;\allowbreak{} folio\allowbreak{} 18\allowbreak{}}\par

\par\medskip\noindent\hypertarget{app-E7II-X-014.1}{}\textbf{E7II-X-014.1}\quad printed folio 18.\par
\noindent Printed reading: {\ttfamily\footnotesize S\allowbreak{}}\par
\noindent Corrected reading: {\ttfamily\footnotesize D\allowbreak{}}\par
\noindent Source: {\ttfamily\footnotesize fr/sga7ii/\{diplomati\allowbreak{}c,current\}/expose\_X\_\allowbreak{}body.tex;\allowbreak{} folio\allowbreak{} 18;\allowbreak{} formula\allowbreak{} 582\allowbreak{}}\par

\par\medskip\noindent\hypertarget{app-E7II-X-015.1}{}\textbf{E7II-X-015.1}\quad printed folio 19.\par
\noindent Printed reading: \emph{(empty)}\par
\noindent Corrected reading: {\ttfamily\footnotesize '\allowbreak{}}\par
\noindent Source: {\ttfamily\footnotesize ED-EN-FORMULA-G085;\allowbreak{} fr/sga7ii/\{diplomati\allowbreak{}c,current\}/expose\_X\_\allowbreak{}body.tex;\allowbreak{} folio\allowbreak{} 19\allowbreak{}}\par

\par\medskip\noindent\hypertarget{app-E7II-X-016.1}{}\textbf{E7II-X-016.1}\quad printed folio 20.\par
\noindent Printed reading: {\ttfamily\footnotesize Lf\allowbreak{}}\par
\noindent Corrected reading: {\ttfamily\footnotesize Lp\allowbreak{}}\par
\noindent Source: {\ttfamily\footnotesize fr/sga7ii/\{diplomati\allowbreak{}c,current\}/expose\_X\_\allowbreak{}body.tex;\allowbreak{} folio\allowbreak{} 20;\allowbreak{} formula\allowbreak{} 640\allowbreak{}}\par

\par\medskip\noindent\hypertarget{app-E7II-X-016.2}{}\textbf{E7II-X-016.2}\quad printed folio 20.\par
\noindent Printed reading: {\ttfamily\footnotesize f\allowbreak{}}\par
\noindent Corrected reading: {\ttfamily\footnotesize p\allowbreak{}}\par
\noindent Source: {\ttfamily\footnotesize fr/sga7ii/\{diplomati\allowbreak{}c,current\}/expose\_X\_\allowbreak{}body.tex;\allowbreak{} folio\allowbreak{} 20;\allowbreak{} formula\allowbreak{} 640\allowbreak{}}\par

\par\medskip\noindent\hypertarget{app-E7II-X-017.1}{}\textbf{E7II-X-017.1}\quad printed folio 20.\par
\noindent Printed reading: {\ttfamily\footnotesize f\allowbreak{}}\par
\noindent Corrected reading: {\ttfamily\footnotesize p\allowbreak{}}\par
\noindent Source: {\ttfamily\footnotesize fr/sga7ii/\{diplomati\allowbreak{}c,current\}/expose\_X\_\allowbreak{}body.tex;\allowbreak{} folio\allowbreak{} 20;\allowbreak{} formula\allowbreak{} 645\allowbreak{}}\par

\par\medskip\noindent\hypertarget{app-E7II-X-017.2}{}\textbf{E7II-X-017.2}\quad printed folio 20.\par
\noindent Printed reading: {\ttfamily\footnotesize f\allowbreak{}}\par
\noindent Corrected reading: {\ttfamily\footnotesize p\allowbreak{}}\par
\noindent Source: {\ttfamily\footnotesize fr/sga7ii/\{diplomati\allowbreak{}c,current\}/expose\_X\_\allowbreak{}body.tex;\allowbreak{} folio\allowbreak{} 20;\allowbreak{} formula\allowbreak{} 645\allowbreak{}}\par

\par\medskip\noindent\hypertarget{app-E7II-X-017.3}{}\textbf{E7II-X-017.3}\quad printed folio 20.\par
\noindent Printed reading: {\ttfamily\footnotesize f\allowbreak{}}\par
\noindent Corrected reading: {\ttfamily\footnotesize p\allowbreak{}}\par
\noindent Source: {\ttfamily\footnotesize fr/sga7ii/\{diplomati\allowbreak{}c,current\}/expose\_X\_\allowbreak{}body.tex;\allowbreak{} folio\allowbreak{} 20;\allowbreak{} formula\allowbreak{} 645\allowbreak{}}\par

\par\medskip\noindent\hypertarget{app-E7II-X-017.4}{}\textbf{E7II-X-017.4}\quad printed folio 20.\par
\noindent Printed reading: {\ttfamily\footnotesize f\allowbreak{}}\par
\noindent Corrected reading: {\ttfamily\footnotesize p\_*\textbackslash{}underline\{O\}\_E\allowbreak{} \textbackslash{}longrightarrow\allowbreak{} p\_*\textbackslash{}underline\{O\}\_\{D'\allowbreak{}\textbackslash{}cap\allowbreak{} E\}\allowbreak{} \textbackslash{}longrightarrow\allowbreak{} p\allowbreak{}}\par
\noindent Source: {\ttfamily\footnotesize fr/sga7ii/\{diplomati\allowbreak{}c,current\}/expose\_X\_\allowbreak{}body.tex;\allowbreak{} folio\allowbreak{} 20;\allowbreak{} formula\allowbreak{} 646\allowbreak{}}\par

\par\medskip\noindent\hypertarget{app-E7II-X-017.5}{}\textbf{E7II-X-017.5}\quad printed folio 20.\par
\noindent Printed reading: {\ttfamily\footnotesize  f\_*\textbackslash{}underline\{O\}\_E\allowbreak{} \textbackslash{}longrightarrow\allowbreak{} f\_*\textbackslash{}\allowbreak{} \textbackslash{}underline\{O\}\_\{D'\textbackslash{}ca\allowbreak{}p\allowbreak{} E\}\allowbreak{} \textbackslash{}longrightarrow\allowbreak{}}\par
\noindent Corrected reading: \emph{(empty)}\par
\noindent Source: {\ttfamily\footnotesize fr/sga7ii/\{diplomati\allowbreak{}c,current\}/expose\_X\_\allowbreak{}body.tex;\allowbreak{} folio\allowbreak{} 20;\allowbreak{} formula\allowbreak{} 646\allowbreak{}}\par

\par\medskip\noindent\hypertarget{app-E7II-X-018.1}{}\textbf{E7II-X-018.1}\quad printed folio 21.\par
\noindent Printed reading: {\ttfamily\footnotesize f\allowbreak{}}\par
\noindent Corrected reading: {\ttfamily\footnotesize p\allowbreak{}}\par
\noindent Source: {\ttfamily\footnotesize fr/sga7ii/\{diplomati\allowbreak{}c,current\}/expose\_X\_\allowbreak{}body.tex;\allowbreak{} folio\allowbreak{} 21;\allowbreak{} formula\allowbreak{} 647\allowbreak{}}\par

\par\medskip\noindent\hypertarget{app-E7II-X-018.2}{}\textbf{E7II-X-018.2}\quad printed folio 21.\par
\noindent Printed reading: \emph{(empty)}\par
\noindent Corrected reading: {\ttfamily\footnotesize \textbackslash{}mu\_s(D)\textasciicircum{}2[k(s):k]\allowbreak{} -\allowbreak{}}\par
\noindent Source: {\ttfamily\footnotesize fr/sga7ii/\{diplomati\allowbreak{}c,current\}/expose\_X\_\allowbreak{}body.tex;\allowbreak{} folio\allowbreak{} 21;\allowbreak{} formula\allowbreak{} 647\allowbreak{}}\par

\par\medskip\noindent\hypertarget{app-E7II-X-018.3}{}\textbf{E7II-X-018.3}\quad printed folio 21.\par
\noindent Printed reading: {\ttfamily\footnotesize [k(s):\allowbreak{} k]\allowbreak{} -\allowbreak{} \textbackslash{}mu\_s(D)\textasciicircum{}2\allowbreak{}}\par
\noindent Corrected reading: \emph{(empty)}\par
\noindent Source: {\ttfamily\footnotesize fr/sga7ii/\{diplomati\allowbreak{}c,current\}/expose\_X\_\allowbreak{}body.tex;\allowbreak{} folio\allowbreak{} 21;\allowbreak{} formula\allowbreak{} 647\allowbreak{}}\par

\par\medskip\noindent\hypertarget{app-E7II-X-018.4}{}\textbf{E7II-X-018.4}\quad printed folio 21.\par
\noindent Printed reading: {\ttfamily\footnotesize f\allowbreak{}}\par
\noindent Corrected reading: {\ttfamily\footnotesize p\allowbreak{}}\par
\noindent Source: {\ttfamily\footnotesize fr/sga7ii/\{diplomati\allowbreak{}c,current\}/expose\_X\_\allowbreak{}body.tex;\allowbreak{} folio\allowbreak{} 21;\allowbreak{} formula\allowbreak{} 656\allowbreak{}}\par

\par\medskip\noindent\hypertarget{app-E7II-X-018.5}{}\textbf{E7II-X-018.5}\quad printed folio 21.\par
\noindent Printed reading: {\ttfamily\footnotesize S\allowbreak{}}\par
\noindent Corrected reading: {\ttfamily\footnotesize \textbackslash{}delta\allowbreak{}}\par
\noindent Source: {\ttfamily\footnotesize fr/sga7ii/\{diplomati\allowbreak{}c,current\}/expose\_X\_\allowbreak{}body.tex;\allowbreak{} folio\allowbreak{} 21;\allowbreak{} formula\allowbreak{} 657\allowbreak{}}\par

\par\medskip\noindent\hypertarget{app-E7II-X-019.1}{}\textbf{E7II-X-019.1}\quad printed folio 22.\par
\noindent Printed reading: {\ttfamily\footnotesize \textbackslash{}underline\{let\}\allowbreak{} \$D\$\allowbreak{} \textbackslash{}underline\{be\allowbreak{} a\allowbreak{} divisor\};\allowbreak{} }\par
\noindent Corrected reading: \emph{(empty)}\par
\noindent Source: {\ttfamily\footnotesize ED-EN-NONMATH-G000;\allowbreak{} F7II-X-023.1;\allowbreak{} ADOPT-F7II-X-023.1;\allowbreak{} fr/sga7ii/\{diplomati\allowbreak{}c,current\}/expose\_X\_\allowbreak{}body.tex;\allowbreak{} folio\allowbreak{} 22\allowbreak{}}\par

\par\medskip\noindent\hypertarget{app-E7II-X-020.1}{}\textbf{E7II-X-020.1}\quad printed folio 23.\par
\noindent Printed reading: {\ttfamily\footnotesize holomorphic\allowbreak{}}\par
\noindent Corrected reading: {\ttfamily\footnotesize formal\allowbreak{}}\par
\noindent Source: {\ttfamily\footnotesize ED-EN-NONMATH-G001;\allowbreak{} F7II-X-024.1;\allowbreak{} ADOPT-F7II-X-024.1;\allowbreak{} fr/sga7ii/\{diplomati\allowbreak{}c,current\}/expose\_X\_\allowbreak{}body.tex;\allowbreak{} folio\allowbreak{} 23\allowbreak{}}\par

\par\medskip\noindent\hypertarget{app-E7II-X-021.1}{}\textbf{E7II-X-021.1}\quad printed folio 24.\par
\noindent Printed reading: \emph{(empty)}\par
\noindent Corrected reading: {\ttfamily\footnotesize )\allowbreak{}}\par
\noindent Source: {\ttfamily\footnotesize ED-en7II-C8-source;\allowbreak{} OPEN\_ITEMS\_DECIDED.t\allowbreak{}sv\allowbreak{}}\par

\par\medskip\noindent\hypertarget{app-E7II-X-021.2}{}\textbf{E7II-X-021.2}\quad printed folio 24.\par
\noindent Printed reading: {\ttfamily\footnotesize \textbackslash{}cap\allowbreak{}}\par
\noindent Corrected reading: {\ttfamily\footnotesize \textbackslash{}cup\allowbreak{}}\par
\noindent Source: {\ttfamily\footnotesize ED-en7II-C8-source;\allowbreak{} OPEN\_ITEMS\_DECIDED.t\allowbreak{}sv\allowbreak{}}\par

\par\medskip\noindent\hypertarget{app-E7II-X-021.3}{}\textbf{E7II-X-021.3}\quad printed folio 24.\par
\noindent Printed reading: {\ttfamily\footnotesize \textbackslash{}cap\allowbreak{}}\par
\noindent Corrected reading: {\ttfamily\footnotesize \textbackslash{}cup\allowbreak{}}\par
\noindent Source: {\ttfamily\footnotesize fr/sga7ii/\{diplomati\allowbreak{}c,current\}/expose\_X\_\allowbreak{}body.tex;\allowbreak{} folio\allowbreak{} 24;\allowbreak{} formula\allowbreak{} 761;\allowbreak{} ED-en7II-C8-source;\allowbreak{} OPEN\_ITEMS\_DECIDED.t\allowbreak{}sv\allowbreak{}}\par

\par\medskip\noindent\hypertarget{app-E7II-X-022.1}{}\textbf{E7II-X-022.1}\quad printed folio 24.\par
\noindent Printed reading: \emph{(empty)}\par
\noindent Corrected reading: {\ttfamily\footnotesize \textbackslash{}underline\{O\}\allowbreak{}}\par
\noindent Source: {\ttfamily\footnotesize ED-EN-FORMULA-G087;\allowbreak{} fr/sga7ii/\{diplomati\allowbreak{}c,current\}/expose\_X\_\allowbreak{}body.tex;\allowbreak{} folio\allowbreak{} 24\allowbreak{}}\par

\par\medskip\noindent\hypertarget{app-E7II-X-023.1}{}\textbf{E7II-X-023.1}\quad printed folio 25.\par
\noindent Printed reading: \emph{(empty)}\par
\noindent Corrected reading: {\ttfamily\footnotesize )\textbackslash{}Bigr\allowbreak{}}\par
\noindent Source: {\ttfamily\footnotesize fr/sga7ii/\{diplomati\allowbreak{}c,current\}/expose\_X\_\allowbreak{}body.tex;\allowbreak{} folio\allowbreak{} 25\allowbreak{}}\par

\par\medskip\noindent\hypertarget{app-E7II-X-024.1}{}\textbf{E7II-X-024.1}\quad printed folio 26.\par
\noindent Printed reading: \emph{(empty)}\par
\noindent Corrected reading: {\ttfamily\footnotesize )\textbackslash{}Bigr\allowbreak{}}\par
\noindent Source: {\ttfamily\footnotesize fr/sga7ii/\{diplomati\allowbreak{}c,current\}/expose\_X\_\allowbreak{}body.tex;\allowbreak{} folio\allowbreak{} 26;\allowbreak{} formula\allowbreak{} 820\allowbreak{}}\par

\par\medskip\noindent\hypertarget{app-E7II-X-024.2}{}\textbf{E7II-X-024.2}\quad printed folio 26.\par
\noindent Printed reading: \emph{(empty)}\par
\noindent Corrected reading: {\ttfamily\footnotesize \textbackslash{}Bigr)\allowbreak{}}\par
\noindent Source: {\ttfamily\footnotesize fr/sga7ii/\{diplomati\allowbreak{}c,current\}/expose\_X\_\allowbreak{}body.tex;\allowbreak{} folio\allowbreak{} 26;\allowbreak{} formula\allowbreak{} 822\allowbreak{}}\par

\par\medskip\noindent\hypertarget{app-E7II-X-024.3}{}\textbf{E7II-X-024.3}\quad printed folio 26.\par
\noindent Printed reading: \emph{(empty)}\par
\noindent Corrected reading: {\ttfamily\footnotesize \textbackslash{}substack\allowbreak{}}\par
\noindent Source: {\ttfamily\footnotesize fr/sga7ii/\{diplomati\allowbreak{}c,current\}/expose\_X\_\allowbreak{}body.tex;\allowbreak{} folio\allowbreak{} 26;\allowbreak{} formula\allowbreak{} 822\allowbreak{}}\par

\par\medskip\noindent\hypertarget{app-E7II-X-024.4}{}\textbf{E7II-X-024.4}\quad printed folio 26.\par
\noindent Printed reading: \emph{(empty)}\par
\noindent Corrected reading: {\ttfamily\footnotesize \textbackslash{}\textbackslash{}t\textbackslash{}ne\allowbreak{} s\allowbreak{}}\par
\noindent Source: {\ttfamily\footnotesize fr/sga7ii/\{diplomati\allowbreak{}c,current\}/expose\_X\_\allowbreak{}body.tex;\allowbreak{} folio\allowbreak{} 26;\allowbreak{} formula\allowbreak{} 822\allowbreak{}}\par

\par\medskip\noindent\hypertarget{app-E7II-X-025.1}{}\textbf{E7II-X-025.1}\quad printed folio 26.\par
\noindent Printed reading: {\ttfamily\footnotesize s\allowbreak{}}\par
\noindent Corrected reading: {\ttfamily\footnotesize S\allowbreak{}}\par
\noindent Source: {\ttfamily\footnotesize fr/sga7ii/\{diplomati\allowbreak{}c,current\}/expose\_X\_\allowbreak{}body.tex;\allowbreak{} folio\allowbreak{} 26;\allowbreak{} formula\allowbreak{} 831\allowbreak{}}\par

\par\medskip\noindent\hypertarget{app-E7II-X-025.2}{}\textbf{E7II-X-025.2}\quad printed folio 26.\par
\noindent Printed reading: \emph{(empty)}\par
\noindent Corrected reading: {\ttfamily\footnotesize )\textbackslash{}bigr\allowbreak{}}\par
\noindent Source: {\ttfamily\footnotesize fr/sga7ii/\{diplomati\allowbreak{}c,current\}/expose\_X\_\allowbreak{}body.tex;\allowbreak{} folio\allowbreak{} 26;\allowbreak{} formula\allowbreak{} 831\allowbreak{}}\par

\par\medskip\noindent\hypertarget{app-E7II-X-026.1}{}\textbf{E7II-X-026.1}\quad printed folio 27.\par
\noindent Printed reading: \emph{(empty)}\par
\noindent Corrected reading: {\ttfamily\footnotesize (t)\allowbreak{}}\par
\noindent Source: {\ttfamily\footnotesize ED-EN-FORMULA-G089;\allowbreak{} fr/sga7ii/\{diplomati\allowbreak{}c,current\}/expose\_X\_\allowbreak{}body.tex;\allowbreak{} folio\allowbreak{} 27\allowbreak{}}\par

\par\medskip\noindent\hypertarget{app-E7II-X-027.1}{}\textbf{E7II-X-027.1}\quad printed folio 29.\par
\noindent Printed reading: {\ttfamily\footnotesize 2\allowbreak{}}\par
\noindent Corrected reading: {\ttfamily\footnotesize 1\allowbreak{}}\par
\noindent Source: {\ttfamily\footnotesize fr/sga7ii/\{diplomati\allowbreak{}c,current\}/expose\_X\_\allowbreak{}body.tex;\allowbreak{} folio\allowbreak{} 29;\allowbreak{} formula\allowbreak{} 929\allowbreak{}}\par

\par\medskip\noindent\hypertarget{app-E7II-X-028.1}{}\textbf{E7II-X-028.1}\quad printed folio 31.\par
\noindent Printed reading: {\ttfamily\footnotesize 4\allowbreak{}}\par
\noindent Corrected reading: {\ttfamily\footnotesize 2\allowbreak{}}\par
\noindent Source: {\ttfamily\footnotesize fr/sga7ii/\{diplomati\allowbreak{}c,current\}/expose\_X\_\allowbreak{}body.tex;\allowbreak{} folio\allowbreak{} 31;\allowbreak{} port/PORT\_EN7II\_LOG.\allowbreak{}jsonl\#C510;\allowbreak{} fixes/FIXLOG\_7II\_CUR\allowbreak{}RENT.jsonl\#C510;\allowbreak{} 7II:41\allowbreak{}}\par

\par\medskip\noindent\hypertarget{app-E7II-X-029.1}{}\textbf{E7II-X-029.1}\quad printed folio 32.\par
\noindent Printed reading: \emph{(empty)}\par
\noindent Corrected reading: {\ttfamily\footnotesize \_s\allowbreak{}}\par
\noindent Source: {\ttfamily\footnotesize fr/sga7ii/\{diplomati\allowbreak{}c,current\}/expose\_X\_\allowbreak{}body.tex;\allowbreak{} folio\allowbreak{} 32;\allowbreak{} formula\allowbreak{} 1003\allowbreak{}}\par

\par\medskip\noindent\hypertarget{app-E7II-X-030.1}{}\textbf{E7II-X-030.1}\quad printed folio 33.\par
\noindent Printed reading: {\ttfamily\footnotesize o\allowbreak{}}\par
\noindent Corrected reading: {\ttfamily\footnotesize 0\allowbreak{}}\par
\noindent Source: {\ttfamily\footnotesize fr/sga7ii/\{diplomati\allowbreak{}c,current\}/expose\_X\_\allowbreak{}body.tex;\allowbreak{} folio\allowbreak{} 33;\allowbreak{} formula\allowbreak{} 1048\allowbreak{}}\par

\par\medskip\noindent\hypertarget{app-E7II-X-030.2}{}\textbf{E7II-X-030.2}\quad printed folio 33.\par
\noindent Printed reading: {\ttfamily\footnotesize S\allowbreak{}}\par
\noindent Corrected reading: {\ttfamily\footnotesize T\allowbreak{}}\par
\noindent Source: {\ttfamily\footnotesize ED-EN-FORMULA-G091;\allowbreak{} fr/sga7ii/\{diplomati\allowbreak{}c,current\}/expose\_X\_\allowbreak{}body.tex;\allowbreak{} folio\allowbreak{} 33\allowbreak{}}\par

\par\medskip\noindent\hypertarget{app-E7II-X-030.3}{}\textbf{E7II-X-030.3}\quad printed folio 33.\par
\noindent Printed reading: \emph{(empty)}\par
\noindent Corrected reading: {\ttfamily\footnotesize ,\allowbreak{}}\par
\noindent Source: {\ttfamily\footnotesize fr/sga7ii/\{diplomati\allowbreak{}c,current\}/expose\_X\_\allowbreak{}body.tex;\allowbreak{} folio\allowbreak{} 33;\allowbreak{} formula\allowbreak{} 1066\allowbreak{}}\par

\par\medskip\noindent\hypertarget{app-E7II-X-030.4}{}\textbf{E7II-X-030.4}\quad printed folio 33.\par
\noindent Printed reading: \emph{(1 espace(s), 0 retour(s) de ligne dans le témoin TeX)}\par
\noindent Corrected reading: {\ttfamily\footnotesize ,\allowbreak{}}\par
\noindent Source: {\ttfamily\footnotesize fr/sga7ii/\{diplomati\allowbreak{}c,current\}/expose\_X\_\allowbreak{}body.tex;\allowbreak{} folio\allowbreak{} 33;\allowbreak{} formula\allowbreak{} 1066\allowbreak{}}\par

\par\medskip\noindent\hypertarget{app-E7II-X-030.5}{}\textbf{E7II-X-030.5}\quad printed folio 33.\par
\noindent Printed reading: \emph{(empty)}\par
\noindent Corrected reading: {\ttfamily\footnotesize \_2\allowbreak{}}\par
\noindent Source: {\ttfamily\footnotesize fr/sga7ii/\{diplomati\allowbreak{}c,current\}/expose\_X\_\allowbreak{}body.tex;\allowbreak{} folio\allowbreak{} 33;\allowbreak{} formula\allowbreak{} 1074\allowbreak{}}\par

\par\medskip\noindent\hypertarget{app-E7II-X-031.1}{}\textbf{E7II-X-031.1}\quad printed folio 34.\par
\noindent Printed reading: \emph{(empty)}\par
\noindent Corrected reading: {\ttfamily\footnotesize \_0\allowbreak{}}\par
\noindent Source: {\ttfamily\footnotesize fr/sga7ii/\{diplomati\allowbreak{}c,current\}/expose\_X\_\allowbreak{}body.tex;\allowbreak{} folio\allowbreak{} 34\allowbreak{}}\par

\par\medskip\noindent\hypertarget{app-E7II-X-032.1}{}\textbf{E7II-X-032.1}\quad printed folio 34.\par
\noindent Printed reading: {\ttfamily\footnotesize -\allowbreak{}}\par
\noindent Corrected reading: \emph{(empty)}\par
\noindent Source: {\ttfamily\footnotesize fr/sga7ii/\{diplomati\allowbreak{}c,current\}/expose\_X\_\allowbreak{}body.tex;\allowbreak{} folio\allowbreak{} 34;\allowbreak{} formula\allowbreak{} 1099\allowbreak{}}\par

\par\medskip\noindent\hypertarget{app-E7II-X-033.1}{}\textbf{E7II-X-033.1}\quad printed folio 35.\par
\noindent Printed reading: \emph{(empty)}\par
\noindent Corrected reading: {\ttfamily\footnotesize \textasciicircum{}E\allowbreak{}}\par
\noindent Source: {\ttfamily\footnotesize fr/sga7ii/\{diplomati\allowbreak{}c,current\}/expose\_X\_\allowbreak{}body.tex;\allowbreak{} folio\allowbreak{} 35;\allowbreak{} formula\allowbreak{} 1125\allowbreak{}}\par

\par\medskip\noindent\hypertarget{app-E7II-X-034.1}{}\textbf{E7II-X-034.1}\quad printed folio 36.\par
\noindent Printed reading: \emph{(empty)}\par
\noindent Corrected reading: {\ttfamily\footnotesize \_\{n\textbackslash{}geq0\}\allowbreak{}}\par
\noindent Source: {\ttfamily\footnotesize ED-EN-FORMULA-G093;\allowbreak{} fr/sga7ii/\{diplomati\allowbreak{}c,current\}/expose\_X\_\allowbreak{}body.tex;\allowbreak{} folio\allowbreak{} 36\allowbreak{}}\par

\par\medskip\noindent\hypertarget{app-E7II-X-035.1}{}\textbf{E7II-X-035.1}\quad printed folio 37.\par
\noindent Printed reading: {\ttfamily\footnotesize 15\allowbreak{}}\par
\noindent Corrected reading: {\ttfamily\footnotesize 90\allowbreak{}}\par
\noindent Source: {\ttfamily\footnotesize ED-EN-NONMATH-G013;\allowbreak{} F7II-X-043.3;\allowbreak{} 7II:54;\allowbreak{} fr/sga7ii/\{diplomati\allowbreak{}c,current\}/expose\_X\_\allowbreak{}body.tex;\allowbreak{} folio\allowbreak{} 37\allowbreak{}}\par

\par\endgroup\normalsize\fi
\clearpage\SGAtoc{Exposé XI}
\selectlanguage{english}
\clearpage
\SGApage{II}{39}

\begin{center}
{\large\textbf{SGA 7}}\\[1.5em]
{\Large\underline{EXPOSÉ XI}}\\[2.5em]
{\Large\underline{\underline{COHOMOLOGY OF COMPLETE INTERSECTIONS}}}\\[1.5em]
{\large\underline{by P. DELIGNE}}
\end{center}

\bigskip
\noindent 1. \underline{Cohomology of projective space and complete
intersections}

Let $X$ be a nonsingular complex projective variety, purely of dimension $n$,
and let $\underline O(1)$ be an ample invertible sheaf on $X$.  By Akizuki and
Nakano [1], Kodaira's vanishing theorem [4] implies that

\begin{equation}\tag{*}
H^i(X,\Omega_X^j(-1))=0\qquad\text{for }i+j<n.
\end{equation}

Again by [1], this result implies the ``weak Lefschetz theorem'' in complex
cohomology.  We propose to prove, by brute force, that if $X$ is a projective
space over a field of arbitrary characteristic, then (*) remains true.  This
will allow us, following [1], to compare the Hodge and de Rham cohomology of
complete intersections over a field of arbitrary characteristic with that of
projective space.

We shall then spell out some consequences of the results obtained.  In §2 we
give qualitative results deduced from Hirzebruch's numerical formula (2.2).  I
wish to thank N. Katz for his help in the preparation of that section.

1.0. Let $\underline E$ be a locally free module of rank $r+1$ on a scheme
$S$.  \underline{We assume that $r>0$.}  Recall that
$\mathbb P(\underline E)$ denotes the relative projective space over $S$ of
hyperplanes of $\underline E$.  Let $p:\mathbb P(\underline E)\to S$ be the
projection, and let

\begin{equation*}
\eta\in H^0\bigl(S,R^1p_*\Omega^1_{\mathbb P(\underline E)/S}\bigr)
\end{equation*}

be the first Chern class, in the Hodge sense, of the invertible sheaf
$\underline O(1)$.


\SGApage{II}{\(\approx\)40}
Recall that $\eta$ is the image of the class of $\underline O(1)$ in
$H^0(S,R^1p_*\underline O^*)$ under the logarithmic differential

\begin{equation*}
df/f:\underline O^*\longrightarrow
\Omega^1_{\mathbb P(\underline E)/S}.
\end{equation*}

For every sheaf $\underline F$ of $\underline O$-modules on
$\mathbb P(\underline E)$, put

\begin{equation*}
\underline F(n)=\underline F\otimes\underline O(1)^{\otimes n}
\qquad(n\in\mathbb Z).
\end{equation*}

\underline{Theorem} 1.1. (i) \underline{The sheaves
$R^ip_*(\Omega^j_{\mathbb P(\underline E)/S}(n))$ are locally free, and their
formation is compatible with arbitrary base change.}

(ii) \underline{For $0\leq i\leq r$, the section $\eta^i$ of
$R^ip_*(\Omega^i_{\mathbb P(\underline E)/S})$ defines an isomorphism from
$\underline O_S$ onto this sheaf. Moreover,}

\begin{equation*}
R^ip_*\Omega^j_{\mathbb P(\underline E)/S}=0
\qquad\underline{\text{if }}i\ne j\ \underline{\text{ or }}\ i>r.
\end{equation*}

(iii) \underline{If $n\ne0$, then, outside the following two cases,
$R^ip_*(\Omega^j_{\mathbb P(\underline E)/S}(n))$ vanishes:}

\begin{align}
i&=0\quad\underline{\text{and}}\quad n>j,
  \tag{iii$_a$}\\
i&=r\quad\underline{\text{and}}\quad n<j-r.
  \tag{iii$_b$}
\end{align}

Let

\begin{equation*}
\mathbb V(\underline E)=\operatorname{Spec}
\bigl(\operatorname{Sym}^*\underline E\bigr)
\end{equation*}

be the affine space defined by $\underline E$, let $e$ be its zero section,
and put
$\mathbb V(\underline E)^*=\mathbb V(\underline E)\setminus e(S)$.  Let $a$
and $a'$ be the projections of $\mathbb V(\underline E)$ and
$\mathbb V(\underline E)^*$ onto $S$, and let $q$ be the projection of
$\mathbb V(\underline E)^*$ onto $\mathbb P(\underline E)$:

\begin{center}
\begin{tikzcd}[row sep=large,column sep=large]
\mathbb V(\underline E)
  \arrow[dr,"a"]
\& \mathbb V(\underline E)^*
  \arrow[l,hook']
  \arrow[r,"q"]
  \arrow[d,"a'"']
\& \mathbb P(\underline E)
  \arrow[dl,"p"]
\\
{} \& S\arrow[ul,bend right,"e"'] \& {}
\end{tikzcd}
\qquad .
\end{center}

The punctured affine space $\mathbb V(\underline O(1))^*$ over
$\mathbb P(\underline E)$ is canonically isomorphic to
$\mathbb V(\underline E)^*$.  Consequently, for every quasi-coherent sheaf
$\underline F$ of $\underline O$-modules on $\mathbb P(\underline E)$, there
is a canonical isomorphism


\SGApage{II}{\(\approx\)41}
\begin{equation}\tag{1.1.1}
q_*q^*\underline F\simeq
\bigoplus_{n\in\mathbb Z}\underline F(n).
\end{equation}

This grading of $q_*q^*\underline F$ can also be obtained from the action on
this sheaf of the group $\mathbb G_m$ of homotheties of
$\mathbb V(\underline E)$.

Relative vector fields on $\mathbb V(\underline E)$ are identified with
$S$-morphisms from $\mathbb V(\underline E)$ to itself, the target being
identified with the tangent space at the origin.  The vector field $X$ defined
by the identity map of $\mathbb V(\underline E)$ is invariant under
homotheties.  It is tangent to the fibers of $q$ and defines a trivialization

\begin{equation}\tag{1.1.2}
X:\Omega^1_{\mathbb V(\underline E)^*/\mathbb P(\underline E)}
\longrightarrow\underline O_{\mathbb V(\underline E)^*}.
\end{equation}

If the $x_i$ form a basis of $\underline E$, then

\begin{equation*}
X=\sum_i x_i\,dx^i.
\end{equation*}

Since $\Omega^1_{\mathbb V(\underline E)^*/\mathbb P(\underline E)}$ has
rank one, the exact sequence

\begin{equation*}
0\longrightarrow q^*\Omega^1_{\mathbb P(\underline E)/S}
\longrightarrow\Omega^1_{\mathbb V(\underline E)^*/S}
\longrightarrow
\Omega^1_{\mathbb V(\underline E)^*/\mathbb P(\underline E)}
\longrightarrow0
\end{equation*}

induces exact sequences

\begin{equation*}
0\longrightarrow q^*\Omega^i_{\mathbb P(\underline E)/S}
\longrightarrow\Omega^i_{\mathbb V(\underline E)^*/S}
\longrightarrow
\Omega^1_{\mathbb V(\underline E)^*/\mathbb P(\underline E)}
\otimes q^*\Omega^{i-1}_{\mathbb P(\underline E)/S}
\longrightarrow0.
\end{equation*}

Using (1.1.2), these become

\begin{equation}\tag{1.1.3}
0\longrightarrow q^*\Omega^i_{\mathbb P(\underline E)/S}
\longrightarrow\Omega^i_{\mathbb V(\underline E)^*/S}
\xrightarrow{\ \iota_X\ }
q^*\Omega^{i-1}_{\mathbb P(\underline E)/S}
\longrightarrow0.
\end{equation}

Canonically,

\begin{equation}\tag{1.1.4}
\Omega^i_{\mathbb V(\underline E)/S}
=a^*\bigwedge^i\underline E.
\end{equation}

The short exact sequences (1.1.3) are therefore obtained by splitting the
following long exact sequence on $\mathbb V(\underline E)^*$:

\begin{equation}\tag{1.1.5}
\cdots\longrightarrow {a'}^*\bigwedge^i\underline E
\xrightarrow{\ \iota_X\ }
{a'}^*\bigwedge^{i-1}\underline E
\xrightarrow{\ \iota_X\ }
{a'}^*\bigwedge^{i-2}\underline E\longrightarrow\cdots.
\end{equation}



\SGApage{II}{\(\approx\)42}
\SGAvar{This long exact sequence is simply the restriction to
$\mathbb V(\underline E)^*$ of the Kozul resolution of
$\underline O_{e(S)}$ on $\mathbb V(\underline E)$.

}{This long exact sequence is simply the restriction to
$\mathbb V(\underline E)^*$ of the Koszul resolution of
$\underline O_{e(S)}$ on $\mathbb V(\underline E)$.

}{E7II-XI-001}Push the long exact sequence (1.1.5) forward to $\mathbb P(\underline E)$ by
the affine morphism $q$.  By (1.1.1) and the formula
$q_*{a'}^*=q_*q^*p^*$, one obtains a long exact sequence

\begin{equation}\tag{1.1.6}
\cdots\longrightarrow
\bigoplus_{n\in\mathbb Z}p^*\bigwedge^i\underline E(n)
\xrightarrow{\ \iota_X\ }
\bigoplus_{n\in\mathbb Z}p^*\bigwedge^{i-1}\underline E(n)
\longrightarrow\cdots .
\end{equation}

It splits into short exact sequences, the images of (1.1.3),

\begin{equation}\tag{1.1.7}
0\longrightarrow
\bigoplus_{n\in\mathbb Z}\Omega^i_{\mathbb P(\underline E)/S}(n)
\longrightarrow
\bigoplus_{n\in\mathbb Z}p^*\bigwedge^i\underline E(n)
\xrightarrow{\ \iota_X\ }
\bigoplus_{n\in\mathbb Z}\Omega^{i-1}_{\mathbb P(\underline E)/S}(n)
\longrightarrow0.
\end{equation}

The exact sequences (1.1.7) are direct sums of the exact sequences

\begin{equation}\tag{1.1.8}
0\longrightarrow\Omega^i_{\mathbb P(\underline E)/S}(n)
\longrightarrow p^*\bigwedge^i\underline E(n-i)
\xrightarrow{\ \iota_X\ }
\Omega^{i-1}_{\mathbb P(\underline E)/S}(n)
\longrightarrow0.
\end{equation}

Recall (EGA III, 2.1.15, or FAC) that the sheaves
$R^ip_*\underline O(k)$ are locally free and their formation is compatible
with arbitrary base change, and that

\begin{equation}\tag{1.1.9}
R^ip_*\underline O(k)=0
\quad\text{if }i\ne0,r,
\quad\text{or if }i=0, k<0,
\quad\text{or if }i=r, k>-r-1,
\end{equation}

and hence that $R^ip_*p^*\bigwedge^j\underline E(k)$ vanishes under the
same conditions.

The exact sequences (1.1.8) give a resolution

\begin{equation}\tag{1.1.10}
0\longrightarrow\Omega^i_{\mathbb P(\underline E)/S}(n)
\xrightarrow{\ E\ }p^*\bigwedge^i\underline E(n-i)
\xrightarrow{\ \iota_X\ }\cdots\longrightarrow
p^*\bigwedge^0\underline E(n)\longrightarrow0
\end{equation}

of $\Omega^i_{\mathbb P(\underline E)/S}(n)$.  The spectral sequence
defined by this resolution degenerates ($E_1=E_\infty$) for degree reasons
by (1.1.9); for $0<j<r$, it gives

\begin{equation}\tag{1.1.11}
R^jp_*\Omega^i_{\mathbb P(\underline E)/S}(n)
=\mathcal H\left(\cdots\xrightarrow{\ \iota_X\ }
p_*\bigwedge^{i-j}\underline E(n-i+j)
\xrightarrow{\ \iota_X\ }\cdots\right).
\end{equation}

Since $r\geq1$, a depth argument shows that

\begin{equation}\tag{1.1.12}
\bigoplus_{n\in\mathbb Z}p_*p^*\bigwedge^i\underline E(n)
=p_*q_*q^*p^*\bigwedge^i\underline E
={a'}_*{a'}^*\bigwedge^i\underline E
=a_*a^*\bigwedge^i\underline E.
\end{equation}


\SGApage{II}{\(\approx\)43}
Thus $p_*p^*\bigwedge^i\underline E(n)$ is the homogeneous part of weight
$n+i$ in $a_*a^*\bigwedge^i\underline E$, namely
$\operatorname{Sym}^n(\underline E)\otimes\bigwedge^i\underline E$.
Remembering that (1.1.5) comes from the Koszul resolution, (1.1.11) shows
that, for $0<j<r$, the sheaf
$R^jp_*\Omega^i_{\mathbb P(\underline E)/S}(n)$ can be nonzero only when
$i-j=0$ and $n-i+j=0$, that is, only when $n=0$ and $i=j$.  In that case,
moreover, $R^ip_*\Omega^i_{\mathbb P(\underline E)/S}$ is isomorphic to
$\underline O_S$.

For $j=0$, (1.1.10) also gives

\SGAvar{\begin{equation}\tag{1.1.13}
p_*\Omega^i_{\mathbb P(\underline E)/S}(n)
=\ker\!\left(\iota_X:
p_*\bigwedge^i\underline E(n-i)\longrightarrow
p_*\bigwedge^{i-1}\underline E(n-i+1)\right.,
\end{equation}

}{\begin{equation}\tag{1.1.13}
p_*\Omega^i_{\mathbb P(\underline E)/S}(n)
=\ker\!\left(\iota_X:
p_*\bigwedge^i\underline E(n-i)\longrightarrow
p_*\bigwedge^{i-1}\underline E(n-i+1)\right),
\end{equation}

}{E7II-XI-002}and the same argument shows that $\iota_X$ is injective when $n-i\leq0$.

By (1.1.8), with $n=0$ and $i=r+1$, one has

\begin{equation*}
\Omega^r_{\mathbb P(\underline E)/S}\simeq
p^*\bigwedge^{r+1}\underline E(-r-1).
\end{equation*}

When $S$ is the spectrum of a field, the vanishing assertions in 1.1(ii)
and (iii) now follow from the preceding argument for $j\ne r$, and follow
from it by duality ([2], V, 11.2, p.~211) for $j=r$.  In particular, for
any $j$ and $n$ there is at most one $i$ for which
$R^ip_*\Omega^j_{\mathbb P(\underline E)/S}(n)\ne0$.  This proves (i),
since $\Omega^j_{\mathbb P(\underline E)/S}$ is flat over $S$ (the hardest
part is locating the reference in EGA III; we suggest 6.10.5, 7.4.1, and
7.5.5, followed by a passage to the limit).

\SGAvar{Knowing that $R^ip_*\Omega^j_{\mathbb P(\underline E)/S}$ is an invertible
sheaf for $0\leq i\leq r$, it remains only to prove that, when $S$ is the
spectrum of a field, $\eta^i\ne0$.  Indeed, $\eta^r$ is the cohomology
class of a rational point (the intersection of $r$ hyperplanes), hence a
basis of $H^r(\mathbb P^r,\Omega^r)$ (FGA, Exposé 149, no.~4).

}{Knowing that $R^ip_*\Omega^i_{\mathbb P(\underline E)/S}$ is an invertible
sheaf for $0\leq i\leq r$, it remains only to prove that, when $S$ is the
spectrum of a field, $\eta^i\ne0$.  Indeed, $\eta^r$ is the cohomology
class of a rational point (the intersection of $r$ hyperplanes), hence a
basis of $H^r(\mathbb P^r,\Omega^r)$ (FGA, Exposé 149, no.~4).

}{E7II-XI-003}1.2. The Euler--Poincaré characteristics of the $\Omega^i(n)$ may be
calculated directly from (1.1.10) or (1.1.13).  The rather unappetizing
result is

\SGAvar{\begin{equation*}
\chi\bigl(\mathbb P^r_k,\Omega^j_{\mathbb P^r_k}(n)\bigr)
=\sum_{a\leq i\leq j}(-1)^{i-j}
\binom{r+1}{i}\binom{r+n-i}{r}.
\end{equation*}




}{\begin{equation*}
\chi\bigl(\mathbb P^r_k,\Omega^j_{\mathbb P^r_k}(n)\bigr)
=\sum_{0\leq i\leq j}(-1)^{i-j}
\binom{r+1}{i}\binom{r+n-i}{r}.
\end{equation*}




}{E7II-XI-004}\SGApage{II}{44}
\underline{Proposition} 1.3. \underline{Let $f:X\to S$ be a proper smooth
morphism, let $\underline W$ be a locally free sheaf of rank $c$ on $X$,
let $\underline F$ be a coherent sheaf on $X$, let $d>0$ be an integer,
and let $s:\underline W\to\underline O_X$ be a section of
$\underline W^\vee$.  Assume that $S$ is noetherian and that:}

\SGAvar{\begin{enumerate}[label=(\alph*),leftmargin=2.4em]
\item \underline{The closed subscheme $H$ of $X$ defined by the vanishing
of $s$ is smooth over $S$.}

\item \underline{Locally on $X$, if the $s_i$ are the coordinates of $S$
in a basis of $\underline W^\vee$, the $s_i$ form both an
$\underline O_X$-regular sequence and an $\underline F$-regular sequence.}

\item \underline{For every sequence of integers $k_\ell$, not all zero,
one has}
\end{enumerate}

}{\begin{enumerate}[label=(\alph*),leftmargin=2.4em]
\item \underline{The closed subscheme $H$ of $X$ defined by the vanishing
of $s$ is smooth over $S$.}

\item \underline{Locally on $X$, if the $s_i$ are the coordinates of $s$
in a basis of $\underline W^\vee$, the $s_i$ form both an
$\underline O_X$-regular sequence and an $\underline F$-regular sequence.}

\item \underline{For every sequence of integers $k_\ell$, not all zero,
one has}
\end{enumerate}

}{E7II-XI-005}\begin{equation*}
R^if_*\left(\bigotimes_\ell\bigwedge^{k_\ell}\underline W
\otimes\Omega^j_{X/S}\otimes\underline F\right)=0
\qquad\underline{\text{for }}i+j<d.
\end{equation*}

\underline{Then:}

\SGAvar{\begin{align*}
R^if_*\left(\Omega^j_{H/S}\otimes\underline F\right)
&\xrightarrow{\ \sim\ }
R^if_*\left(\Omega^j_{X/S}\otimes\underline F\right)
&&\underline{\text{for }}i+j<d-c, \tag{i}\\
R^if_*\left(\Omega^j_{H/S}\otimes\underline F\right)
&\lhook\joinrel\longrightarrow
R^if_*\left(\Omega^j_{X/S}\otimes\underline F\right)
&&\underline{\text{for }}i+j=d-c. \tag{ii}
\end{align*}

}{\begin{align*}
R^if_*\left(\Omega^j_{X/S}\otimes\underline F\right)
&\xrightarrow{\ \sim\ }
R^if_*\left(\Omega^j_{H/S}\otimes\underline F\right)
&&\underline{\text{for }}i+j<d-c, \tag{i}\\
R^if_*\left(\Omega^j_{X/S}\otimes\underline F\right)
&\lhook\joinrel\longrightarrow
R^if_*\left(\Omega^j_{H/S}\otimes\underline F\right)
&&\underline{\text{for }}i+j=d-c. \tag{ii}
\end{align*}

}{E7II-XI-006}The conormal sheaf $N_{H/X}$ of $H$ in $X$ is isomorphic to the restriction
of $\underline W$ to $H$, so the exact sequence

\SGAvar{\begin{equation}\tag{1.3.1}
0\longrightarrow N_{H/X}\longrightarrow
\Omega^1_{X/S}\otimes\underline O_Y\longrightarrow
\Omega^1_{H/S}\longrightarrow0
\end{equation}

defines a decreasing filtration on
$\Omega^j_{X/S}\otimes\underline O_Y$ whose successive quotients are

\begin{equation}\tag{1.3.2}
\operatorname{Gr}^k\bigl(\Omega^j_{X/S}\otimes\underline O_Y\bigr)
\simeq\bigwedge^k\underline W\otimes\Omega^{j-k}_{H/S}.
\end{equation}

}{\begin{equation}\tag{1.3.1}
0\longrightarrow N_{H/X}\longrightarrow
\Omega^1_{X/S}\otimes\underline O_H\longrightarrow
\Omega^1_{H/S}\longrightarrow0
\end{equation}

defines a decreasing filtration on
$\Omega^j_{X/S}\otimes\underline O_H$ whose successive quotients are

\begin{equation}\tag{1.3.2}
\operatorname{Gr}^k\bigl(\Omega^j_{X/S}\otimes\underline O_H\bigr)
\simeq\bigwedge^k\underline W\otimes\Omega^{j-k}_{H/S}.
\end{equation}

}{E7II-XI-007}It follows from (b) that one also has

\SGAvar{\begin{equation}\tag{1.3.3}
\operatorname{Gr}^k\bigl(\Omega^j_{X/S}\otimes\underline O_Y
\otimes\underline F\bigr)
=\bigwedge^k\underline W\otimes\Omega^{j-k}_{H/S}
\otimes\underline F.
\end{equation}

}{\begin{equation}\tag{1.3.3}
\operatorname{Gr}^k\bigl(\Omega^j_{X/S}\otimes\underline O_H
\otimes\underline F\bigr)
=\bigwedge^k\underline W\otimes\Omega^{j-k}_{H/S}
\otimes\underline F.
\end{equation}

}{E7II-XI-008}We prove (i) and (ii) by induction on $j$, simultaneously for all
$\underline F$ satisfying (b) and (c).  The assertions are empty for
$j<0$.


\SGApage{II}{\(\approx\)45}
Suppose, then, that (i) and (ii) have been proved for $j'<j$.  The
induction hypothesis applies to the sheaves
$\bigwedge^k\underline W\otimes\underline F$ and, by (c), gives

\begin{equation}\tag{1.3.4}
R^if_*\left(\bigwedge^k\underline W\otimes
\Omega^{j-k}_{H/S}\otimes\underline F\right)=0
\quad\text{for }k>0\text{ and }i+j<d-c+k.
\end{equation}

\SGAvar{The filtration of
$\Omega^j_{X/S}\otimes\underline O_Y\otimes\underline F$ considered
above defines a spectral sequence

\begin{equation*}
E_1^{p,i-p}=R^if_*\left(\bigwedge^p\underline W\otimes
\Omega^{j-p}_{H/S}\otimes\underline F\right)
\Longrightarrow
R^if_*\left(\Omega^j_{X/S}\otimes\underline O_Y
\otimes\underline F\right).
\end{equation*}

}{The filtration of
$\Omega^j_{X/S}\otimes\underline O_H\otimes\underline F$ considered
above defines a spectral sequence

\begin{equation*}
E_1^{p,i-p}=R^if_*\left(\bigwedge^p\underline W\otimes
\Omega^{j-p}_{H/S}\otimes\underline F\right)
\Longrightarrow
R^if_*\left(\Omega^j_{X/S}\otimes\underline O_H
\otimes\underline F\right).
\end{equation*}

}{E7II-XI-009}By (1.3.4), $E_1^{p,i-p}=0$ when $p>0$ and $i+j\leq d-c$; hence

\SGAvar{\begin{equation}\tag{1.3.5}
\left\{
\begin{aligned}
R^if_*\left(\Omega^j_{X/S}\otimes\underline O_Y\otimes\underline F\right)
&\xrightarrow{\ \sim\ }
R^if_*\left(\Omega^j_{H/S}\otimes\underline F\right)
&&\text{if }i+j<d-c,\\
R^if_*\left(\Omega^j_{X/S}\otimes\underline O_Y\otimes\underline F\right)
&\lhook\joinrel\longrightarrow
R^if_*\left(\Omega^j_{H/S}\otimes\underline F\right)
&&\text{if }i+j=d-c.
\end{aligned}
\right.
\end{equation}

By hypothesis (b), with a suitable differential, the sheaves
$\bigwedge^k\underline W\otimes\Omega^j_{X/S}\otimes\underline F$ form a
left Koszul resolution of
$\Omega^j_{X/S}\otimes\underline O_Y\otimes\underline F$.  This
filtration defines a spectral sequence

\begin{equation}\tag{1.3.6}
E_1^{-p,q}=R^qf_*\left(\bigwedge^p\underline W\otimes
\Omega^j_{X/S}\otimes\underline F\right)
\Longrightarrow
R^{q-p}f_*\left(\Omega^j_{X/S}\otimes\underline O_Y
\otimes\underline F\right).
\end{equation}

}{\begin{equation}\tag{1.3.5}
\left\{
\begin{aligned}
R^if_*\left(\Omega^j_{X/S}\otimes\underline O_H\otimes\underline F\right)
&\xrightarrow{\ \sim\ }
R^if_*\left(\Omega^j_{H/S}\otimes\underline F\right)
&&\text{if }i+j<d-c,\\
R^if_*\left(\Omega^j_{X/S}\otimes\underline O_H\otimes\underline F\right)
&\lhook\joinrel\longrightarrow
R^if_*\left(\Omega^j_{H/S}\otimes\underline F\right)
&&\text{if }i+j=d-c.
\end{aligned}
\right.
\end{equation}

By hypothesis (b), with a suitable differential, the sheaves
$\bigwedge^k\underline W\otimes\Omega^j_{X/S}\otimes\underline F$ form a
left Koszul resolution of
$\Omega^j_{X/S}\otimes\underline O_H\otimes\underline F$.  This
filtration defines a spectral sequence

\begin{equation}\tag{1.3.6}
E_1^{-p,q}=R^qf_*\left(\bigwedge^p\underline W\otimes
\Omega^j_{X/S}\otimes\underline F\right)
\Longrightarrow
R^{q-p}f_*\left(\Omega^j_{X/S}\otimes\underline O_H
\otimes\underline F\right).
\end{equation}

}{E7II-XI-010}By hypothesis (c), $E_1^{-p,q}=0$ for $p>0$ and $q+j<d$, and trivially
for $p\notin[0,c]$.  Consequently,

\begin{equation}\tag{1.3.7}
\left\{
\begin{aligned}
R^if_*\left(\Omega^j_{X/S}\otimes\underline F\right)
&\xrightarrow{\ \sim\ }
R^if_*\left(\Omega^j_{H/S}\otimes\underline F\right)
&&\text{if }i+j<d-c,\\
R^if_*\left(\Omega^j_{X/S}\otimes\underline F\right)
&\lhook\joinrel\longrightarrow
R^if_*\left(\Omega^j_{H/S}\otimes\underline F\right)
&&\text{if }i+j=d-c.
\end{aligned}
\right.
\end{equation}

Assertions (i) and (ii) now follow from (1.3.5) and (1.3.7).




\SGApage{II}{46}
1.4. Let $\underline E$ be a locally free sheaf of rank $r+1$ on a scheme
$S$, let $\mathbb P(\underline E)$ be the corresponding projective bundle,
let $d\leq r$ be an integer, and let
$\underline a=(a_1,\ldots,a_d)$ be a sequence of integers.  A relative
\underline{complete intersection} of multidegree $\underline a$ in
$\mathbb P(\underline E)$ is a closed subscheme which has codimension $d$
in every fiber and which, locally on $S$ for the étale topology, is the
intersection of $d$ hypersurfaces of respective degrees
$a_1,\ldots,a_d$.

By 1.1, hypotheses (b) and (c) of 1.3 hold for
$X/S=\mathbb P(\underline E)/S$, $\underline F=\underline O_X$,

\begin{equation*}
\underline W=\bigoplus_i\underline O(-a_i),
\end{equation*}

and a section $s$ of $\underline W^\vee$ defining a relative complete
intersection $Y\subset\mathbb P(\underline E)$ of multidegree
$\underline a$.

\underline{Theorem} 1.5. \underline{Under the hypotheses of 1.4, let
$n=r-d$ be the relative dimension of $Y$, and let
$\eta\in H^0(S,R^1f_*\Omega^1_{Y/S})$ be the first Chern class of
$\underline O(1)$ in the diagram}

\begin{center}
\begin{tikzcd}[column sep=large,row sep=large]
Y\arrow[rr,hook]\arrow[dr,dash,"f"']\&\&
\mathbb P(\underline E)\arrow[dl,"g"]\\
\&S\&{}
\end{tikzcd}
\qquad .
\end{center}

\underline{If $Y$ is smooth over $S$, then:}

\begin{enumerate}[label=(\roman*),leftmargin=2.6em]
\item \underline{The sheaves $R^if_*\Omega^j_{Y/S}$ are locally free,
and their formation is compatible with arbitrary base change.}

\item \underline{The spectral sequence}
\begin{equation*}
E_1^{p,q}=R^qf_*\Omega^p_{Y/S}\Longrightarrow
R^{p+q}f_*(\Omega^*_{Y/S}),
\end{equation*}
\underline{relating Hodge and de Rham cohomology degenerates
($E_1=E_\infty$); its abutment is locally free, and its formation is
compatible with arbitrary base change.}
\end{enumerate}

\noindent(iii)$_a$ \underline{One has}

\SGAvar{\begin{equation*}
R^if_*\Omega^j_{X/S}=0
\qquad\underline{\text{if }}i\ne j\ \underline{\text{ and }}\ i+j\ne n.
\end{equation*}


}{\begin{equation*}
R^if_*\Omega^j_{Y/S}=0
\qquad\underline{\text{if }}i\ne j\ \underline{\text{ and }}\ i+j\ne n.
\end{equation*}


}{E7II-XI-011}\SGApage{II}{\(\approx\)47}
\noindent(iii)$_b$ \underline{$\eta^i$ is a basis of
$R^if_*\Omega^i_{Y/S}$ if $2i<n$.}

\noindent(iii)$_c$ \underline{If $n<2i\leq2n$, and if $\alpha_i$ is the
basis of $R^if_*\Omega^i_{Y/S}$ dual to the basis $\eta^{n-i}$ of
$R^{n-i}f_*\Omega^{n-i}_{Y/S}$, then}

\begin{equation*}
\eta^i=\left(\prod_{h=1}^d a_h\right)\alpha_i.
\end{equation*}

\noindent(iii)$_d$ \underline{If $n=2m$, then $\eta^m\ne0$ (provided
$S\ne\varnothing$), and the quotient
$R^mf_*\Omega^m_{Y/S}/\underline O_S\eta^m$ is locally free.}

It is enough to prove 1.5 in a ``universal'' case, which permits us to
assume that $S$ is smooth over $\operatorname{Spec}(\mathbb Z)$.  In this
case (ii) follows from (i), since it holds over $\mathbb C$, and therefore
over every field of characteristic zero, by Hodge theory.  Moreover, since
$S$ is reduced, (i) follows from (iii), assumed proved when $S$ is the
spectrum of a field.  Indeed, by (iii), for each $j$ the functions

\begin{equation*}
s\longmapsto\dim_{k(s)}H^i(Y_s,\Omega^j_{Y_s})
\end{equation*}

\SGAvar{are constant on $S$, except perhaps when $i=j$; hence they are all constant,
because their alternating sum is constant (EGA III, 7.9.2).  The conclusion
then follows from EGA III, 6.10.5, 7.6.9, and 7.7.5.

}{are constant on $S$, except perhaps when $i+j=n$; hence they are all constant,
because their alternating sum is constant (EGA III, 7.9.2).  The conclusion
then follows from EGA III, 6.10.5, 7.6.9, and 7.7.5.

}{E7II-XI-012}By 1.4 and 1.3, one has

\begin{equation}\tag{1.5.1}
\left\{
\begin{aligned}
R^ig_*\Omega^j_{\mathbb P(\underline E)/S}
&\xrightarrow{\ \sim\ }R^if_*\Omega^j_{Y/S}
&&\text{for }i+j<n,\\
R^ig_*\Omega^j_{\mathbb P(\underline E)/S}
&\lhook\joinrel\longrightarrow R^if_*\Omega^j_{Y/S}
&&\text{for }i+j=n.
\end{aligned}
\right.
\end{equation}

When $S$ is the spectrum of a field, assertions (iii)$_a$, (iii)$_b$, and
(iii)$_d$ follow from (1.5.1) and Serre duality ([2], V, 11.2, p.~211).
The cohomology class


\SGApage{II}{\(\approx\)48}
$\eta^n$ is the class of a zero-cycle of degree
$\prod_{h=1}^d a_h$ (the degree of the complete intersection), so

\begin{equation*}
\eta^n=\left(\prod_{h=1}^d a_h\right)\alpha_n.
\end{equation*}

By definition, $\eta^{n-i}\alpha_i=\alpha_n$, and (iii)$_c$ follows.

Thus (i) and (ii) hold; in the universal case, the assertions in (iii)
follow from (i) and from (iii) in the case already proved, where $S$ is the
spectrum of a field.

\underline{Theorem} 1.6. \underline{Keep the hypotheses and notation of
1.4 and 1.5, assume that $Y$ is smooth, let $\ell$ be a prime invertible in
all residue fields of $S$, and let
$\eta_\ell\in H^0(S,R^2f_*\mathbb Z_\ell(1))$ denote the $\ell$-adic first
Chern class of $\underline O(1)$ (SGA 5, VII).  Then:}

\begin{enumerate}[label=(\roman*),leftmargin=2.6em]
\item \underline{$R^if_*\mathbb Z_\ell=0$ if $i$ is odd and $i\ne n$.}

\item \underline{If $2i<n$, then $R^{2i}f_*\mathbb Z_\ell(i)$ is
isomorphic to $\mathbb Z_\ell$ and has $\eta_\ell^i$ as a basis.}

\item \underline{If $n<2i\leq2n$, then
$R^{2i}f_*\mathbb Z_\ell(i)$ is isomorphic to $\mathbb Z_\ell$ and has}
\begin{equation*}
\eta_\ell^i\bigg/\prod_{h=1}^d a_h
\end{equation*}
\underline{as a basis.}

\item \underline{The $\mathbb Z_\ell$-sheaf $R^nf_*\mathbb Z_\ell$ is
constructible, twisted constant, and torsion-free.  If $n$ is even, then
$\eta_\ell^{n/2}\ne0$ (provided $S\ne\varnothing$), and
$R^nf_*\mathbb Z_\ell(n/2)/\mathbb Z_\ell\eta_\ell^{n/2}$ is torsion-free.}
\end{enumerate}

The $\mathbb Z_\ell$-sheaves $R^if_*\mathbb Z_\ell$ are constructible and
twisted constant, and their formation is compatible with arbitrary base
change, by SGA 4, XVI, 2.2.  It therefore suffices to prove 1.6 when $S$ is
the spectrum of an algebraically closed field.  Repeated application of
SGA 4, XIV, 3.3 then shows, for every $m$, that the canonical map

\SGAvar{\begin{equation*}
H_Y^q\bigl(\mathbb P(\underline E),\mathbb Z/\ell^m\bigr)
\longrightarrow H^q(Y,\mathbb Z/\ell^m)
\end{equation*}


}{\begin{equation*}
H_Y^q\bigl(\mathbb P(\underline E),\mathbb Z/\ell^m\bigr)
\longrightarrow H^q\bigl(\mathbb P(\underline E),\mathbb Z/\ell^m\bigr)
\end{equation*}


}{E7II-XI-013}\SGApage{II}{\(\approx\)49}
is bijective for $q>r+d$ and surjective for $q=r+d$.  From this and
Poincaré duality (SGA 4, XVIII), one deduces (SGA 5, VII) that the
restriction map

\begin{equation*}
H^q\bigl(\mathbb P(\underline E),\mathbb Z_\ell\bigr)
\longrightarrow H^q(Y,\mathbb Z_\ell)
\end{equation*}

is bijective for $q<n$, and injective with torsion-free cokernel for $q=n$.
This proves (i) for $i<n$, as well as (ii) and (iv).  As in 1.5, one checks
that $\eta_\ell^n$ is $\prod_{h=1}^d a_h$ times the canonical class; (i)
and (iii) then follow by Poincaré duality (SGA 4, XVIII).



1.7. If $Y$ is a complete intersection of dimension at least three in a
projective space $\mathbb P^r(k)$ over a field $k$, then
$\operatorname{Pic}(Y)$ is known to be isomorphic to $\mathbb Z$ and to be
generated by $\underline O(1)$ (SGA 2, XII, 3.7).  In dimension two, when
$Y$ is smooth, one can still prove the following.

\underline{Theorem} 1.8. \underline{Let $Y$ be a smooth complete-intersection
surface in the projective space $\mathbb P^r(k)$ over an algebraically
closed field $k$.  Then $\underline{\operatorname{Pic}}^\tau(Y)=0$, and the
class of $\underline O(1)$ is not divisible in the Néron--Severi group
$\operatorname{NS}(Y)$.}

Since $H^1(Y,\underline O)=0$ by 1.5(iii)$_a$, one knows that
$\underline{\operatorname{Pic}}^0(Y)=0$.  The vanishing of
$\underline{\operatorname{Pic}}^\tau(Y)$ then follows from the familiar
lemma below, the fact that $H^2(Y,\mathbb Z_\ell(1))$ is torsion-free by
1.6(iv), and the fact that $H^0(Y,\Omega_Y^1)=0$ by 1.5(iii)$_a$.

\underline{Lemma} 1.9. \underline{Let $Y$ be a complete algebraic variety
over an algebraically closed field $k$.}

\begin{enumerate}[label=(\roman*),leftmargin=2.6em]
\item \underline{If} $\ell$ \underline{is prime to the characteristic of $k$, then the}
$\ell$-\underline{torsion of $\operatorname{NS}(Y)$ identifies with the} $\ell$-\underline{torsion
of $H^2(Y,\mathbb Z_\ell(1))$.}


\item \underline{If $k$ has nonzero characteristic $p$ and $Y$ is smooth,
then the kernel of multiplication by $p$ in $\operatorname{Pic}(Y)$
identifies with the subgroup of $H^0(Y,\Omega_Y^1)$ consisting of locally
logarithmic differentials.}
\end{enumerate}

\SGApage{II}{\(\approx\)50}
The Kummer exact sequences

\begin{equation*}
0\longrightarrow\mu_{\ell^n}\longrightarrow\mathbb G_m
\xrightarrow{\ x\mapsto x^{\ell^n}\ }\mathbb G_m\longrightarrow0
\end{equation*}

give exact sequences

\begin{equation*}
0\longrightarrow
\operatorname{Pic}(Y)/\ell^n\operatorname{Pic}(Y)
\longrightarrow H^2(Y,\mu_{\ell^n})
\longrightarrow H^2(Y,\mathbb G_m)_{\ell^n}
\longrightarrow0,
\end{equation*}

where

\begin{equation*}
\operatorname{Pic}(Y)/\ell^n\operatorname{Pic}(Y)
=\operatorname{NS}(Y)\otimes\mathbb Z/\ell^n.
\end{equation*}

Passing to the inverse limit gives

\begin{equation*}
0\longrightarrow\operatorname{NS}(Y)\otimes\mathbb Z_\ell
\longrightarrow H^2(Y,\mathbb Z_\ell(1))
\longrightarrow
\operatorname{Hom}\bigl(\mathbb Q_\ell/\mathbb Z_\ell,
H^2(Y,\mathbb G_m)\bigr)
\longrightarrow0.
\end{equation*}

The cokernel is torsion-free because $\mathbb Q_\ell/\mathbb Z_\ell$ is
divisible.  This proves (i).

To prove (ii), it is enough to observe that the exact sequence

\begin{equation*}
0\longrightarrow\underline O_Y^*
\xrightarrow{\ x\mapsto x^p\ }\underline O_Y^*
\xrightarrow{\ df/f\ }\Omega_Y^1
\end{equation*}

defines a short exact sequence

\begin{equation*}
0\longrightarrow\underline O_Y^*\longrightarrow\underline O_Y^*
\longrightarrow\Omega^1_{Y,\mathrm{loc}\,\log}\longrightarrow0,
\end{equation*}

and then to take the corresponding long exact cohomology sequence.  This
argument is due to Cartier, who proves moreover that, if $C$ denotes the
Cartier operator on $H^0(Y,\Omega_Y^1)$, then the locally logarithmic
differentials form the kernel of $C-\operatorname{Id}$.

Suppose, for a contradiction, that the class of $\underline O(1)$ in
$\operatorname{NS}(Y)$ has the form $nu$, with $n>1$.  If $\ell$ divides
$n$ and is prime to the characteristic, then in $\ell$-adic cohomology one
again has


\SGApage{II}{\(\approx\)51}
\begin{equation*}
\eta_\ell=n\,c_1(u),
\end{equation*}

contradicting 1.6(iv).

Finally, if the characteristic $p$ of $k$ is nonzero and $p$ divides $n$,
then in Hodge cohomology one again has

\begin{equation*}
\eta=n\,c_1(u)=0,
\end{equation*}

contradicting 1.5(iii)$_d$ and completing the proof.



2. \underline{Numerical results}

2.1. Let $\underline a=(a_i)_{1\leq i\leq d}$ be a sequence of integers
$\geq1$, and let $V_n(\underline a)$ denote a smooth complete intersection
of dimension $n$ and multidegree $\underline a$ in a projective space
$\mathbb P^{n+d}$ over a field $k$.  Put

\begin{equation*}
h^{p,q}(\underline a)
=\dim_k H^q\bigl(V_{p+q}(\underline a),\Omega^p\bigr)
\end{equation*}

and

\begin{equation*}
h_0^{p,q}(\underline a)=h^{p,q}(\underline a)-\delta_{p,q}.
\end{equation*}

Thus the $h_0^{p,q}(\underline a)$ are the Hodge numbers of the
middle-dimensional primitive cohomology of a complete intersection of
multidegree $\underline a$.  They depend only on $\underline a$, $p$, and
$q$, by 1.5.  Set

\begin{equation}\tag{2.1.1}
H(\underline a)=\sum_{p,q\geq0}h_0^{p,q}(\underline a)y^pz^q
\in\mathbb Z[[y,z]].
\end{equation}

For $d=1$ and $\underline a=(a)$, put
$h^{p,q}(a)=h^{p,q}(\underline a)$ and $H(a)=H(\underline a)$.  For $d=0$,
put $\sup(a_i)=1$.

The following formula is the special case $k=0$ of formula (2), p.~160,
in Hirzebruch [3]:


\SGApage{II}{\(\approx\)52}
\SGAvar{\begin{equation}\tag{2.2}
\sum_{n=0}^{\infty}\sum_p
\chi\bigl(V_n(\underline a),\Omega^p\bigr)y^pz^{n+d}
=\frac{1}{(1+yz)(1-z)}
\prod_{i=1}^d
\frac{(1+yz)^{a_i}-(1-z)^{a_i}}
{(i+yz)^{a_i}z+y(1-z)^{a_i}}.
\end{equation}

}{\begin{equation}\tag{2.2}
\sum_{n=0}^{\infty}\sum_p
\chi\bigl(V_n(\underline a),\Omega^p\bigr)y^pz^{n+d}
=\frac{1}{(1+yz)(1-z)}
\prod_{i=1}^d
\frac{(1+yz)^{a_i}-(1-z)^{a_i}}
{(1+yz)^{a_i}z+y(1-z)^{a_i}}.
\end{equation}

}{E7II-XI-014}We shall transform this into the following formula.

\underline{Theorem} 2.3 (Hirzebruch). \underline{With the notation of 2.1,
one has}

\begin{equation*}
H(\underline a)=\frac{1}{(1+y)(1+z)}\left[
\prod_{i=1}^d
\frac{(1+y)^{a_i}-(1+z)^{a_i}}
{-(1+y)^{a_i}z+(1+z)^{a_i}y}-1\right].
\end{equation*}

By 1.5,

\begin{equation*}
\chi\bigl(V_n(\underline a),\Omega^p\bigr)
=(-1)^p+(-1)^{n-p}h_0^{p,n-p}(\underline a).
\end{equation*}

Consequently, the left-hand side $I$ of (2.2) can be rewritten as

\SGAvar{\begin{align*}
I
&=\sum_{n=0}^{\infty}\sum_{p+q=n}
\left((-1)^p+(-1)^q h_0^{p,q}(\underline a)\right)y^pz^{p+q+d}\\
&=\sum_{p,q\geq0}(-1)^p(yz)^pz.^qz^d
+\sum_{p,q\geq0}(-1)^q h_0^{p,q}(\underline a)(yz)^pz.^qz^d\\
&=z^d\left[\frac{1}{(1+yz)(1-z)}
+H(\underline a)(yz,-z)\right].
\end{align*}

}{\begin{align*}
I
&=\sum_{n=0}^{\infty}\sum_{p+q=n}
\left((-1)^p+(-1)^q h_0^{p,q}(\underline a)\right)y^pz^{p+q+d}\\
&=\sum_{p,q\geq0}(-1)^p(yz)^pz^qz^d
+\sum_{p,q\geq0}(-1)^q h_0^{p,q}(\underline a)(yz)^pz^qz^d\\
&=z^d\left[\frac{1}{(1+yz)(1-z)}
+H(\underline a)(yz,-z)\right].
\end{align*}

}{E7II-XI-015}The right-hand side $II$ of (2.2) can be rewritten as

\begin{equation*}
II=\frac{z^d}{(1+yz)(1-z)}
\prod_{i=1}^d
\frac{(1+yz)^{a_i}-(1-z)^{a_i}}
{(1+yz)^{a_i}z+yz(1-z)^{a_i}}.
\end{equation*}

After cancelling $z^d$ in $I=II$, one obtains

\begin{equation*}
H(\underline a)(yz,-z)
=\frac{1}{(1+yz)(1-z)}\left[
\prod_{i=1}^d
\frac{(1+yz)^{a_i}-(1-z)^{a_i}}
{(1+yz)^{a_i}z+(1-z)^{a_i}yz}-1\right],
\end{equation*}

and 2.3 follows by a change of variables.


\SGApage{II}{53}
\underline{Corollary} 2.4. \underline{With the notation of 2.1:}

\SGAvar{\begin{enumerate}[label=(\roman*),leftmargin=2.6em]
\item \underline{For $d=1$, one has}
\begin{equation*}
\begin{aligned}
H(a)
&=-\frac{(1+y)^{a-1}-(1+z)^{a-1}}
{(1+y)^az-(1+z)^ay}\\[1ex]
&=\frac{\displaystyle\sum_{i,j\geq0}
\binom{a-1}{i+j+1}y^iz^j}
{\displaystyle1-\sum_{i,j\geq1}\binom{a}{i+j}y^iz^j}.
\end{aligned}
\end{equation*}

\item \underline{Writing $|P|$ for the number of elements of a subset
$P$ of $[1,d]$, one has}
\begin{equation*}
H(\underline a)=
\sum_{\substack{P\subset[i,d]\\P\ne\varnothing}}
(1+y)^{|P|-1}(1+z)^{|P|-1}
\prod_{i\in P}H(a_i).
\end{equation*}
\end{enumerate}

}{\begin{enumerate}[label=(\roman*),leftmargin=2.6em]
\item \underline{For $d=1$, one has}
\begin{equation*}
\begin{aligned}
H(a)
&=-\frac{(1+y)^{a-1}-(1+z)^{a-1}}
{(1+y)^az-(1+z)^ay}\\[1ex]
&=\frac{\displaystyle\sum_{i,j\geq0}
\binom{a-1}{i+j+1}y^iz^j}
{\displaystyle1-\sum_{i,j\geq1}\binom{a}{i+j}y^iz^j}.
\end{aligned}
\end{equation*}

\item \underline{Writing $|P|$ for the number of elements of a subset
$P$ of $[1,d]$, one has}
\begin{equation*}
H(\underline a)=
\sum_{\substack{P\subset[1,d]\\P\ne\varnothing}}
(1+y)^{|P|-1}(1+z)^{|P|-1}
\prod_{i\in P}H(a_i).
\end{equation*}
\end{enumerate}

}{E7II-XI-016}By 2.3,

\SGAvar{\begin{align*}
H(a)
&=-\frac{1}{(1+y)(1+z)}\left[
\frac{(1+z)^a-(1+z)^a}{(1+z)^az-(1+z)^ay}+1\right]\\
&=-\frac{1}{(1+y)(1+z)}\left[
\frac{(1+y)^a(1+z)-(1+z)^a(1+y)}
{(1+y)^az-(1+z)^ay}\right]\\
&=-\frac{(1+y)^{a-1}-(1+z)^{a-1}}
{(1+y)^az-(1+z)^ay}.
\end{align*}

}{\begin{align*}
H(a)
&=-\frac{1}{(1+y)(1+z)}\left[
\frac{(1+y)^a-(1+z)^a}{(1+y)^az-(1+z)^ay}+1\right]\\
&=-\frac{1}{(1+y)(1+z)}\left[
\frac{(1+y)^a(1+z)-(1+z)^a(1+y)}
{(1+y)^az-(1+z)^ay}\right]\\
&=-\frac{(1+y)^{a-1}-(1+z)^{a-1}}
{(1+y)^az-(1+z)^ay}.
\end{align*}

}{E7II-XI-017}Let $N$ and $D$ be the numerator and denominator in this last expression,
so that $H(a)=-N/D$.


\SGApage{II}{\(\approx\)54}
One has

\begin{equation*}
N=\sum_k\binom{a-1}{k}(y^k-z^k)
=(y-z)\sum_{i,j\geq0}\binom{a-1}{i+j+1}y^iz^j
\end{equation*}

and

\SGAvar{\begin{align*}
D
&=\sum_k\binom ak(y^kz-yz^k)
=(z-y)+yz\sum_{k\geq1}\binom ak(y^{k-1}-z^{k-1})\\
&=(y-z)\left(-1+yz\sum_{i,j\geq0}
\binom a{i+j+2}y^iz^j\right.\\
&=(y-z)\left(-1+\sum_{i,j\geq1}
\binom a{i+j}y^iz^j\right).
\end{align*}

}{\begin{align*}
D
&=\sum_k\binom ak(y^kz-yz^k)
=(z-y)+yz\sum_{k\geq1}\binom ak(y^{k-1}-z^{k-1})\\
&=(y-z)\left(-1+yz\sum_{i,j\geq0}
\binom a{i+j+2}y^iz^j\right)\\
&=(y-z)\left(-1+\sum_{i,j\geq1}
\binom a{i+j}y^iz^j\right).
\end{align*}

}{E7II-XI-018}This proves (i).

Now invert the formula in 2.3 when $d=1$.  One finds

\SGAvar{\begin{equation*}
-\frac{(1+y)^a-(1+z)^a}{(1+y)^az-(1+z)^ay}
=1+(1+x)(1+y)H(a).
\end{equation*}

}{\begin{equation*}
-\frac{(1+y)^a-(1+z)^a}{(1+y)^az-(1+z)^ay}
=1+(1+y)(1+z)H(a).
\end{equation*}

}{E7II-XI-019}Substituting this identity into 2.3 gives

\SGAvar{\begin{align*}
H(\underline a)
&=\frac{1}{(1+y)(1+z)}\left[
\prod_{i=1}^d\left(1+(1+x)(1+y)H(a_i)\right)-1\right]\\
&=\frac{1}{(1+y)(1+z)}\left[
\sum_{\substack{P\subset[i,d]\\P\ne\varnothing}}
(1+x)^{|P|}(1+y)^{|P|}\prod_{i\in P}H(a_i)\right],
\end{align*}

}{\begin{align*}
H(\underline a)
&=\frac{1}{(1+y)(1+z)}\left[
\prod_{i=1}^d\left(1+(1+y)(1+z)H(a_i)\right)-1\right]\\
&=\frac{1}{(1+y)(1+z)}\left[
\sum_{\substack{P\subset[1,d]\\P\ne\varnothing}}
(1+y)^{|P|}(1+z)^{|P|}\prod_{i\in P}H(a_i)\right],
\end{align*}

}{E7II-XI-020}which proves (ii).


\underline{Theorem} 2.5. \underline{With the notation of 2.1:}

\begin{enumerate}[label=(\roman*),leftmargin=2.6em]
\item \underline{If $p+q=p'+q'$ and
$p\leq p'\leq q'\leq q$, then}
\begin{equation*}
h_0^{p,q}(\underline a)\leq h_0^{p',q'}(\underline a).
\end{equation*}

\item \underline{If $p\leq q$, then $h_0^{p,q}(\underline a)\ne0$ if
and only if}
\begin{equation*}
p>\left\lfloor
\frac{p+q+d-\sum_i a_i}{\sup_i(a_i)}
\right\rfloor.
\end{equation*}
\underline{Here $\lfloor\ \rfloor$ denotes the integer part of the
quotient}.
\end{enumerate}


\SGApage{II}{\(\approx\)55}
Let (*) be the following property of a formal power series
$P=\sum a_{p,q}y^pz^q$:

\begin{quote}
(*) the $a_{p,q}$ are nonnegative, $a_{p,q}=a_{q,p}$, and whenever
$p+q=p'+q'$ and $p\leq p'\leq q'\leq q$, one has
$a_{p,q}\leq a_{p',q'}$.
\end{quote}

Assertion (i) follows at once from the formulas in 2.4, the identity
$(1-x)^{-1}=\sum_{n\geq0}x^n$, and the following lemma, whose proof is left
to the reader.

\underline{Lemma} 2.5.1. \underline{The set of formal power series
satisfying (*) is stable under addition and multiplication}.

We first prove (ii) in the special case of hypersurfaces ($d=1$).  For
$p\leq q$, the desired formula can then be rewritten as

\begin{equation}\tag{2.5.2}
h_0^{p,q}(a)\ne0
\quad\Longleftrightarrow\quad
p\geq\left\lfloor\frac{p+q+1}{a}\right\rfloor
\end{equation}

or, equivalently, $(p+1)a\geq p+q+2$.  From the second formula in 2.4(i),
one sees that $h_0^{p,q}(a)$ is nonzero if and only if there exist pairs
$(e,f)$ and $(c_i,d_i)$ such that

\begin{equation}\tag{2.5.3}
\left\{
\begin{aligned}
(e,f)&\geq(0,0),& e+f&\leq a-2,\\
(c_i,d_i)&\geq(1,1),& c_i+d_i&\leq a,\\
(p,q)&=(e,f)+\sum_i(c_i,d_i).&&
\end{aligned}
\right.
\end{equation}


\SGApage{II}{\(\approx\)56}
For fixed $p+q$, the smallest possible $p$ is obtained by taking $e=0$
and every $c_i=1$; it is the number of indices $i$.  Thus it is the lower
bound of those $p$ for which

\begin{equation*}
p\geq\frac{(p+q)-(a-2)}{a},
\qquad\text{i.e.}\qquad (p+1)a\geq p+q+2.
\end{equation*}

We now pass to the general case.  If $\underline a'$ is obtained from
$\underline a$ by deleting the $a_i$ equal to one, then
$H(\underline a)=H(\underline a')$.  In view of the form of 2.5(ii), we may
therefore assume that every $a_i>1$.  If $d=0$, then
$h_0^{p,q}=0$ and the assertion is correct; we lose no generality by
assuming $d>0$ and $a_d=\sup_i(a_i)$.

\SGAvar{By 2.4(ii) and (2.5.2), the condition
$h_0^{p,q}(\underline a)\ne0$ is equivalent to the existence of a nonempty
subset $P$ of $[1,d]$, a pair $(e,f)$, and pairs
$(c_i,d_i)_{i\leqslant i \leqslant d}$ such that

\begin{equation}\tag{2.5.4}
\left\{
\begin{aligned}
(|P|-1,|P|-1)&\geq(e,f)\geq(0,0),\\
c_i,d_i&\geq
\left\lfloor\frac{c_i+d_i+1}{a_i}\right\rfloor,
& (c_i,d_i)&\geq(0,0),\\
(p,q)&=(e,f)+\sum(c_i,d_i).&&
\end{aligned}
\right.
\end{equation}

}{By 2.4(ii) and (2.5.2), the condition
$h_0^{p,q}(\underline a)\ne0$ is equivalent to the existence of a nonempty
subset $P$ of $[1,d]$, a pair $(e,f)$, and pairs
$(c_i,d_i)_{i\in P}$ such that

\begin{equation}\tag{2.5.4}
\left\{
\begin{aligned}
(|P|-1,|P|-1)&\geq(e,f)\geq(0,0),\\
c_i,d_i&\geq
\left\lfloor\frac{c_i+d_i+1}{a_i}\right\rfloor,
& (c_i,d_i)&\geq(0,0),\\
(p,q)&=(e,f)+\sum_{i\in P}(c_i,d_i).&&
\end{aligned}
\right.
\end{equation}

}{E7II-XI-021}For fixed $n=p+q$, seek the smallest possible value of $p$.  Put
$n_i=c_i+d_i$.  For such a $p$, conditions (2.5.4) become


\SGAvar{\begin{equation}\tag{2.5.5}
\left\{
\begin{aligned}
a_i=2&\Longrightarrow n_i\text{ is even}
&&\text{(otherwise (2.5.4) is impossible)},\\
c_i&\geq\left\lfloor\frac{n_i+1}{a_i}\right\rfloor,
&c_i&\leq n_i,\\
(|P|-1,|P|-1)&\geq(e,f)\geq(0,0),&&\\
n&=e+f+\sum n_i,&&\\
p&=e+\sum c_i.&&
\end{aligned}
\right.
\end{equation}

}{\begin{equation}\tag{2.5.5}
\left\{
\begin{aligned}
a_i=2&\Longrightarrow n_i\text{ is even}
&&\text{(otherwise (2.5.4) is impossible)},\\
c_i&\geq\left\lfloor\frac{n_i+1}{a_i}\right\rfloor,
&c_i&\leq n_i,\\
(|P|-1,|P|-1)&\geq(e,f)\geq(0,0),&&\\
n&=e+f+\sum_{i\in P}n_i,&&\\
p&=e+\sum_{i\in P}c_i.&&
\end{aligned}
\right.
\end{equation}

}{E7II-XI-022}\SGApage{II}{\(\approx\)57}
This description first shows that the smallest $p$ is obtained for
$P=[1,d]$, to which case we henceforth restrict.  If

\begin{equation}\tag{2.5.6}
\sum_i(a_i-2)+(d-1)\geq n,
\end{equation}

then a solution exists with $p=c_i=e=0$.  Otherwise, if $a_d\ne2$, the
best choice is $n_i=a_i-2$ for $i\ne d$, so that $c_i=0$ for $i\ne d$,
$e=0$, $f=d-1$, and

\begin{equation*}
n_d=n-(d-1)-\sum_{i\ne d}(a_i-2).
\end{equation*}

Then

\begin{equation}\tag{2.5.7}
p=\left\lfloor
\frac{n-(d-1)-\sum_{i\ne d}(a_i-2)+1}{a_d}
\right\rfloor,
\end{equation}

and hence, for $p\leq q$, one has
$h_0^{p,q}(\underline a)\ne0$ if and only if

\begin{equation}\tag{2.5.8}
p\geq\left\lfloor
\frac{n-\sum_{i\ne d}(a_i-1)+1}{a_d}
\right\rfloor.
\end{equation}

Formula (2.5.8) reduces to (2.5.2) for $d=1$ and is trivially true when
(2.5.6) holds.  If $d\geq2$, (2.5.6) fails, and all the $a_i$ equal two,
the best choice is instead


\SGApage{II}{\(\approx\)58}
$n_i=0$ for $i\ne d$, $e=0$, and $f=d-1$ or $d-2$, chosen congruent to
$n$ modulo two, with $n_d=n-f$.  If $n\equiv d-1\pmod2$, this proves
(2.5.7) as above.  Otherwise, observe that

\begin{equation*}
\left\lfloor\frac{n-(d-2)+1}{2}\right\rfloor
=\left\lfloor\frac{n-(d-1)+1}{2}\right\rfloor,
\end{equation*}

so (2.5.7), and hence (2.5.8), remains valid when $a_d=2$.

It remains to verify that (2.5.8) is equivalent to 2.5(ii).  Indeed,

\begin{equation*}
\left\lfloor\frac{n-\sum_{i\ne d}(a_i-1)+1}{a_d}\right\rfloor
=\left\lfloor\frac{n+d-\sum_{i\ne d}a_i}{a_d}\right\rfloor
=\left\lfloor\frac{p+q+d-\sum_i a_i}{a_d}\right\rfloor+1.
\end{equation*}

2.6. If $X$ is a complete intersection of dimension $n$ in
$\mathbb P^r$, it follows at once from (2.5.8) that the primitive part of
$H^n(X)$ vanishes only when $X$ is an odd-dimensional quadric.



\underline{Definition} 2.7. \underline{If $X$ is a proper smooth algebraic
variety over $\mathbb C$, the Hodge level of $H^n(X)$ is the supremum of
the numbers $q-p$ for which $p+q=n$ and $h^{p,q}(X)\ne0$.}

\SGAvar{The Hodge level $\ell$ is therefore congruent to $n$ modulo two, and is at
most
\[
\min\bigl(n,\dim(X)-n\bigr).
\]
It is nonnegative unless $H^n(X)=0$, in which case it is $-\infty$.


}{The Hodge level $\ell$ is therefore congruent to $n$ modulo two, and is at
most
\[
\min\bigl(n,2\dim(X)-n\bigr).
\]
It is nonnegative unless $H^n(X)=0$, in which case it is $-\infty$.


}{E7II-XI-023}\underline{Corollary} 2.8. \underline{Retain the notation of 2.1, put
$k=\mathbb C$, and suppose that $a_d=\sup_i(a_i)$.  Let $\ell$ be an
integer with $0\leq\ell<n$, congruent to $n$ modulo two.  The Hodge level
of $H^n(V_n(\underline a))$ is at most $\ell$ if and only if}


\begin{equation*}
\sum_{i\ne d}(a_i-1)
\leq (\ell+1)-(a_d-2)\frac{n-\ell}{2}.
\end{equation*}

\SGApage{II}{\(\approx\)59}
By (2.5.8), if the Hodge level is nonnegative, it is the maximum of zero
and

\begin{equation*}
n-2\left\lfloor
\frac{n-\sum_{i\ne d}(a_i-1)+1}{a_d}
\right\rfloor.
\end{equation*}

The inequality

\begin{equation*}
n-2\left\lfloor
\frac{n-\sum_{i\ne d}(a_i-1)+1}{a_d}
\right\rfloor\leq\ell
\end{equation*}

may be rewritten as

\begin{equation*}
\frac{n-\ell}{2}
\leq
\left\lfloor
\frac{n-\sum_{i\ne d}(a_i-1)+1}{a_d}
\right\rfloor,
\end{equation*}

or equivalently as

\begin{equation*}
\sum_{i\ne d}(a_i-1)
\leq 1+n-a_d\frac{n-\ell}{2}
=\ell+1-(a_d-2)\frac{n-\ell}{2}.
\end{equation*}

2.9. Using formula 2.8, it is easy, for small $\ell$, to enumerate the
multidegrees of complete intersections in $\mathbb P^r$ of dimension
$n>1$ whose $H^n$ has level $\ell$.  The table given in the first version
of this exposé was, however, incorrect; the table below is extracted from
the more complete one given by M. Rapoport in his article
``Complément à l'article de P. Deligne \emph{La conjecture de Weil pour
les surfaces K3}'', \emph{Inventiones Math.} \underline{15}, 227--236
(1972).  For $\ell\leq1$, the complete intersections
$V_n(\underline a)$ of dimension $n$ and multidegree
$\underline a=(a_1,\ldots,a_d)$, with $n>\ell$, $d\geq1$, and $a_i\geq2$,
having Hodge level $\ell$, are the following:


\SGApage{II}{\(\approx\)60}
\begin{equation*}
\begin{array}{ll}
\ell=-\infty & V_{2k+1}(2) \qquad (k\geq0),\\[2ex]
\ell=0 & V_{2k}(2)\ (k\geq1),\quad
          V_{2k}(2,2)\ (k\geq1),\quad V_2(3),\\[2ex]
\ell=1 & V_{2k+1}(2,2)\ (k\geq1),\quad
          V_{2k+1}(2,2,2)\ (k\geq1),\\
& V_3(3),\quad V_5(3),\quad V_3(3,2),\quad V_3(4).
\end{array}
\end{equation*}

2.10. For small $\underline a$, the generating functions of the
$h_0^{p,q}(\underline a)$ are, by 2.4,

\begin{equation*}
\begin{array}{rcl}
H(2)&=&\dfrac{1}{1-yz},\\[3ex]
H(3)&=&\dfrac{2+y+z}{1-3yz-y^2z-yz^2},\\[3ex]
H(2,2)&=&\dfrac{y+z}{(1-yz)^2}
          +\dfrac{2}{(1-yz)^2}+\dfrac{1}{1-yz},
\end{array}
\end{equation*}

so that

\begin{equation*}
\begin{array}{rcll}
h_0^{p,p}(2)&=&1&,\\[2ex]
h_0^{1,1}(3)&=&6&,\\[2ex]
h_0^{p,p}(2,2)&=&2p+3&.
\end{array}
\end{equation*}


\bigskip
\begin{center}
\underline{BIBLIOGRAPHY}
\end{center}

\begin{enumerate}[label={[\arabic*]},leftmargin=3em]
\item Y. Akizuki and S. Nakano, ``Note on Kodaira--Spencer's proof of
Lefschetz theorems'', \emph{Proc. Jap. Acad.} 30, 266--272 (1954).

\item R. Hartshorne, \emph{Residues and Duality}, Lecture Notes in
Mathematics 20, Springer, 1966.

\item F. Hirzebruch, \emph{Topological Methods in Algebraic Geometry},
3rd ed., Springer-Verlag, 1966.

\item K. Kodaira, ``On a differential geometric method in the theory of
analytic stacks'', \emph{Proc. Nat. Acad. Sci. USA} 39, 1268--1273
(1953).
\end{enumerate}


\markboth{Exposé XI}{Exposé XI}

\ifnum\SGAform=1
\clearpage\section*{Critical apparatus}\begingroup\small\raggedright\sloppy
Readings are given in the exact TeX notation of the two English witnesses. Decimal suffixes identify the individual differences in each marked passage.

\par\medskip\noindent\hypertarget{app-E7II-XI-001.1}{}\textbf{E7II-XI-001.1}\quad printed folio 42.\par
\noindent Printed reading: {\ttfamily\footnotesize Kozul\allowbreak{}}\par
\noindent Corrected reading: {\ttfamily\footnotesize Koszul\allowbreak{}}\par
\noindent Source: {\ttfamily\footnotesize port/decisions/decis\allowbreak{}ions\_en7II\_r2\_001.js\allowbreak{}on\#C512;\allowbreak{} 7II:59\allowbreak{}}\par

\par\medskip\noindent\hypertarget{app-E7II-XI-002.1}{}\textbf{E7II-XI-002.1}\quad printed folio 43.\par
\noindent Printed reading: {\ttfamily\footnotesize .\allowbreak{}}\par
\noindent Corrected reading: {\ttfamily\footnotesize )\allowbreak{}}\par
\noindent Source: {\ttfamily\footnotesize fr/sga7ii/\{diplomati\allowbreak{}c,current\}/expose\_XI\allowbreak{}\_body.tex;\allowbreak{} folio\allowbreak{} 43;\allowbreak{} formula\allowbreak{} 111\allowbreak{}}\par

\par\medskip\noindent\hypertarget{app-E7II-XI-003.1}{}\textbf{E7II-XI-003.1}\quad printed folio 43.\par
\noindent Printed reading: {\ttfamily\footnotesize j\allowbreak{}}\par
\noindent Corrected reading: {\ttfamily\footnotesize i\allowbreak{}}\par
\noindent Source: {\ttfamily\footnotesize fr/sga7ii/\{diplomati\allowbreak{}c,current\}/expose\_XI\allowbreak{}\_body.tex;\allowbreak{} folio\allowbreak{} 43;\allowbreak{} formula\allowbreak{} 126;\allowbreak{} port/decisions/decis\allowbreak{}ions\_en7II\_r2\_001.js\allowbreak{}on\#C514;\allowbreak{} fixes/FIXLOG\_7II\_CUR\allowbreak{}RENT.jsonl\#C514;\allowbreak{} 7II:62\allowbreak{}}\par

\par\medskip\noindent\hypertarget{app-E7II-XI-004.1}{}\textbf{E7II-XI-004.1}\quad printed folio 43.\par
\noindent Printed reading: {\ttfamily\footnotesize a\allowbreak{}}\par
\noindent Corrected reading: {\ttfamily\footnotesize 0\allowbreak{}}\par
\noindent Source: {\ttfamily\footnotesize fr/sga7ii/\{diplomati\allowbreak{}c,current\}/expose\_XI\allowbreak{}\_body.tex;\allowbreak{} folio\allowbreak{} 43;\allowbreak{} formula\allowbreak{} 135;\allowbreak{} port/decisions/decis\allowbreak{}ions\_en7II\_r2\_001.js\allowbreak{}on\#C515;\allowbreak{} fixes/FIXLOG\_7II\_CUR\allowbreak{}RENT.jsonl\#C515;\allowbreak{} 7II:63\allowbreak{}}\par

\par\medskip\noindent\hypertarget{app-E7II-XI-005.1}{}\textbf{E7II-XI-005.1}\quad printed folio 44.\par
\noindent Printed reading: {\ttfamily\footnotesize S\allowbreak{}}\par
\noindent Corrected reading: {\ttfamily\footnotesize s\allowbreak{}}\par
\noindent Source: {\ttfamily\footnotesize fr/sga7ii/\{diplomati\allowbreak{}c,current\}/expose\_XI\allowbreak{}\_body.tex;\allowbreak{} folio\allowbreak{} 44;\allowbreak{} formula\allowbreak{} 153;\allowbreak{} port/decisions/decis\allowbreak{}ions\_en7II\_r2\_001.js\allowbreak{}on\#C516;\allowbreak{} fixes/FIXLOG\_7II\_CUR\allowbreak{}RENT.jsonl\#C516;\allowbreak{} 7II:64\allowbreak{}}\par

\par\medskip\noindent\hypertarget{app-E7II-XI-006.1}{}\textbf{E7II-XI-006.1}\quad printed folio 44.\par
\noindent Printed reading: {\ttfamily\footnotesize H\allowbreak{}}\par
\noindent Corrected reading: {\ttfamily\footnotesize X\allowbreak{}}\par
\noindent Source: {\ttfamily\footnotesize fr/sga7ii/\{diplomati\allowbreak{}c,current\}/expose\_XI\allowbreak{}\_body.tex;\allowbreak{} folio\allowbreak{} 44\allowbreak{}}\par

\par\medskip\noindent\hypertarget{app-E7II-XI-006.2}{}\textbf{E7II-XI-006.2}\quad printed folio 44.\par
\noindent Printed reading: {\ttfamily\footnotesize X\allowbreak{}}\par
\noindent Corrected reading: {\ttfamily\footnotesize H\allowbreak{}}\par
\noindent Source: {\ttfamily\footnotesize fr/sga7ii/\{diplomati\allowbreak{}c,current\}/expose\_XI\allowbreak{}\_body.tex;\allowbreak{} folio\allowbreak{} 44\allowbreak{}}\par

\par\medskip\noindent\hypertarget{app-E7II-XI-006.3}{}\textbf{E7II-XI-006.3}\quad printed folio 44.\par
\noindent Printed reading: {\ttfamily\footnotesize H\allowbreak{}}\par
\noindent Corrected reading: {\ttfamily\footnotesize X\allowbreak{}}\par
\noindent Source: {\ttfamily\footnotesize ED-EN-FORMULA-G095;\allowbreak{} fr/sga7ii/\{diplomati\allowbreak{}c,current\}/expose\_XI\allowbreak{}\_body.tex;\allowbreak{} folio\allowbreak{} 44\allowbreak{}}\par

\par\medskip\noindent\hypertarget{app-E7II-XI-006.4}{}\textbf{E7II-XI-006.4}\quad printed folio 44.\par
\noindent Printed reading: {\ttfamily\footnotesize X\allowbreak{}}\par
\noindent Corrected reading: {\ttfamily\footnotesize H\allowbreak{}}\par
\noindent Source: {\ttfamily\footnotesize ED-EN-FORMULA-G095;\allowbreak{} fr/sga7ii/\{diplomati\allowbreak{}c,current\}/expose\_XI\allowbreak{}\_body.tex;\allowbreak{} folio\allowbreak{} 44\allowbreak{}}\par

\par\medskip\noindent\hypertarget{app-E7II-XI-007.1}{}\textbf{E7II-XI-007.1}\quad printed folio 44.\par
\noindent Printed reading: {\ttfamily\footnotesize Y\allowbreak{}}\par
\noindent Corrected reading: {\ttfamily\footnotesize H\allowbreak{}}\par
\noindent Source: {\ttfamily\footnotesize fr/sga7ii/\{diplomati\allowbreak{}c,current\}/expose\_XI\allowbreak{}\_body.tex;\allowbreak{} folio\allowbreak{} 44;\allowbreak{} formula\allowbreak{} 167;\allowbreak{} port/decisions/decis\allowbreak{}ions\_en7II\_r2\_001.js\allowbreak{}on\#C518;\allowbreak{} fixes/FIXLOG\_7II\_CUR\allowbreak{}RENT.jsonl\#C518;\allowbreak{} 7II:68\allowbreak{}}\par

\par\medskip\noindent\hypertarget{app-E7II-XI-007.2}{}\textbf{E7II-XI-007.2}\quad printed folio 44.\par
\noindent Printed reading: {\ttfamily\footnotesize Y\allowbreak{}}\par
\noindent Corrected reading: {\ttfamily\footnotesize H\allowbreak{}}\par
\noindent Source: {\ttfamily\footnotesize fr/sga7ii/\{diplomati\allowbreak{}c,current\}/expose\_XI\allowbreak{}\_body.tex;\allowbreak{} folio\allowbreak{} 44;\allowbreak{} formula\allowbreak{} 168;\allowbreak{} port/decisions/decis\allowbreak{}ions\_en7II\_r2\_001.js\allowbreak{}on\#C518;\allowbreak{} fixes/FIXLOG\_7II\_CUR\allowbreak{}RENT.jsonl\#C518;\allowbreak{} 7II:68\allowbreak{}}\par

\par\medskip\noindent\hypertarget{app-E7II-XI-007.3}{}\textbf{E7II-XI-007.3}\quad printed folio 44.\par
\noindent Printed reading: {\ttfamily\footnotesize Y\allowbreak{}}\par
\noindent Corrected reading: {\ttfamily\footnotesize H\allowbreak{}}\par
\noindent Source: {\ttfamily\footnotesize fr/sga7ii/\{diplomati\allowbreak{}c,current\}/expose\_XI\allowbreak{}\_body.tex;\allowbreak{} folio\allowbreak{} 44;\allowbreak{} formula\allowbreak{} 169;\allowbreak{} port/decisions/decis\allowbreak{}ions\_en7II\_r2\_001.js\allowbreak{}on\#C518;\allowbreak{} fixes/FIXLOG\_7II\_CUR\allowbreak{}RENT.jsonl\#C518;\allowbreak{} 7II:68\allowbreak{}}\par

\par\medskip\noindent\hypertarget{app-E7II-XI-008.1}{}\textbf{E7II-XI-008.1}\quad printed folio 44.\par
\noindent Printed reading: {\ttfamily\footnotesize Y\allowbreak{}}\par
\noindent Corrected reading: {\ttfamily\footnotesize H\allowbreak{}}\par
\noindent Source: {\ttfamily\footnotesize fr/sga7ii/\{diplomati\allowbreak{}c,current\}/expose\_XI\allowbreak{}\_body.tex;\allowbreak{} folio\allowbreak{} 44;\allowbreak{} formula\allowbreak{} 170;\allowbreak{} port/decisions/decis\allowbreak{}ions\_en7II\_r2\_001.js\allowbreak{}on\#C518;\allowbreak{} fixes/FIXLOG\_7II\_CUR\allowbreak{}RENT.jsonl\#C518;\allowbreak{} 7II:68\allowbreak{}}\par

\par\medskip\noindent\hypertarget{app-E7II-XI-009.1}{}\textbf{E7II-XI-009.1}\quad printed folio 45.\par
\noindent Printed reading: {\ttfamily\footnotesize Y\allowbreak{}}\par
\noindent Corrected reading: {\ttfamily\footnotesize H\allowbreak{}}\par
\noindent Source: {\ttfamily\footnotesize port/decisions/decis\allowbreak{}ions\_en7II\_r2\_001.js\allowbreak{}on\#C519;\allowbreak{} 7II:69;\allowbreak{} fr/sga7ii/\{diplomati\allowbreak{}c,current\}/expose\_XI\allowbreak{}\_body.tex;\allowbreak{} folio\allowbreak{} 45;\allowbreak{} formula\allowbreak{} 177\allowbreak{}}\par

\par\medskip\noindent\hypertarget{app-E7II-XI-009.2}{}\textbf{E7II-XI-009.2}\quad printed folio 45.\par
\noindent Printed reading: {\ttfamily\footnotesize Y\allowbreak{}}\par
\noindent Corrected reading: {\ttfamily\footnotesize H\allowbreak{}}\par
\noindent Source: {\ttfamily\footnotesize port/decisions/decis\allowbreak{}ions\_en7II\_r2\_001.js\allowbreak{}on\#C519;\allowbreak{} 7II:69;\allowbreak{} fr/sga7ii/\{diplomati\allowbreak{}c,current\}/expose\_XI\allowbreak{}\_body.tex;\allowbreak{} folio\allowbreak{} 45;\allowbreak{} formula\allowbreak{} 178\allowbreak{}}\par

\par\medskip\noindent\hypertarget{app-E7II-XI-010.1}{}\textbf{E7II-XI-010.1}\quad printed folio 45.\par
\noindent Printed reading: {\ttfamily\footnotesize Y\allowbreak{}}\par
\noindent Corrected reading: {\ttfamily\footnotesize H\allowbreak{}}\par
\noindent Source: {\ttfamily\footnotesize port/decisions/decis\allowbreak{}ions\_en7II\_r2\_001.js\allowbreak{}on\#C519;\allowbreak{} 7II:69\allowbreak{}}\par

\par\medskip\noindent\hypertarget{app-E7II-XI-010.2}{}\textbf{E7II-XI-010.2}\quad printed folio 45.\par
\noindent Printed reading: {\ttfamily\footnotesize Y\allowbreak{}}\par
\noindent Corrected reading: {\ttfamily\footnotesize H\allowbreak{}}\par
\noindent Source: {\ttfamily\footnotesize port/decisions/decis\allowbreak{}ions\_en7II\_r2\_001.js\allowbreak{}on\#C519;\allowbreak{} 7II:69\allowbreak{}}\par

\par\medskip\noindent\hypertarget{app-E7II-XI-010.3}{}\textbf{E7II-XI-010.3}\quad printed folio 45.\par
\noindent Printed reading: {\ttfamily\footnotesize Y\allowbreak{}}\par
\noindent Corrected reading: {\ttfamily\footnotesize H\allowbreak{}}\par
\noindent Source: {\ttfamily\footnotesize port/decisions/decis\allowbreak{}ions\_en7II\_r2\_001.js\allowbreak{}on\#C519;\allowbreak{} 7II:69;\allowbreak{} fr/sga7ii/\{diplomati\allowbreak{}c,current\}/expose\_XI\allowbreak{}\_body.tex;\allowbreak{} folio\allowbreak{} 45;\allowbreak{} formula\allowbreak{} 184\allowbreak{}}\par

\par\medskip\noindent\hypertarget{app-E7II-XI-010.4}{}\textbf{E7II-XI-010.4}\quad printed folio 45.\par
\noindent Printed reading: {\ttfamily\footnotesize Y\allowbreak{}}\par
\noindent Corrected reading: {\ttfamily\footnotesize H\allowbreak{}}\par
\noindent Source: {\ttfamily\footnotesize port/decisions/decis\allowbreak{}ions\_en7II\_r2\_001.js\allowbreak{}on\#C519;\allowbreak{} 7II:69;\allowbreak{} fr/sga7ii/\{diplomati\allowbreak{}c,current\}/expose\_XI\allowbreak{}\_body.tex;\allowbreak{} folio\allowbreak{} 45;\allowbreak{} formula\allowbreak{} 185\allowbreak{}}\par

\par\medskip\noindent\hypertarget{app-E7II-XI-011.1}{}\textbf{E7II-XI-011.1}\quad printed folio 46.\par
\noindent Printed reading: {\ttfamily\footnotesize X\allowbreak{}}\par
\noindent Corrected reading: {\ttfamily\footnotesize Y\allowbreak{}}\par
\noindent Source: {\ttfamily\footnotesize ED-EN-FORMULA-G097;\allowbreak{} fr/sga7ii/\{diplomati\allowbreak{}c,current\}/expose\_XI\allowbreak{}\_body.tex;\allowbreak{} folio\allowbreak{} 46;\allowbreak{} port/decisions/decis\allowbreak{}ions\_en7II\_r2\_001.js\allowbreak{}on\#C521;\allowbreak{} fixes/FIXLOG\_7II\_CUR\allowbreak{}RENT.jsonl\#C521;\allowbreak{} 7II:71\allowbreak{}}\par

\par\medskip\noindent\hypertarget{app-E7II-XI-012.1}{}\textbf{E7II-XI-012.1}\quad printed folio 47.\par
\noindent Printed reading: \emph{(empty)}\par
\noindent Corrected reading: {\ttfamily\footnotesize +j\allowbreak{}}\par
\noindent Source: {\ttfamily\footnotesize fr/sga7ii/\{diplomati\allowbreak{}c,current\}/expose\_XI\allowbreak{}\_body.tex;\allowbreak{} folio\allowbreak{} 47;\allowbreak{} formula\allowbreak{} 248;\allowbreak{} port/PORT\_EN7II\_LOG.\allowbreak{}jsonl\#C523;\allowbreak{} fixes/FIXLOG\_7II\_CUR\allowbreak{}RENT.jsonl\#C523;\allowbreak{} 7II:73\allowbreak{}}\par

\par\medskip\noindent\hypertarget{app-E7II-XI-012.2}{}\textbf{E7II-XI-012.2}\quad printed folio 47.\par
\noindent Printed reading: {\ttfamily\footnotesize j\allowbreak{}}\par
\noindent Corrected reading: {\ttfamily\footnotesize n\allowbreak{}}\par
\noindent Source: {\ttfamily\footnotesize fr/sga7ii/\{diplomati\allowbreak{}c,current\}/expose\_XI\allowbreak{}\_body.tex;\allowbreak{} folio\allowbreak{} 47;\allowbreak{} formula\allowbreak{} 248;\allowbreak{} port/PORT\_EN7II\_LOG.\allowbreak{}jsonl\#C523;\allowbreak{} fixes/FIXLOG\_7II\_CUR\allowbreak{}RENT.jsonl\#C523;\allowbreak{} 7II:73\allowbreak{}}\par

\par\medskip\noindent\hypertarget{app-E7II-XI-013.1}{}\textbf{E7II-XI-013.1}\quad printed folio 48.\par
\noindent Printed reading: \emph{(empty)}\par
\noindent Corrected reading: {\ttfamily\footnotesize \textbackslash{}bigl\allowbreak{}}\par
\noindent Source: {\ttfamily\footnotesize ED-EN-FONT-G016;\allowbreak{} F7II-XI-013.1;\allowbreak{} fr/sga7ii/\{diplomati\allowbreak{}c,current\}/expose\_XI\allowbreak{}\_body.tex;\allowbreak{} folio\allowbreak{} 48;\allowbreak{} port/PORT\_EN7II\_LOG.\allowbreak{}jsonl\#C525;\allowbreak{} fixes/FIXLOG\_7II\_CUR\allowbreak{}RENT.jsonl\#C525;\allowbreak{} 7II:75\allowbreak{}}\par

\par\medskip\noindent\hypertarget{app-E7II-XI-013.2}{}\textbf{E7II-XI-013.2}\quad printed folio 48.\par
\noindent Printed reading: {\ttfamily\footnotesize Y\allowbreak{}}\par
\noindent Corrected reading: {\ttfamily\footnotesize \textbackslash{}mathbb\allowbreak{} P(\textbackslash{}underline\allowbreak{} E)\allowbreak{}}\par
\noindent Source: {\ttfamily\footnotesize ED-EN-FONT-G016;\allowbreak{} F7II-XI-013.1;\allowbreak{} fr/sga7ii/\{diplomati\allowbreak{}c,current\}/expose\_XI\allowbreak{}\_body.tex;\allowbreak{} folio\allowbreak{} 48;\allowbreak{} fr/sga7ii/\{diplomati\allowbreak{}c,current\}/expose\_XI\allowbreak{}\_body.tex;\allowbreak{} folio\allowbreak{} 48;\allowbreak{} formula\allowbreak{} 288;\allowbreak{} port/PORT\_EN7II\_LOG.\allowbreak{}jsonl\#C525;\allowbreak{} fixes/FIXLOG\_7II\_CUR\allowbreak{}RENT.jsonl\#C525;\allowbreak{} 7II:75\allowbreak{}}\par

\par\medskip\noindent\hypertarget{app-E7II-XI-013.3}{}\textbf{E7II-XI-013.3}\quad printed folio 48.\par
\noindent Printed reading: \emph{(empty)}\par
\noindent Corrected reading: {\ttfamily\footnotesize \textbackslash{}bigr\allowbreak{}}\par
\noindent Source: {\ttfamily\footnotesize ED-EN-FONT-G016;\allowbreak{} F7II-XI-013.1;\allowbreak{} fr/sga7ii/\{diplomati\allowbreak{}c,current\}/expose\_XI\allowbreak{}\_body.tex;\allowbreak{} folio\allowbreak{} 48;\allowbreak{} port/PORT\_EN7II\_LOG.\allowbreak{}jsonl\#C525;\allowbreak{} fixes/FIXLOG\_7II\_CUR\allowbreak{}RENT.jsonl\#C525;\allowbreak{} 7II:75\allowbreak{}}\par

\par\medskip\noindent\hypertarget{app-E7II-XI-014.1}{}\textbf{E7II-XI-014.1}\quad printed folio 52.\par
\noindent Printed reading: {\ttfamily\footnotesize i\allowbreak{}}\par
\noindent Corrected reading: {\ttfamily\footnotesize 1\allowbreak{}}\par
\noindent Source: {\ttfamily\footnotesize fr/sga7ii/\{diplomati\allowbreak{}c,current\}/expose\_XI\allowbreak{}\_body.tex;\allowbreak{} folio\allowbreak{} 52;\allowbreak{} formula\allowbreak{} 378;\allowbreak{} port/decisions/decis\allowbreak{}ions\_en7II\_r2\_001.js\allowbreak{}on\#C526;\allowbreak{} fixes/FIXLOG\_7II\_CUR\allowbreak{}RENT.jsonl\#C526;\allowbreak{} 7II:77\allowbreak{}}\par

\par\medskip\noindent\hypertarget{app-E7II-XI-015.1}{}\textbf{E7II-XI-015.1}\quad printed folio 52.\par
\noindent Printed reading: {\ttfamily\footnotesize .\allowbreak{}}\par
\noindent Corrected reading: \emph{(empty)}\par
\noindent Source: {\ttfamily\footnotesize fr/sga7ii/\{diplomati\allowbreak{}c,current\}/expose\_XI\allowbreak{}\_body.tex;\allowbreak{} folio\allowbreak{} 52;\allowbreak{} formula\allowbreak{} 382\allowbreak{}}\par

\par\medskip\noindent\hypertarget{app-E7II-XI-015.2}{}\textbf{E7II-XI-015.2}\quad printed folio 52.\par
\noindent Printed reading: {\ttfamily\footnotesize .\allowbreak{}}\par
\noindent Corrected reading: \emph{(empty)}\par
\noindent Source: {\ttfamily\footnotesize fr/sga7ii/\{diplomati\allowbreak{}c,current\}/expose\_XI\allowbreak{}\_body.tex;\allowbreak{} folio\allowbreak{} 52;\allowbreak{} formula\allowbreak{} 382\allowbreak{}}\par

\par\medskip\noindent\hypertarget{app-E7II-XI-016.1}{}\textbf{E7II-XI-016.1}\quad printed folio 54.\par
\noindent Printed reading: {\ttfamily\footnotesize i\allowbreak{}}\par
\noindent Corrected reading: {\ttfamily\footnotesize 1\allowbreak{}}\par
\noindent Source: {\ttfamily\footnotesize port/decisions/decis\allowbreak{}ions\_en7II\_r2\_001.js\allowbreak{}on\#C533;\allowbreak{} fixes/FIXLOG\_7II\_CUR\allowbreak{}RENT.jsonl\#C533;\allowbreak{} 7II:88\allowbreak{}}\par

\par\medskip\noindent\hypertarget{app-E7II-XI-017.1}{}\textbf{E7II-XI-017.1}\quad printed folio 53.\par
\noindent Printed reading: {\ttfamily\footnotesize z\allowbreak{}}\par
\noindent Corrected reading: {\ttfamily\footnotesize y\allowbreak{}}\par
\noindent Source: {\ttfamily\footnotesize port/decisions/decis\allowbreak{}ions\_en7II\_r2\_001.js\allowbreak{}on\#C528;\allowbreak{} fixes/FIXLOG\_7II\_CUR\allowbreak{}RENT.jsonl\#C528;\allowbreak{} 7II:83\allowbreak{}}\par

\par\medskip\noindent\hypertarget{app-E7II-XI-017.2}{}\textbf{E7II-XI-017.2}\quad printed folio 53.\par
\noindent Printed reading: {\ttfamily\footnotesize z\allowbreak{}}\par
\noindent Corrected reading: {\ttfamily\footnotesize y\allowbreak{}}\par
\noindent Source: {\ttfamily\footnotesize fr/sga7ii/\{diplomati\allowbreak{}c,current\}/expose\_XI\allowbreak{}\_body.tex;\allowbreak{} folio\allowbreak{} 53;\allowbreak{} formula\allowbreak{} 394\allowbreak{}}\par

\par\medskip\noindent\hypertarget{app-E7II-XI-018.1}{}\textbf{E7II-XI-018.1}\quad printed folio 54.\par
\noindent Printed reading: {\ttfamily\footnotesize .\allowbreak{}}\par
\noindent Corrected reading: {\ttfamily\footnotesize )\allowbreak{}}\par
\noindent Source: {\ttfamily\footnotesize fr/sga7ii/\{diplomati\allowbreak{}c,current\}/expose\_XI\allowbreak{}\_body.tex;\allowbreak{} folio\allowbreak{} 54;\allowbreak{} formula\allowbreak{} 399;\allowbreak{} port/decisions/decis\allowbreak{}ions\_en7II\_r2\_001.js\allowbreak{}on\#C529;\allowbreak{} fixes/FIXLOG\_7II\_CUR\allowbreak{}RENT.jsonl\#C529;\allowbreak{} 7II:84\allowbreak{}}\par

\par\medskip\noindent\hypertarget{app-E7II-XI-019.1}{}\textbf{E7II-XI-019.1}\quad printed folio 54.\par
\noindent Printed reading: {\ttfamily\footnotesize x\allowbreak{}}\par
\noindent Corrected reading: {\ttfamily\footnotesize y\allowbreak{}}\par
\noindent Source: {\ttfamily\footnotesize fr/sga7ii/\{diplomati\allowbreak{}c,current\}/expose\_XI\allowbreak{}\_body.tex;\allowbreak{} folio\allowbreak{} 54;\allowbreak{} formula\allowbreak{} 401;\allowbreak{} port/decisions/decis\allowbreak{}ions\_en7II\_r2\_001.js\allowbreak{}on\#C530;\allowbreak{} fixes/FIXLOG\_7II\_CUR\allowbreak{}RENT.jsonl\#C530;\allowbreak{} 7II:85\allowbreak{}}\par

\par\medskip\noindent\hypertarget{app-E7II-XI-019.2}{}\textbf{E7II-XI-019.2}\quad printed folio 54.\par
\noindent Printed reading: {\ttfamily\footnotesize y\allowbreak{}}\par
\noindent Corrected reading: {\ttfamily\footnotesize z\allowbreak{}}\par
\noindent Source: {\ttfamily\footnotesize fr/sga7ii/\{diplomati\allowbreak{}c,current\}/expose\_XI\allowbreak{}\_body.tex;\allowbreak{} folio\allowbreak{} 54;\allowbreak{} formula\allowbreak{} 401;\allowbreak{} port/decisions/decis\allowbreak{}ions\_en7II\_r2\_001.js\allowbreak{}on\#C530;\allowbreak{} fixes/FIXLOG\_7II\_CUR\allowbreak{}RENT.jsonl\#C530;\allowbreak{} 7II:85\allowbreak{}}\par

\par\medskip\noindent\hypertarget{app-E7II-XI-020.1}{}\textbf{E7II-XI-020.1}\quad printed folio 54.\par
\noindent Printed reading: {\ttfamily\footnotesize x\allowbreak{}}\par
\noindent Corrected reading: {\ttfamily\footnotesize y\allowbreak{}}\par
\noindent Source: {\ttfamily\footnotesize fr/sga7ii/\{diplomati\allowbreak{}c,current\}/expose\_XI\allowbreak{}\_body.tex;\allowbreak{} folio\allowbreak{} 54;\allowbreak{} formula\allowbreak{} 402;\allowbreak{} port/decisions/decis\allowbreak{}ions\_en7II\_r2\_001.js\allowbreak{}on\#C531;\allowbreak{} fixes/FIXLOG\_7II\_CUR\allowbreak{}RENT.jsonl\#C531;\allowbreak{} 7II:86\allowbreak{}}\par

\par\medskip\noindent\hypertarget{app-E7II-XI-020.2}{}\textbf{E7II-XI-020.2}\quad printed folio 54.\par
\noindent Printed reading: {\ttfamily\footnotesize y\allowbreak{}}\par
\noindent Corrected reading: {\ttfamily\footnotesize z\allowbreak{}}\par
\noindent Source: {\ttfamily\footnotesize fr/sga7ii/\{diplomati\allowbreak{}c,current\}/expose\_XI\allowbreak{}\_body.tex;\allowbreak{} folio\allowbreak{} 54;\allowbreak{} formula\allowbreak{} 402;\allowbreak{} port/decisions/decis\allowbreak{}ions\_en7II\_r2\_001.js\allowbreak{}on\#C531;\allowbreak{} fixes/FIXLOG\_7II\_CUR\allowbreak{}RENT.jsonl\#C531;\allowbreak{} 7II:86\allowbreak{}}\par

\par\medskip\noindent\hypertarget{app-E7II-XI-020.3}{}\textbf{E7II-XI-020.3}\quad printed folio 54.\par
\noindent Printed reading: {\ttfamily\footnotesize i\allowbreak{}}\par
\noindent Corrected reading: {\ttfamily\footnotesize 1\allowbreak{}}\par
\noindent Source: {\ttfamily\footnotesize fr/sga7ii/\{diplomati\allowbreak{}c,current\}/expose\_XI\allowbreak{}\_body.tex;\allowbreak{} folio\allowbreak{} 54;\allowbreak{} formula\allowbreak{} 402\allowbreak{}}\par

\par\medskip\noindent\hypertarget{app-E7II-XI-020.4}{}\textbf{E7II-XI-020.4}\quad printed folio 54.\par
\noindent Printed reading: {\ttfamily\footnotesize x\allowbreak{}}\par
\noindent Corrected reading: {\ttfamily\footnotesize y\allowbreak{}}\par
\noindent Source: {\ttfamily\footnotesize fr/sga7ii/\{diplomati\allowbreak{}c,current\}/expose\_XI\allowbreak{}\_body.tex;\allowbreak{} folio\allowbreak{} 54;\allowbreak{} formula\allowbreak{} 402\allowbreak{}}\par

\par\medskip\noindent\hypertarget{app-E7II-XI-020.5}{}\textbf{E7II-XI-020.5}\quad printed folio 54.\par
\noindent Printed reading: {\ttfamily\footnotesize y\allowbreak{}}\par
\noindent Corrected reading: {\ttfamily\footnotesize z\allowbreak{}}\par
\noindent Source: {\ttfamily\footnotesize fr/sga7ii/\{diplomati\allowbreak{}c,current\}/expose\_XI\allowbreak{}\_body.tex;\allowbreak{} folio\allowbreak{} 54;\allowbreak{} formula\allowbreak{} 402\allowbreak{}}\par

\par\medskip\noindent\hypertarget{app-E7II-XI-021.1}{}\textbf{E7II-XI-021.1}\quad printed folio 56.\par
\noindent Printed reading: {\ttfamily\footnotesize \textbackslash{}leqslant\allowbreak{}}\par
\noindent Corrected reading: {\ttfamily\footnotesize \textbackslash{}in\allowbreak{}}\par
\noindent Source: {\ttfamily\footnotesize fr/sga7ii/\{diplomati\allowbreak{}c,current\}/expose\_XI\allowbreak{}\_body.tex;\allowbreak{} folio\allowbreak{} 56;\allowbreak{} formula\allowbreak{} 446\allowbreak{}}\par

\par\medskip\noindent\hypertarget{app-E7II-XI-021.2}{}\textbf{E7II-XI-021.2}\quad printed folio 56.\par
\noindent Printed reading: {\ttfamily\footnotesize i\allowbreak{} \textbackslash{}leqslant\allowbreak{} d\allowbreak{}}\par
\noindent Corrected reading: {\ttfamily\footnotesize P\allowbreak{}}\par
\noindent Source: {\ttfamily\footnotesize fr/sga7ii/\{diplomati\allowbreak{}c,current\}/expose\_XI\allowbreak{}\_body.tex;\allowbreak{} folio\allowbreak{} 56;\allowbreak{} formula\allowbreak{} 446\allowbreak{}}\par

\par\medskip\noindent\hypertarget{app-E7II-XI-021.3}{}\textbf{E7II-XI-021.3}\quad printed folio 56.\par
\noindent Printed reading: \emph{(empty)}\par
\noindent Corrected reading: {\ttfamily\footnotesize \_\{i\textbackslash{}in\allowbreak{} P\}\allowbreak{}}\par
\noindent Source: {\ttfamily\footnotesize ED-EN-FORMULA-G099;\allowbreak{} fr/sga7ii/\{diplomati\allowbreak{}c,current\}/expose\_XI\allowbreak{}\_body.tex;\allowbreak{} folio\allowbreak{} 56\allowbreak{}}\par

\par\medskip\noindent\hypertarget{app-E7II-XI-022.1}{}\textbf{E7II-XI-022.1}\quad printed folio 57.\par
\noindent Printed reading: \emph{(empty)}\par
\noindent Corrected reading: {\ttfamily\footnotesize \_\{i\textbackslash{}in\allowbreak{}}\par
\noindent Source: {\ttfamily\footnotesize fr/sga7ii/\{diplomati\allowbreak{}c,current\}/expose\_XI\allowbreak{}\_body.tex;\allowbreak{} folio\allowbreak{} 57\allowbreak{}}\par

\par\medskip\noindent\hypertarget{app-E7II-XI-022.2}{}\textbf{E7II-XI-022.2}\quad printed folio 57.\par
\noindent Printed reading: \emph{(empty)}\par
\noindent Corrected reading: {\ttfamily\footnotesize P\}\allowbreak{}}\par
\noindent Source: {\ttfamily\footnotesize fr/sga7ii/\{diplomati\allowbreak{}c,current\}/expose\_XI\allowbreak{}\_body.tex;\allowbreak{} folio\allowbreak{} 57\allowbreak{}}\par

\par\medskip\noindent\hypertarget{app-E7II-XI-022.3}{}\textbf{E7II-XI-022.3}\quad printed folio 57.\par
\noindent Printed reading: \emph{(empty)}\par
\noindent Corrected reading: {\ttfamily\footnotesize \_\{i\textbackslash{}in\allowbreak{}}\par
\noindent Source: {\ttfamily\footnotesize fr/sga7ii/\{diplomati\allowbreak{}c,current\}/expose\_XI\allowbreak{}\_body.tex;\allowbreak{} folio\allowbreak{} 57\allowbreak{}}\par

\par\medskip\noindent\hypertarget{app-E7II-XI-022.4}{}\textbf{E7II-XI-022.4}\quad printed folio 57.\par
\noindent Printed reading: \emph{(empty)}\par
\noindent Corrected reading: {\ttfamily\footnotesize P\}\allowbreak{}}\par
\noindent Source: {\ttfamily\footnotesize fr/sga7ii/\{diplomati\allowbreak{}c,current\}/expose\_XI\allowbreak{}\_body.tex;\allowbreak{} folio\allowbreak{} 57\allowbreak{}}\par

\par\medskip\noindent\hypertarget{app-E7II-XI-023.1}{}\textbf{E7II-XI-023.1}\quad printed folio 58.\par
\noindent Printed reading: \emph{(empty)}\par
\noindent Corrected reading: {\ttfamily\footnotesize 2\allowbreak{}}\par
\noindent Source: {\ttfamily\footnotesize fr/sga7ii/\{diplomati\allowbreak{}c,current\}/expose\_XI\allowbreak{}\_body.tex;\allowbreak{} folio\allowbreak{} 58;\allowbreak{} formula\allowbreak{} 504\allowbreak{}}\par

\par\endgroup\normalsize\fi
\clearpage\SGAtoc{Exposé XII}
\selectlanguage{english}
\clearpage
\SGApage{II}{62}

\begin{center}
S G A 7

\bigskip
\underline{E X P O S É XII}

\bigskip
\underline{\underline{QUADRICS}}

\bigskip
\underline{by P. DELIGNE}
\end{center}

\bigskip
0.0. The present exposé collects several well-known results about
quadrics which we shall need in order to prove the Picard--Lefschetz
formula.

0.1. Let $m$ be a section of an $\mathcal O_S$-module $M$ on a scheme
$S$.  We shall say that $m$ is \underline{nowhere zero} on $S$ if, for
every $s\in S$, the image of $m$ in the $k(s)$-vector space $M_s$ obtained
from $M$ by the base change $\{s\}\longrightarrow S$ is nonzero.  We use
the same terminology when $S=\Spec(A)$ and $m$ is an element of an
$A$-module $M$.

\medskip
1. \underline{Quadratic forms}

1.1. Let $S=\Spec(A)$ be an affine scheme, let $V$ be a locally free
$A$-module of rank $n$, and let $Q$ be a quadratic form on $V$
(Bourbaki, \emph{Algèbre}, Ch.~9, \S3, no.~4).  Recall that $Q$ is called
\underline{nondegenerate} if the associated bilinear form

\begin{equation*}
\Phi(x,y)=Q(x+y)-Q(x)-Q(y)
\end{equation*}

is nondegenerate, that is, if it induces an isomorphism from $V$ to its
dual $V^*$.  If $2$ is not invertible in $A$, this can occur only when
$n$ is even (loc. cit.).


\SGApage{II}{\(\approx\)63}
The quadratic forms on $V$ form a projective $A$-module; a quadratic form
$Q$ is ``nowhere zero'' on $S$ (0.1) if and only if the elements $Q(v)$
generate the unit ideal of $A$.

The notation $\mathbb P(V)$ is used in the sense of EGA; recall that if
$A$ is a field, then $\mathbb P(V)$ is the projective space whose set of
rational points is the quotient $(V^*\setminus\{0\})/A^*$.

The quadratic form $Q$ identifies with a section of the invertible sheaf
$\mathcal O(2)$ on the projective space $\mathbb P(V^*)$ over $S$.  When
$Q$ is nowhere zero, the (``homogeneous'') equation $Q=0$ defines a
subscheme of $\mathbb P(V^*)$, flat and purely of relative dimension
$n-2$ over $S$.  This scheme is called the \underline{quadric defined by
$Q$}.  The form $Q$ is called \underline{ordinary} if it is nowhere zero
on $S$ and the quadric that it defines is smooth over $S$.

The form $Q$ is ordinary if and only if every form obtained from it by
extending the scalars from $A$ to a field is ordinary.  If $A$ is a
field, then:

\SGAvar{\begin{enumerate}[label=\alph*),leftmargin=2.5em]
\item if $n$ is even or $\operatorname{char}(A)\ne2$, then $Q$ is ordinary
if and only if $Q$ is nondegenerate;

\item if $n$ is odd and $\operatorname{card}(A)=2$, then $Q$ is ordinary
if and only if the kernel $N$ of the associated (alternating) bilinear
form $\Phi$ has dimension one and $Q$ does not vanish on $N$.
\end{enumerate}

}{\begin{enumerate}[label=\alph*),leftmargin=2.5em]
\item if $n$ is even or $\operatorname{char}(A)\ne2$, then $Q$ is ordinary
if and only if $Q$ is nondegenerate;

\item if $n$ is odd and $\operatorname{char}(A)=2$, then $Q$ is ordinary
if and only if the kernel $N$ of the associated (alternating) bilinear
form $\Phi$ has dimension one and $Q$ does not vanish on $N$.
\end{enumerate}

}{E7II-XII-001}\underline{Proposition} 1.2. \underline{If $Q$ is ordinary and
$n=2m$ (respectively, $n=2m+1$), then, locally for the étale topology on
$S$, $V$ has a basis $e_1,\ldots,e_n$ such that}


\SGAvar{\begin{equation*}
Q\!\left(\sum_{i=1}^{n}x_i e_i\right)
=\sum_{i=1}^{m-1}x_i x_{i+m}
\end{equation*}

}{\begin{equation*}
Q\!\left(\sum_{i=1}^{n}x_i e_i\right)
=\sum_{i=1}^{m}x_i x_{i+m}
\end{equation*}

}{E7II-XII-002}\noindent\underline{(respectively,}

\SGApage{II}{\(\approx\)64}
\SGAvar{\begin{equation*}
Q\!\left(\sum_{i=1}^{n}x_i e_i\right)
=\sum_{i=1}^{m-1}x_i x_{i+m}+\lambda x_{2m+1}^{2},
\qquad (\lambda\text{ invertible}).
\end{equation*}



}{\begin{equation*}
Q\!\left(\sum_{i=1}^{n}x_i e_i\right)
=\sum_{i=1}^{m}x_i x_{i+m}+\lambda x_{2m+1}^{2},
\qquad (\lambda\text{ invertible})).
\end{equation*}



}{E7II-XII-003}We prove 1.2 by induction on $m$.  If $m=0$, the assertion is clear.  If
$m>0$, the quadric with equation $Q=0$ is smooth over $S$ and has
nonempty geometric fibers; it therefore has sections locally for the
étale topology.  Locally, such a section is defined by an element
$e\in V$, nowhere zero on $\Spec(A)$, with $Q(e)=0$.  At no point of $S$
does $e$ belong to the kernel of $\Phi$; hence, locally, one can find
$f'$ such that $\Phi(e,f')=1$.  If

\SGAvar{\begin{equation*}
f=-Q(f')e+f,
\end{equation*}

}{\begin{equation*}
f=-Q(f')e+f',
\end{equation*}

}{E7II-XII-004}then

\begin{equation*}
Q(e)=Q(f)=0,\qquad \Phi(e,f)=1.
\end{equation*}



If $V_2$ is the orthogonal complement of the submodule $V_1$ generated
by $e$ and $f$, then $V=V_1\oplus V_2$, and the conclusion follows by
applying the induction hypothesis to $V_2$.

1.3. For the definition and principal properties of the Clifford algebra
of $Q$, we refer to Bourbaki, \emph{Algèbre}, Ch.~9, \S9.  We denote by
$\rho$ the canonical map from $V$ to $C(Q)$.

\underline{Lemma} 1.3.1. \underline{Under the hypotheses of 1.1, for
every $\lambda\in A$ there exists a unique homomorphism}

\begin{equation*}
[\lambda]\colon C^+(\lambda Q)\longrightarrow C^+(Q)
\end{equation*}

\underline{such that, for $x,y\in V$,}

\begin{equation*}
[\lambda]\bigl(\rho(x)\rho(y)\bigr)
=\lambda\rho(x)\rho(y).
\end{equation*}


\SGApage{II}{\(\approx\)65}
\SGAvar{This homomorphism is induced on passing to the quotient by the
endomorphism $[\lambda]$ of the even part $T^+(E)$ of the tensor algebra
of $E$, which acts as $\lambda^r$ on $T^{2r}(E)$.  Indeed,

}{This homomorphism is induced on passing to the quotient by the
endomorphism $[\lambda]$ of the even part $T^+(V)$ of the tensor algebra
of $V$, which acts as $\lambda^r$ on $T^{2r}(V)$.  Indeed,

}{E7II-XII-005}\begin{equation*}
[\lambda]\bigl(x\otimes x-\lambda Q(x)\bigr)
=\lambda\bigl(x\otimes x-Q(x)\bigr).
\end{equation*}

\SGAvar{Suppose that $Q$ is nowhere zero on $S$.  Lemma 1.2.1 then defines a
transitive system of isomorphisms between the algebras $C^+(\lambda Q)$
for $\lambda\in A^*$.

}{Suppose that $Q$ is nowhere zero on $S$.  Lemma 1.3.1 then defines a
transitive system of isomorphisms between the algebras $C^+(\lambda Q)$
for $\lambda\in A^*$.

}{E7II-XII-006}Consequently, the algebra $C^+(Q)$ depends (up to unique isomorphism)
only on $V$ and on $Q$ up to multiplication by a factor; in other words,
$C^+(Q)$ depends only on $V$ and on the subscheme of $\mathbb P(V^*)$
with equation $Q=0$.  Every homothety with invertible ratio is an
automorphism of the pair $(V,X\subset\mathbb P(V^*))$; one checks that
these automorphisms act trivially on $C^+(V,X)$.  It follows that the
algebra $C^+(V,X)$ depends only on the pair of schemes
$X\subset\mathbb P(V^*)$ over $A$.  We shall denote it by
$C^+(X,\mathbb P(V^*))$.


1.4. Let $V$ be a locally free $A$-module of rank $2m>0$, and let $Q$
be a nondegenerate quadratic form on $V$.  Suppose given a decomposition
$V=W_1\oplus W_2$ into two totally singular subspaces.  The form
$Q(w_1+w_2)$ then places $W_1$ and $W_2$ in perfect duality.  Recall that,
for $f\in W_2$ (identified with $W_1^*$), contraction by $f$ is the
endomorphism of degree $-1$ of $\bigwedge W_1$ defined by

\SGAvar{\begin{equation*}
f\mathbin{\lrcorner}(e_1\wedge\cdots\wedge e_p)
=\sum_{i=1}^{p}(-1)^i\langle e_i,f\rangle
e_0\wedge\cdots\wedge\widehat e_i\wedge\cdots\wedge e_p.
\end{equation*}



}{\begin{equation*}
f\mathbin{\lrcorner}(e_1\wedge\cdots\wedge e_p)
=\sum_{i=1}^{p}(-1)^{i-1}\langle e_i,f\rangle
e_1\wedge\cdots\wedge\widehat e_i\wedge\cdots\wedge e_p.
\end{equation*}



}{E7II-XII-007}Endow the $A$-module $\bigwedge W_1$ with the $\mathbb Z/2$-grading for
which


\SGApage{II}{\(\approx\)66}
\begin{equation*}
(\bigwedge W_1)^+=\bigoplus_i\bigwedge^{2i}W_1,
\qquad
(\bigwedge W_1)^-=\bigoplus_i\bigwedge^{2i+1}W_1.
\end{equation*}

The algebra $\operatorname{End}(\bigwedge W_1)$ is then
$\mathbb Z/2$-graded: its even (respectively, odd) endomorphisms are those
that preserve the grading (respectively, have degree one modulo two).

We refer to Bourbaki, \emph{Algèbre}, Ch.~9, \S9, no.~4, for the proof of
the existence of a unique isomorphism $s$ of $\mathbb Z/2$-graded algebras
from $C(Q)$ to $\operatorname{End}(\bigwedge W_1)$ such that

\SGAvar{\begin{equation*}
s\bigl(\rho((e,0))=e\wedge(-),
\qquad
s\bigl(\rho((0,f))=f\mathbin{\lrcorner}(-).
\end{equation*}

Temporarily denote by $e(W_1,W_2)$ the element of $C^+(Q)$ for which
$s(e(W_1,W_2))$ is the identity on $(\bigwedge W_1)^+$ and zero on
$(\bigwedge W_1)^-$.  The center $Z(Q)$ of $C^+(Q)$ is isomorphic to
$A\oplus A$, this decomposition being defined by the idempotents
$e(W_1,W_2)$ and $1-e(W_2,W_1)$.  Moreover, $C^+(Q)$ is a matrix algebra
over $Z(Q)$.

}{\begin{equation*}
s\bigl(\rho((e,0))\bigr)=e\wedge(-),
\qquad
s\bigl(\rho((0,f))\bigr)=f\mathbin{\lrcorner}(-).
\end{equation*}

Temporarily denote by $e(W_1,W_2)$ the element of $C^+(Q)$ for which
$s(e(W_1,W_2))$ is the identity on $(\bigwedge W_1)^+$ and zero on
$(\bigwedge W_1)^-$.  The center $Z(Q)$ of $C^+(Q)$ is isomorphic to
$A\oplus A$, this decomposition being defined by the idempotents
$e(W_1,W_2)$ and $1-e(W_1,W_2)$.  Moreover, $C^+(Q)$ is a matrix algebra
over $Z(Q)$.

}{E7II-XII-008}From this and 1.2 one deduces, as in loc. cit.:

\underline{Proposition} 1.5. \underline{Under the hypotheses of 1.1, if
$n=2m>0$ and $Q$ is nondegenerate, then the center $Z(Q)$ of $C^+(Q)$ is
an étale algebra, locally free of rank two over $A$, and $C^+(Q)$ is an
Azumaya algebra over $Z(Q)$.}

1.6. Let $V$ and $Q$ be as in 1.4, and let $W$ be a totally singular
direct summand of $V$, locally free of rank $m$.  Let $W'_1$ be a
complement of $W$, so that $V=W\oplus W'_1$.  Every complement $W'_2$ of
$W$ identifies with the graph of a map $u:W'_1\to W$.  The form $\Phi$
places $W$ and $W'_1$ in duality, so $u$ is uniquely determined by the
bilinear form on $W'_1$


\SGApage{II}{\(\approx\)67}
\begin{equation*}
B(x,y)=\Phi(x,u(y)),
\end{equation*}

and $W'_2$ is totally singular if and only if

\begin{equation*}
B(x,x)=-Q(x)
\end{equation*}

on $W'_1$.  Since such forms $B$ always exist, $W$ has totally singular
complements.  Having chosen one of them, say $W'$, the set of totally
singular complements $W'_0$ of $W$ identifies with the set of alternating
forms on $W'$, or equivalently with the set of sections of the affine
space $\mathbb E(\bigwedge^2W')$ over $S$.

For each section $x$ of $\mathbb E(\bigwedge^2W')$, let $W'_0(x)$ be the
corresponding complement of $W$ and let
$e(x)=e(W,W'_0(x))$.  The function $x\mapsto e(x)$ is compatible with
every extension of scalars and hence defines a morphism of $S$-schemes

\begin{equation*}
e:\mathbb E(\bigwedge^2W')\longrightarrow\Spec(Z(Q)).
\end{equation*}

\SGAvar{The $S$-scheme $\mathbb E(\bigwedge^2W')$ is flat with geometrically
connected fibers, whereas $\Spec(Z(Q))$ is étale over $S$.  The morphism
$e$ therefore factors through a section of
$\mathbb E(\bigwedge^2W')$ over $S$; equivalently, $e(x)$ is independent of $x$.

}{The $S$-scheme $\mathbb E(\bigwedge^2W')$ is flat with geometrically
connected fibers, whereas $\Spec(Z(Q))$ is étale over $S$.  The morphism
$e$ therefore factors through $S$, hence through a section of
$\Spec(Z(Q))$ over $S$; equivalently, $e(x)$ is independent of $x$.

}{E7II-XII-009}This justifies the following notation.

\underline{Notation} 1.7. \underline{For $W$ as above and any totally
singular complement $W'$ of $W$, denote by $e(W)$ the idempotent
$e(W,W')\in Z(Q)$.}

\SGAvar{1.8. We saw in 1.6 that, for invertible $\lambda$, $C^+(Q)$ identifies
with $C^+(\lambda Q)$.  A continuity argument analogous to the preceding
one shows that this identification carries

}{1.8. We saw in 1.3 that, for invertible $\lambda$, $C^+(Q)$ identifies
with $C^+(\lambda Q)$.  A continuity argument analogous to the preceding
one shows that this identification carries

}{E7II-XII-010}\begin{equation*}
e(W)\in Z(Q)\qquad\text{to}\qquad e(W)\in Z(\lambda Q).
\end{equation*}




\SGApage{II}{68}
1.9. Recall that if $B$ and $C$ are two $\mathbb Z/2$-graded algebras,
then multiplication in their graded tensor product $B\otimes' C$ is
defined by the Koszul rule

\begin{equation*}
(b_1\otimes c_1)(b_2\otimes c_2)
=(-1)^{\deg(c_1)\deg(b_2)}b_1b_2\otimes c_1c_2
\end{equation*}

when the $b_i$ and $c_i$ are homogeneous.

If $(V_1,Q_1)$ and $(V_2,Q_2)$ are two quadratic modules, then

\begin{equation*}
C(Q_1\oplus Q_2)=C(Q_1)\otimes' C(Q_2)
\end{equation*}

(Bourbaki, loc. cit.).

If $P$ and $Q$ are two finite-type, locally free, $\mathbb Z/2$-graded
$A$-modules, we identify the $\mathbb Z/2$-graded algebras
$\operatorname{End}(P)\otimes'\operatorname{End}(Q)$ and
$\operatorname{End}(P\otimes Q)$ by setting, for homogeneous
$u,v,p,q$,

\begin{equation*}
(u\otimes v)(p\otimes q)
=(-1)^{\deg(v)\deg(p)}u(p)\otimes v(q).
\end{equation*}


\underline{Lemma} 1.10. \underline{Let $(V,Q)$ be the quadratic module
which is the sum of two quadratic modules $(V_1,Q_1)$ and $(V_2,Q_2)$,
where $V_i$ is locally free of rank $2m_i$ and $Q_i$ is nondegenerate.
Let $V_i=W_i\oplus W'_i$ be a decomposition of $V_i$ into two totally
singular subspaces $(i=1,2)$.  Then the diagram}

\begin{center}
\begin{tikzcd}[column sep=huge,row sep=large]
C(Q) \arrow[rr,dash,"\sim"] \arrow[d,dash,"\wr"]
  \& {} \& C(Q_1)\otimes' C(Q_2) \arrow[d,equal] \\
\operatorname{End}(\bigwedge(W_1\oplus W_2))
  \arrow[r,phantom,"\simeq"]
  \& \operatorname{End}(\bigwedge W_1\otimes\bigwedge W_2)
  \arrow[r,phantom,"\simeq"]
  \& \operatorname{End}(\bigwedge W_1)\otimes'
    \operatorname{End}(\bigwedge W_2)
\end{tikzcd}
\end{center}

\underline{is commutative.}


The verification is left to the reader.  One deduces at once from this
lemma that

\begin{equation}\tag{1.10.1}
e(W_1\oplus W_2)
=e(W_1)e(W_2)+(1-e(W_1))(1-e(W_2)).
\end{equation}

\SGApage{II}{69}
\SGAvar{1.11. The following construction is borrowed from Micali and Villamayor.
If $p:Z\longrightarrow S$ is a double étale covering of $S$, then $Z$ may
be regarded as a $\mathbb Z/2$-torsor over $S$.  Denote the addition of
torsors by $*$; for a section $s$ of $Z$, let $e(s)$ be the idempotent of
$\rho_*\mathcal O_Z$ that equals one on $s$ and zero elsewhere.  With the
notation of 1.10, there is a unique isomorphism

}{1.11. The following construction is borrowed from Micali and Villamayor.
If $p:Z\longrightarrow S$ is a double étale covering of $S$, then $Z$ may
be regarded as a $\mathbb Z/2$-torsor over $S$.  Denote the addition of
torsors by $*$; for a section $s$ of $Z$, let $e(s)$ be the idempotent of
$p_*\mathcal O_Z$ that equals one on $s$ and zero elsewhere.  With the
notation of 1.10, there is a unique isomorphism

}{E7II-XII-011}\begin{equation}\tag{1.11.1}
\Spec(Z(Q_1))*\Spec(Z(Q_2))\longrightarrow\Spec(Z(Q)),
\end{equation}

whose underlying addition morphism is

\begin{equation*}
+:\Spec(Z(Q_1))\times_S\Spec(Z(Q_2))
\longrightarrow\Spec(Z(Q)),
\end{equation*}

and which, after every base change, satisfies

\begin{equation*}
e(s+t)=e(s)e(t)+(1-e(s))(1-e(t)).
\end{equation*}

If $s(W)$ is the section of $\Spec(Z(Q))$ such that $e(s(W))=e(W)$, then
formula (1.10.1) becomes

\begin{equation*}
s(W_1\oplus W_2)=s(W_1)+s(W_2).
\end{equation*}

\SGAvar{\underline{Proposition} 1.12. \underline{Let $A$ be a field, let $V$ be
a vector space of dimension $2m>0$ over $V$, let $Q$ be a nondegenerate
quadratic form on $V$, and let $W_1,W_2$ be two totally isotropic
subspaces of $V$ of dimension $m$.  Then $e(W_1)=e(W_2)$ if and only if
$\dim(W_1/(W_1\cap W_2))$ is even.}




}{\underline{Proposition} 1.12. \underline{Let $A$ be a field, let $V$ be
a vector space of dimension $2m>0$ over $A$, let $Q$ be a nondegenerate
quadratic form on $V$, and let $W_1,W_2$ be two totally isotropic
subspaces of $V$ of dimension $m$.  Then $e(W_1)=e(W_2)$ if and only if
$\dim(W_1/(W_1\cap W_2))$ is even.}




}{E7II-XII-012}\SGApage{II}{\(\approx\)70}
(Cf. Bourbaki, \emph{Algèbre}, Ch.~9, \S6, Exercise~28.)  Let
$e_1,\ldots,e_m,f_1,\ldots,f_m$ be a basis of $V$ such that

\begin{enumerate}[label=\alph*),leftmargin=2.5em]
\item $Q(\sum_i x_i e_i+\sum_i y_i f_i)=\sum_i x_i y_i$;
\item $W_1$ is generated by the $e_i$;
\item $W_2$ is generated by $(e_i)_{1\leq i\leq s}$ and
$(f_i)_{s<i\leq m}$.
\end{enumerate}

Then $V$ is the sum of the $V_i=Ae_i+Af_i$, and

\begin{equation*}
W_\alpha=\bigoplus_i(W_\alpha\cap V_i),
\qquad \alpha=1,2.
\end{equation*}

The addition formula 1.11 therefore reduces the proof to the case
$\dim V=2$.  If $W_1=W_2$, the assertion is clear.  Otherwise, if
$W_1=Ae$ and $W_2=Af$ with $\Phi(e,f)=1$, then
$e(W_1)=fe$ and $e(W_2)=ef=1-fe$.

\medskip
2. \underline{Quadrics}

2.1. Let $k$ be an algebraically closed field and $V$ a vector space of
dimension $n+2$ over $k$.  A smooth quadric of dimension $n$ over $k$ is
a $k$-scheme $X$ isomorphic to the subscheme of $\mathbb P(V^*)$ defined
by the equation $Q=0$, for some ordinary quadratic form $Q$ over $k$.

\underline{Remark} 2.2. For $n=0,1,2$, $X$ is a smooth quadric of
dimension $n$ over $k$ if and only if

\begin{equation*}
\begin{array}{rcl}
n=0 &:& X\simeq\Spec(k)\amalg\Spec(k),\\[1ex]
n=1 &:& X\simeq\mathbb P^1_k,\\[1ex]
n=2 &:& X\simeq\mathbb P^1_k\times\mathbb P^1_k.
\end{array}
\end{equation*}



\SGApage{II}{71}
\underline{Lemma} 2.3. \underline{With the notation of 2.1, suppose that
$X$ is defined by $Q=0$.}

\begin{enumerate}[label=(\roman*),leftmargin=2.6em]
\item $\Omega_X^n\simeq\mathcal O_X(-n)$.

\item
\begin{equation*}
\operatorname{Pic}(X)\simeq
\begin{cases}
0, & n=0,\\
\mathbb Z, & n=1;\ \mathcal O_X(1)\text{ is the square of a generator},\\
\mathbb Z\oplus\mathbb Z, & n=2,\\
\mathbb Z, & n\geq 3;\ \mathcal O_X(1)\text{ is a generator}.
\end{cases}
\end{equation*}

\item $H^i(X,\mathcal O_X)=0$ for $i>0$, and
$H^1(X,\mathcal O_X(1))=0$.

\item The natural map
\begin{equation*}
H^0\bigl(\mathbb P(V^*),\mathcal O_{\mathbb P(V^*)}(1)\bigr)
=V^*\xrightarrow{\sim}H^0(X,\mathcal O_X(1))
\end{equation*}
is an isomorphism; and, for $n>0$,
$H^0(X,\mathcal O_X)=k$.
\end{enumerate}

Assertion (i) follows from the fact that the conormal sheaf of $X$ in
$\mathbb P(V^*)$ is $\mathcal O_X(-2)$ and that the canonical sheaf of
$\mathbb P^r$ is $\mathcal O(-r-1)$.  Assertion (ii) follows from 2.2
when $n\leq 2$, and from SGA~2, XII, 3.7 when $n\geq 3$.  Finally, (iii)
and (iv) follow from FAC, no.~78.

\underline{Definition} 2.4. \underline{A smooth quadric of dimension
$n$ over a scheme $S$ is an} $S$-\underline{scheme $f:X\longrightarrow S$, proper
and smooth over $S$, whose geometric fibres are smooth quadrics.}

For $n=0$, a smooth quadric of dimension zero over $S$ is simply a
double \'{e}tale covering of $S$.  For $n=1$, it is simply a
Severi--Brauer variety of dimension one over $S$.

2.5. Let $\rho:X\longrightarrow S$ be a smooth quadric of dimension $n$
over $S$.  By 2.3(i), the inverse of the invertible sheaf
$\Omega^n_{X/S}$ is ample.


\SGApage{II}{\(\approx\)72}
\SGAvar{By 2.3(ii), the Picard scheme $\operatorname{Pic}_{X/S}$ is smooth over
$S$ (because $H^2(X,\mathcal O_X)=0$), unramified over $S$ (because
$H^1(X,\mathcal O_X)=0$), and essentially proper over $S$ (because
$X/S$ is smooth).  Thus, locally on $S$ for the \'{e}tale topology, it
is a constant discrete group having the structure described in
2.3(iii).

In particular, if $X$ has a section, there is an invertible sheaf $L$ on
$X$ such that $L^{\otimes(-n)}$ and $\Omega^n_{X/S}$ have the same image
in $\operatorname{Pic}_{X/S}$; any two such sheaves are locally
isomorphic on $S$.  By 2.3(iii) and [1], 0.5, p.~19,
$p_*L$ is locally free and its formation commutes with every base
change.  Hence $\rho_*L$ is locally free of rank $n+2$, the map
$\rho^*\rho_*L\longrightarrow L$ is surjective, and the canonical map
\begin{equation*}
X\longrightarrow\mathbb P(\rho_*L)
\end{equation*}
is a closed immersion.  It identifies $X$ with a relative divisor of
degree two in $\mathbb P(\rho_*L)$; locally on $S$, $X$ is the quadric in
$\mathbb P(\rho_*L)$ defined by an ordinary quadratic form on
$(\rho_*L)^*$, unique up to an invertible factor.



}{By 2.3(iii), the Picard scheme $\operatorname{Pic}_{X/S}$ is smooth over
$S$ (because $H^2(X,\mathcal O_X)=0$), unramified over $S$ (because
$H^1(X,\mathcal O_X)=0$), and essentially proper over $S$ (because
$X/S$ is smooth).  Thus, locally on $S$ for the \'{e}tale topology, it
is a constant discrete group having the structure described in
2.3(ii).

In particular, if $X$ has a section, there is an invertible sheaf $L$ on
$X$ such that $L^{\otimes(-n)}$ and $\Omega^n_{X/S}$ have the same image
in $\operatorname{Pic}_{X/S}$; any two such sheaves are locally
isomorphic on $S$.  By 2.3(iii) and [1], 0.5, p.~19,
$\rho_*L$ is locally free and its formation commutes with every base
change.  Hence $\rho_*L$ is locally free of rank $n+2$, the map
$\rho^*\rho_*L\longrightarrow L$ is surjective, and the canonical map
\begin{equation*}
X\longrightarrow\mathbb P(\rho_*L)
\end{equation*}
is a closed immersion.  It identifies $X$ with a relative divisor of
degree two in $\mathbb P(\rho_*L)$; locally on $S$, $X$ is the quadric in
$\mathbb P(\rho_*L)$ defined by an ordinary quadratic form on
$(\rho_*L)^*$, unique up to an invertible factor.



}{E7II-XII-013}2.6. Suppose that $n>0$.  If $L_1$ and $L_2$ are two invertible sheaves
of the type considered above, then locally on $S$ there is an
isomorphism $u:L_1\xrightarrow{\sim}L_2$.  It is unique up to an
invertible factor, since $\rho_*\mathcal O_X=\mathcal O_S$ by 2.3(iv).
Consequently all the isomorphisms $u:L_1\xrightarrow{\sim}L_2$ induce
the same isomorphism
\begin{equation*}
\mathbb P(u):\mathbb P(\rho_*L_1)\xrightarrow{\sim}
\mathbb P(\rho_*L_2).
\end{equation*}
Thus $\mathbb P(\rho_*L_1)$ is canonically isomorphic to
$\mathbb P(\rho_*L_2)$ and depends only on $X/S$, not on the choice of
$L$.  We denote it by $P(X)$.


\SGApage{II}{\(\approx\)73}
The projective space $P(X)$ is defined locally on $S$ for the \'{e}tale
topology.  By descent it defines a Severi--Brauer variety over $S$,
again denoted $P(X)$, in which $X$ is a relative divisor of degree two.

Locally on $S$, the pair $X\hookrightarrow P(X)$ defines a Clifford
algebra $C^+(X,P(X))$ (1.3).  Descent yields an algebra over $S$, which
we denote simply by $C^+(X)$.

2.7. Suppose that $n=2m>0$.  The centre $Z$ of $C^+(X)$ is then an
$\mathcal O_S$-algebra defining a double \'{e}tale covering $Z(X)$ of
$S$.

A \underline{generator} of $X$ is a subscheme $D$ of $X$ which is a
linear subspace of dimension $m$ in $P(X)$.

Locally on $S$, when $P(X)\simeq\mathbb P(V^*)$ and $X$ is defined by
$Q=0$, a generator $D$ of $X$ is, by definition, identified with a
locally direct-summand submodule $W(D)$ of $V$, totally singular and of
rank $m+1$.  By 1.8, the idempotent $e(W(D))$ in the centre of $C^+(X)$
depends only on $D\subset X$.  Let $e(D)$ denote the corresponding
section of $Z(X)$, namely the one on which $e(W(D))$ has value one.

The generators of $X/S$ are the $S$-points of a scheme
$\operatorname{Gen}(X)$ projective over $S$, and $e$ defines a morphism
of $S$-schemes

\begin{equation}\tag{2.7.1}
e:\operatorname{Gen}(X)\longrightarrow Z(X).
\end{equation}

When $n=0$, we set $Z(X)=X$, call a section of $X$ a generator of $X$,
and define $e$ to be the identity.


\SGApage{II}{74}
\underline{Proposition} 2.8. \underline{Under the hypotheses of 2.7,
$p:\operatorname{Gen}(X)\longrightarrow S$ is projective and smooth,
and (2.7.1) is its Stein factorization.}

The question is local for the \'{e}tale topology on $S$.  Locally, if
$D$ is a generator, then by 1.6 the generators $D'$ disjoint from $D$
are identified with the $S$-points of a suitable affine space over $S$.
In this way $\operatorname{Gen}(X)$ is covered by open subschemes smooth
over $S$.

To prove that the geometric fibres of $e$ are geometrically connected,
it is enough to consider the case in which $S$ is the spectrum of an
algebraically closed field $k$ and $X$ is defined by $Q=0$ in
$\mathbb P(V^*)$.

\SGAvar{Let $W_1$ and $W^*_2$ be two totally isotropic subspaces of rank $m+1$ in
$V$ such that $e(W)=e(W')$, that is, such that
$W/(W\cap W')$ has even dimension (1.12).  Proceeding as in 1.12,
one obtains an orthogonal decomposition
\begin{equation*}
V=V_0\oplus\bigoplus_i V_i
\end{equation*}
such that
\begin{enumerate}[label=\alph*),leftmargin=2.5em]
\item $W_\alpha=(W_\alpha\cap V_0)\oplus
  \displaystyle\bigoplus_i(W_\alpha\cap V_i)$ for $\alpha=1,2$;
\item $W_1\cap V_0=W_2\cap V_0$;
\item $\dim V_i=4$ and
  $V_i=(W_1\cap V_i)\oplus(W_2\cap V_i)$.
\end{enumerate}
It remains to prove that $W_1\cap V_i$ and $W_2\cap V_i$ belong to the
same connected family of maximal totally singular subspaces of $V_i$.
In other words, the assertion reduces to the case of quadrics of
dimension two.  Such a quadric is isomorphic to
$\mathbb P^1\times\mathbb P^1$, and its generators are the subschemes
$\{t\}\times\mathbb P^1$ and $\mathbb P^1\times\{t\}$; the proposition
is then evident.


}{Let $W_1$ and $W_2$ be two totally isotropic subspaces of rank $m+1$ in
$V$ such that $e(W_1)=e(W_2)$, that is, such that
$W_1/(W_1\cap W_2)$ has even dimension (1.12).  Proceeding as in 1.12,
one obtains an orthogonal decomposition
\begin{equation*}
V=V_0\oplus\bigoplus_i V_i
\end{equation*}
such that
\begin{enumerate}[label=\alph*),leftmargin=2.5em]
\item $W_\alpha=(W_\alpha\cap V_0)\oplus
  \displaystyle\bigoplus_i(W_\alpha\cap V_i)$ for $\alpha=1,2$;
\item $W_1\cap V_0=W_2\cap V_0$;
\item $\dim V_i=4$ and
  $V_i=(W_1\cap V_i)\oplus(W_2\cap V_i)$.
\end{enumerate}
It remains to prove that $W_1\cap V_i$ and $W_2\cap V_i$ belong to the
same connected family of maximal totally singular subspaces of $V_i$.
In other words, the assertion reduces to the case of quadrics of
dimension two.  Such a quadric is isomorphic to
$\mathbb P^1\times\mathbb P^1$, and its generators are the subschemes
$\{t\}\times\mathbb P^1$ and $\mathbb P^1\times\{t\}$; the proposition
is then evident.


}{E7II-XII-014}\SGApage{II}{75}
\medskip
3. \underline{Cohomology of quadrics}

\SGAvar{3.1. Let $p:X\longrightarrow S$ be a smooth quadric of dimension $n$
over $S$, and let $\ell$ be a prime number invertible on $S$.  We write
\begin{equation*}
\eta\in H^0\bigl(S,R^1p_*\mathbb Z_\ell(1)\bigr)
\end{equation*}
for the ``cohomology class of a hyperplane section.''  Locally on $S$
for the \'{e}tale topology, $P(X)\simeq\mathbb P(V^*)$; by definition,
$\eta$ is then the image in
$H^0(S,R^2p_*\mathbb Z_\ell(1))$ of the first Chern class
\begin{equation*}
c_1(\mathcal O(1))\in H^2(X,\mathbb Z_\ell(1)).
\end{equation*}


}{3.1. Let $p:X\longrightarrow S$ be a smooth quadric of dimension $n$
over $S$, and let $\ell$ be a prime number invertible on $S$.  We write
\begin{equation*}
\eta\in H^0\bigl(S,R^2p_*\mathbb Z_\ell(1)\bigr)
\end{equation*}
for the ``cohomology class of a hyperplane section.''  Locally on $S$
for the \'{e}tale topology, $P(X)\simeq\mathbb P(V^*)$; by definition,
$\eta$ is then the image in
$H^0(S,R^2p_*\mathbb Z_\ell(1))$ of the first Chern class
\begin{equation*}
c_1(\mathcal O(1))\in H^2(X,\mathbb Z_\ell(1)).
\end{equation*}


}{E7II-XII-015}3.2. Suppose that $n=2m$.  A generator $D$ of $X$ then defines a
cohomology class
$c\ell'(D)\in H^n(X,\mathbb Z_\ell(m))$, whose image in
$H^0(S,R^np_*\mathbb Z_\ell(m))$ is denoted $c\ell(D)$.  By 2.8, the
latter depends only on the section $e(D)$ of $Z(X)$.  We therefore obtain
a morphism of \'{e}tale sheaves on $S$
\begin{equation}\tag{3.2.1}
c\ell:Z(X)\longrightarrow R^np_*\mathbb Z_\ell(m).
\end{equation}

\underline{Theorem} 3.3. \underline{Let $p:X\longrightarrow S$ be a
smooth quadric of dimension $n$ over $S$, and let $\ell$ be a prime
number invertible on $S$.}

\SGAvar{\begin{enumerate}[label=(\roman*),leftmargin=2.6em]
\item $R^{2i+1}p_*\mathbb Z_\ell=0$.


\item If $0\leq 2i<n$ (respectively, if $n<2i\leq 2n$), then
$R^{2i}p_*\mathbb Z_\ell(i)$ is canonically isomorphic to the constant
sheaf $\mathbb Z_\ell$, and is generated by $\eta^i$ (respectively, by
$\eta^i/2$).

\item If $n=2m$, the map deduced from (3.2.1),
\begin{equation}\tag{3.2.1}
\mathbb Z_\ell^{Z(X)}\longrightarrow R^np_*\mathbb Z_\ell(m),
\end{equation}
is an isomorphism.  If $\alpha$ and $\beta$ are two disjoint sections
of $Z(X)$ (assumed isomorphic to $S\amalg S$, as is locally the case),
then
\begin{align*}
\text{(a)}\quad &\eta^m=c\ell(\alpha)+c\ell(\beta),\\[1ex]
\text{(b)}\quad &
\begin{cases}
\operatorname{Tr}(c\ell(\alpha)^2)
=\operatorname{Tr}(c\ell(\beta)^2)=1,
& c\ell(\alpha)c\ell(\beta)=0,\quad m\text{ even},\\[1ex]
c\ell(\alpha)^2=c\ell(\beta)^2=0,
& \operatorname{Tr}(c\ell(\alpha)c\ell(\beta))=1,
  \quad m\text{ odd},
\end{cases}\\[1ex]
\text{(c)}\quad &\eta\bigl(c\ell(\alpha)-c\ell(\beta)\bigr)=0.
\end{align*}
\end{enumerate}



}{\begin{enumerate}[label=(\roman*),leftmargin=2.6em]
\item $R^{2i+1}p_*\mathbb Z_\ell=0$.


\item If $0\leq 2i<n$ (respectively, if $n<2i\leq 2n$), then
$R^{2i}p_*\mathbb Z_\ell(i)$ is canonically isomorphic to the constant
sheaf $\mathbb Z_\ell$, and is generated by $\eta^i$ (respectively, by
$\eta^i/2$).

\item If $n=2m$, the map deduced from (3.2.1),
\begin{equation}\tag{3.3.1}
\mathbb Z_\ell^{Z(X)}\longrightarrow R^np_*\mathbb Z_\ell(m),
\end{equation}
is an isomorphism.  If $\alpha$ and $\beta$ are two disjoint sections
of $Z(X)$ (assumed isomorphic to $S\amalg S$, as is locally the case),
then
\begin{align*}
\text{(a)}\quad &\eta^m=c\ell(\alpha)+c\ell(\beta),\\[1ex]
\text{(b)}\quad &
\begin{cases}
\operatorname{Tr}(c\ell(\alpha)^2)
=\operatorname{Tr}(c\ell(\beta)^2)=1,
& c\ell(\alpha)c\ell(\beta)=0,\quad m\text{ even},\\[1ex]
c\ell(\alpha)^2=c\ell(\beta)^2=0,
& \operatorname{Tr}(c\ell(\alpha)c\ell(\beta))=1,
  \quad m\text{ odd},
\end{cases}\\[1ex]
\text{(c)}\quad &\eta\bigl(c\ell(\alpha)-c\ell(\beta)\bigr)=0.
\end{align*}
\end{enumerate}



}{E7II-XII-016}\SGApage{II}{\(\approx\)76}
Assertions (i) and (ii) are recalled only for reference (XI, 1.6,
supplemented by XI, 2.6 when $2i+1=n$ in (i)).

It is enough to verify (iii) when $S$ is the spectrum of an
algebraically closed field $k$ (one could even take
$S=\Spec(\mathbb C)$).

\SGAvar{\underline{Formula (a).}  We may suppose that $X$ is the quadric in
$\mathbb P^{m+1}(k)$ with equation
\begin{equation*}
\sum_{i=0}^{m}x_i x_{i+m+1}=0.
\end{equation*}
The class $\eta^m$ is the cohomology class of the plane section defined
by
\begin{equation*}
x_i=0\qquad(0\leq i<m),
\end{equation*}
and this cycle is the sum of the generators $D_1,D_2$ defined by


}{\underline{Formula (a).}  We may suppose that $X$ is the quadric in
$\mathbb P^{2m+1}(k)$ with equation
\begin{equation*}
\sum_{i=0}^{m}x_i x_{i+m+1}=0.
\end{equation*}
The class $\eta^m$ is the cohomology class of the plane section defined
by
\begin{equation*}
x_i=0\qquad(0\leq i<m),
\end{equation*}
and this cycle is the sum of the generators $D_1,D_2$ defined by


}{E7II-XII-017}\SGApage{II}{\(\approx\)77}
\begin{align*}
D_1:&\quad x_i=0 &&(0\leq i\leq m),\\
D_2:&\quad x_i=0 &&(0\leq i<m),\qquad x_{2m+1}=0.
\end{align*}
By 1.12, $e(D_1)\neq e(D_2)$, and hence
\begin{equation*}
\eta^m=c\ell(D_1)+c\ell(D_2)
=c\ell(\alpha)+c\ell(\beta).
\end{equation*}



\underline{Formula (b).}  If generators $D_1,D_2$ are disjoint
(respectively, meet transversally in one point), then
$c\ell(D_1)c\ell(D_2)=0$ (respectively,
$\operatorname{Tr}(c\ell(D_1)c\ell(D_2))=1$).  By 1.12, when $m$ is
even we have $e(D_1)\neq e(D_2)$ (respectively,
$e(D_1)=e(D_2)$); when $m$ is odd we have $e(D_1)=e(D_2)$
(respectively, $e(D_1)\neq e(D_2)$).  Formula (b) follows.

\SGAvar{\underline{Formula (c).}  This is clear when $n=0$.  Otherwise, by
(ii), it is enough to prove
\begin{equation*}
\eta^m\bigl(c\ell(\alpha)-c\ell(\beta)=0.
\end{equation*}
Indeed,
\begin{equation*}
\operatorname{Tr}\!\left(
\eta^m(c\ell(\alpha)-c\ell(\beta)\right.
=\operatorname{Tr}\!\left(c\ell(\alpha)^2-c\ell(\beta)^2\right)=0.
\end{equation*}

\underline{Verification.}  When $m$ is even,
\begin{equation*}
2=\operatorname{Tr}(\eta^{2})
=\operatorname{Tr}\!\left(c\ell(\alpha^2
)+2c\ell(\alpha)c\ell(\beta)+c\ell(\beta)^2\right)=1+1;
\end{equation*}
when $m$ is odd,
\begin{equation*}
2=\operatorname{Tr}(\eta^{2})
=\operatorname{Tr}\!\left(c\ell(\alpha)^2
+2c\ell(\alpha)c\ell(\beta)+c\ell(\beta)^2\right)=2.
\end{equation*}


}{\underline{Formula (c).}  This is clear when $n=0$.  Otherwise, by
(ii), it is enough to prove
\begin{equation*}
\eta^m\bigl(c\ell(\alpha)-c\ell(\beta)\bigr)=0.
\end{equation*}
Indeed,
\begin{equation*}
\operatorname{Tr}\!\left(
\eta^m(c\ell(\alpha)-c\ell(\beta))\right)
=\operatorname{Tr}\!\left(c\ell(\alpha)^2-c\ell(\beta)^2\right)=0.
\end{equation*}

\underline{Verification.}  When $m$ is even,
\begin{equation*}
2=\operatorname{Tr}(\eta^{2m})
=\operatorname{Tr}\!\left(c\ell(\alpha)^2
+2c\ell(\alpha)c\ell(\beta)+c\ell(\beta)^2\right)=1+1;
\end{equation*}
when $m$ is odd,
\begin{equation*}
2=\operatorname{Tr}(\eta^{2m})
=\operatorname{Tr}\!\left(c\ell(\alpha)^2
+2c\ell(\alpha)c\ell(\beta)+c\ell(\beta)^2\right)=2.
\end{equation*}


}{E7II-XII-018}That (3.3.1) is an isomorphism now follows from the fact that the two
bilinear forms in two variables
\begin{equation*}
B(x,y)=x_1y_2+x_2y_1,
\qquad
B(x,y)=x_1y_1+x_2y_2
\end{equation*}
have discriminant $\pm1$.


\SGApage{II}{\(\approx\)78}
\underline{Verification} 3.4. Let $X$ be a smooth quadric of dimension
$n$ over a finite field $\mathbb F_q$.  In view of 3.3, the Lefschetz
trace formula gives the classical formula
\begin{equation*}
\#X(\mathbb F_q)=
\begin{cases}
\displaystyle\sum_{i=0}^{n}q^i,
  & n\text{ odd},\\[2ex]
\displaystyle\sum_{i=0}^{n}q^i+\varepsilon q^m,
  & n=2m,\quad \varepsilon=\pm1.
\end{cases}
\end{equation*}
When $n$ is even, $\varepsilon=1$ if $X$ is split, that is, if it has a
rational generator, and $\varepsilon=-1$ otherwise.

3.5. Let $f:X\longrightarrow S$ be a smooth quadric of dimension
$n=2m$.  The \underline{primitive part} of
$R^nf_*\mathbb Z_\ell(m)$ is the orthogonal complement of $\eta^m$;
with the notation of 3.3, it is generated by
$c\ell(\alpha)-c\ell(\beta)$.  The \underline{primitive quotient} of
$R^nf_*\mathbb Z_\ell(m)$ is the quotient by
$\mathbb Z_\ell\eta^m$.  The primitive part and primitive quotient of
$R^nf_*\mathbb Z_\ell$ are obtained by tensoring with
$\mathbb Z_\ell(-m)$.



3.6. Let $V$ be a locally free vector bundle of rank $n+2$ on $S$, and
let $X$ be the smooth quadric over $S$ defined by an ordinary quadratic
form $Q$ on $V$.  Let $H$ moreover be a hyperplane in
$\mathbb P(V^*)$ meeting $X$ transversally, and put $X^o=X-H$.  The
schemes $X$ and $Y=X\cap H$ are smooth quadrics over $S$, of dimensions
$n$ and $n-1$.  We may therefore use 3.3 to compute the cohomology of
$X^o=X-Y$.


\SGApage{II}{\(\approx\)79}
\begin{equation}\tag{3.6.1}
\begin{tikzcd}[column sep=huge,row sep=huge]
X^o \arrow[r,hook] \arrow[rd,"f"']
  \& X \arrow[d,"p"]
  \& Y \arrow[l,hook] \arrow[ld,"q"]\\
{} \& S \& \mathord{.}
\end{tikzcd}
\end{equation}

If $n=0$, then $Y=\varnothing$ and $X=X^o$.  Suppose that $n>0$.  We
have two long exact sequences dual to one another:
\begin{equation}\tag{3.6.2}
\cdots\xrightarrow{\partial}R^if_!\mathbb Z_\ell
\longrightarrow R^ip_*\mathbb Z_\ell
\xrightarrow{r_i}R^iq_*\mathbb Z_\ell\longrightarrow\cdots
\end{equation}
and
\begin{equation}\tag{3.6.3}
\cdots\xrightarrow{\partial}R^{i-2}q_*\mathbb Z_\ell(-1)
\longrightarrow R^ip_*\mathbb Z_\ell
\longrightarrow R^if_*\mathbb Z_\ell\longrightarrow\cdots.
\end{equation}

\SGAvar{For $n$ even, the restriction homomorphism $r_i$ is an isomorphism for
$i\neq n,2n$; for $n$ odd, it is an isomorphism for $i\neq n-1$.  If
$n=2m$ and $\alpha$ is the cohomology class of a generator, then
\begin{equation*}
r_n(\alpha)=\tfrac12\eta^m,
\end{equation*}
so $r_n$ is surjective and its kernel is the primitive part of
$R^np_*\mathbb Z_\ell$.  If $r=2m+1$, then
$r_{2m}(\eta^m)=\eta^m$, so $r_{2m}$ is injective and its cokernel is the
primitive quotient of $R^{2m}q_*\mathbb Z_\ell$.  It follows that
$R^if_!\mathbb Z_\ell=0$ for $i\neq n,2n$, and that
$R^if_!\mathbb Z_\ell$ is a twisted constant free
$\mathbb Z_\ell$-sheaf of rank one for $i=n,2n$.  By duality, or from
(3.5.3), one likewise shows that $R^if_*\mathbb Z_\ell=0$ for
$i\neq0,n$, that $f_*\mathbb Z_\ell=\mathbb Z_\ell$, and that
$R^nf_*\mathbb Z_\ell$ has rank one.



}{For $n$ even, the restriction homomorphism $r_i$ is an isomorphism for
$i\neq n,2n$; for $n$ odd, it is an isomorphism for $i\neq n-1,2n$.  If
$n=2m$ and $\alpha$ is the cohomology class of a generator, then
\begin{equation*}
r_n(\alpha)=\tfrac12\eta^m,
\end{equation*}
so $r_n$ is surjective and its kernel is the primitive part of
$R^np_*\mathbb Z_\ell$.  If $n=2m+1$, then
$r_{2m}(\eta^m)=\eta^m$, so $r_{2m}$ is injective and its cokernel is the
primitive quotient of $R^{2m}q_*\mathbb Z_\ell$.  It follows that
$R^if_!\mathbb Z_\ell=0$ for $i\neq n,2n$, and that
$R^if_!\mathbb Z_\ell$ is a twisted constant free
$\mathbb Z_\ell$-sheaf of rank one for $i=n,2n$.  By duality, or from
(3.6.3), one likewise shows that $R^if_*\mathbb Z_\ell=0$ for
$i\neq0,n$, that $f_*\mathbb Z_\ell=\mathbb Z_\ell$, and that
$R^nf_*\mathbb Z_\ell$ has rank one.



}{E7II-XII-019}Suppose that $n=2m$ and retain the notation of 3.3.


\SGApage{II}{\(\approx\)80}
\begin{center}
\begin{tikzcd}[column sep=large,row sep=large]
0 \arrow[r]
 \& R^nf_!\mathbb Z_\ell(m) \arrow[r] \arrow[d,"\varphi"]
 \& R^np_*\mathbb Z_\ell(m) \arrow[r] \arrow[d,equal]
 \& R^nq_*\mathbb Z_\ell(m) \arrow[r]
 \& 0\\
0
 \& R^nf_*\mathbb Z_\ell(m) \arrow[l]
 \& R^np_*\mathbb Z_\ell(m) \arrow[l]
 \& R^{n-2}q_*\mathbb Z_\ell(m-1) \arrow[l]
 \& 0\mathord{.} \arrow[l]
\end{tikzcd}
\end{center}

The sheaf $R^nf_!\mathbb Z_\ell(m)$ then has two natural opposite
generators $\pm\delta$, whose images in
$R^np_*\mathbb Z_\ell(m)$ are
$\pm(c\ell(\alpha)-c\ell(\beta))$.  By 3.3(iii)(b),
\begin{equation*}
\operatorname{Tr}(\delta^2)
=(c\ell(\alpha)-c\ell(\beta))^2=(-1)^m\cdot2,
\end{equation*}
and $\pm\varphi(\delta)$ is twice the natural dual generators
$\pm\delta'$ of $R^nf_*\mathbb Z_\ell(m)$, namely the opposite images
of $c\ell(\alpha)$ and $c\ell(\beta)$.

Suppose now that $n=2m+1$ and apply the notation of 3.3 to $Y$.

\begin{center}
\begin{tikzcd}[column sep=large,row sep=large]
0 \arrow[r]
 \& R^{2m}p_*\mathbb Z_\ell(m) \arrow[r]
 \& R^{2m}q_*\mathbb Z_\ell(m) \arrow[r,"\partial"]
 \& R^nf_!\mathbb Z_\ell(m) \arrow[r] \arrow[d,"\varphi"]
 \& 0\\
{}
 \& R^{2m+2}p_*\mathbb Z_\ell(m+1)
 \& R^{2m}q_*\mathbb Z_\ell(m) \arrow[l]
 \& R^nf_*\mathbb Z_\ell(m+1) \arrow[l]
 \& 0\mathord{.} \arrow[l]
\end{tikzcd}
\end{center}



The sheaf $R^nf_!\mathbb Z_\ell(m)$ then has two natural opposite
generators $\pm\delta$, namely
$\partial c\ell(\alpha)$ and $\partial c\ell(\beta)$.  We have
\begin{equation*}
\delta\,\varphi(\delta)=\delta^2
=(\partial c\ell(\alpha))^2
=\partial\bigl(c\ell(\alpha)\,\partial c\ell(\alpha)\bigr)=0,
\end{equation*}
and therefore $\varphi(\delta)=0$.  On the other hand,
$R^nf_*\mathbb Z_\ell(m+1)$ has two natural opposite generators
$\pm\delta'$, whose images in $R^{2m}q_*\mathbb Z_\ell(m)$ are
$\pm(c\ell(\alpha)-c\ell(\beta))$, and
$\operatorname{Tr}(\delta\delta')=\pm1$.


These results are assembled in the following table.

\SGApage{II}{\(\approx\)81}
\underline{Table} 3.7. \underline{Cohomology of affine quadrics of
dimension $n$, $X^o=X-Y$.}

\medskip
\noindent\textbf{A.} The cohomology is torsion-free.  The nonzero Betti
numbers are
\begin{align*}
\text{ordinary cohomology:}\quad &b_0=b_n=1,\\
\text{cohomology with compact supports:}\quad &b_{2n}=b_n=1,
\end{align*}
except that $b_0=2$ when $n=0$.

\medskip
\noindent\textbf{B.} If $n=2m>0$, then
\begin{align*}
R^nf_!\mathbb Z_\ell(m)
  &=\text{the primitive part of }R^{2m}p_*\mathbb Z_\ell(m),\\
R^nf_*\mathbb Z_\ell(m)
  &=\text{the primitive quotient of }R^{2m}p_*\mathbb Z_\ell(m).
\end{align*}
If $n=2m+1$, then
\begin{align*}
R^nf_!\mathbb Z_\ell(m)
  &=\text{the primitive quotient of }R^{2m}q_*\mathbb Z_\ell(m),\\
R^nf_*\mathbb Z_\ell(m+1)
  &=\text{the primitive part of }R^{2m}q_*\mathbb Z_\ell(m).
\end{align*}
Locally, these groups have natural generators defined up to sign;
$\delta$ denotes those for cohomology with compact supports, and
$\delta'$ those for ordinary cohomology.

\medskip
\noindent\textbf{C.} The map
\begin{equation*}
\varphi:R^nf_!\mathbb Z_\ell\longrightarrow R^nf_*\mathbb Z_\ell
\end{equation*}
is zero when $n$ is odd, and sends $\pm\delta$ to $\pm2\delta'$ when
$n>0$ is even.  Moreover,
\begin{equation*}
\operatorname{Tr}(\delta\delta')=\pm1,
\qquad
\operatorname{Tr}(\delta^2)=
\begin{cases}
0, & n=2m+1,\\
(-1)^m\cdot2, & n=2m.
\end{cases}
\end{equation*}

\underline{Verification} 3.8. The complex affine quadric $S$ in
$\mathbb C^n$ with equation $\sum_i z_i^2=1$ is diffeomorphic to the
tangent bundle of the real sphere $S\cap\mathbb R^n$.  This agrees with
A and C.

\medskip
\underline{BIBLIOGRAPHY}

\begin{thebibliography}{9}
\bibitem{MumfordGIT}
D. Mumford, \emph{Geometric Invariant Theory}, Ergebnisse \textbf{34},
Springer-Verlag, 1965.
\end{thebibliography}


\markboth{Exposé XII}{Exposé XII}

\ifnum\SGAform=1
\clearpage\section*{Critical apparatus}\begingroup\small\raggedright\sloppy
Readings are given in the exact TeX notation of the two English witnesses. Decimal suffixes identify the individual differences in each marked passage.

\par\medskip\noindent\hypertarget{app-E7II-XII-001.1}{}\textbf{E7II-XII-001.1}\quad printed folio 63.\par
\noindent Printed reading: {\ttfamily\footnotesize card\allowbreak{}}\par
\noindent Corrected reading: {\ttfamily\footnotesize char\allowbreak{}}\par
\noindent Source: {\ttfamily\footnotesize port/decisions/decis\allowbreak{}ions\_en7II\_r2\_001.js\allowbreak{}on\#C534;\allowbreak{} 7II:93;\allowbreak{} fr/sga7ii/\{diplomati\allowbreak{}c,current\}/expose\_XI\allowbreak{}I\_body.tex;\allowbreak{} folio\allowbreak{} 63;\allowbreak{} formula\allowbreak{} 63\allowbreak{}}\par

\par\medskip\noindent\hypertarget{app-E7II-XII-002.1}{}\textbf{E7II-XII-002.1}\quad printed folio 64.\par
\noindent Printed reading: {\ttfamily\footnotesize -1\allowbreak{}}\par
\noindent Corrected reading: \emph{(empty)}\par
\noindent Source: {\ttfamily\footnotesize fr/sga7ii/\{diplomati\allowbreak{}c,current\}/expose\_XI\allowbreak{}I\_body.tex;\allowbreak{} folio\allowbreak{} 64;\allowbreak{} formula\allowbreak{} 75\allowbreak{}}\par

\par\medskip\noindent\hypertarget{app-E7II-XII-003.1}{}\textbf{E7II-XII-003.1}\quad printed folio 64.\par
\noindent Printed reading: {\ttfamily\footnotesize -1\allowbreak{}}\par
\noindent Corrected reading: \emph{(empty)}\par
\noindent Source: {\ttfamily\footnotesize fr/sga7ii/\{diplomati\allowbreak{}c,current\}/expose\_XI\allowbreak{}I\_body.tex;\allowbreak{} folio\allowbreak{} 64;\allowbreak{} formula\allowbreak{} 76\allowbreak{}}\par

\par\medskip\noindent\hypertarget{app-E7II-XII-003.2}{}\textbf{E7II-XII-003.2}\quad printed folio 64.\par
\noindent Printed reading: \emph{(empty)}\par
\noindent Corrected reading: {\ttfamily\footnotesize )\allowbreak{}}\par
\noindent Source: {\ttfamily\footnotesize fr/sga7ii/\{diplomati\allowbreak{}c,current\}/expose\_XI\allowbreak{}I\_body.tex;\allowbreak{} folio\allowbreak{} 64;\allowbreak{} formula\allowbreak{} 76;\allowbreak{} port/PORT\_EN7II\_LOG.\allowbreak{}jsonl\#C535;\allowbreak{} fixes/FIXLOG\_7II\_CUR\allowbreak{}RENT.jsonl\#C535;\allowbreak{} 7II:95\allowbreak{}}\par

\par\medskip\noindent\hypertarget{app-E7II-XII-004.1}{}\textbf{E7II-XII-004.1}\quad printed folio 64.\par
\noindent Printed reading: \emph{(empty)}\par
\noindent Corrected reading: {\ttfamily\footnotesize '\allowbreak{}}\par
\noindent Source: {\ttfamily\footnotesize fr/sga7ii/\{diplomati\allowbreak{}c,current\}/expose\_XI\allowbreak{}I\_body.tex;\allowbreak{} folio\allowbreak{} 64;\allowbreak{} formula\allowbreak{} 90\allowbreak{}}\par

\par\medskip\noindent\hypertarget{app-E7II-XII-005.1}{}\textbf{E7II-XII-005.1}\quad printed folio 65.\par
\noindent Printed reading: {\ttfamily\footnotesize E\allowbreak{}}\par
\noindent Corrected reading: {\ttfamily\footnotesize V\allowbreak{}}\par
\noindent Source: {\ttfamily\footnotesize fr/sga7ii/\{diplomati\allowbreak{}c,current\}/expose\_XI\allowbreak{}I\_body.tex;\allowbreak{} folio\allowbreak{} 65;\allowbreak{} formula\allowbreak{} 108;\allowbreak{} port/decisions/decis\allowbreak{}ions\_en7II\_r2\_001.js\allowbreak{}on\#C537;\allowbreak{} fixes/FIXLOG\_7II\_CUR\allowbreak{}RENT.jsonl\#C537;\allowbreak{} 7II:99\allowbreak{}}\par

\par\medskip\noindent\hypertarget{app-E7II-XII-005.2}{}\textbf{E7II-XII-005.2}\quad printed folio 65.\par
\noindent Printed reading: {\ttfamily\footnotesize E\allowbreak{}}\par
\noindent Corrected reading: {\ttfamily\footnotesize V\allowbreak{}}\par
\noindent Source: {\ttfamily\footnotesize fr/sga7ii/\{diplomati\allowbreak{}c,current\}/expose\_XI\allowbreak{}I\_body.tex;\allowbreak{} folio\allowbreak{} 65;\allowbreak{} formula\allowbreak{} 109;\allowbreak{} port/decisions/decis\allowbreak{}ions\_en7II\_r2\_001.js\allowbreak{}on\#C537;\allowbreak{} fixes/FIXLOG\_7II\_CUR\allowbreak{}RENT.jsonl\#C537;\allowbreak{} 7II:99\allowbreak{}}\par

\par\medskip\noindent\hypertarget{app-E7II-XII-005.3}{}\textbf{E7II-XII-005.3}\quad printed folio 65.\par
\noindent Printed reading: {\ttfamily\footnotesize E\allowbreak{}}\par
\noindent Corrected reading: {\ttfamily\footnotesize V\allowbreak{}}\par
\noindent Source: {\ttfamily\footnotesize ED-EN-FORMULA-G103;\allowbreak{} fr/sga7ii/\{diplomati\allowbreak{}c,current\}/expose\_XI\allowbreak{}I\_body.tex;\allowbreak{} folio\allowbreak{} 65;\allowbreak{} port/decisions/decis\allowbreak{}ions\_en7II\_r2\_001.js\allowbreak{}on\#C537;\allowbreak{} fixes/FIXLOG\_7II\_CUR\allowbreak{}RENT.jsonl\#C537;\allowbreak{} 7II:99\allowbreak{}}\par

\par\medskip\noindent\hypertarget{app-E7II-XII-006.1}{}\textbf{E7II-XII-006.1}\quad printed folio 65.\par
\noindent Printed reading: {\ttfamily\footnotesize 2\allowbreak{}}\par
\noindent Corrected reading: {\ttfamily\footnotesize 3\allowbreak{}}\par
\noindent Source: {\ttfamily\footnotesize port/decisions/decis\allowbreak{}ions\_en7II\_r2\_001.js\allowbreak{}on\#C536;\allowbreak{} fixes/FIXLOG\_7II\_CUR\allowbreak{}RENT.jsonl\#C536;\allowbreak{} 7II:97\allowbreak{}}\par

\par\medskip\noindent\hypertarget{app-E7II-XII-007.1}{}\textbf{E7II-XII-007.1}\quad printed folio 65.\par
\noindent Printed reading: \emph{(empty)}\par
\noindent Corrected reading: {\ttfamily\footnotesize \{\allowbreak{}}\par
\noindent Source: {\ttfamily\footnotesize fr/sga7ii/\{diplomati\allowbreak{}c,current\}/expose\_XI\allowbreak{}I\_body.tex;\allowbreak{} folio\allowbreak{} 65\allowbreak{}}\par

\par\medskip\noindent\hypertarget{app-E7II-XII-007.2}{}\textbf{E7II-XII-007.2}\quad printed folio 65.\par
\noindent Printed reading: \emph{(empty)}\par
\noindent Corrected reading: {\ttfamily\footnotesize -1\}\allowbreak{}}\par
\noindent Source: {\ttfamily\footnotesize fr/sga7ii/\{diplomati\allowbreak{}c,current\}/expose\_XI\allowbreak{}I\_body.tex;\allowbreak{} folio\allowbreak{} 65\allowbreak{}}\par

\par\medskip\noindent\hypertarget{app-E7II-XII-007.3}{}\textbf{E7II-XII-007.3}\quad printed folio 65.\par
\noindent Printed reading: {\ttfamily\footnotesize 0\allowbreak{}}\par
\noindent Corrected reading: {\ttfamily\footnotesize 1\allowbreak{}}\par
\noindent Source: {\ttfamily\footnotesize fr/sga7ii/\{diplomati\allowbreak{}c,current\}/expose\_XI\allowbreak{}I\_body.tex;\allowbreak{} folio\allowbreak{} 65\allowbreak{}}\par

\par\medskip\noindent\hypertarget{app-E7II-XII-008.1}{}\textbf{E7II-XII-008.1}\quad printed folio 66.\par
\noindent Printed reading: \emph{(empty)}\par
\noindent Corrected reading: {\ttfamily\footnotesize )\textbackslash{}bigr\allowbreak{}}\par
\noindent Source: {\ttfamily\footnotesize ED-EN-FORMULA-G105;\allowbreak{} fr/sga7ii/\{diplomati\allowbreak{}c,current\}/expose\_XI\allowbreak{}I\_body.tex;\allowbreak{} folio\allowbreak{} 66\allowbreak{}}\par

\par\medskip\noindent\hypertarget{app-E7II-XII-008.2}{}\textbf{E7II-XII-008.2}\quad printed folio 66.\par
\noindent Printed reading: \emph{(empty)}\par
\noindent Corrected reading: {\ttfamily\footnotesize )\textbackslash{}bigr\allowbreak{}}\par
\noindent Source: {\ttfamily\footnotesize ED-EN-FORMULA-G105;\allowbreak{} fr/sga7ii/\{diplomati\allowbreak{}c,current\}/expose\_XI\allowbreak{}I\_body.tex;\allowbreak{} folio\allowbreak{} 66\allowbreak{}}\par

\par\medskip\noindent\hypertarget{app-E7II-XII-008.3}{}\textbf{E7II-XII-008.3}\quad printed folio 66.\par
\noindent Printed reading: {\ttfamily\footnotesize 2\allowbreak{}}\par
\noindent Corrected reading: {\ttfamily\footnotesize 1\allowbreak{}}\par
\noindent Source: {\ttfamily\footnotesize fr/sga7ii/\{diplomati\allowbreak{}c,current\}/expose\_XI\allowbreak{}I\_body.tex;\allowbreak{} folio\allowbreak{} 66;\allowbreak{} formula\allowbreak{} 168;\allowbreak{} port/PORT\_EN7II\_LOG.\allowbreak{}jsonl\#C540;\allowbreak{} fixes/FIXLOG\_7II\_CUR\allowbreak{}RENT.jsonl\#C540;\allowbreak{} 7II:104\allowbreak{}}\par

\par\medskip\noindent\hypertarget{app-E7II-XII-008.4}{}\textbf{E7II-XII-008.4}\quad printed folio 66.\par
\noindent Printed reading: {\ttfamily\footnotesize 1\allowbreak{}}\par
\noindent Corrected reading: {\ttfamily\footnotesize 2\allowbreak{}}\par
\noindent Source: {\ttfamily\footnotesize fr/sga7ii/\{diplomati\allowbreak{}c,current\}/expose\_XI\allowbreak{}I\_body.tex;\allowbreak{} folio\allowbreak{} 66;\allowbreak{} formula\allowbreak{} 168;\allowbreak{} port/PORT\_EN7II\_LOG.\allowbreak{}jsonl\#C540;\allowbreak{} fixes/FIXLOG\_7II\_CUR\allowbreak{}RENT.jsonl\#C540;\allowbreak{} 7II:104\allowbreak{}}\par

\par\medskip\noindent\hypertarget{app-E7II-XII-009.1}{}\textbf{E7II-XII-009.1}\quad printed folio 67.\par
\noindent Printed reading: \emph{(empty)}\par
\noindent Corrected reading: {\ttfamily\footnotesize \$S\$,\allowbreak{} hence\allowbreak{} through\allowbreak{} }\par
\noindent Source: {\ttfamily\footnotesize fr/sga7ii/\{diplomati\allowbreak{}c,current\}/expose\_XI\allowbreak{}I\_body.tex;\allowbreak{} folio\allowbreak{} 67;\allowbreak{} ED-EN-FORMULA-G101;\allowbreak{} fr/sga7ii/\{diplomati\allowbreak{}c,current\}/expose\_XI\allowbreak{}I\_body.tex;\allowbreak{} folio\allowbreak{} 67\allowbreak{}}\par

\par\medskip\noindent\hypertarget{app-E7II-XII-009.2}{}\textbf{E7II-XII-009.2}\quad printed folio 67.\par
\noindent Printed reading: {\ttfamily\footnotesize \textbackslash{}mathbb\allowbreak{} E\allowbreak{}}\par
\noindent Corrected reading: {\ttfamily\footnotesize \textbackslash{}Spec\allowbreak{}}\par
\noindent Source: {\ttfamily\footnotesize fr/sga7ii/\{diplomati\allowbreak{}c,current\}/expose\_XI\allowbreak{}I\_body.tex;\allowbreak{} folio\allowbreak{} 67;\allowbreak{} ED-EN-FORMULA-G101;\allowbreak{} fr/sga7ii/\{diplomati\allowbreak{}c,current\}/expose\_XI\allowbreak{}I\_body.tex;\allowbreak{} folio\allowbreak{} 67\allowbreak{}}\par

\par\medskip\noindent\hypertarget{app-E7II-XII-009.3}{}\textbf{E7II-XII-009.3}\quad printed folio 67.\par
\noindent Printed reading: {\ttfamily\footnotesize \textbackslash{}bigwedge\textasciicircum{}2W'\allowbreak{}}\par
\noindent Corrected reading: {\ttfamily\footnotesize Z(Q)\allowbreak{}}\par
\noindent Source: {\ttfamily\footnotesize fr/sga7ii/\{diplomati\allowbreak{}c,current\}/expose\_XI\allowbreak{}I\_body.tex;\allowbreak{} folio\allowbreak{} 67;\allowbreak{} ED-EN-FORMULA-G101;\allowbreak{} fr/sga7ii/\{diplomati\allowbreak{}c,current\}/expose\_XI\allowbreak{}I\_body.tex;\allowbreak{} folio\allowbreak{} 67\allowbreak{}}\par

\par\medskip\noindent\hypertarget{app-E7II-XII-010.1}{}\textbf{E7II-XII-010.1}\quad printed folio 67.\par
\noindent Printed reading: {\ttfamily\footnotesize 6\allowbreak{}}\par
\noindent Corrected reading: {\ttfamily\footnotesize 3\allowbreak{}}\par
\noindent Source: {\ttfamily\footnotesize ED-EN-NONMATH-G014;\allowbreak{} F7II-XII-006.1;\allowbreak{} 7II:109;\allowbreak{} fr/sga7ii/\{diplomati\allowbreak{}c,current\}/expose\_XI\allowbreak{}I\_body.tex;\allowbreak{} folio\allowbreak{} 67\allowbreak{}}\par

\par\medskip\noindent\hypertarget{app-E7II-XII-011.1}{}\textbf{E7II-XII-011.1}\quad printed folio 69.\par
\noindent Printed reading: {\ttfamily\footnotesize \textbackslash{}rho\allowbreak{}}\par
\noindent Corrected reading: {\ttfamily\footnotesize p\allowbreak{}}\par
\noindent Source: {\ttfamily\footnotesize fr/sga7ii/\{diplomati\allowbreak{}c,current\}/expose\_XI\allowbreak{}I\_body.tex;\allowbreak{} folio\allowbreak{} 69;\allowbreak{} formula\allowbreak{} 281;\allowbreak{} port/decisions/decis\allowbreak{}ions\_en7II\_r2\_001.js\allowbreak{}on\#C548;\allowbreak{} fixes/FIXLOG\_7II\_CUR\allowbreak{}RENT.jsonl\#C548;\allowbreak{} 7II:114\allowbreak{}}\par

\par\medskip\noindent\hypertarget{app-E7II-XII-012.1}{}\textbf{E7II-XII-012.1}\quad printed folio 69.\par
\noindent Printed reading: {\ttfamily\footnotesize V\allowbreak{}}\par
\noindent Corrected reading: {\ttfamily\footnotesize A\allowbreak{}}\par
\noindent Source: {\ttfamily\footnotesize fr/sga7ii/\{diplomati\allowbreak{}c,current\}/expose\_XI\allowbreak{}I\_body.tex;\allowbreak{} folio\allowbreak{} 69;\allowbreak{} formula\allowbreak{} 295\allowbreak{}}\par

\par\medskip\noindent\hypertarget{app-E7II-XII-013.1}{}\textbf{E7II-XII-013.1}\quad printed folio 72.\par
\noindent Printed reading: {\ttfamily\footnotesize ii\allowbreak{}}\par
\noindent Corrected reading: {\ttfamily\footnotesize iii\allowbreak{}}\par
\noindent Source: {\ttfamily\footnotesize ED-EN-NONMATH-G015;\allowbreak{} F7II-XII-012.1;\allowbreak{} 7II:119;\allowbreak{} fr/sga7ii/\{diplomati\allowbreak{}c,current\}/expose\_XI\allowbreak{}I\_body.tex;\allowbreak{} folio\allowbreak{} 72\allowbreak{}}\par

\par\medskip\noindent\hypertarget{app-E7II-XII-013.2}{}\textbf{E7II-XII-013.2}\quad printed folio 72.\par
\noindent Printed reading: {\ttfamily\footnotesize iii\allowbreak{}}\par
\noindent Corrected reading: {\ttfamily\footnotesize ii\allowbreak{}}\par
\noindent Source: {\ttfamily\footnotesize ED-EN-NONMATH-G016;\allowbreak{} F7II-XII-012.2;\allowbreak{} 7II:120;\allowbreak{} fr/sga7ii/\{diplomati\allowbreak{}c,current\}/expose\_XI\allowbreak{}I\_body.tex;\allowbreak{} folio\allowbreak{} 72\allowbreak{}}\par

\par\medskip\noindent\hypertarget{app-E7II-XII-013.3}{}\textbf{E7II-XII-013.3}\quad printed folio 72.\par
\noindent Printed reading: {\ttfamily\footnotesize p\allowbreak{}}\par
\noindent Corrected reading: {\ttfamily\footnotesize \textbackslash{}rho\allowbreak{}}\par
\noindent Source: {\ttfamily\footnotesize fr/sga7ii/\{diplomati\allowbreak{}c,current\}/expose\_XI\allowbreak{}I\_body.tex;\allowbreak{} folio\allowbreak{} 72;\allowbreak{} formula\allowbreak{} 403\allowbreak{}}\par

\par\medskip\noindent\hypertarget{app-E7II-XII-014.1}{}\textbf{E7II-XII-014.1}\quad printed folio 74.\par
\noindent Printed reading: {\ttfamily\footnotesize \textasciicircum{}*\allowbreak{}}\par
\noindent Corrected reading: \emph{(empty)}\par
\noindent Source: {\ttfamily\footnotesize fr/sga7ii/\{diplomati\allowbreak{}c,current\}/expose\_XI\allowbreak{}I\_body.tex;\allowbreak{} folio\allowbreak{} 74;\allowbreak{} formula\allowbreak{} 493;\allowbreak{} port/decisions/decis\allowbreak{}ions\_en7II\_r2\_001.js\allowbreak{}on\#C551;\allowbreak{} fixes/FIXLOG\_7II\_CUR\allowbreak{}RENT.jsonl\#C551;\allowbreak{} 7II:124\allowbreak{}}\par

\par\medskip\noindent\hypertarget{app-E7II-XII-014.2}{}\textbf{E7II-XII-014.2}\quad printed folio 74.\par
\noindent Printed reading: \emph{(empty)}\par
\noindent Corrected reading: {\ttfamily\footnotesize \_1\allowbreak{}}\par
\noindent Source: {\ttfamily\footnotesize fr/sga7ii/\{diplomati\allowbreak{}c,current\}/expose\_XI\allowbreak{}I\_body.tex;\allowbreak{} folio\allowbreak{} 74;\allowbreak{} formula\allowbreak{} 496;\allowbreak{} port/decisions/decis\allowbreak{}ions\_en7II\_r2\_001.js\allowbreak{}on\#C551;\allowbreak{} fixes/FIXLOG\_7II\_CUR\allowbreak{}RENT.jsonl\#C551;\allowbreak{} 7II:124\allowbreak{}}\par

\par\medskip\noindent\hypertarget{app-E7II-XII-014.3}{}\textbf{E7II-XII-014.3}\quad printed folio 74.\par
\noindent Printed reading: {\ttfamily\footnotesize '\allowbreak{}}\par
\noindent Corrected reading: {\ttfamily\footnotesize \_2\allowbreak{}}\par
\noindent Source: {\ttfamily\footnotesize fr/sga7ii/\{diplomati\allowbreak{}c,current\}/expose\_XI\allowbreak{}I\_body.tex;\allowbreak{} folio\allowbreak{} 74;\allowbreak{} formula\allowbreak{} 496;\allowbreak{} port/decisions/decis\allowbreak{}ions\_en7II\_r2\_001.js\allowbreak{}on\#C551;\allowbreak{} fixes/FIXLOG\_7II\_CUR\allowbreak{}RENT.jsonl\#C551;\allowbreak{} 7II:124\allowbreak{}}\par

\par\medskip\noindent\hypertarget{app-E7II-XII-014.4}{}\textbf{E7II-XII-014.4}\quad printed folio 74.\par
\noindent Printed reading: \emph{(empty)}\par
\noindent Corrected reading: {\ttfamily\footnotesize \_1\allowbreak{}}\par
\noindent Source: {\ttfamily\footnotesize fr/sga7ii/\{diplomati\allowbreak{}c,current\}/expose\_XI\allowbreak{}I\_body.tex;\allowbreak{} folio\allowbreak{} 74;\allowbreak{} formula\allowbreak{} 497;\allowbreak{} port/decisions/decis\allowbreak{}ions\_en7II\_r2\_001.js\allowbreak{}on\#C551;\allowbreak{} fixes/FIXLOG\_7II\_CUR\allowbreak{}RENT.jsonl\#C551;\allowbreak{} 7II:124\allowbreak{}}\par

\par\medskip\noindent\hypertarget{app-E7II-XII-014.5}{}\textbf{E7II-XII-014.5}\quad printed folio 74.\par
\noindent Printed reading: \emph{(empty)}\par
\noindent Corrected reading: {\ttfamily\footnotesize \_1\allowbreak{}}\par
\noindent Source: {\ttfamily\footnotesize fr/sga7ii/\{diplomati\allowbreak{}c,current\}/expose\_XI\allowbreak{}I\_body.tex;\allowbreak{} folio\allowbreak{} 74;\allowbreak{} formula\allowbreak{} 497;\allowbreak{} port/decisions/decis\allowbreak{}ions\_en7II\_r2\_001.js\allowbreak{}on\#C551;\allowbreak{} fixes/FIXLOG\_7II\_CUR\allowbreak{}RENT.jsonl\#C551;\allowbreak{} 7II:124\allowbreak{}}\par

\par\medskip\noindent\hypertarget{app-E7II-XII-014.6}{}\textbf{E7II-XII-014.6}\quad printed folio 74.\par
\noindent Printed reading: {\ttfamily\footnotesize '\allowbreak{}}\par
\noindent Corrected reading: {\ttfamily\footnotesize \_2\allowbreak{}}\par
\noindent Source: {\ttfamily\footnotesize fr/sga7ii/\{diplomati\allowbreak{}c,current\}/expose\_XI\allowbreak{}I\_body.tex;\allowbreak{} folio\allowbreak{} 74;\allowbreak{} formula\allowbreak{} 497;\allowbreak{} port/decisions/decis\allowbreak{}ions\_en7II\_r2\_001.js\allowbreak{}on\#C551;\allowbreak{} fixes/FIXLOG\_7II\_CUR\allowbreak{}RENT.jsonl\#C551;\allowbreak{} 7II:124\allowbreak{}}\par

\par\medskip\noindent\hypertarget{app-E7II-XII-015.1}{}\textbf{E7II-XII-015.1}\quad printed folio 75.\par
\noindent Printed reading: {\ttfamily\footnotesize 1p\allowbreak{}}\par
\noindent Corrected reading: {\ttfamily\footnotesize 2p\allowbreak{}}\par
\noindent Source: {\ttfamily\footnotesize fr/sga7ii/\{diplomati\allowbreak{}c,current\}/expose\_XI\allowbreak{}I\_body.tex;\allowbreak{} folio\allowbreak{} 75;\allowbreak{} formula\allowbreak{} 516\allowbreak{}}\par

\par\medskip\noindent\hypertarget{app-E7II-XII-016.1}{}\textbf{E7II-XII-016.1}\quad printed folio 76.\par
\noindent Printed reading: {\ttfamily\footnotesize 2\allowbreak{}}\par
\noindent Corrected reading: {\ttfamily\footnotesize 3\allowbreak{}}\par
\noindent Source: {\ttfamily\footnotesize ED-EN-NORM-G030;\allowbreak{} complete\allowbreak{} normalized-equal\allowbreak{} formula\allowbreak{} audit;\allowbreak{} fr/sga7ii/\{diplomati\allowbreak{}c,current\}/expose\_XI\allowbreak{}I\_body.tex;\allowbreak{} folio\allowbreak{} 76\allowbreak{}}\par

\par\medskip\noindent\hypertarget{app-E7II-XII-017.1}{}\textbf{E7II-XII-017.1}\quad printed folio 76.\par
\noindent Printed reading: {\ttfamily\footnotesize m\allowbreak{}}\par
\noindent Corrected reading: {\ttfamily\footnotesize 2m\allowbreak{}}\par
\noindent Source: {\ttfamily\footnotesize fr/sga7ii/\{diplomati\allowbreak{}c,current\}/expose\_XI\allowbreak{}I\_body.tex;\allowbreak{} folio\allowbreak{} 76;\allowbreak{} formula\allowbreak{} 560\allowbreak{}}\par

\par\medskip\noindent\hypertarget{app-E7II-XII-018.1}{}\textbf{E7II-XII-018.1}\quad printed folio 77.\par
\noindent Printed reading: \emph{(empty)}\par
\noindent Corrected reading: {\ttfamily\footnotesize \textbackslash{}bigr)\allowbreak{}}\par
\noindent Source: {\ttfamily\footnotesize fr/sga7ii/\{diplomati\allowbreak{}c,current\}/expose\_XI\allowbreak{}I\_body.tex;\allowbreak{} folio\allowbreak{} 77;\allowbreak{} formula\allowbreak{} 585;\allowbreak{} port/decisions/decis\allowbreak{}ions\_en7II\_r2\_001.js\allowbreak{}on\#C552;\allowbreak{} fixes/FIXLOG\_7II\_CUR\allowbreak{}RENT.jsonl\#C552;\allowbreak{} 7II:129\allowbreak{}}\par

\par\medskip\noindent\hypertarget{app-E7II-XII-018.2}{}\textbf{E7II-XII-018.2}\quad printed folio 77.\par
\noindent Printed reading: \emph{(empty)}\par
\noindent Corrected reading: {\ttfamily\footnotesize )\allowbreak{}}\par
\noindent Source: {\ttfamily\footnotesize fr/sga7ii/\{diplomati\allowbreak{}c,current\}/expose\_XI\allowbreak{}I\_body.tex;\allowbreak{} folio\allowbreak{} 77;\allowbreak{} formula\allowbreak{} 586\allowbreak{}}\par

\par\medskip\noindent\hypertarget{app-E7II-XII-018.3}{}\textbf{E7II-XII-018.3}\quad printed folio 77.\par
\noindent Printed reading: {\ttfamily\footnotesize .\allowbreak{}}\par
\noindent Corrected reading: {\ttfamily\footnotesize )\allowbreak{}}\par
\noindent Source: {\ttfamily\footnotesize port/decisions/decis\allowbreak{}ions\_en7II\_r2\_001.js\allowbreak{}on\#C553;\allowbreak{} fixes/FIXLOG\_7II\_CUR\allowbreak{}RENT.jsonl\#C553;\allowbreak{} 7II:130\allowbreak{}}\par

\par\medskip\noindent\hypertarget{app-E7II-XII-018.4}{}\textbf{E7II-XII-018.4}\quad printed folio 77.\par
\noindent Printed reading: {\ttfamily\footnotesize 2\allowbreak{}}\par
\noindent Corrected reading: {\ttfamily\footnotesize 2m\allowbreak{}}\par
\noindent Source: {\ttfamily\footnotesize fr/sga7ii/\{diplomati\allowbreak{}c,current\}/expose\_XI\allowbreak{}I\_body.tex;\allowbreak{} folio\allowbreak{} 77;\allowbreak{} formula\allowbreak{} 588\allowbreak{}}\par

\par\medskip\noindent\hypertarget{app-E7II-XII-018.5}{}\textbf{E7II-XII-018.5}\quad printed folio 77.\par
\noindent Printed reading: \emph{(empty)}\par
\noindent Corrected reading: {\ttfamily\footnotesize )\allowbreak{}}\par
\noindent Source: {\ttfamily\footnotesize fr/sga7ii/\{diplomati\allowbreak{}c,current\}/expose\_XI\allowbreak{}I\_body.tex;\allowbreak{} folio\allowbreak{} 77;\allowbreak{} formula\allowbreak{} 588\allowbreak{}}\par

\par\medskip\noindent\hypertarget{app-E7II-XII-018.6}{}\textbf{E7II-XII-018.6}\quad printed folio 77.\par
\noindent Printed reading: {\ttfamily\footnotesize )\allowbreak{}}\par
\noindent Corrected reading: \emph{(empty)}\par
\noindent Source: {\ttfamily\footnotesize fr/sga7ii/\{diplomati\allowbreak{}c,current\}/expose\_XI\allowbreak{}I\_body.tex;\allowbreak{} folio\allowbreak{} 77;\allowbreak{} formula\allowbreak{} 588\allowbreak{}}\par

\par\medskip\noindent\hypertarget{app-E7II-XII-018.7}{}\textbf{E7II-XII-018.7}\quad printed folio 77.\par
\noindent Printed reading: {\ttfamily\footnotesize 2\allowbreak{}}\par
\noindent Corrected reading: {\ttfamily\footnotesize 2m\allowbreak{}}\par
\noindent Source: {\ttfamily\footnotesize fr/sga7ii/\{diplomati\allowbreak{}c,current\}/expose\_XI\allowbreak{}I\_body.tex;\allowbreak{} folio\allowbreak{} 77;\allowbreak{} formula\allowbreak{} 590\allowbreak{}}\par

\par\medskip\noindent\hypertarget{app-E7II-XII-019.1}{}\textbf{E7II-XII-019.1}\quad printed folio 79.\par
\noindent Printed reading: \emph{(empty)}\par
\noindent Corrected reading: {\ttfamily\footnotesize ,2n\allowbreak{}}\par
\noindent Source: {\ttfamily\footnotesize fr/sga7ii/\{diplomati\allowbreak{}c,current\}/expose\_XI\allowbreak{}I\_body.tex;\allowbreak{} folio\allowbreak{} 79;\allowbreak{} formula\allowbreak{} 641;\allowbreak{} port/PORT\_EN7II\_LOG.\allowbreak{}jsonl\#C554;\allowbreak{} fixes/FIXLOG\_7II\_CUR\allowbreak{}RENT.jsonl\#C554;\allowbreak{} 7II:136\allowbreak{}}\par

\par\medskip\noindent\hypertarget{app-E7II-XII-019.2}{}\textbf{E7II-XII-019.2}\quad printed folio 79.\par
\noindent Printed reading: {\ttfamily\footnotesize r\allowbreak{}}\par
\noindent Corrected reading: {\ttfamily\footnotesize n\allowbreak{}}\par
\noindent Source: {\ttfamily\footnotesize fr/sga7ii/\{diplomati\allowbreak{}c,current\}/expose\_XI\allowbreak{}I\_body.tex;\allowbreak{} folio\allowbreak{} 79;\allowbreak{} formula\allowbreak{} 647\allowbreak{}}\par

\par\medskip\noindent\hypertarget{app-E7II-XII-019.3}{}\textbf{E7II-XII-019.3}\quad printed folio 79.\par
\noindent Printed reading: {\ttfamily\footnotesize 5\allowbreak{}}\par
\noindent Corrected reading: {\ttfamily\footnotesize 6\allowbreak{}}\par
\noindent Source: {\ttfamily\footnotesize ED-EN-NONMATH-G017;\allowbreak{} F7II-XII-020.3;\allowbreak{} 7II:135;\allowbreak{} fr/sga7ii/\{diplomati\allowbreak{}c,current\}/expose\_XI\allowbreak{}I\_body.tex;\allowbreak{} folio\allowbreak{} 79\allowbreak{}}\par

\par\endgroup\normalsize\fi
\clearpage\SGAtoc{Exposé XIII}
\selectlanguage{english}
\clearpage
\SGApage{II}{82}

\begin{center}
{\large SGA 7 --- EXPOS\'E XIII}\par
\vspace{2em}
{\Large\underline{\underline{The Formalism of Vanishing Cycles}}}\par
\vspace{1.5em}
{\large by P. Deligne}
\end{center}

\SGApage{II}{\(\approx\)83}
\medskip
\noindent\textbf{Source contents}
\begin{center}
\begin{tabular}{@{}llr@{}}
Introduction && 2\\
\S1. & Sheaves of sets & 6\\
1.1 & Galois sheaves & 6\\
1.2 & The topos $Y\times_s S$ & 8\\
1.3 & The functor $\Psi$ & 13\\
1.4 & Variation & 16\\
\S2. & Vanishing cycles & 17\\
2.1 & The functor $R\Psi$ & 17\\
2.2 & A compatibility (trace morphisms) & 26\\
2.3 & Finiteness theorem (equal characteristic zero) & 29\\
2.4 & Isolated singularities & 30
\end{tabular}
\end{center}


\medskip
\underline{Introduction}

In this expos\'e we take up again the formalism introduced in no.~2 of
Grothendieck's Expos\'e I.  The only new point is the introduction of the
algebraic analogue of \underline{variation} in the classical theory.

The topological picture I have in mind is as follows.  Let $X$ be an
analytic space, let $D\subset\mathbb C$ be a disc centred at $0$, put
$D^*=D-\{0\}$, and let $f:X\longrightarrow D$ be a proper morphism (the
nonproper case will be considered later).  Replace $D$ by a sufficiently
small disc centred at $0$, replace $X$ by the inverse image of this small
disc, and continue to denote the result by $f:X\longrightarrow D$.
According to the complex-analytic variant of Thom's isotopy theorem,
$f^{-1}(D^*)$ is a topological fibre bundle over $D^*$.  I do not know
whether one can always construct a pair $(\Gamma,r)$ of the kind
described below, but it seems plausible that one can.

\begin{enumerate}[label=(\alph*),leftmargin=2.6em]
\item $\Gamma$ is a locally coherent system of local trivializations of
the bundle $f^{-1}(D^*)$ over $D^*$.  In other words, $\Gamma$ defines a
$\pi_1(D^*)$-equivariant trivialization of the inverse image of $X$ on a
universal covering $\widetilde D^*$ of $D^*$.

\item $r$ is a retraction of $X$ onto the special fibre
$X_0=f^{-1}(0)$, compatible with $\Gamma$: if, on a connected open
subset $U$ of $D^*$, $\Gamma$ identifies
$X_U=f^{-1}(U)$ with $X_1\times U$, then $r=r(x_1,u)$ is independent of
$u$.
\end{enumerate}

Let $(X,f,\Gamma,r)$ be given.  Let $t\in D$ and
$X_t=f^{-1}(t)$ be the ``general fibre.''  The bundle $\varphi^*X$,
obtained from $X/D$ by pullback along
\begin{equation*}
\varphi:[0,1]\longrightarrow D^*,
\qquad x\longmapsto \exp(2\pi i x)\,t,
\end{equation*}
is trivialized by $\Gamma$.  The \underline{monodromy}
$T:X_t\longrightarrow X_t$ is the homeomorphism
\begin{equation*}
X_t=(\varphi^*X)_0\xrightarrow[\Gamma]{\sim}
(\varphi^*X)_1=X_t.
\end{equation*}
Let $r_t:X_t\longrightarrow X_0$ be the restriction of $r$ to $X_t$.
Then $r_tT=r_t$.  One reconstructs $(X,f)$ from
$(X_t,X_0,T,r_t)$ as follows:
\begin{enumerate}[label=(\alph*),leftmargin=2.6em]
\item identify the ends $X_t\times\{0\}$ and $X_t\times\{1\}$ of
$X_t\times[0,1]$ using $T$, obtaining $f':X'\longrightarrow S^1$;
\item $r_t$ defines $r':X'\longrightarrow X_0$, and $(X,f)$ is
identified with the map of mapping cones
\begin{equation*}
\operatorname{cone}(f'):
\operatorname{cone}(r':X'\longrightarrow X_0)
\longrightarrow D=\operatorname{cone}(S^1\longrightarrow\mathrm{pt}).
\end{equation*}
\end{enumerate}


\SGApage{II}{\(\approx\)84}
Put $\Psi=(r,f):X\longrightarrow X_0\times D$.  The higher Leray direct
images $R^i\Psi_*\mathbb Z$ satisfy the following condition:

\begin{quote}
(*) The space $D^*$ is covered by open subsets $U$ such that the
restriction of the sheaf $F$ to $X_0\times U$ is pulled back from a sheaf
on $X_0$.
\end{quote}

\underline{Construction} 0.1. \underline{The sheaves $F$ on
$X_0\times D$ satisfying (*) are identified with systems
$(F_0,F_t,T,\alpha)$ in which}
\begin{enumerate}[label=\alph*),leftmargin=2.5em]
\item \underline{$F_0$ is a sheaf on $X_0$};
\item \underline{$F_t$ is a sheaf on $X_0$, endowed with an action of
the generator $T$ of $\pi_1(D^*,t)$};
\item \underline{$\alpha$ is a morphism from $F_0$ to the sheaf of}
$T$-\underline{invariants of $F_t$}.
\end{enumerate}

Let $G$ be a sheaf on $X_0\times D^*$ satisfying (*), and let
$\widetilde D^*$ be the universal covering of $D^*$, with base point
$t$.  The pullback of $G$ to $X_0\times\widetilde D^*$ is, in one and
only one way, pulled back from a sheaf $G_t$ on $X_0$.  The sheaf $G_t$
is endowed with an action of $\pi_1(D^*,t)$, and $G$ is obtained from
$\operatorname{pr}_1^*G_t$ on $X_0\times\widetilde D^*$ by taking the
quotient by $\pi_1(D^*,t)$.

A sheaf $F$ on $X_0\times D$ is determined by
\begin{enumerate}[label=\arabic*),leftmargin=2.5em]
\item its restriction $F_0$ to $X_0$;
\item its restriction $G$ to $X_0\times D^*$;
\item a morphism $\alpha:F_0\longrightarrow i^*j_*G$, where
$i:X_0\hookrightarrow X_0\times D$ and
$j:X_0\times D^*\hookrightarrow X_0\times D$.
\end{enumerate}
If $F$ satisfies (*), then $G$ is described by $(G_t,T)$ on $X_0$, and
the asserted description of $F$ follows.

For the sheaves $R^i\Psi_*\mathbb Z$ we have
\begin{equation*}
(R^i\Psi_*\mathbb Z)_0=
\begin{cases}
\mathbb Z, & i=0,\\
0, & i>0,
\end{cases}
\end{equation*}
whereas, for $i>0$, the $(R^i\Psi_*\mathbb Z)_t$ deserve the name
\underline{sheaves of vanishing cycles} (cf. the discussion at the
beginning of paragraph 2 of Expos\'e I).


\SGApage{II}{\(\approx\)85}
This discussion is useful only after it is refined by passing to the
derived category: the $R^i\Psi_*\mathbb Z$ are obtained from an object
of $D^+(X_0\times D)$, and indeed from an object
$R\Psi_*\mathbb Z$ of the derived category of the category of abelian
sheaves on $X_0\times D$ satisfying (*).  This object determines not
only the sheaves of vanishing cycles and their monodromy, but also the
classical ``variation'' morphisms (cf. 1.4 and XIV, 3.1).

We shall define its analogue in algebraic geometry.  If $S$ is a
strictly henselian trait and $f:X\longrightarrow S$ is a morphism of
schemes, we shall define an object $R\Psi(\mathbb Z/\ell^n)$ of the
derived category of the category of abelian sheaves on the product
topos $(X_0)_{\mathrm{et}}\times S_{\mathrm{et}}$.  A sheaf on this
topos may be interpreted as a system $(F_0,F_t,T,\alpha)$ in which
\begin{enumerate}[label=\alph*),leftmargin=2.5em]
\item[$\alpha$)] $F_0$ is a sheaf on $X_0$;
\item[$\beta$)] $F_t$ is a sheaf on $X_0$, endowed with a continuous action $T$
of the inertia group $I$;
\item[$\gamma$)] $\alpha$ is a morphism from $F_0$ to the sheaf of $I$-invariants
of $F_t$.
\end{enumerate}

We shall have to work in a slightly more general framework than that
described above:
\begin{enumerate}[label=\alph*),leftmargin=2.5em]
\item The construction of $R\Psi(\mathbb Z/\ell^n)$ does not require
$f$ to be proper.  For a topological interpretation, one would repeat
the preceding construction after replacing $X$ by a neighbourhood of
the special fibre $X_0$.

\item To construct $R\Psi(\mathbb Z/\ell^n)$ and prove its properties,
one must consider not only the constant sheaf $\mathbb Z/\ell^n$, but
also general sheaves and even complexes of sheaves.

\item When the residue field $k$ of the henselian trait $S$ is not
assumed separably closed, the product topos
$(X_0)_{\mathrm{et}}\times S_{\mathrm{et}}$ must be replaced by the
fibred product topos
\begin{equation*}
(X_0)_{\mathrm{et}}
\times_{\Spec(k)_{\mathrm{et}}}S_{\mathrm{et}}.
\end{equation*}
The reader who dislikes $2$-fibred products of topoi may take the Galois
description of this topos in 1.2.4 as its definition.
\end{enumerate}


\SGApage{II}{86}
0.2. \underline{Notation and terminology}.

\begin{description}
\item[(0.2.1) \underline{trait}:]
the spectrum of a discrete valuation ring.

\item[(0.2.2) \underline{henselian}:]
for a scheme, the spectrum of a henselian local ring.

\item[(0.2.3) \underline{strictly henselian}:]
for a ring, henselian with separably closed residue field; for a scheme,
the spectrum of such a ring.

\item[(0.2.4) \underline{geometric point} of $S$:]
a morphism $\bar x:\Spec(\bar k)\longrightarrow S$ with image $x\in S$,
where $k(\bar x)=_{\mathrm{def}}\bar k$ is a separable closure of
$\kappa(x)$.  We say that $\bar x$ is localized at $x$.  By abuse of
language, we shall sometimes also call a morphism
$\bar x:\Spec(\bar k)\longrightarrow S$ a geometric point when $\bar k$
is a separably closed extension of $\kappa(x)$; such a morphism factors
uniquely through a geometric point of $S$ in the proper sense.

\item[(0.2.5) $s,\eta,\bar s,\bar\eta,\bar S$:]
If $S=\Spec(V)$ is a henselian trait, we shall normally denote by $s$
the closed point of $S$, by $\eta$ its generic point, by $\bar\eta$ a
generic geometric point of $S$ (localized at $\eta$), and by $\bar s$
the corresponding geometric point localized at $s$ (I, 0.0.3).  Let
$\bar S$ denote the spectrum of the normalization of $V$ in
$k(\bar\eta)$; its residue field is an inseparable extension of
$k(\bar s)$.  Write $\operatorname{Gal}(\bar\eta/\eta)$ for the Galois
group of $k(\bar\eta)$ over $k(\eta)$, and similarly for $s$.  There is
an exact sequence
\begin{equation*}
0\longrightarrow I\longrightarrow
\operatorname{Gal}(\bar\eta/\eta)\longrightarrow
\operatorname{Gal}(\bar s/s)\longrightarrow0,
\end{equation*}
where $I$ is the inertia group.

\item[(0.2.6) \underline{sheaf}:]
on a scheme $X$, always means a sheaf on the \'{e}tale site
$X_{\mathrm{et}}$ of $X$.

\item[(0.2.7) \underline{prime to}:]
a torsion sheaf $F$ on a scheme $S$ (or a torsion abelian group $F$) is
prime to the residue characteristics of $S$ if, for every $s\in S$,
multiplication by the characteristic exponent of $k(s)$ is an
automorphism of $F$.
\end{description}



\SGApage{II}{87}
\medskip
1. \underline{Sheaves of sets}

The sheaves of vanishing cycles will be defined as values of derived
functors of the left exact functor $\Psi$ defined in no.~3 for every
sheaf of sets.

\medskip
1.1. \underline{Galois sheaves.}

1.1.1. Let $Y$ be a scheme over a field $k$ with separable closure
$\bar{k}$, and put $\bar{Y}=Y\otimes_k\bar{k}$.  Let $G$ be a profinite
group and let
\begin{equation*}
u:G\longrightarrow\operatorname{Gal}(\bar{k}/k)
\end{equation*}
be a continuous homomorphism from $G$ to the Galois group of $\bar{k}$
over $k$.  The Galois group $\operatorname{Gal}(\bar{k}/k)$ acts on the
left on $\bar{Y}$ by transport of structure.  The group $G$ acts on
$\bar{Y}$ through $u$.

Let $\mathfrak F$ be a sheaf of sets on (the \'{e}tale site of) $\bar Y$.
An action of $G$ on $\mathfrak F$ \underline{compatible} with the action
of $G$ on $\bar Y$ is an action of $G$ by automorphisms of the pair
$(\bar Y,\mathfrak F)$ which induces the given action on $\bar Y$.  In
other words, it is a system of isomorphisms
\begin{equation*}
\sigma(g):u(g)_*\mathfrak F\longrightarrow\mathfrak F
\end{equation*}
satisfying $\sigma(gh)=\sigma(g)\,\sigma(h)$.  If $G$ acts on
$\mathfrak F$, then for every $U$ \'{e}tale over $Y$, with
$\bar U=U\otimes_k\bar k$ \'{e}tale over $\bar Y$, the group $G$ acts on
$\mathfrak F(\bar U)$.

\underline{Definition} 1.1.2. \underline{We say that $G$ acts
continuously on $\mathfrak F$ if, for every quasi-compact $U$ \'{e}tale
over $Y$, the action of $G$ on the discrete set $\mathfrak F(\bar U)$ is
continuous.}

It would amount to the same thing to consider only affine $U$.

If $\mathfrak F$ is a sheaf of sets on $Y$, with inverse image
$\bar{\mathfrak F}$ on $\bar Y$, then
$\operatorname{Gal}(\bar k/k)$ acts on $\bar{\mathfrak F}$ by transport
of structure, compatibly with its action on $\bar Y$.

\underline{Reminder 1.1.3.}
\begin{enumerate}[label=(\roman*),leftmargin=2.7em]
\item \underline{The action of $\operatorname{Gal}(\bar k/k)$ on
$\bar{\mathfrak F}$ is continuous.}


\item \underline{The functor}
\begin{equation*}
\mathfrak F\longmapsto\bar{\mathfrak F}
\end{equation*}
\underline{where $\bar{\mathfrak F}$ is endowed with the action of
$\operatorname{Gal}(\bar k/k)$, is an equivalence from the category of sheaves of sets on
$Y$ to the category of sheaves of sets on $\bar Y$ endowed with a
continuous action of $\operatorname{Gal}(\bar k/k)$ compatible with the
action of $\operatorname{Gal}(\bar k/k)$ on $\bar Y$.}
\end{enumerate}

\SGApage{II}{\(\approx\)88}
Write $\bar k$ as the direct limit of finite Galois extensions $k_i$.
Let
\begin{equation*}
\varphi_i:Y_i=Y\otimes_k k_i\longrightarrow Y,
\qquad
\bar\varphi:\bar Y\longrightarrow Y
\end{equation*}
be the projections.  For $U$ \'{e}tale over $Y$, assumed affine, put
$U_i=U\otimes_k k_i$.  Then
\begin{equation*}
\bar{\mathfrak F}(\bar U)
=\varinjlim_i(\varphi_i^*\mathfrak F)(U_i).
\end{equation*}
On $(\varphi_i^*\mathfrak F)(U_i)$, the group
$\operatorname{Gal}(\bar k/k)$ acts through
$\operatorname{Gal}(k_i/k)$, which proves (i).

The functor in (ii) has as right adjoint the functor
\begin{equation*}
\mathfrak G\longmapsto
(\bar\varphi_*\mathfrak G)^{\operatorname{Gal}(\bar k/k)}.
\end{equation*}
For every $\mathfrak F$ on $Y$, one has
\begin{equation*}
\bar\varphi_*\bar\varphi^*\mathfrak F
=\varinjlim_i\varphi_{i*}\varphi_i^*\mathfrak F.
\end{equation*}
Thus, to prove that
\begin{equation*}
\mathfrak F\xrightarrow{\sim}
(\bar\varphi_*\bar\varphi^*\mathfrak F)^{
  \operatorname{Gal}(\bar k/k)},
\end{equation*}
it is enough to prove, for every $i$, that
\begin{equation*}
\mathfrak F\xrightarrow{\sim}
(\varphi_{i*}\varphi_i^*\mathfrak F)^{
  \operatorname{Gal}(k_i/k)}.
\end{equation*}
It is even enough to prove that this arrow becomes an isomorphism after
extension of scalars from $k$ to $k_i$; then $\operatorname{Spec}(k_i)$
becomes a sum of copies of $\operatorname{Spec}(k)$ permuted simply
transitively by $\operatorname{Gal}(k_i/k)$, and the assertion becomes
trivial.

The exact functor in (ii) is therefore fully faithful.  To conclude, it
remains to verify that the arrow
\begin{equation*}
\bar\varphi^*(\bar\varphi_*\mathfrak G)^{
  \operatorname{Gal}(\bar k/k)}\longrightarrow\mathfrak G
\end{equation*}
is an epimorphism.  Let $\bar\varphi_i:\bar Y\longrightarrow Y_i$.
The hypothesis implies that
\begin{equation*}
\varinjlim_i
\bar\varphi_i^*(\bar\varphi_{i*}\mathfrak G)^{
  \operatorname{Gal}(\bar k/k_i)}
\longrightarrow\mathfrak G
\end{equation*}

\SGApage{II}{\(\approx\)89}
is an epimorphism.

\SGAvar{We have
\begin{equation*}
\bar\varphi^*(\varphi_*\mathfrak G)^{
  \operatorname{Gal}(\bar k/k)}
=
\bar\varphi_i^*\!\left(
  \varphi_i^*\varphi_{i*}
  \left((\bar\varphi_{i*}\mathfrak G)^{
    \operatorname{Gal}(\bar k/k_i)}\right)^{
    \operatorname{Gal}(k_i/k)}
\right),
\end{equation*}
and the conclusion follows by observing, by the same Galois-descent
argument as above, that for $\mathfrak G'$ on $Y_i$, endowed with an
action of $\operatorname{Gal}(k_i/k)$ compatible with its action on
$Y_i$, one has
\begin{equation*}
\varphi_i^*(\varphi_{i*}\mathfrak G')^{
  \operatorname{Gal}(k_i/k)}
\xrightarrow{\sim}\mathfrak G'.
\end{equation*}


}{We have
\begin{equation*}
\bar\varphi^*(\bar\varphi_*\mathfrak G)^{
  \operatorname{Gal}(\bar k/k)}
=
\bar\varphi_i^*\!\left(
  \varphi_i^*\varphi_{i*}
  \left((\bar\varphi_{i*}\mathfrak G)^{
    \operatorname{Gal}(\bar k/k_i)}\right)^{
    \operatorname{Gal}(k_i/k)}
\right),
\end{equation*}
and the conclusion follows by observing, by the same Galois-descent
argument as above, that for $\mathfrak G'$ on $Y_i$, endowed with an
action of $\operatorname{Gal}(k_i/k)$ compatible with its action on
$Y_i$, one has
\begin{equation*}
\varphi_i^*(\varphi_{i*}\mathfrak G')^{
  \operatorname{Gal}(k_i/k)}
\xrightarrow{\sim}\mathfrak G'.
\end{equation*}


}{E7II-XIII-001}\subsection*{1.2. The topos $Y\times_s S$}

\subsubsection*{1.2.1.}
Let $S$ be a henselian local scheme with closed point
$i:s\hookrightarrow S$.  The functor
\begin{equation*}
S'\longmapsto S'\times_Ss=i^*S'
\end{equation*}
is an equivalence from the category of finite \'{e}tale $S$-schemes
$S'$ to the category of finite \'{e}tale $s$-schemes $s'$ (spectra of
\'{e}tale $k(s)$-algebras).  The inverse functor $\mathrm{sp}^*$, from
the \'{e}tale site of $s$ to the \'{e}tale site of $S$, is a morphism
of sites
\begin{equation*}
\mathrm{sp}:S\longrightarrow s,
\end{equation*}
the \underline{specialization morphism} (cf. SGA~4, VIII, 7).  If
$s'$ is \'{e}tale over $s$, one says that $\mathrm{sp}^*(s')$ is
obtained from $S$ by \underline{extension of the residue field}.  For
a sheaf $F$ on $S$, there is a functorial isomorphism
\begin{equation}\tag{1.2.1.1}
\mathrm{sp}_*F=i^*F.
\end{equation}
For every open subscheme of $S$, $j:U\hookrightarrow S$, or for every
point of $S$, $j:\eta\hookrightarrow S$, we shall again write
$\mathrm{sp}$ for the composite morphism $\mathrm{sp}\circ j$.  For a
sheaf $F$ on $U$ (respectively on $\eta$), one has
\begin{equation}\tag{1.2.1.2}
\mathrm{sp}_*F=i^*j_*F.
\end{equation}

\subsubsection*{1.2.2.}
Let $S$ be a henselian trait, and let
$s,\eta,\bar s,\bar\eta,I$ be as in 0.2.5.
\begin{center}
\begin{tikzcd}
s \arrow[r, hook, "i"] \& S \& \eta \arrow[l, hook, "j"']
\end{tikzcd}
.\end{center}
Here is the Galois-theoretic interpretation of sheaves on $s$, $\eta$,
and $S$, and of the specialization morphism.

\begin{enumerate}[label=(\alph*)]
\item Sheaves $F$ on $s$ (respectively on $\eta$) are identified, by
$F\mapsto F_{\bar s}$ (respectively $F\mapsto F_{\bar\eta}$), with
sets endowed with a continuous action of $\operatorname{Gal}(\bar s/s)$
(respectively $\operatorname{Gal}(\bar\eta/\eta)$) (SGA~4, VIII, 2.1,
or 1.1.3).

\item A sheaf $F$ on $S$ defines sheaves $F_s=i^*F$ and
$F_\eta=j^*F$ on $s$ and $\eta$, and an adjunction arrow
$F\longrightarrow j_*j^*F$ inducing
\begin{equation*}
\varphi:F_s=i^*F\longrightarrow i^*j_*j^*F
=\mathrm{sp}_*F_\eta.
\end{equation*}
It is known (SGA~4, IV, 9.5.4) that the functor
\begin{equation*}
F\longmapsto(F_s,F_\eta,\varphi)
\end{equation*}
is an equivalence from the category of sheaves on $S$ to the category
of triples $(F_s,F_\eta,\varphi)$, where $F_s$ is a sheaf on $s$,
$F_\eta$ is a sheaf on $\eta$, and
$\varphi:F_s\longrightarrow i^*j_*F_\eta$.

Under the equivalence in (a), the functor $i^*j_*$ is identified with
the functor of ``invariants under $I$''.  Thus sheaves on $S$ are
identified with triples consisting of a set $F_{\bar s}$ endowed with
an action of $\operatorname{Gal}(\bar s/s)$, a set $F_{\bar\eta}$
endowed with an action of $\operatorname{Gal}(\bar\eta/\eta)$, and an
equivariant morphism
$\varphi:F_{\bar s}\longrightarrow F_{\bar\eta}$.

\item Under (a) and (b), the morphisms $\mathrm{sp}$, $j$, $i$, and
$\mathrm{sp}=\mathrm{sp}\circ j$ are expressed as follows:
\begin{equation*}
\begin{array}{llcl}
\mathrm{sp}:S\longrightarrow s:\;&
 \mathrm{sp}^*(F_{\bar s})&=&(F_{\bar s},F_{\bar s},\mathrm{Id}),\\[2pt]
&\mathrm{sp}_*(F_{\bar s},F_{\bar\eta},\varphi)&=&F_{\bar s},\\[3pt]
j:\eta\longrightarrow S:\;&
 j^*(F_{\bar s},F_{\bar\eta},\varphi)&=&F_{\bar\eta},\\[2pt]
&j_*(F_{\bar\eta})&=&(F_{\bar\eta}^{I},F_{\bar\eta},
 \text{inclusion}),\\[3pt]
i:s\longrightarrow S:\;&
 i^*(F_{\bar s},F_{\bar\eta},\varphi)&=&F_{\bar s},\\[2pt]
&i_*(F_{\bar s})&=&(F_{\bar s},\{0\},\text{projection}),\\[3pt]
\mathrm{sp}=\mathrm{sp}\circ j:\eta\longrightarrow s:\;&
 \mathrm{sp}^*(F_{\bar s})&=&F_{\bar s},\\[2pt]
&\mathrm{sp}_*(F_{\bar\eta})&=&F_{\bar\eta}^{I}.
\end{array}
\end{equation*}
\end{enumerate}

\SGApage{II}{\(\approx\)90}
\subsubsection*{1.2.3.}
Let $S$ be a henselian trait as in 1.2.2, and let $Y$ be a scheme over
$s$.  We shall have to work with the 2-fiber product $Y\times_sS$ of
the \'{e}tale topoi of $Y$ and $S$ over the \'{e}tale topos of $s$.
By Giraud~[1], this 2-fiber product exists, and it is readily seen to
admit the Galois-theoretic description given below.  Those who dislike
2-fiber products may take this description as the definition.

\SGApage{II}{\(\approx\)91}
\subsubsection*{Construction 1.2.4.}
\underline{Sheaves $F$ on $Y\times_sS$ are identified with triples
$(F_s,F_\eta,\varphi)$ in which:}
\begin{enumerate}[label=(\alph*)]
\item $F_s$ is a sheaf on $Y$, or equivalently a sheaf $F_{\bar s}$
on $\bar Y=Y\otimes_s\bar s$, endowed with a continuous action of
$\operatorname{Gal}(\bar s/s)$ compatible with the action of
$\operatorname{Gal}(\bar s/s)$ on $\bar Y$;

\item $F_\eta$ is a sheaf $F_{\bar\eta}$ on $\bar Y$, endowed with a
continuous action of $\operatorname{Gal}(\bar\eta/\eta)$ compatible
with the action (via 1.1.1) of $\operatorname{Gal}(\bar\eta/\eta)$ on
$\bar Y$;

\item $\varphi$ is an equivariant morphism from $F_{\bar s}$ to
$F_{\bar\eta}$.
\end{enumerate}
Similarly, a \underline{sheaf $F_\eta$ on $Y\times_s\eta$} is an
object as in 1.2.4(b).  There is a diagram of topoi
\begin{equation}\tag{1.2.4.1}
\begin{tikzcd}[column sep=large, row sep=small]
Y \arrow[dd] \& {} \& Y \arrow[r, hook, "i"] \arrow[dd]
  \& Y\times_sS \arrow[dd]
  \& Y\times_s\eta \arrow[l, hook', "j"'] \arrow[dd] \& {} \\
{} \& {} \arrow[l, "\mathrm{sp}"'] \& {} \& {} \& {} \& {} \\
s \& {} \& s \arrow[r, hook, "i"] \& S
  \& \eta \arrow[l, hook', "j"'] \& {.}
\end{tikzcd}
\end{equation}
The topos $Y\times_sS$ is the union of the open subtopos
$Y\times_s\eta$ and its closed complement $Y$.  The formulas in
1.2.2(c) remain valid for
$\mathrm{sp}:Y\times_sS\longrightarrow Y$ (also denoted
$\mathrm{pr}_1$), $j:Y\times_s\eta\hookrightarrow Y\times_sS$,
$i:Y\hookrightarrow Y\times_sS$, and
$\mathrm{sp}:Y\times_s\eta\longrightarrow Y$ (also denoted
$\mathrm{pr}_1$).

\paragraph{1.2.4.2.}
If 1.2.4 is taken as the definition of $Y\times_sS$, one must show by
a direct construction that, up to unique isomorphism, $Y\times_sS$
does not depend on the choice of a separable closure $k(\bar\eta)$ of
$k(\eta)$.  This is done by interpreting the $F_\eta$ of 1.2.4(b) as
a functor which, to each separable closure $k(\bar\eta)$ of $k(\eta)$
defining a separable closure $k(\bar s)$ of $k(s)$, associates a sheaf
on $Y\otimes_{k(s)}k(\bar s)$; the $F_\eta$ of 1.2.4(b) is the
restriction of this functor to $\{\bar\eta\}$.



\SGApage{II}{\(\approx\)92}
\subsubsection*{1.2.5.}
Every geometric point $\bar x$ of $Y_{\bar s}$ defines a geometric
point $(\bar x,\bar\eta)$ of $Y\times_s\eta$, whose corresponding
fiber functor is
\begin{equation*}
F_\eta\longmapsto(F_{\bar\eta})_{\bar x}.
\end{equation*}
It defines geometric points $(\bar x,\bar\eta)$ and
$(\bar x,\bar s)$ of $Y\times_sS$ by
\begin{equation*}
F_{(\bar x,\bar\eta)}=(F_{\bar\eta})_{\bar x}
\qquad\text{and}\qquad
F_{(\bar x,\bar s)}=(F_{\bar s})_{\bar x}.
\end{equation*}
One verifies that all the points of the topoi $Y\times_s\eta$ and
$Y\times_sS$ are obtained in this way.  We shall not use this.  One
also verifies at once that these sets of points are conservative.

\subsubsection*{${}^{*}$Remark 1.2.6.}
The construction in 1.2.5 defines an equivalence from the category
$\operatorname{Point}(Y\times_s\eta)$ (respectively
$\operatorname{Point}(Y\times_sS)$) of points of $Y\times_s\eta$
(respectively of $Y\times_sS$) to the 2-fiber product
\begin{equation*}
\operatorname{Point}(Y)\times_{\operatorname{Point}(s)}
\operatorname{Point}(\eta).
\end{equation*}
This equivalence also follows from the fact that $Y\times_s\eta$
(respectively $Y\times_sS$) is a 2-fiber product of topoi.${}^{*}$

\subsubsection*{1.2.7.}
By the functoriality of 2-fiber products, diagram (1.2.4.1) is
functorial in $Y$ and $S$.  In particular:
\begin{enumerate}[label=\alph*)]
\item Every morphism $f:Y\longrightarrow Y'$ of schemes over $s$
induces a morphism of topoi from $Y\times_sS$ to $Y'\times_sS$.  One
has
\begin{equation}\tag{1.2.7.1}
f^*(F_{\bar s},F_{\bar\eta},\varphi)
=(f^*F_{\bar s},f^*F_{\bar\eta},f^*(\varphi)),
\end{equation}
and, when $f$ is quasi-compact,
\begin{equation}\tag{1.2.7.2}
f_*(F_{\bar s},F_{\bar\eta},\varphi)
=(f_*F_{\bar s},f_*F_{\bar\eta},f_*(\varphi)).
\end{equation}
Analogous formulas hold for
$f:Y\times_s\eta\longrightarrow Y'\times_s\eta$.

\item Every surjective morphism of henselian traits
$f:S'\longrightarrow S$ induces a morphism of topoi from
$Y_{s'}\times_{s'}S'$ to $Y\times_sS$.  Again writing $f$ for the
natural morphism
$Y_{s'}=Y\otimes_{k(s)}k(s')\longrightarrow Y$, one has
\begin{equation}\tag{1.2.7.3}
f^*(F_{\bar s},F_{\bar\eta},\varphi)
=(f^*F_{\bar s},f^*F_{\bar\eta},f^*(\varphi)).
\end{equation}
When $f$ is finite, $f_*$ is expressed as an induced representation.
\end{enumerate}

\SGApage{II}{\(\approx\)93}
\subsubsection*{1.2.8.}
Let $f:Y\longrightarrow Y'$ be a separated morphism, locally of finite
type, between $s$-schemes.  Define a ``direct image with proper support''
functor
\begin{equation*}
f_!: (\text{sheaves of pointed sets on }Y\times_sS)
\longrightarrow
(\text{sheaves of pointed sets on }Y'\times_sS)
\end{equation*}
by the formula
\begin{equation}\tag{1.2.8.1}
f_!(F_{\bar s},F_{\bar\eta},\varphi)
=(f_!F_{\bar s},f_!F_{\bar\eta},f_!(\varphi)).
\end{equation}
When $f$ is a locally closed immersion, $f_!$ is called the
extension-by-zero functor.

An analogous functor is defined for $Y\times_s\eta$.

\subsubsection*{1.2.9.}
Let $f:Y\hookrightarrow Y'$ be a locally closed immersion of
$s$-schemes.  The functor $f_!$ has a right adjoint $f^!$, ``sections
with support in $Y$'', satisfying
\begin{equation}\tag{1.2.9.1}
f^!(F_{\bar s},F_{\bar\eta},\varphi)
=(f^!F_{\bar s},f^!F_{\bar\eta},f^!(\varphi)).
\end{equation}
Suppose only that $f$ is quasi-finite.  The functor $f_!$ from
abelian sheaves on $Y\times_sS$ to abelian sheaves on
$Y'\times_sS$ still has a right adjoint $f^!$, given by (1.2.9.1)
(cf. SGA~4, XVIII, 3.1.8).  If $f$ is \'{e}tale, then $f^!=f^*$.

Analogous functors are defined for $Y\times_s\eta$.


\SGApage{II}{94}
\subsection*{1.3. The functor $\Psi$}

\subsubsection*{1.3.1.}
Let $S$ be a henselian trait as in 0.2.5, and let
$p:X\longrightarrow S$.  We propose to define a left exact functor
$\Psi$ from sheaves on $X$ to sheaves on the topos $X_s\times_sS$
constructed in 1.2.  In general, this functor is not the direct-image
functor of a morphism of topoi.

We shall also define
\begin{equation*}
\Psi_\eta:(\text{sheaves on }X_\eta)
\longrightarrow(\text{sheaves on }X_s\times_s\eta).
\end{equation*}

\subsubsection*{1.3.2.}
Put $\bar X=X\times_S\bar S$.  There is a Cartesian diagram
\begin{equation}\tag{1.3.2.1}
\begin{tikzcd}
X_{\bar s} \arrow[r, hook, "\bar i"] \arrow[d]
  \& \bar X \arrow[d]
  \& X_{\bar\eta} \arrow[l, hook', "\bar j"'] \arrow[d] \\
X_s \arrow[r, hook] \& X \& X_\eta \arrow[l, hook']
\end{tikzcd}
\end{equation}
over the diagram
\begin{center}
\begin{tikzcd}
\bar s \arrow[r, hook] \arrow[d]
  \& \bar S \arrow[d]
  \& \bar\eta \arrow[l, hook'] \arrow[d] \\
s \arrow[r, hook] \& S \& \eta \arrow[l, hook']
\end{tikzcd}
.\end{center}
Let $F$ be a sheaf on $X_\eta$, with inverse image $F_{\bar\eta}$
on $X_{\bar\eta}$.  Put
\begin{equation}\tag{1.3.2.2}
\Psi_\eta(F)=\bar i^*\bar j_*F_{\bar\eta}.
\end{equation}
By transport of structure, this sheaf is endowed with a continuous
action of $\operatorname{Gal}(\bar\eta/\eta)$, compatible with the
action of $\operatorname{Gal}(\bar\eta/\eta)$ on $X_{\bar s}$.  The
functor
\begin{equation}\tag{1.3.2.3}
\Psi_\eta:(X_{\eta,\mathrm{et}})^{\sim}
\longrightarrow(X_s\times_s\eta)^{\sim}
\end{equation}
is left exact.

\subsubsection*{1.3.3.}
Let $F$ be a sheaf on $X$, and let $F_\eta$ and $F_s$ be its
restrictions to $X_\eta$ and $X_s$.  Put
\begin{equation}\tag{1.3.3.1}
(\Psi(F))_\eta=\Psi_\eta(F_\eta)
\end{equation}
and
\begin{equation}\tag{1.3.3.2}
(\Psi(F))_s=F_s.
\end{equation}
Let $\bar F$ be the inverse image of $F$ on $\bar X$, and let
$\varphi'$ be the adjunction morphism
\begin{equation*}
\varphi':\bar F\longrightarrow\bar j_*\bar j^*\bar F.
\end{equation*}
This morphism induces
\begin{equation*}
\varphi:F_{\bar s}=\bar i^*\bar F
\longrightarrow
\bar i^*\bar j_*\bar j^*\bar F
=(\Psi_\eta(F_\eta))_{\bar\eta}.
\end{equation*}
The triple
\begin{equation*}
\Psi(F)=((\Psi(F))_s,(\Psi(F))_\eta,\varphi)
\end{equation*}
is a sheaf on $X_s\times_sS$.  The functor
\begin{equation}\tag{1.3.3.3}
\Psi:(X_{\mathrm{et}})^{\sim}
\longrightarrow(X_s\times_sS)^{\sim}
\end{equation}
is left exact.

\SGApage{II}{\(\approx\)95}
\subsubsection*{1.3.4.}
When there is no risk of confusion, we shall write
\begin{equation*}
\Psi_\eta(F)=(\Psi(F))_\eta,
\qquad
\Psi_{\bar\eta}(F)=(\Psi(F))_{\bar\eta},
\qquad\text{and}\qquad
\Psi_{\bar\eta}(F)=(\Psi_\eta(F))_{\bar\eta}.
\end{equation*}

\subsubsection*{1.3.5.}
For the remainder of this section, we list functoriality properties of
$\Psi$ (cf. 1.2.7--1.2.9).  They will be useful chiefly in their
derived form (2.1).

\SGAvar{Let $f:X\longrightarrow X'$ be a morphism of $S$-schemes, giving a
diagram
\begin{equation}\tag{1.3.5.1}
\begin{tikzcd}
X_{\bar s} \arrow[r, hook, "\bar i"] \arrow[d, "f"]
  \&  X \arrow[d, "f"]
  \& X_{\bar\eta} \arrow[l, hook', "\bar j"'] \arrow[d, "f"] \\
X'_{\bar s} \arrow[r, hook, "\bar i"]
  \& \bar X' \& X'_{\bar\eta} \arrow[l, hook', "\bar j"']
\end{tikzcd}
\end{equation}
  The ambiguity in the notation for the arrows simplifies
the formulas and causes no confusion: only one interpretation makes
sense of the composites below.  Likewise, we write $\Psi$ for the
functor (1.3.3.3) for both $X$ and $X'$.

}{Let $f:X\longrightarrow X'$ be a morphism of $S$-schemes, giving a
diagram
\begin{equation}\tag{1.3.5.1}
\begin{tikzcd}
X_{\bar s} \arrow[r, hook, "\bar i"] \arrow[d, "f"]
  \& \bar X \arrow[d, "f"]
  \& X_{\bar\eta} \arrow[l, hook', "\bar j"'] \arrow[d, "f"] \\
X'_{\bar s} \arrow[r, hook, "\bar i"]
  \& \bar X' \& X'_{\bar\eta} \arrow[l, hook', "\bar j"']
\end{tikzcd}
\end{equation}
  The ambiguity in the notation for the arrows simplifies
the formulas and causes no confusion: only one interpretation makes
sense of the composites below.  Likewise, we write $\Psi$ for the
functor (1.3.3.3) for both $X$ and $X'$.

}{E7II-XIII-002}The morphisms of functors described below for $\Psi$ have variants for
$\Psi_\eta$, which we generally omit to state explicitly.

\SGApage{II}{\(\approx\)96}
\subsubsection*{1.3.6.}
The base-change morphism
$\bar i^*f_*\longrightarrow f_*\bar i^*$ induces a morphism of
functors
\begin{equation}\tag{1.3.6.1}
\Psi f_*\longrightarrow f_*\Psi,
\end{equation}
which is an isomorphism when $f$ is proper (SGA~4, XII, 5.1(i)).

The essential case is $X'=S$.  One obtains a morphism
\begin{equation}\tag{1.3.6.2}
f_*\longrightarrow f_*\Psi
\end{equation}
whose essential component is
\begin{equation}\tag{1.3.6.3}
\Gamma(X_{\bar\eta},F)
\longrightarrow
\Gamma(X_{\bar s},\Psi_{\bar\eta}(F));
\end{equation}
this morphism is an isomorphism when $f$ is proper.

\subsubsection*{1.3.7.}
The base-change morphism
$f^*\bar j_*\longrightarrow\bar j_*f^*$ induces a morphism of
functors
\begin{equation}\tag{1.3.7.1}
f^*\Psi\longrightarrow\Psi f^*.
\end{equation}

\subsubsection*{1.3.8.}
Suppose that $f$ is quasi-finite.  The functors $f_!$ commute with base
change; the morphism
$f_!\bar j_*\longrightarrow\bar j_*f_!$ induces a morphism of
functors
\begin{equation}\tag{1.3.8.1}
f_!\Psi\longrightarrow\Psi f_!.
\end{equation}
When $f$ is finite, this is the inverse of (1.3.6.1).

\subsubsection*{1.3.9.}
The situation is dual for $f^!$ (we suppose that $f$ is quasi-finite
when working with abelian sheaves, or that it is an immersion when
working with sheaves of pointed sets).  This functor commutes with
direct images, and the morphism
$\bar i^*f^!\longrightarrow f^!\bar i^*$ induces a morphism of
functors
\begin{equation}\tag{1.3.9.1}
\Psi f^!\longrightarrow f^!\Psi.
\end{equation}
When $f$ is \'{e}tale, this is the inverse of (1.3.7.1).

\subsubsection*{1.3.10.}
Finally, a base change $S'\xrightarrow{f}S$ gives a base-change
morphism
\begin{equation}\tag{1.3.10.1}
f^*\Psi\longrightarrow\Psi f^*.
\end{equation}



\SGApage{II}{97}
\subsection*{1.4. The variation}

\subsubsection*{1.4.1.}
Let $S$ be a henselian trait as in 0.2.5, and let $Y$ be a scheme over
$s$.  Let $\Lambda$ be a ring.  More generally, one could take for
$\Lambda$ a sheaf of rings on $Y$.  For every object $K$ of the
derived category $D(Y\times_sS,\Lambda)$ of the category of sheaves of
$\Lambda$-modules on $Y\times_sS$, we propose to define a
\underline{variation} morphism.

\SGAvar{\subsubsection*{1.4.2.}
A complex of sheaves of $\Lambda$-modules $K$ on $X\times_sS$ is
identified with a triple $(K_s,K_\eta,\varphi)$ in which
\begin{enumerate}[label=\alph*)]
\item $K_s$ (respectively $K_\eta$) is a complex of sheaves
$K_{\bar s}$ (respectively $K_{\bar\eta}$) on $\bar Y$, endowed with a
continuous action of $\operatorname{Gal}(\bar s/\bar s)$ (respectively
$\operatorname{Gal}(\bar\eta/\eta)$);

\item $\varphi$ is an equivariant morphism
$\varphi:K_{\bar s}\longrightarrow K_{\bar\eta}$.
\end{enumerate}


}{\subsubsection*{1.4.2.}
A complex of sheaves of $\Lambda$-modules $K$ on $Y\times_sS$ is
identified with a triple $(K_s,K_\eta,\varphi)$ in which
\begin{enumerate}[label=\alph*)]
\item $K_s$ (respectively $K_\eta$) is a complex of sheaves
$K_{\bar s}$ (respectively $K_{\bar\eta}$) on $\bar Y$, endowed with a
continuous action of $\operatorname{Gal}(\bar s/s)$ (respectively
$\operatorname{Gal}(\bar\eta/\eta)$);

\item $\varphi$ is an equivariant morphism
$\varphi:K_{\bar s}\longrightarrow K_{\bar\eta}$.
\end{enumerate}


}{E7II-XIII-003}Such a complex $K$ is always homotopic to a complex
$K'=(K'_s,K'_\eta,\varphi')$ such that $\varphi'$ is injective and the
exact sequence
\begin{equation}\tag{1.4.2.1}
0\longrightarrow K'_{\bar s}
\longrightarrow K'_{\bar\eta}
\longrightarrow\operatorname{coker}(\varphi')
\longrightarrow0
\end{equation}
is split degree by degree.  It is enough to take $K'_s=K_s$ and to
take for $K'_\eta$ the direct sum of $K_\eta$ and the cone of
$\mathrm{sp}^*K_s$.  Put
\begin{equation*}
\Phi(K)=\operatorname{coker}(\varphi').
\end{equation*}
The exact sequence (1.4.2.1) defines a distinguished triangle
\begin{equation}\tag{1.4.2.2}
\begin{tikzcd}[column sep=large, row sep=large]
{} \& \Phi(K) \arrow[dl] \& {} \\
\mathrm{sp}^*K_s \arrow[rr] \& {} \& K_\eta \arrow[ul]
\end{tikzcd}
\end{equation}
in the category $K(Y\times_s\eta,\Lambda)$ of complexes of sheaves of
$\Lambda$-modules on $Y\times_s\eta$, up to homotopy.  This triangle
depends only on $K\in\operatorname{Ob}K(Y\times_sS,\Lambda)$.

\SGApage{II}{\(\approx\)98}
\SGAvar{The construction passes to the derived category and associates with
every object $K$ of $D(X\times_sS,\Lambda)$ a distinguished triangle
of $D(X\times_s\eta,\Lambda)$.

}{The construction passes to the derived category and associates with
every object $K$ of $D(Y\times_sS,\Lambda)$ a distinguished triangle
of $D(Y\times_s\eta,\Lambda)$.

}{E7II-XIII-004}Triangle (1.4.2.2) is useful chiefly through the long exact
hypercohomology sequence that it induces:
\begin{equation*}
\cdots\longrightarrow H^i(\bar Y,K_{\bar s})
\longrightarrow H^i(\bar Y,K_{\bar\eta})
\longrightarrow H^i(\bar Y,\Phi(K)_{\bar\eta})
\longrightarrow\cdots.
\end{equation*}

\subsubsection*{1.4.3.}
The inertia group $I$ acts trivially on $K'_{\bar s}$.  Hence, for
every $\sigma\in I$, the endomorphism $\sigma-1$ of
$K'_{\bar\eta}$ (notation (1.4.2.1)) factors through a morphism
\begin{equation*}
\operatorname{Var}(\sigma):
\operatorname{coker}(\varphi')\longrightarrow K'_{\bar\eta}.
\end{equation*}
This construction passes to the derived category: for
$K\in\operatorname{Ob}D(Y\times_sS,\Lambda)$ and $\sigma\in I$, it
defines the \underline{variation}
\begin{equation}\tag{1.4.3.1}
\operatorname{Var}(\sigma):
\Phi(K)_{\bar\eta}\longrightarrow K_{\bar\eta}.
\end{equation}
Writing $q$ for the arrow from $K_\eta$ to $\Phi(K)$, one has
\begin{align}\tag{1.4.3.2}
\sigma&=1+\operatorname{Var}(\sigma)q
&&\text{on }K_{\bar\eta},\\
\tag{1.4.3.3}
\sigma&=1+q\operatorname{Var}(\sigma)
&&\text{on }\Phi(K)_{\bar\eta},
\end{align}
and
\begin{equation}\tag{1.4.3.4}
\operatorname{Var}(\sigma\tau)
=\operatorname{Var}(\sigma)q\operatorname{Var}(\tau)
 +\operatorname{Var}(\sigma)+\operatorname{Var}(\tau).
\end{equation}



\section*{2. Vanishing cycles}

\subsection*{2.1. The functor $R\Psi$}

\subsubsection*{2.1.1.}
Let $S$ be a henselian trait as in 0.2.5, and let $\Lambda$ be a
torsion ring whose torsion is prime to the residual characteristic
$p$ of $S$ (0.2.7).  More generally, one could take for $\Lambda$ a
torsion sheaf whose torsion is prime to the residual characteristics;
the most important case is
$\Lambda=\mathbb Z/\ell^n$ with $\ell$ an invertible prime in
$\mathcal O_S$.

\SGApage{II}{\(\approx\)99}
Let $X$ be a scheme over $S$.  The functor $\Psi$ (respectively
$\Psi_\eta$) carries sheaves of $\Lambda$-modules on $X$
(respectively on $X_\eta$) to sheaves of $\Lambda$-modules on
$X_s\times_sS$ (respectively on $X_s\times_s\eta$).  Let $R\Psi$ and
$R\Psi_\eta$ be the derived functors
\begin{equation*}
R\Psi:D^+(X,\Lambda)\longrightarrow D^+(X_s\times_sS,\Lambda),
\end{equation*}
\begin{equation*}
R\Psi_\eta:D^+(X_\eta,\Lambda)
\longrightarrow D^+(X_s\times_s\eta,\Lambda),
\end{equation*}
and let $R\Phi$ be the composite functor $\Phi\circ R\Psi$:
\begin{equation*}
R\Phi:D^+(X,\Lambda)
\longrightarrow D^+(X_s\times_s\eta,\Lambda).
\end{equation*}

\SGAvar{\subsubsection*{2.1.2.}
Consider the commutative diagram (1.3.2.1)
\begin{equation}\notag
\begin{tikzcd}
X_{\bar s} \arrow[r, hook, "\bar i"] \arrow[d]
  \& \bar X \arrow[d]
  \& X_{\bar\eta} \arrow[l, hook', "\bar j"'] \arrow[d] \\
X_s \arrow[r, hook, "i"] \& X \& X_\eta \arrow[l, hook', "j"']
\end{tikzcd}
\end{equation}
If $F$ is an injective sheaf of modules on $X$, its restriction to the
open subscheme $X_\eta$ is still injective.  If $F_\eta$ is injective
on $X_\eta$, its inverse image $F_{\bar\eta}$ is acyclic.  Thus the
derived forms of (1.3.3.1) and (1.3.3.2) are
\begin{equation}\tag{2.1.2.1}
R\Psi(K)_s=i^*K
\qquad(K\in\operatorname{Ob}D^+(X,\Lambda)),
\end{equation}
\begin{equation}\tag{2.1.2.2}
R\Psi(K)_\eta=R\Psi_\eta(K_\eta)
\qquad(K\in\operatorname{Ob}D^+(X,\Lambda)),
\end{equation}
and
\begin{equation}\tag{2.1.2.3}
R\Psi_\eta(K)_{\bar\eta}
=\bar i^*R\bar j_*K_{\bar\eta}
\qquad(K\in\operatorname{Ob}D^+(X_\eta,\Lambda)).
\end{equation}
This allows us to write, without conflict,
\begin{equation*}
R\Psi_\eta(K)=(R\Psi(K))_\eta,
\qquad
R\Psi_{\bar\eta}(K)=(R\Psi(K))_{\bar\eta},
\qquad\text{and}\qquad
R\Psi_{\bar\eta}(K)=(R\Psi_\eta(K))_{\bar\eta}.
\end{equation*}
With this notation, the distinguished triangle (1.4.2.2) becomes
\begin{equation}\tag{2.1.2.4}
\begin{tikzcd}[row sep=large, column sep=huge]
{} \& R\Phi(K) \arrow[dl] \& {} \\
\mathrm{sp}^*i^*K \arrow[rr] \& {}
  \& R\Psi_\eta(K_s) \arrow[ul]
\end{tikzcd}
.\end{equation}


}{\subsubsection*{2.1.2.}
Consider the commutative diagram (1.3.2.1)
\begin{equation}\notag
\begin{tikzcd}
X_{\bar s} \arrow[r, hook, "\bar i"] \arrow[d]
  \& \bar X \arrow[d]
  \& X_{\bar\eta} \arrow[l, hook', "\bar j"'] \arrow[d] \\
X_s \arrow[r, hook, "i"] \& X \& X_\eta \arrow[l, hook', "j"']
\end{tikzcd}
\end{equation}
If $F$ is an injective sheaf of modules on $X$, its restriction to the
open subscheme $X_\eta$ is still injective.  If $F_\eta$ is injective
on $X_\eta$, its inverse image $F_{\bar\eta}$ is acyclic.  Thus the
derived forms of (1.3.3.1) and (1.3.3.2) are
\begin{equation}\tag{2.1.2.1}
R\Psi(K)_s=i^*K
\qquad(K\in\operatorname{Ob}D^+(X,\Lambda)),
\end{equation}
\begin{equation}\tag{2.1.2.2}
R\Psi(K)_\eta=R\Psi_\eta(K_\eta)
\qquad(K\in\operatorname{Ob}D^+(X,\Lambda)),
\end{equation}
and
\begin{equation}\tag{2.1.2.3}
R\Psi_\eta(K)_{\bar\eta}
=\bar i^*R\bar j_*K_{\bar\eta}
\qquad(K\in\operatorname{Ob}D^+(X_\eta,\Lambda)).
\end{equation}
This allows us to write, without conflict,
\begin{equation*}
R\Psi_\eta(K)=(R\Psi(K))_\eta,
\qquad
R\Psi_{\bar\eta}(K)=(R\Psi(K))_{\bar\eta},
\qquad\text{and}\qquad
R\Psi_{\bar\eta}(K)=(R\Psi_\eta(K))_{\bar\eta}.
\end{equation*}
With this notation, the distinguished triangle (1.4.2.2) becomes
\begin{equation}\tag{2.1.2.4}
\begin{tikzcd}[row sep=large, column sep=huge]
{} \& R\Phi(K) \arrow[dl] \& {} \\
\mathrm{sp}^*i^*K \arrow[rr] \& {}
  \& R\Psi_\eta(K_\eta) \arrow[ul]
\end{tikzcd}
.\end{equation}


}{E7II-XIII-005}\SGApage{II}{\(\approx\)100}
For $\sigma\in I$, there is a \underline{variation} morphism
\begin{equation}\tag{2.1.2.5}
\operatorname{Var}(\sigma):R\Phi(K)_{\bar\eta}
\longrightarrow R\Psi_{\bar\eta}(K).
\end{equation}
The sheaves $R^i\Phi(\Lambda)$, or more generally $R^i\Phi(K)$, are
the \underline{sheaves of vanishing cycles}.

\SGAvar{\subsubsection*{2.1.3.}
Let $(\bar x,\bar\eta)$ be the geometric point of
$X_s\times_s\eta$ defined by a geometric point $\bar x$ of
$X_{\bar s}$ (1.2.5).  We compute the fiber of
$R\Psi_\eta(K_\eta)$ at $(\bar x,\bar\eta)$ using (2.1.2.3)
(cf. I, 2.3).  The strict henselization $X_{(\bar x)}$ of $X$ at
$\bar x$ is a scheme over the strict henselization $S_{\mathrm{nr}}$
of $S$, with geometric point $\eta_{\mathrm{nr}}$.

\subsubsection*{Proposition 2.1.4.}
\underline{One has}
\begin{equation*}
(R\Psi(K_\eta))_{(\bar x,\bar\eta)}
=R\Gamma\!\left(
X_{(\bar x)}\times_{\eta_{\mathrm{xr}}}\bar\eta,K
\right);
\end{equation*}
\underline{in particular,}
\begin{equation*}
R^i\Psi(K_\eta)_{(\bar x,\bar\eta)}
=\mathbb H^i\!\left(
X_{(\bar x)}\times_{\eta_{\mathrm{xr}}}\bar\eta,K
\right).
\end{equation*}

}{\subsubsection*{2.1.3.}
Let $(\bar x,\bar\eta)$ be the geometric point of
$X_s\times_s\eta$ defined by a geometric point $\bar x$ of
$X_{\bar s}$ (1.2.5).  We compute the fiber of
$R\Psi_\eta(K_\eta)$ at $(\bar x,\bar\eta)$ using (2.1.2.3)
(cf. I, 2.3).  The strict henselization $X_{(\bar x)}$ of $X$ at
$\bar x$ is a scheme over the strict henselization $S_{\mathrm{nr}}$
of $S$, with generic point $\eta_{\mathrm{nr}}$.

\subsubsection*{Proposition 2.1.4.}
\underline{One has}
\begin{equation*}
(R\Psi(K_\eta))_{(\bar x,\bar\eta)}
=R\Gamma\!\left(
X_{(\bar x)}\times_{\eta_{\mathrm{nr}}}\bar\eta,K
\right);
\end{equation*}
\underline{in particular,}
\begin{equation*}
R^i\Psi(K_\eta)_{(\bar x,\bar\eta)}
=\mathbb H^i\!\left(
X_{(\bar x)}\times_{\eta_{\mathrm{nr}}}\bar\eta,K
\right).
\end{equation*}

}{E7II-XIII-006}Apply SGA~4, XV, 2.1 (the local-acyclicity theorem for smooth
morphisms).  This gives the following restatement.

\subsubsection*{Restatement 2.1.5.}
\underline{If $f$ is smooth and $F$ is a locally constant sheaf, then}
\begin{equation*}
R\Phi(F)=0.
\end{equation*}
More generally, if $K\in\operatorname{Ob}D^+(X,\Lambda)$, then
$R\Phi(K)=0$ wherever $f$ is smooth and the sheaves
$\mathcal H^i(K)$ are locally constant.  See also 2.1.11.

\subsubsection*{2.1.6.}
Let $f:Y\longrightarrow Y'$ be a morphism of schemes over $s$.
The functors $f^*$, $f_*$, $f^!$, and $f_!$ (1.2.7--1.2.9) have
natural analogues in the derived category.
\begin{enumerate}[label=\alph*)]
\item The functor $f^*$ is exact and therefore derives trivially.

\item The functor $f_*$ has a right derived functor $Rf_*$.  When $f$
is quasi-compact, one has
\begin{equation}\tag{2.1.6.1}
(Rf_*K)_{\bar s}=Rf_{\bar s*}(K_{\bar s})
\end{equation}
and
\begin{equation}\tag{2.1.6.2}
(Rf_*K)_{\bar\eta}=Rf_{\bar s*}(K_{\bar\eta}).
\end{equation}

\item When $f$ is quasi-finite, the functor $f_!$ is exact and
therefore derives trivially.  In general, the derived functor of $f_!$
is pathological, and one proceeds as in SGA~4, XVII.  Suppose that $f$
is separated and of finite type and that $Y'$ is noetherian.  There is
then [2] a factorization $f=\widetilde f\,j$, where $j$ is an open
immersion and $\widetilde f$ is proper, and one defines
\begin{equation*}
Rf_!=R\widetilde f_*\circ j_!.
\end{equation*}
The arguments of SGA~4, XVII show that $Rf_!$ is independent of the
choice of the factorization $\widetilde f\,j$.  The analogues of
(2.1.6.1) and (2.1.6.2) hold.

\item We shall use $Rf^!$ only for a quasi-finite morphism, in which
case it is the derived functor of $f^!$.  Again, the analogues of
(2.1.6.1) and (2.1.6.2) hold.
\end{enumerate}



\SGApage{II}{\(\approx\)101}
\subsubsection*{2.1.7.}
We now review the functoriality properties of $R\Psi$, that is, we
derive 1.3.5--1.3.10.  Let $f:X\longrightarrow X'$ be a morphism of
$S$-schemes, and retain the notation of 1.3.5.

The morphism (1.3.6.1) defines
\begin{equation}\tag{2.1.7.1}
R\Psi\,Rf_*\longrightarrow Rf_*\,R\Psi.
\end{equation}
By the proper base-change theorem (SGA~4, XII, 5.1),
\underline{this morphism is an isomorphism if $f$ is proper.}

The morphism (1.3.7.1) defines
\begin{equation}\tag{2.1.7.2}
f^*R\Psi\longrightarrow R\Psi f^*.
\end{equation}
By the smooth base-change theorem (SGA~4, XVI, 1.2),
\underline{this morphism is an isomorphism if $f$ is smooth.}

\SGAvar{The isomorphism (2.1.7.1) for $f$ proper and the morphism (1.3.8.1)
define a morphism
\begin{equation}\tag{2.1.7.3}
Rf_!R\Psi\longrightarrow R\Psi Rf_!,
\end{equation}
which, when $f$ is proper, is the inverse of (2.1.8.1).  

When $f$ is quasi-finite, (1.3.9.1) derives to
\begin{equation}\tag{2.1.7.4}
R\Psi Rf^!\longrightarrow Rf^!R\Psi,
\end{equation}
which, when $f$ is \'{e}tale, is the inverse of (2.1.7.1).  

}{The isomorphism (2.1.7.1) for $f$ proper and the morphism (1.3.8.1)
define a morphism
\begin{equation}\tag{2.1.7.3}
Rf_!R\Psi\longrightarrow R\Psi Rf_!,
\end{equation}
which, when $f$ is proper, is the inverse of (2.1.7.1).  

When $f$ is quasi-finite, (1.3.9.1) derives to
\begin{equation}\tag{2.1.7.4}
R\Psi Rf^!\longrightarrow Rf^!R\Psi,
\end{equation}
which, when $f$ is \'{e}tale, is the inverse of (2.1.7.2).  

}{E7II-XIII-007}\SGApage{II}{\(\approx\)102}
Let $f:S'\longrightarrow S$ be a change of trait, and put
$X'=X\times_SS'$:
\begin{center}
\begin{tikzcd}
X' \arrow[r, "f"] \arrow[d, "p"'] \& X \arrow[d, "p"] \\
S' \arrow[r, "f"] \& S
\end{tikzcd}
\end{center}
(for the ambiguity in the notation, see 1.3.5).  Suppose that $f$ is
surjective.  The base-change arrow (1.3.10.1) derives to
\begin{equation}\tag{2.1.7.5}
f^*R\Psi\longrightarrow R\Psi f^*.
\end{equation}
The same statements hold for $R\Psi_\eta$.

\subsubsection*{2.1.8.}
The essential case of (2.1.7.1) and (2.1.7.3) is $X'=S$.  For
$K\in\operatorname{Ob}D^+(X_{\bar\eta},\Lambda)$ one then obtains
morphisms
\begin{equation}\tag{2.1.8.1}
Rf_*\longrightarrow Rf_*R\Psi
\end{equation}
and
\begin{equation}\tag{2.1.8.2}
Rf_!R\Psi\longrightarrow Rf_!,
\end{equation}
hence
\begin{equation}\tag{2.1.8.3}
\mathbb H^i(X_{\bar\eta},K)
\longrightarrow
\mathbb H^i(X_{\bar s},R\Psi_{\bar\eta}(K))
\end{equation}
and
\begin{equation}\tag{2.1.8.4}
\mathbb H_c^i(X_{\bar s},R\Psi_{\bar\eta}(K))
\longrightarrow
\mathbb H_c^i(X_{\bar\eta},K),
\end{equation}
and the commutativity of the diagrams
\begin{equation}\tag{2.1.8.5}
\begin{tikzcd}[column sep=large]
(R^if_*K)_{\bar s} \arrow[r, "\mathrm{sp}"] \arrow[d]
  \& \mathbb H^i(X_{\bar\eta},K) \arrow[d] \\
\mathbb H^i(X_{\bar s},K) \arrow[r]
  \& \mathbb H^i(X_{\bar s},R\Psi_{\bar\eta}(K))
\end{tikzcd}
\end{equation}
and
\begin{equation}\tag{2.1.8.6}
\begin{tikzcd}[column sep=large]
\mathbb H_c^i(X_{\bar s},K) \arrow[rr, "\mathrm{sp}"] \arrow[dr]
  \& {} \& \mathbb H_c^i(X_{\bar\eta},K) \\
{} \& \mathbb H_c^i(X_{\bar s},R\Psi_{\bar\eta}(K)) \arrow[ur] \& {}
\end{tikzcd}
.\end{equation}

\SGApage{II}{\(\approx\)103}
\SGAvar{For $K\in\operatorname{Ob}D^+(X_\eta,\Lambda)$ and $\sigma\in I$,
the following diagrams are commutative:
\begin{equation}\tag{2.1.8.7}
\begin{tikzcd}[column sep=small, cells={nodes={font=\small}}]
\mathbb H^i(X_{\bar\eta},K)
  \arrow[rr, dash, "\sigma-1"] \arrow[d]
  \& {} \& \mathbb H^i(X_{\bar\eta},K) \arrow[d] \\
\mathbb H^i(X_{\bar s},R\Psi_{\bar\eta}(K)) \arrow[r]
  \& \mathbb H^i(X_{\bar s},R\Phi_{\bar\eta}(K))
      \arrow[r, "\operatorname{Var}(\sigma)"]
  \& \mathbb H^i(X_{\bar s},R\Psi_{\bar\eta}(K))
\end{tikzcd}
\end{equation}
and
\begin{equation}\tag{2.1.8.8}
\begin{tikzcd}[column sep=small, cells={nodes={font=\small}}]
\mathbb H_c^i(X_{\bar\eta},K)
  \arrow[rr, "\sigma-1"]
  \& {} \& \mathbb H_c^i(X_{\bar\eta},K) \\
\mathbb H_c^i(X_s,R\Psi_{\bar\eta}(K)) \arrow[r] \arrow[u]
  \& \mathbb H_c^i(X_{\bar s},R\Phi_{\bar\eta}(K)
      \arrow[r, "\operatorname{Var}(\sigma)"]
  \& \mathbb H_c^i(X_{\bar s},R\Psi_{\bar\eta}(K)) \arrow[u]
\end{tikzcd}
.\end{equation}


}{For $K\in\operatorname{Ob}D^+(X_\eta,\Lambda)$ and $\sigma\in I$,
the following diagrams are commutative:
\begin{equation}\tag{2.1.8.7}
\begin{tikzcd}[column sep=small, cells={nodes={font=\small}}]
\mathbb H^i(X_{\bar\eta},K)
  \arrow[rr, "\sigma-1"] \arrow[d]
  \& {} \& \mathbb H^i(X_{\bar\eta},K) \arrow[d] \\
\mathbb H^i(X_{\bar s},R\Psi_{\bar\eta}(K)) \arrow[r]
  \& \mathbb H^i(X_{\bar s},R\Phi_{\bar\eta}(K))
      \arrow[r, "\operatorname{Var}(\sigma)"]
  \& \mathbb H^i(X_{\bar s},R\Psi_{\bar\eta}(K))
\end{tikzcd}
\end{equation}
and
\begin{equation}\tag{2.1.8.8}
\begin{tikzcd}[column sep=small, cells={nodes={font=\small}}]
\mathbb H_c^i(X_{\bar\eta},K)
  \arrow[rr, "\sigma-1"]
  \& {} \& \mathbb H_c^i(X_{\bar\eta},K) \\
\mathbb H_c^i(X_{\bar s},R\Psi_{\bar\eta}(K)) \arrow[r] \arrow[u]
  \& \mathbb H_c^i(X_{\bar s},R\Phi_{\bar\eta}(K))
      \arrow[r, "\operatorname{Var}(\sigma)"]
  \& \mathbb H_c^i(X_{\bar s},R\Psi_{\bar\eta}(K)) \arrow[u]
\end{tikzcd}
.\end{equation}


}{E7II-XIII-008}When $f$ is proper, (2.1.8.1) and (2.1.8.2), as well as (2.1.8.3) and
(2.1.8.4), are mutually inverse isomorphisms.  The distinguished
triangle (2.1.2.4) gives a long exact sequence
\begin{equation}\tag{2.1.8.9}
\cdots\longrightarrow\mathbb H^i(X_{\bar s},K)
\longrightarrow\mathbb H^i(X_{\bar\eta},K)
\longrightarrow\mathbb H^i(X_{\bar s},R\Phi_{\bar\eta}(K))
\longrightarrow\cdots.
\end{equation}
Thus the vanishing complex $R\Phi(K)$ expresses the difference between
the cohomologies of the geometric fibers $X_{\bar s}$ and
$X_{\bar\eta}$, with coefficients in the inverse image of $K$.
By (2.1.8.7), the variation $\operatorname{Var}(\sigma)$ determines
the action of the inertia group $I$ on
$\mathbb H^i(X_{\bar\eta},K)$.  These facts underlie all applications
of the method of vanishing cycles.



\SGApage{II}{\(\approx\)104}
The properness hypothesis imposed above on $f$ may be replaced by
regularity-at-infinity conditions on $f$ and $K$.  For example, one has
the following result (for
$K\in\operatorname{Ob}D(X_{\bar\eta},\Lambda)$).

\subsubsection*{Proposition 2.1.9.}
\underline{Suppose that $f$ admits a factorization $f=f_1k$}
\begin{center}
\begin{tikzcd}[column sep=huge]
X \arrow[r, hook, "k"] \& X_1 \arrow[r, "f_1"] \& S,
\end{tikzcd}
\end{center}
\underline{where $f_1$ is proper and $k$ identifies $X$ with the
complement in $X_1$ of a relative normal-crossings divisor $D$}
(thus, by hypothesis, $f_1$ is smooth on a neighborhood of $D$).
\underline{If, on a neighborhood $U$ of $D$, the sheaves
$\mathcal H^i(K)$ are locally constant on $U-(D\cup X_0)$ and tamely
ramified along the normal-crossings divisor $D\cup X_0$, then
(2.1.8.1) and (2.1.8.2) are isomorphisms.}

\SGAvar{The hypercohomology spectral sequence for $K$ shows that it is enough
to treat the case in which $K$ consists of a sheaf $F_\eta$ in degree
zero.  After replacing $S$ by its normalization in a finite extension
(which is permissible, as one sees directly from the definitions) and
applying Abhyankhar's lemma, we may suppose that $F_\eta$ is the
restriction to $X_\eta$ of a sheaf $F$ on $X$, locally constant on
$U\cap X$ and tamely ramified along $D$.  With the notation of 2.1.2,
let $\overline F$ be the inverse image of $F$ on $\overline X$, and
write $\bar f:\overline X\longrightarrow\overline S$.

}{The hypercohomology spectral sequence for $K$ shows that it is enough
to treat the case in which $K$ consists of a sheaf $F_\eta$ in degree
zero.  After replacing $S$ by its normalization in a finite extension
(which is permissible, as one sees directly from the definitions) and
applying Abhyankar's lemma, we may suppose that $F_\eta$ is the
restriction to $X_\eta$ of a sheaf $F$ on $X$, locally constant on
$U\cap X$ and tamely ramified along $D$.  With the notation of 2.1.2,
let $\overline F$ be the inverse image of $F$ on $\overline X$, and
write $\bar f:\overline X\longrightarrow\overline S$.

}{E7II-XIII-009}To prove that (2.1.8.1) is an isomorphism, it is enough to prove that
\[
R\bar f_*\bigl(R\bar j_*\bar j^*\overline F\bigr)
\]
commutes with base change by $\bar s\longrightarrow\overline S$.  In
the distinguished triangle
\[
\longrightarrow\overline F\longrightarrow
R\bar j_*\bar j^*\overline F\longrightarrow\Delta\longrightarrow,
\]
the sheaves $\mathcal H^i(\Delta)$ are supported on $X_{\bar s}$; for
$\Delta$, the base-change statement in question is tautological.  It
therefore remains to prove that $R\bar f_*(\overline F)$ commutes with
base change by $\bar s\longrightarrow\overline S$, or, in view of the
base-change theorem for the proper morphism $f_1$, that
$Rk_*(\overline F)$ commutes with base change
$X_{1\bar s}\longrightarrow X_1$.  This is Lemma 2.1.10 below.

\SGApage{II}{\(\approx\)105}
To prove that (2.1.8.2) is an isomorphism, it is enough to prove, for
the sheaf $F$, that extension by zero $k_!$ commutes with the functor
$R\bar j_*$; that is, on a neighborhood of $D$,
\[
R\Phi(k_!F)=0\qquad\hbox{on }X_{1\bar s}.
\]
This is Lemma 2.1.11 below.

\subsubsection*{Lemma 2.1.10.}
\underline{Let} $f:X\longrightarrow S$ \underline{be smooth, let}
$D\subset X$ \underline{be a relative normal-crossings divisor, put}
$U=X-D\mathrel{\mathop{\hookrightarrow}^{j}}X$, \underline{and let}
$F$ \underline{be a sheaf of $\Lambda$-modules on $U$ which is locally
constant and tamely ramified along} $D$.  \underline{Formation of}
$Rj_*F$ \underline{commutes with arbitrary base change}
$T\longrightarrow S$.

The question is local on $X$ for the \'{e}tale topology.  We may
therefore suppose that there is a morphism $g:X'\longrightarrow X$,
where $X'$ is the normalization of $X$ in a finite \'{e}tale
surjective covering $g':U'\longrightarrow U$, tamely ramified along
$D$, such that $g'^*F$ is constant on $U'$.  By Abhyankar's lemma, we
may further suppose that $X'$ is smooth over $S$ and that $U'$ is the
complement in $X'$ of a relative normal-crossings divisor $D'$.

The sheaf $F$ injects into $g'_*g'^*F$, and the cokernel again
satisfies the hypotheses of the lemma.  Iterating this construction
(the same $g$ continues to work) gives a resolution of $F$ by sheaves
of the form $g'_*G$, where $G$ is constant and $g'$ is as above.  It
therefore suffices to prove the assertion for such a sheaf, or,
equivalently, for $(X'/S,D',G)$ with $G$ constant on $U'$.  This is a
special case of SGA~5, II, 4.4.

\subsubsection*{Lemma 2.1.11.}
\underline{Let $S$ be a henselian trait, let $f:X\longrightarrow S$
be smooth, let $D\subset X$ be a relative normal-crossings divisor,
put $U=X-D\mathrel{\mathop{\hookrightarrow}^{j}}X$, and let $F$ be a
sheaf of $\Lambda$-modules on $U$ which is locally constant and tamely
ramified along $D$.  Then}
\[
\underline{R\Phi(j_!F)=0.}
\]

\SGAvar{The d\'{e}vissage used in 2.1.11 reduces us to the case in which $F$ is
constant, with value $G$.  We may suppose that $D$ is the union of
smooth divisors $(D_i)_{i\in I}$.  For $J\subset I$, let $D_J$ be the
intersection of the $D_i$ for $i\in J$.  The sheaf $F$ has a
resolution $F^\bullet$ whose terms are sums of sheaves $G_J$, constant
on $D_J$ and extended by zero to $X$.  It is therefore enough to prove
the lemma for one such sheaf, and this follows from (2.1.5), applied to
$D_J/S$.

}{The d\'{e}vissage just used reduces us to the case in which $F$ is
constant, with value $G$.  We may suppose that $D$ is the union of
smooth divisors $(D_i)_{i\in I}$.  For $J\subset I$, let $D_J$ be the
intersection of the $D_i$ for $i\in J$.  The sheaf $F$ has a
resolution $F^\bullet$ whose terms are sums of sheaves $G_J$, constant
on $D_J$ and extended by zero to $X$.  It is therefore enough to prove
the lemma for one such sheaf, and this follows from (2.1.5), applied to
$D_J/S$.

}{E7II-XIII-010}\SGApage{II}{\(\approx\)106}
I conjecture that (2.1.7.5) is an isomorphism when $S$ and $S'$ are
excellent.  Here is a partial result (see also 2.4.2).

\subsubsection*{Proposition 2.1.12.}
\underline{The morphism (2.1.7.5) is an isomorphism in the following
cases:}

\medskip
\noindent(a) $S$ is excellent and $S'$ is its completion;

\noindent(b) $S$ and $S'$ have equal characteristic zero.

In case (b), we reduce to the situation in which $S$ and $S'$ are
strictly local and a uniformizer of $S$ is still a uniformizer of
$S'$ (extend the scalars without changing
$\bar S'\longrightarrow\bar S$).  In both cases (a) and (b), writing
an overbar for normalization in a separable closure of the fraction
field, one then has
\[
\bar S'\simeq S'\times_S\bar S.
\]
Returning to the definitions, it is enough to prove that
$S'\longrightarrow S$ is universally locally acyclic (apply SGA~4,
XVI, 1.1).  This follows from I, 0.5.1 (based on Artin's version of
N\'{e}ron desingularization) and SGA~4, XV, 2.1 (the local-acyclicity
theorem for smooth morphisms).

\subsubsection*{2.1.13.}
Suppose that $f:X\longrightarrow S$ is of finite type.  The
cohomological dimension (SGA~4, XVII, 1.2.9) of $R\Psi$ is then at most
$\dim(X_\eta)$, as follows from 2.1.4 and SGA~4, XIV, 3.2 (cf. I,
4.2(i)).  Thus $R\Psi$ extends to the entire derived category,
\[
R\Psi:D(X,\Lambda)\longrightarrow D(X_s\times_sS,\Lambda),
\]
and it carries $D^-$ into $D^-$.  Let
$L\in\operatorname{Ob}D^-(\Lambda)$ be a bounded-above complex of
right $\Lambda$-modules, and let
$K\in\operatorname{Ob}D^-(X,\Lambda)$.  There is then an isomorphism
\begin{equation}\tag{2.1.13.1}
L\overset{\mathbf L}{\otimes}_\Lambda R\Psi(K)
\xrightarrow{\ \sim\ }
R\Psi\bigl(L\overset{\mathbf L}{\otimes}_\Lambda K\bigr).
\end{equation}
To verify (2.1.13.1), first replace $L$ by a complex of free
$\Lambda$-modules.  Since $R\Psi$ is triangulated and has finite
cohomological dimension, it is enough to verify the assertion after
replacing $L$ by one of its terms.  Passage to the filtered colimit
reduces us to the case in which $L$ is a finite free module, and then
to the tautological case $L=\Lambda$.

\SGApage{II}{\(\approx\)107}
It follows from (2.1.13.1), by the usual argument (cf. SGA~4, XVII,
5.2.11), that if $K\in\operatorname{Ob}D(X_\eta,\Lambda)$ has finite
Tor-dimension at most $d$, where $d\in\mathbb Z$ (SGA~4, XVII, 4.1.9),
then $R\Psi_\eta(K)$ likewise has Tor-dimension at most $d$.


\subsection*{2.2. A compatibility (trace morphisms)}

\SGAvar{\subsubsection*{2.2.1.}
Let $S$ be a henselian trait, and let
$u:X_s\longrightarrow Y_s$ be a separated, quasi-finite, flat
morphism of finite presentation between $s$-schemes.  For every
abelian sheaf $F$ on $Y_s$, there is a trace morphism
\[
\operatorname{Tr}_u:u_!uF\longrightarrow F
\]
(SGA~4, XVII, 6.2.3).  If $F$ is an abelian sheaf on
$Y_s\times_sS$ or on $Y_s\times_s\eta$, this gives an analogous
morphism $\operatorname{Tr}_u:u_!u^*F\longrightarrow F$.

}{\subsubsection*{2.2.1.}
Let $S$ be a henselian trait, and let
$u:X_s\longrightarrow Y_s$ be a separated, quasi-finite, flat
morphism of finite presentation between $s$-schemes.  For every
abelian sheaf $F$ on $Y_s$, there is a trace morphism
\[
\operatorname{Tr}_u:u_!u^*F\longrightarrow F
\]
(SGA~4, XVII, 6.2.3).  If $F$ is an abelian sheaf on
$Y_s\times_sS$ or on $Y_s\times_s\eta$, this gives an analogous
morphism $\operatorname{Tr}_u:u_!u^*F\longrightarrow F$.

}{E7II-XIII-011}Now let $u:X\longrightarrow Y$ be a separated, quasi-finite, flat
morphism of finite presentation between $S$-schemes.  We again write
$u$ for the induced morphism
\[
X_s\times_sS\longrightarrow Y_s\times_sS,
\]
and we use $\Psi$ for the functors of 1.3.3, both for $X$ and for $Y$.

\subsubsection*{Proposition 2.2.2.}
\underline{For every abelian sheaf $F$ on $Y$, the following diagram
is commutative:}
\begin{center}
\begin{tikzcd}[column sep=large, row sep=large]
u_!u^*\Psi F
  \arrow[r, "{(1.3.7.1)}"] \arrow[d, "\operatorname{Tr}_u"']
  \& u_!\Psi u^*F \arrow[r, "{(1.3.8.1)}"]
  \& \Psi u_!u^*F \arrow[d, "\operatorname{Tr}_u"] \\
\Psi F \arrow[rr, equal] \& {} \& \Psi F.
\end{tikzcd}
\end{center}
The verification is left to the reader.  There is a variant for
$\Psi_\eta$ and for $F$ on $Y_\eta$.

\SGApage{II}{\(\approx\)108}
\SGAvar{\subsubsection*{2.2.3.}
Let $\Lambda$ be as in 2.2.1, and let $u:X\longrightarrow Y$ be a
quasi-finite flat morphism between separated, finite-type, flat
$S$-schemes:
\begin{center}
\begin{tikzcd}[column sep=large, row sep=large]
X \arrow[rr, "u"] \arrow[dr, "f"'] \& {} \& Y \arrow[dl, "g"] \\
{} \& S \& {}
\end{tikzcd}
\end{center}
  We again write $f$ and $g$ for the morphisms
\[
X_s\times_sS\longrightarrow S,
\qquad
Y_s\times_sS\longrightarrow S,
\]
and suppose that $\dim Y_\eta\leq n$.  For every
$K\in\operatorname{Ob}D(S,\Lambda)$, there is a trace morphism
(SGA~4, XVIII, 2.9)
\[
\operatorname{Tr}_f:Rf_!f^*K(n)[2n]\longrightarrow K.
\]

}{\subsubsection*{2.2.3.}
Let $\Lambda$ be as in 2.1.1, and let $u:X\longrightarrow Y$ be a
quasi-finite flat morphism between separated, finite-type, flat
$S$-schemes:
\begin{center}
\begin{tikzcd}[column sep=large, row sep=large]
X \arrow[rr, "u"] \arrow[dr, "f"'] \& {} \& Y \arrow[dl, "g"] \\
{} \& S \& {}
\end{tikzcd}
\end{center}
  We again write $f$ and $g$ for the morphisms
\[
X_s\times_sS\longrightarrow S,
\qquad
Y_s\times_sS\longrightarrow S,
\]
and suppose that $\dim Y_\eta\leq n$.  For every
$K\in\operatorname{Ob}D(S,\Lambda)$, there is a trace morphism
(SGA~4, XVIII, 2.9)
\[
\operatorname{Tr}_f:Rf_!f^*K(n)[2n]\longrightarrow K.
\]

}{E7II-XIII-012}\subsubsection*{Corollary 2.2.4.}
\underline{The following diagram is commutative:}
\begin{center}
\resizebox{\textwidth}{!}{
\begin{tikzcd}[ampersand replacement=\&, column sep=large, row sep=large]
Rg_!u_!u^*R\Psi\,g^*K(n)[2n]
  \arrow[r, "{(1.3.7.1)}"] \arrow[d, "\operatorname{Tr}_u"']
  \& Rf_!R\Psi\,f^*K(n)[2n]
      \arrow[r, "{(2.1.10.1)}"]
  \& Rf_!f^*K(n)[2n] \arrow[d, "\operatorname{Tr}_f"] \\
Rg_!R\Psi\,g^*K(n)[2n]
  \arrow[r, "{(2.1.10.1)}"]
  \& Rg_!g^*K(n)[2n] \arrow[r, "\operatorname{Tr}_g"]
  \& K
\end{tikzcd}}
\end{center}
\underline{If $u$ is \'{e}tale, (1.3.7.1) is an isomorphism.}

With the arguments suppressed, the diagram (2.2.4) expands as
follows:
\begin{center}
\resizebox{\textwidth}{!}{
\begin{tikzcd}[ampersand replacement=\&, column sep=large, row sep=huge]
Rg_!u_!u^*R\Psi\,g^*
  \arrow[r] \arrow[d, "\operatorname{Tr}_u"']
  \arrow[rd, phantom, "{(1)}" description]
  \& Rg_!R\Psi\,u_!u^*g^*
      \arrow[r] \arrow[d, "\operatorname{Tr}_u"']
      \arrow[rd, phantom, "{(2)}" description]
  \& Rg_!u_!u^*g^*
      \arrow[r, "\sim"] \arrow[d, "\operatorname{Tr}_u"']
      \arrow[rd, phantom, "{(3)}" description]
  \& Rf_!f^* \arrow[d, "\operatorname{Tr}_f"] \\
Rg_!R\Psi\,g^* \arrow[r, equal]
  \& Rg_!R\Psi\,g^* \arrow[r]
  \& Rg_!g^* \arrow[r, "\operatorname{Tr}_g"]
  \& K.
\end{tikzcd}}
\end{center}
The commutativity of (1) follows from 2.2.2; that of (2) is the
functoriality of (2.1.10.2); and that of (3) is the transitivity of the
trace (SGA~4, XVIII, 2.9, Variant 3).

\SGApage{II}{\(\approx\)109}
\subsubsection*{2.2.4.1.}
There is an analogous result for $\Psi_\eta$ and
$K\in\operatorname{Ob}D(\eta,\Lambda)$.

\SGAvar{\subsubsection*{2.2.5.}
Let $\Phi$ be a subset of $Y_s$ which is proper over $s$, and write
$Rg_\Phi$ for the ``direct image with support in $\Phi$'' functor.
Let $\operatorname{Tr}_g$ denote the composite
\begin{equation}\tag{2.2.5.1}
\begin{split}
\operatorname{Tr}_g:
Rf_\Phi R\Psi g^*K(n)[2n]
&\longrightarrow Rg_!R\Psi g^*K(n)[2n] \\
&\longrightarrow Rg_!g^*K(n)[2n]
\longrightarrow K,
\end{split}
\end{equation}
and similarly for $\Psi_\eta$ and $K\in D(\eta,\Lambda)$.  

}{\subsubsection*{2.2.5.}
Let $\Phi$ be a subset of $Y_s$ which is proper over $s$, and write
$Rg_\Phi$ for the ``direct image with support in $\Phi$'' functor.
Let $\operatorname{Tr}_g$ denote the composite
\begin{equation}\tag{2.2.5.1}
\begin{split}
\operatorname{Tr}_g:
Rg_\Phi R\Psi g^*K(n)[2n]
&\longrightarrow Rg_!R\Psi g^*K(n)[2n] \\
&\longrightarrow Rg_!g^*K(n)[2n]
\longrightarrow K,
\end{split}
\end{equation}
and similarly for $\Psi_\eta$ and $K\in D(\eta,\Lambda)$.  

}{E7II-XIII-013}If $\Phi$ is a finite set of closed points, then
\[
Rg_\Phi\circ u_!=Rf_{u^{-1}\Phi},
\]
and hence there is a trace morphism
\[
\operatorname{Tr}_u:
Rf_{u^{-1}\Phi}u^*=Rg_\Phi u_!u^*\longrightarrow Rg_\Phi.
\]
By 2.2.4, the following diagram is commutative:
\begin{center}
\resizebox{\textwidth}{!}{
\begin{tikzcd}[ampersand replacement=\&, column sep=large, row sep=large]
Rf_{u^{-1}\Phi}u^*R\Psi\,g^*K(n)[2n]
  \arrow[r] \arrow[d, "\operatorname{Tr}_u"']
  \& Rf_{u^{-1}\Phi}R\Psi\,f^*K(n)[2n]
      \arrow[r, "\operatorname{Tr}_f"]
  \& K \arrow[d, equal] \\
Rg_\Phi R\Psi\,g^*K(n)[2n]
  \arrow[rr, "\operatorname{Tr}_g"']
  \& {} \& K.
\end{tikzcd}}
\end{center}

\SGAvar{Suppose that $u$ is \'{e}tale, let $x\in X_s$, and suppose that $x$
and $y=u(x)$ have the same residue field.  Then
\[
Rf_{\{x\}}\xrightarrow{\ \sim\ }Rg_{\{x\}}.
\]
  This
isomorphism, given by $\operatorname{Tr}_u$, expresses the \'{e}tale
locality of cohomology with support at $x$.  The diagram
\begin{equation}\tag{2.2.5.2}
\begin{tikzcd}[column sep=huge, row sep=large]
Rf_{\{x\}}R\Psi\,f^*K(n)[2n]
  \arrow[r, "\operatorname{Tr}_f"] \arrow[d, equal]
  \& K \arrow[d, equal] \\
Rg_{\{y\}}R\Psi\,g^*K(n)[2n]
  \arrow[r, "\operatorname{Tr}_g"]
  \& K
\end{tikzcd}
\end{equation}
is commutative.

\subsubsection*{Scholium 2.2.6.}
The compatibility (2.3.5.2) says that the trace morphism (2.3.5.1),
for $\Phi=\{x\}$, is a local invariant: it depends only on the
henselization of the $S$-scheme $X$ at $x$.  The same holds for
$R\Psi_\eta$.  


}{Suppose that $u$ is \'{e}tale, let $x\in X_s$, and suppose that $x$
and $y=u(x)$ have the same residue field.  Then
\[
Rf_{\{x\}}\xrightarrow{\ \sim\ }Rg_{\{y\}}.
\]
  This
isomorphism, given by $\operatorname{Tr}_u$, expresses the \'{e}tale
locality of cohomology with support at $x$.  The diagram
\begin{equation}\tag{2.2.5.2}
\begin{tikzcd}[column sep=huge, row sep=large]
Rf_{\{x\}}R\Psi\,f^*K(n)[2n]
  \arrow[r, "\operatorname{Tr}_f"] \arrow[d, equal]
  \& K \arrow[d, equal] \\
Rg_{\{y\}}R\Psi\,g^*K(n)[2n]
  \arrow[r, "\operatorname{Tr}_g"]
  \& K
\end{tikzcd}
\end{equation}
is commutative.

\subsubsection*{Scholium 2.2.6.}
The compatibility (2.2.5.2) says that the trace morphism (2.2.5.1),
for $\Phi=\{y\}$, is a local invariant: it depends only on the
henselization of the $S$-scheme $X$ at $x$.  The same holds for
$R\Psi_\eta$.  


}{E7II-XIII-014}\SGApage{II}{110}
\subsection*{2.3. Finiteness theorem (in equal characteristic zero)}

\subsubsection*{Theorem 2.3.1.}
\underline{Let $S$ be a henselian trait as in 0.2.5, let $\Lambda$ be
a noetherian torsion ring, and let $f:X\longrightarrow S$ be a
morphism of finite type.  Suppose that $S$ has equal characteristic
zero.  Then, for every constructible sheaf of $\Lambda$-modules $F$
on $X_\eta$, the sheaves $R^i\Psi_{\bar\eta}(F)$ on $X_{\bar s}$ are
constructible.}

We begin by repeating the opening of the proof of SGA~4, XIV, 1.10.
Apply SGA~4, IX, 2.14(ii).  There is a resolution $F^\bullet$ of $F$
\begin{center}
\begin{tikzcd}[column sep=large, row sep=large]
{} \& F \arrow[d] \& {} \& {} \\
0 \arrow[r] \& F^0 \arrow[r] \& F^1 \arrow[r] \& \cdots
\end{tikzcd}
\end{center}
by sheaves of the form $\bigoplus_i p'_{i*}G_i$, where
$p'_i:Y_{i\eta}\longrightarrow X_\eta$ is finite and $G_i$ is a
constructible constant sheaf.  We may suppose that $Y_{i\eta}$ is
reduced; let $p_i:Y_i\longrightarrow X$ be the normalization of $X$
in $Y_{i\eta}$.  Using the spectral sequence of the resolution
$F^\bullet$ and applying (2.1.7.1) in the elementary case of the
finite morphism $p_i$, we see that it suffices to prove the theorem
for each pair $(Y_i,G_i)$.

We next prove the theorem when $F$ is the constant sheaf defined by a
$\Lambda$-module $M$.  Recall that $\Lambda$ is a
$\mathbb Z/n$-algebra for a suitable $n$.  The module $M$ may
therefore be d\'{e}viss\'{e} into modules $M_i$ over
$\Lambda/\ell_i\Lambda$, where each $\ell_i\mid n$ is a suitable
prime.  This reduces us to the case in which $\Lambda$ is a
$\mathbb Z/\ell$-algebra, with $\ell$ prime.  Formula (2.1.13.1) then
gives
\[
M\otimes_{\mathbb Z/\ell}
R^i\Psi_{\bar\eta}(\mathbb Z/\ell)
\xrightarrow{\ \sim\ }
R^i\Psi_{\bar\eta}(M).
\]
It remains to treat the constant sheaf $F=\mathbb Z/\ell$.

\SGApage{II}{\(\approx\)111}
Choose a descending sequence of closed subsets
\[
X_\eta=X'_0\supset\cdots\supset X'_n
\]
such that $(X'_i-X'_{i+1})_{\mathrm{red}}$ is smooth over $\eta$, and
let $X_i=\overline{X'_i}$ be the closure in $X$.  The constant sheaf
$\mathbb Z/\ell$ is d\'{e}viss\'{e} into the constant sheaves
$\mathbb Z/\ell$ on the strata $X'_i-X'_{i+1}$, extended by zero.
Thus:

\medskip
\noindent\underline{Reduction 1.}
It is enough to prove Theorem 2.3.1 for a sheaf
$j_!\mathbb Z/\ell$, where
$j:U\hookrightarrow X_\eta$ is the inclusion of a smooth open subset.

By resolution of singularities there is a morphism
$g:X_1\longrightarrow X$ such that

\medskip
\noindent(a) $g$ is proper and induces an isomorphism
$g^{-1}(U)\xrightarrow{\sim}U$;

\noindent(b) $X_1$ is regular and $X_1-g^{-1}(U)$ is a
normal-crossings divisor in $X_1$, which is a union of regular
divisors $D_\alpha$.

Let $j_1:g^{-1}(U)\hookrightarrow X_1$.  Then
\[
Rg_*j_{1!}\mathbb Z/\ell=j_!\mathbb Z/\ell.
\]
Applying (2.1.7.1) and the finiteness theorem SGA~4, XIV, 1.1, we see
that it suffices to prove Theorem 2.3.1 for
$(X_1,j_{1!}\mathbb Z/\ell)$.

There is a resolution
\[
0\longrightarrow j_{1!}\mathbb Z/\ell
\longrightarrow\mathbb Z/\ell
\longrightarrow\bigoplus_\alpha(\mathbb Z/\ell)_{D_\alpha}
\longrightarrow
\bigoplus_{\alpha,\beta}
(\mathbb Z/\ell)_{D_\alpha\cap D_\beta}
\longrightarrow\cdots,
\]
and hence:

\medskip
\noindent\underline{Reduction 2.}
It is enough to prove Theorem 2.3.1 for the constant sheaf
$\mathbb Z/\ell$ on a regular $S$-scheme $X$ such that $X_s$ is a
normal-crossings divisor in $X$.

In this final case one uses the explicit computation I, 3.3, which in
turn uses the purity theorem SGA~4, XIX, 3.2 and the characteristic-zero
hypothesis.


\subsection*{2.4. Isolated singularities}

\SGAvar{\subsubsection*{2.4.1.}
Throughout this section, for simplicity, we suppose that $S$ is a
strictly henselian trait.  Let $\Lambda$ be a noetherian torsion ring
whose torsion is prime to the residual characteristic of $S$, let
$f:X\longrightarrow S$ be a morphism of finite type, and let $F$ be a
constructible sheaf of $\Lambda$-modules on $X$.  A point $x\in X$ is
called a smooth point of $(X,F)$ if $f$ is smooth at $X$ and $F$ is
locally constant in a neighborhood of $x$.  Let $\Sigma$ denote the
set of points where $(X,F)$ is not smooth, and let $\Sigma_s$ be its
special fiber.  We know that $R\Phi(F)=0$ outside $\Sigma_s$.

}{\subsubsection*{2.4.1.}
Throughout this section, for simplicity, we suppose that $S$ is a
strictly henselian trait.  Let $\Lambda$ be a noetherian torsion ring
whose torsion is prime to the residual characteristic of $S$, let
$f:X\longrightarrow S$ be a morphism of finite type, and let $F$ be a
constructible sheaf of $\Lambda$-modules on $X$.  A point $x\in X$ is
called a smooth point of $(X,F)$ if $f$ is smooth at $x$ and $F$ is
locally constant in a neighborhood of $x$.  Let $\Sigma$ denote the
set of points where $(X,F)$ is not smooth, and let $\Sigma_s$ be its
special fiber.  We know that $R\Phi(F)=0$ outside $\Sigma_s$.

}{E7II-XIII-015}\SGApage{II}{112}
\subsubsection*{Theorem 2.4.2.}
\underline{At every isolated point $x$ of $\Sigma_s$, the modules
$R^i\Phi_{\bar\eta}(F)_x$ are $\Lambda$-modules of finite type and
are invariant under change of trait (via 2.1.7.1).}

In equal characteristic zero, these assertions were proved in
2.1.12 and 2.3.1 without assuming that $x$ is isolated.

Let $x$ be an isolated point of $\Sigma_s$.  The theorem is local in
a neighborhood of $x$, so we may suppose that $f$ is proper.  We then
have the long exact sequence (2.1.8.9)
\[
\cdots\longrightarrow H^i(X_s,F)
\longrightarrow H^i(X_{\bar\eta},F)
\longrightarrow
\mathbb H^i\bigl(X_s,R\Phi_{\bar\eta}(F)\bigr)
\longrightarrow\cdots .
\]
In this sequence, $H^i(X_s,F)$ and $H^i(X_{\bar\eta},F)$ are
invariant under change of trait and are $\Lambda$-modules of finite
type (SGA~4, XIV, 1.1).  The modules
$\mathbb H^i(X_s,R\Phi_{\bar\eta}(F))$ therefore have the same
properties.  Recall the following lemma.

\subsubsection*{Lemma 2.4.3.}
\underline{Let $(T,\Lambda)$ be a ringed topos, let $Z$ be a closed
subtopos of $T$, and let $i:Z\hookrightarrow T$ be the inclusion.}
\underline{Then the functor}
\[
i_*:D(Z,\Lambda)\longrightarrow D(T,\Lambda)
\]
\underline{is an equivalence from $D(Z,\Lambda)$ onto the full
subcategory of $D(T,\Lambda)$ consisting of the complexes $K$ for
which $\mathcal H^i(K)$ is supported on $Z$.}

Indeed, for such a complex $K$, the morphism
$K\longrightarrow i_*i^*K$ is a quasi-isomorphism.

We may consequently regard $R\Phi_{\bar\eta}(F)$ as an object of
\[
D^b(\Sigma_s,\Lambda)
=D^b(\Sigma_s-\{x\},\Lambda)\oplus D^b(\{x\},\Lambda).
\]
Passing to a direct summand, we find that
\[
\mathbb H^i\bigl(\{x\},R\Phi_{\bar\eta}(F)\bigr)
=R^i\Phi_{\bar\eta}(F)_x
\]
has the properties stated in 2.4.2.  This completes the proof.

\subsubsection*{Remark 2.4.4.}
The same proof shows that, if $K$ is constructible in
$D^+(X,\Lambda)$ and $x$ is isolated in the support of
$R\Phi_{\bar\eta}(K)$, then the cohomology modules of the stalk
$R\Phi_{\bar\eta}(K)_x$ are of finite type over $\Lambda$.  If, after
a change of trait, $x$ remains an isolated point of the support, these
modules are invariant under that change of trait.


\SGApage{II}{\(\approx\)113}
\SGAvar{\subsubsection*{2.4.5.}
Let $x\in X_s$ be an isolated point where $(X,F)$ is not smooth.
More generally, let $K\in\operatorname{Ob}D^b(X,\Lambda)$, and let
$x$ be an isolated point of the support of $R\Phi(K)$.  Let
$i:\{x\}\hookrightarrow X_s$ be the inclusion.  The complex
$i_*i^*R\Phi(K)$, which is isomorphic to $R\Phi(K)$ in a neighborhood
of $x$, is a direct summand of $R\Phi(F)$ by 2.4.3.  Thus, for
$\sigma\in I$, the variation morphism induces
\[
\operatorname{Var}(\sigma):
i_*i^*R\Phi_{\bar\eta}(K)\longrightarrow R\Psi_{\bar\eta}(K),
\]
that is,
\begin{equation}\tag{2.4.5.1}
\operatorname{Var}(\sigma):
\bigl(R\Phi_{\bar\eta}(K)\bigr)_x
\longrightarrow
R\Gamma_{\{x\}}\bigl(R\Psi_{\bar\eta}(K)\bigr).
\end{equation}


\subsubsection*{2.4.6.}
Suppose that the support $\Sigma$ of $R\Phi(K)$ consists of finitely
many isolated points.  For every $\sigma\in I$, the endomorphism
$\sigma-1$ of $R\Psi_{\bar\eta}(K)$ is then the composite displayed
in the following diagram:
\begin{equation}\tag{2.4.6.1}
\resizebox{\textwidth}{!}{
\begin{tikzcd}[ampersand replacement=\&, column sep=huge, row sep=huge]
R\Psi_{\bar\eta}(K)
  \arrow[r] \arrow[d, "\sigma-1"']
  \& R\Phi_{\bar\eta}(K)
     =\displaystyle\bigoplus_{x\in\Sigma}
       i_{x*}\bigl(R\Phi_{\bar\eta}(K)\bigr)_x
     \arrow[d, "{\bigoplus i_{x*}\operatorname{Var}(\sigma)}"'] \\
R\Psi_{\bar\eta}(K)
  \& \displaystyle\bigoplus_{x\in\Sigma}
       i_{x}R\Gamma_{\{x\}}R\Psi_{\bar\eta}(K)
     \arrow[l]
\end{tikzcd}}
\end{equation}

If $f$ is proper, applying $H^i$ to (2.4.6.1) shows that the
endomorphism $\sigma-1$ of $\mathbb H^i(X_{\bar\eta},K)$ is the
composite
\begin{equation}\tag{2.4.6.2}
\resizebox{\textwidth}{!}{
\begin{tikzcd}[ampersand replacement=\&, column sep=large, row sep=huge]
\mathbb H^i(X_{\bar\eta},K)
  \arrow[r, "\sim"] \arrow[d, "\sigma-1"']
  \& \mathbb H^i\bigl(X_s,R\Psi_{\bar\eta}(K)\bigr)
     \arrow[r]
  \& \mathbb H^i\bigl(X_s,R\Phi_{\bar\eta}(K)\bigr)
     =\displaystyle\bigoplus_{x\in\Sigma}R^i\Phi_{\bar\eta}(K)_x
     \arrow[d, "{\bigoplus\operatorname{Var}(\sigma)}"'] \\
\mathbb H^i(X_{\bar\eta},K)
  \& \mathbb H^i\bigl(X_s,R\Psi_{\bar\eta}(K)\bigr)
     \arrow[l, "\sim"']
  \& \displaystyle\bigoplus_{x\in\Sigma}
     \mathbb H^i_{\{x\}}\bigl(X_s,R\Psi_{\bar\eta}(K)\bigr)
     \arrow[l]
\end{tikzcd}}
\end{equation}
This describes the action of $I$ on $\mathbb H^i(X_{\bar\eta},K)$;
its essential ingredients are the variation morphisms (2.4.4.1),
which are local on $X$.  We also have the long exact sequence
\begin{equation}\tag{2.4.6.3}
\longrightarrow \mathbb H^i(X_s,K)
\longrightarrow \mathbb H^i(X_{\bar\eta},K)
\longrightarrow \bigoplus_{x\in\Sigma}R^i\Phi_{\bar\eta}(K)_x
\longrightarrow\cdots .
\end{equation}

}{\subsubsection*{2.4.5.}
Let $x\in X_s$ be an isolated point where $(X,F)$ is not smooth.
More generally, let $K\in\operatorname{Ob}D^b(X,\Lambda)$, and let
$x$ be an isolated point of the support of $R\Phi(K)$.  Let
$i:\{x\}\hookrightarrow X_s$ be the inclusion.  The complex
$i_*i^*R\Phi(K)$, which is isomorphic to $R\Phi(K)$ in a neighborhood
of $x$, is a direct summand of $R\Phi(K)$ by 2.4.3.  Thus, for
$\sigma\in I$, the variation morphism induces
\[
\operatorname{Var}(\sigma):
i_*i^*R\Phi_{\bar\eta}(K)\longrightarrow R\Psi_{\bar\eta}(K),
\]
that is,
\begin{equation}\tag{2.4.5.1}
\operatorname{Var}(\sigma):
\bigl(R\Phi_{\bar\eta}(K)\bigr)_x
\longrightarrow
R\Gamma_{\{x\}}\bigl(R\Psi_{\bar\eta}(K)\bigr).
\end{equation}


\subsubsection*{2.4.6.}
Suppose that the support $\Sigma$ of $R\Phi(K)$ consists of finitely
many isolated points.  For every $\sigma\in I$, the endomorphism
$\sigma-1$ of $R\Psi_{\bar\eta}(K)$ is then the composite displayed
in the following diagram:
\begin{equation}\tag{2.4.6.1}
\resizebox{\textwidth}{!}{
\begin{tikzcd}[ampersand replacement=\&, column sep=huge, row sep=huge]
R\Psi_{\bar\eta}(K)
  \arrow[r] \arrow[d, "\sigma-1"']
  \& R\Phi_{\bar\eta}(K)
     =\displaystyle\bigoplus_{x\in\Sigma}
       i_{x*}\bigl(R\Phi_{\bar\eta}(K)\bigr)_x
     \arrow[d, "{\bigoplus i_{x*}\operatorname{Var}(\sigma)}"'] \\
R\Psi_{\bar\eta}(K)
  \& \displaystyle\bigoplus_{x\in\Sigma}
       i_{x*}R\Gamma_{\{x\}}R\Psi_{\bar\eta}(K)
     \arrow[l]
\end{tikzcd}}
\end{equation}

If $f$ is proper, applying $H^i$ to (2.4.6.1) shows that the
endomorphism $\sigma-1$ of $\mathbb H^i(X_{\bar\eta},K)$ is the
composite
\begin{equation}\tag{2.4.6.2}
\resizebox{\textwidth}{!}{
\begin{tikzcd}[ampersand replacement=\&, column sep=large, row sep=huge]
\mathbb H^i(X_{\bar\eta},K)
  \arrow[r, "\sim"] \arrow[d, "\sigma-1"']
  \& \mathbb H^i\bigl(X_s,R\Psi_{\bar\eta}(K)\bigr)
     \arrow[r]
  \& \mathbb H^i\bigl(X_s,R\Phi_{\bar\eta}(K)\bigr)
     =\displaystyle\bigoplus_{x\in\Sigma}R^i\Phi_{\bar\eta}(K)_x
     \arrow[d, "{\bigoplus\operatorname{Var}(\sigma)}"'] \\
\mathbb H^i(X_{\bar\eta},K)
  \& \mathbb H^i\bigl(X_s,R\Psi_{\bar\eta}(K)\bigr)
     \arrow[l, "\sim"']
  \& \displaystyle\bigoplus_{x\in\Sigma}
     \mathbb H^i_{\{x\}}\bigl(X_s,R\Psi_{\bar\eta}(K)\bigr)
     \arrow[l]
\end{tikzcd}}
\end{equation}
This describes the action of $I$ on $\mathbb H^i(X_{\bar\eta},K)$;
its essential ingredients are the variation morphisms (2.4.5.1),
which are local on $X$.  We also have the long exact sequence
\begin{equation}\tag{2.4.6.3}
\longrightarrow \mathbb H^i(X_s,K)
\longrightarrow \mathbb H^i(X_{\bar\eta},K)
\longrightarrow \bigoplus_{x\in\Sigma}R^i\Phi_{\bar\eta}(K)_x
\longrightarrow\cdots .
\end{equation}

}{E7II-XIII-016}The same argument applies when $f$ is not necessarily proper,
provided that $(X,K)$ satisfies 2.1.9.  The diagram (2.4.6.5) below
defines a morphism
\begin{equation}\tag{2.4.6.4}
\operatorname{Var}(\sigma):
R\Gamma(X_{\bar\eta},K)\longrightarrow R\Gamma_c(X_{\bar\eta},K)
\end{equation}
such that the endomorphism $\sigma-1$ of
$R\Gamma(X_{\bar\eta},K)$ (respectively,
$R\Gamma_c(X_{\bar\eta},K)$) is the composite
$\alpha\operatorname{Var}(\sigma)$ (respectively,
$\operatorname{Var}(\sigma)\alpha$), where
\[
\alpha:R\Gamma_c(X_{\bar\eta},K)\longrightarrow
R\Gamma(X_{\bar\eta},K).
\]
\begin{equation}\tag{2.4.6.5}
\resizebox{\textwidth}{!}{
\begin{tikzcd}[ampersand replacement=\&, column sep=huge, row sep=huge]
R\Gamma(X_{\bar\eta},K)
  \arrow[r, "\sim"] \arrow[d, "\operatorname{Var}(\sigma)"']
  \& R\Gamma\bigl(X_s,R\Psi_{\bar\eta}(K)\bigr)
     \arrow[r]
  \& \displaystyle\bigoplus_{x\in\Sigma}
     \bigl(R\Phi_{\bar\eta}(K)\bigr)_x
     \arrow[d, "\bigoplus"']
  \& {} \\
R\Gamma_c(X_{\bar\eta},K)
  \arrow[r, "\sim"']
  \& R\Gamma_c\bigl(X_s,R\Psi_{\bar\eta}(K)\bigr)
  \& \displaystyle\bigoplus_{x\in\Sigma}
     R\Gamma_{\{x\}}\bigl(X_s,R\Psi_{\bar\eta}(K)\bigr)
     \arrow[l]
  \& .
\end{tikzcd}}
\end{equation}

\SGApage{II}{\(\approx\)114}
In this case we also have the long exact sequences
\begin{equation}\tag{2.4.6.6}
\resizebox{\textwidth}{!}{
\begin{tikzcd}[ampersand replacement=\&, column sep=large, row sep=large]
\cdots \arrow[r]
  \& \mathbb H^i_c(X_s,K) \arrow[r] \arrow[d]
  \& \mathbb H^i_c(X_{\bar\eta},K) \arrow[r] \arrow[d]
  \& \displaystyle\bigoplus_{x\in\Sigma}
     \bigl(R^i\Phi_{\bar\eta}(K)\bigr)_x
     \arrow[r] \arrow[d, equal]
  \& \cdots \\
\cdots \arrow[r]
  \& \mathbb H^i(X_s,K) \arrow[r]
  \& \mathbb H^i(X_{\bar\eta},K) \arrow[r]
  \& \displaystyle\bigoplus_{x\in\Sigma}
     R^i\Phi_{\bar\eta}(K)_x \arrow[r]
  \& \cdots
\end{tikzcd}}
\end{equation}




\SGAvar{\subsubsection*{2.4.6.}
Let $X_{s(x)}$ be the henselization of $X_s$ at $x$, and put
\[
X^*_{s(x)}=X_{s(x)}-\{x\}.
\]
The long hypercohomology exact sequences of
$(X_{s(x)},\{x\})$, with coefficients in $K_s$,
$R\Psi_{\bar\eta}(K)$, and $R\Phi_{\bar\eta}(K)$, give the following
commutative diagram of long exact sequences:
\begin{equation}\tag{2.4.6.1}
\resizebox{\textwidth}{!}{
\begin{tikzcd}[ampersand replacement=\&, column sep=large, row sep=large]
{}
  \& 0 \arrow[r] \arrow[d]
  \& R^{n-1}\Phi_{\bar\eta}(K)_x \arrow[r] \arrow[d]
  \& R^{n-1}\Phi_{\bar\eta}(K)_x \arrow[r] \arrow[d]
  \& 0 \arrow[d]
  \& {} \\
{} \arrow[r]
  \& \mathbb H^{n-1}(X^*_{s(x)},K)
     \arrow[r, "j"] \arrow[d, equal]
  \& \mathbb H^n_{\{x\}}(X_s,K) \arrow[r] \arrow[d]
  \& \underline{H}^{n}(K)_x \arrow[r] \arrow[d]
  \& \mathbb H^n(X^*_{s(x)},K)
     \arrow[r] \arrow[d, equal]
  \& {} \\
{} \arrow[r]
  \& \mathbb H^{n-1}(X^*_{s(x)},K)
     \arrow[r, "j", "{\mbox{\textcircled{1}}}"'] \arrow[d]
  \& \mathbb H^n_{\{x\}}\bigl(X_s,R\Psi_{\bar\eta}(K)
     \arrow[r] \arrow[d]
  \& R^n\Psi_{\bar\eta}(K)_x \arrow[r] \arrow[d]
  \& \mathbb H^n(X_{s(x)},K) \arrow[r] \arrow[d]
  \& {} \\
{}
  \& 0 \arrow[r]
  \& R^n\Phi_{\bar\eta}(K)_x \arrow[r, equal]
  \& R^n\Phi_{\bar\eta}(K)_x \arrow[r]
  \& 0
  \& {}
\end{tikzcd}}
\end{equation}
The image of the arrow \textcircled{1} is contained in the inertia
invariants.



}{\subsubsection*{2.4.7.}
Let $X_{s(x)}$ be the henselization of $X_s$ at $x$, and put
\[
X^*_{s(x)}=X_{s(x)}-\{x\}.
\]
The long hypercohomology exact sequences of
$(X_{s(x)},\{x\})$, with coefficients in $K_s$,
$R\Psi_{\bar\eta}(K)$, and $R\Phi_{\bar\eta}(K)$, give the following
commutative diagram of long exact sequences:
\begin{equation}\tag{2.4.7.1}
\resizebox{\textwidth}{!}{
\begin{tikzcd}[ampersand replacement=\&, column sep=large, row sep=large]
{}
  \& 0 \arrow[r] \arrow[d]
  \& R^{n-1}\Phi_{\bar\eta}(K)_x \arrow[r] \arrow[d]
  \& R^{n-1}\Phi_{\bar\eta}(K)_x \arrow[r] \arrow[d]
  \& 0 \arrow[d]
  \& {} \\
{} \arrow[r]
  \& \mathbb H^{n-1}(X^*_{s(x)},K)
     \arrow[r, "j"] \arrow[d, equal]
  \& \mathbb H^n_{\{x\}}(X_s,K) \arrow[r] \arrow[d]
  \& \underline{H}^{n}(K)_x \arrow[r] \arrow[d]
  \& \mathbb H^n(X^*_{s(x)},K)
     \arrow[r] \arrow[d, equal]
  \& {} \\
{} \arrow[r]
  \& \mathbb H^{n-1}(X^*_{s(x)},K)
     \arrow[r, "j", "{\mbox{\textcircled{1}}}"'] \arrow[d]
  \& \mathbb H^n_{\{x\}}\bigl(X_s,R\Psi_{\bar\eta}(K)\bigr)
     \arrow[r] \arrow[d]
  \& R^n\Psi_{\bar\eta}(K)_x \arrow[r] \arrow[d]
  \& \mathbb H^n(X^*_{s(x)},K) \arrow[r] \arrow[d]
  \& {} \\
{}
  \& 0 \arrow[r]
  \& R^n\Phi_{\bar\eta}(K)_x \arrow[r, equal]
  \& R^n\Phi_{\bar\eta}(K)_x \arrow[r]
  \& 0
  \& {}
\end{tikzcd}}
\end{equation}
The image of the arrow \textcircled{1} is contained in the inertia
invariants.



}{E7II-XIII-017}\renewcommand{\refname}{Bibliography}
\begin{thebibliography}{9}
\bibitem{GiraudClassifyingTopos}
J.~Giraud, ``Classifying topos,'' pp.~43--56 in
\emph{Toposes, Algebraic Geometry and Logic},
Lecture Notes in Mathematics 274, Springer-Verlag, 1971.

\bibitem{NagataImbedding}
M.~Nagata, ``A generalization of the imbedding problem,''
\emph{J. Math. Kyoto} \textbf{3}, no.~1 (1963), pp.~89--102.
\end{thebibliography}


\markboth{Exposé XIII}{Exposé XIII}

\ifnum\SGAform=1
\clearpage\section*{Critical apparatus}\begingroup\small\raggedright\sloppy
Readings are given in the exact TeX notation of the two English witnesses. Decimal suffixes identify the individual differences in each marked passage.

\par\medskip\noindent\hypertarget{app-E7II-XIII-001.1}{}\textbf{E7II-XIII-001.1}\quad printed folio 89.\par
\noindent Printed reading: \emph{(empty)}\par
\noindent Corrected reading: {\ttfamily\footnotesize \textbackslash{}bar\allowbreak{}}\par
\noindent Source: {\ttfamily\footnotesize fr/sga7ii/\{diplomati\allowbreak{}c,current\}/expose\_XI\allowbreak{}II\_body.tex;\allowbreak{} folio\allowbreak{} 89;\allowbreak{} formula\allowbreak{} 280\allowbreak{}}\par

\par\medskip\noindent\hypertarget{app-E7II-XIII-002.1}{}\textbf{E7II-XIII-002.1}\quad printed folio 95.\par
\noindent Printed reading: \emph{(2 espace(s), 0 retour(s) de ligne dans le témoin TeX)}\par
\noindent Corrected reading: {\ttfamily\footnotesize  \textbackslash{}bar\allowbreak{} }\par
\noindent Source: {\ttfamily\footnotesize fr/sga7ii/\{diplomati\allowbreak{}c,current\}/expose\_XI\allowbreak{}II\_body.tex;\allowbreak{} folio\allowbreak{} 95;\allowbreak{} formula\allowbreak{} 574\allowbreak{}}\par

\par\medskip\noindent\hypertarget{app-E7II-XIII-003.1}{}\textbf{E7II-XIII-003.1}\quad printed folio 97.\par
\noindent Printed reading: {\ttfamily\footnotesize X\allowbreak{}}\par
\noindent Corrected reading: {\ttfamily\footnotesize Y\allowbreak{}}\par
\noindent Source: {\ttfamily\footnotesize fr/sga7ii/\{diplomati\allowbreak{}c,current\}/expose\_XI\allowbreak{}II\_body.tex;\allowbreak{} folio\allowbreak{} 97;\allowbreak{} formula\allowbreak{} 637\allowbreak{}}\par

\par\medskip\noindent\hypertarget{app-E7II-XIII-003.2}{}\textbf{E7II-XIII-003.2}\quad printed folio 97.\par
\noindent Printed reading: {\ttfamily\footnotesize \textbackslash{}bar\allowbreak{} }\par
\noindent Corrected reading: \emph{(empty)}\par
\noindent Source: {\ttfamily\footnotesize fr/sga7ii/\{diplomati\allowbreak{}c,current\}/expose\_XI\allowbreak{}II\_body.tex;\allowbreak{} folio\allowbreak{} 97;\allowbreak{} formula\allowbreak{} 644\allowbreak{}}\par

\par\medskip\noindent\hypertarget{app-E7II-XIII-004.1}{}\textbf{E7II-XIII-004.1}\quad printed folio 98.\par
\noindent Printed reading: {\ttfamily\footnotesize X\allowbreak{}}\par
\noindent Corrected reading: {\ttfamily\footnotesize Y\allowbreak{}}\par
\noindent Source: {\ttfamily\footnotesize fr/sga7ii/\{diplomati\allowbreak{}c,current\}/expose\_XI\allowbreak{}II\_body.tex;\allowbreak{} folio\allowbreak{} 98;\allowbreak{} formula\allowbreak{} 663;\allowbreak{} port/decisions/decis\allowbreak{}ions\_en7II\_r2\_001.js\allowbreak{}on\#C563;\allowbreak{} fixes/FIXLOG\_7II\_CUR\allowbreak{}RENT.jsonl\#C563;\allowbreak{} 7II:149\allowbreak{}}\par

\par\medskip\noindent\hypertarget{app-E7II-XIII-004.2}{}\textbf{E7II-XIII-004.2}\quad printed folio 98.\par
\noindent Printed reading: {\ttfamily\footnotesize X\allowbreak{}}\par
\noindent Corrected reading: {\ttfamily\footnotesize Y\allowbreak{}}\par
\noindent Source: {\ttfamily\footnotesize fr/sga7ii/\{diplomati\allowbreak{}c,current\}/expose\_XI\allowbreak{}II\_body.tex;\allowbreak{} folio\allowbreak{} 98;\allowbreak{} formula\allowbreak{} 664;\allowbreak{} port/decisions/decis\allowbreak{}ions\_en7II\_r2\_001.js\allowbreak{}on\#C564;\allowbreak{} fixes/FIXLOG\_7II\_CUR\allowbreak{}RENT.jsonl\#C564;\allowbreak{} 7II:150\allowbreak{}}\par

\par\medskip\noindent\hypertarget{app-E7II-XIII-005.1}{}\textbf{E7II-XIII-005.1}\quad printed folio 100.\par
\noindent Printed reading: {\ttfamily\footnotesize s\allowbreak{}}\par
\noindent Corrected reading: {\ttfamily\footnotesize \textbackslash{}eta\allowbreak{}}\par
\noindent Source: {\ttfamily\footnotesize ED-EN-PORT-WITNESS-7\allowbreak{}II-C566;\allowbreak{} fr/sga7ii/diplomatic\allowbreak{}/expose\_XIII\_body.te\allowbreak{}x;\allowbreak{} folio\allowbreak{} 100;\allowbreak{} fr/sga7ii/\{diplomati\allowbreak{}c,current\}/expose\_XI\allowbreak{}II\_body.tex;\allowbreak{} folio\allowbreak{} 100;\allowbreak{} formula\allowbreak{} 745\allowbreak{}}\par

\par\medskip\noindent\hypertarget{app-E7II-XIII-006.1}{}\textbf{E7II-XIII-006.1}\quad printed folio 100.\par
\noindent Printed reading: {\ttfamily\footnotesize geometric\allowbreak{}}\par
\noindent Corrected reading: {\ttfamily\footnotesize generic\allowbreak{}}\par
\noindent Source: {\ttfamily\footnotesize ED-EN-NONMATH-G046;\allowbreak{} F7II-XIII-013.1;\allowbreak{} 7II:153;\allowbreak{} fr/sga7ii/\{diplomati\allowbreak{}c,current\}/expose\_XI\allowbreak{}II\_body.tex;\allowbreak{} folio\allowbreak{} 100\allowbreak{}}\par

\par\medskip\noindent\hypertarget{app-E7II-XIII-006.2}{}\textbf{E7II-XIII-006.2}\quad printed folio 100.\par
\noindent Printed reading: {\ttfamily\footnotesize xr\allowbreak{}}\par
\noindent Corrected reading: {\ttfamily\footnotesize nr\allowbreak{}}\par
\noindent Source: {\ttfamily\footnotesize fr/sga7ii/\{diplomati\allowbreak{}c,current\}/expose\_XI\allowbreak{}II\_body.tex;\allowbreak{} folio\allowbreak{} 100;\allowbreak{} formula\allowbreak{} 762\allowbreak{}}\par

\par\medskip\noindent\hypertarget{app-E7II-XIII-006.3}{}\textbf{E7II-XIII-006.3}\quad printed folio 100.\par
\noindent Printed reading: {\ttfamily\footnotesize xr\allowbreak{}}\par
\noindent Corrected reading: {\ttfamily\footnotesize nr\allowbreak{}}\par
\noindent Source: {\ttfamily\footnotesize port/decisions/decis\allowbreak{}ions\_en7II\_r2\_001.js\allowbreak{}on\#C568;\allowbreak{} fixes/FIXLOG\_7II\_CUR\allowbreak{}RENT.jsonl\#C568;\allowbreak{} 7II:155\allowbreak{}}\par

\par\medskip\noindent\hypertarget{app-E7II-XIII-007.1}{}\textbf{E7II-XIII-007.1}\quad printed folio 102.\par
\noindent Printed reading: {\ttfamily\footnotesize 8\allowbreak{}}\par
\noindent Corrected reading: {\ttfamily\footnotesize 7\allowbreak{}}\par
\noindent Source: {\ttfamily\footnotesize ED-EN-PORT-WITNESS-7\allowbreak{}II-C569;\allowbreak{} fr/sga7ii/diplomatic\allowbreak{}/expose\_XIII\_body.te\allowbreak{}x;\allowbreak{} folio\allowbreak{} 102\allowbreak{}}\par

\par\medskip\noindent\hypertarget{app-E7II-XIII-007.2}{}\textbf{E7II-XIII-007.2}\quad printed folio 102.\par
\noindent Printed reading: {\ttfamily\footnotesize 1\allowbreak{}}\par
\noindent Corrected reading: {\ttfamily\footnotesize 2\allowbreak{}}\par
\noindent Source: {\ttfamily\footnotesize ED-EN-NONMATH-G018;\allowbreak{} F7II-XIII-015.1;\allowbreak{} 7II:157;\allowbreak{} fr/sga7ii/\{diplomati\allowbreak{}c,current\}/expose\_XI\allowbreak{}II\_body.tex;\allowbreak{} folio\allowbreak{} 102\allowbreak{}}\par

\par\medskip\noindent\hypertarget{app-E7II-XIII-008.1}{}\textbf{E7II-XIII-008.1}\quad printed folio 103.\par
\noindent Printed reading: {\ttfamily\footnotesize ,\allowbreak{} dash\allowbreak{}}\par
\noindent Corrected reading: \emph{(empty)}\par
\noindent Source: {\ttfamily\footnotesize port/PORT\_EN7II\_LOG.\allowbreak{}jsonl\#C570\#C570;\allowbreak{} 7II:159;\allowbreak{} fr/sga7ii/\{diplomati\allowbreak{}c,current\}/expose\_XI\allowbreak{}II\_body.tex;\allowbreak{} folio\allowbreak{} 103;\allowbreak{} formula\allowbreak{} 835\allowbreak{}}\par

\par\medskip\noindent\hypertarget{app-E7II-XIII-008.2}{}\textbf{E7II-XIII-008.2}\quad printed folio 103.\par
\noindent Printed reading: \emph{(empty)}\par
\noindent Corrected reading: {\ttfamily\footnotesize \{\textbackslash{}bar\allowbreak{} }\par
\noindent Source: {\ttfamily\footnotesize fr/sga7ii/\{diplomati\allowbreak{}c,current\}/expose\_XI\allowbreak{}II\_body.tex;\allowbreak{} folio\allowbreak{} 103\allowbreak{}}\par

\par\medskip\noindent\hypertarget{app-E7II-XIII-008.3}{}\textbf{E7II-XIII-008.3}\quad printed folio 103.\par
\noindent Printed reading: \emph{(empty)}\par
\noindent Corrected reading: {\ttfamily\footnotesize \}\allowbreak{}}\par
\noindent Source: {\ttfamily\footnotesize fr/sga7ii/\{diplomati\allowbreak{}c,current\}/expose\_XI\allowbreak{}II\_body.tex;\allowbreak{} folio\allowbreak{} 103\allowbreak{}}\par

\par\medskip\noindent\hypertarget{app-E7II-XIII-008.4}{}\textbf{E7II-XIII-008.4}\quad printed folio 103.\par
\noindent Printed reading: \emph{(empty)}\par
\noindent Corrected reading: {\ttfamily\footnotesize )\allowbreak{}}\par
\noindent Source: {\ttfamily\footnotesize fr/sga7ii/\{diplomati\allowbreak{}c,current\}/expose\_XI\allowbreak{}II\_body.tex;\allowbreak{} folio\allowbreak{} 103\allowbreak{}}\par

\par\medskip\noindent\hypertarget{app-E7II-XIII-009.1}{}\textbf{E7II-XIII-009.1}\quad printed folio 104.\par
\noindent Printed reading: {\ttfamily\footnotesize Abhyankhar\allowbreak{}}\par
\noindent Corrected reading: {\ttfamily\footnotesize Abhyankar\allowbreak{}}\par
\noindent Source: {\ttfamily\footnotesize port/decisions/decis\allowbreak{}ions\_en7II\_r2\_001.js\allowbreak{}on\#C571;\allowbreak{} fixes/FIXLOG\_7II\_CUR\allowbreak{}RENT.jsonl\#C571;\allowbreak{} 7II:162;\allowbreak{} port/decisions/decis\allowbreak{}ions\_en7II\_r2\_001.js\allowbreak{}on\#C572;\allowbreak{} fixes/FIXLOG\_7II\_CUR\allowbreak{}RENT.jsonl\#C572;\allowbreak{} 7II:163\allowbreak{}}\par

\par\medskip\noindent\hypertarget{app-E7II-XIII-010.1}{}\textbf{E7II-XIII-010.1}\quad printed folio 106.\par
\noindent Printed reading: \emph{(empty)}\par
\noindent Corrected reading: {\ttfamily\footnotesize just\allowbreak{} }\par
\noindent Source: {\ttfamily\footnotesize ED-EN-NONMATH-G019;\allowbreak{} F7II-XIII-020.1;\allowbreak{} 7II:165;\allowbreak{} fr/sga7ii/\{diplomati\allowbreak{}c,current\}/expose\_XI\allowbreak{}II\_body.tex;\allowbreak{} folio\allowbreak{} 106\allowbreak{}}\par

\par\medskip\noindent\hypertarget{app-E7II-XIII-010.2}{}\textbf{E7II-XIII-010.2}\quad printed folio 106.\par
\noindent Printed reading: {\ttfamily\footnotesize  in\allowbreak{} 2.1.11\allowbreak{}}\par
\noindent Corrected reading: \emph{(empty)}\par
\noindent Source: {\ttfamily\footnotesize ED-EN-NONMATH-G019;\allowbreak{} F7II-XIII-020.1;\allowbreak{} 7II:165;\allowbreak{} fr/sga7ii/\{diplomati\allowbreak{}c,current\}/expose\_XI\allowbreak{}II\_body.tex;\allowbreak{} folio\allowbreak{} 106\allowbreak{}}\par

\par\medskip\noindent\hypertarget{app-E7II-XIII-011.1}{}\textbf{E7II-XIII-011.1}\quad printed folio 107.\par
\noindent Printed reading: {\ttfamily\footnotesize uF\allowbreak{}}\par
\noindent Corrected reading: {\ttfamily\footnotesize u\textasciicircum{}*F\allowbreak{}}\par
\noindent Source: {\ttfamily\footnotesize fr/sga7ii/\{diplomati\allowbreak{}c,current\}/expose\_XI\allowbreak{}II\_body.tex;\allowbreak{} folio\allowbreak{} 107;\allowbreak{} formula\allowbreak{} 1008;\allowbreak{} port/decisions/decis\allowbreak{}ions\_en7II\_r2\_001.js\allowbreak{}on\#C577;\allowbreak{} fixes/FIXLOG\_7II\_CUR\allowbreak{}RENT.jsonl\#C577;\allowbreak{} 7II:168\allowbreak{}}\par

\par\medskip\noindent\hypertarget{app-E7II-XIII-012.1}{}\textbf{E7II-XIII-012.1}\quad printed folio 108.\par
\noindent Printed reading: {\ttfamily\footnotesize 2\allowbreak{}}\par
\noindent Corrected reading: {\ttfamily\footnotesize 1\allowbreak{}}\par
\noindent Source: {\ttfamily\footnotesize ED-EN-NONMATH-G002;\allowbreak{} F7II-XIII-023.1;\allowbreak{} ADOPT-F7II-XIII-023.\allowbreak{}1;\allowbreak{} fr/sga7ii/\{diplomati\allowbreak{}c,current\}/expose\_XI\allowbreak{}II\_body.tex;\allowbreak{} folio\allowbreak{} 108\allowbreak{}}\par

\par\medskip\noindent\hypertarget{app-E7II-XIII-013.1}{}\textbf{E7II-XIII-013.1}\quad printed folio 109.\par
\noindent Printed reading: {\ttfamily\footnotesize Rf\allowbreak{}}\par
\noindent Corrected reading: {\ttfamily\footnotesize Rg\allowbreak{}}\par
\noindent Source: {\ttfamily\footnotesize fr/sga7ii/\{diplomati\allowbreak{}c,current\}/expose\_XI\allowbreak{}II\_body.tex;\allowbreak{} folio\allowbreak{} 109;\allowbreak{} formula\allowbreak{} 1066\allowbreak{}}\par

\par\medskip\noindent\hypertarget{app-E7II-XIII-014.1}{}\textbf{E7II-XIII-014.1}\quad printed folio 109.\par
\noindent Printed reading: {\ttfamily\footnotesize x\allowbreak{}}\par
\noindent Corrected reading: {\ttfamily\footnotesize y\allowbreak{}}\par
\noindent Source: {\ttfamily\footnotesize ED-EN-FORMULA-G107;\allowbreak{} fr/sga7ii/\{diplomati\allowbreak{}c,current\}/expose\_XI\allowbreak{}II\_body.tex;\allowbreak{} folio\allowbreak{} 109\allowbreak{}}\par

\par\medskip\noindent\hypertarget{app-E7II-XIII-014.2}{}\textbf{E7II-XIII-014.2}\quad printed folio 110.\par
\noindent Printed reading: {\ttfamily\footnotesize 3\allowbreak{}}\par
\noindent Corrected reading: {\ttfamily\footnotesize 2\allowbreak{}}\par
\noindent Source: {\ttfamily\footnotesize ED-EN-NONMATH-G020;\allowbreak{} F7II-XIII-027.1;\allowbreak{} 7II:172;\allowbreak{} fr/sga7ii/\{diplomati\allowbreak{}c,current\}/expose\_XI\allowbreak{}II\_body.tex;\allowbreak{} folio\allowbreak{} 110\allowbreak{}}\par

\par\medskip\noindent\hypertarget{app-E7II-XIII-014.3}{}\textbf{E7II-XIII-014.3}\quad printed folio 110.\par
\noindent Printed reading: {\ttfamily\footnotesize 3\allowbreak{}}\par
\noindent Corrected reading: {\ttfamily\footnotesize 2\allowbreak{}}\par
\noindent Source: {\ttfamily\footnotesize ED-EN-NONMATH-G020;\allowbreak{} F7II-XIII-027.1;\allowbreak{} 7II:172;\allowbreak{} fr/sga7ii/\{diplomati\allowbreak{}c,current\}/expose\_XI\allowbreak{}II\_body.tex;\allowbreak{} folio\allowbreak{} 110\allowbreak{}}\par

\par\medskip\noindent\hypertarget{app-E7II-XIII-014.4}{}\textbf{E7II-XIII-014.4}\quad printed folio 110.\par
\noindent Printed reading: {\ttfamily\footnotesize x\allowbreak{}}\par
\noindent Corrected reading: {\ttfamily\footnotesize y\allowbreak{}}\par
\noindent Source: {\ttfamily\footnotesize fr/sga7ii/\{diplomati\allowbreak{}c,current\}/expose\_XI\allowbreak{}II\_body.tex;\allowbreak{} folio\allowbreak{} 110;\allowbreak{} formula\allowbreak{} 1089\allowbreak{}}\par

\par\medskip\noindent\hypertarget{app-E7II-XIII-015.1}{}\textbf{E7II-XIII-015.1}\quad printed folio 112.\par
\noindent Printed reading: {\ttfamily\footnotesize X\allowbreak{}}\par
\noindent Corrected reading: {\ttfamily\footnotesize x\allowbreak{}}\par
\noindent Source: {\ttfamily\footnotesize fr/sga7ii/\{diplomati\allowbreak{}c,current\}/expose\_XI\allowbreak{}II\_body.tex;\allowbreak{} folio\allowbreak{} 112;\allowbreak{} formula\allowbreak{} 1168;\allowbreak{} port/decisions/decis\allowbreak{}ions\_en7II\_r2\_001.js\allowbreak{}on\#C579;\allowbreak{} fixes/FIXLOG\_7II\_CUR\allowbreak{}RENT.jsonl\#C579;\allowbreak{} 7II:174\allowbreak{}}\par

\par\medskip\noindent\hypertarget{app-E7II-XIII-016.1}{}\textbf{E7II-XIII-016.1}\quad printed folio 113.\par
\noindent Printed reading: {\ttfamily\footnotesize F\allowbreak{}}\par
\noindent Corrected reading: {\ttfamily\footnotesize K\allowbreak{}}\par
\noindent Source: {\ttfamily\footnotesize fr/sga7ii/\{diplomati\allowbreak{}c,current\}/expose\_XI\allowbreak{}II\_body.tex;\allowbreak{} folio\allowbreak{} 113;\allowbreak{} formula\allowbreak{} 1226\allowbreak{}}\par

\par\medskip\noindent\hypertarget{app-E7II-XIII-016.2}{}\textbf{E7II-XIII-016.2}\quad printed folio 113.\par
\noindent Printed reading: \emph{(empty)}\par
\noindent Corrected reading: {\ttfamily\footnotesize *\allowbreak{}}\par
\noindent Source: {\ttfamily\footnotesize fr/sga7ii/\{diplomati\allowbreak{}c,current\}/expose\_XI\allowbreak{}II\_body.tex;\allowbreak{} folio\allowbreak{} 113;\allowbreak{} formula\allowbreak{} 1235;\allowbreak{} port/decisions/decis\allowbreak{}ions\_en7II\_r2\_001.js\allowbreak{}on\#C580;\allowbreak{} fixes/FIXLOG\_7II\_CUR\allowbreak{}RENT.jsonl\#C580;\allowbreak{} 7II:176\allowbreak{}}\par

\par\medskip\noindent\hypertarget{app-E7II-XIII-016.3}{}\textbf{E7II-XIII-016.3}\quad printed folio 113.\par
\noindent Printed reading: {\ttfamily\footnotesize 4\allowbreak{}}\par
\noindent Corrected reading: {\ttfamily\footnotesize 5\allowbreak{}}\par
\noindent Source: {\ttfamily\footnotesize ED-EN-NONMATH-G021;\allowbreak{} F7II-XIII-031.1;\allowbreak{} 7II:177;\allowbreak{} fr/sga7ii/\{diplomati\allowbreak{}c,current\}/expose\_XI\allowbreak{}II\_body.tex;\allowbreak{} folio\allowbreak{} 113\allowbreak{}}\par

\par\medskip\noindent\hypertarget{app-E7II-XIII-017.1}{}\textbf{E7II-XIII-017.1}\quad printed folio 114.\par
\noindent Printed reading: {\ttfamily\footnotesize 6\allowbreak{}}\par
\noindent Corrected reading: {\ttfamily\footnotesize 7\allowbreak{}}\par
\noindent Source: {\ttfamily\footnotesize port/PORT\_EN7II\_LOG.\allowbreak{}jsonl\#C581\#C581;\allowbreak{} 7II:178\allowbreak{}}\par

\par\medskip\noindent\hypertarget{app-E7II-XIII-017.2}{}\textbf{E7II-XIII-017.2}\quad printed folio 115.\par
\noindent Printed reading: {\ttfamily\footnotesize 6\allowbreak{}}\par
\noindent Corrected reading: {\ttfamily\footnotesize 7\allowbreak{}}\par
\noindent Source: {\ttfamily\footnotesize port/PORT\_EN7II\_LOG.\allowbreak{}jsonl\#C582\#C582;\allowbreak{} 7II:179\allowbreak{}}\par

\par\medskip\noindent\hypertarget{app-E7II-XIII-017.3}{}\textbf{E7II-XIII-017.3}\quad printed folio 115.\par
\noindent Printed reading: \emph{(empty)}\par
\noindent Corrected reading: {\ttfamily\footnotesize \textbackslash{}bigr)\allowbreak{}}\par
\noindent Source: {\ttfamily\footnotesize fr/sga7ii/\{diplomati\allowbreak{}c,current\}/expose\_XI\allowbreak{}II\_body.tex;\allowbreak{} folio\allowbreak{} 115;\allowbreak{} formula\allowbreak{} 1285;\allowbreak{} port/decisions/decis\allowbreak{}ions\_en7II\_r2\_001.js\allowbreak{}on\#C583;\allowbreak{} fixes/FIXLOG\_7II\_CUR\allowbreak{}RENT.jsonl\#C583;\allowbreak{} 7II:180\allowbreak{}}\par

\par\medskip\noindent\hypertarget{app-E7II-XIII-017.4}{}\textbf{E7II-XIII-017.4}\quad printed folio 115.\par
\noindent Printed reading: \emph{(empty)}\par
\noindent Corrected reading: {\ttfamily\footnotesize \textasciicircum{}*\allowbreak{}}\par
\noindent Source: {\ttfamily\footnotesize port/PORT\_EN7II\_LOG.\allowbreak{}jsonl\#C582\#C582;\allowbreak{} 7II:179;\allowbreak{} fr/sga7ii/\{diplomati\allowbreak{}c,current\}/expose\_XI\allowbreak{}II\_body.tex;\allowbreak{} folio\allowbreak{} 115;\allowbreak{} formula\allowbreak{} 1285\allowbreak{}}\par

\par\endgroup\normalsize\fi
\clearpage\SGAtoc{Exposé XIV}
\selectlanguage{english}
\clearpage
\SGApage{II}{116}

\begin{center}
{\large SGA 7 --- EXPOS\'E XIV}\par
\vspace{2em}
{\Large\underline{\underline{Comparison with Transcendental Theory}}}\par
\vspace{1.5em}
{\large by P. Deligne}
\end{center}

\medskip
\noindent\textbf{Source contents}
\begin{center}
\begin{tabular}{@{}llr@{}}
Introduction && 1\\
\S1. & Transcendental formalism of vanishing cycles & 1\\
1.1 & Fundamental group of the punctured disc: sign conventions & 2\\
1.2 & Sheaves of sets & 5\\
1.3 & The functor $R\Psi$ & 10\\
1.4 & A special case & 12\\
\S2. & The comparison theorem & 16\\
\S3. & Isolated singularities & 20\\
3.1 & The vanishing-cycle variety & 20\\
3.2 & The Picard--Lefschetz formula (transcendental case) & 31\\
\S4. & de Rham cohomology & 38
\end{tabular}
\end{center}

\medskip
\underline{Introduction}

To compare the theory of vanishing cycles in Expos\'e XIII with its
transcendental analogue over $\mathbb C$, we proceed in two stages.

\begin{enumerate}[label=(\alph*),leftmargin=2.6em]
\item We give a transcendental formalism modeled on that of
Expos\'e XIII and prove a comparison theorem by the usual methods
(already used to prove the finiteness theorem in XIII, 2.3).

\item In the case of isolated singularities, we explain in more
concrete terms the meaning of the transcendental formalism introduced
here.
\end{enumerate}

Besides the comparison theorem, this expos\'e contains a proof of the
Picard--Lefschetz formula in the transcendental case.

\SGApage{II}{\(\approx\)117}
In \S4 we give a result preliminary to the Hodge theory of vanishing
cycles.  That section will not be used later in the seminar.

\section*{1. Transcendental formalism of vanishing cycles}

\subsection*{1.1. The fundamental group of the punctured disc:
sign conventions}

We shall have to consider punctured discs $D^*$ and their fundamental
group $\mathbb Z$.  Since this group is abelian, sign ambiguities may
arise.  We eliminate them by forgetting that it is abelian.

\subsubsection*{1.1.1.}
Let $X$ be a nice connected topological space and let $x\in X$.
The elements of the fundamental group are homotopy classes of loops
based at $x$.  If $\alpha,\beta\in\pi_1(X,x)$ are represented by
loops $a,b$, then $\alpha\beta$ is represented by the loop obtained
by traversing $b$ and then $a$.

\subsubsection*{1.1.2.}
Let $F$ be a locally constant sheaf on $X$.  For every path
$a:[0,1]\longrightarrow X$, the inverse image $a^*F$ on $[0,1]$ is
locally constant, hence constant, and therefore gives a bijection
\[
a:F_{a(0)}
\xleftarrow{\ \sim\ }H^0([0,1],a^*F)
\xrightarrow{\ \sim\ }F_{a(1)}.
\]
This bijection depends only on the homotopy class of $a$.  Through
this construction, $\pi_1(X,x)$ acts on the left on $F_x$.  The
functor $F\longmapsto F_x$ is an equivalence from the category of
locally constant sheaves on $X$ to the category of sets equipped with
an action of $\pi_1(X,x)$.

\subsubsection*{1.1.3.}
As $x$ varies, the groups $\pi_1(X,x)$ are the fibers of a locally
constant sheaf of groups $\Pi_1(X)$, which acts on every locally
constant sheaf $F$.  If $\pi_1(X,x)$ is abelian, then $\Pi_1(X)$ is
constant; its value is denoted by $\pi_1(X)$.

\subsubsection*{1.1.4.}
The group $\pi_1(X,x)$ acts on the universal covering
\[
p:(\widetilde X,\widetilde x)\longrightarrow(X,x).
\]
Regarded as an \'etal\'e space over $X$, and hence as a sheaf (the
sheaf of its local sections), this covering corresponds under 1.1.2
to the set $\pi_1(X,x)$ with distinguished element $e$, on which
$\pi_1(X,x)$ acts by left translation.  The group also acts on this
$\pi_1$-set by automorphisms of $\pi_1$-sets: an element
$\sigma\in\pi_1(X,x)$ acts by
\[
g\longmapsto g\sigma^{-1}.
\]
This construction identifies
\[
\pi_1(X,x)=\operatorname{Aut}_X(\widetilde X).
\]

\SGApage{II}{\(\approx\)118}
If $F$ is locally constant on $X$, the bijection
\begin{equation}\tag{1.1.4.1}
H^0(\widetilde X,p^*F)
\longrightarrow(p^*F)_{\widetilde x}=F_x
\end{equation}
commutes with the action of $\pi_1(X,x)$: on the left the action is
transport of structure through its action on $\widetilde X$, while on
the right it is the action of 1.1.2.  The equivalence of 1.1.2 may
therefore be reformulated by saying that
\[
F\longmapsto
\bigl(\widetilde X\longmapsto
H^0(\widetilde X,p^*F)\bigr)
\]
is an equivalence from the category of locally constant sheaves on
$X$ to the category of functors
\[
(\text{universal coverings of }X)\longrightarrow(\mathrm{Sets}).
\]

\subsubsection*{1.1.5.}
If $\pi_1(X,x)$ is abelian and hence independent of $x$, the action
of 1.1.4 of $\pi_1(X)$ on $\widetilde X$ is independent of the base
point: $\pi_1(X)$ acts on every universal covering $\widetilde X$ of
$X$.

One must bear in mind that this action, which is the only one we shall
use, corresponds to the inverse of the action in 1.1.3.


\subsubsection*{1.1.6.}
Let $D$ be a disc, that is, a topological space homeomorphic to the
standard disc
\[
\{z\in\mathbb C:|z|<1\}.
\]
Let $0\in D$ and put $D^*=D-\{0\}$.  Suppose that $D$ is equipped
with the structure of an analytic manifold, and hence is oriented.
The fundamental group $\pi_1(D^*)$ is then canonically isomorphic to
$\mathbb Z$.  For the standard disc, with center $0$ and
$t\in D^*$, the positive generator of $\pi_1(D^*,t)$ is represented
by the loop
\[
x\longmapsto \exp(2\pi i x)t.
\]

\SGAvar{\subsubsection*{1.1.7.}
Let $z:D\longrightarrow\mathbb C$ be a uniformizing parameter, so
that $(dz)_0\ne0$, and let $\widetilde D^*$ be a universal covering
of $D$.  Write $\log z$ for any function on $\widetilde D^*$ such
that
\[
\exp(\log z)=z.
\]
The pair $(\widetilde D^*,\log z)$ is unique up to unique
isomorphism.

}{\subsubsection*{1.1.7.}
Let $z:D\longrightarrow\mathbb C$ be a uniformizing parameter, so
that $(dz)_0\ne0$, and let $\widetilde D^*$ be a universal covering
of $D^*$.  Write $\log z$ for any function on $\widetilde D^*$ such
that
\[
\exp(\log z)=z.
\]
The pair $(\widetilde D^*,\log z)$ is unique up to unique
isomorphism.

}{E7II-XIV-001}\SGApage{II}{\(\approx\)119}
For $\alpha\in\mathbb C$, put
\[
z^\alpha=\exp(\alpha\log z).
\]
Let $T$ be the positive generator of $\pi_1(D^*)$, acting on
$\widetilde D^*$ as in 1.1.5 and on functions on
$\widetilde D^*$ by transport of structure.  Then
\begin{equation}\tag{1.1.7.1}
T\log z=\log z+2\pi i,
\end{equation}
and
\begin{equation}\tag{1.1.7.2}
Tz^\alpha=\exp(2\pi i\alpha)z^\alpha.
\end{equation}

\subsubsection*{1.1.8.}
Let $D$ be a disc, let $0\in D$, and put $D^*=D-\{0\}$.  If
$D'\subset D$ is a smaller disc containing $0$, restriction to
$D'^{*}$ is an equivalence from the category of locally constant
sheaves on $D^*$ to the category of locally constant sheaves on
$D'^{*}$.  Indeed, for $t\in D'^{*}$,
\[
\pi_1(D'^{*},t)\xrightarrow{\ \sim\ }\pi_1(D^*,t).
\]

Suppose now that $D$ is equipped with a differentiable-manifold
structure.  Let $T$ be the tangent space to $D$ at $0$, and put
$T^*=T-\{0\}$.  For $D'\subset D$ as above and an embedding
$f:D'\hookrightarrow T$ satisfying $f(0)=0$ and $(df)_0=\mathrm{id}$,
the composite
\[
(\text{restriction to }D')^{-1}\circ f^*
\]
is an equivalence $[f]$ from the category of locally constant sheaves
on $T^*$ to the category of locally constant sheaves on $D^*$.
Up to canonical isomorphism it depends only on the germ of $f$ at
$0$, subject to $(df)_0=\mathrm{id}$.  A path from one such germ to
another defines an isomorphism of equivalences, and homotopic paths
define the same isomorphism.  Since the space of such germs is
contractible---by the retraction to a fixed germ $f_0$ given by
$(1-\lambda)f+\lambda f_0$---we obtain an equivalence, well defined
up to unique isomorphism, from locally constant sheaves on $T^*$ to
locally constant sheaves on $D^*$.

\subsubsection*{1.1.9.}
Let $D$ be a disc, let $0\in D$, put $D^*=D-\{0\}$, and let
$p:\widetilde D^*\longrightarrow D^*$ be a universal covering.  We
write
\[
p:\widetilde D\longrightarrow D
\]
for the topological space over $D$ obtained from
$\widetilde D^*$ by adjoining a point $0$, in such a way that:
\begin{enumerate}[label=(\alph*),leftmargin=2.6em]
\item $\widetilde D^*$ is open in $\widetilde D$;
\item $p:\widetilde D\longrightarrow D$ extends
  $p:\widetilde D^*\longrightarrow D^*$ and $p(0)=0$;
\item the sets $p^{-1}(U)$, as $U$ ranges over the neighborhoods of
  $0$ in $D$, form a neighborhood basis of $0$ in $\widetilde D$.
\end{enumerate}

\SGApage{II}{\(\approx\)120}
\subsubsection*{1.1.10.}
Let $X$ be a nonsingular complex algebraic curve, let $0\in X$, and
let $D$ be a disc neighborhood of $0$ in $X$.  The ring of germs at
$0$ of holomorphic functions on $D$ which are algebraic over the
field $K(X)$ of rational functions on $X$ is the henselization
$\mathcal O_0^h$ of the local ring $\mathcal O_0$ of $X$ at $0$.
Let $\widetilde D^*$ be a universal covering of
$D^*=D-\{0\}$.  The field of germs at $0\in\widetilde D$ of
meromorphic functions on $\widetilde D^*$ which are algebraic over
$K(X)$ is an algebraic closure $\overline K$ of $K(X)$, and hence
also of the fraction field $K_0^h$ of $\mathcal O_0^h$.

The group $\pi_1(D^*)$ of $D$-automorphisms of
$\widetilde D^*$ (1.1.5) acts by transport of structure on
$\overline K$, giving
\begin{equation}\tag{1.1.10.1}
\begin{tikzcd}[column sep=huge, row sep=large]
\pi_1(D^*) \arrow[r] \arrow[d, equal]
  \& \operatorname{Gal}(\overline K/K_0^h) \arrow[d, equal]
  \& {} \\
\mathbb Z
  \& \widehat{\mathbb Z}(1)
  \& .
\end{tikzcd}
\end{equation}
It follows from (1.1.7.2) that this diagram is compatible with the
natural inclusion, over $\mathbb C$, of $\mathbb Z$ into
$\widehat{\mathbb Z}(1)$.  This inclusion is the composite of
$\mathbb Z\hookrightarrow\widehat{\mathbb Z}$ and the isomorphism
\[
\widehat{\mathbb Z}\xrightarrow{\ \sim\ }
\widehat{\mathbb Z}(1),
\]
which is the inverse limit of the isomorphisms
\[
\mathbb Z/n\longrightarrow\mu_n,
\qquad
k\longmapsto\exp\left(\frac{2\pi i}{n}k\right).
\]

\subsubsection*{1.1.11.}
Every sheaf $F$ on $\Spec(K_0^h)$ defines a locally constant sheaf on
$D^*$.  We use the equivalence of 1.1.4 and, to each universal
covering $\widetilde D^*$ of $D^*$, associate $F_{\overline K}$ for
the field $\overline K$ constructed in 1.1.10.


\subsection*{1.2. Sheaves of sets}

\subsubsection*{1.2.1.}
Let $D$ be a disc (1.1.6), let $0\in D$, and put
$D^*=D-\{0\}$.  We denote by $\mathcal D^*$ the topos of locally
constant sheaves of sets on $D^*$.  By 1.1.8 this topos does not
``change'' if $D$ is replaced by a smaller disc $D'$ containing $0$,
nor, when $D$ is endowed with the structure of a differentiable
manifold, if $D$ is replaced by its tangent space at $0$.

\SGApage{II}{\(\approx\)121}
For $t\in D^*$, the functor $F\longmapsto F_t$ identifies this topos
with the topos of sets equipped with an action of $\pi_1(D^*)$
(1.1.2).  By 1.1.4, the functor
\[
F\longmapsto
\bigl(\widetilde D^*\longmapsto
H^0(\widetilde D^*,p^*F)\bigr)
\]
identifies it with the topos of functors
\[
(\text{universal coverings of }D^*)
\longrightarrow(\mathrm{Sets}),
\]
and these identifications are related by (1.1.4.1).

\subsubsection*{1.2.2.}
Let $F$ be a sheaf on $D$ whose restriction to $D^*$ is locally
constant.  There is a unique morphism of sheaves
\[
\underline{F_0}\xrightarrow{\ \alpha\ }F
\]
whose map on the stalk at $0$ is the identity.  As explained in
XIII, 0.1 (with $X_0=(P^t)$), the functor
\[
F\longmapsto(F_0,F|_{D^*},\alpha)
\]
identifies these sheaves with triples consisting of a set $F_0$, a
locally constant sheaf $F_\eta$ on $D^*$, and a morphism
\[
\alpha:
(\text{constant sheaf with value }F_0\text{ on }D^*)
\longrightarrow F_\eta.
\]

\SGAvar{\subsubsection*{1.2.3.}
We denote by $\mathcal D$ the topos of sheaves of sets on $D$ whose
restriction to $D^*$ is locally constant.  The inclusion of
$\mathcal D$ into the topos of sheaves on $D$ is the inverse-image
functor of a morphism of topoi $D\longrightarrow\mathcal D$, which
fits into the commutative diagram
\[
\begin{tikzcd}
\{0\} \arrow[r, hook] \arrow[d, equal]
  \& D \arrow[d, "\beta"]
  \& D^* \arrow[l, hook'] \arrow[d, "\beta"] \& {} \\
\{0\} \arrow[r, hook]
  \& \mathcal D
  \& \mathcal D^* \arrow[l, hook'] \& {}.
\end{tikzcd}
\]
As in 1.1.1, one checks that this topos does not change when $D$ is
replaced by a smaller disc or by the tangent space to $D$ at $0$.

}{\subsubsection*{1.2.3.}
We denote by $\mathcal D$ the topos of sheaves of sets on $D$ whose
restriction to $D^*$ is locally constant.  The inclusion of
$\mathcal D$ into the topos of sheaves on $D$ is the inverse-image
functor of a morphism of topoi $D\longrightarrow\mathcal D$, which
fits into the commutative diagram
\[
\begin{tikzcd}
\{0\} \arrow[r, hook] \arrow[d, equal]
  \& D \arrow[d, "\beta"]
  \& D^* \arrow[l, hook'] \arrow[d, "\beta"] \& {} \\
\{0\} \arrow[r, hook]
  \& \mathcal D
  \& \mathcal D^* \arrow[l, hook'] \& {}.
\end{tikzcd}
\]
As in 1.2.1, one checks that this topos does not change when $D$ is
replaced by a smaller disc or by the tangent space to $D$ at $0$.

}{E7II-XIV-002}\subsubsection*{1.2.4.}
Let $Y$ be a topological space.  The product topos
$Y\times\mathcal D$ has the following three descriptions (for their
equivalence, see XIII, 0.1).

\SGApage{II}{\(\approx\)122}
\medskip
\noindent\textit{(a)} $Y\times\mathcal D$ is the topos of sheaves
$F$ on $Y\times D$ satisfying

\smallskip
\noindent $(*)$ $D^*$ is covered by open subsets $U$ such that
$F|_{Y\times U}$ is the inverse image of a sheaf on $Y$.

\medskip
\noindent\textit{(b)} Let $t\in D^*$.  The topos
$Y\times\mathcal D$ is the topos of triples $(F_0,F_t,\alpha)$ in
which

\smallskip
\noindent $\alpha)$ $F_0$ is a sheaf on $Y$;

\noindent $\beta)$ $F_t$ is a sheaf on $Y$ equipped with an action
of $\pi_1(D^*)$;

\noindent $\gamma)$ $\alpha:F_0\longrightarrow F_t$ is a morphism
whose image is contained in the invariants.

\medskip
\noindent\textit{(c)} The topos $Y\times\mathcal D$ is the topos of
triples $(F_0,F_\eta,\alpha)$ in which

\smallskip
\noindent $\alpha')$ $F_0$ is a sheaf on $Y$;

\noindent $\beta')$ $F_\eta$ is a functor, with values denoted by
$F_{\bar\eta}$,
\[
(\text{universal coverings of }D^*)
\longrightarrow(\text{sheaves on }Y);
\]

\noindent $\gamma')$ $\alpha$ is a morphism from $F_0$ (regarded as
the constant functor with value $F_0$ from the universal coverings
of $D^*$ to the sheaves on $Y$) to $F_\eta$.

It is immaterial here whether this topos is literally the product
topos of $Y$ and $\mathcal D$.  Likewise, we define
$Y\times\mathcal D^*$, at our choice, as

\medskip
\noindent\textit{(a$^*$)} the topos of sheaves on $Y\times D^*$
satisfying $(*)$;

\noindent\textit{(b$^*$)} for $t\in D^*$, the topos of sheaves $F_t$
on $Y$ equipped with an action of $\pi_1(D^*)$;

\noindent\textit{(c$^*$)} the topos of functors
\[
(\text{universal coverings of }D^*)
\longrightarrow(\text{sheaves on }Y).
\]



\subsubsection*{1.2.5.}
The topos $Y\times\mathcal D$ has the same functoriality properties
in $Y$ as the topoi $Y\times_s S$ of Expos\'e XIII, and we shall not
list them.  To fix notation, we merely record the diagram
\[
\begin{tikzcd}[row sep=huge, column sep=huge]
Y \arrow[d, ""{name=A}]
  \& Y \arrow[d, ""{name=B}] \arrow[r, hook, "i"]
  \& Y\times\mathcal D \arrow[d]
  \& Y\times\mathcal D^* \arrow[d] \arrow[l, hook', "j"'] \\
(P^t)
  \& \{0\} \arrow[r, hook, "i"]
  \& \mathcal D
  \& \mathcal D^* \arrow[l, hook', "j"']
    \arrow[from=B, to=A, "\mathrm{sp}"']
\end{tikzcd}
\]

\SGApage{II}{\(\approx\)123}
\subsubsection*{1.2.6.}
Let $K$ be a complex of sheaves of abelian groups on
$Y\times\mathcal D$.  As in XIII, 1.4.2, one obtains from $K$ a
distinguished triangle in $K(Y\times\mathcal D^*)$:
\[
\begin{tikzcd}[row sep=large, column sep=large]
{} \& \Phi(K) \arrow[dl] \& {} \\
\mathrm{sp}^*K_0 \arrow[rr] \& {} \& K_\eta \arrow[ul]
\end{tikzcd}
\]
This construction passes to the derived category and defines a
functor
\[
D(Y\times\mathcal D)
\longrightarrow
(\text{distinguished triangles of }D(Y\times\mathcal D^*)).
\]

Suppose that $D$ is oriented (for example, complex analytic), and
let $T$ be the positive generator of $\pi_1(D^*)$ (1.1.6).  For
$K\in\operatorname{Ob}D(Y\times\mathcal D)$, we define, as in XIII,
1.4.3, the \underline{variation} $V=\operatorname{Var}(T)$,
\[
V:\Phi(K)\longrightarrow K_\eta.
\]
The endomorphism $T-\mathrm{Id}$ of $K_\eta$ is the composite
\[
T-\mathrm{Id}:K_\eta\longrightarrow\Phi(K)
\longrightarrow K_\eta.
\]



\SGAvar{\subsubsection*{1.2.7.}
Let $D$ be a disc, let $0\in D$, put $D^*=D-\{0\}$, let
$f:X\longrightarrow D$ be a continuous map, and put
$X_0=f^{-1}(0)$.  We shall define a functor
\[
\Psi:(\text{sheaves on }X)
\longrightarrow(\text{sheaves on }X_0\times\mathcal D).
\]
Let $\widetilde D^*$ be a universal covering of $D^*$, let
$\widetilde D$ be as in 1.1.9, and put
\[
\overline X=X\times_D\widetilde D,
\qquad
X^*=X-X_0=X\times_DD^*,
\qquad
\overline X^*=\overline X-X_0=X\times_D\widetilde D^*.
\]
There is a commutative diagram
\begin{equation}\tag{1.2.7.1}
\begin{tikzcd}
X_0 \arrow[r, hook, "\overline i"] \arrow[d, equal]
  \& \overline X \arrow[d, "p"]
  \&  X^* \arrow[l, hook, "\overline j"']
      \arrow[d, "p'"] \\
X_0 \arrow[r, hook, "i"]
  \& X
  \& X^* \arrow[l, hook, "j"']
\end{tikzcd}
\end{equation}
cartesian over the diagram
\begin{equation}\tag{1.2.7.2}
\begin{tikzcd}
\{0\} \arrow[r, hook, "\overline i"] \arrow[d, equal]
  \& \widetilde D \arrow[d, "p"]
  \& \widetilde D^* \arrow[l, hook, "\overline j"']
      \arrow[d, "p'"] \\
\{0\} \arrow[r, hook, "i"]
  \& D
  \& D^* \arrow[l, hook, "j"']
\end{tikzcd}
\end{equation}


}{\subsubsection*{1.2.7.}
Let $D$ be a disc, let $0\in D$, put $D^*=D-\{0\}$, let
$f:X\longrightarrow D$ be a continuous map, and put
$X_0=f^{-1}(0)$.  We shall define a functor
\[
\Psi:(\text{sheaves on }X)
\longrightarrow(\text{sheaves on }X_0\times\mathcal D).
\]
Let $\widetilde D^*$ be a universal covering of $D^*$, let
$\widetilde D$ be as in 1.1.9, and put
\[
\overline X=X\times_D\widetilde D,
\qquad
X^*=X-X_0=X\times_DD^*,
\qquad
\overline X^*=\overline X-X_0=X\times_D\widetilde D^*.
\]
There is a commutative diagram
\begin{equation}\tag{1.2.7.1}
\begin{tikzcd}
X_0 \arrow[r, hook, "\overline i"] \arrow[d, equal]
  \& \overline X \arrow[d, "p"]
  \& \overline X^* \arrow[l, hook, "\overline j"']
      \arrow[d, "p'"] \\
X_0 \arrow[r, hook, "i"]
  \& X
  \& X^* \arrow[l, hook, "j"']
\end{tikzcd}
\end{equation}
cartesian over the diagram
\begin{equation}\tag{1.2.7.2}
\begin{tikzcd}
\{0\} \arrow[r, hook, "\overline i"] \arrow[d, equal]
  \& \widetilde D \arrow[d, "p"]
  \& \widetilde D^* \arrow[l, hook, "\overline j"']
      \arrow[d, "p'"] \\
\{0\} \arrow[r, hook, "i"]
  \& D
  \& D^* \arrow[l, hook, "j"']
\end{tikzcd}
\end{equation}


}{E7II-XIV-003}\SGApage{II}{\(\approx\)124}
For a sheaf $F$ on $X$, we set (with the notation of 1.2.4(c))
\[
\Psi(F)=
\bigl(i^*F,\ \overline i^*\overline j_*((jp')^*F),\ \alpha\bigr),
\]
where
\[
\begin{split}
\alpha:i^*F=\overline i^*p^*F
&\longrightarrow
\overline i^*\overline j_*((jp')^*F)\\
&=\overline i^*\overline j_*\overline j^*p^*F
\end{split}
\]
is obtained, after applying $\overline i^*$, from the adjunction map
$\mathrm{Id}\longrightarrow\overline j_*\overline j^*$.

We also define a functor
\[
\Psi_\eta:(\text{sheaves on }X^*)
\longrightarrow(\text{sheaves on }X_0\times\mathcal D^*)
\]
by
\[
\Psi_\eta(F)=\overline i^*\overline j_*({p'}^*F).
\]

\subsubsection*{1.2.8.}
The functoriality properties of $\Psi$ are the same as those of its
analogue studied in XIII, 1.3.5--1.3.9.



\SGApage{II}{125}
\subsection*{1.3. The functor $R\Psi$}

\subsubsection*{1.3.1.}
Let $D$ be a disc, let $0\in D$, put $D^*=D-\{0\}$, let
$f:X\longrightarrow D$ be a continuous map, and put
$X_0=f^{-1}(0)$.  The functors $R\Psi$ and $R\Psi_\eta$ are the
derived functors of $\Psi$ and $\Psi_\eta$:
\[
R\Psi:D^+(X)\longrightarrow D^+(X_0\times\mathcal D),
\]
\[
R\Psi_\eta:D^+(X^*)\longrightarrow D^+(X_0\times\mathcal D^*).
\]
We denote by $R\Phi$ the composite functor $\Phi\circ R\Psi$:
\[
R\Phi:D^+(X)\longrightarrow D^+(X_0\times\mathcal D^*).
\]
The properties of these functors are exactly analogous to those of
the functors with the same names in XIII, 2.1.  Their proofs are much
more elementary, since in the topological setting the base-change
theorems for a proper morphism or for a smooth morphism (that is,
one locally isomorphic to $\mathbb R^n\times S\longrightarrow S$)
are easy to prove.  The minor changes required are as follows.

\subsubsection*{1.3.2.}
The functor $Rf_!$ is defined only for a separated morphism between
locally compact spaces, and is the derived functor of the direct-image
functor with proper supports $f_!$.

\subsubsection*{1.3.3.}
For a proper (and separated) morphism $f:X\longrightarrow D$, the
fundamental isomorphism (XIII, (2.1.8.3)) must be rewritten as follows.
Let $p:\widetilde D^*\longrightarrow D^*$ be a universal covering of
$D^*$.  For every neighborhood $U$ of $0$ in $D$, put
\[
\widetilde U^*=p^{-1}(U),
\qquad
X_{\widetilde U^*}=X\times_D\widetilde U^*.
\]
For $K\in\operatorname{Ob}D^+(f^{-1}(D^*))$, one has
\begin{equation}\tag{1.3.3.1}
\varinjlim_U\mathbb H^i(X_{\widetilde U^*},K)
\xrightarrow{\ \sim\ }
\mathbb H^i(X_0,R\Psi_{\bar\eta}(K)).
\end{equation}
This isomorphism is especially useful when
$X^*=f^{-1}(D^*)$ is a topological fiber bundle over $D^*$.  As a
disc $U\subset D$ shrinks to $0$, the Leray spectral sequence for
$X_{\widetilde U^*}\longrightarrow\widetilde U^*$ then shows that,
for every locally constant sheaf $\Lambda$ on $X^*$ and every
$t\in\widetilde U^*$,
\[
H^i(X_{\widetilde U^*},\Lambda)
\xrightarrow{\ \sim\ }H^i(X_t,\Lambda).
\]
Thus, for $t\in D^*$,
\begin{equation}\tag{1.3.3.2}
H^i(X_t,\Lambda)
\xrightarrow{\ \sim\ }
H^i(X_0,R\Psi_t(\Lambda)),
\end{equation}
provided that $X/D$ is proper, $X^*$ is a fiber bundle, and
$\Lambda$ is locally constant.

\SGApage{II}{\(\approx\)126}
In practice, to verify that $X^*$ is a fiber bundle over $D^*$ one
uses Thom's isotopy theorem, one form of which is stated below.  For
the proof of this theorem, none of whose consequences will be used in
the subsequent expos\'es, we refer to Whitney and Mather.



\subsubsection*{1.3.4. Preliminary.}
Let $f:X\longrightarrow Y$ be a smooth morphism (by definition, one
with no critical point) of differentiable manifolds.  A
\underline{connection} $\nabla$ on $X/Y$ is a vector subbundle of the
tangent bundle of $X$ complementary to $\ker(df)$.  For a vector
field $u$ on $Y$, there is a unique vector field $\nabla(u)$ on $X$
which lifts $u$ and lies in this subbundle.  The connection is called
integrable if
\[
\nabla([u,v])=[\nabla(u),\nabla(v)].
\]
When $f$ is proper, an integrable connection gives local
trivializations over $Y$ of $X/Y$,
\[
\alpha:f^{-1}(U)\xrightarrow{\ \sim\ }X_0\times U,
\]
compatible with projection to $U$ and satisfying
$d\alpha(\nabla(u))=(0,u)$.  When $f$ is not proper, the same result
holds under suitable conditions at infinity; for example, it is
enough that $\nabla(u)g$ vanish for a suitable exhaustion function
$g$.

\subsubsection*{1.3.5. Thom's isotopy theorem.}
\textit{Let
\[
D=\{z\in\mathbb C:|z|<1\},
\qquad D^*=D-\{0\},
\]
and let $f:X\longrightarrow D$ be a proper morphism of analytic
spaces.  After replacing $D$ by a smaller disc and $X$ by its inverse
image, there exists a stratification $\Sigma$ of
$X^*=f^{-1}(D^*)$ having the following properties:}

\begin{enumerate}[label=(\alph*),leftmargin=2.7em]
\item \textit{$\Sigma$ is a filtration
\[
X^*=F_0\supset F_1\supset\cdots\supset F_n
\supset F_{n+1}=\varnothing
\]
of $X^*$ by closed subsets such that the closure $\overline F_i$ of
$F_i$ in $X$ is an analytic subspace of $X$.}

\item \textit{The open strata $F_i-F_{i+1}$ are smooth over $D^*$.}

\item \textit{$\Sigma$ satisfies Whitney's conditions (a) and (b).}
\end{enumerate}

\SGApage{II}{\(\approx\)127}
\textit{For such a stratification $\Sigma$, there are integrable
connections $\nabla_i$ on the $F_i-F_{i+1}$ such that}

\begin{enumerate}[label=(\roman*),leftmargin=2.7em]
\item \textit{locally on $D^*$, $\nabla_i$ defines a trivialization
of $F_i-F_{i+1}$;}

\item \textit{these trivializations glue to give local
trivializations of $X^*/D^*$.}
\end{enumerate}

\textit{Finally, if $(Y_\alpha)$ is a finite family of closed
analytic subspaces of $X$, then, after shrinking $D$ further, one can
choose $\Sigma$ so that every $Y_\alpha$ is a union of connected
components of open strata $(F_i-F_{i+1})$.  A connection $\nabla$ as
above then locally trivializes $(X,(Y_\alpha))$.}

\subsubsection*{1.3.6.}
The isomorphism (1.3.3.2) remains valid for hypercohomology with
values in a complex of sheaves $K$ when $X$ is proper over $D$ and
the following condition holds:

\medskip
\noindent $(*)$ For every $i\in\mathbb Z$, the space $D^*$ is
covered by open subsets $U$ for which there exist an isomorphism
compatible with projection to $U$,
\[
\varphi:f^{-1}(U)\xrightarrow{\ \sim\ }Y\times U,
\]
and a sheaf $G$ on $Y$ such that, on $f^{-1}(U)$,
$\underline H^i(K)$ is isomorphic to $\varphi^*\operatorname{pr}_1^*G$.

\SGAvar{Under this hypothesis, the triangle of 1.2.1 gives a long exact
sequence
\begin{equation}\tag{1.3.6.1}
\cdots\longrightarrow
\mathbb H^i(X_0,K)\longrightarrow
\mathbb H^i(X_t,K)\longrightarrow
\mathbb H^i(X_0,R\Phi_t(K))\longrightarrow\cdots.
\end{equation}


}{Under this hypothesis, the triangle of 1.2.6 gives a long exact
sequence
\begin{equation}\tag{1.3.6.1}
\cdots\longrightarrow
\mathbb H^i(X_0,K)\longrightarrow
\mathbb H^i(X_t,K)\longrightarrow
\mathbb H^i(X_0,R\Phi_t(K))\longrightarrow\cdots.
\end{equation}


}{E7II-XIV-004}By Thom's theorem, condition $(*)$ holds (after shrinking $D$) whenever
$X$ admits a complex analytic stratification such that all the
$\underline H^i(K)$ are locally constant on every open stratum.



\subsection*{1.4. A special case}

This section will not be used later in the seminar.  It formalizes a
construction sketched in the Introduction to Expos\'e XIII.

\SGApage{II}{\(\approx\)128}
\subsubsection*{1.4.1.}
Let $D$ be an oriented disc, let $0\in D$, and put
$D^*=D-\{0\}$.  A \underline{topological space} $\mathfrak X^*$ over
$\mathfrak D^*$ is a functor
\[
(\text{universal coverings of }D^*)
\longrightarrow(\text{topological spaces}).
\]
After a universal covering $\widetilde D^*$ has been chosen, this is
a topological space $X_\eta$ equipped with an action of
$\pi_1(D^*)$; when $\widetilde D^*$ is the universal covering of
$(D^*,t)$, we denote it by $\mathfrak X_t^*$.

A \underline{sheaf} on $\mathfrak X^*$ is a functor
\[
(\text{universal coverings of }D^*)
\longrightarrow
(\text{topological spaces equipped with a sheaf}),
\]
or, equivalently, a $\pi_1(D^*)$-equivariant sheaf on $X_\eta$.

To $\mathfrak X^*$ we associate the topological space
\[
\beta^*\mathfrak X^*
=X_\eta\times\widetilde D^*/\pi_1(D^*)
\]
over $D^*$.  Every sheaf $F$ on $\mathfrak X^*$ descends to a sheaf
$\beta^*F$ on $\beta^*\mathfrak X^*$, giving a commutative diagram of
topoi
\[
\begin{tikzcd}
\beta^*\mathfrak X^* \arrow[r] \arrow[d]
  \& \mathfrak X^* \arrow[d] \\
D^* \arrow[r, "\beta"]
  \& \mathfrak D^*.
\end{tikzcd}
\]

\subsubsection*{1.4.2.}
A \underline{topological space} $\mathfrak X$ over $\mathfrak D$
consists of
\begin{enumerate}[label=(\alph*),leftmargin=2.7em]
\item a topological space $\mathfrak X^*$ over $\mathfrak D^*$;
\item a topological space $X_0$;
\item a $\pi_1(D^*)$-invariant morphism
\[
\mathrm{sp}:X_\eta\longrightarrow X_0,
\]
where $\pi_1(D^*)$ acts trivially on $X_0$.
\end{enumerate}

A \underline{sheaf} on $\mathfrak X$ consists of a sheaf $F_0$ on
$X_0$, a sheaf on $\mathfrak X^*$ given by $F_\eta$, and an equivariant
morphism
\[
\mathrm{sp}:\mathrm{sp}^*F_0\longrightarrow F_\eta
\qquad
(\text{or }F_0\longrightarrow\mathrm{sp}_*F_\eta).
\]
There is a commutative diagram
\[
\begin{tikzcd}
X_0 \arrow[r, hook] \arrow[d]
  \& \mathfrak X \arrow[d]
  \& \mathfrak X^* \arrow[l, hook] \arrow[d] \\
\{0\} \arrow[r, hook]
  \& \mathfrak D
  \& \mathfrak D^* \arrow[l, hook]
\end{tikzcd}
\]
and, in addition, a morphism of topoi
\[
\Psi:\mathfrak X\longrightarrow X_0\times\mathfrak D.
\]
It is obtained from the functoriality of the preceding construction,
applied to
\[
(X_\eta,X_0,\mathrm{sp})
\longrightarrow(X_0,X_0,\mathrm{Id}),
\]
and satisfies
\[
\Psi(F_0,F_\eta,\mathrm{sp})
=(F_0,\mathrm{sp}_*F_\eta,\mathrm{sp}).
\]
This morphism induces
\[
\Psi_\eta:\mathfrak X^*\longrightarrow X_0\times\mathfrak D^*.
\]



\SGApage{II}{\(\approx\)129}
\subsubsection*{1.4.3.}
To a topological space $\mathfrak X$ over $\mathfrak D$ we associate a
topological space $X=\beta^*\mathfrak X$ over $D$ as follows.  To
$X^*=\beta^*\mathfrak X^*$ we adjoin the fiber $X_0$ over $0$.  A
fundamental system of neighborhoods of $x_0\in X_0$ is formed by the
sets
\[
V\cup
\bigl(\mathrm{sp}^{-1}(V)\times\widetilde U^*/\pi_1(D^*)\bigr),
\]
where $V$ is a neighborhood of $x_0$ in $X_0$, $U$ is a neighborhood
of $0$ in $D$, and
$\widetilde U^*=U\times_D\widetilde D^*$.  Let $j$ (respectively $i$)
denote the inclusion of $X^*$ (respectively $X_0$) into $X$.  For a
sheaf $F$ on $\mathfrak X^*$, one has
\[
i^*j_*\beta^*F
=(\mathrm{sp}_*F_\eta)^{\pi_1(D^*)}.
\]
It follows that there is a commutative diagram of topoi
\[
\begin{tikzcd}
X \arrow[r, "\beta"] \arrow[d]
  \& \mathfrak X \arrow[d] \\
D \arrow[r, "\beta"]
  \& \mathfrak D
\end{tikzcd}
\]
such that, for every sheaf $F$ on $\mathfrak X$, there is an
isomorphism
\begin{equation}\tag{1.4.3.1}
\Psi(F)\xrightarrow{\ \sim\ }\Psi(\beta^*F).
\end{equation}

\subsubsection*{Proposition 1.4.4.}
The map (1.4.3.1) derives to an isomorphism
\[
R\Psi\xrightarrow{\ \sim\ }R\Psi\circ\beta^*
\]
between functors from $D^+(\mathfrak X)$ to
$D^+(X_0\times\mathfrak D)$.

\SGApage{II}{\(\approx\)130}
Indeed, let $\overline j$ be the inclusion of
\[
\overline X^*=X\times_D\widetilde D^*
\simeq X_\eta\times\widetilde D^*
\]
into $\overline X=X\times_D\widetilde D$.  The proposition reduces to
checking that, for every flabby sheaf of abelian groups $F$ on
$X_\eta$, the sheaf $\operatorname{pr}_1^*F$ is acyclic for
$\overline j_*$.  This in turn reduces (with $B=\widetilde D^*$) to
the following familiar assertion, proved by a homotopy argument.  If
$F$ is a sheaf of abelian groups on a space $Y$ and $B$ is a disc,
then
\[
H^i(Y,F)\xrightarrow{\ \sim\ }
H^i(Y\times B,\operatorname{pr}_1^*F).
\]

\SGApage{II}{129}
\subsubsection*{1.4.5.}
Proposition 1.4.4 gives a method for computing $R\Psi$ for sheaves on
$X$ of the form $\beta^*F$, reducing the computation of $R\Psi$
essentially to that of $R\,\mathrm{sp}_*$.



\SGApage{II}{\(\approx\)130}
\subsubsection*{1.4.6.}
Later we shall consider the case in which $X_\eta$ is locally compact
and $X_0$ is obtained from $X_\eta$ by contracting a compact subset
$V\subset X_\eta$ to a point $0$, with $\mathrm{sp}$ the contraction
map and the action of $\pi_1(D^*)$ trivial outside $V$.  Proper base
change for $\mathrm{sp}$ then gives, for
$K\in\operatorname{Ob}D^+(\mathfrak X^*)$,
\begin{equation}\tag{1.4.6.1}
R\Psi_\eta(\beta^*K)_0=R\Gamma(V,K)
\end{equation}
and
\begin{equation}\tag{1.4.6.2}
R\Gamma_{\{0\}}R\Psi_\eta(\beta^*K)
=R\Gamma_V(X_\eta,K),
\end{equation}
the latter being cohomology with supports.

For $K$ on $\mathfrak X$, suppose that
\[
\mathrm{sp}:\mathrm{sp}^*K_0\longrightarrow K_\eta
\]
is a quasi-isomorphism outside $V$.  Then $R\Phi(K)=0$ outside $0$.
In particular, for a constant sheaf $\Lambda$ on $\mathcal X$,
\[
R^i\Phi(\Lambda)=\widetilde H^i(V,\Lambda),
\qquad\text{supported at }0,
\]
where the right-hand side is reduced cohomology.

\SGApage{II}{\(\approx\)131}
\subsubsection*{1.4.7.}
The case of interest to us is that in which $X_\eta$ is a
differentiable manifold and $V\subset X_\eta$ is a submanifold with
boundary and interior $V^\circ$.  Then $X-\{0\}$ is a differentiable
manifold obtained by gluing $\beta^*X_\eta$ to the constant bundle
over $D$ with fiber $X_\eta-V$.  The map
\begin{equation}\tag{1.4.7.1}
H_c^i(V^\circ,\Lambda)
\xrightarrow{\ \sim\ }H_V^i(X_\eta,\Lambda)
\end{equation}
is an isomorphism.  These groups are also denoted by
$H^i(V\bmod\partial V,\Lambda)$.  The isomorphism (1.4.7.1) defines
an isomorphism
\begin{equation}\tag{1.4.7.2}
R\Gamma_c(V^\circ,\Lambda)
\xrightarrow{\ \sim\ }
R\Gamma_{\{0\}}R\Psi_\eta(\Lambda),
\end{equation}
and the \underline{variation} is defined, by the procedure of XIII,
2.4.5 and through (1.4.7.2), as a morphism whose source is the reduced
cohomology of $V$:
\begin{equation}\tag{1.4.7.3}
\operatorname{Var}:
R\widetilde\Gamma(V,\Lambda)
\longrightarrow R\Gamma_c(V^\circ,\Lambda).
\end{equation}

The automorphism induced by $T$ on the pair $(V,\partial V)$ is again
denoted by $T$, and its restriction $\partial T$ to $\partial V$ is
the identity.  The variation (1.4.7.3) should coincide with the
classical variation morphism (3.1.9) attached to such an automorphism
$T$.  A statement of this kind will be proved in 3.1.10.



\section*{2. The comparison theorem}

\subsubsection*{2.1.}
For every scheme $X$ of finite type over $\mathbb C$, we denote by
$X_{\mathrm{cl}}$ the space $X(\mathbb C)$ endowed with its usual
topology.

Let $S$ be a smooth curve over $\mathbb C$, that is, a smooth
separated scheme of finite type over $\mathbb C$, pure of dimension
$1$.  Let $0\in S(\mathbb C)$, let $f:X\longrightarrow S$ be a
separated $S$-scheme of finite type, and let $\Lambda$ be a
noetherian torsion ring (for example $\mathbb Z/n\mathbb Z$).  Put
$X_0=f^{-1}(0)$.

\SGApage{II}{\(\approx\)132}
\subsubsection*{2.2.}
We denote by $S^h$ the henselized henselian trait of $S$ at $0$, and
by $\eta$ the generic point of $S^h$.  There is a functor
\[
\begin{aligned}
\Psi:(\text{sheaves on }X)
&\xrightarrow{\ \mathrm{restriction}\ }
(\text{sheaves on }X\times_SS^h)\\
&\xrightarrow{\ \Psi\ }
(\text{sheaves on }X_0\times S^h).
\end{aligned}
\]
Its derived functor is the composite
\[
R\Psi:D(X,\Lambda)
\longrightarrow D(X\times_SS^h,\Lambda)
\xrightarrow{\ R\Psi\ }D(X_0\times S^h,\Lambda).
\]

\subsubsection*{2.3.}
Let $D\subset S_{\mathrm{cl}}$ be a disc containing $0$.  The
composite functor $\Psi_{\mathrm{cl}}$ (or simply $\Psi$)
\[
\begin{aligned}
\Psi_{\mathrm{cl}}:(\text{sheaves on }X_{\mathrm{cl}})
&\xrightarrow{\ \mathrm{restriction}\ }
(\text{sheaves on }X_{\mathrm{cl}}\times_{S_{\mathrm{cl}}}D)\\
&\xrightarrow{\ \Psi\ }
(\text{sheaves on }(X_0)_{\mathrm{cl}}\times\mathcal D)
\end{aligned}
\]
is independent, up to unique isomorphism, of the choice of $D$;
recall that $\mathcal D$ does not change when $D$ is replaced by a
smaller disc.  Its derived functor is the composite
\[
R\Psi_{\mathrm{cl}}:D(X_{\mathrm{cl}},\Lambda)
\longrightarrow D(f_{\mathrm{cl}}^{-1}(D),\Lambda)
\xrightarrow{\ R\Psi\ }
D((X_0)_{\mathrm{cl}}\times\mathcal D,\Lambda).
\]

\subsubsection*{2.4.}
Likewise, we define
\[
\begin{aligned}
R\Psi_\eta &:D(X-X_0,\Lambda)
  \longrightarrow D(X_0\times\eta,\Lambda),\\
R\Psi_{\eta\,\mathrm{cl}} &:D((X-X_0)_{\mathrm{cl}},\Lambda)
  \longrightarrow D((X_0)_{\mathrm{cl}}\times\mathcal D^*,\Lambda).
\end{aligned}
\]

\subsubsection*{2.5.}
The construction of 1.1.11 defines a morphism of topoi from
$\mathcal D^*$ to $\eta$, which extends to
\[
\varepsilon:\mathcal D\longrightarrow S^h.
\]
Together with
$\varepsilon:(X_0)_{\mathrm{cl}}\longrightarrow X_0$ (SGA 4, XI,
4), this morphism defines
\[
\varepsilon:(X_0)_{\mathrm{cl}}\times\mathcal D
\longrightarrow X_0\times S^h.
\]
If $\widetilde D^*$ is a universal covering of $D^*$ and
$k(\bar\eta)$ is the corresponding algebraic closure (1.1.10) of
$k(\eta)$, then
\[
\varepsilon^*(F_0,F_\eta,\alpha)
=\bigl(\varepsilon^*F_0,\varepsilon^*F_{\bar\eta},
\varepsilon^*\alpha\bigr).
\]



\SGApage{II}{\(\approx\)133}
\subsubsection*{2.6.}
Let $\widetilde D^*$ be a universal covering of $D^*$, let
$k(\bar\eta)$ be the corresponding algebraic closure of $k(\eta)$,
and let $\overline S$ be the normalization of $S^h$ in
$k(\bar\eta)$.  Write $\overline S$ as an inverse limit of affine
$S$-schemes $S_i$.  For every $i$, after replacing $D$ by a smaller
disc $D_i$, there is a natural morphism of $S_{\mathrm{cl}}$-spaces
\[
a_i:\widetilde D_i\longrightarrow(S_i)_{\mathrm{cl}}.
\]
Let $j_i:U_i\hookrightarrow S_i$ be the open complement in $S_i$ of
the image $0$ of the closed point of $\overline S$, and let
$i:0\hookrightarrow S_i$.

Let $F$ be a sheaf on $X$, and put $F_{\mathrm{cl}}=\varepsilon^*F$.
We again denote by $j_i$ the inclusion of $X\times_SU_i$ into
$X\times_SS_i$, by $i$ the inclusion of $X_0$ into
$X\times_SS_i$, and by $F$ the inverse image of $F$ on
$X\times_SS_i$.  Then
\[
\Psi(F)=
\left(i^*F,\ \varinjlim_i i^*j_{i*}j_i^*F,\
\varinjlim_i\varphi_i\right).
\]

The functors $\varepsilon^*$ and inverse image commute.  The morphism
\[
\varepsilon^*j_{i*}\longrightarrow(j_{i,\mathrm{cl}})_*\varepsilon^*
\]
therefore gives
\[
c_i:
(i^*j_{i*}j_i^*F)_{\mathrm{cl}}
\longrightarrow i^*j_{i*}j_i^*(F_{\mathrm{cl}}).
\]

Let $\overline j$ (respectively $\overline j_i$) be the inclusion of
$\widetilde D^*$ (respectively $\widetilde D_i^*$) into
$\widetilde D$ (respectively $\widetilde D_i$), as well as the
induced inclusions
\[
X\times_S\widetilde D^*\hookrightarrow X\times_S\widetilde D,
\qquad
X\times_S\widetilde D_i^*\hookrightarrow X\times_S\widetilde D_i.
\]
The adjunction map $\mathrm{Id}\longrightarrow a_{i*}a_i^*$ defines
\[
i^*j_{i*}j_i^*F_{\mathrm{cl}}
\longrightarrow
i^*\overline j_{i*}\overline j_i^*\overline F_{\mathrm{cl}}
=i^*\overline j_*\overline j^*\overline F_{\mathrm{cl}}.
\]
Passing to the limit gives
\[
c_\eta:\Psi_{\bar\eta}(F)_{\mathrm{cl}}
\longrightarrow\Psi_{\mathrm{cl}}(F_{\mathrm{cl}}).
\]
This morphism defines the \underline{comparison morphism}
\begin{equation}\tag{2.6.1}
c:\varepsilon^*\Psi(F)
\longrightarrow\Psi_{\mathrm{cl}}(\varepsilon^*F).
\end{equation}
Similarly, one defines
\begin{equation}\tag{2.6.2}
c_\eta:\varepsilon^*\Psi_\eta(F)
\longrightarrow\Psi_\eta(\varepsilon^*F).
\end{equation}



\SGApage{II}{\(\approx\)134}
\subsubsection*{2.7.}
The morphisms (2.6.1) and (2.6.2) derive to
\[
c:\varepsilon^*R\Psi
\longrightarrow R\Psi_{\mathrm{cl}}\circ\varepsilon^*
\]
and
\[
c_\eta:\varepsilon^*R\Psi_\eta
\longrightarrow R\Psi_{\eta\,\mathrm{cl}}\circ\varepsilon^*.
\]

\subsubsection*{Theorem 2.8.}
Let $K\in\operatorname{Ob}D(X,\Lambda)$ (respectively,
$K\in\operatorname{Ob}D(X-X_0,\Lambda)$) be a complex of sheaves of
$\Lambda$-modules whose cohomology sheaves $\underline H^i(K)$ are
constructible.  Then the comparison morphism
\[
c:\varepsilon^*R\Psi(K)
\longrightarrow R\Psi_{\mathrm{cl}}(\varepsilon^*K)
\]
(respectively,
\[
c:\varepsilon^*R\Psi_\eta(K)
\longrightarrow R\Psi_{\eta\,\mathrm{cl}}(\varepsilon^*K))
\]
is an isomorphism.

It suffices to prove the assertion for $R\Psi_\eta$.  Moreover, both
sides of 2.7 are triangulated functors of finite cohomological
dimension, so it is enough to treat the case in which $K$ is a single
sheaf $F$ placed in degree $0$.  The proof is exactly analogous to
that of the finiteness theorem XIII, 2.3.1.  Since $R\Psi$ and
$R\Psi_{\mathrm{cl}}$ have the same variance properties, the
d\'evissages used there (with permission to shrink $S$) reduce the
problem to the case in which $X$ is regular, $f^{-1}(0)$ is a normal
crossings divisor in $X$, and $F=\mathbb Z/\ell\mathbb Z$.  One then
compares Expos\'e I, 3.3 with its transcendental analogue.



\SGApage{II}{\(\approx\)135}
\section*{3. Isolated singularities}

\subsection*{3.1. The manifold of vanishing cycles}

\subsubsection*{3.1.1.}
Let
\[
D=\{z\in\mathbb C\mid |z|<1\},
\]
let $X$ be an analytic space with $0\in X$, and let
$f:X\to D$ be a morphism such that $f(0)=0$.  We suppose that $f$ is
smooth away from $0$ (in particular, $X$ is smooth away from $0$).
We intend to study the vanishing cycles near $0$; hence there is no
loss of generality in supposing that $X$ is embedded in
$\mathbb C^n\times D\subset\mathbb C^{n+1}$ compatibly with the
projection to $D$, and in such a way that $0$ maps to $0$.

\subsubsection*{3.1.2.}
Let $\delta(x)$ be a real-analytic function on $X$ (the restriction to
$X$ of a real-analytic function on $\mathbb C^{n+1}$) such that, for
$x\in X_0=f^{-1}(0)$,

\begin{enumerate}[label=(\alph*)]
\item $\delta(x)\geq 0$;
\item $\delta(x)=0$ if and only if $x=0$.
\end{enumerate}

For example, we may take $\delta$ to be the square of the Euclidean
distance from $x$ to $0$ in $\mathbb C^{n+1}$.  The critical locus of
the restriction of $\delta$ to $X_0$ is real analytic and therefore,
in a suitable neighborhood of $0$, has only finitely many connected
components.  Replacing $X$ by an open neighborhood of $0$ and applying
Sard's theorem, we may consequently obtain the following situation:

\begin{enumerate}[label=(\alph*),start=3]
\item for a suitable $r>0$, the space
  $X_0\cap\delta^{-1}([0,r])$ is compact;
\item on $X_0$, $\delta$ has no critical value in $(0,r]$.
\end{enumerate}

After shrinking $X$ further, we then have:

\begin{enumerate}[label=(\alph*),start=5]
\item there exists a $C^\infty$ function $\tau(s)$, $0\leq s\leq r$,
  with $\tau(0)=0$ and $\tau'(s)>0$ for $s>0$, such that, for
  $0<s\leq r$ and $0\leq |t|^2\leq\tau(s)$, the fiber
  $X_t=f^{-1}(t)$ is transverse to the real hypersurface
  $\delta(x)=s$ in $X$, and the box
  \[
  B_{t,d}=\{x\in X\mid \delta(x)\leq d\ \text{and}\ |f(x)|^2<t\}
  \]
  is proper over the open disc $D_t$ of radius $t^{1/2}$.
\end{enumerate}

\SGApage{II}{\(\approx\)136}
Let $L=X_0\cap\delta^{-1}(r)$.  By (c) and (d), there is a
homeomorphism from $X_0\cap\delta^{-1}([0,r])$ to the cone
\[
L\times[0,r]\big/\bigl(L\times\{0\}\text{ contracted to a point}\bigr)
\]
which transforms $\delta$ into the second coordinate and is
$C^\infty$ away from the origin.

\SGAvar{\subsubsection*{3.1.3.}
Let $d\leq r$ and $t\leq r(d)$.  Away from $0$, $B_{s,t}$ is a
manifold with corners, and the projection $f:B_{t,d}\to D_t$ is a
fibration.  In particular:

}{\subsubsection*{3.1.3.}
Let $d\leq r$ and $t\leq\tau(d)$.  Away from $0$, $B_{t,d}$ is a
manifold with corners, and the projection $f:B_{t,d}\to D_t$ is a
fibration.  In particular:

}{E7II-XIV-005}\begin{enumerate}[label=(\alph*)]
\item over $D_t^*=D_t-\{0\}$, $B_{t,d}$ induces a $C^\infty$ bundle
  of compact manifolds with boundary;
\item the boundary $\delta(x)=d$ of this bundle extends to a
  $C^\infty$ bundle $\partial B_{t,d}$ of compact manifolds over
  $D_t$.
\end{enumerate}

Let $y\in D_t^*$.  We call the manifold with boundary
$V=f^{-1}(y)$ the \underline{manifold of vanishing cycles}.  Since the
disc is contractible, the bundle $\partial B_{t,d}$ over $D_t$ is
isomorphic to the constant bundle with fiber $\partial V$, and the
space of bundle isomorphisms
\[
\alpha:\partial B_{t,d}\xrightarrow{\ \sim\ }\partial V\times D_t
\]
that restrict to the identity on the fiber over $y$ is contractible.
Choose one, denoted by $\alpha$.  Lift the vector field
$\partial/\partial\theta$ on $S_{|y|}$, the circle of radius $|y|$,
to a vector field $u$ on $f^{-1}(S_{|y|})$ which is compatible with
$\alpha$ on the boundary, so that
\[
d\alpha(u)=(0,\partial/\partial\theta).
\]
The space of such lifts is contractible.  Integrating $u$ once around
the circle in the positive direction defines a diffeomorphism $T$ of
$V$ onto itself which is the identity on $\partial V$; this is the
\underline{monodromy transformation}.  It is well defined up to
isotopy relative to $\partial V$.

\SGAvar{In a sense whose precise formulation is left to the reader, $(V,T)$
is independent of the choices of $t$ and $s$.  I do not know whether
it is independent of the choice of $\delta$.  This is nevertheless
the case if we restrict $\delta$ to be the restriction to $X$ of a
real-analytic function $g$ on $\mathbb C^{n+1}$ such that, on the
Zariski tangent space of $X_0$ at $0$, one has $dg=0$ and $d^2g>0$
(positive definite).  Indeed, for $\delta$ and $\delta'$ of this
kind, $r$ can be chosen as in 3.1.2(c),(d), uniformly for every
$\lambda\delta+(1-\lambda)\delta'$, $0\leq\lambda\leq1$.



}{In a sense whose precise formulation is left to the reader, $(V,T)$
is independent of the choices of $t$ and $d$.  I do not know whether
it is independent of the choice of $\delta$.  This is nevertheless
the case if we restrict $\delta$ to be the restriction to $X$ of a
real-analytic function $g$ on $\mathbb C^{n+1}$ such that, on the
Zariski tangent space of $X_0$ at $0$, one has $dg=0$ and $d^2g>0$
(positive definite).  Indeed, for $\delta$ and $\delta'$ of this
kind, $r$ can be chosen as in 3.1.2(c),(d), uniformly for every
$\lambda\delta+(1-\lambda)\delta'$, $0\leq\lambda\leq1$.



}{E7II-XIV-006}\subsubsection*{3.1.4.}
Let $(V_1,T_1)$ be obtained from $(V,T)$ by adjoining a collar:
\[
V_1-V^\circ\simeq\partial V\times[0,1],
\qquad
T_1|_{V_1-V^\circ}=\operatorname{id}.
\]
Define the following space over $D_t$ (cf. 1.4.3):
\[
f':W\longrightarrow D_t.
\]

\SGApage{II}{\(\approx\)137}
\begin{enumerate}[label=(\alph*)]
\item Let $\widetilde D_t^*$ be a universal covering of $D_t^*$ with
  base point $y$, and let $W^*$ be the bundle over $D_t^*$ obtained as
  the quotient of $V_1\times\widetilde D_t^*$ by the equivalence
  relation generated by
  \[
  (v,y)\sim(T_1v,e^{2\pi i}y).
  \]
\item Since $T_1|_{V_1-V^\circ}$ is trivial, $W^*$ can be glued to
  the trivial bundle over $D_t$ with fiber $V_1-V$, giving a space
  $W_1$ over $D_t$.
\item Put $W=W_1\cup\{0\}$ and $f'(0)=0$.  Give $W$ the unique
  topology inducing that of $W_1$, for which $f'$ is continuous and,
  for $t'<t$, ${f'}^{-1}(\overline D_{t'})$ is the one-point
  compactification of
  $W_1\cap {f'}^{-1}(\overline D_{t'})$.
\end{enumerate}

Equip $W-\{0\}$ with the integrable connection of 1.3.4 for which
the trivialization, supplied by (a), of the inverse image of $W$ on
$\widetilde D^*$ is horizontal.

\SGAvar{\subsubsection*{Proposition 3.1.5.}
There exists a homeomorphism
\[
\Psi:B_{s,t}\xrightarrow{\ \sim\ }W
\]
compatible with the projections to $D_t$ and carrying $0$ to $0$;
this homeomorphism may be chosen to be a diffeomorphism away from
$0$.

}{\subsubsection*{Proposition 3.1.5.}
There exists a homeomorphism
\[
\Psi:B_{t,d}\xrightarrow{\ \sim\ }W
\]
compatible with the projections to $D_t$ and carrying $0$ to $0$;
this homeomorphism may be chosen to be a diffeomorphism away from
$0$.

}{E7II-XIV-007}This result will not be used in the remainder of the seminar.

\SGAvar{Let $d_0<d_1\leq r$, let $t_0\leq\tau(d_0)$, and let
$y\in D_{t_0}^*$.  We have
\[
B_{t_0,d_0}\subset B_{t_0,d_1},
\]
and $B_{t_0,d_1}-B_{t_0,d_0}^{\circ}$ is a bundle over $D_{|t|_0}$.
Trivialize it and, over the circle $S_{|y|}$, lift
$\partial/\partial\theta$ compatibly with this trivialization.  Let
$T_1$ be the corresponding endomorphism of
\[
V_1=f^{-1}(y)\cap B_{t_0,d_1},
\qquad
V_0=f^{-1}(y)\cap B_{t_0,d_0}.
\]
Then $(V_0\subset V_1,T_1)$ is isomorphic to
$(V\subset V_1,T_1)$ of 3.1.4.  Let $W$ be constructed from
$(V_0\subset V_1,T_1)$.  We have ${f'}^{-1}(y)=V_1$ canonically.
The proof will show that $\Psi$ can be chosen to induce the identity
\[
f^{-1}(y)\xrightarrow{\ \sim\ }{f'}^{-1}(y).
\]

}{Let $d_0<d_1\leq r$, let $t_0\leq\tau(d_0)$, and let
$y\in D_{t_0}^*$.  We have
\[
B_{t_0,d_0}\subset B_{t_0,d_1},
\]
and $B_{t_0,d_1}-B_{t_0,d_0}^{\circ}$ is a bundle over $D_{t_0}$.
Trivialize it and, over the circle $S_{|y|}$, lift
$\partial/\partial\theta$ compatibly with this trivialization.  Let
$T_1$ be the corresponding endomorphism of
\[
V_1=f^{-1}(y)\cap B_{t_0,d_1},
\qquad
V_0=f^{-1}(y)\cap B_{t_0,d_0}.
\]
Then $(V_0\subset V_1,T_1)$ is isomorphic to
$(V\subset V_1,T_1)$ of 3.1.4.  Let $W$ be constructed from
$(V_0\subset V_1,T_1)$.  We have ${f'}^{-1}(y)=V_1$ canonically.
The proof will show that $\Psi$ can be chosen to induce the identity
\[
f^{-1}(y)\xrightarrow{\ \sim\ }{f'}^{-1}(y).
\]

}{E7II-XIV-008}The following drawings of $f$ and of its restriction to
$f^{-1}(\mathbb R)$ summarize the proof.

\begin{center}
\begin{tikzpicture}[x=1cm,y=1cm,line width=.45pt]
  \draw (-4.2,2.6) arc[start angle=180,end angle=360,x radius=1.05,y radius=.28];
  \draw[dashed] (-2.1,2.6) arc[start angle=0,end angle=180,x radius=1.05,y radius=.28];
  \draw (-4.2,2.6)--(-4.2,.15);
  \draw (-2.1,2.6)--(-2.1,.15);
  \draw (-4.2,.15) arc[start angle=180,end angle=360,x radius=1.05,y radius=.28];
  \draw[dashed] (-2.1,.15) arc[start angle=0,end angle=180,x radius=1.05,y radius=.28];
  \draw (-3.15,2.75)--(-3.15,.05);
  \fill (-3.15,1.35) circle (1.8pt);
  \node[left] at (-3.15,1.35) {$0$};
  \draw (-3.65,-1.45) circle (.58);
  \fill (-3.65,-1.45) circle (1.2pt);
  \node[right] at (-3.08,-1.78) {$D_t$};
  \draw (1.2,.1) rectangle (4.0,2.75);
  \foreach \yy in {.28,.48,2.36,2.56}{\draw (1.2,\yy)--(4.0,\yy);}
  \draw (1.2,2.15) .. controls (2.05,2.10) and (2.25,1.75) .. (2.60,1.48)
                    .. controls (2.95,1.75) and (3.18,2.10) .. (4.0,2.15);
  \draw (1.2,.72) .. controls (2.05,.76) and (2.28,1.15) .. (2.60,1.48)
                   .. controls (2.93,1.15) and (3.17,.76) .. (4.0,.72);
  \draw (1.2,1.92)--(2.42,1.54);
  \draw (4.0,1.92)--(2.78,1.54);
  \draw (1.2,1.05)--(2.42,1.42);
  \draw (4.0,1.05)--(2.78,1.42);
  \draw (1.2,1.48)--(2.42,1.48);
  \draw (4.0,1.48)--(2.78,1.48);
  \fill (2.60,1.48) circle (1.8pt);
  \draw (1.15,-1.45)--(4.05,-1.45);
  \fill (2.60,-1.45) circle (1.8pt);
\end{tikzpicture}
\end{center}

\SGApage{II}{\(\approx\)138}
Let $B^+\subset B_{t_0,d_1}$ be the set of points $x$ such that
\[
|f(x)|^2\leq\tau(\delta(x)).
\]
Let $\varphi_n$ be a sequence of $C^\infty$ functions from
$[0,t_0]$ to $(0,d_1]$ such that, for a sequence
$t_0>r_1>r_2>\cdots$ tending to $0$,

\begin{enumerate}[label=(\alph*)]
\item for $|t|^2\geq r_i$, $\varphi_i$ is the inverse function of
  $\tau$;
\item for $|t|^2<r_i$ and $j>i$,
  $\varphi_j(t)<\varphi_i(t)$.
\end{enumerate}

\begin{center}
\begin{tikzpicture}[x=1cm,y=1cm,>=Latex,line width=.5pt]
  \draw[->] (0,0)--(4.5,0) node[right] {$t$};
  \draw[->] (0,0)--(0,3.4) node[above] {$d$};
  \draw (0,3.0)--(3.05,3.0)--(3.05,0);
  \node[left] at (0,3.0) {$d_1$};
  \node[below] at (3.05,0) {$t_0$};
  \draw (0,0) .. controls (.35,1.05) and (1.35,2.15) .. (3.05,2.45);
  \node[right] at (3.05,2.45) {$t=\tau(d)$};
  \draw (0,1.45) .. controls (.45,1.42) and (.80,1.40) .. (1.18,1.68)
        .. controls (1.45,1.88) and (1.68,2.00) .. (1.92,2.10);
  \node at (.55,1.72) {$\varphi_0$};
  \draw (0,.58) .. controls (.23,.56) and (.35,.65) .. (.50,.90)
        .. controls (.67,1.18) and (.84,1.40) .. (1.18,1.68);
\end{tikzpicture}
\end{center}

\SGAvar{The space $B^+-\{0\}$ is the union of the bundles $B_i^+$ over
$D_t$ defined by
\[
\delta(x)\geq\varphi_i(|f(x)|^2).
\]
A standard argument, using the fact that $\delta$ has no critical
points on the fibers of $B^+-\{0\}\to D_t$, shows that there is a
diffeomorphism $\alpha_i:B_i^+\to B_{i+1}^+$, compatible with the
projection to $D_t$ and equal to the identity outside a small
neighborhood of the subspace
\[
\delta(x)=\varphi_i(|f(x)|^2),
\qquad |f(x)|^2\leq r_i,
\]
of $B_i^+$.  Composing the $\alpha_i$ and passing to the limit gives
a diffeomorphism
\[
\alpha:B_0^+-\{x\mid (x)=0,\ \delta(x)=\varphi_0(0)\}
\xrightarrow{\ \sim\ }B^+-\{0\}.
\]
Transport a trivialization of the bundle $B_0^+$ over $D_t$ by
$\alpha$ to obtain an integrable connection $\nabla'$ on
$B^+-\{0\}$.  We may then:

}{The space $B^+-\{0\}$ is the union of the bundles $B_i^+$ over
$D_t$ defined by
\[
\delta(x)\geq\varphi_i(|f(x)|^2).
\]
A standard argument, using the fact that $\delta$ has no critical
points on the fibers of $B^+-\{0\}\to D_t$, shows that there is a
diffeomorphism $\alpha_i:B_i^+\to B_{i+1}^+$, compatible with the
projection to $D_t$ and equal to the identity outside a small
neighborhood of the subspace
\[
\delta(x)=\varphi_i(|f(x)|^2),
\qquad |f(x)|^2\leq r_i,
\]
of $B_i^+$.  Composing the $\alpha_i$ and passing to the limit gives
a diffeomorphism
\[
\alpha:B_0^+-\{x\mid f(x)=0,\ \delta(x)=\varphi_0(0)\}
\xrightarrow{\ \sim\ }B^+-\{0\}.
\]
Transport a trivialization of the bundle $B_0^+$ over $D_t$ by
$\alpha$ to obtain an integrable connection $\nabla'$ on
$B^+-\{0\}$.  We may then:

}{E7II-XIV-009}\SGApage{II}{\(\approx\)139}
\begin{enumerate}[label=(\alph*)]
\item extend the restriction of this connection to
  $f^{-1}(S_{|y|})\cap B^+$ to a connection on
  $f^{-1}(S_{|y|})/S_{|y|}$; this extension gives $T_1$;
\item extend to $B_{t_0,d_1}$ the lift, supplied by $\nabla'$ on
  $B^+-\{0\}$, of the radial vector field on $D_t$.
\end{enumerate}

There is then a unique integrable connection $\nabla$ on
$B_{t_0,d_1}$ extending $\nabla'$, (a), and (b), and a unique
horizontal diffeomorphism from $B_{t_0,d_1}-\{0\}$ to $W-\{0\}$
that is the identity on $f^{-1}(y)$.  This is the required
diffeomorphism.



\subsubsection*{3.1.6. A generalization.}
The preceding discussion generalizes to the case in which
$f:X\to D$ is a morphism such that $0\in X$ is an isolated
non-smooth point \underline{in} $X_0$.  After shrinking $X$ and $D$,
the non-smooth locus of $f$ is then finite over $D$.  The manifold of
vanishing cycles $V$ has a finite set $\Sigma$ of critical points;
$V-\Sigma$ is a differentiable manifold with boundary (its boundary
is compact, although the manifold is not).  The monodromy
transformation $T$ is the identity on the boundary and stabilizes
$\Sigma$.  Finally, Proposition 3.1.5 remains valid: the fibered
homeomorphism $\Psi$ is a diffeomorphism away from the critical locus,
which it carries to the part of $W$ arising from $\Sigma$.

\SGApage{II}{137}
\subsubsection*{3.1.7.}
Proposition 3.1.5 allows us to use the calculations of 1.4.  We obtain:

\SGApage{II}{\(\approx\)139}
\begin{enumerate}[label=(\Alph*)]
\item The vanishing-cycle groups $(R^i\Phi(\Lambda))_0$ are the
  reduced cohomology groups, with coefficients in $\Lambda$, of the
  manifold of vanishing cycles $V$.
\item The groups
  $H^i_{\{0\}}(X_0,R\Psi(\Lambda))$ are the cohomology groups with
  proper support of $V^\circ$, with coefficients in $\Lambda$.
\item The variation
  \[
  \operatorname{Var}:\widetilde H^i(V,\Lambda)
  \longrightarrow H_c^i(V^\circ,\Lambda)
  =H^i(V\bmod\partial V,\Lambda)
  \]
  is the classical variation defined by the automorphism $T$ of $V$,
  which is the identity on $\partial V$.
\end{enumerate}

\SGApage{II}{137}
\subsubsection*{3.1.8.}
Assertions (A), (B), and (C) are in fact more elementary than
Proposition 3.1.5.  Here is a direct proof and construction which
spares us from checking that the isomorphisms in (A) and (B) are
independent of the choices made in the construction of 3.1.5.

\SGApage{II}{\(\approx\)140}
\SGAvar{Let $d\leq r$ and $t\leq\tau(d)$, and let $B$ be the open box
\[
B=\{x\mid\delta(n)<d\ \text{and}\ |f(n)|^2<t\}.
\]
Let $\widetilde D^*$ be a universal covering of $D^*$, and consider
the diagram
\begin{equation}\tag{3.1.8.1}
\begin{tikzcd}[column sep=large,row sep=large]
\{0\} \arrow[r,hook,"\bar\imath"] \arrow[d,equal]
  \& \widetilde D_t \arrow[d,"p"]
  \& \widetilde D_t^* \arrow[l,hook',"\bar\jmath"'] \arrow[d,"p'"] \\
\{0\} \arrow[r,hook,"i"]
  \& D_t
  \& D_t^* \arrow[l,hook',"j"']
\end{tikzcd}
\end{equation}
obtained from (1.2.7.2) by restriction to $D_t$.  Let
\begin{equation}\tag{3.1.8.2}
\begin{tikzcd}[column sep=large,row sep=large]
B_0 \arrow[r,hook,"\bar\imath"] \arrow[d,equal]
  \& \overline B \arrow[d,"p"]
  \& \overline B^* \arrow[l,hook',"\bar\jmath"'] \arrow[d,"p'"] \\
B_0 \arrow[r,hook,"i"]
  \& B
  \& B^* \arrow[l,hook',"j"']
\end{tikzcd}
\end{equation}
be the inverse image of $f:B\to D$ over (3.1.8.1).  We again write
$f$ for the projection from (3.1.8.2) to (3.1.8.1).

\medskip
\noindent\textup{(A)}
By definition, $R^i\Psi_\eta(\Lambda)_0$ is the direct limit, as $U$
runs through the neighborhoods of $0$ in $X$, of the cohomology groups
\[
H^i(U\times_D\widetilde D^*,\Lambda).
\]
Let $U$ run through the fundamental system of neighborhoods of $0$
formed by boxes $B$ as above.  The fiber product
$U\times_D D^*=\overline B^*$ is a bundle with fiber
$V^\circ$ over the contractible space $D_t^*$.  Hence the
direct system of the groups $H^i(\overline B^*,\Lambda)$ is constant,
with value
\[
H^i(V^\circ,\Lambda)=H^i(V,\Lambda).
\]
The resulting isomorphisms
\begin{equation}\tag{3.1.8.3}
R^i\Psi_\eta(\Lambda)_0\simeq H^i(V,\Lambda)
\end{equation}
may be refined to the following composite isomorphism (cf. (B)
below):
\begin{equation}\tag{3.1.8.4}
\begin{tikzcd}[column sep=small,row sep=large,cells={nodes={font=\small}}]
R\Psi_\eta(\Lambda)_0
  \& R\Gamma(B_0,R\Psi_\eta(\Lambda)) \arrow[l,"\sim"']
  \& (Rf_*R\bar\jmath_*\Lambda)_0 \arrow[l,"\sim"']
  \& R\Gamma(\widetilde D_t,Rf_*R\bar\jmath_*\Lambda)
      \arrow[l,"\sim"'] \arrow[d,equal] \\
{}
  \& R\Gamma(V,\Lambda) \arrow[r,"\sim"]
  \& R\Gamma(V^\circ,\Lambda)
  \& R\Gamma(\overline B^*,\Lambda) \arrow[l,"\sim"']
\end{tikzcd}
\end{equation}
The complex $R\Phi(\Lambda)_0$ is defined by the distinguished
triangle
\[
\Lambda\longrightarrow(R\Psi_\eta\Lambda)_0
\longrightarrow R\Phi(\Lambda)_0\longrightarrow{}.
\]
Consequently, the groups $R^i\Phi(\Lambda)_0$ are the reduced
cohomology groups of $V$.

}{Let $d\leq r$ and $t\leq\tau(d)$, and let $B$ be the open box
\[
B=\{x\mid\delta(x)<d\ \text{and}\ |f(x)|^2<t\}.
\]
Let $\widetilde D^*$ be a universal covering of $D^*$, and consider
the diagram
\begin{equation}\tag{3.1.8.1}
\begin{tikzcd}[column sep=large,row sep=large]
\{0\} \arrow[r,hook,"\bar\imath"] \arrow[d,equal]
  \& \widetilde D_t \arrow[d,"p"]
  \& \widetilde D_t^* \arrow[l,hook',"\bar\jmath"'] \arrow[d,"p'"] \\
\{0\} \arrow[r,hook,"i"]
  \& D_t
  \& D_t^* \arrow[l,hook',"j"']
\end{tikzcd}
\end{equation}
obtained from (1.2.7.2) by restriction to $D_t$.  Let
\begin{equation}\tag{3.1.8.2}
\begin{tikzcd}[column sep=large,row sep=large]
B_0 \arrow[r,hook,"\bar\imath"] \arrow[d,equal]
  \& \overline B \arrow[d,"p"]
  \& \overline B^* \arrow[l,hook',"\bar\jmath"'] \arrow[d,"p'"] \\
B_0 \arrow[r,hook,"i"]
  \& B
  \& B^* \arrow[l,hook',"j"']
\end{tikzcd}
\end{equation}
be the inverse image of $f:B\to D$ over (3.1.8.1).  We again write
$f$ for the projection from (3.1.8.2) to (3.1.8.1).

\medskip
\noindent\textup{(A)}
By definition, $R^i\Psi_\eta(\Lambda)_0$ is the direct limit, as $U$
runs through the neighborhoods of $0$ in $X$, of the cohomology groups
\[
H^i(U\times_D\widetilde D^*,\Lambda).
\]
Let $U$ run through the fundamental system of neighborhoods of $0$
formed by boxes $B$ as above.  The fiber product
$U\times_D\widetilde D^*=\overline B^*$ is a bundle with fiber
$V^\circ$ over the contractible space $\widetilde D_t^*$.  Hence the
direct system of the groups $H^i(\overline B^*,\Lambda)$ is constant,
with value
\[
H^i(V^\circ,\Lambda)=H^i(V,\Lambda).
\]
The resulting isomorphisms
\begin{equation}\tag{3.1.8.3}
R^i\Psi_\eta(\Lambda)_0\simeq H^i(V,\Lambda)
\end{equation}
may be refined to the following composite isomorphism (cf. (B)
below):
\begin{equation}\tag{3.1.8.4}
\begin{tikzcd}[column sep=small,row sep=large,cells={nodes={font=\small}}]
R\Psi_\eta(\Lambda)_0
  \& R\Gamma(B_0,R\Psi_\eta(\Lambda)) \arrow[l,"\sim"']
  \& (Rf_*R\bar\jmath_*\Lambda)_0 \arrow[l,"\sim"']
  \& R\Gamma(\widetilde D_t,Rf_*R\bar\jmath_*\Lambda)
      \arrow[l,"\sim"'] \arrow[d,equal] \\
{}
  \& R\Gamma(V,\Lambda) \arrow[r,"\sim"]
  \& R\Gamma(V^\circ,\Lambda)
  \& R\Gamma(\overline B^*,\Lambda) \arrow[l,"\sim"']
\end{tikzcd}
\end{equation}
The complex $R\Phi(\Lambda)_0$ is defined by the distinguished
triangle
\[
\Lambda\longrightarrow(R\Psi_\eta\Lambda)_0
\longrightarrow R\Phi(\Lambda)_0\longrightarrow{}.
\]
Consequently, the groups $R^i\Phi(\Lambda)_0$ are the reduced
cohomology groups of $V$.

}{E7II-XIV-010}\SGApage{II}{\(\approx\)141}
\medskip
\noindent\textup{(B)}
Let $\mathfrak F$ be the family of closed subsets of $\overline B^*$
whose closure in $\overline B$ is proper over $\widetilde D$.  The
arrows in the following diagram are isomorphisms:
\begin{equation}\tag{3.1.8.5}
\begin{tikzcd}[column sep=small,row sep=large,cells={nodes={font=\small}}]
R\Gamma_{\{0\}}(R\Psi_\eta(\Lambda))
  \arrow[r,"\sim","\text{\textcircled{1}}"']
  \& R\Gamma_c(B_0,R\Psi_\eta(\Lambda))
  \& (Rf_!R\bar\jmath_*\Lambda)_0
      \arrow[l,"\sim"',"\text{\textcircled{2}}"] \\
R\Gamma_c(V^\circ,\Lambda)
  \& R\Gamma_{\mathfrak F}(\overline B^*,\Lambda)
      \arrow[l,"\sim"',"\text{\textcircled{5}}"]
      \arrow[r,equal,"\text{\textcircled{4}}"']
  \& R\Gamma(\widetilde D_t,Rf_!R\bar\jmath_*\Lambda)
      \arrow[u,"\text{\textcircled{3}}","\wr"']
\end{tikzcd}
\end{equation}
Indeed, \textcircled{1} is an isomorphism because $B_0$ is a cone
(3.1.2) and $R^i\Psi_\eta(\Lambda)$ is constant away from $0$;
\textcircled{2} by the proper base-change theorem;
\textcircled{3} because $Rf_!R\bar\jmath_*\Lambda$ is constant away
from $0$; \textcircled{4} is the derived-category version of the
spectral sequence of a composite functor; and \textcircled{5} because
$\widetilde D^*$ is contractible.

The isomorphism (3.1.8.5) supplies the desired isomorphisms
\begin{equation}\tag{3.1.8.6}
R^i\Gamma_{\{0\}}(R\Psi_\eta(\Lambda))
=H_c^i(V^\circ,\Lambda).
\end{equation}
We shall verify (C) in 3.1.10.



\SGApage{II}{\(\approx\)142}
\subsubsection*{3.1.9.}
Cf. [3].  Let $X$ be a topological space, let
$i:Y\hookrightarrow X$ be a closed subset, let $T$ be an
automorphism of $X$ inducing the identity on $Y$, and let $\Lambda$ be
an abelian group.  To define the variation $\operatorname{Var}$,
choose a $T$-equivariant resolution $\Lambda_X^*$ of $\Lambda$ on
$X$, a resolution $\Lambda_Y^*$ of the constant sheaf $\Lambda$ on
$Y$ on which, by convention, $T$ acts trivially, and an equivariant
morphism
\[
\varphi:\Lambda_X^*\longrightarrow i_*\Lambda_Y^*
\]
resolving the evident morphism
$\varphi_0:\Lambda_X\to i_*\Lambda_Y$.  Suppose that the sheaves
$\Lambda_X^i$ and $\Lambda_Y^i$ are acyclic for the functors to be
considered.  Passing to global sections gives
\[
\Gamma(\varphi):\Gamma(\Lambda_X^*)
\longrightarrow\Gamma(\Lambda_Y^*).
\]
For simplicity, suppose that every $\Gamma(\varphi)^i$ is surjective;
otherwise, make it surjective by adding the cone of
$\Gamma(\Lambda_Y^*)$ to $\Gamma(\Lambda_X^*)$.  Put
\[
R\Gamma(X\bmod Y,\Lambda)=\operatorname{Ker}(\Gamma(\varphi)).
\]
Under reasonable hypotheses, if $\mathfrak F$ is the family of closed
subsets of $X$ disjoint from $Y$, there is a quasi-isomorphism
\[
\Gamma_{\mathfrak F}(X-Y,\Lambda_X^*)
\xrightarrow{\ \approx\ }R\Gamma(X\bmod Y,\Lambda).
\]
The endomorphism $T-1$ of $\Gamma(\Lambda_X^*)$, where $T$ acts by
transport of structure, is zero on $\Lambda$ (the constants) and
factors through $\operatorname{Ker}(\Gamma(\varphi))$.  Hence, in the
derived category and under the same hypotheses, it gives a variation
morphism
\[
\operatorname{Var}:R\widetilde\Gamma(X,\Lambda)
\longrightarrow R\Gamma_{\mathfrak F}(X-Y,\Lambda).
\]
Without these hypotheses, one obtains instead
\[
\operatorname{Var}:R\widetilde\Gamma(X,\Lambda)
\longrightarrow R\Gamma(X\bmod Y,\Lambda).
\]

\SGApage{II}{\(\approx\)143}
This description simplifies when $T$ is the identity on a
neighborhood of $Y$ and acts as the identity on $\Lambda_X^*$ in that
neighborhood.  The endomorphism $T-1$ of $\Gamma(\Lambda_X^*)$ then
factors through $\Gamma_{\mathfrak F}(\Lambda_X^*)$ and supplies the
variation.

\subsubsection*{Proposition 3.1.10.}
Under the isomorphisms of 3.1.8(A),(B), the variation
\[
(R\Phi_\eta(\Lambda))_0
\longrightarrow R\Gamma_{\{0\}}(R\Psi_\eta(\Lambda))
\]
is equal to the variation of 3.1.9 defined by the automorphism $T$ of
$V$, which is the identity on $\partial V$.

The reader is invited to draw a small picture rather than read the
proof.  Let ${}_fB$ be the closed box of 3.1.2 and let
${}_fC\subset{}_fB$ be a collar along the boundary; thus
${}_fC\simeq\partial_fB\times[0,1]$ as a bundle over $D_t$, through
an isomorphism carrying $\partial_fB\times\{1\}$ identically onto
$\partial_fB$.  Let
\[
{}_f\alpha:{}_fC\xrightarrow{\ \sim\ }
(\partial V\times[0,1])\times D_t
\]
be a trivialization of the bundle ${}_fC$ over $D_t$ whose fiber at
$y\in D_t$ induces the identity
\[
\partial V=\partial(f^{-1}(y))
\xrightarrow{\ \sim\ }\partial V\times\{1\}.
\]
Put $C={}_fC\cap B$ and let
\[
\alpha:C\xrightarrow{\ \sim\ }
(\partial V\times[0,1))\times D_t
\]
be induced by ${}_f\alpha$.  We retain the notation of (3.1.8.2).

\SGAvar{Let $R$ be the equivalence relation on $\overline B$ defined by
$xRy$ if either $x=y$, or $p(x),p(y)\in C$ and
$\alpha(p(x))$, $\alpha(p(y))$ have the same image in
$\partial V\times[0,1)$.  Put
\[
{}_1\overline B=\overline B/R,
\qquad
{}_1\overline B^*=\overline B^*/R,
\qquad
{}_1\overline C^*=p^{-1}(C)/R.
\]
There is a commutative diagram
\[
\begin{tikzcd}[column sep=large,row sep=large]
B_0 \arrow[r,hook,"\bar\imath"] \arrow[d,equal]
  \& \overline B \arrow[d,"q"]
  \& {}_1\overline B^* \arrow[l,hook',"\bar\jmath"'] \arrow[d,"q'"] \\
B_0 \arrow[r,hook,"\bar\imath"]
  \& {}_1\overline B
  \& {}_1\overline B^* \arrow[l,hook',"{}_1\bar\jmath"']
\end{tikzcd}
\]
Let $\Lambda^*$ be a $\pi_1(D^*)$-equivariant resolution of
$\Lambda$ by acyclic sheaves on $\overline B^*$.  By definition, the
variation is computed from the morphism of complexes
\[
\Lambda[0]\longrightarrow
\bar\imath^*\bar\jmath_*(\Lambda^*).
\]
Since $\Lambda\xrightarrow{\approx}Rq_*\Lambda$ and
$\Lambda\xrightarrow{\approx}Rq'_*\Lambda$, it may equally be computed
using an equivariant resolution of $\Lambda$ on
${}_1\overline B^*$ and the morphism
\[
\Lambda[0]\longrightarrow
\bar\imath^*{}_1\bar\jmath_*(\Lambda).
\]
Under (A) and (B), the variation is thereby identified with the
composite morphism
\begin{equation}\tag{3.1.10.1}
\begin{tikzcd}[column sep=small,row sep=large,cells={nodes={font=\small}}]
R\Gamma(V,\Lambda) \arrow[r,"\sim"]
  \& R\Gamma(V^\circ,\Lambda)
  \& R\Gamma({}_1\overline B^*,\Lambda)
      \arrow[l,"\sim"',"\text{\textcircled{1}}"]
      \arrow[d,"\operatorname{Var}"] \\
{}
  \& R\Gamma_c(V^\circ,\Lambda)
  \& R\Gamma({}_1\overline B^*\bmod{}_1\overline C^*,\Lambda)
      \arrow[l,"\sim"',"\text{\textcircled{2}}"]
\end{tikzcd}
\end{equation}
where \textcircled{1} and \textcircled{2} are induced by the inclusion
of $V^\circ$ in ${}_1\overline B^*$, and where
$\operatorname{Var}$ is the variation of 3.1.9 for
${}_1\overline B^*$ and the action of the positive generator of
$\pi_1(D^*)$.  This group acts on ${}_1\overline B^*$ through the
action described in 1.15 and acts trivially on
${}_1\overline C^*$.

}{Let $R$ be the equivalence relation on $\overline B$ defined by
$xRy$ if either $x=y$, or $p(x),p(y)\in C$ and
$\alpha(p(x))$, $\alpha(p(y))$ have the same image in
$\partial V\times[0,1)$.  Put
\[
{}_1\overline B=\overline B/R,
\qquad
{}_1\overline B^*=\overline B^*/R,
\qquad
{}_1\overline C^*=p^{-1}(C)/R.
\]
There is a commutative diagram
\[
\begin{tikzcd}[column sep=large,row sep=large]
B_0 \arrow[r,hook,"\bar\imath"] \arrow[d,equal]
  \& \overline B \arrow[d,"q"]
  \& \overline B^* \arrow[l,hook',"\bar\jmath"'] \arrow[d,"q'"] \\
B_0 \arrow[r,hook,"\bar\imath"]
  \& {}_1\overline B
  \& {}_1\overline B^* \arrow[l,hook',"{}_1\bar\jmath"']
\end{tikzcd}
\]
Let $\Lambda^*$ be a $\pi_1(D^*)$-equivariant resolution of
$\Lambda$ by acyclic sheaves on $\overline B^*$.  By definition, the
variation is computed from the morphism of complexes
\[
\Lambda[0]\longrightarrow
\bar\imath^*\bar\jmath_*(\Lambda^*).
\]
Since $\Lambda\xrightarrow{\approx}Rq_*\Lambda$ and
$\Lambda\xrightarrow{\approx}Rq'_*\Lambda$, it may equally be computed
using an equivariant resolution of $\Lambda$ on
${}_1\overline B^*$ and the morphism
\[
\Lambda[0]\longrightarrow
\bar\imath^*{}_1\bar\jmath_*(\Lambda^*).
\]
Under (A) and (B), the variation is thereby identified with the
composite morphism
\begin{equation}\tag{3.1.10.1}
\begin{tikzcd}[column sep=small,row sep=large,cells={nodes={font=\small}}]
R\Gamma(V,\Lambda) \arrow[r,"\sim"]
  \& R\Gamma(V^\circ,\Lambda)
  \& R\Gamma({}_1\overline B^*,\Lambda)
      \arrow[l,"\sim"',"\text{\textcircled{1}}"]
      \arrow[d,"\operatorname{Var}"] \\
{}
  \& R\Gamma_c(V^\circ,\Lambda)
  \& R\Gamma({}_1\overline B^*\bmod{}_1\overline C^*,\Lambda)
      \arrow[l,"\sim"',"\text{\textcircled{2}}"]
\end{tikzcd}
\end{equation}
where \textcircled{1} and \textcircled{2} are induced by the inclusion
of $V^\circ$ in ${}_1\overline B^*$, and where
$\operatorname{Var}$ is the variation of 3.1.9 for
${}_1\overline B^*$ and the action of the positive generator of
$\pi_1(D^*)$.  This group acts on ${}_1\overline B^*$ through the
action described in 1.1.5 and acts trivially on
${}_1\overline C^*$.

}{E7II-XIV-011}\SGApage{II}{\(\approx\)144}
Choose a connection $\nabla$ on $B^*/D^*$ compatible with the
trivialization of $C$.  This connection trivializes $\overline B^*$
and, on passing to the quotient, defines a map
$u:{}_1\overline B^*\to V$.  For $T$ defined as in 3.1.3, the diagram
\[
\begin{tikzcd}[column sep=large,row sep=large]
{}_1\overline B^* \arrow[r,"u"] \arrow[d,"T"']
  \& V \arrow[d,"T"] \\
{}_1\overline B^* \arrow[r,"u"]
  \& V
\end{tikzcd}
\]
commutes.  The map $u$ induces isomorphisms in cohomology, so
Proposition 3.1.10 follows from (3.1.10.1) and the functoriality of
variation.

\SGApage{II}{143}
Proposition 3.1.10 is meant to justify an explicit calculation to be
made later in which, as here, $T$ is the identity on a neighborhood
of $\partial V$.  We shall take a cohomology class $c$ represented by
a current $C$ and use the fact that $\operatorname{Var}(c)$ is
represented by the current $TC-C$, compactly supported in $V^\circ$.



\SGApage{II}{\(\approx\)145}
\subsubsection*{3.1.11.}
In practice, $T$ is often calculated in two steps.

\begin{enumerate}[label=(\alph*)]
\item Fix $t$ and $d$ with $t\leq\tau(d)$, choose
  $y\in D_t^*$, and choose a connection $\nabla$ on
  $f^{-1}(S_{|y|})\cap B_{t,d}$.  This connection defines a
  diffeomorphism $T_0$ of
  \[
  V_0=f^{-1}(y)\cap B_{t,d}
  \]
  onto itself, which need not be the identity on the boundary.
\item Choose a trivialization of $\partial B_{t,d}$ over $D_t$.
  Integrating $\nabla$ gives a diffeomorphism $\varphi(\theta)$ from
  $\partial V$ to
  \[
  f^{-1}(e^{2\pi i\theta}y)\cap\partial B_{t,d},
  \]
  which is identified with $\partial V$ by the trivialization.  For
  $0\leq\theta\leq1$, $\varphi(\theta)$ is a path in
  $\operatorname{Diff}(\partial V)$ joining the identity to
  $\partial T_0$.
\end{enumerate}

The resulting pair $(T_0,\varphi)$ depends on the arbitrary choices
made.  Since the space of these choices is contractible, it is
independent of them up to homotopy.  For a transformation $T$ of the
kind considered in 3.1.3, $(T,\text{constant path})$ is a pair of the
kind considered here.

\subsubsection*{3.1.12.}
The map
\[
\operatorname{Diff}(V)\longrightarrow\operatorname{Diff}(\partial V)
\]
is a fibration.  Consequently, the assignment
$T\mapsto(T,\text{constant path})$ induces a bijection between:

\SGAvar{\begin{enumerate}[label=\alph*)]
\item diffeomorphisms $T:V\to V$ satisfying
  $\partial T=\operatorname{id}$, up to isotopy; and
\item pairs $(T_0,\varphi)$, where $T_0:V\xrightarrow{\sim}V$ and
  $\varphi$ is a homotopy class of paths from
  $\operatorname{id}$ to $\partial T$, up to isotopy.
\end{enumerate}

}{\begin{enumerate}[label=\alph*)]
\item diffeomorphisms $T:V\to V$ satisfying
  $\partial T=\operatorname{id}$, up to isotopy; and
\item pairs $(T_0,\varphi)$, where $T_0:V\xrightarrow{\sim}V$ and
  $\varphi$ is a homotopy class of paths from
  $\operatorname{id}$ to $\partial T_0$, up to isotopy.
\end{enumerate}

}{E7II-XIV-012}\SGApage{II}{\(\approx\)146}
A pair $(T_0,\varphi)$ calculated by 3.1.11(a),(b) therefore
determines $T$.  Concretely, calculate $T$ from $(T_0,\varphi)$ as
follows.

\medskip
\noindent\textup{($\alpha$)} Choose a collar embedding
  \[
  [0,1]\times\partial V\xhookrightarrow{\ \alpha\ }V,
  \qquad \alpha(1,x)=x,
  \]
  and a representative
  $\varphi:[0,1]\to\operatorname{Diff}(\partial V)$ such that its
  extension by $\operatorname{id}$ for negative parameters is
  $C^\infty$.  On the collar put
  \[
  \beta(x)=\operatorname{pr}_1\alpha^{-1}(x).
  \]
  Thus $\beta$ has no critical point, takes values in $[0,1]$, and
  equals $1$ on $\partial V$.

\medskip
\noindent\textup{($\beta$)} Take $T=T_0$ away from the collar and, on the collar, set
  \[
  T(x)=T_0\bigl(\varphi(\beta(x))^{-1}(x)\bigr).
  \]

\subsubsection*{3.1.13.}
Here is a variant of 3.1.11(b) for calculating the homotopy class of
the path $\varphi$.  On every ray $\arg(z)=\theta$, choose a
trivialization of the bundle $\partial B_{t,d}$ which varies
continuously with $\theta$.  Use the trivialization over the ray
through $y$ to identify
\[
(\partial V)_0=f^{-1}(0)\cap\partial B_{t,d}
\]
with $\partial V$.  Let
\[
\Psi_\theta:V\longrightarrow
f^{-1}(e^{i\theta}y)\cap B_{t,d}
\]
be obtained by integrating $\nabla$.  Transporting
$\partial\Psi_\theta$ along the rays gives a diffeomorphism of
$(\partial V)_0$, hence of $\partial V$, onto itself.  This is
$\varphi(\theta)$.


\subsection*{3.2. The Picard--Lefschetz formula (transcendental case)}

\subsubsection*{3.2.1.}
We shall make the results of 3.1 explicit for the map
\[
f=\sum_{i=1}^{n+1}z_i^2:\mathbb C^{n+1}\longrightarrow\mathbb C.
\]
For $\delta$ take the function $\sum |z_i|^2$, and take the box
$B=B_{a,b}$ defined by
\[
\left|\sum z_i^2\right|\leq a^2,
\qquad
\sum |z_i|^2\leq b^2,
\]
where $0<a<b$.  It is only to simplify the notation that we use a
closed box rather than a box as in 3.1.2.

\SGApage{II}{\(\approx\)147}
\SGAvar{\subsubsection*{3.2.2.}
The manifold of vanishing cycles $V$ has equations
\begin{equation}\tag{3.2.2.1}
\begin{array}{rcl}
\displaystyle\sum z_i^2&=&a^2,\\
\displaystyle\sum |z_i|^2&\leq&b^2.
\end{array}
\end{equation}
Write $z=(z_i)$ in terms of its real and imaginary parts,
$z=x+iy$.  The preceding equations become
\begin{equation}\tag{3.2.2.2}
\begin{array}{rcl}
x^2-y^2&=&a^2,\\
x\mathbin{\cdot}y&=&0
  \qquad\bigl(\text{i.e. }x\perp y\text{ in }\mathbb R^{n}\bigr),\\
x^2+y^2&\leq&b^2.
\end{array}
\end{equation}

}{\subsubsection*{3.2.2.}
The manifold of vanishing cycles $V$ has equations
\begin{equation}\tag{3.2.2.1}
\begin{array}{rcl}
\displaystyle\sum z_i^2&=&a^2,\\
\displaystyle\sum |z_i|^2&\leq&b^2.
\end{array}
\end{equation}
Write $z=(z_i)$ in terms of its real and imaginary parts,
$z=x+iy$.  The preceding equations become
\begin{equation}\tag{3.2.2.2}
\begin{array}{rcl}
x^2-y^2&=&a^2,\\
x\mathbin{\cdot}y&=&0
  \qquad\bigl(\text{i.e. }x\perp y\text{ in }\mathbb R^{n+1}\bigr),\\
x^2+y^2&\leq&b^2.
\end{array}
\end{equation}

}{E7II-XIV-013}Introduce on $V$ the new coordinate system
\[
X=\frac{x}{|x|},
\qquad
Y=y\left(\frac{b^2-a^2}{2}\right)^{-1/2}.
\]
Then $x=X(a^2+y^2)^{1/2}$, and (3.2.2.2) becomes
\begin{equation}\tag{3.2.2.3}
\begin{array}{rcl}
X^2&=&1,\\
X\mathbin{\cdot}Y&=&0,\\
Y^2&\leq&1.
\end{array}
\end{equation}
This identifies $V$ with the tangent ball bundle of the sphere $S^n$.

\subsubsection*{3.2.3.}
The map
\[
\theta\longmapsto \exp(2\pi i\theta)a
\]
identifies $\mathbb R$ with the universal covering of the circle
$S_a$ in $\mathbb C$ having radius $a$ and center $0$.  The pullback
to $\mathbb R$ of the bundle $B$ is isomorphic to $\mathbb R\times V$.
We trivialize it by means of the isomorphism
\[
v\longmapsto \exp(\pi i\theta)v
\]
from $V$ to the fiber of $B$ over $\exp(2\pi i\theta)a$.  Let $T_0$
be the map furnished by this trivialization as in 3.1.11.  In the
coordinates (3.2.2.3),
\[
T_0(X,Y)=(-X,-Y).
\]

\subsubsection*{3.2.4.}
The bundle over the disc $D_a$ of radius $a$ which forms the boundary
of the ``bundle'' $B$ has equations
\[
\sum |z_i|^2=b^2,
\qquad
f(z)\in D_a,
\]
the projection to the disc being $f$.  We shall describe a
trivialization of the inverse image of this bundle over each ray
\[
e^{2\pi i\theta}\mathbb R^+\cap D_a.
\]

\SGApage{II}{\(\approx\)148}
Over the ray $\mathbb R^+\cap D_a$, take the isomorphism with
$(\mathbb R^+\cap D_a)\times\partial V$ which, in the coordinates
(3.2.2.3) on $\partial V$, is given by
\[
z=x+iy\longmapsto(f(z),X,Y)
  =\left(f(z),\frac{x}{|x|},\frac{y}{|y|}\right).
\]
Over $(e^{2\pi i\theta}\mathbb R^+)\cap D_a$, use the image of this
trivialization under multiplication by $e^{\pi i\theta}$.  Apply
3.1.13 and use on $\partial V$ the coordinates $(X,Y)$ satisfying
\begin{equation}\tag{3.2.4.1}
X^2=Y^2=1,
\qquad
X\mathbin{\cdot}Y=0.
\end{equation}
If $\varphi(\theta)$ takes $(X,Y)$ to $(X_\theta,Y_\theta)$, one finds
\begin{equation}\tag{3.2.4.2}
X_\theta+iY_\theta
=e^{\pi i\theta}(X+iY)
\qquad(0\leq\theta\leq1).
\end{equation}

\subsubsection*{3.2.5.}
In summary, the manifold of vanishing cycles $V$ is the tangent ball
bundle of $S^n$.  The monodromy transformation is represented by
$(T_0,\varphi)$, where $T_0$ is the differential of the antipodal map
and $\varphi$ is the following path from $\operatorname{Id}$ to
$\partial T_0$: under $\varphi(\theta)$, each unit tangent vector is
parallel-transported along the geodesic which it determines through a
distance $+\pi\theta$ (the plus sign means ``in the direction in which
the vector points'').  At $\theta=1$ one reaches the antipodal vector.

\SGAvar{By 3.1.13 this gives the following description of $T$.  Let
$\varphi:[0,1]\to[0,1]$ be a $C^\infty$ map which is $0$ near $0$ and
$1$ near $1$.  The map $T$ moves $u$, along the geodesic which $u$
determines, through the distance
\[
-\pi\bigl(1-\varphi(|u|).
\]
Near $|u|=1$ this is the identity; near $|u|=0$ it is the antipodal
map.


\subsubsection*{3.2.7.}
Let $W$ be a differentiable manifold of dimension $x$, and let
\[
p:T(W)\longrightarrow W
\]
be its tangent bundle, regarded as a manifold.  If
$(x_1,\ldots,x_n)$ is a local coordinate system on $W$, then
\[
(x_1,\ldots,x_n,dx_1,\ldots,dx_n)
\]
is a local coordinate system on $T(W)$.  Let $\varepsilon'$ be the
orientation of $T(W)$ defined by this coordinate system.  It is
independent of the choice of the $x_i$, and is therefore globally
defined.

}{By 3.1.12 this gives the following description of $T$.  Let
$\varphi:[0,1]\to[0,1]$ be a $C^\infty$ map which is $0$ near $0$ and
$1$ near $1$.  The map $T$ moves $u$, along the geodesic which $u$
determines, through the distance
\[
-\pi\bigl(1-\varphi(|u|)\bigr).
\]
Near $|u|=1$ this is the identity; near $|u|=0$ it is the antipodal
map.


\subsubsection*{3.2.7.}
Let $W$ be a differentiable manifold of dimension $d$, and let
\[
p:T(W)\longrightarrow W
\]
be its tangent bundle, regarded as a manifold.  If
$(x_1,\ldots,x_d)$ is a local coordinate system on $W$, then
\[
(x_1,\ldots,x_d,dx_1,\ldots,dx_d)
\]
is a local coordinate system on $T(W)$.  Let $\varepsilon'$ be the
orientation of $T(W)$ defined by this coordinate system.  It is
independent of the choice of the $x_i$, and is therefore globally
defined.

}{E7II-XIV-014}\SGApage{II}{\(\approx\)149}
Suppose that $W$ is orientable.  The submanifold $W\subset T(W)$ (the
zero section), endowed with an orientation (respectively, with an
orientation of its normal bundle), defines an element $w$ of
$H_d(T(W))$ (respectively, of $H_c^d(T(W))$).  The self-intersection
of $W\subset T(W)$ is the Euler--Poincare characteristic:
\[
\int w\wedge w=\chi(W)
\]
for the orientation $\varepsilon'$.



For the manifold of vanishing cycles $V$, its complex structure
defines an orientation $\varepsilon$.  It is related to the
orientation $\varepsilon'$ of the tangent bundle of $S^n$ by
\[
\varepsilon=(-1)^{\frac{n(n-1)}2}\varepsilon'.
\]

\subsubsection*{3.2.8.}
The \underline{vanishing cycle} $\delta$ is well defined up to sign.
For $n\ne0$, it is the element of $H_c^n(V^\circ)$ defined by the
submanifold $S^n\subset V^\circ$.  For $n=0$, the space $V$ consists
of two points, and $\delta$ is their difference.  With the complex
orientation of $V$, 3.2.7 gives
\begin{equation}\tag{3.2.8.1}
(\delta,\delta)
=(-1)^{\frac{n(n-1)}2}\bigl(1+(-1)^n\bigr).
\end{equation}

\subsubsection*{3.2.9.}
Put
\[
\widetilde H_c^i(V^\circ)
=\operatorname{Ker}\!\left(
H_c^i(V^\circ)\xrightarrow{\ \operatorname{Tr}\ }\mathbb Z
\right);
\]
this is $H_c^i(V^\circ)$ for $i\ne2n$.  Then
\[
\begin{aligned}
\widetilde H_c^i(V^\circ)&=0 &&\text{for }i\ne n,\\
\widetilde H_c^n(V^\circ)&=\mathbb Z,
  &&\text{generated by }\delta.
\end{aligned}
\]
Dually,
\[
\begin{aligned}
\widetilde H^i(V)&=0 &&\text{for }i\ne n,\\
\widetilde H^n(V)&\xrightarrow[\sim]{\ (\,\cdot\,,\delta)\ }\mathbb Z.
\end{aligned}
\]
For $n\ne0$, the image of $\delta\in H_c^n(V^\circ)$ in $H^n(V)$
is zero when $n$ is odd, and is divisible by two (and no further)
when $n$ is even.


\SGApage{II}{\(\approx\)150}
\SGAvar{\subsubsection*{3.2.10.}
Let us calculate the variation
\[
\operatorname{Var}:\widetilde H^n(V)
\longrightarrow H_c^n(V^\circ).
\]
We first use the orientation $\varepsilon'$ and suppose that $n\ne0$.
Let $\alpha$ be an orientation of the tangent bundle of $S^n$, and
let $\delta\in H_c^n(V^\circ)$ be defined by
$S^n\subset V^\circ$ together with the orientation $\alpha$ of its
normal bundle.  The cocycle $\delta'\in H^n(V)$ such that
$(\delta',\delta)=1$ is represented by a fiber of the fibration of
$V^\circ$ over $S^n$.  The normal bundle of this bundle is the constant
bundle with value the tangent space at one point of $S^n$, and it is
endowed with the orientation $\alpha$.

This assertion is independent of the sign convention chosen to define
the map
\[
(\text{submanifold + orientation of its normal bundle})
\longrightarrow\text{cohomology}.
\]
Indeed, changing the convention changes $\delta$ to $-\delta$ and
$\delta'$ to $-\delta'$.  The assertion is valid when
\[
\operatorname{Tr}:H_c^{2n}(V,\mathbb R)\longrightarrow\mathbb R
\]
is defined as follows:
\begin{description}[leftmargin=2.8em,labelsep=.5em]
\item[\textup{($\alpha$)}] since $\Omega_V^*$ is a soft resolution of $\underline{\mathbb R}$,
  $\Gamma_c(\Omega_V^*)$ represents $R\Gamma(V,\mathbb R)$, whence
  an isomorphism
  \[
  H^i\bigl(\Gamma_c(V,\Omega_V^*)\bigr)
  \simeq H_c^i(V,\mathbb R)
  \]
  (the iterated coboundary is not used);
\item[\textup{($\beta$)}] $\operatorname{Tr}$ is integration.
\end{description}

}{\subsubsection*{3.2.10.}
Let us calculate the variation
\[
\operatorname{Var}:\widetilde H^n(V)
\longrightarrow H_c^n(V^\circ).
\]
We first use the orientation $\varepsilon'$ and suppose that $n\ne0$.
Let $\alpha$ be an orientation of the tangent bundle of $S^n$, and
let $\delta\in H_c^n(V^\circ)$ be defined by
$S^n\subset V^\circ$ together with the orientation $\alpha$ of its
normal bundle.  The cocycle $\delta'\in H^n(V)$ such that
$(\delta',\delta)=1$ is represented by a fiber of the fibration of
$V^\circ$ over $S^n$.  The normal bundle of this fiber is the constant
bundle with value the tangent space at one point of $S^n$, and it is
endowed with the orientation $\alpha$.

This assertion is independent of the sign convention chosen to define
the map
\[
(\text{submanifold + orientation of its normal bundle})
\longrightarrow\text{cohomology}.
\]
Indeed, changing the convention changes $\delta$ to $-\delta$ and
$\delta'$ to $-\delta'$.  The assertion is valid when
\[
\operatorname{Tr}:H_c^{2n}(V,\mathbb R)\longrightarrow\mathbb R
\]
is defined as follows:
\begin{description}[leftmargin=2.8em,labelsep=.5em]
\item[\textup{($\alpha$)}] since $\Omega_V^*$ is a soft resolution of $\underline{\mathbb R}$,
  $\Gamma_c(\Omega_V^*)$ represents $R\Gamma_c(V,\mathbb R)$, whence
  an isomorphism
  \[
  H^i\bigl(\Gamma_c(V,\Omega_V^*)\bigr)
  \simeq H_c^i(V,\mathbb R)
  \]
  (the iterated coboundary is not used);
\item[\textup{($\beta$)}] $\operatorname{Tr}$ is integration.
\end{description}

}{E7II-XIV-015}The class $T(\delta')$ coincides with $\delta'$ in a neighborhood of
$\partial V$, so that
\(
\operatorname{Var}(\delta')=T(\delta')-\delta'
\)
has compact support in $V^\circ$.  We have
\[
T(\delta')-\delta'\sim-\delta,
\qquad\text{and hence}
\]
\[
\operatorname{Var}(\delta')
=-(\delta',\delta)\,\delta.
\]

Here is the picture for $n=1$.  We draw the tangent bundle of $S^1$
so that vectors pointing in the chosen counterclockwise orientation of
$S^1$ point outward.

\SGApage{II}{\(\approx\)151}
\begin{center}
\begin{tikzpicture}[>=Latex,line width=.45pt,scale=.92]
  \begin{scope}[shift={(-2.8,0)}]
    \draw (0,0) circle (2.15);
    \draw (0,0) circle (1.20);
    \draw (0,0) circle (.43);
    \draw (0,.43)--(0,2.15);
    \draw[->] (-.02,1.88)--(-.48,1.88);
    \draw[->] (-.02,1.13)--(-.48,1.13);
    \node[right] at (.10,1.10) {$\delta'$};
    \draw[->] (-.85,.85)--(-1.16,1.16);
    \draw[->] (-.85,-.85)--(-1.17,-1.17);
    \node[above left] at (-1.02,-.98) {$\delta$};
    \draw[->] (.85,-.85)--(1.17,-1.17);
    \draw[->] (1.43,1.12)--(1.43,1.72)
      node[right] {$2$};
    \draw[->] (1.43,1.12)--(.83,1.12)
      node[below] {$1$};
    \node at (1.03,1.63) {$\varepsilon$};
    \node[below=5pt] at (0,-2.15)
      {$\delta,\ \delta'\ \text{and}\ \varepsilon$};
  \end{scope}
  \begin{scope}[shift={(2.8,0)}]
    \draw (0,0) circle (2.15);
    \draw (0,0) circle (.43);
    \draw (0,.43)--(0,2.15);
    \draw
      (0,2.15)
      .. controls (.02,1.68) and (1.18,1.42) .. (1.28,.30)
      .. controls (1.40,-.72) and (.62,-1.35) .. (-.12,-1.31)
      .. controls (-1.05,-1.27) and (-1.42,-.42) .. (-1.10,.48)
      .. controls (-.83,1.10) and (-.34,1.25) .. (0,1.02);
    \draw[->] (1.08,.78)--(.74,.48);
    \node[below=5pt] at (0,-2.15)
      {$\operatorname{Var}(\delta')$};
  \end{scope}
\end{tikzpicture}
\end{center}

For the orientation $\varepsilon$, we therefore have
\[
\operatorname{Var}(x)
=-\frac{\varepsilon}{\varepsilon'}(x,\delta)\,\delta.
\]

\subsubsection*{Theorem 3.2.11.}
\underline{The group $\widetilde H_c^n(V^\circ)$ is isomorphic to
$\mathbb Z$, generated by the vanishing cycle $\delta$.  The group
$\widetilde H^n(V)$ is also isomorphic to $\mathbb Z$ and is dual to
the preceding group.  All the other groups
$\widetilde H_c^i(V^\circ)$ and $\widetilde H^i(V)$ vanish.}

\medskip
\noindent\underline{The variation is given by}
\[
\operatorname{Var}(x)
=-(-1)^{\frac{n(n-1)}2}(x,\delta)\,\delta,
\qquad\underline{\text{and}}
\]
\[
(\delta,\delta)
=(-1)^{\frac{n(n-1)}2}\bigl(1+(-1)^n\bigr).
\]
According to the congruence class of $n$ modulo $4$, we have the
following table:
\[
\begin{array}{lcccc}
n\bmod4&0&1&2&3\\[2pt]
(\delta,\delta)&2&0&-2&0\\[2pt]
\operatorname{Var}(x)=\pm(x,\delta)\delta&-&-&+&+\\[2pt]
T(\delta)=(-1)^{n+1}\delta&-\delta&\delta&-\delta&\delta
\end{array}
\]
As a check, we indeed have
\[
T(\delta)=\delta+\operatorname{Var}(\delta).
\]


\SGApage{II}{152}
\section*{4. De Rham cohomology}

\SGAvar{\subsubsection*{4.1.}
Let $D$ be a complex analytic disc, let $0\in D$, put
$D^*=D-\{0\}$, and let $\widetilde D^*$ be a universal covering of
$D^*$.  Let $X$ be an analytic space, let
$f:X\longrightarrow D$ be a morphism, and consider the diagram of
(1.2.7.1):
\begin{equation}\tag{4.1.1}
\begin{tikzcd}[column sep=large,row sep=large]
X_0 \arrow[r,hook,"\bar i"] \arrow[d,equal]
  \& \bar X \arrow[d,"p"']
  \& \bar X^* \arrow[l,hook',"\bar j"'] \arrow[d,"p'"'] \\
X_0 \arrow[r,hook,"i"]
  \& X
  \& X^* \arrow[l,hook',"j"']
\end{tikzcd}
\end{equation}
Suppose that $X^*=X-X_0$ is nonsingular (that is, smooth).  Its
covering $\bar X^*$ is therefore smooth.  Denote by
$\Omega_{\bar X^*}^*$ the De Rham complex of sheaves of holomorphic
differential forms on $\Omega^*_{\bar X^*}$.  It is a resolution of the constant
sheaf $\mathbb C$ (the holomorphic Poincare lemma), and hence gives a
morphism
\begin{equation}\tag{4.1.2}
\bar j_*\Omega_{\bar X^*}^*
\longrightarrow R\bar j_*\mathbb C
\end{equation}
in $D^+(\bar X,\mathbb C)$.

}{\subsubsection*{4.1.}
Let $D$ be a complex analytic disc, let $0\in D$, put
$D^*=D-\{0\}$, and let $\widetilde D^*$ be a universal covering of
$D^*$.  Let $X$ be an analytic space, let
$f:X\longrightarrow D$ be a morphism, and consider the diagram of
(1.2.7.1):
\begin{equation}\tag{4.1.1}
\begin{tikzcd}[column sep=large,row sep=large]
X_0 \arrow[r,hook,"\bar i"] \arrow[d,equal]
  \& \bar X \arrow[d,"p"']
  \& \bar X^* \arrow[l,hook',"\bar j"'] \arrow[d,"p'"'] \\
X_0 \arrow[r,hook,"i"]
  \& X
  \& X^* \arrow[l,hook',"j"']
\end{tikzcd}
\end{equation}
Suppose that $X^*=X-X_0$ is nonsingular (that is, smooth).  Its
covering $\bar X^*$ is therefore smooth.  Denote by
$\Omega_{\bar X^*}^*$ the De Rham complex of sheaves of holomorphic
differential forms on $\bar X^*$.  It is a resolution of the constant
sheaf $\mathbb C$ (the holomorphic Poincare lemma), and hence gives a
morphism
\begin{equation}\tag{4.1.2}
\bar j_*\Omega_{\bar X^*}^*
\longrightarrow R\bar j_*\mathbb C
\end{equation}
in $D^+(\bar X,\mathbb C)$.

}{E7II-XIV-016}\subsubsection*{Proposition 4.2.}
\underline{The morphism (4.1.2) is an isomorphism.}

\SGAvar{It is clear that (4.1.1) is a quasi-isomorphism at every point of
$\bar X^*\subset\bar X$.
It is therefore enough to show that every point $x\in X_0$ has a
fundamental system of neighborhoods $U$ in $\bar X$ such that
\[
H^q\bigl(\bar j^{-1}(U),\Omega^p\bigr)=0
\qquad\text{for }q>0.
\]
As $V$ runs through a fundamental system of neighborhoods of $x$ in
$X$, the sets $p^{-1}(V)$ form a fundamental system of neighborhoods
of $X$ in $\bar X$.  In particular, the $p^{-1}(V)$ for
$V$ a Stein neighborhood of $x$ in $X$ form a fundamental system of neighborhoods of $X$.  If $V$ is Stein, then
$V-X_0$ is still Stein (since $X_0\cap V$ is defined by one equation
in $V$), as is its covering
\[
p'^{-1}(V-X_0)=\bar j^{-1}p^{-1}(V).
\]
Consequently, for every coherent sheaf $\mathfrak F$ on $\bar X^*$,
\[
H^q\bigl(\bar j^{-1}p^{-1}(V),\mathfrak F\bigr)=0
\qquad\text{for }q>0,
\]
which proves the proposition.

}{It is clear that (4.1.2) is a quasi-isomorphism at every point of
$\bar X^*\subset\bar X$.
It is therefore enough to show that every point $x\in X_0$ has a
fundamental system of neighborhoods $U$ in $\bar X$ such that
\[
H^q\bigl(\bar j^{-1}(U),\Omega^p\bigr)=0
\qquad\text{for }q>0.
\]
As $V$ runs through a fundamental system of neighborhoods of $x$ in
$X$, the sets $p^{-1}(V)$ form a fundamental system of neighborhoods
of $x$ in $\bar X$.  In particular, one may take the $p^{-1}(V)$ for
$V$ a Stein neighborhood of $x$ in $X$.  If $V$ is Stein, then
$V-X_0$ is still Stein (since $X_0\cap V$ is defined by one equation
in $V$), as is its covering
\[
p'^{-1}(V-X_0)=\bar j^{-1}p^{-1}(V).
\]
Consequently, for every coherent sheaf $\mathfrak F$ on $\bar X^*$,
\[
H^q\bigl(\bar j^{-1}p^{-1}(V),\mathfrak F\bigr)=0
\qquad\text{for }q>0,
\]
which proves the proposition.

}{E7II-XIV-017}\SGApage{II}{153}
\subsubsection*{Corollary 4.3.}
\underline{The isomorphism (4.1.2) induces an isomorphism}
\begin{equation}\tag{4.3.1}
\bar i^*\bar j_*\Omega_{\bar X^*}^*
\xrightarrow{\ \sim\ }R\Psi_\eta(\mathbb C).
\end{equation}

\subsubsection*{Variant 4.4.}
\underline{Let $V$ be a locally constant sheaf of finite-dimensional
complex vector spaces on $X^*$} (\underline{a complex local system}).  \underline{The
quasi-isomorphism}
\[
V\longrightarrow\Omega_{X^*}^*(V)
\]
\underline{of the Poincare lemma induces a quasi-isomorphism}
\begin{equation}\tag{4.4.1}
\bar i^*\bar j_*\Omega_{\bar X^*}^*(p'^*V)
\xrightarrow{\ \sim\ }R\Psi_\eta(V).
\end{equation}

We shall show that (4.3.1) remains true when
$\bar i^*\bar j_*\Omega_{\bar X^*}^*$ is replaced by a subcomplex
defined by growth and monodromy conditions.


\SGAvar{\subsubsection*{4.5.}
Let $f:X^*\longrightarrow D^*$ be a topological space over $D^*$,
let $F$ be a sheaf of $f^\circ\mathcal O_{D^*}$-modules on $X^*$,
and let $\bar F$ be its inverse image on
\[
\bar X^*=X^*\times_{D^*} D^*.
\]

}{\subsubsection*{4.5.}
Let $f:X^*\longrightarrow D^*$ be a topological space over $D^*$,
let $F$ be a sheaf of $f^\circ\mathcal O_{D^*}$-modules on $X^*$,
and let $\bar F$ be its inverse image on
\[
\bar X^*=X^*\times_{D^*}\widetilde D^*.
\]

}{E7II-XIV-018}Let $z:D\hookrightarrow\mathbb C$ be a holomorphic embedding with
$z(0)=0$.  We shall always be free to shrink $D$ below; for a
sufficiently small disc the existence of $z$ is trivial.  Choose on
$\widetilde D^*$ a function $\log z$ such that
$\exp(\log z)=z$, and put
\[
z^\alpha=\exp(\alpha\log z)
\qquad(\alpha\in\mathbb C).
\]

\SGAvar{The group $\pi_1(D^*)$ acts, by transport of structure, on
$H^0(\bar X^*,\bar F)$.  An element
$f\in H^0(\bar X^*,\bar F)$ is of \underline{finite determination}
if its transforms under $\pi_1(D^*)$ span a finite-dimensional
complex vector space $V(F)$.  Let $T$ be the positive generator of
$\pi_1(D^*)$.  The section $f$ is of finite determination if and only
if there is a nonzero polynomial $P\in\mathbb C[X]$ such that
$P(T)(f)=0$.  The polynomial of least degree with this property is
then both the minimal polynomial and the characteristic polynomial of
$T$ acting on $V(f)$.

}{The group $\pi_1(D^*)$ acts, by transport of structure, on
$H^0(\bar X^*,\bar F)$.  An element
$f\in H^0(\bar X^*,\bar F)$ is of \underline{finite determination}
if its transforms under $\pi_1(D^*)$ span a finite-dimensional
complex vector space $V(f)$.  Let $T$ be the positive generator of
$\pi_1(D^*)$.  The section $f$ is of finite determination if and only
if there is a nonzero polynomial $P\in\mathbb C[X]$ such that
$P(T)(f)=0$.  The polynomial of least degree with this property is
then both the minimal polynomial and the characteristic polynomial of
$T$ acting on $V(f)$.

}{E7II-XIV-019}\SGApage{II}{\(\approx\)154}
We say that $f$ is of \underline{unipotent finite determination} if
there is an $A$ such that
\[
(T-1)^A f=0.
\]
On the space of these sections define
\[
N=\frac{\log T}{2\pi i}
\]
by the finite sum
\[
N(f)=-\frac1{2\pi i}
\sum_{n>0}\frac{(1-T)^n}{n}\,f.
\]
If $f$ (respectively, $g$) is a section of unipotent finite
determination in $H^0(D^*,\mathcal O)$ (respectively, in
$H^0(\bar X^*,\bar F)$), then
\[
N(fg)=N(f)g+fN(g),
\]
since the logarithm of an automorphism is a derivation.

We say that $f$ is of
\underline{quasi-unipotent finite determination} if there are
$A$ and $B>0$ such that
\[
(T^B-1)^A f=0.
\]

\SGAvar{\subsubsection*{Examples 4.6.}
For $X=D^*$ and $F=\mathcal O$, one has
\[
\bigl(T-\exp(2\pi i\alpha)\bigr)^{k+1}
  \bigl(z^\alpha(\log z)^k\bigr)=0,
\]
\[
N(\log z)=1,
\]
and
\[
N\left.\frac{(\log z)^k}{k!}\right)
=\frac{(\log z)^{k-1}}{(k-1)!},
\qquad
(0\text{ for }k=0).
\]

}{\subsubsection*{Examples 4.6.}
For $X=D^*$ and $F=\mathcal O$, one has
\[
\bigl(T-\exp(2\pi i\alpha)\bigr)^{k+1}
  \bigl(z^\alpha(\log z)^k\bigr)=0,
\]
\[
N(\log z)=1,
\]
and
\[
N\left(\frac{(\log z)^k}{k!}\right)
=\frac{(\log z)^{k-1}}{(k-1)!},
\qquad
(0\text{ for }k=0).
\]

}{E7II-XIV-020}\subsubsection*{Proposition 4.7.}
\underline{A section $f\in H^0(\bar X^*,\bar F)$ is of finite
determination if and only if it can be written}
\begin{equation}\tag{4.7.1}
f=\sum_{\alpha,k}f_{\alpha,k}\,
z^\alpha\frac{(\log z)^k}{k!},
\end{equation}
\underline{where $0\leq\operatorname{Re}(\alpha)<1$,
$k\in\mathbb N$, and the
$f_{\alpha,k}\in H^0(X^*,F)$ are almost all zero.  For $f$ of finite
determination, the decomposition (4.7.1) is unique.}

\SGApage{II}{\(\approx\)155}
By 4.6, every section of the form (4.7.1) is of finite determination.
Conversely, let $f$ be of finite determination, with $P(T)(f)=0$.
Write
\[
P(X)=\prod_{\beta\in\mathbb C}(X-\beta)^{n'(\beta)}
=\prod_{0\leq\operatorname{Re}(\alpha)<1}
 \bigl(X-\exp(2\pi i\alpha)\bigr)^{n(\alpha)},
\]
where
$n(\alpha)=n'(\exp(2\pi i\alpha))$ is almost always zero.
Decomposing the $\mathbb C[X]$-module $V(f)$ into its isotypic
components, one finds a unique decomposition
\[
f=\sum_\alpha f'_\alpha
\]
such that
\[
\bigl(T-\exp(2\pi i\alpha)\bigr)^{n(\alpha)}f'_\alpha=0.
\]

Choose a polynomial $Q_\alpha\in\mathbb C[X]$ congruent to $1$
modulo
$\bigl(X-\exp(2\pi i\alpha)\bigr)^{n(\alpha)}$ and to $0$ modulo
$\bigl(X-\exp(2\pi i\alpha')\bigr)^{n(\alpha')}$ for every
$\alpha'\ne\alpha$.  Then
$f'_\alpha=Q_\alpha(T)(f)$.  Put
\begin{equation}\tag{4.7.2}
f_\alpha=z^{-\alpha}f'_\alpha
=z^{-\alpha}Q_\alpha(T)(f).
\end{equation}
Since
\[
T(z^{-\alpha}f'_\alpha)
=\exp(-2\pi i\alpha)z^{-\alpha}T(f'_\alpha),
\]
we have
\[
(T-1)^{n(\alpha)}f_\alpha=0.
\]

If
\[
f_\alpha=\sum_k f_{\alpha,k}\frac{(\log z)^k}{k!},
\qquad f_{\alpha,k}\in H^0(X^*,F),
\]
then 4.6 and the identity
$\sum_j(-1)^j\binom{k}{j}=\delta_{0,k}$ give
\begin{equation}\tag{4.7.3}
f_{\alpha,k}
=\sum_{j\geq0}(-1)^j\frac{(\log z)^j}{j!}
  N^{j+k}f_\alpha,
\end{equation}
which proves uniqueness.

Conversely, if $f_\alpha$ is a section of unipotent finite
determination, the $f_{\alpha,k}$ defined by (4.7.3) vanish for large
$k$ (because $N^k f_\alpha=0$ for large $k$) and are invariant under
$T$; hence they identify with sections of $F$ on $X^*$.  Equivalently,
we must show that $Nf_{\alpha,k}=0$.  Apply to $N^k f_\alpha$ the
identity
\[
\begin{aligned}
N\sum_{j\geq0}(-1)^j\frac{(\log z)^j}{j!}N^j
&=\sum_{j>0}(-1)^j\frac{(\log z)^{j-1}}{(j-1)!}N^j\\
&\quad+\sum_{j\geq0}(-1)^j\frac{(\log z)^j}{j!}N^{j+1}=0.
\end{aligned}
\]
Finally,
\[
\sum_{k\geq0}f_{\alpha,k}\frac{(\log z)^k}{k!}
=\sum_{j,k\geq0}(-1)^j
 \frac{(\log z)^{j+k}}{j!k!}N^{j+k}f_\alpha
\]
\[
=\sum_{n\geq0}
 \left(\sum_j(-1)^j\binom nj\right)
 \frac{(\log z)^n}{n!}N^n f_\alpha
=f_\alpha,
\]
which proves 4.7.

\SGApage{II}{156}
\subsubsection*{Remark 4.8.1.}
The condition $\operatorname{Re}(\alpha)\in[0,1)$ in 4.7 may be
replaced by $\alpha\in C$, where $C$ is any fundamental domain for
$\mathbb C\longrightarrow\mathbb C/\mathbb Z$.

\subsubsection*{Remark 4.8.2.}
Let $f$ be of finite determination, with decomposition (4.7.1).  For
$0\leq\operatorname{Re}(\alpha)<1$ and for $\alpha$ such that some
$f_{\alpha,k}\ne0$, put
\[
m(\alpha)=1+\sup\{k\mid f_{\alpha,k}\ne0\}.
\]
If all the $f_{\alpha,k}$ vanish, put $m(\alpha)=0$.  The
characteristic polynomial of $T$ acting on $V(f)$ is
\[
\operatorname{Pc}(T,V(f))
=\prod_\alpha\bigl(X-\exp(2\pi i\alpha)\bigr)^{m(\alpha)}.
\]
In particular, $f$ is of quasi-unipotent finite determination if and
only if the $\alpha$ in (4.7.1) may all be taken in $\mathbb Q$.


\subsubsection*{4.9.}
Return to the more restrictive hypotheses of 4.1, and let $F$ be a
coherent sheaf on $X$ whose restriction to $X^*$ is locally free.

\subsubsection*{Proposition 4.10.}
\underline{Let $f$ be a section of finite determination of the inverse
image $\bar F$ of $F$ on $\bar X^*$.  Then $f$ has moderate growth
along $X_0$ if and only if the $f_{\alpha,k}$ in the decomposition
(4.7.1) are meromorphic along $X_0$.}

For the definition of ``moderate growth,'' we refer to [1], 2.10.  It
is clear from (4.7.1), (4.7.2), and (4.7.3) that $f$ has moderate
growth if and only if the $f_{\alpha,k}$ do.  For sections of $F$ on
$X^*$, moderate growth along $X_0$ is known to be equivalent to being
meromorphic along $X_0$.

\SGApage{II}{\(\approx\)157}
For $X=D$ and $F=\mathcal O$, 4.10 has the following elementary
corollary.

\SGAvar{\subsubsection*{Corollary 4.11.}
\underline{Let $K$ be the field of germs at $0$ of meromorphic
functions on $D$, and let $K^m$} (\underline{respectively, $K^{mu}$ and $K^{mqu}$})
\underline{be the} $K$-\underline{algebra of germs at $0$ of holomorphic functions on
$\widetilde D^*$ which have moderate growth at $0$ and finite
determination (respectively, unipotent finite determination and
quasi-unipotent finite determination}).  \underline{Then the following are bases
over $K$:}
\[
\begin{array}{lll}
K^m:&z^\alpha(\log z)^k,
  &0\leq\operatorname{Re}(\alpha)<1,\ k\in\mathbb N,\\[2pt]
K^{mu}:&(\log z)^k,&k\in\mathbb N,\\[2pt]
K^{muq}:&z^\alpha(\log z)^k,
  &\alpha\in\mathbb Q,\ 0\leq\alpha<1,\ k\in\mathbb N.
\end{array}
\]

Under the hypotheses of 4.9, denote by $\Psi_\eta^m(F)$
(respectively, $\Psi_\eta^{mu}(F)$ and $\Psi_\eta^{muq}(F)$) the
subsheaf of $\bar i^*\bar j_*\bar F$ consisting of sections of
$\bar F$, near points of $X_0$, which have moderate growth along
$X_0$ and finite determination (respectively, unipotent finite
determination and quasi-unipotent finite determination).  Let
\[
F^m\subset j_*j^*F
\]
be the sheaf on $X$ of sections of $j_*j^*F$ which are meromorphic
along $X_0$.  Proposition 4.10 has the following corollary, independent
of the arbitrary choices of $z$ and $\log z$.

\subsubsection*{Corollary 4.12.}
\underline{The natural morphisms}
\[
\textup{(i)}\qquad
K^m\otimes_K i^*F^m\longrightarrow\Psi_\eta^m(F),
\]
\[
\textup{(ii)}\qquad
K^{mu}\otimes_K i^*F^m\longrightarrow\Psi_\eta^{mu}(F),
\]
\[
\textup{(iii)}\qquad
K^{muq}\otimes_K i^*F^m\longrightarrow\Psi_\eta^{muq}(F)
\]
\underline{are isomorphisms.}

}{\subsubsection*{Corollary 4.11.}
\underline{Let $K$ be the field of germs at $0$ of meromorphic
functions on $D$, and let $K^m$} (\underline{respectively, $K^{mu}$ and $K^{mqu}$})
\underline{be the} $K$-\underline{algebra of germs at $0$ of holomorphic functions on
$\widetilde D^*$ which have moderate growth at $0$ and finite
determination (respectively, unipotent finite determination and
quasi-unipotent finite determination}).  \underline{Then the following are bases
over $K$:}
\[
\begin{array}{lll}
K^m:&z^\alpha(\log z)^k,
  &0\leq\operatorname{Re}(\alpha)<1,\ k\in\mathbb N,\\[2pt]
K^{mu}:&(\log z)^k,&k\in\mathbb N,\\[2pt]
K^{mqu}:&z^\alpha(\log z)^k,
  &\alpha\in\mathbb Q,\ 0\leq\alpha<1,\ k\in\mathbb N.
\end{array}
\]

Under the hypotheses of 4.9, denote by $\Psi_\eta^m(F)$
(respectively, $\Psi_\eta^{mu}(F)$ and $\Psi_\eta^{mqu}(F)$) the
subsheaf of $\bar i^*\bar j_*\bar F$ consisting of sections of
$\bar F$, near points of $X_0$, which have moderate growth along
$X_0$ and finite determination (respectively, unipotent finite
determination and quasi-unipotent finite determination).  Let
\[
F^m\subset j_*j^*F
\]
be the sheaf on $X$ of sections of $j_*j^*F$ which are meromorphic
along $X_0$.  Proposition 4.10 has the following corollary, independent
of the arbitrary choices of $z$ and $\log z$.

\subsubsection*{Corollary 4.12.}
\underline{The natural morphisms}
\[
\textup{(i)}\qquad
K^m\otimes_K i^*F^m\longrightarrow\Psi_\eta^m(F),
\]
\[
\textup{(ii)}\qquad
K^{mu}\otimes_K i^*F^m\longrightarrow\Psi_\eta^{mu}(F),
\]
\[
\textup{(iii)}\qquad
K^{mqu}\otimes_K i^*F^m\longrightarrow\Psi_\eta^{mqu}(F)
\]
\underline{are isomorphisms.}

}{E7II-XIV-021}The main result of this section is the following.

\subsubsection*{Theorem 4.13.}
\underline{The natural morphism induced by (4.3.1),}
\begin{equation}\tag{4.13.1}
\Psi_\eta^{mqu}(\Omega_X^*)
\longrightarrow R\Psi_\eta(\mathbb C),
\end{equation}
\underline{is an isomorphism in $D^+(X_0,\mathbb C)$.}

\SGApage{II}{\(\approx\)158}
We shall in fact prove the more general results 4.15 and 4.17.  Their
statements require a few preliminaries which a reader interested only
in 4.13 may omit.

\subsubsection*{4.14.}
Let $V$ be a locally constant sheaf of finite-dimensional complex
vector spaces on $X^*$ (a local system).  Let
\[
\mathcal V=V\otimes_{\mathbb C}\mathcal O
\]
be the corresponding locally free $\mathcal O$-module, and let
$\nabla$ be the integrable connection on $\mathcal V$ for which the
local sections of $V$ are horizontal.  It is known that $\mathcal V$
admits an extension $\widetilde{\mathcal V}$ to a coherent sheaf on
$X$ such that $\nabla$ is \underline{regular} along $X_0$ relative to
this extension.  It is also known that any two such extensions are
meromorphically equivalent ([1], II, 5.7).  Thus the sheaf
$\Psi_\eta^m(\widetilde{\mathcal V})$, and its variants
$\Psi^{mu}$ and $\Psi^{mqu}$, does not depend on the chosen extension;
we denote it by $\Psi_\eta^m(\mathcal V)$.  Likewise define
\[
\Psi_\eta^m\bigl(\Omega_{X^*}^p(V)\bigr)
=\Psi_\eta^m(\Omega_X^p\otimes\mathcal V),
\]
and its variants $\Psi^{mu}$ and $\Psi^{mqu}$.

\subsubsection*{Proposition 4.15.}
\underline{The natural morphism induced by (4.4.1),}
\begin{equation}\tag{4.15.1}
\Psi_\eta^m\bigl(\Omega_{X^*}^*(V)\bigr)
\longrightarrow R\Psi_\eta(V),
\end{equation}
\underline{is an isomorphism in $D^+(X_0,\mathbb C)$.}

\subsubsection*{4.16.}
A local system $V$ as in 4.14 is said to be \underline{unipotent}
(respectively, \underline{quasi-unipotent}) along $X_0$ if, for every
morphism $u$ from the standard disc $D_0$ to $X$ such that
$u^{-1}(X_0)=\{0\}$, the monodromy transformation of the local system
$u^*V$ on $D_0^*$ is unipotent (respectively, quasi-unipotent).

\subsubsection*{Proposition 4.17.}
\underline{If $V$ is quasi-unipotent, the natural morphism induced by
(4.4.1),}
\begin{equation}\tag{4.17.1}
\Psi_\eta^{mqu}\bigl(\Omega_{X^*}^*(V)\bigr)
\longrightarrow R\Psi_\eta(V),
\end{equation}
\underline{is an isomorphism in $D^+(X_0,\mathbb C)$.}

\SGApage{II}{157}
Theorem 4.13 is the special case of 4.17 in which
$V=\underline{\mathbb C}$.


\SGApage{II}{\(\approx\)159}
\subsubsection*{4.18.}
We shall first treat the case in which $X$ is smooth and
$(X_0)_{\mathrm{red}}$ is a normal-crossings divisor in $X$.  Since
the theorems are local, we may suppose that $f$ is the map
\[
f(\underline z)=\prod_{i=1}^d z_i^{n_i},
\qquad 0<d\leq e,\quad n_i>0,
\]
from $X=D^e$ to $D$, where this time $D$ denotes the standard disc
$D=\{z\in\mathbb C\mid |z|<1\}$.  It is enough to prove that the
fiber at $0\in D^e$ of (4.15.1) (respectively, (4.17.1)) is a
quasi-isomorphism.

Let $U$ be a neighborhood of $0$ in $D^e$ of the form
$|z_i|<a$, where $0<a\leq1$.  Define $U_0$, $\bar U$, $U^*$, and
$\bar U^*$ by taking the inverse image of $U$ in (4.1.1).  The
cohomology of $U^*$ or $\bar U^*$ is independent of $a$, and hence
\[
H^*(\bar U^*,V)\xrightarrow{\ \sim\ }
  \bigl(R^*\Psi_\eta(V)\bigr)_0.
\]
The fundamental group $\mathbb Z^d$ of
$X^*=(D^*)^d\times D^{e-d}$ (which is also that of $U^*$) is
abelian.  Therefore $V$ has a finite filtration whose successive
quotients are rank-one local systems.  By devissage, we may thus
suppose that $V$ has rank one.

\SGAvar{For $\widetilde{\mathcal V}$ as in 4.14, write
$j_*^m\Omega_{X^*}^*(V)$ for the subcomplex of
$j_*\Omega_{X^*}^*(V)$ whose degree-$p$ term is
$(\Omega_X^p\otimes\widetilde{\mathcal V})^m$.  By [1], II 6.4.1
and 6.8, the natural map induced by the quasi-isomorphism
$V\xrightarrow{\ \approx\ }\Omega_{X^*}^*(V)$ (the holomorphic
Poincare lemma),
\begin{equation}\tag{4.18.1}
j_*^m\Omega_{X^*}^*(V)\longrightarrow Rj_*V,
\end{equation}
is a quasi-isomorphism.  Thus
\begin{equation}\tag{4.18.2}
H^*\bigl((j_*^m\Omega_{X^*}^*(V))_0\bigr)
 \xrightarrow{\ \sim\ }(R^*j_*V)_0
 \xleftarrow{\ \sim\ }H^*(U,V).
\end{equation}

}{For $\widetilde{\mathcal V}$ as in 4.14, write
$j_*^m\Omega_{X^*}^*(V)$ for the subcomplex of
$j_*\Omega_{X^*}^*(V)$ whose degree-$p$ term is
$(\Omega_X^p\otimes\widetilde{\mathcal V})^m$.  By [1], II 6.4.1
and 6.8, the natural map induced by the quasi-isomorphism
$V\xrightarrow{\ \approx\ }\Omega_{X^*}^*(V)$ (the holomorphic
Poincare lemma),
\begin{equation}\tag{4.18.1}
j_*^m\Omega_{X^*}^*(V)\longrightarrow Rj_*V,
\end{equation}
is a quasi-isomorphism.  Thus
\begin{equation}\tag{4.18.2}
H^*\bigl((j_*^m\Omega_{X^*}^*(V))_0\bigr)
 \xrightarrow{\ \sim\ }(R^*j_*V)_0
 \xleftarrow{\ \sim\ }H^*(U^*,V).
\end{equation}

}{E7II-XIV-022}This will be the only analytic ingredient in this part 4.18 of the
proof.

\subsubsection*{Lemma 4.18.3.}
\underline{For a rank-one $V$, one has}

\noindent\textup{(i)}\qquad
\uline{if $V$ is nontrivial:}\quad $H^*(U^*,V)=0$,

\SGApage{II}{\(\approx\)160}
\medskip
\noindent\textup{(ii)}\qquad
\uline{if $V=\underline{\mathbb C}$:}\quad
$H^*(U^*,V)=\overset{*}{\Lambda}\mathbb C^d$.

Indeed, $U^*$ has the homotopy type of a torus.

Let $\eta$ be the generator of $H^1(D^*,\mathbb Z)$, and also its
inverse image in $H^1(U^*,\mathbb Z)$.

\SGAvar{\subsubsection*{Lemma 4.18.4.}
\underline{For a rank-one $V$, the natural map}
\[
H^*(U^*,\mathbb C)
 \xrightarrow[(4.18.1)]{\ \sim\ }
H^*\bigl((j_*^m\Omega_{X^*}^*(V))_0\bigr)
 \longrightarrow
H^*\bigl(\Psi_\eta^{mu}(\Omega_{X^*}^*(V))_0\bigr)
\]
\underline{identifies
$H^*(\Psi_\eta^{mu}(\Omega_{X^*}^*(V))_0)$ with
$H^*(U^*,V)/\eta H^*(U^*,V)$.}


}{\subsubsection*{Lemma 4.18.4.}
\underline{For a rank-one $V$, the natural map}
\[
H^*(U^*,V)
 \xrightarrow[(4.18.1)]{\ \sim\ }
H^*\bigl((j_*^m\Omega_{X^*}^*(V))_0\bigr)
 \longrightarrow
H^*\bigl(\Psi_\eta^{mu}(\Omega_{X^*}^*(V))_0\bigr)
\]
\underline{identifies
$H^*(\Psi_\eta^{mu}(\Omega_{X^*}^*(V))_0)$ with
$H^*(U^*,V)/\eta H^*(U^*,V)$.}


}{E7II-XIV-023}Use 4.12(ii) and the identity
$d\log f=\sum_i n_i\,dz_i/z_i$.  The complex
$\Psi_\eta^{mu}(\Omega_{X^*}^*(V))_0$ identifies with the simple
complex $\underline{\underline{\mathrm{s}}}K$ associated (using direct sums) with the double
complex $K$ whose terms are
\[
K^{-p,n+p}=\frac{(\log f)^p}{p!}\cdot
  (j_*^m\Omega_{X^*}^n(V))_0.
\]
This complex lies in a bad quadrant
$(p\leq0,q\geq0,p+q\geq0)$.  Nevertheless, $\underline{\underline{\mathrm{s}}}K$ is the
inductive limit of the subcomplexes
$\underline{\underline{\mathrm{s}}}\,\sigma_{p\geq p_0}(K)$ whose terms are
\[
(\underline{\underline{\mathrm{s}}}\,\sigma_{p\geq p_0}K)^n
=\bigoplus_{\substack{p+q=n\\p\geq p_0}}K^{pq},
\]
and this suffices to justify the argument below.

The spectral sequence for the filtration by the first degree is
\[
E_1^{-p,n+p}=H^n(U^*,V)\qquad(p\geq0),
\]
with
\[
d_1=2\pi i\,\eta\wedge:
H^n(U^*,V)\longrightarrow H^{n+1}(U^*,V).
\]
If $V$ is nontrivial, then $E_1=0$ by 4.18.3, and 4.18.4 reduces
to $0=0$.  Suppose that $V=\underline{\mathbb C}$.  The complexes
$(E_1,d_1)$ are then Koszul complexes, possibly truncated, so that,
by acyclicity of the Koszul complex,
\[
E_2^{pq}=0\qquad\text{for }p\ne0,
\]
and
\[
E_2^{0n}=H^n(U^*,V)/\eta H^{n-1}(U^*,V),
\]
which proves 4.18.4.

\SGApage{II}{\(\approx\)161}
For $\beta\in\mathbb C^*$, let $V_\beta$ denote the rank-one local
system on $D^*$ whose inverse image on $\widetilde D^*$ is
trivialized and whose monodromy is $\beta$.  We use the same notation
for its inverse image on $U^*$.  On $\bar U^*$, the inverse image of
$V\otimes V_\beta$ is canonically isomorphic to that of $V$.

\subsubsection*{Lemma 4.18.5.}
\underline{For a rank-one $V$, the map}
\[
\bigoplus_{\beta\in\mathbb C^*}
H^*(U^*,V\otimes V_\beta)/
 \eta H^*(U^*,V\otimes V_\beta)
\longrightarrow H^*(\bar U^*,V)
\]
\underline{is bijective.}

\SGAvar{Let $k$ be the greatest common divisor of the $n_i$.  The space
$\bar U^*$ has $k$ connected components, indexed by the $k$th roots
of unity.  On the component $\bar U_\xi^*$ indexed by $\xi$, the
inverse image of the function $z^{1/k}$ on $\widetilde D^*$
coincides with the inverse image of the function
\[
\xi:\prod_{i=1}^d z_i^{n_i/k}
\]
on $U^*$.  The component $\bar U_\xi^*$ has the homotopy type of a
torus, and the sequence
\[
0\longrightarrow\pi_1(\bar U_\xi^*)
 \longrightarrow\pi_1(U^*)\longrightarrow\pi_1(D^*)
\]
is exact.  Write $\varphi:\mathbb Z^d\to\mathbb Z$ for the map
$\pi_1(U^*)\to\pi_1(D^*)$.  If $\ker(\varphi)$ acts nontrivially on
$V$, both sides of 4.18.4 vanish.
Otherwise, after replacing $V$ by a suitable $V\otimes V_\beta$, we
may suppose that $V$ is trivial, that is, $V=\mathbb C$.  In complex
cohomology we then have
\[
H^0(\bar U^*)\otimes
 H^*(U^*)/\eta H^*(U^*)
 \xrightarrow{\ \sim\ }H^*(\bar U^*).
\]
The result follows on observing that $V_\beta$ is trivial on $U^*$
when $\beta^k=1$, that
\[
\bigoplus_{\beta^k=1}H^0(U^*,V_\beta)
=\bigoplus_{\beta^k=1}\mathbb C
 \xrightarrow{\ \sim\ }H^0(\bar U^*,\mathbb C),
\]
and that $H^*(U^*,V_\beta)=0$ when $\beta^k\ne1$.

}{Let $k$ be the greatest common divisor of the $n_i$.  The space
$\bar U^*$ has $k$ connected components, indexed by the $k$th roots
of unity.  On the component $\bar U_\xi^*$ indexed by $\xi$, the
inverse image of the function $z^{1/k}$ on $\widetilde D^*$
coincides with the inverse image of the function
\[
\xi\prod_{i=1}^d z_i^{n_i/k}
\]
on $U^*$.  The component $\bar U_\xi^*$ has the homotopy type of a
torus, and the sequence
\[
0\longrightarrow\pi_1(\bar U_\xi^*)
 \longrightarrow\pi_1(U^*)\longrightarrow\pi_1(D^*)
\]
is exact.  Write $\varphi:\mathbb Z^d\to\mathbb Z$ for the map
$\pi_1(U^*)\to\pi_1(D^*)$.  If $\ker(\varphi)$ acts nontrivially on
$V$, both sides of 4.18.5 vanish.
Otherwise, after replacing $V$ by a suitable $V\otimes V_\beta$, we
may suppose that $V$ is trivial, that is, $V=\mathbb C$.  In complex
cohomology we then have
\[
H^0(\bar U^*)\otimes
 H^*(U^*)/\eta H^*(U^*)
 \xrightarrow{\ \sim\ }H^*(\bar U^*).
\]
The result follows on observing that $V_\beta$ is trivial on $U^*$
when $\beta^k=1$, that
\[
\bigoplus_{\beta^k=1}H^0(U^*,V_\beta)
=\bigoplus_{\beta^k=1}\mathbb C
 \xrightarrow{\ \sim\ }H^0(\bar U^*,\mathbb C),
\]
and that $H^*(U^*,V_\beta)=0$ when $\beta^k\ne1$.

}{E7II-XIV-024}\SGApage{II}{162}
\subsubsection*{Remark 4.18.6.}
If $V$ is quasi-unipotent, the sum over $\mathbb C^*$ in 4.18.5 may
be replaced by a sum over the roots of unity; the proof is the same.

Applying 4.12 gives a commutative diagram
\begin{center}
\begin{tikzcd}[
  column sep=2.2em,
  row sep=large,
  cells={nodes={font=\scriptsize}}
]
H^*\Psi_\eta^m(\Omega_{X^*}^*(V))_0
  \arrow[d]
\& \displaystyle\bigoplus_\beta
  H^*\Psi_\eta^{mu}(\Omega_{X^*}^*(V\otimes V_\beta))
  \arrow[l,"\sim"',"4.12"]
\& \displaystyle\bigoplus_\beta
  H^*(U^*,V\otimes V_\beta)/
    \eta H^*(U^*,V\otimes V_\beta)
  \arrow[l,"\sim"',"4.18.4"]
  \arrow[d,"\wr"',"4.18.5"]
\\
H^*R\Psi(V)_0
\& {}
\& H^*(\bar U^*,V)\arrow[ll,"\sim"']
\end{tikzcd}
\end{center}
from which 4.15 follows.  Proposition 4.17 is obtained by replacing
$\Psi^m$ with $\Psi^{mqu}$, the sum over $\beta$ with a sum over the
roots of unity, and 4.18.5 with 4.18.6.


\subsubsection*{4.19.}
Let us deduce 4.15 and 4.17 from the special case 4.18.  Let
$u:Y\longrightarrow X$ be a resolution of $(X,X_0)$, that is, a
proper morphism inducing an isomorphism
$Y^*\xrightarrow{\ \sim\ }X^*$, such that $Y$ is smooth and
$(Y_0)_{\mathrm{red}}$ is a normal-crossings divisor in $Y$.  If $V$
is quasi-unipotent along $X_0$, its inverse image $u^*V$ is
quasi-unipotent along $Y_0$.

For every coherent sheaf $F$ on $X$, denote by $F^m$, or by
$j_*^m j^*F$, the subsheaf of $j_*j^*F$ formed by the meromorphic
sections.  If $f$ is an equation for $X_0$, then $j_*^m j^*F$ is the
inductive limit of the system
\begin{equation}\tag{4.19.1}
F\xrightarrow{f}F\xrightarrow{f}F\longrightarrow\cdots.
\end{equation}
The following lemma is proved in [2].

\subsubsection*{Lemma 4.19.2.}
\underline{For every coherent sheaf $F$ on $X$, one has}
\[
F^m\xrightarrow{\ \sim\ }u_*(u^*F)^m
\]
\underline{and $R^iu_*(u^*F)^m=0$ for $i>0$.}

\SGApage{II}{\(\approx\)163}
\SGAvar{By Grauert's theorem, the $R^iu_*u^*F$ are coherent sheaves.  The
kernels and cokernels of $F\longrightarrow u_*u^*F$, as well as the
$R^iu_*F$ for $i>0$, vanish outside $X_0$ and are therefore
locally annihilated by a power of $f$.  Since $u$ is proper,
$R^iu_*$ commutes with filtered inductive limits, in particular with
the inductive limit (4.19.1), and (4.19.2) follows.

}{By Grauert's theorem, the $R^iu_*u^*F$ are coherent sheaves.  The
kernels and cokernels of $F\longrightarrow u_*u^*F$, as well as the
$R^iu_*u^*F$ for $i>0$, vanish outside $X_0$ and are therefore
locally annihilated by a power of $f$.  Since $u$ is proper,
$R^iu_*$ commutes with filtered inductive limits, in particular with
the inductive limit (4.19.1), and (4.19.2) follows.

}{E7II-XIV-025}From 4.19.2 and 4.12 one deduces that, under the hypotheses just
stated and still denoting by $u$ the map $Y_0\to X_0$ induced by
$u$, one has the following.

\subsubsection*{Lemma 4.19.3.}
\underline{The morphisms}
\[
\Psi_\eta^m(F)\longrightarrow
u_*\Psi_\eta^m(u^*F)\longrightarrow
Ru_*\Psi_\eta^m(u^*F)
\]
\underline{are isomorphisms.}

That the second arrow is an isomorphism means that
\[
R^iu_*\Psi_\eta^m(u^*F)=0\qquad(i>0).
\]
This follows from 4.12 and from
$R^iu_*(i^*u^*F^m)=0$ for $i>0$, which in turn follows from
$R^iu_*(u^*F^m)=0$ for $i>0$ and the proper base-change theorem for
$u$.

To prove 4.15, it remains to consider the commutative diagram
\begin{center}
\begin{tikzcd}[
  column sep=large,
  row sep=large,
  cells={nodes={font=\small}}
]
\Psi^m(\Omega_{X^*}^*(V))
  \arrow[d]
  \arrow[r,equal,"(4.19.3)"']
\& u_*\Psi^m(\Omega_{Y^*}^*(u^*V))
  \arrow[r,"\sim","(4.19.3)"']
\& Ru_*\Psi^m(\Omega_{Y^*}^*(u^*V))
  \arrow[d,"\wr"',"(4.18)"]
\\
R\Psi(V)
  \arrow[rr,"(1.3.1;\ \mathrm{XIII}\ 2.1.7.1)"]
\& {}
\& Ru_*R\Psi(u^*V)
\end{tikzcd}
\end{center}
The proof of 4.17 is the same, with $\Psi^m$ replaced everywhere by
$\Psi^{mqu}$.

\subsubsection*{Remark 4.20.}
Theorem 4.13 gives an incarnation of the monodromy theorem at the
level of complexes.

\subsection*{Bibliography}

\begin{enumerate}[label={[\arabic*]},leftmargin=3em]
\item P. Deligne, \emph{Equations diff\'erentielles \`a points
  singuliers r\'eguliers}, Lecture Notes in Mathematics 163,
  Springer-Verlag.
\item A. Grothendieck, \emph{On the De Rham cohomology of algebraic
  varieties}, Publ. Math. IHES \underline{29}, 95--103 (1966).
\item K. Lamotke, \emph{Isolated critical points and monodromy}.
\end{enumerate}


\markboth{Exposé XIV}{Exposé XIV}

\ifnum\SGAform=1
\clearpage\section*{Critical apparatus}\begingroup\small\raggedright\sloppy
Readings are given in the exact TeX notation of the two English witnesses. Decimal suffixes identify the individual differences in each marked passage.

\par\medskip\noindent\hypertarget{app-E7II-XIV-001.1}{}\textbf{E7II-XIV-001.1}\quad printed folio 118.\par
\noindent Printed reading: \emph{(empty)}\par
\noindent Corrected reading: {\ttfamily\footnotesize \textasciicircum{}*\allowbreak{}}\par
\noindent Source: {\ttfamily\footnotesize fr/sga7ii/\{diplomati\allowbreak{}c,current\}/expose\_XI\allowbreak{}V\_body.tex;\allowbreak{} folio\allowbreak{} 118;\allowbreak{} formula\allowbreak{} 76;\allowbreak{} port/decisions/decis\allowbreak{}ions\_en7II\_r2\_001.js\allowbreak{}on\#C587;\allowbreak{} fixes/FIXLOG\_7II\_CUR\allowbreak{}RENT.jsonl\#C587;\allowbreak{} 7II:184\allowbreak{}}\par

\par\medskip\noindent\hypertarget{app-E7II-XIV-002.1}{}\textbf{E7II-XIV-002.1}\quad printed folio 121.\par
\noindent Printed reading: {\ttfamily\footnotesize 1\allowbreak{}}\par
\noindent Corrected reading: {\ttfamily\footnotesize 2\allowbreak{}}\par
\noindent Source: {\ttfamily\footnotesize port/PORT\_EN7II\_LOG.\allowbreak{}jsonl\#C591;\allowbreak{} fixes/FIXLOG\_7II\_CUR\allowbreak{}RENT.jsonl\#C591;\allowbreak{} 7II:188\allowbreak{}}\par

\par\medskip\noindent\hypertarget{app-E7II-XIV-003.1}{}\textbf{E7II-XIV-003.1}\quad printed folio 124.\par
\noindent Printed reading: \emph{(2 espace(s), 0 retour(s) de ligne dans le témoin TeX)}\par
\noindent Corrected reading: {\ttfamily\footnotesize  \textbackslash{}overline\allowbreak{} }\par
\noindent Source: {\ttfamily\footnotesize fr/sga7ii/\{diplomati\allowbreak{}c,current\}/expose\_XI\allowbreak{}V\_body.tex;\allowbreak{} folio\allowbreak{} 124;\allowbreak{} formula\allowbreak{} 375\allowbreak{}}\par

\par\medskip\noindent\hypertarget{app-E7II-XIV-004.1}{}\textbf{E7II-XIV-004.1}\quad printed folio 127.\par
\noindent Printed reading: {\ttfamily\footnotesize 1\allowbreak{}}\par
\noindent Corrected reading: {\ttfamily\footnotesize 6\allowbreak{}}\par
\noindent Source: {\ttfamily\footnotesize ED-EN-NONMATH-G003;\allowbreak{} F7II-XIV-010.1;\allowbreak{} ADOPT-F7II-XIV-010.1\allowbreak{};\allowbreak{} fr/sga7ii/\{diplomati\allowbreak{}c,current\}/expose\_XI\allowbreak{}V\_body.tex;\allowbreak{} folio\allowbreak{} 127\allowbreak{}}\par

\par\medskip\noindent\hypertarget{app-E7II-XIV-005.1}{}\textbf{E7II-XIV-005.1}\quad printed folio 136.\par
\noindent Printed reading: {\ttfamily\footnotesize  r\allowbreak{}}\par
\noindent Corrected reading: {\ttfamily\footnotesize \textbackslash{}tau\allowbreak{}}\par
\noindent Source: {\ttfamily\footnotesize fr/sga7ii/\{diplomati\allowbreak{}c,current\}/expose\_XI\allowbreak{}V\_body.tex;\allowbreak{} folio\allowbreak{} 136;\allowbreak{} formula\allowbreak{} 833\allowbreak{}}\par

\par\medskip\noindent\hypertarget{app-E7II-XIV-005.2}{}\textbf{E7II-XIV-005.2}\quad printed folio 137.\par
\noindent Printed reading: {\ttfamily\footnotesize s\allowbreak{}}\par
\noindent Corrected reading: {\ttfamily\footnotesize t\allowbreak{}}\par
\noindent Source: {\ttfamily\footnotesize port/decisions/decis\allowbreak{}ions\_en7II\_r2\_001.js\allowbreak{}on\#C607;\allowbreak{} fixes/FIXLOG\_7II\_CUR\allowbreak{}RENT.jsonl\#C607;\allowbreak{} 7II:206\allowbreak{}}\par

\par\medskip\noindent\hypertarget{app-E7II-XIV-005.3}{}\textbf{E7II-XIV-005.3}\quad printed folio 137.\par
\noindent Printed reading: {\ttfamily\footnotesize t\allowbreak{}}\par
\noindent Corrected reading: {\ttfamily\footnotesize d\allowbreak{}}\par
\noindent Source: {\ttfamily\footnotesize port/decisions/decis\allowbreak{}ions\_en7II\_r2\_001.js\allowbreak{}on\#C607;\allowbreak{} fixes/FIXLOG\_7II\_CUR\allowbreak{}RENT.jsonl\#C607;\allowbreak{} 7II:206\allowbreak{}}\par

\par\medskip\noindent\hypertarget{app-E7II-XIV-006.1}{}\textbf{E7II-XIV-006.1}\quad printed folio 136.\par
\noindent Printed reading: {\ttfamily\footnotesize s\allowbreak{}}\par
\noindent Corrected reading: {\ttfamily\footnotesize d\allowbreak{}}\par
\noindent Source: {\ttfamily\footnotesize fr/sga7ii/\{diplomati\allowbreak{}c,current\}/expose\_XI\allowbreak{}V\_body.tex;\allowbreak{} folio\allowbreak{} 136;\allowbreak{} formula\allowbreak{} 868;\allowbreak{} port/decisions/decis\allowbreak{}ions\_en7II\_r2\_001.js\allowbreak{}on\#C603;\allowbreak{} fixes/FIXLOG\_7II\_CUR\allowbreak{}RENT.jsonl\#C603;\allowbreak{} 7II:202\allowbreak{}}\par

\par\medskip\noindent\hypertarget{app-E7II-XIV-007.1}{}\textbf{E7II-XIV-007.1}\quad printed folio 137.\par
\noindent Printed reading: {\ttfamily\footnotesize s\allowbreak{}}\par
\noindent Corrected reading: {\ttfamily\footnotesize t\allowbreak{}}\par
\noindent Source: {\ttfamily\footnotesize fr/sga7ii/\{diplomati\allowbreak{}c,current\}/expose\_XI\allowbreak{}V\_body.tex;\allowbreak{} folio\allowbreak{} 137\allowbreak{}}\par

\par\medskip\noindent\hypertarget{app-E7II-XIV-007.2}{}\textbf{E7II-XIV-007.2}\quad printed folio 137.\par
\noindent Printed reading: {\ttfamily\footnotesize t\allowbreak{}}\par
\noindent Corrected reading: {\ttfamily\footnotesize d\allowbreak{}}\par
\noindent Source: {\ttfamily\footnotesize fr/sga7ii/\{diplomati\allowbreak{}c,current\}/expose\_XI\allowbreak{}V\_body.tex;\allowbreak{} folio\allowbreak{} 137\allowbreak{}}\par

\par\medskip\noindent\hypertarget{app-E7II-XIV-008.1}{}\textbf{E7II-XIV-008.1}\quad printed folio 137.\par
\noindent Printed reading: {\ttfamily\footnotesize |\allowbreak{}}\par
\noindent Corrected reading: \emph{(empty)}\par
\noindent Source: {\ttfamily\footnotesize ED-EN-FORMULA-G111;\allowbreak{} fr/sga7ii/\{diplomati\allowbreak{}c,current\}/expose\_XI\allowbreak{}V\_body.tex;\allowbreak{} folio\allowbreak{} 137;\allowbreak{} port/decisions/decis\allowbreak{}ions\_en7II\_r2\_001.js\allowbreak{}on\#C606;\allowbreak{} fixes/FIXLOG\_7II\_CUR\allowbreak{}RENT.jsonl\#C606;\allowbreak{} 7II:205\allowbreak{}}\par

\par\medskip\noindent\hypertarget{app-E7II-XIV-008.2}{}\textbf{E7II-XIV-008.2}\quad printed folio 137.\par
\noindent Printed reading: {\ttfamily\footnotesize |\allowbreak{}}\par
\noindent Corrected reading: \emph{(empty)}\par
\noindent Source: {\ttfamily\footnotesize ED-EN-FORMULA-G111;\allowbreak{} fr/sga7ii/\{diplomati\allowbreak{}c,current\}/expose\_XI\allowbreak{}V\_body.tex;\allowbreak{} folio\allowbreak{} 137;\allowbreak{} port/decisions/decis\allowbreak{}ions\_en7II\_r2\_001.js\allowbreak{}on\#C606;\allowbreak{} fixes/FIXLOG\_7II\_CUR\allowbreak{}RENT.jsonl\#C606;\allowbreak{} 7II:205\allowbreak{}}\par

\par\medskip\noindent\hypertarget{app-E7II-XIV-009.1}{}\textbf{E7II-XIV-009.1}\quad printed folio 139.\par
\noindent Printed reading: \emph{(empty)}\par
\noindent Corrected reading: {\ttfamily\footnotesize f\allowbreak{}}\par
\noindent Source: {\ttfamily\footnotesize fr/sga7ii/\{diplomati\allowbreak{}c,current\}/expose\_XI\allowbreak{}V\_body.tex;\allowbreak{} folio\allowbreak{} 139;\allowbreak{} port/decisions/decis\allowbreak{}ions\_en7II\_r2\_001.js\allowbreak{}on\#C610;\allowbreak{} fixes/FIXLOG\_7II\_CUR\allowbreak{}RENT.jsonl\#C610;\allowbreak{} 7II:210\allowbreak{}}\par

\par\medskip\noindent\hypertarget{app-E7II-XIV-010.1}{}\textbf{E7II-XIV-010.1}\quad printed folio 140.\par
\noindent Printed reading: {\ttfamily\footnotesize n\allowbreak{}}\par
\noindent Corrected reading: {\ttfamily\footnotesize x\allowbreak{}}\par
\noindent Source: {\ttfamily\footnotesize fr/sga7ii/\{diplomati\allowbreak{}c,current\}/expose\_XI\allowbreak{}V\_body.tex;\allowbreak{} folio\allowbreak{} 140;\allowbreak{} formula\allowbreak{} 1033;\allowbreak{} port/decisions/decis\allowbreak{}ions\_en7II\_r2\_001.js\allowbreak{}on\#C611;\allowbreak{} fixes/FIXLOG\_7II\_CUR\allowbreak{}RENT.jsonl\#C611;\allowbreak{} 7II:211\allowbreak{}}\par

\par\medskip\noindent\hypertarget{app-E7II-XIV-010.2}{}\textbf{E7II-XIV-010.2}\quad printed folio 140.\par
\noindent Printed reading: {\ttfamily\footnotesize n\allowbreak{}}\par
\noindent Corrected reading: {\ttfamily\footnotesize x\allowbreak{}}\par
\noindent Source: {\ttfamily\footnotesize fr/sga7ii/\{diplomati\allowbreak{}c,current\}/expose\_XI\allowbreak{}V\_body.tex;\allowbreak{} folio\allowbreak{} 140;\allowbreak{} formula\allowbreak{} 1033;\allowbreak{} port/decisions/decis\allowbreak{}ions\_en7II\_r2\_001.js\allowbreak{}on\#C611;\allowbreak{} fixes/FIXLOG\_7II\_CUR\allowbreak{}RENT.jsonl\#C611;\allowbreak{} 7II:211\allowbreak{}}\par

\par\medskip\noindent\hypertarget{app-E7II-XIV-010.3}{}\textbf{E7II-XIV-010.3}\quad printed folio 140.\par
\noindent Printed reading: \emph{(empty)}\par
\noindent Corrected reading: {\ttfamily\footnotesize \textbackslash{}widetilde\allowbreak{}}\par
\noindent Source: {\ttfamily\footnotesize port/decisions/decis\allowbreak{}ions\_en7II\_r2\_001.js\allowbreak{}on\#C612;\allowbreak{} fixes/FIXLOG\_7II\_CUR\allowbreak{}RENT.jsonl\#C612;\allowbreak{} 7II:212\allowbreak{}}\par

\par\medskip\noindent\hypertarget{app-E7II-XIV-010.4}{}\textbf{E7II-XIV-010.4}\quad printed folio 140.\par
\noindent Printed reading: \emph{(empty)}\par
\noindent Corrected reading: {\ttfamily\footnotesize \textbackslash{}widetilde\allowbreak{} }\par
\noindent Source: {\ttfamily\footnotesize ED-EN-FORMULA-G113;\allowbreak{} fr/sga7ii/\{diplomati\allowbreak{}c,current\}/expose\_XI\allowbreak{}V\_body.tex;\allowbreak{} folio\allowbreak{} 140\allowbreak{}}\par

\par\medskip\noindent\hypertarget{app-E7II-XIV-011.1}{}\textbf{E7II-XIV-011.1}\quad printed folio 144.\par
\noindent Printed reading: {\ttfamily\footnotesize \{\}\_1\allowbreak{}}\par
\noindent Corrected reading: \emph{(empty)}\par
\noindent Source: {\ttfamily\footnotesize port/decisions/decis\allowbreak{}ions\_en7II\_r2\_001.js\allowbreak{}on\#C617;\allowbreak{} 7II:217\allowbreak{}}\par

\par\medskip\noindent\hypertarget{app-E7II-XIV-011.2}{}\textbf{E7II-XIV-011.2}\quad printed folio 144.\par
\noindent Printed reading: \emph{(empty)}\par
\noindent Corrected reading: {\ttfamily\footnotesize \textasciicircum{}*\allowbreak{}}\par
\noindent Source: {\ttfamily\footnotesize fr/sga7ii/\{diplomati\allowbreak{}c,current\}/expose\_XI\allowbreak{}V\_body.tex;\allowbreak{} folio\allowbreak{} 144;\allowbreak{} formula\allowbreak{} 1152;\allowbreak{} port/PORT\_EN7II\_LOG.\allowbreak{}jsonl\#C619;\allowbreak{} fixes/FIXLOG\_7II\_CUR\allowbreak{}RENT.jsonl\#C619;\allowbreak{} 7II:219\allowbreak{}}\par

\par\medskip\noindent\hypertarget{app-E7II-XIV-011.3}{}\textbf{E7II-XIV-011.3}\quad printed folio 144.\par
\noindent Printed reading: {\ttfamily\footnotesize 15\allowbreak{}}\par
\noindent Corrected reading: {\ttfamily\footnotesize 1.5\allowbreak{}}\par
\noindent Source: {\ttfamily\footnotesize port/decisions/decis\allowbreak{}ions\_en7II\_r2\_001.js\allowbreak{}on\#C620;\allowbreak{} fixes/FIXLOG\_7II\_CUR\allowbreak{}RENT.jsonl\#C620;\allowbreak{} 7II:220\allowbreak{}}\par

\par\medskip\noindent\hypertarget{app-E7II-XIV-012.1}{}\textbf{E7II-XIV-012.1}\quad printed folio 146.\par
\noindent Printed reading: \emph{(empty)}\par
\noindent Corrected reading: {\ttfamily\footnotesize \_0\allowbreak{}}\par
\noindent Source: {\ttfamily\footnotesize port/decisions/decis\allowbreak{}ions\_en7II\_r2\_001.js\allowbreak{}on\#C621;\allowbreak{} 7II:221\allowbreak{}}\par

\par\medskip\noindent\hypertarget{app-E7II-XIV-013.1}{}\textbf{E7II-XIV-013.1}\quad printed folio 147.\par
\noindent Printed reading: \emph{(empty)}\par
\noindent Corrected reading: {\ttfamily\footnotesize +1\allowbreak{}}\par
\noindent Source: {\ttfamily\footnotesize fr/sga7ii/\{diplomati\allowbreak{}c,current\}/expose\_XI\allowbreak{}V\_body.tex;\allowbreak{} folio\allowbreak{} 147;\allowbreak{} formula\allowbreak{} 1257\allowbreak{}}\par

\par\medskip\noindent\hypertarget{app-E7II-XIV-014.1}{}\textbf{E7II-XIV-014.1}\quad printed folio 148.\par
\noindent Printed reading: {\ttfamily\footnotesize 13\allowbreak{}}\par
\noindent Corrected reading: {\ttfamily\footnotesize 12\allowbreak{}}\par
\noindent Source: {\ttfamily\footnotesize port/PORT\_EN7II\_LOG.\allowbreak{}jsonl\#C625;\allowbreak{} fixes/FIXLOG\_7II\_CUR\allowbreak{}RENT.jsonl\#C625;\allowbreak{} 7II:227\allowbreak{}}\par

\par\medskip\noindent\hypertarget{app-E7II-XIV-014.2}{}\textbf{E7II-XIV-014.2}\quad printed folio 148.\par
\noindent Printed reading: \emph{(empty)}\par
\noindent Corrected reading: {\ttfamily\footnotesize \textbackslash{}bigr)\allowbreak{}}\par
\noindent Source: {\ttfamily\footnotesize fr/sga7ii/\{diplomati\allowbreak{}c,current\}/expose\_XI\allowbreak{}V\_body.tex;\allowbreak{} folio\allowbreak{} 148;\allowbreak{} formula\allowbreak{} 1321;\allowbreak{} port/decisions/decis\allowbreak{}ions\_en7II\_r2\_002.js\allowbreak{}on\#C627;\allowbreak{} fixes/FIXLOG\_7II\_CUR\allowbreak{}RENT.jsonl\#C627;\allowbreak{} 7II:229\allowbreak{}}\par

\par\medskip\noindent\hypertarget{app-E7II-XIV-014.3}{}\textbf{E7II-XIV-014.3}\quad printed folio 148.\par
\noindent Printed reading: {\ttfamily\footnotesize x\allowbreak{}}\par
\noindent Corrected reading: {\ttfamily\footnotesize d\allowbreak{}}\par
\noindent Source: {\ttfamily\footnotesize fr/sga7ii/\{diplomati\allowbreak{}c,current\}/expose\_XI\allowbreak{}V\_body.tex;\allowbreak{} folio\allowbreak{} 148;\allowbreak{} formula\allowbreak{} 1325\allowbreak{}}\par

\par\medskip\noindent\hypertarget{app-E7II-XIV-014.4}{}\textbf{E7II-XIV-014.4}\quad printed folio 149.\par
\noindent Printed reading: {\ttfamily\footnotesize n\allowbreak{}}\par
\noindent Corrected reading: {\ttfamily\footnotesize d\allowbreak{}}\par
\noindent Source: {\ttfamily\footnotesize fr/sga7ii/\{diplomati\allowbreak{}c,current\}/expose\_XI\allowbreak{}V\_body.tex;\allowbreak{} folio\allowbreak{} 149;\allowbreak{} formula\allowbreak{} 1327\allowbreak{}}\par

\par\medskip\noindent\hypertarget{app-E7II-XIV-014.5}{}\textbf{E7II-XIV-014.5}\quad printed folio 149.\par
\noindent Printed reading: {\ttfamily\footnotesize n\allowbreak{}}\par
\noindent Corrected reading: {\ttfamily\footnotesize d\allowbreak{}}\par
\noindent Source: {\ttfamily\footnotesize fr/sga7ii/\{diplomati\allowbreak{}c,current\}/expose\_XI\allowbreak{}V\_body.tex;\allowbreak{} folio\allowbreak{} 149;\allowbreak{} formula\allowbreak{} 1329\allowbreak{}}\par

\par\medskip\noindent\hypertarget{app-E7II-XIV-014.6}{}\textbf{E7II-XIV-014.6}\quad printed folio 149.\par
\noindent Printed reading: {\ttfamily\footnotesize n\allowbreak{}}\par
\noindent Corrected reading: {\ttfamily\footnotesize d\allowbreak{}}\par
\noindent Source: {\ttfamily\footnotesize fr/sga7ii/\{diplomati\allowbreak{}c,current\}/expose\_XI\allowbreak{}V\_body.tex;\allowbreak{} folio\allowbreak{} 149;\allowbreak{} formula\allowbreak{} 1329\allowbreak{}}\par

\par\medskip\noindent\hypertarget{app-E7II-XIV-015.1}{}\textbf{E7II-XIV-015.1}\quad printed folio 150.\par
\noindent Printed reading: {\ttfamily\footnotesize bundle\allowbreak{}}\par
\noindent Corrected reading: {\ttfamily\footnotesize fiber\allowbreak{}}\par
\noindent Source: {\ttfamily\footnotesize ED-EN-NONMATH-G047;\allowbreak{} F7II-XIV-039.1;\allowbreak{} 7II:231;\allowbreak{} fr/sga7ii/\{diplomati\allowbreak{}c,current\}/expose\_XI\allowbreak{}V\_body.tex;\allowbreak{} folio\allowbreak{} 150\allowbreak{}}\par

\par\medskip\noindent\hypertarget{app-E7II-XIV-015.2}{}\textbf{E7II-XIV-015.2}\quad printed folio 150.\par
\noindent Printed reading: \emph{(empty)}\par
\noindent Corrected reading: {\ttfamily\footnotesize \_c\allowbreak{}}\par
\noindent Source: {\ttfamily\footnotesize fr/sga7ii/\{diplomati\allowbreak{}c,current\}/expose\_XI\allowbreak{}V\_body.tex;\allowbreak{} folio\allowbreak{} 150;\allowbreak{} formula\allowbreak{} 1390;\allowbreak{} port/decisions/decis\allowbreak{}ions\_en7II\_r2\_002.js\allowbreak{}on\#C629;\allowbreak{} fixes/FIXLOG\_7II\_CUR\allowbreak{}RENT.jsonl\#C629;\allowbreak{} 7II:232\allowbreak{}}\par

\par\medskip\noindent\hypertarget{app-E7II-XIV-016.1}{}\textbf{E7II-XIV-016.1}\quad printed folio 152.\par
\noindent Printed reading: {\ttfamily\footnotesize \textbackslash{}Omega\textasciicircum{}*\_\{\allowbreak{}}\par
\noindent Corrected reading: \emph{(empty)}\par
\noindent Source: {\ttfamily\footnotesize fr/sga7ii/\{diplomati\allowbreak{}c,current\}/expose\_XI\allowbreak{}V\_body.tex;\allowbreak{} folio\allowbreak{} 152;\allowbreak{} formula\allowbreak{} 1469\allowbreak{}}\par

\par\medskip\noindent\hypertarget{app-E7II-XIV-016.2}{}\textbf{E7II-XIV-016.2}\quad printed folio 152.\par
\noindent Printed reading: {\ttfamily\footnotesize \}\allowbreak{}}\par
\noindent Corrected reading: \emph{(empty)}\par
\noindent Source: {\ttfamily\footnotesize fr/sga7ii/\{diplomati\allowbreak{}c,current\}/expose\_XI\allowbreak{}V\_body.tex;\allowbreak{} folio\allowbreak{} 152;\allowbreak{} formula\allowbreak{} 1469\allowbreak{}}\par

\par\medskip\noindent\hypertarget{app-E7II-XIV-017.1}{}\textbf{E7II-XIV-017.1}\quad printed folio 152.\par
\noindent Printed reading: {\ttfamily\footnotesize 1\allowbreak{}}\par
\noindent Corrected reading: {\ttfamily\footnotesize 2\allowbreak{}}\par
\noindent Source: {\ttfamily\footnotesize ED-EN-NONMATH-G022;\allowbreak{} F7II-XIV-043.1;\allowbreak{} 7II:244;\allowbreak{} fr/sga7ii/\{diplomati\allowbreak{}c,current\}/expose\_XI\allowbreak{}V\_body.tex;\allowbreak{} folio\allowbreak{} 152\allowbreak{}}\par

\par\medskip\noindent\hypertarget{app-E7II-XIV-017.2}{}\textbf{E7II-XIV-017.2}\quad printed folio 152.\par
\noindent Printed reading: {\ttfamily\footnotesize X\allowbreak{}}\par
\noindent Corrected reading: {\ttfamily\footnotesize x\allowbreak{}}\par
\noindent Source: {\ttfamily\footnotesize port/decisions/decis\allowbreak{}ions\_en7II\_r2\_002.js\allowbreak{}on\#C633;\allowbreak{} fixes/FIXLOG\_7II\_CUR\allowbreak{}RENT.jsonl\#C633;\allowbreak{} 7II:246\allowbreak{}}\par

\par\medskip\noindent\hypertarget{app-E7II-XIV-017.3}{}\textbf{E7II-XIV-017.3}\quad printed folio 152.\par
\noindent Printed reading: \emph{(empty)}\par
\noindent Corrected reading: {\ttfamily\footnotesize  one\allowbreak{} may\allowbreak{} take\allowbreak{}}\par
\noindent Source: {\ttfamily\footnotesize ED-EN-FORMULA-G115;\allowbreak{} fr/sga7ii/\{diplomati\allowbreak{}c,current\}/expose\_XI\allowbreak{}V\_body.tex;\allowbreak{} folio\allowbreak{} 152\allowbreak{}}\par

\par\medskip\noindent\hypertarget{app-E7II-XIV-017.4}{}\textbf{E7II-XIV-017.4}\quad printed folio 152.\par
\noindent Printed reading: {\ttfamily\footnotesize  \$X\$\allowbreak{} form\allowbreak{} a\allowbreak{} fundamental\allowbreak{} system\allowbreak{} of\allowbreak{} neighborhoods\allowbreak{} of\allowbreak{}}\par
\noindent Corrected reading: \emph{(empty)}\par
\noindent Source: {\ttfamily\footnotesize ED-EN-FORMULA-G115;\allowbreak{} fr/sga7ii/\{diplomati\allowbreak{}c,current\}/expose\_XI\allowbreak{}V\_body.tex;\allowbreak{} folio\allowbreak{} 152\allowbreak{}}\par

\par\medskip\noindent\hypertarget{app-E7II-XIV-018.1}{}\textbf{E7II-XIV-018.1}\quad printed folio 153.\par
\noindent Printed reading: \emph{(empty)}\par
\noindent Corrected reading: {\ttfamily\footnotesize \textbackslash{}widetilde\allowbreak{}}\par
\noindent Source: {\ttfamily\footnotesize fr/sga7ii/\{diplomati\allowbreak{}c,current\}/expose\_XI\allowbreak{}V\_body.tex;\allowbreak{} folio\allowbreak{} 153;\allowbreak{} formula\allowbreak{} 1509;\allowbreak{} port/decisions/decis\allowbreak{}ions\_en7II\_r2\_002.js\allowbreak{}on\#C634;\allowbreak{} fixes/FIXLOG\_7II\_CUR\allowbreak{}RENT.jsonl\#C634;\allowbreak{} 7II:247\allowbreak{}}\par

\par\medskip\noindent\hypertarget{app-E7II-XIV-019.1}{}\textbf{E7II-XIV-019.1}\quad printed folio 153.\par
\noindent Printed reading: {\ttfamily\footnotesize F\allowbreak{}}\par
\noindent Corrected reading: {\ttfamily\footnotesize f\allowbreak{}}\par
\noindent Source: {\ttfamily\footnotesize fr/sga7ii/\{diplomati\allowbreak{}c,current\}/expose\_XI\allowbreak{}V\_body.tex;\allowbreak{} folio\allowbreak{} 153;\allowbreak{} formula\allowbreak{} 1524\allowbreak{}}\par

\par\medskip\noindent\hypertarget{app-E7II-XIV-020.1}{}\textbf{E7II-XIV-020.1}\quad printed folio 154.\par
\noindent Printed reading: {\ttfamily\footnotesize .\allowbreak{}}\par
\noindent Corrected reading: {\ttfamily\footnotesize (\allowbreak{}}\par
\noindent Source: {\ttfamily\footnotesize fr/sga7ii/\{diplomati\allowbreak{}c,current\}/expose\_XI\allowbreak{}V\_body.tex;\allowbreak{} folio\allowbreak{} 154;\allowbreak{} formula\allowbreak{} 1553\allowbreak{}}\par

\par\medskip\noindent\hypertarget{app-E7II-XIV-021.1}{}\textbf{E7II-XIV-021.1}\quad printed folio 157.\par
\noindent Printed reading: {\ttfamily\footnotesize muq\allowbreak{}}\par
\noindent Corrected reading: {\ttfamily\footnotesize mqu\allowbreak{}}\par
\noindent Source: {\ttfamily\footnotesize fr/sga7ii/\{diplomati\allowbreak{}c,current\}/expose\_XI\allowbreak{}V\_body.tex;\allowbreak{} folio\allowbreak{} 157;\allowbreak{} formula\allowbreak{} 1650;\allowbreak{} port/decisions/decis\allowbreak{}ions\_en7II\_r2\_002.js\allowbreak{}on\#C638;\allowbreak{} fixes/FIXLOG\_7II\_CUR\allowbreak{}RENT.jsonl\#C638;\allowbreak{} 7II:235\allowbreak{}}\par

\par\medskip\noindent\hypertarget{app-E7II-XIV-021.2}{}\textbf{E7II-XIV-021.2}\quad printed folio 157.\par
\noindent Printed reading: {\ttfamily\footnotesize muq\allowbreak{}}\par
\noindent Corrected reading: {\ttfamily\footnotesize mqu\allowbreak{}}\par
\noindent Source: {\ttfamily\footnotesize fr/sga7ii/\{diplomati\allowbreak{}c,current\}/expose\_XI\allowbreak{}V\_body.tex;\allowbreak{} folio\allowbreak{} 157;\allowbreak{} port/decisions/decis\allowbreak{}ions\_en7II\_r2\_002.js\allowbreak{}on\#C639;\allowbreak{} fixes/FIXLOG\_7II\_CUR\allowbreak{}RENT.jsonl\#C639;\allowbreak{} 7II:236\allowbreak{}}\par

\par\medskip\noindent\hypertarget{app-E7II-XIV-021.3}{}\textbf{E7II-XIV-021.3}\quad printed folio 157.\par
\noindent Printed reading: {\ttfamily\footnotesize muq\allowbreak{}}\par
\noindent Corrected reading: {\ttfamily\footnotesize mqu\allowbreak{}}\par
\noindent Source: {\ttfamily\footnotesize ED-EN-FORMULA-G117;\allowbreak{} fr/sga7ii/\{diplomati\allowbreak{}c,current\}/expose\_XI\allowbreak{}V\_body.tex;\allowbreak{} folio\allowbreak{} 157;\allowbreak{} port/decisions/decis\allowbreak{}ions\_en7II\_r2\_002.js\allowbreak{}on\#C640;\allowbreak{} fixes/FIXLOG\_7II\_CUR\allowbreak{}RENT.jsonl\#C640;\allowbreak{} 7II:237\allowbreak{}}\par

\par\medskip\noindent\hypertarget{app-E7II-XIV-021.4}{}\textbf{E7II-XIV-021.4}\quad printed folio 157.\par
\noindent Printed reading: {\ttfamily\footnotesize muq\allowbreak{}}\par
\noindent Corrected reading: {\ttfamily\footnotesize mqu\allowbreak{}}\par
\noindent Source: {\ttfamily\footnotesize ED-EN-FORMULA-G117;\allowbreak{} fr/sga7ii/\{diplomati\allowbreak{}c,current\}/expose\_XI\allowbreak{}V\_body.tex;\allowbreak{} folio\allowbreak{} 157;\allowbreak{} port/decisions/decis\allowbreak{}ions\_en7II\_r2\_002.js\allowbreak{}on\#C640;\allowbreak{} fixes/FIXLOG\_7II\_CUR\allowbreak{}RENT.jsonl\#C640;\allowbreak{} 7II:237\allowbreak{}}\par

\par\medskip\noindent\hypertarget{app-E7II-XIV-022.1}{}\textbf{E7II-XIV-022.1}\quad printed folio 159.\par
\noindent Printed reading: \emph{(empty)}\par
\noindent Corrected reading: {\ttfamily\footnotesize \textasciicircum{}*\allowbreak{}}\par
\noindent Source: {\ttfamily\footnotesize fr/sga7ii/\{diplomati\allowbreak{}c,current\}/expose\_XI\allowbreak{}V\_body.tex;\allowbreak{} folio\allowbreak{} 159;\allowbreak{} formula\allowbreak{} 1744\allowbreak{}}\par

\par\medskip\noindent\hypertarget{app-E7II-XIV-023.1}{}\textbf{E7II-XIV-023.1}\quad printed folio 160.\par
\noindent Printed reading: {\ttfamily\footnotesize \textbackslash{}mathbb\allowbreak{} C\allowbreak{}}\par
\noindent Corrected reading: {\ttfamily\footnotesize V\allowbreak{}}\par
\noindent Source: {\ttfamily\footnotesize ED-EN-FONT-G017;\allowbreak{} F7II-XIV-053.1;\allowbreak{} fr/sga7ii/\{diplomati\allowbreak{}c,current\}/expose\_XI\allowbreak{}V\_body.tex;\allowbreak{} folio\allowbreak{} 160\allowbreak{}}\par

\par\medskip\noindent\hypertarget{app-E7II-XIV-024.1}{}\textbf{E7II-XIV-024.1}\quad printed folio 161.\par
\noindent Printed reading: {\ttfamily\footnotesize :\allowbreak{}}\par
\noindent Corrected reading: \emph{(empty)}\par
\noindent Source: {\ttfamily\footnotesize ED-EN-FORMULA-G119;\allowbreak{} fr/sga7ii/\{diplomati\allowbreak{}c,current\}/expose\_XI\allowbreak{}V\_body.tex;\allowbreak{} folio\allowbreak{} 161;\allowbreak{} port/decisions/decis\allowbreak{}ions\_en7II\_r2\_002.js\allowbreak{}on\#C646;\allowbreak{} 7II:255\allowbreak{}}\par

\par\medskip\noindent\hypertarget{app-E7II-XIV-024.2}{}\textbf{E7II-XIV-024.2}\quad printed folio 161.\par
\noindent Printed reading: {\ttfamily\footnotesize 4\allowbreak{}}\par
\noindent Corrected reading: {\ttfamily\footnotesize 5\allowbreak{}}\par
\noindent Source: {\ttfamily\footnotesize ED-EN-NONMATH-G023;\allowbreak{} F7II-XIV-055.4;\allowbreak{} 7II:256;\allowbreak{} fr/sga7ii/\{diplomati\allowbreak{}c,current\}/expose\_XI\allowbreak{}V\_body.tex;\allowbreak{} folio\allowbreak{} 161\allowbreak{}}\par

\par\medskip\noindent\hypertarget{app-E7II-XIV-025.1}{}\textbf{E7II-XIV-025.1}\quad printed folio 163.\par
\noindent Printed reading: \emph{(empty)}\par
\noindent Corrected reading: {\ttfamily\footnotesize *u\textasciicircum{}\allowbreak{}}\par
\noindent Source: {\ttfamily\footnotesize fr/sga7ii/\{diplomati\allowbreak{}c,current\}/expose\_XI\allowbreak{}V\_body.tex;\allowbreak{} folio\allowbreak{} 163;\allowbreak{} formula\allowbreak{} 1847\allowbreak{}}\par

\par\endgroup\normalsize\fi
\clearpage\SGAtoc{Exposé XV}
\selectlanguage{english}
\clearpage
\SGApage{II}{165}

\begin{center}
{\large SGA 7 --- EXPOS\'E XV}\par
\vspace{2em}
{\Large\underline{The Picard--Lefschetz Formula}}\par
\vspace{1.5em}
{\large by P. Deligne}
\end{center}

\SGApage{II}{\(\approx\)166}
\medskip
\noindent\textbf{Source contents}
\begin{center}
\begin{tabular}{@{}llr@{}}
\S1. & Ordinary quadratic singularities (normal forms) & 2\\
1.1 & Reminders & 2\\
1.2 & Ordinary quadratic singularities & 4\\
1.3 & Moduli of quadratic singularities & 10\\
\S2. & Computation of vanishing cycles in a standard ordinary
  quadratic case & 13\\
2.1 & Preliminary: the cohomology of a cone & 13\\
2.2 & Vanishing cycles & 17\\
\S3. & The Picard--Lefschetz formula & 23\\
3.1 & Preliminary results & 23\\
3.2 & Even relative dimension & 24\\
3.3 & Odd relative dimension & 26\\
3.4 & Summary & 31
\end{tabular}
\end{center}

\section*{1. Ordinary quadratic singularities (normal forms)}

In this section we take up the results of VI 6 again, adopting a
``henselian'' rather than a ``formal'' point of view.  We give normal
forms for ordinary quadratic singularities and their deformations,
into which they may be put locally for the etale topology.

These results follow easily from VI 6, the implicit-function theorem
(in Tougeron's form), and R. Elkik's theorem on versal henselian
moduli of isolated singularities.  At the time of the oral seminar,
the latter was not yet available; a weaker theorem, covering only the
case of complete intersections (which was enough for us), had been
proved and will not be published.

\subsection*{1.1. Reminders}

\SGAvar{\subsubsection*{1.1.1.}
Let $f:X\longrightarrow S$ be a morphism, with $X$ locally
describable by $p$ equations in relative affine space
$\mathbb A_S^{n+p}$ of dimension $n+p$ over $S$.  The Jacobian ideal
$J^n(X/S)$ was defined in VI 5.  If $X$ is the closed subscheme of
$\mathbb A_S^{n+p}$ defined by the equations
$f_i=0$ $(1< i\leq p)$, it is the ideal sheaf generated by the
$p\times p$ minors of the Jacobian matrix
$(\partial f_i/\partial x_j)$.  The following theorem is well known.
A special case of its formal analogue is proved in Bourbaki,
\emph{Alg\`ebre commutative}, III, \S4, no.~6, Corollary 3 to
Theorem 2; the $C^\infty$ analogue is proved in Tougeron,
\emph{Id\'eaux de fonctions diff\'erentiables} (thesis, Rennes,
May 1967); the case considered here is treated in Artin,
\emph{Algebraic approximation of structures over complete local
rings}, Publ. Math. IHES \underline{36} (Zariski's volume),
Lemma 5.10.

}{\subsubsection*{1.1.1.}
Let $f:X\longrightarrow S$ be a morphism, with $X$ locally
describable by $p$ equations in relative affine space
$\mathbb A_S^{n+p}$ of dimension $n+p$ over $S$.  The Jacobian ideal
$J^n(X/S)$ was defined in VI 5.  If $X$ is the closed subscheme of
$\mathbb A_S^{n+p}$ defined by the equations
$f_i=0$ $(1\leq i\leq p)$, it is the ideal sheaf generated by the
$p\times p$ minors of the Jacobian matrix
$(\partial f_i/\partial x_j)$.  The following theorem is well known.
A special case of its formal analogue is proved in Bourbaki,
\emph{Alg\`ebre commutative}, III, \S4, no.~6, Corollary 3 to
Theorem 2; the $C^\infty$ analogue is proved in Tougeron,
\emph{Id\'eaux de fonctions diff\'erentiables} (thesis, Rennes,
May 1967); the case considered here is treated in Artin,
\emph{Algebraic approximation of structures over complete local
rings}, Publ. Math. IHES \underline{36} (Zariski's volume),
Lemma 5.10.

}{E7II-XV-001}\subsubsection*{Theorem 1.1.2 (Tougeron--Artin).}
\underline{Let $S$ be the spectrum of a henselian local ring $A$,
let $f:X\to S$ be a morphism as in 1.1.1, let $a$ and $I$ be two
proper ideals of $A$}, \underline{and let $\delta'$ be an ideal sheaf of $X$
contained in $J^n(X/S)$}.  \underline{Suppose that $I$ is of finite type.  Let
$s$ be a section of $f$ over $S/I$, that is, $s\in X(A/I)$}.  \underline{Denote
by $s^*\delta'^2$ the ideal of $A$ inverse to the ideal
$s^*\delta'^2$ of $A/I$.}

\underline{If $(s^*\delta'^2)\,a\supset I$, then there is a section
$s_1\in X(A)$ of $f$ congruent to $s$ modulo
$(s^*\delta')\,a$.}

\SGApage{II}{\(\approx\)167}
\subsubsection*{1.1.3.}
Let $S$ be the spectrum of a henselian local ring, with closed point
$s$.  Let $X'_0$ be a scheme of finite type over $k(s)$, smooth away
from a closed point $x_0$, and let $X_0$ be the henselization of
$X'_0$ at $x_0$.  We are interested in systems
\begin{center}
\begin{tikzcd}[column sep=large,row sep=large]
X_0 \arrow[r,hook] \arrow[d]
  \& X \arrow[d]
  \& {} \\
s \arrow[r,hook]
  \& T \arrow[r]
  \& S
\end{tikzcd}
\end{center}
of the following kind:
\begin{enumerate}[label=(\alph*),leftmargin=2.6em]
\item $T$ is henselian local and essentially of finite presentation
over $S$ (I 0.0.5); the morphism $T\to S$ is local and induces an
isomorphism of residue fields;

\item $X$ is henselian local, flat, and essentially of finite
presentation over $T$;

\item the special fiber of $X/T$ is equipped with a
$k(s)$-isomorphism to $X_0$.
\end{enumerate}
These are the \underline{henselian $S$-deformations} of $X_0$.
A \underline{morphism} from $(X,T)$ to $(X',T')$ is a commutative
diagram of $S$-schemes
\begin{center}
\begin{tikzcd}[column sep=large,row sep=large]
X \arrow[r] \arrow[d] \& X' \arrow[d] \\
T \arrow[r] \& T'
\end{tikzcd}
\end{center}
which induces the identity of $X_0$ after passage to the special
fibers.  A deformation $X/T$ is called \underline{versal}
(minimally versal) if:
\begin{enumerate}[label=(\alph*),leftmargin=2.6em]
\item for every morphism $X''/T''\longrightarrow X'/T'$ for which
$T''\to T'$ is a closed immersion, the map
\[
\operatorname{Hom}(X'/T',X/T)
\longrightarrow\operatorname{Hom}(X''/T'',X/T)
\]
is surjective.  In particular, on taking $X''/T''=X_0/s$, one has
$\operatorname{Hom}(X'/T',X/T)\ne\varnothing$;

\item for
$T'=\operatorname{Spec}(k(s)[\varepsilon]/(\varepsilon^2))$, the
``inverse-image'' map
\[
\operatorname{Hom}(T',T)\longrightarrow
\{\text{isomorphism classes of deformations of $X_0$ over $T'$}\}
\]
is bijective.
\end{enumerate}

\SGApage{II}{\(\approx\)168}
\subsubsection*{Theorem 1.1.4 (R. Elkik).}
\underline{There exists a versal henselian} $S$-\underline{deformation of $X_0$
(and it is unique up to a nonunique isomorphism).  If $S$ is
noetherian, a deformation $X/T$ is such a deformation if and only if
$X$ induces over the completion $\widehat T$ of $T$ a versal formal}
$S$-\underline{deformation of $X_0$ (a minimally formal versal deformation;
VI 1.9 and VI 4).}

For the proof of this delicate theorem, we refer to R. Elkik
(ENS seminar, 1971/72).


\subsection*{1.2. Ordinary quadratic singularities}

Let $y$ be a closed point of a scheme $Y$ of finite type over a field
$k$ of characteristic $p$, and let $n$ be the dimension of $Y$ at
$y$.

\SGAvar{\subsubsection*{Definition 1.2.1 (= VI 6.6).}
\begin{enumerate}[label=(\roman*),leftmargin=2.6em]
\item[(1)] \underline{If $k$ is algebraically closed, $y$ is called an
ordinary quadratic point of $Y$ if the completion
$\widehat{\mathcal O}_{Y,y}$ of the local ring of $Y$ at $y$ is
isomorphic to the quotient of $k[[x_1,\ldots,x_{n+1}]]$ by the ideal
generated by a single formal power series $f$ which can be written}
\begin{equation}\tag{4.2.1}
f(\underline X)=Q(\underline X)=+\text{terms of order }>2,
\end{equation}
\underline{where $Q(\underline X)$ is an ordinary quadratic form.}


\item \underline{If $ k$ is an algebraic closure of $k$, $y$ is
called an ordinary quadratic point of $Y$ if the points of
$\bar Y=Y\otimes_k\bar k$ lying over $y$ are ordinary quadratic
points of $\bar Y$.}
\end{enumerate}

}{\subsubsection*{Definition 1.2.1 (= VI 6.6).}
\begin{enumerate}[label=(\roman*),leftmargin=2.6em]
\item \underline{If $k$ is algebraically closed, $y$ is called an
ordinary quadratic point of $Y$ if the completion
$\widehat{\mathcal O}_{Y,y}$ of the local ring of $Y$ at $y$ is
isomorphic to the quotient of $k[[x_1,\ldots,x_{n+1}]]$ by the ideal
generated by a single formal power series $f$ which can be written}
\begin{equation}\tag{1.2.1.1}
f(\underline X)=Q(\underline X)+\text{terms of order }>2,
\end{equation}
\underline{where $Q(\underline X)$ is an ordinary quadratic form.}


\item \underline{If $\bar k$ is an algebraic closure of $k$, $y$ is
called an ordinary quadratic point of $Y$ if the points of
$\bar Y=Y\otimes_k\bar k$ lying over $y$ are ordinary quadratic
points of $\bar Y$.}
\end{enumerate}

}{E7II-XV-002}\subsubsection*{1.2.2.}
Replacing ``ordinary quadratic form'' by ``nondegenerate quadratic
form'' in this definition gives the notion of a
\underline{nondegenerate quadratic point} (VI 6.0).  A point $y$ is
a nondegenerate quadratic point of $Y$ if and only if it is an
ordinary quadratic point and either $p\ne2$ or $n$ is odd.

\SGApage{II}{169}
\subsubsection*{Example 1.2.3.}
Let $Q(\underline X)$ be a nondegenerate quadratic form in $n+1$
variables over $k$.  Let $Y_0$ be the quadratic cone in
$\mathbb A_k^{n+1}$ with equation $Q=0$.  Then the origin, denoted by
$y_0$, is a nondegenerate quadratic point of $Y_0$.

\subsubsection*{Example 1.2.4.}
Suppose that $p=2$ and $n=2m$.  Let $a\in k$, and let $Q$ be a
polynomial
\[
Q(\underline X)=(X_0^2-a)
 +\sum_{0<i\leq j\leq2m}a_{ij}x_ix_j.
\]
Suppose that the quadratic form in $2m$ variables
$\sum_{0<i\leq j}a_{ij}x_ix_j$ is nondegenerate, and denote by $y_0$
the closed point of $\mathbb A_k^{n+1}$ with coordinates
$(\sqrt a,0,\ldots,0)$.  Then $y_0$ is an ordinary quadratic point
of the closed subscheme $Y_0\subset\mathbb A_k^{n+1}$ with equation
$Q=0$.

\subsubsection*{Example 1.2.5.}
Let $k'$ be a finite separable extension of $k$, and let $Y_0$ be a
$k'$-scheme of the kind considered in 1.2.3 (respectively, 1.2.4).
If $Y_0$ is regarded as a $k$-scheme, then $y_0$ is a nondegenerate
(respectively, ordinary) quadratic point of $Y_0$.  The extension
$k'/k$ identifies with the largest separable subextension of
$k(y_0)/k$.

\subsubsection*{Theorem 1.2.6.}
\underline{Suppose that $y$ is an ordinary quadratic point of the
}$k$-\underline{scheme $Y$, and let $k'$ be the largest separable subextension
of $k(y)/k$.  There then exist a} $k'$-\underline{scheme
$Y_0\subset\mathbb A_{k'}^{n+1}$ of the kind in 1.2.3 or 1.2.4 and a
}$k$-\underline{isomorphism $\varphi$ between the henselization $Y_{(y)}$ of $Y$
at $y$ and the henselization $Y_{0(y_0)}$ of $Y_0$ at $y_0$.}

First observe that there is an etale neighborhood $Y'$ of $y$ in
$Y$
\begin{center}
\begin{tikzcd}[row sep=large,column sep=large]
{} \& Y'\arrow[d]\\
y\arrow[ur,hook]\arrow[r,hook]\&Y
\end{tikzcd}
\end{center}
which is a $k'$-scheme.  Indeed, $Y_{(y)}$ is canonically a
$k'$-scheme.  Replacing $(k,Y,y)$ by $(k',Y',y)$ reduces us to the
case $k=k'$, that is, to the case in which $k(y)/k$ is radicial.


\SGApage{II}{\(\approx\)170}
The $k(y)$-vector space $(\Omega^1_{Y/k})_y$ has dimension $n+1$.
It is enough to check this after extending scalars from $k(y)$ to its
algebraic closure and completing at $y$, where it is immediate.  A
neighborhood of $y$ in $Y$ is therefore a closed subscheme of a
smooth $k$-scheme $Z$ of dimension $n+1$.  Moreover, after replacing
$Y$ and $Z$ by suitable neighborhoods of $y$, the ideal in
$\mathcal O_Z$ defining $Y$ is principal (the same argument as
above), say with equation $f$.

\subsubsection*{Lemma 1.2.7.}
\underline{Suppose that $y$ is a nondegenerate quadratic point of
$Y$ and that $k(y)/k$ is radicial.  Then, in a neighborhood of $y$,
$J^n_{Y/k}$ (1.1.1) is generated by one element, supported at $y$,
and of rank one over $k$.  In particular, $k(y)=k$.}

It is enough to prove the first assertion after passing to an
algebraic closure of $k$ and completing at $y$.  It is then immediate
by direct calculation: with the notation of 1.2.1, Nakayama's lemma
shows that the $\partial f/\partial x_i$ (and hence, a fortiori, $f$
and the $\partial f/\partial x_i$) generate the maximal ideal.

\SGAvar{Let us prove 1.2.6 when $y$ is nondegenerate.  We may and shall
suppose that $k(y)=k$, that $Y$ is defined by an equation $f$ in a
smooth scheme $Z$, and that there is an etale morphism
\[
x:Z\longrightarrow\mathbb A_k^{n+1}
\]
(a ``system of local coordinates $x_i$'') taking $y$ to $0$.  There
is then a nondegenerate quadratic form $Q$ in $n+1$ variables such
that, at $y$,
\[
f=Q(x_i)+\text{terms of order }>2.
\]
Let $Y_0$ be defined by $Q$ as in 4.3.

}{Let us prove 1.2.6 when $y$ is nondegenerate.  We may and shall
suppose that $k(y)=k$, that $Y$ is defined by an equation $f$ in a
smooth scheme $Z$, and that there is an etale morphism
\[
x:Z\longrightarrow\mathbb A_k^{n+1}
\]
(a ``system of local coordinates $x_i$'') taking $y$ to $0$.  There
is then a nondegenerate quadratic form $Q$ in $n+1$ variables such
that, at $y$,
\[
f=Q(x_i)+\text{terms of order }>2.
\]
Let $Y_0$ be defined by $Q$ as in 1.2.3.

}{E7II-XV-003}\underline{On $Y$}, the system $(x_i)$ is an ``approximate solution''
of the equation $Q(x_i)=0$: if $\mathfrak m$ is the ideal defining
$y$, then
\[
Q(x_i)\in\mathfrak m^3\qquad\text{on }Y.
\]
Since $Q$ is nondegenerate, the ideal of $\mathbb A_k^{n+1}$
generated by the $\partial Q/\partial x_i$ is the maximal ideal at
$0$, and its inverse image under
$x=(x_i):Y\to\mathbb A_k^{n+1}$ is $\mathfrak m$.  By the
implicit-function theorem 1.1.2 (applied with $a=\mathfrak m$ and
$\delta'=J$), there are functions $x'_i$ on the henselization
$Y_{(y)}$, congruent to the $x_i$ modulo $\mathfrak m^2$, which
satisfy $Q(x'_i)=0$.  They give the required isomorphism between
$Y_{(y)}$ and $Y_{0(y_0)}$.

\SGApage{II}{171}
\subsubsection*{Lemma 1.2.8.}
\underline{Suppose that $y$ is a degenerate ordinary quadratic point
($p=2$ and $n$ even), and that $k(y)/k$ is radicial.  Then, in a
neighborhood of $y$, $T^1_{Y/k}$ is generated by one element,
supported at $y$, and of rank two over $k$.  In particular, either
$k(y)=k$, or $k(y)$ is isomorphic to $k(\sqrt a)$ for some
$a\in k$, $a\notin k^2$.}

\SGAvar{Proceed as in 1.2.7.  With the notation of 1.2.1, it is enough to
check that
\[
\dim_k\left(
k[[x_1,\ldots,x_{n+1}]]\big/
  (f,\partial f/\partial x_i)
\right))=2.
\]
In suitable coordinates $x_0,x_1,\ldots,x_{2m}$, the formal power
series $f$ can be written
\[
f=x_0^2+\sum_{i=1}^m x_ix_{i+m}+R,
\qquad\operatorname{ord}(R)\geq3.
\]
Let $\mathfrak n'$ be the ideal generated by $f$ and the
$\partial f/\partial x_i$, let $\mathfrak n$ be the ideal generated
by $x_0^2$ and the $x_i$ for $i\ne0$, and let $\mathfrak q$ be the
maximal ideal.  We prove that $\mathfrak n=\mathfrak n'$.  Plainly
$\mathfrak n'\subset\mathfrak n$, and by Nakayama it is enough to
check that
\[
\mathfrak n'+\mathfrak n\mathfrak q=\mathfrak n.
\]
Modulo $\mathfrak n\mathfrak q$, one has $f\equiv x_0^2$; modulo
$\mathfrak n\mathfrak q+(x_0^2)$, one has
\[
\frac{\partial f}{\partial x_i}\equiv x_{i+m},
\qquad
\frac{\partial f}{\partial x_{i+m}}\equiv x_i
\qquad(1\leq i\leq m),
\]
which proves the assertion.

}{Proceed as in 1.2.7.  With the notation of 1.2.1, it is enough to
check that
\[
\dim_k\left(
k[[x_1,\ldots,x_{n+1}]]\big/
  (f,\partial f/\partial x_i)
\right)=2.
\]
In suitable coordinates $x_0,x_1,\ldots,x_{2m}$, the formal power
series $f$ can be written
\[
f=x_0^2+\sum_{i=1}^m x_ix_{i+m}+R,
\qquad\operatorname{ord}(R)\geq3.
\]
Let $\mathfrak n'$ be the ideal generated by $f$ and the
$\partial f/\partial x_i$, let $\mathfrak n$ be the ideal generated
by $x_0^2$ and the $x_i$ for $i\ne0$, and let $\mathfrak q$ be the
maximal ideal.  We prove that $\mathfrak n=\mathfrak n'$.  Plainly
$\mathfrak n'\subset\mathfrak n$, and by Nakayama it is enough to
check that
\[
\mathfrak n'+\mathfrak n\mathfrak q=\mathfrak n.
\]
Modulo $\mathfrak n\mathfrak q$, one has $f\equiv x_0^2$; modulo
$\mathfrak n\mathfrak q+(x_0^2)$, one has
\[
\frac{\partial f}{\partial x_i}\equiv x_{i+m},
\qquad
\frac{\partial f}{\partial x_{i+m}}\equiv x_i
\qquad(1\leq i\leq m),
\]
which proves the assertion.

}{E7II-XV-004}Under the hypotheses of 1.2.8, if $k=k(y)$, we shall let $a$ denote
any square in $k$.

\SGApage{II}{172}
\subsubsection*{Lemma 1.2.9.}
\underline{Let $V$ be a vector bundle on a scheme $S$, and let $T$
be a closed subscheme of $\mathbb A_S(V)$ locally free of rank two
over $S$.  Then $T$ is contained in one and only one line in
$\mathbb A_S(V)$.}

The question is local for the fppf topology on $S$, so we may suppose
that $T/S$ has a section.  By translation, we may suppose that $T$
contains the zero section.  Let $\mathfrak n$ be the ideal of $T$
defining $0$; it is locally free of rank one over $S$.  Every
homogeneous linear form $v\in V$ on $\mathbb A_S(V)$ defines an
element of $\mathfrak n$, and hence an epimorphism
$V\twoheadrightarrow\mathfrak n$.  Its kernel defines the required
line.

\subsubsection*{1.2.10.}
In particular, if $A$ is a radicial rank-two closed subscheme of
affine space $\mathbb A_k^{n+1}$ over a field $k$ of characteristic
$2$, there is an $a\in k$ such that, in suitable affine coordinates,
$A$ has equations
\[
x_0^2-a=0,
\qquad x_i=0\quad(1\leq i\leq n).
\]
It is enough to take the $x_i$ which vanish on the line defined by
$A$.

\subsubsection*{1.2.11.}
Let us prove 1.2.6 for a degenerate $y$ ($p=2$ and $n=2m$).  We may
and shall suppose that $k(y)=k(\sqrt a)$, that $Y$ is defined by an
equation $f$ in a smooth scheme $Z$, and that there is a system of
local coordinates at $y$
\[
x=(x_i)_{0\leq i\leq2m}:
Z\longrightarrow\mathbb A_k^{2m+1}
\]
taking $y$ to the closed point $(\sqrt a,0,\ldots,0)$, such that, in
a neighborhood of $y$,
\[
J^n_{Y/k}\simeq
\mathcal O_Y/(x_0^2-a,x_1,\ldots,x_{2m}).
\]
After replacing $f$ by a suitable multiple, one has
\[
f=(x_0^2-a)+
 \sum_{1\leq i\leq j\leq2m}a_{ij}x_ix_j+R,
\]
where
\begin{enumerate}[label=(\alph*),leftmargin=2.6em]
\item the quadratic form in $2m$ variables
$\sum a_{ij}x_ix_j$ is nondegenerate;
\item after extending scalars to $\bar k$, $R$ has order at least
$3$ in $x_0-\sqrt a$ and the $x_i$.
\end{enumerate}
Let $Y_0$ be defined, as in 1.2.4, by
\[
Q=x_0^2-a+\sum a_{ij}x_ix_j.
\]

\SGApage{II}{\(\approx\)173}
\SGAvar{\underline{On $Y$}, the system $(x_i)$ may again be regarded as an
approximate solution of $Q(X_i)=0$.  This time, the ideal $\delta$,
the inverse image under $x$ of the ideal generated by the
$\partial Q/\partial X_i$, satisfies
$\mathcal O_Y/\delta\simeq J^n_{Y/k}$; at $y$ it is
\[
\delta=(x_0^2-a,(x_i)_{1\leq i}).
\]
Let $\mathfrak q$ be the maximal ideal defining $y$.  To conclude as
in the nondegenerate case by applying the implicit-function theorem,
it remains to verify that
\begin{equation}\tag{1.2.11.1}
Q(x_i)\equiv0\pmod{\delta^2\mathfrak q}
\qquad\underline{\text{on }Y}.
\end{equation}
Let $\bar{\mathfrak q}$ be the ideal defining $y$ after scalar
extension to an algebraic closure of $k$.  One checks that
\[
\delta^2\bar{\mathfrak q}\cap\mathcal O_{Y,y}
=\delta^2\mathfrak q.
\]
It is therefore enough to verify (4.11.1) after extending scalars
to $\bar k$; we may suppose that $a=0$.
At $y$ one then has
\[
\delta^2\mathfrak q=
(x_0^5,(x_0^3x_i)_{1\leq i},
 (x_0x_ix_j)_{1\leq i,j},
 (x_ix_jx_k)_{1\leq i,j,k}).
\]
It follows that $\delta^2\mathfrak q\supset\mathfrak q^5$ and that
\[
x_0^2f\equiv x_0^4\pmod{\delta^2\mathfrak q},
\qquad
\delta^2\mathfrak q+(x_0^2f)\supset\mathfrak q^4,
\]
\[
x_0f\equiv x_0^3
 \pmod{\delta^2\mathfrak q+\mathfrak q^4},
\qquad
\delta^2\mathfrak q+(x_0f)\supset\mathfrak q^3.
\]
Finally,
\[
Q(\chi_i)\equiv0\pmod{\delta^2\mathfrak q+(f)}
\qquad\underline{\text{on }Z},
\]
which is (1.2.11.1).

}{\underline{On $Y$}, the system $(x_i)$ may again be regarded as an
approximate solution of $Q(X_i)=0$.  This time, the ideal $\delta$,
the inverse image under $x$ of the ideal generated by the
$\partial Q/\partial X_i$, satisfies
$\mathcal O_Y/\delta\simeq J^n_{Y/k}$; at $y$ it is
\[
\delta=(x_0^2-a,(x_i)_{1\leq i}).
\]
Let $\mathfrak q$ be the maximal ideal defining $y$.  To conclude as
in the nondegenerate case by applying the implicit-function theorem,
it remains to verify that
\begin{equation}\tag{1.2.11.1}
Q(x_i)\equiv0\pmod{\delta^2\mathfrak q}
\qquad\underline{\text{on }Y}.
\end{equation}
Let $\bar{\mathfrak q}$ be the ideal defining $y$ after scalar
extension to an algebraic closure of $k$.  One checks that
\[
\delta^2\bar{\mathfrak q}\cap\mathcal O_{Y,y}
=\delta^2\mathfrak q.
\]
It is therefore enough to verify (1.2.11.1) after extending scalars
to $\bar k$; we may suppose that $a=0$.
At $y$ one then has
\[
\delta^2\mathfrak q=
(x_0^5,(x_0^3x_i)_{1\leq i},
 (x_0x_ix_j)_{1\leq i,j},
 (x_ix_jx_k)_{1\leq i,j,k}).
\]
It follows that $\delta^2\mathfrak q\supset\mathfrak q^5$ and that
\[
x_0^2f\equiv x_0^4\pmod{\delta^2\mathfrak q},
\qquad
\delta^2\mathfrak q+(x_0^2f)\supset\mathfrak q^4,
\]
\[
x_0f\equiv x_0^3
 \pmod{\delta^2\mathfrak q+\mathfrak q^4},
\qquad
\delta^2\mathfrak q+(x_0f)\supset\mathfrak q^3.
\]
Finally,
\[
Q(\chi_i)\equiv0\pmod{\delta^2\mathfrak q+(f)}
\qquad\underline{\text{on }Z},
\]
which is (1.2.11.1).

}{E7II-XV-005}\SGApage{II}{174}
\subsubsection*{Remark 1.2.12.}
Let $Q$ be a not necessarily homogeneous, nonzero quadratic form in
$n+1$ variables, and suppose that the closed subscheme
$Y\subset\mathbb A_k^{n+1}$ with equation $Q$ has an ordinary
quadratic singularity at a point $y$.  A nonhomogeneous linear change
of variables then puts $Y$ into the form 1.2.3 or 1.2.4.  In the
nondegenerate case $k(y)=k$, and it is enough to put $y$ at the
origin.  In the degenerate case it is enough that the support of
$T^1_{Y/k}$ lie on the line $x_i=0$ $(1\leq i\leq2m)$
(cf. 1.2.9).


\subsection*{1.3. Moduli of quadratic singularities}

The following proposition follows from R. Elkik's theorem 1.1.4 and
from explicit calculations left to the reader (and essentially
contained in VI 6).

\subsubsection*{Proposition 1.3.1.}
\underline{Let $S=\operatorname{Spec}(A)$ be a henselian local scheme
with closed point $s$, and let
$Y_0\subset\mathbb A_{k(s)}^{n+1}$ be a scheme of the kind described
in 1.2.3 or 1.2.4.}

\begin{enumerate}[label=(\roman*),leftmargin=2.6em]
\item \underline{Suppose that $Y_0$ is of the kind described in
1.2.3, defined by the equation}
\[
Q(\underline X)=\sum_{i\leq j}\bar a_{ij}x_i x_j=0
\qquad(\bar a_{ij}\in k(s)),
\]
\underline{and lift the $\bar a_{ij}$ to elements $a_{ij}$ of $A$.
Then the versal henselian $S$-scheme of deformations of the
singularity of $Y_0$ at $y_0$ is}
\[
T=\operatorname{Spec}(A\{b\}),
\]
\underline{the versal deformation being the henselization at $y_0$
of the $T$-scheme $Y\subset\mathbb A_T^{n+1}$ with equation}
\[
\sum a_{ij}x_i x_j-b=0.
\]

\item \underline{Suppose that $Y_0$ is of the kind described in
1.2.4, defined by the equation}
\[
Q(\underline x)=(x_0^2-\bar a)+\sum\bar a_{ij}x_i x_j=0
\qquad(\bar a,\bar a_{ij}\in k(s)),
\]
\underline{and lift the $\bar a_{ij}$ and $\bar a$ to elements
$a_{ij}$ and $a$ of $A$.  Then the versal henselian $S$-scheme of
deformations of the singularity of $Y_0$ at $y_0$ is}
\[
T=\operatorname{Spec}(A\{b,c\}),
\]
\underline{the versal deformation being the henselization at $y_0$
of the $T$-scheme $Y\subset\mathbb A_T^{n+1}$ with equation}
\[
(x_0^2-a)+\sum a_{ij}x_i x_j+b x_0+c=0.
\]
\end{enumerate}

\SGApage{II}{\(\approx\)175}
The following corollary follows at once from 1.2.6 and 1.3.1.

\subsubsection*{Corollary 1.3.2.}
\underline{Let $S=\operatorname{Spec}(A)$ be a henselian local
scheme with closed point $s$, and let $f:X\to S$ be a flat morphism
of finite presentation.  Suppose that, at the closed point $x$ of
$X_s$, the scheme $X_s$ has an ordinary quadratic singularity, that
$k(x)$ is radicial over $k(s)$, and that $X_s$ has dimension $n$.}

\SGAvar{\begin{enumerate}[label=(\roman*),leftmargin=2.6em]
\item If $x$ is nondegenerate, \underline{there exist a
nondegenerate quadratic form $Q$ in $n+1$ variables with
coefficients in $A$, and an element $b$ in the maximal ideal of $A$,
such that the henselized} $S$-\underline{scheme of $X$ at $x$ is isomorphic to
the henselization at the origin of the closed subscheme of
$\mathbb A_S^{n+1}$ with equation $Q-b=0$, as in 1.3.1(i).}

\item If $x$ is degenerate
($n=2m$, $\operatorname{char}(k(s))=2$), \underline{there exists a
nonhomogeneous quadratic form in $n+1$ variables with coefficients
in $A$}
\begin{equation}\tag{1.3.2.1}
Q(\underline X)=x_0+b x_0+c+
 \sum_{1\leq i\leq j\leq 2m}a_{ij}x_i x_j,
\end{equation}
\underline{such that $b$ lies in the maximal ideal, the quadratic
form $\sum a_{ij}x_i x_j$ in $2m$ variables is nondegenerate, and
the henselized} $S$-\underline{scheme of $X$ at $x$ is isomorphic to the
henselization at $(\sqrt c,0,\ldots,0)$ of the closed subscheme of
$\mathbb A_S^{n+1}$ with equation $Q=0$.}

\end{enumerate}

}{\begin{enumerate}[label=(\roman*),leftmargin=2.6em]
\item If $x$ is nondegenerate, \underline{there exist a
nondegenerate quadratic form $Q$ in $n+1$ variables with
coefficients in $A$, and an element $b$ in the maximal ideal of $A$,
such that the henselized} $S$-\underline{scheme of $X$ at $x$ is isomorphic to
the henselization at the origin of the closed subscheme of
$\mathbb A_S^{n+1}$ with equation $Q-b=0$, as in 1.3.1(i).}

\item If $x$ is degenerate
($n=2m$, $\operatorname{char}(k(s))=2$), \underline{there exists a
nonhomogeneous quadratic form in $n+1$ variables with coefficients
in $A$}
\begin{equation}\tag{1.3.2.1}
Q(\underline X)=x_0^2+b x_0+c+
 \sum_{1\leq i\leq j\leq 2m}a_{ij}x_i x_j,
\end{equation}
\underline{such that $b$ lies in the maximal ideal, the quadratic
form $\sum a_{ij}x_i x_j$ in $2m$ variables is nondegenerate, and
the henselized} $S$-\underline{scheme of $X$ at $x$ is isomorphic to the
henselization at $(\sqrt c,0,\ldots,0)$ of the closed subscheme of
$\mathbb A_S^{n+1}$ with equation $Q=0$.}

\end{enumerate}

}{E7II-XV-006}\SGApage{II}{176}
\subsubsection*{Remark 1.3.3.}
In case (i), suppose that the henselization of $X_s$ at $x$ is
isomorphic, by an isomorphism $\varphi$, to the henselization at $O$
of the closed subscheme
$X_0\subset\mathbb A_k^{n+1}$ with equation
\[
\sum\bar a_{ij}x_i x_j=0,
\]
and suppose that the $a_{ij}$ lift the $\bar a_{ij}$.  One may then
take
\[
Q=\sum a_{ij}x_i x_j.
\]

In case (ii), suppose that the henselization of $X_s$ at $x$ is
isomorphic, by an isomorphism $\varphi$, to the henselization at
$(\sqrt a,0,\ldots,0)$ of the closed subscheme
$X_0\subset\mathbb A_k^{n+1}$ with equation
\[
(x_0^2+\bar a)+\sum\bar a_{ij}x_i x_j=0,
\]
and suppose that the $a_{ij}$ lift the $\bar a_{ij}$.  One may then
take
\[
Q=x_0^2+b x_0+c+\sum a_{ij}x_i x_j,
\qquad c\equiv\bar a\pmod{\mathfrak m_A}.
\]
In both cases, the isomorphism promised in 1.3.2 may be chosen to
extend $\varphi$.

\subsubsection*{Corollary 1.3.4.}
\underline{Let $f:X\to S$ be as in 1.3.2.  There is a neighborhood
$U$ of $x$ in $X$ such that the nonsmooth points of $f$ contained in
$U$ are ordinary quadratic points of their respective fibers.}


\section*{2. Computation of vanishing cycles in a standard ordinary
quadratic case}

\SGApage{II}{177}
\subsection*{2.1. Preliminary: the cohomology of a cone}

\subsubsection*{2.1.1.}
Let $Y\subset\mathbb P^r$ be a projective variety over an
algebraically closed field $k$.  We use the following notation.

The affine cone over $Y$ is denoted by
$X\subset\mathbb A^{r+1}$, and its vertex by $O$;

$X_{(O)}$ is the henselization of $X$ at $O$;

$X^*=X-\{O\}$ and $X^*_{(O)}=X_{(O)}-\{O\}$;

$X_1\subset\mathbb P^{r+1}$ is the projective cone over $Y$, so
$X=X_1-Y$;

$\widetilde X$, $\widetilde X_{(O)}$, and $\widetilde X_1$ are the
blowups of $X$, $X_{(O)}$, and $X_1$ at $O$;

$Y_0$ is the inverse image of $O$ in $\widetilde X$,
$\widetilde X_{(O)}$, or $\widetilde X_1$;

$h:\widetilde X_1\to Y$ is the natural projection;

$y_0,y_\infty:Y\to\widetilde X_1$ are the sections of $h$ whose
images are $Y_0$ and $Y$;

$F$ is an abelian torsion group of order prime to the characteristic
of $k$.  All cohomology groups considered below have coefficients in
$F$, which will not be displayed.

Our purpose in this subsection is to relate the cohomology of
$X^*_{(O)}$ to that of $Y$.

\subsubsection*{Proposition 2.1.2.}
\underline{There are isomorphisms}
\[
\text{(i)}\qquad
H^i(X)\xrightarrow{\ \sim\ }H^i(\{O\})=
\begin{cases}
F,& i=0,\\
0,& i>0;
\end{cases}
\]
\[
\text{(ii)}\qquad
H^i_{\{O\}}(X)\xrightarrow{\ \sim\ }H^i_c(X).
\]
This proposition is proved by a homotopy argument, which we now
formalize.

\SGAvar{Let $k$ be an algebraically closed field, let $U$ and $V$ be two
$k$-schemes, let $\Lambda$ be a torsion ring of order prime to the
characteristic of $k$, and let
\[
K\in\operatorname{Ob}D^+(U,\Lambda),
\qquad L\in\operatorname{Ob}D^+(V,X).
\]
By a ``morphism'' $f$ from $(U,K)$ to $(V,L)$ we mean a pair
consisting of a morphism $f:U\to V$ and a morphism
$\varphi:f^*L\to K$.  Such a morphism induces
\[
f^*:H^*(V,L)\longrightarrow H^*(U,K).
\]

}{Let $k$ be an algebraically closed field, let $U$ and $V$ be two
$k$-schemes, let $\Lambda$ be a torsion ring of order prime to the
characteristic of $k$, and let
\[
K\in\operatorname{Ob}D^+(U,\Lambda),
\qquad L\in\operatorname{Ob}D^+(V,\Lambda).
\]
By a ``morphism'' $f$ from $(U,K)$ to $(V,L)$ we mean a pair
consisting of a morphism $f:U\to V$ and a morphism
$\varphi:f^*L\to K$.  Such a morphism induces
\[
f^*:H^*(V,L)\longrightarrow H^*(U,K).
\]

}{E7II-XV-007}\SGApage{II}{\(\approx\)178}
We say that two ``morphisms'' $f_0$ and $f_1$ are homotopic if there
exist a connected $k$-scheme $T$ of finite type, two points $0$ and
$1$ of $T$, and a ``morphism'' $H$ from
$(U\times T,\operatorname{pr}_1^*K)$ to $(V,L)$ whose fibers at $0$
and $1$ induce $f_0$ and $f_1$.

\subsubsection*{Lemma 2.1.3.}
\underline{If $f_0$ is homotopic to $f_1$, then
$f_0^*=f_1^*$.}

There is a sequence of points
$0=x_0,x_1,\ldots,x_n=1$ of $T$, smooth connected curves
$\Gamma_i$, and morphisms $a_i:\Gamma_i\to T$ such that $x_i$ and
$x_{i+1}$ belong to the image of $a_i$.  To construct the
$\Gamma_i$, normalize a connected one-dimensional closed subscheme
of $T$ containing $0$ and $1$.  This reduces us to the case in which
$T$ is a smooth curve.  Consider the commutative square
\begin{center}
\begin{tikzcd}[row sep=large,column sep=large]
U\times T \arrow[r,"\operatorname{pr}_1"]
  \arrow[d,"\operatorname{pr}_2"']
  \& U \arrow[d,"f"] \\
T \arrow[r,"t"'] \& k .
\end{tikzcd}
\end{center}
By smooth base change, applied to $t$, one has
\[
t^*Rf_*K\xrightarrow{\ \sim\ }
R\operatorname{pr}_{2*}(\operatorname{pr}_1^*K).
\]
The map $f_i^*$ $(i=0,1)$ therefore factors through
\begin{center}
\begin{tikzcd}[row sep=large,column sep=large]
H^n(V,L) \arrow[r,"H^*"]
  \& H^n(U\times T,\operatorname{pr}_1^*K) \arrow[r]
  \& H^0\bigl(T,R^n\operatorname{pr}_{2*}
      (\operatorname{pr}_1^*K)\bigr)
      \arrow[dd,"\text{fiber at }i"] \\
{} \& H^0\bigl(T,t^*H^n(U,K)\bigr) \arrow[ur,"\wr"] \& {} \\
{} \& {} \& H^n(U,K) \arrow[ul,"\wr"']
\end{tikzcd}
\end{center}
and is consequently independent of $i$.


\subsubsection*{2.1.3.}
Let us prove 2.1.2.  Scalar multiplication gives a homotopy between
the identity map of $X$ and the constant map from $X$ to $X$ with
value $O$; this proves (i).

\SGApage{II}{\(\approx\)179}
Using scalar multiplications whose factors tend to infinity, one
sees that $Y$ is a deformation retract of $X_1^*$.  Assertion (ii)
then follows from the five lemma and the diagram
\begin{center}
\begin{tikzcd}[row sep=large,column sep=large]
{} \arrow[r,"\partial"]
  \& H^i_{\{O\}}(X) \arrow[r] \arrow[d]
  \& H^i(X_1) \arrow[r] \arrow[d,equal]
  \& H^i(X_1^*) \arrow[r,"\partial"] \arrow[d,"\wr"']
  \& {} \\
{} \arrow[r,"\partial"]
  \& H^i_c(X) \arrow[r]
  \& H^i(X_1) \arrow[r]
  \& H^i(Y) \arrow[r,"\partial"]
  \& {} .
\end{tikzcd}
\end{center}

\subsubsection*{Corollary 2.1.4.}
\[
H^i(X^*)\xrightarrow{\ \sim\ }H^i(X^*_{(O)}).
\]
Indeed, apply the five lemma to
\begin{center}
\begin{tikzcd}[row sep=large,column sep=large]
{} \arrow[r]
  \& H^i_{\{O\}}(X) \arrow[r] \arrow[d,"\wr"]
  \& H^i(X) \arrow[r]
      \arrow[d,"\wr"',"{(2.1.2(i))}"]
  \& H^i(X^*) \arrow[r] \arrow[d]
  \& {} \\
{} \arrow[r]
  \& H^i_{\{O\}}(X_{(O)}) \arrow[r]
  \& H^i(X_{(O)}) \arrow[r]
  \& H^i(X^*_{(O)}) \arrow[r]
  \& {} .
\end{tikzcd}
\end{center}

\subsubsection*{2.1.5.}
The Leray spectral sequences for the restrictions of $h$ to
$\widetilde X_1-Y_0$ and $\widetilde X_1-Y$ show that the restriction
maps
\[
H^i(\widetilde X_1-Y_0)\xrightarrow{\ \sim\ }H^i(Y),
\qquad
H^i(\widetilde X)\xrightarrow{\ \sim\ }H^i(Y_0)=H^i(Y)
\]
are isomorphisms.  Under these isomorphisms, the long exact
cohomology sequences of the pairs
$(\widetilde X_1-Y_0,Y)$ and $(\widetilde X_1-Y,Y_0)$ form the rows
of the diagram
\begin{equation}\tag{2.1.5.1}
\begin{tikzcd}[column sep=large]
H^{i-1}(X^*) \arrow[r] \arrow[d,equal]
  \& H^{i-2}(Y)(-1) \arrow[r,"\eta"'] \arrow[d,equal]
  \& H^i(Y) \arrow[r] \arrow[d,dash,"-1"]
  \& H^i(X^*) \arrow[r] \arrow[d,equal]
  \& {} \\
H^{i-1}(X^*) \arrow[r]
  \& H^{i-2}(Y)(-1) \arrow[r,"-\eta"']
  \& H^i(Y) \arrow[r]
  \& H^i(X^*) \arrow[r]
  \& {} .
\end{tikzcd}
\end{equation}

\SGApage{II}{180}
\subsubsection*{Lemma 2.1.6.}
\underline{The arrows denoted by $\eta$ and $-\eta$ in (2.1.5.1)
are cup products with the cohomology class of a hyperplane section.
Diagram (2.1.5.1) is commutative.}

The first assertion follows because $\eta$ (respectively $-\eta$)
is the restriction to $Y$ (respectively to $Y_0\simeq Y$) of the
cohomology class of the invertible sheaf $\mathcal O(Y)$
(respectively $\mathcal O(Y_0)$) on $\widetilde X_1-Y_0$
(respectively $\widetilde X_1-Y$).

Use the Leray spectral sequence for $h$.  The two rows of (2.1.5.1)
are obtained by applying the functor $H(Y,-)$ to the distinguished
triangles
\begin{center}
\begin{tikzcd}[column sep=large]
{} \& y_{\infty*}Ry_\infty^!F \arrow[r]
  \& R(h|\widetilde X_1-Y_0)_*F
      \arrow[r,"\textcircled{\scriptsize 1}"]
  \& R(h|X^*)_*F \arrow[r,"\textcircled{\scriptsize 2}"]
  \& {} \\
\text{and}
  \& y_{0*}Ry_0^!F \arrow[r]
  \& R(h|\widetilde X_1-Y)_*F \arrow[r]
  \& R(h|X^*)_*F \arrow[r]
  \& {} .
\end{tikzcd}
\end{center}
The nonzero cohomology sheaves of the complexes at the vertices of
these triangles occur in degrees $(2)$, $(0)$, and $(0$ and $1)$.
The arrows \textcircled{\scriptsize 1} and
\textcircled{\scriptsize 2} are therefore determined by their effect
on cohomology sheaves.  This reduces us to the case in which $Y$ is
a point, which is left to the reader.


\SGApage{II}{\(\approx\)181}
\SGAvar{\subsubsection*{2.1.7.}
The long exact sequence of the pair
$(\widetilde X_{(O)},Y_0)$ is a local analogue of the second row of
(2.1.5.1).  One has
\[
H^i(\widetilde X_{(O)})\xrightarrow{\ \sim\ }H^i(Y_0)
\]
by proper base change for the morphism
$\widetilde X_{(O)}\to X_{(O)}$.  The long exact sequence in question
may therefore be written
\begin{equation}\tag{2.1 7.1}
\begin{tikzcd}[column sep=large]
H^{i-1}(X^*_{(O)}) \arrow[r]
  \& H^{i-2}(Y_0)(-1) \arrow[r,"-\eta"]
  \& H^i(Y_0) \arrow[r]
  \& H^i(X^*_{(O)}) \& {} \\
H^{i-1}(X^*) \arrow[r] \arrow[u,"\wr"]
  \& H^{i-2}(Y)(-1) \arrow[r,"-\eta"] \arrow[u,equal]
  \& H^i(Y) \arrow[r] \arrow[u,equal]
  \& H^i(X^*) \arrow[u,dash,"\wr"]
  \& {} ,
\end{tikzcd}
\end{equation}
the second row of this commutative diagram being the second row of
(2.1.5.1).


}{\subsubsection*{2.1.7.}
The long exact sequence of the pair
$(\widetilde X_{(O)},Y_0)$ is a local analogue of the second row of
(2.1.5.1).  One has
\[
H^i(\widetilde X_{(O)})\xrightarrow{\ \sim\ }H^i(Y_0)
\]
by proper base change for the morphism
$\widetilde X_{(O)}\to X_{(O)}$.  The long exact sequence in question
may therefore be written
\begin{equation}\tag{2.1.7.1}
\begin{tikzcd}[column sep=large]
H^{i-1}(X^*_{(O)}) \arrow[r]
  \& H^{i-2}(Y_0)(-1) \arrow[r,"-\eta"]
  \& H^i(Y_0) \arrow[r]
  \& H^i(X^*_{(O)}) \& {} \\
H^{i-1}(X^*) \arrow[r] \arrow[u,"\wr"]
  \& H^{i-2}(Y)(-1) \arrow[r,"-\eta"] \arrow[u,equal]
  \& H^i(Y) \arrow[r] \arrow[u,equal]
  \& H^i(X^*) \arrow[u,dash,"\wr"]
  \& {} ,
\end{tikzcd}
\end{equation}
the second row of this commutative diagram being the second row of
(2.1.5.1).


}{E7II-XV-008}The following technical lemma will be used in 2.2.7.

\SGAvar{\subsubsection*{Lemma 2.7.8.}
\begin{center}
\begin{tikzcd}[column sep=large]
H^{n-1}(Y_0) \arrow[d,equal]
  \& H^{n-1}(\widetilde X) \arrow[l,"\sim"'] \arrow[r]
  \& H^{n-1}(X^*) \arrow[r]
  \& H^n_{\{x\}}(X) \arrow[r]
  \& H^n_c(X_s) \arrow[d,equal] \\
H^{n-1}(Y) \arrow[rrrr,"\partial"]
  \& {} \& {} \& {} \& H^n_c(X_s)
\end{tikzcd}
\end{center}
\underline{is anticommutative.}


}{\subsubsection*{Lemma 2.1.8.}
\begin{center}
\begin{tikzcd}[column sep=large]
H^{n-1}(Y_0) \arrow[d,equal]
  \& H^{n-1}(\widetilde X) \arrow[l,"\sim"'] \arrow[r]
  \& H^{n-1}(X^*) \arrow[r]
  \& H^n_{\{x\}}(X) \arrow[r]
  \& H^n_c(X_s) \arrow[d,equal] \\
H^{n-1}(Y) \arrow[rrrr,"\partial"]
  \& {} \& {} \& {} \& H^n_c(X_s)
\end{tikzcd}
\end{center}
\underline{is anticommutative.}


}{E7II-XV-009}By 2.1.7, this is equivalent to proving the commutativity of
\begin{center}
\begin{tikzcd}[column sep=large,row sep=large]
H^{n-1}(X_1-\{O\}) \arrow[r,"\partial"] \arrow[d,"\wr"]
  \& H^n_{\{O\}}(X_1)=H^n_{\{O\}}(X) \arrow[d] \\
H^{n-1}(Y) \arrow[r,"\partial"]
  \& H^n_c(X).
\end{tikzcd}
\end{center}
This is general nonsense.


\subsection*{2.2. Vanishing cycles}

\subsubsection*{2.2.1.}
Let $S$ be a henselian trait, and let
$s,\eta,\bar s,\bar\eta$ have the meaning of XIII, 0.2.5.  Put
$\Lambda=\mathbb Z/k$, where $k$ is prime to the residual
characteristic of $S$.  Let $\mathbb A^{n+1}_S$ be standard affine
$(n+1)$-space over $S$, and let $X$ be the closed subscheme of
$\mathbb A^{n+1}_S$ defined by a quadratic equation $Q=0$ which is
nonzero modulo a uniformizer:
\[
Q(\underline{\mathbf{x}})=\sum_{i\leq j}a_{ij}x_ix_j+\sum_i b_ix_i+c.
\]
Regard $\mathbb A^{n+1}_S$ as the complement in
$\mathbb P^{n+1}_S$ of the hyperplane at infinity $H$.  Then $X$ is
obtained from the quadric $X_1\subset\mathbb P^{n+1}_S$ with
homogeneous equation
\begin{equation}\tag{2.2.1.1}
\sum_{i\leq j}a_{ij}x_ix_j+\sum_i b_ix_iz+cz^2=0
\end{equation}
by removing $Y=X_1\cap H$.

\SGApage{II}{\(\approx\)182}
We make the following assumptions:

\SGAvar{\begin{enumerate}
\item[(a)] $Y$ is a smooth quadric over $S$, that is, the quadratic
form $\sum_{i\leq j}a_{ij}x_ix_j$ is ordinary;

\item[(b)] the quadric $X'_s$ is singular, that is,
$X_{\bar s}$ is a quadratic cone.
\end{enumerate}


}{\begin{enumerate}
\item[(a)] $Y$ is a smooth quadric over $S$, that is, the quadratic
form $\sum_{i\leq j}a_{ij}x_ix_j$ is ordinary;

\item[(b)] the quadric $X_s$ is singular, that is,
$X_{\bar s}$ is a quadratic cone.
\end{enumerate}


}{E7II-XV-010}Let $x_0$ denote the singular point of $X_s$, and let
$f:X\to S$ be the projection.

\subsubsection*{Remark 2.2.2.}
If $x_0$ is a rational point, a translation puts it at the origin of
the coordinates.  If $x_0=0$, the $b_i$ and $c$ are divisible by a
uniformizer.

The subscheme $A$ of $X_s$ defined by the equations
$\partial Q/\partial x_i$ is concentrated at $x_0$.  Except in the
exceptional case where $k(s)$ has characteristic $2$ and $n+1$ is
even, this scheme has degree one over $k(s)$, and hence $x_0$ is
rational.  In the exceptional case it is ramified over $k(s)$ and
has rank $2$, so that $k(x_0)$ is either $k(s)$ or a purely
inseparable quadratic extension of $k(s)$.  The assertions about
$A$ are checked most easily after passing to the algebraic closure
$k(\bar s)$ of $k(s)$.

\subsubsection*{Proposition 2.2.3.}
\underline{Under the assumptions of 2.2.1, the morphisms}
\[
H^i(X_{\bar\eta},\Lambda)
 \xrightarrow{\ \sim\ }
\mathbb H^i\bigl(X_{\bar s},R\psi_{\bar\eta}(\Lambda)\bigr)
 \xrightarrow{\ \sim\ }
R^i\psi_{\bar\eta}(\Lambda)_{x_0}
\]
\[
H_c^i(X_{\bar\eta},\Lambda)
 \xleftarrow{\ \sim\ }
\mathbb H_c^i\bigl(X_{\bar s},R\psi_{\bar\eta}(\Lambda)\bigr)
 \xleftarrow{\ \sim\ }
\mathbb H^i_{\{x_0\}}\bigl(X_{\bar s},R\psi_{\bar\eta}(\Lambda)\bigr)
\]
\underline{are isomorphisms.}

\SGApage{II}{\(\approx\)183}
That the left-hand arrows are isomorphisms follows from 2.1.8.6 and
2.1.10.5; the assertion for the right-hand arrows is a corollary of
2.1.2.  Indeed, there is a distinguished triangle
\[
(\Lambda\text{ on }X_{\bar s})[0]
 \longrightarrow R\psi_{\bar\eta}(\Lambda)
 \longrightarrow R\Phi_{\bar\eta}(\Lambda),
\]
and 2.1.2 applies to $(\Lambda\text{ on }X_{\bar s})$, whereas
$R\Phi_{\bar\eta}(\Lambda)$ is supported at $x_0$.

\subsubsection*{Corollary 2.2.4.}
\underline{If the geometric generic fiber $X_{\bar\eta}$ is
singular, that is, if it is still a quadratic cone, then all the
sheaves $R^i\Phi(\Lambda)$ vanish.}

\SGAvar{By 2.1.2,
\[
\Lambda=H^o(X_{\bar s},\Lambda)
 \xrightarrow{\ \sim\ }H^o(X_{\bar\eta},\Lambda),
\qquad
0=H^i(X_{\bar s},\Lambda)=H^i(X_{\bar\eta},\Lambda)
\quad(i>0).
\]
We then apply the long exact sequence XIII, 2.1.8.9 (valid by
XIII, 2.1.9).



}{By 2.1.2,
\[
\Lambda=H^0(X_{\bar s},\Lambda)
 \xrightarrow{\ \sim\ }H^0(X_{\bar\eta},\Lambda),
\qquad
0=H^i(X_{\bar s},\Lambda)=H^i(X_{\bar\eta},\Lambda)
\quad(i>0).
\]
We then apply the long exact sequence XIII, 2.1.8.9 (valid by
XIII, 2.1.9).



}{E7II-XV-011}\subsubsection*{2.2.5.}
Suppose now that
\[
(*)\qquad\text{the generic fiber $X_\eta$ is smooth.}
\]
For simplicity we shall also assume that $S$ is strictly henselian.
The sheaves of vanishing cycles are easily computed with the aid of
2.2.3 and XII, 3.7.  The cohomology with coefficients in $\Lambda$
considered here follows from the $\ell$-adic cohomology of the cited
passage by the universal-coefficient formula (cf. 2.1.13).  The
results are as follows.

\medskip
\noindent\underline{A. Cohomology without supports.}

\SGAvar{One has $R^i\psi_{\eta}(\Lambda)=0$ for $i\ne0,n$.

\begin{enumerate}
\item[1)] \underline{If $n\ne0$:}
$\psi_{\bar\eta}(\Lambda)=\Lambda$  Moreover,
$R^n\psi_{\bar\eta}(\Lambda)$ is (noncanonically) isomorphic to the
sheaf $\Lambda$ on $x_0$, extended by zero to $X_s$.

\item[2)] \underline{For arbitrary $n$:}
even when $n=0$, one checks that
$R^i\Phi_{\bar\eta}(\Lambda)=0$ for $i\ne n$, and that
$R^n\Phi_{\bar\eta}(\Lambda)$ is (noncanonically) isomorphic to the
sheaf $\Lambda$ on $x_0$, extended by zero to $X_s$.
\end{enumerate}

}{One has $R^i\psi_{\bar\eta}(\Lambda)=0$ for $i\ne0,n$.

\begin{enumerate}
\item[1)] \underline{If $n\ne0$:}
$\psi_{\bar\eta}(\Lambda)=\Lambda$.  Moreover,
$R^n\psi_{\bar\eta}(\Lambda)$ is (noncanonically) isomorphic to the
sheaf $\Lambda$ on $x_0$, extended by zero to $X_s$.

\item[2)] \underline{For arbitrary $n$:}
even when $n=0$, one checks that
$R^i\Phi_{\bar\eta}(\Lambda)=0$ for $i\ne n$, and that
$R^n\Phi_{\bar\eta}(\Lambda)$ is (noncanonically) isomorphic to the
sheaf $\Lambda$ on $x_0$, extended by zero to $X_s$.
\end{enumerate}

}{E7II-XV-012}\SGApage{II}{\(\approx\)184}
\medskip
\noindent\underline{B. Cohomology with support at $x_0$.}

One has
\[
H^i_{\{x_0\}}\bigl(X_s,R\psi_{\bar\eta}(\Lambda)\bigr)=0
\qquad(i\ne n,2n).
\]
There is, moreover, a trace morphism
\begin{equation}\tag{2.2.5.1}
\operatorname{tr}:
H^{2n}_{\{x_0\}}\bigl(X_s,R\psi_{\bar\eta}(\Lambda(n))\bigr)
\longrightarrow\Lambda.
\end{equation}
\underline{If $n\ne0$:} the morphism (2.2.5.1) is an isomorphism.
Moreover,
\[
H^n_{\{x_0\}}\bigl(X_s,R\psi_{\bar\eta}(\Lambda(n))\bigr)
\]
is (noncanonically) isomorphic to $\Lambda$.

\SGAvar{\medskip
\noindent\underline{C. Duality.}
For
\[
a\in R^n\psi_{\bar\eta}(\Lambda)_{x_0},
\qquad
b\in H^n_{\{x_0\}}\bigl(X_s,R\psi_{\bar\eta}(\Lambda(n))\bigr),
\]
put $(a,b)=\operatorname{Tr}(a\wedge b)$, where $\operatorname{Tr}$
is the morphism XIII, 2.2.6
\[
\operatorname{Tr}:
H^{2n}_{\{x_0\}}\bigl(X_s,R\psi_{\eta}(\Lambda(n))\bigr)
\longrightarrow\Lambda.
\]
It follows from Poincar\'e duality on $X_{\bar\eta}$ and the
isomorphisms 2.2.3 that $(a,b)$ gives a perfect duality between the
free $\Lambda$-modules
\[
R^n\psi_{\bar\eta}(\Lambda)_{x_0}
\quad\text{and}\quad
H^n_{\{x_0\}}\bigl(X_s,R\psi_{\bar\eta}(\Lambda(n)).
\]



}{\medskip
\noindent\underline{C. Duality.}
For
\[
a\in R^n\psi_{\bar\eta}(\Lambda)_{x_0},
\qquad
b\in H^n_{\{x_0\}}\bigl(X_s,R\psi_{\bar\eta}(\Lambda(n))\bigr),
\]
put $(a,b)=\operatorname{Tr}(a\wedge b)$, where $\operatorname{Tr}$
is the morphism XIII, 2.2.6
\[
\operatorname{Tr}:
H^{2n}_{\{x_0\}}\bigl(X_s,R\psi_{\bar\eta}(\Lambda(n))\bigr)
\longrightarrow\Lambda.
\]
It follows from Poincar\'e duality on $X_{\bar\eta}$ and the
isomorphisms 2.2.3 that $(a,b)$ gives a perfect duality between the
free $\Lambda$-modules
\[
R^n\psi_{\bar\eta}(\Lambda)_{x_0}
\quad\text{and}\quad
H^n_{\{x_0\}}\bigl(X_s,R\psi_{\bar\eta}(\Lambda(n))\bigr).
\]



}{E7II-XV-013}\medskip
\noindent\underline{D. $n=2m>0$.}

By XII, 3.7, the groups
\[
H^n_{\{x_0\}}\bigl(X_s,R\psi_{\bar\eta}(\Lambda(m))\bigr)
\quad\text{and}\quad
R^n\psi_{\bar\eta}(\Lambda(m))_{x_0}
\]
have natural generators, defined up to sign; denote them by
$\delta$ and $\delta'$, respectively.  Let $\varphi$ be the natural map
from the first of these groups to the second.  The signs of $\delta$
and $\delta'$ may be normalized so that
\begin{equation}\tag{2.2.5.2}
(\delta',\delta)=1;
\end{equation}
then
\begin{equation}\tag{2.2.5.3}
\varphi(\delta)=(-1)^m\,2\,\delta'.
\end{equation}

\SGApage{II}{\(\approx\)185}
\SGAvar{Let $Z$ be the separable quadratic extension of $k(\eta)$ which is
the center of the even part of the Clifford algebra of the quadratic
form (2.2.1.1), and let $\epsilon$ be the corresponding character,
of order $1$ or $2$, of the inertia group $I$.  For $\sigma\in I$,
one has
\begin{equation}\tag{2.2.5.4}
\sigma\delta=\epsilon(\sigma)\delta
\end{equation}
and
\begin{equation}\tag{2.2.5.5}
\sigma\delta'=\epsilon(\sigma)\delta'.
\end{equation}
Formulas (2.2.5.2), (2.2.5.3), and (2.2.5.4), in cohomology with
coefficients in $\mathbb Z/2k$, imply that the variation is
\begin{equation}\tag{2.2.5.6}
\operatorname{Var}(\sigma)(a)=
\frac{\epsilon(\sigma)-1}{2}\,(-1)^m\,(a\delta)\,\delta.
\end{equation}

}{Let $Z$ be the separable quadratic extension of $k(\eta)$ which is
the center of the even part of the Clifford algebra of the quadratic
form (2.2.1.1), and let $\epsilon$ be the corresponding character,
of order $1$ or $2$, of the inertia group $I$.  For $\sigma\in I$,
one has
\begin{equation}\tag{2.2.5.4}
\sigma\delta=\epsilon(\sigma)\delta
\end{equation}
and
\begin{equation}\tag{2.2.5.5}
\sigma\delta'=\epsilon(\sigma)\delta'.
\end{equation}
Formulas (2.2.5.2), (2.2.5.3), and (2.2.5.4), in cohomology with
coefficients in $\mathbb Z/2k$, imply that the variation is
\begin{equation}\tag{2.2.5.6}
\operatorname{Var}(\sigma)(a)=
\frac{\epsilon(\sigma)-1}{2}\,(-1)^m\,(a,\delta)\,\delta.
\end{equation}

}{E7II-XV-014}\medskip
\noindent\underline{E. $n=2m+1$---Vanishing cycles.}

By XII, 3.7, the groups
\[
H^n_{\{x_0\}}\bigl(X_s,R\psi_{\bar\eta}(\Lambda(m))\bigr)
\quad\text{and}\quad
R^n\psi_{\bar\eta}(\Lambda(m+1))_{x_0}
\]
have natural generators, defined up to sign; denote them by
$\delta$ and $\delta'$, respectively.  With the notation of D, their
signs may be normalized so that
\begin{equation}\tag{2.2.5.7}
(\delta',\delta)=1;
\end{equation}
then
\begin{equation}\tag{2.2.5.8}
\varphi(\delta)=0.
\end{equation}

\medskip
\noindent\underline{F. $n=2m+1$---The variation.}

The inertia group $I$ acts trivially on the cohomology of $X_\eta$.
This follows from XII, 3.7 and from the fact that $I$ acts trivially
on the cohomology of $Y_\eta$, since $Y$ is proper and smooth over
$S$.

\SGAvar{For $\sigma$ in the inertia group one necessarily has
\begin{equation}\tag{2.2.5.9}
D(\sigma)(\delta')=\lambda(\sigma)\delta,
\qquad \lambda(\sigma)\in\Lambda(1).
\end{equation}
The additivity formula (XIII, 1.4.3.4) becomes
\[
\lambda(\sigma\tau)=\lambda(\sigma)+\lambda(\tau).
\]
Let
\[
\epsilon:I\longrightarrow\Lambda(1)=\mu_k
\]
be the homomorphism such that
\[
\sigma\bigl(\sqrt[k]{t}\bigr)
=\epsilon(\sigma)\sqrt[k]{t}
\]
for a uniformizer $t$.  In view of the structure of the tame inertia
group, $\lambda$ is necessarily a multiple of $\epsilon$.  Formula
(2.2.5.9) may therefore be rewritten
\begin{equation}\tag{2.2.5.10}
\operatorname{Var}(\sigma)(a)
=\lambda_X\epsilon(\sigma)(a\delta)\delta
\end{equation}
for a suitable $\lambda_X\in\Lambda$ (depending on $X/S$).  We shall
determine $\lambda_X$ later.



}{For $\sigma$ in the inertia group one necessarily has
\begin{equation}\tag{2.2.5.9}
D(\sigma)(\delta')=\lambda(\sigma)\delta,
\qquad \lambda(\sigma)\in\Lambda(1).
\end{equation}
The additivity formula (XIII, 1.4.3.4) becomes
\[
\lambda(\sigma\tau)=\lambda(\sigma)+\lambda(\tau).
\]
Let
\[
\epsilon:I\longrightarrow\Lambda(1)=\mu_k
\]
be the homomorphism such that
\[
\sigma\bigl(\sqrt[k]{t}\bigr)
=\epsilon(\sigma)\sqrt[k]{t}
\]
for a uniformizer $t$.  In view of the structure of the tame inertia
group, $\lambda$ is necessarily a multiple of $\epsilon$.  Formula
(2.2.5.9) may therefore be rewritten
\begin{equation}\tag{2.2.5.10}
\operatorname{Var}(\sigma)(a)
=\lambda_X\epsilon(\sigma)(a,\delta)\delta
\end{equation}
for a suitable $\lambda_X\in\Lambda$ (depending on $X/S$).  We shall
determine $\lambda_X$ later.



}{E7II-XV-015}\SGApage{II}{\(\approx\)186}
\subsubsection*{2.2.6.}
We shall give the cycles $\delta$ and $\delta'$ of D and E a
description which is local for the \`etale topology.

The case of even $n$ is easy.  If $2\nmid k$
($\Lambda=\mathbb Z/k$), then $\pm\delta$ is characterized by
(2.2.5.3) and (2.2.5.4), that is, by
\begin{equation}\tag{2.2.6.1}
(\delta,\delta)=(-1)^m\,2.
\end{equation}
In general, $\pm\delta$ is the image of any cycle
$\widetilde\delta$, in cohomology with coefficients in
$\mathbb Z/2^ak$ for $a$ sufficiently large, which satisfies
(2.2.6.1).

\SGAvar{Suppose that $n$ is odd.  Then $x_0$ is a rational point (2.2.2).
Let $X_{s(0)}$ be the henselization of $X_s$ at $x_0$, let
$\widetilde X_{s(0)}$ be its blowup at $x_0$, let $Y_0$ be the
inverse image of $x_0$ in $X_{s(0)}$ (a quadric), and put
\[
X^*_{s(0)}=X_{s(0)}-\{x_0\}
=\widetilde X_{s(0)}-Y_0.
\]
Consider the composite map (cf. XIII, 2.4.6.1)
\begin{equation}\tag{2.2.6.2}
\begin{tikzcd}[column sep=huge,row sep=huge]
H^{n-1}(Y_0,\Lambda(m))
  \& H^{n-1}(\widetilde X_{s(0)},\Lambda(m))
      \arrow[l,"\sim"'] \arrow[r]
  \& H^{n-1}(X^*_{s(0)},\Lambda(m)) \arrow[d] \\
{}
  \& H^n_{\{x\}}\bigl(X_s,R\psi_{\bar\eta}(\Lambda(m))\bigr)
  \& H^n_{\{x\}}(X_s,\Lambda(m)) \arrow[l]
\end{tikzcd}
\end{equation}

}{Suppose that $n$ is odd.  Then $x_0$ is a rational point (2.2.2).
Let $X_{s(0)}$ be the henselization of $X_s$ at $x_0$, let
$\widetilde X_{s(0)}$ be its blowup at $x_0$, let $Y_0$ be the
inverse image of $x_0$ in $\widetilde X_{s(0)}$ (a quadric), and put
\[
X^*_{s(0)}=X_{s(0)}-\{x_0\}
=\widetilde X_{s(0)}-Y_0.
\]
Consider the composite map (cf. XIII, 2.4.6.1)
\begin{equation}\tag{2.2.6.2}
\begin{tikzcd}[column sep=huge,row sep=huge]
H^{n-1}(Y_0,\Lambda(m))
  \& H^{n-1}(\widetilde X_{s(0)},\Lambda(m))
      \arrow[l,"\sim"'] \arrow[r]
  \& H^{n-1}(X^*_{s(0)},\Lambda(m)) \arrow[d] \\
{}
  \& H^n_{\{x_0\}}\bigl(X_s,R\psi_{\bar\eta}(\Lambda(m))\bigr)
  \& H^n_{\{x_0\}}(X_s,\Lambda(m)) \arrow[l]
\end{tikzcd}
\end{equation}

}{E7II-XV-016}\subsubsection*{Lemma 2.2.7.}
\underline{The map (2.2.6.2) identifies
$\mathbb H^n_{\{x_0\}}(X_s,R\psi_{\bar\eta}(\Lambda(m)))$ with the
primitive quotient of the middle-dimensional cohomology of the
$2m$-dimensional quadric $Y_0$; the two classes $\pm\delta$ are the
images of the natural generators of this primitive quotient.}

\SGApage{II}{\(\approx\)187}
\SGAvar{The cycles $\pm\delta$ are characterized by their images in
$H_c^i(X_{\bar\eta},\Lambda(m))$.  Using (2.1.7.1), it is therefore
enough to prove that the composite map
\begin{center}
\begin{tikzcd}[column sep=huge,row sep=huge]
H^{n-1}(Y_0,\Lambda(m))
  \& H^{n-1}(\widetilde X_s,\Lambda(m))
      \arrow[l,"\sim"'] \arrow[r]
  \& H^{n-1}(X_s^*,\Lambda(m)) \arrow[d] \\
H_c^n(X_{\bar\eta},\Lambda(m))
  \& H_c^n(X_s,\Lambda(m)) \arrow[l,"\operatorname{sp}"']
  \& H^n_{\{x_0\}}(X_s,\Lambda(m)) \arrow[l]
\end{tikzcd}
\end{center}
sends the distinguished generators of the primitive quotient of the
first group to the distinguished generators of the last group.

}{The cycles $\pm\delta$ are characterized by their images in
$H_c^n(X_{\bar\eta},\Lambda(m))$.  Using (2.1.7.1), it is therefore
enough to prove that the composite map
\begin{center}
\begin{tikzcd}[column sep=huge,row sep=huge]
H^{n-1}(Y_0,\Lambda(m))
  \& H^{n-1}(\widetilde X_s,\Lambda(m))
      \arrow[l,"\sim"'] \arrow[r]
  \& H^{n-1}(X_s^*,\Lambda(m)) \arrow[d] \\
H_c^n(X_{\bar\eta},\Lambda(m))
  \& H_c^n(X_s,\Lambda(m)) \arrow[l,"\operatorname{sp}"']
  \& H^n_{\{x_0\}}(X_s,\Lambda(m)) \arrow[l]
\end{tikzcd}
\end{center}
sends the distinguished generators of the primitive quotient of the
first group to the distinguished generators of the last group.

}{E7II-XV-017}Apply 2.1.8.  Up to sign, the map may be rewritten
\begin{center}
\begin{tikzcd}[column sep=huge,row sep=huge]
H^{n-1}(Y_s,\Lambda(m)) \arrow[r,"\partial"] \arrow[d,equal]
  \& H_c^n(X_s,\Lambda(m)) \arrow[d] \\
H^{n-1}(Y_{\bar\eta},\Lambda(m)) \arrow[r]
  \& H_c^n(X_{\bar\eta},\Lambda(m)),
\end{tikzcd}
\end{center}
and 2.2.7 follows.



\section*{3. The Picard--Lefschetz formula}

\subsection*{3.1. Preliminary results}

\subsubsection*{3.1.1.}
Let $S$ be a henselian trait and let $f:X\to S$ be a flat morphism
of finite type, purely of relative dimension $n$.  Suppose that the
special fiber $X_s$ is smooth except at a finite set of points
$\Sigma'$, at which $X_s$ has an ordinary quadratic singularity.

\SGApage{II}{\(\approx\)188}
Let $\Sigma$ be the set of points $x\in\Sigma'$ such that the
generic fiber $X_\eta$ is smooth in a neighborhood of $x$.  Finally,
let $\Lambda$ be a torsion ring whose order is prime to the residual
characteristic of $S$.

\SGAvar{\subsubsection*{Proposition 3.1.2.}
\begin{enumerate}
\item[(i)] \underline{The sheaves $R^i\Phi(\Lambda)$ vanish for
$i\ne n$.}

\item[(ii)] \underline{The sheaf $R^n\Phi(\Lambda)$ is alone outside
$\Sigma$.  Its restriction to $\Sigma$ is a sheaf of
$\Lambda$-modules of rank one.}

\item[(iii)] \underline{For every geometric point $\bar x$ of
$\Sigma$, the pairing $(a,b)$ of 2.2.5(C) puts the $\Lambda$-modules}
\[
\bigl(R^n\psi(\Lambda)\bigr)_{\bar x}
\quad\text{and}\quad
H^n_{\{\bar x\}}\bigl(R^n\psi_\eta(\Lambda)\bigr)
\]
\underline{in duality.  They are free of rank one if $n\ne0$, and
of rank two if $n=0$.}
\end{enumerate}


}{\subsubsection*{Proposition 3.1.2.}
\begin{enumerate}
\item[(i)] \underline{The sheaves $R^i\Phi(\Lambda)$ vanish for
$i\ne n$.}

\item[(ii)] \underline{The sheaf $R^n\Phi(\Lambda)$ is zero outside
$\Sigma$.  Its restriction to $\Sigma$ is a sheaf of
$\Lambda$-modules of rank one.}

\item[(iii)] \underline{For every geometric point $\bar x$ of
$\Sigma$, the pairing $(a,b)$ of 2.2.5(C) puts the $\Lambda$-modules}
\[
\bigl(R^n\psi(\Lambda)\bigr)_{\bar x}
\quad\text{and}\quad
H^n_{\{\bar x\}}\bigl(X_{\bar s},R\psi_{\bar\eta}(\Lambda)\bigr)
\]
\underline{in duality.  They are free of rank one if $n\ne0$, and
of rank two if $n=0$.}
\end{enumerate}


}{E7II-XV-018}The assertions of 3.1.2 are local for the \`etale topology.  We may
therefore treat only the case where $S$ is strictly henselian and
$X$ is the hypersurface in relative affine space
$\mathbb A^{n+1}_S$ defined by a quadratic equation satisfying
2.2.1(a) and (b) (apply 1.3.2).

The morphism $f$ is smooth at every point of $X_s-\Sigma'$, hence
$R\Phi=0$ outside $\Sigma'$.  By 2.2.4, $R\Phi$ even vanishes outside
$\Sigma$.  We may now assume 2.2.5(*); assertions (ii) and (iii) then
follow from 2.2.5(A), (B), and (C).


\subsection*{3.2. Even relative dimension}

We retain the notation of 3.1 and take $n=2m$ to be even.  For
simplicity we suppose that $S$ is strictly henselian (the general
case would follow by descent).

\SGAvar{\subsubsection*{Proposition 3.2.1.}
\begin{enumerate}
\item[(i)] \underline{For $x\in\Sigma$, the group
$H^n_{\{x\}}(R^n\psi_\eta(\Lambda(m))$ has a natural
generator $\delta$, well defined up to sign, characterized by being
natural in $\Lambda$ and by satisfying}
\[
(\delta,\delta)=(-1)^m\,2.
\]
\underline{For $n=0$, add the condition
$\operatorname{Tr}(\delta)=0$.}

\item[(ii)] \underline{For a suitable character, independent of
$\Lambda$,}
\[
\varepsilon_x:I\longrightarrow\{\pm1\}
\]
\underline{of the inertia group $I$, the variation is}
\[
\operatorname{Var}(\sigma)(a)
=(-1)^m\left(\frac{\varepsilon_x(\sigma)-1}{2}\right)
(a\delta)\delta
\]
\underline{for $a\in R^n\Phi_{\bar\eta}(\Lambda(m))$.  Consequently,}
\[
\sigma(\delta)=\varepsilon_x(\sigma)\delta.
\]
\end{enumerate}


}{\subsubsection*{Proposition 3.2.1.}
\begin{enumerate}
\item[(i)] \underline{For $x\in\Sigma$, the group
$H^n_{\{x\}}(X_s,R\psi_{\bar\eta}(\Lambda(m)))$ has a natural
generator $\delta$, well defined up to sign, characterized by being
natural in $\Lambda$ and by satisfying}
\[
(\delta,\delta)=(-1)^m\,2.
\]
\underline{For $n=0$, add the condition
$\operatorname{Tr}(\delta)=0$.}

\item[(ii)] \underline{For a suitable character, independent of
$\Lambda$,}
\[
\varepsilon_x:I\longrightarrow\{\pm1\}
\]
\underline{of the inertia group $I$, the variation is}
\[
\operatorname{Var}(\sigma)(a)
=(-1)^m\left(\frac{\varepsilon_x(\sigma)-1}{2}\right)
(a,\delta)\delta
\]
\underline{for $a\in R^n\Phi_{\bar\eta}(\Lambda(m))$.  Consequently,}
\[
\sigma(\delta)=\varepsilon_x(\sigma)\delta.
\]
\end{enumerate}


}{E7II-XV-019}\SGApage{II}{\(\approx\)189}
Passing to the inverse limit, it is enough to treat the universal
case $\Lambda=\mathbb Z_\ell$.  Up to sign, only one $\delta$ can
then satisfy the conditions in (i).  As in 3.1.2, we reduce to the
case 2.2.5(*) and apply 2.2.5(D).

\subsubsection*{Complement 3.2.2.}
Suppose that, at $x\in\Sigma$, the henselization $X_{(x)}$ of $X$ at
$x$ is isomorphic to the henselization at $x_0$ of the projective
scheme $Q$ with equation
\[
\sum_{i\leq j}a_{ij}x_ix_j=0.
\]
Then the character $\varepsilon_x$ is defined by the separable
quadratic extension $Z(C^+(Q))$ of the fraction field.  This follows
from 2.2.5(D).

\subsubsection*{3.2.3.}
In residual characteristic different from $2$, Complement 3.2.2 can
be made more explicit.  In that case $I$ has a unique nontrivial
character $\varepsilon$ of order $2$.  For a uniformizer $t$, one
has
\[
\sigma(\sqrt t)=\varepsilon(\sigma)\sqrt t.
\]
On the other hand, $X_{(x)}$ is then isomorphic to the henselization
at $0$ of an affine scheme
\[
\sum a_{ij}x_ix_j=b,
\]
where $b$ lies in the maximal ideal and
$Q(x)=\sum a_{ij}x_ix_j$ is nondegenerate.  The ideal $(b)$ defines
the image of the subscheme of $X_{(x)}$ which is the support of
$J^n_{X_{(x)}/S}$.  The center of the Clifford algebra is isomorphic
to
\[
k(\eta)\left(\left((-1)^{m+1}\cdot2b\det(a_{ij})\right)^{1/2}\right)
\]
(Bourbaki, \emph{Alg\`ebre}, Ch.~9, \S9, no.~4, final remark), whence
\[
\varepsilon_x=\varepsilon^{v(b)}.
\]
The variation is zero if $b$ has even valuation; otherwise
$\varepsilon_x=\varepsilon$.


\SGApage{II}{190}
\subsection*{3.3. Odd relative dimension}

To prove the Picard--Lefschetz formula in odd relative dimension, we
shall have to reduce to the transcendental case treated in XIV,
3.2.1.1.  The specialization argument required for this forces us to
work over a base which is no longer one-dimensional.  We shall use
ad hoc arguments adapted to the precise case we need: I do not have
a usable theory over a base of dimension greater than one, and I have
preferred not to cover my ignorance with generalities.

\subsubsection*{3.3.1.}
Let $S=\operatorname{Spec}(A)$ be a strictly local noetherian scheme,
and let $f:X\to S$ be a flat morphism of finite type, purely of
relative dimension $n$, whose special fiber $X_s$ has no nonsmooth
points other than nondegenerate quadratic points.  Let $s$ be the
closed point of $S$, let $X_s$ be the special fiber, and let
$x_0\in X_s$ be a singular point.

By 1.3.2(i), in a neighborhood of $x_0$ the scheme $X/S$ is locally
isomorphic to an $S$-scheme defined in $A^{n+1}_S$ by an
equation $Q-b=0$, where $b$ belongs to the maximal ideal and $Q$ is a
nondegenerate quadratic form.  The ideal $(b)$ of $A$ depends only
on $(X/S,x_0)$: in a neighborhood of $x$ there is a section
\[
s\in X(A/(b))
\]
of $X$ over $\operatorname{Spec}(A/(b))$ such that
\[
J^n_{X/S}\simeq
\mathcal O_{s(\operatorname{Spec}(A/b))}.
\]
We denote this section by $s_x$, and denote by $b(x)$ any generator
of $(b)$.

\SGApage{II}{\(\approx\)191}
\subsubsection*{3.3.2.}
Suppose that $S$ is normal, and let $b\in A$.  If $n$ is invertible
in $A$ and $b\ne0$, write $\varepsilon_b^{(n)}$ for the following
homomorphism from the Galois group of the fraction field of $A$ to
$\mathbb Z/(n)(1)$:
\[
\varepsilon_b^{(n)}(\sigma)
= (\sigma\sqrt[n]{b})/\sqrt[n]{b}.
\]
One has
\[
\bigl(\varepsilon_b^{(nm)}\bigr)^m=\varepsilon_b^{(n)}.
\]
If $\Lambda$ is a $\mathbb Z/n$-algebra, we again denote by
$\varepsilon_b$ the composite
\[
\operatorname{Gal}\longrightarrow\mathbb Z/(n)(1)
\longrightarrow\Lambda(1).
\]
It is independent of $n$.

The character $\varepsilon_b$ factors through
\[
\pi_1\bigl(S-(\text{the subscheme }b=0)\bigr).
\]

\subsubsection*{3.3.3.}
Suppose that $S$ is a strictly henselian trait and that $n$ is odd,
$n=2m+1$.  Let $x$ be a singular point of $X_s$ in a neighborhood
of which $X_\eta$ is smooth, and let $X_1/S$ be a subscheme of
$\mathbb A^{n+1}_S$ with equation $Q-b=0$, as in 1.3.2(i).  The
assumption on $x$ means that $b$, although zero modulo the maximal
ideal, is not zero.

\SGAvar{Let $X_{s(x)}$ be the henselization of $X_s$ at $x$, let
$\widetilde X_{s(x)}$ be the blowup of $\widetilde X_s$ at $x$, let $Y$ be the
quadric which is the inverse image of $x$, and put
\[
\widetilde X^*_{s(x)}=X_{s(x)}-\{x\}=\widetilde X_{s(x)}-Y.
\]
The composite morphism
\begin{center}
\begin{tikzcd}[column sep=huge,row sep=huge]
H^{n-1}(Y,\Lambda(m))
  \& H^{n-1}(\widetilde X_{s(x)},\Lambda(m))
      \arrow[l,"\sim"'] \arrow[r]
  \& H^{n-1}(X^*_{s(x)},\Lambda(m)) \arrow[d] \\
{}
  \& H^n_{\{x\}}\bigl(X_s,R\psi_{\bar\eta}(\Lambda(m))\bigr)
  \& H^n_{\{x\}}(X_s,\Lambda(m)) \arrow[l]
\end{tikzcd}
\end{center}
identifies the latter group with the primitive quotient (XII, 3.5)
of the cohomology of the quadric $Y$.  Indeed, this follows from
2.2.7 in a locally isomorphic situation.  The images of the two
natural generators of this quotient (the cohomology classes of the
rulings of $Y$) are the two basis vectors
\[
\pm\delta\in
H^n_{\{x\}}\bigl(X_s,R\psi_{\bar\eta}(\Lambda(m))\bigr).
\]



}{Let $X_{s(x)}$ be the henselization of $X_s$ at $x$, let
$\widetilde X_{s(x)}$ be the blowup of $X_s$ at $x$, let $Y$ be the
quadric which is the inverse image of $x$, and put
\[
X^*_{s(x)}=X_{s(x)}-\{x\}=\widetilde X_{s(x)}-Y.
\]
The composite morphism
\begin{center}
\begin{tikzcd}[column sep=huge,row sep=huge]
H^{n-1}(Y,\Lambda(m))
  \& H^{n-1}(\widetilde X_{s(x)},\Lambda(m))
      \arrow[l,"\sim"'] \arrow[r]
  \& H^{n-1}(X^*_{s(x)},\Lambda(m)) \arrow[d] \\
{}
  \& H^n_{\{x\}}\bigl(X_s,R\psi_{\bar\eta}(\Lambda(m))\bigr)
  \& H^n_{\{x\}}(X_s,\Lambda(m)) \arrow[l]
\end{tikzcd}
\end{center}
identifies the latter group with the primitive quotient (XII, 3.5)
of the cohomology of the quadric $Y$.  Indeed, this follows from
2.2.7 in a locally isomorphic situation.  The images of the two
natural generators of this quotient (the cohomology classes of the
rulings of $Y$) are the two basis vectors
\[
\pm\delta\in
H^n_{\{x\}}\bigl(X_s,R\psi_{\bar\eta}(\Lambda(m))\bigr).
\]



}{E7II-XV-020}\SGApage{II}{\(\approx\)192}
\subsubsection*{3.3.4.}
Return to the more general assumptions of 3.3.1, with $n=2m+1$.
For every singular point $x$ of $X_s$, we again denote by
$\pm\delta_x$ the sections of $R^nf_!\Lambda(m)$ over $S$ which are
the images of the generators of the primitive quotient of
$H^{n-1}(Y,\Lambda(m))$ (with the same notation as above) under the
composite morphism
\begin{center}
\begin{tikzcd}[column sep=huge,row sep=huge]
H^{n-1}(Y,\Lambda(m))
  \& H^{n-1}(\widetilde X_{s(x)},\Lambda(m))
      \arrow[l,"\sim"'] \arrow[r]
  \& H^{n-1}(X^*_{s(x)},\Lambda(m)) \arrow[d] \\
H^0\bigl(S,R^nf_!\Lambda(m)\bigr)
  \& H_c^n(X_s,\Lambda(m)) \arrow[l,"\sim"']
  \& H^n_{\{x\}}(X_s,\Lambda(m)) \arrow[l]
\end{tikzcd}
\end{center}
The class $\pm\delta_x$ is the \underline{vanishing cycle} at $x$.

When $S$ is a trait, the class
\[
\pm\delta\in H_c^n(X_{\bar\eta},\Lambda(m))
\]
is the image under XIII, 2.1.8.4 of the class of 3.3.3
\[
\pm\delta\in
H^n_{\{x\}}\bigl(X_s,R\psi_{\bar\eta}(\Lambda(m))\bigr).
\]

We shall prove the following two results simultaneously.

\subsubsection*{Proposition 3.3.5.}
\underline{Let $S$ be a henselian trait and let $X/S$ be as in
3.3.3.  The variation is given by}
\[
\operatorname{Var}(\sigma)(a)
=(-1)^{m+1}\varepsilon_{b(x)}(\sigma)(a,\delta)\delta.
\]

\subsubsection*{Proposition 3.3.6.}
\underline{Let $S$ be strictly local and regular, and let
$f:X\to S$ be a proper flat morphism of relative dimension
$n=2m+1$ such that $X_s$ has exactly one nonsmooth point $x$, assumed
to be a nondegenerate quadratic point.  Suppose also that $b(x)$ is
a parameter.  Then Galois acts on the cohomology of the geometric
generic fiber by}
\[
\sigma(a)=a+(-1)^{m+1}
\varepsilon_{b(x)}(\sigma)(a,\delta_x)\delta_x.
\]


\medskip
\noindent\textup{(A)}
Let $k$ be an algebraically closed field of characteristic $0$, and
let $S$ be the henselization of $\operatorname{Spec}(k[T])$ at $0$.
We prove 3.3.5 when $b(x)$ is a uniformizer.

\SGApage{II}{\(\approx\)193}
By the Lefschetz principle, we reduce to the case $k=\mathbb C$.
Since the question is local, we may suppose that $f:X\to S$ is
obtained by base change from the map
\[
\sum_i z_i^2:\mathbb C^{n+1}\longrightarrow\mathbb C.
\]
The construction of $\pm\delta$ given above has a transcendental
analogue in integral cohomology.  Via XIV, 2.1, the class
$\pm\delta$ therefore corresponds to a transcendental vanishing
cycle, the image of
\[
\pm\delta_{\mathrm{cl}}
\in H^n_{\{0\}}\bigl(X_0,R\psi_{\mathrm{cl}\eta}(\mathbb Z)\bigr).
\]
By XIV, 2.1 and the corresponding property of the algebraic cycles
$\pm\delta$, the image of $\pm\delta_{\mathrm{cl}}$ in
\[
H^n_{\{0\}}\bigl(X_0,R\psi_{\mathrm{cl}\eta}(\mathbb Z)\bigr)
\otimes\mathbb Z_\ell
\]
is a generator of this group.  Thus $\pm\delta_{\mathrm{cl}}$ are the
two generators of
$H^n_{\{0\}}(X_0,R\psi_{\mathrm{cl}\eta}(\mathbb Z))$, and coincide
with the vanishing cycles of XIV, 3.2.  Formula 3.3.5 consequently
follows from XIV, 3.2.11.  Strictly speaking, one should have checked
compatibility between the isomorphism of XIV, 2.1 and cup products
and traces.

\medskip
\noindent\textup{(B)}
We prove 3.3.6 for $X/S$ as in (A), with $X/S$ proper in addition.
This follows from XIII, 2.4.6.2.

\medskip
\noindent\textup{(C)}
Let $p$ be a prime number, and let $S_0$ be the strict henselization
of $\operatorname{Spec}(\mathbb Z[T])$ at $(p,T)$.  We prove 3.3.6
for $X/S_0$ such that $b(x)=T$.

\SGAvar{The restriction of the sheaf $R^if_*\Lambda$ to the strict
henselization of $S$ at $(T)$ is
governed by (B).  This sheaf is
locally constant on the complement of the smooth divisor $T=0$,
which has mixed characteristic.  Assertion (C) therefore follows
from (B) and Abhyankar's lemma.


}{The restriction of the sheaf $R^if_*\Lambda$ to the strict
henselization of $S_0$ at $(T)$ is
governed by (B).  This sheaf is
locally constant on the complement of the smooth divisor $T=0$,
which has mixed characteristic.  Assertion (C) therefore follows
from (B) and Abhyankar's lemma.


}{E7II-XV-021}\par\medskip
\noindent\textup{(D)}
In at least one case to which (C) applies, $\delta_x$ is not a torsion
element in $\ell$-adic cohomology.

\SGApage{II}{\(\approx\)194}
Let $S_1$ be the strict henselization of
$\operatorname{Spec}(\mathbb Z)$ at $(p)$, and let
$V\subset\mathbb P^r/S_1$ be a smooth projective variety over $S_1$,
purely of relative dimension $n+1$.  Suppose that there is a
Lefschetz pencil of hyperplane sections for the special fiber $V_p$
of $V/S_1$ (XVII, 2.5).  This pencil is defined by a line $d'$ in the
special fiber of the projective space $\check{\mathbb P}^{,r}/S_1$
dual to $\mathbb P^r/S_1$.  Lift $d'$ to a line $d$ in
$\check{\mathbb P}^{,r}/S_1$, and let $X\to d$ be the pencil of
hyperplane sections of $V$ parametrized by $d$ (XVIII, 3.1).  The
exceptional locus $\Delta$ of $d$ is \`etale over $S_1$ (XVII, 6).
After \`etale localization at a point of $\Delta_p\subset d$, we
therefore obtain a situation of type (C).

To check that $\delta_x$ is not torsion, it is enough to do so after
extending scalars from $S_1$ to $\mathbb C$.  Take
$V=\mathbb P^{n+1}$ with the projective embedding defined by
$\mathcal O(N)$.  For $N\ge2$, XVII, 2.5 applies (it is valid for
every $N$ in this particular case).  We then use the following lemma.

\subsubsection*{Lemma 3.3.7.}
\underline{For a Lefschetz pencil of hypersurfaces of sufficiently
large degree $N$ in $\mathbb P^{n+1}(\mathbb C)$ ($N\ge3$ suffices),
a vanishing cycle in $\mathbb Q_\ell$-cohomology is nonzero.}

This lemma follows from the theory of Lefschetz pencils, for example
from XVIII, 6.6.1.  There is no circularity in using here this result
of XVIII, based on the ``Picard--Lefschetz formula,'' since (A) is
already available.  One may instead use the more direct, purely
transcendental variant XIX, no.~4, which depends only on XIV, 3.2.

\medskip
\noindent\textup{(E)}
We prove 3.3.5 in residual characteristic $p$ for
$(X/S,x\in X_s)$.

Let $X_0/S_0$ be as in (D).  By 1.3.2(i), there is a morphism
$g:S\to S_0$ such that the henselization at $x_0$ of
\[
g^*X_0=X_0\times_{S_0}S
\]
is $S$-isomorphic to the henselization of $X$ at $x$.  It is
therefore enough to prove 3.3.5 for $g^*X_0$.  By base change from
(C) for $X_0/S_0$, one obtains
\[
\sigma(a)=a-(-1)^m
\varepsilon_{b(x)}(\sigma)(a,\delta_x)\delta_x
\]
in $H^n(X_{0\bar\eta})$, for $\sigma$ in the inertia group.  Let
$\pi$ be a uniformizer and compare this formula with (2.2.5.9), where
$\varepsilon=\varepsilon_\pi$.  We find
\[
\lambda_X\varepsilon_\pi(\sigma)(a,\delta_x)\delta_x
=-(-1)^m\varepsilon_{b(x)}(\sigma)(a,\delta_x)\delta_x
\]
in $H^n(X_{0\bar\eta})$.  Since $\delta_x$ is not torsion, we may
cancel $(a,\delta_x)\delta_x$ from this identity.  Formula (2.2.5.9),
with the resulting value of $\lambda_X$, is precisely 3.3.5.

\SGApage{II}{\(\approx\)195}
\medskip
\noindent\textup{(F) Proof of 3.3.6 in the general case.}

The general case of 3.3.6 reduces to the general case of 3.3.5 just
as (B) and (C) reduce to (A).


\subsection*{3.4. Summary}

Propositions 3.1.2, 3.2.1, 3.2.2, 3.2.3, and 3.3.5 constitute
the ``Picard--Lefschetz formula.''  This formula is the algebraic
analogue of XIV, 3.2.11, but takes very different forms according to
the parity of the relative dimension.

In the most interesting case, where $X/S$ is proper, the results may
be summarized as follows.

\subsubsection*{Theorem 3.4 (Picard--Lefschetz formula).}
\underline{Let $S$ be a strictly henselian trait and let
$f:X\to S$ be a proper flat morphism, purely of relative dimension
$n=2m$ or $2m+1$.  Suppose that $X_\eta$ is smooth and that $X_s$
has only ordinary quadratic singularities.  Let $\Sigma$ be the
singular locus of $X_s$.}

\SGAvar{\begin{enumerate}
\item[(i)] \underline{To every $x\in\Sigma$ is attached a vanishing
cycle}
\[
\pm\delta_x\in H_n(X_{\bar\eta})(m).
\]
\underline{The $\delta_x$ are pairwise orthogonal.  When $n$ is even,}
\[
(\delta_x,\delta_x)=(-1)^m,2.
\]

\item[(ii)] \underline{There are isomorphisms}
\[
H^i(X_s)\xrightarrow{\ \sim\ }H^i(X_{\bar\eta})
\]
\underline{for $i\ne n,n+1$, and an exact sequence}
\[
0\longrightarrow H^n(X_s)
 \longrightarrow H^n(X_{\bar\eta})
 \xrightarrow{\,((\,\cdot\,,\delta_x))_{x\in\Sigma}\,}
 \bigoplus_{x\in\Sigma}\Lambda(n-m)
 \longrightarrow H^{n+1}(X_s)
 \longrightarrow H^{n+1}(X_{\eta})
 \longrightarrow0.
\]

\item[(iii)] \underline{There are defined local characters of the
inertia group $I$:}
\[
\begin{array}{lll}
n\text{ even}:&\varepsilon_x:I\longrightarrow\{\pm1\},\\[2pt]
n\text{ odd}:&\varepsilon_x:I\longrightarrow\Lambda(1).
\end{array}
\]
\underline{When $X$ is regular at $x$, these characters are
surjective; when $x$ is nondegenerate, they are respectively the
unique character of order $2$ and the canonical character.}

\underline{The action of the inertia group} $I$ \underline{is given by}
\[
n=2m\qquad
\sigma(a)=a+(-1)^m\sum_{x\in\Sigma}
 \frac{\varepsilon_x(\sigma)-1}{2}
 (a\delta_x)\delta_x,
\]
\[
n=2m+1\qquad
\sigma(a)=a-(-1)^m\sum_{x\in\Sigma}
 \varepsilon_x(\sigma)(a\delta_x)\delta_x.
\]
\underline{In particular,}
\[
\begin{array}{lll}
n\text{ even}:&\sigma(\delta_x)=\varepsilon_x(\sigma)\delta_x,\\[2pt]
n\text{ odd}:&\sigma(\delta_x)=\delta_x.
\end{array}
\]
\end{enumerate}


}{\begin{enumerate}
\item[(i)] \underline{To every $x\in\Sigma$ is attached a vanishing
cycle}
\[
\pm\delta_x\in H^n(X_{\bar\eta})(m).
\]
\underline{The $\delta_x$ are pairwise orthogonal.  When $n$ is even,}
\[
(\delta_x,\delta_x)=(-1)^m,2.
\]

\item[(ii)] \underline{There are isomorphisms}
\[
H^i(X_s)\xrightarrow{\ \sim\ }H^i(X_{\bar\eta})
\]
\underline{for $i\ne n,n+1$, and an exact sequence}
\[
0\longrightarrow H^n(X_s)
 \longrightarrow H^n(X_{\bar\eta})
 \xrightarrow{\,((\,\cdot\,,\delta_x))_{x\in\Sigma}\,}
 \bigoplus_{x\in\Sigma}\Lambda(n-m)
 \longrightarrow H^{n+1}(X_s)
 \longrightarrow H^{n+1}(X_{\bar\eta})
 \longrightarrow0.
\]

\item[(iii)] \underline{There are defined local characters of the
inertia group $I$:}
\[
\begin{array}{lll}
n\text{ even}:&\varepsilon_x:I\longrightarrow\{\pm1\},\\[2pt]
n\text{ odd}:&\varepsilon_x:I\longrightarrow\Lambda(1).
\end{array}
\]
\underline{When $X$ is regular at $x$, these characters are
surjective; when $x$ is nondegenerate, they are respectively the
unique character of order $2$ and the canonical character.}

\underline{The action of the inertia group} $I$ \underline{is given by}
\[
n=2m\qquad
\sigma(a)=a+(-1)^m\sum_{x\in\Sigma}
 \frac{\varepsilon_x(\sigma)-1}{2}
 (a,\delta_x)\delta_x,
\]
\[
n=2m+1\qquad
\sigma(a)=a-(-1)^m\sum_{x\in\Sigma}
 \varepsilon_x(\sigma)(a,\delta_x)\delta_x.
\]
\underline{In particular,}
\[
\begin{array}{lll}
n\text{ even}:&\sigma(\delta_x)=\varepsilon_x(\sigma)\delta_x,\\[2pt]
n\text{ odd}:&\sigma(\delta_x)=\delta_x.
\end{array}
\]
\end{enumerate}


}{E7II-XV-022}\markboth{Exposé XV}{Exposé XV}

\ifnum\SGAform=1
\clearpage\section*{Critical apparatus}\begingroup\small\raggedright\sloppy
Readings are given in the exact TeX notation of the two English witnesses. Decimal suffixes identify the individual differences in each marked passage.

\par\medskip\noindent\hypertarget{app-E7II-XV-001.1}{}\textbf{E7II-XV-001.1}\quad printed folio 166.\par
\noindent Printed reading: {\ttfamily\footnotesize <\allowbreak{}}\par
\noindent Corrected reading: {\ttfamily\footnotesize \textbackslash{}leq\allowbreak{}}\par
\noindent Source: {\ttfamily\footnotesize fr/sga7ii/\{diplomati\allowbreak{}c,current\}/expose\_XV\allowbreak{}\_body.tex;\allowbreak{} folio\allowbreak{} 166;\allowbreak{} formula\allowbreak{} 11;\allowbreak{} port/decisions/decis\allowbreak{}ions\_en7II\_r2\_002.js\allowbreak{}on\#C650;\allowbreak{} fixes/FIXLOG\_7II\_CUR\allowbreak{}RENT.jsonl\#C650;\allowbreak{} 7II:260\allowbreak{}}\par

\par\medskip\noindent\hypertarget{app-E7II-XV-002.1}{}\textbf{E7II-XV-002.1}\quad printed folio 168.\par
\noindent Printed reading: {\ttfamily\footnotesize [(1)]\allowbreak{}}\par
\noindent Corrected reading: \emph{(empty)}\par
\noindent Source: {\ttfamily\footnotesize port/decisions/decis\allowbreak{}ions\_en7II\_r2\_002.js\allowbreak{}on\#C653;\allowbreak{} 7II:263\allowbreak{}}\par

\par\medskip\noindent\hypertarget{app-E7II-XV-002.2}{}\textbf{E7II-XV-002.2}\quad printed folio 168.\par
\noindent Printed reading: {\ttfamily\footnotesize 4\allowbreak{}}\par
\noindent Corrected reading: {\ttfamily\footnotesize 1\allowbreak{}}\par
\noindent Source: {\ttfamily\footnotesize fr/sga7ii/\{diplomati\allowbreak{}c,current\}/expose\_XV\allowbreak{}\_body.tex;\allowbreak{} folio\allowbreak{} 168\allowbreak{}}\par

\par\medskip\noindent\hypertarget{app-E7II-XV-002.3}{}\textbf{E7II-XV-002.3}\quad printed folio 168.\par
\noindent Printed reading: \emph{(empty)}\par
\noindent Corrected reading: {\ttfamily\footnotesize .1\allowbreak{}}\par
\noindent Source: {\ttfamily\footnotesize fr/sga7ii/\{diplomati\allowbreak{}c,current\}/expose\_XV\allowbreak{}\_body.tex;\allowbreak{} folio\allowbreak{} 168\allowbreak{}}\par

\par\medskip\noindent\hypertarget{app-E7II-XV-002.4}{}\textbf{E7II-XV-002.4}\quad printed folio 168.\par
\noindent Printed reading: {\ttfamily\footnotesize =\allowbreak{}}\par
\noindent Corrected reading: \emph{(empty)}\par
\noindent Source: {\ttfamily\footnotesize fr/sga7ii/\{diplomati\allowbreak{}c,current\}/expose\_XV\allowbreak{}\_body.tex;\allowbreak{} folio\allowbreak{} 168\allowbreak{}}\par

\par\medskip\noindent\hypertarget{app-E7II-XV-002.5}{}\textbf{E7II-XV-002.5}\quad printed folio 168.\par
\noindent Printed reading: \emph{(empty)}\par
\noindent Corrected reading: {\ttfamily\footnotesize \textbackslash{}bar\allowbreak{}}\par
\noindent Source: {\ttfamily\footnotesize fr/sga7ii/\{diplomati\allowbreak{}c,current\}/expose\_XV\allowbreak{}\_body.tex;\allowbreak{} folio\allowbreak{} 168;\allowbreak{} formula\allowbreak{} 99;\allowbreak{} port/decisions/decis\allowbreak{}ions\_en7II\_r2\_002.js\allowbreak{}on\#C654;\allowbreak{} fixes/FIXLOG\_7II\_CUR\allowbreak{}RENT.jsonl\#C654;\allowbreak{} 7II:265\allowbreak{}}\par

\par\medskip\noindent\hypertarget{app-E7II-XV-003.1}{}\textbf{E7II-XV-003.1}\quad printed folio 170.\par
\noindent Printed reading: {\ttfamily\footnotesize 4\allowbreak{}}\par
\noindent Corrected reading: {\ttfamily\footnotesize 1.2\allowbreak{}}\par
\noindent Source: {\ttfamily\footnotesize ED-EN-NONMATH-G024;\allowbreak{} F7II-XV-005.2;\allowbreak{} 7II:268;\allowbreak{} fr/sga7ii/\{diplomati\allowbreak{}c,current\}/expose\_XV\allowbreak{}\_body.tex;\allowbreak{} folio\allowbreak{} 170\allowbreak{}}\par

\par\medskip\noindent\hypertarget{app-E7II-XV-004.1}{}\textbf{E7II-XV-004.1}\quad printed folio 171.\par
\noindent Printed reading: {\ttfamily\footnotesize )\allowbreak{}}\par
\noindent Corrected reading: \emph{(empty)}\par
\noindent Source: {\ttfamily\footnotesize fr/sga7ii/\{diplomati\allowbreak{}c,current\}/expose\_XV\allowbreak{}\_body.tex;\allowbreak{} folio\allowbreak{} 171;\allowbreak{} formula\allowbreak{} 252;\allowbreak{} port/decisions/decis\allowbreak{}ions\_en7II\_r2\_002.js\allowbreak{}on\#C656;\allowbreak{} fixes/FIXLOG\_7II\_CUR\allowbreak{}RENT.jsonl\#C656;\allowbreak{} 7II:269\allowbreak{}}\par

\par\medskip\noindent\hypertarget{app-E7II-XV-005.1}{}\textbf{E7II-XV-005.1}\quad printed folio 173.\par
\noindent Printed reading: {\ttfamily\footnotesize 4\allowbreak{}}\par
\noindent Corrected reading: {\ttfamily\footnotesize 1.2\allowbreak{}}\par
\noindent Source: {\ttfamily\footnotesize ED-EN-NONMATH-G025;\allowbreak{} F7II-XV-008.1;\allowbreak{} 7II:272;\allowbreak{} fr/sga7ii/\{diplomati\allowbreak{}c,current\}/expose\_XV\allowbreak{}\_body.tex;\allowbreak{} folio\allowbreak{} 173\allowbreak{}}\par

\par\medskip\noindent\hypertarget{app-E7II-XV-006.1}{}\textbf{E7II-XV-006.1}\quad printed folio 175.\par
\noindent Printed reading: \emph{(empty)}\par
\noindent Corrected reading: {\ttfamily\footnotesize \textasciicircum{}2\allowbreak{}}\par
\noindent Source: {\ttfamily\footnotesize fr/sga7ii/\{diplomati\allowbreak{}c,current\}/expose\_XV\allowbreak{}\_body.tex;\allowbreak{} folio\allowbreak{} 175;\allowbreak{} formula\allowbreak{} 424\allowbreak{}}\par

\par\medskip\noindent\hypertarget{app-E7II-XV-007.1}{}\textbf{E7II-XV-007.1}\quad printed folio 177.\par
\noindent Printed reading: {\ttfamily\footnotesize X\allowbreak{}}\par
\noindent Corrected reading: {\ttfamily\footnotesize \textbackslash{}Lambda\allowbreak{}}\par
\noindent Source: {\ttfamily\footnotesize fr/sga7ii/\{diplomati\allowbreak{}c,current\}/expose\_XV\allowbreak{}\_body.tex;\allowbreak{} folio\allowbreak{} 177;\allowbreak{} formula\allowbreak{} 504;\allowbreak{} port/decisions/decis\allowbreak{}ions\_en7II\_r2\_002.js\allowbreak{}on\#C660;\allowbreak{} fixes/FIXLOG\_7II\_CUR\allowbreak{}RENT.jsonl\#C660;\allowbreak{} 7II:278\allowbreak{}}\par

\par\medskip\noindent\hypertarget{app-E7II-XV-008.1}{}\textbf{E7II-XV-008.1}\quad printed folio 181.\par
\noindent Printed reading: \emph{(1 espace(s), 0 retour(s) de ligne dans le témoin TeX)}\par
\noindent Corrected reading: {\ttfamily\footnotesize .\allowbreak{}}\par
\noindent Source: {\ttfamily\footnotesize ED-EN-NORM-G033;\allowbreak{} complete\allowbreak{} normalized-equal\allowbreak{} formula\allowbreak{} audit;\allowbreak{} fr/sga7ii/\{diplomati\allowbreak{}c,current\}/expose\_XV\allowbreak{}\_body.tex;\allowbreak{} folio\allowbreak{} 181\allowbreak{}}\par

\par\medskip\noindent\hypertarget{app-E7II-XV-009.1}{}\textbf{E7II-XV-009.1}\quad printed folio 181.\par
\noindent Printed reading: {\ttfamily\footnotesize 7\allowbreak{}}\par
\noindent Corrected reading: {\ttfamily\footnotesize 1\allowbreak{}}\par
\noindent Source: {\ttfamily\footnotesize port/decisions/decis\allowbreak{}ions\_en7II\_r2\_002.js\allowbreak{}on\#C661;\allowbreak{} fixes/FIXLOG\_7II\_CUR\allowbreak{}RENT.jsonl\#C661;\allowbreak{} 7II:280\allowbreak{}}\par

\par\medskip\noindent\hypertarget{app-E7II-XV-010.1}{}\textbf{E7II-XV-010.1}\quad printed folio 182.\par
\noindent Printed reading: {\ttfamily\footnotesize '\allowbreak{}}\par
\noindent Corrected reading: \emph{(empty)}\par
\noindent Source: {\ttfamily\footnotesize fr/sga7ii/\{diplomati\allowbreak{}c,current\}/expose\_XV\allowbreak{}\_body.tex;\allowbreak{} folio\allowbreak{} 182;\allowbreak{} formula\allowbreak{} 664\allowbreak{}}\par

\par\medskip\noindent\hypertarget{app-E7II-XV-011.1}{}\textbf{E7II-XV-011.1}\quad printed folio 183.\par
\noindent Printed reading: {\ttfamily\footnotesize o\allowbreak{}}\par
\noindent Corrected reading: {\ttfamily\footnotesize 0\allowbreak{}}\par
\noindent Source: {\ttfamily\footnotesize ED-EN-FORMULA-G121;\allowbreak{} fr/sga7ii/\{diplomati\allowbreak{}c,current\}/expose\_XV\allowbreak{}\_body.tex;\allowbreak{} folio\allowbreak{} 183\allowbreak{}}\par

\par\medskip\noindent\hypertarget{app-E7II-XV-011.2}{}\textbf{E7II-XV-011.2}\quad printed folio 183.\par
\noindent Printed reading: {\ttfamily\footnotesize o\allowbreak{}}\par
\noindent Corrected reading: {\ttfamily\footnotesize 0\allowbreak{}}\par
\noindent Source: {\ttfamily\footnotesize ED-EN-FORMULA-G121;\allowbreak{} fr/sga7ii/\{diplomati\allowbreak{}c,current\}/expose\_XV\allowbreak{}\_body.tex;\allowbreak{} folio\allowbreak{} 183\allowbreak{}}\par

\par\medskip\noindent\hypertarget{app-E7II-XV-012.1}{}\textbf{E7II-XV-012.1}\quad printed folio 183.\par
\noindent Printed reading: \emph{(empty)}\par
\noindent Corrected reading: {\ttfamily\footnotesize \textbackslash{}bar\allowbreak{}}\par
\noindent Source: {\ttfamily\footnotesize fr/sga7ii/\{diplomati\allowbreak{}c,current\}/expose\_XV\allowbreak{}\_body.tex;\allowbreak{} folio\allowbreak{} 183;\allowbreak{} formula\allowbreak{} 706\allowbreak{}}\par

\par\medskip\noindent\hypertarget{app-E7II-XV-012.2}{}\textbf{E7II-XV-012.2}\quad printed folio 184.\par
\noindent Printed reading: \emph{(empty)}\par
\noindent Corrected reading: {\ttfamily\footnotesize .\allowbreak{}}\par
\noindent Source: {\ttfamily\footnotesize port/decisions/decis\allowbreak{}ions\_en7II\_r2\_002.js\allowbreak{}on\#C662;\allowbreak{} fixes/FIXLOG\_7II\_CUR\allowbreak{}RENT.jsonl\#C662;\allowbreak{} 7II:283\allowbreak{}}\par

\par\medskip\noindent\hypertarget{app-E7II-XV-013.1}{}\textbf{E7II-XV-013.1}\quad printed folio 184.\par
\noindent Printed reading: \emph{(empty)}\par
\noindent Corrected reading: {\ttfamily\footnotesize \textbackslash{}bar\allowbreak{}}\par
\noindent Source: {\ttfamily\footnotesize fr/sga7ii/\{diplomati\allowbreak{}c,current\}/expose\_XV\allowbreak{}\_body.tex;\allowbreak{} folio\allowbreak{} 184;\allowbreak{} formula\allowbreak{} 735\allowbreak{}}\par

\par\medskip\noindent\hypertarget{app-E7II-XV-013.2}{}\textbf{E7II-XV-013.2}\quad printed folio 184.\par
\noindent Printed reading: \emph{(empty)}\par
\noindent Corrected reading: {\ttfamily\footnotesize \textbackslash{}bigr)\allowbreak{}}\par
\noindent Source: {\ttfamily\footnotesize fr/sga7ii/\{diplomati\allowbreak{}c,current\}/expose\_XV\allowbreak{}\_body.tex;\allowbreak{} folio\allowbreak{} 184;\allowbreak{} port/decisions/decis\allowbreak{}ions\_en7II\_r2\_002.js\allowbreak{}on\#C663;\allowbreak{} fixes/FIXLOG\_7II\_CUR\allowbreak{}RENT.jsonl\#C663;\allowbreak{} 7II:285\allowbreak{}}\par

\par\medskip\noindent\hypertarget{app-E7II-XV-014.1}{}\textbf{E7II-XV-014.1}\quad printed folio 185.\par
\noindent Printed reading: \emph{(empty)}\par
\noindent Corrected reading: {\ttfamily\footnotesize ,\allowbreak{}}\par
\noindent Source: {\ttfamily\footnotesize fr/sga7ii/\{diplomati\allowbreak{}c,current\}/expose\_XV\allowbreak{}\_body.tex;\allowbreak{} folio\allowbreak{} 185;\allowbreak{} formula\allowbreak{} 759\allowbreak{}}\par

\par\medskip\noindent\hypertarget{app-E7II-XV-015.1}{}\textbf{E7II-XV-015.1}\quad printed folio 186.\par
\noindent Printed reading: \emph{(empty)}\par
\noindent Corrected reading: {\ttfamily\footnotesize ,\allowbreak{}}\par
\noindent Source: {\ttfamily\footnotesize fr/sga7ii/\{diplomati\allowbreak{}c,current\}/expose\_XV\allowbreak{}\_body.tex;\allowbreak{} folio\allowbreak{} 186;\allowbreak{} formula\allowbreak{} 784\allowbreak{}}\par

\par\medskip\noindent\hypertarget{app-E7II-XV-016.1}{}\textbf{E7II-XV-016.1}\quad printed folio 186.\par
\noindent Printed reading: \emph{(empty)}\par
\noindent Corrected reading: {\ttfamily\footnotesize \textbackslash{}widetilde\allowbreak{} }\par
\noindent Source: {\ttfamily\footnotesize ED-EN-FORMULA-G123;\allowbreak{} fr/sga7ii/\{diplomati\allowbreak{}c,current\}/expose\_XV\allowbreak{}\_body.tex;\allowbreak{} folio\allowbreak{} 186\allowbreak{}}\par

\par\medskip\noindent\hypertarget{app-E7II-XV-016.2}{}\textbf{E7II-XV-016.2}\quad printed folio 186.\par
\noindent Printed reading: \emph{(empty)}\par
\noindent Corrected reading: {\ttfamily\footnotesize \_0\allowbreak{}}\par
\noindent Source: {\ttfamily\footnotesize fr/sga7ii/\{diplomati\allowbreak{}c,current\}/expose\_XV\allowbreak{}\_body.tex;\allowbreak{} folio\allowbreak{} 186\allowbreak{}}\par

\par\medskip\noindent\hypertarget{app-E7II-XV-016.3}{}\textbf{E7II-XV-016.3}\quad printed folio 186.\par
\noindent Printed reading: \emph{(empty)}\par
\noindent Corrected reading: {\ttfamily\footnotesize \_0\allowbreak{}}\par
\noindent Source: {\ttfamily\footnotesize fr/sga7ii/\{diplomati\allowbreak{}c,current\}/expose\_XV\allowbreak{}\_body.tex;\allowbreak{} folio\allowbreak{} 186\allowbreak{}}\par

\par\medskip\noindent\hypertarget{app-E7II-XV-017.1}{}\textbf{E7II-XV-017.1}\quad printed folio 187.\par
\noindent Printed reading: {\ttfamily\footnotesize i\allowbreak{}}\par
\noindent Corrected reading: {\ttfamily\footnotesize n\allowbreak{}}\par
\noindent Source: {\ttfamily\footnotesize fr/sga7ii/\{diplomati\allowbreak{}c,current\}/expose\_XV\allowbreak{}\_body.tex;\allowbreak{} folio\allowbreak{} 187;\allowbreak{} formula\allowbreak{} 818\allowbreak{}}\par

\par\medskip\noindent\hypertarget{app-E7II-XV-018.1}{}\textbf{E7II-XV-018.1}\quad printed folio 188.\par
\noindent Printed reading: {\ttfamily\footnotesize alone\allowbreak{}}\par
\noindent Corrected reading: {\ttfamily\footnotesize zero\allowbreak{}}\par
\noindent Source: {\ttfamily\footnotesize ED-EN-NONMATH-G048;\allowbreak{} F7II-XV-023.1;\allowbreak{} 7II:290;\allowbreak{} fr/sga7ii/\{diplomati\allowbreak{}c,current\}/expose\_XV\allowbreak{}\_body.tex;\allowbreak{} folio\allowbreak{} 188\allowbreak{}}\par

\par\medskip\noindent\hypertarget{app-E7II-XV-018.2}{}\textbf{E7II-XV-018.2}\quad printed folio 188.\par
\noindent Printed reading: \emph{(empty)}\par
\noindent Corrected reading: {\ttfamily\footnotesize X\_\{\textbackslash{}bar\allowbreak{} s\},\allowbreak{}}\par
\noindent Source: {\ttfamily\footnotesize ED-EN-FORMULA-G125;\allowbreak{} fr/sga7ii/\{diplomati\allowbreak{}c,current\}/expose\_XV\allowbreak{}\_body.tex;\allowbreak{} folio\allowbreak{} 188\allowbreak{}}\par

\par\medskip\noindent\hypertarget{app-E7II-XV-018.3}{}\textbf{E7II-XV-018.3}\quad printed folio 188.\par
\noindent Printed reading: {\ttfamily\footnotesize \textasciicircum{}n\allowbreak{}}\par
\noindent Corrected reading: \emph{(empty)}\par
\noindent Source: {\ttfamily\footnotesize ED-EN-FORMULA-G125;\allowbreak{} fr/sga7ii/\{diplomati\allowbreak{}c,current\}/expose\_XV\allowbreak{}\_body.tex;\allowbreak{} folio\allowbreak{} 188\allowbreak{}}\par

\par\medskip\noindent\hypertarget{app-E7II-XV-018.4}{}\textbf{E7II-XV-018.4}\quad printed folio 188.\par
\noindent Printed reading: \emph{(empty)}\par
\noindent Corrected reading: {\ttfamily\footnotesize \{\textbackslash{}bar\allowbreak{}}\par
\noindent Source: {\ttfamily\footnotesize ED-EN-FORMULA-G125;\allowbreak{} fr/sga7ii/\{diplomati\allowbreak{}c,current\}/expose\_XV\allowbreak{}\_body.tex;\allowbreak{} folio\allowbreak{} 188\allowbreak{}}\par

\par\medskip\noindent\hypertarget{app-E7II-XV-018.5}{}\textbf{E7II-XV-018.5}\quad printed folio 188.\par
\noindent Printed reading: \emph{(empty)}\par
\noindent Corrected reading: {\ttfamily\footnotesize \}\allowbreak{}}\par
\noindent Source: {\ttfamily\footnotesize ED-EN-FORMULA-G125;\allowbreak{} fr/sga7ii/\{diplomati\allowbreak{}c,current\}/expose\_XV\allowbreak{}\_body.tex;\allowbreak{} folio\allowbreak{} 188\allowbreak{}}\par

\par\medskip\noindent\hypertarget{app-E7II-XV-019.1}{}\textbf{E7II-XV-019.1}\quad printed folio 188.\par
\noindent Printed reading: \emph{(empty)}\par
\noindent Corrected reading: {\ttfamily\footnotesize X\_s,\allowbreak{}}\par
\noindent Source: {\ttfamily\footnotesize fr/sga7ii/\{diplomati\allowbreak{}c,current\}/expose\_XV\allowbreak{}\_body.tex;\allowbreak{} folio\allowbreak{} 188\allowbreak{}}\par

\par\medskip\noindent\hypertarget{app-E7II-XV-019.2}{}\textbf{E7II-XV-019.2}\quad printed folio 188.\par
\noindent Printed reading: {\ttfamily\footnotesize \textasciicircum{}n\allowbreak{}}\par
\noindent Corrected reading: \emph{(empty)}\par
\noindent Source: {\ttfamily\footnotesize fr/sga7ii/\{diplomati\allowbreak{}c,current\}/expose\_XV\allowbreak{}\_body.tex;\allowbreak{} folio\allowbreak{} 188\allowbreak{}}\par

\par\medskip\noindent\hypertarget{app-E7II-XV-019.3}{}\textbf{E7II-XV-019.3}\quad printed folio 188.\par
\noindent Printed reading: \emph{(empty)}\par
\noindent Corrected reading: {\ttfamily\footnotesize \{\textbackslash{}bar\allowbreak{}}\par
\noindent Source: {\ttfamily\footnotesize fr/sga7ii/\{diplomati\allowbreak{}c,current\}/expose\_XV\allowbreak{}\_body.tex;\allowbreak{} folio\allowbreak{} 188\allowbreak{}}\par

\par\medskip\noindent\hypertarget{app-E7II-XV-019.4}{}\textbf{E7II-XV-019.4}\quad printed folio 188.\par
\noindent Printed reading: \emph{(empty)}\par
\noindent Corrected reading: {\ttfamily\footnotesize \}\allowbreak{}}\par
\noindent Source: {\ttfamily\footnotesize fr/sga7ii/\{diplomati\allowbreak{}c,current\}/expose\_XV\allowbreak{}\_body.tex;\allowbreak{} folio\allowbreak{} 188\allowbreak{}}\par

\par\medskip\noindent\hypertarget{app-E7II-XV-019.5}{}\textbf{E7II-XV-019.5}\quad printed folio 188.\par
\noindent Printed reading: \emph{(empty)}\par
\noindent Corrected reading: {\ttfamily\footnotesize )\allowbreak{}}\par
\noindent Source: {\ttfamily\footnotesize fr/sga7ii/\{diplomati\allowbreak{}c,current\}/expose\_XV\allowbreak{}\_body.tex;\allowbreak{} folio\allowbreak{} 188\allowbreak{}}\par

\par\medskip\noindent\hypertarget{app-E7II-XV-019.6}{}\textbf{E7II-XV-019.6}\quad printed folio 189.\par
\noindent Printed reading: \emph{(empty)}\par
\noindent Corrected reading: {\ttfamily\footnotesize ,\allowbreak{}}\par
\noindent Source: {\ttfamily\footnotesize fr/sga7ii/\{diplomati\allowbreak{}c,current\}/expose\_XV\allowbreak{}\_body.tex;\allowbreak{} folio\allowbreak{} 189;\allowbreak{} formula\allowbreak{} 866\allowbreak{}}\par

\par\medskip\noindent\hypertarget{app-E7II-XV-020.1}{}\textbf{E7II-XV-020.1}\quad printed folio 191.\par
\noindent Printed reading: {\ttfamily\footnotesize \textbackslash{}widetilde\allowbreak{} }\par
\noindent Corrected reading: \emph{(empty)}\par
\noindent Source: {\ttfamily\footnotesize fr/sga7ii/\{diplomati\allowbreak{}c,current\}/expose\_XV\allowbreak{}\_body.tex;\allowbreak{} folio\allowbreak{} 191;\allowbreak{} formula\allowbreak{} 959\allowbreak{}}\par

\par\medskip\noindent\hypertarget{app-E7II-XV-020.2}{}\textbf{E7II-XV-020.2}\quad printed folio 191.\par
\noindent Printed reading: {\ttfamily\footnotesize \textbackslash{}widetilde\allowbreak{} }\par
\noindent Corrected reading: \emph{(empty)}\par
\noindent Source: {\ttfamily\footnotesize fr/sga7ii/\{diplomati\allowbreak{}c,current\}/expose\_XV\allowbreak{}\_body.tex;\allowbreak{} folio\allowbreak{} 191;\allowbreak{} formula\allowbreak{} 963\allowbreak{}}\par

\par\medskip\noindent\hypertarget{app-E7II-XV-021.1}{}\textbf{E7II-XV-021.1}\quad printed folio 193.\par
\noindent Printed reading: \emph{(empty)}\par
\noindent Corrected reading: {\ttfamily\footnotesize \_0\allowbreak{}}\par
\noindent Source: {\ttfamily\footnotesize ED-EN-FORMULA-G127;\allowbreak{} fr/sga7ii/\{diplomati\allowbreak{}c,current\}/expose\_XV\allowbreak{}\_body.tex;\allowbreak{} folio\allowbreak{} 193\allowbreak{}}\par

\par\medskip\noindent\hypertarget{app-E7II-XV-022.1}{}\textbf{E7II-XV-022.1}\quad printed folio 195.\par
\noindent Printed reading: {\ttfamily\footnotesize \_\allowbreak{}}\par
\noindent Corrected reading: {\ttfamily\footnotesize \textasciicircum{}\allowbreak{}}\par
\noindent Source: {\ttfamily\footnotesize fr/sga7ii/\{diplomati\allowbreak{}c,current\}/expose\_XV\allowbreak{}\_body.tex;\allowbreak{} folio\allowbreak{} 195;\allowbreak{} formula\allowbreak{} 1086;\allowbreak{} port/PORT\_EN7II\_LOG.\allowbreak{}jsonl\#C673;\allowbreak{} fixes/FIXLOG\_7II\_CUR\allowbreak{}RENT.jsonl\#C673;\allowbreak{} 7II:301\allowbreak{}}\par

\par\medskip\noindent\hypertarget{app-E7II-XV-022.2}{}\textbf{E7II-XV-022.2}\quad printed folio 196.\par
\noindent Printed reading: \emph{(empty)}\par
\noindent Corrected reading: {\ttfamily\footnotesize \textbackslash{}bar\allowbreak{}}\par
\noindent Source: {\ttfamily\footnotesize fr/sga7ii/\{diplomati\allowbreak{}c,current\}/expose\_XV\allowbreak{}\_body.tex;\allowbreak{} folio\allowbreak{} 196;\allowbreak{} formula\allowbreak{} 1093\allowbreak{}}\par

\par\medskip\noindent\hypertarget{app-E7II-XV-022.3}{}\textbf{E7II-XV-022.3}\quad printed folio 196.\par
\noindent Printed reading: \emph{(empty)}\par
\noindent Corrected reading: {\ttfamily\footnotesize ,\allowbreak{}}\par
\noindent Source: {\ttfamily\footnotesize fr/sga7ii/\{diplomati\allowbreak{}c,current\}/expose\_XV\allowbreak{}\_body.tex;\allowbreak{} folio\allowbreak{} 196;\allowbreak{} formula\allowbreak{} 1102\allowbreak{}}\par

\par\medskip\noindent\hypertarget{app-E7II-XV-022.4}{}\textbf{E7II-XV-022.4}\quad printed folio 196.\par
\noindent Printed reading: \emph{(empty)}\par
\noindent Corrected reading: {\ttfamily\footnotesize ,\allowbreak{}}\par
\noindent Source: {\ttfamily\footnotesize fr/sga7ii/\{diplomati\allowbreak{}c,current\}/expose\_XV\allowbreak{}\_body.tex;\allowbreak{} folio\allowbreak{} 196;\allowbreak{} formula\allowbreak{} 1103\allowbreak{}}\par

\par\endgroup\normalsize\fi
\clearpage\SGAtoc{Exposé XVI}
\selectlanguage{english}
\clearpage
\SGApage{II}{197}

\begin{center}
{\large SGA 7 --- EXPOS\'E XVI}\par
\vspace{2em}
{\Large\underline{\underline{The Milnor Formula}}}\par
\vspace{1.5em}
{\large by P. Deligne}
\end{center}

\medskip
In this expos\'e, Milnor's formula [3] (see also [2]) is transposed
to characteristic $p$.  In certain cases this formula makes it
possible to compute the dimensions of spaces of vanishing cycles.
The proof is global in nature and treats only the ``geometric'' case
in which the base is a henselian trait of equal characteristic with
perfect residue field.

\section*{1. Statement of the problem}

\subsubsection*{1.1.}
Let $S$ be a regular scheme, purely of dimension one, let $s$ be a
closed point of $S$ \underline{with perfect residue field}, let
$f:X\to S$ be a flat morphism of finite type, let $x_0$ be a closed
point of the fiber $X_s$, and let $\bar x_0$ be a geometric point of
$X$ localized at $x_0$.  Suppose that
\begin{enumerate}
\item[(a)] $x_0$ is an \underline{isolated point at which $f$ is not
smooth}: there are an integer $n$ and a neighborhood $U$ of $x_0$,
purely of relative dimension $n$, such that $U-\{x_0\}$ is smooth
over $S$;
\item[(b)] $X$ is \underline{regular} at $x_0$.
\end{enumerate}

Under these assumptions, $X$ is a relative complete intersection
over $S$ at $x_0$.  We refer to VI, 3.10 for the definition of
$D^1_{X/S}$.  On the neighborhood $U$ one has
\[
D^1_{X/S}
=\mathcal H^1\bigl(R\mathcal Hom(L\Omega_{X/S},\mathcal O_X)\bigr)
=\mathcal Ext^1(\Omega^1_{X/S},\mathcal O_X);
\]
on this neighborhood, $D^1_{X/S}$ has finite length and support
contained in $\{x_0\}$.

\SGApage{II}{198}
\subsubsection*{Definition 1.2.}
\underline{The Milnor number
$\mu=\mu(X/S,\bar x_0)$ of $X/S$ at $\bar x_0$ is the length of the
$\mathcal O_{X,x_0}$-module defined by $D^1_{X/S}$.}

\subsubsection*{Example 1.3.}
Take $S$ to be a smooth curve over an algebraically closed field
$k$.  After replacing $X$ by a neighborhood of $x_0$, we may suppose
that $X$ is smooth over $k$.  The cotangent complex then reduces to
the following resolution of $\Omega^1_{X/S}$:
\[
L\Omega_{X/S}
=\bigl[f^*\Omega^1_{S/k}\longrightarrow\Omega^1_{X/k}\bigr].
\]
Choose a local coordinate $t:S\to\mathbb E^1_k$ on $S$ at $s$
(with $t$ \'etale at $s$), again denote by $f$ the composite
$t\circ f$, and let $(x_i)_{1\le i\le n+1}$ be a system of local
coordinates on $X$ at $x_0$.  At $x_0$, the module
$f^*\Omega^1_{S/k}$ then has basis $df$, while $\Omega^1_{X/k}$ has
basis the $dx_i$.  Thus $D^1_{X/S}$ is identified with the quotient
of $\mathcal O_X$ by the ideal generated by the
$\partial f/\partial x_i$, and
\begin{equation}\tag{1.3.1}
\mu=\dim_k\!\left(
\mathcal O_{X,x_0}/
\left(\frac{\partial f}{\partial x_1},\ldots,
      \frac{\partial f}{\partial x_{n+1}}\right)\right).
\end{equation}
This is the definition given by Milnor.

\subsubsection*{Example 1.4.}
Suppose that $X$ is the closed subscheme of $\mathbb E^{n+1}_S$
defined by an equation $g$, and let
$i:X\hookrightarrow\mathbb E^{n+1}_S$ be the inclusion.  The
cotangent complex then reduces to
\[
\bigl[(g)/(g)^2\longrightarrow
i^*\Omega^1_{\mathbb E^{n+1}/S}\bigr],
\]
and $D^1_{X/S}$ is isomorphic to the quotient of $\mathcal O_X$ by
the ideal generated by
$(\partial g/\partial x_i)_{1\le i\le n+1}$.  Equivalently, it is
the quotient of $\mathcal O_{\mathbb E^{n+1}}$ by the ideal generated
by $g$ and the $\partial g/\partial x_i$.

\subsubsection*{Example 1.5.}
Let $S$ be the spectrum of a discrete valuation ring $V$ with
perfect residue field and fraction field $K$, and let $X$ be the
normalization of $S$ in a separable extension $K'$ of $K$.  By 1.4
and [5], III, 7.14, p.~68, $D^1_{X/S}$ is then isomorphic to the
quotient of $\mathcal O_X$ by the different $\mathfrak d_{X/S}$.  If
$\bar x_0$ is a geometric point of $X_s$, localized at
$x_0\in X_s$, its Milnor number is therefore
\begin{equation}\tag{1.5.1}
\mu=\operatorname{length}
\bigl(\mathcal O_{X,x_0}/\mathfrak d_{X/S}\bigr).
\end{equation}
Let $\bar s$ be a geometric point of $S$ localized at the closed
point $s$, and let $d_{K'/K}\in V$ generate the discriminant ideal,
with valuation $v(d_{K'/K})$.  It follows from (1.5.1) that
\begin{equation}\tag{1.5.2}
v(d_{K'/K})=\sum_{x\in X(\bar s)}\mu(X/S,x).
\end{equation}


\SGApage{II}{\(\approx\)199}
\subsubsection*{1.6.}
Retain the assumptions of 1.1.  Let $\bar s$ be the geometric point
of $S$ which is the image of $\bar x_0$, let $S_{(\bar s)}$ be the
spectrum of the strict henselization of $\mathcal O_{S,s}$ at
$\bar s$, and let $X_{(\bar s)}$ be the inverse image of $X$ over
$S_{(\bar s)}$.  Let $\bar\eta$ be the spectrum of an algebraic
closure of the fraction field $k(\eta)$ of $S_{(\bar s)}$, and put
\[
I=\operatorname{Gal}(k(\bar\eta)/k(\eta)).
\]

Let $\ell$ be a prime number different from the residual
characteristic of $S$.  In a neighborhood of $\bar x_0$ in
$X_{\bar s}$, the sheaves of vanishing cycles
\[
\Phi^i_{\bar\eta}(\mathbb Z/\ell)
=R^i\Phi_{\bar\eta}(\mathbb Z/\ell)
\]
relative to the constant sheaf $\mathbb Z/\ell$ on
$X_{(\bar s)}$ are concentrated at $\bar x_0$.  Moreover, the
inertia group $I$ acts on the $\mathbb Z/\ell$-vector spaces
\[
\Phi^i_{\bar\eta}(\mathbb Z/\ell)_{\bar x_0}.
\]
These spaces are finite-dimensional (XIII, 2.4.2).

\subsubsection*{1.7.}
Let $I_1$ be a finite quotient of the inertia group.  Recall that
the \underline{Swan character} $\operatorname{Sw}$ of $I_1$ is the
difference $a-u$ between the Artin character $a$ and the character
$u$ of the augmentation representation.  Recall also that, up to
isomorphism, there is a unique projective $\mathbb Z_\ell[I_1]$-module
$S(I_1)$ with character $\operatorname{Sw}$ (see [6]).

\SGApage{II}{\(\approx\)200}
\SGAvar{Let $V$ be a finite-dimensional $\mathbb Z/\ell$-vector space with
an action $\rho:I\to\operatorname{GL}(V)$, and suppose that $\rho$
factors through $I_1$.  The \underline{wild dimension} of $V$ is
then
\begin{equation}\tag{1.7.1}
\dim\operatorname{Sw}(V)
=\dim_{\mathbb Z/\ell}
\operatorname{Hom}_{I_1}\!\left(
S(I_1)\otimes_{\mathbb Z/\ell}\mathbb Z/\ell,,V\right).
\end{equation}
This integer is independent of the choice of $I_1$, since for a
normal subgroup $H$ of $I_1$ one has
$S(I_1/H)\sim S(I_1)_H$.

}{Let $V$ be a finite-dimensional $\mathbb Z/\ell$-vector space with
an action $\rho:I\to\operatorname{GL}(V)$, and suppose that $\rho$
factors through $I_1$.  The \underline{wild dimension} of $V$ is
then
\begin{equation}\tag{1.7.1}
\dim\operatorname{Sw}(V)
=\dim_{\mathbb Z/\ell}
\operatorname{Hom}_{I_1}\!\left(
S(I_1)\otimes_{\mathbb Z_\ell}\mathbb Z/\ell,,V\right).
\end{equation}
This integer is independent of the choice of $I_1$, since for a
normal subgroup $H$ of $I_1$ one has
$S(I_1/H)\sim S(I_1)_H$.

}{E7II-XVI-001}We further put
\[
\dim\operatorname{tot}(V)
=\dim\operatorname{Sw}(V)+\dim_{\mathbb Z/\ell}(V)
\]
and call this number the \underline{total dimension} of $V$.

\subsubsection*{Definition 1.8.}
\underline{With the terminology of 1.6 and 1.7, the number of
vanishing cycles of $X/S$ at $\bar x_0$ (respectively the wild
number of vanishing cycles, respectively the total number of
vanishing cycles) is the integer}
\[
(-1)^n\dim\Phi_{\bar x_0}
\underset{\mathrm{def}}{=}
(-1)^n\sum_i(-1)^i
\dim\Phi^i_{\bar\eta}(\mathbb Z/\ell)_{\bar x_0},
\]
\[
(-1)^n\dim\operatorname{Sw}\Phi_{\bar x_0}
\underset{\mathrm{def}}{=}
(-1)^n\sum_i(-1)^i
\dim\operatorname{Sw}
\Phi^i_{\bar\eta}(\mathbb Z/\ell)_{\bar x_0},
\]
\[
\begin{aligned}
(-1)^n\dim\operatorname{tot}\Phi_{\bar x_0}
&\underset{\mathrm{def}}{=}
(-1)^n\sum_i(-1)^i\\[-2pt]
&\qquad{}\cdot\dim\operatorname{tot}
\Phi^i_{\bar\eta}(\mathbb Z/\ell)_{\bar x_0}\\
&=(-1)^n\bigl(
\dim\Phi_{\bar x_0}+\dim\operatorname{Sw}\Phi_{\bar x_0}
\bigr).
\end{aligned}
\]

\subsubsection*{Conjecture 1.9.}
\underline{Under the assumptions of 1.1, the Milnor number $\mu$ of
$X/S$ at $\bar x_0$ is equal to the total number of vanishing cycles
of $X/S$ at $\bar x_0$.}

\subsubsection*{Remark 1.10.}
Let $S$ be a henselian trait with imperfect residue field and let
$f:X\to S$ be a morphism which is smooth except at
$x_0\in X_s$.  Suppose that $X$ is regular at $x_0$.  Definition
1.2 of the Milnor number $\mu$ still makes sense under these
assumptions, but I no longer have a geometric interpretation of
$\mu$ to propose.


\SGApage{II}{\(\approx\)201}
\subsubsection*{1.11.}
For the remainder of this section, suppose that
$S=\operatorname{Spec}(V)$ is a henselian trait with
\underline{algebraically closed} residue field and closed point $s$,
that $f:X\to S$ is a flat morphism of finite type, purely of relative
dimension $n$, that $X$ is regular, that $x_0$ is a closed point of
$X_s$, and that $X-\{x_0\}$ is smooth over $S$.

This entails no loss of generality.

\subsubsection*{Proposition 1.12.}
\underline{If $n=0$ (that is, if $X/S$ is finite), Conjecture 1.9
holds for $(X/S,x_0)$.}

We may suppose that $X$ is the normalization of $S$ in a finite
extension $K'$ of the fraction field $K$.  By 1.5,
$\mu=v(d_{K'/K})$.

\SGAvar{The extension $K'$ is defined by an open subgroup $H$ of the inertia
group $I$.  Let $H'\subset H$ be an open normal subgroup of $I$, and
put $I'=I/H'$.  Let $s_{I/H}$ be the character of the permutation
representation of $I'$ on $I/H$, let $a_{I}$ be the Artin
character, let $\operatorname{Sw}_{I}$ be the Swan character, and
let $u_{I'}$ be the character of the augmentation representation.
By Artin's ``F\"uhrerdiskriminantenproduktformel,'' the discriminant
ideal is the Artin conductor of $s_{I/H}$ ([5], VI, 3, Corollary 1,
p.~111):
\[
\mu=(a_{I'},s_{I/H})
=|I'|^{-1}\sum_{i\in I'}a_{I'}(i)s_{I/H}(i).
\]
One has
\[
a_{I'}=u_{I'}+\operatorname{Sw}_{I'},
\qquad
(u_{I'},s_{I/H})=\#(I/H)-1,
\]
and
\begin{align*}
(\operatorname{Sw}_{I'},s_{I/H})
&=\dim_{\mathbb Z/\ell}\operatorname{Hom}_I\!\left(
S(I')\otimes_{\mathbb Z_\ell}\mathbb Z/\ell,
(\mathbb Z/\ell)^{I/H}\right)\\
&=\dim_{\mathbb Z/\ell}\operatorname{Hom}_I\!\left(
S(I')\otimes_{\mathbb Z_\ell}\mathbb Z/\ell,
\operatorname{Coker}\bigl(
\mathbb Z/\ell\longrightarrow(\mathbb Z/\ell)^{I/H}\bigr)
\right).
\end{align*}
Consequently,
\[
\mu=\#(I/H)-1+
\dim\operatorname{Sw}\operatorname{Coker}\bigl(
\mathbb Z/\ell\longrightarrow(\mathbb Z/\ell)^{I/H}\bigr).
\]
One checks immediately that $\Phi^i_{x_0}=0$ for $i\ne0$ and that
\[
\Phi^0_{x_0}=\operatorname{Coker}\bigl(
\mathbb Z/\ell\longrightarrow(\mathbb Z/\ell)^{I/H}\bigr);
\]
the formula follows.


}{The extension $K'$ is defined by an open subgroup $H$ of the inertia
group $I$.  Let $H'\subset H$ be an open normal subgroup of $I$, and
put $I'=I/H'$.  Let $s_{I/H}$ be the character of the permutation
representation of $I'$ on $I/H$, let $a_{I'}$ be the Artin
character, let $\operatorname{Sw}_{I'}$ be the Swan character, and
let $u_{I'}$ be the character of the augmentation representation.
By Artin's ``F\"uhrerdiskriminantenproduktformel,'' the discriminant
ideal is the Artin conductor of $s_{I/H}$ ([5], VI, 3, Corollary 1,
p.~111):
\[
\mu=(a_{I'},s_{I/H})
=|I'|^{-1}\sum_{i\in I'}a_{I'}(i)s_{I/H}(i).
\]
One has
\[
a_{I'}=u_{I'}+\operatorname{Sw}_{I'},
\qquad
(u_{I'},s_{I/H})=\#(I/H)-1,
\]
and
\begin{align*}
(\operatorname{Sw}_{I'},s_{I/H})
&=\dim_{\mathbb Z/\ell}\operatorname{Hom}_I\!\left(
S(I')\otimes_{\mathbb Z_\ell}\mathbb Z/\ell,
(\mathbb Z/\ell)^{I/H}\right)\\
&=\dim_{\mathbb Z/\ell}\operatorname{Hom}_I\!\left(
S(I')\otimes_{\mathbb Z_\ell}\mathbb Z/\ell,
\operatorname{Coker}\bigl(
\mathbb Z/\ell\longrightarrow(\mathbb Z/\ell)^{I/H}\bigr)
\right).
\end{align*}
Consequently,
\[
\mu=\#(I/H)-1+
\dim\operatorname{Sw}\operatorname{Coker}\bigl(
\mathbb Z/\ell\longrightarrow(\mathbb Z/\ell)^{I/H}\bigr).
\]
One checks immediately that $\Phi^i_{x_0}=0$ for $i\ne0$ and that
\[
\Phi^0_{x_0}=\operatorname{Coker}\bigl(
\mathbb Z/\ell\longrightarrow(\mathbb Z/\ell)^{I/H}\bigr);
\]
the formula follows.


}{E7II-XVI-002}\SGApage{II}{202}
\subsubsection*{Proposition 1.13.}
\underline{If $X_s$ has an ordinary quadratic singularity at $x_0$,
Conjecture 1.9 holds for $(X/S,x_0)$.}

\medskip
\noindent\underline{Case 1: $n=\dim(X_s)=2m-1$.}

\SGAvar{The problem is local for the \'etale topology on $X$.  We may
therefore suppose that $X_s$ is the closed subscheme of
$E_S^{n+1}$ defined by an equation
\[
f=\sum_{i=1}^{m}x_i x_{i+m}+a=0,
\]
where $a$ lies in the maximal ideal of $V$ and $x_0$ is the origin
of the coordinates.  Moreover, $X$ is regular at $x_0$ if and only
if $a$ is a uniformizer.  By 1.4,
\[
\mu=\operatorname{length}\!\left(
\mathcal O_{\mathbb E_S^{2m},x_0}/
\left(f,\frac{\partial f}{\partial x_i}\right)\right)=1.
\]
By the ordinary-quadratic vanishing-cycle calculation of Expos\'e XV,
$\Phi^i_{x_0}=0$ for $i\ne n$, while $\Phi^n_{x_0}$ has dimension
one and inertia acts trivially on it.
The wild number of vanishing cycles is therefore zero and the total
number is $1$, which verifies 1.9.

}{The problem is local for the \'etale topology on $X$.  We may
therefore suppose that $X$ is the closed subscheme of
$\mathbb E_S^{n+1}$ defined by an equation
\[
f=\sum_{i=1}^{m}x_i x_{i+m}+a=0,
\]
where $a$ lies in the maximal ideal of $V$ and $x_0$ is the origin
of the coordinates.  Moreover, $X$ is regular at $x_0$ if and only
if $a$ is a uniformizer.  By 1.4,
\[
\mu=\operatorname{length}\!\left(
\mathcal O_{\mathbb E_S^{2m},x_0}/
\left(f,\frac{\partial f}{\partial x_i}\right)\right)=1.
\]
By the ordinary-quadratic vanishing-cycle calculation of Expos\'e XV,
$\Phi^i_{x_0}=0$ for $i\ne n$, while $\Phi^n_{x_0}$ has dimension
one and inertia acts trivially on it.
The wild number of vanishing cycles is therefore zero and the total
number is $1$, which verifies 1.9.

}{E7II-XVI-003}\SGApage{II}{\(\approx\)203}
\medskip
\noindent\underline{Case 2: $n=\dim(X_s)=2m$ and
$\operatorname{char}(k(s))\ne2$.}

\SGAvar{Under these assumptions we may take $X\subset\mathbb E_S^{n+1}$ to
be defined by
\[
f=\sum_{i=1}^{m}x_i x_{i+m}+x_{2m+1}^2+a=0,
\]
where $a$ is a uniformizer.  Again $\mu=1$,
$\Phi^i_{x_0}=0$ for $i\ne n$, and $\Phi^i_{x_0}$ has dimension
one.  This time the inertia group acts by $\pm1$ on
$\Phi^n_{x_0}$, hence tamely.  The wild number of vanishing cycles is
again zero, and 1.9 is verified.

}{Under these assumptions we may take $X\subset\mathbb E_S^{n+1}$ to
be defined by
\[
f=\sum_{i=1}^{m}x_i x_{i+m}+x_{2m+1}^2+a=0,
\]
where $a$ is a uniformizer.  Again $\mu=1$,
$\Phi^i_{x_0}=0$ for $i\ne n$, and $\Phi^n_{x_0}$ has dimension
one.  This time the inertia group acts by $\pm1$ on
$\Phi^n_{x_0}$, hence tamely.  The wild number of vanishing cycles is
again zero, and 1.9 is verified.

}{E7II-XVI-004}\medskip
\noindent\underline{Case 3: $n=\dim(X_s)=2m$ and
$\operatorname{char}(k(s))=2$.}

This time we may take $X\subset\mathbb E_S^{n+1}$ to be defined by
\[
f=\sum_{i=1}^{m}x_i x_{i+m}
 +(x_{2m+1}^2+a x_{2m+1}+b)=0,
\]
where $a$ and $b$ lie in the maximal ideal.  The scheme $X$ is
regular if and only if $b$ is a uniformizer, and its generic fiber
is smooth if and only if $a^2-4b\ne0$.

The Milnor number of $X/S$ is the same as that of
$X_0\subset\mathbb E_S^1$ defined by
\begin{equation}\tag{1}
x^2+a x+b=0.
\end{equation}
Likewise, $\Phi^i_{x_0}=0$ for $i\ne n$,
$\Phi^n_{x_0}$ has dimension one, and $I$ acts on it by $\pm1$; the
index-two subgroup of $I$ acting trivially corresponds to the
extension of $K$ defined by (1).  The total number of vanishing
cycles of $X/S$ at $x_0$ therefore coincides with the one for
$X_0/S$, and 1.9 in this case follows from 1.12.  The Milnor number
is given by the discriminant formula of 1.5:
\begin{equation}\tag{1.13.1}
\mu=v(a^2-4b).
\end{equation}


\SGApage{II}{204}
\section*{2. The ``geometric'' case}

We shall prove that Conjecture 1.9 holds in equal characteristic.
The key geometric result is the following.  Its possible use was
suggested to me by conversations with B. Saint-Donat.

\subsubsection*{Proposition 2.1.}
\underline{Let} $S$ \underline{be a smooth connected projective curve over an
algebraically closed field} $k$, \underline{let} $X$ \underline{be a smooth proper variety
over} $k$, \underline{purely of dimension} $n+1$, \underline{and let} $f:X\to S$ \underline{be a
morphism which is smooth away from a finite set} $F$ \underline{of points of}
$X$.  \underline{Then the Euler--Poincar\'e characteristics in \'etale
cohomology of} $S$, $X$, \underline{and the geometric generic fiber}
$X_{\bar\eta}$ \underline{of} $f$ \underline{satisfy}
\begin{equation}\tag{2.1.1}
\chi(X)=\chi(S)\chi(X_{\bar\eta})
-(-1)^n\sum_{x\in F}\mu(X/S,x).
\end{equation}

The Euler--Poincar\'e characteristic of $X$ is the trace (or integral
over $X$) of the top Chern class of the tangent bundle:
\[
\chi(X)=\int c_{n+1}(\Theta_X)
=(-1)^{n+1}\int c_{n+1}(\Omega_X^1).
\]
Write
\[
\Omega_X^1=\left(\Omega_X^1\otimes
(f^*\Omega_S^1)^{\otimes-1}\right)\otimes f^*\Omega_S^1.
\]
The total Chern class of the invertible sheaf $f^*\Omega_S^1$ is
the inverse image of the total Chern class of $\Omega_S^1$ on $S$:
\[
c^*(f^*\Omega_S^1)=1+f^*c^1(S,\Omega_S^1),
\]
and the cup powers of $c^1$ vanish:
\[
\bigl(f^*c^1(S,\Omega_S^1)\bigr)^n
=f^*\bigl(c^1(S,\Omega_S^1)^n\bigr)=0
\qquad(n>1).
\]
It follows that
\begin{align*}
c^{n+1}(\Omega_X^1)
={}&c^n\!\left(\Omega_X^1\otimes
(f^*\Omega_S^1)^{\otimes-1}\right)
\smile f^*c_1(S,\Omega_S^1)\\
&+c^{n+1}\!\left(\Omega_X^1\otimes
(f^*\Omega_S^1)^{\otimes-1}\right).
\end{align*}

\SGApage{II}{\(\approx\)205}
If $S$ has genus $g$, then $c_1(S,\Omega_S^1)$ is $(2g-2)$
times the cohomology class of any closed point $s$ of $S$.
Consequently, when $X_s$ is nonsingular,
\begin{align*}
\mathrm I
&=\int_Xc^n\!\left(\Omega_X^1\otimes
(f^*\Omega_S^1)^{\otimes-1}\right)
\smile f^*c_1(S,\Omega_S^1)\\
&=(2g-2)\int_Xc^n\!\left(\Omega_X^1\otimes
(f^*\Omega_S^1)^{\otimes-1}\right)
\smile f^*\operatorname{cl}(s)\\
&=(2g-2)\int_{X_s}c^n\!\left(
X_s,\Omega_X^1\otimes(f^*\Omega_S^1)^{\otimes-1}
\otimes\mathcal O_{X_s}\right).
\end{align*}
Moreover, on $X_s$ the sheaf $f^*\Omega_S^1$ is isomorphic to
$\mathcal O_{X_s}$, and
$\Omega_X^1\otimes_{\mathcal O_X}\mathcal O_{X_s}$ occurs in the
exact sequence
\[
0\longrightarrow f^*\Omega_S^1\otimes\mathcal O_{X_s}
\longrightarrow\Omega_X^1\otimes\mathcal O_{X_s}
\longrightarrow\Omega_{X_s}^1\longrightarrow0.
\]
Hence
\[
c^n\!\left(X_s,\Omega_X^1\otimes
(f^*\Omega_S^1)^{\otimes-1}\otimes\mathcal O_{X_s}\right)
=c^n(X_s,\Omega_{X_s}^1),
\]
and
\[
\mathrm I=(2g-2)\int_{X_s}c^n(X_s,\Omega_{X_s}^1)
=(-1)^{n+1}\chi(S)\chi(X_s).
\]

\SGApage{II}{\(\approx\)206}
The canonical map
\[
f^*:f^*\Omega_S^1\longrightarrow\Omega_X^1
\]
is identified with a section $a$ of
$\Omega_X^1\otimes(f^*\Omega_S^1)^{\otimes-1}$.  After choosing
systems of local coordinates on $S$ and $X$, thereby identifying
this bundle with $\mathcal O_X^{n+1}$, the coordinates of $a$ are
the partial derivatives $\partial f/\partial x_i$.  Thus $a$
vanishes only on $F$, and by 1.3 the multiplicity of its zero at
$x\in F$ is the Milnor number $\mu(X/S,x)$.

Recall that if a section $a$ of a vector bundle $G$ of rank $d$ on
a smooth proper $k$-scheme $X$, purely of dimension $d$, vanishes
only at isolated points, then $\int_Xc^d(G)$ is the number of zeros
of $a$, counted with multiplicity.  Here
\[
\mathrm{II}=\int_Xc^{n+1}\!\left(
\Omega_X^1\otimes(f^*\Omega_S^1)^{\otimes-1}\right)
=\sum_{x\in F}\mu(X/S,x).
\]
Collecting terms gives
\[
\chi(X)=(-1)^{n+1}(\mathrm I+\mathrm{II})
=\chi(S)\chi(X_s)-(-1)^n\sum_{x\in F}\mu(X/S,x),
\]
which proves 2.1, since
$\chi(X_s)=\chi(X_{\bar\eta})$.


\subsubsection*{Corollary 2.2.}
\underline{Under the preceding assumptions, if $n(X/S,x)$ is the
total number of vanishing cycles of $X/S$ at $x$ (1.8), then}
\[
\sum_{x\in F}\mu(X/S,x)=\sum_{x\in F}n(X/S,x).
\]

\SGAvar{The Leray spectral sequence for $f$ gives
\begin{align*}
\chi(X)
&=\sum_{p,q}(-1)^{p+q}
\dim_{\mathbb Z/\ell}H^p(S,R^qf_*\mathbb Z/\ell)\\
&=\sum_q(-1)^q\chi(S,R^qf_*\mathbb Z/\ell).
\end{align*}
For every constructible sheaf $G$ of $\mathbb Z/\ell$-modules on
$S$, one has ([4] or SGA 5, X, 7.2)
\[
\chi(S,G)=\chi(S)\dim_{\mathbb Z/\ell}(G_{\bar\eta})
-\sum_{s\in S}\varepsilon_s(G),
\]
where, for each $s\in S$, after identifying $\bar\eta$ with the
algebraic closure of the fraction field of the henselization of $S$
at $s$,
\[
\varepsilon_s(G)=
\dim(G_{\bar\eta})-\dim(G_s)-\dim\operatorname{Sw}(G_{\bar\eta}).
\]


}{The Leray spectral sequence for $f$ gives
\begin{align*}
\chi(X)
&=\sum_{p,q}(-1)^{p+q}
\dim_{\mathbb Z/\ell}H^p(S,R^qf_*\mathbb Z/\ell)\\
&=\sum_q(-1)^q\chi(S,R^qf_*\mathbb Z/\ell).
\end{align*}
For every constructible sheaf $G$ of $\mathbb Z/\ell$-modules on
$S$, one has ([4] or SGA 5, X, 7.2)
\[
\chi(S,G)=\chi(S)\dim_{\mathbb Z/\ell}(G_{\bar\eta})
-\sum_{s\in S}\varepsilon_s(G),
\]
where, for each $s\in S$, after identifying $\bar\eta$ with the
algebraic closure of the fraction field of the henselization of $S$
at $s$,
\[
\varepsilon_s(G)=
\dim(G_{\bar\eta})-\dim(G_s)+\dim\operatorname{Sw}(G_{\bar\eta}).
\]


}{E7II-XVI-005}\SGApage{II}{\(\approx\)207}
For $s\in S$, the exact sequences XIII, 2.4.6.3,
\[
\cdots\longrightarrow(R^qf_*\mathbb Z/\ell)_s
\longrightarrow(R^qf_*\mathbb Z/\ell)_{\bar\eta}
\longrightarrow\bigoplus_{x\in F(s)}
\Phi^q_{\bar\eta}(\mathbb Z/\ell)_x
\longrightarrow\cdots
\]
give
\begin{equation}\tag{a}
\begin{aligned}
&\sum_q(-1)^q\left(
\dim(R^qf_*\mathbb Z/\ell)_{\bar\eta}
-\dim(R^qf_*\mathbb Z/\ell)_s\right)\\
&\hspace{4em}=\sum_{x\in F(s)}\sum_q(-1)^q
\dim\Phi^q_{\bar\eta}(\mathbb Z/\ell)_x,
\end{aligned}
\end{equation}
and
\begin{equation}\tag{b}
\sum_q(-1)^q\dim\operatorname{Sw}
(R^qf_*\mathbb Z/\ell)_{\bar\eta}
=\sum_{x\in F(s)}\sum_q(-1)^q
\dim\operatorname{Sw}\Phi^q_{\bar\eta}(\mathbb Z/\ell)_x.
\end{equation}
Consequently,
\begin{equation}\tag{c}
\sum_q(-1)^q\varepsilon_s(R^qf_*\mathbb Z/\ell)
=\sum_{x\in F(s)}\dim\operatorname{tot}\Phi_x.
\end{equation}
The Euler--Poincar\'e formula therefore becomes
\begin{align*}
\chi(X)
&=\sum_q(-1)^q\chi(S,R^qf_*\mathbb Z/\ell)\\
&=\sum_q(-1)^q\left(
\chi(S)\dim(R^qf_*\mathbb Z/\ell)_{\bar\eta}
-\sum_s\varepsilon_s(R^qf_*\mathbb Z/\ell)\right)\\
&=\chi(S)\sum_q(-1)^q
\dim(R^qf_*\mathbb Z/\ell)_{\bar\eta}
-\sum_s\sum_q(-1)^q\varepsilon_s(R^qf_*\mathbb Z/\ell)\\
&=\chi(S)\chi(X_{\bar\eta})
-\sum_s\sum_{x\in F(s)}\dim\operatorname{tot}\Phi_x\\
&=\chi(S)\chi(X_{\bar\eta})
-(-1)^n\sum_{x\in F}n(X/S,x).
\end{align*}
Comparison with (2.1.1) proves 2.2.


\SGApage{II}{208}
\subsubsection*{Corollary 2.3.}
\underline{Under the preceding assumptions and with the preceding
notation, suppose that $x_0\in F$ and that, for every
$x\in F$ distinct from $x_0$, the fiber of $f$ through $x$ has an
ordinary quadratic singularity at $x$.  Then}
\[
\mu(X/S,x_0)=n(X/S,x_0).
\]
For $x\ne x_0$, Proposition 1.13 gives
\[
\mu(X/S,x)=n(X/S,x),
\]
and it remains to apply 2.2.

\subsubsection*{Theorem 2.4.}
\underline{Let $S$ be a henselian trait of equal characteristic with
perfect residue field, let $f:X\to S$ be a flat morphism of finite
type, and let $\bar x$ be a geometric point of $X$ whose image in
$X$ is a closed point $x$ of $X_s$ at which $X/S$ is not smooth,
and which is isolated as such.  Suppose that $X$ is regular at $x$.
Then the} total \underline{number (1.8) of vanishing cycles of $X/S$ at $\bar x$
coincides with the Milnor number $\mu(X/S,x)$.}

\SGAvar{Let $S'$ be the henselian trait obtained from $S$ by extending the
residue field $k(s)$ to its separable closure, and put
$X'=X\times_s S'$.  The assertion of 2.4 for $X/S$ is equivalent to
the assertion for $X'/S'$.  Since $k(s)$ is already assumed perfect,
we may therefore suppose that the residue field of $S$ is
algebraically closed.

For such an $S$, let $S'$ be its completion and put
$X'=X\times_s S'$.  The Milnor number at $x$ is the same for $X/S$
and $X'/S'$, and by XIII, 2.4.2 the same is true of the total number
of vanishing cycles.  We may thus suppose in addition that $S$ is
complete, hence isomorphic to $k[[t]]$, with $k$ algebraically
closed.  If $X_s$ has dimension $n$ at $x$, we then have:

}{Let $S'$ be the henselian trait obtained from $S$ by extending the
residue field $k(s)$ to its separable closure, and put
$X'=X\times_S S'$.  The assertion of 2.4 for $X/S$ is equivalent to
the assertion for $X'/S'$.  Since $k(s)$ is already assumed perfect,
we may therefore suppose that the residue field of $S$ is
algebraically closed.

For such an $S$, let $S'$ be its completion and put
$X'=X\times_S S'$.  The Milnor number at $x$ is the same for $X/S$
and $X'/S'$, and by XIII, 2.4.2 the same is true of the total number
of vanishing cycles.  We may thus suppose in addition that $S$ is
complete, hence isomorphic to $k[[t]]$, with $k$ algebraically
closed.  If $X_s$ has dimension $n$ at $x$, we then have:

}{E7II-XVI-006}\subsubsection*{Proposition 2.5.}
\underline{There exist a smooth projective variety $Y$ over $k$,
purely of dimension $n+1$, a morphism
$f:Y\to\mathbb P^1_k$ which is smooth outside a finite set $F$ of
points of $Y$, and a point $y_0\in F$ such that:}
\begin{enumerate}
\item[(a)] \underline{if $y\in F$ and $y\ne y_0$, then $y$ is an
ordinary quadratic singular point of the fiber of $f$ through $y$;}
\item[(b)] \underline{$f(y_0)=0$, and if
$a:S\to\mathbb P^1_k$ is the natural identification of $S$ with the
completion of $\mathbb P^1_k$ at $0$, there is an} $S$-\underline{isomorphism
between the henselization of $X$ at $x$ and the henselization of
$a^*Y$ at $y_0$.}
\end{enumerate}

\SGApage{II}{\(\approx\)209}
\SGAvar{Let us show that 2.5 implies 2.4.  By 2.3 and 2.5(a), the Milnor
number of $Y/\mathbb T$ at $y_0$ coincides with the total number
of vanishing cycles.
As above, XIII, 2.4.2 gives the same result for $a^*Y/S$, hence also
for $(X/S,x)$, which by 2.5(b) is locally isomorphic to
$(a^*Y/S,y_0)$ for the \'etale topology.


}{Let us show that 2.5 implies 2.4.  By 2.3 and 2.5(a), the Milnor
number of $Y/\mathbb P^1_k$ at $y_0$ coincides with the total number
of vanishing cycles.
As above, XIII, 2.4.2 gives the same result for $a^*Y/S$, hence also
for $(X/S,x)$, which by 2.5(b) is locally isomorphic to
$(a^*Y/S,y_0)$ for the \'etale topology.


}{E7II-XVI-007}We prove 2.5.  As a $k[[t]]$-algebra, the completion of the local
ring $\mathcal O_{X,x}$ is isomorphic to a quotient
\[
\widehat{\mathcal O}_{X,x}
=k[[t,x_1,\ldots,x_{n+1}]]/(f).
\]
The implicit-function theorem XV, 1.1.2 (that is, [1], Lemma 5.10)
then shows that there is an integer $N$ such that condition 2.5(b)
holds whenever
\[
\widehat{\mathcal O}_{Y,y_0}
\simeq k[[t,x_1,\ldots,x_{n+1}]]/(f_1)
\]
as $k[[t]]$-algebras, with $f_1-f$ in the $N$th power $\mathfrak m^N$
of the maximal ideal of $k[[t,x_1,\ldots,x_{n+1}]]$.

Let $(t_0,t_1)$ and $(x_0,\ldots,x_{n+1})$ be the usual homogeneous
coordinate systems on $\mathbb P^1$ and $\mathbb P^{n+1}$.  We take
$Y/\mathbb P^1$ to be the subvariety of
$\mathbb P^1\times\mathbb P^{n+1}$ defined by an equation
\begin{equation}\tag{2.5.1}
g(t_0,t_1;x_0,\ldots,x_{n+1})=0,
\end{equation}
homogeneous of degree $d'$ in $(t_0,t_1)$ and of degree $d''$ in the
$x_i$.  More intrinsically, $g$ is a section of
\[
\mathcal O(d',d'')=\operatorname{pr}_1^*\mathcal O(d')
\otimes\operatorname{pr}_2^*\mathcal O(d'').
\]
Condition (b) holds if
\begin{equation}\tag{2.5.2}
g(1,t;1,x_1,\ldots,x_{n+1})\equiv f\pmod{\mathfrak m^N}.
\end{equation}

\SGApage{II}{\(\approx\)210}
\SGAvar{Fix $d'$ and $d''$.  The sections $g$ of type (2.5.1) are the
$k$-rational points of an affine space $A$.  For
$u\in\mathbb P^1\times\mathbb P^{n+1}$, let $B(u,k)$ denote the
affine space whose points are the $k$-jets at $u$ of sections of
$\mathcal O(d',d'')$.  Put
$u_0=(1,0;1,0,\ldots,0)$.  If $u\ne u_0$, then for sufficiently
large $d'$ and $d''$ the always-linear restriction map
\begin{equation}\tag{2.5.3}
A\longrightarrow B(u_0,N)\times B(u,2)
\end{equation}
is surjective.

}{Fix $d'$ and $d''$.  The sections $g$ of type (2.5.1) are the
$k$-rational points of an affine space $A$.  For
$u\in\mathbb P^1\times\mathbb P^{n+1}$, let $B(u,r)$ denote the
affine space whose points are the $r$-jets at $u$ of sections of
$\mathcal O(d',d'')$.  Put
$u_0=(1,0;1,0,\ldots,0)$.  If $u\ne u_0$, then for sufficiently
large $d'$ and $d''$ the always-linear restriction map
\begin{equation}\tag{2.5.3}
A\longrightarrow B(u_0,N)\times B(u,2)
\end{equation}
is surjective.

}{E7II-XVI-008}Let $C(u)\subset B(u,2)$ be the set of $2$-jets which are the
$2$-jets either of the zero germ or of a germ $\widetilde h$ such
that
\begin{enumerate}
\item[(a)] $u$ lies on the variety $V$ defined by $\widetilde h=0$;
\item[(b)] $u$ is neither a smooth point nor an ordinary quadratic
point of the fiber through $u$ of $V/\mathbb P^1$.
\end{enumerate}
It is readily checked that $C(u)$ is closed of codimension $n+3$ in
$B(u,2)$.

Let $D$ be the affine subspace of $A$ consisting of the $g$ which
satisfy (2.5.2).  For $u\ne u_0$, let $D(u)$ be the subspace of $D$
consisting of the zero $g$ and of those $g$ such that $g(u)=0$ and
$u$ is neither a smooth point nor an ordinary quadratic point of the
fiber through $u$ of the variety $V$ defined by $g=0$.  If $d'$ and
$d''$ are sufficiently large for (2.5.3) to be surjective, then
$D(u)$ has codimension $n+3$ in $D$.  The union of the $D(u)$ for
$u\ne u_0$ therefore has codimension at least
\[
1=(n+3)-(n+2),
\]
and a general element $g$ of $D$ lies in none of the $D(u)$.  It
solves the problem.

\subsection*{Bibliography}

\begin{enumerate}[label={[\arabic*]},leftmargin=3em]
\item M. Artin, ``Algebraic approximation of structures over complete
local rings,'' \emph{Publ. Math. IHES} \underline{36} (1969)
(Zariski's volume).
\item L\^e D\~ung Tr\'ang, \emph{Singularit\'es isol\'ees des
hypersurfaces complexes}, Publ. Centre Math. \'Ecole Polytechnique,
May 1969.
\item J. Milnor, \emph{Singular Points of Complex Hypersurfaces},
Princeton University Press, 1968.
\item M. Raynaud, ``Caract\'eristique d'Euler--Poincar\'e d'un
faisceau et cohomologie des vari\'et\'es ab\'eliennes,'' S\'eminaire
Bourbaki 286, February 1965.
\item J.-P. Serre, \emph{Corps locaux}, Publ. Inst. Math. Univ.
Nancago VIII, Hermann, 1962.
\item J.-P. Serre, ``SGA 5, IX. Introduction \`a la th\'eorie de
Brauer,'' reproduced in \emph{Repr\'esentations lin\'eaires des
groupes}, 3rd part, Hermann, Paris, 1967.
\end{enumerate}


\markboth{Exposé XVI}{Exposé XVI}

\ifnum\SGAform=1
\clearpage\section*{Critical apparatus}\begingroup\small\raggedright\sloppy
Readings are given in the exact TeX notation of the two English witnesses. Decimal suffixes identify the individual differences in each marked passage.

\par\medskip\noindent\hypertarget{app-E7II-XVI-001.1}{}\textbf{E7II-XVI-001.1}\quad printed folio 200.\par
\noindent Printed reading: {\ttfamily\footnotesize /\allowbreak{}}\par
\noindent Corrected reading: {\ttfamily\footnotesize \_\allowbreak{}}\par
\noindent Source: {\ttfamily\footnotesize fr/sga7ii/\{diplomati\allowbreak{}c,current\}/expose\_XV\allowbreak{}I\_body.tex;\allowbreak{} folio\allowbreak{} 200;\allowbreak{} formula\allowbreak{} 137;\allowbreak{} port/decisions/decis\allowbreak{}ions\_en7II\_r2\_002.js\allowbreak{}on\#C675;\allowbreak{} fixes/FIXLOG\_7II\_CUR\allowbreak{}RENT.jsonl\#C675;\allowbreak{} 7II:304\allowbreak{}}\par

\par\medskip\noindent\hypertarget{app-E7II-XVI-002.1}{}\textbf{E7II-XVI-002.1}\quad printed folio 201.\par
\noindent Printed reading: \emph{(empty)}\par
\noindent Corrected reading: {\ttfamily\footnotesize '\allowbreak{}}\par
\noindent Source: {\ttfamily\footnotesize fr/sga7ii/\{diplomati\allowbreak{}c,current\}/expose\_XV\allowbreak{}I\_body.tex;\allowbreak{} folio\allowbreak{} 201;\allowbreak{} formula\allowbreak{} 189;\allowbreak{} port/decisions/decis\allowbreak{}ions\_en7II\_r2\_002.js\allowbreak{}on\#C676;\allowbreak{} fixes/FIXLOG\_7II\_CUR\allowbreak{}RENT.jsonl\#C676;\allowbreak{} 7II:305\allowbreak{}}\par

\par\medskip\noindent\hypertarget{app-E7II-XVI-002.2}{}\textbf{E7II-XVI-002.2}\quad printed folio 201.\par
\noindent Printed reading: \emph{(empty)}\par
\noindent Corrected reading: {\ttfamily\footnotesize '\allowbreak{}}\par
\noindent Source: {\ttfamily\footnotesize fr/sga7ii/\{diplomati\allowbreak{}c,current\}/expose\_XV\allowbreak{}I\_body.tex;\allowbreak{} folio\allowbreak{} 201;\allowbreak{} formula\allowbreak{} 190;\allowbreak{} port/decisions/decis\allowbreak{}ions\_en7II\_r2\_002.js\allowbreak{}on\#C676;\allowbreak{} fixes/FIXLOG\_7II\_CUR\allowbreak{}RENT.jsonl\#C676;\allowbreak{} 7II:305\allowbreak{}}\par

\par\medskip\noindent\hypertarget{app-E7II-XVI-003.1}{}\textbf{E7II-XVI-003.1}\quad printed folio 202.\par
\noindent Printed reading: {\ttfamily\footnotesize \_s\allowbreak{}}\par
\noindent Corrected reading: \emph{(empty)}\par
\noindent Source: {\ttfamily\footnotesize fr/sga7ii/\{diplomati\allowbreak{}c,current\}/expose\_XV\allowbreak{}I\_body.tex;\allowbreak{} folio\allowbreak{} 202;\allowbreak{} formula\allowbreak{} 207\allowbreak{}}\par

\par\medskip\noindent\hypertarget{app-E7II-XVI-003.2}{}\textbf{E7II-XVI-003.2}\quad printed folio 202.\par
\noindent Printed reading: \emph{(empty)}\par
\noindent Corrected reading: {\ttfamily\footnotesize \textbackslash{}mathbb\allowbreak{} }\par
\noindent Source: {\ttfamily\footnotesize ED-EN-FONT-G007;\allowbreak{} F7II-XVI-004.2;\allowbreak{} fr/sga7ii/\{diplomati\allowbreak{}c,current\}/expose\_XV\allowbreak{}I\_body.tex;\allowbreak{} folio\allowbreak{} 202\allowbreak{}}\par

\par\medskip\noindent\hypertarget{app-E7II-XVI-004.1}{}\textbf{E7II-XVI-004.1}\quad printed folio 203.\par
\noindent Printed reading: {\ttfamily\footnotesize i\allowbreak{}}\par
\noindent Corrected reading: {\ttfamily\footnotesize n\allowbreak{}}\par
\noindent Source: {\ttfamily\footnotesize fr/sga7ii/\{diplomati\allowbreak{}c,current\}/expose\_XV\allowbreak{}I\_body.tex;\allowbreak{} folio\allowbreak{} 203\allowbreak{}}\par

\par\medskip\noindent\hypertarget{app-E7II-XVI-005.1}{}\textbf{E7II-XVI-005.1}\quad printed folio 207.\par
\noindent Printed reading: {\ttfamily\footnotesize -\allowbreak{}}\par
\noindent Corrected reading: {\ttfamily\footnotesize +\allowbreak{}}\par
\noindent Source: {\ttfamily\footnotesize fr/sga7ii/\{diplomati\allowbreak{}c,current\}/expose\_XV\allowbreak{}I\_body.tex;\allowbreak{} folio\allowbreak{} 207;\allowbreak{} formula\allowbreak{} 339\allowbreak{}}\par

\par\medskip\noindent\hypertarget{app-E7II-XVI-006.1}{}\textbf{E7II-XVI-006.1}\quad printed folio 208.\par
\noindent Printed reading: {\ttfamily\footnotesize s\allowbreak{}}\par
\noindent Corrected reading: {\ttfamily\footnotesize S\allowbreak{}}\par
\noindent Source: {\ttfamily\footnotesize fr/sga7ii/\{diplomati\allowbreak{}c,current\}/expose\_XV\allowbreak{}I\_body.tex;\allowbreak{} folio\allowbreak{} 208;\allowbreak{} formula\allowbreak{} 371\allowbreak{}}\par

\par\medskip\noindent\hypertarget{app-E7II-XVI-006.2}{}\textbf{E7II-XVI-006.2}\quad printed folio 208.\par
\noindent Printed reading: {\ttfamily\footnotesize s\allowbreak{}}\par
\noindent Corrected reading: {\ttfamily\footnotesize S\allowbreak{}}\par
\noindent Source: {\ttfamily\footnotesize fr/sga7ii/\{diplomati\allowbreak{}c,current\}/expose\_XV\allowbreak{}I\_body.tex;\allowbreak{} folio\allowbreak{} 208;\allowbreak{} formula\allowbreak{} 379\allowbreak{}}\par

\par\medskip\noindent\hypertarget{app-E7II-XVI-007.1}{}\textbf{E7II-XVI-007.1}\quad printed folio 209.\par
\noindent Printed reading: {\ttfamily\footnotesize T\allowbreak{}}\par
\noindent Corrected reading: {\ttfamily\footnotesize P\textasciicircum{}1\_k\allowbreak{}}\par
\noindent Source: {\ttfamily\footnotesize fr/sga7ii/\{diplomati\allowbreak{}c,current\}/expose\_XV\allowbreak{}I\_body.tex;\allowbreak{} folio\allowbreak{} 209;\allowbreak{} formula\allowbreak{} 413\allowbreak{}}\par

\par\medskip\noindent\hypertarget{app-E7II-XVI-008.1}{}\textbf{E7II-XVI-008.1}\quad printed folio 210.\par
\noindent Printed reading: {\ttfamily\footnotesize k\allowbreak{}}\par
\noindent Corrected reading: {\ttfamily\footnotesize r\allowbreak{}}\par
\noindent Source: {\ttfamily\footnotesize fr/sga7ii/\{diplomati\allowbreak{}c,current\}/expose\_XV\allowbreak{}I\_body.tex;\allowbreak{} folio\allowbreak{} 210;\allowbreak{} formula\allowbreak{} 450\allowbreak{}}\par

\par\medskip\noindent\hypertarget{app-E7II-XVI-008.2}{}\textbf{E7II-XVI-008.2}\quad printed folio 210.\par
\noindent Printed reading: {\ttfamily\footnotesize k\allowbreak{}}\par
\noindent Corrected reading: {\ttfamily\footnotesize r\allowbreak{}}\par
\noindent Source: {\ttfamily\footnotesize fr/sga7ii/\{diplomati\allowbreak{}c,current\}/expose\_XV\allowbreak{}I\_body.tex;\allowbreak{} folio\allowbreak{} 210;\allowbreak{} formula\allowbreak{} 451\allowbreak{}}\par

\par\endgroup\normalsize\fi
\clearpage\SGAtoc{Exposé XVII}
\selectlanguage{english}
\clearpage
\SGApage{II}{212}

\begin{center}
{\large SGA 7 --- EXPOS\'E XVII}\par
\vspace{2em}
{\Large\underline{\underline{Lefschetz Pencils: Existence Theorem}}}\par
\vspace{1.5em}
{\large by N. M. Katz}
\end{center}

\section*{0. Introduction}

\SGApage{II}{\(\approx\)213}
In the present expos\'e (which, in the oral seminar, came before
Expos\'es XIV and XV devoted to the Picard--Lefschetz formula), we
introduce the ``Lefschetz pencils'' of hyperplane sections of a smooth
projective scheme and give conditions for their existence.  In the
following expos\'e we make an elementary cohomological study of them,
based on the known structure of the cohomology of blown-up varieties.
The fact that certain hypotheses on Lefschetz pencils which we shall
introduce are often satisfied will follow from the Picard--Lefschetz
formula of Expos\'e XV, which is less elementary in nature than the
results of the present expos\'e and the next.

Some developments of the present expos\'e and the next will be given
elsewhere in a more general setting (general base schemes,
quasi-projectivity hypotheses in place of projectivity, and so forth):
in particular, the precohomological study of pencils of hyperplane
sections, and especially of Lefschetz pencils, in EGA V, and the
cohomological study of blown-up schemes in Jouanolou's Expos\'e VII of
SGA 5.  The reader will also find in EGA V proofs in which intrinsic
diagram chasing replaces the use of coordinates favored by the writer
of the present expos\'e.

I wish to thank P. Deligne and P. Griffiths for the help they gave me
in conceiving this expos\'e.  Throughout it, $k$ denotes an
algebraically closed field unless expressly stated otherwise.

\section*{1. A reminder on quadratic singularities (cf. VI, 6)}

\SGAvar{\subsubsection*{1.1.}
Let $Y$ be a $k$-scheme pure of dimension $n$.

}{\subsubsection*{1.1.}
Let $Y$ be a $k$-scheme pure of dimension $n-1$.

}{E7II-XVII-001}A closed point $y\in Y$ is called an \underline{ordinary quadratic
singularity} of $Y$ if:
\begin{enumerate}
\item[(a)] the completion $\widehat{\underline O}_{Y,y}$ of the local
ring of $Y$ at $y$ is isomorphic to
\[
k[[T_1,\ldots,T_n]]/(Q(T)),
\]
the quotient of the formal power-series ring in
$n=1+\dim Y$ variables by the ideal generated by an element $Q(T)$;
\item[(b)] $Q$ begins in degree two, and the subvariety of
$\mathbb P^{n-1}_k$ defined by the vanishing of the homogeneous
degree-two part of $Q(T)$ is smooth over $k$ when $n\geq2$, while the
degree-two part of $Q(T)$ is nonzero when $n=1$.
\end{enumerate}

If the characteristic of $k$ is different from two, \underline{or} if
$n$ is even, condition (b) is equivalent to:
\begin{enumerate}
\item[(b$'$)] $Q$ begins in degree two and its Hessian matrix
\[
\left(
  \frac{\partial^2Q}{\partial T_i\partial T_j}(0)
\right)_{1\leq i,j\leq n}
\]
is invertible.
\end{enumerate}
This follows from the Jacobian criterion for smoothness.

\SGApage{II}{\(\approx\)214}
\SGAvar{The standard form for the homogeneous degree-two part $Q_2$ of a
series $Q$ satisfying (b), obtainable after a linear change of
variables, is
\begin{equation}\tag{1.1.1}
\left\{
\begin{aligned}
Q_2(T)&=\sum_{i=1}^{n/2}T_{2i-1}T_{2i}
  &&\text{if $n$ is even},\\[1ex]
Q_2(T)&=T_n^2+\sum_{i=1}^{(n-1)/2}T_{2i-1}T_{2i}
  &&\text{if $n$ is odd}.
\end{aligned}
\right.
\end{equation}
This follows from the theory of quadratic forms (Bourbaki,
\emph{Alg\`ebre}, Ch.~X).

\subsubsection*{1.2.}
A closed point $y\in Y$ satisfying (a) and (b) is called a
\underline{nondegenerate quadratic singularity}.  Since the
determinant of the Hessian matrix is identically zero in
characteristic two when $n$ is odd, nondegenerate quadratic
singularities of dimension $n$ exist only when $n$ is even
or the characteristic of $k$ is different from two; in those cases
they are precisely the ordinary quadratic singularities.


}{The standard form for the homogeneous degree-two part $Q_2$ of a
series $Q$ satisfying (b), obtainable after a linear change of
variables, is
\begin{equation}\tag{1.1.1}
\left\{
\begin{aligned}
Q_2(T)&=\sum_{i=1}^{n/2}T_{2i-1}T_{2i}
  &&\text{if $n$ is even},\\[1ex]
Q_2(T)&=T_n^2+\sum_{i=1}^{(n-1)/2}T_{2i-1}T_{2i}
  &&\text{if $n$ is odd}.
\end{aligned}
\right.
\end{equation}
This follows from the theory of quadratic forms (Bourbaki,
\emph{Alg\`ebre}, Ch.~IX).

\subsubsection*{1.2.}
A closed point $y\in Y$ satisfying (a) and (b$'$) is called a
\underline{nondegenerate quadratic singularity}.  Since the
determinant of the Hessian matrix is identically zero in
characteristic two when $n$ is odd, nondegenerate quadratic
singularities presented in $n$ variables exist only when $n$ is even
or the characteristic of $k$ is different from two; in those cases
they are precisely the ordinary quadratic singularities.


}{E7II-XVII-002}\SGApage{II}{215}
\section*{2. Lefschetz pencils: statement of the results}

\subsubsection*{2.1.}
Let $P^r$ denote projective $r$-space over $k$,
$\check{\mathbb P}^{r}$ the \underline{dual} projective space of
hyperplanes in $\mathbb P^r$, and
$\operatorname{Gr}(1,\check{\mathbb P}^{r})$ the variety of lines in
$\check{\mathbb P}^{r}$.  A point
$D\in\operatorname{Gr}(1,\check{\mathbb P}^{r})$ is called a
\underline{pencil of hyperplanes} in $\mathbb P^r$.  Regarding $D$ as
a line in $\check{\mathbb P}^{r}$, we write, by abuse of notation,
$D=\{H_t\}_{t\in D}$, where the $H_t$ are the hyperplanes of
$\mathbb P^r$ which are the points of the line $D$.

The \underline{axis of $D$} is the codimension-two linear subvariety
of $\mathbb P^r$ which is the intersection $H_0\cap H_\infty$ of any
two distinct members of the pencil.  Regarding this axis as a point
of $\operatorname{Gr}(r-2,\mathbb P^r)$, the variety of
$(r-2)$-dimensional linear subvarieties of $\mathbb P^r$, the map
sending $D$ to its axis is the familiar orthogonality isomorphism
\[
\operatorname{Gr}(1,\check{\mathbb P}^{r})
\xrightarrow{\ \sim\ }
\operatorname{Gr}(r-2,\mathbb P^r).
\]

\SGAvar{\subsubsection*{2.2.}
Let $X$ be a proper, smooth, irreducible $k$-scheme of dimension
$n\geq1$, equipped with an embedding $X\hookrightarrow\mathbb P^r$.
A pencil $D=\{H_t\}_{t\in\mathbb P}$ of hyperplanes in $\mathbb P^r$ is
called a \underline{Lefschetz pencil} (relative to this embedding) if:
\begin{enumerate}
\item[(a)] the axis of $D$ meets $X$ transversely (EGA IV, 17.13.7);
\item[(b)] there is a dense open subset $U\subset D$ such that, for
every $t\in U$, the hyperplane $H_t$ meets $X$ transversely;
\item[(c)] for $t_0\in D-U$, the hyperplane $H_{t_0}$ meets $X$
transversely except at one point, which is an ordinary quadratic
singular point of $X\cap H_{t_0}$.
\end{enumerate}
We shall see below in 3.2.1 that the Lefschetz pencils form an
\underline{open subset} of
$\operatorname{Gr}(1,{\mathbb P}^{r})$.

}{\subsubsection*{2.2.}
Let $X$ be a proper, smooth, irreducible $k$-scheme of dimension
$n\geq1$, equipped with an embedding $X\hookrightarrow\mathbb P^r$.
A pencil $D=\{H_t\}_{t\in D}$ of hyperplanes in $\mathbb P^r$ is
called a \underline{Lefschetz pencil} (relative to this embedding) if:
\begin{enumerate}
\item[(a)] the axis of $D$ meets $X$ transversely (EGA IV, 17.13.7);
\item[(b)] there is a dense open subset $U\subset D$ such that, for
every $t\in U$, the hyperplane $H_t$ meets $X$ transversely;
\item[(c)] for $t_0\in D-U$, the hyperplane $H_{t_0}$ meets $X$
transversely except at one point, which is an ordinary quadratic
singular point of $X\cap H_{t_0}$.
\end{enumerate}
We shall see below in 3.2.1 that the Lefschetz pencils form an
\underline{open subset} of
$\operatorname{Gr}(1,\check{\mathbb P}^{r})$.

}{E7II-XVII-003}\SGApage{II}{\(\approx\)216}
\subsubsection*{2.3.}
The embedding $X\hookrightarrow\mathbb P^r$ is called a
\underline{Lefschetz embedding} if the Lefschetz pencils form a dense
open subset of $\operatorname{Gr}(1,\check{\mathbb P}^{r})$.
\footnote{More precisely, they form a dense open subset in the set of
closed points of $\operatorname{Gr}(1,\check{\mathbb P}^{r})$.  We
shall permit ourselves such abuses of language below.}

\SGAvar{\subsubsection*{2.4.}
Recall the notion of a \underline{multiple of a given embedding}
$X\hookrightarrow\mathbb P^r$.  It is the composite of that embedding
with the $d$-th Segre embedding

\[
S_d:\mathbb P^r\hookrightarrow
\mathbb P^{\binom{r+d}{d}-1},
\qquad
(X_0,\ldots,X_r)\longmapsto(,\ldots,X^W,\ldots),
\]
where
\[
X^W=X_0^{W_0}\cdots X_r^{W_o}
\quad\text{and}\quad
W_0+\cdots+W_r=d.
\]
Thus a ``section of $X$ by a hypersurface of degree $d$'' for the
given embedding---in English, a hypersurface section of degree
$d$---is simply a hyperplane section for the $d$-th multiple of that
embedding.

}{\subsubsection*{2.4.}
Recall the notion of a \underline{multiple of a given embedding}
$X\hookrightarrow\mathbb P^r$.  It is the composite of that embedding
with the $d$-th Veronese embedding

\[
S_d:\mathbb P^r\hookrightarrow
\mathbb P^{\binom{r+d}{d}-1},
\qquad
(X_0,\ldots,X_r)\longmapsto(\ldots,X^W,\ldots),
\]
where
\[
X^W=X_0^{W_0}\cdots X_r^{W_r}
\quad\text{and}\quad
W_0+\cdots+W_r=d.
\]
Thus a ``section of $X$ by a hypersurface of degree $d$'' for the
given embedding---in English, a hypersurface section of degree
$d$---is simply a hyperplane section for the $d$-th multiple of that
embedding.

}{E7II-XVII-004}\SGApage{II}{217}
\subsubsection*{Theorem 2.5.}
\underline{Let $X$ be irreducible, proper, and smooth over $k$, of
dimension $n\geq1$, and equipped with an embedding
$i:X\hookrightarrow\mathbb P^r$.  Then:}
\begin{enumerate}
\item[2.5.1.] \underline{every $d$-th multiple of $i$, with $d\geq2$,
is a Lefschetz embedding;}
\item[2.5.2.] \underline{if $k$ has characteristic zero, the given
embedding $i$ is a Lefschetz embedding.}
\end{enumerate}
The proof is divided into several pieces, which follow.


\section*{3. The dual variety}

\SGAvar{\subsubsection*{3.1.}
Let $I$ be the ideal in $\mathbb P^r$ defining $X$.  Write
$N=I/I^2$ for the conormal sheaf, locally free of rank $r-n$ on $X$,
and let $\mathbb P(N)$ be the projective bundle over $X$ formed by the
lines in $\check N$.  (When $X=\mathbb P^r$, the bundle
$\mathbb P(N)$ is empty.)  The bundle $\mathbb P(N)$ identifies with
the subvariety of $X\times\check{\mathbb P}^{r}$ formed by the pairs
$(x,H)$ for which $H$ is tangent to $X$ at $x$, by the map
\[
\nu:\mathbb P(N)\longrightarrow X\times\check{\mathbb P}^{r}.
\]
More explicitly,
\begin{equation}\tag{3.1.1}
\mathbb P(I/I^2)_x\ni F\longmapsto
\left(x,\text{ the hyperplane with equation }
\sum_i\frac{\partial F}{\partial X_i}(x)T_i=0\right),
\end{equation}
where $F$ denotes a local section of $I/I^2$.  Define
\begin{equation}\tag{3.1.2}
\varphi:\mathbb P(N)\longrightarrow\check{\mathbb P}^{r}
\end{equation}
to be the composite
\[
\mathbb P(N)\xrightarrow{\nu}
X\times\check{\mathbb P}^{r}
\stackrel{\operatorname{pr}_2}{\relbar\joinrel\relbar\joinrel\relbar}\check{\mathbb P}^{r}.
\]

}{\subsubsection*{3.1.}
Let $I$ be the ideal in $\mathbb P^r$ defining $X$.  Write
$N=I/I^2$ for the conormal sheaf, locally free of rank $r-n$ on $X$,
and let $\mathbb P(N)$ be the projective bundle over $X$ formed by the
lines in $\check N$.  (When $X=\mathbb P^r$, the bundle
$\mathbb P(N)$ is empty.)  The bundle $\mathbb P(N)$ identifies with
the subvariety of $X\times\check{\mathbb P}^{r}$ formed by the pairs
$(x,H)$ for which $H$ is tangent to $X$ at $x$, by the map
\[
\nu:\mathbb P(N)\longrightarrow X\times\check{\mathbb P}^{r}.
\]
More explicitly,
\begin{equation}\tag{3.1.1}
\mathbb P(I/I^2)_x\ni F\longmapsto
\left(x,\text{ the hyperplane with equation }
\sum_i\frac{\partial F}{\partial X_i}(x)T_i=0\right),
\end{equation}
where $F$ denotes a local section of $I/I^2$.  Define
\begin{equation}\tag{3.1.2}
\varphi:\mathbb P(N)\longrightarrow\check{\mathbb P}^{r}
\end{equation}
to be the composite
\[
\mathbb P(N)\xrightarrow{\nu}
X\times\check{\mathbb P}^{r}
\xrightarrow{\operatorname{pr}_2}\check{\mathbb P}^{r}.
\]

}{E7II-XVII-005}\SGApage{II}{\(\approx\)218}
\subsubsection*{3.1.3.}
Let $\check X$ denote the image of $\varphi$.  It is the set of
hyperplanes which do not meet $X$ transversely, and is called the
\underline{dual variety} of $X$ for the embedding
$i:X\hookrightarrow\mathbb P^r$.

\subsubsection*{Proposition 3.1.4.}
\underline{The dual variety $\check X$ is irreducible and proper, of
dimension at most $r-1$.}

\noindent\underline{Proof.}
It is the image under $\varphi$ of $\mathbb P(N)$, which is itself
irreducible of dimension $r-1$ if $X\ne\mathbb P^r$, and empty if
$X=\mathbb P^r$.

\SGAvar{\subsubsection*{Remark 3.1.5.}
The interpretation of $\mathbb P(N)$ can be made more precise as
follows.  Let
$\mathcal Y\subset X\times\check{\mathbb P}^{r}$ be the incidence
subscheme whose fiber over every point
$H\in\check{\mathbb P}^{r}(k)$ is $X\cap H$.  The scheme $\mathcal Y$
is plainly smooth of dimension $n+r-1$, being a projective bundle of
relative dimension $r-1$ over $X$.  Thus
$Y\to\check{\mathbb P}^{r}$ makes $\mathcal Y$ a relative
complete intersection of virtual relative dimension
$(n+r-1)-r=n-1$.

}{\subsubsection*{Remark 3.1.5.}
The interpretation of $\mathbb P(N)$ can be made more precise as
follows.  Let
$\mathcal Y\subset X\times\check{\mathbb P}^{r}$ be the incidence
subscheme whose fiber over every point
$H\in\check{\mathbb P}^{r}(k)$ is $X\cap H$.  The scheme $\mathcal Y$
is plainly smooth of dimension $n+r-1$, being a projective bundle of
relative dimension $r-1$ over $X$.  Thus
$\mathcal Y\to\check{\mathbb P}^{r}$ makes $\mathcal Y$ a relative
complete intersection of virtual relative dimension
$(n+r-1)-r=n-1$.

}{E7II-XVII-006}The \underline{Jacobian subscheme}
\[
J_{n-1}(\mathcal Y/\check P^{r})
\]
(VI, 5.2) is then precisely $\mathbb P(N)$.  Formation of the
Jacobian subscheme of a relative complete intersection is compatible
with arbitrary base change.  Consequently, for every line
$D\subset\check{\mathbb P}^{r}$, the Jacobian subscheme
$J_{n-1}(\widetilde X/D)$, where
\[
\widetilde X=\mathcal Y\times_{\check{\mathbb P}^{r}}D,
\]
is precisely
$\mathbb P(N)\times_{\check{\mathbb P}^{r}}D$.


\subsubsection*{Proposition 3.2.}
\underline{The subset $B\subset\check X$ formed by the hyperplanes
$H$ such that $X\cap H$ has exactly one singular point, and that point
is an ordinary quadratic singularity, is open.}

\SGApage{II}{\(\approx\)219}
This is a special case of XV, 1.3.2, based on R. Elkik's theory of
``henselian modules'' (XV, 1.1.4).  The case in which one restricts to
\underline{nondegenerate} quadratic points---that is, where
$\operatorname{char}(k)\ne2$ \underline{or} $n=\dim X$ is even---is
also immediate from 3.3 below, whose proof does not use 3.2.

\subsubsection*{Corollary 3.2.1.}
\underline{Under the hypotheses of 3.2, let $F=\check X-B$, a closed
subset of $\check X$.  Then the Lefschetz pencils form an open subset
of $\operatorname{Gr}(1,\check{\mathbb P}^{r})$, and the following
conditions are equivalent:}
\begin{enumerate}
\item[3.2.2.] \underline{the given embedding
$i:X\hookrightarrow\mathbb P^r$ is a Lefschetz embedding;}
\item[3.2.3.] \underline{there exists a Lefschetz pencil of
hyperplanes for the given embedding;}
\item[3.2.4.] \underline{the codimension of $F$ in
$\check{\mathbb P}^{r}$ is at least two.}
\end{enumerate}

\SGAvar{\noindent\underline{Proof.}
A line $D\subset\check{\mathbb P}^{r}$ is a Lefschetz pencil if and
only if:
\begin{enumerate}
\item[(3.2.5)] the axis of $D$ meets $X$ transversely;
\item[(3.2.6)] $D$ is not contained in $ X$;
\item[(3.2.7)] $D$ does not meet $F$.
\end{enumerate}
Conditions (3.2.5) and (3.2.6) define open subsets of
$\operatorname{Gr}(1,\check{\mathbb P}^{r})$; since $F$ is closed, so
does (3.2.7).  Hence the Lefschetz pencils form an open subset.

Because $\dim X\leq r-1$, the open subset in (3.2.6) is nonempty and
therefore dense.  The open subset in (3.2.5) is likewise nonempty and
dense.  Indeed, choose a hyperplane $H$ meeting $X$ transversely; such
an $H$ exists because $\dim\check X\leq r-1$.  Since
$\dim\check{(X\mathbin{\cdot}G)}\leq r-1$, there is a hyperplane $G$ meeting
$X\cap H$ transversely, and the line $D$ through $H$ and $G$ then has
an axis meeting $X$ transversely.  Thus the embedding is Lefschetz if
and only if the open subset in (3.2.7) is nonempty, hence dense.  Since
$F$ is closed, this occurs if and only if $F$ has codimension at least
two in $\check{\mathbb P}^{r}$ (cf. EGA V, or Mumford,
\emph{Introduction to Algebraic Geometry}, p.~88, Cor.~3).

}{\noindent\underline{Proof.}
A line $D\subset\check{\mathbb P}^{r}$ is a Lefschetz pencil if and
only if:
\begin{enumerate}
\item[(3.2.5)] the axis of $D$ meets $X$ transversely;
\item[(3.2.6)] $D$ is not contained in $\check X$;
\item[(3.2.7)] $D$ does not meet $F$.
\end{enumerate}
Conditions (3.2.5) and (3.2.6) define open subsets of
$\operatorname{Gr}(1,\check{\mathbb P}^{r})$; since $F$ is closed, so
does (3.2.7).  Hence the Lefschetz pencils form an open subset.

Because $\dim\check X\leq r-1$, the open subset in (3.2.6) is nonempty and
therefore dense.  The open subset in (3.2.5) is likewise nonempty and
dense.  Indeed, choose a hyperplane $H$ meeting $X$ transversely; such
an $H$ exists because $\dim\check X\leq r-1$.  Since
$\dim\check{(X\mathbin{\cdot}H)}\leq r-1$, there is a hyperplane $G$ meeting
$X\cap H$ transversely, and the line $D$ through $H$ and $G$ then has
an axis meeting $X$ transversely.  Thus the embedding is Lefschetz if
and only if the open subset in (3.2.7) is nonempty, hence dense.  Since
$F$ is closed, this occurs if and only if $F$ has codimension at least
two in $\check{\mathbb P}^{r}$ (cf. EGA V, or Mumford,
\emph{Introduction to Algebraic Geometry}, p.~88, Cor.~3).

}{E7II-XVII-007}\SGApage{II}{220}
\subsubsection*{Corollary 3.2.8.}
\underline{Under the hypotheses of 3.2, suppose that
$i:X\hookrightarrow\mathbb P^r$ is a Lefschetz embedding.  Let $H_0$
be a hyperplane meeting $X$ transversely, and let
$\operatorname{Gr}(1,\check{\mathbb P}^{r})_{H_0}$ be the subvariety
of $\operatorname{Gr}(1,\check{\mathbb P}^{r})$ formed by the lines
through $H_0$.  Then the Lefschetz pencils through $H_0$ form a dense
open subset of $\operatorname{Gr}(1,\check{\mathbb P}^{r})_{H_0}$.}

\noindent\underline{Proof.}
By 3.2.1 these pencils form an open subset of
$\operatorname{Gr}(1,\check{\mathbb P}^{r})_{H_0}$.  The argument
just given to prove that the open subsets (3.2.5) and (3.2.6) are
nonempty also shows that their intersections with this subvariety are
nonempty.  It remains to prove that some lines through $H_0$ do not
meet $F$.

\SGApage{II}{\(\approx\)221}
Choose a hyperplane $T\subset\check{\mathbb P}^{r}$ not containing
$H_0$.  There is an isomorphism
\begin{center}
\begin{tikzcd}[column sep=huge]
\operatorname{Gr}(1,\check{\mathbb P}^{r})_{H_0}
  \& T \arrow[l,"\sim"']\\
\substack{\text{the line through}\\ H_0\text{ and }\tau}
  \& \tau \arrow[l,mapsto]
\end{tikzcd}
\end{center}
so $\operatorname{Gr}(1,\check{\mathbb P}^{r})_{H_0}$ ``is'' a
projective space of dimension $r-1$.  The lines in it which pass
through a point of $F$ form a closed subset of dimension at most
$r-2$, being the image of the morphism
\begin{center}
\begin{tikzcd}[column sep=huge]
F \arrow[r] \& \operatorname{Gr}(1,\check{\mathbb P}^{r})_{H_0}\\
f \arrow[r,mapsto] \&
  \substack{\text{the line through }H_0\\\text{and through }f}.
\end{tikzcd}
\end{center}

\subsubsection*{Proposition 3.3.}
\underline{The morphism
$\varphi:\mathbb P(N)\to\check{\mathbb P}^{r}$ is unramified at a
closed point $(x_0,H)$ if and only if $x_0$ is a nondegenerate
quadratic singularity (1.2) of $X\cap H$.  In particular, the subset
of $\mathbb P(N)$ formed by the points $(x,H)$ for which $x$ is a
nondegenerate quadratic singularity of $X\cap H$ is open.}


\SGAvar{\noindent\underline{Proof of Proposition 3.3.}
We may restrict to the case $X\ne\mathbb P^r$.  The two properties
asserted to be equivalent are independent of the particular choice of
coordinates, so we may describe them using coordinates chosen in a
very special way.  Since $X$ is a smooth subvariety of $\mathbb P^r$
of dimension $n\geq1$, choose homogeneous coordinates
$X_0,\ldots,X_r$ so that:
\begin{enumerate}
\item[3.3.1.] the point $x_0\in\mathbb P^r$ is $(1,0,\ldots,0)$;
\item[3.3.2.] putting $x_i=X_i/X_0$ for $i=1,\ldots,r$, so that
$\widehat{\mathcal O}_{\mathbb P^r,x_0}$ identifies with
$k[[x_1,\ldots,x_r]]$, the ideal defining the trace of $X$ in
$\operatorname{Spec}\widehat{\mathcal O}_{\mathbb P^r,x_0}$ is
generated by elements $g_i$ $(i=1,\ldots,r-n)$ which, in the
completion, have the form
\[
g_i=x_{n+i}-h_i(x_1,\ldots,x_n),
\qquad i=1,\ldots,r-n,
\]
where the $h_i$ are power series in $x_1,\ldots,x_n$ alone and lie in
the square of the maximal ideal of $k[[x_1,\ldots,x_n]]$;
\item[3.3.3.] under (3.1.1), the hyperplane $H$ comes from
$g_{n-r}$.
\end{enumerate}

Put $V=\operatorname{Spec}(\widehat{\mathcal O}_{X,x_0})$.  It follows
from 3.3.2 that there is a trivialization
\[
(I/I^2)|_V\simeq\mathcal O_V^{r-n}
\]
given by
\begin{equation}\tag{3.3.4}
(\lambda_1,\ldots,\lambda_{r-n})\longmapsto
\sum_1^{n-r}\lambda_i g_{n+i}.
\end{equation}
Together with 3.3.3 this gives a neighborhood $W$ of $(x_0,H)$ in
$\mathbb P(N)|_V$ and an isomorphism
$V\times\mathbb A^{r-n-1}\simeq W$, given by
\begin{equation}\tag{3.3.5}
\begin{aligned}
V\times\mathbb A^{r-n-1}&\longrightarrow\mathbb P(I/I^2)|_V,\\
(v,(\lambda_1,\ldots,\lambda_{r-n-1}))&\longmapsto
(v,(g_{r-n}+\sum_1^{n-r-1}\lambda_i g_i),
\end{aligned}
\end{equation}
where $\lambda_i=\Lambda_i/\Lambda_{n-r}$.

In particular, the completion of the local ring of $\mathbb P(\check N)$ at
$(x_0,H)$ identifies with
\[
k[[x_1,\ldots,x_n,\lambda_1,\ldots,\lambda_{r-n-1}]].
\]
We now express $\varphi$ in these coordinates at the level of
completions:
\begin{equation}\tag{3.3.6}
\varphi(x,\lambda)
=\text{the hyperplane with equation }
\sum_{i=0}^{n+r}\varphi_i(x,\lambda)T_i=0.
\end{equation}
Since this image hyperplane contains the point $x$, whose homogeneous
coordinates are
\begin{equation}\tag{3.3.7}
x=(1,x_1,\ldots,x_n,h_1(x),\ldots,h_{r-n}(x),
\end{equation}
we have
\begin{equation}\tag{3.3.8}
\varphi_0(x,\lambda)
=-\sum_{i=1}^{n}x_i\varphi_i(x,\lambda)
 -\sum_{i=1}^{r-n}h_i(x)\varphi_{n+i}(x,\lambda)
\end{equation}
in $\widehat{\mathcal O}_{\mathbb P(N),(x_0,H)}$.  Moreover, by
(3.1.1), if
\[
D_i=\frac{\partial}{\partial x_i}\in
\operatorname{Der}_k(
\widehat{\mathcal O}_{\mathbb P^r,x_0},
\widehat{\mathcal O}_{\mathbb P^r,x_0}),
\]
then for $1\leq i\leq r$,
\begin{equation}\tag{3.3.9}
\varphi_i(x,\lambda)
=D_i\!\left(g_{r-n}+
\sum_{i=1}^{r-n-1}\lambda_ig_i\right)(x)
\quad\text{in }\widehat{\mathcal O}_{\mathbb P(N),(x_0H)}.
\end{equation}

Recalling that $g_i=x_{n+i}-h_i(x_1,\ldots,x_n)$, one obtains
\begin{equation}\tag{3.3.10}
\varphi_i(x,\lambda)
=-\frac{\partial}{\partial x_i}
\left(h_{r-n}+\sum_{i=1}^{r-n-1}\lambda_ih_i\right),
\qquad 1\leq i\leq n,
\end{equation}
\begin{equation}\tag{3.3.11}
\varphi_{i}(x,\lambda)=\lambda_i,
\qquad 1\leq i\leq n,
\end{equation}
and
\begin{equation}\tag{3.3.12}
\varphi_r(x,\lambda)=1.
\end{equation}


}{\noindent\underline{Proof of Proposition 3.3.}
We may restrict to the case $X\ne\mathbb P^r$.  The two properties
asserted to be equivalent are independent of the particular choice of
coordinates, so we may describe them using coordinates chosen in a
very special way.  Since $X$ is a smooth subvariety of $\mathbb P^r$
of dimension $n\geq1$, choose homogeneous coordinates
$X_0,\ldots,X_r$ so that:
\begin{enumerate}
\item[3.3.1.] the point $x_0\in\mathbb P^r$ is $(1,0,\ldots,0)$;
\item[3.3.2.] putting $x_i=X_i/X_0$ for $i=1,\ldots,r$, so that
$\widehat{\mathcal O}_{\mathbb P^r,x_0}$ identifies with
$k[[x_1,\ldots,x_r]]$, the ideal defining the trace of $X$ in
$\operatorname{Spec}\widehat{\mathcal O}_{\mathbb P^r,x_0}$ is
generated by elements $g_i$ $(i=1,\ldots,r-n)$ which, in the
completion, have the form
\[
g_i=x_{n+i}-h_i(x_1,\ldots,x_n),
\qquad i=1,\ldots,r-n,
\]
where the $h_i$ are power series in $x_1,\ldots,x_n$ alone and lie in
the square of the maximal ideal of $k[[x_1,\ldots,x_n]]$;
\item[3.3.3.] under (3.1.1), the hyperplane $H$ comes from
$g_{r-n}$.
\end{enumerate}

Put $V=\operatorname{Spec}(\widehat{\mathcal O}_{X,x_0})$.  It follows
from 3.3.2 that there is a trivialization
\[
(I/I^2)|_V\simeq\mathcal O_V^{r-n}
\]
given by
\begin{equation}\tag{3.3.4}
(\Lambda_1,\ldots,\Lambda_{r-n})\longmapsto
\sum_{i=1}^{r-n}\Lambda_i g_i.
\end{equation}
Together with 3.3.3 this gives a neighborhood $W$ of $(x_0,H)$ in
$\mathbb P(N)|_V$ and an isomorphism
$V\times\mathbb A^{r-n-1}\simeq W$, given by
\begin{equation}\tag{3.3.5}
\begin{aligned}
V\times\mathbb A^{r-n-1}&\longrightarrow\mathbb P(I/I^2)|_V,\\
(v,(\lambda_1,\ldots,\lambda_{r-n-1}))&\longmapsto
\left(v,\left[g_{r-n}+\sum_{i=1}^{r-n-1}\lambda_i g_i\right]\right),
\end{aligned}
\end{equation}
where $\lambda_i=\Lambda_i/\Lambda_{r-n}$.

In particular, the completion of the local ring of $\mathbb P(N)$ at
$(x_0,H)$ identifies with
\[
k[[x_1,\ldots,x_n,\lambda_1,\ldots,\lambda_{r-n-1}]].
\]
We now express $\varphi$ in these coordinates at the level of
completions:
\begin{equation}\tag{3.3.6}
\varphi(x,\lambda)
=\text{the hyperplane with equation }
\sum_{i=0}^{r}\varphi_i(x,\lambda)T_i=0.
\end{equation}
Since this image hyperplane contains the point $x$, whose homogeneous
coordinates are
\begin{equation}\tag{3.3.7}
x=(1,x_1,\ldots,x_n,h_1(x),\ldots,h_{r-n}(x)),
\end{equation}
we have
\begin{equation}\tag{3.3.8}
\varphi_0(x,\lambda)
=-\sum_{i=1}^{n}x_i\varphi_i(x,\lambda)
 -\sum_{j=1}^{r-n}h_j(x)\varphi_{n+j}(x,\lambda)
\end{equation}
in $\widehat{\mathcal O}_{\mathbb P(N),(x_0,H)}$.  Moreover, by
(3.1.1), if
\[
D_i=\frac{\partial}{\partial x_i}\in
\operatorname{Der}_k(
\widehat{\mathcal O}_{\mathbb P^r,x_0},
\widehat{\mathcal O}_{\mathbb P^r,x_0}),
\]
then for $1\leq i\leq r$,
\begin{equation}\tag{3.3.9}
\varphi_i(x,\lambda)
=D_i\!\left(g_{r-n}+
\sum_{j=1}^{r-n-1}\lambda_jg_j\right)(x)
\quad\text{in }\widehat{\mathcal O}_{\mathbb P(N),(x_0,H)}.
\end{equation}

Recalling that $g_i=x_{n+i}-h_i(x_1,\ldots,x_n)$, one obtains
\begin{equation}\tag{3.3.10}
\varphi_i(x,\lambda)
=-\frac{\partial}{\partial x_i}
\left(h_{r-n}+\sum_{j=1}^{r-n-1}\lambda_jh_j\right),
\qquad 1\leq i\leq n,
\end{equation}
\begin{equation}\tag{3.3.11}
\varphi_{n+j}(x,\lambda)=\lambda_j,
\qquad 1\leq j\leq r-n-1,
\end{equation}
and
\begin{equation}\tag{3.3.12}
\varphi_r(x,\lambda)=1.
\end{equation}


}{E7II-XVII-008}\SGApage{II}{\(\approx\)222}
Let $\mathfrak m$ be the maximal ideal of
$\widehat{\mathcal O}_{\mathbb P(N),(x_0,H)}$.  Since
$h_i\equiv0\pmod{\mathfrak m^2}$ for $i=1,\ldots,r-n$, we find
\begin{equation}\tag{3.3.13}
\varphi_0(x,\lambda)\equiv0\pmod{\mathfrak m^2}
\end{equation}
and
\begin{equation}\tag{3.3.14}
\varphi_i(x,\lambda)
\equiv-\frac{\partial h_{r-n}}{\partial x_i}
\pmod{\mathfrak m^2},
\qquad i=1,\ldots,n.
\end{equation}


Putting (3.3.11)--(3.3.14) together, the Jacobian matrix of
$\varphi$ at $(x_0,H)$ has the block form
\begin{equation}\tag{3.3.15}
J_{(x_0,H)}(\varphi)=
\left(
\begin{array}{c|c}
0_{1\times n} & 0_{1\times(r-n-1)}\\ \hline
-\operatorname{Hess}(h_{r-n})(0) & 0_{n\times(r-n-1)}\\ \hline
0_{(r-n-1)\times n} & I_{r-n-1}\\ \hline
0_{1\times n} & 0_{1\times(r-n-1)}
\end{array}
\right).
\end{equation}
Thus $\varphi$ is unramified at $(x_0,H)$ if and only if the Hessian
matrix
\[
\left(
\frac{\partial^2h_{r-n}}{\partial x_i\partial x_j}(0)
\right)_{1\leq i,j\leq n}
\]
is invertible.  On the other hand, 3.3.2 and 3.3.3 show that the
completion of the local ring of $X\cap H$ at $x_0$ is
\begin{equation}\tag{3.3.16}
k[[x_1,\ldots,x_n]]/
\bigl(h_{r-n}(x_1,\ldots,x_n)\bigr),
\end{equation}
and the result follows from Definition 1.2.

\SGApage{II}{226}
\subsubsection*{Example 3.4.}
Suppose that $X$ is a hypersurface in $\mathbb P^r$, defined by a
homogeneous equation $F(X_0,\ldots,X_r)=0$.  Then
$\mathbb P(N)=X$, and in homogeneous coordinates
$\varphi:X\to\check{\mathbb P}^{r}$ is given by
\begin{equation}\tag{3.4.1}
\varphi(X_0,\ldots,X_r)=
\left(
\frac{\partial F}{\partial X_0},\ldots,
\frac{\partial F}{\partial X_r}
\right).
\end{equation}
This is the \underline{Gauss morphism}.

Here is an example in which $\varphi$ is everywhere ramified.  Let
$p$ be the characteristic of $k$, let $r\geq1$, and take the
hypersurface $X$ with equation
\begin{equation}\tag{3.4.2}
\sum_{i=0}^{r}X_i^{p^r+1}=0.
\end{equation}
Then
\begin{equation}\tag{3.4.3}
\varphi(X_0,\ldots,X_r)
=(X_0^{p^r},\ldots,X_r^{p^r}).
\end{equation}
Thus $\check X$ is reduced to a point if $p=0$, while if $p>0$ one
has $\check X\simeq X$ and $\varphi$ is the Frobenius morphism.  It
follows from 3.3 (or from 3.5.0 below) that if $p>2$, this embedding is
not a Lefschetz embedding.

We now return to the general case.


\SGApage{II}{227}
\SGAvar{\subsubsection*{Proposition 3.5.}
\underline{Suppose that the morphism
$\varphi:\mathbb P(N)\to\check{\mathbb P}^{r}$ is generically
unramified (that is, is not everywhere ramified).  Then $\varphi$
induces a birational morphism from $\mathbb P(N)$ onto $\check X$.
More precisely, the following three open subsets of $ X$ are
identical}:
\begin{enumerate}
\item[(a)] \underline{the largest open subset on which $\check X$ is
smooth;}
\item[(b)] \underline{the largest open subset $V\subset\check X$ such
that $\varphi$ induces an isomorphism
$\varphi^{-1}(V)\xrightarrow{\sim}V$;}
\item[(c)] \underline{the open subset (cf. 3.2) formed by those
hyperplanes $H$ for which $X\cap H$ has exactly one singular point,
that point being a nondegenerate quadratic singularity.}
\end{enumerate}

}{\subsubsection*{Proposition 3.5.}
\underline{Suppose that the morphism
$\varphi:\mathbb P(N)\to\check{\mathbb P}^{r}$ is generically
unramified (that is, is not everywhere ramified).  Then $\varphi$
induces a birational morphism from $\mathbb P(N)$ onto $\check X$.
More precisely, the following three open subsets of $\check X$ are
identical}:
\begin{enumerate}
\item[(a)] \underline{the largest open subset on which $\check X$ is
smooth;}
\item[(b)] \underline{the largest open subset $V\subset\check X$ such
that $\varphi$ induces an isomorphism
$\varphi^{-1}(V)\xrightarrow{\sim}V$;}
\item[(c)] \underline{the open subset (cf. 3.2) formed by those
hyperplanes $H$ for which $X\cap H$ has exactly one singular point,
that point being a nondegenerate quadratic singularity.}
\end{enumerate}

}{E7II-XVII-009}\subsubsection*{Corollary 3.5.0.}
\underline{Suppose that $\operatorname{char}(k)\ne2$ or that $\dim X$
is even.  The given embedding is a Lefschetz embedding if and only if
one of the following two conditions holds}:
\begin{enumerate}
\item[(a)] $\varphi$ \underline{is generically unramified;}
\item[(b)] $\dim\check X\leq r-2$.
\end{enumerate}

Indeed, sufficiency follows from criterion (3.2.4).  In case (a),
Proposition 3.5 identifies $F$ with the singular locus of $\check X$,
so
\[
\dim F<\dim\check X=r-1;
\]
in case (b), the criterion is immediate.  Conversely, suppose that
$D$ is a Lefschetz pencil and $\dim\check X=r-1$.  There is a point
$H\in D\cap\check X$, and 3.3 together with Definition 2.2(c) shows
that $\varphi$ is unramified at the points of $\varphi^{-1}(H)$.
Hence it is generically unramified.

\noindent\underline{Proof of Proposition 3.5.}
The hypothesis implies
\[
\dim\check X=\dim\mathbb P(N)=r-1,
\]
and therefore 3.2 shows that $\check X$ is an irreducible hypersurface
in $\check{\mathbb P}^{r}$.

\SGApage{II}{\(\approx\)228}
Let $R$ be the ramification locus of $\varphi$.  By hypothesis it is
a closed subset of $\mathbb P(N)$ of dimension at most $r-2$; hence
$\varphi(R)$ is a closed subset of $\check X$ of dimension at most
$r-2$.  Let $V'$ be the dense open subset
\begin{equation}\tag{3.5.1}
V'=(\check X-\varphi(R))\cap
   (\check X-\operatorname{Sing}(\check X))
\end{equation}
of $\check X$, and denote its complement by $F$:
\begin{equation}\tag{3.5.1.1}
F=\varphi(R)\cup\operatorname{Sing}(\check X).
\end{equation}

We first prove that $\varphi$ induces an isomorphism
$\varphi^{-1}(V')\simeq V'$ by constructing an inverse
$\psi:V'\to\varphi^{-1}(V')$.  Regard $\varphi^{-1}(V')$ as embedded
in $X\times V'$ by (3.1.1), so that $\varphi$ is the second
projection.  Constructing $\psi$ is then equivalent to constructing a
morphism $\psi_1:V'\to\mathbb P^r$ which makes the diagram
\begin{equation}\tag{3.5.2}
\begin{tikzcd}[column sep=small]
\mathbb P(N) \arrow[dd]
  \& {}\supset\varphi^{-1}(V') \arrow[r,"\varphi"]
  \& V' \arrow[dd] \& \\
  \& \& \& \psi_1 \\
X \arrow[rr,hook] \& \& \mathbb P^r \&
\end{tikzcd}
\end{equation}
commute; one then puts
$\psi(v)=(\psi_1(v),v)\in\mathbb P^r\times V'$.

For any hypersurface in projective space, the Gauss morphism (3.4.1)
is defined at every smooth point.  Apply this to the hypersurface
$\check X\subset\check{\mathbb P}^{r}$.  Since $V'$ lies in the
smooth locus $V$ of $\check X$, the Gauss morphism gives
\begin{equation}\tag{3.5.3}
\psi_1:V'\longrightarrow
\check{\check{\mathbb P}}^{r}\simeq\mathbb P^r.
\end{equation}
It remains to check that (3.5.2) commutes.  Let $(x_0,H_0)$ be a
closed point of $\varphi^{-1}(V')$.  We must show that $x_0$, viewed
as a hyperplane in $\check{\mathbb P}^{r}$, is tangent to $\check X$
at $H_0$: in other words, tangency of $H_0$ to $X$ at $x_0$ implies
tangency of $x_0$ to $\check X$ at $H_0$.  Since $\varphi$ is
unramified at $(x_0,H_0)$ and $\check X$ is smooth at
$H_0=\varphi(x_0,H_0)$, its differential is an isomorphism
\begin{equation}\tag{3.5.4}
d\varphi:
T_{\mathbb P(N),(x_0,H_0)}\xrightarrow{\sim}T_{\check X,H_0}.
\end{equation}
The first part of the proof is therefore completed by the following
lemma.


\SGApage{II}{229}
\subsubsection*{Lemma 3.6 (Euler).}
\underline{Let $(x_0,H_0)$ be a closed point of $\mathbb P(N)$}, \underline{and
let $T_L$ denote the tangent bundle of the hyperplane
$L\subset\check{\mathbb P}^{r}$ defined by $x_0$.  Then, inside
$T_{\check{\mathbb P}^{r},H_0}$,}
\[
d\varphi\bigl(T_{\mathbb P(N),(x_0,H_0)}\bigr)
\subset T_{L,H_0}.
\]

\noindent\underline{Proof.}
Choose homogeneous coordinates $X_0,\ldots,X_r$ on $\mathbb P^r$ in
which $x_0$ is the point
$\underline A=(A_0,\ldots,A_r)$.  A point of $T_{X,x_0}$ is a point
of $\mathbb P^r$ with values in $k[\epsilon]/(\epsilon^2)$ of the
form
\begin{equation}\tag{3.6.1}
\underline A+\epsilon\underline B,
\end{equation}
such that, for every local section $H$ of $I$,
\begin{equation}\tag{3.6.2}
\sum_{i=0}^{r}B_i\frac{\partial H}{\partial X_i}(\underline A)=0.
\end{equation}
Choose also a local section $H$ of $I$ which gives rise to $H_0$ by
(3.1.1).  Since $\mathbb P(N)$ is locally trivial over $X$, a point
of $T_{\mathbb P(N),(x_0,H_0)}$ can be represented as
\begin{equation}\tag{3.6.3}
(\underline A+\epsilon\underline B,\ H+\epsilon F),
\end{equation}
where $F$ is a local section of $I$ and $\underline A,\underline B$
are as above.  Its image under $\varphi$ is the hyperplane over
$k[\epsilon]/(\epsilon^2)$ with equation
\begin{equation}\tag{3.6.4}
\sum_{i=0}^{r}
\frac{\partial(H+\epsilon F)}{\partial X_i}
(\underline A+\epsilon\underline B)T_i=0.
\end{equation}
Thus the assertion is equivalent to
\begin{equation}\tag{3.6.5}
\sum_{i=0}^{r}A_i
\frac{\partial(H+\epsilon F)}{\partial X_i}
(\underline A+\epsilon\underline B)=0.
\end{equation}
Because $\epsilon^2=0$, Taylor's formula gives
\begin{equation}\tag{3.6.6}
\frac{\partial(H+\epsilon F)}{\partial X_i}
(\underline A+\epsilon\underline B)
=\frac{\partial H}{\partial X_i}(\underline A)
+\epsilon\frac{\partial F}{\partial X_i}(\underline A)
+\epsilon\sum_{j=0}^{r}B_j
\frac{\partial^2H}{\partial X_j\partial X_i}(\underline A).
\end{equation}
It is therefore enough to verify
\begin{equation}\tag{3.6.7}
\sum_{i=0}^{r}A_i\frac{\partial H}{\partial X_i}(\underline A)=0,
\end{equation}
\begin{equation}\tag{3.6.8}
\sum_{i=0}^{r}A_i\frac{\partial F}{\partial X_i}(\underline A)=0,
\end{equation}
and
\begin{equation}\tag{3.6.9}
\sum_{i,j=0}^{r}A_iB_j
\frac{\partial^2H}{\partial X_i\partial X_j}(\underline A)=0.
\end{equation}
Here $H$ and $F$ are represented by homogeneous functions of degree
zero, so Euler's formula gives
\begin{equation}\tag{3.6.10}
\sum_iX_i\frac{\partial H}{\partial X_i}=0,
\qquad
\sum_iX_i\frac{\partial F}{\partial X_i}=0.
\end{equation}
Since $\partial H/\partial X_j$ is homogeneous of degree $-1$, the
same formula gives
\[
\sum_iX_i\frac{\partial^2H}{\partial X_i\partial X_j}
=-\frac{\partial H}{\partial X_j}.
\]
Consequently, (3.6.9) becomes
\begin{equation}\tag{3.6.11}
-\sum_jB_j\frac{\partial H}{\partial X_j}(\underline A)=0,
\end{equation}
which is precisely (3.6.2).

\SGApage{II}{\(\approx\)231}
This proves that $\varphi$ induces an isomorphism
\begin{equation}\tag{3.6.12}
\varphi^{-1}(V')\xrightarrow{\sim}V',
\end{equation}
whose inverse is $\psi$.


\SGApage{II}{\(\approx\)232}
Now observe that the morphism
\begin{equation}\tag{3.6.13}
\left\{
\begin{aligned}
\psi &:V'\longrightarrow
  \mathbb P^r\times V'\subset\mathbb P^r\times\check X,\\
\psi(v)&=(\psi_1(v),v)
\end{aligned}
\right.
\end{equation}
extends to a morphism
\begin{equation}\tag{3.6.14}
\left\{
\begin{aligned}
\widetilde\psi &:V\longrightarrow\mathbb P^r\times V,\\
\widetilde\psi(v)&=(\widetilde\psi_1(v),v),
\end{aligned}
\right.
\end{equation}
where
\begin{equation}\tag{3.6.15}
\widetilde\psi_1:V\longrightarrow\mathbb P^r
\end{equation}
is the Gauss morphism, now regarded as defined on all of $V$.

We have just proved
\begin{equation}\tag{3.6.16}
\psi(V')\subset\mathbb P(N).
\end{equation}
Since $V'$ is dense in $V$ and $\mathbb P(N)$ is closed in
$\mathbb P^r\times\check X$, this implies
\begin{equation}\tag{3.6.17}
\widetilde\psi(V)\subset\mathbb P(N),
\end{equation}
and in fact plainly
\begin{equation}\tag{3.6.18}
\widetilde\psi(V)\subset\varphi^{-1}(V).
\end{equation}
Again by continuity, the identities
\[
\psi\circ\varphi=\operatorname{id}_{\varphi^{-1}(V')},
\qquad
\varphi\circ\psi=\operatorname{id}_{V'}
\]
extend to
\[
\widetilde\psi\circ\varphi=\operatorname{id}_{\varphi^{-1}(V)},
\qquad
\varphi\circ\widetilde\psi=\operatorname{id}_{V}.
\]
Thus the open subsets in 3.5(a) and 3.5(b) are identical.

\SGApage{II}{\(\approx\)233}
We next show that the open subsets in 3.5(a) and 3.5(c) are
identical.  By 3.3, the latter is the subset $U\subset\check X$
defined by
\begin{equation}\tag{3.6.19}
U=\left\{H\in\check X\ \middle|\
\begin{array}{l}
\varphi^{-1}(H)\text{ consists of exactly one point, and}\\
\varphi\text{ is unramified at that point}
\end{array}
\right\}.
\end{equation}
The scheme $\varphi^{-1}(U)$ is normal, since it is smooth, and the
induced morphism
\[
\varphi:\varphi^{-1}(U)\longrightarrow U
\]
is birational, as already proved, radicial by the definition of $U$,
and proper.  It is therefore the normalization of $U$.  In
particular, $U$ is geometrically unibranch.  Since this morphism is
unramified, EGA IV, 18.10.1 shows that it is étale; EGA IV, 18.10.18
then shows that it is an open immersion.  It is consequently an
isomorphism, and $U$ is smooth.  This proves
$3.5(c)\subset 3.5(a)$.  Formula (3.6.19) makes the inclusion
$3.5(b)\subset 3.5(c)$ clear, and the equality
$3.5(a)=3.5(b)$ completes the proof.


\SGApage{II}{217}
\subsubsection*{3.7.}
Theorem 2.5 is true when $\operatorname{char}(k)\ne2$, and also when
$X$ has even dimension.

\SGApage{II}{\(\approx\)234}
\noindent\underline{Proof.}
In view of 3.5.0 and 3.3, it is enough to verify the following
assertion:

\medskip
\noindent(3.7.1) For every closed point $x_0\in X$ and every integer
$d\geq2$, there is a hypersurface $H$ of degree $d$, tangent to $X$
at $x_0$, such that $x_0$ is an ordinary quadratic singularity of
$X\cap H$.

Choose homogeneous coordinates $X_0,\ldots,X_r$ so that
$x_0=(1,0,\ldots,0)$ and the functions $X_i/X_0$, $i=1,\ldots,n$,
are local coordinates on $X$ near $x_0$.  One may take $H$ to be the
hypersurface with equation
\begin{equation}\tag{3.7.1.1}
H=X_0^{d-2}\sum_{i=1}^{n}X_i^2
\qquad\text{if }\operatorname{char}(k)\ne2,
\end{equation}
or
\begin{equation}\tag{3.7.1.2}
H=X_0^{d-2}\sum_{i=1}^{n/2}X_iX_{n/2+i}
\qquad\text{if }n\text{ is even}.
\end{equation}
When $n$ is odd (under the first hypothesis), the equivalent standard
choice recorded in the source is
\begin{equation}\tag{3.7.1.3}
H=X_0^{d-2}
\left(X_n^2+
\sum_{i=1}^{\lfloor n/2\rfloor}
X_iX_{\lfloor n/2\rfloor+i}\right).
\end{equation}

\SGAvar{For assertion 2.5.2, note that either $\varphi$ is generically
unramified, in which case 3.5.0 applies, or it is everywhere ramified.
In characteristic zero the latter case implies
$\dim X\leq r-2$, so criterion (3.2.4) is then automatically
satisfied.


}{For assertion 2.5.2, note that either $\varphi$ is generically
unramified, in which case 3.5.0 applies, or it is everywhere ramified.
In characteristic zero the latter case implies
$\dim\check X\leq r-2$, so criterion (3.2.4) is then automatically
satisfied.


}{E7II-XVII-010}\SGApage{II}{235}
\section*{4. The general case}

We shall sketch another method of proving 2.5, valid in every
characteristic, assuming the following statement, which is close to
3.1 and is a special case of XV, 1.1.4.

\medskip
\noindent(4.1) \underline{The subset of $\mathbb P(N)$ formed by the
points $(x,H)$ such that $x$ is an ordinary quadratic singularity of
$X\cap H$ is open.}

Observe that 4.1 follows from 3.3, except when
$\operatorname{char}(k)=2$ and $\dim X$ is odd.

\subsubsection*{4.1.1.}
Let $R_1$ be the closed complement in $\mathbb P(N)$ of the open set
4.1, and let
\[
F_1=\varphi(R_1)\subset\check X\subset\check{\mathbb P}^{r}
\]
be its image.

\subsubsection*{Lemma 4.1.2.}
\underline{If the open set 4.1 is nonempty, then
$\operatorname{codim}_{\check{\mathbb P}^{r}}(F_1)\geq2$.}

\noindent\underline{Proof.}
By hypothesis,
\[
\dim R_1<\dim\mathbb P(N)=r-1,
\]
and $F_1=\varphi(R_1)$.

\SGAvar{\subsubsection*{4.1.3.}
Let $F_2$ be the subset of $\check X$ formed by the hyperplanes which
are tangent to $X$ at at least two points.  Notice that $F_2$ is not
closed, but that
\begin{equation}\tag{4.1.4}
F_2\subset F_1\cup F_2.
\end{equation}
This follows from the fact that $F_1\cup F_2=F$ (cf. 3.2) is closed.
(Inclusion (4.1.4) ``says'' that a \underline{confluence} of two
singularities is not an ordinary quadratic singularity.)  In any
case, the construction in 4.2.4 will show that $F_2$ is constructible.

To prove 2.5.1 we shall apply criterion 3.2.4.  For this we need to
know that
\[
\operatorname{codim}_{\check{\mathbb P}^{r}}(F_1)\geq2
\quad\text{and}\quad
\operatorname{codim}_{\check{\mathbb P}^{r}}(F_2)\geq2.
\]
On the one hand, 3.8.1 and 4.1.2 show that, after replacing the given
embedding by its $d$-th multiple, $d\geq2$, one has
$\operatorname{codim}_{\check{\mathbb P}^{r}}(F_2)\geq2$.

}{\subsubsection*{4.1.3.}
Let $F_2$ be the subset of $\check X$ formed by the hyperplanes which
are tangent to $X$ at at least two points.  Notice that $F_2$ is not
closed, but that
\begin{equation}\tag{4.1.4}
\overline{F_2}\subset F_1\cup F_2.
\end{equation}
This follows from the fact that $F_1\cup F_2=F$ (cf. 3.2) is closed.
(Inclusion (4.1.4) ``says'' that a \underline{confluence} of two
singularities is not an ordinary quadratic singularity.)  In any
case, the construction in 4.2.4 will show that $F_2$ is constructible.

To prove 2.5.1 we shall apply criterion 3.2.4.  For this we need to
know that
\[
\operatorname{codim}_{\check{\mathbb P}^{r}}(F_1)\geq2
\quad\text{and}\quad
\operatorname{codim}_{\check{\mathbb P}^{r}}(F_2)\geq2.
\]
On the one hand, 3.7.1 and 4.1.2 show that, after replacing the given
embedding by its $d$-th multiple, $d\geq2$, one has
$\operatorname{codim}_{\check{\mathbb P}^{r}}(F_1)\geq2$.

}{E7II-XVII-011}\SGApage{II}{\(\approx\)236}
Thus 2.5.1 will follow once we know the following proposition.

\SGAvar{\subsubsection*{Proposition 4.2.}
\underline{Let $X$ be a smooth irreducible subvariety of
$\mathbb P^r$, of dimension $n\geq1$.  For every integer $d\geq2$,
the subset $F_2$ of}
\[
\check{\mathbb P}^{N}
=\check{\mathbb P}^{\binom{n+d}{d}-1}
\]
\underline{formed by the hypersurfaces of degree $d$ in
$\mathbb P^r$ which are tangent to $X$ at at least two points has
codimension at least $2$ in $\check{\mathbb P}^{N}$.}



}{\subsubsection*{Proposition 4.2.}
\underline{Let $X$ be a smooth irreducible subvariety of
$\mathbb P^r$, of dimension $n\geq1$.  For every integer $d\geq2$,
the subset $F_2$ of}
\[
\check{\mathbb P}^{N}
=\check{\mathbb P}^{\binom{r+d}{d}-1}
\]
\underline{formed by the hypersurfaces of degree $d$ in
$\mathbb P^r$ which are tangent to $X$ at at least two points has
codimension at least $2$ in $\check{\mathbb P}^{N}$.}



}{E7II-XVII-012}\subsubsection*{4.2.1.}
Begin with the special case $d=2$ and $X$ a linear subvariety of
$\mathbb P^r$.  It is enough to find a line in
$\check{\mathbb P}^{N}$ which does not meet $F$ (cf. 3.2.4).  Choose
homogeneous coordinates $X_0,\ldots,X_r$ on $\mathbb P^r$ so that
$X$ is defined by
\[
X_{n+1}=\cdots=X_r=0,
\]
and homogeneous coordinates $(\lambda,\mu)$ on $\mathbb P^1$.

\SGApage{II}{\(\approx\)237}
For $n$ even, $n=2k$, take the pencil of quadrics
\begin{equation}\tag{4.2.2}
H_{(\lambda,\mu)}
=\lambda X_0^2+(\lambda X_k+\mu X_0)X_{2k}
 +\sum_{i=1}^{k-1}(\lambda X_i+\mu X_{i+1})X_{k+i}.
\end{equation}
For $n$ odd, $n=2k+1$, take the pencil
\begin{equation}\tag{4.2.3}
H_{(\lambda,\mu)}
=(\lambda X_1+\mu X_0)X_0
 +\sum_{i=1}^{k}(\lambda X_{2i}+\mu X_i)X_{2i+1}.
\end{equation}
In both cases one checks that $X\cap H_{(\lambda,\mu)}$ is smooth for
$(\lambda,\mu)\ne(0,1)$, while $X\cap H_{(0,1)}$ has just one
singular point, which is moreover an ordinary quadratic singularity.
Thus the pencil does not meet $F$.

\subsubsection*{4.2.4.}
Let $\operatorname{Btg}^{d}(X)$ be the subvariety of
\[
\check{\mathbb P}^{N}\times
\bigl((X\times X)-\Delta(X\times X)\bigr)
\]
formed by triples $(H,x_1,x_2)$ such that $H$ is a hypersurface of
degree $d$ in $\mathbb P^r$, tangent to $X$ at $x_1$ and at $x_2$,
with $x_1\ne x_2$.  The image of $\operatorname{Btg}^{d}(X)$ under
$\operatorname{pr}_1$ is the set $F_2$ of Proposition 4.2, so it is
enough to prove
\begin{equation}\tag{4.2.5}
\dim\operatorname{Btg}^{d}(X)\leq N-2
\quad\text{if $d\geq3$, or if $d=2$ and $X$ is not linear.}
\end{equation}

\SGAvar{For this purpose, let $\operatorname{Lin}(X\times X)$ be the
subvariety of $(X\times X)-\Delta(X\times X)$ given by
\begin{equation}\tag{4.2.6}
\begin{aligned}
\operatorname{Lin}(X\times X)
&=\bigl\{(x_1,x_2)\in X\times X\mid x_1\ne x_2,\\[-1ex]
&\hspace{8em}\text{the line spanned by $x_1,x_2$ is tangent to $X$ at $x_1$}
\bigr\}\\
&\bigl\{(x_1,x_2)\in X\times X\mid x_1\ne x_2,\\[-1ex]
&\hspace{8em}x_2\in X\cap\operatorname{Tg}_X(x_1)\bigr\}.
\end{aligned}
\end{equation}
Clearly:

}{For this purpose, let $\operatorname{Lin}(X\times X)$ be the
subvariety of $(X\times X)-\Delta(X\times X)$ given by
\begin{equation}\tag{4.2.6}
\begin{aligned}
\operatorname{Lin}(X\times X)
&=\bigl\{(x_1,x_2)\in X\times X\mid x_1\ne x_2,\\[-1ex]
&\hspace{8em}\text{the line spanned by $x_1,x_2$ is tangent to $X$ at $x_1$}
\bigr\}\\
&=\bigl\{(x_1,x_2)\in X\times X\mid x_1\ne x_2,\\[-1ex]
&\hspace{8em}x_2\in X\cap\operatorname{Tg}_X(x_1)\bigr\}.
\end{aligned}
\end{equation}
Clearly:

}{E7II-XVII-013}\SGApage{II}{\(\approx\)238}
\medskip
\noindent(4.2.7) \underline{If $X$ is not a linear subvariety of
$\mathbb P^r$, then}
\[
\dim\operatorname{Lin}(X\times X)<2\dim X.
\]
Consequently, to prove (4.2.5) it is enough to show that the fibers of
the morphism
\begin{equation}\tag{4.2.8}
\operatorname{Btg}^{d}(X)
\xrightarrow{\ \operatorname{pr}_{2,3}\ }
(X\times X)-\Delta(X\times X),
\end{equation}
regarded as subvarieties of $\check{\mathbb P}^{N}$, have
\begin{equation}\tag{4.2.9}
\left\{
\begin{array}{ll}
\text{codimension at least $2\dim X+1$}
 &\text{over $\operatorname{Lin}(X\times X)$},\\[1ex]
\text{codimension at least $2\dim X+2$}
 &\text{over $(X\times X)-\Delta(X\times X)
                 -\operatorname{Lin}(X\times X)$.}
\end{array}
\right.
\end{equation}
Statement (4.2.9) follows from the next proposition by taking $L_i$
to be the tangent space to $X$ at $x_i$, for $i=1,2$.


\SGAvar{\subsubsection*{Proposition 4.3.}
\underline{Let $x_1$ and $x_2$ be two distinct closed points of
$\mathbb P^r$, and let $L_1$ and $L_2$ be two linear subspaces of
$\mathbb P^r$ such that $x_1\in L_1$ and $x_2\in L_2$.  In the space}
\[
\check{\mathbb P}^{N}
=\check{\mathbb P}^{\binom{r+d}{d}-d}
\]
\underline{of hypersurfaces of degree $d$, the linear subspace formed
by the hypersurfaces tangent to $L_1$ at $x_1$ and to $L_2$ at $x_2$
has codimension}
\[
\left\{
\begin{aligned}
\dim L_1+\dim L_2+2,
 &\quad\text{if $d\geq3$, or if $d=2$ and $L_1\cap L_2$ does not}\\[-0.2ex]
 &\quad\text{contain the line $(x_1,x_2)$ through $x_1,x_2$},\\[1ex]
\dim L_1+\dim L_2+1,
 &\quad\text{if $d=2$ and $L_1\cap L_2$ contains the line $(x_1,x_2)$.}
\end{aligned}
\right.
\]

\noindent\underline{Proof.}
For $d\geq3$, one must prove that
$\dim L_1+\dim L_2+2$ linear conditions are independent.  It is
therefore enough to treat the case
$L_1=L_2=\mathbb P^r$, that is, to determine how many independent
conditions make a hypersurface singular at both $x_1$ and $x_2$.
Choose homogeneous coordinates so that
\[
x_1=(1,0,\ldots,0),\qquad x_1=(0,1,0,\ldots,0).
\]
The hypersurface with equation $\sum_W A_WX^W=0$ is singular at
$(1,0,\ldots,0)$ if and only if all the coefficients of the monomials
$X_0^{d-1}X_i$, $i=0,\ldots,r$, vanish.  It is singular at
$(0,1,0,\ldots,0)$ if and only if all the coefficients of the
monomials $X_1^{d-1}X_i$, $i=0,\ldots,r$, vanish.  For $d\geq3$ all
these monomials are distinct, which proves the assertion in this case.

}{\subsubsection*{Proposition 4.3.}
\underline{Let $x_1$ and $x_2$ be two distinct closed points of
$\mathbb P^r$, and let $L_1$ and $L_2$ be two linear subspaces of
$\mathbb P^r$ such that $x_1\in L_1$ and $x_2\in L_2$.  In the space}
\[
\check{\mathbb P}^{N}
=\check{\mathbb P}^{\binom{r+d}{d}-1}
\]
\underline{of hypersurfaces of degree $d$, the linear subspace formed
by the hypersurfaces tangent to $L_1$ at $x_1$ and to $L_2$ at $x_2$
has codimension}
\[
\left\{
\begin{aligned}
\dim L_1+\dim L_2+2,
 &\quad\text{if $d\geq3$, or if $d=2$ and $L_1\cap L_2$ does not}\\[-0.2ex]
 &\quad\text{contain the line $(x_1,x_2)$ through $x_1,x_2$},\\[1ex]
\dim L_1+\dim L_2+1,
 &\quad\text{if $d=2$ and $L_1\cap L_2$ contains the line $(x_1,x_2)$.}
\end{aligned}
\right.
\]

\noindent\underline{Proof.}
For $d\geq3$, one must prove that
$\dim L_1+\dim L_2+2$ linear conditions are independent.  It is
therefore enough to treat the case
$L_1=L_2=\mathbb P^r$, that is, to determine how many independent
conditions make a hypersurface singular at both $x_1$ and $x_2$.
Choose homogeneous coordinates so that
\[
x_1=(1,0,\ldots,0),\qquad x_2=(0,1,0,\ldots,0).
\]
The hypersurface with equation $\sum_W A_WX^W=0$ is singular at
$(1,0,\ldots,0)$ if and only if all the coefficients of the monomials
$X_0^{d-1}X_i$, $i=0,\ldots,r$, vanish.  It is singular at
$(0,1,0,\ldots,0)$ if and only if all the coefficients of the
monomials $X_1^{d-1}X_i$, $i=0,\ldots,r$, vanish.  For $d\geq3$ all
these monomials are distinct, which proves the assertion in this case.

}{E7II-XVII-014}\SGApage{II}{\(\approx\)239}
For $d=2$, the same calculation gives
\[
\operatorname{codim}\geq\dim L_1+\dim L_2+1,
\]
because the monomial $X_0X_1$ is counted twice.

\SGAvar{Suppose first that $L_1\cap L_2$ does not contain the line
$(x_1,x_2)$.  We reduce to the case $L_1=\mathbb P^r$ and $L_2$ is a
hyperplane in $\mathbb P^r$ which does not contain $(x_1,x_2)$.
Choose homogeneous coordinates $X_0,\ldots,X_r$ so that
\[
x_1=(1,0,\ldots,0),\qquad x_2=(0,1,0,\ldots,0),
\qquad L_2=(X_0=0).
\]
A quadric $\sum_{i\leq j}A_{ij}X_iX_j$ is singular at
$(1,0,\ldots,0)$ if and only if $A_{0j}=0$ for $j=0,\ldots,r$, and
is tangent to $X_0=0$ at $(0,1,\ldots,0)$ if and only if
$A_{ij}=0$ for $j=1,\ldots,r$.  Together these are
\[
2r+1=\dim L_1+\dim L_2+2
\]
independent conditions.

}{Suppose first that $L_1\cap L_2$ does not contain the line
$(x_1,x_2)$.  We reduce to the case $L_1=\mathbb P^r$ and $L_2$ is a
hyperplane in $\mathbb P^r$ which does not contain $(x_1,x_2)$.
Choose homogeneous coordinates $X_0,\ldots,X_r$ so that
\[
x_1=(1,0,\ldots,0),\qquad x_2=(0,1,0,\ldots,0),
\qquad L_2=(X_0=0).
\]
A quadric $\sum_{i\leq j}A_{ij}X_iX_j$ is singular at
$(1,0,\ldots,0)$ if and only if $A_{0j}=0$ for $j=0,\ldots,r$, and
is tangent to $X_0=0$ at $(0,1,\ldots,0)$ if and only if
$A_{1j}=0$ for $j=1,\ldots,r$.  Together these are
\[
2r+1=\dim L_1+\dim L_2+2
\]
independent conditions.

}{E7II-XVII-015}\SGApage{II}{\(\approx\)240}
Finally, suppose that $L_1\cap L_2$ contains the line $(x_1,x_2)$.
We reduce to the case $L_1=L_2=(x_1,x_2)$.  In homogeneous
coordinates with
\[
x_1=(1,0,\ldots,0),\qquad x_2=(0,1,0,\ldots,0),
\]
the quadric $\sum_{i\leq j}A_{ij}X_iX_j$ passes through
$(1,0,\ldots,0)$ and is tangent there to $L_1$ if and only if
$A_{00}=A_{01}=0$.  It passes through $(0,1,0,\ldots,0)$ and is
tangent there to $L_2$ if and only if $A_{11}=A_{01}=0$.  Thus there
are altogether only
\[
3=\dim L_1+\dim L_2+1
\]
conditions.


\section*{5. The degree of the dual variety by an
\lq\lq elementary\rq\rq\ method}

(Cf. XVIII, 3.2 for another method when $\dim X$ is even
\underline{or} $\operatorname{char}(k)\ne2$.)

We continue to use the notation of 2.2: thus
$X\hookrightarrow\mathbb P^r$ is a smooth irreducible subvariety of
dimension $n$.

\SGAvar{\subsubsection*{5.1.}
We first consider the morphism (3.1.1)
\begin{equation}\tag{5.1.0}
\nu:\mathbb P(N)\longrightarrow X\times\check{\mathbb P}^{r}.
\end{equation}
By definition, the conormal sheaf $\check N$ occurs in the exact
sequence of coherent sheaves on $X$
\begin{equation}\tag{5.1.1}
0\longrightarrow\check N
\longrightarrow\Omega^1_{\mathbb P^r/k}|_X
\longrightarrow\Omega^1_X
\longrightarrow0.
\end{equation}
On the other hand, writing differential forms on $\mathbb P^r$ in
homogeneous coordinates gives an exact sequence of sheaves on
$\mathbb P^r$
\begin{equation}\tag{5.1.2}
0\longrightarrow\Omega^1_{\mathbb P^r/k}
\longrightarrow\mathcal O_{\mathbb P^r}(-1)^{r+1}
\longrightarrow\mathcal O_{\mathbb P^r}
\longrightarrow0.
\end{equation}
Its restriction to $X$ is therefore the exact sequence
\begin{equation}\tag{5.1.3}
0\longrightarrow\Omega^1_{\mathbb P^r/k}|_X
\longrightarrow\mathcal O_X(-1)^{r+1}
\longrightarrow\mathcal O_{\mathbb P}
\longrightarrow0.
\end{equation}

}{\subsubsection*{5.1.}
We first consider the morphism (3.1.1)
\begin{equation}\tag{5.1.0}
\nu:\mathbb P(N)\longrightarrow X\times\check{\mathbb P}^{r}.
\end{equation}
By definition, the conormal sheaf $\check N$ occurs in the exact
sequence of coherent sheaves on $X$
\begin{equation}\tag{5.1.1}
0\longrightarrow\check N
\longrightarrow\Omega^1_{\mathbb P^r/k}|_X
\longrightarrow\Omega^1_X
\longrightarrow0.
\end{equation}
On the other hand, writing differential forms on $\mathbb P^r$ in
homogeneous coordinates gives an exact sequence of sheaves on
$\mathbb P^r$
\begin{equation}\tag{5.1.2}
0\longrightarrow\Omega^1_{\mathbb P^r/k}
\longrightarrow\mathcal O_{\mathbb P^r}(-1)^{r+1}
\longrightarrow\mathcal O_{\mathbb P^r}
\longrightarrow0.
\end{equation}
Its restriction to $X$ is therefore the exact sequence
\begin{equation}\tag{5.1.3}
0\longrightarrow\Omega^1_{\mathbb P^r/k}|_X
\longrightarrow\mathcal O_X(-1)^{r+1}
\longrightarrow\mathcal O_X
\longrightarrow0.
\end{equation}

}{E7II-XVII-016}Combining (5.1.1) and (5.1.3) gives an inclusion
\begin{equation}\tag{5.1.4}
\check N\lhook\joinrel\longrightarrow\mathcal O_X(-1)^{r+1},
\end{equation}
which, by transposition, gives a surjection
\begin{equation}\tag{5.1.5}
\mathcal O_X(1)^{r+1}\longrightarrow N.
\end{equation}
Hence there is a closed immersion
\begin{equation}\tag{5.1.6}
\alpha:\mathbb P(N)
\longrightarrow\mathbb P\bigl(\mathcal O_X(1)^{r+1}\bigr)
\simeq X\times\check{\mathbb P}^{r},
\end{equation}
which is precisely (5.1.0).



\SGApage{II}{242}
\subsubsection*{5.2. The idea of the calculation.}

From now on suppose that
$\check X=\varphi(\mathbb P(N))$ has dimension $r-1$, so that it is a
hypersurface in $\check{\mathbb P}^{r}$.  Its degree is naturally
calculated by means of Chern classes [1].

For every smooth projective variety $T$, let $\operatorname{CH}(T)$
denote its Chow ring, graded by codimension.  For
$z\in\operatorname{CH}(T)$, write
\[
\int_T z
\]
for the degree of its zero-dimensional component.  The class in
$\operatorname{CH}^1(T)$ of an invertible sheaf $\mathcal L$ on $T$
will also be denoted by $\mathcal L$ whenever no confusion can arise
(for example, under the integral sign).

\SGAvar{Let $L$ be the class of a line in
$\operatorname{CH}(\check{\mathbb P}^{r})$.  Then
\begin{equation}\tag{5.2.1}
\deg(\check X)
=\int_{\check{\mathbb P}^{r}}L\cdot\check X.
\end{equation}
Set-theoretically, $\check X$ is the image of $\mathbb P(N)$ under
$\varphi$ (3.1.2).  From the viewpoint of the proper pushforward
\[
\varphi_*:\operatorname{CH}(\mathbb P(N))
\longrightarrow\operatorname{CH}(\check{\mathbb P}^{r}),
\]
however, one has
\[
\varphi_*[\mathbb P(N)]=\deg(\varphi)[\check X].
\]
Consequently, (5.2.1) becomes
\begin{equation}\tag{5.2.2}
\deg(\check X)
=\frac{1}{\deg(\varphi)}
\int_{\check{\mathbb P}^{r}}L\cdot\varphi_*[\mathbb P(N)].
\end{equation}
Here, by abuse of notation, $\varphi$ in $\deg(\varphi)$ denotes the
morphism $\mathbb P(N)\to\check X$ induced by (3.1.2); its degree is a
nonnegative integer, taken to be zero precisely when
$\dim X<r-1=\dim\mathbb P(N)$.

}{Let $L$ be the class of a line in
$\operatorname{CH}(\check{\mathbb P}^{r})$.  Then
\begin{equation}\tag{5.2.1}
\deg(\check X)
=\int_{\check{\mathbb P}^{r}}L\cdot\check X.
\end{equation}
Set-theoretically, $\check X$ is the image of $\mathbb P(N)$ under
$\varphi$ (3.1.2).  From the viewpoint of the proper pushforward
\[
\varphi_*:\operatorname{CH}(\mathbb P(N))
\longrightarrow\operatorname{CH}(\check{\mathbb P}^{r}),
\]
however, one has
\[
\varphi_*[\mathbb P(N)]=\deg(\varphi)[\check X].
\]
Consequently, (5.2.1) becomes
\begin{equation}\tag{5.2.2}
\deg(\check X)
=\frac{1}{\deg(\varphi)}
\int_{\check{\mathbb P}^{r}}L\cdot\varphi_*[\mathbb P(N)].
\end{equation}
Here, by abuse of notation, $\varphi$ in $\deg(\varphi)$ denotes the
morphism $\mathbb P(N)\to\check X$ induced by (3.1.2); its degree is a
nonnegative integer, taken to be zero precisely when
$\dim\check X<r-1=\dim\mathbb P(N)$.

}{E7II-XVII-017}The projection formula rewrites (5.2.2) as
\begin{equation}\tag{5.2.3}
\deg(\check X)
=\frac{1}{\deg(\varphi)}
\int_{\mathbb P(N)}\varphi^*(L).
\end{equation}

\SGApage{II}{\(\approx\)243}
Let $\eta\in\operatorname{CH}(\check{\mathbb P}^{r})$ be the class of
a hyperplane.  Since $L=\eta^{r-1}$, equation (5.2.3) becomes
\begin{equation}\tag{5.2.4}
\deg(\check X)
=\frac{1}{\deg(\varphi)}
\int_{\mathbb P(N)}\bigl(\varphi^*(\eta)\bigr)^{r-1}.
\end{equation}
We now calculate $\varphi^*(\eta)$ concretely.

\SGAvar{Let $\xi$ (respectively $\tau$) denote the canonical invertible sheaf
$\mathcal O_{\mathbb P(N)}(1)$ on $\mathbb P(N)$ (respectively
$\mathcal O_{\mathbb P(\mathcal O_X(1)^{r+1})}$ on
$\mathbb P(\mathcal O_X(1)^{r+1})$).  Functoriality gives, by
(5.1.6),
\begin{equation}\tag{5.2.5}
\alpha^*(\tau)=\xi.
\end{equation}
Under the identification
\[
\mathbb P\bigl(\mathcal O_X(1)^{r+1}\bigr)
\simeq X\times\check{\mathbb P}^{r},
\]
and writing $H=\mathcal O_X(1)$, one plainly has
\begin{equation}\tag{5.2.6}
\tau=\operatorname{pr}_1^*(H)\otimes\operatorname{pr}_2^*(\eta).
\end{equation}
Let $\pi:\mathbb P(N)\to X$ be the structural morphism.  Then
\[
\pi=\operatorname{pr}_1\circ\alpha,
\qquad
\varphi=\operatorname{pr}_2\circ\alpha.
\]
Applying $\alpha^*$ to (5.2.6) and using (5.2.5), we obtain
\begin{equation}\tag{5.2.6}
\xi=\alpha^*(\tau)=\pi^*(H)\otimes\varphi^*(\eta).
\end{equation}

}{Let $\xi$ (respectively $\tau$) denote the canonical invertible sheaf
$\mathcal O_{\mathbb P(N)}(1)$ on $\mathbb P(N)$ (respectively
$\mathcal O_{\mathbb P(\mathcal O_X(1)^{r+1})}(1)$ on
$\mathbb P(\mathcal O_X(1)^{r+1})$).  Functoriality gives, by
(5.1.6),
\begin{equation}\tag{5.2.5}
\alpha^*(\tau)=\xi.
\end{equation}
Under the identification
\[
\mathbb P\bigl(\mathcal O_X(1)^{r+1}\bigr)
\simeq X\times\check{\mathbb P}^{r},
\]
and writing $H=\mathcal O_X(1)$, one plainly has
\begin{equation}\tag{5.2.6}
\tau=\operatorname{pr}_1^*(H)\otimes\operatorname{pr}_2^*(\eta).
\end{equation}
Let $\pi:\mathbb P(N)\to X$ be the structural morphism.  Then
\[
\pi=\operatorname{pr}_1\circ\alpha,
\qquad
\varphi=\operatorname{pr}_2\circ\alpha.
\]
Applying $\alpha^*$ to (5.2.6) and using (5.2.5), we obtain
\begin{equation}\tag{5.2.6$'$}
\xi=\alpha^*(\tau)=\pi^*(H)\otimes\varphi^*(\eta).
\end{equation}

}{E7II-XVII-018}Thus, in $\operatorname{CH}(\mathbb P(N))$,
\begin{equation}\tag{5.2.7}
\varphi^*(\eta)=\xi-\pi^*(H),
\end{equation}
and (5.2.3) becomes
\begin{equation}\tag{5.2.8}
\deg(\check X)
=\frac{1}{\deg(\varphi)}
\int_{\mathbb P(N)}\bigl(\xi-\pi^*(H)\bigr)^{r-1}.
\end{equation}
The right-hand side of (5.2.8) is calculated by means of the
following lemma (cf. 5.3.3).


\SGApage{II}{244}
\subsubsection*{Lemma 5.3.}
\underline{Let $X$ be a smooth connected quasi-projective variety of
dimension $n$, and let $M$ be a locally free sheaf of finite rank $a$
on $X$.  Let $\xi$ denote the invertible sheaf
$\mathcal O_{\mathbb P(M)}(1)$ on $\mathbb P(M)$, and let
$\pi:\mathbb P(M)\to X$ be the structural morphism.  Recall that
$\operatorname{CH}(\mathbb P(M))$ is a free
$\operatorname{CH}(X)$-module (via $\pi^*$), with basis
$1,\xi,\ldots,\xi^{a-1}$, where $\xi$ satisfies}
\begin{equation}\tag{5.3.1}
\xi^a-C_1\xi^{a-1}+C_2\xi^{a-2}-C_3\xi^{a-3}+\cdots=0,
\end{equation}
\underline{the $C_i$ being the Chern classes of $M$ in
$\operatorname{CH}(X)$.  Put
$C(M)=1+C_1+C_2+\cdots$, the total Chern class of $M$; it is an
invertible element of $\operatorname{CH}(X)$.  Then:}

\SGApage{II}{\(\approx\)245}
\begin{enumerate}
\item[(i)] \underline{one has}
\begin{equation}\tag{5.3.2}
\pi_*\!\left(\frac{\xi^{a-1}}{1+\xi}\right)
=\frac{1}{C(M)};
\end{equation}
\item[(ii)] \underline{let
$F(T)\in\operatorname{CH}(X)[T]$ be a polynomial such that the
coefficient of $T^i$ is homogeneous of degree $n+a-1-i$.  Thus
$F(\xi)$ has maximal degree
$n+a-1=\dim\mathbb P(M)$ in
$\operatorname{CH}(\mathbb P(M))$.  Then}
\begin{equation}\tag{5.3.3}
\int_{\mathbb P(M)}F(\xi)
=(-1)^{a-1}\int_X\frac{F(-1)}{C(M)}.
\end{equation}
\end{enumerate}

\noindent\underline{Proof.}
We first deduce (5.3.3) from (5.3.2).  Write
\begin{equation}\tag{5.3.4}
\frac{1}{C(M)}=\sum_{j\geq0}b_j,\qquad
b_0=1,\qquad
b_j\in\operatorname{CH}^j(X),\qquad
b_j=0\ \text{for }j<0.
\end{equation}
Then (5.3.2) is equivalent to
\begin{equation}\tag{5.3.5}
\pi_*(\xi^{a-1+j})=(-1)^jb_j
\qquad\text{whenever }a-1+j\geq0,
\end{equation}
since $\pi_*(\xi^q)=0$ for $q<a-1$ for dimensional reasons.

\SGApage{II}{\(\approx\)246}
By linearity it suffices to prove (5.3.3) when
$F(T)=uT^{a-1+j}$, with $u\in\operatorname{CH}^{n-j}(X)$
homogeneous.  By (5.3.5) and the projection formula,
\begin{equation}\tag{5.3.6}
\begin{aligned}
\int_{\mathbb P(M)}F(\xi)
&=\int_{\mathbb P(M)}\pi^*(u)\xi^{a-1+j}\\
&=\int_Xu\,\pi_*(\xi^{a-1+j})
=(-1)^j\int_Xu\,b_j.
\end{aligned}
\end{equation}
On the other hand,
\begin{equation}\tag*{5.3.7.}
\begin{aligned}
(-1)^{a-1}\int_X\frac{F(-1)}{C(M)}
&=(-1)^{a-1}\int_X
  \frac{(-1)^{a-1+j}u}{C(M)}\\
&=(-1)^j\int_X\sum_i u\,b_i.
\end{aligned}
\end{equation}
Because $\deg u=n-j$ and $\deg b_i=i$, one has
\[
\int_Xu\,b_i=0\qquad\text{for }i\ne j.
\]
Thus (5.3.7) reduces to
\begin{equation}\tag{5.3.8}
(-1)^{a-1}\int_X\frac{F(-1)}{C(M)}
=(-1)^j\int_Xu\,b_j,
\end{equation}
which proves the asserted implication.

\SGApage{II}{\(\approx\)247}
It remains to prove (5.3.2), or equivalently (5.3.5).



\SGAvar{For dimensional reasons,
\[
\pi_*(\xi^j)=0\qquad\text{if }j<a-1,
\]
whereas
\[
\pi_*(\xi^{a-1})=1_X
\]
for the geometric reason that $\xi^{a-1}$ meets each fiber in one
point.  The projection formula therefore gives
\begin{equation}\tag{5.3.9}
\pi_*\!\left(
\sum_{j=0}^{a-1}\pi^*(U_i)\xi^i
\right)=U_{j-i}.
\end{equation}

}{For dimensional reasons,
\[
\pi_*(\xi^j)=0\qquad\text{if }j<a-1,
\]
whereas
\[
\pi_*(\xi^{a-1})=1_X
\]
for the geometric reason that $\xi^{a-1}$ meets each fiber in one
point.  The projection formula therefore gives
\begin{equation}\tag{5.3.9}
\pi_*\!\left(
\sum_{i=0}^{a-1}\pi^*(U_i)\xi^i
\right)=U_{a-1}.
\end{equation}

}{E7II-XVII-019}By (5.3.9), assertion (5.3.5) follows from the next lemma, in which
$\xi$ is replaced by $T$ and $C_i$ by $S_i$, in view of (5.3.1).

\SGAvar{\subsubsection*{Lemma 5.4.}
\underline{Let} $a>0$ \underline{be an integer and let
$A=Z[[X_1,\ldots,X_a]]$ be the ring of formal power series in
$a$ indeterminates.  Let $B$ be the $A$-algebra}
\[
B=A[[T]]\bigg/\left(\prod_{i=1}^{a}(T-X_i)\right).
\]
\underline{Thus $B$ is a free $A$-module with basis
$1,T,\ldots,T^{a-1}$:}
\begin{equation}\tag{5.4.1}
B=AT^{a-1}\oplus AT^{a-2}\oplus\cdots\oplus A.
\end{equation}
\underline{Let $S_i\in A$ be the} $i$-\underline{th elementary symmetric
function}
\[
S_i=\sum_{1\leq\alpha_1<\cdots<\alpha_i\leq a}
X_{\alpha_1}\cdots X_{\alpha_i},
\]
\underline{so that}
\begin{equation}\tag{5.4.2}
\sum_{i=0}^{a}(-1)^iS_iT^{a-i}=0
\qquad\text{in }B.
\end{equation}
\underline{The element $1+S_1+\cdots+S_a$ is invertible in $A$;
write its inverse as}
\begin{equation}\tag{5.4.3}
\frac{1}{\sum_{i=0}^{a}S_i}
=\sum_{j\geq0}b_j,
\qquad
b_j\text{ homogeneous of degree }j\text{ in }A,
\qquad b_0=1.
\end{equation}
\underline{For every integer $j\geq0$, define a homogeneous element
$a_j\in A$ of degree $j$ by writing, using (5.4.1),}
\begin{equation}\tag{5.4.4}
T^{a-1+j}
=a_jT^{a-1}+{}?\,T^{a-2}+{}?\,T^{a-3}+\cdots+{}?,
\qquad ?\in A.
\end{equation}
\underline{Then, for every $j\geq0$,}
\begin{equation}\tag{5.4.5}
a_j=(-1)^jb_j.
\end{equation}



\noindent\underline{Proof.}
In view of the definition (5.4.3) of the $b_j$, it is enough to prove
\begin{equation}\tag{5.4.6}
\left\{
\begin{aligned}
a_0&=1,\\
\sum_{i=0}^{a}(-1)^iS_i\,a_{j-i}&=0
&&\text{if }j>0.
\end{aligned}
\right.
\end{equation}
Clearly $a_0=1$.  For a fixed $j>0$, the second relation is obtained
by taking the coefficient of $T^{a-1}$, using (5.4.1), in the
relation
\begin{equation}\tag{5.4.7}
\sum_{i=0}(-1)^iS_iT^{a-1+j-i}=0
\qquad\text{in }B.
\end{equation}

}{\subsubsection*{Lemma 5.4.}
\underline{Let} $a>0$ \underline{be an integer and let
$A=\mathbb Z[[X_1,\ldots,X_a]]$ be the ring of formal power series in
$a$ indeterminates.  Let $B$ be the $A$-algebra}
\[
B=A[[T]]\bigg/\left(\prod_{i=1}^{a}(T-X_i)\right).
\]
\underline{Thus $B$ is a free $A$-module with basis
$1,T,\ldots,T^{a-1}$:}
\begin{equation}\tag{5.4.1}
B=AT^{a-1}\oplus AT^{a-2}\oplus\cdots\oplus A.
\end{equation}
\underline{Let $S_i\in A$ be the} $i$-\underline{th elementary symmetric
function}
\[
S_i=\sum_{1\leq\alpha_1<\cdots<\alpha_i\leq a}
X_{\alpha_1}\cdots X_{\alpha_i},
\]
\underline{so that}
\begin{equation}\tag{5.4.2}
\sum_{i=0}^{a}(-1)^iS_iT^{a-i}=0
\qquad\text{in }B.
\end{equation}
\underline{The element $1+S_1+\cdots+S_a$ is invertible in $A$;
write its inverse as}
\begin{equation}\tag{5.4.3}
\frac{1}{\sum_{i=0}^{a}S_i}
=\sum_{j\geq0}b_j,
\qquad
b_j\text{ homogeneous of degree }j\text{ in }A,
\qquad b_0=1.
\end{equation}
\underline{For every integer $j\geq0$, define a homogeneous element
$a_j\in A$ of degree $j$ by writing, using (5.4.1),}
\begin{equation}\tag{5.4.4}
T^{a-1+j}
=a_jT^{a-1}+{}?\,T^{a-2}+{}?\,T^{a-3}+\cdots+{}?,
\qquad ?\in A.
\end{equation}
\underline{Then, for every $j\geq0$,}
\begin{equation}\tag{5.4.5}
a_j=(-1)^jb_j.
\end{equation}



\noindent\underline{Proof.}
In view of the definition (5.4.3) of the $b_j$, it is enough to prove
\begin{equation}\tag{5.4.6}
\left\{
\begin{aligned}
a_0&=1,\\
\sum_{i=0}^{a}(-1)^iS_i\,a_{j-i}&=0
&&\text{if }j>0.
\end{aligned}
\right.
\end{equation}
Clearly $a_0=1$.  For a fixed $j>0$, the second relation is obtained
by taking the coefficient of $T^{a-1}$, using (5.4.1), in the
relation
\begin{equation}\tag{5.4.7}
\sum_{i=0}^{a}(-1)^iS_iT^{a-1+j-i}=0
\qquad\text{in }B.
\end{equation}

}{E7II-XVII-020}\SGApage{II}{\(\approx\)248}
This relation is obtained by multiplying (5.4.2) by $T^{j-1}$,
which proves the lemma.

\SGApage{II}{249}
\subsubsection*{Corollary 5.5.}
\underline{One has}
\begin{equation}\tag{5.5.1}
\deg(\check X)
=\frac{(-1)^n}{\deg(\varphi)}
\int_X\frac{(1+H)^{r-1}}{C(N)}.
\end{equation}

\noindent\underline{Proof.}
Apply (5.3.3), with $M=N$ the normal sheaf of
$X\hookrightarrow\mathbb P^r$, to (5.2.8).

For a smooth quasi-projective variety $S$, write
\begin{equation}\tag{5.5.2}
C(S)\stackrel{\mathrm{def}}{=}C(T_S)
\end{equation}
for its total Chern class, where $T_S$ denotes the tangent bundle of
$S$.

Dualizing the exact sequence (5.1.2) gives
\begin{equation}\tag{5.5.3}
C(\mathbb P^r)=(1+H)^{r+1}
\qquad\text{in }\operatorname{CH}(\mathbb P^r).
\end{equation}
Dualizing the exact sequence (5.1.1) gives
\begin{equation}\tag{5.5.4}
C(N)\,C(X)=C(\mathbb P^r)|_X
\qquad\text{in }\operatorname{CH}(X).
\end{equation}
Combining (5.5.3) and (5.5.4), formula (5.5.1) becomes:

\SGApage{II}{250}
\SGAvar{\subsubsection*{Corollary 5.6.}
\underline{One has}
\begin{equation}\tag{5.6.1}
\deg(\check X)
=\frac{(-1)^r}{\deg(\varphi)}
\int_X\frac{C(X)}{(1+H)^2}.
\end{equation}


}{\subsubsection*{Corollary 5.6.}
\underline{One has}
\begin{equation}\tag{5.6.1}
\deg(\check X)
=\frac{(-1)^n}{\deg(\varphi)}
\int_X\frac{C(X)}{(1+H)^2}.
\end{equation}


}{E7II-XVII-021}\subsubsection*{5.7.}
For a smooth projective variety $S$, define its
\underline{Euler--Poincaré characteristic} by
\begin{equation}\tag{5.7.1}
\chi(S)=\int_SC(S).
\end{equation}

\subsubsection*{Proposition 5.7.2.}
\underline{Let $Y$ be a smooth hyperplane section of $X$, and let
$\Delta$ be a smooth section of $X$ by a codimension-two linear
subspace of $\mathbb P^r$.  Then}
\begin{equation}\tag{5.7.3}
\deg(\check X)
=\frac{(-1)^n}{\deg(\varphi)}
\bigl[\chi(X)+\chi(\Delta)-2\chi(Y)\bigr].
\end{equation}



\SGApage{II}{\(\approx\)251}
\noindent\underline{Proof.}
Let $i:Y\hookrightarrow X$ be the inclusion.  One has
\[
N_{Y/X}=H|_Y=i^*H.
\]
The exact sequence
\[
0\longrightarrow T_Y
\longrightarrow i^*T_X
\longrightarrow i^*H
\longrightarrow0
\]
therefore gives
\begin{equation}\tag{5.7.4}
C(Y)=i^*\!\left(\frac{C(X)}{1+H}\right).
\end{equation}
Hence
\[
\chi(Y)=\int_Yi^*\!\left(\frac{C(X)}{1+H}\right),
\]
and the projection formula gives
\begin{equation}\tag{5.7.5}
\chi(Y)
=\int_Xi_*i^*\!\left(\frac{C(X)}{1+H}\right)
=\int_X\frac{H\,C(X)}{1+H}.
\end{equation}
Iterating this calculation for $\Delta$, one obtains
\begin{equation}\tag{5.7.6}
\begin{aligned}
\chi(\Delta)
&=\int_Y\frac{H\,C(Y)}{1+H}\\
&=\int_Yi^*\!\left(
  \frac{H}{1+H}\cdot\frac{C(X)}{1+H}
  \right)
=\int_X\frac{H^2C(X)}{(1+H)^2}.
\end{aligned}
\end{equation}
It follows that
\[
\begin{aligned}
\chi(X)-2\chi(Y)+\chi(\Delta)
&=\int_XC(X)\left[
1-\frac{2H}{1+H}+\frac{H^2}{(1+H)^2}
\right]\\
&=\int_X\frac{C(X)}{(1+H)^2}.
\end{aligned}
\]
Together with (5.6.1), this proves the proposition.

\SGApage{II}{252}
\section*{6. The case of a general base}

\subsubsection*{6.1.}
Let $S$ be a scheme and let $\bar s$ be a geometric point of $S$,
defined by an algebraic closure of the residue field $k(s)$ for a
point $s\in S$.  Let $\mathbb P^r_S$ be projective $r$-space over
$S$, let $\check{\mathbb P}^{r}_S$ be its dual projective space, and
let
\[
I\subset\mathbb P^r_S\times_S\check{\mathbb P}^{r}_S
\]
be the incidence variety.  Let
$D\subset\check{\mathbb P}^{r}_S$ be a line, defined by an
$S$-valued point of the appropriate Grassmannian, and let $\Delta$ be
the axis of the corresponding pencil of hyperplanes (cf. 2.1).

Let $X$ be a proper smooth $S$-scheme, pure of dimension $n\geq1$,
equipped with a projective embedding $X\hookrightarrow\mathbb P^r_S$.
Put
\[
\widetilde X=(X\times_SD)\cap I.
\]
Thus $\widetilde X$ is the scheme over $D$ whose fibers are the
hyperplane sections of $X$ parametrized by $D$; let
$u:\widetilde X\to D$ be the projection.

We use a subscript $\bar s$ to denote base change by
$\bar s\to S$.  Suppose that $\Delta_{\bar s}$ is transverse to
$X_{\bar s}$.  Then, over a neighborhood $U$ of $s$ in $S$,
$\Delta$ is transverse to $X$.  In what follows we replace $S$ by
this neighborhood and write $U=S$.  The scheme $\widetilde X$ is then
smooth over $S$, and a neighborhood in $\widetilde X$ of the inverse
image of $\Delta$ is smooth over $D$.

Suppose moreover that one of the hyperplane sections of $X_{\bar s}$
belonging to the pencil $D_{\bar s}$ is smooth.  After again replacing
$S$ by a suitable neighborhood of $s$, the closed subscheme
$F\subset\widetilde X$ defined by the Jacobian ideal
$J^{n-1}(\widetilde X/D)$ is finite and flat over $S$.

\SGApage{II}{\(\approx\)253}
\subsubsection*{6.2.}
The \underline{exceptional locus} of the pencil $D$ is the relative
Cartier divisor $E$ of $D/S$ defined by the finite flat
$\mathcal O_D$-module $u_*\mathcal O_F$.  Locally on $D$, if
\[
\mathcal O_D^m\xrightarrow{\ v\ }\mathcal O_D^m
\longrightarrow u_*\mathcal O_F\longrightarrow0
\]
is a presentation of $u_*\mathcal O_F$, then $E$ has equation
$\det(v)=0$.  Formation of $E$ commutes with arbitrary base change.

The divisor $E_{\bar s}$ on $D_{\bar s}$ can be written
\[
E_{\bar s}=\sum_P n(P)[P],
\]
where $n(P)$ is the sum of the colengths of the Jacobian ideal
$J^{n-1}(\widetilde X/D)$ at the critical points of $u$ mapping to
$P$.

\subsubsection*{Proposition 6.3.}
\underline{Suppose that $D_{\bar s}$ is a Lefschetz pencil of
hyperplane sections of $X_{\bar s}$ and that
$\operatorname{char}k(s)\ne2$ or $n$ is even.  There is a
neighborhood $U$ of $s$ in $S$ such that:}
\begin{enumerate}
\item[(a)] \underline{for every geometric point $\bar t$ of $U$,
$D_{\bar t}$ is a Lefschetz pencil of hyperplane sections of
$X_{\bar t}$;}
\item[(b)] \underline{over $U$, the exceptional locus $E$ is étale
over $S$.}
\end{enumerate}
To prove (a) and (b), it is enough to observe that, under these
hypotheses, $D_{\bar s}$ satisfies condition 2.2(c) if and only if
the divisor $E_{\bar s}$ is reduced.  Invoking 3.3, one could use
only the easier implication from left to right.

\SGAvar{\subsubsection*{6.3 (exceptional case).}

The proposition is false in the exceptional case just mentioned.
Nevertheless, if $S$ has characteristic two and $n$ is odd, there is
a neighborhood $U$ of $S$ over which (a) holds, and one can define a
divisor $E'$ étale over $U$ such that $E=2E'$.  We shall not use this
result.

}{\subsubsection*{6.4 (exceptional case).}

The proposition is false in the exceptional case just mentioned.
Nevertheless, if $S$ has characteristic two and $n$ is odd, there is
a neighborhood $U$ of $s$ over which (a) holds, and one can define a
divisor $E'$ étale over $U$ such that $E=2E'$.  We shall not use this
result.

}{E7II-XVII-022}\subsubsection*{References}

\noindent[1] A. Grothendieck, \emph{La théorie des classes de Chern},
Bull. Soc. Math. France \textbf{86} (1958), 137--154.



\markboth{Exposé XVII}{Exposé XVII}

\ifnum\SGAform=1
\clearpage\section*{Critical apparatus}\begingroup\small\raggedright\sloppy
Readings are given in the exact TeX notation of the two English witnesses. Decimal suffixes identify the individual differences in each marked passage.

\par\medskip\noindent\hypertarget{app-E7II-XVII-001.1}{}\textbf{E7II-XVII-001.1}\quad printed folio 213.\par
\noindent Printed reading: \emph{(empty)}\par
\noindent Corrected reading: {\ttfamily\footnotesize -1\allowbreak{}}\par
\noindent Source: {\ttfamily\footnotesize fr/sga7ii/\{diplomati\allowbreak{}c,current\}/expose\_XV\allowbreak{}II\_body.tex;\allowbreak{} folio\allowbreak{} 213;\allowbreak{} formula\allowbreak{} 4\allowbreak{}}\par

\par\medskip\noindent\hypertarget{app-E7II-XVII-002.1}{}\textbf{E7II-XVII-002.1}\quad printed folio 214.\par
\noindent Printed reading: {\ttfamily\footnotesize X\allowbreak{}}\par
\noindent Corrected reading: {\ttfamily\footnotesize IX\allowbreak{}}\par
\noindent Source: {\ttfamily\footnotesize port/PORT\_EN7II\_LOG.\allowbreak{}jsonl\#C683;\allowbreak{} fixes/FIXLOG\_7II\_CUR\allowbreak{}RENT.jsonl\#C683;\allowbreak{} 7II:321\allowbreak{}}\par

\par\medskip\noindent\hypertarget{app-E7II-XVII-002.2}{}\textbf{E7II-XVII-002.2}\quad printed folio 214.\par
\noindent Printed reading: \emph{(empty)}\par
\noindent Corrected reading: {\ttfamily\footnotesize \$'\$\allowbreak{}}\par
\noindent Source: {\ttfamily\footnotesize port/PORT\_EN7II\_LOG.\allowbreak{}jsonl\#C684;\allowbreak{} fixes/FIXLOG\_7II\_CUR\allowbreak{}RENT.jsonl\#C684;\allowbreak{} 7II:322\allowbreak{}}\par

\par\medskip\noindent\hypertarget{app-E7II-XVII-002.3}{}\textbf{E7II-XVII-002.3}\quad printed folio 214.\par
\noindent Printed reading: {\ttfamily\footnotesize of\allowbreak{}}\par
\noindent Corrected reading: {\ttfamily\footnotesize presented\allowbreak{}}\par
\noindent Source: {\ttfamily\footnotesize ED-EN-FORMULA-G131;\allowbreak{} fr/sga7ii/\{diplomati\allowbreak{}c,current\}/expose\_XV\allowbreak{}II\_body.tex;\allowbreak{} folio\allowbreak{} 214\allowbreak{}}\par

\par\medskip\noindent\hypertarget{app-E7II-XVII-002.4}{}\textbf{E7II-XVII-002.4}\quad printed folio 214.\par
\noindent Printed reading: {\ttfamily\footnotesize dimension\allowbreak{}}\par
\noindent Corrected reading: {\ttfamily\footnotesize in\allowbreak{}}\par
\noindent Source: {\ttfamily\footnotesize ED-EN-FORMULA-G131;\allowbreak{} fr/sga7ii/\{diplomati\allowbreak{}c,current\}/expose\_XV\allowbreak{}II\_body.tex;\allowbreak{} folio\allowbreak{} 214\allowbreak{}}\par

\par\medskip\noindent\hypertarget{app-E7II-XVII-002.5}{}\textbf{E7II-XVII-002.5}\quad printed folio 214.\par
\noindent Printed reading: \emph{(empty)}\par
\noindent Corrected reading: {\ttfamily\footnotesize  variables\allowbreak{}}\par
\noindent Source: {\ttfamily\footnotesize ED-EN-FORMULA-G131;\allowbreak{} fr/sga7ii/\{diplomati\allowbreak{}c,current\}/expose\_XV\allowbreak{}II\_body.tex;\allowbreak{} folio\allowbreak{} 214\allowbreak{}}\par

\par\medskip\noindent\hypertarget{app-E7II-XVII-003.1}{}\textbf{E7II-XVII-003.1}\quad printed folio 215.\par
\noindent Printed reading: {\ttfamily\footnotesize \textbackslash{}mathbb\allowbreak{}}\par
\noindent Corrected reading: \emph{(empty)}\par
\noindent Source: {\ttfamily\footnotesize ED-EN-FONT-G018;\allowbreak{} F7II-XVII-004.1;\allowbreak{} fr/sga7ii/\{diplomati\allowbreak{}c,current\}/expose\_XV\allowbreak{}II\_body.tex;\allowbreak{} folio\allowbreak{} 215\allowbreak{}}\par

\par\medskip\noindent\hypertarget{app-E7II-XVII-003.2}{}\textbf{E7II-XVII-003.2}\quad printed folio 215.\par
\noindent Printed reading: {\ttfamily\footnotesize P\allowbreak{}}\par
\noindent Corrected reading: {\ttfamily\footnotesize D\allowbreak{}}\par
\noindent Source: {\ttfamily\footnotesize ED-EN-FONT-G018;\allowbreak{} F7II-XVII-004.1;\allowbreak{} fr/sga7ii/\{diplomati\allowbreak{}c,current\}/expose\_XV\allowbreak{}II\_body.tex;\allowbreak{} folio\allowbreak{} 215;\allowbreak{} fr/sga7ii/\{diplomati\allowbreak{}c,current\}/expose\_XV\allowbreak{}II\_body.tex;\allowbreak{} folio\allowbreak{} 215;\allowbreak{} formula\allowbreak{} 62\allowbreak{}}\par

\par\medskip\noindent\hypertarget{app-E7II-XVII-003.3}{}\textbf{E7II-XVII-003.3}\quad printed folio 216.\par
\noindent Printed reading: \emph{(empty)}\par
\noindent Corrected reading: {\ttfamily\footnotesize \textbackslash{}check\allowbreak{}}\par
\noindent Source: {\ttfamily\footnotesize fr/sga7ii/\{diplomati\allowbreak{}c,current\}/expose\_XV\allowbreak{}II\_body.tex;\allowbreak{} folio\allowbreak{} 216;\allowbreak{} formula\allowbreak{} 76;\allowbreak{} port/decisions/decis\allowbreak{}ions\_en7II\_r2\_002.js\allowbreak{}on\#C687;\allowbreak{} fixes/FIXLOG\_7II\_CUR\allowbreak{}RENT.jsonl\#C687;\allowbreak{} 7II:327\allowbreak{}}\par

\par\medskip\noindent\hypertarget{app-E7II-XVII-004.1}{}\textbf{E7II-XVII-004.1}\quad printed folio 216.\par
\noindent Printed reading: {\ttfamily\footnotesize Segre\allowbreak{}}\par
\noindent Corrected reading: {\ttfamily\footnotesize Veronese\allowbreak{}}\par
\noindent Source: {\ttfamily\footnotesize ED-EN-NONMATH-G004;\allowbreak{} F7II-XVII-005.4;\allowbreak{} ADOPT-F7II-XVII-005.\allowbreak{}4;\allowbreak{} fr/sga7ii/\{diplomati\allowbreak{}c,current\}/expose\_XV\allowbreak{}II\_body.tex;\allowbreak{} folio\allowbreak{} 216\allowbreak{}}\par

\par\medskip\noindent\hypertarget{app-E7II-XVII-004.2}{}\textbf{E7II-XVII-004.2}\quad printed folio 216.\par
\noindent Printed reading: {\ttfamily\footnotesize ,\allowbreak{}}\par
\noindent Corrected reading: \emph{(empty)}\par
\noindent Source: {\ttfamily\footnotesize fr/sga7ii/\{diplomati\allowbreak{}c,current\}/expose\_XV\allowbreak{}II\_body.tex;\allowbreak{} folio\allowbreak{} 216;\allowbreak{} port/decisions/decis\allowbreak{}ions\_en7II\_r2\_002.js\allowbreak{}on\#C690;\allowbreak{} fixes/FIXLOG\_7II\_CUR\allowbreak{}RENT.jsonl\#C690;\allowbreak{} 7II:330\allowbreak{}}\par

\par\medskip\noindent\hypertarget{app-E7II-XVII-004.3}{}\textbf{E7II-XVII-004.3}\quad printed folio 216.\par
\noindent Printed reading: {\ttfamily\footnotesize o\allowbreak{}}\par
\noindent Corrected reading: {\ttfamily\footnotesize r\allowbreak{}}\par
\noindent Source: {\ttfamily\footnotesize ED-EN-FORMULA-G133;\allowbreak{} fr/sga7ii/\{diplomati\allowbreak{}c,current\}/expose\_XV\allowbreak{}II\_body.tex;\allowbreak{} folio\allowbreak{} 216\allowbreak{}}\par

\par\medskip\noindent\hypertarget{app-E7II-XVII-005.1}{}\textbf{E7II-XVII-005.1}\quad printed folio 217.\par
\noindent Printed reading: {\ttfamily\footnotesize \textbackslash{}stackrel\allowbreak{}}\par
\noindent Corrected reading: {\ttfamily\footnotesize \textbackslash{}xrightarrow\allowbreak{}}\par
\noindent Source: {\ttfamily\footnotesize fr/sga7ii/\{diplomati\allowbreak{}c,current\}/expose\_XV\allowbreak{}II\_body.tex;\allowbreak{} folio\allowbreak{} 217;\allowbreak{} formula\allowbreak{} 130\allowbreak{}}\par

\par\medskip\noindent\hypertarget{app-E7II-XVII-005.2}{}\textbf{E7II-XVII-005.2}\quad printed folio 217.\par
\noindent Printed reading: {\ttfamily\footnotesize \}\{\textbackslash{}relbar\textbackslash{}joinrel\textbackslash{}re\allowbreak{}lbar\textbackslash{}joinrel\textbackslash{}relbar\allowbreak{}}\par
\noindent Corrected reading: \emph{(empty)}\par
\noindent Source: {\ttfamily\footnotesize fr/sga7ii/\{diplomati\allowbreak{}c,current\}/expose\_XV\allowbreak{}II\_body.tex;\allowbreak{} folio\allowbreak{} 217;\allowbreak{} formula\allowbreak{} 130\allowbreak{}}\par

\par\medskip\noindent\hypertarget{app-E7II-XVII-006.1}{}\textbf{E7II-XVII-006.1}\quad printed folio 218.\par
\noindent Printed reading: \emph{(empty)}\par
\noindent Corrected reading: {\ttfamily\footnotesize \textbackslash{}mathcal\allowbreak{} }\par
\noindent Source: {\ttfamily\footnotesize ED-EN-FONT-G025;\allowbreak{} F7II-XVII-008.1;\allowbreak{} fr/sga7ii/\{diplomati\allowbreak{}c,current\}/expose\_XV\allowbreak{}II\_body.tex;\allowbreak{} folio\allowbreak{} 218;\allowbreak{} fr/sga7ii/\{diplomati\allowbreak{}c,current\}/expose\_XV\allowbreak{}II\_body.tex;\allowbreak{} folio\allowbreak{} 218;\allowbreak{} port/decisions/decis\allowbreak{}ions\_en7II\_r2\_002.js\allowbreak{}on\#C692;\allowbreak{} 7II:335\allowbreak{}}\par

\par\medskip\noindent\hypertarget{app-E7II-XVII-007.1}{}\textbf{E7II-XVII-007.1}\quad printed folio 219.\par
\noindent Printed reading: \emph{(empty)}\par
\noindent Corrected reading: {\ttfamily\footnotesize \textbackslash{}check\allowbreak{}}\par
\noindent Source: {\ttfamily\footnotesize fr/sga7ii/\{diplomati\allowbreak{}c,current\}/expose\_XV\allowbreak{}II\_body.tex;\allowbreak{} folio\allowbreak{} 219;\allowbreak{} formula\allowbreak{} 183\allowbreak{}}\par

\par\medskip\noindent\hypertarget{app-E7II-XVII-007.2}{}\textbf{E7II-XVII-007.2}\quad printed folio 220.\par
\noindent Printed reading: \emph{(empty)}\par
\noindent Corrected reading: {\ttfamily\footnotesize \textbackslash{}check\allowbreak{}}\par
\noindent Source: {\ttfamily\footnotesize fr/sga7ii/\{diplomati\allowbreak{}c,current\}/expose\_XV\allowbreak{}II\_body.tex;\allowbreak{} folio\allowbreak{} 220;\allowbreak{} formula\allowbreak{} 189;\allowbreak{} port/PORT\_EN7II\_LOG.\allowbreak{}jsonl\#C696;\allowbreak{} fixes/FIXLOG\_7II\_CUR\allowbreak{}RENT.jsonl\#C696;\allowbreak{} 7II:341\allowbreak{}}\par

\par\medskip\noindent\hypertarget{app-E7II-XVII-007.3}{}\textbf{E7II-XVII-007.3}\quad printed folio 220.\par
\noindent Printed reading: {\ttfamily\footnotesize G\allowbreak{}}\par
\noindent Corrected reading: {\ttfamily\footnotesize H\allowbreak{}}\par
\noindent Source: {\ttfamily\footnotesize fr/sga7ii/\{diplomati\allowbreak{}c,current\}/expose\_XV\allowbreak{}II\_body.tex;\allowbreak{} folio\allowbreak{} 220;\allowbreak{} formula\allowbreak{} 194;\allowbreak{} port/PORT\_EN7II\_LOG.\allowbreak{}jsonl\#C697;\allowbreak{} fixes/FIXLOG\_7II\_CUR\allowbreak{}RENT.jsonl\#C697;\allowbreak{} 7II:342\allowbreak{}}\par

\par\medskip\noindent\hypertarget{app-E7II-XVII-008.1}{}\textbf{E7II-XVII-008.1}\quad printed folio 222.\par
\noindent Printed reading: \emph{(empty)}\par
\noindent Corrected reading: {\ttfamily\footnotesize r-\allowbreak{}}\par
\noindent Source: {\ttfamily\footnotesize fr/sga7ii/\{diplomati\allowbreak{}c,current\}/expose\_XV\allowbreak{}II\_body.tex;\allowbreak{} folio\allowbreak{} 222;\allowbreak{} formula\allowbreak{} 276\allowbreak{}}\par

\par\medskip\noindent\hypertarget{app-E7II-XVII-008.2}{}\textbf{E7II-XVII-008.2}\quad printed folio 222.\par
\noindent Printed reading: {\ttfamily\footnotesize -r\allowbreak{}}\par
\noindent Corrected reading: \emph{(empty)}\par
\noindent Source: {\ttfamily\footnotesize fr/sga7ii/\{diplomati\allowbreak{}c,current\}/expose\_XV\allowbreak{}II\_body.tex;\allowbreak{} folio\allowbreak{} 222;\allowbreak{} formula\allowbreak{} 276\allowbreak{}}\par

\par\medskip\noindent\hypertarget{app-E7II-XVII-008.3}{}\textbf{E7II-XVII-008.3}\quad printed folio 222.\par
\noindent Printed reading: {\ttfamily\footnotesize \textbackslash{}lambda\allowbreak{}}\par
\noindent Corrected reading: {\ttfamily\footnotesize \textbackslash{}Lambda\allowbreak{}}\par
\noindent Source: {\ttfamily\footnotesize ED-EN-FORMULA-G135;\allowbreak{} fr/sga7ii/\{diplomati\allowbreak{}c,current\}/expose\_XV\allowbreak{}II\_body.tex;\allowbreak{} folio\allowbreak{} 222\allowbreak{}}\par

\par\medskip\noindent\hypertarget{app-E7II-XVII-008.4}{}\textbf{E7II-XVII-008.4}\quad printed folio 222.\par
\noindent Printed reading: {\ttfamily\footnotesize \textbackslash{}lambda\allowbreak{}}\par
\noindent Corrected reading: {\ttfamily\footnotesize \textbackslash{}Lambda\allowbreak{}}\par
\noindent Source: {\ttfamily\footnotesize ED-EN-FORMULA-G135;\allowbreak{} fr/sga7ii/\{diplomati\allowbreak{}c,current\}/expose\_XV\allowbreak{}II\_body.tex;\allowbreak{} folio\allowbreak{} 222\allowbreak{}}\par

\par\medskip\noindent\hypertarget{app-E7II-XVII-008.5}{}\textbf{E7II-XVII-008.5}\quad printed folio 222.\par
\noindent Printed reading: \emph{(empty)}\par
\noindent Corrected reading: {\ttfamily\footnotesize \{i=\allowbreak{}}\par
\noindent Source: {\ttfamily\footnotesize ED-EN-FORMULA-G135;\allowbreak{} fr/sga7ii/\{diplomati\allowbreak{}c,current\}/expose\_XV\allowbreak{}II\_body.tex;\allowbreak{} folio\allowbreak{} 222\allowbreak{}}\par

\par\medskip\noindent\hypertarget{app-E7II-XVII-008.6}{}\textbf{E7II-XVII-008.6}\quad printed folio 222.\par
\noindent Printed reading: \emph{(empty)}\par
\noindent Corrected reading: {\ttfamily\footnotesize \}\allowbreak{}}\par
\noindent Source: {\ttfamily\footnotesize ED-EN-FORMULA-G135;\allowbreak{} fr/sga7ii/\{diplomati\allowbreak{}c,current\}/expose\_XV\allowbreak{}II\_body.tex;\allowbreak{} folio\allowbreak{} 222\allowbreak{}}\par

\par\medskip\noindent\hypertarget{app-E7II-XVII-008.7}{}\textbf{E7II-XVII-008.7}\quad printed folio 222.\par
\noindent Printed reading: \emph{(empty)}\par
\noindent Corrected reading: {\ttfamily\footnotesize r-\allowbreak{}}\par
\noindent Source: {\ttfamily\footnotesize ED-EN-FORMULA-G135;\allowbreak{} fr/sga7ii/\{diplomati\allowbreak{}c,current\}/expose\_XV\allowbreak{}II\_body.tex;\allowbreak{} folio\allowbreak{} 222\allowbreak{}}\par

\par\medskip\noindent\hypertarget{app-E7II-XVII-008.8}{}\textbf{E7II-XVII-008.8}\quad printed folio 222.\par
\noindent Printed reading: {\ttfamily\footnotesize -r\allowbreak{}}\par
\noindent Corrected reading: \emph{(empty)}\par
\noindent Source: {\ttfamily\footnotesize ED-EN-FORMULA-G135;\allowbreak{} fr/sga7ii/\{diplomati\allowbreak{}c,current\}/expose\_XV\allowbreak{}II\_body.tex;\allowbreak{} folio\allowbreak{} 222\allowbreak{}}\par

\par\medskip\noindent\hypertarget{app-E7II-XVII-008.9}{}\textbf{E7II-XVII-008.9}\quad printed folio 222.\par
\noindent Printed reading: {\ttfamily\footnotesize \textbackslash{}lambda\allowbreak{}}\par
\noindent Corrected reading: {\ttfamily\footnotesize \textbackslash{}Lambda\allowbreak{}}\par
\noindent Source: {\ttfamily\footnotesize ED-EN-FORMULA-G135;\allowbreak{} fr/sga7ii/\{diplomati\allowbreak{}c,current\}/expose\_XV\allowbreak{}II\_body.tex;\allowbreak{} folio\allowbreak{} 222\allowbreak{}}\par

\par\medskip\noindent\hypertarget{app-E7II-XVII-008.10}{}\textbf{E7II-XVII-008.10}\quad printed folio 222.\par
\noindent Printed reading: {\ttfamily\footnotesize \{n+\allowbreak{}}\par
\noindent Corrected reading: \emph{(empty)}\par
\noindent Source: {\ttfamily\footnotesize ED-EN-FORMULA-G135;\allowbreak{} fr/sga7ii/\{diplomati\allowbreak{}c,current\}/expose\_XV\allowbreak{}II\_body.tex;\allowbreak{} folio\allowbreak{} 222\allowbreak{}}\par

\par\medskip\noindent\hypertarget{app-E7II-XVII-008.11}{}\textbf{E7II-XVII-008.11}\quad printed folio 222.\par
\noindent Printed reading: {\ttfamily\footnotesize \}\allowbreak{}}\par
\noindent Corrected reading: \emph{(empty)}\par
\noindent Source: {\ttfamily\footnotesize ED-EN-FORMULA-G135;\allowbreak{} fr/sga7ii/\{diplomati\allowbreak{}c,current\}/expose\_XV\allowbreak{}II\_body.tex;\allowbreak{} folio\allowbreak{} 222\allowbreak{}}\par

\par\medskip\noindent\hypertarget{app-E7II-XVII-008.12}{}\textbf{E7II-XVII-008.12}\quad printed folio 223.\par
\noindent Printed reading: \emph{(empty)}\par
\noindent Corrected reading: {\ttfamily\footnotesize \textbackslash{}left\allowbreak{}}\par
\noindent Source: {\ttfamily\footnotesize fr/sga7ii/\{diplomati\allowbreak{}c,current\}/expose\_XV\allowbreak{}II\_body.tex;\allowbreak{} folio\allowbreak{} 223\allowbreak{}}\par

\par\medskip\noindent\hypertarget{app-E7II-XVII-008.13}{}\textbf{E7II-XVII-008.13}\quad printed folio 223.\par
\noindent Printed reading: {\ttfamily\footnotesize (\allowbreak{}}\par
\noindent Corrected reading: {\ttfamily\footnotesize \textbackslash{}left[\allowbreak{}}\par
\noindent Source: {\ttfamily\footnotesize fr/sga7ii/\{diplomati\allowbreak{}c,current\}/expose\_XV\allowbreak{}II\_body.tex;\allowbreak{} folio\allowbreak{} 223\allowbreak{}}\par

\par\medskip\noindent\hypertarget{app-E7II-XVII-008.14}{}\textbf{E7II-XVII-008.14}\quad printed folio 223.\par
\noindent Printed reading: \emph{(empty)}\par
\noindent Corrected reading: {\ttfamily\footnotesize \{i=\allowbreak{}}\par
\noindent Source: {\ttfamily\footnotesize fr/sga7ii/\{diplomati\allowbreak{}c,current\}/expose\_XV\allowbreak{}II\_body.tex;\allowbreak{} folio\allowbreak{} 223\allowbreak{}}\par

\par\medskip\noindent\hypertarget{app-E7II-XVII-008.15}{}\textbf{E7II-XVII-008.15}\quad printed folio 223.\par
\noindent Printed reading: \emph{(empty)}\par
\noindent Corrected reading: {\ttfamily\footnotesize \}\allowbreak{}}\par
\noindent Source: {\ttfamily\footnotesize fr/sga7ii/\{diplomati\allowbreak{}c,current\}/expose\_XV\allowbreak{}II\_body.tex;\allowbreak{} folio\allowbreak{} 223\allowbreak{}}\par

\par\medskip\noindent\hypertarget{app-E7II-XVII-008.16}{}\textbf{E7II-XVII-008.16}\quad printed folio 223.\par
\noindent Printed reading: \emph{(empty)}\par
\noindent Corrected reading: {\ttfamily\footnotesize r-\allowbreak{}}\par
\noindent Source: {\ttfamily\footnotesize fr/sga7ii/\{diplomati\allowbreak{}c,current\}/expose\_XV\allowbreak{}II\_body.tex;\allowbreak{} folio\allowbreak{} 223\allowbreak{}}\par

\par\medskip\noindent\hypertarget{app-E7II-XVII-008.17}{}\textbf{E7II-XVII-008.17}\quad printed folio 223.\par
\noindent Printed reading: {\ttfamily\footnotesize -r\allowbreak{}}\par
\noindent Corrected reading: \emph{(empty)}\par
\noindent Source: {\ttfamily\footnotesize fr/sga7ii/\{diplomati\allowbreak{}c,current\}/expose\_XV\allowbreak{}II\_body.tex;\allowbreak{} folio\allowbreak{} 223\allowbreak{}}\par

\par\medskip\noindent\hypertarget{app-E7II-XVII-008.18}{}\textbf{E7II-XVII-008.18}\quad printed folio 223.\par
\noindent Printed reading: \emph{(empty)}\par
\noindent Corrected reading: {\ttfamily\footnotesize \textbackslash{}right]\textbackslash{}right\allowbreak{}}\par
\noindent Source: {\ttfamily\footnotesize fr/sga7ii/\{diplomati\allowbreak{}c,current\}/expose\_XV\allowbreak{}II\_body.tex;\allowbreak{} folio\allowbreak{} 223\allowbreak{}}\par

\par\medskip\noindent\hypertarget{app-E7II-XVII-008.19}{}\textbf{E7II-XVII-008.19}\quad printed folio 223.\par
\noindent Printed reading: \emph{(empty)}\par
\noindent Corrected reading: {\ttfamily\footnotesize r-\allowbreak{}}\par
\noindent Source: {\ttfamily\footnotesize fr/sga7ii/\{diplomati\allowbreak{}c,current\}/expose\_XV\allowbreak{}II\_body.tex;\allowbreak{} folio\allowbreak{} 223;\allowbreak{} formula\allowbreak{} 285\allowbreak{}}\par

\par\medskip\noindent\hypertarget{app-E7II-XVII-008.20}{}\textbf{E7II-XVII-008.20}\quad printed folio 223.\par
\noindent Printed reading: {\ttfamily\footnotesize -r\allowbreak{}}\par
\noindent Corrected reading: \emph{(empty)}\par
\noindent Source: {\ttfamily\footnotesize fr/sga7ii/\{diplomati\allowbreak{}c,current\}/expose\_XV\allowbreak{}II\_body.tex;\allowbreak{} folio\allowbreak{} 223;\allowbreak{} formula\allowbreak{} 285\allowbreak{}}\par

\par\medskip\noindent\hypertarget{app-E7II-XVII-008.21}{}\textbf{E7II-XVII-008.21}\quad printed folio 223.\par
\noindent Printed reading: {\ttfamily\footnotesize \textbackslash{}check\allowbreak{} }\par
\noindent Corrected reading: \emph{(empty)}\par
\noindent Source: {\ttfamily\footnotesize ED-EN-FORMULA-G137;\allowbreak{} fr/sga7ii/\{diplomati\allowbreak{}c,current\}/expose\_XV\allowbreak{}II\_body.tex;\allowbreak{} folio\allowbreak{} 223\allowbreak{}}\par

\par\medskip\noindent\hypertarget{app-E7II-XVII-008.22}{}\textbf{E7II-XVII-008.22}\quad printed folio 223.\par
\noindent Printed reading: {\ttfamily\footnotesize n+\allowbreak{}}\par
\noindent Corrected reading: \emph{(empty)}\par
\noindent Source: {\ttfamily\footnotesize fr/sga7ii/\{diplomati\allowbreak{}c,current\}/expose\_XV\allowbreak{}II\_body.tex;\allowbreak{} folio\allowbreak{} 223;\allowbreak{} formula\allowbreak{} 291\allowbreak{}}\par

\par\medskip\noindent\hypertarget{app-E7II-XVII-008.23}{}\textbf{E7II-XVII-008.23}\quad printed folio 223.\par
\noindent Printed reading: \emph{(empty)}\par
\noindent Corrected reading: {\ttfamily\footnotesize )\allowbreak{}}\par
\noindent Source: {\ttfamily\footnotesize fr/sga7ii/\{diplomati\allowbreak{}c,current\}/expose\_XV\allowbreak{}II\_body.tex;\allowbreak{} folio\allowbreak{} 223;\allowbreak{} formula\allowbreak{} 294\allowbreak{}}\par

\par\medskip\noindent\hypertarget{app-E7II-XVII-008.24}{}\textbf{E7II-XVII-008.24}\quad printed folio 223.\par
\noindent Printed reading: {\ttfamily\footnotesize i\allowbreak{}}\par
\noindent Corrected reading: {\ttfamily\footnotesize j\allowbreak{}}\par
\noindent Source: {\ttfamily\footnotesize fr/sga7ii/\{diplomati\allowbreak{}c,current\}/expose\_XV\allowbreak{}II\_body.tex;\allowbreak{} folio\allowbreak{} 223;\allowbreak{} formula\allowbreak{} 295\allowbreak{}}\par

\par\medskip\noindent\hypertarget{app-E7II-XVII-008.25}{}\textbf{E7II-XVII-008.25}\quad printed folio 223.\par
\noindent Printed reading: {\ttfamily\footnotesize i\allowbreak{}}\par
\noindent Corrected reading: {\ttfamily\footnotesize j\allowbreak{}}\par
\noindent Source: {\ttfamily\footnotesize fr/sga7ii/\{diplomati\allowbreak{}c,current\}/expose\_XV\allowbreak{}II\_body.tex;\allowbreak{} folio\allowbreak{} 223;\allowbreak{} formula\allowbreak{} 295\allowbreak{}}\par

\par\medskip\noindent\hypertarget{app-E7II-XVII-008.26}{}\textbf{E7II-XVII-008.26}\quad printed folio 223.\par
\noindent Printed reading: {\ttfamily\footnotesize i\allowbreak{}}\par
\noindent Corrected reading: {\ttfamily\footnotesize j\allowbreak{}}\par
\noindent Source: {\ttfamily\footnotesize fr/sga7ii/\{diplomati\allowbreak{}c,current\}/expose\_XV\allowbreak{}II\_body.tex;\allowbreak{} folio\allowbreak{} 223;\allowbreak{} formula\allowbreak{} 295\allowbreak{}}\par

\par\medskip\noindent\hypertarget{app-E7II-XVII-008.27}{}\textbf{E7II-XVII-008.27}\quad printed folio 223.\par
\noindent Printed reading: {\ttfamily\footnotesize i\allowbreak{}}\par
\noindent Corrected reading: {\ttfamily\footnotesize j\allowbreak{}}\par
\noindent Source: {\ttfamily\footnotesize fr/sga7ii/\{diplomati\allowbreak{}c,current\}/expose\_XV\allowbreak{}II\_body.tex;\allowbreak{} folio\allowbreak{} 223;\allowbreak{} formula\allowbreak{} 299\allowbreak{}}\par

\par\medskip\noindent\hypertarget{app-E7II-XVII-008.28}{}\textbf{E7II-XVII-008.28}\quad printed folio 223.\par
\noindent Printed reading: {\ttfamily\footnotesize ig\allowbreak{}}\par
\noindent Corrected reading: {\ttfamily\footnotesize jg\allowbreak{}}\par
\noindent Source: {\ttfamily\footnotesize fr/sga7ii/\{diplomati\allowbreak{}c,current\}/expose\_XV\allowbreak{}II\_body.tex;\allowbreak{} folio\allowbreak{} 223;\allowbreak{} formula\allowbreak{} 299\allowbreak{}}\par

\par\medskip\noindent\hypertarget{app-E7II-XVII-008.29}{}\textbf{E7II-XVII-008.29}\quad printed folio 223.\par
\noindent Printed reading: {\ttfamily\footnotesize i\allowbreak{}}\par
\noindent Corrected reading: {\ttfamily\footnotesize j\allowbreak{}}\par
\noindent Source: {\ttfamily\footnotesize fr/sga7ii/\{diplomati\allowbreak{}c,current\}/expose\_XV\allowbreak{}II\_body.tex;\allowbreak{} folio\allowbreak{} 223;\allowbreak{} formula\allowbreak{} 299\allowbreak{}}\par

\par\medskip\noindent\hypertarget{app-E7II-XVII-008.30}{}\textbf{E7II-XVII-008.30}\quad printed folio 223.\par
\noindent Printed reading: {\ttfamily\footnotesize 0H\allowbreak{}}\par
\noindent Corrected reading: {\ttfamily\footnotesize 0,H\allowbreak{}}\par
\noindent Source: {\ttfamily\footnotesize fr/sga7ii/\{diplomati\allowbreak{}c,current\}/expose\_XV\allowbreak{}II\_body.tex;\allowbreak{} folio\allowbreak{} 223;\allowbreak{} formula\allowbreak{} 299\allowbreak{}}\par

\par\medskip\noindent\hypertarget{app-E7II-XVII-008.31}{}\textbf{E7II-XVII-008.31}\quad printed folio 224.\par
\noindent Printed reading: {\ttfamily\footnotesize i\allowbreak{}}\par
\noindent Corrected reading: {\ttfamily\footnotesize j\allowbreak{}}\par
\noindent Source: {\ttfamily\footnotesize fr/sga7ii/\{diplomati\allowbreak{}c,current\}/expose\_XV\allowbreak{}II\_body.tex;\allowbreak{} folio\allowbreak{} 224;\allowbreak{} formula\allowbreak{} 301\allowbreak{}}\par

\par\medskip\noindent\hypertarget{app-E7II-XVII-008.32}{}\textbf{E7II-XVII-008.32}\quad printed folio 224.\par
\noindent Printed reading: {\ttfamily\footnotesize ih\allowbreak{}}\par
\noindent Corrected reading: {\ttfamily\footnotesize jh\allowbreak{}}\par
\noindent Source: {\ttfamily\footnotesize fr/sga7ii/\{diplomati\allowbreak{}c,current\}/expose\_XV\allowbreak{}II\_body.tex;\allowbreak{} folio\allowbreak{} 224;\allowbreak{} formula\allowbreak{} 301\allowbreak{}}\par

\par\medskip\noindent\hypertarget{app-E7II-XVII-008.33}{}\textbf{E7II-XVII-008.33}\quad printed folio 224.\par
\noindent Printed reading: {\ttfamily\footnotesize i\allowbreak{}}\par
\noindent Corrected reading: {\ttfamily\footnotesize j\allowbreak{}}\par
\noindent Source: {\ttfamily\footnotesize fr/sga7ii/\{diplomati\allowbreak{}c,current\}/expose\_XV\allowbreak{}II\_body.tex;\allowbreak{} folio\allowbreak{} 224;\allowbreak{} formula\allowbreak{} 301\allowbreak{}}\par

\par\medskip\noindent\hypertarget{app-E7II-XVII-008.34}{}\textbf{E7II-XVII-008.34}\quad printed folio 224.\par
\noindent Printed reading: {\ttfamily\footnotesize i\allowbreak{}}\par
\noindent Corrected reading: {\ttfamily\footnotesize n+j\allowbreak{}}\par
\noindent Source: {\ttfamily\footnotesize fr/sga7ii/\{diplomati\allowbreak{}c,current\}/expose\_XV\allowbreak{}II\_body.tex;\allowbreak{} folio\allowbreak{} 224;\allowbreak{} formula\allowbreak{} 302\allowbreak{}}\par

\par\medskip\noindent\hypertarget{app-E7II-XVII-008.35}{}\textbf{E7II-XVII-008.35}\quad printed folio 224.\par
\noindent Printed reading: {\ttfamily\footnotesize i\allowbreak{}}\par
\noindent Corrected reading: {\ttfamily\footnotesize j\allowbreak{}}\par
\noindent Source: {\ttfamily\footnotesize fr/sga7ii/\{diplomati\allowbreak{}c,current\}/expose\_XV\allowbreak{}II\_body.tex;\allowbreak{} folio\allowbreak{} 224;\allowbreak{} formula\allowbreak{} 302\allowbreak{}}\par

\par\medskip\noindent\hypertarget{app-E7II-XVII-008.36}{}\textbf{E7II-XVII-008.36}\quad printed folio 224.\par
\noindent Printed reading: {\ttfamily\footnotesize i\allowbreak{}}\par
\noindent Corrected reading: {\ttfamily\footnotesize j\allowbreak{}}\par
\noindent Source: {\ttfamily\footnotesize fr/sga7ii/\{diplomati\allowbreak{}c,current\}/expose\_XV\allowbreak{}II\_body.tex;\allowbreak{} folio\allowbreak{} 224;\allowbreak{} formula\allowbreak{} 302\allowbreak{}}\par

\par\medskip\noindent\hypertarget{app-E7II-XVII-008.37}{}\textbf{E7II-XVII-008.37}\quad printed folio 224.\par
\noindent Printed reading: \emph{(empty)}\par
\noindent Corrected reading: {\ttfamily\footnotesize r-\allowbreak{}}\par
\noindent Source: {\ttfamily\footnotesize fr/sga7ii/\{diplomati\allowbreak{}c,current\}/expose\_XV\allowbreak{}II\_body.tex;\allowbreak{} folio\allowbreak{} 224;\allowbreak{} formula\allowbreak{} 302\allowbreak{}}\par

\par\medskip\noindent\hypertarget{app-E7II-XVII-008.38}{}\textbf{E7II-XVII-008.38}\quad printed folio 224.\par
\noindent Printed reading: \emph{(empty)}\par
\noindent Corrected reading: {\ttfamily\footnotesize -1\allowbreak{}}\par
\noindent Source: {\ttfamily\footnotesize fr/sga7ii/\{diplomati\allowbreak{}c,current\}/expose\_XV\allowbreak{}II\_body.tex;\allowbreak{} folio\allowbreak{} 224;\allowbreak{} formula\allowbreak{} 302\allowbreak{}}\par

\par\medskip\noindent\hypertarget{app-E7II-XVII-009.1}{}\textbf{E7II-XVII-009.1}\quad printed folio 227.\par
\noindent Printed reading: \emph{(empty)}\par
\noindent Corrected reading: {\ttfamily\footnotesize \textbackslash{}check\allowbreak{}}\par
\noindent Source: {\ttfamily\footnotesize fr/sga7ii/\{diplomati\allowbreak{}c,current\}/expose\_XV\allowbreak{}II\_body.tex;\allowbreak{} folio\allowbreak{} 227;\allowbreak{} formula\allowbreak{} 343\allowbreak{}}\par

\par\medskip\noindent\hypertarget{app-E7II-XVII-010.1}{}\textbf{E7II-XVII-010.1}\quad printed folio 235.\par
\noindent Printed reading: \emph{(empty)}\par
\noindent Corrected reading: {\ttfamily\footnotesize \textbackslash{}check\allowbreak{}}\par
\noindent Source: {\ttfamily\footnotesize fr/sga7ii/\{diplomati\allowbreak{}c,current\}/expose\_XV\allowbreak{}II\_body.tex;\allowbreak{} folio\allowbreak{} 235;\allowbreak{} formula\allowbreak{} 529\allowbreak{}}\par

\par\medskip\noindent\hypertarget{app-E7II-XVII-011.1}{}\textbf{E7II-XVII-011.1}\quad printed folio 236.\par
\noindent Printed reading: \emph{(empty)}\par
\noindent Corrected reading: {\ttfamily\footnotesize \textbackslash{}overline\{\allowbreak{}}\par
\noindent Source: {\ttfamily\footnotesize fr/sga7ii/\{diplomati\allowbreak{}c,current\}/expose\_XV\allowbreak{}II\_body.tex;\allowbreak{} folio\allowbreak{} 236\allowbreak{}}\par

\par\medskip\noindent\hypertarget{app-E7II-XVII-011.2}{}\textbf{E7II-XVII-011.2}\quad printed folio 236.\par
\noindent Printed reading: \emph{(empty)}\par
\noindent Corrected reading: {\ttfamily\footnotesize \}\allowbreak{}}\par
\noindent Source: {\ttfamily\footnotesize fr/sga7ii/\{diplomati\allowbreak{}c,current\}/expose\_XV\allowbreak{}II\_body.tex;\allowbreak{} folio\allowbreak{} 236\allowbreak{}}\par

\par\medskip\noindent\hypertarget{app-E7II-XVII-011.3}{}\textbf{E7II-XVII-011.3}\quad printed folio 236.\par
\noindent Printed reading: {\ttfamily\footnotesize 8\allowbreak{}}\par
\noindent Corrected reading: {\ttfamily\footnotesize 7\allowbreak{}}\par
\noindent Source: {\ttfamily\footnotesize ED-EN-NONMATH-G026;\allowbreak{} F7II-XVII-027.6;\allowbreak{} 7II:375;\allowbreak{} fr/sga7ii/\{diplomati\allowbreak{}c,current\}/expose\_XV\allowbreak{}II\_body.tex;\allowbreak{} folio\allowbreak{} 236\allowbreak{}}\par

\par\medskip\noindent\hypertarget{app-E7II-XVII-011.4}{}\textbf{E7II-XVII-011.4}\quad printed folio 236.\par
\noindent Printed reading: {\ttfamily\footnotesize 2\allowbreak{}}\par
\noindent Corrected reading: {\ttfamily\footnotesize 1\allowbreak{}}\par
\noindent Source: {\ttfamily\footnotesize fr/sga7ii/\{diplomati\allowbreak{}c,current\}/expose\_XV\allowbreak{}II\_body.tex;\allowbreak{} folio\allowbreak{} 236;\allowbreak{} formula\allowbreak{} 555\allowbreak{}}\par

\par\medskip\noindent\hypertarget{app-E7II-XVII-012.1}{}\textbf{E7II-XVII-012.1}\quad printed folio 236.\par
\noindent Printed reading: {\ttfamily\footnotesize n\allowbreak{}}\par
\noindent Corrected reading: {\ttfamily\footnotesize r\allowbreak{}}\par
\noindent Source: {\ttfamily\footnotesize ED-EN-FORMULA-G139;\allowbreak{} fr/sga7ii/\{diplomati\allowbreak{}c,current\}/expose\_XV\allowbreak{}II\_body.tex;\allowbreak{} folio\allowbreak{} 236\allowbreak{}}\par

\par\medskip\noindent\hypertarget{app-E7II-XVII-013.1}{}\textbf{E7II-XVII-013.1}\quad printed folio 237.\par
\noindent Printed reading: \emph{(empty)}\par
\noindent Corrected reading: {\ttfamily\footnotesize =\allowbreak{}}\par
\noindent Source: {\ttfamily\footnotesize fr/sga7ii/\{diplomati\allowbreak{}c,current\}/expose\_XV\allowbreak{}II\_body.tex;\allowbreak{} folio\allowbreak{} 237\allowbreak{}}\par

\par\medskip\noindent\hypertarget{app-E7II-XVII-014.1}{}\textbf{E7II-XVII-014.1}\quad printed folio 238.\par
\noindent Printed reading: {\ttfamily\footnotesize d\allowbreak{}}\par
\noindent Corrected reading: {\ttfamily\footnotesize 1\allowbreak{}}\par
\noindent Source: {\ttfamily\footnotesize fr/sga7ii/\{diplomati\allowbreak{}c,current\}/expose\_XV\allowbreak{}II\_body.tex;\allowbreak{} folio\allowbreak{} 238;\allowbreak{} formula\allowbreak{} 623\allowbreak{}}\par

\par\medskip\noindent\hypertarget{app-E7II-XVII-014.2}{}\textbf{E7II-XVII-014.2}\quad printed folio 239.\par
\noindent Printed reading: {\ttfamily\footnotesize 1\allowbreak{}}\par
\noindent Corrected reading: {\ttfamily\footnotesize 2\allowbreak{}}\par
\noindent Source: {\ttfamily\footnotesize fr/sga7ii/\{diplomati\allowbreak{}c,current\}/expose\_XV\allowbreak{}II\_body.tex;\allowbreak{} folio\allowbreak{} 239\allowbreak{}}\par

\par\medskip\noindent\hypertarget{app-E7II-XVII-015.1}{}\textbf{E7II-XVII-015.1}\quad printed folio 240.\par
\noindent Printed reading: {\ttfamily\footnotesize ij\allowbreak{}}\par
\noindent Corrected reading: {\ttfamily\footnotesize 1j\allowbreak{}}\par
\noindent Source: {\ttfamily\footnotesize fr/sga7ii/\{diplomati\allowbreak{}c,current\}/expose\_XV\allowbreak{}II\_body.tex;\allowbreak{} folio\allowbreak{} 240;\allowbreak{} formula\allowbreak{} 670\allowbreak{}}\par

\par\medskip\noindent\hypertarget{app-E7II-XVII-016.1}{}\textbf{E7II-XVII-016.1}\quad printed folio 241.\par
\noindent Printed reading: {\ttfamily\footnotesize \{\textbackslash{}mathbb\allowbreak{} P\}\allowbreak{}}\par
\noindent Corrected reading: {\ttfamily\footnotesize X\allowbreak{}}\par
\noindent Source: {\ttfamily\footnotesize ED-EN-FORMULA-G141;\allowbreak{} fr/sga7ii/\{diplomati\allowbreak{}c,current\}/expose\_XV\allowbreak{}II\_body.tex;\allowbreak{} folio\allowbreak{} 241\allowbreak{}}\par

\par\medskip\noindent\hypertarget{app-E7II-XVII-017.1}{}\textbf{E7II-XVII-017.1}\quad printed folio 243.\par
\noindent Printed reading: \emph{(empty)}\par
\noindent Corrected reading: {\ttfamily\footnotesize \textbackslash{}check\allowbreak{}}\par
\noindent Source: {\ttfamily\footnotesize fr/sga7ii/\{diplomati\allowbreak{}c,current\}/expose\_XV\allowbreak{}II\_body.tex;\allowbreak{} folio\allowbreak{} 243;\allowbreak{} formula\allowbreak{} 728\allowbreak{}}\par

\par\medskip\noindent\hypertarget{app-E7II-XVII-018.1}{}\textbf{E7II-XVII-018.1}\quad printed folio 243.\par
\noindent Printed reading: \emph{(empty)}\par
\noindent Corrected reading: {\ttfamily\footnotesize (1)\allowbreak{}}\par
\noindent Source: {\ttfamily\footnotesize fr/sga7ii/\{diplomati\allowbreak{}c,current\}/expose\_XV\allowbreak{}II\_body.tex;\allowbreak{} folio\allowbreak{} 243;\allowbreak{} formula\allowbreak{} 740;\allowbreak{} port/decisions/decis\allowbreak{}ions\_en7II\_r2\_002.js\allowbreak{}on\#C730;\allowbreak{} fixes/FIXLOG\_7II\_CUR\allowbreak{}RENT.jsonl\#C730;\allowbreak{} 7II:398\allowbreak{}}\par

\par\medskip\noindent\hypertarget{app-E7II-XVII-018.2}{}\textbf{E7II-XVII-018.2}\quad printed folio 244.\par
\noindent Printed reading: \emph{(empty)}\par
\noindent Corrected reading: {\ttfamily\footnotesize \$'\$\allowbreak{}}\par
\noindent Source: {\ttfamily\footnotesize ED-EN-NORM-G032;\allowbreak{} complete\allowbreak{} normalized-equal\allowbreak{} formula\allowbreak{} audit;\allowbreak{} fr/sga7ii/\{diplomati\allowbreak{}c,current\}/expose\_XV\allowbreak{}II\_body.tex;\allowbreak{} folio\allowbreak{} 244\allowbreak{}}\par

\par\medskip\noindent\hypertarget{app-E7II-XVII-019.1}{}\textbf{E7II-XVII-019.1}\quad printed folio 247.\par
\noindent Printed reading: {\ttfamily\footnotesize j\allowbreak{}}\par
\noindent Corrected reading: {\ttfamily\footnotesize i\allowbreak{}}\par
\noindent Source: {\ttfamily\footnotesize fr/sga7ii/\{diplomati\allowbreak{}c,current\}/expose\_XV\allowbreak{}II\_body.tex;\allowbreak{} folio\allowbreak{} 247\allowbreak{}}\par

\par\medskip\noindent\hypertarget{app-E7II-XVII-019.2}{}\textbf{E7II-XVII-019.2}\quad printed folio 247.\par
\noindent Printed reading: {\ttfamily\footnotesize j\allowbreak{}}\par
\noindent Corrected reading: {\ttfamily\footnotesize a\allowbreak{}}\par
\noindent Source: {\ttfamily\footnotesize fr/sga7ii/\{diplomati\allowbreak{}c,current\}/expose\_XV\allowbreak{}II\_body.tex;\allowbreak{} folio\allowbreak{} 247\allowbreak{}}\par

\par\medskip\noindent\hypertarget{app-E7II-XVII-019.3}{}\textbf{E7II-XVII-019.3}\quad printed folio 247.\par
\noindent Printed reading: {\ttfamily\footnotesize i\allowbreak{}}\par
\noindent Corrected reading: {\ttfamily\footnotesize 1\allowbreak{}}\par
\noindent Source: {\ttfamily\footnotesize fr/sga7ii/\{diplomati\allowbreak{}c,current\}/expose\_XV\allowbreak{}II\_body.tex;\allowbreak{} folio\allowbreak{} 247\allowbreak{}}\par

\par\medskip\noindent\hypertarget{app-E7II-XVII-020.1}{}\textbf{E7II-XVII-020.1}\quad printed folio 247.\par
\noindent Printed reading: \emph{(empty)}\par
\noindent Corrected reading: {\ttfamily\footnotesize \textbackslash{}mathbb\allowbreak{} }\par
\noindent Source: {\ttfamily\footnotesize ED-EN-FONT-G006;\allowbreak{} F7II-XVII-042.5;\allowbreak{} fr/sga7ii/\{diplomati\allowbreak{}c,current\}/expose\_XV\allowbreak{}II\_body.tex;\allowbreak{} folio\allowbreak{} 247\allowbreak{}}\par

\par\medskip\noindent\hypertarget{app-E7II-XVII-020.2}{}\textbf{E7II-XVII-020.2}\quad printed folio 249.\par
\noindent Printed reading: \emph{(empty)}\par
\noindent Corrected reading: {\ttfamily\footnotesize \textasciicircum{}\{a\}\allowbreak{}}\par
\noindent Source: {\ttfamily\footnotesize ED-EN-FORMULA-G143;\allowbreak{} fr/sga7ii/\{diplomati\allowbreak{}c,current\}/expose\_XV\allowbreak{}II\_body.tex;\allowbreak{} folio\allowbreak{} 249\allowbreak{}}\par

\par\medskip\noindent\hypertarget{app-E7II-XVII-021.1}{}\textbf{E7II-XVII-021.1}\quad printed folio 250.\par
\noindent Printed reading: {\ttfamily\footnotesize r\allowbreak{}}\par
\noindent Corrected reading: {\ttfamily\footnotesize n\allowbreak{}}\par
\noindent Source: {\ttfamily\footnotesize fr/sga7ii/\{diplomati\allowbreak{}c,current\}/expose\_XV\allowbreak{}II\_body.tex;\allowbreak{} folio\allowbreak{} 250\allowbreak{}}\par

\par\medskip\noindent\hypertarget{app-E7II-XVII-022.1}{}\textbf{E7II-XVII-022.1}\quad printed folio 253.\par
\noindent Printed reading: {\ttfamily\footnotesize 3\allowbreak{}}\par
\noindent Corrected reading: {\ttfamily\footnotesize 4\allowbreak{}}\par
\noindent Source: {\ttfamily\footnotesize port/PORT\_EN7II\_LOG.\allowbreak{}jsonl\#C738;\allowbreak{} fixes/FIXLOG\_7II\_CUR\allowbreak{}RENT.jsonl\#C738;\allowbreak{} 7II:410\allowbreak{}}\par

\par\medskip\noindent\hypertarget{app-E7II-XVII-022.2}{}\textbf{E7II-XVII-022.2}\quad printed folio 253.\par
\noindent Printed reading: {\ttfamily\footnotesize S\allowbreak{}}\par
\noindent Corrected reading: {\ttfamily\footnotesize s\allowbreak{}}\par
\noindent Source: {\ttfamily\footnotesize fr/sga7ii/\{diplomati\allowbreak{}c,current\}/expose\_XV\allowbreak{}II\_body.tex;\allowbreak{} folio\allowbreak{} 253;\allowbreak{} formula\allowbreak{} 952;\allowbreak{} port/decisions/decis\allowbreak{}ions\_en7II\_r2\_002.js\allowbreak{}on\#C739;\allowbreak{} fixes/FIXLOG\_7II\_CUR\allowbreak{}RENT.jsonl\#C739;\allowbreak{} 7II:411\allowbreak{}}\par

\par\endgroup\normalsize\fi
\clearpage\SGAtoc{Exposé XVIII}
\selectlanguage{english}
\clearpage
\SGApage{II}{254}

\begin{center}
{\large SGA 7 --- EXPOS\'E XVIII}\par
\vspace{2em}
{\Large\underline{Cohomological Study of Lefschetz Pencils}}\par
\vspace{1.5em}
{\large by N. Katz\footnote{Based on succinct notes of Grothendieck.}}
\end{center}

\section*{Contents}

\noindent 1. The cohomology of a projective bundle. \hfill 2

\noindent 2. The cohomology of blown-up varieties. \hfill 4

\noindent 3. The blow-up associated with a pencil. \hfill 13

\noindent 4. The cohomology of the blow-up associated with a pencil. \hfill 17

\noindent 5. Cohomological study of the morphism
$\rho:\widetilde X\longrightarrow\mathbb P^1$ associated with a pencil.
\hfill 24

\noindent\quad 5.1. The restriction homomorphism
$H^*(\widetilde X)\longrightarrow H^*(Y)$. \hfill 24

\noindent\quad 5.2. The ``Lefschetz vache theorem'' (LV). \hfill 27

\noindent\quad 5.3. Condition (A). \hfill 33

\noindent\quad 5.8. The Griffiths homomorphism. \hfill 47

\noindent\quad 5.9. An example in which hypothesis (A) is not satisfied.
\hfill 57

\noindent 6. Applications of the local theory of vanishing cycles
(Picard--Lefschetz). \hfill 58

\section*{Introduction}

In the present expos\'e we give the elementary cohomological structure of
Lefschetz pencils. Sections 1--4 give the ``well-known'' facts concerning the
cohomological structure of projective bundles and blow-ups. In Section 5 we
carry out the study named in the title, assuming that $X$ satisfies the
``Lefschetz vache theorem'' (5.2) and that the pencil $D$ in question satisfies
a certain condition (A), local in nature on $D$ (5.3). The principal results
are 5.6, 5.7, and 5.8.7; the last plays an essential role in the proof of
Griffiths's theorem in Expos\'e XX. In the final section, which is distinctly
less elementary, we use Mme Raynaud's Expos\'e XIII of SGA 1 on the
\underline{tame fundamental group}, and the \underline{Picard--Lefschetz
formula} (XVI), to prove:
\begin{enumerate}[label=(\arabic*)]
\item condition (A) is often satisfied (6.3, 6.4);
\item except when $\dim(X)$ is odd and the characteristic is $2$, the
monodromy group of a Lefschetz pencil acts irreducibly on the ``vanishing
cohomology,'' and transitively (up to sign) on the set of ``local vanishing
cycles'' associated, via the Picard--Lefschetz formula, with the various
``critical points'' of the pencil (6.6 and 6.7).
\end{enumerate}

\SGApage{II}{\(\approx\)256}
\section*{1. The cohomology of a projective bundle}

\SGAvar{\subsubsection*{Theorem 1.1 (Deligne).}
\underline{Let $f:X\longrightarrow S$ be a proper morphism, let $\ell$ be a
prime number invertible on $S$, let
$u\in H^2(X,\mathbb Z_\ell(1))$ (respectively,
$u\in H^2(X,\mathbb Q_\ell(1))$ be a cohomology class on $X$, and let
$\mathfrak F$ be a} $\mathbb Z_\ell$-\underline{sheaf} (respectively, a
$\mathbb Q_\ell$-\underline{sheaf}) \underline{on $X$. Suppose that there is an integer $n$ such
that, for every integer $j$ and every integer $i>0$, the map}
\begin{equation}\tag{1.1.1}
R^{n-i}f_*\bigl(\mathfrak F(j)\bigr)
\xrightarrow{\ \wedge u^i\ }
R^{n+i}f_*\bigl(\mathfrak F(j+i)\bigr)
\end{equation}
\underline{is an isomorphism of} $\mathbb Z_\ell$-\underline{sheaves} (respectively,
$\mathbb Q_\ell$-\underline{sheaves}) \underline{on $S$. Then, for every morphism
$g:S\longrightarrow T$, the Leray spectral sequence}
\[
E_2^{p,q}=R^pg_*R^qf_*(\mathfrak F)
\Longrightarrow R^{p+q}(gf)_*(\mathfrak F)
\]
\underline{degenerates.}

}{\subsubsection*{Theorem 1.1 (Deligne).}
\underline{Let $f:X\longrightarrow S$ be a proper morphism, let $\ell$ be a
prime number invertible on $S$, let
$u\in H^2(X,\mathbb Z_\ell(1))$ (respectively,
$u\in H^2(X,\mathbb Q_\ell(1))$) be a cohomology class on $X$, and let
$\mathfrak F$ be a} $\mathbb Z_\ell$-\underline{sheaf} (respectively, a
$\mathbb Q_\ell$-\underline{sheaf}) \underline{on $X$. Suppose that there is an integer $n$ such
that, for every integer $j$ and every integer $i>0$, the map}
\begin{equation}\tag{1.1.1}
R^{n-i}f_*\bigl(\mathfrak F(j)\bigr)
\xrightarrow{\ \wedge u^i\ }
R^{n+i}f_*\bigl(\mathfrak F(j+i)\bigr)
\end{equation}
\underline{is an isomorphism of} $\mathbb Z_\ell$-\underline{sheaves} (respectively,
$\mathbb Q_\ell$-\underline{sheaves}) \underline{on $S$. Then, for every morphism
$g:S\longrightarrow T$, the Leray spectral sequence}
\[
E_2^{p,q}=R^pg_*R^qf_*(\mathfrak F)
\Longrightarrow R^{p+q}(gf)_*(\mathfrak F)
\]
\underline{degenerates.}

}{E7II-XVIII-001}The proof is found in [1, Prop.~2.4]. Notice that, by the proper base-change
theorem (SGA 4 XII, 5.2), the hypothesis can be checked fibre by fibre.

\SGApage{II}{257}
\SGAvar{\subsubsection*{Theorem 1.2 (cf. SGA 5 VII, 2.2.1).}
\underline{Let $\mathcal E\longrightarrow X$ be a vector bundle, locally
free of rank $1+r$. Let $f:\mathbb P(\mathcal E)\longrightarrow X$ denote
the projective bundle defined by $\mathcal E$, and let
$u\in H^2(\mathbb P(\mathcal E),\mathbb Z_\ell(1))$ be the first Chern
class of the canonical invertible sheaf
$\mathcal O_{\mathbb P(\mathcal E)}(1)$. For $\ell$ a prime number
invertible on $X$, one has:}
\begin{enumerate}[label=\alph*)]
\item \underline{the spectral sequence}
\begin{equation}\tag{1.2.1}
E_2^{p,q}=H^p\bigl(X,R^qf_*\mathbb Z_\ell\bigr)
\Longrightarrow
H^{p+q}\bigl(\mathbb P(\mathcal E),\mathbb Z_\ell\bigr)
\end{equation}
\underline{degenerates;}
\item \underline{the map}
\begin{equation}\tag{1.2.2}
\bigoplus_{0\leq j\leq n}
H^{q-2-j}\bigl(X,\mathbb Z_\ell(-j)
\longrightarrow
H^q\bigl(\mathbb P(\mathcal E),\mathbb Z_\ell\bigr),
\qquad
(\tau_{q-2j})_j\longmapsto
\sum_j f^*(\tau_{q-2j})\wedge u^j
\end{equation}
\underline{is an isomorphism.}
\end{enumerate}


\noindent\underline{Proof.}
The fibres of $f$ are projective spaces of dimension $r$, and one knows that
\begin{equation}\tag{1.2.3}
R^if_*\mathbb Z_\ell
\xleftarrow[\ \sim\ ]{\ \wedge^{i/2}u\ }
\begin{cases}
\mathbb Z_\ell\left(-\dfrac{i}{2}\right),
  & \text{if $i$ is even and $0\leq i\leq2r$},\\[2ex]
0, & \text{otherwise}.
\end{cases}
\end{equation}
Theorem 1.1 applies to $(f,\ell,u,\mathbb Z_\ell,r)$, which proves (a).
By (1.2.3),
\begin{equation}\tag{1.2.4}
E_2^{a,b}=
\begin{cases}
H^a\left(X,\mathbb Z_\ell\left(-\dfrac b2\right)\right),
  & \text{if $b$ is even and $0\leq b\leq2r$},\\[2ex]
0, & \text{otherwise}.
\end{cases}
\end{equation}
Thus the graded object associated with
$H^q(\mathbb P(\mathcal E),\mathbb Z_\ell)$ for the filtration coming
from the spectral sequence is
\[
\bigoplus_{0\leq j\leq r}
H^{q-2j}\bigl(X,\mathbb Z_\ell(-j)\bigr).
\]
If this direct sum is equipped with the filtration
\[
F^i=
\bigoplus_{\substack{0\leq j\leq r\\q-2j\leq i}}
H^{q-2j}\bigl(X,\mathbb Z_\ell(-j)\bigr),
\]
the map (1.2.2) is compatible with the filtrations and induces the identity on
the associated graded objects. This proves (b).



}{\subsubsection*{Theorem 1.2 (cf. SGA 5 VII, 2.2.1).}
\underline{Let $\mathcal E\longrightarrow X$ be a vector bundle, locally
free of rank $1+r$. Let $f:\mathbb P(\mathcal E)\longrightarrow X$ denote
the projective bundle defined by $\mathcal E$, and let
$u\in H^2(\mathbb P(\mathcal E),\mathbb Z_\ell(1))$ be the first Chern
class of the canonical invertible sheaf
$\mathcal O_{\mathbb P(\mathcal E)}(1)$. For $\ell$ a prime number
invertible on $X$, one has:}
\begin{enumerate}[label=\alph*)]
\item \underline{the spectral sequence}
\begin{equation}\tag{1.2.1}
E_2^{p,q}=H^p\bigl(X,R^qf_*\mathbb Z_\ell\bigr)
\Longrightarrow
H^{p+q}\bigl(\mathbb P(\mathcal E),\mathbb Z_\ell\bigr)
\end{equation}
\underline{degenerates;}
\item \underline{the map}
\begin{equation}\tag{1.2.2}
\bigoplus_{0\leq j\leq r}
H^{q-2j}\bigl(X,\mathbb Z_\ell(-j)\bigr)
\longrightarrow
H^q\bigl(\mathbb P(\mathcal E),\mathbb Z_\ell\bigr),
\qquad
(\tau_{q-2j})_j\longmapsto
\sum_j f^*(\tau_{q-2j})\wedge u^j
\end{equation}
\underline{is an isomorphism.}
\end{enumerate}


\noindent\underline{Proof.}
The fibres of $f$ are projective spaces of dimension $r$, and one knows that
\begin{equation}\tag{1.2.3}
R^if_*\mathbb Z_\ell
\xleftarrow[\ \sim\ ]{\ \wedge^{i/2}u\ }
\begin{cases}
\mathbb Z_\ell\left(-\dfrac{i}{2}\right),
  & \text{if $i$ is even and $0\leq i\leq2r$},\\[2ex]
0, & \text{otherwise}.
\end{cases}
\end{equation}
Theorem 1.1 applies to $(f,\ell,u,\mathbb Z_\ell,r)$, which proves (a).
By (1.2.3),
\begin{equation}\tag{1.2.4}
E_2^{a,b}=
\begin{cases}
H^a\left(X,\mathbb Z_\ell\left(-\dfrac b2\right)\right),
  & \text{if $b$ is even and $0\leq b\leq2r$},\\[2ex]
0, & \text{otherwise}.
\end{cases}
\end{equation}
Thus the graded object associated with
$H^q(\mathbb P(\mathcal E),\mathbb Z_\ell)$ for the filtration coming
from the spectral sequence is
\[
\bigoplus_{0\leq j\leq r}
H^{q-2j}\bigl(X,\mathbb Z_\ell(-j)\bigr).
\]
If this direct sum is equipped with the filtration
\[
F^i=
\bigoplus_{\substack{0\leq j\leq r\\q-2j\geq i}}
H^{q-2j}\bigl(X,\mathbb Z_\ell(-j)\bigr),
\]
the map (1.2.2) is compatible with the filtrations and induces the identity on
the associated graded objects. This proves (b).



}{E7II-XVIII-002}\SGApage{II}{\(\approx\)258}
\section*{2. The cohomology of blown-up varieties (cf. SGA 5 VII, 8)}

\subsubsection*{2.1.}
Let $X$ be a smooth scheme over a field $k$, and let
$\Delta\subset X$ be an everywhere smooth subscheme of codimension $1+r$.
It is known (EGA IV, 19.4.9) that there is a \underline{smooth} scheme
$\widetilde X$, equipped with a projective morphism
\begin{equation}\tag{2.1.1}
f:\widetilde X\longrightarrow X,
\end{equation}
such that, on setting
\begin{equation}\tag{2.1.2}
\widetilde\Delta=f^{-1}(\Delta)
=\Delta\times_X\widetilde X,
\end{equation}
$f$ induces an isomorphism
\begin{equation}\tag{2.1.3}
f|_{\widetilde X-\widetilde\Delta}:
\widetilde X-\widetilde\Delta
\xrightarrow{\ \sim\ }X-\Delta,
\end{equation}
and there is a $\Delta$-isomorphism between $\widetilde\Delta$ and
$\check{\mathbb P}(\underline N_{\Delta/X})$, the projective bundle of
lines in the normal bundle $\underline N_{\Delta/X}$ of $\Delta$ in $X$.
Thus one has a commutative diagram
\begin{equation}\tag{2.1.4}
\begin{tikzcd}[row sep=huge,column sep=small]
\widetilde\Delta
  \arrow[rd,"f|_{\widetilde\Delta}"']
  \& \simeq \&
\check{\mathbb P}(\underline N_{\Delta|X})
  \arrow[ld,"\text{canonical projection}"]\\
{} \& \Delta \& {}
\end{tikzcd}.
\end{equation}

\SGApage{II}{\(\approx\)259}
\SGAvar{\subsubsection*{2.1.5.}
In the special case
\[
X=\operatorname{Spec}(A),\qquad
\Delta=\operatorname{Spec}(A/I),\qquad
I=(g_o,\ldots,g_r),
\]
$\widetilde X$ is obtained by the Hopf construction (cf. EGA IV, 19.4.11).
One defines a morphism from $X-\Delta$ to $\mathbb P^r$ by
\[
x\in X-\Delta\longmapsto
\bigl(g_o(x),\ldots,g_r(x)\in\mathbb P^r
\]
and regards its graph $\Gamma$ as embedded in $X\times\mathbb P^r$. Then
$\widetilde X$ is defined to be the schematic closure of $\Gamma$ in
$X\times\mathbb P^r$, and the morphism
$f:\widetilde X\longrightarrow X$ is the composite
\[
\widetilde X\lhook\joinrel\longrightarrow
X\times\mathbb P^r\xrightarrow{\ \operatorname{pr}_1\ }X.
\]

\subsubsection*{2.1.6.}
Similarly (cf. again EGA IV, 19.4.11), if $X$ is projective and $\Delta$ is
the intersection of $X$ with $1+r$ hypersurfaces, say with equations
$G_i=0$, $i=0,\ldots,r$, one again considers the morphism
\[
X-\Delta\longrightarrow\mathbb P^r,\qquad
x\longmapsto\bigl(G_o(x),\ldots,G_r(x)\bigr),
\]
and obtains $\widetilde X$ as the closure of its graph in
$X\times\mathbb P^r$. Concretely, in terms of homogeneous coordinates
$\lambda_o,\ldots,\lambda_r$ on $\mathbb P^r$, $\widetilde X$ is the
subscheme of $X\times\mathbb P^r$ defined by the equations
\begin{equation}\tag{2.1.6.1}
\lambda_iG_j(X)-\lambda_jG_i(X),
\qquad 0\leq i<j\leq r.
\end{equation}
These equations show that the point $(x,\lambda)$ lies in
$\widetilde\Delta$ if and only if $x\in\Delta$; in other words,
$\check{\mathbb P}(\underline N_{\Delta/X})$ is the trivial projective
bundle
\begin{equation}\tag{2.1.6.2}
\begin{tikzcd}[column sep=large,row sep=large]
\check{\mathbb P}(\underline N_{\Delta/X})
  \arrow[rr,phantom,"\simeq"]
  \arrow[dr,"\text{canonical projection}"']
  \& {} \&
\Delta\times\mathbb P^r
  \arrow[dl,"\operatorname{pr}_i"]\\
{} \& \Delta \& {}
\end{tikzcd}.
\end{equation}
One has the following commutative diagram:
\begin{equation}\tag{2.1.6.3}
\begin{tikzcd}[column sep=large,row sep=large]
\Delta\times\mathbb P^r
  \arrow[r,"\sim"]
  \arrow[dr,"\operatorname{pr}_1"]
\&
\widetilde\Delta
  \arrow[r,hook]
  \arrow[d,"f|_{\widetilde\Delta}"]
\&
\widetilde X
  \arrow[r,hook]
  \arrow[d,"f"]
\&
X\times\mathbb P^r
  \arrow[dl,"\operatorname{pr}_1"]
\\
{} \& \Delta \arrow[r,hook] \& X \& {}
\end{tikzcd},
\end{equation}
and the composite inclusion
$\Delta\times\mathbb P^r\hookrightarrow X\times\mathbb P^r$ is obtained
from $\Delta\hookrightarrow X$ by base change.


}{\subsubsection*{2.1.5.}
In the special case
\[
X=\operatorname{Spec}(A),\qquad
\Delta=\operatorname{Spec}(A/I),\qquad
I=(g_0,\ldots,g_r),
\]
$\widetilde X$ is obtained by the Hopf construction (cf. EGA IV, 19.4.11).
One defines a morphism from $X-\Delta$ to $\mathbb P^r$ by
\[
x\in X-\Delta\longmapsto
\bigl(g_0(x),\ldots,g_r(x)\bigr)\in\mathbb P^r
\]
and regards its graph $\Gamma$ as embedded in $X\times\mathbb P^r$. Then
$\widetilde X$ is defined to be the schematic closure of $\Gamma$ in
$X\times\mathbb P^r$, and the morphism
$f:\widetilde X\longrightarrow X$ is the composite
\[
\widetilde X\lhook\joinrel\longrightarrow
X\times\mathbb P^r\xrightarrow{\ \operatorname{pr}_1\ }X.
\]

\subsubsection*{2.1.6.}
Similarly (cf. again EGA IV, 19.4.11), if $X$ is projective and $\Delta$ is
the intersection of $X$ with $1+r$ hypersurfaces, say with equations
$G_i=0$, $i=0,\ldots,r$, one again considers the morphism
\[
X-\Delta\longrightarrow\mathbb P^r,\qquad
x\longmapsto\bigl(G_0(x),\ldots,G_r(x)\bigr),
\]
and obtains $\widetilde X$ as the closure of its graph in
$X\times\mathbb P^r$. Concretely, in terms of homogeneous coordinates
$\lambda_0,\ldots,\lambda_r$ on $\mathbb P^r$, $\widetilde X$ is the
subscheme of $X\times\mathbb P^r$ defined by the equations
\begin{equation}\tag{2.1.6.1}
\lambda_iG_j(X)-\lambda_jG_i(X),
\qquad 0\leq i<j\leq r.
\end{equation}
These equations show that the point $(x,\lambda)$ lies in
$\widetilde\Delta$ if and only if $x\in\Delta$; in other words,
$\check{\mathbb P}(\underline N_{\Delta/X})$ is the trivial projective
bundle
\begin{equation}\tag{2.1.6.2}
\begin{tikzcd}[column sep=large,row sep=large]
\check{\mathbb P}(\underline N_{\Delta/X})
  \arrow[rr,phantom,"\simeq"]
  \arrow[dr,"\text{canonical projection}"']
  \& {} \&
\Delta\times\mathbb P^r
  \arrow[dl,"\operatorname{pr}_1"]\\
{} \& \Delta \& {}
\end{tikzcd}.
\end{equation}
One has the following commutative diagram:
\begin{equation}\tag{2.1.6.3}
\begin{tikzcd}[column sep=large,row sep=large]
\Delta\times\mathbb P^r
  \arrow[r,"\sim"]
  \arrow[dr,"\operatorname{pr}_1"]
\&
\widetilde\Delta
  \arrow[r,hook]
  \arrow[d,"f|_{\widetilde\Delta}"]
\&
\widetilde X
  \arrow[r,hook]
  \arrow[d,"f"]
\&
X\times\mathbb P^r
  \arrow[dl,"\operatorname{pr}_1"]
\\
{} \& \Delta \arrow[r,hook] \& X \& {}
\end{tikzcd},
\end{equation}
and the composite inclusion
$\Delta\times\mathbb P^r\hookrightarrow X\times\mathbb P^r$ is obtained
from $\Delta\hookrightarrow X$ by base change.


}{E7II-XVIII-003}\SGApage{II}{260}
\subsubsection*{Theorem 2.2.}
\underline{Let $X$ be smooth over a field $k$, let $\Delta$ be an everywhere
smooth subscheme of codimension $1+r$, and let
$f:\widetilde X\longrightarrow X$ be the blow-up of $X$ along $\Delta$.
Then:}

\SGAvar{\medskip
\noindent(2.2.1) \underline{The Leray spectral sequence}
\[
E_2^{p,q}=H^p\bigl(X,R^qf_*\mathbb Z_\ell
\Longrightarrow
H^{p+q}\bigl(\widetilde X_o\;\mathbb Z_\ell\bigr)
\]
\underline{degenerates.}

}{\medskip
\noindent(2.2.1) \underline{The Leray spectral sequence}
\[
E_2^{p,q}=H^p\bigl(X,R^qf_*\mathbb Z_\ell\bigr)
\Longrightarrow
H^{p+q}\bigl(\widetilde X,\mathbb Z_\ell\bigr)
\]
\underline{degenerates.}

}{E7II-XVIII-004}\medskip
\noindent(2.2.2) \underline{There is an additive isomorphism}
\[
\begin{aligned}
H^q(\widetilde X,\mathbb Z_\ell)
&\xleftarrow{\ \sim\ }
H^q(X,\mathbb Z_\ell)\oplus
\frac{H^q(\widetilde\Delta,\mathbb Z_\ell)}
     {\operatorname{Im}H^q(\Delta,\mathbb Z_\ell)}
\\
&\simeq
\left(\bigoplus_{j=1}^{r}
H^{q-2j}(\Delta,\mathbb Z_\ell(-j))\right)
\oplus H^q(X,\mathbb Z_\ell).
\end{aligned}
\]

\SGApage{II}{\(\approx\)261}
\noindent\underline{Proof.}
Let
\begin{equation}\tag{2.2.3}
g:\widetilde\Delta\longrightarrow\Delta
\end{equation}
be the restriction of $f$ to $\widetilde\Delta$, and let
\begin{equation}\tag{2.2.4}
\left\{
\begin{aligned}
j&:\Delta\hookrightarrow X,\\
i&:\widetilde\Delta\hookrightarrow\widetilde X
\end{aligned}
\right.
\end{equation}
be the inclusions, so that one has the commutative diagram
\begin{equation}\tag{2.2.5}
\begin{tikzcd}[column sep=huge,row sep=large]
\widetilde X \arrow[d,"f"'] \&
\widetilde\Delta \arrow[l,"i"'] \arrow[d,"g"]\\
X \& \Delta \arrow[l,"j"']
\end{tikzcd}.
\end{equation}
It induces a homomorphism between the Leray spectral sequences for the
morphisms $f$ and $g$, with the sheaf $\mathbb Z_\ell$ on $\widetilde X$ and
$\widetilde\Delta$:
\begin{equation}\tag{2.2.6}
E_s^{p,q}(f)\xrightarrow{\ (i^*,j^*)\ }E_s^{p,q}(g).
\end{equation}
On the other hand, since $f$ induces an isomorphism from
$\widetilde X-\widetilde\Delta$ to $X-\Delta$, the proper base-change theorem
gives
\begin{equation}\tag{2.2.7}
R^qf_*\mathbb Z_\ell\simeq
\begin{cases}
\mathbb Z_\ell,&q=0,\\[1ex]
j_*R^qg_*\mathbb Z_\ell,&q>0.
\end{cases}
\end{equation}
Hence
\begin{equation}\tag{2.2.8}
E_2^{p,q}(f)
\xrightarrow[\ \sim\ ]{\ (i^*,j^*)\ }
E_2^{p,q}(g)
\qquad(q\ne0).
\end{equation}
We shall prove by induction on $s$ that
\begin{equation}\tag{2.2.9}
E_s^{p,q}(f)
\xrightarrow[\ \sim\ ]{\ (i^*,j^*)\ }
E_s^{p,q}(g)
\qquad(q\ne0,\ s\geq2)
\end{equation}
and
\begin{equation}\tag{2.2.10}
d_s^{p,q}(f)=0\qquad(s\geq2).
\end{equation}
For the moment, admit the formula
\begin{equation}\tag{2.2.11}
d_s^{p,s-1}(f)=0\qquad(s\geq2,\ \text{for every }p).
\end{equation}

\SGApage{II}{\(\approx\)262}
Begin with $s=2$. We have a commutative diagram
\begin{equation}\tag{2.2.12}
\begin{tikzcd}[column sep=huge,row sep=large]
E_2^{p,q}(f)
  \arrow[r,"{(i^*,j^*)}"]
  \arrow[d,"{d_2^{p,q}(f)}"]
\&
E_2^{p,q}(g)
  \arrow[d,"{0=d_2^{p,q}(g)}"]
\\
E_2^{p+2,q-1}(f)
  \arrow[r,"{(i^*,j^*)}"]
\&
E_2^{p+2,q-1}(g)
\end{tikzcd}.
\end{equation}
For $q\leq0$, one has $E_2^{p+2,q-1}(f)=0$, hence
$d_2^{p,q}(f)=0$. For $q=1$, formula (2.2.11) gives
$d_2^{p,1}(f)=0$. For $q>1$, the horizontal arrows in (2.2.12) are
isomorphisms, while $d_2^{p,q}(g)=0$ by 1.2 and 2.1.4. This proves
(2.2.10) for $s=2$, and (2.2.9) for $s=2$ is exactly (2.2.8).

\SGApage{II}{\(\approx\)263}
\SGAvar{Suppose (2.2.9) and (2.2.10) are known for some $s\geq2$. Hypothesis
(2.2.10) implies
$E_s^{p,q}(f)\simeq E_{s+1}^{p,q}(f)$, and by 1.2 one has
$E_s^{p,q}(g)\simeq E_{s+1}^{p,q}(g)$, so there is a commutative diagram
\begin{equation}\tag{2.2.13}
\begin{tikzcd}[column sep=huge,row sep=large]
E_s^{p,q}(f)
  \arrow[r,"{(i^*,j^*)}"]
  \arrow[d,"\wr"']
\&
E_s^{p,q}(g)
  \arrow[d,"\wr"]
\\
R_{s+1}^{p,q}(f)
  \arrow[r,"{(i^*,j^*)}"]
\&
E_{s+1}^{p,q}(g)
\end{tikzcd},
\end{equation}
which proves (2.2.9) for $s+1$, in view of (2.2.9) for $s$. To deduce
(2.2.10) for $s+1$, consider the commutative diagram
\begin{equation}\tag{2.2.14}
\begin{tikzcd}[column sep=huge,row sep=large]
E_{s+1}^{p,q}(f)
  \arrow[r,"{(i^*,j^*)}"]
  \arrow[d,"{d_{s+1}^{p,q}(f)}"]
\&
E_{s+1}^{p,q}(g)
  \arrow[d,"{0=d_{s+1}^{p,q}(g)}"]
\\
E_{s+1}^{p+s+1,q-s}(f)
  \arrow[r,"{(i^*,j^*)}"]
\&
E_{s+1}^{p+s+1,q-s}(g)
\end{tikzcd}.
\end{equation}
For $q\leq s-1$, one has $E_{s+1}^{p+s+1,q-s}(f)=0$, hence
$d_s^{p,q}(f)=0$. For $q=s$, one has (2.2.11), and for $q>s$ the
horizontal arrows in (2.2.14) are isomorphisms and
$d_{s+1}^{p,q}(g)=0$ by 1.2. This gives the desired conclusion.


}{Suppose (2.2.9) and (2.2.10) are known for some $s\geq2$. Hypothesis
(2.2.10) implies
$E_s^{p,q}(f)\simeq E_{s+1}^{p,q}(f)$, and by 1.2 one has
$E_s^{p,q}(g)\simeq E_{s+1}^{p,q}(g)$, so there is a commutative diagram
\begin{equation}\tag{2.2.13}
\begin{tikzcd}[column sep=huge,row sep=large]
E_s^{p,q}(f)
  \arrow[r,"{(i^*,j^*)}"]
  \arrow[d,"\wr"']
\&
E_s^{p,q}(g)
  \arrow[d,"\wr"]
\\
E_{s+1}^{p,q}(f)
  \arrow[r,"{(i^*,j^*)}"]
\&
E_{s+1}^{p,q}(g)
\end{tikzcd},
\end{equation}
which proves (2.2.9) for $s+1$, in view of (2.2.9) for $s$. To deduce
(2.2.10) for $s+1$, consider the commutative diagram
\begin{equation}\tag{2.2.14}
\begin{tikzcd}[column sep=huge,row sep=large]
E_{s+1}^{p,q}(f)
  \arrow[r,"{(i^*,j^*)}"]
  \arrow[d,"{d_{s+1}^{p,q}(f)}"]
\&
E_{s+1}^{p,q}(g)
  \arrow[d,"{0=d_{s+1}^{p,q}(g)}"]
\\
E_{s+1}^{p+s+1,q-s}(f)
  \arrow[r,"{(i^*,j^*)}"]
\&
E_{s+1}^{p+s+1,q-s}(g)
\end{tikzcd}.
\end{equation}
For $q\leq s-1$, one has $E_{s+1}^{p+s+1,q-s}(f)=0$, hence
$d_{s+1}^{p,q}(f)=0$. For $q=s$, one has (2.2.11), and for $q>s$ the
horizontal arrows in (2.2.14) are isomorphisms and
$d_{s+1}^{p,q}(g)=0$ by 1.2. This gives the desired conclusion.


}{E7II-XVIII-005}It remains to verify (2.2.11), which says that the composite edge
homomorphism
\begin{equation}\tag{2.2.15}
E_2^{p,0}(f)\longrightarrow E_3^{p,0}(f)
\longrightarrow E_4^{p,0}(f)\longrightarrow\cdots
\longrightarrow E_\infty^{p,0}(f)
\end{equation}
is an isomorphism; that is, that
\begin{equation}\tag{2.2.16}
f^*:H^p(X,\mathbb Z_\ell)
\longrightarrow H^p(\widetilde X,\mathbb Z_\ell)
\end{equation}
is injective. But it has a left inverse, namely the
\underline{Gysin homomorphism} (SGA 5 IV)
\begin{equation}\tag{2.2.17}
f_*:H^p(\widetilde X,\mathbb Z_\ell)
\longrightarrow H^p(X,\mathbb Z_\ell),
\end{equation}
which indeed satisfies
\begin{equation}\tag{2.2.18}
f_*f^*=\mathrm{id}.
\end{equation}
This follows from the projection formula and the fact that $f_*(1)=1$, since
$f$ is birational and therefore has degree one. This proves (2.2.1).



\SGApage{II}{\(\approx\)264}
\SGAvar{As for (2.2.2), passing to cohomology in diagram (2.2.5) gives a
commutative diagram
\begin{equation}\tag{2.2.19}
\begin{tikzcd}[row sep=huge,column sep=huge]
H^p(\widetilde X,\mathbb Z_\ell)
  \arrow[r,"i^*"]
\&
H^p(\widetilde\Delta,\mathbb Z_\ell)
\\
H^p(X,\mathbb Z_\ell)
  \arrow[r,"j^*"]
  \arrow[u,"f^*"']
\&
H^p(\Delta,\mathbb Z_\ell)
  \arrow[u,"g^*"']
\end{tikzcd},
\end{equation}
hence a morphism
\begin{equation}\tag{2.2.20}
\frac{H^p(\widetilde X,\mathbb Z_\ell)}
     {f^*H^p(X,\mathbb Z_\ell)}
\xrightarrow{\ (i^*,j^*)\ }
\frac{H^p(\widetilde\Delta,\mathbb Z_\ell)}
     {g^*H^p(\Delta,\mathbb Z_\ell)}.
\end{equation}
Let $F^i$ denote the filtrations on
$H^\bullet(\widetilde X,\mathbb Z_\ell)$ and
$H^\bullet(\widetilde\Delta,\mathbb Z_\ell)$ coming from the Leray spectral
sequences. Then (2.2.20) can be rewritten as
\begin{equation}\tag{2.2.21}
\frac{H^p(\widetilde X,\mathbb Z_\ell)}
     {F^pH^p(\widetilde X,\mathbb Z_\ell)}
\xrightarrow{\ (i^*,j^*)\ }
\frac{H^p(\widetilde\Delta,\mathbb Z_\ell)}
     {H^pF^p(\widetilde\Delta,\mathbb Z_\ell)}.
\end{equation}
By (2.2.9), the map induced by (2.2.21) on associated graded objects is an
isomorphism, so (2.2.20) is an isomorphism. On the other hand, (2.2.18) gives
a direct-sum decomposition
\begin{equation}\tag{2.2.22}
\begin{tikzcd}[column sep=7em]
H^\bullet(\widetilde X,\mathbb Z_\ell)
  \arrow[r,shift right=1.2ex,
    "f_*\oplus(1-f^*f_*)"']
\&
H^\bullet(X,\mathbb Z_\ell)\oplus\Ker f_*
  \arrow[l,shift right=1.2ex,
    "f^*+\text{inclusion}"']
\end{tikzcd},
\end{equation}
and (1.2.2) gives the direct-sum decomposition
\begin{equation}\tag{2.2.23}
\frac{H^\bullet(\widetilde\Delta,\mathbb Z_\ell)}
     {g^*H^\bullet(\Delta,\mathbb Z_\ell)}
\xleftarrow{\ \sim\ }
\bigoplus_{1\leq j\leq r}
H^{\bullet-2j}(\Delta,\mathbb Z_\ell(-j)).
\end{equation}
Combining the isomorphism (2.2.20) and decompositions (2.2.22) and (2.2.23)
gives (2.2.2), as required.

\subsubsection*{Remark 2.2.24.}
Let us sketch a second, distinctly simpler proof of (2.2.2), valid when $X$
is \underline{proper} over $k$. More generally, without any hypothesis on
$X$, one obtains the corresponding result (2.2.26) in cohomology with proper
supports. The exact sequences of cohomology with proper supports, denoted
$H_c^\bullet$, for the pairs
$(\widetilde X,\Delta)$ and $(X,\Delta)$ form the commutative
diagram
\begin{center}
\resizebox{\textwidth}{!}{
\begin{tikzcd}[ampersand replacement=\&,column sep=small,row sep=large]
(2.2.25) \& {} \arrow[r] \&
H_c^i(\widetilde X-\widetilde\Delta,\mathbb Z_\ell)
  \arrow[r] \&
H_c^i(\widetilde X,\mathbb Z_\ell)
  \arrow[r] \&
H_c^i(\widetilde\Delta,\mathbb Z_\ell)
  \arrow[r] \&
H_c^{i+1}(\widetilde X-\widetilde\Delta,\mathbb Z_\ell)
\\
(2.2.26) \& {} \arrow[r] \&
H_c^i(X-\Delta,\mathbb Z_\ell)
  \arrow[r]
  \arrow[u,"\wr" left] \&
H_c^i(X,\mathbb Z_\ell)
  \arrow[r]
  \arrow[u,hook,"f^*" right] \&
H_c^i(\Delta,\mathbb Z_\ell)
  \arrow[r]
  \arrow[u,hook,"g^*" right] \&
H_c^{i+1}(X-\Delta,\mathbb Z_\ell)
  \arrow[u,"\wr" left]
\end{tikzcd}}
\end{center}
The quotient of the acyclic complex (2.2.25) by the acyclic subcomplex
(2.2.26) gives the \underline{acyclic} complex
\begin{equation}\tag{2.2.27}
\longrightarrow 0\longrightarrow
\frac{H_c^i(\widetilde X,\mathbb Z_\ell)}
     {f^*H_c^i(X,\mathbb Z_\ell)}
\longrightarrow
\frac{H_C^i(\widetilde\Delta,\mathbb Z_\ell)}
     {g^*H_c^i(\Delta,\mathbb Z_\ell)}
\longrightarrow0\longrightarrow.
\end{equation}
Thus, without considering spectral sequences, one finds that (2.2.20) is an
isomorphism. The arguments just given then apply and yield (2.2.2).

\subsubsection*{Corollary 2.3.}
\underline{Set}
\[
b_q(X)=\dim_{\mathbb Q_\ell}H^q(\chi,\mathbb Q_\ell),
\qquad
\chi(X)=\sum_q(-1)^qb_q(X).
\]
\underline{Then}
\begin{equation}\tag{2.3.1}
b_q(\widetilde X)
=b_q(X)+\sum_{j=1}^{r}b_{q-j}(\Delta),
\end{equation}
\begin{equation}\tag{2.3.2}
\chi(\widetilde X)=\chi(X)+r\,\chi(\Delta).
\end{equation}


}{As for (2.2.2), passing to cohomology in diagram (2.2.5) gives a
commutative diagram
\begin{equation}\tag{2.2.19}
\begin{tikzcd}[row sep=huge,column sep=huge]
H^p(\widetilde X,\mathbb Z_\ell)
  \arrow[r,"i^*"]
\&
H^p(\widetilde\Delta,\mathbb Z_\ell)
\\
H^p(X,\mathbb Z_\ell)
  \arrow[r,"j^*"]
  \arrow[u,"f^*"']
\&
H^p(\Delta,\mathbb Z_\ell)
  \arrow[u,"g^*"']
\end{tikzcd},
\end{equation}
hence a morphism
\begin{equation}\tag{2.2.20}
\frac{H^p(\widetilde X,\mathbb Z_\ell)}
     {f^*H^p(X,\mathbb Z_\ell)}
\xrightarrow{\ (i^*,j^*)\ }
\frac{H^p(\widetilde\Delta,\mathbb Z_\ell)}
     {g^*H^p(\Delta,\mathbb Z_\ell)}.
\end{equation}
Let $F^i$ denote the filtrations on
$H^\bullet(\widetilde X,\mathbb Z_\ell)$ and
$H^\bullet(\widetilde\Delta,\mathbb Z_\ell)$ coming from the Leray spectral
sequences. Then (2.2.20) can be rewritten as
\begin{equation}\tag{2.2.21}
\frac{H^p(\widetilde X,\mathbb Z_\ell)}
     {F^pH^p(\widetilde X,\mathbb Z_\ell)}
\xrightarrow{\ (i^*,j^*)\ }
\frac{H^p(\widetilde\Delta,\mathbb Z_\ell)}
     {F^pH^p(\widetilde\Delta,\mathbb Z_\ell)}.
\end{equation}
By (2.2.9), the map induced by (2.2.21) on associated graded objects is an
isomorphism, so (2.2.20) is an isomorphism. On the other hand, (2.2.18) gives
a direct-sum decomposition
\begin{equation}\tag{2.2.22}
\begin{tikzcd}[column sep=7em]
H^\bullet(\widetilde X,\mathbb Z_\ell)
  \arrow[r,shift right=1.2ex,
    "f_*\oplus(1-f^*f_*)"']
\&
H^\bullet(X,\mathbb Z_\ell)\oplus\Ker f_*
  \arrow[l,shift right=1.2ex,
    "f^*+\text{inclusion}"']
\end{tikzcd},
\end{equation}
and (1.2.2) gives the direct-sum decomposition
\begin{equation}\tag{2.2.23}
\frac{H^\bullet(\widetilde\Delta,\mathbb Z_\ell)}
     {g^*H^\bullet(\Delta,\mathbb Z_\ell)}
\xleftarrow{\ \sim\ }
\bigoplus_{1\leq j\leq r}
H^{\bullet-2j}(\Delta,\mathbb Z_\ell(-j)).
\end{equation}
Combining the isomorphism (2.2.20) and decompositions (2.2.22) and (2.2.23)
gives (2.2.2), as required.

\subsubsection*{Remark 2.2.24.}
Let us sketch a second, distinctly simpler proof of (2.2.2), valid when $X$
is \underline{proper} over $k$. More generally, without any hypothesis on
$X$, one obtains the corresponding result (2.2.26) in cohomology with proper
supports. The exact sequences of cohomology with proper supports, denoted
$H_c^\bullet$, for the pairs
$(\widetilde X,\widetilde\Delta)$ and $(X,\Delta)$ form the commutative
diagram
\begin{center}
\resizebox{\textwidth}{!}{
\begin{tikzcd}[ampersand replacement=\&,column sep=small,row sep=large]
(2.2.25) \& {} \arrow[r] \&
H_c^i(\widetilde X-\widetilde\Delta,\mathbb Z_\ell)
  \arrow[r] \&
H_c^i(\widetilde X,\mathbb Z_\ell)
  \arrow[r] \&
H_c^i(\widetilde\Delta,\mathbb Z_\ell)
  \arrow[r] \&
H_c^{i+1}(\widetilde X-\widetilde\Delta,\mathbb Z_\ell)
\\
(2.2.26) \& {} \arrow[r] \&
H_c^i(X-\Delta,\mathbb Z_\ell)
  \arrow[r]
  \arrow[u,"\wr" left] \&
H_c^i(X,\mathbb Z_\ell)
  \arrow[r]
  \arrow[u,hook,"f^*" right] \&
H_c^i(\Delta,\mathbb Z_\ell)
  \arrow[r]
  \arrow[u,hook,"g^*" right] \&
H_c^{i+1}(X-\Delta,\mathbb Z_\ell)
  \arrow[u,"\wr" left]
\end{tikzcd}}
\end{center}
The quotient of the acyclic complex (2.2.25) by the acyclic subcomplex
(2.2.26) gives the \underline{acyclic} complex
\begin{equation}\tag{2.2.27}
\longrightarrow 0\longrightarrow
\frac{H_c^i(\widetilde X,\mathbb Z_\ell)}
     {f^*H_c^i(X,\mathbb Z_\ell)}
\longrightarrow
\frac{H_c^i(\widetilde\Delta,\mathbb Z_\ell)}
     {g^*H_c^i(\Delta,\mathbb Z_\ell)}
\longrightarrow0\longrightarrow.
\end{equation}
Thus, without considering spectral sequences, one finds that (2.2.20) is an
isomorphism. The arguments just given then apply and yield (2.2.2).

\subsubsection*{Corollary 2.3.}
\underline{Set}
\[
b_q(X)=\dim_{\mathbb Q_\ell}H^q(X,\mathbb Q_\ell),
\qquad
\chi(X)=\sum_q(-1)^qb_q(X).
\]
\underline{Then}
\begin{equation}\tag{2.3.1}
b_q(\widetilde X)
=b_q(X)+\sum_{j=1}^{r}b_{q-2j}(\Delta),
\end{equation}
\begin{equation}\tag{2.3.2}
\chi(\widetilde X)=\chi(X)+r\,\chi(\Delta).
\end{equation}


}{E7II-XVIII-006}\SGApage{II}{\(\approx\)267}
\section*{3. The blow-up associated with a pencil}

\SGAvar{\subsection*{3.1.}
Let $X$ be projective, smooth, and irreducible over an algebraically closed
field $k$, of dimension $u$, and let
$D=\{H_\tau\}_{\tau\in\mathbb P^1}$ be a pencil of hypersurfaces whose axis
meets $X$ transversely, in $\Delta$.
We consider the blow-up of $X$ along $\Delta$, which we shall now describe
explicitly.

In terms of equations, let $(\lambda,\mu)$ be homogeneous coordinates on
$\mathbb P^1$, and let $F$ and $G$ be two linear forms such that $F=G=0$
defines the axis of the pencil.  The pencil is written
\begin{equation}\tag{3.1.1}
\lambda F-\mu G.
\end{equation}
By definition, the subscheme $\Delta$ is the intersection of $X$ with the
linear subvariety $G=F=0$.  By 2.1.6,
\begin{equation}\tag{3.1.2}
\widetilde X
=\{(x,(\lambda,\mu)\in X\times\mathbb P^1
       \mid \lambda F-\mu G=0\};
\end{equation}
it is the scheme-theoretic closure, in $X\times\mathbb P^1$, of the graph of
the map
\begin{equation}\tag{3.1.3}
\begin{tikzcd}[row sep=large,column sep=huge]
X-\Delta \arrow[r] \& \mathbb P^1 \\
x \arrow[r] \& (G(x),F(x)).
\end{tikzcd}
\end{equation}
The great advantage of $\widetilde X$ is that it is projective (being
embedded in $X\times\mathbb P^1$) and smooth (because $\Delta$ is smooth)
over $k$, and admits a morphism $\rho$ to $\mathbb P^1$, namely
\begin{equation}\tag{3.1.4}
\rho:\widetilde X\longrightarrow X\times\mathbb P^1
\xrightarrow{\operatorname{pr}_2}\mathbb P^1,
\end{equation}
whose fibers are the hyperplane sections of $X$ by the members of our
pencil:
\[
\rho^{-1}(t)=X_t\overset{\mathrm{def}}{=}X\cdot H_t.
\]

}{\subsection*{3.1.}
Let $X$ be projective, smooth, and irreducible over an algebraically closed
field $k$, of dimension $n$, and let
$D=\{H_\tau\}_{\tau\in\mathbb P^1}$ be a pencil of hypersurfaces whose axis
meets $X$ transversely, in $\Delta$.
We consider the blow-up of $X$ along $\Delta$, which we shall now describe
explicitly.

In terms of equations, let $(\lambda,\mu)$ be homogeneous coordinates on
$\mathbb P^1$, and let $F$ and $G$ be two linear forms such that $F=G=0$
defines the axis of the pencil.  The pencil is written
\begin{equation}\tag{3.1.1}
\lambda F-\mu G.
\end{equation}
By definition, the subscheme $\Delta$ is the intersection of $X$ with the
linear subvariety $G=F=0$.  By 2.1.6,
\begin{equation}\tag{3.1.2}
\widetilde X
=\{(x,(\lambda,\mu))\in X\times\mathbb P^1
       \mid \lambda F-\mu G=0\};
\end{equation}
it is the scheme-theoretic closure, in $X\times\mathbb P^1$, of the graph of
the map
\begin{equation}\tag{3.1.3}
\begin{tikzcd}[row sep=large,column sep=huge]
X-\Delta \arrow[r] \& \mathbb P^1 \\
x \arrow[r] \& (G(x),F(x)).
\end{tikzcd}
\end{equation}
The great advantage of $\widetilde X$ is that it is projective (being
embedded in $X\times\mathbb P^1$) and smooth (because $\Delta$ is smooth)
over $k$, and admits a morphism $\rho$ to $\mathbb P^1$, namely
\begin{equation}\tag{3.1.4}
\rho:\widetilde X\longrightarrow X\times\mathbb P^1
\xrightarrow{\operatorname{pr}_2}\mathbb P^1,
\end{equation}
whose fibers are the hyperplane sections of $X$ by the members of our
pencil:
\[
\rho^{-1}(t)=X_t\overset{\mathrm{def}}{=}X\cdot H_t.
\]

}{E7II-XVIII-007}\SGApage{II}{268}
\subsection*{3.2. The degree of the dual variety: topological method}

Suppose in addition that the pencil of 3.1 is such that
\begin{equation}\tag{3.2.1}
X_t\text{ is smooth for every }t\text{ in a nonempty open subset of }
\mathbb P^1,
\end{equation}
and
\begin{equation}\tag{3.2.2}
\text{if $X_t$ is singular, its only singularities are nondegenerate
quadratic singularities (XVII 1.1).}
\end{equation}
These conditions hold in particular for a \underline{Lefschetz pencil}
(XVII 2.21) if $\operatorname{char}k\ne2$ or if $X$ has even dimension.

\SGAvar{Let $\eta$ be the generic point of $\mathbb P^1$, let $\bar\eta$ be a
geometric point above it, let $n=\dim X$, and let $\chi$ be the
Euler--Poincar\'e characteristic with coefficients in $\mathbb Q_\ell$.
The Euler--Poincar\'e formula XVI 2.1 applied to $S$ gives
\begin{equation}\tag{3.2.3}
\chi(\widetilde X)
=\chi(\mathbb P^1)\,\chi(X_{\bar\eta})
 +(-1)^{\eta}\sum_{\substack{\text{singular points}\\
                         \text{of the fibers}}}1,
\end{equation}
which, in view of (2.2.2), can be rewritten
\begin{equation}\tag{3.2.4}
(-1)^\eta
\left[\begin{array}{c}
\text{total number of}\\[3pt]
\text{singular points}
\end{array}\right]
=\chi(X)+\chi(\Delta)-2\chi(X_{\bar\eta}).
\end{equation}
Notice that the quantity
$\chi(X)+\chi(\Delta)-2\chi(X_{\bar\eta})$ is independent of the particular
choice of pencil.  It depends only on the embedding of $X$ in some
$\mathbb P^r$, and indeed only on the Chern class of
$\underline{\mathcal O}_X(1)$; see XVII 5.7.5 and 5.7.6.

\subsubsection*{3.2.5.}
Let us make this explicit in a concrete case.  Take $X=\mathbb P^2$ and let
$\{H_\tau\}$ be a pencil of plane curves of degree $d$ satisfying 3.1,
(3.2.1), and (3.2.2).  Then
\begin{equation}\tag{3.2.6}
\begin{array}{ll}
\chi(X)=3 & \text{(by 1.2)},\\[4pt]
\chi(X_{\bar\eta})=-d(d-3) & \text{(see SGA 5 VII 7.3)},\\[4pt]
\chi(\Delta)=d^2,&
\end{array}
\end{equation}
whence
\begin{equation}\tag{3.2.7}
\text{the total number of singular points is }3(d-1)^2.
\end{equation}
Take the special case $d=3$ and consider the family of elliptic curves
(in characteristic different from $3$)
\begin{equation}\tag{3.2.8}
X^3+Y^3+Z^3+3\mu XYZ.
\end{equation}
The four singular fibers lie over $\mu=$ a cube root of unity and $\mu=\infty$,
respectively,
and each singular fiber is the union of three lines as follows:
\begin{equation}\tag{3.2.9}
\begin{tikzcd}[row sep=4.5em,column sep=4.5em]
{} \& {} \arrow[dl,dash,shorten <=-9mm,shorten >=-9mm]
        \arrow[dr,dash,shorten <=-9mm,shorten >=-9mm] \& {} \& {} \\
{} \arrow[rr,dash,shorten <=-11mm,shorten >=-11mm] \& {} \& {} \& {,}
\end{tikzcd}
\end{equation}
It has exactly three singular points, all nondegenerate quadratic
singularities.  Altogether this gives twelve singular points and verifies
(3.2.7) in this example.

\subsubsection*{Proposition 3.2.10.}
\underline{Let $X\hookrightarrow\mathbb P^r$ be a Lefschetz embedding}
(XVII 2.3), \underline{and suppose that $n=\dim X$ is even, or that
$\operatorname{char}k\ne2$.  Then every Lefschetz pencil, regarded as a
line in $\check{\mathbb P}^{r}$, meets $\check X$ in}
\begin{equation}\tag{3.2.11}
(-1)^\eta\bigl[\chi(X)+\chi(\Delta)-2\chi(X_{\bar\eta})\bigr]
\end{equation}
\underline{points, and meets it transversely} (EGA IV 17.13.7).

}{Let $\eta$ be the generic point of $\mathbb P^1$, let $\bar\eta$ be a
geometric point above it, let $n=\dim X$, and let $\chi$ be the
Euler--Poincar\'e characteristic with coefficients in $\mathbb Q_\ell$.
The Euler--Poincar\'e formula XVI 2.1 applied to $\rho$ gives
\begin{equation}\tag{3.2.3}
\chi(\widetilde X)
=\chi(\mathbb P^1)\,\chi(X_{\bar\eta})
 +(-1)^n\sum_{\substack{\text{singular points}\\
                         \text{of the fibers}}}1,
\end{equation}
which, in view of (2.2.2), can be rewritten
\begin{equation}\tag{3.2.4}
(-1)^n
\left[\begin{array}{c}
\text{total number of}\\[3pt]
\text{singular points}
\end{array}\right]
=\chi(X)+\chi(\Delta)-2\chi(X_{\bar\eta}).
\end{equation}
Notice that the quantity
$\chi(X)+\chi(\Delta)-2\chi(X_{\bar\eta})$ is independent of the particular
choice of pencil.  It depends only on the embedding of $X$ in some
$\mathbb P^r$, and indeed only on the Chern class of
$\underline{\mathcal O}_X(1)$; see XVII 5.7.5 and 5.7.6.

\subsubsection*{3.2.5.}
Let us make this explicit in a concrete case.  Take $X=\mathbb P^2$ and let
$\{H_\tau\}$ be a pencil of plane curves of degree $d$ satisfying 3.1,
(3.2.1), and (3.2.2).  Then
\begin{equation}\tag{3.2.6}
\begin{array}{ll}
\chi(X)=3 & \text{(by 1.2)},\\[4pt]
\chi(X_{\bar\eta})=-d(d-3) & \text{(see SGA 5 VII 7.3)},\\[4pt]
\chi(\Delta)=d^2,&
\end{array}
\end{equation}
whence
\begin{equation}\tag{3.2.7}
\text{the total number of singular points is }3(d-1)^2.
\end{equation}
Take the special case $d=3$ and consider the family of elliptic curves
(in characteristic different from $3$)
\begin{equation}\tag{3.2.8}
X^3+Y^3+Z^3+3\mu XYZ.
\end{equation}
The four singular fibers lie over $\mu^3=-1$ and $\mu=\infty$,
respectively,
and each singular fiber is the union of three lines as follows:
\begin{equation}\tag{3.2.9}
\begin{tikzcd}[row sep=4.5em,column sep=4.5em]
{} \& {} \arrow[dl,dash,shorten <=-9mm,shorten >=-9mm]
        \arrow[dr,dash,shorten <=-9mm,shorten >=-9mm] \& {} \& {} \\
{} \arrow[rr,dash,shorten <=-11mm,shorten >=-11mm] \& {} \& {} \& {,}
\end{tikzcd}
\end{equation}
It has exactly three singular points, all nondegenerate quadratic
singularities.  Altogether this gives twelve singular points and verifies
(3.2.7) in this example.

\subsubsection*{Proposition 3.2.10.}
\underline{Let $X\hookrightarrow\mathbb P^r$ be a Lefschetz embedding}
(XVII 2.3), \underline{and suppose that $n=\dim X$ is even, or that
$\operatorname{char}k\ne2$.  Then every Lefschetz pencil, regarded as a
line in $\check{\mathbb P}^{r}$, meets $\check X$ in}
\begin{equation}\tag{3.2.11}
(-1)^n\bigl[\chi(X)+\chi(\Delta)-2\chi(X_{\bar\eta})\bigr]
\end{equation}
\underline{points, and meets it transversely} (EGA IV 17.13.7).

}{E7II-XVIII-008}\SGApage{II}{\(\approx\)270}
\underline{Proof.}
By hypothesis, every singularity of a fiber of the Lefschetz pencil is a
nondegenerate quadratic singularity.  Formula (3.2.4) therefore shows that
the number of singular fibers in the pencil $D$---each containing exactly
one singular point---is the number (3.2.11).  Now the fiber $X_\tau$ is
singular if and only if $\tau\in\check X\cap D$.

\SGAvar{To see that the intersection is transverse, observe that (3.2.11) is
\underline{independent} of the particular Lefschetz pencil.  Since the
generic pencil is a Lefschetz pencil, (3.2.11) is the degree of $\check X$;
that is, its degree as a hypersurface in $\check{\mathbb P}^{r}$, taken by
convention to be $0$ if $\dim\check X\le r-2$.  If a line $D$ meets a
hypersurface of degree $d>0$ in $d$ distinct points $x_i$
($i\le i\le d$), the intersection must be transverse.  Indeed, the local
intersection multiplicity
\[
\mu_i=\operatorname{rank}_{x_i}(D\cap\check X)
\]
is at least $1$, while $\sum^d\mu_i=d$.  Hence $\mu_i=1$ for every
$i$; equivalently, $D\cap\check X$ is smooth, which is precisely
transversality.

}{To see that the intersection is transverse, observe that (3.2.11) is
\underline{independent} of the particular Lefschetz pencil.  Since the
generic pencil is a Lefschetz pencil, (3.2.11) is the degree of $\check X$;
that is, its degree as a hypersurface in $\check{\mathbb P}^{r}$, taken by
convention to be $0$ if $\dim\check X\le r-2$.  If a line $D$ meets a
hypersurface of degree $d>0$ in $d$ distinct points $x_i$
($1\le i\le d$), the intersection must be transverse.  Indeed, the local
intersection multiplicity
\[
\mu_i=\operatorname{length}_{x_i}(D\cap\check X)
\]
is at least $1$, while $\sum_{i=1}^d\mu_i=d$.  Hence $\mu_i=1$ for every
$i$; equivalently, $D\cap\check X$ is smooth, which is precisely
transversality.

}{E7II-XVIII-009}\SGApage{II}{271}
\subsubsection*{Remark 3.2.12.}
For a Lefschetz pencil $D$ already assumed to meet $\check X$ transversely
(a condition satisfied in any case by sufficiently general $D$), the
result 3.2.10 may also be regarded as the conjunction of XVII 5.7.2 and
the fact that $\deg\varphi=1$ contained in XVII 3.5.  Of course, the fact
proved here that \underline{every} Lefschetz pencil $D$ meets $\check X$
transversely---by means of the Euler--Poincar\'e formula XVI 2.1---may also
be verified directly and elementarily, using the fact that the axis of
$D$ meets $X$ transversely.


\section*{4. Cohomology of the blow-up associated with a pencil}

\subsection*{4.1.}
We return to the situation of 3.1, summarized by the commutative diagram
(see 2.1.6.3):
\begin{equation}\tag{4.1.1}
\resizebox{\textwidth}{!}{
\begin{tikzcd}[ampersand replacement=\&,row sep=huge,column sep=huge]
X\times\mathbb P^1
  \& \widetilde X
      \arrow[l,hook',"f\times\rho"]
      \arrow[d,"f"]
      \arrow[dl,"\rho"']
  \& \widetilde\Delta
      \arrow[l,hook',"i"]
      \arrow[d,"g"]
      \arrow[r,"{(g,\rho_i)}","\sim"']
  \& \Delta\times\mathbb P^1 \\
\mathbb P^1 \& X \& \Delta\arrow[l,hook',"j"] \& {}
\end{tikzcd}}
\end{equation}
We begin by making (2.2.2) more explicit in this case.

\subsection*{Proposition 4.2.}
\underline{Under the hypotheses of 3.1:}

\smallskip
\noindent
\textup{(4.2.1)} \underline{The homomorphisms of} $\mathbb Z_\ell$\underline{-modules}
\[
\begin{tikzcd}[column sep=huge]
H^\bullet(\widetilde X,\mathbb Z_\ell)
  \arrow[r,shift left=.7ex,"f_*\oplus(1-f^*f_*)"]
\& H^\bullet(X,\mathbb Z_\ell)\oplus
   \operatorname{Ker}\!\left(f_*:H^\bullet(\widetilde X,\mathbb Z_\ell)
       \to H^\bullet(X,\mathbb Z_\ell)\right)
  \arrow[l,shift left=.7ex,"f^*+\operatorname{incl}"']
\end{tikzcd}
\]
\underline{are mutually inverse isomorphisms.}

\SGApage{II}{\(\approx\)272}
\smallskip
\noindent
\textup{(4.2.2)} \underline{The homomorphisms of $\mathbb Z_\ell$-modules}
\[
\begin{tikzcd}[column sep=huge]
H^\bullet(\widetilde X,\mathbb Z_\ell)
  \arrow[r,shift left=.7ex,"f_*\oplus(-g_*i^*)"]
\& H^\bullet(X,\mathbb Z_\ell)\oplus
   H^{\bullet-2}(\Delta,\mathbb Z_\ell(-1))
  \arrow[l,shift left=.7ex,"f^*+i_*g^*"']
\end{tikzcd}
\]
\underline{are mutually inverse isomorphisms.}

\SGAvar{\smallskip
\noindent
\textup{(4.2.3)} \underline{The algebra structure on}
\[
H^\bullet(X,\mathbb Z_\ell)\oplus
H^{\bullet-2}(\Delta,\mathbb Z_\ell(-1)),
\]
\underline{transported by (4.2.2) from the cup product on}
$H^\bullet(\widetilde X,\mathbb Z_\ell)$,
\underline{is expressed, for}
\[
a,b\in H^\bullet(X,\mathbb Z_\ell),
\qquad
x,y\in H^{\bullet-2}(\Delta,\mathbb Z_\ell(-1),
\]
\underline{by the formulas}
\begin{equation}\tag{4.2.3.1}
(0\oplus x)\wedge(0\oplus y)
=-j_*(xy)\oplus2xy\,c_1(\underline{\mathcal O}_\Delta(1)),
\end{equation}
\begin{equation}\tag{4.2.3.2}
(a\oplus0)\wedge(b\oplus0=ab\oplus0,
\end{equation}
\begin{equation}\tag{4.2.3.3}
(a\oplus0)\wedge(0\oplus y)=0\oplus j^*(a)y,
\end{equation}
\begin{equation}\tag{4.2.3.4}
(0\oplus x)\wedge(b\oplus0)=0\oplus xj^*(b).
\end{equation}
\underline{In particular, if}
\[
a\oplus x\in H^\bullet(X,\mathbb Z_\ell)
  \oplus H^{\bullet-2}(\Delta,\mathbb Z_\ell(-1)
\]
and
\[
b\oplus y\in H^{2n-\bullet}(X,\mathbb Z_\ell)
  \oplus H^{2n-2-\bullet}(\Delta,\mathbb Z_\ell(-1)),
\]
\underline{then}
\begin{equation}\tag{4.2.3.5}
(a\oplus x)\wedge(b\oplus y)=ab-j_*(xy).
\end{equation}
In other words, the total intersection matrix of
$H^\bullet(\widetilde X,\mathbb Z_\ell)$ is
\[
\left(
\begin{array}{c|c}
\begin{array}{c}
\text{total intersection}\\[-1pt]
\text{matrix on $X$}
\end{array} & 0\\ \hline
0 &
-\begin{array}{c}
\text{total intersection}\\[-1pt]
\text{matrix on $\Delta$}
\end{array}
\end{array}
\right).
\]

\underline{Proof.}
We have already proved (4.2.1); see (2.2.22).  For (4.2.2), first observe
that, since $g:\widetilde\Delta\to\Delta$ is a projective-line bundle, the
sequence
\begin{equation}\tag{4.2.4}
0\longrightarrow H^\bullet(\Delta,\mathbb Z_\ell)
\xrightarrow{g^*}H^\bullet(\widetilde\Delta,\mathbb Z_\ell)
\xrightarrow{g_*}H^{\bullet-2}(\Delta,\mathbb Z_\ell(-1))
\longrightarrow0
\end{equation}
is exact, by (1.2.2).  Together with (2.2.26), this gives the commutative
diagram of exact sequences
\begin{equation}\tag{4.2.5}
\resizebox{\textwidth}{!}{
\begin{tikzcd}[ampersand replacement=\&,column sep=large,row sep=huge]
0\arrow[r]
  \&H^\bullet(X,\mathbb Z_\ell)
       \arrow[r,"f^*"]\arrow[d,leftrightarrow,"j^*" right]
  \&H^\bullet(\widetilde X,\mathbb Z_\ell)
       \arrow[r,"g_*i^*"]\arrow[d,"i^*" right]
  \&H^{\bullet-2}(\Delta,\mathbb Z_\ell(-1))
       \arrow[r]\arrow[d,equal]
  \&0\\
0\arrow[r]
  \&H^\bullet(\Delta,\mathbb Z_\ell)
       \arrow[r,"g^*"]
  \&H^\bullet(\widetilde\Delta,\mathbb Z_\ell)
       \arrow[r,"g_*"]
  \&H^{\bullet-2}(\Delta,\mathbb Z_\ell(-1))
       \arrow[r]
  \&0.
\end{tikzcd}}
\end{equation}
In view of the splitting (4.2.1), this yields the isomorphism
\begin{equation}\tag{4.2.6}
f_*\oplus g_*i^*:
H^\bullet(\widetilde X,\mathbb Z_\ell)
\xrightarrow{\sim}
H^\bullet(X,\mathbb Z_\ell)\oplus
H^{\bullet-2}(\Delta,\mathbb Z_\ell(-1)).
\end{equation}
The transpose of (4.2.5), after the degree reversal
$\bullet\mapsto2n-\bullet$, gives the commutative diagram of exact sequences
\begin{equation}\tag{4.2.7}
\resizebox{\textwidth}{!}{
\begin{tikzcd}[ampersand replacement=\&,column sep=large,row sep=large]
0
  \&H^\bullet(X,\mathbb Z_\ell)\arrow[l]
  \&H^\bullet(\widetilde X,\mathbb Z_\ell)\arrow[l,"f_*"']
  \&H^{\bullet-2}(\Delta,\mathbb Z_\ell(-1))
       \arrow[l,"i_*g^*"']
  \&0\arrow[l]\\
0
  \&H^{\bullet-2}(\Delta,\mathbb Z_\ell(-2))
       \arrow[l]\arrow[u,"j_*"']
  \&H^{\bullet-2}(\widetilde\Delta,\mathbb Z_\ell(-1))
       \arrow[l,"g_*"']\arrow[u,"i_*"']
  \&H^{\bullet-2}(\Delta,\mathbb Z_\ell(-1))
       \arrow[l,"g_*"']\arrow[u,equal]
  \&0\arrow[l]
\end{tikzcd}}
\end{equation}
and hence the isomorphism
\begin{equation}\tag{4.2.8}
H^\bullet(X,\mathbb Z_\ell)\oplus
H^{\bullet-2}(\Delta,\mathbb Z_\ell(-1))
\xrightarrow{\ f^*+i_*g^*\ }
H^\bullet(\widetilde X,\mathbb Z_\ell).
\end{equation}


It remains to prove that the composite
\begin{equation}\tag{4.2.9}
\begin{aligned}
&H^\bullet(X,\mathbb Z_\ell)\oplus
 H^{\bullet-2}(\Delta,\mathbb Z_\ell(-1))
 \xrightarrow{\ f^*+i_*g^*\ }
 H^\bullet(\widetilde X,\mathbb Z_\ell)\\[-2pt]
&\hspace{55mm}\xrightarrow{\ f_*\oplus(g_*i^*)\ }
 H^\bullet(X,\mathbb Z_\ell)\oplus
 H^{\bullet-2}(\Delta,\mathbb Z_\ell(-1))
\end{aligned}
\end{equation}
is the identity.  Now
\[
f_*f^*=\mathrm{id}\quad\text{on }H^\bullet(X,\mathbb Z_\ell)
\quad\text{by (2.2.18)},
\]
\[
g_*i^*f^*=0\quad\text{by (4.2.5)},
\qquad
f_*i_*g^*=0\quad\text{by (4.2.7)}.
\]
It therefore suffices to prove
\begin{equation}\tag{4.2.10}
g_*i^*i_*g^*(X)=X
\quad\text{for }X\in H^\bullet(\Delta,\mathbb Z_\ell(-1)).
\end{equation}

To verify (4.2.10), first observe that
$i:\widetilde\Delta\hookrightarrow\widetilde X$ is the inclusion of a
smooth subvariety of codimension one.  Thus (SGA 5 VII)
\begin{equation}\tag{4.2.11}
i^*i_*(x)=c_1(\underline N_{\widetilde\Delta/\widetilde X})\,x,
\end{equation}
and
\begin{equation}\tag{4.2.12}
i_*(x)i_*(y)
=i_*\!\left(xy\,c_1(\underline N_{\widetilde\Delta/\widetilde X})\right).
\end{equation}
Write
$\xi_X=\underline{\mathcal O}_X(1)$,
$\xi_\Delta=\underline{\mathcal O}_\Delta(1)$, and
$\xi_{\mathbb P^1}=\underline{\mathcal O}_{\mathbb P^1}(1)$.
By EGA II 8.1.7, if $\widetilde{\mathfrak J}$ is the ideal of
$\widetilde\Delta$ in $\widetilde X$, then
\begin{equation}\tag{4.2.13}
\begin{aligned}
\underline{\check N}_{\widetilde\Delta/\widetilde X}
&\simeq\widetilde{\mathfrak J}/\widetilde{\mathfrak J}^{,2}
 \simeq\widetilde{\mathfrak J}
       \otimes_{\underline{\mathcal O}_{\widetilde X}}
       \underline{\mathcal O}_{\widetilde\Delta}\\
&\simeq\underline{\mathcal O}_{\widetilde X}(1)
       \otimes_{\underline{\mathcal O}_{\widetilde X}}
       \underline{\mathcal O}_{\widetilde\Delta}
 =\underline{\mathcal O}_{\widetilde\Delta}(1).
\end{aligned}
\end{equation}
Here $\widetilde\Delta$ is viewed as
$\mathbb P(\mathfrak J/\mathfrak J^2)$, where $\mathfrak J$ is the ideal of
$\Delta$ in $X$.  Since $\Delta$ is the intersection of $X$ with two given
hyperplanes, there is a canonical isomorphism
\[
\mathfrak J/\mathfrak J^2\simeq
\underline{\mathcal O}_\Delta^{,2}(-1),
\]
which induces
\[
\widetilde\Delta=\mathbb P(\mathfrak J/\mathfrak J^2)
\simeq\mathbb P(\underline{\mathcal O}_\Delta^{,2}(-1))
\xrightarrow[\sim]{\alpha}
\mathbb P(\underline{\mathcal O}_\Delta^{,2})
\simeq\Delta\times\mathbb P^1.
\]
The canonical isomorphism $\alpha$ of EGA II 4.1.4 carries the tautological
line bundle to
\[
\underline{\mathcal O}_{\mathbb P(\underline{\mathcal O}_\Delta^{,2})}(1)
\otimes_{\underline{\mathcal O}_\Delta^{,2}}\underline{\mathcal O}_\Delta(-1)
=\underline{\mathcal O}_{\mathbb P^1}(1)
\otimes_K\underline{\mathcal O}_\Delta(-1).
\]
Consequently, (4.2.13) becomes
\begin{equation}\tag{4.2.14}
\underline{\check N}_{\widetilde\Delta/\widetilde X}
\simeq(\rho i)^*\underline{\mathcal O}_{\mathbb P}(1)
       \otimes g^*\underline{\mathcal O}_\Delta(-1).
\end{equation}

\subsection*{Lemma 4.3.}
\underline{One has}
\[
c_1(\underline N_{\widetilde\Delta/\widetilde X})
=g^*c_1(\xi_\Delta)-(\rho i)^*c_1(\xi_{\mathbb P^1}).
\]
Return now to the proof of (4.2.10).  By (4.2.11) and Lemma 4.3, it is enough
to prove
\[
g_*\!\left(g^*(x)\left[(\rho i)^*c_1(\xi_{\mathbb P^1})
                       -g^*c_1(\xi_\Delta\right]\right)=x.
\]
By (4.2.4), $g_*g^*=0$, while
\[
\begin{tikzcd}[column sep=huge]
\widetilde\Delta\arrow[r,"g\times\rho_i","\sim"']
  \&\Delta\times\mathbb P^1,
\end{tikzcd}
\]
so
\begin{equation}\tag{4.3.1}
g_*\!\left(g^*(x)(\rho i)^*c_1(\xi_{\mathbb P^1})\right.=x,
\end{equation}
and hence
\begin{equation}\tag{4.3.2}
g_*\!\left(g^*(x)c_1(\underline N_{\widetilde\Delta/\widetilde X})\right)
=-x.
\end{equation}

We now compute the algebra structure on
\[
H^\bullet(X,\mathbb Z_\ell)\oplus
H^{\bullet-2}(\Delta,\mathbb Z_\ell(-1))
\]
transported from the cup product by (4.2.2).  Let $a,b$ belong to
$H^\bullet(X,\mathbb Z_\ell)$ and $x,y$ to
$H^{\bullet-2}(\Delta,\mathbb Z_\ell(-1))$.  Begin with (4.2.3.1), and write
\[
(0\oplus x)\Lambda(0\oplus y)=A\oplus B.
\]
By (4.2.12) and (4.3.2),
\[
\begin{aligned}
A
&=f_*\!\left((i_*g^*x)(i_*g^*y)\right.\\
&=f_*i_*\!\left(g^*(xy)c_1(\underline N_{\widetilde\Delta/\widetilde X})\right)\\
&=j_*g_*\!\left(g^*(xy)c_1(\underline N_{\widetilde\Delta/\widetilde X})\right)
=-j_*(xy),
\end{aligned}
\]
and
\[
\begin{aligned}
B
&=-g_*i^*\!\left((i_*g^*x)(i_*g^*y)\right)\\
&=-g_*i^*i_*\!\left(g^*(xy)c_1(\underline N_{\widetilde\Delta/\widetilde X})\right)\qquad\text{by (4.3.2)}\\
&=-g_*\!\left(g^*(xy)c_1(\underline N_{\widetilde\Delta/\widetilde X})^2\right)
\qquad\text{by (4.3.2)}.
\end{aligned}
\]
Lemma 4.3 gives
\[
\begin{aligned}
c_1(\underline N_{\widetilde\Delta/\widetilde X})^2
&=\left(g^*c_1(\underline{\mathcal O}_{\Delta(1)})
 -(\rho i)^*c_1(\xi_{\mathbb P^1})\right)^2\\
&=g^*\!\left(c_1(\xi_\Delta)^2\right)
 -2g^*c_1(\xi_\Delta)(\rho i)^*c_1(\xi_{\mathbb P^1}),
\end{aligned}
\]
because $c_1(\xi_{\mathbb P^1})^2=0$.  Therefore, by (4.2.4) and (4.3.1),
\[
\begin{aligned}
B&=-g_*(g^*(xy)g^*(c_1(\xi_\Delta)^2))
 +g_*(2g^*(xy)g^*(c_1(\xi_\Delta)(\rho i)^*(c_1(\xi_{\mathbb P^1}))\\
A&=0+2xy\,c_1(\xi_\Delta),
\end{aligned}
\]
which proves (4.2.3.1).  Formula (4.2.3.2), $(a\oplus0)\Lambda(b\oplus0)=ab\oplus0$, expresses the fact that $f^*$ is
an algebra homomorphism.

To prove (4.2.3.3), the projection formula shows that
$\operatorname{Ker}f_*\cap\operatorname{im}f^*\subset\operatorname{Ker}f_*$.  Hence
\[
(a\oplus0)\Lambda(0\oplus y)=0\oplus C,
\]
where
\begin{equation}\tag{4.3.3}
C=-g_*i^*\!\left(f^*(a)i_*g^*(y)\right).
\end{equation}
Using (4.3.2),
\[
\begin{aligned}
C
&=-g_*\!\left(i^*f^*(a)g^*(y)
      c_1(\underline N_{\widetilde\Delta/\widetilde X})\right)\\
&=-g_*\!\left(g^*j^*(a)g^*(y)
      c_1(\underline N_{\widetilde\Delta/\widetilde X})\right)\\
&=j^*(a)y.
\end{aligned}
\]
This proves (4.2.3.3).  Finally, (4.2.3.4) follows from (4.2.3.3) by graded
commutativity of the cup product.


}{\smallskip
\noindent
\textup{(4.2.3)} \underline{The algebra structure on}
\[
H^\bullet(X,\mathbb Z_\ell)\oplus
H^{\bullet-2}(\Delta,\mathbb Z_\ell(-1)),
\]
\underline{transported by (4.2.2) from the cup product on}
$H^\bullet(\widetilde X,\mathbb Z_\ell)$,
\underline{is expressed, for}
\[
a,b\in H^\bullet(X,\mathbb Z_\ell),
\qquad
x,y\in H^{\bullet-2}(\Delta,\mathbb Z_\ell(-1)),
\]
\underline{by the formulas}
\begin{equation}\tag{4.2.3.1}
(0\oplus x)\wedge(0\oplus y)
=-j_*(xy)\oplus2xy\,c_1(\underline{\mathcal O}_\Delta(1)),
\end{equation}
\begin{equation}\tag{4.2.3.2}
(a\oplus0)\wedge(b\oplus0)=ab\oplus0,
\end{equation}
\begin{equation}\tag{4.2.3.3}
(a\oplus0)\wedge(0\oplus y)=0\oplus j^*(a)y,
\end{equation}
\begin{equation}\tag{4.2.3.4}
(0\oplus x)\wedge(b\oplus0)=0\oplus xj^*(b).
\end{equation}
\underline{In particular, if}
\[
a\oplus x\in H^\bullet(X,\mathbb Z_\ell)
  \oplus H^{\bullet-2}(\Delta,\mathbb Z_\ell(-1))
\]
and
\[
b\oplus y\in H^{2n-\bullet}(X,\mathbb Z_\ell)
  \oplus H^{2n-2-\bullet}(\Delta,\mathbb Z_\ell(-1)),
\]
\underline{then}
\begin{equation}\tag{4.2.3.5}
(a\oplus x)\wedge(b\oplus y)=ab-j_*(xy).
\end{equation}
In other words, the total intersection matrix of
$H^\bullet(\widetilde X,\mathbb Z_\ell)$ is
\[
\left(
\begin{array}{c|c}
\begin{array}{c}
\text{total intersection}\\[-1pt]
\text{matrix on $X$}
\end{array} & 0\\ \hline
0 &
-\begin{array}{c}
\text{total intersection}\\[-1pt]
\text{matrix on $\Delta$}
\end{array}
\end{array}
\right).
\]

\underline{Proof.}
We have already proved (4.2.1); see (2.2.22).  For (4.2.2), first observe
that, since $g:\widetilde\Delta\to\Delta$ is a projective-line bundle, the
sequence
\begin{equation}\tag{4.2.4}
0\longrightarrow H^\bullet(\Delta,\mathbb Z_\ell)
\xrightarrow{g^*}H^\bullet(\widetilde\Delta,\mathbb Z_\ell)
\xrightarrow{g_*}H^{\bullet-2}(\Delta,\mathbb Z_\ell(-1))
\longrightarrow0
\end{equation}
is exact, by (1.2.2).  Together with (2.2.26), this gives the commutative
diagram of exact sequences
\begin{equation}\tag{4.2.5}
\resizebox{\textwidth}{!}{
\begin{tikzcd}[ampersand replacement=\&,column sep=large,row sep=huge]
0\arrow[r]
  \&H^\bullet(X,\mathbb Z_\ell)
       \arrow[r,"f^*"]\arrow[d,"j^*" right]
  \&H^\bullet(\widetilde X,\mathbb Z_\ell)
       \arrow[r,"g_*i^*"]\arrow[d,"i^*" right]
  \&H^{\bullet-2}(\Delta,\mathbb Z_\ell(-1))
       \arrow[r]\arrow[d,equal]
  \&0\\
0\arrow[r]
  \&H^\bullet(\Delta,\mathbb Z_\ell)
       \arrow[r,"g^*"]
  \&H^\bullet(\widetilde\Delta,\mathbb Z_\ell)
       \arrow[r,"g_*"]
  \&H^{\bullet-2}(\Delta,\mathbb Z_\ell(-1))
       \arrow[r]
  \&0.
\end{tikzcd}}
\end{equation}
In view of the splitting (4.2.1), this yields the isomorphism
\begin{equation}\tag{4.2.6}
f_*\oplus(-g_*i^*):
H^\bullet(\widetilde X,\mathbb Z_\ell)
\xrightarrow{\sim}
H^\bullet(X,\mathbb Z_\ell)\oplus
H^{\bullet-2}(\Delta,\mathbb Z_\ell(-1)).
\end{equation}
The transpose of (4.2.5), after the degree reversal
$\bullet\mapsto2n-\bullet$, gives the commutative diagram of exact sequences
\begin{equation}\tag{4.2.7}
\resizebox{\textwidth}{!}{
\begin{tikzcd}[ampersand replacement=\&,column sep=large,row sep=large]
0
  \&H^\bullet(X,\mathbb Z_\ell)\arrow[l]
  \&H^\bullet(\widetilde X,\mathbb Z_\ell)\arrow[l,"f_*"']
  \&H^{\bullet-2}(\Delta,\mathbb Z_\ell(-1))
       \arrow[l,"i_*g^*"']
  \&0\arrow[l]\\
0
  \&H^{\bullet-4}(\Delta,\mathbb Z_\ell(-2))
       \arrow[l]\arrow[u,"j_*"']
  \&H^{\bullet-2}(\widetilde\Delta,\mathbb Z_\ell(-1))
       \arrow[l,"g_*"']\arrow[u,"i_*"']
  \&H^{\bullet-2}(\Delta,\mathbb Z_\ell(-1))
       \arrow[l,"g^*"']\arrow[u,equal]
  \&0\arrow[l]
\end{tikzcd}}
\end{equation}
and hence the isomorphism
\begin{equation}\tag{4.2.8}
H^\bullet(X,\mathbb Z_\ell)\oplus
H^{\bullet-2}(\Delta,\mathbb Z_\ell(-1))
\xrightarrow[\sim]{\ f^*+i_*g^*\ }
H^\bullet(\widetilde X,\mathbb Z_\ell).
\end{equation}


It remains to prove that the composite
\begin{equation}\tag{4.2.9}
\begin{aligned}
&H^\bullet(X,\mathbb Z_\ell)\oplus
 H^{\bullet-2}(\Delta,\mathbb Z_\ell(-1))
 \xrightarrow{\ f^*+i_*g^*\ }
 H^\bullet(\widetilde X,\mathbb Z_\ell)\\[-2pt]
&\hspace{55mm}\xrightarrow{\ f_*\oplus(-g_*i^*)\ }
 H^\bullet(X,\mathbb Z_\ell)\oplus
 H^{\bullet-2}(\Delta,\mathbb Z_\ell(-1))
\end{aligned}
\end{equation}
is the identity.  Now
\[
f_*f^*=\mathrm{id}\quad\text{on }H^\bullet(X,\mathbb Z_\ell)
\quad\text{by (2.2.18)},
\]
\[
g_*i^*f^*=0\quad\text{by (4.2.5)},
\qquad
f_*i_*g^*=0\quad\text{by (4.2.7)}.
\]
It therefore suffices to prove
\begin{equation}\tag{4.2.10}
-g_*i^*i_*g^*(x)=x
\quad\text{for }x\in H^\bullet(\Delta,\mathbb Z_\ell(-1)).
\end{equation}

To verify (4.2.10), first observe that
$i:\widetilde\Delta\hookrightarrow\widetilde X$ is the inclusion of a
smooth subvariety of codimension one.  Thus (SGA 5 VII)
\begin{equation}\tag{4.2.11}
i^*i_*(x)=c_1(\underline N_{\widetilde\Delta/\widetilde X})\,x,
\end{equation}
and
\begin{equation}\tag{4.2.12}
i_*(x)i_*(y)
=i_*\!\left(xy\,c_1(\underline N_{\widetilde\Delta/\widetilde X})\right).
\end{equation}
Write
$\xi_X=\underline{\mathcal O}_X(1)$,
$\xi_\Delta=\underline{\mathcal O}_\Delta(1)$, and
$\xi_{\mathbb P^1}=\underline{\mathcal O}_{\mathbb P^1}(1)$.
By EGA II 8.1.7, if $\widetilde{\mathfrak J}$ is the ideal of
$\widetilde\Delta$ in $\widetilde X$, then
\begin{equation}\tag{4.2.13}
\begin{aligned}
\underline{\check N}_{\widetilde\Delta/\widetilde X}
&\simeq\widetilde{\mathfrak J}/\widetilde{\mathfrak J}^{,2}
 \simeq\widetilde{\mathfrak J}
       \otimes_{\underline{\mathcal O}_{\widetilde X}}
       \underline{\mathcal O}_{\widetilde\Delta}\\
&\simeq\underline{\mathcal O}_{\widetilde X}(1)
       \otimes_{\underline{\mathcal O}_{\widetilde X}}
       \underline{\mathcal O}_{\widetilde\Delta}
 =\underline{\mathcal O}_{\widetilde\Delta}(1).
\end{aligned}
\end{equation}
Here $\widetilde\Delta$ is viewed as
$\mathbb P(\mathfrak J/\mathfrak J^2)$, where $\mathfrak J$ is the ideal of
$\Delta$ in $X$.  Since $\Delta$ is the intersection of $X$ with two given
hyperplanes, there is a canonical isomorphism
\[
\mathfrak J/\mathfrak J^2\simeq
\underline{\mathcal O}_\Delta^{,2}(-1),
\]
which induces
\[
\widetilde\Delta=\mathbb P(\mathfrak J/\mathfrak J^2)
\simeq\mathbb P(\underline{\mathcal O}_\Delta^{,2}(-1))
\xrightarrow[\sim]{\alpha}
\mathbb P(\underline{\mathcal O}_\Delta^{,2})
\simeq\Delta\times\mathbb P^1.
\]
The canonical isomorphism $\alpha$ of EGA II 4.1.4 carries the tautological
line bundle to
\[
(\rho i)^*\underline{\mathcal O}_{\mathbb P^1}(1)
\otimes g^*\underline{\mathcal O}_\Delta(-1).
\]
Consequently, (4.2.13) becomes
\begin{equation}\tag{4.2.14}
\underline{\check N}_{\widetilde\Delta/\widetilde X}
\simeq(\rho i)^*\underline{\mathcal O}_{\mathbb P^1}(1)
       \otimes g^*\underline{\mathcal O}_\Delta(-1).
\end{equation}

\subsection*{Lemma 4.3.}
\underline{One has}
\[
c_1(\underline N_{\widetilde\Delta/\widetilde X})
=g^*c_1(\xi_\Delta)-(\rho i)^*c_1(\xi_{\mathbb P^1}).
\]
Return now to the proof of (4.2.10).  By (4.2.11) and Lemma 4.3, it is enough
to prove
\[
g_*\!\left(g^*(x)\left[(\rho i)^*c_1(\xi_{\mathbb P^1})
                       -g^*c_1(\xi_\Delta)\right]\right)=x.
\]
By (4.2.4), $g_*g^*=0$, while
\[
\begin{tikzcd}[column sep=huge]
\widetilde\Delta\arrow[r,"g\times\rho_i","\sim"']
  \&\Delta\times\mathbb P^1,
\end{tikzcd}
\]
so
\begin{equation}\tag{4.3.1}
g_*\!\left(g^*(x)(\rho i)^*c_1(\xi_{\mathbb P^1})\right)=x,
\end{equation}
and hence
\begin{equation}\tag{4.3.2}
g_*\!\left(g^*(x)c_1(\underline N_{\widetilde\Delta/\widetilde X})\right)
=-x.
\end{equation}

We now compute the algebra structure on
\[
H^\bullet(X,\mathbb Z_\ell)\oplus
H^{\bullet-2}(\Delta,\mathbb Z_\ell(-1))
\]
transported from the cup product by (4.2.2).  Let $a,b$ belong to
$H^\bullet(X,\mathbb Z_\ell)$ and $x,y$ to
$H^{\bullet-2}(\Delta,\mathbb Z_\ell(-1))$.  Begin with (4.2.3.1), and write
\[
(0\oplus x)\wedge(0\oplus y)=A\oplus B.
\]
By (4.2.12) and (4.3.2),
\[
\begin{aligned}
A
&=f_*\!\left((i_*g^*x)(i_*g^*y)\right)\\
&=f_*i_*\!\left(g^*(xy)c_1(\underline N_{\widetilde\Delta/\widetilde X})\right)\\
&=j_*g_*\!\left(g^*(xy)c_1(\underline N_{\widetilde\Delta/\widetilde X})\right)
=-j_*(xy),
\end{aligned}
\]
and
\[
\begin{aligned}
B
&=-g_*i^*\!\left((i_*g^*x)(i_*g^*y)\right)\\
&=-g_*i^*i_*\!\left(g^*(xy)c_1(\underline N_{\widetilde\Delta/\widetilde X})\right)\\
&=-g_*\!\left(g^*(xy)c_1(\underline N_{\widetilde\Delta/\widetilde X})^2\right)
\qquad\text{by (4.2.11)}.
\end{aligned}
\]
Lemma 4.3 gives
\[
\begin{aligned}
c_1(\underline N_{\widetilde\Delta/\widetilde X})^2
&=g^*\!\left(c_1(\xi_\Delta)^2\right)
 -2g^*c_1(\xi_\Delta)(\rho i)^*c_1(\xi_{\mathbb P^1}),
\end{aligned}
\]
because $c_1(\xi_{\mathbb P^1})^2=0$.  Therefore, by (4.2.4) and (4.3.1),
\[
B=0+2xy\,c_1(\xi_\Delta),
\]
which proves (4.2.3.1).  Formula (4.2.3.2) expresses the fact that $f^*$ is
an algebra homomorphism.

To prove (4.2.3.3), the projection formula shows that
$\operatorname{Ker}f_*$ is stable under multiplication by
$\operatorname{im}f^*$.  Hence
\[
(a\oplus0)\wedge(0\oplus y)=0\oplus C,
\]
where
\begin{equation}\tag{4.3.3}
C=-g_*i^*\!\left(f^*(a)i_*g^*(y)\right).
\end{equation}
Using (4.2.11) and (4.3.2),
\[
\begin{aligned}
C
&=-g_*\!\left(i^*f^*(a)g^*(y)
      c_1(\underline N_{\widetilde\Delta/\widetilde X})\right)\\
&=-g_*\!\left(g^*j^*(a)g^*(y)
      c_1(\underline N_{\widetilde\Delta/\widetilde X})\right)\\
&=j^*(a)y.
\end{aligned}
\]
This proves (4.2.3.3).  Finally, (4.2.3.4) follows from (4.2.3.3) by graded
commutativity of the cup product.


}{E7II-XVIII-010}\SGApage{II}{\(\approx\)278}
\section*{5. Cohomological study of the morphism
$\rho:\widetilde X\longrightarrow\mathbb P^1$ associated with a pencil}

\subsection*{5.1. The restriction homomorphism
$H^*(\widetilde X)\longrightarrow H^*(Y)$}

Let $Y$ denote one of the smooth fibers of
$\rho:\widetilde X\to\mathbb P^1$, say
\[
Y=\rho^{-1}(t_0)=X\mathbin{\cdot}H_{t_0}.
\]
We have the commutative diagram
\begin{center}
\resizebox{\textwidth}{!}{
\begin{tikzcd}[ampersand replacement=\&,column sep=large,row sep=huge]
\mathbb P^1
  \& {} \& \widetilde X\arrow[ll,"\rho"']\arrow[d,"f"]
  \& {} \& \widetilde\Delta
       \arrow[ll,hook',"i"]\arrow[d,"g"]
       \arrow[r,phantom,"\simeq"]
  \& \Delta\times\mathbb P^1 \\
\mathrm{Sp}(k)=\{t_0\}\arrow[u,hook]
  \& {} \& X
  \& {} \& \Delta
       \arrow[ll,hook',"j"]
       \arrow[ur,hook,"{\ell=\mathrm{id}\times\{t_0\}}"']
       \arrow[dlll,hook,"h"]
  \& {} \\
{} \& Y\arrow[ul]\arrow[uur,"k"]\arrow[ur,"m"']
  \& {} \& {} \& {} \& {}
\end{tikzcd}}
\end{center}

\SGApage{II}{279}
\subsection*{Proposition 5.1.1.}
\underline{Under the hypotheses of (3.1), with $Y$ denoting a smooth
fiber:}

\SGApage{II}{\(\approx\)278}
\SGAvar{\smallskip
\noindent
\textup{(5.1.2)} \underline{The restriction homomorphism}
\[
k^*:H^\bullet(\widetilde X,\mathbb Z_\ell)
   \longrightarrow H^\bullet(Y,\mathbb Z_\ell)
\]
\underline{is expressed (via (4.2.2)) by}
\[
H^\bullet(X,\mathbb Z_\ell)\oplus
H^{\bullet-2}(\Delta,\mathbb Z_\ell(-1))
\xrightarrow{\ m^*+h^*\ }
H^\bullet(Y,\mathbb Z_\ell).
\]

}{\smallskip
\noindent
\textup{(5.1.2)} \underline{The restriction homomorphism}
\[
k^*:H^\bullet(\widetilde X,\mathbb Z_\ell)
   \longrightarrow H^\bullet(Y,\mathbb Z_\ell)
\]
\underline{is expressed (via (4.2.2)) by}
\[
H^\bullet(X,\mathbb Z_\ell)\oplus
H^{\bullet-2}(\Delta,\mathbb Z_\ell(-1))
\xrightarrow{\ m^*+h_*\ }
H^\bullet(Y,\mathbb Z_\ell).
\]

}{E7II-XVIII-011}\SGApage{II}{\(\approx\)279}
\noindent
\textup{(5.1.3)} \underline{The Gysin map}
\[
k_*:H^{\bullet-2}(Y,\mathbb Z_\ell(-1))
   \longrightarrow H^\bullet(\widetilde X,\mathbb Z_\ell)
\]
\underline{is expressed by}
\[
H^{\bullet-2}(Y,\mathbb Z_\ell(-1))
\xrightarrow{\ m_*\oplus(-h^*)\ }
H^\bullet(X,\mathbb Z_\ell)\oplus
H^{\bullet-2}(\Delta,\mathbb Z_\ell(-1)).
\]

\underline{Proof.}
The smooth subschemes $\widetilde\Delta$ and $Y$ of $\widetilde X$
intersect transversely in $\Delta$; in other words, the square of smooth
schemes
\begin{equation}\tag{5.1.4}
\begin{tikzcd}[column sep=large,row sep=large]
\Delta\arrow[r,"\ell"]\arrow[d,hook,"h"]
  \&\widetilde\Delta\arrow[d,hook,"i"]\\
Y\arrow[r,"k"']\&\widetilde X
\end{tikzcd}
\end{equation}
is Cartesian.  It therefore implies the formula [SGA 5, IV]
\begin{equation}\tag{5.1.5}
h_*\ell^*=k^*i_* ,\qquad \ell_*h^*=i^*k_*.
\end{equation}
The composite
\[
H^\bullet(X,\mathbb Z_\ell)\oplus
H^{\bullet-2}(\Delta,\mathbb Z_\ell(-1))
\xrightarrow{\ f^*+i_*g^*\ }
H^\bullet(\widetilde X,\mathbb Z_\ell)
\xrightarrow{\ k^*\ }
H^\bullet(Y,\mathbb Z_\ell)
\]
is the sum of
$k^*f^*=m^*$ and
$k^*i_*g^*=h_*\ell^*g^*$ by (5.1.5).  Since
$g\ell=\mathrm{id}_\Delta$, this proves (5.1.2).

\SGApage{II}{\(\approx\)280}
\SGAvar{The composite
\[
H^{\bullet-2}(Y,\mathbb Z_\ell(-1))
\xrightarrow{\ k_*\ }
H^\bullet(\widetilde X,\mathbb Z_\ell)
\xrightarrow{\ f_*\oplus(g_*i^*)\ }
H^\bullet(X,\mathbb Z_\ell)\oplus
H^{\bullet-2}(\Delta,\mathbb Z_\ell(-1))
\]
is the sum of
\[
f_*k_*=m_*
\qquad\text{and}\qquad
-g_*i^*k_*=g_*\ell_*h^*=-h^*,
\]
which proves (5.1.3).

}{The composite
\[
H^{\bullet-2}(Y,\mathbb Z_\ell(-1))
\xrightarrow{\ k_*\ }
H^\bullet(\widetilde X,\mathbb Z_\ell)
\xrightarrow{\ f_*\oplus(-g_*i^*)\ }
H^\bullet(X,\mathbb Z_\ell)\oplus
H^{\bullet-2}(\Delta,\mathbb Z_\ell(-1))
\]
is the sum of
\[
f_*k_*=m_*
\qquad\text{and}\qquad
-g_*i^*k_*=-g_*\ell_*h^*=-h^*,
\]
which proves (5.1.3).

}{E7II-XVIII-012}\subsection*{Corollary 5.1.6.}
\underline{Under the hypotheses of 5.1.1, for $i\leq n-1$, the images
of}
\[
m^*:H^i(X,\mathbb Z_\ell)\longrightarrow H^i(Y,\mathbb Z_\ell)
\]
\underline{and of}
\[
k^*:H^i(\widetilde X,\mathbb Z_\ell)
   \longrightarrow H^i(Y,\mathbb Z_\ell)
\]
\underline{are equal.}

\SGAvar{\underline{Proof.}
By the weak Lefschetz theorem for hyperplane sections (SGA 5, VII 7.1),
for $i\leq n-1$ the map
\begin{equation}\tag{5.1.7}
H^{i-2}(X,\mathbb Z_\ell(-1))
\xrightarrow{\ j^*\ }
H^{i-2}(\Delta,\mathbb Z_\ell(-1))
\end{equation}
is an isomorphism.  Thus, by (5.1.2), the image of $h^*$ is the sum of
the image of $m^*$ and the image of the composite
\[
H^{i-2}(X,\mathbb Z_\ell(-1))
\xrightarrow{\ j^*\ }
H^{i-2}(\Delta,\mathbb Z_\ell(-1))
\xrightarrow{\ h^*\ }
H^i(Y,\mathbb Z_\ell).
\]
But the projection formula gives
\[
\begin{aligned}
h_*j^*(x)
 &=h_*\bigl(h^*m^*(x)
  =m^*(x)h_*(1)
  =\xi_Y\,m^*(x)\\
 &=m^*(x\xi_X),
 \qquad\text{since }m^*(\xi_X)=\xi_Y.
\end{aligned}
\]
Consequently $\operatorname{Im}k^*\subset\operatorname{Im}m^*$, and
the reverse inclusion follows directly from (5.1.2).  This proves the
corollary.


}{\underline{Proof.}
By the weak Lefschetz theorem for hyperplane sections (SGA 5, VII 7.1),
for $i\leq n-1$ the map
\begin{equation}\tag{5.1.7}
H^{i-2}(X,\mathbb Z_\ell(-1))
\xrightarrow{\ j^*\ }
H^{i-2}(\Delta,\mathbb Z_\ell(-1))
\end{equation}
is an isomorphism.  Thus, by (5.1.2), the image of $k^*$ is the sum of
the image of $m^*$ and the image of the composite
\[
H^{i-2}(X,\mathbb Z_\ell(-1))
\xrightarrow{\ j^*\ }
H^{i-2}(\Delta,\mathbb Z_\ell(-1))
\xrightarrow{\ h_*\ }
H^i(Y,\mathbb Z_\ell).
\]
But the projection formula gives
\[
\begin{aligned}
h_*j^*(x)
 &=h_*\bigl(h^*m^*(x)\bigr)
  =m^*(x)h_*(1)
  =\xi_Y\,m^*(x)\\
 &=m^*(x\xi_X),
 \qquad\text{since }m^*(\xi_X)=\xi_Y.
\end{aligned}
\]
Consequently $\operatorname{Im}k^*\subset\operatorname{Im}m^*$, and
the reverse inclusion follows directly from (5.1.2).  This proves the
corollary.


}{E7II-XVIII-013}\SGApage{II}{281}
\subsection*{5.2. The hard Lefschetz theorem (LV)}

\SGAvar{\subsubsection*{5.2.1.}
Let $S$ be projective of dimension $n$ over the algebraically closed field
$k$, and let $S\hookrightarrow\mathbb P$ be a projective embedding.  Denote by
\[
u\in H^2(S,\mathbb Z_\ell(1))
\]
the cohomology class of a hyperplane section.  Thus $u$ is the pullback along
$S\hookrightarrow\mathbb P$ of the first Chern class of
$\underline{\mathcal O}_{\mathbb P}(1)$ in
$H^2(\mathbb P,\mathbb Z_\ell(1))$.  Write
\[
L_S:H^\bullet(S,\mathbb Q_\ell(j)
   \longrightarrow H^{\bullet+2}(S,\mathbb Q_\ell(j+1))
\]
for multiplication by $u$.

}{\subsubsection*{5.2.1.}
Let $S$ be projective of dimension $n$ over the algebraically closed field
$k$, and let $S\hookrightarrow\mathbb P$ be a projective embedding.  Denote by
\[
u\in H^2(S,\mathbb Z_\ell(1))
\]
the cohomology class of a hyperplane section.  Thus $u$ is the pullback along
$S\hookrightarrow\mathbb P$ of the first Chern class of
$\underline{\mathcal O}_{\mathbb P}(1)$ in
$H^2(\mathbb P,\mathbb Z_\ell(1))$.  Write
\[
L_S:H^\bullet(S,\mathbb Q_\ell(j))
   \longrightarrow H^{\bullet+2}(S,\mathbb Q_\ell(j+1))
\]
for multiplication by $u$.

}{E7II-XVIII-014}\subsubsection*{5.2.2.}
\underline{We say that $S$ (equipped with} $u$\underline{) satisfies (LV), the hard
Lefschetz theorem, if, for every (or for one) integer $j$ and every integer
$i>0$, the maps}
\[
L_S^i:H^{n-i}(S,\mathbb Q_\ell(j))
\longrightarrow H^{n+i}(S,\mathbb Q_\ell(i+j))
\]
\underline{are isomorphisms.}

\subsubsection*{5.2.2.1.}
It is conjectured that this condition holds when $S$ is smooth.  It is true
when $k$ has characteristic zero [6].  The specialization theorem
[SGA 4, XVI 2.2] therefore gives the same conclusion when $S$, together with
its projective embedding, lifts to characteristic zero.  In particular, it
holds for a smooth complete intersection.  It is also known for a smooth $S$
of dimension at most two [2], and for an abelian variety, either directly
(Lieberman [4] and [2]) or by Mumford's theorem [5] that a polarized abelian
variety lifts, together with its polarization, to characteristic zero.

\SGApage{II}{\(\approx\)282}
\SGAvar{Recall that, for a projective variety $S$ and $i\leq n$, the primitive part of
cohomology is defined by
\begin{equation}\tag{5.2.2.2}
\operatorname{Prim}^i(S,\mathbb Q_\ell(j))
=\operatorname{Ker}\!\left(
L^{n-i+1}:H^i(S,\mathbb Q_\ell(j))
\longrightarrow
H^{2n-i+2}(S,\mathbb Q_\ell(j+2n-2i+2))
\right..
\end{equation}
If $S$ satisfies (LV), then for $i\leq n$ one has the decomposition
\begin{equation}\tag{5.2.2.3}
H^i(S,\mathbb Q_\ell(j))
\xleftarrow[\sim]{\ 1+L\ }
\operatorname{Prim}^i(S,\mathbb Q_\ell(j))
\oplus H^{i-2}(S,\mathbb Q_\ell(j-1)).
\end{equation}

}{Recall that, for a projective variety $S$ and $i\leq n$, the primitive part of
cohomology is defined by
\begin{equation}\tag{5.2.2.2}
\operatorname{Prim}^i(S,\mathbb Q_\ell(j))
=\operatorname{Ker}\!\left(
L^{n-i+1}:H^i(S,\mathbb Q_\ell(j))
\longrightarrow
H^{2n-i+2}(S,\mathbb Q_\ell(j+n-i+1))
\right).
\end{equation}
If $S$ satisfies (LV), then for $i\leq n$ one has the decomposition
\begin{equation}\tag{5.2.2.3}
H^i(S,\mathbb Q_\ell(j))
\xleftarrow[\sim]{\ 1+L\ }
\operatorname{Prim}^i(S,\mathbb Q_\ell(j))
\oplus H^{i-2}(S,\mathbb Q_\ell(j-1)).
\end{equation}

}{E7II-XVIII-015}\subsection*{Proposition 5.2.3.}
\underline{Let $X$ be smooth and projective and suppose that it satisfies
(LV).  Let $\tau:Y\hookrightarrow X$ be the inclusion of a smooth hyperplane
section.  Then}
\begin{equation}\tag{5.2.3.1}
\operatorname{Prim}^n(X,\mathbb Q_\ell(j))
=\operatorname{Ker}\!\left(
\tau^*:H^n(X,\mathbb Q_\ell(j))
\longrightarrow H^n(Y,\mathbb Q_\ell(j))
\right).
\end{equation}

\SGAvar{\underline{Proof.}
The projection formula for $m:Y\to X$ gives the commutative diagram
\begin{equation}\tag{5.2.3.2}
\begin{tikzcd}[column sep=huge,row sep=huge]
H^n(X,\mathbb Q_\ell(j))
  \arrow[rr,"L"]\arrow[dr,"\tau^*"']
  \& {} \&
H^{n+2}(X,\mathbb Q_\ell(j+1))\arrow[dl, "\tau_*"]\\
{} \& H^n(Y,\mathbb Q_\ell(j))  \& {}
\end{tikzcd}
\end{equation}
By duality and weak Lefschetz,
\[
\tau_*:H^n(Y,\mathbb Q_\ell(j))
\longrightarrow H^{n+2}(X,\mathbb Q_\ell(j+1))
\]
is the transpose of an isomorphism and hence is itself an isomorphism.  The
assertion follows.

}{\underline{Proof.}
The projection formula for $\tau:Y\to X$ gives the commutative diagram
\begin{equation}\tag{5.2.3.2}
\begin{tikzcd}[column sep=huge,row sep=huge]
H^n(X,\mathbb Q_\ell(j))
  \arrow[rr,"L"]\arrow[dr,"\tau^*"']
  \& {} \&
H^{n+2}(X,\mathbb Q_\ell(j+1))\\
{} \& H^n(Y,\mathbb Q_\ell(j)) \arrow[ur,"\tau_*"'] \& {}
\end{tikzcd}
\end{equation}
By duality and weak Lefschetz,
\[
\tau_*:H^n(Y,\mathbb Q_\ell(j))
\longrightarrow H^{n+2}(X,\mathbb Q_\ell(j+1))
\]
is the transpose of an isomorphism and hence is itself an isomorphism.  The
assertion follows.

}{E7II-XVIII-016}\SGApage{II}{283}
\subsection*{Lemma 5.2.4.}
\underline{Let $X$ be smooth and projective, and let
$\tau:Y\hookrightarrow X$ be the inclusion of a smooth hyperplane section.
Then $X$ satisfies (LV) if and only if the following two conditions hold:}

\smallskip
\noindent
\textup{5.2.4.1.} \underline{$Y$ satisfies (LV).}

\smallskip
\noindent
\textup{5.2.4.2.} \underline{The restriction of the cup product on
$H^{n-1}(Y,\mathbb Q_\ell)$ to the subspace
$\operatorname{Im}H^{n-1}(X,\mathbb Q_\ell)$ is nondegenerate.}

\underline{Proof.}
We have
\begin{equation}\tag{5.2.4.3}
\left\{
\begin{aligned}
L_X&=\tau_*\tau^*,\\[1ex]
L_Y&=\tau^*\tau_* \qquad\text{(by 4.2.11),}
\end{aligned}
\right.
\end{equation}
and consequently
\begin{equation}\tag{5.2.4.4}
L_X^{i+1}=\tau_*L_Y^i\tau^*.
\end{equation}
In the commutative diagram
\begin{equation}\tag{5.2.4.5}
\begin{tikzcd}[column sep=huge,row sep=huge]
H^{n-i-1}(X,\mathbb Q_\ell(j))
  \arrow[r,"L_X^{i+1}"]\arrow[d,"\tau^*"]
  \&H^{n+i+1}(X,\mathbb Q_\ell(j+i+1))\\
H^{n-1-i}(Y,\mathbb Q_\ell(j))
  \arrow[r,"L_Y^i"]
  \&H^{n-1+i}(Y,\mathbb Q_\ell(j+i))\arrow[u,"\tau_*"']
\end{tikzcd}
\end{equation}
$\tau^*$ is an isomorphism for $i\geq1$ by weak Lefschetz, and $\tau_*$ is
an isomorphism by duality.  Hence $X$ satisfies (LV) if and only if $Y$
satisfies (LV) and
\begin{equation}\tag{5.2.4.6}
L_X:H^{n-1}(X,\mathbb Q_\ell(j))
\longrightarrow H^{n+1}(X,\mathbb Q_\ell(j+1))
\end{equation}
is an isomorphism.

\SGApage{II}{\(\approx\)284}
\SGAvar{It remains to identify (5.2.4.6) with 5.2.4.2.  Let
$a,b\in H^{n-1}(X,\mathbb Q_\ell(j))$.  By Poincaré duality on $X$,
(5.2.4.6) is equivalent to nondegeneracy of the bilinear form
\begin{equation}\notag
\begin{aligned}
H^{n-1}(X,\mathbb Q_\ell(j))\times
H^{n-1}(X,\mathbb Q_\ell(j))
&\longrightarrow H^{2n}(X,\mathbb Q_\ell(2j+2),\\
(a,b)&\longmapsto a\,L_X(b).
\end{aligned}
\end{equation}
But by (5.2.5) and the projection formula,
\[
aL_X(b)=a\tau_*\tau^*(b)
=\tau_*\bigl(\tau^*(a)\tau^*(b)\bigr).
\]
This proves the equivalence of (5.2.4.7) and 5.2.4.2.

}{It remains to identify (5.2.4.6) with 5.2.4.2.  Let
$a,b\in H^{n-1}(X,\mathbb Q_\ell(j))$.  By Poincaré duality on $X$,
(5.2.4.6) is equivalent to nondegeneracy of the bilinear form
\begin{equation}\tag{5.2.4.7}
\begin{aligned}
H^{n-1}(X,\mathbb Q_\ell(j))\times
H^{n-1}(X,\mathbb Q_\ell(j))
&\longrightarrow H^{2n}(X,\mathbb Q_\ell(2j+1)),\\
(a,b)&\longmapsto a\,L_X(b).
\end{aligned}
\end{equation}
But by (5.2.4.3) and the projection formula,
\[
aL_X(b)=a\tau_*\tau^*(b)
=\tau_*\bigl(\tau^*(a)\tau^*(b)\bigr).
\]
This proves the equivalence of (5.2.4.7) and 5.2.4.2.

}{E7II-XVIII-017}\SGApage{II}{285}
\subsection*{Corollary 5.2.5.}
\underline{Let $X$ be smooth and projective of dimension $n$, and let
$\tau:Y\hookrightarrow X$ be the inclusion of a smooth hyperplane section.
Suppose that $X$ satisfies (LV).  Then, for $q\geq n$, the map}
\[
\tau^*:H^q(X,\mathbb Q_\ell(j))
\longrightarrow H^q(Y,\mathbb Q_\ell(j))
\]
\underline{is surjective.}

\underline{Proof.}
By (5.2.4.3), for $i\geq0$ one has
\[
L_Y^{i+1}=\tau^*L_X^i\tau_*,
\]
and hence the commutative diagram
\begin{center}
\begin{tikzcd}[column sep=huge,row sep=huge]
H^{n-i}(X,\mathbb Q_\ell(j+1))
  \arrow[r,"L_X^i"]
  \&H^{n+i}(X,\mathbb Q_\ell(j+i+1))\arrow[d,"\tau^*"]\\
H^{n-i-2}(Y,\mathbb Q_\ell(j))
  \arrow[u,"\tau_*"']\arrow[r,"L_Y^{i+1}"{yshift=7pt},"\sim"]
  \&H^{n+i}(Y,\mathbb Q_\ell(j+i+1)).
\end{tikzcd}
\end{center}
This gives the required result.

\SGApage{II}{286}
\subsection*{Lemma 5.2.6.}
\underline{Let $f:F\to B$ be a smooth proper morphism with $B$ connected.
Let $\underline{\mathcal O}_F(1)$ be a relatively ample invertible sheaf on
$F$, and equip each geometric fiber $F_{\bar b}$ with the ample invertible
sheaf induced by $\underline{\mathcal O}_F(1)$.  If one geometric fiber
satisfies (LV), then so do all the others.}

\SGAvar{\underline{Proof.}
Let
$\xi_F\in H^2(F,\mathbb Z_\ell(1))$ be the first Chern class of
$\underline{\mathcal O}_F(1)$.  This class and its powers
$\xi_F^i\in H^{ii}(F,\mathbb Z_\ell(i))$ induce, by cup product, the
homomorphisms
\begin{equation}\tag{5.2.6.1}
\xi_F^i:R^{n-i}f_*\mathbb Q_\ell(j)
\longrightarrow R^{n+i}f_*\mathbb Q_\ell(j+i),
\end{equation}
where $n$ is the relative dimension of the fibers.  By proper base change
[SGA 4, XII 5.2], all geometric fibers satisfy (LV) if and only if all the
maps (5.2.6.1) are isomorphisms.  By the specialization theorem
[SGA 4, XVI 2.2], the sheaves $R^if_*\mathbb Q_\ell(j)$ are twisted constant
sheaves.  Thus the maps (5.2.6.1) are isomorphisms if and only if they are so
over one given geometric point.

}{\underline{Proof.}
Let
$\xi_F\in H^2(F,\mathbb Z_\ell(1))$ be the first Chern class of
$\underline{\mathcal O}_F(1)$.  This class and its powers
$\xi_F^i\in H^{2i}(F,\mathbb Z_\ell(i))$ induce, by cup product, the
homomorphisms
\begin{equation}\tag{5.2.6.1}
\xi_F^i:R^{n-i}f_*\mathbb Q_\ell(j)
\longrightarrow R^{n+i}f_*\mathbb Q_\ell(j+i),
\end{equation}
where $n$ is the relative dimension of the fibers.  By proper base change
[SGA 4, XII 5.2], all geometric fibers satisfy (LV) if and only if all the
maps (5.2.6.1) are isomorphisms.  By the specialization theorem
[SGA 4, XVI 2.2], the sheaves $R^if_*\mathbb Q_\ell(j)$ are twisted constant
sheaves.  Thus the maps (5.2.6.1) are isomorphisms if and only if they are so
over one given geometric point.

}{E7II-XVIII-018}\subsection*{Corollary 5.2.7.}
\underline{Let $X$ be smooth and projective.  If one smooth hyperplane
section satisfies (LV), then every smooth hyperplane section does.}

\SGAvar{\underline{Proof.}
The universal family of smooth hyperplane sections of $X$ satisfies the
hypotheses of Lemma 5.2.5.


}{\underline{Proof.}
The universal family of smooth hyperplane sections of $X$ satisfies the
hypotheses of Lemma 5.2.6.


}{E7II-XVIII-019}\SGApage{II}{287}
\subsection*{5.3. Condition (A)}

\subsubsection*{5.3.1.}
Let
\[
\rho:\widetilde X\longrightarrow\mathbb P^1
\]
be the ``fibration'' associated with a pencil whose axis meets $X$
transversely and whose generic fiber is assumed to be smooth.

Let $U$ be a nonempty open subset of $\mathbb P^1$ over which $\rho$ is
smooth, and let
\begin{equation}\tag{5.3.1.1}
\nu:U\hookrightarrow\mathbb P^1
\end{equation}
be the inclusion.  Since the sheaves
$R^i\rho_*\mathbb Q_\ell(j)$ are \underline{twisted constant over $U$}, the
sheaves
\[
\nu_*\nu^*R^i\rho_*\mathbb Q_\ell(j)
\]
are seen to be essentially \underline{independent} of the choice of $U$;
the canonical ``adjunction'' homomorphisms
\begin{equation}\tag{5.3.2}
R^i\rho_*\mathbb Q_\ell(j)
\longrightarrow
\nu_*\nu^*R^i\rho_*\mathbb Q_\ell(j)
\end{equation}
are likewise independent of the choice of $U$.

\subsubsection*{5.3.3.}
Observe that, for an \underline{arbitrary} pencil of hyperplane sections,
the weak Lefschetz theorem implies that, for $i\leq n-2$, the immersion
\[
\widetilde X\overset{(f\times\rho)}{\hookrightarrow}
X\times\mathbb P^1
\]
gives an isomorphism of sheaves on $\mathbb P^1$
\[
R^i\rho_*\mathbb Q_\ell(j)
\xleftarrow{\ \sim\ }
H^i(X,\mathbb Q_\ell(j))_{\mathbb P^1},
\qquad i\leq n-2.
\]
In other words, $R^i\rho_*\mathbb Q_\ell(j)$ is constant for
$i\leq n-2$ and, a fortiori, the maps (5.3.2) are isomorphisms for
$i\leq n-2$.

\SGApage{II}{\(\approx\)288}
\subsubsection*{5.3.4.}
The local theory of vanishing cycles, and in particular the
\underline{Picard--Lefschetz formula} [SGA 7, XV 3.4], shows that,
\underline{for a Lefschetz pencil, the homomorphisms (5.3.2) are
isomorphisms for $i\neq n$, and are surjective for $i=n$}.

\subsubsection*{5.3.5.}
We say that $\rho$, or the pencil from which it arises,
\underline{satisfies condition (A)} if \underline{all} the homomorphisms
(5.3.2) are isomorphisms (equivalently, if this holds for every $i$ with
$j=0$).

We shall see below (6.3, 6.4) that condition (A) has a strong tendency to
hold.


\subsection*{Lemma 5.4.}
\underline{Suppose that one smooth hyperplane section of $X$} (\underline{and hence
every such section, by Corollary 5.2.7}) \underline{satisfies (LV)---for example, suppose
that $X$ satisfies (LV), cf. Lemma 5.2.4---and suppose that $\rho$ satisfies
condition (A).  Then every fiber} (\underline{including every singular fiber}) \underline{satisfies
(LV), and the sheaves $R^i\rho_*\mathbb Q_\ell(j)$ are constant for
$i\neq n-1$.}

\SGAvar{\underline{Proof.}
Let
\[
L:R^i\rho_*\mathbb Q_\ell(j)
\longrightarrow R^{i+2}\rho_*\mathbb Q_\ell(j+1)
\]
denote multiplication by
$f^*c_1(\mathcal O_X(1)\in
H^2(\widetilde X,\mathbb Z_\ell(1))$; fiber by fiber, this is the operation
considered in (5.2.1).  By hypothesis, over $U$ we have an isomorphism
\begin{equation}\tag{5.4.1}
\nu^*(L^i):
\nu^*R^{n-1-i}\rho_*\mathbb Q_\ell(j)
\overset{\sim}{\longrightarrow}
\nu^*R^{n-1+i}\rho_*\mathbb Q_\ell(j+i)
\end{equation}
for every $i>0$ and every $j$.  Applying the functor $\nu_*$ gives the
\underline{isomorphisms}
\begin{equation}\tag{5.4.2}
(\nu_*\nu^*)L^i:
\nu_*\nu^*R^{n-1-i}\rho_*\mathbb Q_\ell(j)
\overset{\sim}{\longrightarrow}
\nu_*\nu^*R^{n-1+i}\rho_*\mathbb Q_\ell(j+i).
\end{equation}
In view of (A), these identify with the morphisms
\begin{equation}\tag{5.4.3}
L^i:R^{n-1-i}\rho^*\mathbb Q_\ell(j)
\overset{\sim}{\longrightarrow}
R^{n-1+i}\rho_*\mathbb Q_\ell(j+i).
\end{equation}
The conclusion follows from 5.3.3.

}{\underline{Proof.}
Let
\[
L:R^i\rho_*\mathbb Q_\ell(j)
\longrightarrow R^{i+2}\rho_*\mathbb Q_\ell(j+1)
\]
denote multiplication by
$f^*c_1(\mathcal O_X(1))\in
H^2(\widetilde X,\mathbb Z_\ell(1))$; fiber by fiber, this is the operation
considered in (5.2.1).  By hypothesis, over $U$ we have an isomorphism
\begin{equation}\tag{5.4.1}
\nu^*(L^i):
\nu^*R^{n-1-i}\rho_*\mathbb Q_\ell(j)
\overset{\sim}{\longrightarrow}
\nu^*R^{n-1+i}\rho_*\mathbb Q_\ell(j+i)
\end{equation}
for every $i>0$ and every $j$.  Applying the functor $\nu_*$ gives the
\underline{isomorphisms}
\begin{equation}\tag{5.4.2}
(\nu_*\nu^*)L^i:
\nu_*\nu^*R^{n-1-i}\rho_*\mathbb Q_\ell(j)
\overset{\sim}{\longrightarrow}
\nu_*\nu^*R^{n-1+i}\rho_*\mathbb Q_\ell(j+i).
\end{equation}
In view of (A), these identify with the morphisms
\begin{equation}\tag{5.4.3}
L^i:R^{n-1-i}\rho_*\mathbb Q_\ell(j)
\overset{\sim}{\longrightarrow}
R^{n-1+i}\rho_*\mathbb Q_\ell(j+i).
\end{equation}
The conclusion follows from 5.3.3.

}{E7II-XVIII-020}\SGApage{II}{\(\approx\)289}
\subsubsection*{5.4.4.}
Let $Y$ be a smooth fiber of
$\rho:\widetilde X\to\mathbb P^1$.  By weak Lefschetz, the restriction map
\[
H^{n-1}(X,\mathbb Q_\ell(j))
\overset{m^*}{\hookrightarrow}
H^{n-1}(Y,\mathbb Q_\ell(j))
\]
is injective.  We saw in 5.2.4 that, if $X$ satisfies (LV), then the cup
product on $H^{n-1}(Y,\mathbb Q_\ell(j))$ remains nondegenerate when
restricted to $m^*H^{n-1}(X,\mathbb Q_\ell(j))$.  The orthogonal complement
of $H^{n-1}(X,\mathbb Q_\ell(j))$ in
$H^{n-1}(Y,\mathbb Q_\ell(j))$ is therefore a complementary subspace to
$H^{n-1}(X,\mathbb Q_\ell(j))$ in
$H^{n-1}(Y,\mathbb Q_\ell(j))$; following the classical terminology, it is
called the \underline{vanishing-cohomology space of $Y$}.

Let $\nu:U\hookrightarrow\mathbb P^1$ be any nonempty open subset over
which $\rho$ is smooth.  The corresponding subsheaf of
$\nu_*\nu^*R^{n-1}\rho_*\mathbb Q_\ell(j)$, the
``\underline{vanishing-cohomology sheaf}''
$E^{n-1}\rho_*\mathbb Q_\ell(j)$ on $\mathbb P^1$, is defined by
\begin{equation}\tag{5.4.4.1}
E^{n-1}\rho_*\mathbb Q_\ell(j)
=\nu_*\left\{
\begin{array}{l}
\text{the orthogonal complement in }
\nu^*R^{n-1}\rho_*\mathbb Q_\ell(j)\text{ of the}\\[1ex]
\text{constant subsheaf }H^{n-1}(X,\mathbb Q_\ell(j))_U
\end{array}
\right\}.
\end{equation}
It is clear that this sheaf is independent of the choice of $U$.

\SGApage{II}{290}
\subsection*{Lemma 5.5.}
\underline{Suppose that $X$ satisfies (LV) and that the maps (5.3.2) are
isomorphisms for $i=n-1$---for example, suppose that $\rho$ arises from a
Lefschetz pencil, cf. 5.3.4.  Then there is a direct-sum decomposition}
\begin{equation}\tag{5.5.1}
R^{n-1}\rho_*\mathbb Q_\ell(j)
\simeq
H^{n-1}(X,\mathbb Q_\ell(j))_{\mathbb P^1}
\oplus E^{n-1}\rho_*\mathbb Q_\ell(j).
\end{equation}

\underline{Proof.}
By (LV), over $U$ there is a decomposition
\begin{equation}\tag{5.5.2}
\nu^*R^{n-1}\rho_*\mathbb Q_\ell(j)
\simeq
H^{n-1}(X,\mathbb Q_\ell(j))_U
\oplus \nu^*E^{n-1}\rho_*\mathbb Q_\ell(j),
\end{equation}
to which we apply $\nu_*$.  This proves the assertion.


\SGAvar{\subsection*{Theorem 5.6.}
\underline{Let $X$ be a smooth, irreducible, projective variety satisfying
(LV) (5.2.2).  Let $\{H_\tau\}$ be a pencil of hyperplanes whose generic
member and axis meet $X$ transversely, in $X_\eta$ and $\Delta$ respectively.
Suppose that the projection
$\rho:\widetilde X\to\mathbb P^1$ satisfies condition (5.3.2).  Then:}


\smallskip
\noindent
\textup{(5.6.1)}\quad
\underline{the Leray spectral sequence}
\[
E_2^{p,q}=H^p(\mathbb P^1,R^q\rho_*\mathbb Q_\ell)
\Longrightarrow H^{m+q}(\widetilde X,\mathbb Q_\ell)
\]

}{\subsection*{Theorem 5.6.}
\underline{Let $X$ be a smooth, irreducible, projective variety satisfying
(LV) (5.2.2).  Let $\{H_\tau\}$ be a pencil of hyperplanes whose generic
member and axis meet $X$ transversely, in $X_\eta$ and $\Delta$ respectively.
Suppose that the projection
$\rho:\widetilde X\to\mathbb P^1$ satisfies condition (A).  Then:}


\smallskip
\noindent
\textup{(5.6.1)}\quad
\underline{the Leray spectral sequence}
\[
E_2^{p,q}=H^p(\mathbb P^1,R^q\rho_*\mathbb Q_\ell)
\Longrightarrow H^{p+q}(\widetilde X,\mathbb Q_\ell)
\]

}{E7II-XVIII-021}\underline{degenerates.}

\noindent
\textup{(Note.)}\quad
\underline{This remains valid if, instead of assuming that $X$ satisfies
(LV), one assumes only that a smooth hyperplane section $Y$ does.}

\SGApage{II}{\(\approx\)291}
\subsubsection*{5.6.2.}
\underline{Let $\eta$ denote the generic point of $\mathbb P^1$,
$\bar\eta$ the geometric generic point, $X_{\bar\eta}$ the geometric generic
fiber of $\rho$, and $\pi$ the Galois group of
$k(\bar\eta)/k(\eta)$.}

\SGAvar{\underline{For $q\neq n-1$ there are canonical isomorphisms}
\begin{align}
E_2^{o,q}
 &=H^o(\mathbb P^1,R^q\rho_*\mathbb Q_\ell)
 \simeq H^q(X_{\bar\eta},\mathbb Q_\ell),
 \tag{5.6.2.1}\\
E_2^{2,q}
 &=H^2(\mathbb P^1,R^q\rho_*\mathbb Q_\ell)
 \simeq H^q(X_{\bar\eta},\mathbb Q_\ell(-1)),
 \tag{5.6.2.2}\\
E_2^{p,q}
 &=H^p(\mathbb P^1,R^q\rho_*\mathbb Q_\ell)=0,
 \qquad p\neq0,2.
 \tag{5.6.2.3}
\end{align}
\underline{For $q=n-1$ one has}
\begin{align}
E_2^{0,n-1}
 &=H^o(\mathbb P^1,R^{n-1}\rho_*\mathbb Q_\ell)
 \simeq H^{n-1}(X_{\bar\eta},\mathbb Q_\ell)^\pi
 \simeq H^{n-1}(X,\mathbb Q_\ell),
 \tag{5.6.2.4}\\
E_2^{2,n-1}(
 &=H^2(\mathbb P^1,R^{n-1}\rho_*\mathbb Q_\ell)
 \simeq H^{n-1}(X_{\bar\eta},\mathbb Q_\ell(-1))^\pi
 \simeq H^{n-1}(X,\mathbb Q_\ell(-1)).
 \tag{5.6.2.5}
\end{align}
(\underline{The missing value $E_2^{1,n-1}$ will be found in 5.7 below}).

}{\underline{For $q\neq n-1$ there are canonical isomorphisms}
\begin{align}
E_2^{0,q}
 &=H^0(\mathbb P^1,R^q\rho_*\mathbb Q_\ell)
 \simeq H^q(X_{\bar\eta},\mathbb Q_\ell),
 \tag{5.6.2.1}\\
E_2^{2,q}
 &=H^2(\mathbb P^1,R^q\rho_*\mathbb Q_\ell)
 \simeq H^q(X_{\bar\eta},\mathbb Q_\ell(-1)),
 \tag{5.6.2.2}\\
E_2^{p,q}
 &=H^p(\mathbb P^1,R^q\rho_*\mathbb Q_\ell)=0,
 \qquad p\neq0,2.
 \tag{5.6.2.3}
\end{align}
\underline{For $q=n-1$ one has}
\begin{align}
E_2^{0,n-1}
 &=H^0(\mathbb P^1,R^{n-1}\rho_*\mathbb Q_\ell)
 \simeq H^{n-1}(X_{\bar\eta},\mathbb Q_\ell)^\pi
 \simeq H^{n-1}(X,\mathbb Q_\ell),
 \tag{5.6.2.4}\\
E_2^{2,n-1}
 &=H^2(\mathbb P^1,R^{n-1}\rho_*\mathbb Q_\ell)
 \simeq H^{n-1}(X_{\bar\eta},\mathbb Q_\ell(-1))^\pi
 \simeq H^{n-1}(X,\mathbb Q_\ell(-1)).
 \tag{5.6.2.5}
\end{align}
(\underline{The missing value $E_2^{1,n-1}$ will be found in 5.7 below}).

}{E7II-XVIII-022}\underline{Proof.}
By (5.4.3), Theorem 1.1 applies and gives (5.6.1).  The first three
formulas (5.6.2.1)--(5.6.2.3) follow from the constancy of
$R^q\rho_*\mathbb Q_\ell$ for $q\neq n-1$ (Lemma 5.4); for (5.6.2.3) one
also uses, of course, the simple connectedness of $\mathbb P^1$.

\SGAvar{The first isomorphism in (5.6.2.4) follows from
\[
R^{n-1}\rho_*\mathbb Q_\ell
\overset{\sim}{\longrightarrow}
\nu_*\nu^*R^{n-1}\rho_*\mathbb Q_\ell.
\]
Thus
\[
H^0(\mathbb P^1,R^{n-1}\rho_*\mathbb Q_\ell)
\simeq
H^0(U,\nu^*R^{n-1}f_*\mathbb Q_\ell),
\]
and $\nu^*R^{n-1}f_*\mathbb Q_\ell$ is a
\textquotedblleft twisted constant\textquotedblright{} sheaf on $U$, hence
corresponds to the $\pi$-module
$H^{n-1}(X_{\bar\eta},\mathbb Q_\ell)$.  The second isomorphism in
(5.6.2.4) follows from degeneration together with Corollary 5.1.6.

We turn to (5.6.2.5).  Let $F^\bullet$ denote the filtration on
$H^\bullet(\widetilde X,\mathbb Q_\ell)$ defined by the spectral sequence.
By degeneration,
$F^2H^{n+1}(\widetilde X,\mathbb Q_\ell)=E_2^{2,n-1}$ and
$\operatorname{gr}^1H^{n+1}(\widetilde X,\mathbb Q_\ell)=E_2^{1,n}=0$
by (5.6.2.3).  Consequently (cf. 5.1.0),
\begin{align}
E_2^{2,n-1}
 &=F^1H^{n+1}(\widetilde X,\mathbb Q_\ell) \notag\\
 &=\operatorname{Ker}\!\left(
 H^{n+1}(\widetilde X,\mathbb Q_\ell)
 \longrightarrow H^0(\mathbb P^1,R^{n+1}\rho_*\mathbb Q_\ell)
 \right. \notag\\
 &\overset{(A)}{=}
 \operatorname{Ker}\!\left(
 k^*:H^{n+1}(\widetilde X,\mathbb Q_\ell)
 \longrightarrow H^{n+1}(X_{\bar\eta},\mathbb Q_\ell)
 \right).
\tag{5.6.3}
\end{align}
By (5.1.2), $k^*$ may be written as
\begin{equation}\tag{5.6.4}
H^{n+1}(X,\mathbb Q_\ell)
\oplus H^{n-1}(\Delta,\mathbb Q_\ell(-1))
\xrightarrow{\ m^*+h_*\ }
H^{n+1}(X_{\bar\eta},\mathbb Q_\ell).
\end{equation}
Using (5.2.2.3), this becomes
\begin{equation}\tag{5.6.5}
\begin{split}
&\operatorname{Prim}^{n-1}(X,\mathbb Q_\ell(-1)
\oplus H^{n-3}(X,\mathbb Q_\ell(-2))
\oplus H^{n-1}(\Delta,\mathbb Q_\ell(-1))\\
&\hspace{35mm}\xrightarrow{\ m^*L+m^*L^2h_*\ }
H^{n+1}(X_{\bar\eta},\mathbb Q_\ell.
\end{split}
\end{equation}

Now the map
\[
H^{n-3}(X,\mathbb Q_\ell(-2)
\xrightarrow{\ m^*L^2\ }
H^{n+1}(X_{\bar\eta},\mathbb Q_\ell)
\]
is an isomorphism by (LV), since
$m^*L_X^2=L_{X_{\bar\eta}}^2m^*$ (cf. 5.2.4 and 5.2.5).  Hence the kernel
of (5.6.8) identifies, by $\operatorname{pr}_1\oplus\operatorname{pr}_3$,
with
\[
\operatorname{Prim}^{n-1}(X,\mathbb Q_\ell(-1))
\oplus H^{n-1}(\Delta,\mathbb Q_\ell(-1)).
\]
By (LV) and weak Lefschetz,
\[
H^{n-1}(\Delta,\mathbb Q_\ell(-1))
\xleftarrow[\sim]{\,L\,}
H^{n-3}(\Delta,\mathbb Q_\ell(-2))
\xleftarrow[\sim]{\ }
H^{n-3}(X,\mathbb Q_\ell(-2)).
\]
Together with (5.2.2.3), this gives
\begin{equation}\tag{5.6.6}
E_2^{2,n-1}\simeq H^{n-1}(X,\mathbb Q_\ell(-1)).
\end{equation}
The other isomorphism in (5.6.2.5) follows by twisting (5.6.2.4).

}{The first isomorphism in (5.6.2.4) follows from
\[
R^{n-1}\rho_*\mathbb Q_\ell
\overset{\sim}{\longrightarrow}
\nu_*\nu^*R^{n-1}\rho_*\mathbb Q_\ell.
\]
Thus
\[
H^0(\mathbb P^1,R^{n-1}\rho_*\mathbb Q_\ell)
\simeq
H^0(U,\nu^*R^{n-1}\rho_*\mathbb Q_\ell),
\]
and $\nu^*R^{n-1}\rho_*\mathbb Q_\ell$ is a
\textquotedblleft twisted constant\textquotedblright{} sheaf on $U$, hence
corresponds to the $\pi$-module
$H^{n-1}(X_{\bar\eta},\mathbb Q_\ell)$.  The second isomorphism in
(5.6.2.4) follows from degeneration together with Corollary 5.1.6.

We turn to (5.6.2.5).  Let $F^\bullet$ denote the filtration on
$H^\bullet(\widetilde X,\mathbb Q_\ell)$ defined by the spectral sequence.
By degeneration,
$F^2H^{n+1}(\widetilde X,\mathbb Q_\ell)=E_2^{2,n-1}$ and
$\operatorname{gr}^1H^{n+1}(\widetilde X,\mathbb Q_\ell)=E_2^{1,n}=0$
by (5.6.2.3).  Consequently (cf. 5.1.0),
\begin{align}
E_2^{2,n-1}
 &=F^1H^{n+1}(\widetilde X,\mathbb Q_\ell) \notag\\
 &=\operatorname{Ker}\!\left(
 H^{n+1}(\widetilde X,\mathbb Q_\ell)
 \longrightarrow H^0(\mathbb P^1,R^{n+1}\rho_*\mathbb Q_\ell)
 \right) \notag\\
 &\overset{(A)}{=}
 \operatorname{Ker}\!\left(
 k^*:H^{n+1}(\widetilde X,\mathbb Q_\ell)
 \longrightarrow H^{n+1}(X_{\bar\eta},\mathbb Q_\ell)
 \right).
\tag{5.6.3}
\end{align}
By (5.1.2), $k^*$ may be written as
\begin{equation}\tag{5.6.4}
H^{n+1}(X,\mathbb Q_\ell)
\oplus H^{n-1}(\Delta,\mathbb Q_\ell(-1))
\xrightarrow{\ m^*+h_*\ }
H^{n+1}(X_{\bar\eta},\mathbb Q_\ell).
\end{equation}
Using (5.2.2.3), this becomes
\begin{equation}\tag{5.6.5}
\begin{split}
&\operatorname{Prim}^{n-1}(X,\mathbb Q_\ell(-1))
\oplus H^{n-3}(X,\mathbb Q_\ell(-2))
\oplus H^{n-1}(\Delta,\mathbb Q_\ell(-1))\\
&\hspace{35mm}\xrightarrow{\ m^*L+m^*L^2+h_*\ }
H^{n+1}(X_{\bar\eta},\mathbb Q_\ell).
\end{split}
\end{equation}

Now the map
\[
H^{n-3}(X,\mathbb Q_\ell(-2))
\xrightarrow{\ m^*L^2\ }
H^{n+1}(X_{\bar\eta},\mathbb Q_\ell)
\]
is an isomorphism by (LV), since
$m^*L_X^2=L_{X_{\bar\eta}}^2m^*$ (cf. 5.2.4 and 5.2.5).  Hence the kernel
of (5.6.5) identifies, by $\operatorname{pr}_1\oplus\operatorname{pr}_3$,
with
\[
\operatorname{Prim}^{n-1}(X,\mathbb Q_\ell(-1))
\oplus H^{n-1}(\Delta,\mathbb Q_\ell(-1)).
\]
By (LV) and weak Lefschetz,
\[
H^{n-1}(\Delta,\mathbb Q_\ell(-1))
\xleftarrow[\sim]{\,L\,}
H^{n-3}(\Delta,\mathbb Q_\ell(-2))
\xleftarrow[\sim]{\ }
H^{n-3}(X,\mathbb Q_\ell(-2)).
\]
Together with (5.2.2.3), this gives
\begin{equation}\tag{5.6.6}
E_2^{2,n-1}\simeq H^{n-1}(X,\mathbb Q_\ell(-1)).
\end{equation}
The other isomorphism in (5.6.2.5) follows by twisting (5.6.2.4).

}{E7II-XVIII-023}\SGApage{II}{293}
\subsubsection*{Remark 5.6.7.}
The assertion (5.6.1) can be proved under substantially weaker hypotheses,
as Deligne explained to me:


\subsection*{Theorem 5.6.8.}
\underline{Let $X$ be smooth of dimension $n$ over an algebraically closed
field $k$, and let}
\[
\rho:X\longrightarrow\mathbb P^1_k
\]
\underline{be a projective morphism whose generic fiber is smooth.  Suppose
that:}

\smallskip
\noindent
\textup{(5.6.8.1)}\quad
\underline{the geometric generic fiber, equipped with the restriction of the
relatively ample invertible sheaf $\underline{\mathcal O}_X(1)$, satisfies
(LV);}

\smallskip
\noindent
\textup{(5.6.8.2)}\quad
\underline{the maps (5.3.2) are isomorphisms for $i\neq n$ and are
surjective for $i=n$ (a condition satisfied in particular when $\rho$ arises
from a Lefschetz pencil, cf. 5.3.4).  Then the Leray spectral sequence}
\[
E_2^{p,q}=H^p(\mathbb P^1,R^q\rho_*\mathbb Q_\ell)
\Longrightarrow H^{p+q}(X,\mathbb Q_\ell)
\]
\underline{degenerates.}

\smallskip
\noindent
\underline{Remark.}
The same theorem, with the same proof, holds with $\mathbb P^1_k$ replaced
by any smooth, irreducible, projective curve.

\SGApage{II}{\(\approx\)294}
\SGAvar{\underline{Proof.}
For the proof we identify the sheaf $\mathbb Q_\ell(1)$ with
$\mathbb Q_\ell$; this is possible because $k$ is algebraically closed.
We may plainly suppose that $X$ is connected and that $\rho$ is
\underline{surjective}, hence flat with fibers of dimension $n-1$.  Let
\[
\xi_X\in H^2(X,\mathbb Q_\ell)
\]
be the first Chern class of $\underline{\mathcal O}_{X(1)}$, and let
\[
\nu:U\hookrightarrow\mathbb P^1
\]
be the inclusion of an open subset over which $\rho$ is smooth.  By
(5.6.8.1) and Lemma 5.2.6, the cup-product maps
\[
\xi_X^i:
\nu^*R^{n-1-i}\rho_*\mathbb Q_\ell
\longrightarrow
\nu^*R^{n-1+i}\rho_*\mathbb Q_\ell
\]
are isomorphisms for $i\geq1$.  Applying $\nu_*$ and using (5.6.8.2) gives,
for $i\geq2$, the isomorphisms
\begin{equation}\tag{5.6.8.3}
\xi_X^i:R^{n-1-i}\rho_*\mathbb Q_\ell
\overset{\sim}{\longrightarrow}
R^{n-1+i}\rho_*\mathbb Q_\ell,
\end{equation}
and, for $i=1$, the isomorphism
\begin{equation}\tag{5.6.8.4}
\xi_X:R^{n-2}\rho_*\mathbb Q_\ell
\overset{\sim}{\longrightarrow}
\nu_*\nu^*R^n\rho_*\mathbb Q_\ell.
\end{equation}
The latter factors as
\begin{equation}\tag{5.6.8.5}
R^{n-2}\rho_*\mathbb Q_\ell
\xrightarrow{\ \xi_X\ }
R^n\rho_*\mathbb Q_\ell
\xrightarrow{\ \text{adjunction}\ }
\nu_*\nu^*R^n\rho_*\mathbb Q_\ell.
\end{equation}
Consequently there is a canonical decomposition
\begin{equation}\tag{5.6.8.6}
R^n\rho_*\mathbb Q_\ell
\xleftarrow[\sim]{\ \xi_X+\mathrm{incl}\ }
R^{n-2}\rho_*\mathbb Q_\ell\oplus
\operatorname{Local}^n\mathbb Q_\ell,
\end{equation}
where
\begin{equation}\tag{5.6.8.7}
\operatorname{Local}^n\mathbb Q_\ell
=\operatorname{Ker}\!\left(
R^n\rho_*\mathbb Q_\ell
\longrightarrow\nu_*\nu^*R^n\rho_*\mathbb Q_\ell
\right).
\end{equation}
The sheaf $\operatorname{Local}^n\mathbb Q_\ell$ is concentrated at the
points of $\mathbb P^1$ where the fiber of $\rho$ is not smooth.  Combining
this fact with (5.6.8.3), we obtain
\begin{equation}\tag{5.6.8.8}
\xi_X:
H^j(\mathbb P^1,R^{n-2-i}\rho_*\mathbb Q_\ell)
\overset{\sim}{\longrightarrow}
H^j(\mathbb P^1,R^{n-1+i}\mathbb Q_\ell),
\qquad i\geq1,\quad j\geq1.
\end{equation}


We now prove degeneration.  Since $E_2^{p,q}=0$ unless $p=0,1,2$, it is
enough to prove that all the differentials
\[
d_2:H^0(\mathbb P^1,R^i\rho_*\mathbb Q_\ell)
\longrightarrow H^2(\mathbb P^1,R^{i-1},\rho_*\mathbb Q_\ell)
\]
vanish.  For $q\leq n-1$, set
\begin{equation}\tag{5.6.8.9}
\operatorname{Prim}(E_2^{p,q})
=\operatorname{Ker}\!\left(
\xi_X^{n-q}:E_2^{p,q}\longrightarrow E_2^{p,2n-q}
\right).
\end{equation}
Equations (5.6.8.3) and (5.6.8.6), in particular the injectivity of
\[
\xi_X:H^p(\mathbb P^1,R^{n-2}\rho_*\mathbb Q_\ell)
\longrightarrow H^p(\mathbb P^1,\mathbb R^1\rho_*\mathbb Q_\ell),
\]
give, for $q\leq n-1$, the decomposition
\begin{equation}\tag{5.6.8.10}
E_2^{p,q}
\xleftarrow[\sim]{\ 1+\xi_X+\xi_X^2+\cdots\ }
\operatorname{Prim}(E_2^{p,q})
\oplus\operatorname{Prim}(E_2^{p,q-2})
\oplus\operatorname{Prim}(E_2^{p,q-4})
.
\end{equation}



To prove that
\[
d_2^{p,q}=0\qquad(q\leq n-1),
\]
consider the commutative diagram
\[
\begin{tikzcd}[column sep=huge,row sep=huge]
\operatorname{Prim}(E_2^{p,q})
  \arrow[r,"\xi_X^{n-q}=0"]
  \arrow[d,"d_2^{p,q}"]
\&
E_2^{p,2n-q}
  \arrow[d,"d_2^{p,2n-q}"]
\\
E_2^{p+2,q-1}
  \arrow[r,"\xi_X^{n-1-(q-1)}","\sim"']
\&
E_2^{p+2,2n-q-1}.
\end{tikzcd}
\]
The lower arrow is an isomorphism by (5.6.8.8).  The diagram therefore shows
that $d_2^{p,q}$ kills $\operatorname{Prim}(E_2^{p,q})$ for $q\leq n-1$,
and (5.6.8.10) then shows that $d_2$ kills $E_2^{p,q}$ for $q\leq n-1$.
By (5.6.8.3), it also kills $E_2^{p,q}$ for $q\geq n+1$.  It remains only
to prove that
\[
d_2^{0,n}:E_2^{0,n}\longrightarrow E_2^{2,n-1}
\]
is zero, or equivalently that
\begin{equation}\tag{5.6.8.12}
\dim E_3^{2,n-1}=\dim E_2^{2,n-1}.
\end{equation}
We have already proved
\begin{equation}\tag{5.6.8.13}
\dim E_2^{p,q}=\dim E_3^{p,q}
\qquad\text{for }(p,q)\neq(0,n),(2,n-1),
\end{equation}
and for degree reasons
\begin{equation}\tag{5.6.8.14}
E_3^{p,q}=E_\infty^{p,q}
\qquad\text{for every pair }(p,q).
\end{equation}
It follows that
\[
b_i(X)=\dim H^i(X,\mathbb Q_\ell)
=\sum_{p+q=i}\dim E_3^{p,q}.
\]
In view of (5.6.8.13), we therefore have
\begin{equation}\tag{5.6.8.15}
b_{n-1}(X)
=\dim E_2^{c,n-1}
 +\dim E_2^{1,n-2}
 +\dim _2^{2,n-3}
\end{equation}
and
\begin{equation}\tag{5.6.8.16}
b_{n+1}(X)
=\dim E_2^{0,n+1}
 +\dim E_2^{1,n}
 +\dim E_3^{2,n-1}.
\end{equation}
Poincaré duality gives
\begin{equation}\tag{5.6.8.17}
b_{n-1}(X)=b_{n+1}(X).
\end{equation}
First, by (5.6.8.8),
\[
\dim E_2^{1,n}=\dim E_2^{1,n-2}.
\]
Thus (5.6.8.17) shows that the proof of (5.6.8.12) is complete once we know
\begin{align}
\dim E_2^{2,n-3}&=\dim E_2^{0,n+1},
\tag{5.6.8.18}\\
\dim E_2^{0,n-1}&=\dim E_2^{2,n-1}.
\tag{5.6.8.19}
\end{align}
These two formulas follow from (5.6.8.2), Poincaré duality on the smooth
fibers of $\rho$, and the following lemma applied to $\mathbb P^1$ and the
sheaves $R^{n-3}\rho_*\mathbb Q_\ell$ and
$R^{n-1}\rho_*\mathbb Q_\ell$.

}{\underline{Proof.}
For the proof we identify the sheaf $\mathbb Q_\ell(1)$ with
$\mathbb Q_\ell$; this is possible because $k$ is algebraically closed.
We may plainly suppose that $X$ is connected and that $\rho$ is
\underline{surjective}, hence flat with fibers of dimension $n-1$.  Let
\[
\xi_X\in H^2(X,\mathbb Q_\ell)
\]
be the first Chern class of $\underline{\mathcal O}_X(1)$, and let
\[
\nu:U\hookrightarrow\mathbb P^1
\]
be the inclusion of an open subset over which $\rho$ is smooth.  By
(5.6.8.1) and Lemma 5.2.6, the cup-product maps
\[
\xi_X^i:
\nu^*R^{n-1-i}\rho_*\mathbb Q_\ell
\longrightarrow
\nu^*R^{n-1+i}\rho_*\mathbb Q_\ell
\]
are isomorphisms for $i\geq1$.  Applying $\nu_*$ and using (5.6.8.2) gives,
for $i\geq2$, the isomorphisms
\begin{equation}\tag{5.6.8.3}
\xi_X^i:R^{n-1-i}\rho_*\mathbb Q_\ell
\overset{\sim}{\longrightarrow}
R^{n-1+i}\rho_*\mathbb Q_\ell,
\end{equation}
and, for $i=1$, the isomorphism
\begin{equation}\tag{5.6.8.4}
\xi_X:R^{n-2}\rho_*\mathbb Q_\ell
\overset{\sim}{\longrightarrow}
\nu_*\nu^*R^n\rho_*\mathbb Q_\ell.
\end{equation}
The latter factors as
\begin{equation}\tag{5.6.8.5}
R^{n-2}\rho_*\mathbb Q_\ell
\xrightarrow{\ \xi_X\ }
R^n\rho_*\mathbb Q_\ell
\xrightarrow{\ \text{adjunction}\ }
\nu_*\nu^*R^n\rho_*\mathbb Q_\ell.
\end{equation}
Consequently there is a canonical decomposition
\begin{equation}\tag{5.6.8.6}
R^n\rho_*\mathbb Q_\ell
\xleftarrow[\sim]{\ \xi_X+\mathrm{incl}\ }
R^{n-2}\rho_*\mathbb Q_\ell\oplus
\operatorname{Local}^n\mathbb Q_\ell,
\end{equation}
where
\begin{equation}\tag{5.6.8.7}
\operatorname{Local}^n\mathbb Q_\ell
=\operatorname{Ker}\!\left(
R^n\rho_*\mathbb Q_\ell
\longrightarrow\nu_*\nu^*R^n\rho_*\mathbb Q_\ell
\right).
\end{equation}
The sheaf $\operatorname{Local}^n\mathbb Q_\ell$ is concentrated at the
points of $\mathbb P^1$ where the fiber of $\rho$ is not smooth.  Combining
this fact with (5.6.8.3), we obtain
\begin{equation}\tag{5.6.8.8}
\xi_X^{i+1}:
H^j(\mathbb P^1,R^{n-2-i}\rho_*\mathbb Q_\ell)
\overset{\sim}{\longrightarrow}
H^j(\mathbb P^1,R^{n+i}\rho_*\mathbb Q_\ell),
\qquad i\geq0,\quad j\geq1.
\end{equation}


We now prove degeneration.  Since $E_2^{p,q}=0$ unless $p=0,1,2$, it is
enough to prove that all the differentials
\[
d_2:H^0(\mathbb P^1,R^i\rho_*\mathbb Q_\ell)
\longrightarrow H^2(\mathbb P^1,R^{i-1}\rho_*\mathbb Q_\ell)
\]
vanish.  For $q\leq n-1$, set
\begin{equation}\tag{5.6.8.9}
\operatorname{Prim}(E_2^{p,q})
=\operatorname{Ker}\!\left(
\xi_X^{n-q}:E_2^{p,q}\longrightarrow E_2^{p,2n-q}
\right).
\end{equation}
Equations (5.6.8.3) and (5.6.8.6), in particular the injectivity of
\[
\xi_X:H^p(\mathbb P^1,R^{n-2}\rho_*\mathbb Q_\ell)
\longrightarrow H^p(\mathbb P^1,R^n\rho_*\mathbb Q_\ell),
\]
give, for $q\leq n-1$, the decomposition
\begin{equation}\tag{5.6.8.10}
E_2^{p,q}
\xleftarrow[\sim]{\ 1+\xi_X+\xi_X^2+\cdots\ }
\operatorname{Prim}(E_2^{p,q})
\oplus\operatorname{Prim}(E_2^{p,q-2})
\oplus\operatorname{Prim}(E_2^{p,q-4})
\oplus\cdots.
\end{equation}



To prove that
\[
d_2^{p,q}=0\qquad(q\leq n-1),
\]
consider the commutative diagram
\[
\begin{tikzcd}[column sep=huge,row sep=huge]
\operatorname{Prim}(E_2^{p,q})
  \arrow[r,"\xi_X^{n-q}=0"]
  \arrow[d,"d_2^{p,q}"]
\&
E_2^{p,2n-q}
  \arrow[d,"d_2^{p,2n-q}"]
\\
E_2^{p+2,q-1}
  \arrow[r,"\xi_X^{n-1-(q-1)}","\sim"']
\&
E_2^{p+2,2n-q-1}.
\end{tikzcd}
\]
The lower arrow is an isomorphism by (5.6.8.8).  The diagram therefore shows
that $d_2^{p,q}$ kills $\operatorname{Prim}(E_2^{p,q})$ for $q\leq n-1$,
and (5.6.8.10) then shows that $d_2$ kills $E_2^{p,q}$ for $q\leq n-1$.
By (5.6.8.3), it also kills $E_2^{p,q}$ for $q\geq n+1$.  It remains only
to prove that
\[
d_2^{0,n}:E_2^{0,n}\longrightarrow E_2^{2,n-1}
\]
is zero, or equivalently that
\begin{equation}\tag{5.6.8.12}
\dim E_3^{2,n-1}=\dim E_2^{2,n-1}.
\end{equation}
We have already proved
\begin{equation}\tag{5.6.8.13}
\dim E_2^{p,q}=\dim E_3^{p,q}
\qquad\text{for }(p,q)\neq(0,n),(2,n-1),
\end{equation}
and for degree reasons
\begin{equation}\tag{5.6.8.14}
E_3^{p,q}=E_\infty^{p,q}
\qquad\text{for every pair }(p,q).
\end{equation}
It follows that
\[
b_i(X)=\dim H^i(X,\mathbb Q_\ell)
=\sum_{p+q=i}\dim E_3^{p,q}.
\]
In view of (5.6.8.13), we therefore have
\begin{equation}\tag{5.6.8.15}
b_{n-1}(X)
=\dim E_2^{0,n-1}
 +\dim E_2^{1,n-2}
 +\dim E_2^{2,n-3}
\end{equation}
and
\begin{equation}\tag{5.6.8.16}
b_{n+1}(X)
=\dim E_2^{0,n+1}
 +\dim E_2^{1,n}
 +\dim E_3^{2,n-1}.
\end{equation}
Poincaré duality gives
\begin{equation}\tag{5.6.8.17}
b_{n-1}(X)=b_{n+1}(X).
\end{equation}
First, by (5.6.8.8),
\[
\dim E_2^{1,n}=\dim E_2^{1,n-2}.
\]
Thus (5.6.8.17) shows that the proof of (5.6.8.12) is complete once we know
\begin{align}
\dim E_2^{2,n-3}&=\dim E_2^{0,n+1},
\tag{5.6.8.18}\\
\dim E_2^{0,n-1}&=\dim E_2^{2,n-1}.
\tag{5.6.8.19}
\end{align}
These two formulas follow from (5.6.8.2), Poincaré duality on the smooth
fibers of $\rho$, and the following lemma applied to $\mathbb P^1$ and the
sheaves $R^{n-3}\rho_*\mathbb Q_\ell$ and
$R^{n-1}\rho_*\mathbb Q_\ell$.

}{E7II-XVIII-024}\SGApage{II}{298}
\subsection*{Lemma 5.6.9.}
\underline{Let $D$ be a smooth, connected, proper curve over an
algebraically closed field $k$.  Let $i:\eta\hookrightarrow D$ be the
inclusion of the generic point, and let $F$ be a constructible
$\mathbb Q_\ell$-sheaf on $D$ such that
$F\overset{\sim}{\longrightarrow}i_*i^*F$.  Then}
\[
H^2(D,F)\simeq H^0(D,\check F(-1))^\vee,
\]
\underline{where $\check F\simeq i_*((i^*F)^\vee)$.}

\subsection*{Corollary 5.6.10.}
\textup{(Uses Picard--Lefschetz.)}
\underline{Let $\rho:\widetilde X\to\mathbb P^1$ be the
\textquotedblleft fibration\textquotedblright{} arising from a Lefschetz
pencil, and suppose that the smooth fibers satisfy (LV).  Then}
\begin{equation}\tag{5.6.10.1}
H^i(X_{\bar\eta},\mathbb Q_\ell)^\pi
=\operatorname{Im}\!\left(
H^i(X,\mathbb Q_\ell)\longrightarrow
H^i(X_{\bar\eta},\mathbb Q_\ell)
\right).
\end{equation}

\underline{Proof.}
By 5.3.4, which uses the Picard--Lefschetz formula, Theorem 5.6.8 applies
and gives
\[
H^i(X_{\bar\eta},\mathbb Q_\ell)^\pi
=\operatorname{Im}\!\left(
H^i(\widetilde X,\mathbb Q_\ell)\longrightarrow
H^i(X_{\bar\eta},\mathbb Q_\ell)
\right).
\]
It remains to check that $H^i(X,\mathbb Q_\ell)$ and
$H^i(\widetilde X,\mathbb Q_\ell)$ have the same image in
$H^i(X_{\bar\eta},\mathbb Q_\ell)$.  This holds for $i\leq n-1$ by
Corollary 5.1.6, and for $i\geq n$ by Corollary 5.2.5.

We now turn to the most important computation:


\subsection*{Theorem 5.7.}
\underline{Under the hypotheses of 5.6, set}
\[
\begin{aligned}
E^{n-2}(\Delta,\mathbb Q_\ell(j))
&=\underline{\text{the orthogonal complement in }}
H^{n-2}(\Delta,\mathbb Q_\ell(j))\\
&\qquad\underline{\text{of the image of }}
H^{n-2}(X,\mathbb Q_\ell(j)).
\end{aligned}
\]
(\underline{This is the} \textquotedblleft\underline{vanishing}\textquotedblright{}
\underline{part of the cohomology of $\Delta$; cf. 5.4.4}.)  \underline{Then there is a
direct-sum decomposition} (cf. 4.2)
\begin{equation}\tag{5.7.1}
E_2^{1,n-1}
=H^1(\mathbb P^1,R^{n-1}\rho_*\mathbb Q_\ell)
\xleftarrow[\;f^*+i_*g^*\;]{\sim}
\operatorname{Prim}^n(X,\mathbb Q_\ell)
\oplus E^{n-2}(\Delta,\mathbb Q_\ell(-1)).
\end{equation}

\SGApage{II}{\(\approx\)299}
\SGAvar{\underline{Proof.}
By degeneration,
\[
E_2^{1,n-1}\simeq
F^1H^n(\widetilde X,\mathbb Q_\ell)
\big/F^2H^n(\widetilde X,\mathbb Q_\ell).
\]
By (A),
\begin{equation}\tag{5.7.2}
F^1H^n(\widetilde X,\mathbb Q_\ell)
=\ker\!\left(
k^*:H^n(\widetilde X,\mathbb Q_\ell)\longrightarrow
H^n(X_{\bar\eta},\mathbb Q_\ell)
\right),
\end{equation}
which, by (5.1.2), identifies with
\[
\ker:\!\left(
H^n(X,\mathbb Q_\ell)
\oplus H^{n-2}(\Delta,\mathbb Q_\ell(-1))
\xrightarrow{\;m^*+h_*\;}
H^n(X_{\bar\eta},\mathbb Q_\ell)
\right).
\]
Under (LV), $k^*$ can also be written as
\begin{equation}\tag{5.7.3}
\operatorname{Prim}^n(X,\mathbb Q_\ell)
\oplus H^{n-2}(X,\mathbb Q_\ell(-1))
\oplus H^{n-2}(\Delta,\mathbb Q_\ell(-1))
\xrightarrow{\;m^*+m^*L+h_*\;}
H^n(X_{\bar\eta},\mathbb Q_\ell),
\end{equation}
and the middle arrow
\begin{equation}\tag{5.7.4}
H^{n-2}(X,\mathbb Q_\ell(-1))
\xrightarrow{\;m^*L\;}
H^n(X_{\bar\eta},\mathbb Q_\ell)
\end{equation}
is an isomorphism (cf. 5.2.3 and (5.2.2.3)).  Thus, denoting by
\begin{equation}\tag{5.7.5}
\wp:H^n(X,\mathbb Q_\ell)\longrightarrow
\operatorname{Prim}^n(X,\mathbb Q_\ell)
\end{equation}
the projection onto the primitive part, we obtain

\begin{equation}\tag{5.7.6}
F^1H^n(\widetilde X,\mathbb Q_\ell)
\xrightarrow[\;\wp f_*\oplus(-g_*2^*)\;]{\sim}
\operatorname{Prim}^n(X,\mathbb Q_\ell)
\oplus H^{n-2}(\Delta,\mathbb Q_\ell(-1)).
\end{equation}
By (5.6.2.2) and degeneration,
\begin{equation}\tag{5.7.7}
F^2H^n(\widetilde X,\mathbb Q_\ell)
\simeq E_2^{2,n-2}
\simeq H^{n-2}(X_{\bar\eta},\mathbb Q_\ell(-1)).
\end{equation}

On the other hand, the inclusion of
$F^2H^N(\widetilde X,\mathbb Q_\ell)
=H^{n-2}(X_{\bar\eta},\mathbb Q_\ell(-1))$
in $H^n(\widetilde X,\mathbb Q_\ell)$ is the Gysin map $k_*$.

Strictly speaking, the Gysin map is defined only as a map
\[
H^{i-2}(X_{\bar\eta},\mathbb Q_\ell))
\longrightarrow
H^i(X\times_k k(\bar\eta),\mathbb Q_\ell),
\]
but, in view of the canonical isomorphism (SGA 4 XII 5.4)
\[
H^i(X,\mathbb Q_\ell)\xrightarrow{\sim}
H^i(X\otimes_kK,\mathbb Q_\ell)
\]
for every algebraically closed field $K/k$, we hope the reader will allow
us such abuses.
By (5.1.3), the Gysin map is expressed as
\begin{equation}\tag{5.7.8}
H^{n-2}(X_{\bar\eta},\mathbb Q_\ell(-1))
\xrightarrow{\;m_*\oplus(-h^*)\;}
H^n(X,\mathbb Q_\ell)
\oplus H^{n-2}(\Delta,\mathbb Q_\ell(-1)).
\end{equation}
By weak Lefschetz,
\[
m_*H^{n-2}(X_{\bar\eta},\mathbb Q_\ell(-1))
=m_*m^*H^{n-2}(X,\mathbb Q_\ell(-1))
=L_XH^{n-2}(X,\mathbb Q_\ell(-1))
\]
is the orthogonal complement of the primitive part.  Hence we have the
commutative diagram
\begin{equation}\tag{5.7.9}
\begin{tikzcd}[column sep=huge,row sep=huge]
F^2H^n(\widetilde X,\mathbb Q_\ell)
  \arrow[r,hook]
\& F^1H^n(\widetilde X,\mathbb Q_\ell)
  \arrow[d,"\wp f_*\oplus(-g_*i^*)","\wr"']
\\
H^{n-2}(X_{\bar\eta},\mathbb Q_\ell(-1))
  \arrow[u,"k_*"',"\wr"]
  \arrow[r,"0\oplus h^*"]
\& \operatorname{Prim}(X,\mathbb Q_\ell)
  \oplus H^{n-2}(\Delta,\mathbb Q_\ell(-1)).
\end{tikzcd}
\end{equation}
This proves (5.7.1), as required.

}{\underline{Proof.}
By degeneration,
\[
E_2^{1,n-1}\simeq
F^1H^n(\widetilde X,\mathbb Q_\ell)
\big/F^2H^n(\widetilde X,\mathbb Q_\ell).
\]
By (A),
\begin{equation}\tag{5.7.2}
F^1H^n(\widetilde X,\mathbb Q_\ell)
=\ker\!\left(
k^*:H^n(\widetilde X,\mathbb Q_\ell)\longrightarrow
H^n(X_{\bar\eta},\mathbb Q_\ell)
\right),
\end{equation}
which, by (5.1.2), identifies with
\[
\ker\!\left(
H^n(X,\mathbb Q_\ell)
\oplus H^{n-2}(\Delta,\mathbb Q_\ell(-1))
\xrightarrow{\;m^*+h_*\;}
H^n(X_{\bar\eta},\mathbb Q_\ell)
\right).
\]
Under (LV), $k^*$ can also be written as
\begin{equation}\tag{5.7.3}
\operatorname{Prim}^n(X,\mathbb Q_\ell)
\oplus H^{n-2}(X,\mathbb Q_\ell(-1))
\oplus H^{n-2}(\Delta,\mathbb Q_\ell(-1))
\xrightarrow{\;m^*+m^*L+h_*\;}
H^n(X_{\bar\eta},\mathbb Q_\ell),
\end{equation}
and the middle arrow
\begin{equation}\tag{5.7.4}
H^{n-2}(X,\mathbb Q_\ell(-1))
\xrightarrow{\;m^*L\;}
H^n(X_{\bar\eta},\mathbb Q_\ell)
\end{equation}
is an isomorphism (cf. 5.2.3 and (5.2.2.3)).  Thus, denoting by
\begin{equation}\tag{5.7.5}
\wp:H^n(X,\mathbb Q_\ell)\longrightarrow
\operatorname{Prim}^n(X,\mathbb Q_\ell)
\end{equation}
the projection onto the primitive part, we obtain

\begin{equation}\tag{5.7.6}
F^1H^n(\widetilde X,\mathbb Q_\ell)
\xrightarrow[\;\wp f_*\oplus(-g_*i^*)\;]{\sim}
\operatorname{Prim}^n(X,\mathbb Q_\ell)
\oplus H^{n-2}(\Delta,\mathbb Q_\ell(-1)).
\end{equation}
By (5.6.2.2) and degeneration,
\begin{equation}\tag{5.7.7}
F^2H^n(\widetilde X,\mathbb Q_\ell)
\simeq E_2^{2,n-2}
\simeq H^{n-2}(X_{\bar\eta},\mathbb Q_\ell(-1)).
\end{equation}

On the other hand, the inclusion of
$F^2H^n(\widetilde X,\mathbb Q_\ell)
=H^{n-2}(X_{\bar\eta},\mathbb Q_\ell(-1))$
in $H^n(\widetilde X,\mathbb Q_\ell)$ is the Gysin map $k_*$.

Strictly speaking, the Gysin map is defined only as a map
\[
H^{i-2}(X_{\bar\eta},\mathbb Q_\ell(-1))
\longrightarrow
H^i(X\times_k k(\bar\eta),\mathbb Q_\ell),
\]
but, in view of the canonical isomorphism (SGA 4 XII 5.4)
\[
H^i(X,\mathbb Q_\ell)\xrightarrow{\sim}
H^i(X\otimes_kK,\mathbb Q_\ell)
\]
for every algebraically closed field $K/k$, we hope the reader will allow
us such abuses.
By (5.1.3), the Gysin map is expressed as
\begin{equation}\tag{5.7.8}
H^{n-2}(X_{\bar\eta},\mathbb Q_\ell(-1))
\xrightarrow{\;m_*\oplus(-h^*)\;}
H^n(X,\mathbb Q_\ell)
\oplus H^{n-2}(\Delta,\mathbb Q_\ell(-1)).
\end{equation}
By weak Lefschetz,
\[
m_*H^{n-2}(X_{\bar\eta},\mathbb Q_\ell(-1))
=m_*m^*H^{n-2}(X,\mathbb Q_\ell(-1))
=L_XH^{n-2}(X,\mathbb Q_\ell(-1))
\]
is the orthogonal complement of the primitive part.  Hence we have the
commutative diagram
\begin{equation}\tag{5.7.9}
\begin{tikzcd}[column sep=huge,row sep=huge]
F^2H^n(\widetilde X,\mathbb Q_\ell)
  \arrow[r,hook]
\& F^1H^n(\widetilde X,\mathbb Q_\ell)
  \arrow[d,"\wp f_*\oplus(-g_*i^*)","\wr"']
\\
H^{n-2}(X_{\bar\eta},\mathbb Q_\ell(-1))
  \arrow[u,"k_*"',"\wr"]
  \arrow[r,"0\oplus h^*"]
\& \operatorname{Prim}^n(X,\mathbb Q_\ell)
  \oplus H^{n-2}(\Delta,\mathbb Q_\ell(-1)).
\end{tikzcd}
\end{equation}
This proves (5.7.1), as required.

}{E7II-XVIII-025}\SGApage{II}{301}
\subsection*{Remark 5.7.10.}
In 5.7, hypothesis (A) may be replaced by the weaker assumption that, for
$q=n$, the maps
\begin{equation}\tag{5.7.11}
H^q(\widetilde X,\mathbb Q_\ell)\longrightarrow
H^0(\mathbb P^1,R^q\rho_*\mathbb Q_\ell)
\end{equation}
and
\begin{equation}\tag{5.7.12}
H^q(\widetilde X,\mathbb Q_\ell)\longrightarrow
H^q(X_{\bar\eta},\mathbb Q_\ell)
\end{equation}
have the same kernel, as our proof shows.  Without this condition, the
conclusion can fail; cf. 5.9.


\section*{5.8. The Griffiths homomorphism}

\SGAvar{\subsection*{5.8.1.}
Let $f:T\to S$ be a morphism of schemes.  The Leray spectral sequence
defines a filtration $(F^i)$ on the groups
$H^\bullet(T,\mathbb Q_\ell(j))$.  Define
\begin{equation}\tag{5.8.1.1}
\operatorname{Prim}^\bullet(T/S,\mathbb Q_\ell(j))
=F^1H^\bullet(T,\mathbb Q_\ell(j)).
\end{equation}
Because the Leray spectral sequence lies in the first quadrant, we have
\begin{equation}\tag{5.8.1.2}
E_\infty^{0,\bullet}\lhook\joinrel\longrightarrow E_2^{0,\bullet},
\end{equation}
and hence
\begin{equation}\tag{5.8.1.3}
\operatorname{Prim}^\bullet(T_{/S},\mathbb Q_\ell(j)
=\ker\!\left(
H^\bullet(T,\mathbb Q_\ell(j))
\longrightarrow H^0(S,R^\bullet f_*\mathbb Q_\ell(j))
\right..
\end{equation}
For the same reason,
\begin{equation}\tag{5.8.1.4}
E_\infty^{1,\bullet}\lhook\joinrel\longrightarrow E_2^{1,\bullet}.
\end{equation}
The composite
\begin{equation}\tag{5.8.1.5}
\begin{tikzcd}[column sep=large,row sep=large]
\operatorname{Prim}^i(T/S,\mathbb Q_\ell(j))
=F^1H^i(T,\mathbb Q_\ell(j))
  \arrow[r]
\&
\operatorname{gr}^1H^i(T,\mathbb Q_\ell(j))
=E_\infty^{1,i-1}
  \arrow[d,hook]
\\
{}
\&
H^1(S,R^{i-1}_{f_*}\mathbb Q_\ell(j))
=E_2^{1,i-1}
\end{tikzcd}
\end{equation}
will be denoted

\begin{equation}\tag{5.8.1.6}
U^i:\operatorname{Prim}(T_{/S},\mathbb Q_\ell(j))
\longrightarrow H^1(S,R^{\cdot-1}f_*\mathbb Q_\ell(j)).
\end{equation}

\subsection*{5.8.2.}
Now let $X$ be projective, smooth, and irreducible, and let
$D=\{H_\tau\}$ be a pencil of hyperplanes whose generic member and axis
meet $X$ transversely.  Let
$U\xrightarrow{\nu}\mathbb P^1$ be a nonempty open subset over which the
projection $\rho:\widetilde X\to\mathbb P^1$ is smooth, and put
$\widetilde X|_U=\widetilde X\times_{\mathbb P^1}U$.  We have the
commutative diagram
\begin{equation}\tag{5.8.2.1}
\begin{tikzcd}[column sep=large,row sep=large]
{}
\&
\widetilde X
  \arrow[dl,"f"]
  \arrow[d,"\rho"]
\&
\widetilde X|_U
  \arrow[l,"\alpha"']
  \arrow[d,"\rho|_U"]
\\
X
\&
\mathbb P^1
\&
U\arrow[l,"\nu"].
\end{tikzcd}
\end{equation}
Since $\rho|_U:\widetilde X|_U\to U$ is proper and smooth, the
specialization theorem for cohomology (SGA 4 XVI 2.2) gives
\begin{equation}\tag{5.8.2.2}
\operatorname{Prim}^\bullet(\widetilde X|_U/U,\mathbb Q_\ell(j)
=\ker\!\left(
H^\bullet(\widetilde X|_U,\mathbb Q_\ell(j))
\longrightarrow H^\bullet(X_{\bar\eta},\mathbb Q_\ell(j))
\right).
\end{equation}
Consequently, by (5.2.3.1), we have:

}{\subsection*{5.8.1.}
Let $f:T\to S$ be a morphism of schemes.  The Leray spectral sequence
defines a filtration $(F^i)$ on the groups
$H^\bullet(T,\mathbb Q_\ell(j))$.  Define
\begin{equation}\tag{5.8.1.1}
\operatorname{Prim}^\bullet(T/S,\mathbb Q_\ell(j))
=F^1H^\bullet(T,\mathbb Q_\ell(j)).
\end{equation}
Because the Leray spectral sequence lies in the first quadrant, we have
\begin{equation}\tag{5.8.1.2}
E_\infty^{0,\bullet}\lhook\joinrel\longrightarrow E_2^{0,\bullet},
\end{equation}
and hence
\begin{equation}\tag{5.8.1.3}
\operatorname{Prim}^\bullet(T/S,\mathbb Q_\ell(j))
=\ker\!\left(
H^\bullet(T,\mathbb Q_\ell(j))
\longrightarrow H^0(S,R^\bullet f_*\mathbb Q_\ell(j))
\right).
\end{equation}
For the same reason,
\begin{equation}\tag{5.8.1.4}
E_\infty^{1,\bullet}\lhook\joinrel\longrightarrow E_2^{1,\bullet}.
\end{equation}
The composite
\begin{equation}\tag{5.8.1.5}
\begin{tikzcd}[column sep=large,row sep=large]
\operatorname{Prim}^i(T/S,\mathbb Q_\ell(j))
=F^1H^i(T,\mathbb Q_\ell(j))
  \arrow[r]
\&
\operatorname{gr}^1H^i(T,\mathbb Q_\ell(j))
=E_\infty^{1,i-1}
  \arrow[d,hook]
\\
{}
\&
H^1(S,R^{i-1}f_*\mathbb Q_\ell(j))
=E_2^{1,i-1}
\end{tikzcd}
\end{equation}
will be denoted

\begin{equation}\tag{5.8.1.6}
U^i:\operatorname{Prim}^i(T/S,\mathbb Q_\ell(j))
\longrightarrow H^1(S,R^{i-1}f_*\mathbb Q_\ell(j)).
\end{equation}

\subsection*{5.8.2.}
Now let $X$ be projective, smooth, and irreducible, and let
$D=\{H_\tau\}$ be a pencil of hyperplanes whose generic member and axis
meet $X$ transversely.  Let
$U\xrightarrow{\nu}\mathbb P^1$ be a nonempty open subset over which the
projection $\rho:\widetilde X\to\mathbb P^1$ is smooth, and put
$\widetilde X|_U=\widetilde X\times_{\mathbb P^1}U$.  We have the
commutative diagram
\begin{equation}\tag{5.8.2.1}
\begin{tikzcd}[column sep=large,row sep=large]
{}
\&
\widetilde X
  \arrow[dl,"f"]
  \arrow[d,"\rho"]
\&
\widetilde X|_U
  \arrow[l,"\alpha"']
  \arrow[d,"\rho|_U"]
\\
X
\&
\mathbb P^1
\&
U\arrow[l,"\nu"].
\end{tikzcd}
\end{equation}
Since $\rho|_U:\widetilde X|_U\to U$ is proper and smooth, the
specialization theorem for cohomology (SGA 4 XVI 2.2) gives
\begin{equation}\tag{5.8.2.2}
\operatorname{Prim}^\bullet(\widetilde X|_U/U,\mathbb Q_\ell(j))
=\ker\!\left(
H^\bullet(\widetilde X|_U,\mathbb Q_\ell(j))
\longrightarrow H^\bullet(X_{\bar\eta},\mathbb Q_\ell(j))
\right).
\end{equation}
Consequently, by (5.2.3.1), we have:

}{E7II-XVIII-026}\SGApage{II}{303}
\subsection*{Proposition 5.8.2.3.}
\underline{Under the hypotheses of 5.8.2, for
$x\in H^n(X,\mathbb Q_\ell(j))$, the following conditions are equivalent:}
\begin{equation}\tag{5.8.2.4}
x\in\operatorname{Prim}^n(X,\mathbb Q_\ell(j))
\quad\Longleftrightarrow\quad
\alpha^*f^*(x)\in
\operatorname{Prim}^n(\widetilde X|_U/U,\mathbb Q_\ell(j)).
\end{equation}
\underline{If, in addition, $\rho$ satisfies hypothesis (A) of 5.3, then}
\begin{equation}\tag{5.8.2.5}
x\in\operatorname{Prim}^n(X,\mathbb Q_\ell(j))
\quad\Longleftrightarrow\quad
f^*(x)\in
\operatorname{Prim}^n(\widetilde X/\mathbb P^1,\mathbb Q_\ell(j)).
\end{equation}

\SGAvar{\underline{Proof.}
Formula (5.8.2.4) is proved.  For (5.8.2.5), hypothesis (A) implies that
\[
H^o(\mathbb P^1,R^n\rho_*\mathbb Q_\ell(j))
\xrightarrow[\nu^*]{\sim}
H^o(U,\nu^*R^n\rho_*\mathbb Q_\ell(j)).
\]
Thus[$f^*(x)$ is primitive relative to $\mathbb P^1\alpha$] $\Longrightarrow
\alpha^*f^*(\alpha)$ is primitive relative to $U$.


}{\underline{Proof.}
Formula (5.8.2.4) is proved.  For (5.8.2.5), hypothesis (A) implies that
\[
H^0(\mathbb P^1,R^n\rho_*\mathbb Q_\ell(j))
\xrightarrow[\nu^*]{\sim}
H^0(U,\nu^*R^n\rho_*\mathbb Q_\ell(j)).
\]
Thus $f^*(x)$ is primitive relative to $\mathbb P^1$ if and only if
$\alpha^*f^*(x)$ is primitive relative to $U$.


}{E7II-XVIII-027}\subsection*{Definition 5.8.3.}
\underline{In the situation of 5.8.2, the composite homomorphism}

\SGApage{II}{\(\approx\)304}
\SGAvar{\begin{equation}\tag{5.8.3.1}
\begin{tikzcd}[column sep=huge,row sep=huge]
\operatorname{Prim}^n(X,\mathbb Q_\ell(j))
  \arrow[r,"\alpha^*f^*"]
  \arrow[dr,"\operatorname{Grif}"]
\&
\operatorname{Prim}^n(\widetilde X|_U/U,\mathbb Q_\ell(j))
  \arrow[d,"{[\mathrm{U}^n\,?]}"]
\\
{}
\&
H^1(U,\nu^*R^{n-1}\rho_*\mathbb Q_\ell(j))
\end{tikzcd}
\end{equation}
\underline{is called the Griffiths homomorphism (for $X$, relative to the
given pencil $D$).}

}{\begin{equation}\tag{5.8.3.1}
\begin{tikzcd}[column sep=huge,row sep=huge]
\operatorname{Prim}^n(X,\mathbb Q_\ell(j))
  \arrow[r,"\alpha^*f^*"]
  \arrow[dr,"\operatorname{Grif}"]
\&
\operatorname{Prim}^n(\widetilde X|_U/U,\mathbb Q_\ell(j))
  \arrow[d,"U^n"]
\\
{}
\&
H^1(U,\nu^*R^{n-1}\rho_*\mathbb Q_\ell(j))
\end{tikzcd}
\end{equation}
\underline{is called the Griffiths homomorphism (for $X$, relative to the
given pencil $D$).}

}{E7II-XVIII-028}Adjoin to diagram (5.8.2.1) the constant family over $U$ with fiber $X$:
\begin{equation}\tag{5.8.3.2}
\begin{tikzcd}[row sep=large,column sep=large]
{}
\&
\widetilde X
  \arrow[dl,"f"']
  \arrow[dd]
\&
\widetilde X|_U
  \arrow[l,"\alpha"']
  \arrow[r,"f\times\rho|_U"]
  \arrow[dd,"\rho|_U"]
\&
X\times U
  \arrow[ddl,"\operatorname{pr}_2"]
\\
X\&{}\&{}\&{}
\\
{}
\&
\mathbb P^1
\&
U\arrow[l,"\nu"']
\&
.
\end{tikzcd}
\end{equation}



\SGAvar{\subsection*{Theorem 5.8.4.}
\underline{Suppose that $X$ satisfies (LV) (5.2.2), and that $\rho$
satisfies hypothesis (A) (5.3).  Then, for every nonempty smoothness open
$U\xrightarrow{\nu}\mathbb P^1$ of $\rho$, the Griffiths homomorphism
\textquotedblleft modulo the fixed part\textquotedblright{}, which is by
definition the composite}
\begin{equation}\tag{5.8.4.1}
\resizebox{\textwidth}{!}{
\begin{tikzcd}[ampersand replacement=\&,row sep=huge,column sep=huge]
\operatorname{Prim}^n(X,\mathbb Q_\ell(j))
  \arrow[r,"\operatorname{Grif}"]
\&
H^1(U,\nu^*R^{n-1}\rho_*\mathbb Q_\ell(j))
  \arrow[d]
\\
{}
\&
H^1(U,E^{n-1}\rho_*(\mathbb Q_\ell(j))
\simeq
\dfrac{H^1(U,\nu^*R^{n-1}\rho_*\mathbb Q_\ell(j))}
{(f\times\rho|_U)^*
 H^1(U,H^{n-1}(X,\mathbb Q_\ell(j))_U)}
\end{tikzcd}}
\end{equation}
\underline{is injective.}


\underline{Proof.}
By (5.5.1), over $\mathbb P^1$ we have
\begin{equation}\tag{5.8.4.2}
R^{n-1}\rho_*\mathbb Q_\ell(j)
\simeq
H^{n-1}(X,\mathbb Q_\ell(j))_{\mathbb P^1}
\oplus E^{n-1}\rho_*\mathbb Q_\ell(j),
\end{equation}
and over $U$ we have
\begin{equation}\tag{5.8.4.3}
\nu^*R^{n-1}\rho_*\mathbb Q_\ell(j))
\simeq
H^{n-1}(X,\mathbb Q_\ell(j))_U
\oplus \nu^*E^{n-1}\rho_*\mathbb Q_\ell(j).
\end{equation}
This permits us to rewrite (5.8.4.1) as indicated:
\begin{equation}\tag{5.8.4.4}
\begin{tikzcd}[row sep=huge,column sep=large]
\operatorname{Prim}^n(X,\mathbb Q_\ell(j))
  \arrow[r,"\operatorname{Grif}"]
\&
H^1(U,\nu^*R^{n-1}\rho_*\mathbb Q_\ell(j))
  \arrow[d]
\\
{}
\&
H^1(U,\nu^*E^{n-1}\rho_*\mathbb Q_\ell(j)).
\end{tikzcd}
\end{equation}
By (5.7.1), the map (defined only by virtue of hypothesis (A))
\begin{equation}\tag{5.8.4.5}
u:\operatorname{Prim}^n(X,\mathbb Q_\ell(j))
\longrightarrow
H^1(\mathbb P^1,R^{n-1}\rho_*\mathbb Q_\ell)
\end{equation}
is injective.  Since $\mathbb P^1$ is simply connected, decomposition
(5.8.4.2) gives

\begin{equation}\tag{5.8.4.6}
H^1(\mathbb P^1,R^{n-1}\rho_*\mathbb Q_\ell)
\simeq
H^1(\mathbb P^1,E^{n-1}\rho_*\mathbb Q_\ell(j)).
\end{equation}
We can therefore complete (5.8.4.4) to a commutative square:
\begin{equation}\tag{5.8.4.7}
\begin{tikzcd}[column sep=huge,row sep=huge]
\operatorname{Prim}^n(X,\mathbb Q_\ell(j))
  \arrow[r,"\operatorname{Grif}"]
  \arrow[d,hook,"u"]
\&
H^1(U,\nu^*R^{n-1}\rho_*\mathbb Q_\ell(j))
  \arrow[d]
\\
H^1(\mathbb P^1,E^{n-1}\rho_*\mathbb Q_\ell(j))
  \arrow[r,"\nu^*"]
\&
H^1(U,\nu^*E^{n-1}\rho_*\mathbb Q_\ell(j)).
\end{tikzcd}
\end{equation}
It remains only to see that the lower arrow is injective.  This follows from
\begin{equation}\tag{5.8.4.8}
E^{n-1}\rho_*\mathbb Q_\ell(j)
\simeq
\nu_*\nu^*E^{n-1}\rho_*\mathbb Q_\ell(j)
\end{equation}
and the following well-known lemma.

}{\subsection*{Theorem 5.8.4.}
\underline{Suppose that $X$ satisfies (LV) (5.2.2), and that $\rho$
satisfies hypothesis (A) (5.3).  Then, for every nonempty smoothness open
$U\xrightarrow{\nu}\mathbb P^1$ of $\rho$, the Griffiths homomorphism
\textquotedblleft modulo the fixed part\textquotedblright{}, which is by
definition the composite}
\begin{equation}\tag{5.8.4.1}
\resizebox{\textwidth}{!}{
\begin{tikzcd}[ampersand replacement=\&,row sep=huge,column sep=huge]
\operatorname{Prim}^n(X,\mathbb Q_\ell(j))
  \arrow[r,"\operatorname{Grif}"]
\&
H^1(U,\nu^*R^{n-1}\rho_*\mathbb Q_\ell(j))
  \arrow[d]
\\
{}
\&
H^1(U,\nu^*E^{n-1}\rho_*\mathbb Q_\ell(j))
\simeq
\dfrac{H^1(U,\nu^*R^{n-1}\rho_*\mathbb Q_\ell(j))}
{(f\times\rho|_U)^*
 H^1(U,H^{n-1}(X,\mathbb Q_\ell(j))_U)}
\end{tikzcd}}
\end{equation}
\underline{is injective.}


\underline{Proof.}
By (5.5.1), over $\mathbb P^1$ we have
\begin{equation}\tag{5.8.4.2}
R^{n-1}\rho_*\mathbb Q_\ell(j)
\simeq
H^{n-1}(X,\mathbb Q_\ell(j))_{\mathbb P^1}
\oplus E^{n-1}\rho_*\mathbb Q_\ell(j),
\end{equation}
and over $U$ we have
\begin{equation}\tag{5.8.4.3}
\nu^*R^{n-1}\rho_*\mathbb Q_\ell(j)
\simeq
H^{n-1}(X,\mathbb Q_\ell(j))_U
\oplus \nu^*E^{n-1}\rho_*\mathbb Q_\ell(j).
\end{equation}
This permits us to rewrite (5.8.4.1) as indicated:
\begin{equation}\tag{5.8.4.4}
\begin{tikzcd}[row sep=huge,column sep=large]
\operatorname{Prim}^n(X,\mathbb Q_\ell(j))
  \arrow[r,"\operatorname{Grif}"]
\&
H^1(U,\nu^*R^{n-1}\rho_*\mathbb Q_\ell(j))
  \arrow[d]
\\
{}
\&
H^1(U,\nu^*E^{n-1}\rho_*\mathbb Q_\ell(j)).
\end{tikzcd}
\end{equation}
By (5.7.1), the map (defined only by virtue of hypothesis (A))
\begin{equation}\tag{5.8.4.5}
u:\operatorname{Prim}^n(X,\mathbb Q_\ell(j))
\longrightarrow
H^1(\mathbb P^1,R^{n-1}\rho_*\mathbb Q_\ell(j))
\end{equation}
is injective.  Since $\mathbb P^1$ is simply connected, decomposition
(5.8.4.2) gives

\begin{equation}\tag{5.8.4.6}
H^1(\mathbb P^1,R^{n-1}\rho_*\mathbb Q_\ell(j))
\simeq
H^1(\mathbb P^1,E^{n-1}\rho_*\mathbb Q_\ell(j)).
\end{equation}
We can therefore complete (5.8.4.4) to a commutative square:
\begin{equation}\tag{5.8.4.7}
\begin{tikzcd}[column sep=huge,row sep=huge]
\operatorname{Prim}^n(X,\mathbb Q_\ell(j))
  \arrow[r,"\operatorname{Grif}"]
  \arrow[d,hook,"u"]
\&
H^1(U,\nu^*R^{n-1}\rho_*\mathbb Q_\ell(j))
  \arrow[d]
\\
H^1(\mathbb P^1,E^{n-1}\rho_*\mathbb Q_\ell(j))
  \arrow[r,"\nu^*"]
\&
H^1(U,\nu^*E^{n-1}\rho_*\mathbb Q_\ell(j)).
\end{tikzcd}
\end{equation}
It remains only to see that the lower arrow is injective.  This follows from
\begin{equation}\tag{5.8.4.8}
E^{n-1}\rho_*\mathbb Q_\ell(j)
\simeq
\nu_*\nu^*E^{n-1}\rho_*\mathbb Q_\ell(j)
\end{equation}
and the following well-known lemma.

}{E7II-XVIII-029}\SGApage{II}{306}
\subsection*{Lemma 5.8.5.}
\underline{Let $\nu:A\to B$ be a morphism of topoi, and let $F$ be a
sheaf on $A$.  Then the canonical map}
\begin{equation}\tag{5.8.5.1}
H^1(B,\nu_*F)\longrightarrow H^1(A,F)
\end{equation}
\underline{is injective.}

\underline{Proof.}
In a first-quadrant spectral sequence,
\begin{equation}\tag{5.8.5.2}
E_2^{1,0}\simeq E_\infty^{1,0}.
\end{equation}
Thus, for the Leray spectral sequence
\[
E_2^{p,q}=H^p(B,R^q\nu_*F)
\Longrightarrow H^{p+q}(A,F),
\]
we have
\begin{equation}\tag{5.8.5.3}
E_2^{1,0}=H^1(B,\nu_*F)
\simeq E_\infty^{1,0}
=F^1H^1(A,F).
\end{equation}
Taking $F=\nu^*E$, where $E$ is a sheaf on $B$, gives
\begin{equation}\tag{5.8.5.4}
H^1(B,\nu_*\nu^*E)
\lhook\joinrel\longrightarrow
H^1(A,\nu^*E).
\end{equation}


\SGApage{II}{\(\approx\)307}
\SGAvar{\subsection*{5.8.6.}
Let $L/k(\eta)$ be a finite extension of the rational function field of
$\mathbb P^1$, and put
\begin{equation}\tag{5.8.6.1}
S=\text{the normalization of $\mathbb P^1$ in $L$}.
\end{equation}
Make the base change $S\to\mathbb P^1$ in situation (5.8.3.2):
\begin{equation}\tag{5.8.6.1 bis}
\begin{tikzcd}[column sep=huge,row sep=large]
{} \& \widetilde X \arrow[dl,"f"] \arrow[dd,"\rho"]
\& \widetilde X_S \arrow[l,"\alpha_S"'] \arrow[dd,"\rho_S"] \\
X\&{}\&{} \\
{} \& \mathbb P^1 \& S\arrow[l,"\pi"'] .
\end{tikzcd}
\end{equation}
For every open immersion $V\xhookrightarrow{\mu}S$ over which $\rho_S$
is smooth, base change by $V\xrightarrow{\pi\mu}\mathbb P^1$ in
(5.8.3.2) gives
\begin{equation}\tag{5.8.6.2}
\begin{tikzcd}[column sep=huge,row sep=large]
{} \& \widetilde X \arrow[dl,"f"'] \arrow[dd,"\rho"]
\& \widetilde X_V \arrow[l,"\alpha_V"'] \arrow[dd,"\rho_V"]
  \arrow[dr,"f\times\rho_V"] \&{} \\
X\&{}\&{}\&X\times V\arrow[dl,"\operatorname{pr}_2"] \\
{} \& \mathbb P^1 \& V\arrow[l,"\pi\mu"'] \&{} .
\end{tikzcd}
\end{equation}
We still have (cf. (5.8.2.3))
\begin{equation}\tag{5.8.6.3}
(\alpha_Vf)^*\operatorname{Prim}^n(X,\mathbb Q_\ell(j))
\subset
\operatorname{Prim}^n(\widetilde X_V/_V,\mathbb Q_\ell(j)),
\end{equation}
so that we may again speak of the Griffiths homomorphism
\begin{equation}\tag{5.8.6.4}
\operatorname{Prim}^n(X,\mathbb Q_\ell(j))
\xrightarrow{\operatorname{Grif}}
H^1(V,R^{n-1}\rho_{V*}\mathbb Q_\ell(j)).
\end{equation}

\subsection*{Corollary 5.8.7.}
\underline{With the preceding notation and under the hypotheses of 5.8.4,
the Griffiths homomorphism ``modulo the fixed part,'' that is, the
composite}
\begin{equation}\tag{5.8.7.1}
\resizebox{\textwidth}{!}{
\begin{tikzcd}[ampersand replacement=\&,column sep=huge,row sep=huge]
\operatorname{Prim}^n(X,\mathbb Q_\ell(j))
  \arrow[r,"\operatorname{Grif}"]
\& H^1(V,R^{n-1}\rho_{V_*}\mathbb Q_\ell(j)) \arrow[d] \\
{}
\& H^1(V,(\pi\mu)^*E^{n-1}\rho_*(\mathbb Q_\ell(j)))
\simeq
\dfrac{H^1(V,R^{n-1}\rho_{V_*}\mathbb Q_\ell(j))}
{(f\times\rho_V)^*H^1(V,H^{n-1}(X,\mathbb Q_\ell(j))_V)}
\end{tikzcd}}
\end{equation}
\underline{is injective.}

\underline{Proof.}
By the specialization theorem for cohomology,
\[
R^{n-1}\rho_{V_*}\mathbb Q_\ell(j))
=(\pi\mu)^*R^{n-1}\rho_*\mathbb Q_\ell(j).
\]
Thus decomposition (5.8.4.2) rewrites (5.8.7.1) as indicated:
\[
\begin{tikzcd}[column sep=huge,row sep=huge]
\operatorname{Prim}^n(X,\mathbb Q_\ell(j))
  \arrow[r,"\operatorname{Grif}"]
\& H^1(V,(\pi\mu)^*R^{n-1}\rho_*\mathbb Q_\ell(j)) \arrow[d] \\
{}
\& H^1(V,(\pi\mu)^*E^{n-1}\rho_*\mathbb Q_\ell(j)).
\end{tikzcd}
\]
It is therefore enough to prove the injectivity of
\begin{equation}\notag
\operatorname{Prim}^n(X,\mathbb Q_\ell(j))
\longrightarrow
H^1(V,\mu^*\pi^*E^{n-1}\rho_*\mathbb Q_\ell(j))
\end{equation}
for arbitrarily small open immersions $V\xhookrightarrow{\mu}S$.  We are
reduced to the case
\[
\left\{
\begin{aligned}
U&\xhookrightarrow{\nu}\mathbb P^1
   &&\text{is an open subset over which $\rho$ is smooth},\\
V&\xrightarrow{\pi_U}U
   &&\text{is the normalization of $U$ in $L$}.
\end{aligned}
\right.
\]
From the cartesian square
\[
\begin{tikzcd}[column sep=huge,row sep=huge]
V\arrow[r,hook,"\mu"]\arrow[d,"\pi"]
\&S\arrow[d,"\pi"]\\
U\arrow[r,hook,"\nu"]
\&\mathbb P^1
\end{tikzcd}
\]
we obtain
\[
\mu^*\pi^*E^{n-1}\rho_*\mathbb Q_\ell(j)
=\pi_U^*\nu^*E^{n-1}\rho_*\mathbb Q_\ell.
\]
Consequently, (5.8.7.3) factors as
\[
\begin{tikzcd}[column sep=huge,row sep=huge]
\operatorname{Prim}^n(X,\mathbb Q_\ell(j)) \arrow[r]
\&H^1(U,\nu^*E^{n-1}\rho_*\mathbb Q_\ell(j))
  \arrow[d,"\pi_U^*"] \\
{}
\&H^1(V,\pi_U^*\nu^*E^{n-1}\rho_*\mathbb Q_\ell(j)).
\end{tikzcd}
\]
By (5.8.4), it is enough to verify that
\[
\pi_U^*:H^1(U,\nu^*E^{n-1}\rho_*\mathbb Q_\ell(j))
\longrightarrow
H^1(V,\pi_U^*\nu^*E^{n-1}\rho_*\mathbb Q_\ell(j))
\]
is injective.  Since $\pi_U:V\to U$ is finite and flat, trace theory
[SGA 4 XVII 6.2] supplies, for every $\mathbb Q_\ell$-sheaf $F$ on $U$
(for example $F=\nu^*E^{n-1}\rho_*\mathbb Q_\ell(j)$), a map
\[
\operatorname{tr}_{\pi_U}:H^1(V,\pi_U^*)\longrightarrow H^1(U,F(j))
\]
such that
\[
\operatorname{tr}_{\pi}\circ\pi^*
=\deg(\pi_U).\operatorname{id}.
\]
It follows that $H^*(U,F)\to H^*(V,\pi_U^*F)$ is injective, completing
the proof.



}{\subsection*{5.8.6.}
Let $L/k(\eta)$ be a finite extension of the rational function field of
$\mathbb P^1$, and put
\begin{equation}\tag{5.8.6.1}
S=\text{the normalization of $\mathbb P^1$ in $L$}.
\end{equation}
Make the base change $S\to\mathbb P^1$ in situation (5.8.3.2):
\begin{equation}\tag{5.8.6.1 bis}
\begin{tikzcd}[column sep=huge,row sep=large]
{} \& \widetilde X \arrow[dl,"f"] \arrow[dd,"\rho"]
\& \widetilde X_S \arrow[l,"\alpha_S"'] \arrow[dd,"\rho_S"] \\
X\&{}\&{} \\
{} \& \mathbb P^1 \& S\arrow[l,"\pi"'] .
\end{tikzcd}
\end{equation}
For every open immersion $V\xhookrightarrow{\mu}S$ over which $\rho_S$
is smooth, base change by $V\xrightarrow{\pi\mu}\mathbb P^1$ in
(5.8.3.2) gives
\begin{equation}\tag{5.8.6.2}
\begin{tikzcd}[column sep=huge,row sep=large]
{} \& \widetilde X \arrow[dl,"f"'] \arrow[dd,"\rho"]
\& \widetilde X_V \arrow[l,"\alpha_V"'] \arrow[dd,"\rho_V"]
  \arrow[dr,"f\times\rho_V"] \&{} \\
X\&{}\&{}\&X\times V\arrow[dl,"\operatorname{pr}_2"] \\
{} \& \mathbb P^1 \& V\arrow[l,"\pi\mu"'] \&{} .
\end{tikzcd}
\end{equation}
We still have (cf. (5.8.2.3))
\begin{equation}\tag{5.8.6.3}
\alpha_V^*f^*\operatorname{Prim}^n(X,\mathbb Q_\ell(j))
\subset
\operatorname{Prim}^n(\widetilde X_V/V,\mathbb Q_\ell(j)),
\end{equation}
so that we may again speak of the Griffiths homomorphism
\begin{equation}\tag{5.8.6.4}
\operatorname{Prim}^n(X,\mathbb Q_\ell(j))
\xrightarrow{\operatorname{Grif}}
H^1(V,R^{n-1}\rho_{V*}\mathbb Q_\ell(j)).
\end{equation}

\subsection*{Corollary 5.8.7.}
\underline{With the preceding notation and under the hypotheses of 5.8.4,
the Griffiths homomorphism ``modulo the fixed part,'' that is, the
composite}
\begin{equation}\tag{5.8.7.1}
\resizebox{\textwidth}{!}{
\begin{tikzcd}[ampersand replacement=\&,column sep=huge,row sep=huge]
\operatorname{Prim}^n(X,\mathbb Q_\ell(j))
  \arrow[r,"\operatorname{Grif}"]
\& H^1(V,R^{n-1}\rho_{V*}\mathbb Q_\ell(j)) \arrow[d] \\
{}
\& H^1(V,(\pi\mu)^*E^{n-1}\rho_*\mathbb Q_\ell(j))
\simeq
\dfrac{H^1(V,R^{n-1}\rho_{V*}\mathbb Q_\ell(j))}
{(f\times\rho_V)^*H^1(V,H^{n-1}(X,\mathbb Q_\ell(j))_V)}
\end{tikzcd}}
\end{equation}
\underline{is injective.}

\underline{Proof.}
By the specialization theorem for cohomology,
\[
R^{n-1}\rho_{V*}\mathbb Q_\ell(j)
=(\pi\mu)^*R^{n-1}\rho_*\mathbb Q_\ell(j).
\]
Thus decomposition (5.8.4.2) rewrites (5.8.7.1) as indicated:
\[
\begin{tikzcd}[column sep=huge,row sep=huge]
\operatorname{Prim}^n(X,\mathbb Q_\ell(j))
  \arrow[r,"\operatorname{Grif}"]
\& H^1(V,(\pi\mu)^*R^{n-1}\rho_*\mathbb Q_\ell(j)) \arrow[d] \\
{}
\& H^1(V,(\pi\mu)^*E^{n-1}\rho_*\mathbb Q_\ell(j)).
\end{tikzcd}
\]
It is therefore enough to prove the injectivity of
\begin{equation}\tag{5.8.7.3}
\operatorname{Prim}^n(X,\mathbb Q_\ell(j))
\longrightarrow
H^1(V,\mu^*\pi^*E^{n-1}\rho_*\mathbb Q_\ell(j))
\end{equation}
for arbitrarily small open immersions $V\xhookrightarrow{\mu}S$.  We are
reduced to the case
\[
\left\{
\begin{aligned}
U&\xhookrightarrow{\nu}\mathbb P^1
   &&\text{is an open subset over which $\rho$ is smooth},\\
V&\xrightarrow{\pi_U}U
   &&\text{is the normalization of $U$ in $L$}.
\end{aligned}
\right.
\]
From the cartesian square
\[
\begin{tikzcd}[column sep=huge,row sep=huge]
V\arrow[r,hook,"\mu"]\arrow[d,"\pi_U"]
\&S\arrow[d,"\pi"]\\
U\arrow[r,hook,"\nu"]
\&\mathbb P^1
\end{tikzcd}
\]
we obtain
\[
\mu^*\pi^*E^{n-1}\rho_*\mathbb Q_\ell(j)
=\pi_U^*\nu^*E^{n-1}\rho_*\mathbb Q_\ell(j).
\]
Consequently, (5.8.7.3) factors as
\[
\begin{tikzcd}[column sep=huge,row sep=huge]
\operatorname{Prim}^n(X,\mathbb Q_\ell(j)) \arrow[r]
\&H^1(U,\nu^*E^{n-1}\rho_*\mathbb Q_\ell(j))
  \arrow[d,"\pi_U^*"] \\
{}
\&H^1(V,\pi_U^*\nu^*E^{n-1}\rho_*\mathbb Q_\ell(j)).
\end{tikzcd}
\]
By (5.8.4), it is enough to verify that
\[
\pi_U^*:H^1(U,\nu^*E^{n-1}\rho_*\mathbb Q_\ell(j))
\longrightarrow
H^1(V,\pi_U^*\nu^*E^{n-1}\rho_*\mathbb Q_\ell(j))
\]
is injective.  Since $\pi_U:V\to U$ is finite and flat, trace theory
[SGA 4 XVII 6.2] supplies, for every $\mathbb Q_\ell$-sheaf $F$ on $U$
(for example $F=\nu^*E^{n-1}\rho_*\mathbb Q_\ell(j)$), a map
\[
\operatorname{tr}_{\pi_U}:H^1(V,\pi_U^*F)\longrightarrow H^1(U,F)
\]
such that
\[
\operatorname{tr}_{\pi_U}\circ\pi_U^*
=\deg(\pi_U)\,\operatorname{id}.
\]
It follows that $H^*(U,F)\to H^*(V,\pi_U^*F)$ is injective, completing
the proof.



}{E7II-XVIII-030}\SGApage{II}{311}
\subsection*{5.9. \underline{An example in which hypothesis} (A)
\underline{is not satisfied}.}

Let $X$ be a smooth quadric surface in $\mathbb P^3$.  Then $X$ satisfies
(LV) (5.2.4).  Take a pencil (a Lefschetz pencil, if desired) of
\underline{hyperplanes} whose generic member and axis meet $X$
transversely.  Then (as Deligne pointed out to me) \underline{the projection}
\[
\rho:\widetilde X\longrightarrow\mathbb P^1
\]
\underline{does not satisfy} (A).  Indeed,
$R^2\rho_*\mathbb Q_\ell$ is not constant: the singular fibers are singular
conics, hence reducible, and in a Lefschetz pencil each is the union of two
lines meeting at one point (this is the definition of a Lefschetz pencil in
this simple case).  Thus
\[
\dim_{\mathbb Q_\ell}H^2(\text{singular fiber},\mathbb Q_\ell)=2,
\]
whereas
\[
\dim_{\mathbb Q_\ell}H^2(\text{smooth fiber},\mathbb Q_\ell)=1.
\]

Moreover, the conclusion of Theorem 5.8.4 is false here.  Let
$U\xhookrightarrow{\nu}\mathbb P^1$ be a nonempty open subset over which
$\rho$ is smooth.  Since the fibers of $\rho$ over $U$ are smooth curves of
genus zero, we have
\[
\nu^*R^1\rho_*\mathbb Q_\ell=0
\]
and hence
\[
H^1(U,\nu^*R^1\rho_*\mathbb Q_\ell)=0.
\]
Consequently, the Griffiths homomorphism
\[
\operatorname{Prim}^2(X,\mathbb Q_\ell)
\xrightarrow{\operatorname{Grif}}
H^1(U,\nu^*R^1\rho_*\mathbb Q_\ell)
\]
is zero.  But on a quadric surface,
\[
\dim_{\mathbb Q_\ell}\operatorname{Prim}^2(X,\mathbb Q_\ell)=1.
\]
Likewise, the conclusion of 5.7 fails here.


\SGApage{II}{\(\approx\)312}
\section*{6. \underline{Application of the local theory of vanishing cycles}\\
(\underline{Picard}--\underline{Lefschetz})}

\subsection*{Proposition 6.1.}
\underline{Let $Z$ be a hypersurface in projective space $\mathbb P^r_k$,
$r\geq2$, where $k$ is an algebraically closed field.  Let}
\[
i_o:D_o\lhook\joinrel\longrightarrow\mathbb P^r
\]
\underline{be the inclusion of a line which meets $Z$ only in its smooth
locus and meets it transversely.  Choose a geometric point $\bar a$ of
$D_o-D_o\cap Z$, and denote by}
\[
\pi_1^{\mathrm{mod}}(D_o-D_o\cap Z,\bar a)
\]
\underline{the quotient of $\pi_1(D_o-D_o\cap Z,\bar a)$ classifying the
étale coverings of $D_o-D_o\cap Z$ which are tamely ramified at every point
of $D_o\cap Z$.  Then:}

\SGApage{II}{\(\approx\)313}
\subsubsection*{6.1.1.}
\underline{The canonical homomorphism}
\begin{equation}\tag{6.1.1.1}
i_{o*}:\pi_1^{\mathrm{mod}}(D_o-D_o\cap Z,\bar a)
\longrightarrow
\pi_1^{\mathrm{mod}}(\mathbb P^r-Z,\bar a)
\end{equation}
\underline{is surjective.}

\subsubsection*{6.1.2.}
\underline{Every choice of a geometric point $\bar x_o\in D_o\cap Z$ and
of a} ``\underline{path}'' \underline{in $D_o$ from $\bar x_o$ to $\bar a$ determines a
homomorphism}
\begin{equation}\tag{6.1.2.1}
\prod_{\ell\ne\operatorname{char}(k)}\mathbb Z_\ell(1)
\longrightarrow
\pi_1^{\mathrm{mod}}(D_o-D_o\cap Z,\bar a)
\end{equation}
\underline{whose image is} ``\underline{the}'' \underline{inertia group at $\bar x_o$.  Composing
this morphism with $i_{o*}$ gives a homomorphism}
\begin{equation}\tag{6.1.2.2}
\prod_{\ell\ne\operatorname{char}(k)}\mathbb Z_\ell(1)
\longrightarrow
\pi_1^{\mathrm{mod}}(\mathbb P^r-Z,\bar a).
\end{equation}
\underline{If $Z$ is irreducible, all these homomorphisms are conjugate to
one another.}

\underline{Proof.}
Let $T$ be the subscheme of the Grassmannian formed by the lines in
$\mathbb P^r$ which
\begin{enumerate}[label=(\roman*)]
\item pass through $\bar a$;
\item meet $Z$ only in its smooth locus, and meet it there transversely.
\end{enumerate}
Let $\eta$ be its generic point, and let
\[
i_\eta:D_\eta\longrightarrow\mathbb P^r\otimes_k k(\eta)
\]
be the corresponding line (thus a line ``joining $\bar a$ to the generic
point of $\mathbb P^r_k$'').  Choose also a geometric point $\bar\eta$ over
$\eta$, and write $D_{\bar\eta}$ for the corresponding ``geometric'' line.
Let $o\in T$ be the point corresponding to the given line $D_o$.


\SGAvar{By [SGA 1, XIII 2.8], there is a specialization homomorphism (depending on
a path in $T$ from $\bar\eta$ to $O$, and hence defined up to conjugacy)
\begin{equation}\tag{6.1.3}
\pi_1^{\mathrm{mod}}(D_{\bar\eta}-D_{\bar\eta}\cap Z),\bar a)
\xrightarrow{\ \bar\eta\longrightarrow o\ }
\pi_1^{\mathrm{mod}}(D_o-D_o\cap Z,\bar a)
\end{equation}
such that the following assertions hold.

\subsubsection*{6.1.4.}
The local-monodromy homomorphisms (6.1.2.1)
\[
\prod_{\ell\ne\operatorname{char}(k)}\mathbb Z_\ell(1)
\longrightarrow
\pi_1(D_o-D_o\cap Z,\bar a)
\]
are obtained by composing the specialization arrow (6.1.3) with the local
monodromy homomorphisms for $D_{\bar\eta}$:
\begin{equation}\tag{6.1.4.1}
\prod_{\ell\ne\operatorname{char}(k)}\mathbb Z_\ell(1)
\longrightarrow
\pi_1^{\mathrm{mod}}(D_{\bar\eta}-D_{\bar\eta}\cap Z,\bar a)
\xrightarrow{\ \bar\eta\longrightarrow o\ }
\pi_1^{\mathrm{mod}}(D_o-D_o\cap Z,\bar a).
\end{equation}

\subsubsection*{6.1.5.}
The following diagram is commutative:
\begin{equation}\notag
\begin{tikzcd}[column sep=huge,row sep=huge]
{} \& \pi_1^{\mathrm{mod}}(\mathbb P^r-Z,\bar a) \& {} \\
\pi_1^{\mathrm{mod}}(D_\eta-D_\eta\cap Z,\bar a)
  \arrow[ur,"i_{\eta*}"]
\& {}
\& \pi_1^{\mathrm{mod}}(D_o-D_o\cap Z,\bar a)
  \arrow[ul,"i_{o*}"'] \\
{} \& \pi_1^{\mathrm{mod}}(D_{\bar\eta}-D_{\bar\eta}\cap Z,\bar a)
  \arrow[uu,"i_{\bar\eta_*}"']
  \arrow[ul]
  \arrow[ur,"\bar\eta\longrightarrow o"'] \& {} .
\end{tikzcd}
\end{equation}

Assume this for the moment, and first prove 6.1.1.  By the commutativity of
6.1.5, it is enough to prove that
\[
i_{\bar\eta*}:
\pi_1^{\mathrm{mod}}(D_{\bar\eta}-D_{\bar\eta}\cap Z,\bar a)
\longrightarrow
\pi_1^{\mathrm{mod}}(\mathbb P^r-Z,\bar a)
\]
is surjective.  For this, it is enough to prove that the analogous
homomorphism of the \underline{ordinary} fundamental groups
\begin{equation}\tag{6.1.6}
i_{\bar\eta}:
\pi_1(D_{\bar\eta}-D_{\bar\eta}\cap Z,\bar a)
\longrightarrow
\pi_1(\mathbb P^r-Z,\bar a)
\end{equation}
is surjective.  In other words, for every étale covering
$E\to\mathbb P^r-Z$ with $E$ connected (and hence irreducible, since
$\mathbb P^r-Z$ is smooth and therefore unibranch [EGA IV, 18.10.3]), its
inverse image
\[
E_{\bar\eta}
=E\times_{\mathbb P^r-Z}(D_{\bar\eta}-D_{\bar\eta}\cap Z)
\]
over $D_{\bar\eta}-D_{\bar\eta}\cap Z$ must remain geometrically
irreducible.  It is enough to use the following variant of Bertini's theorem
([EGA V, citation left blank in the source]) for the composite
$E\to\mathbb P^r-Z\lhook\joinrel\longrightarrow\mathbb P^r$ and $s=1$.

}{By [SGA 1, XIII 2.8], there is a specialization homomorphism (depending on
a path in $T$ from $\bar\eta$ to $o$, and hence defined up to conjugacy)
\begin{equation}\tag{6.1.3}
\pi_1^{\mathrm{mod}}(D_{\bar\eta}-D_{\bar\eta}\cap Z,\bar a)
\xrightarrow{\ \bar\eta\longrightarrow o\ }
\pi_1^{\mathrm{mod}}(D_o-D_o\cap Z,\bar a)
\end{equation}
such that the following assertions hold.

\subsubsection*{6.1.4.}
The local-monodromy homomorphisms (6.1.2.1)
\[
\prod_{\ell\ne\operatorname{char}(k)}\mathbb Z_\ell(1)
\longrightarrow
\pi_1^{\mathrm{mod}}(D_o-D_o\cap Z,\bar a)
\]
are obtained by composing the specialization arrow (6.1.3) with the local
monodromy homomorphisms for $D_{\bar\eta}$:
\begin{equation}\tag{6.1.4.1}
\prod_{\ell\ne\operatorname{char}(k)}\mathbb Z_\ell(1)
\longrightarrow
\pi_1^{\mathrm{mod}}(D_{\bar\eta}-D_{\bar\eta}\cap Z,\bar a)
\xrightarrow{\ \bar\eta\longrightarrow o\ }
\pi_1^{\mathrm{mod}}(D_o-D_o\cap Z,\bar a).
\end{equation}

\subsubsection*{6.1.5.}
The following diagram is commutative:
\begin{equation}\notag
\begin{tikzcd}[column sep=huge,row sep=huge]
{} \& \pi_1^{\mathrm{mod}}(\mathbb P^r-Z,\bar a) \& {} \\
\pi_1^{\mathrm{mod}}(D_\eta-D_\eta\cap Z,\bar a)
  \arrow[ur,"i_{\eta*}"]
\& {}
\& \pi_1^{\mathrm{mod}}(D_o-D_o\cap Z,\bar a)
  \arrow[ul,"i_{o*}"'] \\
{} \& \pi_1^{\mathrm{mod}}(D_{\bar\eta}-D_{\bar\eta}\cap Z,\bar a)
  \arrow[uu,"i_{\bar\eta*}"']
  \arrow[ul]
  \arrow[ur,"\bar\eta\longrightarrow o"'] \& {} .
\end{tikzcd}
\end{equation}

Assume this for the moment, and first prove 6.1.1.  By the commutativity of
6.1.5, it is enough to prove that
\[
i_{\bar\eta*}:
\pi_1^{\mathrm{mod}}(D_{\bar\eta}-D_{\bar\eta}\cap Z,\bar a)
\longrightarrow
\pi_1^{\mathrm{mod}}(\mathbb P^r-Z,\bar a)
\]
is surjective.  For this, it is enough to prove that the analogous
homomorphism of the \underline{ordinary} fundamental groups
\begin{equation}\tag{6.1.6}
i_{\bar\eta*}:
\pi_1(D_{\bar\eta}-D_{\bar\eta}\cap Z,\bar a)
\longrightarrow
\pi_1(\mathbb P^r-Z,\bar a)
\end{equation}
is surjective.  In other words, for every étale covering
$E\to\mathbb P^r-Z$ with $E$ connected (and hence irreducible, since
$\mathbb P^r-Z$ is smooth and therefore unibranch [EGA IV, 18.10.3]), its
inverse image
\[
E_{\bar\eta}
=E\times_{\mathbb P^r-Z}(D_{\bar\eta}-D_{\bar\eta}\cap Z)
\]
over $D_{\bar\eta}-D_{\bar\eta}\cap Z$ must remain geometrically
irreducible.  It is enough to use the following variant of Bertini's theorem
([EGA V, citation left blank in the source]) for the composite
$E\to\mathbb P^r-Z\lhook\joinrel\longrightarrow\mathbb P^r$ and $s=1$.

}{E7II-XVIII-031}\SGApage{II}{315}
\subsubsection*{6.1.6.1. Bertini's theorem.}
Let $k$ be a field, let $\mathbb P^r=\mathbb P^r_k$, and let
$f:X\to\mathbb P^r$ be a nonconstant morphism, where $X$ is a geometrically
irreducible algebraic scheme.  For every integer
$s<\dim\overline{f(X)}$, let $L_t$ denote the \underline{generic} linear
subspace of codimension $s$ in $\mathbb P^r$.  Then
$X_t=f^{-1}(L_t)$ is geometrically irreducible.


\SGAvar{We now prove 6.1.2.  By 6.1.4 and the commutativity of (6.1.5), it is
enough to prove that all the local-monodromy homomorphisms (6.1.3)
\underline{for} $D_{\bar\eta}$,
\begin{equation}\tag{6.1.7.1}
\prod_{\ell\ne\operatorname{char}(k)}\mathbb Z_\ell(1)
\longrightarrow
\pi_1^{\mathrm{mod}}(D_{\bar\eta}-D_{\bar\eta}\cap Z,\bar a),
\end{equation}
become conjugate after composition with the homomorphism
\begin{equation}\tag{6.1.7.2}
\pi_1^{\mathrm{mod}}(D_{\bar\eta}-D_{\bar\eta}\cap Z,\bar a)
\longrightarrow
\pi_1^{\mathrm{mod}}(D_\eta-D_\eta\cap Z,\bar a).
\end{equation}
To prove this, observe that for every finite Galois extension $K$ of
$k(\eta)$, with
\[
k(\eta)\subset K\subset K(\bar\eta),
\]
the scheme
\[
D_K\stackrel{\mathrm{def}}{=}D_\eta\times_{k(\eta)}K
\]
is an étale covering of $D_\eta$.  Thus
$D_K-D_K\cap Z$ is a tame covering of $D_\eta-D_\eta\cap Z$, and hence
$\pi_1^{\mathrm{mod}}(D_{\overline\eta}-D_\eta\cap Z,\bar a)$ acts on the scheme
$D_K\cap Z$.  We take $K$ sufficiently large that every point of
$D_K\cap Z$ is $K$-rational; this is possible because, by hypothesis,
$D_\eta\cap Z$ is étale over $k(\eta)$.

}{We now prove 6.1.2.  By 6.1.4 and the commutativity of (6.1.5), it is
enough to prove that all the local-monodromy homomorphisms (6.1.3)
\underline{for} $D_{\bar\eta}$,
\begin{equation}\tag{6.1.7.1}
\prod_{\ell\ne\operatorname{char}(k)}\mathbb Z_\ell(1)
\longrightarrow
\pi_1^{\mathrm{mod}}(D_{\bar\eta}-D_{\bar\eta}\cap Z,\bar a),
\end{equation}
become conjugate after composition with the homomorphism
\begin{equation}\tag{6.1.7.2}
\pi_1^{\mathrm{mod}}(D_{\bar\eta}-D_{\bar\eta}\cap Z,\bar a)
\longrightarrow
\pi_1^{\mathrm{mod}}(D_\eta-D_\eta\cap Z,\bar a).
\end{equation}
To prove this, observe that for every finite Galois extension $K$ of
$k(\eta)$, with
\[
k(\eta)\subset K\subset k(\bar\eta),
\]
the scheme
\[
D_K\stackrel{\mathrm{def}}{=}D_\eta\times_{k(\eta)}K
\]
is an étale covering of $D_\eta$.  Thus
$D_K-D_K\cap Z$ is a tame covering of $D_\eta-D_\eta\cap Z$, and hence
$\pi_1^{\mathrm{mod}}(D_\eta-D_\eta\cap Z,\bar a)$ acts on the scheme
$D_K\cap Z$.  We take $K$ sufficiently large that every point of
$D_K\cap Z$ is $K$-rational; this is possible because, by hypothesis,
$D_\eta\cap Z$ is étale over $k(\eta)$.

}{E7II-XVIII-032}\SGApage{II}{\(\approx\)316}
It remains to prove that
$\pi_1^{\mathrm{mod}}(D_\eta-D_\eta\cap Z,\bar a)$ acts
\underline{transitively} on the points of $D_K\cap Z$.  This will prove
that the composite local-inertia homomorphisms
\[
\prod_{\ell\ne\operatorname{char}(k)}\mathbb Z_\ell(1)
\longrightarrow
\pi_1^{\mathrm{mod}}(D_\eta-D_\eta\cap Z,\bar a)
\]
attached to the various points of $D_K\cap Z$ are conjugate.  Now the
action on $D_K\cap Z$ factors through the canonical quotient

\SGAvar{\[
\pi_1^{\mathrm{mod}}(D_\eta-D_\eta\cap Z,\bar a)
\longrightarrow
\pi_1(\eta,\bar\eta)
=\operatorname{Gal}\bigl(/k(\eta).
\]
The action of $\operatorname{Gal}(/k(\eta))$ on
$K_K\cap Z$ is the one compatible with the inclusion
\[
D_K\cap Z\lhook\joinrel\longrightarrow\mathbb P^r_k\times_k K,
\]
where the Galois group acts on $\mathbb P^r_k\times_k K$ only through
the second factor.  It follows that the desired assertion is equivalent
to saying that $D_\eta\cap Z$ has only one point, that is, that it is the
spectrum of a \underline{field}.  In elementary terms, this is precisely
the following Bertini-type fact (take an affine equation $F=0$ for $Z$
and suppose that $\bar a$ is the origin).

}{\[
\pi_1^{\mathrm{mod}}(D_\eta-D_\eta\cap Z,\bar a)
\longrightarrow
\pi_1(\eta,\bar\eta)
=\operatorname{Gal}\bigl(k(\bar\eta)/k(\eta)\bigr).
\]
The action of $\operatorname{Gal}(k(\bar\eta)/k(\eta))$ on
$D_K\cap Z$ is the one compatible with the inclusion
\[
D_K\cap Z\lhook\joinrel\longrightarrow\mathbb P^r_k\times_k K,
\]
where the Galois group acts on $\mathbb P^r_k\times_k K$ only through
the second factor.  It follows that the desired assertion is equivalent
to saying that $D_\eta\cap Z$ has only one point, that is, that it is the
spectrum of a \underline{field}.  In elementary terms, this is precisely
the following Bertini-type fact (take an affine equation $F=0$ for $Z$
and suppose that $\bar a$ is the origin).

}{E7II-XVIII-033}\SGApage{II}{\(\approx\)317}
\subsubsection*{Fact 6.1.8.}
Let $k$ be an algebraically closed field, let $r\ge2$, and let
$F(X_1,\ldots,X_r)\in k[X_1,\ldots,X_r]$ be irreducible.  Then
\[
G(T)=F(X_1T,\ldots,X_rT)\in k(X_1,\ldots,X_r)[T]
\]
is irreducible as a polynomial in $T$ with coefficients in
$k(X_1,\ldots,X_r)$.

This completes the proof of 6.1.2.

\subsubsection*{Remark 6.1.9.}
Another method of proving 6.1.2, valid under considerably more general
conditions, would be simply to establish a general compatibility for
``monodromy homomorphisms'': the composite homomorphisms in 6.1.2 are
conjugate to the monodromy homomorphism at the generic point of $Z$, and
therefore are conjugate to one another.

\SGApage{II}{318}
\SGAvar{\subsubsection*{Corollary 6.1.10.}
Let $X$ be a smooth connected projective variety of dimension $n$ over an
algebraically closed field $k$.  Suppose that $n=\dim X$ is even, or that
$\operatorname{card}(k)\ne2$.
Let $D$ be a Lefschetz pencil for $X$, let $\eta$ be its generic point,
and let $\bar\eta$ be a geometric point over $\eta$.  Then the natural
representation of
$\pi_1(D-D\cap\check X,\bar\eta)$ on
$H^{n-1}(X_{\bar\eta},\mathbb Q_\ell)$ factors through
$\pi_1^{\mathrm{mod}}(D-D\cap\check X,\bar\eta)$.  If this
representation is nontrivial, then $D\cap\check X\ne\varnothing$, and
the inertia group at every point of $D\cap\check X$ acts nontrivially.

}{\subsubsection*{Corollary 6.1.10.}
Let $X$ be a smooth connected projective variety of dimension $n$ over an
algebraically closed field $k$.  Suppose that $n=\dim X$ is even, or that
$\operatorname{char}(k)\ne2$.
Let $D$ be a Lefschetz pencil for $X$, let $\eta$ be its generic point,
and let $\bar\eta$ be a geometric point over $\eta$.  Then the natural
representation of
$\pi_1(D-D\cap\check X,\bar\eta)$ on
$H^{n-1}(X_{\bar\eta},\mathbb Q_\ell)$ factors through
$\pi_1^{\mathrm{mod}}(D-D\cap\check X,\bar\eta)$.  If this
representation is nontrivial, then $D\cap\check X\ne\varnothing$, and
the inertia group at every point of $D\cap\check X$ acts nontrivially.

}{E7II-XVIII-034}\subsubsection*{Proof.}
If $\dim(\check X)\le r-2$, then $D\cap\check X$ is empty and there is
nothing to prove, since $\mathbb P^1$ is simply connected.  Otherwise
$\check X$ is a hypersurface in $\mathbb P^r$ [XVII, 3.1.4], and 6.1.1
applies by 3.2.10 and XVII 3.5, with $Z=\check X$ and $D_o=D$.  The fact
that the action of $\pi_1$ on
$H^{n-1}(X_{\bar\eta},\mathbb Q_\ell)$ is tame was established in the
local theory by the Picard--Lefschetz formula [XV, 3.4], recalled below.
Finally, the nontrivial fact that
$\pi_1^{\mathrm{mod}}(D-D\cap\check X,\bar\eta)$ is generated by the
inertia groups at the various points of $D\cap\check X$ [SGA 1, XIII
2.10] gives the required conclusion.

We shall now use the preceding result to study condition (A) of (5.3.5).
First recall the following consequence of the Picard--Lefschetz formula
[XV, 3.4].


\SGApage{II}{319}
\subsubsection*{Proposition 6.2.}
Let $X$ be a smooth connected projective variety of dimension $n$ over
the algebraically closed field $k$.  Let $D$ be a Lefschetz pencil, let
$\eta$ be its generic point, and let $\bar\eta$ be a geometric point
over $\eta$.  Then the pencil $D$ satisfies condition (A) if and only if
the inertia group at every point of $D\cap\check X$ acts nontrivially on
$H^{n-1}(X_{\bar\eta},\mathbb Q_\ell)$.

\subsubsection*{Theorem 6.3.}
Under the hypotheses of 6.2, the following assertions hold.

\medskip
\noindent
\textup{(6.3.1)} If $n=\dim X$ is odd, condition (A) is satisfied by
every Lefschetz pencil for $X$.

\medskip
\noindent
\textup{(6.3.2)} Suppose that $n=\dim X$ is even, or that
$\operatorname{char}(k)\ne2$.  Then the following conditions are
equivalent:
\begin{enumerate}
\item[(i)] $D$ does not satisfy condition (A) of (5.3.5);
\item[(ii)] $\pi_1(D-D\cap\check X)$ acts trivially on
  $H^{n-1}(X_{\bar\eta},\mathbb Q_\ell)$, and
  $D\cap\check X\ne\varnothing$;
\item[(iii)] there is a point of $D\cap\check X$ whose inertia group
  acts trivially on $H^{n-1}(X_{\bar\eta},\mathbb Q_\ell)$;
\item[(iv)] the sheaf $R^{n-1}\rho_*\mathbb Q_\ell$ is constant on
  $D$, and $D\cap\check X$ is nonempty;
\item[(v)] the restriction
  $R^{n-1}\rho_*\mathbb Q_\ell\vert_U$ is constant on $U$, where $U$
  is a nonempty open subset of $D$ over which $\rho$ is smooth, and
  $\rho$ is not smooth.
\end{enumerate}

\SGApage{II}{\(\approx\)320}
\noindent
\textup{(6.3.3)} If, in addition, $X$ satisfies condition (LV) of
(5.2.2), conditions (i)--(v) are also equivalent to:
\begin{enumerate}
\item[(vi)] the restriction map
\[
H^{n-1}(X,\mathbb Q_\ell)
\longrightarrow H^{n-1}(X_{\bar\eta},\mathbb Q_\ell)
\]
is surjective (and hence is an isomorphism by the weak Lefschetz
theorem).
\end{enumerate}

\SGAvar{\subsubsection*{Proof.}
Assertion 6.3.1 follows from XV 3.4.  For the equivalences in 6.3.2,
(i) is equivalent to (iii) by 6.2, (ii) to (iii) by 6.1.6, (ii) to (v)
tautologically, and (iv) to (v) by 5.3.4.  Finally, if $X$ satisfies
condition (LV), 5.6.10 gives the equivalence of (ii) and (vi).

\subsubsection*{Remark 6.3.4.}
Condition (A) does not depend on the choice of the Lefschetz pencil
$D$, but only on the chosen (Lefschetz) immersion
$i:X\longrightarrow\mathbb P^r$.  This is clear from (6.3.1) when
$\dim X$ is odd.  When $\dim X$ is even, condition (A) also says that
the local system on $\mathbb P^r- X$ formed by the groups
$H^{n-1}(X\cap H_s,Q_\ell)$ is constant.  This follows from
criterion (ii), from the fact that the local system in question is
tamely ramified at the generic point of $\check X$, and from the
surjectivity in (6.1.2).

\subsubsection*{Corollary 6.4.}
Under the hypotheses of 6.2, suppose that
$n=\dim X$ is even, or $\operatorname{char}(k)\ne2$ and that $X$
satisfies condition (LV).
Then there is an integer $d_o$ such that, for $d\ge d_o$, every
Lefschetz pencil of sections of $X$ by hypersurfaces of degree $d$
satisfies condition (A).

}{\subsubsection*{Proof.}
Assertion 6.3.1 follows from XV 3.4.  For the equivalences in 6.3.2,
(i) is equivalent to (iii) by 6.2, (ii) to (iii) by 6.1.10, (ii) to (v)
tautologically, and (iv) to (v) by 5.3.4.  Finally, if $X$ satisfies
condition (LV), 5.6.10 gives the equivalence of (ii) and (vi).

\subsubsection*{Remark 6.3.4.}
Condition (A) does not depend on the choice of the Lefschetz pencil
$D$, but only on the chosen (Lefschetz) immersion
$i:X\longrightarrow\mathbb P^r$.  This is clear from (6.3.1) when
$\dim X$ is odd.  When $\dim X$ is even, condition (A) also says that
the local system on $\mathbb P^r-\check X$ formed by the groups
$H^{n-1}(X\cap H_s,\mathbb Q_\ell)$ is constant.  This follows from
criterion (ii), from the fact that the local system in question is
tamely ramified at the generic point of $\check X$, and from the
surjectivity in (6.1.1).

\subsubsection*{Corollary 6.4.}
Under the hypotheses of 6.2, suppose that
$n=\dim X$ is even or $\operatorname{char}(k)\ne2$, and that $X$
satisfies condition (LV).
Then there is an integer $d_o$ such that, for $d\ge d_o$, every
Lefschetz pencil of sections of $X$ by hypersurfaces of degree $d$
satisfies condition (A).

}{E7II-XVIII-035}\SGApage{II}{\(\approx\)321}
\subsubsection*{Proof.}
By the equivalence (i)$\Longleftrightarrow$(vi) of (6.3.2), it is
enough to prove that, if $Y(d)$ denotes a smooth hypersurface section of
$X$ of degree $d$, then
\begin{equation}\tag{6.4.1}
\dim H^{n-1}(Y(d),\mathbb Q_\ell)\longrightarrow\infty
\qquad\text{as }d\longrightarrow\infty.
\end{equation}
Indeed, we have the following precise result.

\subsubsection*{Lemma 6.4.2.}
Let $X$ be a smooth connected projective variety of dimension $n\ge1$
over an algebraically closed field $k$, and let $Y(d)$ denote a
hypersurface section of degree $d$.  Then there is a polynomial
$f(T)\in\mathbb Z[T]$ such that
\begin{align}
f(T)&=\deg(X)T^n+\text{a polynomial of degree at most }n-1,
  \tag{6.4.2.1}\\
f(d)&=\dim H^{n-1}(Y(d),\mathbb Q_\ell)\qquad(d\ge1).
  \tag{6.4.2.2}
\end{align}

\SGAvar{\subsubsection*{Proof.}
By XVII 5.7.5, if $C(X)$ denotes the total Chern class of $X$ and $H$
the first Chern class of a hyperplane, the Euler--Poincaré
characteristic of $Y(d)$ is
\begin{equation}\tag{6.4.2.3}
\chi(Y(d)
=\int_X\frac{dH\,C(X)}{1+dH}
=\sum_{i=1}^{n}(-1)^{i-1}d^i\int_X H^iC_{n-i}.
\end{equation}
On the other hand, Poincaré duality and weak Lefschetz give
\begin{equation}\tag{6.4.2.4}
\chi(Y(d))
=(-1)^{n-1}\dim H^{n-1}(Y(d),\mathbb Q_\ell)
+2\sum_{i\ge1}(-1)^{n-1-i}
  \dim H^{n-1-i}(X,\mathbb Q_\ell).
\end{equation}
Combining (6.4.2.3) and (6.4.2.4) gives the result.



}{\subsubsection*{Proof.}
By XVII 5.7.5, if $C(X)$ denotes the total Chern class of $X$ and $H$
the first Chern class of a hyperplane, the Euler--Poincaré
characteristic of $Y(d)$ is
\begin{equation}\tag{6.4.2.3}
\chi(Y(d))
=\int_X\frac{dH\,C(X)}{1+dH}
=\sum_{i=1}^{n}(-1)^{i-1}d^i\int_X H^iC_{n-i}.
\end{equation}
On the other hand, Poincaré duality and weak Lefschetz give
\begin{equation}\tag{6.4.2.4}
\chi(Y(d))
=(-1)^{n-1}\dim H^{n-1}(Y(d),\mathbb Q_\ell)
+2\sum_{i\ge1}(-1)^{n-1-i}
  \dim H^{n-1-i}(X,\mathbb Q_\ell).
\end{equation}
Combining (6.4.2.3) and (6.4.2.4) gives the result.



}{E7II-XVIII-036}\SGApage{II}{322}
\subsection*{6.5. A reminder on the Picard--Lefschetz formula (XV 3.4)}

Let $S$ be the spectrum of a strictly henselian valuation ring, with
generic point $\eta$, geometric generic point $\bar\eta$, and closed
point $s$.  Let
\[
\rho:X\longrightarrow S
\]
be a proper morphism, where $X$ is a connected regular scheme of
dimension $n$.  Suppose that $\rho$ is smooth away from a single point
$x$ in the special fiber and that $x$ is an ordinary quadratic
singularity of that fiber.

Let $I$ denote the inertia group of $S$ at $s$.  By XV 3.4, $I$ acts
trivially on $H^i(X_{\bar\eta},\mathbb Q_\ell)$ for $i\ne n-1$, and its
action on $H^{n-1}(X_{\bar\eta},\mathbb Q_\ell)$ is described as follows.

\SGAvar{\subsubsection*{6.5.1.}
If $n=2m+1$, then $I$ acts through a quotient isomorphic to
$\mathbb Z/2\mathbb Z$ (if $\operatorname{char}(k)\ne2$, there is only
one such quotient).  There is an element
\[
\delta\in H^{n-1}(X_{\bar\eta},\mathbb Q_\ell(m)),
\]
determined uniquely up to sign and called the \underline{vanishing
cycle at $s$}, such that
\begin{equation}\tag{6.5.1.1}
\delta\cdot\delta=(-1)^m2
\quad\text{in}\quad
H^{2n-2}(X_{\bar\eta},\mathbb Q_\ell(n-1))\simeq\mathbb Q_\ell.
\end{equation}
The nontrivial element $\sigma$ of $\mathbb Z/2\mathbb Z$ acts on
$H^{i}(X_{\bar\eta},\mathbb Q_\ell)$ by the reflection
\begin{equation}\tag{6.5.2.2}
\sigma x=x-(-1)^m(x\cdot\delta)\delta.
\end{equation}

Here $x\cdot\delta$ belongs to
$H^{2n-2}(X_{\bar\eta},\mathbb Q_\ell(m)\simeq\mathbb Q_\ell(-m)$,
so $(x\cdot\delta)\delta$ defines an element of
$H^{n-1}(X_{\bar\eta},\mathbb Q_\ell)$.

\subsubsection*{6.5.2.}
If $n=2m+2$, then $I$ acts through its canonical quotient
$\mathbb Z_\ell(1)$.  There is an element
\[
\delta\in H^{n-1}(X_{\bar\eta},\mathbb Q_\ell(m),
\]
canonically determined up to sign and called the \underline{vanishing
cycle at $s$}, such that $\sigma\in\mathbb Z_\ell(1)$ acts on
$\chi\in H^{n-1}(X_{\bar\eta},\mathbb Q_\ell)$ by the transvection
\begin{equation}\tag{6.5.2.1}
\sigma(x)=x-(-1)^m\sigma\,(x\cdot\delta)\delta.
\end{equation}
Indeed, $x\cdot\delta$ belongs to
$H^{2n-2}(X_{\bar\eta},\mathbb Q_\ell(m)\simeq\mathbb Q_\ell(-m-1)$.
Multiplication by $\sigma\in\mathbb Z_\ell(1)$ therefore gives an
element $\sigma(x\cdot\delta)\in\mathbb Q_\ell(-m)$, whose product with
$\delta\in H^{n-1}(X_{\bar\eta},\mathbb Q_\ell(m)$ lies in
$H^{n-1}(X_{\bar\eta},\mathbb Q_\ell)$.

}{\subsubsection*{6.5.1.}
If $n=2m+1$, then $I$ acts through a quotient isomorphic to
$\mathbb Z/2\mathbb Z$ (if $\operatorname{char}(k)\ne2$, there is only
one such quotient).  There is an element
\[
\delta\in H^{n-1}(X_{\bar\eta},\mathbb Q_\ell(m)),
\]
determined uniquely up to sign and called the \underline{vanishing
cycle at $s$}, such that
\begin{equation}\tag{6.5.1.1}
\delta\cdot\delta=(-1)^m2
\quad\text{in}\quad
H^{2n-2}(X_{\bar\eta},\mathbb Q_\ell(n-1))\simeq\mathbb Q_\ell.
\end{equation}
The nontrivial element $\sigma$ of $\mathbb Z/2\mathbb Z$ acts on
$H^{n-1}(X_{\bar\eta},\mathbb Q_\ell)$ by the reflection
\begin{equation}\tag{6.5.1.2}
\sigma x=x-(-1)^m(x\cdot\delta)\delta.
\end{equation}

Here $x\cdot\delta$ belongs to
$H^{2n-2}(X_{\bar\eta},\mathbb Q_\ell(m))\simeq\mathbb Q_\ell(-m)$,
so $(x\cdot\delta)\delta$ defines an element of
$H^{n-1}(X_{\bar\eta},\mathbb Q_\ell)$.

\subsubsection*{6.5.2.}
If $n=2m+2$, then $I$ acts through its canonical quotient
$\mathbb Z_\ell(1)$.  There is an element
\[
\delta\in H^{n-1}(X_{\bar\eta},\mathbb Q_\ell(m)),
\]
canonically determined up to sign and called the \underline{vanishing
cycle at $s$}, such that $\sigma\in\mathbb Z_\ell(1)$ acts on
$x\in H^{n-1}(X_{\bar\eta},\mathbb Q_\ell)$ by the transvection
\begin{equation}\tag{6.5.2.1}
\sigma(x)=x-(-1)^m\sigma\,(x\cdot\delta)\delta.
\end{equation}
Indeed, $x\cdot\delta$ belongs to
$H^{2n-2}(X_{\bar\eta},\mathbb Q_\ell(m))\simeq\mathbb Q_\ell(-m-1)$.
Multiplication by $\sigma\in\mathbb Z_\ell(1)$ therefore gives an
element $\sigma(x\cdot\delta)\in\mathbb Q_\ell(-m)$, whose product with
$\delta\in H^{n-1}(X_{\bar\eta},\mathbb Q_\ell(m))$ lies in
$H^{n-1}(X_{\bar\eta},\mathbb Q_\ell)$.

}{E7II-XVIII-037}\SGApage{II}{\(\approx\)323}
We return to the global situation.

\subsubsection*{Theorem 6.6.}
Let $X$ be a smooth connected projective variety of dimension $n$ over
an algebraically closed field $k$.  Suppose that $n$ is even or that
$\operatorname{char}(k)\ne2$.  Let $D$ be a Lefschetz pencil for $X$,
let $\eta$ be its generic point, let $\bar\eta$ be a geometric point
over $\eta$, and let $s_1,\ldots,s_N$ be the points of
$D\cap\check X$.

\SGApage{II}{\(\approx\)324}
\SGAvar{For each $i$, choose an isomorphism $\ell_i$ between $\bar\eta$ and the
geometric generic point of $S_i$, the strict henselization of $D$ at
$s_i$.  Let $I_i$ be the inertia group of $\delta_i$ at $S_i$ acting on
$H^{n-1}(X_{\bar\eta},\mathbb Q_\ell(j))$, and let $\delta_i$ be the
vanishing cycle at $s_i$ defined in 6.5.
Thus
\[
\delta_i\in H^{n-1}(X,\mathbb Q_\ell(m),
\qquad
m=\begin{cases}
(n-1)/2,& n\text{ odd},\\
n/2-1,& n\text{ even}.
\end{cases}
\]

Choose the $\ell_i$ so that the subgroups
\[
I_i\subset
\pi=\pi_1^{\mathrm{mod}}(D-D\cap\check X,\bar\eta)
\]
topologically generate $\pi_1$, as is possible by SGA 1, XIII 2.10.
Then:

\medskip
\noindent
\textup{6.6.1.} If $X$ satisfies condition (LV), the vanishing
cohomology space
\[
E^{n-1}(X_{\bar\eta},\mathbb Q_\ell(m)
\stackrel{\mathrm{def}}{=}
H^{n-1}(X,\mathbb Q_\ell(m))^\perp
\subset H^{n-1}(X_{\bar\eta},\mathbb Q_\ell(m)
\]
(cf. 5.4.4) is generated by the vanishing cycles
$\delta_i$, $i=1,\ldots,N$.

}{For each $i$, choose an isomorphism $\ell_i$ between $\bar\eta$ and the
geometric generic point of $S_i$, the strict henselization of $D$ at
$s_i$.  Via $\ell_i$, let $I_i$ be the inertia group at $s_i$ acting on
$H^{n-1}(X_{\bar\eta},\mathbb Q_\ell(j))$, and let $\delta_i$ be the
vanishing cycle at $s_i$ defined in 6.5.
Thus
\[
\delta_i\in H^{n-1}(X_{\bar\eta},\mathbb Q_\ell(m)),
\qquad
m=\begin{cases}
(n-1)/2,& n\text{ odd},\\
n/2-1,& n\text{ even}.
\end{cases}
\]

Choose the $\ell_i$ so that the subgroups
\[
I_i\subset
\pi=\pi_1^{\mathrm{mod}}(D-D\cap\check X,\bar\eta)
\]
topologically generate $\pi$, as is possible by SGA 1, XIII 2.10.
Then:

\medskip
\noindent
\textup{6.6.1.} If $X$ satisfies condition (LV), the vanishing
cohomology space
\[
E^{n-1}(X_{\bar\eta},\mathbb Q_\ell(m))
\stackrel{\mathrm{def}}{=}
H^{n-1}(X,\mathbb Q_\ell(m))^\perp
\subset H^{n-1}(X_{\bar\eta},\mathbb Q_\ell(m))
\]
(cf. 5.4.4) is generated by the vanishing cycles
$\delta_i$, $i=1,\ldots,N$.

}{E7II-XVIII-038}\medskip
\noindent
\textup{6.6.2.} The vanishing cycles $\delta_i$ are conjugate, up to
sign, under the global monodromy group
$\pi_1(D-D\cap\check X)$.

\SGAvar{\subsubsection*{Proof.}
For 6.6.1, the Picard--Lefschetz formula 6.5 and the fact that the
$I_i$ generate $\pi=\pi_1(D-D\cap\check X)$ give
\begin{equation}\tag{6.6.3}
H^{n-1}(X_{\bar\eta},\mathbb Q_\ell(m)^\pi
=\bigl\langle\delta_1,\ldots,\delta_N\bigr\rangle^\perp
\quad\text{in }H^{n-1}(X_{\bar\eta},\mathbb Q_\ell(m)).
\end{equation}
By 5.6.10,
\begin{equation}\tag{6.6.4}
H^{n-1}(X_{\bar\eta},\mathbb Q_\ell(m))^\pi
=\operatorname{Im}\bigl(
H^{n-1}(X,\mathbb Q_\ell(i))
\bigr).
\end{equation}
Taking orthogonals in (6.6.3) proves 6.6.1.

For 6.6.2, note simply that the vanishing cycle $\delta_i$ is
determined up to sign by the representation of the inertia group at
$s_i$.  Apply 6.1.2, which gives conjugacy in
$\pi_1^{\mathrm{mod}}({\mathbb P}^{r}- X,\bar\eta)$ of the
composite inertia homomorphisms
\begin{center}
\begin{tikzcd}[column sep=huge,row sep=huge]
\mathbb Z_\ell(1) \arrow[r]
  \& \pi_1^{\mathrm{mod}}(D-D\cap\check X,\bar\eta) \arrow[d] \& {}\\
{}
  \& \pi_1^{\mathrm{mod}}(\check{\mathbb P}^{r}-\check X,\bar\eta) \& .
\end{tikzcd}
\end{center}
Consequently, the vanishing cycles determined by these homomorphisms
are conjugate up to sign by elements of
$\pi_1^{\mathrm{mod}}(\check{\mathbb P}^{r}-\check X,\bar\eta)$.
The proof is completed by 6.1.1, which gives the surjection
\[
\pi_1^{\mathrm{mod}}(D-D\cap\check X)
\longrightarrow
\pi_1^{\mathrm{mod}}(\check{\mathbb P}^{r}-\check X).
\]



}{\subsubsection*{Proof.}
For 6.6.1, the Picard--Lefschetz formula 6.5 and the fact that the
$I_i$ generate $\pi=\pi_1(D-D\cap\check X)$ give
\begin{equation}\tag{6.6.3}
H^{n-1}(X_{\bar\eta},\mathbb Q_\ell(m))^\pi
=\bigl\langle\delta_1,\ldots,\delta_N\bigr\rangle^\perp
\quad\text{in }H^{n-1}(X_{\bar\eta},\mathbb Q_\ell(m)).
\end{equation}
By 5.6.10,
\begin{equation}\tag{6.6.4}
H^{n-1}(X_{\bar\eta},\mathbb Q_\ell(m))^\pi
=\operatorname{Im}\bigl(
H^{n-1}(X,\mathbb Q_\ell(m))
\longrightarrow H^{n-1}(X_{\bar\eta},\mathbb Q_\ell(m))
\bigr).
\end{equation}
Taking orthogonals in (6.6.3) proves 6.6.1.

For 6.6.2, note simply that the vanishing cycle $\delta_i$ is
determined up to sign by the representation of the inertia group at
$s_i$.  Apply 6.1.2, which gives conjugacy in
$\pi_1^{\mathrm{mod}}(\check{\mathbb P}^{r}-\check X,\bar\eta)$ of the
composite inertia homomorphisms
\begin{center}
\begin{tikzcd}[column sep=huge,row sep=huge]
\mathbb Z_\ell(1) \arrow[r]
  \& \pi_1^{\mathrm{mod}}(D-D\cap\check X,\bar\eta) \arrow[d] \& {}\\
{}
  \& \pi_1^{\mathrm{mod}}(\check{\mathbb P}^{r}-\check X,\bar\eta) \& .
\end{tikzcd}
\end{center}
Consequently, the vanishing cycles determined by these homomorphisms
are conjugate up to sign by elements of
$\pi_1^{\mathrm{mod}}(\check{\mathbb P}^{r}-\check X,\bar\eta)$.
The proof is completed by 6.1.1, which gives the surjection
\[
\pi_1^{\mathrm{mod}}(D-D\cap\check X)
\longrightarrow
\pi_1^{\mathrm{mod}}(\check{\mathbb P}^{r}-\check X).
\]



}{E7II-XVIII-039}\SGApage{II}{326}
\subsubsection*{Corollary 6.7.}
Under the hypotheses of 6.6, suppose that $X$ satisfies condition
(LV).  Then:

\medskip
\noindent
\textup{6.7.1.} The representation of
$\pi_1(D-D\cap\check X)$ on
$E^{n-1}(X_{\bar\eta},\mathbb Q_\ell)$ is absolutely irreducible.

\medskip
\noindent
\textup{6.7.2.} If $n$ is even, every finite-index subgroup of
$\pi_1(D-D\cap\check X)$ acts absolutely irreducibly on
$E^{n-1}(X_{\bar\eta},\mathbb Q_\ell)$.

\SGAvar{\subsubsection*{Proof.}
We may replace the module in question by its twist
\[
E=E^{n-1}(X_{\bar\eta},\mathbb Q_\ell(m).
\]
Let
$\bar E=E\otimes_{\mathbb Q_\ell}\overline{\mathbb Q_\ell}$, and choose
$0\ne\chi\in\bar E$.  To prove 6.7.1, we must show that the
$\overline{\mathbb Q_\ell}$-linear subspace of $\bar E$ generated by $\chi$
and its conjugates under $\pi$ is all of $\bar E$.  By 6.6.1 and 6.6.2,
it is enough to prove that this subspace contains one of the vanishing
cycles $\delta_i$.

Since $\chi$ satisfies (LV), the restriction of the cup-product form to
\[
H^{n-1}(X,\mathbb Q_\ell(m))
\lhook\joinrel\longrightarrow
H^{n-1}(X_{\bar\eta},\mathbb Q_\ell(m))
\]
is nondegenerate; hence so is its restriction to the orthogonal
$E^{n-1}(X_{\bar\eta},\mathbb Q_\ell(m))$.  By 6.6.1, there is therefore
at least one $i$ such that
\[
x\cdot\delta_i\ne0.
\]
For this $i$, 6.5 supplies an element $\sigma_i$ of the inertia group at
$s_i$ such that
\begin{equation}\tag{6.7.3}
\sigma_i(x)-x
=\text{a nonzero multiple of }\delta_i.
\end{equation}
This proves 6.7.1.

For 6.7.2, $n-1$ is odd, so $\delta\cdot\delta=0$.  Iterating
6.5.2 gives, for every $k$,
\begin{equation}\tag{6.7.4}
\sigma_i^k(x)
=X+k\,\sigma_i(x\cdot\delta_i)\delta_i.
\end{equation}
Thus a $Q_\ell$-subspace $M$ of $E$ is stable under the action
of $\pi$ if and only if it is stable under a finite-index subgroup:
both conditions are equivalent to
\[
x\in M
\quad\Longrightarrow\quad
\delta_i\in M\quad\text{for every $i$ with }
x\cdot\delta_i\ne0.
\]
The assertion therefore reduces to 6.7.1.

}{\subsubsection*{Proof.}
We may replace the module in question by its twist
\[
E=E^{n-1}(X_{\bar\eta},\mathbb Q_\ell(m)).
\]
Let
$\bar E=E\otimes_{\mathbb Q_\ell}\bar{\mathbb Q}_\ell$, and choose
$0\ne\chi\in\bar E$.  To prove 6.7.1, we must show that the
$\bar{\mathbb Q}_\ell$-linear subspace of $\bar E$ generated by $\chi$
and its conjugates under $\pi$ is all of $\bar E$.  By 6.6.1 and 6.6.2,
it is enough to prove that this subspace contains one of the vanishing
cycles $\delta_i$.

Since $X$ satisfies (LV), the restriction of the cup-product form to
\[
H^{n-1}(X,\mathbb Q_\ell(m))
\lhook\joinrel\longrightarrow
H^{n-1}(X_{\bar\eta},\mathbb Q_\ell(m))
\]
is nondegenerate; hence so is its restriction to the orthogonal
$E^{n-1}(X_{\bar\eta},\mathbb Q_\ell(m))$.  By 6.6.1, there is therefore
at least one $i$ such that
\[
\chi\cdot\delta_i\ne0.
\]
For this $i$, 6.5 supplies an element $\sigma_i$ of the inertia group at
$s_i$ such that
\begin{equation}\tag{6.7.3}
\sigma_i(\chi)-\chi
=\text{a nonzero multiple of }\delta_i.
\end{equation}
This proves 6.7.1.

For 6.7.2, $n-1$ is odd, so $\delta_i\cdot\delta_i=0$.  Iterating
6.5.2 gives, for every $k$,
\begin{equation}\tag{6.7.4}
\sigma_i^k(\chi)
=\chi+k\,\sigma_i(\chi\cdot\delta_i)\delta_i.
\end{equation}
Thus a $\mathbb Q_\ell$-subspace $M$ of $E$ is stable under the action
of $\pi$ if and only if it is stable under a finite-index subgroup:
both conditions are equivalent to
\[
\chi\in M
\quad\Longrightarrow\quad
\delta_i\in M\quad\text{for every $i$ with }
\chi\cdot\delta_i\ne0.
\]
The assertion therefore reduces to 6.7.1.

}{E7II-XVIII-040}\SGApage{II}{\(\approx\)327}
\subsection*{Bibliography}

\SGAvar{\begin{description}
\item[{[1]}] P. Deligne, \emph{Théorème de Lefschetz et critères de
dégénérescence de suites spectrales}, Publ. Math. IHÉS \textbf{35}
(1968), 259--278.

\item[{[2]}] S. Kleiman, \emph{Algebraic cycles and the Weil
conjectures}, in \emph{Dix exposés sur la cohomologie des schémas},
North-Holland, 1968.

\item[{[3]}] S. Lefschetz, \emph{L'analyse situs et la géométrie
algébrique}, Gauthier-Villars, Paris, 1924.

\item[{[4]}] D. Lieberman, \emph{Higher Picard Varieties}, Amer. J.
Math. \textbf{15} (4) (1968), 1165--1199.

\item[{[5]}] D. Mumford, unpublished.

\item[{[6]}] A. Weil, \emph{Variétés kählériennes}, Hermann, Paris,
1958.
\end{description}



}{\begin{description}
\item[{[1]}] P. Deligne, \emph{Théorème de Lefschetz et critères de
dégénérescence de suites spectrales}, Publ. Math. IHÉS \textbf{35}
(1968), 259--278.

\item[{[2]}] S. Kleiman, \emph{Algebraic cycles and the Weil
conjectures}, in \emph{Dix exposés sur la cohomologie des schémas},
North-Holland, 1968.

\item[{[3]}] S. Lefschetz, \emph{L'analyse situs et la géométrie
algébrique}, Gauthier-Villars, Paris, 1924.

\item[{[4]}] D. Lieberman, \emph{Higher Picard Varieties}, Amer. J.
Math. \textbf{90} (4) (1968), 1165--1199.

\item[{[5]}] D. Mumford, unpublished.

\item[{[6]}] A. Weil, \emph{Variétés kählériennes}, Hermann, Paris,
1958.
\end{description}



}{E7II-XVIII-041}\markboth{Exposé XVIII}{Exposé XVIII}

\ifnum\SGAform=1
\clearpage\section*{Critical apparatus}\begingroup\small\raggedright\sloppy
Readings are given in the exact TeX notation of the two English witnesses. Decimal suffixes identify the individual differences in each marked passage.

\par\medskip\noindent\hypertarget{app-E7II-XVIII-001.1}{}\textbf{E7II-XVIII-001.1}\quad printed folio 256.\par
\noindent Printed reading: \emph{(empty)}\par
\noindent Corrected reading: {\ttfamily\footnotesize )\allowbreak{}}\par
\noindent Source: {\ttfamily\footnotesize port/decisions/decis\allowbreak{}ions\_en7II\_r2\_002.js\allowbreak{}on\#C740;\allowbreak{} fixes/FIXLOG\_7II\_CUR\allowbreak{}RENT.jsonl\#C740;\allowbreak{} 7II:414\allowbreak{}}\par

\par\medskip\noindent\hypertarget{app-E7II-XVIII-002.1}{}\textbf{E7II-XVIII-002.1}\quad printed folio 257.\par
\noindent Printed reading: {\ttfamily\footnotesize n\allowbreak{}}\par
\noindent Corrected reading: {\ttfamily\footnotesize r\allowbreak{}}\par
\noindent Source: {\ttfamily\footnotesize fr/sga7ii/\{diplomati\allowbreak{}c,current\}/expose\_XV\allowbreak{}III\_body.tex;\allowbreak{} folio\allowbreak{} 257\allowbreak{}}\par

\par\medskip\noindent\hypertarget{app-E7II-XVIII-002.2}{}\textbf{E7II-XVIII-002.2}\quad printed folio 257.\par
\noindent Printed reading: {\ttfamily\footnotesize 2-j\allowbreak{}}\par
\noindent Corrected reading: {\ttfamily\footnotesize 2j\allowbreak{}}\par
\noindent Source: {\ttfamily\footnotesize fr/sga7ii/\{diplomati\allowbreak{}c,current\}/expose\_XV\allowbreak{}III\_body.tex;\allowbreak{} folio\allowbreak{} 257\allowbreak{}}\par

\par\medskip\noindent\hypertarget{app-E7II-XVIII-002.3}{}\textbf{E7II-XVIII-002.3}\quad printed folio 257.\par
\noindent Printed reading: \emph{(empty)}\par
\noindent Corrected reading: {\ttfamily\footnotesize )\textbackslash{}bigr\allowbreak{}}\par
\noindent Source: {\ttfamily\footnotesize fr/sga7ii/\{diplomati\allowbreak{}c,current\}/expose\_XV\allowbreak{}III\_body.tex;\allowbreak{} folio\allowbreak{} 257\allowbreak{}}\par

\par\medskip\noindent\hypertarget{app-E7II-XVIII-002.4}{}\textbf{E7II-XVIII-002.4}\quad printed folio 258.\par
\noindent Printed reading: {\ttfamily\footnotesize \textbackslash{}leq\allowbreak{}}\par
\noindent Corrected reading: {\ttfamily\footnotesize \textbackslash{}geq\allowbreak{}}\par
\noindent Source: {\ttfamily\footnotesize fr/sga7ii/\{diplomati\allowbreak{}c,current\}/expose\_XV\allowbreak{}III\_body.tex;\allowbreak{} folio\allowbreak{} 258;\allowbreak{} formula\allowbreak{} 46;\allowbreak{} port/PORT\_EN7II\_LOG.\allowbreak{}jsonl\#C742;\allowbreak{} fixes/FIXLOG\_7II\_CUR\allowbreak{}RENT.jsonl\#C742;\allowbreak{} 7II:419\allowbreak{}}\par

\par\medskip\noindent\hypertarget{app-E7II-XVIII-003.1}{}\textbf{E7II-XVIII-003.1}\quad printed folio 259.\par
\noindent Printed reading: {\ttfamily\footnotesize o\allowbreak{}}\par
\noindent Corrected reading: {\ttfamily\footnotesize 0\allowbreak{}}\par
\noindent Source: {\ttfamily\footnotesize fr/sga7ii/\{diplomati\allowbreak{}c,current\}/expose\_XV\allowbreak{}III\_body.tex;\allowbreak{} folio\allowbreak{} 259;\allowbreak{} formula\allowbreak{} 78\allowbreak{}}\par

\par\medskip\noindent\hypertarget{app-E7II-XVIII-003.2}{}\textbf{E7II-XVIII-003.2}\quad printed folio 259.\par
\noindent Printed reading: {\ttfamily\footnotesize o\allowbreak{}}\par
\noindent Corrected reading: {\ttfamily\footnotesize 0\allowbreak{}}\par
\noindent Source: {\ttfamily\footnotesize fr/sga7ii/\{diplomati\allowbreak{}c,current\}/expose\_XV\allowbreak{}III\_body.tex;\allowbreak{} folio\allowbreak{} 259;\allowbreak{} formula\allowbreak{} 82\allowbreak{}}\par

\par\medskip\noindent\hypertarget{app-E7II-XVIII-003.3}{}\textbf{E7II-XVIII-003.3}\quad printed folio 259.\par
\noindent Printed reading: \emph{(empty)}\par
\noindent Corrected reading: {\ttfamily\footnotesize )\textbackslash{}bigr\allowbreak{}}\par
\noindent Source: {\ttfamily\footnotesize fr/sga7ii/\{diplomati\allowbreak{}c,current\}/expose\_XV\allowbreak{}III\_body.tex;\allowbreak{} folio\allowbreak{} 259;\allowbreak{} formula\allowbreak{} 82\allowbreak{}}\par

\par\medskip\noindent\hypertarget{app-E7II-XVIII-003.4}{}\textbf{E7II-XVIII-003.4}\quad printed folio 259.\par
\noindent Printed reading: {\ttfamily\footnotesize o\allowbreak{}}\par
\noindent Corrected reading: {\ttfamily\footnotesize 0\allowbreak{}}\par
\noindent Source: {\ttfamily\footnotesize ED-EN-FORMULA-G147;\allowbreak{} fr/sga7ii/\{diplomati\allowbreak{}c,current\}/expose\_XV\allowbreak{}III\_body.tex;\allowbreak{} folio\allowbreak{} 259\allowbreak{}}\par

\par\medskip\noindent\hypertarget{app-E7II-XVIII-003.5}{}\textbf{E7II-XVIII-003.5}\quad printed folio 260.\par
\noindent Printed reading: {\ttfamily\footnotesize o\allowbreak{}}\par
\noindent Corrected reading: {\ttfamily\footnotesize 0\allowbreak{}}\par
\noindent Source: {\ttfamily\footnotesize fr/sga7ii/\{diplomati\allowbreak{}c,current\}/expose\_XV\allowbreak{}III\_body.tex;\allowbreak{} folio\allowbreak{} 260;\allowbreak{} formula\allowbreak{} 100\allowbreak{}}\par

\par\medskip\noindent\hypertarget{app-E7II-XVIII-003.6}{}\textbf{E7II-XVIII-003.6}\quad printed folio 260.\par
\noindent Printed reading: {\ttfamily\footnotesize i\allowbreak{}}\par
\noindent Corrected reading: {\ttfamily\footnotesize 1\allowbreak{}}\par
\noindent Source: {\ttfamily\footnotesize fr/sga7ii/\{diplomati\allowbreak{}c,current\}/expose\_XV\allowbreak{}III\_body.tex;\allowbreak{} folio\allowbreak{} 260\allowbreak{}}\par

\par\medskip\noindent\hypertarget{app-E7II-XVIII-004.1}{}\textbf{E7II-XVIII-004.1}\quad printed folio 260.\par
\noindent Printed reading: \emph{(empty)}\par
\noindent Corrected reading: {\ttfamily\footnotesize \textbackslash{}bigr)\allowbreak{}}\par
\noindent Source: {\ttfamily\footnotesize ED-EN-FORMULA-G149;\allowbreak{} fr/sga7ii/\{diplomati\allowbreak{}c,current\}/expose\_XV\allowbreak{}III\_body.tex;\allowbreak{} folio\allowbreak{} 260\allowbreak{}}\par

\par\medskip\noindent\hypertarget{app-E7II-XVIII-004.2}{}\textbf{E7II-XVIII-004.2}\quad printed folio 260.\par
\noindent Printed reading: {\ttfamily\footnotesize \_o\textbackslash{};\allowbreak{}}\par
\noindent Corrected reading: {\ttfamily\footnotesize ,\allowbreak{}}\par
\noindent Source: {\ttfamily\footnotesize ED-EN-FORMULA-G149;\allowbreak{} fr/sga7ii/\{diplomati\allowbreak{}c,current\}/expose\_XV\allowbreak{}III\_body.tex;\allowbreak{} folio\allowbreak{} 260\allowbreak{}}\par

\par\medskip\noindent\hypertarget{app-E7II-XVIII-005.1}{}\textbf{E7II-XVIII-005.1}\quad printed folio 263.\par
\noindent Printed reading: {\ttfamily\footnotesize R\allowbreak{}}\par
\noindent Corrected reading: {\ttfamily\footnotesize E\allowbreak{}}\par
\noindent Source: {\ttfamily\footnotesize fr/sga7ii/\{diplomati\allowbreak{}c,current\}/expose\_XV\allowbreak{}III\_body.tex;\allowbreak{} folio\allowbreak{} 263\allowbreak{}}\par

\par\medskip\noindent\hypertarget{app-E7II-XVIII-005.2}{}\textbf{E7II-XVIII-005.2}\quad printed folio 263.\par
\noindent Printed reading: \emph{(empty)}\par
\noindent Corrected reading: {\ttfamily\footnotesize \{\allowbreak{}}\par
\noindent Source: {\ttfamily\footnotesize ED-EN-FORMULA-G151;\allowbreak{} fr/sga7ii/\{diplomati\allowbreak{}c,current\}/expose\_XV\allowbreak{}III\_body.tex;\allowbreak{} folio\allowbreak{} 263\allowbreak{}}\par

\par\medskip\noindent\hypertarget{app-E7II-XVIII-005.3}{}\textbf{E7II-XVIII-005.3}\quad printed folio 263.\par
\noindent Printed reading: \emph{(empty)}\par
\noindent Corrected reading: {\ttfamily\footnotesize +1\}\allowbreak{}}\par
\noindent Source: {\ttfamily\footnotesize ED-EN-FORMULA-G151;\allowbreak{} fr/sga7ii/\{diplomati\allowbreak{}c,current\}/expose\_XV\allowbreak{}III\_body.tex;\allowbreak{} folio\allowbreak{} 263\allowbreak{}}\par

\par\medskip\noindent\hypertarget{app-E7II-XVIII-006.1}{}\textbf{E7II-XVIII-006.1}\quad printed folio 265.\par
\noindent Printed reading: {\ttfamily\footnotesize H\allowbreak{}}\par
\noindent Corrected reading: {\ttfamily\footnotesize F\allowbreak{}}\par
\noindent Source: {\ttfamily\footnotesize fr/sga7ii/\{diplomati\allowbreak{}c,current\}/expose\_XV\allowbreak{}III\_body.tex;\allowbreak{} folio\allowbreak{} 265\allowbreak{}}\par

\par\medskip\noindent\hypertarget{app-E7II-XVIII-006.2}{}\textbf{E7II-XVIII-006.2}\quad printed folio 265.\par
\noindent Printed reading: {\ttfamily\footnotesize pF\allowbreak{}}\par
\noindent Corrected reading: {\ttfamily\footnotesize pH\allowbreak{}}\par
\noindent Source: {\ttfamily\footnotesize fr/sga7ii/\{diplomati\allowbreak{}c,current\}/expose\_XV\allowbreak{}III\_body.tex;\allowbreak{} folio\allowbreak{} 265\allowbreak{}}\par

\par\medskip\noindent\hypertarget{app-E7II-XVIII-006.3}{}\textbf{E7II-XVIII-006.3}\quad printed folio 265.\par
\noindent Printed reading: \emph{(empty)}\par
\noindent Corrected reading: {\ttfamily\footnotesize \textbackslash{}widetilde\allowbreak{}}\par
\noindent Source: {\ttfamily\footnotesize fr/sga7ii/\{diplomati\allowbreak{}c,current\}/expose\_XV\allowbreak{}III\_body.tex;\allowbreak{} folio\allowbreak{} 265;\allowbreak{} formula\allowbreak{} 246\allowbreak{}}\par

\par\medskip\noindent\hypertarget{app-E7II-XVIII-006.4}{}\textbf{E7II-XVIII-006.4}\quad printed folio 266.\par
\noindent Printed reading: {\ttfamily\footnotesize C\allowbreak{}}\par
\noindent Corrected reading: {\ttfamily\footnotesize c\allowbreak{}}\par
\noindent Source: {\ttfamily\footnotesize ED-EN-FORMULA-G153;\allowbreak{} fr/sga7ii/\{diplomati\allowbreak{}c,current\}/expose\_XV\allowbreak{}III\_body.tex;\allowbreak{} folio\allowbreak{} 266\allowbreak{}}\par

\par\medskip\noindent\hypertarget{app-E7II-XVIII-006.5}{}\textbf{E7II-XVIII-006.5}\quad printed folio 266.\par
\noindent Printed reading: {\ttfamily\footnotesize \textbackslash{}chi\allowbreak{}}\par
\noindent Corrected reading: {\ttfamily\footnotesize X\allowbreak{}}\par
\noindent Source: {\ttfamily\footnotesize fr/sga7ii/\{diplomati\allowbreak{}c,current\}/expose\_XV\allowbreak{}III\_body.tex;\allowbreak{} folio\allowbreak{} 266\allowbreak{}}\par

\par\medskip\noindent\hypertarget{app-E7II-XVIII-006.6}{}\textbf{E7II-XVIII-006.6}\quad printed folio 266.\par
\noindent Printed reading: {\ttfamily\footnotesize j\allowbreak{}}\par
\noindent Corrected reading: {\ttfamily\footnotesize 2j\allowbreak{}}\par
\noindent Source: {\ttfamily\footnotesize ED-EN-FORMULA-G155;\allowbreak{} fr/sga7ii/\{diplomati\allowbreak{}c,current\}/expose\_XV\allowbreak{}III\_body.tex;\allowbreak{} folio\allowbreak{} 266\allowbreak{}}\par

\par\medskip\noindent\hypertarget{app-E7II-XVIII-007.1}{}\textbf{E7II-XVIII-007.1}\quad printed folio 267.\par
\noindent Printed reading: {\ttfamily\footnotesize u\allowbreak{}}\par
\noindent Corrected reading: {\ttfamily\footnotesize n\allowbreak{}}\par
\noindent Source: {\ttfamily\footnotesize fr/sga7ii/\{diplomati\allowbreak{}c,current\}/expose\_XV\allowbreak{}III\_body.tex;\allowbreak{} folio\allowbreak{} 267;\allowbreak{} formula\allowbreak{} 254\allowbreak{}}\par

\par\medskip\noindent\hypertarget{app-E7II-XVIII-007.2}{}\textbf{E7II-XVIII-007.2}\quad printed folio 267.\par
\noindent Printed reading: \emph{(empty)}\par
\noindent Corrected reading: {\ttfamily\footnotesize )\allowbreak{}}\par
\noindent Source: {\ttfamily\footnotesize fr/sga7ii/\{diplomati\allowbreak{}c,current\}/expose\_XV\allowbreak{}III\_body.tex;\allowbreak{} folio\allowbreak{} 267\allowbreak{}}\par

\par\medskip\noindent\hypertarget{app-E7II-XVIII-008.1}{}\textbf{E7II-XVIII-008.1}\quad printed folio 268.\par
\noindent Printed reading: {\ttfamily\footnotesize S\allowbreak{}}\par
\noindent Corrected reading: {\ttfamily\footnotesize \textbackslash{}rho\allowbreak{}}\par
\noindent Source: {\ttfamily\footnotesize fr/sga7ii/\{diplomati\allowbreak{}c,current\}/expose\_XV\allowbreak{}III\_body.tex;\allowbreak{} folio\allowbreak{} 268;\allowbreak{} formula\allowbreak{} 292\allowbreak{}}\par

\par\medskip\noindent\hypertarget{app-E7II-XVIII-008.2}{}\textbf{E7II-XVIII-008.2}\quad printed folio 268.\par
\noindent Printed reading: {\ttfamily\footnotesize \{\textbackslash{}eta\}\allowbreak{}}\par
\noindent Corrected reading: {\ttfamily\footnotesize n\allowbreak{}}\par
\noindent Source: {\ttfamily\footnotesize ED-EN-FORMULA-G157;\allowbreak{} fr/sga7ii/\{diplomati\allowbreak{}c,current\}/expose\_XV\allowbreak{}III\_body.tex;\allowbreak{} folio\allowbreak{} 268\allowbreak{}}\par

\par\medskip\noindent\hypertarget{app-E7II-XVIII-008.3}{}\textbf{E7II-XVIII-008.3}\quad printed folio 268.\par
\noindent Printed reading: {\ttfamily\footnotesize \textbackslash{}eta\allowbreak{}}\par
\noindent Corrected reading: {\ttfamily\footnotesize n\allowbreak{}}\par
\noindent Source: {\ttfamily\footnotesize fr/sga7ii/\{diplomati\allowbreak{}c,current\}/expose\_XV\allowbreak{}III\_body.tex;\allowbreak{} folio\allowbreak{} 268\allowbreak{}}\par

\par\medskip\noindent\hypertarget{app-E7II-XVIII-008.4}{}\textbf{E7II-XVIII-008.4}\quad printed folio 269.\par
\noindent Printed reading: \emph{(empty)}\par
\noindent Corrected reading: {\ttfamily\footnotesize \textasciicircum{}3\allowbreak{}}\par
\noindent Source: {\ttfamily\footnotesize ED-EN-NONMATH-G005;\allowbreak{} F7II-XVIII-019.4;\allowbreak{} ADOPT-F7II-XVIII-019\allowbreak{}.4;\allowbreak{} fr/sga7ii/\{diplomati\allowbreak{}c,current\}/expose\_XV\allowbreak{}III\_body.tex;\allowbreak{} folio\allowbreak{} 269;\allowbreak{} fr/sga7ii/\{diplomati\allowbreak{}c,current\}/expose\_XV\allowbreak{}III\_body.tex;\allowbreak{} folio\allowbreak{} 269;\allowbreak{} formula\allowbreak{} 308\allowbreak{}}\par

\par\medskip\noindent\hypertarget{app-E7II-XVIII-008.5}{}\textbf{E7II-XVIII-008.5}\quad printed folio 269.\par
\noindent Printed reading: \emph{(empty)}\par
\noindent Corrected reading: {\ttfamily\footnotesize -1\allowbreak{}}\par
\noindent Source: {\ttfamily\footnotesize ED-EN-NONMATH-G005;\allowbreak{} F7II-XVIII-019.4;\allowbreak{} ADOPT-F7II-XVIII-019\allowbreak{}.4;\allowbreak{} fr/sga7ii/\{diplomati\allowbreak{}c,current\}/expose\_XV\allowbreak{}III\_body.tex;\allowbreak{} folio\allowbreak{} 269;\allowbreak{} fr/sga7ii/\{diplomati\allowbreak{}c,current\}/expose\_XV\allowbreak{}III\_body.tex;\allowbreak{} folio\allowbreak{} 269;\allowbreak{} formula\allowbreak{} 308\allowbreak{}}\par

\par\medskip\noindent\hypertarget{app-E7II-XVIII-008.6}{}\textbf{E7II-XVIII-008.6}\quad printed folio 269.\par
\noindent Printed reading: {\ttfamily\footnotesize  a\allowbreak{} cube\allowbreak{} root\allowbreak{} of\allowbreak{} unity\allowbreak{}}\par
\noindent Corrected reading: \emph{(empty)}\par
\noindent Source: {\ttfamily\footnotesize ED-EN-NONMATH-G005;\allowbreak{} F7II-XVIII-019.4;\allowbreak{} ADOPT-F7II-XVIII-019\allowbreak{}.4;\allowbreak{} fr/sga7ii/\{diplomati\allowbreak{}c,current\}/expose\_XV\allowbreak{}III\_body.tex;\allowbreak{} folio\allowbreak{} 269\allowbreak{}}\par

\par\medskip\noindent\hypertarget{app-E7II-XVIII-008.7}{}\textbf{E7II-XVIII-008.7}\quad printed folio 270.\par
\noindent Printed reading: {\ttfamily\footnotesize \textbackslash{}eta\allowbreak{}}\par
\noindent Corrected reading: {\ttfamily\footnotesize n\allowbreak{}}\par
\noindent Source: {\ttfamily\footnotesize fr/sga7ii/\{diplomati\allowbreak{}c,current\}/expose\_XV\allowbreak{}III\_body.tex;\allowbreak{} folio\allowbreak{} 270;\allowbreak{} formula\allowbreak{} 316\allowbreak{}}\par

\par\medskip\noindent\hypertarget{app-E7II-XVIII-009.1}{}\textbf{E7II-XVIII-009.1}\quad printed folio 270.\par
\noindent Printed reading: {\ttfamily\footnotesize i\allowbreak{}}\par
\noindent Corrected reading: {\ttfamily\footnotesize 1\allowbreak{}}\par
\noindent Source: {\ttfamily\footnotesize fr/sga7ii/\{diplomati\allowbreak{}c,current\}/expose\_XV\allowbreak{}III\_body.tex;\allowbreak{} folio\allowbreak{} 270;\allowbreak{} formula\allowbreak{} 328\allowbreak{}}\par

\par\medskip\noindent\hypertarget{app-E7II-XVIII-009.2}{}\textbf{E7II-XVIII-009.2}\quad printed folio 270.\par
\noindent Printed reading: {\ttfamily\footnotesize rank\allowbreak{}}\par
\noindent Corrected reading: {\ttfamily\footnotesize length\allowbreak{}}\par
\noindent Source: {\ttfamily\footnotesize ED-EN-FORMULA-G159;\allowbreak{} fr/sga7ii/\{diplomati\allowbreak{}c,current\}/expose\_XV\allowbreak{}III\_body.tex;\allowbreak{} folio\allowbreak{} 270\allowbreak{}}\par

\par\medskip\noindent\hypertarget{app-E7II-XVIII-009.3}{}\textbf{E7II-XVIII-009.3}\quad printed folio 270.\par
\noindent Printed reading: \emph{(empty)}\par
\noindent Corrected reading: {\ttfamily\footnotesize \_\{i=1\}\allowbreak{}}\par
\noindent Source: {\ttfamily\footnotesize fr/sga7ii/\{diplomati\allowbreak{}c,current\}/expose\_XV\allowbreak{}III\_body.tex;\allowbreak{} folio\allowbreak{} 270\allowbreak{}}\par

\par\medskip\noindent\hypertarget{app-E7II-XVIII-010.1}{}\textbf{E7II-XVIII-010.1}\quad printed folio 272.\par
\noindent Printed reading: \emph{(empty)}\par
\noindent Corrected reading: {\ttfamily\footnotesize )\allowbreak{}}\par
\noindent Source: {\ttfamily\footnotesize fr/sga7ii/\{diplomati\allowbreak{}c,current\}/expose\_XV\allowbreak{}III\_body.tex;\allowbreak{} folio\allowbreak{} 272\allowbreak{}}\par

\par\medskip\noindent\hypertarget{app-E7II-XVIII-010.2}{}\textbf{E7II-XVIII-010.2}\quad printed folio 272.\par
\noindent Printed reading: \emph{(empty)}\par
\noindent Corrected reading: {\ttfamily\footnotesize )\allowbreak{}}\par
\noindent Source: {\ttfamily\footnotesize ED-EN-FORMULA-G163;\allowbreak{} fr/sga7ii/\{diplomati\allowbreak{}c,current\}/expose\_XV\allowbreak{}III\_body.tex;\allowbreak{} folio\allowbreak{} 272\allowbreak{}}\par

\par\medskip\noindent\hypertarget{app-E7II-XVIII-010.3}{}\textbf{E7II-XVIII-010.3}\quad printed folio 272.\par
\noindent Printed reading: \emph{(empty)}\par
\noindent Corrected reading: {\ttfamily\footnotesize )\allowbreak{}}\par
\noindent Source: {\ttfamily\footnotesize fr/sga7ii/\{diplomati\allowbreak{}c,current\}/expose\_XV\allowbreak{}III\_body.tex;\allowbreak{} folio\allowbreak{} 272;\allowbreak{} formula\allowbreak{} 384\allowbreak{}}\par

\par\medskip\noindent\hypertarget{app-E7II-XVIII-010.4}{}\textbf{E7II-XVIII-010.4}\quad printed folio 273.\par
\noindent Printed reading: {\ttfamily\footnotesize leftrightarrow,\allowbreak{}}\par
\noindent Corrected reading: \emph{(empty)}\par
\noindent Source: {\ttfamily\footnotesize fr/sga7ii/\{diplomati\allowbreak{}c,current\}/expose\_XV\allowbreak{}III\_body.tex;\allowbreak{} folio\allowbreak{} 273\allowbreak{}}\par

\par\medskip\noindent\hypertarget{app-E7II-XVIII-010.5}{}\textbf{E7II-XVIII-010.5}\quad printed folio 274.\par
\noindent Printed reading: \emph{(1 espace(s), 0 retour(s) de ligne dans le témoin TeX)}\par
\noindent Corrected reading: {\ttfamily\footnotesize (-\allowbreak{}}\par
\noindent Source: {\ttfamily\footnotesize ED-EN-FORMULA-G165;\allowbreak{} fr/sga7ii/\{diplomati\allowbreak{}c,current\}/expose\_XV\allowbreak{}III\_body.tex;\allowbreak{} folio\allowbreak{} 274\allowbreak{}}\par

\par\medskip\noindent\hypertarget{app-E7II-XVIII-010.6}{}\textbf{E7II-XVIII-010.6}\quad printed folio 274.\par
\noindent Printed reading: \emph{(empty)}\par
\noindent Corrected reading: {\ttfamily\footnotesize )\allowbreak{}}\par
\noindent Source: {\ttfamily\footnotesize ED-EN-FORMULA-G165;\allowbreak{} fr/sga7ii/\{diplomati\allowbreak{}c,current\}/expose\_XV\allowbreak{}III\_body.tex;\allowbreak{} folio\allowbreak{} 274\allowbreak{}}\par

\par\medskip\noindent\hypertarget{app-E7II-XVIII-010.7}{}\textbf{E7II-XVIII-010.7}\quad printed folio 274.\par
\noindent Printed reading: {\ttfamily\footnotesize 2\allowbreak{}}\par
\noindent Corrected reading: {\ttfamily\footnotesize 4\allowbreak{}}\par
\noindent Source: {\ttfamily\footnotesize fr/sga7ii/\{diplomati\allowbreak{}c,current\}/expose\_XV\allowbreak{}III\_body.tex;\allowbreak{} folio\allowbreak{} 274;\allowbreak{} formula\allowbreak{} 394;\allowbreak{} port/PORT\_EN7II\_LOG.\allowbreak{}jsonl\#C746;\allowbreak{} fixes/FIXLOG\_7II\_CUR\allowbreak{}RENT.jsonl\#C746;\allowbreak{} 7II:450\allowbreak{}}\par

\par\medskip\noindent\hypertarget{app-E7II-XVIII-010.8}{}\textbf{E7II-XVIII-010.8}\quad printed folio 274.\par
\noindent Printed reading: {\ttfamily\footnotesize \_\allowbreak{}}\par
\noindent Corrected reading: {\ttfamily\footnotesize \textasciicircum{}\allowbreak{}}\par
\noindent Source: {\ttfamily\footnotesize fr/sga7ii/\{diplomati\allowbreak{}c,current\}/expose\_XV\allowbreak{}III\_body.tex;\allowbreak{} folio\allowbreak{} 274;\allowbreak{} formula\allowbreak{} 394;\allowbreak{} port/PORT\_EN7II\_LOG.\allowbreak{}jsonl\#C747;\allowbreak{} fixes/FIXLOG\_7II\_CUR\allowbreak{}RENT.jsonl\#C747;\allowbreak{} 7II:451\allowbreak{}}\par

\par\medskip\noindent\hypertarget{app-E7II-XVIII-010.9}{}\textbf{E7II-XVIII-010.9}\quad printed folio 274.\par
\noindent Printed reading: \emph{(empty)}\par
\noindent Corrected reading: {\ttfamily\footnotesize [\textbackslash{}sim]\allowbreak{}}\par
\noindent Source: {\ttfamily\footnotesize fr/sga7ii/\{diplomati\allowbreak{}c,current\}/expose\_XV\allowbreak{}III\_body.tex;\allowbreak{} folio\allowbreak{} 274\allowbreak{}}\par

\par\medskip\noindent\hypertarget{app-E7II-XVIII-010.10}{}\textbf{E7II-XVIII-010.10}\quad printed folio 274.\par
\noindent Printed reading: \emph{(empty)}\par
\noindent Corrected reading: {\ttfamily\footnotesize -\allowbreak{}}\par
\noindent Source: {\ttfamily\footnotesize ED-EN-FORMULA-G167;\allowbreak{} fr/sga7ii/\{diplomati\allowbreak{}c,current\}/expose\_XV\allowbreak{}III\_body.tex;\allowbreak{} folio\allowbreak{} 274\allowbreak{}}\par

\par\medskip\noindent\hypertarget{app-E7II-XVIII-010.11}{}\textbf{E7II-XVIII-010.11}\quad printed folio 274.\par
\noindent Printed reading: \emph{(empty)}\par
\noindent Corrected reading: {\ttfamily\footnotesize -\allowbreak{}}\par
\noindent Source: {\ttfamily\footnotesize fr/sga7ii/\{diplomati\allowbreak{}c,current\}/expose\_XV\allowbreak{}III\_body.tex;\allowbreak{} folio\allowbreak{} 274\allowbreak{}}\par

\par\medskip\noindent\hypertarget{app-E7II-XVIII-010.12}{}\textbf{E7II-XVIII-010.12}\quad printed folio 274.\par
\noindent Printed reading: {\ttfamily\footnotesize X\allowbreak{}}\par
\noindent Corrected reading: {\ttfamily\footnotesize x\allowbreak{}}\par
\noindent Source: {\ttfamily\footnotesize fr/sga7ii/\{diplomati\allowbreak{}c,current\}/expose\_XV\allowbreak{}III\_body.tex;\allowbreak{} folio\allowbreak{} 274\allowbreak{}}\par

\par\medskip\noindent\hypertarget{app-E7II-XVIII-010.13}{}\textbf{E7II-XVIII-010.13}\quad printed folio 274.\par
\noindent Printed reading: {\ttfamily\footnotesize X\allowbreak{}}\par
\noindent Corrected reading: {\ttfamily\footnotesize x\allowbreak{}}\par
\noindent Source: {\ttfamily\footnotesize fr/sga7ii/\{diplomati\allowbreak{}c,current\}/expose\_XV\allowbreak{}III\_body.tex;\allowbreak{} folio\allowbreak{} 274\allowbreak{}}\par

\par\medskip\noindent\hypertarget{app-E7II-XVIII-010.14}{}\textbf{E7II-XVIII-010.14}\quad printed folio 274.\par
\noindent Printed reading: {\ttfamily\footnotesize X\allowbreak{}}\par
\noindent Corrected reading: {\ttfamily\footnotesize x\allowbreak{}}\par
\noindent Source: {\ttfamily\footnotesize fr/sga7ii/\{diplomati\allowbreak{}c,current\}/expose\_XV\allowbreak{}III\_body.tex;\allowbreak{} folio\allowbreak{} 274\allowbreak{}}\par

\par\medskip\noindent\hypertarget{app-E7II-XVIII-010.15}{}\textbf{E7II-XVIII-010.15}\quad printed folio 275.\par
\noindent Printed reading: {\ttfamily\footnotesize \textbackslash{}underline\{\textbackslash{}mathcal\allowbreak{}}\par
\noindent Corrected reading: {\ttfamily\footnotesize (\textbackslash{}rho\allowbreak{}}\par
\noindent Source: {\ttfamily\footnotesize ED-EN-FORMULA-G169;\allowbreak{} fr/sga7ii/\{diplomati\allowbreak{}c,current\}/expose\_XV\allowbreak{}III\_body.tex;\allowbreak{} folio\allowbreak{} 275\allowbreak{}}\par

\par\medskip\noindent\hypertarget{app-E7II-XVIII-010.16}{}\textbf{E7II-XVIII-010.16}\quad printed folio 275.\par
\noindent Printed reading: {\ttfamily\footnotesize O\}\_\{\textbackslash{}mathbb\allowbreak{} P(\textbackslash{}underline\{\textbackslash{}mathca\allowbreak{}l\allowbreak{} O\}\_\textbackslash{}Delta\allowbreak{}}\par
\noindent Corrected reading: {\ttfamily\footnotesize i)\allowbreak{}}\par
\noindent Source: {\ttfamily\footnotesize ED-EN-FORMULA-G169;\allowbreak{} fr/sga7ii/\{diplomati\allowbreak{}c,current\}/expose\_XV\allowbreak{}III\_body.tex;\allowbreak{} folio\allowbreak{} 275\allowbreak{}}\par

\par\medskip\noindent\hypertarget{app-E7II-XVIII-010.17}{}\textbf{E7II-XVIII-010.17}\quad printed folio 275.\par
\noindent Printed reading: {\ttfamily\footnotesize \{,2\})\}(1)\allowbreak{} \textbackslash{}otimes\_\{\textbackslash{}underline\{\allowbreak{}\textbackslash{}mathcal\allowbreak{} O\}\_\textbackslash{}Delta\textasciicircum{}\{,2\}\}\textbackslash{}unde\allowbreak{}rline\{\textbackslash{}mathcal\allowbreak{} O\}\_\textbackslash{}Delta(-1)\allowbreak{} =\allowbreak{}}\par
\noindent Corrected reading: {\ttfamily\footnotesize *\allowbreak{}}\par
\noindent Source: {\ttfamily\footnotesize ED-EN-FORMULA-G169;\allowbreak{} fr/sga7ii/\{diplomati\allowbreak{}c,current\}/expose\_XV\allowbreak{}III\_body.tex;\allowbreak{} folio\allowbreak{} 275\allowbreak{}}\par

\par\medskip\noindent\hypertarget{app-E7II-XVIII-010.18}{}\textbf{E7II-XVIII-010.18}\quad printed folio 275.\par
\noindent Printed reading: {\ttfamily\footnotesize \_K\allowbreak{}}\par
\noindent Corrected reading: {\ttfamily\footnotesize  g\textasciicircum{}*\allowbreak{}}\par
\noindent Source: {\ttfamily\footnotesize ED-EN-FORMULA-G169;\allowbreak{} fr/sga7ii/\{diplomati\allowbreak{}c,current\}/expose\_XV\allowbreak{}III\_body.tex;\allowbreak{} folio\allowbreak{} 275\allowbreak{}}\par

\par\medskip\noindent\hypertarget{app-E7II-XVIII-010.19}{}\textbf{E7II-XVIII-010.19}\quad printed folio 275.\par
\noindent Printed reading: \emph{(empty)}\par
\noindent Corrected reading: {\ttfamily\footnotesize \textasciicircum{}1\allowbreak{}}\par
\noindent Source: {\ttfamily\footnotesize fr/sga7ii/\{diplomati\allowbreak{}c,current\}/expose\_XV\allowbreak{}III\_body.tex;\allowbreak{} folio\allowbreak{} 275\allowbreak{}}\par

\par\medskip\noindent\hypertarget{app-E7II-XVIII-010.20}{}\textbf{E7II-XVIII-010.20}\quad printed folio 276.\par
\noindent Printed reading: \emph{(empty)}\par
\noindent Corrected reading: {\ttfamily\footnotesize )\allowbreak{}}\par
\noindent Source: {\ttfamily\footnotesize ED-EN-FORMULA-G171;\allowbreak{} fr/sga7ii/\{diplomati\allowbreak{}c,current\}/expose\_XV\allowbreak{}III\_body.tex;\allowbreak{} folio\allowbreak{} 276\allowbreak{}}\par

\par\medskip\noindent\hypertarget{app-E7II-XVIII-010.21}{}\textbf{E7II-XVIII-010.21}\quad printed folio 276.\par
\noindent Printed reading: {\ttfamily\footnotesize .\allowbreak{}}\par
\noindent Corrected reading: {\ttfamily\footnotesize )\allowbreak{}}\par
\noindent Source: {\ttfamily\footnotesize fr/sga7ii/\{diplomati\allowbreak{}c,current\}/expose\_XV\allowbreak{}III\_body.tex;\allowbreak{} folio\allowbreak{} 276\allowbreak{}}\par

\par\medskip\noindent\hypertarget{app-E7II-XVIII-010.22}{}\textbf{E7II-XVIII-010.22}\quad printed folio 276.\par
\noindent Printed reading: {\ttfamily\footnotesize \textbackslash{}Lambda\allowbreak{}}\par
\noindent Corrected reading: {\ttfamily\footnotesize \textbackslash{}wedge\allowbreak{}}\par
\noindent Source: {\ttfamily\footnotesize ED-EN-FORMULA-G173;\allowbreak{} fr/sga7ii/\{diplomati\allowbreak{}c,current\}/expose\_XV\allowbreak{}III\_body.tex;\allowbreak{} folio\allowbreak{} 276\allowbreak{}}\par

\par\medskip\noindent\hypertarget{app-E7II-XVIII-010.23}{}\textbf{E7II-XVIII-010.23}\quad printed folio 276.\par
\noindent Printed reading: {\ttfamily\footnotesize .\allowbreak{}}\par
\noindent Corrected reading: {\ttfamily\footnotesize )\allowbreak{}}\par
\noindent Source: {\ttfamily\footnotesize fr/sga7ii/\{diplomati\allowbreak{}c,current\}/expose\_XV\allowbreak{}III\_body.tex;\allowbreak{} folio\allowbreak{} 276\allowbreak{}}\par

\par\medskip\noindent\hypertarget{app-E7II-XVIII-010.24}{}\textbf{E7II-XVIII-010.24}\quad printed folio 277.\par
\noindent Printed reading: {\ttfamily\footnotesize \textbackslash{}qquad\textbackslash{}text\{by\allowbreak{} (4.3.2)\}\allowbreak{}}\par
\noindent Corrected reading: \emph{(empty)}\par
\noindent Source: {\ttfamily\footnotesize ED-EN-FORMULA-G175;\allowbreak{} fr/sga7ii/\{diplomati\allowbreak{}c,current\}/expose\_XV\allowbreak{}III\_body.tex;\allowbreak{} folio\allowbreak{} 277\allowbreak{}}\par

\par\medskip\noindent\hypertarget{app-E7II-XVIII-010.25}{}\textbf{E7II-XVIII-010.25}\quad printed folio 277.\par
\noindent Printed reading: {\ttfamily\footnotesize 3\allowbreak{}}\par
\noindent Corrected reading: {\ttfamily\footnotesize 2\allowbreak{}}\par
\noindent Source: {\ttfamily\footnotesize ED-EN-FORMULA-G175;\allowbreak{} fr/sga7ii/\{diplomati\allowbreak{}c,current\}/expose\_XV\allowbreak{}III\_body.tex;\allowbreak{} folio\allowbreak{} 277\allowbreak{}}\par

\par\medskip\noindent\hypertarget{app-E7II-XVIII-010.26}{}\textbf{E7II-XVIII-010.26}\quad printed folio 277.\par
\noindent Printed reading: {\ttfamily\footnotesize 2\allowbreak{}}\par
\noindent Corrected reading: {\ttfamily\footnotesize 11\allowbreak{}}\par
\noindent Source: {\ttfamily\footnotesize ED-EN-FORMULA-G175;\allowbreak{} fr/sga7ii/\{diplomati\allowbreak{}c,current\}/expose\_XV\allowbreak{}III\_body.tex;\allowbreak{} folio\allowbreak{} 277\allowbreak{}}\par

\par\medskip\noindent\hypertarget{app-E7II-XVIII-010.27}{}\textbf{E7II-XVIII-010.27}\quad printed folio 277.\par
\noindent Printed reading: {\ttfamily\footnotesize  \&=\textbackslash{}left(g\textasciicircum{}*c\_1(\textbackslash{}unde\allowbreak{}rline\{\textbackslash{}mathcal\allowbreak{} O\}\_\{\textbackslash{}Delta(1)\})\allowbreak{} -(\textbackslash{}rho\allowbreak{} i)\textasciicircum{}*c\_1(\textbackslash{}xi\_\{\textbackslash{}mathbb\allowbreak{} P\textasciicircum{}1\})\textbackslash{}right)\textasciicircum{}2\textbackslash{}\textbackslash{}\allowbreak{}}\par
\noindent Corrected reading: \emph{(empty)}\par
\noindent Source: {\ttfamily\footnotesize ED-EN-NORM-G039;\allowbreak{} normalized-equal\allowbreak{} subscript-scope\allowbreak{} audit;\allowbreak{} fr/sga7ii/\{diplomati\allowbreak{}c,current\}/expose\_XV\allowbreak{}III\_body.tex;\allowbreak{} folio\allowbreak{} 277\allowbreak{}}\par

\par\medskip\noindent\hypertarget{app-E7II-XVIII-010.28}{}\textbf{E7II-XVIII-010.28}\quad printed folio 277.\par
\noindent Printed reading: {\ttfamily\footnotesize \textbackslash{}begin\{aligned\}\allowbreak{} }\par
\noindent Corrected reading: \emph{(empty)}\par
\noindent Source: {\ttfamily\footnotesize fr/sga7ii/\{diplomati\allowbreak{}c,current\}/expose\_XV\allowbreak{}III\_body.tex;\allowbreak{} folio\allowbreak{} 277\allowbreak{}}\par

\par\medskip\noindent\hypertarget{app-E7II-XVIII-010.29}{}\textbf{E7II-XVIII-010.29}\quad printed folio 277.\par
\noindent Printed reading: {\ttfamily\footnotesize \&=-g\_*(g\textasciicircum{}*(xy)g\textasciicircum{}*(c\_\allowbreak{}1(\textbackslash{}xi\_\textbackslash{}Delta)\textasciicircum{}2))\allowbreak{} +g\_*(2g\textasciicircum{}*(xy)g\textasciicircum{}*(c\_1\allowbreak{}(\textbackslash{}xi\_\textbackslash{}Delta)(\textbackslash{}rho\allowbreak{} i)\textasciicircum{}*(c\_1(\textbackslash{}xi\_\{\textbackslash{}mathb\allowbreak{}b\allowbreak{} P\textasciicircum{}1\}))\textbackslash{}\textbackslash{}\allowbreak{} A\&\allowbreak{}}\par
\noindent Corrected reading: \emph{(empty)}\par
\noindent Source: {\ttfamily\footnotesize fr/sga7ii/\{diplomati\allowbreak{}c,current\}/expose\_XV\allowbreak{}III\_body.tex;\allowbreak{} folio\allowbreak{} 277\allowbreak{}}\par

\par\medskip\noindent\hypertarget{app-E7II-XVIII-010.30}{}\textbf{E7II-XVIII-010.30}\quad printed folio 277.\par
\noindent Printed reading: {\ttfamily\footnotesize  \textbackslash{}end\{aligned\}\allowbreak{}}\par
\noindent Corrected reading: \emph{(empty)}\par
\noindent Source: {\ttfamily\footnotesize fr/sga7ii/\{diplomati\allowbreak{}c,current\}/expose\_XV\allowbreak{}III\_body.tex;\allowbreak{} folio\allowbreak{} 277\allowbreak{}}\par

\par\medskip\noindent\hypertarget{app-E7II-XVIII-010.31}{}\textbf{E7II-XVIII-010.31}\quad printed folio 277.\par
\noindent Printed reading: {\ttfamily\footnotesize ,\allowbreak{} \$(a\textbackslash{}oplus0)\textbackslash{}Lambda(b\allowbreak{}\textbackslash{}oplus0)=ab\textbackslash{}oplus0\$,\allowbreak{}}\par
\noindent Corrected reading: \emph{(empty)}\par
\noindent Source: {\ttfamily\footnotesize ED-EN-FORMULA-G177;\allowbreak{} fr/sga7ii/\{diplomati\allowbreak{}c,current\}/expose\_XV\allowbreak{}III\_body.tex;\allowbreak{} folio\allowbreak{} 277\allowbreak{}}\par

\par\medskip\noindent\hypertarget{app-E7II-XVIII-010.32}{}\textbf{E7II-XVIII-010.32}\quad printed folio 277.\par
\noindent Printed reading: {\ttfamily\footnotesize \textbackslash{}cap\allowbreak{}}\par
\noindent Corrected reading: {\ttfamily\footnotesize \$\allowbreak{} is\allowbreak{} stable\allowbreak{} under\allowbreak{} multiplication\allowbreak{} by\allowbreak{} \$\allowbreak{}}\par
\noindent Source: {\ttfamily\footnotesize fr/sga7ii/\{diplomati\allowbreak{}c,current\}/expose\_XV\allowbreak{}III\_body.tex;\allowbreak{} folio\allowbreak{} 277\allowbreak{}}\par

\par\medskip\noindent\hypertarget{app-E7II-XVIII-010.33}{}\textbf{E7II-XVIII-010.33}\quad printed folio 277.\par
\noindent Printed reading: {\ttfamily\footnotesize *\textbackslash{}subset\textbackslash{}operatornam\allowbreak{}e\{Ker\}f\_\allowbreak{}}\par
\noindent Corrected reading: \emph{(empty)}\par
\noindent Source: {\ttfamily\footnotesize fr/sga7ii/\{diplomati\allowbreak{}c,current\}/expose\_XV\allowbreak{}III\_body.tex;\allowbreak{} folio\allowbreak{} 277\allowbreak{}}\par

\par\medskip\noindent\hypertarget{app-E7II-XVIII-010.34}{}\textbf{E7II-XVIII-010.34}\quad printed folio 277.\par
\noindent Printed reading: {\ttfamily\footnotesize \textbackslash{}Lambda\allowbreak{}}\par
\noindent Corrected reading: {\ttfamily\footnotesize \textbackslash{}wedge\allowbreak{}}\par
\noindent Source: {\ttfamily\footnotesize fr/sga7ii/\{diplomati\allowbreak{}c,current\}/expose\_XV\allowbreak{}III\_body.tex;\allowbreak{} folio\allowbreak{} 277;\allowbreak{} formula\allowbreak{} 447\allowbreak{}}\par

\par\medskip\noindent\hypertarget{app-E7II-XVIII-010.35}{}\textbf{E7II-XVIII-010.35}\quad printed folio 278.\par
\noindent Printed reading: \emph{(empty)}\par
\noindent Corrected reading: {\ttfamily\footnotesize  (4.2.11)\allowbreak{} and\allowbreak{}}\par
\noindent Source: {\ttfamily\footnotesize ED-EN-FORMULA-G179;\allowbreak{} fr/sga7ii/\{diplomati\allowbreak{}c,current\}/expose\_XV\allowbreak{}III\_body.tex;\allowbreak{} folio\allowbreak{} 278;\allowbreak{} port/PORT\_EN7II\_LOG.\allowbreak{}jsonl\#C824;\allowbreak{} fixes/FIXLOG\_7II\_CUR\allowbreak{}RENT.jsonl\#C824;\allowbreak{} 7II:820\allowbreak{}}\par

\par\medskip\noindent\hypertarget{app-E7II-XVIII-011.1}{}\textbf{E7II-XVIII-011.1}\quad printed folio 279.\par
\noindent Printed reading: {\ttfamily\footnotesize \textasciicircum{}\allowbreak{}}\par
\noindent Corrected reading: {\ttfamily\footnotesize \_\allowbreak{}}\par
\noindent Source: {\ttfamily\footnotesize fr/sga7ii/\{diplomati\allowbreak{}c,current\}/expose\_XV\allowbreak{}III\_body.tex;\allowbreak{} folio\allowbreak{} 279;\allowbreak{} formula\allowbreak{} 508\allowbreak{}}\par

\par\medskip\noindent\hypertarget{app-E7II-XVIII-012.1}{}\textbf{E7II-XVIII-012.1}\quad printed folio 280.\par
\noindent Printed reading: \emph{(empty)}\par
\noindent Corrected reading: {\ttfamily\footnotesize -\allowbreak{}}\par
\noindent Source: {\ttfamily\footnotesize fr/sga7ii/\{diplomati\allowbreak{}c,current\}/expose\_XV\allowbreak{}III\_body.tex;\allowbreak{} folio\allowbreak{} 280;\allowbreak{} formula\allowbreak{} 540\allowbreak{}}\par

\par\medskip\noindent\hypertarget{app-E7II-XVIII-012.2}{}\textbf{E7II-XVIII-012.2}\quad printed folio 280.\par
\noindent Printed reading: \emph{(empty)}\par
\noindent Corrected reading: {\ttfamily\footnotesize -\allowbreak{}}\par
\noindent Source: {\ttfamily\footnotesize fr/sga7ii/\{diplomati\allowbreak{}c,current\}/expose\_XV\allowbreak{}III\_body.tex;\allowbreak{} folio\allowbreak{} 280;\allowbreak{} formula\allowbreak{} 542\allowbreak{}}\par

\par\medskip\noindent\hypertarget{app-E7II-XVIII-013.1}{}\textbf{E7II-XVIII-013.1}\quad printed folio 280.\par
\noindent Printed reading: {\ttfamily\footnotesize h\allowbreak{}}\par
\noindent Corrected reading: {\ttfamily\footnotesize k\allowbreak{}}\par
\noindent Source: {\ttfamily\footnotesize fr/sga7ii/\{diplomati\allowbreak{}c,current\}/expose\_XV\allowbreak{}III\_body.tex;\allowbreak{} folio\allowbreak{} 280;\allowbreak{} formula\allowbreak{} 549\allowbreak{}}\par

\par\medskip\noindent\hypertarget{app-E7II-XVIII-013.2}{}\textbf{E7II-XVIII-013.2}\quad printed folio 280.\par
\noindent Printed reading: {\ttfamily\footnotesize \textasciicircum{}\allowbreak{}}\par
\noindent Corrected reading: {\ttfamily\footnotesize \_\allowbreak{}}\par
\noindent Source: {\ttfamily\footnotesize fr/sga7ii/\{diplomati\allowbreak{}c,current\}/expose\_XV\allowbreak{}III\_body.tex;\allowbreak{} folio\allowbreak{} 280;\allowbreak{} formula\allowbreak{} 551\allowbreak{}}\par

\par\medskip\noindent\hypertarget{app-E7II-XVIII-013.3}{}\textbf{E7II-XVIII-013.3}\quad printed folio 281.\par
\noindent Printed reading: \emph{(empty)}\par
\noindent Corrected reading: {\ttfamily\footnotesize )\textbackslash{}bigr\allowbreak{}}\par
\noindent Source: {\ttfamily\footnotesize fr/sga7ii/\{diplomati\allowbreak{}c,current\}/expose\_XV\allowbreak{}III\_body.tex;\allowbreak{} folio\allowbreak{} 281\allowbreak{}}\par

\par\medskip\noindent\hypertarget{app-E7II-XVIII-014.1}{}\textbf{E7II-XVIII-014.1}\quad printed folio 281.\par
\noindent Printed reading: \emph{(empty)}\par
\noindent Corrected reading: {\ttfamily\footnotesize )\allowbreak{}}\par
\noindent Source: {\ttfamily\footnotesize fr/sga7ii/\{diplomati\allowbreak{}c,current\}/expose\_XV\allowbreak{}III\_body.tex;\allowbreak{} folio\allowbreak{} 281;\allowbreak{} formula\allowbreak{} 564\allowbreak{}}\par

\par\medskip\noindent\hypertarget{app-E7II-XVIII-015.1}{}\textbf{E7II-XVIII-015.1}\quad printed folio 282.\par
\noindent Printed reading: {\ttfamily\footnotesize 2n\allowbreak{}}\par
\noindent Corrected reading: {\ttfamily\footnotesize n\allowbreak{}}\par
\noindent Source: {\ttfamily\footnotesize ED-EN-PORT-WITNESS-7\allowbreak{}II-C752;\allowbreak{} fr/sga7ii/diplomatic\allowbreak{}/expose\_XVIII\_body.t\allowbreak{}ex;\allowbreak{} folio\allowbreak{} 282\allowbreak{}}\par

\par\medskip\noindent\hypertarget{app-E7II-XVIII-015.2}{}\textbf{E7II-XVIII-015.2}\quad printed folio 282.\par
\noindent Printed reading: {\ttfamily\footnotesize 2i\allowbreak{}}\par
\noindent Corrected reading: {\ttfamily\footnotesize i\allowbreak{}}\par
\noindent Source: {\ttfamily\footnotesize ED-EN-PORT-WITNESS-7\allowbreak{}II-C752;\allowbreak{} fr/sga7ii/diplomatic\allowbreak{}/expose\_XVIII\_body.t\allowbreak{}ex;\allowbreak{} folio\allowbreak{} 282\allowbreak{}}\par

\par\medskip\noindent\hypertarget{app-E7II-XVIII-015.3}{}\textbf{E7II-XVIII-015.3}\quad printed folio 282.\par
\noindent Printed reading: {\ttfamily\footnotesize 2\allowbreak{}}\par
\noindent Corrected reading: {\ttfamily\footnotesize 1\allowbreak{}}\par
\noindent Source: {\ttfamily\footnotesize ED-EN-PORT-WITNESS-7\allowbreak{}II-C752;\allowbreak{} fr/sga7ii/diplomatic\allowbreak{}/expose\_XVIII\_body.t\allowbreak{}ex;\allowbreak{} folio\allowbreak{} 282\allowbreak{}}\par

\par\medskip\noindent\hypertarget{app-E7II-XVIII-015.4}{}\textbf{E7II-XVIII-015.4}\quad printed folio 282.\par
\noindent Printed reading: {\ttfamily\footnotesize .\allowbreak{}}\par
\noindent Corrected reading: {\ttfamily\footnotesize )\allowbreak{}}\par
\noindent Source: {\ttfamily\footnotesize ED-EN-PORT-WITNESS-7\allowbreak{}II-C752;\allowbreak{} fr/sga7ii/diplomatic\allowbreak{}/expose\_XVIII\_body.t\allowbreak{}ex;\allowbreak{} folio\allowbreak{} 282\allowbreak{}}\par

\par\medskip\noindent\hypertarget{app-E7II-XVIII-016.1}{}\textbf{E7II-XVIII-016.1}\quad printed folio 282.\par
\noindent Printed reading: {\ttfamily\footnotesize m\allowbreak{}}\par
\noindent Corrected reading: {\ttfamily\footnotesize \textbackslash{}tau\allowbreak{}}\par
\noindent Source: {\ttfamily\footnotesize fr/sga7ii/\{diplomati\allowbreak{}c,current\}/expose\_XV\allowbreak{}III\_body.tex;\allowbreak{} folio\allowbreak{} 282;\allowbreak{} formula\allowbreak{} 586\allowbreak{}}\par

\par\medskip\noindent\hypertarget{app-E7II-XVIII-016.2}{}\textbf{E7II-XVIII-016.2}\quad printed folio 283.\par
\noindent Printed reading: {\ttfamily\footnotesize \textbackslash{}arrow[dl,\allowbreak{} "\textbackslash{}tau\_*"]\allowbreak{}}\par
\noindent Corrected reading: \emph{(empty)}\par
\noindent Source: {\ttfamily\footnotesize fr/sga7ii/\{diplomati\allowbreak{}c,current\}/expose\_XV\allowbreak{}III\_body.tex;\allowbreak{} folio\allowbreak{} 283;\allowbreak{} formula\allowbreak{} 587\allowbreak{}}\par

\par\medskip\noindent\hypertarget{app-E7II-XVIII-016.3}{}\textbf{E7II-XVIII-016.3}\quad printed folio 283.\par
\noindent Printed reading: \emph{(2 espace(s), 0 retour(s) de ligne dans le témoin TeX)}\par
\noindent Corrected reading: {\ttfamily\footnotesize  \textbackslash{}arrow[ur,"\textbackslash{}tau\_*"']\allowbreak{} }\par
\noindent Source: {\ttfamily\footnotesize fr/sga7ii/\{diplomati\allowbreak{}c,current\}/expose\_XV\allowbreak{}III\_body.tex;\allowbreak{} folio\allowbreak{} 283;\allowbreak{} formula\allowbreak{} 587\allowbreak{}}\par

\par\medskip\noindent\hypertarget{app-E7II-XVIII-017.1}{}\textbf{E7II-XVIII-017.1}\quad printed folio 284.\par
\noindent Printed reading: {\ttfamily\footnotesize \textbackslash{}notag\allowbreak{}}\par
\noindent Corrected reading: {\ttfamily\footnotesize \textbackslash{}tag\{5.2.4.7\}\allowbreak{}}\par
\noindent Source: {\ttfamily\footnotesize ED-EN-FORMULA-G181;\allowbreak{} fr/sga7ii/\{diplomati\allowbreak{}c,current\}/expose\_XV\allowbreak{}III\_body.tex;\allowbreak{} folio\allowbreak{} 284\allowbreak{}}\par

\par\medskip\noindent\hypertarget{app-E7II-XVIII-017.2}{}\textbf{E7II-XVIII-017.2}\quad printed folio 284.\par
\noindent Printed reading: {\ttfamily\footnotesize 2\allowbreak{}}\par
\noindent Corrected reading: {\ttfamily\footnotesize 1)\allowbreak{}}\par
\noindent Source: {\ttfamily\footnotesize ED-EN-FORMULA-G181;\allowbreak{} fr/sga7ii/\{diplomati\allowbreak{}c,current\}/expose\_XV\allowbreak{}III\_body.tex;\allowbreak{} folio\allowbreak{} 284\allowbreak{}}\par

\par\medskip\noindent\hypertarget{app-E7II-XVIII-017.3}{}\textbf{E7II-XVIII-017.3}\quad printed folio 285.\par
\noindent Printed reading: {\ttfamily\footnotesize 5\allowbreak{}}\par
\noindent Corrected reading: {\ttfamily\footnotesize 4.3\allowbreak{}}\par
\noindent Source: {\ttfamily\footnotesize ED-EN-NONMATH-G027;\allowbreak{} F7II-XVIII-047.1;\allowbreak{} 7II:790;\allowbreak{} fr/sga7ii/\{diplomati\allowbreak{}c,current\}/expose\_XV\allowbreak{}III\_body.tex;\allowbreak{} folio\allowbreak{} 285\allowbreak{}}\par

\par\medskip\noindent\hypertarget{app-E7II-XVIII-018.1}{}\textbf{E7II-XVIII-018.1}\quad printed folio 286.\par
\noindent Printed reading: {\ttfamily\footnotesize ii\allowbreak{}}\par
\noindent Corrected reading: {\ttfamily\footnotesize 2i\allowbreak{}}\par
\noindent Source: {\ttfamily\footnotesize fr/sga7ii/\{diplomati\allowbreak{}c,current\}/expose\_XV\allowbreak{}III\_body.tex;\allowbreak{} folio\allowbreak{} 286;\allowbreak{} formula\allowbreak{} 628\allowbreak{}}\par

\par\medskip\noindent\hypertarget{app-E7II-XVIII-019.1}{}\textbf{E7II-XVIII-019.1}\quad printed folio 286.\par
\noindent Printed reading: {\ttfamily\footnotesize 5\allowbreak{}}\par
\noindent Corrected reading: {\ttfamily\footnotesize 6\allowbreak{}}\par
\noindent Source: {\ttfamily\footnotesize ED-EN-NONMATH-G028;\allowbreak{} F7II-XVIII-050.4;\allowbreak{} 7II:481;\allowbreak{} fr/sga7ii/\{diplomati\allowbreak{}c,current\}/expose\_XV\allowbreak{}III\_body.tex;\allowbreak{} folio\allowbreak{} 286\allowbreak{}}\par

\par\medskip\noindent\hypertarget{app-E7II-XVIII-020.1}{}\textbf{E7II-XVIII-020.1}\quad printed folio 288.\par
\noindent Printed reading: \emph{(empty)}\par
\noindent Corrected reading: {\ttfamily\footnotesize )\allowbreak{}}\par
\noindent Source: {\ttfamily\footnotesize fr/sga7ii/\{diplomati\allowbreak{}c,current\}/expose\_XV\allowbreak{}III\_body.tex;\allowbreak{} folio\allowbreak{} 288;\allowbreak{} formula\allowbreak{} 664\allowbreak{}}\par

\par\medskip\noindent\hypertarget{app-E7II-XVIII-020.2}{}\textbf{E7II-XVIII-020.2}\quad printed folio 289.\par
\noindent Printed reading: {\ttfamily\footnotesize \textasciicircum{}\allowbreak{}}\par
\noindent Corrected reading: {\ttfamily\footnotesize \_\allowbreak{}}\par
\noindent Source: {\ttfamily\footnotesize fr/sga7ii/\{diplomati\allowbreak{}c,current\}/expose\_XV\allowbreak{}III\_body.tex;\allowbreak{} folio\allowbreak{} 289;\allowbreak{} formula\allowbreak{} 671\allowbreak{}}\par

\par\medskip\noindent\hypertarget{app-E7II-XVIII-021.1}{}\textbf{E7II-XVIII-021.1}\quad printed folio 290.\par
\noindent Printed reading: {\ttfamily\footnotesize 5.3.2\allowbreak{}}\par
\noindent Corrected reading: {\ttfamily\footnotesize A\allowbreak{}}\par
\noindent Source: {\ttfamily\footnotesize ED-EN-NONMATH-G029;\allowbreak{} F7II-XVIII-054.2;\allowbreak{} WITNESS-F7II-XVIII-0\allowbreak{}54.2;\allowbreak{} fr/sga7ii/\{diplomati\allowbreak{}c,current\}/expose\_XV\allowbreak{}III\_body.tex;\allowbreak{} folio\allowbreak{} 290\allowbreak{}}\par

\par\medskip\noindent\hypertarget{app-E7II-XVIII-021.2}{}\textbf{E7II-XVIII-021.2}\quad printed folio 290.\par
\noindent Printed reading: {\ttfamily\footnotesize m\allowbreak{}}\par
\noindent Corrected reading: {\ttfamily\footnotesize p\allowbreak{}}\par
\noindent Source: {\ttfamily\footnotesize fr/sga7ii/\{diplomati\allowbreak{}c,current\}/expose\_XV\allowbreak{}III\_body.tex;\allowbreak{} folio\allowbreak{} 290;\allowbreak{} formula\allowbreak{} 704\allowbreak{}}\par

\par\medskip\noindent\hypertarget{app-E7II-XVIII-022.1}{}\textbf{E7II-XVIII-022.1}\quad printed folio 291.\par
\noindent Printed reading: {\ttfamily\footnotesize o\allowbreak{}}\par
\noindent Corrected reading: {\ttfamily\footnotesize 0\allowbreak{}}\par
\noindent Source: {\ttfamily\footnotesize fr/sga7ii/\{diplomati\allowbreak{}c,current\}/expose\_XV\allowbreak{}III\_body.tex;\allowbreak{} folio\allowbreak{} 291;\allowbreak{} formula\allowbreak{} 715\allowbreak{}}\par

\par\medskip\noindent\hypertarget{app-E7II-XVIII-022.2}{}\textbf{E7II-XVIII-022.2}\quad printed folio 291.\par
\noindent Printed reading: {\ttfamily\footnotesize o\allowbreak{}}\par
\noindent Corrected reading: {\ttfamily\footnotesize 0\allowbreak{}}\par
\noindent Source: {\ttfamily\footnotesize fr/sga7ii/\{diplomati\allowbreak{}c,current\}/expose\_XV\allowbreak{}III\_body.tex;\allowbreak{} folio\allowbreak{} 291;\allowbreak{} formula\allowbreak{} 715\allowbreak{}}\par

\par\medskip\noindent\hypertarget{app-E7II-XVIII-022.3}{}\textbf{E7II-XVIII-022.3}\quad printed folio 291.\par
\noindent Printed reading: {\ttfamily\footnotesize o\allowbreak{}}\par
\noindent Corrected reading: {\ttfamily\footnotesize 0\allowbreak{}}\par
\noindent Source: {\ttfamily\footnotesize fr/sga7ii/\{diplomati\allowbreak{}c,current\}/expose\_XV\allowbreak{}III\_body.tex;\allowbreak{} folio\allowbreak{} 291\allowbreak{}}\par

\par\medskip\noindent\hypertarget{app-E7II-XVIII-022.4}{}\textbf{E7II-XVIII-022.4}\quad printed folio 291.\par
\noindent Printed reading: {\ttfamily\footnotesize (\allowbreak{}}\par
\noindent Corrected reading: \emph{(empty)}\par
\noindent Source: {\ttfamily\footnotesize ED-EN-FORMULA-G183;\allowbreak{} fr/sga7ii/\{diplomati\allowbreak{}c,current\}/expose\_XV\allowbreak{}III\_body.tex;\allowbreak{} folio\allowbreak{} 291\allowbreak{}}\par

\par\medskip\noindent\hypertarget{app-E7II-XVIII-023.1}{}\textbf{E7II-XVIII-023.1}\quad printed folio 292.\par
\noindent Printed reading: {\ttfamily\footnotesize f\allowbreak{}}\par
\noindent Corrected reading: {\ttfamily\footnotesize \textbackslash{}rho\allowbreak{}}\par
\noindent Source: {\ttfamily\footnotesize fr/sga7ii/\{diplomati\allowbreak{}c,current\}/expose\_XV\allowbreak{}III\_body.tex;\allowbreak{} folio\allowbreak{} 292;\allowbreak{} formula\allowbreak{} 724\allowbreak{}}\par

\par\medskip\noindent\hypertarget{app-E7II-XVIII-023.2}{}\textbf{E7II-XVIII-023.2}\quad printed folio 292.\par
\noindent Printed reading: {\ttfamily\footnotesize f\allowbreak{}}\par
\noindent Corrected reading: {\ttfamily\footnotesize \textbackslash{}rho\allowbreak{}}\par
\noindent Source: {\ttfamily\footnotesize fr/sga7ii/\{diplomati\allowbreak{}c,current\}/expose\_XV\allowbreak{}III\_body.tex;\allowbreak{} folio\allowbreak{} 292;\allowbreak{} formula\allowbreak{} 725\allowbreak{}}\par

\par\medskip\noindent\hypertarget{app-E7II-XVIII-023.3}{}\textbf{E7II-XVIII-023.3}\quad printed folio 292.\par
\noindent Printed reading: {\ttfamily\footnotesize .\allowbreak{}}\par
\noindent Corrected reading: {\ttfamily\footnotesize )\allowbreak{}}\par
\noindent Source: {\ttfamily\footnotesize fr/sga7ii/\{diplomati\allowbreak{}c,current\}/expose\_XV\allowbreak{}III\_body.tex;\allowbreak{} folio\allowbreak{} 292;\allowbreak{} port/decisions/decis\allowbreak{}ions\_en7II\_r2\_003.js\allowbreak{}on\#C755;\allowbreak{} fixes/FIXLOG\_7II\_CUR\allowbreak{}RENT.jsonl\#C755;\allowbreak{} 7II:490\allowbreak{}}\par

\par\medskip\noindent\hypertarget{app-E7II-XVIII-023.4}{}\textbf{E7II-XVIII-023.4}\quad printed folio 292.\par
\noindent Printed reading: \emph{(empty)}\par
\noindent Corrected reading: {\ttfamily\footnotesize )\allowbreak{}}\par
\noindent Source: {\ttfamily\footnotesize fr/sga7ii/\{diplomati\allowbreak{}c,current\}/expose\_XV\allowbreak{}III\_body.tex;\allowbreak{} folio\allowbreak{} 292;\allowbreak{} formula\allowbreak{} 737\allowbreak{}}\par

\par\medskip\noindent\hypertarget{app-E7II-XVIII-023.5}{}\textbf{E7II-XVIII-023.5}\quad printed folio 292.\par
\noindent Printed reading: {\ttfamily\footnotesize 2h\allowbreak{}}\par
\noindent Corrected reading: {\ttfamily\footnotesize 2+h\allowbreak{}}\par
\noindent Source: {\ttfamily\footnotesize fr/sga7ii/\{diplomati\allowbreak{}c,current\}/expose\_XV\allowbreak{}III\_body.tex;\allowbreak{} folio\allowbreak{} 292;\allowbreak{} formula\allowbreak{} 737\allowbreak{}}\par

\par\medskip\noindent\hypertarget{app-E7II-XVIII-023.6}{}\textbf{E7II-XVIII-023.6}\quad printed folio 292.\par
\noindent Printed reading: \emph{(empty)}\par
\noindent Corrected reading: {\ttfamily\footnotesize )\allowbreak{}}\par
\noindent Source: {\ttfamily\footnotesize fr/sga7ii/\{diplomati\allowbreak{}c,current\}/expose\_XV\allowbreak{}III\_body.tex;\allowbreak{} folio\allowbreak{} 292;\allowbreak{} formula\allowbreak{} 737\allowbreak{}}\par

\par\medskip\noindent\hypertarget{app-E7II-XVIII-023.7}{}\textbf{E7II-XVIII-023.7}\quad printed folio 292.\par
\noindent Printed reading: \emph{(empty)}\par
\noindent Corrected reading: {\ttfamily\footnotesize )\allowbreak{}}\par
\noindent Source: {\ttfamily\footnotesize fr/sga7ii/\{diplomati\allowbreak{}c,current\}/expose\_XV\allowbreak{}III\_body.tex;\allowbreak{} folio\allowbreak{} 292;\allowbreak{} formula\allowbreak{} 738\allowbreak{}}\par

\par\medskip\noindent\hypertarget{app-E7II-XVIII-023.8}{}\textbf{E7II-XVIII-023.8}\quad printed folio 292.\par
\noindent Printed reading: {\ttfamily\footnotesize 8\allowbreak{}}\par
\noindent Corrected reading: {\ttfamily\footnotesize 5\allowbreak{}}\par
\noindent Source: {\ttfamily\footnotesize ED-EN-NONMATH-G030;\allowbreak{} F7II-XVIII-058.5;\allowbreak{} 7II:494;\allowbreak{} fr/sga7ii/\{diplomati\allowbreak{}c,current\}/expose\_XV\allowbreak{}III\_body.tex;\allowbreak{} folio\allowbreak{} 292\allowbreak{}}\par

\par\medskip\noindent\hypertarget{app-E7II-XVIII-024.1}{}\textbf{E7II-XVIII-024.1}\quad printed folio 294.\par
\noindent Printed reading: {\ttfamily\footnotesize \{\allowbreak{}}\par
\noindent Corrected reading: \emph{(empty)}\par
\noindent Source: {\ttfamily\footnotesize ED-EN-NORM-G040;\allowbreak{} normalized-equal\allowbreak{} subscript-scope\allowbreak{} audit;\allowbreak{} fr/sga7ii/\{diplomati\allowbreak{}c,current\}/expose\_XV\allowbreak{}III\_body.tex;\allowbreak{} folio\allowbreak{} 294\allowbreak{}}\par

\par\medskip\noindent\hypertarget{app-E7II-XVIII-024.2}{}\textbf{E7II-XVIII-024.2}\quad printed folio 294.\par
\noindent Printed reading: {\ttfamily\footnotesize \}\allowbreak{}}\par
\noindent Corrected reading: \emph{(empty)}\par
\noindent Source: {\ttfamily\footnotesize ED-EN-NORM-G040;\allowbreak{} normalized-equal\allowbreak{} subscript-scope\allowbreak{} audit;\allowbreak{} fr/sga7ii/\{diplomati\allowbreak{}c,current\}/expose\_XV\allowbreak{}III\_body.tex;\allowbreak{} folio\allowbreak{} 294\allowbreak{}}\par

\par\medskip\noindent\hypertarget{app-E7II-XVIII-024.3}{}\textbf{E7II-XVIII-024.3}\quad printed folio 295.\par
\noindent Printed reading: \emph{(empty)}\par
\noindent Corrected reading: {\ttfamily\footnotesize \textasciicircum{}\{i+1\}\allowbreak{}}\par
\noindent Source: {\ttfamily\footnotesize ED-EN-FORMULA-G185;\allowbreak{} fr/sga7ii/\{diplomati\allowbreak{}c,current\}/expose\_XV\allowbreak{}III\_body.tex;\allowbreak{} folio\allowbreak{} 295\allowbreak{}}\par

\par\medskip\noindent\hypertarget{app-E7II-XVIII-024.4}{}\textbf{E7II-XVIII-024.4}\quad printed folio 295.\par
\noindent Printed reading: {\ttfamily\footnotesize -1\allowbreak{}}\par
\noindent Corrected reading: \emph{(empty)}\par
\noindent Source: {\ttfamily\footnotesize ED-EN-FORMULA-G185;\allowbreak{} fr/sga7ii/\{diplomati\allowbreak{}c,current\}/expose\_XV\allowbreak{}III\_body.tex;\allowbreak{} folio\allowbreak{} 295\allowbreak{}}\par

\par\medskip\noindent\hypertarget{app-E7II-XVIII-024.5}{}\textbf{E7II-XVIII-024.5}\quad printed folio 295.\par
\noindent Printed reading: \emph{(empty)}\par
\noindent Corrected reading: {\ttfamily\footnotesize \textbackslash{}rho\_*\allowbreak{}}\par
\noindent Source: {\ttfamily\footnotesize ED-EN-FORMULA-G185;\allowbreak{} fr/sga7ii/\{diplomati\allowbreak{}c,current\}/expose\_XV\allowbreak{}III\_body.tex;\allowbreak{} folio\allowbreak{} 295\allowbreak{}}\par

\par\medskip\noindent\hypertarget{app-E7II-XVIII-024.6}{}\textbf{E7II-XVIII-024.6}\quad printed folio 295.\par
\noindent Printed reading: {\ttfamily\footnotesize 1\allowbreak{}}\par
\noindent Corrected reading: {\ttfamily\footnotesize 0\allowbreak{}}\par
\noindent Source: {\ttfamily\footnotesize ED-EN-FORMULA-G185;\allowbreak{} fr/sga7ii/\{diplomati\allowbreak{}c,current\}/expose\_XV\allowbreak{}III\_body.tex;\allowbreak{} folio\allowbreak{} 295\allowbreak{}}\par

\par\medskip\noindent\hypertarget{app-E7II-XVIII-024.7}{}\textbf{E7II-XVIII-024.7}\quad printed folio 295.\par
\noindent Printed reading: {\ttfamily\footnotesize ,\allowbreak{}}\par
\noindent Corrected reading: \emph{(empty)}\par
\noindent Source: {\ttfamily\footnotesize fr/sga7ii/\{diplomati\allowbreak{}c,current\}/expose\_XV\allowbreak{}III\_body.tex;\allowbreak{} folio\allowbreak{} 295;\allowbreak{} formula\allowbreak{} 781\allowbreak{}}\par

\par\medskip\noindent\hypertarget{app-E7II-XVIII-024.8}{}\textbf{E7II-XVIII-024.8}\quad printed folio 295.\par
\noindent Printed reading: {\ttfamily\footnotesize \textbackslash{}mathbb\allowbreak{} }\par
\noindent Corrected reading: \emph{(empty)}\par
\noindent Source: {\ttfamily\footnotesize ED-EN-FONT-G019;\allowbreak{} F7II-XVIII-060.7;\allowbreak{} fr/sga7ii/\{diplomati\allowbreak{}c,current\}/expose\_XV\allowbreak{}III\_body.tex;\allowbreak{} folio\allowbreak{} 295\allowbreak{}}\par

\par\medskip\noindent\hypertarget{app-E7II-XVIII-024.9}{}\textbf{E7II-XVIII-024.9}\quad printed folio 295.\par
\noindent Printed reading: {\ttfamily\footnotesize 1\allowbreak{}}\par
\noindent Corrected reading: {\ttfamily\footnotesize n\allowbreak{}}\par
\noindent Source: {\ttfamily\footnotesize ED-EN-FONT-G019;\allowbreak{} F7II-XVIII-060.7;\allowbreak{} fr/sga7ii/\{diplomati\allowbreak{}c,current\}/expose\_XV\allowbreak{}III\_body.tex;\allowbreak{} folio\allowbreak{} 295;\allowbreak{} fr/sga7ii/\{diplomati\allowbreak{}c,current\}/expose\_XV\allowbreak{}III\_body.tex;\allowbreak{} folio\allowbreak{} 295;\allowbreak{} formula\allowbreak{} 784\allowbreak{}}\par

\par\medskip\noindent\hypertarget{app-E7II-XVIII-024.10}{}\textbf{E7II-XVIII-024.10}\quad printed folio 295.\par
\noindent Printed reading: \emph{(empty)}\par
\noindent Corrected reading: {\ttfamily\footnotesize \textbackslash{}oplus\textbackslash{}cdots\allowbreak{}}\par
\noindent Source: {\ttfamily\footnotesize fr/sga7ii/\{diplomati\allowbreak{}c,current\}/expose\_XV\allowbreak{}III\_body.tex;\allowbreak{} folio\allowbreak{} 295;\allowbreak{} formula\allowbreak{} 786\allowbreak{}}\par

\par\medskip\noindent\hypertarget{app-E7II-XVIII-024.11}{}\textbf{E7II-XVIII-024.11}\quad printed folio 297.\par
\noindent Printed reading: {\ttfamily\footnotesize c\allowbreak{}}\par
\noindent Corrected reading: {\ttfamily\footnotesize 0\allowbreak{}}\par
\noindent Source: {\ttfamily\footnotesize fr/sga7ii/\{diplomati\allowbreak{}c,current\}/expose\_XV\allowbreak{}III\_body.tex;\allowbreak{} folio\allowbreak{} 297;\allowbreak{} formula\allowbreak{} 802\allowbreak{}}\par

\par\medskip\noindent\hypertarget{app-E7II-XVIII-024.12}{}\textbf{E7II-XVIII-024.12}\quad printed folio 297.\par
\noindent Printed reading: \emph{(empty)}\par
\noindent Corrected reading: {\ttfamily\footnotesize E\allowbreak{}}\par
\noindent Source: {\ttfamily\footnotesize fr/sga7ii/\{diplomati\allowbreak{}c,current\}/expose\_XV\allowbreak{}III\_body.tex;\allowbreak{} folio\allowbreak{} 297;\allowbreak{} formula\allowbreak{} 802\allowbreak{}}\par

\par\medskip\noindent\hypertarget{app-E7II-XVIII-025.1}{}\textbf{E7II-XVIII-025.1}\quad printed folio 299.\par
\noindent Printed reading: {\ttfamily\footnotesize :\allowbreak{}}\par
\noindent Corrected reading: \emph{(empty)}\par
\noindent Source: {\ttfamily\footnotesize fr/sga7ii/\{diplomati\allowbreak{}c,current\}/expose\_XV\allowbreak{}III\_body.tex;\allowbreak{} folio\allowbreak{} 299;\allowbreak{} formula\allowbreak{} 834\allowbreak{}}\par

\par\medskip\noindent\hypertarget{app-E7II-XVIII-025.2}{}\textbf{E7II-XVIII-025.2}\quad printed folio 299.\par
\noindent Printed reading: {\ttfamily\footnotesize 2\allowbreak{}}\par
\noindent Corrected reading: {\ttfamily\footnotesize i\allowbreak{}}\par
\noindent Source: {\ttfamily\footnotesize fr/sga7ii/\{diplomati\allowbreak{}c,current\}/expose\_XV\allowbreak{}III\_body.tex;\allowbreak{} folio\allowbreak{} 299;\allowbreak{} formula\allowbreak{} 839\allowbreak{}}\par

\par\medskip\noindent\hypertarget{app-E7II-XVIII-025.3}{}\textbf{E7II-XVIII-025.3}\quad printed folio 300.\par
\noindent Printed reading: {\ttfamily\footnotesize N\allowbreak{}}\par
\noindent Corrected reading: {\ttfamily\footnotesize n\allowbreak{}}\par
\noindent Source: {\ttfamily\footnotesize fr/sga7ii/\{diplomati\allowbreak{}c,current\}/expose\_XV\allowbreak{}III\_body.tex;\allowbreak{} folio\allowbreak{} 300;\allowbreak{} formula\allowbreak{} 841\allowbreak{}}\par

\par\medskip\noindent\hypertarget{app-E7II-XVIII-025.4}{}\textbf{E7II-XVIII-025.4}\quad printed folio 300.\par
\noindent Printed reading: \emph{(empty)}\par
\noindent Corrected reading: {\ttfamily\footnotesize (-1\allowbreak{}}\par
\noindent Source: {\ttfamily\footnotesize fr/sga7ii/\{diplomati\allowbreak{}c,current\}/expose\_XV\allowbreak{}III\_body.tex;\allowbreak{} folio\allowbreak{} 300;\allowbreak{} formula\allowbreak{} 844\allowbreak{}}\par

\par\medskip\noindent\hypertarget{app-E7II-XVIII-025.5}{}\textbf{E7II-XVIII-025.5}\quad printed folio 300.\par
\noindent Printed reading: \emph{(empty)}\par
\noindent Corrected reading: {\ttfamily\footnotesize \textasciicircum{}n\allowbreak{}}\par
\noindent Source: {\ttfamily\footnotesize fr/sga7ii/\{diplomati\allowbreak{}c,current\}/expose\_XV\allowbreak{}III\_body.tex;\allowbreak{} folio\allowbreak{} 300;\allowbreak{} formula\allowbreak{} 849\allowbreak{}}\par

\par\medskip\noindent\hypertarget{app-E7II-XVIII-026.1}{}\textbf{E7II-XVIII-026.1}\quad printed folio 301.\par
\noindent Printed reading: {\ttfamily\footnotesize \_\{\allowbreak{}}\par
\noindent Corrected reading: \emph{(empty)}\par
\noindent Source: {\ttfamily\footnotesize fr/sga7ii/\{diplomati\allowbreak{}c,current\}/expose\_XV\allowbreak{}III\_body.tex;\allowbreak{} folio\allowbreak{} 301;\allowbreak{} formula\allowbreak{} 867\allowbreak{}}\par

\par\medskip\noindent\hypertarget{app-E7II-XVIII-026.2}{}\textbf{E7II-XVIII-026.2}\quad printed folio 301.\par
\noindent Printed reading: {\ttfamily\footnotesize \}\allowbreak{}}\par
\noindent Corrected reading: \emph{(empty)}\par
\noindent Source: {\ttfamily\footnotesize fr/sga7ii/\{diplomati\allowbreak{}c,current\}/expose\_XV\allowbreak{}III\_body.tex;\allowbreak{} folio\allowbreak{} 301;\allowbreak{} formula\allowbreak{} 867\allowbreak{}}\par

\par\medskip\noindent\hypertarget{app-E7II-XVIII-026.3}{}\textbf{E7II-XVIII-026.3}\quad printed folio 301.\par
\noindent Printed reading: \emph{(empty)}\par
\noindent Corrected reading: {\ttfamily\footnotesize )\allowbreak{}}\par
\noindent Source: {\ttfamily\footnotesize fr/sga7ii/\{diplomati\allowbreak{}c,current\}/expose\_XV\allowbreak{}III\_body.tex;\allowbreak{} folio\allowbreak{} 301;\allowbreak{} formula\allowbreak{} 867\allowbreak{}}\par

\par\medskip\noindent\hypertarget{app-E7II-XVIII-026.4}{}\textbf{E7II-XVIII-026.4}\quad printed folio 301.\par
\noindent Printed reading: {\ttfamily\footnotesize .\allowbreak{}}\par
\noindent Corrected reading: {\ttfamily\footnotesize )\allowbreak{}}\par
\noindent Source: {\ttfamily\footnotesize fr/sga7ii/\{diplomati\allowbreak{}c,current\}/expose\_XV\allowbreak{}III\_body.tex;\allowbreak{} folio\allowbreak{} 301;\allowbreak{} formula\allowbreak{} 867\allowbreak{}}\par

\par\medskip\noindent\hypertarget{app-E7II-XVIII-026.5}{}\textbf{E7II-XVIII-026.5}\quad printed folio 302.\par
\noindent Printed reading: {\ttfamily\footnotesize \_\{\allowbreak{}}\par
\noindent Corrected reading: \emph{(empty)}\par
\noindent Source: {\ttfamily\footnotesize fr/sga7ii/\{diplomati\allowbreak{}c,current\}/expose\_XV\allowbreak{}III\_body.tex;\allowbreak{} folio\allowbreak{} 302;\allowbreak{} formula\allowbreak{} 869\allowbreak{}}\par

\par\medskip\noindent\hypertarget{app-E7II-XVIII-026.6}{}\textbf{E7II-XVIII-026.6}\quad printed folio 302.\par
\noindent Printed reading: {\ttfamily\footnotesize \}\allowbreak{}}\par
\noindent Corrected reading: \emph{(empty)}\par
\noindent Source: {\ttfamily\footnotesize fr/sga7ii/\{diplomati\allowbreak{}c,current\}/expose\_XV\allowbreak{}III\_body.tex;\allowbreak{} folio\allowbreak{} 302;\allowbreak{} formula\allowbreak{} 869\allowbreak{}}\par

\par\medskip\noindent\hypertarget{app-E7II-XVIII-026.7}{}\textbf{E7II-XVIII-026.7}\quad printed folio 302.\par
\noindent Printed reading: \emph{(empty)}\par
\noindent Corrected reading: {\ttfamily\footnotesize \textasciicircum{}i\allowbreak{}}\par
\noindent Source: {\ttfamily\footnotesize fr/sga7ii/\{diplomati\allowbreak{}c,current\}/expose\_XV\allowbreak{}III\_body.tex;\allowbreak{} folio\allowbreak{} 302;\allowbreak{} formula\allowbreak{} 870\allowbreak{}}\par

\par\medskip\noindent\hypertarget{app-E7II-XVIII-026.8}{}\textbf{E7II-XVIII-026.8}\quad printed folio 302.\par
\noindent Printed reading: {\ttfamily\footnotesize \_\{\allowbreak{}}\par
\noindent Corrected reading: \emph{(empty)}\par
\noindent Source: {\ttfamily\footnotesize fr/sga7ii/\{diplomati\allowbreak{}c,current\}/expose\_XV\allowbreak{}III\_body.tex;\allowbreak{} folio\allowbreak{} 302;\allowbreak{} formula\allowbreak{} 870\allowbreak{}}\par

\par\medskip\noindent\hypertarget{app-E7II-XVIII-026.9}{}\textbf{E7II-XVIII-026.9}\quad printed folio 302.\par
\noindent Printed reading: {\ttfamily\footnotesize \}\allowbreak{}}\par
\noindent Corrected reading: \emph{(empty)}\par
\noindent Source: {\ttfamily\footnotesize fr/sga7ii/\{diplomati\allowbreak{}c,current\}/expose\_XV\allowbreak{}III\_body.tex;\allowbreak{} folio\allowbreak{} 302;\allowbreak{} formula\allowbreak{} 870\allowbreak{}}\par

\par\medskip\noindent\hypertarget{app-E7II-XVIII-026.10}{}\textbf{E7II-XVIII-026.10}\quad printed folio 302.\par
\noindent Printed reading: {\ttfamily\footnotesize \textbackslash{}cdot\allowbreak{}}\par
\noindent Corrected reading: {\ttfamily\footnotesize i\allowbreak{}}\par
\noindent Source: {\ttfamily\footnotesize fr/sga7ii/\{diplomati\allowbreak{}c,current\}/expose\_XV\allowbreak{}III\_body.tex;\allowbreak{} folio\allowbreak{} 302;\allowbreak{} formula\allowbreak{} 870\allowbreak{}}\par

\par\medskip\noindent\hypertarget{app-E7II-XVIII-026.11}{}\textbf{E7II-XVIII-026.11}\quad printed folio 303.\par
\noindent Printed reading: \emph{(empty)}\par
\noindent Corrected reading: {\ttfamily\footnotesize )\allowbreak{}}\par
\noindent Source: {\ttfamily\footnotesize fr/sga7ii/\{diplomati\allowbreak{}c,current\}/expose\_XV\allowbreak{}III\_body.tex;\allowbreak{} folio\allowbreak{} 303;\allowbreak{} formula\allowbreak{} 877\allowbreak{}}\par

\par\medskip\noindent\hypertarget{app-E7II-XVIII-027.1}{}\textbf{E7II-XVIII-027.1}\quad printed folio 303.\par
\noindent Printed reading: {\ttfamily\footnotesize o\allowbreak{}}\par
\noindent Corrected reading: {\ttfamily\footnotesize 0\allowbreak{}}\par
\noindent Source: {\ttfamily\footnotesize ED-EN-FORMULA-G145;\allowbreak{} fr/sga7ii/\{diplomati\allowbreak{}c,current\}/expose\_XV\allowbreak{}III\_body.tex;\allowbreak{} folio\allowbreak{} 303\allowbreak{}}\par

\par\medskip\noindent\hypertarget{app-E7II-XVIII-027.2}{}\textbf{E7II-XVIII-027.2}\quad printed folio 303.\par
\noindent Printed reading: {\ttfamily\footnotesize o\allowbreak{}}\par
\noindent Corrected reading: {\ttfamily\footnotesize 0\allowbreak{}}\par
\noindent Source: {\ttfamily\footnotesize ED-EN-FORMULA-G145;\allowbreak{} fr/sga7ii/\{diplomati\allowbreak{}c,current\}/expose\_XV\allowbreak{}III\_body.tex;\allowbreak{} folio\allowbreak{} 303\allowbreak{}}\par

\par\medskip\noindent\hypertarget{app-E7II-XVIII-027.3}{}\textbf{E7II-XVIII-027.3}\quad printed folio 303.\par
\noindent Printed reading: {\ttfamily\footnotesize [\allowbreak{}}\par
\noindent Corrected reading: \emph{(1 espace(s), 0 retour(s) de ligne dans le témoin TeX)}\par
\noindent Source: {\ttfamily\footnotesize ED-EN-FORMULA-G145;\allowbreak{} fr/sga7ii/\{diplomati\allowbreak{}c,current\}/expose\_XV\allowbreak{}III\_body.tex;\allowbreak{} folio\allowbreak{} 303\allowbreak{}}\par

\par\medskip\noindent\hypertarget{app-E7II-XVIII-027.4}{}\textbf{E7II-XVIII-027.4}\quad printed folio 303.\par
\noindent Printed reading: {\ttfamily\footnotesize \textbackslash{}alpha\allowbreak{}}\par
\noindent Corrected reading: \emph{(empty)}\par
\noindent Source: {\ttfamily\footnotesize ED-EN-FORMULA-G145;\allowbreak{} fr/sga7ii/\{diplomati\allowbreak{}c,current\}/expose\_XV\allowbreak{}III\_body.tex;\allowbreak{} folio\allowbreak{} 303\allowbreak{}}\par

\par\medskip\noindent\hypertarget{app-E7II-XVIII-027.5}{}\textbf{E7II-XVIII-027.5}\quad printed folio 303.\par
\noindent Printed reading: {\ttfamily\footnotesize ]\allowbreak{}}\par
\noindent Corrected reading: \emph{(empty)}\par
\noindent Source: {\ttfamily\footnotesize ED-EN-FORMULA-G145;\allowbreak{} fr/sga7ii/\{diplomati\allowbreak{}c,current\}/expose\_XV\allowbreak{}III\_body.tex;\allowbreak{} folio\allowbreak{} 303\allowbreak{}}\par

\par\medskip\noindent\hypertarget{app-E7II-XVIII-027.6}{}\textbf{E7II-XVIII-027.6}\quad printed folio 303.\par
\noindent Printed reading: \emph{(empty)}\par
\noindent Corrected reading: {\ttfamily\footnotesize if\allowbreak{} and\allowbreak{} only\allowbreak{} if\allowbreak{} }\par
\noindent Source: {\ttfamily\footnotesize ED-EN-FORMULA-G145;\allowbreak{} fr/sga7ii/\{diplomati\allowbreak{}c,current\}/expose\_XV\allowbreak{}III\_body.tex;\allowbreak{} folio\allowbreak{} 303\allowbreak{}}\par

\par\medskip\noindent\hypertarget{app-E7II-XVIII-027.7}{}\textbf{E7II-XVIII-027.7}\quad printed folio 303.\par
\noindent Printed reading: {\ttfamily\footnotesize \textbackslash{}Longrightarrow\allowbreak{} }\par
\noindent Corrected reading: \emph{(empty)}\par
\noindent Source: {\ttfamily\footnotesize ED-EN-FORMULA-G145;\allowbreak{} fr/sga7ii/\{diplomati\allowbreak{}c,current\}/expose\_XV\allowbreak{}III\_body.tex;\allowbreak{} folio\allowbreak{} 303\allowbreak{}}\par

\par\medskip\noindent\hypertarget{app-E7II-XVIII-027.8}{}\textbf{E7II-XVIII-027.8}\quad printed folio 303.\par
\noindent Printed reading: {\ttfamily\footnotesize \textbackslash{}alpha\allowbreak{}}\par
\noindent Corrected reading: {\ttfamily\footnotesize x\allowbreak{}}\par
\noindent Source: {\ttfamily\footnotesize ED-EN-FORMULA-G145;\allowbreak{} fr/sga7ii/\{diplomati\allowbreak{}c,current\}/expose\_XV\allowbreak{}III\_body.tex;\allowbreak{} folio\allowbreak{} 303\allowbreak{}}\par

\par\medskip\noindent\hypertarget{app-E7II-XVIII-028.1}{}\textbf{E7II-XVIII-028.1}\quad printed folio 303.\par
\noindent Printed reading: {\ttfamily\footnotesize \{[\textbackslash{}mathrm\{\allowbreak{}}\par
\noindent Corrected reading: \emph{(empty)}\par
\noindent Source: {\ttfamily\footnotesize fr/sga7ii/\{diplomati\allowbreak{}c,current\}/expose\_XV\allowbreak{}III\_body.tex;\allowbreak{} folio\allowbreak{} 303;\allowbreak{} formula\allowbreak{} 887\allowbreak{}}\par
\noindent Editorial explanation: Retain the supplied diplomatic witness’s explicitly bracketed conjectural label, including roman U and question mark; its source comment identifies an illegible scan label. The unbracketed automatic U\textasciicircum{}n? would misrepresent the witness. Corrected English retains its released U\textasciicircum{}n; the apparatus and inherited uncertainty note distinguish the supplied conjecture from a secure printed reading.\par

\par\medskip\noindent\hypertarget{app-E7II-XVIII-028.2}{}\textbf{E7II-XVIII-028.2}\quad printed folio 303.\par
\noindent Printed reading: {\ttfamily\footnotesize \}\allowbreak{}}\par
\noindent Corrected reading: \emph{(empty)}\par
\noindent Source: {\ttfamily\footnotesize fr/sga7ii/\{diplomati\allowbreak{}c,current\}/expose\_XV\allowbreak{}III\_body.tex;\allowbreak{} folio\allowbreak{} 303;\allowbreak{} formula\allowbreak{} 887\allowbreak{}}\par
\noindent Editorial explanation: Retain the supplied diplomatic witness’s explicitly bracketed conjectural label, including roman U and question mark; its source comment identifies an illegible scan label. The unbracketed automatic U\textasciicircum{}n? would misrepresent the witness. Corrected English retains its released U\textasciicircum{}n; the apparatus and inherited uncertainty note distinguish the supplied conjecture from a secure printed reading.\par

\par\medskip\noindent\hypertarget{app-E7II-XVIII-028.3}{}\textbf{E7II-XVIII-028.3}\quad printed folio 303.\par
\noindent Printed reading: {\ttfamily\footnotesize \textbackslash{},?]\}\allowbreak{}}\par
\noindent Corrected reading: \emph{(empty)}\par
\noindent Source: {\ttfamily\footnotesize fr/sga7ii/\{diplomati\allowbreak{}c,current\}/expose\_XV\allowbreak{}III\_body.tex;\allowbreak{} folio\allowbreak{} 303;\allowbreak{} formula\allowbreak{} 887\allowbreak{}}\par
\noindent Editorial explanation: Retain the supplied diplomatic witness’s explicitly bracketed conjectural label, including roman U and question mark; its source comment identifies an illegible scan label. The unbracketed automatic U\textasciicircum{}n? would misrepresent the witness. Corrected English retains its released U\textasciicircum{}n; the apparatus and inherited uncertainty note distinguish the supplied conjecture from a secure printed reading.\par

\par\medskip\noindent\hypertarget{app-E7II-XVIII-029.1}{}\textbf{E7II-XVIII-029.1}\quad printed folio 304.\par
\noindent Printed reading: \emph{(empty)}\par
\noindent Corrected reading: {\ttfamily\footnotesize \textbackslash{}nu\textasciicircum{}*\allowbreak{}}\par
\noindent Source: {\ttfamily\footnotesize fr/sga7ii/\{diplomati\allowbreak{}c,current\}/expose\_XV\allowbreak{}III\_body.tex;\allowbreak{} folio\allowbreak{} 304;\allowbreak{} formula\allowbreak{} 917\allowbreak{}}\par

\par\medskip\noindent\hypertarget{app-E7II-XVIII-029.2}{}\textbf{E7II-XVIII-029.2}\quad printed folio 304.\par
\noindent Printed reading: {\ttfamily\footnotesize (\allowbreak{}}\par
\noindent Corrected reading: \emph{(empty)}\par
\noindent Source: {\ttfamily\footnotesize fr/sga7ii/\{diplomati\allowbreak{}c,current\}/expose\_XV\allowbreak{}III\_body.tex;\allowbreak{} folio\allowbreak{} 304;\allowbreak{} formula\allowbreak{} 917\allowbreak{}}\par

\par\medskip\noindent\hypertarget{app-E7II-XVIII-029.3}{}\textbf{E7II-XVIII-029.3}\quad printed folio 305.\par
\noindent Printed reading: {\ttfamily\footnotesize )\allowbreak{}}\par
\noindent Corrected reading: \emph{(empty)}\par
\noindent Source: {\ttfamily\footnotesize fr/sga7ii/\{diplomati\allowbreak{}c,current\}/expose\_XV\allowbreak{}III\_body.tex;\allowbreak{} folio\allowbreak{} 305;\allowbreak{} formula\allowbreak{} 931\allowbreak{}}\par

\par\medskip\noindent\hypertarget{app-E7II-XVIII-029.4}{}\textbf{E7II-XVIII-029.4}\quad printed folio 305.\par
\noindent Printed reading: \emph{(empty)}\par
\noindent Corrected reading: {\ttfamily\footnotesize (j)\allowbreak{}}\par
\noindent Source: {\ttfamily\footnotesize fr/sga7ii/\{diplomati\allowbreak{}c,current\}/expose\_XV\allowbreak{}III\_body.tex;\allowbreak{} folio\allowbreak{} 305;\allowbreak{} formula\allowbreak{} 942\allowbreak{}}\par

\par\medskip\noindent\hypertarget{app-E7II-XVIII-029.5}{}\textbf{E7II-XVIII-029.5}\quad printed folio 305.\par
\noindent Printed reading: \emph{(empty)}\par
\noindent Corrected reading: {\ttfamily\footnotesize (j)\allowbreak{}}\par
\noindent Source: {\ttfamily\footnotesize fr/sga7ii/\{diplomati\allowbreak{}c,current\}/expose\_XV\allowbreak{}III\_body.tex;\allowbreak{} folio\allowbreak{} 305;\allowbreak{} formula\allowbreak{} 944\allowbreak{}}\par

\par\medskip\noindent\hypertarget{app-E7II-XVIII-030.1}{}\textbf{E7II-XVIII-030.1}\quad printed folio 308.\par
\noindent Printed reading: {\ttfamily\footnotesize (\allowbreak{}}\par
\noindent Corrected reading: \emph{(empty)}\par
\noindent Source: {\ttfamily\footnotesize fr/sga7ii/\{diplomati\allowbreak{}c,current\}/expose\_XV\allowbreak{}III\_body.tex;\allowbreak{} folio\allowbreak{} 308;\allowbreak{} formula\allowbreak{} 967\allowbreak{}}\par

\par\medskip\noindent\hypertarget{app-E7II-XVIII-030.2}{}\textbf{E7II-XVIII-030.2}\quad printed folio 308.\par
\noindent Printed reading: {\ttfamily\footnotesize Vf)\allowbreak{}}\par
\noindent Corrected reading: {\ttfamily\footnotesize V\textasciicircum{}*f\allowbreak{}}\par
\noindent Source: {\ttfamily\footnotesize fr/sga7ii/\{diplomati\allowbreak{}c,current\}/expose\_XV\allowbreak{}III\_body.tex;\allowbreak{} folio\allowbreak{} 308;\allowbreak{} formula\allowbreak{} 967\allowbreak{}}\par

\par\medskip\noindent\hypertarget{app-E7II-XVIII-030.3}{}\textbf{E7II-XVIII-030.3}\quad printed folio 308.\par
\noindent Printed reading: {\ttfamily\footnotesize \_\allowbreak{}}\par
\noindent Corrected reading: \emph{(empty)}\par
\noindent Source: {\ttfamily\footnotesize fr/sga7ii/\{diplomati\allowbreak{}c,current\}/expose\_XV\allowbreak{}III\_body.tex;\allowbreak{} folio\allowbreak{} 308;\allowbreak{} formula\allowbreak{} 967\allowbreak{}}\par

\par\medskip\noindent\hypertarget{app-E7II-XVIII-030.4}{}\textbf{E7II-XVIII-030.4}\quad printed folio 308.\par
\noindent Printed reading: {\ttfamily\footnotesize \_\allowbreak{}}\par
\noindent Corrected reading: \emph{(empty)}\par
\noindent Source: {\ttfamily\footnotesize fr/sga7ii/\{diplomati\allowbreak{}c,current\}/expose\_XV\allowbreak{}III\_body.tex;\allowbreak{} folio\allowbreak{} 308;\allowbreak{} formula\allowbreak{} 969\allowbreak{}}\par

\par\medskip\noindent\hypertarget{app-E7II-XVIII-030.5}{}\textbf{E7II-XVIII-030.5}\quad printed folio 308.\par
\noindent Printed reading: {\ttfamily\footnotesize (\allowbreak{}}\par
\noindent Corrected reading: \emph{(empty)}\par
\noindent Source: {\ttfamily\footnotesize fr/sga7ii/\{diplomati\allowbreak{}c,current\}/expose\_XV\allowbreak{}III\_body.tex;\allowbreak{} folio\allowbreak{} 308;\allowbreak{} formula\allowbreak{} 969\allowbreak{}}\par

\par\medskip\noindent\hypertarget{app-E7II-XVIII-030.6}{}\textbf{E7II-XVIII-030.6}\quad printed folio 308.\par
\noindent Printed reading: {\ttfamily\footnotesize )\allowbreak{}}\par
\noindent Corrected reading: \emph{(empty)}\par
\noindent Source: {\ttfamily\footnotesize fr/sga7ii/\{diplomati\allowbreak{}c,current\}/expose\_XV\allowbreak{}III\_body.tex;\allowbreak{} folio\allowbreak{} 308;\allowbreak{} formula\allowbreak{} 969\allowbreak{}}\par

\par\medskip\noindent\hypertarget{app-E7II-XVIII-030.7}{}\textbf{E7II-XVIII-030.7}\quad printed folio 308.\par
\noindent Printed reading: {\ttfamily\footnotesize \_\allowbreak{}}\par
\noindent Corrected reading: \emph{(empty)}\par
\noindent Source: {\ttfamily\footnotesize fr/sga7ii/\{diplomati\allowbreak{}c,current\}/expose\_XV\allowbreak{}III\_body.tex;\allowbreak{} folio\allowbreak{} 308;\allowbreak{} formula\allowbreak{} 969\allowbreak{}}\par

\par\medskip\noindent\hypertarget{app-E7II-XVIII-030.8}{}\textbf{E7II-XVIII-030.8}\quad printed folio 308.\par
\noindent Printed reading: {\ttfamily\footnotesize \_\allowbreak{}}\par
\noindent Corrected reading: \emph{(empty)}\par
\noindent Source: {\ttfamily\footnotesize fr/sga7ii/\{diplomati\allowbreak{}c,current\}/expose\_XV\allowbreak{}III\_body.tex;\allowbreak{} folio\allowbreak{} 308;\allowbreak{} formula\allowbreak{} 970\allowbreak{}}\par

\par\medskip\noindent\hypertarget{app-E7II-XVIII-030.9}{}\textbf{E7II-XVIII-030.9}\quad printed folio 308.\par
\noindent Printed reading: {\ttfamily\footnotesize )\allowbreak{}}\par
\noindent Corrected reading: \emph{(empty)}\par
\noindent Source: {\ttfamily\footnotesize fr/sga7ii/\{diplomati\allowbreak{}c,current\}/expose\_XV\allowbreak{}III\_body.tex;\allowbreak{} folio\allowbreak{} 308;\allowbreak{} formula\allowbreak{} 970\allowbreak{}}\par

\par\medskip\noindent\hypertarget{app-E7II-XVIII-030.10}{}\textbf{E7II-XVIII-030.10}\quad printed folio 309.\par
\noindent Printed reading: {\ttfamily\footnotesize \textbackslash{}notag\allowbreak{}}\par
\noindent Corrected reading: {\ttfamily\footnotesize \textbackslash{}tag\{5.8.7.3\}\allowbreak{}}\par
\noindent Source: {\ttfamily\footnotesize ED-EN-NORM-G031;\allowbreak{} complete\allowbreak{} normalized-equal\allowbreak{} formula\allowbreak{} audit;\allowbreak{} fr/sga7ii/\{diplomati\allowbreak{}c,current\}/expose\_XV\allowbreak{}III\_body.tex;\allowbreak{} folio\allowbreak{} 309\allowbreak{}}\par

\par\medskip\noindent\hypertarget{app-E7II-XVIII-030.11}{}\textbf{E7II-XVIII-030.11}\quad printed folio 309.\par
\noindent Printed reading: \emph{(empty)}\par
\noindent Corrected reading: {\ttfamily\footnotesize \_U\allowbreak{}}\par
\noindent Source: {\ttfamily\footnotesize ED-EN-PORT-WITNESS-7\allowbreak{}II-C889;\allowbreak{} fr/sga7ii/diplomatic\allowbreak{}/expose\_XVIII\_body.t\allowbreak{}ex;\allowbreak{} folio\allowbreak{} 309;\allowbreak{} port/PORT\_EN7II\_LOG.\allowbreak{}jsonl\#C889\#C889;\allowbreak{} ;\allowbreak{} fr/sga7ii/\{diplomati\allowbreak{}c,current\}/expose\_XV\allowbreak{}III\_body.tex;\allowbreak{} folio\allowbreak{} 309;\allowbreak{} formula\allowbreak{} 975\allowbreak{}}\par

\par\medskip\noindent\hypertarget{app-E7II-XVIII-030.12}{}\textbf{E7II-XVIII-030.12}\quad printed folio 310.\par
\noindent Printed reading: \emph{(empty)}\par
\noindent Corrected reading: {\ttfamily\footnotesize (j)\allowbreak{}}\par
\noindent Source: {\ttfamily\footnotesize fr/sga7ii/\{diplomati\allowbreak{}c,current\}/expose\_XV\allowbreak{}III\_body.tex;\allowbreak{} folio\allowbreak{} 310;\allowbreak{} formula\allowbreak{} 976\allowbreak{}}\par

\par\medskip\noindent\hypertarget{app-E7II-XVIII-030.13}{}\textbf{E7II-XVIII-030.13}\quad printed folio 310.\par
\noindent Printed reading: \emph{(empty)}\par
\noindent Corrected reading: {\ttfamily\footnotesize F\allowbreak{}}\par
\noindent Source: {\ttfamily\footnotesize fr/sga7ii/\{diplomati\allowbreak{}c,current\}/expose\_XV\allowbreak{}III\_body.tex;\allowbreak{} folio\allowbreak{} 310;\allowbreak{} formula\allowbreak{} 983\allowbreak{}}\par

\par\medskip\noindent\hypertarget{app-E7II-XVIII-030.14}{}\textbf{E7II-XVIII-030.14}\quad printed folio 310.\par
\noindent Printed reading: {\ttfamily\footnotesize (j)\allowbreak{}}\par
\noindent Corrected reading: \emph{(empty)}\par
\noindent Source: {\ttfamily\footnotesize fr/sga7ii/\{diplomati\allowbreak{}c,current\}/expose\_XV\allowbreak{}III\_body.tex;\allowbreak{} folio\allowbreak{} 310;\allowbreak{} formula\allowbreak{} 983\allowbreak{}}\par

\par\medskip\noindent\hypertarget{app-E7II-XVIII-030.15}{}\textbf{E7II-XVIII-030.15}\quad printed folio 310.\par
\noindent Printed reading: \emph{(empty)}\par
\noindent Corrected reading: {\ttfamily\footnotesize \_U\allowbreak{}}\par
\noindent Source: {\ttfamily\footnotesize fr/sga7ii/\{diplomati\allowbreak{}c,current\}/expose\_XV\allowbreak{}III\_body.tex;\allowbreak{} folio\allowbreak{} 310;\allowbreak{} formula\allowbreak{} 984\allowbreak{}}\par

\par\medskip\noindent\hypertarget{app-E7II-XVIII-030.16}{}\textbf{E7II-XVIII-030.16}\quad printed folio 310.\par
\noindent Printed reading: \emph{(empty)}\par
\noindent Corrected reading: {\ttfamily\footnotesize \_U\allowbreak{}}\par
\noindent Source: {\ttfamily\footnotesize fr/sga7ii/\{diplomati\allowbreak{}c,current\}/expose\_XV\allowbreak{}III\_body.tex;\allowbreak{} folio\allowbreak{} 310;\allowbreak{} formula\allowbreak{} 984\allowbreak{}}\par

\par\medskip\noindent\hypertarget{app-E7II-XVIII-030.17}{}\textbf{E7II-XVIII-030.17}\quad printed folio 310.\par
\noindent Printed reading: {\ttfamily\footnotesize .\allowbreak{}}\par
\noindent Corrected reading: {\ttfamily\footnotesize \textbackslash{},\allowbreak{}}\par
\noindent Source: {\ttfamily\footnotesize fr/sga7ii/\{diplomati\allowbreak{}c,current\}/expose\_XV\allowbreak{}III\_body.tex;\allowbreak{} folio\allowbreak{} 310;\allowbreak{} formula\allowbreak{} 984\allowbreak{}}\par

\par\medskip\noindent\hypertarget{app-E7II-XVIII-031.1}{}\textbf{E7II-XVIII-031.1}\quad printed folio 314.\par
\noindent Printed reading: {\ttfamily\footnotesize O\allowbreak{}}\par
\noindent Corrected reading: {\ttfamily\footnotesize o\allowbreak{}}\par
\noindent Source: {\ttfamily\footnotesize fr/sga7ii/\{diplomati\allowbreak{}c,current\}/expose\_XV\allowbreak{}III\_body.tex;\allowbreak{} folio\allowbreak{} 314;\allowbreak{} formula\allowbreak{} 1032;\allowbreak{} port/decisions/decis\allowbreak{}ions\_en7II\_r2\_003.js\allowbreak{}on\#C759;\allowbreak{} fixes/FIXLOG\_7II\_CUR\allowbreak{}RENT.jsonl\#C759;\allowbreak{} 7II:533\allowbreak{}}\par

\par\medskip\noindent\hypertarget{app-E7II-XVIII-031.2}{}\textbf{E7II-XVIII-031.2}\quad printed folio 314.\par
\noindent Printed reading: {\ttfamily\footnotesize )\allowbreak{}}\par
\noindent Corrected reading: \emph{(empty)}\par
\noindent Source: {\ttfamily\footnotesize fr/sga7ii/\{diplomati\allowbreak{}c,current\}/expose\_XV\allowbreak{}III\_body.tex;\allowbreak{} folio\allowbreak{} 314;\allowbreak{} formula\allowbreak{} 1034\allowbreak{}}\par

\par\medskip\noindent\hypertarget{app-E7II-XVIII-031.3}{}\textbf{E7II-XVIII-031.3}\quad printed folio 314.\par
\noindent Printed reading: \emph{(empty)}\par
\noindent Corrected reading: {\ttfamily\footnotesize \textasciicircum{}\{\textbackslash{}mathrm\{mod\}\}\allowbreak{}}\par
\noindent Source: {\ttfamily\footnotesize ED-EN-FORMULA-G187;\allowbreak{} fr/sga7ii/\{diplomati\allowbreak{}c,current\}/expose\_XV\allowbreak{}III\_body.tex;\allowbreak{} folio\allowbreak{} 314\allowbreak{}}\par

\par\medskip\noindent\hypertarget{app-E7II-XVIII-031.4}{}\textbf{E7II-XVIII-031.4}\quad printed folio 314.\par
\noindent Printed reading: {\ttfamily\footnotesize \_\allowbreak{}}\par
\noindent Corrected reading: \emph{(empty)}\par
\noindent Source: {\ttfamily\footnotesize fr/sga7ii/\{diplomati\allowbreak{}c,current\}/expose\_XV\allowbreak{}III\_body.tex;\allowbreak{} folio\allowbreak{} 314;\allowbreak{} formula\allowbreak{} 1038\allowbreak{}}\par

\par\medskip\noindent\hypertarget{app-E7II-XVIII-031.5}{}\textbf{E7II-XVIII-031.5}\quad printed folio 315.\par
\noindent Printed reading: \emph{(empty)}\par
\noindent Corrected reading: {\ttfamily\footnotesize *\allowbreak{}}\par
\noindent Source: {\ttfamily\footnotesize fr/sga7ii/\{diplomati\allowbreak{}c,current\}/expose\_XV\allowbreak{}III\_body.tex;\allowbreak{} folio\allowbreak{} 315\allowbreak{}}\par

\par\medskip\noindent\hypertarget{app-E7II-XVIII-032.1}{}\textbf{E7II-XVIII-032.1}\quad printed folio 316.\par
\noindent Printed reading: {\ttfamily\footnotesize K\allowbreak{}}\par
\noindent Corrected reading: {\ttfamily\footnotesize k\allowbreak{}}\par
\noindent Source: {\ttfamily\footnotesize fr/sga7ii/\{diplomati\allowbreak{}c,current\}/expose\_XV\allowbreak{}III\_body.tex;\allowbreak{} folio\allowbreak{} 316;\allowbreak{} formula\allowbreak{} 1063\allowbreak{}}\par

\par\medskip\noindent\hypertarget{app-E7II-XVIII-032.2}{}\textbf{E7II-XVIII-032.2}\quad printed folio 316.\par
\noindent Printed reading: {\ttfamily\footnotesize \{\textbackslash{}overline\allowbreak{}}\par
\noindent Corrected reading: \emph{(empty)}\par
\noindent Source: {\ttfamily\footnotesize fr/sga7ii/\{diplomati\allowbreak{}c,current\}/expose\_XV\allowbreak{}III\_body.tex;\allowbreak{} folio\allowbreak{} 316;\allowbreak{} formula\allowbreak{} 1068\allowbreak{}}\par

\par\medskip\noindent\hypertarget{app-E7II-XVIII-032.3}{}\textbf{E7II-XVIII-032.3}\quad printed folio 316.\par
\noindent Printed reading: {\ttfamily\footnotesize \}\allowbreak{}}\par
\noindent Corrected reading: \emph{(empty)}\par
\noindent Source: {\ttfamily\footnotesize fr/sga7ii/\{diplomati\allowbreak{}c,current\}/expose\_XV\allowbreak{}III\_body.tex;\allowbreak{} folio\allowbreak{} 316;\allowbreak{} formula\allowbreak{} 1068\allowbreak{}}\par

\par\medskip\noindent\hypertarget{app-E7II-XVIII-033.1}{}\textbf{E7II-XVIII-033.1}\quad printed folio 316.\par
\noindent Printed reading: \emph{(empty)}\par
\noindent Corrected reading: {\ttfamily\footnotesize k(\textbackslash{}bar\textbackslash{}eta)\allowbreak{}}\par
\noindent Source: {\ttfamily\footnotesize fr/sga7ii/\{diplomati\allowbreak{}c,current\}/expose\_XV\allowbreak{}III\_body.tex;\allowbreak{} folio\allowbreak{} 316;\allowbreak{} formula\allowbreak{} 1081\allowbreak{}}\par

\par\medskip\noindent\hypertarget{app-E7II-XVIII-033.2}{}\textbf{E7II-XVIII-033.2}\quad printed folio 316.\par
\noindent Printed reading: \emph{(empty)}\par
\noindent Corrected reading: {\ttfamily\footnotesize )\textbackslash{}bigr\allowbreak{}}\par
\noindent Source: {\ttfamily\footnotesize fr/sga7ii/\{diplomati\allowbreak{}c,current\}/expose\_XV\allowbreak{}III\_body.tex;\allowbreak{} folio\allowbreak{} 316;\allowbreak{} formula\allowbreak{} 1081\allowbreak{}}\par

\par\medskip\noindent\hypertarget{app-E7II-XVIII-033.3}{}\textbf{E7II-XVIII-033.3}\quad printed folio 317.\par
\noindent Printed reading: \emph{(empty)}\par
\noindent Corrected reading: {\ttfamily\footnotesize k(\textbackslash{}bar\textbackslash{}eta)\allowbreak{}}\par
\noindent Source: {\ttfamily\footnotesize ED-EN-FORMULA-G189;\allowbreak{} fr/sga7ii/\{diplomati\allowbreak{}c,current\}/expose\_XV\allowbreak{}III\_body.tex;\allowbreak{} folio\allowbreak{} 317;\allowbreak{} fr/sga7ii/\{diplomati\allowbreak{}c,current\}/expose\_XV\allowbreak{}III\_body.tex;\allowbreak{} folio\allowbreak{} 317;\allowbreak{} formula\allowbreak{} 1082\allowbreak{}}\par

\par\medskip\noindent\hypertarget{app-E7II-XVIII-033.4}{}\textbf{E7II-XVIII-033.4}\quad printed folio 317.\par
\noindent Printed reading: {\ttfamily\footnotesize K\allowbreak{}}\par
\noindent Corrected reading: {\ttfamily\footnotesize D\allowbreak{}}\par
\noindent Source: {\ttfamily\footnotesize fr/sga7ii/\{diplomati\allowbreak{}c,current\}/expose\_XV\allowbreak{}III\_body.tex;\allowbreak{} folio\allowbreak{} 317;\allowbreak{} formula\allowbreak{} 1083\allowbreak{}}\par

\par\medskip\noindent\hypertarget{app-E7II-XVIII-034.1}{}\textbf{E7II-XVIII-034.1}\quad printed folio 318.\par
\noindent Printed reading: {\ttfamily\footnotesize card\allowbreak{}}\par
\noindent Corrected reading: {\ttfamily\footnotesize char\allowbreak{}}\par
\noindent Source: {\ttfamily\footnotesize fr/sga7ii/\{diplomati\allowbreak{}c,current\}/expose\_XV\allowbreak{}III\_body.tex;\allowbreak{} folio\allowbreak{} 318\allowbreak{}}\par

\par\medskip\noindent\hypertarget{app-E7II-XVIII-035.1}{}\textbf{E7II-XVIII-035.1}\quad printed folio 320.\par
\noindent Printed reading: {\ttfamily\footnotesize 6\allowbreak{}}\par
\noindent Corrected reading: {\ttfamily\footnotesize 10\allowbreak{}}\par
\noindent Source: {\ttfamily\footnotesize port/PORT\_EN7II\_LOG.\allowbreak{}jsonl\#C903;\allowbreak{} fixes/FIXLOG\_7II\_CUR\allowbreak{}RENT.jsonl\#C903;\allowbreak{} 7II:945\allowbreak{}}\par

\par\medskip\noindent\hypertarget{app-E7II-XVIII-035.2}{}\textbf{E7II-XVIII-035.2}\quad printed folio 320.\par
\noindent Printed reading: \emph{(empty)}\par
\noindent Corrected reading: {\ttfamily\footnotesize \textbackslash{}check\allowbreak{}}\par
\noindent Source: {\ttfamily\footnotesize fr/sga7ii/\{diplomati\allowbreak{}c,current\}/expose\_XV\allowbreak{}III\_body.tex;\allowbreak{} folio\allowbreak{} 320;\allowbreak{} formula\allowbreak{} 1167\allowbreak{}}\par

\par\medskip\noindent\hypertarget{app-E7II-XVIII-035.3}{}\textbf{E7II-XVIII-035.3}\quad printed folio 320.\par
\noindent Printed reading: \emph{(empty)}\par
\noindent Corrected reading: {\ttfamily\footnotesize \textbackslash{}mathbb\allowbreak{} }\par
\noindent Source: {\ttfamily\footnotesize ED-EN-FONT-G004;\allowbreak{} F7II-XVIII-103.4;\allowbreak{} fr/sga7ii/\{diplomati\allowbreak{}c,current\}/expose\_XV\allowbreak{}III\_body.tex;\allowbreak{} folio\allowbreak{} 320\allowbreak{}}\par

\par\medskip\noindent\hypertarget{app-E7II-XVIII-035.4}{}\textbf{E7II-XVIII-035.4}\quad printed folio 320.\par
\noindent Printed reading: {\ttfamily\footnotesize 2\allowbreak{}}\par
\noindent Corrected reading: {\ttfamily\footnotesize 1\allowbreak{}}\par
\noindent Source: {\ttfamily\footnotesize port/PORT\_EN7II\_LOG.\allowbreak{}jsonl\#C904;\allowbreak{} fixes/FIXLOG\_7II\_CUR\allowbreak{}RENT.jsonl\#C904;\allowbreak{} 7II:948\allowbreak{}}\par

\par\medskip\noindent\hypertarget{app-E7II-XVIII-035.5}{}\textbf{E7II-XVIII-035.5}\quad printed folio 320.\par
\noindent Printed reading: {\ttfamily\footnotesize ,\allowbreak{}}\par
\noindent Corrected reading: \emph{(empty)}\par
\noindent Source: {\ttfamily\footnotesize ED-EN-NONMATH-G052;\allowbreak{} F7II-XVIII-103.7;\allowbreak{} ADOPT-F7II-XVIII-103\allowbreak{}.7;\allowbreak{} fr/sga7ii/\{diplomati\allowbreak{}c,current\}/expose\_XV\allowbreak{}III\_body.tex;\allowbreak{} folio\allowbreak{} 320\allowbreak{}}\par

\par\medskip\noindent\hypertarget{app-E7II-XVIII-035.6}{}\textbf{E7II-XVIII-035.6}\quad printed folio 320.\par
\noindent Printed reading: \emph{(empty)}\par
\noindent Corrected reading: {\ttfamily\footnotesize ,\allowbreak{}}\par
\noindent Source: {\ttfamily\footnotesize ED-EN-NONMATH-G052;\allowbreak{} F7II-XVIII-103.7;\allowbreak{} ADOPT-F7II-XVIII-103\allowbreak{}.7;\allowbreak{} fr/sga7ii/\{diplomati\allowbreak{}c,current\}/expose\_XV\allowbreak{}III\_body.tex;\allowbreak{} folio\allowbreak{} 320\allowbreak{}}\par

\par\medskip\noindent\hypertarget{app-E7II-XVIII-036.1}{}\textbf{E7II-XVIII-036.1}\quad printed folio 321.\par
\noindent Printed reading: \emph{(empty)}\par
\noindent Corrected reading: {\ttfamily\footnotesize )\allowbreak{}}\par
\noindent Source: {\ttfamily\footnotesize fr/sga7ii/\{diplomati\allowbreak{}c,current\}/expose\_XV\allowbreak{}III\_body.tex;\allowbreak{} folio\allowbreak{} 321;\allowbreak{} formula\allowbreak{} 1196\allowbreak{}}\par

\par\medskip\noindent\hypertarget{app-E7II-XVIII-037.1}{}\textbf{E7II-XVIII-037.1}\quad printed folio 323.\par
\noindent Printed reading: {\ttfamily\footnotesize i\allowbreak{}}\par
\noindent Corrected reading: {\ttfamily\footnotesize n-1\allowbreak{}}\par
\noindent Source: {\ttfamily\footnotesize fr/sga7ii/\{diplomati\allowbreak{}c,current\}/expose\_XV\allowbreak{}III\_body.tex;\allowbreak{} folio\allowbreak{} 323;\allowbreak{} formula\allowbreak{} 1224\allowbreak{}}\par

\par\medskip\noindent\hypertarget{app-E7II-XVIII-037.2}{}\textbf{E7II-XVIII-037.2}\quad printed folio 323.\par
\noindent Printed reading: {\ttfamily\footnotesize 2\allowbreak{}}\par
\noindent Corrected reading: {\ttfamily\footnotesize 1\allowbreak{}}\par
\noindent Source: {\ttfamily\footnotesize ED-EN-NONMATH-G031;\allowbreak{} F7II-XVIII-106.1;\allowbreak{} 7II:551;\allowbreak{} fr/sga7ii/\{diplomati\allowbreak{}c,current\}/expose\_XV\allowbreak{}III\_body.tex;\allowbreak{} folio\allowbreak{} 323\allowbreak{}}\par

\par\medskip\noindent\hypertarget{app-E7II-XVIII-037.3}{}\textbf{E7II-XVIII-037.3}\quad printed folio 323.\par
\noindent Printed reading: \emph{(empty)}\par
\noindent Corrected reading: {\ttfamily\footnotesize )\allowbreak{}}\par
\noindent Source: {\ttfamily\footnotesize fr/sga7ii/\{diplomati\allowbreak{}c,current\}/expose\_XV\allowbreak{}III\_body.tex;\allowbreak{} folio\allowbreak{} 323;\allowbreak{} formula\allowbreak{} 1227\allowbreak{}}\par

\par\medskip\noindent\hypertarget{app-E7II-XVIII-037.4}{}\textbf{E7II-XVIII-037.4}\quad printed folio 323.\par
\noindent Printed reading: \emph{(empty)}\par
\noindent Corrected reading: {\ttfamily\footnotesize )\allowbreak{}}\par
\noindent Source: {\ttfamily\footnotesize fr/sga7ii/\{diplomati\allowbreak{}c,current\}/expose\_XV\allowbreak{}III\_body.tex;\allowbreak{} folio\allowbreak{} 323;\allowbreak{} formula\allowbreak{} 1233\allowbreak{}}\par

\par\medskip\noindent\hypertarget{app-E7II-XVIII-037.5}{}\textbf{E7II-XVIII-037.5}\quad printed folio 323.\par
\noindent Printed reading: {\ttfamily\footnotesize \textbackslash{}chi\allowbreak{}}\par
\noindent Corrected reading: {\ttfamily\footnotesize x\allowbreak{}}\par
\noindent Source: {\ttfamily\footnotesize fr/sga7ii/\{diplomati\allowbreak{}c,current\}/expose\_XV\allowbreak{}III\_body.tex;\allowbreak{} folio\allowbreak{} 323;\allowbreak{} formula\allowbreak{} 1236\allowbreak{}}\par

\par\medskip\noindent\hypertarget{app-E7II-XVIII-037.6}{}\textbf{E7II-XVIII-037.6}\quad printed folio 323.\par
\noindent Printed reading: \emph{(empty)}\par
\noindent Corrected reading: {\ttfamily\footnotesize )\allowbreak{}}\par
\noindent Source: {\ttfamily\footnotesize fr/sga7ii/\{diplomati\allowbreak{}c,current\}/expose\_XV\allowbreak{}III\_body.tex;\allowbreak{} folio\allowbreak{} 323;\allowbreak{} formula\allowbreak{} 1239\allowbreak{}}\par

\par\medskip\noindent\hypertarget{app-E7II-XVIII-037.7}{}\textbf{E7II-XVIII-037.7}\quad printed folio 323.\par
\noindent Printed reading: \emph{(empty)}\par
\noindent Corrected reading: {\ttfamily\footnotesize )\allowbreak{}}\par
\noindent Source: {\ttfamily\footnotesize fr/sga7ii/\{diplomati\allowbreak{}c,current\}/expose\_XV\allowbreak{}III\_body.tex;\allowbreak{} folio\allowbreak{} 323;\allowbreak{} formula\allowbreak{} 1242\allowbreak{}}\par

\par\medskip\noindent\hypertarget{app-E7II-XVIII-038.1}{}\textbf{E7II-XVIII-038.1}\quad printed folio 324.\par
\noindent Printed reading: {\ttfamily\footnotesize Let\allowbreak{}}\par
\noindent Corrected reading: {\ttfamily\footnotesize Via\allowbreak{} \$\textbackslash{}ell\_i\$,\allowbreak{} let\allowbreak{}}\par
\noindent Source: {\ttfamily\footnotesize fr/sga7ii/\{diplomati\allowbreak{}c,current\}/expose\_XV\allowbreak{}III\_body.tex;\allowbreak{} folio\allowbreak{} 324\allowbreak{}}\par

\par\medskip\noindent\hypertarget{app-E7II-XVIII-038.2}{}\textbf{E7II-XVIII-038.2}\quad printed folio 324.\par
\noindent Printed reading: {\ttfamily\footnotesize of\allowbreak{} \$\textbackslash{}delta\_i\$\allowbreak{} }\par
\noindent Corrected reading: \emph{(empty)}\par
\noindent Source: {\ttfamily\footnotesize fr/sga7ii/\{diplomati\allowbreak{}c,current\}/expose\_XV\allowbreak{}III\_body.tex;\allowbreak{} folio\allowbreak{} 324\allowbreak{}}\par

\par\medskip\noindent\hypertarget{app-E7II-XVIII-038.3}{}\textbf{E7II-XVIII-038.3}\quad printed folio 324.\par
\noindent Printed reading: {\ttfamily\footnotesize S\allowbreak{}}\par
\noindent Corrected reading: {\ttfamily\footnotesize s\allowbreak{}}\par
\noindent Source: {\ttfamily\footnotesize fr/sga7ii/\{diplomati\allowbreak{}c,current\}/expose\_XV\allowbreak{}III\_body.tex;\allowbreak{} folio\allowbreak{} 324\allowbreak{}}\par

\par\medskip\noindent\hypertarget{app-E7II-XVIII-038.4}{}\textbf{E7II-XVIII-038.4}\quad printed folio 324.\par
\noindent Printed reading: \emph{(empty)}\par
\noindent Corrected reading: {\ttfamily\footnotesize \_\{\textbackslash{}bar\textbackslash{}eta\}\allowbreak{}}\par
\noindent Source: {\ttfamily\footnotesize ED-EN-FORMULA-G193;\allowbreak{} fr/sga7ii/\{diplomati\allowbreak{}c,current\}/expose\_XV\allowbreak{}III\_body.tex;\allowbreak{} folio\allowbreak{} 324\allowbreak{}}\par

\par\medskip\noindent\hypertarget{app-E7II-XVIII-038.5}{}\textbf{E7II-XVIII-038.5}\quad printed folio 324.\par
\noindent Printed reading: \emph{(empty)}\par
\noindent Corrected reading: {\ttfamily\footnotesize )\allowbreak{}}\par
\noindent Source: {\ttfamily\footnotesize ED-EN-FORMULA-G193;\allowbreak{} fr/sga7ii/\{diplomati\allowbreak{}c,current\}/expose\_XV\allowbreak{}III\_body.tex;\allowbreak{} folio\allowbreak{} 324\allowbreak{}}\par

\par\medskip\noindent\hypertarget{app-E7II-XVIII-038.6}{}\textbf{E7II-XVIII-038.6}\quad printed folio 324.\par
\noindent Printed reading: {\ttfamily\footnotesize \_1\allowbreak{}}\par
\noindent Corrected reading: \emph{(empty)}\par
\noindent Source: {\ttfamily\footnotesize fr/sga7ii/\{diplomati\allowbreak{}c,current\}/expose\_XV\allowbreak{}III\_body.tex;\allowbreak{} folio\allowbreak{} 324\allowbreak{}}\par

\par\medskip\noindent\hypertarget{app-E7II-XVIII-038.7}{}\textbf{E7II-XVIII-038.7}\quad printed folio 324.\par
\noindent Printed reading: \emph{(empty)}\par
\noindent Corrected reading: {\ttfamily\footnotesize )\allowbreak{}}\par
\noindent Source: {\ttfamily\footnotesize fr/sga7ii/\{diplomati\allowbreak{}c,current\}/expose\_XV\allowbreak{}III\_body.tex;\allowbreak{} folio\allowbreak{} 324;\allowbreak{} formula\allowbreak{} 1272\allowbreak{}}\par

\par\medskip\noindent\hypertarget{app-E7II-XVIII-038.8}{}\textbf{E7II-XVIII-038.8}\quad printed folio 324.\par
\noindent Printed reading: \emph{(empty)}\par
\noindent Corrected reading: {\ttfamily\footnotesize )\allowbreak{}}\par
\noindent Source: {\ttfamily\footnotesize fr/sga7ii/\{diplomati\allowbreak{}c,current\}/expose\_XV\allowbreak{}III\_body.tex;\allowbreak{} folio\allowbreak{} 324;\allowbreak{} formula\allowbreak{} 1275\allowbreak{}}\par

\par\medskip\noindent\hypertarget{app-E7II-XVIII-039.1}{}\textbf{E7II-XVIII-039.1}\quad printed folio 325.\par
\noindent Printed reading: \emph{(empty)}\par
\noindent Corrected reading: {\ttfamily\footnotesize )\allowbreak{}}\par
\noindent Source: {\ttfamily\footnotesize ED-EN-FORMULA-G195;\allowbreak{} fr/sga7ii/\{diplomati\allowbreak{}c,current\}/expose\_XV\allowbreak{}III\_body.tex;\allowbreak{} folio\allowbreak{} 325\allowbreak{}}\par

\par\medskip\noindent\hypertarget{app-E7II-XVIII-039.2}{}\textbf{E7II-XVIII-039.2}\quad printed folio 325.\par
\noindent Printed reading: {\ttfamily\footnotesize i\allowbreak{}}\par
\noindent Corrected reading: {\ttfamily\footnotesize m))\allowbreak{} \textbackslash{}longrightarrow\allowbreak{} H\textasciicircum{}\{n-1\}(X\_\{\textbackslash{}bar\textbackslash{}eta\}\allowbreak{},\textbackslash{}mathbb\allowbreak{} Q\_\textbackslash{}ell(m\allowbreak{}}\par
\noindent Source: {\ttfamily\footnotesize fr/sga7ii/\{diplomati\allowbreak{}c,current\}/expose\_XV\allowbreak{}III\_body.tex;\allowbreak{} folio\allowbreak{} 325;\allowbreak{} formula\allowbreak{} 1284\allowbreak{}}\par

\par\medskip\noindent\hypertarget{app-E7II-XVIII-039.3}{}\textbf{E7II-XVIII-039.3}\quad printed folio 325.\par
\noindent Printed reading: \emph{(empty)}\par
\noindent Corrected reading: {\ttfamily\footnotesize \textbackslash{}check\allowbreak{}}\par
\noindent Source: {\ttfamily\footnotesize fr/sga7ii/\{diplomati\allowbreak{}c,current\}/expose\_XV\allowbreak{}III\_body.tex;\allowbreak{} folio\allowbreak{} 325;\allowbreak{} formula\allowbreak{} 1287\allowbreak{}}\par

\par\medskip\noindent\hypertarget{app-E7II-XVIII-039.4}{}\textbf{E7II-XVIII-039.4}\quad printed folio 325.\par
\noindent Printed reading: \emph{(empty)}\par
\noindent Corrected reading: {\ttfamily\footnotesize \textbackslash{}check\allowbreak{}}\par
\noindent Source: {\ttfamily\footnotesize fr/sga7ii/\{diplomati\allowbreak{}c,current\}/expose\_XV\allowbreak{}III\_body.tex;\allowbreak{} folio\allowbreak{} 325;\allowbreak{} formula\allowbreak{} 1287\allowbreak{}}\par

\par\medskip\noindent\hypertarget{app-E7II-XVIII-040.1}{}\textbf{E7II-XVIII-040.1}\quad printed folio 326.\par
\noindent Printed reading: \emph{(empty)}\par
\noindent Corrected reading: {\ttfamily\footnotesize )\allowbreak{}}\par
\noindent Source: {\ttfamily\footnotesize fr/sga7ii/\{diplomati\allowbreak{}c,current\}/expose\_XV\allowbreak{}III\_body.tex;\allowbreak{} folio\allowbreak{} 326;\allowbreak{} formula\allowbreak{} 1299\allowbreak{}}\par

\par\medskip\noindent\hypertarget{app-E7II-XVIII-040.2}{}\textbf{E7II-XVIII-040.2}\quad printed folio 326.\par
\noindent Printed reading: {\ttfamily\footnotesize \textbackslash{}overline\allowbreak{}}\par
\noindent Corrected reading: {\ttfamily\footnotesize \textbackslash{}bar\allowbreak{}}\par
\noindent Source: {\ttfamily\footnotesize ED-EN-NORM-G041;\allowbreak{} normalized-equal\allowbreak{} accent-scope\allowbreak{} audit;\allowbreak{} fr/sga7ii/\{diplomati\allowbreak{}c,current\}/expose\_XV\allowbreak{}III\_body.tex;\allowbreak{} folio\allowbreak{} 326\allowbreak{}}\par

\par\medskip\noindent\hypertarget{app-E7II-XVIII-040.3}{}\textbf{E7II-XVIII-040.3}\quad printed folio 326.\par
\noindent Printed reading: \emph{(empty)}\par
\noindent Corrected reading: {\ttfamily\footnotesize \}\allowbreak{}}\par
\noindent Source: {\ttfamily\footnotesize ED-EN-NORM-G041;\allowbreak{} normalized-equal\allowbreak{} accent-scope\allowbreak{} audit;\allowbreak{} fr/sga7ii/\{diplomati\allowbreak{}c,current\}/expose\_XV\allowbreak{}III\_body.tex;\allowbreak{} folio\allowbreak{} 326\allowbreak{}}\par

\par\medskip\noindent\hypertarget{app-E7II-XVIII-040.4}{}\textbf{E7II-XVIII-040.4}\quad printed folio 326.\par
\noindent Printed reading: {\ttfamily\footnotesize \}\allowbreak{}}\par
\noindent Corrected reading: \emph{(empty)}\par
\noindent Source: {\ttfamily\footnotesize ED-EN-NORM-G041;\allowbreak{} normalized-equal\allowbreak{} accent-scope\allowbreak{} audit;\allowbreak{} fr/sga7ii/\{diplomati\allowbreak{}c,current\}/expose\_XV\allowbreak{}III\_body.tex;\allowbreak{} folio\allowbreak{} 326\allowbreak{}}\par

\par\medskip\noindent\hypertarget{app-E7II-XVIII-040.5}{}\textbf{E7II-XVIII-040.5}\quad printed folio 326.\par
\noindent Printed reading: {\ttfamily\footnotesize \textbackslash{}overline\allowbreak{}}\par
\noindent Corrected reading: {\ttfamily\footnotesize \textbackslash{}bar\allowbreak{}}\par
\noindent Source: {\ttfamily\footnotesize ED-EN-NORM-G042;\allowbreak{} normalized-equal\allowbreak{} accent-scope\allowbreak{} audit;\allowbreak{} fr/sga7ii/\{diplomati\allowbreak{}c,current\}/expose\_XV\allowbreak{}III\_body.tex;\allowbreak{} folio\allowbreak{} 326\allowbreak{}}\par

\par\medskip\noindent\hypertarget{app-E7II-XVIII-040.6}{}\textbf{E7II-XVIII-040.6}\quad printed folio 326.\par
\noindent Printed reading: \emph{(empty)}\par
\noindent Corrected reading: {\ttfamily\footnotesize \}\allowbreak{}}\par
\noindent Source: {\ttfamily\footnotesize ED-EN-NORM-G042;\allowbreak{} normalized-equal\allowbreak{} accent-scope\allowbreak{} audit;\allowbreak{} fr/sga7ii/\{diplomati\allowbreak{}c,current\}/expose\_XV\allowbreak{}III\_body.tex;\allowbreak{} folio\allowbreak{} 326\allowbreak{}}\par

\par\medskip\noindent\hypertarget{app-E7II-XVIII-040.7}{}\textbf{E7II-XVIII-040.7}\quad printed folio 326.\par
\noindent Printed reading: {\ttfamily\footnotesize \}\allowbreak{}}\par
\noindent Corrected reading: \emph{(empty)}\par
\noindent Source: {\ttfamily\footnotesize ED-EN-NORM-G042;\allowbreak{} normalized-equal\allowbreak{} accent-scope\allowbreak{} audit;\allowbreak{} fr/sga7ii/\{diplomati\allowbreak{}c,current\}/expose\_XV\allowbreak{}III\_body.tex;\allowbreak{} folio\allowbreak{} 326\allowbreak{}}\par

\par\medskip\noindent\hypertarget{app-E7II-XVIII-040.8}{}\textbf{E7II-XVIII-040.8}\quad printed folio 326.\par
\noindent Printed reading: {\ttfamily\footnotesize \textbackslash{}chi\allowbreak{}}\par
\noindent Corrected reading: {\ttfamily\footnotesize X\allowbreak{}}\par
\noindent Source: {\ttfamily\footnotesize fr/sga7ii/\{diplomati\allowbreak{}c,current\}/expose\_XV\allowbreak{}III\_body.tex;\allowbreak{} folio\allowbreak{} 326;\allowbreak{} formula\allowbreak{} 1310\allowbreak{}}\par

\par\medskip\noindent\hypertarget{app-E7II-XVIII-040.9}{}\textbf{E7II-XVIII-040.9}\quad printed folio 326.\par
\noindent Printed reading: {\ttfamily\footnotesize x\allowbreak{}}\par
\noindent Corrected reading: {\ttfamily\footnotesize \textbackslash{}chi\allowbreak{}}\par
\noindent Source: {\ttfamily\footnotesize fr/sga7ii/\{diplomati\allowbreak{}c,current\}/expose\_XV\allowbreak{}III\_body.tex;\allowbreak{} folio\allowbreak{} 326\allowbreak{}}\par

\par\medskip\noindent\hypertarget{app-E7II-XVIII-040.10}{}\textbf{E7II-XVIII-040.10}\quad printed folio 326.\par
\noindent Printed reading: {\ttfamily\footnotesize x\allowbreak{}}\par
\noindent Corrected reading: {\ttfamily\footnotesize \textbackslash{}chi\allowbreak{}}\par
\noindent Source: {\ttfamily\footnotesize ED-EN-FORMULA-G197;\allowbreak{} fr/sga7ii/\{diplomati\allowbreak{}c,current\}/expose\_XV\allowbreak{}III\_body.tex;\allowbreak{} folio\allowbreak{} 326\allowbreak{}}\par

\par\medskip\noindent\hypertarget{app-E7II-XVIII-040.11}{}\textbf{E7II-XVIII-040.11}\quad printed folio 326.\par
\noindent Printed reading: {\ttfamily\footnotesize x\allowbreak{}}\par
\noindent Corrected reading: {\ttfamily\footnotesize \textbackslash{}chi\allowbreak{}}\par
\noindent Source: {\ttfamily\footnotesize ED-EN-FORMULA-G197;\allowbreak{} fr/sga7ii/\{diplomati\allowbreak{}c,current\}/expose\_XV\allowbreak{}III\_body.tex;\allowbreak{} folio\allowbreak{} 326\allowbreak{}}\par

\par\medskip\noindent\hypertarget{app-E7II-XVIII-040.12}{}\textbf{E7II-XVIII-040.12}\quad printed folio 327.\par
\noindent Printed reading: \emph{(empty)}\par
\noindent Corrected reading: {\ttfamily\footnotesize \_i\allowbreak{}}\par
\noindent Source: {\ttfamily\footnotesize fr/sga7ii/\{diplomati\allowbreak{}c,current\}/expose\_XV\allowbreak{}III\_body.tex;\allowbreak{} folio\allowbreak{} 327;\allowbreak{} formula\allowbreak{} 1320\allowbreak{}}\par

\par\medskip\noindent\hypertarget{app-E7II-XVIII-040.13}{}\textbf{E7II-XVIII-040.13}\quad printed folio 327.\par
\noindent Printed reading: \emph{(empty)}\par
\noindent Corrected reading: {\ttfamily\footnotesize \_i\allowbreak{}}\par
\noindent Source: {\ttfamily\footnotesize fr/sga7ii/\{diplomati\allowbreak{}c,current\}/expose\_XV\allowbreak{}III\_body.tex;\allowbreak{} folio\allowbreak{} 327;\allowbreak{} formula\allowbreak{} 1320\allowbreak{}}\par

\par\medskip\noindent\hypertarget{app-E7II-XVIII-040.14}{}\textbf{E7II-XVIII-040.14}\quad printed folio 327.\par
\noindent Printed reading: {\ttfamily\footnotesize x\allowbreak{}}\par
\noindent Corrected reading: {\ttfamily\footnotesize \textbackslash{}chi\allowbreak{}}\par
\noindent Source: {\ttfamily\footnotesize fr/sga7ii/\{diplomati\allowbreak{}c,current\}/expose\_XV\allowbreak{}III\_body.tex;\allowbreak{} folio\allowbreak{} 327;\allowbreak{} formula\allowbreak{} 1323\allowbreak{}}\par

\par\medskip\noindent\hypertarget{app-E7II-XVIII-040.15}{}\textbf{E7II-XVIII-040.15}\quad printed folio 327.\par
\noindent Printed reading: {\ttfamily\footnotesize X\allowbreak{}}\par
\noindent Corrected reading: {\ttfamily\footnotesize \textbackslash{}chi\allowbreak{}}\par
\noindent Source: {\ttfamily\footnotesize fr/sga7ii/\{diplomati\allowbreak{}c,current\}/expose\_XV\allowbreak{}III\_body.tex;\allowbreak{} folio\allowbreak{} 327;\allowbreak{} formula\allowbreak{} 1323\allowbreak{}}\par

\par\medskip\noindent\hypertarget{app-E7II-XVIII-040.16}{}\textbf{E7II-XVIII-040.16}\quad printed folio 327.\par
\noindent Printed reading: {\ttfamily\footnotesize x\allowbreak{}}\par
\noindent Corrected reading: {\ttfamily\footnotesize \textbackslash{}chi\allowbreak{}}\par
\noindent Source: {\ttfamily\footnotesize fr/sga7ii/\{diplomati\allowbreak{}c,current\}/expose\_XV\allowbreak{}III\_body.tex;\allowbreak{} folio\allowbreak{} 327;\allowbreak{} formula\allowbreak{} 1323\allowbreak{}}\par

\par\medskip\noindent\hypertarget{app-E7II-XVIII-040.17}{}\textbf{E7II-XVIII-040.17}\quad printed folio 327.\par
\noindent Printed reading: \emph{(empty)}\par
\noindent Corrected reading: {\ttfamily\footnotesize \textbackslash{}mathbb\allowbreak{} }\par
\noindent Source: {\ttfamily\footnotesize ED-EN-FONT-G005;\allowbreak{} F7II-XVIII-110.10;\allowbreak{} fr/sga7ii/\{diplomati\allowbreak{}c,current\}/expose\_XV\allowbreak{}III\_body.tex;\allowbreak{} folio\allowbreak{} 327\allowbreak{}}\par

\par\medskip\noindent\hypertarget{app-E7II-XVIII-040.18}{}\textbf{E7II-XVIII-040.18}\quad printed folio 327.\par
\noindent Printed reading: {\ttfamily\footnotesize x\allowbreak{}}\par
\noindent Corrected reading: {\ttfamily\footnotesize \textbackslash{}chi\allowbreak{}}\par
\noindent Source: {\ttfamily\footnotesize fr/sga7ii/\{diplomati\allowbreak{}c,current\}/expose\_XV\allowbreak{}III\_body.tex;\allowbreak{} folio\allowbreak{} 327\allowbreak{}}\par

\par\medskip\noindent\hypertarget{app-E7II-XVIII-040.19}{}\textbf{E7II-XVIII-040.19}\quad printed folio 327.\par
\noindent Printed reading: {\ttfamily\footnotesize x\allowbreak{}}\par
\noindent Corrected reading: {\ttfamily\footnotesize \textbackslash{}chi\allowbreak{}}\par
\noindent Source: {\ttfamily\footnotesize fr/sga7ii/\{diplomati\allowbreak{}c,current\}/expose\_XV\allowbreak{}III\_body.tex;\allowbreak{} folio\allowbreak{} 327\allowbreak{}}\par

\par\medskip\noindent\hypertarget{app-E7II-XVIII-041.1}{}\textbf{E7II-XVIII-041.1}\quad printed folio 327.\par
\noindent Printed reading: {\ttfamily\footnotesize 15\allowbreak{}}\par
\noindent Corrected reading: {\ttfamily\footnotesize 90\allowbreak{}}\par
\noindent Source: {\ttfamily\footnotesize ED-EN-NONMATH-G032;\allowbreak{} F7II-XVIII-111.2;\allowbreak{} 7II:574;\allowbreak{} 7II:977;\allowbreak{} fr/sga7ii/\{diplomati\allowbreak{}c,current\}/expose\_XV\allowbreak{}III\_body.tex;\allowbreak{} folio\allowbreak{} 327\allowbreak{}}\par

\par\endgroup\normalsize\fi
\clearpage\SGAtoc{Exposé XIX}
\selectlanguage{english}
\clearpage
\SGApage{II}{328}

\begin{center}
{\large SGA 7 --- EXPOS\'E XIX}\par
\vspace{2em}
{\Large\underline{The Noether Theorem}}\par
\vspace{1.5em}
{\large by P. Deligne}
\end{center}

\section*{Contents}

\noindent 1. Statement of the theorems. \hfill 1

\noindent 2. A complement to XVIII. \hfill 4

\noindent 3. Proof of 1.3. \hfill 5

\noindent 4. A second complement to XVIII. \hfill 8

\noindent 5. The exceptional cases of the Noether theorem. \hfill 11

\section*{1. Statement of the theorems}

The present expos\'e updates Lefschetz's proof [3, p.~108] of the Noether
theorem (1.2) stated below.

\subsubsection*{1.1.}
Let $k$ be an algebraically closed field of characteristic $p$.  Let
$d\geq1$ and let
\[
\underline a=(a_i)_{1\leq i\leq d}
\]
be a sequence of integers at least $1$.  A complete intersection $V$ of
dimension $n$ and multidegree $\underline a$ in the projective space
$\mathbb P^{n+d}$ over $k$ will be called \underline{generic} if, after an
extension of $k$, it is projectively isomorphic to a complete intersection
defined by equations whose coefficients are algebraically independent over
the prime field.  In what follows, unless expressly stated otherwise, we
suppose that $d\geq1$ and $a_i\geq2$.

\subsubsection*{Theorem 1.2.}
Let $V\subset\mathbb P^{2+d}$ be a generic complete-intersection surface of
multidegree $\underline a$.  Suppose that
\[
\underline a\ne(2),(3),(2,2).
\]
Then every (reduced) curve on $V$ is the scheme-theoretic intersection of
$V$ with a hypersurface of $\mathbb P^{2+d}$.

Let us show that the theorem follows from its corollary
\begin{equation}\tag{1.2.1}
\operatorname{Pic}(\mathbb P^{2+d})=\mathbb Z
\xrightarrow{\ \sim\ }\operatorname{Pic}(V).
\end{equation}
If $C$ is a curve on $V$, (1.2.1) implies that
$\mathcal O(C)\simeq\mathcal O(n)$ for a suitable $n$.  If $H$ is a
hyperplane section, then
\[
n=\bigl(\textstyle\prod_i a_i\bigr)^{-1}(C,H)>0.
\]
It remains to prove that the linear system induced on $V$ by $|nH|$ is
complete.  Write $V$ as the intersection of hypersurfaces $H_i$ of degree
$a_i$, and put $V_i=H_1\cap\cdots\cap H_i$.  Induction on $i$ shows that the
linear system induced by $|nH|$ on $V_i$ is complete: in the exact sequence
\[
H^0(V_i,\mathcal O(n))\longrightarrow
H^0(V_{i+1},\mathcal O(n))\longrightarrow
H^1(V_i,\mathcal O(n-a_{i+1})),
\]
the last group vanishes (FAC).

\SGApage{II}{329}
\subsubsection*{Theorem 1.3.}
Let $V\subset\mathbb P^{2n+d}$ be a generic complete intersection of
nonzero even dimension $2n$ and multidegree $\underline a$.  Suppose that
one of the following conditions holds:
\begin{enumerate}[label=(\alph*)]
\item $k$ is not of characteristic $2$, and $(2n;\underline a)$ is distinct
  from $(2n;2)$, $(2n;2,2)$, and $(2;3)$;
\item the Hodge number $h^{2n,0}$ is nonzero, i.e.
  $\sum_i(a_i-1)>2n$.
\end{enumerate}
Then every algebraic cohomology class in
$H^{2n}(V,\mathbb Q_\ell(n))$ is a rational multiple of the class
$\eta^n$ of a linear section.

\SGAvar{Let us deduce (1.2.1) from 1.3.  The exceptional cases in 1.2 are exactly
those for which $\sum_i(d_i-1)\leq2$; hence 1.3(b) applies to the surface
$V$ of 1.2.  We know that
$\operatorname{Pic}^{\tau}(V)=0$ (XI 1.8), so
$\operatorname{Pic}(V)$ injects into $H^2(V,\mathbb Q_\ell(1))$.
By 1.3 every element of $\operatorname{Pic}(V)$ is a rational multiple of
the class $\eta$ of $\mathcal O(1)$, and (XI 1.8) gives the conclusion.

}{Let us deduce (1.2.1) from 1.3.  The exceptional cases in 1.2 are exactly
those for which $\sum_i(a_i-1)\leq2$; hence 1.3(b) applies to the surface
$V$ of 1.2.  We know that
$\operatorname{Pic}^{\tau}(V)=0$ (XI 1.8), so
$\operatorname{Pic}(V)$ injects into $H^2(V,\mathbb Q_\ell(1))$.
By 1.3 every element of $\operatorname{Pic}(V)$ is a rational multiple of
the class $\eta$ of $\mathcal O(1)$, and (XI 1.8) gives the conclusion.

}{E7II-XIX-001}It is clear that the validity of 1.2 and 1.3 does not depend on the chosen
generic complete intersection, but only on $p$, $n$, and $\underline a$.

\SGApage{II}{330}
\subsubsection*{Theorem 1.4.}
Let $X\subset\mathbb P$ be a smooth projective variety, pure of dimension
$2n+1$ over $k$.  Let $\ell\subset\check{\mathbb P}$ be a Lefschetz pencil
of hyperplane sections of $X$.  If $k$ has characteristic $2$, suppose in
addition that $\ell$ is a generic Lefschetz pencil.

Suppose that $X$ satisfies condition (LV) (XVIII 5.2.2), and let $V$ be the
generic hyperplane section in the pencil $\ell$.  One of the following two
conditions holds:
\begin{enumerate}[label=(\alph*)]
\item no nonzero cohomology class in the vanishing part
  $\operatorname{Ev}(V,\mathbb Q_\ell(n))$ of
  $H^{2n}(V,\mathbb Q_\ell(n))$ is algebraic;
\item the algebraic cohomology classes generate
  $\operatorname{Ev}(V,\mathbb Q_\ell(n))$.
\end{enumerate}

Let us make precise the meaning of ``generic'' in 1.4.

\begin{enumerate}[label=(\alph*)]
\item Let $G$ be the Grassmannian of lines in the projective space
  $\check{\mathbb P}$ dual to $\mathbb P$.  By a ``geometric point'' of
  $G$ we mean a morphism from the spectrum of an algebraically closed field
  to $G$.  Every geometric point $\bar s\to G$ defines a pencil $\ell$ of
  hyperplane sections of the variety obtained from $X$ by extending scalars
  from $k$ to $k(\bar s)$.  We call $\ell$ generic if $\bar s$ lies over
  the generic point of $G$.
\item Let $\ell$ be a pencil of hyperplane sections (a line in
  $\check{\mathbb P}$).  Every geometric point $\bar s\to\ell$ defines a
  hyperplane section of the variety obtained from $X$ by extending scalars
  from $k$ to $k(\bar s)$.  We call this section generic if $\bar s$ lies
  over the generic point of $\ell$.
\end{enumerate}

Let $\bar s$ be as above, with image the generic point $s$ of $\ell$, and
let $S\subset\ell$ be the finite set of points giving singular hyperplane
sections.  The monodromy group $\pi_1(\ell-S,\bar s)$, a quotient of
$\operatorname{Gal}(k(\bar s)/k(s))$, acts on
$\operatorname{Ev}(V,\mathbb Q_\ell(n))$.  If $k$ is not of characteristic
$2$, this action is irreducible (XVIII 6.7); in Section 2 below we shall show
that the assertion remains true without a characteristic hypothesis for a
generic pencil.  The subspace of $\operatorname{Ev}(V,\mathbb Q_\ell(n))$
generated by algebraic cohomology classes is stable under monodromy (that is,
under Galois).  It is therefore either zero or the whole space, which proves
1.4.

\SGApage{II}{331}
\subsubsection*{Corollary 1.5.}
Let $V\subset\mathbb P^{2n+d}$ be a generic complete intersection of even
dimension $2n$ and multidegree $\underline a$.  Then either
\begin{enumerate}[label=(\alph*)]
\item every algebraic cohomology class in
  $H^{2n}(V,\mathbb Q_\ell(n))$ is a rational multiple of $\eta^n$; or
\item $H^{2n}(V,\mathbb Q_\ell(n))$ is generated by algebraic cohomology
  classes.
\end{enumerate}
If (b) holds, then for every complete intersection of dimension $2n$ and
multidegree $\underline a$ over an algebraically closed field of
characteristic $p$, the group $H^{2n}(V,\mathbb Q_\ell(n))$ is generated by
algebraic classes.

Apply 1.4 with $X$ a generic complete intersection of dimension $2n+1$ and
multidegree $(a_1,\ldots,a_{d-1})$, and with $\ell$ a generic pencil of
sections of $X$ by hypersurfaces of degree $a_d$ (cf. XVII 2.4).  One could
equally take $X$ of multidegree $(a_1,\ldots,a_d)$ and $\ell$ a generic
pencil of hyperplane sections.  The final assertion follows because an
algebraic cohomology class remains algebraic under specialization.



\section*{2. A complement to XVIII}

\SGAvar{\subsubsection*{Proposition 2.1.}
Let $X\subset\mathbb P$ be a smooth projective variety, pure of even
dimension $2n$ over an algebraically closed field $k$, and let
$\ell\subset\check{\mathbb P}$ be a generic Lefschetz pencil of hyperplane
sections of $X$.  Let $s$ be a geometric point of $\ell$, let $H_s$ be the
corresponding hyperplane section, let $S\subset\ell$ be the set of exceptional
values of the pencil, and let $\pi_1(\ell-S,s)$ be the monodromy group.
Suppose that $X$ satisfies condition (LV).  Then the representation of
$\pi_1(\ell-S,s)$ on the vanishing cohomology
\[
\operatorname{Ev}(H_s,\mathbb Q_\ell)
\subset H^{2n}(H_s,\mathbb Q_\ell)
\]
is absolutely irreducible.

Recall that $\operatorname{Ev}(H_s,\mathbb Q_\ell)$ is the orthogonal
complement of the image of $H^{2n}(X,\mathbb Q_\ell)$ in
$H^{2n}(H_{s,}\mathbb Q_\ell)$, and that this image equals the subgroup of
monodromy invariants (XVIII 5.6.10).

\subsubsection*{Proof.}
The projective line $\mathbb P^1$ is simply connected.  Hence
$\pi_1(\ell-S,s)$ is topologically generated by the inertia subgroups at the
points of $S$ and their conjugates.  Its image in
$\operatorname{GL}(H^{n}(H_s,\mathbb Q_\ell))$ is therefore topologically
generated by the transformations
\[
x\longmapsto x-(-1)^n(x\cdot\delta)\delta,
\]
where $\delta$ is conjugate to a vanishing cycle.
Consequently, the image of $H^{2n}(X,\mathbb Q_\ell)$ in
$H^{2n}(H_s,\mathbb Q_\ell)$ is the orthogonal complement of the subspace
generated by the vanishing cycles and their conjugates; thus
$\operatorname{Ev}(H_s,\mathbb Q_\ell)$ is generated by those cycles and
their conjugates.

}{\subsubsection*{Proposition 2.1.}
Let $X\subset\mathbb P$ be a smooth projective variety, pure of odd
dimension $2n+1$ over an algebraically closed field $k$, and let
$\ell\subset\check{\mathbb P}$ be a generic Lefschetz pencil of hyperplane
sections of $X$.  Let $s$ be a geometric point of $\ell$, let $H_s$ be the
corresponding hyperplane section, let $S\subset\ell$ be the set of exceptional
values of the pencil, and let $\pi_1(\ell-S,s)$ be the monodromy group.
Suppose that $X$ satisfies condition (LV).  Then the representation of
$\pi_1(\ell-S,s)$ on the vanishing cohomology
\[
\operatorname{Ev}(H_s,\mathbb Q_\ell)
\subset H^{2n}(H_s,\mathbb Q_\ell)
\]
is absolutely irreducible.

Recall that $\operatorname{Ev}(H_s,\mathbb Q_\ell)$ is the orthogonal
complement of the image of $H^{2n}(X,\mathbb Q_\ell)$ in
$H^{2n}(H_s,\mathbb Q_\ell)$, and that this image equals the subgroup of
monodromy invariants (XVIII 5.6.10).

\subsubsection*{Proof.}
The projective line $\mathbb P^1$ is simply connected.  Hence
$\pi_1(\ell-S,s)$ is topologically generated by the inertia subgroups at the
points of $S$ and their conjugates.  Its image in
$\operatorname{GL}(H^{2n}(H_s,\mathbb Q_\ell))$ is therefore topologically
generated by the transformations
\[
x\longmapsto x-(-1)^n(x\cdot\delta)\delta,
\]
where $\delta$ is conjugate to a vanishing cycle.
Consequently, the image of $H^{2n}(X,\mathbb Q_\ell)$ in
$H^{2n}(H_s,\mathbb Q_\ell)$ is the orthogonal complement of the subspace
generated by the vanishing cycles and their conjugates; thus
$\operatorname{Ev}(H_s,\mathbb Q_\ell)$ is generated by those cycles and
their conjugates.

}{E7II-XIX-002}\SGApage{II}{\(\approx\)332}
For a \underline{generic} Lefschetz pencil, the arguments of XVIII §6 show
that the vanishing cycles attached to the different points of $S$ are
conjugate to one another.

\SGAvar{Let $W\subset\operatorname{Ev}(H_s,\mathbb Q_\ell)$ be a nonzero invariant
subspace.  There exist a conjugate $\delta$ of a vanishing cycle and an
element $v\in W$ such that $(v,\delta)\ne0$; this is where condition (LV) is
used.  Since
\[
v-(-1)^m(v\cdot\delta)\delta\in W,
\]
we have $\delta\in W$.  The subspace $W$, containing one
vanishing cycle and all its conjugates, is therefore
$\operatorname{Ev}(H_s,\mathbb Q_\ell)$.  The argument remains valid after
extending the coefficient field $\mathbb Q_\ell$, proving absolute
irreducibility.

}{Let $W\subset\operatorname{Ev}(H_s,\mathbb Q_\ell)$ be a nonzero invariant
subspace.  There exist a conjugate $\delta$ of a vanishing cycle and an
element $v\in W$ such that $(v,\delta)\ne0$; this is where condition (LV) is
used.  Since
\[
v-(-1)^n(v\cdot\delta)\delta\in W,
\]
we have $\delta\in W$.  The subspace $W$, containing one
vanishing cycle and all its conjugates, is therefore
$\operatorname{Ev}(H_s,\mathbb Q_\ell)$.  The argument remains valid after
extending the coefficient field $\mathbb Q_\ell$, proving absolute
irreducibility.

}{E7II-XIX-003}\subsubsection*{Remark 2.2.}
If condition (LV) is omitted, the preceding arguments still show that every
invariant subspace of $H^{2n}(H_s,\mathbb Q_\ell)$ either contains
$\operatorname{Ev}(H_s,\mathbb Q_\ell)$ or is contained in the image of
$H^{2n}(X,\mathbb Q_\ell)$.  The same assertion holds for $X$ of even
dimension (and for every Lefschetz pencil).



\section*{3. Proof of 1.3}

\SGAvar{\subsubsection*{3.1.}
Under the hypotheses of the theorem, it remains to prove that 1.5(b) does
not hold.  The case $h^{0,2n}\ne0$ will be treated by Katz in XXI 4.1 (the
condition $h^{0,2n}\ne0$ is equivalent to
\[
  \sum_i d_i-(2n+r+1)\geq0,
\]
this number being the integer $m$ such that
$\Omega_V^{2n}\simeq\mathcal O(m)$).

}{\subsubsection*{3.1.}
Under the hypotheses of the theorem, it remains to prove that 1.5(b) does
not hold.  The case $h^{0,2n}\ne0$ will be treated by Katz in XXI 4.1 (the
condition $h^{0,2n}\ne0$ is equivalent to
\[
  \sum_i a_i-(2n+r+1)\geq0,
\]
this number being the integer $m$ such that
$\Omega_V^{2n}\simeq\mathcal O(m)$).

}{E7II-XIX-004}The results of the present expos\'e will not be used again in the remainder
of the seminar, so there is no vicious circle.

\SGApage{II}{\(\approx\)333}
\subsubsection*{3.2.}
When $k=\mathbb C$, one proves 1.3 by observing that if 1.5(b) holds, the
cohomology of $V$ is purely of type $(n,n)$, since algebraic cohomology is
of type $(n,n)$.  This occurs only in the exceptional cases excluded in
1.3 (XI 2.9).  By the Lefschetz principle, 1.3 is therefore true in
characteristic zero.

\subsubsection*{Proposition 3.3.}
Assume the hypotheses of 1.4, and suppose that (b) holds.  Let
$\operatorname{Ev}(V,\mathbb Q)$ be the $\mathbb Q$-vector subspace of
$\operatorname{Ev}(V,\mathbb Q_\ell(n))$ generated by the algebraic
cohomology classes.

\SGAvar{\begin{enumerate}[label=(\alph*)]
\item There is an isomorphism
  \[
  \operatorname{Ev}(V,\mathbb Q)\otimes\mathbb Q_\ell
  \xrightarrow{\ \sim\ }
  \operatorname{Ev}(V,\mathbb Q_\ell(n)).
  \]
  On $\operatorname{Ev}(V,\mathbb Q)$ the intersection form is rational
  and definite, of sign $(-1)^n$.
\item The vanishing cycles and their conjugates belong to
  $\operatorname{Ev}(V,\mathbb Q)$; there they form an irreducible root
  system of type $A$, $D$, or $E$.
\item The monodromy group (the image of $\pi_1(\ell-S,\bar s)$ in
  $\operatorname{GL}(\operatorname{Ev}(V,\mathbb Q))$) is finite; it is the
  Weil group of the root system in (b).
\end{enumerate}

}{\begin{enumerate}[label=(\alph*)]
\item There is an isomorphism
  \[
  \operatorname{Ev}(V,\mathbb Q)\otimes\mathbb Q_\ell
  \xrightarrow{\ \sim\ }
  \operatorname{Ev}(V,\mathbb Q_\ell(n)).
  \]
  On $\operatorname{Ev}(V,\mathbb Q)$ the intersection form is rational
  and definite, of sign $(-1)^n$.
\item The vanishing cycles and their conjugates belong to
  $\operatorname{Ev}(V,\mathbb Q)$; there they form an irreducible root
  system of type $A$, $D$, or $E$.
\item The monodromy group (the image of $\pi_1(\ell-S,\bar s)$ in
  $\operatorname{GL}(\operatorname{Ev}(V,\mathbb Q))$) is finite; it is the
  Weyl group of the root system in (b).
\end{enumerate}

}{E7II-XIX-005}Apart from definiteness of the intersection form, (a) follows because the
intersection form is nondegenerate on
$\operatorname{Ev}(V,\mathbb Q_\ell(n))$ (by (LV)) and rational on
$\operatorname{Ev}(V,\mathbb Q)$.

Condition 1.4(b) says equivalently that there are algebraic cycles on $V$,
modulo numerical equivalence, which are not invariant under monodromy; this
condition is therefore independent of $\ell$.  Let $x$ be such an algebraic
cycle, and let $\delta$ be the vanishing cycle associated with $s_0\in S$
and with a path $c$ from $\bar s$ (which defines $V$) to $s_0$.  For a
suitable $c$, $(x,\delta)\ne0$.  The Galois transform
\[
  \sigma(x)=x-(-1)^n(x,\delta)\delta
\]
is again algebraic; hence $(x,\delta)\delta$ is algebraic, the image of a
cycle $y$ independent of $\ell$.  The rational number
\[
  \frac{(-1)^n}{2}(y,y)=(x,\delta)^2
\]
is consequently independent of $\ell$.  It is a square in every
$\mathbb Q_\ell$, hence a square in $\mathbb Q$; thus, up to sign,
$(x,\delta)$ belongs to $\mathbb Q$ and is independent of $\ell$.  This
shows not only that $\delta\in\operatorname{Ev}(V,\mathbb Q)$, but also that
the algebraic cycle modulo numerical equivalence
\[
  \pm\delta=\frac{\sigma(x)-x}{(x,\delta)}
\]
is independent of $\ell$.

\SGApage{II}{\(\approx\)334}
\SGAvar{An algebraic cycle on a variety $V$ is always defined over a finite
extension of any field of definition of $V$.  It follows that
$\pi_1(\ell-S,s)$ acts on $\operatorname{Ev}(V,\mathbb Q)$ through a finite
group.  This group is generated by the reflections associated with the
vanishing cycles $\delta$ and their conjugates (XVIII 6.6.1 and Section 2).
The cycles $\delta$ and their conjugates therefore form a root system whose
Weil group is the monodromy group.  Moreover, all roots of this root system
are conjugate under the Weil group; it is of type $A$, $D$, or $E$.

Up to a scalar, the intersection form is the unique invariant form under
the Weil group.  It is therefore definite.  Its sign is determined by
computing
\[
  (\delta,\delta)=2(-1)^n.
\]

}{An algebraic cycle on a variety $V$ is always defined over a finite
extension of any field of definition of $V$.  It follows that
$\pi_1(\ell-S,s)$ acts on $\operatorname{Ev}(V,\mathbb Q)$ through a finite
group.  This group is generated by the reflections associated with the
vanishing cycles $\delta$ and their conjugates (XVIII 6.6.1 and Section 2).
The cycles $\delta$ and their conjugates therefore form a root system whose
Weyl group is the monodromy group.  Moreover, all roots of this root system
are conjugate under the Weyl group; it is of type $A$, $D$, or $E$.

Up to a scalar, the intersection form is the unique invariant form under
the Weyl group.  It is therefore definite.  Its sign is determined by
computing
\[
  (\delta,\delta)=2(-1)^n.
\]

}{E7II-XIX-006}\subsubsection*{Proposition 3.4.}
If $k=\mathbb C$, then under the hypotheses of 1.4 the following conditions
are equivalent:
\begin{enumerate}[label=(\roman*)]
\item the monodromy group is finite;
\item the vanishing cohomology $\operatorname{Ev}(V)$ is purely of type
  $(n,n)$.
\end{enumerate}

Let $\widetilde X\to\ell$ be the family whose fibers are the hyperplane
sections in the pencil $\ell$, and let $V_s$ be the fiber at
$s\in\ell(\mathbb C)$.  The $\operatorname{Ev}(V_s)$ form an algebraic
variation of Hodge structures on $\ell$.  By the fixed-part theorem
[2, 4.1.2], applied to a finite covering of $\ell$ (cf. [2, 4.1.3.3]), (i)
implies that the Hodge structure on $\operatorname{Ev}(V_s)$ is locally
constant.  The monodromy group therefore acts by automorphisms of the Hodge
structure.  For a Picard--Lefschetz transformation $T$,
$\operatorname{Im}(T-1)$ is a rank-one Hodge substructure, hence is of type
$(n,n)$ and is generated by a vanishing cycle $\delta$.  Since the
vanishing cycles are of type $(n,n)$, all of
$\operatorname{Ev}(V_s)$ is of that type.

Conversely, if $\operatorname{Ev}(V)$ is purely of type $(n,n)$, the
intersection form on $\operatorname{Ev}(V)$ is definite and the monodromy
group is finite (cf. [2, 4.1.3.3]).

\SGApage{II}{335}
\SGAvar{\subsubsection*{Remark 3.5.}
If the equivalent conditions of 3.3 hold, the argument of 3.3 shows that
the vanishing cycles and their conjugates form a root system of type $A$,
$D$, or $E$, whose Weil group is the monodromy group.

}{\subsubsection*{Remark 3.5.}
If the equivalent conditions of 3.4 hold, the argument of 3.3 shows that
the vanishing cycles and their conjugates form a root system of type $A$,
$D$, or $E$, whose Weyl group is the monodromy group.

}{E7II-XIX-007}\subsubsection*{3.6.}
Let us prove 1.3 when $k$ has finite characteristic different from $2$.
Let $S=\operatorname{Spec}(\mathcal V)$ be the spectrum of a discrete
valuation ring of mixed characteristic with residue field $k$, and let
$X\subset\mathbb P_S^{2n+d}$ be a complete intersection of relative
dimension $2n+1$ and multidegree $(a_1,\ldots,a_{d-1})$ over $S$.  Let
$\ell$ be a Lefschetz pencil of sections of $X$ by hypersurfaces of degree
$a_d$ (in the special fiber, $\ell_s$ is a Lefschetz pencil for $X_s$).
Let
\[
  f:\widetilde X\longrightarrow\ell\simeq\mathbb P^1_S
\]
be the family whose fibers are the sections in the pencil.  The exceptional
locus $T\subset\ell$ is \`etale over $S$ (XVII 6; here one uses
$p\ne2$).  The sheaf $R^{2n}f_*\mathbb Q_\ell$ is locally constant on
$\ell-T$; it is therefore tamely ramified, and its monodromy is the same in
characteristic $p$ as in characteristic zero.  If 1.5(b) holds, it follows
from 3.3(c) and 3.4 that the cohomology group $H^{2n}(V)$ of complex
complete intersections of dimension $2n$ and multidegree
$\underline a$ is purely of type $(n,n)$.  This occurs only in the cases
excluded by 1.3 (XI 2.9).



\section*{4. A second complement to XVIII}

In this section we work over $\mathbb C$ and with integral, classical
cohomology.  Using methods closer to Lefschetz's, we revisit some of the
results of XVIII.  There they were proved only after tensoring with
$\mathbb Q$ (or rather with $\mathbb Q_\ell$, which amounts to the same
thing).

\subsubsection*{4.1.}
Let $X\subset\mathbb P$ be a smooth complex projective variety, pure of
dimension $n+1$, and let $\ell$ be a Lefschetz pencil of hyperplane
sections of $X$, with axis $A$.  We also write $\ell$ for the parameter
space of the pencil, which is isomorphic to $\mathbb P^1$.  Let
$S\subset\ell$ be the set of exceptional values.  For $s\in\ell$, write
$H_s$ for the corresponding hyperplane section of $X$ (of dimension $n$).

\SGApage{II}{\(\approx\)336}
Let $T(A)$ be a closed tubular neighborhood of $A\cap X$ in $X$, whose
boundary is transverse to every $H_s$.

\SGAvar{Choose two points $0,\infty\in\ell-S$, and let
$\alpha:\ell-S\to\mathbb P^1(\mathbb C)$ be an orientation-preserving
diffeomorphism that sends $0$ to $0$, $\infty$ to $\infty$, and $S$ to the
circle of radius one centered at $0$.

}{Choose two points $0,\infty\in\ell-S$, and let
$\alpha:\ell\to\mathbb P^1(\mathbb C)$ be an orientation-preserving
diffeomorphism that sends $0$ to $0$, $\infty$ to $\infty$, and $S$ to the
circle of radius one centered at $0$.

}{E7II-XIX-008}The composite $\alpha f$ of $\alpha$ with the projection to $\ell$ makes
$X-T(A)^\circ$ a singular fibration in manifolds with boundary over
$\mathbb P^1(\mathbb C)$ (one singularity above each point of
$\alpha(S)$).  Write $c_s$ for the straight segment joining $0$ to
$\alpha(s)$.

Suppressing $\alpha$ from the notation, the picture is
\begin{center}
\begin{tikzpicture}[x=1.2cm,y=1.2cm]
  \coordinate (O) at (0,0);
  \coordinate (L) at (-1.8,1.8);
  \coordinate (S2) at (0,2.8);
  \coordinate (S1) at (2.2,2.0);
  \draw (O)--(L);
  \draw (O)--(S2);
  \draw (O)--(S1);
  \node[right] at (O) {$0$};
  \node[above right] at (S2) {$s_2$};
  \node[above right] at (S1) {$s_1$};
  \node[right] at (0,1.65) {$c_2$};
  \node[below right] at (1.25,1.14) {$c_1$};
  \fill (-0.73,0.76) circle (0.035);
  \fill (-0.55,0.91) circle (0.035);
  \fill (-0.37,1.03) circle (0.035);
\end{tikzpicture}
\end{center}

Near $s\in S$ and the singular point $x_s$ of $H_s\subset X$, suitable
local coordinates put us in the situation studied in XIV 3.2.  If
$\alpha$ is chosen holomorphic near $s$, one can even choose local
coordinates on $X$ and $\ell$ in which $f$ is written
$\sum z_i^2$ and $\alpha^{-1}(c_s)$ is given by ``the $z_i$ are real and
positive.''  Above the intersection of a neighborhood of $s$ with
$\alpha^{-1}(c_s)-\{s\}$, consider the bundle of real spheres
$\Delta'_s\subset X$ defined by ``the $z_i$ are real.''  As $t\to s$, the
fiber $\Delta'_{st}$ of $\Delta'_s$ at $t$ tends to $x_s$.

Let $\Delta_s$ denote the union of $\{x_s\}$ and a sphere bundle over
$\alpha^{-1}(c_s)-\{s\}$, contained in $X-T(A)$ and extending
$\Delta'_s$.  Then $\Delta_s$ is an embedded disk in $X$ whose boundary
$H_0\cap\Delta_s$ is a sphere representing the vanishing cycle $\delta_s$
at $s$, transported to $H_0$ along the path $\alpha^{-1}(c_s)$.

Apply Morse theory to the function $|\alpha f|^2$ to study how the
homotopy type of
\[
 \{x\in X\mid x\in T(A)\ \text{or}\ |\alpha f(x)|^2\leq r\}
\]
changes as $r$ increases.  Put
\[
 X^- = \{x\in X\mid x\in T(A)\ \text{or}\ \alpha f(x)\ne\infty\}.
\]
One finds that the inclusions
\[
 H_0\cup\bigcup_{s\in S}\Delta_s
 \hookrightarrow
 T(A)\cup H_0\cup\bigcup_{s\in S}\Delta_s
 \hookrightarrow X^-
\]
are homotopy equivalences.  Morse theory also shows that $X-X^-$ can be
deformation-retracted onto a subcomplex of real dimension at most $n$.
Altogether this gives the following result.

\SGApage{II}{337}
\subsubsection*{Theorem 4.2.}
There exists a compact subcomplex $\Sigma$ of real dimension at most $n$ in
$H_\infty^{\mathrm{aff}}=H_\infty-A$ such that the inclusion
\[
 H_0\cup\bigcup_{s\in S}\Delta_s\longrightarrow X-\Sigma
\]
is a homotopy equivalence.

\SGAvar{\subsubsection*{4.3.}
Removing $\Sigma$ has no effect on the cohomology, homology, or homotopy
groups of $X$ in degrees $i\leq n$.  Since
$H_0\cup\bigcup_{s\in S}\Delta_s$ is obtained from $H_0$ by attaching
$(n+1)$-disks along spheres representing the vanishing cycles, homology
gives
\[
 \mathbb Z^S\longrightarrow H_n(V_0)\longrightarrow H_n(X)
 \longrightarrow0.
\]
The analogous exact sequence in cohomology, together with the sequence
deduced from it by Poincar\'e duality, forms the commutative diagram
\begin{equation}\tag{4.3.1}
\begin{tikzcd}[row sep=large,column sep=large]
{} \& \mathbb Z^S \& {} \& {} \\
\mathbb Z^S \arrow[r,"{\sum n_s\delta_s}"] \&
H^n(V_0) \arrow[u,"{(\ ,\delta_s)}"'] \arrow[r] \&
H^{n+2}(X) \arrow[r] \& 0 \\
{} \& H^n(X) \arrow[u] \arrow[ur,"{\eta\Lambda}"'] \& {} \& {} \\
{} \& 0 \arrow[u] \& {} \& {}
\end{tikzcd}
\end{equation}


}{\subsubsection*{4.3.}
Removing $\Sigma$ has no effect on the cohomology, homology, or homotopy
groups of $X$ in degrees $i\leq n$.  Since
$H_0\cup\bigcup_{s\in S}\Delta_s$ is obtained from $H_0$ by attaching
$(n+1)$-disks along spheres representing the vanishing cycles, homology
gives
\[
 \mathbb Z^S\longrightarrow H_n(H_0)\longrightarrow H_n(X)
 \longrightarrow0.
\]
The analogous exact sequence in cohomology, together with the sequence
deduced from it by Poincar\'e duality, forms the commutative diagram
\begin{equation}\tag{4.3.1}
\begin{tikzcd}[row sep=large,column sep=large]
{} \& \mathbb Z^S \& {} \& {} \\
\mathbb Z^S \arrow[r,"{\sum n_s\delta_s}"] \&
H^n(H_0) \arrow[u,"{(\ ,\delta_s)}"'] \arrow[r] \&
H^{n+2}(X) \arrow[r] \& 0 \\
{} \& H^n(X) \arrow[u] \arrow[ur,"{\eta\wedge}"'] \& {} \& {} \\
{} \& 0 \arrow[u] \& {} \& {}
\end{tikzcd}
\end{equation}


}{E7II-XIX-009}\section*{5. The exceptional cases of the Noether theorem}

In this section we continue to work over $\mathbb C$, with integral
cohomology.

\SGApage{II}{\(\approx\)338}
\subsubsection*{5.1.}
Let $2n$, $d$, and $\underline a=(a_1,\ldots,a_d)$, with $d\geq1$ and
$a_i\geq2$, be such that the middle cohomology of smooth complete
intersections $V$ of dimension $2n$ and multidegree $\underline a$ is
purely of type $(n,n)$.  In other words, either $n=0$, or
$(2n;\underline a)$ is one of $(2n;2)$, $(2n;2,2)$, and $(2;3)$
(XI 2.9).

Suppose that $V$ is obtained from a complete intersection
$X\subset\mathbb P^{2n+d}$ (respectively,
$X\subset\mathbb P^{2n+d+1}$) of dimension $2n+1$ and multidegree
$(a_1,\ldots,a_{d-1})$ (respectively, $(a_1,\ldots,a_d)$) by taking a
section by a hypersurface of degree $a_d$ (respectively, a hyperplane
section) belonging to a Lefschetz pencil $\ell$.

By (4.3.1), the vanishing part $\operatorname{Ev}(V)$ of $H^{2n}(V)$,
which equals its primitive part, is generated over $\mathbb Z$ by the
vanishing cycles.  It is the orthogonal complement of the power $\eta^n$
of the cohomology class $\eta$ of a hyperplane section.

\SGAvar{Since the intersection form on $H^{2n}(V)$ is unimodular and
$\eta^{2n}$ is not $b y$, with $b>1$, for
another cohomology class $y$ (XI 1.6(iv)), the
discriminant of the intersection form on $\operatorname{Ev}(V)$ is
\[
 (\eta^n,\eta^n)=\prod_i a_i.
\]

}{Since the intersection form on $H^{2n}(V)$ is unimodular and
$\eta^n$ is not $b y$, with $b>1$, for
another cohomology class $y$ (XI 1.6(iv)), the
discriminant of the intersection form on $\operatorname{Ev}(V)$ is
\[
 (\eta^n,\eta^n)=\prod_i a_i.
\]

}{E7II-XIX-010}Let $R$ be the root system in $H^{2n}(V)$ formed by the vanishing cycles
and their conjugates.  We have the following information about $R$:
\begin{enumerate}[label=(\alph*)]
\item $R$ is empty or is of type $A$, $D$, or $E$;
\item the rank of $R$ is given in XI 1.10;
\item its discriminant is $\prod_i a_i$.  In the tables, this discriminant
  is the order of the center of the corresponding simply connected complex
  simple group.  The argument by which we computed the discriminant shows
  that this center, dual to the quotient of the weight lattice by the root
  lattice, is cyclic.
\end{enumerate}

\SGApage{II}{\(\approx\)339}
In every case this information determines $R$ completely if $n\ne0$.  One
finds:

\subsubsection*{Proposition 5.2.}
\begin{enumerate}[label=(\roman*)]
\item If $V$ is an even-dimensional quadric, $R$ is of type $A_1$.
\item If $V$ is a cubic surface, $R$ is of type $E_6$.
\item If $V$ has dimension $2n$ and is the intersection of two quadrics,
  $R$ is of type $D_{2n+3}$.
\end{enumerate}

\SGAvar{\subsubsection*{5.3.}
Suppose that $n\ne0$.  The set $R^1$ of vectors of length $\pm2$ in
$\operatorname{Ev}(V)$ is also a root system and has the same properties
(a), (b), and (c) as $R$.  Hence $R=R^1$.  Let $\pi_1$ be the fundamental
group of the parameter space of all complete intersections
$V\subset\mathbb P^{2n+d}$ of the type under consideration.  The image of
$\pi_1$ in $\operatorname{GL}(\operatorname{Ev}(V))$ is the Weil group of
$R$.  In every case it is the automorphism group of the triple
\[
 \bigl(H^{2n}(V),\eta^n,\text{intersection form}\bigr).
\]

\subsubsection*{5.4.}
With $R$ replaced by $R^1$ from 5.3, assertion 4.2(ii) is well known [1],
and for small values of $n$, 4.2(iii) was known, at least to Manin
($n=0$ is trivial; a surface that is the intersection of two quadrics is a
del Pezzo surface).  A detailed study of the intersection of two quadrics was
made by M. Reid [4].  He deduces 4.2(iii) from a study of the geometry of
the $n$-dimensional linear subspaces on an intersection of two quadrics
$V\subset\mathbb P^{2n+2}$.

}{\subsubsection*{5.3.}
Suppose that $n\ne0$.  The set $R^1$ of vectors of length $\pm2$ in
$\operatorname{Ev}(V)$ is also a root system and has the same properties
(a), (b), and (c) as $R$.  Hence $R=R^1$.  Let $\pi_1$ be the fundamental
group of the parameter space of all complete intersections
$V\subset\mathbb P^{2n+d}$ of the type under consideration.  The image of
$\pi_1$ in $\operatorname{GL}(\operatorname{Ev}(V))$ is the Weyl group of
$R$.  In every case it is the automorphism group of the triple
\[
 \bigl(H^{2n}(V),\eta^n,\text{intersection form}\bigr).
\]

\subsubsection*{5.4.}
With $R$ replaced by $R^1$ from 5.3, assertion 5.2(ii) is well known [1],
and for small values of $n$, 5.2(iii) was known, at least to Manin
($n=0$ is trivial; a surface that is the intersection of two quadrics is a
del Pezzo surface).  A detailed study of the intersection of two quadrics was
made by M. Reid [4].  He deduces 5.2(iii) from a study of the geometry of
the $n$-dimensional linear subspaces on an intersection of two quadrics
$V\subset\mathbb P^{2n+2}$.

}{E7II-XIX-011}\section*{Bibliography}

\begin{enumerate}[label={[\arabic*]}]
\item E. Cartan, ``Quelques remarques sur les 28 bitangentes d'une quartique
  plane et les 27 droites d'une surface cubique,'' \emph{Bull. Sci. Math.}
  \textbf{70} (1946), 42--45; \emph{Oeuvres}, part I, vol. 2, p.~1353.
\item P. Deligne, ``Th\'eorie de Hodge II,'' \emph{Publ. Math. IHES}
  \textbf{40} (1971), 5--58.
\item S. Lefschetz, \emph{L'analysis situs et la g\'eom\'etrie alg\'ebrique},
  Gauthier--Villars, 1924 (reprinted in \emph{Selected Papers}, Chelsea).
\item M. Reid, Ph.D. thesis, Cambridge, 1972.
\item J.-P. Serre, ``Faisceaux alg\'ebriques coh\'erents,''
  \emph{Ann. of Math.} \textbf{61} (1955), 197--278 (cited as FAC).
\end{enumerate}


\markboth{Exposé XIX}{Exposé XIX}

\ifnum\SGAform=1
\clearpage\section*{Critical apparatus}\begingroup\small\raggedright\sloppy
Readings are given in the exact TeX notation of the two English witnesses. Decimal suffixes identify the individual differences in each marked passage.

\par\medskip\noindent\hypertarget{app-E7II-XIX-001.1}{}\textbf{E7II-XIX-001.1}\quad printed folio 329.\par
\noindent Printed reading: {\ttfamily\footnotesize d\allowbreak{}}\par
\noindent Corrected reading: {\ttfamily\footnotesize a\allowbreak{}}\par
\noindent Source: {\ttfamily\footnotesize fr/sga7ii/\{diplomati\allowbreak{}c,current\}/expose\_XI\allowbreak{}X\_body.tex;\allowbreak{} folio\allowbreak{} 329\allowbreak{}}\par

\par\medskip\noindent\hypertarget{app-E7II-XIX-002.1}{}\textbf{E7II-XIX-002.1}\quad printed folio 331.\par
\noindent Printed reading: {\ttfamily\footnotesize even\allowbreak{}}\par
\noindent Corrected reading: {\ttfamily\footnotesize odd\allowbreak{}}\par
\noindent Source: {\ttfamily\footnotesize ED-EN-NONMATH-G049;\allowbreak{} F7II-XIX-003.1;\allowbreak{} 7II:576;\allowbreak{} 7II:980;\allowbreak{} fr/sga7ii/\{diplomati\allowbreak{}c,current\}/expose\_XI\allowbreak{}X\_body.tex;\allowbreak{} folio\allowbreak{} 331\allowbreak{}}\par

\par\medskip\noindent\hypertarget{app-E7II-XIX-002.2}{}\textbf{E7II-XIX-002.2}\quad printed folio 331.\par
\noindent Printed reading: \emph{(empty)}\par
\noindent Corrected reading: {\ttfamily\footnotesize +1\allowbreak{}}\par
\noindent Source: {\ttfamily\footnotesize fr/sga7ii/\{diplomati\allowbreak{}c,current\}/expose\_XI\allowbreak{}X\_body.tex;\allowbreak{} folio\allowbreak{} 331;\allowbreak{} formula\allowbreak{} 126\allowbreak{}}\par

\par\medskip\noindent\hypertarget{app-E7II-XIX-002.3}{}\textbf{E7II-XIX-002.3}\quad printed folio 331.\par
\noindent Printed reading: {\ttfamily\footnotesize \{\allowbreak{}}\par
\noindent Corrected reading: \emph{(empty)}\par
\noindent Source: {\ttfamily\footnotesize ED-EN-NORM-G038;\allowbreak{} normalized-equal\allowbreak{} subscript-scope\allowbreak{} audit;\allowbreak{} fr/sga7ii/\{diplomati\allowbreak{}c,current\}/expose\_XI\allowbreak{}X\_body.tex;\allowbreak{} folio\allowbreak{} 331\allowbreak{}}\par

\par\medskip\noindent\hypertarget{app-E7II-XIX-002.4}{}\textbf{E7II-XIX-002.4}\quad printed folio 331.\par
\noindent Printed reading: {\ttfamily\footnotesize \}\allowbreak{}}\par
\noindent Corrected reading: \emph{(empty)}\par
\noindent Source: {\ttfamily\footnotesize ED-EN-NORM-G038;\allowbreak{} normalized-equal\allowbreak{} subscript-scope\allowbreak{} audit;\allowbreak{} fr/sga7ii/\{diplomati\allowbreak{}c,current\}/expose\_XI\allowbreak{}X\_body.tex;\allowbreak{} folio\allowbreak{} 331\allowbreak{}}\par

\par\medskip\noindent\hypertarget{app-E7II-XIX-002.5}{}\textbf{E7II-XIX-002.5}\quad printed folio 332.\par
\noindent Printed reading: {\ttfamily\footnotesize n\allowbreak{}}\par
\noindent Corrected reading: {\ttfamily\footnotesize 2n\allowbreak{}}\par
\noindent Source: {\ttfamily\footnotesize fr/sga7ii/\{diplomati\allowbreak{}c,current\}/expose\_XI\allowbreak{}X\_body.tex;\allowbreak{} folio\allowbreak{} 332;\allowbreak{} formula\allowbreak{} 144\allowbreak{}}\par

\par\medskip\noindent\hypertarget{app-E7II-XIX-003.1}{}\textbf{E7II-XIX-003.1}\quad printed folio 332.\par
\noindent Printed reading: {\ttfamily\footnotesize m\allowbreak{}}\par
\noindent Corrected reading: {\ttfamily\footnotesize n\allowbreak{}}\par
\noindent Source: {\ttfamily\footnotesize fr/sga7ii/\{diplomati\allowbreak{}c,current\}/expose\_XI\allowbreak{}X\_body.tex;\allowbreak{} folio\allowbreak{} 332;\allowbreak{} formula\allowbreak{} 157\allowbreak{}}\par

\par\medskip\noindent\hypertarget{app-E7II-XIX-004.1}{}\textbf{E7II-XIX-004.1}\quad printed folio 332.\par
\noindent Printed reading: {\ttfamily\footnotesize d\allowbreak{}}\par
\noindent Corrected reading: {\ttfamily\footnotesize a\allowbreak{}}\par
\noindent Source: {\ttfamily\footnotesize ED-EN-FORMULA-G205;\allowbreak{} fr/sga7ii/\{diplomati\allowbreak{}c,current\}/expose\_XI\allowbreak{}X\_body.tex;\allowbreak{} folio\allowbreak{} 332\allowbreak{}}\par

\par\medskip\noindent\hypertarget{app-E7II-XIX-005.1}{}\textbf{E7II-XIX-005.1}\quad printed folio 333.\par
\noindent Printed reading: {\ttfamily\footnotesize Weil\allowbreak{}}\par
\noindent Corrected reading: {\ttfamily\footnotesize Weyl\allowbreak{}}\par
\noindent Source: {\ttfamily\footnotesize ED-EN-NONMATH-G006;\allowbreak{} F7II-XIX-009.3;\allowbreak{} ADOPT-F7II-XIX-009.3\allowbreak{};\allowbreak{} fr/sga7ii/\{diplomati\allowbreak{}c,current\}/expose\_XI\allowbreak{}X\_body.tex;\allowbreak{} folio\allowbreak{} 333\allowbreak{}}\par

\par\medskip\noindent\hypertarget{app-E7II-XIX-006.1}{}\textbf{E7II-XIX-006.1}\quad printed folio 334.\par
\noindent Printed reading: {\ttfamily\footnotesize Weil\allowbreak{}}\par
\noindent Corrected reading: {\ttfamily\footnotesize Weyl\allowbreak{}}\par
\noindent Source: {\ttfamily\footnotesize ED-EN-NONMATH-G007;\allowbreak{} F7II-XIX-010.4;\allowbreak{} WITNESS-F7II-XIX-010\allowbreak{}.4;\allowbreak{} fr/sga7ii/\{diplomati\allowbreak{}c,current\}/expose\_XI\allowbreak{}X\_body.tex;\allowbreak{} folio\allowbreak{} 334\allowbreak{}}\par

\par\medskip\noindent\hypertarget{app-E7II-XIX-006.2}{}\textbf{E7II-XIX-006.2}\quad printed folio 334.\par
\noindent Printed reading: {\ttfamily\footnotesize Weil\allowbreak{}}\par
\noindent Corrected reading: {\ttfamily\footnotesize Weyl\allowbreak{}}\par
\noindent Source: {\ttfamily\footnotesize ED-EN-NONMATH-G008;\allowbreak{} F7II-XIX-010.5;\allowbreak{} ADOPT-F7II-XIX-010.5\allowbreak{};\allowbreak{} fr/sga7ii/\{diplomati\allowbreak{}c,current\}/expose\_XI\allowbreak{}X\_body.tex;\allowbreak{} folio\allowbreak{} 334\allowbreak{}}\par

\par\medskip\noindent\hypertarget{app-E7II-XIX-006.3}{}\textbf{E7II-XIX-006.3}\quad printed folio 334.\par
\noindent Printed reading: {\ttfamily\footnotesize Weil\allowbreak{}}\par
\noindent Corrected reading: {\ttfamily\footnotesize Weyl\allowbreak{}}\par
\noindent Source: {\ttfamily\footnotesize ED-EN-NONMATH-G009;\allowbreak{} F7II-XIX-010.6;\allowbreak{} ADOPT-F7II-XIX-010.6\allowbreak{};\allowbreak{} fr/sga7ii/\{diplomati\allowbreak{}c,current\}/expose\_XI\allowbreak{}X\_body.tex;\allowbreak{} folio\allowbreak{} 334\allowbreak{}}\par

\par\medskip\noindent\hypertarget{app-E7II-XIX-007.1}{}\textbf{E7II-XIX-007.1}\quad printed folio 335.\par
\noindent Printed reading: {\ttfamily\footnotesize 3\allowbreak{}}\par
\noindent Corrected reading: {\ttfamily\footnotesize 4\allowbreak{}}\par
\noindent Source: {\ttfamily\footnotesize port/PORT\_EN7II\_LOG.\allowbreak{}jsonl\#C763;\allowbreak{} fixes/FIXLOG\_7II\_CUR\allowbreak{}RENT.jsonl\#C763;\allowbreak{} 7II:588\allowbreak{}}\par

\par\medskip\noindent\hypertarget{app-E7II-XIX-007.2}{}\textbf{E7II-XIX-007.2}\quad printed folio 335.\par
\noindent Printed reading: {\ttfamily\footnotesize Weil\allowbreak{}}\par
\noindent Corrected reading: {\ttfamily\footnotesize Weyl\allowbreak{}}\par
\noindent Source: {\ttfamily\footnotesize ED-EN-NONMATH-G010;\allowbreak{} F7II-XIX-011.2;\allowbreak{} ADOPT-F7II-XIX-011.2\allowbreak{};\allowbreak{} fr/sga7ii/\{diplomati\allowbreak{}c,current\}/expose\_XI\allowbreak{}X\_body.tex;\allowbreak{} folio\allowbreak{} 335\allowbreak{}}\par

\par\medskip\noindent\hypertarget{app-E7II-XIX-008.1}{}\textbf{E7II-XIX-008.1}\quad printed folio 336.\par
\noindent Printed reading: {\ttfamily\footnotesize -S\allowbreak{}}\par
\noindent Corrected reading: \emph{(empty)}\par
\noindent Source: {\ttfamily\footnotesize fr/sga7ii/\{diplomati\allowbreak{}c,current\}/expose\_XI\allowbreak{}X\_body.tex;\allowbreak{} folio\allowbreak{} 336\allowbreak{}}\par

\par\medskip\noindent\hypertarget{app-E7II-XIX-009.1}{}\textbf{E7II-XIX-009.1}\quad printed folio 337.\par
\noindent Printed reading: {\ttfamily\footnotesize V\allowbreak{}}\par
\noindent Corrected reading: {\ttfamily\footnotesize H\allowbreak{}}\par
\noindent Source: {\ttfamily\footnotesize fr/sga7ii/\{diplomati\allowbreak{}c,current\}/expose\_XI\allowbreak{}X\_body.tex;\allowbreak{} folio\allowbreak{} 337;\allowbreak{} formula\allowbreak{} 375\allowbreak{}}\par

\par\medskip\noindent\hypertarget{app-E7II-XIX-009.2}{}\textbf{E7II-XIX-009.2}\quad printed folio 337.\par
\noindent Printed reading: {\ttfamily\footnotesize V\allowbreak{}}\par
\noindent Corrected reading: {\ttfamily\footnotesize H\allowbreak{}}\par
\noindent Source: {\ttfamily\footnotesize fr/sga7ii/\{diplomati\allowbreak{}c,current\}/expose\_XI\allowbreak{}X\_body.tex;\allowbreak{} folio\allowbreak{} 337;\allowbreak{} formula\allowbreak{} 376\allowbreak{}}\par

\par\medskip\noindent\hypertarget{app-E7II-XIX-009.3}{}\textbf{E7II-XIX-009.3}\quad printed folio 337.\par
\noindent Printed reading: {\ttfamily\footnotesize \textbackslash{}Lambda\allowbreak{}}\par
\noindent Corrected reading: {\ttfamily\footnotesize \textbackslash{}wedge\allowbreak{}}\par
\noindent Source: {\ttfamily\footnotesize fr/sga7ii/\{diplomati\allowbreak{}c,current\}/expose\_XI\allowbreak{}X\_body.tex;\allowbreak{} folio\allowbreak{} 337;\allowbreak{} formula\allowbreak{} 376\allowbreak{}}\par

\par\medskip\noindent\hypertarget{app-E7II-XIX-010.1}{}\textbf{E7II-XIX-010.1}\quad printed folio 338.\par
\noindent Printed reading: {\ttfamily\footnotesize \{2n\}\allowbreak{}}\par
\noindent Corrected reading: {\ttfamily\footnotesize n\allowbreak{}}\par
\noindent Source: {\ttfamily\footnotesize fr/sga7ii/\{diplomati\allowbreak{}c,current\}/expose\_XI\allowbreak{}X\_body.tex;\allowbreak{} folio\allowbreak{} 338;\allowbreak{} formula\allowbreak{} 440\allowbreak{}}\par

\par\medskip\noindent\hypertarget{app-E7II-XIX-011.1}{}\textbf{E7II-XIX-011.1}\quad printed folio 339.\par
\noindent Printed reading: {\ttfamily\footnotesize Weil\allowbreak{}}\par
\noindent Corrected reading: {\ttfamily\footnotesize Weyl\allowbreak{}}\par
\noindent Source: {\ttfamily\footnotesize ED-EN-NONMATH-G011;\allowbreak{} F7II-XIX-018.2;\allowbreak{} 7II:1001;\allowbreak{} fr/sga7ii/\{diplomati\allowbreak{}c,current\}/expose\_XI\allowbreak{}X\_body.tex;\allowbreak{} folio\allowbreak{} 339\allowbreak{}}\par

\par\medskip\noindent\hypertarget{app-E7II-XIX-011.2}{}\textbf{E7II-XIX-011.2}\quad printed folio 339.\par
\noindent Printed reading: {\ttfamily\footnotesize 4\allowbreak{}}\par
\noindent Corrected reading: {\ttfamily\footnotesize 5\allowbreak{}}\par
\noindent Source: {\ttfamily\footnotesize ED-EN-NONMATH-G033;\allowbreak{} F7II-XIX-018.3;\allowbreak{} 7II:597;\allowbreak{} 7II:1002;\allowbreak{} fr/sga7ii/\{diplomati\allowbreak{}c,current\}/expose\_XI\allowbreak{}X\_body.tex;\allowbreak{} folio\allowbreak{} 339\allowbreak{}}\par

\par\medskip\noindent\hypertarget{app-E7II-XIX-011.3}{}\textbf{E7II-XIX-011.3}\quad printed folio 339.\par
\noindent Printed reading: {\ttfamily\footnotesize 4\allowbreak{}}\par
\noindent Corrected reading: {\ttfamily\footnotesize 5\allowbreak{}}\par
\noindent Source: {\ttfamily\footnotesize ED-EN-NONMATH-G034;\allowbreak{} F7II-XIX-018.4;\allowbreak{} 7II:598;\allowbreak{} 7II:1003;\allowbreak{} fr/sga7ii/\{diplomati\allowbreak{}c,current\}/expose\_XI\allowbreak{}X\_body.tex;\allowbreak{} folio\allowbreak{} 339\allowbreak{}}\par

\par\medskip\noindent\hypertarget{app-E7II-XIX-011.4}{}\textbf{E7II-XIX-011.4}\quad printed folio 339.\par
\noindent Printed reading: {\ttfamily\footnotesize 4\allowbreak{}}\par
\noindent Corrected reading: {\ttfamily\footnotesize 5\allowbreak{}}\par
\noindent Source: {\ttfamily\footnotesize ED-EN-NONMATH-G035;\allowbreak{} F7II-XIX-018.5;\allowbreak{} 7II:599;\allowbreak{} 7II:1004;\allowbreak{} fr/sga7ii/\{diplomati\allowbreak{}c,current\}/expose\_XI\allowbreak{}X\_body.tex;\allowbreak{} folio\allowbreak{} 339\allowbreak{}}\par

\par\endgroup\normalsize\fi
\clearpage\SGAtoc{Exposé XX}
\selectlanguage{english}
\clearpage
\SGApage{II}{341}

\begin{center}
{\large SGA 7 --- EXPOS\'E XX}\par
\vspace{2em}
{\Large\underline{\underline{The Griffiths Theorem}}}\par
\vspace{1.5em}
{\large by N. Katz}
\end{center}

\section*{0. Introduction}

We prove Griffiths's theorem~[1], which gives a systematic (though
nonconstructive) way of ``finding'' algebraic cycles that are
$\mathbb Q_\ell$-cohomologous to zero, but no multiple of which is
algebraically equivalent to zero.  The proof given here (due to
Grothendieck) is the translation into purely algebraic terms of Griffiths's
original, more or less transcendental, proof.  In both proofs, the
irreducibility theorem XVIII~6.7 is a key point (cf.~3.1.6 and~3.2).

\section*{1. The formalism of primitive classes}

Throughout what follows, a prime number $\ell$ is fixed.  All schemes under
consideration are assumed to have characteristic different from $\ell$.
Consider the situation

\SGAvar{\begin{equation}\tag{1.0.1}
\begin{tikzcd}[column sep=huge, row sep=huge]
X \arrow[rr, "\pi"] \arrow[dr, "f"'] \& {} \& Y \arrow[dl, "g"] \\
{} \& S \& {}
\end{tikzcd}
\end{equation}
in which we assume
\begin{equation}\tag{1.0.3}
f,\ g \quad\text{smooth of relative dimensions}\quad\ell\quad d_1\ \text{and}\ d_2,
\end{equation}
and
\begin{equation}\tag{1.0.4}
\pi\quad\text{proper}.
\end{equation}
Since $X$ and $Y$ are smooth over $S$, the properness of $\pi$ gives a
\underline{Gysin morphism} (SGA~5, Expos\'e~IV)
\begin{equation}\tag{1.0.5}
\pi_*:H^q(X,\mathbb Q_\ell(i))\longrightarrow
H^{q-2(d_1-d_2)}(Y,\mathbb Q_\ell(i-d_1+d_2)).
\end{equation}
Localizing on $S$ in the \`etale topology, we likewise obtain a morphism of
sheaves on $S$
\begin{equation}\tag{1.0.6}
\pi_*:R^qf_*(\mathbb Q_\ell(i)\longrightarrow
R^{q-2(d_1-d_2)}g_*(\mathbb Q_\ell(i-d_1-d_2)),
\end{equation}
and, on applying the functor $H^p(S,-)$, a morphism
\begin{equation}\tag{1.0.7}
\pi_*:H^p(S,R^qf_*\mathbb Q_\ell(i))\longrightarrow
H^p(S,R^{q-2(d_1-d_2)}g_*\mathbb Q_\ell(i-d_1+d_2)).
\end{equation}

\subsubsection*{Definition 1.0.8.}
Let $(E_r^{p,q},d_r^{p,q})$ and
$({}'E_r^{p,q},{}'d_r^{p,q})$ be two spectral sequences in a given abelian
category, and let $(a,b)\in\mathbb Z\times\mathbb Z$.  The shift of
$({}'E_r^{p,q},{}'d_r^{p,q})$ by $(a,b)$ is the spectral sequence
$({}''E_r^{p,q},{}''d^{p,q})$ given by
\begin{equation}\tag{1.0.9}
\left\{
\begin{aligned}
{}''E_r^{p,q}&={}'E_r^{p+a,q+b},\\
{}''d_r^{p,q}&={}d_r^{p+a,q+b}.
\end{aligned}
\right.
\end{equation}

}{\begin{equation}\tag{1.0.1}
\begin{tikzcd}[column sep=huge, row sep=huge]
X \arrow[rr, "\pi"] \arrow[dr, "f"'] \& {} \& Y \arrow[dl, "g"] \\
{} \& S \& {}
\end{tikzcd}
\end{equation}
in which we assume
\begin{equation}\tag{1.0.3}
f,\ g \quad\text{smooth of relative dimensions}\quad d_1\ \text{and}\ d_2,
\end{equation}
and
\begin{equation}\tag{1.0.4}
\pi\quad\text{proper}.
\end{equation}
Since $X$ and $Y$ are smooth over $S$, the properness of $\pi$ gives a
\underline{Gysin morphism} (SGA~5, Expos\'e~IV)
\begin{equation}\tag{1.0.5}
\pi_*:H^q(X,\mathbb Q_\ell(i))\longrightarrow
H^{q-2(d_1-d_2)}(Y,\mathbb Q_\ell(i-d_1+d_2)).
\end{equation}
Localizing on $S$ in the \`etale topology, we likewise obtain a morphism of
sheaves on $S$
\begin{equation}\tag{1.0.6}
\pi_*:R^qf_*\mathbb Q_\ell(i)\longrightarrow
R^{q-2(d_1-d_2)}g_*\mathbb Q_\ell(i-d_1+d_2),
\end{equation}
and, on applying the functor $H^p(S,-)$, a morphism
\begin{equation}\tag{1.0.7}
\pi_*:H^p(S,R^qf_*\mathbb Q_\ell(i))\longrightarrow
H^p(S,R^{q-2(d_1-d_2)}g_*\mathbb Q_\ell(i-d_1+d_2)).
\end{equation}

\subsubsection*{Definition 1.0.8.}
Let $(E_r^{p,q},d_r^{p,q})$ and
$({}'E_r^{p,q},{}'d_r^{p,q})$ be two spectral sequences in a given abelian
category, and let $(a,b)\in\mathbb Z\times\mathbb Z$.  The shift of
$({}'E_r^{p,q},{}'d_r^{p,q})$ by $(a,b)$ is the spectral sequence
$({}''E_r^{p,q},{}''d_r^{p,q})$ given by
\begin{equation}\tag{1.0.9}
\left\{
\begin{aligned}
{}''E_r^{p,q}&={}'E_r^{p+a,q+b},\\
{}''d_r^{p,q}&={}'d_r^{p+a,q+b}.
\end{aligned}
\right.
\end{equation}

}{E7II-XX-001}\SGApage{II}{\(\approx\)343}
\subsubsection*{1.0.10.}
A morphism of spectral sequences from
$(E_r^{p,q},d_r^{p,q})$ to $({}''E_r^{p,q},{}''d_r^{p,q})$ is called a
morphism of bidegree $(a,b)$ from $(E_r^{p,q},d_r^{p,q})$ to
$({}'E_r^{p,q},{}'d_r^{p,q})$.

The following lemma is an immediate consequence of the theory of the
\underline{Gysin morphism} (SGA~5, Expos\'e~IV).

\SGAvar{\subsubsection*{Lemma 1.1.}
There is a canonical morphism $\pi_*$ of bidegree
$(a,-2(d_1-d_2))$ between the Leray spectral sequences
\begin{equation}\tag{1.1.0}
R^p\alpha_*R^qf_*\mathbb Q_\ell(i)\Longrightarrow
R^{p+q}(\alpha f)_*\mathbb Q_\ell(i)
\end{equation}
and
\begin{equation}\tag{1.1.1}
R^p\alpha_*R^qg_*\mathbb Q_\ell(i-d_1+d_2)\Longrightarrow
R^{p+q}(\alpha g)_*\mathbb Q_\ell(i-d_1+d_2),
\end{equation}
which is given on the $E_2$ term by (1.0.7), and on the abutment by
(1.0.5).

\subsubsection*{1.2.}
Now consider the situation
\begin{equation}\tag{1.2.1}
\begin{tikzcd}[column sep=large, row sep=large]
{} \& X\times_S Y \arrow[dl, "\mathrm{pr}_1"']
                  \arrow[dr, "\mathrm{pr}_2"] \& {} \\
X \arrow[dr, "f"'] \& {} \& Y \arrow[dl, "g"] \\
{} \& S \arrow[d, "\alpha"] \& {} \\
{} \& T \& {}
\end{tikzcd}
\end{equation}
in which $S,X,f,Y,g$ are as in (1.0.2)--(1.0.4), but now $f$ is also
assumed proper.  Let
\begin{equation}\tag{1.2.2}
z\in H^a(X\times_S Y,\mathbb Q_\ell(b)).
\end{equation}
Then $z$ induces a map, again denoted by $z$,
\begin{equation}\tag{1.2.3}
z:H^q(X,\mathbb Q_\ell(i))\longrightarrow
H^{q+a-2d_1}(Y,\mathbb Q_\ell(i+b-d_1))
\end{equation}
by the formula
\begin{equation}\tag{1.2.4}
z(x)=\mathrm{pr}_{2*}\bigl(z\wedge\mathrm{pr}_1^*(x)\bigr),
\end{equation}
where $\mathrm{pr}_{2*}$ is the Gysin homomorphism.  By \`etale localization
on $S$, the class $z$ likewise defines
\begin{equation}\tag{1.2.5}
z:R^qf_*\mathbb Q_\ell(i)\longrightarrow
R^{q+a-2d_1}g_*\mathbb Q_\ell(i+b-d_1),
\end{equation}
and, after applying $H^p(S,-)$, maps
\begin{equation}\tag{1.2.6}
H^p(S,R^qf_*\mathbb Q_\ell(i))\longrightarrow
H^p(S,R^{q+a-2d_1}g_*\mathbb Q_\ell(i+b-d_1).
\end{equation}

\subsubsection*{Lemma 1.3.}
The class $z$ defines a morphism of bidegree
$(0,b-2d_1)$ between the Leray spectral sequences
\begin{equation}\tag{1.2.7}
H^p(S,R^qf_*\mathbb Q_\ell(i))\Longrightarrow
H^{p+q}(X,\mathbb Q_\ell(i))
\end{equation}
and
\begin{equation}\tag{1.2.8}
H^p(S,R^qg_*\mathbb Q_\ell(i+b-d_1))\Longrightarrow
H^{p+q}(Y,\mathbb Q_\ell(i+b-d_1)),
\end{equation}
which is given on the $E_2$ terms by (1.2.6), and on the abutments by
(1.2.3).

\noindent\underline{Proof.}
In view of (1.2.4), and since ``multiplication by $z$'' defines a morphism
of bidegree $(0,a)$ between the Leray spectral sequences
\begin{equation}\tag{1.2.9}
H^p\bigl(S,R^q(\mathrm{pr}_1f)_*\mathbb Q_\ell(i)\bigr)
\Longrightarrow H^{p+q}(X\times_S Y,\mathbb Q_\ell(i))
\end{equation}
\begin{center}
\begin{tikzcd}
{} \arrow[d, "z"'] \\
{}
\end{tikzcd}
\end{center}
\begin{equation}\tag{1.2.10}
H^p\bigl(S,R^{q+a}(\mathrm{pr}_1f)_*\mathbb Q_\ell(i+b)\bigr)
\Longrightarrow H^{p+q+a}(X\times_S Y,\mathbb Q_\ell(i+b)),
\end{equation}
it suffices to apply Lemma~1.1 to the situation
\begin{equation}\tag{1.2.11}
\begin{tikzcd}[column sep=large, row sep=large]
X\times_S Y \arrow[rr, "\mathrm{pr}_2"]
             \arrow[dr, "f\mathrm{pr}_1"]
  \& {} \& Y \arrow[dl, "g"] \\
{} \& S \& {}
\end{tikzcd}
\qquad .
\end{equation}

}{\subsubsection*{Lemma 1.1.}
There is a canonical morphism $\pi_*$ of bidegree
$(0,-2(d_1-d_2))$ between the Leray spectral sequences
\begin{equation}\tag{1.1.0}
R^p\alpha_*R^qf_*\mathbb Q_\ell(i)\Longrightarrow
R^{p+q}(\alpha f)_*\mathbb Q_\ell(i)
\end{equation}
and
\begin{equation}\tag{1.1.1}
R^p\alpha_*R^qg_*\mathbb Q_\ell(i-d_1+d_2)\Longrightarrow
R^{p+q}(\alpha g)_*\mathbb Q_\ell(i-d_1+d_2),
\end{equation}
which is given on the $E_2$ term by (1.0.7), and on the abutment by
(1.0.5).

\subsubsection*{1.2.}
Now consider the situation
\begin{equation}\tag{1.2.1}
\begin{tikzcd}[column sep=large, row sep=large]
{} \& X\times_S Y \arrow[dl, "\mathrm{pr}_1"']
                  \arrow[dr, "\mathrm{pr}_2"] \& {} \\
X \arrow[dr, "f"'] \& {} \& Y \arrow[dl, "g"] \\
{} \& S \arrow[d, "\alpha"] \& {} \\
{} \& T \& {}
\end{tikzcd}
\end{equation}
in which $S,X,f,Y,g$ are as in (1.0.1)--(1.0.4), but now $f$ is also
assumed proper.  Let
\begin{equation}\tag{1.2.2}
z\in H^a(X\times_S Y,\mathbb Q_\ell(b)).
\end{equation}
Then $z$ induces a map, again denoted by $z$,
\begin{equation}\tag{1.2.3}
z:H^q(X,\mathbb Q_\ell(i))\longrightarrow
H^{q+a-2d_1}(Y,\mathbb Q_\ell(i+b-d_1))
\end{equation}
by the formula
\begin{equation}\tag{1.2.4}
z(x)=\mathrm{pr}_{2*}\bigl(z\wedge\mathrm{pr}_1^*(x)\bigr),
\end{equation}
where $\mathrm{pr}_{2*}$ is the Gysin homomorphism.  By \`etale localization
on $S$, the class $z$ likewise defines
\begin{equation}\tag{1.2.5}
z:R^qf_*\mathbb Q_\ell(i)\longrightarrow
R^{q+a-2d_1}g_*\mathbb Q_\ell(i+b-d_1),
\end{equation}
and, after applying $H^p(S,-)$, maps
\begin{equation}\tag{1.2.6}
H^p(S,R^qf_*\mathbb Q_\ell(i))\longrightarrow
H^p(S,R^{q+a-2d_1}g_*\mathbb Q_\ell(i+b-d_1)).
\end{equation}

\subsubsection*{Lemma 1.3.}
The class $z$ defines a morphism of bidegree
$(0,a-2d_1)$ between the Leray spectral sequences
\begin{equation}\tag{1.2.7}
H^p(S,R^qf_*\mathbb Q_\ell(i))\Longrightarrow
H^{p+q}(X,\mathbb Q_\ell(i))
\end{equation}
and
\begin{equation}\tag{1.2.8}
H^p(S,R^qg_*\mathbb Q_\ell(i+b-d_1))\Longrightarrow
H^{p+q}(Y,\mathbb Q_\ell(i+b-d_1)),
\end{equation}
which is given on the $E_2$ terms by (1.2.6), and on the abutments by
(1.2.3).

\noindent\underline{Proof.}
In view of (1.2.4), and since ``multiplication by $z$'' defines a morphism
of bidegree $(0,a)$ between the Leray spectral sequences
\begin{equation}\tag{1.2.9}
H^p\bigl(S,R^q(f\circ\mathrm{pr}_1)_*\mathbb Q_\ell(i)\bigr)
\Longrightarrow H^{p+q}(X\times_S Y,\mathbb Q_\ell(i))
\end{equation}
\begin{center}
\begin{tikzcd}
{} \arrow[d, "z"'] \\
{}
\end{tikzcd}
\end{center}
\begin{equation}\tag{1.2.10}
H^p\bigl(S,R^{q+a}(f\circ\mathrm{pr}_1)_*\mathbb Q_\ell(i+b)\bigr)
\Longrightarrow H^{p+q+a}(X\times_S Y,\mathbb Q_\ell(i+b)),
\end{equation}
it suffices to apply Lemma~1.1 to the situation
\begin{equation}\tag{1.2.11}
\begin{tikzcd}[column sep=large, row sep=large]
X\times_S Y \arrow[rr, "\mathrm{pr}_2"]
             \arrow[dr, "f\circ\mathrm{pr}_1"]
  \& {} \& Y \arrow[dl, "g"] \\
{} \& S \& {}
\end{tikzcd}
\qquad .
\end{equation}

}{E7II-XX-002}\SGApage{II}{345}
\subsubsection*{Corollary 1.4.}
In the situation (1.2.1)--(1.2.2), the morphism (1.2.3)
\begin{equation}\tag{1.4.1}
z:H^q(X,\mathbb Q_\ell(i))\longrightarrow
H^{q+a-2d_1}(Y,\mathbb Q_\ell(i+b-d_1))
\end{equation}
takes primitive classes (XVIII~5.8.1) to primitive classes, and the following
diagram is commutative:
\begin{equation}\tag{1.4.2}
\begin{tikzcd}[column sep=large, row sep=large]
\mathrm{Prim}^q(X/S,\mathbb Q_\ell(i))
  \arrow[r, "z"] \arrow[d, "u_f"]
\& \mathrm{Prim}^{q+a-2d_1}(Y/S,\mathbb Q_\ell(i+b-d_1))
  \arrow[d, "u_g"] \\
H^1(S,R^{q-1}f_*\mathbb Q_\ell(i))
  \arrow[r, "z"]
\& H^1(S,R^{q+a-2d_1-1}g_*\mathbb Q_\ell(i+b-d_1)).
\end{tikzcd}
\end{equation}
Here $u_f$ and $u_g$ are the canonical homomorphisms of XVIII~5.8.1.


\SGApage{II}{\(\approx\)346}
\subsubsection*{1.5.}
Let $\bar\eta$ be a geometric point of $S$ (lying over a point $\eta$ of
$S$), and let
\begin{equation}\tag{1.5.1}
z_{\bar\eta}\in
H^a\bigl(X_{\bar\eta}\times_{k(\bar\eta)}Y_{\bar\eta},
\mathbb Q_\ell(b)\bigr)
\end{equation}
be the class induced by $z$ on the corresponding geometric fiber of
\begin{equation}\tag{1.5.2}
X\times_S Y\longrightarrow S.
\end{equation}

\SGAvar{\subsubsection*{Proposition 1.5.3.}
Suppose that $f$ and $g$ are proper, that $S$ is \underline{annexe}, and let
$q\in\mathbb Z$.  If the class $z_{\bar\eta}$ induces the zero
homomorphism
\begin{equation}\tag{1.5.3}
0=z_{\bar\eta}:H^{q-1}(X_{\bar\eta},\mathbb Q_\ell(i))
\longrightarrow
H^{q+a-2d_1-1}(Y_{\bar\eta},\mathbb Q_\ell(i+b-d_1)),
\end{equation}
then, for every
$X\in\mathrm{Prim}^q(X/S,\mathbb Q_\ell(i))$, one has
\begin{equation}\tag{1.5.4}
u_g(z(x))=0\quad\text{in}\quad
H^1\bigl(S,R^{q+a-2d_1-1}g_*\mathbb Q_\ell(i+b-d_1)\bigr).
\end{equation}

}{\subsubsection*{Proposition 1.5.3.}
Suppose that $f$ and $g$ are proper, that $S$ is connected, and let
$q\in\mathbb Z$.  If the class $z_{\bar\eta}$ induces the zero
homomorphism
\begin{equation}\tag{1.5.3}
0=z_{\bar\eta}:H^{q-1}(X_{\bar\eta},\mathbb Q_\ell(i))
\longrightarrow
H^{q+a-2d_1-1}(Y_{\bar\eta},\mathbb Q_\ell(i+b-d_1)),
\end{equation}
then, for every
$x\in\mathrm{Prim}^q(X/S,\mathbb Q_\ell(i))$, one has
\begin{equation}\tag{1.5.4}
u_g(z(x))=0\quad\text{in}\quad
H^1\bigl(S,R^{q+a-2d_1-1}g_*\mathbb Q_\ell(i+b-d_1)\bigr).
\end{equation}

}{E7II-XX-003}\SGApage{II}{\(\approx\)347}
\noindent\underline{Proof.}
By the commutativity of (1.4.2), it is enough to see that $z$ induces the
zero homomorphism
\begin{equation}\tag{1.5.5}
z:R^{q-1}f_*\mathbb Q_\ell(i)
\stackrel{0}{\longrightarrow}
R^{q+a-2d_1-1}g_*\mathbb Q_\ell(i+b-d_1).
\end{equation}
Indeed, since $f$ and $g$ are proper and smooth, the sheaves
$R^\bullet f_*\mathbb Q_\ell(\mathord\cdot)$ and
$R^\bullet g_*\mathbb Q_\ell(\mathord\cdot)$ are twisted constant sheaves
(SGA~4, XVI~2.2), hence are equivalent to their values at $\bar\eta$,
regarded as modules under $\pi_1(S,\bar\eta)$.

\SGAvar{\subsubsection*{1.6. A construction of Manin}
Suppose that two sections of $f:X\longrightarrow S$ are given:
\begin{equation}\tag{1.6.0}
\begin{tikzcd}[column sep=huge]
S \arrow[r, bend left=45, "t_1"]
  \arrow[r, bend right=45, "t_2"']
\& X \arrow[l, "f" description]
\end{tikzcd}
\qquad ,
\end{equation}
and associate with each its cohomology class
\begin{equation}\tag{1.6.1}
\mathrm{cl}(t_i)=t_{i*}(1)
\in H^{2d_1}(X,\mathbb Q_\ell(d_1)).
\end{equation}
If the fibers of $f:X\longrightarrow S$ are geometrically connected
(more generally, if, for $s\in S$, $f(s)$ and $g(s)$ belong to the same
connected component of $X_s$), it is clear that
\begin{equation}\tag{1.6.2}
\mathrm{cl}(t_1)-\mathrm{cl}(t_2)
\in\mathrm{Prim}^{2d_1}(X/S,\mathbb Q_\ell(d_1)).
\end{equation}
We denote by
\begin{equation}\tag{1.6.3}
u_f(t_1-t_2)\in
H^1(S,R^{2d_1-1}f_*\mathbb Q_\ell(d_1))
\end{equation}
its image under
\begin{equation}\tag{1.6.4}
u_f:\mathrm{Prim}^{2d_1}(X/S,\mathbb Q_\ell(d_1))
\longrightarrow H^1(S,R^{2d_1-1}f_*\mathbb Q_\ell(d_1)).
\end{equation}
(It is an analogue $\mu(t_1-t_2)$ of this element that Manin studies
in~[5].)

}{\subsubsection*{1.6. A construction of Manin}
Suppose that two sections of $f:X\longrightarrow S$ are given:
\begin{equation}\tag{1.6.0}
\begin{tikzcd}[column sep=huge]
S \arrow[r, bend left=45, "t_1"]
  \arrow[r, bend right=45, "t_2"']
\& X \arrow[l, "f" description]
\end{tikzcd}
\qquad ,
\end{equation}
and associate with each its cohomology class
\begin{equation}\tag{1.6.1}
\mathrm{cl}(t_i)=t_{i*}(1)
\in H^{2d_1}(X,\mathbb Q_\ell(d_1)).
\end{equation}
If the fibers of $f:X\longrightarrow S$ are geometrically connected
(more generally, if, for $s\in S$, $t_1(s)$ and $t_2(s)$ belong to the same
connected component of $X_s$), it is clear that
\begin{equation}\tag{1.6.2}
\mathrm{cl}(t_1)-\mathrm{cl}(t_2)
\in\mathrm{Prim}^{2d_1}(X/S,\mathbb Q_\ell(d_1)).
\end{equation}
We denote by
\begin{equation}\tag{1.6.3}
u_f(t_1-t_2)\in
H^1(S,R^{2d_1-1}f_*\mathbb Q_\ell(d_1))
\end{equation}
its image under
\begin{equation}\tag{1.6.4}
u_f:\mathrm{Prim}^{2d_1}(X/S,\mathbb Q_\ell(d_1))
\longrightarrow H^1(S,R^{2d_1-1}f_*\mathbb Q_\ell(d_1)).
\end{equation}
(It is an analogue $\mu(t_1-t_2)$ of this element that Manin studies
in~[5].)

}{E7II-XX-004}\SGApage{II}{\(\approx\)348}
\subsubsection*{1.6.5.}
Now let $z$ be an algebraic cycle on $X\times_S Y$, of codimension at least
$b$ on each fiber of $X\times_S Y$,\footnote{If $S$ is smooth over a scheme
$T$ (for example $T=\operatorname{Spec}k$, with $k$ a field), it would be
enough to suppose that $z$ has codimension at least $b$ on each fiber over
$T$.} which determines a cohomology class (SGA~5, Expos\'e~IV)
\begin{equation}\tag{1.6.6}
\mathrm{cl}(z)\in H^{2b}(X\times_S Y,\mathbb Q_\ell(b)).
\end{equation}
This class defines a map (cf.~(1.2.3))
\begin{equation}\tag{1.6.7}
\mathrm{cl}(z):H^{2d_1}(X,\mathbb Q_\ell(d_1))
\longrightarrow H^{2b}(Y,\mathbb Q_\ell(b)),
\end{equation}
and, if the cycles $z(t_1)$ and $z(t_2)$ are defined, the compatibility of
the cup product with intersections of cycles (SGA~5, Expos\'e~IV) gives
\begin{equation}\tag{1.6.8}
\mathrm{cl}(z)\bigl(\mathrm{cl}(t_1)-\mathrm{cl}(t_2)\bigr)
=\mathrm{cl}\bigl(z(t_1)-z(t_2)\bigr).
\end{equation}
The commutativity of (1.4.2) gives us
\begin{equation}\tag{1.6.9}
\mathrm{cl}(z)\bigl(u_f(t_1-t_2)\bigr)
=u_g\bigl(\mathrm{cl}(z(t_1)-z(t_2))\bigr).
\end{equation}
As a corollary of Proposition~1.5.3, we therefore have:

\SGApage{II}{349}
\SGAvar{\subsubsection*{Proposition 1.7.}
Let $t_1$ and $t_2$ be two sections of $f:X\longrightarrow S$, with $f$
assumed to have connected fibers, and let $z$ be an algebraic cycle on
$X\times_S Y$ of codimension $b$, such that the algebraic cycle
\begin{equation}\tag{1.7.1}
z(t_1)-z(t_2)
\end{equation}
is defined.  Then its cohomology class is primitive (XVIII~5.8.1).  If the
algebraic cycle $Z_{\bar\eta}$ on
$X_{\bar\eta}\times Y_{\bar\eta}$ induces the zero homomorphism
\begin{equation}\tag{1.7.2}
z_{\bar\eta}:H^{2d_1-1}(X_{\bar\eta},\mathbb Q_\ell(d_1))
\stackrel{0}{\longrightarrow}
H^{2b-1}(Y_{\bar\eta},\mathbb Q_\ell(b)),
\end{equation}
then
\begin{equation}\tag{1.7.3}
u_g\bigl(z(t_1)-z(t_2)\bigr)=0\quad\text{in}\quad
H^1(S,R^{2b-1}g_*\mathbb Q_\ell(b)),
\end{equation}
where $u_g$ is the homomorphism of XVIII~5.8.1 occurring in~1.4.



}{\subsubsection*{Proposition 1.7.}
Let $t_1$ and $t_2$ be two sections of $f:X\longrightarrow S$, with $f$
assumed to have connected fibers, and let $z$ be an algebraic cycle on
$X\times_S Y$ of codimension $b$, such that the algebraic cycle
\begin{equation}\tag{1.7.1}
z(t_1)-z(t_2)
\end{equation}
is defined.  Then its cohomology class is primitive (XVIII~5.8.1).  If the
algebraic cycle $z_{\bar\eta}$ on
$X_{\bar\eta}\times Y_{\bar\eta}$ induces the zero homomorphism
\begin{equation}\tag{1.7.2}
z_{\bar\eta}:H^{2d_1-1}(X_{\bar\eta},\mathbb Q_\ell(d_1))
\stackrel{0}{\longrightarrow}
H^{2b-1}(Y_{\bar\eta},\mathbb Q_\ell(b)),
\end{equation}
then
\begin{equation}\tag{1.7.3}
u_g\bigl(z(t_1)-z(t_2)\bigr)=0\quad\text{in}\quad
H^1(S,R^{2b-1}g_*\mathbb Q_\ell(b)),
\end{equation}
where $u_g$ is the homomorphism of XVIII~5.8.1 occurring in~1.4.



}{E7II-XX-005}\SGApage{II}{350}
\section*{2. A reminder on level}
\noindent(cf.~[2], 9.7 and 10.1).

\SGAvar{\subsubsection*{2.0.1.}
Let $X$ be proper and smooth over an algebraically closed field $k$.  We saw
in (1.2.3) that, for every
\begin{equation}\tag{2.0.2}
T\quad\text{proper, smooth, and connected over $k$, of dimension $d$},
\end{equation}
and every
\begin{equation}\tag{2.0.3}
\text{algebraic cycle $z$ of codimension $b$ on $T\times X$},
\end{equation}
the cohomology class
\begin{equation}\tag{2.0.4}
\mathrm{cl}(z)\in H^{2b}(T\times X,\mathbb Q_\ell(b))
\end{equation}
gives rise to maps
\begin{equation}\tag{2.0.5}
\mathrm{cl}(z):H^q(T,\mathbb Q_\ell(i))\longrightarrow
H^{q+2b-2d}(X,\mathbb Q_\ell(i+b-d)).
\end{equation}
The \underline{coniveau filtration} $(N^j)$ on
$H^q(X,\mathbb Q_\ell(i))$ is defined by
\begin{equation}\tag{2.0.6}
N^jH^q(X,\mathbb Q_\ell(i))
=\sum_{(T,z)}\operatorname{Im}\!\left(
\mathrm{cl}(z):H^{q-2j}(T,\mathbb Q_\ell(i-j))
\longrightarrow H^q(X,\mathbb Q_\ell(i))\right),
\end{equation}
where the sum extends over the pairs
\begin{equation}\tag{2.0.7}
(T,z),\qquad
\begin{array}{l}
T/k\text{ proper, smooth, and connected, of dimension }d,\\
z\text{ an algebraic cycle on }T\times X\text{ of codimension }j+d.
\end{array}
\end{equation}
It should be noted that one sometimes speaks of the increasing filtration
$F^j$ ``by level'', defined by
\begin{equation}\tag{2.0.8}
F^jH^q(X,\mathbb Q_\ell(i))
=N^{\left[\frac{g-j}{2}\right]}H^q(X,\mathbb Q_\ell(i)).
\end{equation}
Replacing $(T,z)$ by $(\mathbb P^1\times T,\{\mathrm{pt}\}\times z)$, one
sees that
\begin{equation}\tag{2.0.9}
N^jH^q(X,\mathbb Q_\ell(i))
\subset N^{j-1}H^q(X,\mathbb Q_\ell(i)).
\end{equation}
Obviously, one has:

}{\subsubsection*{2.0.1.}
Let $X$ be proper and smooth over an algebraically closed field $k$.  We saw
in (1.2.3) that, for every
\begin{equation}\tag{2.0.2}
T\quad\text{proper, smooth, and connected over $k$, of dimension $d$},
\end{equation}
and every
\begin{equation}\tag{2.0.3}
\text{algebraic cycle $z$ of codimension $b$ on $T\times X$},
\end{equation}
the cohomology class
\begin{equation}\tag{2.0.4}
\mathrm{cl}(z)\in H^{2b}(T\times X,\mathbb Q_\ell(b))
\end{equation}
gives rise to maps
\begin{equation}\tag{2.0.5}
\mathrm{cl}(z):H^q(T,\mathbb Q_\ell(i))\longrightarrow
H^{q+2b-2d}(X,\mathbb Q_\ell(i+b-d)).
\end{equation}
The \underline{coniveau filtration} $(N^j)$ on
$H^q(X,\mathbb Q_\ell(i))$ is defined by
\begin{equation}\tag{2.0.6}
N^jH^q(X,\mathbb Q_\ell(i))
=\sum_{(T,z)}\operatorname{Im}\!\left(
\mathrm{cl}(z):H^{q-2j}(T,\mathbb Q_\ell(i-j))
\longrightarrow H^q(X,\mathbb Q_\ell(i))\right),
\end{equation}
where the sum extends over the pairs
\begin{equation}\tag{2.0.7}
(T,z),\qquad
\begin{array}{l}
T/k\text{ proper, smooth, and connected, of dimension }d,\\
z\text{ an algebraic cycle on }T\times X\text{ of codimension }j+d.
\end{array}
\end{equation}
It should be noted that one sometimes speaks of the increasing filtration
$F^j$ ``by level'', defined by
\begin{equation}\tag{2.0.8}
F^jH^q(X,\mathbb Q_\ell(i))
=N^{\left[\frac{q-j}{2}\right]}H^q(X,\mathbb Q_\ell(i)).
\end{equation}
Replacing $(T,z)$ by $(\mathbb P^1\times T,\{\mathrm{pt}\}\times z)$, one
sees that
\begin{equation}\tag{2.0.9}
N^jH^q(X,\mathbb Q_\ell(i))
\subset N^{j-1}H^q(X,\mathbb Q_\ell(i)).
\end{equation}
Obviously, one has:

}{E7II-XX-006}\SGApage{II}{351}
\subsubsection*{Proposition 2.1.}
Let $X$ be proper, smooth, and geometrically connected over a field $k$, let
$\bar k$ be an algebraic closure of $k$, and put
$X_{\bar k}=X\otimes_k\bar k$.  Then $\operatorname{Gal}(\bar k/k)$,
acting on $H^q(X_{\bar k},\mathbb Q_\ell(i))$, preserves the coniveau
filtration.

Standard arguments (cf.~[2]) prove the following statements.

\subsubsection*{Proposition 2.2.}
The level filtration is invariant under an extension $k'/k$ of
algebraically closed fields.

\subsubsection*{Proposition 2.3.}
Let $S$ be a reduced and irreducible scheme, with generic point $\eta$, and
let
\begin{equation*}
\pi:X\longrightarrow S
\end{equation*}
be a proper and smooth morphism with geometrically connected fibers.  If,
for given integers $j$ and $q$, one has on the geometric generic fiber
$X_{\bar\eta}$
\begin{equation}\tag{2.3.1}
H^q(X_{\bar\eta},\mathbb Q_\ell)
=N^jH^q(X_{\bar\eta},\mathbb Q_\ell),
\end{equation}
then there exists a nonempty open subset $U\subset S$ such that, at every
point $u\in U$, one has on the geometric fiber $X_{\bar u}$
\begin{equation}\tag{2.3.2}
H^q(X_{\bar u},\mathbb Q_\ell)
=N^jH^q(X_{\bar u},\mathbb Q_\ell).
\end{equation}



\SGApage{II}{352}
\section*{3. Application to Lefschetz pencils}

\SGAvar{\subsubsection*{3.1.}
From now on, we fix
\begin{equation}\tag{3.1.0}
k,\quad\text{an algebraically closed field},
\end{equation}
and
\begin{equation}\tag{3.1.1}
\begin{aligned}
X\subset\mathbb P^r,\quad
&\text{a projective, smooth, connected }k\text{-scheme}\\
&\text{of even dimension }n\text{, satisfying (LV) (XVIII~5.2.2).}
\end{aligned}
\end{equation}
Take
\begin{equation}\tag{3.1.2}
D=\{H_t\}_{t\in\mathbb P^1},
\quad\text{a Lefschetz pencil of hypersurfaces of degree $d$},
\end{equation}
with generic fiber
\begin{equation}\tag{3.1.3}
X_\eta,\qquad \eta=\text{the generic point of }\mathbb P^1.
\end{equation}
Put (cf.~XVIII~5.4.4)
\begin{equation}\tag{3.1.4}
m:X_{\bar\eta}\hookrightarrow X\times_k k(\bar\eta),
\quad\text{the inclusion},
\end{equation}
and
\begin{equation}\tag{3.1.5}
E^{n-1}(X_{\bar\eta},\mathbb Q_\ell(i))
=\bigl(m^*H^{n-1}(X,\mathbb Q_\ell(i))\bigr)^\perp
\subset H^{n-1}(X_{\bar\eta},\mathbb Q_\ell(i)).
\end{equation}
We saw in XVIII~6.7.2 that the \underline{Picard--Lefschetz formula}
implies
\begin{equation}\tag{3.1.6}
\begin{gathered}
\text{If }n\text{ is even, every finite-index subgroup of }
\operatorname{Gal}(k(\bar\eta)/k(\eta))\\
\text{acts irreducibly on }E^{2n-1}(X_{\bar\eta},\mathbb Q_\ell(i)).
\end{gathered}
\end{equation}

}{\subsubsection*{3.1.}
From now on, we fix
\begin{equation}\tag{3.1.0}
k,\quad\text{an algebraically closed field},
\end{equation}
and
\begin{equation}\tag{3.1.1}
\begin{aligned}
X\subset\mathbb P^r,\quad
&\text{a projective, smooth, connected }k\text{-scheme}\\
&\text{of even dimension }n\text{, satisfying (LV) (XVIII~5.2.2).}
\end{aligned}
\end{equation}
Take
\begin{equation}\tag{3.1.2}
D=\{H_t\}_{t\in\mathbb P^1},
\quad\text{a Lefschetz pencil of hypersurfaces of degree $d$},
\end{equation}
with generic fiber
\begin{equation}\tag{3.1.3}
X_\eta,\qquad \eta=\text{the generic point of }\mathbb P^1.
\end{equation}
Put (cf.~XVIII~5.4.4)
\begin{equation}\tag{3.1.4}
m:X_{\bar\eta}\hookrightarrow X\times_k k(\bar\eta),
\quad\text{the inclusion},
\end{equation}
and
\begin{equation}\tag{3.1.5}
E^{n-1}(X_{\bar\eta},\mathbb Q_\ell(i))
=\bigl(m^*H^{n-1}(X,\mathbb Q_\ell(i))\bigr)^\perp
\subset H^{n-1}(X_{\bar\eta},\mathbb Q_\ell(i)).
\end{equation}
We saw in XVIII~6.7.2 that the \underline{Picard--Lefschetz formula}
implies
\begin{equation}\tag{3.1.6}
\begin{gathered}
\text{If }n\text{ is even, every finite-index subgroup of }
\operatorname{Gal}(k(\bar\eta)/k(\eta))\\
\text{acts irreducibly on }E^{n-1}(X_{\bar\eta},\mathbb Q_\ell(i)).
\end{gathered}
\end{equation}

}{E7II-XX-007}\subsubsection*{Proposition 3.2.}
Let $D=\{H_t\}$ be a Lefschetz pencil, and suppose that $n=2m$ is even.
Then either
\begin{equation}\tag{3.2.1}
H^{n-1}(X_{\bar\eta},\mathbb Q_\ell(i))
\subset m^*H^{n-1}(X,\mathbb Q_\ell(i))
+N^{m-1}H^{n-1}(X_{\bar\eta},\mathbb Q_\ell(i)),
\end{equation}
or else:

\SGApage{II}{\(\approx\)353}
\subsubsection*{(3.2.2)}
For every curve $C$ proper and smooth over $k(\bar\eta)$ and every
algebraic cycle $z$ on $C\times X_{\bar\eta}$ of codimension $m$, the
induced homomorphism (1.2.3)
\begin{equation}\tag{3.2.3}
\mathrm{cl}(z):H^1(C,\mathbb Q_\ell(1+i-m))
\longrightarrow H^{2m-1}(X_{\bar\eta},\mathbb Q_\ell(i))
\end{equation}
has its image in $m^*H^{2m-1}(X,\mathbb Q_\ell(i))$.

\SGAvar{\noindent\underline{Proof.}
Suppose that the image of such a homomorphism (3.2.3) is not contained in
$m^*H^{n-1}(X,\mathbb Q_\ell(i))$, or, equivalently, that its image in
$E^{n-1}(X_{\bar\eta},\mathbb Q_\ell(i))=H^{n-1}(X_{\bar\eta},\mathbb Q_\ell(i))\quad H^{n-1}(X,Q_\ell(i))$ is nonzero.  This image is
evidently stable under the action of a finite-index subgroup of
$\operatorname{Gal}(k(\bar\eta)/k(\eta))$.  The irreducibility (3.1.6) of
$E^{n-1}(X_{\bar\eta},\mathbb Q_\ell(i)$ therefore says that
\begin{equation*}
H^{n-1}(X_{\bar\eta},\mathbb Q_\ell(i))
=m^*H^{n-1}(X,\mathbb Q_\ell(i))+\operatorname{Im}(3.2.3).
\end{equation*}
By definition,
\begin{equation*}
\operatorname{Im}(3.2.3)
=\mathrm{cl}(z)\bigl(H^1(C,\mathbb Q_\ell(1+i-m))\bigr)
\subset N^{m-1}H^{n-1}(X_{\bar\eta},\mathbb Q_\ell(i)).
\end{equation*}
Consequently,
\begin{equation*}
H^{n-1}(X_{\bar\eta},\mathbb Q_\ell(i))
=m^*H^{n-1}(X,\mathbb Q_\ell(i))
+N^{m-1}H^{n-1}(X_{\bar\eta},\mathbb Q_\ell(i)),
\end{equation*}
which is precisely the first alternative.


}{\noindent\underline{Proof.}
Suppose that the image of such a homomorphism (3.2.3) is not contained in
$m^*H^{n-1}(X,\mathbb Q_\ell(i))$, or, equivalently, that its image in
$E^{n-1}(X_{\bar\eta},\mathbb Q_\ell(i))=H^{n-1}(X_{\bar\eta},\mathbb Q_\ell(i))/H^{n-1}(X,\mathbb Q_\ell(i))$ is nonzero.  This image is
evidently stable under the action of a finite-index subgroup of
$\operatorname{Gal}(k(\bar\eta)/k(\eta))$.  The irreducibility (3.1.6) of
$E^{n-1}(X_{\bar\eta},\mathbb Q_\ell(i))$ therefore says that
\begin{equation*}
H^{n-1}(X_{\bar\eta},\mathbb Q_\ell(i))
=m^*H^{n-1}(X,\mathbb Q_\ell(i))+\operatorname{Im}(3.2.3).
\end{equation*}
By definition,
\begin{equation*}
\operatorname{Im}(3.2.3)
=\mathrm{cl}(z)\bigl(H^1(C,\mathbb Q_\ell(1+i-m))\bigr)
\subset N^{m-1}H^{n-1}(X_{\bar\eta},\mathbb Q_\ell(i)).
\end{equation*}
Consequently,
\begin{equation*}
H^{n-1}(X_{\bar\eta},\mathbb Q_\ell(i))
=m^*H^{n-1}(X,\mathbb Q_\ell(i))
+N^{m-1}H^{n-1}(X_{\bar\eta},\mathbb Q_\ell(i)),
\end{equation*}
which is precisely the first alternative.


}{E7II-XX-008}\SGApage{II}{354}
\subsubsection*{3.3. Condition (B)}

\SGAvar{\paragraph{3.3.1.}
We shall say that the inclusion
\begin{equation*}
m:X_{\bar\eta}\hookrightarrow X\times_k k(\bar\eta)
\end{equation*}
satisfies condition (B) if there exists
\begin{equation}\tag{3.3.2}
V_{\bar\eta},\quad
\text{an algebraic cycle of codimension $n-1$ on }
X_{\bar\eta}\times\bigl(X\times_k k(\bar\eta)\bigr),
\end{equation}
such that the map it defines
\begin{equation}\tag{3.3.3}
V_{\bar\eta};H^{n-1}(X_{\bar\eta},\mathbb Q_\ell(i))
\longrightarrow H^{n-1}(X,\mathbb Q_\ell(i))
\end{equation}
is a left inverse of
\begin{equation}\tag{3.3.4}
m^*:H^{n-1}(X,\mathbb Q_\ell(i))
\longrightarrow H^{n-1}(X_{\eta},\mathbb Q_\ell(i)),
\end{equation}
that is, such that
\begin{equation}\tag{3.3.5}
V_{\bar\eta}\,m^*=\mathrm{id}.
\end{equation}

}{\paragraph{3.3.1.}
We shall say that the inclusion
\begin{equation*}
m:X_{\bar\eta}\hookrightarrow X\times_k k(\bar\eta)
\end{equation*}
satisfies condition (B) if there exists
\begin{equation}\tag{3.3.2}
V_{\bar\eta},\quad
\text{an algebraic cycle of codimension $n-1$ on }
X_{\bar\eta}\times\bigl(X\times_k k(\bar\eta)\bigr),
\end{equation}
such that the map it defines
\begin{equation}\tag{3.3.3}
V_{\bar\eta}:H^{n-1}(X_{\bar\eta},\mathbb Q_\ell(i))
\longrightarrow H^{n-1}(X,\mathbb Q_\ell(i))
\end{equation}
is a left inverse of
\begin{equation}\tag{3.3.4}
m^*:H^{n-1}(X,\mathbb Q_\ell(i))
\longrightarrow H^{n-1}(X_{\bar\eta},\mathbb Q_\ell(i)),
\end{equation}
that is, such that
\begin{equation}\tag{3.3.5}
V_{\bar\eta}\,m^*=\mathrm{id}.
\end{equation}

}{E7II-XX-009}\subsubsection*{3.4.}
Condition (B) is satisfied if the operation $\Lambda$ of [4, 1.4.2.1]
is algebraic (cf.~[4, 2.12(ii)]).  It is conjectured that $\Lambda$ is
always algebraic, and hence that condition (B) always holds; this is
known in many cases [4].  It holds trivially when $X$ is a complete
intersection of even dimension, since in that case we know (IX~1.6)
that
\begin{equation}\tag{3.4.1}
H^{n-1}(X,\mathbb Q_\ell(i))=0.
\end{equation}


\SGApage{II}{355}
\section*{4. Griffiths's theorem}

\subsubsection*{Theorem 4.1.}
Let $X$ be projective, smooth, and connected over $k$, of even dimension
$n=2m\geq4$.  Let $D$ be a Lefschetz pencil of sections by hypersurfaces
of degree $d$.  Assume that:

\medskip
\noindent\textbf{(4.1.0)} $X$ satisfies (LV) (XVIII~5.2.2);

\noindent\textbf{(4.1.1)} the pencil $D$ satisfies hypothesis (A)
(XVIII~5.3);

\noindent\textbf{(4.1.2)} the pencil $D$ satisfies alternative (3.2.2),
that is, in the setting of Proposition~3.2, the vanishing cohomology
$E^{n-1}(X_{\bar\eta},\mathbb Q_\ell)$ is not of level $\leq1$;

\SGAvar{\noindent\textbf{(4.1.3)} the inclusion
$m^*:X_{\bar\eta}\hookrightarrow X\times_k k(\bar\eta)$ satisfies
hypothesis (B) (3.3).

}{\noindent\textbf{(4.1.3)} the inclusion
$m:X_{\bar\eta}\hookrightarrow X\times_k k(\bar\eta)$ satisfies
hypothesis (B) (3.3).

}{E7II-XX-010}\medskip
Then, if an algebraic cycle $z$ of codimension $m$ on $X$ has a
restriction to $X_{\bar\eta}$ which, as an algebraic cycle, satisfies
\begin{equation}\tag{4.1.4}
z_{\bar\eta}=z|_{X_{\bar\eta}}
\quad\text{is algebraically equivalent to zero},
\end{equation}
then
\begin{equation}\tag{4.1.5}
\operatorname{cl}(z)=0
\quad\text{in }H^{2m}(X,\mathbb Q_\ell(m));
\end{equation}
in other words, $z$ is $\mathbb Q_\ell$-homologous to zero.

\subsubsection*{4.2. A remark on the case where $X$ is a complete intersection}
As already noted (cf.~XVIII~5.2.5 and (3.4)), hypotheses (4.1.0) and
(4.1.3) then hold automatically.  We saw in XVIII~6.3.4 and 6.4 that
hypothesis (4.1.1) holds for a pencil of hypersurfaces of degree $d>2n$.
Moreover, in the next exposé we shall prove condition (4.1.2) for a
``general'' pencil of hypersurfaces of degree $d>2n$ (XXI~5.2).
Consequently, a ``general'' pencil of hypersurfaces of degree $d>2n$
satisfies all the hypotheses of Theorem~4.1.

\SGAvar{\subsubsection*{4.3.}
Returning to the case of arbitrary $X$, it seems very plausible that, for
$d$ sufficiently large, all the conditions of 4.1 hold for $D$, at least
when $D$ is sufficiently general.  We have already seen this for (4.1.1)
(XVIII~6.4), and we noted the conjecture that (4.1.0) and (4.1.3) always
hold (with (4.1.0) known in every case when $k=0$).
As for condition (4.1.2), Hodge theory proves it for large $d$ when
$k=0$, while heuristic arguments, given in the
appendix to the next exposé, make the same conclusion plausible when
$k>0$.

}{\subsubsection*{4.3.}
Returning to the case of arbitrary $X$, it seems very plausible that, for
$d$ sufficiently large, all the conditions of 4.1 hold for $D$, at least
when $D$ is sufficiently general.  We have already seen this for (4.1.1)
(XVIII~6.4), and we noted the conjecture that (4.1.0) and (4.1.3) always
hold (with (4.1.0) known in every case when $\operatorname{char}(k)=0$).
As for condition (4.1.2), Hodge theory proves it for large $d$ when
$\operatorname{char}(k)=0$, while heuristic arguments, given in the
appendix to the next exposé, make the same conclusion plausible when
$\operatorname{char}(k)>0$.

}{E7II-XX-011}\SGApage{II}{\(\approx\)356}
More generally, these arguments---Hodge theory and, respectively,
crystalline cohomology---prove, and respectively make plausible, that for
a smooth projective variety $X$ of dimension $n$, and for every Lefschetz
pencil $D$ of hypersurfaces of sufficiently large degree $d$ (respectively,
every sufficiently general pencil $D$ of hypersurfaces of degree $d$), the
vanishing cohomology
$E^{n-1}(X_{\bar\eta},\mathbb Q_\ell)$ has level exactly $n-1$; that is, it
is not of level $\leq n-3$, or equivalently is not of coniveau $1$.


\subsubsection*{4.4. Proof of Griffiths's theorem}

\SGAvar{Let $z$ be an algebraic cycle of codimension $n$ on $X$ such that
\begin{equation}\tag{4.4.1}
z_{\bar\eta}=z|_{X_{\bar\eta}}
\quad\text{is algebraically equivalent to zero}.
\end{equation}
By definition, there exist
\begin{equation}\tag{4.4.2}
\text{a proper, smooth, connected curve $C_{\bar\eta}$ over $k(\bar\eta)$},
\end{equation}
\begin{equation}\tag{4.4.3}
t_i(\bar\eta):\Spec k(\bar\eta)\longrightarrow C,
\qquad i=1,2,
\end{equation}
and
\begin{equation}\tag{4.4.4}
\text{an algebraic cycle $P_{\bar\eta}$ of codimension $n$ on
$C_{\bar\eta}\times X_{\bar\eta}$},
\end{equation}
such that
\begin{equation}\tag{4.4.5}
z_{\bar\eta}
=P_{\bar\eta}\bigl(t_1(\bar\eta)\bigr)
-P_{\bar\eta}\bigl(t_2(\bar\eta)\bigr)
\quad\text{on }X_{\bar\eta}.
\end{equation}


Regard $P_{\bar\eta}$ as an algebraic correspondence from
$C_{\bar\eta}$ to $X_{\bar\eta}$.  It defines a cohomological
correspondence, again denoted by $P_{\bar\eta}$, or equivalently a
homomorphism
\begin{equation}\tag{4.4.6}
P_{\bar\eta}:H^*(C_{\bar\eta},\mathbb Q_\ell(i))
\longrightarrow
H^{*+2(m-1)}(X_{\bar\eta},\mathbb Q_\ell((i+m-1)).
\end{equation}


Likewise, consider the algebraic correspondence $V_{\bar\eta}$ of
(3.3.2), from $X_{\bar\eta}$ to $X\otimes_k k(\bar\eta)$, supplied by
hypothesis (B) in (4.1.3).  It defines
\begin{equation}\tag{4.4.7}
V_{\bar\eta}:H^*(X_{\bar\eta},\mathbb Q_\ell(i+m-1))
\longrightarrow
H^*(X\otimes_k k(\bar\eta),\mathbb Q_\ell(i+m-1)).
\end{equation}
The composite algebraic correspondence
$Q_{\bar\eta}=V_{\bar\eta}\circ P$ from $C_{\bar\eta}$ to
$X\otimes_k k(\bar\eta)$ therefore defines
\begin{equation}\tag{4.4.8}
Q_{\bar\eta}=V_{\bar\eta}\circ P_{\bar\eta}:
H^*(C_{\bar\eta},\mathbb Q_\ell(i))
\longrightarrow
H^{*+2(m-1)}(X\otimes_k k(\bar\eta),\mathbb Q_\ell(i+m-1)).
\end{equation}

Consider the inclusion
\begin{equation*}
m:X_{\bar\eta}\longrightarrow X\otimes_k k(\bar\eta).
\end{equation*}
It gives the algebraic correspondence $m^*$ from
$X\otimes_k k(\bar\eta)$ to $X_{\bar\eta}$, hence the composite
correspondence $m^*Q_{\bar\eta}\;{}_{\bar\eta}$ from $C_{\bar\eta}$ to
$X_{\bar\eta}$, defining
\begin{equation}\tag{4.4.9}
m^*Q_{\bar\eta}=m^*V_{\bar\eta}P_{\bar\eta}:
H^*(C_{\bar\eta},\mathbb Q_\ell(i))
\longrightarrow
H^{*+2(m-1)}(X_{\bar\eta},\mathbb Q_\ell(i+m-1)).
\end{equation}


Now consider the algebraic correspondence
\begin{equation}\tag{4.4.10}
\begin{split}
R_{\bar\eta}
&=P_{\bar\eta}m^*Q_{\bar\eta}
=(\id-m^*V_{\bar\eta})P_{\bar\eta}:\\
&\hspace{2em}H^*(C_{\bar\eta},\mathbb Q_\ell(i))
\longrightarrow
H^{*+2(m-1)}(X_{\bar\eta},\mathbb Q_\ell(i+m-1)).
\end{split}
\end{equation}
Since $V_{\bar\eta}$ is a left inverse of $m^*$,
$\id-m^*V_{\bar\eta}$ is a projector on $H^*(X_{\bar\eta},-)$ whose
image is a complement of
\begin{equation*}
\operatorname{Im}\bigl(m^*:H^*(X\otimes_k k(\bar\eta),-)
\longrightarrow H^*(X_{\bar\eta},-)\bigr).
\end{equation*}
By hypothesis (4.1.2), the image of
\begin{equation}\tag{4.4.11}
R_{\bar\eta}:H^1(C_{\bar\eta},\mathbb Q_\ell(i))
\longrightarrow
H^{2m-1}(X_{\bar\eta},\mathbb Q_\ell(i+m-1))
\end{equation}
lies in $\operatorname{Im}m^*$.  Consequently,
\begin{equation}\tag{4.4.12}
\text{the homomorphism (4.4.11) is zero.}
\end{equation}

}{Let $z$ be an algebraic cycle of codimension $m$ on $X$ such that
\begin{equation}\tag{4.4.1}
z_{\bar\eta}=z|_{X_{\bar\eta}}
\quad\text{is algebraically equivalent to zero}.
\end{equation}
By definition, there exist
\begin{equation}\tag{4.4.2}
\text{a proper, smooth, connected curve $C_{\bar\eta}$ over $k(\bar\eta)$},
\end{equation}
\begin{equation}\tag{4.4.3}
t_i(\bar\eta):\Spec k(\bar\eta)\longrightarrow C_{\bar\eta},
\qquad i=1,2,
\end{equation}
and
\begin{equation}\tag{4.4.4}
\text{an algebraic cycle $P_{\bar\eta}$ of codimension $m$ on
$C_{\bar\eta}\times X_{\bar\eta}$},
\end{equation}
such that
\begin{equation}\tag{4.4.5}
z_{\bar\eta}
=P_{\bar\eta}\bigl(t_1(\bar\eta)\bigr)
-P_{\bar\eta}\bigl(t_2(\bar\eta)\bigr)
\quad\text{on }X_{\bar\eta}.
\end{equation}


Regard $P_{\bar\eta}$ as an algebraic correspondence from
$C_{\bar\eta}$ to $X_{\bar\eta}$.  It defines a cohomological
correspondence, again denoted by $P_{\bar\eta}$, or equivalently a
homomorphism
\begin{equation}\tag{4.4.6}
P_{\bar\eta}:H^*(C_{\bar\eta},\mathbb Q_\ell(i))
\longrightarrow
H^{*+2(m-1)}(X_{\bar\eta},\mathbb Q_\ell(i+m-1)).
\end{equation}


Likewise, consider the algebraic correspondence $V_{\bar\eta}$ of
(3.3.2), from $X_{\bar\eta}$ to $X\otimes_k k(\bar\eta)$, supplied by
hypothesis (B) in (4.1.3).  It defines
\begin{equation}\tag{4.4.7}
V_{\bar\eta}:H^*(X_{\bar\eta},\mathbb Q_\ell(i+m-1))
\longrightarrow
H^*(X\otimes_k k(\bar\eta),\mathbb Q_\ell(i+m-1)).
\end{equation}
The composite algebraic correspondence
$Q_{\bar\eta}=V_{\bar\eta}\circ P_{\bar\eta}$ from $C_{\bar\eta}$ to
$X\otimes_k k(\bar\eta)$ therefore defines
\begin{equation}\tag{4.4.8}
Q_{\bar\eta}=V_{\bar\eta}\circ P_{\bar\eta}:
H^*(C_{\bar\eta},\mathbb Q_\ell(i))
\longrightarrow
H^{*+2(m-1)}(X\otimes_k k(\bar\eta),\mathbb Q_\ell(i+m-1)).
\end{equation}

Consider the inclusion
\begin{equation*}
m:X_{\bar\eta}\longrightarrow X\otimes_k k(\bar\eta).
\end{equation*}
It gives the algebraic correspondence $m^*$ from
$X\otimes_k k(\bar\eta)$ to $X_{\bar\eta}$, hence the composite
correspondence $m^*Q_{\bar\eta}$ from $C_{\bar\eta}$ to
$X_{\bar\eta}$, defining
\begin{equation}\tag{4.4.9}
m^*Q_{\bar\eta}=m^*V_{\bar\eta}P_{\bar\eta}:
H^*(C_{\bar\eta},\mathbb Q_\ell(i))
\longrightarrow
H^{*+2(m-1)}(X_{\bar\eta},\mathbb Q_\ell(i+m-1)).
\end{equation}


Now consider the algebraic correspondence
\begin{equation}\tag{4.4.10}
\begin{split}
R_{\bar\eta}
&=P_{\bar\eta}-m^*Q_{\bar\eta}
=(\id-m^*V_{\bar\eta})P_{\bar\eta}:\\
&\hspace{2em}H^*(C_{\bar\eta},\mathbb Q_\ell(i))
\longrightarrow
H^{*+2(m-1)}(X_{\bar\eta},\mathbb Q_\ell(i+m-1)).
\end{split}
\end{equation}
Since $V_{\bar\eta}$ is a left inverse of $m^*$,
$\id-m^*V_{\bar\eta}$ is a projector on $H^*(X_{\bar\eta},-)$ whose
image is a complement of
\begin{equation*}
\operatorname{Im}\bigl(m^*:H^*(X\otimes_k k(\bar\eta),-)
\longrightarrow H^*(X_{\bar\eta},-)\bigr).
\end{equation*}
By hypothesis (4.1.2), the image of
\begin{equation}\tag{4.4.11}
R_{\bar\eta}:H^1(C_{\bar\eta},\mathbb Q_\ell(i))
\longrightarrow
H^{2m-1}(X_{\bar\eta},\mathbb Q_\ell(i+m-1))
\end{equation}
lies in $\operatorname{Im}m^*$.  Consequently,
\begin{equation}\tag{4.4.12}
\text{the homomorphism (4.4.11) is zero.}
\end{equation}

}{E7II-XX-012}\SGApage{II}{\(\approx\)358}
Choose a finite extension
\begin{equation*}
L/k(\eta)
\end{equation*}
over which the objects
\begin{equation*}
C_{\bar\eta},\quad t_1(\bar\eta),\quad t_2(\bar\eta),\quad
Q_{\bar\eta},\quad R_{\bar\eta}
\end{equation*}
are already defined, and let $U$ be a nonempty open subset of the
normalization of $\mathbb P^1$ in $L$.  After shrinking $U$, there is a
proper smooth relative curve
\begin{equation*}
C\longrightarrow U
\end{equation*}
with geometrically connected fibers and geometric generic fiber
$C_{\bar\eta}$, together with two sections
\begin{equation*}
t_i:U\longrightarrow C,\qquad i=1,2,
\end{equation*}
whose values at $\bar\eta$ are $t_i(\bar\eta)$.

\SGApage{II}{\(\approx\)359}
\SGAvar{Let
\begin{equation*}
\rho:\widetilde X\longrightarrow\mathbb P^1
\end{equation*}
be the projection of XVIII~3.1.4 whose fibers are the members of the
pencil, put
\begin{equation*}
\widetilde X_U=\widetilde X\times_{\mathbb P^1}U,
\end{equation*}
and let
\begin{equation*}
M:\widetilde X_U\hookrightarrow X\times U
\end{equation*}
be the canonical immersion into the constant family over $U$ with fiber
$X$.  Set
\begin{equation*}
z_U=M^*\pr_i^*(z),
\end{equation*}
the algebraic cycle on $\widetilde X_U$ induced by $z$.

}{Let
\begin{equation*}
\rho:\widetilde X\longrightarrow\mathbb P^1
\end{equation*}
be the projection of XVIII~3.1.4 whose fibers are the members of the
pencil, put
\begin{equation*}
\widetilde X_U=\widetilde X\times_{\mathbb P^1}U,
\end{equation*}
and let
\begin{equation*}
M:\widetilde X_U\hookrightarrow X\times U
\end{equation*}
be the canonical immersion into the constant family over $U$ with fiber
$X$.  Set
\begin{equation*}
z_U=M^*\pr_1^*(z),
\end{equation*}
the algebraic cycle on $\widetilde X_U$ induced by $z$.

}{E7II-XX-013}After shrinking $U$ once more, there are an algebraic cycle $R$ on
$C\times_U\widetilde X_U$ whose value on
$C_{\bar\eta}\times X_{\bar\eta}$ is $R_{\bar\eta}$, and an algebraic
cycle $Q$ on $C\times X\times U$ whose value on
$C_{\bar\eta}\times(X\otimes_k k(\bar\eta))$ is $Q_{\bar\eta}$, such
that
\begin{equation}\tag{4.4.13}
z_U=R(t_1)-R(t_2)
+(\id\times M)^*\bigl[Q(t_1)-Q(t_2)\bigr].
\end{equation}

\SGAvar{To prove the theorem it suffices, by the injectivity of the Griffiths
homomorphism ``modulo the fixed part'' (XVIII~5.8.7),
\begin{equation*}
\operatorname{Prim}^{2m}(X,\mathbb Q_\ell(n))
\xrightarrow{\ \operatorname{grif}\ }
\frac{H^1(U,R^{2m-1}\rho_{U*}\mathbb Q_\ell(m))}
{M^*H^1(U,H^{2m-1}(X,\mathbb Q_\ell(m))},
\end{equation*}
to show that the image of $z_U$ under
\begin{equation*}
u_0:\operatorname{Prim}^{2m}(\widetilde X_U/U,\mathbb Q_\ell(m))
\longrightarrow H^1(U,R^{2m-1}\rho_{U*}\mathbb Q_\ell(m))
\end{equation*}
lies in
$M^*(H^1(U,H^{2m-1}(X,\mathbb Q_\ell(m))$.


In view of (4.4.13) and the commutativity of (1.4.2), we have the
commutative diagram
\begin{center}
\begin{tikzcd}[column sep=4.2em,row sep=4.2em]
\operatorname{Prim}^{2}(C/U,\mathbb Q_\ell(1))
  \arrow[r,"Q"]
\& \operatorname{Prim}^{2m}(X\times U/U,\mathbb Q_\ell(m))
  \arrow[r,"M^*"] \arrow[d,"u_{\pr_2}"]
\& \operatorname{Prim}^{2m}(\widetilde X_U/U,\mathbb Q_\ell(m))
  \arrow[d,"u_0"] \\
{}
\& H^1(U,H^{2m-1}(X,\mathbb Q_\ell(m)))
  \arrow[r,"M^*"]
\& H^1(U,R^{2m-1}\rho_{U*}\mathbb Q_\ell(m)),
\end{tikzcd}
\end{center}
so it is enough to prove
\begin{equation*}
u_0\bigl(R(t_1)-R(t_2)\bigr)=0
\quad\text{in}\quad
H^1(U,R^{2n-1}\rho_{U*}\mathbb Q_\ell(n)).
\end{equation*}
This follows from (4.3.12), by 1.7.  


}{To prove the theorem it suffices, by the injectivity of the Griffiths
homomorphism ``modulo the fixed part'' (XVIII~5.8.7),
\begin{equation*}
\operatorname{Prim}^{2m}(X,\mathbb Q_\ell(m))
\xrightarrow{\ \operatorname{grif}\ }
\frac{H^1(U,R^{2m-1}\rho_{U*}\mathbb Q_\ell(m))}
{M^*H^1(U,H^{2m-1}(X,\mathbb Q_\ell(m)))},
\end{equation*}
to show that the image of $z_U$ under
\begin{equation*}
u_0:\operatorname{Prim}^{2m}(\widetilde X_U/U,\mathbb Q_\ell(m))
\longrightarrow H^1(U,R^{2m-1}\rho_{U*}\mathbb Q_\ell(m))
\end{equation*}
lies in
$M^*H^1(U,H^{2m-1}(X,\mathbb Q_\ell(m)))$.


In view of (4.4.13) and the commutativity of (1.4.2), we have the
commutative diagram
\begin{center}
\begin{tikzcd}[column sep=4.2em,row sep=4.2em]
\operatorname{Prim}^{2}(C/U,\mathbb Q_\ell(1))
  \arrow[r,"Q"]
\& \operatorname{Prim}^{2m}(X\times U/U,\mathbb Q_\ell(m))
  \arrow[r,"M^*"] \arrow[d,"u_{\pr_2}"]
\& \operatorname{Prim}^{2m}(\widetilde X_U/U,\mathbb Q_\ell(m))
  \arrow[d,"u_0"] \\
{}
\& H^1(U,H^{2m-1}(X,\mathbb Q_\ell(m)))
  \arrow[r,"M^*"]
\& H^1(U,R^{2m-1}\rho_{U*}\mathbb Q_\ell(m)),
\end{tikzcd}
\end{center}
so it is enough to prove
\begin{equation*}
u_0\bigl(R(t_1)-R(t_2)\bigr)=0
\quad\text{in}\quad
H^1(U,R^{2m-1}\rho_{U*}\mathbb Q_\ell(m)).
\end{equation*}
This follows from (4.4.12), by 1.7.  


}{E7II-XX-014}\SGApage{II}{360}
\section*{5. The Griffiths group}

\subsubsection*{5.0.}
Let $S$ be proper and smooth over an algebraically closed field $k$.
Denote by

\SGApage{II}{\(\approx\)361}
\medskip
\noindent\textbf{(5.0.1)} $Z^i(S)$ the group of algebraic cycles of
codimension $i$ on $S$;

\noindent\textbf{(5.0.2)} $Z^i_{\mathrm{alg}}(S)$ the subgroup of
$Z^i(S)$ consisting of cycles algebraically equivalent to zero; and, for
every prime number $\ell\ne\operatorname{char}(k)$,

\noindent\textbf{(5.0.3)} $Z^i_{\ell\text{-}\mathrm{coh}}(S)$ the subgroup
of $Z^i(S)$ consisting of cycles whose cohomology class in
$H^{2i}(S,\mathbb Q_\ell(i))$ is zero.

\SGAvar{It is conjectured in \textup{[4]} that this group is independent of $\ell$;
this is known at least when $k=0$.  Set
\begin{equation}\tag{5.0.4}
\mathrm{Grif}^i_\ell(S)
 =Z^i_{\ell\text{-}\mathrm{coh}}(S)/Z^i_{\mathrm{alg}}(S).
\end{equation}
If, moreover, $S$ is projective and is equipped with a projective embedding,
set

}{It is conjectured in \textup{[4]} that this group is independent of $\ell$;
this is known at least when $\operatorname{char}(k)=0$.  Set
\begin{equation}\tag{5.0.4}
\mathrm{Grif}^i_\ell(S)
 =Z^i_{\ell\text{-}\mathrm{coh}}(S)/Z^i_{\mathrm{alg}}(S).
\end{equation}
If, moreover, $S$ is projective and is equipped with a projective embedding,
set

}{E7II-XX-015}\medskip
\noindent\textbf{(5.0.5)} $Z^{i,\mathrm{prim}}(S)$ equal to the subgroup
of $Z^i(S)$ consisting of primitive cycles (that is, cycles whose
cohomology class is primitive), and

\noindent\textbf{(5.0.6)}
$\mathrm{Prim}^{2i}_{\mathrm{alg}}(S,\mathbb Q_\ell(i))$ equal to the
$\mathbb Q$-vector subspace of $H^{2i}(S,\mathbb Q_\ell(i))$ spanned by
the cohomology classes of primitive algebraic cycles.

\SGAvar{\subsubsection*{Corollary 5.1.}
Under the hypotheses of Theorem~4.1, one has
\begin{equation}\tag{5.1.1}
\dim_{\mathbb Q}\!\left(
  \mathrm{Grif}^n(X_{\bar\eta})\otimes_{\mathbb Z}\mathbb Q
\right)
\geq
\dim_Q\!\left(
  \mathrm{Prim}^{2n}_{\mathrm{alg}}(X,\mathbb Q_\ell(n))
\right).
\end{equation}
Indeed, by (4.1.5), the kernel of the composite map
\begin{equation*}
\mathbb Q\otimes_{\mathbb Z}Z^{n,\mathrm{prim}}(X)
 \xrightarrow{\ \mathrm{res.}\ \text{to }X_{\bar\eta}\ }
\mathbb Q\otimes_{\mathbb Z}Z^n_{\ell\text{-}\mathrm{coh}}
  (X_{\bar\eta})
 \longrightarrow \mathrm{Grif}^n(X_{\bar\eta})
\end{equation*}
is contained in
  $Z^n_{\ell\text{-}\mathrm{coh}}(X)\otimes_{\mathbb Z}\mathbb Q$.
Consequently,
  $\mathrm{Prim}^{2n}_{\mathrm{alg}}(X,\mathbb Q_\ell(n))$ is a quotient
of a subspace of
  $\mathrm{Grif}^n(X_{\bar\eta})\otimes_{\mathbb Z}\mathbb Q$,
which proves the assertion.

\subsubsection*{Remark 5.2.}
For $X$ a quadric hypersurface in $\mathbb P^{2m+1}$, one has
\begin{equation*}
\dim_{\mathbb Q_\ell}\!\left(
  \mathrm{Prim}^{2m}_{\mathrm{alg}}(X,\mathbb Q_\ell(n))
\right.=1.
\end{equation*}


In particular (4.2), for $X$ a quadric in $\mathbb P^5$, the section
of $X$ by a general hypersurface of degree $>5$ contains a curve which is
$Q_\ell$-homologous to zero but of which no multiple is
algebraically equivalent to zero.

}{\subsubsection*{Corollary 5.1.}
Under the hypotheses of Theorem~4.1, one has
\begin{equation}\tag{5.1.1}
\dim_{\mathbb Q}\!\left(
  \mathrm{Grif}^m_\ell(X_{\bar\eta})\otimes_{\mathbb Z}\mathbb Q
\right)
\geq
\dim_{\mathbb Q}\!\left(
  \mathrm{Prim}^{2m}_{\mathrm{alg}}(X,\mathbb Q_\ell(m))
\right).
\end{equation}
Indeed, by (4.1.5), the kernel of the composite map
\begin{equation*}
\mathbb Q\otimes_{\mathbb Z}Z^{m,\mathrm{prim}}(X)
 \xrightarrow{\ \mathrm{res.}\ \text{to }X_{\bar\eta}\ }
\mathbb Q\otimes_{\mathbb Z}Z^m_{\ell\text{-}\mathrm{coh}}
  (X_{\bar\eta})
 \longrightarrow \mathrm{Grif}^m_\ell(X_{\bar\eta})
\end{equation*}
is contained in
  $Z^m_{\ell\text{-}\mathrm{coh}}(X)\otimes_{\mathbb Z}\mathbb Q$.
Consequently,
  $\mathrm{Prim}^{2m}_{\mathrm{alg}}(X,\mathbb Q_\ell(m))$ is a quotient
of a subspace of
  $\mathrm{Grif}^m_\ell(X_{\bar\eta})\otimes_{\mathbb Z}\mathbb Q$,
which proves the assertion.

\subsubsection*{Remark 5.2.}
For $X$ a quadric hypersurface in $\mathbb P^{2m+1}$, one has
\begin{equation*}
\dim_{\mathbb Q_\ell}\!\left(
  \mathrm{Prim}^{2m}_{\mathrm{alg}}(X,\mathbb Q_\ell(m))
\right)=1.
\end{equation*}


In particular, when $m=2$, for $X$ a quadric in $\mathbb P^5$, the section
of $X$ by a general hypersurface of degree $>5$ contains a curve which is
$\mathbb Q_\ell$-homologous to zero but of which no nonzero multiple is
algebraically equivalent to zero.

}{E7II-XX-016}\SGApage{II}{\(\approx\)362}
\begin{thebibliography}{9}
\bibitem{GriffithsPeriodsXX}
P. Griffiths, ``On the Periods of Certain Rational Integrals III,''
forthcoming in \emph{Annals of Mathematics}.

\bibitem{GrothendieckBrauerIIIXX}
A. Grothendieck, ``Le Groupe de Brauer III'' (especially \S\S 9--10), in
\emph{Dix Exposés sur la Cohomologie des Schémas}, North-Holland, 1968.

\bibitem{JouanolouCohomologieXX}
J.-P. Jouanolou, \emph{Cohomologie $\ell$-adique}, thesis, Paris, 1969.

\bibitem{KleimanCyclesXX}
S. Kleiman, ``Algebraic Cycles and the Weil Conjectures,'' in
\emph{Dix Exposés sur la Cohomologie des Schémas}, North-Holland, 1968.

\bibitem{ManinCurvesXX}
J. Manin, ``Rational Points of Algebraic Curves over Function Fields,''
\emph{Translations of the American Mathematical Society}.
\end{thebibliography}


\markboth{Exposé XX}{Exposé XX}

\ifnum\SGAform=1
\clearpage\section*{Critical apparatus}\begingroup\small\raggedright\sloppy
Readings are given in the exact TeX notation of the two English witnesses. Decimal suffixes identify the individual differences in each marked passage.

\par\medskip\noindent\hypertarget{app-E7II-XX-001.1}{}\textbf{E7II-XX-001.1}\quad printed folio 342.\par
\noindent Printed reading: {\ttfamily\footnotesize \textbackslash{}quad\textbackslash{}ell\allowbreak{}}\par
\noindent Corrected reading: \emph{(empty)}\par
\noindent Source: {\ttfamily\footnotesize fr/sga7ii/\{diplomati\allowbreak{}c,current\}/expose\_XX\allowbreak{}\_body.tex;\allowbreak{} folio\allowbreak{} 342\allowbreak{}}\par

\par\medskip\noindent\hypertarget{app-E7II-XX-001.2}{}\textbf{E7II-XX-001.2}\quad printed folio 342.\par
\noindent Printed reading: {\ttfamily\footnotesize (\allowbreak{}}\par
\noindent Corrected reading: \emph{(empty)}\par
\noindent Source: {\ttfamily\footnotesize fr/sga7ii/\{diplomati\allowbreak{}c,current\}/expose\_XX\allowbreak{}\_body.tex;\allowbreak{} folio\allowbreak{} 342;\allowbreak{} formula\allowbreak{} 15\allowbreak{}}\par

\par\medskip\noindent\hypertarget{app-E7II-XX-001.3}{}\textbf{E7II-XX-001.3}\quad printed folio 342.\par
\noindent Printed reading: {\ttfamily\footnotesize (\allowbreak{}}\par
\noindent Corrected reading: \emph{(empty)}\par
\noindent Source: {\ttfamily\footnotesize fr/sga7ii/\{diplomati\allowbreak{}c,current\}/expose\_XX\allowbreak{}\_body.tex;\allowbreak{} folio\allowbreak{} 342;\allowbreak{} formula\allowbreak{} 15\allowbreak{}}\par

\par\medskip\noindent\hypertarget{app-E7II-XX-001.4}{}\textbf{E7II-XX-001.4}\quad printed folio 342.\par
\noindent Printed reading: {\ttfamily\footnotesize -\allowbreak{}}\par
\noindent Corrected reading: {\ttfamily\footnotesize +\allowbreak{}}\par
\noindent Source: {\ttfamily\footnotesize fr/sga7ii/\{diplomati\allowbreak{}c,current\}/expose\_XX\allowbreak{}\_body.tex;\allowbreak{} folio\allowbreak{} 342;\allowbreak{} formula\allowbreak{} 15\allowbreak{}}\par

\par\medskip\noindent\hypertarget{app-E7II-XX-001.5}{}\textbf{E7II-XX-001.5}\quad printed folio 342.\par
\noindent Printed reading: {\ttfamily\footnotesize )\allowbreak{}}\par
\noindent Corrected reading: \emph{(empty)}\par
\noindent Source: {\ttfamily\footnotesize fr/sga7ii/\{diplomati\allowbreak{}c,current\}/expose\_XX\allowbreak{}\_body.tex;\allowbreak{} folio\allowbreak{} 342;\allowbreak{} formula\allowbreak{} 15\allowbreak{}}\par

\par\medskip\noindent\hypertarget{app-E7II-XX-001.6}{}\textbf{E7II-XX-001.6}\quad printed folio 342.\par
\noindent Printed reading: \emph{(empty)}\par
\noindent Corrected reading: {\ttfamily\footnotesize \_r\allowbreak{}}\par
\noindent Source: {\ttfamily\footnotesize fr/sga7ii/\{diplomati\allowbreak{}c,current\}/expose\_XX\allowbreak{}\_body.tex;\allowbreak{} folio\allowbreak{} 342;\allowbreak{} formula\allowbreak{} 23\allowbreak{}}\par

\par\medskip\noindent\hypertarget{app-E7II-XX-001.7}{}\textbf{E7II-XX-001.7}\quad printed folio 342.\par
\noindent Printed reading: \emph{(empty)}\par
\noindent Corrected reading: {\ttfamily\footnotesize '\allowbreak{}}\par
\noindent Source: {\ttfamily\footnotesize fr/sga7ii/\{diplomati\allowbreak{}c,current\}/expose\_XX\allowbreak{}\_body.tex;\allowbreak{} folio\allowbreak{} 342;\allowbreak{} formula\allowbreak{} 24\allowbreak{}}\par

\par\medskip\noindent\hypertarget{app-E7II-XX-002.1}{}\textbf{E7II-XX-002.1}\quad printed folio 343.\par
\noindent Printed reading: {\ttfamily\footnotesize a\allowbreak{}}\par
\noindent Corrected reading: {\ttfamily\footnotesize 0\allowbreak{}}\par
\noindent Source: {\ttfamily\footnotesize fr/sga7ii/\{diplomati\allowbreak{}c,current\}/expose\_XX\allowbreak{}\_body.tex;\allowbreak{} folio\allowbreak{} 343;\allowbreak{} formula\allowbreak{} 31\allowbreak{}}\par

\par\medskip\noindent\hypertarget{app-E7II-XX-002.2}{}\textbf{E7II-XX-002.2}\quad printed folio 344.\par
\noindent Printed reading: {\ttfamily\footnotesize 2\allowbreak{}}\par
\noindent Corrected reading: {\ttfamily\footnotesize 1\allowbreak{}}\par
\noindent Source: {\ttfamily\footnotesize ED-EN-NONMATH-G036;\allowbreak{} F7II-XX-006.1;\allowbreak{} WITNESS-F7II-XX-006.\allowbreak{}1;\allowbreak{} fr/sga7ii/\{diplomati\allowbreak{}c,current\}/expose\_XX\allowbreak{}\_body.tex;\allowbreak{} folio\allowbreak{} 344\allowbreak{}}\par

\par\medskip\noindent\hypertarget{app-E7II-XX-002.3}{}\textbf{E7II-XX-002.3}\quad printed folio 344.\par
\noindent Printed reading: \emph{(empty)}\par
\noindent Corrected reading: {\ttfamily\footnotesize )\allowbreak{}}\par
\noindent Source: {\ttfamily\footnotesize ED-EN-FORMULA-G209;\allowbreak{} fr/sga7ii/\{diplomati\allowbreak{}c,current\}/expose\_XX\allowbreak{}\_body.tex;\allowbreak{} folio\allowbreak{} 344\allowbreak{}}\par

\par\medskip\noindent\hypertarget{app-E7II-XX-002.4}{}\textbf{E7II-XX-002.4}\quad printed folio 344.\par
\noindent Printed reading: {\ttfamily\footnotesize b\allowbreak{}}\par
\noindent Corrected reading: {\ttfamily\footnotesize a\allowbreak{}}\par
\noindent Source: {\ttfamily\footnotesize fr/sga7ii/\{diplomati\allowbreak{}c,current\}/expose\_XX\allowbreak{}\_body.tex;\allowbreak{} folio\allowbreak{} 344;\allowbreak{} formula\allowbreak{} 54\allowbreak{}}\par

\par\medskip\noindent\hypertarget{app-E7II-XX-002.5}{}\textbf{E7II-XX-002.5}\quad printed folio 345.\par
\noindent Printed reading: \emph{(empty)}\par
\noindent Corrected reading: {\ttfamily\footnotesize f\textbackslash{}circ\allowbreak{}}\par
\noindent Source: {\ttfamily\footnotesize fr/sga7ii/\{diplomati\allowbreak{}c,current\}/expose\_XX\allowbreak{}\_body.tex;\allowbreak{} folio\allowbreak{} 345;\allowbreak{} formula\allowbreak{} 60\allowbreak{}}\par

\par\medskip\noindent\hypertarget{app-E7II-XX-002.6}{}\textbf{E7II-XX-002.6}\quad printed folio 345.\par
\noindent Printed reading: {\ttfamily\footnotesize 1f\allowbreak{}}\par
\noindent Corrected reading: {\ttfamily\footnotesize 1\allowbreak{}}\par
\noindent Source: {\ttfamily\footnotesize fr/sga7ii/\{diplomati\allowbreak{}c,current\}/expose\_XX\allowbreak{}\_body.tex;\allowbreak{} folio\allowbreak{} 345;\allowbreak{} formula\allowbreak{} 60\allowbreak{}}\par

\par\medskip\noindent\hypertarget{app-E7II-XX-002.7}{}\textbf{E7II-XX-002.7}\quad printed folio 345.\par
\noindent Printed reading: \emph{(empty)}\par
\noindent Corrected reading: {\ttfamily\footnotesize f\textbackslash{}circ\allowbreak{}}\par
\noindent Source: {\ttfamily\footnotesize fr/sga7ii/\{diplomati\allowbreak{}c,current\}/expose\_XX\allowbreak{}\_body.tex;\allowbreak{} folio\allowbreak{} 345;\allowbreak{} formula\allowbreak{} 61\allowbreak{}}\par

\par\medskip\noindent\hypertarget{app-E7II-XX-002.8}{}\textbf{E7II-XX-002.8}\quad printed folio 345.\par
\noindent Printed reading: {\ttfamily\footnotesize 1f\allowbreak{}}\par
\noindent Corrected reading: {\ttfamily\footnotesize 1\allowbreak{}}\par
\noindent Source: {\ttfamily\footnotesize fr/sga7ii/\{diplomati\allowbreak{}c,current\}/expose\_XX\allowbreak{}\_body.tex;\allowbreak{} folio\allowbreak{} 345;\allowbreak{} formula\allowbreak{} 61\allowbreak{}}\par

\par\medskip\noindent\hypertarget{app-E7II-XX-002.9}{}\textbf{E7II-XX-002.9}\quad printed folio 345.\par
\noindent Printed reading: \emph{(empty)}\par
\noindent Corrected reading: {\ttfamily\footnotesize \textbackslash{}circ\allowbreak{}}\par
\noindent Source: {\ttfamily\footnotesize fr/sga7ii/\{diplomati\allowbreak{}c,current\}/expose\_XX\allowbreak{}\_body.tex;\allowbreak{} folio\allowbreak{} 345;\allowbreak{} formula\allowbreak{} 62\allowbreak{}}\par

\par\medskip\noindent\hypertarget{app-E7II-XX-003.1}{}\textbf{E7II-XX-003.1}\quad printed folio 346.\par
\noindent Printed reading: {\ttfamily\footnotesize \textbackslash{}underline\{annexe\}\allowbreak{}}\par
\noindent Corrected reading: {\ttfamily\footnotesize connected\allowbreak{}}\par
\noindent Source: {\ttfamily\footnotesize ED-EN-XX-ANNEXE;\allowbreak{} F7II-XX-011.1;\allowbreak{} 7II:613;\allowbreak{} 7II:1005;\allowbreak{} fr/sga7ii/\{diplomati\allowbreak{}c,current\}/expose\_XX\allowbreak{}\_body.tex;\allowbreak{} folio\allowbreak{} 346\allowbreak{}}\par
\noindent Editorial explanation: ED-EN-XX-ANNEXE: the French printed adjective annexe is corrected to connexe, rendered connected in native English. Retain the corrupt source adjective untranslated and underlined in diplomatic English, because selecting a determinate English predicate would add meaning; corrected English remains connected. No sic or conjectured English predicate is added.\par

\par\medskip\noindent\hypertarget{app-E7II-XX-003.2}{}\textbf{E7II-XX-003.2}\quad printed folio 346.\par
\noindent Printed reading: {\ttfamily\footnotesize X\allowbreak{}}\par
\noindent Corrected reading: {\ttfamily\footnotesize x\allowbreak{}}\par
\noindent Source: {\ttfamily\footnotesize fr/sga7ii/\{diplomati\allowbreak{}c,current\}/expose\_XX\allowbreak{}\_body.tex;\allowbreak{} folio\allowbreak{} 346;\allowbreak{} formula\allowbreak{} 80\allowbreak{}}\par

\par\medskip\noindent\hypertarget{app-E7II-XX-004.1}{}\textbf{E7II-XX-004.1}\quad printed folio 347.\par
\noindent Printed reading: {\ttfamily\footnotesize f\allowbreak{}}\par
\noindent Corrected reading: {\ttfamily\footnotesize t\_1\allowbreak{}}\par
\noindent Source: {\ttfamily\footnotesize fr/sga7ii/\{diplomati\allowbreak{}c,current\}/expose\_XX\allowbreak{}\_body.tex;\allowbreak{} folio\allowbreak{} 347;\allowbreak{} formula\allowbreak{} 95\allowbreak{}}\par

\par\medskip\noindent\hypertarget{app-E7II-XX-004.2}{}\textbf{E7II-XX-004.2}\quad printed folio 347.\par
\noindent Printed reading: {\ttfamily\footnotesize g\allowbreak{}}\par
\noindent Corrected reading: {\ttfamily\footnotesize t\_2\allowbreak{}}\par
\noindent Source: {\ttfamily\footnotesize fr/sga7ii/\{diplomati\allowbreak{}c,current\}/expose\_XX\allowbreak{}\_body.tex;\allowbreak{} folio\allowbreak{} 347;\allowbreak{} formula\allowbreak{} 96\allowbreak{}}\par

\par\medskip\noindent\hypertarget{app-E7II-XX-005.1}{}\textbf{E7II-XX-005.1}\quad printed folio 349.\par
\noindent Printed reading: {\ttfamily\footnotesize Z\allowbreak{}}\par
\noindent Corrected reading: {\ttfamily\footnotesize z\allowbreak{}}\par
\noindent Source: {\ttfamily\footnotesize fr/sga7ii/\{diplomati\allowbreak{}c,current\}/expose\_XX\allowbreak{}\_body.tex;\allowbreak{} folio\allowbreak{} 349;\allowbreak{} formula\allowbreak{} 128\allowbreak{}}\par

\par\medskip\noindent\hypertarget{app-E7II-XX-006.1}{}\textbf{E7II-XX-006.1}\quad printed folio 350.\par
\noindent Printed reading: {\ttfamily\footnotesize g\allowbreak{}}\par
\noindent Corrected reading: {\ttfamily\footnotesize q\allowbreak{}}\par
\noindent Source: {\ttfamily\footnotesize fr/sga7ii/\{diplomati\allowbreak{}c,current\}/expose\_XX\allowbreak{}\_body.tex;\allowbreak{} folio\allowbreak{} 350;\allowbreak{} formula\allowbreak{} 144\allowbreak{}}\par

\par\medskip\noindent\hypertarget{app-E7II-XX-007.1}{}\textbf{E7II-XX-007.1}\quad printed folio 352.\par
\noindent Printed reading: {\ttfamily\footnotesize 2n\allowbreak{}}\par
\noindent Corrected reading: {\ttfamily\footnotesize n\allowbreak{}}\par
\noindent Source: {\ttfamily\footnotesize fr/sga7ii/\{diplomati\allowbreak{}c,current\}/expose\_XX\allowbreak{}\_body.tex;\allowbreak{} folio\allowbreak{} 352\allowbreak{}}\par

\par\medskip\noindent\hypertarget{app-E7II-XX-008.1}{}\textbf{E7II-XX-008.1}\quad printed folio 353.\par
\noindent Printed reading: {\ttfamily\footnotesize \textbackslash{}quad\allowbreak{} }\par
\noindent Corrected reading: {\ttfamily\footnotesize /\allowbreak{}}\par
\noindent Source: {\ttfamily\footnotesize ED-7II-1018-EN;\allowbreak{} fixes/MISPRINTS\_7II.\allowbreak{}json:626,1018;\allowbreak{} fr/sga7ii/\{diplomati\allowbreak{}c,current\}/expose\_XX\allowbreak{}\_body.tex;\allowbreak{} folio\allowbreak{} 353;\allowbreak{} ED-7II-1018\allowbreak{}}\par

\par\medskip\noindent\hypertarget{app-E7II-XX-008.2}{}\textbf{E7II-XX-008.2}\quad printed folio 353.\par
\noindent Printed reading: \emph{(empty)}\par
\noindent Corrected reading: {\ttfamily\footnotesize \textbackslash{}mathbb\allowbreak{} }\par
\noindent Source: {\ttfamily\footnotesize ED-7II-1018-EN;\allowbreak{} fixes/MISPRINTS\_7II.\allowbreak{}json:626,1018;\allowbreak{} fr/sga7ii/\{diplomati\allowbreak{}c,current\}/expose\_XX\allowbreak{}\_body.tex;\allowbreak{} folio\allowbreak{} 353;\allowbreak{} ED-7II-1018\allowbreak{}}\par

\par\medskip\noindent\hypertarget{app-E7II-XX-008.3}{}\textbf{E7II-XX-008.3}\quad printed folio 353.\par
\noindent Printed reading: \emph{(empty)}\par
\noindent Corrected reading: {\ttfamily\footnotesize )\allowbreak{}}\par
\noindent Source: {\ttfamily\footnotesize fr/sga7ii/\{diplomati\allowbreak{}c,current\}/expose\_XX\allowbreak{}\_body.tex;\allowbreak{} folio\allowbreak{} 353;\allowbreak{} formula\allowbreak{} 196\allowbreak{}}\par

\par\medskip\noindent\hypertarget{app-E7II-XX-009.1}{}\textbf{E7II-XX-009.1}\quad printed folio 354.\par
\noindent Printed reading: {\ttfamily\footnotesize ;\allowbreak{}}\par
\noindent Corrected reading: {\ttfamily\footnotesize :\allowbreak{}}\par
\noindent Source: {\ttfamily\footnotesize fr/sga7ii/\{diplomati\allowbreak{}c,current\}/expose\_XX\allowbreak{}\_body.tex;\allowbreak{} folio\allowbreak{} 354;\allowbreak{} formula\allowbreak{} 204\allowbreak{}}\par

\par\medskip\noindent\hypertarget{app-E7II-XX-009.2}{}\textbf{E7II-XX-009.2}\quad printed folio 354.\par
\noindent Printed reading: \emph{(empty)}\par
\noindent Corrected reading: {\ttfamily\footnotesize \textbackslash{}bar\allowbreak{}}\par
\noindent Source: {\ttfamily\footnotesize fr/sga7ii/\{diplomati\allowbreak{}c,current\}/expose\_XX\allowbreak{}\_body.tex;\allowbreak{} folio\allowbreak{} 354;\allowbreak{} formula\allowbreak{} 205\allowbreak{}}\par

\par\medskip\noindent\hypertarget{app-E7II-XX-010.1}{}\textbf{E7II-XX-010.1}\quad printed folio 355.\par
\noindent Printed reading: {\ttfamily\footnotesize \textasciicircum{}*\allowbreak{}}\par
\noindent Corrected reading: \emph{(empty)}\par
\noindent Source: {\ttfamily\footnotesize fr/sga7ii/\{diplomati\allowbreak{}c,current\}/expose\_XX\allowbreak{}\_body.tex;\allowbreak{} folio\allowbreak{} 355;\allowbreak{} formula\allowbreak{} 221\allowbreak{}}\par

\par\medskip\noindent\hypertarget{app-E7II-XX-011.1}{}\textbf{E7II-XX-011.1}\quad printed folio 356.\par
\noindent Printed reading: \emph{(empty)}\par
\noindent Corrected reading: {\ttfamily\footnotesize \textbackslash{}operatorname\{char\}(\allowbreak{}}\par
\noindent Source: {\ttfamily\footnotesize fr/sga7ii/\{diplomati\allowbreak{}c,current\}/expose\_XX\allowbreak{}\_body.tex;\allowbreak{} folio\allowbreak{} 356;\allowbreak{} formula\allowbreak{} 238\allowbreak{}}\par

\par\medskip\noindent\hypertarget{app-E7II-XX-011.2}{}\textbf{E7II-XX-011.2}\quad printed folio 356.\par
\noindent Printed reading: \emph{(empty)}\par
\noindent Corrected reading: {\ttfamily\footnotesize )\allowbreak{}}\par
\noindent Source: {\ttfamily\footnotesize fr/sga7ii/\{diplomati\allowbreak{}c,current\}/expose\_XX\allowbreak{}\_body.tex;\allowbreak{} folio\allowbreak{} 356;\allowbreak{} formula\allowbreak{} 238\allowbreak{}}\par

\par\medskip\noindent\hypertarget{app-E7II-XX-011.3}{}\textbf{E7II-XX-011.3}\quad printed folio 356.\par
\noindent Printed reading: \emph{(empty)}\par
\noindent Corrected reading: {\ttfamily\footnotesize \textbackslash{}operatorname\{char\}(\allowbreak{}}\par
\noindent Source: {\ttfamily\footnotesize fr/sga7ii/\{diplomati\allowbreak{}c,current\}/expose\_XX\allowbreak{}\_body.tex;\allowbreak{} folio\allowbreak{} 356;\allowbreak{} formula\allowbreak{} 240\allowbreak{}}\par

\par\medskip\noindent\hypertarget{app-E7II-XX-011.4}{}\textbf{E7II-XX-011.4}\quad printed folio 356.\par
\noindent Printed reading: \emph{(empty)}\par
\noindent Corrected reading: {\ttfamily\footnotesize )\allowbreak{}}\par
\noindent Source: {\ttfamily\footnotesize fr/sga7ii/\{diplomati\allowbreak{}c,current\}/expose\_XX\allowbreak{}\_body.tex;\allowbreak{} folio\allowbreak{} 356;\allowbreak{} formula\allowbreak{} 240\allowbreak{}}\par

\par\medskip\noindent\hypertarget{app-E7II-XX-011.5}{}\textbf{E7II-XX-011.5}\quad printed folio 356.\par
\noindent Printed reading: \emph{(empty)}\par
\noindent Corrected reading: {\ttfamily\footnotesize \textbackslash{}operatorname\{char\}(\allowbreak{}}\par
\noindent Source: {\ttfamily\footnotesize fr/sga7ii/\{diplomati\allowbreak{}c,current\}/expose\_XX\allowbreak{}\_body.tex;\allowbreak{} folio\allowbreak{} 356;\allowbreak{} formula\allowbreak{} 241\allowbreak{}}\par

\par\medskip\noindent\hypertarget{app-E7II-XX-011.6}{}\textbf{E7II-XX-011.6}\quad printed folio 356.\par
\noindent Printed reading: \emph{(empty)}\par
\noindent Corrected reading: {\ttfamily\footnotesize )\allowbreak{}}\par
\noindent Source: {\ttfamily\footnotesize fr/sga7ii/\{diplomati\allowbreak{}c,current\}/expose\_XX\allowbreak{}\_body.tex;\allowbreak{} folio\allowbreak{} 356;\allowbreak{} formula\allowbreak{} 241\allowbreak{}}\par

\par\medskip\noindent\hypertarget{app-E7II-XX-012.1}{}\textbf{E7II-XX-012.1}\quad printed folio 356.\par
\noindent Printed reading: {\ttfamily\footnotesize n\allowbreak{}}\par
\noindent Corrected reading: {\ttfamily\footnotesize m\allowbreak{}}\par
\noindent Source: {\ttfamily\footnotesize fr/sga7ii/\{diplomati\allowbreak{}c,current\}/expose\_XX\allowbreak{}\_body.tex;\allowbreak{} folio\allowbreak{} 356;\allowbreak{} formula\allowbreak{} 252\allowbreak{}}\par

\par\medskip\noindent\hypertarget{app-E7II-XX-012.2}{}\textbf{E7II-XX-012.2}\quad printed folio 357.\par
\noindent Printed reading: \emph{(empty)}\par
\noindent Corrected reading: {\ttfamily\footnotesize \_\{\textbackslash{}bar\textbackslash{}eta\}\allowbreak{}}\par
\noindent Source: {\ttfamily\footnotesize ED-EN-FORMULA-G211;\allowbreak{} fr/sga7ii/\{diplomati\allowbreak{}c,current\}/expose\_XX\allowbreak{}\_body.tex;\allowbreak{} folio\allowbreak{} 357\allowbreak{}}\par

\par\medskip\noindent\hypertarget{app-E7II-XX-012.3}{}\textbf{E7II-XX-012.3}\quad printed folio 357.\par
\noindent Printed reading: {\ttfamily\footnotesize n\allowbreak{}}\par
\noindent Corrected reading: {\ttfamily\footnotesize m\allowbreak{}}\par
\noindent Source: {\ttfamily\footnotesize fr/sga7ii/\{diplomati\allowbreak{}c,current\}/expose\_XX\allowbreak{}\_body.tex;\allowbreak{} folio\allowbreak{} 357\allowbreak{}}\par

\par\medskip\noindent\hypertarget{app-E7II-XX-012.4}{}\textbf{E7II-XX-012.4}\quad printed folio 357.\par
\noindent Printed reading: {\ttfamily\footnotesize (\allowbreak{}}\par
\noindent Corrected reading: \emph{(empty)}\par
\noindent Source: {\ttfamily\footnotesize fr/sga7ii/\{diplomati\allowbreak{}c,current\}/expose\_XX\allowbreak{}\_body.tex;\allowbreak{} folio\allowbreak{} 357;\allowbreak{} formula\allowbreak{} 262\allowbreak{}}\par

\par\medskip\noindent\hypertarget{app-E7II-XX-012.5}{}\textbf{E7II-XX-012.5}\quad printed folio 357.\par
\noindent Printed reading: \emph{(empty)}\par
\noindent Corrected reading: {\ttfamily\footnotesize \_\{\textbackslash{}bar\textbackslash{}eta\}\allowbreak{}}\par
\noindent Source: {\ttfamily\footnotesize ED-EN-FORMULA-G213;\allowbreak{} fr/sga7ii/\{diplomati\allowbreak{}c,current\}/expose\_XX\allowbreak{}\_body.tex;\allowbreak{} folio\allowbreak{} 357\allowbreak{}}\par

\par\medskip\noindent\hypertarget{app-E7II-XX-012.6}{}\textbf{E7II-XX-012.6}\quad printed folio 357.\par
\noindent Printed reading: {\ttfamily\footnotesize \textbackslash{};\{\}\_\{\textbackslash{}bar\textbackslash{}eta\}\allowbreak{}}\par
\noindent Corrected reading: \emph{(empty)}\par
\noindent Source: {\ttfamily\footnotesize fr/sga7ii/\{diplomati\allowbreak{}c,current\}/expose\_XX\allowbreak{}\_body.tex;\allowbreak{} folio\allowbreak{} 357\allowbreak{}}\par

\par\medskip\noindent\hypertarget{app-E7II-XX-012.7}{}\textbf{E7II-XX-012.7}\quad printed folio 358.\par
\noindent Printed reading: \emph{(empty)}\par
\noindent Corrected reading: {\ttfamily\footnotesize -\allowbreak{}}\par
\noindent Source: {\ttfamily\footnotesize ED-EN-FORMULA-G215;\allowbreak{} fr/sga7ii/\{diplomati\allowbreak{}c,current\}/expose\_XX\allowbreak{}\_body.tex;\allowbreak{} folio\allowbreak{} 358\allowbreak{}}\par

\par\medskip\noindent\hypertarget{app-E7II-XX-013.1}{}\textbf{E7II-XX-013.1}\quad printed folio 359.\par
\noindent Printed reading: {\ttfamily\footnotesize i\allowbreak{}}\par
\noindent Corrected reading: {\ttfamily\footnotesize 1\allowbreak{}}\par
\noindent Source: {\ttfamily\footnotesize fr/sga7ii/\{diplomati\allowbreak{}c,current\}/expose\_XX\allowbreak{}\_body.tex;\allowbreak{} folio\allowbreak{} 359;\allowbreak{} formula\allowbreak{} 302\allowbreak{}}\par

\par\medskip\noindent\hypertarget{app-E7II-XX-014.1}{}\textbf{E7II-XX-014.1}\quad printed folio 360.\par
\noindent Printed reading: {\ttfamily\footnotesize n\allowbreak{}}\par
\noindent Corrected reading: {\ttfamily\footnotesize m\allowbreak{}}\par
\noindent Source: {\ttfamily\footnotesize fr/sga7ii/\{diplomati\allowbreak{}c,current\}/expose\_XX\allowbreak{}\_body.tex;\allowbreak{} folio\allowbreak{} 360;\allowbreak{} formula\allowbreak{} 315\allowbreak{}}\par

\par\medskip\noindent\hypertarget{app-E7II-XX-014.2}{}\textbf{E7II-XX-014.2}\quad printed folio 360.\par
\noindent Printed reading: \emph{(empty)}\par
\noindent Corrected reading: {\ttfamily\footnotesize )\allowbreak{}}\par
\noindent Source: {\ttfamily\footnotesize fr/sga7ii/\{diplomati\allowbreak{}c,current\}/expose\_XX\allowbreak{}\_body.tex;\allowbreak{} folio\allowbreak{} 360;\allowbreak{} formula\allowbreak{} 315\allowbreak{}}\par

\par\medskip\noindent\hypertarget{app-E7II-XX-014.3}{}\textbf{E7II-XX-014.3}\quad printed folio 360.\par
\noindent Printed reading: {\ttfamily\footnotesize (\allowbreak{}}\par
\noindent Corrected reading: \emph{(empty)}\par
\noindent Source: {\ttfamily\footnotesize fr/sga7ii/\{diplomati\allowbreak{}c,current\}/expose\_XX\allowbreak{}\_body.tex;\allowbreak{} folio\allowbreak{} 360\allowbreak{}}\par

\par\medskip\noindent\hypertarget{app-E7II-XX-014.4}{}\textbf{E7II-XX-014.4}\quad printed folio 360.\par
\noindent Printed reading: \emph{(empty)}\par
\noindent Corrected reading: {\ttfamily\footnotesize )\allowbreak{}}\par
\noindent Source: {\ttfamily\footnotesize fr/sga7ii/\{diplomati\allowbreak{}c,current\}/expose\_XX\allowbreak{}\_body.tex;\allowbreak{} folio\allowbreak{} 360\allowbreak{}}\par

\par\medskip\noindent\hypertarget{app-E7II-XX-014.5}{}\textbf{E7II-XX-014.5}\quad printed folio 360.\par
\noindent Printed reading: {\ttfamily\footnotesize 2n\allowbreak{}}\par
\noindent Corrected reading: {\ttfamily\footnotesize 2m\allowbreak{}}\par
\noindent Source: {\ttfamily\footnotesize fr/sga7ii/\{diplomati\allowbreak{}c,current\}/expose\_XX\allowbreak{}\_body.tex;\allowbreak{} folio\allowbreak{} 360;\allowbreak{} formula\allowbreak{} 329\allowbreak{}}\par

\par\medskip\noindent\hypertarget{app-E7II-XX-014.6}{}\textbf{E7II-XX-014.6}\quad printed folio 360.\par
\noindent Printed reading: {\ttfamily\footnotesize n\allowbreak{}}\par
\noindent Corrected reading: {\ttfamily\footnotesize m\allowbreak{}}\par
\noindent Source: {\ttfamily\footnotesize fr/sga7ii/\{diplomati\allowbreak{}c,current\}/expose\_XX\allowbreak{}\_body.tex;\allowbreak{} folio\allowbreak{} 360;\allowbreak{} formula\allowbreak{} 329\allowbreak{}}\par

\par\medskip\noindent\hypertarget{app-E7II-XX-014.7}{}\textbf{E7II-XX-014.7}\quad printed folio 360.\par
\noindent Printed reading: {\ttfamily\footnotesize 3\allowbreak{}}\par
\noindent Corrected reading: {\ttfamily\footnotesize 4\allowbreak{}}\par
\noindent Source: {\ttfamily\footnotesize ED-EN-NONMATH-G037;\allowbreak{} F7II-XX-034.3;\allowbreak{} 7II:645;\allowbreak{} 7II:1081;\allowbreak{} fr/sga7ii/\{diplomati\allowbreak{}c,current\}/expose\_XX\allowbreak{}\_body.tex;\allowbreak{} folio\allowbreak{} 360\allowbreak{}}\par

\par\medskip\noindent\hypertarget{app-E7II-XX-015.1}{}\textbf{E7II-XX-015.1}\quad printed folio 361.\par
\noindent Printed reading: \emph{(empty)}\par
\noindent Corrected reading: {\ttfamily\footnotesize \textbackslash{}operatorname\{char\}(\allowbreak{}}\par
\noindent Source: {\ttfamily\footnotesize fr/sga7ii/\{diplomati\allowbreak{}c,current\}/expose\_XX\allowbreak{}\_body.tex;\allowbreak{} folio\allowbreak{} 361;\allowbreak{} formula\allowbreak{} 342\allowbreak{}}\par

\par\medskip\noindent\hypertarget{app-E7II-XX-015.2}{}\textbf{E7II-XX-015.2}\quad printed folio 361.\par
\noindent Printed reading: \emph{(empty)}\par
\noindent Corrected reading: {\ttfamily\footnotesize )\allowbreak{}}\par
\noindent Source: {\ttfamily\footnotesize fr/sga7ii/\{diplomati\allowbreak{}c,current\}/expose\_XX\allowbreak{}\_body.tex;\allowbreak{} folio\allowbreak{} 361;\allowbreak{} formula\allowbreak{} 342\allowbreak{}}\par

\par\medskip\noindent\hypertarget{app-E7II-XX-016.1}{}\textbf{E7II-XX-016.1}\quad printed folio 361.\par
\noindent Printed reading: {\ttfamily\footnotesize n\allowbreak{}}\par
\noindent Corrected reading: {\ttfamily\footnotesize m\_\textbackslash{}ell\allowbreak{}}\par
\noindent Source: {\ttfamily\footnotesize fr/sga7ii/\{diplomati\allowbreak{}c,current\}/expose\_XX\allowbreak{}\_body.tex;\allowbreak{} folio\allowbreak{} 361;\allowbreak{} formula\allowbreak{} 350\allowbreak{}}\par

\par\medskip\noindent\hypertarget{app-E7II-XX-016.2}{}\textbf{E7II-XX-016.2}\quad printed folio 361.\par
\noindent Printed reading: \emph{(empty)}\par
\noindent Corrected reading: {\ttfamily\footnotesize \{\textbackslash{}mathbb\allowbreak{} }\par
\noindent Source: {\ttfamily\footnotesize ED-EN-FONT-G020;\allowbreak{} F7II-XX-038.4;\allowbreak{} fr/sga7ii/\{diplomati\allowbreak{}c,current\}/expose\_XX\allowbreak{}\_body.tex;\allowbreak{} folio\allowbreak{} 361\allowbreak{}}\par

\par\medskip\noindent\hypertarget{app-E7II-XX-016.3}{}\textbf{E7II-XX-016.3}\quad printed folio 361.\par
\noindent Printed reading: \emph{(empty)}\par
\noindent Corrected reading: {\ttfamily\footnotesize \}\allowbreak{}}\par
\noindent Source: {\ttfamily\footnotesize ED-EN-FONT-G020;\allowbreak{} F7II-XX-038.4;\allowbreak{} fr/sga7ii/\{diplomati\allowbreak{}c,current\}/expose\_XX\allowbreak{}\_body.tex;\allowbreak{} folio\allowbreak{} 361\allowbreak{}}\par

\par\medskip\noindent\hypertarget{app-E7II-XX-016.4}{}\textbf{E7II-XX-016.4}\quad printed folio 361.\par
\noindent Printed reading: {\ttfamily\footnotesize 2n\allowbreak{}}\par
\noindent Corrected reading: {\ttfamily\footnotesize 2m\allowbreak{}}\par
\noindent Source: {\ttfamily\footnotesize fr/sga7ii/\{diplomati\allowbreak{}c,current\}/expose\_XX\allowbreak{}\_body.tex;\allowbreak{} folio\allowbreak{} 361;\allowbreak{} formula\allowbreak{} 350\allowbreak{}}\par

\par\medskip\noindent\hypertarget{app-E7II-XX-016.5}{}\textbf{E7II-XX-016.5}\quad printed folio 361.\par
\noindent Printed reading: {\ttfamily\footnotesize n\allowbreak{}}\par
\noindent Corrected reading: {\ttfamily\footnotesize m\allowbreak{}}\par
\noindent Source: {\ttfamily\footnotesize fr/sga7ii/\{diplomati\allowbreak{}c,current\}/expose\_XX\allowbreak{}\_body.tex;\allowbreak{} folio\allowbreak{} 361;\allowbreak{} formula\allowbreak{} 350\allowbreak{}}\par

\par\medskip\noindent\hypertarget{app-E7II-XX-016.6}{}\textbf{E7II-XX-016.6}\quad printed folio 361.\par
\noindent Printed reading: {\ttfamily\footnotesize n\allowbreak{}}\par
\noindent Corrected reading: {\ttfamily\footnotesize m\allowbreak{}}\par
\noindent Source: {\ttfamily\footnotesize fr/sga7ii/\{diplomati\allowbreak{}c,current\}/expose\_XX\allowbreak{}\_body.tex;\allowbreak{} folio\allowbreak{} 361;\allowbreak{} formula\allowbreak{} 351\allowbreak{}}\par

\par\medskip\noindent\hypertarget{app-E7II-XX-016.7}{}\textbf{E7II-XX-016.7}\quad printed folio 361.\par
\noindent Printed reading: {\ttfamily\footnotesize n\allowbreak{}}\par
\noindent Corrected reading: {\ttfamily\footnotesize m\allowbreak{}}\par
\noindent Source: {\ttfamily\footnotesize fr/sga7ii/\{diplomati\allowbreak{}c,current\}/expose\_XX\allowbreak{}\_body.tex;\allowbreak{} folio\allowbreak{} 361;\allowbreak{} formula\allowbreak{} 351\allowbreak{}}\par

\par\medskip\noindent\hypertarget{app-E7II-XX-016.8}{}\textbf{E7II-XX-016.8}\quad printed folio 361.\par
\noindent Printed reading: {\ttfamily\footnotesize n\allowbreak{}}\par
\noindent Corrected reading: {\ttfamily\footnotesize m\_\textbackslash{}ell\allowbreak{}}\par
\noindent Source: {\ttfamily\footnotesize fr/sga7ii/\{diplomati\allowbreak{}c,current\}/expose\_XX\allowbreak{}\_body.tex;\allowbreak{} folio\allowbreak{} 361;\allowbreak{} formula\allowbreak{} 351\allowbreak{}}\par

\par\medskip\noindent\hypertarget{app-E7II-XX-016.9}{}\textbf{E7II-XX-016.9}\quad printed folio 362.\par
\noindent Printed reading: {\ttfamily\footnotesize n\allowbreak{}}\par
\noindent Corrected reading: {\ttfamily\footnotesize m\allowbreak{}}\par
\noindent Source: {\ttfamily\footnotesize fr/sga7ii/\{diplomati\allowbreak{}c,current\}/expose\_XX\allowbreak{}\_body.tex;\allowbreak{} folio\allowbreak{} 362;\allowbreak{} formula\allowbreak{} 352\allowbreak{}}\par

\par\medskip\noindent\hypertarget{app-E7II-XX-016.10}{}\textbf{E7II-XX-016.10}\quad printed folio 362.\par
\noindent Printed reading: {\ttfamily\footnotesize 2n\allowbreak{}}\par
\noindent Corrected reading: {\ttfamily\footnotesize 2m\allowbreak{}}\par
\noindent Source: {\ttfamily\footnotesize fr/sga7ii/\{diplomati\allowbreak{}c,current\}/expose\_XX\allowbreak{}\_body.tex;\allowbreak{} folio\allowbreak{} 362;\allowbreak{} formula\allowbreak{} 353\allowbreak{}}\par

\par\medskip\noindent\hypertarget{app-E7II-XX-016.11}{}\textbf{E7II-XX-016.11}\quad printed folio 362.\par
\noindent Printed reading: {\ttfamily\footnotesize n\allowbreak{}}\par
\noindent Corrected reading: {\ttfamily\footnotesize m\allowbreak{}}\par
\noindent Source: {\ttfamily\footnotesize fr/sga7ii/\{diplomati\allowbreak{}c,current\}/expose\_XX\allowbreak{}\_body.tex;\allowbreak{} folio\allowbreak{} 362;\allowbreak{} formula\allowbreak{} 353\allowbreak{}}\par

\par\medskip\noindent\hypertarget{app-E7II-XX-016.12}{}\textbf{E7II-XX-016.12}\quad printed folio 362.\par
\noindent Printed reading: {\ttfamily\footnotesize n\allowbreak{}}\par
\noindent Corrected reading: {\ttfamily\footnotesize m\_\textbackslash{}ell\allowbreak{}}\par
\noindent Source: {\ttfamily\footnotesize fr/sga7ii/\{diplomati\allowbreak{}c,current\}/expose\_XX\allowbreak{}\_body.tex;\allowbreak{} folio\allowbreak{} 362;\allowbreak{} formula\allowbreak{} 354\allowbreak{}}\par

\par\medskip\noindent\hypertarget{app-E7II-XX-016.13}{}\textbf{E7II-XX-016.13}\quad printed folio 362.\par
\noindent Printed reading: {\ttfamily\footnotesize n\allowbreak{}}\par
\noindent Corrected reading: {\ttfamily\footnotesize m\allowbreak{}}\par
\noindent Source: {\ttfamily\footnotesize fr/sga7ii/\{diplomati\allowbreak{}c,current\}/expose\_XX\allowbreak{}\_body.tex;\allowbreak{} folio\allowbreak{} 362;\allowbreak{} formula\allowbreak{} 357\allowbreak{}}\par

\par\medskip\noindent\hypertarget{app-E7II-XX-016.14}{}\textbf{E7II-XX-016.14}\quad printed folio 362.\par
\noindent Printed reading: {\ttfamily\footnotesize .\allowbreak{}}\par
\noindent Corrected reading: {\ttfamily\footnotesize )\allowbreak{}}\par
\noindent Source: {\ttfamily\footnotesize fr/sga7ii/\{diplomati\allowbreak{}c,current\}/expose\_XX\allowbreak{}\_body.tex;\allowbreak{} folio\allowbreak{} 362;\allowbreak{} formula\allowbreak{} 357\allowbreak{}}\par

\par\medskip\noindent\hypertarget{app-E7II-XX-016.15}{}\textbf{E7II-XX-016.15}\quad printed folio 362.\par
\noindent Printed reading: \emph{(empty)}\par
\noindent Corrected reading: {\ttfamily\footnotesize ,\allowbreak{}}\par
\noindent Source: {\ttfamily\footnotesize ED-EN-FORMULA-G207;\allowbreak{} fr/sga7ii/\{diplomati\allowbreak{}c,current\}/expose\_XX\allowbreak{}\_body.tex;\allowbreak{} folio\allowbreak{} 362\allowbreak{}}\par

\par\medskip\noindent\hypertarget{app-E7II-XX-016.16}{}\textbf{E7II-XX-016.16}\quad printed folio 362.\par
\noindent Printed reading: {\ttfamily\footnotesize (4.\allowbreak{}}\par
\noindent Corrected reading: {\ttfamily\footnotesize when\allowbreak{} \$m=\allowbreak{}}\par
\noindent Source: {\ttfamily\footnotesize ED-EN-FORMULA-G207;\allowbreak{} fr/sga7ii/\{diplomati\allowbreak{}c,current\}/expose\_XX\allowbreak{}\_body.tex;\allowbreak{} folio\allowbreak{} 362\allowbreak{}}\par

\par\medskip\noindent\hypertarget{app-E7II-XX-016.17}{}\textbf{E7II-XX-016.17}\quad printed folio 362.\par
\noindent Printed reading: {\ttfamily\footnotesize )\allowbreak{}}\par
\noindent Corrected reading: {\ttfamily\footnotesize \$\allowbreak{}}\par
\noindent Source: {\ttfamily\footnotesize ED-EN-FORMULA-G207;\allowbreak{} fr/sga7ii/\{diplomati\allowbreak{}c,current\}/expose\_XX\allowbreak{}\_body.tex;\allowbreak{} folio\allowbreak{} 362\allowbreak{}}\par

\par\medskip\noindent\hypertarget{app-E7II-XX-016.18}{}\textbf{E7II-XX-016.18}\quad printed folio 362.\par
\noindent Printed reading: \emph{(empty)}\par
\noindent Corrected reading: {\ttfamily\footnotesize \textbackslash{}mathbb\allowbreak{} }\par
\noindent Source: {\ttfamily\footnotesize ED-EN-FONT-G011;\allowbreak{} F7II-XX-040.6;\allowbreak{} fr/sga7ii/\{diplomati\allowbreak{}c,current\}/expose\_XX\allowbreak{}\_body.tex;\allowbreak{} folio\allowbreak{} 362\allowbreak{}}\par

\par\medskip\noindent\hypertarget{app-E7II-XX-016.19}{}\textbf{E7II-XX-016.19}\quad printed folio 362.\par
\noindent Printed reading: \emph{(empty)}\par
\noindent Corrected reading: {\ttfamily\footnotesize nonzero\allowbreak{} }\par
\noindent Source: {\ttfamily\footnotesize ED-EN-NONMATH-G050;\allowbreak{} F7II-XX-040.9;\allowbreak{} WITNESS-F7II-XX-040.\allowbreak{}9;\allowbreak{} fr/sga7ii/\{diplomati\allowbreak{}c,current\}/expose\_XX\allowbreak{}\_body.tex;\allowbreak{} folio\allowbreak{} 362\allowbreak{}}\par

\par\endgroup\normalsize\fi
\clearpage\SGAtoc{Exposé XXI}
\selectlanguage{english}
\clearpage
\SGApage{II}{363}

\begin{center}
{\Large\bfseries EXPOSÉ XXI}\par
\vspace{1.5cm}
{\Large\bfseries THE LEVEL OF THE COHOMOLOGY OF COMPLETE
INTERSECTIONS}\par
\vspace{0.8cm}
{\large by N. Katz}\par
\vspace{0.5cm}
with an Appendix by P. Deligne (\S 5)
\end{center}

\SGApage{II}{\(\approx\)364}
\vspace{1.5cm}
\section*{Contents}
\begin{description}
\item[0.] Introduction \hfill 1
\item[1.] The Hasse--Witt matrix of a complete intersection \hfill 2
\item[2.] Coniveau over a finite field of definition \hfill 14
\item[3.] Application to complete intersections \hfill 16
\item[4.] Conclusions \hfill 19
\item[5.] Integrality theorems \hfill 21
\end{description}

\clearpage
\section*{0. Introduction}

Let $X$ be a smooth projective variety of dimension $n\geq3$, let $Y(d)$
be a section by a hypersurface of degree $d$, and let
$E^{n-1}(Y(d),\mathbb Q_\ell)$ be the ``vanishing'' cohomology of $Y(d)$
(XVIII~5.4.4).  The aim of this exposé is to prove that, for $d$
sufficiently large and $Y(d)$ sufficiently general, it is not true that the
whole of $E^{n-1}(Y(d),\mathbb Q_\ell)$ ``comes'' from the groups $H^i$,
$i<n-1$, of smooth proper algebraic varieties by means of algebraic
correspondences.  In characteristic zero this follows trivially from the
Hodge bigrading on complex cohomology and from the fact that, in coherent
cohomology,
\begin{equation*}
\dim H^{n-1}(Y(d),\mathcal O_{Y(d)})\longrightarrow\infty
\qquad\text{as }d\longrightarrow\infty.
\end{equation*}


In characteristic $p$, the method again relates the coherent group
$H^{n-1}(Y(d),\mathcal O_{Y(d)})$ to the étale group
$H^{n-1}_{\mathrm{et}}(Y(d),\mathbb Q_\ell)$.  To do so, one works over a
finite base field.  The two are then linked by the zeta function of $Y(d)$:
this function is expressed in terms of $\ell$-adic cohomology, while its
``reduction modulo $p$'' is expressed in terms of coherent cohomology
(cf.~XXII~3.9).  For complete intersections this method works without
difficulty, owing to the detailed knowledge of both kinds of cohomology
(étale and coherent) in that case.  It yields the main result, Theorem~4.2,
as a consequence of a result on Hasse--Witt operations, Theorem~1.4 (see
1.0).

\SGApage{II}{\(\approx\)365}
There are heuristic arguments due to A. Grothendieck, based on the yoga of
crystalline cohomology, which make the general statement plausible for every
smooth projective $X$, by essentially the same method.

\SGAvar{Even in the special case of complete intersections, the proof makes
essential use in 2.2 of Theorem~2.1 on the integrality of the eigenvalues of
Frobenius, recently proved by P. Deligne and proved in Appendix I.  

}{Even in the special case of complete intersections, the proof makes
essential use in 2.2 of Theorem~2.1 on the integrality of the eigenvalues of
Frobenius, recently proved by P. Deligne and proved here in the Appendix
(\S 5).  

}{E7II-XXI-001}\section*{1. The Hasse--Witt matrix of a complete intersection}

\SGAvar{\subsubsection*{Definition 1.0.}
A proper scheme over a field $k$ of characteristic exponent $p$ is called
\emph{special} if the additive $p$-linear operation
\begin{equation}\tag{1.0.1}
\mathfrak F_p:\mathcal O_X\longrightarrow\mathcal O_X,
\qquad g\longmapsto g^p,
\end{equation}
induces a nilpotent operation
\begin{equation}\tag{1.0.2}
\mathfrak F_p:H^{\dim X}(X,\mathcal O_X)
\longrightarrow H^{\dim X}(X,\mathcal O_X).
\end{equation}
Equivalently, if
$M=\dim_k H^{\dim X}(X,\mathcal O_X)$, the $p^M$-linear operation
\begin{equation}\tag{1.0.3}
\mathfrak F_p^M=\mathfrak F_{p^M}:\mathcal O_X\longrightarrow\mathcal O_X,
\qquad g\longmapsto g^{p^M},
\end{equation}
annihilates $H^{\dim X}(X,\mathcal O_X)$.

}{\subsubsection*{Definition 1.0.}
A proper scheme $X$ over a field $k$ of characteristic exponent $p$ is called
\emph{special} if the additive $p$-linear operation
\begin{equation}\tag{1.0.1}
\mathfrak F_p:\mathcal O_X\longrightarrow\mathcal O_X,
\qquad g\longmapsto g^p,
\end{equation}
induces a nilpotent operation
\begin{equation}\tag{1.0.2}
\mathfrak F_p:H^{\dim X}(X,\mathcal O_X)
\longrightarrow H^{\dim X}(X,\mathcal O_X).
\end{equation}
Equivalently, if
$M=\dim_k H^{\dim X}(X,\mathcal O_X)$, the $p^M$-linear operation
\begin{equation}\tag{1.0.3}
\mathfrak F_p^M=\mathfrak F_{p^M}:\mathcal O_X\longrightarrow\mathcal O_X,
\qquad g\longmapsto g^{p^M},
\end{equation}
annihilates $H^{\dim X}(X,\mathcal O_X)$.

}{E7II-XXI-002}\SGApage{II}{\(\approx\)366}
Recall from \textup{[2]} that, for every $i\in\mathbb Z$, the homomorphism
\begin{equation*}
H^i(X,\mathcal O_X)\xrightarrow{\ \mathfrak F_p\ }H^i(X,\mathcal O_X)
\end{equation*}
induced by (1.0.1) is classically called the \emph{Hasse--Witt
homomorphism} on $H^i(X,\mathcal O_X)$.  Notice that, when $p=1$, saying
that $X$ is non-special means simply that
$H^{\dim X}(X,\mathcal O_X)\ne0$.


\subsubsection*{Theorem 1.1.}
Fix a characteristic exponent $p$, a dimension $n\geq1$, an integer
$r\geq1$, and a multidegree $(d_1,\dots,d_r)$.  If
\begin{equation}\tag{1.1.0}
\sum_{i=1}^{r}d_i\geq(n+r)+1,
\end{equation}
then, among all complete intersections of dimension $n$ in
$\mathbb P^{n+r}$ having multidegree $(d_1,\dots,d_r)$ in characteristic
$p$, a ``sufficiently general'' member is non-special in the sense of 1.0.

\medskip
\noindent\emph{Proof.}
Let
\begin{equation}\tag{1.1.1}
f:X_U\longrightarrow U
\end{equation}
denote the ``universal family'' of complete intersections in
$\mathbb P^{n+r}$ of multidegree $(d_1,\dots,d_r)$ in characteristic $p$,
together with their defining equations.  Thus $U$ is an open subset of a
product of $r$ projective spaces over $\mathbb F_p$, of impressive
dimension.

\SGApage{II}{\(\approx\)367}
The $\mathcal O_U$-module
\begin{equation}\tag{1.1.2}
R^nf_*\mathcal O_{X_U}
\end{equation}
is locally free, and its formation commutes with every base change
(XI~1.5(i)).  By (1.1.0), its rank $M$ is nonzero (XI~2.5).  When $p=1$,
this proves the theorem for every complete intersection satisfying (1.1.0).

Suppose now that $p\geq2$.  The iterated Hasse--Witt operation
\begin{equation}\tag{1.1.3}
\mathfrak F_p^M=\mathfrak F_{p^M}:
R^nf_*\mathcal O_{X_U}\longrightarrow R^nf_*\mathcal O_{X_U}
\end{equation}
is $p^M$-linear over $\mathcal O_U$.  Since
$R^nf_*\mathcal O_{X_U}$ is locally free, (1.1.3) is represented locally on
$U$ by an $(M,M)$-matrix $H$ with entries in $\mathcal O_U$.  Consequently,
\begin{equation}\tag{1.1.4}
\{u\in U\mid X_u\text{ is non-special}\}\quad\text{is open in }U.
\end{equation}
Indeed, this is the set of $u$ for which $H(u)\ne0$.  Since $U$ is
irreducible, it is therefore enough to find one value of $u$ for which
$X_u$ is non-special; we shall do so over a finite field.

\subsubsection*{Proposition 1.1.5.}
Let $X$ be a complete intersection over $\mathbb F_q$.  If $X$ has no
$\mathbb F_q$-rational points, then $X$ is non-special.

\medskip
\noindent\emph{Proof.}
It is plainly enough to prove, for a complete intersection $X$, the
implication
\begin{equation}\tag{1.1.5.1}
X\text{ special}\quad\Longrightarrow\quad
\operatorname{card}\bigl(X(\mathbb F_q)\bigr)\equiv1\pmod p.
\end{equation}
By the trace formula, proved in the following Exposé XXII, 3.2,
\begin{equation*}
\operatorname{card}\bigl(X(\mathbb F_q)\bigr)
\equiv\sum_{i\geq0}(-1)^i
\operatorname{tr}\bigl(\mathfrak F_q\mid H^i(X,\mathcal O_X)\bigr)
\pmod p.
\end{equation*}
Now $H^i(X,\mathcal O_X)=0$ unless $i=0$ or $i=n=\dim X$, while
$H^0(X,\mathcal O_X)\simeq\mathbb F_q$ (XI~1.5(iii)); hence this becomes
\begin{equation*}
\operatorname{card}\bigl(X(\mathbb F_q)\bigr)
\equiv1+(-1)^n
\operatorname{tr}\bigl(\mathfrak F_q\mid H^n(X,\mathcal O_X)\bigr)
\pmod p.
\end{equation*}
If $X$ is special, the operation of $\mathfrak F_q$ on
$H^n(X,\mathcal O_X)$ is nilpotent by definition and therefore has trace
zero.  This proves (1.1.5.1), and hence the proposition.



\SGApage{II}{368}
\SGAvar{\subsubsection*{Proposition 1.2.}
For every power $q$ of $p$, if
\begin{equation*}
\sum_{1}^{r}d_i\geq n+r+1,
\end{equation*}
there exists a complete intersection in $\mathbb P^{n+r}_{\mathbb F_q}$ of
multidegree $(d_1,\dots,d_r)$ having no $\mathbb F_q$-rational point.

\medskip
\noindent\emph{Proof.}
Begin with the case $r=1$.  For every $d\geq n+2$, we must find a
homogeneous form of degree $d$,
\begin{equation}\tag{1.2.1}
H(X_1),\dots,X_{n+2})\in
\mathbb F_q[X_1,\dots,X_{n+2}],
\end{equation}
such that
\begin{equation}\tag{1.2.2}
\begin{aligned}
&(\alpha_1,\dots,\alpha_{n+2})\in(\mathbb F_q)^{n+2},\qquad
H(\alpha_1,\dots,\alpha_{n+2})=0\quad\Longrightarrow\\
&\hspace{7em}\alpha_1=\cdots=\alpha_{n+2}=0.
\end{aligned}
\end{equation}
Set $\mathbb E=\mathbb F_{q^d}$, and choose elements
$\beta_1,\dots,\beta_{n+2}$ of $\mathbb E$ that are linearly independent
over $\mathbb F_q$.  Take
\begin{equation}\tag{1.2.3}
H(X_1,\dots,X_{n+2})=
N_{\mathbb E/\mathbb F_q}
\left(\sum_{i=1}^{n+2}\beta_iX_i\right).
\end{equation}

In the general case, since
\begin{equation}\tag{1.2.4}
\sum_{i}^{r}d_i\geq n+r+1,
\end{equation}
we can choose $r$ integers
\begin{equation}\tag{1.2.5}
1\leq a_1<a_2<\cdots<a_r=n+r+1
\end{equation}
such that
\begin{equation}\tag{1.2.6}
\left\{\begin{aligned}
d_1&\geq a_o,\\
d_2&\geq a_2-a_1,\\
d_r&\geq a_r-a_{r-1}.
\end{aligned}\right.
\end{equation}
For each $i=1,\dots,r$, choose a homogeneous form of degree $d_i$,
\begin{equation*}
H_i(X_{a_{i-1}+1},\dots,X_{a_i})
\in\mathbb F_q[X_{a_{i-1}+1},\dots,X_{a_i}],
\end{equation*}
in the $a_i-a_{i-1}$ displayed variables and having property (1.2.2).
The required complete intersection is defined by the homogeneous forms
\begin{equation*}
H_1(X_1,\dots,X_{a_1}),\quad\dots,\quad
H_r(X_{a_{r-1}},\dots,X_{a_r}).
\end{equation*}



}{\subsubsection*{Proposition 1.2.}
For every power $q$ of $p$, if
\begin{equation*}
\sum_{i=1}^{r}d_i\geq n+r+1,
\end{equation*}
there exists a complete intersection in $\mathbb P^{n+r}_{\mathbb F_q}$ of
multidegree $(d_1,\dots,d_r)$ having no $\mathbb F_q$-rational point.

\medskip
\noindent\emph{Proof.}
Begin with the case $r=1$.  For every $d\geq n+2$, we must find a
homogeneous form of degree $d$,
\begin{equation}\tag{1.2.1}
H(X_1,\dots,X_{n+2})\in
\mathbb F_q[X_1,\dots,X_{n+2}],
\end{equation}
such that
\begin{equation}\tag{1.2.2}
\begin{aligned}
&(\alpha_1,\dots,\alpha_{n+2})\in(\mathbb F_q)^{n+2},\qquad
H(\alpha_1,\dots,\alpha_{n+2})=0\quad\Longrightarrow\\
&\hspace{7em}\alpha_1=\cdots=\alpha_{n+2}=0.
\end{aligned}
\end{equation}
Set $\mathbb E=\mathbb F_{q^d}$, and choose elements
$\beta_1,\dots,\beta_{n+2}$ of $\mathbb E$ that are linearly independent
over $\mathbb F_q$.  Take
\begin{equation}\tag{1.2.3}
H(X_1,\dots,X_{n+2})=
N_{\mathbb E/\mathbb F_q}
\left(\sum_{i=1}^{n+2}\beta_iX_i\right).
\end{equation}

In the general case, since
\begin{equation}\tag{1.2.4}
\sum_{i=1}^{r}d_i\geq n+r+1,
\end{equation}
we can set $a_0=0$ and choose $r$ integers
\begin{equation}\tag{1.2.5}
1\leq a_1<a_2<\cdots<a_r=n+r+1
\end{equation}
such that
\begin{equation}\tag{1.2.6}
d_i\geq a_i-a_{i-1}\qquad(1\leq i\leq r).
\end{equation}
For each $i=1,\dots,r$, choose a homogeneous form of degree $d_i$,
\begin{equation*}
H_i(X_{a_{i-1}+1},\dots,X_{a_i})
\in\mathbb F_q[X_{a_{i-1}+1},\dots,X_{a_i}],
\end{equation*}
in the $a_i-a_{i-1}$ displayed variables and having property (1.2.2).
The required complete intersection is defined by the homogeneous forms
\begin{equation*}
H_1(X_1,\dots,X_{a_1}),\quad\dots,\quad
H_r(X_{a_{r-1}+1},\dots,X_{a_r}).
\end{equation*}



}{E7II-XXI-003}\SGApage{II}{370}
\subsubsection*{Theorem 1.3.}
Let $X$ be an irreducible complete intersection of dimension $n\geq1$ over
a field $k$ of characteristic exponent $p$.  For every integer $d\geq n+1$,
a sufficiently general section of $X$ by a hypersurface of degree $d$ is
non-special in the sense of 1.0.

\SGAvar{\medskip
\noindent\emph{Proof.}
The sections $X.H_u$ of $X$ by hypersurfaces of degree $d$ are parametrized
by the lines in $H^0(X,\mathcal O_X(d))$, that is, by
\begin{equation}\tag{1.3.1}
U=\mathbb P\bigl(H^0(X,\mathcal O_X(d)).
\end{equation}
Over $U$ we have the ``universal section of $X$ by a hypersurface of degree
$d$,''
\begin{equation}\tag{1.3.2}
f:X.H_u\longrightarrow U.
\end{equation}
Let $U'$ be the open subset of those $u\in U$ for which
$\dim(X_u)=n-1$.  By XI~1.5(i) and 2.5, the restriction
\begin{equation}\tag{1.3.3}
\left.R^{n-1}f_*\mathcal O_{X.H_u}\right|_{U'}
\end{equation}
is locally free, its formation commutes with every base change, and its rank
$M$ is at least one.  If $p=1$, this proves the theorem for every section by
a hypersurface of degree at least $n+1$.

Suppose therefore that $p>0$.  The iterated Hasse--Witt operation
$\mathfrak F_{p^n}$ is again represented locally by an $M\times M$ matrix
with entries in $\mathcal O_u$.  Consequently,
\begin{equation}\tag{1.3.4}
\{u\in U'\mid X.H_u\text{ is non-special}\}
\quad\text{is open in }U'.
\end{equation}
It remains only to prove that the set in (1.3.4) is nonempty.

First consider the case
\begin{equation}\tag{1.3.5}
k=\mathbb F_q.
\end{equation}
After replacing $\mathbb F_q$ by a finite extension, we may suppose that
\begin{equation}\tag{1.3.6}
\operatorname{card}\bigl(X(\mathbb F_q)\bigr)>0
\end{equation}
and that
\begin{equation}\tag{1.3.7}
\begin{aligned}
&\text{there are homogeneous coordinates}\\
&\qquad X_1,\dots,X_{n+r+1}\text{ in the ambient space }\mathbb P_k^{n+r}\\
&\text{such that }X\text{ does not meet the linear subvariety}\\
&\qquad X_1=\cdots=X_{n+1}=0.
\end{aligned}
\end{equation}
Choose a homogeneous form of degree $d$,
\begin{equation}\tag{1.3.8}
H(X_1,\dots,X_{n+1})\in
\mathbb F_q[X_1,\dots,X_{n+1}],
\end{equation}
having property (1.2.2), which is possible because $d\geq n+1$.  Let $H$
also denote the hypersurface in $\mathbb P^{n+r}$ defined by
$H(X_1,\dots,X_{n+1})=0$.  By construction, $X.H$ has no
$\mathbb F_q$-rational point.  In view of (1.3.6), this also shows that
$X\not\subset H$; since $X$ is irreducible, $X.H$ has dimension $n-1$, so
$H\in U'$.  Proposition~1.5 now shows that $X.H$ is non-special.



}{\medskip
\noindent\emph{Proof.}
The sections $X.H_u$ of $X$ by hypersurfaces of degree $d$ are parametrized
by the lines in $H^0(X,\mathcal O_X(d))$, that is, by
\begin{equation}\tag{1.3.1}
U=\mathbb P\bigl(H^0(X,\mathcal O_X(d))\bigr).
\end{equation}
Over $U$ we have the ``universal section of $X$ by a hypersurface of degree
$d$,''
\begin{equation}\tag{1.3.2}
f:X.H_u\longrightarrow U.
\end{equation}
Let $U'$ be the open subset of those $u\in U$ for which
$\dim(X_u)=n-1$.  By XI~1.5(i) and 2.5, the restriction
\begin{equation}\tag{1.3.3}
\left.R^{n-1}f_*\mathcal O_{X.H_u}\right|_{U'}
\end{equation}
is locally free, its formation commutes with every base change, and its rank
$M$ is at least one.  If $p=1$, this proves the theorem for every section by
a hypersurface of degree at least $n+1$.

Suppose therefore that $p>1$.  The iterated Hasse--Witt operation
$\mathfrak F_{p^M}$ is again represented locally by an $M\times M$ matrix
with entries in $\mathcal O_U$.  Consequently,
\begin{equation}\tag{1.3.4}
\{u\in U'\mid X.H_u\text{ is non-special}\}
\quad\text{is open in }U'.
\end{equation}
It remains only to prove that the set in (1.3.4) is nonempty.

First consider the case
\begin{equation}\tag{1.3.5}
k=\mathbb F_q.
\end{equation}
After replacing $\mathbb F_q$ by a finite extension, we may suppose that
\begin{equation}\tag{1.3.6}
\operatorname{card}\bigl(X(\mathbb F_q)\bigr)>0
\end{equation}
and that
\begin{equation}\tag{1.3.7}
\begin{aligned}
&\text{there are homogeneous coordinates}\\
&\qquad X_1,\dots,X_{n+r+1}\text{ in the ambient space }\mathbb P_k^{n+r}\\
&\text{such that }X\text{ does not meet the linear subvariety}\\
&\qquad X_1=\cdots=X_{n+1}=0.
\end{aligned}
\end{equation}
Choose a homogeneous form of degree $d$,
\begin{equation}\tag{1.3.8}
H(X_1,\dots,X_{n+1})\in
\mathbb F_q[X_1,\dots,X_{n+1}],
\end{equation}
having property (1.2.2), which is possible because $d\geq n+1$.  Let $H$
also denote the hypersurface in $\mathbb P^{n+r}$ defined by
$H(X_1,\dots,X_{n+1})=0$.  By construction, $X.H$ has no
$\mathbb F_q$-rational point.  In view of (1.3.6), this also shows that
$X\not\subset H$; since $X$ is irreducible, $X.H$ has dimension $n-1$, so
$H\in U'$.  Proposition~1.1.5 now shows that $X.H$ is non-special.



}{E7II-XXI-004}\SGApage{II}{\(\approx\)372}
\par\medskip
We now pass to the general case.  Looking at the equations defining $X$, we
may suppose that $k$ is finitely generated over $\mathbb F_p$.  There is
then a subring $B\subset k$, finitely generated as an algebra over
$\mathbb F_p$ and having fraction field $k$, such that
\begin{equation}\tag{1.3.9}
X\subset\mathbb P_k^{n+r}
\end{equation}
is the generic fiber of a relative complete intersection
$\widetilde X\subset\mathbb P_B^{n+r}$ which, fiber by fiber, has the same
multidegree as $X$ and is geometrically irreducible.

\SGAvar{Write
\begin{equation}\tag{1.3.10}
\rho:\widetilde X\longrightarrow S=\operatorname{Spec}(B)
\end{equation}
for the projection,
\begin{equation}\tag{1.3.11}
E=\rho_*\mathcal O_{\widetilde X}(d)
\end{equation}
for the locally free $\mathcal O_S$-module, and
\begin{equation}\tag{1.3.12}
\pi:T=\check{\mathbb P}(E)\longrightarrow S
\end{equation}
for the associated bundle of lines in $\mathbb E$.  The generic fiber of (1.3.12)
is precisely
\[
U=\mathbb P\bigl(H^0(X,\mathcal O_X(d)).
\]
Over $T$ we have the universal section of
$\widetilde X\subset\mathbb P_B^{n+r}$ by hypersurfaces of degree $d$,
denoted
\begin{equation}\tag{1.3.13}
\alpha:\widetilde X.H_T\longrightarrow T.
\end{equation}
These constructions fit into the following diagram of cartesian squares:
\begin{equation}\tag{1.3.14}
\begin{tikzcd}[column sep=huge,row sep=large]
\widetilde X.H_T \arrow[d,"\alpha"]
  \& X.H_u=(\widetilde X.H_T)\times_U T \arrow[l]\arrow[d] \\
T \arrow[d,"\pi"]
  \& U=T\times_S\operatorname{Sp}(k) \arrow[l]\arrow[d] \\
S \& \operatorname{Sp}(k)\arrow[l]
\end{tikzcd}
\end{equation}
The $\mathcal O_T$-module
\begin{equation}\tag{1.3.15}
R^{n-1}\alpha_*\mathcal O_{\widetilde X.H_T}
\end{equation}
is locally free over the open subset $T'\subset T$ consisting of the
$t$ for which $\dim\alpha^{-1}(t)=n-1$, and its formation there commutes
with every base change.  Since $T$ and $U$ have the same generic point, it
is enough, in order to show that (1.3.4) is nonempty, to prove that the open
subset of $T'$
\begin{equation}\tag{1.3.16}
\bigl\{t\in T'\mid\widetilde X_{\pi(t)}.H_t
   \text{ is non-special}\bigr\}
\end{equation}
is nonempty.  But the already established case (1.3.5) shows that, for
every closed point $s\in S$, the fiber $\pi^{-1}(s)$ has nonempty
intersection with the open set (1.3.16).  This proves the theorem.



\subsubsection*{Theorem 1.4.}
Let $X\subset\mathbb P^N$ be a projective Cohen--Macaulay scheme,
equidimensional of dimension $n\geq1$, over a field $k$.  There is an
integer $d_o$ such that, for every $d\geq d_o$, if
$i:X.H\longrightarrow X$ denotes the inclusion of a section by a general
hypersurface of degree $d$, then the action of
\begin{equation}\tag{1.4.0}
\mathfrak F_p\quad\text{on}\quad
\frac{H^{n-1}(X.H,\mathcal O_{X.H})}
     {i^*H^{n-1}(X,\mathcal O_X)}
\end{equation}
is non-nilpotent.

\medskip
\noindent\emph{Proof.}
As above, we reduce to the case $k=\mathbb F_q$ and to the problem of
finding just one $H$---for any fixed sufficiently large degree for which
$\dim H^{n-1}(X.H,\mathcal O_{X.H})$ is independent of $H$; see
(1.4.4)---for which (1.4.0) holds.  After enlarging $q$, we may suppose
that the irreducible components $X_i$ of $X$ are geometrically irreducible,
that
\begin{equation}\tag{1.4.1}
\operatorname{card}\bigl(X_i(\mathbb F_q)\bigr)\geq2,
\end{equation}
and that
\begin{equation}\tag{1.4.2}
\begin{aligned}
&\text{there are projective coordinates }X_o,\dots,X_N\text{ such that}\\
&\qquad X\text{ does not meet the linear subvariety }
X_o=\cdots=X_n=0.
\end{aligned}
\end{equation}
By a vanishing theorem of Serre ([1], or SGA~2 XII~1.4), there is an
integer $d_1$ such that, for every $d\geq d_1$,
\begin{equation}\tag{1.4.3}
H^i(X,\mathcal O_X(-d))=0\qquad\text{if }i\ne n=\dim X.
\end{equation}
Thus, if $H$ is a hypersurface of degree $d\geq d_1$, then (1.4.3) gives
\begin{equation}\tag{1.4.4}
\begin{aligned}
H^i(X,\mathcal O_X)&\xrightarrow{\ \sim\ }
 H^i(X.H,\mathcal O_{X.H}) &&\text{if }i\leq n-2,\\[1ex]
H^{n-1}(X,\mathcal O_X)&\lhook\joinrel\longrightarrow
 H^{n-1}(X.H,\mathcal O_{X.H}) &&\text{injectively}.
\end{aligned}
\end{equation}

We prove the theorem with $d_o=\max(N+1,d_1)$.  For every $d\geq d_o$,
we must find a hypersurface $H$ of degree $d$, defined over $\mathbb F_q$,
which contains no $X_i$ and for which
\begin{equation}\tag{1.4.5}
\operatorname{tr}\!\left(\mathfrak F_q\ \middle|\
\frac{H^{n-1}(X.H,\mathcal O_{X.H})}
     {H^{n-1}(X,\mathcal O_X)}\right)\ne0.
\end{equation}
By the trace formula (XX~3.2.1), modulo $p$ we have
\begin{equation}\tag{1.4.6}
\begin{aligned}
\operatorname{card}\bigl((X.H)(\mathbb F_q)\bigr)
&=\sum_{i=0}^{n-1}(-1)^i
 \operatorname{tr}\bigl(\mathfrak F_q\mid
 H^i(X.H,\mathcal O_{X.H})\bigr)\\
&\stackrel{(1.4.4)}{=}
 \sum_{i=0}^{n-1}(-1)^i
 \operatorname{tr}\bigl(\mathfrak F_q\mid H^i(X,\mathcal O_X)\bigr)\\
&\qquad+(1)^{n-1}\operatorname{tr}\!\left(\mathfrak F_q\ \middle|\
 \left(\frac{H^{n-1}(X.H,\mathcal O_{X.H})}
      {H^{n-1}(X,\mathcal O_X)}\right.\right..
\end{aligned}
\end{equation}
Let $a\in\mathbb Z/p\mathbb Z$ denote the quantity
\[
a=\sum_{i=0}^{n-1}(-1)^i
\operatorname{tr}\bigl(\mathfrak F_q\mid H^i(X,\mathcal O_X)\bigr).
\]
It is therefore enough, for every $d\geq d_o$, to find a hypersurface $H$
of degree $d$, defined over $\mathbb F_q$, which does not contain $X$ and
such that
\begin{equation}\tag{1.4.7}
\operatorname{card}\bigl((X.H)(\mathbb F_q)\bigr)\not\equiv a\pmod p.
\end{equation}

}{Write
\begin{equation}\tag{1.3.10}
\rho:\widetilde X\longrightarrow S=\operatorname{Spec}(B)
\end{equation}
for the projection,
\begin{equation}\tag{1.3.11}
E=\rho_*\mathcal O_{\widetilde X}(d)
\end{equation}
for the locally free $\mathcal O_S$-module, and
\begin{equation}\tag{1.3.12}
\pi:T=\check{\mathbb P}(E)\longrightarrow S
\end{equation}
for the associated bundle of lines in $\mathbb E$.  The generic fiber of (1.3.12)
is precisely
\[
U=\mathbb P\bigl(H^0(X,\mathcal O_X(d))\bigr).
\]
Over $T$ we have the universal section of
$\widetilde X\subset\mathbb P_B^{n+r}$ by hypersurfaces of degree $d$,
denoted
\begin{equation}\tag{1.3.13}
\alpha:\widetilde X.H_T\longrightarrow T.
\end{equation}
These constructions fit into the following diagram of cartesian squares:
\begin{equation}\tag{1.3.14}
\begin{tikzcd}[column sep=huge,row sep=large]
\widetilde X.H_T \arrow[d,"\alpha"]
  \& X.H_u=(\widetilde X.H_T)\times_T U \arrow[l]\arrow[d] \\
T \arrow[d,"\pi"]
  \& U=T\times_S\operatorname{Sp}(k) \arrow[l]\arrow[d] \\
S \& \operatorname{Sp}(k)\arrow[l]
\end{tikzcd}
\end{equation}
The $\mathcal O_T$-module
\begin{equation}\tag{1.3.15}
R^{n-1}\alpha_*\mathcal O_{\widetilde X.H_T}
\end{equation}
is locally free over the open subset $T'\subset T$ consisting of the
$t$ for which $\dim\alpha^{-1}(t)=n-1$, and its formation there commutes
with every base change.  Since $T$ and $U$ have the same generic point, it
is enough, in order to show that (1.3.4) is nonempty, to prove that the open
subset of $T'$
\begin{equation}\tag{1.3.16}
\bigl\{t\in T'\mid\widetilde X_{\pi(t)}.H_t
   \text{ is non-special}\bigr\}
\end{equation}
is nonempty.  But the already established case (1.3.5) shows that, for
every closed point $s\in S$, the fiber $\pi^{-1}(s)$ has nonempty
intersection with the open set (1.3.16).  This proves the theorem.



\subsubsection*{Theorem 1.4.}
Let $X\subset\mathbb P^N$ be a projective Cohen--Macaulay scheme,
equidimensional of dimension $n\geq1$, over a field $k$.  There is an
integer $d_0$ such that, for every $d\geq d_0$, if
$i:X.H\longrightarrow X$ denotes the inclusion of a section by a general
hypersurface of degree $d$, then the action of
\begin{equation}\tag{1.4.0}
\mathfrak F_p\quad\text{on}\quad
\frac{H^{n-1}(X.H,\mathcal O_{X.H})}
     {i^*H^{n-1}(X,\mathcal O_X)}
\end{equation}
is non-nilpotent.

\medskip
\noindent\emph{Proof.}
As above, we reduce to the case $k=\mathbb F_q$ and to the problem of
finding just one $H$---for any fixed sufficiently large degree for which
$\dim H^{n-1}(X.H,\mathcal O_{X.H})$ is independent of $H$; see
(1.4.4)---for which (1.4.0) holds.  After enlarging $q$, we may suppose
that the irreducible components $X_i$ of $X$ are geometrically irreducible,
that
\begin{equation}\tag{1.4.1}
\operatorname{card}\bigl(X_i(\mathbb F_q)\bigr)\geq2,
\end{equation}
and that
\begin{equation}\tag{1.4.2}
\begin{aligned}
&\text{there are projective coordinates }X_0,\dots,X_N\text{ such that}\\
&\qquad X\text{ does not meet the linear subvariety }
X_0=\cdots=X_n=0.
\end{aligned}
\end{equation}
By a vanishing theorem of Serre ([1], or SGA~2 XII~1.4), there is an
integer $d_1$ such that, for every $d\geq d_1$,
\begin{equation}\tag{1.4.3}
H^i(X,\mathcal O_X(-d))=0\qquad\text{if }i\ne n=\dim X.
\end{equation}
Thus, if $H$ is a hypersurface of degree $d\geq d_1$, then (1.4.3) gives
\begin{equation}\tag{1.4.4}
\begin{aligned}
H^i(X,\mathcal O_X)&\xrightarrow{\ \sim\ }
 H^i(X.H,\mathcal O_{X.H}) &&\text{if }i\leq n-2,\\[1ex]
H^{n-1}(X,\mathcal O_X)&\lhook\joinrel\longrightarrow
 H^{n-1}(X.H,\mathcal O_{X.H}) &&\text{injectively}.
\end{aligned}
\end{equation}

We prove the theorem with $d_0=\max(N+1,d_1)$.  For every $d\geq d_0$,
we must find a hypersurface $H$ of degree $d$, defined over $\mathbb F_q$,
which contains no $X_i$ and for which
\begin{equation}\tag{1.4.5}
\operatorname{tr}\!\left(\mathfrak F_q\ \middle|\
\frac{H^{n-1}(X.H,\mathcal O_{X.H})}
     {H^{n-1}(X,\mathcal O_X)}\right)\ne0.
\end{equation}
By the trace formula (XXII~3.2.1), modulo $p$ we have
\begin{equation}\tag{1.4.6}
\begin{aligned}
\operatorname{card}\bigl((X.H)(\mathbb F_q)\bigr)
&\equiv\sum_{i=0}^{n-1}(-1)^i
 \operatorname{tr}\bigl(\mathfrak F_q\mid
 H^i(X.H,\mathcal O_{X.H})\bigr)\\
&\stackrel{(1.4.4)}{\equiv}
 \sum_{i=0}^{n-1}(-1)^i
 \operatorname{tr}\bigl(\mathfrak F_q\mid H^i(X,\mathcal O_X)\bigr)\\
&\qquad+(-1)^{n-1}\operatorname{tr}\!\left(\mathfrak F_q\ \middle|\
 \left(\frac{H^{n-1}(X.H,\mathcal O_{X.H})}
      {H^{n-1}(X,\mathcal O_X)}\right)\right).
\end{aligned}
\end{equation}
Let $a\in\mathbb Z/p\mathbb Z$ denote the quantity
\[
a=\sum_{i=0}^{n-1}(-1)^i
\operatorname{tr}\bigl(\mathfrak F_q\mid H^i(X,\mathcal O_X)\bigr).
\]
It is therefore enough, for every $d\geq d_0$, to find a hypersurface $H$
of degree $d$, defined over $\mathbb F_q$, which does not contain $X$ and
such that
\begin{equation}\tag{1.4.7}
\operatorname{card}\bigl((X.H)(\mathbb F_q)\bigr)\not\equiv a\pmod p.
\end{equation}

}{E7II-XXI-005}\SGApage{II}{\(\approx\)376}
There are two cases.  If $a\ne0$, it is enough to take an $H$ over
$\mathbb F_q$ for which $(X.H)(\mathbb F_q)$ is empty, just as in (1.3.8).
Such an $H$ satisfies (1.4.7) and contains no $X_i$, because
$\operatorname{card}(X_i(\mathbb F_q))\geq2$ by hypothesis.

\SGAvar{If $a=0$, it is enough to find an $H$ such that
$\operatorname{card}(X.H)(\mathbb F_q)=1$.  Then (1.4.7) holds and
$X_i\not\subset H$, since $\operatorname{card}(X_i)(\mathbb F_q)\geq2$.
Choose projective coordinates $X_o,\dots,X_N$ such that $(1,0,\dots,0)$ is
a point of $X$.  Take $H$ to have equation $H(X_1,\dots,X_N)=0$, where $H$
is a homogeneous form satisfying (cf. (1.2.2))
\begin{equation}\tag{1.4.8}
\alpha_1,\dots,\alpha_N\in\mathbb F_q,\quad
H(\alpha_1,.,\alpha_N)=0
\quad\Longrightarrow\quad
\alpha_1=\cdots=\alpha_N=0.
\end{equation}
The hypersurface defined by $H$ in $\mathbb P^N$ then has only the point
$(1,0,\dots,0)$ over $\mathbb F_q$; a fortiori, $X.H$ has only this point
over $\mathbb F_q$.  This proves the theorem.



}{If $a=0$, it is enough to find an $H$ such that
$\operatorname{card}((X.H)(\mathbb F_q))=1$.  Then (1.4.7) holds and
$X_i\not\subset H$, since $\operatorname{card}(X_i(\mathbb F_q))\geq2$.
Choose projective coordinates $X_0,\dots,X_N$ such that $(1,0,\dots,0)$ is
a point of $X$.  Take $H$ to have equation $H(X_1,\dots,X_N)=0$, where $H$
is a homogeneous form satisfying (cf. (1.2.2))
\begin{equation}\tag{1.4.8}
\alpha_1,\dots,\alpha_N\in\mathbb F_q,\quad
H(\alpha_1,\dots,\alpha_N)=0
\quad\Longrightarrow\quad
\alpha_1=\cdots=\alpha_N=0.
\end{equation}
The hypersurface defined by $H$ in $\mathbb P^N$ then has only the point
$(1,0,\dots,0)$ over $\mathbb F_q$; a fortiori, $X.H$ has only this point
over $\mathbb F_q$.  This proves the theorem.



}{E7II-XXI-006}\SGApage{II}{\(\approx\)377}
\section*{2. Coniveau over a finite field of definition}

\SGAvar{\subsection*{2.0.}
Let $X$ be a separated scheme of finite type over $\mathbb F_q$.  For every
prime number $\ell\ne\operatorname{char}(\mathbb F_q)$, write
\begin{equation}\tag{2.0.1}
H_c^i(\overline X,\mathbb Q_\ell)
\end{equation}
for the compactly supported cohomology groups of
$\overline X=X\otimes_{\mathbb F_q}\overline{\mathbb F}_q$.  These are
finite-dimensional $\mathbb Q_\ell$-vector spaces on which the Galois group
$\operatorname{Gal}(/\mathbb F_q)$ acts.  Let
\begin{equation}\tag{2.0.2}
\varphi_q\in\operatorname{Gal}(/\mathbb F_q)
\end{equation}
denote the Frobenius automorphism.  It is known (SGA~5 XIII~3.2.1) that the
zeta function of $X$ (XXII~3.0.1) is given by
\begin{equation}\tag{2.0.3}
Z(t,X/\mathbb F_q)=
\prod_{i\geq0}
\det\!\left(1-t\varphi_q^{-1}\mid
H_c^i(\overline X,\mathbb Q_\ell)\right)^{(-1)^{i+1}}.
\end{equation}
Deligne has recently proved, generalizing a result of Washnitzer for smooth
projective $X$, the following theorem, which is proved in the appendix below
(5.5.3).

}{\subsection*{2.0.}
Let $X$ be a separated scheme of finite type over $\mathbb F_q$.  For every
prime number $\ell\ne\operatorname{char}(\mathbb F_q)$, write
\begin{equation}\tag{2.0.1}
H_c^i(\overline X,\mathbb Q_\ell)
\end{equation}
for the compactly supported cohomology groups of
$\overline X=X\otimes_{\mathbb F_q}\overline{\mathbb F}_q$.  These are
finite-dimensional $\mathbb Q_\ell$-vector spaces on which the Galois group
$\operatorname{Gal}(\overline{\mathbb F}_q/\mathbb F_q)$ acts.  Let
\begin{equation}\tag{2.0.2}
\varphi_q\in\operatorname{Gal}(\overline{\mathbb F}_q/\mathbb F_q)
\end{equation}
denote the Frobenius automorphism.  It is known (SGA~5 XIII~3.2.1) that the
zeta function of $X$ (XXII~3.0.1) is given by
\begin{equation}\tag{2.0.3}
Z(t,X/\mathbb F_q)=
\prod_{i\geq0}
\det\!\left(1-t\varphi_q^{-1}\mid
H_c^i(\overline X,\mathbb Q_\ell)\right)^{(-1)^{i+1}}.
\end{equation}
Deligne has recently proved, generalizing a result of Washnitzer for smooth
projective $X$, the following theorem, which is proved in the appendix below
(5.5.3).

}{E7II-XXI-007}\subsubsection*{Theorem 2.1.}
The eigenvalues of $\varphi_q^{-1}$ on
$H_c^i(\overline X,\mathbb Q_\ell)$ are algebraic integers; if
$i>\dim X$, they are divisible, as algebraic integers, by
$q^{,i-\dim X}$.

In particular, tautologically, we obtain:

\SGApage{II}{\(\approx\)378}
\medskip
\noindent\textbf{(2.1 bis).}
The eigenvalues of $\varphi_q^{-1}$ on
$H_c^i(\overline X,\mathbb Q_\ell(-j))$ are algebraic integers divisible by
$q^j$.

\subsubsection*{Corollary 2.2.}
Let $X$ be proper, smooth, and geometrically connected over $\mathbb F_q$.
Then the eigenvalues of $\varphi_q^{-1}$ on
$N^jH^q(\overline X,\mathbb Q_\ell)$ (notation of XX~2.0.6) are algebraic
integers divisible by $q^j$.

\SGAvar{\medskip
\noindent\emph{Proof.}
The assertion is invariant under a finite extension
$\mathbb F_{q^M}/\mathbb F_q$, which amounts to replacing the eigenvalues
under consideration by their $M$-th powers.  After making such a base
change, the definition gives finitely many projective, smooth, geometrically
connected schemes $T_i$ over $\mathbb F_q$, of dimensions $d_i$, and
algebraic cycles $Z_i$ on $T_i\times_{\mathbb F_q}X$ of codimension
$j+d_i$, such that the homomorphism of
$(/\mathbb F_q)$-modules
\begin{equation*}
 H^{q-2j}(\overline T_i,\mathbb Q_\ell(-j)
\xrightarrow{\ \oplus\operatorname{Cl}(Z_i)\ }
N^jH^q(\overline X,\mathbb Q_\ell)
\end{equation*}
is surjective.  The conclusion follows from (2.1 bis).



}{\medskip
\noindent\emph{Proof.}
The assertion is invariant under a finite extension
$\mathbb F_{q^M}/\mathbb F_q$, which amounts to replacing the eigenvalues
under consideration by their $M$-th powers.  After making such a base
change, the definition gives finitely many projective, smooth, geometrically
connected schemes $T_i$ over $\mathbb F_q$, of dimensions $d_i$, and
algebraic cycles $Z_i$ on $T_i\times_{\mathbb F_q}X$ of codimension
$j+d_i$, such that the homomorphism of
$\operatorname{Gal}(\overline{\mathbb F}_q/\mathbb F_q)$-modules
\begin{equation*}
\bigoplus_i H^{q-2j}(\overline T_i,\mathbb Q_\ell(-j))
\xrightarrow{\ \oplus_i\operatorname{Cl}(Z_i)\ }
N^jH^q(\overline X,\mathbb Q_\ell)
\end{equation*}
is surjective.  The conclusion follows from (2.1 bis).



}{E7II-XXI-008}\SGApage{II}{\(\approx\)379}
\section*{3. Application to complete intersections}

\subsection*{3.0.}
Let $X$ be a smooth complete intersection of dimension $n\geq1$ in
$\mathbb P^N$ over a field $k$.  For
$\ell\ne\operatorname{char}(k)$, weak Lefschetz and duality give
(XI~1.6)
\begin{equation}\tag{3.0.1}
H^i(\overline X,\mathbb Q_\ell)=
\begin{cases}
0 & \text{if $i$ is odd and $i\ne n$},\\[1.5ex]
\mathbb Q_\ell\!\left(-\dfrac{i}{2}\right)
  & \text{if $i$ is even, $i\leq2n$, and $i\ne n$},\\[1.5ex]
H^n(\overline X,\mathbb Q_\ell) & \text{if $i=n$}.
\end{cases}
\end{equation}
For $k=\mathbb F_q$ this gives the following formula for the zeta function:
\begin{equation}\tag{3.0.2}
Z(t,X/\mathbb F_q)\prod_{i=0}^{n}(1-q^it)
=\det\!\left(1-t\varphi_q^{-1}\mid
\operatorname{Prim}^n(\overline X,\mathbb Q_\ell)\right)^{(-1)^{n+1}}.
\end{equation}
Here we have used the primitive decomposition and (3.0.1).

\subsubsection*{Corollary 3.0.3.}
Under the hypotheses of 3.0, with $k=\mathbb F_q$, there is a congruence
modulo $p$
\[
\det\!\left(1-t\varphi_q^{-1}\mid
H^n(\overline X,\mathbb Q_\ell)\right)
\equiv
\det\!\left(1-t\mathfrak F_q\mid H^n(X,\mathcal O_X)\right).
\]

\SGAvar{\medskip
\noindent\emph{Proof.}
By the congruence formula (XXII~3.2.1), and since
\begin{equation}\tag{3.0.4}
H^i(X,\mathcal O_X)\simeq
\begin{cases}
\mathbb F_q & i=0,\\[1.5ex]
0 & 1\leq\;\leq n-1,\\[1.5ex]
0 & i\geq n+1,
\end{cases}
\end{equation}
we have
\begin{equation}\tag{3.0.5}
Z(t,X/\mathbb F_q)\equiv
(1-t)^{-1}
\det\!\left(1-t\mathfrak F_q\mid H^n(X,\mathcal O_X)\right)^{(-1)^{n+1}}
\pmod p.
\end{equation}
Consequently, (3.0.5) and (3.0.2) give
\begin{equation}\tag{3.0.6}
\det\!\left(1-t\varphi_q^{-1}\mid
\operatorname{Prim}^n(\overline X,\mathbb Q_\ell)\right)
\equiv
\det\!\left(1-t\mathfrak F_q\mid H^n(X,\mathcal O_X)\right)
\pmod p.
\end{equation}
The desired conclusion follows because
\begin{equation}\tag{3.0.7}
H^n(\overline X,\mathbb Q_\ell)\simeq
\begin{cases}
\operatorname{Prim}^n(\overline X,\mathbb Q_\ell)
  & \text{if $n$ is odd},\\[2mm]
\operatorname{Prim}^n(\overline X,\mathbb Q_\ell)
 \oplus\mathbb Q_\ell\!\left(-\dfrac n2\right)
  & \text{if $n$ is even}.
\end{cases}
\end{equation}

}{\medskip
\noindent\emph{Proof.}
By the congruence formula (XXII~3.2.1), and since
\begin{equation}\tag{3.0.4}
H^i(X,\mathcal O_X)\simeq
\begin{cases}
\mathbb F_q & i=0,\\[1.5ex]
0 & 1\leq i\leq n-1,\\[1.5ex]
0 & i\geq n+1,
\end{cases}
\end{equation}
we have
\begin{equation}\tag{3.0.5}
Z(t,X/\mathbb F_q)\equiv
(1-t)^{-1}
\det\!\left(1-t\mathfrak F_q\mid H^n(X,\mathcal O_X)\right)^{(-1)^{n+1}}
\pmod p.
\end{equation}
Consequently, (3.0.5) and (3.0.2) give
\begin{equation}\tag{3.0.6}
\det\!\left(1-t\varphi_q^{-1}\mid
\operatorname{Prim}^n(\overline X,\mathbb Q_\ell)\right)
\equiv
\det\!\left(1-t\mathfrak F_q\mid H^n(X,\mathcal O_X)\right)
\pmod p.
\end{equation}
The desired conclusion follows because
\begin{equation}\tag{3.0.7}
H^n(\overline X,\mathbb Q_\ell)\simeq
\begin{cases}
\operatorname{Prim}^n(\overline X,\mathbb Q_\ell)
  & \text{if $n$ is odd},\\[2mm]
\operatorname{Prim}^n(\overline X,\mathbb Q_\ell)
 \oplus\mathbb Q_\ell\!\left(-\dfrac n2\right)
  & \text{if $n$ is even}.
\end{cases}
\end{equation}

}{E7II-XXI-009}\SGApage{II}{380}
\subsubsection*{Corollary 3.1.}
Let $X$ be a smooth complete intersection of dimension $n\geq1$ over
$\mathbb F_q$.  If
\begin{equation}\tag{3.1.0}
H^n(\overline X,\mathbb Q_\ell)=N^1H^n(\overline X,\mathbb Q_\ell),
\end{equation}
then
\begin{equation}\tag{3.1.1}
\mathfrak F_p\quad\text{is nilpotent on}\quad H^n(X,\mathcal O_X).
\end{equation}

\medskip
\noindent\emph{Proof.}
By Corollary~2.2, (3.1.0) implies
\begin{equation}\tag{3.1.2}
\det\!\left(1-t\varphi_q^{-1}\mid
H^n(\overline X,\mathbb Q_\ell)\right)
\equiv1\pmod p
\qquad\text{(indeed, even modulo $q$)}.
\end{equation}
Together with Corollary~3.0.3, this gives
\begin{equation}\tag{3.1.3}
\det\!\left(1-t\mathfrak F_q\mid H^n(X,\mathcal O_X)\right)=1.
\end{equation}
Thus $\mathfrak F_q$ is nilpotent on $H^n(X,\mathcal O_X)$, and hence so is
$\mathfrak F_p$.  This proves the corollary.



\SGApage{II}{381}
\subsubsection*{Proposition 3.2.}
Let $X$ be a smooth complete intersection of dimension $n\geq1$ over a
field $k$.  If
\begin{equation}\tag{3.2.1}
H^n(\overline X,\mathbb Q_\ell)=N^1H^n(\overline X,\mathbb Q_\ell),
\end{equation}
then $X$ is special; that is,
\begin{equation}\tag{3.2.2}
\mathfrak F_p\quad\text{is nilpotent on}\quad H^n(X,\mathcal O_X).
\end{equation}

\SGAvar{\medskip
\noindent\emph{Proof.}
We may suppose that $k$ is finitely generated over $\mathbb F_p$ (the case
$p=0$ being trivial), and that there is a finitely generated
$\mathbb F_p$-subalgebra $B$ of $k$ such that $X/k$ is the generic fiber of a
smooth relative complete intersection $\widetilde X$:
\begin{equation}\tag{3.2.3}
\begin{tikzcd}[column sep=4.5em,row sep=3em]
\widetilde X \arrow[d,"\pi"'] \&
X=\widetilde X/_\eta \arrow[l] \arrow[d] \\
\operatorname{Spec}(B)=S \& \eta=\operatorname{Sp}(k)\ . \arrow[l]
\end{tikzcd}
\end{equation}
The $B$-module $R^n\pi_*(\mathcal O_{\widetilde X})$ is locally free, and its
formation commutes with every base change.  Consequently the set of points
\begin{equation}\tag{3.2.4}
\{\,s\in S\mid X_s\text{ is special}\,\}
\end{equation}
is closed in $S$.

}{\medskip
\noindent\emph{Proof.}
We may suppose that $k$ is finitely generated over $\mathbb F_p$ (the case
$p=0$ being trivial), and that there is a finitely generated
$\mathbb F_p$-subalgebra $B$ of $k$ such that $X/k$ is the generic fiber of a
smooth relative complete intersection $\widetilde X$:
\begin{equation}\tag{3.2.3}
\begin{tikzcd}[column sep=4.5em,row sep=3em]
\widetilde X \arrow[d,"\pi"'] \&
X=\widetilde X_\eta \arrow[l] \arrow[d] \\
\operatorname{Spec}(B)=S \& \eta=\operatorname{Sp}(k)\ . \arrow[l]
\end{tikzcd}
\end{equation}
The $B$-module $R^n\pi_*(\mathcal O_{\widetilde X})$ is locally free, and its
formation commutes with every base change.  Consequently the set of points
\begin{equation}\tag{3.2.4}
\{\,s\in S\mid X_s\text{ is special}\,\}
\end{equation}
is closed in $S$.

}{E7II-XXI-010}\SGApage{II}{\(\approx\)382}
By (XX~2.3), after shrinking $S$ we may suppose that
\begin{equation}\tag{3.2.5}
\forall s\in S,\qquad
H^n(X_{\bar s},\mathbb Q_\ell)=N^1H^n(X_{\bar s},\mathbb Q_\ell).
\end{equation}
By Corollary~3.1, $X_s$ is special for every closed point $s$ of $S$.
It follows that the closed subset (3.2.4) is all of $S$ and hence, in
particular, that $X=\widetilde X_\eta$ is special.



\section*{4. Conclusions}

Combining Theorem~1.1 with Proposition~3.2 gives the following result.

\SGAvar{\subsubsection*{Theorem 4.1.}
Fix a characteristic $p$, a dimension $n$, and a multidegree
$(d_1,\ldots,d_r)$.  Let
\begin{equation}\tag{4.1.0}
X_V\longrightarrow V
\end{equation}
denote the universal smooth complete intersection of dimension $n$ and
multidegree $(d_1,\ldots,d_r)$ in characteristic $p$.  If
\begin{equation}\tag{4.1.1}
\sum_{1}^{r}d_i\geq n+r+1,
\end{equation}
then there is a dense open subset $U\subset V$ such that
\begin{equation}\tag{4.1.2}
\text{for every }u\in U,\qquad
H^n(X_{\bar u},\mathbb Q_\ell)\ne N^1H^n(X_{\bar u},\mathbb Q_\ell).
\end{equation}

}{\subsubsection*{Theorem 4.1.}
Fix a characteristic $p$, a dimension $n$, and a multidegree
$(d_1,\ldots,d_r)$.  Let
\begin{equation}\tag{4.1.0}
X_V\longrightarrow V
\end{equation}
denote the universal smooth complete intersection of dimension $n$ and
multidegree $(d_1,\ldots,d_r)$ in characteristic $p$.  If
\begin{equation}\tag{4.1.1}
\sum_{i=1}^{r}d_i\geq n+r+1,
\end{equation}
then there is a dense open subset $U\subset V$ such that
\begin{equation}\tag{4.1.2}
\text{for every }u\in U,\qquad
H^n(X_{\bar u},\mathbb Q_\ell)\ne N^1H^n(X_{\bar u},\mathbb Q_\ell).
\end{equation}

}{E7II-XXI-011}\SGApage{II}{\(\approx\)383}
Combining Theorem~1.3 with Proposition~3.2 likewise gives the following.

\subsubsection*{Theorem 4.2.}
Let $X$ be a smooth complete intersection of dimension $n\geq1$ over a
field $k$.  Fix an integer
\begin{equation}\tag{4.2.1}
d\geq n+1.
\end{equation}
Let
\begin{equation}\tag{4.2.2}
X\!\cdot\!H_u\longrightarrow U
\end{equation}
denote the universal family of sections of $X$ by hypersurfaces of degree
$d$ that are smooth of dimension $n-1$.  There is a dense open subset
$T\subset U$ such that
\begin{equation}\tag{4.2.3}
\text{for every }t\in T,\qquad
H^{n-1}(X\!\cdot\!H_{\bar t},\mathbb Q_\ell)
\ne N^1H^{n-1}(X\!\cdot\!H_{\bar t},\mathbb Q_\ell).
\end{equation}

\subsection*{Bibliography}

\begin{enumerate}[label={\textup{[\arabic*]}},leftmargin=2.6em]
\item J.-P. Serre, \emph{Faisceaux algébriques cohérents}, Annals of
  Mathematics \textbf{61} (1955), 197--278.
\item J.-P. Serre, \emph{Sur la topologie des variétés algébriques en
  caractéristique $p$}, in \emph{Symposium Internacional de Topología
  Algebraica}, Mexico, 1958.
\end{enumerate}



\SGApage{II}{363}
\section*{5. Appendix: Integrality Theorems}

\SGApage{II}{\(\approx\)384}
\begin{center}
By P. Deligne
\end{center}

\subsection*{5.0.}
Let $p$ be a prime number, $q$ a power of $p$, $\mathbb F_q$ a field with
$q$ elements, $E$ the completion of a number field at a place dividing a
prime number $\ell\ne p$, and $T$ a set of prime numbers (for example,
$T=\varnothing$).

An element $\alpha$ of an algebraic closure of $E$ is called
\emph{$T$-integral} if it is algebraic over $\mathbb Q$ and integral over
$\mathbb Z[(1/r)_{r\in T}]$.  An endomorphism $F$ of a finite-dimensional
$E$-vector space is called $T$-integral if its eigenvalues are $T$-integral.

Let $k$ be a finite field with $q'$ elements, let $\bar k$ be an algebraic
closure of $k$, and let
\[
\mathfrak G=\operatorname{Gal}(\bar k/k)
\]
be the Galois group of $k$.  We denote by $f\in\mathfrak G$ the Frobenius
substitution
\[
f(x)=x^{q'},
\]
and by $F$ the ``geometric'' Frobenius, the inverse of $f$ in $\mathfrak G$.
A representation $\rho$ of $\mathfrak G$ on a finite-dimensional
$E$-vector space is called $T$-integral if $\rho(F)$ is $T$-integral.

If $X$ is a scheme of finite type over $\mathbb F_q$ and $x$ is a closed
point of $X$, we denote by $F_x$ the geometric Frobenius element of the
Galois group of $k(x)$.  If $\bar x$ is a geometric point of $X$ lying over
$x$ and $\underline G$ is an $E$-sheaf on $X$, then $F_x$ acts, by transport
of structure, on the fiber $G_{\bar x}$.

\subsubsection*{Definition 5.1.}
A constructible $E$-sheaf $\underline G$ on a scheme $X$ of finite type over
$\mathbb F_q$ is called $T$-integral if, for every geometric point $\bar x$
of $X$ lying over a closed point $x$ of $X$, $F_x$ is a $T$-integral
automorphism of $\underline G_{\bar x}$.

\SGApage{II}{\(\approx\)385}
If $a:X\longrightarrow\operatorname{Spec}(\mathbb F_q)$ is a scheme of
finite type over $\mathbb F_q$ and $\underline G$ is a constructible
$E$-sheaf on $X$, we write $\underline H^*(X,\underline G)$ for the
$E$-sheaves $R^*a_*\underline G$ on
$\operatorname{Spec}(\mathbb F_q)_{\acute e t}$.  If
$\overline{\mathbb F}_q$ is an algebraic closure of $\mathbb F_q$ and
$\overline X=X\times_{\mathbb F_q}\overline{\mathbb F}_q$, these sheaves
identify with the cohomology groups $H^*(\overline X,\underline G)$ regarded
as Galois modules, that is, endowed with an action of the geometric
Frobenius $F$.  We use the same convention for cohomology with proper
support.

\subsubsection*{Lemma 5.2.1.}
Let $\underline G$ be a constructible $E$-sheaf on a scheme $X$ of finite
type over $\mathbb F_q$, of dimension at most $n$.

\begin{enumerate}[label=(\roman*)]
\item There is a closed subset $Y$ of $X$, of dimension $0$, such that the
  evident map
  \[
  \underline H^0(X,\underline G)\longrightarrow
  \underline H^0(Y,\underline G)
  \]
  is injective.
\item If $X$ is separated and $U$ is an open subset of $X$ whose complement
  $Z$ has dimension less than $n$, there is a closed subset $Y$ of $U$, of
  dimension $0$, and a surjective map
  \[
  \underline H^0(Y,\underline G)(-n)\longrightarrow
  \underline H_c^{2n}(X,\underline G).
  \]
\end{enumerate}

\medskip
\noindent\emph{Proof.}
Assertion (i) is evident.  We prove (ii).  After replacing $X$ by
$X_{\mathrm{red}}$ and shrinking $U$, we may suppose that $U$ is smooth and
that $\underline G|_U$ is twisted constant.  Since $\dim Z<n$, we have
\[
0=\underline H_c^{2n-1}(Z,\underline G)
\longrightarrow \underline H_c^{2n}(U,\underline G)
\xrightarrow{\sim}\underline H_c^{2n}(X,\underline G)
\longrightarrow \underline H_c^{2n}(Z,\underline G)=0,
\]
and Poincar\'e duality (SGA~4, Expos\'e~XVIII, and [2]) gives
\[
\underline H_c^{2n}(U,\underline G)\simeq
\underline H^0(U,\underline G^\vee(n))^\vee.
\]
The conclusion follows from (i).


\SGApage{II}{386}
\SGAvar{\subsubsection*{Theorem 5.2.2.}
Let $X$ be a separated scheme of finite type over $\mathbb F_q$, of
dimension at most $n$, and let $\underline G$ be a constructible
$T$-integral $E$-sheaf.  Then the endomorphisms $F$ and $q^{n-i}F$ of
$H_c^i(X,\underline G)$ are $T$-integral.

\medskip
\noindent\emph{Proof.}
The case $\dim X\leq0$ is trivial, and the cases $i=0$ and $i=2n$ follow
from Lemma~5.2.1.  Let $Z(X,\underline G,t)$ be the formal power series
\[
Z(X,\underline G,t)=
\prod_{x\in X}\det\!\left(1-F_xt^{\deg(x)}\,;\,\underline G_x\right)^{-1}
\in F[[t]].
\]
By SGA~5, Expos\'e~XV, 3.2, one has
\begin{equation}\tag{5.2.3}
Z(X,\underline G,t)=
\prod_i\det\!\left(1-Ft\,;\,\underline H_c^i(X,\underline G)\right)^{(-1)^{i+1}}.
\end{equation}

}{\subsubsection*{Theorem 5.2.2.}
Let $X$ be a separated scheme of finite type over $\mathbb F_q$, of
dimension at most $n$, and let $\underline G$ be a constructible
$T$-integral $E$-sheaf.  Then the endomorphisms $F$ and $q^{n-i}F$ of
$\underline H_c^i(X,\underline G)$ are $T$-integral.

\medskip
\noindent\emph{Proof.}
The case $\dim X\leq0$ is trivial, and the cases $i=0$ and $i=2n$ follow
from Lemma~5.2.1.  Let $Z(X,\underline G,t)$ be the formal power series
\[
Z(X,\underline G,t)=
\prod_{x\in X}\det\!\left(1-F_xt^{\deg(x)}\,;\,\underline G_x\right)^{-1}
\in E[[t]].
\]
By SGA~5, Expos\'e~XV, 3.2, one has
\begin{equation}\tag{5.2.3}
Z(X,\underline G,t)=
\prod_i\det\!\left(1-Ft\,;\,\underline H_c^i(X,\underline G)\right)^{(-1)^{i+1}}.
\end{equation}

}{E7II-XXI-012}We first prove the theorem for $n=1$.  Then
$\underline H_c^i(X,\underline G)=0$ for $i>2$ (SGA~4, Expos\'e~X, 4.3,
or Expos\'e~XVII, 5.2.8.1).  By hypothesis the coefficients of
$Z(X,\underline G,t)$ are $T$-integral.  By the cases $i=0$ and $i=2n=2$
already established, the product of $Z(X,\underline G,t)$ with
\[
\det\!\left(1-Ft\,;\,\underline H_c^0(X,\underline G)\right)
\det\!\left(1-Ft\,;\,\underline H_c^2(X,\underline G)\right)
\]
still has $T$-integral coefficients.  This product is precisely
$\det(1-Ft\,;\,\underline H_c^1(X,\underline G))$, proving the theorem for
$n=1$.

We now proceed by induction on $n=\dim X$.  Let $U$ be a Zariski-open subset
of $X$ whose complement $Z$ has dimension less than $n$, and suppose there
is a morphism $f:U\longrightarrow Y$ to a curve, with fibers of dimension
at most $n-1$:
\begin{center}
\begin{tikzcd}
U \arrow[r,hook,"j"] \arrow[d,"f"'] \& X \& Z \arrow[l,hook',"i"'] \\
Y
\end{tikzcd}
\qquad .
\end{center}
The exact sequence
\[
0\longrightarrow j_!j^*\underline G\longrightarrow\underline G
\longrightarrow i_!i^*\underline G\longrightarrow0
\]
induces exact sequences
\[
\underline H_c^i(U,\underline G)\longrightarrow
\underline H_c^i(X,\underline G)\longrightarrow
\underline H_c^i(Z,\underline G).
\]
By induction it is therefore enough to verify the theorem for
$\underline G$ on $U$.  Applying the induction hypothesis to the fibers of
$f$, together with proper base change, shows that the sheaves
$R^if_!\underline G$ and $R^if_!\underline G(i-(n-1))$ on $Y$ are
$T$-integral.  The Leray spectral sequence for $f$ is
\[
E_2^{pq}=\underline H_c^p(Y,R^qf_!\underline G)
\Longrightarrow\underline H_c^{p+q}(U,\underline G).
\]
By what precedes and by the theorem already proved for $Y$, both
$E_2^{pq}$ and $E_2^{pq}(p+q-n)$ are $T$-integral.  The theorem follows.

\SGApage{II}{387}
\subsubsection*{Corollary 5.3.}
Let $X$ be a scheme of finite type over $\mathbb F_q$, of dimension at most
$1$, let $U$ be an open subset of $X$, let $j:U\hookrightarrow X$ be the
inclusion, and let $\underline G$ be a constructible $E$-sheaf on $U$.

\begin{enumerate}[label=(\alph*)]
\item If $\underline G$ is $T$-integral, then $j_*\underline G$ is
  $T$-integral.
\item If $G_x^\vee$ is $T$-integral for every $x\in U$, then
  $(j_*\underline G)_x^\vee$ is $T$-integral for every $x\in X$.
\end{enumerate}

\SGApage{II}{\(\approx\)388}
\SGAvar{\medskip
\noindent\emph{Proof.}
We reduce first to the case where $X$ is affine, and then to the case where
$X$ is projective.  Define $\underline H$ by the exact sequence
\[
0\longrightarrow j_!\underline G\longrightarrow j_*\underline G
\longrightarrow\underline H\longrightarrow0.
\]
If $Y$ is the finite support of $\underline H$, the corresponding
cohomology sequence gives an exact sequence
\[
\underline H^0(U,\underline G)\longrightarrow
\underline H^0(Y,\underline H)\longrightarrow
\underline H_c^1(U,\underline G).
\]
Lemma~5.2.1 and Theorem~5.2 show that the two outer terms are
$T$-integral.  Hence $\underline H^0(Y,\underline H)$ is $T$-integral,
which proves (a).

Let $X'$ be the normalization of the pure one-dimensional part of
$X_{\mathrm{red}}$, and choose an open subset $U'$ of $U$, dense in the
smooth one-dimensional locus of $U_{\mathrm{red}}$, such that
$\underline G|_U$ is twisted constant:
\begin{center}
\begin{tikzcd}
U' \arrow[r,hook,"j'"'] \arrow[d,hook] \& X' \arrow[d,hook,"p"] \& {} \\
U \arrow[r,hook,"j"'] \& X \& Y \arrow[l,hook']
\end{tikzcd}
\qquad ,
\end{center}
Then
\[
(p_*j'_*\underline G)|_Y\simeq(j_*\underline G)|_Y.
\]
It is therefore enough to treat the case in which $X$ is smooth of
dimension one and $\underline G|_U$ is twisted constant.  The sheaf
$\underline G^\vee$ is then $T$-integral, so we already know that
$(j_*\underline G^\vee)_x$ is $T$-integral for every $x\in X$.

}{\medskip
\noindent\emph{Proof.}
We reduce first to the case where $X$ is affine, and then to the case where
$X$ is projective.  Define $\underline H$ by the exact sequence
\[
0\longrightarrow j_!\underline G\longrightarrow j_*\underline G
\longrightarrow\underline H\longrightarrow0.
\]
If $Y$ is the finite support of $\underline H$, the corresponding
cohomology sequence gives an exact sequence
\[
\underline H^0(U,\underline G)\longrightarrow
\underline H^0(Y,\underline H)\longrightarrow
\underline H_c^1(U,\underline G).
\]
Lemma~5.2.1 and Theorem~5.2.2 show that the two outer terms are
$T$-integral.  Hence $\underline H^0(Y,\underline H)$ is $T$-integral,
which proves (a).

Let $X'$ be the normalization of the pure one-dimensional part of
$X_{\mathrm{red}}$, and choose an open subset $U'$ of $U$, dense in the
smooth one-dimensional locus of $U_{\mathrm{red}}$, such that
$\underline G|_{U'}$ is twisted constant:
\begin{center}
\begin{tikzcd}
U' \arrow[r,hook,"j'"'] \arrow[d,hook] \& X' \arrow[d,hook,"p"] \& {} \\
U \arrow[r,hook,"j"'] \& X \& Y \arrow[l,hook']
\end{tikzcd}
\qquad ,
\end{center}
Then
\[
(p_*j'_*\underline G)|_Y\simeq(j_*\underline G)|_Y.
\]
It is therefore enough to treat the case in which $X$ is smooth of
dimension one and $\underline G|_U$ is twisted constant.  The sheaf
$\underline G^\vee$ is then $T$-integral, so we already know that
$(j_*\underline G^\vee)_x$ is $T$-integral for every $x\in X$.

}{E7II-XXI-013}\SGApage{II}{\(\approx\)389}
Let $\mathfrak G_x$ denote the Galois group of the fraction field of
$\widehat{\mathcal O}_{X,x}$, and let $\mathfrak I_x$ be the inertia group.
The sheaf $\underline G$ defines a representation of $\mathfrak G_x$, again
denoted by $\underline G$.  The $\mathfrak G_x/\mathfrak I_x=
\operatorname{Gal}(\overline{k(x)}/k(x))$-module
$(j_*\underline G^\vee)_x$ (respectively,
$(j_*\underline G)_x^\vee$) identifies with the module of invariants
(respectively, coinvariants) in $\underline G^\vee$.  It remains to prove:





\SGAvar{\subsubsection*{Lemma 5.3.1.}
Let $K$ be a local field with residue field $\mathbb F_q$, let $G$ be its
Galois group, let $I$ be the inertia group, and let $\rho$ be a continuous
representation of $G$ on an $n$-dimensional $E$-vector space $V$ (cf. 5.0).
Let $V^I$ and $V_I$ denote respectively the space of invariants and the
space of coinvariants under $I$, regarded as $G/I$-modules.  Let $\alpha_1,\ldots,\alpha_n$ be the eigenvalues of
$F$ on $V^I$.  Then there are integers $m_i\geq0$ such that the eigenvalues
of $F$ on $V_I$ are
\[
\alpha_iq^{m_i}.
\]

}{\subsubsection*{Lemma 5.3.1.}
Let $K$ be a local field with residue field $\mathbb F_q$, let $G$ be its
Galois group, let $I$ be the inertia group, and let $\rho$ be a continuous
representation of $G$ on an $n$-dimensional $E$-vector space $V$ (cf. 5.0).
Let $V^I$ and $V_I$ denote respectively the space of invariants and the
space of coinvariants under $I$, regarded as $G/I$-modules.  Put
$r=\dim_E V^I$, and let $\alpha_1,\ldots,\alpha_r$ be the eigenvalues of
$F$ on $V^I$.  Then there are integers $m_i\geq0$ such that the eigenvalues
of $F$ on $V_I$ are
\[
\alpha_iq^{m_i}\qquad(1\leq i\leq r).
\]

}{E7II-XXI-014}\medskip
\noindent\emph{Proof.}
Let $I'$ be the largest subgroup of $I$ of order prime to $\ell$.  The group
$\rho(I')$ is finite, since it is a compact subgroup of $\operatorname{GL}(V)$
of order prime to $\ell$.  Hence
\[
V_I\simeq\bigl(V^{\rho(I')}\bigr)_I.
\]
After replacing $V$ by $V^{I'}$, we may therefore suppose that $I'$ acts
trivially on $V$.  As $G/I$-modules one has
\[
I/I'\simeq\mathbb Z_\ell(1).
\]
Let $T$ be a topological generator of $I/I'$.  For each $\alpha$, let
$V^{(\alpha)}$ be the generalized eigenspace of $T$ on $V$ with eigenvalue
$\alpha$.  Then
\[
V=\bigoplus_\alpha V^{(\alpha)},
\]
and the decomposition
\[
V=V^{(1)}\oplus\bigoplus_{\alpha\ne1}V^{(\alpha)}
\]
is stable under $G$ and compatible with duality.  We may thus replace $V$
by $V^{(1)}$ and suppose that $I$ acts unipotently.

\SGApage{II}{390}
\SGAvar{\subsubsection*{Lemma 5.3.2.}
Let $V$ be a finite-dimensional vector space over a field $k$ of
characteristic $0$, let $I$ be a one-parameter additive group, let
$\rho$ $I\longrightarrow\operatorname{GL}(V)$ be a representation, and let
$N$ be an infinitesimal generator of $I$.  Set
\[
V^{I,k}=V^I\cap\operatorname{Im}(N^k)
\]
and
\[
V_{I,k}=\operatorname{Im}\!\left(\operatorname{Ker}(N^{k+1})
\longrightarrow V_I\right).
\]
Then $N^k$ induces an isomorphism
\[
N^k:V_{I,k}/V_{I,k-1}\xrightarrow{\sim}
V^{I,k}/V^{I,k+1},
\]
and defines canonical isomorphisms
\[
V_{I,k}/V_{I,k-1}\otimes I^{\otimes k}
\simeq V^{I,k}/V^{I,k+1}.
\]

\medskip
\noindent\emph{Proof.}
Via $N$, representations $V$ of $I$ identify with $k[T]$-modules killed by
a power of $T$.  It is enough to verify Lemma~3.2 when $V$ corresponds to
an irreducible $k[T]$-module $k[T]/T^k$, where the assertion is immediate.

}{\subsubsection*{Lemma 5.3.2.}
Let $V$ be a finite-dimensional vector space over a field $k$ of
characteristic $0$, let $I$ be a one-parameter additive group, let
$\rho:I\longrightarrow\operatorname{GL}(V)$ be a representation, and let
$N$ be an infinitesimal generator of $I$.  Set
\[
V^{I,k}=V^I\cap\operatorname{Im}(N^k)
\]
and
\[
V_{I,k}=\operatorname{Im}\!\left(\operatorname{Ker}(N^{k+1})
\longrightarrow V_I\right).
\]
Then $N^k$ induces an isomorphism
\[
N^k:V_{I,k}/V_{I,k-1}\xrightarrow{\sim}
V^{I,k}/V^{I,k+1},
\]
and defines canonical isomorphisms
\[
V_{I,k}/V_{I,k-1}\otimes I^{\otimes k}
\simeq V^{I,k}/V^{I,k+1}.
\]

\medskip
\noindent\emph{Proof.}
Via $N$, representations $V$ of $I$ identify with $k[T]$-modules killed by
a power of $T$.  It is enough to verify Lemma~5.3.2 when $V$ corresponds to
an indecomposable $k[T]$-module $k[T]/T^k$, where the assertion is immediate.

}{E7II-XXI-015}We now finish the proof of Lemma~5.3.1.  Lemma~5.3.2 supplies increasing and
decreasing filtrations $V_{I,k}$ and $V^{I,k}$ of $V_I$ and $V^I$, together
with isomorphisms of $G/I$-modules
\[
V^{I,k}/V^{I,k+1}\simeq
\bigl(V_{I,k}/V_{I,k-1}\bigr)(k).
\]
The assertion follows.



\SGApage{II}{391}
\SGAvar{\subsubsection*{Theorem 5.4.}
Let $X$ be a separated scheme of finite type over $\mathbb F_q$, of
dimension $n$, and let $\underline G$ be a constructible $E$-sheaf on $X$.
Suppose that, for every $x\in X$, the endomorphism $F_x$ of
$G_x^\vee$ is $T$-integral.  Then the endomorphisms $q^iF^\vee$ and
$q^nF^\vee$ of $H_c^i(X,\underline G)^\vee$ are $T$-integral.

}{\subsubsection*{Theorem 5.4.}
Let $X$ be a separated scheme of finite type over $\mathbb F_q$, of
dimension $n$, and let $\underline G$ be a constructible $E$-sheaf on $X$.
Suppose that, for every $x\in X$, the endomorphism $F_x$ of
$G_x^\vee$ is $T$-integral.  Then the endomorphisms $q^iF^\vee$ and
$q^nF^\vee$ of $\underline H_c^i(X,\underline G)^\vee$ are $T$-integral.

}{E7II-XXI-016}\medskip
\noindent\emph{Proof.}
As in Theorem~5.2.2, we reduce to the case in which $X$ is a smooth curve
and $\underline G$ is twisted constant on $X$.  Let $j$ be the inclusion of
$X$ in a smooth projective curve $X'$ compactifying it.

By Poincar\'e duality,
$\underline H^i(X',j_*\underline G)$ is dual to
$\underline H^{2-i}(X',j_*\underline G^\vee(1))$.  By Corollary~5.3(a),
Theorem~5.2.2 applies to $j_*\underline G^\vee$:
\begin{align*}
\underline H^0(X',j_*\underline G^\vee)
  &\simeq \underline H^2(X',j_*\underline G)^\vee(-1)
  &&\text{is \(T\)-integral},\\
\underline H^1(X',j_*\underline G^\vee)
  &\simeq \underline H^1(X',j_*\underline G)^\vee(-1)
  &&\text{is \(T\)-integral},\\
\underline H^2(X',j_*\underline G^\vee)(1)
  &\simeq \underline H^0(X',j_*\underline G)^\vee
  &&\text{is \(T\)-integral}.
\end{align*}
This proves the theorem for $j_*\underline G$.

Define $\underline H$ by the exact sequence
\[
0\longrightarrow j_!\underline G\longrightarrow j_*\underline G
\longrightarrow\underline H\longrightarrow0.
\]
If $Y$ is the support of $\underline H$, the long exact cohomology sequence
provides
\begin{center}
\begin{tikzcd}[row sep=large,column sep=large]
{} \& \underline H_c^0(X,\underline G)
  \arrow[r,hook] \& \underline H^0(X',j_*\underline G) \& ,\\
\underline H^0(Y,\underline H) \arrow[r] \&
  \underline H_c^1(X,\underline G) \arrow[r] \&
  \underline H^1(X',j_*\underline G) \& ,\\
{} \& \underline H_c^2(X,\underline G)
  \arrow[r,"\sim"] \& \underline H^2(X',j_*\underline G) \& ,
\end{tikzcd}
\end{center}
Thus the theorem for $(X,\underline G)$ follows from the theorem for
$(X',j_*\underline G)$.


\subsection*{5.5.}
The preceding method of fibration into curves also gives Lang--Weil
estimates~[3] for the complex absolute values of the eigenvalues of
Frobenius.

\SGApage{II}{\(\approx\)392}
Let $X$ be a scheme of finite type over $\mathbb F_q$, let $\underline G$
be a constructible $E$-sheaf on $X$, and let $\lambda\in\mathbb R$.  Suppose
that, for every point $x\in X$, with $\#k(x)=q_x$, the eigenvalues of $F_x$
acting on $G_x$ are algebraic numbers $\alpha$ all of whose complex absolute
values satisfy
\[
|\alpha|\leq q_x^\lambda.
\]
We shall then write
\[
\log|\underline G|\leq\lambda.
\]

\subsubsection*{Lemma 5.5.1.}
Let $X$ be a curve over $\mathbb F_q$, by which we mean here a separated
scheme $X$ of finite type over $\mathbb F_q$ and of dimension at most $1$.
Let $\underline G$ be a constructible $E$-sheaf on $X$ such that
$\log|\underline G|\leq\lambda$, with $\lambda\in\mathbb R$.  Then the
complex absolute values of the eigenvalues $\alpha$ of $F$ acting on
$\underline H_c^1(X,G)$ satisfy
\[
|\alpha|\leq q^{\lambda+1}.
\]

\medskip
\noindent\emph{Proof.}
Let $\sigma$ be a complex embedding of $E$, or simply a complex embedding of
the algebraic closure of $\mathbb Q$ in $E$.  It is enough to verify that,
for every $\sigma$, the zeros $\beta$ of the rational function
$\sigma Z(X,\underline G,t)$ satisfy
\[
|\beta^{-1}|\leq q^{\lambda+1}.
\]
Indeed, Lemma~5.2.1 already shows that, if $\alpha$ is an eigenvalue of
Frobenius acting on $\underline H_c^0(X,\underline G)$ or on
$\underline H_c^2(X,\underline G)$, then
$|\sigma\alpha|\leq q^{\lambda+1}$.

\SGApage{II}{\(\approx\)393}
\SGAvar{We verify that the series
\[
\log\sigma Z(X,\underline G,t)
=\sum_k\frac{t^k}{k}
  \sum_{x\in X(\mathbb F_{q^k})}
  \sigma\operatorname{Tr}(F,x^*\underline G)
\]
converges for $|t|<q^{-\lambda-1}$.  Representing $X$ as a ramified covering
of $\mathbb P^1$, one obtains the elementary estimate
\[
\#X(\mathbb F_{q^k})\leq A q^k.
\]
On the other hand, the hypothesis gives
\[
\left|\sigma\operatorname{Tr}(F,x^*\underline G)\right|
\leq B(q^k)^\lambda.
\]
Thus the coefficient $c_k$ of $t^k$ in
$\sigma Z(X,\underline G,t)$ satisfies
\[
|c_k|\leq\frac{1}{k}\,C(q^{1+\lambda})^k,
\]
and the assertion follows.

\subsubsection*{Theorem 5.5.2.}
Let $X$ be a separated scheme of finite type over $\mathbb F_q$, of
dimension at most $n$, and let $G$ be a constructible $E$-sheaf
on $X$ such that $\log|G|\leq\lambda$, with
$\lambda\in\mathbb R$.  Then the complex absolute values of the eigenvalues
$\alpha$ of Frobenius acting on $H_c^i(X,\underline G)$ satisfy
\[
|\alpha|\leq q^{\lambda+i}
\qquad\text{and}\qquad
|\alpha|\leq q^{\lambda+n}.
\]

\medskip
\noindent\emph{Proof.}
If $n\leq0$, the theorem is trivial.  If $n=1$, its cases $i=0$ and $i=2$
follow from Lemma~5.2.1, while the case $i=1$ follows from
Lemma~5.5.1.  We then proceed by induction on $n$, as in the proof of
Theorem~5.5.2.

}{We verify that the series
\[
\log\sigma Z(X,\underline G,t)
=\sum_k\frac{t^k}{k}
  \sum_{x\in X(\mathbb F_{q^k})}
  \sigma\operatorname{Tr}(F,x^*\underline G)
\]
converges for $|t|<q^{-\lambda-1}$.  Representing $X$ as a ramified covering
of $\mathbb P^1$, one obtains the elementary estimate
\[
\#X(\mathbb F_{q^k})\leq A q^k.
\]
On the other hand, the hypothesis gives
\[
\left|\sigma\operatorname{Tr}(F,x^*\underline G)\right|
\leq B(q^k)^\lambda.
\]
Thus the coefficient $c_k$ of $t^k$ in
$\log\sigma Z(X,\underline G,t)$ satisfies
\[
|c_k|\leq\frac{1}{k}\,C(q^{1+\lambda})^k,
\]
and the assertion follows.

\subsubsection*{Theorem 5.5.2.}
Let $X$ be a separated scheme of finite type over $\mathbb F_q$, of
dimension at most $n$, and let $G$ be a constructible $E$-sheaf
on $X$ such that $\log|G|\leq\lambda$, with
$\lambda\in\mathbb R$.  Then the complex absolute values of the eigenvalues
$\alpha$ of Frobenius acting on $\underline H_c^i(X,\underline G)$ satisfy
\[
|\alpha|\leq q^{\lambda+i}
\qquad\text{and}\qquad
|\alpha|\leq q^{\lambda+n}.
\]

\medskip
\noindent\emph{Proof.}
If $n\leq0$, the theorem is trivial.  If $n=1$, its cases $i=0$ and $i=2$
follow from Lemma~5.2.1, while the case $i=1$ follows from
Lemma~5.5.1.  We then proceed by induction on $n$, as in the proof of
Theorem~5.2.2.

}{E7II-XXI-017}\SGApage{II}{394}
\subsubsection*{Corollary 5.5.3.}
Let $X$ be a separated scheme of finite type and dimension $n$ over
$\mathbb F_q$.  Then the eigenvalues $\alpha$ of the geometric Frobenius
endomorphism $F$ of $\underline H_c^i(X,\mathbb Q_\ell)$ are algebraic
integers.  Moreover:
\begin{enumerate}[label=(\roman*)]
\item for $\ell'\ne p$, $\alpha$ is an $\ell'$-adic unit;
\item for $i>n$, $\alpha$ is divisible by $q^{i-n}$;
\item for $i\leq n$ (respectively, $i\geq n$),
  $q^i\alpha^{-1}$ (respectively, $q^n\alpha^{-1}$) is an algebraic integer;
\item for $0<i$, the complex absolute values of $\alpha$ satisfy
  \[
  |\alpha|\leq q^{i-\frac12}
  \qquad\text{and}\qquad
  |\alpha|\leq q^n.
  \]
\end{enumerate}

\medskip
\noindent\emph{Proof.}
By Theorem~5.2.2, $\alpha$ and $q^{n-i}\alpha$ are algebraic integers.
Assertions (i)--(iii) then follow from Theorem~5.4.  For $i>n$, assertion
(iv) follows from Theorem~5.5.2.  To obtain (iv) when $i\leq n$, we must
combine Theorem~5.5.2 with the Riemann hypothesis for curves.

\SGAvar{We argue by induction on $n$, the case $n=0$ being trivial.  If $U$ is an
open subset of $X$ whose complement $Z$ has dimension less than $n$, the
exact sequence
\[
\underline H_c^i(U,\mathbb Q_\ell)\longrightarrow
\underline H_c^i(X,\mathbb Q_\ell)\longrightarrow
\underline H_c^i(Z,\mathbb Q_\ell)
\]
shows that it is enough to prove 5.2.2 for $U$.  After
first replacing $X$ by $X_{\mathrm{red}}$, this reduces us to the case in
which $X$ is smooth and affine and there is a smooth morphism
$f:X\longrightarrow S$ satisfying:
\begin{enumerate}[label=(\alph*)]
\item $\dim S\leq n-1$;
\item the fibers of $f$ are purely one-dimensional.
\end{enumerate}
Since the fibers of $f$ are affine curves, one has
\[
R^0f_!\mathbb Q_\ell=0.
\]
By Weil~[4],
\begin{equation}\tag{5.5.3.1}
\log\bigl|R^1f_!\mathbb Q_\ell\bigr|\leq\frac12,
\end{equation}
and Lemma~5.2.1 applied to the fibers of $f$ gives
\begin{equation}\tag{5.5.3.2}
\log\bigl|R^2f_!\mathbb Q_\ell\bigr|\leq1.
\end{equation}
The Leray spectral sequence for $f$ reduces here to a long exact sequence;
the relevant part is
\[
\underline H_c^{i-2}(S,R^2f_!\mathbb Q_\ell)\longrightarrow
\underline H_c^i(X,\mathbb Q_\ell)\longrightarrow
\underline H_c^{i-1}(S,R^1f_!\mathbb Q_\ell).
\]
By (5.5.3.1) (respectively, (5.5.3.2)) and Theorem~5.5.2, the complex
absolute values of the eigenvalues $\alpha$ of Frobenius on
$\underline H_c^{i-1}(S,R^1f_!\mathbb Q_\ell)$ (respectively, on
$\underline H_c^{i-2}(S,R^2f_!\mathbb Q_\ell)$) satisfy
\[
|\alpha|\leq q^{i-1+\frac12}=q^{i-\frac12}
\]
(respectively,
\[
|\alpha|\leq q^{i-2+1}=q^{i-1}\leq q^{i-\frac12}.
\]
The corollary follows.



}{We argue by induction on $n$, the case $n=0$ being trivial.  If $U$ is an
open subset of $X$ whose complement $Z$ has dimension less than $n$, the
exact sequence
\[
\underline H_c^i(U,\mathbb Q_\ell)\longrightarrow
\underline H_c^i(X,\mathbb Q_\ell)\longrightarrow
\underline H_c^i(Z,\mathbb Q_\ell)
\]
shows that it is enough to prove the desired assertion for $U$.  After
first replacing $X$ by $X_{\mathrm{red}}$, this reduces us to the case in
which $X$ is smooth and affine and there is a smooth morphism
$f:X\longrightarrow S$ satisfying:
\begin{enumerate}[label=(\alph*)]
\item $\dim S\leq n-1$;
\item the fibers of $f$ are purely one-dimensional.
\end{enumerate}
Since the fibers of $f$ are affine curves, one has
\[
R^0f_!\mathbb Q_\ell=0.
\]
By Weil~[4],
\begin{equation}\tag{5.5.3.1}
\log\bigl|R^1f_!\mathbb Q_\ell\bigr|\leq\frac12,
\end{equation}
and Lemma~5.2.1 applied to the fibers of $f$ gives
\begin{equation}\tag{5.5.3.2}
\log\bigl|R^2f_!\mathbb Q_\ell\bigr|\leq1.
\end{equation}
The Leray spectral sequence for $f$ reduces here to a long exact sequence;
the relevant part is
\[
\underline H_c^{i-2}(S,R^2f_!\mathbb Q_\ell)\longrightarrow
\underline H_c^i(X,\mathbb Q_\ell)\longrightarrow
\underline H_c^{i-1}(S,R^1f_!\mathbb Q_\ell).
\]
By (5.5.3.1) (respectively, (5.5.3.2)) and Theorem~5.5.2, the complex
absolute values of the eigenvalues $\alpha$ of Frobenius on
$\underline H_c^{i-1}(S,R^1f_!\mathbb Q_\ell)$ (respectively, on
$\underline H_c^{i-2}(S,R^2f_!\mathbb Q_\ell)$) satisfy
\[
|\alpha|\leq q^{i-1+\frac12}=q^{i-\frac12}
\]
(respectively,
\[
|\alpha|\leq q^{i-2+1}=q^{i-1}\leq q^{i-\frac12}).
\]
The corollary follows.



}{E7II-XXI-018}\SGApage{II}{\(\approx\)396}
\subsection*{5.6.}
To obtain results analogous to the preceding ones for cohomology without
supports, it seems necessary to have resolution of singularities available.
In dimension at most $1$, the problem is solved by Corollary~5.3.

\subsubsection*{Theorem 5.6.}
Let $N$ be an integer, and suppose that strong resolution of singularities
in the sense of Hironaka is available for schemes of finite type over
$\mathbb F_q$ of dimension at most $N$ (at present one may therefore take
$N=2$~[1]).  Let $f:X\longrightarrow Y$ be a morphism of schemes of finite
type over $\mathbb F_q$, and let $\underline G$ be a constructible
$T$-integral $E$-sheaf on $X$.  If $\dim X\leq N$, then the $E$-sheaves
$R^if_*\underline G$ on $Y$ are $T$-integral.

\SGAvar{\medskip
\noindent\emph{Proof.}
We argue by induction on the dimension $n\leq N$ of $X$.  Let
$j:U\hookrightarrow X$ be an open subset of $X$, smooth and affine, such
that $\underline G|_U$ is twisted constant and the dimension of its
complement $Z$ is less than $n$.  By resolution of singularities, the
sheaves $R^ij_*(j^*\underline G)$ are constructible
(SGA~5, Expos\'e~I; [2]); for $i>0$ they are supported on $Z$, and for
$i=0$, $j_*(j^*\underline G)$ differs from $\underline G$ only on $Z$.
The induction hypothesis therefore shows that it is enough to prove the
theorem for the sheaf $j^*\underline G$ and the morphisms $j$ and
$f j$.  Thus we may further suppose that $X$ is smooth and affine and
that $\underline G$ is twisted constant, that is, comes from a continuous
representation of the fundamental group of $X$ on a finite-dimensional
$E$-vector space.

}{\medskip
\noindent\emph{Proof.}
We argue by induction on the dimension $n\leq N$ of $X$.  Let
$j:U\hookrightarrow X$ be an open subset of $X$, smooth and affine, such
that $\underline G|_U$ is twisted constant and the dimension of its
complement $Z$ is less than $n$.  By resolution of singularities, the
sheaves $R^ij_*(j^*\underline G)$ are constructible
(SGA~5, Expos\'e~I; [2]); for $i>0$ they are supported on $Z$, and for
$i=0$, $j_*(j^*\underline G)$ differs from $\underline G$ only on $Z$.
The induction hypothesis therefore shows that it is enough to prove the
theorem for the sheaf $j^*\underline G$ and the morphisms $j$ and
$f\circ j$.  Thus we may further suppose that $X$ is smooth and affine and
that $\underline G$ is twisted constant, that is, comes from a continuous
representation of the fundamental group of $X$ on a finite-dimensional
$E$-vector space.

}{E7II-XXI-019}\SGApage{II}{\(\approx\)397}
Let $R$ be the ring of integers of $E$, let $\lambda$ be a uniformizer, let
$e$ be the absolute ramification index, and choose an integer
$b>e/(\ell-1)$.  Since $X$ is smooth, hence normal, and its fundamental
group is profinite, there is a twisted-constant $R$-sheaf
$\underline G^\circ$ such that
\[
\underline G\simeq\underline G^\circ\otimes_R E.
\]
Let $p:X'\longrightarrow X$ be the \'etale covering that trivializes the
locally constant sheaf
$\underline G^\circ\otimes_R R/\lambda^b$.  The sheaf $\underline G$ is a
direct summand of $p_*p^*\underline G$, so it is enough to prove the theorem
for $p^*\underline G$ and $f\circ p$.  It remains to treat the case in which
the following conditions hold:
\[
(*)\quad
\begin{minipage}{0.82\linewidth}
$X$ is smooth and affine, and there is a twisted-constant $R$-sheaf
$\underline G^\circ$ such that
$\underline G\simeq\underline G^\circ\otimes_R E$ and
$\underline G^\circ\otimes_R R/\lambda^b$ is constant.
\end{minipage}
\]

Factor $f$ as a proper morphism followed by an open immersion,
$f=\bar f\circ j$:
\begin{center}
\begin{tikzcd}
X \arrow[r,hook,"j",swap] \arrow[d,"f"'] \&
X' \arrow[dl,"\bar f"]\\
Y \& {}
\end{tikzcd}
\qquad .
\end{center}
The theorem has already been proved for $\bar f$ (5.2), so it is enough to
prove it for $j$.  Applying resolution of singularities to $X'$, we may
suppose that $X'$ is smooth and that $X$ is the complement in $X'$ of a
normal-crossings divisor which is a union of smooth divisors.  Hypothesis
$(*)$ implies that $\underline G$ is tamely ramified along this divisor.
We then have:

\SGApage{II}{398}
\subsubsection*{Lemma 5.6.1.}
Let $X$ be a smooth scheme of finite type over $\mathbb F_q$, let $D$ be a
normal-crossings divisor in $X$ which is a union of smooth divisors $D_i$,
let $j:U\hookrightarrow X$ be the complement of $D$ in $X$, and let
$\underline G$ be a twisted-constant $E$-sheaf on $U$, tamely ramified
along $D$.  If $\underline G$ is $T$-integral, then the sheaves
$R^ij_*\underline G$ are $T$-integral.

We shall admit the following result (cf. SGA~1, Expos\'e~XIII).

\subsubsection*{Lemma 5.6.2.}
\begin{enumerate}[label=(\roman*)]
\item Under the hypotheses of Lemma~5.6.1, let $D'$ be a component of $D$,
  let $D''$ be the union of the other components, and put
  $U'=X-D''$ and $D^{\prime\circ}=D'\cap U'$:
  \begin{center}
  \begin{tikzcd}
  {} \& D^{\prime\circ}
    \arrow[r,hook,"j'_D",swap] \arrow[d,"i^\circ"] \&
    D' \arrow[d,"i"]\\
  U \arrow[r,hook,"j_1",swap] \&
    U' \arrow[r,hook,"j'",swap] \& X
  \end{tikzcd}
  \qquad .
  \end{center}
  Then the base-change map
  \[
  i^*Rj'_*(Rj_{1*}\underline G)
  \longrightarrow
  Rj'_{D*}\,i^{\circ *}(Rj_{1*}\underline G)
  \]
  is an isomorphism, and the sheaves
  $i^{\circ *}Rj_{1*}\underline G$ are twisted constant and tamely ramified
  along $D'-D^{\prime\circ}$.

\item If $C$ is a smooth curve on $X$, transverse to $D$ at a smooth point
  $x$ of $D$,
  \begin{center}
  \begin{tikzcd}
  C\cap U \arrow[r,"j_c",swap] \arrow[d,"k'"] \&
    C \arrow[d,"k"]\\
  U \arrow[r,"j",swap] \& X
  \end{tikzcd}
  \qquad ,
  \end{center}
  then the base-change map
  \[
  k^*Rj_*\underline G\longrightarrow Rj_{c*}\,k^{\prime *}\underline G
  \]
  is an isomorphism at $x$.
\end{enumerate}

\SGApage{II}{\(\approx\)399}
\SGAvar{Let us deduce Lemma~5.6.1 from Lemma~5.6.2 by induction on the dimension,
temporarily admitting Lemma~5.6.1 for curves.  Lemma~5.6.2(ii) and this
hypothesis show that the sheaves $R^ij_*\underline G$ are $T$-integral at
every smooth point of $D$.  In particular, in the notation of
Lemma~5.6.2(i), the sheaves $R^ij_{1*}\underline G$ are $T$-integral.
Lemma~6.2(i), together with the induction hypothesis applied to
$i^{\circ *}Rj_{1*}\underline G$ and to $j'_D$, therefore shows that the
sheaves $i^*R^ij_*\underline G$ are $T$-integral, whichever component $D'$
of $D$ is chosen.  This proves Lemma~5.6.1.

}{Let us deduce Lemma~5.6.1 from Lemma~5.6.2 by induction on the dimension,
temporarily admitting Lemma~5.6.1 for curves.  Lemma~5.6.2(ii) and this
hypothesis show that the sheaves $R^ij_*\underline G$ are $T$-integral at
every smooth point of $D$.  In particular, in the notation of
Lemma~5.6.2(i), the sheaves $R^ij_{1*}\underline G$ are $T$-integral.
Lemma~5.6.2(i), together with the induction hypothesis applied to
$i^{\circ *}Rj_{1*}\underline G$ and to $j'_D$, therefore shows that the
sheaves $i^*R^ij_*\underline G$ are $T$-integral, whichever component $D'$
of $D$ is chosen.  This proves Lemma~5.6.1.

}{E7II-XXI-020}It remains to handle curves.  Under the hypotheses of Lemma~5.6.1, let $x$
be a point of $X-U$, let $\mathfrak G$ be the Galois group of the fraction
field of $\widehat{\mathcal O}_{x,X}$, and let $\mathfrak I$ be the inertia
group.  Denote again by $\underline G$ the representation of $\mathfrak G$
defined by the sheaf, and let $\bar x$ be a geometric point of $X$ localized
at $x$.  There are isomorphisms of Galois
$\mathfrak G/\mathfrak I$-modules
\[
(j_*\underline G)_{\bar x}\simeq\underline G^{\mathfrak I},
\]
\[
(R^1j_*\underline G)_{\bar x}\simeq\underline G_{\mathfrak I}(-1),
\]
and
\[
(R^ij_*\underline G)_{\bar x}=0\qquad\text{for }i\geq2.
\]
The assertion follows from Corollary~5.3(a) and Lemma~5.3.1.

\section*{Bibliography}
\begin{enumerate}[label={\textup{[\arabic*]}},leftmargin=2.4em]
\item S.~S. Abhyankar, \emph{Resolution of Singularities of Embedded Algebraic Surfaces}, Academic Press, New York, 1966.
\item J.-P.~Jouanolou, \emph{Cohomologie $\ell$-adique des schémas}, forthcoming in \emph{Lecture Notes in Mathematics}.
\item S.~Lang and A.~Weil, \emph{Number of Points of Varieties in Finite Fields}, \emph{Amer. J. Math.} \textbf{76} (1954), 819--827.
\item A.~Weil, \emph{Sur les courbes algébriques et les variétés qui s'en déduisent}, \emph{Actualités Sci. Indust.} no.~1041, Hermann, Paris, 1948.
\end{enumerate}



\markboth{Exposé XXI}{Exposé XXI}

\ifnum\SGAform=1
\clearpage\section*{Critical apparatus}\begingroup\small\raggedright\sloppy
Readings are given in the exact TeX notation of the two English witnesses. Decimal suffixes identify the individual differences in each marked passage.

\par\medskip\noindent\hypertarget{app-E7II-XXI-001.1}{}\textbf{E7II-XXI-001.1}\quad printed folio 365.\par
\noindent Printed reading: \emph{(empty)}\par
\noindent Corrected reading: {\ttfamily\footnotesize here\allowbreak{} }\par
\noindent Source: {\ttfamily\footnotesize fr/sga7ii/\{diplomati\allowbreak{}c,current\}/expose\_XX\allowbreak{}I\_body.tex;\allowbreak{} folio\allowbreak{} 365\allowbreak{}}\par

\par\medskip\noindent\hypertarget{app-E7II-XXI-001.2}{}\textbf{E7II-XXI-001.2}\quad printed folio 365.\par
\noindent Printed reading: \emph{(empty)}\par
\noindent Corrected reading: {\ttfamily\footnotesize the\allowbreak{} }\par
\noindent Source: {\ttfamily\footnotesize fr/sga7ii/\{diplomati\allowbreak{}c,current\}/expose\_XX\allowbreak{}I\_body.tex;\allowbreak{} folio\allowbreak{} 365\allowbreak{}}\par

\par\medskip\noindent\hypertarget{app-E7II-XXI-001.3}{}\textbf{E7II-XXI-001.3}\quad printed folio 365.\par
\noindent Printed reading: \emph{(empty)}\par
\noindent Corrected reading: {\ttfamily\footnotesize  (\textbackslash{}S\allowbreak{}}\par
\noindent Source: {\ttfamily\footnotesize fr/sga7ii/\{diplomati\allowbreak{}c,current\}/expose\_XX\allowbreak{}I\_body.tex;\allowbreak{} folio\allowbreak{} 365\allowbreak{}}\par

\par\medskip\noindent\hypertarget{app-E7II-XXI-001.4}{}\textbf{E7II-XXI-001.4}\quad printed folio 365.\par
\noindent Printed reading: {\ttfamily\footnotesize I\allowbreak{}}\par
\noindent Corrected reading: {\ttfamily\footnotesize 5)\allowbreak{}}\par
\noindent Source: {\ttfamily\footnotesize fr/sga7ii/\{diplomati\allowbreak{}c,current\}/expose\_XX\allowbreak{}I\_body.tex;\allowbreak{} folio\allowbreak{} 365\allowbreak{}}\par

\par\medskip\noindent\hypertarget{app-E7II-XXI-002.1}{}\textbf{E7II-XXI-002.1}\quad printed folio 365.\par
\noindent Printed reading: \emph{(empty)}\par
\noindent Corrected reading: {\ttfamily\footnotesize  \$X\$\allowbreak{}}\par
\noindent Source: {\ttfamily\footnotesize ED-EN-FORMULA-G219;\allowbreak{} fr/sga7ii/\{diplomati\allowbreak{}c,current\}/expose\_XX\allowbreak{}I\_body.tex;\allowbreak{} folio\allowbreak{} 365;\allowbreak{} port/PORT\_EN7II\_LOG.\allowbreak{}jsonl\#C765;\allowbreak{} fixes/FIXLOG\_7II\_CUR\allowbreak{}RENT.jsonl\#C765;\allowbreak{} 7II:654\allowbreak{}}\par

\par\medskip\noindent\hypertarget{app-E7II-XXI-003.1}{}\textbf{E7II-XXI-003.1}\quad printed folio 368.\par
\noindent Printed reading: \emph{(empty)}\par
\noindent Corrected reading: {\ttfamily\footnotesize i=\allowbreak{}}\par
\noindent Source: {\ttfamily\footnotesize fr/sga7ii/\{diplomati\allowbreak{}c,current\}/expose\_XX\allowbreak{}I\_body.tex;\allowbreak{} folio\allowbreak{} 368;\allowbreak{} formula\allowbreak{} 92\allowbreak{}}\par

\par\medskip\noindent\hypertarget{app-E7II-XXI-003.2}{}\textbf{E7II-XXI-003.2}\quad printed folio 369.\par
\noindent Printed reading: {\ttfamily\footnotesize )\allowbreak{}}\par
\noindent Corrected reading: \emph{(empty)}\par
\noindent Source: {\ttfamily\footnotesize fr/sga7ii/\{diplomati\allowbreak{}c,current\}/expose\_XX\allowbreak{}I\_body.tex;\allowbreak{} folio\allowbreak{} 369;\allowbreak{} formula\allowbreak{} 99\allowbreak{}}\par

\par\medskip\noindent\hypertarget{app-E7II-XXI-003.3}{}\textbf{E7II-XXI-003.3}\quad printed folio 369.\par
\noindent Printed reading: \emph{(empty)}\par
\noindent Corrected reading: {\ttfamily\footnotesize =1\allowbreak{}}\par
\noindent Source: {\ttfamily\footnotesize fr/sga7ii/\{diplomati\allowbreak{}c,current\}/expose\_XX\allowbreak{}I\_body.tex;\allowbreak{} folio\allowbreak{} 369;\allowbreak{} formula\allowbreak{} 106\allowbreak{}}\par

\par\medskip\noindent\hypertarget{app-E7II-XXI-003.4}{}\textbf{E7II-XXI-003.4}\quad printed folio 369.\par
\noindent Printed reading: \emph{(empty)}\par
\noindent Corrected reading: {\ttfamily\footnotesize  set\allowbreak{} \$a\_0=0\$\allowbreak{} and\allowbreak{}}\par
\noindent Source: {\ttfamily\footnotesize ED-EN-FORMULA-G221;\allowbreak{} fr/sga7ii/\{diplomati\allowbreak{}c,current\}/expose\_XX\allowbreak{}I\_body.tex;\allowbreak{} folio\allowbreak{} 369\allowbreak{}}\par

\par\medskip\noindent\hypertarget{app-E7II-XXI-003.5}{}\textbf{E7II-XXI-003.5}\quad printed folio 369.\par
\noindent Printed reading: {\ttfamily\footnotesize \textbackslash{}left\textbackslash{}\{\textbackslash{}begin\{aligne\allowbreak{}d\}\allowbreak{} }\par
\noindent Corrected reading: \emph{(empty)}\par
\noindent Source: {\ttfamily\footnotesize ED-EN-FORMULA-G225;\allowbreak{} fr/sga7ii/\{diplomati\allowbreak{}c,current\}/expose\_XX\allowbreak{}I\_body.tex;\allowbreak{} folio\allowbreak{} 369\allowbreak{}}\par

\par\medskip\noindent\hypertarget{app-E7II-XXI-003.6}{}\textbf{E7II-XXI-003.6}\quad printed folio 369.\par
\noindent Printed reading: {\ttfamily\footnotesize 1\&\allowbreak{}}\par
\noindent Corrected reading: {\ttfamily\footnotesize i\allowbreak{}}\par
\noindent Source: {\ttfamily\footnotesize ED-EN-FORMULA-G225;\allowbreak{} fr/sga7ii/\{diplomati\allowbreak{}c,current\}/expose\_XX\allowbreak{}I\_body.tex;\allowbreak{} folio\allowbreak{} 369\allowbreak{}}\par

\par\medskip\noindent\hypertarget{app-E7II-XXI-003.7}{}\textbf{E7II-XXI-003.7}\quad printed folio 369.\par
\noindent Printed reading: {\ttfamily\footnotesize o,\textbackslash{}\textbackslash{}\allowbreak{} d\_2\&\textbackslash{}geq\allowbreak{} a\_2-a\_1,\textbackslash{}\textbackslash{}\allowbreak{} d\_r\&\textbackslash{}geq\allowbreak{} a\_r\allowbreak{}}\par
\noindent Corrected reading: {\ttfamily\footnotesize i\allowbreak{}}\par
\noindent Source: {\ttfamily\footnotesize ED-EN-FORMULA-G225;\allowbreak{} fr/sga7ii/\{diplomati\allowbreak{}c,current\}/expose\_XX\allowbreak{}I\_body.tex;\allowbreak{} folio\allowbreak{} 369\allowbreak{}}\par

\par\medskip\noindent\hypertarget{app-E7II-XXI-003.8}{}\textbf{E7II-XXI-003.8}\quad printed folio 369.\par
\noindent Printed reading: {\ttfamily\footnotesize r\allowbreak{}}\par
\noindent Corrected reading: {\ttfamily\footnotesize i\allowbreak{}}\par
\noindent Source: {\ttfamily\footnotesize ED-EN-FORMULA-G225;\allowbreak{} fr/sga7ii/\{diplomati\allowbreak{}c,current\}/expose\_XX\allowbreak{}I\_body.tex;\allowbreak{} folio\allowbreak{} 369\allowbreak{}}\par

\par\medskip\noindent\hypertarget{app-E7II-XXI-003.9}{}\textbf{E7II-XXI-003.9}\quad printed folio 369.\par
\noindent Printed reading: {\ttfamily\footnotesize .\allowbreak{} \textbackslash{}end\{aligned\}\textbackslash{}right\allowbreak{}}\par
\noindent Corrected reading: {\ttfamily\footnotesize \textbackslash{}qquad(1\textbackslash{}leq\allowbreak{} i\textbackslash{}leq\allowbreak{} r)\allowbreak{}}\par
\noindent Source: {\ttfamily\footnotesize ED-EN-FORMULA-G225;\allowbreak{} fr/sga7ii/\{diplomati\allowbreak{}c,current\}/expose\_XX\allowbreak{}I\_body.tex;\allowbreak{} folio\allowbreak{} 369\allowbreak{}}\par

\par\medskip\noindent\hypertarget{app-E7II-XXI-003.10}{}\textbf{E7II-XXI-003.10}\quad printed folio 370.\par
\noindent Printed reading: \emph{(empty)}\par
\noindent Corrected reading: {\ttfamily\footnotesize +1\allowbreak{}}\par
\noindent Source: {\ttfamily\footnotesize fr/sga7ii/\{diplomati\allowbreak{}c,current\}/expose\_XX\allowbreak{}I\_body.tex;\allowbreak{} folio\allowbreak{} 370;\allowbreak{} formula\allowbreak{} 115\allowbreak{}}\par

\par\medskip\noindent\hypertarget{app-E7II-XXI-004.1}{}\textbf{E7II-XXI-004.1}\quad printed folio 370.\par
\noindent Printed reading: \emph{(empty)}\par
\noindent Corrected reading: {\ttfamily\footnotesize )\textbackslash{}bigr\allowbreak{}}\par
\noindent Source: {\ttfamily\footnotesize fr/sga7ii/\{diplomati\allowbreak{}c,current\}/expose\_XX\allowbreak{}I\_body.tex;\allowbreak{} folio\allowbreak{} 370\allowbreak{}}\par

\par\medskip\noindent\hypertarget{app-E7II-XXI-004.2}{}\textbf{E7II-XXI-004.2}\quad printed folio 371.\par
\noindent Printed reading: {\ttfamily\footnotesize 0\allowbreak{}}\par
\noindent Corrected reading: {\ttfamily\footnotesize 1\allowbreak{}}\par
\noindent Source: {\ttfamily\footnotesize fr/sga7ii/\{diplomati\allowbreak{}c,current\}/expose\_XX\allowbreak{}I\_body.tex;\allowbreak{} folio\allowbreak{} 371;\allowbreak{} formula\allowbreak{} 140;\allowbreak{} port/PORT\_EN7II\_LOG.\allowbreak{}jsonl\#C893;\allowbreak{} fixes/FIXLOG\_7II\_CUR\allowbreak{}RENT.jsonl\#C893;\allowbreak{} 7II:1105\allowbreak{}}\par

\par\medskip\noindent\hypertarget{app-E7II-XXI-004.3}{}\textbf{E7II-XXI-004.3}\quad printed folio 371.\par
\noindent Printed reading: {\ttfamily\footnotesize n\allowbreak{}}\par
\noindent Corrected reading: {\ttfamily\footnotesize M\allowbreak{}}\par
\noindent Source: {\ttfamily\footnotesize fr/sga7ii/\{diplomati\allowbreak{}c,current\}/expose\_XX\allowbreak{}I\_body.tex;\allowbreak{} folio\allowbreak{} 371;\allowbreak{} formula\allowbreak{} 141\allowbreak{}}\par

\par\medskip\noindent\hypertarget{app-E7II-XXI-004.4}{}\textbf{E7II-XXI-004.4}\quad printed folio 371.\par
\noindent Printed reading: {\ttfamily\footnotesize u\allowbreak{}}\par
\noindent Corrected reading: {\ttfamily\footnotesize U\allowbreak{}}\par
\noindent Source: {\ttfamily\footnotesize fr/sga7ii/\{diplomati\allowbreak{}c,current\}/expose\_XX\allowbreak{}I\_body.tex;\allowbreak{} folio\allowbreak{} 371;\allowbreak{} formula\allowbreak{} 143\allowbreak{}}\par

\par\medskip\noindent\hypertarget{app-E7II-XXI-004.5}{}\textbf{E7II-XXI-004.5}\quad printed folio 372.\par
\noindent Printed reading: \emph{(empty)}\par
\noindent Corrected reading: {\ttfamily\footnotesize 1.\allowbreak{}}\par
\noindent Source: {\ttfamily\footnotesize ED-EN-NONMATH-G038;\allowbreak{} F7II-XXI-014.1;\allowbreak{} WITNESS-F7II-XXI-014\allowbreak{}.1;\allowbreak{} fr/sga7ii/\{diplomati\allowbreak{}c,current\}/expose\_XX\allowbreak{}I\_body.tex;\allowbreak{} folio\allowbreak{} 372\allowbreak{}}\par

\par\medskip\noindent\hypertarget{app-E7II-XXI-005.1}{}\textbf{E7II-XXI-005.1}\quad printed folio 372.\par
\noindent Printed reading: \emph{(empty)}\par
\noindent Corrected reading: {\ttfamily\footnotesize \textbackslash{}bigr)\allowbreak{}}\par
\noindent Source: {\ttfamily\footnotesize fr/sga7ii/\{diplomati\allowbreak{}c,current\}/expose\_XX\allowbreak{}I\_body.tex;\allowbreak{} folio\allowbreak{} 372;\allowbreak{} formula\allowbreak{} 177\allowbreak{}}\par

\par\medskip\noindent\hypertarget{app-E7II-XXI-005.2}{}\textbf{E7II-XXI-005.2}\quad printed folio 373.\par
\noindent Printed reading: \emph{(empty)}\par
\noindent Corrected reading: {\ttfamily\footnotesize T\allowbreak{} }\par
\noindent Source: {\ttfamily\footnotesize fr/sga7ii/\{diplomati\allowbreak{}c,current\}/expose\_XX\allowbreak{}I\_body.tex;\allowbreak{} folio\allowbreak{} 373;\allowbreak{} formula\allowbreak{} 182\allowbreak{}}\par

\par\medskip\noindent\hypertarget{app-E7II-XXI-005.3}{}\textbf{E7II-XXI-005.3}\quad printed folio 373.\par
\noindent Printed reading: {\ttfamily\footnotesize  T\allowbreak{}}\par
\noindent Corrected reading: \emph{(empty)}\par
\noindent Source: {\ttfamily\footnotesize fr/sga7ii/\{diplomati\allowbreak{}c,current\}/expose\_XX\allowbreak{}I\_body.tex;\allowbreak{} folio\allowbreak{} 373;\allowbreak{} formula\allowbreak{} 182\allowbreak{}}\par

\par\medskip\noindent\hypertarget{app-E7II-XXI-005.4}{}\textbf{E7II-XXI-005.4}\quad printed folio 374.\par
\noindent Printed reading: {\ttfamily\footnotesize o\allowbreak{}}\par
\noindent Corrected reading: {\ttfamily\footnotesize 0\allowbreak{}}\par
\noindent Source: {\ttfamily\footnotesize fr/sga7ii/\{diplomati\allowbreak{}c,current\}/expose\_XX\allowbreak{}I\_body.tex;\allowbreak{} folio\allowbreak{} 374;\allowbreak{} formula\allowbreak{} 219\allowbreak{}}\par

\par\medskip\noindent\hypertarget{app-E7II-XXI-005.5}{}\textbf{E7II-XXI-005.5}\quad printed folio 374.\par
\noindent Printed reading: {\ttfamily\footnotesize o\allowbreak{}}\par
\noindent Corrected reading: {\ttfamily\footnotesize 0\allowbreak{}}\par
\noindent Source: {\ttfamily\footnotesize fr/sga7ii/\{diplomati\allowbreak{}c,current\}/expose\_XX\allowbreak{}I\_body.tex;\allowbreak{} folio\allowbreak{} 374;\allowbreak{} formula\allowbreak{} 220\allowbreak{}}\par

\par\medskip\noindent\hypertarget{app-E7II-XXI-005.6}{}\textbf{E7II-XXI-005.6}\quad printed folio 374.\par
\noindent Printed reading: {\ttfamily\footnotesize o\allowbreak{}}\par
\noindent Corrected reading: {\ttfamily\footnotesize 0\allowbreak{}}\par
\noindent Source: {\ttfamily\footnotesize ED-EN-FORMULA-G227;\allowbreak{} fr/sga7ii/\{diplomati\allowbreak{}c,current\}/expose\_XX\allowbreak{}I\_body.tex;\allowbreak{} folio\allowbreak{} 374\allowbreak{}}\par

\par\medskip\noindent\hypertarget{app-E7II-XXI-005.7}{}\textbf{E7II-XXI-005.7}\quad printed folio 374.\par
\noindent Printed reading: {\ttfamily\footnotesize o\allowbreak{}}\par
\noindent Corrected reading: {\ttfamily\footnotesize 0\allowbreak{}}\par
\noindent Source: {\ttfamily\footnotesize fr/sga7ii/\{diplomati\allowbreak{}c,current\}/expose\_XX\allowbreak{}I\_body.tex;\allowbreak{} folio\allowbreak{} 374\allowbreak{}}\par

\par\medskip\noindent\hypertarget{app-E7II-XXI-005.8}{}\textbf{E7II-XXI-005.8}\quad printed folio 375.\par
\noindent Printed reading: {\ttfamily\footnotesize o\allowbreak{}}\par
\noindent Corrected reading: {\ttfamily\footnotesize 0\allowbreak{}}\par
\noindent Source: {\ttfamily\footnotesize ED-EN-FORMULA-G229;\allowbreak{} fr/sga7ii/\{diplomati\allowbreak{}c,current\}/expose\_XX\allowbreak{}I\_body.tex;\allowbreak{} folio\allowbreak{} 375\allowbreak{}}\par

\par\medskip\noindent\hypertarget{app-E7II-XXI-005.9}{}\textbf{E7II-XXI-005.9}\quad printed folio 375.\par
\noindent Printed reading: {\ttfamily\footnotesize o\allowbreak{}}\par
\noindent Corrected reading: {\ttfamily\footnotesize 0\allowbreak{}}\par
\noindent Source: {\ttfamily\footnotesize fr/sga7ii/\{diplomati\allowbreak{}c,current\}/expose\_XX\allowbreak{}I\_body.tex;\allowbreak{} folio\allowbreak{} 375;\allowbreak{} formula\allowbreak{} 242\allowbreak{}}\par

\par\medskip\noindent\hypertarget{app-E7II-XXI-005.10}{}\textbf{E7II-XXI-005.10}\quad printed folio 375.\par
\noindent Printed reading: {\ttfamily\footnotesize XX\allowbreak{}}\par
\noindent Corrected reading: {\ttfamily\footnotesize XXII\allowbreak{}}\par
\noindent Source: {\ttfamily\footnotesize port/PORT\_EN7II\_LOG.\allowbreak{}jsonl\#C894;\allowbreak{} fixes/FIXLOG\_7II\_CUR\allowbreak{}RENT.jsonl\#C894;\allowbreak{} 7II:1112\allowbreak{}}\par

\par\medskip\noindent\hypertarget{app-E7II-XXI-005.11}{}\textbf{E7II-XXI-005.11}\quad printed folio 375.\par
\noindent Printed reading: {\ttfamily\footnotesize =\allowbreak{}}\par
\noindent Corrected reading: {\ttfamily\footnotesize \textbackslash{}equiv\allowbreak{}}\par
\noindent Source: {\ttfamily\footnotesize fr/sga7ii/\{diplomati\allowbreak{}c,current\}/expose\_XX\allowbreak{}I\_body.tex;\allowbreak{} folio\allowbreak{} 375;\allowbreak{} formula\allowbreak{} 248\allowbreak{}}\par

\par\medskip\noindent\hypertarget{app-E7II-XXI-005.12}{}\textbf{E7II-XXI-005.12}\quad printed folio 375.\par
\noindent Printed reading: {\ttfamily\footnotesize =\allowbreak{}}\par
\noindent Corrected reading: {\ttfamily\footnotesize \textbackslash{}equiv\allowbreak{}}\par
\noindent Source: {\ttfamily\footnotesize fr/sga7ii/\{diplomati\allowbreak{}c,current\}/expose\_XX\allowbreak{}I\_body.tex;\allowbreak{} folio\allowbreak{} 375;\allowbreak{} formula\allowbreak{} 248\allowbreak{}}\par

\par\medskip\noindent\hypertarget{app-E7II-XXI-005.13}{}\textbf{E7II-XXI-005.13}\quad printed folio 375.\par
\noindent Printed reading: \emph{(empty)}\par
\noindent Corrected reading: {\ttfamily\footnotesize -\allowbreak{}}\par
\noindent Source: {\ttfamily\footnotesize fr/sga7ii/\{diplomati\allowbreak{}c,current\}/expose\_XX\allowbreak{}I\_body.tex;\allowbreak{} folio\allowbreak{} 375;\allowbreak{} formula\allowbreak{} 248\allowbreak{}}\par

\par\medskip\noindent\hypertarget{app-E7II-XXI-005.14}{}\textbf{E7II-XXI-005.14}\quad printed folio 375.\par
\noindent Printed reading: {\ttfamily\footnotesize .\allowbreak{}}\par
\noindent Corrected reading: {\ttfamily\footnotesize )\allowbreak{}}\par
\noindent Source: {\ttfamily\footnotesize fr/sga7ii/\{diplomati\allowbreak{}c,current\}/expose\_XX\allowbreak{}I\_body.tex;\allowbreak{} folio\allowbreak{} 375;\allowbreak{} formula\allowbreak{} 248\allowbreak{}}\par

\par\medskip\noindent\hypertarget{app-E7II-XXI-005.15}{}\textbf{E7II-XXI-005.15}\quad printed folio 375.\par
\noindent Printed reading: {\ttfamily\footnotesize .\allowbreak{}}\par
\noindent Corrected reading: {\ttfamily\footnotesize )\allowbreak{}}\par
\noindent Source: {\ttfamily\footnotesize fr/sga7ii/\{diplomati\allowbreak{}c,current\}/expose\_XX\allowbreak{}I\_body.tex;\allowbreak{} folio\allowbreak{} 375;\allowbreak{} formula\allowbreak{} 248\allowbreak{}}\par

\par\medskip\noindent\hypertarget{app-E7II-XXI-005.16}{}\textbf{E7II-XXI-005.16}\quad printed folio 376.\par
\noindent Printed reading: {\ttfamily\footnotesize o\allowbreak{}}\par
\noindent Corrected reading: {\ttfamily\footnotesize 0\allowbreak{}}\par
\noindent Source: {\ttfamily\footnotesize fr/sga7ii/\{diplomati\allowbreak{}c,current\}/expose\_XX\allowbreak{}I\_body.tex;\allowbreak{} folio\allowbreak{} 376;\allowbreak{} formula\allowbreak{} 251\allowbreak{}}\par

\par\medskip\noindent\hypertarget{app-E7II-XXI-006.1}{}\textbf{E7II-XXI-006.1}\quad printed folio 376.\par
\noindent Printed reading: \emph{(empty)}\par
\noindent Corrected reading: {\ttfamily\footnotesize (\allowbreak{}}\par
\noindent Source: {\ttfamily\footnotesize fr/sga7ii/\{diplomati\allowbreak{}c,current\}/expose\_XX\allowbreak{}I\_body.tex;\allowbreak{} folio\allowbreak{} 376;\allowbreak{} formula\allowbreak{} 266\allowbreak{}}\par

\par\medskip\noindent\hypertarget{app-E7II-XXI-006.2}{}\textbf{E7II-XXI-006.2}\quad printed folio 376.\par
\noindent Printed reading: \emph{(empty)}\par
\noindent Corrected reading: {\ttfamily\footnotesize )\allowbreak{}}\par
\noindent Source: {\ttfamily\footnotesize fr/sga7ii/\{diplomati\allowbreak{}c,current\}/expose\_XX\allowbreak{}I\_body.tex;\allowbreak{} folio\allowbreak{} 376;\allowbreak{} formula\allowbreak{} 266\allowbreak{}}\par

\par\medskip\noindent\hypertarget{app-E7II-XXI-006.3}{}\textbf{E7II-XXI-006.3}\quad printed folio 376.\par
\noindent Printed reading: {\ttfamily\footnotesize )\allowbreak{}}\par
\noindent Corrected reading: \emph{(empty)}\par
\noindent Source: {\ttfamily\footnotesize fr/sga7ii/\{diplomati\allowbreak{}c,current\}/expose\_XX\allowbreak{}I\_body.tex;\allowbreak{} folio\allowbreak{} 376;\allowbreak{} formula\allowbreak{} 268\allowbreak{}}\par

\par\medskip\noindent\hypertarget{app-E7II-XXI-006.4}{}\textbf{E7II-XXI-006.4}\quad printed folio 376.\par
\noindent Printed reading: \emph{(empty)}\par
\noindent Corrected reading: {\ttfamily\footnotesize )\allowbreak{}}\par
\noindent Source: {\ttfamily\footnotesize fr/sga7ii/\{diplomati\allowbreak{}c,current\}/expose\_XX\allowbreak{}I\_body.tex;\allowbreak{} folio\allowbreak{} 376;\allowbreak{} formula\allowbreak{} 268\allowbreak{}}\par

\par\medskip\noindent\hypertarget{app-E7II-XXI-006.5}{}\textbf{E7II-XXI-006.5}\quad printed folio 376.\par
\noindent Printed reading: {\ttfamily\footnotesize o\allowbreak{}}\par
\noindent Corrected reading: {\ttfamily\footnotesize 0\allowbreak{}}\par
\noindent Source: {\ttfamily\footnotesize fr/sga7ii/\{diplomati\allowbreak{}c,current\}/expose\_XX\allowbreak{}I\_body.tex;\allowbreak{} folio\allowbreak{} 376;\allowbreak{} formula\allowbreak{} 269\allowbreak{}}\par

\par\medskip\noindent\hypertarget{app-E7II-XXI-006.6}{}\textbf{E7II-XXI-006.6}\quad printed folio 376.\par
\noindent Printed reading: {\ttfamily\footnotesize .\allowbreak{}}\par
\noindent Corrected reading: {\ttfamily\footnotesize \textbackslash{}dots\allowbreak{}}\par
\noindent Source: {\ttfamily\footnotesize fr/sga7ii/\{diplomati\allowbreak{}c,current\}/expose\_XX\allowbreak{}I\_body.tex;\allowbreak{} folio\allowbreak{} 376;\allowbreak{} formula\allowbreak{} 275\allowbreak{}}\par

\par\medskip\noindent\hypertarget{app-E7II-XXI-007.1}{}\textbf{E7II-XXI-007.1}\quad printed folio 377.\par
\noindent Printed reading: \emph{(empty)}\par
\noindent Corrected reading: {\ttfamily\footnotesize \textbackslash{}overline\{\textbackslash{}mathbb\allowbreak{} F\}\_q\allowbreak{}}\par
\noindent Source: {\ttfamily\footnotesize fr/sga7ii/\{diplomati\allowbreak{}c,current\}/expose\_XX\allowbreak{}I\_body.tex;\allowbreak{} folio\allowbreak{} 377;\allowbreak{} formula\allowbreak{} 288\allowbreak{}}\par

\par\medskip\noindent\hypertarget{app-E7II-XXI-007.2}{}\textbf{E7II-XXI-007.2}\quad printed folio 377.\par
\noindent Printed reading: \emph{(empty)}\par
\noindent Corrected reading: {\ttfamily\footnotesize \textbackslash{}overline\{\textbackslash{}mathbb\allowbreak{} F\}\_q\allowbreak{}}\par
\noindent Source: {\ttfamily\footnotesize fr/sga7ii/\{diplomati\allowbreak{}c,current\}/expose\_XX\allowbreak{}I\_body.tex;\allowbreak{} folio\allowbreak{} 377\allowbreak{}}\par

\par\medskip\noindent\hypertarget{app-E7II-XXI-008.1}{}\textbf{E7II-XXI-008.1}\quad printed folio 378.\par
\noindent Printed reading: \emph{(empty)}\par
\noindent Corrected reading: {\ttfamily\footnotesize \textbackslash{}operatorname\{Gal\}\allowbreak{}}\par
\noindent Source: {\ttfamily\footnotesize fr/sga7ii/\{diplomati\allowbreak{}c,current\}/expose\_XX\allowbreak{}I\_body.tex;\allowbreak{} folio\allowbreak{} 378;\allowbreak{} formula\allowbreak{} 314\allowbreak{}}\par

\par\medskip\noindent\hypertarget{app-E7II-XXI-008.2}{}\textbf{E7II-XXI-008.2}\quad printed folio 378.\par
\noindent Printed reading: \emph{(empty)}\par
\noindent Corrected reading: {\ttfamily\footnotesize \textbackslash{}overline\{\textbackslash{}mathbb\allowbreak{} F\}\_q\allowbreak{}}\par
\noindent Source: {\ttfamily\footnotesize fr/sga7ii/\{diplomati\allowbreak{}c,current\}/expose\_XX\allowbreak{}I\_body.tex;\allowbreak{} folio\allowbreak{} 378;\allowbreak{} formula\allowbreak{} 314\allowbreak{}}\par

\par\medskip\noindent\hypertarget{app-E7II-XXI-008.3}{}\textbf{E7II-XXI-008.3}\quad printed folio 378.\par
\noindent Printed reading: \emph{(1 espace(s), 1 retour(s) de ligne dans le témoin TeX)}\par
\noindent Corrected reading: {\ttfamily\footnotesize  \textbackslash{}bigoplus\_i\allowbreak{} }\par
\noindent Source: {\ttfamily\footnotesize fr/sga7ii/\{diplomati\allowbreak{}c,current\}/expose\_XX\allowbreak{}I\_body.tex;\allowbreak{} folio\allowbreak{} 378;\allowbreak{} formula\allowbreak{} 315\allowbreak{}}\par

\par\medskip\noindent\hypertarget{app-E7II-XXI-008.4}{}\textbf{E7II-XXI-008.4}\quad printed folio 378.\par
\noindent Printed reading: \emph{(empty)}\par
\noindent Corrected reading: {\ttfamily\footnotesize )\allowbreak{}}\par
\noindent Source: {\ttfamily\footnotesize fr/sga7ii/\{diplomati\allowbreak{}c,current\}/expose\_XX\allowbreak{}I\_body.tex;\allowbreak{} folio\allowbreak{} 378;\allowbreak{} formula\allowbreak{} 315\allowbreak{}}\par

\par\medskip\noindent\hypertarget{app-E7II-XXI-008.5}{}\textbf{E7II-XXI-008.5}\quad printed folio 378.\par
\noindent Printed reading: \emph{(empty)}\par
\noindent Corrected reading: {\ttfamily\footnotesize \_i\allowbreak{}}\par
\noindent Source: {\ttfamily\footnotesize fr/sga7ii/\{diplomati\allowbreak{}c,current\}/expose\_XX\allowbreak{}I\_body.tex;\allowbreak{} folio\allowbreak{} 378;\allowbreak{} formula\allowbreak{} 315\allowbreak{}}\par

\par\medskip\noindent\hypertarget{app-E7II-XXI-009.1}{}\textbf{E7II-XXI-009.1}\quad printed folio 379.\par
\noindent Printed reading: {\ttfamily\footnotesize \textbackslash{};\allowbreak{}}\par
\noindent Corrected reading: {\ttfamily\footnotesize  i\allowbreak{}}\par
\noindent Source: {\ttfamily\footnotesize ED-EN-FORMULA-G231;\allowbreak{} fr/sga7ii/\{diplomati\allowbreak{}c,current\}/expose\_XX\allowbreak{}I\_body.tex;\allowbreak{} folio\allowbreak{} 379\allowbreak{}}\par

\par\medskip\noindent\hypertarget{app-E7II-XXI-010.1}{}\textbf{E7II-XXI-010.1}\quad printed folio 381.\par
\noindent Printed reading: {\ttfamily\footnotesize /\allowbreak{}}\par
\noindent Corrected reading: \emph{(empty)}\par
\noindent Source: {\ttfamily\footnotesize fr/sga7ii/\{diplomati\allowbreak{}c,current\}/expose\_XX\allowbreak{}I\_body.tex;\allowbreak{} folio\allowbreak{} 381\allowbreak{}}\par

\par\medskip\noindent\hypertarget{app-E7II-XXI-011.1}{}\textbf{E7II-XXI-011.1}\quad printed folio 382.\par
\noindent Printed reading: \emph{(empty)}\par
\noindent Corrected reading: {\ttfamily\footnotesize i=\allowbreak{}}\par
\noindent Source: {\ttfamily\footnotesize fr/sga7ii/\{diplomati\allowbreak{}c,current\}/expose\_XX\allowbreak{}I\_body.tex;\allowbreak{} folio\allowbreak{} 382;\allowbreak{} formula\allowbreak{} 381\allowbreak{}}\par

\par\medskip\noindent\hypertarget{app-E7II-XXI-012.1}{}\textbf{E7II-XXI-012.1}\quad printed folio 386.\par
\noindent Printed reading: \emph{(empty)}\par
\noindent Corrected reading: {\ttfamily\footnotesize \textbackslash{}underline\allowbreak{} }\par
\noindent Source: {\ttfamily\footnotesize ED-EN-FORMULA-G233;\allowbreak{} fr/sga7ii/\{diplomati\allowbreak{}c,current\}/expose\_XX\allowbreak{}I\_body.tex;\allowbreak{} folio\allowbreak{} 386\allowbreak{}}\par

\par\medskip\noindent\hypertarget{app-E7II-XXI-012.2}{}\textbf{E7II-XXI-012.2}\quad printed folio 386.\par
\noindent Printed reading: {\ttfamily\footnotesize F\allowbreak{}}\par
\noindent Corrected reading: {\ttfamily\footnotesize E\allowbreak{}}\par
\noindent Source: {\ttfamily\footnotesize fr/sga7ii/\{diplomati\allowbreak{}c,current\}/expose\_XX\allowbreak{}I\_body.tex;\allowbreak{} folio\allowbreak{} 386;\allowbreak{} formula\allowbreak{} 502\allowbreak{}}\par

\par\medskip\noindent\hypertarget{app-E7II-XXI-013.1}{}\textbf{E7II-XXI-013.1}\quad printed folio 388.\par
\noindent Printed reading: \emph{(empty)}\par
\noindent Corrected reading: {\ttfamily\footnotesize .2\allowbreak{}}\par
\noindent Source: {\ttfamily\footnotesize port/decisions/decis\allowbreak{}ions\_en7II\_r4\_001.js\allowbreak{}on\#C895;\allowbreak{} fixes/FIXLOG\_7II\_CUR\allowbreak{}RENT.jsonl\#C895;\allowbreak{} 7II:1117\allowbreak{}}\par

\par\medskip\noindent\hypertarget{app-E7II-XXI-013.2}{}\textbf{E7II-XXI-013.2}\quad printed folio 388.\par
\noindent Printed reading: \emph{(empty)}\par
\noindent Corrected reading: {\ttfamily\footnotesize \{\allowbreak{}}\par
\noindent Source: {\ttfamily\footnotesize fr/sga7ii/\{diplomati\allowbreak{}c,current\}/expose\_XX\allowbreak{}I\_body.tex;\allowbreak{} folio\allowbreak{} 388\allowbreak{}}\par

\par\medskip\noindent\hypertarget{app-E7II-XXI-013.3}{}\textbf{E7II-XXI-013.3}\quad printed folio 388.\par
\noindent Printed reading: \emph{(empty)}\par
\noindent Corrected reading: {\ttfamily\footnotesize '\}\allowbreak{}}\par
\noindent Source: {\ttfamily\footnotesize fr/sga7ii/\{diplomati\allowbreak{}c,current\}/expose\_XX\allowbreak{}I\_body.tex;\allowbreak{} folio\allowbreak{} 388\allowbreak{}}\par

\par\medskip\noindent\hypertarget{app-E7II-XXI-014.1}{}\textbf{E7II-XXI-014.1}\quad printed folio 389.\par
\noindent Printed reading: {\ttfamily\footnotesize Let\allowbreak{}}\par
\noindent Corrected reading: {\ttfamily\footnotesize Put\allowbreak{} \$r=\textbackslash{}dim\_E\allowbreak{} V\textasciicircum{}I\$,\allowbreak{} and\allowbreak{} let\allowbreak{}}\par
\noindent Source: {\ttfamily\footnotesize ED-EN-FORMULA-G223;\allowbreak{} fr/sga7ii/\{diplomati\allowbreak{}c,current\}/expose\_XX\allowbreak{}I\_body.tex;\allowbreak{} folio\allowbreak{} 389\allowbreak{}}\par

\par\medskip\noindent\hypertarget{app-E7II-XXI-014.2}{}\textbf{E7II-XXI-014.2}\quad printed folio 389.\par
\noindent Printed reading: {\ttfamily\footnotesize n\allowbreak{}}\par
\noindent Corrected reading: {\ttfamily\footnotesize r\allowbreak{}}\par
\noindent Source: {\ttfamily\footnotesize ED-EN-FORMULA-G223;\allowbreak{} fr/sga7ii/\{diplomati\allowbreak{}c,current\}/expose\_XX\allowbreak{}I\_body.tex;\allowbreak{} folio\allowbreak{} 389\allowbreak{}}\par

\par\medskip\noindent\hypertarget{app-E7II-XXI-014.3}{}\textbf{E7II-XXI-014.3}\quad printed folio 389.\par
\noindent Printed reading: \emph{(empty)}\par
\noindent Corrected reading: {\ttfamily\footnotesize \textbackslash{}qquad(1\textbackslash{}leq\allowbreak{} i\textbackslash{}leq\allowbreak{} r)\allowbreak{}}\par
\noindent Source: {\ttfamily\footnotesize ED-EN-FORMULA-G223;\allowbreak{} fr/sga7ii/\{diplomati\allowbreak{}c,current\}/expose\_XX\allowbreak{}I\_body.tex;\allowbreak{} folio\allowbreak{} 389\allowbreak{}}\par

\par\medskip\noindent\hypertarget{app-E7II-XXI-015.1}{}\textbf{E7II-XXI-015.1}\quad printed folio 390.\par
\noindent Printed reading: {\ttfamily\footnotesize \$\allowbreak{} \$\allowbreak{}}\par
\noindent Corrected reading: {\ttfamily\footnotesize :\allowbreak{}}\par
\noindent Source: {\ttfamily\footnotesize port/decisions/decis\allowbreak{}ions\_en7II\_r4\_001.js\allowbreak{}on\#C897;\allowbreak{} fixes/FIXLOG\_7II\_CUR\allowbreak{}RENT.jsonl\#C897;\allowbreak{} 7II:1122\allowbreak{}}\par

\par\medskip\noindent\hypertarget{app-E7II-XXI-015.2}{}\textbf{E7II-XXI-015.2}\quad printed folio 390.\par
\noindent Printed reading: \emph{(empty)}\par
\noindent Corrected reading: {\ttfamily\footnotesize 5.\allowbreak{}}\par
\noindent Source: {\ttfamily\footnotesize ED-EN-NONMATH-G051;\allowbreak{} F7II-XXI-038.1;\allowbreak{} 7II:684;\allowbreak{} 7II:1123;\allowbreak{} fr/sga7ii/\{diplomati\allowbreak{}c,current\}/expose\_XX\allowbreak{}I\_body.tex;\allowbreak{} folio\allowbreak{} 390\allowbreak{}}\par

\par\medskip\noindent\hypertarget{app-E7II-XXI-015.3}{}\textbf{E7II-XXI-015.3}\quad printed folio 390.\par
\noindent Printed reading: {\ttfamily\footnotesize irreducible\allowbreak{}}\par
\noindent Corrected reading: {\ttfamily\footnotesize indecomposable\allowbreak{}}\par
\noindent Source: {\ttfamily\footnotesize ED-EN-NONMATH-G051;\allowbreak{} F7II-XXI-038.1;\allowbreak{} 7II:684;\allowbreak{} 7II:1123;\allowbreak{} fr/sga7ii/\{diplomati\allowbreak{}c,current\}/expose\_XX\allowbreak{}I\_body.tex;\allowbreak{} folio\allowbreak{} 390\allowbreak{}}\par

\par\medskip\noindent\hypertarget{app-E7II-XXI-016.1}{}\textbf{E7II-XXI-016.1}\quad printed folio 391.\par
\noindent Printed reading: \emph{(empty)}\par
\noindent Corrected reading: {\ttfamily\footnotesize \textbackslash{}underline\allowbreak{} }\par
\noindent Source: {\ttfamily\footnotesize ED-EN-FORMULA-G235;\allowbreak{} fr/sga7ii/\{diplomati\allowbreak{}c,current\}/expose\_XX\allowbreak{}I\_body.tex;\allowbreak{} folio\allowbreak{} 391\allowbreak{}}\par

\par\medskip\noindent\hypertarget{app-E7II-XXI-017.1}{}\textbf{E7II-XXI-017.1}\quad printed folio 393.\par
\noindent Printed reading: \emph{(empty)}\par
\noindent Corrected reading: {\ttfamily\footnotesize \textbackslash{}log\allowbreak{}}\par
\noindent Source: {\ttfamily\footnotesize fr/sga7ii/\{diplomati\allowbreak{}c,current\}/expose\_XX\allowbreak{}I\_body.tex;\allowbreak{} folio\allowbreak{} 393;\allowbreak{} formula\allowbreak{} 745\allowbreak{}}\par

\par\medskip\noindent\hypertarget{app-E7II-XXI-017.2}{}\textbf{E7II-XXI-017.2}\quad printed folio 393.\par
\noindent Printed reading: \emph{(empty)}\par
\noindent Corrected reading: {\ttfamily\footnotesize \textbackslash{}underline\allowbreak{} }\par
\noindent Source: {\ttfamily\footnotesize fr/sga7ii/\{diplomati\allowbreak{}c,current\}/expose\_XX\allowbreak{}I\_body.tex;\allowbreak{} folio\allowbreak{} 393\allowbreak{}}\par

\par\medskip\noindent\hypertarget{app-E7II-XXI-017.3}{}\textbf{E7II-XXI-017.3}\quad printed folio 394.\par
\noindent Printed reading: {\ttfamily\footnotesize 5\allowbreak{}}\par
\noindent Corrected reading: {\ttfamily\footnotesize 2\allowbreak{}}\par
\noindent Source: {\ttfamily\footnotesize ED-EN-NONMATH-G039;\allowbreak{} F7II-XXI-042.1;\allowbreak{} WITNESS-F7II-XXI-042\allowbreak{}.1;\allowbreak{} fr/sga7ii/\{diplomati\allowbreak{}c,current\}/expose\_XX\allowbreak{}I\_body.tex;\allowbreak{} folio\allowbreak{} 394\allowbreak{}}\par

\par\medskip\noindent\hypertarget{app-E7II-XXI-018.1}{}\textbf{E7II-XXI-018.1}\quad printed folio 394.\par
\noindent Printed reading: {\ttfamily\footnotesize 5.2.2\allowbreak{}}\par
\noindent Corrected reading: {\ttfamily\footnotesize the\allowbreak{} desired\allowbreak{} assertion\allowbreak{}}\par
\noindent Source: {\ttfamily\footnotesize ED-EN-NONMATH-G040;\allowbreak{} F7II-XXI-043.1;\allowbreak{} WITNESS-F7II-XXI-043\allowbreak{}.1;\allowbreak{} fr/sga7ii/\{diplomati\allowbreak{}c,current\}/expose\_XX\allowbreak{}I\_body.tex;\allowbreak{} folio\allowbreak{} 394\allowbreak{}}\par

\par\medskip\noindent\hypertarget{app-E7II-XXI-018.2}{}\textbf{E7II-XXI-018.2}\quad printed folio 395.\par
\noindent Printed reading: \emph{(empty)}\par
\noindent Corrected reading: {\ttfamily\footnotesize )\allowbreak{}}\par
\noindent Source: {\ttfamily\footnotesize fr/sga7ii/\{diplomati\allowbreak{}c,current\}/expose\_XX\allowbreak{}I\_body.tex;\allowbreak{} folio\allowbreak{} 395;\allowbreak{} formula\allowbreak{} 813;\allowbreak{} port/decisions/decis\allowbreak{}ions\_en7II\_r2\_003.js\allowbreak{}on\#C766;\allowbreak{} fixes/FIXLOG\_7II\_CUR\allowbreak{}RENT.jsonl\#C766;\allowbreak{} 7II:688\allowbreak{}}\par

\par\medskip\noindent\hypertarget{app-E7II-XXI-019.1}{}\textbf{E7II-XXI-019.1}\quad printed folio 396.\par
\noindent Printed reading: \emph{(empty)}\par
\noindent Corrected reading: {\ttfamily\footnotesize \textbackslash{}circ\allowbreak{}}\par
\noindent Source: {\ttfamily\footnotesize fr/sga7ii/\{diplomati\allowbreak{}c,current\}/expose\_XX\allowbreak{}I\_body.tex;\allowbreak{} folio\allowbreak{} 396;\allowbreak{} formula\allowbreak{} 847\allowbreak{}}\par

\par\medskip\noindent\hypertarget{app-E7II-XXI-020.1}{}\textbf{E7II-XXI-020.1}\quad printed folio 399.\par
\noindent Printed reading: \emph{(empty)}\par
\noindent Corrected reading: {\ttfamily\footnotesize 5.\allowbreak{}}\par
\noindent Source: {\ttfamily\footnotesize ED-EN-NONMATH-G041;\allowbreak{} F7II-XXI-048.1;\allowbreak{} 7II:692;\allowbreak{} 7II:831;\allowbreak{} fr/sga7ii/\{diplomati\allowbreak{}c,current\}/expose\_XX\allowbreak{}I\_body.tex;\allowbreak{} folio\allowbreak{} 399\allowbreak{}}\par

\par\endgroup\normalsize\fi
\clearpage\SGAtoc{Exposé XXII}
\selectlanguage{english}
\clearpage
\SGApage{II}{401}

\begin{center}
{\Large\bfseries EXPOSÉ XXII}\par
\vspace{1.5cm}
{\Large\bfseries A CONGRUENCE FORMULA FOR THE ZETA FUNCTION}\par
\vspace{0.8cm}
{\large by N. Katz}\par
\end{center}

\vspace{1.5cm}
\section*{0. Introduction}

Let $X$ be a proper scheme over $\mathbb F_q$. Its zeta function, regarded as
an element of $\mathbb Z[[T]]$, may be reduced modulo $p$; we wish to compute
it modulo $p$. We obtain an expression of the kind one might have expected
from fixed-point formulas of ``Woods Hole'' type: an alternating product of
characteristic polynomials of a suitable ``Frobenius'' acting on any one of
the following three cohomology spaces:

\begin{enumerate}[label=\arabic*)]
\item coherent cohomology $H^i(X,\mathcal O_X)$ (3.1.1);
\item de Rham cohomology
  $H^i_{\mathrm{DR}}(X)=\mathbb H^i(X,\Omega_X^\bullet)$ (3.1.1);
\item étale cohomology $H^i_{\mathrm{et}}(X,\mathbb Z/p\mathbb Z)$ (3.1.2).
\end{enumerate}

This expression should be viewed as an improvement of a special case of the
Woods Hole trace formula, namely that of the structure sheaf and the
Frobenius endomorphism. That case is indeed a corollary of the congruence
formula (cf. 3.2.1).

The method of this exposé is completely naïve. It consists in reducing to the
case of a hypersurface and then ``reducing modulo $p$'' Dwork's computation
of the actual zeta function of a hypersurface. By ``reducing'' carefully, one
obtains the improvement 5.0.6, due to Ax~[1] (and also to Dwork in the smooth
case [4, \S 7]), of the well-known Warning--Chevalley theorem~[8].

\SGApage{II}{\(\approx\)402}
Another method, entirely cohomological and in the spirit of SGA~5, which also
applies to ``twisted coefficients,'' has just been found by Deligne and ought
to appear in the new edition of SGA~5. It rests on the trace formula of Woods
Hole type and on Deligne's ``symmetric Künneth formula''
(SGA~4, XVII~5.5.21).

\section*{1. A reminder on ``$p$-linear'' operations}

\subsection*{1.0.}
Let $V$ be a finite-dimensional vector space over a field $k$ of
characteristic $p>0$, and let $q$ be a power of $p$. A map

\begin{equation}\tag{1.0.1}
\varphi:V\longrightarrow V
\end{equation}

is called $q$-\emph{linear} if it is additive and satisfies

\begin{equation}\tag{1.0.2}
\varphi(\lambda v)=\lambda^q\varphi(v)
\qquad\text{for every }\lambda\in k,\ v\in V.
\end{equation}

\subsection*{1.0.3.}
Notice that the $n$-th iterate $\varphi^n$ of a $q$-linear endomorphism is
$q^n$-linear and that, if $k$ is not perfect, the image $\varphi(V)$ of a
$q$-linear endomorphism need not be a $k$-vector subspace.

\subsection*{1.0.4.}
For every extension $K/k$, a $q$-linear endomorphism $\varphi$ of $V$ gives
rise to a unique $q$-linear endomorphism $\varphi_K$ of
$V_K=V\otimes_kK$ satisfying

\begin{equation}\tag{1.0.5}
\varphi_K(\mu\otimes v)=\mu^q\varphi(v)
\qquad\text{for }\mu\in K,\ v\in V.
\end{equation}

\SGApage{II}{\(\approx\)403}
\subsection*{1.0.6.}
A $q$-linear endomorphism $\varphi$ of $V$ is called \emph{semisimple} if
its image spans $V$ over $k$, and \emph{nilpotent} if some iterate is zero.
Define the following $\varphi$-stable $k$-subspaces:

\begin{equation}\tag{1.0.7}
V_{\mathrm{ss}}=
\bigcap_{n\geq1}\langle\varphi^n(V)\rangle_k,
\end{equation}

\begin{equation}\tag{1.0.8}
V_{\mathrm{nilp}}=
\bigcup_{n\geq1}\Ker(\varphi^n:V\longrightarrow V).
\end{equation}

Clearly, the restriction of $\varphi$ to $V_{\mathrm{ss}}$ is semisimple,
while its restriction to $V_{\mathrm{nilp}}$ is nilpotent. \underline{If} $k$ \underline{is
perfect}, then

\begin{equation}\tag{1.0.9}
V=V_{\mathrm{ss}}\oplus V_{\mathrm{nilp}}.
\end{equation}

If $k$ is not perfect, however, (1.0.9) is false, since there exist injective
$q$-linear endomorphisms which are not semisimple. For an arbitrary extension
$K/k$, one has

\begin{equation}\tag{1.0.10}
(V_K)_{\mathrm{ss}}=V_{\mathrm{ss}}\otimes_kK.
\end{equation}



\subsubsection*{Proposition 1.1.}
Let $\varphi$ be a $q$-linear endomorphism of a finite-dimensional vector
space $V$ over a separably closed field $k$. Suppose that $\varphi$ is
semisimple, and put

\[
V^{1-\varphi}
=\text{\underline{the $\mathbb F_q$-vector space of fixed points of} $\varphi$}
=\Ker(1-\varphi).
\]

Then

\begin{equation}\tag{1.1.1}
V\xleftarrow{\ \sim\ }
(V^{1-\varphi})\otimes_{\mathbb F_q}k.
\end{equation}

\SGApage{II}{404}
\subsubsection*{Corollary 1.1.2.}
There is an isomorphism

\[
V^{1-\varphi^n}\xleftarrow{\ \sim\ }
(V^{1-\varphi})\otimes_{\mathbb F_q}\mathbb F_{q^n}.
\]

\medskip
\noindent\emph{Proof of Proposition 1.1.}
(Lang, cf.~[6], p.~119.)
Let $\underline{e}$ be a $k$-basis of $V$, and write

\begin{equation}\tag{1.1.3}
\varphi(\underline{e})=A\underline{e},\qquad A\in\mathrm{GL}(n,k).
\end{equation}

Let

\begin{equation}\tag{1.1.4}
F:\mathrm{GL}(n,k)\longrightarrow\mathrm{GL}(n,k)
\end{equation}

denote the map which raises every coordinate to the $q$-th power. For
$g\in\mathrm{GL}(n,k)$, the basis $g\underline{e}$ of $V$ consists of fixed points of
$\varphi$ if and only if

\begin{equation}\tag{1.1.5}
F(g)Ag^{-1}=1.
\end{equation}

Let $\bar k$ be an algebraic closure of $k$. The group
$\mathrm{GL}(n,\bar k)$ acts on the variety $\mathrm{GL}(n,\bar k)$ by

\begin{equation}\tag{1.1.6}
\left\{
\begin{aligned}
\mathrm{GL}(n,\bar k)\times\mathrm{GL}(n,\bar k)
&\longrightarrow\mathrm{GL}(n,\bar k),\\
(g,A)&\longmapsto F(g)Ag^{-1}.
\end{aligned}
\right.
\end{equation}

For fixed $A$, the map

\begin{equation}\tag{1.1.7}
g\longmapsto F(g)Ag^{-1}
\end{equation}

is étale, as one sees by computing its tangent map. Its orbits are therefore
open. Since $\mathrm{GL}(n,\bar k)$ is irreducible, there is only one orbit;
in other words, the map

\SGAvar{\begin{equation}\tag{1.1.8}
\begin{aligned}
\mathrm{Gl}(n,\bar k)&\longrightarrow\mathrm{Gl}(n,\bar k),\\
g&\longrightarrow F(g^{-1})\,g
\end{aligned}
\end{equation}

}{\begin{equation}\tag{1.1.8}
\begin{aligned}
\mathrm{GL}(n,\bar k)&\longrightarrow\mathrm{GL}(n,\bar k),\\
g&\longrightarrow F(g^{-1})\,g
\end{aligned}
\end{equation}

}{E7II-XXII-001}\SGApage{II}{\(\approx\)405}
is surjective. Since it is also étale and $k$ is separably closed, the map

\begin{equation}\tag{1.1.9}
\begin{aligned}
\mathrm{GL}(n,k)&\longrightarrow\mathrm{GL}(n,k),\\
g&\longrightarrow F(g^{-1})\,g
\end{aligned}
\end{equation}

is surjective. Thus $V$ has a $k$-basis consisting of fixed points of
$\varphi$. Since $\varphi$ is $q$-linear, every fixed point of $\varphi$ lies
in the $\mathbb F_q$-vector space spanned by such a $k$-basis. This proves
Proposition~1.1.

\subsubsection*{Corollary 1.1.10.}
Let $\varphi$ be a $q$-linear endomorphism of a finite-dimensional vector
space $V$ over a separably closed field $k$. There is a canonical
isomorphism

\begin{equation}\tag{1.1.11}
V_{\mathrm{ss}}\xleftarrow{\ \sim\ }
V^{1-\varphi}\otimes_{\mathbb F_q}k.
\end{equation}

Indeed, it is enough to apply Proposition~1.1 to $V_{\mathrm{ss}}$.

\SGAvar{\subsubsection*{Proposition 1.2.}
Let $\varphi$ be a $q$-linear endomorphism of a finite-dimensional vector
space over a separably closed field $k$. Then the additive map

}{\subsubsection*{Proposition 1.2.}
Let $\varphi$ be a $q$-linear endomorphism of a finite-dimensional vector
space $V$ over a separably closed field $k$. Then the additive map

}{E7II-XXII-002}\[
1-\varphi:V\longrightarrow V
\]

is surjective.

\medskip
\noindent\emph{Proof.}
There is an exact sequence of abelian groups

\[
0\longrightarrow V_{\mathrm{ss}}\longrightarrow V
\longrightarrow V/V_{\mathrm{ss}}\longrightarrow0.
\]


\SGApage{II}{\(\approx\)406}
on which $1-\varphi$ acts. By the snake lemma, it is enough to treat
$V_{\mathrm{ss}}$, which is immediate from Proposition~1.1 by using a basis
consisting of fixed points, and $V/V_{\mathrm{ss}}$, on which $\varphi$ is
nilpotent and hence $1-\varphi$ is invertible.



\section*{2. Coherent cohomology and étale cohomology}

\subsubsection*{Proposition 2.0.}
Let $\bar X$ be a proper scheme over a separably closed field $k$, and let

\begin{equation}\tag{2.0.1}
\overline{\mathfrak F}_p:H^{\cdot}(\bar X,\mathcal O_{\bar X})
\longrightarrow H^{\cdot}(\bar X,\mathcal O_{\bar X})
\end{equation}

denote the $p$-linear map induced by

\begin{equation}\tag{2.0.2}
\overline{\mathfrak F}_p:\mathcal O_{\bar X}\longrightarrow\mathcal O_{\bar X},
\qquad \overline{\mathfrak F}_p(f)=f^p.
\end{equation}

Denote by

\[
H^{\cdot}_{\mathrm{et}}(\bar X,\mathbb Z/p\mathbb Z)
\]

the étale cohomology groups with coefficients in the constant sheaf
$\mathbb Z/p\mathbb Z$. Then there are canonical isomorphisms

\begin{equation}\tag{2.0.3}
H^i_{\mathrm{et}}(\bar X,\mathbb Z/p\mathbb Z)
\simeq
(H^i(\bar X,\mathcal O_{\bar X}))^{1-\overline{\mathfrak F}_p}.
\end{equation}

\medskip
\noindent\emph{Proof.}
There is an exact sequence of sheaves on $\bar X_{\mathrm{et}}$:

\begin{equation}\tag{2.0.4}
0\longrightarrow\mathbb Z/p\mathbb Z\longrightarrow\mathcal O_{\bar X}
\xrightarrow{\ 1-\mathfrak F_p\ }\mathcal O_{\bar X}\longrightarrow0,
\end{equation}

which gives an exact cohomology sequence

\begin{equation}\tag{2.0.5}
\cdots\longrightarrow H^i_{\mathrm{et}}(\bar X,\mathbb Z/p\mathbb Z)
\longrightarrow H^i_{\mathrm{et}}(\bar X,\mathcal O_{\bar X})
\xrightarrow{\ 1-\mathfrak F_p\ }
H^i_{\mathrm{et}}(\bar X,\mathcal O_{\bar X})\longrightarrow\cdots .
\end{equation}

One has
$H^i_{\mathrm{et}}(\bar X,\mathcal O_{\bar X})\simeq
H^i(\bar X,\mathcal O_{\bar X})$; this is a \underline{finite-dimensional}
$k$-vector space because $\bar X$ is proper. Applying Proposition~1.2 gives
exact sequences

\SGApage{II}{\(\approx\)407}
\begin{equation}\tag{2.0.6}
0\longrightarrow H^i_{\mathrm{et}}(\bar X,\mathbb Z/p\mathbb Z)
\longrightarrow H^i(\bar X,\mathcal O_{\bar X})
\xrightarrow{\ 1-\mathfrak F_p\ }
H^i(\bar X,\mathcal O_{\bar X})\longrightarrow0,
\end{equation}

which proves the assertion. Proposition~1.1.10 now gives the following
corollary of Proposition~2.0.

\subsubsection*{Corollary 2.1.}
For $\bar X$ proper over a separably closed field $k$, there is a canonical
isomorphism

\begin{equation}\tag{2.1.1}
H^i(\bar X,\mathcal O_{\bar X})_{\mathrm{ss}}
\xleftarrow{\ \sim\ }
H^i_{\mathrm{et}}(\bar X,\mathbb Z/p\mathbb Z)
\otimes_{\mathbb F_p}k.
\end{equation}


\subsection*{2.2.}
Let now $f:X\longrightarrow\Spec(\mathbb F_q)$ be a proper scheme over
$\mathbb F_q$. Let $k$ be an algebraic closure of $\mathbb F_q$, and put

\begin{equation}\tag{2.2.1}
\bar X=X\times_{\mathbb F_q}k.
\end{equation}

The étale sheaf $R^if_*(\mathbb Z/p\mathbb Z)$ on
$\Spec(\mathbb F_q)$ identifies with
$H^i_{\mathrm{et}}(\bar X,\mathbb Z/p\mathbb Z)$, viewed as a
$\Gal(k/\mathbb F_q)$-module. Let

\begin{equation}\tag{2.2.2}
\mathfrak F_q:H^i(X,\mathcal O_X)\longrightarrow H^i(X,\mathcal O_X)
\end{equation}

denote the linear operation (it is $q$-linear, but the ground field is
$\mathbb F_q$) induced by

\begin{equation}\tag{2.2.3}
\mathfrak F_q:\mathcal O_X\longrightarrow\mathcal O_X,
\qquad \mathfrak F_q(f)=f^q.
\end{equation}

\subsubsection*{Proposition 2.2.4.}
Let $\varphi_q\in\Gal(k/\mathbb F_q)$ be the ``Frobenius'' element
$\varphi_q(\lambda)=\lambda^q$, and let $i\in\mathbb Z$. There is an
identity of characteristic polynomials

\begin{equation}\tag{2.2.4.1}
\det_{\mathbb F_p}\!\left(
1-t\varphi_q^{-1}\mid H^i_{\mathrm{et}}(\bar X,\mathbb Z/p\mathbb Z)
\right)
=
\det_{\mathbb F_q}\!\left(
1-t\mathfrak F_q\mid H^i(X,\mathcal O_X)
\right).
\end{equation}

Here the left-hand side is the characteristic polynomial of the element
$\varphi_q^{-1}$ of $\pi_1(\Spec(\mathbb F_q))$ acting on the Galois module
$H^i_{\mathrm{et}}(\bar X,\mathbb Z/p\mathbb Z)$.

Proposition~2.2.4 is a corollary of the following result.

\SGApage{II}{408}
\subsubsection*{Proposition 2.2.5.}
Let $X$ be proper over $\Spec(\mathbb F_q)$. There is a canonical
isomorphism

\begin{equation}\tag{2.2.6}
H^i(X,\mathcal O_X)_{\mathrm{ss}}\otimes_{\mathbb F_q}k
\simeq
H^i_{\mathrm{et}}(\bar X,\mathbb Z/p\mathbb Z)
\otimes_{\mathbb F_p}k
\end{equation}

\SGAvar{which carries the automorphism induced by $\mathfrak F_q$ on the left-hand
side to the automorphism $\varphi_q^{-1}\otimes_{\mathbb F_q}k$ of the right-hand side.
The latter may also be interpreted as
$\mathrm{fr}_q^{X*}\otimes_{\mathbb F_p}k$, where
$\mathrm{fr}_q^X:\bar X\longrightarrow\bar X$ is the Frobenius
$k$-endomorphism (cf. SGA~5, XVIII~2.3.4).

}{which carries the automorphism induced by $\mathfrak F_q$ on the left-hand
side to the automorphism induced by $\varphi_q^{-1}$ on the right-hand side.
The latter may also be interpreted as
$\mathrm{fr}_q^{X*}\otimes_{\mathbb F_p}k$, where
$\mathrm{fr}_q^X:\bar X\longrightarrow\bar X$ is the Frobenius
$k$-endomorphism (cf. SGA~5, XVIII~2.3.4).

}{E7II-XXII-003}\medskip
\noindent\emph{Proof.}
Since

\begin{equation}\tag{2.2.7}
H^i(\bar X,\mathcal O_{\bar X})\simeq
H^i(X,\mathcal O_X)\otimes_{\mathbb F_q}k,
\end{equation}

Proposition~1.0.10 gives

\begin{equation}\tag{2.2.8}
H^i(\bar X,\mathcal O_{\bar X})_{\mathrm{ss}}\simeq
H^i(X,\mathcal O_X)_{\mathrm{ss}}\otimes_{\mathbb F_q}k.
\end{equation}

Thus the isomorphism (2.2.6) is obtained by composing (2.2.8) with
(2.1.1).

In view of (2.2.7), the $q$-linear operation

\begin{equation}\tag{2.2.9}
\mathfrak F_q:H^i(X,\mathcal O_X)\longrightarrow H^i(X,\mathcal O_X)
\end{equation}

has a $q$-linear extension (cf. 1.0.5)

\begin{equation}\tag{2.2.10}
\overline{\mathfrak F}_q:H^i(\bar X,\mathcal O_{\bar X})
\longrightarrow H^i(\bar X,\mathcal O_{\bar X}),
\end{equation}

and also a $k$-linear extension, again denoted by
$\mathfrak F_q$,

\begin{equation}\tag{2.2.11}
\mathfrak F_q:H^i(\bar X,\mathcal O_{\bar X})
\longrightarrow H^i(\bar X,\mathcal O_{\bar X}).
\end{equation}

\SGApage{II}{\(\approx\)409}
Let

\begin{equation}\tag{2.2.12}
\varphi_q:H^i(\bar X,\mathcal O_{\bar X})
\longrightarrow H^i(\bar X,\mathcal O_{\bar X})
\end{equation}

\SGAvar{denote the $q$-linear operation whose fixed points are
$H^i(X,\mathcal O_X)$. Under (2.2.7), the operators
$\mathfrak F_q$, $\varphi_q$, and $\overline{\mathfrak F}_q$ may also be
written respectively as
$\mathfrak F_q\otimes\mathrm{id}_k$,
$\mathrm{id}\otimes\varphi_q$, and
$\mathfrak F_k\otimes\varphi_k$. They therefore commute with one another,
and

}{denote the $q$-linear operation whose fixed points are
$H^i(X,\mathcal O_X)$. Under (2.2.7), the operators
$\mathfrak F_q$, $\varphi_q$, and $\overline{\mathfrak F}_q$ may also be
written respectively as
$\mathfrak F_q\otimes\mathrm{id}_k$,
$\mathrm{id}\otimes\varphi_q$, and
$\mathfrak F_q\otimes\varphi_q$. They therefore commute with one another,
and

}{E7II-XXII-004}\begin{equation}\tag{2.2.13}
\overline{\mathfrak F}_q
=\varphi_q\mathfrak F_q
=\mathfrak F_q\circ\varphi_q.
\end{equation}

They all commute with $\overline{\mathfrak F}_p$, so $\mathfrak F_q$ and $\varphi_q$
act on both
$H^i(\bar X,\mathcal O_{\bar X})^{1-\overline{\mathfrak F}_q}$ and
$H^i(\bar X,\mathcal O_{\bar X})^{1-\overline{\mathfrak F}_p}$. By (2.2.13),
$\mathfrak F_q$ and $\varphi_q$ are inverse to one another \underline{on}
$H^i(\bar X,\mathcal O_{\bar X})^{1-\overline{\mathfrak F}_q}$. Hence the
action of $\mathfrak F_q$ on

\[
H^i(\bar X,\mathcal O_{\bar X})_{\mathrm{ss}}
\simeq
H^i(\bar X,\mathcal O_{\bar X})^{1-\overline{\mathfrak F}_q}
\otimes_{\mathbb F_q}k
\qquad\text{(cf. 1.1.11)}
\]

is inverse to the action of $\varphi_q\otimes\mathrm{id}_k$ on

\[
H^i(\bar X,\mathcal O_{\bar X})^{1-\overline{\mathfrak F}_q}
\otimes_{\mathbb F_q}k.
\]

Writing

\SGAvar{\begin{equation}\tag{2.2.14}
H^i(\bar X,\mathcal O_{\bar X})^{1-{\mathfrak F}_q}
\simeq
H^i(\bar X,\mathcal O_{\bar X})^{1-{\mathfrak F}_p}
\otimes_{\mathbb F_p}\mathbb F_q,
\end{equation}

the automorphism $\varphi_q$ of the left-hand side of (2.2.14) becomes,
via (2.2.14), the automorphism induced by $\varphi_q$ on
$H^i(\bar X,\mathcal O_{\bar X})^{1-\overline{\mathfrak F}_p}
\otimes_{\mathbb F_p}\mathbb F_q$. Consequently the automorphism induced by
$\varphi_q$ on
$H^i(\bar X,\mathcal O_{\bar X})_{\mathrm{ss}}$ becomes, via (2.1.1), the
automorphism induced by $\varphi_q$ on
$H^i_{\mathrm{et}}(\bar X,\mathbb Z/p\mathbb Z)
\otimes_{\mathbb F_p}k$. This proves Proposition~2.2.5.



}{\begin{equation}\tag{2.2.14}
H^i(\bar X,\mathcal O_{\bar X})^{1-\overline{\mathfrak F}_q}
\simeq
H^i(\bar X,\mathcal O_{\bar X})^{1-\overline{\mathfrak F}_p}
\otimes_{\mathbb F_p}\mathbb F_q,
\end{equation}

the automorphism $\varphi_q$ of the left-hand side of (2.2.14) becomes,
via (2.2.14), the automorphism induced by $\varphi_q$ on
$H^i(\bar X,\mathcal O_{\bar X})^{1-\overline{\mathfrak F}_p}
\otimes_{\mathbb F_p}\mathbb F_q$. Consequently the automorphism induced by
$\varphi_q$ on
$H^i(\bar X,\mathcal O_{\bar X})_{\mathrm{ss}}$ becomes, via (2.1.1), the
automorphism induced by $\varphi_q^{-1}$ on
$H^i_{\mathrm{et}}(\bar X,\mathbb Z/p\mathbb Z)
\otimes_{\mathbb F_p}k$. This proves Proposition~2.2.5.



}{E7II-XXII-005}\SGApage{II}{410}
\section*{3. The congruence formula: statements and elementary equivalences}

\SGAvar{I wish to thank Grothendieck and Deligne for explaining to me the variants
(6.1.2), (3.1.1.2), and (3.1.3) of (3.1.1).

}{I wish to thank Grothendieck and Deligne for explaining to me the variants
(3.1.2), (3.2.1)--(3.2.2), and (3.3.1) of (3.1.1).

}{E7II-XXII-006}\subsection*{3.0.}
Let $X$ be a scheme of finite type over $\mathbb F_q$. The degree of a closed
point $y\in X$ is defined to be the degree of its residue field
$\mathbb F_q(y)/\mathbb F_q$. The zeta function of $X/\mathbb F_q$, as an
element of $\mathbb Z[[t]]$, is defined by the formula (cf. also SGA~5,
XIII~3.1)

\begin{equation}\tag{3.0.1}
Z(t,X/\mathbb F_q)=\prod_y(1-t^{\deg(y)})^{-1}.
\end{equation}

By the ``zeta function modulo $p$'' we shall always mean the image of
$Z(t,X/\mathbb F_q)$ in $\mathbb F_q[[t]]$ under the ``reduction modulo $p\mathbb Z[[t]]$''
map

\[
\mathbb Z[[t]]\longrightarrow\mathbb F_q[[t]].
\]

\subsubsection*{Theorem 3.1.}
Let $X$ be a proper scheme over $\mathbb F_q$. Then

\begin{equation}\tag{3.1.1}
Z(t,X/\mathbb F_q)
\equiv
\prod_{i\geq0}
\det\!\left(1-t\mathfrak F_q\mid H^i(X,\mathcal O_X)\right)^{(-1)^{i+1}}
\pmod p,
\end{equation}

or, equivalently by Proposition~2.2.4,

\begin{equation}\tag{3.1.2}
Z(t,X/\mathbb F_q)
\equiv
\prod_{i\geq0}
\det\!\left(1-t\varphi_q^{-1}\mid
H^i_{\mathrm{et}}(\bar X,\mathbb Z/p\mathbb Z)\right)^{(-1)^{i+1}}
\pmod p.
\end{equation}

\subsubsection*{Corollary 3.2.}
Let $X$ be a proper scheme over $\mathbb F_q$. Then

\SGAvar{\begin{equation}\tag{3.2.1}
\sum_{i\geq0}(-1)^i
\operatorname{Tr}\!\left(\mathfrak F_q\mid H^i(X,\mathcal O_X)\right.
\equiv \#X(\mathbb F_q)\pmod p,
\end{equation}

}{\begin{equation}\tag{3.2.1}
\sum_{i\geq0}(-1)^i
\operatorname{Tr}\!\left(\mathfrak F_q\mid H^i(X,\mathcal O_X)\right)
\equiv \#X(\mathbb F_q)\pmod p,
\end{equation}

}{E7II-XXII-007}or, equivalently by Proposition~2.2.4,

\SGApage{II}{\(\approx\)411}
\begin{equation}\tag{3.2.2}
\sum_{i\geq0}(-1)^i
\operatorname{Tr}\!\left(\varphi_q^{-1}\mid
H^i_{\mathrm{et}}(\bar X,\mathbb Z/p\mathbb Z)\right)
\equiv \#X(\mathbb F_q)\pmod p.
\end{equation}

Indeed, take the coefficients of $t$ in the congruences (3.1.1) and (3.1.2).

\subsubsection*{Theorem 3.3.}
Let $X$ be a proper scheme over $\mathbb F_q$, and denote by
$H^i_{\mathrm{DR}}(X/\mathbb F_q)$ the $\mathbb F_q$-vector spaces
$\mathbb H^i(X,\Omega^\bullet_{X/\mathbb F_q})$, on which $\mathfrak F_q$ acts. Then

\begin{equation}\tag{3.3.1}
Z(t,X/\mathbb F_q)
\equiv
\prod_{i\geq0}
\det\!\left(1-t\mathfrak F_q\mid
H^i_{\mathrm{DR}}(X/\mathbb F_q)\right)^{(-1)^{i+1}}
\pmod p.
\end{equation}

Indeed, consider the spectral sequence

\begin{equation}\tag{3.3.2}
E_1^{p,q}=H^q(X,\Omega^p_{X/\mathbb F_q})
\Longrightarrow H^{p+q}_{\mathrm{DR}}(X/\mathbb F_q).
\end{equation}

\SGAvar{The operator $\mathfrak F_q$ acts on it through the endomorphism
$\mathfrak F_q=\mathrm{Fr}_p^{X*}$ induced by
$\mathrm{Fr}_q^X:X\longrightarrow X$, and kills $E_1^{p,q}$ for $q\ne0$.
It follows that (3.3.1) is equivalent to (3.1.1).

}{The operator $\mathfrak F_q$ acts on it through the endomorphism
$\mathfrak F_q=(\mathrm{Fr}_q^X)^*$ induced by
$\mathrm{Fr}_q^X:X\longrightarrow X$, and kills $E_1^{p,q}$ for $p\ne0$.
It follows that (3.3.1) is equivalent to (3.1.1).

}{E7II-XXII-008}Let us add a small complement. Since formation of the ``semisimple part''
(cf. 1.0.7) is an exact functor, (3.3.2) gives a spectral sequence

\begin{equation}\tag{3.3.3}
(E_1^{p,q})_{\mathrm{ss}}
\Longrightarrow H^{p+q}_{\mathrm{DR}}(X/\mathbb F_q)_{\mathrm{ss}}.
\end{equation}

\SGAvar{Since $(E_1^{p,q})_{\mathrm{ss}}=0$ for $q\ne0$, we obtain for every $i$ an
isomorphism

}{Since $(E_1^{p,q})_{\mathrm{ss}}=0$ for $p\ne0$, we obtain for every $i$ an
isomorphism

}{E7II-XXII-009}\begin{equation}\tag{3.3.4}
H^i_{\mathrm{DR}}(X/\mathbb F_q)_{\mathrm{ss}}
\xrightarrow{\ \sim\ }
H^i(X,\mathcal O_X)_{\mathrm{ss}}
\end{equation}

that commutes with $\mathfrak F_q$. Hence

\begin{equation}\tag{3.3.5}
\det\!\left(1-t\mathfrak F_q\mid H^i_{\mathrm{DR}}(X/\mathbb F_q)\right)
=
\det\!\left(1-t\mathfrak F_q\mid H^i(X,\mathcal O_X)\right).
\end{equation}


\SGApage{II}{412}
\subsubsection*{Theorem 3.4.}
Let $X$ be a separated scheme of finite type over $\mathbb F_q$. Denoting by
$H_c^i$ étale cohomology with proper supports (SGA~4, XVII~5.1.9(iii)), one
has

\begin{equation}\tag{3.4.1}
Z(t,X/\mathbb F_q)
\equiv
\prod_{i\geq0}
\det\!\left(1-t\varphi_q^{-1}\mid
H_c^i(\bar X,\mathbb Z/p\mathbb Z)\right)^{(-1)^{i+1}}
\pmod p.
\end{equation}

We begin by proving the equivalence of (3.4.1) and (3.1.1). For the moment,
write

\begin{equation}\tag{3.4.2}
Z(X)\quad\text{in place of }Z(t,X/\mathbb F_q)
\end{equation}

\begin{equation}\tag{3.4.3}
Z'(X)\quad\text{in place of }\quad
\prod\det\!\left(1-t\varphi_q^{-1}\mid
H_c^i(\bar X,\mathbb Z/p\mathbb Z)\right)^{(-1)^{i+1}}.
\end{equation}

\subsubsection*{Lemma 3.5.}
If $\{U_i\}$ is a finite covering of $X$ by open subschemes, then

\begin{equation}\tag{3.5.1}
Z'(X)=
\prod_{j\geq0}
\prod_{1\leq i_0<\cdots<i_j\leq n}
Z'(U_{i_0}\cap\cdots\cap U_{i_j})^{(-1)^j}.
\end{equation}

\medskip
\noindent\emph{Proof.}
This follows immediately from the localization spectral sequence (SGA~4,
XVII~6.2.10)

\SGAvar{\[
E_1^{-a,b}=\bigoplus_{i_0<\cdots<i_b}
H_c^a(\bar U_{i_0}\cap\cdots\cap\bar U_{i_b},\mathbb Z/p\mathbb Z)
\Longrightarrow H_c^{-a+b}(\bar X,\mathbb Z/p\mathbb Z).
\]

}{\[
E_1^{-a,b}=\bigoplus_{i_0<\cdots<i_a}
H_c^b(\bar U_{i_0}\cap\cdots\cap\bar U_{i_a},\mathbb Z/p\mathbb Z)
\Longrightarrow H_c^{-a+b}(\bar X,\mathbb Z/p\mathbb Z).
\]

}{E7II-XXII-010}It may also be obtained by induction on $n$, using Lemma~3.6 below.

\subsubsection*{Lemma 3.6.}
If $U$ is an open subscheme of $X$ and $Y$ is the closed subscheme $X-U$, then

\begin{equation}\tag{3.6.1}
Z'(X)=Z'(U)Z'(Y).
\end{equation}

\medskip
\noindent\emph{Proof.}
This follows from the exact cohomology sequence

\begin{equation}\tag{3.6.2}
\longrightarrow H_c^i(\bar U,\mathbb Z/p\mathbb Z)
\longrightarrow H_c^i(\bar X,\mathbb Z/p\mathbb Z)
\longrightarrow H_c^i(\bar Y,\mathbb Z/p\mathbb Z)\longrightarrow.
\end{equation}

\SGApage{II}{413}
\subsubsection*{Lemma 3.7 (M\"obius).}
Let $B$ be a finite set and let $\{A_i\}$ be a finite covering by subsets.
Then

\begin{equation}\tag{3.7.1}
\operatorname{card}(B)=
\sum_{j\geq0}(-1)^j
\sum_{i_0<\cdots<i_j}
\operatorname{card}(A_{i_0}\cap\cdots\cap A_{i_j}).
\end{equation}

\medskip
\noindent\emph{Proof.}
There is a Leray spectral sequence for the rational cohomology of the discrete
space $B$ and the open covering $\{A_i\}$,

\SGAvar{\[
E_1^{p,q}=\bigoplus_q
H^p(A_{i_0}\cap\cdots\cap A_{i_q},\mathbb Q)
\Longrightarrow H^{p+q}(B,\mathbb Q),
\]

}{\[
E_1^{p,q}=\bigoplus_{i_0<\cdots<i_q}
H^p(A_{i_0}\cap\cdots\cap A_{i_q},\mathbb Q)
\Longrightarrow H^{p+q}(B,\mathbb Q),
\]

}{E7II-XXII-011}and one computes the Euler--Poincar\'e characteristic of $B$.

\subsubsection*{Lemma 3.8.}
If $\{U_i\}$ is a finite covering of $X$ by subschemes, then

\begin{equation}\tag{3.8.1}
Z(X)=
\prod_{j\geq0}
\prod_{i_0<\cdots<i_j}
Z(U_{i_0}\cap\cdots\cap U_{i_j})^{(-1)^j}.
\end{equation}

\medskip
\noindent\emph{Proof.}
One has

\begin{equation}\tag{3.8.2}
\operatorname{card}(X(\mathbb F_{q^n}))
=\sum_{r\mid n}r.\operatorname{card}\{y\in X\mid\deg(y)=r\},
\end{equation}

which gives

\begin{equation}\tag{3.8.3}
Z(X)=\exp\!\left(
\sum_{n\geq1}\operatorname{card}(X(\mathbb F_{q^n}))\frac{t^n}{n}
\right).
\end{equation}

The conclusion follows from (3.7.1).

\subsubsection*{Lemma 3.9.}
Suppose that formula (3.4.1), or equivalently (3.1.2) and hence (3.1.1), holds
when $X$ is a projective hypersurface. Then it holds for every separated
scheme $X$ of finite type.

\SGApage{II}{\(\approx\)414}
\medskip
\noindent\emph{Proof.}
By Lemmas~3.8 and 3.6, (3.4.1) holds for the complement
$\mathbb P^n-H$ of a hypersurface. Since

\[
\bigcap(\mathbb P^n-H_i)=\mathbb P^n-\bigcup_iH_i
\]

\SGAvar{is again the complement of a hypersurface, the open covering
$\{\mathbb P^n-H_i\}$ of $\mathbb P^n-\bigcap H_i$ allows us to apply
Lemmas~3.8 and 3.5 and deduce (3.1.3) for
$\mathbb P^n-\bigcap H_i$. By Lemmas~3.8 and 3.6, (3.4.1) then holds for
$\bigcap H_i$, and hence for every projective $X$. Another application of
Lemmas~3.8 and 3.5 gives (3.1.3) for every quasi-projective $X$, in particular
for every affine $X$. Finally, if $X$ is separated and of finite type, choose
a finite covering by affine open subschemes and apply Lemmas~3.8 and 3.5.



}{is again the complement of a hypersurface, the open covering
$\{\mathbb P^n-H_i\}$ of $\mathbb P^n-\bigcap H_i$ allows us to apply
Lemmas~3.8 and 3.5 and deduce (3.4.1) for
$\mathbb P^n-\bigcap H_i$. By Lemmas~3.8 and 3.6, (3.4.1) then holds for
$\bigcap H_i$, and hence for every projective $X$. Another application of
Lemmas~3.8 and 3.5 gives (3.4.1) for every quasi-projective $X$, in particular
for every affine $X$. Finally, if $X$ is separated and of finite type, choose
a finite covering by affine open subschemes and apply Lemmas~3.8 and 3.5.



}{E7II-XXII-012}\section*{4. The calculation of the zeta function of a hypersurface by
Dwork's method}

\subsection*{4.0.}
Let

\begin{equation}\tag{4.0.1}
f(X)=f(X_1,\ldots,X_{n+2})\in\mathbb F_q[X]
\end{equation}

be homogeneous of degree $d$, and let its zero locus define the hypersurface
$V\subset\mathbb P^{n+1}$. Denote by $V^*$ the open subset of $V$ on which
$\prod_{i=1}^{n+2}X_i$ is invertible, and by $V^*_{\mathrm{aff}}$ the affine
cone over $V^*$.

Choose a nontrivial additive character

\begin{equation}\tag{4.0.2}
\chi:\mathbb F_{q^\nu}\longrightarrow\mu_p(C),
\end{equation}

where $C$ is a fixed algebraically closed field of characteristic zero (for
example $C=\mathbb C$). Thus

\SGAvar{\begin{equation}\tag{4.0.3}
\text{for }x\in\mathbb F_{q^\nu},\qquad
\sum_{y\in\mathbb F_{q^\nu}}\chi(xy)=
\begin{cases}
q^\nu,&\chi=0,\\
0,&\chi\ne0.
\end{cases}
\end{equation}

For $\chi=(\chi_1,\ldots,\chi_{n+2})\in(\mathbb F_{q^\nu}^*)^{n+2}$ we therefore
have

\begin{equation}\tag{4.0.4}
1+\sum_{X_0\in\mathbb F_{q^\nu}^*}\chi(x_0f(x))
=\sum_{x_0\in\mathbb F_{q^\nu}}\chi(x_0f(x))
=
\begin{cases}
q^\nu,&f(x)=0,\\
0,&\text{otherwise}.
\end{cases}
\end{equation}

}{\begin{equation}\tag{4.0.3}
\text{for }x\in\mathbb F_{q^\nu},\qquad
\sum_{y\in\mathbb F_{q^\nu}}\chi(xy)=
\begin{cases}
q^\nu,&x=0,\\
0,&x\ne0.
\end{cases}
\end{equation}

For $x=(x_1,\ldots,x_{n+2})\in(\mathbb F_{q^\nu}^*)^{n+2}$ we therefore
have

\begin{equation}\tag{4.0.4}
1+\sum_{x_0\in\mathbb F_{q^\nu}^*}\chi(x_0f(x))
=\sum_{x_0\in\mathbb F_{q^\nu}}\chi(x_0f(x))
=
\begin{cases}
q^\nu,&f(x)=0,\\
0,&\text{otherwise}.
\end{cases}
\end{equation}

}{E7II-XXII-013}\SGApage{II}{\(\approx\)415}
Summing (4.0.4) over $x\in(\mathbb F_{q^\nu}^*)^{n+2}$ gives

\begin{equation}\tag{4.0.5}
\operatorname{card}V^*_{\mathrm{aff}}(\mathbb F_{q^\nu})
=\frac1{q^\nu}
\sum_{x\in(\mathbb F_{q^\nu}^*)^{n+2}}
\left\{1+\sum_{x_0\in\mathbb F_{q^\nu}^*}\chi(x_0f(x))\right\}.
\end{equation}

Since

\begin{equation}\tag{4.0.6}
\operatorname{card}V^*_{\mathrm{aff}}(\mathbb F_{q^\nu})
=(q^\nu-1)\operatorname{card}V^*(\mathbb F_{q^\nu}),
\end{equation}

formula (4.0.5) becomes

\begin{equation}\tag{4.0.7}
(q^\nu-1)\operatorname{card}V^*(\mathbb F_{q^\nu})
=\frac{(q^\nu-1)^{n+2}}{q^\nu}
+\frac1{q^\nu}
\sum_{(x_0,x)\in(\mathbb F_{q^\nu}^*)^{n+3}}\chi(x_0f(x)).
\end{equation}

Write $f(X)$ explicitly as a sum of monomials,

\begin{equation}\tag{4.0.8}
f(X)=\sum a_wX^w,
\end{equation}

where $w=(w_1,\ldots,w_{n+2})$, $\sum w_i=d$, and
$X^w\mathrel{\mathop:}=X_1^{w_1}\cdots X_{n+2}^{w_{n+2}}$. Then (4.0.7)
may be rewritten as

\begin{equation}\tag{4.0.9}
(q^\nu-1)\operatorname{card}V^*(\mathbb F_{q^\nu})
=\frac{(q^\nu-1)^{n+2}}{q^\nu}
+\frac1{q^\nu}
\sum_{(x_0,x)\in(\mathbb F_{q^\nu}^*)^{n+3}}
\prod_w\chi(a_wx_0x^w).
\end{equation}

\subsection*{4.1.}
We must now make explicit our ``additive character with values in
characteristic zero.'' One would like simply to take

\SGAvar{\begin{equation}\tag{4.1.0}
\chi:x\longmapsto
\exp\!\left(2\pi i\,\operatorname{tr}_{\mathbb F_{q^\nu}/\mathbb F_p}(x)\right.
=\xi_p^{\operatorname{tr}_{\mathbb F_{q^\nu}/\mathbb H_p}(x)}.
\end{equation}

}{\begin{equation}\tag{4.1.0}
\chi:x\longmapsto
\exp\!\left(2\pi i\,\operatorname{tr}_{\mathbb F_{q^\nu}/\mathbb F_p}(x)/p\right)
=\xi_p^{\operatorname{tr}_{\mathbb F_{q^\nu}/\mathbb F_p}(x)}.
\end{equation}

}{E7II-XXII-014}\SGApage{II}{\(\approx\)416}
Let

\begin{align}
\tag{4.1.1}
\Omega&=\text{an algebraic closure of }\mathbb Q_p,\\
\tag{4.1.2}
\operatorname{ord}:\Omega^*&\longrightarrow\mathbb Q
  &&\text{be the valuation normalized by }\operatorname{ord}(p)=1,\\
\tag{4.1.3}
\xi_p&\in\Omega
  &&\text{be a primitive $p$-th root of unity.}
\end{align}

One checks (cf. [2]) that there exists

\begin{equation}\tag{4.1.4}
\pi\in\mathbb Q_p(\xi_p),\qquad
\operatorname{ord}(\pi)=\frac1{p-1},
\end{equation}

such that

\begin{equation}\tag{4.1.5}
\sum_{n\geq0}\frac{\pi^{p^n}}{p^n}=0.
\end{equation}

\SGAvar{(Notice also that the left-hand side converges in $Q_p(\xi_p)$.)

}{(Notice also that the left-hand side converges in $\mathbb Q_p(\xi_p)$.)

}{E7II-XXII-015}Define \underline{the Artin--Hasse exponential}

\SGAvar{\begin{equation}\tag{4.1.6}
E(X)=\exp\!\left(\sum_{n\geq0}\frac{X^{p^n}}{p^n}\right)
\in Q_p[[X]],
\end{equation}

}{\begin{equation}\tag{4.1.6}
E(X)=\exp\!\left(\sum_{n\geq0}\frac{X^{p^n}}{p^n}\right)
\in\mathbb Q_p[[X]],
\end{equation}

}{E7II-XXII-016}which, miraculously, lies in $\mathbb Z_p[[X]]$. Define

\begin{equation}\tag{4.1.7}
\Theta(X)=E(\pi X).
\end{equation}

Its fundamental properties are summarized in the following lemma.

\subsubsection*{Lemma 4.2.}

\SGAvar{\begin{equation}\tag{4.2.1}
\Theta(X)=\sum B_n X^n,\qquad B_0=1,\quad A_1=\pi,\quad
\operatorname{ord}(B_n)\geq\frac n{p-1}.
\end{equation}

}{\begin{equation}\tag{4.2.1}
\Theta(X)=\sum B_n X^n,\qquad B_0=1,\quad B_1=\pi,\quad
\operatorname{ord}(B_n)\geq\frac n{p-1}.
\end{equation}

}{E7II-XXII-017}\begin{equation}\tag{4.2.2}
\Theta(1)=\xi_p,\qquad\text{a primitive $p$-th root of unity.}
\end{equation}

Let $t\in\Omega$, and suppose that, for a given integer $a>0$,

\begin{equation}\tag{4.2.3}
t=t^{p^a};
\end{equation}

\SGAvar{that is, $t$ is a Teichmüller representative in
$\mathbb W(\mathbb F_{q^a})\subset\Omega$. Then
$\sum_{j=0}^{a-1}t^{p^j}\in\mathbb Z_p$, so the formal series

}{that is, $t$ is a Teichmüller representative in
$\mathbb W(\mathbb F_{p^a})\subset\Omega$. Then
$\sum_{j=0}^{a-1}t^{p^j}\in\mathbb Z_p$, so the formal series

}{E7II-XXII-018}\[
(\Theta(X))^{\sum_{j=0}^{a-1}t^{p^j}}
\]

\SGApage{II}{\(\approx\)417}
is defined, and one has the identity of formal power series

\SGAvar{\begin{equation}\tag{4.2.3.1}
\prod_{j=0}^{a-1}\Theta(t^{p^j}X)
=(\Theta(X))^{\sum_{j=0}^{a-1}t^{j}}.
\end{equation}


}{\begin{equation}\tag{4.2.3.1}
\prod_{j=0}^{a-1}\Theta(t^{p^j}X)
=(\Theta(X))^{\sum_{j=0}^{a-1}t^{p^j}}.
\end{equation}


}{E7II-XXII-019}\medskip
\noindent\emph{Proof.}
Since $E(X)\in\mathbb Z_p[[X]]$, (4.2.1) follows from the definition. For
(4.2.2), we have

\begin{equation}\tag{4.2.4}
\log(\Theta(1))=0
\qquad\text{by the choice of $\pi$ in (4.1.5),}
\end{equation}

and

\begin{equation}\tag{4.2.5}
\operatorname{ord}(\Theta(1)-1)=\frac1{p-1}
\qquad\text{by (4.2.1).}
\end{equation}

By (4.2.4), there is an integer $n\geq0$ such that
$\Theta(1)^{p^n}=1$ (this is the characterization of the kernel of the
logarithm), and (4.2.5) implies that $n=1$.

\SGAvar{As for (4.2.3.1), after taking logarithms it follows by an immediate
calculation, once one knows that
$\sum_{j=0}^{a-1}t^{p^j}\in\mathbb Z_p$. But the $t^{p^j}$,
$0\leq j\leq a-1$, are precisely the conjugates over $\mathbb Z_p$ of
$t\in\mathbb W(\mathbb F_{q^a})$; hence

\[
\sum_{j=0}^{a-1}t^{p^j}
=\operatorname{Tr}_{\mathbb W(\mathbb F_{q^a})/\mathbb Z_p}(t)\in\mathbb Z_p.
\]

}{As for (4.2.3.1), after taking logarithms it follows by an immediate
calculation, once one knows that
$\sum_{j=0}^{a-1}t^{p^j}\in\mathbb Z_p$. But the $t^{p^j}$,
$0\leq j\leq a-1$, are precisely the conjugates over $\mathbb Z_p$ of
$t\in\mathbb W(\mathbb F_{p^a})$; hence

\[
\sum_{j=0}^{a-1}t^{p^j}
=\operatorname{Tr}_{\mathbb W(\mathbb F_{p^a})/\mathbb Z_p}(t)\in\mathbb Z_p.
\]

}{E7II-XXII-020}Evaluating the two sides of (4.2.3.1) at $X=1$ gives the following result.

\subsubsection*{Corollary 4.3.}
Let $t\in\Omega$ satisfy $t^{p^a}=t$, and denote its residue class by
$\bar t\in\mathbb F_{p^a}$. Then

\begin{equation}\tag{4.3.1}
\prod_{j=0}^{a-1}\Theta(t^{p^j})
=\xi_p^{\operatorname{tr}_{\mathbb F_{p^a}/\mathbb F_p}(\bar t)}.
\end{equation}



\SGApage{II}{\(\approx\)418}
\subsection*{4.4.}
Recall that the residue field of the ring of integers of
$\mathbb Q_p(\xi_p)$ is $\mathbb F_p$. Let

\begin{equation}\tag{4.4.0}
K=\text{the unramified extension of $\mathbb Q_p(\xi_p)$ with residue
field $\mathbb F_q$},
\end{equation}

and let

\SGAvar{\begin{equation}\tag{4.4.1}
\tau\in\operatorname{Gal}(K/Q_p(\xi_p))
\end{equation}

}{\begin{equation}\tag{4.4.1}
\tau\in\operatorname{Gal}(K/\mathbb Q_p(\xi_p))
\end{equation}

}{E7II-XXII-021}be the ``Frobenius'', that is, the automorphism reducing to the Frobenius
automorphism of $\mathbb F_q/\mathbb F_p$, $x\longmapsto x^p$.

Let

\begin{equation}\tag{4.4.2}
F(X)=\sum A_WX^W\in K[X_1,\ldots,X_{n+2}]
\end{equation}

\SGAvar{be the lifting of $f(X)$ whose coefficients satisfy $A_W^q=A_W$.
Putting (4.3.1) and (4.0.9) together, and taking into account the
definition (4.1.0) of $X$, one obtains

}{be the lifting of $f(X)$ whose coefficients satisfy $A_W^q=A_W$.
Putting (4.3.1) and (4.0.9) together, and taking into account the
definition (4.1.0) of $\chi$, one obtains

}{E7II-XXII-022}\begin{equation}\tag{4.4.3}
\left\{
\begin{aligned}
&(q^\nu-1)\operatorname{card}V^*(\mathbb F_{q^\nu})
-\frac{(q^\nu-1)^{n+2}}{q^\nu}\\
&\quad=\frac1{q^\nu}
\sum_{(X_0,X)\in(\mu_{q^\nu-1}(\Omega))^{n+3}}
\prod_W\prod_{j=0}^{\operatorname{ord}(q^\nu)-1}
\Theta\!\left(A_W^{p^j}X_0^{p^j}X^{p^jW}\right).
\end{aligned}
\right.
\end{equation}

Define

\[
H(X)=H(X_0,\ldots,X_{n+2})\in K[[X_0,\ldots,X_{n+2}]]
\]

by

\begin{equation}\tag{4.4.4}
H(X)=\prod_W
\Theta\!\left(A_WX_0X_1^{W_1}\cdots X_{n+2}^{W_{n+2}}\right).
\end{equation}

We let $\operatorname{Gal}(K/\mathbb Q_p(\xi_p))$ act on
$K[[X_0,\ldots,X_{n+2}]]$ through its action on the coefficients; the
componentwise action of
$\sigma\in\operatorname{Gal}(K/\mathbb Q_p(\xi_p))$ will be denoted

\begin{equation}\tag{4.4.5}
G(X)\longmapsto G^\sigma(X).
\end{equation}

\SGApage{II}{\(\approx\)419}
In particular,

\SGAvar{\[
H^{\tau^j}(X)=\prod_W
\Theta\!\left(A_W^{\tau^j}X_0X_1^{W}\cdots
X_{n+2}^{W_{n+2}}\right).
\]

}{\[
H^{\tau^j}(X)=\prod_W
\Theta\!\left(A_W^{\tau^j}X_0X_1^{W_1}\cdots
X_{n+2}^{W_{n+2}}\right).
\]

}{E7II-XXII-023}Since $A_W^{\tau^j}=A_W^{p^j}$, formula (4.4.3) may therefore be rewritten
as

\begin{equation}\tag{4.4.6}
\begin{aligned}
(q^\nu-1)\operatorname{Card}V^*(\mathbb F_{q^\nu})
&=\frac{(q^\nu-1)^{n+2}}{q^\nu}\\
&\quad+\frac1{q^\nu}
\sum_{(X_0,X)\in(\mu_{q^\nu-1}(\Omega))^{n+3}}
\prod_{j=0}^{\operatorname{ord}(q^\nu)-1}H^{\tau^j}(X^{p^j}).
\end{aligned}
\end{equation}


\subsection*{4.5.}
We shall now express the sum in (4.4.6) as a trace.

For each $b\in\mathbb Q$, define a $p$-adic Banach $K$-algebra

\begin{equation}\tag{4.5.1}
L(b)=\text{the subring of $K[[X_0,\ldots,X_{n+2}]]$ consisting of the
series}
\end{equation}

which can be written

\begin{equation}\tag{4.5.2}
\sum_{dU_0=\sum_{i\geq1}U_i}
B_UX_0^{U_0}X_1^{U_1}\cdots X_{n+2}^{U_{n+2}}
\qquad (U=(U_0,\ldots,U_{n+2}))
\end{equation}

and which satisfy

\begin{equation}\tag{4.5.3}
\inf_U\bigl(\operatorname{ord}(B_U)-bU_0\bigr)>-\infty.
\end{equation}

It is a $p$-adic Banach algebra for the order function (the negative of the
logarithm of the norm)

\begin{equation}\tag{4.5.4}
\inf_U\bigl(\operatorname{ord}(B_U)-bU_0\bigr)
=\operatorname{ord}\!\left(\sum B_UX^U\right).
\end{equation}

It follows from (4.4.4) and (4.2.1) that

\[
H(X)\in L\!\left(\frac1{p-1}\right).
\]

\SGApage{II}{\(\approx\)420}
Let us introduce some operators. For every integer $k>0$, define a
continuous operator between Banach spaces over $K$ by

\begin{equation}\tag{4.5.5}
\left\{
\begin{aligned}
\psi_k:L(b)&\longrightarrow L(kb),\\
\psi_k\!\left(\sum B_UX^U\right)&=\sum B_{kU}X^U.
\end{aligned}
\right.
\end{equation}

For every pair $(b,b')$ with $b'<b$, there is a ``restriction'' operator,
which is completely continuous (that is, a uniform limit of finite-rank
operators):

\begin{equation}\tag{4.5.6}
\left\{
\begin{aligned}
L(b)&\xrightarrow{\ \text{restriction}\ }L(b'),\\
\sum B_UX^U&\longmapsto\sum B_UX^U.
\end{aligned}
\right.
\end{equation}

Each element $G$ of the algebra $L(b)$ gives the continuous endomorphism
of multiplication by $G$,

\begin{equation}\tag{4.5.7}
\begin{aligned}
G:L(b)&\longrightarrow L(b)\\
\eta&\longmapsto\eta G.
\end{aligned}
\end{equation}

These operators are related by the following lemma, whose proof is found in
[2] (see also [4]).


\SGApage{II}{421}
\subsubsection*{Lemma 4.5.8 (Dwork trace lemma).}
Let $b\in\mathbb Q$, $b>0$, let $k\in\mathbb Z$, $k\geq2$, and let
$G\in L(b/k)$. Then the composite

\[
L(b)\xrightarrow{\ \text{restriction}\ }L(b/k)
\xrightarrow{\ \times G\ }L(b/k)
\xrightarrow{\ \psi_k\ }L(b),
\]

denoted by $\psi_k\circ G$, is completely continuous, and its trace
(cf. [4]) is given by

\begin{equation}\tag{4.5.9}
\operatorname{tr}(\psi_k\circ G)
=\frac1{(k-1)^{n+2}}
\sum_{X\in(\mu_{k-1}(\Omega))^{n+3}}G(X).
\end{equation}

Apply this lemma with

\[
\left\{
\begin{aligned}
k&=q^\nu,\\
G(X)&=\prod_{j=0}^{\operatorname{ord}(q^\nu)-1}
H^{\tau^j}(X^{p^j})
\in L\!\left(\frac{p}{(p-1)q^\nu}\right),\\
b&=\frac1{p-1}.
\end{aligned}
\right.
\]

Formula (4.4.6) becomes

\SGAvar{\begin{equation}\tag{4.5.10}
\begin{aligned}
(q-1)\operatorname{card}V^*(\mathbb F_{q^\nu})
&=\frac{(q^\nu-1)^{n+2}}{q^\nu}\\
&\quad+\frac{(q^\nu-1)^{n+3}}{q^\nu}
\operatorname{tr}\!\left(
\left.\psi_{q}
\prod_{j=0}^{\operatorname{ord}(q^\nu)-1}H^{\tau^j}(X^{p^j})
\right|L\!\left(\frac1{p-1}\right)\right).
\end{aligned}
\end{equation}

Let $L^0(b)$ be the ideal of the elements of $L(b)$ which ``vanish at
zero'', that is, which have no constant term. It is stable under
$\psi\circ G$, and evidently

}{\begin{equation}\tag{4.5.10}
\begin{aligned}
(q^\nu-1)\operatorname{card}V^*(\mathbb F_{q^\nu})
&=\frac{(q^\nu-1)^{n+2}}{q^\nu}\\
&\quad+\frac{(q^\nu-1)^{n+3}}{q^\nu}
\operatorname{tr}\!\left(
\left.\psi_{q^\nu}
\prod_{j=0}^{\operatorname{ord}(q^\nu)-1}H^{\tau^j}(X^{p^j})
\right|L\!\left(\frac1{p-1}\right)\right).
\end{aligned}
\end{equation}

Let $L^0(b)$ be the ideal of the elements of $L(b)$ which ``vanish at
zero'', that is, which have no constant term. It is stable under
$\psi_k\circ G$, and evidently

}{E7II-XXII-024}\begin{equation}\tag{4.5.11}
\operatorname{tr}(\psi_k\circ G)
=G(0)+\operatorname{tr}\!\left(\left.\psi_k\circ G\right|L^0\right).
\end{equation}

Since $H(0)=1$, formula (4.5.10) may be rewritten as

\SGAvar{\begin{equation}\tag{4.5.12}
\begin{aligned}
\operatorname{card}V^*(\mathbb F_{q^\nu})
&=(q^\nu-1)^{n+1}\\
&\quad+\frac{(q^\nu-1)^{n+2}}{q}
\operatorname{tr}\!\left(
\left.\psi_{q}\circ
\prod_{j=0}^{\operatorname{ord}(q^\nu)-1}H^{\tau^j}(X^{p^j})
\right|L^0\!\left(\frac1{p-1}\right)\right).
\end{aligned}
\end{equation}


}{\begin{equation}\tag{4.5.12}
\begin{aligned}
\operatorname{card}V^*(\mathbb F_{q^\nu})
&=(q^\nu-1)^{n+1}\\
&\quad+\frac{(q^\nu-1)^{n+2}}{q^\nu}
\operatorname{tr}\!\left(
\left.\psi_{q^\nu}\circ
\prod_{j=0}^{\operatorname{ord}(q^\nu)-1}H^{\tau^j}(X^{p^j})
\right|L^0\!\left(\frac1{p-1}\right)\right).
\end{aligned}
\end{equation}


}{E7II-XXII-025}\SGApage{II}{\(\approx\)422}
Consider now the $\tau^{-1}$-linear operator

\begin{equation}\tag{4.5.13}
\tau^{-1}\circ\psi_p\circ H(X):L(b)\longrightarrow L(b).
\end{equation}

Its iterates are

\SGAvar{\begin{equation}\tag{5.4.14}
(\tau^{-1}\circ\psi_p\circ H(X))^a
=\tau^{-a}\circ\psi_{p^a}\circ
H(X)H^\tau(X^p)\cdots H^{\tau^{a-1}}(X^{p^{a-1}}),
\end{equation}

}{\begin{equation}\tag{4.5.14}
(\tau^{-1}\circ\psi_p\circ H(X))^a
=\tau^{-a}\circ\psi_{p^a}\circ
H(X)H^\tau(X^p)\cdots H^{\tau^{a-1}}(X^{p^{a-1}}),
\end{equation}

}{E7II-XXII-026}and, if $\operatorname{ord}(q)\mid a$, then $\tau^{-a}=\operatorname{id}$.
Define the $\tau^{-1}$-linear operation

\begin{equation}\tag{4.5.15}
\left\{
\begin{aligned}
\alpha:L^0\!\left(\frac1{p-1}\right)
&\longrightarrow L^0\!\left(\frac1{p-1}\right),\\
\alpha&=\frac1p\,\tau^{-1}\circ\psi_p\circ H(X).
\end{aligned}
\right.
\end{equation}

Then (4.5.12) takes the form

\SGAvar{\begin{equation}\tag{4.5.16}
\operatorname{card}V^*(\mathbb F_{q^\nu})
=(q^\nu-1)^{n+1}
+(q^\nu-1)^{n+2}
\operatorname{tr}\!\left(
\left.(\alpha)^{\operatorname{ord}(q^\nu)}
\right|L^0\!\left(\frac1{p-1}\right)\right..
\end{equation}

}{\begin{equation}\tag{4.5.16}
\operatorname{card}V^*(\mathbb F_{q^\nu})
=(q^\nu-1)^{n+1}
+(q^\nu-1)^{n+2}
\operatorname{tr}\!\left(
\left.(\alpha)^{\operatorname{ord}(q^\nu)}
\right|L^0\!\left(\frac1{p-1}\right)\right).
\end{equation}

}{E7II-XXII-027}Put

\begin{equation}\tag{4.5.17}
S=\{1,\ldots,n+2\}.
\end{equation}

For every nonempty subset $A\subset S$, define the subspace

\begin{equation}\tag{4.5.18}
\begin{aligned}
L^A(b)={}&\text{the set of series $\sum B_UX^U\in L(b)$ such that, for every
monomial $X^U$ with $B_U\ne0$,}\\
&\text{the indices $i$ for which $U_i>0$ are exactly $0$ and the
$j\in A$}.
\end{aligned}
\end{equation}

Then

\begin{equation}\tag{4.5.19}
L^0(b)=\bigoplus_{\substack{A\subset S\\A\ne\varnothing}}L^A(b).
\end{equation}



\SGApage{II}{\(\approx\)423}
\par\medskip
And, for every nonempty subset $B\subset S$, one has

\begin{equation}\tag{4.5.20}
\bigoplus_{A\supset B}L^A(b)
=L(b)\cap\left(\prod_{i\in B}X_i\right)K[[X]].
\end{equation}

\SGAvar{Clearly every subspace $\bigoplus_{A\subseteq B}L^A(b)$ is stable under
$\alpha$. It follows that the same is true of
$\bigoplus_{A\supsetneq B}L^A(b)$, which is a sum of subspaces of the
preceding type. Define

}{Clearly every subspace $\bigoplus_{A\supset B}L^A(b)$ is stable under
$\alpha$. It follows that the same is true of
$\bigoplus_{A\supsetneq B}L^A(b)$, which is a sum of subspaces of the
preceding type. Define

}{E7II-XXII-028}\begin{equation}\tag{4.5.21}
\alpha_B=\text{the $\tau^{-1}$-linear endomorphism induced by $\alpha$ on}
\quad
\frac{\displaystyle\bigoplus_{A\supset B}
L^A\!\left(\frac1{p-1}\right)}
{\displaystyle\bigoplus_{A\supsetneq B}
L^A\!\left(\frac1{p-1}\right)}
\simeq L^B\!\left(\frac1{p-1}\right).
\end{equation}

Explicitly, $\alpha_B$ is given by

\begin{equation}\tag{4.5.22}
\alpha_B=\frac{\tau^{-1}}p\,\psi_p\,H_B(X),
\end{equation}

an operation under which $L^B$ is stable, where $H_B(X)$ is obtained from
$H(X)$ by substituting $X_i=0$ for $i\notin\{0\}\cup B$.

Rewriting (4.5.16) once more gives

\SGAvar{\begin{equation}\tag{4.5.2.3}
\operatorname{card}V^*(\mathbb F_{q^\nu})
=(q^\nu-1)^{n+1}
+(q^\nu-1)^{n+2}
\sum_{\substack{A\subset S\\A\ne\varnothing}}
\operatorname{tr}\!\left((\alpha_A)^{\operatorname{ord}(q^\nu)}\right).
\end{equation}

}{\begin{equation}\tag{4.5.23}
\operatorname{card}V^*(\mathbb F_{q^\nu})
=(q^\nu-1)^{n+1}
+(q^\nu-1)^{n+2}
\sum_{\substack{A\subset S\\A\ne\varnothing}}
\operatorname{tr}\!\left((\alpha_A)^{\operatorname{ord}(q^\nu)}\right).
\end{equation}

}{E7II-XXII-029}For $A\subset S$, $A\ne\varnothing$, define the hypersurface

\begin{equation}\tag{4.5.24}
V_A\subset\mathbb P^{\operatorname{card}(A)-1}
\simeq
\mathbb P^{n+1}\cap
\{X_i=0\text{ for every }i\notin A\}
\end{equation}

by $f_A(X)$, the polynomial obtained from $f(X)$ by substituting $X_i=0$
for $i\notin A$.

\SGApage{II}{\(\approx\)424}
By compatibility of the formation of the series $H$ from $f(X)$ with
substitution of $0$ for some of the variables, (4.5.23) gives, for every
$A\subseteq S$, $A\ne\varnothing$,

\SGAvar{\begin{equation}\tag{4.5.25}
\operatorname{card}(V_A^*(\mathbb F_{q^\nu}))
=(q^\nu-1)^{\operatorname{car}(A)-1}
+(q^\nu-1)^{\operatorname{card}(A)}
\sum_{B\subset A}
\operatorname{tr}\!\left((\alpha_B))^{\operatorname{ord}(q^\nu)}\right).
\end{equation}

}{\begin{equation}\tag{4.5.25}
\operatorname{card}(V_A^*(\mathbb F_{q^\nu}))
=(q^\nu-1)^{\operatorname{card}(A)-1}
+(q^\nu-1)^{\operatorname{card}(A)}
\sum_{B\subset A}
\operatorname{tr}\!\left((\alpha_B)^{\operatorname{ord}(q^\nu)}\right).
\end{equation}

}{E7II-XXII-030}On the other hand,

\begin{equation}\tag{4.5.26}
\operatorname{card}V(\mathbb F_{q^\nu})
=\sum_{\substack{A\subset S\\A\ne\varnothing}}
\operatorname{card}V_A^*(\mathbb F_{q^\nu}).
\end{equation}

After an easy summation, (4.5.25) therefore gives

\begin{equation}\tag{4.5.27}
\begin{aligned}
\operatorname{card}V(\mathbb F_{q^\nu})
={}&1+q^\nu+q^{2\nu}+\cdots
+q^{\nu(\operatorname{card}(S)-1)}\\
&+\sum_{\substack{B\subseteq S\\B\ne\varnothing}}
(q^\nu-1)^{\operatorname{card}(B)}
q^{\nu\operatorname{card}(S-B)}
\operatorname{tr}\!\left((\alpha_B)^{\operatorname{ord}(q^\nu)}\right).
\end{aligned}
\end{equation}

Since $\alpha$, and hence the $\alpha_B$, are $\tau^{-1}$-linear, the
operators $\alpha_B^{\operatorname{ord}(q)}$ are $K$-linear. Define the
characteristic series

\[
\Delta_B(t)\in1+tK[[t]]
\]

by

\begin{equation}\tag{4.5.28}
\Delta_B(t)=\det\!\left(
1-t\alpha_B^{\operatorname{ord}(q)}
\mathrel{\Big|}L^B\!\left(\frac1{p-1}\right)\right).
\end{equation}

Introduce the following endomorphisms of $1+tK[[t]]$:

\begin{equation}\tag{4.5.29}
\varphi:g(t)\longmapsto g^\varphi(t)=g(qt)
\end{equation}

and

\begin{equation}\tag{4.5.30}
\delta=1-\varphi:g(t)\longmapsto
g^\delta(t)=\frac{g(t)}{g(qt)}.
\end{equation}


\SGApage{II}{\(\approx\)425}
Taking account of the identities

\begin{equation}\tag{4.5.31}
Z(t,V/\mathbb F_q)
=\exp\!\left(
\sum_{\nu\geq1}
\operatorname{card}(V(\mathbb F_{q^\nu}))\frac{t^\nu}{\nu}
\right)
\end{equation}

and

\begin{equation}\tag{4.5.32}
\det\!\left(1-t\alpha_B^{\operatorname{ord}(q)}\right)
=\exp\!\left(
-\sum_{\nu\geq1}
\operatorname{tr}\!\left((\alpha_B)^{\nu\operatorname{ord}q}\right)
\frac{t^\nu}{\nu}
\right),
\end{equation}

formula (4.5.27) can be rewritten as

\SGAvar{\begin{equation}\tag{4.5.23}
Z(t,V/\mathbb F_q)\prod_{i=0}^{n+1}(1-q^it)
=\prod_{j=1}^{n+2}
\left(
\prod_{\substack{B\subset S\\\operatorname{card}(B)=j}}
\Delta_B(t)
\right)^{-(-\delta)^j\varphi^{n+2-j}}.
\end{equation}



}{\begin{equation}\tag{4.5.33}
Z(t,V/\mathbb F_q)\prod_{i=0}^{n+1}(1-q^it)
=\prod_{j=1}^{n+2}
\left(
\prod_{\substack{B\subset S\\\operatorname{card}(B)=j}}
\Delta_B(t)
\right)^{-(-\delta)^j\varphi^{n+2-j}}.
\end{equation}



}{E7II-XXII-031}\subsection*{4.6.}
We now study the integrality of the series $\Delta_B$. Put

\begin{equation}\tag{4.6.1}
\mathcal O=\text{the ring of integers of $K$},
\end{equation}

and define a \underline{filtration} on $1+t\mathcal O[[t]]$ by

\begin{equation}\tag{4.6.2}
F^i(1+t\mathcal O[[t]])
=1+q^it\mathcal O[[q^it]]
=\varphi^i(1+t\mathcal O[[t]]).
\end{equation}

One has

\begin{align}
\varphi(F^i)&\subset F^{i+1}, \tag{4.6.3}\\
\delta(F^i)&\subset F^i, \tag{4.6.4}\\
\delta&\text{ induces the identity on }F^i/F^{i+1}. \tag{4.6.5}
\end{align}

\subsubsection*{Lemma 4.6.6.}
One has $\Delta_B(t)\in1+t\mathcal O[[t]]$; moreover, if an integer $j$
satisfies

\begin{equation}\tag{4.6.7}
jd<\operatorname{card}(B),
\end{equation}

then

\begin{equation}\tag{4.6.8}
\Delta_B(t)\in F^j(1+t\mathcal O[[t]]).
\end{equation}



\SGApage{II}{\(\approx\)426}
\noindent\emph{Proof.}
The Banach spaces
$L^B\!\left(\frac1{p-1}\right)$ are all of type $B(I)$ in the sense of
Serre~[7]; that is, they admit orthonormal bases, namely the monomials

\begin{equation}\tag{4.6.9}
\pi^{U_0}X_0^{U_0}X_1^{U_1}\cdots X_{n+2}^{U_{n+2}}
\end{equation}

\noindent
which belong to them.  To say that this is an orthonormal basis means, by
definition, that

\begin{equation}\tag{4.6.10}
\operatorname{ord}\!\left(
\sum C_U\pi^{U_0}X_0^{U_0}\cdots X_{n+2}^{U_{n+2}}
\right)
=\min_U\operatorname{ord}(C_U).
\end{equation}

Let

\SGAvar{\begin{equation}\tag{4.6.11}
\mathbf R_B=\text{the $\mathcal O$-submodule of
$L^B\!\left(\frac1{p-1}\right)$ formed by the elements of ordinal $\geq0$}.
\end{equation}

}{\begin{equation}\tag{4.6.11}
\mathbf R_B=\text{the $\mathcal O$-submodule of
$L^B\!\left(\frac1{p-1}\right)$ formed by the elements of order $\geq0$}.
\end{equation}

}{E7II-XXII-032}To prove the lemma, it is enough to show

\begin{equation}\tag{4.6.12}
jd<\operatorname{card}(B)
\quad\Longrightarrow\quad
(\alpha_B)^{\operatorname{ord}(q)}R_B\subset q^jR_B,
\end{equation}

or, a fortiori, to show

\begin{equation}\tag{4.6.13}
jd<\operatorname{card}(B)
\quad\Longrightarrow\quad
\alpha_B(R_B)\subset p^jR_B.
\end{equation}

By (4.5.22),

\begin{equation}\tag{4.6.14}
\alpha_B=\frac{\tau^{-1}}p\,\psi_p\circ H_B(X).
\end{equation}

Since

\begin{equation}\tag{4.6.15}
H(X)\in L\!\left(\frac1{p-1}\right),
\qquad \operatorname{ord}(H(X))\geq0,
\end{equation}

we have

\begin{equation}\tag{4.6.16}
H_B(X)\cdot R_B\subset R_B
\end{equation}

and, plainly,

\begin{equation}\tag{4.6.17}
\tau^{-1}(R_B)=R_B.
\end{equation}

\SGApage{II}{\(\approx\)427}
Consequently, (4.6.13) follows from

\begin{equation}\tag{4.6.18}
jd<\operatorname{card}(B)
\quad\Longrightarrow\quad
\psi_p(R_B)\subset p^{j+1}R_B.
\end{equation}

We verify (4.6.18) monomial by monomial.  We have

\SGAvar{\begin{equation}\tag{4.6.19}
\begin{cases}
\psi_p\!\left(\pi^{U_0}X^{U_0}\cdots X_{n+2}^{U_{n+2}}\right)=0,
&\text{unless }p\mid U_i\text{ for every }i,\\
\text{and}\\[3pt]
\psi_p\!\left(
\pi^{pU_0}X_0^{pU_0}\cdots X_{n+2}^{pU_{n+2}}
\right)
=\pi^{(p-1)U_0}\cdot
\underbrace{\pi^{U_0}X_0^{U_0}\cdots X_{n+2}^{U_{n+2}}}_{\in R_B}.
\end{cases}
\end{equation}

}{\begin{equation}\tag{4.6.19}
\begin{cases}
\psi_p\!\left(\pi^{U_0}X_0^{U_0}\cdots X_{n+2}^{U_{n+2}}\right)=0,
&\text{unless }p\mid U_i\text{ for every }i,\\
\text{and}\\[3pt]
\psi_p\!\left(
\pi^{pU_0}X_0^{pU_0}\cdots X_{n+2}^{pU_{n+2}}
\right)
=\pi^{(p-1)U_0}\cdot
\underbrace{\pi^{U_0}X_0^{U_0}\cdots X_{n+2}^{U_{n+2}}}_{\in R_B}.
\end{cases}
\end{equation}

}{E7II-XXII-033}The conditions for
$X_0^{U_0}\cdots X_{n+2}^{U_{n+2}}$ to belong to
$L^B\!\left(\frac1{p-1}\right)$ are

\begin{equation}\tag{4.6.20}
\begin{cases}
dU_0=\displaystyle\sum_{i\geq1}U_i,\\
U_i>0\Longleftrightarrow i\in B,\qquad i\geq1.
\end{cases}
\end{equation}

In particular,

\begin{equation}\tag{4.6.21}
dU_0\geq\operatorname{card}(B).
\end{equation}

Hence

\begin{equation}\tag{4.6.22}
\operatorname{ord}\!\left(\pi^{(p-1)U_0}\right)
=U_0\geq\frac{\operatorname{card}(B)}d>j,
\end{equation}

and, since $U_0\in\mathbb Z$,

\begin{equation}\tag{4.6.23}
\operatorname{ord}\!\left(\pi^{(p-1)U_0}\right)
=U_0\geq j+1.
\end{equation}

Together with (4.6.19), this proves (4.6.18).



\SGApage{II}{428}
\section*{5. \underline{A theorem of Ax} [1] (see also [4], Section 7)}

\subsubsection*{Theorem 5.0.}
\underline{Let} $V$ \underline{be a hypersurface in} $\mathbb P^{n+1}$ \underline{of degree} $d>0$, \underline{defined
over} $\mathbb F_q$, \underline{and let} $j$ \underline{be an integer}.  \underline{If}

\begin{equation}\tag{5.0.1}
jd<n+2,
\end{equation}

\underline{then}

\begin{equation}\tag{5.0.2}
Z(t,V/\mathbb F_q)\prod_{i=0}^{n+1}(1-q^it)
\in1+q^jt\mathbb Z[[q^jt]],
\end{equation}

and the reciprocal of every zero or pole of

\begin{equation}\tag{5.0.3}
Z(t,V/\mathbb F_q)\prod_{i=0}^{n+1}(1-q^it)
\end{equation}

\noindent
\underline{is of the form} $q^j$ times \underline{an algebraic integer}.

\subsubsection*{Theorem 5.0.bis.}
\underline{Let} $V_{\mathrm{aff}}$ \underline{be a hypersurface in} $\mathbb A^{n+2}$ \underline{defined by a
homogeneous form of degree} $d$, \underline{over} $\mathbb F_q$.  \underline{If}

\SGAvar{\begin{equation}\tag{5.0.7}
jd<n+2,
\end{equation}

}{\begin{equation}\tag{5.0.4}
jd<n+2,
\end{equation}

}{E7II-XXII-034}\underline{then}

\begin{equation}\tag{5.0.5}
Z(t,V_{\mathrm{aff}}/\mathbb F_q)
\in1+q^jt\mathbb Z[[q^jt]].
\end{equation}

\underline{In particular},

\begin{equation}\tag{5.0.6}
\operatorname{card}\bigl(V_{\mathrm{aff}}(\mathbb F_q)\bigr)
\equiv0\mod{q^j},
\end{equation}

\underline{and consequently}

\begin{equation}\tag{5.0.6 bis}
\operatorname{card}\bigl(V(\mathbb F_q)\bigr)
\geq\frac{q^j-1}{q-1}>0,
\end{equation}

\noindent
\underline{which improves the Chevalley--Warning theorem}~[8].



\SGApage{II}{\(\approx\)429}
\SGAvar{\medskip
\noindent\emph{Proof.}
The implication (5.0.5) $\Rightarrow$ (5.0.6) is clear on taking the
coefficient of $t$ in the zeta function.  Assertions (5.0.4) and (5.0.2)
are equivalent, since

}{\medskip
\noindent\emph{Proof.}
The implication (5.0.5) $\Rightarrow$ (5.0.6) is clear on taking the
coefficient of $t$ in the zeta function.  Assertions (5.0.5) and (5.0.2)
are equivalent, since

}{E7II-XXII-035}\begin{equation}\tag{5.0.7}
\operatorname{card}V_{\mathrm{aff}}(\mathbb F_{q^\nu})
=1+(q^\nu-1)\operatorname{card}V(\mathbb F_{q^\nu}),
\end{equation}

whence

\SGAvar{\begin{equation}\tag{5.0.8}
Z(t,V_{\mathrm{aff}}/\mathbb F_q)
=\frac{Z(t,V/\mathbb F_q)}
{Z(qt,V/\mathbb F_q)(1-t)}
\end{equation}

}{\begin{equation}\tag{5.0.8}
Z(t,V_{\mathrm{aff}}/\mathbb F_q)
=\frac{Z(qt,V/\mathbb F_q)}
{Z(t,V/\mathbb F_q)(1-t)}
\end{equation}

}{E7II-XXII-036}and

\SGAvar{\begin{equation}\tag{5.0.9}
Z(t,V/\mathbb F_q)
=\prod_{i\geq0}
\left[(1-q^if)Z(q^it,V_{\mathrm{aff}}/\mathbb F_q)\right].
\end{equation}

}{\begin{equation}\tag{5.0.9}
Z(t,V/\mathbb F_q)
=\prod_{i\geq0}
\left[(1-q^it)Z(q^it,V_{\mathrm{aff}}/\mathbb F_q)\right]^{-1}.
\end{equation}

}{E7II-XXII-037}\noindent
Assertions (5.0.2) and (5.0.3) are likewise equivalent, by the theorem of
Fatou~[3] and the rationality of the zeta function.  Since

\begin{equation}\tag{5.0.10}
Z(t,V/\mathbb F_q)\in\mathbb Z[[t]],
\end{equation}

it is enough to prove

\begin{equation}\tag{5.0.11}
Z(t,V/\mathbb F_q)\prod_{i=0}^{n+1}(1-q^it)
\in F^j(1+t\mathcal O[[t]]).
\end{equation}

In view of (4.5.33) and (4.6.4)--(4.6.5), it is enough to show, for
$\varnothing\ne B\subset\{1,\ldots,n+2\}$, that

\SGAvar{\begin{equation}\tag{5.0.12}
\Lambda_B(t)\in F^{j+\operatorname{card}(B)-(n+2)}.
\end{equation}

We may finish the proof by applying Lemma~4.4.6 once we have proved

}{\begin{equation}\tag{5.0.12}
\Delta_B(t)\in F^{j+\operatorname{card}(B)-(n+2)}.
\end{equation}

We may finish the proof by applying Lemma~4.6.6 once we have proved

}{E7II-XXII-038}\begin{equation}\tag{5.0.13}
jd<n+2
\quad\Longrightarrow\quad
\bigl(j+\operatorname{card}(B)-(n+2)\bigr)d
<\operatorname{card}(B).
\end{equation}

This may be rewritten as

\begin{equation}\tag{5.0.14}
jd<n+2
\quad\Longrightarrow\quad
(j-(n+2))d<\operatorname{card}(B)(1-d).
\end{equation}

\SGApage{II}{\(\approx\)430}
But

\[
\inf_{\varnothing\ne B\subset\{1,\ldots,n+2\}}
\operatorname{card}(B)(1-d)
=(n+2)(1-d),
\]

so it is enough to prove

\begin{equation}\tag{5.0.15}
jd<n+2
\quad\Longrightarrow\quad
(j-(n+2))d<(n+2)(1-d),
\end{equation}

which is true.





\subsubsection*{Remark 5.1.}
Condition

\begin{equation}\tag{5.1.1}
jd<n+2
\end{equation}

admits a cohomological interpretation.  Let $V$ be a \underline{smooth} hypersurface of
degree $d$ and dimension $n$.  As in XI, 2.1, put, for $a+b=n$,

\begin{equation}\tag{5.1.2}
h_0^{a,b}
=\dim H^b(V,\underline{\Omega}_V^a)-\delta_{a,b}
=\dim E^{a,b}(V),
\end{equation}

where $E^{a,b}(V)$ is the bidegree-$(a,b)$ component of the vanishing Hodge
cohomology (compare XVII, 5.4.4)

\SGAvar{\begin{equation}\tag{5.1.3}
E^{a,b}(V)
=\operatorname{Coker}\!\left(
H^b(\mathbb P^{n+1},\underline{\Omega}_{\mathbb P^{n+1}}^a)
\hookrightarrow H^b(V,\underline{\Omega}_V^a)\right..
\end{equation}

}{\begin{equation}\tag{5.1.3}
E^{a,b}(V)
=\operatorname{Coker}\!\left(
H^b(\mathbb P^{n+1},\underline{\Omega}_{\mathbb P^{n+1}}^a)
\hookrightarrow H^b(V,\underline{\Omega}_V^a)\right).
\end{equation}

}{E7II-XXII-039}Then, by XI, 2.8,

\begin{equation}\tag{5.1.4}
jd<n+2
\quad\Longleftrightarrow\quad
\left[h_0^{a,n-a}=0\quad\text{for }a<j\text{ and }a>n-j\right].
\end{equation}

In other words, the condition says that the Hodge level (XI, 2.7) of the
vanishing de Rham cohomology of $V$,

\SGApage{II}{\(\approx\)431}
\[
E_{\mathrm{DR}}^n(V)
=\operatorname{Coker}\!\left(
H_{\mathrm{Hdg}}^n(\mathbb P^{n+1})
\hookrightarrow H_{\mathrm{Hdg}}^n(V)\right),
\]

is at most $n-2j$, or equivalently that its coniveau is at least $j$.
On the other hand, the conclusion also says that the eigenvalues of the
Frobenius $\mathfrak F_q$ acting on the vanishing $\ell$-adic cohomology
($\ell\ne p$) are algebraic integers \underline{divisible by} $q^j$ (XXI, A.2).
Thus Ax's result 5.0 supports the following conjecture, suggested by the
philosophy of the level of motives and relating the classical Hodge and
Tate conjectures~[5].

\subsubsection*{Conjecture 5.2.}
\underline{Let} $f:X\to S$ \underline{be a proper smooth morphism}, \underline{with} $S$ \underline{integral}, \underline{of finite
type over} $\operatorname{Spec}(\mathbb Z)$, \underline{and dominant over}
$\operatorname{Spec}(\mathbb Z)$; \underline{let} $i,j$ \underline{be integers}.  \underline{Suppose that the
Hodge coniveau of} $H^i$ \underline{of the generic fiber of} $f$ \underline{is} at least $j$.  \underline{Then,
for every point} $s\in S$, \underline{the coniveau} (XX, 2.0.6) \underline{of}
$H^i(X_s,\mathbb Q_\ell)$ \underline{is} at least $j$.  \underline{Hence} (XX, 2.2), \underline{if} $s$ \underline{is a
closed point of} $S$, \underline{the eigenvalues of Frobenius acting on}
$H^i(X_s,\mathbb Q_\ell)$ \underline{are algebraic integers divisible by} $q^j$, \underline{where}
$q=\operatorname{card}(k(s))$.

\section*{6. End of the proof of the congruence formula}

\SGAvar{Taking account of (4.6.6) and (4.6.3)--(4.6.4), formula (4.5.24) gives

}{Taking account of (4.6.6) and (4.6.3)--(4.6.4), formula (4.5.33) gives

}{E7II-XXII-040}\begin{equation}\tag{6.0.1}
Z(t,V/\mathbb F_q)(1-t)
\equiv\Delta_S(t)^{(-1)^{n+1}}
\mod{1+qt\mathcal O[[t]]}.
\end{equation}

We now calculate $\Delta_S(t)$ modulo $\pi\mathcal O[[t]]$.  Define, as in
(4.6.11),

\SGAvar{\begin{equation}\tag{6.0.2}
W_S=R_B\otimes_{\mathcal O}\mathbb F_q.
\end{equation}



}{\begin{equation}\tag{6.0.2}
W_S=R_S\otimes_{\mathcal O}\mathbb F_q.
\end{equation}



}{E7II-XXII-041}\SGApage{II}{\(\approx\)432}
\begin{equation}\tag{6.0.3}
\bar\alpha_S:W_S\longrightarrow W_S
\end{equation}

\noindent
denotes the $\tau^{-1}$-linear operation induced by $\alpha_S$;

\SGAvar{\begin{equation}\tag{6.0.4}
R_S(t)
\text{ the submodule of $R_S$ formed by the elements of degree $i$ in $X_0$},
\end{equation}

}{\begin{equation}\tag{6.0.4}
R_S(i)
\text{ the submodule of $R_S$ formed by the elements of degree $i$ in $X_0$},
\end{equation}

}{E7II-XXII-042}and

\begin{equation}\tag{6.0.5}
W_S(i)=R_S(i)\otimes_{\mathcal O}\mathbb F_q.
\end{equation}

\subsubsection*{Proposition 6.0.6.}
\underline{One has}

\[
\alpha_S(R_S)\subset R_S(1)+pR_S.
\]

\noindent\emph{Proof.}
We have the factorization (4.6.14)

\begin{equation}\tag{6.0.7}
\alpha_S=\tau^{-1}\circ\frac1p\,\psi_p\circ H(X).
\end{equation}

We know (see (4.6.15)) that

\begin{equation}\tag{6.0.8}
H(X)\cdot R_S\subset R_S,
\end{equation}

and plainly

\begin{equation}\tag{6.0.9}
\tau^{-1}\bigl(R_S(1)+pR_S\bigr)=R_S(1)+pR_S.
\end{equation}

It is therefore enough to verify

\begin{equation}\tag{6.0.10}
\frac1p\,\psi_p(R_S)\subset R_S(1)+pR_S,
\end{equation}

which is clear from (4.6.19).

\subsubsection*{Corollary 6.0.11.}
\underline{One has}
\[
\bar\alpha_S W_S\subset W_S(1).
\]

\SGApage{II}{433}
\subsubsection*{Corollary 6.0.12.}
\underline{One has the congruence}

\[
\Delta_S(t)\equiv
\det\!\left(
1-t(\bar\alpha_S)^{\operatorname{ord}(q)}
\mathrel{\Big|}W_S(1)\right)
\pmod{\pi\mathcal O[[t]]}.
\]


\SGAvar{\subsubsection*{Proposition 6.1.}
\underline{The} $\tau^{1}$-\underline{linear operation}
$\bar\alpha_S:W_S(1)\to W_S(1)$ \underline{is equal to}

}{\subsubsection*{Proposition 6.1.}
\underline{The} $\tau^{-1}$-\underline{linear operation}
$\bar\alpha_S:W_S(1)\to W_S(1)$ \underline{is equal to}

}{E7II-XXII-043}\begin{equation}\tag{6.1.1}
\tau^{-1}\psi_p\circ(X_0f(X))^{p-1}.
\end{equation}

\noindent\emph{Proof.}
Write

\begin{equation}\tag{6.1.2}
H(X)=1+\sum_{i\geq1}H_i(X),
\end{equation}

where

\begin{equation}\tag{6.1.3}
H_i(X)\text{ has degree $i$ in }X_0.
\end{equation}

Again by (4.6.19),

\begin{equation}\tag{6.1.4}
\frac1p\,\psi_pR_S(i)\subset pR_S
\qquad\text{if }i\ne p.
\end{equation}

Since

\begin{equation}\tag{6.1.5}
H_i(X)\cdot R_S(j)\subset R_S(i+j),
\end{equation}

we have

\begin{equation}\tag{6.1.6}
\left(
\alpha_S-\frac{\tau^{-1}}p\,\psi_pH_{p-1}(X)
\right)(R_S(1))
\subset pR_S.
\end{equation}

Consequently,

\begin{equation}\tag{6.1.7}
\left[
(\alpha_S)^{\operatorname{ord}(q)}
-
\left(
\frac{\tau^{-1}}p\,\psi_p\cdot H_{p-1}(X)
\right)^{\operatorname{ord}(q)}
\right]R_S(1)
\subset pR_S.
\end{equation}

It remains to calculate $H_{p-1}(X)$.  From (4.4.4),

\SGAvar{\begin{equation}\tag{6.1.8}
\begin{aligned}
H(X)
&=\prod_W\Theta(A_WX_0X^W)\\
&=\prod_W\exp\!\left(
\sum_{n\geq0}\frac{\pi A_WX_0X^W)^{p^n}}{p^n}
\right)\\
&=\prod_W\prod_{n\geq0}
\exp\!\left(\frac{(\pi A_WX_0X^W)^{p^n}}{p^n}\right)\\
&\equiv\prod_W\exp(\pi A_WX_0X^W)
\pmod{\text{terms of degree }\geq p}\\
&=\exp\!\left(\pi\sum A_WX_0X^W\right)
\pmod{\text{terms of degree }\geq p}\\
&=\exp(\pi X_0F(X))
\pmod{\text{terms of degree }\geq p}.
\end{aligned}
\end{equation}

}{\begin{equation}\tag{6.1.8}
\begin{aligned}
H(X)
&=\prod_W\Theta(A_WX_0X^W)\\
&=\prod_W\exp\!\left(
\sum_{n\geq0}\frac{(\pi A_WX_0X^W)^{p^n}}{p^n}
\right)\\
&=\prod_W\prod_{n\geq0}
\exp\!\left(\frac{(\pi A_WX_0X^W)^{p^n}}{p^n}\right)\\
&\equiv\prod_W\exp(\pi A_WX_0X^W)
\pmod{\text{terms of degree }\geq p}\\
&=\exp\!\left(\pi\sum A_WX_0X^W\right)
\pmod{\text{terms of degree }\geq p}\\
&=\exp(\pi X_0F(X))
\pmod{\text{terms of degree }\geq p}.
\end{aligned}
\end{equation}

}{E7II-XXII-044}\SGApage{II}{\(\approx\)434}
It follows that

\begin{equation}\tag{6.1.9}
H_{p-1}(X)
\equiv\frac{\pi^{p-1}X_0^{p-1}}{(p-1)!}\,F(X)^{p-1}.
\end{equation}

One easily checks that

\begin{equation}\tag{6.1.10}
\frac{\pi^{p-1}}p\equiv-1\pmod p,
\end{equation}

and hence

\SGAvar{\begin{equation}\tag{6.1.11}
\frac1p\,H_{p-1}(X)
\equiv(X_0f(X)^{p-1}\pmod p.
\end{equation}

}{\begin{equation}\tag{6.1.11}
\frac1p\,H_{p-1}(X)
\equiv(X_0f(X))^{p-1}\pmod p.
\end{equation}

}{E7II-XXII-045}This completes the proof of Proposition~6.1.

\subsubsection*{Corollary 6.1.12.}
\underline{The linear operator}
\[
(\bar\alpha_S)^{\operatorname{ord}(q)}:W_S(1)\longrightarrow W_S(1)
\]
\underline{is precisely}
\[
\psi_q\circ(X_0f(X))^{q-1}.
\]

\SGApage{II}{\(\approx\)435}
\noindent\emph{Proof.}
By (6.1.1),

\SGAvar{\begin{equation*}
\begin{aligned}
(\bar\alpha_S)^{\operatorname{ord}(q)}
&=
\left(
\tau^{-1}\psi_p\circ(X_0f(X))^{p-1}
\right)^{\operatorname{ord}(q)}\\
&=
\psi_q
\prod_{j=0}^{\operatorname{ord}(a)-1}
\left(X_0^{p^j}f^{\tau^j}(X^{p^j})\right)^{p-1}\\
&=
\psi_q
\prod_{j=0}^{\operatorname{ord}(q)-1}
(X_0f(X))^{p^j(p-1)}\\
&=\psi_q\circ(X_0f(X))^{q-1}.
\end{aligned}
\end{equation*}



}{\begin{equation*}
\begin{aligned}
(\bar\alpha_S)^{\operatorname{ord}(q)}
&=
\left(
\tau^{-1}\psi_p\circ(X_0f(X))^{p-1}
\right)^{\operatorname{ord}(q)}\\
&=
\psi_q\circ
\prod_{j=0}^{\operatorname{ord}(q)-1}
\left(X_0^{p^j}f^{\tau^j}(X^{p^j})\right)^{p-1}\\
&=
\psi_q\circ
\prod_{j=0}^{\operatorname{ord}(q)-1}
(X_0f(X))^{p^j(p-1)}\\
&=\psi_q\circ(X_0f(X))^{q-1}.
\end{aligned}
\end{equation*}



}{E7II-XXII-046}\subsubsection*{Corollary 6.1.13.}
\underline{Modulo} $p$, \underline{one has the congruence}

\begin{equation*}
Z(t,V/\mathbb F_q)
\equiv
\frac{
\det\!\left(
1-t\,\psi_q(X_0f(X))^{q-1}
\mathrel{\Big|}W_S(1)
\right)^{(-1)^{n+1}}
}{1-t}.
\end{equation*}



\subsection*{6.2.}
It remains to calculate the cohomology groups of $\mathcal O_V$.
Put

\begin{equation}\tag{6.2.1}
\mathfrak F_q
=\text{the operation of raising to the $q$-th power in $\mathcal O_V$}.
\end{equation}

There is an exact sequence of sheaves on
$\mathbb P^{n+1}=\mathbb P$,

\SGAvar{\begin{equation}\tag{6.2.2}
0\longrightarrow\mathcal O_{\mathbb P}(-d)
\xrightarrow{\,f\,}\mathcal O_{\mathbb P}
\longrightarrow\mathcal O_V\longrightarrow\mathcal O,
\end{equation}

}{\begin{equation}\tag{6.2.2}
0\longrightarrow\mathcal O_{\mathbb P}(-d)
\xrightarrow{\,f\,}\mathcal O_{\mathbb P}
\longrightarrow\mathcal O_V\longrightarrow0,
\end{equation}

}{E7II-XXII-047}and the following diagram is commutative:

\SGAvar{\begin{equation}\tag{6.2.3}
\begin{tikzcd}[column sep=large,row sep=large]
0 \arrow[r] \&
\mathcal O_{\mathbb P}(-d) \arrow[r,"f"] \arrow[d,"f^{q-1}\mathfrak F_q"'] \&
\mathcal O_{\mathbb P} \arrow[r] \arrow[d,"\mathfrak F_q"'] \&
\mathcal O_V \arrow[r] \arrow[d,"\mathfrak F_q"'] \&
\mathcal O\\
0 \arrow[r] \&
\mathcal O_{\mathbb P}(-d) \arrow[r,"f"] \&
\mathcal O_{\mathbb P} \arrow[r] \&
\mathcal O_V \arrow[r] \&
\mathcal O .
\end{tikzcd}
\end{equation}

}{\begin{equation}\tag{6.2.3}
\begin{tikzcd}[column sep=large,row sep=large]
0 \arrow[r] \&
\mathcal O_{\mathbb P}(-d) \arrow[r,"f"] \arrow[d,"f^{q-1}\mathfrak F_q"'] \&
\mathcal O_{\mathbb P} \arrow[r] \arrow[d,"\mathfrak F_q"'] \&
\mathcal O_V \arrow[r] \arrow[d,"\mathfrak F_q"'] \&
0\\
0 \arrow[r] \&
\mathcal O_{\mathbb P}(-d) \arrow[r,"f"] \&
\mathcal O_{\mathbb P} \arrow[r] \&
\mathcal O_V \arrow[r] \&
0 .
\end{tikzcd}
\end{equation}

}{E7II-XXII-048}\SGApage{II}{\(\approx\)436}
Taking account of

\begin{equation}\tag{6.2.4}
\begin{cases}
H^0(\mathbb P^{n+1},\mathcal O_{\mathbb P})\simeq\mathbb F_q,\\
H^i(\mathbb P^{n+1},\mathcal O_{\mathbb P})=0,
&i\ne0,\\
H^i(\mathbb P^{n+1},\mathcal O_{\mathbb P}(-d))=0,
&i\ne0,n+1,
\end{cases}
\end{equation}

the long exact cohomology sequence associated with (6.2.2) gives

\begin{equation}\tag{6.2.5}
\begin{cases}
H^0(V,\mathcal O_V)\simeq\mathbb F_q,\\
H^i(V,\mathcal O_V)=0,&i\ne0,n,\\
H^n(V,\mathcal O_V)
\simeq H^{n+1}(\mathbb P^{n+1},\mathcal O_{\mathbb P}(-d)).
\end{cases}
\end{equation}

The commutative diagram (6.2.3) therefore gives

\begin{equation}\tag{6.2.6}
\det\!\left(1-t\mathfrak F_q
\mathrel{\Big|}H^0(V,\mathcal O_V)\right)=1-t
\end{equation}

and

\SGAvar{\begin{equation}\tag{6.2.7}
\det\!\left(1-t\mathfrak F_q
\mathrel{\Big|}H^n(V,\mathcal O_V)\right)
=
\det\!\left(
1-t(f^{q-1}\mathfrak F_q)
\mathrel{\Big|}
H^{n+1}(\mathbb P^{n+1},\mathcal O_{\mathbb P}(-d))
\right..
\end{equation}

}{\begin{equation}\tag{6.2.7}
\det\!\left(1-t\mathfrak F_q
\mathrel{\Big|}H^n(V,\mathcal O_V)\right)
=
\det\!\left(
1-t(f^{q-1}\mathfrak F_q)
\mathrel{\Big|}
H^{n+1}(\mathbb P^{n+1},\mathcal O_{\mathbb P}(-d))
\right).
\end{equation}

}{E7II-XXII-049}It remains only to prove

\SGAvar{\begin{equation}\tag{6.2.8}
\begin{aligned}
&\det\!\left(
1-t(f^{q-1}\mathfrak F_q)
\mathrel{\Big|}
H^{n+1}(\mathbb P^{n+1},\mathcal O_{\mathbb P}(-d))
\right.\\
&\qquad=
\det\!\left(
1-t\,(\psi_q\cdot(X_0f(X))^{q-1}
\mathrel{\Big|}W_S(1)
\right)
\end{aligned}
\end{equation}

}{\begin{equation}\tag{6.2.8}
\begin{aligned}
&\det\!\left(
1-t(f^{q-1}\mathfrak F_q)
\mathrel{\Big|}
H^{n+1}(\mathbb P^{n+1},\mathcal O_{\mathbb P}(-d))
\right)\\
&\qquad=
\det\!\left(
1-t\,\psi_q\cdot(X_0f(X))^{q-1}
\mathrel{\Big|}W_S(1)
\right)
\end{aligned}
\end{equation}

}{E7II-XXII-050}which is a consequence of the following proposition.



\subsubsection*{\underline{Proposition} 6.3.}
\underline{There is an isomorphism}

\SGAvar{\begin{equation}\tag{6.3.1}
W_S(1)\simeq
\operatorname{Hom}_{\mathbb F_q}\!\left(
H^{n+1}(\mathbb P^{n+1}.\mathcal O_{\mathbb P}(-d),
\mathbb F_q\right)
\end{equation}

}{\begin{equation}\tag{6.3.1}
W_S(1)\simeq
\operatorname{Hom}_{\mathbb F_q}\!\left(
H^{n+1}(\mathbb P^{n+1},\mathcal O_{\mathbb P}(-d)),
\mathbb F_q\right)
\end{equation}

}{E7II-XXII-051}\underline{with respect to which the operator}

\begin{equation}\tag{6.3.2}
\psi_q\circ(X_0f(X))^{q-1}:W_S(1)\longrightarrow W_S(1)
\end{equation}

\SGApage{II}{\(\approx\)437}
\underline{is the transpose of}

\begin{equation}\tag{6.3.3}
f(X)^{q-1}\mathfrak F_q:
H^{n+1}(\mathbb P^{n+1},\mathcal O_{\mathbb P}(-d))
\longrightarrow
H^{n+1}(\mathbb P^{n+1},\mathcal O_{\mathbb P}(-d)).
\end{equation}

\noindent\emph{Proof.}
Computing
$H^{n+1}(\mathbb P^{n+1},\mathcal O_{\mathbb P}(-d))$
by means of the standard covering of $\mathbb P^{n+1}$ gives

\SGAvar{\begin{equation}\tag{6.3.4}
\begin{aligned}
H^{i+1}(\mathbb P^{n+1},\mathcal O_{\mathbb P}(-d))
&=\text{the $\mathbb F_q$-vector space with basis}\\
&\quad X_1^{W_1}\cdots X_{n+2}^{W_{n+2}}
\qquad\text{with }W_i\in\mathbb Z\text{ satisfying }\sum W_i=-d,\\
&=\text{the $\mathbb F_q$-subspace with basis those monomials}\\
&\quad X_1^{W_1}\cdots X_{n+2}^{W_{n+2}}
\text{ for which }\exists i\text{ with }W_i\geq0.
\end{aligned}
\end{equation}

}{\begin{equation}\tag{6.3.4}
\begin{aligned}
H^{n+1}(\mathbb P^{n+1},\mathcal O_{\mathbb P}(-d))
&=\text{the quotient of the $\mathbb F_q$-vector space with basis}\\
&\quad X_1^{W_1}\cdots X_{n+2}^{W_{n+2}}
\qquad\text{with }W_i\in\mathbb Z\text{ satisfying }\sum W_i=-d,\\
&\quad\text{by the $\mathbb F_q$-subspace with basis those monomials}\\
&\quad X_1^{W_1}\cdots X_{n+2}^{W_{n+2}}
\text{ for which }\exists i\text{ with }W_i\geq0.
\end{aligned}
\end{equation}

}{E7II-XXII-052}On the other hand, $W_S(1)=R_S(1)\otimes_{\mathcal O}\mathbb F_q$, and

\begin{equation}\tag{6.3.5}
\begin{aligned}
R_S(1)
&=\text{the free $\mathcal O$-module with basis the monomials}\\
&\quad
(\pi X_0)X_1^{U_1}\cdots X_{n+2}^{U_{n+2}}
\text{ satisfying }
\sum U_i=d,\quad U_i>0\ \forall_i.
\end{aligned}
\end{equation}

The pairing

\begin{equation}\tag{6.3.6}
(\ ,\ ):
W_S(1)\times
H^{n+1}(\mathbb P^{n+1},\mathcal O_{\mathbb P}(-d))
\longrightarrow\mathbb F_q
\end{equation}

defined by

\SGAvar{\begin{equation}\tag{6.3.7}
\left(
\text{class of }\pi X_0X_1^{U_1}\cdots X_{n+2}^{U_{n+2}},
X_1^{W_1}\cdots X_{n+2}^{W_{n+2}}
\right)
=
\begin{cases}
1,&U_1=-W_i\text{ for every }i,\\
0,&\text{otherwise},
\end{cases}
\end{equation}

is plainly a perfect duality of $\mathbb F_q$-vector spaces.  One checks
immediately that
$\psi_q\cdot(X_0f(X))^{q-1}$ is the transpose of
$f(X)^{q-1}\cdot f_q$, which proves Proposition~6.3.



}{\begin{equation}\tag{6.3.7}
\left(
\text{class of }\pi X_0X_1^{U_1}\cdots X_{n+2}^{U_{n+2}},
X_1^{W_1}\cdots X_{n+2}^{W_{n+2}}
\right)
=
\begin{cases}
1,&U_i=-W_i\text{ for every }i,\\
0,&\text{otherwise},
\end{cases}
\end{equation}

is plainly a perfect duality of $\mathbb F_q$-vector spaces.  One checks
immediately that
$\psi_q\cdot(X_0f(X))^{q-1}$ is the transpose of
$f(X)^{q-1}\cdot \mathfrak F_q$, which proves Proposition~6.3.



}{E7II-XXII-053}\section*{Bibliography}

\begin{enumerate}[label={[\arabic*]},leftmargin=2.4em]
\item J.~Ax, \emph{Zeroes of Polynomials Over Finite Fields},
  \emph{Amer. J. Math.} \textbf{86} (1964), 255--261.
\item B.~Dwork, \emph{On the Rationality of the Zeta Function of an Algebraic
  Variety}, \emph{Amer. J. Math.} \textbf{82} (1960), 631--648.
\item B.~Dwork, \emph{An Algebraic Variety}, in \emph{Proceedings of the
  Woods Hole Summer Institute 1964} (mimeographed).
\item B.~Dwork, \emph{On the Zeta Function of a Hypersurface II},
  \emph{Ann. of Math.} (2) \textbf{80} (1964), 227--299.
\item A.~Grothendieck, \emph{Le Groupe de Brauer III}, in
  \emph{Dix Exposés sur la Cohomologie des Schémas},
  North-Holland Publishing Co., 1968.
\item J.-P.~Serre, \emph{Groupes algébriques et corps de classes},
  Hermann, Paris, 1959.
\item J.-P.~Serre, \emph{Endomorphismes complètement continus des espaces de
  Banach $p$-adiques}, \emph{Publ. Math. I.H.E.S.} no.~12 (1962).
\item C.~Chevalley, \emph{Démonstration d'une hypothèse de M.~Artin},
  \emph{Abh. Math. Sem. Univ. Hamburg} \textbf{11} (1936), 73--75.
\end{enumerate}


\markboth{Exposé XXII}{Exposé XXII}

\ifnum\SGAform=1
\clearpage\section*{Critical apparatus}\begingroup\small\raggedright\sloppy
Readings are given in the exact TeX notation of the two English witnesses. Decimal suffixes identify the individual differences in each marked passage.

\par\medskip\noindent\hypertarget{app-E7II-XXII-001.1}{}\textbf{E7II-XXII-001.1}\quad printed folio 404.\par
\noindent Printed reading: {\ttfamily\footnotesize Gl\allowbreak{}}\par
\noindent Corrected reading: {\ttfamily\footnotesize GL\allowbreak{}}\par
\noindent Source: {\ttfamily\footnotesize fr/sga7ii/\{diplomati\allowbreak{}c,current\}/expose\_XX\allowbreak{}II\_body.tex;\allowbreak{} folio\allowbreak{} 404;\allowbreak{} formula\allowbreak{} 82\allowbreak{}}\par

\par\medskip\noindent\hypertarget{app-E7II-XXII-001.2}{}\textbf{E7II-XXII-001.2}\quad printed folio 404.\par
\noindent Printed reading: {\ttfamily\footnotesize Gl\allowbreak{}}\par
\noindent Corrected reading: {\ttfamily\footnotesize GL\allowbreak{}}\par
\noindent Source: {\ttfamily\footnotesize fr/sga7ii/\{diplomati\allowbreak{}c,current\}/expose\_XX\allowbreak{}II\_body.tex;\allowbreak{} folio\allowbreak{} 404;\allowbreak{} formula\allowbreak{} 82\allowbreak{}}\par

\par\medskip\noindent\hypertarget{app-E7II-XXII-002.1}{}\textbf{E7II-XXII-002.1}\quad printed folio 405.\par
\noindent Printed reading: \emph{(empty)}\par
\noindent Corrected reading: {\ttfamily\footnotesize \$V\$\allowbreak{} }\par
\noindent Source: {\ttfamily\footnotesize ED-EN-FORMULA-G237;\allowbreak{} fr/sga7ii/\{diplomati\allowbreak{}c,current\}/expose\_XX\allowbreak{}II\_body.tex;\allowbreak{} folio\allowbreak{} 405;\allowbreak{} port/PORT\_EN7II\_LOG.\allowbreak{}jsonl\#C826;\allowbreak{} fixes/FIXLOG\_7II\_CUR\allowbreak{}RENT.jsonl\#C826;\allowbreak{} 7II:840\allowbreak{}}\par

\par\medskip\noindent\hypertarget{app-E7II-XXII-003.1}{}\textbf{E7II-XXII-003.1}\quad printed folio 408.\par
\noindent Printed reading: \emph{(empty)}\par
\noindent Corrected reading: {\ttfamily\footnotesize induced\allowbreak{} by\allowbreak{} }\par
\noindent Source: {\ttfamily\footnotesize ED-EN-FORMULA-G239;\allowbreak{} fr/sga7ii/\{diplomati\allowbreak{}c,current\}/expose\_XX\allowbreak{}II\_body.tex;\allowbreak{} folio\allowbreak{} 408\allowbreak{}}\par

\par\medskip\noindent\hypertarget{app-E7II-XXII-003.2}{}\textbf{E7II-XXII-003.2}\quad printed folio 408.\par
\noindent Printed reading: {\ttfamily\footnotesize \textbackslash{}otimes\_\{\textbackslash{}mathbb\allowbreak{} F\_q\}k\allowbreak{}}\par
\noindent Corrected reading: \emph{(empty)}\par
\noindent Source: {\ttfamily\footnotesize ED-EN-FORMULA-G239;\allowbreak{} fr/sga7ii/\{diplomati\allowbreak{}c,current\}/expose\_XX\allowbreak{}II\_body.tex;\allowbreak{} folio\allowbreak{} 408\allowbreak{}}\par

\par\medskip\noindent\hypertarget{app-E7II-XXII-003.3}{}\textbf{E7II-XXII-003.3}\quad printed folio 408.\par
\noindent Printed reading: {\ttfamily\footnotesize of\allowbreak{}}\par
\noindent Corrected reading: {\ttfamily\footnotesize on\allowbreak{}}\par
\noindent Source: {\ttfamily\footnotesize ED-EN-FORMULA-G239;\allowbreak{} fr/sga7ii/\{diplomati\allowbreak{}c,current\}/expose\_XX\allowbreak{}II\_body.tex;\allowbreak{} folio\allowbreak{} 408\allowbreak{}}\par

\par\medskip\noindent\hypertarget{app-E7II-XXII-004.1}{}\textbf{E7II-XXII-004.1}\quad printed folio 409.\par
\noindent Printed reading: {\ttfamily\footnotesize k\allowbreak{}}\par
\noindent Corrected reading: {\ttfamily\footnotesize q\allowbreak{}}\par
\noindent Source: {\ttfamily\footnotesize fr/sga7ii/\{diplomati\allowbreak{}c,current\}/expose\_XX\allowbreak{}II\_body.tex;\allowbreak{} folio\allowbreak{} 409;\allowbreak{} formula\allowbreak{} 174\allowbreak{}}\par

\par\medskip\noindent\hypertarget{app-E7II-XXII-004.2}{}\textbf{E7II-XXII-004.2}\quad printed folio 409.\par
\noindent Printed reading: {\ttfamily\footnotesize k\allowbreak{}}\par
\noindent Corrected reading: {\ttfamily\footnotesize q\allowbreak{}}\par
\noindent Source: {\ttfamily\footnotesize fr/sga7ii/\{diplomati\allowbreak{}c,current\}/expose\_XX\allowbreak{}II\_body.tex;\allowbreak{} folio\allowbreak{} 409;\allowbreak{} formula\allowbreak{} 174\allowbreak{}}\par

\par\medskip\noindent\hypertarget{app-E7II-XXII-005.1}{}\textbf{E7II-XXII-005.1}\quad printed folio 409.\par
\noindent Printed reading: \emph{(empty)}\par
\noindent Corrected reading: {\ttfamily\footnotesize \textbackslash{}overline\allowbreak{}}\par
\noindent Source: {\ttfamily\footnotesize fr/sga7ii/\{diplomati\allowbreak{}c,current\}/expose\_XX\allowbreak{}II\_body.tex;\allowbreak{} folio\allowbreak{} 409;\allowbreak{} formula\allowbreak{} 188\allowbreak{}}\par

\par\medskip\noindent\hypertarget{app-E7II-XXII-005.2}{}\textbf{E7II-XXII-005.2}\quad printed folio 409.\par
\noindent Printed reading: \emph{(empty)}\par
\noindent Corrected reading: {\ttfamily\footnotesize \textbackslash{}overline\allowbreak{}}\par
\noindent Source: {\ttfamily\footnotesize fr/sga7ii/\{diplomati\allowbreak{}c,current\}/expose\_XX\allowbreak{}II\_body.tex;\allowbreak{} folio\allowbreak{} 409;\allowbreak{} formula\allowbreak{} 188\allowbreak{}}\par

\par\medskip\noindent\hypertarget{app-E7II-XXII-005.3}{}\textbf{E7II-XXII-005.3}\quad printed folio 409.\par
\noindent Printed reading: \emph{(empty)}\par
\noindent Corrected reading: {\ttfamily\footnotesize \textasciicircum{}\{-1\}\allowbreak{}}\par
\noindent Source: {\ttfamily\footnotesize ED-EN-FORMULA-G241;\allowbreak{} fr/sga7ii/\{diplomati\allowbreak{}c,current\}/expose\_XX\allowbreak{}II\_body.tex;\allowbreak{} folio\allowbreak{} 409\allowbreak{}}\par

\par\medskip\noindent\hypertarget{app-E7II-XXII-006.1}{}\textbf{E7II-XXII-006.1}\quad printed folio 410.\par
\noindent Printed reading: {\ttfamily\footnotesize 6\allowbreak{}}\par
\noindent Corrected reading: {\ttfamily\footnotesize 3\allowbreak{}}\par
\noindent Source: {\ttfamily\footnotesize ED-EN-NONMATH-G042;\allowbreak{} F7II-XXII-019.1;\allowbreak{} WITNESS-F7II-XXII-01\allowbreak{}9.1;\allowbreak{} fr/sga7ii/\{diplomati\allowbreak{}c,current\}/expose\_XX\allowbreak{}II\_body.tex;\allowbreak{} folio\allowbreak{} 410\allowbreak{}}\par

\par\medskip\noindent\hypertarget{app-E7II-XXII-006.2}{}\textbf{E7II-XXII-006.2}\quad printed folio 410.\par
\noindent Printed reading: {\ttfamily\footnotesize 1\allowbreak{}}\par
\noindent Corrected reading: {\ttfamily\footnotesize 2\allowbreak{}}\par
\noindent Source: {\ttfamily\footnotesize ED-EN-NONMATH-G042;\allowbreak{} F7II-XXII-019.1;\allowbreak{} WITNESS-F7II-XXII-01\allowbreak{}9.1;\allowbreak{} fr/sga7ii/\{diplomati\allowbreak{}c,current\}/expose\_XX\allowbreak{}II\_body.tex;\allowbreak{} folio\allowbreak{} 410\allowbreak{}}\par

\par\medskip\noindent\hypertarget{app-E7II-XXII-006.3}{}\textbf{E7II-XXII-006.3}\quad printed folio 410.\par
\noindent Printed reading: \emph{(empty)}\par
\noindent Corrected reading: {\ttfamily\footnotesize )--(3.2\allowbreak{}}\par
\noindent Source: {\ttfamily\footnotesize ED-EN-NONMATH-G042;\allowbreak{} F7II-XXII-019.1;\allowbreak{} WITNESS-F7II-XXII-01\allowbreak{}9.1;\allowbreak{} fr/sga7ii/\{diplomati\allowbreak{}c,current\}/expose\_XX\allowbreak{}II\_body.tex;\allowbreak{} folio\allowbreak{} 410\allowbreak{}}\par

\par\medskip\noindent\hypertarget{app-E7II-XXII-006.4}{}\textbf{E7II-XXII-006.4}\quad printed folio 410.\par
\noindent Printed reading: \emph{(empty)}\par
\noindent Corrected reading: {\ttfamily\footnotesize 3.\allowbreak{}}\par
\noindent Source: {\ttfamily\footnotesize ED-EN-NONMATH-G042;\allowbreak{} F7II-XXII-019.1;\allowbreak{} WITNESS-F7II-XXII-01\allowbreak{}9.1;\allowbreak{} fr/sga7ii/\{diplomati\allowbreak{}c,current\}/expose\_XX\allowbreak{}II\_body.tex;\allowbreak{} folio\allowbreak{} 410\allowbreak{}}\par

\par\medskip\noindent\hypertarget{app-E7II-XXII-006.5}{}\textbf{E7II-XXII-006.5}\quad printed folio 410.\par
\noindent Printed reading: {\ttfamily\footnotesize .3\allowbreak{}}\par
\noindent Corrected reading: \emph{(empty)}\par
\noindent Source: {\ttfamily\footnotesize ED-EN-NONMATH-G042;\allowbreak{} F7II-XXII-019.1;\allowbreak{} WITNESS-F7II-XXII-01\allowbreak{}9.1;\allowbreak{} fr/sga7ii/\{diplomati\allowbreak{}c,current\}/expose\_XX\allowbreak{}II\_body.tex;\allowbreak{} folio\allowbreak{} 410\allowbreak{}}\par

\par\medskip\noindent\hypertarget{app-E7II-XXII-007.1}{}\textbf{E7II-XXII-007.1}\quad printed folio 410.\par
\noindent Printed reading: {\ttfamily\footnotesize .\allowbreak{}}\par
\noindent Corrected reading: {\ttfamily\footnotesize )\allowbreak{}}\par
\noindent Source: {\ttfamily\footnotesize fr/sga7ii/\{diplomati\allowbreak{}c,current\}/expose\_XX\allowbreak{}II\_body.tex;\allowbreak{} folio\allowbreak{} 410;\allowbreak{} formula\allowbreak{} 215\allowbreak{}}\par

\par\medskip\noindent\hypertarget{app-E7II-XXII-008.1}{}\textbf{E7II-XXII-008.1}\quad printed folio 411.\par
\noindent Printed reading: \emph{(empty)}\par
\noindent Corrected reading: {\ttfamily\footnotesize (\allowbreak{}}\par
\noindent Source: {\ttfamily\footnotesize fr/sga7ii/\{diplomati\allowbreak{}c,current\}/expose\_XX\allowbreak{}II\_body.tex;\allowbreak{} folio\allowbreak{} 411;\allowbreak{} formula\allowbreak{} 227\allowbreak{}}\par

\par\medskip\noindent\hypertarget{app-E7II-XXII-008.2}{}\textbf{E7II-XXII-008.2}\quad printed folio 411.\par
\noindent Printed reading: {\ttfamily\footnotesize p\allowbreak{}}\par
\noindent Corrected reading: {\ttfamily\footnotesize q\allowbreak{}}\par
\noindent Source: {\ttfamily\footnotesize fr/sga7ii/\{diplomati\allowbreak{}c,current\}/expose\_XX\allowbreak{}II\_body.tex;\allowbreak{} folio\allowbreak{} 411;\allowbreak{} formula\allowbreak{} 227\allowbreak{}}\par

\par\medskip\noindent\hypertarget{app-E7II-XXII-008.3}{}\textbf{E7II-XXII-008.3}\quad printed folio 411.\par
\noindent Printed reading: {\ttfamily\footnotesize \{\allowbreak{}}\par
\noindent Corrected reading: \emph{(empty)}\par
\noindent Source: {\ttfamily\footnotesize fr/sga7ii/\{diplomati\allowbreak{}c,current\}/expose\_XX\allowbreak{}II\_body.tex;\allowbreak{} folio\allowbreak{} 411;\allowbreak{} formula\allowbreak{} 227\allowbreak{}}\par

\par\medskip\noindent\hypertarget{app-E7II-XXII-008.4}{}\textbf{E7II-XXII-008.4}\quad printed folio 411.\par
\noindent Printed reading: \emph{(empty)}\par
\noindent Corrected reading: {\ttfamily\footnotesize )\textasciicircum{}\allowbreak{}}\par
\noindent Source: {\ttfamily\footnotesize fr/sga7ii/\{diplomati\allowbreak{}c,current\}/expose\_XX\allowbreak{}II\_body.tex;\allowbreak{} folio\allowbreak{} 411;\allowbreak{} formula\allowbreak{} 227\allowbreak{}}\par

\par\medskip\noindent\hypertarget{app-E7II-XXII-008.5}{}\textbf{E7II-XXII-008.5}\quad printed folio 411.\par
\noindent Printed reading: {\ttfamily\footnotesize \}\allowbreak{}}\par
\noindent Corrected reading: \emph{(empty)}\par
\noindent Source: {\ttfamily\footnotesize fr/sga7ii/\{diplomati\allowbreak{}c,current\}/expose\_XX\allowbreak{}II\_body.tex;\allowbreak{} folio\allowbreak{} 411;\allowbreak{} formula\allowbreak{} 227\allowbreak{}}\par

\par\medskip\noindent\hypertarget{app-E7II-XXII-008.6}{}\textbf{E7II-XXII-008.6}\quad printed folio 411.\par
\noindent Printed reading: {\ttfamily\footnotesize q\allowbreak{}}\par
\noindent Corrected reading: {\ttfamily\footnotesize p\allowbreak{}}\par
\noindent Source: {\ttfamily\footnotesize fr/sga7ii/\{diplomati\allowbreak{}c,current\}/expose\_XX\allowbreak{}II\_body.tex;\allowbreak{} folio\allowbreak{} 411;\allowbreak{} formula\allowbreak{} 230\allowbreak{}}\par

\par\medskip\noindent\hypertarget{app-E7II-XXII-009.1}{}\textbf{E7II-XXII-009.1}\quad printed folio 411.\par
\noindent Printed reading: {\ttfamily\footnotesize q\allowbreak{}}\par
\noindent Corrected reading: {\ttfamily\footnotesize p\allowbreak{}}\par
\noindent Source: {\ttfamily\footnotesize fr/sga7ii/\{diplomati\allowbreak{}c,current\}/expose\_XX\allowbreak{}II\_body.tex;\allowbreak{} folio\allowbreak{} 411;\allowbreak{} formula\allowbreak{} 233\allowbreak{}}\par

\par\medskip\noindent\hypertarget{app-E7II-XXII-010.1}{}\textbf{E7II-XXII-010.1}\quad printed folio 412.\par
\noindent Printed reading: {\ttfamily\footnotesize b\allowbreak{}}\par
\noindent Corrected reading: {\ttfamily\footnotesize a\allowbreak{}}\par
\noindent Source: {\ttfamily\footnotesize fr/sga7ii/\{diplomati\allowbreak{}c,current\}/expose\_XX\allowbreak{}II\_body.tex;\allowbreak{} folio\allowbreak{} 412;\allowbreak{} formula\allowbreak{} 247;\allowbreak{} port/PORT\_EN7II\_LOG.\allowbreak{}jsonl\#C899;\allowbreak{} fixes/FIXLOG\_7II\_CUR\allowbreak{}RENT.jsonl\#C899;\allowbreak{} 7II:888\allowbreak{}}\par

\par\medskip\noindent\hypertarget{app-E7II-XXII-010.2}{}\textbf{E7II-XXII-010.2}\quad printed folio 412.\par
\noindent Printed reading: {\ttfamily\footnotesize a\allowbreak{}}\par
\noindent Corrected reading: {\ttfamily\footnotesize b\allowbreak{}}\par
\noindent Source: {\ttfamily\footnotesize fr/sga7ii/\{diplomati\allowbreak{}c,current\}/expose\_XX\allowbreak{}II\_body.tex;\allowbreak{} folio\allowbreak{} 412;\allowbreak{} formula\allowbreak{} 247;\allowbreak{} port/PORT\_EN7II\_LOG.\allowbreak{}jsonl\#C899;\allowbreak{} fixes/FIXLOG\_7II\_CUR\allowbreak{}RENT.jsonl\#C899;\allowbreak{} 7II:888\allowbreak{}}\par

\par\medskip\noindent\hypertarget{app-E7II-XXII-010.3}{}\textbf{E7II-XXII-010.3}\quad printed folio 412.\par
\noindent Printed reading: {\ttfamily\footnotesize b\allowbreak{}}\par
\noindent Corrected reading: {\ttfamily\footnotesize a\allowbreak{}}\par
\noindent Source: {\ttfamily\footnotesize fr/sga7ii/\{diplomati\allowbreak{}c,current\}/expose\_XX\allowbreak{}II\_body.tex;\allowbreak{} folio\allowbreak{} 412;\allowbreak{} formula\allowbreak{} 247;\allowbreak{} port/PORT\_EN7II\_LOG.\allowbreak{}jsonl\#C899;\allowbreak{} fixes/FIXLOG\_7II\_CUR\allowbreak{}RENT.jsonl\#C899;\allowbreak{} 7II:888\allowbreak{}}\par

\par\medskip\noindent\hypertarget{app-E7II-XXII-011.1}{}\textbf{E7II-XXII-011.1}\quad printed folio 413.\par
\noindent Printed reading: \emph{(empty)}\par
\noindent Corrected reading: {\ttfamily\footnotesize \{i\_0<\textbackslash{}cdots<i\_\allowbreak{}}\par
\noindent Source: {\ttfamily\footnotesize fr/sga7ii/\{diplomati\allowbreak{}c,current\}/expose\_XX\allowbreak{}II\_body.tex;\allowbreak{} folio\allowbreak{} 413\allowbreak{}}\par

\par\medskip\noindent\hypertarget{app-E7II-XXII-011.2}{}\textbf{E7II-XXII-011.2}\quad printed folio 413.\par
\noindent Printed reading: \emph{(empty)}\par
\noindent Corrected reading: {\ttfamily\footnotesize \}\allowbreak{}}\par
\noindent Source: {\ttfamily\footnotesize fr/sga7ii/\{diplomati\allowbreak{}c,current\}/expose\_XX\allowbreak{}II\_body.tex;\allowbreak{} folio\allowbreak{} 413\allowbreak{}}\par

\par\medskip\noindent\hypertarget{app-E7II-XXII-012.1}{}\textbf{E7II-XXII-012.1}\quad printed folio 414.\par
\noindent Printed reading: \emph{(empty)}\par
\noindent Corrected reading: {\ttfamily\footnotesize 4.\allowbreak{}}\par
\noindent Source: {\ttfamily\footnotesize ED-EN-NONMATH-G043;\allowbreak{} F7II-XXII-030.1;\allowbreak{} 7II:890;\allowbreak{} fr/sga7ii/\{diplomati\allowbreak{}c,current\}/expose\_XX\allowbreak{}II\_body.tex;\allowbreak{} folio\allowbreak{} 414\allowbreak{}}\par

\par\medskip\noindent\hypertarget{app-E7II-XXII-012.2}{}\textbf{E7II-XXII-012.2}\quad printed folio 414.\par
\noindent Printed reading: {\ttfamily\footnotesize .3\allowbreak{}}\par
\noindent Corrected reading: \emph{(empty)}\par
\noindent Source: {\ttfamily\footnotesize ED-EN-NONMATH-G043;\allowbreak{} F7II-XXII-030.1;\allowbreak{} 7II:890;\allowbreak{} fr/sga7ii/\{diplomati\allowbreak{}c,current\}/expose\_XX\allowbreak{}II\_body.tex;\allowbreak{} folio\allowbreak{} 414\allowbreak{}}\par

\par\medskip\noindent\hypertarget{app-E7II-XXII-012.3}{}\textbf{E7II-XXII-012.3}\quad printed folio 414.\par
\noindent Printed reading: \emph{(empty)}\par
\noindent Corrected reading: {\ttfamily\footnotesize 4.\allowbreak{}}\par
\noindent Source: {\ttfamily\footnotesize ED-EN-NONMATH-G044;\allowbreak{} F7II-XXII-030.8;\allowbreak{} WITNESS-F7II-XXII-03\allowbreak{}0.8;\allowbreak{} fr/sga7ii/\{diplomati\allowbreak{}c,current\}/expose\_XX\allowbreak{}II\_body.tex;\allowbreak{} folio\allowbreak{} 414\allowbreak{}}\par

\par\medskip\noindent\hypertarget{app-E7II-XXII-012.4}{}\textbf{E7II-XXII-012.4}\quad printed folio 414.\par
\noindent Printed reading: {\ttfamily\footnotesize .3\allowbreak{}}\par
\noindent Corrected reading: \emph{(empty)}\par
\noindent Source: {\ttfamily\footnotesize ED-EN-NONMATH-G044;\allowbreak{} F7II-XXII-030.8;\allowbreak{} WITNESS-F7II-XXII-03\allowbreak{}0.8;\allowbreak{} fr/sga7ii/\{diplomati\allowbreak{}c,current\}/expose\_XX\allowbreak{}II\_body.tex;\allowbreak{} folio\allowbreak{} 414\allowbreak{}}\par

\par\medskip\noindent\hypertarget{app-E7II-XXII-013.1}{}\textbf{E7II-XXII-013.1}\quad printed folio 414.\par
\noindent Printed reading: {\ttfamily\footnotesize \textbackslash{}chi\allowbreak{}}\par
\noindent Corrected reading: {\ttfamily\footnotesize x\allowbreak{}}\par
\noindent Source: {\ttfamily\footnotesize fr/sga7ii/\{diplomati\allowbreak{}c,current\}/expose\_XX\allowbreak{}II\_body.tex;\allowbreak{} folio\allowbreak{} 414;\allowbreak{} formula\allowbreak{} 289\allowbreak{}}\par

\par\medskip\noindent\hypertarget{app-E7II-XXII-013.2}{}\textbf{E7II-XXII-013.2}\quad printed folio 414.\par
\noindent Printed reading: {\ttfamily\footnotesize \textbackslash{}chi\allowbreak{}}\par
\noindent Corrected reading: {\ttfamily\footnotesize x\allowbreak{}}\par
\noindent Source: {\ttfamily\footnotesize fr/sga7ii/\{diplomati\allowbreak{}c,current\}/expose\_XX\allowbreak{}II\_body.tex;\allowbreak{} folio\allowbreak{} 414;\allowbreak{} formula\allowbreak{} 289\allowbreak{}}\par

\par\medskip\noindent\hypertarget{app-E7II-XXII-013.3}{}\textbf{E7II-XXII-013.3}\quad printed folio 414.\par
\noindent Printed reading: {\ttfamily\footnotesize \textbackslash{}chi\allowbreak{}}\par
\noindent Corrected reading: {\ttfamily\footnotesize x\allowbreak{}}\par
\noindent Source: {\ttfamily\footnotesize fr/sga7ii/\{diplomati\allowbreak{}c,current\}/expose\_XX\allowbreak{}II\_body.tex;\allowbreak{} folio\allowbreak{} 414;\allowbreak{} formula\allowbreak{} 290\allowbreak{}}\par

\par\medskip\noindent\hypertarget{app-E7II-XXII-013.4}{}\textbf{E7II-XXII-013.4}\quad printed folio 414.\par
\noindent Printed reading: {\ttfamily\footnotesize \textbackslash{}chi\allowbreak{}}\par
\noindent Corrected reading: {\ttfamily\footnotesize x\allowbreak{}}\par
\noindent Source: {\ttfamily\footnotesize fr/sga7ii/\{diplomati\allowbreak{}c,current\}/expose\_XX\allowbreak{}II\_body.tex;\allowbreak{} folio\allowbreak{} 414;\allowbreak{} formula\allowbreak{} 290\allowbreak{}}\par

\par\medskip\noindent\hypertarget{app-E7II-XXII-013.5}{}\textbf{E7II-XXII-013.5}\quad printed folio 414.\par
\noindent Printed reading: {\ttfamily\footnotesize \textbackslash{}chi\allowbreak{}}\par
\noindent Corrected reading: {\ttfamily\footnotesize x\allowbreak{}}\par
\noindent Source: {\ttfamily\footnotesize fr/sga7ii/\{diplomati\allowbreak{}c,current\}/expose\_XX\allowbreak{}II\_body.tex;\allowbreak{} folio\allowbreak{} 414;\allowbreak{} formula\allowbreak{} 290\allowbreak{}}\par

\par\medskip\noindent\hypertarget{app-E7II-XXII-013.6}{}\textbf{E7II-XXII-013.6}\quad printed folio 415.\par
\noindent Printed reading: {\ttfamily\footnotesize X\allowbreak{}}\par
\noindent Corrected reading: {\ttfamily\footnotesize x\allowbreak{}}\par
\noindent Source: {\ttfamily\footnotesize fr/sga7ii/\{diplomati\allowbreak{}c,current\}/expose\_XX\allowbreak{}II\_body.tex;\allowbreak{} folio\allowbreak{} 415;\allowbreak{} formula\allowbreak{} 291\allowbreak{}}\par

\par\medskip\noindent\hypertarget{app-E7II-XXII-014.1}{}\textbf{E7II-XXII-014.1}\quad printed folio 415.\par
\noindent Printed reading: \emph{(empty)}\par
\noindent Corrected reading: {\ttfamily\footnotesize /p\allowbreak{}}\par
\noindent Source: {\ttfamily\footnotesize port/PORT\_EN7II\_LOG.\allowbreak{}jsonl\#C900\#C900;\allowbreak{} 7II:896;\allowbreak{} fr/sga7ii/\{diplomati\allowbreak{}c,current\}/expose\_XX\allowbreak{}II\_body.tex;\allowbreak{} folio\allowbreak{} 415;\allowbreak{} formula\allowbreak{} 302\allowbreak{}}\par

\par\medskip\noindent\hypertarget{app-E7II-XXII-014.2}{}\textbf{E7II-XXII-014.2}\quad printed folio 415.\par
\noindent Printed reading: {\ttfamily\footnotesize .\allowbreak{}}\par
\noindent Corrected reading: {\ttfamily\footnotesize )\allowbreak{}}\par
\noindent Source: {\ttfamily\footnotesize port/PORT\_EN7II\_LOG.\allowbreak{}jsonl\#C900\#C900;\allowbreak{} 7II:896;\allowbreak{} fr/sga7ii/\{diplomati\allowbreak{}c,current\}/expose\_XX\allowbreak{}II\_body.tex;\allowbreak{} folio\allowbreak{} 415;\allowbreak{} formula\allowbreak{} 302\allowbreak{}}\par

\par\medskip\noindent\hypertarget{app-E7II-XXII-014.3}{}\textbf{E7II-XXII-014.3}\quad printed folio 415.\par
\noindent Printed reading: {\ttfamily\footnotesize H\allowbreak{}}\par
\noindent Corrected reading: {\ttfamily\footnotesize F\allowbreak{}}\par
\noindent Source: {\ttfamily\footnotesize fr/sga7ii/\{diplomati\allowbreak{}c,current\}/expose\_XX\allowbreak{}II\_body.tex;\allowbreak{} folio\allowbreak{} 415;\allowbreak{} formula\allowbreak{} 302\allowbreak{}}\par

\par\medskip\noindent\hypertarget{app-E7II-XXII-015.1}{}\textbf{E7II-XXII-015.1}\quad printed folio 416.\par
\noindent Printed reading: \emph{(empty)}\par
\noindent Corrected reading: {\ttfamily\footnotesize \textbackslash{}mathbb\allowbreak{} }\par
\noindent Source: {\ttfamily\footnotesize ED-EN-FONT-G008;\allowbreak{} F7II-XXII-034.1;\allowbreak{} fr/sga7ii/\{diplomati\allowbreak{}c,current\}/expose\_XX\allowbreak{}II\_body.tex;\allowbreak{} folio\allowbreak{} 416\allowbreak{}}\par

\par\medskip\noindent\hypertarget{app-E7II-XXII-016.1}{}\textbf{E7II-XXII-016.1}\quad printed folio 416.\par
\noindent Printed reading: \emph{(empty)}\par
\noindent Corrected reading: {\ttfamily\footnotesize \textbackslash{}mathbb\allowbreak{}}\par
\noindent Source: {\ttfamily\footnotesize ED-EN-FONT-G009;\allowbreak{} F7II-XXII-035.1;\allowbreak{} fr/sga7ii/\{diplomati\allowbreak{}c,current\}/expose\_XX\allowbreak{}II\_body.tex;\allowbreak{} folio\allowbreak{} 416\allowbreak{}}\par

\par\medskip\noindent\hypertarget{app-E7II-XXII-017.1}{}\textbf{E7II-XXII-017.1}\quad printed folio 416.\par
\noindent Printed reading: {\ttfamily\footnotesize A\allowbreak{}}\par
\noindent Corrected reading: {\ttfamily\footnotesize B\allowbreak{}}\par
\noindent Source: {\ttfamily\footnotesize fr/sga7ii/\{diplomati\allowbreak{}c,current\}/expose\_XX\allowbreak{}II\_body.tex;\allowbreak{} folio\allowbreak{} 416;\allowbreak{} formula\allowbreak{} 310\allowbreak{}}\par

\par\medskip\noindent\hypertarget{app-E7II-XXII-018.1}{}\textbf{E7II-XXII-018.1}\quad printed folio 416.\par
\noindent Printed reading: {\ttfamily\footnotesize q\allowbreak{}}\par
\noindent Corrected reading: {\ttfamily\footnotesize p\allowbreak{}}\par
\noindent Source: {\ttfamily\footnotesize fr/sga7ii/\{diplomati\allowbreak{}c,current\}/expose\_XX\allowbreak{}II\_body.tex;\allowbreak{} folio\allowbreak{} 416;\allowbreak{} formula\allowbreak{} 316\allowbreak{}}\par

\par\medskip\noindent\hypertarget{app-E7II-XXII-019.1}{}\textbf{E7II-XXII-019.1}\quad printed folio 417.\par
\noindent Printed reading: \emph{(empty)}\par
\noindent Corrected reading: {\ttfamily\footnotesize p\textasciicircum{}\allowbreak{}}\par
\noindent Source: {\ttfamily\footnotesize fr/sga7ii/\{diplomati\allowbreak{}c,current\}/expose\_XX\allowbreak{}II\_body.tex;\allowbreak{} folio\allowbreak{} 417;\allowbreak{} formula\allowbreak{} 319\allowbreak{}}\par

\par\medskip\noindent\hypertarget{app-E7II-XXII-020.1}{}\textbf{E7II-XXII-020.1}\quad printed folio 417.\par
\noindent Printed reading: {\ttfamily\footnotesize q\allowbreak{}}\par
\noindent Corrected reading: {\ttfamily\footnotesize p\allowbreak{}}\par
\noindent Source: {\ttfamily\footnotesize fr/sga7ii/\{diplomati\allowbreak{}c,current\}/expose\_XX\allowbreak{}II\_body.tex;\allowbreak{} folio\allowbreak{} 417;\allowbreak{} formula\allowbreak{} 329\allowbreak{}}\par

\par\medskip\noindent\hypertarget{app-E7II-XXII-020.2}{}\textbf{E7II-XXII-020.2}\quad printed folio 417.\par
\noindent Printed reading: {\ttfamily\footnotesize q\allowbreak{}}\par
\noindent Corrected reading: {\ttfamily\footnotesize p\allowbreak{}}\par
\noindent Source: {\ttfamily\footnotesize fr/sga7ii/\{diplomati\allowbreak{}c,current\}/expose\_XX\allowbreak{}II\_body.tex;\allowbreak{} folio\allowbreak{} 417;\allowbreak{} formula\allowbreak{} 331\allowbreak{}}\par

\par\medskip\noindent\hypertarget{app-E7II-XXII-021.1}{}\textbf{E7II-XXII-021.1}\quad printed folio 418.\par
\noindent Printed reading: \emph{(empty)}\par
\noindent Corrected reading: {\ttfamily\footnotesize \textbackslash{}mathbb\allowbreak{} }\par
\noindent Source: {\ttfamily\footnotesize ED-EN-FONT-G010;\allowbreak{} F7II-XXII-039.1;\allowbreak{} fr/sga7ii/\{diplomati\allowbreak{}c,current\}/expose\_XX\allowbreak{}II\_body.tex;\allowbreak{} folio\allowbreak{} 418\allowbreak{}}\par

\par\medskip\noindent\hypertarget{app-E7II-XXII-022.1}{}\textbf{E7II-XXII-022.1}\quad printed folio 418.\par
\noindent Printed reading: {\ttfamily\footnotesize X\allowbreak{}}\par
\noindent Corrected reading: {\ttfamily\footnotesize \textbackslash{}chi\allowbreak{}}\par
\noindent Source: {\ttfamily\footnotesize fr/sga7ii/\{diplomati\allowbreak{}c,current\}/expose\_XX\allowbreak{}II\_body.tex;\allowbreak{} folio\allowbreak{} 418;\allowbreak{} formula\allowbreak{} 345\allowbreak{}}\par

\par\medskip\noindent\hypertarget{app-E7II-XXII-023.1}{}\textbf{E7II-XXII-023.1}\quad printed folio 419.\par
\noindent Printed reading: \emph{(empty)}\par
\noindent Corrected reading: {\ttfamily\footnotesize \_1\allowbreak{}}\par
\noindent Source: {\ttfamily\footnotesize ED-EN-PORT-WITNESS-7\allowbreak{}II-C819;\allowbreak{} fr/sga7ii/diplomatic\allowbreak{}/expose\_XXII\_body.te\allowbreak{}x;\allowbreak{} folio\allowbreak{} 419;\allowbreak{} port/PORT\_EN7II\_LOG.\allowbreak{}jsonl\#C819\#C819;\allowbreak{} ;\allowbreak{} fr/sga7ii/\{diplomati\allowbreak{}c,current\}/expose\_XX\allowbreak{}II\_body.tex;\allowbreak{} folio\allowbreak{} 419;\allowbreak{} formula\allowbreak{} 353\allowbreak{}}\par

\par\medskip\noindent\hypertarget{app-E7II-XXII-024.1}{}\textbf{E7II-XXII-024.1}\quad printed folio 421.\par
\noindent Printed reading: \emph{(empty)}\par
\noindent Corrected reading: {\ttfamily\footnotesize \textasciicircum{}\textbackslash{}nu\allowbreak{}}\par
\noindent Source: {\ttfamily\footnotesize fr/sga7ii/\{diplomati\allowbreak{}c,current\}/expose\_XX\allowbreak{}II\_body.tex;\allowbreak{} folio\allowbreak{} 421;\allowbreak{} formula\allowbreak{} 384\allowbreak{}}\par

\par\medskip\noindent\hypertarget{app-E7II-XXII-024.2}{}\textbf{E7II-XXII-024.2}\quad printed folio 421.\par
\noindent Printed reading: \emph{(empty)}\par
\noindent Corrected reading: {\ttfamily\footnotesize \textasciicircum{}\textbackslash{}nu\allowbreak{}}\par
\noindent Source: {\ttfamily\footnotesize fr/sga7ii/\{diplomati\allowbreak{}c,current\}/expose\_XX\allowbreak{}II\_body.tex;\allowbreak{} folio\allowbreak{} 421;\allowbreak{} formula\allowbreak{} 384\allowbreak{}}\par

\par\medskip\noindent\hypertarget{app-E7II-XXII-024.3}{}\textbf{E7II-XXII-024.3}\quad printed folio 421.\par
\noindent Printed reading: \emph{(empty)}\par
\noindent Corrected reading: {\ttfamily\footnotesize \_k\allowbreak{}}\par
\noindent Source: {\ttfamily\footnotesize fr/sga7ii/\{diplomati\allowbreak{}c,current\}/expose\_XX\allowbreak{}II\_body.tex;\allowbreak{} folio\allowbreak{} 421;\allowbreak{} formula\allowbreak{} 388\allowbreak{}}\par

\par\medskip\noindent\hypertarget{app-E7II-XXII-025.1}{}\textbf{E7II-XXII-025.1}\quad printed folio 421.\par
\noindent Printed reading: \emph{(empty)}\par
\noindent Corrected reading: {\ttfamily\footnotesize \textasciicircum{}\textbackslash{}nu\allowbreak{}}\par
\noindent Source: {\ttfamily\footnotesize fr/sga7ii/\{diplomati\allowbreak{}c,current\}/expose\_XX\allowbreak{}II\_body.tex;\allowbreak{} folio\allowbreak{} 421;\allowbreak{} formula\allowbreak{} 391\allowbreak{}}\par

\par\medskip\noindent\hypertarget{app-E7II-XXII-025.2}{}\textbf{E7II-XXII-025.2}\quad printed folio 421.\par
\noindent Printed reading: \emph{(empty)}\par
\noindent Corrected reading: {\ttfamily\footnotesize \textasciicircum{}\textbackslash{}nu\allowbreak{}}\par
\noindent Source: {\ttfamily\footnotesize fr/sga7ii/\{diplomati\allowbreak{}c,current\}/expose\_XX\allowbreak{}II\_body.tex;\allowbreak{} folio\allowbreak{} 421;\allowbreak{} formula\allowbreak{} 391\allowbreak{}}\par

\par\medskip\noindent\hypertarget{app-E7II-XXII-026.1}{}\textbf{E7II-XXII-026.1}\quad printed folio 422.\par
\noindent Printed reading: \emph{(empty)}\par
\noindent Corrected reading: {\ttfamily\footnotesize 4.\allowbreak{}}\par
\noindent Source: {\ttfamily\footnotesize ED-EN-NORM-G034;\allowbreak{} complete\allowbreak{} normalized-equal\allowbreak{} formula\allowbreak{} audit;\allowbreak{} fr/sga7ii/\{diplomati\allowbreak{}c,current\}/expose\_XX\allowbreak{}II\_body.tex;\allowbreak{} folio\allowbreak{} 422\allowbreak{}}\par

\par\medskip\noindent\hypertarget{app-E7II-XXII-026.2}{}\textbf{E7II-XXII-026.2}\quad printed folio 422.\par
\noindent Printed reading: {\ttfamily\footnotesize .4\allowbreak{}}\par
\noindent Corrected reading: \emph{(empty)}\par
\noindent Source: {\ttfamily\footnotesize ED-EN-NORM-G034;\allowbreak{} complete\allowbreak{} normalized-equal\allowbreak{} formula\allowbreak{} audit;\allowbreak{} fr/sga7ii/\{diplomati\allowbreak{}c,current\}/expose\_XX\allowbreak{}II\_body.tex;\allowbreak{} folio\allowbreak{} 422\allowbreak{}}\par

\par\medskip\noindent\hypertarget{app-E7II-XXII-027.1}{}\textbf{E7II-XXII-027.1}\quad printed folio 422.\par
\noindent Printed reading: {\ttfamily\footnotesize .\allowbreak{}}\par
\noindent Corrected reading: {\ttfamily\footnotesize )\allowbreak{}}\par
\noindent Source: {\ttfamily\footnotesize fr/sga7ii/\{diplomati\allowbreak{}c,current\}/expose\_XX\allowbreak{}II\_body.tex;\allowbreak{} folio\allowbreak{} 422;\allowbreak{} formula\allowbreak{} 399\allowbreak{}}\par

\par\medskip\noindent\hypertarget{app-E7II-XXII-028.1}{}\textbf{E7II-XXII-028.1}\quad printed folio 423.\par
\noindent Printed reading: {\ttfamily\footnotesize \textbackslash{}subseteq\allowbreak{}}\par
\noindent Corrected reading: {\ttfamily\footnotesize \textbackslash{}supset\allowbreak{}}\par
\noindent Source: {\ttfamily\footnotesize fr/sga7ii/\{diplomati\allowbreak{}c,current\}/expose\_XX\allowbreak{}II\_body.tex;\allowbreak{} folio\allowbreak{} 423;\allowbreak{} formula\allowbreak{} 408\allowbreak{}}\par

\par\medskip\noindent\hypertarget{app-E7II-XXII-029.1}{}\textbf{E7II-XXII-029.1}\quad printed folio 423.\par
\noindent Printed reading: {\ttfamily\footnotesize 2.3\allowbreak{}}\par
\noindent Corrected reading: {\ttfamily\footnotesize 23\allowbreak{}}\par
\noindent Source: {\ttfamily\footnotesize ED-EN-NORM-G035;\allowbreak{} complete\allowbreak{} normalized-equal\allowbreak{} formula\allowbreak{} audit;\allowbreak{} fr/sga7ii/\{diplomati\allowbreak{}c,current\}/expose\_XX\allowbreak{}II\_body.tex;\allowbreak{} folio\allowbreak{} 423\allowbreak{}}\par

\par\medskip\noindent\hypertarget{app-E7II-XXII-030.1}{}\textbf{E7II-XXII-030.1}\quad printed folio 424.\par
\noindent Printed reading: {\ttfamily\footnotesize car\allowbreak{}}\par
\noindent Corrected reading: {\ttfamily\footnotesize card\allowbreak{}}\par
\noindent Source: {\ttfamily\footnotesize fr/sga7ii/\{diplomati\allowbreak{}c,current\}/expose\_XX\allowbreak{}II\_body.tex;\allowbreak{} folio\allowbreak{} 424;\allowbreak{} formula\allowbreak{} 431\allowbreak{}}\par

\par\medskip\noindent\hypertarget{app-E7II-XXII-030.2}{}\textbf{E7II-XXII-030.2}\quad printed folio 424.\par
\noindent Printed reading: {\ttfamily\footnotesize )\allowbreak{}}\par
\noindent Corrected reading: \emph{(empty)}\par
\noindent Source: {\ttfamily\footnotesize fr/sga7ii/\{diplomati\allowbreak{}c,current\}/expose\_XX\allowbreak{}II\_body.tex;\allowbreak{} folio\allowbreak{} 424;\allowbreak{} formula\allowbreak{} 431\allowbreak{}}\par

\par\medskip\noindent\hypertarget{app-E7II-XXII-031.1}{}\textbf{E7II-XXII-031.1}\quad printed folio 425.\par
\noindent Printed reading: {\ttfamily\footnotesize 23\allowbreak{}}\par
\noindent Corrected reading: {\ttfamily\footnotesize 33\allowbreak{}}\par
\noindent Source: {\ttfamily\footnotesize ED-EN-NORM-G036;\allowbreak{} complete\allowbreak{} normalized-equal\allowbreak{} formula\allowbreak{} audit;\allowbreak{} fr/sga7ii/\{diplomati\allowbreak{}c,current\}/expose\_XX\allowbreak{}II\_body.tex;\allowbreak{} folio\allowbreak{} 425\allowbreak{}}\par

\par\medskip\noindent\hypertarget{app-E7II-XXII-032.1}{}\textbf{E7II-XXII-032.1}\quad printed folio 426.\par
\noindent Printed reading: {\ttfamily\footnotesize ordinal\allowbreak{}}\par
\noindent Corrected reading: {\ttfamily\footnotesize order\allowbreak{}}\par
\noindent Source: {\ttfamily\footnotesize fr/sga7ii/\{diplomati\allowbreak{}c,current\}/expose\_XX\allowbreak{}II\_body.tex;\allowbreak{} folio\allowbreak{} 426\allowbreak{}}\par

\par\medskip\noindent\hypertarget{app-E7II-XXII-033.1}{}\textbf{E7II-XXII-033.1}\quad printed folio 427.\par
\noindent Printed reading: \emph{(empty)}\par
\noindent Corrected reading: {\ttfamily\footnotesize \_0\allowbreak{}}\par
\noindent Source: {\ttfamily\footnotesize fr/sga7ii/\{diplomati\allowbreak{}c,current\}/expose\_XX\allowbreak{}II\_body.tex;\allowbreak{} folio\allowbreak{} 427;\allowbreak{} formula\allowbreak{} 468\allowbreak{}}\par

\par\medskip\noindent\hypertarget{app-E7II-XXII-034.1}{}\textbf{E7II-XXII-034.1}\quad printed folio 428.\par
\noindent Printed reading: {\ttfamily\footnotesize 7\allowbreak{}}\par
\noindent Corrected reading: {\ttfamily\footnotesize 4\allowbreak{}}\par
\noindent Source: {\ttfamily\footnotesize ED-EN-NORM-G037;\allowbreak{} complete\allowbreak{} normalized-equal\allowbreak{} formula\allowbreak{} audit;\allowbreak{} fr/sga7ii/\{diplomati\allowbreak{}c,current\}/expose\_XX\allowbreak{}II\_body.tex;\allowbreak{} folio\allowbreak{} 428\allowbreak{}}\par

\par\medskip\noindent\hypertarget{app-E7II-XXII-035.1}{}\textbf{E7II-XXII-035.1}\quad printed folio 429.\par
\noindent Printed reading: {\ttfamily\footnotesize 4\allowbreak{}}\par
\noindent Corrected reading: {\ttfamily\footnotesize 5\allowbreak{}}\par
\noindent Source: {\ttfamily\footnotesize port/PORT\_EN7II\_LOG.\allowbreak{}jsonl\#C768;\allowbreak{} fixes/FIXLOG\_7II\_CUR\allowbreak{}RENT.jsonl\#C768;\allowbreak{} 7II:740\allowbreak{}}\par

\par\medskip\noindent\hypertarget{app-E7II-XXII-036.1}{}\textbf{E7II-XXII-036.1}\quad printed folio 429.\par
\noindent Printed reading: {\ttfamily\footnotesize t\allowbreak{}}\par
\noindent Corrected reading: {\ttfamily\footnotesize qt\allowbreak{}}\par
\noindent Source: {\ttfamily\footnotesize ED-EN-PORT-WITNESS-7\allowbreak{}II-C901;\allowbreak{} fr/sga7ii/diplomatic\allowbreak{}/expose\_XXII\_body.te\allowbreak{}x;\allowbreak{} folio\allowbreak{} 429;\allowbreak{} port/PORT\_EN7II\_LOG.\allowbreak{}jsonl\#C901\#C901;\allowbreak{} 7II:881;\allowbreak{} fr/sga7ii/\{diplomati\allowbreak{}c,current\}/expose\_XX\allowbreak{}II\_body.tex;\allowbreak{} folio\allowbreak{} 429;\allowbreak{} formula\allowbreak{} 495\allowbreak{}}\par

\par\medskip\noindent\hypertarget{app-E7II-XXII-036.2}{}\textbf{E7II-XXII-036.2}\quad printed folio 429.\par
\noindent Printed reading: {\ttfamily\footnotesize qt\allowbreak{}}\par
\noindent Corrected reading: {\ttfamily\footnotesize t\allowbreak{}}\par
\noindent Source: {\ttfamily\footnotesize ED-EN-PORT-WITNESS-7\allowbreak{}II-C901;\allowbreak{} fr/sga7ii/diplomatic\allowbreak{}/expose\_XXII\_body.te\allowbreak{}x;\allowbreak{} folio\allowbreak{} 429;\allowbreak{} port/PORT\_EN7II\_LOG.\allowbreak{}jsonl\#C901\#C901;\allowbreak{} 7II:881;\allowbreak{} fr/sga7ii/\{diplomati\allowbreak{}c,current\}/expose\_XX\allowbreak{}II\_body.tex;\allowbreak{} folio\allowbreak{} 429;\allowbreak{} formula\allowbreak{} 495\allowbreak{}}\par

\par\medskip\noindent\hypertarget{app-E7II-XXII-037.1}{}\textbf{E7II-XXII-037.1}\quad printed folio 429.\par
\noindent Printed reading: {\ttfamily\footnotesize if\allowbreak{}}\par
\noindent Corrected reading: {\ttfamily\footnotesize it\allowbreak{}}\par
\noindent Source: {\ttfamily\footnotesize ED-EN-PORT-WITNESS-7\allowbreak{}II-C901;\allowbreak{} fr/sga7ii/diplomatic\allowbreak{}/expose\_XXII\_body.te\allowbreak{}x;\allowbreak{} folio\allowbreak{} 429;\allowbreak{} ED-EN-FORMULA-G245;\allowbreak{} fr/sga7ii/\{diplomati\allowbreak{}c,current\}/expose\_XX\allowbreak{}II\_body.tex;\allowbreak{} folio\allowbreak{} 429;\allowbreak{} port/PORT\_EN7II\_LOG.\allowbreak{}jsonl\#C901\#C901;\allowbreak{} 7II:881\allowbreak{}}\par

\par\medskip\noindent\hypertarget{app-E7II-XXII-037.2}{}\textbf{E7II-XXII-037.2}\quad printed folio 429.\par
\noindent Printed reading: \emph{(empty)}\par
\noindent Corrected reading: {\ttfamily\footnotesize \textasciicircum{}\{-1\}\allowbreak{}}\par
\noindent Source: {\ttfamily\footnotesize ED-EN-PORT-WITNESS-7\allowbreak{}II-C901;\allowbreak{} fr/sga7ii/diplomatic\allowbreak{}/expose\_XXII\_body.te\allowbreak{}x;\allowbreak{} folio\allowbreak{} 429;\allowbreak{} ED-EN-FORMULA-G245;\allowbreak{} fr/sga7ii/\{diplomati\allowbreak{}c,current\}/expose\_XX\allowbreak{}II\_body.tex;\allowbreak{} folio\allowbreak{} 429;\allowbreak{} port/PORT\_EN7II\_LOG.\allowbreak{}jsonl\#C901\#C901;\allowbreak{} 7II:881\allowbreak{}}\par

\par\medskip\noindent\hypertarget{app-E7II-XXII-038.1}{}\textbf{E7II-XXII-038.1}\quad printed folio 429.\par
\noindent Printed reading: {\ttfamily\footnotesize \textbackslash{}Lambda\allowbreak{}}\par
\noindent Corrected reading: {\ttfamily\footnotesize \textbackslash{}Delta\allowbreak{}}\par
\noindent Source: {\ttfamily\footnotesize fr/sga7ii/\{diplomati\allowbreak{}c,current\}/expose\_XX\allowbreak{}II\_body.tex;\allowbreak{} folio\allowbreak{} 429;\allowbreak{} formula\allowbreak{} 501\allowbreak{}}\par

\par\medskip\noindent\hypertarget{app-E7II-XXII-038.2}{}\textbf{E7II-XXII-038.2}\quad printed folio 429.\par
\noindent Printed reading: {\ttfamily\footnotesize 4\allowbreak{}}\par
\noindent Corrected reading: {\ttfamily\footnotesize 6\allowbreak{}}\par
\noindent Source: {\ttfamily\footnotesize ED-EN-NONMATH-G045;\allowbreak{} F7II-XXII-059.2;\allowbreak{} 7II:884;\allowbreak{} fr/sga7ii/\{diplomati\allowbreak{}c,current\}/expose\_XX\allowbreak{}II\_body.tex;\allowbreak{} folio\allowbreak{} 429\allowbreak{}}\par

\par\medskip\noindent\hypertarget{app-E7II-XXII-039.1}{}\textbf{E7II-XXII-039.1}\quad printed folio 430.\par
\noindent Printed reading: {\ttfamily\footnotesize .\allowbreak{}}\par
\noindent Corrected reading: {\ttfamily\footnotesize )\allowbreak{}}\par
\noindent Source: {\ttfamily\footnotesize fr/sga7ii/\{diplomati\allowbreak{}c,current\}/expose\_XX\allowbreak{}II\_body.tex;\allowbreak{} folio\allowbreak{} 430;\allowbreak{} formula\allowbreak{} 514\allowbreak{}}\par

\par\medskip\noindent\hypertarget{app-E7II-XXII-040.1}{}\textbf{E7II-XXII-040.1}\quad printed folio 431.\par
\noindent Printed reading: {\ttfamily\footnotesize 24\allowbreak{}}\par
\noindent Corrected reading: {\ttfamily\footnotesize 33\allowbreak{}}\par
\noindent Source: {\ttfamily\footnotesize ED-EN-NONMATH-G012;\allowbreak{} F7II-XXII-062.1;\allowbreak{} ADOPT-F7II-XXII-062.\allowbreak{}1;\allowbreak{} fr/sga7ii/\{diplomati\allowbreak{}c,current\}/expose\_XX\allowbreak{}II\_body.tex;\allowbreak{} folio\allowbreak{} 431\allowbreak{}}\par

\par\medskip\noindent\hypertarget{app-E7II-XXII-041.1}{}\textbf{E7II-XXII-041.1}\quad printed folio 431.\par
\noindent Printed reading: {\ttfamily\footnotesize B\allowbreak{}}\par
\noindent Corrected reading: {\ttfamily\footnotesize S\allowbreak{}}\par
\noindent Source: {\ttfamily\footnotesize fr/sga7ii/\{diplomati\allowbreak{}c,current\}/expose\_XX\allowbreak{}II\_body.tex;\allowbreak{} folio\allowbreak{} 431;\allowbreak{} formula\allowbreak{} 543\allowbreak{}}\par

\par\medskip\noindent\hypertarget{app-E7II-XXII-042.1}{}\textbf{E7II-XXII-042.1}\quad printed folio 432.\par
\noindent Printed reading: {\ttfamily\footnotesize t\allowbreak{}}\par
\noindent Corrected reading: {\ttfamily\footnotesize i\allowbreak{}}\par
\noindent Source: {\ttfamily\footnotesize fr/sga7ii/\{diplomati\allowbreak{}c,current\}/expose\_XX\allowbreak{}II\_body.tex;\allowbreak{} folio\allowbreak{} 432\allowbreak{}}\par

\par\medskip\noindent\hypertarget{app-E7II-XXII-043.1}{}\textbf{E7II-XXII-043.1}\quad printed folio 433.\par
\noindent Printed reading: \emph{(empty)}\par
\noindent Corrected reading: {\ttfamily\footnotesize -\allowbreak{}}\par
\noindent Source: {\ttfamily\footnotesize fr/sga7ii/\{diplomati\allowbreak{}c,current\}/expose\_XX\allowbreak{}II\_body.tex;\allowbreak{} folio\allowbreak{} 433;\allowbreak{} formula\allowbreak{} 556\allowbreak{}}\par

\par\medskip\noindent\hypertarget{app-E7II-XXII-044.1}{}\textbf{E7II-XXII-044.1}\quad printed folio 434.\par
\noindent Printed reading: \emph{(empty)}\par
\noindent Corrected reading: {\ttfamily\footnotesize (\allowbreak{}}\par
\noindent Source: {\ttfamily\footnotesize fr/sga7ii/\{diplomati\allowbreak{}c,current\}/expose\_XX\allowbreak{}II\_body.tex;\allowbreak{} folio\allowbreak{} 434;\allowbreak{} formula\allowbreak{} 566\allowbreak{}}\par

\par\medskip\noindent\hypertarget{app-E7II-XXII-045.1}{}\textbf{E7II-XXII-045.1}\quad printed folio 434.\par
\noindent Printed reading: \emph{(empty)}\par
\noindent Corrected reading: {\ttfamily\footnotesize )\allowbreak{}}\par
\noindent Source: {\ttfamily\footnotesize fr/sga7ii/\{diplomati\allowbreak{}c,current\}/expose\_XX\allowbreak{}II\_body.tex;\allowbreak{} folio\allowbreak{} 434;\allowbreak{} formula\allowbreak{} 569\allowbreak{}}\par

\par\medskip\noindent\hypertarget{app-E7II-XXII-046.1}{}\textbf{E7II-XXII-046.1}\quad printed folio 435.\par
\noindent Printed reading: \emph{(empty)}\par
\noindent Corrected reading: {\ttfamily\footnotesize \textbackslash{}circ\allowbreak{}}\par
\noindent Source: {\ttfamily\footnotesize fr/sga7ii/\{diplomati\allowbreak{}c,current\}/expose\_XX\allowbreak{}II\_body.tex;\allowbreak{} folio\allowbreak{} 435;\allowbreak{} formula\allowbreak{} 572\allowbreak{}}\par

\par\medskip\noindent\hypertarget{app-E7II-XXII-046.2}{}\textbf{E7II-XXII-046.2}\quad printed folio 435.\par
\noindent Printed reading: {\ttfamily\footnotesize a\allowbreak{}}\par
\noindent Corrected reading: {\ttfamily\footnotesize q\allowbreak{}}\par
\noindent Source: {\ttfamily\footnotesize fr/sga7ii/\{diplomati\allowbreak{}c,current\}/expose\_XX\allowbreak{}II\_body.tex;\allowbreak{} folio\allowbreak{} 435;\allowbreak{} formula\allowbreak{} 572\allowbreak{}}\par

\par\medskip\noindent\hypertarget{app-E7II-XXII-046.3}{}\textbf{E7II-XXII-046.3}\quad printed folio 435.\par
\noindent Printed reading: \emph{(empty)}\par
\noindent Corrected reading: {\ttfamily\footnotesize \textbackslash{}circ\allowbreak{}}\par
\noindent Source: {\ttfamily\footnotesize fr/sga7ii/\{diplomati\allowbreak{}c,current\}/expose\_XX\allowbreak{}II\_body.tex;\allowbreak{} folio\allowbreak{} 435;\allowbreak{} formula\allowbreak{} 572\allowbreak{}}\par

\par\medskip\noindent\hypertarget{app-E7II-XXII-047.1}{}\textbf{E7II-XXII-047.1}\quad printed folio 435.\par
\noindent Printed reading: {\ttfamily\footnotesize \textbackslash{}mathcal\allowbreak{} O\allowbreak{}}\par
\noindent Corrected reading: {\ttfamily\footnotesize 0\allowbreak{}}\par
\noindent Source: {\ttfamily\footnotesize ED-EN-FONT-G021;\allowbreak{} F7II-XXII-069.1;\allowbreak{} fr/sga7ii/\{diplomati\allowbreak{}c,current\}/expose\_XX\allowbreak{}II\_body.tex;\allowbreak{} folio\allowbreak{} 435;\allowbreak{} fr/sga7ii/\{diplomati\allowbreak{}c,current\}/expose\_XX\allowbreak{}II\_body.tex;\allowbreak{} folio\allowbreak{} 435;\allowbreak{} formula\allowbreak{} 578\allowbreak{}}\par

\par\medskip\noindent\hypertarget{app-E7II-XXII-048.1}{}\textbf{E7II-XXII-048.1}\quad printed folio 435.\par
\noindent Printed reading: {\ttfamily\footnotesize \textbackslash{}mathcal\allowbreak{} O\allowbreak{}}\par
\noindent Corrected reading: {\ttfamily\footnotesize 0\allowbreak{}}\par
\noindent Source: {\ttfamily\footnotesize ED-EN-FONT-G022;\allowbreak{} F7II-XXII-070.1;\allowbreak{} fr/sga7ii/\{diplomati\allowbreak{}c,current\}/expose\_XX\allowbreak{}II\_body.tex;\allowbreak{} folio\allowbreak{} 435;\allowbreak{} fr/sga7ii/\{diplomati\allowbreak{}c,current\}/expose\_XX\allowbreak{}II\_body.tex;\allowbreak{} folio\allowbreak{} 435;\allowbreak{} formula\allowbreak{} 579\allowbreak{}}\par

\par\medskip\noindent\hypertarget{app-E7II-XXII-048.2}{}\textbf{E7II-XXII-048.2}\quad printed folio 435.\par
\noindent Printed reading: {\ttfamily\footnotesize \textbackslash{}mathcal\allowbreak{} O\allowbreak{}}\par
\noindent Corrected reading: {\ttfamily\footnotesize 0\allowbreak{}}\par
\noindent Source: {\ttfamily\footnotesize ED-EN-FONT-G023;\allowbreak{} F7II-XXII-070.2;\allowbreak{} fr/sga7ii/\{diplomati\allowbreak{}c,current\}/expose\_XX\allowbreak{}II\_body.tex;\allowbreak{} folio\allowbreak{} 435;\allowbreak{} fr/sga7ii/\{diplomati\allowbreak{}c,current\}/expose\_XX\allowbreak{}II\_body.tex;\allowbreak{} folio\allowbreak{} 435;\allowbreak{} formula\allowbreak{} 579\allowbreak{}}\par

\par\medskip\noindent\hypertarget{app-E7II-XXII-049.1}{}\textbf{E7II-XXII-049.1}\quad printed folio 436.\par
\noindent Printed reading: {\ttfamily\footnotesize .\allowbreak{}}\par
\noindent Corrected reading: {\ttfamily\footnotesize )\allowbreak{}}\par
\noindent Source: {\ttfamily\footnotesize fr/sga7ii/\{diplomati\allowbreak{}c,current\}/expose\_XX\allowbreak{}II\_body.tex;\allowbreak{} folio\allowbreak{} 436;\allowbreak{} formula\allowbreak{} 583\allowbreak{}}\par

\par\medskip\noindent\hypertarget{app-E7II-XXII-050.1}{}\textbf{E7II-XXII-050.1}\quad printed folio 436.\par
\noindent Printed reading: {\ttfamily\footnotesize .\allowbreak{}}\par
\noindent Corrected reading: {\ttfamily\footnotesize )\allowbreak{}}\par
\noindent Source: {\ttfamily\footnotesize fr/sga7ii/\{diplomati\allowbreak{}c,current\}/expose\_XX\allowbreak{}II\_body.tex;\allowbreak{} folio\allowbreak{} 436;\allowbreak{} formula\allowbreak{} 584\allowbreak{}}\par

\par\medskip\noindent\hypertarget{app-E7II-XXII-050.2}{}\textbf{E7II-XXII-050.2}\quad printed folio 436.\par
\noindent Printed reading: {\ttfamily\footnotesize (\allowbreak{}}\par
\noindent Corrected reading: \emph{(empty)}\par
\noindent Source: {\ttfamily\footnotesize fr/sga7ii/\{diplomati\allowbreak{}c,current\}/expose\_XX\allowbreak{}II\_body.tex;\allowbreak{} folio\allowbreak{} 436;\allowbreak{} formula\allowbreak{} 584\allowbreak{}}\par

\par\medskip\noindent\hypertarget{app-E7II-XXII-051.1}{}\textbf{E7II-XXII-051.1}\quad printed folio 436.\par
\noindent Printed reading: {\ttfamily\footnotesize .\allowbreak{}}\par
\noindent Corrected reading: {\ttfamily\footnotesize ,\allowbreak{}}\par
\noindent Source: {\ttfamily\footnotesize fr/sga7ii/\{diplomati\allowbreak{}c,current\}/expose\_XX\allowbreak{}II\_body.tex;\allowbreak{} folio\allowbreak{} 436;\allowbreak{} formula\allowbreak{} 585\allowbreak{}}\par

\par\medskip\noindent\hypertarget{app-E7II-XXII-051.2}{}\textbf{E7II-XXII-051.2}\quad printed folio 436.\par
\noindent Printed reading: \emph{(empty)}\par
\noindent Corrected reading: {\ttfamily\footnotesize )\allowbreak{}}\par
\noindent Source: {\ttfamily\footnotesize fr/sga7ii/\{diplomati\allowbreak{}c,current\}/expose\_XX\allowbreak{}II\_body.tex;\allowbreak{} folio\allowbreak{} 436;\allowbreak{} formula\allowbreak{} 585\allowbreak{}}\par

\par\medskip\noindent\hypertarget{app-E7II-XXII-052.1}{}\textbf{E7II-XXII-052.1}\quad printed folio 437.\par
\noindent Printed reading: {\ttfamily\footnotesize i\allowbreak{}}\par
\noindent Corrected reading: {\ttfamily\footnotesize n\allowbreak{}}\par
\noindent Source: {\ttfamily\footnotesize ED-EN-FORMULA-G247;\allowbreak{} fr/sga7ii/\{diplomati\allowbreak{}c,current\}/expose\_XX\allowbreak{}II\_body.tex;\allowbreak{} folio\allowbreak{} 437\allowbreak{}}\par

\par\medskip\noindent\hypertarget{app-E7II-XXII-052.2}{}\textbf{E7II-XXII-052.2}\quad printed folio 437.\par
\noindent Printed reading: \emph{(empty)}\par
\noindent Corrected reading: {\ttfamily\footnotesize the\allowbreak{} quotient\allowbreak{} of\allowbreak{} }\par
\noindent Source: {\ttfamily\footnotesize ED-EN-FORMULA-G247;\allowbreak{} fr/sga7ii/\{diplomati\allowbreak{}c,current\}/expose\_XX\allowbreak{}II\_body.tex;\allowbreak{} folio\allowbreak{} 437\allowbreak{}}\par

\par\medskip\noindent\hypertarget{app-E7II-XXII-052.3}{}\textbf{E7II-XXII-052.3}\quad printed folio 437.\par
\noindent Printed reading: {\ttfamily\footnotesize =\allowbreak{}}\par
\noindent Corrected reading: {\ttfamily\footnotesize \textbackslash{}quad\allowbreak{}}\par
\noindent Source: {\ttfamily\footnotesize ED-EN-FORMULA-G247;\allowbreak{} fr/sga7ii/\{diplomati\allowbreak{}c,current\}/expose\_XX\allowbreak{}II\_body.tex;\allowbreak{} folio\allowbreak{} 437\allowbreak{}}\par

\par\medskip\noindent\hypertarget{app-E7II-XXII-052.4}{}\textbf{E7II-XXII-052.4}\quad printed folio 437.\par
\noindent Printed reading: \emph{(empty)}\par
\noindent Corrected reading: {\ttfamily\footnotesize by\allowbreak{} }\par
\noindent Source: {\ttfamily\footnotesize ED-EN-FORMULA-G247;\allowbreak{} fr/sga7ii/\{diplomati\allowbreak{}c,current\}/expose\_XX\allowbreak{}II\_body.tex;\allowbreak{} folio\allowbreak{} 437\allowbreak{}}\par

\par\medskip\noindent\hypertarget{app-E7II-XXII-053.1}{}\textbf{E7II-XXII-053.1}\quad printed folio 437.\par
\noindent Printed reading: {\ttfamily\footnotesize 1\allowbreak{}}\par
\noindent Corrected reading: {\ttfamily\footnotesize i\allowbreak{}}\par
\noindent Source: {\ttfamily\footnotesize fr/sga7ii/\{diplomati\allowbreak{}c,current\}/expose\_XX\allowbreak{}II\_body.tex;\allowbreak{} folio\allowbreak{} 437\allowbreak{}}\par

\par\medskip\noindent\hypertarget{app-E7II-XXII-053.2}{}\textbf{E7II-XXII-053.2}\quad printed folio 437.\par
\noindent Printed reading: {\ttfamily\footnotesize f\allowbreak{}}\par
\noindent Corrected reading: {\ttfamily\footnotesize \textbackslash{}mathfrak\allowbreak{} F\allowbreak{}}\par
\noindent Source: {\ttfamily\footnotesize ED-EN-FONT-G024;\allowbreak{} F7II-XXII-076.2;\allowbreak{} fr/sga7ii/\{diplomati\allowbreak{}c,current\}/expose\_XX\allowbreak{}II\_body.tex;\allowbreak{} folio\allowbreak{} 437\allowbreak{}}\par

\par\endgroup\normalsize\fi

\ifnum\SGAform=1\clearpage\chapter*{Inherited editorial notes}\small\raggedright\sloppy
These notes are inherited from the released English. They are preserved here as provenance records, separate from the variant apparatus and the authorial text.
\par\medskip\noindent\textbf{EN-NOTE-001} 7I IX; approximate folio 324\par
{\ttfamily\footnotesize The\allowbreak{} printed\allowbreak{} formula,\allowbreak{} confirmed\allowbreak{} in\allowbreak{} a\allowbreak{} direct\allowbreak{} high-resolution\allowbreak{} check,\allowbreak{} reuses\allowbreak{} \$\textbackslash{}Phi'\_o\$\allowbreak{} for\allowbreak{} the\allowbreak{} first\allowbreak{} subgroup\allowbreak{} and\allowbreak{} writes\allowbreak{} \$\textbackslash{}Psi'\_o\textbackslash{}subset\textbackslash{}Psi\_\allowbreak{}o\$\allowbreak{} for\allowbreak{} the\allowbreak{} second,\allowbreak{} colliding\allowbreak{} with\allowbreak{} the\allowbreak{} established\allowbreak{} notation\allowbreak{} for\allowbreak{} the\allowbreak{} component\allowbreak{} group\allowbreak{} of\allowbreak{} \$A'\$.\allowbreak{} The\allowbreak{} neutral\allowbreak{} local\allowbreak{} names\allowbreak{} \$F\_o,F'\_o\$\allowbreak{} above\allowbreak{} make\allowbreak{} the\allowbreak{} intended\allowbreak{} two\allowbreak{} subgroups\allowbreak{} unambiguous\allowbreak{} without\allowbreak{} changing\allowbreak{} the\allowbreak{} assertion.\allowbreak{}}\par
\par\medskip\noindent\textbf{EN-NOTE-002} 7I IX; approximate folio 387\par
{\ttfamily\footnotesize The\allowbreak{} source\allowbreak{} leaves\allowbreak{} the\allowbreak{} precise\allowbreak{} internal\allowbreak{} reference\allowbreak{} blank.\allowbreak{}}\par
\par\medskip\noindent\textbf{EN-NOTE-003} 7I IX; approximate folio 401\par
{\ttfamily\footnotesize The\allowbreak{} printed\allowbreak{} text\allowbreak{} changes\allowbreak{} incoherently\allowbreak{} here\allowbreak{} from\allowbreak{} an\allowbreak{} invertible\allowbreak{} sheaf\allowbreak{} \$\textbackslash{}mathbb\allowbreak{} M\$\allowbreak{} on\allowbreak{} \$G\$\allowbreak{} to\allowbreak{} the\allowbreak{} restriction\allowbreak{} \$\textbackslash{}mathbb\allowbreak{} L\_n\$\allowbreak{} of\allowbreak{} \$\textbackslash{}mathbb\allowbreak{} L\$\allowbreak{} to\allowbreak{} \$G\_n\$.\allowbreak{} The\allowbreak{} notation\allowbreak{} above\allowbreak{} is\allowbreak{} the\allowbreak{} type-correct\allowbreak{} reading\allowbreak{} fixed\allowbreak{} by\allowbreak{} the\allowbreak{} surrounding\allowbreak{} construction\allowbreak{} and\allowbreak{} by\allowbreak{} Proposition\textasciitilde{}\textbackslash{}SGARefL\allowbreak{}ink\{sga7.sga7i.eix.n\allowbreak{}7-1-5\}\{7.1.5\}.\allowbreak{}}\par
\par\medskip\noindent\textbf{EN-NOTE-004} 7I IX; approximate folio 402\par
{\ttfamily\footnotesize At\allowbreak{} both\allowbreak{} occurrences\allowbreak{} in\allowbreak{} this\allowbreak{} statement\allowbreak{} the\allowbreak{} print\allowbreak{} writes\allowbreak{} \$\textbackslash{}operatorname\{Specf\allowbreak{}\}(A)\$,\allowbreak{} and\allowbreak{} it\allowbreak{} also\allowbreak{} writes\allowbreak{} \$\textbackslash{}operatorname\{Spec\}\allowbreak{}(A/\textbackslash{}underline\allowbreak{} m\textasciicircum{}\{n+1\})\$\allowbreak{} below.\allowbreak{} Since\allowbreak{} \$V\$\allowbreak{} is\allowbreak{} the\allowbreak{} adic\allowbreak{} base\allowbreak{} ring\allowbreak{} and\allowbreak{} \$A\$\allowbreak{} denotes\allowbreak{} the\allowbreak{} varying\allowbreak{} group\allowbreak{} scheme,\allowbreak{} the\allowbreak{} type-correct\allowbreak{} base\allowbreak{} is\allowbreak{} \$\textbackslash{}operatorname\{Spf\}(\allowbreak{}V)\$\allowbreak{} and\allowbreak{} \$S\_n=\textbackslash{}operatorname\{S\allowbreak{}pec\}(V/\textbackslash{}underline\allowbreak{} m\textasciicircum{}\{n+1\})\$.\allowbreak{}}\par
\par\medskip\noindent\textbf{EN-NOTE-005} 7I IX; approximate folio 405\par
{\ttfamily\footnotesize The\allowbreak{} print\allowbreak{} writes\allowbreak{} \$S\$\allowbreak{} rather\allowbreak{} than\allowbreak{} the\allowbreak{} fiber\allowbreak{} bases\allowbreak{} \$S\_o\$\allowbreak{} and\allowbreak{} \$S\_n\$\allowbreak{} in\allowbreak{} these\allowbreak{} sentences,\allowbreak{} and\allowbreak{} leaves\allowbreak{} blank\allowbreak{} the\allowbreak{} parenthetical\allowbreak{} citation\allowbreak{} from\allowbreak{} which\allowbreak{} representability\allowbreak{} of\allowbreak{} \$\textbackslash{}Phi\_n\$\allowbreak{} is\allowbreak{} said\allowbreak{} to\allowbreak{} follow.\allowbreak{} The\allowbreak{} bases\allowbreak{} are\allowbreak{} restored\allowbreak{} from\allowbreak{} the\allowbreak{} statement\allowbreak{} of\allowbreak{} the\allowbreak{} lemma;\allowbreak{} the\allowbreak{} missing\allowbreak{} citation\allowbreak{} is\allowbreak{} not\allowbreak{} guessed.\allowbreak{}}\par
\par\medskip\noindent\textbf{EN-NOTE-006} 7I IX; approximate folio 421\par
{\ttfamily\footnotesize The\allowbreak{} printed\allowbreak{} source\allowbreak{} assigns\allowbreak{} \$P,R,Q\$,\allowbreak{} in\allowbreak{} that\allowbreak{} order,\allowbreak{} to\allowbreak{} these\allowbreak{} three\allowbreak{} graded\allowbreak{} pieces.\allowbreak{} The\allowbreak{} immediately\allowbreak{} following\allowbreak{} exact\allowbreak{} sequences\allowbreak{} and\allowbreak{} every\allowbreak{} subsequent\allowbreak{} construction\allowbreak{} require\allowbreak{} \$Q,R,P\$;\allowbreak{} the\allowbreak{} evident\allowbreak{} interchange\allowbreak{} of\allowbreak{} \$P\$\allowbreak{} and\allowbreak{} \$Q\$\allowbreak{} has\allowbreak{} therefore\allowbreak{} been\allowbreak{} corrected\allowbreak{} here.\allowbreak{}}\par
\par\medskip\noindent\textbf{EN-NOTE-007} 7I IX; approximate folio 423\par
{\ttfamily\footnotesize At\allowbreak{} both\allowbreak{} occurrences\allowbreak{} in\allowbreak{} this\allowbreak{} paragraph\allowbreak{} the\allowbreak{} printed\allowbreak{} source\allowbreak{} has\allowbreak{} \$F\$\allowbreak{} where\allowbreak{} the\allowbreak{} objects\allowbreak{} involved\allowbreak{} are\allowbreak{} extensions\allowbreak{} of\allowbreak{} \$Q\$\allowbreak{} by\allowbreak{} \$R\$\allowbreak{} and\allowbreak{} hence\allowbreak{} must\allowbreak{} be\allowbreak{} identified\allowbreak{} with\allowbreak{} \$E\$.\allowbreak{}}\par
\par\medskip\noindent\textbf{EN-NOTE-008} 7I IX; approximate folio 424\par
{\ttfamily\footnotesize The\allowbreak{} printed\allowbreak{} source\allowbreak{} says\allowbreak{} ``prove\allowbreak{} (d)'',\allowbreak{} although\allowbreak{} Proposition\allowbreak{} \textbackslash{}SGARefLink\{sga7.sga\allowbreak{}7i.eix.n9-3-8\}\{9.3.8\allowbreak{}\}\allowbreak{} has\allowbreak{} only\allowbreak{} parts\allowbreak{} (a),\allowbreak{} (b),\allowbreak{} and\allowbreak{} (c).\allowbreak{}}\par
\par\medskip\noindent\textbf{EN-NOTE-009} 7I IX; approximate folio 425\par
{\ttfamily\footnotesize The\allowbreak{} printed\allowbreak{} final\allowbreak{} term\allowbreak{} is\allowbreak{} \$\textbackslash{}operatorname\{Ext\}\textasciicircum{}\allowbreak{}2(R,Q)\$;\allowbreak{} applying\allowbreak{} \$\textbackslash{}operatorname\{Ext\}(\allowbreak{}-,P)\$\allowbreak{} to\allowbreak{} \$0\textbackslash{}to\allowbreak{} R\textbackslash{}to\allowbreak{} E\textbackslash{}to\allowbreak{} Q\textbackslash{}to0\$\allowbreak{} gives\allowbreak{} the\allowbreak{} displayed\allowbreak{} \$\textbackslash{}operatorname\{Ext\}\textasciicircum{}\allowbreak{}2(Q,P)\$,\allowbreak{} which\allowbreak{} is\allowbreak{} also\allowbreak{} the\allowbreak{} group\allowbreak{} specified\allowbreak{} in\allowbreak{} (\textbackslash{}SGARefLink\{sga7.sg\allowbreak{}a7i.eix.n9-3-9\}\{9.3.\allowbreak{}9\}).\allowbreak{}}\par
\par\medskip\noindent\textbf{EN-NOTE-010} 7I IX; approximate folio 426\par
{\ttfamily\footnotesize The\allowbreak{} printed\allowbreak{} right-hand\allowbreak{} side\allowbreak{} of\allowbreak{} the\allowbreak{} following\allowbreak{} formula\allowbreak{} has\allowbreak{} \$c\_X(Y)\$;\allowbreak{} the\allowbreak{} fixed\allowbreak{} object\allowbreak{} in\allowbreak{} this\allowbreak{} construction\allowbreak{} is\allowbreak{} \$Y'\$.\allowbreak{}}\par
\par\medskip\noindent\textbf{EN-NOTE-011} 7I IX; approximate folio 430\par
{\ttfamily\footnotesize The\allowbreak{} printed\allowbreak{} target\allowbreak{} of\allowbreak{} (\textbackslash{}SGARefLink\{sga7.sg\allowbreak{}a7i.eix.n9-4-7-2\}\{9.\allowbreak{}4.7.2\})\allowbreak{} is\allowbreak{} \$\textbackslash{}Lambda\_s\$;\allowbreak{} the\allowbreak{} French\allowbreak{} transcription's\allowbreak{} \$\textbackslash{}wedge\_s\$\allowbreak{} is\allowbreak{} a\allowbreak{} character-recognitio\allowbreak{}n\allowbreak{} error.\allowbreak{}}\par
\par\medskip\noindent\textbf{EN-NOTE-012} 7I IX; approximate folio 431\par
{\ttfamily\footnotesize The\allowbreak{} printed\allowbreak{} proof\allowbreak{} has\allowbreak{} \$\textbackslash{}underline\{\textbackslash{}operato\allowbreak{}rname\{Hom\}\}(P,Q)\$\allowbreak{} here.\allowbreak{} Proposition\allowbreak{} 9.3.8(a),\allowbreak{} formula\allowbreak{} (\textbackslash{}SGARefLink\{sga7.sg\allowbreak{}a7i.eix.n9-3-4\}\{9.3.\allowbreak{}4\}),\allowbreak{} and\allowbreak{} the\allowbreak{} displayed\allowbreak{} identification\allowbreak{} with\allowbreak{} \$D(Q\textbackslash{}otimes\allowbreak{} Q\textasciicircum{}*)\$\allowbreak{} all\allowbreak{} require\allowbreak{} \$\textbackslash{}underline\{\textbackslash{}operato\allowbreak{}rname\{Hom\}\}(Q,P)\$.\allowbreak{}}\par
\par\medskip\noindent\textbf{EN-NOTE-013} 7I IX; approximate folio 431\par
{\ttfamily\footnotesize The\allowbreak{} printed\allowbreak{} Ext\allowbreak{} groups\allowbreak{} in\allowbreak{} this\allowbreak{} paragraph\allowbreak{} have\allowbreak{} the\allowbreak{} arguments\allowbreak{} \$(P,Q)\$.\allowbreak{} The\allowbreak{} action\allowbreak{} of\allowbreak{} \$\textbackslash{}operatorname\{Ext\}\textasciicircum{}\allowbreak{}1(Q,P)\$\allowbreak{} on\allowbreak{} mixed\allowbreak{} extensions\allowbreak{} in\allowbreak{} Proposition\allowbreak{} \textbackslash{}SGARefLink\{sga7.sga\allowbreak{}7i.eix.n9-3-8\}\{9.3.8\allowbreak{}\}(b),\allowbreak{} and\allowbreak{} Lemma\allowbreak{} \textbackslash{}SGARefLink\{sga7.sga\allowbreak{}7i.eix.n9-4-3\}\{9.4.3\allowbreak{}\}\allowbreak{} itself,\allowbreak{} require\allowbreak{} \$(Q,P)\$.\allowbreak{}}\par
\par\medskip\noindent\textbf{EN-NOTE-014} 7I IX; approximate folio 432\par
{\ttfamily\footnotesize The\allowbreak{} printed\allowbreak{} proof\allowbreak{} cites\allowbreak{} the\allowbreak{} nonexistent\allowbreak{} pairings\allowbreak{} ``(9.4.6.2)'';\allowbreak{} these\allowbreak{} are\allowbreak{} the\allowbreak{} pairings\allowbreak{} (\textbackslash{}SGARefLink\{sga7.sg\allowbreak{}a7i.eix.n9-4-7-2\}\{9.\allowbreak{}4.7.2\}).\allowbreak{}}\par
\par\medskip\noindent\textbf{EN-NOTE-015} 7I IX; approximate folio 434\par
{\ttfamily\footnotesize The\allowbreak{} printed\allowbreak{} proof\allowbreak{} cites\allowbreak{} \textbackslash{}SGARefLink\{sga7.sga\allowbreak{}7i.eix.n9-4-6\}\{9.4.6\allowbreak{}\},\allowbreak{} which\allowbreak{} only\allowbreak{} performs\allowbreak{} the\allowbreak{} henselian\allowbreak{} descent\allowbreak{} construction;\allowbreak{} the\allowbreak{} full-faithfulness\allowbreak{} and\allowbreak{} essential-image\allowbreak{} statement\allowbreak{} being\allowbreak{} applied\allowbreak{} componentwise\allowbreak{} is\allowbreak{} Proposition\allowbreak{} \textbackslash{}SGARefLink\{sga7.sga\allowbreak{}7i.eix.n9-4-7\}\{9.4.7\allowbreak{}\}.\allowbreak{}}\par
\par\medskip\noindent\textbf{EN-NOTE-016} 7I IX; approximate folio 434\par
{\ttfamily\footnotesize At\allowbreak{} this\allowbreak{} point\allowbreak{} the\allowbreak{} printed\allowbreak{} source\allowbreak{} gives\allowbreak{} the\allowbreak{} extension\allowbreak{} over\allowbreak{} \$S\$\allowbreak{} the\allowbreak{} name\allowbreak{} \$X\_\textbackslash{}eta\$\allowbreak{} as\allowbreak{} well.\allowbreak{} We\allowbreak{} reserve\allowbreak{} \$X\_\textbackslash{}eta\$\allowbreak{} for\allowbreak{} the\allowbreak{} generic\allowbreak{} fibre\allowbreak{} and\allowbreak{} write\allowbreak{} \$X\$\allowbreak{} for\allowbreak{} its\allowbreak{} extension\allowbreak{} over\allowbreak{} \$S\$.\allowbreak{}}\par
\par\medskip\noindent\textbf{EN-NOTE-017} 7I IX; approximate folio 434\par
{\ttfamily\footnotesize The\allowbreak{} printed\allowbreak{} sentence\allowbreak{} says\allowbreak{} that\allowbreak{} \$Q(n)\$\allowbreak{} is\allowbreak{} étale\allowbreak{} over\allowbreak{} \$\textbackslash{}eta\$.\allowbreak{} The\allowbreak{} group\allowbreak{} \$Q(n)\$\allowbreak{} was\allowbreak{} defined\allowbreak{} over\allowbreak{} \$S\$\allowbreak{} and\allowbreak{} is\allowbreak{} assumed\allowbreak{} of\allowbreak{} étale\allowbreak{} type\allowbreak{} there;\allowbreak{} the\allowbreak{} base\allowbreak{} required\allowbreak{} for\allowbreak{} this\allowbreak{} flatness\allowbreak{} argument\allowbreak{} is\allowbreak{} \$S\$.\allowbreak{}}\par
\par\medskip\noindent\textbf{EN-NOTE-018} 7I IX; approximate folio 435\par
{\ttfamily\footnotesize The\allowbreak{} printed\allowbreak{} source\allowbreak{} refers\allowbreak{} here\allowbreak{} to\allowbreak{} ``the\allowbreak{} homomorphism\allowbreak{} (\textbackslash{}SGARefLink\{sga7.sg\allowbreak{}a7i.eix.n9-5-6\}\{9.5.\allowbreak{}6\})'';\allowbreak{} the\allowbreak{} homomorphism\allowbreak{} constructed\allowbreak{} in\allowbreak{} this\allowbreak{} section\allowbreak{} is\allowbreak{} (\textbackslash{}SGARefLink\{sga7.sg\allowbreak{}a7i.eix.n9-5-5\}\{9.5.\allowbreak{}5\}),\allowbreak{} whereas\allowbreak{} \textbackslash{}SGARefLink\{sga7.sga\allowbreak{}7i.eix.n9-5-6\}\{9.5.6\allowbreak{}\}\allowbreak{} is\allowbreak{} the\allowbreak{} proposition.\allowbreak{}}\par
\par\medskip\noindent\textbf{EN-NOTE-019} 7I IX; approximate folio 437\par
{\ttfamily\footnotesize The\allowbreak{} printed\allowbreak{} formula\allowbreak{} has\allowbreak{} \$F\_\textbackslash{}eta=D(F'\_\textbackslash{}eta)\$.\allowbreak{} Since\allowbreak{} \$E=D(F')\$\allowbreak{} and\allowbreak{} the\allowbreak{} displayed\allowbreak{} quotient\allowbreak{} is\allowbreak{} the\allowbreak{} extension\allowbreak{} of\allowbreak{} \$Q\_\textbackslash{}eta\$\allowbreak{} by\allowbreak{} \$R\_\textbackslash{}eta\$,\allowbreak{} the\allowbreak{} required\allowbreak{} term\allowbreak{} is\allowbreak{} \$E\_\textbackslash{}eta\$.\allowbreak{}}\par
\par\medskip\noindent\textbf{EN-NOTE-020} 7I IX; approximate folio 437\par
{\ttfamily\footnotesize The\allowbreak{} printed\allowbreak{} source\allowbreak{} calls\allowbreak{} this\allowbreak{} a\allowbreak{} pairing\allowbreak{} ``(\textbackslash{}SGARefLink\{sga7.\allowbreak{}sga7i.eix.n9-5-6\}\{9.\allowbreak{}5.6\})'';\allowbreak{} the\allowbreak{} pairing\allowbreak{} is\allowbreak{} the\allowbreak{} homomorphism\allowbreak{} (\textbackslash{}SGARefLink\{sga7.sg\allowbreak{}a7i.eix.n9-5-5\}\{9.5.\allowbreak{}5\}),\allowbreak{} whereas\allowbreak{} \textbackslash{}SGARefLink\{sga7.sga\allowbreak{}7i.eix.n9-5-6\}\{9.5.6\allowbreak{}\}\allowbreak{} is\allowbreak{} the\allowbreak{} subsequent\allowbreak{} proposition.\allowbreak{}}\par
\par\medskip\noindent\textbf{EN-NOTE-021} 7I IX; approximate folio 437\par
{\ttfamily\footnotesize The\allowbreak{} printed\allowbreak{} sentence\allowbreak{} says\allowbreak{} that\allowbreak{} \$X\_\textbackslash{}eta\$\allowbreak{} is\allowbreak{} ``over\allowbreak{} \$S\$'';\allowbreak{} its\allowbreak{} notation\allowbreak{} and\allowbreak{} definition\allowbreak{} make\allowbreak{} it\allowbreak{} the\allowbreak{} generic-fibre\allowbreak{} group\allowbreak{} over\allowbreak{} \$\textbackslash{}eta\$.\allowbreak{} It\allowbreak{} also\allowbreak{} lacks\allowbreak{} its\allowbreak{} closing\allowbreak{} parenthesis.\allowbreak{}}\par
\par\medskip\noindent\textbf{EN-NOTE-022} 7I IX; approximate folio 441\par
{\ttfamily\footnotesize In\allowbreak{} the\allowbreak{} printed\allowbreak{} right-hand\allowbreak{} side\allowbreak{} the\allowbreak{} underlying\allowbreak{} pairing\allowbreak{} is\allowbreak{} denoted\allowbreak{} only\allowbreak{} by\allowbreak{} \$\textbackslash{}varphi\textasciicircum{}t\$,\allowbreak{} without\allowbreak{} a\allowbreak{} definition.\allowbreak{} It\allowbreak{} is\allowbreak{} the\allowbreak{} monodromy\allowbreak{} pairing\allowbreak{} \$u\_\textbackslash{}ell\textasciicircum{}\{A\_K\}\$\allowbreak{} used\allowbreak{} in\allowbreak{} the\allowbreak{} identical\allowbreak{} construction\allowbreak{} (\textbackslash{}SGARefLink\{sga7.sg\allowbreak{}a7i.eix.n10-2-4\}\{10.\allowbreak{}2.4\}).\allowbreak{}}\par
\par\medskip\noindent\textbf{EN-NOTE-023} 7I IX; approximate folio 441\par
{\ttfamily\footnotesize The\allowbreak{} printed\allowbreak{} source\allowbreak{} refers\allowbreak{} here\allowbreak{} to\allowbreak{} (\textbackslash{}SGARefLink\{sga7.sg\allowbreak{}a7i.eix.n9-2-1\}\{9.2.\allowbreak{}1\}),\allowbreak{} which\allowbreak{} is\allowbreak{} the\allowbreak{} inertia\allowbreak{} orthogonality\allowbreak{} formula\allowbreak{} \$gx-x\textbackslash{}in\allowbreak{} V\textasciicircum{}\textbackslash{}perp\$,\allowbreak{} not\allowbreak{} the\allowbreak{} monodromy\allowbreak{} pairing.\allowbreak{} The\allowbreak{} pairing\allowbreak{} is\allowbreak{} numbered\allowbreak{} (\textbackslash{}SGARefLink\{sga7.sg\allowbreak{}a7i.eix.n9-1-2\}\{9.1.\allowbreak{}2\}).\allowbreak{}}\par
\par\medskip\noindent\textbf{EN-NOTE-024} 7I IX; approximate folio 442\par
{\ttfamily\footnotesize The\allowbreak{} printed\allowbreak{} source\allowbreak{} leaves\allowbreak{} the\allowbreak{} internal\allowbreak{} reference\allowbreak{} following\allowbreak{} ``Expos\textbackslash{}'e\allowbreak{} I''\allowbreak{} blank.\allowbreak{}}\par
\par\medskip\noindent\textbf{EN-NOTE-025} 7I IX; approximate folio 442\par
{\ttfamily\footnotesize The\allowbreak{} printed\allowbreak{} source\allowbreak{} cites\allowbreak{} (\textbackslash{}SGARefLink\{sga7.sg\allowbreak{}a7i.eix.n10-3-3\}\{10.\allowbreak{}3.3\}),\allowbreak{} the\allowbreak{} displayed\allowbreak{} map\allowbreak{} \$g\_*\$;\allowbreak{} the\allowbreak{} multiplication\allowbreak{} formula\allowbreak{} it\allowbreak{} uses\allowbreak{} is\allowbreak{} (\textbackslash{}SGARefLink\{sga7.sg\allowbreak{}a7i.eix.n10-3-4\}\{10.\allowbreak{}3.4\}).\allowbreak{}}\par
\par\medskip\noindent\textbf{EN-NOTE-026} 7I IX; approximate folio 443\par
{\ttfamily\footnotesize The\allowbreak{} printed\allowbreak{} source\allowbreak{} says\allowbreak{} that\allowbreak{} \$K\textasciicircum{}*/V\textasciicircum{}*\textbackslash{}simeq\textbackslash{}mathb\allowbreak{}b\allowbreak{} Z\$\allowbreak{} is\allowbreak{} deduced\allowbreak{} from\allowbreak{} the\allowbreak{} ``norm.''\allowbreak{} The\allowbreak{} displayed\allowbreak{} map\allowbreak{} is\allowbreak{} the\allowbreak{} normalized\allowbreak{} discrete\allowbreak{} valuation.\allowbreak{}}\par
\par\medskip\noindent\textbf{EN-NOTE-027} 7I IX; approximate folio 445\par
{\ttfamily\footnotesize The\allowbreak{} printed\allowbreak{} source\allowbreak{} says\allowbreak{} that\allowbreak{} \$\textbackslash{}overline\allowbreak{} A\_K\$\allowbreak{} is\allowbreak{} ``over\allowbreak{} \$k\$'';\allowbreak{} the\allowbreak{} notation\allowbreak{} and\allowbreak{} the\allowbreak{} direct-factor\allowbreak{} argument\allowbreak{} require\allowbreak{} ``over\allowbreak{} \$K\$.''\allowbreak{}}\par
\par\medskip\noindent\textbf{EN-NOTE-028} 7I IX; approximate folio 446\par
{\ttfamily\footnotesize The\allowbreak{} printed\allowbreak{} source\allowbreak{} leaves\allowbreak{} the\allowbreak{} pinpoint\allowbreak{} following\allowbreak{} citation\allowbreak{} [15]\allowbreak{} blank.\allowbreak{}}\par
\par\medskip\noindent\textbf{EN-NOTE-029} 7I IX; approximate folio 446\par
{\ttfamily\footnotesize The\allowbreak{} printed\allowbreak{} source\allowbreak{} says\allowbreak{} that\allowbreak{} the\allowbreak{} curve\allowbreak{} \$X\_K\$\allowbreak{} is\allowbreak{} ``over\allowbreak{} \$X\$'';\allowbreak{} the\allowbreak{} displayed\allowbreak{} relative\allowbreak{} Picard\allowbreak{} scheme\allowbreak{} and\allowbreak{} the\allowbreak{} argument\allowbreak{} require\allowbreak{} ``over\allowbreak{} \$K\$.''\allowbreak{}}\par
\par\medskip\noindent\textbf{EN-NOTE-030} 7I IX; approximate folio 446\par
{\ttfamily\footnotesize The\allowbreak{} printed\allowbreak{} source\allowbreak{} leaves\allowbreak{} the\allowbreak{} internal\allowbreak{} reference\allowbreak{} following\allowbreak{} ``Expos\textbackslash{}'e\allowbreak{} V''\allowbreak{} blank.\allowbreak{}}\par
\par\medskip\noindent\textbf{EN-NOTE-031} 7I IX; approximate folio 451\par
{\ttfamily\footnotesize The\allowbreak{} printed\allowbreak{} source\allowbreak{} leaves\allowbreak{} the\allowbreak{} internal\allowbreak{} reference\allowbreak{} following\allowbreak{} the\allowbreak{} first\allowbreak{} occurrence\allowbreak{} of\allowbreak{} \$I\$\allowbreak{} blank.\allowbreak{}}\par
\par\medskip\noindent\textbf{EN-NOTE-032} 7I IX; approximate folio 454\par
{\ttfamily\footnotesize The\allowbreak{} printed\allowbreak{} source\allowbreak{} reads\allowbreak{} ``as\allowbreak{} in\allowbreak{} \textbackslash{}SGARefLink\{sga7.sga\allowbreak{}7i.eix.n11-4\}\{11.4\}'\allowbreak{}'\allowbreak{} here,\allowbreak{} within\allowbreak{} section\allowbreak{} \textbackslash{}SGARefLink\{sga7.sga\allowbreak{}7i.eix.n11-4\}\{11.4\}\allowbreak{} itself.\allowbreak{} The\allowbreak{} reference\allowbreak{} is\allowbreak{} corrected\allowbreak{} to\allowbreak{} the\allowbreak{} immediately\allowbreak{} preceding\allowbreak{} construction\allowbreak{} in\allowbreak{} \textbackslash{}SGARefLink\{sga7.sga\allowbreak{}7i.eix.n11-3\}\{11.3\}.\allowbreak{}}\par
\par\medskip\noindent\textbf{EN-NOTE-033} 7I IX; approximate folio 456\par
{\ttfamily\footnotesize At\allowbreak{} this\allowbreak{} point\allowbreak{} the\allowbreak{} printed\allowbreak{} source\allowbreak{} labels\allowbreak{} \$V/M\_\textbackslash{}ell\$\allowbreak{} as\allowbreak{} \$\textbackslash{}Phi(\textbackslash{}ell)\$\allowbreak{} and\allowbreak{} \$V/\textbackslash{}check\allowbreak{} M'\_\textbackslash{}ell\$\allowbreak{} as\allowbreak{} \$\textbackslash{}Phi'(\textbackslash{}ell)\$.\allowbreak{} Those\allowbreak{} quotients\allowbreak{} are\allowbreak{} instead\allowbreak{} the\allowbreak{} source\allowbreak{} and\allowbreak{} target\allowbreak{} of\allowbreak{} the\allowbreak{} map\allowbreak{} whose\allowbreak{} kernel\allowbreak{} is\allowbreak{} the\allowbreak{} finite\allowbreak{} cokernel\allowbreak{} \$\textbackslash{}check\allowbreak{} M'\_\textbackslash{}ell/M\_\textbackslash{}ell\$;\allowbreak{} the\allowbreak{} contradictory\allowbreak{} labels\allowbreak{} are\allowbreak{} not\allowbreak{} reproduced.\allowbreak{}}\par
\par\medskip\noindent\textbf{EN-NOTE-034} 7I IX; approximate folio 460\par
{\ttfamily\footnotesize The\allowbreak{} printed\allowbreak{} source\allowbreak{} refers\allowbreak{} here\allowbreak{} to\allowbreak{} nonexistent\allowbreak{} formula\allowbreak{} (11.7.5);\allowbreak{} the\allowbreak{} filtration\allowbreak{} whose\allowbreak{} generic\allowbreak{} fibre\allowbreak{} is\allowbreak{} being\allowbreak{} used\allowbreak{} is\allowbreak{} (\textbackslash{}SGARefLink\{sga7.sg\allowbreak{}a7i.eix.n11-6-5\}\{11.\allowbreak{}6.5\}).\allowbreak{}}\par
\par\medskip\noindent\textbf{EN-NOTE-035} 7I IX; approximate folio 460\par
{\ttfamily\footnotesize The\allowbreak{} printed\allowbreak{} source\allowbreak{} cites\allowbreak{} \textbackslash{}SGARefLink\{sga7.sga\allowbreak{}7i.eix.n9-4-6\}\{9.4.6\allowbreak{}\},\allowbreak{} but\allowbreak{} the\allowbreak{} required\allowbreak{} extension\allowbreak{} criterion\allowbreak{} is\allowbreak{} Proposition\allowbreak{} \textbackslash{}SGARefLink\{sga7.sga\allowbreak{}7i.eix.n9-4-7\}\{9.4.7\allowbreak{}\}.\allowbreak{}}\par
\par\medskip\noindent\textbf{EN-NOTE-036} 7I IX; approximate folio 460\par
{\ttfamily\footnotesize The\allowbreak{} printed\allowbreak{} source\allowbreak{} says\allowbreak{} ``of\allowbreak{} \$K\_n\$''\allowbreak{} in\allowbreak{} this\allowbreak{} general\allowbreak{} criterion,\allowbreak{} although\allowbreak{} the\allowbreak{} subgroup\allowbreak{} being\allowbreak{} varied\allowbreak{} at\allowbreak{} this\allowbreak{} point\allowbreak{} is\allowbreak{} \$K'\_n\$;\allowbreak{} the\allowbreak{} specialization\allowbreak{} \$K'\_n=K\_n\$\allowbreak{} is\allowbreak{} made\allowbreak{} only\allowbreak{} below.\allowbreak{}}\par
\par\medskip\noindent\textbf{EN-NOTE-037} 7I IX; approximate folio 461\par
{\ttfamily\footnotesize The\allowbreak{} printed\allowbreak{} source\allowbreak{} says\allowbreak{} ``the\allowbreak{} construction\allowbreak{} made\allowbreak{} in\allowbreak{} \textbackslash{}SGARefLink\{sga7.sga\allowbreak{}7i.eix.n11-7\}\{11.7\}'\allowbreak{}'\allowbreak{} at\allowbreak{} the\allowbreak{} opening\allowbreak{} of\allowbreak{} \textbackslash{}SGARefLink\{sga7.sga\allowbreak{}7i.eix.n11-7\}\{11.7\};\allowbreak{} the\allowbreak{} reference\allowbreak{} is\allowbreak{} to\allowbreak{} the\allowbreak{} construction\allowbreak{} just\allowbreak{} completed\allowbreak{} in\allowbreak{} \textbackslash{}SGARefLink\{sga7.sga\allowbreak{}7i.eix.n11-6\}\{11.6\}.\allowbreak{}}\par
\par\medskip\noindent\textbf{EN-NOTE-038} 7I IX; approximate folio 464\par
{\ttfamily\footnotesize The\allowbreak{} printed\allowbreak{} source\allowbreak{} appends\allowbreak{} ``cf.\allowbreak{} I''\allowbreak{} followed\allowbreak{} by\allowbreak{} a\allowbreak{} blank\allowbreak{} internal\allowbreak{} reference.\allowbreak{} No\allowbreak{} target\allowbreak{} is\allowbreak{} supplied,\allowbreak{} so\allowbreak{} none\allowbreak{} is\allowbreak{} invented\allowbreak{} here.\allowbreak{}}\par
\par\medskip\noindent\textbf{EN-NOTE-039} 7I IX; approximate folio 467\par
{\ttfamily\footnotesize The\allowbreak{} printed\allowbreak{} source\allowbreak{} says\allowbreak{} ``on\allowbreak{} \$S\$''\allowbreak{} here\allowbreak{} and\allowbreak{} later\allowbreak{} writes\allowbreak{} \$X\_i\textasciicircum{}o\$\allowbreak{} for\allowbreak{} the\allowbreak{} same\allowbreak{} divisor.\allowbreak{} The\allowbreak{} sheaf\allowbreak{} is\allowbreak{} the\allowbreak{} invertible\allowbreak{} sheaf\allowbreak{} on\allowbreak{} \$X\$\allowbreak{} associated\allowbreak{} with\allowbreak{} the\allowbreak{} component\allowbreak{} \$X\_o\textasciicircum{}i\$;\allowbreak{} the\allowbreak{} English\allowbreak{} text\allowbreak{} keeps\allowbreak{} that\allowbreak{} notation\allowbreak{} consistently.\allowbreak{}}\par
\par\medskip\noindent\textbf{EN-NOTE-040} 7I IX; approximate folio 468\par
{\ttfamily\footnotesize The\allowbreak{} printed\allowbreak{} source\allowbreak{} leaves\allowbreak{} the\allowbreak{} internal\allowbreak{} locator\allowbreak{} following\allowbreak{} ``III''\allowbreak{} blank;\allowbreak{} no\allowbreak{} target\allowbreak{} is\allowbreak{} invented\allowbreak{} here.\allowbreak{}}\par
\par\medskip\noindent\textbf{EN-NOTE-041} 7I IX; approximate folio 470\par
{\ttfamily\footnotesize The\allowbreak{} printed\allowbreak{} source\allowbreak{} has\allowbreak{} \$\textbackslash{}underline\{\textbackslash{}operato\allowbreak{}rname\{Pic\}\}\textasciicircum{}o\_\{X/k\}\$\allowbreak{} at\allowbreak{} the\allowbreak{} end\allowbreak{} of\allowbreak{} this\allowbreak{} argument.\allowbreak{} The\allowbreak{} curve\allowbreak{} throughout\allowbreak{} \textbackslash{}SGARefLink\{sga7.sga\allowbreak{}7i.eix.n12-3\}\{12.3\}\allowbreak{} is\allowbreak{} \$C\$,\allowbreak{} and\allowbreak{} the\allowbreak{} preceding\allowbreak{} and\allowbreak{} following\allowbreak{} occurrences\allowbreak{} both\allowbreak{} read\allowbreak{} \$\textbackslash{}underline\{\textbackslash{}operato\allowbreak{}rname\{Pic\}\}\textasciicircum{}o\_\{C/k\}\$\allowbreak{};\allowbreak{} the\allowbreak{} English\allowbreak{} text\allowbreak{} restores\allowbreak{} \$C\$.\allowbreak{}}\par
\par\medskip\noindent\textbf{EN-NOTE-042} 7I IX; approximate folio 471\par
{\ttfamily\footnotesize The\allowbreak{} printed\allowbreak{} source\allowbreak{} numbers\allowbreak{} these\allowbreak{} two\allowbreak{} displays\allowbreak{} (12.2.3)\allowbreak{} and\allowbreak{} (12.2.4).\allowbreak{} They\allowbreak{} belong\allowbreak{} to\allowbreak{} section\allowbreak{} \textbackslash{}SGARefLink\{sga7.sga\allowbreak{}7i.eix.n12-3\}\{12.3\}\allowbreak{} and\allowbreak{} are\allowbreak{} cited\allowbreak{} below\allowbreak{} as\allowbreak{} (\textbackslash{}SGARefLink\{sga7.sg\allowbreak{}a7i.eix.n12-3-3\}\{12.\allowbreak{}3.3\})\allowbreak{} and\allowbreak{} (\textbackslash{}SGARefLink\{sga7.sg\allowbreak{}a7i.eix.n12-3-4\}\{12.\allowbreak{}3.4\});\allowbreak{} those\allowbreak{} coherent\allowbreak{} numbers\allowbreak{} are\allowbreak{} used\allowbreak{} here.\allowbreak{}}\par
\par\medskip\noindent\textbf{EN-NOTE-043} 7I IX; approximate folio 473\par
{\ttfamily\footnotesize The\allowbreak{} printed\allowbreak{} display\allowbreak{} defining\allowbreak{} (\textbackslash{}SGARefLink\{sga7.sg\allowbreak{}a7i.eix.n12-3-8\}\{12.\allowbreak{}3.8\})\allowbreak{} carries\allowbreak{} no\allowbreak{} number\allowbreak{} even\allowbreak{} though\allowbreak{} the\allowbreak{} following\allowbreak{} text\allowbreak{} cites\allowbreak{} it,\allowbreak{} and\allowbreak{} (\textbackslash{}SGARefLink\{sga7.sg\allowbreak{}a7i.eix.n12-3-13\}\{12\allowbreak{}.3.13\})\allowbreak{} writes\allowbreak{} \$\textbackslash{}mathbb\allowbreak{} Z\textasciicircum{}\{(\textbackslash{}Gamma)\}\$\allowbreak{} where\allowbreak{} the\allowbreak{} edge\allowbreak{} chain\allowbreak{} group\allowbreak{} is\allowbreak{} \$\textbackslash{}mathbb\allowbreak{} Z\textasciicircum{}\{(\textbackslash{}Gamma\_1)\}\$.\allowbreak{} The\allowbreak{} English\allowbreak{} text\allowbreak{} supplies\allowbreak{} the\allowbreak{} cited\allowbreak{} label\allowbreak{} and\allowbreak{} the\allowbreak{} required\allowbreak{} subscript.\allowbreak{} The\allowbreak{} working\allowbreak{} transcription's\allowbreak{} \$\textbackslash{}widetilde\allowbreak{} C\_i\textbackslash{}times\_C\textbackslash{}widetild\allowbreak{}e\allowbreak{} C\_1\$\allowbreak{} in\allowbreak{} the\allowbreak{} loop-edge\allowbreak{} description\allowbreak{} is\allowbreak{} also\allowbreak{} restored\allowbreak{} from\allowbreak{} the\allowbreak{} scan\allowbreak{} as\allowbreak{} \$\textbackslash{}widetilde\allowbreak{} C\_i\textbackslash{}times\_C\textbackslash{}widetild\allowbreak{}e\allowbreak{} C\_i\$.\allowbreak{}}\par
\par\medskip\noindent\textbf{EN-NOTE-044} 7I IX; approximate folio 475\par
{\ttfamily\footnotesize The\allowbreak{} printed\allowbreak{} citation\allowbreak{} following\allowbreak{} the\allowbreak{} canonical\allowbreak{} polarization\allowbreak{} is\allowbreak{} ``[5]''\allowbreak{} followed\allowbreak{} by\allowbreak{} a\allowbreak{} blank\allowbreak{} pinpoint.\allowbreak{} The\allowbreak{} English\allowbreak{} text\allowbreak{} preserves\allowbreak{} the\allowbreak{} reference\allowbreak{} without\allowbreak{} inventing\allowbreak{} a\allowbreak{} missing\allowbreak{} locator.\allowbreak{}}\par
\par\medskip\noindent\textbf{EN-NOTE-045} 7I IX; approximate folio 476\par
{\ttfamily\footnotesize The\allowbreak{} working\allowbreak{} transcription\allowbreak{} adds\allowbreak{} a\allowbreak{} dual\allowbreak{} superscript\allowbreak{} to\allowbreak{} \$T\_\textbackslash{}ell(A\_\{\textbackslash{}bar\textbackslash{}eta\}\allowbreak{})\$\allowbreak{} in\allowbreak{} (\textbackslash{}SGARefLink\{sga7.sg\allowbreak{}a7i.eix.n12-7-1\}\{12.\allowbreak{}7.1\}).\allowbreak{} A\allowbreak{} direct\allowbreak{} 5,000-dpi\allowbreak{} inspection\allowbreak{} of\allowbreak{} the\allowbreak{} printed\allowbreak{} formula\allowbreak{} shows\allowbreak{} no\allowbreak{} such\allowbreak{} superscript,\allowbreak{} in\allowbreak{} agreement\allowbreak{} with\allowbreak{} the\allowbreak{} standard\allowbreak{} identification\allowbreak{} \$H\textasciicircum{}1(A\_\{\textbackslash{}bar\textbackslash{}eta\},\textbackslash{}m\allowbreak{}athbb\allowbreak{} Z\_\textbackslash{}ell)=\allowbreak{} \textbackslash{}operatorname\{Hom\}(T\allowbreak{}\_\textbackslash{}ell(A\_\{\textbackslash{}bar\textbackslash{}eta\}),\allowbreak{}\textbackslash{}mathbb\allowbreak{} Z\_\textbackslash{}ell)\$.\allowbreak{}}\par
\par\medskip\noindent\textbf{EN-NOTE-046} 7I IX; approximate folio 477\par
{\ttfamily\footnotesize The\allowbreak{} printed\allowbreak{} source\allowbreak{} leaves\allowbreak{} the\allowbreak{} internal\allowbreak{} pinpoint\allowbreak{} following\allowbreak{} ``V''\allowbreak{} blank;\allowbreak{} no\allowbreak{} section\allowbreak{} target\allowbreak{} is\allowbreak{} supplied.\allowbreak{}}\par
\par\medskip\noindent\textbf{EN-NOTE-047} 7I IX; approximate folio 477\par
{\ttfamily\footnotesize The\allowbreak{} printed\allowbreak{} source\allowbreak{} likewise\allowbreak{} leaves\allowbreak{} the\allowbreak{} internal\allowbreak{} pinpoint\allowbreak{} following\allowbreak{} ``VI''\allowbreak{} blank.\allowbreak{}}\par
\par\medskip\noindent\textbf{EN-NOTE-048} 7I IX; approximate folio 478\par
{\ttfamily\footnotesize The\allowbreak{} working\allowbreak{} transcription\allowbreak{} reads\allowbreak{} \$\textbackslash{}underline\allowbreak{} F\$\allowbreak{} here.\allowbreak{} Direct\allowbreak{} inspection\allowbreak{} of\allowbreak{} the\allowbreak{} printed\allowbreak{} page\allowbreak{} shows\allowbreak{} \$\textbackslash{}underline\allowbreak{} f\$,\allowbreak{} as\allowbreak{} also\allowbreak{} required\allowbreak{} by\allowbreak{} the\allowbreak{} sentence's\allowbreak{} reference\allowbreak{} to\allowbreak{} the\allowbreak{} special\allowbreak{} fibre\allowbreak{} of\allowbreak{} the\allowbreak{} universal\allowbreak{} curve.\allowbreak{}}\par
\par\medskip\noindent\textbf{EN-NOTE-049} 7I IX; approximate folio 480\par
{\ttfamily\footnotesize The\allowbreak{} printed\allowbreak{} source\allowbreak{} reads\allowbreak{} \$T\_\textbackslash{}ell(B\_U)\$\allowbreak{} here,\allowbreak{} although\allowbreak{} no\allowbreak{} \$B\_U\$\allowbreak{} has\allowbreak{} been\allowbreak{} introduced\allowbreak{} and\allowbreak{} the\allowbreak{} pairing\allowbreak{} and\allowbreak{} every\allowbreak{} adjacent\allowbreak{} term\allowbreak{} involve\allowbreak{} \$P\_U\$.\allowbreak{} The\allowbreak{} English\allowbreak{} text\allowbreak{} restores\allowbreak{} the\allowbreak{} internally\allowbreak{} forced\allowbreak{} \$P\_U\$.\allowbreak{}}\par
\par\medskip\noindent\textbf{EN-NOTE-050} 7I IX; approximate folio 480\par
{\ttfamily\footnotesize The\allowbreak{} printed\allowbreak{} upper\allowbreak{} identification\allowbreak{} in\allowbreak{} (\textbackslash{}SGARefLink\{sga7.sg\allowbreak{}a7i.eix.n12-8-9\}\{12.\allowbreak{}8.9\})\allowbreak{} omits\allowbreak{} the\allowbreak{} outer\allowbreak{} closing\allowbreak{} parenthesis\allowbreak{} after\allowbreak{} \$\textbackslash{}underline\{\textbackslash{}operato\allowbreak{}rname\{Hom\}\}(M,T\_\textbackslash{}ell\allowbreak{}(\textbackslash{}mathbb\allowbreak{} G\_\{m\textbackslash{}underline\allowbreak{} S\}))\$;\allowbreak{} the\allowbreak{} balanced\allowbreak{} expression\allowbreak{} has\allowbreak{} been\allowbreak{} supplied\allowbreak{} here.\allowbreak{}}\par
\par\medskip\noindent\textbf{EN-NOTE-051} 7I IX; approximate folio 482\par
{\ttfamily\footnotesize Both\allowbreak{} printed\allowbreak{} citations\allowbreak{} to\allowbreak{} Expos\textbackslash{}'e\allowbreak{} VI\allowbreak{} in\allowbreak{} this\allowbreak{} sentence\allowbreak{} leave\allowbreak{} their\allowbreak{} internal\allowbreak{} section\allowbreak{} pinpoints\allowbreak{} blank.\allowbreak{}}\par
\par\medskip\noindent\textbf{EN-NOTE-052} 7I IX; approximate folio 482\par
{\ttfamily\footnotesize The\allowbreak{} printed\allowbreak{} source\allowbreak{} has\allowbreak{} the\allowbreak{} lowercase\allowbreak{} letter\allowbreak{} \$x\$\allowbreak{} as\allowbreak{} the\allowbreak{} subject\allowbreak{} of\allowbreak{} ``is\allowbreak{} regular.''\allowbreak{} The\allowbreak{} following\allowbreak{} sentence\allowbreak{} and\allowbreak{} the\allowbreak{} mathematical\allowbreak{} statement\allowbreak{} require\allowbreak{} the\allowbreak{} scheme\allowbreak{} \$X\$.\allowbreak{}}\par
\par\medskip\noindent\textbf{EN-NOTE-053} 7I IX; approximate folio 482\par
{\ttfamily\footnotesize In\allowbreak{} this\allowbreak{} reduction\allowbreak{} the\allowbreak{} printed\allowbreak{} source\allowbreak{} drops\allowbreak{} the\allowbreak{} subscript\allowbreak{} \$\textbackslash{}lambda\$\allowbreak{} from\allowbreak{} \$D\_\textbackslash{}lambda\$\allowbreak{} and\allowbreak{} prints\allowbreak{} an\allowbreak{} un-underlined\allowbreak{} \$X\$\allowbreak{} as\allowbreak{} the\allowbreak{} right-hand\allowbreak{} factor\allowbreak{} of\allowbreak{} the\allowbreak{} following\allowbreak{} fibre\allowbreak{} product.\allowbreak{} Both\allowbreak{} marks\allowbreak{} are\allowbreak{} restored:\allowbreak{} the\allowbreak{} localization\allowbreak{} is\allowbreak{} at\allowbreak{} the\allowbreak{} component\allowbreak{} \$D\_\textbackslash{}lambda\$,\allowbreak{} and\allowbreak{} the\allowbreak{} fibre\allowbreak{} product\allowbreak{} is\allowbreak{} formed\allowbreak{} from\allowbreak{} the\allowbreak{} universal\allowbreak{} curve\allowbreak{} \$\textbackslash{}underline\allowbreak{} X\$\allowbreak{} over\allowbreak{} \$\textbackslash{}underline\allowbreak{} S\$.\allowbreak{}}\par
\par\medskip\noindent\textbf{EN-NOTE-054} 7I IX; approximate folio 483\par
{\ttfamily\footnotesize The\allowbreak{} working\allowbreak{} transcription\allowbreak{} misreads\allowbreak{} the\allowbreak{} right\allowbreak{} diagonal\allowbreak{} label\allowbreak{} in\allowbreak{} (\textbackslash{}SGARefLink\{sga7.sg\allowbreak{}a7i.eix.n12-8-15\}\{12\allowbreak{}.8.15\})\allowbreak{} as\allowbreak{} \$\textbackslash{}varphi\_\{\textbackslash{}lambda\_\textbackslash{}l\allowbreak{}ambda\}\$.\allowbreak{} The\allowbreak{} printed\allowbreak{} diagram,\allowbreak{} the\allowbreak{} immediately\allowbreak{} preceding\allowbreak{} definition,\allowbreak{} and\allowbreak{} the\allowbreak{} formula\allowbreak{} below\allowbreak{} all\allowbreak{} give\allowbreak{} \$\textbackslash{}varphi\_\{y\_\textbackslash{}lambda\}\allowbreak{}\$.\allowbreak{}}\par
\par\medskip\noindent\textbf{EN-NOTE-055} 7I IX; approximate folio 483\par
{\ttfamily\footnotesize At\allowbreak{} this\allowbreak{} final\allowbreak{} step\allowbreak{} the\allowbreak{} printed\allowbreak{} source\allowbreak{} writes\allowbreak{} \$\textbackslash{}underline\allowbreak{} S\$\allowbreak{} rather\allowbreak{} than\allowbreak{} the\allowbreak{} localized\allowbreak{} trait\allowbreak{} \$\textbackslash{}underline\allowbreak{} S\textasciicircum{}\textbackslash{}lambda\$\allowbreak{} used\allowbreak{} throughout\allowbreak{} the\allowbreak{} reduction.\allowbreak{} The\allowbreak{} English\allowbreak{} text\allowbreak{} restores\allowbreak{} the\allowbreak{} mathematically\allowbreak{} required\allowbreak{} superscript.\allowbreak{}}\par
\par\medskip\noindent\textbf{EN-NOTE-056} 7I IX; approximate folio 489\par
{\ttfamily\footnotesize The\allowbreak{} printed\allowbreak{} source\allowbreak{} refers\allowbreak{} here\allowbreak{} to\allowbreak{} the\allowbreak{} conditions\allowbreak{} of\allowbreak{} \textbackslash{}SGARefLink\{sga7.sga\allowbreak{}7i.eix.n13-2-2\}\{13.2\allowbreak{}.2\}\allowbreak{} itself;\allowbreak{} the\allowbreak{} three\allowbreak{} conditions\allowbreak{} are\allowbreak{} introduced\allowbreak{} in\allowbreak{} Lemma\allowbreak{} \textbackslash{}SGARefLink\{sga7.sga\allowbreak{}7i.eix.n13-2-1\}\{13.2\allowbreak{}.1\}.\allowbreak{}}\par
\par\medskip\noindent\textbf{EN-NOTE-057} 7I IX; approximate folio 489\par
{\ttfamily\footnotesize The\allowbreak{} printed\allowbreak{} source\allowbreak{} has\allowbreak{} \$A=A\_o\$\allowbreak{} here.\allowbreak{} The\allowbreak{} required\allowbreak{} hypothesis,\allowbreak{} and\allowbreak{} the\allowbreak{} one\allowbreak{} used\allowbreak{} below,\allowbreak{} is\allowbreak{} \$A=A\textasciicircum{}o\$;\allowbreak{} compare\allowbreak{} \textbackslash{}SGARefLink\{sga7.sga\allowbreak{}7i.eix.n13-1\}\{13.1\}\allowbreak{} and\allowbreak{} Lemma\allowbreak{} \textbackslash{}SGARefLink\{sga7.sga\allowbreak{}7i.eix.n13-2-1\}\{13.2\allowbreak{}.1\}.\allowbreak{}}\par
\par\medskip\noindent\textbf{EN-NOTE-058} 7I IX; approximate folio 496\par
{\ttfamily\footnotesize The\allowbreak{} printed\allowbreak{} source\allowbreak{} says\allowbreak{} here\allowbreak{} that\allowbreak{} \$B\$\allowbreak{} is\allowbreak{} a\allowbreak{} quotient\allowbreak{} of\allowbreak{} \$A\$;\allowbreak{} the\allowbreak{} displayed\allowbreak{} definition\allowbreak{} \$B=A\textasciicircum{}\textbackslash{}natural/T\$\allowbreak{} gives\allowbreak{} the\allowbreak{} intended\allowbreak{} group.\allowbreak{}}\par
\par\medskip\noindent\textbf{EN-NOTE-059} 7I IX; approximate folio 502\par
{\ttfamily\footnotesize The\allowbreak{} printed\allowbreak{} source\allowbreak{} says\allowbreak{} ``projective\allowbreak{} over\allowbreak{} \$S\$''\allowbreak{} here;\allowbreak{} the\allowbreak{} hypothesis\allowbreak{} of\allowbreak{} Lemma\allowbreak{} \textbackslash{}SGARefLink\{sga7.sga\allowbreak{}7i.eix.n13-5-2\}\{13.5\allowbreak{}.2\}\allowbreak{} and\allowbreak{} the\allowbreak{} ensuing\allowbreak{} normality\allowbreak{} argument\allowbreak{} require\allowbreak{} ``projective\allowbreak{} over\allowbreak{} \$\textbackslash{}hat\{\textbackslash{}underline\allowbreak{} O\}\$''.\allowbreak{}}\par
\par\medskip\noindent\textbf{EN-NOTE-060} 7II XI; approximate folio 58\par
{\ttfamily\footnotesize The\allowbreak{} printed\allowbreak{} source\allowbreak{} reads\allowbreak{} \$\textbackslash{}min(n,\textbackslash{}dim(X)-n)\$.\allowbreak{} The\allowbreak{} factor\allowbreak{} \$2\$\allowbreak{} is\allowbreak{} required\allowbreak{} by\allowbreak{} the\allowbreak{} Hodge\allowbreak{} ranges\allowbreak{} \$0\textbackslash{}leq\allowbreak{} p,q\textbackslash{}leq\textbackslash{}dim(X)\$;\allowbreak{} the\allowbreak{} English\allowbreak{} text\allowbreak{} corrects\allowbreak{} this\allowbreak{} evident\allowbreak{} source\allowbreak{} typo.\allowbreak{}}\par
\par\medskip\noindent\textbf{EN-NOTE-061} 7II XII; approximate folio 64\par
{\ttfamily\footnotesize In\allowbreak{} both\allowbreak{} displayed\allowbreak{} sums\allowbreak{} the\allowbreak{} printed\allowbreak{} source\allowbreak{} has\allowbreak{} upper\allowbreak{} limit\allowbreak{} \$m-1\$.\allowbreak{} That\allowbreak{} leaves\allowbreak{} a\allowbreak{} two-dimensional\allowbreak{} radical\allowbreak{} and\allowbreak{} contradicts\allowbreak{} the\allowbreak{} statement;\allowbreak{} the\allowbreak{} standard\allowbreak{} ordinary\allowbreak{} normal\allowbreak{} form,\allowbreak{} and\allowbreak{} the\allowbreak{} induction\allowbreak{} that\allowbreak{} follows,\allowbreak{} require\allowbreak{} the\allowbreak{} upper\allowbreak{} limit\allowbreak{} \$m\$.\allowbreak{}}\par
\par\medskip\noindent\textbf{EN-NOTE-062} 7II XII; approximate folio 64\par
{\ttfamily\footnotesize The\allowbreak{} final\allowbreak{} term\allowbreak{} in\allowbreak{} the\allowbreak{} printed\allowbreak{} definition\allowbreak{} of\allowbreak{} \$f\$\allowbreak{} is\allowbreak{} \$f\$;\allowbreak{} it\allowbreak{} is\allowbreak{} the\allowbreak{} evident\allowbreak{} source\allowbreak{} typo\allowbreak{} \$f'\$,\allowbreak{} as\allowbreak{} is\allowbreak{} confirmed\allowbreak{} by\allowbreak{} the\allowbreak{} displayed\allowbreak{} conclusion.\allowbreak{}}\par
\par\medskip\noindent\textbf{EN-NOTE-063} 7II XII; approximate folio 65\par
{\ttfamily\footnotesize The\allowbreak{} printed\allowbreak{} formula\allowbreak{} has\allowbreak{} \$(-1)\textasciicircum{}i\$\allowbreak{} and\allowbreak{} begins\allowbreak{} the\allowbreak{} remaining\allowbreak{} wedge\allowbreak{} with\allowbreak{} an\allowbreak{} undefined\allowbreak{} \$e\_0\$.\allowbreak{} The\allowbreak{} English\allowbreak{} restores\allowbreak{} the\allowbreak{} standard\allowbreak{} contraction\allowbreak{} formula\allowbreak{} \$(-1)\textasciicircum{}\{i-1\}e\_1\textbackslash{}wedge\allowbreak{}\textbackslash{}cdots\textbackslash{}widehat\allowbreak{} e\_i\textbackslash{}cdots\textbackslash{}wedge\allowbreak{} e\_p\$,\allowbreak{} required\allowbreak{} by\allowbreak{} the\allowbreak{} Clifford\allowbreak{} action\allowbreak{} displayed\allowbreak{} immediately\allowbreak{} below.\allowbreak{} The\allowbreak{} reading\allowbreak{} was\allowbreak{} checked\allowbreak{} in\allowbreak{} a\allowbreak{} targeted\allowbreak{} high-detail\allowbreak{} enlargement\allowbreak{} of\allowbreak{} folio\allowbreak{} 65.\allowbreak{}}\par
\par\medskip\noindent\textbf{EN-NOTE-064} 7II XII; approximate folio 67\par
{\ttfamily\footnotesize The\allowbreak{} printed\allowbreak{} sentence\allowbreak{} says\allowbreak{} that\allowbreak{} \$e\$\allowbreak{} factors\allowbreak{} through\allowbreak{} a\allowbreak{} section\allowbreak{} of\allowbreak{} the\allowbreak{} affine\allowbreak{} space\allowbreak{} \$\textbackslash{}mathbb\allowbreak{} E(\textbackslash{}bigwedge\textasciicircum{}2W')\$.\allowbreak{} The\allowbreak{} connected-to-étale\allowbreak{} argument\allowbreak{} instead\allowbreak{} gives\allowbreak{} the\allowbreak{} factorization\allowbreak{} through\allowbreak{} \$S\$\allowbreak{} and\allowbreak{} a\allowbreak{} section\allowbreak{} of\allowbreak{} \$\textbackslash{}Spec(Z(Q))\$,\allowbreak{} exactly\allowbreak{} as\allowbreak{} the\allowbreak{} following\allowbreak{} conclusion\allowbreak{} requires.\allowbreak{}}\par
\par\medskip\noindent\textbf{EN-NOTE-065} 7II XII; approximate folio 67\par
{\ttfamily\footnotesize The\allowbreak{} printed\allowbreak{} cross-reference\allowbreak{} is\allowbreak{} 1.6;\allowbreak{} the\allowbreak{} scalar-change\allowbreak{} identification\allowbreak{} is\allowbreak{} the\allowbreak{} construction\allowbreak{} of\allowbreak{} 1.3.1.\allowbreak{}}\par
\par\medskip\noindent\textbf{EN-NOTE-066} 7II XII; approximate folio 69\par
{\ttfamily\footnotesize The\allowbreak{} printed\allowbreak{} proposition\allowbreak{} says\allowbreak{} that\allowbreak{} \$V\$\allowbreak{} has\allowbreak{} dimension\allowbreak{} \$2m\$\allowbreak{} ``over\allowbreak{} \$V\$'';\allowbreak{} the\allowbreak{} base\allowbreak{} field\allowbreak{} introduced\allowbreak{} in\allowbreak{} the\allowbreak{} same\allowbreak{} sentence\allowbreak{} is\allowbreak{} \$A\$.\allowbreak{}}\par
\par\medskip\noindent\textbf{EN-NOTE-067} 7II XII; approximate folio 72\par
{\ttfamily\footnotesize The\allowbreak{} print\allowbreak{} cites\allowbreak{} 2.3(ii)\allowbreak{} for\allowbreak{} the\allowbreak{} cohomological\allowbreak{} vanishing\allowbreak{} and\allowbreak{} 2.3(iii)\allowbreak{} for\allowbreak{} the\allowbreak{} Picard-group\allowbreak{} computation.\allowbreak{} The\allowbreak{} two\allowbreak{} references\allowbreak{} are\allowbreak{} interchanged\allowbreak{} here.\allowbreak{}}\par
\par\medskip\noindent\textbf{EN-NOTE-068} 7II XII; approximate folio 72\par
{\ttfamily\footnotesize After\allowbreak{} introducing\allowbreak{} the\allowbreak{} projection\allowbreak{} as\allowbreak{} \$\textbackslash{}rho\$,\allowbreak{} the\allowbreak{} print\allowbreak{} changes\allowbreak{} its\allowbreak{} symbol\allowbreak{} to\allowbreak{} \$p\$\allowbreak{} in\allowbreak{} the\allowbreak{} direct-image\allowbreak{} formulas.\allowbreak{} The\allowbreak{} notation\allowbreak{} is\allowbreak{} regularized\allowbreak{} here;\allowbreak{} the\allowbreak{} reference\allowbreak{} to\allowbreak{} 2.3(iii)\allowbreak{} is\allowbreak{} retained.\allowbreak{}}\par
\par\medskip\noindent\textbf{EN-NOTE-069} 7II XII; approximate folio 75\par
{\ttfamily\footnotesize The\allowbreak{} displayed\allowbreak{} definition\allowbreak{} in\allowbreak{} the\allowbreak{} print\allowbreak{} has\allowbreak{} \$R\textasciicircum{}1p\_*\$,\allowbreak{} whereas\allowbreak{} the\allowbreak{} following\allowbreak{} sentence\allowbreak{} and\allowbreak{} the\allowbreak{} degree\allowbreak{} of\allowbreak{} \$c\_1(\textbackslash{}mathcal\allowbreak{} O(1))\$\allowbreak{} both\allowbreak{} give\allowbreak{} \$R\textasciicircum{}2p\_*\$.\allowbreak{}}\par
\par\medskip\noindent\textbf{EN-NOTE-070} 7II XII; approximate folio 76\par
{\ttfamily\footnotesize The\allowbreak{} print\allowbreak{} repeats\allowbreak{} the\allowbreak{} tag\allowbreak{} (3.2.1)\allowbreak{} on\allowbreak{} the\allowbreak{} map\allowbreak{} in\allowbreak{} part\allowbreak{} (iii),\allowbreak{} but\allowbreak{} the\allowbreak{} proof\allowbreak{} itself\allowbreak{} refers\allowbreak{} to\allowbreak{} this\allowbreak{} map\allowbreak{} as\allowbreak{} (3.3.1).\allowbreak{}}\par
\par\medskip\noindent\textbf{EN-NOTE-071} 7II XII; approximate folio 77\par
{\ttfamily\footnotesize The\allowbreak{} printed\allowbreak{} ambient\allowbreak{} projective\allowbreak{} space\allowbreak{} is\allowbreak{} \$\textbackslash{}mathbb\allowbreak{} P\textasciicircum{}\{m+1\}\$,\allowbreak{} although\allowbreak{} the\allowbreak{} displayed\allowbreak{} nondegenerate\allowbreak{} equation\allowbreak{} has\allowbreak{} \$2m+2\$\allowbreak{} coordinates.\allowbreak{} The\allowbreak{} required\allowbreak{} ambient\allowbreak{} space\allowbreak{} is\allowbreak{} \$\textbackslash{}mathbb\allowbreak{} P\textasciicircum{}\{2m+1\}\$.\allowbreak{}}\par
\par\medskip\noindent\textbf{EN-NOTE-072} 7II XII; approximate folio 77\par
{\ttfamily\footnotesize In\allowbreak{} these\allowbreak{} two\allowbreak{} top-degree\allowbreak{} intersection\allowbreak{} computations\allowbreak{} the\allowbreak{} print\allowbreak{} has\allowbreak{} \$\textbackslash{}eta\textasciicircum{}2\$\allowbreak{} in\allowbreak{} place\allowbreak{} of\allowbreak{} \$\textbackslash{}eta\textasciicircum{}\{2m\}\$,\allowbreak{} and\allowbreak{} once\allowbreak{} prints\allowbreak{} \$c\textbackslash{}ell(\textbackslash{}alpha\textasciicircum{}2)\$\allowbreak{} in\allowbreak{} place\allowbreak{} of\allowbreak{} \$c\textbackslash{}ell(\textbackslash{}alpha)\textasciicircum{}2\$.\allowbreak{}}\par
\par\medskip\noindent\textbf{EN-NOTE-073} 7II XII; approximate folio 79\par
{\ttfamily\footnotesize This\allowbreak{} paragraph\allowbreak{} prints\allowbreak{} \$r=2m+1\$\allowbreak{} instead\allowbreak{} of\allowbreak{} \$n=2m+1\$\allowbreak{} and\allowbreak{} cites\allowbreak{} 3.5.3\allowbreak{} instead\allowbreak{} of\allowbreak{} the\allowbreak{} displayed\allowbreak{} sequence\allowbreak{} (3.6.3).\allowbreak{}}\par
\par\medskip\noindent\textbf{EN-NOTE-074} 7II XII; approximate folio 80\par
{\ttfamily\footnotesize The\allowbreak{} leftmost\allowbreak{} term\allowbreak{} in\allowbreak{} the\allowbreak{} printed\allowbreak{} lower\allowbreak{} row\allowbreak{} has\allowbreak{} twist\allowbreak{} \$(m)\$;\allowbreak{} duality\allowbreak{} and\allowbreak{} the\allowbreak{} other\allowbreak{} three\allowbreak{} terms\allowbreak{} require\allowbreak{} the\allowbreak{} displayed\allowbreak{} twist\allowbreak{} \$(m+1)\$.\allowbreak{}}\par
\par\medskip\noindent\textbf{EN-NOTE-075} 7II XIII; approximate folio 95\par
{\ttfamily\footnotesize (the\allowbreak{} printed\allowbreak{} diagram\allowbreak{} inadvertently\allowbreak{} omits\allowbreak{} the\allowbreak{} bar\allowbreak{} from\allowbreak{} its\allowbreak{} top-middle\allowbreak{} \$\textbackslash{}bar\allowbreak{} X\$).\allowbreak{}}\par
\par\medskip\noindent\textbf{EN-NOTE-076} 7II XIII; approximate folio 97\par
{\ttfamily\footnotesize Here\allowbreak{} and\allowbreak{} below\allowbreak{} we\allowbreak{} retain\allowbreak{} the\allowbreak{} scheme\allowbreak{} \$Y\$\allowbreak{} fixed\allowbreak{} in\allowbreak{} 1.4.1;\allowbreak{} the\allowbreak{} printed\allowbreak{} opening\allowbreak{} sentence\allowbreak{} of\allowbreak{} 1.4.2\allowbreak{} inadvertently\allowbreak{} writes\allowbreak{} \$X\textbackslash{}times\_sS\$.\allowbreak{}}\par
\par\medskip\noindent\textbf{EN-NOTE-077} 7II XIII; approximate folio 99\par
{\ttfamily\footnotesize The\allowbreak{} printed\allowbreak{} bottom-right\allowbreak{} term\allowbreak{} has\allowbreak{} \$K\_s\$;\allowbreak{} the\allowbreak{} defining\allowbreak{} identity\allowbreak{} (2.1.2.2)\allowbreak{} requires\allowbreak{} \$K\_\textbackslash{}eta\$,\allowbreak{} as\allowbreak{} written\allowbreak{} here.\allowbreak{}}\par
\par\medskip\noindent\textbf{EN-NOTE-078} 7II XIII; approximate folio 101\par
{\ttfamily\footnotesize The\allowbreak{} print\allowbreak{} mistakenly\allowbreak{} refers\allowbreak{} here\allowbreak{} to\allowbreak{} the\allowbreak{} later\allowbreak{} special\allowbreak{} case\allowbreak{} (2.1.8.1).\allowbreak{}}\par
\par\medskip\noindent\textbf{EN-NOTE-079} 7II XIII; approximate folio 102\par
{\ttfamily\footnotesize The\allowbreak{} printed\allowbreak{} reference\allowbreak{} to\allowbreak{} (2.1.7.1)\allowbreak{} is\allowbreak{} incompatible\allowbreak{} with\allowbreak{} the\allowbreak{} two\allowbreak{} functors\allowbreak{} involved.\allowbreak{}}\par
\par\medskip\noindent\textbf{EN-NOTE-080} 7II XIII; approximate folio 103\par
{\ttfamily\footnotesize In\allowbreak{} (2.1.8.8)\allowbreak{} the\allowbreak{} English\allowbreak{} restores\allowbreak{} the\allowbreak{} bar\allowbreak{} omitted\allowbreak{} from\allowbreak{} the\allowbreak{} printed\allowbreak{} bottom-left\allowbreak{} \$X\_\{\textbackslash{}bar\allowbreak{} s\}\$\allowbreak{} and\allowbreak{} the\allowbreak{} missing\allowbreak{} closing\allowbreak{} parenthesis\allowbreak{} in\allowbreak{} the\allowbreak{} bottom-middle\allowbreak{} term.\allowbreak{} In\allowbreak{} (2.1.8.7),\allowbreak{} the\allowbreak{} upper\allowbreak{} connector\allowbreak{} is\allowbreak{} corrected\allowbreak{} to\allowbreak{} an\allowbreak{} arrow,\allowbreak{} as\allowbreak{} in\allowbreak{} (2.1.8.8).\allowbreak{}}\par
\par\medskip\noindent\textbf{EN-NOTE-081} 7II XIII; approximate folio 108\par
{\ttfamily\footnotesize The\allowbreak{} print\allowbreak{} refers\allowbreak{} here\allowbreak{} to\allowbreak{} 2.2.1,\allowbreak{} but\allowbreak{} the\allowbreak{} coefficient\allowbreak{} ring\allowbreak{} \$\textbackslash{}Lambda\$\allowbreak{} is\allowbreak{} fixed\allowbreak{} in\allowbreak{} 2.1.1.\allowbreak{}}\par
\par\medskip\noindent\textbf{EN-NOTE-082} 7II XIII; approximate folio 109\par
{\ttfamily\footnotesize The\allowbreak{} print\allowbreak{} has\allowbreak{} \$Rf\_\textbackslash{}Phi\$\allowbreak{} as\allowbreak{} the\allowbreak{} first\allowbreak{} functor\allowbreak{} in\allowbreak{} (2.2.5.1);\allowbreak{} the\allowbreak{} immediately\allowbreak{} preceding\allowbreak{} definition\allowbreak{} and\allowbreak{} the\allowbreak{} remainder\allowbreak{} of\allowbreak{} the\allowbreak{} composite\allowbreak{} require\allowbreak{} \$Rg\_\textbackslash{}Phi\$,\allowbreak{} which\allowbreak{} is\allowbreak{} used\allowbreak{} here.\allowbreak{}}\par
\par\medskip\noindent\textbf{EN-NOTE-083} 7II XIII; approximate folio 109\par
{\ttfamily\footnotesize The\allowbreak{} print\allowbreak{} has\allowbreak{} \$\textbackslash{}\{x\textbackslash{}\}\$\allowbreak{} also\allowbreak{} on\allowbreak{} the\allowbreak{} right,\allowbreak{} but\allowbreak{} the\allowbreak{} support\allowbreak{} for\allowbreak{} \$g\$\allowbreak{} lies\allowbreak{} in\allowbreak{} \$Y\_s\$\allowbreak{} and\allowbreak{} is\allowbreak{} \$\textbackslash{}\{y\textbackslash{}\}\$,\allowbreak{} as\allowbreak{} the\allowbreak{} following\allowbreak{} diagram\allowbreak{} confirms.\allowbreak{}}\par
\par\medskip\noindent\textbf{EN-NOTE-084} 7II XIII; approximate folio 109\par
{\ttfamily\footnotesize The\allowbreak{} print's\allowbreak{} references\allowbreak{} to\allowbreak{} (2.3.5.2)\allowbreak{} and\allowbreak{} (2.3.5.1)\allowbreak{} are\allowbreak{} the\allowbreak{} immediately\allowbreak{} preceding\allowbreak{} (2.2.5.2)\allowbreak{} and\allowbreak{} (2.2.5.1);\allowbreak{} likewise,\allowbreak{} \$\textbackslash{}Phi\textbackslash{}subset\allowbreak{} Y\_s\$\allowbreak{} makes\allowbreak{} \$\textbackslash{}\{y\textbackslash{}\}\$,\allowbreak{} rather\allowbreak{} than\allowbreak{} the\allowbreak{} printed\allowbreak{} \$\textbackslash{}\{x\textbackslash{}\}\$,\allowbreak{} the\allowbreak{} relevant\allowbreak{} support.\allowbreak{}}\par
\par\medskip\noindent\textbf{EN-NOTE-085} 7II XIII; approximate folio 113\par
{\ttfamily\footnotesize The\allowbreak{} printed\allowbreak{} \$R\textbackslash{}Phi(F)\$\allowbreak{} in\allowbreak{} the\allowbreak{} direct-summand\allowbreak{} sentence\allowbreak{} is\allowbreak{} \$R\textbackslash{}Phi(K)\$\allowbreak{} in\allowbreak{} this\allowbreak{} general\allowbreak{} formulation.\allowbreak{}}\par
\par\medskip\noindent\textbf{EN-NOTE-086} 7II XIII; approximate folio 114\par
{\ttfamily\footnotesize The\allowbreak{} printed\allowbreak{} reference\allowbreak{} (2.4.4.1)\allowbreak{} in\allowbreak{} the\allowbreak{} discussion\allowbreak{} following\allowbreak{} (2.4.6.2)\allowbreak{} is\allowbreak{} the\allowbreak{} variation\allowbreak{} morphism\allowbreak{} (2.4.5.1).\allowbreak{}}\par
\par\medskip\noindent\textbf{EN-NOTE-087} 7II XIII; approximate folio 114\par
{\ttfamily\footnotesize The\allowbreak{} source\allowbreak{} prints\allowbreak{} both\allowbreak{} this\allowbreak{} paragraph\allowbreak{} and\allowbreak{} the\allowbreak{} preceding\allowbreak{} paragraph\allowbreak{} as\allowbreak{} 2.4.6,\allowbreak{} and\allowbreak{} prints\allowbreak{} both\allowbreak{} principal\allowbreak{} diagrams\allowbreak{} as\allowbreak{} (2.4.6.1).\allowbreak{} These\allowbreak{} repeated\allowbreak{} numbers\allowbreak{} are\allowbreak{} corrected\allowbreak{} here\allowbreak{} to\allowbreak{} 2.4.7\allowbreak{} and\allowbreak{} (2.4.7.1)\allowbreak{} for\allowbreak{} this\allowbreak{} paragraph\allowbreak{} and\allowbreak{} its\allowbreak{} diagram,\allowbreak{} the\allowbreak{} latter\allowbreak{} being\allowbreak{} the\allowbreak{} one\allowbreak{} cited\allowbreak{} in\allowbreak{} Expos\textbackslash{}'e\allowbreak{} XV,\allowbreak{} 2.2.6.\allowbreak{}}\par
\par\medskip\noindent\textbf{EN-NOTE-088} 7II XIV; approximate folio 123\par
{\ttfamily\footnotesize \textbackslash{}noindent\textbackslash{}emph\{Edito\allowbreak{}rial\allowbreak{} note.\}\allowbreak{} In\allowbreak{} (1.2.7.1)\allowbreak{} the\allowbreak{} printed\allowbreak{} top-right\allowbreak{} node\allowbreak{} lacks\allowbreak{} the\allowbreak{} overline.\allowbreak{} The\allowbreak{} preceding\allowbreak{} definition\allowbreak{} and\allowbreak{} the\allowbreak{} map\allowbreak{} \$p'\$\allowbreak{} force\allowbreak{} that\allowbreak{} node\allowbreak{} to\allowbreak{} be\allowbreak{} \$\textbackslash{}overline\allowbreak{} X\textasciicircum{}*\$,\allowbreak{} which\allowbreak{} is\allowbreak{} the\allowbreak{} notation\allowbreak{} used\allowbreak{} here.\allowbreak{}}\par
\par\medskip\noindent\textbf{EN-NOTE-089} 7II XIV; approximate folio 127\par
{\ttfamily\footnotesize \textbackslash{}noindent\textbackslash{}emph\{Edito\allowbreak{}rial\allowbreak{} note.\}\allowbreak{} The\allowbreak{} print\allowbreak{} cites\allowbreak{} the\allowbreak{} triangle\allowbreak{} ``1.2.1''\allowbreak{} here;\allowbreak{} the\allowbreak{} distinguished\allowbreak{} triangle\allowbreak{} used\allowbreak{} is\allowbreak{} the\allowbreak{} one\allowbreak{} constructed\allowbreak{} in\allowbreak{} 1.2.6.\allowbreak{}}\par
\par\medskip\noindent\textbf{EN-NOTE-090} 7II XIV; approximate folio 147\par
{\ttfamily\footnotesize The\allowbreak{} printed\allowbreak{} source\allowbreak{} has\allowbreak{} \$\textbackslash{}mathbb\allowbreak{} R\textasciicircum{}n\$\allowbreak{} in\allowbreak{} (3.2.2.2),\allowbreak{} while\allowbreak{} \$z\$\allowbreak{} has\allowbreak{} \$n+1\$\allowbreak{} coordinates.\allowbreak{} The\allowbreak{} English\allowbreak{} text\allowbreak{} corrects\allowbreak{} this\allowbreak{} evident\allowbreak{} dimension\allowbreak{} typo\allowbreak{} to\allowbreak{} \$\textbackslash{}mathbb\allowbreak{} R\textasciicircum{}\{n+1\}\$.\allowbreak{}}\par
\par\medskip\noindent\textbf{EN-NOTE-091} 7II XIV; approximate folio 149\par
{\ttfamily\footnotesize The\allowbreak{} printed\allowbreak{} sentence\allowbreak{} mixes\allowbreak{} three\allowbreak{} dimension\allowbreak{} symbols:\allowbreak{} it\allowbreak{} calls\allowbreak{} the\allowbreak{} dimension\allowbreak{} \$x\$,\allowbreak{} indexes\allowbreak{} the\allowbreak{} coordinates\allowbreak{} by\allowbreak{} \$n\$,\allowbreak{} and\allowbreak{} places\allowbreak{} the\allowbreak{} classes\allowbreak{} in\allowbreak{} degree\allowbreak{} \$d\$.\allowbreak{} The\allowbreak{} English\allowbreak{} text\allowbreak{} consistently\allowbreak{} uses\allowbreak{} \$d\$\allowbreak{} in\allowbreak{} the\allowbreak{} general\allowbreak{} statement;\allowbreak{} the\allowbreak{} specialization\allowbreak{} below\allowbreak{} uses\allowbreak{} \$d=n\$.\allowbreak{}}\par
\par\medskip\noindent\textbf{EN-NOTE-092} 7II XIV; approximate folio 152\par
{\ttfamily\footnotesize The\allowbreak{} printed\allowbreak{} proof\allowbreak{} says\allowbreak{} (4.1.1),\allowbreak{} which\allowbreak{} is\allowbreak{} the\allowbreak{} preceding\allowbreak{} diagram\allowbreak{} rather\allowbreak{} than\allowbreak{} the\allowbreak{} morphism\allowbreak{} under\allowbreak{} discussion;\allowbreak{} the\allowbreak{} English\allowbreak{} restores\allowbreak{} the\allowbreak{} evident\allowbreak{} reference\allowbreak{} to\allowbreak{} (4.1.2).\allowbreak{}}\par
\par\medskip\noindent\textbf{EN-NOTE-093} 7II XIV; approximate folio 159\par
{\ttfamily\footnotesize The\allowbreak{} printed\allowbreak{} right-hand\allowbreak{} term\allowbreak{} in\allowbreak{} (4.18.2)\allowbreak{} is\allowbreak{} \$H\textasciicircum{}*(U,V)\$,\allowbreak{} although\allowbreak{} \$V\$\allowbreak{} is\allowbreak{} defined\allowbreak{} on\allowbreak{} \$U\textasciicircum{}*\$;\allowbreak{} the\allowbreak{} English\allowbreak{} restores\allowbreak{} the\allowbreak{} evident\allowbreak{} \$U\textasciicircum{}*\$.\allowbreak{}}\par
\par\medskip\noindent\textbf{EN-NOTE-094} 7II XIV; approximate folio 160\par
{\ttfamily\footnotesize The\allowbreak{} printed\allowbreak{} left-hand\allowbreak{} term\allowbreak{} is\allowbreak{} \$H\textasciicircum{}*(U\textasciicircum{}*,\textbackslash{}mathbb\allowbreak{} C)\$,\allowbreak{} while\allowbreak{} both\allowbreak{} the\allowbreak{} hypothesis\allowbreak{} and\allowbreak{} the\allowbreak{} conclusion\allowbreak{} concern\allowbreak{} the\allowbreak{} general\allowbreak{} rank-one\allowbreak{} system\allowbreak{} \$V\$;\allowbreak{} the\allowbreak{} English\allowbreak{} restores\allowbreak{} \$V\$.\allowbreak{}}\par
\par\medskip\noindent\textbf{EN-NOTE-095} 7II XIV; approximate folio 161\par
{\ttfamily\footnotesize The\allowbreak{} print\allowbreak{} refers\allowbreak{} here\allowbreak{} to\allowbreak{} 4.18.4;\allowbreak{} the\allowbreak{} two\allowbreak{} sides\allowbreak{} under\allowbreak{} discussion\allowbreak{} are\allowbreak{} those\allowbreak{} of\allowbreak{} 4.18.5.\allowbreak{}}\par
\par\medskip\noindent\textbf{EN-NOTE-096} 7II XV; approximate folio 168\par
{\ttfamily\footnotesize The\allowbreak{} printed\allowbreak{} display\allowbreak{} is\allowbreak{} numbered\allowbreak{} (4.2.1)\allowbreak{} and\allowbreak{} has\allowbreak{} an\allowbreak{} extra\allowbreak{} equals\allowbreak{} sign\allowbreak{} after\allowbreak{} \$Q(\textbackslash{}underline\allowbreak{} X)\$;\allowbreak{} the\allowbreak{} English\allowbreak{} restores\allowbreak{} the\allowbreak{} local\allowbreak{} numbering\allowbreak{} and\allowbreak{} the\allowbreak{} asserted\allowbreak{} expansion\allowbreak{} \$f=Q+\$\allowbreak{} higher-order\allowbreak{} terms.\allowbreak{}}\par
\par\medskip\noindent\textbf{EN-NOTE-097} 7II XV; approximate folio 170\par
{\ttfamily\footnotesize The\allowbreak{} print\allowbreak{} refers\allowbreak{} to\allowbreak{} 4.3\allowbreak{} here;\allowbreak{} the\allowbreak{} normal\allowbreak{} form\allowbreak{} just\allowbreak{} invoked\allowbreak{} is\allowbreak{} Example\allowbreak{} 1.2.3.\allowbreak{}}\par
\par\medskip\noindent\textbf{EN-NOTE-098} 7II XV; approximate folio 173\par
{\ttfamily\footnotesize The\allowbreak{} print\allowbreak{} refers\allowbreak{} at\allowbreak{} this\allowbreak{} point\allowbreak{} to\allowbreak{} (4.11.1);\allowbreak{} the\allowbreak{} displayed\allowbreak{} congruence\allowbreak{} under\allowbreak{} proof\allowbreak{} is\allowbreak{} (1.2.11.1).\allowbreak{}}\par
\par\medskip\noindent\textbf{EN-NOTE-099} 7II XV; approximate folio 175\par
{\ttfamily\footnotesize The\allowbreak{} print\allowbreak{} labels\allowbreak{} this\allowbreak{} display\allowbreak{} (4.14.2)\allowbreak{} and\allowbreak{} writes\allowbreak{} \$x\_0\$\allowbreak{} in\allowbreak{} its\allowbreak{} first\allowbreak{} term.\allowbreak{} The\allowbreak{} surrounding\allowbreak{} normal\allowbreak{} forms\allowbreak{} and\allowbreak{} Remark\allowbreak{} 1.3.3\allowbreak{} require\allowbreak{} the\allowbreak{} local\allowbreak{} tag\allowbreak{} (1.3.2.1)\allowbreak{} and\allowbreak{} the\allowbreak{} quadratic\allowbreak{} term\allowbreak{} \$x\_0\textasciicircum{}2\$;\allowbreak{} the\allowbreak{} English\allowbreak{} restores\allowbreak{} both.\allowbreak{}}\par
\par\medskip\noindent\textbf{EN-NOTE-100} 7II XV; approximate folio 181\par
{\ttfamily\footnotesize The\allowbreak{} printed\allowbreak{} tag\allowbreak{} omits\allowbreak{} the\allowbreak{} dot\allowbreak{} between\allowbreak{} 1\allowbreak{} and\allowbreak{} 7;\allowbreak{} the\allowbreak{} English\allowbreak{} uses\allowbreak{} the\allowbreak{} local\allowbreak{} number\allowbreak{} (2.1.7.1).\allowbreak{}}\par
\par\medskip\noindent\textbf{EN-NOTE-101} 7II XV; approximate folio 181\par
{\ttfamily\footnotesize The\allowbreak{} print\allowbreak{} numbers\allowbreak{} this\allowbreak{} lemma\allowbreak{} 2.7.8.\allowbreak{} Its\allowbreak{} position\allowbreak{} and\allowbreak{} the\allowbreak{} forward\allowbreak{} reference\allowbreak{} in\allowbreak{} 2.2.7\allowbreak{} require\allowbreak{} 2.1.8.\allowbreak{}}\par
\par\medskip\noindent\textbf{EN-NOTE-102} 7II XV; approximate folio 182\par
{\ttfamily\footnotesize The\allowbreak{} print\allowbreak{} has\allowbreak{} \$X'\_s\$\allowbreak{} in\allowbreak{} (b).\allowbreak{} The\allowbreak{} equivalent\allowbreak{} clause\allowbreak{} and\allowbreak{} the\allowbreak{} immediately\allowbreak{} following\allowbreak{} definition\allowbreak{} of\allowbreak{} the\allowbreak{} singular\allowbreak{} point\allowbreak{} show\allowbreak{} that\allowbreak{} \$X\_s\$\allowbreak{} is\allowbreak{} intended.\allowbreak{}}\par
\par\medskip\noindent\textbf{EN-NOTE-103} 7II XV; approximate folio 183\par
{\ttfamily\footnotesize The\allowbreak{} print\allowbreak{} uses\allowbreak{} the\allowbreak{} letter\allowbreak{} \$o\$\allowbreak{} in\allowbreak{} the\allowbreak{} two\allowbreak{} occurrences\allowbreak{} of\allowbreak{} \$H\textasciicircum{}o\$;\allowbreak{} the\allowbreak{} cohomological\allowbreak{} degree\allowbreak{} is\allowbreak{} \$0\$.\allowbreak{}}\par
\par\medskip\noindent\textbf{EN-NOTE-104} 7II XV; approximate folio 184\par
{\ttfamily\footnotesize The\allowbreak{} print\allowbreak{} drops\allowbreak{} the\allowbreak{} bar\allowbreak{} from\allowbreak{} \$\textbackslash{}bar\textbackslash{}eta\$\allowbreak{} in\allowbreak{} the\allowbreak{} displayed\allowbreak{} trace\allowbreak{} map;\allowbreak{} the\allowbreak{} surrounding\allowbreak{} formulas\allowbreak{} determine\allowbreak{} the\allowbreak{} indicated\allowbreak{} term.\allowbreak{}}\par
\par\medskip\noindent\textbf{EN-NOTE-105} 7II XV; approximate folio 185\par
{\ttfamily\footnotesize In\allowbreak{} (2.2.5.6)\allowbreak{} and\allowbreak{} (2.2.5.10),\allowbreak{} the\allowbreak{} print\allowbreak{} suppresses\allowbreak{} the\allowbreak{} comma\allowbreak{} in\allowbreak{} \$(a,\textbackslash{}delta)\$;\allowbreak{} this\allowbreak{} is\allowbreak{} the\allowbreak{} pairing\allowbreak{} defined\allowbreak{} in\allowbreak{} C.\allowbreak{}}\par
\par\medskip\noindent\textbf{EN-NOTE-106} 7II XV; approximate folio 187\par
{\ttfamily\footnotesize In\allowbreak{} this\allowbreak{} proof\allowbreak{} the\allowbreak{} print\allowbreak{} has\allowbreak{} \$H\_c\textasciicircum{}i\$\allowbreak{} for\allowbreak{} the\allowbreak{} group\allowbreak{} in\allowbreak{} which\allowbreak{} the\allowbreak{} degree-\$n\$\allowbreak{} cycle\allowbreak{} \$\textbackslash{}delta\$\allowbreak{} is\allowbreak{} characterized;\allowbreak{} the\allowbreak{} displayed\allowbreak{} maps\allowbreak{} and\allowbreak{} all\allowbreak{} surrounding\allowbreak{} terms\allowbreak{} require\allowbreak{} \$H\_c\textasciicircum{}n\$.\allowbreak{} The\allowbreak{} print\allowbreak{} also\allowbreak{} abbreviates\allowbreak{} the\allowbreak{} support\allowbreak{} point\allowbreak{} in\allowbreak{} (2.2.6.2)\allowbreak{} as\allowbreak{} \$x\$;\allowbreak{} the\allowbreak{} English\allowbreak{} uses\allowbreak{} the\allowbreak{} fixed\allowbreak{} singular\allowbreak{} point\allowbreak{} \$x\_0\$.\allowbreak{}}\par
\par\medskip\noindent\textbf{EN-NOTE-107} 7II XV; approximate folio 188\par
{\ttfamily\footnotesize The\allowbreak{} print\allowbreak{} says\allowbreak{} that\allowbreak{} \$R\textasciicircum{}n\textbackslash{}Phi(\textbackslash{}Lambda)\$\allowbreak{} is\allowbreak{} ``alone''\allowbreak{} outside\allowbreak{} \$\textbackslash{}Sigma\$\allowbreak{} (\textbackslash{}emph\{seul\});\allowbreak{} the\allowbreak{} proof\allowbreak{} and\allowbreak{} the\allowbreak{} asserted\allowbreak{} support\allowbreak{} on\allowbreak{} \$\textbackslash{}Sigma\$\allowbreak{} require\allowbreak{} ``zero''\allowbreak{} (\textbackslash{}emph\{nul\}).\allowbreak{} In\allowbreak{} (iii)\allowbreak{} it\allowbreak{} prints\allowbreak{} \$H\textasciicircum{}n\_\{\textbackslash{}\{\textbackslash{}bar\allowbreak{} x\textbackslash{}\}\}(R\textasciicircum{}n\textbackslash{}psi\_\{\textbackslash{}bar\textbackslash{}e\allowbreak{}ta\}(\textbackslash{}Lambda))\$;\allowbreak{} the\allowbreak{} pairing\allowbreak{} cited\allowbreak{} from\allowbreak{} 2.2.5(C)\allowbreak{} has\allowbreak{} the\allowbreak{} second\allowbreak{} term\allowbreak{} \$H\textasciicircum{}n\_\{\textbackslash{}\{\textbackslash{}bar\allowbreak{} x\textbackslash{}\}\}(X\_\{\textbackslash{}bar\allowbreak{} s\},R\textbackslash{}psi\_\{\textbackslash{}bar\textbackslash{}eta\}(\allowbreak{}\textbackslash{}Lambda))\$.\allowbreak{}}\par
\par\medskip\noindent\textbf{EN-NOTE-108} 7II XV; approximate folio 189\par
{\ttfamily\footnotesize The\allowbreak{} print\allowbreak{} has\allowbreak{} \$H\textasciicircum{}n\_\{\textbackslash{}\{x\textbackslash{}\}\}(R\textasciicircum{}n\textbackslash{}psi\allowbreak{}\_\{\textbackslash{}bar\textbackslash{}eta\}(\textbackslash{}Lambda(\allowbreak{}m)))\$\allowbreak{} in\allowbreak{} (i).\allowbreak{} The\allowbreak{} proof\allowbreak{} explicitly\allowbreak{} invokes\allowbreak{} 2.2.5(D),\allowbreak{} where\allowbreak{} the\allowbreak{} natural\allowbreak{} generator\allowbreak{} belongs\allowbreak{} to\allowbreak{} \$H\textasciicircum{}n\_\{\textbackslash{}\{x\textbackslash{}\}\}(X\_s,R\textbackslash{}p\allowbreak{}si\_\{\textbackslash{}bar\textbackslash{}eta\}(\textbackslash{}Lambd\allowbreak{}a(m)))\$.\allowbreak{}}\par
\par\medskip\noindent\textbf{EN-NOTE-109} 7II XV; approximate folio 191\par
{\ttfamily\footnotesize At\allowbreak{} the\allowbreak{} definition\allowbreak{} of\allowbreak{} the\allowbreak{} punctured\allowbreak{} henselization,\allowbreak{} the\allowbreak{} print\allowbreak{} places\allowbreak{} a\allowbreak{} tilde\allowbreak{} on\allowbreak{} the\allowbreak{} first\allowbreak{} \$X\textasciicircum{}*\_\{s(x)\}\$.\allowbreak{} The\allowbreak{} equality\allowbreak{} with\allowbreak{} \$X\_\{s(x)\}-\textbackslash{}\{x\textbackslash{}\}\$\allowbreak{} and\allowbreak{} the\allowbreak{} identical\allowbreak{} construction\allowbreak{} in\allowbreak{} 2.2.6\allowbreak{} determine\allowbreak{} the\allowbreak{} untilded\allowbreak{} term\allowbreak{} used\allowbreak{} here.\allowbreak{}}\par
\par\medskip\noindent\textbf{EN-NOTE-110} 7II XVI; approximate folio 202\par
{\ttfamily\footnotesize After\allowbreak{} the\allowbreak{} printed\allowbreak{} word\allowbreak{} ``By,''\allowbreak{} the\allowbreak{} source\allowbreak{} leaves\allowbreak{} the\allowbreak{} expected\allowbreak{} reference\allowbreak{} blank.\allowbreak{} The\allowbreak{} stated\allowbreak{} concentration\allowbreak{} and\allowbreak{} inertia\allowbreak{} action\allowbreak{} are\allowbreak{} the\allowbreak{} ordinary-quadratic\allowbreak{} vanishing-cycle\allowbreak{} calculation\allowbreak{} established\allowbreak{} in\allowbreak{} Expos\textbackslash{}'e\allowbreak{} XV.\allowbreak{}}\par
\par\medskip\noindent\textbf{EN-NOTE-111} 7II XVI; approximate folio 206\par
{\ttfamily\footnotesize The\allowbreak{} print\allowbreak{} has\allowbreak{} a\allowbreak{} minus\allowbreak{} sign\allowbreak{} before\allowbreak{} the\allowbreak{} Swan\allowbreak{} term.\allowbreak{} The\allowbreak{} plus\allowbreak{} sign\allowbreak{} is\allowbreak{} forced\allowbreak{} both\allowbreak{} by\allowbreak{} the\allowbreak{} standard\allowbreak{} Artin-conductor\allowbreak{} formula\allowbreak{} cited\allowbreak{} in\allowbreak{} the\allowbreak{} preceding\allowbreak{} sentence\allowbreak{} and\allowbreak{} by\allowbreak{} the\allowbreak{} printed\allowbreak{} identity\allowbreak{} (c),\allowbreak{} which\allowbreak{} identifies\allowbreak{} the\allowbreak{} alternating\allowbreak{} sum\allowbreak{} with\allowbreak{} the\allowbreak{} \textbackslash{}emph\{total\}\allowbreak{} dimension\allowbreak{} \$\textbackslash{}dim+\textbackslash{}dim\textbackslash{}operatorn\allowbreak{}ame\{Sw\}\$.\allowbreak{}}\par
\par\medskip\noindent\textbf{EN-NOTE-112} 7II XVI; approximate folio 209\par
{\ttfamily\footnotesize The\allowbreak{} print\allowbreak{} says\allowbreak{} \$Y/T\$\allowbreak{} here,\allowbreak{} but\allowbreak{} no\allowbreak{} base\allowbreak{} \$T\$\allowbreak{} is\allowbreak{} introduced;\allowbreak{} Proposition\allowbreak{} 2.5\allowbreak{} defines\allowbreak{} the\allowbreak{} morphism\allowbreak{} \$Y\textbackslash{}to\textbackslash{}mathbb\allowbreak{} P\textasciicircum{}1\_k\$.\allowbreak{}}\par
\par\medskip\noindent\textbf{EN-NOTE-113} 7II XVII; approximate folio 213\par
{\ttfamily\footnotesize The\allowbreak{} print\allowbreak{} first\allowbreak{} says\allowbreak{} ``of\allowbreak{} dimension\allowbreak{} \$n\$''\allowbreak{} and\allowbreak{} then,\allowbreak{} in\allowbreak{} the\allowbreak{} hypersurface\allowbreak{} presentation\allowbreak{} below,\allowbreak{} imposes\allowbreak{} \$n=1+\textbackslash{}dim\allowbreak{} Y\$.\allowbreak{} The\allowbreak{} latter\allowbreak{} is\allowbreak{} the\allowbreak{} convention\allowbreak{} used\allowbreak{} by\allowbreak{} the\allowbreak{} displayed\allowbreak{} Hessian\allowbreak{} and\allowbreak{} normal\allowbreak{} form.\allowbreak{}}\par
\par\medskip\noindent\textbf{EN-NOTE-114} 7II XVII; approximate folio 216\par
{\ttfamily\footnotesize The\allowbreak{} print\allowbreak{} calls\allowbreak{} this\allowbreak{} the\allowbreak{} \$d\$-th\allowbreak{} Segre\allowbreak{} embedding;\allowbreak{} the\allowbreak{} displayed\allowbreak{} monomial\allowbreak{} map\allowbreak{} is\allowbreak{} the\allowbreak{} Veronese\allowbreak{} embedding.\allowbreak{}}\par
\par\medskip\noindent\textbf{EN-NOTE-115} 7II XVII; approximate folio 222\par
{\ttfamily\footnotesize The\allowbreak{} print\allowbreak{} gives\allowbreak{} incompatible\allowbreak{} ranges\allowbreak{} in\allowbreak{} (3.3.11).\allowbreak{} The\allowbreak{} coordinates\allowbreak{} in\allowbreak{} (3.3.5),\allowbreak{} the\allowbreak{} coefficient\allowbreak{} formula\allowbreak{} (3.3.9),\allowbreak{} and\allowbreak{} the\allowbreak{} identity\allowbreak{} block\allowbreak{} of\allowbreak{} the\allowbreak{} printed\allowbreak{} Jacobian\allowbreak{} matrix\allowbreak{} force\allowbreak{} the\allowbreak{} indices\allowbreak{} shown\allowbreak{} here.\allowbreak{}}\par
\par\medskip\noindent\textbf{EN-NOTE-116} 7II XVII; approximate folio 234\par
{\ttfamily\footnotesize The\allowbreak{} print\allowbreak{} has\allowbreak{} \$\textbackslash{}dim\allowbreak{} X\textbackslash{}leq\allowbreak{} r-2\$\allowbreak{} here.\allowbreak{} The\allowbreak{} alternative\allowbreak{} in\allowbreak{} Corollary\allowbreak{} 3.5.0\allowbreak{} and\allowbreak{} the\allowbreak{} criterion\allowbreak{} being\allowbreak{} invoked\allowbreak{} require\allowbreak{} \$\textbackslash{}dim\textbackslash{}check\allowbreak{} X\textbackslash{}leq\allowbreak{} r-2\$.\allowbreak{}}\par
\par\medskip\noindent\textbf{EN-NOTE-117} 7II XVII; approximate folio 236\par
{\ttfamily\footnotesize The\allowbreak{} print\allowbreak{} reads\allowbreak{} ``3.8.1''\allowbreak{} and\allowbreak{} \$F\_2\$\allowbreak{} in\allowbreak{} this\allowbreak{} sentence.\allowbreak{} The\allowbreak{} ordinary-quadratic\allowbreak{} construction\allowbreak{} is\allowbreak{} (3.7.1),\allowbreak{} and\allowbreak{} Lemma\allowbreak{} 4.1.2\allowbreak{} applies\allowbreak{} to\allowbreak{} \$F\_1\$;\allowbreak{} Proposition\allowbreak{} 4.2\allowbreak{} supplies\allowbreak{} the\allowbreak{} required\allowbreak{} assertion\allowbreak{} for\allowbreak{} \$F\_2\$.\allowbreak{}}\par
\par\medskip\noindent\textbf{EN-NOTE-118} 7II XVII; approximate folio 236\par
{\ttfamily\footnotesize The\allowbreak{} print\allowbreak{} has\allowbreak{} \$\textbackslash{}binom\{n+d\}\{d\}-1\$\allowbreak{} here\allowbreak{} and\allowbreak{} \$\textbackslash{}binom\{r+d\}\{d\}-d\$\allowbreak{} in\allowbreak{} Proposition\allowbreak{} 4.3.\allowbreak{} The\allowbreak{} degree-\$d\$\allowbreak{} Veronese\allowbreak{} embedding\allowbreak{} defined\allowbreak{} in\allowbreak{} 2.4\allowbreak{} gives\allowbreak{} \$N=\textbackslash{}binom\{r+d\}\{d\}-1\$\allowbreak{} in\allowbreak{} both\allowbreak{} places.\allowbreak{}}\par
\par\medskip\noindent\textbf{EN-NOTE-119} 7II XVII; approximate folio 240\par
{\ttfamily\footnotesize The\allowbreak{} print\allowbreak{} leaves\allowbreak{} the\allowbreak{} last\allowbreak{} term\allowbreak{} as\allowbreak{} \$\textbackslash{}mathcal\allowbreak{} O\_\{\textbackslash{}mathbb\allowbreak{} P\textasciicircum{}r\}\$\allowbreak{} after\allowbreak{} restricting\allowbreak{} the\allowbreak{} sequence\allowbreak{} to\allowbreak{} \$X\$;\allowbreak{} the\allowbreak{} restricted\allowbreak{} term\allowbreak{} is\allowbreak{} \$\textbackslash{}mathcal\allowbreak{} O\_X\$.\allowbreak{}}\par
\par\medskip\noindent\textbf{EN-NOTE-120} 7II XVII; approximate folio 242\par
{\ttfamily\footnotesize The\allowbreak{} print\allowbreak{} has\allowbreak{} \$\textbackslash{}dim\allowbreak{} X<r-1\$.\allowbreak{} The\allowbreak{} condition\allowbreak{} for\allowbreak{} the\allowbreak{} induced\allowbreak{} map\allowbreak{} to\allowbreak{} the\allowbreak{} dual\allowbreak{} variety\allowbreak{} to\allowbreak{} be\allowbreak{} generically\allowbreak{} finite\allowbreak{} is\allowbreak{} \$\textbackslash{}dim\textbackslash{}check\allowbreak{} X=r-1\$.\allowbreak{}}\par
\par\medskip\noindent\textbf{EN-NOTE-121} 7II XVII; approximate folio 243\par
{\ttfamily\footnotesize The\allowbreak{} print\allowbreak{} numbers\allowbreak{} both\allowbreak{} this\allowbreak{} display\allowbreak{} and\allowbreak{} the\allowbreak{} preceding\allowbreak{} one\allowbreak{} (5.2.6);\allowbreak{} the\allowbreak{} prime\allowbreak{} distinguishes\allowbreak{} the\allowbreak{} second\allowbreak{} occurrence.\allowbreak{}}\par
\par\medskip\noindent\textbf{EN-NOTE-122} 7II XVII; approximate folio 247\par
{\ttfamily\footnotesize The\allowbreak{} printed\allowbreak{} display\allowbreak{} mixes\allowbreak{} the\allowbreak{} indices\allowbreak{} \$i,j\$\allowbreak{} and\allowbreak{} gives\allowbreak{} \$U\_\{j-i\}\$\allowbreak{} on\allowbreak{} the\allowbreak{} right.\allowbreak{} The\allowbreak{} two\allowbreak{} immediately\allowbreak{} preceding\allowbreak{} pushforward\allowbreak{} identities\allowbreak{} force\allowbreak{} the\allowbreak{} coefficient\allowbreak{} \$U\_\{a-1\}\$.\allowbreak{}}\par
\par\medskip\noindent\textbf{EN-NOTE-123} 7II XVII; approximate folio 248\par
{\ttfamily\footnotesize The\allowbreak{} print\allowbreak{} omits\allowbreak{} the\allowbreak{} upper\allowbreak{} summation\allowbreak{} bound\allowbreak{} \$a\$\allowbreak{} in\allowbreak{} (5.4.7);\allowbreak{} the\allowbreak{} relation\allowbreak{} is\allowbreak{} (5.4.2)\allowbreak{} multiplied\allowbreak{} by\allowbreak{} \$T\textasciicircum{}\{j-1\}\$.\allowbreak{}}\par
\par\medskip\noindent\textbf{EN-NOTE-124} 7II XVII; approximate folio 250\par
{\ttfamily\footnotesize The\allowbreak{} print\allowbreak{} has\allowbreak{} \$(-1)\textasciicircum{}r\$.\allowbreak{} Substitution\allowbreak{} of\allowbreak{} (5.5.3)\allowbreak{} and\allowbreak{} (5.5.4)\allowbreak{} into\allowbreak{} (5.5.1)\allowbreak{} does\allowbreak{} not\allowbreak{} alter\allowbreak{} its\allowbreak{} sign\allowbreak{} \$(-1)\textasciicircum{}n\$;\allowbreak{} formula\allowbreak{} (5.7.3)\allowbreak{} immediately\allowbreak{} below\allowbreak{} confirms\allowbreak{} the\allowbreak{} same\allowbreak{} exponent.\allowbreak{}}\par
\par\medskip\noindent\textbf{EN-NOTE-125} 7II XVII; approximate folio 253\par
{\ttfamily\footnotesize The\allowbreak{} print\allowbreak{} repeats\allowbreak{} the\allowbreak{} number\allowbreak{} 6.3\allowbreak{} after\allowbreak{} Proposition\allowbreak{} 6.3.\allowbreak{}}\par
\par\medskip\noindent\textbf{EN-NOTE-126} 7II XVIII; approximate folio 257\par
{\ttfamily\footnotesize In\allowbreak{} (1.2.2)\allowbreak{} the\allowbreak{} print\allowbreak{} has\allowbreak{} the\allowbreak{} upper\allowbreak{} bound\allowbreak{} \$n\$,\allowbreak{} the\allowbreak{} degree\allowbreak{} \$q-2-j\$,\allowbreak{} and\allowbreak{} one\allowbreak{} fewer\allowbreak{} closing\allowbreak{} parenthesis.\allowbreak{} The\allowbreak{} rank\allowbreak{} \$1+r\$,\allowbreak{} the\allowbreak{} defining\allowbreak{} map,\allowbreak{} and\allowbreak{} the\allowbreak{} proof\allowbreak{} immediately\allowbreak{} below\allowbreak{} require\allowbreak{} \$0\textbackslash{}leq\allowbreak{} j\textbackslash{}leq\allowbreak{} r\$\allowbreak{} and\allowbreak{} \$q-2j\$;\allowbreak{} the\allowbreak{} corrected\allowbreak{} formula\allowbreak{} is\allowbreak{} printed\allowbreak{} here.\allowbreak{}}\par
\par\medskip\noindent\textbf{EN-NOTE-127} 7II XVIII; approximate folio 259\par
{\ttfamily\footnotesize The\allowbreak{} print\allowbreak{} labels\allowbreak{} the\allowbreak{} right-hand\allowbreak{} projection\allowbreak{} in\allowbreak{} (2.1.6.2)\allowbreak{} \$\textbackslash{}operatorname\{pr\}\_i\allowbreak{}\$.\allowbreak{} Its\allowbreak{} target\allowbreak{} is\allowbreak{} \$\textbackslash{}Delta\$,\allowbreak{} and\allowbreak{} the\allowbreak{} same\allowbreak{} projection\allowbreak{} is\allowbreak{} labelled\allowbreak{} \$\textbackslash{}operatorname\{pr\}\_1\allowbreak{}\$\allowbreak{} in\allowbreak{} (2.1.6.3);\allowbreak{} the\allowbreak{} corrected\allowbreak{} label\allowbreak{} is\allowbreak{} used\allowbreak{} here.\allowbreak{}}\par
\par\medskip\noindent\textbf{EN-NOTE-128} 7II XVIII; approximate folio 263\par
{\ttfamily\footnotesize At\allowbreak{} this\allowbreak{} point\allowbreak{} the\allowbreak{} print\allowbreak{} writes\allowbreak{} \$d\_s\textasciicircum{}\{p,q\}(f)=0\$,\allowbreak{} although\allowbreak{} the\allowbreak{} inductive\allowbreak{} step\allowbreak{} and\allowbreak{} diagram\allowbreak{} (2.2.14)\allowbreak{} concern\allowbreak{} \$d\_\{s+1\}\textasciicircum{}\{p,q\}(f)\$.\allowbreak{}}\par
\par\medskip\noindent\textbf{EN-NOTE-129} 7II XVIII; approximate folio 264\par
{\ttfamily\footnotesize In\allowbreak{} the\allowbreak{} right-hand\allowbreak{} denominator\allowbreak{} of\allowbreak{} (2.2.21),\allowbreak{} the\allowbreak{} print\allowbreak{} reverses\allowbreak{} the\allowbreak{} notation\allowbreak{} to\allowbreak{} \$H\textasciicircum{}pF\textasciicircum{}p\$.\allowbreak{} The\allowbreak{} filtration\allowbreak{} used\allowbreak{} in\allowbreak{} the\allowbreak{} preceding\allowbreak{} sentence\allowbreak{} and\allowbreak{} on\allowbreak{} the\allowbreak{} left\allowbreak{} is\allowbreak{} \$F\textasciicircum{}pH\textasciicircum{}p\$;\allowbreak{} that\allowbreak{} notation\allowbreak{} is\allowbreak{} restored\allowbreak{} here.\allowbreak{}}\par
\par\medskip\noindent\textbf{EN-NOTE-130} 7II XVIII; approximate folio 266\par
{\ttfamily\footnotesize The\allowbreak{} print\allowbreak{} has\allowbreak{} \$b\_\{q-j\}(\textbackslash{}Delta)\$\allowbreak{} in\allowbreak{} (2.3.1).\allowbreak{} The\allowbreak{} decomposition\allowbreak{} (2.2.2)\allowbreak{} shifts\allowbreak{} cohomological\allowbreak{} degree\allowbreak{} by\allowbreak{} \$2j\$,\allowbreak{} so\allowbreak{} the\allowbreak{} correct\allowbreak{} summand\allowbreak{} is\allowbreak{} \$b\_\{q-2j\}(\textbackslash{}Delta)\$.\allowbreak{}}\par
\par\medskip\noindent\textbf{EN-NOTE-131} 7II XVIII; approximate folio 267\par
{\ttfamily\footnotesize The\allowbreak{} print\allowbreak{} writes\allowbreak{} \$u\$\allowbreak{} for\allowbreak{} the\allowbreak{} dimension\allowbreak{} in\allowbreak{} this\allowbreak{} sentence.\allowbreak{} Section\allowbreak{} 3.2\allowbreak{} immediately\allowbreak{} denotes\allowbreak{} the\allowbreak{} same\allowbreak{} dimension\allowbreak{} by\allowbreak{} \$n\$,\allowbreak{} which\allowbreak{} also\allowbreak{} governs\allowbreak{} the\allowbreak{} signs\allowbreak{} in\allowbreak{} (3.2.3)--(3.2.4);\allowbreak{} that\allowbreak{} notation\allowbreak{} is\allowbreak{} used\allowbreak{} here.\allowbreak{}}\par
\par\medskip\noindent\textbf{EN-NOTE-132} 7II XVIII; approximate folio 269\par
{\ttfamily\footnotesize The\allowbreak{} print\allowbreak{} says\allowbreak{} that\allowbreak{} the\allowbreak{} finite\allowbreak{} parameters\allowbreak{} are\allowbreak{} the\allowbreak{} cube\allowbreak{} roots\allowbreak{} of\allowbreak{} unity.\allowbreak{} With\allowbreak{} the\allowbreak{} printed\allowbreak{} plus\allowbreak{} sign\allowbreak{} in\allowbreak{} (3.2.8),\allowbreak{} the\allowbreak{} Hesse\allowbreak{} factorization\allowbreak{} gives\allowbreak{} \$\textbackslash{}mu\textasciicircum{}3=-1\$;\allowbreak{} cube\allowbreak{} roots\allowbreak{} of\allowbreak{} \$+1\$\allowbreak{} correspond\allowbreak{} to\allowbreak{} the\allowbreak{} polynomial\allowbreak{} with\allowbreak{} coefficient\allowbreak{} \$-3\textbackslash{}mu\$.\allowbreak{}}\par
\par\medskip\noindent\textbf{EN-NOTE-133} 7II XVIII; approximate folio 270\par
{\ttfamily\footnotesize The\allowbreak{} print\allowbreak{} denotes\allowbreak{} the\allowbreak{} local\allowbreak{} length\allowbreak{} by\allowbreak{} \$\textbackslash{}operatorname\{rg\}\$;\allowbreak{} the\allowbreak{} standard\allowbreak{} local-length\allowbreak{} notation\allowbreak{} and\allowbreak{} the\allowbreak{} stated\allowbreak{} intersection-multipl\allowbreak{}icity\allowbreak{} argument\allowbreak{} require\allowbreak{} \$\textbackslash{}operatorname\{lg\}\$\allowbreak{} (or,\allowbreak{} equivalently,\allowbreak{} \$\textbackslash{}operatorname\{lengt\allowbreak{}h\}\$).\allowbreak{}}\par
\par\medskip\noindent\textbf{EN-NOTE-134} 7II XVIII; approximate folio 272\par
{\ttfamily\footnotesize The\allowbreak{} print\allowbreak{} omits\allowbreak{} the\allowbreak{} minus\allowbreak{} sign\allowbreak{} before\allowbreak{} \$g\_*i\textasciicircum{}*\$\allowbreak{} in\allowbreak{} (4.2.2),\allowbreak{} (4.2.6),\allowbreak{} (4.2.9),\allowbreak{} and\allowbreak{} (4.2.10).\allowbreak{} But\allowbreak{} Lemma\allowbreak{} 4.3\allowbreak{} gives\allowbreak{} \$g\_*i\textasciicircum{}*i\_*g\textasciicircum{}*=-\textbackslash{}math\allowbreak{}rm\{id\}\$,\allowbreak{} and\allowbreak{} the\allowbreak{} printed\allowbreak{} product\allowbreak{} calculation\allowbreak{} itself\allowbreak{} uses\allowbreak{} \$-g\_*i\textasciicircum{}*\$.\allowbreak{} The\allowbreak{} sign\allowbreak{} is\allowbreak{} therefore\allowbreak{} restored\allowbreak{} consistently.\allowbreak{}}\par
\par\medskip\noindent\textbf{EN-NOTE-135} 7II XVIII; approximate folio 290\par
{\ttfamily\footnotesize The\allowbreak{} print\allowbreak{} says\allowbreak{} ``condition\allowbreak{} (5.3.2)'',\allowbreak{} but\allowbreak{} (5.3.2)\allowbreak{} is\allowbreak{} the\allowbreak{} displayed\allowbreak{} family\allowbreak{} of\allowbreak{} adjunction\allowbreak{} maps;\allowbreak{} condition\allowbreak{} (A)\allowbreak{} is\allowbreak{} defined\allowbreak{} from\allowbreak{} those\allowbreak{} maps\allowbreak{} in\allowbreak{} 5.3.5.\allowbreak{}}\par
\par\medskip\noindent\textbf{EN-NOTE-136} 7II XVIII; approximate folio 290\par
{\ttfamily\footnotesize The\allowbreak{} printed\allowbreak{} abutment\allowbreak{} has\allowbreak{} \$H\textasciicircum{}\{m+q\}\$;\allowbreak{} the\allowbreak{} bidegree\allowbreak{} and\allowbreak{} the\allowbreak{} parallel\allowbreak{} statement\allowbreak{} in\allowbreak{} Theorem\allowbreak{} 5.6.8\allowbreak{} require\allowbreak{} \$H\textasciicircum{}\{p+q\}\$.\allowbreak{}}\par
\par\medskip\noindent\textbf{EN-NOTE-137} 7II XVIII; approximate folio 292\par
{\ttfamily\footnotesize The\allowbreak{} printed\allowbreak{} (5.6.5)\allowbreak{} omits\allowbreak{} several\allowbreak{} closing\allowbreak{} parentheses\allowbreak{} and\allowbreak{} the\allowbreak{} plus\allowbreak{} sign\allowbreak{} before\allowbreak{} the\allowbreak{} third-summand\allowbreak{} map\allowbreak{} \$h\_*\$.\allowbreak{} The\allowbreak{} three\allowbreak{} displayed\allowbreak{} source\allowbreak{} summands\allowbreak{} and\allowbreak{} (5.6.4)\allowbreak{} force\allowbreak{} the\allowbreak{} corrected\allowbreak{} three-term\allowbreak{} map\allowbreak{} shown\allowbreak{} here.\allowbreak{}}\par
\par\medskip\noindent\textbf{EN-NOTE-138} 7II XVIII; approximate folio 294\par
{\ttfamily\footnotesize The\allowbreak{} printed\allowbreak{} (5.6.8.8)\allowbreak{} has\allowbreak{} an\allowbreak{} odd\allowbreak{} degree\allowbreak{} shift,\allowbreak{} labels\allowbreak{} the\allowbreak{} map\allowbreak{} only\allowbreak{} \$\textbackslash{}xi\_X\$,\allowbreak{} and\allowbreak{} starts\allowbreak{} at\allowbreak{} \$i=1\$.\allowbreak{} The\allowbreak{} corrected\allowbreak{} form\allowbreak{} is\allowbreak{} forced\allowbreak{} by\allowbreak{} (5.6.8.3),\allowbreak{} (5.6.8.6),\allowbreak{} the\allowbreak{} lower\allowbreak{} arrow\allowbreak{} of\allowbreak{} the\allowbreak{} next\allowbreak{} commutative\allowbreak{} square,\allowbreak{} and\allowbreak{} the\allowbreak{} later\allowbreak{} specialization\allowbreak{} \$E\_2\textasciicircum{}\{1,n-2\}\textbackslash{}simeq\allowbreak{} E\_2\textasciicircum{}\{1,n\}\$.\allowbreak{}}\par
\par\medskip\noindent\textbf{EN-NOTE-139} 7II XVIII; approximate folio 295\par
{\ttfamily\footnotesize The\allowbreak{} print\allowbreak{} displays\allowbreak{} the\allowbreak{} first\allowbreak{} three\allowbreak{} primitive\allowbreak{} summands\allowbreak{} while\allowbreak{} its\allowbreak{} arrow\allowbreak{} label\allowbreak{} already\allowbreak{} contains\allowbreak{} an\allowbreak{} ellipsis;\allowbreak{} the\allowbreak{} final\allowbreak{} \$\textbackslash{}oplus\textbackslash{}cdots\$\allowbreak{} makes\allowbreak{} the\allowbreak{} same\allowbreak{} finite\allowbreak{} decomposition\allowbreak{} explicit.\allowbreak{}}\par
\par\medskip\noindent\textbf{EN-NOTE-140} 7II XVIII; approximate folio 299\par
{\ttfamily\footnotesize The\allowbreak{} printed\allowbreak{} label\allowbreak{} in\allowbreak{} (5.7.6)\allowbreak{} has\allowbreak{} \$2\textasciicircum{}*\$\allowbreak{} after\allowbreak{} \$g\_*\$.\allowbreak{} The\allowbreak{} inverse\allowbreak{} blow-up\allowbreak{} decomposition\allowbreak{} (4.2.2),\allowbreak{} used\allowbreak{} here\allowbreak{} verbatim,\allowbreak{} requires\allowbreak{} \$i\textasciicircum{}*\$.\allowbreak{}}\par
\par\medskip\noindent\textbf{EN-NOTE-141} 7II XVIII; approximate folio 300\par
{\ttfamily\footnotesize The\allowbreak{} print\allowbreak{} has\allowbreak{} \$F\textasciicircum{}2H\textasciicircum{}N\$\allowbreak{} in\allowbreak{} the\allowbreak{} first\allowbreak{} sentence\allowbreak{} and\allowbreak{} omits\allowbreak{} the\allowbreak{} Tate\allowbreak{} twist\allowbreak{} \$(-1)\$\allowbreak{} from\allowbreak{} the\allowbreak{} displayed\allowbreak{} source\allowbreak{} of\allowbreak{} the\allowbreak{} Gysin\allowbreak{} map.\allowbreak{} The\allowbreak{} surrounding\allowbreak{} identifications\allowbreak{} and\allowbreak{} the\allowbreak{} degree-two\allowbreak{} Gysin\allowbreak{} shift\allowbreak{} force\allowbreak{} the\allowbreak{} readings\allowbreak{} shown\allowbreak{} here.\allowbreak{}}\par
\par\medskip\noindent\textbf{EN-NOTE-142} 7II XVIII; approximate folio 301\par
{\ttfamily\footnotesize The\allowbreak{} print\allowbreak{} has\allowbreak{} a\allowbreak{} baseline\allowbreak{} dot\allowbreak{} in\allowbreak{} place\allowbreak{} of\allowbreak{} the\allowbreak{} degree\allowbreak{} index\allowbreak{} in\allowbreak{} (5.8.1.6).\allowbreak{} The\allowbreak{} source\allowbreak{} and\allowbreak{} target\allowbreak{} of\allowbreak{} (5.8.1.5)\allowbreak{} force\allowbreak{} \$R\textasciicircum{}\{i-1\}f\_*\$.\allowbreak{}}\par
\par\medskip\noindent\textbf{EN-NOTE-143} 7II XVIII; approximate folio 303\par
{\ttfamily\footnotesize The\allowbreak{} last\allowbreak{} printed\allowbreak{} implication\allowbreak{} has\allowbreak{} a\allowbreak{} stray\allowbreak{} \$\textbackslash{}alpha\$\allowbreak{} after\allowbreak{} \$\textbackslash{}mathbb\allowbreak{} P\textasciicircum{}1\$\allowbreak{} and\allowbreak{} substitutes\allowbreak{} \$\textbackslash{}alpha\$\allowbreak{} for\allowbreak{} the\allowbreak{} class\allowbreak{} \$x\$\allowbreak{} in\allowbreak{} \$\textbackslash{}alpha\textasciicircum{}*f\textasciicircum{}*(x)\$.\allowbreak{} The\allowbreak{} displayed\allowbreak{} equivalence\allowbreak{} and\allowbreak{} the\allowbreak{} preceding\allowbreak{} restriction\allowbreak{} isomorphism\allowbreak{} determine\allowbreak{} the\allowbreak{} corrected\allowbreak{} sentence.\allowbreak{}}\par
\par\medskip\noindent\textbf{EN-NOTE-144} 7II XVIII; approximate folio 304\par
{\ttfamily\footnotesize The\allowbreak{} ink\allowbreak{} beside\allowbreak{} the\allowbreak{} vertical\allowbreak{} arrow\allowbreak{} is\allowbreak{} not\allowbreak{} independently\allowbreak{} legible\allowbreak{} in\allowbreak{} the\allowbreak{} scan.\allowbreak{} The\allowbreak{} map\allowbreak{} is\allowbreak{} nevertheless\allowbreak{} determined\allowbreak{} uniquely:\allowbreak{} (5.8.1.6)\allowbreak{} has\allowbreak{} just\allowbreak{} defined\allowbreak{} \$U\textasciicircum{}i\$\allowbreak{} with\allowbreak{} exactly\allowbreak{} this\allowbreak{} source\allowbreak{} and\allowbreak{} target,\allowbreak{} so\allowbreak{} the\allowbreak{} degree-\$n\$\allowbreak{} label\allowbreak{} is\allowbreak{} \$U\textasciicircum{}n\$.\allowbreak{}}\par
\par\medskip\noindent\textbf{EN-NOTE-145} 7II XVIII; approximate folio 304\par
{\ttfamily\footnotesize The\allowbreak{} print\allowbreak{} omits\allowbreak{} \$\textbackslash{}nu\textasciicircum{}*\$\allowbreak{} from\allowbreak{} the\allowbreak{} lower-left\allowbreak{} sheaf\allowbreak{} in\allowbreak{} (5.8.4.1)\allowbreak{} and\allowbreak{} leaves\allowbreak{} its\allowbreak{} coefficient\allowbreak{} parenthesis\allowbreak{} unclosed.\allowbreak{} Formula\allowbreak{} (5.8.4.3)\allowbreak{} and\allowbreak{} the\allowbreak{} same\allowbreak{} lower\allowbreak{} node\allowbreak{} in\allowbreak{} (5.8.4.4)\allowbreak{} give\allowbreak{} the\allowbreak{} corrected\allowbreak{} reading.\allowbreak{}}\par
\par\medskip\noindent\textbf{EN-NOTE-146} 7II XVIII; approximate folio 305\par
{\ttfamily\footnotesize The\allowbreak{} print\allowbreak{} omits\allowbreak{} the\allowbreak{} twist\allowbreak{} \$(j)\$\allowbreak{} from\allowbreak{} the\allowbreak{} targets\allowbreak{} in\allowbreak{} (5.8.4.5)\allowbreak{} and\allowbreak{} from\allowbreak{} the\allowbreak{} left-hand\allowbreak{} side\allowbreak{} of\allowbreak{} (5.8.4.6),\allowbreak{} although\allowbreak{} the\allowbreak{} domain\allowbreak{} and\allowbreak{} the\allowbreak{} decomposition\allowbreak{} (5.8.4.2)\allowbreak{} retain\allowbreak{} it.\allowbreak{}}\par
\par\medskip\noindent\textbf{EN-NOTE-147} 7II XVIII; approximate folio 309\par
{\ttfamily\footnotesize The\allowbreak{} print\allowbreak{} drops\allowbreak{} the\allowbreak{} twist\allowbreak{} \$(j)\$\allowbreak{} from\allowbreak{} one\allowbreak{} base-change\allowbreak{} display\allowbreak{} and\allowbreak{} the\allowbreak{} frozen\allowbreak{} transcription\allowbreak{} corrupts\allowbreak{} both\allowbreak{} the\allowbreak{} source\allowbreak{} and\allowbreak{} target\allowbreak{} of\allowbreak{} the\allowbreak{} trace\allowbreak{} map.\allowbreak{} The\allowbreak{} corrected\allowbreak{} formulas\allowbreak{} are\allowbreak{} the\allowbreak{} standard\allowbreak{} finite-flat\allowbreak{} pullback--trace\allowbreak{} identities\allowbreak{} used\allowbreak{} by\allowbreak{} the\allowbreak{} concluding\allowbreak{} argument.\allowbreak{}}\par
\par\medskip\noindent\textbf{EN-NOTE-148} 7II XVIII; approximate folio 316\par
{\ttfamily\footnotesize The\allowbreak{} numerator\allowbreak{} of\allowbreak{} the\allowbreak{} printed\allowbreak{} Galois\allowbreak{} group\allowbreak{} is\allowbreak{} blank.\allowbreak{} The\allowbreak{} field\allowbreak{} tower\allowbreak{} immediately\allowbreak{} preceding\allowbreak{} the\allowbreak{} display\allowbreak{} determines\allowbreak{} it\allowbreak{} as\allowbreak{} \$k(\textbackslash{}bar\textbackslash{}eta)\$.\allowbreak{}}\par
\par\medskip\noindent\textbf{EN-NOTE-149} 7II XVIII; approximate folio 318\par
{\ttfamily\footnotesize The\allowbreak{} print\allowbreak{} has\allowbreak{} \$\textbackslash{}operatorname\{card\}\allowbreak{}(k)\textbackslash{}ne2\$;\allowbreak{} the\allowbreak{} parallel\allowbreak{} hypothesis\allowbreak{} in\allowbreak{} (6.3.2)\allowbreak{} and\allowbreak{} the\allowbreak{} tame-monodromy\allowbreak{} argument\allowbreak{} require\allowbreak{} the\allowbreak{} characteristic.\allowbreak{}}\par
\par\medskip\noindent\textbf{EN-NOTE-150} 7II XVIII; approximate folio 320\par
{\ttfamily\footnotesize In\allowbreak{} the\allowbreak{} print,\allowbreak{} the\allowbreak{} placement\allowbreak{} of\allowbreak{} the\allowbreak{} conjunction\allowbreak{} attaches\allowbreak{} condition\allowbreak{} (LV)\allowbreak{} only\allowbreak{} to\allowbreak{} the\allowbreak{} second\allowbreak{} alternative.\allowbreak{} The\allowbreak{} proof\allowbreak{} invokes\allowbreak{} the\allowbreak{} equivalence\allowbreak{} (i)\$\textbackslash{}Longleftrightar\allowbreak{}row\$(vi)\allowbreak{} of\allowbreak{} (6.3.2),\allowbreak{} so\allowbreak{} condition\allowbreak{} (LV)\allowbreak{} is\allowbreak{} required\allowbreak{} in\allowbreak{} either\allowbreak{} case.\allowbreak{}}\par
\par\medskip\noindent\textbf{EN-NOTE-151} 7II XVIII; approximate folio 322\par
{\ttfamily\footnotesize The\allowbreak{} print\allowbreak{} numbers\allowbreak{} the\allowbreak{} reflection\allowbreak{} formula\allowbreak{} (6.5.2.2),\allowbreak{} although\allowbreak{} it\allowbreak{} is\allowbreak{} the\allowbreak{} second\allowbreak{} formula\allowbreak{} of\allowbreak{} 6.5.1\allowbreak{} and\allowbreak{} is\allowbreak{} followed\allowbreak{} by\allowbreak{} (6.5.2.1).\allowbreak{}}\par
\par\medskip\noindent\textbf{EN-NOTE-152} 7II XVIII; approximate folio 324\par
{\ttfamily\footnotesize The\allowbreak{} print\allowbreak{} calls\allowbreak{} \$I\_i\$\allowbreak{} the\allowbreak{} inertia\allowbreak{} group\allowbreak{} ``of\allowbreak{} \$\textbackslash{}delta\_i\$''\allowbreak{} before\allowbreak{} \$\textbackslash{}delta\_i\$\allowbreak{} is\allowbreak{} defined.\allowbreak{} The\allowbreak{} standard\allowbreak{} construction\allowbreak{} is\allowbreak{} the\allowbreak{} inertia\allowbreak{} group\allowbreak{} at\allowbreak{} \$s\_i\$,\allowbreak{} embedded\allowbreak{} by\allowbreak{} the\allowbreak{} chosen\allowbreak{} path\allowbreak{} \$\textbackslash{}ell\_i\$.\allowbreak{}}\par
\par\medskip\noindent\textbf{EN-NOTE-153} 7II XVIII; approximate folio 324\par
{\ttfamily\footnotesize The\allowbreak{} print\allowbreak{} omits\allowbreak{} the\allowbreak{} generic-fiber\allowbreak{} subscript\allowbreak{} in\allowbreak{} this\allowbreak{} display;\allowbreak{} the\allowbreak{} preceding\allowbreak{} action\allowbreak{} and\allowbreak{} the\allowbreak{} definition\allowbreak{} in\allowbreak{} 6.5\allowbreak{} place\allowbreak{} every\allowbreak{} \$\textbackslash{}delta\_i\$\allowbreak{} in\allowbreak{} the\allowbreak{} cohomology\allowbreak{} of\allowbreak{} \$X\_\{\textbackslash{}bar\textbackslash{}eta\}\$.\allowbreak{}}\par
\par\medskip\noindent\textbf{EN-NOTE-154} 7II XVIII; approximate folio 327\par
{\ttfamily\footnotesize The\allowbreak{} print\allowbreak{} gives\allowbreak{} the\allowbreak{} journal\allowbreak{} volume\allowbreak{} as\allowbreak{} 15;\allowbreak{} the\allowbreak{} article\allowbreak{} appeared\allowbreak{} in\allowbreak{} volume\allowbreak{} 90.\allowbreak{}}\par
\par\medskip\noindent\textbf{EN-NOTE-155} 7II XIX; approximate folio 329\par
{\ttfamily\footnotesize The\allowbreak{} print\allowbreak{} writes\allowbreak{} \$d\_i\$\allowbreak{} in\allowbreak{} this\allowbreak{} sum,\allowbreak{} although\allowbreak{} the\allowbreak{} multidegree\allowbreak{} has\allowbreak{} been\allowbreak{} denoted\allowbreak{} throughout\allowbreak{} by\allowbreak{} \$(a\_i)\$.\allowbreak{}}\par
\par\medskip\noindent\textbf{EN-NOTE-156} 7II XIX; approximate folio 331\par
{\ttfamily\footnotesize The\allowbreak{} print\allowbreak{} says\allowbreak{} that\allowbreak{} \$X\$\allowbreak{} has\allowbreak{} even\allowbreak{} dimension\allowbreak{} \$2n\$.\allowbreak{} The\allowbreak{} displayed\allowbreak{} degree\allowbreak{} \$H\textasciicircum{}\{2n\}(H\_s)\$,\allowbreak{} the\allowbreak{} reflection\allowbreak{} formula,\allowbreak{} Remark\allowbreak{} 2.2,\allowbreak{} and\allowbreak{} the\allowbreak{} application\allowbreak{} to\allowbreak{} Theorem\allowbreak{} 1.4\allowbreak{} all\allowbreak{} force\allowbreak{} odd\allowbreak{} dimension\allowbreak{} \$2n+1\$.\allowbreak{}}\par
\par\medskip\noindent\textbf{EN-NOTE-157} 7II XIX; approximate folio 332\par
{\ttfamily\footnotesize The\allowbreak{} print\allowbreak{} gives\allowbreak{} \$H\textasciicircum{}n(H\_s)\$\allowbreak{} at\allowbreak{} this\allowbreak{} point,\allowbreak{} although\allowbreak{} every\allowbreak{} surrounding\allowbreak{} occurrence\allowbreak{} and\allowbreak{} the\allowbreak{} statement\allowbreak{} of\allowbreak{} the\allowbreak{} proposition\allowbreak{} use\allowbreak{} the\allowbreak{} middle\allowbreak{} degree\allowbreak{} \$H\textasciicircum{}\{2n\}(H\_s)\$.\allowbreak{}}\par
\par\medskip\noindent\textbf{EN-NOTE-158} 7II XIX; approximate folio 332\par
{\ttfamily\footnotesize The\allowbreak{} print\allowbreak{} uses\allowbreak{} the\allowbreak{} undefined\allowbreak{} exponent\allowbreak{} \$m\$\allowbreak{} in\allowbreak{} this\allowbreak{} second\allowbreak{} occurrence;\allowbreak{} the\allowbreak{} immediately\allowbreak{} preceding\allowbreak{} reflection\allowbreak{} and\allowbreak{} the\allowbreak{} Picard--Lefschetz\allowbreak{} formula\allowbreak{} give\allowbreak{} \$n\$.\allowbreak{}}\par
\par\medskip\noindent\textbf{EN-NOTE-159} 7II XIX; approximate folio 332\par
{\ttfamily\footnotesize The\allowbreak{} print\allowbreak{} writes\allowbreak{} \$d\_i\$\allowbreak{} in\allowbreak{} this\allowbreak{} sum,\allowbreak{} although\allowbreak{} the\allowbreak{} multidegree\allowbreak{} is\allowbreak{} denoted\allowbreak{} by\allowbreak{} \$(a\_i)\$.\allowbreak{}}\par
\par\medskip\noindent\textbf{EN-NOTE-160} 7II XIX; approximate folio 333\par
{\ttfamily\footnotesize The\allowbreak{} print\allowbreak{} repeatedly\allowbreak{} says\allowbreak{} ``groupe\allowbreak{} de\allowbreak{} Weil'';\allowbreak{} the\allowbreak{} group\allowbreak{} attached\allowbreak{} to\allowbreak{} a\allowbreak{} root\allowbreak{} system\allowbreak{} is\allowbreak{} its\allowbreak{} Weyl\allowbreak{} group.\allowbreak{}}\par
\par\medskip\noindent\textbf{EN-NOTE-161} 7II XIX; approximate folio 336\par
{\ttfamily\footnotesize The\allowbreak{} print\allowbreak{} gives\allowbreak{} \$\textbackslash{}ell-S\$\allowbreak{} as\allowbreak{} the\allowbreak{} domain\allowbreak{} of\allowbreak{} \$\textbackslash{}alpha\$,\allowbreak{} although\allowbreak{} the\allowbreak{} same\allowbreak{} sentence\allowbreak{} evaluates\allowbreak{} \$\textbackslash{}alpha\$\allowbreak{} on\allowbreak{} \$S\$;\allowbreak{} the\allowbreak{} required\allowbreak{} domain\allowbreak{} is\allowbreak{} \$\textbackslash{}ell\$.\allowbreak{}}\par
\par\medskip\noindent\textbf{EN-NOTE-162} 7II XIX; approximate folio 337\par
{\ttfamily\footnotesize The\allowbreak{} print\allowbreak{} writes\allowbreak{} \$V\_0\$\allowbreak{} in\allowbreak{} this\allowbreak{} sequence\allowbreak{} and\allowbreak{} in\allowbreak{} the\allowbreak{} diagram\allowbreak{} below,\allowbreak{} but\allowbreak{} Section\allowbreak{} 4\allowbreak{} defines\allowbreak{} the\allowbreak{} fiber\allowbreak{} as\allowbreak{} \$H\_0\$.\allowbreak{}}\par
\par\medskip\noindent\textbf{EN-NOTE-163} 7II XIX; approximate folio 338\par
{\ttfamily\footnotesize The\allowbreak{} print\allowbreak{} has\allowbreak{} \$\textbackslash{}eta\textasciicircum{}\{2n\}\$\allowbreak{} in\allowbreak{} the\allowbreak{} primitivity\allowbreak{} assertion,\allowbreak{} but\allowbreak{} that\allowbreak{} class\allowbreak{} lies\allowbreak{} in\allowbreak{} degree\allowbreak{} \$4n\$;\allowbreak{} both\allowbreak{} the\allowbreak{} surrounding\allowbreak{} discussion\allowbreak{} and\allowbreak{} the\allowbreak{} next\allowbreak{} display\allowbreak{} require\allowbreak{} \$\textbackslash{}eta\textasciicircum{}n\textbackslash{}in\allowbreak{} H\textasciicircum{}\{2n\}(V)\$.\allowbreak{}}\par
\par\medskip\noindent\textbf{EN-NOTE-164} 7II XIX; approximate folio 339\par
{\ttfamily\footnotesize The\allowbreak{} print\allowbreak{} refers\allowbreak{} here\allowbreak{} to\allowbreak{} 4.2(ii)\allowbreak{} and\allowbreak{} 4.2(iii),\allowbreak{} but\allowbreak{} the\allowbreak{} assertions\allowbreak{} cited\allowbreak{} are\allowbreak{} the\allowbreak{} immediately\allowbreak{} preceding\allowbreak{} cases\allowbreak{} of\allowbreak{} Proposition\allowbreak{} 5.2.\allowbreak{}}\par
\par\medskip\noindent\textbf{EN-NOTE-165} 7II XX; approximate folio 341\par
{\ttfamily\footnotesize ;\allowbreak{} the\allowbreak{} further\allowbreak{} citation\allowbreak{} is\allowbreak{} left\allowbreak{} blank\allowbreak{} in\allowbreak{} the\allowbreak{} print\allowbreak{}}\par
\par\medskip\noindent\textbf{EN-NOTE-166} 7II XX; approximate folio 341\par
{\ttfamily\footnotesize In\allowbreak{} (1.0.6)\allowbreak{} the\allowbreak{} print\allowbreak{} has\allowbreak{} an\allowbreak{} apostrophe\allowbreak{} in\allowbreak{} place\allowbreak{} of\allowbreak{} the\allowbreak{} closing\allowbreak{} parenthesis\allowbreak{} in\allowbreak{} the\allowbreak{} exponent\allowbreak{} and\allowbreak{} has\allowbreak{} \$i-d\_1-d\_2\$\allowbreak{} in\allowbreak{} the\allowbreak{} target\allowbreak{} twist.\allowbreak{} Formulae\allowbreak{} (1.0.5)\allowbreak{} and\allowbreak{} (1.0.7),\allowbreak{} as\allowbreak{} well\allowbreak{} as\allowbreak{} the\allowbreak{} Gysin\allowbreak{} degree,\allowbreak{} give\allowbreak{} the\allowbreak{} reading\allowbreak{} displayed\allowbreak{} here.\allowbreak{}}\par
\par\medskip\noindent\textbf{EN-NOTE-167} 7II XX; approximate folio 342\par
{\ttfamily\footnotesize The\allowbreak{} print\allowbreak{} omits\allowbreak{} the\allowbreak{} prime\allowbreak{} on\allowbreak{} the\allowbreak{} right-hand\allowbreak{} differential\allowbreak{} in\allowbreak{} (1.0.9).\allowbreak{} The\allowbreak{} definition\allowbreak{} of\allowbreak{} the\allowbreak{} shift\allowbreak{} of\allowbreak{} the\allowbreak{} primed\allowbreak{} spectral\allowbreak{} sequence\allowbreak{} requires\allowbreak{} the\allowbreak{} primed\allowbreak{} differential.\allowbreak{}}\par
\par\medskip\noindent\textbf{EN-NOTE-168} 7II XX; approximate folio 343\par
{\ttfamily\footnotesize The\allowbreak{} print\allowbreak{} has\allowbreak{} \$(a,-2(d\_1-d\_2))\$.\allowbreak{} The\allowbreak{} map\allowbreak{} (1.0.7)\allowbreak{} leaves\allowbreak{} the\allowbreak{} first\allowbreak{} Leray\allowbreak{} index\allowbreak{} \$p\$\allowbreak{} unchanged,\allowbreak{} so\allowbreak{} the\allowbreak{} first\allowbreak{} component\allowbreak{} of\allowbreak{} the\allowbreak{} bidegree\allowbreak{} is\allowbreak{} \$0\$.\allowbreak{}}\par
\par\medskip\noindent\textbf{EN-NOTE-169} 7II XX; approximate folio 343\par
{\ttfamily\footnotesize The\allowbreak{} print\allowbreak{} begins\allowbreak{} this\allowbreak{} range\allowbreak{} at\allowbreak{} the\allowbreak{} nonexistent\allowbreak{} (1.0.2);\allowbreak{} the\allowbreak{} displayed\allowbreak{} situation\allowbreak{} is\allowbreak{} (1.0.1).\allowbreak{}}\par
\par\medskip\noindent\textbf{EN-NOTE-170} 7II XX; approximate folio 344\par
{\ttfamily\footnotesize The\allowbreak{} print\allowbreak{} has\allowbreak{} \$(0,b-2d\_1)\$.\allowbreak{} The\allowbreak{} maps\allowbreak{} (1.2.3),\allowbreak{} (1.2.5),\allowbreak{} and\allowbreak{} (1.2.6)\allowbreak{} change\allowbreak{} cohomological\allowbreak{} degree\allowbreak{} by\allowbreak{} \$a-2d\_1\$;\allowbreak{} \$b\$\allowbreak{} controls\allowbreak{} the\allowbreak{} Tate\allowbreak{} twist.\allowbreak{}}\par
\par\medskip\noindent\textbf{EN-NOTE-171} 7II XX; approximate folio 345\par
{\ttfamily\footnotesize In\allowbreak{} (1.2.9)\allowbreak{} and\allowbreak{} (1.2.10)\allowbreak{} the\allowbreak{} print\allowbreak{} writes\allowbreak{} \$(\textbackslash{}mathrm\{pr\}\_1\allowbreak{} f)\_*\$,\allowbreak{} whereas\allowbreak{} the\allowbreak{} map\allowbreak{} to\allowbreak{} \$S\$\allowbreak{} is\allowbreak{} \$f\textbackslash{}circ\textbackslash{}mathrm\{pr\}\_1\allowbreak{}\$,\allowbreak{} as\allowbreak{} the\allowbreak{} printed\allowbreak{} diagram\allowbreak{} (1.2.11)\allowbreak{} itself\allowbreak{} shows.\allowbreak{}}\par
\par\medskip\noindent\textbf{EN-NOTE-172} 7II XX; approximate folio 346\par
{\ttfamily\footnotesize The\allowbreak{} print\allowbreak{} has\allowbreak{} the\allowbreak{} nonsensical\allowbreak{} adjective\allowbreak{} ``annexe''\allowbreak{} in\allowbreak{} place\allowbreak{} of\allowbreak{} ``connexe'',\allowbreak{} and\allowbreak{} uses\allowbreak{} the\allowbreak{} capital\allowbreak{} \$X\$\allowbreak{} for\allowbreak{} the\allowbreak{} element\allowbreak{} in\allowbreak{} the\allowbreak{} conclusion.\allowbreak{} The\allowbreak{} proof\allowbreak{} requires\allowbreak{} connectedness\allowbreak{} of\allowbreak{} \$S\$,\allowbreak{} while\allowbreak{} (1.5.4)\allowbreak{} itself\allowbreak{} uses\allowbreak{} the\allowbreak{} element\allowbreak{} \$x\$.\allowbreak{}}\par
\par\medskip\noindent\textbf{EN-NOTE-173} 7II XX; approximate folio 347\par
{\ttfamily\footnotesize The\allowbreak{} print\allowbreak{} says\allowbreak{} \$f(s)\$\allowbreak{} and\allowbreak{} \$g(s)\$\allowbreak{} here,\allowbreak{} although\allowbreak{} the\allowbreak{} two\allowbreak{} maps\allowbreak{} under\allowbreak{} discussion\allowbreak{} are\allowbreak{} the\allowbreak{} sections\allowbreak{} \$t\_1\$\allowbreak{} and\allowbreak{} \$t\_2\$.\allowbreak{}}\par
\par\medskip\noindent\textbf{EN-NOTE-174} 7II XX; approximate folio 348\par
{\ttfamily\footnotesize ;\allowbreak{} the\allowbreak{} further\allowbreak{} citation\allowbreak{} is\allowbreak{} blank\allowbreak{} in\allowbreak{} the\allowbreak{} print\allowbreak{}}\par
\par\medskip\noindent\textbf{EN-NOTE-175} 7II XX; approximate folio 348\par
{\ttfamily\footnotesize ;\allowbreak{} the\allowbreak{} further\allowbreak{} citation\allowbreak{} is\allowbreak{} blank\allowbreak{} in\allowbreak{} the\allowbreak{} print\allowbreak{}}\par
\par\medskip\noindent\textbf{EN-NOTE-176} 7II XX; approximate folio 349\par
{\ttfamily\footnotesize The\allowbreak{} sentence\allowbreak{} in\allowbreak{} the\allowbreak{} print\allowbreak{} uses\allowbreak{} an\allowbreak{} upper-case\allowbreak{} \$Z\_\{\textbackslash{}bar\textbackslash{}eta\}\$,\allowbreak{} whereas\allowbreak{} (1.7.2)\allowbreak{} and\allowbreak{} the\allowbreak{} cycle\allowbreak{} used\allowbreak{} throughout\allowbreak{} the\allowbreak{} section\allowbreak{} are\allowbreak{} denoted\allowbreak{} by\allowbreak{} \$z\_\{\textbackslash{}bar\textbackslash{}eta\}\$.\allowbreak{}}\par
\par\medskip\noindent\textbf{EN-NOTE-177} 7II XX; approximate folio 352\par
{\ttfamily\footnotesize The\allowbreak{} print\allowbreak{} has\allowbreak{} \$E\textasciicircum{}\{2n-1\}\$\allowbreak{} in\allowbreak{} (3.1.6).\allowbreak{} The\allowbreak{} space\allowbreak{} defined\allowbreak{} in\allowbreak{} (3.1.5),\allowbreak{} and\allowbreak{} used\allowbreak{} throughout\allowbreak{} the\allowbreak{} proof\allowbreak{} of\allowbreak{} Proposition\textasciitilde{}3.2,\allowbreak{} is\allowbreak{} \$E\textasciicircum{}\{n-1\}\$.\allowbreak{}}\par
\par\medskip\noindent\textbf{EN-NOTE-178} 7II XX; approximate folio 355\par
{\ttfamily\footnotesize The\allowbreak{} print\allowbreak{} writes\allowbreak{} \$m\textasciicircum{}*\$\allowbreak{} before\allowbreak{} this\allowbreak{} geometric\allowbreak{} inclusion.\allowbreak{} Condition\allowbreak{} (B)\allowbreak{} in\allowbreak{} 3.3\allowbreak{} is\allowbreak{} imposed\allowbreak{} on\allowbreak{} \$m\$;\allowbreak{} \$m\textasciicircum{}*\$\allowbreak{} denotes\allowbreak{} its\allowbreak{} induced\allowbreak{} map\allowbreak{} in\allowbreak{} cohomology.\allowbreak{}}\par
\par\medskip\noindent\textbf{EN-NOTE-179} 7II XX; approximate folio 355\par
{\ttfamily\footnotesize The\allowbreak{} print\allowbreak{} has\allowbreak{} \$\textbackslash{}operatorname\{char\}\allowbreak{}(k)\textbackslash{}ne0\$\allowbreak{} in\allowbreak{} this\allowbreak{} Hodge-theoretic\allowbreak{} clause.\allowbreak{} The\allowbreak{} following\allowbreak{} clause\allowbreak{} treats\allowbreak{} positive\allowbreak{} characteristic,\allowbreak{} so\allowbreak{} the\allowbreak{} required\allowbreak{} reading\allowbreak{} here\allowbreak{} is\allowbreak{} \$\textbackslash{}operatorname\{char\}\allowbreak{}(k)=0\$.\allowbreak{}}\par
\par\medskip\noindent\textbf{EN-NOTE-180} 7II XX; approximate folio 356\par
{\ttfamily\footnotesize The\allowbreak{} print\allowbreak{} writes\allowbreak{} target\allowbreak{} \$C\$\allowbreak{} rather\allowbreak{} than\allowbreak{} \$C\_\{\textbackslash{}bar\textbackslash{}eta\}\$\allowbreak{} in\allowbreak{} (4.4.3),\allowbreak{} and\allowbreak{} both\allowbreak{} above\allowbreak{} and\allowbreak{} in\allowbreak{} (4.4.4)\allowbreak{} it\allowbreak{} writes\allowbreak{} codimension\allowbreak{} \$n\$\allowbreak{} rather\allowbreak{} than\allowbreak{} \$m\$;\allowbreak{} the\allowbreak{} defined\allowbreak{} curve,\allowbreak{} Theorem\textasciitilde{}4.1,\allowbreak{} and\allowbreak{} the\allowbreak{} degree\allowbreak{} shift\allowbreak{} below\allowbreak{} force\allowbreak{} the\allowbreak{} readings\allowbreak{} used\allowbreak{} here.\allowbreak{}}\par
\par\medskip\noindent\textbf{EN-NOTE-181} 7II XX; approximate folio 357\par
{\ttfamily\footnotesize The\allowbreak{} print\allowbreak{} omits\allowbreak{} the\allowbreak{} final\allowbreak{} closing\allowbreak{} parenthesis\allowbreak{} in\allowbreak{} the\allowbreak{} target\allowbreak{} of\allowbreak{} (4.4.6).\allowbreak{}}\par
\par\medskip\noindent\textbf{EN-NOTE-182} 7II XX; approximate folio 357\par
{\ttfamily\footnotesize The\allowbreak{} print\allowbreak{} repeats\allowbreak{} an\allowbreak{} extraneous\allowbreak{} \$\textbackslash{}bar\textbackslash{}eta\$\allowbreak{} after\allowbreak{} \$m\textasciicircum{}*Q\_\{\textbackslash{}bar\textbackslash{}eta\}\$\allowbreak{} in\allowbreak{} the\allowbreak{} preceding\allowbreak{} sentence.\allowbreak{}}\par
\par\medskip\noindent\textbf{EN-NOTE-183} 7II XX; approximate folio 359\par
{\ttfamily\footnotesize The\allowbreak{} print\allowbreak{} gives\allowbreak{} the\allowbreak{} twist\allowbreak{} \$n\$\allowbreak{} in\allowbreak{} the\allowbreak{} domain\allowbreak{} of\allowbreak{} the\allowbreak{} Griffiths\allowbreak{} map;\allowbreak{} the\allowbreak{} codimension\allowbreak{} and\allowbreak{} all\allowbreak{} adjacent\allowbreak{} terms\allowbreak{} require\allowbreak{} \$m\$.\allowbreak{}}\par
\par\medskip\noindent\textbf{EN-NOTE-184} 7II XX; approximate folio 360\par
{\ttfamily\footnotesize The\allowbreak{} print\allowbreak{} has\allowbreak{} \$2n-1\$\allowbreak{} and\allowbreak{} twist\allowbreak{} \$n\$\allowbreak{} in\allowbreak{} the\allowbreak{} last\allowbreak{} group\allowbreak{} and\allowbreak{} cites\allowbreak{} the\allowbreak{} nonexistent\allowbreak{} (4.3.12);\allowbreak{} the\allowbreak{} coherent\allowbreak{} readings\allowbreak{} are\allowbreak{} \$2m-1\$,\allowbreak{} twist\allowbreak{} \$m\$,\allowbreak{} and\allowbreak{} (4.4.12).\allowbreak{}}\par
\par\medskip\noindent\textbf{EN-NOTE-185} 7II XX; approximate folio 362\par
{\ttfamily\footnotesize The\allowbreak{} print\allowbreak{} has\allowbreak{} twist\allowbreak{} \$n\$\allowbreak{} here\allowbreak{} and\allowbreak{} omits\allowbreak{} the\allowbreak{} outer\allowbreak{} closing\allowbreak{} parenthesis;\allowbreak{} the\allowbreak{} middle\allowbreak{} degree\allowbreak{} \$2m\$\allowbreak{} forces\allowbreak{} the\allowbreak{} twist\allowbreak{} \$m\$.\allowbreak{}}\par
\par\medskip\noindent\textbf{EN-NOTE-186} 7II XX; approximate folio 362\par
{\ttfamily\footnotesize The\allowbreak{} print\allowbreak{} inserts\allowbreak{} the\allowbreak{} literal\allowbreak{} parenthetical\allowbreak{} ``(4.2)''\allowbreak{} at\allowbreak{} this\allowbreak{} point.\allowbreak{} Direct\allowbreak{} inspection\allowbreak{} of\allowbreak{} the\allowbreak{} native\allowbreak{} scan\allowbreak{} at\allowbreak{} 9000-dpi-equivalent\allowbreak{} enlargement\allowbreak{} confirms\allowbreak{} the\allowbreak{} glyphs;\allowbreak{} the\allowbreak{} numerical\allowbreak{} specialization\allowbreak{} from\allowbreak{} \$\textbackslash{}mathbb\allowbreak{} P\textasciicircum{}\{2m+1\}\$\allowbreak{} to\allowbreak{} \$\textbackslash{}mathbb\allowbreak{} P\textasciicircum{}5\$\allowbreak{} is\allowbreak{} \$m=2\$.\allowbreak{}}\par
\par\medskip\noindent\textbf{EN-NOTE-187} 7II XXI; approximate folio 364\par
{\ttfamily\footnotesize The\allowbreak{} print\allowbreak{} has\allowbreak{} the\allowbreak{} evident\allowbreak{} typographical\allowbreak{} forms\allowbreak{} \textbackslash{}emph\{et\allowbreak{} de\allowbreak{} démontrer\}\allowbreak{} and\allowbreak{} \textbackslash{}emph\{suffisament\};\allowbreak{} the\allowbreak{} readings\allowbreak{} used\allowbreak{} above\allowbreak{} are\allowbreak{} \textbackslash{}emph\{est\allowbreak{} de\allowbreak{} démontrer\}\allowbreak{} and\allowbreak{} \textbackslash{}emph\{suffisamment\}.\allowbreak{}}\par
\par\medskip\noindent\textbf{EN-NOTE-188} 7II XXI; approximate folio 365\par
{\ttfamily\footnotesize The\allowbreak{} print\allowbreak{} calls\allowbreak{} this\allowbreak{} ``Appendix\allowbreak{} I,''\allowbreak{} while\allowbreak{} the\allowbreak{} title\allowbreak{} and\allowbreak{} the\allowbreak{} document\allowbreak{} structure\allowbreak{} identify\allowbreak{} Deligne's\allowbreak{} Appendix\allowbreak{} as\allowbreak{} \textbackslash{}S\allowbreak{} 5.\allowbreak{}}\par
\par\medskip\noindent\textbf{EN-NOTE-189} 7II XXI; approximate folio 368\par
{\ttfamily\footnotesize \textbackslash{}medskip\allowbreak{} \textbackslash{}noindent\textbackslash{}emph\{Sourc\allowbreak{}e\allowbreak{} note.\}\allowbreak{} The\allowbreak{} diplomatic\allowbreak{} French\allowbreak{} retains\allowbreak{} the\allowbreak{} printed\allowbreak{} forms\allowbreak{} \textbackslash{}emph\{parmi\allowbreak{} toute\allowbreak{} les\},\allowbreak{} \textbackslash{}emph\{théroème\},\allowbreak{} \textbackslash{}emph\{intersection\allowbreak{} complète\allowbreak{} vérifiante\},\allowbreak{} and\allowbreak{} \textbackslash{}emph\{n'a\allowbreak{} pas\allowbreak{} des\allowbreak{} points\};\allowbreak{} the\allowbreak{} editorial\allowbreak{} French\allowbreak{} and\allowbreak{} the\allowbreak{} translation\allowbreak{} use\allowbreak{} the\allowbreak{} grammatically\allowbreak{} coherent\allowbreak{} readings\allowbreak{} recorded\allowbreak{} in\allowbreak{} the\allowbreak{} source-defect\allowbreak{} register.\allowbreak{} }\par
\par\medskip\noindent\textbf{EN-NOTE-190} 7II XXI; approximate folio 370\par
{\ttfamily\footnotesize \textbackslash{}medskip\allowbreak{} \textbackslash{}noindent\textbackslash{}emph\{Sourc\allowbreak{}e\allowbreak{} note.\}\allowbreak{} The\allowbreak{} diplomatic\allowbreak{} French\allowbreak{} retains\allowbreak{} the\allowbreak{} printed\allowbreak{} misplaced\allowbreak{} parenthesis\allowbreak{} in\allowbreak{} (1.2.1),\allowbreak{} the\allowbreak{} literal\allowbreak{} bound\allowbreak{} \$d\_1\textbackslash{}geq\allowbreak{} a\_o\$\allowbreak{} in\allowbreak{} (1.2.6),\allowbreak{} and\allowbreak{} the\allowbreak{} missing\allowbreak{} \$+1\$\allowbreak{} in\allowbreak{} the\allowbreak{} terminal\allowbreak{} \$H\_r\$\allowbreak{} block.\allowbreak{} Direct\allowbreak{} native-pixel\allowbreak{} inspection\allowbreak{} confirms\allowbreak{} these\allowbreak{} as\allowbreak{} print\allowbreak{} features;\allowbreak{} the\allowbreak{} construction\allowbreak{} forces\allowbreak{} the\allowbreak{} coherent\allowbreak{} forms\allowbreak{} displayed\allowbreak{} above.\allowbreak{} }\par
\par\medskip\noindent\textbf{EN-NOTE-191} 7II XXI; approximate folio 372\par
{\ttfamily\footnotesize \textbackslash{}medskip\allowbreak{} \textbackslash{}noindent\textbackslash{}emph\{Sourc\allowbreak{}e\allowbreak{} note.\}\allowbreak{} The\allowbreak{} diplomatic\allowbreak{} French\allowbreak{} retains\allowbreak{} the\allowbreak{} printed\allowbreak{} exponent\allowbreak{} \$p\textasciicircum{}n\$,\allowbreak{} the\allowbreak{} base\allowbreak{} sheaf\allowbreak{} \$\textbackslash{}mathcal\allowbreak{} O\_u\$,\allowbreak{} and\allowbreak{} the\allowbreak{} shortened\allowbreak{} citation\allowbreak{} ``(1.5).''\allowbreak{} The\allowbreak{} rank-\$M\$\allowbreak{} nilpotence\allowbreak{} test,\allowbreak{} the\allowbreak{} base\allowbreak{} \$U\$,\allowbreak{} and\allowbreak{} Proposition\textasciitilde{}1.1.5\allowbreak{} force\allowbreak{} the\allowbreak{} coherent\allowbreak{} readings\allowbreak{} used\allowbreak{} above.\allowbreak{} The\allowbreak{} closing\allowbreak{} parenthesis\allowbreak{} in\allowbreak{} (1.3.1),\allowbreak{} absent\allowbreak{} only\allowbreak{} from\allowbreak{} the\allowbreak{} inherited\allowbreak{} TeX,\allowbreak{} has\allowbreak{} been\allowbreak{} restored\allowbreak{} directly\allowbreak{} from\allowbreak{} the\allowbreak{} scan.\allowbreak{} }\par
\par\medskip\noindent\textbf{EN-NOTE-192} 7II XXI; approximate folio 374\par
{\ttfamily\footnotesize \textbackslash{}medskip\allowbreak{} \textbackslash{}noindent\textbackslash{}emph\{Sourc\allowbreak{}e\allowbreak{} note.\}\allowbreak{} The\allowbreak{} diplomatic\allowbreak{} French\allowbreak{} retains\allowbreak{} the\allowbreak{} printed\allowbreak{} top-right\allowbreak{} object\allowbreak{} \$(\textbackslash{}widetilde\allowbreak{} X.H\_T)\textbackslash{}times\_U\allowbreak{} T\$\allowbreak{} and\allowbreak{} the\allowbreak{} phrase\allowbreak{} \textbackslash{}emph\{se\allowbreak{} formation\}.\allowbreak{} The\allowbreak{} cartesian\allowbreak{} square\allowbreak{} and\allowbreak{} the\allowbreak{} definitions\allowbreak{} of\allowbreak{} \$T\$\allowbreak{} and\allowbreak{} its\allowbreak{} generic\allowbreak{} fiber\allowbreak{} \$U\$\allowbreak{} force\allowbreak{} \$(\textbackslash{}widetilde\allowbreak{} X.H\_T)\textbackslash{}times\_T\allowbreak{} U\$\allowbreak{} and\allowbreak{} \textbackslash{}emph\{sa\allowbreak{} formation\},\allowbreak{} as\allowbreak{} used\allowbreak{} above.\allowbreak{} The\allowbreak{} plain\allowbreak{} \$E\$,\allowbreak{} the\allowbreak{} exponent\allowbreak{} \$0\$\allowbreak{} in\allowbreak{} \$H\textasciicircum{}0\$,\allowbreak{} and\allowbreak{} the\allowbreak{} outer\allowbreak{} parenthesis\allowbreak{} defining\allowbreak{} \$U\$\allowbreak{} are\allowbreak{} transcription\allowbreak{} repairs\allowbreak{} made\allowbreak{} directly\allowbreak{} from\allowbreak{} the\allowbreak{} scan\allowbreak{} rather\allowbreak{} than\allowbreak{} editorial\allowbreak{} emendations.\allowbreak{} }\par
\par\medskip\noindent\textbf{EN-NOTE-193} 7II XXI; approximate folio 376\par
{\ttfamily\footnotesize \textbackslash{}medskip\allowbreak{} \textbackslash{}noindent\textbackslash{}emph\{Sourc\allowbreak{}e\allowbreak{} note.\}\allowbreak{} The\allowbreak{} diplomatic\allowbreak{} French\allowbreak{} restores\allowbreak{} the\allowbreak{} scan's\allowbreak{} \textbackslash{}emph\{coordonnées\},\allowbreak{} the\allowbreak{} two\allowbreak{} printed\allowbreak{} equality\allowbreak{} signs\allowbreak{} in\allowbreak{} (1.4.6),\allowbreak{} and\allowbreak{} the\allowbreak{} three-dot\allowbreak{} ellipsis\allowbreak{} in\allowbreak{} (1.4.8),\allowbreak{} all\allowbreak{} of\allowbreak{} which\allowbreak{} differ\allowbreak{} from\allowbreak{} the\allowbreak{} inherited\allowbreak{} transcription.\allowbreak{} The\allowbreak{} print\allowbreak{} uses\allowbreak{} the\allowbreak{} letter-like\allowbreak{} subscript\allowbreak{} \$o\$\allowbreak{} for\allowbreak{} the\allowbreak{} zero\allowbreak{} indices\allowbreak{} in\allowbreak{} \$d\_o\$\allowbreak{} and\allowbreak{} \$X\_o\$,\allowbreak{} says\allowbreak{} \textbackslash{}emph\{une\allowbreak{} \$H\$\allowbreak{} défini\},\allowbreak{} and\allowbreak{} uses\allowbreak{} irregular\allowbreak{} point-count\allowbreak{} parentheses.\allowbreak{} The\allowbreak{} corrected\allowbreak{} French\allowbreak{} and\allowbreak{} English\allowbreak{} regularize\allowbreak{} these\allowbreak{} as\allowbreak{} \$d\_0\$,\allowbreak{} \$X\_0\$,\allowbreak{} masculine\allowbreak{} agreement,\allowbreak{} and\allowbreak{} ordinary\allowbreak{} evaluation\allowbreak{} notation\allowbreak{} without\allowbreak{} changing\allowbreak{} the\allowbreak{} argument.\allowbreak{} }\par
\par\medskip\noindent\textbf{EN-NOTE-194} 7II XXI; approximate folio 378\par
{\ttfamily\footnotesize \textbackslash{}medskip\allowbreak{} \textbackslash{}noindent\textbackslash{}emph\{Sourc\allowbreak{}e\allowbreak{} note.\}\allowbreak{} The\allowbreak{} diplomatic\allowbreak{} French\allowbreak{} preserves\allowbreak{} the\allowbreak{} printed\allowbreak{} terms\allowbreak{} \textbackslash{}emph\{\$\textbackslash{}mathbb\allowbreak{} Q\_\textbackslash{}ell\$-vectoriels\},\allowbreak{} \textbackslash{}emph\{fonction\allowbreak{} zêta\allowbreak{} ...\allowbreak{} donné\},\allowbreak{} the\allowbreak{} overstruck\allowbreak{} phrase\allowbreak{} following\allowbreak{} the\allowbreak{} extension\allowbreak{} \$\textbackslash{}mathbb\allowbreak{} F\_\{q\textasciicircum{}M\}/\textbackslash{}mathbb\allowbreak{} F\_q\$,\allowbreak{} the\allowbreak{} Galois\allowbreak{} group\allowbreak{} with\allowbreak{} a\allowbreak{} blank\allowbreak{} numerator,\allowbreak{} and\allowbreak{} the\allowbreak{} displayed\allowbreak{} domain\allowbreak{} without\allowbreak{} an\allowbreak{} explicit\allowbreak{} direct-sum\allowbreak{} sign.\allowbreak{} The\allowbreak{} corrected\allowbreak{} French\allowbreak{} and\allowbreak{} English\allowbreak{} use\allowbreak{} standard\allowbreak{} vector-space\allowbreak{} terminology\allowbreak{} and\allowbreak{} agreement,\allowbreak{} state\allowbreak{} the\allowbreak{} \$M\$-th-power\allowbreak{} operation\allowbreak{} plainly,\allowbreak{} restore\allowbreak{} \$\textbackslash{}overline\{\textbackslash{}mathbb\allowbreak{} F\}\_q\$\allowbreak{} in\allowbreak{} the\allowbreak{} Galois\allowbreak{} group,\allowbreak{} and\allowbreak{} display\allowbreak{} the\allowbreak{} direct\allowbreak{} sum\allowbreak{} required\allowbreak{} by\allowbreak{} \$\textbackslash{}oplus\_i\textbackslash{}operatorna\allowbreak{}me\{Cl\}(Z\_i)\$.\allowbreak{} }\par
\par\medskip\noindent\textbf{EN-NOTE-195} 7II XXI; approximate folio 381\par
{\ttfamily\footnotesize \textbackslash{}medskip\allowbreak{} \textbackslash{}noindent\textbackslash{}emph\{Sourc\allowbreak{}e\allowbreak{} note.\}\allowbreak{} The\allowbreak{} diplomatic\allowbreak{} French\allowbreak{} preserves\allowbreak{} the\allowbreak{} printed\allowbreak{} word\allowbreak{} order\allowbreak{} \textbackslash{}emph\{intersection\allowbreak{} complète\allowbreak{} dans\allowbreak{} \$\textbackslash{}mathbb\allowbreak{} P\textasciicircum{}N\$\allowbreak{} lisse\},\allowbreak{} the\allowbreak{} typo\allowbreak{} \textbackslash{}emph\{da\allowbreak{} décomposition\allowbreak{} primitive\},\allowbreak{} and\allowbreak{} the\allowbreak{} omitted\allowbreak{} index\allowbreak{} in\allowbreak{} \$1\textbackslash{}leq\{\}\textbackslash{}leq\allowbreak{} n-1\$\allowbreak{} in\allowbreak{} (3.0.4).\allowbreak{} The\allowbreak{} corrected\allowbreak{} French\allowbreak{} and\allowbreak{} English\allowbreak{} use\allowbreak{} normal\allowbreak{} word\allowbreak{} order,\allowbreak{} \textbackslash{}emph\{la\allowbreak{} décomposition\allowbreak{} primitive\},\allowbreak{} and\allowbreak{} restore\allowbreak{} the\allowbreak{} forced\allowbreak{} index\allowbreak{} \$i\$\allowbreak{} in\allowbreak{} the\allowbreak{} cohomology\allowbreak{} range.\allowbreak{} }\par
\par\medskip\noindent\textbf{EN-NOTE-196} 7II XXI; approximate folio 382\par
{\ttfamily\footnotesize \textbackslash{}medskip\allowbreak{} \textbackslash{}noindent\textbackslash{}emph\{Sourc\allowbreak{}e\allowbreak{} note.\}\allowbreak{} In\allowbreak{} diagram\allowbreak{} (3.2.3)\allowbreak{} the\allowbreak{} diplomatic\allowbreak{} French\allowbreak{} preserves\allowbreak{} the\allowbreak{} printed\allowbreak{} \$X=\textbackslash{}widetilde\allowbreak{} X/\_\{\textbackslash{}eta\}\$\allowbreak{} notation.\allowbreak{} The\allowbreak{} corrected\allowbreak{} French\allowbreak{} and\allowbreak{} English\allowbreak{} use\allowbreak{} the\allowbreak{} standard\allowbreak{} generic-fiber\allowbreak{} notation\allowbreak{} \$X=\textbackslash{}widetilde\allowbreak{} X\_\textbackslash{}eta\$,\allowbreak{} agreeing\allowbreak{} with\allowbreak{} the\allowbreak{} concluding\allowbreak{} sentence\allowbreak{} of\allowbreak{} the\allowbreak{} proof.\allowbreak{} }\par
\par\medskip\noindent\textbf{EN-NOTE-197} 7II XXI; approximate folio 383\par
{\ttfamily\footnotesize \textbackslash{}medskip\allowbreak{} \textbackslash{}noindent\textbackslash{}emph\{Sourc\allowbreak{}e\allowbreak{} note.\}\allowbreak{} The\allowbreak{} diplomatic\allowbreak{} French\allowbreak{} preserves\allowbreak{} the\allowbreak{} printed\allowbreak{} lower\allowbreak{} limit\allowbreak{} \$1\$\allowbreak{} in\allowbreak{} (4.1.1)\allowbreak{} and\allowbreak{} the\allowbreak{} unaccented\allowbreak{} phrase\allowbreak{} \textbackslash{}emph\{de\allowbreak{} meme\}.\allowbreak{} The\allowbreak{} corrected\allowbreak{} French\allowbreak{} and\allowbreak{} English\allowbreak{} make\allowbreak{} the\allowbreak{} summation\allowbreak{} index\allowbreak{} explicit\allowbreak{} as\allowbreak{} \$i=1\$\allowbreak{} and\allowbreak{} use\allowbreak{} the\allowbreak{} ordinary\allowbreak{} accented\allowbreak{} form.\allowbreak{} }\par
\par\medskip\noindent\textbf{EN-NOTE-198} 7II XXI; approximate folio 389\par
{\ttfamily\footnotesize \textbackslash{}medskip\allowbreak{} \textbackslash{}noindent\textbackslash{}emph\{Sourc\allowbreak{}e\allowbreak{} note.\}\allowbreak{} The\allowbreak{} diplomatic\allowbreak{} French\allowbreak{} preserves\allowbreak{} the\allowbreak{} printed\allowbreak{} coefficient\allowbreak{} field\allowbreak{} \$F[[t]]\$\allowbreak{} in\allowbreak{} the\allowbreak{} zeta\allowbreak{} series\allowbreak{} and\allowbreak{} the\allowbreak{} printed\allowbreak{} restriction\allowbreak{} of\allowbreak{} \$\textbackslash{}underline\allowbreak{} G\$\allowbreak{} to\allowbreak{} \$U\$\allowbreak{} in\allowbreak{} the\allowbreak{} normalization\allowbreak{} reduction.\allowbreak{} The\allowbreak{} corrected\allowbreak{} French\allowbreak{} and\allowbreak{} English\allowbreak{} use\allowbreak{} the\allowbreak{} coefficient\allowbreak{} field\allowbreak{} \$E\$\allowbreak{} fixed\allowbreak{} in\allowbreak{} 5.0\allowbreak{} and\allowbreak{} the\allowbreak{} newly\allowbreak{} chosen\allowbreak{} open\allowbreak{} \$U'\$.\allowbreak{} Direct\allowbreak{} scan\allowbreak{} comparison\allowbreak{} also\allowbreak{} repairs\allowbreak{} the\allowbreak{} inherited\allowbreak{} joined\allowbreak{} words,\allowbreak{} Fraktur\allowbreak{} Galois\allowbreak{} symbols,\allowbreak{} and\allowbreak{} doubled\allowbreak{} consonant\allowbreak{} in\allowbreak{} \textbackslash{}emph\{D\textbackslash{}'esignons\}\allowbreak{} to\allowbreak{} the\allowbreak{} printed\allowbreak{} forms.\allowbreak{}}\par
\par\medskip\noindent\textbf{EN-NOTE-199} 7II XXI; approximate folio 390\par
{\ttfamily\footnotesize \textbackslash{}medskip\allowbreak{} \textbackslash{}noindent\textbackslash{}emph\{Sourc\allowbreak{}e\allowbreak{} note.\}\allowbreak{} The\allowbreak{} diplomatic\allowbreak{} French\allowbreak{} enumerates\allowbreak{} the\allowbreak{} eigenvalues\allowbreak{} on\allowbreak{} \$V\textasciicircum{}I\$\allowbreak{} by\allowbreak{} the\allowbreak{} dimension\allowbreak{} \$n\$\allowbreak{} of\allowbreak{} \$V\$,\allowbreak{} cites\allowbreak{} the\allowbreak{} lemma\allowbreak{} as\allowbreak{} ``3.2'',\allowbreak{} and\allowbreak{} calls\allowbreak{} \$k[T]/T\textasciicircum{}k\$\allowbreak{} irreducible.\allowbreak{} The\allowbreak{} corrected\allowbreak{} French\allowbreak{} and\allowbreak{} English\allowbreak{} use\allowbreak{} \$r=\textbackslash{}dim\_E\allowbreak{} V\textasciicircum{}I\$,\allowbreak{} the\allowbreak{} exact\allowbreak{} locator\allowbreak{} 5.3.2,\allowbreak{} and\allowbreak{} the\allowbreak{} mathematically\allowbreak{} correct\allowbreak{} term\allowbreak{} \textbackslash{}emph\{indecomposable\allowbreak{}\}.\allowbreak{} }\par
\par\medskip\noindent\textbf{EN-NOTE-200} 7II XXI; approximate folio 395\par
{\ttfamily\footnotesize \textbackslash{}medskip\allowbreak{} \textbackslash{}noindent\textbackslash{}emph\{Sourc\allowbreak{}e\allowbreak{} note.\}\allowbreak{} The\allowbreak{} diplomatic\allowbreak{} French\allowbreak{} preserves\allowbreak{} three\allowbreak{} printed\allowbreak{} slips:\allowbreak{} the\allowbreak{} coefficient\allowbreak{} \$c\_k\$\allowbreak{} is\allowbreak{} said\allowbreak{} to\allowbreak{} belong\allowbreak{} to\allowbreak{} \$\textbackslash{}sigma\allowbreak{} Z\$\allowbreak{} although\allowbreak{} the\allowbreak{} displayed\allowbreak{} series\allowbreak{} and\allowbreak{} the\allowbreak{} factor\allowbreak{} \$1/k\$\allowbreak{} make\allowbreak{} it\allowbreak{} a\allowbreak{} coefficient\allowbreak{} of\allowbreak{} \$\textbackslash{}log\textbackslash{}sigma\allowbreak{} Z\$;\allowbreak{} the\allowbreak{} proof\allowbreak{} of\allowbreak{} Theorem\allowbreak{} 5.5.2\allowbreak{} refers\allowbreak{} to\allowbreak{} itself\allowbreak{} rather\allowbreak{} than\allowbreak{} to\allowbreak{} the\allowbreak{} parallel\allowbreak{} induction\allowbreak{} in\allowbreak{} Theorem\allowbreak{} 5.2.2;\allowbreak{} and\allowbreak{} the\allowbreak{} proof\allowbreak{} of\allowbreak{} Corollary\allowbreak{} 5.5.3\allowbreak{} says\allowbreak{} to\allowbreak{} prove\allowbreak{} Theorem\allowbreak{} 5.2.2\allowbreak{} for\allowbreak{} \$U\$\allowbreak{} although\allowbreak{} the\allowbreak{} argument\allowbreak{} is\allowbreak{} proving\allowbreak{} assertion\allowbreak{} (iv)\allowbreak{} of\allowbreak{} the\allowbreak{} corollary.\allowbreak{} The\allowbreak{} corrected\allowbreak{} French\allowbreak{} and\allowbreak{} English\allowbreak{} give\allowbreak{} the\allowbreak{} coherent\allowbreak{} readings.\allowbreak{} Direct\allowbreak{} scan\allowbreak{} comparison\allowbreak{} also\allowbreak{} restores\allowbreak{} the\allowbreak{} notation\allowbreak{} underlines\allowbreak{} omitted\allowbreak{} by\allowbreak{} the\allowbreak{} inherited\allowbreak{} French\allowbreak{} TeX\allowbreak{} in\allowbreak{} the\allowbreak{} statement\allowbreak{} of\allowbreak{} Theorem\allowbreak{} 5.5.2.\allowbreak{} }\par
\par\medskip\noindent\textbf{EN-NOTE-201} 7II XXI; approximate folio 399\par
{\ttfamily\footnotesize \textbackslash{}medskip\allowbreak{} \textbackslash{}noindent\textbackslash{}emph\{Sourc\allowbreak{}e\allowbreak{} note.\}\allowbreak{} Direct\allowbreak{} authority\allowbreak{} comparison\allowbreak{} restores\allowbreak{} the\allowbreak{} factorization\allowbreak{} diagram's\allowbreak{} \$j\$\allowbreak{} label\allowbreak{} to\allowbreak{} the\allowbreak{} lower\allowbreak{} side\allowbreak{} of\allowbreak{} its\allowbreak{} arrow\allowbreak{} and\allowbreak{} restores\allowbreak{} the\allowbreak{} printed\allowbreak{} Fraktur\allowbreak{} \$\textbackslash{}mathfrak\allowbreak{} G\$\allowbreak{} in\allowbreak{} the\allowbreak{} curve\allowbreak{} case.\allowbreak{} The\allowbreak{} diplomatic\allowbreak{} French\allowbreak{} preserves\allowbreak{} the\allowbreak{} printed\allowbreak{} singular\allowbreak{} \textbackslash{}emph\{morphisme\},\allowbreak{} the\allowbreak{} punctuation\allowbreak{} collision\allowbreak{} before\allowbreak{} hypothesis\allowbreak{} \$(*)\$,\allowbreak{} and\allowbreak{} the\allowbreak{} shortened\allowbreak{} reference\allowbreak{} 6.2(i).\allowbreak{} Corrected\allowbreak{} French\allowbreak{} and\allowbreak{} English\allowbreak{} use\allowbreak{} plural\allowbreak{} \textbackslash{}emph\{morphismes\},\allowbreak{} ordinary\allowbreak{} sentence\allowbreak{} punctuation,\allowbreak{} and\allowbreak{} the\allowbreak{} exact\allowbreak{} reference\allowbreak{} Lemma\allowbreak{} 5.6.2(i).\allowbreak{} }\par
\par\medskip\noindent\textbf{EN-NOTE-202} 7II XXII; approximate folio 406\par
{\ttfamily\footnotesize \textbackslash{}medskip\allowbreak{} \textbackslash{}noindent\textbackslash{}emph\{Sourc\allowbreak{}e\allowbreak{} note.\}\allowbreak{} The\allowbreak{} diplomatic\allowbreak{} French\allowbreak{} preserves\allowbreak{} the\allowbreak{} print's\allowbreak{} hyphen\allowbreak{} in\allowbreak{} \textbackslash{}emph\{peut-être\},\allowbreak{} the\allowbreak{} singular\allowbreak{} phrase\allowbreak{} \textbackslash{}emph\{coefficient\allowbreak{} tordus\},\allowbreak{} \textbackslash{}emph\{Una\allowbreak{} application\},\allowbreak{} the\allowbreak{} omitted\allowbreak{} verb\allowbreak{} in\allowbreak{} \textbackslash{}emph\{si\allowbreak{} elle\allowbreak{} additive\},\allowbreak{} and\allowbreak{} two\allowbreak{} spacing/punctuation\allowbreak{} collisions.\allowbreak{} The\allowbreak{} corrected\allowbreak{} French\allowbreak{} and\allowbreak{} English\allowbreak{} regularize\allowbreak{} those\allowbreak{} surface\allowbreak{} defects\allowbreak{} without\allowbreak{} changing\allowbreak{} the\allowbreak{} mathematical\allowbreak{} content.\allowbreak{} Formula\allowbreak{} (1.0.7)\allowbreak{} was\allowbreak{} read\allowbreak{} directly\allowbreak{} from\allowbreak{} a\allowbreak{} high-detail\allowbreak{} authority\allowbreak{} crop:\allowbreak{} its\allowbreak{} large\allowbreak{} operator\allowbreak{} is\allowbreak{} an\allowbreak{} intersection.\allowbreak{}}\par
\par\medskip\noindent\textbf{EN-NOTE-203} 7II XXII; approximate folio 409\par
{\ttfamily\footnotesize \textbackslash{}medskip\allowbreak{} \textbackslash{}noindent\textbackslash{}emph\{Sourc\allowbreak{}e\allowbreak{} note.\}\allowbreak{} The\allowbreak{} final\allowbreak{} printed\allowbreak{} sentence\allowbreak{} of\allowbreak{} the\allowbreak{} proof\allowbreak{} says\allowbreak{} \$\textbackslash{}varphi\_q\$\allowbreak{} where\allowbreak{} the\allowbreak{} proposition\allowbreak{} and\allowbreak{} the\allowbreak{} immediately\allowbreak{} preceding\allowbreak{} inverse\allowbreak{} relation\allowbreak{} require\allowbreak{} \$\textbackslash{}varphi\_q\textasciicircum{}\{-1\}\$.\allowbreak{} The\allowbreak{} English\allowbreak{} follows\allowbreak{} the\allowbreak{} mathematically\allowbreak{} coherent\allowbreak{} reading\allowbreak{} and\allowbreak{} the\allowbreak{} corrected\allowbreak{} French\allowbreak{} edition;\allowbreak{} the\allowbreak{} diplomatic\allowbreak{} French\allowbreak{} text\allowbreak{} retains\allowbreak{} the\allowbreak{} printed\allowbreak{} wording.\allowbreak{} }\par
\par\medskip\noindent\textbf{EN-NOTE-204} 7II XXII; approximate folio 414\par
{\ttfamily\footnotesize \textbackslash{}medskip\allowbreak{} \textbackslash{}noindent\textbackslash{}emph\{Sourc\allowbreak{}e\allowbreak{} note.\}\allowbreak{} The\allowbreak{} printed\allowbreak{} opening\allowbreak{} references\allowbreak{} three\allowbreak{} nonexistent\allowbreak{} or\allowbreak{} malformed\allowbreak{} variants\allowbreak{} of\allowbreak{} (3.1.1);\allowbreak{} the\allowbreak{} corrected\allowbreak{} French\allowbreak{} and\allowbreak{} English\allowbreak{} cite\allowbreak{} the\allowbreak{} actual\allowbreak{} formulas\allowbreak{} (3.1.2),\allowbreak{} (3.2.1)--(3.2.2),\allowbreak{} and\allowbreak{} (3.3.1).\allowbreak{} In\allowbreak{} the\allowbreak{} proof\allowbreak{} of\allowbreak{} Theorem\textasciitilde{}3.3\allowbreak{} the\allowbreak{} print\allowbreak{} writes\allowbreak{} \$\textbackslash{}mathrm\{fr\}\_p\textasciicircum{}\{X*\}\$\allowbreak{} although\allowbreak{} the\allowbreak{} map\allowbreak{} displayed\allowbreak{} beside\allowbreak{} it\allowbreak{} is\allowbreak{} \$\textbackslash{}mathrm\{fr\}\_q\textasciicircum{}X\$,\allowbreak{} and\allowbreak{} twice\allowbreak{} says\allowbreak{} that\allowbreak{} Frobenius\allowbreak{} kills\allowbreak{} spectral-sequence\allowbreak{} terms\allowbreak{} for\allowbreak{} \$q\textbackslash{}ne0\$\allowbreak{} rather\allowbreak{} than\allowbreak{} for\allowbreak{} positive\allowbreak{} differential-form\allowbreak{} degree\allowbreak{} \$p\textbackslash{}ne0\$.\allowbreak{} Finally,\allowbreak{} the\allowbreak{} proof\allowbreak{} of\allowbreak{} Lemma\textasciitilde{}3.9\allowbreak{} twice\allowbreak{} cites\allowbreak{} nonexistent\allowbreak{} (3.1.3)\allowbreak{} where\allowbreak{} its\allowbreak{} own\allowbreak{} argument\allowbreak{} is\allowbreak{} proving\allowbreak{} (3.4.1).\allowbreak{} The\allowbreak{} diplomatic\allowbreak{} French\allowbreak{} retains\allowbreak{} all\allowbreak{} these\allowbreak{} printed\allowbreak{} forms;\allowbreak{} corrected\allowbreak{} French\allowbreak{} and\allowbreak{} English\allowbreak{} use\allowbreak{} the\allowbreak{} coherent\allowbreak{} readings\allowbreak{} just\allowbreak{} stated.\allowbreak{} }\par
\par\medskip\noindent\textbf{EN-NOTE-205} 7II XXII; approximate folio 417\par
{\ttfamily\footnotesize \textbackslash{}medskip\allowbreak{} \textbackslash{}noindent\textbackslash{}emph\{Sourc\allowbreak{}e\allowbreak{} note.\}\allowbreak{} The\allowbreak{} print\allowbreak{} has\allowbreak{} three\allowbreak{} mutually\allowbreak{} exposed\allowbreak{} notation\allowbreak{} slips\allowbreak{} in\allowbreak{} Lemma\textasciitilde{}4.2:\allowbreak{} it\allowbreak{} writes\allowbreak{} \$A\_1=\textbackslash{}pi\$\allowbreak{} in\allowbreak{} the\allowbreak{} coefficient\allowbreak{} expansion\allowbreak{} in\allowbreak{} (4.2.1),\allowbreak{} places\allowbreak{} the\allowbreak{} Teichmüller\allowbreak{} representative\allowbreak{} in\allowbreak{} \$W(\textbackslash{}mathbb\allowbreak{} F\_\{q\textasciicircum{}a\})\$\allowbreak{} although\allowbreak{} both\allowbreak{} \$t\textasciicircum{}\{p\textasciicircum{}a\}=t\$\allowbreak{} and\allowbreak{} Corollary\textasciitilde{}4.3\allowbreak{} use\allowbreak{} \$\textbackslash{}mathbb\allowbreak{} F\_\{p\textasciicircum{}a\}\$,\allowbreak{} and\allowbreak{} prints\allowbreak{} \$\textbackslash{}sum\allowbreak{} t\textasciicircum{}j\$\allowbreak{} in\allowbreak{} the\allowbreak{} exponent\allowbreak{} of\allowbreak{} (4.2.3.1)\allowbreak{} although\allowbreak{} the\allowbreak{} product\allowbreak{} and\allowbreak{} the\allowbreak{} following\allowbreak{} proof\allowbreak{} use\allowbreak{} \$\textbackslash{}sum\allowbreak{} t\textasciicircum{}\{p\textasciicircum{}j\}\$.\allowbreak{} The\allowbreak{} diplomatic\allowbreak{} French\allowbreak{} preserves\allowbreak{} those\allowbreak{} forms;\allowbreak{} corrected\allowbreak{} French\allowbreak{} and\allowbreak{} English\allowbreak{} use\allowbreak{} \$B\_1\$,\allowbreak{} \$\textbackslash{}mathbb\allowbreak{} F\_\{p\textasciicircum{}a\}\$,\allowbreak{} and\allowbreak{} \$t\textasciicircum{}\{p\textasciicircum{}j\}\$.\allowbreak{} }\par
\par\medskip\noindent\textbf{EN-NOTE-206} 7II XXII; approximate folio 422\par
{\ttfamily\footnotesize \textbackslash{}medskip\allowbreak{} \textbackslash{}noindent\textbackslash{}emph\{Sourc\allowbreak{}e\allowbreak{} note.\}\allowbreak{} The\allowbreak{} print\allowbreak{} drops\allowbreak{} \$\textbackslash{}nu\$\allowbreak{} from\allowbreak{} the\allowbreak{} factor\allowbreak{} \$q\textasciicircum{}\textbackslash{}nu-1\$\allowbreak{} in\allowbreak{} (4.5.10)\allowbreak{} and\allowbreak{} from\allowbreak{} the\allowbreak{} denominator\allowbreak{} \$q\textasciicircum{}\textbackslash{}nu\$\allowbreak{} in\allowbreak{} (4.5.12),\allowbreak{} although\allowbreak{} both\allowbreak{} formulas\allowbreak{} are\allowbreak{} direct\allowbreak{} rewrites\allowbreak{} of\allowbreak{} (4.4.6).\allowbreak{} It\allowbreak{} also\allowbreak{} transposes\allowbreak{} the\allowbreak{} formula\allowbreak{} number\allowbreak{} (4.5.14)\allowbreak{} as\allowbreak{} (5.4.14).\allowbreak{} The\allowbreak{} diplomatic\allowbreak{} French\allowbreak{} preserves\allowbreak{} all\allowbreak{} three\allowbreak{} printed\allowbreak{} forms;\allowbreak{} corrected\allowbreak{} French\allowbreak{} and\allowbreak{} English\allowbreak{} use\allowbreak{} the\allowbreak{} algebraically\allowbreak{} forced\allowbreak{} readings\allowbreak{} displayed\allowbreak{} above.\allowbreak{} }\par
\par\medskip\noindent\textbf{EN-NOTE-207} 7II XXII; approximate folio 425\par
{\ttfamily\footnotesize \textbackslash{}medskip\allowbreak{} \textbackslash{}noindent\textbackslash{}emph\{Sourc\allowbreak{}e\allowbreak{} note.\}\allowbreak{} The\allowbreak{} print\allowbreak{} numbers\allowbreak{} (4.5.23)\allowbreak{} as\allowbreak{} (4.5.2.3),\allowbreak{} abbreviates\allowbreak{} \$\textbackslash{}operatorname\{card\}\allowbreak{}(A)\$\allowbreak{} as\allowbreak{} \$\textbackslash{}operatorname\{car\}(\allowbreak{}A)\$\allowbreak{} in\allowbreak{} the\allowbreak{} first\allowbreak{} exponent\allowbreak{} of\allowbreak{} (4.5.25),\allowbreak{} and\allowbreak{} reuses\allowbreak{} the\allowbreak{} number\allowbreak{} (4.5.23)\allowbreak{} for\allowbreak{} the\allowbreak{} final\allowbreak{} zeta\allowbreak{} identity.\allowbreak{} The\allowbreak{} corrected\allowbreak{} French\allowbreak{} and\allowbreak{} English\allowbreak{} layers\allowbreak{} use\allowbreak{} the\allowbreak{} sequential\allowbreak{} numbers\allowbreak{} and\allowbreak{} the\allowbreak{} full\allowbreak{} cardinality\allowbreak{} notation\allowbreak{} displayed\allowbreak{} above;\allowbreak{} the\allowbreak{} diplomatic\allowbreak{} French\allowbreak{} preserves\allowbreak{} the\allowbreak{} printed\allowbreak{} forms.\allowbreak{}}\par
\par\medskip\noindent\textbf{EN-NOTE-208} 7II XXII; approximate folio 428\par
{\ttfamily\footnotesize \textbackslash{}medskip\allowbreak{} \textbackslash{}noindent\textbackslash{}emph\{Sourc\allowbreak{}e\allowbreak{} note.\}\allowbreak{} The\allowbreak{} print\allowbreak{} numbers\allowbreak{} the\allowbreak{} hypothesis\allowbreak{} in\allowbreak{} Theorem\textasciitilde{}5.0\allowbreak{} bis\allowbreak{} as\allowbreak{} (5.0.7),\allowbreak{} duplicating\allowbreak{} the\allowbreak{} number\allowbreak{} subsequently\allowbreak{} used\allowbreak{} for\allowbreak{} the\allowbreak{} point-count\allowbreak{} identity.\allowbreak{} The\allowbreak{} corrected\allowbreak{} French\allowbreak{} and\allowbreak{} English\allowbreak{} layers\allowbreak{} use\allowbreak{} the\allowbreak{} sequential\allowbreak{} number\allowbreak{} (5.0.4);\allowbreak{} the\allowbreak{} diplomatic\allowbreak{} French\allowbreak{} preserves\allowbreak{} the\allowbreak{} printed\allowbreak{} number.\allowbreak{}}\par
\par\medskip\noindent\textbf{EN-NOTE-209} 7II XXII; approximate folio 430\par
{\ttfamily\footnotesize \textbackslash{}medskip\allowbreak{} \textbackslash{}noindent\textbackslash{}emph\{Sourc\allowbreak{}e\allowbreak{} note.\}\allowbreak{} In\allowbreak{} (5.0.9)\allowbreak{} the\allowbreak{} print\allowbreak{} has\allowbreak{} the\allowbreak{} impossible\allowbreak{} factor\allowbreak{} \$(1-q\textasciicircum{}if)\$\allowbreak{} rather\allowbreak{} than\allowbreak{} \$(1-q\textasciicircum{}it)\$,\allowbreak{} and\allowbreak{} the\allowbreak{} last\allowbreak{} paragraph\allowbreak{} cites\allowbreak{} Lemma\textasciitilde{}4.4.6\allowbreak{} rather\allowbreak{} than\allowbreak{} the\allowbreak{} integrality\allowbreak{} lemma\allowbreak{} 4.6.6\allowbreak{} used\allowbreak{} there.\allowbreak{} Both\allowbreak{} are\allowbreak{} corrected\allowbreak{} above;\allowbreak{} the\allowbreak{} diplomatic\allowbreak{} French\allowbreak{} retains\allowbreak{} the\allowbreak{} printed\allowbreak{} forms.\allowbreak{}}\par
\par\medskip\noindent\textbf{EN-NOTE-210} 7II XXII; approximate folio 431\par
{\ttfamily\footnotesize \textbackslash{}medskip\allowbreak{} \textbackslash{}noindent\textbackslash{}emph\{Sourc\allowbreak{}e\allowbreak{} note.\}\allowbreak{} The\allowbreak{} print\allowbreak{} cites\allowbreak{} (4.5.24),\allowbreak{} but\allowbreak{} the\allowbreak{} congruence\allowbreak{} follows\allowbreak{} from\allowbreak{} the\allowbreak{} terminal\allowbreak{} zeta\allowbreak{} identity\allowbreak{} (4.5.33).\allowbreak{} It\allowbreak{} also\allowbreak{} defines\allowbreak{} \$W\_S\$\allowbreak{} using\allowbreak{} \$R\_B\$,\allowbreak{} although\allowbreak{} the\allowbreak{} entire\allowbreak{} section\allowbreak{} and\allowbreak{} every\allowbreak{} following\allowbreak{} formula\allowbreak{} concern\allowbreak{} \$R\_S\$.\allowbreak{} The\allowbreak{} corrected\allowbreak{} French\allowbreak{} and\allowbreak{} English\allowbreak{} layers\allowbreak{} use\allowbreak{} (4.5.33)\allowbreak{} and\allowbreak{} \$R\_S\$;\allowbreak{} the\allowbreak{} diplomatic\allowbreak{} French\allowbreak{} retains\allowbreak{} the\allowbreak{} printed\allowbreak{} forms.\allowbreak{} }\par
\par\medskip\noindent\textbf{EN-NOTE-211} 7II XXII; approximate folio 435\par
{\ttfamily\footnotesize \textbackslash{}medskip\allowbreak{} \textbackslash{}noindent\textbackslash{}emph\{Sourc\allowbreak{}e\allowbreak{} note.\}\allowbreak{} The\allowbreak{} first\allowbreak{} product\allowbreak{} in\allowbreak{} the\allowbreak{} printed\allowbreak{} proof\allowbreak{} has\allowbreak{} the\allowbreak{} undefined\allowbreak{} upper\allowbreak{} bound\allowbreak{} \$\textbackslash{}operatorname\{ord\}(\allowbreak{}a)-1\$;\allowbreak{} the\allowbreak{} next\allowbreak{} line\allowbreak{} and\allowbreak{} the\allowbreak{} definition\allowbreak{} of\allowbreak{} \$q\$\allowbreak{} require\allowbreak{} \$\textbackslash{}operatorname\{ord\}(\allowbreak{}q)-1\$.\allowbreak{} The\allowbreak{} corrected\allowbreak{} French\allowbreak{} and\allowbreak{} English\allowbreak{} layers\allowbreak{} use\allowbreak{} the\allowbreak{} latter;\allowbreak{} the\allowbreak{} diplomatic\allowbreak{} French\allowbreak{} preserves\allowbreak{} the\allowbreak{} print.\allowbreak{}}\par
\par\medskip\noindent\textbf{EN-NOTE-212} 7II XXII; approximate folio 437\par
{\ttfamily\footnotesize \textbackslash{}medskip\allowbreak{} \textbackslash{}noindent\textbackslash{}emph\{Sourc\allowbreak{}e\allowbreak{} note.\}\allowbreak{} In\allowbreak{} (6.3.4)\allowbreak{} the\allowbreak{} print\allowbreak{} has\allowbreak{} \$H\textasciicircum{}\{i+1\}\$,\allowbreak{} while\allowbreak{} Proposition\textasciitilde{}6.3,\allowbreak{} (6.3.3),\allowbreak{} and\allowbreak{} the\allowbreak{} sentence\allowbreak{} introducing\allowbreak{} the\allowbreak{} computation\allowbreak{} all\allowbreak{} require\allowbreak{} \$H\textasciicircum{}\{n+1\}\$.\allowbreak{} The\allowbreak{} corrected\allowbreak{} French\allowbreak{} and\allowbreak{} English\allowbreak{} layers\allowbreak{} use\allowbreak{} \$n+1\$;\allowbreak{} the\allowbreak{} diplomatic\allowbreak{} French\allowbreak{} preserves\allowbreak{} the\allowbreak{} printed\allowbreak{} exponent.\allowbreak{}}\par

\normalsize
\clearpage

\section*{Retained readings and unresolved editorial questions}

These notes record alternatives that have not been inserted as new corrections. They are separate from the numbered variant apparatus: a note without a diplomatic/corrected difference adds no apparatus entry. “Needs the scan” concerns a physical reading that this build cannot verify. The attached diplomatic layer remains the printed-reading witness.

\subsection*{ED-G27-x-hat — 7I:123}

Retained/source reading: \textbackslash{}mathfrak X\_\{(\textbackslash{}hat x)\} in the supplied source layers

Alternative: earlier transcription \textbackslash{}mathcal X\_\{(\textbackslash{}mathfrak x)\}

Keep the supplied source-layer hat reading; do not invent a further glyph correction. The two current source pages both use the completion hat and the same completion-to-localization isomorphism. The earlier convention entry questioned the physical overstrike; without the scan that paleographic question cannot be independently resolved. The authoritative attached layers remain the printed-reading witness.

Status: needs the scan. Source: \nolinkurl{conventions/GLOBAL_ISSUES_SGA.md:27}

\subsection*{ED-G27-missing-reference — 7I:128}

Retained/source reading: is smooth by and therefore

Alternative: is smooth by 4.13(a) and therefore

Retain the printed blank reference in every body; record the suggested reference only as a note. The statement lacks a number in the recorded print. More than one preceding smoothness result could support the argument, and adding 4.13(a) would manufacture a citation. No supplied source establishes the omitted numeral.

Status: retained; authorial reference unresolved. Source: \nolinkurl{conventions/GLOBAL_ISSUES_SGA.md:27}

\subsection*{ED-G32-run-in-headings — 7I:163,165 (and analogous numbered headings)}

Retained/source reading: Attached layers use separate paragraphs after numbered headings

Alternative: possible printed run-in headings with following prose on the same paragraph

Keep the supplied paragraph boundaries; do not make an unchecked global reflow. The two specific changes were held at low confidence and the other headings were not verified. A text-only layer cannot establish a printed indentation or run-in break. This is a layout uncertainty, not grounds for altering any authorial wording.

Status: needs the scan. Source: \nolinkurl{conventions/GLOBAL_ISSUES_SGA.md:32}

\subsection*{ED-G36-E-or-X — 7I:180}

Retained/source reading: d\_1\textasciicircum{}*d\_0\textasciicircum{}*(E)\textbackslash{}simeq(d\_0d\_1)\textasciicircum{}*(E)

Alternative: d\_1\textasciicircum{}*d\_0\textasciicircum{}*(X)\textbackslash{}simeq(d\_0d\_1)\textasciicircum{}*(X)

Retain E in all bodies; record the proposed X reading without adopting it. The preceding concrete extension is called X, but the transitivity isomorphism is valid for any extension E. Consequently the typography does not force a mathematical error. Both supplied layers retain E, and an authorial generic E remains possible.

Status: retained; no determinate emendation. Source: \nolinkurl{conventions/GLOBAL_ISSUES_SGA.md:36}

\subsection*{ED-G36-psi-footnote — 7I:180}

Retained/source reading: \textbackslash{}psi\_\{L.\}(A)

Alternative: \textbackslash{}underline\textbackslash{}Psi\_\{L.\}(A)

Retain the source-layer lowercase psi in the body. The footnote names a category also denoted by underlined capital Psi in the main text, but the recorded low-confidence observation does not establish whether the glyph is a typist variant or an accidental change. No scan is supplied to adjudicate its shape.

Status: needs the scan. Source: \nolinkurl{conventions/GLOBAL_ISSUES_SGA.md:36}

\subsection*{ED-G37-duplicate-number — 7I:183}

Retained/source reading: two displays tagged (3.2.3)

Alternative: second display (3.2.3\textbackslash{}prime) or another new tag

Retain both printed (3.2.3) tags; distinguish their locations by description in the records. The following source already uses (3.2.4); no unoccupied intended number is established. An editorial prime would be a new citation system, not recovery of the author’s number. The diagram and c(X,alpha) display remain distinguishable by their content.

Status: retained; intended numbering unresolved. Source: \nolinkurl{conventions/GLOBAL_ISSUES_SGA.md:37}

\subsection*{ED-G38-speck — 7II:202}

Retained/source reading: (\textbackslash{}mathrm\{Sw\}\_\{I\textbackslash{}prime\},s\_\{I/H\}) without an extra lower dot

Alternative: a possible printed dot below the opening parenthesis

Keep the supplied formula without an additional dot. No mathematical operation is defined by the possible mark. Its status as ink, damage or punctuation is a visual-source question; neither attached layer encodes it.

Status: needs the scan. Source: \nolinkurl{conventions/GLOBAL_ISSUES_SGA.md:38}

\subsection*{ED-G41-exton-sion — 7I:195}

Retained/source reading: extension (both attached layers)

Alternative: exton-sion (held low-confidence visual reading)

Retain the attached diplomatic reading extension, and list the alternative only in the note. The two source layers agree. The earlier candidate was explicitly held and supplies no established printed lexical error. The scan is needed to distinguish an actual o from a damaged e and a line-end hyphen.

Status: needs the scan. Source: \nolinkurl{conventions/GLOBAL_ISSUES_SGA.md:41}

\subsection*{ED-G45-mathematical-repair — 7II:237}

Retained/source reading: the printed pencil in (4.2.3)

Alternative: a proposed replacement pencil with an isolated ordinary quadratic singularity

Preserve the supplied mathematical formula; do not construct a new proof or pencil. The global record explicitly calls the proposed replacement a repair of the mathematics rather than a recovered reading. The detailed proposed formula is not supplied in the attached findings file. The recorded concern (a singular plane for n=2k+1 at least 7) is disclosed, but no conjectural replacement is inserted.

Status: retained; proposed mathematical replacement unavailable. Source: \nolinkurl{conventions/GLOBAL_ISSUES_SGA.md:45}

\subsection*{ED-G48-inner-underlines — 7II:248 and other internal commas}

Retained/source reading: source-layer continuous underlines through some inner commas

Alternative: stop and restart each underline around the comma

Keep the source-layer markup pending visual evidence at each occurrence. The global record states that the typist sometimes underlined through commas and sometimes did not. Therefore a string-wide rule is not justified; each of the potential cases requires its scan. No authorial wording is changed.

Status: needs the scan. Source: \nolinkurl{conventions/GLOBAL_ISSUES_SGA.md:48}

\subsection*{ED-G58-comma — 7II:332}

Retained/source reading: H\_s,

Alternative: H\_\{s,\}

Keep the comma as argument punctuation; do not move it into the subscript. The lower comma was a held visual observation, not a new index. The H\_s followed by coefficient argument consistently denotes the fibre. Both layers retain ordinary argument punctuation; a physical baseline question requires the scan.

Status: needs the scan. Source: \nolinkurl{conventions/GLOBAL_ISSUES_SGA.md:58}

\subsection*{ED-G63-miracle — 7II:416}

Retained/source reading: oh miracle (both attached layers)

Alternative: ô miracle (earlier held reading)

Retain oh miracle as the attached diplomatic reading. The later supplied layers agree on oh. The earlier low-confidence record does not establish the visual reading independently; changing both layers back to ô would discard the current witness without evidence.

Status: needs the scan. Source: \nolinkurl{conventions/GLOBAL_ISSUES_SGA.md:63}

\subsection*{ED-G64-quotient — 7II:423}

Retained/source reading: the displayed fraction in (4.5.21)

Alternative: a slash presentation of the same quotient

Keep the displayed fraction in both layers. Both attached layers encode the same numerator and denominator as a quotient. The proposed slash concerns the original graphic layout alone. No mathematical change follows, and without a scan the alternative typography cannot be established.

Status: needs the scan. Source: \nolinkurl{conventions/GLOBAL_ISSUES_SGA.md:64}

\subsection*{ED-G64-subset-stroke — 7II:424}

Retained/source reading: \textbackslash{}subseteq in the supplied source

Alternative: \textbackslash{}subset without the lower stroke

Keep the attached diplomatic subseteq reading. The issue reports a possible overstrike, not a verified symbol. Historical subset notation may also be non-strict, so semantics do not settle the physical stroke. No inequality is silently modernized.

Status: needs the scan. Source: \nolinkurl{conventions/GLOBAL_ISSUES_SGA.md:64}

\subsection*{ED-G64-Delta-Lambda — 7II:424–429}

Retained/source reading: \textbackslash{}Delta\_B throughout the defining discussion; \textbackslash{}Lambda\_B in printed (5.0.12)

Alternative: change all seven \textbackslash{}Delta\_B to \textbackslash{}Lambda\_B, or retain current \textbackslash{}Delta\_B at (5.0.12)

Keep the supplied diplomatic glyphs and the local current Delta correction; reject an unchecked all-seven replacement. The attached layers define Delta\_B in (4.5.28) and use it in the intervening argument. Current (5.0.12) refers back to that same characteristic series. The earlier global all-seven proposal requires scans of all occurrences and cannot override the later source-layer distinction. The physical glyph family remains a scan question.

Status: needs the scan. Source: \nolinkurl{conventions/GLOBAL_ISSUES_SGA.md:64}

\subsection*{ED-G67-Dwork-psi — 7II:436–438}

Retained/source reading: \textbackslash{}psi\_q in the attached layers

Alternative: \textbackslash{}Psi\_q suggested by the typewriter capital-form glyph

Retain psi consistently with the attached source notation; do not recase the operator. Issue 78 excludes glyph-only substitutions from textual diplomatic variants. The operator is the same Dwork operator defined earlier in Exposé XXII; its cap-height typescript form does not introduce a second operator. No scan-backed case distinction is available.

Status: needs the scan. Source: \nolinkurl{conventions/GLOBAL_ISSUES_SGA.md:67}

\subsection*{ED-G73-statement-underlines — 7II:403,405 and similar statements}

Retained/source reading: source underlines sometimes include mathematical symbols and final punctuation

Alternative: split each underline around symbols and punctuation according to the print

Retain the supplied markup except for exact separately documented repairs, including ED-XV-1.3.2-i. The global record defers these cases to the individual visual review required by issue 48; the typist did not follow one mechanically inferable rule. A broad replacement would silently modify authorial presentation at unverified loci.

Status: needs the scan. Source: \nolinkurl{conventions/GLOBAL_ISSUES_SGA.md:73}

\fi

\end{document}
